% docs/DEEPEST_de.tex -- generated by tools/deepest_booklet.py --lang de.
% Do not edit: regenerate from docs/DEEPEST.json instead.
% Build: lualatex DEEPEST_de.tex (twice -- the index carries page refs).
% Needs Source Serif Pro, Linux Biolinum O, chessboard, tcolorbox, pgfplots, qrcode.

\documentclass[a5paper,10pt,twoside,openany,headings=small]{scrbook}

\usepackage[ngerman]{babel}
\usepackage[a5paper,inner=15mm,outer=15mm,top=14mm,bottom=18mm]{geometry}
\usepackage{fontspec}
\setmainfont{Source Serif Pro}[Numbers=Lining]
\setsansfont{Linux Biolinum O}
\usepackage{microtype}
\usepackage{booktabs}
\usepackage{longtable}
\usepackage{array}
\usepackage[most]{tcolorbox}
\usepackage{pgfplots}
\pgfplotsset{compat=1.16}
\usepackage{chessfss}
\usepackage{chessboard}
\usepackage{qrcode}
\usepackage[hidelinks]{hyperref}
\usepackage[automark]{scrlayer-scrpage}

\definecolor{ink}{HTML}{24476B}
\setchessboard{boardfontsize=17pt,labelfontsize=5pt,showmover=false,borderwidth=0.3pt,margin=false}
\addtokomafont{disposition}{\color{ink}}

\clearpairofpagestyles
\ofoot{\sffamily\footnotesize\pagemark}
\ifoot{\sffamily\footnotesize\color{ink}Längste Hilfsmatts mit maximal 6 Steinen}
\pagestyle{scrheadings}

\raggedbottom
\setlength{\parindent}{0pt}
\setlength{\parskip}{0.5\baselineskip}
\setcounter{tocdepth}{0}
\setcounter{secnumdepth}{0}
\pgfplotsset{every axis/.append style={/pgf/number format/.cd,use comma,1000 sep={\,}}}

% A theme as a small tinted tag.
\newcommand{\themetag}[1]{\tcbox[on line,colback=ink!8,colframe=ink!8,boxrule=0pt,arc=1.5pt,
  left=2pt,right=2pt,top=0.5pt,bottom=0.5pt,fontupper=\sffamily\scriptsize\color{ink}]{#1}}
% The diagram, measured so the caption under it can end at its right edge.
\newsavebox{\hmboard}
% A labelled key figure in the statistics column.
\newcommand{\hmfigure}[2]{\textcolor{ink}{\textsc{#1}}\\ #2\\[1mm]}

\begin{document}
\frontmatter
\begin{titlepage}
\centering\sffamily
\vspace*{2cm}
{\Huge\bfseries\color{ink} Längste Hilfsmatts\\[0.3em] mit maximal 6 Steinen\par}
\vspace{1.2cm}
{\large Eine Aufgabe je Materialklasse:\\ das längste Hilfsmatt, das noch eine einzige Lösung hat\par}
\vspace{2cm}
{\large Alle 861 Materialklassen von drei bis sechs Steinen\par}
\vfill
{\small Erzeugt aus der Hilfsmatt-Tablebase.\\ Jedes Diagramm, jede Länge und jede Lösung in diesem Heft ist maschinell ermittelt.\par}
\vspace{1cm}
\end{titlepage}

\chapter*{Die Tablebase in Kürze}
\markboth{Die Tablebase in Kürze}{}
\addcontentsline{toc}{chapter}{Die Tablebase in Kürze}

Ein \emph{Hilfsmatt} ist eine Schachaufgabe, in der beide Seiten zusammenarbeiten: Schwarz zieht an und hilft Weiß, matt zu setzen. \textbf{h\#n} verlangt Matt im \(n\)-ten Zug von Weiß; bei \textbf{h\#n.5} beginnt Weiß, die Lösung beginnt also mit \(1\ldots\) und dauert einen Halbzug länger. \emph{Korrekt} ist eine Aufgabe nur, wenn ihre Lösung eindeutig ist: eine zweite Lösung ist eine Nebenlösung.

Die Hilfsmatt-Tablebase löst jede legale Stellung jeder Materialklasse mit bis zu sechs Steinen. Für jede Stellung und jedes Zugrecht speichert sie die Entfernung zum Matt bei Zusammenarbeit und die Zahl der verschiedenen kürzesten Lösungen; eine korrekte Aufgabe ist also einfach eine Stellung, deren Zahl eins ist.

\begin{tcolorbox}[enhanced,colback=ink!5,colframe=ink!5,arc=1.5pt,boxrule=0pt]
\sffamily\small
\begin{tabular}{@{}l@{\quad}l@{}}
Klassen in diesem Heft & 861 (drei bis sechs Steine)\\
Stellungen darin & 18,7 Billionen (beide Zugrechte)\\
Mit Hilfsmatt & 7,07 Billionen\\
Mit eindeutiger Lösung & 438 Milliarden\\
Längste eindeutige Aufgabe & h\#13: KBvkbbp, KBvkrp, KQBvkrp\\
Längstes Hilfsmatt & h\#17.5: KBvkqqp\\
Gefundene veröffentlichte Aufgaben & 26\\
\end{tabular}
\end{tcolorbox}

Die Tabellen erzeugt ein C++-Generator, der jede Klasse Tiefe um Tiefe vom Matt aus nach außen durcharbeitet, nachdem er zuvor jede Tabelle gebaut hat, in die ein Schlag oder eine Umwandlung führt. Geprüft werden sie durch einen unabhängigen Löser an Stichproben, durch einen vollständigen Abgleich von KQvk mit python-chess und durch die Forderung, dass ein Lauf auf vielen Threads Byte für Byte dieselbe Datei schreibt wie ein Lauf auf einem. Jede von außen beigesteuerte Tabelle durchläuft vor der Veröffentlichung eine Prüfung in sieben Schritten.

\section*{So liest man eine Seite}

Jede Klasse hat eine Seite. Der \textcolor{ink}{Kopf} nennt das Material, die weißen Steine zuerst, und die Forderung der Aufgabe. Darunter steht die Klasse im Überblick: ihre längste eindeutige Aufgabe, wie viele Stellungen diese Länge mit eindeutiger Lösung erreichen, ihr längstes Hilfsmatt und ihre Größe; daneben ein Schaubild aller ihrer Stellungen nach Entfernung zum Matt (Balken), mit der Zahl der eindeutigen Lösungen je Entfernung auf logarithmischer Skala (Linie) und der Länge der Aufgabe selbst gepunktet.

Unter dem Diagramm stehen die Forderung und die Steinzahl, weiße plus schwarze Steine. Die Lösung ist in Hilfsmatt-Reihenfolge nummeriert, der schwarze Zug zuerst. Ist die Stellung oder eine nahe Verwandte als Komposition veröffentlicht, nennt der Kasten unter der Lösung Autor und Quellen. Der Link \emph{lösen} im Kopf und der QR-Code öffnen die Stellung im Helpmate Analyzer (\texttt{helpman.komtera.lt}) mit bereits gesetzter Forderung.

\section*{Urheberschaft und Mitwirkende}

Die Hilfsmatt-Tablebase und dieses Heft sind ein Werk von Oliver Sick, der das Projekt betreut: Generator, Verifikation und Auswahl der Aufgaben.

Die meisten Sechssteiner-Tabellen haben T31M (584 Tabellen) und popeye37 (77 Tabellen) auf eigenen Rechnern berechnet; vor der Veröffentlichung wurden sie geprüft. T31M hat außerdem den Generator und die Suche nach Aufgaben um ein Vielfaches beschleunigt.

\begin{tcolorbox}[enhanced,colback=ink!5,colframe=ink!5,arc=1.5pt,boxrule=0pt]
\sffamily\small
\begin{tabular}{@{}l@{\quad}l@{}}
Betreuer & Oliver Sick (\texttt{osick})\\
Mitwirkende & T31M (584 Tabellen) und popeye37 (77 Tabellen)\\
Code & {\footnotesize\url{https://github.com/osick/helpmate-tablebase}}\\
Tabellen & {\footnotesize\url{https://huggingface.co/datasets/osick/helpmate-tables}}\\
\end{tabular}
\end{tcolorbox}

\section*{Frühere Veröffentlichungen}

Die Suche nach früheren Veröffentlichungen ist nicht vollständig. Jede Stellung wurde mit einer Sammlung veröffentlichter Hilfsmatts abgeglichen; gefundene Treffer stehen auf der Seite der Aufgabe. Fehlt dort ein solcher Kasten, heißt das \emph{nicht}, dass die Stellung unveröffentlicht ist: Für keine Stellung in diesem Heft wird Originalität oder Priorität beansprucht. Veröffentlichte Aufgaben bleiben geistiges Eigentum ihrer Autoren. Hinweise auf übersehene Vorgänger nehmen wir gern als Issue auf GitHub entgegen.

\tableofcontents
\mainmatter

\chapter{Der Bestand in Zahlen}
\markboth{Der Bestand in Zahlen}{}

Dieses Heft umfasst 861 Materialklassen mit 9\,353\,470\,023\,072 indizierten Stellungen. Die längste eindeutige Aufgabe des ganzen Bestands ist \textbf{h\#13} in KQBvkrp (Seite~\pageref{hm:KQBvkrp}).

\section{Abdeckung}

Jede Materialklasse von drei bis sechs Steinen, in der ein Matt möglich ist, ist gelöst. Beide Seitenverteilungen zählen getrennt -- Weiß setzt matt, also sind KRPvkb und KBvkrp verschiedene Aufgaben --, das ergibt auf dem Papier 1\,000 Klassen. Davon haben 125 auf weißer Seite nur den nackten König und können kein Matt enthalten; es bleiben die hier gezählten 875.

\begin{center}
\begin{tabular}{lrrr}
\toprule
Steine & mit eindeutiger Aufgabe & Matt möglich & Anteil \\
\midrule
3 & 3 & 5 & 60,0\,\% \\
4 & 37 & 40 & 92,5\,\% \\
5 & 181 & 185 & 97,8\,\% \\
6 & 640 & 645 & 99,2\,\% \\
\midrule
gesamt & 861 & 875 & 98,4\,\% \\
\bottomrule
\end{tabular}
\end{center}

Eine Klasse kann vollständig gelöst sein und trotzdem keine Seite haben: entweder gibt es in ihr überhaupt kein Matt, oder keine Tiefe enthält eine Stellung mit nur einer Lösung. Eine solche Klasse hat keine eindeutige Aufgabe zu zeigen.

\section{Wie lang eine eindeutige Aufgabe wird}

Verteilung der größten eindeutigen Länge über alle 861 Klassen.
\begin{center}
\begin{tikzpicture}
\begin{axis}[
  width=0.95\linewidth, height=45mm,
  ybar, bar width=5pt,
  xlabel={größte eindeutige Länge (h\#n)}, ylabel={Klassen},
  ylabel near ticks, xlabel near ticks,
  ymin=0, ymax=198,
  tick label style={font=\scriptsize},
  label style={font=\scriptsize},
  axis lines=left, enlarge x limits=0.06,
]
\addplot[fill=ink!55,draw=none] coordinates {(2.0,10) (2.5,3) (3.0,17) (3.5,4) (4.0,26) (4.5,5) (5.0,52) (5.5,27) (6.0,158) (6.5,58) (7.0,172) (7.5,73) (8.0,69) (8.5,52) (9.0,39) (9.5,41) (10.0,21) (10.5,9) (11.0,5) (11.5,6) (12.0,5) (12.5,6) (13.0,3)};
\end{axis}
\end{tikzpicture}
\end{center}

\section{Der Abstand zum Klassenmaximum}

Was Eindeutigkeit an Länge kostet: das längste Hilfsmatt der Klasse minus ihre längste eindeutige Aufgabe, in Halbzügen. Er ist nie null.
\begin{center}
\begin{tikzpicture}
\begin{axis}[
  width=0.95\linewidth, height=45mm,
  ybar, bar width=5pt,
  xlabel={Abstand (Halbzüge)}, ylabel={Klassen},
  ylabel near ticks, xlabel near ticks,
  ymin=0, ymax=351,
  tick label style={font=\scriptsize},
  label style={font=\scriptsize},
  axis lines=left, enlarge x limits=0.06,
]
\addplot[fill=ink!55,draw=none] coordinates {(1,1) (2,305) (3,122) (4,185) (5,23) (6,63) (7,25) (8,35) (9,12) (10,34) (11,7) (12,20) (13,8) (14,11) (16,5) (17,1) (18,3) (20,1)};
\end{axis}
\end{tikzpicture}
\end{center}

Mittel 4,5 Halbzüge, Median 4, größter 20 -- KRBvkp, das bis h\#16 reicht, dessen längste eindeutige Aufgabe aber nur h\#6 ist (Seite~\pageref{hm:KRBvkp}).

\section{Die zehn längsten}

\begin{center}
\begin{tabular}{lrrrr}
\toprule
Material & eindeutig & Klassenmax. & Abstand & Stellungen \\
\midrule
KBvkbbp & h\#13 & h\#16.5 & 7 & 2 \\
KBvkrp & h\#13 & h\#16.5 & 7 & 2 \\
KQBvkrp & h\#13 & h\#17 & 8 & 2 \\
KBvkqbp & h\#12.5 & h\#16.5 & 8 & 59 \\
KBvkqqp & h\#12.5 & h\#17.5 & 10 & 210 \\
KBvkqrp & h\#12.5 & h\#16.5 & 8 & 495 \\
KBvkrpp & h\#12.5 & h\#16.5 & 8 & 14 \\
KBvkrrp & h\#12.5 & h\#16.5 & 8 & 79 \\
KPvkbpp & h\#12.5 & h\#15.5 & 6 & 7 \\
KBBPvkp & h\#12 & h\#15 & 6 & 4 \\
\bottomrule
\end{tabular}
\end{center}

\section{Seltenheit}

In ihrer größten eindeutigen Länge kann eine Klasse eine einzige Stellung enthalten und sonst nichts. 81 Klassen sind so, und 335 enthalten zehn oder weniger. Das sind die Stellungen ganz ohne Spielraum: ein Stein anders, und die Aufgabe ist dahin.

{\footnotesize
\begin{longtable}{lll}
\toprule
\multicolumn{3}{c}{Klassen mit genau einer Stellung in ihrer größten eindeutigen Länge} \\
\midrule
\endhead
KQBvkqp~h\#12 & KBBvkrp~h\#11 & KBPPvkp~h\#11 \\
KBPvkrp~h\#11 & KBvkppp~h\#11 & KQBPvkp~h\#11 \\
KBPvkp~h\#10.5 & KBPvkqp~h\#10.5 & KBvkbnp~h\#10.5 \\
KBvknpp~h\#10.5 & KPPvkpp~h\#10.5 & KBNPvkp~h\#10 \\
KQNvkpp~h\#10 & KRBvkbp~h\#10 & KRBvknp~h\#10 \\
KNvkqb~h\#9.5 & KNvkqrb~h\#9.5 & KNvkrpp~h\#9.5 \\
KPvknnp~h\#9.5 & KBBvkp~h\#9 & KQNvknp~h\#9 \\
KRNPvkp~h\#9 & KNPvknp~h\#8.5 & KNvknp~h\#8.5 \\
KNvkrn~h\#8.5 & KNvkrnn~h\#8.5 & KPvkqn~h\#8.5 \\
KPvkqrn~h\#8.5 & KBNPvkq~h\#8 & KBNvkqb~h\#8 \\
KNNvkpp~h\#8 & KNNvkqp~h\#8 & KNvkqrr~h\#8 \\
KQPvkqr~h\#8 & KRPvkqb~h\#8 & KBBNvkq~h\#7.5 \\
KBNvkq~h\#7.5 & KBNvkrr~h\#7.5 & KBvkrrb~h\#7.5 \\
KPPvkb~h\#7.5 & KPPvkn~h\#7.5 & KPPvkq~h\#7.5 \\
KPPvkrb~h\#7.5 & KPPvkrn~h\#7.5 & KPvkr~h\#7.5 \\
KQNPvkn~h\#7.5 & KRPPvkq~h\#7.5 & KBBNvk~h\#7 \\
KBNNvkq~h\#7 & KBPPPvk~h\#7 & KNNNvkp~h\#7 \\
KNNPvkb~h\#7 & KNNvkq~h\#7 & KNNvkr~h\#7 \\
KQBBNvk~h\#7 & KQBBvk~h\#7 & KRBBvk~h\#7 \\
KRBvkrb~h\#7 & KRNNvkb~h\#7 & KRNNvkq~h\#7 \\
KRPPPvk~h\#7 & KRPPvkb~h\#7 & KRvkqbn~h\#7 \\
KRvkqnn~h\#7 & KRvkqnp~h\#7 & KRvkrpp~h\#7 \\
KBBNNvk~h\#6.5 & KBBvk~h\#6.5 & KNNvkrb~h\#6.5 \\
KNNvkrn~h\#6.5 & KRvknnn~h\#6.5 & KRvkqb~h\#6.5 \\
KNNNvkr~h\#6 & KQvkbn~h\#6 & KRRvkr~h\#6 \\
KQQvkqb~h\#5.5 & KRRNvkq~h\#5.5 & KQvkqrr~h\#5 \\
KRRNvkb~h\#5 & KRvk~h\#4.5 & KQQBvkr~h\#3.5 \\
\bottomrule
\end{longtable}
}

\chapter{Drei Steine}
\markboth{Drei Steine}{}

3 der 5 Materialklassen mit 3 Steinen, in denen ein Matt möglich ist, längste eindeutige Aufgabe zuerst.
\clearpage

\phantomsection\label{hm:KPvk}
\begin{tcolorbox}[enhanced,colback=ink,colframe=ink,arc=0pt,boxrule=0pt,left=4mm,right=4mm,top=2.5mm,bottom=2.5mm]
\color{white}\sffamily{\LARGE\bfseries KPvk}\hfill{\Large h\#6.5}\\[0.5mm]
{\small 3 Steine\quad Weiß am Zug}\hfill{\small Nr.~1\qquad\href{https://helpman.komtera.lt/?fen=6k1/8/8/8/8/8/4P3/2K5&moves=6.5}{\color{white}\bfseries lösen\,\ensuremath{\nearrow}}}
\end{tcolorbox}
\vspace{2mm}
\noindent\begin{minipage}[t]{36mm}\vspace{0pt}\sffamily\small\raggedright
\hmfigure{längste eindeutige}{{\Large h\#6.5}}\hmfigure{eindeutig bei dieser Länge}{{\Large 13} Stellungen}\hmfigure{längstes Matt}{{\Large h\#8.5}}\hmfigure{in der Klasse}{173 Tausend Stellungen,\\ 166 Tausend mit Hilfsmatt}
\end{minipage}\hfill
\begin{minipage}[t]{80mm}\vspace{0pt}\begin{tikzpicture}\begin{axis}[width=72mm,height=42mm,ybar,bar width=1.6pt,xmin=-1,xmax=18,ymin=0,ymax=24.21,axis y line*=left,axis x line*=bottom,xlabel={\scriptsize Halbzüge bis Matt},ylabel={\scriptsize Stellungen (Tsd.)},tick label style={font=\tiny},label style={font=\scriptsize},enlarge x limits=false,ylabel shift=-3pt,xlabel shift=-3pt]\addplot[fill=ink!55,draw=none] coordinates {(1,0.04) (2,0.069) (3,0.515) (4,1.194) (5,4.396) (6,7.289) (7,13.13) (8,16.03) (9,19.32) (10,19.3) (11,21.81) (12,22.01) (13,18.59) (14,16.26) (15,3.849) (16,1.838) (17,0.014)};\draw[densely dotted] (axis cs:13,0) -- (axis cs:13,24.21);\end{axis}\begin{axis}[width=72mm,height=42mm,axis y line*=right,axis x line=none,ymode=log,xmin=-1,xmax=18,ylabel={\scriptsize eindeutig (log)},tick label style={font=\tiny},label style={font=\scriptsize},ylabel shift=-3pt]\addplot[mark=*,mark size=0.7pt,draw=black,thin] coordinates {(1,6) (2,16) (3,159) (4,239) (5,411) (6,308) (7,181) (8,108) (9,102) (10,87) (11,87) (12,65) (13,13)};\end{axis}\end{tikzpicture}\end{minipage}
\vspace{3mm}

\noindent\begin{minipage}[t]{60mm}\vspace{0pt}\sbox\hmboard{\chessboard[setfen=6k1/8/8/8/8/8/4P3/2K5 w - - 0 1]}\usebox\hmboard\\[1mm]
\makebox[\wd\hmboard]{\sffamily\hspace{3mm}\textbf{h\#6.5}\hfill(2+1)}\end{minipage}\hfill
\begin{minipage}[t]{52mm}\vspace{0pt}
\textcolor{ink}{\sffamily\textsc{Lösung}}\par\smallskip\small
\begin{tabular}{@{}r@{\enspace}l@{\enspace}l@{}}
\textbf{1.} & \ldots & e3\\
\textbf{2.} & \symking{}f7 & e4\\
\textbf{3.} & \symking{}e6 & e5\\
\textbf{4.} & \symking{}d5 & e6\\
\textbf{5.} & \symking{}c4 & e7\\
\textbf{6.} & \symking{}b3 & e8=\symqueen{}\\
\textbf{7.} & \symking{}a2 & \symqueen{}a4\#\\
\end{tabular}\par\vspace{3mm}
\qrcode[height=12mm]{https://helpman.komtera.lt/?fen=6k1/8/8/8/8/8/4P3/2K5&moves=6.5}\\{\sffamily\tiny\color{ink}scannen zum Lösen}
\end{minipage}
\vspace{2mm}

{\raggedright\noindent\textcolor{ink}{\sffamily\textsc{Themen}}\enspace \themetag{Satzspiel} \themetag{reines Matt} \themetag{Mustermatt} \themetag{Idealmatt} \themetag{Spiegelmatt} \themetag{Umwandlung} \themetag{Excelsior} \themetag{ein Stein zieht} \themetag{ohne Schlag} \themetag{ohne Schach} \themetag{Umnov}\par}

\begin{tcolorbox}[enhanced,colback=white,colframe=ink!35,arc=1.5pt,boxrule=0.4pt,left=3mm,right=3mm,top=1.5mm,bottom=1.5mm]
\textcolor{ink}{\sffamily\textsc{veröffentlicht}}\par\smallskip\small \textbf{Niemann, John}, h\#6\quad{\color{ink}\scriptsize diese Stellung}\\{\scriptsize Schachmatt, No. 427, 13/07/1947; Catalogue h\# (CHM) avec 6 pieces, No. 61, 1995; Wenigsteiner im Hilfsmatt, No. 77, 1977; \ldots}\par\smallskip\textbf{Ban, Jenö}, h\#7 0.1...\quad{\color{ink}\scriptsize verwandte Stellung}\\{\scriptsize Magyar Sakkelet, 1961; Ideal-Mate Encyclopedia Vol.1, No. 12, 1999}
\end{tcolorbox}
\clearpage
\phantomsection\label{hm:KQvk}
\begin{tcolorbox}[enhanced,colback=ink,colframe=ink,arc=0pt,boxrule=0pt,left=4mm,right=4mm,top=2.5mm,bottom=2.5mm]
\color{white}\sffamily{\LARGE\bfseries KQvk}\hfill{\Large h\#6}\\[0.5mm]
{\small 3 Steine\quad Schwarz am Zug}\hfill{\small Nr.~2\qquad\href{https://helpman.komtera.lt/?fen=8/8/7k/6Q1/8/8/8/K7&moves=6}{\color{white}\bfseries lösen\,\ensuremath{\nearrow}}}
\end{tcolorbox}
\vspace{2mm}
\noindent\begin{minipage}[t]{36mm}\vspace{0pt}\sffamily\small\raggedright
\hmfigure{längste eindeutige}{{\Large h\#6}}\hmfigure{eindeutig bei dieser Länge}{{\Large 3} Stellungen}\hmfigure{längstes Matt}{{\Large h\#7}}\hmfigure{in der Klasse}{59,1 Tausend Stellungen,\\ 45,7 Tausend mit Hilfsmatt}
\end{minipage}\hfill
\begin{minipage}[t]{80mm}\vspace{0pt}\begin{tikzpicture}\begin{axis}[width=72mm,height=42mm,ybar,bar width=1.6pt,xmin=-1,xmax=15,ymin=0,ymax=12.33,axis y line*=left,axis x line*=bottom,xlabel={\scriptsize Halbzüge bis Matt},ylabel={\scriptsize Stellungen (Tsd.)},tick label style={font=\tiny},label style={font=\scriptsize},enlarge x limits=false,ylabel shift=-3pt,xlabel shift=-3pt]\addplot[fill=ink!55,draw=none] coordinates {(0,0.046) (1,0.306) (2,0.58) (3,1.751) (4,4.336) (5,5.851) (6,11.21) (7,7.267) (8,9.197) (9,2.551) (10,1.987) (11,0.353) (12,0.284) (13,0.002) (14,0.004)};\draw[densely dotted] (axis cs:12,0) -- (axis cs:12,12.33);\end{axis}\begin{axis}[width=72mm,height=42mm,axis y line*=right,axis x line=none,ymode=log,xmin=-1,xmax=15,ylabel={\scriptsize eindeutig (log)},tick label style={font=\tiny},label style={font=\scriptsize},ylabel shift=-3pt]\addplot[mark=*,mark size=0.7pt,draw=black,thin] coordinates {(0,46) (1,182) (2,356) (3,582) (4,878) (5,469) (6,373) (7,110) (8,57) (9,3) (10,3) (11,2) (12,3)};\end{axis}\end{tikzpicture}\end{minipage}
\vspace{3mm}

\noindent\begin{minipage}[t]{60mm}\vspace{0pt}\sbox\hmboard{\chessboard[setfen=8/8/7k/6Q1/8/8/8/K7 b - - 0 1]}\usebox\hmboard\\[1mm]
\makebox[\wd\hmboard]{\sffamily\hspace{3mm}\textbf{h\#6}\hfill(2+1)}\end{minipage}\hfill
\begin{minipage}[t]{52mm}\vspace{0pt}
\textcolor{ink}{\sffamily\textsc{Lösung}}\par\smallskip\small
\begin{tabular}{@{}r@{\enspace}l@{\enspace}l@{}}
\textbf{1.} & \symking{}h7 & \symking{}b2\\
\textbf{2.} & \symking{}h8 & \symking{}c3\\
\textbf{3.} & \symking{}h7 & \symking{}d4\\
\textbf{4.} & \symking{}h8 & \symking{}e5\\
\textbf{5.} & \symking{}h7 & \symking{}f6\\
\textbf{6.} & \symking{}h8 & \symqueen{}g7\#\\
\end{tabular}\par\vspace{3mm}
\qrcode[height=12mm]{https://helpman.komtera.lt/?fen=8/8/7k/6Q1/8/8/8/K7&moves=6}\\{\sffamily\tiny\color{ink}scannen zum Lösen}
\end{minipage}
\vspace{2mm}

{\raggedright\noindent\textcolor{ink}{\sffamily\textsc{Themen}}\enspace \themetag{reines Matt} \themetag{Mustermatt} \themetag{Idealmatt} \themetag{Switchback} \themetag{ein Stein zieht} \themetag{ohne Schlag} \themetag{ohne Schach}\par\smallskip{\footnotesize\textit{Hinweis:} die Partei am Zug steht im Diagramm im Schach.}\par}

\begin{tcolorbox}[enhanced,colback=white,colframe=ink!35,arc=1.5pt,boxrule=0.4pt,left=3mm,right=3mm,top=1.5mm,bottom=1.5mm]
\textcolor{ink}{\sffamily\textsc{veröffentlicht}}\par\smallskip\small \textbf{Csaszar, Imre}, h\#6\quad{\color{ink}\scriptsize verwandte Stellung}\\{\scriptsize Magyar Sakkvilag, 1945; Ideal-Mate Review, 1990; Catalogue h\# (CHM) avec 6 pieces, No. 165, 1995}
\end{tcolorbox}
\clearpage
\phantomsection\label{hm:KRvk}
\begin{tcolorbox}[enhanced,colback=ink,colframe=ink,arc=0pt,boxrule=0pt,left=4mm,right=4mm,top=2.5mm,bottom=2.5mm]
\color{white}\sffamily{\LARGE\bfseries KRvk}\hfill{\Large h\#4.5}\\[0.5mm]
{\small 3 Steine\quad Weiß am Zug}\hfill{\small Nr.~3\qquad\href{https://helpman.komtera.lt/?fen=4k3/8/8/8/8/8/2R5/1K6&moves=4.5}{\color{white}\bfseries lösen\,\ensuremath{\nearrow}}}
\end{tcolorbox}
\vspace{2mm}
\noindent\begin{minipage}[t]{36mm}\vspace{0pt}\sffamily\small\raggedright
\hmfigure{längste eindeutige}{{\Large h\#4.5}}\hmfigure{eindeutig bei dieser Länge}{{\Large 1} Stellung}\hmfigure{längstes Matt}{{\Large h\#7}}\hmfigure{in der Klasse}{59,1 Tausend Stellungen,\\ 50 Tausend mit Hilfsmatt}
\end{minipage}\hfill
\begin{minipage}[t]{80mm}\vspace{0pt}\begin{tikzpicture}\begin{axis}[width=72mm,height=42mm,ybar,bar width=1.6pt,xmin=-1,xmax=15,ymin=0,ymax=11.14,axis y line*=left,axis x line*=bottom,xlabel={\scriptsize Halbzüge bis Matt},ylabel={\scriptsize Stellungen (Tsd.)},tick label style={font=\tiny},label style={font=\scriptsize},enlarge x limits=false,ylabel shift=-3pt,xlabel shift=-3pt]\addplot[fill=ink!55,draw=none] coordinates {(0,0.027) (1,0.189) (2,0.294) (3,1.372) (4,3.27) (5,5.248) (6,9.039) (7,7.893) (8,10.13) (9,6.181) (10,4.473) (11,1.041) (12,0.734) (13,0.035) (14,0.027)};\draw[densely dotted] (axis cs:9,0) -- (axis cs:9,11.14);\end{axis}\begin{axis}[width=72mm,height=42mm,axis y line*=right,axis x line=none,ymode=log,xmin=-1,xmax=15,ylabel={\scriptsize eindeutig (log)},tick label style={font=\tiny},label style={font=\scriptsize},ylabel shift=-3pt]\addplot[mark=*,mark size=0.7pt,draw=black,thin] coordinates {(0,27) (1,189) (2,294) (3,893) (4,1405) (5,664) (6,379) (7,36) (8,22) (9,1)};\end{axis}\end{tikzpicture}\end{minipage}
\vspace{3mm}

\noindent\begin{minipage}[t]{60mm}\vspace{0pt}\sbox\hmboard{\chessboard[setfen=4k3/8/8/8/8/8/2R5/1K6 w - - 0 1]}\usebox\hmboard\\[1mm]
\makebox[\wd\hmboard]{\sffamily\hspace{3mm}\textbf{h\#4.5}\hfill(2+1)}\end{minipage}\hfill
\begin{minipage}[t]{52mm}\vspace{0pt}
\textcolor{ink}{\sffamily\textsc{Lösung}}\par\smallskip\small
\begin{tabular}{@{}r@{\enspace}l@{\enspace}l@{}}
\textbf{1.} & \ldots & \symking{}c1\\
\textbf{2.} & \symking{}f7 & \symking{}d2\\
\textbf{3.} & \symking{}g6 & \symking{}e3\\
\textbf{4.} & \symking{}h5 & \symking{}f4\\
\textbf{5.} & \symking{}h4 & \symrook{}h2\#\\
\end{tabular}\par\vspace{3mm}
\qrcode[height=12mm]{https://helpman.komtera.lt/?fen=4k3/8/8/8/8/8/2R5/1K6&moves=4.5}\\{\sffamily\tiny\color{ink}scannen zum Lösen}
\end{minipage}
\vspace{2mm}

{\raggedright\noindent\textcolor{ink}{\sffamily\textsc{Themen}}\enspace \themetag{reines Matt} \themetag{Mustermatt} \themetag{Idealmatt} \themetag{Spiegelmatt} \themetag{ein Stein zieht} \themetag{ohne Schlag} \themetag{ohne Schach}\par}

\begin{tcolorbox}[enhanced,colback=white,colframe=ink!35,arc=1.5pt,boxrule=0.4pt,left=3mm,right=3mm,top=1.5mm,bottom=1.5mm]
\textcolor{ink}{\sffamily\textsc{veröffentlicht}}\par\smallskip\small \textbf{Ban, Jenö}, h\#5 0.1...\quad{\color{ink}\scriptsize verwandte Stellung}\\{\scriptsize Diagramme und Figuren, No. 1295, 27/11/1965; Catalogue h\# (CHM) avec 6 pieces, No. 122, 1995; Ideal-Mate Encyclopedia Vol.1, No. 1848, 1999}
\end{tcolorbox}
\clearpage

\chapter{Vier Steine}
\markboth{Vier Steine}{}

37 der 40 Materialklassen mit 4 Steinen, in denen ein Matt möglich ist, längste eindeutige Aufgabe zuerst.
\clearpage

\phantomsection\label{hm:KPvkp}
\begin{tcolorbox}[enhanced,colback=ink,colframe=ink,arc=0pt,boxrule=0pt,left=4mm,right=4mm,top=2.5mm,bottom=2.5mm]
\color{white}\sffamily{\LARGE\bfseries KPvkp}\hfill{\Large h\#8.5}\\[0.5mm]
{\small 4 Steine\quad Weiß am Zug}\hfill{\small Nr.~4\qquad\href{https://helpman.komtera.lt/?fen=8/8/p7/8/P7/8/3K4/7k&moves=8.5}{\color{white}\bfseries lösen\,\ensuremath{\nearrow}}}
\end{tcolorbox}
\vspace{2mm}
\noindent\begin{minipage}[t]{36mm}\vspace{0pt}\sffamily\small\raggedright
\hmfigure{längste eindeutige}{{\Large h\#8.5}}\hmfigure{eindeutig bei dieser Länge}{{\Large 2} Stellungen}\hmfigure{längstes Matt}{{\Large h\#14.5}}\hmfigure{in der Klasse}{8,32 Millionen Stellungen,\\ 7,44 Millionen mit Hilfsmatt}
\end{minipage}\hfill
\begin{minipage}[t]{80mm}\vspace{0pt}\begin{tikzpicture}\begin{axis}[width=72mm,height=42mm,ybar,bar width=1.6pt,xmin=-1,xmax=30,ymin=0,ymax=1072,axis y line*=left,axis x line*=bottom,xlabel={\scriptsize Halbzüge bis Matt},ylabel={\scriptsize Stellungen (Tsd.)},tick label style={font=\tiny},label style={font=\scriptsize},enlarge x limits=false,ylabel shift=-3pt,xlabel shift=-3pt]\addplot[fill=ink!55,draw=none] coordinates {(0,0.002) (1,1.863) (2,3.908) (3,25.98) (4,59.05) (5,210.2) (6,344.4) (7,615.3) (8,742.4) (9,895.5) (10,858.5) (11,974.2) (12,935.4) (13,739.8) (14,572.2) (15,107.1) (16,55.51) (17,43.37) (18,44.31) (19,47.19) (20,44.77) (21,29.95) (22,29.75) (23,14.29) (24,14.29) (25,9.156) (26,9.294) (27,4.036) (28,4.068) (29,0.03)};\draw[densely dotted] (axis cs:17,0) -- (axis cs:17,1072);\end{axis}\begin{axis}[width=72mm,height=42mm,axis y line*=right,axis x line=none,ymode=log,xmin=-1,xmax=30,ylabel={\scriptsize eindeutig (log)},tick label style={font=\tiny},label style={font=\scriptsize},ylabel shift=-3pt]\addplot[mark=*,mark size=0.7pt,draw=black,thin] coordinates {(0,2) (1,332) (2,623) (3,5739) (4,9935) (5,18250) (6,14116) (7,9142) (8,5997) (9,5909) (10,4674) (11,5251) (12,3850) (13,442) (14,58) (15,16) (16,6) (17,2)};\end{axis}\end{tikzpicture}\end{minipage}
\vspace{3mm}

\noindent\begin{minipage}[t]{60mm}\vspace{0pt}\sbox\hmboard{\chessboard[setfen=8/8/p7/8/P7/8/3K4/7k w - - 0 1]}\usebox\hmboard\\[1mm]
\makebox[\wd\hmboard]{\sffamily\hspace{3mm}\textbf{h\#8.5}\hfill(2+2)}\end{minipage}\hfill
\begin{minipage}[t]{52mm}\vspace{0pt}
\textcolor{ink}{\sffamily\textsc{Lösung}}\par\smallskip\small
\begin{tabular}{@{}r@{\enspace}l@{\enspace}l@{}}
\textbf{1.} & \ldots & \symking{}c3\\
\textbf{2.} & \symking{}g2 & \symking{}b4\\
\textbf{3.} & a5+ & \symking{}xa5\\
\textbf{4.} & \symking{}f3 & \symking{}b4\\
\textbf{5.} & \symking{}e4 & a5\\
\textbf{6.} & \symking{}d5 & a6\\
\textbf{7.} & \symking{}c6 & a7\\
\textbf{8.} & \symking{}b7 & \symking{}c5\\
\textbf{9.} & \symking{}a6 & a8=\symqueen{}\#\\
\end{tabular}\par\vspace{3mm}
\qrcode[height=12mm]{https://helpman.komtera.lt/?fen=8/8/p7/8/P7/8/3K4/7k&moves=8.5}\\{\sffamily\tiny\color{ink}scannen zum Lösen}
\end{minipage}
\vspace{2mm}

{\raggedright\noindent\textcolor{ink}{\sffamily\textsc{Themen}}\enspace \themetag{reines Matt} \themetag{Mustermatt} \themetag{Idealmatt} \themetag{Spiegelmatt} \themetag{Umwandlung} \themetag{Switchback}\par}

\clearpage
\phantomsection\label{hm:KBvkn}
\begin{tcolorbox}[enhanced,colback=ink,colframe=ink,arc=0pt,boxrule=0pt,left=4mm,right=4mm,top=2.5mm,bottom=2.5mm]
\color{white}\sffamily{\LARGE\bfseries KBvkn}\hfill{\Large h\#8}\\[0.5mm]
{\small 4 Steine\quad Schwarz am Zug}\hfill{\small Nr.~5\qquad\href{https://helpman.komtera.lt/?fen=8/8/8/8/8/1k6/5n2/1K5B&moves=8}{\color{white}\bfseries lösen\,\ensuremath{\nearrow}}}
\end{tcolorbox}
\vspace{2mm}
\noindent\begin{minipage}[t]{36mm}\vspace{0pt}\sffamily\small\raggedright
\hmfigure{längste eindeutige}{{\Large h\#8}}\hmfigure{eindeutig bei dieser Länge}{{\Large 19} Stellungen}\hmfigure{längstes Matt}{{\Large h\#12}}\hmfigure{in der Klasse}{3,78 Millionen Stellungen,\\ 3,05 Millionen mit Hilfsmatt}
\end{minipage}\hfill
\begin{minipage}[t]{80mm}\vspace{0pt}\begin{tikzpicture}\begin{axis}[width=72mm,height=42mm,ybar,bar width=1.6pt,xmin=-1,xmax=25,ymin=0,ymax=507.7,axis y line*=left,axis x line*=bottom,xlabel={\scriptsize Halbzüge bis Matt},ylabel={\scriptsize Stellungen (Tsd.)},tick label style={font=\tiny},label style={font=\scriptsize},enlarge x limits=false,ylabel shift=-3pt,xlabel shift=-3pt]\addplot[fill=ink!55,draw=none] coordinates {(0,0.001) (1,0.002) (2,0.003) (3,0.029) (4,0.199) (5,1.022) (6,3.221) (7,9.96) (8,25.81) (9,49.69) (10,99.44) (11,151.7) (12,258.9) (13,322.3) (14,460.1) (15,434.6) (16,461.6) (17,278.9) (18,185) (19,159.5) (20,64.95) (21,63.06) (22,8.045) (23,8.377) (24,0.004)};\draw[densely dotted] (axis cs:16,0) -- (axis cs:16,507.7);\end{axis}\begin{axis}[width=72mm,height=42mm,axis y line*=right,axis x line=none,ymode=log,xmin=-1,xmax=25,ylabel={\scriptsize eindeutig (log)},tick label style={font=\tiny},label style={font=\scriptsize},ylabel shift=-3pt]\addplot[mark=*,mark size=0.7pt,draw=black,thin] coordinates {(0,1) (1,2) (2,3) (3,27) (4,165) (5,305) (6,341) (7,364) (8,528) (9,596) (10,856) (11,454) (12,486) (13,204) (14,74) (15,45) (16,19)};\end{axis}\end{tikzpicture}\end{minipage}
\vspace{3mm}

\noindent\begin{minipage}[t]{60mm}\vspace{0pt}\sbox\hmboard{\chessboard[setfen=8/8/8/8/8/1k6/5n2/1K5B b - - 0 1]}\usebox\hmboard\\[1mm]
\makebox[\wd\hmboard]{\sffamily\hspace{3mm}\textbf{h\#8}\hfill(2+2)}\end{minipage}\hfill
\begin{minipage}[t]{52mm}\vspace{0pt}
\textcolor{ink}{\sffamily\textsc{Lösung}}\par\smallskip\small
\begin{tabular}{@{}r@{\enspace}l@{\enspace}l@{}}
\textbf{1.} & \symking{}c3 & \symbishop{}g2\\
\textbf{2.} & \symking{}d2 & \symbishop{}f1\\
\textbf{3.} & \symking{}e1 & \symking{}c2\\
\textbf{4.} & \symknight{}h3 & \symking{}d3\\
\textbf{5.} & \symking{}f2 & \symking{}e4\\
\textbf{6.} & \symking{}g1 & \symking{}f3\\
\textbf{7.} & \symking{}h1 & \symking{}g3\\
\textbf{8.} & \symknight{}g1 & \symbishop{}g2\#\\
\end{tabular}\par\vspace{3mm}
\qrcode[height=12mm]{https://helpman.komtera.lt/?fen=8/8/8/8/8/1k6/5n2/1K5B&moves=8}\\{\sffamily\tiny\color{ink}scannen zum Lösen}
\end{minipage}
\vspace{2mm}

{\raggedright\noindent\textcolor{ink}{\sffamily\textsc{Themen}}\enspace \themetag{reines Matt} \themetag{Mustermatt} \themetag{Idealmatt} \themetag{Switchback} \themetag{Selbstblock} \themetag{ohne Schlag} \themetag{ohne Schach}\par}

\clearpage
\phantomsection\label{hm:KBvkp}
\begin{tcolorbox}[enhanced,colback=ink,colframe=ink,arc=0pt,boxrule=0pt,left=4mm,right=4mm,top=2.5mm,bottom=2.5mm]
\color{white}\sffamily{\LARGE\bfseries KBvkp}\hfill{\Large h\#8}\\[0.5mm]
{\small 4 Steine\quad Schwarz am Zug}\hfill{\small Nr.~6\qquad\href{https://helpman.komtera.lt/?fen=8/8/8/4k3/2p5/8/8/K3B3&moves=8}{\color{white}\bfseries lösen\,\ensuremath{\nearrow}}}
\end{tcolorbox}
\vspace{2mm}
\noindent\begin{minipage}[t]{36mm}\vspace{0pt}\sffamily\small\raggedright
\hmfigure{längste eindeutige}{{\Large h\#8}}\hmfigure{eindeutig bei dieser Länge}{{\Large 16} Stellungen}\hmfigure{längstes Matt}{{\Large h\#16}}\hmfigure{in der Klasse}{11,1 Millionen Stellungen,\\ 9,43 Millionen mit Hilfsmatt}
\end{minipage}\hfill
\begin{minipage}[t]{80mm}\vspace{0pt}\begin{tikzpicture}\begin{axis}[width=72mm,height=42mm,ybar,bar width=1.6pt,xmin=-1,xmax=33,ymin=0,ymax=847.1,axis y line*=left,axis x line*=bottom,xlabel={\scriptsize Halbzüge bis Matt},ylabel={\scriptsize Stellungen (Tsd.)},tick label style={font=\tiny},label style={font=\scriptsize},enlarge x limits=false,ylabel shift=-3pt,xlabel shift=-3pt]\addplot[fill=ink!55,draw=none] coordinates {(0,0.028) (1,0.104) (2,0.112) (3,0.483) (4,1.528) (5,4.113) (6,13.79) (7,26.22) (8,66.13) (9,95.9) (10,202.5) (11,253.4) (12,472.5) (13,415.5) (14,673.5) (15,555.7) (16,770.1) (17,641) (18,742.9) (19,634.7) (20,687.3) (21,610) (22,579.4) (23,514.1) (24,386.7) (25,344) (26,239.6) (27,215.7) (28,103.7) (29,95.53) (30,41.35) (31,39.38) (32,0.109)};\draw[densely dotted] (axis cs:16,0) -- (axis cs:16,847.1);\end{axis}\begin{axis}[width=72mm,height=42mm,axis y line*=right,axis x line=none,ymode=log,xmin=-1,xmax=33,ylabel={\scriptsize eindeutig (log)},tick label style={font=\tiny},label style={font=\scriptsize},ylabel shift=-3pt]\addplot[mark=*,mark size=0.7pt,draw=black,thin] coordinates {(0,28) (1,104) (2,110) (3,229) (4,721) (5,932) (6,2085) (7,1954) (8,1920) (9,1682) (10,1813) (11,1286) (12,1215) (13,132) (14,116) (15,17) (16,16)};\end{axis}\end{tikzpicture}\end{minipage}
\vspace{3mm}

\noindent\begin{minipage}[t]{60mm}\vspace{0pt}\sbox\hmboard{\chessboard[setfen=8/8/8/4k3/2p5/8/8/K3B3 b - - 0 1]}\usebox\hmboard\\[1mm]
\makebox[\wd\hmboard]{\sffamily\hspace{3mm}\textbf{h\#8}\hfill(2+2)}\end{minipage}\hfill
\begin{minipage}[t]{52mm}\vspace{0pt}
\textcolor{ink}{\sffamily\textsc{Lösung}}\par\smallskip\small
\begin{tabular}{@{}r@{\enspace}l@{\enspace}l@{}}
\textbf{1.} & c3 & \symbishop{}h4\\
\textbf{2.} & c2 & \symking{}b2\\
\textbf{3.} & c1=\symknight{} & \symking{}c2\\
\textbf{4.} & \symking{}d4 & \symking{}d1\\
\textbf{5.} & \symking{}c3 & \symbishop{}e7\\
\textbf{6.} & \symking{}b2 & \symbishop{}a3+\\
\textbf{7.} & \symking{}a1 & \symking{}c2\\
\textbf{8.} & \symknight{}a2 & \symbishop{}b2\#\\
\end{tabular}\par\vspace{3mm}
\qrcode[height=12mm]{https://helpman.komtera.lt/?fen=8/8/8/4k3/2p5/8/8/K3B3&moves=8}\\{\sffamily\tiny\color{ink}scannen zum Lösen}
\end{minipage}
\vspace{2mm}

{\raggedright\noindent\textcolor{ink}{\sffamily\textsc{Themen}}\enspace \themetag{reines Matt} \themetag{Mustermatt} \themetag{Idealmatt} \themetag{Umwandlung} \themetag{Unterverwandlung} \themetag{Switchback} \themetag{Selbstblock} \themetag{ohne Schlag} \themetag{Umnov}\par}

\clearpage
\phantomsection\label{hm:KNvkb}
\begin{tcolorbox}[enhanced,colback=ink,colframe=ink,arc=0pt,boxrule=0pt,left=4mm,right=4mm,top=2.5mm,bottom=2.5mm]
\color{white}\sffamily{\LARGE\bfseries KNvkb}\hfill{\Large h\#8}\\[0.5mm]
{\small 4 Steine\quad Schwarz am Zug}\hfill{\small Nr.~7\qquad\href{https://helpman.komtera.lt/?fen=8/8/7k/8/8/1b6/4N3/K7&moves=8}{\color{white}\bfseries lösen\,\ensuremath{\nearrow}}}
\end{tcolorbox}
\vspace{2mm}
\noindent\begin{minipage}[t]{36mm}\vspace{0pt}\sffamily\small\raggedright
\hmfigure{längste eindeutige}{{\Large h\#8}}\hmfigure{eindeutig bei dieser Länge}{{\Large 5} Stellungen}\hmfigure{längstes Matt}{{\Large h\#10.5}}\hmfigure{in der Klasse}{3,78 Millionen Stellungen,\\ 3,05 Millionen mit Hilfsmatt}
\end{minipage}\hfill
\begin{minipage}[t]{80mm}\vspace{0pt}\begin{tikzpicture}\begin{axis}[width=72mm,height=42mm,ybar,bar width=1.6pt,xmin=-1,xmax=22,ymin=0,ymax=520.7,axis y line*=left,axis x line*=bottom,xlabel={\scriptsize Halbzüge bis Matt},ylabel={\scriptsize Stellungen (Tsd.)},tick label style={font=\tiny},label style={font=\scriptsize},enlarge x limits=false,ylabel shift=-3pt,xlabel shift=-3pt]\addplot[fill=ink!55,draw=none] coordinates {(0,0.001) (1,0.005) (2,0.035) (3,0.213) (4,1.145) (5,4.812) (6,10.67) (7,30.93) (8,51.16) (9,109.6) (10,174) (11,283.3) (12,375.7) (13,464.5) (14,473.4) (15,428.7) (16,320.5) (17,206.5) (18,72.57) (19,38.62) (20,0.053) (21,0.006)};\draw[densely dotted] (axis cs:16,0) -- (axis cs:16,520.7);\end{axis}\begin{axis}[width=72mm,height=42mm,axis y line*=right,axis x line=none,ymode=log,xmin=-1,xmax=22,ylabel={\scriptsize eindeutig (log)},tick label style={font=\tiny},label style={font=\scriptsize},ylabel shift=-3pt]\addplot[mark=*,mark size=0.7pt,draw=black,thin] coordinates {(0,1) (1,5) (2,35) (3,159) (4,629) (5,899) (6,585) (7,690) (8,955) (9,1368) (10,1703) (11,1011) (12,557) (13,139) (14,112) (15,6) (16,5)};\end{axis}\end{tikzpicture}\end{minipage}
\vspace{3mm}

\noindent\begin{minipage}[t]{60mm}\vspace{0pt}\sbox\hmboard{\chessboard[setfen=8/8/7k/8/8/1b6/4N3/K7 b - - 0 1]}\usebox\hmboard\\[1mm]
\makebox[\wd\hmboard]{\sffamily\hspace{3mm}\textbf{h\#8}\hfill(2+2)}\end{minipage}\hfill
\begin{minipage}[t]{52mm}\vspace{0pt}
\textcolor{ink}{\sffamily\textsc{Lösung}}\par\smallskip\small
\begin{tabular}{@{}r@{\enspace}l@{\enspace}l@{}}
\textbf{1.} & \symking{}g5 & \symknight{}d4\\
\textbf{2.} & \symking{}f4 & \symking{}b2\\
\textbf{3.} & \symking{}e3 & \symking{}a3\\
\textbf{4.} & \symking{}d2 & \symking{}b4\\
\textbf{5.} & \symking{}c1 & \symking{}c3\\
\textbf{6.} & \symking{}b1 & \symking{}d2\\
\textbf{7.} & \symking{}a1 & \symking{}c1\\
\textbf{8.} & \symbishop{}a2 & \symknight{}c2\#\\
\end{tabular}\par\vspace{3mm}
\qrcode[height=12mm]{https://helpman.komtera.lt/?fen=8/8/7k/8/8/1b6/4N3/K7&moves=8}\\{\sffamily\tiny\color{ink}scannen zum Lösen}
\end{minipage}
\vspace{2mm}

{\raggedright\noindent\textcolor{ink}{\sffamily\textsc{Themen}}\enspace \themetag{reines Matt} \themetag{Mustermatt} \themetag{Idealmatt} \themetag{Selbstblock} \themetag{ohne Schlag} \themetag{ohne Schach}\par}

\clearpage
\phantomsection\label{hm:KNvkp}
\begin{tcolorbox}[enhanced,colback=ink,colframe=ink,arc=0pt,boxrule=0pt,left=4mm,right=4mm,top=2.5mm,bottom=2.5mm]
\color{white}\sffamily{\LARGE\bfseries KNvkp}\hfill{\Large h\#8}\\[0.5mm]
{\small 4 Steine\quad Schwarz am Zug}\hfill{\small Nr.~8\qquad\href{https://helpman.komtera.lt/?fen=8/8/3p4/8/8/1k6/N7/K7&moves=8}{\color{white}\bfseries lösen\,\ensuremath{\nearrow}}}
\end{tcolorbox}
\vspace{2mm}
\noindent\begin{minipage}[t]{36mm}\vspace{0pt}\sffamily\small\raggedright
\hmfigure{längste eindeutige}{{\Large h\#8}}\hmfigure{eindeutig bei dieser Länge}{{\Large 4} Stellungen}\hmfigure{längstes Matt}{{\Large h\#11.5}}\hmfigure{in der Klasse}{11,1 Millionen Stellungen,\\ 9,69 Millionen mit Hilfsmatt}
\end{minipage}\hfill
\begin{minipage}[t]{80mm}\vspace{0pt}\begin{tikzpicture}\begin{axis}[width=72mm,height=42mm,ybar,bar width=1.6pt,xmin=-1,xmax=24,ymin=0,ymax=1.448,axis y line*=left,axis x line*=bottom,xlabel={\scriptsize Halbzüge bis Matt},ylabel={\scriptsize Stellungen (Mio.)},tick label style={font=\tiny},label style={font=\scriptsize},enlarge x limits=false,ylabel shift=-3pt,xlabel shift=-3pt]\addplot[fill=ink!55,draw=none] coordinates {(0,4e-06) (1,1.8e-05) (2,2.5e-05) (3,0.000157) (4,0.000591) (5,0.002408) (6,0.009652) (7,0.02658) (8,0.0724) (9,0.1418) (10,0.3082) (11,0.4708) (12,0.831) (13,0.9817) (14,1.286) (15,1.194) (16,1.317) (17,1.026) (18,0.882) (19,0.6582) (20,0.2724) (21,0.2023) (22,0.00284) (23,1.5e-05)};\draw[densely dotted] (axis cs:16,0) -- (axis cs:16,1.448);\end{axis}\begin{axis}[width=72mm,height=42mm,axis y line*=right,axis x line=none,ymode=log,xmin=-1,xmax=24,ylabel={\scriptsize eindeutig (log)},tick label style={font=\tiny},label style={font=\scriptsize},ylabel shift=-3pt]\addplot[mark=*,mark size=0.7pt,draw=black,thin] coordinates {(0,4) (1,17) (2,23) (3,96) (4,418) (5,481) (6,1179) (7,990) (8,1719) (9,1580) (10,2654) (11,1583) (12,1651) (13,286) (14,274) (15,14) (16,4)};\end{axis}\end{tikzpicture}\end{minipage}
\vspace{3mm}

\noindent\begin{minipage}[t]{60mm}\vspace{0pt}\sbox\hmboard{\chessboard[setfen=8/8/3p4/8/8/1k6/N7/K7 b - - 0 1]}\usebox\hmboard\\[1mm]
\makebox[\wd\hmboard]{\sffamily\hspace{3mm}\textbf{h\#8}\hfill(2+2)}\end{minipage}\hfill
\begin{minipage}[t]{52mm}\vspace{0pt}
\textcolor{ink}{\sffamily\textsc{Lösung}}\par\smallskip\small
\begin{tabular}{@{}r@{\enspace}l@{\enspace}l@{}}
\textbf{1.} & d5 & \symking{}b1\\
\textbf{2.} & d4 & \symking{}c1\\
\textbf{3.} & d3 & \symking{}d2\\
\textbf{4.} & \symking{}b2 & \symknight{}b4\\
\textbf{5.} & \symking{}a1 & \symking{}c3\\
\textbf{6.} & d2 & \symking{}b3\\
\textbf{7.} & d1=\symrook{} & \symking{}a3\\
\textbf{8.} & \symrook{}b1 & \symknight{}c2\#\\
\end{tabular}\par\vspace{3mm}
\qrcode[height=12mm]{https://helpman.komtera.lt/?fen=8/8/3p4/8/8/1k6/N7/K7&moves=8}\\{\sffamily\tiny\color{ink}scannen zum Lösen}
\end{minipage}
\vspace{2mm}

{\raggedright\noindent\textcolor{ink}{\sffamily\textsc{Themen}}\enspace \themetag{reines Matt} \themetag{Mustermatt} \themetag{Idealmatt} \themetag{Umwandlung} \themetag{Unterverwandlung} \themetag{Selbstblock} \themetag{ohne Schlag} \themetag{ohne Schach} \themetag{Umnov}\par}

\clearpage
\phantomsection\label{hm:KNvkr}
\begin{tcolorbox}[enhanced,colback=ink,colframe=ink,arc=0pt,boxrule=0pt,left=4mm,right=4mm,top=2.5mm,bottom=2.5mm]
\color{white}\sffamily{\LARGE\bfseries KNvkr}\hfill{\Large h\#8}\\[0.5mm]
{\small 4 Steine\quad Schwarz am Zug}\hfill{\small Nr.~9\qquad\href{https://helpman.komtera.lt/?fen=8/8/8/8/8/7k/5r1N/K7&moves=8}{\color{white}\bfseries lösen\,\ensuremath{\nearrow}}}
\end{tcolorbox}
\vspace{2mm}
\noindent\begin{minipage}[t]{36mm}\vspace{0pt}\sffamily\small\raggedright
\hmfigure{längste eindeutige}{{\Large h\#8}}\hmfigure{eindeutig bei dieser Länge}{{\Large 9} Stellungen}\hmfigure{längstes Matt}{{\Large h\#9.5}}\hmfigure{in der Klasse}{3,78 Millionen Stellungen,\\ 2,91 Millionen mit Hilfsmatt}
\end{minipage}\hfill
\begin{minipage}[t]{80mm}\vspace{0pt}\begin{tikzpicture}\begin{axis}[width=72mm,height=42mm,ybar,bar width=1.6pt,xmin=-1,xmax=20,ymin=0,ymax=559.5,axis y line*=left,axis x line*=bottom,xlabel={\scriptsize Halbzüge bis Matt},ylabel={\scriptsize Stellungen (Tsd.)},tick label style={font=\tiny},label style={font=\scriptsize},enlarge x limits=false,ylabel shift=-3pt,xlabel shift=-3pt]\addplot[fill=ink!55,draw=none] coordinates {(0,0.001) (1,0.004) (2,0.039) (3,0.274) (4,1.573) (5,7.263) (6,16.4) (7,51.82) (8,82.65) (9,188.4) (10,283) (11,488.8) (12,465.5) (13,508.7) (14,349) (15,265.1) (16,137) (17,53.04) (18,12.74) (19,1.705)};\draw[densely dotted] (axis cs:16,0) -- (axis cs:16,559.5);\end{axis}\begin{axis}[width=72mm,height=42mm,axis y line*=right,axis x line=none,ymode=log,xmin=-1,xmax=20,ylabel={\scriptsize eindeutig (log)},tick label style={font=\tiny},label style={font=\scriptsize},ylabel shift=-3pt]\addplot[mark=*,mark size=0.7pt,draw=black,thin] coordinates {(0,1) (1,4) (2,39) (3,229) (4,661) (5,1677) (6,1108) (7,1707) (8,1544) (9,2600) (10,3363) (11,2531) (12,638) (13,310) (14,135) (15,24) (16,9)};\end{axis}\end{tikzpicture}\end{minipage}
\vspace{3mm}

\noindent\begin{minipage}[t]{60mm}\vspace{0pt}\sbox\hmboard{\chessboard[setfen=8/8/8/8/8/7k/5r1N/K7 b - - 0 1]}\usebox\hmboard\\[1mm]
\makebox[\wd\hmboard]{\sffamily\hspace{3mm}\textbf{h\#8}\hfill(2+2)}\end{minipage}\hfill
\begin{minipage}[t]{52mm}\vspace{0pt}
\textcolor{ink}{\sffamily\textsc{Lösung}}\par\smallskip\small
\begin{tabular}{@{}r@{\enspace}l@{\enspace}l@{}}
\textbf{1.} & \symking{}g2 & \symknight{}f3\\
\textbf{2.} & \symking{}f1 & \symking{}b1\\
\textbf{3.} & \symking{}e2 & \symking{}b2\\
\textbf{4.} & \symking{}d3+ & \symking{}c1\\
\textbf{5.} & \symking{}c3 & \symking{}d1\\
\textbf{6.} & \symking{}b2 & \symknight{}d4\\
\textbf{7.} & \symking{}a1 & \symking{}c1\\
\textbf{8.} & \symrook{}a2 & \symknight{}b3\#\\
\end{tabular}\par\vspace{3mm}
\qrcode[height=12mm]{https://helpman.komtera.lt/?fen=8/8/8/8/8/7k/5r1N/K7&moves=8}\\{\sffamily\tiny\color{ink}scannen zum Lösen}
\end{minipage}
\vspace{2mm}

{\raggedright\noindent\textcolor{ink}{\sffamily\textsc{Themen}}\enspace \themetag{reines Matt} \themetag{Mustermatt} \themetag{Idealmatt} \themetag{Switchback} \themetag{Selbstblock} \themetag{ohne Schlag}\par}

\begin{tcolorbox}[enhanced,colback=white,colframe=ink!35,arc=1.5pt,boxrule=0.4pt,left=3mm,right=3mm,top=1.5mm,bottom=1.5mm]
\textcolor{ink}{\sffamily\textsc{veröffentlicht}}\par\smallskip\small \textbf{Elkies, Noam D.}, h\#8\quad{\color{ink}\scriptsize verwandte Stellung}\\{\scriptsize Variantim, No. 546, 11/1995}\par\smallskip\textbf{Mertes, Helmut}, h\#8\quad{\color{ink}\scriptsize diese Stellung}\\{\scriptsize feenschach, 01/1977; Ideal-Mate Encyclopedia Vol.1, No. 400, 1999}
\end{tcolorbox}
\clearpage
\phantomsection\label{hm:KNvkn}
\begin{tcolorbox}[enhanced,colback=ink,colframe=ink,arc=0pt,boxrule=0pt,left=4mm,right=4mm,top=2.5mm,bottom=2.5mm]
\color{white}\sffamily{\LARGE\bfseries KNvkn}\hfill{\Large h\#7.5}\\[0.5mm]
{\small 4 Steine\quad Weiß am Zug}\hfill{\small Nr.~10\qquad\href{https://helpman.komtera.lt/?fen=5n2/8/8/8/8/8/3k2N1/K7&moves=7.5}{\color{white}\bfseries lösen\,\ensuremath{\nearrow}}}
\end{tcolorbox}
\vspace{2mm}
\noindent\begin{minipage}[t]{36mm}\vspace{0pt}\sffamily\small\raggedright
\hmfigure{längste eindeutige}{{\Large h\#7.5}}\hmfigure{eindeutig bei dieser Länge}{{\Large 14} Stellungen}\hmfigure{längstes Matt}{{\Large h\#10}}\hmfigure{in der Klasse}{3,78 Millionen Stellungen,\\ 3,13 Millionen mit Hilfsmatt}
\end{minipage}\hfill
\begin{minipage}[t]{80mm}\vspace{0pt}\begin{tikzpicture}\begin{axis}[width=72mm,height=42mm,ybar,bar width=1.6pt,xmin=-1,xmax=21,ymin=0,ymax=636.4,axis y line*=left,axis x line*=bottom,xlabel={\scriptsize Halbzüge bis Matt},ylabel={\scriptsize Stellungen (Tsd.)},tick label style={font=\tiny},label style={font=\scriptsize},enlarge x limits=false,ylabel shift=-3pt,xlabel shift=-3pt]\addplot[fill=ink!55,draw=none] coordinates {(0,0.001) (1,0.005) (2,0.009) (3,0.073) (4,0.436) (5,2.188) (6,6.909) (7,20.12) (8,50.27) (9,100.6) (10,192.7) (11,309) (12,469.6) (13,578.5) (14,564.8) (15,455) (16,250.6) (17,98.22) (18,30.97) (19,3.546) (20,0.966)};\draw[densely dotted] (axis cs:15,0) -- (axis cs:15,636.4);\end{axis}\begin{axis}[width=72mm,height=42mm,axis y line*=right,axis x line=none,ymode=log,xmin=-1,xmax=21,ylabel={\scriptsize eindeutig (log)},tick label style={font=\tiny},label style={font=\scriptsize},ylabel shift=-3pt]\addplot[mark=*,mark size=0.7pt,draw=black,thin] coordinates {(0,1) (1,5) (2,9) (3,59) (4,304) (5,583) (6,733) (7,521) (8,619) (9,545) (10,765) (11,519) (12,392) (13,164) (14,40) (15,14)};\end{axis}\end{tikzpicture}\end{minipage}
\vspace{3mm}

\noindent\begin{minipage}[t]{60mm}\vspace{0pt}\sbox\hmboard{\chessboard[setfen=5n2/8/8/8/8/8/3k2N1/K7 w - - 0 1]}\usebox\hmboard\\[1mm]
\makebox[\wd\hmboard]{\sffamily\hspace{3mm}\textbf{h\#7.5}\hfill(2+2)}\end{minipage}\hfill
\begin{minipage}[t]{52mm}\vspace{0pt}
\textcolor{ink}{\sffamily\textsc{Lösung}}\par\smallskip\small
\begin{tabular}{@{}r@{\enspace}l@{\enspace}l@{}}
\textbf{1.} & \ldots & \symknight{}h4\\
\textbf{2.} & \symking{}e3 & \symking{}b2\\
\textbf{3.} & \symking{}f4 & \symking{}c3\\
\textbf{4.} & \symking{}g5 & \symking{}d4\\
\textbf{5.} & \symking{}h6 & \symking{}e5\\
\textbf{6.} & \symking{}h7 & \symking{}f6\\
\textbf{7.} & \symking{}h8 & \symking{}f7\\
\textbf{8.} & \symknight{}h7 & \symknight{}g6\#\\
\end{tabular}\par\vspace{3mm}
\qrcode[height=12mm]{https://helpman.komtera.lt/?fen=5n2/8/8/8/8/8/3k2N1/K7&moves=7.5}\\{\sffamily\tiny\color{ink}scannen zum Lösen}
\end{minipage}
\vspace{2mm}

{\raggedright\noindent\textcolor{ink}{\sffamily\textsc{Themen}}\enspace \themetag{reines Matt} \themetag{Mustermatt} \themetag{Idealmatt} \themetag{Selbstblock} \themetag{ohne Schlag} \themetag{ohne Schach}\par}

\begin{tcolorbox}[enhanced,colback=white,colframe=ink!35,arc=1.5pt,boxrule=0.4pt,left=3mm,right=3mm,top=1.5mm,bottom=1.5mm]
\textcolor{ink}{\sffamily\textsc{veröffentlicht}}\par\smallskip\small \textbf{Kotesovec, Vaclav}, h\#8 0.1...\quad{\color{ink}\scriptsize verwandte Stellung}\\{\scriptsize Sachova skladba, No. 1065, 06/1987}\par\smallskip\textbf{?, ?}, h\#8 0.1...\quad{\color{ink}\scriptsize verwandte Stellung}\\{\scriptsize ?, 03/1975}
\end{tcolorbox}
\clearpage
\phantomsection\label{hm:KPvkb}
\begin{tcolorbox}[enhanced,colback=ink,colframe=ink,arc=0pt,boxrule=0pt,left=4mm,right=4mm,top=2.5mm,bottom=2.5mm]
\color{white}\sffamily{\LARGE\bfseries KPvkb}\hfill{\Large h\#7.5}\\[0.5mm]
{\small 4 Steine\quad Weiß am Zug}\hfill{\small Nr.~11\qquad\href{https://helpman.komtera.lt/?fen=8/8/1K6/8/2b5/2P5/7k/8&moves=7.5}{\color{white}\bfseries lösen\,\ensuremath{\nearrow}}}
\end{tcolorbox}
\vspace{2mm}
\noindent\begin{minipage}[t]{36mm}\vspace{0pt}\sffamily\small\raggedright
\hmfigure{längste eindeutige}{{\Large h\#7.5}}\hmfigure{eindeutig bei dieser Länge}{{\Large 13} Stellungen}\hmfigure{längstes Matt}{{\Large h\#8.5}}\hmfigure{in der Klasse}{11,1 Millionen Stellungen,\\ 9,43 Millionen mit Hilfsmatt}
\end{minipage}\hfill
\begin{minipage}[t]{80mm}\vspace{0pt}\begin{tikzpicture}\begin{axis}[width=72mm,height=42mm,ybar,bar width=1.6pt,xmin=-1,xmax=18,ymin=0,ymax=1.536,axis y line*=left,axis x line*=bottom,xlabel={\scriptsize Halbzüge bis Matt},ylabel={\scriptsize Stellungen (Mio.)},tick label style={font=\tiny},label style={font=\scriptsize},enlarge x limits=false,ylabel shift=-3pt,xlabel shift=-3pt]\addplot[fill=ink!55,draw=none] coordinates {(0,4e-06) (1,0.001283) (2,0.003574) (3,0.02802) (4,0.064) (5,0.2559) (6,0.3963) (7,0.7995) (8,0.8766) (9,1.218) (10,1.082) (11,1.397) (12,1.24) (13,1.147) (14,0.7683) (15,0.133) (16,0.0159) (17,0.002214)};\draw[densely dotted] (axis cs:15,0) -- (axis cs:15,1.536);\end{axis}\begin{axis}[width=72mm,height=42mm,axis y line*=right,axis x line=none,ymode=log,xmin=-1,xmax=18,ylabel={\scriptsize eindeutig (log)},tick label style={font=\tiny},label style={font=\scriptsize},ylabel shift=-3pt]\addplot[mark=*,mark size=0.7pt,draw=black,thin] coordinates {(0,4) (1,260) (2,512) (3,6031) (4,10252) (5,22888) (6,13735) (7,12916) (8,6370) (9,7550) (10,5121) (11,6391) (12,3807) (13,1453) (14,84) (15,13)};\end{axis}\end{tikzpicture}\end{minipage}
\vspace{3mm}

\noindent\begin{minipage}[t]{60mm}\vspace{0pt}\sbox\hmboard{\chessboard[setfen=8/8/1K6/8/2b5/2P5/7k/8 w - - 0 1]}\usebox\hmboard\\[1mm]
\makebox[\wd\hmboard]{\sffamily\hspace{3mm}\textbf{h\#7.5}\hfill(2+2)}\end{minipage}\hfill
\begin{minipage}[t]{52mm}\vspace{0pt}
\textcolor{ink}{\sffamily\textsc{Lösung}}\par\smallskip\small
\begin{tabular}{@{}r@{\enspace}l@{\enspace}l@{}}
\textbf{1.} & \ldots & \symking{}a5\\
\textbf{2.} & \symbishop{}d5 & c4\\
\textbf{3.} & \symking{}g3 & cxd5\\
\textbf{4.} & \symking{}f4 & d6\\
\textbf{5.} & \symking{}e5 & d7\\
\textbf{6.} & \symking{}d6 & \symking{}b5\\
\textbf{7.} & \symking{}c7 & \symking{}a6\\
\textbf{8.} & \symking{}b8 & d8=\symqueen{}\#\\
\end{tabular}\par\vspace{3mm}
\qrcode[height=12mm]{https://helpman.komtera.lt/?fen=8/8/1K6/8/2b5/2P5/7k/8&moves=7.5}\\{\sffamily\tiny\color{ink}scannen zum Lösen}
\end{minipage}
\vspace{2mm}

{\raggedright\noindent\textcolor{ink}{\sffamily\textsc{Themen}}\enspace \themetag{Satzspiel} \themetag{reines Matt} \themetag{Mustermatt} \themetag{Idealmatt} \themetag{Spiegelmatt} \themetag{Umwandlung} \themetag{ohne Schach} \themetag{Umnov}\par}

\clearpage
\phantomsection\label{hm:KPvkn}
\begin{tcolorbox}[enhanced,colback=ink,colframe=ink,arc=0pt,boxrule=0pt,left=4mm,right=4mm,top=2.5mm,bottom=2.5mm]
\color{white}\sffamily{\LARGE\bfseries KPvkn}\hfill{\Large h\#7.5}\\[0.5mm]
{\small 4 Steine\quad Weiß am Zug}\hfill{\small Nr.~12\qquad\href{https://helpman.komtera.lt/?fen=K7/8/8/n7/2P5/8/8/7k&moves=7.5}{\color{white}\bfseries lösen\,\ensuremath{\nearrow}}}
\end{tcolorbox}
\vspace{2mm}
\noindent\begin{minipage}[t]{36mm}\vspace{0pt}\sffamily\small\raggedright
\hmfigure{längste eindeutige}{{\Large h\#7.5}}\hmfigure{eindeutig bei dieser Länge}{{\Large 8} Stellungen}\hmfigure{längstes Matt}{{\Large h\#8.5}}\hmfigure{in der Klasse}{11,1 Millionen Stellungen,\\ 9,69 Millionen mit Hilfsmatt}
\end{minipage}\hfill
\begin{minipage}[t]{80mm}\vspace{0pt}\begin{tikzpicture}\begin{axis}[width=72mm,height=42mm,ybar,bar width=1.6pt,xmin=-1,xmax=18,ymin=0,ymax=1.496,axis y line*=left,axis x line*=bottom,xlabel={\scriptsize Halbzüge bis Matt},ylabel={\scriptsize Stellungen (Mio.)},tick label style={font=\tiny},label style={font=\scriptsize},enlarge x limits=false,ylabel shift=-3pt,xlabel shift=-3pt]\addplot[fill=ink!55,draw=none] coordinates {(0,2e-06) (1,0.001766) (2,0.004197) (3,0.03012) (4,0.06625) (5,0.2466) (6,0.3809) (7,0.7671) (8,0.8894) (9,1.187) (10,1.118) (11,1.36) (12,1.266) (13,1.159) (14,0.8928) (15,0.2274) (16,0.08652) (17,0.002358)};\draw[densely dotted] (axis cs:15,0) -- (axis cs:15,1.496);\end{axis}\begin{axis}[width=72mm,height=42mm,axis y line*=right,axis x line=none,ymode=log,xmin=-1,xmax=18,ylabel={\scriptsize eindeutig (log)},tick label style={font=\tiny},label style={font=\scriptsize},ylabel shift=-3pt]\addplot[mark=*,mark size=0.7pt,draw=black,thin] coordinates {(0,2) (1,287) (2,564) (3,5835) (4,9281) (5,20725) (6,14175) (7,10762) (8,6557) (9,6958) (10,4950) (11,5977) (12,3648) (13,1316) (14,111) (15,8)};\end{axis}\end{tikzpicture}\end{minipage}
\vspace{3mm}

\noindent\begin{minipage}[t]{60mm}\vspace{0pt}\sbox\hmboard{\chessboard[setfen=K7/8/8/n7/2P5/8/8/7k w - - 0 1]}\usebox\hmboard\\[1mm]
\makebox[\wd\hmboard]{\sffamily\hspace{3mm}\textbf{h\#7.5}\hfill(2+2)}\end{minipage}\hfill
\begin{minipage}[t]{52mm}\vspace{0pt}
\textcolor{ink}{\sffamily\textsc{Lösung}}\par\smallskip\small
\begin{tabular}{@{}r@{\enspace}l@{\enspace}l@{}}
\textbf{1.} & \ldots & c5\\
\textbf{2.} & \symking{}g2 & c6\\
\textbf{3.} & \symking{}f3 & c7\\
\textbf{4.} & \symking{}e4 & c8=\symqueen{}\\
\textbf{5.} & \symking{}d5 & \symqueen{}f5+\\
\textbf{6.} & \symking{}c6 & \symking{}b8\\
\textbf{7.} & \symking{}b6 & \symking{}c8\\
\textbf{8.} & \symking{}a7 & \symqueen{}xa5\#\\
\end{tabular}\par\vspace{3mm}
\qrcode[height=12mm]{https://helpman.komtera.lt/?fen=K7/8/8/n7/2P5/8/8/7k&moves=7.5}\\{\sffamily\tiny\color{ink}scannen zum Lösen}
\end{minipage}
\vspace{2mm}

{\raggedright\noindent\textcolor{ink}{\sffamily\textsc{Themen}}\enspace \themetag{reines Matt} \themetag{Mustermatt} \themetag{Idealmatt} \themetag{Spiegelmatt} \themetag{Umwandlung} \themetag{ein Stein zieht}\par}

\clearpage
\phantomsection\label{hm:KPvkq}
\begin{tcolorbox}[enhanced,colback=ink,colframe=ink,arc=0pt,boxrule=0pt,left=4mm,right=4mm,top=2.5mm,bottom=2.5mm]
\color{white}\sffamily{\LARGE\bfseries KPvkq}\hfill{\Large h\#7.5}\\[0.5mm]
{\small 4 Steine\quad Weiß am Zug}\hfill{\small Nr.~13\qquad\href{https://helpman.komtera.lt/?fen=8/K7/8/8/8/8/3P4/2q4k&moves=7.5}{\color{white}\bfseries lösen\,\ensuremath{\nearrow}}}
\end{tcolorbox}
\vspace{2mm}
\noindent\begin{minipage}[t]{36mm}\vspace{0pt}\sffamily\small\raggedright
\hmfigure{längste eindeutige}{{\Large h\#7.5}}\hmfigure{eindeutig bei dieser Länge}{{\Large 32} Stellungen}\hmfigure{längstes Matt}{{\Large h\#9}}\hmfigure{in der Klasse}{11,1 Millionen Stellungen,\\ 8,34 Millionen mit Hilfsmatt}
\end{minipage}\hfill
\begin{minipage}[t]{80mm}\vspace{0pt}\begin{tikzpicture}\begin{axis}[width=72mm,height=42mm,ybar,bar width=1.6pt,xmin=-1,xmax=19,ymin=0,ymax=1.519,axis y line*=left,axis x line*=bottom,xlabel={\scriptsize Halbzüge bis Matt},ylabel={\scriptsize Stellungen (Mio.)},tick label style={font=\tiny},label style={font=\scriptsize},enlarge x limits=false,ylabel shift=-3pt,xlabel shift=-3pt]\addplot[fill=ink!55,draw=none] coordinates {(1,0.000543) (2,0.002283) (3,0.02037) (4,0.04664) (5,0.2155) (6,0.2931) (7,0.7483) (8,0.6602) (9,1.2) (10,0.8259) (11,1.381) (12,0.9621) (13,1.206) (14,0.566) (15,0.1948) (16,0.01481) (17,0.004737) (18,8e-06)};\draw[densely dotted] (axis cs:15,0) -- (axis cs:15,1.519);\end{axis}\begin{axis}[width=72mm,height=42mm,axis y line*=right,axis x line=none,ymode=log,xmin=-1,xmax=19,ylabel={\scriptsize eindeutig (log)},tick label style={font=\tiny},label style={font=\scriptsize},ylabel shift=-3pt]\addplot[mark=*,mark size=0.7pt,draw=black,thin] coordinates {(1,134) (2,321) (3,4529) (4,6264) (5,19251) (6,7733) (7,12559) (8,4177) (9,7285) (10,3249) (11,6055) (12,2273) (13,1932) (14,69) (15,32)};\end{axis}\end{tikzpicture}\end{minipage}
\vspace{3mm}

\noindent\begin{minipage}[t]{60mm}\vspace{0pt}\sbox\hmboard{\chessboard[setfen=8/K7/8/8/8/8/3P4/2q4k w - - 0 1]}\usebox\hmboard\\[1mm]
\makebox[\wd\hmboard]{\sffamily\hspace{3mm}\textbf{h\#7.5}\hfill(2+2)}\end{minipage}\hfill
\begin{minipage}[t]{52mm}\vspace{0pt}
\textcolor{ink}{\sffamily\textsc{Lösung}}\par\smallskip\small
\begin{tabular}{@{}r@{\enspace}l@{\enspace}l@{}}
\textbf{1.} & \ldots & d4\\
\textbf{2.} & \symking{}g2 & d5\\
\textbf{3.} & \symking{}f3 & d6\\
\textbf{4.} & \symking{}e4 & d7\\
\textbf{5.} & \symking{}d5 & d8=\symqueen{}+\\
\textbf{6.} & \symking{}c6 & \symqueen{}d5+\\
\textbf{7.} & \symking{}c7 & \symking{}a6\\
\textbf{8.} & \symking{}b8 & \symqueen{}b7\#\\
\end{tabular}\par\vspace{3mm}
\qrcode[height=12mm]{https://helpman.komtera.lt/?fen=8/K7/8/8/8/8/3P4/2q4k&moves=7.5}\\{\sffamily\tiny\color{ink}scannen zum Lösen}
\end{minipage}
\vspace{2mm}

{\raggedright\noindent\textcolor{ink}{\sffamily\textsc{Themen}}\enspace \themetag{Umwandlung} \themetag{Excelsior} \themetag{Rundlauf} \themetag{ein Stein zieht} \themetag{ohne Schlag} \themetag{Umnov}\par}

\clearpage
\phantomsection\label{hm:KPvkr}
\begin{tcolorbox}[enhanced,colback=ink,colframe=ink,arc=0pt,boxrule=0pt,left=4mm,right=4mm,top=2.5mm,bottom=2.5mm]
\color{white}\sffamily{\LARGE\bfseries KPvkr}\hfill{\Large h\#7.5}\\[0.5mm]
{\small 4 Steine\quad Weiß am Zug}\hfill{\small Nr.~14\qquad\href{https://helpman.komtera.lt/?fen=8/7k/8/8/8/7r/1K5P/8&moves=7.5}{\color{white}\bfseries lösen\,\ensuremath{\nearrow}}}
\end{tcolorbox}
\vspace{2mm}
\noindent\begin{minipage}[t]{36mm}\vspace{0pt}\sffamily\small\raggedright
\hmfigure{längste eindeutige}{{\Large h\#7.5}}\hmfigure{eindeutig bei dieser Länge}{{\Large 1} Stellung}\hmfigure{längstes Matt}{{\Large h\#8.5}}\hmfigure{in der Klasse}{11,1 Millionen Stellungen,\\ 9,03 Millionen mit Hilfsmatt}
\end{minipage}\hfill
\begin{minipage}[t]{80mm}\vspace{0pt}\begin{tikzpicture}\begin{axis}[width=72mm,height=42mm,ybar,bar width=1.6pt,xmin=-1,xmax=18,ymin=0,ymax=1.479,axis y line*=left,axis x line*=bottom,xlabel={\scriptsize Halbzüge bis Matt},ylabel={\scriptsize Stellungen (Mio.)},tick label style={font=\tiny},label style={font=\scriptsize},enlarge x limits=false,ylabel shift=-3pt,xlabel shift=-3pt]\addplot[fill=ink!55,draw=none] coordinates {(1,0.001025) (2,0.003763) (3,0.02941) (4,0.07967) (5,0.3111) (6,0.5241) (7,0.9961) (8,0.9584) (9,1.345) (10,0.9543) (11,1.34) (12,1.076) (13,0.9003) (14,0.4515) (15,0.05363) (16,0.002327) (17,0.000154)};\draw[densely dotted] (axis cs:15,0) -- (axis cs:15,1.479);\end{axis}\begin{axis}[width=72mm,height=42mm,axis y line*=right,axis x line=none,ymode=log,xmin=-1,xmax=18,ylabel={\scriptsize eindeutig (log)},tick label style={font=\tiny},label style={font=\scriptsize},ylabel shift=-3pt]\addplot[mark=*,mark size=0.7pt,draw=black,thin] coordinates {(1,259) (2,937) (3,7812) (4,16296) (5,28734) (6,14045) (7,15266) (8,6530) (9,8362) (10,4063) (11,5800) (12,2421) (13,699) (14,4) (15,1)};\end{axis}\end{tikzpicture}\end{minipage}
\vspace{3mm}

\noindent\begin{minipage}[t]{60mm}\vspace{0pt}\sbox\hmboard{\chessboard[setfen=8/7k/8/8/8/7r/1K5P/8 w - - 0 1]}\usebox\hmboard\\[1mm]
\makebox[\wd\hmboard]{\sffamily\hspace{3mm}\textbf{h\#7.5}\hfill(2+2)}\end{minipage}\hfill
\begin{minipage}[t]{52mm}\vspace{0pt}
\textcolor{ink}{\sffamily\textsc{Lösung}}\par\smallskip\small
\begin{tabular}{@{}r@{\enspace}l@{\enspace}l@{}}
\textbf{1.} & \ldots & \symking{}a2\\
\textbf{2.} & \symrook{}a3+ & \symking{}xa3\\
\textbf{3.} & \symking{}g6 & h4\\
\textbf{4.} & \symking{}f5 & h5\\
\textbf{5.} & \symking{}e4 & h6\\
\textbf{6.} & \symking{}d3 & h7\\
\textbf{7.} & \symking{}c2 & h8=\symqueen{}\\
\textbf{8.} & \symking{}b1 & \symqueen{}b2\#\\
\end{tabular}\par\vspace{3mm}
\qrcode[height=12mm]{https://helpman.komtera.lt/?fen=8/7k/8/8/8/7r/1K5P/8&moves=7.5}\\{\sffamily\tiny\color{ink}scannen zum Lösen}
\end{minipage}
\vspace{2mm}

{\raggedright\noindent\textcolor{ink}{\sffamily\textsc{Themen}}\enspace \themetag{Satzspiel} \themetag{Umwandlung} \themetag{Excelsior}\par}

\begin{tcolorbox}[enhanced,colback=white,colframe=ink!35,arc=1.5pt,boxrule=0.4pt,left=3mm,right=3mm,top=1.5mm,bottom=1.5mm]
\textcolor{ink}{\sffamily\textsc{veröffentlicht}}\par\smallskip\small \textbf{Maslar, Zdravko}, h\#8 0.1...\quad{\color{ink}\scriptsize diese Stellung}\\{\scriptsize ?,}
\end{tcolorbox}
\clearpage
\phantomsection\label{hm:KBPvk}
\begin{tcolorbox}[enhanced,colback=ink,colframe=ink,arc=0pt,boxrule=0pt,left=4mm,right=4mm,top=2.5mm,bottom=2.5mm]
\color{white}\sffamily{\LARGE\bfseries KBPvk}\hfill{\Large h\#7}\\[0.5mm]
{\small 4 Steine\quad Schwarz am Zug}\hfill{\small Nr.~15\qquad\href{https://helpman.komtera.lt/?fen=7k/8/8/8/B7/8/P7/3K4&moves=7}{\color{white}\bfseries lösen\,\ensuremath{\nearrow}}}
\end{tcolorbox}
\vspace{2mm}
\noindent\begin{minipage}[t]{36mm}\vspace{0pt}\sffamily\small\raggedright
\hmfigure{längste eindeutige}{{\Large h\#7}}\hmfigure{eindeutig bei dieser Länge}{{\Large 12} Stellungen}\hmfigure{längstes Matt}{{\Large h\#8}}\hmfigure{in der Klasse}{11,1 Millionen Stellungen,\\ 9,44 Millionen mit Hilfsmatt}
\end{minipage}\hfill
\begin{minipage}[t]{80mm}\vspace{0pt}\begin{tikzpicture}\begin{axis}[width=72mm,height=42mm,ybar,bar width=1.6pt,xmin=-1,xmax=17,ymin=0,ymax=1.494,axis y line*=left,axis x line*=bottom,xlabel={\scriptsize Halbzüge bis Matt},ylabel={\scriptsize Stellungen (Mio.)},tick label style={font=\tiny},label style={font=\scriptsize},enlarge x limits=false,ylabel shift=-3pt,xlabel shift=-3pt]\addplot[fill=ink!55,draw=none] coordinates {(0,0.000117) (1,0.003118) (2,0.006206) (3,0.05643) (4,0.1671) (5,0.6887) (6,0.9375) (7,0.8846) (8,0.9984) (9,0.888) (10,1.057) (11,1.123) (12,1.359) (13,0.6596) (14,0.5877) (15,0.01359) (16,0.007124)};\draw[densely dotted] (axis cs:14,0) -- (axis cs:14,1.494);\end{axis}\begin{axis}[width=72mm,height=42mm,axis y line*=right,axis x line=none,ymode=log,xmin=-1,xmax=17,ylabel={\scriptsize eindeutig (log)},tick label style={font=\tiny},label style={font=\scriptsize},ylabel shift=-3pt]\addplot[mark=*,mark size=0.7pt,draw=black,thin] coordinates {(0,117) (1,1071) (2,2736) (3,30260) (4,54754) (5,69846) (6,36972) (7,27207) (8,18560) (9,14540) (10,11833) (11,10480) (12,6920) (13,81) (14,12)};\end{axis}\end{tikzpicture}\end{minipage}
\vspace{3mm}

\noindent\begin{minipage}[t]{60mm}\vspace{0pt}\sbox\hmboard{\chessboard[setfen=7k/8/8/8/B7/8/P7/3K4 b - - 0 1]}\usebox\hmboard\\[1mm]
\makebox[\wd\hmboard]{\sffamily\hspace{3mm}\textbf{h\#7}\hfill(3+1)}\end{minipage}\hfill
\begin{minipage}[t]{52mm}\vspace{0pt}
\textcolor{ink}{\sffamily\textsc{Lösung}}\par\smallskip\small
\begin{tabular}{@{}r@{\enspace}l@{\enspace}l@{}}
\textbf{1.} & \symking{}g7 & \symbishop{}b3\\
\textbf{2.} & \symking{}f6 & a4\\
\textbf{3.} & \symking{}e5 & a5\\
\textbf{4.} & \symking{}d4 & a6\\
\textbf{5.} & \symking{}c3 & a7\\
\textbf{6.} & \symking{}b2 & a8=\symqueen{}\\
\textbf{7.} & \symking{}b1 & \symqueen{}a2\#\\
\end{tabular}\par\vspace{3mm}
\qrcode[height=12mm]{https://helpman.komtera.lt/?fen=7k/8/8/8/B7/8/P7/3K4&moves=7}\\{\sffamily\tiny\color{ink}scannen zum Lösen}
\end{minipage}
\vspace{2mm}

{\raggedright\noindent\textcolor{ink}{\sffamily\textsc{Themen}}\enspace \themetag{Umwandlung} \themetag{Excelsior} \themetag{Rundlauf} \themetag{ein Stein zieht} \themetag{ohne Schlag} \themetag{ohne Schach}\par}

\begin{tcolorbox}[enhanced,colback=white,colframe=ink!35,arc=1.5pt,boxrule=0.4pt,left=3mm,right=3mm,top=1.5mm,bottom=1.5mm]
\textcolor{ink}{\sffamily\textsc{veröffentlicht}}\par\smallskip\small \textbf{Elkies, Noam D.}, h\#7\quad{\color{ink}\scriptsize diese Stellung}\\{\scriptsize The Problemist, No. H2794, 07/2004}
\end{tcolorbox}
\clearpage
\phantomsection\label{hm:KNPvk}
\begin{tcolorbox}[enhanced,colback=ink,colframe=ink,arc=0pt,boxrule=0pt,left=4mm,right=4mm,top=2.5mm,bottom=2.5mm]
\color{white}\sffamily{\LARGE\bfseries KNPvk}\hfill{\Large h\#7}\\[0.5mm]
{\small 4 Steine\quad Schwarz am Zug}\hfill{\small Nr.~16\qquad\href{https://helpman.komtera.lt/?fen=8/8/8/8/8/N7/P7/3K3k&moves=7}{\color{white}\bfseries lösen\,\ensuremath{\nearrow}}}
\end{tcolorbox}
\vspace{2mm}
\noindent\begin{minipage}[t]{36mm}\vspace{0pt}\sffamily\small\raggedright
\hmfigure{längste eindeutige}{{\Large h\#7}}\hmfigure{eindeutig bei dieser Länge}{{\Large 49} Stellungen}\hmfigure{längstes Matt}{{\Large h\#8.5}}\hmfigure{in der Klasse}{11,1 Millionen Stellungen,\\ 9,7 Millionen mit Hilfsmatt}
\end{minipage}\hfill
\begin{minipage}[t]{80mm}\vspace{0pt}\begin{tikzpicture}\begin{axis}[width=72mm,height=42mm,ybar,bar width=1.6pt,xmin=-1,xmax=18,ymin=0,ymax=1.45,axis y line*=left,axis x line*=bottom,xlabel={\scriptsize Halbzüge bis Matt},ylabel={\scriptsize Stellungen (Mio.)},tick label style={font=\tiny},label style={font=\scriptsize},enlarge x limits=false,ylabel shift=-3pt,xlabel shift=-3pt]\addplot[fill=ink!55,draw=none] coordinates {(0,2.6e-05) (1,0.002591) (2,0.004813) (3,0.04594) (4,0.1254) (5,0.5173) (6,0.7779) (7,0.9545) (8,1.03) (9,0.9719) (10,1.044) (11,1.156) (12,1.319) (13,0.8793) (14,0.8034) (15,0.04696) (16,0.01796) (17,1.8e-05)};\draw[densely dotted] (axis cs:14,0) -- (axis cs:14,1.45);\end{axis}\begin{axis}[width=72mm,height=42mm,axis y line*=right,axis x line=none,ymode=log,xmin=-1,xmax=18,ylabel={\scriptsize eindeutig (log)},tick label style={font=\tiny},label style={font=\scriptsize},ylabel shift=-3pt]\addplot[mark=*,mark size=0.7pt,draw=black,thin] coordinates {(0,26) (1,558) (2,1475) (3,22342) (4,38108) (5,50476) (6,26775) (7,19630) (8,12108) (9,9541) (10,6416) (11,8571) (12,5032) (13,167) (14,49)};\end{axis}\end{tikzpicture}\end{minipage}
\vspace{3mm}

\noindent\begin{minipage}[t]{60mm}\vspace{0pt}\sbox\hmboard{\chessboard[setfen=8/8/8/8/8/N7/P7/3K3k b - - 0 1]}\usebox\hmboard\\[1mm]
\makebox[\wd\hmboard]{\sffamily\hspace{3mm}\textbf{h\#7}\hfill(3+1)}\end{minipage}\hfill
\begin{minipage}[t]{52mm}\vspace{0pt}
\textcolor{ink}{\sffamily\textsc{Lösung}}\par\smallskip\small
\begin{tabular}{@{}r@{\enspace}l@{\enspace}l@{}}
\textbf{1.} & \symking{}g2 & \symknight{}c2\\
\textbf{2.} & \symking{}f3 & a4\\
\textbf{3.} & \symking{}e4 & a5\\
\textbf{4.} & \symking{}d3 & a6\\
\textbf{5.} & \symking{}c3 & a7\\
\textbf{6.} & \symking{}b2 & a8=\symqueen{}\\
\textbf{7.} & \symking{}b1 & \symqueen{}a1\#\\
\end{tabular}\par\vspace{3mm}
\qrcode[height=12mm]{https://helpman.komtera.lt/?fen=8/8/8/8/8/N7/P7/3K3k&moves=7}\\{\sffamily\tiny\color{ink}scannen zum Lösen}
\end{minipage}
\vspace{2mm}

{\raggedright\noindent\textcolor{ink}{\sffamily\textsc{Themen}}\enspace \themetag{Umwandlung} \themetag{Excelsior} \themetag{ein Stein zieht} \themetag{ohne Schlag} \themetag{ohne Schach}\par}

\begin{tcolorbox}[enhanced,colback=white,colframe=ink!35,arc=1.5pt,boxrule=0.4pt,left=3mm,right=3mm,top=1.5mm,bottom=1.5mm]
\textcolor{ink}{\sffamily\textsc{veröffentlicht}}\par\smallskip\small \textbf{Elkies, Noam D.}, h\#7\quad{\color{ink}\scriptsize diese Stellung}\\{\scriptsize The Problemist, No. H2793, 07/2004}
\end{tcolorbox}
\clearpage
\phantomsection\label{hm:KPPvk}
\begin{tcolorbox}[enhanced,colback=ink,colframe=ink,arc=0pt,boxrule=0pt,left=4mm,right=4mm,top=2.5mm,bottom=2.5mm]
\color{white}\sffamily{\LARGE\bfseries KPPvk}\hfill{\Large h\#7}\\[0.5mm]
{\small 4 Steine\quad Schwarz am Zug}\hfill{\small Nr.~17\qquad\href{https://helpman.komtera.lt/?fen=8/7k/8/8/8/6P1/K5P1/8&moves=7}{\color{white}\bfseries lösen\,\ensuremath{\nearrow}}}
\end{tcolorbox}
\vspace{2mm}
\noindent\begin{minipage}[t]{36mm}\vspace{0pt}\sffamily\small\raggedright
\hmfigure{längste eindeutige}{{\Large h\#7}}\hmfigure{eindeutig bei dieser Länge}{{\Large 2} Stellungen}\hmfigure{längstes Matt}{{\Large h\#8}}\hmfigure{in der Klasse}{8,32 Millionen Stellungen,\\ 3,72 Millionen mit Hilfsmatt}
\end{minipage}\hfill
\begin{minipage}[t]{80mm}\vspace{0pt}\begin{tikzpicture}\begin{axis}[width=72mm,height=42mm,ybar,bar width=1.6pt,xmin=-1,xmax=17,ymin=0,ymax=638.2,axis y line*=left,axis x line*=bottom,xlabel={\scriptsize Halbzüge bis Matt},ylabel={\scriptsize Stellungen (Tsd.)},tick label style={font=\tiny},label style={font=\scriptsize},enlarge x limits=false,ylabel shift=-3pt,xlabel shift=-3pt]\addplot[fill=ink!55,draw=none] coordinates {(0,0.012) (1,1.863) (2,3.369) (3,27.26) (4,68.17) (5,255.1) (6,397.4) (7,540.8) (8,580.1) (9,512.9) (10,468.6) (11,342.6) (12,303.3) (13,120.4) (14,89.47) (15,5.716) (16,1.375)};\draw[densely dotted] (axis cs:14,0) -- (axis cs:14,638.2);\end{axis}\begin{axis}[width=72mm,height=42mm,axis y line*=right,axis x line=none,ymode=log,xmin=-1,xmax=17,ylabel={\scriptsize eindeutig (log)},tick label style={font=\tiny},label style={font=\scriptsize},ylabel shift=-3pt]\addplot[mark=*,mark size=0.7pt,draw=black,thin] coordinates {(0,12) (1,352) (2,938) (3,10467) (4,17948) (5,22896) (6,13662) (7,7903) (8,4058) (9,2716) (10,1840) (11,1281) (12,692) (13,36) (14,2)};\end{axis}\end{tikzpicture}\end{minipage}
\vspace{3mm}

\noindent\begin{minipage}[t]{60mm}\vspace{0pt}\sbox\hmboard{\chessboard[setfen=8/7k/8/8/8/6P1/K5P1/8 b - - 0 1]}\usebox\hmboard\\[1mm]
\makebox[\wd\hmboard]{\sffamily\hspace{3mm}\textbf{h\#7}\hfill(3+1)}\end{minipage}\hfill
\begin{minipage}[t]{52mm}\vspace{0pt}
\textcolor{ink}{\sffamily\textsc{Lösung}}\par\smallskip\small
\begin{tabular}{@{}r@{\enspace}l@{\enspace}l@{}}
\textbf{1.} & \symking{}g6 & \symking{}a3\\
\textbf{2.} & \symking{}f5 & g4+\\
\textbf{3.} & \symking{}e4 & g5\\
\textbf{4.} & \symking{}d3 & g6\\
\textbf{5.} & \symking{}c2 & g7\\
\textbf{6.} & \symking{}b1 & g8=\symqueen{}\\
\textbf{7.} & \symking{}a1 & \symqueen{}a2\#\\
\end{tabular}\par\vspace{3mm}
\qrcode[height=12mm]{https://helpman.komtera.lt/?fen=8/7k/8/8/8/6P1/K5P1/8&moves=7}\\{\sffamily\tiny\color{ink}scannen zum Lösen}
\end{minipage}
\vspace{2mm}

{\raggedright\noindent\textcolor{ink}{\sffamily\textsc{Themen}}\enspace \themetag{Umwandlung} \themetag{ein Stein zieht} \themetag{ohne Schlag}\par}

\begin{tcolorbox}[enhanced,colback=white,colframe=ink!35,arc=1.5pt,boxrule=0.4pt,left=3mm,right=3mm,top=1.5mm,bottom=1.5mm]
\textcolor{ink}{\sffamily\textsc{veröffentlicht}}\par\smallskip\small \textbf{Salazar Lopez, Francisco}, h\#7\quad{\color{ink}\scriptsize verwandte Stellung}\\{\scriptsize British Chess Magazine, 03/1973; Catalogue h\# (CHM) avec 6 pieces, No. 213, 1995; The Problemist, 05/2005; \ldots}
\end{tcolorbox}
\clearpage
\phantomsection\label{hm:KQPvk}
\begin{tcolorbox}[enhanced,colback=ink,colframe=ink,arc=0pt,boxrule=0pt,left=4mm,right=4mm,top=2.5mm,bottom=2.5mm]
\color{white}\sffamily{\LARGE\bfseries KQPvk}\hfill{\Large h\#7}\\[0.5mm]
{\small 4 Steine\quad Schwarz am Zug}\hfill{\small Nr.~18\qquad\href{https://helpman.komtera.lt/?fen=6Qk/8/8/8/8/8/4P3/2K5&moves=7}{\color{white}\bfseries lösen\,\ensuremath{\nearrow}}}
\end{tcolorbox}
\vspace{2mm}
\noindent\begin{minipage}[t]{36mm}\vspace{0pt}\sffamily\small\raggedright
\hmfigure{längste eindeutige}{{\Large h\#7}}\hmfigure{eindeutig bei dieser Länge}{{\Large 12} Stellungen}\hmfigure{längstes Matt}{{\Large h\#8}}\hmfigure{in der Klasse}{11,1 Millionen Stellungen,\\ 8,38 Millionen mit Hilfsmatt}
\end{minipage}\hfill
\begin{minipage}[t]{80mm}\vspace{0pt}\begin{tikzpicture}\begin{axis}[width=72mm,height=42mm,ybar,bar width=1.6pt,xmin=-1,xmax=17,ymin=0,ymax=2.462,axis y line*=left,axis x line*=bottom,xlabel={\scriptsize Halbzüge bis Matt},ylabel={\scriptsize Stellungen (Mio.)},tick label style={font=\tiny},label style={font=\scriptsize},enlarge x limits=false,ylabel shift=-3pt,xlabel shift=-3pt]\addplot[fill=ink!55,draw=none] coordinates {(0,0.01023) (1,0.07721) (2,0.1821) (3,0.6692) (4,1.684) (5,1.511) (6,2.238) (7,0.8701) (8,0.8592) (9,0.1414) (10,0.094) (11,0.004873) (12,0.01722) (13,1.1e-05) (14,0.01408) (16,0.002299)};\draw[densely dotted] (axis cs:14,0) -- (axis cs:14,2.462);\end{axis}\begin{axis}[width=72mm,height=42mm,axis y line*=right,axis x line=none,ymode=log,xmin=-1,xmax=17,ylabel={\scriptsize eindeutig (log)},tick label style={font=\tiny},label style={font=\scriptsize},ylabel shift=-3pt]\addplot[mark=*,mark size=0.7pt,draw=black,thin] coordinates {(0,10232) (1,47073) (2,111848) (3,153118) (4,219190) (5,70286) (6,48465) (7,6915) (8,3101) (9,83) (10,116) (11,6) (12,46) (14,12)};\end{axis}\end{tikzpicture}\end{minipage}
\vspace{3mm}

\noindent\begin{minipage}[t]{60mm}\vspace{0pt}\sbox\hmboard{\chessboard[setfen=6Qk/8/8/8/8/8/4P3/2K5 b - - 0 1]}\usebox\hmboard\\[1mm]
\makebox[\wd\hmboard]{\sffamily\hspace{3mm}\textbf{h\#7}\hfill(3+1)}\end{minipage}\hfill
\begin{minipage}[t]{52mm}\vspace{0pt}
\textcolor{ink}{\sffamily\textsc{Lösung}}\par\smallskip\small
\begin{tabular}{@{}r@{\enspace}l@{\enspace}l@{}}
\textbf{1.} & \symking{}xg8 & e3\\
\textbf{2.} & \symking{}f7 & e4\\
\textbf{3.} & \symking{}e6 & e5\\
\textbf{4.} & \symking{}d5 & e6\\
\textbf{5.} & \symking{}c4 & e7\\
\textbf{6.} & \symking{}b3 & e8=\symqueen{}\\
\textbf{7.} & \symking{}a2 & \symqueen{}a4\#\\
\end{tabular}\par\vspace{3mm}
\qrcode[height=12mm]{https://helpman.komtera.lt/?fen=6Qk/8/8/8/8/8/4P3/2K5&moves=7}\\{\sffamily\tiny\color{ink}scannen zum Lösen}
\end{minipage}
\vspace{2mm}

{\raggedright\noindent\textcolor{ink}{\sffamily\textsc{Themen}}\enspace \themetag{reines Matt} \themetag{Mustermatt} \themetag{Idealmatt} \themetag{Spiegelmatt} \themetag{Umwandlung} \themetag{Excelsior} \themetag{ein Stein zieht} \themetag{Phönix} \themetag{ohne Schach} \themetag{Umnov}\par\smallskip{\footnotesize\textit{Hinweis:} die Partei am Zug steht im Diagramm im Schach; die Lösung beginnt mit einem Schlag.}\par}

\clearpage
\phantomsection\label{hm:KBBvk}
\begin{tcolorbox}[enhanced,colback=ink,colframe=ink,arc=0pt,boxrule=0pt,left=4mm,right=4mm,top=2.5mm,bottom=2.5mm]
\color{white}\sffamily{\LARGE\bfseries KBBvk}\hfill{\Large h\#6.5}\\[0.5mm]
{\small 4 Steine\quad Weiß am Zug}\hfill{\small Nr.~19\qquad\href{https://helpman.komtera.lt/?fen=7k/8/7B/8/8/8/2B5/K7&moves=6.5}{\color{white}\bfseries lösen\,\ensuremath{\nearrow}}}
\end{tcolorbox}
\vspace{2mm}
\noindent\begin{minipage}[t]{36mm}\vspace{0pt}\sffamily\small\raggedright
\hmfigure{längste eindeutige}{{\Large h\#6.5}}\hmfigure{eindeutig bei dieser Länge}{{\Large 1} Stellung}\hmfigure{längstes Matt}{{\Large h\#8}}\hmfigure{in der Klasse}{3,78 Millionen Stellungen,\\ 745 Tausend mit Hilfsmatt}
\end{minipage}\hfill
\begin{minipage}[t]{80mm}\vspace{0pt}\begin{tikzpicture}\begin{axis}[width=72mm,height=42mm,ybar,bar width=1.6pt,xmin=-1,xmax=17,ymin=0,ymax=191.8,axis y line*=left,axis x line*=bottom,xlabel={\scriptsize Halbzüge bis Matt},ylabel={\scriptsize Stellungen (Tsd.)},tick label style={font=\tiny},label style={font=\scriptsize},enlarge x limits=false,ylabel shift=-3pt,xlabel shift=-3pt]\addplot[fill=ink!55,draw=none] coordinates {(0,0.194) (1,0.789) (2,1.587) (3,9.013) (4,34.63) (5,89.84) (6,174.4) (7,125) (8,158) (9,69.18) (10,52.46) (11,16.93) (12,9.959) (13,2.258) (14,0.845) (15,0.024) (16,0.001)};\draw[densely dotted] (axis cs:13,0) -- (axis cs:13,191.8);\end{axis}\begin{axis}[width=72mm,height=42mm,axis y line*=right,axis x line=none,ymode=log,xmin=-1,xmax=17,ylabel={\scriptsize eindeutig (log)},tick label style={font=\tiny},label style={font=\scriptsize},ylabel shift=-3pt]\addplot[mark=*,mark size=0.7pt,draw=black,thin] coordinates {(0,194) (1,779) (2,1549) (3,7407) (4,13459) (5,8466) (6,8301) (7,1350) (8,597) (9,230) (10,45) (11,23) (12,2) (13,1)};\end{axis}\end{tikzpicture}\end{minipage}
\vspace{3mm}

\noindent\begin{minipage}[t]{60mm}\vspace{0pt}\sbox\hmboard{\chessboard[setfen=7k/8/7B/8/8/8/2B5/K7 w - - 0 1]}\usebox\hmboard\\[1mm]
\makebox[\wd\hmboard]{\sffamily\hspace{3mm}\textbf{h\#6.5}\hfill(3+1)}\end{minipage}\hfill
\begin{minipage}[t]{52mm}\vspace{0pt}
\textcolor{ink}{\sffamily\textsc{Lösung}}\par\smallskip\small
\begin{tabular}{@{}r@{\enspace}l@{\enspace}l@{}}
\textbf{1.} & \ldots & \symking{}b2\\
\textbf{2.} & \symking{}g8 & \symking{}c3\\
\textbf{3.} & \symking{}f7 & \symking{}d3\\
\textbf{4.} & \symking{}g6 & \symking{}e4\\
\textbf{5.} & \symking{}h5 & \symking{}f5\\
\textbf{6.} & \symking{}h4 & \symbishop{}g5+\\
\textbf{7.} & \symking{}h5 & \symbishop{}d1\#\\
\end{tabular}\par\vspace{3mm}
\qrcode[height=12mm]{https://helpman.komtera.lt/?fen=7k/8/7B/8/8/8/2B5/K7&moves=6.5}\\{\sffamily\tiny\color{ink}scannen zum Lösen}
\end{minipage}
\vspace{2mm}

{\raggedright\noindent\textcolor{ink}{\sffamily\textsc{Themen}}\enspace \themetag{Switchback} \themetag{ein Stein zieht} \themetag{ohne Schlag}\par}

\begin{tcolorbox}[enhanced,colback=white,colframe=ink!35,arc=1.5pt,boxrule=0.4pt,left=3mm,right=3mm,top=1.5mm,bottom=1.5mm]
\textcolor{ink}{\sffamily\textsc{veröffentlicht}}\par\smallskip\small \textbf{Maslar, Zdravko}, h\#7 0.1...\quad{\color{ink}\scriptsize verwandte Stellung}\\{\scriptsize feenschach, 01-03/1978; Catalogue h\# (CHM) avec 6 pieces, No. 899, 1995; Problemkiste, No. A1, 03/2012}
\end{tcolorbox}
\clearpage
\phantomsection\label{hm:KBvkb}
\begin{tcolorbox}[enhanced,colback=ink,colframe=ink,arc=0pt,boxrule=0pt,left=4mm,right=4mm,top=2.5mm,bottom=2.5mm]
\color{white}\sffamily{\LARGE\bfseries KBvkb}\hfill{\Large h\#6.5}\\[0.5mm]
{\small 4 Steine\quad Weiß am Zug}\hfill{\small Nr.~20\qquad\href{https://helpman.komtera.lt/?fen=8/8/6b1/8/8/6k1/3B4/K7&moves=6.5}{\color{white}\bfseries lösen\,\ensuremath{\nearrow}}}
\end{tcolorbox}
\vspace{2mm}
\noindent\begin{minipage}[t]{36mm}\vspace{0pt}\sffamily\small\raggedright
\hmfigure{längste eindeutige}{{\Large h\#6.5}}\hmfigure{eindeutig bei dieser Länge}{{\Large 8} Stellungen}\hmfigure{längstes Matt}{{\Large h\#10}}\hmfigure{in der Klasse}{3,78 Millionen Stellungen,\\ 1,49 Millionen mit Hilfsmatt}
\end{minipage}\hfill
\begin{minipage}[t]{80mm}\vspace{0pt}\begin{tikzpicture}\begin{axis}[width=72mm,height=42mm,ybar,bar width=1.6pt,xmin=-1,xmax=21,ymin=0,ymax=290.9,axis y line*=left,axis x line*=bottom,xlabel={\scriptsize Halbzüge bis Matt},ylabel={\scriptsize Stellungen (Tsd.)},tick label style={font=\tiny},label style={font=\scriptsize},enlarge x limits=false,ylabel shift=-3pt,xlabel shift=-3pt]\addplot[fill=ink!55,draw=none] coordinates {(0,0.014) (1,0.052) (2,0.321) (3,1.158) (4,5.189) (5,12.17) (6,26.85) (7,48.73) (8,78.41) (9,119.7) (10,178.6) (11,228.6) (12,264.4) (13,172.5) (14,127.6) (15,97.82) (16,50.87) (17,50.11) (18,14.32) (19,15.81) (20,0.01)};\draw[densely dotted] (axis cs:13,0) -- (axis cs:13,290.9);\end{axis}\begin{axis}[width=72mm,height=42mm,axis y line*=right,axis x line=none,ymode=log,xmin=-1,xmax=21,ylabel={\scriptsize eindeutig (log)},tick label style={font=\tiny},label style={font=\scriptsize},ylabel shift=-3pt]\addplot[mark=*,mark size=0.7pt,draw=black,thin] coordinates {(0,14) (1,52) (2,321) (3,527) (4,1148) (5,1909) (6,1471) (7,1108) (8,698) (9,841) (10,516) (11,256) (12,107) (13,8)};\end{axis}\end{tikzpicture}\end{minipage}
\vspace{3mm}

\noindent\begin{minipage}[t]{60mm}\vspace{0pt}\sbox\hmboard{\chessboard[setfen=8/8/6b1/8/8/6k1/3B4/K7 w - - 0 1]}\usebox\hmboard\\[1mm]
\makebox[\wd\hmboard]{\sffamily\hspace{3mm}\textbf{h\#6.5}\hfill(2+2)}\end{minipage}\hfill
\begin{minipage}[t]{52mm}\vspace{0pt}
\textcolor{ink}{\sffamily\textsc{Lösung}}\par\smallskip\small
\begin{tabular}{@{}r@{\enspace}l@{\enspace}l@{}}
\textbf{1.} & \ldots & \symking{}b2\\
\textbf{2.} & \symking{}g4 & \symking{}c3\\
\textbf{3.} & \symking{}f5 & \symking{}d4\\
\textbf{4.} & \symking{}f6 & \symbishop{}c3\\
\textbf{5.} & \symking{}g7 & \symking{}e5\\
\textbf{6.} & \symking{}h8 & \symking{}f6\\
\textbf{7.} & \symbishop{}h7 & \symking{}f7\#\\
\end{tabular}\par\vspace{3mm}
\qrcode[height=12mm]{https://helpman.komtera.lt/?fen=8/8/6b1/8/8/6k1/3B4/K7&moves=6.5}\\{\sffamily\tiny\color{ink}scannen zum Lösen}
\end{minipage}
\vspace{2mm}

{\raggedright\noindent\textcolor{ink}{\sffamily\textsc{Themen}}\enspace \themetag{Selbstblock} \themetag{ohne Schlag} \themetag{ohne Schach}\par}

\clearpage
\phantomsection\label{hm:KNNvk}
\begin{tcolorbox}[enhanced,colback=ink,colframe=ink,arc=0pt,boxrule=0pt,left=4mm,right=4mm,top=2.5mm,bottom=2.5mm]
\color{white}\sffamily{\LARGE\bfseries KNNvk}\hfill{\Large h\#6.5}\\[0.5mm]
{\small 4 Steine\quad Weiß am Zug}\hfill{\small Nr.~21\qquad\href{https://helpman.komtera.lt/?fen=1k5N/8/8/8/8/8/7N/3K4&moves=6.5}{\color{white}\bfseries lösen\,\ensuremath{\nearrow}}}
\end{tcolorbox}
\vspace{2mm}
\noindent\begin{minipage}[t]{36mm}\vspace{0pt}\sffamily\small\raggedright
\hmfigure{längste eindeutige}{{\Large h\#6.5}}\hmfigure{eindeutig bei dieser Länge}{{\Large 9} Stellungen}\hmfigure{längstes Matt}{{\Large h\#8}}\hmfigure{in der Klasse}{3,78 Millionen Stellungen,\\ 1,57 Millionen mit Hilfsmatt}
\end{minipage}\hfill
\begin{minipage}[t]{80mm}\vspace{0pt}\begin{tikzpicture}\begin{axis}[width=72mm,height=42mm,ybar,bar width=1.6pt,xmin=-1,xmax=17,ymin=0,ymax=389.2,axis y line*=left,axis x line*=bottom,xlabel={\scriptsize Halbzüge bis Matt},ylabel={\scriptsize Stellungen (Tsd.)},tick label style={font=\tiny},label style={font=\scriptsize},enlarge x limits=false,ylabel shift=-3pt,xlabel shift=-3pt]\addplot[fill=ink!55,draw=none] coordinates {(0,0.015) (1,0.077) (2,0.165) (3,1.05) (4,4.559) (5,23.95) (6,55.72) (7,137.9) (8,248.8) (9,287.1) (10,353.8) (11,207.3) (12,169) (13,55.1) (14,21.37) (15,6.634) (16,0.003)};\draw[densely dotted] (axis cs:13,0) -- (axis cs:13,389.2);\end{axis}\begin{axis}[width=72mm,height=42mm,axis y line*=right,axis x line=none,ymode=log,xmin=-1,xmax=17,ylabel={\scriptsize eindeutig (log)},tick label style={font=\tiny},label style={font=\scriptsize},ylabel shift=-3pt]\addplot[mark=*,mark size=0.7pt,draw=black,thin] coordinates {(0,15) (1,75) (2,163) (3,897) (4,2661) (5,3924) (6,4458) (7,3077) (8,3709) (9,983) (10,846) (11,92) (12,27) (13,9)};\end{axis}\end{tikzpicture}\end{minipage}
\vspace{3mm}

\noindent\begin{minipage}[t]{60mm}\vspace{0pt}\sbox\hmboard{\chessboard[setfen=1k5N/8/8/8/8/8/7N/3K4 w - - 0 1]}\usebox\hmboard\\[1mm]
\makebox[\wd\hmboard]{\sffamily\hspace{3mm}\textbf{h\#6.5}\hfill(3+1)}\end{minipage}\hfill
\begin{minipage}[t]{52mm}\vspace{0pt}
\textcolor{ink}{\sffamily\textsc{Lösung}}\par\smallskip\small
\begin{tabular}{@{}r@{\enspace}l@{\enspace}l@{}}
\textbf{1.} & \ldots & \symking{}e2\\
\textbf{2.} & \symking{}c7 & \symking{}f3\\
\textbf{3.} & \symking{}d6 & \symking{}g4\\
\textbf{4.} & \symking{}e5 & \symking{}h5\\
\textbf{5.} & \symking{}f4 & \symknight{}g6+\\
\textbf{6.} & \symking{}g3 & \symknight{}f1+\\
\textbf{7.} & \symking{}h3 & \symknight{}f4\#\\
\end{tabular}\par\vspace{3mm}
\qrcode[height=12mm]{https://helpman.komtera.lt/?fen=1k5N/8/8/8/8/8/7N/3K4&moves=6.5}\\{\sffamily\tiny\color{ink}scannen zum Lösen}
\end{minipage}
\vspace{2mm}

{\raggedright\noindent\textcolor{ink}{\sffamily\textsc{Themen}}\enspace \themetag{reines Matt} \themetag{Mustermatt} \themetag{Idealmatt} \themetag{Spiegelmatt} \themetag{ein Stein zieht} \themetag{ohne Schlag}\par}

\begin{tcolorbox}[enhanced,colback=white,colframe=ink!35,arc=1.5pt,boxrule=0.4pt,left=3mm,right=3mm,top=1.5mm,bottom=1.5mm]
\textcolor{ink}{\sffamily\textsc{veröffentlicht}}\par\smallskip\small \textbf{Bebesi, Gyula}, h\#7 0.1...\quad{\color{ink}\scriptsize diese Stellung}\\{\scriptsize Diagramme und Figuren, No. 1910, 12/12/1966; Catalogue h\# (CHM) avec 6 pieces, No. 624, 1995; Ideal-Mate Encyclopedia Vol.1, No. 930, 1999}
\end{tcolorbox}
\clearpage
\phantomsection\label{hm:KBNvk}
\begin{tcolorbox}[enhanced,colback=ink,colframe=ink,arc=0pt,boxrule=0pt,left=4mm,right=4mm,top=2.5mm,bottom=2.5mm]
\color{white}\sffamily{\LARGE\bfseries KBNvk}\hfill{\Large h\#6}\\[0.5mm]
{\small 4 Steine\quad Schwarz am Zug}\hfill{\small Nr.~22\qquad\href{https://helpman.komtera.lt/?fen=8/6N1/8/8/8/2B5/8/K2k4&moves=6}{\color{white}\bfseries lösen\,\ensuremath{\nearrow}}}
\end{tcolorbox}
\vspace{2mm}
\noindent\begin{minipage}[t]{36mm}\vspace{0pt}\sffamily\small\raggedright
\hmfigure{längste eindeutige}{{\Large h\#6}}\hmfigure{eindeutig bei dieser Länge}{{\Large 23} Stellungen}\hmfigure{längstes Matt}{{\Large h\#8}}\hmfigure{in der Klasse}{3,78 Millionen Stellungen,\\ 3,06 Millionen mit Hilfsmatt}
\end{minipage}\hfill
\begin{minipage}[t]{80mm}\vspace{0pt}\begin{tikzpicture}\begin{axis}[width=72mm,height=42mm,ybar,bar width=1.6pt,xmin=-1,xmax=17,ymin=0,ymax=828,axis y line*=left,axis x line*=bottom,xlabel={\scriptsize Halbzüge bis Matt},ylabel={\scriptsize Stellungen (Tsd.)},tick label style={font=\tiny},label style={font=\scriptsize},enlarge x limits=false,ylabel shift=-3pt,xlabel shift=-3pt]\addplot[fill=ink!55,draw=none] coordinates {(0,0.058) (1,0.23) (2,0.417) (3,2.538) (4,8.084) (5,48.11) (6,109.9) (7,285.3) (8,523.3) (9,594.1) (10,752.8) (11,362.2) (12,283) (13,64.12) (14,27.49) (15,2.989) (16,0.096)};\draw[densely dotted] (axis cs:12,0) -- (axis cs:12,828);\end{axis}\begin{axis}[width=72mm,height=42mm,axis y line*=right,axis x line=none,ymode=log,xmin=-1,xmax=17,ylabel={\scriptsize eindeutig (log)},tick label style={font=\tiny},label style={font=\scriptsize},ylabel shift=-3pt]\addplot[mark=*,mark size=0.7pt,draw=black,thin] coordinates {(0,58) (1,217) (2,392) (3,2210) (4,5863) (5,15118) (6,19105) (7,10247) (8,12152) (9,3015) (10,2251) (11,94) (12,23)};\end{axis}\end{tikzpicture}\end{minipage}
\vspace{3mm}

\noindent\begin{minipage}[t]{60mm}\vspace{0pt}\sbox\hmboard{\chessboard[setfen=8/6N1/8/8/8/2B5/8/K2k4 b - - 0 1]}\usebox\hmboard\\[1mm]
\makebox[\wd\hmboard]{\sffamily\hspace{3mm}\textbf{h\#6}\hfill(3+1)}\end{minipage}\hfill
\begin{minipage}[t]{52mm}\vspace{0pt}
\textcolor{ink}{\sffamily\textsc{Lösung}}\par\smallskip\small
\begin{tabular}{@{}r@{\enspace}l@{\enspace}l@{}}
\textbf{1.} & \symking{}c2 & \symking{}a2\\
\textbf{2.} & \symking{}c1 & \symking{}b3\\
\textbf{3.} & \symking{}b1 & \symking{}c4\\
\textbf{4.} & \symking{}a2 & \symknight{}e6\\
\textbf{5.} & \symking{}a3 & \symbishop{}b4+\\
\textbf{6.} & \symking{}a4 & \symknight{}c5\#\\
\end{tabular}\par\vspace{3mm}
\qrcode[height=12mm]{https://helpman.komtera.lt/?fen=8/6N1/8/8/8/2B5/8/K2k4&moves=6}\\{\sffamily\tiny\color{ink}scannen zum Lösen}
\end{minipage}
\vspace{2mm}

{\raggedright\noindent\textcolor{ink}{\sffamily\textsc{Themen}}\enspace \themetag{Satzspiel} \themetag{ein Stein zieht} \themetag{ohne Schlag}\par}

\clearpage
\phantomsection\label{hm:KQBvk}
\begin{tcolorbox}[enhanced,colback=ink,colframe=ink,arc=0pt,boxrule=0pt,left=4mm,right=4mm,top=2.5mm,bottom=2.5mm]
\color{white}\sffamily{\LARGE\bfseries KQBvk}\hfill{\Large h\#6}\\[0.5mm]
{\small 4 Steine\quad Schwarz am Zug}\hfill{\small Nr.~23\qquad\href{https://helpman.komtera.lt/?fen=7B/7k/8/6Q1/8/8/8/K7&moves=6}{\color{white}\bfseries lösen\,\ensuremath{\nearrow}}}
\end{tcolorbox}
\vspace{2mm}
\noindent\begin{minipage}[t]{36mm}\vspace{0pt}\sffamily\small\raggedright
\hmfigure{längste eindeutige}{{\Large h\#6}}\hmfigure{eindeutig bei dieser Länge}{{\Large 3} Stellungen}\hmfigure{längstes Matt}{{\Large h\#7}}\hmfigure{in der Klasse}{3,78 Millionen Stellungen,\\ 2,63 Millionen mit Hilfsmatt}
\end{minipage}\hfill
\begin{minipage}[t]{80mm}\vspace{0pt}\begin{tikzpicture}\begin{axis}[width=72mm,height=42mm,ybar,bar width=1.6pt,xmin=-1,xmax=15,ymin=0,ymax=1.392,axis y line*=left,axis x line*=bottom,xlabel={\scriptsize Halbzüge bis Matt},ylabel={\scriptsize Stellungen (Mio.)},tick label style={font=\tiny},label style={font=\scriptsize},enlarge x limits=false,ylabel shift=-3pt,xlabel shift=-3pt]\addplot[fill=ink!55,draw=none] coordinates {(0,0.008228) (1,0.05076) (2,0.1603) (3,0.5118) (4,1.265) (5,0.3989) (6,0.2309) (7,0.000893) (8,0.002146) (10,0.00072) (12,0.000114) (14,4e-06)};\draw[densely dotted] (axis cs:12,0) -- (axis cs:12,1.392);\end{axis}\begin{axis}[width=72mm,height=42mm,axis y line*=right,axis x line=none,ymode=log,xmin=-1,xmax=15,ylabel={\scriptsize eindeutig (log)},tick label style={font=\tiny},label style={font=\scriptsize},ylabel shift=-3pt]\addplot[mark=*,mark size=0.7pt,draw=black,thin] coordinates {(0,8228) (1,40448) (2,125768) (3,165738) (4,107474) (5,202) (6,95) (8,21) (10,4) (12,3)};\end{axis}\end{tikzpicture}\end{minipage}
\vspace{3mm}

\noindent\begin{minipage}[t]{60mm}\vspace{0pt}\sbox\hmboard{\chessboard[setfen=7B/7k/8/6Q1/8/8/8/K7 b - - 0 1]}\usebox\hmboard\\[1mm]
\makebox[\wd\hmboard]{\sffamily\hspace{3mm}\textbf{h\#6}\hfill(3+1)}\end{minipage}\hfill
\begin{minipage}[t]{52mm}\vspace{0pt}
\textcolor{ink}{\sffamily\textsc{Lösung}}\par\smallskip\small
\begin{tabular}{@{}r@{\enspace}l@{\enspace}l@{}}
\textbf{1.} & \symking{}xh8 & \symking{}b2\\
\textbf{2.} & \symking{}h7 & \symking{}c3\\
\textbf{3.} & \symking{}h8 & \symking{}d4\\
\textbf{4.} & \symking{}h7 & \symking{}e5\\
\textbf{5.} & \symking{}h8 & \symking{}f6\\
\textbf{6.} & \symking{}h7 & \symqueen{}g7\#\\
\end{tabular}\par\vspace{3mm}
\qrcode[height=12mm]{https://helpman.komtera.lt/?fen=7B/7k/8/6Q1/8/8/8/K7&moves=6}\\{\sffamily\tiny\color{ink}scannen zum Lösen}
\end{minipage}
\vspace{2mm}

{\raggedright\noindent\textcolor{ink}{\sffamily\textsc{Themen}}\enspace \themetag{Switchback} \themetag{ein Stein zieht} \themetag{Pendel} \themetag{ohne Schach}\par\smallskip{\footnotesize\textit{Hinweis:} die Lösung beginnt mit einem Schlag.}\par}

\clearpage
\phantomsection\label{hm:KQNvk}
\begin{tcolorbox}[enhanced,colback=ink,colframe=ink,arc=0pt,boxrule=0pt,left=4mm,right=4mm,top=2.5mm,bottom=2.5mm]
\color{white}\sffamily{\LARGE\bfseries KQNvk}\hfill{\Large h\#6}\\[0.5mm]
{\small 4 Steine\quad Schwarz am Zug}\hfill{\small Nr.~24\qquad\href{https://helpman.komtera.lt/?fen=7N/7k/8/6Q1/8/8/8/K7&moves=6}{\color{white}\bfseries lösen\,\ensuremath{\nearrow}}}
\end{tcolorbox}
\vspace{2mm}
\noindent\begin{minipage}[t]{36mm}\vspace{0pt}\sffamily\small\raggedright
\hmfigure{längste eindeutige}{{\Large h\#6}}\hmfigure{eindeutig bei dieser Länge}{{\Large 3} Stellungen}\hmfigure{längstes Matt}{{\Large h\#7}}\hmfigure{in der Klasse}{3,78 Millionen Stellungen,\\ 2,71 Millionen mit Hilfsmatt}
\end{minipage}\hfill
\begin{minipage}[t]{80mm}\vspace{0pt}\begin{tikzpicture}\begin{axis}[width=72mm,height=42mm,ybar,bar width=1.6pt,xmin=-1,xmax=15,ymin=0,ymax=922.1,axis y line*=left,axis x line*=bottom,xlabel={\scriptsize Halbzüge bis Matt},ylabel={\scriptsize Stellungen (Tsd.)},tick label style={font=\tiny},label style={font=\scriptsize},enlarge x limits=false,ylabel shift=-3pt,xlabel shift=-3pt]\addplot[fill=ink!55,draw=none] coordinates {(0,4.75) (1,35.98) (2,103.5) (3,309.9) (4,838.2) (5,600) (6,679.1) (7,84.74) (8,49.75) (9,0.199) (10,0.32) (12,0.053) (14,0.002)};\draw[densely dotted] (axis cs:12,0) -- (axis cs:12,922.1);\end{axis}\begin{axis}[width=72mm,height=42mm,axis y line*=right,axis x line=none,ymode=log,xmin=-1,xmax=15,ylabel={\scriptsize eindeutig (log)},tick label style={font=\tiny},label style={font=\scriptsize},ylabel shift=-3pt]\addplot[mark=*,mark size=0.7pt,draw=black,thin] coordinates {(0,4750) (1,27785) (2,79125) (3,77957) (4,103775) (5,9404) (6,4821) (7,17) (8,11) (10,3) (12,3)};\end{axis}\end{tikzpicture}\end{minipage}
\vspace{3mm}

\noindent\begin{minipage}[t]{60mm}\vspace{0pt}\sbox\hmboard{\chessboard[setfen=7N/7k/8/6Q1/8/8/8/K7 b - - 0 1]}\usebox\hmboard\\[1mm]
\makebox[\wd\hmboard]{\sffamily\hspace{3mm}\textbf{h\#6}\hfill(3+1)}\end{minipage}\hfill
\begin{minipage}[t]{52mm}\vspace{0pt}
\textcolor{ink}{\sffamily\textsc{Lösung}}\par\smallskip\small
\begin{tabular}{@{}r@{\enspace}l@{\enspace}l@{}}
\textbf{1.} & \symking{}xh8 & \symking{}b2\\
\textbf{2.} & \symking{}h7 & \symking{}c3\\
\textbf{3.} & \symking{}h8 & \symking{}d4\\
\textbf{4.} & \symking{}h7 & \symking{}e5\\
\textbf{5.} & \symking{}h8 & \symking{}f6\\
\textbf{6.} & \symking{}h7 & \symqueen{}g7\#\\
\end{tabular}\par\vspace{3mm}
\qrcode[height=12mm]{https://helpman.komtera.lt/?fen=7N/7k/8/6Q1/8/8/8/K7&moves=6}\\{\sffamily\tiny\color{ink}scannen zum Lösen}
\end{minipage}
\vspace{2mm}

{\raggedright\noindent\textcolor{ink}{\sffamily\textsc{Themen}}\enspace \themetag{Switchback} \themetag{ein Stein zieht} \themetag{Pendel} \themetag{ohne Schach}\par\smallskip{\footnotesize\textit{Hinweis:} die Lösung beginnt mit einem Schlag.}\par}

\clearpage
\phantomsection\label{hm:KQQvk}
\begin{tcolorbox}[enhanced,colback=ink,colframe=ink,arc=0pt,boxrule=0pt,left=4mm,right=4mm,top=2.5mm,bottom=2.5mm]
\color{white}\sffamily{\LARGE\bfseries KQQvk}\hfill{\Large h\#6}\\[0.5mm]
{\small 4 Steine\quad Schwarz am Zug}\hfill{\small Nr.~25\qquad\href{https://helpman.komtera.lt/?fen=7Q/7k/8/6Q1/8/8/8/K7&moves=6}{\color{white}\bfseries lösen\,\ensuremath{\nearrow}}}
\end{tcolorbox}
\vspace{2mm}
\noindent\begin{minipage}[t]{36mm}\vspace{0pt}\sffamily\small\raggedright
\hmfigure{längste eindeutige}{{\Large h\#6}}\hmfigure{eindeutig bei dieser Länge}{{\Large 2} Stellungen}\hmfigure{längstes Matt}{{\Large h\#7}}\hmfigure{in der Klasse}{3,78 Millionen Stellungen,\\ 1,19 Millionen mit Hilfsmatt}
\end{minipage}\hfill
\begin{minipage}[t]{80mm}\vspace{0pt}\begin{tikzpicture}\begin{axis}[width=72mm,height=42mm,ybar,bar width=1.6pt,xmin=-1,xmax=15,ymin=0,ymax=472.3,axis y line*=left,axis x line*=bottom,xlabel={\scriptsize Halbzüge bis Matt},ylabel={\scriptsize Stellungen (Tsd.)},tick label style={font=\tiny},label style={font=\scriptsize},enlarge x limits=false,ylabel shift=-3pt,xlabel shift=-3pt]\addplot[fill=ink!55,draw=none] coordinates {(0,31.52) (1,95.65) (2,347.1) (3,256.7) (4,429.4) (5,1.45) (6,11.79) (8,12.25) (10,3.834) (12,0.664) (14,0.004)};\draw[densely dotted] (axis cs:12,0) -- (axis cs:12,472.3);\end{axis}\begin{axis}[width=72mm,height=42mm,axis y line*=right,axis x line=none,ymode=log,xmin=-1,xmax=15,ylabel={\scriptsize eindeutig (log)},tick label style={font=\tiny},label style={font=\scriptsize},ylabel shift=-3pt]\addplot[mark=*,mark size=0.7pt,draw=black,thin] coordinates {(0,31523) (1,47189) (2,133257) (3,2198) (4,1153) (6,743) (8,166) (10,3) (12,2)};\end{axis}\end{tikzpicture}\end{minipage}
\vspace{3mm}

\noindent\begin{minipage}[t]{60mm}\vspace{0pt}\sbox\hmboard{\chessboard[setfen=7Q/7k/8/6Q1/8/8/8/K7 b - - 0 1]}\usebox\hmboard\\[1mm]
\makebox[\wd\hmboard]{\sffamily\hspace{3mm}\textbf{h\#6}\hfill(3+1)}\end{minipage}\hfill
\begin{minipage}[t]{52mm}\vspace{0pt}
\textcolor{ink}{\sffamily\textsc{Lösung}}\par\smallskip\small
\begin{tabular}{@{}r@{\enspace}l@{\enspace}l@{}}
\textbf{1.} & \symking{}xh8 & \symking{}b2\\
\textbf{2.} & \symking{}h7 & \symking{}c3\\
\textbf{3.} & \symking{}h8 & \symking{}d4\\
\textbf{4.} & \symking{}h7 & \symking{}e5\\
\textbf{5.} & \symking{}h8 & \symking{}f6\\
\textbf{6.} & \symking{}h7 & \symqueen{}g7\#\\
\end{tabular}\par\vspace{3mm}
\qrcode[height=12mm]{https://helpman.komtera.lt/?fen=7Q/7k/8/6Q1/8/8/8/K7&moves=6}\\{\sffamily\tiny\color{ink}scannen zum Lösen}
\end{minipage}
\vspace{2mm}

{\raggedright\noindent\textcolor{ink}{\sffamily\textsc{Themen}}\enspace \themetag{Switchback} \themetag{ein Stein zieht} \themetag{Pendel} \themetag{ohne Schach}\par\smallskip{\footnotesize\textit{Hinweis:} die Partei am Zug steht im Diagramm im Schach; die Lösung beginnt mit einem Schlag.}\par}

\clearpage
\phantomsection\label{hm:KQRvk}
\begin{tcolorbox}[enhanced,colback=ink,colframe=ink,arc=0pt,boxrule=0pt,left=4mm,right=4mm,top=2.5mm,bottom=2.5mm]
\color{white}\sffamily{\LARGE\bfseries KQRvk}\hfill{\Large h\#6}\\[0.5mm]
{\small 4 Steine\quad Schwarz am Zug}\hfill{\small Nr.~26\qquad\href{https://helpman.komtera.lt/?fen=7R/7k/8/6Q1/8/8/8/K7&moves=6}{\color{white}\bfseries lösen\,\ensuremath{\nearrow}}}
\end{tcolorbox}
\vspace{2mm}
\noindent\begin{minipage}[t]{36mm}\vspace{0pt}\sffamily\small\raggedright
\hmfigure{längste eindeutige}{{\Large h\#6}}\hmfigure{eindeutig bei dieser Länge}{{\Large 2} Stellungen}\hmfigure{längstes Matt}{{\Large h\#7}}\hmfigure{in der Klasse}{3,78 Millionen Stellungen,\\ 2,55 Millionen mit Hilfsmatt}
\end{minipage}\hfill
\begin{minipage}[t]{80mm}\vspace{0pt}\begin{tikzpicture}\begin{axis}[width=72mm,height=42mm,ybar,bar width=1.6pt,xmin=-1,xmax=15,ymin=0,ymax=1.345,axis y line*=left,axis x line*=bottom,xlabel={\scriptsize Halbzüge bis Matt},ylabel={\scriptsize Stellungen (Mio.)},tick label style={font=\tiny},label style={font=\scriptsize},enlarge x limits=false,ylabel shift=-3pt,xlabel shift=-3pt]\addplot[fill=ink!55,draw=none] coordinates {(0,0.02845) (1,0.1113) (2,0.308) (3,0.6007) (4,1.223) (5,0.1507) (6,0.1052) (7,0.001378) (8,0.0152) (10,0.00862) (12,0.001574) (14,5.8e-05)};\draw[densely dotted] (axis cs:12,0) -- (axis cs:12,1.345);\end{axis}\begin{axis}[width=72mm,height=42mm,axis y line*=right,axis x line=none,ymode=log,xmin=-1,xmax=15,ylabel={\scriptsize eindeutig (log)},tick label style={font=\tiny},label style={font=\scriptsize},ylabel shift=-3pt]\addplot[mark=*,mark size=0.7pt,draw=black,thin] coordinates {(0,28454) (1,69157) (2,138399) (3,59987) (4,37796) (5,360) (6,1323) (8,108) (10,7) (12,2)};\end{axis}\end{tikzpicture}\end{minipage}
\vspace{3mm}

\noindent\begin{minipage}[t]{60mm}\vspace{0pt}\sbox\hmboard{\chessboard[setfen=7R/7k/8/6Q1/8/8/8/K7 b - - 0 1]}\usebox\hmboard\\[1mm]
\makebox[\wd\hmboard]{\sffamily\hspace{3mm}\textbf{h\#6}\hfill(3+1)}\end{minipage}\hfill
\begin{minipage}[t]{52mm}\vspace{0pt}
\textcolor{ink}{\sffamily\textsc{Lösung}}\par\smallskip\small
\begin{tabular}{@{}r@{\enspace}l@{\enspace}l@{}}
\textbf{1.} & \symking{}xh8 & \symking{}b2\\
\textbf{2.} & \symking{}h7 & \symking{}c3\\
\textbf{3.} & \symking{}h8 & \symking{}d4\\
\textbf{4.} & \symking{}h7 & \symking{}e5\\
\textbf{5.} & \symking{}h8 & \symking{}f6\\
\textbf{6.} & \symking{}h7 & \symqueen{}g7\#\\
\end{tabular}\par\vspace{3mm}
\qrcode[height=12mm]{https://helpman.komtera.lt/?fen=7R/7k/8/6Q1/8/8/8/K7&moves=6}\\{\sffamily\tiny\color{ink}scannen zum Lösen}
\end{minipage}
\vspace{2mm}

{\raggedright\noindent\textcolor{ink}{\sffamily\textsc{Themen}}\enspace \themetag{Switchback} \themetag{ein Stein zieht} \themetag{Pendel} \themetag{ohne Schach}\par\smallskip{\footnotesize\textit{Hinweis:} die Partei am Zug steht im Diagramm im Schach; die Lösung beginnt mit einem Schlag.}\par}

\clearpage
\phantomsection\label{hm:KQvkn}
\begin{tcolorbox}[enhanced,colback=ink,colframe=ink,arc=0pt,boxrule=0pt,left=4mm,right=4mm,top=2.5mm,bottom=2.5mm]
\color{white}\sffamily{\LARGE\bfseries KQvkn}\hfill{\Large h\#6}\\[0.5mm]
{\small 4 Steine\quad Schwarz am Zug}\hfill{\small Nr.~27\qquad\href{https://helpman.komtera.lt/?fen=7k/8/8/6Q1/8/8/1n6/K7&moves=6}{\color{white}\bfseries lösen\,\ensuremath{\nearrow}}}
\end{tcolorbox}
\vspace{2mm}
\noindent\begin{minipage}[t]{36mm}\vspace{0pt}\sffamily\small\raggedright
\hmfigure{längste eindeutige}{{\Large h\#6}}\hmfigure{eindeutig bei dieser Länge}{{\Large 4} Stellungen}\hmfigure{längstes Matt}{{\Large h\#7}}\hmfigure{in der Klasse}{3,78 Millionen Stellungen,\\ 2,67 Millionen mit Hilfsmatt}
\end{minipage}\hfill
\begin{minipage}[t]{80mm}\vspace{0pt}\begin{tikzpicture}\begin{axis}[width=72mm,height=42mm,ybar,bar width=1.6pt,xmin=-1,xmax=15,ymin=0,ymax=675.4,axis y line*=left,axis x line*=bottom,xlabel={\scriptsize Halbzüge bis Matt},ylabel={\scriptsize Stellungen (Tsd.)},tick label style={font=\tiny},label style={font=\scriptsize},enlarge x limits=false,ylabel shift=-3pt,xlabel shift=-3pt]\addplot[fill=ink!55,draw=none] coordinates {(0,2.096) (1,14.43) (2,35.62) (3,105.9) (4,239.3) (5,353.7) (6,614) (7,448.5) (8,527.3) (9,177.4) (10,118.4) (11,19.3) (12,12.61) (13,0.019) (14,0.012)};\draw[densely dotted] (axis cs:12,0) -- (axis cs:12,675.4);\end{axis}\begin{axis}[width=72mm,height=42mm,axis y line*=right,axis x line=none,ymode=log,xmin=-1,xmax=15,ylabel={\scriptsize eindeutig (log)},tick label style={font=\tiny},label style={font=\scriptsize},ylabel shift=-3pt]\addplot[mark=*,mark size=0.7pt,draw=black,thin] coordinates {(0,2096) (1,9343) (2,14346) (3,25797) (4,34491) (5,22810) (6,20088) (7,7137) (8,4018) (9,365) (10,82) (11,4) (12,4)};\end{axis}\end{tikzpicture}\end{minipage}
\vspace{3mm}

\noindent\begin{minipage}[t]{60mm}\vspace{0pt}\sbox\hmboard{\chessboard[setfen=7k/8/8/6Q1/8/8/1n6/K7 b - - 0 1]}\usebox\hmboard\\[1mm]
\makebox[\wd\hmboard]{\sffamily\hspace{3mm}\textbf{h\#6}\hfill(2+2)}\end{minipage}\hfill
\begin{minipage}[t]{52mm}\vspace{0pt}
\textcolor{ink}{\sffamily\textsc{Lösung}}\par\smallskip\small
\begin{tabular}{@{}r@{\enspace}l@{\enspace}l@{}}
\textbf{1.} & \symking{}h7 & \symking{}xb2\\
\textbf{2.} & \symking{}h8 & \symking{}c3\\
\textbf{3.} & \symking{}h7 & \symking{}d4\\
\textbf{4.} & \symking{}h8 & \symking{}e5\\
\textbf{5.} & \symking{}h7 & \symking{}f6\\
\textbf{6.} & \symking{}h8 & \symqueen{}g7\#\\
\end{tabular}\par\vspace{3mm}
\qrcode[height=12mm]{https://helpman.komtera.lt/?fen=7k/8/8/6Q1/8/8/1n6/K7&moves=6}\\{\sffamily\tiny\color{ink}scannen zum Lösen}
\end{minipage}
\vspace{2mm}

{\raggedright\noindent\textcolor{ink}{\sffamily\textsc{Themen}}\enspace \themetag{Satzspiel} \themetag{reines Matt} \themetag{Mustermatt} \themetag{Idealmatt} \themetag{Switchback} \themetag{ein Stein zieht} \themetag{Pendel} \themetag{ohne Schach}\par}

\clearpage
\phantomsection\label{hm:KQvkp}
\begin{tcolorbox}[enhanced,colback=ink,colframe=ink,arc=0pt,boxrule=0pt,left=4mm,right=4mm,top=2.5mm,bottom=2.5mm]
\color{white}\sffamily{\LARGE\bfseries KQvkp}\hfill{\Large h\#6}\\[0.5mm]
{\small 4 Steine\quad Schwarz am Zug}\hfill{\small Nr.~28\qquad\href{https://helpman.komtera.lt/?fen=K7/1p6/8/8/8/8/4Q3/5k2&moves=6}{\color{white}\bfseries lösen\,\ensuremath{\nearrow}}}
\end{tcolorbox}
\vspace{2mm}
\noindent\begin{minipage}[t]{36mm}\vspace{0pt}\sffamily\small\raggedright
\hmfigure{längste eindeutige}{{\Large h\#6}}\hmfigure{eindeutig bei dieser Länge}{{\Large 4} Stellungen}\hmfigure{längstes Matt}{{\Large h\#7}}\hmfigure{in der Klasse}{11,1 Millionen Stellungen,\\ 8,29 Millionen mit Hilfsmatt}
\end{minipage}\hfill
\begin{minipage}[t]{80mm}\vspace{0pt}\begin{tikzpicture}\begin{axis}[width=72mm,height=42mm,ybar,bar width=1.6pt,xmin=-1,xmax=15,ymin=0,ymax=2.232,axis y line*=left,axis x line*=bottom,xlabel={\scriptsize Halbzüge bis Matt},ylabel={\scriptsize Stellungen (Mio.)},tick label style={font=\tiny},label style={font=\scriptsize},enlarge x limits=false,ylabel shift=-3pt,xlabel shift=-3pt]\addplot[fill=ink!55,draw=none] coordinates {(0,0.008156) (1,0.05441) (2,0.1241) (3,0.3545) (4,0.8799) (5,1.124) (6,2.029) (7,1.333) (8,1.555) (9,0.457) (10,0.2997) (11,0.04837) (12,0.02679) (13,8.9e-05) (14,0.000126)};\draw[densely dotted] (axis cs:12,0) -- (axis cs:12,2.232);\end{axis}\begin{axis}[width=72mm,height=42mm,axis y line*=right,axis x line=none,ymode=log,xmin=-1,xmax=15,ylabel={\scriptsize eindeutig (log)},tick label style={font=\tiny},label style={font=\scriptsize},ylabel shift=-3pt]\addplot[mark=*,mark size=0.7pt,draw=black,thin] coordinates {(0,8156) (1,34532) (2,62533) (3,87058) (4,124844) (5,68251) (6,54498) (7,17747) (8,8860) (9,259) (10,106) (11,6) (12,4)};\end{axis}\end{tikzpicture}\end{minipage}
\vspace{3mm}

\noindent\begin{minipage}[t]{60mm}\vspace{0pt}\sbox\hmboard{\chessboard[setfen=K7/1p6/8/8/8/8/4Q3/5k2 b - - 0 1]}\usebox\hmboard\\[1mm]
\makebox[\wd\hmboard]{\sffamily\hspace{3mm}\textbf{h\#6}\hfill(2+2)}\end{minipage}\hfill
\begin{minipage}[t]{52mm}\vspace{0pt}
\textcolor{ink}{\sffamily\textsc{Lösung}}\par\smallskip\small
\begin{tabular}{@{}r@{\enspace}l@{\enspace}l@{}}
\textbf{1.} & \symking{}g1 & \symking{}xb7\\
\textbf{2.} & \symking{}h1 & \symking{}c6\\
\textbf{3.} & \symking{}g1 & \symking{}d5\\
\textbf{4.} & \symking{}h1 & \symking{}e4\\
\textbf{5.} & \symking{}g1 & \symking{}f3\\
\textbf{6.} & \symking{}h1 & \symqueen{}g2\#\\
\end{tabular}\par\vspace{3mm}
\qrcode[height=12mm]{https://helpman.komtera.lt/?fen=K7/1p6/8/8/8/8/4Q3/5k2&moves=6}\\{\sffamily\tiny\color{ink}scannen zum Lösen}
\end{minipage}
\vspace{2mm}

{\raggedright\noindent\textcolor{ink}{\sffamily\textsc{Themen}}\enspace \themetag{reines Matt} \themetag{Mustermatt} \themetag{Idealmatt} \themetag{Switchback} \themetag{ein Stein zieht} \themetag{ohne Schach}\par\smallskip{\footnotesize\textit{Hinweis:} die Partei am Zug steht im Diagramm im Schach.}\par}

\clearpage
\phantomsection\label{hm:KRBvk}
\begin{tcolorbox}[enhanced,colback=ink,colframe=ink,arc=0pt,boxrule=0pt,left=4mm,right=4mm,top=2.5mm,bottom=2.5mm]
\color{white}\sffamily{\LARGE\bfseries KRBvk}\hfill{\Large h\#6}\\[0.5mm]
{\small 4 Steine\quad Schwarz am Zug}\hfill{\small Nr.~29\qquad\href{https://helpman.komtera.lt/?fen=6R1/8/7k/8/8/8/8/K2B4&moves=6}{\color{white}\bfseries lösen\,\ensuremath{\nearrow}}}
\end{tcolorbox}
\vspace{2mm}
\noindent\begin{minipage}[t]{36mm}\vspace{0pt}\sffamily\small\raggedright
\hmfigure{längste eindeutige}{{\Large h\#6}}\hmfigure{eindeutig bei dieser Länge}{{\Large 20} Stellungen}\hmfigure{längstes Matt}{{\Large h\#7}}\hmfigure{in der Klasse}{3,78 Millionen Stellungen,\\ 2,87 Millionen mit Hilfsmatt}
\end{minipage}\hfill
\begin{minipage}[t]{80mm}\vspace{0pt}\begin{tikzpicture}\begin{axis}[width=72mm,height=42mm,ybar,bar width=1.6pt,xmin=-1,xmax=15,ymin=0,ymax=691.1,axis y line*=left,axis x line*=bottom,xlabel={\scriptsize Halbzüge bis Matt},ylabel={\scriptsize Stellungen (Tsd.)},tick label style={font=\tiny},label style={font=\scriptsize},enlarge x limits=false,ylabel shift=-3pt,xlabel shift=-3pt]\addplot[fill=ink!55,draw=none] coordinates {(0,2.042) (1,11.5) (2,19.57) (3,76.91) (4,200.8) (5,336.8) (6,628.3) (7,459.8) (8,605.5) (9,248.7) (10,214.4) (11,36.02) (12,26.5) (13,1.195) (14,0.746)};\draw[densely dotted] (axis cs:12,0) -- (axis cs:12,691.1);\end{axis}\begin{axis}[width=72mm,height=42mm,axis y line*=right,axis x line=none,ymode=log,xmin=-1,xmax=15,ylabel={\scriptsize eindeutig (log)},tick label style={font=\tiny},label style={font=\scriptsize},ylabel shift=-3pt]\addplot[mark=*,mark size=0.7pt,draw=black,thin] coordinates {(0,2042) (1,11174) (2,18417) (3,44230) (4,90618) (5,53934) (6,46527) (7,6307) (8,3673) (9,387) (10,176) (11,14) (12,20)};\end{axis}\end{tikzpicture}\end{minipage}
\vspace{3mm}

\noindent\begin{minipage}[t]{60mm}\vspace{0pt}\sbox\hmboard{\chessboard[setfen=6R1/8/7k/8/8/8/8/K2B4 b - - 0 1]}\usebox\hmboard\\[1mm]
\makebox[\wd\hmboard]{\sffamily\hspace{3mm}\textbf{h\#6}\hfill(3+1)}\end{minipage}\hfill
\begin{minipage}[t]{52mm}\vspace{0pt}
\textcolor{ink}{\sffamily\textsc{Lösung}}\par\smallskip\small
\begin{tabular}{@{}r@{\enspace}l@{\enspace}l@{}}
\textbf{1.} & \symking{}h7 & \symking{}b2\\
\textbf{2.} & \symking{}h6 & \symking{}c3\\
\textbf{3.} & \symking{}h7 & \symking{}d4\\
\textbf{4.} & \symking{}h6 & \symking{}e5\\
\textbf{5.} & \symking{}h7 & \symking{}f6\\
\textbf{6.} & \symking{}h6 & \symrook{}h8\#\\
\end{tabular}\par\vspace{3mm}
\qrcode[height=12mm]{https://helpman.komtera.lt/?fen=6R1/8/7k/8/8/8/8/K2B4&moves=6}\\{\sffamily\tiny\color{ink}scannen zum Lösen}
\end{minipage}
\vspace{2mm}

{\raggedright\noindent\textcolor{ink}{\sffamily\textsc{Themen}}\enspace \themetag{Satzspiel} \themetag{Spiegelmatt} \themetag{Switchback} \themetag{ein Stein zieht} \themetag{Pendel} \themetag{ohne Schlag} \themetag{ohne Schach}\par}

\clearpage
\phantomsection\label{hm:KRPvk}
\begin{tcolorbox}[enhanced,colback=ink,colframe=ink,arc=0pt,boxrule=0pt,left=4mm,right=4mm,top=2.5mm,bottom=2.5mm]
\color{white}\sffamily{\LARGE\bfseries KRPvk}\hfill{\Large h\#6}\\[0.5mm]
{\small 4 Steine\quad Schwarz am Zug}\hfill{\small Nr.~30\qquad\href{https://helpman.komtera.lt/?fen=8/7k/6R1/8/8/8/P7/K7&moves=6}{\color{white}\bfseries lösen\,\ensuremath{\nearrow}}}
\end{tcolorbox}
\vspace{2mm}
\noindent\begin{minipage}[t]{36mm}\vspace{0pt}\sffamily\small\raggedright
\hmfigure{längste eindeutige}{{\Large h\#6}}\hmfigure{eindeutig bei dieser Länge}{{\Large 25} Stellungen}\hmfigure{längstes Matt}{{\Large h\#8}}\hmfigure{in der Klasse}{11,1 Millionen Stellungen,\\ 9,06 Millionen mit Hilfsmatt}
\end{minipage}\hfill
\begin{minipage}[t]{80mm}\vspace{0pt}\begin{tikzpicture}\begin{axis}[width=72mm,height=42mm,ybar,bar width=1.6pt,xmin=-1,xmax=17,ymin=0,ymax=2.227,axis y line*=left,axis x line*=bottom,xlabel={\scriptsize Halbzüge bis Matt},ylabel={\scriptsize Stellungen (Mio.)},tick label style={font=\tiny},label style={font=\scriptsize},enlarge x limits=false,ylabel shift=-3pt,xlabel shift=-3pt]\addplot[fill=ink!55,draw=none] coordinates {(0,0.005051) (1,0.04032) (2,0.07291) (3,0.3979) (4,0.8997) (5,1.381) (6,2.025) (7,1.355) (8,1.57) (9,0.6839) (10,0.4908) (11,0.07927) (12,0.05528) (13,0.000871) (14,0.002824) (16,0.000306)};\draw[densely dotted] (axis cs:12,0) -- (axis cs:12,2.227);\end{axis}\begin{axis}[width=72mm,height=42mm,axis y line*=right,axis x line=none,ymode=log,xmin=-1,xmax=17,ylabel={\scriptsize eindeutig (log)},tick label style={font=\tiny},label style={font=\scriptsize},ylabel shift=-3pt]\addplot[mark=*,mark size=0.7pt,draw=black,thin] coordinates {(0,5051) (1,31699) (2,50616) (3,150201) (4,233403) (5,106050) (6,66071) (7,7521) (8,4068) (9,424) (10,181) (11,24) (12,25)};\end{axis}\end{tikzpicture}\end{minipage}
\vspace{3mm}

\noindent\begin{minipage}[t]{60mm}\vspace{0pt}\sbox\hmboard{\chessboard[setfen=8/7k/6R1/8/8/8/P7/K7 b - - 0 1]}\usebox\hmboard\\[1mm]
\makebox[\wd\hmboard]{\sffamily\hspace{3mm}\textbf{h\#6}\hfill(3+1)}\end{minipage}\hfill
\begin{minipage}[t]{52mm}\vspace{0pt}
\textcolor{ink}{\sffamily\textsc{Lösung}}\par\smallskip\small
\begin{tabular}{@{}r@{\enspace}l@{\enspace}l@{}}
\textbf{1.} & \symking{}h8 & a4\\
\textbf{2.} & \symking{}h7 & a5\\
\textbf{3.} & \symking{}h8 & a6\\
\textbf{4.} & \symking{}h7 & a7\\
\textbf{5.} & \symking{}h8 & a8=\symqueen{}+\\
\textbf{6.} & \symking{}h7 & \symqueen{}g8\#\\
\end{tabular}\par\vspace{3mm}
\qrcode[height=12mm]{https://helpman.komtera.lt/?fen=8/7k/6R1/8/8/8/P7/K7&moves=6}\\{\sffamily\tiny\color{ink}scannen zum Lösen}
\end{minipage}
\vspace{2mm}

{\raggedright\noindent\textcolor{ink}{\sffamily\textsc{Themen}}\enspace \themetag{Umwandlung} \themetag{Excelsior} \themetag{Switchback} \themetag{ein Stein zieht} \themetag{Pendel} \themetag{ohne Schlag}\par}

\clearpage
\phantomsection\label{hm:KRvkp}
\begin{tcolorbox}[enhanced,colback=ink,colframe=ink,arc=0pt,boxrule=0pt,left=4mm,right=4mm,top=2.5mm,bottom=2.5mm]
\color{white}\sffamily{\LARGE\bfseries KRvkp}\hfill{\Large h\#6}\\[0.5mm]
{\small 4 Steine\quad Schwarz am Zug}\hfill{\small Nr.~31\qquad\href{https://helpman.komtera.lt/?fen=8/8/7k/6R1/8/1p6/8/K7&moves=6}{\color{white}\bfseries lösen\,\ensuremath{\nearrow}}}
\end{tcolorbox}
\vspace{2mm}
\noindent\begin{minipage}[t]{36mm}\vspace{0pt}\sffamily\small\raggedright
\hmfigure{längste eindeutige}{{\Large h\#6}}\hmfigure{eindeutig bei dieser Länge}{{\Large 84} Stellungen}\hmfigure{längstes Matt}{{\Large h\#7}}\hmfigure{in der Klasse}{11,1 Millionen Stellungen,\\ 9,03 Millionen mit Hilfsmatt}
\end{minipage}\hfill
\begin{minipage}[t]{80mm}\vspace{0pt}\begin{tikzpicture}\begin{axis}[width=72mm,height=42mm,ybar,bar width=1.6pt,xmin=-1,xmax=15,ymin=0,ymax=2.016,axis y line*=left,axis x line*=bottom,xlabel={\scriptsize Halbzüge bis Matt},ylabel={\scriptsize Stellungen (Mio.)},tick label style={font=\tiny},label style={font=\scriptsize},enlarge x limits=false,ylabel shift=-3pt,xlabel shift=-3pt]\addplot[fill=ink!55,draw=none] coordinates {(0,0.004672) (1,0.03052) (2,0.06531) (3,0.2485) (4,0.5687) (5,0.9588) (6,1.558) (7,1.448) (8,1.833) (9,1.159) (10,0.8176) (11,0.2011) (12,0.1274) (13,0.003985) (14,0.002503)};\draw[densely dotted] (axis cs:12,0) -- (axis cs:12,2.016);\end{axis}\begin{axis}[width=72mm,height=42mm,axis y line*=right,axis x line=none,ymode=log,xmin=-1,xmax=15,ylabel={\scriptsize eindeutig (log)},tick label style={font=\tiny},label style={font=\scriptsize},ylabel shift=-3pt]\addplot[mark=*,mark size=0.7pt,draw=black,thin] coordinates {(0,4672) (1,30522) (2,49698) (3,111619) (4,208806) (5,94540) (6,50666) (7,9205) (8,5514) (9,997) (10,427) (11,61) (12,84)};\end{axis}\end{tikzpicture}\end{minipage}
\vspace{3mm}

\noindent\begin{minipage}[t]{60mm}\vspace{0pt}\sbox\hmboard{\chessboard[setfen=8/8/7k/6R1/8/1p6/8/K7 b - - 0 1]}\usebox\hmboard\\[1mm]
\makebox[\wd\hmboard]{\sffamily\hspace{3mm}\textbf{h\#6}\hfill(2+2)}\end{minipage}\hfill
\begin{minipage}[t]{52mm}\vspace{0pt}
\textcolor{ink}{\sffamily\textsc{Lösung}}\par\smallskip\small
\begin{tabular}{@{}r@{\enspace}l@{\enspace}l@{}}
\textbf{1.} & \symking{}h7 & \symking{}b2\\
\textbf{2.} & \symking{}h8 & \symking{}c3\\
\textbf{3.} & b2 & \symking{}d4\\
\textbf{4.} & b1=\symrook{} & \symking{}e5\\
\textbf{5.} & \symrook{}b8 & \symking{}f6\\
\textbf{6.} & \symrook{}g8 & \symrook{}h5\#\\
\end{tabular}\par\vspace{3mm}
\qrcode[height=12mm]{https://helpman.komtera.lt/?fen=8/8/7k/6R1/8/1p6/8/K7&moves=6}\\{\sffamily\tiny\color{ink}scannen zum Lösen}
\end{minipage}
\vspace{2mm}

{\raggedright\noindent\textcolor{ink}{\sffamily\textsc{Themen}}\enspace \themetag{Satzspiel} \themetag{reines Matt} \themetag{Mustermatt} \themetag{Idealmatt} \themetag{Umwandlung} \themetag{Unterverwandlung} \themetag{Selbstblock} \themetag{ohne Schlag} \themetag{ohne Schach} \themetag{Umnov}\par}

\clearpage
\phantomsection\label{hm:KQvkq}
\begin{tcolorbox}[enhanced,colback=ink,colframe=ink,arc=0pt,boxrule=0pt,left=4mm,right=4mm,top=2.5mm,bottom=2.5mm]
\color{white}\sffamily{\LARGE\bfseries KQvkq}\hfill{\Large h\#5.5}\\[0.5mm]
{\small 4 Steine\quad Weiß am Zug}\hfill{\small Nr.~32\qquad\href{https://helpman.komtera.lt/?fen=8/8/5k2/8/6Q1/q7/8/K7&moves=5.5}{\color{white}\bfseries lösen\,\ensuremath{\nearrow}}}
\end{tcolorbox}
\vspace{2mm}
\noindent\begin{minipage}[t]{36mm}\vspace{0pt}\sffamily\small\raggedright
\hmfigure{längste eindeutige}{{\Large h\#5.5}}\hmfigure{eindeutig bei dieser Länge}{{\Large 7} Stellungen}\hmfigure{längstes Matt}{{\Large h\#7}}\hmfigure{in der Klasse}{3,78 Millionen Stellungen,\\ 2,22 Millionen mit Hilfsmatt}
\end{minipage}\hfill
\begin{minipage}[t]{80mm}\vspace{0pt}\begin{tikzpicture}\begin{axis}[width=72mm,height=42mm,ybar,bar width=1.6pt,xmin=-1,xmax=15,ymin=0,ymax=481.7,axis y line*=left,axis x line*=bottom,xlabel={\scriptsize Halbzüge bis Matt},ylabel={\scriptsize Stellungen (Tsd.)},tick label style={font=\tiny},label style={font=\scriptsize},enlarge x limits=false,ylabel shift=-3pt,xlabel shift=-3pt]\addplot[fill=ink!55,draw=none] coordinates {(0,0.801) (1,7.223) (2,25.43) (3,97.04) (4,195.2) (5,341.8) (6,421.3) (7,437.9) (8,365.2) (9,211.7) (10,89.4) (11,22.3) (12,8.965) (13,0.358) (14,0.052)};\draw[densely dotted] (axis cs:11,0) -- (axis cs:11,481.7);\end{axis}\begin{axis}[width=72mm,height=42mm,axis y line*=right,axis x line=none,ymode=log,xmin=-1,xmax=15,ylabel={\scriptsize eindeutig (log)},tick label style={font=\tiny},label style={font=\scriptsize},ylabel shift=-3pt]\addplot[mark=*,mark size=0.7pt,draw=black,thin] coordinates {(0,801) (1,5995) (2,11934) (3,22622) (4,21248) (5,18143) (6,11093) (7,7242) (8,2529) (9,561) (10,51) (11,7)};\end{axis}\end{tikzpicture}\end{minipage}
\vspace{3mm}

\noindent\begin{minipage}[t]{60mm}\vspace{0pt}\sbox\hmboard{\chessboard[setfen=8/8/5k2/8/6Q1/q7/8/K7 w - - 0 1]}\usebox\hmboard\\[1mm]
\makebox[\wd\hmboard]{\sffamily\hspace{3mm}\textbf{h\#5.5}\hfill(2+2)}\end{minipage}\hfill
\begin{minipage}[t]{52mm}\vspace{0pt}
\textcolor{ink}{\sffamily\textsc{Lösung}}\par\smallskip\small
\begin{tabular}{@{}r@{\enspace}l@{\enspace}l@{}}
\textbf{1.} & \ldots & \symking{}b1\\
\textbf{2.} & \symqueen{}d3+ & \symking{}c1\\
\textbf{3.} & \symqueen{}g6 & \symking{}d2\\
\textbf{4.} & \symking{}g7 & \symking{}e3\\
\textbf{5.} & \symking{}h6 & \symking{}f4\\
\textbf{6.} & \symqueen{}h7 & \symqueen{}g5\#\\
\end{tabular}\par\vspace{3mm}
\qrcode[height=12mm]{https://helpman.komtera.lt/?fen=8/8/5k2/8/6Q1/q7/8/K7&moves=5.5}\\{\sffamily\tiny\color{ink}scannen zum Lösen}
\end{minipage}
\vspace{2mm}

{\raggedright\noindent\textcolor{ink}{\sffamily\textsc{Themen}}\enspace \themetag{reines Matt} \themetag{Mustermatt} \themetag{Idealmatt} \themetag{Selbstblock} \themetag{ohne Schlag}\par\smallskip{\footnotesize\textit{Hinweis:} die Partei am Zug steht im Diagramm im Schach.}\par}

\begin{tcolorbox}[enhanced,colback=white,colframe=ink!35,arc=1.5pt,boxrule=0.4pt,left=3mm,right=3mm,top=1.5mm,bottom=1.5mm]
\textcolor{ink}{\sffamily\textsc{veröffentlicht}}\par\smallskip\small \textbf{Sheglow, Wiktor S.}, h\#6 0.1...\quad{\color{ink}\scriptsize verwandte Stellung}\\{\scriptsize Suomen Tehtäväniekat, No. (257), 05-06/1998}
\end{tcolorbox}
\clearpage
\phantomsection\label{hm:KRvkq}
\begin{tcolorbox}[enhanced,colback=ink,colframe=ink,arc=0pt,boxrule=0pt,left=4mm,right=4mm,top=2.5mm,bottom=2.5mm]
\color{white}\sffamily{\LARGE\bfseries KRvkq}\hfill{\Large h\#5.5}\\[0.5mm]
{\small 4 Steine\quad Weiß am Zug}\hfill{\small Nr.~33\qquad\href{https://helpman.komtera.lt/?fen=8/8/5k2/8/8/6R1/5q2/K7&moves=5.5}{\color{white}\bfseries lösen\,\ensuremath{\nearrow}}}
\end{tcolorbox}
\vspace{2mm}
\noindent\begin{minipage}[t]{36mm}\vspace{0pt}\sffamily\small\raggedright
\hmfigure{längste eindeutige}{{\Large h\#5.5}}\hmfigure{eindeutig bei dieser Länge}{{\Large 6} Stellungen}\hmfigure{längstes Matt}{{\Large h\#7.5}}\hmfigure{in der Klasse}{3,78 Millionen Stellungen,\\ 2,46 Millionen mit Hilfsmatt}
\end{minipage}\hfill
\begin{minipage}[t]{80mm}\vspace{0pt}\begin{tikzpicture}\begin{axis}[width=72mm,height=42mm,ybar,bar width=1.6pt,xmin=-1,xmax=16,ymin=0,ymax=530.2,axis y line*=left,axis x line*=bottom,xlabel={\scriptsize Halbzüge bis Matt},ylabel={\scriptsize Stellungen (Tsd.)},tick label style={font=\tiny},label style={font=\scriptsize},enlarge x limits=false,ylabel shift=-3pt,xlabel shift=-3pt]\addplot[fill=ink!55,draw=none] coordinates {(0,0.306) (1,2.381) (2,10.9) (3,57.15) (4,104.6) (5,272.9) (6,329.5) (7,482) (8,431.1) (9,428.2) (10,216) (11,101.5) (12,25.71) (13,2.45) (14,0.213) (15,0.001)};\draw[densely dotted] (axis cs:11,0) -- (axis cs:11,530.2);\end{axis}\begin{axis}[width=72mm,height=42mm,axis y line*=right,axis x line=none,ymode=log,xmin=-1,xmax=16,ylabel={\scriptsize eindeutig (log)},tick label style={font=\tiny},label style={font=\scriptsize},ylabel shift=-3pt]\addplot[mark=*,mark size=0.7pt,draw=black,thin] coordinates {(0,306) (1,2381) (2,5207) (3,20529) (4,19018) (5,24613) (6,12134) (7,7025) (8,3485) (9,1203) (10,199) (11,6)};\end{axis}\end{tikzpicture}\end{minipage}
\vspace{3mm}

\noindent\begin{minipage}[t]{60mm}\vspace{0pt}\sbox\hmboard{\chessboard[setfen=8/8/5k2/8/8/6R1/5q2/K7 w - - 0 1]}\usebox\hmboard\\[1mm]
\makebox[\wd\hmboard]{\sffamily\hspace{3mm}\textbf{h\#5.5}\hfill(2+2)}\end{minipage}\hfill
\begin{minipage}[t]{52mm}\vspace{0pt}
\textcolor{ink}{\sffamily\textsc{Lösung}}\par\smallskip\small
\begin{tabular}{@{}r@{\enspace}l@{\enspace}l@{}}
\textbf{1.} & \ldots & \symking{}b1\\
\textbf{2.} & \symqueen{}f5+ & \symking{}c1\\
\textbf{3.} & \symqueen{}g6 & \symking{}d2\\
\textbf{4.} & \symking{}g7 & \symking{}e3\\
\textbf{5.} & \symking{}h6 & \symking{}f4\\
\textbf{6.} & \symking{}h5 & \symrook{}h3\#\\
\end{tabular}\par\vspace{3mm}
\qrcode[height=12mm]{https://helpman.komtera.lt/?fen=8/8/5k2/8/8/6R1/5q2/K7&moves=5.5}\\{\sffamily\tiny\color{ink}scannen zum Lösen}
\end{minipage}
\vspace{2mm}

{\raggedright\noindent\textcolor{ink}{\sffamily\textsc{Themen}}\enspace \themetag{reines Matt} \themetag{Mustermatt} \themetag{Idealmatt} \themetag{Selbstblock} \themetag{ohne Schlag}\par}

\clearpage
\phantomsection\label{hm:KRvkr}
\begin{tcolorbox}[enhanced,colback=ink,colframe=ink,arc=0pt,boxrule=0pt,left=4mm,right=4mm,top=2.5mm,bottom=2.5mm]
\color{white}\sffamily{\LARGE\bfseries KRvkr}\hfill{\Large h\#5.5}\\[0.5mm]
{\small 4 Steine\quad Weiß am Zug}\hfill{\small Nr.~34\qquad\href{https://helpman.komtera.lt/?fen=8/8/5k2/8/6R1/8/7r/K7&moves=5.5}{\color{white}\bfseries lösen\,\ensuremath{\nearrow}}}
\end{tcolorbox}
\vspace{2mm}
\noindent\begin{minipage}[t]{36mm}\vspace{0pt}\sffamily\small\raggedright
\hmfigure{längste eindeutige}{{\Large h\#5.5}}\hmfigure{eindeutig bei dieser Länge}{{\Large 9} Stellungen}\hmfigure{längstes Matt}{{\Large h\#7}}\hmfigure{in der Klasse}{3,78 Millionen Stellungen,\\ 2,69 Millionen mit Hilfsmatt}
\end{minipage}\hfill
\begin{minipage}[t]{80mm}\vspace{0pt}\begin{tikzpicture}\begin{axis}[width=72mm,height=42mm,ybar,bar width=1.6pt,xmin=-1,xmax=15,ymin=0,ymax=562.6,axis y line*=left,axis x line*=bottom,xlabel={\scriptsize Halbzüge bis Matt},ylabel={\scriptsize Stellungen (Tsd.)},tick label style={font=\tiny},label style={font=\scriptsize},enlarge x limits=false,ylabel shift=-3pt,xlabel shift=-3pt]\addplot[fill=ink!55,draw=none] coordinates {(0,0.603) (1,4.454) (2,15.65) (3,69.2) (4,138.5) (5,284.2) (6,406) (7,483.4) (8,511.4) (9,416.8) (10,245.5) (11,87.76) (12,29.09) (13,1.406) (14,0.227)};\draw[densely dotted] (axis cs:11,0) -- (axis cs:11,562.6);\end{axis}\begin{axis}[width=72mm,height=42mm,axis y line*=right,axis x line=none,ymode=log,xmin=-1,xmax=15,ylabel={\scriptsize eindeutig (log)},tick label style={font=\tiny},label style={font=\scriptsize},ylabel shift=-3pt]\addplot[mark=*,mark size=0.7pt,draw=black,thin] coordinates {(0,603) (1,4454) (2,7661) (3,23365) (4,29905) (5,23422) (6,13659) (7,6133) (8,3714) (9,991) (10,297) (11,9)};\end{axis}\end{tikzpicture}\end{minipage}
\vspace{3mm}

\noindent\begin{minipage}[t]{60mm}\vspace{0pt}\sbox\hmboard{\chessboard[setfen=8/8/5k2/8/6R1/8/7r/K7 w - - 0 1]}\usebox\hmboard\\[1mm]
\makebox[\wd\hmboard]{\sffamily\hspace{3mm}\textbf{h\#5.5}\hfill(2+2)}\end{minipage}\hfill
\begin{minipage}[t]{52mm}\vspace{0pt}
\textcolor{ink}{\sffamily\textsc{Lösung}}\par\smallskip\small
\begin{tabular}{@{}r@{\enspace}l@{\enspace}l@{}}
\textbf{1.} & \ldots & \symking{}b1\\
\textbf{2.} & \symrook{}h6 & \symking{}c2\\
\textbf{3.} & \symrook{}g6 & \symking{}d3\\
\textbf{4.} & \symking{}g7 & \symking{}e4\\
\textbf{5.} & \symking{}h6 & \symking{}f5\\
\textbf{6.} & \symrook{}g7 & \symrook{}h4\#\\
\end{tabular}\par\vspace{3mm}
\qrcode[height=12mm]{https://helpman.komtera.lt/?fen=8/8/5k2/8/6R1/8/7r/K7&moves=5.5}\\{\sffamily\tiny\color{ink}scannen zum Lösen}
\end{minipage}
\vspace{2mm}

{\raggedright\noindent\textcolor{ink}{\sffamily\textsc{Themen}}\enspace \themetag{reines Matt} \themetag{Mustermatt} \themetag{Idealmatt} \themetag{Selbstblock} \themetag{ohne Schlag} \themetag{ohne Schach}\par}

\clearpage
\phantomsection\label{hm:KQvkb}
\begin{tcolorbox}[enhanced,colback=ink,colframe=ink,arc=0pt,boxrule=0pt,left=4mm,right=4mm,top=2.5mm,bottom=2.5mm]
\color{white}\sffamily{\LARGE\bfseries KQvkb}\hfill{\Large h\#5}\\[0.5mm]
{\small 4 Steine\quad Schwarz am Zug}\hfill{\small Nr.~35\qquad\href{https://helpman.komtera.lt/?fen=8/8/7k/8/8/b5Q1/8/K7&moves=5}{\color{white}\bfseries lösen\,\ensuremath{\nearrow}}}
\end{tcolorbox}
\vspace{2mm}
\noindent\begin{minipage}[t]{36mm}\vspace{0pt}\sffamily\small\raggedright
\hmfigure{längste eindeutige}{{\Large h\#5}}\hmfigure{eindeutig bei dieser Länge}{{\Large 64} Stellungen}\hmfigure{längstes Matt}{{\Large h\#6}}\hmfigure{in der Klasse}{3,78 Millionen Stellungen,\\ 2,58 Millionen mit Hilfsmatt}
\end{minipage}\hfill
\begin{minipage}[t]{80mm}\vspace{0pt}\begin{tikzpicture}\begin{axis}[width=72mm,height=42mm,ybar,bar width=1.6pt,xmin=-1,xmax=13,ymin=0,ymax=633.8,axis y line*=left,axis x line*=bottom,xlabel={\scriptsize Halbzüge bis Matt},ylabel={\scriptsize Stellungen (Tsd.)},tick label style={font=\tiny},label style={font=\scriptsize},enlarge x limits=false,ylabel shift=-3pt,xlabel shift=-3pt]\addplot[fill=ink!55,draw=none] coordinates {(0,1.692) (1,12.74) (2,34.96) (3,112.7) (4,263.5) (5,364.5) (6,576.2) (7,437.6) (8,470.1) (9,174.2) (10,105) (11,17.51) (12,10.92)};\draw[densely dotted] (axis cs:10,0) -- (axis cs:10,633.8);\end{axis}\begin{axis}[width=72mm,height=42mm,axis y line*=right,axis x line=none,ymode=log,xmin=-1,xmax=13,ylabel={\scriptsize eindeutig (log)},tick label style={font=\tiny},label style={font=\scriptsize},ylabel shift=-3pt]\addplot[mark=*,mark size=0.7pt,draw=black,thin] coordinates {(0,1692) (1,9101) (2,15194) (3,24220) (4,31664) (5,20810) (6,16405) (7,7027) (8,3284) (9,299) (10,64)};\end{axis}\end{tikzpicture}\end{minipage}
\vspace{3mm}

\noindent\begin{minipage}[t]{60mm}\vspace{0pt}\sbox\hmboard{\chessboard[setfen=8/8/7k/8/8/b5Q1/8/K7 b - - 0 1]}\usebox\hmboard\\[1mm]
\makebox[\wd\hmboard]{\sffamily\hspace{3mm}\textbf{h\#5}\hfill(2+2)}\end{minipage}\hfill
\begin{minipage}[t]{52mm}\vspace{0pt}
\textcolor{ink}{\sffamily\textsc{Lösung}}\par\smallskip\small
\begin{tabular}{@{}r@{\enspace}l@{\enspace}l@{}}
\textbf{1.} & \symking{}h5 & \symqueen{}xa3\\
\textbf{2.} & \symking{}g4 & \symqueen{}a2\\
\textbf{3.} & \symking{}f3 & \symking{}b2\\
\textbf{4.} & \symking{}e2 & \symking{}c3+\\
\textbf{5.} & \symking{}d1 & \symqueen{}d2\#\\
\end{tabular}\par\vspace{3mm}
\qrcode[height=12mm]{https://helpman.komtera.lt/?fen=8/8/7k/8/8/b5Q1/8/K7&moves=5}\\{\sffamily\tiny\color{ink}scannen zum Lösen}
\end{minipage}
\vspace{2mm}

{\raggedright\noindent\textcolor{ink}{\sffamily\textsc{Themen}}\enspace \themetag{ein Stein zieht}\par}

\clearpage
\phantomsection\label{hm:KQvkr}
\begin{tcolorbox}[enhanced,colback=ink,colframe=ink,arc=0pt,boxrule=0pt,left=4mm,right=4mm,top=2.5mm,bottom=2.5mm]
\color{white}\sffamily{\LARGE\bfseries KQvkr}\hfill{\Large h\#5}\\[0.5mm]
{\small 4 Steine\quad Schwarz am Zug}\hfill{\small Nr.~36\qquad\href{https://helpman.komtera.lt/?fen=8/8/8/8/7k/1r4Q1/8/K7&moves=5}{\color{white}\bfseries lösen\,\ensuremath{\nearrow}}}
\end{tcolorbox}
\vspace{2mm}
\noindent\begin{minipage}[t]{36mm}\vspace{0pt}\sffamily\small\raggedright
\hmfigure{längste eindeutige}{{\Large h\#5}}\hmfigure{eindeutig bei dieser Länge}{{\Large 7} Stellungen}\hmfigure{längstes Matt}{{\Large h\#6}}\hmfigure{in der Klasse}{3,78 Millionen Stellungen,\\ 2,45 Millionen mit Hilfsmatt}
\end{minipage}\hfill
\begin{minipage}[t]{80mm}\vspace{0pt}\begin{tikzpicture}\begin{axis}[width=72mm,height=42mm,ybar,bar width=1.6pt,xmin=-1,xmax=13,ymin=0,ymax=657.5,axis y line*=left,axis x line*=bottom,xlabel={\scriptsize Halbzüge bis Matt},ylabel={\scriptsize Stellungen (Tsd.)},tick label style={font=\tiny},label style={font=\scriptsize},enlarge x limits=false,ylabel shift=-3pt,xlabel shift=-3pt]\addplot[fill=ink!55,draw=none] coordinates {(0,1.374) (1,11.29) (2,39.91) (3,132.2) (4,357.5) (5,435.8) (6,597.7) (7,416.3) (8,313.7) (9,122.2) (10,22.05) (11,1.178) (12,0.061)};\draw[densely dotted] (axis cs:10,0) -- (axis cs:10,657.5);\end{axis}\begin{axis}[width=72mm,height=42mm,axis y line*=right,axis x line=none,ymode=log,xmin=-1,xmax=13,ylabel={\scriptsize eindeutig (log)},tick label style={font=\tiny},label style={font=\scriptsize},ylabel shift=-3pt]\addplot[mark=*,mark size=0.7pt,draw=black,thin] coordinates {(0,1374) (1,8752) (2,21234) (3,24414) (4,26434) (5,15184) (6,8603) (7,4016) (8,1270) (9,109) (10,7)};\end{axis}\end{tikzpicture}\end{minipage}
\vspace{3mm}

\noindent\begin{minipage}[t]{60mm}\vspace{0pt}\sbox\hmboard{\chessboard[setfen=8/8/8/8/7k/1r4Q1/8/K7 b - - 0 1]}\usebox\hmboard\\[1mm]
\makebox[\wd\hmboard]{\sffamily\hspace{3mm}\textbf{h\#5}\hfill(2+2)}\end{minipage}\hfill
\begin{minipage}[t]{52mm}\vspace{0pt}
\textcolor{ink}{\sffamily\textsc{Lösung}}\par\smallskip\small
\begin{tabular}{@{}r@{\enspace}l@{\enspace}l@{}}
\textbf{1.} & \symking{}h5 & \symqueen{}xb3\\
\textbf{2.} & \symking{}g4 & \symqueen{}a2\\
\textbf{3.} & \symking{}f3 & \symking{}b2\\
\textbf{4.} & \symking{}e2 & \symking{}c3+\\
\textbf{5.} & \symking{}d1 & \symqueen{}d2\#\\
\end{tabular}\par\vspace{3mm}
\qrcode[height=12mm]{https://helpman.komtera.lt/?fen=8/8/8/8/7k/1r4Q1/8/K7&moves=5}\\{\sffamily\tiny\color{ink}scannen zum Lösen}
\end{minipage}
\vspace{2mm}

{\raggedright\noindent\textcolor{ink}{\sffamily\textsc{Themen}}\enspace \themetag{ein Stein zieht}\par\smallskip{\footnotesize\textit{Hinweis:} die Partei am Zug steht im Diagramm im Schach.}\par}

\clearpage
\phantomsection\label{hm:KRNvk}
\begin{tcolorbox}[enhanced,colback=ink,colframe=ink,arc=0pt,boxrule=0pt,left=4mm,right=4mm,top=2.5mm,bottom=2.5mm]
\color{white}\sffamily{\LARGE\bfseries KRNvk}\hfill{\Large h\#5}\\[0.5mm]
{\small 4 Steine\quad Schwarz am Zug}\hfill{\small Nr.~37\qquad\href{https://helpman.komtera.lt/?fen=8/8/8/8/6k1/N4R2/8/K7&moves=5}{\color{white}\bfseries lösen\,\ensuremath{\nearrow}}}
\end{tcolorbox}
\vspace{2mm}
\noindent\begin{minipage}[t]{36mm}\vspace{0pt}\sffamily\small\raggedright
\hmfigure{längste eindeutige}{{\Large h\#5}}\hmfigure{eindeutig bei dieser Länge}{{\Large 19} Stellungen}\hmfigure{längstes Matt}{{\Large h\#7}}\hmfigure{in der Klasse}{3,78 Millionen Stellungen,\\ 2,94 Millionen mit Hilfsmatt}
\end{minipage}\hfill
\begin{minipage}[t]{80mm}\vspace{0pt}\begin{tikzpicture}\begin{axis}[width=72mm,height=42mm,ybar,bar width=1.6pt,xmin=-1,xmax=15,ymin=0,ymax=761.9,axis y line*=left,axis x line*=bottom,xlabel={\scriptsize Halbzüge bis Matt},ylabel={\scriptsize Stellungen (Tsd.)},tick label style={font=\tiny},label style={font=\scriptsize},enlarge x limits=false,ylabel shift=-3pt,xlabel shift=-3pt]\addplot[fill=ink!55,draw=none] coordinates {(0,1.835) (1,11.46) (2,19.45) (3,81.6) (4,211.9) (5,384.3) (6,692.6) (7,550.7) (8,644.6) (9,203.3) (10,125.7) (11,7.058) (12,4.287) (13,0.005) (14,0.005)};\draw[densely dotted] (axis cs:10,0) -- (axis cs:10,761.9);\end{axis}\begin{axis}[width=72mm,height=42mm,axis y line*=right,axis x line=none,ymode=log,xmin=-1,xmax=15,ylabel={\scriptsize eindeutig (log)},tick label style={font=\tiny},label style={font=\scriptsize},ylabel shift=-3pt]\addplot[mark=*,mark size=0.7pt,draw=black,thin] coordinates {(0,1835) (1,11261) (2,18965) (3,47524) (4,90466) (5,53428) (6,39774) (7,6321) (8,1552) (9,26) (10,19)};\end{axis}\end{tikzpicture}\end{minipage}
\vspace{3mm}

\noindent\begin{minipage}[t]{60mm}\vspace{0pt}\sbox\hmboard{\chessboard[setfen=8/8/8/8/6k1/N4R2/8/K7 b - - 0 1]}\usebox\hmboard\\[1mm]
\makebox[\wd\hmboard]{\sffamily\hspace{3mm}\textbf{h\#5}\hfill(3+1)}\end{minipage}\hfill
\begin{minipage}[t]{52mm}\vspace{0pt}
\textcolor{ink}{\sffamily\textsc{Lösung}}\par\smallskip\small
\begin{tabular}{@{}r@{\enspace}l@{\enspace}l@{}}
\textbf{1.} & \symking{}g5 & \symrook{}b3\\
\textbf{2.} & \symking{}f4 & \symrook{}b2\\
\textbf{3.} & \symking{}e3 & \symknight{}c2+\\
\textbf{4.} & \symking{}d2 & \symknight{}e3+\\
\textbf{5.} & \symking{}c1 & \symrook{}c2\#\\
\end{tabular}\par\vspace{3mm}
\qrcode[height=12mm]{https://helpman.komtera.lt/?fen=8/8/8/8/6k1/N4R2/8/K7&moves=5}\\{\sffamily\tiny\color{ink}scannen zum Lösen}
\end{minipage}
\vspace{2mm}

{\raggedright\noindent\textcolor{ink}{\sffamily\textsc{Themen}}\enspace \themetag{Satzspiel} \themetag{ein Stein zieht} \themetag{ohne Schlag} \themetag{Umnov}\par}

\begin{tcolorbox}[enhanced,colback=white,colframe=ink!35,arc=1.5pt,boxrule=0.4pt,left=3mm,right=3mm,top=1.5mm,bottom=1.5mm]
\textcolor{ink}{\sffamily\textsc{veröffentlicht}}\par\smallskip\small \textbf{Maslar, Zdravko}, h\#5\quad{\color{ink}\scriptsize verwandte Stellung}\\{\scriptsize Stern, 15/01/1981; Catalogue h\# (CHM) avec 6 pieces, No. 784, 1995}
\end{tcolorbox}
\clearpage
\phantomsection\label{hm:KRvkb}
\begin{tcolorbox}[enhanced,colback=ink,colframe=ink,arc=0pt,boxrule=0pt,left=4mm,right=4mm,top=2.5mm,bottom=2.5mm]
\color{white}\sffamily{\LARGE\bfseries KRvkb}\hfill{\Large h\#5}\\[0.5mm]
{\small 4 Steine\quad Schwarz am Zug}\hfill{\small Nr.~38\qquad\href{https://helpman.komtera.lt/?fen=8/8/8/8/8/7k/6R1/K4b2&moves=5}{\color{white}\bfseries lösen\,\ensuremath{\nearrow}}}
\end{tcolorbox}
\vspace{2mm}
\noindent\begin{minipage}[t]{36mm}\vspace{0pt}\sffamily\small\raggedright
\hmfigure{längste eindeutige}{{\Large h\#5}}\hmfigure{eindeutig bei dieser Länge}{{\Large 365} Stellungen}\hmfigure{längstes Matt}{{\Large h\#7}}\hmfigure{in der Klasse}{3,78 Millionen Stellungen,\\ 2,83 Millionen mit Hilfsmatt}
\end{minipage}\hfill
\begin{minipage}[t]{80mm}\vspace{0pt}\begin{tikzpicture}\begin{axis}[width=72mm,height=42mm,ybar,bar width=1.6pt,xmin=-1,xmax=15,ymin=0,ymax=616.8,axis y line*=left,axis x line*=bottom,xlabel={\scriptsize Halbzüge bis Matt},ylabel={\scriptsize Stellungen (Tsd.)},tick label style={font=\tiny},label style={font=\scriptsize},enlarge x limits=false,ylabel shift=-3pt,xlabel shift=-3pt]\addplot[fill=ink!55,draw=none] coordinates {(0,0.816) (1,5.694) (2,16.8) (3,73.05) (4,155.7) (5,311.8) (6,456) (7,488.1) (8,560.7) (9,398.9) (10,254.8) (11,69.43) (12,33.04) (13,0.905) (14,0.356)};\draw[densely dotted] (axis cs:10,0) -- (axis cs:10,616.8);\end{axis}\begin{axis}[width=72mm,height=42mm,axis y line*=right,axis x line=none,ymode=log,xmin=-1,xmax=15,ylabel={\scriptsize eindeutig (log)},tick label style={font=\tiny},label style={font=\scriptsize},ylabel shift=-3pt]\addplot[mark=*,mark size=0.7pt,draw=black,thin] coordinates {(0,816) (1,5694) (2,7994) (3,24641) (4,39647) (5,26747) (6,14990) (7,4370) (8,3077) (9,805) (10,365)};\end{axis}\end{tikzpicture}\end{minipage}
\vspace{3mm}

\noindent\begin{minipage}[t]{60mm}\vspace{0pt}\sbox\hmboard{\chessboard[setfen=8/8/8/8/8/7k/6R1/K4b2 b - - 0 1]}\usebox\hmboard\\[1mm]
\makebox[\wd\hmboard]{\sffamily\hspace{3mm}\textbf{h\#5}\hfill(2+2)}\end{minipage}\hfill
\begin{minipage}[t]{52mm}\vspace{0pt}
\textcolor{ink}{\sffamily\textsc{Lösung}}\par\smallskip\small
\begin{tabular}{@{}r@{\enspace}l@{\enspace}l@{}}
\textbf{1.} & \symbishop{}d3 & \symrook{}a2\\
\textbf{2.} & \symking{}g3 & \symking{}b2\\
\textbf{3.} & \symking{}f2 & \symking{}c3+\\
\textbf{4.} & \symking{}e1 & \symking{}xd3\\
\textbf{5.} & \symking{}d1 & \symrook{}a1\#\\
\end{tabular}\par\vspace{3mm}
\qrcode[height=12mm]{https://helpman.komtera.lt/?fen=8/8/8/8/8/7k/6R1/K4b2&moves=5}\\{\sffamily\tiny\color{ink}scannen zum Lösen}
\end{minipage}
\vspace{2mm}

{\raggedright\noindent\textcolor{ink}{\sffamily\textsc{Themen}}\enspace \themetag{reines Matt} \themetag{Mustermatt} \themetag{Idealmatt} \themetag{Spiegelmatt} \themetag{Inder}\par}

\clearpage
\phantomsection\label{hm:KRvkn}
\begin{tcolorbox}[enhanced,colback=ink,colframe=ink,arc=0pt,boxrule=0pt,left=4mm,right=4mm,top=2.5mm,bottom=2.5mm]
\color{white}\sffamily{\LARGE\bfseries KRvkn}\hfill{\Large h\#5}\\[0.5mm]
{\small 4 Steine\quad Schwarz am Zug}\hfill{\small Nr.~39\qquad\href{https://helpman.komtera.lt/?fen=8/8/8/8/8/n6k/6R1/K7&moves=5}{\color{white}\bfseries lösen\,\ensuremath{\nearrow}}}
\end{tcolorbox}
\vspace{2mm}
\noindent\begin{minipage}[t]{36mm}\vspace{0pt}\sffamily\small\raggedright
\hmfigure{längste eindeutige}{{\Large h\#5}}\hmfigure{eindeutig bei dieser Länge}{{\Large 247} Stellungen}\hmfigure{längstes Matt}{{\Large h\#7}}\hmfigure{in der Klasse}{3,78 Millionen Stellungen,\\ 2,91 Millionen mit Hilfsmatt}
\end{minipage}\hfill
\begin{minipage}[t]{80mm}\vspace{0pt}\begin{tikzpicture}\begin{axis}[width=72mm,height=42mm,ybar,bar width=1.6pt,xmin=-1,xmax=15,ymin=0,ymax=641.9,axis y line*=left,axis x line*=bottom,xlabel={\scriptsize Halbzüge bis Matt},ylabel={\scriptsize Stellungen (Tsd.)},tick label style={font=\tiny},label style={font=\scriptsize},enlarge x limits=false,ylabel shift=-3pt,xlabel shift=-3pt]\addplot[fill=ink!55,draw=none] coordinates {(0,1.166) (1,8.028) (2,19.64) (3,77.1) (4,169.1) (5,310.2) (6,479.3) (7,482.8) (8,583.5) (9,395.6) (10,273.7) (11,72.8) (12,39.3) (13,1.427) (14,0.703)};\draw[densely dotted] (axis cs:10,0) -- (axis cs:10,641.9);\end{axis}\begin{axis}[width=72mm,height=42mm,axis y line*=right,axis x line=none,ymode=log,xmin=-1,xmax=15,ylabel={\scriptsize eindeutig (log)},tick label style={font=\tiny},label style={font=\scriptsize},ylabel shift=-3pt]\addplot[mark=*,mark size=0.7pt,draw=black,thin] coordinates {(0,1166) (1,8028) (2,10483) (3,30272) (4,51989) (5,28420) (6,16712) (7,3799) (8,2623) (9,567) (10,247)};\end{axis}\end{tikzpicture}\end{minipage}
\vspace{3mm}

\noindent\begin{minipage}[t]{60mm}\vspace{0pt}\sbox\hmboard{\chessboard[setfen=8/8/8/8/8/n6k/6R1/K7 b - - 0 1]}\usebox\hmboard\\[1mm]
\makebox[\wd\hmboard]{\sffamily\hspace{3mm}\textbf{h\#5}\hfill(2+2)}\end{minipage}\hfill
\begin{minipage}[t]{52mm}\vspace{0pt}
\textcolor{ink}{\sffamily\textsc{Lösung}}\par\smallskip\small
\begin{tabular}{@{}r@{\enspace}l@{\enspace}l@{}}
\textbf{1.} & \symknight{}b5 & \symrook{}a2\\
\textbf{2.} & \symking{}g3 & \symking{}b2\\
\textbf{3.} & \symking{}f2 & \symking{}c2\\
\textbf{4.} & \symking{}e1 & \symking{}d3\\
\textbf{5.} & \symking{}d1 & \symrook{}a1\#\\
\end{tabular}\par\vspace{3mm}
\qrcode[height=12mm]{https://helpman.komtera.lt/?fen=8/8/8/8/8/n6k/6R1/K7&moves=5}\\{\sffamily\tiny\color{ink}scannen zum Lösen}
\end{minipage}
\vspace{2mm}

{\raggedright\noindent\textcolor{ink}{\sffamily\textsc{Themen}}\enspace \themetag{reines Matt} \themetag{Mustermatt} \themetag{Spiegelmatt} \themetag{ohne Schlag} \themetag{ohne Schach}\par}

\clearpage
\phantomsection\label{hm:KRRvk}
\begin{tcolorbox}[enhanced,colback=ink,colframe=ink,arc=0pt,boxrule=0pt,left=4mm,right=4mm,top=2.5mm,bottom=2.5mm]
\color{white}\sffamily{\LARGE\bfseries KRRvk}\hfill{\Large h\#4}\\[0.5mm]
{\small 4 Steine\quad Schwarz am Zug}\hfill{\small Nr.~40\qquad\href{https://helpman.komtera.lt/?fen=8/8/8/8/8/8/1R6/K4kR1&moves=4}{\color{white}\bfseries lösen\,\ensuremath{\nearrow}}}
\end{tcolorbox}
\vspace{2mm}
\noindent\begin{minipage}[t]{36mm}\vspace{0pt}\sffamily\small\raggedright
\hmfigure{längste eindeutige}{{\Large h\#4}}\hmfigure{eindeutig bei dieser Länge}{{\Large 27} Stellungen}\hmfigure{längstes Matt}{{\Large h\#7}}\hmfigure{in der Klasse}{3,78 Millionen Stellungen,\\ 1,37 Millionen mit Hilfsmatt}
\end{minipage}\hfill
\begin{minipage}[t]{80mm}\vspace{0pt}\begin{tikzpicture}\begin{axis}[width=72mm,height=42mm,ybar,bar width=1.6pt,xmin=-1,xmax=15,ymin=0,ymax=657.9,axis y line*=left,axis x line*=bottom,xlabel={\scriptsize Halbzüge bis Matt},ylabel={\scriptsize Stellungen (Tsd.)},tick label style={font=\tiny},label style={font=\scriptsize},enlarge x limits=false,ylabel shift=-3pt,xlabel shift=-3pt]\addplot[fill=ink!55,draw=none] coordinates {(0,9.052) (1,41.61) (2,113) (3,310.1) (4,598) (5,164) (6,124) (7,4.964) (8,4.863) (9,0.001) (10,2.367) (12,0.475) (14,0.023)};\draw[densely dotted] (axis cs:8,0) -- (axis cs:8,657.9);\end{axis}\begin{axis}[width=72mm,height=42mm,axis y line*=right,axis x line=none,ymode=log,xmin=-1,xmax=15,ylabel={\scriptsize eindeutig (log)},tick label style={font=\tiny},label style={font=\scriptsize},ylabel shift=-3pt]\addplot[mark=*,mark size=0.7pt,draw=black,thin] coordinates {(0,9052) (1,35985) (2,69998) (3,30066) (4,30852) (5,506) (6,382) (8,27)};\end{axis}\end{tikzpicture}\end{minipage}
\vspace{3mm}

\noindent\begin{minipage}[t]{60mm}\vspace{0pt}\sbox\hmboard{\chessboard[setfen=8/8/8/8/8/8/1R6/K4kR1 b - - 0 1]}\usebox\hmboard\\[1mm]
\makebox[\wd\hmboard]{\sffamily\hspace{3mm}\textbf{h\#4}\hfill(3+1)}\end{minipage}\hfill
\begin{minipage}[t]{52mm}\vspace{0pt}
\textcolor{ink}{\sffamily\textsc{Lösung}}\par\smallskip\small
\begin{tabular}{@{}r@{\enspace}l@{\enspace}l@{}}
\textbf{1.} & \symking{}xg1 & \symking{}b1\\
\textbf{2.} & \symking{}f1 & \symking{}c2\\
\textbf{3.} & \symking{}e1 & \symking{}d3\\
\textbf{4.} & \symking{}d1 & \symrook{}b1\#\\
\end{tabular}\par\vspace{3mm}
\qrcode[height=12mm]{https://helpman.komtera.lt/?fen=8/8/8/8/8/8/1R6/K4kR1&moves=4}\\{\sffamily\tiny\color{ink}scannen zum Lösen}
\end{minipage}
\vspace{2mm}

{\raggedright\noindent\textcolor{ink}{\sffamily\textsc{Themen}}\enspace \themetag{reines Matt} \themetag{Mustermatt} \themetag{Idealmatt} \themetag{Spiegelmatt} \themetag{Switchback} \themetag{ein Stein zieht} \themetag{ohne Schach}\par\smallskip{\footnotesize\textit{Hinweis:} die Partei am Zug steht im Diagramm im Schach; die Lösung beginnt mit einem Schlag.}\par}

\clearpage

\chapter{Fünf Steine}
\markboth{Fünf Steine}{}

181 der 185 Materialklassen mit 5 Steinen, in denen ein Matt möglich ist, längste eindeutige Aufgabe zuerst.
\clearpage

\phantomsection\label{hm:KBvkrp}
\begin{tcolorbox}[enhanced,colback=ink,colframe=ink,arc=0pt,boxrule=0pt,left=4mm,right=4mm,top=2.5mm,bottom=2.5mm]
\color{white}\sffamily{\LARGE\bfseries KBvkrp}\hfill{\Large h\#13}\\[0.5mm]
{\small 5 Steine\quad Schwarz am Zug}\hfill{\small Nr.~41\qquad\href{https://helpman.komtera.lt/?fen=K7/8/3p4/8/8/2k5/8/5Br1&moves=13}{\color{white}\bfseries lösen\,\ensuremath{\nearrow}}}
\end{tcolorbox}
\vspace{2mm}
\noindent\begin{minipage}[t]{36mm}\vspace{0pt}\sffamily\small\raggedright
\hmfigure{längste eindeutige}{{\Large h\#13}}\hmfigure{eindeutig bei dieser Länge}{{\Large 2} Stellungen}\hmfigure{längstes Matt}{{\Large h\#16.5}}\hmfigure{in der Klasse}{710 Millionen Stellungen,\\ 506 Millionen mit Hilfsmatt}
\end{minipage}\hfill
\begin{minipage}[t]{80mm}\vspace{0pt}\begin{tikzpicture}\begin{axis}[width=72mm,height=42mm,ybar,bar width=1.6pt,xmin=-1,xmax=34,ymin=0,ymax=44.41,axis y line*=left,axis x line*=bottom,xlabel={\scriptsize Halbzüge bis Matt},ylabel={\scriptsize Stellungen (Mio.)},tick label style={font=\tiny},label style={font=\scriptsize},enlarge x limits=false,ylabel shift=-3pt,xlabel shift=-3pt]\addplot[fill=ink!55,draw=none] coordinates {(0,0.000504) (1,0.00216) (2,0.008019) (3,0.03317) (4,0.1257) (5,0.4084) (6,1.234) (7,2.923) (8,6.914) (9,12.25) (10,21.47) (11,30.45) (12,38.65) (13,40.37) (14,38.66) (15,39.51) (16,37.08) (17,39.67) (18,32.88) (19,34.8) (20,21.84) (21,23.38) (22,14.47) (23,15.93) (24,11.89) (25,13.06) (26,8.003) (27,8.875) (28,3.531) (29,3.991) (30,1.555) (31,1.821) (32,0.003993) (33,0.000316)};\draw[densely dotted] (axis cs:26,0) -- (axis cs:26,44.41);\end{axis}\begin{axis}[width=72mm,height=42mm,axis y line*=right,axis x line=none,ymode=log,xmin=-1,xmax=34,ylabel={\scriptsize eindeutig (log)},tick label style={font=\tiny},label style={font=\scriptsize},ylabel shift=-3pt]\addplot[mark=*,mark size=0.7pt,draw=black,thin] coordinates {(0,504) (1,2160) (2,4094) (3,12931) (4,48188) (5,71008) (6,113310) (7,139615) (8,202164) (9,141712) (10,155152) (11,102577) (12,78574) (13,31527) (14,26138) (15,10526) (16,5649) (17,2139) (18,1276) (19,923) (20,399) (21,466) (22,454) (23,569) (24,123) (25,22) (26,2)};\end{axis}\end{tikzpicture}\end{minipage}
\vspace{3mm}

\noindent\begin{minipage}[t]{60mm}\vspace{0pt}\sbox\hmboard{\chessboard[setfen=K7/8/3p4/8/8/2k5/8/5Br1 b - - 0 1]}\usebox\hmboard\\[1mm]
\makebox[\wd\hmboard]{\sffamily\hspace{3mm}\textbf{h\#13}\hfill(2+3)}\end{minipage}\hfill
\begin{minipage}[t]{52mm}\vspace{0pt}
\textcolor{ink}{\sffamily\textsc{Lösung}}\par\smallskip\small
\begin{tabular}{@{}r@{\enspace}l@{\enspace}l@{}}
\textbf{1.} & d5 & \symking{}b7\\
\textbf{2.} & d4 & \symking{}c6\\
\textbf{3.} & d3 & \symking{}d5\\
\textbf{4.} & d2 & \symking{}e4\\
\textbf{5.} & d1=\symknight{} & \symking{}f3\\
\textbf{6.} & \symknight{}e3 & \symking{}f2\\
\textbf{7.} & \symking{}d2 & \symking{}xg1\\
\textbf{8.} & \symking{}e1 & \symking{}h2\\
\textbf{9.} & \symking{}f2 & \symking{}h3\\
\textbf{10.} & \symking{}g1 & \symking{}g3\\
\textbf{11.} & \symking{}h1 & \symbishop{}h3\\
\textbf{12.} & \symknight{}f1+ & \symking{}f2\\
\textbf{13.} & \symknight{}h2 & \symbishop{}g2\#\\
\end{tabular}\par\vspace{3mm}
\qrcode[height=12mm]{https://helpman.komtera.lt/?fen=K7/8/3p4/8/8/2k5/8/5Br1&moves=13}\\{\sffamily\tiny\color{ink}scannen zum Lösen}
\end{minipage}
\vspace{2mm}

{\raggedright\noindent\textcolor{ink}{\sffamily\textsc{Themen}}\enspace \themetag{reines Matt} \themetag{Mustermatt} \themetag{Idealmatt} \themetag{Umwandlung} \themetag{Unterverwandlung} \themetag{Rundlauf} \themetag{Selbstblock} \themetag{Umnov}\par}

\clearpage
\phantomsection\label{hm:KBvkqp}
\begin{tcolorbox}[enhanced,colback=ink,colframe=ink,arc=0pt,boxrule=0pt,left=4mm,right=4mm,top=2.5mm,bottom=2.5mm]
\color{white}\sffamily{\LARGE\bfseries KBvkqp}\hfill{\Large h\#12}\\[0.5mm]
{\small 5 Steine\quad Schwarz am Zug}\hfill{\small Nr.~42\qquad\href{https://helpman.komtera.lt/?fen=K7/1B6/8/4q3/8/1p6/k7/8&moves=12}{\color{white}\bfseries lösen\,\ensuremath{\nearrow}}}
\end{tcolorbox}
\vspace{2mm}
\noindent\begin{minipage}[t]{36mm}\vspace{0pt}\sffamily\small\raggedright
\hmfigure{längste eindeutige}{{\Large h\#12}}\hmfigure{eindeutig bei dieser Länge}{{\Large 3} Stellungen}\hmfigure{längstes Matt}{{\Large h\#17}}\hmfigure{in der Klasse}{710 Millionen Stellungen,\\ 466 Millionen mit Hilfsmatt}
\end{minipage}\hfill
\begin{minipage}[t]{80mm}\vspace{0pt}\begin{tikzpicture}\begin{axis}[width=72mm,height=42mm,ybar,bar width=1.6pt,xmin=-1,xmax=35,ymin=0,ymax=43.34,axis y line*=left,axis x line*=bottom,xlabel={\scriptsize Halbzüge bis Matt},ylabel={\scriptsize Stellungen (Mio.)},tick label style={font=\tiny},label style={font=\scriptsize},enlarge x limits=false,ylabel shift=-3pt,xlabel shift=-3pt]\addplot[fill=ink!55,draw=none] coordinates {(0,0.000332) (1,0.001436) (2,0.007107) (3,0.03179) (4,0.1104) (5,0.4142) (6,1.097) (7,3.098) (8,6.153) (9,12.96) (10,18.77) (11,31.24) (12,32.48) (13,39.4) (14,31.31) (15,38.32) (16,30.37) (17,39.18) (18,26.74) (19,33.98) (20,17.3) (21,22.05) (22,12.08) (23,16.03) (24,10.26) (25,13.57) (26,7.266) (27,9.622) (28,3.634) (29,4.897) (30,1.545) (31,2.142) (32,0.1797) (33,0.233) (34,2.8e-05)};\draw[densely dotted] (axis cs:24,0) -- (axis cs:24,43.34);\end{axis}\begin{axis}[width=72mm,height=42mm,axis y line*=right,axis x line=none,ymode=log,xmin=-1,xmax=35,ylabel={\scriptsize eindeutig (log)},tick label style={font=\tiny},label style={font=\scriptsize},ylabel shift=-3pt]\addplot[mark=*,mark size=0.7pt,draw=black,thin] coordinates {(0,332) (1,1436) (2,3585) (3,13640) (4,28786) (5,65937) (6,84800) (7,139532) (8,147886) (9,164340) (10,127901) (11,116646) (12,67142) (13,38779) (14,21453) (15,11537) (16,4838) (17,1585) (18,685) (19,187) (20,100) (21,36) (22,18) (23,10) (24,3)};\end{axis}\end{tikzpicture}\end{minipage}
\vspace{3mm}

\noindent\begin{minipage}[t]{60mm}\vspace{0pt}\sbox\hmboard{\chessboard[setfen=K7/1B6/8/4q3/8/1p6/k7/8 b - - 0 1]}\usebox\hmboard\\[1mm]
\makebox[\wd\hmboard]{\sffamily\hspace{3mm}\textbf{h\#12}\hfill(2+3)}\end{minipage}\hfill
\begin{minipage}[t]{52mm}\vspace{0pt}
\textcolor{ink}{\sffamily\textsc{Lösung}}\par\smallskip\small
\begin{tabular}{@{}r@{\enspace}l@{\enspace}l@{}}
\textbf{1.} & b2 & \symking{}a7\\
\textbf{2.} & \symking{}b3 & \symking{}b6\\
\textbf{3.} & \symking{}c4 & \symbishop{}a6+\\
\textbf{4.} & \symking{}d5 & \symking{}b5\\
\textbf{5.} & \symking{}e4+ & \symking{}c4\\
\textbf{6.} & \symking{}f3 & \symking{}d3\\
\textbf{7.} & \symking{}g2 & \symking{}d2\\
\textbf{8.} & \symking{}h1 & \symbishop{}e2\\
\textbf{9.} & b1=\symknight{}+ & \symking{}e1\\
\textbf{10.} & \symknight{}d2 & \symking{}f2\\
\textbf{11.} & \symknight{}f3 & \symbishop{}f1\\
\textbf{12.} & \symknight{}h2 & \symbishop{}g2\#\\
\end{tabular}\par\vspace{3mm}
\qrcode[height=12mm]{https://helpman.komtera.lt/?fen=K7/1B6/8/4q3/8/1p6/k7/8&moves=12}\\{\sffamily\tiny\color{ink}scannen zum Lösen}
\end{minipage}
\vspace{2mm}

{\raggedright\noindent\textcolor{ink}{\sffamily\textsc{Themen}}\enspace \themetag{reines Matt} \themetag{Mustermatt} \themetag{Umwandlung} \themetag{Unterverwandlung} \themetag{Selbstblock} \themetag{ohne Schlag} \themetag{Umnov}\par}

\clearpage
\phantomsection\label{hm:KBvkbp}
\begin{tcolorbox}[enhanced,colback=ink,colframe=ink,arc=0pt,boxrule=0pt,left=4mm,right=4mm,top=2.5mm,bottom=2.5mm]
\color{white}\sffamily{\LARGE\bfseries KBvkbp}\hfill{\Large h\#11.5}\\[0.5mm]
{\small 5 Steine\quad Weiß am Zug}\hfill{\small Nr.~43\qquad\href{https://helpman.komtera.lt/?fen=8/8/2k2p2/8/3K4/8/8/1B5b&moves=11.5}{\color{white}\bfseries lösen\,\ensuremath{\nearrow}}}
\end{tcolorbox}
\vspace{2mm}
\noindent\begin{minipage}[t]{36mm}\vspace{0pt}\sffamily\small\raggedright
\hmfigure{längste eindeutige}{{\Large h\#11.5}}\hmfigure{eindeutig bei dieser Länge}{{\Large 8} Stellungen}\hmfigure{längstes Matt}{{\Large h\#16}}\hmfigure{in der Klasse}{710 Millionen Stellungen,\\ 528 Millionen mit Hilfsmatt}
\end{minipage}\hfill
\begin{minipage}[t]{80mm}\vspace{0pt}\begin{tikzpicture}\begin{axis}[width=72mm,height=42mm,ybar,bar width=1.6pt,xmin=-1,xmax=33,ymin=0,ymax=63.32,axis y line*=left,axis x line*=bottom,xlabel={\scriptsize Halbzüge bis Matt},ylabel={\scriptsize Stellungen (Mio.)},tick label style={font=\tiny},label style={font=\scriptsize},enlarge x limits=false,ylabel shift=-3pt,xlabel shift=-3pt]\addplot[fill=ink!55,draw=none] coordinates {(0,0.003861) (1,0.0148) (2,0.08468) (3,0.269) (4,1.164) (5,2.724) (6,6.372) (7,11.73) (8,20.52) (9,31.12) (10,45.06) (11,56.45) (12,57.57) (13,42.28) (14,33.35) (15,27.59) (16,23.81) (17,22.75) (18,20.36) (19,19.93) (20,18.21) (21,18.47) (22,14.95) (23,15.15) (24,9.772) (25,10.05) (26,5.651) (27,5.861) (28,2.255) (29,2.413) (30,0.7678) (31,0.8242) (32,0.001736)};\draw[densely dotted] (axis cs:23,0) -- (axis cs:23,63.32);\end{axis}\begin{axis}[width=72mm,height=42mm,axis y line*=right,axis x line=none,ymode=log,xmin=-1,xmax=33,ylabel={\scriptsize eindeutig (log)},tick label style={font=\tiny},label style={font=\scriptsize},ylabel shift=-3pt]\addplot[mark=*,mark size=0.7pt,draw=black,thin] coordinates {(0,3861) (1,14732) (2,77008) (3,101342) (4,248731) (5,376084) (6,310942) (7,273673) (8,209087) (9,206788) (10,146149) (11,93998) (12,51356) (13,14480) (14,11691) (15,2784) (16,1834) (17,330) (18,242) (19,25) (20,41) (21,5) (22,9) (23,8)};\end{axis}\end{tikzpicture}\end{minipage}
\vspace{3mm}

\noindent\begin{minipage}[t]{60mm}\vspace{0pt}\sbox\hmboard{\chessboard[setfen=8/8/2k2p2/8/3K4/8/8/1B5b w - - 0 1]}\usebox\hmboard\\[1mm]
\makebox[\wd\hmboard]{\sffamily\hspace{3mm}\textbf{h\#11.5}\hfill(2+3)}\end{minipage}\hfill
\begin{minipage}[t]{52mm}\vspace{0pt}
\textcolor{ink}{\sffamily\textsc{Lösung}}\par\smallskip\small
\begin{tabular}{@{}r@{\enspace}l@{\enspace}l@{}}
\textbf{1.} & \ldots & \symking{}e3\\
\textbf{2.} & f5 & \symking{}f2\\
\textbf{3.} & f4 & \symking{}g1\\
\textbf{4.} & f3 & \symking{}xh1\\
\textbf{5.} & f2 & \symking{}g2\\
\textbf{6.} & f1=\symknight{} & \symking{}f2\\
\textbf{7.} & \symking{}d5 & \symbishop{}a2+\\
\textbf{8.} & \symking{}e4 & \symking{}e1\\
\textbf{9.} & \symking{}f3 & \symbishop{}e6\\
\textbf{10.} & \symking{}g2 & \symbishop{}h3+\\
\textbf{11.} & \symking{}h1 & \symking{}f2\\
\textbf{12.} & \symknight{}h2 & \symbishop{}g2\#\\
\end{tabular}\par\vspace{3mm}
\qrcode[height=12mm]{https://helpman.komtera.lt/?fen=8/8/2k2p2/8/3K4/8/8/1B5b&moves=11.5}\\{\sffamily\tiny\color{ink}scannen zum Lösen}
\end{minipage}
\vspace{2mm}

{\raggedright\noindent\textcolor{ink}{\sffamily\textsc{Themen}}\enspace \themetag{reines Matt} \themetag{Mustermatt} \themetag{Idealmatt} \themetag{Umwandlung} \themetag{Unterverwandlung} \themetag{Switchback} \themetag{Rundlauf} \themetag{Selbstblock} \themetag{Kniest} \themetag{Umnov}\par}

\clearpage
\phantomsection\label{hm:KBPvkp}
\begin{tcolorbox}[enhanced,colback=ink,colframe=ink,arc=0pt,boxrule=0pt,left=4mm,right=4mm,top=2.5mm,bottom=2.5mm]
\color{white}\sffamily{\LARGE\bfseries KBPvkp}\hfill{\Large h\#10.5}\\[0.5mm]
{\small 5 Steine\quad Weiß am Zug}\hfill{\small Nr.~44\qquad\href{https://helpman.komtera.lt/?fen=5B2/8/8/7p/7P/1k6/8/1K6&moves=10.5}{\color{white}\bfseries lösen\,\ensuremath{\nearrow}}}
\end{tcolorbox}
\vspace{2mm}
\noindent\begin{minipage}[t]{36mm}\vspace{0pt}\sffamily\small\raggedright
\hmfigure{längste eindeutige}{{\Large h\#10.5}}\hmfigure{eindeutig bei dieser Länge}{{\Large 1} Stellung}\hmfigure{längstes Matt}{{\Large h\#16}}\hmfigure{in der Klasse}{533 Millionen Stellungen,\\ 417 Millionen mit Hilfsmatt}
\end{minipage}\hfill
\begin{minipage}[t]{80mm}\vspace{0pt}\begin{tikzpicture}\begin{axis}[width=72mm,height=42mm,ybar,bar width=1.6pt,xmin=-1,xmax=33,ymin=0,ymax=59.42,axis y line*=left,axis x line*=bottom,xlabel={\scriptsize Halbzüge bis Matt},ylabel={\scriptsize Stellungen (Mio.)},tick label style={font=\tiny},label style={font=\scriptsize},enlarge x limits=false,ylabel shift=-3pt,xlabel shift=-3pt]\addplot[fill=ink!55,draw=none] coordinates {(0,0.006703) (1,0.1514) (2,0.3783) (3,3.111) (4,8.469) (5,31.38) (6,42.79) (7,41.34) (8,47.3) (9,42.97) (10,50.75) (11,49.86) (12,54.02) (13,20.39) (14,14.76) (15,3.452) (16,3.157) (17,0.9723) (18,0.7628) (19,0.332) (20,0.3732) (21,0.2002) (22,0.2013) (23,0.07897) (24,0.07647) (25,0.01593) (26,0.01643) (27,0.0102) (28,0.01086) (29,0.000333) (30,0.000364) (32,2.7e-05)};\draw[densely dotted] (axis cs:21,0) -- (axis cs:21,59.42);\end{axis}\begin{axis}[width=72mm,height=42mm,axis y line*=right,axis x line=none,ymode=log,xmin=-1,xmax=33,ylabel={\scriptsize eindeutig (log)},tick label style={font=\tiny},label style={font=\scriptsize},ylabel shift=-3pt]\addplot[mark=*,mark size=0.7pt,draw=black,thin] coordinates {(0,6703) (1,55837) (2,154301) (3,1393163) (4,2003993) (5,2462222) (6,1495997) (7,1117681) (8,814576) (9,618629) (10,457613) (11,341913) (12,199140) (13,15398) (14,14440) (15,4763) (16,1487) (17,300) (18,93) (19,38) (20,7) (21,1)};\end{axis}\end{tikzpicture}\end{minipage}
\vspace{3mm}

\noindent\begin{minipage}[t]{60mm}\vspace{0pt}\sbox\hmboard{\chessboard[setfen=5B2/8/8/7p/7P/1k6/8/1K6 w - - 0 1]}\usebox\hmboard\\[1mm]
\makebox[\wd\hmboard]{\sffamily\hspace{3mm}\textbf{h\#10.5}\hfill(3+2)}\end{minipage}\hfill
\begin{minipage}[t]{52mm}\vspace{0pt}
\textcolor{ink}{\sffamily\textsc{Lösung}}\par\smallskip\small
\begin{tabular}{@{}r@{\enspace}l@{\enspace}l@{}}
\textbf{1.} & \ldots & \symking{}c1\\
\textbf{2.} & \symking{}c4 & \symking{}d2\\
\textbf{3.} & \symking{}d5 & \symking{}e3\\
\textbf{4.} & \symking{}e6 & \symking{}f4\\
\textbf{5.} & \symking{}f7 & \symking{}g5\\
\textbf{6.} & \symking{}g8 & \symking{}xh5\\
\textbf{7.} & \symking{}h8 & \symking{}g6\\
\textbf{8.} & \symking{}g8 & h5\\
\textbf{9.} & \symking{}h8 & h6\\
\textbf{10.} & \symking{}g8 & h7+\\
\textbf{11.} & \symking{}h8 & \symbishop{}g7\#\\
\end{tabular}\par\vspace{3mm}
\qrcode[height=12mm]{https://helpman.komtera.lt/?fen=5B2/8/8/7p/7P/1k6/8/1K6&moves=10.5}\\{\sffamily\tiny\color{ink}scannen zum Lösen}
\end{minipage}
\vspace{2mm}

{\raggedright\noindent\textcolor{ink}{\sffamily\textsc{Themen}}\enspace \themetag{reines Matt} \themetag{Mustermatt} \themetag{Idealmatt} \themetag{Switchback} \themetag{ein Stein zieht}\par}

\begin{tcolorbox}[enhanced,colback=white,colframe=ink!35,arc=1.5pt,boxrule=0.4pt,left=3mm,right=3mm,top=1.5mm,bottom=1.5mm]
\textcolor{ink}{\sffamily\textsc{veröffentlicht}}\par\smallskip\small \textbf{Szentai, Endre}, h\#11 0.1...\quad{\color{ink}\scriptsize verwandte Stellung}\\{\scriptsize Magyar Sakkelet, 1958; Suomen Tehtäväniekat, No. 320, 18/10/1999; Ideal-Mate Encyclopedia Vol.1, No. 1304, 1999}
\end{tcolorbox}
\clearpage
\phantomsection\label{hm:KPvkpp}
\begin{tcolorbox}[enhanced,colback=ink,colframe=ink,arc=0pt,boxrule=0pt,left=4mm,right=4mm,top=2.5mm,bottom=2.5mm]
\color{white}\sffamily{\LARGE\bfseries KPvkpp}\hfill{\Large h\#10}\\[0.5mm]
{\small 5 Steine\quad Schwarz am Zug}\hfill{\small Nr.~45\qquad\href{https://helpman.komtera.lt/?fen=6k1/8/8/8/1p2p3/8/4P3/K7&moves=10}{\color{white}\bfseries lösen\,\ensuremath{\nearrow}}}
\end{tcolorbox}
\vspace{2mm}
\noindent\begin{minipage}[t]{36mm}\vspace{0pt}\sffamily\small\raggedright
\hmfigure{längste eindeutige}{{\Large h\#10}}\hmfigure{eindeutig bei dieser Länge}{{\Large 2} Stellungen}\hmfigure{längstes Matt}{{\Large h\#15.5}}\hmfigure{in der Klasse}{399 Millionen Stellungen,\\ 163 Millionen mit Hilfsmatt}
\end{minipage}\hfill
\begin{minipage}[t]{80mm}\vspace{0pt}\begin{tikzpicture}\begin{axis}[width=72mm,height=42mm,ybar,bar width=1.6pt,xmin=-1,xmax=32,ymin=0,ymax=23.68,axis y line*=left,axis x line*=bottom,xlabel={\scriptsize Halbzüge bis Matt},ylabel={\scriptsize Stellungen (Mio.)},tick label style={font=\tiny},label style={font=\scriptsize},enlarge x limits=false,ylabel shift=-3pt,xlabel shift=-3pt]\addplot[fill=ink!55,draw=none] coordinates {(0,9.5e-05) (1,0.04302) (2,0.0925) (3,0.5879) (4,1.347) (5,5.028) (6,8.303) (7,14.5) (8,16.87) (9,20.22) (10,18.62) (11,21.53) (12,20.27) (13,14.31) (14,10.09) (15,2.934) (16,2.181) (17,1.751) (18,1.451) (19,1.164) (20,0.7904) (21,0.4959) (22,0.3111) (23,0.1645) (24,0.09321) (25,0.04261) (26,0.02642) (27,0.01682) (28,0.0155) (29,0.007124) (30,0.007104) (31,0.0001)};\draw[densely dotted] (axis cs:20,0) -- (axis cs:20,23.68);\end{axis}\begin{axis}[width=72mm,height=42mm,axis y line*=right,axis x line=none,ymode=log,xmin=-1,xmax=32,ylabel={\scriptsize eindeutig (log)},tick label style={font=\tiny},label style={font=\scriptsize},ylabel shift=-3pt]\addplot[mark=*,mark size=0.7pt,draw=black,thin] coordinates {(0,95) (1,8930) (2,19149) (3,135301) (4,233528) (5,444930) (6,330470) (7,247655) (8,155569) (9,159822) (10,108994) (11,123710) (12,73496) (13,8016) (14,2801) (15,1150) (16,611) (17,262) (18,205) (19,25) (20,2)};\end{axis}\end{tikzpicture}\end{minipage}
\vspace{3mm}

\noindent\begin{minipage}[t]{60mm}\vspace{0pt}\sbox\hmboard{\chessboard[setfen=6k1/8/8/8/1p2p3/8/4P3/K7 b - - 0 1]}\usebox\hmboard\\[1mm]
\makebox[\wd\hmboard]{\sffamily\hspace{3mm}\textbf{h\#10}\hfill(2+3)}\end{minipage}\hfill
\begin{minipage}[t]{52mm}\vspace{0pt}
\textcolor{ink}{\sffamily\textsc{Lösung}}\par\smallskip\small
\begin{tabular}{@{}r@{\enspace}l@{\enspace}l@{}}
\textbf{1.} & b3 & \symking{}b1\\
\textbf{2.} & b2 & \symking{}c2\\
\textbf{3.} & b1=\symbishop{}+ & \symking{}c1\\
\textbf{4.} & \symbishop{}d3 & exd3\\
\textbf{5.} & \symking{}f7 & dxe4\\
\textbf{6.} & \symking{}e6 & e5\\
\textbf{7.} & \symking{}d5 & e6\\
\textbf{8.} & \symking{}c4 & e7\\
\textbf{9.} & \symking{}b3 & e8=\symqueen{}\\
\textbf{10.} & \symking{}a2 & \symqueen{}a4\#\\
\end{tabular}\par\vspace{3mm}
\qrcode[height=12mm]{https://helpman.komtera.lt/?fen=6k1/8/8/8/1p2p3/8/4P3/K7&moves=10}\\{\sffamily\tiny\color{ink}scannen zum Lösen}
\end{minipage}
\vspace{2mm}

{\raggedright\noindent\textcolor{ink}{\sffamily\textsc{Themen}}\enspace \themetag{reines Matt} \themetag{Mustermatt} \themetag{Idealmatt} \themetag{Spiegelmatt} \themetag{Umwandlung} \themetag{Unterverwandlung} \themetag{Excelsior} \themetag{Umnov}\par}

\begin{tcolorbox}[enhanced,colback=white,colframe=ink!35,arc=1.5pt,boxrule=0.4pt,left=3mm,right=3mm,top=1.5mm,bottom=1.5mm]
\textcolor{ink}{\sffamily\textsc{veröffentlicht}}\par\smallskip\small \textbf{Ban, Jenö}, h\#10\quad{\color{ink}\scriptsize verwandte Stellung}\\{\scriptsize Stella Polaris, No. 1072, 06/1967; Ideal-Mate Encyclopedia Vol.1, No. 70, 1999}
\end{tcolorbox}
\clearpage
\phantomsection\label{hm:KBvknp}
\begin{tcolorbox}[enhanced,colback=ink,colframe=ink,arc=0pt,boxrule=0pt,left=4mm,right=4mm,top=2.5mm,bottom=2.5mm]
\color{white}\sffamily{\LARGE\bfseries KBvknp}\hfill{\Large h\#9.5}\\[0.5mm]
{\small 5 Steine\quad Weiß am Zug}\hfill{\small Nr.~46\qquad\href{https://helpman.komtera.lt/?fen=8/k1B2p2/8/5n2/8/8/K7/8&moves=9.5}{\color{white}\bfseries lösen\,\ensuremath{\nearrow}}}
\end{tcolorbox}
\vspace{2mm}
\noindent\begin{minipage}[t]{36mm}\vspace{0pt}\sffamily\small\raggedright
\hmfigure{längste eindeutige}{{\Large h\#9.5}}\hmfigure{eindeutig bei dieser Länge}{{\Large 2} Stellungen}\hmfigure{längstes Matt}{{\Large h\#15.5}}\hmfigure{in der Klasse}{710 Millionen Stellungen,\\ 542 Millionen mit Hilfsmatt}
\end{minipage}\hfill
\begin{minipage}[t]{80mm}\vspace{0pt}\begin{tikzpicture}\begin{axis}[width=72mm,height=42mm,ybar,bar width=1.6pt,xmin=-1,xmax=32,ymin=0,ymax=91.82,axis y line*=left,axis x line*=bottom,xlabel={\scriptsize Halbzüge bis Matt},ylabel={\scriptsize Stellungen (Mio.)},tick label style={font=\tiny},label style={font=\scriptsize},enlarge x limits=false,ylabel shift=-3pt,xlabel shift=-3pt]\addplot[fill=ink!55,draw=none] coordinates {(0,0.001248) (1,0.004426) (2,0.01069) (3,0.04269) (4,0.1434) (5,0.4782) (6,1.366) (7,3.391) (8,7.63) (9,13.78) (10,24.79) (11,36.83) (12,58.15) (13,65.65) (14,83.47) (15,73.68) (16,66.23) (17,43.9) (18,25.05) (19,22.26) (20,7.228) (21,7.197) (22,0.407) (23,0.4286) (24,0.000237) (25,9.7e-05) (27,5.1e-05) (29,1.4e-05) (31,9e-06)};\draw[densely dotted] (axis cs:19,0) -- (axis cs:19,91.82);\end{axis}\begin{axis}[width=72mm,height=42mm,axis y line*=right,axis x line=none,ymode=log,xmin=-1,xmax=32,ylabel={\scriptsize eindeutig (log)},tick label style={font=\tiny},label style={font=\scriptsize},ylabel shift=-3pt]\addplot[mark=*,mark size=0.7pt,draw=black,thin] coordinates {(0,1248) (1,4426) (2,5955) (3,17673) (4,55600) (5,102779) (6,137799) (7,183985) (8,217330) (9,236217) (10,262904) (11,185662) (12,159924) (13,57680) (14,21681) (15,8187) (16,3073) (17,211) (18,35) (19,2)};\end{axis}\end{tikzpicture}\end{minipage}
\vspace{3mm}

\noindent\begin{minipage}[t]{60mm}\vspace{0pt}\sbox\hmboard{\chessboard[setfen=8/k1B2p2/8/5n2/8/8/K7/8 w - - 0 1]}\usebox\hmboard\\[1mm]
\makebox[\wd\hmboard]{\sffamily\hspace{3mm}\textbf{h\#9.5}\hfill(2+3)}\end{minipage}\hfill
\begin{minipage}[t]{52mm}\vspace{0pt}
\textcolor{ink}{\sffamily\textsc{Lösung}}\par\smallskip\small
\begin{tabular}{@{}r@{\enspace}l@{\enspace}l@{}}
\textbf{1.} & \ldots & \symking{}b3\\
\textbf{2.} & \symking{}b7 & \symking{}c4\\
\textbf{3.} & \symking{}c8 & \symking{}d5\\
\textbf{4.} & \symking{}d7 & \symking{}e5\\
\textbf{5.} & \symking{}e8 & \symking{}f6\\
\textbf{6.} & \symking{}f8 & \symbishop{}d6+\\
\textbf{7.} & \symking{}g8 & \symbishop{}f8\\
\textbf{8.} & \symking{}h8 & \symking{}xf7\\
\textbf{9.} & \symknight{}h6+ & \symking{}g6\\
\textbf{10.} & \symknight{}g8 & \symbishop{}g7\#\\
\end{tabular}\par\vspace{3mm}
\qrcode[height=12mm]{https://helpman.komtera.lt/?fen=8/k1B2p2/8/5n2/8/8/K7/8&moves=9.5}\\{\sffamily\tiny\color{ink}scannen zum Lösen}
\end{minipage}
\vspace{2mm}

{\raggedright\noindent\textcolor{ink}{\sffamily\textsc{Themen}}\enspace \themetag{reines Matt} \themetag{Mustermatt} \themetag{Idealmatt} \themetag{Selbstblock} \themetag{Umnov}\par}

\clearpage
\phantomsection\label{hm:KBvkqn}
\begin{tcolorbox}[enhanced,colback=ink,colframe=ink,arc=0pt,boxrule=0pt,left=4mm,right=4mm,top=2.5mm,bottom=2.5mm]
\color{white}\sffamily{\LARGE\bfseries KBvkqn}\hfill{\Large h\#9.5}\\[0.5mm]
{\small 5 Steine\quad Weiß am Zug}\hfill{\small Nr.~47\qquad\href{https://helpman.komtera.lt/?fen=8/8/5q2/8/8/1k6/8/1K1n1B2&moves=9.5}{\color{white}\bfseries lösen\,\ensuremath{\nearrow}}}
\end{tcolorbox}
\vspace{2mm}
\noindent\begin{minipage}[t]{36mm}\vspace{0pt}\sffamily\small\raggedright
\hmfigure{längste eindeutige}{{\Large h\#9.5}}\hmfigure{eindeutig bei dieser Länge}{{\Large 46} Stellungen}\hmfigure{längstes Matt}{{\Large h\#13}}\hmfigure{in der Klasse}{242 Millionen Stellungen,\\ 152 Millionen mit Hilfsmatt}
\end{minipage}\hfill
\begin{minipage}[t]{80mm}\vspace{0pt}\begin{tikzpicture}\begin{axis}[width=72mm,height=42mm,ybar,bar width=1.6pt,xmin=-1,xmax=27,ymin=0,ymax=27.61,axis y line*=left,axis x line*=bottom,xlabel={\scriptsize Halbzüge bis Matt},ylabel={\scriptsize Stellungen (Mio.)},tick label style={font=\tiny},label style={font=\scriptsize},enlarge x limits=false,ylabel shift=-3pt,xlabel shift=-3pt]\addplot[fill=ink!55,draw=none] coordinates {(0,2.8e-05) (1,6.4e-05) (2,0.000207) (3,0.001828) (4,0.008009) (5,0.04499) (6,0.1264) (7,0.4656) (8,0.9202) (9,2.484) (10,3.578) (11,8.033) (12,9.867) (13,18.35) (14,19.09) (15,25.1) (16,17.05) (17,18.85) (18,8.589) (19,10.76) (20,3.104) (21,4.143) (22,0.4095) (23,0.5614) (24,0.003201) (25,0.003833) (26,3e-06)};\draw[densely dotted] (axis cs:19,0) -- (axis cs:19,27.61);\end{axis}\begin{axis}[width=72mm,height=42mm,axis y line*=right,axis x line=none,ymode=log,xmin=-1,xmax=27,ylabel={\scriptsize eindeutig (log)},tick label style={font=\tiny},label style={font=\scriptsize},ylabel shift=-3pt]\addplot[mark=*,mark size=0.7pt,draw=black,thin] coordinates {(0,28) (1,64) (2,130) (3,974) (4,4543) (5,14381) (6,13692) (7,27140) (8,24962) (9,42806) (10,40279) (11,46227) (12,33051) (13,34460) (14,13111) (15,8268) (16,1665) (17,942) (18,163) (19,46)};\end{axis}\end{tikzpicture}\end{minipage}
\vspace{3mm}

\noindent\begin{minipage}[t]{60mm}\vspace{0pt}\sbox\hmboard{\chessboard[setfen=8/8/5q2/8/8/1k6/8/1K1n1B2 w - - 0 1]}\usebox\hmboard\\[1mm]
\makebox[\wd\hmboard]{\sffamily\hspace{3mm}\textbf{h\#9.5}\hfill(2+3)}\end{minipage}\hfill
\begin{minipage}[t]{52mm}\vspace{0pt}
\textcolor{ink}{\sffamily\textsc{Lösung}}\par\smallskip\small
\begin{tabular}{@{}r@{\enspace}l@{\enspace}l@{}}
\textbf{1.} & \ldots & \symking{}c1\\
\textbf{2.} & \symknight{}f2 & \symking{}d2\\
\textbf{3.} & \symknight{}h3 & \symking{}e3\\
\textbf{4.} & \symking{}c2 & \symbishop{}d3+\\
\textbf{5.} & \symking{}d1 & \symbishop{}f5\\
\textbf{6.} & \symking{}e1 & \symking{}f3\\
\textbf{7.} & \symking{}f1 & \symking{}g3\\
\textbf{8.} & \symking{}g1 & \symbishop{}d3\\
\textbf{9.} & \symking{}h1 & \symbishop{}f1\\
\textbf{10.} & \symknight{}g1 & \symbishop{}g2\#\\
\end{tabular}\par\vspace{3mm}
\qrcode[height=12mm]{https://helpman.komtera.lt/?fen=8/8/5q2/8/8/1k6/8/1K1n1B2&moves=9.5}\\{\sffamily\tiny\color{ink}scannen zum Lösen}
\end{minipage}
\vspace{2mm}

{\raggedright\noindent\textcolor{ink}{\sffamily\textsc{Themen}}\enspace \themetag{reines Matt} \themetag{Mustermatt} \themetag{Switchback} \themetag{Selbstblock} \themetag{ohne Schlag}\par}

\clearpage
\phantomsection\label{hm:KBvkrn}
\begin{tcolorbox}[enhanced,colback=ink,colframe=ink,arc=0pt,boxrule=0pt,left=4mm,right=4mm,top=2.5mm,bottom=2.5mm]
\color{white}\sffamily{\LARGE\bfseries KBvkrn}\hfill{\Large h\#9.5}\\[0.5mm]
{\small 5 Steine\quad Weiß am Zug}\hfill{\small Nr.~48\qquad\href{https://helpman.komtera.lt/?fen=8/8/5r2/8/8/1k6/4B3/1K1n4&moves=9.5}{\color{white}\bfseries lösen\,\ensuremath{\nearrow}}}
\end{tcolorbox}
\vspace{2mm}
\noindent\begin{minipage}[t]{36mm}\vspace{0pt}\sffamily\small\raggedright
\hmfigure{längste eindeutige}{{\Large h\#9.5}}\hmfigure{eindeutig bei dieser Länge}{{\Large 24} Stellungen}\hmfigure{längstes Matt}{{\Large h\#12.5}}\hmfigure{in der Klasse}{242 Millionen Stellungen,\\ 164 Millionen mit Hilfsmatt}
\end{minipage}\hfill
\begin{minipage}[t]{80mm}\vspace{0pt}\begin{tikzpicture}\begin{axis}[width=72mm,height=42mm,ybar,bar width=1.6pt,xmin=-1,xmax=26,ymin=0,ymax=28.53,axis y line*=left,axis x line*=bottom,xlabel={\scriptsize Halbzüge bis Matt},ylabel={\scriptsize Stellungen (Mio.)},tick label style={font=\tiny},label style={font=\scriptsize},enlarge x limits=false,ylabel shift=-3pt,xlabel shift=-3pt]\addplot[fill=ink!55,draw=none] coordinates {(0,4.1e-05) (1,8.9e-05) (2,0.000258) (3,0.002284) (4,0.01042) (5,0.05729) (6,0.1692) (7,0.5665) (8,1.237) (9,2.836) (10,4.661) (11,8.768) (12,12.36) (13,19.45) (14,23.2) (15,25.94) (16,20.56) (17,17.99) (18,9.456) (19,9.703) (20,3.012) (21,3.392) (22,0.2506) (23,0.2785) (24,0.001074) (25,0.000472)};\draw[densely dotted] (axis cs:19,0) -- (axis cs:19,28.53);\end{axis}\begin{axis}[width=72mm,height=42mm,axis y line*=right,axis x line=none,ymode=log,xmin=-1,xmax=26,ylabel={\scriptsize eindeutig (log)},tick label style={font=\tiny},label style={font=\scriptsize},ylabel shift=-3pt]\addplot[mark=*,mark size=0.7pt,draw=black,thin] coordinates {(0,41) (1,89) (2,164) (3,1235) (4,6464) (5,17886) (6,18809) (7,28239) (8,31753) (9,43457) (10,45903) (11,39876) (12,35185) (13,26038) (14,12275) (15,5681) (16,1629) (17,680) (18,121) (19,24)};\end{axis}\end{tikzpicture}\end{minipage}
\vspace{3mm}

\noindent\begin{minipage}[t]{60mm}\vspace{0pt}\sbox\hmboard{\chessboard[setfen=8/8/5r2/8/8/1k6/4B3/1K1n4 w - - 0 1]}\usebox\hmboard\\[1mm]
\makebox[\wd\hmboard]{\sffamily\hspace{3mm}\textbf{h\#9.5}\hfill(2+3)}\end{minipage}\hfill
\begin{minipage}[t]{52mm}\vspace{0pt}
\textcolor{ink}{\sffamily\textsc{Lösung}}\par\smallskip\small
\begin{tabular}{@{}r@{\enspace}l@{\enspace}l@{}}
\textbf{1.} & \ldots & \symking{}c1\\
\textbf{2.} & \symknight{}f2 & \symking{}d2\\
\textbf{3.} & \symknight{}h3 & \symking{}e3\\
\textbf{4.} & \symking{}c2 & \symbishop{}d3+\\
\textbf{5.} & \symking{}d1 & \symbishop{}f5\\
\textbf{6.} & \symking{}e1 & \symking{}f3\\
\textbf{7.} & \symking{}f1 & \symking{}g3\\
\textbf{8.} & \symking{}g1 & \symbishop{}d3\\
\textbf{9.} & \symking{}h1 & \symbishop{}f1\\
\textbf{10.} & \symknight{}g1 & \symbishop{}g2\#\\
\end{tabular}\par\vspace{3mm}
\qrcode[height=12mm]{https://helpman.komtera.lt/?fen=8/8/5r2/8/8/1k6/4B3/1K1n4&moves=9.5}\\{\sffamily\tiny\color{ink}scannen zum Lösen}
\end{minipage}
\vspace{2mm}

{\raggedright\noindent\textcolor{ink}{\sffamily\textsc{Themen}}\enspace \themetag{reines Matt} \themetag{Mustermatt} \themetag{Switchback} \themetag{Selbstblock} \themetag{ohne Schlag}\par}

\clearpage
\phantomsection\label{hm:KNvkpp}
\begin{tcolorbox}[enhanced,colback=ink,colframe=ink,arc=0pt,boxrule=0pt,left=4mm,right=4mm,top=2.5mm,bottom=2.5mm]
\color{white}\sffamily{\LARGE\bfseries KNvkpp}\hfill{\Large h\#9.5}\\[0.5mm]
{\small 5 Steine\quad Weiß am Zug}\hfill{\small Nr.~49\qquad\href{https://helpman.komtera.lt/?fen=8/6p1/1K2k1p1/8/8/8/8/2N5&moves=9.5}{\color{white}\bfseries lösen\,\ensuremath{\nearrow}}}
\end{tcolorbox}
\vspace{2mm}
\noindent\begin{minipage}[t]{36mm}\vspace{0pt}\sffamily\small\raggedright
\hmfigure{längste eindeutige}{{\Large h\#9.5}}\hmfigure{eindeutig bei dieser Länge}{{\Large 2} Stellungen}\hmfigure{längstes Matt}{{\Large h\#11.5}}\hmfigure{in der Klasse}{533 Millionen Stellungen,\\ 214 Millionen mit Hilfsmatt}
\end{minipage}\hfill
\begin{minipage}[t]{80mm}\vspace{0pt}\begin{tikzpicture}\begin{axis}[width=72mm,height=42mm,ybar,bar width=1.6pt,xmin=-1,xmax=24,ymin=0,ymax=35.53,axis y line*=left,axis x line*=bottom,xlabel={\scriptsize Halbzüge bis Matt},ylabel={\scriptsize Stellungen (Mio.)},tick label style={font=\tiny},label style={font=\scriptsize},enlarge x limits=false,ylabel shift=-3pt,xlabel shift=-3pt]\addplot[fill=ink!55,draw=none] coordinates {(0,0.000192) (1,0.000882) (2,0.001974) (3,0.01199) (4,0.0442) (5,0.1701) (6,0.5979) (7,1.571) (8,4.005) (9,7.365) (10,14.18) (11,20.12) (12,29.32) (13,31.73) (14,32.3) (15,26.42) (16,20.03) (17,12.98) (18,6.85) (19,4.299) (20,1.062) (21,0.6997) (22,0.007057) (23,3e-06)};\draw[densely dotted] (axis cs:19,0) -- (axis cs:19,35.53);\end{axis}\begin{axis}[width=72mm,height=42mm,axis y line*=right,axis x line=none,ymode=log,xmin=-1,xmax=24,ylabel={\scriptsize eindeutig (log)},tick label style={font=\tiny},label style={font=\scriptsize},ylabel shift=-3pt]\addplot[mark=*,mark size=0.7pt,draw=black,thin] coordinates {(0,192) (1,831) (2,1583) (3,6493) (4,22702) (5,29310) (6,58659) (7,54293) (8,78945) (9,73766) (10,97940) (11,62706) (12,49601) (13,9913) (14,5842) (15,546) (16,179) (17,18) (18,13) (19,2)};\end{axis}\end{tikzpicture}\end{minipage}
\vspace{3mm}

\noindent\begin{minipage}[t]{60mm}\vspace{0pt}\sbox\hmboard{\chessboard[setfen=8/6p1/1K2k1p1/8/8/8/8/2N5 w - - 0 1]}\usebox\hmboard\\[1mm]
\makebox[\wd\hmboard]{\sffamily\hspace{3mm}\textbf{h\#9.5}\hfill(2+3)}\end{minipage}\hfill
\begin{minipage}[t]{52mm}\vspace{0pt}
\textcolor{ink}{\sffamily\textsc{Lösung}}\par\smallskip\small
\begin{tabular}{@{}r@{\enspace}l@{\enspace}l@{}}
\textbf{1.} & \ldots & \symking{}c7\\
\textbf{2.} & g5 & \symking{}d8\\
\textbf{3.} & g4 & \symking{}e8\\
\textbf{4.} & g3 & \symking{}f8\\
\textbf{5.} & g2 & \symking{}xg7\\
\textbf{6.} & g1=\symrook{}+ & \symking{}h7\\
\textbf{7.} & \symking{}f7 & \symking{}h6\\
\textbf{8.} & \symking{}g8 & \symknight{}d3\\
\textbf{9.} & \symking{}h8 & \symknight{}e5\\
\textbf{10.} & \symrook{}g8 & \symknight{}f7\#\\
\end{tabular}\par\vspace{3mm}
\qrcode[height=12mm]{https://helpman.komtera.lt/?fen=8/6p1/1K2k1p1/8/8/8/8/2N5&moves=9.5}\\{\sffamily\tiny\color{ink}scannen zum Lösen}
\end{minipage}
\vspace{2mm}

{\raggedright\noindent\textcolor{ink}{\sffamily\textsc{Themen}}\enspace \themetag{reines Matt} \themetag{Mustermatt} \themetag{Idealmatt} \themetag{Umwandlung} \themetag{Unterverwandlung} \themetag{Selbstblock}\par}

\clearpage
\phantomsection\label{hm:KNvkqb}
\begin{tcolorbox}[enhanced,colback=ink,colframe=ink,arc=0pt,boxrule=0pt,left=4mm,right=4mm,top=2.5mm,bottom=2.5mm]
\color{white}\sffamily{\LARGE\bfseries KNvkqb}\hfill{\Large h\#9.5}\\[0.5mm]
{\small 5 Steine\quad Weiß am Zug}\hfill{\small Nr.~50\qquad\href{https://helpman.komtera.lt/?fen=7k/6N1/8/8/8/8/b2q4/1K6&moves=9.5}{\color{white}\bfseries lösen\,\ensuremath{\nearrow}}}
\end{tcolorbox}
\vspace{2mm}
\noindent\begin{minipage}[t]{36mm}\vspace{0pt}\sffamily\small\raggedright
\hmfigure{längste eindeutige}{{\Large h\#9.5}}\hmfigure{eindeutig bei dieser Länge}{{\Large 1} Stellung}\hmfigure{längstes Matt}{{\Large h\#11.5}}\hmfigure{in der Klasse}{242 Millionen Stellungen,\\ 152 Millionen mit Hilfsmatt}
\end{minipage}\hfill
\begin{minipage}[t]{80mm}\vspace{0pt}\begin{tikzpicture}\begin{axis}[width=72mm,height=42mm,ybar,bar width=1.6pt,xmin=-1,xmax=24,ymin=0,ymax=31,axis y line*=left,axis x line*=bottom,xlabel={\scriptsize Halbzüge bis Matt},ylabel={\scriptsize Stellungen (Mio.)},tick label style={font=\tiny},label style={font=\scriptsize},enlarge x limits=false,ylabel shift=-3pt,xlabel shift=-3pt]\addplot[fill=ink!55,draw=none] coordinates {(0,3e-05) (1,0.000153) (2,0.001195) (3,0.007457) (4,0.03429) (5,0.1614) (6,0.3282) (7,1.216) (8,1.646) (9,4.86) (10,5.694) (11,13.73) (12,13.38) (13,26.22) (14,18.98) (15,28.18) (16,14.55) (17,15.72) (18,4.188) (19,3.466) (20,0.03979) (21,0.01197) (22,5e-06) (23,4e-06)};\draw[densely dotted] (axis cs:19,0) -- (axis cs:19,31);\end{axis}\begin{axis}[width=72mm,height=42mm,axis y line*=right,axis x line=none,ymode=log,xmin=-1,xmax=24,ylabel={\scriptsize eindeutig (log)},tick label style={font=\tiny},label style={font=\scriptsize},ylabel shift=-3pt]\addplot[mark=*,mark size=0.7pt,draw=black,thin] coordinates {(0,30) (1,153) (2,986) (3,4674) (4,15746) (5,32521) (6,14775) (7,33464) (8,25005) (9,50216) (10,51958) (11,70402) (12,33767) (13,28689) (14,11315) (15,6300) (16,1901) (17,618) (18,4) (19,1)};\end{axis}\end{tikzpicture}\end{minipage}
\vspace{3mm}

\noindent\begin{minipage}[t]{60mm}\vspace{0pt}\sbox\hmboard{\chessboard[setfen=7k/6N1/8/8/8/8/b2q4/1K6 w - - 0 1]}\usebox\hmboard\\[1mm]
\makebox[\wd\hmboard]{\sffamily\hspace{3mm}\textbf{h\#9.5}\hfill(2+3)}\end{minipage}\hfill
\begin{minipage}[t]{52mm}\vspace{0pt}
\textcolor{ink}{\sffamily\textsc{Lösung}}\par\smallskip\small
\begin{tabular}{@{}r@{\enspace}l@{\enspace}l@{}}
\textbf{1.} & \ldots & \symking{}a1\\
\textbf{2.} & \symking{}g8 & \symknight{}f5\\
\textbf{3.} & \symking{}f7 & \symknight{}e3\\
\textbf{4.} & \symking{}e6 & \symknight{}c2\\
\textbf{5.} & \symking{}d5 & \symking{}b2\\
\textbf{6.} & \symking{}c4 & \symking{}a3\\
\textbf{7.} & \symking{}c3 & \symking{}a4\\
\textbf{8.} & \symking{}b2 & \symknight{}d4\\
\textbf{9.} & \symking{}a1 & \symking{}a3\\
\textbf{10.} & \symbishop{}b1 & \symknight{}b3\#\\
\end{tabular}\par\vspace{3mm}
\qrcode[height=12mm]{https://helpman.komtera.lt/?fen=7k/6N1/8/8/8/8/b2q4/1K6&moves=9.5}\\{\sffamily\tiny\color{ink}scannen zum Lösen}
\end{minipage}
\vspace{2mm}

{\raggedright\noindent\textcolor{ink}{\sffamily\textsc{Themen}}\enspace \themetag{reines Matt} \themetag{Mustermatt} \themetag{Switchback} \themetag{Selbstblock} \themetag{ohne Schlag} \themetag{ohne Schach}\par\smallskip{\footnotesize\textit{Hinweis:} die Partei am Zug steht im Diagramm im Schach.}\par}

\begin{tcolorbox}[enhanced,colback=white,colframe=ink!35,arc=1.5pt,boxrule=0.4pt,left=3mm,right=3mm,top=1.5mm,bottom=1.5mm]
\textcolor{ink}{\sffamily\textsc{veröffentlicht}}\par\smallskip\small \textbf{Sheglow, Wiktor S.}, h\#10 0.1...\quad{\color{ink}\scriptsize verwandte Stellung}\\{\scriptsize Suomen Tehtäväniekat, No. (291), 05-06/1998}
\end{tcolorbox}
\clearpage
\phantomsection\label{hm:KNvkqp}
\begin{tcolorbox}[enhanced,colback=ink,colframe=ink,arc=0pt,boxrule=0pt,left=4mm,right=4mm,top=2.5mm,bottom=2.5mm]
\color{white}\sffamily{\LARGE\bfseries KNvkqp}\hfill{\Large h\#9.5}\\[0.5mm]
{\small 5 Steine\quad Weiß am Zug}\hfill{\small Nr.~51\qquad\href{https://helpman.komtera.lt/?fen=8/8/4p3/4k3/8/8/2q5/K3N3&moves=9.5}{\color{white}\bfseries lösen\,\ensuremath{\nearrow}}}
\end{tcolorbox}
\vspace{2mm}
\noindent\begin{minipage}[t]{36mm}\vspace{0pt}\sffamily\small\raggedright
\hmfigure{längste eindeutige}{{\Large h\#9.5}}\hmfigure{eindeutig bei dieser Länge}{{\Large 194} Stellungen}\hmfigure{längstes Matt}{{\Large h\#11.5}}\hmfigure{in der Klasse}{710 Millionen Stellungen,\\ 481 Millionen mit Hilfsmatt}
\end{minipage}\hfill
\begin{minipage}[t]{80mm}\vspace{0pt}\begin{tikzpicture}\begin{axis}[width=72mm,height=42mm,ybar,bar width=1.6pt,xmin=-1,xmax=24,ymin=0,ymax=84.58,axis y line*=left,axis x line*=bottom,xlabel={\scriptsize Halbzüge bis Matt},ylabel={\scriptsize Stellungen (Mio.)},tick label style={font=\tiny},label style={font=\scriptsize},enlarge x limits=false,ylabel shift=-3pt,xlabel shift=-3pt]\addplot[fill=ink!55,draw=none] coordinates {(0,0.000138) (1,0.00067) (2,0.002844) (3,0.01833) (4,0.08879) (5,0.3915) (6,1.228) (7,3.938) (8,8.429) (9,20.41) (10,29.72) (11,54.87) (12,52.76) (13,76.89) (14,53.38) (15,64.57) (16,37.23) (17,41.99) (18,14.88) (19,16.76) (20,1.585) (21,1.849) (22,0.000268) (23,0.001315)};\draw[densely dotted] (axis cs:19,0) -- (axis cs:19,84.58);\end{axis}\begin{axis}[width=72mm,height=42mm,axis y line*=right,axis x line=none,ymode=log,xmin=-1,xmax=24,ylabel={\scriptsize eindeutig (log)},tick label style={font=\tiny},label style={font=\scriptsize},ylabel shift=-3pt]\addplot[mark=*,mark size=0.7pt,draw=black,thin] coordinates {(0,138) (1,653) (2,2033) (3,9630) (4,29077) (5,62017) (6,89188) (7,165925) (8,167420) (9,249045) (10,170429) (11,207873) (12,107984) (13,102401) (14,40027) (15,29563) (16,8055) (17,5306) (18,368) (19,194)};\end{axis}\end{tikzpicture}\end{minipage}
\vspace{3mm}

\noindent\begin{minipage}[t]{60mm}\vspace{0pt}\sbox\hmboard{\chessboard[setfen=8/8/4p3/4k3/8/8/2q5/K3N3 w - - 0 1]}\usebox\hmboard\\[1mm]
\makebox[\wd\hmboard]{\sffamily\hspace{3mm}\textbf{h\#9.5}\hfill(2+3)}\end{minipage}\hfill
\begin{minipage}[t]{52mm}\vspace{0pt}
\textcolor{ink}{\sffamily\textsc{Lösung}}\par\smallskip\small
\begin{tabular}{@{}r@{\enspace}l@{\enspace}l@{}}
\textbf{1.} & \ldots & \symknight{}f3+\\
\textbf{2.} & \symking{}f4 & \symknight{}d2\\
\textbf{3.} & \symqueen{}d1+ & \symking{}b2\\
\textbf{4.} & e5 & \symking{}c3\\
\textbf{5.} & e4 & \symking{}d4\\
\textbf{6.} & \symking{}g3 & \symking{}e5\\
\textbf{7.} & \symking{}f2 & \symking{}f4\\
\textbf{8.} & \symking{}e1 & \symking{}g3\\
\textbf{9.} & e3 & \symking{}g2\\
\textbf{10.} & e2 & \symknight{}f3\#\\
\end{tabular}\par\vspace{3mm}
\qrcode[height=12mm]{https://helpman.komtera.lt/?fen=8/8/4p3/4k3/8/8/2q5/K3N3&moves=9.5}\\{\sffamily\tiny\color{ink}scannen zum Lösen}
\end{minipage}
\vspace{2mm}

{\raggedright\noindent\textcolor{ink}{\sffamily\textsc{Themen}}\enspace \themetag{reines Matt} \themetag{Mustermatt} \themetag{Idealmatt} \themetag{Switchback} \themetag{Selbstblock} \themetag{ohne Schlag}\par}

\clearpage
\phantomsection\label{hm:KPvkbp}
\begin{tcolorbox}[enhanced,colback=ink,colframe=ink,arc=0pt,boxrule=0pt,left=4mm,right=4mm,top=2.5mm,bottom=2.5mm]
\color{white}\sffamily{\LARGE\bfseries KPvkbp}\hfill{\Large h\#9.5}\\[0.5mm]
{\small 5 Steine\quad Weiß am Zug}\hfill{\small Nr.~52\qquad\href{https://helpman.komtera.lt/?fen=8/8/6p1/8/6P1/k7/8/Kb6&moves=9.5}{\color{white}\bfseries lösen\,\ensuremath{\nearrow}}}
\end{tcolorbox}
\vspace{2mm}
\noindent\begin{minipage}[t]{36mm}\vspace{0pt}\sffamily\small\raggedright
\hmfigure{längste eindeutige}{{\Large h\#9.5}}\hmfigure{eindeutig bei dieser Länge}{{\Large 9} Stellungen}\hmfigure{längstes Matt}{{\Large h\#14.5}}\hmfigure{in der Klasse}{533 Millionen Stellungen,\\ 417 Millionen mit Hilfsmatt}
\end{minipage}\hfill
\begin{minipage}[t]{80mm}\vspace{0pt}\begin{tikzpicture}\begin{axis}[width=72mm,height=42mm,ybar,bar width=1.6pt,xmin=-1,xmax=30,ymin=0,ymax=70.94,axis y line*=left,axis x line*=bottom,xlabel={\scriptsize Halbzüge bis Matt},ylabel={\scriptsize Stellungen (Mio.)},tick label style={font=\tiny},label style={font=\scriptsize},enlarge x limits=false,ylabel shift=-3pt,xlabel shift=-3pt]\addplot[fill=ink!55,draw=none] coordinates {(0,0.000282) (1,0.06189) (2,0.1716) (3,1.282) (4,3.054) (5,12.55) (6,20.16) (7,39.28) (8,41.89) (9,56.82) (10,49.15) (11,64.49) (12,55.38) (13,42.16) (14,21.62) (15,4.564) (16,1.615) (17,0.8817) (18,0.4145) (19,0.4467) (20,0.3756) (21,0.3223) (22,0.2768) (23,0.135) (24,0.1138) (25,0.04263) (26,0.03478) (27,0.02711) (28,0.02407) (29,0.000782)};\draw[densely dotted] (axis cs:19,0) -- (axis cs:19,70.94);\end{axis}\begin{axis}[width=72mm,height=42mm,axis y line*=right,axis x line=none,ymode=log,xmin=-1,xmax=30,ylabel={\scriptsize eindeutig (log)},tick label style={font=\tiny},label style={font=\scriptsize},ylabel shift=-3pt]\addplot[mark=*,mark size=0.7pt,draw=black,thin] coordinates {(0,282) (1,14535) (2,31091) (3,298976) (4,544264) (5,1153728) (6,714092) (7,696395) (8,396766) (9,469854) (10,291109) (11,361573) (12,195013) (13,45051) (14,9327) (15,2109) (16,435) (17,151) (18,10) (19,9)};\end{axis}\end{tikzpicture}\end{minipage}
\vspace{3mm}

\noindent\begin{minipage}[t]{60mm}\vspace{0pt}\sbox\hmboard{\chessboard[setfen=8/8/6p1/8/6P1/k7/8/Kb6 w - - 0 1]}\usebox\hmboard\\[1mm]
\makebox[\wd\hmboard]{\sffamily\hspace{3mm}\textbf{h\#9.5}\hfill(2+3)}\end{minipage}\hfill
\begin{minipage}[t]{52mm}\vspace{0pt}
\textcolor{ink}{\sffamily\textsc{Lösung}}\par\smallskip\small
\begin{tabular}{@{}r@{\enspace}l@{\enspace}l@{}}
\textbf{1.} & \ldots & g5\\
\textbf{2.} & \symking{}b4 & \symking{}b2\\
\textbf{3.} & \symking{}c5 & \symking{}c3\\
\textbf{4.} & \symking{}d6 & \symking{}d4\\
\textbf{5.} & \symking{}e7 & \symking{}e5\\
\textbf{6.} & \symking{}f8 & \symking{}f6\\
\textbf{7.} & \symbishop{}a2 & \symking{}xg6\\
\textbf{8.} & \symking{}g8 & \symking{}h6\\
\textbf{9.} & \symking{}h8 & g6\\
\textbf{10.} & \symbishop{}g8 & g7\#\\
\end{tabular}\par\vspace{3mm}
\qrcode[height=12mm]{https://helpman.komtera.lt/?fen=8/8/6p1/8/6P1/k7/8/Kb6&moves=9.5}\\{\sffamily\tiny\color{ink}scannen zum Lösen}
\end{minipage}
\vspace{2mm}

{\raggedright\noindent\textcolor{ink}{\sffamily\textsc{Themen}}\enspace \themetag{reines Matt} \themetag{Mustermatt} \themetag{Idealmatt} \themetag{Selbstblock} \themetag{ohne Schach}\par}

\begin{tcolorbox}[enhanced,colback=white,colframe=ink!35,arc=1.5pt,boxrule=0.4pt,left=3mm,right=3mm,top=1.5mm,bottom=1.5mm]
\textcolor{ink}{\sffamily\textsc{veröffentlicht}}\par\smallskip\small \textbf{Locker, Miklos S.}, h\#10 0.1...\quad{\color{ink}\scriptsize diese Stellung}\\{\scriptsize Magyar Sakkelet, 02/1966}
\end{tcolorbox}
\clearpage
\phantomsection\label{hm:KBBvkn}
\begin{tcolorbox}[enhanced,colback=ink,colframe=ink,arc=0pt,boxrule=0pt,left=4mm,right=4mm,top=2.5mm,bottom=2.5mm]
\color{white}\sffamily{\LARGE\bfseries KBBvkn}\hfill{\Large h\#9}\\[0.5mm]
{\small 5 Steine\quad Schwarz am Zug}\hfill{\small Nr.~53\qquad\href{https://helpman.komtera.lt/?fen=8/8/4n3/8/3B4/8/1B4k1/3K4&moves=9}{\color{white}\bfseries lösen\,\ensuremath{\nearrow}}}
\end{tcolorbox}
\vspace{2mm}
\noindent\begin{minipage}[t]{36mm}\vspace{0pt}\sffamily\small\raggedright
\hmfigure{längste eindeutige}{{\Large h\#9}}\hmfigure{eindeutig bei dieser Länge}{{\Large 2} Stellungen}\hmfigure{längstes Matt}{{\Large h\#12}}\hmfigure{in der Klasse}{242 Millionen Stellungen,\\ 85,4 Millionen mit Hilfsmatt}
\end{minipage}\hfill
\begin{minipage}[t]{80mm}\vspace{0pt}\begin{tikzpicture}\begin{axis}[width=72mm,height=42mm,ybar,bar width=1.6pt,xmin=-1,xmax=25,ymin=0,ymax=10.21,axis y line*=left,axis x line*=bottom,xlabel={\scriptsize Halbzüge bis Matt},ylabel={\scriptsize Stellungen (Mio.)},tick label style={font=\tiny},label style={font=\scriptsize},enlarge x limits=false,ylabel shift=-3pt,xlabel shift=-3pt]\addplot[fill=ink!55,draw=none] coordinates {(0,0.007958) (1,0.03261) (2,0.1101) (3,0.6909) (4,2.101) (5,5.117) (6,9.28) (7,7.597) (8,9.203) (9,5.253) (10,4.619) (11,3.123) (12,4.53) (13,4.428) (14,7.029) (15,5.372) (16,6.315) (17,3.658) (18,2.723) (19,2.144) (20,0.9676) (21,0.8592) (22,0.1237) (23,0.128) (24,0.000516)};\draw[densely dotted] (axis cs:18,0) -- (axis cs:18,10.21);\end{axis}\begin{axis}[width=72mm,height=42mm,axis y line*=right,axis x line=none,ymode=log,xmin=-1,xmax=25,ylabel={\scriptsize eindeutig (log)},tick label style={font=\tiny},label style={font=\scriptsize},ylabel shift=-3pt]\addplot[mark=*,mark size=0.7pt,draw=black,thin] coordinates {(0,7958) (1,32283) (2,73812) (3,351827) (4,480403) (5,335590) (6,255806) (7,68297) (8,39237) (9,12869) (10,11720) (11,4849) (12,6557) (13,2188) (14,1144) (15,353) (16,207) (17,4) (18,2)};\end{axis}\end{tikzpicture}\end{minipage}
\vspace{3mm}

\noindent\begin{minipage}[t]{60mm}\vspace{0pt}\sbox\hmboard{\chessboard[setfen=8/8/4n3/8/3B4/8/1B4k1/3K4 b - - 0 1]}\usebox\hmboard\\[1mm]
\makebox[\wd\hmboard]{\sffamily\hspace{3mm}\textbf{h\#9}\hfill(3+2)}\end{minipage}\hfill
\begin{minipage}[t]{52mm}\vspace{0pt}
\textcolor{ink}{\sffamily\textsc{Lösung}}\par\smallskip\small
\begin{tabular}{@{}r@{\enspace}l@{\enspace}l@{}}
\textbf{1.} & \symking{}f3 & \symbishop{}c1\\
\textbf{2.} & \symking{}e4 & \symking{}e2\\
\textbf{3.} & \symking{}xd4 & \symking{}f3\\
\textbf{4.} & \symking{}e5 & \symking{}g4\\
\textbf{5.} & \symking{}f6 & \symking{}h5\\
\textbf{6.} & \symking{}g7 & \symbishop{}h6+\\
\textbf{7.} & \symking{}h8 & \symking{}g6\\
\textbf{8.} & \symknight{}f8+ & \symking{}f7\\
\textbf{9.} & \symknight{}h7 & \symbishop{}g7\#\\
\end{tabular}\par\vspace{3mm}
\qrcode[height=12mm]{https://helpman.komtera.lt/?fen=8/8/4n3/8/3B4/8/1B4k1/3K4&moves=9}\\{\sffamily\tiny\color{ink}scannen zum Lösen}
\end{minipage}
\vspace{2mm}

{\raggedright\noindent\textcolor{ink}{\sffamily\textsc{Themen}}\enspace \themetag{reines Matt} \themetag{Mustermatt} \themetag{Idealmatt} \themetag{Selbstblock}\par}

\clearpage
\phantomsection\label{hm:KBBvkp}
\begin{tcolorbox}[enhanced,colback=ink,colframe=ink,arc=0pt,boxrule=0pt,left=4mm,right=4mm,top=2.5mm,bottom=2.5mm]
\color{white}\sffamily{\LARGE\bfseries KBBvkp}\hfill{\Large h\#9}\\[0.5mm]
{\small 5 Steine\quad Schwarz am Zug}\hfill{\small Nr.~54\qquad\href{https://helpman.komtera.lt/?fen=3K4/1B6/2B5/8/1k6/8/5p2/8&moves=9}{\color{white}\bfseries lösen\,\ensuremath{\nearrow}}}
\end{tcolorbox}
\vspace{2mm}
\noindent\begin{minipage}[t]{36mm}\vspace{0pt}\sffamily\small\raggedright
\hmfigure{längste eindeutige}{{\Large h\#9}}\hmfigure{eindeutig bei dieser Länge}{{\Large 1} Stellung}\hmfigure{längstes Matt}{{\Large h\#16}}\hmfigure{in der Klasse}{710 Millionen Stellungen,\\ 265 Millionen mit Hilfsmatt}
\end{minipage}\hfill
\begin{minipage}[t]{80mm}\vspace{0pt}\begin{tikzpicture}\begin{axis}[width=72mm,height=42mm,ybar,bar width=1.6pt,xmin=-1,xmax=33,ymin=0,ymax=34.24,axis y line*=left,axis x line*=bottom,xlabel={\scriptsize Halbzüge bis Matt},ylabel={\scriptsize Stellungen (Mio.)},tick label style={font=\tiny},label style={font=\scriptsize},enlarge x limits=false,ylabel shift=-3pt,xlabel shift=-3pt]\addplot[fill=ink!55,draw=none] coordinates {(0,0.03603) (1,0.145) (2,0.4201) (3,2.47) (4,7.685) (5,16.79) (6,31.13) (7,22.75) (8,28.07) (9,13.91) (10,11.8) (11,6.668) (12,9.805) (13,7.386) (14,13.08) (15,9.609) (16,14.25) (17,10.15) (18,12.13) (19,9.029) (20,9.354) (21,7.364) (22,6.605) (23,5.125) (24,3.232) (25,2.669) (26,1.439) (27,1.161) (28,0.3848) (29,0.3825) (30,0.000745) (32,7.3e-05)};\draw[densely dotted] (axis cs:18,0) -- (axis cs:18,34.24);\end{axis}\begin{axis}[width=72mm,height=42mm,axis y line*=right,axis x line=none,ymode=log,xmin=-1,xmax=33,ylabel={\scriptsize eindeutig (log)},tick label style={font=\tiny},label style={font=\scriptsize},ylabel shift=-3pt]\addplot[mark=*,mark size=0.7pt,draw=black,thin] coordinates {(0,36030) (1,142202) (2,380487) (3,1719724) (4,2360286) (5,1198315) (6,957263) (7,214917) (8,115227) (9,42940) (10,42996) (11,20511) (12,28808) (13,6557) (14,8583) (15,465) (16,369) (17,33) (18,1)};\end{axis}\end{tikzpicture}\end{minipage}
\vspace{3mm}

\noindent\begin{minipage}[t]{60mm}\vspace{0pt}\sbox\hmboard{\chessboard[setfen=3K4/1B6/2B5/8/1k6/8/5p2/8 b - - 0 1]}\usebox\hmboard\\[1mm]
\makebox[\wd\hmboard]{\sffamily\hspace{3mm}\textbf{h\#9}\hfill(3+2)}\end{minipage}\hfill
\begin{minipage}[t]{52mm}\vspace{0pt}
\textcolor{ink}{\sffamily\textsc{Lösung}}\par\smallskip\small
\begin{tabular}{@{}r@{\enspace}l@{\enspace}l@{}}
\textbf{1.} & \symking{}c5 & \symbishop{}c8\\
\textbf{2.} & \symking{}xc6 & \symking{}e7\\
\textbf{3.} & \symking{}d5 & \symking{}f6\\
\textbf{4.} & \symking{}e4 & \symking{}g5\\
\textbf{5.} & \symking{}f3 & \symking{}h4\\
\textbf{6.} & \symking{}g2 & \symbishop{}h3+\\
\textbf{7.} & \symking{}h1 & \symking{}g3\\
\textbf{8.} & f1=\symknight{}+ & \symking{}f2\\
\textbf{9.} & \symknight{}h2 & \symbishop{}g2\#\\
\end{tabular}\par\vspace{3mm}
\qrcode[height=12mm]{https://helpman.komtera.lt/?fen=3K4/1B6/2B5/8/1k6/8/5p2/8&moves=9}\\{\sffamily\tiny\color{ink}scannen zum Lösen}
\end{minipage}
\vspace{2mm}

{\raggedright\noindent\textcolor{ink}{\sffamily\textsc{Themen}}\enspace \themetag{reines Matt} \themetag{Mustermatt} \themetag{Idealmatt} \themetag{Umwandlung} \themetag{Unterverwandlung} \themetag{Selbstblock} \themetag{Umnov}\par}

\clearpage
\phantomsection\label{hm:KBvkbn}
\begin{tcolorbox}[enhanced,colback=ink,colframe=ink,arc=0pt,boxrule=0pt,left=4mm,right=4mm,top=2.5mm,bottom=2.5mm]
\color{white}\sffamily{\LARGE\bfseries KBvkbn}\hfill{\Large h\#9}\\[0.5mm]
{\small 5 Steine\quad Schwarz am Zug}\hfill{\small Nr.~55\qquad\href{https://helpman.komtera.lt/?fen=8/8/1n6/3b4/8/8/8/KBk5&moves=9}{\color{white}\bfseries lösen\,\ensuremath{\nearrow}}}
\end{tcolorbox}
\vspace{2mm}
\noindent\begin{minipage}[t]{36mm}\vspace{0pt}\sffamily\small\raggedright
\hmfigure{längste eindeutige}{{\Large h\#9}}\hmfigure{eindeutig bei dieser Länge}{{\Large 57} Stellungen}\hmfigure{längstes Matt}{{\Large h\#12}}\hmfigure{in der Klasse}{242 Millionen Stellungen,\\ 171 Millionen mit Hilfsmatt}
\end{minipage}\hfill
\begin{minipage}[t]{80mm}\vspace{0pt}\begin{tikzpicture}\begin{axis}[width=72mm,height=42mm,ybar,bar width=1.6pt,xmin=-1,xmax=25,ymin=0,ymax=24.63,axis y line*=left,axis x line*=bottom,xlabel={\scriptsize Halbzüge bis Matt},ylabel={\scriptsize Stellungen (Mio.)},tick label style={font=\tiny},label style={font=\scriptsize},enlarge x limits=false,ylabel shift=-3pt,xlabel shift=-3pt]\addplot[fill=ink!55,draw=none] coordinates {(0,0.000581) (1,0.00212) (2,0.01518) (3,0.05681) (4,0.2448) (5,0.6426) (6,1.4) (7,2.928) (8,4.677) (9,8.342) (10,12.2) (11,18.39) (12,22.39) (13,21.88) (14,21.16) (15,18.81) (16,13.19) (17,11.28) (18,5.13) (19,5.163) (20,1.357) (21,1.462) (22,0.08162) (23,0.09614) (24,4.2e-05)};\draw[densely dotted] (axis cs:18,0) -- (axis cs:18,24.63);\end{axis}\begin{axis}[width=72mm,height=42mm,axis y line*=right,axis x line=none,ymode=log,xmin=-1,xmax=25,ylabel={\scriptsize eindeutig (log)},tick label style={font=\tiny},label style={font=\scriptsize},ylabel shift=-3pt]\addplot[mark=*,mark size=0.7pt,draw=black,thin] coordinates {(0,581) (1,2120) (2,12766) (3,23553) (4,58121) (5,98662) (6,74276) (7,76841) (8,63322) (9,78561) (10,65201) (11,45234) (12,25701) (13,12804) (14,5374) (15,2084) (16,673) (17,208) (18,57)};\end{axis}\end{tikzpicture}\end{minipage}
\vspace{3mm}

\noindent\begin{minipage}[t]{60mm}\vspace{0pt}\sbox\hmboard{\chessboard[setfen=8/8/1n6/3b4/8/8/8/KBk5 b - - 0 1]}\usebox\hmboard\\[1mm]
\makebox[\wd\hmboard]{\sffamily\hspace{3mm}\textbf{h\#9}\hfill(2+3)}\end{minipage}\hfill
\begin{minipage}[t]{52mm}\vspace{0pt}
\textcolor{ink}{\sffamily\textsc{Lösung}}\par\smallskip\small
\begin{tabular}{@{}r@{\enspace}l@{\enspace}l@{}}
\textbf{1.} & \symking{}d2 & \symbishop{}d3\\
\textbf{2.} & \symking{}c3 & \symbishop{}a6\\
\textbf{3.} & \symking{}b4 & \symking{}b2\\
\textbf{4.} & \symking{}a5 & \symking{}c3\\
\textbf{5.} & \symknight{}c8 & \symking{}d4\\
\textbf{6.} & \symking{}b6 & \symking{}xd5\\
\textbf{7.} & \symking{}a7 & \symking{}c6\\
\textbf{8.} & \symking{}a8 & \symking{}c7\\
\textbf{9.} & \symknight{}a7 & \symbishop{}b7\#\\
\end{tabular}\par\vspace{3mm}
\qrcode[height=12mm]{https://helpman.komtera.lt/?fen=8/8/1n6/3b4/8/8/8/KBk5&moves=9}\\{\sffamily\tiny\color{ink}scannen zum Lösen}
\end{minipage}
\vspace{2mm}

{\raggedright\noindent\textcolor{ink}{\sffamily\textsc{Themen}}\enspace \themetag{reines Matt} \themetag{Mustermatt} \themetag{Idealmatt} \themetag{Selbstblock} \themetag{ohne Schach}\par}

\clearpage
\phantomsection\label{hm:KBvknn}
\begin{tcolorbox}[enhanced,colback=ink,colframe=ink,arc=0pt,boxrule=0pt,left=4mm,right=4mm,top=2.5mm,bottom=2.5mm]
\color{white}\sffamily{\LARGE\bfseries KBvknn}\hfill{\Large h\#9}\\[0.5mm]
{\small 5 Steine\quad Schwarz am Zug}\hfill{\small Nr.~56\qquad\href{https://helpman.komtera.lt/?fen=8/8/8/8/8/8/5n2/kBK4n&moves=9}{\color{white}\bfseries lösen\,\ensuremath{\nearrow}}}
\end{tcolorbox}
\vspace{2mm}
\noindent\begin{minipage}[t]{36mm}\vspace{0pt}\sffamily\small\raggedright
\hmfigure{längste eindeutige}{{\Large h\#9}}\hmfigure{eindeutig bei dieser Länge}{{\Large 34} Stellungen}\hmfigure{längstes Matt}{{\Large h\#12.5}}\hmfigure{in der Klasse}{242 Millionen Stellungen,\\ 87,6 Millionen mit Hilfsmatt}
\end{minipage}\hfill
\begin{minipage}[t]{80mm}\vspace{0pt}\begin{tikzpicture}\begin{axis}[width=72mm,height=42mm,ybar,bar width=1.6pt,xmin=-1,xmax=26,ymin=0,ymax=16.24,axis y line*=left,axis x line*=bottom,xlabel={\scriptsize Halbzüge bis Matt},ylabel={\scriptsize Stellungen (Mio.)},tick label style={font=\tiny},label style={font=\scriptsize},enlarge x limits=false,ylabel shift=-3pt,xlabel shift=-3pt]\addplot[fill=ink!55,draw=none] coordinates {(0,5.1e-05) (1,0.000109) (2,0.000295) (3,0.002729) (4,0.01247) (5,0.06474) (6,0.1867) (7,0.5812) (8,1.196) (9,2.423) (10,3.983) (11,6.501) (12,9.635) (13,12.55) (14,14.76) (15,12.72) (16,9.892) (17,6.425) (18,2.964) (19,2.745) (20,0.4882) (21,0.5144) (22,0.002776) (23,0.002774) (24,2e-06) (25,1e-06)};\draw[densely dotted] (axis cs:18,0) -- (axis cs:18,16.24);\end{axis}\begin{axis}[width=72mm,height=42mm,axis y line*=right,axis x line=none,ymode=log,xmin=-1,xmax=26,ylabel={\scriptsize eindeutig (log)},tick label style={font=\tiny},label style={font=\scriptsize},ylabel shift=-3pt]\addplot[mark=*,mark size=0.7pt,draw=black,thin] coordinates {(0,51) (1,109) (2,185) (3,1488) (4,7390) (5,16194) (6,14471) (7,19044) (8,26660) (9,34472) (10,39059) (11,29172) (12,23025) (13,11192) (14,3163) (15,1230) (16,377) (17,109) (18,34)};\end{axis}\end{tikzpicture}\end{minipage}
\vspace{3mm}

\noindent\begin{minipage}[t]{60mm}\vspace{0pt}\sbox\hmboard{\chessboard[setfen=8/8/8/8/8/8/5n2/kBK4n b - - 0 1]}\usebox\hmboard\\[1mm]
\makebox[\wd\hmboard]{\sffamily\hspace{3mm}\textbf{h\#9}\hfill(2+3)}\end{minipage}\hfill
\begin{minipage}[t]{52mm}\vspace{0pt}
\textcolor{ink}{\sffamily\textsc{Lösung}}\par\smallskip\small
\begin{tabular}{@{}r@{\enspace}l@{\enspace}l@{}}
\textbf{1.} & \symknight{}h3 & \symking{}d2\\
\textbf{2.} & \symking{}b2 & \symking{}e3\\
\textbf{3.} & \symking{}c1 & \symbishop{}e4\\
\textbf{4.} & \symking{}d1 & \symbishop{}xh1\\
\textbf{5.} & \symking{}e1 & \symking{}f3\\
\textbf{6.} & \symking{}f1 & \symbishop{}g2+\\
\textbf{7.} & \symking{}g1 & \symbishop{}f1\\
\textbf{8.} & \symking{}h1 & \symking{}g3\\
\textbf{9.} & \symknight{}g1 & \symbishop{}g2\#\\
\end{tabular}\par\vspace{3mm}
\qrcode[height=12mm]{https://helpman.komtera.lt/?fen=8/8/8/8/8/8/5n2/kBK4n&moves=9}\\{\sffamily\tiny\color{ink}scannen zum Lösen}
\end{minipage}
\vspace{2mm}

{\raggedright\noindent\textcolor{ink}{\sffamily\textsc{Themen}}\enspace \themetag{reines Matt} \themetag{Mustermatt} \themetag{Idealmatt} \themetag{Switchback} \themetag{Selbstblock} \themetag{Kniest} \themetag{Umnov}\par}

\clearpage
\phantomsection\label{hm:KBvkpp}
\begin{tcolorbox}[enhanced,colback=ink,colframe=ink,arc=0pt,boxrule=0pt,left=4mm,right=4mm,top=2.5mm,bottom=2.5mm]
\color{white}\sffamily{\LARGE\bfseries KBvkpp}\hfill{\Large h\#9}\\[0.5mm]
{\small 5 Steine\quad Schwarz am Zug}\hfill{\small Nr.~57\qquad\href{https://helpman.komtera.lt/?fen=8/6k1/8/p3p3/8/8/1B6/K7&moves=9}{\color{white}\bfseries lösen\,\ensuremath{\nearrow}}}
\end{tcolorbox}
\vspace{2mm}
\noindent\begin{minipage}[t]{36mm}\vspace{0pt}\sffamily\small\raggedright
\hmfigure{längste eindeutige}{{\Large h\#9}}\hmfigure{eindeutig bei dieser Länge}{{\Large 51} Stellungen}\hmfigure{längstes Matt}{{\Large h\#16}}\hmfigure{in der Klasse}{533 Millionen Stellungen,\\ 208 Millionen mit Hilfsmatt}
\end{minipage}\hfill
\begin{minipage}[t]{80mm}\vspace{0pt}\begin{tikzpicture}\begin{axis}[width=72mm,height=42mm,ybar,bar width=1.6pt,xmin=-1,xmax=33,ymin=0,ymax=25.77,axis y line*=left,axis x line*=bottom,xlabel={\scriptsize Halbzüge bis Matt},ylabel={\scriptsize Stellungen (Mio.)},tick label style={font=\tiny},label style={font=\scriptsize},enlarge x limits=false,ylabel shift=-3pt,xlabel shift=-3pt]\addplot[fill=ink!55,draw=none] coordinates {(0,0.00114) (1,0.004387) (2,0.009944) (3,0.03381) (4,0.109) (5,0.271) (6,0.8116) (7,1.46) (8,3.516) (9,5.057) (10,10.08) (11,12.47) (12,20.67) (13,18.02) (14,23.43) (15,18.94) (16,19.6) (17,16.49) (18,13.78) (19,12.08) (20,8.69) (21,7.99) (22,4.681) (23,4.335) (24,2.049) (25,1.938) (26,0.764) (27,0.7277) (28,0.1749) (29,0.1701) (30,0.03874) (31,0.03667) (32,6e-06)};\draw[densely dotted] (axis cs:18,0) -- (axis cs:18,25.77);\end{axis}\begin{axis}[width=72mm,height=42mm,axis y line*=right,axis x line=none,ymode=log,xmin=-1,xmax=33,ylabel={\scriptsize eindeutig (log)},tick label style={font=\tiny},label style={font=\scriptsize},ylabel shift=-3pt]\addplot[mark=*,mark size=0.7pt,draw=black,thin] coordinates {(0,1140) (1,4367) (2,8815) (3,14371) (4,39197) (5,46816) (6,76962) (7,64119) (8,71636) (9,62791) (10,67053) (11,51494) (12,41130) (13,7906) (14,6108) (15,920) (16,653) (17,37) (18,51)};\end{axis}\end{tikzpicture}\end{minipage}
\vspace{3mm}

\noindent\begin{minipage}[t]{60mm}\vspace{0pt}\sbox\hmboard{\chessboard[setfen=8/6k1/8/p3p3/8/8/1B6/K7 b - - 0 1]}\usebox\hmboard\\[1mm]
\makebox[\wd\hmboard]{\sffamily\hspace{3mm}\textbf{h\#9}\hfill(2+3)}\end{minipage}\hfill
\begin{minipage}[t]{52mm}\vspace{0pt}
\textcolor{ink}{\sffamily\textsc{Lösung}}\par\smallskip\small
\begin{tabular}{@{}r@{\enspace}l@{\enspace}l@{}}
\textbf{1.} & \symking{}h8 & \symbishop{}a3\\
\textbf{2.} & e4 & \symking{}b2\\
\textbf{3.} & e3 & \symking{}c3\\
\textbf{4.} & e2 & \symking{}d4\\
\textbf{5.} & e1=\symknight{} & \symking{}e5\\
\textbf{6.} & \symknight{}f3+ & \symking{}f6\\
\textbf{7.} & \symknight{}g5 & \symbishop{}f8\\
\textbf{8.} & \symknight{}h7+ & \symking{}f7\\
\textbf{9.} & a4 & \symbishop{}g7\#\\
\end{tabular}\par\vspace{3mm}
\qrcode[height=12mm]{https://helpman.komtera.lt/?fen=8/6k1/8/p3p3/8/8/1B6/K7&moves=9}\\{\sffamily\tiny\color{ink}scannen zum Lösen}
\end{minipage}
\vspace{2mm}

{\raggedright\noindent\textcolor{ink}{\sffamily\textsc{Themen}}\enspace \themetag{reines Matt} \themetag{Mustermatt} \themetag{Umwandlung} \themetag{Unterverwandlung} \themetag{Selbstblock} \themetag{ohne Schlag}\par}

\clearpage
\phantomsection\label{hm:KNPvkp}
\begin{tcolorbox}[enhanced,colback=ink,colframe=ink,arc=0pt,boxrule=0pt,left=4mm,right=4mm,top=2.5mm,bottom=2.5mm]
\color{white}\sffamily{\LARGE\bfseries KNPvkp}\hfill{\Large h\#9}\\[0.5mm]
{\small 5 Steine\quad Schwarz am Zug}\hfill{\small Nr.~58\qquad\href{https://helpman.komtera.lt/?fen=N7/8/8/8/6k1/4p3/4P3/K7&moves=9}{\color{white}\bfseries lösen\,\ensuremath{\nearrow}}}
\end{tcolorbox}
\vspace{2mm}
\noindent\begin{minipage}[t]{36mm}\vspace{0pt}\sffamily\small\raggedright
\hmfigure{längste eindeutige}{{\Large h\#9}}\hmfigure{eindeutig bei dieser Länge}{{\Large 9} Stellungen}\hmfigure{längstes Matt}{{\Large h\#14}}\hmfigure{in der Klasse}{533 Millionen Stellungen,\\ 428 Millionen mit Hilfsmatt}
\end{minipage}\hfill
\begin{minipage}[t]{80mm}\vspace{0pt}\begin{tikzpicture}\begin{axis}[width=72mm,height=42mm,ybar,bar width=1.6pt,xmin=-1,xmax=29,ymin=0,ymax=61.3,axis y line*=left,axis x line*=bottom,xlabel={\scriptsize Halbzüge bis Matt},ylabel={\scriptsize Stellungen (Mio.)},tick label style={font=\tiny},label style={font=\scriptsize},enlarge x limits=false,ylabel shift=-3pt,xlabel shift=-3pt]\addplot[fill=ink!55,draw=none] coordinates {(0,0.001539) (1,0.124) (2,0.2851) (3,2.463) (4,6.322) (5,24.36) (6,36.22) (7,44.06) (8,48.01) (9,46.36) (10,50.84) (11,53.06) (12,55.73) (13,28.5) (14,20.75) (15,4.375) (16,3.685) (17,1.468) (18,1.112) (19,0.1693) (20,0.1113) (21,0.008614) (22,0.004389) (23,2.3e-05) (24,3.1e-05) (26,1e-05) (28,4e-06)};\draw[densely dotted] (axis cs:18,0) -- (axis cs:18,61.3);\end{axis}\begin{axis}[width=72mm,height=42mm,axis y line*=right,axis x line=none,ymode=log,xmin=-1,xmax=29,ylabel={\scriptsize eindeutig (log)},tick label style={font=\tiny},label style={font=\scriptsize},ylabel shift=-3pt]\addplot[mark=*,mark size=0.7pt,draw=black,thin] coordinates {(0,1539) (1,30915) (2,78150) (3,1018167) (4,1366404) (5,1814923) (6,1196678) (7,906162) (8,621261) (9,503640) (10,352129) (11,352761) (12,188633) (13,15456) (14,12238) (15,5542) (16,1296) (17,79) (18,9)};\end{axis}\end{tikzpicture}\end{minipage}
\vspace{3mm}

\noindent\begin{minipage}[t]{60mm}\vspace{0pt}\sbox\hmboard{\chessboard[setfen=N7/8/8/8/6k1/4p3/4P3/K7 b - - 0 1]}\usebox\hmboard\\[1mm]
\makebox[\wd\hmboard]{\sffamily\hspace{3mm}\textbf{h\#9}\hfill(3+2)}\end{minipage}\hfill
\begin{minipage}[t]{52mm}\vspace{0pt}
\textcolor{ink}{\sffamily\textsc{Lösung}}\par\smallskip\small
\begin{tabular}{@{}r@{\enspace}l@{\enspace}l@{}}
\textbf{1.} & \symking{}g3 & \symknight{}b6\\
\textbf{2.} & \symking{}f2 & \symknight{}c4\\
\textbf{3.} & \symking{}e1 & \symknight{}d2\\
\textbf{4.} & exd2 & e4\\
\textbf{5.} & \symking{}d1 & e5\\
\textbf{6.} & \symking{}c1 & e6\\
\textbf{7.} & d1=\symrook{} & e7\\
\textbf{8.} & \symrook{}d2 & e8=\symqueen{}\\
\textbf{9.} & \symrook{}c2 & \symqueen{}e1\#\\
\end{tabular}\par\vspace{3mm}
\qrcode[height=12mm]{https://helpman.komtera.lt/?fen=N7/8/8/8/6k1/4p3/4P3/K7&moves=9}\\{\sffamily\tiny\color{ink}scannen zum Lösen}
\end{minipage}
\vspace{2mm}

{\raggedright\noindent\textcolor{ink}{\sffamily\textsc{Themen}}\enspace \themetag{Umwandlung} \themetag{Unterverwandlung} \themetag{Excelsior} \themetag{Switchback} \themetag{Selbstblock} \themetag{ohne Schach}\par}

\clearpage
\phantomsection\label{hm:KPvknp}
\begin{tcolorbox}[enhanced,colback=ink,colframe=ink,arc=0pt,boxrule=0pt,left=4mm,right=4mm,top=2.5mm,bottom=2.5mm]
\color{white}\sffamily{\LARGE\bfseries KPvknp}\hfill{\Large h\#9}\\[0.5mm]
{\small 5 Steine\quad Schwarz am Zug}\hfill{\small Nr.~59\qquad\href{https://helpman.komtera.lt/?fen=6k1/8/7n/8/8/p7/P7/K7&moves=9}{\color{white}\bfseries lösen\,\ensuremath{\nearrow}}}
\end{tcolorbox}
\vspace{2mm}
\noindent\begin{minipage}[t]{36mm}\vspace{0pt}\sffamily\small\raggedright
\hmfigure{längste eindeutige}{{\Large h\#9}}\hmfigure{eindeutig bei dieser Länge}{{\Large 4} Stellungen}\hmfigure{längstes Matt}{{\Large h\#14.5}}\hmfigure{in der Klasse}{533 Millionen Stellungen,\\ 428 Millionen mit Hilfsmatt}
\end{minipage}\hfill
\begin{minipage}[t]{80mm}\vspace{0pt}\begin{tikzpicture}\begin{axis}[width=72mm,height=42mm,ybar,bar width=1.6pt,xmin=-1,xmax=30,ymin=0,ymax=68.67,axis y line*=left,axis x line*=bottom,xlabel={\scriptsize Halbzüge bis Matt},ylabel={\scriptsize Stellungen (Mio.)},tick label style={font=\tiny},label style={font=\scriptsize},enlarge x limits=false,ylabel shift=-3pt,xlabel shift=-3pt]\addplot[fill=ink!55,draw=none] coordinates {(0,0.000194) (1,0.08141) (2,0.1942) (3,1.336) (4,2.956) (5,11.43) (6,17.98) (7,35.68) (8,41.06) (9,54.81) (10,50.55) (11,62.43) (12,57.03) (13,48.31) (14,31.87) (15,7.771) (16,2.929) (17,0.9934) (18,0.3415) (19,0.2102) (20,0.04479) (21,0.02456) (22,0.001526) (23,0.000728) (24,4e-06) (25,4.8e-05) (27,2e-05) (29,8e-06)};\draw[densely dotted] (axis cs:18,0) -- (axis cs:18,68.67);\end{axis}\begin{axis}[width=72mm,height=42mm,axis y line*=right,axis x line=none,ymode=log,xmin=-1,xmax=30,ylabel={\scriptsize eindeutig (log)},tick label style={font=\tiny},label style={font=\scriptsize},ylabel shift=-3pt]\addplot[mark=*,mark size=0.7pt,draw=black,thin] coordinates {(0,194) (1,15722) (2,30795) (3,264224) (4,447317) (5,979177) (6,712453) (7,605360) (8,390689) (9,426524) (10,296740) (11,368203) (12,230814) (13,54052) (14,10647) (15,1796) (16,424) (17,79) (18,4)};\end{axis}\end{tikzpicture}\end{minipage}
\vspace{3mm}

\noindent\begin{minipage}[t]{60mm}\vspace{0pt}\sbox\hmboard{\chessboard[setfen=6k1/8/7n/8/8/p7/P7/K7 b - - 0 1]}\usebox\hmboard\\[1mm]
\makebox[\wd\hmboard]{\sffamily\hspace{3mm}\textbf{h\#9}\hfill(2+3)}\end{minipage}\hfill
\begin{minipage}[t]{52mm}\vspace{0pt}
\textcolor{ink}{\sffamily\textsc{Lösung}}\par\smallskip\small
\begin{tabular}{@{}r@{\enspace}l@{\enspace}l@{}}
\textbf{1.} & \symknight{}f5 & \symking{}b1\\
\textbf{2.} & \symknight{}d4 & \symking{}c1\\
\textbf{3.} & \symknight{}b3+ & axb3\\
\textbf{4.} & \symking{}f7 & b4\\
\textbf{5.} & \symking{}e6 & b5\\
\textbf{6.} & \symking{}d5 & b6\\
\textbf{7.} & \symking{}c4 & b7\\
\textbf{8.} & \symking{}b3 & b8=\symqueen{}+\\
\textbf{9.} & \symking{}a2 & \symqueen{}b1\#\\
\end{tabular}\par\vspace{3mm}
\qrcode[height=12mm]{https://helpman.komtera.lt/?fen=6k1/8/7n/8/8/p7/P7/K7&moves=9}\\{\sffamily\tiny\color{ink}scannen zum Lösen}
\end{minipage}
\vspace{2mm}

{\raggedright\noindent\textcolor{ink}{\sffamily\textsc{Themen}}\enspace \themetag{Umwandlung} \themetag{Excelsior}\par}

\clearpage
\phantomsection\label{hm:KNvkbb}
\begin{tcolorbox}[enhanced,colback=ink,colframe=ink,arc=0pt,boxrule=0pt,left=4mm,right=4mm,top=2.5mm,bottom=2.5mm]
\color{white}\sffamily{\LARGE\bfseries KNvkbb}\hfill{\Large h\#8.5}\\[0.5mm]
{\small 5 Steine\quad Weiß am Zug}\hfill{\small Nr.~60\qquad\href{https://helpman.komtera.lt/?fen=8/8/N7/7k/b1b5/8/8/K7&moves=8.5}{\color{white}\bfseries lösen\,\ensuremath{\nearrow}}}
\end{tcolorbox}
\vspace{2mm}
\noindent\begin{minipage}[t]{36mm}\vspace{0pt}\sffamily\small\raggedright
\hmfigure{längste eindeutige}{{\Large h\#8.5}}\hmfigure{eindeutig bei dieser Länge}{{\Large 69} Stellungen}\hmfigure{längstes Matt}{{\Large h\#10.5}}\hmfigure{in der Klasse}{242 Millionen Stellungen,\\ 85,4 Millionen mit Hilfsmatt}
\end{minipage}\hfill
\begin{minipage}[t]{80mm}\vspace{0pt}\begin{tikzpicture}\begin{axis}[width=72mm,height=42mm,ybar,bar width=1.6pt,xmin=-1,xmax=22,ymin=0,ymax=16.54,axis y line*=left,axis x line*=bottom,xlabel={\scriptsize Halbzüge bis Matt},ylabel={\scriptsize Stellungen (Mio.)},tick label style={font=\tiny},label style={font=\scriptsize},enlarge x limits=false,ylabel shift=-3pt,xlabel shift=-3pt]\addplot[fill=ink!55,draw=none] coordinates {(0,4.4e-05) (1,0.000222) (2,0.001744) (3,0.01072) (4,0.04498) (5,0.1983) (6,0.396) (7,1.271) (8,1.901) (9,4.638) (10,6.357) (11,11.82) (12,11.73) (13,15.03) (14,11.26) (15,10.21) (16,5.918) (17,3.474) (18,0.797) (19,0.3444) (20,0.000336) (21,4e-06)};\draw[densely dotted] (axis cs:17,0) -- (axis cs:17,16.54);\end{axis}\begin{axis}[width=72mm,height=42mm,axis y line*=right,axis x line=none,ymode=log,xmin=-1,xmax=22,ylabel={\scriptsize eindeutig (log)},tick label style={font=\tiny},label style={font=\scriptsize},ylabel shift=-3pt]\addplot[mark=*,mark size=0.7pt,draw=black,thin] coordinates {(0,44) (1,222) (2,1484) (3,6875) (4,19179) (5,37414) (6,20456) (7,37073) (8,36036) (9,61063) (10,63450) (11,52333) (12,20700) (13,7973) (14,4191) (15,882) (16,522) (17,69)};\end{axis}\end{tikzpicture}\end{minipage}
\vspace{3mm}

\noindent\begin{minipage}[t]{60mm}\vspace{0pt}\sbox\hmboard{\chessboard[setfen=8/8/N7/7k/b1b5/8/8/K7 w - - 0 1]}\usebox\hmboard\\[1mm]
\makebox[\wd\hmboard]{\sffamily\hspace{3mm}\textbf{h\#8.5}\hfill(2+3)}\end{minipage}\hfill
\begin{minipage}[t]{52mm}\vspace{0pt}
\textcolor{ink}{\sffamily\textsc{Lösung}}\par\smallskip\small
\begin{tabular}{@{}r@{\enspace}l@{\enspace}l@{}}
\textbf{1.} & \ldots & \symking{}b2\\
\textbf{2.} & \symking{}g4 & \symking{}a3\\
\textbf{3.} & \symking{}f3 & \symking{}xa4\\
\textbf{4.} & \symking{}e2 & \symking{}b4\\
\textbf{5.} & \symking{}d1 & \symking{}c3\\
\textbf{6.} & \symking{}c1 & \symknight{}b4\\
\textbf{7.} & \symking{}b1 & \symking{}d2\\
\textbf{8.} & \symking{}a1 & \symking{}c1\\
\textbf{9.} & \symbishop{}a2 & \symknight{}c2\#\\
\end{tabular}\par\vspace{3mm}
\qrcode[height=12mm]{https://helpman.komtera.lt/?fen=8/8/N7/7k/b1b5/8/8/K7&moves=8.5}\\{\sffamily\tiny\color{ink}scannen zum Lösen}
\end{minipage}
\vspace{2mm}

{\raggedright\noindent\textcolor{ink}{\sffamily\textsc{Themen}}\enspace \themetag{reines Matt} \themetag{Mustermatt} \themetag{Idealmatt} \themetag{Selbstblock} \themetag{ohne Schach}\par}

\clearpage
\phantomsection\label{hm:KNvkbp}
\begin{tcolorbox}[enhanced,colback=ink,colframe=ink,arc=0pt,boxrule=0pt,left=4mm,right=4mm,top=2.5mm,bottom=2.5mm]
\color{white}\sffamily{\LARGE\bfseries KNvkbp}\hfill{\Large h\#8.5}\\[0.5mm]
{\small 5 Steine\quad Weiß am Zug}\hfill{\small Nr.~61\qquad\href{https://helpman.komtera.lt/?fen=k7/8/8/7p/8/8/5b2/3K3N&moves=8.5}{\color{white}\bfseries lösen\,\ensuremath{\nearrow}}}
\end{tcolorbox}
\vspace{2mm}
\noindent\begin{minipage}[t]{36mm}\vspace{0pt}\sffamily\small\raggedright
\hmfigure{längste eindeutige}{{\Large h\#8.5}}\hmfigure{eindeutig bei dieser Länge}{{\Large 12} Stellungen}\hmfigure{längstes Matt}{{\Large h\#10.5}}\hmfigure{in der Klasse}{710 Millionen Stellungen,\\ 542 Millionen mit Hilfsmatt}
\end{minipage}\hfill
\begin{minipage}[t]{80mm}\vspace{0pt}\begin{tikzpicture}\begin{axis}[width=72mm,height=42mm,ybar,bar width=1.6pt,xmin=-1,xmax=22,ymin=0,ymax=105,axis y line*=left,axis x line*=bottom,xlabel={\scriptsize Halbzüge bis Matt},ylabel={\scriptsize Stellungen (Mio.)},tick label style={font=\tiny},label style={font=\scriptsize},enlarge x limits=false,ylabel shift=-3pt,xlabel shift=-3pt]\addplot[fill=ink!55,draw=none] coordinates {(0,0.00039) (1,0.001891) (2,0.00968) (3,0.05956) (4,0.277) (5,1.165) (6,2.972) (7,8.584) (8,17.14) (9,34.71) (10,55.82) (11,82.26) (12,90.75) (13,95.44) (14,70.25) (15,51.19) (16,21.39) (17,8.639) (18,1.292) (19,0.1941) (20,0.00054) (21,1e-06)};\draw[densely dotted] (axis cs:17,0) -- (axis cs:17,105);\end{axis}\begin{axis}[width=72mm,height=42mm,axis y line*=right,axis x line=none,ymode=log,xmin=-1,xmax=22,ylabel={\scriptsize eindeutig (log)},tick label style={font=\tiny},label style={font=\scriptsize},ylabel shift=-3pt]\addplot[mark=*,mark size=0.7pt,draw=black,thin] coordinates {(0,390) (1,1848) (2,8320) (3,37800) (4,139650) (5,216255) (6,205181) (7,245999) (8,318491) (9,348141) (10,367704) (11,214437) (12,109705) (13,36013) (14,20035) (15,2924) (16,485) (17,12)};\end{axis}\end{tikzpicture}\end{minipage}
\vspace{3mm}

\noindent\begin{minipage}[t]{60mm}\vspace{0pt}\sbox\hmboard{\chessboard[setfen=k7/8/8/7p/8/8/5b2/3K3N w - - 0 1]}\usebox\hmboard\\[1mm]
\makebox[\wd\hmboard]{\sffamily\hspace{3mm}\textbf{h\#8.5}\hfill(2+3)}\end{minipage}\hfill
\begin{minipage}[t]{52mm}\vspace{0pt}
\textcolor{ink}{\sffamily\textsc{Lösung}}\par\smallskip\small
\begin{tabular}{@{}r@{\enspace}l@{\enspace}l@{}}
\textbf{1.} & \ldots & \symking{}e2\\
\textbf{2.} & \symking{}b7 & \symking{}f3\\
\textbf{3.} & \symking{}c6 & \symking{}f4\\
\textbf{4.} & \symking{}d5 & \symking{}g5\\
\textbf{5.} & \symking{}e4 & \symknight{}g3+\\
\textbf{6.} & \symking{}f3 & \symknight{}xh5\\
\textbf{7.} & \symking{}g2 & \symking{}g4\\
\textbf{8.} & \symking{}h1 & \symking{}h3\\
\textbf{9.} & \symbishop{}g1 & \symknight{}g3\#\\
\end{tabular}\par\vspace{3mm}
\qrcode[height=12mm]{https://helpman.komtera.lt/?fen=k7/8/8/7p/8/8/5b2/3K3N&moves=8.5}\\{\sffamily\tiny\color{ink}scannen zum Lösen}
\end{minipage}
\vspace{2mm}

{\raggedright\noindent\textcolor{ink}{\sffamily\textsc{Themen}}\enspace \themetag{reines Matt} \themetag{Mustermatt} \themetag{Idealmatt} \themetag{Switchback} \themetag{Selbstblock}\par}

\clearpage
\phantomsection\label{hm:KNvknn}
\begin{tcolorbox}[enhanced,colback=ink,colframe=ink,arc=0pt,boxrule=0pt,left=4mm,right=4mm,top=2.5mm,bottom=2.5mm]
\color{white}\sffamily{\LARGE\bfseries KNvknn}\hfill{\Large h\#8.5}\\[0.5mm]
{\small 5 Steine\quad Weiß am Zug}\hfill{\small Nr.~62\qquad\href{https://helpman.komtera.lt/?fen=7k/8/8/8/8/8/8/KNn1n3&moves=8.5}{\color{white}\bfseries lösen\,\ensuremath{\nearrow}}}
\end{tcolorbox}
\vspace{2mm}
\noindent\begin{minipage}[t]{36mm}\vspace{0pt}\sffamily\small\raggedright
\hmfigure{längste eindeutige}{{\Large h\#8.5}}\hmfigure{eindeutig bei dieser Länge}{{\Large 3} Stellungen}\hmfigure{längstes Matt}{{\Large h\#9.5}}\hmfigure{in der Klasse}{242 Millionen Stellungen,\\ 90,1 Millionen mit Hilfsmatt}
\end{minipage}\hfill
\begin{minipage}[t]{80mm}\vspace{0pt}\begin{tikzpicture}\begin{axis}[width=72mm,height=42mm,ybar,bar width=1.6pt,xmin=-1,xmax=20,ymin=0,ymax=18.18,axis y line*=left,axis x line*=bottom,xlabel={\scriptsize Halbzüge bis Matt},ylabel={\scriptsize Stellungen (Mio.)},tick label style={font=\tiny},label style={font=\scriptsize},enlarge x limits=false,ylabel shift=-3pt,xlabel shift=-3pt]\addplot[fill=ink!55,draw=none] coordinates {(0,5.6e-05) (1,0.000293) (2,0.000981) (3,0.00801) (4,0.03872) (5,0.1933) (6,0.6553) (7,1.936) (8,4.338) (9,8.181) (10,12.22) (11,16.53) (12,16.29) (13,15.76) (14,8.522) (15,4.24) (16,1.055) (17,0.164) (18,0.01115) (19,0.000361)};\draw[densely dotted] (axis cs:17,0) -- (axis cs:17,18.18);\end{axis}\begin{axis}[width=72mm,height=42mm,axis y line*=right,axis x line=none,ymode=log,xmin=-1,xmax=20,ylabel={\scriptsize eindeutig (log)},tick label style={font=\tiny},label style={font=\scriptsize},ylabel shift=-3pt]\addplot[mark=*,mark size=0.7pt,draw=black,thin] coordinates {(0,56) (1,293) (2,715) (3,4610) (4,19674) (5,41056) (6,54790) (7,51067) (8,48444) (9,41770) (10,32227) (11,24817) (12,7923) (13,2888) (14,550) (15,32) (16,4) (17,3)};\end{axis}\end{tikzpicture}\end{minipage}
\vspace{3mm}

\noindent\begin{minipage}[t]{60mm}\vspace{0pt}\sbox\hmboard{\chessboard[setfen=7k/8/8/8/8/8/8/KNn1n3 w - - 0 1]}\usebox\hmboard\\[1mm]
\makebox[\wd\hmboard]{\sffamily\hspace{3mm}\textbf{h\#8.5}\hfill(2+3)}\end{minipage}\hfill
\begin{minipage}[t]{52mm}\vspace{0pt}
\textcolor{ink}{\sffamily\textsc{Lösung}}\par\smallskip\small
\begin{tabular}{@{}r@{\enspace}l@{\enspace}l@{}}
\textbf{1.} & \ldots & \symking{}b2\\
\textbf{2.} & \symking{}g7 & \symking{}c3\\
\textbf{3.} & \symking{}f6 & \symking{}d2\\
\textbf{4.} & \symking{}e5 & \symking{}d1\\
\textbf{5.} & \symking{}d4 & \symknight{}d2\\
\textbf{6.} & \symking{}c3 & \symking{}xe1\\
\textbf{7.} & \symking{}b2 & \symking{}d1\\
\textbf{8.} & \symking{}a1 & \symking{}c2\\
\textbf{9.} & \symknight{}a2 & \symknight{}b3\#\\
\end{tabular}\par\vspace{3mm}
\qrcode[height=12mm]{https://helpman.komtera.lt/?fen=7k/8/8/8/8/8/8/KNn1n3&moves=8.5}\\{\sffamily\tiny\color{ink}scannen zum Lösen}
\end{minipage}
\vspace{2mm}

{\raggedright\noindent\textcolor{ink}{\sffamily\textsc{Themen}}\enspace \themetag{reines Matt} \themetag{Mustermatt} \themetag{Idealmatt} \themetag{Switchback} \themetag{Selbstblock} \themetag{ohne Schach}\par}

\clearpage
\phantomsection\label{hm:KNvknp}
\begin{tcolorbox}[enhanced,colback=ink,colframe=ink,arc=0pt,boxrule=0pt,left=4mm,right=4mm,top=2.5mm,bottom=2.5mm]
\color{white}\sffamily{\LARGE\bfseries KNvknp}\hfill{\Large h\#8.5}\\[0.5mm]
{\small 5 Steine\quad Weiß am Zug}\hfill{\small Nr.~63\qquad\href{https://helpman.komtera.lt/?fen=N5k1/8/8/8/8/1p6/8/Kn6&moves=8.5}{\color{white}\bfseries lösen\,\ensuremath{\nearrow}}}
\end{tcolorbox}
\vspace{2mm}
\noindent\begin{minipage}[t]{36mm}\vspace{0pt}\sffamily\small\raggedright
\hmfigure{längste eindeutige}{{\Large h\#8.5}}\hmfigure{eindeutig bei dieser Länge}{{\Large 1} Stellung}\hmfigure{längstes Matt}{{\Large h\#10}}\hmfigure{in der Klasse}{710 Millionen Stellungen,\\ 557 Millionen mit Hilfsmatt}
\end{minipage}\hfill
\begin{minipage}[t]{80mm}\vspace{0pt}\begin{tikzpicture}\begin{axis}[width=72mm,height=42mm,ybar,bar width=1.6pt,xmin=-1,xmax=21,ymin=0,ymax=114.8,axis y line*=left,axis x line*=bottom,xlabel={\scriptsize Halbzüge bis Matt},ylabel={\scriptsize Stellungen (Mio.)},tick label style={font=\tiny},label style={font=\scriptsize},enlarge x limits=false,ylabel shift=-3pt,xlabel shift=-3pt]\addplot[fill=ink!55,draw=none] coordinates {(0,0.000412) (1,0.002013) (2,0.005573) (3,0.03946) (4,0.1722) (5,0.7762) (6,2.55) (7,7.062) (8,16.59) (9,31.36) (10,53.78) (11,78.1) (12,96.21) (13,104.3) (14,79.77) (15,53.39) (16,23.66) (17,7.041) (18,1.72) (19,0.1456) (20,0.01634)};\draw[densely dotted] (axis cs:17,0) -- (axis cs:17,114.8);\end{axis}\begin{axis}[width=72mm,height=42mm,axis y line*=right,axis x line=none,ymode=log,xmin=-1,xmax=21,ylabel={\scriptsize eindeutig (log)},tick label style={font=\tiny},label style={font=\scriptsize},ylabel shift=-3pt]\addplot[mark=*,mark size=0.7pt,draw=black,thin] coordinates {(0,412) (1,1963) (2,4181) (3,23114) (4,92023) (5,165829) (6,241734) (7,206811) (8,235504) (9,225198) (10,249418) (11,171123) (12,91107) (13,27784) (14,6473) (15,849) (16,136) (17,1)};\end{axis}\end{tikzpicture}\end{minipage}
\vspace{3mm}

\noindent\begin{minipage}[t]{60mm}\vspace{0pt}\sbox\hmboard{\chessboard[setfen=N5k1/8/8/8/8/1p6/8/Kn6 w - - 0 1]}\usebox\hmboard\\[1mm]
\makebox[\wd\hmboard]{\sffamily\hspace{3mm}\textbf{h\#8.5}\hfill(2+3)}\end{minipage}\hfill
\begin{minipage}[t]{52mm}\vspace{0pt}
\textcolor{ink}{\sffamily\textsc{Lösung}}\par\smallskip\small
\begin{tabular}{@{}r@{\enspace}l@{\enspace}l@{}}
\textbf{1.} & \ldots & \symking{}b2\\
\textbf{2.} & \symknight{}c3 & \symking{}a3\\
\textbf{3.} & \symking{}f7 & \symking{}b4\\
\textbf{4.} & \symking{}e6 & \symking{}a5\\
\textbf{5.} & \symking{}d5 & \symking{}b6\\
\textbf{6.} & \symking{}c4 & \symking{}a6\\
\textbf{7.} & \symking{}b4 & \symknight{}b6\\
\textbf{8.} & \symking{}a3 & \symking{}a5\\
\textbf{9.} & \symknight{}a2 & \symknight{}c4\#\\
\end{tabular}\par\vspace{3mm}
\qrcode[height=12mm]{https://helpman.komtera.lt/?fen=N5k1/8/8/8/8/1p6/8/Kn6&moves=8.5}\\{\sffamily\tiny\color{ink}scannen zum Lösen}
\end{minipage}
\vspace{2mm}

{\raggedright\noindent\textcolor{ink}{\sffamily\textsc{Themen}}\enspace \themetag{Satzspiel} \themetag{reines Matt} \themetag{Mustermatt} \themetag{Idealmatt} \themetag{Rundlauf} \themetag{Selbstblock} \themetag{ohne Schlag} \themetag{ohne Schach}\par}

\begin{tcolorbox}[enhanced,colback=white,colframe=ink!35,arc=1.5pt,boxrule=0.4pt,left=3mm,right=3mm,top=1.5mm,bottom=1.5mm]
\textcolor{ink}{\sffamily\textsc{veröffentlicht}}\par\smallskip\small \textbf{Sheglow, Wiktor S.}, h\#9 0.1...\quad{\color{ink}\scriptsize diese Stellung}\\{\scriptsize Suomen Tehtäväniekat, No. (286), 05-06/1998}
\end{tcolorbox}
\clearpage
\phantomsection\label{hm:KNvkqr}
\begin{tcolorbox}[enhanced,colback=ink,colframe=ink,arc=0pt,boxrule=0pt,left=4mm,right=4mm,top=2.5mm,bottom=2.5mm]
\color{white}\sffamily{\LARGE\bfseries KNvkqr}\hfill{\Large h\#8.5}\\[0.5mm]
{\small 5 Steine\quad Weiß am Zug}\hfill{\small Nr.~64\qquad\href{https://helpman.komtera.lt/?fen=8/7k/N7/8/8/8/4r3/Kq6&moves=8.5}{\color{white}\bfseries lösen\,\ensuremath{\nearrow}}}
\end{tcolorbox}
\vspace{2mm}
\noindent\begin{minipage}[t]{36mm}\vspace{0pt}\sffamily\small\raggedright
\hmfigure{längste eindeutige}{{\Large h\#8.5}}\hmfigure{eindeutig bei dieser Länge}{{\Large 8} Stellungen}\hmfigure{längstes Matt}{{\Large h\#9.5}}\hmfigure{in der Klasse}{242 Millionen Stellungen,\\ 146 Millionen mit Hilfsmatt}
\end{minipage}\hfill
\begin{minipage}[t]{80mm}\vspace{0pt}\begin{tikzpicture}\begin{axis}[width=72mm,height=42mm,ybar,bar width=1.6pt,xmin=-1,xmax=20,ymin=0,ymax=41.52,axis y line*=left,axis x line*=bottom,xlabel={\scriptsize Halbzüge bis Matt},ylabel={\scriptsize Stellungen (Mio.)},tick label style={font=\tiny},label style={font=\scriptsize},enlarge x limits=false,ylabel shift=-3pt,xlabel shift=-3pt]\addplot[fill=ink!55,draw=none] coordinates {(0,3.6e-05) (1,0.000149) (2,0.001902) (3,0.01314) (4,0.08877) (5,0.454) (6,1.637) (7,6.34) (8,11.98) (9,29.53) (10,22.22) (11,37.75) (12,14.63) (13,16.34) (14,2.54) (15,1.719) (16,0.0588) (17,0.1912) (19,0.01492)};\draw[densely dotted] (axis cs:17,0) -- (axis cs:17,41.52);\end{axis}\begin{axis}[width=72mm,height=42mm,axis y line*=right,axis x line=none,ymode=log,xmin=-1,xmax=20,ylabel={\scriptsize eindeutig (log)},tick label style={font=\tiny},label style={font=\scriptsize},ylabel shift=-3pt]\addplot[mark=*,mark size=0.7pt,draw=black,thin] coordinates {(0,36) (1,149) (2,1737) (3,9942) (4,31684) (5,109248) (6,163460) (7,410778) (8,258218) (9,415253) (10,122803) (11,164302) (12,18451) (13,14469) (14,231) (15,237) (17,8)};\end{axis}\end{tikzpicture}\end{minipage}
\vspace{3mm}

\noindent\begin{minipage}[t]{60mm}\vspace{0pt}\sbox\hmboard{\chessboard[setfen=8/7k/N7/8/8/8/4r3/Kq6 w - - 0 1]}\usebox\hmboard\\[1mm]
\makebox[\wd\hmboard]{\sffamily\hspace{3mm}\textbf{h\#8.5}\hfill(2+3)}\end{minipage}\hfill
\begin{minipage}[t]{52mm}\vspace{0pt}
\textcolor{ink}{\sffamily\textsc{Lösung}}\par\smallskip\small
\begin{tabular}{@{}r@{\enspace}l@{\enspace}l@{}}
\textbf{1.} & \ldots & \symking{}xb1\\
\textbf{2.} & \symking{}g6 & \symknight{}b4\\
\textbf{3.} & \symking{}f5 & \symknight{}c2\\
\textbf{4.} & \symking{}e4 & \symking{}b2\\
\textbf{5.} & \symking{}d3 & \symking{}c1\\
\textbf{6.} & \symking{}c3 & \symking{}d1\\
\textbf{7.} & \symking{}b2 & \symknight{}d4\\
\textbf{8.} & \symking{}a1 & \symking{}c1\\
\textbf{9.} & \symrook{}a2 & \symknight{}b3\#\\
\end{tabular}\par\vspace{3mm}
\qrcode[height=12mm]{https://helpman.komtera.lt/?fen=8/7k/N7/8/8/8/4r3/Kq6&moves=8.5}\\{\sffamily\tiny\color{ink}scannen zum Lösen}
\end{minipage}
\vspace{2mm}

{\raggedright\noindent\textcolor{ink}{\sffamily\textsc{Themen}}\enspace \themetag{reines Matt} \themetag{Mustermatt} \themetag{Idealmatt} \themetag{Switchback} \themetag{Selbstblock} \themetag{ohne Schach}\par\smallskip{\footnotesize\textit{Hinweis:} die Partei am Zug steht im Diagramm im Schach; die Lösung beginnt mit einem Schlag.}\par}

\clearpage
\phantomsection\label{hm:KNvkrn}
\begin{tcolorbox}[enhanced,colback=ink,colframe=ink,arc=0pt,boxrule=0pt,left=4mm,right=4mm,top=2.5mm,bottom=2.5mm]
\color{white}\sffamily{\LARGE\bfseries KNvkrn}\hfill{\Large h\#8.5}\\[0.5mm]
{\small 5 Steine\quad Weiß am Zug}\hfill{\small Nr.~65\qquad\href{https://helpman.komtera.lt/?fen=6k1/8/8/8/8/n7/4r3/1KN5&moves=8.5}{\color{white}\bfseries lösen\,\ensuremath{\nearrow}}}
\end{tcolorbox}
\vspace{2mm}
\noindent\begin{minipage}[t]{36mm}\vspace{0pt}\sffamily\small\raggedright
\hmfigure{längste eindeutige}{{\Large h\#8.5}}\hmfigure{eindeutig bei dieser Länge}{{\Large 1} Stellung}\hmfigure{längstes Matt}{{\Large h\#10.5}}\hmfigure{in der Klasse}{242 Millionen Stellungen,\\ 169 Millionen mit Hilfsmatt}
\end{minipage}\hfill
\begin{minipage}[t]{80mm}\vspace{0pt}\begin{tikzpicture}\begin{axis}[width=72mm,height=42mm,ybar,bar width=1.6pt,xmin=-1,xmax=22,ymin=0,ymax=41.29,axis y line*=left,axis x line*=bottom,xlabel={\scriptsize Halbzüge bis Matt},ylabel={\scriptsize Stellungen (Mio.)},tick label style={font=\tiny},label style={font=\scriptsize},enlarge x limits=false,ylabel shift=-3pt,xlabel shift=-3pt]\addplot[fill=ink!55,draw=none] coordinates {(0,0.000101) (1,0.000468) (2,0.003654) (3,0.02612) (4,0.1445) (5,0.6744) (6,2.003) (7,6.131) (8,12.04) (9,23.54) (10,27.61) (11,37.54) (12,24.83) (13,22.53) (14,7.693) (15,3.434) (16,0.5849) (17,0.05794) (18,0.001718) (19,0.00059) (21,1e-06)};\draw[densely dotted] (axis cs:17,0) -- (axis cs:17,41.29);\end{axis}\begin{axis}[width=72mm,height=42mm,axis y line*=right,axis x line=none,ymode=log,xmin=-1,xmax=22,ylabel={\scriptsize eindeutig (log)},tick label style={font=\tiny},label style={font=\scriptsize},ylabel shift=-3pt]\addplot[mark=*,mark size=0.7pt,draw=black,thin] coordinates {(0,101) (1,467) (2,3181) (3,18406) (4,62398) (5,148360) (6,153324) (7,205001) (8,165366) (9,188259) (10,118721) (11,81495) (12,16662) (13,9603) (14,948) (15,78) (16,3) (17,1)};\end{axis}\end{tikzpicture}\end{minipage}
\vspace{3mm}

\noindent\begin{minipage}[t]{60mm}\vspace{0pt}\sbox\hmboard{\chessboard[setfen=6k1/8/8/8/8/n7/4r3/1KN5 w - - 0 1]}\usebox\hmboard\\[1mm]
\makebox[\wd\hmboard]{\sffamily\hspace{3mm}\textbf{h\#8.5}\hfill(2+3)}\end{minipage}\hfill
\begin{minipage}[t]{52mm}\vspace{0pt}
\textcolor{ink}{\sffamily\textsc{Lösung}}\par\smallskip\small
\begin{tabular}{@{}r@{\enspace}l@{\enspace}l@{}}
\textbf{1.} & \ldots & \symking{}a1\\
\textbf{2.} & \symking{}f7 & \symknight{}xe2\\
\textbf{3.} & \symking{}e6 & \symking{}b2\\
\textbf{4.} & \symking{}d5 & \symking{}c1\\
\textbf{5.} & \symking{}c4 & \symking{}d2\\
\textbf{6.} & \symking{}b3 & \symknight{}d4+\\
\textbf{7.} & \symking{}a2 & \symking{}c3\\
\textbf{8.} & \symking{}a1 & \symking{}b3\\
\textbf{9.} & \symknight{}b1 & \symknight{}c2\#\\
\end{tabular}\par\vspace{3mm}
\qrcode[height=12mm]{https://helpman.komtera.lt/?fen=6k1/8/8/8/8/n7/4r3/1KN5&moves=8.5}\\{\sffamily\tiny\color{ink}scannen zum Lösen}
\end{minipage}
\vspace{2mm}

{\raggedright\noindent\textcolor{ink}{\sffamily\textsc{Themen}}\enspace \themetag{reines Matt} \themetag{Mustermatt} \themetag{Idealmatt} \themetag{Selbstblock}\par\smallskip{\footnotesize\textit{Hinweis:} die Partei am Zug steht im Diagramm im Schach.}\par}

\clearpage
\phantomsection\label{hm:KNvkrp}
\begin{tcolorbox}[enhanced,colback=ink,colframe=ink,arc=0pt,boxrule=0pt,left=4mm,right=4mm,top=2.5mm,bottom=2.5mm]
\color{white}\sffamily{\LARGE\bfseries KNvkrp}\hfill{\Large h\#8.5}\\[0.5mm]
{\small 5 Steine\quad Weiß am Zug}\hfill{\small Nr.~66\qquad\href{https://helpman.komtera.lt/?fen=6k1/8/8/1p6/8/8/1r6/K6N&moves=8.5}{\color{white}\bfseries lösen\,\ensuremath{\nearrow}}}
\end{tcolorbox}
\vspace{2mm}
\noindent\begin{minipage}[t]{36mm}\vspace{0pt}\sffamily\small\raggedright
\hmfigure{längste eindeutige}{{\Large h\#8.5}}\hmfigure{eindeutig bei dieser Länge}{{\Large 46} Stellungen}\hmfigure{längstes Matt}{{\Large h\#11.5}}\hmfigure{in der Klasse}{710 Millionen Stellungen,\\ 520 Millionen mit Hilfsmatt}
\end{minipage}\hfill
\begin{minipage}[t]{80mm}\vspace{0pt}\begin{tikzpicture}\begin{axis}[width=72mm,height=42mm,ybar,bar width=1.6pt,xmin=-1,xmax=24,ymin=0,ymax=113.2,axis y line*=left,axis x line*=bottom,xlabel={\scriptsize Halbzüge bis Matt},ylabel={\scriptsize Stellungen (Mio.)},tick label style={font=\tiny},label style={font=\scriptsize},enlarge x limits=false,ylabel shift=-3pt,xlabel shift=-3pt]\addplot[fill=ink!55,draw=none] coordinates {(0,0.000388) (1,0.00171) (2,0.01146) (3,0.07765) (4,0.409) (5,1.833) (6,4.716) (7,13.89) (8,25.14) (9,50.34) (10,69.14) (11,102.9) (12,84.87) (13,82.37) (14,43.96) (15,27.56) (16,9.73) (17,2.977) (18,0.3349) (19,0.06497) (21,0.005521) (23,0.000123)};\draw[densely dotted] (axis cs:17,0) -- (axis cs:17,113.2);\end{axis}\begin{axis}[width=72mm,height=42mm,axis y line*=right,axis x line=none,ymode=log,xmin=-1,xmax=24,ylabel={\scriptsize eindeutig (log)},tick label style={font=\tiny},label style={font=\scriptsize},ylabel shift=-3pt]\addplot[mark=*,mark size=0.7pt,draw=black,thin] coordinates {(0,388) (1,1666) (2,10156) (3,55809) (4,179320) (5,384974) (6,324034) (7,459356) (8,464858) (9,605160) (10,519436) (11,348652) (12,95190) (13,46563) (14,16500) (15,3645) (16,520) (17,46)};\end{axis}\end{tikzpicture}\end{minipage}
\vspace{3mm}

\noindent\begin{minipage}[t]{60mm}\vspace{0pt}\sbox\hmboard{\chessboard[setfen=6k1/8/8/1p6/8/8/1r6/K6N w - - 0 1]}\usebox\hmboard\\[1mm]
\makebox[\wd\hmboard]{\sffamily\hspace{3mm}\textbf{h\#8.5}\hfill(2+3)}\end{minipage}\hfill
\begin{minipage}[t]{52mm}\vspace{0pt}
\textcolor{ink}{\sffamily\textsc{Lösung}}\par\smallskip\small
\begin{tabular}{@{}r@{\enspace}l@{\enspace}l@{}}
\textbf{1.} & \ldots & \symknight{}g3\\
\textbf{2.} & \symrook{}h2 & \symking{}b1\\
\textbf{3.} & \symrook{}h8 & \symking{}c2\\
\textbf{4.} & b4 & \symking{}d3\\
\textbf{5.} & b3 & \symking{}e4\\
\textbf{6.} & b2 & \symking{}f5\\
\textbf{7.} & b1=\symrook{} & \symking{}g6\\
\textbf{8.} & \symrook{}b8 & \symknight{}f5\\
\textbf{9.} & \symrook{}f8 & \symknight{}e7\#\\
\end{tabular}\par\vspace{3mm}
\qrcode[height=12mm]{https://helpman.komtera.lt/?fen=6k1/8/8/1p6/8/8/1r6/K6N&moves=8.5}\\{\sffamily\tiny\color{ink}scannen zum Lösen}
\end{minipage}
\vspace{2mm}

{\raggedright\noindent\textcolor{ink}{\sffamily\textsc{Themen}}\enspace \themetag{reines Matt} \themetag{Mustermatt} \themetag{Idealmatt} \themetag{Umwandlung} \themetag{Unterverwandlung} \themetag{Selbstblock} \themetag{ohne Schlag} \themetag{ohne Schach}\par}

\clearpage
\phantomsection\label{hm:KNvkrr}
\begin{tcolorbox}[enhanced,colback=ink,colframe=ink,arc=0pt,boxrule=0pt,left=4mm,right=4mm,top=2.5mm,bottom=2.5mm]
\color{white}\sffamily{\LARGE\bfseries KNvkrr}\hfill{\Large h\#8.5}\\[0.5mm]
{\small 5 Steine\quad Weiß am Zug}\hfill{\small Nr.~67\qquad\href{https://helpman.komtera.lt/?fen=8/7k/N7/8/8/8/4r3/Kr6&moves=8.5}{\color{white}\bfseries lösen\,\ensuremath{\nearrow}}}
\end{tcolorbox}
\vspace{2mm}
\noindent\begin{minipage}[t]{36mm}\vspace{0pt}\sffamily\small\raggedright
\hmfigure{längste eindeutige}{{\Large h\#8.5}}\hmfigure{eindeutig bei dieser Länge}{{\Large 5} Stellungen}\hmfigure{längstes Matt}{{\Large h\#9.5}}\hmfigure{in der Klasse}{242 Millionen Stellungen,\\ 78,4 Millionen mit Hilfsmatt}
\end{minipage}\hfill
\begin{minipage}[t]{80mm}\vspace{0pt}\begin{tikzpicture}\begin{axis}[width=72mm,height=42mm,ybar,bar width=1.6pt,xmin=-1,xmax=20,ymin=0,ymax=20.4,axis y line*=left,axis x line*=bottom,xlabel={\scriptsize Halbzüge bis Matt},ylabel={\scriptsize Stellungen (Mio.)},tick label style={font=\tiny},label style={font=\scriptsize},enlarge x limits=false,ylabel shift=-3pt,xlabel shift=-3pt]\addplot[fill=ink!55,draw=none] coordinates {(0,4.5e-05) (1,0.000184) (2,0.002122) (3,0.01454) (4,0.07807) (5,0.3829) (6,1.263) (7,4.236) (8,7.808) (9,15.78) (10,13.27) (11,18.54) (12,7.978) (13,7.022) (14,1.376) (15,0.6046) (16,0.03368) (17,0.03848) (19,0.001425)};\draw[densely dotted] (axis cs:17,0) -- (axis cs:17,20.4);\end{axis}\begin{axis}[width=72mm,height=42mm,axis y line*=right,axis x line=none,ymode=log,xmin=-1,xmax=20,ylabel={\scriptsize eindeutig (log)},tick label style={font=\tiny},label style={font=\scriptsize},ylabel shift=-3pt]\addplot[mark=*,mark size=0.7pt,draw=black,thin] coordinates {(0,45) (1,184) (2,1880) (3,10574) (4,25628) (5,84965) (6,111776) (7,205185) (8,170475) (9,200841) (10,71811) (11,69563) (12,8936) (13,5762) (14,152) (15,126) (17,5)};\end{axis}\end{tikzpicture}\end{minipage}
\vspace{3mm}

\noindent\begin{minipage}[t]{60mm}\vspace{0pt}\sbox\hmboard{\chessboard[setfen=8/7k/N7/8/8/8/4r3/Kr6 w - - 0 1]}\usebox\hmboard\\[1mm]
\makebox[\wd\hmboard]{\sffamily\hspace{3mm}\textbf{h\#8.5}\hfill(2+3)}\end{minipage}\hfill
\begin{minipage}[t]{52mm}\vspace{0pt}
\textcolor{ink}{\sffamily\textsc{Lösung}}\par\smallskip\small
\begin{tabular}{@{}r@{\enspace}l@{\enspace}l@{}}
\textbf{1.} & \ldots & \symking{}xb1\\
\textbf{2.} & \symking{}g6 & \symknight{}b4\\
\textbf{3.} & \symking{}f5 & \symknight{}c2\\
\textbf{4.} & \symking{}e4 & \symking{}b2\\
\textbf{5.} & \symking{}d3 & \symking{}c1\\
\textbf{6.} & \symking{}c3 & \symking{}d1\\
\textbf{7.} & \symking{}b2 & \symknight{}d4\\
\textbf{8.} & \symking{}a1 & \symking{}c1\\
\textbf{9.} & \symrook{}a2 & \symknight{}b3\#\\
\end{tabular}\par\vspace{3mm}
\qrcode[height=12mm]{https://helpman.komtera.lt/?fen=8/7k/N7/8/8/8/4r3/Kr6&moves=8.5}\\{\sffamily\tiny\color{ink}scannen zum Lösen}
\end{minipage}
\vspace{2mm}

{\raggedright\noindent\textcolor{ink}{\sffamily\textsc{Themen}}\enspace \themetag{reines Matt} \themetag{Mustermatt} \themetag{Idealmatt} \themetag{Switchback} \themetag{Selbstblock} \themetag{ohne Schach}\par\smallskip{\footnotesize\textit{Hinweis:} die Partei am Zug steht im Diagramm im Schach; die Lösung beginnt mit einem Schlag.}\par}

\clearpage
\phantomsection\label{hm:KPPvkp}
\begin{tcolorbox}[enhanced,colback=ink,colframe=ink,arc=0pt,boxrule=0pt,left=4mm,right=4mm,top=2.5mm,bottom=2.5mm]
\color{white}\sffamily{\LARGE\bfseries KPPvkp}\hfill{\Large h\#8.5}\\[0.5mm]
{\small 5 Steine\quad Weiß am Zug}\hfill{\small Nr.~68\qquad\href{https://helpman.komtera.lt/?fen=8/8/p7/8/P7/P7/K7/7k&moves=8.5}{\color{white}\bfseries lösen\,\ensuremath{\nearrow}}}
\end{tcolorbox}
\vspace{2mm}
\noindent\begin{minipage}[t]{36mm}\vspace{0pt}\sffamily\small\raggedright
\hmfigure{längste eindeutige}{{\Large h\#8.5}}\hmfigure{eindeutig bei dieser Länge}{{\Large 6} Stellungen}\hmfigure{längstes Matt}{{\Large h\#14.5}}\hmfigure{in der Klasse}{399 Millionen Stellungen,\\ 163 Millionen mit Hilfsmatt}
\end{minipage}\hfill
\begin{minipage}[t]{80mm}\vspace{0pt}\begin{tikzpicture}\begin{axis}[width=72mm,height=42mm,ybar,bar width=1.6pt,xmin=-1,xmax=30,ymin=0,ymax=28.3,axis y line*=left,axis x line*=bottom,xlabel={\scriptsize Halbzüge bis Matt},ylabel={\scriptsize Stellungen (Mio.)},tick label style={font=\tiny},label style={font=\scriptsize},enlarge x limits=false,ylabel shift=-3pt,xlabel shift=-3pt]\addplot[fill=ink!55,draw=none] coordinates {(0,0.000637) (1,0.08735) (2,0.1933) (3,1.367) (4,3.303) (5,11.87) (6,18.26) (7,24.65) (8,25.73) (9,22.71) (10,19.48) (11,14.57) (12,12.33) (13,4.798) (14,3.224) (15,0.2641) (16,0.1159) (17,0.05577) (18,0.057) (19,0.05212) (20,0.053) (21,0.03198) (22,0.03222) (23,0.0118) (24,0.01186) (25,0.005907) (26,0.006128) (27,0.001742) (28,0.001776) (29,1e-05)};\draw[densely dotted] (axis cs:17,0) -- (axis cs:17,28.3);\end{axis}\begin{axis}[width=72mm,height=42mm,axis y line*=right,axis x line=none,ymode=log,xmin=-1,xmax=30,ylabel={\scriptsize eindeutig (log)},tick label style={font=\tiny},label style={font=\scriptsize},ylabel shift=-3pt]\addplot[mark=*,mark size=0.7pt,draw=black,thin] coordinates {(0,637) (1,18997) (2,44547) (3,453168) (4,672505) (5,929670) (6,609775) (7,387675) (8,224503) (9,153344) (10,95655) (11,66120) (12,36404) (13,1371) (14,209) (15,21) (16,10) (17,6)};\end{axis}\end{tikzpicture}\end{minipage}
\vspace{3mm}

\noindent\begin{minipage}[t]{60mm}\vspace{0pt}\sbox\hmboard{\chessboard[setfen=8/8/p7/8/P7/P7/K7/7k w - - 0 1]}\usebox\hmboard\\[1mm]
\makebox[\wd\hmboard]{\sffamily\hspace{3mm}\textbf{h\#8.5}\hfill(3+2)}\end{minipage}\hfill
\begin{minipage}[t]{52mm}\vspace{0pt}
\textcolor{ink}{\sffamily\textsc{Lösung}}\par\smallskip\small
\begin{tabular}{@{}r@{\enspace}l@{\enspace}l@{}}
\textbf{1.} & \ldots & \symking{}b3\\
\textbf{2.} & \symking{}g2 & \symking{}b4\\
\textbf{3.} & a5+ & \symking{}xa5\\
\textbf{4.} & \symking{}f3 & \symking{}b4\\
\textbf{5.} & \symking{}e4 & a5\\
\textbf{6.} & \symking{}d5 & a6\\
\textbf{7.} & \symking{}c6 & a7\\
\textbf{8.} & \symking{}b7 & \symking{}c5\\
\textbf{9.} & \symking{}a6 & a8=\symqueen{}\#\\
\end{tabular}\par\vspace{3mm}
\qrcode[height=12mm]{https://helpman.komtera.lt/?fen=8/8/p7/8/P7/P7/K7/7k&moves=8.5}\\{\sffamily\tiny\color{ink}scannen zum Lösen}
\end{minipage}
\vspace{2mm}

{\raggedright\noindent\textcolor{ink}{\sffamily\textsc{Themen}}\enspace \themetag{reines Matt} \themetag{Mustermatt} \themetag{Spiegelmatt} \themetag{Umwandlung} \themetag{Switchback}\par}

\clearpage
\phantomsection\label{hm:KPvkqn}
\begin{tcolorbox}[enhanced,colback=ink,colframe=ink,arc=0pt,boxrule=0pt,left=4mm,right=4mm,top=2.5mm,bottom=2.5mm]
\color{white}\sffamily{\LARGE\bfseries KPvkqn}\hfill{\Large h\#8.5}\\[0.5mm]
{\small 5 Steine\quad Weiß am Zug}\hfill{\small Nr.~69\qquad\href{https://helpman.komtera.lt/?fen=K7/8/8/8/8/2n5/2P5/2q4k&moves=8.5}{\color{white}\bfseries lösen\,\ensuremath{\nearrow}}}
\end{tcolorbox}
\vspace{2mm}
\noindent\begin{minipage}[t]{36mm}\vspace{0pt}\sffamily\small\raggedright
\hmfigure{längste eindeutige}{{\Large h\#8.5}}\hmfigure{eindeutig bei dieser Länge}{{\Large 1} Stellung}\hmfigure{längstes Matt}{{\Large h\#9.5}}\hmfigure{in der Klasse}{710 Millionen Stellungen,\\ 487 Millionen mit Hilfsmatt}
\end{minipage}\hfill
\begin{minipage}[t]{80mm}\vspace{0pt}\begin{tikzpicture}\begin{axis}[width=72mm,height=42mm,ybar,bar width=1.6pt,xmin=-1,xmax=20,ymin=0,ymax=93.03,axis y line*=left,axis x line*=bottom,xlabel={\scriptsize Halbzüge bis Matt},ylabel={\scriptsize Stellungen (Mio.)},tick label style={font=\tiny},label style={font=\scriptsize},enlarge x limits=false,ylabel shift=-3pt,xlabel shift=-3pt]\addplot[fill=ink!55,draw=none] coordinates {(0,5.6e-05) (1,0.0269) (2,0.1155) (3,1.041) (4,2.39) (5,11.67) (6,15.65) (7,42.8) (8,36.76) (9,71.95) (10,47.78) (11,84.57) (12,55.82) (13,72.68) (14,29.47) (15,12.74) (16,1.043) (17,0.391) (18,0.001644) (19,0.000343)};\draw[densely dotted] (axis cs:17,0) -- (axis cs:17,93.03);\end{axis}\begin{axis}[width=72mm,height=42mm,axis y line*=right,axis x line=none,ymode=log,xmin=-1,xmax=20,ylabel={\scriptsize eindeutig (log)},tick label style={font=\tiny},label style={font=\scriptsize},ylabel shift=-3pt]\addplot[mark=*,mark size=0.7pt,draw=black,thin] coordinates {(0,56) (1,6625) (2,16494) (3,224517) (4,318374) (5,1051043) (6,449779) (7,813788) (8,272967) (9,495730) (10,195409) (11,388284) (12,134292) (13,126331) (14,5531) (15,2569) (16,4) (17,1)};\end{axis}\end{tikzpicture}\end{minipage}
\vspace{3mm}

\noindent\begin{minipage}[t]{60mm}\vspace{0pt}\sbox\hmboard{\chessboard[setfen=K7/8/8/8/8/2n5/2P5/2q4k w - - 0 1]}\usebox\hmboard\\[1mm]
\makebox[\wd\hmboard]{\sffamily\hspace{3mm}\textbf{h\#8.5}\hfill(2+3)}\end{minipage}\hfill
\begin{minipage}[t]{52mm}\vspace{0pt}
\textcolor{ink}{\sffamily\textsc{Lösung}}\par\smallskip\small
\begin{tabular}{@{}r@{\enspace}l@{\enspace}l@{}}
\textbf{1.} & \ldots & \symking{}a7\\
\textbf{2.} & \symknight{}d5 & c4\\
\textbf{3.} & \symking{}g2 & cxd5\\
\textbf{4.} & \symking{}f3 & d6\\
\textbf{5.} & \symking{}e4 & d7\\
\textbf{6.} & \symking{}d5 & d8=\symqueen{}+\\
\textbf{7.} & \symking{}c6 & \symqueen{}d5+\\
\textbf{8.} & \symking{}c7 & \symking{}a6\\
\textbf{9.} & \symking{}b8 & \symqueen{}b7\#\\
\end{tabular}\par\vspace{3mm}
\qrcode[height=12mm]{https://helpman.komtera.lt/?fen=K7/8/8/8/8/2n5/2P5/2q4k&moves=8.5}\\{\sffamily\tiny\color{ink}scannen zum Lösen}
\end{minipage}
\vspace{2mm}

{\raggedright\noindent\textcolor{ink}{\sffamily\textsc{Themen}}\enspace \themetag{Umwandlung} \themetag{Excelsior} \themetag{Rundlauf} \themetag{Umnov}\par}

\clearpage
\phantomsection\label{hm:KPvkqp}
\begin{tcolorbox}[enhanced,colback=ink,colframe=ink,arc=0pt,boxrule=0pt,left=4mm,right=4mm,top=2.5mm,bottom=2.5mm]
\color{white}\sffamily{\LARGE\bfseries KPvkqp}\hfill{\Large h\#8.5}\\[0.5mm]
{\small 5 Steine\quad Weiß am Zug}\hfill{\small Nr.~70\qquad\href{https://helpman.komtera.lt/?fen=6kq/8/8/8/8/p7/P7/1K6&moves=8.5}{\color{white}\bfseries lösen\,\ensuremath{\nearrow}}}
\end{tcolorbox}
\vspace{2mm}
\noindent\begin{minipage}[t]{36mm}\vspace{0pt}\sffamily\small\raggedright
\hmfigure{längste eindeutige}{{\Large h\#8.5}}\hmfigure{eindeutig bei dieser Länge}{{\Large 9} Stellungen}\hmfigure{längstes Matt}{{\Large h\#14.5}}\hmfigure{in der Klasse}{533 Millionen Stellungen,\\ 372 Millionen mit Hilfsmatt}
\end{minipage}\hfill
\begin{minipage}[t]{80mm}\vspace{0pt}\begin{tikzpicture}\begin{axis}[width=72mm,height=42mm,ybar,bar width=1.6pt,xmin=-1,xmax=30,ymin=0,ymax=70.62,axis y line*=left,axis x line*=bottom,xlabel={\scriptsize Halbzüge bis Matt},ylabel={\scriptsize Stellungen (Mio.)},tick label style={font=\tiny},label style={font=\scriptsize},enlarge x limits=false,ylabel shift=-3pt,xlabel shift=-3pt]\addplot[fill=ink!55,draw=none] coordinates {(0,5.9e-05) (1,0.02905) (2,0.1189) (3,0.9855) (4,2.345) (5,10.95) (6,15.79) (7,38.02) (8,32.61) (9,56.36) (10,38.38) (11,64.2) (12,43.64) (13,45.66) (14,15.39) (15,5.711) (16,0.9341) (17,0.5077) (18,0.03676) (19,0.04294) (20,6.8e-05) (21,0.03471) (23,0.0175) (25,0.01151) (27,0.00831) (29,0.001656)};\draw[densely dotted] (axis cs:17,0) -- (axis cs:17,70.62);\end{axis}\begin{axis}[width=72mm,height=42mm,axis y line*=right,axis x line=none,ymode=log,xmin=-1,xmax=30,ylabel={\scriptsize eindeutig (log)},tick label style={font=\tiny},label style={font=\scriptsize},ylabel shift=-3pt]\addplot[mark=*,mark size=0.7pt,draw=black,thin] coordinates {(0,59) (1,7807) (2,21494) (3,230188) (4,358866) (5,996904) (6,450724) (7,669347) (8,263446) (9,418348) (10,191492) (11,313129) (12,122834) (13,60123) (14,3470) (15,1739) (16,23) (17,9)};\end{axis}\end{tikzpicture}\end{minipage}
\vspace{3mm}

\noindent\begin{minipage}[t]{60mm}\vspace{0pt}\sbox\hmboard{\chessboard[setfen=6kq/8/8/8/8/p7/P7/1K6 w - - 0 1]}\usebox\hmboard\\[1mm]
\makebox[\wd\hmboard]{\sffamily\hspace{3mm}\textbf{h\#8.5}\hfill(2+3)}\end{minipage}\hfill
\begin{minipage}[t]{52mm}\vspace{0pt}
\textcolor{ink}{\sffamily\textsc{Lösung}}\par\smallskip\small
\begin{tabular}{@{}r@{\enspace}l@{\enspace}l@{}}
\textbf{1.} & \ldots & \symking{}c2\\
\textbf{2.} & \symqueen{}h3 & \symking{}c1\\
\textbf{3.} & \symqueen{}b3 & axb3\\
\textbf{4.} & \symking{}f7 & b4\\
\textbf{5.} & \symking{}e6 & b5\\
\textbf{6.} & \symking{}d5 & b6\\
\textbf{7.} & \symking{}c4 & b7\\
\textbf{8.} & \symking{}b3 & b8=\symqueen{}+\\
\textbf{9.} & \symking{}a2 & \symqueen{}b1\#\\
\end{tabular}\par\vspace{3mm}
\qrcode[height=12mm]{https://helpman.komtera.lt/?fen=6kq/8/8/8/8/p7/P7/1K6&moves=8.5}\\{\sffamily\tiny\color{ink}scannen zum Lösen}
\end{minipage}
\vspace{2mm}

{\raggedright\noindent\textcolor{ink}{\sffamily\textsc{Themen}}\enspace \themetag{Satzspiel} \themetag{Umwandlung} \themetag{Excelsior}\par}

\clearpage
\phantomsection\label{hm:KPvkrp}
\begin{tcolorbox}[enhanced,colback=ink,colframe=ink,arc=0pt,boxrule=0pt,left=4mm,right=4mm,top=2.5mm,bottom=2.5mm]
\color{white}\sffamily{\LARGE\bfseries KPvkrp}\hfill{\Large h\#8.5}\\[0.5mm]
{\small 5 Steine\quad Weiß am Zug}\hfill{\small Nr.~71\qquad\href{https://helpman.komtera.lt/?fen=8/8/8/8/6k1/5p2/2r2P2/K7&moves=8.5}{\color{white}\bfseries lösen\,\ensuremath{\nearrow}}}
\end{tcolorbox}
\vspace{2mm}
\noindent\begin{minipage}[t]{36mm}\vspace{0pt}\sffamily\small\raggedright
\hmfigure{längste eindeutige}{{\Large h\#8.5}}\hmfigure{eindeutig bei dieser Länge}{{\Large 101} Stellungen}\hmfigure{längstes Matt}{{\Large h\#14.5}}\hmfigure{in der Klasse}{533 Millionen Stellungen,\\ 401 Millionen mit Hilfsmatt}
\end{minipage}\hfill
\begin{minipage}[t]{80mm}\vspace{0pt}\begin{tikzpicture}\begin{axis}[width=72mm,height=42mm,ybar,bar width=1.6pt,xmin=-1,xmax=30,ymin=0,ymax=67.97,axis y line*=left,axis x line*=bottom,xlabel={\scriptsize Halbzüge bis Matt},ylabel={\scriptsize Stellungen (Mio.)},tick label style={font=\tiny},label style={font=\scriptsize},enlarge x limits=false,ylabel shift=-3pt,xlabel shift=-3pt]\addplot[fill=ink!55,draw=none] coordinates {(0,8.2e-05) (1,0.05184) (2,0.1873) (3,1.377) (4,3.906) (5,15.57) (6,26.1) (7,48.01) (8,42.99) (9,59.41) (10,43.61) (11,61.79) (12,48.11) (13,32.51) (14,11.88) (15,3.52) (16,1.216) (17,0.5721) (18,0.06091) (19,0.01607) (20,0.00015) (21,0.0045) (22,5e-05) (23,0.002614) (24,3.2e-05) (25,0.001052) (26,1e-05) (27,0.000742) (28,4e-06) (29,0.000296)};\draw[densely dotted] (axis cs:17,0) -- (axis cs:17,67.97);\end{axis}\begin{axis}[width=72mm,height=42mm,axis y line*=right,axis x line=none,ymode=log,xmin=-1,xmax=30,ylabel={\scriptsize eindeutig (log)},tick label style={font=\tiny},label style={font=\scriptsize},ylabel shift=-3pt]\addplot[mark=*,mark size=0.7pt,draw=black,thin] coordinates {(0,82) (1,14298) (2,52506) (3,391072) (4,860155) (5,1529414) (6,765029) (7,815257) (8,343989) (9,436960) (10,221859) (11,247427) (12,93153) (13,27683) (14,6915) (15,2724) (16,279) (17,101)};\end{axis}\end{tikzpicture}\end{minipage}
\vspace{3mm}

\noindent\begin{minipage}[t]{60mm}\vspace{0pt}\sbox\hmboard{\chessboard[setfen=8/8/8/8/6k1/5p2/2r2P2/K7 w - - 0 1]}\usebox\hmboard\\[1mm]
\makebox[\wd\hmboard]{\sffamily\hspace{3mm}\textbf{h\#8.5}\hfill(2+3)}\end{minipage}\hfill
\begin{minipage}[t]{52mm}\vspace{0pt}
\textcolor{ink}{\sffamily\textsc{Lösung}}\par\smallskip\small
\begin{tabular}{@{}r@{\enspace}l@{\enspace}l@{}}
\textbf{1.} & \ldots & \symking{}b1\\
\textbf{2.} & \symrook{}e2 & \symking{}c1\\
\textbf{3.} & \symrook{}e3 & fxe3\\
\textbf{4.} & \symking{}f5 & e4+\\
\textbf{5.} & \symking{}e6 & e5\\
\textbf{6.} & \symking{}d5 & e6\\
\textbf{7.} & \symking{}c4 & e7\\
\textbf{8.} & \symking{}b3 & e8=\symqueen{}\\
\textbf{9.} & \symking{}a2 & \symqueen{}a4\#\\
\end{tabular}\par\vspace{3mm}
\qrcode[height=12mm]{https://helpman.komtera.lt/?fen=8/8/8/8/6k1/5p2/2r2P2/K7&moves=8.5}\\{\sffamily\tiny\color{ink}scannen zum Lösen}
\end{minipage}
\vspace{2mm}

{\raggedright\noindent\textcolor{ink}{\sffamily\textsc{Themen}}\enspace \themetag{reines Matt} \themetag{Mustermatt} \themetag{Spiegelmatt} \themetag{Umwandlung} \themetag{Excelsior} \themetag{Umnov}\par}

\clearpage
\phantomsection\label{hm:KBvkbb}
\begin{tcolorbox}[enhanced,colback=ink,colframe=ink,arc=0pt,boxrule=0pt,left=4mm,right=4mm,top=2.5mm,bottom=2.5mm]
\color{white}\sffamily{\LARGE\bfseries KBvkbb}\hfill{\Large h\#8}\\[0.5mm]
{\small 5 Steine\quad Schwarz am Zug}\hfill{\small Nr.~72\qquad\href{https://helpman.komtera.lt/?fen=5b2/8/8/8/7B/7b/7k/K7&moves=8}{\color{white}\bfseries lösen\,\ensuremath{\nearrow}}}
\end{tcolorbox}
\vspace{2mm}
\noindent\begin{minipage}[t]{36mm}\vspace{0pt}\sffamily\small\raggedright
\hmfigure{längste eindeutige}{{\Large h\#8}}\hmfigure{eindeutig bei dieser Länge}{{\Large 8} Stellungen}\hmfigure{längstes Matt}{{\Large h\#10}}\hmfigure{in der Klasse}{242 Millionen Stellungen,\\ 62,7 Millionen mit Hilfsmatt}
\end{minipage}\hfill
\begin{minipage}[t]{80mm}\vspace{0pt}\begin{tikzpicture}\begin{axis}[width=72mm,height=42mm,ybar,bar width=1.6pt,xmin=-1,xmax=21,ymin=0,ymax=10.75,axis y line*=left,axis x line*=bottom,xlabel={\scriptsize Halbzüge bis Matt},ylabel={\scriptsize Stellungen (Mio.)},tick label style={font=\tiny},label style={font=\scriptsize},enlarge x limits=false,ylabel shift=-3pt,xlabel shift=-3pt]\addplot[fill=ink!55,draw=none] coordinates {(0,0.000646) (1,0.002456) (2,0.02021) (3,0.06606) (4,0.2935) (5,0.7282) (6,1.811) (7,3.482) (8,5.309) (9,8.07) (10,7.755) (11,9.772) (12,7.897) (13,6.206) (14,3.894) (15,3.394) (16,1.515) (17,1.578) (18,0.3821) (19,0.4826) (20,0.000189)};\draw[densely dotted] (axis cs:16,0) -- (axis cs:16,10.75);\end{axis}\begin{axis}[width=72mm,height=42mm,axis y line*=right,axis x line=none,ymode=log,xmin=-1,xmax=21,ylabel={\scriptsize eindeutig (log)},tick label style={font=\tiny},label style={font=\scriptsize},ylabel shift=-3pt]\addplot[mark=*,mark size=0.7pt,draw=black,thin] coordinates {(0,646) (1,2431) (2,17722) (3,26679) (4,62973) (5,93051) (6,55879) (7,65542) (8,39329) (9,44969) (10,26151) (11,16965) (12,5353) (13,1442) (14,444) (15,18) (16,8)};\end{axis}\end{tikzpicture}\end{minipage}
\vspace{3mm}

\noindent\begin{minipage}[t]{60mm}\vspace{0pt}\sbox\hmboard{\chessboard[setfen=5b2/8/8/8/7B/7b/7k/K7 b - - 0 1]}\usebox\hmboard\\[1mm]
\makebox[\wd\hmboard]{\sffamily\hspace{3mm}\textbf{h\#8}\hfill(2+3)}\end{minipage}\hfill
\begin{minipage}[t]{52mm}\vspace{0pt}
\textcolor{ink}{\sffamily\textsc{Lösung}}\par\smallskip\small
\begin{tabular}{@{}r@{\enspace}l@{\enspace}l@{}}
\textbf{1.} & \symbishop{}f5 & \symking{}b2\\
\textbf{2.} & \symking{}h3 & \symking{}c3\\
\textbf{3.} & \symking{}g4 & \symking{}d4\\
\textbf{4.} & \symking{}h5 & \symking{}e5\\
\textbf{5.} & \symking{}h6 & \symking{}f6\\
\textbf{6.} & \symking{}h7 & \symking{}f7\\
\textbf{7.} & \symking{}h8 & \symking{}xf8\\
\textbf{8.} & \symbishop{}h7 & \symbishop{}f6\#\\
\end{tabular}\par\vspace{3mm}
\qrcode[height=12mm]{https://helpman.komtera.lt/?fen=5b2/8/8/8/7B/7b/7k/K7&moves=8}\\{\sffamily\tiny\color{ink}scannen zum Lösen}
\end{minipage}
\vspace{2mm}

{\raggedright\noindent\textcolor{ink}{\sffamily\textsc{Themen}}\enspace \themetag{Selbstblock} \themetag{ohne Schach}\par}

\clearpage
\phantomsection\label{hm:KNvkbn}
\begin{tcolorbox}[enhanced,colback=ink,colframe=ink,arc=0pt,boxrule=0pt,left=4mm,right=4mm,top=2.5mm,bottom=2.5mm]
\color{white}\sffamily{\LARGE\bfseries KNvkbn}\hfill{\Large h\#8}\\[0.5mm]
{\small 5 Steine\quad Schwarz am Zug}\hfill{\small Nr.~73\qquad\href{https://helpman.komtera.lt/?fen=7k/8/8/8/8/1b6/8/Kn2N3&moves=8}{\color{white}\bfseries lösen\,\ensuremath{\nearrow}}}
\end{tcolorbox}
\vspace{2mm}
\noindent\begin{minipage}[t]{36mm}\vspace{0pt}\sffamily\small\raggedright
\hmfigure{längste eindeutige}{{\Large h\#8}}\hmfigure{eindeutig bei dieser Länge}{{\Large 5} Stellungen}\hmfigure{längstes Matt}{{\Large h\#10}}\hmfigure{in der Klasse}{242 Millionen Stellungen,\\ 176 Millionen mit Hilfsmatt}
\end{minipage}\hfill
\begin{minipage}[t]{80mm}\vspace{0pt}\begin{tikzpicture}\begin{axis}[width=72mm,height=42mm,ybar,bar width=1.6pt,xmin=-1,xmax=21,ymin=0,ymax=39.88,axis y line*=left,axis x line*=bottom,xlabel={\scriptsize Halbzüge bis Matt},ylabel={\scriptsize Stellungen (Mio.)},tick label style={font=\tiny},label style={font=\scriptsize},enlarge x limits=false,ylabel shift=-3pt,xlabel shift=-3pt]\addplot[fill=ink!55,draw=none] coordinates {(0,0.000104) (1,0.000533) (2,0.002946) (3,0.01983) (4,0.09845) (5,0.4451) (6,1.264) (7,3.786) (8,8.327) (9,16.95) (10,25.35) (11,36.25) (12,31.51) (13,30) (14,13.66) (15,6.322) (16,1.605) (17,0.2337) (18,0.0293) (19,0.000754) (20,2e-06)};\draw[densely dotted] (axis cs:16,0) -- (axis cs:16,39.88);\end{axis}\begin{axis}[width=72mm,height=42mm,axis y line*=right,axis x line=none,ymode=log,xmin=-1,xmax=21,ylabel={\scriptsize eindeutig (log)},tick label style={font=\tiny},label style={font=\scriptsize},ylabel shift=-3pt]\addplot[mark=*,mark size=0.7pt,draw=black,thin] coordinates {(0,104) (1,533) (2,2416) (3,12297) (4,51108) (5,90572) (6,94522) (7,111312) (8,132497) (9,144264) (10,109514) (11,69265) (12,21505) (13,6215) (14,924) (15,42) (16,5)};\end{axis}\end{tikzpicture}\end{minipage}
\vspace{3mm}

\noindent\begin{minipage}[t]{60mm}\vspace{0pt}\sbox\hmboard{\chessboard[setfen=7k/8/8/8/8/1b6/8/Kn2N3 b - - 0 1]}\usebox\hmboard\\[1mm]
\makebox[\wd\hmboard]{\sffamily\hspace{3mm}\textbf{h\#8}\hfill(2+3)}\end{minipage}\hfill
\begin{minipage}[t]{52mm}\vspace{0pt}
\textcolor{ink}{\sffamily\textsc{Lösung}}\par\smallskip\small
\begin{tabular}{@{}r@{\enspace}l@{\enspace}l@{}}
\textbf{1.} & \symking{}g7 & \symking{}xb1\\
\textbf{2.} & \symking{}f6 & \symking{}c1\\
\textbf{3.} & \symking{}e5 & \symking{}d2\\
\textbf{4.} & \symking{}d4 & \symking{}e2\\
\textbf{5.} & \symking{}c3 & \symking{}e3\\
\textbf{6.} & \symking{}b2 & \symking{}d2\\
\textbf{7.} & \symking{}a1 & \symking{}c1\\
\textbf{8.} & \symbishop{}a2 & \symknight{}c2\#\\
\end{tabular}\par\vspace{3mm}
\qrcode[height=12mm]{https://helpman.komtera.lt/?fen=7k/8/8/8/8/1b6/8/Kn2N3&moves=8}\\{\sffamily\tiny\color{ink}scannen zum Lösen}
\end{minipage}
\vspace{2mm}

{\raggedright\noindent\textcolor{ink}{\sffamily\textsc{Themen}}\enspace \themetag{reines Matt} \themetag{Mustermatt} \themetag{Idealmatt} \themetag{Rundlauf} \themetag{Selbstblock} \themetag{ohne Schach}\par}

\clearpage
\phantomsection\label{hm:KNvkqn}
\begin{tcolorbox}[enhanced,colback=ink,colframe=ink,arc=0pt,boxrule=0pt,left=4mm,right=4mm,top=2.5mm,bottom=2.5mm]
\color{white}\sffamily{\LARGE\bfseries KNvkqn}\hfill{\Large h\#8}\\[0.5mm]
{\small 5 Steine\quad Schwarz am Zug}\hfill{\small Nr.~74\qquad\href{https://helpman.komtera.lt/?fen=8/7k/8/8/8/N7/1nq5/K7&moves=8}{\color{white}\bfseries lösen\,\ensuremath{\nearrow}}}
\end{tcolorbox}
\vspace{2mm}
\noindent\begin{minipage}[t]{36mm}\vspace{0pt}\sffamily\small\raggedright
\hmfigure{längste eindeutige}{{\Large h\#8}}\hmfigure{eindeutig bei dieser Länge}{{\Large 2} Stellungen}\hmfigure{längstes Matt}{{\Large h\#10.5}}\hmfigure{in der Klasse}{242 Millionen Stellungen,\\ 156 Millionen mit Hilfsmatt}
\end{minipage}\hfill
\begin{minipage}[t]{80mm}\vspace{0pt}\begin{tikzpicture}\begin{axis}[width=72mm,height=42mm,ybar,bar width=1.6pt,xmin=-1,xmax=22,ymin=0,ymax=40.27,axis y line*=left,axis x line*=bottom,xlabel={\scriptsize Halbzüge bis Matt},ylabel={\scriptsize Stellungen (Mio.)},tick label style={font=\tiny},label style={font=\scriptsize},enlarge x limits=false,ylabel shift=-3pt,xlabel shift=-3pt]\addplot[fill=ink!55,draw=none] coordinates {(0,3.1e-05) (1,0.000166) (2,0.000926) (3,0.007424) (4,0.04409) (5,0.237) (6,0.9774) (7,3.449) (8,7.941) (9,18.6) (10,21.71) (11,36.61) (12,23.24) (13,28.99) (14,8.208) (15,5.557) (16,0.6365) (17,0.2696) (18,0.002897) (19,0.01686) (21,0.000194)};\draw[densely dotted] (axis cs:16,0) -- (axis cs:16,40.27);\end{axis}\begin{axis}[width=72mm,height=42mm,axis y line*=right,axis x line=none,ymode=log,xmin=-1,xmax=22,ylabel={\scriptsize eindeutig (log)},tick label style={font=\tiny},label style={font=\scriptsize},ylabel shift=-3pt]\addplot[mark=*,mark size=0.7pt,draw=black,thin] coordinates {(0,31) (1,166) (2,726) (3,4560) (4,18400) (5,51151) (6,84031) (7,138022) (8,95342) (9,123322) (10,48013) (11,61889) (12,12446) (13,11258) (14,650) (15,170) (16,2)};\end{axis}\end{tikzpicture}\end{minipage}
\vspace{3mm}

\noindent\begin{minipage}[t]{60mm}\vspace{0pt}\sbox\hmboard{\chessboard[setfen=8/7k/8/8/8/N7/1nq5/K7 b - - 0 1]}\usebox\hmboard\\[1mm]
\makebox[\wd\hmboard]{\sffamily\hspace{3mm}\textbf{h\#8}\hfill(2+3)}\end{minipage}\hfill
\begin{minipage}[t]{52mm}\vspace{0pt}
\textcolor{ink}{\sffamily\textsc{Lösung}}\par\smallskip\small
\begin{tabular}{@{}r@{\enspace}l@{\enspace}l@{}}
\textbf{1.} & \symqueen{}a4 & \symking{}b1\\
\textbf{2.} & \symking{}g6 & \symking{}c1\\
\textbf{3.} & \symking{}f5 & \symking{}d2\\
\textbf{4.} & \symking{}e4 & \symking{}e1\\
\textbf{5.} & \symking{}d3 & \symknight{}b5\\
\textbf{6.} & \symking{}c2 & \symking{}e2\\
\textbf{7.} & \symking{}b1 & \symking{}d2\\
\textbf{8.} & \symqueen{}a1 & \symknight{}c3\#\\
\end{tabular}\par\vspace{3mm}
\qrcode[height=12mm]{https://helpman.komtera.lt/?fen=8/7k/8/8/8/N7/1nq5/K7&moves=8}\\{\sffamily\tiny\color{ink}scannen zum Lösen}
\end{minipage}
\vspace{2mm}

{\raggedright\noindent\textcolor{ink}{\sffamily\textsc{Themen}}\enspace \themetag{reines Matt} \themetag{Mustermatt} \themetag{Idealmatt} \themetag{Rundlauf} \themetag{Selbstblock} \themetag{ohne Schlag} \themetag{ohne Schach}\par}

\begin{tcolorbox}[enhanced,colback=white,colframe=ink!35,arc=1.5pt,boxrule=0.4pt,left=3mm,right=3mm,top=1.5mm,bottom=1.5mm]
\textcolor{ink}{\sffamily\textsc{veröffentlicht}}\par\smallskip\small \textbf{Sheglow, Wiktor S.}, h\#8\quad{\color{ink}\scriptsize verwandte Stellung}\\{\scriptsize Suomen Tehtäväniekat, No. (289), 05-06/1998}
\end{tcolorbox}
\clearpage
\phantomsection\label{hm:KQBvkn}
\begin{tcolorbox}[enhanced,colback=ink,colframe=ink,arc=0pt,boxrule=0pt,left=4mm,right=4mm,top=2.5mm,bottom=2.5mm]
\color{white}\sffamily{\LARGE\bfseries KQBvkn}\hfill{\Large h\#8}\\[0.5mm]
{\small 5 Steine\quad Schwarz am Zug}\hfill{\small Nr.~75\qquad\href{https://helpman.komtera.lt/?fen=8/8/8/7B/8/2Q5/3k1n2/1K6&moves=8}{\color{white}\bfseries lösen\,\ensuremath{\nearrow}}}
\end{tcolorbox}
\vspace{2mm}
\noindent\begin{minipage}[t]{36mm}\vspace{0pt}\sffamily\small\raggedright
\hmfigure{längste eindeutige}{{\Large h\#8}}\hmfigure{eindeutig bei dieser Länge}{{\Large 25} Stellungen}\hmfigure{längstes Matt}{{\Large h\#12}}\hmfigure{in der Klasse}{242 Millionen Stellungen,\\ 153 Millionen mit Hilfsmatt}
\end{minipage}\hfill
\begin{minipage}[t]{80mm}\vspace{0pt}\begin{tikzpicture}\begin{axis}[width=72mm,height=42mm,ybar,bar width=1.6pt,xmin=-1,xmax=25,ymin=0,ymax=76.35,axis y line*=left,axis x line*=bottom,xlabel={\scriptsize Halbzüge bis Matt},ylabel={\scriptsize Stellungen (Mio.)},tick label style={font=\tiny},label style={font=\scriptsize},enlarge x limits=false,ylabel shift=-3pt,xlabel shift=-3pt]\addplot[fill=ink!55,draw=none] coordinates {(0,0.3796) (1,2.383) (2,9.493) (3,31.53) (4,69.41) (5,24.67) (6,13.13) (7,0.2556) (8,0.1053) (9,2e-06) (10,0.0985) (12,0.1818) (14,0.3021) (16,0.3531) (18,0.2866) (20,0.1936) (22,0.07338) (24,0.008247)};\draw[densely dotted] (axis cs:16,0) -- (axis cs:16,76.35);\end{axis}\begin{axis}[width=72mm,height=42mm,axis y line*=right,axis x line=none,ymode=log,xmin=-1,xmax=25,ylabel={\scriptsize eindeutig (log)},tick label style={font=\tiny},label style={font=\scriptsize},ylabel shift=-3pt]\addplot[mark=*,mark size=0.7pt,draw=black,thin] coordinates {(0,379608) (1,1949857) (2,5694265) (3,6965823) (4,3580336) (5,44077) (6,4510) (8,1225) (10,753) (12,335) (14,85) (16,25)};\end{axis}\end{tikzpicture}\end{minipage}
\vspace{3mm}

\noindent\begin{minipage}[t]{60mm}\vspace{0pt}\sbox\hmboard{\chessboard[setfen=8/8/8/7B/8/2Q5/3k1n2/1K6 b - - 0 1]}\usebox\hmboard\\[1mm]
\makebox[\wd\hmboard]{\sffamily\hspace{3mm}\textbf{h\#8}\hfill(3+2)}\end{minipage}\hfill
\begin{minipage}[t]{52mm}\vspace{0pt}
\textcolor{ink}{\sffamily\textsc{Lösung}}\par\smallskip\small
\begin{tabular}{@{}r@{\enspace}l@{\enspace}l@{}}
\textbf{1.} & \symking{}xc3 & \symbishop{}e2\\
\textbf{2.} & \symking{}d2 & \symbishop{}f1\\
\textbf{3.} & \symking{}e1 & \symking{}c2\\
\textbf{4.} & \symknight{}h3 & \symking{}d3\\
\textbf{5.} & \symking{}f2 & \symking{}e4\\
\textbf{6.} & \symking{}g1 & \symking{}f3\\
\textbf{7.} & \symking{}h1 & \symking{}g3\\
\textbf{8.} & \symknight{}g1 & \symbishop{}g2\#\\
\end{tabular}\par\vspace{3mm}
\qrcode[height=12mm]{https://helpman.komtera.lt/?fen=8/8/8/7B/8/2Q5/3k1n2/1K6&moves=8}\\{\sffamily\tiny\color{ink}scannen zum Lösen}
\end{minipage}
\vspace{2mm}

{\raggedright\noindent\textcolor{ink}{\sffamily\textsc{Themen}}\enspace \themetag{reines Matt} \themetag{Mustermatt} \themetag{Idealmatt} \themetag{Switchback} \themetag{Selbstblock} \themetag{ohne Schach}\par\smallskip{\footnotesize\textit{Hinweis:} die Partei am Zug steht im Diagramm im Schach; die Lösung beginnt mit einem Schlag.}\par}

\clearpage
\phantomsection\label{hm:KQNvkb}
\begin{tcolorbox}[enhanced,colback=ink,colframe=ink,arc=0pt,boxrule=0pt,left=4mm,right=4mm,top=2.5mm,bottom=2.5mm]
\color{white}\sffamily{\LARGE\bfseries KQNvkb}\hfill{\Large h\#8}\\[0.5mm]
{\small 5 Steine\quad Schwarz am Zug}\hfill{\small Nr.~76\qquad\href{https://helpman.komtera.lt/?fen=8/8/8/6Qk/8/1b6/4N3/K7&moves=8}{\color{white}\bfseries lösen\,\ensuremath{\nearrow}}}
\end{tcolorbox}
\vspace{2mm}
\noindent\begin{minipage}[t]{36mm}\vspace{0pt}\sffamily\small\raggedright
\hmfigure{längste eindeutige}{{\Large h\#8}}\hmfigure{eindeutig bei dieser Länge}{{\Large 5} Stellungen}\hmfigure{längstes Matt}{{\Large h\#11}}\hmfigure{in der Klasse}{242 Millionen Stellungen,\\ 152 Millionen mit Hilfsmatt}
\end{minipage}\hfill
\begin{minipage}[t]{80mm}\vspace{0pt}\begin{tikzpicture}\begin{axis}[width=72mm,height=42mm,ybar,bar width=1.6pt,xmin=-1,xmax=23,ymin=0,ymax=54.49,axis y line*=left,axis x line*=bottom,xlabel={\scriptsize Halbzüge bis Matt},ylabel={\scriptsize Stellungen (Mio.)},tick label style={font=\tiny},label style={font=\scriptsize},enlarge x limits=false,ylabel shift=-3pt,xlabel shift=-3pt]\addplot[fill=ink!55,draw=none] coordinates {(0,0.1998) (1,1.601) (2,6.29) (3,20.73) (4,49.53) (5,35.57) (6,30.5) (7,4.822) (8,1.401) (9,0.0269) (10,0.1118) (11,1e-06) (12,0.1967) (14,0.2635) (16,0.3262) (17,1.2e-05) (18,0.2191) (20,0.01348) (22,1e-06)};\draw[densely dotted] (axis cs:16,0) -- (axis cs:16,54.49);\end{axis}\begin{axis}[width=72mm,height=42mm,axis y line*=right,axis x line=none,ymode=log,xmin=-1,xmax=23,ylabel={\scriptsize eindeutig (log)},tick label style={font=\tiny},label style={font=\scriptsize},ylabel shift=-3pt]\addplot[mark=*,mark size=0.7pt,draw=black,thin] coordinates {(0,199802) (1,1328518) (2,3664551) (3,3993792) (4,4175507) (5,549169) (6,134704) (7,2289) (8,861) (10,769) (12,355) (14,11) (16,5)};\end{axis}\end{tikzpicture}\end{minipage}
\vspace{3mm}

\noindent\begin{minipage}[t]{60mm}\vspace{0pt}\sbox\hmboard{\chessboard[setfen=8/8/8/6Qk/8/1b6/4N3/K7 b - - 0 1]}\usebox\hmboard\\[1mm]
\makebox[\wd\hmboard]{\sffamily\hspace{3mm}\textbf{h\#8}\hfill(3+2)}\end{minipage}\hfill
\begin{minipage}[t]{52mm}\vspace{0pt}
\textcolor{ink}{\sffamily\textsc{Lösung}}\par\smallskip\small
\begin{tabular}{@{}r@{\enspace}l@{\enspace}l@{}}
\textbf{1.} & \symking{}xg5 & \symknight{}d4\\
\textbf{2.} & \symking{}f4 & \symking{}b2\\
\textbf{3.} & \symking{}e3 & \symking{}a3\\
\textbf{4.} & \symking{}d2 & \symking{}b4\\
\textbf{5.} & \symking{}c1 & \symking{}c3\\
\textbf{6.} & \symking{}b1 & \symking{}d2\\
\textbf{7.} & \symking{}a1 & \symking{}c1\\
\textbf{8.} & \symbishop{}a2 & \symknight{}c2\#\\
\end{tabular}\par\vspace{3mm}
\qrcode[height=12mm]{https://helpman.komtera.lt/?fen=8/8/8/6Qk/8/1b6/4N3/K7&moves=8}\\{\sffamily\tiny\color{ink}scannen zum Lösen}
\end{minipage}
\vspace{2mm}

{\raggedright\noindent\textcolor{ink}{\sffamily\textsc{Themen}}\enspace \themetag{reines Matt} \themetag{Mustermatt} \themetag{Idealmatt} \themetag{Selbstblock} \themetag{ohne Schach}\par\smallskip{\footnotesize\textit{Hinweis:} die Partei am Zug steht im Diagramm im Schach; die Lösung beginnt mit einem Schlag.}\par}

\clearpage
\phantomsection\label{hm:KQNvkn}
\begin{tcolorbox}[enhanced,colback=ink,colframe=ink,arc=0pt,boxrule=0pt,left=4mm,right=4mm,top=2.5mm,bottom=2.5mm]
\color{white}\sffamily{\LARGE\bfseries KQNvkn}\hfill{\Large h\#8}\\[0.5mm]
{\small 5 Steine\quad Schwarz am Zug}\hfill{\small Nr.~77\qquad\href{https://helpman.komtera.lt/?fen=5n2/8/8/8/8/8/3Q2N1/K1k5&moves=8}{\color{white}\bfseries lösen\,\ensuremath{\nearrow}}}
\end{tcolorbox}
\vspace{2mm}
\noindent\begin{minipage}[t]{36mm}\vspace{0pt}\sffamily\small\raggedright
\hmfigure{längste eindeutige}{{\Large h\#8}}\hmfigure{eindeutig bei dieser Länge}{{\Large 16} Stellungen}\hmfigure{längstes Matt}{{\Large h\#10}}\hmfigure{in der Klasse}{242 Millionen Stellungen,\\ 157 Millionen mit Hilfsmatt}
\end{minipage}\hfill
\begin{minipage}[t]{80mm}\vspace{0pt}\begin{tikzpicture}\begin{axis}[width=72mm,height=42mm,ybar,bar width=1.6pt,xmin=-1,xmax=21,ymin=0,ymax=53.08,axis y line*=left,axis x line*=bottom,xlabel={\scriptsize Halbzüge bis Matt},ylabel={\scriptsize Stellungen (Mio.)},tick label style={font=\tiny},label style={font=\scriptsize},enlarge x limits=false,ylabel shift=-3pt,xlabel shift=-3pt]\addplot[fill=ink!55,draw=none] coordinates {(0,0.2295) (1,1.741) (2,6.285) (3,20.04) (4,48.26) (5,35.43) (6,35.88) (7,5.498) (8,2.145) (9,0.04093) (10,0.1396) (12,0.3054) (14,0.4197) (16,0.283) (18,0.06724) (20,0.002411)};\draw[densely dotted] (axis cs:16,0) -- (axis cs:16,53.08);\end{axis}\begin{axis}[width=72mm,height=42mm,axis y line*=right,axis x line=none,ymode=log,xmin=-1,xmax=21,ylabel={\scriptsize eindeutig (log)},tick label style={font=\tiny},label style={font=\scriptsize},ylabel shift=-3pt]\addplot[mark=*,mark size=0.7pt,draw=black,thin] coordinates {(0,229504) (1,1395649) (2,3590985) (3,3865402) (4,4142903) (5,495260) (6,143655) (7,2304) (8,1117) (10,390) (12,249) (14,39) (16,16)};\end{axis}\end{tikzpicture}\end{minipage}
\vspace{3mm}

\noindent\begin{minipage}[t]{60mm}\vspace{0pt}\sbox\hmboard{\chessboard[setfen=5n2/8/8/8/8/8/3Q2N1/K1k5 b - - 0 1]}\usebox\hmboard\\[1mm]
\makebox[\wd\hmboard]{\sffamily\hspace{3mm}\textbf{h\#8}\hfill(3+2)}\end{minipage}\hfill
\begin{minipage}[t]{52mm}\vspace{0pt}
\textcolor{ink}{\sffamily\textsc{Lösung}}\par\smallskip\small
\begin{tabular}{@{}r@{\enspace}l@{\enspace}l@{}}
\textbf{1.} & \symking{}xd2 & \symknight{}h4\\
\textbf{2.} & \symking{}e3 & \symking{}b2\\
\textbf{3.} & \symking{}f4 & \symking{}c3\\
\textbf{4.} & \symking{}g5 & \symking{}d4\\
\textbf{5.} & \symking{}h6 & \symking{}e5\\
\textbf{6.} & \symking{}h7 & \symking{}f6\\
\textbf{7.} & \symking{}h8 & \symking{}f7\\
\textbf{8.} & \symknight{}h7 & \symknight{}g6\#\\
\end{tabular}\par\vspace{3mm}
\qrcode[height=12mm]{https://helpman.komtera.lt/?fen=5n2/8/8/8/8/8/3Q2N1/K1k5&moves=8}\\{\sffamily\tiny\color{ink}scannen zum Lösen}
\end{minipage}
\vspace{2mm}

{\raggedright\noindent\textcolor{ink}{\sffamily\textsc{Themen}}\enspace \themetag{reines Matt} \themetag{Mustermatt} \themetag{Idealmatt} \themetag{Selbstblock} \themetag{ohne Schach}\par\smallskip{\footnotesize\textit{Hinweis:} die Partei am Zug steht im Diagramm im Schach; die Lösung beginnt mit einem Schlag.}\par}

\clearpage
\phantomsection\label{hm:KQNvkp}
\begin{tcolorbox}[enhanced,colback=ink,colframe=ink,arc=0pt,boxrule=0pt,left=4mm,right=4mm,top=2.5mm,bottom=2.5mm]
\color{white}\sffamily{\LARGE\bfseries KQNvkp}\hfill{\Large h\#8}\\[0.5mm]
{\small 5 Steine\quad Schwarz am Zug}\hfill{\small Nr.~78\qquad\href{https://helpman.komtera.lt/?fen=8/8/8/3p4/8/kQ6/N7/K7&moves=8}{\color{white}\bfseries lösen\,\ensuremath{\nearrow}}}
\end{tcolorbox}
\vspace{2mm}
\noindent\begin{minipage}[t]{36mm}\vspace{0pt}\sffamily\small\raggedright
\hmfigure{längste eindeutige}{{\Large h\#8}}\hmfigure{eindeutig bei dieser Länge}{{\Large 2} Stellungen}\hmfigure{längstes Matt}{{\Large h\#11}}\hmfigure{in der Klasse}{710 Millionen Stellungen,\\ 488 Millionen mit Hilfsmatt}
\end{minipage}\hfill
\begin{minipage}[t]{80mm}\vspace{0pt}\begin{tikzpicture}\begin{axis}[width=72mm,height=42mm,ybar,bar width=1.6pt,xmin=-1,xmax=23,ymin=0,ymax=173.8,axis y line*=left,axis x line*=bottom,xlabel={\scriptsize Halbzüge bis Matt},ylabel={\scriptsize Stellungen (Mio.)},tick label style={font=\tiny},label style={font=\scriptsize},enlarge x limits=false,ylabel shift=-3pt,xlabel shift=-3pt]\addplot[fill=ink!55,draw=none] coordinates {(0,0.8632) (1,6.372) (2,21.64) (3,65.24) (4,158) (5,103.4) (6,108.1) (7,14.02) (8,6.206) (9,0.04266) (10,0.2242) (11,1e-06) (12,0.4945) (14,0.8413) (15,1e-06) (16,1.016) (18,0.8058) (20,0.3909) (22,0.1206)};\draw[densely dotted] (axis cs:16,0) -- (axis cs:16,173.8);\end{axis}\begin{axis}[width=72mm,height=42mm,axis y line*=right,axis x line=none,ymode=log,xmin=-1,xmax=23,ylabel={\scriptsize eindeutig (log)},tick label style={font=\tiny},label style={font=\scriptsize},ylabel shift=-3pt]\addplot[mark=*,mark size=0.7pt,draw=black,thin] coordinates {(0,863243) (1,4989856) (2,14534874) (3,14041921) (4,14020679) (5,1337052) (6,468790) (7,3001) (8,1914) (10,1469) (12,821) (14,106) (16,2)};\end{axis}\end{tikzpicture}\end{minipage}
\vspace{3mm}

\noindent\begin{minipage}[t]{60mm}\vspace{0pt}\sbox\hmboard{\chessboard[setfen=8/8/8/3p4/8/kQ6/N7/K7 b - - 0 1]}\usebox\hmboard\\[1mm]
\makebox[\wd\hmboard]{\sffamily\hspace{3mm}\textbf{h\#8}\hfill(3+2)}\end{minipage}\hfill
\begin{minipage}[t]{52mm}\vspace{0pt}
\textcolor{ink}{\sffamily\textsc{Lösung}}\par\smallskip\small
\begin{tabular}{@{}r@{\enspace}l@{\enspace}l@{}}
\textbf{1.} & \symking{}xb3 & \symking{}b1\\
\textbf{2.} & d4 & \symking{}c1\\
\textbf{3.} & d3 & \symking{}d2\\
\textbf{4.} & \symking{}b2 & \symknight{}b4\\
\textbf{5.} & \symking{}a1 & \symking{}c3\\
\textbf{6.} & d2 & \symking{}b3\\
\textbf{7.} & d1=\symrook{} & \symking{}a3\\
\textbf{8.} & \symrook{}b1 & \symknight{}c2\#\\
\end{tabular}\par\vspace{3mm}
\qrcode[height=12mm]{https://helpman.komtera.lt/?fen=8/8/8/3p4/8/kQ6/N7/K7&moves=8}\\{\sffamily\tiny\color{ink}scannen zum Lösen}
\end{minipage}
\vspace{2mm}

{\raggedright\noindent\textcolor{ink}{\sffamily\textsc{Themen}}\enspace \themetag{reines Matt} \themetag{Mustermatt} \themetag{Idealmatt} \themetag{Umwandlung} \themetag{Unterverwandlung} \themetag{Selbstblock} \themetag{ohne Schach} \themetag{Umnov}\par\smallskip{\footnotesize\textit{Hinweis:} die Partei am Zug steht im Diagramm im Schach; die Lösung beginnt mit einem Schlag.}\par}

\clearpage
\phantomsection\label{hm:KQNvkr}
\begin{tcolorbox}[enhanced,colback=ink,colframe=ink,arc=0pt,boxrule=0pt,left=4mm,right=4mm,top=2.5mm,bottom=2.5mm]
\color{white}\sffamily{\LARGE\bfseries KQNvkr}\hfill{\Large h\#8}\\[0.5mm]
{\small 5 Steine\quad Schwarz am Zug}\hfill{\small Nr.~79\qquad\href{https://helpman.komtera.lt/?fen=8/8/8/8/8/8/5r1N/K5Qk&moves=8}{\color{white}\bfseries lösen\,\ensuremath{\nearrow}}}
\end{tcolorbox}
\vspace{2mm}
\noindent\begin{minipage}[t]{36mm}\vspace{0pt}\sffamily\small\raggedright
\hmfigure{längste eindeutige}{{\Large h\#8}}\hmfigure{eindeutig bei dieser Länge}{{\Large 11} Stellungen}\hmfigure{längstes Matt}{{\Large h\#10}}\hmfigure{in der Klasse}{242 Millionen Stellungen,\\ 144 Millionen mit Hilfsmatt}
\end{minipage}\hfill
\begin{minipage}[t]{80mm}\vspace{0pt}\begin{tikzpicture}\begin{axis}[width=72mm,height=42mm,ybar,bar width=1.6pt,xmin=-1,xmax=21,ymin=0,ymax=54.75,axis y line*=left,axis x line*=bottom,xlabel={\scriptsize Halbzüge bis Matt},ylabel={\scriptsize Stellungen (Mio.)},tick label style={font=\tiny},label style={font=\scriptsize},enlarge x limits=false,ylabel shift=-3pt,xlabel shift=-3pt]\addplot[fill=ink!55,draw=none] coordinates {(0,0.17) (1,1.425) (2,6.288) (3,20.87) (4,49.77) (5,36.72) (6,23.85) (7,3.724) (8,0.461) (9,0.004524) (10,0.1847) (12,0.3321) (14,0.2785) (16,0.1682) (18,0.0396) (20,0.000247)};\draw[densely dotted] (axis cs:16,0) -- (axis cs:16,54.75);\end{axis}\begin{axis}[width=72mm,height=42mm,axis y line*=right,axis x line=none,ymode=log,xmin=-1,xmax=21,ylabel={\scriptsize eindeutig (log)},tick label style={font=\tiny},label style={font=\scriptsize},ylabel shift=-3pt]\addplot[mark=*,mark size=0.7pt,draw=black,thin] coordinates {(0,169999) (1,1218654) (2,3791840) (3,3829412) (4,3455272) (5,541211) (6,60158) (7,482) (8,1554) (10,1467) (12,1402) (14,176) (16,11)};\end{axis}\end{tikzpicture}\end{minipage}
\vspace{3mm}

\noindent\begin{minipage}[t]{60mm}\vspace{0pt}\sbox\hmboard{\chessboard[setfen=8/8/8/8/8/8/5r1N/K5Qk b - - 0 1]}\usebox\hmboard\\[1mm]
\makebox[\wd\hmboard]{\sffamily\hspace{3mm}\textbf{h\#8}\hfill(3+2)}\end{minipage}\hfill
\begin{minipage}[t]{52mm}\vspace{0pt}
\textcolor{ink}{\sffamily\textsc{Lösung}}\par\smallskip\small
\begin{tabular}{@{}r@{\enspace}l@{\enspace}l@{}}
\textbf{1.} & \symking{}xg1 & \symknight{}f3+\\
\textbf{2.} & \symking{}f1 & \symking{}b1\\
\textbf{3.} & \symking{}e2 & \symking{}b2\\
\textbf{4.} & \symking{}d3+ & \symking{}c1\\
\textbf{5.} & \symking{}c3 & \symking{}d1\\
\textbf{6.} & \symking{}b2 & \symknight{}d4\\
\textbf{7.} & \symking{}a1 & \symking{}c1\\
\textbf{8.} & \symrook{}a2 & \symknight{}b3\#\\
\end{tabular}\par\vspace{3mm}
\qrcode[height=12mm]{https://helpman.komtera.lt/?fen=8/8/8/8/8/8/5r1N/K5Qk&moves=8}\\{\sffamily\tiny\color{ink}scannen zum Lösen}
\end{minipage}
\vspace{2mm}

{\raggedright\noindent\textcolor{ink}{\sffamily\textsc{Themen}}\enspace \themetag{reines Matt} \themetag{Mustermatt} \themetag{Idealmatt} \themetag{Switchback} \themetag{Selbstblock}\par\smallskip{\footnotesize\textit{Hinweis:} die Partei am Zug steht im Diagramm im Schach; die Lösung beginnt mit einem Schlag.}\par}

\clearpage
\phantomsection\label{hm:KQPvkb}
\begin{tcolorbox}[enhanced,colback=ink,colframe=ink,arc=0pt,boxrule=0pt,left=4mm,right=4mm,top=2.5mm,bottom=2.5mm]
\color{white}\sffamily{\LARGE\bfseries KQPvkb}\hfill{\Large h\#8}\\[0.5mm]
{\small 5 Steine\quad Schwarz am Zug}\hfill{\small Nr.~80\qquad\href{https://helpman.komtera.lt/?fen=8/8/1K6/8/2b5/2P5/7Q/7k&moves=8}{\color{white}\bfseries lösen\,\ensuremath{\nearrow}}}
\end{tcolorbox}
\vspace{2mm}
\noindent\begin{minipage}[t]{36mm}\vspace{0pt}\sffamily\small\raggedright
\hmfigure{längste eindeutige}{{\Large h\#8}}\hmfigure{eindeutig bei dieser Länge}{{\Large 4} Stellungen}\hmfigure{längstes Matt}{{\Large h\#9}}\hmfigure{in der Klasse}{710 Millionen Stellungen,\\ 467 Millionen mit Hilfsmatt}
\end{minipage}\hfill
\begin{minipage}[t]{80mm}\vspace{0pt}\begin{tikzpicture}\begin{axis}[width=72mm,height=42mm,ybar,bar width=1.6pt,xmin=-1,xmax=19,ymin=0,ymax=125.6,axis y line*=left,axis x line*=bottom,xlabel={\scriptsize Halbzüge bis Matt},ylabel={\scriptsize Stellungen (Mio.)},tick label style={font=\tiny},label style={font=\scriptsize},enlarge x limits=false,ylabel shift=-3pt,xlabel shift=-3pt]\addplot[fill=ink!55,draw=none] coordinates {(0,0.4076) (1,3.249) (2,10.78) (3,40.76) (4,98.83) (5,95.57) (6,114.2) (7,51.23) (8,38.04) (9,8.296) (10,3.757) (11,0.1854) (12,0.8941) (13,3.7e-05) (14,0.7146) (16,0.07246) (18,0.000188)};\draw[densely dotted] (axis cs:16,0) -- (axis cs:16,125.6);\end{axis}\begin{axis}[width=72mm,height=42mm,axis y line*=right,axis x line=none,ymode=log,xmin=-1,xmax=19,ylabel={\scriptsize eindeutig (log)},tick label style={font=\tiny},label style={font=\scriptsize},ylabel shift=-3pt]\addplot[mark=*,mark size=0.7pt,draw=black,thin] coordinates {(0,407642) (1,2301270) (2,5127291) (3,6923727) (4,8761083) (5,3193143) (6,1853885) (7,415403) (8,130886) (9,4487) (10,4172) (12,2213) (14,775) (16,4)};\end{axis}\end{tikzpicture}\end{minipage}
\vspace{3mm}

\noindent\begin{minipage}[t]{60mm}\vspace{0pt}\sbox\hmboard{\chessboard[setfen=8/8/1K6/8/2b5/2P5/7Q/7k b - - 0 1]}\usebox\hmboard\\[1mm]
\makebox[\wd\hmboard]{\sffamily\hspace{3mm}\textbf{h\#8}\hfill(3+2)}\end{minipage}\hfill
\begin{minipage}[t]{52mm}\vspace{0pt}
\textcolor{ink}{\sffamily\textsc{Lösung}}\par\smallskip\small
\begin{tabular}{@{}r@{\enspace}l@{\enspace}l@{}}
\textbf{1.} & \symking{}xh2 & \symking{}a5\\
\textbf{2.} & \symbishop{}d5 & c4\\
\textbf{3.} & \symking{}g3 & cxd5\\
\textbf{4.} & \symking{}f4 & d6\\
\textbf{5.} & \symking{}e5 & d7\\
\textbf{6.} & \symking{}d6 & \symking{}b5\\
\textbf{7.} & \symking{}c7 & \symking{}a6\\
\textbf{8.} & \symking{}b8 & d8=\symqueen{}\#\\
\end{tabular}\par\vspace{3mm}
\qrcode[height=12mm]{https://helpman.komtera.lt/?fen=8/8/1K6/8/2b5/2P5/7Q/7k&moves=8}\\{\sffamily\tiny\color{ink}scannen zum Lösen}
\end{minipage}
\vspace{2mm}

{\raggedright\noindent\textcolor{ink}{\sffamily\textsc{Themen}}\enspace \themetag{reines Matt} \themetag{Mustermatt} \themetag{Idealmatt} \themetag{Spiegelmatt} \themetag{Umwandlung} \themetag{Phönix} \themetag{ohne Schach} \themetag{Umnov}\par\smallskip{\footnotesize\textit{Hinweis:} die Partei am Zug steht im Diagramm im Schach; die Lösung beginnt mit einem Schlag.}\par}

\clearpage
\phantomsection\label{hm:KQPvkn}
\begin{tcolorbox}[enhanced,colback=ink,colframe=ink,arc=0pt,boxrule=0pt,left=4mm,right=4mm,top=2.5mm,bottom=2.5mm]
\color{white}\sffamily{\LARGE\bfseries KQPvkn}\hfill{\Large h\#8}\\[0.5mm]
{\small 5 Steine\quad Schwarz am Zug}\hfill{\small Nr.~81\qquad\href{https://helpman.komtera.lt/?fen=K7/8/2n5/8/8/8/2P4Q/7k&moves=8}{\color{white}\bfseries lösen\,\ensuremath{\nearrow}}}
\end{tcolorbox}
\vspace{2mm}
\noindent\begin{minipage}[t]{36mm}\vspace{0pt}\sffamily\small\raggedright
\hmfigure{längste eindeutige}{{\Large h\#8}}\hmfigure{eindeutig bei dieser Länge}{{\Large 7} Stellungen}\hmfigure{längstes Matt}{{\Large h\#9}}\hmfigure{in der Klasse}{710 Millionen Stellungen,\\ 482 Millionen mit Hilfsmatt}
\end{minipage}\hfill
\begin{minipage}[t]{80mm}\vspace{0pt}\begin{tikzpicture}\begin{axis}[width=72mm,height=42mm,ybar,bar width=1.6pt,xmin=-1,xmax=19,ymin=0,ymax=136,axis y line*=left,axis x line*=bottom,xlabel={\scriptsize Halbzüge bis Matt},ylabel={\scriptsize Stellungen (Mio.)},tick label style={font=\tiny},label style={font=\scriptsize},enlarge x limits=false,ylabel shift=-3pt,xlabel shift=-3pt]\addplot[fill=ink!55,draw=none] coordinates {(0,0.4741) (1,3.595) (2,10.79) (3,39.34) (4,92.67) (5,91.27) (6,123.6) (7,55.32) (8,47.84) (9,9.521) (10,4.961) (11,0.2536) (12,0.9541) (13,5.9e-05) (14,0.7885) (15,6.3e-05) (16,0.1363) (17,2.7e-05) (18,0.000267)};\draw[densely dotted] (axis cs:16,0) -- (axis cs:16,136);\end{axis}\begin{axis}[width=72mm,height=42mm,axis y line*=right,axis x line=none,ymode=log,xmin=-1,xmax=19,ylabel={\scriptsize eindeutig (log)},tick label style={font=\tiny},label style={font=\scriptsize},ylabel shift=-3pt]\addplot[mark=*,mark size=0.7pt,draw=black,thin] coordinates {(0,474081) (1,2346185) (2,4786059) (3,6738880) (4,8304580) (5,3395926) (6,2354870) (7,452174) (8,176643) (9,5314) (10,4439) (12,2490) (14,931) (16,7)};\end{axis}\end{tikzpicture}\end{minipage}
\vspace{3mm}

\noindent\begin{minipage}[t]{60mm}\vspace{0pt}\sbox\hmboard{\chessboard[setfen=K7/8/2n5/8/8/8/2P4Q/7k b - - 0 1]}\usebox\hmboard\\[1mm]
\makebox[\wd\hmboard]{\sffamily\hspace{3mm}\textbf{h\#8}\hfill(3+2)}\end{minipage}\hfill
\begin{minipage}[t]{52mm}\vspace{0pt}
\textcolor{ink}{\sffamily\textsc{Lösung}}\par\smallskip\small
\begin{tabular}{@{}r@{\enspace}l@{\enspace}l@{}}
\textbf{1.} & \symking{}xh2 & c3\\
\textbf{2.} & \symknight{}d4 & cxd4\\
\textbf{3.} & \symking{}g3 & d5\\
\textbf{4.} & \symking{}f4 & d6\\
\textbf{5.} & \symking{}e5 & d7\\
\textbf{6.} & \symking{}d6 & \symking{}a7\\
\textbf{7.} & \symking{}c7 & \symking{}a6\\
\textbf{8.} & \symking{}b8 & d8=\symqueen{}\#\\
\end{tabular}\par\vspace{3mm}
\qrcode[height=12mm]{https://helpman.komtera.lt/?fen=K7/8/2n5/8/8/8/2P4Q/7k&moves=8}\\{\sffamily\tiny\color{ink}scannen zum Lösen}
\end{minipage}
\vspace{2mm}

{\raggedright\noindent\textcolor{ink}{\sffamily\textsc{Themen}}\enspace \themetag{reines Matt} \themetag{Mustermatt} \themetag{Idealmatt} \themetag{Spiegelmatt} \themetag{Umwandlung} \themetag{Excelsior} \themetag{Phönix} \themetag{ohne Schach} \themetag{Umnov}\par\smallskip{\footnotesize\textit{Hinweis:} die Partei am Zug steht im Diagramm im Schach; die Lösung beginnt mit einem Schlag.}\par}

\clearpage
\phantomsection\label{hm:KQPvkp}
\begin{tcolorbox}[enhanced,colback=ink,colframe=ink,arc=0pt,boxrule=0pt,left=4mm,right=4mm,top=2.5mm,bottom=2.5mm]
\color{white}\sffamily{\LARGE\bfseries KQPvkp}\hfill{\Large h\#8}\\[0.5mm]
{\small 5 Steine\quad Schwarz am Zug}\hfill{\small Nr.~82\qquad\href{https://helpman.komtera.lt/?fen=8/8/p7/8/P7/K7/6Q1/6k1&moves=8}{\color{white}\bfseries lösen\,\ensuremath{\nearrow}}}
\end{tcolorbox}
\vspace{2mm}
\noindent\begin{minipage}[t]{36mm}\vspace{0pt}\sffamily\small\raggedright
\hmfigure{längste eindeutige}{{\Large h\#8}}\hmfigure{eindeutig bei dieser Länge}{{\Large 42} Stellungen}\hmfigure{längstes Matt}{{\Large h\#15}}\hmfigure{in der Klasse}{533 Millionen Stellungen,\\ 372 Millionen mit Hilfsmatt}
\end{minipage}\hfill
\begin{minipage}[t]{80mm}\vspace{0pt}\begin{tikzpicture}\begin{axis}[width=72mm,height=42mm,ybar,bar width=1.6pt,xmin=-1,xmax=31,ymin=0,ymax=105.9,axis y line*=left,axis x line*=bottom,xlabel={\scriptsize Halbzüge bis Matt},ylabel={\scriptsize Stellungen (Mio.)},tick label style={font=\tiny},label style={font=\scriptsize},enlarge x limits=false,ylabel shift=-3pt,xlabel shift=-3pt]\addplot[fill=ink!55,draw=none] coordinates {(0,0.456) (1,3.378) (2,9.448) (3,32.73) (4,79.32) (5,69.44) (6,96.29) (7,37.96) (8,32.96) (9,5.608) (10,3.115) (11,0.1409) (12,0.719) (13,0.000254) (14,0.527) (16,0.06487) (18,0.03221) (20,0.03396) (22,0.02071) (24,0.009569) (26,0.005906) (28,0.002446) (30,1.2e-05)};\draw[densely dotted] (axis cs:16,0) -- (axis cs:16,105.9);\end{axis}\begin{axis}[width=72mm,height=42mm,axis y line*=right,axis x line=none,ymode=log,xmin=-1,xmax=31,ylabel={\scriptsize eindeutig (log)},tick label style={font=\tiny},label style={font=\scriptsize},ylabel shift=-3pt]\addplot[mark=*,mark size=0.7pt,draw=black,thin] coordinates {(0,456045) (1,2151527) (2,5064444) (3,6127116) (4,7837730) (5,2492143) (6,1587630) (7,255260) (8,100760) (9,1924) (10,3249) (11,34) (12,2684) (14,384) (16,42)};\end{axis}\end{tikzpicture}\end{minipage}
\vspace{3mm}

\noindent\begin{minipage}[t]{60mm}\vspace{0pt}\sbox\hmboard{\chessboard[setfen=8/8/p7/8/P7/K7/6Q1/6k1 b - - 0 1]}\usebox\hmboard\\[1mm]
\makebox[\wd\hmboard]{\sffamily\hspace{3mm}\textbf{h\#8}\hfill(3+2)}\end{minipage}\hfill
\begin{minipage}[t]{52mm}\vspace{0pt}
\textcolor{ink}{\sffamily\textsc{Lösung}}\par\smallskip\small
\begin{tabular}{@{}r@{\enspace}l@{\enspace}l@{}}
\textbf{1.} & \symking{}xg2 & \symking{}b4\\
\textbf{2.} & a5+ & \symking{}xa5\\
\textbf{3.} & \symking{}f3 & \symking{}b4\\
\textbf{4.} & \symking{}e4 & a5\\
\textbf{5.} & \symking{}d5 & a6\\
\textbf{6.} & \symking{}c6 & a7\\
\textbf{7.} & \symking{}b7 & \symking{}c5\\
\textbf{8.} & \symking{}a6 & a8=\symqueen{}\#\\
\end{tabular}\par\vspace{3mm}
\qrcode[height=12mm]{https://helpman.komtera.lt/?fen=8/8/p7/8/P7/K7/6Q1/6k1&moves=8}\\{\sffamily\tiny\color{ink}scannen zum Lösen}
\end{minipage}
\vspace{2mm}

{\raggedright\noindent\textcolor{ink}{\sffamily\textsc{Themen}}\enspace \themetag{reines Matt} \themetag{Mustermatt} \themetag{Idealmatt} \themetag{Spiegelmatt} \themetag{Umwandlung} \themetag{Switchback} \themetag{Phönix}\par\smallskip{\footnotesize\textit{Hinweis:} die Partei am Zug steht im Diagramm im Schach; die Lösung beginnt mit einem Schlag.}\par}

\clearpage
\phantomsection\label{hm:KRBvkn}
\begin{tcolorbox}[enhanced,colback=ink,colframe=ink,arc=0pt,boxrule=0pt,left=4mm,right=4mm,top=2.5mm,bottom=2.5mm]
\color{white}\sffamily{\LARGE\bfseries KRBvkn}\hfill{\Large h\#8}\\[0.5mm]
{\small 5 Steine\quad Schwarz am Zug}\hfill{\small Nr.~83\qquad\href{https://helpman.komtera.lt/?fen=8/8/7n/8/8/8/1K1R1B2/3k4&moves=8}{\color{white}\bfseries lösen\,\ensuremath{\nearrow}}}
\end{tcolorbox}
\vspace{2mm}
\noindent\begin{minipage}[t]{36mm}\vspace{0pt}\sffamily\small\raggedright
\hmfigure{längste eindeutige}{{\Large h\#8}}\hmfigure{eindeutig bei dieser Länge}{{\Large 4} Stellungen}\hmfigure{längstes Matt}{{\Large h\#12}}\hmfigure{in der Klasse}{242 Millionen Stellungen,\\ 165 Millionen mit Hilfsmatt}
\end{minipage}\hfill
\begin{minipage}[t]{80mm}\vspace{0pt}\begin{tikzpicture}\begin{axis}[width=72mm,height=42mm,ybar,bar width=1.6pt,xmin=-1,xmax=25,ymin=0,ymax=44.36,axis y line*=left,axis x line*=bottom,xlabel={\scriptsize Halbzüge bis Matt},ylabel={\scriptsize Stellungen (Mio.)},tick label style={font=\tiny},label style={font=\scriptsize},enlarge x limits=false,ylabel shift=-3pt,xlabel shift=-3pt]\addplot[fill=ink!55,draw=none] coordinates {(0,0.09333) (1,0.4997) (2,1.275) (3,4.776) (4,12.17) (5,22.36) (6,40.33) (7,30.16) (8,33.24) (9,12.37) (10,6.693) (11,0.9196) (12,0.1548) (13,0.001442) (14,0.01031) (16,0.01142) (18,0.01514) (20,0.01599) (22,0.009635) (24,0.001227)};\draw[densely dotted] (axis cs:16,0) -- (axis cs:16,44.36);\end{axis}\begin{axis}[width=72mm,height=42mm,axis y line*=right,axis x line=none,ymode=log,xmin=-1,xmax=25,ylabel={\scriptsize eindeutig (log)},tick label style={font=\tiny},label style={font=\scriptsize},ylabel shift=-3pt]\addplot[mark=*,mark size=0.7pt,draw=black,thin] coordinates {(0,93330) (1,487824) (2,689106) (3,1702723) (4,3661567) (5,2792724) (6,2175796) (7,358455) (8,171941) (9,23461) (10,4144) (11,15) (12,18) (16,4)};\end{axis}\end{tikzpicture}\end{minipage}
\vspace{3mm}

\noindent\begin{minipage}[t]{60mm}\vspace{0pt}\sbox\hmboard{\chessboard[setfen=8/8/7n/8/8/8/1K1R1B2/3k4 b - - 0 1]}\usebox\hmboard\\[1mm]
\makebox[\wd\hmboard]{\sffamily\hspace{3mm}\textbf{h\#8}\hfill(3+2)}\end{minipage}\hfill
\begin{minipage}[t]{52mm}\vspace{0pt}
\textcolor{ink}{\sffamily\textsc{Lösung}}\par\smallskip\small
\begin{tabular}{@{}r@{\enspace}l@{\enspace}l@{}}
\textbf{1.} & \symking{}xd2 & \symbishop{}c5\\
\textbf{2.} & \symking{}d3 & \symbishop{}f8\\
\textbf{3.} & \symking{}e4 & \symking{}c3\\
\textbf{4.} & \symking{}f5 & \symking{}d4\\
\textbf{5.} & \symking{}g6 & \symking{}e5\\
\textbf{6.} & \symking{}h7 & \symking{}f6\\
\textbf{7.} & \symking{}h8 & \symking{}g6\\
\textbf{8.} & \symknight{}g8 & \symbishop{}g7\#\\
\end{tabular}\par\vspace{3mm}
\qrcode[height=12mm]{https://helpman.komtera.lt/?fen=8/8/7n/8/8/8/1K1R1B2/3k4&moves=8}\\{\sffamily\tiny\color{ink}scannen zum Lösen}
\end{minipage}
\vspace{2mm}

{\raggedright\noindent\textcolor{ink}{\sffamily\textsc{Themen}}\enspace \themetag{reines Matt} \themetag{Mustermatt} \themetag{Idealmatt} \themetag{Selbstblock} \themetag{ohne Schach}\par\smallskip{\footnotesize\textit{Hinweis:} die Partei am Zug steht im Diagramm im Schach; die Lösung beginnt mit einem Schlag.}\par}

\clearpage
\phantomsection\label{hm:KRNvkn}
\begin{tcolorbox}[enhanced,colback=ink,colframe=ink,arc=0pt,boxrule=0pt,left=4mm,right=4mm,top=2.5mm,bottom=2.5mm]
\color{white}\sffamily{\LARGE\bfseries KRNvkn}\hfill{\Large h\#8}\\[0.5mm]
{\small 5 Steine\quad Schwarz am Zug}\hfill{\small Nr.~84\qquad\href{https://helpman.komtera.lt/?fen=5n2/8/8/8/8/3N4/3R4/1K1k4&moves=8}{\color{white}\bfseries lösen\,\ensuremath{\nearrow}}}
\end{tcolorbox}
\vspace{2mm}
\noindent\begin{minipage}[t]{36mm}\vspace{0pt}\sffamily\small\raggedright
\hmfigure{längste eindeutige}{{\Large h\#8}}\hmfigure{eindeutig bei dieser Länge}{{\Large 2} Stellungen}\hmfigure{längstes Matt}{{\Large h\#10}}\hmfigure{in der Klasse}{242 Millionen Stellungen,\\ 169 Millionen mit Hilfsmatt}
\end{minipage}\hfill
\begin{minipage}[t]{80mm}\vspace{0pt}\begin{tikzpicture}\begin{axis}[width=72mm,height=42mm,ybar,bar width=1.6pt,xmin=-1,xmax=21,ymin=0,ymax=52.13,axis y line*=left,axis x line*=bottom,xlabel={\scriptsize Halbzüge bis Matt},ylabel={\scriptsize Stellungen (Mio.)},tick label style={font=\tiny},label style={font=\scriptsize},enlarge x limits=false,ylabel shift=-3pt,xlabel shift=-3pt]\addplot[fill=ink!55,draw=none] coordinates {(0,0.08203) (1,0.4975) (2,1.285) (3,5.174) (4,13.24) (5,26.67) (6,47.39) (7,35.96) (8,30.38) (9,6.569) (10,1.582) (11,0.0482) (12,0.0262) (14,0.02215) (16,0.01023) (18,0.002116) (20,2.5e-05)};\draw[densely dotted] (axis cs:16,0) -- (axis cs:16,52.13);\end{axis}\begin{axis}[width=72mm,height=42mm,axis y line*=right,axis x line=none,ymode=log,xmin=-1,xmax=21,ylabel={\scriptsize eindeutig (log)},tick label style={font=\tiny},label style={font=\scriptsize},ylabel shift=-3pt]\addplot[mark=*,mark size=0.7pt,draw=black,thin] coordinates {(0,82025) (1,488801) (2,716009) (3,1905065) (4,3977137) (5,2862713) (6,1965442) (7,288143) (8,44850) (9,1116) (10,42) (12,16) (16,2)};\end{axis}\end{tikzpicture}\end{minipage}
\vspace{3mm}

\noindent\begin{minipage}[t]{60mm}\vspace{0pt}\sbox\hmboard{\chessboard[setfen=5n2/8/8/8/8/3N4/3R4/1K1k4 b - - 0 1]}\usebox\hmboard\\[1mm]
\makebox[\wd\hmboard]{\sffamily\hspace{3mm}\textbf{h\#8}\hfill(3+2)}\end{minipage}\hfill
\begin{minipage}[t]{52mm}\vspace{0pt}
\textcolor{ink}{\sffamily\textsc{Lösung}}\par\smallskip\small
\begin{tabular}{@{}r@{\enspace}l@{\enspace}l@{}}
\textbf{1.} & \symking{}xd2 & \symknight{}e5\\
\textbf{2.} & \symking{}e3 & \symking{}c2\\
\textbf{3.} & \symking{}f4 & \symking{}d3\\
\textbf{4.} & \symking{}g5 & \symking{}e4\\
\textbf{5.} & \symking{}h6 & \symking{}f5\\
\textbf{6.} & \symking{}h7 & \symking{}f6\\
\textbf{7.} & \symking{}h8 & \symking{}f7\\
\textbf{8.} & \symknight{}h7 & \symknight{}g6\#\\
\end{tabular}\par\vspace{3mm}
\qrcode[height=12mm]{https://helpman.komtera.lt/?fen=5n2/8/8/8/8/3N4/3R4/1K1k4&moves=8}\\{\sffamily\tiny\color{ink}scannen zum Lösen}
\end{minipage}
\vspace{2mm}

{\raggedright\noindent\textcolor{ink}{\sffamily\textsc{Themen}}\enspace \themetag{reines Matt} \themetag{Mustermatt} \themetag{Idealmatt} \themetag{Selbstblock} \themetag{ohne Schach}\par\smallskip{\footnotesize\textit{Hinweis:} die Partei am Zug steht im Diagramm im Schach; die Lösung beginnt mit einem Schlag.}\par}

\clearpage
\phantomsection\label{hm:KBNvkq}
\begin{tcolorbox}[enhanced,colback=ink,colframe=ink,arc=0pt,boxrule=0pt,left=4mm,right=4mm,top=2.5mm,bottom=2.5mm]
\color{white}\sffamily{\LARGE\bfseries KBNvkq}\hfill{\Large h\#7.5}\\[0.5mm]
{\small 5 Steine\quad Weiß am Zug}\hfill{\small Nr.~85\qquad\href{https://helpman.komtera.lt/?fen=1N6/7k/8/B7/8/1q6/8/1K6&moves=7.5}{\color{white}\bfseries lösen\,\ensuremath{\nearrow}}}
\end{tcolorbox}
\vspace{2mm}
\noindent\begin{minipage}[t]{36mm}\vspace{0pt}\sffamily\small\raggedright
\hmfigure{längste eindeutige}{{\Large h\#7.5}}\hmfigure{eindeutig bei dieser Länge}{{\Large 1} Stellung}\hmfigure{längstes Matt}{{\Large h\#8.5}}\hmfigure{in der Klasse}{242 Millionen Stellungen,\\ 150 Millionen mit Hilfsmatt}
\end{minipage}\hfill
\begin{minipage}[t]{80mm}\vspace{0pt}\begin{tikzpicture}\begin{axis}[width=72mm,height=42mm,ybar,bar width=1.6pt,xmin=-1,xmax=18,ymin=0,ymax=38.03,axis y line*=left,axis x line*=bottom,xlabel={\scriptsize Halbzüge bis Matt},ylabel={\scriptsize Stellungen (Mio.)},tick label style={font=\tiny},label style={font=\scriptsize},enlarge x limits=false,ylabel shift=-3pt,xlabel shift=-3pt]\addplot[fill=ink!55,draw=none] coordinates {(0,0.001137) (1,0.004793) (2,0.01943) (3,0.1537) (4,0.4838) (5,2.651) (6,5.456) (7,16.16) (8,21.75) (9,34.57) (10,28.51) (11,23.73) (12,10.84) (13,4.328) (14,1.06) (15,0.1751) (16,0.007834) (17,0.000203)};\draw[densely dotted] (axis cs:15,0) -- (axis cs:15,38.03);\end{axis}\begin{axis}[width=72mm,height=42mm,axis y line*=right,axis x line=none,ymode=log,xmin=-1,xmax=18,ylabel={\scriptsize eindeutig (log)},tick label style={font=\tiny},label style={font=\scriptsize},ylabel shift=-3pt]\addplot[mark=*,mark size=0.7pt,draw=black,thin] coordinates {(0,1137) (1,4568) (2,10488) (3,55897) (4,110158) (5,377516) (6,355828) (7,512804) (8,276238) (9,196883) (10,57678) (11,16565) (12,2939) (13,370) (14,19) (15,1)};\end{axis}\end{tikzpicture}\end{minipage}
\vspace{3mm}

\noindent\begin{minipage}[t]{60mm}\vspace{0pt}\sbox\hmboard{\chessboard[setfen=1N6/7k/8/B7/8/1q6/8/1K6 w - - 0 1]}\usebox\hmboard\\[1mm]
\makebox[\wd\hmboard]{\sffamily\hspace{3mm}\textbf{h\#7.5}\hfill(3+2)}\end{minipage}\hfill
\begin{minipage}[t]{52mm}\vspace{0pt}
\textcolor{ink}{\sffamily\textsc{Lösung}}\par\smallskip\small
\begin{tabular}{@{}r@{\enspace}l@{\enspace}l@{}}
\textbf{1.} & \ldots & \symking{}a1\\
\textbf{2.} & \symking{}g6 & \symknight{}c6\\
\textbf{3.} & \symking{}f5 & \symknight{}d4+\\
\textbf{4.} & \symking{}e4 & \symknight{}xb3\\
\textbf{5.} & \symking{}d3 & \symking{}a2\\
\textbf{6.} & \symking{}c2 & \symking{}a3\\
\textbf{7.} & \symking{}b1 & \symknight{}d2+\\
\textbf{8.} & \symking{}a1 & \symbishop{}c3\#\\
\end{tabular}\par\vspace{3mm}
\qrcode[height=12mm]{https://helpman.komtera.lt/?fen=1N6/7k/8/B7/8/1q6/8/1K6&moves=7.5}\\{\sffamily\tiny\color{ink}scannen zum Lösen}
\end{minipage}
\vspace{2mm}

{\raggedright\noindent\textcolor{ink}{\sffamily\textsc{Themen}}\enspace \themetag{Spiegelmatt} \themetag{ein Stein zieht}\par\smallskip{\footnotesize\textit{Hinweis:} die Partei am Zug steht im Diagramm im Schach.}\par}

\clearpage
\phantomsection\label{hm:KNPvkb}
\begin{tcolorbox}[enhanced,colback=ink,colframe=ink,arc=0pt,boxrule=0pt,left=4mm,right=4mm,top=2.5mm,bottom=2.5mm]
\color{white}\sffamily{\LARGE\bfseries KNPvkb}\hfill{\Large h\#7.5}\\[0.5mm]
{\small 5 Steine\quad Weiß am Zug}\hfill{\small Nr.~86\qquad\href{https://helpman.komtera.lt/?fen=8/8/2K5/8/N7/8/P5b1/7k&moves=7.5}{\color{white}\bfseries lösen\,\ensuremath{\nearrow}}}
\end{tcolorbox}
\vspace{2mm}
\noindent\begin{minipage}[t]{36mm}\vspace{0pt}\sffamily\small\raggedright
\hmfigure{längste eindeutige}{{\Large h\#7.5}}\hmfigure{eindeutig bei dieser Länge}{{\Large 4} Stellungen}\hmfigure{längstes Matt}{{\Large h\#11}}\hmfigure{in der Klasse}{710 Millionen Stellungen,\\ 542 Millionen mit Hilfsmatt}
\end{minipage}\hfill
\begin{minipage}[t]{80mm}\vspace{0pt}\begin{tikzpicture}\begin{axis}[width=72mm,height=42mm,ybar,bar width=1.6pt,xmin=-1,xmax=23,ymin=0,ymax=82.41,axis y line*=left,axis x line*=bottom,xlabel={\scriptsize Halbzüge bis Matt},ylabel={\scriptsize Stellungen (Mio.)},tick label style={font=\tiny},label style={font=\scriptsize},enlarge x limits=false,ylabel shift=-3pt,xlabel shift=-3pt]\addplot[fill=ink!55,draw=none] coordinates {(0,0.001685) (1,0.09008) (2,0.2774) (3,2.897) (4,8.292) (5,32.81) (6,47.64) (7,64.18) (8,63.75) (9,69.41) (10,68) (11,74.92) (12,66.51) (13,29.6) (14,13.14) (15,0.5819) (16,0.03358) (17,8e-05) (18,0.001484) (20,7.2e-05) (22,2e-06)};\draw[densely dotted] (axis cs:15,0) -- (axis cs:15,82.41);\end{axis}\begin{axis}[width=72mm,height=42mm,axis y line*=right,axis x line=none,ymode=log,xmin=-1,xmax=23,ylabel={\scriptsize eindeutig (log)},tick label style={font=\tiny},label style={font=\scriptsize},ylabel shift=-3pt]\addplot[mark=*,mark size=0.7pt,draw=black,thin] coordinates {(0,1685) (1,26748) (2,77505) (3,1087207) (4,1942822) (5,2707993) (6,1853266) (7,1559385) (8,927177) (9,800869) (10,408896) (11,390653) (12,141310) (13,9854) (14,429) (15,4)};\end{axis}\end{tikzpicture}\end{minipage}
\vspace{3mm}

\noindent\begin{minipage}[t]{60mm}\vspace{0pt}\sbox\hmboard{\chessboard[setfen=8/8/2K5/8/N7/8/P5b1/7k w - - 0 1]}\usebox\hmboard\\[1mm]
\makebox[\wd\hmboard]{\sffamily\hspace{3mm}\textbf{h\#7.5}\hfill(3+2)}\end{minipage}\hfill
\begin{minipage}[t]{52mm}\vspace{0pt}
\textcolor{ink}{\sffamily\textsc{Lösung}}\par\smallskip\small
\begin{tabular}{@{}r@{\enspace}l@{\enspace}l@{}}
\textbf{1.} & \ldots & \symking{}d6\\
\textbf{2.} & \symking{}g1 & \symknight{}c5\\
\textbf{3.} & \symking{}f2 & a4\\
\textbf{4.} & \symking{}e3 & a5\\
\textbf{5.} & \symking{}d4 & a6\\
\textbf{6.} & \symking{}c4 & a7\\
\textbf{7.} & \symking{}b5 & a8=\symqueen{}\\
\textbf{8.} & \symking{}b6 & \symqueen{}a6\#\\
\end{tabular}\par\vspace{3mm}
\qrcode[height=12mm]{https://helpman.komtera.lt/?fen=8/8/2K5/8/N7/8/P5b1/7k&moves=7.5}\\{\sffamily\tiny\color{ink}scannen zum Lösen}
\end{minipage}
\vspace{2mm}

{\raggedright\noindent\textcolor{ink}{\sffamily\textsc{Themen}}\enspace \themetag{Umwandlung} \themetag{Excelsior} \themetag{Rundlauf} \themetag{ein Stein zieht} \themetag{ohne Schlag} \themetag{ohne Schach}\par\smallskip{\footnotesize\textit{Hinweis:} die Partei am Zug steht im Diagramm im Schach.}\par}

\clearpage
\phantomsection\label{hm:KNPvkn}
\begin{tcolorbox}[enhanced,colback=ink,colframe=ink,arc=0pt,boxrule=0pt,left=4mm,right=4mm,top=2.5mm,bottom=2.5mm]
\color{white}\sffamily{\LARGE\bfseries KNPvkn}\hfill{\Large h\#7.5}\\[0.5mm]
{\small 5 Steine\quad Weiß am Zug}\hfill{\small Nr.~87\qquad\href{https://helpman.komtera.lt/?fen=N7/K7/8/7n/8/8/P7/7k&moves=7.5}{\color{white}\bfseries lösen\,\ensuremath{\nearrow}}}
\end{tcolorbox}
\vspace{2mm}
\noindent\begin{minipage}[t]{36mm}\vspace{0pt}\sffamily\small\raggedright
\hmfigure{längste eindeutige}{{\Large h\#7.5}}\hmfigure{eindeutig bei dieser Länge}{{\Large 2} Stellungen}\hmfigure{längstes Matt}{{\Large h\#9}}\hmfigure{in der Klasse}{710 Millionen Stellungen,\\ 557 Millionen mit Hilfsmatt}
\end{minipage}\hfill
\begin{minipage}[t]{80mm}\vspace{0pt}\begin{tikzpicture}\begin{axis}[width=72mm,height=42mm,ybar,bar width=1.6pt,xmin=-1,xmax=19,ymin=0,ymax=84.83,axis y line*=left,axis x line*=bottom,xlabel={\scriptsize Halbzüge bis Matt},ylabel={\scriptsize Stellungen (Mio.)},tick label style={font=\tiny},label style={font=\scriptsize},enlarge x limits=false,ylabel shift=-3pt,xlabel shift=-3pt]\addplot[fill=ink!55,draw=none] coordinates {(0,0.001655) (1,0.1158) (2,0.2958) (3,2.834) (4,7.202) (5,30.05) (6,44.96) (7,60.58) (8,65.81) (9,69.01) (10,74.15) (11,77.12) (12,72.52) (13,33.97) (14,17.23) (15,0.7988) (16,0.08053) (17,0.000237) (18,0.00019)};\draw[densely dotted] (axis cs:15,0) -- (axis cs:15,84.83);\end{axis}\begin{axis}[width=72mm,height=42mm,axis y line*=right,axis x line=none,ymode=log,xmin=-1,xmax=19,ylabel={\scriptsize eindeutig (log)},tick label style={font=\tiny},label style={font=\scriptsize},ylabel shift=-3pt]\addplot[mark=*,mark size=0.7pt,draw=black,thin] coordinates {(0,1655) (1,28181) (2,66101) (3,987661) (4,1460016) (5,2317535) (6,1574227) (7,1284833) (8,886515) (9,732926) (10,435557) (11,451657) (12,158582) (13,10132) (14,846) (15,2)};\end{axis}\end{tikzpicture}\end{minipage}
\vspace{3mm}

\noindent\begin{minipage}[t]{60mm}\vspace{0pt}\sbox\hmboard{\chessboard[setfen=N7/K7/8/7n/8/8/P7/7k w - - 0 1]}\usebox\hmboard\\[1mm]
\makebox[\wd\hmboard]{\sffamily\hspace{3mm}\textbf{h\#7.5}\hfill(3+2)}\end{minipage}\hfill
\begin{minipage}[t]{52mm}\vspace{0pt}
\textcolor{ink}{\sffamily\textsc{Lösung}}\par\smallskip\small
\begin{tabular}{@{}r@{\enspace}l@{\enspace}l@{}}
\textbf{1.} & \ldots & a4\\
\textbf{2.} & \symknight{}f6 & a5\\
\textbf{3.} & \symknight{}d7 & a6\\
\textbf{4.} & \symknight{}b8 & \symking{}b6\\
\textbf{5.} & \symking{}g2 & a7\\
\textbf{6.} & \symking{}f3 & axb8=\symqueen{}\\
\textbf{7.} & \symking{}e4 & \symqueen{}f4+\\
\textbf{8.} & \symking{}d5 & \symknight{}c7\#\\
\end{tabular}\par\vspace{3mm}
\qrcode[height=12mm]{https://helpman.komtera.lt/?fen=N7/K7/8/7n/8/8/P7/7k&moves=7.5}\\{\sffamily\tiny\color{ink}scannen zum Lösen}
\end{minipage}
\vspace{2mm}

{\raggedright\noindent\textcolor{ink}{\sffamily\textsc{Themen}}\enspace \themetag{reines Matt} \themetag{Mustermatt} \themetag{Idealmatt} \themetag{Spiegelmatt} \themetag{Umwandlung} \themetag{Excelsior}\par}

\begin{tcolorbox}[enhanced,colback=white,colframe=ink!35,arc=1.5pt,boxrule=0.4pt,left=3mm,right=3mm,top=1.5mm,bottom=1.5mm]
\textcolor{ink}{\sffamily\textsc{veröffentlicht}}\par\smallskip\small \textbf{Paliulionis, Viktoras}, h\#8 0.1.1...\quad{\color{ink}\scriptsize diese Stellung}\\{\scriptsize Sachmatija, No. H5293, 01-03/2024}
\end{tcolorbox}
\clearpage
\phantomsection\label{hm:KNPvkq}
\begin{tcolorbox}[enhanced,colback=ink,colframe=ink,arc=0pt,boxrule=0pt,left=4mm,right=4mm,top=2.5mm,bottom=2.5mm]
\color{white}\sffamily{\LARGE\bfseries KNPvkq}\hfill{\Large h\#7.5}\\[0.5mm]
{\small 5 Steine\quad Weiß am Zug}\hfill{\small Nr.~88\qquad\href{https://helpman.komtera.lt/?fen=8/8/8/8/8/7N/1q5P/k1K5&moves=7.5}{\color{white}\bfseries lösen\,\ensuremath{\nearrow}}}
\end{tcolorbox}
\vspace{2mm}
\noindent\begin{minipage}[t]{36mm}\vspace{0pt}\sffamily\small\raggedright
\hmfigure{längste eindeutige}{{\Large h\#7.5}}\hmfigure{eindeutig bei dieser Länge}{{\Large 17} Stellungen}\hmfigure{längstes Matt}{{\Large h\#8.5}}\hmfigure{in der Klasse}{710 Millionen Stellungen,\\ 479 Millionen mit Hilfsmatt}
\end{minipage}\hfill
\begin{minipage}[t]{80mm}\vspace{0pt}\begin{tikzpicture}\begin{axis}[width=72mm,height=42mm,ybar,bar width=1.6pt,xmin=-1,xmax=18,ymin=0,ymax=81.48,axis y line*=left,axis x line*=bottom,xlabel={\scriptsize Halbzüge bis Matt},ylabel={\scriptsize Stellungen (Mio.)},tick label style={font=\tiny},label style={font=\scriptsize},enlarge x limits=false,ylabel shift=-3pt,xlabel shift=-3pt]\addplot[fill=ink!55,draw=none] coordinates {(0,0.000812) (1,0.04336) (2,0.1965) (3,2.264) (4,6.87) (5,29.02) (6,37.27) (7,61.66) (8,45.8) (9,63.35) (10,46.69) (11,74.07) (12,56.98) (13,41.71) (14,11.39) (15,1.986) (16,0.02013) (17,0.009923)};\draw[densely dotted] (axis cs:15,0) -- (axis cs:15,81.48);\end{axis}\begin{axis}[width=72mm,height=42mm,axis y line*=right,axis x line=none,ymode=log,xmin=-1,xmax=18,ylabel={\scriptsize eindeutig (log)},tick label style={font=\tiny},label style={font=\scriptsize},ylabel shift=-3pt]\addplot[mark=*,mark size=0.7pt,draw=black,thin] coordinates {(0,812) (1,15234) (2,59640) (3,848577) (4,1484037) (5,2455200) (6,992537) (7,1147602) (8,457013) (9,585101) (10,243526) (11,334303) (12,99973) (13,24362) (14,144) (15,17)};\end{axis}\end{tikzpicture}\end{minipage}
\vspace{3mm}

\noindent\begin{minipage}[t]{60mm}\vspace{0pt}\sbox\hmboard{\chessboard[setfen=8/8/8/8/8/7N/1q5P/k1K5 w - - 0 1]}\usebox\hmboard\\[1mm]
\makebox[\wd\hmboard]{\sffamily\hspace{3mm}\textbf{h\#7.5}\hfill(3+2)}\end{minipage}\hfill
\begin{minipage}[t]{52mm}\vspace{0pt}
\textcolor{ink}{\sffamily\textsc{Lösung}}\par\smallskip\small
\begin{tabular}{@{}r@{\enspace}l@{\enspace}l@{}}
\textbf{1.} & \ldots & \symking{}d1\\
\textbf{2.} & \symqueen{}g7 & \symknight{}f4\\
\textbf{3.} & \symking{}b2 & h4\\
\textbf{4.} & \symking{}c3 & h5\\
\textbf{5.} & \symking{}d4 & h6\\
\textbf{6.} & \symking{}e3 & hxg7\\
\textbf{7.} & \symking{}f2 & g8=\symqueen{}\\
\textbf{8.} & \symking{}f1 & \symqueen{}g2\#\\
\end{tabular}\par\vspace{3mm}
\qrcode[height=12mm]{https://helpman.komtera.lt/?fen=8/8/8/8/8/7N/1q5P/k1K5&moves=7.5}\\{\sffamily\tiny\color{ink}scannen zum Lösen}
\end{minipage}
\vspace{2mm}

{\raggedright\noindent\textcolor{ink}{\sffamily\textsc{Themen}}\enspace \themetag{Umwandlung} \themetag{Excelsior} \themetag{ohne Schach}\par\smallskip{\footnotesize\textit{Hinweis:} die Partei am Zug steht im Diagramm im Schach.}\par}

\clearpage
\phantomsection\label{hm:KNPvkr}
\begin{tcolorbox}[enhanced,colback=ink,colframe=ink,arc=0pt,boxrule=0pt,left=4mm,right=4mm,top=2.5mm,bottom=2.5mm]
\color{white}\sffamily{\LARGE\bfseries KNPvkr}\hfill{\Large h\#7.5}\\[0.5mm]
{\small 5 Steine\quad Weiß am Zug}\hfill{\small Nr.~89\qquad\href{https://helpman.komtera.lt/?fen=N7/8/8/8/r7/K1k5/P7/8&moves=7.5}{\color{white}\bfseries lösen\,\ensuremath{\nearrow}}}
\end{tcolorbox}
\vspace{2mm}
\noindent\begin{minipage}[t]{36mm}\vspace{0pt}\sffamily\small\raggedright
\hmfigure{längste eindeutige}{{\Large h\#7.5}}\hmfigure{eindeutig bei dieser Länge}{{\Large 3} Stellungen}\hmfigure{längstes Matt}{{\Large h\#9}}\hmfigure{in der Klasse}{710 Millionen Stellungen,\\ 519 Millionen mit Hilfsmatt}
\end{minipage}\hfill
\begin{minipage}[t]{80mm}\vspace{0pt}\begin{tikzpicture}\begin{axis}[width=72mm,height=42mm,ybar,bar width=1.6pt,xmin=-1,xmax=19,ymin=0,ymax=82.51,axis y line*=left,axis x line*=bottom,xlabel={\scriptsize Halbzüge bis Matt},ylabel={\scriptsize Stellungen (Mio.)},tick label style={font=\tiny},label style={font=\scriptsize},enlarge x limits=false,ylabel shift=-3pt,xlabel shift=-3pt]\addplot[fill=ink!55,draw=none] coordinates {(0,0.001362) (1,0.07788) (2,0.3171) (3,3.148) (4,10.71) (5,37.75) (6,52.79) (7,71.94) (8,63.54) (9,75.01) (10,67.5) (11,68.87) (12,45.62) (13,17.16) (14,4.557) (15,0.3302) (16,0.01065) (17,0.000683) (18,0.000135)};\draw[densely dotted] (axis cs:15,0) -- (axis cs:15,82.51);\end{axis}\begin{axis}[width=72mm,height=42mm,axis y line*=right,axis x line=none,ymode=log,xmin=-1,xmax=19,ylabel={\scriptsize eindeutig (log)},tick label style={font=\tiny},label style={font=\scriptsize},ylabel shift=-3pt]\addplot[mark=*,mark size=0.7pt,draw=black,thin] coordinates {(0,1362) (1,26343) (2,111408) (3,1200779) (4,2472189) (5,3211292) (6,1603450) (7,1395046) (8,641119) (9,613546) (10,309045) (11,184874) (12,39413) (13,4806) (14,80) (15,3)};\end{axis}\end{tikzpicture}\end{minipage}
\vspace{3mm}

\noindent\begin{minipage}[t]{60mm}\vspace{0pt}\sbox\hmboard{\chessboard[setfen=N7/8/8/8/r7/K1k5/P7/8 w - - 0 1]}\usebox\hmboard\\[1mm]
\makebox[\wd\hmboard]{\sffamily\hspace{3mm}\textbf{h\#7.5}\hfill(3+2)}\end{minipage}\hfill
\begin{minipage}[t]{52mm}\vspace{0pt}
\textcolor{ink}{\sffamily\textsc{Lösung}}\par\smallskip\small
\begin{tabular}{@{}r@{\enspace}l@{\enspace}l@{}}
\textbf{1.} & \ldots & \symking{}xa4\\
\textbf{2.} & \symking{}d4 & \symking{}b4\\
\textbf{3.} & \symking{}d5 & a4\\
\textbf{4.} & \symking{}c6 & a5\\
\textbf{5.} & \symking{}b7 & a6+\\
\textbf{6.} & \symking{}xa8 & a7\\
\textbf{7.} & \symking{}b7 & \symking{}c5\\
\textbf{8.} & \symking{}a6 & a8=\symqueen{}\#\\
\end{tabular}\par\vspace{3mm}
\qrcode[height=12mm]{https://helpman.komtera.lt/?fen=N7/8/8/8/r7/K1k5/P7/8&moves=7.5}\\{\sffamily\tiny\color{ink}scannen zum Lösen}
\end{minipage}
\vspace{2mm}

{\raggedright\noindent\textcolor{ink}{\sffamily\textsc{Themen}}\enspace \themetag{reines Matt} \themetag{Mustermatt} \themetag{Idealmatt} \themetag{Spiegelmatt} \themetag{Umwandlung} \themetag{Excelsior} \themetag{Switchback} \themetag{ein Stein zieht}\par\smallskip{\footnotesize\textit{Hinweis:} die Partei am Zug steht im Diagramm im Schach; die Lösung beginnt mit einem Schlag.}\par}

\clearpage
\phantomsection\label{hm:KNvkrb}
\begin{tcolorbox}[enhanced,colback=ink,colframe=ink,arc=0pt,boxrule=0pt,left=4mm,right=4mm,top=2.5mm,bottom=2.5mm]
\color{white}\sffamily{\LARGE\bfseries KNvkrb}\hfill{\Large h\#7.5}\\[0.5mm]
{\small 5 Steine\quad Weiß am Zug}\hfill{\small Nr.~90\qquad\href{https://helpman.komtera.lt/?fen=8/8/8/8/6k1/1b6/1N6/Kr6&moves=7.5}{\color{white}\bfseries lösen\,\ensuremath{\nearrow}}}
\end{tcolorbox}
\vspace{2mm}
\noindent\begin{minipage}[t]{36mm}\vspace{0pt}\sffamily\small\raggedright
\hmfigure{längste eindeutige}{{\Large h\#7.5}}\hmfigure{eindeutig bei dieser Länge}{{\Large 12} Stellungen}\hmfigure{längstes Matt}{{\Large h\#10.5}}\hmfigure{in der Klasse}{242 Millionen Stellungen,\\ 165 Millionen mit Hilfsmatt}
\end{minipage}\hfill
\begin{minipage}[t]{80mm}\vspace{0pt}\begin{tikzpicture}\begin{axis}[width=72mm,height=42mm,ybar,bar width=1.6pt,xmin=-1,xmax=22,ymin=0,ymax=43.55,axis y line*=left,axis x line*=bottom,xlabel={\scriptsize Halbzüge bis Matt},ylabel={\scriptsize Stellungen (Mio.)},tick label style={font=\tiny},label style={font=\scriptsize},enlarge x limits=false,ylabel shift=-3pt,xlabel shift=-3pt]\addplot[fill=ink!55,draw=none] coordinates {(0,0.000105) (1,0.000476) (2,0.004613) (3,0.03083) (4,0.171) (5,0.7864) (6,2.431) (7,7.616) (8,14.98) (9,30.26) (10,29.71) (11,39.59) (12,20.25) (13,14.73) (14,3.434) (15,0.8838) (16,0.09138) (17,0.02255) (19,0.001881) (21,1e-06)};\draw[densely dotted] (axis cs:15,0) -- (axis cs:15,43.55);\end{axis}\begin{axis}[width=72mm,height=42mm,axis y line*=right,axis x line=none,ymode=log,xmin=-1,xmax=22,ylabel={\scriptsize eindeutig (log)},tick label style={font=\tiny},label style={font=\scriptsize},ylabel shift=-3pt]\addplot[mark=*,mark size=0.7pt,draw=black,thin] coordinates {(0,105) (1,471) (2,4095) (3,21966) (4,72211) (5,165805) (6,177950) (7,270232) (8,269307) (9,314086) (10,120512) (11,80735) (12,6991) (13,1537) (14,57) (15,12)};\end{axis}\end{tikzpicture}\end{minipage}
\vspace{3mm}

\noindent\begin{minipage}[t]{60mm}\vspace{0pt}\sbox\hmboard{\chessboard[setfen=8/8/8/8/6k1/1b6/1N6/Kr6 w - - 0 1]}\usebox\hmboard\\[1mm]
\makebox[\wd\hmboard]{\sffamily\hspace{3mm}\textbf{h\#7.5}\hfill(2+3)}\end{minipage}\hfill
\begin{minipage}[t]{52mm}\vspace{0pt}
\textcolor{ink}{\sffamily\textsc{Lösung}}\par\smallskip\small
\begin{tabular}{@{}r@{\enspace}l@{\enspace}l@{}}
\textbf{1.} & \ldots & \symking{}xb1\\
\textbf{2.} & \symking{}f3 & \symknight{}c4\\
\textbf{3.} & \symking{}e2 & \symking{}b2\\
\textbf{4.} & \symking{}d1 & \symking{}c3\\
\textbf{5.} & \symking{}c1 & \symknight{}e3\\
\textbf{6.} & \symking{}b1 & \symking{}d2\\
\textbf{7.} & \symking{}a1 & \symking{}c1\\
\textbf{8.} & \symbishop{}a2 & \symknight{}c2\#\\
\end{tabular}\par\vspace{3mm}
\qrcode[height=12mm]{https://helpman.komtera.lt/?fen=8/8/8/8/6k1/1b6/1N6/Kr6&moves=7.5}\\{\sffamily\tiny\color{ink}scannen zum Lösen}
\end{minipage}
\vspace{2mm}

{\raggedright\noindent\textcolor{ink}{\sffamily\textsc{Themen}}\enspace \themetag{reines Matt} \themetag{Mustermatt} \themetag{Idealmatt} \themetag{Selbstblock} \themetag{ohne Schach}\par\smallskip{\footnotesize\textit{Hinweis:} die Partei am Zug steht im Diagramm im Schach; die Lösung beginnt mit einem Schlag.}\par}

\clearpage
\phantomsection\label{hm:KPPvkb}
\begin{tcolorbox}[enhanced,colback=ink,colframe=ink,arc=0pt,boxrule=0pt,left=4mm,right=4mm,top=2.5mm,bottom=2.5mm]
\color{white}\sffamily{\LARGE\bfseries KPPvkb}\hfill{\Large h\#7.5}\\[0.5mm]
{\small 5 Steine\quad Weiß am Zug}\hfill{\small Nr.~91\qquad\href{https://helpman.komtera.lt/?fen=K7/8/8/b7/8/P7/P7/7k&moves=7.5}{\color{white}\bfseries lösen\,\ensuremath{\nearrow}}}
\end{tcolorbox}
\vspace{2mm}
\noindent\begin{minipage}[t]{36mm}\vspace{0pt}\sffamily\small\raggedright
\hmfigure{längste eindeutige}{{\Large h\#7.5}}\hmfigure{eindeutig bei dieser Länge}{{\Large 1} Stellung}\hmfigure{längstes Matt}{{\Large h\#8.5}}\hmfigure{in der Klasse}{533 Millionen Stellungen,\\ 208 Millionen mit Hilfsmatt}
\end{minipage}\hfill
\begin{minipage}[t]{80mm}\vspace{0pt}\begin{tikzpicture}\begin{axis}[width=72mm,height=42mm,ybar,bar width=1.6pt,xmin=-1,xmax=18,ymin=0,ymax=37.29,axis y line*=left,axis x line*=bottom,xlabel={\scriptsize Halbzüge bis Matt},ylabel={\scriptsize Stellungen (Mio.)},tick label style={font=\tiny},label style={font=\scriptsize},enlarge x limits=false,ylabel shift=-3pt,xlabel shift=-3pt]\addplot[fill=ink!55,draw=none] coordinates {(0,0.000771) (1,0.06291) (2,0.1821) (3,1.567) (4,4.04) (5,15.31) (6,22.86) (7,33.9) (8,31.89) (9,31.79) (10,24.16) (11,20.4) (12,14.68) (13,5.273) (14,2.225) (15,0.09906) (16,0.001882) (17,1.1e-05)};\draw[densely dotted] (axis cs:15,0) -- (axis cs:15,37.29);\end{axis}\begin{axis}[width=72mm,height=42mm,axis y line*=right,axis x line=none,ymode=log,xmin=-1,xmax=18,ylabel={\scriptsize eindeutig (log)},tick label style={font=\tiny},label style={font=\scriptsize},ylabel shift=-3pt]\addplot[mark=*,mark size=0.7pt,draw=black,thin] coordinates {(0,771) (1,15695) (2,40140) (3,487277) (4,850745) (5,1278465) (6,797237) (7,623333) (8,341934) (9,260844) (10,103646) (11,77157) (12,23712) (13,1752) (14,38) (15,1)};\end{axis}\end{tikzpicture}\end{minipage}
\vspace{3mm}

\noindent\begin{minipage}[t]{60mm}\vspace{0pt}\sbox\hmboard{\chessboard[setfen=K7/8/8/b7/8/P7/P7/7k w - - 0 1]}\usebox\hmboard\\[1mm]
\makebox[\wd\hmboard]{\sffamily\hspace{3mm}\textbf{h\#7.5}\hfill(3+2)}\end{minipage}\hfill
\begin{minipage}[t]{52mm}\vspace{0pt}
\textcolor{ink}{\sffamily\textsc{Lösung}}\par\smallskip\small
\begin{tabular}{@{}r@{\enspace}l@{\enspace}l@{}}
\textbf{1.} & \ldots & \symking{}b8\\
\textbf{2.} & \symbishop{}b4 & axb4\\
\textbf{3.} & \symking{}g2 & a4\\
\textbf{4.} & \symking{}f3 & a5\\
\textbf{5.} & \symking{}e4 & a6\\
\textbf{6.} & \symking{}d5 & a7\\
\textbf{7.} & \symking{}c6 & a8=\symqueen{}+\\
\textbf{8.} & \symking{}b6 & \symqueen{}b7\#\\
\end{tabular}\par\vspace{3mm}
\qrcode[height=12mm]{https://helpman.komtera.lt/?fen=K7/8/8/b7/8/P7/P7/7k&moves=7.5}\\{\sffamily\tiny\color{ink}scannen zum Lösen}
\end{minipage}
\vspace{2mm}

{\raggedright\noindent\textcolor{ink}{\sffamily\textsc{Themen}}\enspace \themetag{Umwandlung} \themetag{Excelsior}\par}

\clearpage
\phantomsection\label{hm:KPPvkn}
\begin{tcolorbox}[enhanced,colback=ink,colframe=ink,arc=0pt,boxrule=0pt,left=4mm,right=4mm,top=2.5mm,bottom=2.5mm]
\color{white}\sffamily{\LARGE\bfseries KPPvkn}\hfill{\Large h\#7.5}\\[0.5mm]
{\small 5 Steine\quad Weiß am Zug}\hfill{\small Nr.~92\qquad\href{https://helpman.komtera.lt/?fen=1K6/8/n7/8/8/P7/P7/7k&moves=7.5}{\color{white}\bfseries lösen\,\ensuremath{\nearrow}}}
\end{tcolorbox}
\vspace{2mm}
\noindent\begin{minipage}[t]{36mm}\vspace{0pt}\sffamily\small\raggedright
\hmfigure{längste eindeutige}{{\Large h\#7.5}}\hmfigure{eindeutig bei dieser Länge}{{\Large 1} Stellung}\hmfigure{längstes Matt}{{\Large h\#8.5}}\hmfigure{in der Klasse}{533 Millionen Stellungen,\\ 214 Millionen mit Hilfsmatt}
\end{minipage}\hfill
\begin{minipage}[t]{80mm}\vspace{0pt}\begin{tikzpicture}\begin{axis}[width=72mm,height=42mm,ybar,bar width=1.6pt,xmin=-1,xmax=18,ymin=0,ymax=36.17,axis y line*=left,axis x line*=bottom,xlabel={\scriptsize Halbzüge bis Matt},ylabel={\scriptsize Stellungen (Mio.)},tick label style={font=\tiny},label style={font=\scriptsize},enlarge x limits=false,ylabel shift=-3pt,xlabel shift=-3pt]\addplot[fill=ink!55,draw=none] coordinates {(0,0.000698) (1,0.08184) (2,0.2025) (3,1.588) (4,3.759) (5,14.18) (6,21.34) (7,32.17) (8,32.88) (9,31.99) (10,26.7) (11,21.37) (12,16.69) (13,6.738) (14,3.747) (15,0.2703) (16,0.03765) (17,7.2e-05)};\draw[densely dotted] (axis cs:15,0) -- (axis cs:15,36.17);\end{axis}\begin{axis}[width=72mm,height=42mm,axis y line*=right,axis x line=none,ymode=log,xmin=-1,xmax=18,ylabel={\scriptsize eindeutig (log)},tick label style={font=\tiny},label style={font=\scriptsize},ylabel shift=-3pt]\addplot[mark=*,mark size=0.7pt,draw=black,thin] coordinates {(0,698) (1,16697) (2,37432) (3,442040) (4,676225) (5,1075687) (6,720483) (7,528559) (8,318315) (9,233054) (10,122611) (11,96053) (12,36670) (13,2995) (14,59) (15,1)};\end{axis}\end{tikzpicture}\end{minipage}
\vspace{3mm}

\noindent\begin{minipage}[t]{60mm}\vspace{0pt}\sbox\hmboard{\chessboard[setfen=1K6/8/n7/8/8/P7/P7/7k w - - 0 1]}\usebox\hmboard\\[1mm]
\makebox[\wd\hmboard]{\sffamily\hspace{3mm}\textbf{h\#7.5}\hfill(3+2)}\end{minipage}\hfill
\begin{minipage}[t]{52mm}\vspace{0pt}
\textcolor{ink}{\sffamily\textsc{Lösung}}\par\smallskip\small
\begin{tabular}{@{}r@{\enspace}l@{\enspace}l@{}}
\textbf{1.} & \ldots & \symking{}c8\\
\textbf{2.} & \symknight{}b4 & axb4\\
\textbf{3.} & \symking{}g2 & a4\\
\textbf{4.} & \symking{}f3 & a5\\
\textbf{5.} & \symking{}e4 & a6\\
\textbf{6.} & \symking{}d5 & a7\\
\textbf{7.} & \symking{}c6 & a8=\symqueen{}+\\
\textbf{8.} & \symking{}b6 & \symqueen{}b7\#\\
\end{tabular}\par\vspace{3mm}
\qrcode[height=12mm]{https://helpman.komtera.lt/?fen=1K6/8/n7/8/8/P7/P7/7k&moves=7.5}\\{\sffamily\tiny\color{ink}scannen zum Lösen}
\end{minipage}
\vspace{2mm}

{\raggedright\noindent\textcolor{ink}{\sffamily\textsc{Themen}}\enspace \themetag{Umwandlung} \themetag{Excelsior}\par\smallskip{\footnotesize\textit{Hinweis:} die Partei am Zug steht im Diagramm im Schach.}\par}

\clearpage
\phantomsection\label{hm:KPPvkq}
\begin{tcolorbox}[enhanced,colback=ink,colframe=ink,arc=0pt,boxrule=0pt,left=4mm,right=4mm,top=2.5mm,bottom=2.5mm]
\color{white}\sffamily{\LARGE\bfseries KPPvkq}\hfill{\Large h\#7.5}\\[0.5mm]
{\small 5 Steine\quad Weiß am Zug}\hfill{\small Nr.~93\qquad\href{https://helpman.komtera.lt/?fen=8/7k/8/8/8/6P1/q5P1/K7&moves=7.5}{\color{white}\bfseries lösen\,\ensuremath{\nearrow}}}
\end{tcolorbox}
\vspace{2mm}
\noindent\begin{minipage}[t]{36mm}\vspace{0pt}\sffamily\small\raggedright
\hmfigure{längste eindeutige}{{\Large h\#7.5}}\hmfigure{eindeutig bei dieser Länge}{{\Large 1} Stellung}\hmfigure{längstes Matt}{{\Large h\#8.5}}\hmfigure{in der Klasse}{533 Millionen Stellungen,\\ 185 Millionen mit Hilfsmatt}
\end{minipage}\hfill
\begin{minipage}[t]{80mm}\vspace{0pt}\begin{tikzpicture}\begin{axis}[width=72mm,height=42mm,ybar,bar width=1.6pt,xmin=-1,xmax=18,ymin=0,ymax=37.05,axis y line*=left,axis x line*=bottom,xlabel={\scriptsize Halbzüge bis Matt},ylabel={\scriptsize Stellungen (Mio.)},tick label style={font=\tiny},label style={font=\scriptsize},enlarge x limits=false,ylabel shift=-3pt,xlabel shift=-3pt]\addplot[fill=ink!55,draw=none] coordinates {(0,0.000376) (1,0.03056) (2,0.1294) (3,1.223) (4,3.255) (5,13.46) (6,17.96) (7,33.68) (8,24.55) (9,32.22) (10,18.34) (11,20.87) (12,11.15) (13,6.511) (14,1.451) (15,0.2285) (16,0.001308) (17,0.000359)};\draw[densely dotted] (axis cs:15,0) -- (axis cs:15,37.05);\end{axis}\begin{axis}[width=72mm,height=42mm,axis y line*=right,axis x line=none,ymode=log,xmin=-1,xmax=18,ylabel={\scriptsize eindeutig (log)},tick label style={font=\tiny},label style={font=\scriptsize},ylabel shift=-3pt]\addplot[mark=*,mark size=0.7pt,draw=black,thin] coordinates {(0,376) (1,8827) (2,30015) (3,376599) (4,619412) (5,1193615) (6,497574) (7,570060) (8,175815) (9,192529) (10,59135) (11,61071) (12,12797) (13,3203) (14,17) (15,1)};\end{axis}\end{tikzpicture}\end{minipage}
\vspace{3mm}

\noindent\begin{minipage}[t]{60mm}\vspace{0pt}\sbox\hmboard{\chessboard[setfen=8/7k/8/8/8/6P1/q5P1/K7 w - - 0 1]}\usebox\hmboard\\[1mm]
\makebox[\wd\hmboard]{\sffamily\hspace{3mm}\textbf{h\#7.5}\hfill(3+2)}\end{minipage}\hfill
\begin{minipage}[t]{52mm}\vspace{0pt}
\textcolor{ink}{\sffamily\textsc{Lösung}}\par\smallskip\small
\begin{tabular}{@{}r@{\enspace}l@{\enspace}l@{}}
\textbf{1.} & \ldots & \symking{}xa2\\
\textbf{2.} & \symking{}g6 & \symking{}a3\\
\textbf{3.} & \symking{}f5 & g4+\\
\textbf{4.} & \symking{}e4 & g5\\
\textbf{5.} & \symking{}d3 & g6\\
\textbf{6.} & \symking{}c2 & g7\\
\textbf{7.} & \symking{}b1 & g8=\symqueen{}\\
\textbf{8.} & \symking{}a1 & \symqueen{}a2\#\\
\end{tabular}\par\vspace{3mm}
\qrcode[height=12mm]{https://helpman.komtera.lt/?fen=8/7k/8/8/8/6P1/q5P1/K7&moves=7.5}\\{\sffamily\tiny\color{ink}scannen zum Lösen}
\end{minipage}
\vspace{2mm}

{\raggedright\noindent\textcolor{ink}{\sffamily\textsc{Themen}}\enspace \themetag{Umwandlung} \themetag{ein Stein zieht}\par\smallskip{\footnotesize\textit{Hinweis:} die Partei am Zug steht im Diagramm im Schach; die Lösung beginnt mit einem Schlag.}\par}

\clearpage
\phantomsection\label{hm:KPvkbb}
\begin{tcolorbox}[enhanced,colback=ink,colframe=ink,arc=0pt,boxrule=0pt,left=4mm,right=4mm,top=2.5mm,bottom=2.5mm]
\color{white}\sffamily{\LARGE\bfseries KPvkbb}\hfill{\Large h\#7.5}\\[0.5mm]
{\small 5 Steine\quad Weiß am Zug}\hfill{\small Nr.~94\qquad\href{https://helpman.komtera.lt/?fen=8/7k/8/3b4/8/bP6/8/K7&moves=7.5}{\color{white}\bfseries lösen\,\ensuremath{\nearrow}}}
\end{tcolorbox}
\vspace{2mm}
\noindent\begin{minipage}[t]{36mm}\vspace{0pt}\sffamily\small\raggedright
\hmfigure{längste eindeutige}{{\Large h\#7.5}}\hmfigure{eindeutig bei dieser Länge}{{\Large 448} Stellungen}\hmfigure{längstes Matt}{{\Large h\#9.5}}\hmfigure{in der Klasse}{710 Millionen Stellungen,\\ 265 Millionen mit Hilfsmatt}
\end{minipage}\hfill
\begin{minipage}[t]{80mm}\vspace{0pt}\begin{tikzpicture}\begin{axis}[width=72mm,height=42mm,ybar,bar width=1.6pt,xmin=-1,xmax=20,ymin=0,ymax=49.37,axis y line*=left,axis x line*=bottom,xlabel={\scriptsize Halbzüge bis Matt},ylabel={\scriptsize Stellungen (Mio.)},tick label style={font=\tiny},label style={font=\scriptsize},enlarge x limits=false,ylabel shift=-3pt,xlabel shift=-3pt]\addplot[fill=ink!55,draw=none] coordinates {(0,0.000194) (1,0.02221) (2,0.07239) (3,0.65) (4,1.625) (5,7.361) (6,11.42) (7,25.37) (8,25) (9,38.78) (10,30.43) (11,44.88) (12,34.52) (13,29.36) (14,12.27) (15,2.91) (16,0.279) (17,0.05972) (18,7.2e-05) (19,2e-06)};\draw[densely dotted] (axis cs:15,0) -- (axis cs:15,49.37);\end{axis}\begin{axis}[width=72mm,height=42mm,axis y line*=right,axis x line=none,ymode=log,xmin=-1,xmax=20,ylabel={\scriptsize eindeutig (log)},tick label style={font=\tiny},label style={font=\scriptsize},ylabel shift=-3pt]\addplot[mark=*,mark size=0.7pt,draw=black,thin] coordinates {(0,194) (1,5581) (2,11913) (3,164152) (4,275514) (5,694702) (6,371337) (7,485792) (8,217775) (9,291158) (10,130811) (11,187842) (12,70062) (13,37257) (14,1953) (15,448)};\end{axis}\end{tikzpicture}\end{minipage}
\vspace{3mm}

\noindent\begin{minipage}[t]{60mm}\vspace{0pt}\sbox\hmboard{\chessboard[setfen=8/7k/8/3b4/8/bP6/8/K7 w - - 0 1]}\usebox\hmboard\\[1mm]
\makebox[\wd\hmboard]{\sffamily\hspace{3mm}\textbf{h\#7.5}\hfill(2+3)}\end{minipage}\hfill
\begin{minipage}[t]{52mm}\vspace{0pt}
\textcolor{ink}{\sffamily\textsc{Lösung}}\par\smallskip\small
\begin{tabular}{@{}r@{\enspace}l@{\enspace}l@{}}
\textbf{1.} & \ldots & \symking{}a2\\
\textbf{2.} & \symking{}g6 & \symking{}xa3\\
\textbf{3.} & \symking{}f5 & b4\\
\textbf{4.} & \symking{}e4 & b5\\
\textbf{5.} & \symking{}d3 & b6\\
\textbf{6.} & \symking{}c2 & b7\\
\textbf{7.} & \symking{}b1 & b8=\symqueen{}+\\
\textbf{8.} & \symking{}a1 & \symqueen{}b2\#\\
\end{tabular}\par\vspace{3mm}
\qrcode[height=12mm]{https://helpman.komtera.lt/?fen=8/7k/8/3b4/8/bP6/8/K7&moves=7.5}\\{\sffamily\tiny\color{ink}scannen zum Lösen}
\end{minipage}
\vspace{2mm}

{\raggedright\noindent\textcolor{ink}{\sffamily\textsc{Themen}}\enspace \themetag{Umwandlung} \themetag{ein Stein zieht}\par}

\clearpage
\phantomsection\label{hm:KPvkbn}
\begin{tcolorbox}[enhanced,colback=ink,colframe=ink,arc=0pt,boxrule=0pt,left=4mm,right=4mm,top=2.5mm,bottom=2.5mm]
\color{white}\sffamily{\LARGE\bfseries KPvkbn}\hfill{\Large h\#7.5}\\[0.5mm]
{\small 5 Steine\quad Weiß am Zug}\hfill{\small Nr.~95\qquad\href{https://helpman.komtera.lt/?fen=8/6k1/8/8/2n5/b7/P7/K7&moves=7.5}{\color{white}\bfseries lösen\,\ensuremath{\nearrow}}}
\end{tcolorbox}
\vspace{2mm}
\noindent\begin{minipage}[t]{36mm}\vspace{0pt}\sffamily\small\raggedright
\hmfigure{längste eindeutige}{{\Large h\#7.5}}\hmfigure{eindeutig bei dieser Länge}{{\Large 1\,222} Stellungen}\hmfigure{längstes Matt}{{\Large h\#9.5}}\hmfigure{in der Klasse}{710 Millionen Stellungen,\\ 545 Millionen mit Hilfsmatt}
\end{minipage}\hfill
\begin{minipage}[t]{80mm}\vspace{0pt}\begin{tikzpicture}\begin{axis}[width=72mm,height=42mm,ybar,bar width=1.6pt,xmin=-1,xmax=20,ymin=0,ymax=93.87,axis y line*=left,axis x line*=bottom,xlabel={\scriptsize Halbzüge bis Matt},ylabel={\scriptsize Stellungen (Mio.)},tick label style={font=\tiny},label style={font=\scriptsize},enlarge x limits=false,ylabel shift=-3pt,xlabel shift=-3pt]\addplot[fill=ink!55,draw=none] coordinates {(0,0.000296) (1,0.05813) (2,0.1736) (3,1.407) (4,3.204) (5,13.82) (6,20.73) (7,46.03) (8,48.14) (9,73.29) (10,61.58) (11,85.33) (12,71.02) (13,69.38) (14,40.01) (15,9.311) (16,1.228) (17,0.2141) (18,0.000415) (19,1.1e-05)};\draw[densely dotted] (axis cs:15,0) -- (axis cs:15,93.87);\end{axis}\begin{axis}[width=72mm,height=42mm,axis y line*=right,axis x line=none,ymode=log,xmin=-1,xmax=20,ylabel={\scriptsize eindeutig (log)},tick label style={font=\tiny},label style={font=\scriptsize},ylabel shift=-3pt]\addplot[mark=*,mark size=0.7pt,draw=black,thin] coordinates {(0,296) (1,12411) (2,25537) (3,300630) (4,493940) (5,1254489) (6,723487) (7,843263) (8,396669) (9,515256) (10,290314) (11,419503) (12,207750) (13,104120) (14,6822) (15,1222)};\end{axis}\end{tikzpicture}\end{minipage}
\vspace{3mm}

\noindent\begin{minipage}[t]{60mm}\vspace{0pt}\sbox\hmboard{\chessboard[setfen=8/6k1/8/8/2n5/b7/P7/K7 w - - 0 1]}\usebox\hmboard\\[1mm]
\makebox[\wd\hmboard]{\sffamily\hspace{3mm}\textbf{h\#7.5}\hfill(2+3)}\end{minipage}\hfill
\begin{minipage}[t]{52mm}\vspace{0pt}
\textcolor{ink}{\sffamily\textsc{Lösung}}\par\smallskip\small
\begin{tabular}{@{}r@{\enspace}l@{\enspace}l@{}}
\textbf{1.} & \ldots & \symking{}b1\\
\textbf{2.} & \symbishop{}b4 & a4\\
\textbf{3.} & \symking{}f6 & a5\\
\textbf{4.} & \symking{}e5 & a6\\
\textbf{5.} & \symking{}d4 & a7\\
\textbf{6.} & \symking{}c3 & a8=\symqueen{}\\
\textbf{7.} & \symking{}b3 & \symqueen{}g2\\
\textbf{8.} & \symking{}a3 & \symqueen{}a2\#\\
\end{tabular}\par\vspace{3mm}
\qrcode[height=12mm]{https://helpman.komtera.lt/?fen=8/6k1/8/8/2n5/b7/P7/K7&moves=7.5}\\{\sffamily\tiny\color{ink}scannen zum Lösen}
\end{minipage}
\vspace{2mm}

{\raggedright\noindent\textcolor{ink}{\sffamily\textsc{Themen}}\enspace \themetag{Satzspiel} \themetag{Umwandlung} \themetag{Excelsior} \themetag{Rundlauf} \themetag{Selbstblock} \themetag{ohne Schlag} \themetag{ohne Schach}\par}

\clearpage
\phantomsection\label{hm:KPvknn}
\begin{tcolorbox}[enhanced,colback=ink,colframe=ink,arc=0pt,boxrule=0pt,left=4mm,right=4mm,top=2.5mm,bottom=2.5mm]
\color{white}\sffamily{\LARGE\bfseries KPvknn}\hfill{\Large h\#7.5}\\[0.5mm]
{\small 5 Steine\quad Weiß am Zug}\hfill{\small Nr.~96\qquad\href{https://helpman.komtera.lt/?fen=8/8/8/8/5k2/8/3n1P2/K2n4&moves=7.5}{\color{white}\bfseries lösen\,\ensuremath{\nearrow}}}
\end{tcolorbox}
\vspace{2mm}
\noindent\begin{minipage}[t]{36mm}\vspace{0pt}\sffamily\small\raggedright
\hmfigure{längste eindeutige}{{\Large h\#7.5}}\hmfigure{eindeutig bei dieser Länge}{{\Large 476} Stellungen}\hmfigure{längstes Matt}{{\Large h\#9.5}}\hmfigure{in der Klasse}{710 Millionen Stellungen,\\ 279 Millionen mit Hilfsmatt}
\end{minipage}\hfill
\begin{minipage}[t]{80mm}\vspace{0pt}\begin{tikzpicture}\begin{axis}[width=72mm,height=42mm,ybar,bar width=1.6pt,xmin=-1,xmax=20,ymin=0,ymax=45.28,axis y line*=left,axis x line*=bottom,xlabel={\scriptsize Halbzüge bis Matt},ylabel={\scriptsize Stellungen (Mio.)},tick label style={font=\tiny},label style={font=\scriptsize},enlarge x limits=false,ylabel shift=-3pt,xlabel shift=-3pt]\addplot[fill=ink!55,draw=none] coordinates {(0,0.000102) (1,0.03849) (2,0.101) (3,0.7538) (4,1.637) (5,6.634) (6,9.828) (7,22.09) (8,24.16) (9,35.67) (10,31.43) (11,41.16) (12,35.21) (13,35.67) (14,24.62) (15,7.29) (16,2.517) (17,0.1157) (18,0.000549) (19,7.9e-05)};\draw[densely dotted] (axis cs:15,0) -- (axis cs:15,45.28);\end{axis}\begin{axis}[width=72mm,height=42mm,axis y line*=right,axis x line=none,ymode=log,xmin=-1,xmax=20,ylabel={\scriptsize eindeutig (log)},tick label style={font=\tiny},label style={font=\scriptsize},ylabel shift=-3pt]\addplot[mark=*,mark size=0.7pt,draw=black,thin] coordinates {(0,102) (1,6757) (2,13624) (3,141246) (4,218712) (5,567001) (6,357269) (7,358951) (8,190144) (9,228553) (10,136549) (11,191159) (12,100397) (13,50810) (14,4516) (15,476)};\end{axis}\end{tikzpicture}\end{minipage}
\vspace{3mm}

\noindent\begin{minipage}[t]{60mm}\vspace{0pt}\sbox\hmboard{\chessboard[setfen=8/8/8/8/5k2/8/3n1P2/K2n4 w - - 0 1]}\usebox\hmboard\\[1mm]
\makebox[\wd\hmboard]{\sffamily\hspace{3mm}\textbf{h\#7.5}\hfill(2+3)}\end{minipage}\hfill
\begin{minipage}[t]{52mm}\vspace{0pt}
\textcolor{ink}{\sffamily\textsc{Lösung}}\par\smallskip\small
\begin{tabular}{@{}r@{\enspace}l@{\enspace}l@{}}
\textbf{1.} & \ldots & \symking{}a2\\
\textbf{2.} & \symking{}e4 & f4\\
\textbf{3.} & \symknight{}f3 & \symking{}a3\\
\textbf{4.} & \symknight{}g5 & fxg5\\
\textbf{5.} & \symking{}d3 & g6\\
\textbf{6.} & \symking{}c2 & g7\\
\textbf{7.} & \symking{}b1 & g8=\symqueen{}\\
\textbf{8.} & \symking{}a1 & \symqueen{}a2\#\\
\end{tabular}\par\vspace{3mm}
\qrcode[height=12mm]{https://helpman.komtera.lt/?fen=8/8/8/8/5k2/8/3n1P2/K2n4&moves=7.5}\\{\sffamily\tiny\color{ink}scannen zum Lösen}
\end{minipage}
\vspace{2mm}

{\raggedright\noindent\textcolor{ink}{\sffamily\textsc{Themen}}\enspace \themetag{Umwandlung} \themetag{Excelsior} \themetag{ohne Schach} \themetag{Umnov}\par}

\clearpage
\phantomsection\label{hm:KPvkqb}
\begin{tcolorbox}[enhanced,colback=ink,colframe=ink,arc=0pt,boxrule=0pt,left=4mm,right=4mm,top=2.5mm,bottom=2.5mm]
\color{white}\sffamily{\LARGE\bfseries KPvkqb}\hfill{\Large h\#7.5}\\[0.5mm]
{\small 5 Steine\quad Weiß am Zug}\hfill{\small Nr.~97\qquad\href{https://helpman.komtera.lt/?fen=1q6/8/7k/8/8/7b/7P/K7&moves=7.5}{\color{white}\bfseries lösen\,\ensuremath{\nearrow}}}
\end{tcolorbox}
\vspace{2mm}
\noindent\begin{minipage}[t]{36mm}\vspace{0pt}\sffamily\small\raggedright
\hmfigure{längste eindeutige}{{\Large h\#7.5}}\hmfigure{eindeutig bei dieser Länge}{{\Large 342} Stellungen}\hmfigure{längstes Matt}{{\Large h\#11.5}}\hmfigure{in der Klasse}{710 Millionen Stellungen,\\ 475 Millionen mit Hilfsmatt}
\end{minipage}\hfill
\begin{minipage}[t]{80mm}\vspace{0pt}\begin{tikzpicture}\begin{axis}[width=72mm,height=42mm,ybar,bar width=1.6pt,xmin=-1,xmax=24,ymin=0,ymax=95.13,axis y line*=left,axis x line*=bottom,xlabel={\scriptsize Halbzüge bis Matt},ylabel={\scriptsize Stellungen (Mio.)},tick label style={font=\tiny},label style={font=\scriptsize},enlarge x limits=false,ylabel shift=-3pt,xlabel shift=-3pt]\addplot[fill=ink!55,draw=none] coordinates {(0,0.000119) (1,0.02142) (2,0.09383) (3,0.9271) (4,2.351) (5,12.45) (6,17.6) (7,48.55) (8,37.7) (9,74.59) (10,46.92) (11,86.48) (12,52.58) (13,66.28) (14,19.95) (15,7.915) (16,0.2612) (17,0.05298) (18,4.9e-05) (19,1.9e-05) (20,1.2e-05) (21,2.5e-05) (22,1.1e-05) (23,9e-06)};\draw[densely dotted] (axis cs:15,0) -- (axis cs:15,95.13);\end{axis}\begin{axis}[width=72mm,height=42mm,axis y line*=right,axis x line=none,ymode=log,xmin=-1,xmax=24,ylabel={\scriptsize eindeutig (log)},tick label style={font=\tiny},label style={font=\scriptsize},ylabel shift=-3pt]\addplot[mark=*,mark size=0.7pt,draw=black,thin] coordinates {(0,119) (1,5996) (2,15786) (3,240058) (4,377084) (5,1240858) (6,511286) (7,952211) (8,300813) (9,556616) (10,194093) (11,364811) (12,103006) (13,84747) (14,1366) (15,342)};\end{axis}\end{tikzpicture}\end{minipage}
\vspace{3mm}

\noindent\begin{minipage}[t]{60mm}\vspace{0pt}\sbox\hmboard{\chessboard[setfen=1q6/8/7k/8/8/7b/7P/K7 w - - 0 1]}\usebox\hmboard\\[1mm]
\makebox[\wd\hmboard]{\sffamily\hspace{3mm}\textbf{h\#7.5}\hfill(2+3)}\end{minipage}\hfill
\begin{minipage}[t]{52mm}\vspace{0pt}
\textcolor{ink}{\sffamily\textsc{Lösung}}\par\smallskip\small
\begin{tabular}{@{}r@{\enspace}l@{\enspace}l@{}}
\textbf{1.} & \ldots & \symking{}a2\\
\textbf{2.} & \symbishop{}g4 & h4\\
\textbf{3.} & \symbishop{}d1 & h5\\
\textbf{4.} & \symking{}g5 & h6\\
\textbf{5.} & \symking{}f4 & h7\\
\textbf{6.} & \symking{}e3 & h8=\symqueen{}\\
\textbf{7.} & \symking{}d2 & \symqueen{}xb8\\
\textbf{8.} & \symking{}c1 & \symqueen{}b2\#\\
\end{tabular}\par\vspace{3mm}
\qrcode[height=12mm]{https://helpman.komtera.lt/?fen=1q6/8/7k/8/8/7b/7P/K7&moves=7.5}\\{\sffamily\tiny\color{ink}scannen zum Lösen}
\end{minipage}
\vspace{2mm}

{\raggedright\noindent\textcolor{ink}{\sffamily\textsc{Themen}}\enspace \themetag{Satzspiel} \themetag{Umwandlung} \themetag{Excelsior} \themetag{Selbstblock} \themetag{ohne Schach} \themetag{Umnov}\par}

\clearpage
\phantomsection\label{hm:KPvkqq}
\begin{tcolorbox}[enhanced,colback=ink,colframe=ink,arc=0pt,boxrule=0pt,left=4mm,right=4mm,top=2.5mm,bottom=2.5mm]
\color{white}\sffamily{\LARGE\bfseries KPvkqq}\hfill{\Large h\#7.5}\\[0.5mm]
{\small 5 Steine\quad Weiß am Zug}\hfill{\small Nr.~98\qquad\href{https://helpman.komtera.lt/?fen=8/8/8/1q5k/8/6q1/6P1/K7&moves=7.5}{\color{white}\bfseries lösen\,\ensuremath{\nearrow}}}
\end{tcolorbox}
\vspace{2mm}
\noindent\begin{minipage}[t]{36mm}\vspace{0pt}\sffamily\small\raggedright
\hmfigure{längste eindeutige}{{\Large h\#7.5}}\hmfigure{eindeutig bei dieser Länge}{{\Large 367} Stellungen}\hmfigure{längstes Matt}{{\Large h\#9.5}}\hmfigure{in der Klasse}{710 Millionen Stellungen,\\ 210 Millionen mit Hilfsmatt}
\end{minipage}\hfill
\begin{minipage}[t]{80mm}\vspace{0pt}\begin{tikzpicture}\begin{axis}[width=72mm,height=42mm,ybar,bar width=1.6pt,xmin=-1,xmax=20,ymin=0,ymax=44.71,axis y line*=left,axis x line*=bottom,xlabel={\scriptsize Halbzüge bis Matt},ylabel={\scriptsize Stellungen (Mio.)},tick label style={font=\tiny},label style={font=\scriptsize},enlarge x limits=false,ylabel shift=-3pt,xlabel shift=-3pt]\addplot[fill=ink!55,draw=none] coordinates {(1,0.004855) (2,0.02647) (3,0.3001) (4,0.8345) (5,5.039) (6,6.707) (7,22.34) (8,13.99) (9,34.96) (10,17.04) (11,40.64) (12,20.09) (13,35.23) (14,7.53) (15,5.365) (16,0.07308) (17,0.03639) (18,1e-06) (19,5e-06)};\draw[densely dotted] (axis cs:15,0) -- (axis cs:15,44.71);\end{axis}\begin{axis}[width=72mm,height=42mm,axis y line*=right,axis x line=none,ymode=log,xmin=-1,xmax=20,ylabel={\scriptsize eindeutig (log)},tick label style={font=\tiny},label style={font=\scriptsize},ylabel shift=-3pt]\addplot[mark=*,mark size=0.7pt,draw=black,thin] coordinates {(1,1479) (2,4806) (3,83798) (4,141536) (5,562564) (6,200075) (7,494897) (8,105576) (9,257064) (10,65256) (11,173146) (12,37152) (13,47822) (14,637) (15,367)};\end{axis}\end{tikzpicture}\end{minipage}
\vspace{3mm}

\noindent\begin{minipage}[t]{60mm}\vspace{0pt}\sbox\hmboard{\chessboard[setfen=8/8/8/1q5k/8/6q1/6P1/K7 w - - 0 1]}\usebox\hmboard\\[1mm]
\makebox[\wd\hmboard]{\sffamily\hspace{3mm}\textbf{h\#7.5}\hfill(2+3)}\end{minipage}\hfill
\begin{minipage}[t]{52mm}\vspace{0pt}
\textcolor{ink}{\sffamily\textsc{Lösung}}\par\smallskip\small
\begin{tabular}{@{}r@{\enspace}l@{\enspace}l@{}}
\textbf{1.} & \ldots & \symking{}a2\\
\textbf{2.} & \symqueen{}e3 & g4+\\
\textbf{3.} & \symking{}g6 & g5\\
\textbf{4.} & \symking{}f5 & g6\\
\textbf{5.} & \symking{}e4 & g7\\
\textbf{6.} & \symking{}d3 & g8=\symqueen{}\\
\textbf{7.} & \symqueen{}a5+ & \symking{}b3\\
\textbf{8.} & \symqueen{}ad2 & \symqueen{}c4\#\\
\end{tabular}\par\vspace{3mm}
\qrcode[height=12mm]{https://helpman.komtera.lt/?fen=8/8/8/1q5k/8/6q1/6P1/K7&moves=7.5}\\{\sffamily\tiny\color{ink}scannen zum Lösen}
\end{minipage}
\vspace{2mm}

{\raggedright\noindent\textcolor{ink}{\sffamily\textsc{Themen}}\enspace \themetag{Satzspiel} \themetag{Umwandlung} \themetag{Excelsior} \themetag{Selbstblock} \themetag{ohne Schlag} \themetag{Umnov}\par}

\clearpage
\phantomsection\label{hm:KPvkqr}
\begin{tcolorbox}[enhanced,colback=ink,colframe=ink,arc=0pt,boxrule=0pt,left=4mm,right=4mm,top=2.5mm,bottom=2.5mm]
\color{white}\sffamily{\LARGE\bfseries KPvkqr}\hfill{\Large h\#7.5}\\[0.5mm]
{\small 5 Steine\quad Weiß am Zug}\hfill{\small Nr.~99\qquad\href{https://helpman.komtera.lt/?fen=8/7k/8/1q6/K7/8/Pr6/8&moves=7.5}{\color{white}\bfseries lösen\,\ensuremath{\nearrow}}}
\end{tcolorbox}
\vspace{2mm}
\noindent\begin{minipage}[t]{36mm}\vspace{0pt}\sffamily\small\raggedright
\hmfigure{längste eindeutige}{{\Large h\#7.5}}\hmfigure{eindeutig bei dieser Länge}{{\Large 40} Stellungen}\hmfigure{längstes Matt}{{\Large h\#8.5}}\hmfigure{in der Klasse}{710 Millionen Stellungen,\\ 454 Millionen mit Hilfsmatt}
\end{minipage}\hfill
\begin{minipage}[t]{80mm}\vspace{0pt}\begin{tikzpicture}\begin{axis}[width=72mm,height=42mm,ybar,bar width=1.6pt,xmin=-1,xmax=18,ymin=0,ymax=94.93,axis y line*=left,axis x line*=bottom,xlabel={\scriptsize Halbzüge bis Matt},ylabel={\scriptsize Stellungen (Mio.)},tick label style={font=\tiny},label style={font=\scriptsize},enlarge x limits=false,ylabel shift=-3pt,xlabel shift=-3pt]\addplot[fill=ink!55,draw=none] coordinates {(1,0.0178) (2,0.09538) (3,0.9293) (4,3.021) (5,15.33) (6,25.91) (7,65.7) (8,44.45) (9,85.73) (10,43.71) (11,86.3) (12,41.16) (13,38.68) (14,1.98) (15,1.172) (16,0.000736) (17,0.007602)};\draw[densely dotted] (axis cs:15,0) -- (axis cs:15,94.93);\end{axis}\begin{axis}[width=72mm,height=42mm,axis y line*=right,axis x line=none,ymode=log,xmin=-1,xmax=18,ylabel={\scriptsize eindeutig (log)},tick label style={font=\tiny},label style={font=\scriptsize},ylabel shift=-3pt]\addplot[mark=*,mark size=0.7pt,draw=black,thin] coordinates {(1,5862) (2,27120) (3,289022) (4,679012) (5,1690711) (6,854002) (7,1324068) (8,300942) (9,534854) (10,143471) (11,170707) (12,11286) (13,6100) (14,12) (15,40)};\end{axis}\end{tikzpicture}\end{minipage}
\vspace{3mm}

\noindent\begin{minipage}[t]{60mm}\vspace{0pt}\sbox\hmboard{\chessboard[setfen=8/7k/8/1q6/K7/8/Pr6/8 w - - 0 1]}\usebox\hmboard\\[1mm]
\makebox[\wd\hmboard]{\sffamily\hspace{3mm}\textbf{h\#7.5}\hfill(2+3)}\end{minipage}\hfill
\begin{minipage}[t]{52mm}\vspace{0pt}
\textcolor{ink}{\sffamily\textsc{Lösung}}\par\smallskip\small
\begin{tabular}{@{}r@{\enspace}l@{\enspace}l@{}}
\textbf{1.} & \ldots & \symking{}a3\\
\textbf{2.} & \symqueen{}b3+ & axb3\\
\textbf{3.} & \symking{}g6 & b4\\
\textbf{4.} & \symking{}f5 & b5\\
\textbf{5.} & \symking{}e4 & b6\\
\textbf{6.} & \symking{}d3 & b7\\
\textbf{7.} & \symking{}c2 & b8=\symqueen{}\\
\textbf{8.} & \symking{}b1 & \symqueen{}xb2\#\\
\end{tabular}\par\vspace{3mm}
\qrcode[height=12mm]{https://helpman.komtera.lt/?fen=8/7k/8/1q6/K7/8/Pr6/8&moves=7.5}\\{\sffamily\tiny\color{ink}scannen zum Lösen}
\end{minipage}
\vspace{2mm}

{\raggedright\noindent\textcolor{ink}{\sffamily\textsc{Themen}}\enspace \themetag{Umwandlung} \themetag{Excelsior}\par\smallskip{\footnotesize\textit{Hinweis:} die Partei am Zug steht im Diagramm im Schach.}\par}

\clearpage
\phantomsection\label{hm:KPvkrb}
\begin{tcolorbox}[enhanced,colback=ink,colframe=ink,arc=0pt,boxrule=0pt,left=4mm,right=4mm,top=2.5mm,bottom=2.5mm]
\color{white}\sffamily{\LARGE\bfseries KPvkrb}\hfill{\Large h\#7.5}\\[0.5mm]
{\small 5 Steine\quad Weiß am Zug}\hfill{\small Nr.~100\qquad\href{https://helpman.komtera.lt/?fen=K7/8/8/8/1r6/b7/P6k/8&moves=7.5}{\color{white}\bfseries lösen\,\ensuremath{\nearrow}}}
\end{tcolorbox}
\vspace{2mm}
\noindent\begin{minipage}[t]{36mm}\vspace{0pt}\sffamily\small\raggedright
\hmfigure{längste eindeutige}{{\Large h\#7.5}}\hmfigure{eindeutig bei dieser Länge}{{\Large 6} Stellungen}\hmfigure{längstes Matt}{{\Large h\#8.5}}\hmfigure{in der Klasse}{710 Millionen Stellungen,\\ 512 Millionen mit Hilfsmatt}
\end{minipage}\hfill
\begin{minipage}[t]{80mm}\vspace{0pt}\begin{tikzpicture}\begin{axis}[width=72mm,height=42mm,ybar,bar width=1.6pt,xmin=-1,xmax=18,ymin=0,ymax=97.44,axis y line*=left,axis x line*=bottom,xlabel={\scriptsize Halbzüge bis Matt},ylabel={\scriptsize Stellungen (Mio.)},tick label style={font=\tiny},label style={font=\scriptsize},enlarge x limits=false,ylabel shift=-3pt,xlabel shift=-3pt]\addplot[fill=ink!55,draw=none] coordinates {(0,0.000165) (1,0.03907) (2,0.1608) (3,1.375) (4,4.159) (5,18.51) (6,32.86) (7,67.2) (8,57.3) (9,87.04) (10,58.59) (11,88.59) (12,55.63) (13,35.04) (14,5.167) (15,0.6487) (16,0.003014) (17,0.001026)};\draw[densely dotted] (axis cs:15,0) -- (axis cs:15,97.44);\end{axis}\begin{axis}[width=72mm,height=42mm,axis y line*=right,axis x line=none,ymode=log,xmin=-1,xmax=18,ylabel={\scriptsize eindeutig (log)},tick label style={font=\tiny},label style={font=\scriptsize},ylabel shift=-3pt]\addplot[mark=*,mark size=0.7pt,draw=black,thin] coordinates {(0,165) (1,11650) (2,45673) (3,414516) (4,934599) (5,1890107) (6,1105027) (7,1350753) (8,580654) (9,786505) (10,246497) (11,233150) (12,35224) (13,9473) (14,34) (15,6)};\end{axis}\end{tikzpicture}\end{minipage}
\vspace{3mm}

\noindent\begin{minipage}[t]{60mm}\vspace{0pt}\sbox\hmboard{\chessboard[setfen=K7/8/8/8/1r6/b7/P6k/8 w - - 0 1]}\usebox\hmboard\\[1mm]
\makebox[\wd\hmboard]{\sffamily\hspace{3mm}\textbf{h\#7.5}\hfill(2+3)}\end{minipage}\hfill
\begin{minipage}[t]{52mm}\vspace{0pt}
\textcolor{ink}{\sffamily\textsc{Lösung}}\par\smallskip\small
\begin{tabular}{@{}r@{\enspace}l@{\enspace}l@{}}
\textbf{1.} & \ldots & \symking{}a7\\
\textbf{2.} & \symrook{}e4 & \symking{}b6\\
\textbf{3.} & \symbishop{}e7 & a4\\
\textbf{4.} & \symking{}g3 & a5\\
\textbf{5.} & \symking{}f4 & a6\\
\textbf{6.} & \symking{}e5 & a7\\
\textbf{7.} & \symking{}d6 & a8=\symqueen{}\\
\textbf{8.} & \symrook{}e5 & \symqueen{}c6\#\\
\end{tabular}\par\vspace{3mm}
\qrcode[height=12mm]{https://helpman.komtera.lt/?fen=K7/8/8/8/1r6/b7/P6k/8&moves=7.5}\\{\sffamily\tiny\color{ink}scannen zum Lösen}
\end{minipage}
\vspace{2mm}

{\raggedright\noindent\textcolor{ink}{\sffamily\textsc{Themen}}\enspace \themetag{Satzspiel} \themetag{Umwandlung} \themetag{Excelsior} \themetag{Selbstblock} \themetag{ohne Schlag} \themetag{ohne Schach}\par}

\clearpage
\phantomsection\label{hm:KPvkrn}
\begin{tcolorbox}[enhanced,colback=ink,colframe=ink,arc=0pt,boxrule=0pt,left=4mm,right=4mm,top=2.5mm,bottom=2.5mm]
\color{white}\sffamily{\LARGE\bfseries KPvkrn}\hfill{\Large h\#7.5}\\[0.5mm]
{\small 5 Steine\quad Weiß am Zug}\hfill{\small Nr.~101\qquad\href{https://helpman.komtera.lt/?fen=8/7k/8/8/8/7r/1K5P/n7&moves=7.5}{\color{white}\bfseries lösen\,\ensuremath{\nearrow}}}
\end{tcolorbox}
\vspace{2mm}
\noindent\begin{minipage}[t]{36mm}\vspace{0pt}\sffamily\small\raggedright
\hmfigure{längste eindeutige}{{\Large h\#7.5}}\hmfigure{eindeutig bei dieser Länge}{{\Large 50} Stellungen}\hmfigure{längstes Matt}{{\Large h\#9.5}}\hmfigure{in der Klasse}{710 Millionen Stellungen,\\ 524 Millionen mit Hilfsmatt}
\end{minipage}\hfill
\begin{minipage}[t]{80mm}\vspace{0pt}\begin{tikzpicture}\begin{axis}[width=72mm,height=42mm,ybar,bar width=1.6pt,xmin=-1,xmax=20,ymin=0,ymax=96.94,axis y line*=left,axis x line*=bottom,xlabel={\scriptsize Halbzüge bis Matt},ylabel={\scriptsize Stellungen (Mio.)},tick label style={font=\tiny},label style={font=\scriptsize},enlarge x limits=false,ylabel shift=-3pt,xlabel shift=-3pt]\addplot[fill=ink!55,draw=none] coordinates {(0,8.2e-05) (1,0.04821) (2,0.1854) (3,1.477) (4,4.017) (5,16.87) (6,28.02) (7,57.64) (8,55.25) (9,83.74) (10,59.69) (11,88.13) (12,64.87) (13,48.83) (14,13.13) (15,1.764) (16,0.04686) (17,0.008883) (18,5e-06) (19,1e-06)};\draw[densely dotted] (axis cs:15,0) -- (axis cs:15,96.94);\end{axis}\begin{axis}[width=72mm,height=42mm,axis y line*=right,axis x line=none,ymode=log,xmin=-1,xmax=20,ylabel={\scriptsize eindeutig (log)},tick label style={font=\tiny},label style={font=\scriptsize},ylabel shift=-3pt]\addplot[mark=*,mark size=0.7pt,draw=black,thin] coordinates {(0,82) (1,12467) (2,45583) (3,384450) (4,810809) (5,1635215) (6,925841) (7,1108958) (8,508291) (9,701961) (10,296233) (11,366789) (12,86751) (13,27806) (14,327) (15,50)};\end{axis}\end{tikzpicture}\end{minipage}
\vspace{3mm}

\noindent\begin{minipage}[t]{60mm}\vspace{0pt}\sbox\hmboard{\chessboard[setfen=8/7k/8/8/8/7r/1K5P/n7 w - - 0 1]}\usebox\hmboard\\[1mm]
\makebox[\wd\hmboard]{\sffamily\hspace{3mm}\textbf{h\#7.5}\hfill(2+3)}\end{minipage}\hfill
\begin{minipage}[t]{52mm}\vspace{0pt}
\textcolor{ink}{\sffamily\textsc{Lösung}}\par\smallskip\small
\begin{tabular}{@{}r@{\enspace}l@{\enspace}l@{}}
\textbf{1.} & \ldots & \symking{}a2\\
\textbf{2.} & \symrook{}a3+ & \symking{}xa3\\
\textbf{3.} & \symking{}g6 & h4\\
\textbf{4.} & \symking{}f5 & h5\\
\textbf{5.} & \symking{}e4 & h6\\
\textbf{6.} & \symking{}d3 & h7\\
\textbf{7.} & \symking{}c2 & h8=\symqueen{}\\
\textbf{8.} & \symking{}b1 & \symqueen{}b2\#\\
\end{tabular}\par\vspace{3mm}
\qrcode[height=12mm]{https://helpman.komtera.lt/?fen=8/7k/8/8/8/7r/1K5P/n7&moves=7.5}\\{\sffamily\tiny\color{ink}scannen zum Lösen}
\end{minipage}
\vspace{2mm}

{\raggedright\noindent\textcolor{ink}{\sffamily\textsc{Themen}}\enspace \themetag{Satzspiel} \themetag{Umwandlung} \themetag{Excelsior}\par}

\clearpage
\phantomsection\label{hm:KBBBvk}
\begin{tcolorbox}[enhanced,colback=ink,colframe=ink,arc=0pt,boxrule=0pt,left=4mm,right=4mm,top=2.5mm,bottom=2.5mm]
\color{white}\sffamily{\LARGE\bfseries KBBBvk}\hfill{\Large h\#7}\\[0.5mm]
{\small 5 Steine\quad Schwarz am Zug}\hfill{\small Nr.~102\qquad\href{https://helpman.komtera.lt/?fen=5B2/5B1k/8/8/8/8/7B/K7&moves=7}{\color{white}\bfseries lösen\,\ensuremath{\nearrow}}}
\end{tcolorbox}
\vspace{2mm}
\noindent\begin{minipage}[t]{36mm}\vspace{0pt}\sffamily\small\raggedright
\hmfigure{längste eindeutige}{{\Large h\#7}}\hmfigure{eindeutig bei dieser Länge}{{\Large 3} Stellungen}\hmfigure{längstes Matt}{{\Large h\#8}}\hmfigure{in der Klasse}{242 Millionen Stellungen,\\ 21 Millionen mit Hilfsmatt}
\end{minipage}\hfill
\begin{minipage}[t]{80mm}\vspace{0pt}\begin{tikzpicture}\begin{axis}[width=72mm,height=42mm,ybar,bar width=1.6pt,xmin=-1,xmax=17,ymin=0,ymax=6.956,axis y line*=left,axis x line*=bottom,xlabel={\scriptsize Halbzüge bis Matt},ylabel={\scriptsize Stellungen (Mio.)},tick label style={font=\tiny},label style={font=\scriptsize},enlarge x limits=false,ylabel shift=-3pt,xlabel shift=-3pt]\addplot[fill=ink!55,draw=none] coordinates {(0,0.02149) (1,0.07132) (2,0.1718) (3,0.6649) (4,2.735) (5,3.53) (6,6.324) (7,2.844) (8,2.907) (9,0.7878) (10,0.6427) (11,0.1742) (12,0.1231) (13,0.01763) (14,0.009309) (15,4.4e-05) (16,5e-06)};\draw[densely dotted] (axis cs:14,0) -- (axis cs:14,6.956);\end{axis}\begin{axis}[width=72mm,height=42mm,axis y line*=right,axis x line=none,ymode=log,xmin=-1,xmax=17,ylabel={\scriptsize eindeutig (log)},tick label style={font=\tiny},label style={font=\scriptsize},ylabel shift=-3pt]\addplot[mark=*,mark size=0.7pt,draw=black,thin] coordinates {(0,21487) (1,62657) (2,151832) (3,374652) (4,863478) (5,128782) (6,111965) (7,17471) (8,7168) (9,1851) (10,709) (11,303) (12,56) (13,3) (14,3)};\end{axis}\end{tikzpicture}\end{minipage}
\vspace{3mm}

\noindent\begin{minipage}[t]{60mm}\vspace{0pt}\sbox\hmboard{\chessboard[setfen=5B2/5B1k/8/8/8/8/7B/K7 b - - 0 1]}\usebox\hmboard\\[1mm]
\makebox[\wd\hmboard]{\sffamily\hspace{3mm}\textbf{h\#7}\hfill(4+1)}\end{minipage}\hfill
\begin{minipage}[t]{52mm}\vspace{0pt}
\textcolor{ink}{\sffamily\textsc{Lösung}}\par\smallskip\small
\begin{tabular}{@{}r@{\enspace}l@{\enspace}l@{}}
\textbf{1.} & \symking{}h8 & \symking{}b2\\
\textbf{2.} & \symking{}h7 & \symking{}c3\\
\textbf{3.} & \symking{}h8 & \symking{}d4\\
\textbf{4.} & \symking{}h7 & \symking{}e5\\
\textbf{5.} & \symking{}h8 & \symking{}f6\\
\textbf{6.} & \symking{}h7 & \symbishop{}e5\\
\textbf{7.} & \symking{}h8 & \symking{}g6\#\\
\end{tabular}\par\vspace{3mm}
\qrcode[height=12mm]{https://helpman.komtera.lt/?fen=5B2/5B1k/8/8/8/8/7B/K7&moves=7}\\{\sffamily\tiny\color{ink}scannen zum Lösen}
\end{minipage}
\vspace{2mm}

{\raggedright\noindent\textcolor{ink}{\sffamily\textsc{Themen}}\enspace \themetag{Satzspiel} \themetag{Spiegelmatt} \themetag{Switchback} \themetag{ein Stein zieht} \themetag{Pendel} \themetag{ohne Schlag} \themetag{ohne Schach}\par}

\clearpage
\phantomsection\label{hm:KBBNvk}
\begin{tcolorbox}[enhanced,colback=ink,colframe=ink,arc=0pt,boxrule=0pt,left=4mm,right=4mm,top=2.5mm,bottom=2.5mm]
\color{white}\sffamily{\LARGE\bfseries KBBNvk}\hfill{\Large h\#7}\\[0.5mm]
{\small 5 Steine\quad Schwarz am Zug}\hfill{\small Nr.~103\qquad\href{https://helpman.komtera.lt/?fen=5B1N/7k/8/8/8/1B6/8/K7&moves=7}{\color{white}\bfseries lösen\,\ensuremath{\nearrow}}}
\end{tcolorbox}
\vspace{2mm}
\noindent\begin{minipage}[t]{36mm}\vspace{0pt}\sffamily\small\raggedright
\hmfigure{längste eindeutige}{{\Large h\#7}}\hmfigure{eindeutig bei dieser Länge}{{\Large 1} Stellung}\hmfigure{längstes Matt}{{\Large h\#8}}\hmfigure{in der Klasse}{242 Millionen Stellungen,\\ 86,5 Millionen mit Hilfsmatt}
\end{minipage}\hfill
\begin{minipage}[t]{80mm}\vspace{0pt}\begin{tikzpicture}\begin{axis}[width=72mm,height=42mm,ybar,bar width=1.6pt,xmin=-1,xmax=17,ymin=0,ymax=21.22,axis y line*=left,axis x line*=bottom,xlabel={\scriptsize Halbzüge bis Matt},ylabel={\scriptsize Stellungen (Mio.)},tick label style={font=\tiny},label style={font=\scriptsize},enlarge x limits=false,ylabel shift=-3pt,xlabel shift=-3pt]\addplot[fill=ink!55,draw=none] coordinates {(0,0.0218) (1,0.07924) (2,0.1816) (3,0.875) (4,3.481) (5,8) (6,16.02) (7,13.75) (8,19.29) (9,9.628) (10,10.22) (11,2.666) (12,1.772) (13,0.3409) (14,0.1236) (15,0.01886) (16,4e-05)};\draw[densely dotted] (axis cs:14,0) -- (axis cs:14,21.22);\end{axis}\begin{axis}[width=72mm,height=42mm,axis y line*=right,axis x line=none,ymode=log,xmin=-1,xmax=17,ylabel={\scriptsize eindeutig (log)},tick label style={font=\tiny},label style={font=\scriptsize},ylabel shift=-3pt]\addplot[mark=*,mark size=0.7pt,draw=black,thin] coordinates {(0,21799) (1,72711) (2,165948) (3,608383) (4,1362284) (5,748009) (6,826303) (7,163432) (8,134734) (9,12347) (10,5789) (11,407) (12,48) (13,1) (14,1)};\end{axis}\end{tikzpicture}\end{minipage}
\vspace{3mm}

\noindent\begin{minipage}[t]{60mm}\vspace{0pt}\sbox\hmboard{\chessboard[setfen=5B1N/7k/8/8/8/1B6/8/K7 b - - 0 1]}\usebox\hmboard\\[1mm]
\makebox[\wd\hmboard]{\sffamily\hspace{3mm}\textbf{h\#7}\hfill(4+1)}\end{minipage}\hfill
\begin{minipage}[t]{52mm}\vspace{0pt}
\textcolor{ink}{\sffamily\textsc{Lösung}}\par\smallskip\small
\begin{tabular}{@{}r@{\enspace}l@{\enspace}l@{}}
\textbf{1.} & \symking{}xh8 & \symking{}b2\\
\textbf{2.} & \symking{}h7 & \symking{}c3\\
\textbf{3.} & \symking{}g6 & \symking{}c4\\
\textbf{4.} & \symking{}f7 & \symking{}d5\\
\textbf{5.} & \symking{}e8 & \symking{}e6\\
\textbf{6.} & \symking{}d8 & \symbishop{}e7+\\
\textbf{7.} & \symking{}e8 & \symbishop{}a4\#\\
\end{tabular}\par\vspace{3mm}
\qrcode[height=12mm]{https://helpman.komtera.lt/?fen=5B1N/7k/8/8/8/1B6/8/K7&moves=7}\\{\sffamily\tiny\color{ink}scannen zum Lösen}
\end{minipage}
\vspace{2mm}

{\raggedright\noindent\textcolor{ink}{\sffamily\textsc{Themen}}\enspace \themetag{Switchback} \themetag{ein Stein zieht}\par\smallskip{\footnotesize\textit{Hinweis:} die Lösung beginnt mit einem Schlag.}\par}

\clearpage
\phantomsection\label{hm:KBBPvk}
\begin{tcolorbox}[enhanced,colback=ink,colframe=ink,arc=0pt,boxrule=0pt,left=4mm,right=4mm,top=2.5mm,bottom=2.5mm]
\color{white}\sffamily{\LARGE\bfseries KBBPvk}\hfill{\Large h\#7}\\[0.5mm]
{\small 5 Steine\quad Schwarz am Zug}\hfill{\small Nr.~104\qquad\href{https://helpman.komtera.lt/?fen=1B6/8/8/8/7B/8/7P/K5k1&moves=7}{\color{white}\bfseries lösen\,\ensuremath{\nearrow}}}
\end{tcolorbox}
\vspace{2mm}
\noindent\begin{minipage}[t]{36mm}\vspace{0pt}\sffamily\small\raggedright
\hmfigure{längste eindeutige}{{\Large h\#7}}\hmfigure{eindeutig bei dieser Länge}{{\Large 556} Stellungen}\hmfigure{längstes Matt}{{\Large h\#9}}\hmfigure{in der Klasse}{710 Millionen Stellungen,\\ 266 Millionen mit Hilfsmatt}
\end{minipage}\hfill
\begin{minipage}[t]{80mm}\vspace{0pt}\begin{tikzpicture}\begin{axis}[width=72mm,height=42mm,ybar,bar width=1.6pt,xmin=-1,xmax=19,ymin=0,ymax=54.94,axis y line*=left,axis x line*=bottom,xlabel={\scriptsize Halbzüge bis Matt},ylabel={\scriptsize Stellungen (Mio.)},tick label style={font=\tiny},label style={font=\scriptsize},enlarge x limits=false,ylabel shift=-3pt,xlabel shift=-3pt]\addplot[fill=ink!55,draw=none] coordinates {(0,0.04954) (1,0.2844) (2,0.6687) (3,4.354) (4,14.81) (5,31.89) (6,49.94) (7,31.88) (8,38.42) (9,19.77) (10,21.31) (11,16.33) (12,20.84) (13,7.363) (14,7.211) (15,0.2511) (16,0.175) (17,0.00094) (18,0.000518)};\draw[densely dotted] (axis cs:14,0) -- (axis cs:14,54.94);\end{axis}\begin{axis}[width=72mm,height=42mm,axis y line*=right,axis x line=none,ymode=log,xmin=-1,xmax=19,ylabel={\scriptsize eindeutig (log)},tick label style={font=\tiny},label style={font=\scriptsize},ylabel shift=-3pt]\addplot[mark=*,mark size=0.7pt,draw=black,thin] coordinates {(0,49539) (1,195777) (2,463360) (3,2548704) (4,4659935) (5,2777072) (6,2258871) (7,789040) (8,470551) (9,281422) (10,202724) (11,171889) (12,111205) (13,1431) (14,556)};\end{axis}\end{tikzpicture}\end{minipage}
\vspace{3mm}

\noindent\begin{minipage}[t]{60mm}\vspace{0pt}\sbox\hmboard{\chessboard[setfen=1B6/8/8/8/7B/8/7P/K5k1 b - - 0 1]}\usebox\hmboard\\[1mm]
\makebox[\wd\hmboard]{\sffamily\hspace{3mm}\textbf{h\#7}\hfill(4+1)}\end{minipage}\hfill
\begin{minipage}[t]{52mm}\vspace{0pt}
\textcolor{ink}{\sffamily\textsc{Lösung}}\par\smallskip\small
\begin{tabular}{@{}r@{\enspace}l@{\enspace}l@{}}
\textbf{1.} & \symking{}g2 & \symbishop{}hg3\\
\textbf{2.} & \symking{}f3 & h4\\
\textbf{3.} & \symking{}e4 & h5\\
\textbf{4.} & \symking{}d5 & h6\\
\textbf{5.} & \symking{}c6 & h7\\
\textbf{6.} & \symking{}b7 & h8=\symqueen{}\\
\textbf{7.} & \symking{}a8 & \symqueen{}h1\#\\
\end{tabular}\par\vspace{3mm}
\qrcode[height=12mm]{https://helpman.komtera.lt/?fen=1B6/8/8/8/7B/8/7P/K5k1&moves=7}\\{\sffamily\tiny\color{ink}scannen zum Lösen}
\end{minipage}
\vspace{2mm}

{\raggedright\noindent\textcolor{ink}{\sffamily\textsc{Themen}}\enspace \themetag{reines Matt} \themetag{Mustermatt} \themetag{Umwandlung} \themetag{Excelsior} \themetag{ein Stein zieht} \themetag{ohne Schlag} \themetag{ohne Schach}\par}

\clearpage
\phantomsection\label{hm:KBBvkb}
\begin{tcolorbox}[enhanced,colback=ink,colframe=ink,arc=0pt,boxrule=0pt,left=4mm,right=4mm,top=2.5mm,bottom=2.5mm]
\color{white}\sffamily{\LARGE\bfseries KBBvkb}\hfill{\Large h\#7}\\[0.5mm]
{\small 5 Steine\quad Schwarz am Zug}\hfill{\small Nr.~105\qquad\href{https://helpman.komtera.lt/?fen=8/8/6b1/8/8/8/3B3k/K3B3&moves=7}{\color{white}\bfseries lösen\,\ensuremath{\nearrow}}}
\end{tcolorbox}
\vspace{2mm}
\noindent\begin{minipage}[t]{36mm}\vspace{0pt}\sffamily\small\raggedright
\hmfigure{längste eindeutige}{{\Large h\#7}}\hmfigure{eindeutig bei dieser Länge}{{\Large 95} Stellungen}\hmfigure{längstes Matt}{{\Large h\#10.5}}\hmfigure{in der Klasse}{242 Millionen Stellungen,\\ 62,7 Millionen mit Hilfsmatt}
\end{minipage}\hfill
\begin{minipage}[t]{80mm}\vspace{0pt}\begin{tikzpicture}\begin{axis}[width=72mm,height=42mm,ybar,bar width=1.6pt,xmin=-1,xmax=22,ymin=0,ymax=11.57,axis y line*=left,axis x line*=bottom,xlabel={\scriptsize Halbzüge bis Matt},ylabel={\scriptsize Stellungen (Mio.)},tick label style={font=\tiny},label style={font=\scriptsize},enlarge x limits=false,ylabel shift=-3pt,xlabel shift=-3pt]\addplot[fill=ink!55,draw=none] coordinates {(0,0.009628) (1,0.03888) (2,0.182) (3,1.018) (4,3.554) (5,6.519) (6,10.52) (7,8.331) (8,8.45) (9,5.071) (10,4.069) (11,3.428) (12,4.207) (13,2.257) (14,1.827) (15,1.303) (16,0.7514) (17,0.6702) (18,0.2254) (19,0.2413) (20,0.001651) (21,1e-06)};\draw[densely dotted] (axis cs:14,0) -- (axis cs:14,11.57);\end{axis}\begin{axis}[width=72mm,height=42mm,axis y line*=right,axis x line=none,ymode=log,xmin=-1,xmax=22,ylabel={\scriptsize eindeutig (log)},tick label style={font=\tiny},label style={font=\scriptsize},ylabel shift=-3pt]\addplot[mark=*,mark size=0.7pt,draw=black,thin] coordinates {(0,9628) (1,37830) (2,142716) (3,548312) (4,795381) (5,357085) (6,228809) (7,70314) (8,30116) (9,14967) (10,11782) (11,4684) (12,1988) (13,322) (14,95)};\end{axis}\end{tikzpicture}\end{minipage}
\vspace{3mm}

\noindent\begin{minipage}[t]{60mm}\vspace{0pt}\sbox\hmboard{\chessboard[setfen=8/8/6b1/8/8/8/3B3k/K3B3 b - - 0 1]}\usebox\hmboard\\[1mm]
\makebox[\wd\hmboard]{\sffamily\hspace{3mm}\textbf{h\#7}\hfill(3+2)}\end{minipage}\hfill
\begin{minipage}[t]{52mm}\vspace{0pt}
\textcolor{ink}{\sffamily\textsc{Lösung}}\par\smallskip\small
\begin{tabular}{@{}r@{\enspace}l@{\enspace}l@{}}
\textbf{1.} & \symking{}h3 & \symking{}b2\\
\textbf{2.} & \symking{}g4 & \symking{}c3\\
\textbf{3.} & \symking{}f5 & \symking{}d4\\
\textbf{4.} & \symking{}f6 & \symbishop{}c3\\
\textbf{5.} & \symking{}g7 & \symking{}e5\\
\textbf{6.} & \symking{}h8 & \symking{}f6\\
\textbf{7.} & \symbishop{}h7 & \symking{}f7\#\\
\end{tabular}\par\vspace{3mm}
\qrcode[height=12mm]{https://helpman.komtera.lt/?fen=8/8/6b1/8/8/8/3B3k/K3B3&moves=7}\\{\sffamily\tiny\color{ink}scannen zum Lösen}
\end{minipage}
\vspace{2mm}

{\raggedright\noindent\textcolor{ink}{\sffamily\textsc{Themen}}\enspace \themetag{Selbstblock} \themetag{ohne Schlag} \themetag{ohne Schach}\par}

\clearpage
\phantomsection\label{hm:KBNPvk}
\begin{tcolorbox}[enhanced,colback=ink,colframe=ink,arc=0pt,boxrule=0pt,left=4mm,right=4mm,top=2.5mm,bottom=2.5mm]
\color{white}\sffamily{\LARGE\bfseries KBNPvk}\hfill{\Large h\#7}\\[0.5mm]
{\small 5 Steine\quad Schwarz am Zug}\hfill{\small Nr.~106\qquad\href{https://helpman.komtera.lt/?fen=2K5/8/8/8/8/4N1Bk/4P3/8&moves=7}{\color{white}\bfseries lösen\,\ensuremath{\nearrow}}}
\end{tcolorbox}
\vspace{2mm}
\noindent\begin{minipage}[t]{36mm}\vspace{0pt}\sffamily\small\raggedright
\hmfigure{längste eindeutige}{{\Large h\#7}}\hmfigure{eindeutig bei dieser Länge}{{\Large 9} Stellungen}\hmfigure{längstes Matt}{{\Large h\#8}}\hmfigure{in der Klasse}{710 Millionen Stellungen,\\ 546 Millionen mit Hilfsmatt}
\end{minipage}\hfill
\begin{minipage}[t]{80mm}\vspace{0pt}\begin{tikzpicture}\begin{axis}[width=72mm,height=42mm,ybar,bar width=1.6pt,xmin=-1,xmax=17,ymin=0,ymax=118.7,axis y line*=left,axis x line*=bottom,xlabel={\scriptsize Halbzüge bis Matt},ylabel={\scriptsize Stellungen (Mio.)},tick label style={font=\tiny},label style={font=\scriptsize},enlarge x limits=false,ylabel shift=-3pt,xlabel shift=-3pt]\addplot[fill=ink!55,draw=none] coordinates {(0,0.02354) (1,0.2758) (2,0.6371) (3,5.363) (4,16.99) (5,49.28) (6,73.67) (7,75.57) (8,107.9) (9,77.08) (10,86.34) (11,29.18) (12,20.76) (13,2.289) (14,0.7995) (15,0.01161) (16,0.001114)};\draw[densely dotted] (axis cs:14,0) -- (axis cs:14,118.7);\end{axis}\begin{axis}[width=72mm,height=42mm,axis y line*=right,axis x line=none,ymode=log,xmin=-1,xmax=17,ylabel={\scriptsize eindeutig (log)},tick label style={font=\tiny},label style={font=\scriptsize},ylabel shift=-3pt]\addplot[mark=*,mark size=0.7pt,draw=black,thin] coordinates {(0,23539) (1,129030) (2,330373) (3,2948971) (4,5646510) (5,5866712) (6,5079565) (7,2466592) (8,1992242) (9,482102) (10,270645) (11,49352) (12,10220) (13,47) (14,9)};\end{axis}\end{tikzpicture}\end{minipage}
\vspace{3mm}

\noindent\begin{minipage}[t]{60mm}\vspace{0pt}\sbox\hmboard{\chessboard[setfen=2K5/8/8/8/8/4N1Bk/4P3/8 b - - 0 1]}\usebox\hmboard\\[1mm]
\makebox[\wd\hmboard]{\sffamily\hspace{3mm}\textbf{h\#7}\hfill(4+1)}\end{minipage}\hfill
\begin{minipage}[t]{52mm}\vspace{0pt}
\textcolor{ink}{\sffamily\textsc{Lösung}}\par\smallskip\small
\begin{tabular}{@{}r@{\enspace}l@{\enspace}l@{}}
\textbf{1.} & \symking{}xg3 & \symknight{}d5\\
\textbf{2.} & \symking{}f2 & e4\\
\textbf{3.} & \symking{}e2 & e5\\
\textbf{4.} & \symking{}d3 & e6\\
\textbf{5.} & \symking{}c4 & e7\\
\textbf{6.} & \symking{}b5 & e8=\symqueen{}+\\
\textbf{7.} & \symking{}a6 & \symqueen{}a4\#\\
\end{tabular}\par\vspace{3mm}
\qrcode[height=12mm]{https://helpman.komtera.lt/?fen=2K5/8/8/8/8/4N1Bk/4P3/8&moves=7}\\{\sffamily\tiny\color{ink}scannen zum Lösen}
\end{minipage}
\vspace{2mm}

{\raggedright\noindent\textcolor{ink}{\sffamily\textsc{Themen}}\enspace \themetag{reines Matt} \themetag{Mustermatt} \themetag{Idealmatt} \themetag{Spiegelmatt} \themetag{Umwandlung} \themetag{Excelsior} \themetag{ein Stein zieht} \themetag{Umnov}\par\smallskip{\footnotesize\textit{Hinweis:} die Lösung beginnt mit einem Schlag.}\par}

\clearpage
\phantomsection\label{hm:KBNvkb}
\begin{tcolorbox}[enhanced,colback=ink,colframe=ink,arc=0pt,boxrule=0pt,left=4mm,right=4mm,top=2.5mm,bottom=2.5mm]
\color{white}\sffamily{\LARGE\bfseries KBNvkb}\hfill{\Large h\#7}\\[0.5mm]
{\small 5 Steine\quad Schwarz am Zug}\hfill{\small Nr.~107\qquad\href{https://helpman.komtera.lt/?fen=7k/8/8/7N/8/8/8/K2B3b&moves=7}{\color{white}\bfseries lösen\,\ensuremath{\nearrow}}}
\end{tcolorbox}
\vspace{2mm}
\noindent\begin{minipage}[t]{36mm}\vspace{0pt}\sffamily\small\raggedright
\hmfigure{längste eindeutige}{{\Large h\#7}}\hmfigure{eindeutig bei dieser Länge}{{\Large 7} Stellungen}\hmfigure{längstes Matt}{{\Large h\#11}}\hmfigure{in der Klasse}{242 Millionen Stellungen,\\ 171 Millionen mit Hilfsmatt}
\end{minipage}\hfill
\begin{minipage}[t]{80mm}\vspace{0pt}\begin{tikzpicture}\begin{axis}[width=72mm,height=42mm,ybar,bar width=1.6pt,xmin=-1,xmax=23,ymin=0,ymax=40.59,axis y line*=left,axis x line*=bottom,xlabel={\scriptsize Halbzüge bis Matt},ylabel={\scriptsize Stellungen (Mio.)},tick label style={font=\tiny},label style={font=\scriptsize},enlarge x limits=false,ylabel shift=-3pt,xlabel shift=-3pt]\addplot[fill=ink!55,draw=none] coordinates {(0,0.003993) (1,0.01558) (2,0.07573) (3,0.4329) (4,1.845) (5,6.582) (6,14.88) (7,25.8) (8,36.9) (9,33.77) (10,28.28) (11,13.49) (12,6.566) (13,1.697) (14,0.4945) (15,0.06407) (16,0.002237) (17,2e-06) (18,0.001071) (20,0.000179) (22,1e-06)};\draw[densely dotted] (axis cs:14,0) -- (axis cs:14,40.59);\end{axis}\begin{axis}[width=72mm,height=42mm,axis y line*=right,axis x line=none,ymode=log,xmin=-1,xmax=23,ylabel={\scriptsize eindeutig (log)},tick label style={font=\tiny},label style={font=\scriptsize},ylabel shift=-3pt]\addplot[mark=*,mark size=0.7pt,draw=black,thin] coordinates {(0,3993) (1,14881) (2,63251) (3,212747) (4,567864) (5,810806) (6,714551) (7,396705) (8,295436) (9,81756) (10,29516) (11,3172) (12,1366) (13,74) (14,7)};\end{axis}\end{tikzpicture}\end{minipage}
\vspace{3mm}

\noindent\begin{minipage}[t]{60mm}\vspace{0pt}\sbox\hmboard{\chessboard[setfen=7k/8/8/7N/8/8/8/K2B3b b - - 0 1]}\usebox\hmboard\\[1mm]
\makebox[\wd\hmboard]{\sffamily\hspace{3mm}\textbf{h\#7}\hfill(3+2)}\end{minipage}\hfill
\begin{minipage}[t]{52mm}\vspace{0pt}
\textcolor{ink}{\sffamily\textsc{Lösung}}\par\smallskip\small
\begin{tabular}{@{}r@{\enspace}l@{\enspace}l@{}}
\textbf{1.} & \symking{}h7 & \symking{}b2\\
\textbf{2.} & \symking{}g6 & \symknight{}f4+\\
\textbf{3.} & \symking{}f5 & \symknight{}d3\\
\textbf{4.} & \symking{}e4 & \symbishop{}f3+\\
\textbf{5.} & \symking{}e3 & \symbishop{}xh1\\
\textbf{6.} & \symking{}e2 & \symking{}c3\\
\textbf{7.} & \symking{}d1 & \symbishop{}f3\#\\
\end{tabular}\par\vspace{3mm}
\qrcode[height=12mm]{https://helpman.komtera.lt/?fen=7k/8/8/7N/8/8/8/K2B3b&moves=7}\\{\sffamily\tiny\color{ink}scannen zum Lösen}
\end{minipage}
\vspace{2mm}

{\raggedright\noindent\textcolor{ink}{\sffamily\textsc{Themen}}\enspace \themetag{reines Matt} \themetag{Mustermatt} \themetag{Idealmatt} \themetag{Spiegelmatt} \themetag{Switchback} \themetag{ein Stein zieht}\par}

\clearpage
\phantomsection\label{hm:KBNvkn}
\begin{tcolorbox}[enhanced,colback=ink,colframe=ink,arc=0pt,boxrule=0pt,left=4mm,right=4mm,top=2.5mm,bottom=2.5mm]
\color{white}\sffamily{\LARGE\bfseries KBNvkn}\hfill{\Large h\#7}\\[0.5mm]
{\small 5 Steine\quad Schwarz am Zug}\hfill{\small Nr.~108\qquad\href{https://helpman.komtera.lt/?fen=7k/7N/8/8/2n5/8/8/K4B2&moves=7}{\color{white}\bfseries lösen\,\ensuremath{\nearrow}}}
\end{tcolorbox}
\vspace{2mm}
\noindent\begin{minipage}[t]{36mm}\vspace{0pt}\sffamily\small\raggedright
\hmfigure{längste eindeutige}{{\Large h\#7}}\hmfigure{eindeutig bei dieser Länge}{{\Large 7} Stellungen}\hmfigure{längstes Matt}{{\Large h\#12}}\hmfigure{in der Klasse}{242 Millionen Stellungen,\\ 176 Millionen mit Hilfsmatt}
\end{minipage}\hfill
\begin{minipage}[t]{80mm}\vspace{0pt}\begin{tikzpicture}\begin{axis}[width=72mm,height=42mm,ybar,bar width=1.6pt,xmin=-1,xmax=25,ymin=0,ymax=40.3,axis y line*=left,axis x line*=bottom,xlabel={\scriptsize Halbzüge bis Matt},ylabel={\scriptsize Stellungen (Mio.)},tick label style={font=\tiny},label style={font=\scriptsize},enlarge x limits=false,ylabel shift=-3pt,xlabel shift=-3pt]\addplot[fill=ink!55,draw=none] coordinates {(0,0.00274) (1,0.01106) (2,0.03553) (3,0.2587) (4,0.8609) (5,4.039) (6,9.796) (7,21) (8,35.54) (9,35.56) (10,36.64) (11,18.19) (12,10.3) (13,2.679) (14,0.8385) (15,0.1078) (16,0.003473) (17,5e-06) (18,0.00013) (20,1e-06) (22,4.4e-05) (24,1e-06)};\draw[densely dotted] (axis cs:14,0) -- (axis cs:14,40.3);\end{axis}\begin{axis}[width=72mm,height=42mm,axis y line*=right,axis x line=none,ymode=log,xmin=-1,xmax=25,ylabel={\scriptsize eindeutig (log)},tick label style={font=\tiny},label style={font=\scriptsize},ylabel shift=-3pt]\addplot[mark=*,mark size=0.7pt,draw=black,thin] coordinates {(0,2740) (1,10418) (2,23326) (3,120727) (4,280782) (5,510547) (6,615941) (7,350184) (8,314850) (9,86068) (10,44713) (11,4393) (12,1685) (13,96) (14,7)};\end{axis}\end{tikzpicture}\end{minipage}
\vspace{3mm}

\noindent\begin{minipage}[t]{60mm}\vspace{0pt}\sbox\hmboard{\chessboard[setfen=7k/7N/8/8/2n5/8/8/K4B2 b - - 0 1]}\usebox\hmboard\\[1mm]
\makebox[\wd\hmboard]{\sffamily\hspace{3mm}\textbf{h\#7}\hfill(3+2)}\end{minipage}\hfill
\begin{minipage}[t]{52mm}\vspace{0pt}
\textcolor{ink}{\sffamily\textsc{Lösung}}\par\smallskip\small
\begin{tabular}{@{}r@{\enspace}l@{\enspace}l@{}}
\textbf{1.} & \symking{}g7 & \symknight{}g5\\
\textbf{2.} & \symking{}f6 & \symking{}b1\\
\textbf{3.} & \symking{}e5 & \symking{}c1\\
\textbf{4.} & \symking{}d4 & \symknight{}f3+\\
\textbf{5.} & \symking{}c3 & \symknight{}e1\\
\textbf{6.} & \symking{}b3 & \symknight{}c2\\
\textbf{7.} & \symking{}a2 & \symbishop{}xc4\#\\
\end{tabular}\par\vspace{3mm}
\qrcode[height=12mm]{https://helpman.komtera.lt/?fen=7k/7N/8/8/2n5/8/8/K4B2&moves=7}\\{\sffamily\tiny\color{ink}scannen zum Lösen}
\end{minipage}
\vspace{2mm}

{\raggedright\noindent\textcolor{ink}{\sffamily\textsc{Themen}}\enspace \themetag{reines Matt} \themetag{Mustermatt} \themetag{Idealmatt} \themetag{Spiegelmatt} \themetag{ein Stein zieht}\par}

\clearpage
\phantomsection\label{hm:KBNvkp}
\begin{tcolorbox}[enhanced,colback=ink,colframe=ink,arc=0pt,boxrule=0pt,left=4mm,right=4mm,top=2.5mm,bottom=2.5mm]
\color{white}\sffamily{\LARGE\bfseries KBNvkp}\hfill{\Large h\#7}\\[0.5mm]
{\small 5 Steine\quad Schwarz am Zug}\hfill{\small Nr.~109\qquad\href{https://helpman.komtera.lt/?fen=7N/7k/8/8/2p5/8/8/K5B1&moves=7}{\color{white}\bfseries lösen\,\ensuremath{\nearrow}}}
\end{tcolorbox}
\vspace{2mm}
\noindent\begin{minipage}[t]{36mm}\vspace{0pt}\sffamily\small\raggedright
\hmfigure{längste eindeutige}{{\Large h\#7}}\hmfigure{eindeutig bei dieser Länge}{{\Large 50} Stellungen}\hmfigure{längstes Matt}{{\Large h\#16}}\hmfigure{in der Klasse}{710 Millionen Stellungen,\\ 545 Millionen mit Hilfsmatt}
\end{minipage}\hfill
\begin{minipage}[t]{80mm}\vspace{0pt}\begin{tikzpicture}\begin{axis}[width=72mm,height=42mm,ybar,bar width=1.6pt,xmin=-1,xmax=33,ymin=0,ymax=133.1,axis y line*=left,axis x line*=bottom,xlabel={\scriptsize Halbzüge bis Matt},ylabel={\scriptsize Stellungen (Mio.)},tick label style={font=\tiny},label style={font=\scriptsize},enlarge x limits=false,ylabel shift=-3pt,xlabel shift=-3pt]\addplot[fill=ink!55,draw=none] coordinates {(0,0.01354) (1,0.0532) (2,0.135) (3,0.8661) (4,2.712) (5,12.29) (6,28.74) (7,61.41) (8,107.4) (9,107) (10,121) (11,56.14) (12,36.1) (13,8.042) (14,2.71) (15,0.2769) (16,0.01155) (17,1.4e-05) (18,0.00583) (19,2.9e-05) (20,0.003139) (21,1.8e-05) (22,0.002002) (23,1.4e-05) (24,0.001049) (25,7e-06) (26,0.000427) (27,4e-06) (28,0.000289) (29,3e-06) (30,4e-05) (32,1.6e-05)};\draw[densely dotted] (axis cs:14,0) -- (axis cs:14,133.1);\end{axis}\begin{axis}[width=72mm,height=42mm,axis y line*=right,axis x line=none,ymode=log,xmin=-1,xmax=33,ylabel={\scriptsize eindeutig (log)},tick label style={font=\tiny},label style={font=\scriptsize},ylabel shift=-3pt]\addplot[mark=*,mark size=0.7pt,draw=black,thin] coordinates {(0,13541) (1,50218) (2,117542) (3,571700) (4,1255762) (5,2200446) (6,2861119) (7,1484906) (8,1522924) (9,339215) (10,204000) (11,20136) (12,10874) (13,971) (14,50)};\end{axis}\end{tikzpicture}\end{minipage}
\vspace{3mm}

\noindent\begin{minipage}[t]{60mm}\vspace{0pt}\sbox\hmboard{\chessboard[setfen=7N/7k/8/8/2p5/8/8/K5B1 b - - 0 1]}\usebox\hmboard\\[1mm]
\makebox[\wd\hmboard]{\sffamily\hspace{3mm}\textbf{h\#7}\hfill(3+2)}\end{minipage}\hfill
\begin{minipage}[t]{52mm}\vspace{0pt}
\textcolor{ink}{\sffamily\textsc{Lösung}}\par\smallskip\small
\begin{tabular}{@{}r@{\enspace}l@{\enspace}l@{}}
\textbf{1.} & c3 & \symknight{}f7\\
\textbf{2.} & \symking{}g6 & \symking{}a2\\
\textbf{3.} & \symking{}f5 & \symking{}b3\\
\textbf{4.} & \symking{}e4 & \symknight{}d6+\\
\textbf{5.} & \symking{}d3 & \symknight{}b5\\
\textbf{6.} & \symking{}d2 & \symknight{}xc3\\
\textbf{7.} & \symking{}c1 & \symbishop{}e3\#\\
\end{tabular}\par\vspace{3mm}
\qrcode[height=12mm]{https://helpman.komtera.lt/?fen=7N/7k/8/8/2p5/8/8/K5B1&moves=7}\\{\sffamily\tiny\color{ink}scannen zum Lösen}
\end{minipage}
\vspace{2mm}

{\raggedright\noindent\textcolor{ink}{\sffamily\textsc{Themen}}\enspace \themetag{reines Matt} \themetag{Mustermatt} \themetag{Idealmatt} \themetag{Spiegelmatt}\par}

\clearpage
\phantomsection\label{hm:KBNvkr}
\begin{tcolorbox}[enhanced,colback=ink,colframe=ink,arc=0pt,boxrule=0pt,left=4mm,right=4mm,top=2.5mm,bottom=2.5mm]
\color{white}\sffamily{\LARGE\bfseries KBNvkr}\hfill{\Large h\#7}\\[0.5mm]
{\small 5 Steine\quad Schwarz am Zug}\hfill{\small Nr.~110\qquad\href{https://helpman.komtera.lt/?fen=7k/8/5r2/8/N7/8/8/KB6&moves=7}{\color{white}\bfseries lösen\,\ensuremath{\nearrow}}}
\end{tcolorbox}
\vspace{2mm}
\noindent\begin{minipage}[t]{36mm}\vspace{0pt}\sffamily\small\raggedright
\hmfigure{längste eindeutige}{{\Large h\#7}}\hmfigure{eindeutig bei dieser Länge}{{\Large 2} Stellungen}\hmfigure{längstes Matt}{{\Large h\#9}}\hmfigure{in der Klasse}{242 Millionen Stellungen,\\ 163 Millionen mit Hilfsmatt}
\end{minipage}\hfill
\begin{minipage}[t]{80mm}\vspace{0pt}\begin{tikzpicture}\begin{axis}[width=72mm,height=42mm,ybar,bar width=1.6pt,xmin=-1,xmax=19,ymin=0,ymax=38.48,axis y line*=left,axis x line*=bottom,xlabel={\scriptsize Halbzüge bis Matt},ylabel={\scriptsize Stellungen (Mio.)},tick label style={font=\tiny},label style={font=\scriptsize},enlarge x limits=false,ylabel shift=-3pt,xlabel shift=-3pt]\addplot[fill=ink!55,draw=none] coordinates {(0,0.001839) (1,0.007766) (2,0.03058) (3,0.2361) (4,0.9377) (5,4.299) (6,10.12) (7,21.73) (8,31.64) (9,34.99) (10,29.86) (11,18.03) (12,8.359) (13,2.457) (14,0.5973) (15,0.06798) (16,0.002768) (17,3e-06) (18,7.7e-05)};\draw[densely dotted] (axis cs:14,0) -- (axis cs:14,38.48);\end{axis}\begin{axis}[width=72mm,height=42mm,axis y line*=right,axis x line=none,ymode=log,xmin=-1,xmax=19,ylabel={\scriptsize eindeutig (log)},tick label style={font=\tiny},label style={font=\scriptsize},ylabel shift=-3pt]\addplot[mark=*,mark size=0.7pt,draw=black,thin] coordinates {(0,1839) (1,7297) (2,19860) (3,104290) (4,320232) (5,620447) (6,672876) (7,489921) (8,307758) (9,124492) (10,43136) (11,8824) (12,1846) (13,70) (14,2)};\end{axis}\end{tikzpicture}\end{minipage}
\vspace{3mm}

\noindent\begin{minipage}[t]{60mm}\vspace{0pt}\sbox\hmboard{\chessboard[setfen=7k/8/5r2/8/N7/8/8/KB6 b - - 0 1]}\usebox\hmboard\\[1mm]
\makebox[\wd\hmboard]{\sffamily\hspace{3mm}\textbf{h\#7}\hfill(3+2)}\end{minipage}\hfill
\begin{minipage}[t]{52mm}\vspace{0pt}
\textcolor{ink}{\sffamily\textsc{Lösung}}\par\smallskip\small
\begin{tabular}{@{}r@{\enspace}l@{\enspace}l@{}}
\textbf{1.} & \symking{}g7 & \symbishop{}a2\\
\textbf{2.} & \symking{}g6 & \symking{}b2\\
\textbf{3.} & \symking{}f5 & \symking{}a3\\
\textbf{4.} & \symking{}e4 & \symbishop{}d5+\\
\textbf{5.} & \symking{}d3 & \symknight{}c5+\\
\textbf{6.} & \symking{}c2 & \symknight{}b3\\
\textbf{7.} & \symking{}b1 & \symbishop{}e4\#\\
\end{tabular}\par\vspace{3mm}
\qrcode[height=12mm]{https://helpman.komtera.lt/?fen=7k/8/5r2/8/N7/8/8/KB6&moves=7}\\{\sffamily\tiny\color{ink}scannen zum Lösen}
\end{minipage}
\vspace{2mm}

{\raggedright\noindent\textcolor{ink}{\sffamily\textsc{Themen}}\enspace \themetag{reines Matt} \themetag{Mustermatt} \themetag{Spiegelmatt} \themetag{ein Stein zieht} \themetag{ohne Schlag}\par}

\clearpage
\phantomsection\label{hm:KBPPvk}
\begin{tcolorbox}[enhanced,colback=ink,colframe=ink,arc=0pt,boxrule=0pt,left=4mm,right=4mm,top=2.5mm,bottom=2.5mm]
\color{white}\sffamily{\LARGE\bfseries KBPPvk}\hfill{\Large h\#7}\\[0.5mm]
{\small 5 Steine\quad Schwarz am Zug}\hfill{\small Nr.~111\qquad\href{https://helpman.komtera.lt/?fen=8/8/8/8/2B5/2PK4/2P4k/8&moves=7}{\color{white}\bfseries lösen\,\ensuremath{\nearrow}}}
\end{tcolorbox}
\vspace{2mm}
\noindent\begin{minipage}[t]{36mm}\vspace{0pt}\sffamily\small\raggedright
\hmfigure{längste eindeutige}{{\Large h\#7}}\hmfigure{eindeutig bei dieser Länge}{{\Large 12} Stellungen}\hmfigure{längstes Matt}{{\Large h\#8}}\hmfigure{in der Klasse}{533 Millionen Stellungen,\\ 209 Millionen mit Hilfsmatt}
\end{minipage}\hfill
\begin{minipage}[t]{80mm}\vspace{0pt}\begin{tikzpicture}\begin{axis}[width=72mm,height=42mm,ybar,bar width=1.6pt,xmin=-1,xmax=17,ymin=0,ymax=40.63,axis y line*=left,axis x line*=bottom,xlabel={\scriptsize Halbzüge bis Matt},ylabel={\scriptsize Stellungen (Mio.)},tick label style={font=\tiny},label style={font=\scriptsize},enlarge x limits=false,ylabel shift=-3pt,xlabel shift=-3pt]\addplot[fill=ink!55,draw=none] coordinates {(0,0.007065) (1,0.1595) (2,0.3634) (3,2.901) (4,8.707) (5,27.58) (6,36.93) (7,28.27) (8,30.8) (9,19.87) (10,22.01) (11,13.13) (12,14.06) (13,2.388) (14,1.754) (15,0.004802) (16,0.001843)};\draw[densely dotted] (axis cs:14,0) -- (axis cs:14,40.63);\end{axis}\begin{axis}[width=72mm,height=42mm,axis y line*=right,axis x line=none,ymode=log,xmin=-1,xmax=17,ylabel={\scriptsize eindeutig (log)},tick label style={font=\tiny},label style={font=\scriptsize},ylabel shift=-3pt]\addplot[mark=*,mark size=0.7pt,draw=black,thin] coordinates {(0,7065) (1,56336) (2,155953) (3,1508731) (4,2681544) (5,2632764) (6,1490591) (7,866098) (8,539228) (9,286687) (10,193341) (11,77232) (12,36380) (13,49) (14,12)};\end{axis}\end{tikzpicture}\end{minipage}
\vspace{3mm}

\noindent\begin{minipage}[t]{60mm}\vspace{0pt}\sbox\hmboard{\chessboard[setfen=8/8/8/8/2B5/2PK4/2P4k/8 b - - 0 1]}\usebox\hmboard\\[1mm]
\makebox[\wd\hmboard]{\sffamily\hspace{3mm}\textbf{h\#7}\hfill(4+1)}\end{minipage}\hfill
\begin{minipage}[t]{52mm}\vspace{0pt}
\textcolor{ink}{\sffamily\textsc{Lösung}}\par\smallskip\small
\begin{tabular}{@{}r@{\enspace}l@{\enspace}l@{}}
\textbf{1.} & \symking{}g3 & \symbishop{}a6\\
\textbf{2.} & \symking{}f4 & c4\\
\textbf{3.} & \symking{}e5 & c5\\
\textbf{4.} & \symking{}d5 & c6\\
\textbf{5.} & \symking{}c5 & c7\\
\textbf{6.} & \symking{}b6 & c8=\symqueen{}\\
\textbf{7.} & \symking{}a7 & \symqueen{}b7\#\\
\end{tabular}\par\vspace{3mm}
\qrcode[height=12mm]{https://helpman.komtera.lt/?fen=8/8/8/8/2B5/2PK4/2P4k/8&moves=7}\\{\sffamily\tiny\color{ink}scannen zum Lösen}
\end{minipage}
\vspace{2mm}

{\raggedright\noindent\textcolor{ink}{\sffamily\textsc{Themen}}\enspace \themetag{reines Matt} \themetag{Mustermatt} \themetag{Umwandlung} \themetag{ein Stein zieht} \themetag{ohne Schlag} \themetag{ohne Schach} \themetag{Umnov}\par}

\clearpage
\phantomsection\label{hm:KBPvkb}
\begin{tcolorbox}[enhanced,colback=ink,colframe=ink,arc=0pt,boxrule=0pt,left=4mm,right=4mm,top=2.5mm,bottom=2.5mm]
\color{white}\sffamily{\LARGE\bfseries KBPvkb}\hfill{\Large h\#7}\\[0.5mm]
{\small 5 Steine\quad Schwarz am Zug}\hfill{\small Nr.~112\qquad\href{https://helpman.komtera.lt/?fen=7k/6b1/8/8/3B4/8/3P4/K7&moves=7}{\color{white}\bfseries lösen\,\ensuremath{\nearrow}}}
\end{tcolorbox}
\vspace{2mm}
\noindent\begin{minipage}[t]{36mm}\vspace{0pt}\sffamily\small\raggedright
\hmfigure{längste eindeutige}{{\Large h\#7}}\hmfigure{eindeutig bei dieser Länge}{{\Large 65} Stellungen}\hmfigure{längstes Matt}{{\Large h\#10}}\hmfigure{in der Klasse}{710 Millionen Stellungen,\\ 528 Millionen mit Hilfsmatt}
\end{minipage}\hfill
\begin{minipage}[t]{80mm}\vspace{0pt}\begin{tikzpicture}\begin{axis}[width=72mm,height=42mm,ybar,bar width=1.6pt,xmin=-1,xmax=21,ymin=0,ymax=73.56,axis y line*=left,axis x line*=bottom,xlabel={\scriptsize Halbzüge bis Matt},ylabel={\scriptsize Stellungen (Mio.)},tick label style={font=\tiny},label style={font=\scriptsize},enlarge x limits=false,ylabel shift=-3pt,xlabel shift=-3pt]\addplot[fill=ink!55,draw=none] coordinates {(0,0.008238) (1,0.1235) (2,0.4495) (3,4.018) (4,12.72) (5,43.36) (6,59.32) (7,64.71) (8,66.87) (9,66.28) (10,66.01) (11,64.37) (12,55.76) (13,16.88) (14,6.499) (15,0.1506) (16,0.00431) (18,0.000398) (20,0.000266)};\draw[densely dotted] (axis cs:14,0) -- (axis cs:14,73.56);\end{axis}\begin{axis}[width=72mm,height=42mm,axis y line*=right,axis x line=none,ymode=log,xmin=-1,xmax=21,ylabel={\scriptsize eindeutig (log)},tick label style={font=\tiny},label style={font=\scriptsize},ylabel shift=-3pt]\addplot[mark=*,mark size=0.7pt,draw=black,thin] coordinates {(0,8238) (1,58746) (2,213182) (3,1578871) (4,2977515) (5,3511958) (6,2333861) (7,1727218) (8,1023980) (9,781779) (10,372487) (11,238151) (12,98281) (13,2737) (14,65)};\end{axis}\end{tikzpicture}\end{minipage}
\vspace{3mm}

\noindent\begin{minipage}[t]{60mm}\vspace{0pt}\sbox\hmboard{\chessboard[setfen=7k/6b1/8/8/3B4/8/3P4/K7 b - - 0 1]}\usebox\hmboard\\[1mm]
\makebox[\wd\hmboard]{\sffamily\hspace{3mm}\textbf{h\#7}\hfill(3+2)}\end{minipage}\hfill
\begin{minipage}[t]{52mm}\vspace{0pt}
\textcolor{ink}{\sffamily\textsc{Lösung}}\par\smallskip\small
\begin{tabular}{@{}r@{\enspace}l@{\enspace}l@{}}
\textbf{1.} & \symking{}h7 & \symbishop{}c3\\
\textbf{2.} & \symking{}g6 & d4\\
\textbf{3.} & \symking{}f5 & d5\\
\textbf{4.} & \symking{}e4 & d6\\
\textbf{5.} & \symking{}d3 & d7\\
\textbf{6.} & \symking{}c2 & d8=\symqueen{}\\
\textbf{7.} & \symking{}c1 & \symqueen{}d2\#\\
\end{tabular}\par\vspace{3mm}
\qrcode[height=12mm]{https://helpman.komtera.lt/?fen=7k/6b1/8/8/3B4/8/3P4/K7&moves=7}\\{\sffamily\tiny\color{ink}scannen zum Lösen}
\end{minipage}
\vspace{2mm}

{\raggedright\noindent\textcolor{ink}{\sffamily\textsc{Themen}}\enspace \themetag{Umwandlung} \themetag{Excelsior} \themetag{Rundlauf} \themetag{ein Stein zieht} \themetag{ohne Schlag} \themetag{ohne Schach}\par}

\clearpage
\phantomsection\label{hm:KBPvkn}
\begin{tcolorbox}[enhanced,colback=ink,colframe=ink,arc=0pt,boxrule=0pt,left=4mm,right=4mm,top=2.5mm,bottom=2.5mm]
\color{white}\sffamily{\LARGE\bfseries KBPvkn}\hfill{\Large h\#7}\\[0.5mm]
{\small 5 Steine\quad Schwarz am Zug}\hfill{\small Nr.~113\qquad\href{https://helpman.komtera.lt/?fen=8/8/8/8/8/K1n5/PB6/7k&moves=7}{\color{white}\bfseries lösen\,\ensuremath{\nearrow}}}
\end{tcolorbox}
\vspace{2mm}
\noindent\begin{minipage}[t]{36mm}\vspace{0pt}\sffamily\small\raggedright
\hmfigure{längste eindeutige}{{\Large h\#7}}\hmfigure{eindeutig bei dieser Länge}{{\Large 101} Stellungen}\hmfigure{längstes Matt}{{\Large h\#12}}\hmfigure{in der Klasse}{710 Millionen Stellungen,\\ 542 Millionen mit Hilfsmatt}
\end{minipage}\hfill
\begin{minipage}[t]{80mm}\vspace{0pt}\begin{tikzpicture}\begin{axis}[width=72mm,height=42mm,ybar,bar width=1.6pt,xmin=-1,xmax=25,ymin=0,ymax=81.67,axis y line*=left,axis x line*=bottom,xlabel={\scriptsize Halbzüge bis Matt},ylabel={\scriptsize Stellungen (Mio.)},tick label style={font=\tiny},label style={font=\scriptsize},enlarge x limits=false,ylabel shift=-3pt,xlabel shift=-3pt]\addplot[fill=ink!55,draw=none] coordinates {(0,0.005421) (1,0.1384) (2,0.3912) (3,3.55) (4,9.617) (5,39.02) (6,53.36) (7,56.56) (8,62.03) (9,61.44) (10,70) (11,72.79) (12,74.25) (13,26.17) (14,12.59) (15,0.2252) (16,0.02047) (18,0.001696) (20,0.001729) (22,0.000188) (24,7.8e-05)};\draw[densely dotted] (axis cs:14,0) -- (axis cs:14,81.67);\end{axis}\begin{axis}[width=72mm,height=42mm,axis y line*=right,axis x line=none,ymode=log,xmin=-1,xmax=25,ylabel={\scriptsize eindeutig (log)},tick label style={font=\tiny},label style={font=\scriptsize},ylabel shift=-3pt]\addplot[mark=*,mark size=0.7pt,draw=black,thin] coordinates {(0,5421) (1,49983) (2,127879) (3,1300816) (4,2082636) (5,2841705) (6,1783996) (7,1428209) (8,1098570) (9,890803) (10,578768) (11,442537) (12,188235) (13,3996) (14,101)};\end{axis}\end{tikzpicture}\end{minipage}
\vspace{3mm}

\noindent\begin{minipage}[t]{60mm}\vspace{0pt}\sbox\hmboard{\chessboard[setfen=8/8/8/8/8/K1n5/PB6/7k b - - 0 1]}\usebox\hmboard\\[1mm]
\makebox[\wd\hmboard]{\sffamily\hspace{3mm}\textbf{h\#7}\hfill(3+2)}\end{minipage}\hfill
\begin{minipage}[t]{52mm}\vspace{0pt}
\textcolor{ink}{\sffamily\textsc{Lösung}}\par\smallskip\small
\begin{tabular}{@{}r@{\enspace}l@{\enspace}l@{}}
\textbf{1.} & \symking{}g2 & \symking{}b4\\
\textbf{2.} & \symking{}f3 & a4\\
\textbf{3.} & \symking{}e4 & a5\\
\textbf{4.} & \symking{}d5 & a6\\
\textbf{5.} & \symking{}c6 & a7\\
\textbf{6.} & \symking{}b7 & \symking{}c5\\
\textbf{7.} & \symking{}a6 & a8=\symqueen{}\#\\
\end{tabular}\par\vspace{3mm}
\qrcode[height=12mm]{https://helpman.komtera.lt/?fen=8/8/8/8/8/K1n5/PB6/7k&moves=7}\\{\sffamily\tiny\color{ink}scannen zum Lösen}
\end{minipage}
\vspace{2mm}

{\raggedright\noindent\textcolor{ink}{\sffamily\textsc{Themen}}\enspace \themetag{reines Matt} \themetag{Spiegelmatt} \themetag{Umwandlung} \themetag{Excelsior} \themetag{ein Stein zieht} \themetag{ohne Schlag} \themetag{ohne Schach}\par}

\clearpage
\phantomsection\label{hm:KBPvkq}
\begin{tcolorbox}[enhanced,colback=ink,colframe=ink,arc=0pt,boxrule=0pt,left=4mm,right=4mm,top=2.5mm,bottom=2.5mm]
\color{white}\sffamily{\LARGE\bfseries KBPvkq}\hfill{\Large h\#7}\\[0.5mm]
{\small 5 Steine\quad Schwarz am Zug}\hfill{\small Nr.~114\qquad\href{https://helpman.komtera.lt/?fen=8/8/B7/8/8/K7/P4q2/6k1&moves=7}{\color{white}\bfseries lösen\,\ensuremath{\nearrow}}}
\end{tcolorbox}
\vspace{2mm}
\noindent\begin{minipage}[t]{36mm}\vspace{0pt}\sffamily\small\raggedright
\hmfigure{längste eindeutige}{{\Large h\#7}}\hmfigure{eindeutig bei dieser Länge}{{\Large 75} Stellungen}\hmfigure{längstes Matt}{{\Large h\#8.5}}\hmfigure{in der Klasse}{710 Millionen Stellungen,\\ 465 Millionen mit Hilfsmatt}
\end{minipage}\hfill
\begin{minipage}[t]{80mm}\vspace{0pt}\begin{tikzpicture}\begin{axis}[width=72mm,height=42mm,ybar,bar width=1.6pt,xmin=-1,xmax=18,ymin=0,ymax=75.74,axis y line*=left,axis x line*=bottom,xlabel={\scriptsize Halbzüge bis Matt},ylabel={\scriptsize Stellungen (Mio.)},tick label style={font=\tiny},label style={font=\scriptsize},enlarge x limits=false,ylabel shift=-3pt,xlabel shift=-3pt]\addplot[fill=ink!55,draw=none] coordinates {(0,0.002251) (1,0.0544) (2,0.2734) (3,3.018) (4,9.834) (5,37.06) (6,44.04) (7,60.88) (8,46.36) (9,61.12) (10,48.7) (11,68.85) (12,50.94) (13,27.76) (14,5.041) (15,0.8364) (16,0.001846) (17,0.003074)};\draw[densely dotted] (axis cs:14,0) -- (axis cs:14,75.74);\end{axis}\begin{axis}[width=72mm,height=42mm,axis y line*=right,axis x line=none,ymode=log,xmin=-1,xmax=18,ylabel={\scriptsize eindeutig (log)},tick label style={font=\tiny},label style={font=\scriptsize},ylabel shift=-3pt]\addplot[mark=*,mark size=0.7pt,draw=black,thin] coordinates {(0,2251) (1,25908) (2,110489) (3,1206832) (4,2182037) (5,3037764) (6,1287125) (7,1343414) (8,543294) (9,609918) (10,299492) (11,244056) (12,68061) (13,11263) (14,75)};\end{axis}\end{tikzpicture}\end{minipage}
\vspace{3mm}

\noindent\begin{minipage}[t]{60mm}\vspace{0pt}\sbox\hmboard{\chessboard[setfen=8/8/B7/8/8/K7/P4q2/6k1 b - - 0 1]}\usebox\hmboard\\[1mm]
\makebox[\wd\hmboard]{\sffamily\hspace{3mm}\textbf{h\#7}\hfill(3+2)}\end{minipage}\hfill
\begin{minipage}[t]{52mm}\vspace{0pt}
\textcolor{ink}{\sffamily\textsc{Lösung}}\par\smallskip\small
\begin{tabular}{@{}r@{\enspace}l@{\enspace}l@{}}
\textbf{1.} & \symking{}g2 & \symking{}b4\\
\textbf{2.} & \symking{}f3 & a4\\
\textbf{3.} & \symking{}e4 & a5\\
\textbf{4.} & \symqueen{}b6+ & axb6\\
\textbf{5.} & \symking{}d5 & b7\\
\textbf{6.} & \symking{}c6 & \symbishop{}b5+\\
\textbf{7.} & \symking{}b6 & b8=\symqueen{}\#\\
\end{tabular}\par\vspace{3mm}
\qrcode[height=12mm]{https://helpman.komtera.lt/?fen=8/8/B7/8/8/K7/P4q2/6k1&moves=7}\\{\sffamily\tiny\color{ink}scannen zum Lösen}
\end{minipage}
\vspace{2mm}

{\raggedright\noindent\textcolor{ink}{\sffamily\textsc{Themen}}\enspace \themetag{reines Matt} \themetag{Mustermatt} \themetag{Idealmatt} \themetag{Umwandlung} \themetag{Excelsior} \themetag{Kniest}\par}

\clearpage
\phantomsection\label{hm:KBPvkr}
\begin{tcolorbox}[enhanced,colback=ink,colframe=ink,arc=0pt,boxrule=0pt,left=4mm,right=4mm,top=2.5mm,bottom=2.5mm]
\color{white}\sffamily{\LARGE\bfseries KBPvkr}\hfill{\Large h\#7}\\[0.5mm]
{\small 5 Steine\quad Schwarz am Zug}\hfill{\small Nr.~115\qquad\href{https://helpman.komtera.lt/?fen=8/8/B7/2r5/8/K7/P7/7k&moves=7}{\color{white}\bfseries lösen\,\ensuremath{\nearrow}}}
\end{tcolorbox}
\vspace{2mm}
\noindent\begin{minipage}[t]{36mm}\vspace{0pt}\sffamily\small\raggedright
\hmfigure{längste eindeutige}{{\Large h\#7}}\hmfigure{eindeutig bei dieser Länge}{{\Large 69} Stellungen}\hmfigure{längstes Matt}{{\Large h\#8.5}}\hmfigure{in der Klasse}{710 Millionen Stellungen,\\ 505 Millionen mit Hilfsmatt}
\end{minipage}\hfill
\begin{minipage}[t]{80mm}\vspace{0pt}\begin{tikzpicture}\begin{axis}[width=72mm,height=42mm,ybar,bar width=1.6pt,xmin=-1,xmax=18,ymin=0,ymax=75.78,axis y line*=left,axis x line*=bottom,xlabel={\scriptsize Halbzüge bis Matt},ylabel={\scriptsize Stellungen (Mio.)},tick label style={font=\tiny},label style={font=\scriptsize},enlarge x limits=false,ylabel shift=-3pt,xlabel shift=-3pt]\addplot[fill=ink!55,draw=none] coordinates {(0,0.0033) (1,0.09315) (2,0.4002) (3,3.897) (4,13.8) (5,43.66) (6,55.51) (7,62.67) (8,55.33) (9,62.48) (10,59.47) (11,68.89) (12,56.76) (13,17.87) (14,3.755) (15,0.1666) (16,0.002493) (17,0.00029)};\draw[densely dotted] (axis cs:14,0) -- (axis cs:14,75.78);\end{axis}\begin{axis}[width=72mm,height=42mm,axis y line*=right,axis x line=none,ymode=log,xmin=-1,xmax=18,ylabel={\scriptsize eindeutig (log)},tick label style={font=\tiny},label style={font=\scriptsize},ylabel shift=-3pt]\addplot[mark=*,mark size=0.7pt,draw=black,thin] coordinates {(0,3300) (1,41748) (2,170145) (3,1579600) (4,3352319) (5,3691295) (6,1760510) (7,1387451) (8,653699) (9,627581) (10,376613) (11,181281) (12,44198) (13,2436) (14,69)};\end{axis}\end{tikzpicture}\end{minipage}
\vspace{3mm}

\noindent\begin{minipage}[t]{60mm}\vspace{0pt}\sbox\hmboard{\chessboard[setfen=8/8/B7/2r5/8/K7/P7/7k b - - 0 1]}\usebox\hmboard\\[1mm]
\makebox[\wd\hmboard]{\sffamily\hspace{3mm}\textbf{h\#7}\hfill(3+2)}\end{minipage}\hfill
\begin{minipage}[t]{52mm}\vspace{0pt}
\textcolor{ink}{\sffamily\textsc{Lösung}}\par\smallskip\small
\begin{tabular}{@{}r@{\enspace}l@{\enspace}l@{}}
\textbf{1.} & \symking{}g2 & \symking{}b4\\
\textbf{2.} & \symking{}f3 & a4\\
\textbf{3.} & \symrook{}b5+ & axb5\\
\textbf{4.} & \symking{}e4 & b6\\
\textbf{5.} & \symking{}d5 & b7\\
\textbf{6.} & \symking{}c6 & \symbishop{}b5+\\
\textbf{7.} & \symking{}b6 & b8=\symqueen{}\#\\
\end{tabular}\par\vspace{3mm}
\qrcode[height=12mm]{https://helpman.komtera.lt/?fen=8/8/B7/2r5/8/K7/P7/7k&moves=7}\\{\sffamily\tiny\color{ink}scannen zum Lösen}
\end{minipage}
\vspace{2mm}

{\raggedright\noindent\textcolor{ink}{\sffamily\textsc{Themen}}\enspace \themetag{reines Matt} \themetag{Mustermatt} \themetag{Idealmatt} \themetag{Umwandlung} \themetag{Excelsior}\par}

\clearpage
\phantomsection\label{hm:KNNPvk}
\begin{tcolorbox}[enhanced,colback=ink,colframe=ink,arc=0pt,boxrule=0pt,left=4mm,right=4mm,top=2.5mm,bottom=2.5mm]
\color{white}\sffamily{\LARGE\bfseries KNNPvk}\hfill{\Large h\#7}\\[0.5mm]
{\small 5 Steine\quad Schwarz am Zug}\hfill{\small Nr.~116\qquad\href{https://helpman.komtera.lt/?fen=8/K7/8/8/6N1/4N3/4P3/6k1&moves=7}{\color{white}\bfseries lösen\,\ensuremath{\nearrow}}}
\end{tcolorbox}
\vspace{2mm}
\noindent\begin{minipage}[t]{36mm}\vspace{0pt}\sffamily\small\raggedright
\hmfigure{längste eindeutige}{{\Large h\#7}}\hmfigure{eindeutig bei dieser Länge}{{\Large 5} Stellungen}\hmfigure{längstes Matt}{{\Large h\#8}}\hmfigure{in der Klasse}{710 Millionen Stellungen,\\ 280 Millionen mit Hilfsmatt}
\end{minipage}\hfill
\begin{minipage}[t]{80mm}\vspace{0pt}\begin{tikzpicture}\begin{axis}[width=72mm,height=42mm,ybar,bar width=1.6pt,xmin=-1,xmax=17,ymin=0,ymax=58.87,axis y line*=left,axis x line*=bottom,xlabel={\scriptsize Halbzüge bis Matt},ylabel={\scriptsize Stellungen (Mio.)},tick label style={font=\tiny},label style={font=\scriptsize},enlarge x limits=false,ylabel shift=-3pt,xlabel shift=-3pt]\addplot[fill=ink!55,draw=none] coordinates {(0,0.005409) (1,0.1095) (2,0.2431) (3,2.258) (4,7.078) (5,22.82) (6,35.13) (7,39.55) (8,53.52) (9,40.71) (10,43.83) (11,18.11) (12,13.06) (13,2.403) (14,0.7517) (15,0.03541) (16,0.000105)};\draw[densely dotted] (axis cs:14,0) -- (axis cs:14,58.87);\end{axis}\begin{axis}[width=72mm,height=42mm,axis y line*=right,axis x line=none,ymode=log,xmin=-1,xmax=17,ylabel={\scriptsize eindeutig (log)},tick label style={font=\tiny},label style={font=\scriptsize},ylabel shift=-3pt]\addplot[mark=*,mark size=0.7pt,draw=black,thin] coordinates {(0,5409) (1,41156) (2,110803) (3,1240271) (4,2378089) (5,2273371) (6,1706365) (7,947905) (8,717132) (9,197079) (10,111707) (11,32299) (12,5837) (13,208) (14,5)};\end{axis}\end{tikzpicture}\end{minipage}
\vspace{3mm}

\noindent\begin{minipage}[t]{60mm}\vspace{0pt}\sbox\hmboard{\chessboard[setfen=8/K7/8/8/6N1/4N3/4P3/6k1 b - - 0 1]}\usebox\hmboard\\[1mm]
\makebox[\wd\hmboard]{\sffamily\hspace{3mm}\textbf{h\#7}\hfill(4+1)}\end{minipage}\hfill
\begin{minipage}[t]{52mm}\vspace{0pt}
\textcolor{ink}{\sffamily\textsc{Lösung}}\par\smallskip\small
\begin{tabular}{@{}r@{\enspace}l@{\enspace}l@{}}
\textbf{1.} & \symking{}h1 & \symknight{}f5\\
\textbf{2.} & \symking{}g2 & e4\\
\textbf{3.} & \symking{}f3 & e5\\
\textbf{4.} & \symking{}f4 & e6\\
\textbf{5.} & \symking{}g5 & e7\\
\textbf{6.} & \symking{}g6 & e8=\symqueen{}+\\
\textbf{7.} & \symking{}h7 & \symknight{}f6\#\\
\end{tabular}\par\vspace{3mm}
\qrcode[height=12mm]{https://helpman.komtera.lt/?fen=8/K7/8/8/6N1/4N3/4P3/6k1&moves=7}\\{\sffamily\tiny\color{ink}scannen zum Lösen}
\end{minipage}
\vspace{2mm}

{\raggedright\noindent\textcolor{ink}{\sffamily\textsc{Themen}}\enspace \themetag{Satzspiel} \themetag{Spiegelmatt} \themetag{Umwandlung} \themetag{Excelsior} \themetag{ein Stein zieht} \themetag{ohne Schlag}\par}

\clearpage
\phantomsection\label{hm:KNNvkp}
\begin{tcolorbox}[enhanced,colback=ink,colframe=ink,arc=0pt,boxrule=0pt,left=4mm,right=4mm,top=2.5mm,bottom=2.5mm]
\color{white}\sffamily{\LARGE\bfseries KNNvkp}\hfill{\Large h\#7}\\[0.5mm]
{\small 5 Steine\quad Schwarz am Zug}\hfill{\small Nr.~117\qquad\href{https://helpman.komtera.lt/?fen=8/4N2k/7p/7N/8/8/8/K7&moves=7}{\color{white}\bfseries lösen\,\ensuremath{\nearrow}}}
\end{tcolorbox}
\vspace{2mm}
\noindent\begin{minipage}[t]{36mm}\vspace{0pt}\sffamily\small\raggedright
\hmfigure{längste eindeutige}{{\Large h\#7}}\hmfigure{eindeutig bei dieser Länge}{{\Large 3} Stellungen}\hmfigure{längstes Matt}{{\Large h\#11}}\hmfigure{in der Klasse}{710 Millionen Stellungen,\\ 279 Millionen mit Hilfsmatt}
\end{minipage}\hfill
\begin{minipage}[t]{80mm}\vspace{0pt}\begin{tikzpicture}\begin{axis}[width=72mm,height=42mm,ybar,bar width=1.6pt,xmin=-1,xmax=23,ymin=0,ymax=67.82,axis y line*=left,axis x line*=bottom,xlabel={\scriptsize Halbzüge bis Matt},ylabel={\scriptsize Stellungen (Mio.)},tick label style={font=\tiny},label style={font=\scriptsize},enlarge x limits=false,ylabel shift=-3pt,xlabel shift=-3pt]\addplot[fill=ink!55,draw=none] coordinates {(0,0.003389) (1,0.01679) (2,0.04799) (3,0.3619) (4,1.179) (5,5.505) (6,13.02) (7,29.12) (8,51.52) (9,55.01) (10,61.65) (11,33.23) (12,20.57) (13,5.95) (14,1.442) (15,0.3191) (16,0.001746) (17,8e-06) (18,0.000708) (19,2e-06) (20,0.000225) (22,0.000119)};\draw[densely dotted] (axis cs:14,0) -- (axis cs:14,67.82);\end{axis}\begin{axis}[width=72mm,height=42mm,axis y line*=right,axis x line=none,ymode=log,xmin=-1,xmax=23,ylabel={\scriptsize eindeutig (log)},tick label style={font=\tiny},label style={font=\scriptsize},ylabel shift=-3pt]\addplot[mark=*,mark size=0.7pt,draw=black,thin] coordinates {(0,3389) (1,15989) (2,42213) (3,248425) (4,483764) (5,669624) (6,849781) (7,574553) (8,640208) (9,152498) (10,91068) (11,13218) (12,5956) (13,1749) (14,3)};\end{axis}\end{tikzpicture}\end{minipage}
\vspace{3mm}

\noindent\begin{minipage}[t]{60mm}\vspace{0pt}\sbox\hmboard{\chessboard[setfen=8/4N2k/7p/7N/8/8/8/K7 b - - 0 1]}\usebox\hmboard\\[1mm]
\makebox[\wd\hmboard]{\sffamily\hspace{3mm}\textbf{h\#7}\hfill(3+2)}\end{minipage}\hfill
\begin{minipage}[t]{52mm}\vspace{0pt}
\textcolor{ink}{\sffamily\textsc{Lösung}}\par\smallskip\small
\begin{tabular}{@{}r@{\enspace}l@{\enspace}l@{}}
\textbf{1.} & \symking{}h8 & \symknight{}f4\\
\textbf{2.} & h5 & \symking{}b2\\
\textbf{3.} & h4 & \symking{}c3\\
\textbf{4.} & h3 & \symking{}d4\\
\textbf{5.} & h2 & \symking{}e5\\
\textbf{6.} & h1=\symrook{} & \symking{}f6\\
\textbf{7.} & \symrook{}h7 & \symknight{}fg6\#\\
\end{tabular}\par\vspace{3mm}
\qrcode[height=12mm]{https://helpman.komtera.lt/?fen=8/4N2k/7p/7N/8/8/8/K7&moves=7}\\{\sffamily\tiny\color{ink}scannen zum Lösen}
\end{minipage}
\vspace{2mm}

{\raggedright\noindent\textcolor{ink}{\sffamily\textsc{Themen}}\enspace \themetag{Satzspiel} \themetag{reines Matt} \themetag{Mustermatt} \themetag{Idealmatt} \themetag{Umwandlung} \themetag{Unterverwandlung} \themetag{Selbstblock} \themetag{ohne Schlag} \themetag{ohne Schach} \themetag{Umnov}\par}

\clearpage
\phantomsection\label{hm:KNNvkq}
\begin{tcolorbox}[enhanced,colback=ink,colframe=ink,arc=0pt,boxrule=0pt,left=4mm,right=4mm,top=2.5mm,bottom=2.5mm]
\color{white}\sffamily{\LARGE\bfseries KNNvkq}\hfill{\Large h\#7}\\[0.5mm]
{\small 5 Steine\quad Schwarz am Zug}\hfill{\small Nr.~118\qquad\href{https://helpman.komtera.lt/?fen=k7/8/8/2q5/7N/3K4/7N/8&moves=7}{\color{white}\bfseries lösen\,\ensuremath{\nearrow}}}
\end{tcolorbox}
\vspace{2mm}
\noindent\begin{minipage}[t]{36mm}\vspace{0pt}\sffamily\small\raggedright
\hmfigure{längste eindeutige}{{\Large h\#7}}\hmfigure{eindeutig bei dieser Länge}{{\Large 1} Stellung}\hmfigure{längstes Matt}{{\Large h\#8.5}}\hmfigure{in der Klasse}{242 Millionen Stellungen,\\ 77,1 Millionen mit Hilfsmatt}
\end{minipage}\hfill
\begin{minipage}[t]{80mm}\vspace{0pt}\begin{tikzpicture}\begin{axis}[width=72mm,height=42mm,ybar,bar width=1.6pt,xmin=-1,xmax=18,ymin=0,ymax=19.64,axis y line*=left,axis x line*=bottom,xlabel={\scriptsize Halbzüge bis Matt},ylabel={\scriptsize Stellungen (Mio.)},tick label style={font=\tiny},label style={font=\scriptsize},enlarge x limits=false,ylabel shift=-3pt,xlabel shift=-3pt]\addplot[fill=ink!55,draw=none] coordinates {(0,0.000465) (1,0.002387) (2,0.008542) (3,0.07435) (4,0.2366) (5,1.307) (6,2.736) (7,8.133) (8,11.14) (9,17.86) (10,14.23) (11,12.86) (12,5.231) (13,2.659) (14,0.4921) (15,0.1789) (16,0.000352) (17,6.8e-05)};\draw[densely dotted] (axis cs:14,0) -- (axis cs:14,19.64);\end{axis}\begin{axis}[width=72mm,height=42mm,axis y line*=right,axis x line=none,ymode=log,xmin=-1,xmax=18,ylabel={\scriptsize eindeutig (log)},tick label style={font=\tiny},label style={font=\scriptsize},ylabel shift=-3pt]\addplot[mark=*,mark size=0.7pt,draw=black,thin] coordinates {(0,465) (1,2337) (2,5235) (3,32491) (4,64018) (5,150123) (6,137155) (7,191428) (8,131307) (9,85470) (10,22464) (11,6662) (12,647) (13,103) (14,1)};\end{axis}\end{tikzpicture}\end{minipage}
\vspace{3mm}

\noindent\begin{minipage}[t]{60mm}\vspace{0pt}\sbox\hmboard{\chessboard[setfen=k7/8/8/2q5/7N/3K4/7N/8 b - - 0 1]}\usebox\hmboard\\[1mm]
\makebox[\wd\hmboard]{\sffamily\hspace{3mm}\textbf{h\#7}\hfill(3+2)}\end{minipage}\hfill
\begin{minipage}[t]{52mm}\vspace{0pt}
\textcolor{ink}{\sffamily\textsc{Lösung}}\par\smallskip\small
\begin{tabular}{@{}r@{\enspace}l@{\enspace}l@{}}
\textbf{1.} & \symking{}b7 & \symking{}e4\\
\textbf{2.} & \symking{}c6 & \symking{}f4\\
\textbf{3.} & \symking{}d5 & \symking{}g5\\
\textbf{4.} & \symking{}e6+ & \symking{}h6\\
\textbf{5.} & \symking{}f7 & \symknight{}g4\\
\textbf{6.} & \symking{}g8 & \symknight{}f6+\\
\textbf{7.} & \symking{}h8 & \symknight{}g6\#\\
\end{tabular}\par\vspace{3mm}
\qrcode[height=12mm]{https://helpman.komtera.lt/?fen=k7/8/8/2q5/7N/3K4/7N/8&moves=7}\\{\sffamily\tiny\color{ink}scannen zum Lösen}
\end{minipage}
\vspace{2mm}

{\raggedright\noindent\textcolor{ink}{\sffamily\textsc{Themen}}\enspace \themetag{Spiegelmatt} \themetag{ein Stein zieht} \themetag{ohne Schlag}\par}

\begin{tcolorbox}[enhanced,colback=white,colframe=ink!35,arc=1.5pt,boxrule=0.4pt,left=3mm,right=3mm,top=1.5mm,bottom=1.5mm]
\textcolor{ink}{\sffamily\textsc{veröffentlicht}}\par\smallskip\small \textbf{Sheglow, Wiktor S.}, h\#7\quad{\color{ink}\scriptsize verwandte Stellung}\\{\scriptsize Suomen Tehtäväniekat, No. (269), 05-06/1998}
\end{tcolorbox}
\clearpage
\phantomsection\label{hm:KNNvkr}
\begin{tcolorbox}[enhanced,colback=ink,colframe=ink,arc=0pt,boxrule=0pt,left=4mm,right=4mm,top=2.5mm,bottom=2.5mm]
\color{white}\sffamily{\LARGE\bfseries KNNvkr}\hfill{\Large h\#7}\\[0.5mm]
{\small 5 Steine\quad Schwarz am Zug}\hfill{\small Nr.~119\qquad\href{https://helpman.komtera.lt/?fen=6Nk/7r/8/7N/8/8/8/K7&moves=7}{\color{white}\bfseries lösen\,\ensuremath{\nearrow}}}
\end{tcolorbox}
\vspace{2mm}
\noindent\begin{minipage}[t]{36mm}\vspace{0pt}\sffamily\small\raggedright
\hmfigure{längste eindeutige}{{\Large h\#7}}\hmfigure{eindeutig bei dieser Länge}{{\Large 1} Stellung}\hmfigure{längstes Matt}{{\Large h\#10}}\hmfigure{in der Klasse}{242 Millionen Stellungen,\\ 83,9 Millionen mit Hilfsmatt}
\end{minipage}\hfill
\begin{minipage}[t]{80mm}\vspace{0pt}\begin{tikzpicture}\begin{axis}[width=72mm,height=42mm,ybar,bar width=1.6pt,xmin=-1,xmax=21,ymin=0,ymax=19.99,axis y line*=left,axis x line*=bottom,xlabel={\scriptsize Halbzüge bis Matt},ylabel={\scriptsize Stellungen (Mio.)},tick label style={font=\tiny},label style={font=\scriptsize},enlarge x limits=false,ylabel shift=-3pt,xlabel shift=-3pt]\addplot[fill=ink!55,draw=none] coordinates {(0,0.000778) (1,0.003875) (2,0.01585) (3,0.1298) (4,0.582) (5,2.598) (6,6.074) (7,12.17) (8,16.29) (9,18.17) (10,14.13) (11,8.736) (12,3.443) (13,1.227) (14,0.2408) (15,0.06842) (16,7e-05) (18,3e-06) (20,1e-06)};\draw[densely dotted] (axis cs:14,0) -- (axis cs:14,19.99);\end{axis}\begin{axis}[width=72mm,height=42mm,axis y line*=right,axis x line=none,ymode=log,xmin=-1,xmax=21,ylabel={\scriptsize eindeutig (log)},tick label style={font=\tiny},label style={font=\scriptsize},ylabel shift=-3pt]\addplot[mark=*,mark size=0.7pt,draw=black,thin] coordinates {(0,778) (1,3746) (2,11788) (3,67927) (4,226800) (5,329711) (6,292417) (7,205879) (8,126510) (9,48144) (10,13193) (11,1920) (12,191) (13,12) (14,1)};\end{axis}\end{tikzpicture}\end{minipage}
\vspace{3mm}

\noindent\begin{minipage}[t]{60mm}\vspace{0pt}\sbox\hmboard{\chessboard[setfen=6Nk/7r/8/7N/8/8/8/K7 b - - 0 1]}\usebox\hmboard\\[1mm]
\makebox[\wd\hmboard]{\sffamily\hspace{3mm}\textbf{h\#7}\hfill(3+2)}\end{minipage}\hfill
\begin{minipage}[t]{52mm}\vspace{0pt}
\textcolor{ink}{\sffamily\textsc{Lösung}}\par\smallskip\small
\begin{tabular}{@{}r@{\enspace}l@{\enspace}l@{}}
\textbf{1.} & \symrook{}a7+ & \symking{}b2\\
\textbf{2.} & \symking{}h7 & \symking{}c3\\
\textbf{3.} & \symking{}g6 & \symking{}d4\\
\textbf{4.} & \symking{}f7 & \symking{}e5\\
\textbf{5.} & \symking{}e8 & \symking{}f6\\
\textbf{6.} & \symrook{}a8 & \symking{}g7\\
\textbf{7.} & \symrook{}d8 & \symknight{}hf6\#\\
\end{tabular}\par\vspace{3mm}
\qrcode[height=12mm]{https://helpman.komtera.lt/?fen=6Nk/7r/8/7N/8/8/8/K7&moves=7}\\{\sffamily\tiny\color{ink}scannen zum Lösen}
\end{minipage}
\vspace{2mm}

{\raggedright\noindent\textcolor{ink}{\sffamily\textsc{Themen}}\enspace \themetag{reines Matt} \themetag{Mustermatt} \themetag{Idealmatt} \themetag{Selbstblock} \themetag{ohne Schlag}\par}

\clearpage
\phantomsection\label{hm:KNPPvk}
\begin{tcolorbox}[enhanced,colback=ink,colframe=ink,arc=0pt,boxrule=0pt,left=4mm,right=4mm,top=2.5mm,bottom=2.5mm]
\color{white}\sffamily{\LARGE\bfseries KNPPvk}\hfill{\Large h\#7}\\[0.5mm]
{\small 5 Steine\quad Schwarz am Zug}\hfill{\small Nr.~120\qquad\href{https://helpman.komtera.lt/?fen=8/8/8/8/4N3/K3P3/4P3/k7&moves=7}{\color{white}\bfseries lösen\,\ensuremath{\nearrow}}}
\end{tcolorbox}
\vspace{2mm}
\noindent\begin{minipage}[t]{36mm}\vspace{0pt}\sffamily\small\raggedright
\hmfigure{längste eindeutige}{{\Large h\#7}}\hmfigure{eindeutig bei dieser Länge}{{\Large 18} Stellungen}\hmfigure{längstes Matt}{{\Large h\#8.5}}\hmfigure{in der Klasse}{533 Millionen Stellungen,\\ 214 Millionen mit Hilfsmatt}
\end{minipage}\hfill
\begin{minipage}[t]{80mm}\vspace{0pt}\begin{tikzpicture}\begin{axis}[width=72mm,height=42mm,ybar,bar width=1.6pt,xmin=-1,xmax=18,ymin=0,ymax=36.57,axis y line*=left,axis x line*=bottom,xlabel={\scriptsize Halbzüge bis Matt},ylabel={\scriptsize Stellungen (Mio.)},tick label style={font=\tiny},label style={font=\scriptsize},enlarge x limits=false,ylabel shift=-3pt,xlabel shift=-3pt]\addplot[fill=ink!55,draw=none] coordinates {(0,0.002139) (1,0.1285) (2,0.2671) (3,2.372) (4,6.768) (5,23.19) (6,33.25) (7,31.8) (8,32.68) (9,23.06) (10,23.18) (11,15.14) (12,15.52) (13,3.904) (14,3.02) (15,0.02637) (16,0.004325) (17,1e-06)};\draw[densely dotted] (axis cs:14,0) -- (axis cs:14,36.57);\end{axis}\begin{axis}[width=72mm,height=42mm,axis y line*=right,axis x line=none,ymode=log,xmin=-1,xmax=18,ylabel={\scriptsize eindeutig (log)},tick label style={font=\tiny},label style={font=\scriptsize},ylabel shift=-3pt]\addplot[mark=*,mark size=0.7pt,draw=black,thin] coordinates {(0,2139) (1,32271) (2,90475) (3,1179700) (4,2074753) (5,2134674) (6,1186019) (7,721704) (8,406971) (9,229297) (10,130527) (11,81955) (12,37535) (13,193) (14,18)};\end{axis}\end{tikzpicture}\end{minipage}
\vspace{3mm}

\noindent\begin{minipage}[t]{60mm}\vspace{0pt}\sbox\hmboard{\chessboard[setfen=8/8/8/8/4N3/K3P3/4P3/k7 b - - 0 1]}\usebox\hmboard\\[1mm]
\makebox[\wd\hmboard]{\sffamily\hspace{3mm}\textbf{h\#7}\hfill(4+1)}\end{minipage}\hfill
\begin{minipage}[t]{52mm}\vspace{0pt}
\textcolor{ink}{\sffamily\textsc{Lösung}}\par\smallskip\small
\begin{tabular}{@{}r@{\enspace}l@{\enspace}l@{}}
\textbf{1.} & \symking{}b1 & \symknight{}d6\\
\textbf{2.} & \symking{}c2 & e4\\
\textbf{3.} & \symking{}c3 & e5\\
\textbf{4.} & \symking{}d4 & e6\\
\textbf{5.} & \symking{}c5 & e7\\
\textbf{6.} & \symking{}b6 & e8=\symqueen{}\\
\textbf{7.} & \symking{}a5 & \symqueen{}b5\#\\
\end{tabular}\par\vspace{3mm}
\qrcode[height=12mm]{https://helpman.komtera.lt/?fen=8/8/8/8/4N3/K3P3/4P3/k7&moves=7}\\{\sffamily\tiny\color{ink}scannen zum Lösen}
\end{minipage}
\vspace{2mm}

{\raggedright\noindent\textcolor{ink}{\sffamily\textsc{Themen}}\enspace \themetag{Umwandlung} \themetag{ein Stein zieht} \themetag{ohne Schlag} \themetag{ohne Schach}\par}

\clearpage
\phantomsection\label{hm:KPPPvk}
\begin{tcolorbox}[enhanced,colback=ink,colframe=ink,arc=0pt,boxrule=0pt,left=4mm,right=4mm,top=2.5mm,bottom=2.5mm]
\color{white}\sffamily{\LARGE\bfseries KPPPvk}\hfill{\Large h\#7}\\[0.5mm]
{\small 5 Steine\quad Schwarz am Zug}\hfill{\small Nr.~121\qquad\href{https://helpman.komtera.lt/?fen=8/8/8/8/K7/PP4P1/8/7k&moves=7}{\color{white}\bfseries lösen\,\ensuremath{\nearrow}}}
\end{tcolorbox}
\vspace{2mm}
\noindent\begin{minipage}[t]{36mm}\vspace{0pt}\sffamily\small\raggedright
\hmfigure{längste eindeutige}{{\Large h\#7}}\hmfigure{eindeutig bei dieser Länge}{{\Large 2} Stellungen}\hmfigure{längstes Matt}{{\Large h\#8}}\hmfigure{in der Klasse}{399 Millionen Stellungen,\\ 54,4 Millionen mit Hilfsmatt}
\end{minipage}\hfill
\begin{minipage}[t]{80mm}\vspace{0pt}\begin{tikzpicture}\begin{axis}[width=72mm,height=42mm,ybar,bar width=1.6pt,xmin=-1,xmax=17,ymin=0,ymax=10.8,axis y line*=left,axis x line*=bottom,xlabel={\scriptsize Halbzüge bis Matt},ylabel={\scriptsize Stellungen (Mio.)},tick label style={font=\tiny},label style={font=\scriptsize},enlarge x limits=false,ylabel shift=-3pt,xlabel shift=-3pt]\addplot[fill=ink!55,draw=none] coordinates {(0,0.000568) (1,0.0453) (2,0.09002) (3,0.7027) (4,1.875) (5,6.289) (6,9.336) (7,9.82) (8,9.571) (9,6.323) (10,5.235) (11,2.515) (12,2.035) (13,0.362) (14,0.2361) (15,0.00371) (16,0.000231)};\draw[densely dotted] (axis cs:14,0) -- (axis cs:14,10.8);\end{axis}\begin{axis}[width=72mm,height=42mm,axis y line*=right,axis x line=none,ymode=log,xmin=-1,xmax=17,ylabel={\scriptsize eindeutig (log)},tick label style={font=\tiny},label style={font=\scriptsize},ylabel shift=-3pt]\addplot[mark=*,mark size=0.7pt,draw=black,thin] coordinates {(0,568) (1,10258) (2,28474) (3,308347) (4,545019) (5,549927) (6,318090) (7,161178) (8,74765) (9,34551) (10,17596) (11,7389) (12,2929) (13,47) (14,2)};\end{axis}\end{tikzpicture}\end{minipage}
\vspace{3mm}

\noindent\begin{minipage}[t]{60mm}\vspace{0pt}\sbox\hmboard{\chessboard[setfen=8/8/8/8/K7/PP4P1/8/7k b - - 0 1]}\usebox\hmboard\\[1mm]
\makebox[\wd\hmboard]{\sffamily\hspace{3mm}\textbf{h\#7}\hfill(4+1)}\end{minipage}\hfill
\begin{minipage}[t]{52mm}\vspace{0pt}
\textcolor{ink}{\sffamily\textsc{Lösung}}\par\smallskip\small
\begin{tabular}{@{}r@{\enspace}l@{\enspace}l@{}}
\textbf{1.} & \symking{}g2 & \symking{}b4\\
\textbf{2.} & \symking{}f3 & a4\\
\textbf{3.} & \symking{}e4 & a5\\
\textbf{4.} & \symking{}d5 & a6\\
\textbf{5.} & \symking{}c6 & a7\\
\textbf{6.} & \symking{}b7 & \symking{}c5\\
\textbf{7.} & \symking{}a6 & a8=\symqueen{}\#\\
\end{tabular}\par\vspace{3mm}
\qrcode[height=12mm]{https://helpman.komtera.lt/?fen=8/8/8/8/K7/PP4P1/8/7k&moves=7}\\{\sffamily\tiny\color{ink}scannen zum Lösen}
\end{minipage}
\vspace{2mm}

{\raggedright\noindent\textcolor{ink}{\sffamily\textsc{Themen}}\enspace \themetag{reines Matt} \themetag{Mustermatt} \themetag{Spiegelmatt} \themetag{Umwandlung} \themetag{ein Stein zieht} \themetag{ohne Schlag} \themetag{ohne Schach}\par}

\clearpage
\phantomsection\label{hm:KPvkrr}
\begin{tcolorbox}[enhanced,colback=ink,colframe=ink,arc=0pt,boxrule=0pt,left=4mm,right=4mm,top=2.5mm,bottom=2.5mm]
\color{white}\sffamily{\LARGE\bfseries KPvkrr}\hfill{\Large h\#7}\\[0.5mm]
{\small 5 Steine\quad Schwarz am Zug}\hfill{\small Nr.~122\qquad\href{https://helpman.komtera.lt/?fen=8/7k/8/8/1r6/K7/Pr6/8&moves=7}{\color{white}\bfseries lösen\,\ensuremath{\nearrow}}}
\end{tcolorbox}
\vspace{2mm}
\noindent\begin{minipage}[t]{36mm}\vspace{0pt}\sffamily\small\raggedright
\hmfigure{längste eindeutige}{{\Large h\#7}}\hmfigure{eindeutig bei dieser Länge}{{\Large 5} Stellungen}\hmfigure{längstes Matt}{{\Large h\#8.5}}\hmfigure{in der Klasse}{710 Millionen Stellungen,\\ 244 Millionen mit Hilfsmatt}
\end{minipage}\hfill
\begin{minipage}[t]{80mm}\vspace{0pt}\begin{tikzpicture}\begin{axis}[width=72mm,height=42mm,ybar,bar width=1.6pt,xmin=-1,xmax=18,ymin=0,ymax=47.05,axis y line*=left,axis x line*=bottom,xlabel={\scriptsize Halbzüge bis Matt},ylabel={\scriptsize Stellungen (Mio.)},tick label style={font=\tiny},label style={font=\scriptsize},enlarge x limits=false,ylabel shift=-3pt,xlabel shift=-3pt]\addplot[fill=ink!55,draw=none] coordinates {(1,0.01567) (2,0.07819) (3,0.6794) (4,2.368) (5,10.26) (6,17.69) (7,36.3) (8,25.94) (9,42.77) (10,25.44) (11,42.37) (12,23.27) (13,15.28) (14,0.966) (15,0.2415) (16,0.000389) (17,0.000691)};\draw[densely dotted] (axis cs:14,0) -- (axis cs:14,47.05);\end{axis}\begin{axis}[width=72mm,height=42mm,axis y line*=right,axis x line=none,ymode=log,xmin=-1,xmax=18,ylabel={\scriptsize eindeutig (log)},tick label style={font=\tiny},label style={font=\scriptsize},ylabel shift=-3pt]\addplot[mark=*,mark size=0.7pt,draw=black,thin] coordinates {(1,5415) (2,27912) (3,234133) (4,568973) (5,1052820) (6,537974) (7,667271) (8,178049) (9,266747) (10,106238) (11,72067) (12,6312) (13,1499) (14,5)};\end{axis}\end{tikzpicture}\end{minipage}
\vspace{3mm}

\noindent\begin{minipage}[t]{60mm}\vspace{0pt}\sbox\hmboard{\chessboard[setfen=8/7k/8/8/1r6/K7/Pr6/8 b - - 0 1]}\usebox\hmboard\\[1mm]
\makebox[\wd\hmboard]{\sffamily\hspace{3mm}\textbf{h\#7}\hfill(2+3)}\end{minipage}\hfill
\begin{minipage}[t]{52mm}\vspace{0pt}
\textcolor{ink}{\sffamily\textsc{Lösung}}\par\smallskip\small
\begin{tabular}{@{}r@{\enspace}l@{\enspace}l@{}}
\textbf{1.} & \symrook{}4b3+ & axb3\\
\textbf{2.} & \symking{}g6 & b4\\
\textbf{3.} & \symking{}f5 & b5\\
\textbf{4.} & \symking{}e4 & b6\\
\textbf{5.} & \symking{}d3 & b7\\
\textbf{6.} & \symking{}c2 & b8=\symqueen{}\\
\textbf{7.} & \symking{}b1 & \symqueen{}xb2\#\\
\end{tabular}\par\vspace{3mm}
\qrcode[height=12mm]{https://helpman.komtera.lt/?fen=8/7k/8/8/1r6/K7/Pr6/8&moves=7}\\{\sffamily\tiny\color{ink}scannen zum Lösen}
\end{minipage}
\vspace{2mm}

{\raggedright\noindent\textcolor{ink}{\sffamily\textsc{Themen}}\enspace \themetag{Umwandlung} \themetag{Excelsior} \themetag{ein Stein zieht}\par}

\clearpage
\phantomsection\label{hm:KQBBvk}
\begin{tcolorbox}[enhanced,colback=ink,colframe=ink,arc=0pt,boxrule=0pt,left=4mm,right=4mm,top=2.5mm,bottom=2.5mm]
\color{white}\sffamily{\LARGE\bfseries KQBBvk}\hfill{\Large h\#7}\\[0.5mm]
{\small 5 Steine\quad Schwarz am Zug}\hfill{\small Nr.~123\qquad\href{https://helpman.komtera.lt/?fen=7Q/7k/7B/8/8/8/2B5/K7&moves=7}{\color{white}\bfseries lösen\,\ensuremath{\nearrow}}}
\end{tcolorbox}
\vspace{2mm}
\noindent\begin{minipage}[t]{36mm}\vspace{0pt}\sffamily\small\raggedright
\hmfigure{längste eindeutige}{{\Large h\#7}}\hmfigure{eindeutig bei dieser Länge}{{\Large 1} Stellung}\hmfigure{längstes Matt}{{\Large h\#8}}\hmfigure{in der Klasse}{242 Millionen Stellungen,\\ 75,2 Millionen mit Hilfsmatt}
\end{minipage}\hfill
\begin{minipage}[t]{80mm}\vspace{0pt}\begin{tikzpicture}\begin{axis}[width=72mm,height=42mm,ybar,bar width=1.6pt,xmin=-1,xmax=17,ymin=0,ymax=40.96,axis y line*=left,axis x line*=bottom,xlabel={\scriptsize Halbzüge bis Matt},ylabel={\scriptsize Stellungen (Mio.)},tick label style={font=\tiny},label style={font=\scriptsize},enlarge x limits=false,ylabel shift=-3pt,xlabel shift=-3pt]\addplot[fill=ink!55,draw=none] coordinates {(0,0.6041) (1,2.722) (2,9.644) (3,16.9) (4,37.24) (5,5.547) (6,2.28) (7,0.000914) (8,0.1596) (10,0.08311) (12,0.0213) (14,0.003679) (16,6e-06)};\draw[densely dotted] (axis cs:14,0) -- (axis cs:14,40.96);\end{axis}\begin{axis}[width=72mm,height=42mm,axis y line*=right,axis x line=none,ymode=log,xmin=-1,xmax=17,ylabel={\scriptsize eindeutig (log)},tick label style={font=\tiny},label style={font=\scriptsize},ylabel shift=-3pt]\addplot[mark=*,mark size=0.7pt,draw=black,thin] coordinates {(0,604142) (1,2025115) (2,6336032) (3,2457397) (4,1098453) (5,330) (6,13534) (8,1241) (10,328) (12,23) (14,1)};\end{axis}\end{tikzpicture}\end{minipage}
\vspace{3mm}

\noindent\begin{minipage}[t]{60mm}\vspace{0pt}\sbox\hmboard{\chessboard[setfen=7Q/7k/7B/8/8/8/2B5/K7 b - - 0 1]}\usebox\hmboard\\[1mm]
\makebox[\wd\hmboard]{\sffamily\hspace{3mm}\textbf{h\#7}\hfill(4+1)}\end{minipage}\hfill
\begin{minipage}[t]{52mm}\vspace{0pt}
\textcolor{ink}{\sffamily\textsc{Lösung}}\par\smallskip\small
\begin{tabular}{@{}r@{\enspace}l@{\enspace}l@{}}
\textbf{1.} & \symking{}xh8 & \symking{}b2\\
\textbf{2.} & \symking{}g8 & \symking{}c3\\
\textbf{3.} & \symking{}f7 & \symking{}d3\\
\textbf{4.} & \symking{}g6 & \symking{}e4\\
\textbf{5.} & \symking{}h5 & \symking{}f5\\
\textbf{6.} & \symking{}h4 & \symbishop{}g5+\\
\textbf{7.} & \symking{}h5 & \symbishop{}d1\#\\
\end{tabular}\par\vspace{3mm}
\qrcode[height=12mm]{https://helpman.komtera.lt/?fen=7Q/7k/7B/8/8/8/2B5/K7&moves=7}\\{\sffamily\tiny\color{ink}scannen zum Lösen}
\end{minipage}
\vspace{2mm}

{\raggedright\noindent\textcolor{ink}{\sffamily\textsc{Themen}}\enspace \themetag{Switchback} \themetag{ein Stein zieht}\par\smallskip{\footnotesize\textit{Hinweis:} das Diagramm hat keinen legalen letzten Zug, kann also in keiner Partie entstehen.}\par}

\clearpage
\phantomsection\label{hm:KQBPvk}
\begin{tcolorbox}[enhanced,colback=ink,colframe=ink,arc=0pt,boxrule=0pt,left=4mm,right=4mm,top=2.5mm,bottom=2.5mm]
\color{white}\sffamily{\LARGE\bfseries KQBPvk}\hfill{\Large h\#7}\\[0.5mm]
{\small 5 Steine\quad Schwarz am Zug}\hfill{\small Nr.~124\qquad\href{https://helpman.komtera.lt/?fen=8/6Qk/8/8/8/B7/P7/2K5&moves=7}{\color{white}\bfseries lösen\,\ensuremath{\nearrow}}}
\end{tcolorbox}
\vspace{2mm}
\noindent\begin{minipage}[t]{36mm}\vspace{0pt}\sffamily\small\raggedright
\hmfigure{längste eindeutige}{{\Large h\#7}}\hmfigure{eindeutig bei dieser Länge}{{\Large 73} Stellungen}\hmfigure{längstes Matt}{{\Large h\#8}}\hmfigure{in der Klasse}{710 Millionen Stellungen,\\ 477 Millionen mit Hilfsmatt}
\end{minipage}\hfill
\begin{minipage}[t]{80mm}\vspace{0pt}\begin{tikzpicture}\begin{axis}[width=72mm,height=42mm,ybar,bar width=1.6pt,xmin=-1,xmax=17,ymin=0,ymax=250.2,axis y line*=left,axis x line*=bottom,xlabel={\scriptsize Halbzüge bis Matt},ylabel={\scriptsize Stellungen (Mio.)},tick label style={font=\tiny},label style={font=\scriptsize},enlarge x limits=false,ylabel shift=-3pt,xlabel shift=-3pt]\addplot[fill=ink!55,draw=none] coordinates {(0,1.715) (1,10.32) (2,33.69) (3,95.27) (4,227.4) (5,66.35) (6,37.8) (7,0.1893) (8,1.216) (9,3e-06) (10,0.994) (12,1.207) (14,0.7187) (16,0.01318)};\draw[densely dotted] (axis cs:14,0) -- (axis cs:14,250.2);\end{axis}\begin{axis}[width=72mm,height=42mm,axis y line*=right,axis x line=none,ymode=log,xmin=-1,xmax=17,ylabel={\scriptsize eindeutig (log)},tick label style={font=\tiny},label style={font=\scriptsize},ylabel shift=-3pt]\addplot[mark=*,mark size=0.7pt,draw=black,thin] coordinates {(0,1714611) (1,7807457) (2,24413701) (3,25885531) (4,16873956) (5,35921) (6,90842) (8,29215) (10,15589) (12,9304) (14,73)};\end{axis}\end{tikzpicture}\end{minipage}
\vspace{3mm}

\noindent\begin{minipage}[t]{60mm}\vspace{0pt}\sbox\hmboard{\chessboard[setfen=8/6Qk/8/8/8/B7/P7/2K5 b - - 0 1]}\usebox\hmboard\\[1mm]
\makebox[\wd\hmboard]{\sffamily\hspace{3mm}\textbf{h\#7}\hfill(4+1)}\end{minipage}\hfill
\begin{minipage}[t]{52mm}\vspace{0pt}
\textcolor{ink}{\sffamily\textsc{Lösung}}\par\smallskip\small
\begin{tabular}{@{}r@{\enspace}l@{\enspace}l@{}}
\textbf{1.} & \symking{}xg7 & \symbishop{}b2+\\
\textbf{2.} & \symking{}f7 & a4\\
\textbf{3.} & \symking{}e6 & a5\\
\textbf{4.} & \symking{}d5 & a6\\
\textbf{5.} & \symking{}c4 & a7\\
\textbf{6.} & \symking{}b3 & a8=\symbishop{}\\
\textbf{7.} & \symking{}a2 & \symbishop{}d5\#\\
\end{tabular}\par\vspace{3mm}
\qrcode[height=12mm]{https://helpman.komtera.lt/?fen=8/6Qk/8/8/8/B7/P7/2K5&moves=7}\\{\sffamily\tiny\color{ink}scannen zum Lösen}
\end{minipage}
\vspace{2mm}

{\raggedright\noindent\textcolor{ink}{\sffamily\textsc{Themen}}\enspace \themetag{reines Matt} \themetag{Mustermatt} \themetag{Idealmatt} \themetag{Umwandlung} \themetag{Unterverwandlung} \themetag{Excelsior} \themetag{ein Stein zieht}\par\smallskip{\footnotesize\textit{Hinweis:} die Partei am Zug steht im Diagramm im Schach; die Lösung beginnt mit einem Schlag.}\par}

\clearpage
\phantomsection\label{hm:KQBvkb}
\begin{tcolorbox}[enhanced,colback=ink,colframe=ink,arc=0pt,boxrule=0pt,left=4mm,right=4mm,top=2.5mm,bottom=2.5mm]
\color{white}\sffamily{\LARGE\bfseries KQBvkb}\hfill{\Large h\#7}\\[0.5mm]
{\small 5 Steine\quad Schwarz am Zug}\hfill{\small Nr.~125\qquad\href{https://helpman.komtera.lt/?fen=8/8/6b1/8/8/6Qk/3B4/K7&moves=7}{\color{white}\bfseries lösen\,\ensuremath{\nearrow}}}
\end{tcolorbox}
\vspace{2mm}
\noindent\begin{minipage}[t]{36mm}\vspace{0pt}\sffamily\small\raggedright
\hmfigure{längste eindeutige}{{\Large h\#7}}\hmfigure{eindeutig bei dieser Länge}{{\Large 4} Stellungen}\hmfigure{längstes Matt}{{\Large h\#10}}\hmfigure{in der Klasse}{242 Millionen Stellungen,\\ 147 Millionen mit Hilfsmatt}
\end{minipage}\hfill
\begin{minipage}[t]{80mm}\vspace{0pt}\begin{tikzpicture}\begin{axis}[width=72mm,height=42mm,ybar,bar width=1.6pt,xmin=-1,xmax=21,ymin=0,ymax=74.58,axis y line*=left,axis x line*=bottom,xlabel={\scriptsize Halbzüge bis Matt},ylabel={\scriptsize Stellungen (Mio.)},tick label style={font=\tiny},label style={font=\scriptsize},enlarge x limits=false,ylabel shift=-3pt,xlabel shift=-3pt]\addplot[fill=ink!55,draw=none] coordinates {(0,0.3522) (1,2.286) (2,9.813) (3,31.89) (4,67.8) (5,24.46) (6,9.597) (7,0.1905) (8,0.1476) (9,1e-06) (10,0.1407) (12,0.1873) (14,0.1301) (16,0.1011) (17,3.6e-05) (18,0.07632) (20,0.01137)};\draw[densely dotted] (axis cs:14,0) -- (axis cs:14,74.58);\end{axis}\begin{axis}[width=72mm,height=42mm,axis y line*=right,axis x line=none,ymode=log,xmin=-1,xmax=21,ylabel={\scriptsize eindeutig (log)},tick label style={font=\tiny},label style={font=\scriptsize},ylabel shift=-3pt]\addplot[mark=*,mark size=0.7pt,draw=black,thin] coordinates {(0,352159) (1,1908703) (2,5961044) (3,6610397) (4,3308592) (5,42315) (6,4832) (8,1802) (10,438) (12,184) (14,4)};\end{axis}\end{tikzpicture}\end{minipage}
\vspace{3mm}

\noindent\begin{minipage}[t]{60mm}\vspace{0pt}\sbox\hmboard{\chessboard[setfen=8/8/6b1/8/8/6Qk/3B4/K7 b - - 0 1]}\usebox\hmboard\\[1mm]
\makebox[\wd\hmboard]{\sffamily\hspace{3mm}\textbf{h\#7}\hfill(3+2)}\end{minipage}\hfill
\begin{minipage}[t]{52mm}\vspace{0pt}
\textcolor{ink}{\sffamily\textsc{Lösung}}\par\smallskip\small
\begin{tabular}{@{}r@{\enspace}l@{\enspace}l@{}}
\textbf{1.} & \symking{}xg3 & \symking{}b2\\
\textbf{2.} & \symking{}g4 & \symking{}c3\\
\textbf{3.} & \symking{}f5 & \symking{}d4\\
\textbf{4.} & \symking{}f6 & \symbishop{}c3\\
\textbf{5.} & \symking{}g7 & \symking{}e5\\
\textbf{6.} & \symking{}h8 & \symking{}f6\\
\textbf{7.} & \symbishop{}h7 & \symking{}f7\#\\
\end{tabular}\par\vspace{3mm}
\qrcode[height=12mm]{https://helpman.komtera.lt/?fen=8/8/6b1/8/8/6Qk/3B4/K7&moves=7}\\{\sffamily\tiny\color{ink}scannen zum Lösen}
\end{minipage}
\vspace{2mm}

{\raggedright\noindent\textcolor{ink}{\sffamily\textsc{Themen}}\enspace \themetag{Selbstblock} \themetag{ohne Schach}\par\smallskip{\footnotesize\textit{Hinweis:} die Partei am Zug steht im Diagramm im Schach; die Lösung beginnt mit einem Schlag.}\par}

\clearpage
\phantomsection\label{hm:KQBvkp}
\begin{tcolorbox}[enhanced,colback=ink,colframe=ink,arc=0pt,boxrule=0pt,left=4mm,right=4mm,top=2.5mm,bottom=2.5mm]
\color{white}\sffamily{\LARGE\bfseries KQBvkp}\hfill{\Large h\#7}\\[0.5mm]
{\small 5 Steine\quad Schwarz am Zug}\hfill{\small Nr.~126\qquad\href{https://helpman.komtera.lt/?fen=8/8/8/8/2p5/8/2Q5/K1k1B3&moves=7}{\color{white}\bfseries lösen\,\ensuremath{\nearrow}}}
\end{tcolorbox}
\vspace{2mm}
\noindent\begin{minipage}[t]{36mm}\vspace{0pt}\sffamily\small\raggedright
\hmfigure{längste eindeutige}{{\Large h\#7}}\hmfigure{eindeutig bei dieser Länge}{{\Large 55} Stellungen}\hmfigure{längstes Matt}{{\Large h\#16}}\hmfigure{in der Klasse}{710 Millionen Stellungen,\\ 476 Millionen mit Hilfsmatt}
\end{minipage}\hfill
\begin{minipage}[t]{80mm}\vspace{0pt}\begin{tikzpicture}\begin{axis}[width=72mm,height=42mm,ybar,bar width=1.6pt,xmin=-1,xmax=33,ymin=0,ymax=245.6,axis y line*=left,axis x line*=bottom,xlabel={\scriptsize Halbzüge bis Matt},ylabel={\scriptsize Stellungen (Mio.)},tick label style={font=\tiny},label style={font=\scriptsize},enlarge x limits=false,ylabel shift=-3pt,xlabel shift=-3pt]\addplot[fill=ink!55,draw=none] coordinates {(0,1.438) (1,8.821) (2,32.93) (3,102.4) (4,223.2) (5,65.89) (6,36.11) (7,0.3195) (8,0.3133) (9,1e-05) (10,0.209) (12,0.2898) (14,0.4147) (16,0.5827) (18,0.6494) (20,0.6121) (22,0.6277) (24,0.5315) (26,0.3487) (28,0.2478) (30,0.1242) (32,0.03418)};\draw[densely dotted] (axis cs:14,0) -- (axis cs:14,245.6);\end{axis}\begin{axis}[width=72mm,height=42mm,axis y line*=right,axis x line=none,ymode=log,xmin=-1,xmax=33,ylabel={\scriptsize eindeutig (log)},tick label style={font=\tiny},label style={font=\scriptsize},ylabel shift=-3pt]\addplot[mark=*,mark size=0.7pt,draw=black,thin] coordinates {(0,1438084) (1,7064724) (2,23260243) (3,26920436) (4,12611628) (5,58264) (6,12226) (7,4) (8,3390) (10,1443) (12,644) (14,55)};\end{axis}\end{tikzpicture}\end{minipage}
\vspace{3mm}

\noindent\begin{minipage}[t]{60mm}\vspace{0pt}\sbox\hmboard{\chessboard[setfen=8/8/8/8/2p5/8/2Q5/K1k1B3 b - - 0 1]}\usebox\hmboard\\[1mm]
\makebox[\wd\hmboard]{\sffamily\hspace{3mm}\textbf{h\#7}\hfill(3+2)}\end{minipage}\hfill
\begin{minipage}[t]{52mm}\vspace{0pt}
\textcolor{ink}{\sffamily\textsc{Lösung}}\par\smallskip\small
\begin{tabular}{@{}r@{\enspace}l@{\enspace}l@{}}
\textbf{1.} & \symking{}xc2 & \symking{}a2\\
\textbf{2.} & c3 & \symking{}a3\\
\textbf{3.} & \symking{}b1 & \symking{}b3\\
\textbf{4.} & c2 & \symbishop{}b4\\
\textbf{5.} & \symking{}a1 & \symbishop{}a3\\
\textbf{6.} & c1=\symknight{}+ & \symking{}c2\\
\textbf{7.} & \symknight{}a2 & \symbishop{}b2\#\\
\end{tabular}\par\vspace{3mm}
\qrcode[height=12mm]{https://helpman.komtera.lt/?fen=8/8/8/8/2p5/8/2Q5/K1k1B3&moves=7}\\{\sffamily\tiny\color{ink}scannen zum Lösen}
\end{minipage}
\vspace{2mm}

{\raggedright\noindent\textcolor{ink}{\sffamily\textsc{Themen}}\enspace \themetag{reines Matt} \themetag{Mustermatt} \themetag{Idealmatt} \themetag{Umwandlung} \themetag{Unterverwandlung} \themetag{Selbstblock} \themetag{Umnov}\par\smallskip{\footnotesize\textit{Hinweis:} die Partei am Zug steht im Diagramm im Schach; die Lösung beginnt mit einem Schlag.}\par}

\clearpage
\phantomsection\label{hm:KQNNvk}
\begin{tcolorbox}[enhanced,colback=ink,colframe=ink,arc=0pt,boxrule=0pt,left=4mm,right=4mm,top=2.5mm,bottom=2.5mm]
\color{white}\sffamily{\LARGE\bfseries KQNNvk}\hfill{\Large h\#7}\\[0.5mm]
{\small 5 Steine\quad Schwarz am Zug}\hfill{\small Nr.~127\qquad\href{https://helpman.komtera.lt/?fen=kQ5N/8/8/8/8/8/7N/3K4&moves=7}{\color{white}\bfseries lösen\,\ensuremath{\nearrow}}}
\end{tcolorbox}
\vspace{2mm}
\noindent\begin{minipage}[t]{36mm}\vspace{0pt}\sffamily\small\raggedright
\hmfigure{längste eindeutige}{{\Large h\#7}}\hmfigure{eindeutig bei dieser Länge}{{\Large 3} Stellungen}\hmfigure{längstes Matt}{{\Large h\#8}}\hmfigure{in der Klasse}{242 Millionen Stellungen,\\ 79,7 Millionen mit Hilfsmatt}
\end{minipage}\hfill
\begin{minipage}[t]{80mm}\vspace{0pt}\begin{tikzpicture}\begin{axis}[width=72mm,height=42mm,ybar,bar width=1.6pt,xmin=-1,xmax=17,ymin=0,ymax=36.92,axis y line*=left,axis x line*=bottom,xlabel={\scriptsize Halbzüge bis Matt},ylabel={\scriptsize Stellungen (Mio.)},tick label style={font=\tiny},label style={font=\scriptsize},enlarge x limits=false,ylabel shift=-3pt,xlabel shift=-3pt]\addplot[fill=ink!55,draw=none] coordinates {(0,0.2491) (1,1.722) (2,5.749) (3,13.83) (4,33.56) (5,12.9) (6,10.35) (7,0.4396) (8,0.3274) (10,0.2769) (12,0.178) (14,0.07827) (16,0.001511)};\draw[densely dotted] (axis cs:14,0) -- (axis cs:14,36.92);\end{axis}\begin{axis}[width=72mm,height=42mm,axis y line*=right,axis x line=none,ymode=log,xmin=-1,xmax=17,ylabel={\scriptsize eindeutig (log)},tick label style={font=\tiny},label style={font=\scriptsize},ylabel shift=-3pt]\addplot[mark=*,mark size=0.7pt,draw=black,thin] coordinates {(0,249064) (1,1368281) (2,4206096) (3,2703755) (4,2455957) (5,47812) (6,25796) (8,2783) (10,1095) (12,101) (14,3)};\end{axis}\end{tikzpicture}\end{minipage}
\vspace{3mm}

\noindent\begin{minipage}[t]{60mm}\vspace{0pt}\sbox\hmboard{\chessboard[setfen=kQ5N/8/8/8/8/8/7N/3K4 b - - 0 1]}\usebox\hmboard\\[1mm]
\makebox[\wd\hmboard]{\sffamily\hspace{3mm}\textbf{h\#7}\hfill(4+1)}\end{minipage}\hfill
\begin{minipage}[t]{52mm}\vspace{0pt}
\textcolor{ink}{\sffamily\textsc{Lösung}}\par\smallskip\small
\begin{tabular}{@{}r@{\enspace}l@{\enspace}l@{}}
\textbf{1.} & \symking{}xb8 & \symking{}e2\\
\textbf{2.} & \symking{}c7 & \symking{}f3\\
\textbf{3.} & \symking{}d6 & \symking{}g4\\
\textbf{4.} & \symking{}e5 & \symking{}h5\\
\textbf{5.} & \symking{}f4 & \symknight{}g6+\\
\textbf{6.} & \symking{}g3 & \symknight{}f1+\\
\textbf{7.} & \symking{}h3 & \symknight{}f4\#\\
\end{tabular}\par\vspace{3mm}
\qrcode[height=12mm]{https://helpman.komtera.lt/?fen=kQ5N/8/8/8/8/8/7N/3K4&moves=7}\\{\sffamily\tiny\color{ink}scannen zum Lösen}
\end{minipage}
\vspace{2mm}

{\raggedright\noindent\textcolor{ink}{\sffamily\textsc{Themen}}\enspace \themetag{reines Matt} \themetag{Mustermatt} \themetag{Idealmatt} \themetag{Spiegelmatt} \themetag{ein Stein zieht}\par\smallskip{\footnotesize\textit{Hinweis:} die Partei am Zug steht im Diagramm im Schach; die Lösung beginnt mit einem Schlag.}\par}

\clearpage
\phantomsection\label{hm:KQNPvk}
\begin{tcolorbox}[enhanced,colback=ink,colframe=ink,arc=0pt,boxrule=0pt,left=4mm,right=4mm,top=2.5mm,bottom=2.5mm]
\color{white}\sffamily{\LARGE\bfseries KQNPvk}\hfill{\Large h\#7}\\[0.5mm]
{\small 5 Steine\quad Schwarz am Zug}\hfill{\small Nr.~128\qquad\href{https://helpman.komtera.lt/?fen=8/8/8/8/8/6N1/Q5P1/k2K4&moves=7}{\color{white}\bfseries lösen\,\ensuremath{\nearrow}}}
\end{tcolorbox}
\vspace{2mm}
\noindent\begin{minipage}[t]{36mm}\vspace{0pt}\sffamily\small\raggedright
\hmfigure{längste eindeutige}{{\Large h\#7}}\hmfigure{eindeutig bei dieser Länge}{{\Large 160} Stellungen}\hmfigure{längstes Matt}{{\Large h\#9}}\hmfigure{in der Klasse}{710 Millionen Stellungen,\\ 489 Millionen mit Hilfsmatt}
\end{minipage}\hfill
\begin{minipage}[t]{80mm}\vspace{0pt}\begin{tikzpicture}\begin{axis}[width=72mm,height=42mm,ybar,bar width=1.6pt,xmin=-1,xmax=19,ymin=0,ymax=195.3,axis y line*=left,axis x line*=bottom,xlabel={\scriptsize Halbzüge bis Matt},ylabel={\scriptsize Stellungen (Mio.)},tick label style={font=\tiny},label style={font=\scriptsize},enlarge x limits=false,ylabel shift=-3pt,xlabel shift=-3pt]\addplot[fill=ink!55,draw=none] coordinates {(0,1.043) (1,7.773) (2,23.64) (3,69.98) (4,177.6) (5,96.83) (6,95.14) (7,9.073) (8,5.345) (9,0.0173) (10,0.8514) (12,0.9779) (14,0.7718) (16,0.03224) (18,2e-06)};\draw[densely dotted] (axis cs:14,0) -- (axis cs:14,195.3);\end{axis}\begin{axis}[width=72mm,height=42mm,axis y line*=right,axis x line=none,ymode=log,xmin=-1,xmax=19,ylabel={\scriptsize eindeutig (log)},tick label style={font=\tiny},label style={font=\scriptsize},ylabel shift=-3pt]\addplot[mark=*,mark size=0.7pt,draw=black,thin] coordinates {(0,1043224) (1,5746429) (2,16855890) (3,14900282) (4,17387390) (5,1107815) (6,567790) (7,1384) (8,17633) (10,7890) (12,5327) (14,160)};\end{axis}\end{tikzpicture}\end{minipage}
\vspace{3mm}

\noindent\begin{minipage}[t]{60mm}\vspace{0pt}\sbox\hmboard{\chessboard[setfen=8/8/8/8/8/6N1/Q5P1/k2K4 b - - 0 1]}\usebox\hmboard\\[1mm]
\makebox[\wd\hmboard]{\sffamily\hspace{3mm}\textbf{h\#7}\hfill(4+1)}\end{minipage}\hfill
\begin{minipage}[t]{52mm}\vspace{0pt}
\textcolor{ink}{\sffamily\textsc{Lösung}}\par\smallskip\small
\begin{tabular}{@{}r@{\enspace}l@{\enspace}l@{}}
\textbf{1.} & \symking{}xa2 & \symknight{}e2\\
\textbf{2.} & \symking{}b3 & g4\\
\textbf{3.} & \symking{}c4 & g5\\
\textbf{4.} & \symking{}d3 & g6\\
\textbf{5.} & \symking{}e3 & g7\\
\textbf{6.} & \symking{}f2 & g8=\symqueen{}\\
\textbf{7.} & \symking{}f1 & \symqueen{}g1\#\\
\end{tabular}\par\vspace{3mm}
\qrcode[height=12mm]{https://helpman.komtera.lt/?fen=8/8/8/8/8/6N1/Q5P1/k2K4&moves=7}\\{\sffamily\tiny\color{ink}scannen zum Lösen}
\end{minipage}
\vspace{2mm}

{\raggedright\noindent\textcolor{ink}{\sffamily\textsc{Themen}}\enspace \themetag{Umwandlung} \themetag{Excelsior} \themetag{ein Stein zieht} \themetag{Phönix} \themetag{ohne Schach}\par\smallskip{\footnotesize\textit{Hinweis:} die Partei am Zug steht im Diagramm im Schach; die Lösung beginnt mit einem Schlag.}\par}

\clearpage
\phantomsection\label{hm:KQPPvk}
\begin{tcolorbox}[enhanced,colback=ink,colframe=ink,arc=0pt,boxrule=0pt,left=4mm,right=4mm,top=2.5mm,bottom=2.5mm]
\color{white}\sffamily{\LARGE\bfseries KQPPvk}\hfill{\Large h\#7}\\[0.5mm]
{\small 5 Steine\quad Schwarz am Zug}\hfill{\small Nr.~129\qquad\href{https://helpman.komtera.lt/?fen=6Qk/8/8/8/8/1P6/1P6/2K5&moves=7}{\color{white}\bfseries lösen\,\ensuremath{\nearrow}}}
\end{tcolorbox}
\vspace{2mm}
\noindent\begin{minipage}[t]{36mm}\vspace{0pt}\sffamily\small\raggedright
\hmfigure{längste eindeutige}{{\Large h\#7}}\hmfigure{eindeutig bei dieser Länge}{{\Large 37} Stellungen}\hmfigure{längstes Matt}{{\Large h\#8}}\hmfigure{in der Klasse}{533 Millionen Stellungen,\\ 187 Millionen mit Hilfsmatt}
\end{minipage}\hfill
\begin{minipage}[t]{80mm}\vspace{0pt}\begin{tikzpicture}\begin{axis}[width=72mm,height=42mm,ybar,bar width=1.6pt,xmin=-1,xmax=17,ymin=0,ymax=56.77,axis y line*=left,axis x line*=bottom,xlabel={\scriptsize Halbzüge bis Matt},ylabel={\scriptsize Stellungen (Mio.)},tick label style={font=\tiny},label style={font=\scriptsize},enlarge x limits=false,ylabel shift=-3pt,xlabel shift=-3pt]\addplot[fill=ink!55,draw=none] coordinates {(0,0.2865) (1,2.202) (2,5.854) (3,20.79) (4,51.61) (5,36.45) (6,45.45) (7,11.95) (8,9.797) (9,1.058) (10,0.8233) (11,0.01055) (12,0.2434) (13,8e-06) (14,0.09622) (16,0.002268)};\draw[densely dotted] (axis cs:14,0) -- (axis cs:14,56.77);\end{axis}\begin{axis}[width=72mm,height=42mm,axis y line*=right,axis x line=none,ymode=log,xmin=-1,xmax=17,ylabel={\scriptsize eindeutig (log)},tick label style={font=\tiny},label style={font=\scriptsize},ylabel shift=-3pt]\addplot[mark=*,mark size=0.7pt,draw=black,thin] coordinates {(0,286485) (1,1350390) (2,3548686) (3,4160541) (4,5532290) (5,1112306) (6,670967) (7,56399) (8,25037) (9,220) (10,1748) (11,3) (12,640) (14,37)};\end{axis}\end{tikzpicture}\end{minipage}
\vspace{3mm}

\noindent\begin{minipage}[t]{60mm}\vspace{0pt}\sbox\hmboard{\chessboard[setfen=6Qk/8/8/8/8/1P6/1P6/2K5 b - - 0 1]}\usebox\hmboard\\[1mm]
\makebox[\wd\hmboard]{\sffamily\hspace{3mm}\textbf{h\#7}\hfill(4+1)}\end{minipage}\hfill
\begin{minipage}[t]{52mm}\vspace{0pt}
\textcolor{ink}{\sffamily\textsc{Lösung}}\par\smallskip\small
\begin{tabular}{@{}r@{\enspace}l@{\enspace}l@{}}
\textbf{1.} & \symking{}xg8 & b4\\
\textbf{2.} & \symking{}f7 & b5\\
\textbf{3.} & \symking{}e6 & b6\\
\textbf{4.} & \symking{}d5 & b7\\
\textbf{5.} & \symking{}c4 & b3+\\
\textbf{6.} & \symking{}xb3 & b8=\symqueen{}+\\
\textbf{7.} & \symking{}a2 & \symqueen{}b2\#\\
\end{tabular}\par\vspace{3mm}
\qrcode[height=12mm]{https://helpman.komtera.lt/?fen=6Qk/8/8/8/8/1P6/1P6/2K5&moves=7}\\{\sffamily\tiny\color{ink}scannen zum Lösen}
\end{minipage}
\vspace{2mm}

{\raggedright\noindent\textcolor{ink}{\sffamily\textsc{Themen}}\enspace \themetag{Umwandlung} \themetag{ein Stein zieht} \themetag{Phönix}\par\smallskip{\footnotesize\textit{Hinweis:} die Partei am Zug steht im Diagramm im Schach; die Lösung beginnt mit einem Schlag.}\par}

\clearpage
\phantomsection\label{hm:KQPvkq}
\begin{tcolorbox}[enhanced,colback=ink,colframe=ink,arc=0pt,boxrule=0pt,left=4mm,right=4mm,top=2.5mm,bottom=2.5mm]
\color{white}\sffamily{\LARGE\bfseries KQPvkq}\hfill{\Large h\#7}\\[0.5mm]
{\small 5 Steine\quad Schwarz am Zug}\hfill{\small Nr.~130\qquad\href{https://helpman.komtera.lt/?fen=8/8/kQ6/8/2q5/8/2P5/K7&moves=7}{\color{white}\bfseries lösen\,\ensuremath{\nearrow}}}
\end{tcolorbox}
\vspace{2mm}
\noindent\begin{minipage}[t]{36mm}\vspace{0pt}\sffamily\small\raggedright
\hmfigure{längste eindeutige}{{\Large h\#7}}\hmfigure{eindeutig bei dieser Länge}{{\Large 521} Stellungen}\hmfigure{längstes Matt}{{\Large h\#9}}\hmfigure{in der Klasse}{710 Millionen Stellungen,\\ 404 Millionen mit Hilfsmatt}
\end{minipage}\hfill
\begin{minipage}[t]{80mm}\vspace{0pt}\begin{tikzpicture}\begin{axis}[width=72mm,height=42mm,ybar,bar width=1.6pt,xmin=-1,xmax=19,ymin=0,ymax=106.2,axis y line*=left,axis x line*=bottom,xlabel={\scriptsize Halbzüge bis Matt},ylabel={\scriptsize Stellungen (Mio.)},tick label style={font=\tiny},label style={font=\scriptsize},enlarge x limits=false,ylabel shift=-3pt,xlabel shift=-3pt]\addplot[fill=ink!55,draw=none] coordinates {(0,0.2157) (1,1.92) (2,7.849) (3,34.07) (4,76.52) (5,96.54) (6,87.64) (7,56.09) (8,28.68) (9,10.22) (10,2.991) (11,0.2932) (12,0.7074) (13,0.001226) (14,0.5351) (16,0.06999) (18,0.000359)};\draw[densely dotted] (axis cs:14,0) -- (axis cs:14,106.2);\end{axis}\begin{axis}[width=72mm,height=42mm,axis y line*=right,axis x line=none,ymode=log,xmin=-1,xmax=19,ylabel={\scriptsize eindeutig (log)},tick label style={font=\tiny},label style={font=\scriptsize},ylabel shift=-3pt]\addplot[mark=*,mark size=0.7pt,draw=black,thin] coordinates {(0,215721) (1,1528920) (2,3883562) (3,6499979) (4,6606255) (5,3451230) (6,1262394) (7,452238) (8,96220) (9,9973) (10,2968) (11,107) (12,1670) (14,521)};\end{axis}\end{tikzpicture}\end{minipage}
\vspace{3mm}

\noindent\begin{minipage}[t]{60mm}\vspace{0pt}\sbox\hmboard{\chessboard[setfen=8/8/kQ6/8/2q5/8/2P5/K7 b - - 0 1]}\usebox\hmboard\\[1mm]
\makebox[\wd\hmboard]{\sffamily\hspace{3mm}\textbf{h\#7}\hfill(3+2)}\end{minipage}\hfill
\begin{minipage}[t]{52mm}\vspace{0pt}
\textcolor{ink}{\sffamily\textsc{Lösung}}\par\smallskip\small
\begin{tabular}{@{}r@{\enspace}l@{\enspace}l@{}}
\textbf{1.} & \symking{}xb6 & \symking{}b1\\
\textbf{2.} & \symqueen{}e2 & c4\\
\textbf{3.} & \symking{}a5 & c5\\
\textbf{4.} & \symking{}b4 & c6\\
\textbf{5.} & \symking{}c3 & c7\\
\textbf{6.} & \symking{}d2 & c8=\symqueen{}\\
\textbf{7.} & \symking{}d1 & \symqueen{}c1\#\\
\end{tabular}\par\vspace{3mm}
\qrcode[height=12mm]{https://helpman.komtera.lt/?fen=8/8/kQ6/8/2q5/8/2P5/K7&moves=7}\\{\sffamily\tiny\color{ink}scannen zum Lösen}
\end{minipage}
\vspace{2mm}

{\raggedright\noindent\textcolor{ink}{\sffamily\textsc{Themen}}\enspace \themetag{Umwandlung} \themetag{Excelsior} \themetag{Selbstblock} \themetag{Phönix} \themetag{ohne Schach} \themetag{Umnov}\par\smallskip{\footnotesize\textit{Hinweis:} die Partei am Zug steht im Diagramm im Schach; die Lösung beginnt mit einem Schlag.}\par}

\clearpage
\phantomsection\label{hm:KQPvkr}
\begin{tcolorbox}[enhanced,colback=ink,colframe=ink,arc=0pt,boxrule=0pt,left=4mm,right=4mm,top=2.5mm,bottom=2.5mm]
\color{white}\sffamily{\LARGE\bfseries KQPvkr}\hfill{\Large h\#7}\\[0.5mm]
{\small 5 Steine\quad Schwarz am Zug}\hfill{\small Nr.~131\qquad\href{https://helpman.komtera.lt/?fen=4k3/4Q3/8/8/3r4/8/3P4/K7&moves=7}{\color{white}\bfseries lösen\,\ensuremath{\nearrow}}}
\end{tcolorbox}
\vspace{2mm}
\noindent\begin{minipage}[t]{36mm}\vspace{0pt}\sffamily\small\raggedright
\hmfigure{längste eindeutige}{{\Large h\#7}}\hmfigure{eindeutig bei dieser Länge}{{\Large 105} Stellungen}\hmfigure{längstes Matt}{{\Large h\#9}}\hmfigure{in der Klasse}{710 Millionen Stellungen,\\ 444 Millionen mit Hilfsmatt}
\end{minipage}\hfill
\begin{minipage}[t]{80mm}\vspace{0pt}\begin{tikzpicture}\begin{axis}[width=72mm,height=42mm,ybar,bar width=1.6pt,xmin=-1,xmax=19,ymin=0,ymax=119,axis y line*=left,axis x line*=bottom,xlabel={\scriptsize Halbzüge bis Matt},ylabel={\scriptsize Stellungen (Mio.)},tick label style={font=\tiny},label style={font=\scriptsize},enlarge x limits=false,ylabel shift=-3pt,xlabel shift=-3pt]\addplot[fill=ink!55,draw=none] coordinates {(0,0.3407) (1,2.884) (2,11.34) (3,42.62) (4,108.2) (5,104.2) (6,103) (7,45.38) (8,19.76) (9,4.101) (10,1.163) (11,0.01111) (12,0.7496) (13,4.3e-05) (14,0.4478) (16,0.01465) (18,2.1e-05)};\draw[densely dotted] (axis cs:14,0) -- (axis cs:14,119);\end{axis}\begin{axis}[width=72mm,height=42mm,axis y line*=right,axis x line=none,ymode=log,xmin=-1,xmax=19,ylabel={\scriptsize eindeutig (log)},tick label style={font=\tiny},label style={font=\scriptsize},ylabel shift=-3pt]\addplot[mark=*,mark size=0.7pt,draw=black,thin] coordinates {(0,340705) (1,2157826) (2,5958403) (3,6957971) (4,7929036) (5,2771532) (6,999371) (7,196384) (8,40525) (9,1395) (10,3159) (12,1645) (14,105)};\end{axis}\end{tikzpicture}\end{minipage}
\vspace{3mm}

\noindent\begin{minipage}[t]{60mm}\vspace{0pt}\sbox\hmboard{\chessboard[setfen=4k3/4Q3/8/8/3r4/8/3P4/K7 b - - 0 1]}\usebox\hmboard\\[1mm]
\makebox[\wd\hmboard]{\sffamily\hspace{3mm}\textbf{h\#7}\hfill(3+2)}\end{minipage}\hfill
\begin{minipage}[t]{52mm}\vspace{0pt}
\textcolor{ink}{\sffamily\textsc{Lösung}}\par\smallskip\small
\begin{tabular}{@{}r@{\enspace}l@{\enspace}l@{}}
\textbf{1.} & \symking{}xe7 & d3\\
\textbf{2.} & \symrook{}b4 & d4\\
\textbf{3.} & \symking{}d6 & d5\\
\textbf{4.} & \symking{}c5 & d6\\
\textbf{5.} & \symrook{}b3 & d7\\
\textbf{6.} & \symking{}b4 & d8=\symqueen{}\\
\textbf{7.} & \symking{}a3 & \symqueen{}a5\#\\
\end{tabular}\par\vspace{3mm}
\qrcode[height=12mm]{https://helpman.komtera.lt/?fen=4k3/4Q3/8/8/3r4/8/3P4/K7&moves=7}\\{\sffamily\tiny\color{ink}scannen zum Lösen}
\end{minipage}
\vspace{2mm}

{\raggedright\noindent\textcolor{ink}{\sffamily\textsc{Themen}}\enspace \themetag{Umwandlung} \themetag{Excelsior} \themetag{Selbstblock} \themetag{ein Stein zieht} \themetag{Phönix} \themetag{ohne Schach} \themetag{Umnov}\par\smallskip{\footnotesize\textit{Hinweis:} die Partei am Zug steht im Diagramm im Schach; die Lösung beginnt mit einem Schlag.}\par}

\clearpage
\phantomsection\label{hm:KRBBvk}
\begin{tcolorbox}[enhanced,colback=ink,colframe=ink,arc=0pt,boxrule=0pt,left=4mm,right=4mm,top=2.5mm,bottom=2.5mm]
\color{white}\sffamily{\LARGE\bfseries KRBBvk}\hfill{\Large h\#7}\\[0.5mm]
{\small 5 Steine\quad Schwarz am Zug}\hfill{\small Nr.~132\qquad\href{https://helpman.komtera.lt/?fen=7R/7k/7B/8/8/8/2B5/K7&moves=7}{\color{white}\bfseries lösen\,\ensuremath{\nearrow}}}
\end{tcolorbox}
\vspace{2mm}
\noindent\begin{minipage}[t]{36mm}\vspace{0pt}\sffamily\small\raggedright
\hmfigure{längste eindeutige}{{\Large h\#7}}\hmfigure{eindeutig bei dieser Länge}{{\Large 1} Stellung}\hmfigure{längstes Matt}{{\Large h\#8}}\hmfigure{in der Klasse}{242 Millionen Stellungen,\\ 81,6 Millionen mit Hilfsmatt}
\end{minipage}\hfill
\begin{minipage}[t]{80mm}\vspace{0pt}\begin{tikzpicture}\begin{axis}[width=72mm,height=42mm,ybar,bar width=1.6pt,xmin=-1,xmax=17,ymin=0,ymax=24.18,axis y line*=left,axis x line*=bottom,xlabel={\scriptsize Halbzüge bis Matt},ylabel={\scriptsize Stellungen (Mio.)},tick label style={font=\tiny},label style={font=\scriptsize},enlarge x limits=false,ylabel shift=-3pt,xlabel shift=-3pt]\addplot[fill=ink!55,draw=none] coordinates {(0,0.1522) (1,0.6002) (2,1.61) (3,5.982) (4,18.8) (5,16.45) (6,21.98) (7,6.377) (8,7.215) (9,1.255) (10,0.992) (11,0.08871) (12,0.07416) (13,0.001796) (14,0.002684) (15,3e-06) (16,5e-06)};\draw[densely dotted] (axis cs:14,0) -- (axis cs:14,24.18);\end{axis}\begin{axis}[width=72mm,height=42mm,axis y line*=right,axis x line=none,ymode=log,xmin=-1,xmax=17,ylabel={\scriptsize eindeutig (log)},tick label style={font=\tiny},label style={font=\scriptsize},ylabel shift=-3pt]\addplot[mark=*,mark size=0.7pt,draw=black,thin] coordinates {(0,152151) (1,547667) (2,1418017) (3,3367102) (4,5763092) (5,1304774) (6,1135690) (7,55599) (8,28639) (9,1255) (10,912) (11,13) (12,34) (14,1)};\end{axis}\end{tikzpicture}\end{minipage}
\vspace{3mm}

\noindent\begin{minipage}[t]{60mm}\vspace{0pt}\sbox\hmboard{\chessboard[setfen=7R/7k/7B/8/8/8/2B5/K7 b - - 0 1]}\usebox\hmboard\\[1mm]
\makebox[\wd\hmboard]{\sffamily\hspace{3mm}\textbf{h\#7}\hfill(4+1)}\end{minipage}\hfill
\begin{minipage}[t]{52mm}\vspace{0pt}
\textcolor{ink}{\sffamily\textsc{Lösung}}\par\smallskip\small
\begin{tabular}{@{}r@{\enspace}l@{\enspace}l@{}}
\textbf{1.} & \symking{}xh8 & \symking{}b2\\
\textbf{2.} & \symking{}g8 & \symking{}c3\\
\textbf{3.} & \symking{}f7 & \symking{}d3\\
\textbf{4.} & \symking{}g6 & \symking{}e4\\
\textbf{5.} & \symking{}h5 & \symking{}f5\\
\textbf{6.} & \symking{}h4 & \symbishop{}g5+\\
\textbf{7.} & \symking{}h5 & \symbishop{}d1\#\\
\end{tabular}\par\vspace{3mm}
\qrcode[height=12mm]{https://helpman.komtera.lt/?fen=7R/7k/7B/8/8/8/2B5/K7&moves=7}\\{\sffamily\tiny\color{ink}scannen zum Lösen}
\end{minipage}
\vspace{2mm}

{\raggedright\noindent\textcolor{ink}{\sffamily\textsc{Themen}}\enspace \themetag{Switchback} \themetag{ein Stein zieht}\par\smallskip{\footnotesize\textit{Hinweis:} das Diagramm hat keinen legalen letzten Zug, kann also in keiner Partie entstehen.}\par}

\clearpage
\phantomsection\label{hm:KRBPvk}
\begin{tcolorbox}[enhanced,colback=ink,colframe=ink,arc=0pt,boxrule=0pt,left=4mm,right=4mm,top=2.5mm,bottom=2.5mm]
\color{white}\sffamily{\LARGE\bfseries KRBPvk}\hfill{\Large h\#7}\\[0.5mm]
{\small 5 Steine\quad Schwarz am Zug}\hfill{\small Nr.~133\qquad\href{https://helpman.komtera.lt/?fen=7k/6R1/8/8/8/B7/P7/2K5&moves=7}{\color{white}\bfseries lösen\,\ensuremath{\nearrow}}}
\end{tcolorbox}
\vspace{2mm}
\noindent\begin{minipage}[t]{36mm}\vspace{0pt}\sffamily\small\raggedright
\hmfigure{längste eindeutige}{{\Large h\#7}}\hmfigure{eindeutig bei dieser Länge}{{\Large 17} Stellungen}\hmfigure{längstes Matt}{{\Large h\#8}}\hmfigure{in der Klasse}{710 Millionen Stellungen,\\ 514 Millionen mit Hilfsmatt}
\end{minipage}\hfill
\begin{minipage}[t]{80mm}\vspace{0pt}\begin{tikzpicture}\begin{axis}[width=72mm,height=42mm,ybar,bar width=1.6pt,xmin=-1,xmax=17,ymin=0,ymax=156.7,axis y line*=left,axis x line*=bottom,xlabel={\scriptsize Halbzüge bis Matt},ylabel={\scriptsize Stellungen (Mio.)},tick label style={font=\tiny},label style={font=\scriptsize},enlarge x limits=false,ylabel shift=-3pt,xlabel shift=-3pt]\addplot[fill=ink!55,draw=none] coordinates {(0,0.4229) (1,2.626) (2,5.512) (3,24.79) (4,64.73) (5,90.19) (6,142.5) (7,70.71) (8,78.52) (9,18.46) (10,13.96) (11,1.004) (12,0.8407) (13,0.004338) (14,0.1536) (16,0.002549)};\draw[densely dotted] (axis cs:14,0) -- (axis cs:14,156.7);\end{axis}\begin{axis}[width=72mm,height=42mm,axis y line*=right,axis x line=none,ymode=log,xmin=-1,xmax=17,ylabel={\scriptsize eindeutig (log)},tick label style={font=\tiny},label style={font=\scriptsize},ylabel shift=-3pt]\addplot[mark=*,mark size=0.7pt,draw=black,thin] coordinates {(0,422928) (1,2091263) (2,3928490) (3,10066481) (4,19349519) (5,9999843) (6,7341982) (7,733470) (8,424753) (9,25581) (10,18758) (11,408) (12,2236) (14,17)};\end{axis}\end{tikzpicture}\end{minipage}
\vspace{3mm}

\noindent\begin{minipage}[t]{60mm}\vspace{0pt}\sbox\hmboard{\chessboard[setfen=7k/6R1/8/8/8/B7/P7/2K5 b - - 0 1]}\usebox\hmboard\\[1mm]
\makebox[\wd\hmboard]{\sffamily\hspace{3mm}\textbf{h\#7}\hfill(4+1)}\end{minipage}\hfill
\begin{minipage}[t]{52mm}\vspace{0pt}
\textcolor{ink}{\sffamily\textsc{Lösung}}\par\smallskip\small
\begin{tabular}{@{}r@{\enspace}l@{\enspace}l@{}}
\textbf{1.} & \symking{}xg7 & \symbishop{}b2+\\
\textbf{2.} & \symking{}f7 & a4\\
\textbf{3.} & \symking{}e6 & a5\\
\textbf{4.} & \symking{}d5 & a6\\
\textbf{5.} & \symking{}c4 & a7\\
\textbf{6.} & \symking{}b3 & a8=\symbishop{}\\
\textbf{7.} & \symking{}a2 & \symbishop{}d5\#\\
\end{tabular}\par\vspace{3mm}
\qrcode[height=12mm]{https://helpman.komtera.lt/?fen=7k/6R1/8/8/8/B7/P7/2K5&moves=7}\\{\sffamily\tiny\color{ink}scannen zum Lösen}
\end{minipage}
\vspace{2mm}

{\raggedright\noindent\textcolor{ink}{\sffamily\textsc{Themen}}\enspace \themetag{reines Matt} \themetag{Mustermatt} \themetag{Idealmatt} \themetag{Umwandlung} \themetag{Unterverwandlung} \themetag{Excelsior} \themetag{ein Stein zieht}\par\smallskip{\footnotesize\textit{Hinweis:} die Lösung beginnt mit einem Schlag.}\par}

\begin{tcolorbox}[enhanced,colback=white,colframe=ink!35,arc=1.5pt,boxrule=0.4pt,left=3mm,right=3mm,top=1.5mm,bottom=1.5mm]
\textcolor{ink}{\sffamily\textsc{veröffentlicht}}\par\smallskip\small \textbf{Abdurahmanovic, Fadil; Becker, Richard}, h\#7\quad{\color{ink}\scriptsize diese Stellung}\\{\scriptsize The Problemist, No. (R), 01/2006}
\end{tcolorbox}
\clearpage
\phantomsection\label{hm:KRNPvk}
\begin{tcolorbox}[enhanced,colback=ink,colframe=ink,arc=0pt,boxrule=0pt,left=4mm,right=4mm,top=2.5mm,bottom=2.5mm]
\color{white}\sffamily{\LARGE\bfseries KRNPvk}\hfill{\Large h\#7}\\[0.5mm]
{\small 5 Steine\quad Schwarz am Zug}\hfill{\small Nr.~134\qquad\href{https://helpman.komtera.lt/?fen=8/8/8/8/8/6N1/R5P1/1k1K4&moves=7}{\color{white}\bfseries lösen\,\ensuremath{\nearrow}}}
\end{tcolorbox}
\vspace{2mm}
\noindent\begin{minipage}[t]{36mm}\vspace{0pt}\sffamily\small\raggedright
\hmfigure{längste eindeutige}{{\Large h\#7}}\hmfigure{eindeutig bei dieser Länge}{{\Large 31} Stellungen}\hmfigure{längstes Matt}{{\Large h\#8}}\hmfigure{in der Klasse}{710 Millionen Stellungen,\\ 526 Millionen mit Hilfsmatt}
\end{minipage}\hfill
\begin{minipage}[t]{80mm}\vspace{0pt}\begin{tikzpicture}\begin{axis}[width=72mm,height=42mm,ybar,bar width=1.6pt,xmin=-1,xmax=17,ymin=0,ymax=164.4,axis y line*=left,axis x line*=bottom,xlabel={\scriptsize Halbzüge bis Matt},ylabel={\scriptsize Stellungen (Mio.)},tick label style={font=\tiny},label style={font=\scriptsize},enlarge x limits=false,ylabel shift=-3pt,xlabel shift=-3pt]\addplot[fill=ink!55,draw=none] coordinates {(0,0.3628) (1,2.505) (2,5.049) (3,24.61) (4,62.35) (5,96.97) (6,149.4) (7,80.43) (8,80.93) (9,14.39) (10,8.299) (11,0.199) (12,0.2955) (13,1.2e-05) (14,0.145) (16,0.004811)};\draw[densely dotted] (axis cs:14,0) -- (axis cs:14,164.4);\end{axis}\begin{axis}[width=72mm,height=42mm,axis y line*=right,axis x line=none,ymode=log,xmin=-1,xmax=17,ylabel={\scriptsize eindeutig (log)},tick label style={font=\tiny},label style={font=\scriptsize},ylabel shift=-3pt]\addplot[mark=*,mark size=0.7pt,draw=black,thin] coordinates {(0,362767) (1,1995651) (2,3643042) (3,9889633) (4,18130560) (5,9991851) (6,6611886) (7,568118) (8,184395) (9,2590) (10,3137) (12,667) (14,31)};\end{axis}\end{tikzpicture}\end{minipage}
\vspace{3mm}

\noindent\begin{minipage}[t]{60mm}\vspace{0pt}\sbox\hmboard{\chessboard[setfen=8/8/8/8/8/6N1/R5P1/1k1K4 b - - 0 1]}\usebox\hmboard\\[1mm]
\makebox[\wd\hmboard]{\sffamily\hspace{3mm}\textbf{h\#7}\hfill(4+1)}\end{minipage}\hfill
\begin{minipage}[t]{52mm}\vspace{0pt}
\textcolor{ink}{\sffamily\textsc{Lösung}}\par\smallskip\small
\begin{tabular}{@{}r@{\enspace}l@{\enspace}l@{}}
\textbf{1.} & \symking{}xa2 & \symknight{}e2\\
\textbf{2.} & \symking{}b3 & g4\\
\textbf{3.} & \symking{}c4 & g5\\
\textbf{4.} & \symking{}d3 & g6\\
\textbf{5.} & \symking{}e3 & g7\\
\textbf{6.} & \symking{}f2 & g8=\symqueen{}\\
\textbf{7.} & \symking{}f1 & \symqueen{}g1\#\\
\end{tabular}\par\vspace{3mm}
\qrcode[height=12mm]{https://helpman.komtera.lt/?fen=8/8/8/8/8/6N1/R5P1/1k1K4&moves=7}\\{\sffamily\tiny\color{ink}scannen zum Lösen}
\end{minipage}
\vspace{2mm}

{\raggedright\noindent\textcolor{ink}{\sffamily\textsc{Themen}}\enspace \themetag{Umwandlung} \themetag{Excelsior} \themetag{ein Stein zieht} \themetag{ohne Schach}\par\smallskip{\footnotesize\textit{Hinweis:} die Lösung beginnt mit einem Schlag.}\par}

\clearpage
\phantomsection\label{hm:KRNvkp}
\begin{tcolorbox}[enhanced,colback=ink,colframe=ink,arc=0pt,boxrule=0pt,left=4mm,right=4mm,top=2.5mm,bottom=2.5mm]
\color{white}\sffamily{\LARGE\bfseries KRNvkp}\hfill{\Large h\#7}\\[0.5mm]
{\small 5 Steine\quad Schwarz am Zug}\hfill{\small Nr.~135\qquad\href{https://helpman.komtera.lt/?fen=8/8/8/8/3R4/2kp4/3N4/2K5&moves=7}{\color{white}\bfseries lösen\,\ensuremath{\nearrow}}}
\end{tcolorbox}
\vspace{2mm}
\noindent\begin{minipage}[t]{36mm}\vspace{0pt}\sffamily\small\raggedright
\hmfigure{längste eindeutige}{{\Large h\#7}}\hmfigure{eindeutig bei dieser Länge}{{\Large 6} Stellungen}\hmfigure{längstes Matt}{{\Large h\#11.5}}\hmfigure{in der Klasse}{710 Millionen Stellungen,\\ 524 Millionen mit Hilfsmatt}
\end{minipage}\hfill
\begin{minipage}[t]{80mm}\vspace{0pt}\begin{tikzpicture}\begin{axis}[width=72mm,height=42mm,ybar,bar width=1.6pt,xmin=-1,xmax=24,ymin=0,ymax=156.5,axis y line*=left,axis x line*=bottom,xlabel={\scriptsize Halbzüge bis Matt},ylabel={\scriptsize Stellungen (Mio.)},tick label style={font=\tiny},label style={font=\scriptsize},enlarge x limits=false,ylabel shift=-3pt,xlabel shift=-3pt]\addplot[fill=ink!55,draw=none] coordinates {(0,0.3267) (1,1.908) (2,4.373) (3,16.97) (4,43.58) (5,81.27) (6,142.3) (7,102) (8,97.76) (9,22.67) (10,10.09) (11,0.3913) (12,0.1809) (13,7.2e-05) (14,0.06327) (15,5.4e-05) (16,0.065) (17,6.5e-05) (18,0.04725) (19,2.1e-05) (20,0.02356) (21,1e-06) (22,0.008177) (23,1e-06)};\draw[densely dotted] (axis cs:14,0) -- (axis cs:14,156.5);\end{axis}\begin{axis}[width=72mm,height=42mm,axis y line*=right,axis x line=none,ymode=log,xmin=-1,xmax=24,ylabel={\scriptsize eindeutig (log)},tick label style={font=\tiny},label style={font=\scriptsize},ylabel shift=-3pt]\addplot[mark=*,mark size=0.7pt,draw=black,thin] coordinates {(0,326675) (1,1874134) (2,3297179) (3,7476038) (4,14808826) (5,9073488) (6,6229681) (7,828719) (8,213499) (9,4248) (10,1450) (11,1) (12,63) (14,6)};\end{axis}\end{tikzpicture}\end{minipage}
\vspace{3mm}

\noindent\begin{minipage}[t]{60mm}\vspace{0pt}\sbox\hmboard{\chessboard[setfen=8/8/8/8/3R4/2kp4/3N4/2K5 b - - 0 1]}\usebox\hmboard\\[1mm]
\makebox[\wd\hmboard]{\sffamily\hspace{3mm}\textbf{h\#7}\hfill(3+2)}\end{minipage}\hfill
\begin{minipage}[t]{52mm}\vspace{0pt}
\textcolor{ink}{\sffamily\textsc{Lösung}}\par\smallskip\small
\begin{tabular}{@{}r@{\enspace}l@{\enspace}l@{}}
\textbf{1.} & \symking{}xd4 & \symknight{}f1\\
\textbf{2.} & d2+ & \symking{}b2\\
\textbf{3.} & d1=\symrook{} & \symking{}a3\\
\textbf{4.} & \symking{}c3 & \symking{}a4\\
\textbf{5.} & \symking{}b2 & \symknight{}e3\\
\textbf{6.} & \symking{}a1 & \symking{}a3\\
\textbf{7.} & \symrook{}b1 & \symknight{}c2\#\\
\end{tabular}\par\vspace{3mm}
\qrcode[height=12mm]{https://helpman.komtera.lt/?fen=8/8/8/8/3R4/2kp4/3N4/2K5&moves=7}\\{\sffamily\tiny\color{ink}scannen zum Lösen}
\end{minipage}
\vspace{2mm}

{\raggedright\noindent\textcolor{ink}{\sffamily\textsc{Themen}}\enspace \themetag{reines Matt} \themetag{Mustermatt} \themetag{Idealmatt} \themetag{Umwandlung} \themetag{Unterverwandlung} \themetag{Switchback} \themetag{Selbstblock} \themetag{Umnov}\par\smallskip{\footnotesize\textit{Hinweis:} die Lösung beginnt mit einem Schlag.}\par}

\clearpage
\phantomsection\label{hm:KRNvkr}
\begin{tcolorbox}[enhanced,colback=ink,colframe=ink,arc=0pt,boxrule=0pt,left=4mm,right=4mm,top=2.5mm,bottom=2.5mm]
\color{white}\sffamily{\LARGE\bfseries KRNvkr}\hfill{\Large h\#7}\\[0.5mm]
{\small 5 Steine\quad Schwarz am Zug}\hfill{\small Nr.~136\qquad\href{https://helpman.komtera.lt/?fen=8/8/8/8/8/5N2/1r3R2/K4k2&moves=7}{\color{white}\bfseries lösen\,\ensuremath{\nearrow}}}
\end{tcolorbox}
\vspace{2mm}
\noindent\begin{minipage}[t]{36mm}\vspace{0pt}\sffamily\small\raggedright
\hmfigure{längste eindeutige}{{\Large h\#7}}\hmfigure{eindeutig bei dieser Länge}{{\Large 9} Stellungen}\hmfigure{längstes Matt}{{\Large h\#10}}\hmfigure{in der Klasse}{242 Millionen Stellungen,\\ 156 Millionen mit Hilfsmatt}
\end{minipage}\hfill
\begin{minipage}[t]{80mm}\vspace{0pt}\begin{tikzpicture}\begin{axis}[width=72mm,height=42mm,ybar,bar width=1.6pt,xmin=-1,xmax=21,ymin=0,ymax=57.45,axis y line*=left,axis x line*=bottom,xlabel={\scriptsize Halbzüge bis Matt},ylabel={\scriptsize Stellungen (Mio.)},tick label style={font=\tiny},label style={font=\scriptsize},enlarge x limits=false,ylabel shift=-3pt,xlabel shift=-3pt]\addplot[fill=ink!55,draw=none] coordinates {(0,0.0509) (1,0.3202) (2,1.221) (3,5.814) (4,17.89) (5,36.9) (6,52.23) (7,30.67) (8,10.07) (9,1.186) (10,0.04643) (11,0.001469) (12,0.02391) (13,9e-06) (14,0.00774) (16,0.003291) (18,0.000633) (20,1e-06)};\draw[densely dotted] (axis cs:14,0) -- (axis cs:14,57.45);\end{axis}\begin{axis}[width=72mm,height=42mm,axis y line*=right,axis x line=none,ymode=log,xmin=-1,xmax=21,ylabel={\scriptsize eindeutig (log)},tick label style={font=\tiny},label style={font=\scriptsize},ylabel shift=-3pt]\addplot[mark=*,mark size=0.7pt,draw=black,thin] coordinates {(0,50899) (1,315893) (2,709835) (3,2256603) (4,5140743) (5,3397592) (6,1076232) (7,104065) (8,2270) (9,53) (10,179) (12,95) (14,9)};\end{axis}\end{tikzpicture}\end{minipage}
\vspace{3mm}

\noindent\begin{minipage}[t]{60mm}\vspace{0pt}\sbox\hmboard{\chessboard[setfen=8/8/8/8/8/5N2/1r3R2/K4k2 b - - 0 1]}\usebox\hmboard\\[1mm]
\makebox[\wd\hmboard]{\sffamily\hspace{3mm}\textbf{h\#7}\hfill(3+2)}\end{minipage}\hfill
\begin{minipage}[t]{52mm}\vspace{0pt}
\textcolor{ink}{\sffamily\textsc{Lösung}}\par\smallskip\small
\begin{tabular}{@{}r@{\enspace}l@{\enspace}l@{}}
\textbf{1.} & \symrook{}xf2 & \symking{}b1\\
\textbf{2.} & \symking{}e2 & \symking{}b2\\
\textbf{3.} & \symking{}d3+ & \symking{}c1\\
\textbf{4.} & \symking{}c3 & \symking{}d1\\
\textbf{5.} & \symking{}b2 & \symknight{}d4\\
\textbf{6.} & \symking{}a1 & \symking{}c1\\
\textbf{7.} & \symrook{}a2 & \symknight{}b3\#\\
\end{tabular}\par\vspace{3mm}
\qrcode[height=12mm]{https://helpman.komtera.lt/?fen=8/8/8/8/8/5N2/1r3R2/K4k2&moves=7}\\{\sffamily\tiny\color{ink}scannen zum Lösen}
\end{minipage}
\vspace{2mm}

{\raggedright\noindent\textcolor{ink}{\sffamily\textsc{Themen}}\enspace \themetag{reines Matt} \themetag{Mustermatt} \themetag{Idealmatt} \themetag{Switchback} \themetag{Selbstblock} \themetag{Inder}\par\smallskip{\footnotesize\textit{Hinweis:} die Partei am Zug steht im Diagramm im Schach; die Lösung beginnt mit einem Schlag.}\par}

\clearpage
\phantomsection\label{hm:KRPPvk}
\begin{tcolorbox}[enhanced,colback=ink,colframe=ink,arc=0pt,boxrule=0pt,left=4mm,right=4mm,top=2.5mm,bottom=2.5mm]
\color{white}\sffamily{\LARGE\bfseries KRPPvk}\hfill{\Large h\#7}\\[0.5mm]
{\small 5 Steine\quad Schwarz am Zug}\hfill{\small Nr.~137\qquad\href{https://helpman.komtera.lt/?fen=8/8/6R1/7k/8/6P1/K5P1/8&moves=7}{\color{white}\bfseries lösen\,\ensuremath{\nearrow}}}
\end{tcolorbox}
\vspace{2mm}
\noindent\begin{minipage}[t]{36mm}\vspace{0pt}\sffamily\small\raggedright
\hmfigure{längste eindeutige}{{\Large h\#7}}\hmfigure{eindeutig bei dieser Länge}{{\Large 2} Stellungen}\hmfigure{längstes Matt}{{\Large h\#8}}\hmfigure{in der Klasse}{533 Millionen Stellungen,\\ 201 Millionen mit Hilfsmatt}
\end{minipage}\hfill
\begin{minipage}[t]{80mm}\vspace{0pt}\begin{tikzpicture}\begin{axis}[width=72mm,height=42mm,ybar,bar width=1.6pt,xmin=-1,xmax=17,ymin=0,ymax=55.98,axis y line*=left,axis x line*=bottom,xlabel={\scriptsize Halbzüge bis Matt},ylabel={\scriptsize Stellungen (Mio.)},tick label style={font=\tiny},label style={font=\scriptsize},enlarge x limits=false,ylabel shift=-3pt,xlabel shift=-3pt]\addplot[fill=ink!55,draw=none] coordinates {(0,0.1204) (1,1.04) (2,2.093) (3,11.53) (4,26.32) (5,37.74) (6,50.89) (7,26.7) (8,28.47) (9,8.865) (10,6.264) (11,0.6782) (12,0.4621) (13,0.003661) (14,0.01758) (16,0.000183)};\draw[densely dotted] (axis cs:14,0) -- (axis cs:14,55.98);\end{axis}\begin{axis}[width=72mm,height=42mm,axis y line*=right,axis x line=none,ymode=log,xmin=-1,xmax=17,ylabel={\scriptsize eindeutig (log)},tick label style={font=\tiny},label style={font=\scriptsize},ylabel shift=-3pt]\addplot[mark=*,mark size=0.7pt,draw=black,thin] coordinates {(0,120384) (1,681759) (2,1150011) (3,3444946) (4,5271037) (5,2303424) (6,1450074) (7,188342) (8,86818) (9,7041) (10,3366) (11,175) (12,179) (14,2)};\end{axis}\end{tikzpicture}\end{minipage}
\vspace{3mm}

\noindent\begin{minipage}[t]{60mm}\vspace{0pt}\sbox\hmboard{\chessboard[setfen=8/8/6R1/7k/8/6P1/K5P1/8 b - - 0 1]}\usebox\hmboard\\[1mm]
\makebox[\wd\hmboard]{\sffamily\hspace{3mm}\textbf{h\#7}\hfill(4+1)}\end{minipage}\hfill
\begin{minipage}[t]{52mm}\vspace{0pt}
\textcolor{ink}{\sffamily\textsc{Lösung}}\par\smallskip\small
\begin{tabular}{@{}r@{\enspace}l@{\enspace}l@{}}
\textbf{1.} & \symking{}xg6 & \symking{}a3\\
\textbf{2.} & \symking{}f5 & g4+\\
\textbf{3.} & \symking{}e4 & g5\\
\textbf{4.} & \symking{}d3 & g6\\
\textbf{5.} & \symking{}c2 & g7\\
\textbf{6.} & \symking{}b1 & g8=\symqueen{}\\
\textbf{7.} & \symking{}a1 & \symqueen{}a2\#\\
\end{tabular}\par\vspace{3mm}
\qrcode[height=12mm]{https://helpman.komtera.lt/?fen=8/8/6R1/7k/8/6P1/K5P1/8&moves=7}\\{\sffamily\tiny\color{ink}scannen zum Lösen}
\end{minipage}
\vspace{2mm}

{\raggedright\noindent\textcolor{ink}{\sffamily\textsc{Themen}}\enspace \themetag{Umwandlung} \themetag{ein Stein zieht}\par\smallskip{\footnotesize\textit{Hinweis:} die Lösung beginnt mit einem Schlag.}\par}

\clearpage
\phantomsection\label{hm:KRPvkb}
\begin{tcolorbox}[enhanced,colback=ink,colframe=ink,arc=0pt,boxrule=0pt,left=4mm,right=4mm,top=2.5mm,bottom=2.5mm]
\color{white}\sffamily{\LARGE\bfseries KRPvkb}\hfill{\Large h\#7}\\[0.5mm]
{\small 5 Steine\quad Schwarz am Zug}\hfill{\small Nr.~138\qquad\href{https://helpman.komtera.lt/?fen=6Rk/7b/8/8/8/8/1P6/2K5&moves=7}{\color{white}\bfseries lösen\,\ensuremath{\nearrow}}}
\end{tcolorbox}
\vspace{2mm}
\noindent\begin{minipage}[t]{36mm}\vspace{0pt}\sffamily\small\raggedright
\hmfigure{längste eindeutige}{{\Large h\#7}}\hmfigure{eindeutig bei dieser Länge}{{\Large 7} Stellungen}\hmfigure{längstes Matt}{{\Large h\#8}}\hmfigure{in der Klasse}{710 Millionen Stellungen,\\ 506 Millionen mit Hilfsmatt}
\end{minipage}\hfill
\begin{minipage}[t]{80mm}\vspace{0pt}\begin{tikzpicture}\begin{axis}[width=72mm,height=42mm,ybar,bar width=1.6pt,xmin=-1,xmax=17,ymin=0,ymax=116.6,axis y line*=left,axis x line*=bottom,xlabel={\scriptsize Halbzüge bis Matt},ylabel={\scriptsize Stellungen (Mio.)},tick label style={font=\tiny},label style={font=\scriptsize},enlarge x limits=false,ylabel shift=-3pt,xlabel shift=-3pt]\addplot[fill=ink!55,draw=none] coordinates {(0,0.1597) (1,1.269) (2,3.825) (3,20.25) (4,44.85) (5,82.5) (6,106) (7,85.34) (8,84.59) (9,43.93) (10,26.09) (11,4.992) (12,2.154) (13,0.01966) (14,0.01177) (16,0.000616)};\draw[densely dotted] (axis cs:14,0) -- (axis cs:14,116.6);\end{axis}\begin{axis}[width=72mm,height=42mm,axis y line*=right,axis x line=none,ymode=log,xmin=-1,xmax=17,ylabel={\scriptsize eindeutig (log)},tick label style={font=\tiny},label style={font=\scriptsize},ylabel shift=-3pt]\addplot[mark=*,mark size=0.7pt,draw=black,thin] coordinates {(0,159681) (1,1002276) (2,1453373) (3,4585161) (4,7406220) (5,4692927) (6,2744132) (7,766617) (8,396807) (9,69131) (10,21763) (11,178) (12,22) (14,7)};\end{axis}\end{tikzpicture}\end{minipage}
\vspace{3mm}

\noindent\begin{minipage}[t]{60mm}\vspace{0pt}\sbox\hmboard{\chessboard[setfen=6Rk/7b/8/8/8/8/1P6/2K5 b - - 0 1]}\usebox\hmboard\\[1mm]
\makebox[\wd\hmboard]{\sffamily\hspace{3mm}\textbf{h\#7}\hfill(3+2)}\end{minipage}\hfill
\begin{minipage}[t]{52mm}\vspace{0pt}
\textcolor{ink}{\sffamily\textsc{Lösung}}\par\smallskip\small
\begin{tabular}{@{}r@{\enspace}l@{\enspace}l@{}}
\textbf{1.} & \symking{}xg8 & b3\\
\textbf{2.} & \symking{}f7 & b4\\
\textbf{3.} & \symking{}e6 & b5\\
\textbf{4.} & \symking{}d5 & b6\\
\textbf{5.} & \symking{}c4 & b7\\
\textbf{6.} & \symking{}b3 & b8=\symqueen{}+\\
\textbf{7.} & \symking{}a2 & \symqueen{}b2\#\\
\end{tabular}\par\vspace{3mm}
\qrcode[height=12mm]{https://helpman.komtera.lt/?fen=6Rk/7b/8/8/8/8/1P6/2K5&moves=7}\\{\sffamily\tiny\color{ink}scannen zum Lösen}
\end{minipage}
\vspace{2mm}

{\raggedright\noindent\textcolor{ink}{\sffamily\textsc{Themen}}\enspace \themetag{Umwandlung} \themetag{Excelsior} \themetag{Rundlauf} \themetag{ein Stein zieht}\par\smallskip{\footnotesize\textit{Hinweis:} die Partei am Zug steht im Diagramm im Schach; die Lösung beginnt mit einem Schlag.}\par}

\clearpage
\phantomsection\label{hm:KRPvkn}
\begin{tcolorbox}[enhanced,colback=ink,colframe=ink,arc=0pt,boxrule=0pt,left=4mm,right=4mm,top=2.5mm,bottom=2.5mm]
\color{white}\sffamily{\LARGE\bfseries KRPvkn}\hfill{\Large h\#7}\\[0.5mm]
{\small 5 Steine\quad Schwarz am Zug}\hfill{\small Nr.~139\qquad\href{https://helpman.komtera.lt/?fen=6Rk/7n/8/8/8/8/4P3/2K5&moves=7}{\color{white}\bfseries lösen\,\ensuremath{\nearrow}}}
\end{tcolorbox}
\vspace{2mm}
\noindent\begin{minipage}[t]{36mm}\vspace{0pt}\sffamily\small\raggedright
\hmfigure{längste eindeutige}{{\Large h\#7}}\hmfigure{eindeutig bei dieser Länge}{{\Large 14} Stellungen}\hmfigure{längstes Matt}{{\Large h\#8.5}}\hmfigure{in der Klasse}{710 Millionen Stellungen,\\ 521 Millionen mit Hilfsmatt}
\end{minipage}\hfill
\begin{minipage}[t]{80mm}\vspace{0pt}\begin{tikzpicture}\begin{axis}[width=72mm,height=42mm,ybar,bar width=1.6pt,xmin=-1,xmax=18,ymin=0,ymax=121.4,axis y line*=left,axis x line*=bottom,xlabel={\scriptsize Halbzüge bis Matt},ylabel={\scriptsize Stellungen (Mio.)},tick label style={font=\tiny},label style={font=\scriptsize},enlarge x limits=false,ylabel shift=-3pt,xlabel shift=-3pt]\addplot[fill=ink!55,draw=none] coordinates {(0,0.2169) (1,1.708) (2,4.429) (3,21.64) (4,47.13) (5,81.61) (6,110.3) (7,84.34) (8,89.37) (9,43.87) (10,28.32) (11,5.117) (12,2.419) (13,0.0251) (14,0.01543) (15,1e-06) (16,0.001118) (17,1e-06)};\draw[densely dotted] (axis cs:14,0) -- (axis cs:14,121.4);\end{axis}\begin{axis}[width=72mm,height=42mm,axis y line*=right,axis x line=none,ymode=log,xmin=-1,xmax=18,ylabel={\scriptsize eindeutig (log)},tick label style={font=\tiny},label style={font=\scriptsize},ylabel shift=-3pt]\addplot[mark=*,mark size=0.7pt,draw=black,thin] coordinates {(0,216929) (1,1343421) (2,1809476) (3,5167705) (4,8574622) (5,4685247) (6,2903801) (7,657048) (8,359871) (9,56507) (10,19764) (11,39) (12,15) (14,14)};\end{axis}\end{tikzpicture}\end{minipage}
\vspace{3mm}

\noindent\begin{minipage}[t]{60mm}\vspace{0pt}\sbox\hmboard{\chessboard[setfen=6Rk/7n/8/8/8/8/4P3/2K5 b - - 0 1]}\usebox\hmboard\\[1mm]
\makebox[\wd\hmboard]{\sffamily\hspace{3mm}\textbf{h\#7}\hfill(3+2)}\end{minipage}\hfill
\begin{minipage}[t]{52mm}\vspace{0pt}
\textcolor{ink}{\sffamily\textsc{Lösung}}\par\smallskip\small
\begin{tabular}{@{}r@{\enspace}l@{\enspace}l@{}}
\textbf{1.} & \symking{}xg8 & e3\\
\textbf{2.} & \symking{}f7 & e4\\
\textbf{3.} & \symking{}e6 & e5\\
\textbf{4.} & \symking{}d5 & e6\\
\textbf{5.} & \symking{}c4 & e7\\
\textbf{6.} & \symking{}b3 & e8=\symqueen{}\\
\textbf{7.} & \symking{}a2 & \symqueen{}a4\#\\
\end{tabular}\par\vspace{3mm}
\qrcode[height=12mm]{https://helpman.komtera.lt/?fen=6Rk/7n/8/8/8/8/4P3/2K5&moves=7}\\{\sffamily\tiny\color{ink}scannen zum Lösen}
\end{minipage}
\vspace{2mm}

{\raggedright\noindent\textcolor{ink}{\sffamily\textsc{Themen}}\enspace \themetag{reines Matt} \themetag{Mustermatt} \themetag{Spiegelmatt} \themetag{Umwandlung} \themetag{Excelsior} \themetag{ein Stein zieht} \themetag{ohne Schach} \themetag{Umnov}\par\smallskip{\footnotesize\textit{Hinweis:} die Partei am Zug steht im Diagramm im Schach; die Lösung beginnt mit einem Schlag.}\par}

\clearpage
\phantomsection\label{hm:KRPvkp}
\begin{tcolorbox}[enhanced,colback=ink,colframe=ink,arc=0pt,boxrule=0pt,left=4mm,right=4mm,top=2.5mm,bottom=2.5mm]
\color{white}\sffamily{\LARGE\bfseries KRPvkp}\hfill{\Large h\#7}\\[0.5mm]
{\small 5 Steine\quad Schwarz am Zug}\hfill{\small Nr.~140\qquad\href{https://helpman.komtera.lt/?fen=8/8/8/pK6/P7/8/6R1/7k&moves=7}{\color{white}\bfseries lösen\,\ensuremath{\nearrow}}}
\end{tcolorbox}
\vspace{2mm}
\noindent\begin{minipage}[t]{36mm}\vspace{0pt}\sffamily\small\raggedright
\hmfigure{längste eindeutige}{{\Large h\#7}}\hmfigure{eindeutig bei dieser Länge}{{\Large 41} Stellungen}\hmfigure{längstes Matt}{{\Large h\#14}}\hmfigure{in der Klasse}{533 Millionen Stellungen,\\ 401 Millionen mit Hilfsmatt}
\end{minipage}\hfill
\begin{minipage}[t]{80mm}\vspace{0pt}\begin{tikzpicture}\begin{axis}[width=72mm,height=42mm,ybar,bar width=1.6pt,xmin=-1,xmax=29,ymin=0,ymax=96.08,axis y line*=left,axis x line*=bottom,xlabel={\scriptsize Halbzüge bis Matt},ylabel={\scriptsize Stellungen (Mio.)},tick label style={font=\tiny},label style={font=\scriptsize},enlarge x limits=false,ylabel shift=-3pt,xlabel shift=-3pt]\addplot[fill=ink!55,draw=none] coordinates {(0,0.2171) (1,1.634) (2,3.754) (3,17.7) (4,38.92) (5,62.05) (6,87.35) (7,61.61) (8,69.2) (9,31.35) (10,21.4) (11,3.684) (12,2.167) (13,0.03663) (14,0.03157) (15,7.1e-05) (16,0.002771) (17,1.8e-05) (18,0.003418) (19,3e-06) (20,0.003136) (22,0.002346) (24,0.001142) (26,0.000448) (28,0.00018)};\draw[densely dotted] (axis cs:14,0) -- (axis cs:14,96.08);\end{axis}\begin{axis}[width=72mm,height=42mm,axis y line*=right,axis x line=none,ymode=log,xmin=-1,xmax=29,ylabel={\scriptsize eindeutig (log)},tick label style={font=\tiny},label style={font=\scriptsize},ylabel shift=-3pt]\addplot[mark=*,mark size=0.7pt,draw=black,thin] coordinates {(0,217137) (1,1279027) (2,2096054) (3,4864454) (4,8535722) (5,3912502) (6,2340222) (7,452549) (8,242044) (9,32287) (10,14017) (11,482) (12,586) (14,41)};\end{axis}\end{tikzpicture}\end{minipage}
\vspace{3mm}

\noindent\begin{minipage}[t]{60mm}\vspace{0pt}\sbox\hmboard{\chessboard[setfen=8/8/8/pK6/P7/8/6R1/7k b - - 0 1]}\usebox\hmboard\\[1mm]
\makebox[\wd\hmboard]{\sffamily\hspace{3mm}\textbf{h\#7}\hfill(3+2)}\end{minipage}\hfill
\begin{minipage}[t]{52mm}\vspace{0pt}
\textcolor{ink}{\sffamily\textsc{Lösung}}\par\smallskip\small
\begin{tabular}{@{}r@{\enspace}l@{\enspace}l@{}}
\textbf{1.} & \symking{}xg2 & \symking{}xa5\\
\textbf{2.} & \symking{}f3 & \symking{}b4\\
\textbf{3.} & \symking{}e4 & a5\\
\textbf{4.} & \symking{}d5 & a6\\
\textbf{5.} & \symking{}c6 & a7\\
\textbf{6.} & \symking{}b7 & \symking{}c5\\
\textbf{7.} & \symking{}a6 & a8=\symqueen{}\#\\
\end{tabular}\par\vspace{3mm}
\qrcode[height=12mm]{https://helpman.komtera.lt/?fen=8/8/8/pK6/P7/8/6R1/7k&moves=7}\\{\sffamily\tiny\color{ink}scannen zum Lösen}
\end{minipage}
\vspace{2mm}

{\raggedright\noindent\textcolor{ink}{\sffamily\textsc{Themen}}\enspace \themetag{reines Matt} \themetag{Mustermatt} \themetag{Idealmatt} \themetag{Spiegelmatt} \themetag{Umwandlung} \themetag{ein Stein zieht} \themetag{ohne Schach}\par\smallskip{\footnotesize\textit{Hinweis:} die Lösung beginnt mit einem Schlag.}\par}

\clearpage
\phantomsection\label{hm:KRPvkq}
\begin{tcolorbox}[enhanced,colback=ink,colframe=ink,arc=0pt,boxrule=0pt,left=4mm,right=4mm,top=2.5mm,bottom=2.5mm]
\color{white}\sffamily{\LARGE\bfseries KRPvkq}\hfill{\Large h\#7}\\[0.5mm]
{\small 5 Steine\quad Schwarz am Zug}\hfill{\small Nr.~141\qquad\href{https://helpman.komtera.lt/?fen=6Rk/7q/8/8/8/8/1P6/2K5&moves=7}{\color{white}\bfseries lösen\,\ensuremath{\nearrow}}}
\end{tcolorbox}
\vspace{2mm}
\noindent\begin{minipage}[t]{36mm}\vspace{0pt}\sffamily\small\raggedright
\hmfigure{längste eindeutige}{{\Large h\#7}}\hmfigure{eindeutig bei dieser Länge}{{\Large 7} Stellungen}\hmfigure{längstes Matt}{{\Large h\#9}}\hmfigure{in der Klasse}{710 Millionen Stellungen,\\ 443 Millionen mit Hilfsmatt}
\end{minipage}\hfill
\begin{minipage}[t]{80mm}\vspace{0pt}\begin{tikzpicture}\begin{axis}[width=72mm,height=42mm,ybar,bar width=1.6pt,xmin=-1,xmax=19,ymin=0,ymax=97.79,axis y line*=left,axis x line*=bottom,xlabel={\scriptsize Halbzüge bis Matt},ylabel={\scriptsize Stellungen (Mio.)},tick label style={font=\tiny},label style={font=\scriptsize},enlarge x limits=false,ylabel shift=-3pt,xlabel shift=-3pt]\addplot[fill=ink!55,draw=none] coordinates {(0,0.07029) (1,0.6005) (2,2.612) (3,15.84) (4,33.17) (5,75.94) (6,80.51) (7,88.9) (8,65.67) (9,49.35) (10,21.58) (11,7.334) (12,1.585) (13,0.07895) (14,0.008706) (15,0.000936) (16,0.000631) (17,0.00012) (18,1e-06)};\draw[densely dotted] (axis cs:14,0) -- (axis cs:14,97.79);\end{axis}\begin{axis}[width=72mm,height=42mm,axis y line*=right,axis x line=none,ymode=log,xmin=-1,xmax=19,ylabel={\scriptsize eindeutig (log)},tick label style={font=\tiny},label style={font=\scriptsize},ylabel shift=-3pt]\addplot[mark=*,mark size=0.7pt,draw=black,thin] coordinates {(0,70294) (1,483378) (2,991322) (3,3991046) (4,4293393) (5,4524949) (6,1975615) (7,1059544) (8,399331) (9,116318) (10,15886) (11,423) (12,1) (13,2) (14,7)};\end{axis}\end{tikzpicture}\end{minipage}
\vspace{3mm}

\noindent\begin{minipage}[t]{60mm}\vspace{0pt}\sbox\hmboard{\chessboard[setfen=6Rk/7q/8/8/8/8/1P6/2K5 b - - 0 1]}\usebox\hmboard\\[1mm]
\makebox[\wd\hmboard]{\sffamily\hspace{3mm}\textbf{h\#7}\hfill(3+2)}\end{minipage}\hfill
\begin{minipage}[t]{52mm}\vspace{0pt}
\textcolor{ink}{\sffamily\textsc{Lösung}}\par\smallskip\small
\begin{tabular}{@{}r@{\enspace}l@{\enspace}l@{}}
\textbf{1.} & \symking{}xg8 & b3\\
\textbf{2.} & \symking{}f7 & b4\\
\textbf{3.} & \symking{}e6 & b5\\
\textbf{4.} & \symking{}d5 & b6\\
\textbf{5.} & \symking{}c4 & b7\\
\textbf{6.} & \symking{}b3 & b8=\symqueen{}+\\
\textbf{7.} & \symking{}a2 & \symqueen{}b2\#\\
\end{tabular}\par\vspace{3mm}
\qrcode[height=12mm]{https://helpman.komtera.lt/?fen=6Rk/7q/8/8/8/8/1P6/2K5&moves=7}\\{\sffamily\tiny\color{ink}scannen zum Lösen}
\end{minipage}
\vspace{2mm}

{\raggedright\noindent\textcolor{ink}{\sffamily\textsc{Themen}}\enspace \themetag{Umwandlung} \themetag{Excelsior} \themetag{Rundlauf} \themetag{ein Stein zieht}\par\smallskip{\footnotesize\textit{Hinweis:} die Partei am Zug steht im Diagramm im Schach; die Lösung beginnt mit einem Schlag.}\par}

\clearpage
\phantomsection\label{hm:KRPvkr}
\begin{tcolorbox}[enhanced,colback=ink,colframe=ink,arc=0pt,boxrule=0pt,left=4mm,right=4mm,top=2.5mm,bottom=2.5mm]
\color{white}\sffamily{\LARGE\bfseries KRPvkr}\hfill{\Large h\#7}\\[0.5mm]
{\small 5 Steine\quad Schwarz am Zug}\hfill{\small Nr.~142\qquad\href{https://helpman.komtera.lt/?fen=8/8/8/8/K7/8/kR2P3/r7&moves=7}{\color{white}\bfseries lösen\,\ensuremath{\nearrow}}}
\end{tcolorbox}
\vspace{2mm}
\noindent\begin{minipage}[t]{36mm}\vspace{0pt}\sffamily\small\raggedright
\hmfigure{längste eindeutige}{{\Large h\#7}}\hmfigure{eindeutig bei dieser Länge}{{\Large 2} Stellungen}\hmfigure{längstes Matt}{{\Large h\#8.5}}\hmfigure{in der Klasse}{710 Millionen Stellungen,\\ 483 Millionen mit Hilfsmatt}
\end{minipage}\hfill
\begin{minipage}[t]{80mm}\vspace{0pt}\begin{tikzpicture}\begin{axis}[width=72mm,height=42mm,ybar,bar width=1.6pt,xmin=-1,xmax=18,ymin=0,ymax=109.1,axis y line*=left,axis x line*=bottom,xlabel={\scriptsize Halbzüge bis Matt},ylabel={\scriptsize Stellungen (Mio.)},tick label style={font=\tiny},label style={font=\scriptsize},enlarge x limits=false,ylabel shift=-3pt,xlabel shift=-3pt]\addplot[fill=ink!55,draw=none] coordinates {(0,0.1268) (1,1.06) (2,3.78) (3,19.85) (4,44.79) (5,81.65) (6,99.17) (7,88.1) (8,75.49) (9,42.82) (10,20.53) (11,4.695) (12,1.158) (13,0.01914) (14,0.006082) (15,2.3e-05) (16,0.000171) (17,7e-06)};\draw[densely dotted] (axis cs:14,0) -- (axis cs:14,109.1);\end{axis}\begin{axis}[width=72mm,height=42mm,axis y line*=right,axis x line=none,ymode=log,xmin=-1,xmax=18,ylabel={\scriptsize eindeutig (log)},tick label style={font=\tiny},label style={font=\scriptsize},ylabel shift=-3pt]\addplot[mark=*,mark size=0.7pt,draw=black,thin] coordinates {(0,126791) (1,843391) (2,1549585) (3,4675616) (4,6998890) (5,4855742) (6,2558309) (7,909596) (8,401644) (9,72664) (10,18443) (11,278) (12,7) (14,2)};\end{axis}\end{tikzpicture}\end{minipage}
\vspace{3mm}

\noindent\begin{minipage}[t]{60mm}\vspace{0pt}\sbox\hmboard{\chessboard[setfen=8/8/8/8/K7/8/kR2P3/r7 b - - 0 1]}\usebox\hmboard\\[1mm]
\makebox[\wd\hmboard]{\sffamily\hspace{3mm}\textbf{h\#7}\hfill(3+2)}\end{minipage}\hfill
\begin{minipage}[t]{52mm}\vspace{0pt}
\textcolor{ink}{\sffamily\textsc{Lösung}}\par\smallskip\small
\begin{tabular}{@{}r@{\enspace}l@{\enspace}l@{}}
\textbf{1.} & \symking{}xb2+ & \symking{}b5\\
\textbf{2.} & \symking{}c3 & e4\\
\textbf{3.} & \symking{}d4 & e5\\
\textbf{4.} & \symking{}d5 & e6\\
\textbf{5.} & \symking{}d6 & e7\\
\textbf{6.} & \symking{}d7 & \symking{}b6\\
\textbf{7.} & \symking{}c8 & e8=\symqueen{}\#\\
\end{tabular}\par\vspace{3mm}
\qrcode[height=12mm]{https://helpman.komtera.lt/?fen=8/8/8/8/K7/8/kR2P3/r7&moves=7}\\{\sffamily\tiny\color{ink}scannen zum Lösen}
\end{minipage}
\vspace{2mm}

{\raggedright\noindent\textcolor{ink}{\sffamily\textsc{Themen}}\enspace \themetag{reines Matt} \themetag{Mustermatt} \themetag{Spiegelmatt} \themetag{Umwandlung} \themetag{Excelsior} \themetag{ein Stein zieht}\par\smallskip{\footnotesize\textit{Hinweis:} die Partei am Zug steht im Diagramm im Schach; die Lösung beginnt mit einem Schlag.}\par}

\clearpage
\phantomsection\label{hm:KBBvkq}
\begin{tcolorbox}[enhanced,colback=ink,colframe=ink,arc=0pt,boxrule=0pt,left=4mm,right=4mm,top=2.5mm,bottom=2.5mm]
\color{white}\sffamily{\LARGE\bfseries KBBvkq}\hfill{\Large h\#6.5}\\[0.5mm]
{\small 5 Steine\quad Weiß am Zug}\hfill{\small Nr.~143\qquad\href{https://helpman.komtera.lt/?fen=7k/8/7B/8/8/8/1qB5/K7&moves=6.5}{\color{white}\bfseries lösen\,\ensuremath{\nearrow}}}
\end{tcolorbox}
\vspace{2mm}
\noindent\begin{minipage}[t]{36mm}\vspace{0pt}\sffamily\small\raggedright
\hmfigure{längste eindeutige}{{\Large h\#6.5}}\hmfigure{eindeutig bei dieser Länge}{{\Large 3} Stellungen}\hmfigure{längstes Matt}{{\Large h\#7.5}}\hmfigure{in der Klasse}{242 Millionen Stellungen,\\ 36,2 Millionen mit Hilfsmatt}
\end{minipage}\hfill
\begin{minipage}[t]{80mm}\vspace{0pt}\begin{tikzpicture}\begin{axis}[width=72mm,height=42mm,ybar,bar width=1.6pt,xmin=-1,xmax=16,ymin=0,ymax=8.774,axis y line*=left,axis x line*=bottom,xlabel={\scriptsize Halbzüge bis Matt},ylabel={\scriptsize Stellungen (Mio.)},tick label style={font=\tiny},label style={font=\scriptsize},enlarge x limits=false,ylabel shift=-3pt,xlabel shift=-3pt]\addplot[fill=ink!55,draw=none] coordinates {(0,0.002493) (1,0.01179) (2,0.06276) (3,0.4514) (4,1.404) (5,4.42) (6,6.617) (7,7.976) (8,6.65) (9,5.09) (10,2.276) (11,0.9227) (12,0.2737) (13,0.05736) (14,0.00975) (15,0.000347)};\draw[densely dotted] (axis cs:13,0) -- (axis cs:13,8.774);\end{axis}\begin{axis}[width=72mm,height=42mm,axis y line*=right,axis x line=none,ymode=log,xmin=-1,xmax=16,ylabel={\scriptsize eindeutig (log)},tick label style={font=\tiny},label style={font=\scriptsize},ylabel shift=-3pt]\addplot[mark=*,mark size=0.7pt,draw=black,thin] coordinates {(0,2493) (1,11747) (2,34415) (3,180808) (4,260604) (5,385472) (6,164797) (7,101131) (8,34213) (9,11690) (10,1599) (11,387) (12,11) (13,3)};\end{axis}\end{tikzpicture}\end{minipage}
\vspace{3mm}

\noindent\begin{minipage}[t]{60mm}\vspace{0pt}\sbox\hmboard{\chessboard[setfen=7k/8/7B/8/8/8/1qB5/K7 w - - 0 1]}\usebox\hmboard\\[1mm]
\makebox[\wd\hmboard]{\sffamily\hspace{3mm}\textbf{h\#6.5}\hfill(3+2)}\end{minipage}\hfill
\begin{minipage}[t]{52mm}\vspace{0pt}
\textcolor{ink}{\sffamily\textsc{Lösung}}\par\smallskip\small
\begin{tabular}{@{}r@{\enspace}l@{\enspace}l@{}}
\textbf{1.} & \ldots & \symking{}xb2\\
\textbf{2.} & \symking{}g8 & \symking{}c3\\
\textbf{3.} & \symking{}f7 & \symking{}d3\\
\textbf{4.} & \symking{}g6 & \symking{}e4\\
\textbf{5.} & \symking{}h5 & \symking{}f5\\
\textbf{6.} & \symking{}h4 & \symbishop{}g5+\\
\textbf{7.} & \symking{}h5 & \symbishop{}d1\#\\
\end{tabular}\par\vspace{3mm}
\qrcode[height=12mm]{https://helpman.komtera.lt/?fen=7k/8/7B/8/8/8/1qB5/K7&moves=6.5}\\{\sffamily\tiny\color{ink}scannen zum Lösen}
\end{minipage}
\vspace{2mm}

{\raggedright\noindent\textcolor{ink}{\sffamily\textsc{Themen}}\enspace \themetag{Switchback} \themetag{ein Stein zieht}\par\smallskip{\footnotesize\textit{Hinweis:} die Partei am Zug steht im Diagramm im Schach; die Lösung beginnt mit einem Schlag.}\par}

\clearpage
\phantomsection\label{hm:KBNNvk}
\begin{tcolorbox}[enhanced,colback=ink,colframe=ink,arc=0pt,boxrule=0pt,left=4mm,right=4mm,top=2.5mm,bottom=2.5mm]
\color{white}\sffamily{\LARGE\bfseries KBNNvk}\hfill{\Large h\#6.5}\\[0.5mm]
{\small 5 Steine\quad Weiß am Zug}\hfill{\small Nr.~144\qquad\href{https://helpman.komtera.lt/?fen=7k/8/8/8/8/7N/2B3N1/K7&moves=6.5}{\color{white}\bfseries lösen\,\ensuremath{\nearrow}}}
\end{tcolorbox}
\vspace{2mm}
\noindent\begin{minipage}[t]{36mm}\vspace{0pt}\sffamily\small\raggedright
\hmfigure{längste eindeutige}{{\Large h\#6.5}}\hmfigure{eindeutig bei dieser Länge}{{\Large 8} Stellungen}\hmfigure{längstes Matt}{{\Large h\#8}}\hmfigure{in der Klasse}{242 Millionen Stellungen,\\ 88,7 Millionen mit Hilfsmatt}
\end{minipage}\hfill
\begin{minipage}[t]{80mm}\vspace{0pt}\begin{tikzpicture}\begin{axis}[width=72mm,height=42mm,ybar,bar width=1.6pt,xmin=-1,xmax=17,ymin=0,ymax=27.55,axis y line*=left,axis x line*=bottom,xlabel={\scriptsize Halbzüge bis Matt},ylabel={\scriptsize Stellungen (Mio.)},tick label style={font=\tiny},label style={font=\scriptsize},enlarge x limits=false,ylabel shift=-3pt,xlabel shift=-3pt]\addplot[fill=ink!55,draw=none] coordinates {(0,0.006843) (1,0.0265) (2,0.06014) (3,0.3064) (4,1.214) (5,4.759) (6,10.94) (7,15.57) (8,25.04) (9,13.46) (10,12.1) (11,3.101) (12,1.666) (13,0.3659) (14,0.1173) (15,0.01368) (16,4.5e-05)};\draw[densely dotted] (axis cs:13,0) -- (axis cs:13,27.55);\end{axis}\begin{axis}[width=72mm,height=42mm,axis y line*=right,axis x line=none,ymode=log,xmin=-1,xmax=17,ylabel={\scriptsize eindeutig (log)},tick label style={font=\tiny},label style={font=\scriptsize},ylabel shift=-3pt]\addplot[mark=*,mark size=0.7pt,draw=black,thin] coordinates {(0,6843) (1,24100) (2,54870) (3,226752) (4,666666) (5,854662) (6,1018786) (7,239208) (8,184711) (9,16711) (10,5717) (11,526) (12,115) (13,8)};\end{axis}\end{tikzpicture}\end{minipage}
\vspace{3mm}

\noindent\begin{minipage}[t]{60mm}\vspace{0pt}\sbox\hmboard{\chessboard[setfen=7k/8/8/8/8/7N/2B3N1/K7 w - - 0 1]}\usebox\hmboard\\[1mm]
\makebox[\wd\hmboard]{\sffamily\hspace{3mm}\textbf{h\#6.5}\hfill(4+1)}\end{minipage}\hfill
\begin{minipage}[t]{52mm}\vspace{0pt}
\textcolor{ink}{\sffamily\textsc{Lösung}}\par\smallskip\small
\begin{tabular}{@{}r@{\enspace}l@{\enspace}l@{}}
\textbf{1.} & \ldots & \symking{}b2\\
\textbf{2.} & \symking{}g8 & \symking{}c3\\
\textbf{3.} & \symking{}g7 & \symking{}d4\\
\textbf{4.} & \symking{}h6 & \symking{}e5\\
\textbf{5.} & \symking{}h5 & \symking{}f6\\
\textbf{6.} & \symking{}g4 & \symking{}g7\\
\textbf{7.} & \symking{}h5 & \symbishop{}d1\#\\
\end{tabular}\par\vspace{3mm}
\qrcode[height=12mm]{https://helpman.komtera.lt/?fen=7k/8/8/8/8/7N/2B3N1/K7&moves=6.5}\\{\sffamily\tiny\color{ink}scannen zum Lösen}
\end{minipage}
\vspace{2mm}

{\raggedright\noindent\textcolor{ink}{\sffamily\textsc{Themen}}\enspace \themetag{reines Matt} \themetag{Mustermatt} \themetag{Idealmatt} \themetag{Spiegelmatt} \themetag{Switchback} \themetag{ein Stein zieht} \themetag{ohne Schlag} \themetag{ohne Schach}\par}

\begin{tcolorbox}[enhanced,colback=white,colframe=ink!35,arc=1.5pt,boxrule=0.4pt,left=3mm,right=3mm,top=1.5mm,bottom=1.5mm]
\textcolor{ink}{\sffamily\textsc{veröffentlicht}}\par\smallskip\small \textbf{Sheglow, Wiktor S.}, h\#7 0.1...\quad{\color{ink}\scriptsize verwandte Stellung}\\{\scriptsize Suomen Tehtäväniekat, No. (259), 05-06/1998}
\end{tcolorbox}
\clearpage
\phantomsection\label{hm:KBvkqb}
\begin{tcolorbox}[enhanced,colback=ink,colframe=ink,arc=0pt,boxrule=0pt,left=4mm,right=4mm,top=2.5mm,bottom=2.5mm]
\color{white}\sffamily{\LARGE\bfseries KBvkqb}\hfill{\Large h\#6.5}\\[0.5mm]
{\small 5 Steine\quad Weiß am Zug}\hfill{\small Nr.~145\qquad\href{https://helpman.komtera.lt/?fen=8/B7/1q6/8/2b5/8/8/K6k&moves=6.5}{\color{white}\bfseries lösen\,\ensuremath{\nearrow}}}
\end{tcolorbox}
\vspace{2mm}
\noindent\begin{minipage}[t]{36mm}\vspace{0pt}\sffamily\small\raggedright
\hmfigure{längste eindeutige}{{\Large h\#6.5}}\hmfigure{eindeutig bei dieser Länge}{{\Large 256} Stellungen}\hmfigure{längstes Matt}{{\Large h\#9.5}}\hmfigure{in der Klasse}{242 Millionen Stellungen,\\ 74,5 Millionen mit Hilfsmatt}
\end{minipage}\hfill
\begin{minipage}[t]{80mm}\vspace{0pt}\begin{tikzpicture}\begin{axis}[width=72mm,height=42mm,ybar,bar width=1.6pt,xmin=-1,xmax=20,ymin=0,ymax=20.78,axis y line*=left,axis x line*=bottom,xlabel={\scriptsize Halbzüge bis Matt},ylabel={\scriptsize Stellungen (Mio.)},tick label style={font=\tiny},label style={font=\scriptsize},enlarge x limits=false,ylabel shift=-3pt,xlabel shift=-3pt]\addplot[fill=ink!55,draw=none] coordinates {(0,0.000167) (1,0.000718) (2,0.006768) (3,0.03495) (4,0.1636) (5,0.612) (6,1.605) (7,4.883) (8,7.555) (9,16.77) (10,13.07) (11,18.89) (12,7.108) (13,3.358) (14,0.2048) (15,0.1382) (17,0.04776) (19,0.0142)};\draw[densely dotted] (axis cs:13,0) -- (axis cs:13,20.78);\end{axis}\begin{axis}[width=72mm,height=42mm,axis y line*=right,axis x line=none,ymode=log,xmin=-1,xmax=20,ylabel={\scriptsize eindeutig (log)},tick label style={font=\tiny},label style={font=\scriptsize},ylabel shift=-3pt]\addplot[mark=*,mark size=0.7pt,draw=black,thin] coordinates {(0,167) (1,718) (2,5032) (3,19286) (4,47520) (5,110662) (6,109828) (7,207978) (8,141415) (9,166108) (10,40339) (11,25915) (12,1108) (13,256)};\end{axis}\end{tikzpicture}\end{minipage}
\vspace{3mm}

\noindent\begin{minipage}[t]{60mm}\vspace{0pt}\sbox\hmboard{\chessboard[setfen=8/B7/1q6/8/2b5/8/8/K6k w - - 0 1]}\usebox\hmboard\\[1mm]
\makebox[\wd\hmboard]{\sffamily\hspace{3mm}\textbf{h\#6.5}\hfill(2+3)}\end{minipage}\hfill
\begin{minipage}[t]{52mm}\vspace{0pt}
\textcolor{ink}{\sffamily\textsc{Lösung}}\par\smallskip\small
\begin{tabular}{@{}r@{\enspace}l@{\enspace}l@{}}
\textbf{1.} & \ldots & \symbishop{}b8\\
\textbf{2.} & \symqueen{}g1+ & \symking{}b2\\
\textbf{3.} & \symking{}g2 & \symking{}c3\\
\textbf{4.} & \symking{}f2 & \symking{}d4\\
\textbf{5.} & \symqueen{}d1+ & \symking{}e4\\
\textbf{6.} & \symking{}e1 & \symking{}e3\\
\textbf{7.} & \symbishop{}f1 & \symbishop{}g3\#\\
\end{tabular}\par\vspace{3mm}
\qrcode[height=12mm]{https://helpman.komtera.lt/?fen=8/B7/1q6/8/2b5/8/8/K6k&moves=6.5}\\{\sffamily\tiny\color{ink}scannen zum Lösen}
\end{minipage}
\vspace{2mm}

{\raggedright\noindent\textcolor{ink}{\sffamily\textsc{Themen}}\enspace \themetag{Satzspiel} \themetag{Selbstblock} \themetag{ohne Schlag}\par}

\begin{tcolorbox}[enhanced,colback=white,colframe=ink!35,arc=1.5pt,boxrule=0.4pt,left=3mm,right=3mm,top=1.5mm,bottom=1.5mm]
\textcolor{ink}{\sffamily\textsc{veröffentlicht}}\par\smallskip\small \textbf{Sheglow, Wiktor S.}, h\#7 0.1...\quad{\color{ink}\scriptsize diese Stellung}\\{\scriptsize Suomen Tehtäväniekat, No. (304), 05-06/1998}
\end{tcolorbox}
\clearpage
\phantomsection\label{hm:KBvkrb}
\begin{tcolorbox}[enhanced,colback=ink,colframe=ink,arc=0pt,boxrule=0pt,left=4mm,right=4mm,top=2.5mm,bottom=2.5mm]
\color{white}\sffamily{\LARGE\bfseries KBvkrb}\hfill{\Large h\#6.5}\\[0.5mm]
{\small 5 Steine\quad Weiß am Zug}\hfill{\small Nr.~146\qquad\href{https://helpman.komtera.lt/?fen=8/8/b7/8/4k3/8/4rB2/K7&moves=6.5}{\color{white}\bfseries lösen\,\ensuremath{\nearrow}}}
\end{tcolorbox}
\vspace{2mm}
\noindent\begin{minipage}[t]{36mm}\vspace{0pt}\sffamily\small\raggedright
\hmfigure{längste eindeutige}{{\Large h\#6.5}}\hmfigure{eindeutig bei dieser Länge}{{\Large 128} Stellungen}\hmfigure{längstes Matt}{{\Large h\#9.5}}\hmfigure{in der Klasse}{242 Millionen Stellungen,\\ 80,8 Millionen mit Hilfsmatt}
\end{minipage}\hfill
\begin{minipage}[t]{80mm}\vspace{0pt}\begin{tikzpicture}\begin{axis}[width=72mm,height=42mm,ybar,bar width=1.6pt,xmin=-1,xmax=20,ymin=0,ymax=22.12,axis y line*=left,axis x line*=bottom,xlabel={\scriptsize Halbzüge bis Matt},ylabel={\scriptsize Stellungen (Mio.)},tick label style={font=\tiny},label style={font=\scriptsize},enlarge x limits=false,ylabel shift=-3pt,xlabel shift=-3pt]\addplot[fill=ink!55,draw=none] coordinates {(0,0.000253) (1,0.00108) (2,0.009023) (3,0.04152) (4,0.2028) (5,0.6413) (6,1.809) (7,4.546) (8,8.438) (9,16.05) (10,15.95) (11,20.11) (12,9.156) (13,3.498) (14,0.306) (15,0.02831) (16,2.2e-05) (17,0.001188) (19,0.000209)};\draw[densely dotted] (axis cs:13,0) -- (axis cs:13,22.12);\end{axis}\begin{axis}[width=72mm,height=42mm,axis y line*=right,axis x line=none,ymode=log,xmin=-1,xmax=20,ylabel={\scriptsize eindeutig (log)},tick label style={font=\tiny},label style={font=\scriptsize},ylabel shift=-3pt]\addplot[mark=*,mark size=0.7pt,draw=black,thin] coordinates {(0,253) (1,1080) (2,7099) (3,20630) (4,59767) (5,107688) (6,130963) (7,173868) (8,219904) (9,174809) (10,65753) (11,30285) (12,2422) (13,128)};\end{axis}\end{tikzpicture}\end{minipage}
\vspace{3mm}

\noindent\begin{minipage}[t]{60mm}\vspace{0pt}\sbox\hmboard{\chessboard[setfen=8/8/b7/8/4k3/8/4rB2/K7 w - - 0 1]}\usebox\hmboard\\[1mm]
\makebox[\wd\hmboard]{\sffamily\hspace{3mm}\textbf{h\#6.5}\hfill(2+3)}\end{minipage}\hfill
\begin{minipage}[t]{52mm}\vspace{0pt}
\textcolor{ink}{\sffamily\textsc{Lösung}}\par\smallskip\small
\begin{tabular}{@{}r@{\enspace}l@{\enspace}l@{}}
\textbf{1.} & \ldots & \symbishop{}g1\\
\textbf{2.} & \symrook{}e1+ & \symking{}b2\\
\textbf{3.} & \symking{}f3 & \symking{}c3\\
\textbf{4.} & \symking{}e2 & \symking{}d4\\
\textbf{5.} & \symrook{}d1+ & \symking{}e4\\
\textbf{6.} & \symking{}e1 & \symking{}e3\\
\textbf{7.} & \symbishop{}f1 & \symbishop{}f2\#\\
\end{tabular}\par\vspace{3mm}
\qrcode[height=12mm]{https://helpman.komtera.lt/?fen=8/8/b7/8/4k3/8/4rB2/K7&moves=6.5}\\{\sffamily\tiny\color{ink}scannen zum Lösen}
\end{minipage}
\vspace{2mm}

{\raggedright\noindent\textcolor{ink}{\sffamily\textsc{Themen}}\enspace \themetag{reines Matt} \themetag{Mustermatt} \themetag{Idealmatt} \themetag{Switchback} \themetag{Selbstblock} \themetag{ohne Schlag}\par}

\clearpage
\phantomsection\label{hm:KNNvkb}
\begin{tcolorbox}[enhanced,colback=ink,colframe=ink,arc=0pt,boxrule=0pt,left=4mm,right=4mm,top=2.5mm,bottom=2.5mm]
\color{white}\sffamily{\LARGE\bfseries KNNvkb}\hfill{\Large h\#6.5}\\[0.5mm]
{\small 5 Steine\quad Weiß am Zug}\hfill{\small Nr.~147\qquad\href{https://helpman.komtera.lt/?fen=1k5N/8/8/8/8/8/7N/b2K4&moves=6.5}{\color{white}\bfseries lösen\,\ensuremath{\nearrow}}}
\end{tcolorbox}
\vspace{2mm}
\noindent\begin{minipage}[t]{36mm}\vspace{0pt}\sffamily\small\raggedright
\hmfigure{längste eindeutige}{{\Large h\#6.5}}\hmfigure{eindeutig bei dieser Länge}{{\Large 21} Stellungen}\hmfigure{längstes Matt}{{\Large h\#10}}\hmfigure{in der Klasse}{242 Millionen Stellungen,\\ 87,7 Millionen mit Hilfsmatt}
\end{minipage}\hfill
\begin{minipage}[t]{80mm}\vspace{0pt}\begin{tikzpicture}\begin{axis}[width=72mm,height=42mm,ybar,bar width=1.6pt,xmin=-1,xmax=21,ymin=0,ymax=21.27,axis y line*=left,axis x line*=bottom,xlabel={\scriptsize Halbzüge bis Matt},ylabel={\scriptsize Stellungen (Mio.)},tick label style={font=\tiny},label style={font=\scriptsize},enlarge x limits=false,ylabel shift=-3pt,xlabel shift=-3pt]\addplot[fill=ink!55,draw=none] coordinates {(0,0.00086) (1,0.004295) (2,0.01603) (3,0.1281) (4,0.4953) (5,2.278) (6,5.434) (7,11.98) (8,18.14) (9,19.34) (10,16.56) (11,8.36) (12,3.66) (13,0.9896) (14,0.2223) (15,0.05008) (16,0.00014) (18,9.8e-05) (20,8e-06)};\draw[densely dotted] (axis cs:13,0) -- (axis cs:13,21.27);\end{axis}\begin{axis}[width=72mm,height=42mm,axis y line*=right,axis x line=none,ymode=log,xmin=-1,xmax=21,ylabel={\scriptsize eindeutig (log)},tick label style={font=\tiny},label style={font=\scriptsize},ylabel shift=-3pt]\addplot[mark=*,mark size=0.7pt,draw=black,thin] coordinates {(0,860) (1,4126) (2,11605) (3,66017) (4,195517) (5,270165) (6,306921) (7,194007) (8,163724) (9,41772) (10,13206) (11,910) (12,135) (13,21)};\end{axis}\end{tikzpicture}\end{minipage}
\vspace{3mm}

\noindent\begin{minipage}[t]{60mm}\vspace{0pt}\sbox\hmboard{\chessboard[setfen=1k5N/8/8/8/8/8/7N/b2K4 w - - 0 1]}\usebox\hmboard\\[1mm]
\makebox[\wd\hmboard]{\sffamily\hspace{3mm}\textbf{h\#6.5}\hfill(3+2)}\end{minipage}\hfill
\begin{minipage}[t]{52mm}\vspace{0pt}
\textcolor{ink}{\sffamily\textsc{Lösung}}\par\smallskip\small
\begin{tabular}{@{}r@{\enspace}l@{\enspace}l@{}}
\textbf{1.} & \ldots & \symking{}e2\\
\textbf{2.} & \symking{}c7 & \symking{}f3\\
\textbf{3.} & \symking{}d6 & \symking{}g4\\
\textbf{4.} & \symking{}e5 & \symking{}h5\\
\textbf{5.} & \symking{}f4 & \symknight{}g6+\\
\textbf{6.} & \symking{}g3 & \symknight{}f1+\\
\textbf{7.} & \symking{}h3 & \symknight{}f4\#\\
\end{tabular}\par\vspace{3mm}
\qrcode[height=12mm]{https://helpman.komtera.lt/?fen=1k5N/8/8/8/8/8/7N/b2K4&moves=6.5}\\{\sffamily\tiny\color{ink}scannen zum Lösen}
\end{minipage}
\vspace{2mm}

{\raggedright\noindent\textcolor{ink}{\sffamily\textsc{Themen}}\enspace \themetag{reines Matt} \themetag{Mustermatt} \themetag{Spiegelmatt} \themetag{ein Stein zieht} \themetag{ohne Schlag}\par}

\clearpage
\phantomsection\label{hm:KNNvkn}
\begin{tcolorbox}[enhanced,colback=ink,colframe=ink,arc=0pt,boxrule=0pt,left=4mm,right=4mm,top=2.5mm,bottom=2.5mm]
\color{white}\sffamily{\LARGE\bfseries KNNvkn}\hfill{\Large h\#6.5}\\[0.5mm]
{\small 5 Steine\quad Weiß am Zug}\hfill{\small Nr.~148\qquad\href{https://helpman.komtera.lt/?fen=1k5N/8/8/8/8/8/3K3N/5n2&moves=6.5}{\color{white}\bfseries lösen\,\ensuremath{\nearrow}}}
\end{tcolorbox}
\vspace{2mm}
\noindent\begin{minipage}[t]{36mm}\vspace{0pt}\sffamily\small\raggedright
\hmfigure{längste eindeutige}{{\Large h\#6.5}}\hmfigure{eindeutig bei dieser Länge}{{\Large 2} Stellungen}\hmfigure{längstes Matt}{{\Large h\#7.5}}\hmfigure{in der Klasse}{242 Millionen Stellungen,\\ 90,1 Millionen mit Hilfsmatt}
\end{minipage}\hfill
\begin{minipage}[t]{80mm}\vspace{0pt}\begin{tikzpicture}\begin{axis}[width=72mm,height=42mm,ybar,bar width=1.6pt,xmin=-1,xmax=16,ymin=0,ymax=21.54,axis y line*=left,axis x line*=bottom,xlabel={\scriptsize Halbzüge bis Matt},ylabel={\scriptsize Stellungen (Mio.)},tick label style={font=\tiny},label style={font=\scriptsize},enlarge x limits=false,ylabel shift=-3pt,xlabel shift=-3pt]\addplot[fill=ink!55,draw=none] coordinates {(0,0.000959) (1,0.004953) (2,0.01666) (3,0.1333) (4,0.4852) (5,2.139) (6,5.368) (7,11.17) (8,18.89) (9,19.58) (10,18.7) (11,9.554) (12,3.466) (13,0.5328) (14,0.08256) (15,0.006399)};\draw[densely dotted] (axis cs:13,0) -- (axis cs:13,21.54);\end{axis}\begin{axis}[width=72mm,height=42mm,axis y line*=right,axis x line=none,ymode=log,xmin=-1,xmax=16,ylabel={\scriptsize eindeutig (log)},tick label style={font=\tiny},label style={font=\scriptsize},ylabel shift=-3pt]\addplot[mark=*,mark size=0.7pt,draw=black,thin] coordinates {(0,959) (1,4798) (2,12163) (3,69496) (4,177311) (5,211146) (6,255599) (7,165035) (8,159586) (9,38206) (10,11213) (11,699) (12,60) (13,2)};\end{axis}\end{tikzpicture}\end{minipage}
\vspace{3mm}

\noindent\begin{minipage}[t]{60mm}\vspace{0pt}\sbox\hmboard{\chessboard[setfen=1k5N/8/8/8/8/8/3K3N/5n2 w - - 0 1]}\usebox\hmboard\\[1mm]
\makebox[\wd\hmboard]{\sffamily\hspace{3mm}\textbf{h\#6.5}\hfill(3+2)}\end{minipage}\hfill
\begin{minipage}[t]{52mm}\vspace{0pt}
\textcolor{ink}{\sffamily\textsc{Lösung}}\par\smallskip\small
\begin{tabular}{@{}r@{\enspace}l@{\enspace}l@{}}
\textbf{1.} & \ldots & \symking{}e2\\
\textbf{2.} & \symking{}c7 & \symking{}f3\\
\textbf{3.} & \symking{}d6 & \symking{}g4\\
\textbf{4.} & \symking{}e5 & \symking{}h5\\
\textbf{5.} & \symking{}f4 & \symknight{}g6+\\
\textbf{6.} & \symking{}g3 & \symknight{}xf1+\\
\textbf{7.} & \symking{}h3 & \symknight{}f4\#\\
\end{tabular}\par\vspace{3mm}
\qrcode[height=12mm]{https://helpman.komtera.lt/?fen=1k5N/8/8/8/8/8/3K3N/5n2&moves=6.5}\\{\sffamily\tiny\color{ink}scannen zum Lösen}
\end{minipage}
\vspace{2mm}

{\raggedright\noindent\textcolor{ink}{\sffamily\textsc{Themen}}\enspace \themetag{reines Matt} \themetag{Mustermatt} \themetag{Idealmatt} \themetag{Spiegelmatt} \themetag{ein Stein zieht}\par\smallskip{\footnotesize\textit{Hinweis:} die Partei am Zug steht im Diagramm im Schach.}\par}

\clearpage
\phantomsection\label{hm:KPPvkr}
\begin{tcolorbox}[enhanced,colback=ink,colframe=ink,arc=0pt,boxrule=0pt,left=4mm,right=4mm,top=2.5mm,bottom=2.5mm]
\color{white}\sffamily{\LARGE\bfseries KPPvkr}\hfill{\Large h\#6.5}\\[0.5mm]
{\small 5 Steine\quad Weiß am Zug}\hfill{\small Nr.~149\qquad\href{https://helpman.komtera.lt/?fen=8/8/6r1/6k1/6P1/6P1/8/K7&moves=6.5}{\color{white}\bfseries lösen\,\ensuremath{\nearrow}}}
\end{tcolorbox}
\vspace{2mm}
\noindent\begin{minipage}[t]{36mm}\vspace{0pt}\sffamily\small\raggedright
\hmfigure{längste eindeutige}{{\Large h\#6.5}}\hmfigure{eindeutig bei dieser Länge}{{\Large 784} Stellungen}\hmfigure{längstes Matt}{{\Large h\#8.5}}\hmfigure{in der Klasse}{533 Millionen Stellungen,\\ 200 Millionen mit Hilfsmatt}
\end{minipage}\hfill
\begin{minipage}[t]{80mm}\vspace{0pt}\begin{tikzpicture}\begin{axis}[width=72mm,height=42mm,ybar,bar width=1.6pt,xmin=-1,xmax=18,ymin=0,ymax=42.52,axis y line*=left,axis x line*=bottom,xlabel={\scriptsize Halbzüge bis Matt},ylabel={\scriptsize Stellungen (Mio.)},tick label style={font=\tiny},label style={font=\scriptsize},enlarge x limits=false,ylabel shift=-3pt,xlabel shift=-3pt]\addplot[fill=ink!55,draw=none] coordinates {(0,0.000503) (1,0.05456) (2,0.2059) (3,1.644) (4,4.918) (5,17.48) (6,26.41) (7,38.65) (8,29.57) (9,30.11) (10,18.76) (11,16.57) (12,10.7) (13,3.741) (14,0.9959) (15,0.04578) (16,0.000109) (17,0.000137)};\draw[densely dotted] (axis cs:13,0) -- (axis cs:13,42.52);\end{axis}\begin{axis}[width=72mm,height=42mm,axis y line*=right,axis x line=none,ymode=log,xmin=-1,xmax=18,ylabel={\scriptsize eindeutig (log)},tick label style={font=\tiny},label style={font=\scriptsize},ylabel shift=-3pt]\addplot[mark=*,mark size=0.7pt,draw=black,thin] coordinates {(0,503) (1,15589) (2,60817) (3,540080) (4,1107120) (5,1572484) (6,763911) (7,639394) (8,217943) (9,183216) (10,65342) (11,43893) (12,11132) (13,784)};\end{axis}\end{tikzpicture}\end{minipage}
\vspace{3mm}

\noindent\begin{minipage}[t]{60mm}\vspace{0pt}\sbox\hmboard{\chessboard[setfen=8/8/6r1/6k1/6P1/6P1/8/K7 w - - 0 1]}\usebox\hmboard\\[1mm]
\makebox[\wd\hmboard]{\sffamily\hspace{3mm}\textbf{h\#6.5}\hfill(3+2)}\end{minipage}\hfill
\begin{minipage}[t]{52mm}\vspace{0pt}
\textcolor{ink}{\sffamily\textsc{Lösung}}\par\smallskip\small
\begin{tabular}{@{}r@{\enspace}l@{\enspace}l@{}}
\textbf{1.} & \ldots & \symking{}b2\\
\textbf{2.} & \symrook{}f6 & \symking{}c3\\
\textbf{3.} & \symrook{}f5 & gxf5\\
\textbf{4.} & \symking{}g4 & f6\\
\textbf{5.} & \symking{}f3 & f7\\
\textbf{6.} & \symking{}e2 & f8=\symqueen{}\\
\textbf{7.} & \symking{}d1 & \symqueen{}f1\#\\
\end{tabular}\par\vspace{3mm}
\qrcode[height=12mm]{https://helpman.komtera.lt/?fen=8/8/6r1/6k1/6P1/6P1/8/K7&moves=6.5}\\{\sffamily\tiny\color{ink}scannen zum Lösen}
\end{minipage}
\vspace{2mm}

{\raggedright\noindent\textcolor{ink}{\sffamily\textsc{Themen}}\enspace \themetag{reines Matt} \themetag{Mustermatt} \themetag{Spiegelmatt} \themetag{Umwandlung} \themetag{ohne Schach} \themetag{Umnov}\par}

\clearpage
\phantomsection\label{hm:KRBvkq}
\begin{tcolorbox}[enhanced,colback=ink,colframe=ink,arc=0pt,boxrule=0pt,left=4mm,right=4mm,top=2.5mm,bottom=2.5mm]
\color{white}\sffamily{\LARGE\bfseries KRBvkq}\hfill{\Large h\#6.5}\\[0.5mm]
{\small 5 Steine\quad Weiß am Zug}\hfill{\small Nr.~150\qquad\href{https://helpman.komtera.lt/?fen=7k/3B4/6R1/8/8/8/q7/K7&moves=6.5}{\color{white}\bfseries lösen\,\ensuremath{\nearrow}}}
\end{tcolorbox}
\vspace{2mm}
\noindent\begin{minipage}[t]{36mm}\vspace{0pt}\sffamily\small\raggedright
\hmfigure{längste eindeutige}{{\Large h\#6.5}}\hmfigure{eindeutig bei dieser Länge}{{\Large 3} Stellungen}\hmfigure{längstes Matt}{{\Large h\#7.5}}\hmfigure{in der Klasse}{242 Millionen Stellungen,\\ 139 Millionen mit Hilfsmatt}
\end{minipage}\hfill
\begin{minipage}[t]{80mm}\vspace{0pt}\begin{tikzpicture}\begin{axis}[width=72mm,height=42mm,ybar,bar width=1.6pt,xmin=-1,xmax=16,ymin=0,ymax=50.11,axis y line*=left,axis x line*=bottom,xlabel={\scriptsize Halbzüge bis Matt},ylabel={\scriptsize Stellungen (Mio.)},tick label style={font=\tiny},label style={font=\scriptsize},enlarge x limits=false,ylabel shift=-3pt,xlabel shift=-3pt]\addplot[fill=ink!55,draw=none] coordinates {(0,0.03831) (1,0.1951) (2,0.8876) (3,4.709) (4,14.13) (5,33.44) (6,45.55) (7,30.7) (8,7.458) (9,1.85) (10,0.006034) (11,0.1116) (12,0.000463) (13,0.0135) (14,9e-06) (15,6e-05)};\draw[densely dotted] (axis cs:13,0) -- (axis cs:13,50.11);\end{axis}\begin{axis}[width=72mm,height=42mm,axis y line*=right,axis x line=none,ymode=log,xmin=-1,xmax=16,ylabel={\scriptsize eindeutig (log)},tick label style={font=\tiny},label style={font=\scriptsize},ylabel shift=-3pt]\addplot[mark=*,mark size=0.7pt,draw=black,thin] coordinates {(0,38306) (1,192589) (2,487554) (3,1973315) (4,3924521) (5,2943381) (6,715178) (7,136497) (8,558) (9,1423) (10,5) (11,96) (13,3)};\end{axis}\end{tikzpicture}\end{minipage}
\vspace{3mm}

\noindent\begin{minipage}[t]{60mm}\vspace{0pt}\sbox\hmboard{\chessboard[setfen=7k/3B4/6R1/8/8/8/q7/K7 w - - 0 1]}\usebox\hmboard\\[1mm]
\makebox[\wd\hmboard]{\sffamily\hspace{3mm}\textbf{h\#6.5}\hfill(3+2)}\end{minipage}\hfill
\begin{minipage}[t]{52mm}\vspace{0pt}
\textcolor{ink}{\sffamily\textsc{Lösung}}\par\smallskip\small
\begin{tabular}{@{}r@{\enspace}l@{\enspace}l@{}}
\textbf{1.} & \ldots & \symking{}xa2\\
\textbf{2.} & \symking{}h7 & \symking{}b3\\
\textbf{3.} & \symking{}h8 & \symking{}c4\\
\textbf{4.} & \symking{}h7 & \symking{}d5\\
\textbf{5.} & \symking{}h8 & \symking{}e6\\
\textbf{6.} & \symking{}h7 & \symking{}f7\\
\textbf{7.} & \symking{}h8 & \symrook{}h6\#\\
\end{tabular}\par\vspace{3mm}
\qrcode[height=12mm]{https://helpman.komtera.lt/?fen=7k/3B4/6R1/8/8/8/q7/K7&moves=6.5}\\{\sffamily\tiny\color{ink}scannen zum Lösen}
\end{minipage}
\vspace{2mm}

{\raggedright\noindent\textcolor{ink}{\sffamily\textsc{Themen}}\enspace \themetag{reines Matt} \themetag{Spiegelmatt} \themetag{Switchback} \themetag{ein Stein zieht} \themetag{Pendel} \themetag{ohne Schach}\par\smallskip{\footnotesize\textit{Hinweis:} die Partei am Zug steht im Diagramm im Schach; die Lösung beginnt mit einem Schlag.}\par}

\clearpage
\phantomsection\label{hm:KRvkqb}
\begin{tcolorbox}[enhanced,colback=ink,colframe=ink,arc=0pt,boxrule=0pt,left=4mm,right=4mm,top=2.5mm,bottom=2.5mm]
\color{white}\sffamily{\LARGE\bfseries KRvkqb}\hfill{\Large h\#6.5}\\[0.5mm]
{\small 5 Steine\quad Weiß am Zug}\hfill{\small Nr.~151\qquad\href{https://helpman.komtera.lt/?fen=8/8/5k2/3R4/8/bq6/8/1K6&moves=6.5}{\color{white}\bfseries lösen\,\ensuremath{\nearrow}}}
\end{tcolorbox}
\vspace{2mm}
\noindent\begin{minipage}[t]{36mm}\vspace{0pt}\sffamily\small\raggedright
\hmfigure{längste eindeutige}{{\Large h\#6.5}}\hmfigure{eindeutig bei dieser Länge}{{\Large 1} Stellung}\hmfigure{längstes Matt}{{\Large h\#7.5}}\hmfigure{in der Klasse}{242 Millionen Stellungen,\\ 140 Millionen mit Hilfsmatt}
\end{minipage}\hfill
\begin{minipage}[t]{80mm}\vspace{0pt}\begin{tikzpicture}\begin{axis}[width=72mm,height=42mm,ybar,bar width=1.6pt,xmin=-1,xmax=16,ymin=0,ymax=31.78,axis y line*=left,axis x line*=bottom,xlabel={\scriptsize Halbzüge bis Matt},ylabel={\scriptsize Stellungen (Mio.)},tick label style={font=\tiny},label style={font=\scriptsize},enlarge x limits=false,ylabel shift=-3pt,xlabel shift=-3pt]\addplot[fill=ink!55,draw=none] coordinates {(0,0.0116) (1,0.08649) (2,0.4106) (3,2.533) (4,4.596) (5,14.48) (6,16.05) (7,28.89) (8,23.15) (9,27.42) (10,12.93) (11,7.651) (12,1.617) (13,0.2056) (14,0.01676) (15,0.000419)};\draw[densely dotted] (axis cs:13,0) -- (axis cs:13,31.78);\end{axis}\begin{axis}[width=72mm,height=42mm,axis y line*=right,axis x line=none,ymode=log,xmin=-1,xmax=16,ylabel={\scriptsize eindeutig (log)},tick label style={font=\tiny},label style={font=\scriptsize},ylabel shift=-3pt]\addplot[mark=*,mark size=0.7pt,draw=black,thin] coordinates {(0,11600) (1,86492) (2,177899) (3,800570) (4,759575) (5,1229386) (6,586432) (7,491367) (8,211068) (9,118313) (10,21087) (11,2083) (12,16) (13,1)};\end{axis}\end{tikzpicture}\end{minipage}
\vspace{3mm}

\noindent\begin{minipage}[t]{60mm}\vspace{0pt}\sbox\hmboard{\chessboard[setfen=8/8/5k2/3R4/8/bq6/8/1K6 w - - 0 1]}\usebox\hmboard\\[1mm]
\makebox[\wd\hmboard]{\sffamily\hspace{3mm}\textbf{h\#6.5}\hfill(2+3)}\end{minipage}\hfill
\begin{minipage}[t]{52mm}\vspace{0pt}
\textcolor{ink}{\sffamily\textsc{Lösung}}\par\smallskip\small
\begin{tabular}{@{}r@{\enspace}l@{\enspace}l@{}}
\textbf{1.} & \ldots & \symking{}a1\\
\textbf{2.} & \symbishop{}b2+ & \symking{}b1\\
\textbf{3.} & \symbishop{}e5+ & \symking{}c1\\
\textbf{4.} & \symking{}g5 & \symking{}d2\\
\textbf{5.} & \symking{}h4 & \symking{}e2\\
\textbf{6.} & \symbishop{}c3 & \symking{}f3\\
\textbf{7.} & \symking{}h3 & \symrook{}h5\#\\
\end{tabular}\par\vspace{3mm}
\qrcode[height=12mm]{https://helpman.komtera.lt/?fen=8/8/5k2/3R4/8/bq6/8/1K6&moves=6.5}\\{\sffamily\tiny\color{ink}scannen zum Lösen}
\end{minipage}
\vspace{2mm}

{\raggedright\noindent\textcolor{ink}{\sffamily\textsc{Themen}}\enspace \themetag{reines Matt} \themetag{Mustermatt} \themetag{Spiegelmatt} \themetag{Switchback} \themetag{ohne Schlag}\par\smallskip{\footnotesize\textit{Hinweis:} die Partei am Zug steht im Diagramm im Schach.}\par}

\begin{tcolorbox}[enhanced,colback=white,colframe=ink!35,arc=1.5pt,boxrule=0.4pt,left=3mm,right=3mm,top=1.5mm,bottom=1.5mm]
\textcolor{ink}{\sffamily\textsc{veröffentlicht}}\par\smallskip\small \textbf{Sheglow, Wiktor S.}, h\#7 0.1...\quad{\color{ink}\scriptsize verwandte Stellung}\\{\scriptsize Suomen Tehtäväniekat, No. (312), 05-06/1998}
\end{tcolorbox}
\clearpage
\phantomsection\label{hm:KRvkqp}
\begin{tcolorbox}[enhanced,colback=ink,colframe=ink,arc=0pt,boxrule=0pt,left=4mm,right=4mm,top=2.5mm,bottom=2.5mm]
\color{white}\sffamily{\LARGE\bfseries KRvkqp}\hfill{\Large h\#6.5}\\[0.5mm]
{\small 5 Steine\quad Weiß am Zug}\hfill{\small Nr.~152\qquad\href{https://helpman.komtera.lt/?fen=6k1/5R2/8/8/8/1p6/2q5/1K6&moves=6.5}{\color{white}\bfseries lösen\,\ensuremath{\nearrow}}}
\end{tcolorbox}
\vspace{2mm}
\noindent\begin{minipage}[t]{36mm}\vspace{0pt}\sffamily\small\raggedright
\hmfigure{längste eindeutige}{{\Large h\#6.5}}\hmfigure{eindeutig bei dieser Länge}{{\Large 17} Stellungen}\hmfigure{längstes Matt}{{\Large h\#7.5}}\hmfigure{in der Klasse}{710 Millionen Stellungen,\\ 444 Millionen mit Hilfsmatt}
\end{minipage}\hfill
\begin{minipage}[t]{80mm}\vspace{0pt}\begin{tikzpicture}\begin{axis}[width=72mm,height=42mm,ybar,bar width=1.6pt,xmin=-1,xmax=16,ymin=0,ymax=96.11,axis y line*=left,axis x line*=bottom,xlabel={\scriptsize Halbzüge bis Matt},ylabel={\scriptsize Stellungen (Mio.)},tick label style={font=\tiny},label style={font=\scriptsize},enlarge x limits=false,ylabel shift=-3pt,xlabel shift=-3pt]\addplot[fill=ink!55,draw=none] coordinates {(0,0.0619) (1,0.4371) (2,1.901) (3,10.08) (4,17.95) (5,48.09) (6,57.57) (7,87.37) (8,77.03) (9,78.66) (10,39.58) (11,19.58) (12,5.01) (13,0.4841) (14,0.04403) (15,0.000202)};\draw[densely dotted] (axis cs:13,0) -- (axis cs:13,96.11);\end{axis}\begin{axis}[width=72mm,height=42mm,axis y line*=right,axis x line=none,ymode=log,xmin=-1,xmax=16,ylabel={\scriptsize eindeutig (log)},tick label style={font=\tiny},label style={font=\scriptsize},ylabel shift=-3pt]\addplot[mark=*,mark size=0.7pt,draw=black,thin] coordinates {(0,61900) (1,437139) (2,870061) (3,3384913) (4,3269974) (5,4203062) (6,2104739) (7,1319751) (8,645240) (9,258657) (10,49908) (11,2560) (12,147) (13,17)};\end{axis}\end{tikzpicture}\end{minipage}
\vspace{3mm}

\noindent\begin{minipage}[t]{60mm}\vspace{0pt}\sbox\hmboard{\chessboard[setfen=6k1/5R2/8/8/8/1p6/2q5/1K6 w - - 0 1]}\usebox\hmboard\\[1mm]
\makebox[\wd\hmboard]{\sffamily\hspace{3mm}\textbf{h\#6.5}\hfill(2+3)}\end{minipage}\hfill
\begin{minipage}[t]{52mm}\vspace{0pt}
\textcolor{ink}{\sffamily\textsc{Lösung}}\par\smallskip\small
\begin{tabular}{@{}r@{\enspace}l@{\enspace}l@{}}
\textbf{1.} & \ldots & \symking{}a1\\
\textbf{2.} & \symqueen{}b2+ & \symking{}xb2\\
\textbf{3.} & \symking{}h8 & \symking{}c3\\
\textbf{4.} & b2 & \symking{}d4\\
\textbf{5.} & b1=\symrook{} & \symking{}e5\\
\textbf{6.} & \symrook{}h1 & \symking{}f6\\
\textbf{7.} & \symrook{}h7 & \symrook{}f8\#\\
\end{tabular}\par\vspace{3mm}
\qrcode[height=12mm]{https://helpman.komtera.lt/?fen=6k1/5R2/8/8/8/1p6/2q5/1K6&moves=6.5}\\{\sffamily\tiny\color{ink}scannen zum Lösen}
\end{minipage}
\vspace{2mm}

{\raggedright\noindent\textcolor{ink}{\sffamily\textsc{Themen}}\enspace \themetag{reines Matt} \themetag{Mustermatt} \themetag{Idealmatt} \themetag{Umwandlung} \themetag{Unterverwandlung} \themetag{Selbstblock} \themetag{Umnov}\par\smallskip{\footnotesize\textit{Hinweis:} die Partei am Zug steht im Diagramm im Schach.}\par}

\clearpage
\phantomsection\label{hm:KRvkqq}
\begin{tcolorbox}[enhanced,colback=ink,colframe=ink,arc=0pt,boxrule=0pt,left=4mm,right=4mm,top=2.5mm,bottom=2.5mm]
\color{white}\sffamily{\LARGE\bfseries KRvkqq}\hfill{\Large h\#6.5}\\[0.5mm]
{\small 5 Steine\quad Weiß am Zug}\hfill{\small Nr.~153\qquad\href{https://helpman.komtera.lt/?fen=8/8/5k2/8/8/6R1/2q3q1/1K6&moves=6.5}{\color{white}\bfseries lösen\,\ensuremath{\nearrow}}}
\end{tcolorbox}
\vspace{2mm}
\noindent\begin{minipage}[t]{36mm}\vspace{0pt}\sffamily\small\raggedright
\hmfigure{längste eindeutige}{{\Large h\#6.5}}\hmfigure{eindeutig bei dieser Länge}{{\Large 6} Stellungen}\hmfigure{längstes Matt}{{\Large h\#8}}\hmfigure{in der Klasse}{242 Millionen Stellungen,\\ 61,5 Millionen mit Hilfsmatt}
\end{minipage}\hfill
\begin{minipage}[t]{80mm}\vspace{0pt}\begin{tikzpicture}\begin{axis}[width=72mm,height=42mm,ybar,bar width=1.6pt,xmin=-1,xmax=17,ymin=0,ymax=15.56,axis y line*=left,axis x line*=bottom,xlabel={\scriptsize Halbzüge bis Matt},ylabel={\scriptsize Stellungen (Mio.)},tick label style={font=\tiny},label style={font=\scriptsize},enlarge x limits=false,ylabel shift=-3pt,xlabel shift=-3pt]\addplot[fill=ink!55,draw=none] coordinates {(0,0.002447) (1,0.01999) (2,0.112) (3,0.8374) (4,1.441) (5,5.72) (6,5.449) (7,13.24) (8,8.59) (9,14.14) (10,5.477) (11,5.366) (12,0.8237) (13,0.2868) (14,0.0139) (15,0.001977) (16,1.1e-05)};\draw[densely dotted] (axis cs:13,0) -- (axis cs:13,15.56);\end{axis}\begin{axis}[width=72mm,height=42mm,axis y line*=right,axis x line=none,ymode=log,xmin=-1,xmax=17,ylabel={\scriptsize eindeutig (log)},tick label style={font=\tiny},label style={font=\scriptsize},ylabel shift=-3pt]\addplot[mark=*,mark size=0.7pt,draw=black,thin] coordinates {(0,2447) (1,19986) (2,45378) (3,258028) (4,201283) (5,488376) (6,192448) (7,255798) (8,76375) (9,76791) (10,12305) (11,4477) (12,104) (13,6)};\end{axis}\end{tikzpicture}\end{minipage}
\vspace{3mm}

\noindent\begin{minipage}[t]{60mm}\vspace{0pt}\sbox\hmboard{\chessboard[setfen=8/8/5k2/8/8/6R1/2q3q1/1K6 w - - 0 1]}\usebox\hmboard\\[1mm]
\makebox[\wd\hmboard]{\sffamily\hspace{3mm}\textbf{h\#6.5}\hfill(2+3)}\end{minipage}\hfill
\begin{minipage}[t]{52mm}\vspace{0pt}
\textcolor{ink}{\sffamily\textsc{Lösung}}\par\smallskip\small
\begin{tabular}{@{}r@{\enspace}l@{\enspace}l@{}}
\textbf{1.} & \ldots & \symking{}a1\\
\textbf{2.} & \symqueen{}b1+ & \symking{}xb1\\
\textbf{3.} & \symqueen{}e4+ & \symking{}c1\\
\textbf{4.} & \symqueen{}g6 & \symking{}d2\\
\textbf{5.} & \symking{}g7 & \symking{}e3\\
\textbf{6.} & \symking{}h6 & \symking{}f4\\
\textbf{7.} & \symking{}h5 & \symrook{}h3\#\\
\end{tabular}\par\vspace{3mm}
\qrcode[height=12mm]{https://helpman.komtera.lt/?fen=8/8/5k2/8/8/6R1/2q3q1/1K6&moves=6.5}\\{\sffamily\tiny\color{ink}scannen zum Lösen}
\end{minipage}
\vspace{2mm}

{\raggedright\noindent\textcolor{ink}{\sffamily\textsc{Themen}}\enspace \themetag{reines Matt} \themetag{Mustermatt} \themetag{Idealmatt} \themetag{Switchback} \themetag{Selbstblock} \themetag{Umnov}\par\smallskip{\footnotesize\textit{Hinweis:} die Partei am Zug steht im Diagramm im Schach.}\par}

\clearpage
\phantomsection\label{hm:KRvkqr}
\begin{tcolorbox}[enhanced,colback=ink,colframe=ink,arc=0pt,boxrule=0pt,left=4mm,right=4mm,top=2.5mm,bottom=2.5mm]
\color{white}\sffamily{\LARGE\bfseries KRvkqr}\hfill{\Large h\#6.5}\\[0.5mm]
{\small 5 Steine\quad Weiß am Zug}\hfill{\small Nr.~154\qquad\href{https://helpman.komtera.lt/?fen=8/8/5k2/8/6R1/8/2q4r/1K6&moves=6.5}{\color{white}\bfseries lösen\,\ensuremath{\nearrow}}}
\end{tcolorbox}
\vspace{2mm}
\noindent\begin{minipage}[t]{36mm}\vspace{0pt}\sffamily\small\raggedright
\hmfigure{längste eindeutige}{{\Large h\#6.5}}\hmfigure{eindeutig bei dieser Länge}{{\Large 5} Stellungen}\hmfigure{längstes Matt}{{\Large h\#8}}\hmfigure{in der Klasse}{242 Millionen Stellungen,\\ 134 Millionen mit Hilfsmatt}
\end{minipage}\hfill
\begin{minipage}[t]{80mm}\vspace{0pt}\begin{tikzpicture}\begin{axis}[width=72mm,height=42mm,ybar,bar width=1.6pt,xmin=-1,xmax=17,ymin=0,ymax=30.55,axis y line*=left,axis x line*=bottom,xlabel={\scriptsize Halbzüge bis Matt},ylabel={\scriptsize Stellungen (Mio.)},tick label style={font=\tiny},label style={font=\scriptsize},enlarge x limits=false,ylabel shift=-3pt,xlabel shift=-3pt]\addplot[fill=ink!55,draw=none] coordinates {(0,0.008898) (1,0.07053) (2,0.3581) (3,2.299) (4,4.125) (5,13.37) (6,14.39) (7,27.78) (8,20.54) (9,27.68) (10,12) (11,9.045) (12,1.676) (13,0.3654) (14,0.02221) (15,0.000975) (16,2e-06)};\draw[densely dotted] (axis cs:13,0) -- (axis cs:13,30.55);\end{axis}\begin{axis}[width=72mm,height=42mm,axis y line*=right,axis x line=none,ymode=log,xmin=-1,xmax=17,ylabel={\scriptsize eindeutig (log)},tick label style={font=\tiny},label style={font=\scriptsize},ylabel shift=-3pt]\addplot[mark=*,mark size=0.7pt,draw=black,thin] coordinates {(0,8898) (1,70532) (2,153296) (3,712823) (4,646006) (5,1093493) (6,499330) (7,482280) (8,176678) (9,131268) (10,25306) (11,4796) (12,245) (13,5)};\end{axis}\end{tikzpicture}\end{minipage}
\vspace{3mm}

\noindent\begin{minipage}[t]{60mm}\vspace{0pt}\sbox\hmboard{\chessboard[setfen=8/8/5k2/8/6R1/8/2q4r/1K6 w - - 0 1]}\usebox\hmboard\\[1mm]
\makebox[\wd\hmboard]{\sffamily\hspace{3mm}\textbf{h\#6.5}\hfill(2+3)}\end{minipage}\hfill
\begin{minipage}[t]{52mm}\vspace{0pt}
\textcolor{ink}{\sffamily\textsc{Lösung}}\par\smallskip\small
\begin{tabular}{@{}r@{\enspace}l@{\enspace}l@{}}
\textbf{1.} & \ldots & \symking{}a1\\
\textbf{2.} & \symqueen{}b1+ & \symking{}xb1\\
\textbf{3.} & \symrook{}h6 & \symking{}c2\\
\textbf{4.} & \symrook{}g6 & \symking{}d3\\
\textbf{5.} & \symking{}g7 & \symking{}e4\\
\textbf{6.} & \symking{}h6 & \symking{}f5\\
\textbf{7.} & \symrook{}g7 & \symrook{}h4\#\\
\end{tabular}\par\vspace{3mm}
\qrcode[height=12mm]{https://helpman.komtera.lt/?fen=8/8/5k2/8/6R1/8/2q4r/1K6&moves=6.5}\\{\sffamily\tiny\color{ink}scannen zum Lösen}
\end{minipage}
\vspace{2mm}

{\raggedright\noindent\textcolor{ink}{\sffamily\textsc{Themen}}\enspace \themetag{reines Matt} \themetag{Mustermatt} \themetag{Idealmatt} \themetag{Switchback} \themetag{Selbstblock} \themetag{Umnov}\par\smallskip{\footnotesize\textit{Hinweis:} die Partei am Zug steht im Diagramm im Schach.}\par}

\clearpage
\phantomsection\label{hm:KBBvkr}
\begin{tcolorbox}[enhanced,colback=ink,colframe=ink,arc=0pt,boxrule=0pt,left=4mm,right=4mm,top=2.5mm,bottom=2.5mm]
\color{white}\sffamily{\LARGE\bfseries KBBvkr}\hfill{\Large h\#6}\\[0.5mm]
{\small 5 Steine\quad Schwarz am Zug}\hfill{\small Nr.~155\qquad\href{https://helpman.komtera.lt/?fen=7k/8/7B/8/8/8/1r6/KB6&moves=6}{\color{white}\bfseries lösen\,\ensuremath{\nearrow}}}
\end{tcolorbox}
\vspace{2mm}
\noindent\begin{minipage}[t]{36mm}\vspace{0pt}\sffamily\small\raggedright
\hmfigure{längste eindeutige}{{\Large h\#6}}\hmfigure{eindeutig bei dieser Länge}{{\Large 28} Stellungen}\hmfigure{längstes Matt}{{\Large h\#7.5}}\hmfigure{in der Klasse}{242 Millionen Stellungen,\\ 39,6 Millionen mit Hilfsmatt}
\end{minipage}\hfill
\begin{minipage}[t]{80mm}\vspace{0pt}\begin{tikzpicture}\begin{axis}[width=72mm,height=42mm,ybar,bar width=1.6pt,xmin=-1,xmax=16,ymin=0,ymax=8.882,axis y line*=left,axis x line*=bottom,xlabel={\scriptsize Halbzüge bis Matt},ylabel={\scriptsize Stellungen (Mio.)},tick label style={font=\tiny},label style={font=\scriptsize},enlarge x limits=false,ylabel shift=-3pt,xlabel shift=-3pt]\addplot[fill=ink!55,draw=none] coordinates {(0,0.004014) (1,0.01836) (2,0.07828) (3,0.529) (4,1.691) (5,4.769) (6,8.074) (7,7.851) (8,7.895) (9,4.84) (10,2.605) (11,0.8789) (12,0.3271) (13,0.05353) (14,0.01201) (15,1.4e-05)};\draw[densely dotted] (axis cs:12,0) -- (axis cs:12,8.882);\end{axis}\begin{axis}[width=72mm,height=42mm,axis y line*=right,axis x line=none,ymode=log,xmin=-1,xmax=16,ylabel={\scriptsize eindeutig (log)},tick label style={font=\tiny},label style={font=\scriptsize},ylabel shift=-3pt]\addplot[mark=*,mark size=0.7pt,draw=black,thin] coordinates {(0,4014) (1,18278) (2,46968) (3,231014) (4,387651) (5,388371) (6,232880) (7,88444) (8,40903) (9,11326) (10,1454) (11,300) (12,28)};\end{axis}\end{tikzpicture}\end{minipage}
\vspace{3mm}

\noindent\begin{minipage}[t]{60mm}\vspace{0pt}\sbox\hmboard{\chessboard[setfen=7k/8/7B/8/8/8/1r6/KB6 b - - 0 1]}\usebox\hmboard\\[1mm]
\makebox[\wd\hmboard]{\sffamily\hspace{3mm}\textbf{h\#6}\hfill(3+2)}\end{minipage}\hfill
\begin{minipage}[t]{52mm}\vspace{0pt}
\textcolor{ink}{\sffamily\textsc{Lösung}}\par\smallskip\small
\begin{tabular}{@{}r@{\enspace}l@{\enspace}l@{}}
\textbf{1.} & \symrook{}b8 & \symking{}a2\\
\textbf{2.} & \symrook{}a8+ & \symking{}b3\\
\textbf{3.} & \symking{}g8 & \symking{}c4\\
\textbf{4.} & \symking{}f7 & \symking{}d5\\
\textbf{5.} & \symking{}e8 & \symking{}e6\\
\textbf{6.} & \symrook{}d8 & \symbishop{}g6\#\\
\end{tabular}\par\vspace{3mm}
\qrcode[height=12mm]{https://helpman.komtera.lt/?fen=7k/8/7B/8/8/8/1r6/KB6&moves=6}\\{\sffamily\tiny\color{ink}scannen zum Lösen}
\end{minipage}
\vspace{2mm}

{\raggedright\noindent\textcolor{ink}{\sffamily\textsc{Themen}}\enspace \themetag{Selbstblock} \themetag{ohne Schlag}\par}

\clearpage
\phantomsection\label{hm:KNNNvk}
\begin{tcolorbox}[enhanced,colback=ink,colframe=ink,arc=0pt,boxrule=0pt,left=4mm,right=4mm,top=2.5mm,bottom=2.5mm]
\color{white}\sffamily{\LARGE\bfseries KNNNvk}\hfill{\Large h\#6}\\[0.5mm]
{\small 5 Steine\quad Schwarz am Zug}\hfill{\small Nr.~156\qquad\href{https://helpman.komtera.lt/?fen=5k2/8/8/8/8/3K3N/N7/7N&moves=6}{\color{white}\bfseries lösen\,\ensuremath{\nearrow}}}
\end{tcolorbox}
\vspace{2mm}
\noindent\begin{minipage}[t]{36mm}\vspace{0pt}\sffamily\small\raggedright
\hmfigure{längste eindeutige}{{\Large h\#6}}\hmfigure{eindeutig bei dieser Länge}{{\Large 4} Stellungen}\hmfigure{längstes Matt}{{\Large h\#8}}\hmfigure{in der Klasse}{242 Millionen Stellungen,\\ 30,2 Millionen mit Hilfsmatt}
\end{minipage}\hfill
\begin{minipage}[t]{80mm}\vspace{0pt}\begin{tikzpicture}\begin{axis}[width=72mm,height=42mm,ybar,bar width=1.6pt,xmin=-1,xmax=17,ymin=0,ymax=8.611,axis y line*=left,axis x line*=bottom,xlabel={\scriptsize Halbzüge bis Matt},ylabel={\scriptsize Stellungen (Mio.)},tick label style={font=\tiny},label style={font=\scriptsize},enlarge x limits=false,ylabel shift=-3pt,xlabel shift=-3pt]\addplot[fill=ink!55,draw=none] coordinates {(0,0.001483) (1,0.006728) (2,0.01622) (3,0.08095) (4,0.3517) (5,1.37) (6,3.133) (7,4.831) (8,7.828) (9,5.115) (10,4.813) (11,1.54) (12,0.8518) (13,0.2279) (14,0.06039) (15,0.01272) (16,7e-06)};\draw[densely dotted] (axis cs:12,0) -- (axis cs:12,8.611);\end{axis}\begin{axis}[width=72mm,height=42mm,axis y line*=right,axis x line=none,ymode=log,xmin=-1,xmax=17,ylabel={\scriptsize eindeutig (log)},tick label style={font=\tiny},label style={font=\scriptsize},ylabel shift=-3pt]\addplot[mark=*,mark size=0.7pt,draw=black,thin] coordinates {(0,1483) (1,6017) (2,14539) (3,55313) (4,172999) (5,167402) (6,206132) (7,63227) (8,56950) (9,5828) (10,3310) (11,71) (12,4)};\end{axis}\end{tikzpicture}\end{minipage}
\vspace{3mm}

\noindent\begin{minipage}[t]{60mm}\vspace{0pt}\sbox\hmboard{\chessboard[setfen=5k2/8/8/8/8/3K3N/N7/7N b - - 0 1]}\usebox\hmboard\\[1mm]
\makebox[\wd\hmboard]{\sffamily\hspace{3mm}\textbf{h\#6}\hfill(4+1)}\end{minipage}\hfill
\begin{minipage}[t]{52mm}\vspace{0pt}
\textcolor{ink}{\sffamily\textsc{Lösung}}\par\smallskip\small
\begin{tabular}{@{}r@{\enspace}l@{\enspace}l@{}}
\textbf{1.} & \symking{}g7 & \symking{}e4\\
\textbf{2.} & \symking{}h6 & \symking{}f5\\
\textbf{3.} & \symking{}h5 & \symknight{}g3+\\
\textbf{4.} & \symking{}h4 & \symking{}g6\\
\textbf{5.} & \symking{}g4 & \symknight{}f2+\\
\textbf{6.} & \symking{}h4 & \symknight{}f5\#\\
\end{tabular}\par\vspace{3mm}
\qrcode[height=12mm]{https://helpman.komtera.lt/?fen=5k2/8/8/8/8/3K3N/N7/7N&moves=6}\\{\sffamily\tiny\color{ink}scannen zum Lösen}
\end{minipage}
\vspace{2mm}

{\raggedright\noindent\textcolor{ink}{\sffamily\textsc{Themen}}\enspace \themetag{reines Matt} \themetag{Spiegelmatt} \themetag{Switchback} \themetag{ein Stein zieht} \themetag{ohne Schlag}\par}

\clearpage
\phantomsection\label{hm:KQBNvk}
\begin{tcolorbox}[enhanced,colback=ink,colframe=ink,arc=0pt,boxrule=0pt,left=4mm,right=4mm,top=2.5mm,bottom=2.5mm]
\color{white}\sffamily{\LARGE\bfseries KQBNvk}\hfill{\Large h\#6}\\[0.5mm]
{\small 5 Steine\quad Schwarz am Zug}\hfill{\small Nr.~157\qquad\href{https://helpman.komtera.lt/?fen=8/6N1/8/8/8/2B5/8/K1kQ4&moves=6}{\color{white}\bfseries lösen\,\ensuremath{\nearrow}}}
\end{tcolorbox}
\vspace{2mm}
\noindent\begin{minipage}[t]{36mm}\vspace{0pt}\sffamily\small\raggedright
\hmfigure{längste eindeutige}{{\Large h\#6}}\hmfigure{eindeutig bei dieser Länge}{{\Large 98} Stellungen}\hmfigure{längstes Matt}{{\Large h\#8}}\hmfigure{in der Klasse}{242 Millionen Stellungen,\\ 156 Millionen mit Hilfsmatt}
\end{minipage}\hfill
\begin{minipage}[t]{80mm}\vspace{0pt}\begin{tikzpicture}\begin{axis}[width=72mm,height=42mm,ybar,bar width=1.6pt,xmin=-1,xmax=17,ymin=0,ymax=84.38,axis y line*=left,axis x line*=bottom,xlabel={\scriptsize Halbzüge bis Matt},ylabel={\scriptsize Stellungen (Mio.)},tick label style={font=\tiny},label style={font=\scriptsize},enlarge x limits=false,ylabel shift=-3pt,xlabel shift=-3pt]\addplot[fill=ink!55,draw=none] coordinates {(0,0.7562) (1,4.35) (2,15.03) (3,33.98) (4,76.71) (5,15.93) (6,7.218) (7,0.007076) (8,0.4232) (10,0.7224) (12,0.4035) (14,0.085) (16,0.001245)};\draw[densely dotted] (axis cs:12,0) -- (axis cs:12,84.38);\end{axis}\begin{axis}[width=72mm,height=42mm,axis y line*=right,axis x line=none,ymode=log,xmin=-1,xmax=17,ylabel={\scriptsize eindeutig (log)},tick label style={font=\tiny},label style={font=\scriptsize},ylabel shift=-3pt]\addplot[mark=*,mark size=0.7pt,draw=black,thin] coordinates {(0,756176) (1,3416160) (2,10717852) (3,7581588) (4,3879577) (5,1032) (6,26297) (8,13058) (10,3733) (12,98)};\end{axis}\end{tikzpicture}\end{minipage}
\vspace{3mm}

\noindent\begin{minipage}[t]{60mm}\vspace{0pt}\sbox\hmboard{\chessboard[setfen=8/6N1/8/8/8/2B5/8/K1kQ4 b - - 0 1]}\usebox\hmboard\\[1mm]
\makebox[\wd\hmboard]{\sffamily\hspace{3mm}\textbf{h\#6}\hfill(4+1)}\end{minipage}\hfill
\begin{minipage}[t]{52mm}\vspace{0pt}
\textcolor{ink}{\sffamily\textsc{Lösung}}\par\smallskip\small
\begin{tabular}{@{}r@{\enspace}l@{\enspace}l@{}}
\textbf{1.} & \symking{}xd1 & \symking{}a2\\
\textbf{2.} & \symking{}c1 & \symking{}b3\\
\textbf{3.} & \symking{}b1 & \symking{}c4\\
\textbf{4.} & \symking{}a2 & \symknight{}e6\\
\textbf{5.} & \symking{}a3 & \symbishop{}b4+\\
\textbf{6.} & \symking{}a4 & \symknight{}c5\#\\
\end{tabular}\par\vspace{3mm}
\qrcode[height=12mm]{https://helpman.komtera.lt/?fen=8/6N1/8/8/8/2B5/8/K1kQ4&moves=6}\\{\sffamily\tiny\color{ink}scannen zum Lösen}
\end{minipage}
\vspace{2mm}

{\raggedright\noindent\textcolor{ink}{\sffamily\textsc{Themen}}\enspace \themetag{Switchback} \themetag{ein Stein zieht}\par\smallskip{\footnotesize\textit{Hinweis:} die Partei am Zug steht im Diagramm im Schach; die Lösung beginnt mit einem Schlag.}\par}

\clearpage
\phantomsection\label{hm:KQQPvk}
\begin{tcolorbox}[enhanced,colback=ink,colframe=ink,arc=0pt,boxrule=0pt,left=4mm,right=4mm,top=2.5mm,bottom=2.5mm]
\color{white}\sffamily{\LARGE\bfseries KQQPvk}\hfill{\Large h\#6}\\[0.5mm]
{\small 5 Steine\quad Schwarz am Zug}\hfill{\small Nr.~158\qquad\href{https://helpman.komtera.lt/?fen=K7/8/8/8/8/8/P3Q3/6kQ&moves=6}{\color{white}\bfseries lösen\,\ensuremath{\nearrow}}}
\end{tcolorbox}
\vspace{2mm}
\noindent\begin{minipage}[t]{36mm}\vspace{0pt}\sffamily\small\raggedright
\hmfigure{längste eindeutige}{{\Large h\#6}}\hmfigure{eindeutig bei dieser Länge}{{\Large 6} Stellungen}\hmfigure{längstes Matt}{{\Large h\#7}}\hmfigure{in der Klasse}{710 Millionen Stellungen,\\ 215 Millionen mit Hilfsmatt}
\end{minipage}\hfill
\begin{minipage}[t]{80mm}\vspace{0pt}\begin{tikzpicture}\begin{axis}[width=72mm,height=42mm,ybar,bar width=1.6pt,xmin=-1,xmax=15,ymin=0,ymax=86.64,axis y line*=left,axis x line*=bottom,xlabel={\scriptsize Halbzüge bis Matt},ylabel={\scriptsize Stellungen (Mio.)},tick label style={font=\tiny},label style={font=\scriptsize},enlarge x limits=false,ylabel shift=-3pt,xlabel shift=-3pt]\addplot[fill=ink!55,draw=none] coordinates {(0,5.776) (1,16.8) (2,61.61) (3,46.69) (4,78.76) (5,0.8251) (6,2.808) (8,1.404) (10,0.2061) (12,0.008242) (14,1.8e-05)};\draw[densely dotted] (axis cs:12,0) -- (axis cs:12,86.64);\end{axis}\begin{axis}[width=72mm,height=42mm,axis y line*=right,axis x line=none,ymode=log,xmin=-1,xmax=15,ylabel={\scriptsize eindeutig (log)},tick label style={font=\tiny},label style={font=\scriptsize},ylabel shift=-3pt]\addplot[mark=*,mark size=0.7pt,draw=black,thin] coordinates {(0,5776157) (1,8200270) (2,23420425) (3,739884) (4,314046) (6,106788) (8,9908) (10,112) (12,6)};\end{axis}\end{tikzpicture}\end{minipage}
\vspace{3mm}

\noindent\begin{minipage}[t]{60mm}\vspace{0pt}\sbox\hmboard{\chessboard[setfen=K7/8/8/8/8/8/P3Q3/6kQ b - - 0 1]}\usebox\hmboard\\[1mm]
\makebox[\wd\hmboard]{\sffamily\hspace{3mm}\textbf{h\#6}\hfill(4+1)}\end{minipage}\hfill
\begin{minipage}[t]{52mm}\vspace{0pt}
\textcolor{ink}{\sffamily\textsc{Lösung}}\par\smallskip\small
\begin{tabular}{@{}r@{\enspace}l@{\enspace}l@{}}
\textbf{1.} & \symking{}xh1 & \symking{}b7\\
\textbf{2.} & \symking{}g1 & \symking{}c6\\
\textbf{3.} & \symking{}h1 & \symking{}d5\\
\textbf{4.} & \symking{}g1 & \symking{}e4\\
\textbf{5.} & \symking{}h1 & \symking{}f3\\
\textbf{6.} & \symking{}g1 & \symqueen{}g2\#\\
\end{tabular}\par\vspace{3mm}
\qrcode[height=12mm]{https://helpman.komtera.lt/?fen=K7/8/8/8/8/8/P3Q3/6kQ&moves=6}\\{\sffamily\tiny\color{ink}scannen zum Lösen}
\end{minipage}
\vspace{2mm}

{\raggedright\noindent\textcolor{ink}{\sffamily\textsc{Themen}}\enspace \themetag{Switchback} \themetag{ein Stein zieht} \themetag{Pendel} \themetag{ohne Schach}\par\smallskip{\footnotesize\textit{Hinweis:} die Partei am Zug steht im Diagramm im Schach; die Lösung beginnt mit einem Schlag.}\par}

\clearpage
\phantomsection\label{hm:KQQvkn}
\begin{tcolorbox}[enhanced,colback=ink,colframe=ink,arc=0pt,boxrule=0pt,left=4mm,right=4mm,top=2.5mm,bottom=2.5mm]
\color{white}\sffamily{\LARGE\bfseries KQQvkn}\hfill{\Large h\#6}\\[0.5mm]
{\small 5 Steine\quad Schwarz am Zug}\hfill{\small Nr.~159\qquad\href{https://helpman.komtera.lt/?fen=7Q/7k/8/6Q1/8/8/1n6/K7&moves=6}{\color{white}\bfseries lösen\,\ensuremath{\nearrow}}}
\end{tcolorbox}
\vspace{2mm}
\noindent\begin{minipage}[t]{36mm}\vspace{0pt}\sffamily\small\raggedright
\hmfigure{längste eindeutige}{{\Large h\#6}}\hmfigure{eindeutig bei dieser Länge}{{\Large 4} Stellungen}\hmfigure{längstes Matt}{{\Large h\#7}}\hmfigure{in der Klasse}{242 Millionen Stellungen,\\ 68,9 Millionen mit Hilfsmatt}
\end{minipage}\hfill
\begin{minipage}[t]{80mm}\vspace{0pt}\begin{tikzpicture}\begin{axis}[width=72mm,height=42mm,ybar,bar width=1.6pt,xmin=-1,xmax=15,ymin=0,ymax=27.06,axis y line*=left,axis x line*=bottom,xlabel={\scriptsize Halbzüge bis Matt},ylabel={\scriptsize Stellungen (Mio.)},tick label style={font=\tiny},label style={font=\scriptsize},enlarge x limits=false,ylabel shift=-3pt,xlabel shift=-3pt]\addplot[fill=ink!55,draw=none] coordinates {(0,1.494) (1,4.55) (2,19.23) (3,16.44) (4,24.6) (5,0.9388) (6,0.6665) (7,9.3e-05) (8,0.7162) (9,3e-06) (10,0.26) (12,0.03344) (14,2.7e-05)};\draw[densely dotted] (axis cs:12,0) -- (axis cs:12,27.06);\end{axis}\begin{axis}[width=72mm,height=42mm,axis y line*=right,axis x line=none,ymode=log,xmin=-1,xmax=15,ylabel={\scriptsize eindeutig (log)},tick label style={font=\tiny},label style={font=\scriptsize},ylabel shift=-3pt]\addplot[mark=*,mark size=0.7pt,draw=black,thin] coordinates {(0,1493820) (1,2430552) (2,6163320) (3,238273) (4,52648) (5,5) (6,35658) (7,9) (8,10579) (10,692) (12,4)};\end{axis}\end{tikzpicture}\end{minipage}
\vspace{3mm}

\noindent\begin{minipage}[t]{60mm}\vspace{0pt}\sbox\hmboard{\chessboard[setfen=7Q/7k/8/6Q1/8/8/1n6/K7 b - - 0 1]}\usebox\hmboard\\[1mm]
\makebox[\wd\hmboard]{\sffamily\hspace{3mm}\textbf{h\#6}\hfill(3+2)}\end{minipage}\hfill
\begin{minipage}[t]{52mm}\vspace{0pt}
\textcolor{ink}{\sffamily\textsc{Lösung}}\par\smallskip\small
\begin{tabular}{@{}r@{\enspace}l@{\enspace}l@{}}
\textbf{1.} & \symking{}xh8 & \symking{}xb2\\
\textbf{2.} & \symking{}h7 & \symking{}c3\\
\textbf{3.} & \symking{}h8 & \symking{}d4\\
\textbf{4.} & \symking{}h7 & \symking{}e5\\
\textbf{5.} & \symking{}h8 & \symking{}f6\\
\textbf{6.} & \symking{}h7 & \symqueen{}g7\#\\
\end{tabular}\par\vspace{3mm}
\qrcode[height=12mm]{https://helpman.komtera.lt/?fen=7Q/7k/8/6Q1/8/8/1n6/K7&moves=6}\\{\sffamily\tiny\color{ink}scannen zum Lösen}
\end{minipage}
\vspace{2mm}

{\raggedright\noindent\textcolor{ink}{\sffamily\textsc{Themen}}\enspace \themetag{Switchback} \themetag{ein Stein zieht} \themetag{Pendel} \themetag{ohne Schach}\par\smallskip{\footnotesize\textit{Hinweis:} die Partei am Zug steht im Diagramm im Schach; die Lösung beginnt mit einem Schlag.}\par}

\clearpage
\phantomsection\label{hm:KQQvkp}
\begin{tcolorbox}[enhanced,colback=ink,colframe=ink,arc=0pt,boxrule=0pt,left=4mm,right=4mm,top=2.5mm,bottom=2.5mm]
\color{white}\sffamily{\LARGE\bfseries KQQvkp}\hfill{\Large h\#6}\\[0.5mm]
{\small 5 Steine\quad Schwarz am Zug}\hfill{\small Nr.~160\qquad\href{https://helpman.komtera.lt/?fen=K7/1p6/8/8/8/8/4Q3/6kQ&moves=6}{\color{white}\bfseries lösen\,\ensuremath{\nearrow}}}
\end{tcolorbox}
\vspace{2mm}
\noindent\begin{minipage}[t]{36mm}\vspace{0pt}\sffamily\small\raggedright
\hmfigure{längste eindeutige}{{\Large h\#6}}\hmfigure{eindeutig bei dieser Länge}{{\Large 6} Stellungen}\hmfigure{längstes Matt}{{\Large h\#7}}\hmfigure{in der Klasse}{710 Millionen Stellungen,\\ 216 Millionen mit Hilfsmatt}
\end{minipage}\hfill
\begin{minipage}[t]{80mm}\vspace{0pt}\begin{tikzpicture}\begin{axis}[width=72mm,height=42mm,ybar,bar width=1.6pt,xmin=-1,xmax=15,ymin=0,ymax=82.84,axis y line*=left,axis x line*=bottom,xlabel={\scriptsize Halbzüge bis Matt},ylabel={\scriptsize Stellungen (Mio.)},tick label style={font=\tiny},label style={font=\scriptsize},enlarge x limits=false,ylabel shift=-3pt,xlabel shift=-3pt]\addplot[fill=ink!55,draw=none] coordinates {(0,5.304) (1,15.98) (2,63.59) (3,49.08) (4,75.31) (5,1.232) (6,2.139) (8,2.133) (10,0.664) (12,0.08594) (14,0.000233)};\draw[densely dotted] (axis cs:12,0) -- (axis cs:12,82.84);\end{axis}\begin{axis}[width=72mm,height=42mm,axis y line*=right,axis x line=none,ymode=log,xmin=-1,xmax=15,ylabel={\scriptsize eindeutig (log)},tick label style={font=\tiny},label style={font=\scriptsize},ylabel shift=-3pt]\addplot[mark=*,mark size=0.7pt,draw=black,thin] coordinates {(0,5304427) (1,8123188) (2,21898972) (3,549542) (4,183295) (5,1) (6,105976) (8,25038) (10,364) (12,6)};\end{axis}\end{tikzpicture}\end{minipage}
\vspace{3mm}

\noindent\begin{minipage}[t]{60mm}\vspace{0pt}\sbox\hmboard{\chessboard[setfen=K7/1p6/8/8/8/8/4Q3/6kQ b - - 0 1]}\usebox\hmboard\\[1mm]
\makebox[\wd\hmboard]{\sffamily\hspace{3mm}\textbf{h\#6}\hfill(3+2)}\end{minipage}\hfill
\begin{minipage}[t]{52mm}\vspace{0pt}
\textcolor{ink}{\sffamily\textsc{Lösung}}\par\smallskip\small
\begin{tabular}{@{}r@{\enspace}l@{\enspace}l@{}}
\textbf{1.} & \symking{}xh1 & \symking{}xb7\\
\textbf{2.} & \symking{}g1 & \symking{}c6\\
\textbf{3.} & \symking{}h1 & \symking{}d5\\
\textbf{4.} & \symking{}g1 & \symking{}e4\\
\textbf{5.} & \symking{}h1 & \symking{}f3\\
\textbf{6.} & \symking{}g1 & \symqueen{}g2\#\\
\end{tabular}\par\vspace{3mm}
\qrcode[height=12mm]{https://helpman.komtera.lt/?fen=K7/1p6/8/8/8/8/4Q3/6kQ&moves=6}\\{\sffamily\tiny\color{ink}scannen zum Lösen}
\end{minipage}
\vspace{2mm}

{\raggedright\noindent\textcolor{ink}{\sffamily\textsc{Themen}}\enspace \themetag{Switchback} \themetag{ein Stein zieht} \themetag{Pendel} \themetag{ohne Schach}\par\smallskip{\footnotesize\textit{Hinweis:} die Partei am Zug steht im Diagramm im Schach; die Lösung beginnt mit einem Schlag.}\par}

\clearpage
\phantomsection\label{hm:KQRBvk}
\begin{tcolorbox}[enhanced,colback=ink,colframe=ink,arc=0pt,boxrule=0pt,left=4mm,right=4mm,top=2.5mm,bottom=2.5mm]
\color{white}\sffamily{\LARGE\bfseries KQRBvk}\hfill{\Large h\#6}\\[0.5mm]
{\small 5 Steine\quad Schwarz am Zug}\hfill{\small Nr.~161\qquad\href{https://helpman.komtera.lt/?fen=8/7Q/7k/8/7B/6R1/8/K7&moves=6}{\color{white}\bfseries lösen\,\ensuremath{\nearrow}}}
\end{tcolorbox}
\vspace{2mm}
\noindent\begin{minipage}[t]{36mm}\vspace{0pt}\sffamily\small\raggedright
\hmfigure{längste eindeutige}{{\Large h\#6}}\hmfigure{eindeutig bei dieser Länge}{{\Large 29} Stellungen}\hmfigure{längstes Matt}{{\Large h\#7}}\hmfigure{in der Klasse}{242 Millionen Stellungen,\\ 147 Millionen mit Hilfsmatt}
\end{minipage}\hfill
\begin{minipage}[t]{80mm}\vspace{0pt}\begin{tikzpicture}\begin{axis}[width=72mm,height=42mm,ybar,bar width=1.6pt,xmin=-1,xmax=15,ymin=0,ymax=74.27,axis y line*=left,axis x line*=bottom,xlabel={\scriptsize Halbzüge bis Matt},ylabel={\scriptsize Stellungen (Mio.)},tick label style={font=\tiny},label style={font=\scriptsize},enlarge x limits=false,ylabel shift=-3pt,xlabel shift=-3pt]\addplot[fill=ink!55,draw=none] coordinates {(0,2.487) (1,8.34) (2,28.06) (3,35.15) (4,67.52) (5,2.567) (6,1.457) (7,1e-06) (8,0.7408) (10,0.4047) (12,0.06465) (14,0.002525)};\draw[densely dotted] (axis cs:12,0) -- (axis cs:12,74.27);\end{axis}\begin{axis}[width=72mm,height=42mm,axis y line*=right,axis x line=none,ymode=log,xmin=-1,xmax=15,ylabel={\scriptsize eindeutig (log)},tick label style={font=\tiny},label style={font=\scriptsize},ylabel shift=-3pt]\addplot[mark=*,mark size=0.7pt,draw=black,thin] coordinates {(0,2486646) (1,5065579) (2,13752831) (3,2577232) (4,750815) (6,84955) (8,8631) (10,620) (12,29)};\end{axis}\end{tikzpicture}\end{minipage}
\vspace{3mm}

\noindent\begin{minipage}[t]{60mm}\vspace{0pt}\sbox\hmboard{\chessboard[setfen=8/7Q/7k/8/7B/6R1/8/K7 b - - 0 1]}\usebox\hmboard\\[1mm]
\makebox[\wd\hmboard]{\sffamily\hspace{3mm}\textbf{h\#6}\hfill(4+1)}\end{minipage}\hfill
\begin{minipage}[t]{52mm}\vspace{0pt}
\textcolor{ink}{\sffamily\textsc{Lösung}}\par\smallskip\small
\begin{tabular}{@{}r@{\enspace}l@{\enspace}l@{}}
\textbf{1.} & \symking{}xh7 & \symking{}b2\\
\textbf{2.} & \symking{}h6 & \symking{}c3\\
\textbf{3.} & \symking{}h5 & \symking{}d4\\
\textbf{4.} & \symking{}xh4 & \symking{}e5\\
\textbf{5.} & \symking{}h5 & \symking{}f6\\
\textbf{6.} & \symking{}h6 & \symrook{}h3\#\\
\end{tabular}\par\vspace{3mm}
\qrcode[height=12mm]{https://helpman.komtera.lt/?fen=8/7Q/7k/8/7B/6R1/8/K7&moves=6}\\{\sffamily\tiny\color{ink}scannen zum Lösen}
\end{minipage}
\vspace{2mm}

{\raggedright\noindent\textcolor{ink}{\sffamily\textsc{Themen}}\enspace \themetag{reines Matt} \themetag{Mustermatt} \themetag{Idealmatt} \themetag{Spiegelmatt} \themetag{Switchback} \themetag{ein Stein zieht} \themetag{ohne Schach}\par\smallskip{\footnotesize\textit{Hinweis:} die Partei am Zug steht im Diagramm im Schach; die Lösung beginnt mit einem Schlag.}\par}

\clearpage
\phantomsection\label{hm:KQRPvk}
\begin{tcolorbox}[enhanced,colback=ink,colframe=ink,arc=0pt,boxrule=0pt,left=4mm,right=4mm,top=2.5mm,bottom=2.5mm]
\color{white}\sffamily{\LARGE\bfseries KQRPvk}\hfill{\Large h\#6}\\[0.5mm]
{\small 5 Steine\quad Schwarz am Zug}\hfill{\small Nr.~162\qquad\href{https://helpman.komtera.lt/?fen=7k/7Q/6R1/8/8/8/5P2/K7&moves=6}{\color{white}\bfseries lösen\,\ensuremath{\nearrow}}}
\end{tcolorbox}
\vspace{2mm}
\noindent\begin{minipage}[t]{36mm}\vspace{0pt}\sffamily\small\raggedright
\hmfigure{längste eindeutige}{{\Large h\#6}}\hmfigure{eindeutig bei dieser Länge}{{\Large 14} Stellungen}\hmfigure{längstes Matt}{{\Large h\#7}}\hmfigure{in der Klasse}{710 Millionen Stellungen,\\ 460 Millionen mit Hilfsmatt}
\end{minipage}\hfill
\begin{minipage}[t]{80mm}\vspace{0pt}\begin{tikzpicture}\begin{axis}[width=72mm,height=42mm,ybar,bar width=1.6pt,xmin=-1,xmax=15,ymin=0,ymax=242.9,axis y line*=left,axis x line*=bottom,xlabel={\scriptsize Halbzüge bis Matt},ylabel={\scriptsize Stellungen (Mio.)},tick label style={font=\tiny},label style={font=\scriptsize},enlarge x limits=false,ylabel shift=-3pt,xlabel shift=-3pt]\addplot[fill=ink!55,draw=none] coordinates {(0,5.186) (1,19.85) (2,57.72) (3,109.2) (4,220.8) (5,26.72) (6,17.22) (7,0.0316) (8,2.216) (10,0.876) (12,0.1032) (14,0.001497)};\draw[densely dotted] (axis cs:12,0) -- (axis cs:12,242.9);\end{axis}\begin{axis}[width=72mm,height=42mm,axis y line*=right,axis x line=none,ymode=log,xmin=-1,xmax=15,ylabel={\scriptsize eindeutig (log)},tick label style={font=\tiny},label style={font=\scriptsize},ylabel shift=-3pt]\addplot[mark=*,mark size=0.7pt,draw=black,thin] coordinates {(0,5186183) (1,12246646) (2,26964644) (3,12495740) (4,6735753) (5,11158) (6,191513) (8,11846) (10,598) (12,14)};\end{axis}\end{tikzpicture}\end{minipage}
\vspace{3mm}

\noindent\begin{minipage}[t]{60mm}\vspace{0pt}\sbox\hmboard{\chessboard[setfen=7k/7Q/6R1/8/8/8/5P2/K7 b - - 0 1]}\usebox\hmboard\\[1mm]
\makebox[\wd\hmboard]{\sffamily\hspace{3mm}\textbf{h\#6}\hfill(4+1)}\end{minipage}\hfill
\begin{minipage}[t]{52mm}\vspace{0pt}
\textcolor{ink}{\sffamily\textsc{Lösung}}\par\smallskip\small
\begin{tabular}{@{}r@{\enspace}l@{\enspace}l@{}}
\textbf{1.} & \symking{}xh7 & f4\\
\textbf{2.} & \symking{}h8 & f5\\
\textbf{3.} & \symking{}h7 & f6\\
\textbf{4.} & \symking{}h8 & f7\\
\textbf{5.} & \symking{}h7 & \symrook{}g7+\\
\textbf{6.} & \symking{}h8 & f8=\symqueen{}\#\\
\end{tabular}\par\vspace{3mm}
\qrcode[height=12mm]{https://helpman.komtera.lt/?fen=7k/7Q/6R1/8/8/8/5P2/K7&moves=6}\\{\sffamily\tiny\color{ink}scannen zum Lösen}
\end{minipage}
\vspace{2mm}

{\raggedright\noindent\textcolor{ink}{\sffamily\textsc{Themen}}\enspace \themetag{Umwandlung} \themetag{Excelsior} \themetag{Switchback} \themetag{ein Stein zieht} \themetag{Phönix} \themetag{Pendel}\par\smallskip{\footnotesize\textit{Hinweis:} die Partei am Zug steht im Diagramm im Schach; die Lösung beginnt mit einem Schlag.}\par}

\clearpage
\phantomsection\label{hm:KQRvkn}
\begin{tcolorbox}[enhanced,colback=ink,colframe=ink,arc=0pt,boxrule=0pt,left=4mm,right=4mm,top=2.5mm,bottom=2.5mm]
\color{white}\sffamily{\LARGE\bfseries KQRvkn}\hfill{\Large h\#6}\\[0.5mm]
{\small 5 Steine\quad Schwarz am Zug}\hfill{\small Nr.~163\qquad\href{https://helpman.komtera.lt/?fen=7R/7k/8/6Q1/8/8/1n6/K7&moves=6}{\color{white}\bfseries lösen\,\ensuremath{\nearrow}}}
\end{tcolorbox}
\vspace{2mm}
\noindent\begin{minipage}[t]{36mm}\vspace{0pt}\sffamily\small\raggedright
\hmfigure{längste eindeutige}{{\Large h\#6}}\hmfigure{eindeutig bei dieser Länge}{{\Large 4} Stellungen}\hmfigure{längstes Matt}{{\Large h\#7}}\hmfigure{in der Klasse}{242 Millionen Stellungen,\\ 147 Millionen mit Hilfsmatt}
\end{minipage}\hfill
\begin{minipage}[t]{80mm}\vspace{0pt}\begin{tikzpicture}\begin{axis}[width=72mm,height=42mm,ybar,bar width=1.6pt,xmin=-1,xmax=15,ymin=0,ymax=76.54,axis y line*=left,axis x line*=bottom,xlabel={\scriptsize Halbzüge bis Matt},ylabel={\scriptsize Stellungen (Mio.)},tick label style={font=\tiny},label style={font=\scriptsize},enlarge x limits=false,ylabel shift=-3pt,xlabel shift=-3pt]\addplot[fill=ink!55,draw=none] coordinates {(0,1.291) (1,5.03) (2,16.8) (3,35.68) (4,69.58) (5,12.42) (6,5.027) (7,0.0289) (8,0.7253) (9,2e-05) (10,0.4896) (12,0.09207) (14,0.002394)};\draw[densely dotted] (axis cs:12,0) -- (axis cs:12,76.54);\end{axis}\begin{axis}[width=72mm,height=42mm,axis y line*=right,axis x line=none,ymode=log,xmin=-1,xmax=15,ylabel={\scriptsize eindeutig (log)},tick label style={font=\tiny},label style={font=\scriptsize},ylabel shift=-3pt]\addplot[mark=*,mark size=0.7pt,draw=black,thin] coordinates {(0,1290683) (1,3359840) (2,5916244) (3,3414520) (4,2057195) (5,6799) (6,49675) (7,66) (8,6991) (10,840) (12,4)};\end{axis}\end{tikzpicture}\end{minipage}
\vspace{3mm}

\noindent\begin{minipage}[t]{60mm}\vspace{0pt}\sbox\hmboard{\chessboard[setfen=7R/7k/8/6Q1/8/8/1n6/K7 b - - 0 1]}\usebox\hmboard\\[1mm]
\makebox[\wd\hmboard]{\sffamily\hspace{3mm}\textbf{h\#6}\hfill(3+2)}\end{minipage}\hfill
\begin{minipage}[t]{52mm}\vspace{0pt}
\textcolor{ink}{\sffamily\textsc{Lösung}}\par\smallskip\small
\begin{tabular}{@{}r@{\enspace}l@{\enspace}l@{}}
\textbf{1.} & \symking{}xh8 & \symking{}xb2\\
\textbf{2.} & \symking{}h7 & \symking{}c3\\
\textbf{3.} & \symking{}h8 & \symking{}d4\\
\textbf{4.} & \symking{}h7 & \symking{}e5\\
\textbf{5.} & \symking{}h8 & \symking{}f6\\
\textbf{6.} & \symking{}h7 & \symqueen{}g7\#\\
\end{tabular}\par\vspace{3mm}
\qrcode[height=12mm]{https://helpman.komtera.lt/?fen=7R/7k/8/6Q1/8/8/1n6/K7&moves=6}\\{\sffamily\tiny\color{ink}scannen zum Lösen}
\end{minipage}
\vspace{2mm}

{\raggedright\noindent\textcolor{ink}{\sffamily\textsc{Themen}}\enspace \themetag{Switchback} \themetag{ein Stein zieht} \themetag{Pendel} \themetag{ohne Schach}\par\smallskip{\footnotesize\textit{Hinweis:} die Partei am Zug steht im Diagramm im Schach; die Lösung beginnt mit einem Schlag.}\par}

\clearpage
\phantomsection\label{hm:KQRvkp}
\begin{tcolorbox}[enhanced,colback=ink,colframe=ink,arc=0pt,boxrule=0pt,left=4mm,right=4mm,top=2.5mm,bottom=2.5mm]
\color{white}\sffamily{\LARGE\bfseries KQRvkp}\hfill{\Large h\#6}\\[0.5mm]
{\small 5 Steine\quad Schwarz am Zug}\hfill{\small Nr.~164\qquad\href{https://helpman.komtera.lt/?fen=7Q/7k/8/6R1/p7/8/8/K7&moves=6}{\color{white}\bfseries lösen\,\ensuremath{\nearrow}}}
\end{tcolorbox}
\vspace{2mm}
\noindent\begin{minipage}[t]{36mm}\vspace{0pt}\sffamily\small\raggedright
\hmfigure{längste eindeutige}{{\Large h\#6}}\hmfigure{eindeutig bei dieser Länge}{{\Large 72} Stellungen}\hmfigure{längstes Matt}{{\Large h\#7}}\hmfigure{in der Klasse}{710 Millionen Stellungen,\\ 459 Millionen mit Hilfsmatt}
\end{minipage}\hfill
\begin{minipage}[t]{80mm}\vspace{0pt}\begin{tikzpicture}\begin{axis}[width=72mm,height=42mm,ybar,bar width=1.6pt,xmin=-1,xmax=15,ymin=0,ymax=238.4,axis y line*=left,axis x line*=bottom,xlabel={\scriptsize Halbzüge bis Matt},ylabel={\scriptsize Stellungen (Mio.)},tick label style={font=\tiny},label style={font=\scriptsize},enlarge x limits=false,ylabel shift=-3pt,xlabel shift=-3pt]\addplot[fill=ink!55,draw=none] coordinates {(0,4.77) (1,18.26) (2,56.58) (3,111.1) (4,216.7) (5,30.86) (6,16.58) (7,0.09646) (8,2.274) (9,3.2e-05) (10,1.455) (11,4e-06) (12,0.2645) (14,0.006762)};\draw[densely dotted] (axis cs:12,0) -- (axis cs:12,238.4);\end{axis}\begin{axis}[width=72mm,height=42mm,axis y line*=right,axis x line=none,ymode=log,xmin=-1,xmax=15,ylabel={\scriptsize eindeutig (log)},tick label style={font=\tiny},label style={font=\scriptsize},ylabel shift=-3pt]\addplot[mark=*,mark size=0.7pt,draw=black,thin] coordinates {(0,4769537) (1,11805207) (2,22264648) (3,9642456) (4,5664474) (5,30830) (6,160146) (7,4) (8,16222) (9,1) (10,1179) (12,72)};\end{axis}\end{tikzpicture}\end{minipage}
\vspace{3mm}

\noindent\begin{minipage}[t]{60mm}\vspace{0pt}\sbox\hmboard{\chessboard[setfen=7Q/7k/8/6R1/p7/8/8/K7 b - - 0 1]}\usebox\hmboard\\[1mm]
\makebox[\wd\hmboard]{\sffamily\hspace{3mm}\textbf{h\#6}\hfill(3+2)}\end{minipage}\hfill
\begin{minipage}[t]{52mm}\vspace{0pt}
\textcolor{ink}{\sffamily\textsc{Lösung}}\par\smallskip\small
\begin{tabular}{@{}r@{\enspace}l@{\enspace}l@{}}
\textbf{1.} & \symking{}xh8 & \symking{}b2\\
\textbf{2.} & a3+ & \symking{}c3\\
\textbf{3.} & a2 & \symking{}d4\\
\textbf{4.} & a1=\symrook{} & \symking{}e5\\
\textbf{5.} & \symrook{}a8 & \symking{}f6\\
\textbf{6.} & \symrook{}g8 & \symrook{}h5\#\\
\end{tabular}\par\vspace{3mm}
\qrcode[height=12mm]{https://helpman.komtera.lt/?fen=7Q/7k/8/6R1/p7/8/8/K7&moves=6}\\{\sffamily\tiny\color{ink}scannen zum Lösen}
\end{minipage}
\vspace{2mm}

{\raggedright\noindent\textcolor{ink}{\sffamily\textsc{Themen}}\enspace \themetag{reines Matt} \themetag{Mustermatt} \themetag{Idealmatt} \themetag{Umwandlung} \themetag{Unterverwandlung} \themetag{Selbstblock} \themetag{Kniest}\par\smallskip{\footnotesize\textit{Hinweis:} die Partei am Zug steht im Diagramm im Schach; die Lösung beginnt mit einem Schlag.}\par}

\clearpage
\phantomsection\label{hm:KQRvkq}
\begin{tcolorbox}[enhanced,colback=ink,colframe=ink,arc=0pt,boxrule=0pt,left=4mm,right=4mm,top=2.5mm,bottom=2.5mm]
\color{white}\sffamily{\LARGE\bfseries KQRvkq}\hfill{\Large h\#6}\\[0.5mm]
{\small 5 Steine\quad Schwarz am Zug}\hfill{\small Nr.~165\qquad\href{https://helpman.komtera.lt/?fen=8/1qR5/6Qk/8/8/8/8/K7&moves=6}{\color{white}\bfseries lösen\,\ensuremath{\nearrow}}}
\end{tcolorbox}
\vspace{2mm}
\noindent\begin{minipage}[t]{36mm}\vspace{0pt}\sffamily\small\raggedright
\hmfigure{längste eindeutige}{{\Large h\#6}}\hmfigure{eindeutig bei dieser Länge}{{\Large 2} Stellungen}\hmfigure{längstes Matt}{{\Large h\#7}}\hmfigure{in der Klasse}{242 Millionen Stellungen,\\ 121 Millionen mit Hilfsmatt}
\end{minipage}\hfill
\begin{minipage}[t]{80mm}\vspace{0pt}\begin{tikzpicture}\begin{axis}[width=72mm,height=42mm,ybar,bar width=1.6pt,xmin=-1,xmax=15,ymin=0,ymax=57.72,axis y line*=left,axis x line*=bottom,xlabel={\scriptsize Halbzüge bis Matt},ylabel={\scriptsize Stellungen (Mio.)},tick label style={font=\tiny},label style={font=\scriptsize},enlarge x limits=false,ylabel shift=-3pt,xlabel shift=-3pt]\addplot[fill=ink!55,draw=none] coordinates {(0,0.5957) (1,2.567) (2,11.56) (3,31.31) (4,52.47) (5,19.2) (6,2.462) (7,0.04626) (8,0.5659) (9,0.001107) (10,0.4022) (11,0.000454) (12,0.08088) (13,5.8e-05) (14,0.002302)};\draw[densely dotted] (axis cs:12,0) -- (axis cs:12,57.72);\end{axis}\begin{axis}[width=72mm,height=42mm,axis y line*=right,axis x line=none,ymode=log,xmin=-1,xmax=15,ylabel={\scriptsize eindeutig (log)},tick label style={font=\tiny},label style={font=\scriptsize},ylabel shift=-3pt]\addplot[mark=*,mark size=0.7pt,draw=black,thin] coordinates {(0,595701) (1,1992138) (2,4913797) (3,3909435) (4,1209563) (5,19631) (6,26518) (7,69) (8,7199) (9,1) (10,1297) (12,2)};\end{axis}\end{tikzpicture}\end{minipage}
\vspace{3mm}

\noindent\begin{minipage}[t]{60mm}\vspace{0pt}\sbox\hmboard{\chessboard[setfen=8/1qR5/6Qk/8/8/8/8/K7 b - - 0 1]}\usebox\hmboard\\[1mm]
\makebox[\wd\hmboard]{\sffamily\hspace{3mm}\textbf{h\#6}\hfill(3+2)}\end{minipage}\hfill
\begin{minipage}[t]{52mm}\vspace{0pt}
\textcolor{ink}{\sffamily\textsc{Lösung}}\par\smallskip\small
\begin{tabular}{@{}r@{\enspace}l@{\enspace}l@{}}
\textbf{1.} & \symking{}xg6 & \symking{}a2\\
\textbf{2.} & \symqueen{}d5+ & \symking{}a3\\
\textbf{3.} & \symqueen{}f7 & \symking{}b4\\
\textbf{4.} & \symking{}g7 & \symking{}c5\\
\textbf{5.} & \symking{}f8 & \symking{}d6\\
\textbf{6.} & \symking{}e8 & \symrook{}c8\#\\
\end{tabular}\par\vspace{3mm}
\qrcode[height=12mm]{https://helpman.komtera.lt/?fen=8/1qR5/6Qk/8/8/8/8/K7&moves=6}\\{\sffamily\tiny\color{ink}scannen zum Lösen}
\end{minipage}
\vspace{2mm}

{\raggedright\noindent\textcolor{ink}{\sffamily\textsc{Themen}}\enspace \themetag{reines Matt} \themetag{Mustermatt} \themetag{Idealmatt} \themetag{Selbstblock}\par\smallskip{\footnotesize\textit{Hinweis:} die Partei am Zug steht im Diagramm im Schach; die Lösung beginnt mit einem Schlag.}\par}

\clearpage
\phantomsection\label{hm:KQRvkr}
\begin{tcolorbox}[enhanced,colback=ink,colframe=ink,arc=0pt,boxrule=0pt,left=4mm,right=4mm,top=2.5mm,bottom=2.5mm]
\color{white}\sffamily{\LARGE\bfseries KQRvkr}\hfill{\Large h\#6}\\[0.5mm]
{\small 5 Steine\quad Schwarz am Zug}\hfill{\small Nr.~166\qquad\href{https://helpman.komtera.lt/?fen=8/8/7Q/7k/2R5/1r6/8/K7&moves=6}{\color{white}\bfseries lösen\,\ensuremath{\nearrow}}}
\end{tcolorbox}
\vspace{2mm}
\noindent\begin{minipage}[t]{36mm}\vspace{0pt}\sffamily\small\raggedright
\hmfigure{längste eindeutige}{{\Large h\#6}}\hmfigure{eindeutig bei dieser Länge}{{\Large 6} Stellungen}\hmfigure{längstes Matt}{{\Large h\#7}}\hmfigure{in der Klasse}{242 Millionen Stellungen,\\ 135 Millionen mit Hilfsmatt}
\end{minipage}\hfill
\begin{minipage}[t]{80mm}\vspace{0pt}\begin{tikzpicture}\begin{axis}[width=72mm,height=42mm,ybar,bar width=1.6pt,xmin=-1,xmax=15,ymin=0,ymax=67.41,axis y line*=left,axis x line*=bottom,xlabel={\scriptsize Halbzüge bis Matt},ylabel={\scriptsize Stellungen (Mio.)},tick label style={font=\tiny},label style={font=\scriptsize},enlarge x limits=false,ylabel shift=-3pt,xlabel shift=-3pt]\addplot[fill=ink!55,draw=none] coordinates {(0,0.9348) (1,3.932) (2,15.69) (3,34.9) (4,61.28) (5,14.3) (6,2.493) (7,0.01011) (8,0.6431) (9,6.5e-05) (10,0.4191) (11,7e-06) (12,0.08068) (14,0.001943)};\draw[densely dotted] (axis cs:12,0) -- (axis cs:12,67.41);\end{axis}\begin{axis}[width=72mm,height=42mm,axis y line*=right,axis x line=none,ymode=log,xmin=-1,xmax=15,ylabel={\scriptsize eindeutig (log)},tick label style={font=\tiny},label style={font=\scriptsize},ylabel shift=-3pt]\addplot[mark=*,mark size=0.7pt,draw=black,thin] coordinates {(0,934838) (1,2851372) (2,6317369) (3,3405360) (4,1226778) (5,4489) (6,33012) (7,7) (8,6132) (10,1209) (12,6)};\end{axis}\end{tikzpicture}\end{minipage}
\vspace{3mm}

\noindent\begin{minipage}[t]{60mm}\vspace{0pt}\sbox\hmboard{\chessboard[setfen=8/8/7Q/7k/2R5/1r6/8/K7 b - - 0 1]}\usebox\hmboard\\[1mm]
\makebox[\wd\hmboard]{\sffamily\hspace{3mm}\textbf{h\#6}\hfill(3+2)}\end{minipage}\hfill
\begin{minipage}[t]{52mm}\vspace{0pt}
\textcolor{ink}{\sffamily\textsc{Lösung}}\par\smallskip\small
\begin{tabular}{@{}r@{\enspace}l@{\enspace}l@{}}
\textbf{1.} & \symking{}xh6 & \symking{}a2\\
\textbf{2.} & \symrook{}b7 & \symking{}a3\\
\textbf{3.} & \symrook{}f7 & \symking{}b4\\
\textbf{4.} & \symking{}g7 & \symking{}c5\\
\textbf{5.} & \symking{}f8 & \symking{}d6\\
\textbf{6.} & \symking{}e8 & \symrook{}c8\#\\
\end{tabular}\par\vspace{3mm}
\qrcode[height=12mm]{https://helpman.komtera.lt/?fen=8/8/7Q/7k/2R5/1r6/8/K7&moves=6}\\{\sffamily\tiny\color{ink}scannen zum Lösen}
\end{minipage}
\vspace{2mm}

{\raggedright\noindent\textcolor{ink}{\sffamily\textsc{Themen}}\enspace \themetag{reines Matt} \themetag{Mustermatt} \themetag{Idealmatt} \themetag{Selbstblock} \themetag{ohne Schach}\par\smallskip{\footnotesize\textit{Hinweis:} die Partei am Zug steht im Diagramm im Schach; die Lösung beginnt mit einem Schlag.}\par}

\clearpage
\phantomsection\label{hm:KQvkbn}
\begin{tcolorbox}[enhanced,colback=ink,colframe=ink,arc=0pt,boxrule=0pt,left=4mm,right=4mm,top=2.5mm,bottom=2.5mm]
\color{white}\sffamily{\LARGE\bfseries KQvkbn}\hfill{\Large h\#6}\\[0.5mm]
{\small 5 Steine\quad Schwarz am Zug}\hfill{\small Nr.~167\qquad\href{https://helpman.komtera.lt/?fen=7k/8/8/8/8/b7/7Q/K1n5&moves=6}{\color{white}\bfseries lösen\,\ensuremath{\nearrow}}}
\end{tcolorbox}
\vspace{2mm}
\noindent\begin{minipage}[t]{36mm}\vspace{0pt}\sffamily\small\raggedright
\hmfigure{längste eindeutige}{{\Large h\#6}}\hmfigure{eindeutig bei dieser Länge}{{\Large 1} Stellung}\hmfigure{längstes Matt}{{\Large h\#7}}\hmfigure{in der Klasse}{242 Millionen Stellungen,\\ 149 Millionen mit Hilfsmatt}
\end{minipage}\hfill
\begin{minipage}[t]{80mm}\vspace{0pt}\begin{tikzpicture}\begin{axis}[width=72mm,height=42mm,ybar,bar width=1.6pt,xmin=-1,xmax=15,ymin=0,ymax=34.3,axis y line*=left,axis x line*=bottom,xlabel={\scriptsize Halbzüge bis Matt},ylabel={\scriptsize Stellungen (Mio.)},tick label style={font=\tiny},label style={font=\scriptsize},enlarge x limits=false,ylabel shift=-3pt,xlabel shift=-3pt]\addplot[fill=ink!55,draw=none] coordinates {(0,0.0797) (1,0.6024) (2,1.769) (3,6.261) (4,13.71) (5,21.51) (6,31.19) (7,26.82) (8,26.91) (9,11.77) (10,6.581) (11,1.17) (12,0.6388) (13,0.001426) (14,0.000383)};\draw[densely dotted] (axis cs:12,0) -- (axis cs:12,34.3);\end{axis}\begin{axis}[width=72mm,height=42mm,axis y line*=right,axis x line=none,ymode=log,xmin=-1,xmax=15,ylabel={\scriptsize eindeutig (log)},tick label style={font=\tiny},label style={font=\scriptsize},ylabel shift=-3pt]\addplot[mark=*,mark size=0.7pt,draw=black,thin] coordinates {(0,79698) (1,449539) (2,757933) (3,1316431) (4,1581885) (5,1244297) (6,966619) (7,446633) (8,213035) (9,23226) (10,5041) (11,9) (12,1)};\end{axis}\end{tikzpicture}\end{minipage}
\vspace{3mm}

\noindent\begin{minipage}[t]{60mm}\vspace{0pt}\sbox\hmboard{\chessboard[setfen=7k/8/8/8/8/b7/7Q/K1n5 b - - 0 1]}\usebox\hmboard\\[1mm]
\makebox[\wd\hmboard]{\sffamily\hspace{3mm}\textbf{h\#6}\hfill(2+3)}\end{minipage}\hfill
\begin{minipage}[t]{52mm}\vspace{0pt}
\textcolor{ink}{\sffamily\textsc{Lösung}}\par\smallskip\small
\begin{tabular}{@{}r@{\enspace}l@{\enspace}l@{}}
\textbf{1.} & \symking{}g7 & \symking{}b1\\
\textbf{2.} & \symknight{}e2 & \symking{}c2\\
\textbf{3.} & \symknight{}g1 & \symking{}d3\\
\textbf{4.} & \symknight{}h3 & \symking{}e4\\
\textbf{5.} & \symking{}h6 & \symking{}f5\\
\textbf{6.} & \symking{}h5 & \symqueen{}xh3\#\\
\end{tabular}\par\vspace{3mm}
\qrcode[height=12mm]{https://helpman.komtera.lt/?fen=7k/8/8/8/8/b7/7Q/K1n5&moves=6}\\{\sffamily\tiny\color{ink}scannen zum Lösen}
\end{minipage}
\vspace{2mm}

{\raggedright\noindent\textcolor{ink}{\sffamily\textsc{Themen}}\enspace \themetag{Spiegelmatt} \themetag{ohne Schach}\par\smallskip{\footnotesize\textit{Hinweis:} die Partei am Zug steht im Diagramm im Schach.}\par}

\clearpage
\phantomsection\label{hm:KQvknn}
\begin{tcolorbox}[enhanced,colback=ink,colframe=ink,arc=0pt,boxrule=0pt,left=4mm,right=4mm,top=2.5mm,bottom=2.5mm]
\color{white}\sffamily{\LARGE\bfseries KQvknn}\hfill{\Large h\#6}\\[0.5mm]
{\small 5 Steine\quad Schwarz am Zug}\hfill{\small Nr.~168\qquad\href{https://helpman.komtera.lt/?fen=8/8/7k/6Q1/8/2n5/1n6/K7&moves=6}{\color{white}\bfseries lösen\,\ensuremath{\nearrow}}}
\end{tcolorbox}
\vspace{2mm}
\noindent\begin{minipage}[t]{36mm}\vspace{0pt}\sffamily\small\raggedright
\hmfigure{längste eindeutige}{{\Large h\#6}}\hmfigure{eindeutig bei dieser Länge}{{\Large 9} Stellungen}\hmfigure{längstes Matt}{{\Large h\#7}}\hmfigure{in der Klasse}{242 Millionen Stellungen,\\ 76,7 Millionen mit Hilfsmatt}
\end{minipage}\hfill
\begin{minipage}[t]{80mm}\vspace{0pt}\begin{tikzpicture}\begin{axis}[width=72mm,height=42mm,ybar,bar width=1.6pt,xmin=-1,xmax=15,ymin=0,ymax=18.14,axis y line*=left,axis x line*=bottom,xlabel={\scriptsize Halbzüge bis Matt},ylabel={\scriptsize Stellungen (Mio.)},tick label style={font=\tiny},label style={font=\scriptsize},enlarge x limits=false,ylabel shift=-3pt,xlabel shift=-3pt]\addplot[fill=ink!55,draw=none] coordinates {(0,0.04771) (1,0.3346) (2,0.8902) (3,2.938) (4,6.162) (5,10.41) (6,16.49) (7,13.72) (8,14.97) (9,6.041) (10,3.682) (11,0.6209) (12,0.3613) (13,0.001443) (14,0.000555)};\draw[densely dotted] (axis cs:12,0) -- (axis cs:12,18.14);\end{axis}\begin{axis}[width=72mm,height=42mm,axis y line*=right,axis x line=none,ymode=log,xmin=-1,xmax=15,ylabel={\scriptsize eindeutig (log)},tick label style={font=\tiny},label style={font=\scriptsize},ylabel shift=-3pt]\addplot[mark=*,mark size=0.7pt,draw=black,thin] coordinates {(0,47707) (1,231209) (2,349375) (3,675858) (4,831428) (5,672588) (6,588815) (7,228509) (8,122801) (9,11144) (10,2273) (11,12) (12,9)};\end{axis}\end{tikzpicture}\end{minipage}
\vspace{3mm}

\noindent\begin{minipage}[t]{60mm}\vspace{0pt}\sbox\hmboard{\chessboard[setfen=8/8/7k/6Q1/8/2n5/1n6/K7 b - - 0 1]}\usebox\hmboard\\[1mm]
\makebox[\wd\hmboard]{\sffamily\hspace{3mm}\textbf{h\#6}\hfill(2+3)}\end{minipage}\hfill
\begin{minipage}[t]{52mm}\vspace{0pt}
\textcolor{ink}{\sffamily\textsc{Lösung}}\par\smallskip\small
\begin{tabular}{@{}r@{\enspace}l@{\enspace}l@{}}
\textbf{1.} & \symking{}h7 & \symking{}xb2\\
\textbf{2.} & \symking{}h8 & \symking{}xc3\\
\textbf{3.} & \symking{}h7 & \symking{}d4\\
\textbf{4.} & \symking{}h8 & \symking{}e5\\
\textbf{5.} & \symking{}h7 & \symking{}f6\\
\textbf{6.} & \symking{}h8 & \symqueen{}g7\#\\
\end{tabular}\par\vspace{3mm}
\qrcode[height=12mm]{https://helpman.komtera.lt/?fen=8/8/7k/6Q1/8/2n5/1n6/K7&moves=6}\\{\sffamily\tiny\color{ink}scannen zum Lösen}
\end{minipage}
\vspace{2mm}

{\raggedright\noindent\textcolor{ink}{\sffamily\textsc{Themen}}\enspace \themetag{reines Matt} \themetag{Mustermatt} \themetag{Idealmatt} \themetag{Switchback} \themetag{ein Stein zieht} \themetag{ohne Schach}\par\smallskip{\footnotesize\textit{Hinweis:} die Partei am Zug steht im Diagramm im Schach.}\par}

\clearpage
\phantomsection\label{hm:KQvknp}
\begin{tcolorbox}[enhanced,colback=ink,colframe=ink,arc=0pt,boxrule=0pt,left=4mm,right=4mm,top=2.5mm,bottom=2.5mm]
\color{white}\sffamily{\LARGE\bfseries KQvknp}\hfill{\Large h\#6}\\[0.5mm]
{\small 5 Steine\quad Schwarz am Zug}\hfill{\small Nr.~169\qquad\href{https://helpman.komtera.lt/?fen=8/8/7k/6Q1/4p3/8/8/K2n4&moves=6}{\color{white}\bfseries lösen\,\ensuremath{\nearrow}}}
\end{tcolorbox}
\vspace{2mm}
\noindent\begin{minipage}[t]{36mm}\vspace{0pt}\sffamily\small\raggedright
\hmfigure{längste eindeutige}{{\Large h\#6}}\hmfigure{eindeutig bei dieser Länge}{{\Large 18} Stellungen}\hmfigure{längstes Matt}{{\Large h\#7}}\hmfigure{in der Klasse}{710 Millionen Stellungen,\\ 476 Millionen mit Hilfsmatt}
\end{minipage}\hfill
\begin{minipage}[t]{80mm}\vspace{0pt}\begin{tikzpicture}\begin{axis}[width=72mm,height=42mm,ybar,bar width=1.6pt,xmin=-1,xmax=15,ymin=0,ymax=120.7,axis y line*=left,axis x line*=bottom,xlabel={\scriptsize Halbzüge bis Matt},ylabel={\scriptsize Stellungen (Mio.)},tick label style={font=\tiny},label style={font=\scriptsize},enlarge x limits=false,ylabel shift=-3pt,xlabel shift=-3pt]\addplot[fill=ink!55,draw=none] coordinates {(0,0.3707) (1,2.526) (2,6.444) (3,20.28) (4,45.98) (5,66.48) (6,109.7) (7,81.86) (8,89.03) (9,31.23) (10,18.32) (11,2.838) (12,1.352) (13,0.001476) (14,0.000689)};\draw[densely dotted] (axis cs:12,0) -- (axis cs:12,120.7);\end{axis}\begin{axis}[width=72mm,height=42mm,axis y line*=right,axis x line=none,ymode=log,xmin=-1,xmax=15,ylabel={\scriptsize eindeutig (log)},tick label style={font=\tiny},label style={font=\scriptsize},ylabel shift=-3pt]\addplot[mark=*,mark size=0.7pt,draw=black,thin] coordinates {(0,370737) (1,1718985) (2,2684679) (3,4377568) (4,5651982) (5,3858410) (6,3269327) (7,1194164) (8,614234) (9,51360) (10,9043) (11,32) (12,18)};\end{axis}\end{tikzpicture}\end{minipage}
\vspace{3mm}

\noindent\begin{minipage}[t]{60mm}\vspace{0pt}\sbox\hmboard{\chessboard[setfen=8/8/7k/6Q1/4p3/8/8/K2n4 b - - 0 1]}\usebox\hmboard\\[1mm]
\makebox[\wd\hmboard]{\sffamily\hspace{3mm}\textbf{h\#6}\hfill(2+3)}\end{minipage}\hfill
\begin{minipage}[t]{52mm}\vspace{0pt}
\textcolor{ink}{\sffamily\textsc{Lösung}}\par\smallskip\small
\begin{tabular}{@{}r@{\enspace}l@{\enspace}l@{}}
\textbf{1.} & \symking{}h7 & \symking{}b1\\
\textbf{2.} & e3 & \symking{}c2\\
\textbf{3.} & e2 & \symking{}d3\\
\textbf{4.} & e1=\symbishop{} & \symking{}e4\\
\textbf{5.} & \symbishop{}c3 & \symking{}f5\\
\textbf{6.} & \symbishop{}h8 & \symqueen{}g6\#\\
\end{tabular}\par\vspace{3mm}
\qrcode[height=12mm]{https://helpman.komtera.lt/?fen=8/8/7k/6Q1/4p3/8/8/K2n4&moves=6}\\{\sffamily\tiny\color{ink}scannen zum Lösen}
\end{minipage}
\vspace{2mm}

{\raggedright\noindent\textcolor{ink}{\sffamily\textsc{Themen}}\enspace \themetag{reines Matt} \themetag{Mustermatt} \themetag{Umwandlung} \themetag{Unterverwandlung} \themetag{Selbstblock} \themetag{ohne Schlag} \themetag{ohne Schach}\par\smallskip{\footnotesize\textit{Hinweis:} die Partei am Zug steht im Diagramm im Schach.}\par}

\clearpage
\phantomsection\label{hm:KQvkpp}
\begin{tcolorbox}[enhanced,colback=ink,colframe=ink,arc=0pt,boxrule=0pt,left=4mm,right=4mm,top=2.5mm,bottom=2.5mm]
\color{white}\sffamily{\LARGE\bfseries KQvkpp}\hfill{\Large h\#6}\\[0.5mm]
{\small 5 Steine\quad Schwarz am Zug}\hfill{\small Nr.~170\qquad\href{https://helpman.komtera.lt/?fen=8/7k/5Q2/p4p2/8/8/8/K7&moves=6}{\color{white}\bfseries lösen\,\ensuremath{\nearrow}}}
\end{tcolorbox}
\vspace{2mm}
\noindent\begin{minipage}[t]{36mm}\vspace{0pt}\sffamily\small\raggedright
\hmfigure{längste eindeutige}{{\Large h\#6}}\hmfigure{eindeutig bei dieser Länge}{{\Large 15} Stellungen}\hmfigure{längstes Matt}{{\Large h\#7}}\hmfigure{in der Klasse}{533 Millionen Stellungen,\\ 184 Millionen mit Hilfsmatt}
\end{minipage}\hfill
\begin{minipage}[t]{80mm}\vspace{0pt}\begin{tikzpicture}\begin{axis}[width=72mm,height=42mm,ybar,bar width=1.6pt,xmin=-1,xmax=15,ymin=0,ymax=50.65,axis y line*=left,axis x line*=bottom,xlabel={\scriptsize Halbzüge bis Matt},ylabel={\scriptsize Stellungen (Mio.)},tick label style={font=\tiny},label style={font=\scriptsize},enlarge x limits=false,ylabel shift=-3pt,xlabel shift=-3pt]\addplot[fill=ink!55,draw=none] coordinates {(0,0.1809) (1,1.199) (2,2.93) (3,8.88) (4,22.02) (5,27.14) (6,46.05) (7,29.77) (8,30.83) (9,9.168) (10,4.853) (11,0.6742) (12,0.2464) (13,0.000434) (14,0.000399)};\draw[densely dotted] (axis cs:12,0) -- (axis cs:12,50.65);\end{axis}\begin{axis}[width=72mm,height=42mm,axis y line*=right,axis x line=none,ymode=log,xmin=-1,xmax=15,ylabel={\scriptsize eindeutig (log)},tick label style={font=\tiny},label style={font=\scriptsize},ylabel shift=-3pt]\addplot[mark=*,mark size=0.7pt,draw=black,thin] coordinates {(0,180884) (1,805158) (2,1421608) (3,1821938) (4,2519925) (5,1341529) (6,1004615) (7,339433) (8,148690) (9,7228) (10,1519) (11,11) (12,15)};\end{axis}\end{tikzpicture}\end{minipage}
\vspace{3mm}

\noindent\begin{minipage}[t]{60mm}\vspace{0pt}\sbox\hmboard{\chessboard[setfen=8/7k/5Q2/p4p2/8/8/8/K7 b - - 0 1]}\usebox\hmboard\\[1mm]
\makebox[\wd\hmboard]{\sffamily\hspace{3mm}\textbf{h\#6}\hfill(2+3)}\end{minipage}\hfill
\begin{minipage}[t]{52mm}\vspace{0pt}
\textcolor{ink}{\sffamily\textsc{Lösung}}\par\smallskip\small
\begin{tabular}{@{}r@{\enspace}l@{\enspace}l@{}}
\textbf{1.} & a4 & \symking{}b2\\
\textbf{2.} & a3+ & \symking{}c3\\
\textbf{3.} & a2 & \symking{}d4\\
\textbf{4.} & a1=\symrook{} & \symking{}e5\\
\textbf{5.} & \symrook{}a8 & \symking{}xf5\\
\textbf{6.} & \symrook{}h8 & \symqueen{}g6\#\\
\end{tabular}\par\vspace{3mm}
\qrcode[height=12mm]{https://helpman.komtera.lt/?fen=8/7k/5Q2/p4p2/8/8/8/K7&moves=6}\\{\sffamily\tiny\color{ink}scannen zum Lösen}
\end{minipage}
\vspace{2mm}

{\raggedright\noindent\textcolor{ink}{\sffamily\textsc{Themen}}\enspace \themetag{reines Matt} \themetag{Mustermatt} \themetag{Idealmatt} \themetag{Umwandlung} \themetag{Unterverwandlung} \themetag{Selbstblock} \themetag{ein Stein zieht}\par}

\clearpage
\phantomsection\label{hm:KQvkqn}
\begin{tcolorbox}[enhanced,colback=ink,colframe=ink,arc=0pt,boxrule=0pt,left=4mm,right=4mm,top=2.5mm,bottom=2.5mm]
\color{white}\sffamily{\LARGE\bfseries KQvkqn}\hfill{\Large h\#6}\\[0.5mm]
{\small 5 Steine\quad Schwarz am Zug}\hfill{\small Nr.~171\qquad\href{https://helpman.komtera.lt/?fen=8/8/7k/6Q1/1q1n4/8/8/K7&moves=6}{\color{white}\bfseries lösen\,\ensuremath{\nearrow}}}
\end{tcolorbox}
\vspace{2mm}
\noindent\begin{minipage}[t]{36mm}\vspace{0pt}\sffamily\small\raggedright
\hmfigure{längste eindeutige}{{\Large h\#6}}\hmfigure{eindeutig bei dieser Länge}{{\Large 2} Stellungen}\hmfigure{längstes Matt}{{\Large h\#7}}\hmfigure{in der Klasse}{242 Millionen Stellungen,\\ 130 Millionen mit Hilfsmatt}
\end{minipage}\hfill
\begin{minipage}[t]{80mm}\vspace{0pt}\begin{tikzpicture}\begin{axis}[width=72mm,height=42mm,ybar,bar width=1.6pt,xmin=-1,xmax=15,ymin=0,ymax=29.33,axis y line*=left,axis x line*=bottom,xlabel={\scriptsize Halbzüge bis Matt},ylabel={\scriptsize Stellungen (Mio.)},tick label style={font=\tiny},label style={font=\scriptsize},enlarge x limits=false,ylabel shift=-3pt,xlabel shift=-3pt]\addplot[fill=ink!55,draw=none] coordinates {(0,0.04065) (1,0.3553) (2,1.29) (3,5.32) (4,10.31) (5,20.01) (6,23.05) (7,26.66) (8,21.02) (9,14.04) (10,5.72) (11,1.613) (12,0.5627) (13,0.02519) (14,0.003324)};\draw[densely dotted] (axis cs:12,0) -- (axis cs:12,29.33);\end{axis}\begin{axis}[width=72mm,height=42mm,axis y line*=right,axis x line=none,ymode=log,xmin=-1,xmax=15,ylabel={\scriptsize eindeutig (log)},tick label style={font=\tiny},label style={font=\scriptsize},ylabel shift=-3pt]\addplot[mark=*,mark size=0.7pt,draw=black,thin] coordinates {(0,40654) (1,296770) (2,586192) (3,1185985) (4,1090120) (5,1079912) (6,664285) (7,466111) (8,168669) (9,42385) (10,4864) (11,147) (12,2)};\end{axis}\end{tikzpicture}\end{minipage}
\vspace{3mm}

\noindent\begin{minipage}[t]{60mm}\vspace{0pt}\sbox\hmboard{\chessboard[setfen=8/8/7k/6Q1/1q1n4/8/8/K7 b - - 0 1]}\usebox\hmboard\\[1mm]
\makebox[\wd\hmboard]{\sffamily\hspace{3mm}\textbf{h\#6}\hfill(2+3)}\end{minipage}\hfill
\begin{minipage}[t]{52mm}\vspace{0pt}
\textcolor{ink}{\sffamily\textsc{Lösung}}\par\smallskip\small
\begin{tabular}{@{}r@{\enspace}l@{\enspace}l@{}}
\textbf{1.} & \symking{}h7 & \symking{}a2\\
\textbf{2.} & \symqueen{}b7 & \symking{}a3\\
\textbf{3.} & \symqueen{}g7 & \symking{}b4\\
\textbf{4.} & \symking{}g8 & \symking{}c5\\
\textbf{5.} & \symking{}f8 & \symking{}d6\\
\textbf{6.} & \symqueen{}g8 & \symqueen{}e7\#\\
\end{tabular}\par\vspace{3mm}
\qrcode[height=12mm]{https://helpman.komtera.lt/?fen=8/8/7k/6Q1/1q1n4/8/8/K7&moves=6}\\{\sffamily\tiny\color{ink}scannen zum Lösen}
\end{minipage}
\vspace{2mm}

{\raggedright\noindent\textcolor{ink}{\sffamily\textsc{Themen}}\enspace \themetag{reines Matt} \themetag{Mustermatt} \themetag{Selbstblock} \themetag{ohne Schlag} \themetag{ohne Schach}\par\smallskip{\footnotesize\textit{Hinweis:} die Partei am Zug steht im Diagramm im Schach.}\par}

\clearpage
\phantomsection\label{hm:KQvkqp}
\begin{tcolorbox}[enhanced,colback=ink,colframe=ink,arc=0pt,boxrule=0pt,left=4mm,right=4mm,top=2.5mm,bottom=2.5mm]
\color{white}\sffamily{\LARGE\bfseries KQvkqp}\hfill{\Large h\#6}\\[0.5mm]
{\small 5 Steine\quad Schwarz am Zug}\hfill{\small Nr.~172\qquad\href{https://helpman.komtera.lt/?fen=K7/8/2p5/1q6/8/8/8/5Qk1&moves=6}{\color{white}\bfseries lösen\,\ensuremath{\nearrow}}}
\end{tcolorbox}
\vspace{2mm}
\noindent\begin{minipage}[t]{36mm}\vspace{0pt}\sffamily\small\raggedright
\hmfigure{längste eindeutige}{{\Large h\#6}}\hmfigure{eindeutig bei dieser Länge}{{\Large 4} Stellungen}\hmfigure{längstes Matt}{{\Large h\#7}}\hmfigure{in der Klasse}{710 Millionen Stellungen,\\ 402 Millionen mit Hilfsmatt}
\end{minipage}\hfill
\begin{minipage}[t]{80mm}\vspace{0pt}\begin{tikzpicture}\begin{axis}[width=72mm,height=42mm,ybar,bar width=1.6pt,xmin=-1,xmax=15,ymin=0,ymax=92.96,axis y line*=left,axis x line*=bottom,xlabel={\scriptsize Halbzüge bis Matt},ylabel={\scriptsize Stellungen (Mio.)},tick label style={font=\tiny},label style={font=\scriptsize},enlarge x limits=false,ylabel shift=-3pt,xlabel shift=-3pt]\addplot[fill=ink!55,draw=none] coordinates {(0,0.1581) (1,1.346) (2,4.865) (3,19.46) (4,42.22) (5,72.23) (6,84.51) (7,80.4) (8,55.27) (9,29.51) (10,9.422) (11,2.029) (12,0.4408) (13,0.01811) (14,0.001523)};\draw[densely dotted] (axis cs:12,0) -- (axis cs:12,92.96);\end{axis}\begin{axis}[width=72mm,height=42mm,axis y line*=right,axis x line=none,ymode=log,xmin=-1,xmax=15,ylabel={\scriptsize eindeutig (log)},tick label style={font=\tiny},label style={font=\scriptsize},ylabel shift=-3pt]\addplot[mark=*,mark size=0.7pt,draw=black,thin] coordinates {(0,158059) (1,1111108) (2,2350743) (3,4242336) (4,4223276) (5,3241134) (6,1677305) (7,1013248) (8,300292) (9,61463) (10,5495) (11,96) (12,4)};\end{axis}\end{tikzpicture}\end{minipage}
\vspace{3mm}

\noindent\begin{minipage}[t]{60mm}\vspace{0pt}\sbox\hmboard{\chessboard[setfen=K7/8/2p5/1q6/8/8/8/5Qk1 b - - 0 1]}\usebox\hmboard\\[1mm]
\makebox[\wd\hmboard]{\sffamily\hspace{3mm}\textbf{h\#6}\hfill(2+3)}\end{minipage}\hfill
\begin{minipage}[t]{52mm}\vspace{0pt}
\textcolor{ink}{\sffamily\textsc{Lösung}}\par\smallskip\small
\begin{tabular}{@{}r@{\enspace}l@{\enspace}l@{}}
\textbf{1.} & \symking{}h2 & \symqueen{}f7\\
\textbf{2.} & \symking{}g3 & \symqueen{}b7\\
\textbf{3.} & \symking{}f4 & \symking{}b8\\
\textbf{4.} & \symking{}e5 & \symking{}c8\\
\textbf{5.} & \symking{}d6 & \symking{}d8\\
\textbf{6.} & \symqueen{}d5 & \symqueen{}e7\#\\
\end{tabular}\par\vspace{3mm}
\qrcode[height=12mm]{https://helpman.komtera.lt/?fen=K7/8/2p5/1q6/8/8/8/5Qk1&moves=6}\\{\sffamily\tiny\color{ink}scannen zum Lösen}
\end{minipage}
\vspace{2mm}

{\raggedright\noindent\textcolor{ink}{\sffamily\textsc{Themen}}\enspace \themetag{Selbstblock} \themetag{ohne Schlag} \themetag{ohne Schach}\par\smallskip{\footnotesize\textit{Hinweis:} die Partei am Zug steht im Diagramm im Schach.}\par}

\clearpage
\phantomsection\label{hm:KRBNvk}
\begin{tcolorbox}[enhanced,colback=ink,colframe=ink,arc=0pt,boxrule=0pt,left=4mm,right=4mm,top=2.5mm,bottom=2.5mm]
\color{white}\sffamily{\LARGE\bfseries KRBNvk}\hfill{\Large h\#6}\\[0.5mm]
{\small 5 Steine\quad Schwarz am Zug}\hfill{\small Nr.~173\qquad\href{https://helpman.komtera.lt/?fen=8/6N1/8/8/8/2B5/2R5/K2k4&moves=6}{\color{white}\bfseries lösen\,\ensuremath{\nearrow}}}
\end{tcolorbox}
\vspace{2mm}
\noindent\begin{minipage}[t]{36mm}\vspace{0pt}\sffamily\small\raggedright
\hmfigure{längste eindeutige}{{\Large h\#6}}\hmfigure{eindeutig bei dieser Länge}{{\Large 36} Stellungen}\hmfigure{längstes Matt}{{\Large h\#8}}\hmfigure{in der Klasse}{242 Millionen Stellungen,\\ 167 Millionen mit Hilfsmatt}
\end{minipage}\hfill
\begin{minipage}[t]{80mm}\vspace{0pt}\begin{tikzpicture}\begin{axis}[width=72mm,height=42mm,ybar,bar width=1.6pt,xmin=-1,xmax=17,ymin=0,ymax=67.7,axis y line*=left,axis x line*=bottom,xlabel={\scriptsize Halbzüge bis Matt},ylabel={\scriptsize Stellungen (Mio.)},tick label style={font=\tiny},label style={font=\scriptsize},enlarge x limits=false,ylabel shift=-3pt,xlabel shift=-3pt]\addplot[fill=ink!55,draw=none] coordinates {(0,0.1907) (1,0.9427) (2,2.196) (3,8.578) (4,26.95) (5,41.14) (6,61.55) (7,14.34) (8,10.66) (9,0.399) (10,0.3571) (11,0.000366) (12,0.08727) (13,3e-06) (14,0.02053) (16,0.000128)};\draw[densely dotted] (axis cs:12,0) -- (axis cs:12,67.7);\end{axis}\begin{axis}[width=72mm,height=42mm,axis y line*=right,axis x line=none,ymode=log,xmin=-1,xmax=17,ylabel={\scriptsize eindeutig (log)},tick label style={font=\tiny},label style={font=\scriptsize},ylabel shift=-3pt]\addplot[mark=*,mark size=0.7pt,draw=black,thin] coordinates {(0,190700) (1,890133) (2,1940386) (3,4991471) (4,10949810) (5,3301307) (6,1532048) (7,18553) (8,10079) (9,8) (10,805) (12,36)};\end{axis}\end{tikzpicture}\end{minipage}
\vspace{3mm}

\noindent\begin{minipage}[t]{60mm}\vspace{0pt}\sbox\hmboard{\chessboard[setfen=8/6N1/8/8/8/2B5/2R5/K2k4 b - - 0 1]}\usebox\hmboard\\[1mm]
\makebox[\wd\hmboard]{\sffamily\hspace{3mm}\textbf{h\#6}\hfill(4+1)}\end{minipage}\hfill
\begin{minipage}[t]{52mm}\vspace{0pt}
\textcolor{ink}{\sffamily\textsc{Lösung}}\par\smallskip\small
\begin{tabular}{@{}r@{\enspace}l@{\enspace}l@{}}
\textbf{1.} & \symking{}xc2 & \symking{}a2\\
\textbf{2.} & \symking{}c1 & \symking{}b3\\
\textbf{3.} & \symking{}b1 & \symking{}c4\\
\textbf{4.} & \symking{}a2 & \symknight{}e6\\
\textbf{5.} & \symking{}a3 & \symbishop{}b4+\\
\textbf{6.} & \symking{}a4 & \symknight{}c5\#\\
\end{tabular}\par\vspace{3mm}
\qrcode[height=12mm]{https://helpman.komtera.lt/?fen=8/6N1/8/8/8/2B5/2R5/K2k4&moves=6}\\{\sffamily\tiny\color{ink}scannen zum Lösen}
\end{minipage}
\vspace{2mm}

{\raggedright\noindent\textcolor{ink}{\sffamily\textsc{Themen}}\enspace \themetag{ein Stein zieht}\par\smallskip{\footnotesize\textit{Hinweis:} die Lösung beginnt mit einem Schlag.}\par}

\clearpage
\phantomsection\label{hm:KRBvkb}
\begin{tcolorbox}[enhanced,colback=ink,colframe=ink,arc=0pt,boxrule=0pt,left=4mm,right=4mm,top=2.5mm,bottom=2.5mm]
\color{white}\sffamily{\LARGE\bfseries KRBvkb}\hfill{\Large h\#6}\\[0.5mm]
{\small 5 Steine\quad Schwarz am Zug}\hfill{\small Nr.~174\qquad\href{https://helpman.komtera.lt/?fen=7b/7k/8/4R3/8/8/8/K1B5&moves=6}{\color{white}\bfseries lösen\,\ensuremath{\nearrow}}}
\end{tcolorbox}
\vspace{2mm}
\noindent\begin{minipage}[t]{36mm}\vspace{0pt}\sffamily\small\raggedright
\hmfigure{längste eindeutige}{{\Large h\#6}}\hmfigure{eindeutig bei dieser Länge}{{\Large 28} Stellungen}\hmfigure{längstes Matt}{{\Large h\#10}}\hmfigure{in der Klasse}{242 Millionen Stellungen,\\ 160 Millionen mit Hilfsmatt}
\end{minipage}\hfill
\begin{minipage}[t]{80mm}\vspace{0pt}\begin{tikzpicture}\begin{axis}[width=72mm,height=42mm,ybar,bar width=1.6pt,xmin=-1,xmax=21,ymin=0,ymax=49.22,axis y line*=left,axis x line*=bottom,xlabel={\scriptsize Halbzüge bis Matt},ylabel={\scriptsize Stellungen (Mio.)},tick label style={font=\tiny},label style={font=\scriptsize},enlarge x limits=false,ylabel shift=-3pt,xlabel shift=-3pt]\addplot[fill=ink!55,draw=none] coordinates {(0,0.07724) (1,0.3977) (2,1.273) (3,5.529) (4,17.15) (5,30.16) (6,44.74) (7,25.5) (8,19.46) (9,8.382) (10,5.7) (11,1.095) (12,0.5742) (13,0.01679) (14,0.01235) (16,0.008625) (18,0.008375) (20,0.002074)};\draw[densely dotted] (axis cs:12,0) -- (axis cs:12,49.22);\end{axis}\begin{axis}[width=72mm,height=42mm,axis y line*=right,axis x line=none,ymode=log,xmin=-1,xmax=21,ylabel={\scriptsize eindeutig (log)},tick label style={font=\tiny},label style={font=\scriptsize},ylabel shift=-3pt]\addplot[mark=*,mark size=0.7pt,draw=black,thin] coordinates {(0,77239) (1,388044) (2,697663) (3,2115730) (4,5223778) (5,2725585) (6,1257796) (7,219476) (8,116115) (9,18339) (10,5878) (11,96) (12,28)};\end{axis}\end{tikzpicture}\end{minipage}
\vspace{3mm}

\noindent\begin{minipage}[t]{60mm}\vspace{0pt}\sbox\hmboard{\chessboard[setfen=7b/7k/8/4R3/8/8/8/K1B5 b - - 0 1]}\usebox\hmboard\\[1mm]
\makebox[\wd\hmboard]{\sffamily\hspace{3mm}\textbf{h\#6}\hfill(3+2)}\end{minipage}\hfill
\begin{minipage}[t]{52mm}\vspace{0pt}
\textcolor{ink}{\sffamily\textsc{Lösung}}\par\smallskip\small
\begin{tabular}{@{}r@{\enspace}l@{\enspace}l@{}}
\textbf{1.} & \symking{}g6 & \symbishop{}g5\\
\textbf{2.} & \symking{}h5 & \symking{}b1\\
\textbf{3.} & \symking{}g4 & \symrook{}e8\\
\textbf{4.} & \symking{}f3 & \symbishop{}e7\\
\textbf{5.} & \symking{}e2 & \symbishop{}b4+\\
\textbf{6.} & \symking{}d1 & \symrook{}e1\#\\
\end{tabular}\par\vspace{3mm}
\qrcode[height=12mm]{https://helpman.komtera.lt/?fen=7b/7k/8/4R3/8/8/8/K1B5&moves=6}\\{\sffamily\tiny\color{ink}scannen zum Lösen}
\end{minipage}
\vspace{2mm}

{\raggedright\noindent\textcolor{ink}{\sffamily\textsc{Themen}}\enspace \themetag{ein Stein zieht} \themetag{ohne Schlag} \themetag{Inder}\par}

\clearpage
\phantomsection\label{hm:KRBvkp}
\begin{tcolorbox}[enhanced,colback=ink,colframe=ink,arc=0pt,boxrule=0pt,left=4mm,right=4mm,top=2.5mm,bottom=2.5mm]
\color{white}\sffamily{\LARGE\bfseries KRBvkp}\hfill{\Large h\#6}\\[0.5mm]
{\small 5 Steine\quad Schwarz am Zug}\hfill{\small Nr.~175\qquad\href{https://helpman.komtera.lt/?fen=6R1/8/7k/7p/7B/8/8/K7&moves=6}{\color{white}\bfseries lösen\,\ensuremath{\nearrow}}}
\end{tcolorbox}
\vspace{2mm}
\noindent\begin{minipage}[t]{36mm}\vspace{0pt}\sffamily\small\raggedright
\hmfigure{längste eindeutige}{{\Large h\#6}}\hmfigure{eindeutig bei dieser Länge}{{\Large 175} Stellungen}\hmfigure{längstes Matt}{{\Large h\#16}}\hmfigure{in der Klasse}{710 Millionen Stellungen,\\ 513 Millionen mit Hilfsmatt}
\end{minipage}\hfill
\begin{minipage}[t]{80mm}\vspace{0pt}\begin{tikzpicture}\begin{axis}[width=72mm,height=42mm,ybar,bar width=1.6pt,xmin=-1,xmax=33,ymin=0,ymax=147.1,axis y line*=left,axis x line*=bottom,xlabel={\scriptsize Halbzüge bis Matt},ylabel={\scriptsize Stellungen (Mio.)},tick label style={font=\tiny},label style={font=\scriptsize},enlarge x limits=false,ylabel shift=-3pt,xlabel shift=-3pt]\addplot[fill=ink!55,draw=none] coordinates {(0,0.3739) (1,1.972) (2,4.513) (3,16.48) (4,43.87) (5,73.33) (6,133.7) (7,87.44) (8,94.67) (9,31.9) (10,20.13) (11,2.71) (12,1.253) (13,0.03801) (14,0.033) (15,2.6e-05) (16,0.02967) (17,3.7e-05) (18,0.03692) (19,2.6e-05) (20,0.03759) (21,5.8e-05) (22,0.04233) (23,0.000217) (24,0.04385) (25,1e-05) (26,0.02542) (27,3.1e-05) (28,0.02182) (29,9e-06) (30,0.009949) (31,1.4e-05) (32,0.005458)};\draw[densely dotted] (axis cs:12,0) -- (axis cs:12,147.1);\end{axis}\begin{axis}[width=72mm,height=42mm,axis y line*=right,axis x line=none,ymode=log,xmin=-1,xmax=33,ylabel={\scriptsize eindeutig (log)},tick label style={font=\tiny},label style={font=\scriptsize},ylabel shift=-3pt]\addplot[mark=*,mark size=0.7pt,draw=black,thin] coordinates {(0,373932) (1,1916741) (2,3324059) (3,7255559) (4,15609437) (5,9705271) (6,7190088) (7,938105) (8,467336) (9,37793) (10,13968) (11,312) (12,175)};\end{axis}\end{tikzpicture}\end{minipage}
\vspace{3mm}

\noindent\begin{minipage}[t]{60mm}\vspace{0pt}\sbox\hmboard{\chessboard[setfen=6R1/8/7k/7p/7B/8/8/K7 b - - 0 1]}\usebox\hmboard\\[1mm]
\makebox[\wd\hmboard]{\sffamily\hspace{3mm}\textbf{h\#6}\hfill(3+2)}\end{minipage}\hfill
\begin{minipage}[t]{52mm}\vspace{0pt}
\textcolor{ink}{\sffamily\textsc{Lösung}}\par\smallskip\small
\begin{tabular}{@{}r@{\enspace}l@{\enspace}l@{}}
\textbf{1.} & \symking{}h7 & \symking{}b2\\
\textbf{2.} & \symking{}h6 & \symking{}c3\\
\textbf{3.} & \symking{}h7 & \symking{}d4\\
\textbf{4.} & \symking{}h6 & \symking{}e5\\
\textbf{5.} & \symking{}h7 & \symking{}f6\\
\textbf{6.} & \symking{}h6 & \symrook{}h8\#\\
\end{tabular}\par\vspace{3mm}
\qrcode[height=12mm]{https://helpman.komtera.lt/?fen=6R1/8/7k/7p/7B/8/8/K7&moves=6}\\{\sffamily\tiny\color{ink}scannen zum Lösen}
\end{minipage}
\vspace{2mm}

{\raggedright\noindent\textcolor{ink}{\sffamily\textsc{Themen}}\enspace \themetag{Switchback} \themetag{ein Stein zieht} \themetag{Pendel} \themetag{ohne Schlag} \themetag{ohne Schach}\par}

\clearpage
\phantomsection\label{hm:KRNNvk}
\begin{tcolorbox}[enhanced,colback=ink,colframe=ink,arc=0pt,boxrule=0pt,left=4mm,right=4mm,top=2.5mm,bottom=2.5mm]
\color{white}\sffamily{\LARGE\bfseries KRNNvk}\hfill{\Large h\#6}\\[0.5mm]
{\small 5 Steine\quad Schwarz am Zug}\hfill{\small Nr.~176\qquad\href{https://helpman.komtera.lt/?fen=5k2/4R3/7N/8/6N1/8/2K5/8&moves=6}{\color{white}\bfseries lösen\,\ensuremath{\nearrow}}}
\end{tcolorbox}
\vspace{2mm}
\noindent\begin{minipage}[t]{36mm}\vspace{0pt}\sffamily\small\raggedright
\hmfigure{längste eindeutige}{{\Large h\#6}}\hmfigure{eindeutig bei dieser Länge}{{\Large 27} Stellungen}\hmfigure{längstes Matt}{{\Large h\#8}}\hmfigure{in der Klasse}{242 Millionen Stellungen,\\ 85,5 Millionen mit Hilfsmatt}
\end{minipage}\hfill
\begin{minipage}[t]{80mm}\vspace{0pt}\begin{tikzpicture}\begin{axis}[width=72mm,height=42mm,ybar,bar width=1.6pt,xmin=-1,xmax=17,ymin=0,ymax=33.26,axis y line*=left,axis x line*=bottom,xlabel={\scriptsize Halbzüge bis Matt},ylabel={\scriptsize Stellungen (Mio.)},tick label style={font=\tiny},label style={font=\scriptsize},enlarge x limits=false,ylabel shift=-3pt,xlabel shift=-3pt]\addplot[fill=ink!55,draw=none] coordinates {(0,0.07899) (1,0.4241) (2,0.9268) (3,3.77) (4,11.25) (5,19.56) (6,30.23) (7,10.04) (8,8.174) (9,0.6128) (10,0.3354) (11,0.000905) (12,0.03499) (14,0.01797) (16,4.7e-05)};\draw[densely dotted] (axis cs:12,0) -- (axis cs:12,33.26);\end{axis}\begin{axis}[width=72mm,height=42mm,axis y line*=right,axis x line=none,ymode=log,xmin=-1,xmax=17,ylabel={\scriptsize eindeutig (log)},tick label style={font=\tiny},label style={font=\scriptsize},ylabel shift=-3pt]\addplot[mark=*,mark size=0.7pt,draw=black,thin] coordinates {(0,78988) (1,399874) (2,819673) (3,2054855) (4,4261835) (5,1670934) (6,787974) (7,15583) (8,5980) (9,8) (10,269) (12,27)};\end{axis}\end{tikzpicture}\end{minipage}
\vspace{3mm}

\noindent\begin{minipage}[t]{60mm}\vspace{0pt}\sbox\hmboard{\chessboard[setfen=5k2/4R3/7N/8/6N1/8/2K5/8 b - - 0 1]}\usebox\hmboard\\[1mm]
\makebox[\wd\hmboard]{\sffamily\hspace{3mm}\textbf{h\#6}\hfill(4+1)}\end{minipage}\hfill
\begin{minipage}[t]{52mm}\vspace{0pt}
\textcolor{ink}{\sffamily\textsc{Lösung}}\par\smallskip\small
\begin{tabular}{@{}r@{\enspace}l@{\enspace}l@{}}
\textbf{1.} & \symking{}xe7 & \symknight{}f2\\
\textbf{2.} & \symking{}f6 & \symking{}d3\\
\textbf{3.} & \symking{}g5 & \symking{}e4\\
\textbf{4.} & \symking{}h4 & \symking{}f5\\
\textbf{5.} & \symking{}g3 & \symking{}g6\\
\textbf{6.} & \symking{}h4 & \symknight{}f5\#\\
\end{tabular}\par\vspace{3mm}
\qrcode[height=12mm]{https://helpman.komtera.lt/?fen=5k2/4R3/7N/8/6N1/8/2K5/8&moves=6}\\{\sffamily\tiny\color{ink}scannen zum Lösen}
\end{minipage}
\vspace{2mm}

{\raggedright\noindent\textcolor{ink}{\sffamily\textsc{Themen}}\enspace \themetag{reines Matt} \themetag{Mustermatt} \themetag{Idealmatt} \themetag{Spiegelmatt} \themetag{Switchback} \themetag{ein Stein zieht} \themetag{ohne Schach}\par\smallskip{\footnotesize\textit{Hinweis:} die Lösung beginnt mit einem Schlag.}\par}

\clearpage
\phantomsection\label{hm:KRNvkb}
\begin{tcolorbox}[enhanced,colback=ink,colframe=ink,arc=0pt,boxrule=0pt,left=4mm,right=4mm,top=2.5mm,bottom=2.5mm]
\color{white}\sffamily{\LARGE\bfseries KRNvkb}\hfill{\Large h\#6}\\[0.5mm]
{\small 5 Steine\quad Schwarz am Zug}\hfill{\small Nr.~177\qquad\href{https://helpman.komtera.lt/?fen=8/8/8/8/8/4N3/4R3/Kb2k3&moves=6}{\color{white}\bfseries lösen\,\ensuremath{\nearrow}}}
\end{tcolorbox}
\vspace{2mm}
\noindent\begin{minipage}[t]{36mm}\vspace{0pt}\sffamily\small\raggedright
\hmfigure{längste eindeutige}{{\Large h\#6}}\hmfigure{eindeutig bei dieser Länge}{{\Large 33} Stellungen}\hmfigure{längstes Matt}{{\Large h\#10}}\hmfigure{in der Klasse}{242 Millionen Stellungen,\\ 164 Millionen mit Hilfsmatt}
\end{minipage}\hfill
\begin{minipage}[t]{80mm}\vspace{0pt}\begin{tikzpicture}\begin{axis}[width=72mm,height=42mm,ybar,bar width=1.6pt,xmin=-1,xmax=21,ymin=0,ymax=56.32,axis y line*=left,axis x line*=bottom,xlabel={\scriptsize Halbzüge bis Matt},ylabel={\scriptsize Stellungen (Mio.)},tick label style={font=\tiny},label style={font=\scriptsize},enlarge x limits=false,ylabel shift=-3pt,xlabel shift=-3pt]\addplot[fill=ink!55,draw=none] coordinates {(0,0.0618) (1,0.3733) (2,1.148) (3,5.192) (4,14.21) (5,31.22) (6,51.2) (7,34.72) (8,21.82) (9,3.411) (10,0.5582) (11,0.008298) (12,0.01489) (14,0.0145) (16,0.01422) (18,0.01276) (20,0.002839)};\draw[densely dotted] (axis cs:12,0) -- (axis cs:12,56.32);\end{axis}\begin{axis}[width=72mm,height=42mm,axis y line*=right,axis x line=none,ymode=log,xmin=-1,xmax=21,ylabel={\scriptsize eindeutig (log)},tick label style={font=\tiny},label style={font=\scriptsize},ylabel shift=-3pt]\addplot[mark=*,mark size=0.7pt,draw=black,thin] coordinates {(0,61801) (1,367676) (2,606429) (3,1846171) (4,4218757) (5,3111731) (6,1654393) (7,173845) (8,20413) (9,233) (10,70) (12,33)};\end{axis}\end{tikzpicture}\end{minipage}
\vspace{3mm}

\noindent\begin{minipage}[t]{60mm}\vspace{0pt}\sbox\hmboard{\chessboard[setfen=8/8/8/8/8/4N3/4R3/Kb2k3 b - - 0 1]}\usebox\hmboard\\[1mm]
\makebox[\wd\hmboard]{\sffamily\hspace{3mm}\textbf{h\#6}\hfill(3+2)}\end{minipage}\hfill
\begin{minipage}[t]{52mm}\vspace{0pt}
\textcolor{ink}{\sffamily\textsc{Lösung}}\par\smallskip\small
\begin{tabular}{@{}r@{\enspace}l@{\enspace}l@{}}
\textbf{1.} & \symking{}xe2 & \symking{}b2\\
\textbf{2.} & \symking{}d3 & \symking{}c1\\
\textbf{3.} & \symking{}c3 & \symking{}d1\\
\textbf{4.} & \symking{}b2 & \symking{}d2\\
\textbf{5.} & \symking{}a1 & \symking{}c1\\
\textbf{6.} & \symbishop{}a2 & \symknight{}c2\#\\
\end{tabular}\par\vspace{3mm}
\qrcode[height=12mm]{https://helpman.komtera.lt/?fen=8/8/8/8/8/4N3/4R3/Kb2k3&moves=6}\\{\sffamily\tiny\color{ink}scannen zum Lösen}
\end{minipage}
\vspace{2mm}

{\raggedright\noindent\textcolor{ink}{\sffamily\textsc{Themen}}\enspace \themetag{reines Matt} \themetag{Mustermatt} \themetag{Idealmatt} \themetag{Rundlauf} \themetag{Selbstblock} \themetag{ohne Schach}\par\smallskip{\footnotesize\textit{Hinweis:} die Partei am Zug steht im Diagramm im Schach; die Lösung beginnt mit einem Schlag.}\par}

\clearpage
\phantomsection\label{hm:KRRBvk}
\begin{tcolorbox}[enhanced,colback=ink,colframe=ink,arc=0pt,boxrule=0pt,left=4mm,right=4mm,top=2.5mm,bottom=2.5mm]
\color{white}\sffamily{\LARGE\bfseries KRRBvk}\hfill{\Large h\#6}\\[0.5mm]
{\small 5 Steine\quad Schwarz am Zug}\hfill{\small Nr.~178\qquad\href{https://helpman.komtera.lt/?fen=8/7R/7k/8/7B/6R1/8/K7&moves=6}{\color{white}\bfseries lösen\,\ensuremath{\nearrow}}}
\end{tcolorbox}
\vspace{2mm}
\noindent\begin{minipage}[t]{36mm}\vspace{0pt}\sffamily\small\raggedright
\hmfigure{längste eindeutige}{{\Large h\#6}}\hmfigure{eindeutig bei dieser Länge}{{\Large 17} Stellungen}\hmfigure{längstes Matt}{{\Large h\#7}}\hmfigure{in der Klasse}{242 Millionen Stellungen,\\ 78,7 Millionen mit Hilfsmatt}
\end{minipage}\hfill
\begin{minipage}[t]{80mm}\vspace{0pt}\begin{tikzpicture}\begin{axis}[width=72mm,height=42mm,ybar,bar width=1.6pt,xmin=-1,xmax=15,ymin=0,ymax=38.21,axis y line*=left,axis x line*=bottom,xlabel={\scriptsize Halbzüge bis Matt},ylabel={\scriptsize Stellungen (Mio.)},tick label style={font=\tiny},label style={font=\scriptsize},enlarge x limits=false,ylabel shift=-3pt,xlabel shift=-3pt]\addplot[fill=ink!55,draw=none] coordinates {(0,0.604) (1,2.272) (2,6.372) (3,15.63) (4,34.74) (5,9.725) (6,8.608) (7,0.1938) (8,0.3391) (9,4.6e-05) (10,0.1461) (11,4e-06) (12,0.02612) (14,0.001086)};\draw[densely dotted] (axis cs:12,0) -- (axis cs:12,38.21);\end{axis}\begin{axis}[width=72mm,height=42mm,axis y line*=right,axis x line=none,ymode=log,xmin=-1,xmax=15,ylabel={\scriptsize eindeutig (log)},tick label style={font=\tiny},label style={font=\scriptsize},ylabel shift=-3pt]\addplot[mark=*,mark size=0.7pt,draw=black,thin] coordinates {(0,604000) (1,1882303) (2,3895052) (3,2483972) (4,2574235) (5,48415) (6,42979) (7,3) (8,3806) (10,259) (12,17)};\end{axis}\end{tikzpicture}\end{minipage}
\vspace{3mm}

\noindent\begin{minipage}[t]{60mm}\vspace{0pt}\sbox\hmboard{\chessboard[setfen=8/7R/7k/8/7B/6R1/8/K7 b - - 0 1]}\usebox\hmboard\\[1mm]
\makebox[\wd\hmboard]{\sffamily\hspace{3mm}\textbf{h\#6}\hfill(4+1)}\end{minipage}\hfill
\begin{minipage}[t]{52mm}\vspace{0pt}
\textcolor{ink}{\sffamily\textsc{Lösung}}\par\smallskip\small
\begin{tabular}{@{}r@{\enspace}l@{\enspace}l@{}}
\textbf{1.} & \symking{}xh7 & \symking{}b2\\
\textbf{2.} & \symking{}h6 & \symking{}c3\\
\textbf{3.} & \symking{}h5 & \symking{}d4\\
\textbf{4.} & \symking{}xh4 & \symking{}e5\\
\textbf{5.} & \symking{}h5 & \symking{}f6\\
\textbf{6.} & \symking{}h6 & \symrook{}h3\#\\
\end{tabular}\par\vspace{3mm}
\qrcode[height=12mm]{https://helpman.komtera.lt/?fen=8/7R/7k/8/7B/6R1/8/K7&moves=6}\\{\sffamily\tiny\color{ink}scannen zum Lösen}
\end{minipage}
\vspace{2mm}

{\raggedright\noindent\textcolor{ink}{\sffamily\textsc{Themen}}\enspace \themetag{reines Matt} \themetag{Mustermatt} \themetag{Idealmatt} \themetag{Spiegelmatt} \themetag{Switchback} \themetag{ein Stein zieht} \themetag{ohne Schach}\par\smallskip{\footnotesize\textit{Hinweis:} die Partei am Zug steht im Diagramm im Schach; die Lösung beginnt mit einem Schlag.}\par}

\clearpage
\phantomsection\label{hm:KRRPvk}
\begin{tcolorbox}[enhanced,colback=ink,colframe=ink,arc=0pt,boxrule=0pt,left=4mm,right=4mm,top=2.5mm,bottom=2.5mm]
\color{white}\sffamily{\LARGE\bfseries KRRPvk}\hfill{\Large h\#6}\\[0.5mm]
{\small 5 Steine\quad Schwarz am Zug}\hfill{\small Nr.~179\qquad\href{https://helpman.komtera.lt/?fen=7k/7R/6R1/8/8/8/5P2/K7&moves=6}{\color{white}\bfseries lösen\,\ensuremath{\nearrow}}}
\end{tcolorbox}
\vspace{2mm}
\noindent\begin{minipage}[t]{36mm}\vspace{0pt}\sffamily\small\raggedright
\hmfigure{längste eindeutige}{{\Large h\#6}}\hmfigure{eindeutig bei dieser Länge}{{\Large 5} Stellungen}\hmfigure{längstes Matt}{{\Large h\#7}}\hmfigure{in der Klasse}{710 Millionen Stellungen,\\ 246 Millionen mit Hilfsmatt}
\end{minipage}\hfill
\begin{minipage}[t]{80mm}\vspace{0pt}\begin{tikzpicture}\begin{axis}[width=72mm,height=42mm,ybar,bar width=1.6pt,xmin=-1,xmax=15,ymin=0,ymax=115.3,axis y line*=left,axis x line*=bottom,xlabel={\scriptsize Halbzüge bis Matt},ylabel={\scriptsize Stellungen (Mio.)},tick label style={font=\tiny},label style={font=\scriptsize},enlarge x limits=false,ylabel shift=-3pt,xlabel shift=-3pt]\addplot[fill=ink!55,draw=none] coordinates {(0,1.547) (1,6.801) (2,18.62) (3,51.94) (4,104.8) (5,33.55) (6,27.16) (7,0.8676) (8,0.7937) (9,0.000166) (10,0.2613) (11,1.1e-05) (12,0.03451) (13,2e-06) (14,0.000508)};\draw[densely dotted] (axis cs:12,0) -- (axis cs:12,115.3);\end{axis}\begin{axis}[width=72mm,height=42mm,axis y line*=right,axis x line=none,ymode=log,xmin=-1,xmax=15,ylabel={\scriptsize eindeutig (log)},tick label style={font=\tiny},label style={font=\scriptsize},ylabel shift=-3pt]\addplot[mark=*,mark size=0.7pt,draw=black,thin] coordinates {(0,1547276) (1,5535741) (2,10988512) (3,6512987) (4,6013092) (5,111827) (6,71990) (7,4) (8,3823) (10,263) (12,5)};\end{axis}\end{tikzpicture}\end{minipage}
\vspace{3mm}

\noindent\begin{minipage}[t]{60mm}\vspace{0pt}\sbox\hmboard{\chessboard[setfen=7k/7R/6R1/8/8/8/5P2/K7 b - - 0 1]}\usebox\hmboard\\[1mm]
\makebox[\wd\hmboard]{\sffamily\hspace{3mm}\textbf{h\#6}\hfill(4+1)}\end{minipage}\hfill
\begin{minipage}[t]{52mm}\vspace{0pt}
\textcolor{ink}{\sffamily\textsc{Lösung}}\par\smallskip\small
\begin{tabular}{@{}r@{\enspace}l@{\enspace}l@{}}
\textbf{1.} & \symking{}xh7 & f4\\
\textbf{2.} & \symking{}h8 & f5\\
\textbf{3.} & \symking{}h7 & f6\\
\textbf{4.} & \symking{}h8 & f7\\
\textbf{5.} & \symking{}h7 & \symrook{}g7+\\
\textbf{6.} & \symking{}h8 & f8=\symqueen{}\#\\
\end{tabular}\par\vspace{3mm}
\qrcode[height=12mm]{https://helpman.komtera.lt/?fen=7k/7R/6R1/8/8/8/5P2/K7&moves=6}\\{\sffamily\tiny\color{ink}scannen zum Lösen}
\end{minipage}
\vspace{2mm}

{\raggedright\noindent\textcolor{ink}{\sffamily\textsc{Themen}}\enspace \themetag{Umwandlung} \themetag{Excelsior} \themetag{Switchback} \themetag{ein Stein zieht} \themetag{Pendel}\par\smallskip{\footnotesize\textit{Hinweis:} die Partei am Zug steht im Diagramm im Schach; die Lösung beginnt mit einem Schlag.}\par}

\clearpage
\phantomsection\label{hm:KRRvkp}
\begin{tcolorbox}[enhanced,colback=ink,colframe=ink,arc=0pt,boxrule=0pt,left=4mm,right=4mm,top=2.5mm,bottom=2.5mm]
\color{white}\sffamily{\LARGE\bfseries KRRvkp}\hfill{\Large h\#6}\\[0.5mm]
{\small 5 Steine\quad Schwarz am Zug}\hfill{\small Nr.~180\qquad\href{https://helpman.komtera.lt/?fen=K7/8/8/8/5R1p/7k/6R1/8&moves=6}{\color{white}\bfseries lösen\,\ensuremath{\nearrow}}}
\end{tcolorbox}
\vspace{2mm}
\noindent\begin{minipage}[t]{36mm}\vspace{0pt}\sffamily\small\raggedright
\hmfigure{längste eindeutige}{{\Large h\#6}}\hmfigure{eindeutig bei dieser Länge}{{\Large 43} Stellungen}\hmfigure{längstes Matt}{{\Large h\#7}}\hmfigure{in der Klasse}{710 Millionen Stellungen,\\ 245 Millionen mit Hilfsmatt}
\end{minipage}\hfill
\begin{minipage}[t]{80mm}\vspace{0pt}\begin{tikzpicture}\begin{axis}[width=72mm,height=42mm,ybar,bar width=1.6pt,xmin=-1,xmax=15,ymin=0,ymax=112.3,axis y line*=left,axis x line*=bottom,xlabel={\scriptsize Halbzüge bis Matt},ylabel={\scriptsize Stellungen (Mio.)},tick label style={font=\tiny},label style={font=\scriptsize},enlarge x limits=false,ylabel shift=-3pt,xlabel shift=-3pt]\addplot[fill=ink!55,draw=none] coordinates {(0,1.459) (1,6.296) (2,18.78) (3,54.11) (4,102.1) (5,34.25) (6,26.02) (7,1.101) (8,0.6554) (9,0.000725) (10,0.2907) (11,0.000344) (12,0.06331) (13,3.9e-05) (14,0.001562)};\draw[densely dotted] (axis cs:12,0) -- (axis cs:12,112.3);\end{axis}\begin{axis}[width=72mm,height=42mm,axis y line*=right,axis x line=none,ymode=log,xmin=-1,xmax=15,ylabel={\scriptsize eindeutig (log)},tick label style={font=\tiny},label style={font=\scriptsize},ylabel shift=-3pt]\addplot[mark=*,mark size=0.7pt,draw=black,thin] coordinates {(0,1458510) (1,5467711) (2,9241037) (3,7028962) (4,5049611) (5,113612) (6,43538) (7,26) (8,5020) (10,377) (12,43)};\end{axis}\end{tikzpicture}\end{minipage}
\vspace{3mm}

\noindent\begin{minipage}[t]{60mm}\vspace{0pt}\sbox\hmboard{\chessboard[setfen=K7/8/8/8/5R1p/7k/6R1/8 b - - 0 1]}\usebox\hmboard\\[1mm]
\makebox[\wd\hmboard]{\sffamily\hspace{3mm}\textbf{h\#6}\hfill(3+2)}\end{minipage}\hfill
\begin{minipage}[t]{52mm}\vspace{0pt}
\textcolor{ink}{\sffamily\textsc{Lösung}}\par\smallskip\small
\begin{tabular}{@{}r@{\enspace}l@{\enspace}l@{}}
\textbf{1.} & \symking{}xg2 & \symking{}b7\\
\textbf{2.} & h3 & \symking{}c6\\
\textbf{3.} & h2 & \symking{}d5\\
\textbf{4.} & h1=\symrook{} & \symking{}e4\\
\textbf{5.} & \symking{}h2 & \symking{}f3\\
\textbf{6.} & \symrook{}g1 & \symrook{}h4\#\\
\end{tabular}\par\vspace{3mm}
\qrcode[height=12mm]{https://helpman.komtera.lt/?fen=K7/8/8/8/5R1p/7k/6R1/8&moves=6}\\{\sffamily\tiny\color{ink}scannen zum Lösen}
\end{minipage}
\vspace{2mm}

{\raggedright\noindent\textcolor{ink}{\sffamily\textsc{Themen}}\enspace \themetag{reines Matt} \themetag{Mustermatt} \themetag{Idealmatt} \themetag{Umwandlung} \themetag{Unterverwandlung} \themetag{Selbstblock} \themetag{ohne Schach}\par\smallskip{\footnotesize\textit{Hinweis:} die Lösung beginnt mit einem Schlag.}\par}

\clearpage
\phantomsection\label{hm:KRRvkr}
\begin{tcolorbox}[enhanced,colback=ink,colframe=ink,arc=0pt,boxrule=0pt,left=4mm,right=4mm,top=2.5mm,bottom=2.5mm]
\color{white}\sffamily{\LARGE\bfseries KRRvkr}\hfill{\Large h\#6}\\[0.5mm]
{\small 5 Steine\quad Schwarz am Zug}\hfill{\small Nr.~181\qquad\href{https://helpman.komtera.lt/?fen=4kR2/7R/8/8/8/8/6r1/1K6&moves=6}{\color{white}\bfseries lösen\,\ensuremath{\nearrow}}}
\end{tcolorbox}
\vspace{2mm}
\noindent\begin{minipage}[t]{36mm}\vspace{0pt}\sffamily\small\raggedright
\hmfigure{längste eindeutige}{{\Large h\#6}}\hmfigure{eindeutig bei dieser Länge}{{\Large 1} Stellung}\hmfigure{längstes Matt}{{\Large h\#7}}\hmfigure{in der Klasse}{242 Millionen Stellungen,\\ 72,6 Millionen mit Hilfsmatt}
\end{minipage}\hfill
\begin{minipage}[t]{80mm}\vspace{0pt}\begin{tikzpicture}\begin{axis}[width=72mm,height=42mm,ybar,bar width=1.6pt,xmin=-1,xmax=15,ymin=0,ymax=30.93,axis y line*=left,axis x line*=bottom,xlabel={\scriptsize Halbzüge bis Matt},ylabel={\scriptsize Stellungen (Mio.)},tick label style={font=\tiny},label style={font=\scriptsize},enlarge x limits=false,ylabel shift=-3pt,xlabel shift=-3pt]\addplot[fill=ink!55,draw=none] coordinates {(0,0.2439) (1,1.142) (2,4.528) (3,15.63) (4,28.12) (5,14.33) (6,7.608) (7,0.6982) (8,0.1679) (9,0.000519) (10,0.08829) (11,2.4e-05) (12,0.0205) (14,0.000349)};\draw[densely dotted] (axis cs:12,0) -- (axis cs:12,30.93);\end{axis}\begin{axis}[width=72mm,height=42mm,axis y line*=right,axis x line=none,ymode=log,xmin=-1,xmax=15,ylabel={\scriptsize eindeutig (log)},tick label style={font=\tiny},label style={font=\scriptsize},ylabel shift=-3pt]\addplot[mark=*,mark size=0.7pt,draw=black,thin] coordinates {(0,243885) (1,1058537) (2,1980204) (3,2699164) (4,1266828) (5,75606) (6,12092) (7,44) (8,1569) (10,392) (12,1)};\end{axis}\end{tikzpicture}\end{minipage}
\vspace{3mm}

\noindent\begin{minipage}[t]{60mm}\vspace{0pt}\sbox\hmboard{\chessboard[setfen=4kR2/7R/8/8/8/8/6r1/1K6 b - - 0 1]}\usebox\hmboard\\[1mm]
\makebox[\wd\hmboard]{\sffamily\hspace{3mm}\textbf{h\#6}\hfill(3+2)}\end{minipage}\hfill
\begin{minipage}[t]{52mm}\vspace{0pt}
\textcolor{ink}{\sffamily\textsc{Lösung}}\par\smallskip\small
\begin{tabular}{@{}r@{\enspace}l@{\enspace}l@{}}
\textbf{1.} & \symking{}xf8 & \symking{}c1\\
\textbf{2.} & \symrook{}g7 & \symking{}d2\\
\textbf{3.} & \symking{}f7 & \symking{}e3\\
\textbf{4.} & \symking{}g6 & \symrook{}xg7+\\
\textbf{5.} & \symking{}h5 & \symking{}f4\\
\textbf{6.} & \symking{}h4 & \symrook{}h7\#\\
\end{tabular}\par\vspace{3mm}
\qrcode[height=12mm]{https://helpman.komtera.lt/?fen=4kR2/7R/8/8/8/8/6r1/1K6&moves=6}\\{\sffamily\tiny\color{ink}scannen zum Lösen}
\end{minipage}
\vspace{2mm}

{\raggedright\noindent\textcolor{ink}{\sffamily\textsc{Themen}}\enspace \themetag{reines Matt} \themetag{Mustermatt} \themetag{Idealmatt} \themetag{Spiegelmatt} \themetag{Switchback}\par\smallskip{\footnotesize\textit{Hinweis:} die Partei am Zug steht im Diagramm im Schach; die Lösung beginnt mit einem Schlag.}\par}

\clearpage
\phantomsection\label{hm:KRvkbb}
\begin{tcolorbox}[enhanced,colback=ink,colframe=ink,arc=0pt,boxrule=0pt,left=4mm,right=4mm,top=2.5mm,bottom=2.5mm]
\color{white}\sffamily{\LARGE\bfseries KRvkbb}\hfill{\Large h\#6}\\[0.5mm]
{\small 5 Steine\quad Schwarz am Zug}\hfill{\small Nr.~182\qquad\href{https://helpman.komtera.lt/?fen=8/6R1/6k1/8/8/8/2b5/K1b5&moves=6}{\color{white}\bfseries lösen\,\ensuremath{\nearrow}}}
\end{tcolorbox}
\vspace{2mm}
\noindent\begin{minipage}[t]{36mm}\vspace{0pt}\sffamily\small\raggedright
\hmfigure{längste eindeutige}{{\Large h\#6}}\hmfigure{eindeutig bei dieser Länge}{{\Large 3} Stellungen}\hmfigure{längstes Matt}{{\Large h\#7}}\hmfigure{in der Klasse}{242 Millionen Stellungen,\\ 79,1 Millionen mit Hilfsmatt}
\end{minipage}\hfill
\begin{minipage}[t]{80mm}\vspace{0pt}\begin{tikzpicture}\begin{axis}[width=72mm,height=42mm,ybar,bar width=1.6pt,xmin=-1,xmax=15,ymin=0,ymax=16.56,axis y line*=left,axis x line*=bottom,xlabel={\scriptsize Halbzüge bis Matt},ylabel={\scriptsize Stellungen (Mio.)},tick label style={font=\tiny},label style={font=\scriptsize},enlarge x limits=false,ylabel shift=-3pt,xlabel shift=-3pt]\addplot[fill=ink!55,draw=none] coordinates {(0,0.0131) (1,0.09014) (2,0.3182) (3,1.681) (4,3.398) (5,8.532) (6,11.14) (7,14.93) (8,15.05) (9,12.83) (10,7.526) (11,2.671) (12,0.9187) (13,0.03052) (14,0.008727)};\draw[densely dotted] (axis cs:12,0) -- (axis cs:12,16.56);\end{axis}\begin{axis}[width=72mm,height=42mm,axis y line*=right,axis x line=none,ymode=log,xmin=-1,xmax=15,ylabel={\scriptsize eindeutig (log)},tick label style={font=\tiny},label style={font=\scriptsize},ylabel shift=-3pt]\addplot[mark=*,mark size=0.7pt,draw=black,thin] coordinates {(0,13099) (1,90139) (2,138047) (3,501267) (4,690794) (5,703047) (6,386537) (7,187453) (8,115052) (9,39160) (10,10665) (11,177) (12,3)};\end{axis}\end{tikzpicture}\end{minipage}
\vspace{3mm}

\noindent\begin{minipage}[t]{60mm}\vspace{0pt}\sbox\hmboard{\chessboard[setfen=8/6R1/6k1/8/8/8/2b5/K1b5 b - - 0 1]}\usebox\hmboard\\[1mm]
\makebox[\wd\hmboard]{\sffamily\hspace{3mm}\textbf{h\#6}\hfill(2+3)}\end{minipage}\hfill
\begin{minipage}[t]{52mm}\vspace{0pt}
\textcolor{ink}{\sffamily\textsc{Lösung}}\par\smallskip\small
\begin{tabular}{@{}r@{\enspace}l@{\enspace}l@{}}
\textbf{1.} & \symking{}f6 & \symking{}a2\\
\textbf{2.} & \symbishop{}g6 & \symking{}b3\\
\textbf{3.} & \symbishop{}f7+ & \symking{}b4\\
\textbf{4.} & \symking{}e7 & \symking{}c5\\
\textbf{5.} & \symking{}d8 & \symking{}d6\\
\textbf{6.} & \symbishop{}g8 & \symrook{}xg8\#\\
\end{tabular}\par\vspace{3mm}
\qrcode[height=12mm]{https://helpman.komtera.lt/?fen=8/6R1/6k1/8/8/8/2b5/K1b5&moves=6}\\{\sffamily\tiny\color{ink}scannen zum Lösen}
\end{minipage}
\vspace{2mm}

{\raggedright\noindent\textcolor{ink}{\sffamily\textsc{Themen}}\enspace \themetag{reines Matt} \themetag{Mustermatt} \themetag{Spiegelmatt}\par\smallskip{\footnotesize\textit{Hinweis:} die Partei am Zug steht im Diagramm im Schach.}\par}

\clearpage
\phantomsection\label{hm:KRvkbn}
\begin{tcolorbox}[enhanced,colback=ink,colframe=ink,arc=0pt,boxrule=0pt,left=4mm,right=4mm,top=2.5mm,bottom=2.5mm]
\color{white}\sffamily{\LARGE\bfseries KRvkbn}\hfill{\Large h\#6}\\[0.5mm]
{\small 5 Steine\quad Schwarz am Zug}\hfill{\small Nr.~183\qquad\href{https://helpman.komtera.lt/?fen=8/8/7k/8/2b5/8/7R/K2n4&moves=6}{\color{white}\bfseries lösen\,\ensuremath{\nearrow}}}
\end{tcolorbox}
\vspace{2mm}
\noindent\begin{minipage}[t]{36mm}\vspace{0pt}\sffamily\small\raggedright
\hmfigure{längste eindeutige}{{\Large h\#6}}\hmfigure{eindeutig bei dieser Länge}{{\Large 12} Stellungen}\hmfigure{längstes Matt}{{\Large h\#7}}\hmfigure{in der Klasse}{242 Millionen Stellungen,\\ 163 Millionen mit Hilfsmatt}
\end{minipage}\hfill
\begin{minipage}[t]{80mm}\vspace{0pt}\begin{tikzpicture}\begin{axis}[width=72mm,height=42mm,ybar,bar width=1.6pt,xmin=-1,xmax=15,ymin=0,ymax=34.71,axis y line*=left,axis x line*=bottom,xlabel={\scriptsize Halbzüge bis Matt},ylabel={\scriptsize Stellungen (Mio.)},tick label style={font=\tiny},label style={font=\scriptsize},enlarge x limits=false,ylabel shift=-3pt,xlabel shift=-3pt]\addplot[fill=ink!55,draw=none] coordinates {(0,0.03612) (1,0.2479) (2,0.788) (3,3.753) (4,7.679) (5,17.56) (6,24.14) (7,29.68) (8,31.55) (9,25.18) (10,15.58) (11,5.066) (12,1.994) (13,0.06079) (14,0.01928)};\draw[densely dotted] (axis cs:12,0) -- (axis cs:12,34.71);\end{axis}\begin{axis}[width=72mm,height=42mm,axis y line*=right,axis x line=none,ymode=log,xmin=-1,xmax=15,ylabel={\scriptsize eindeutig (log)},tick label style={font=\tiny},label style={font=\scriptsize},ylabel shift=-3pt]\addplot[mark=*,mark size=0.7pt,draw=black,thin] coordinates {(0,36125) (1,247853) (2,357004) (3,1195283) (4,1774105) (5,1562867) (6,896922) (7,348106) (8,224677) (9,68210) (10,23782) (11,175) (12,12)};\end{axis}\end{tikzpicture}\end{minipage}
\vspace{3mm}

\noindent\begin{minipage}[t]{60mm}\vspace{0pt}\sbox\hmboard{\chessboard[setfen=8/8/7k/8/2b5/8/7R/K2n4 b - - 0 1]}\usebox\hmboard\\[1mm]
\makebox[\wd\hmboard]{\sffamily\hspace{3mm}\textbf{h\#6}\hfill(2+3)}\end{minipage}\hfill
\begin{minipage}[t]{52mm}\vspace{0pt}
\textcolor{ink}{\sffamily\textsc{Lösung}}\par\smallskip\small
\begin{tabular}{@{}r@{\enspace}l@{\enspace}l@{}}
\textbf{1.} & \symking{}g5 & \symrook{}c2\\
\textbf{2.} & \symking{}f4 & \symrook{}xc4+\\
\textbf{3.} & \symking{}e3 & \symking{}a2\\
\textbf{4.} & \symking{}d2 & \symrook{}d4+\\
\textbf{5.} & \symking{}c1 & \symking{}b3\\
\textbf{6.} & \symking{}b1 & \symrook{}xd1\#\\
\end{tabular}\par\vspace{3mm}
\qrcode[height=12mm]{https://helpman.komtera.lt/?fen=8/8/7k/8/2b5/8/7R/K2n4&moves=6}\\{\sffamily\tiny\color{ink}scannen zum Lösen}
\end{minipage}
\vspace{2mm}

{\raggedright\noindent\textcolor{ink}{\sffamily\textsc{Themen}}\enspace \themetag{reines Matt} \themetag{Mustermatt} \themetag{Idealmatt} \themetag{Spiegelmatt} \themetag{ein Stein zieht}\par\smallskip{\footnotesize\textit{Hinweis:} die Partei am Zug steht im Diagramm im Schach.}\par}

\begin{tcolorbox}[enhanced,colback=white,colframe=ink!35,arc=1.5pt,boxrule=0.4pt,left=3mm,right=3mm,top=1.5mm,bottom=1.5mm]
\textcolor{ink}{\sffamily\textsc{veröffentlicht}}\par\smallskip\small \textbf{Sheglow, Wiktor S.}, h\#6\quad{\color{ink}\scriptsize verwandte Stellung}\\{\scriptsize Suomen Tehtäväniekat, No. (306), 05-06/1998}
\end{tcolorbox}
\clearpage
\phantomsection\label{hm:KRvkbp}
\begin{tcolorbox}[enhanced,colback=ink,colframe=ink,arc=0pt,boxrule=0pt,left=4mm,right=4mm,top=2.5mm,bottom=2.5mm]
\color{white}\sffamily{\LARGE\bfseries KRvkbp}\hfill{\Large h\#6}\\[0.5mm]
{\small 5 Steine\quad Schwarz am Zug}\hfill{\small Nr.~184\qquad\href{https://helpman.komtera.lt/?fen=8/8/7k/6p1/6R1/8/8/K1b5&moves=6}{\color{white}\bfseries lösen\,\ensuremath{\nearrow}}}
\end{tcolorbox}
\vspace{2mm}
\noindent\begin{minipage}[t]{36mm}\vspace{0pt}\sffamily\small\raggedright
\hmfigure{längste eindeutige}{{\Large h\#6}}\hmfigure{eindeutig bei dieser Länge}{{\Large 817} Stellungen}\hmfigure{längstes Matt}{{\Large h\#7}}\hmfigure{in der Klasse}{710 Millionen Stellungen,\\ 505 Millionen mit Hilfsmatt}
\end{minipage}\hfill
\begin{minipage}[t]{80mm}\vspace{0pt}\begin{tikzpicture}\begin{axis}[width=72mm,height=42mm,ybar,bar width=1.6pt,xmin=-1,xmax=15,ymin=0,ymax=109.5,axis y line*=left,axis x line*=bottom,xlabel={\scriptsize Halbzüge bis Matt},ylabel={\scriptsize Stellungen (Mio.)},tick label style={font=\tiny},label style={font=\scriptsize},enlarge x limits=false,ylabel shift=-3pt,xlabel shift=-3pt]\addplot[fill=ink!55,draw=none] coordinates {(0,0.1443) (1,0.9373) (2,2.749) (3,12.51) (4,25.93) (5,54.88) (6,78.53) (7,89.04) (8,99.54) (9,73.69) (10,46.54) (11,13.81) (12,6.219) (13,0.1677) (14,0.05779)};\draw[densely dotted] (axis cs:12,0) -- (axis cs:12,109.5);\end{axis}\begin{axis}[width=72mm,height=42mm,axis y line*=right,axis x line=none,ymode=log,xmin=-1,xmax=15,ylabel={\scriptsize eindeutig (log)},tick label style={font=\tiny},label style={font=\scriptsize},ylabel shift=-3pt]\addplot[mark=*,mark size=0.7pt,draw=black,thin] coordinates {(0,144328) (1,937324) (2,1306515) (3,4093017) (4,6409621) (5,4769160) (6,2742405) (7,919724) (8,617201) (9,169104) (10,71113) (11,1388) (12,817)};\end{axis}\end{tikzpicture}\end{minipage}
\vspace{3mm}

\noindent\begin{minipage}[t]{60mm}\vspace{0pt}\sbox\hmboard{\chessboard[setfen=8/8/7k/6p1/6R1/8/8/K1b5 b - - 0 1]}\usebox\hmboard\\[1mm]
\makebox[\wd\hmboard]{\sffamily\hspace{3mm}\textbf{h\#6}\hfill(2+3)}\end{minipage}\hfill
\begin{minipage}[t]{52mm}\vspace{0pt}
\textcolor{ink}{\sffamily\textsc{Lösung}}\par\smallskip\small
\begin{tabular}{@{}r@{\enspace}l@{\enspace}l@{}}
\textbf{1.} & \symking{}h5 & \symrook{}xg5+\\
\textbf{2.} & \symking{}h4 & \symrook{}h5+\\
\textbf{3.} & \symking{}g3 & \symking{}b1\\
\textbf{4.} & \symking{}f2 & \symking{}c2\\
\textbf{5.} & \symking{}e1 & \symking{}d3\\
\textbf{6.} & \symking{}d1 & \symrook{}h1\#\\
\end{tabular}\par\vspace{3mm}
\qrcode[height=12mm]{https://helpman.komtera.lt/?fen=8/8/7k/6p1/6R1/8/8/K1b5&moves=6}\\{\sffamily\tiny\color{ink}scannen zum Lösen}
\end{minipage}
\vspace{2mm}

{\raggedright\noindent\textcolor{ink}{\sffamily\textsc{Themen}}\enspace \themetag{ein Stein zieht} \themetag{Umnov}\par}

\clearpage
\phantomsection\label{hm:KRvknn}
\begin{tcolorbox}[enhanced,colback=ink,colframe=ink,arc=0pt,boxrule=0pt,left=4mm,right=4mm,top=2.5mm,bottom=2.5mm]
\color{white}\sffamily{\LARGE\bfseries KRvknn}\hfill{\Large h\#6}\\[0.5mm]
{\small 5 Steine\quad Schwarz am Zug}\hfill{\small Nr.~185\qquad\href{https://helpman.komtera.lt/?fen=8/8/7k/6R1/8/8/1n5n/K7&moves=6}{\color{white}\bfseries lösen\,\ensuremath{\nearrow}}}
\end{tcolorbox}
\vspace{2mm}
\noindent\begin{minipage}[t]{36mm}\vspace{0pt}\sffamily\small\raggedright
\hmfigure{längste eindeutige}{{\Large h\#6}}\hmfigure{eindeutig bei dieser Länge}{{\Large 38} Stellungen}\hmfigure{längstes Matt}{{\Large h\#7}}\hmfigure{in der Klasse}{242 Millionen Stellungen,\\ 83,9 Millionen mit Hilfsmatt}
\end{minipage}\hfill
\begin{minipage}[t]{80mm}\vspace{0pt}\begin{tikzpicture}\begin{axis}[width=72mm,height=42mm,ybar,bar width=1.6pt,xmin=-1,xmax=15,ymin=0,ymax=18.07,axis y line*=left,axis x line*=bottom,xlabel={\scriptsize Halbzüge bis Matt},ylabel={\scriptsize Stellungen (Mio.)},tick label style={font=\tiny},label style={font=\scriptsize},enlarge x limits=false,ylabel shift=-3pt,xlabel shift=-3pt]\addplot[fill=ink!55,draw=none] coordinates {(0,0.02483) (1,0.1687) (2,0.4625) (3,2.001) (4,4.139) (5,8.767) (6,12.59) (7,14.69) (8,16.43) (9,12.5) (10,8.299) (11,2.609) (12,1.142) (13,0.0417) (14,0.01388)};\draw[densely dotted] (axis cs:12,0) -- (axis cs:12,18.07);\end{axis}\begin{axis}[width=72mm,height=42mm,axis y line*=right,axis x line=none,ymode=log,xmin=-1,xmax=15,ylabel={\scriptsize eindeutig (log)},tick label style={font=\tiny},label style={font=\scriptsize},ylabel shift=-3pt]\addplot[mark=*,mark size=0.7pt,draw=black,thin] coordinates {(0,24834) (1,168666) (2,225705) (3,712138) (4,1102736) (5,828473) (6,498786) (7,165986) (8,111396) (9,30006) (10,11324) (11,131) (12,38)};\end{axis}\end{tikzpicture}\end{minipage}
\vspace{3mm}

\noindent\begin{minipage}[t]{60mm}\vspace{0pt}\sbox\hmboard{\chessboard[setfen=8/8/7k/6R1/8/8/1n5n/K7 b - - 0 1]}\usebox\hmboard\\[1mm]
\makebox[\wd\hmboard]{\sffamily\hspace{3mm}\textbf{h\#6}\hfill(2+3)}\end{minipage}\hfill
\begin{minipage}[t]{52mm}\vspace{0pt}
\textcolor{ink}{\sffamily\textsc{Lösung}}\par\smallskip\small
\begin{tabular}{@{}r@{\enspace}l@{\enspace}l@{}}
\textbf{1.} & \symknight{}d3 & \symking{}b1\\
\textbf{2.} & \symknight{}f4 & \symking{}c2\\
\textbf{3.} & \symknight{}h5 & \symking{}d3\\
\textbf{4.} & \symknight{}g4 & \symking{}e4\\
\textbf{5.} & \symknight{}gf6+ & \symking{}f5\\
\textbf{6.} & \symknight{}h7 & \symrook{}g6\#\\
\end{tabular}\par\vspace{3mm}
\qrcode[height=12mm]{https://helpman.komtera.lt/?fen=8/8/7k/6R1/8/8/1n5n/K7&moves=6}\\{\sffamily\tiny\color{ink}scannen zum Lösen}
\end{minipage}
\vspace{2mm}

{\raggedright\noindent\textcolor{ink}{\sffamily\textsc{Themen}}\enspace \themetag{Selbstblock} \themetag{ohne Schlag}\par}

\clearpage
\phantomsection\label{hm:KRvknp}
\begin{tcolorbox}[enhanced,colback=ink,colframe=ink,arc=0pt,boxrule=0pt,left=4mm,right=4mm,top=2.5mm,bottom=2.5mm]
\color{white}\sffamily{\LARGE\bfseries KRvknp}\hfill{\Large h\#6}\\[0.5mm]
{\small 5 Steine\quad Schwarz am Zug}\hfill{\small Nr.~186\qquad\href{https://helpman.komtera.lt/?fen=8/8/6k1/5R2/p7/8/8/K1n5&moves=6}{\color{white}\bfseries lösen\,\ensuremath{\nearrow}}}
\end{tcolorbox}
\vspace{2mm}
\noindent\begin{minipage}[t]{36mm}\vspace{0pt}\sffamily\small\raggedright
\hmfigure{längste eindeutige}{{\Large h\#6}}\hmfigure{eindeutig bei dieser Länge}{{\Large 1\,731} Stellungen}\hmfigure{längstes Matt}{{\Large h\#7}}\hmfigure{in der Klasse}{710 Millionen Stellungen,\\ 519 Millionen mit Hilfsmatt}
\end{minipage}\hfill
\begin{minipage}[t]{80mm}\vspace{0pt}\begin{tikzpicture}\begin{axis}[width=72mm,height=42mm,ybar,bar width=1.6pt,xmin=-1,xmax=15,ymin=0,ymax=114.1,axis y line*=left,axis x line*=bottom,xlabel={\scriptsize Halbzüge bis Matt},ylabel={\scriptsize Stellungen (Mio.)},tick label style={font=\tiny},label style={font=\scriptsize},enlarge x limits=false,ylabel shift=-3pt,xlabel shift=-3pt]\addplot[fill=ink!55,draw=none] coordinates {(0,0.1992) (1,1.281) (2,3.198) (3,13.28) (4,28.3) (5,55.09) (6,83.16) (7,88.4) (8,103.8) (9,72.86) (10,48.87) (11,13.92) (12,6.731) (13,0.2095) (14,0.07686)};\draw[densely dotted] (axis cs:12,0) -- (axis cs:12,114.1);\end{axis}\begin{axis}[width=72mm,height=42mm,axis y line*=right,axis x line=none,ymode=log,xmin=-1,xmax=15,ylabel={\scriptsize eindeutig (log)},tick label style={font=\tiny},label style={font=\scriptsize},ylabel shift=-3pt]\addplot[mark=*,mark size=0.7pt,draw=black,thin] coordinates {(0,199174) (1,1281254) (2,1649578) (3,4928410) (4,8231507) (5,5071734) (6,3001982) (7,812865) (8,533300) (9,132758) (10,53953) (11,2035) (12,1731)};\end{axis}\end{tikzpicture}\end{minipage}
\vspace{3mm}

\noindent\begin{minipage}[t]{60mm}\vspace{0pt}\sbox\hmboard{\chessboard[setfen=8/8/6k1/5R2/p7/8/8/K1n5 b - - 0 1]}\usebox\hmboard\\[1mm]
\makebox[\wd\hmboard]{\sffamily\hspace{3mm}\textbf{h\#6}\hfill(2+3)}\end{minipage}\hfill
\begin{minipage}[t]{52mm}\vspace{0pt}
\textcolor{ink}{\sffamily\textsc{Lösung}}\par\smallskip\small
\begin{tabular}{@{}r@{\enspace}l@{\enspace}l@{}}
\textbf{1.} & \symking{}h7 & \symking{}b2\\
\textbf{2.} & a3+ & \symking{}c3\\
\textbf{3.} & a2 & \symking{}d4\\
\textbf{4.} & a1=\symrook{} & \symking{}e5\\
\textbf{5.} & \symrook{}a8 & \symking{}f6\\
\textbf{6.} & \symrook{}g8 & \symrook{}h5\#\\
\end{tabular}\par\vspace{3mm}
\qrcode[height=12mm]{https://helpman.komtera.lt/?fen=8/8/6k1/5R2/p7/8/8/K1n5&moves=6}\\{\sffamily\tiny\color{ink}scannen zum Lösen}
\end{minipage}
\vspace{2mm}

{\raggedright\noindent\textcolor{ink}{\sffamily\textsc{Themen}}\enspace \themetag{reines Matt} \themetag{Mustermatt} \themetag{Umwandlung} \themetag{Unterverwandlung} \themetag{Selbstblock} \themetag{ohne Schlag}\par}

\clearpage
\phantomsection\label{hm:KRvkpp}
\begin{tcolorbox}[enhanced,colback=ink,colframe=ink,arc=0pt,boxrule=0pt,left=4mm,right=4mm,top=2.5mm,bottom=2.5mm]
\color{white}\sffamily{\LARGE\bfseries KRvkpp}\hfill{\Large h\#6}\\[0.5mm]
{\small 5 Steine\quad Schwarz am Zug}\hfill{\small Nr.~187\qquad\href{https://helpman.komtera.lt/?fen=8/8/8/6pk/6R1/8/2p5/K7&moves=6}{\color{white}\bfseries lösen\,\ensuremath{\nearrow}}}
\end{tcolorbox}
\vspace{2mm}
\noindent\begin{minipage}[t]{36mm}\vspace{0pt}\sffamily\small\raggedright
\hmfigure{längste eindeutige}{{\Large h\#6}}\hmfigure{eindeutig bei dieser Länge}{{\Large 1\,085} Stellungen}\hmfigure{längstes Matt}{{\Large h\#7}}\hmfigure{in der Klasse}{533 Millionen Stellungen,\\ 200 Millionen mit Hilfsmatt}
\end{minipage}\hfill
\begin{minipage}[t]{80mm}\vspace{0pt}\begin{tikzpicture}\begin{axis}[width=72mm,height=42mm,ybar,bar width=1.6pt,xmin=-1,xmax=15,ymin=0,ymax=44.61,axis y line*=left,axis x line*=bottom,xlabel={\scriptsize Halbzüge bis Matt},ylabel={\scriptsize Stellungen (Mio.)},tick label style={font=\tiny},label style={font=\scriptsize},enlarge x limits=false,ylabel shift=-3pt,xlabel shift=-3pt]\addplot[fill=ink!55,draw=none] coordinates {(0,0.09949) (1,0.607) (2,1.348) (3,5.349) (4,11.78) (5,21.15) (6,33.42) (7,32.93) (8,40.56) (9,26.59) (10,18.3) (11,4.841) (12,2.781) (13,0.08108) (14,0.04221)};\draw[densely dotted] (axis cs:12,0) -- (axis cs:12,44.61);\end{axis}\begin{axis}[width=72mm,height=42mm,axis y line*=right,axis x line=none,ymode=log,xmin=-1,xmax=15,ylabel={\scriptsize eindeutig (log)},tick label style={font=\tiny},label style={font=\scriptsize},ylabel shift=-3pt]\addplot[mark=*,mark size=0.7pt,draw=black,thin] coordinates {(0,99485) (1,606974) (2,784878) (3,2113495) (4,3741938) (5,1884185) (6,1099707) (7,251812) (8,155508) (9,35732) (10,13542) (11,1333) (12,1085)};\end{axis}\end{tikzpicture}\end{minipage}
\vspace{3mm}

\noindent\begin{minipage}[t]{60mm}\vspace{0pt}\sbox\hmboard{\chessboard[setfen=8/8/8/6pk/6R1/8/2p5/K7 b - - 0 1]}\usebox\hmboard\\[1mm]
\makebox[\wd\hmboard]{\sffamily\hspace{3mm}\textbf{h\#6}\hfill(2+3)}\end{minipage}\hfill
\begin{minipage}[t]{52mm}\vspace{0pt}
\textcolor{ink}{\sffamily\textsc{Lösung}}\par\smallskip\small
\begin{tabular}{@{}r@{\enspace}l@{\enspace}l@{}}
\textbf{1.} & c1=\symbishop{} & \symrook{}xg5+\\
\textbf{2.} & \symking{}h4 & \symrook{}h5+\\
\textbf{3.} & \symking{}g3 & \symking{}b1\\
\textbf{4.} & \symking{}f2 & \symking{}c2\\
\textbf{5.} & \symking{}e1 & \symking{}d3\\
\textbf{6.} & \symking{}d1 & \symrook{}h1\#\\
\end{tabular}\par\vspace{3mm}
\qrcode[height=12mm]{https://helpman.komtera.lt/?fen=8/8/8/6pk/6R1/8/2p5/K7&moves=6}\\{\sffamily\tiny\color{ink}scannen zum Lösen}
\end{minipage}
\vspace{2mm}

{\raggedright\noindent\textcolor{ink}{\sffamily\textsc{Themen}}\enspace \themetag{Satzspiel} \themetag{Umwandlung} \themetag{Unterverwandlung} \themetag{Umnov}\par}

\clearpage
\phantomsection\label{hm:KRvkqn}
\begin{tcolorbox}[enhanced,colback=ink,colframe=ink,arc=0pt,boxrule=0pt,left=4mm,right=4mm,top=2.5mm,bottom=2.5mm]
\color{white}\sffamily{\LARGE\bfseries KRvkqn}\hfill{\Large h\#6}\\[0.5mm]
{\small 5 Steine\quad Schwarz am Zug}\hfill{\small Nr.~188\qquad\href{https://helpman.komtera.lt/?fen=1q6/8/6k1/4R3/8/8/8/K3n3&moves=6}{\color{white}\bfseries lösen\,\ensuremath{\nearrow}}}
\end{tcolorbox}
\vspace{2mm}
\noindent\begin{minipage}[t]{36mm}\vspace{0pt}\sffamily\small\raggedright
\hmfigure{längste eindeutige}{{\Large h\#6}}\hmfigure{eindeutig bei dieser Länge}{{\Large 66} Stellungen}\hmfigure{längstes Matt}{{\Large h\#7.5}}\hmfigure{in der Klasse}{242 Millionen Stellungen,\\ 144 Millionen mit Hilfsmatt}
\end{minipage}\hfill
\begin{minipage}[t]{80mm}\vspace{0pt}\begin{tikzpicture}\begin{axis}[width=72mm,height=42mm,ybar,bar width=1.6pt,xmin=-1,xmax=16,ymin=0,ymax=31.85,axis y line*=left,axis x line*=bottom,xlabel={\scriptsize Halbzüge bis Matt},ylabel={\scriptsize Stellungen (Mio.)},tick label style={font=\tiny},label style={font=\scriptsize},enlarge x limits=false,ylabel shift=-3pt,xlabel shift=-3pt]\addplot[fill=ink!55,draw=none] coordinates {(0,0.01517) (1,0.1142) (2,0.528) (3,2.942) (4,5.282) (5,15.22) (6,17.69) (7,28.96) (8,24.34) (9,26.83) (10,13.17) (11,7.13) (12,1.667) (13,0.1798) (14,0.01545) (15,0.000241)};\draw[densely dotted] (axis cs:12,0) -- (axis cs:12,31.85);\end{axis}\begin{axis}[width=72mm,height=42mm,axis y line*=right,axis x line=none,ymode=log,xmin=-1,xmax=16,ylabel={\scriptsize eindeutig (log)},tick label style={font=\tiny},label style={font=\scriptsize},ylabel shift=-3pt]\addplot[mark=*,mark size=0.7pt,draw=black,thin] coordinates {(0,15174) (1,114175) (2,237610) (3,974430) (4,948128) (5,1384306) (6,717039) (7,504043) (8,222438) (9,103588) (10,19847) (11,1139) (12,66)};\end{axis}\end{tikzpicture}\end{minipage}
\vspace{3mm}

\noindent\begin{minipage}[t]{60mm}\vspace{0pt}\sbox\hmboard{\chessboard[setfen=1q6/8/6k1/4R3/8/8/8/K3n3 b - - 0 1]}\usebox\hmboard\\[1mm]
\makebox[\wd\hmboard]{\sffamily\hspace{3mm}\textbf{h\#6}\hfill(2+3)}\end{minipage}\hfill
\begin{minipage}[t]{52mm}\vspace{0pt}
\textcolor{ink}{\sffamily\textsc{Lösung}}\par\smallskip\small
\begin{tabular}{@{}r@{\enspace}l@{\enspace}l@{}}
\textbf{1.} & \symqueen{}d8 & \symking{}b2\\
\textbf{2.} & \symqueen{}g5 & \symking{}c3\\
\textbf{3.} & \symking{}h5 & \symking{}d4\\
\textbf{4.} & \symking{}h4 & \symking{}e4\\
\textbf{5.} & \symqueen{}g2+ & \symking{}f4\\
\textbf{6.} & \symking{}h3 & \symrook{}h5\#\\
\end{tabular}\par\vspace{3mm}
\qrcode[height=12mm]{https://helpman.komtera.lt/?fen=1q6/8/6k1/4R3/8/8/8/K3n3&moves=6}\\{\sffamily\tiny\color{ink}scannen zum Lösen}
\end{minipage}
\vspace{2mm}

{\raggedright\noindent\textcolor{ink}{\sffamily\textsc{Themen}}\enspace \themetag{reines Matt} \themetag{Mustermatt} \themetag{Selbstblock} \themetag{ohne Schlag}\par}

\clearpage
\phantomsection\label{hm:KRvkrb}
\begin{tcolorbox}[enhanced,colback=ink,colframe=ink,arc=0pt,boxrule=0pt,left=4mm,right=4mm,top=2.5mm,bottom=2.5mm]
\color{white}\sffamily{\LARGE\bfseries KRvkrb}\hfill{\Large h\#6}\\[0.5mm]
{\small 5 Steine\quad Schwarz am Zug}\hfill{\small Nr.~189\qquad\href{https://helpman.komtera.lt/?fen=8/7k/8/6R1/8/b7/4r3/K7&moves=6}{\color{white}\bfseries lösen\,\ensuremath{\nearrow}}}
\end{tcolorbox}
\vspace{2mm}
\noindent\begin{minipage}[t]{36mm}\vspace{0pt}\sffamily\small\raggedright
\hmfigure{längste eindeutige}{{\Large h\#6}}\hmfigure{eindeutig bei dieser Länge}{{\Large 55} Stellungen}\hmfigure{längstes Matt}{{\Large h\#7}}\hmfigure{in der Klasse}{242 Millionen Stellungen,\\ 153 Millionen mit Hilfsmatt}
\end{minipage}\hfill
\begin{minipage}[t]{80mm}\vspace{0pt}\begin{tikzpicture}\begin{axis}[width=72mm,height=42mm,ybar,bar width=1.6pt,xmin=-1,xmax=15,ymin=0,ymax=32.24,axis y line*=left,axis x line*=bottom,xlabel={\scriptsize Halbzüge bis Matt},ylabel={\scriptsize Stellungen (Mio.)},tick label style={font=\tiny},label style={font=\scriptsize},enlarge x limits=false,ylabel shift=-3pt,xlabel shift=-3pt]\addplot[fill=ink!55,draw=none] coordinates {(0,0.02149) (1,0.1548) (2,0.609) (3,3.244) (4,6.332) (5,16.02) (6,20.47) (7,29.31) (8,27.52) (9,26.35) (10,14.32) (11,6.306) (12,1.744) (13,0.1045) (14,0.01328)};\draw[densely dotted] (axis cs:12,0) -- (axis cs:12,32.24);\end{axis}\begin{axis}[width=72mm,height=42mm,axis y line*=right,axis x line=none,ymode=log,xmin=-1,xmax=15,ylabel={\scriptsize eindeutig (log)},tick label style={font=\tiny},label style={font=\scriptsize},ylabel shift=-3pt]\addplot[mark=*,mark size=0.7pt,draw=black,thin] coordinates {(0,21489) (1,154833) (2,272509) (3,986677) (4,1194192) (5,1243271) (6,691269) (7,405675) (8,215887) (9,90635) (10,24230) (11,1096) (12,55)};\end{axis}\end{tikzpicture}\end{minipage}
\vspace{3mm}

\noindent\begin{minipage}[t]{60mm}\vspace{0pt}\sbox\hmboard{\chessboard[setfen=8/7k/8/6R1/8/b7/4r3/K7 b - - 0 1]}\usebox\hmboard\\[1mm]
\makebox[\wd\hmboard]{\sffamily\hspace{3mm}\textbf{h\#6}\hfill(2+3)}\end{minipage}\hfill
\begin{minipage}[t]{52mm}\vspace{0pt}
\textcolor{ink}{\sffamily\textsc{Lösung}}\par\smallskip\small
\begin{tabular}{@{}r@{\enspace}l@{\enspace}l@{}}
\textbf{1.} & \symbishop{}b2+ & \symking{}a2\\
\textbf{2.} & \symbishop{}g7+ & \symking{}b3\\
\textbf{3.} & \symking{}g8 & \symking{}c4\\
\textbf{4.} & \symking{}f8 & \symking{}d5\\
\textbf{5.} & \symbishop{}e5 & \symking{}e6\\
\textbf{6.} & \symking{}e8 & \symrook{}g8\#\\
\end{tabular}\par\vspace{3mm}
\qrcode[height=12mm]{https://helpman.komtera.lt/?fen=8/7k/8/6R1/8/b7/4r3/K7&moves=6}\\{\sffamily\tiny\color{ink}scannen zum Lösen}
\end{minipage}
\vspace{2mm}

{\raggedright\noindent\textcolor{ink}{\sffamily\textsc{Themen}}\enspace \themetag{reines Matt} \themetag{Mustermatt} \themetag{Spiegelmatt} \themetag{ohne Schlag}\par}

\clearpage
\phantomsection\label{hm:KRvkrn}
\begin{tcolorbox}[enhanced,colback=ink,colframe=ink,arc=0pt,boxrule=0pt,left=4mm,right=4mm,top=2.5mm,bottom=2.5mm]
\color{white}\sffamily{\LARGE\bfseries KRvkrn}\hfill{\Large h\#6}\\[0.5mm]
{\small 5 Steine\quad Schwarz am Zug}\hfill{\small Nr.~190\qquad\href{https://helpman.komtera.lt/?fen=8/8/6k1/6R1/8/8/3r3n/K7&moves=6}{\color{white}\bfseries lösen\,\ensuremath{\nearrow}}}
\end{tcolorbox}
\vspace{2mm}
\noindent\begin{minipage}[t]{36mm}\vspace{0pt}\sffamily\small\raggedright
\hmfigure{längste eindeutige}{{\Large h\#6}}\hmfigure{eindeutig bei dieser Länge}{{\Large 22} Stellungen}\hmfigure{längstes Matt}{{\Large h\#7.5}}\hmfigure{in der Klasse}{242 Millionen Stellungen,\\ 156 Millionen mit Hilfsmatt}
\end{minipage}\hfill
\begin{minipage}[t]{80mm}\vspace{0pt}\begin{tikzpicture}\begin{axis}[width=72mm,height=42mm,ybar,bar width=1.6pt,xmin=-1,xmax=16,ymin=0,ymax=32.7,axis y line*=left,axis x line*=bottom,xlabel={\scriptsize Halbzüge bis Matt},ylabel={\scriptsize Stellungen (Mio.)},tick label style={font=\tiny},label style={font=\scriptsize},enlarge x limits=false,ylabel shift=-3pt,xlabel shift=-3pt]\addplot[fill=ink!55,draw=none] coordinates {(0,0.02788) (1,0.2014) (2,0.7429) (3,3.596) (4,7.194) (5,16.55) (6,22.6) (7,29.73) (8,29.26) (9,25.75) (10,13.75) (11,5.604) (12,1.288) (13,0.06651) (14,0.004178) (15,6e-06)};\draw[densely dotted] (axis cs:12,0) -- (axis cs:12,32.7);\end{axis}\begin{axis}[width=72mm,height=42mm,axis y line*=right,axis x line=none,ymode=log,xmin=-1,xmax=16,ylabel={\scriptsize eindeutig (log)},tick label style={font=\tiny},label style={font=\scriptsize},ylabel shift=-3pt]\addplot[mark=*,mark size=0.7pt,draw=black,thin] coordinates {(0,27883) (1,201403) (2,347294) (3,1164938) (4,1551113) (5,1376976) (6,822132) (7,418691) (8,224365) (9,77173) (10,18921) (11,801) (12,22)};\end{axis}\end{tikzpicture}\end{minipage}
\vspace{3mm}

\noindent\begin{minipage}[t]{60mm}\vspace{0pt}\sbox\hmboard{\chessboard[setfen=8/8/6k1/6R1/8/8/3r3n/K7 b - - 0 1]}\usebox\hmboard\\[1mm]
\makebox[\wd\hmboard]{\sffamily\hspace{3mm}\textbf{h\#6}\hfill(2+3)}\end{minipage}\hfill
\begin{minipage}[t]{52mm}\vspace{0pt}
\textcolor{ink}{\sffamily\textsc{Lösung}}\par\smallskip\small
\begin{tabular}{@{}r@{\enspace}l@{\enspace}l@{}}
\textbf{1.} & \symking{}h6 & \symking{}b1\\
\textbf{2.} & \symrook{}d7 & \symking{}c2\\
\textbf{3.} & \symrook{}h7 & \symking{}d3\\
\textbf{4.} & \symknight{}g4 & \symking{}e4\\
\textbf{5.} & \symknight{}f6+ & \symking{}f5\\
\textbf{6.} & \symknight{}h5 & \symrook{}g6\#\\
\end{tabular}\par\vspace{3mm}
\qrcode[height=12mm]{https://helpman.komtera.lt/?fen=8/8/6k1/6R1/8/8/3r3n/K7&moves=6}\\{\sffamily\tiny\color{ink}scannen zum Lösen}
\end{minipage}
\vspace{2mm}

{\raggedright\noindent\textcolor{ink}{\sffamily\textsc{Themen}}\enspace \themetag{Selbstblock} \themetag{ohne Schlag}\par\smallskip{\footnotesize\textit{Hinweis:} die Partei am Zug steht im Diagramm im Schach.}\par}

\clearpage
\phantomsection\label{hm:KRvkrp}
\begin{tcolorbox}[enhanced,colback=ink,colframe=ink,arc=0pt,boxrule=0pt,left=4mm,right=4mm,top=2.5mm,bottom=2.5mm]
\color{white}\sffamily{\LARGE\bfseries KRvkrp}\hfill{\Large h\#6}\\[0.5mm]
{\small 5 Steine\quad Schwarz am Zug}\hfill{\small Nr.~191\qquad\href{https://helpman.komtera.lt/?fen=8/8/7k/6R1/8/1p6/1r6/K7&moves=6}{\color{white}\bfseries lösen\,\ensuremath{\nearrow}}}
\end{tcolorbox}
\vspace{2mm}
\noindent\begin{minipage}[t]{36mm}\vspace{0pt}\sffamily\small\raggedright
\hmfigure{längste eindeutige}{{\Large h\#6}}\hmfigure{eindeutig bei dieser Länge}{{\Large 175} Stellungen}\hmfigure{längstes Matt}{{\Large h\#7.5}}\hmfigure{in der Klasse}{710 Millionen Stellungen,\\ 483 Millionen mit Hilfsmatt}
\end{minipage}\hfill
\begin{minipage}[t]{80mm}\vspace{0pt}\begin{tikzpicture}\begin{axis}[width=72mm,height=42mm,ybar,bar width=1.6pt,xmin=-1,xmax=16,ymin=0,ymax=100,axis y line*=left,axis x line*=bottom,xlabel={\scriptsize Halbzüge bis Matt},ylabel={\scriptsize Stellungen (Mio.)},tick label style={font=\tiny},label style={font=\scriptsize},enlarge x limits=false,ylabel shift=-3pt,xlabel shift=-3pt]\addplot[fill=ink!55,draw=none] coordinates {(0,0.1139) (1,0.7736) (2,2.657) (3,12.08) (4,23.9) (5,51.11) (6,71.17) (7,87.96) (8,90.91) (9,76.09) (10,44.01) (11,16.64) (12,5.32) (13,0.2625) (14,0.03864) (15,4e-06)};\draw[densely dotted] (axis cs:12,0) -- (axis cs:12,100);\end{axis}\begin{axis}[width=72mm,height=42mm,axis y line*=right,axis x line=none,ymode=log,xmin=-1,xmax=16,ylabel={\scriptsize eindeutig (log)},tick label style={font=\tiny},label style={font=\scriptsize},ylabel shift=-3pt]\addplot[mark=*,mark size=0.7pt,draw=black,thin] coordinates {(0,113931) (1,773584) (2,1296202) (3,3920342) (4,5095747) (5,4096893) (6,2430697) (7,1128910) (8,663895) (9,203836) (10,62096) (11,2652) (12,175)};\end{axis}\end{tikzpicture}\end{minipage}
\vspace{3mm}

\noindent\begin{minipage}[t]{60mm}\vspace{0pt}\sbox\hmboard{\chessboard[setfen=8/8/7k/6R1/8/1p6/1r6/K7 b - - 0 1]}\usebox\hmboard\\[1mm]
\makebox[\wd\hmboard]{\sffamily\hspace{3mm}\textbf{h\#6}\hfill(2+3)}\end{minipage}\hfill
\begin{minipage}[t]{52mm}\vspace{0pt}
\textcolor{ink}{\sffamily\textsc{Lösung}}\par\smallskip\small
\begin{tabular}{@{}r@{\enspace}l@{\enspace}l@{}}
\textbf{1.} & \symking{}h7 & \symking{}xb2\\
\textbf{2.} & \symking{}h8 & \symking{}c3\\
\textbf{3.} & b2 & \symking{}d4\\
\textbf{4.} & b1=\symrook{} & \symking{}e5\\
\textbf{5.} & \symrook{}b8 & \symking{}f6\\
\textbf{6.} & \symrook{}g8 & \symrook{}h5\#\\
\end{tabular}\par\vspace{3mm}
\qrcode[height=12mm]{https://helpman.komtera.lt/?fen=8/8/7k/6R1/8/1p6/1r6/K7&moves=6}\\{\sffamily\tiny\color{ink}scannen zum Lösen}
\end{minipage}
\vspace{2mm}

{\raggedright\noindent\textcolor{ink}{\sffamily\textsc{Themen}}\enspace \themetag{Satzspiel} \themetag{reines Matt} \themetag{Mustermatt} \themetag{Idealmatt} \themetag{Umwandlung} \themetag{Unterverwandlung} \themetag{Selbstblock} \themetag{Phönix} \themetag{ohne Schach} \themetag{Umnov}\par}

\clearpage
\phantomsection\label{hm:KRvkrr}
\begin{tcolorbox}[enhanced,colback=ink,colframe=ink,arc=0pt,boxrule=0pt,left=4mm,right=4mm,top=2.5mm,bottom=2.5mm]
\color{white}\sffamily{\LARGE\bfseries KRvkrr}\hfill{\Large h\#6}\\[0.5mm]
{\small 5 Steine\quad Schwarz am Zug}\hfill{\small Nr.~192\qquad\href{https://helpman.komtera.lt/?fen=8/8/8/1r4k1/6R1/8/3r4/K7&moves=6}{\color{white}\bfseries lösen\,\ensuremath{\nearrow}}}
\end{tcolorbox}
\vspace{2mm}
\noindent\begin{minipage}[t]{36mm}\vspace{0pt}\sffamily\small\raggedright
\hmfigure{längste eindeutige}{{\Large h\#6}}\hmfigure{eindeutig bei dieser Länge}{{\Large 100} Stellungen}\hmfigure{längstes Matt}{{\Large h\#7.5}}\hmfigure{in der Klasse}{242 Millionen Stellungen,\\ 72,3 Millionen mit Hilfsmatt}
\end{minipage}\hfill
\begin{minipage}[t]{80mm}\vspace{0pt}\begin{tikzpicture}\begin{axis}[width=72mm,height=42mm,ybar,bar width=1.6pt,xmin=-1,xmax=16,ymin=0,ymax=15.52,axis y line*=left,axis x line*=bottom,xlabel={\scriptsize Halbzüge bis Matt},ylabel={\scriptsize Stellungen (Mio.)},tick label style={font=\tiny},label style={font=\scriptsize},enlarge x limits=false,ylabel shift=-3pt,xlabel shift=-3pt]\addplot[fill=ink!55,draw=none] coordinates {(0,0.007904) (1,0.06063) (2,0.2639) (3,1.455) (4,2.757) (5,7.286) (6,8.936) (7,14.11) (8,12.17) (9,13.54) (10,6.732) (11,3.953) (12,0.9118) (13,0.1182) (14,0.01001) (15,3.9e-05)};\draw[densely dotted] (axis cs:12,0) -- (axis cs:12,15.52);\end{axis}\begin{axis}[width=72mm,height=42mm,axis y line*=right,axis x line=none,ymode=log,xmin=-1,xmax=16,ylabel={\scriptsize eindeutig (log)},tick label style={font=\tiny},label style={font=\scriptsize},ylabel shift=-3pt]\addplot[mark=*,mark size=0.7pt,draw=black,thin] coordinates {(0,7904) (1,60633) (2,114246) (3,429249) (4,458300) (5,515936) (6,274439) (7,209129) (8,94023) (9,53680) (10,12744) (11,1546) (12,100)};\end{axis}\end{tikzpicture}\end{minipage}
\vspace{3mm}

\noindent\begin{minipage}[t]{60mm}\vspace{0pt}\sbox\hmboard{\chessboard[setfen=8/8/8/1r4k1/6R1/8/3r4/K7 b - - 0 1]}\usebox\hmboard\\[1mm]
\makebox[\wd\hmboard]{\sffamily\hspace{3mm}\textbf{h\#6}\hfill(2+3)}\end{minipage}\hfill
\begin{minipage}[t]{52mm}\vspace{0pt}
\textcolor{ink}{\sffamily\textsc{Lösung}}\par\smallskip\small
\begin{tabular}{@{}r@{\enspace}l@{\enspace}l@{}}
\textbf{1.} & \symking{}f5 & \symrook{}g5+\\
\textbf{2.} & \symking{}e4 & \symrook{}xb5\\
\textbf{3.} & \symking{}d3 & \symrook{}a5\\
\textbf{4.} & \symking{}c2 & \symking{}a2\\
\textbf{5.} & \symking{}d1+ & \symking{}b3\\
\textbf{6.} & \symking{}c1 & \symrook{}a1\#\\
\end{tabular}\par\vspace{3mm}
\qrcode[height=12mm]{https://helpman.komtera.lt/?fen=8/8/8/1r4k1/6R1/8/3r4/K7&moves=6}\\{\sffamily\tiny\color{ink}scannen zum Lösen}
\end{minipage}
\vspace{2mm}

{\raggedright\noindent\textcolor{ink}{\sffamily\textsc{Themen}}\enspace \themetag{reines Matt} \themetag{Mustermatt} \themetag{Idealmatt} \themetag{ein Stein zieht} \themetag{Umnov}\par\smallskip{\footnotesize\textit{Hinweis:} die Partei am Zug steht im Diagramm im Schach.}\par}

\clearpage
\phantomsection\label{hm:KQBvkq}
\begin{tcolorbox}[enhanced,colback=ink,colframe=ink,arc=0pt,boxrule=0pt,left=4mm,right=4mm,top=2.5mm,bottom=2.5mm]
\color{white}\sffamily{\LARGE\bfseries KQBvkq}\hfill{\Large h\#5.5}\\[0.5mm]
{\small 5 Steine\quad Weiß am Zug}\hfill{\small Nr.~193\qquad\href{https://helpman.komtera.lt/?fen=7B/7k/8/6Q1/8/8/1q6/K7&moves=5.5}{\color{white}\bfseries lösen\,\ensuremath{\nearrow}}}
\end{tcolorbox}
\vspace{2mm}
\noindent\begin{minipage}[t]{36mm}\vspace{0pt}\sffamily\small\raggedright
\hmfigure{längste eindeutige}{{\Large h\#5.5}}\hmfigure{eindeutig bei dieser Länge}{{\Large 6} Stellungen}\hmfigure{längstes Matt}{{\Large h\#7}}\hmfigure{in der Klasse}{242 Millionen Stellungen,\\ 126 Millionen mit Hilfsmatt}
\end{minipage}\hfill
\begin{minipage}[t]{80mm}\vspace{0pt}\begin{tikzpicture}\begin{axis}[width=72mm,height=42mm,ybar,bar width=1.6pt,xmin=-1,xmax=15,ymin=0,ymax=57.18,axis y line*=left,axis x line*=bottom,xlabel={\scriptsize Halbzüge bis Matt},ylabel={\scriptsize Stellungen (Mio.)},tick label style={font=\tiny},label style={font=\scriptsize},enlarge x limits=false,ylabel shift=-3pt,xlabel shift=-3pt]\addplot[fill=ink!55,draw=none] coordinates {(0,0.1843) (1,1.364) (2,6.806) (3,25.7) (4,51.98) (5,31.15) (6,7.978) (7,0.5692) (8,0.06024) (9,0.000332) (10,0.02496) (11,0.00012) (12,0.002992) (13,1.7e-05) (14,6.3e-05)};\draw[densely dotted] (axis cs:11,0) -- (axis cs:11,57.18);\end{axis}\begin{axis}[width=72mm,height=42mm,axis y line*=right,axis x line=none,ymode=log,xmin=-1,xmax=15,ylabel={\scriptsize eindeutig (log)},tick label style={font=\tiny},label style={font=\scriptsize},ylabel shift=-3pt]\addplot[mark=*,mark size=0.7pt,draw=black,thin] coordinates {(0,184332) (1,1212644) (2,4123117) (3,5628200) (4,2609938) (5,156368) (6,3021) (7,12) (8,913) (9,4) (10,49) (11,6)};\end{axis}\end{tikzpicture}\end{minipage}
\vspace{3mm}

\noindent\begin{minipage}[t]{60mm}\vspace{0pt}\sbox\hmboard{\chessboard[setfen=7B/7k/8/6Q1/8/8/1q6/K7 w - - 0 1]}\usebox\hmboard\\[1mm]
\makebox[\wd\hmboard]{\sffamily\hspace{3mm}\textbf{h\#5.5}\hfill(3+2)}\end{minipage}\hfill
\begin{minipage}[t]{52mm}\vspace{0pt}
\textcolor{ink}{\sffamily\textsc{Lösung}}\par\smallskip\small
\begin{tabular}{@{}r@{\enspace}l@{\enspace}l@{}}
\textbf{1.} & \ldots & \symking{}xb2\\
\textbf{2.} & \symking{}xh8 & \symking{}c3\\
\textbf{3.} & \symking{}h7 & \symking{}d4\\
\textbf{4.} & \symking{}h8 & \symking{}e5\\
\textbf{5.} & \symking{}h7 & \symking{}f6\\
\textbf{6.} & \symking{}h8 & \symqueen{}g7\#\\
\end{tabular}\par\vspace{3mm}
\qrcode[height=12mm]{https://helpman.komtera.lt/?fen=7B/7k/8/6Q1/8/8/1q6/K7&moves=5.5}\\{\sffamily\tiny\color{ink}scannen zum Lösen}
\end{minipage}
\vspace{2mm}

{\raggedright\noindent\textcolor{ink}{\sffamily\textsc{Themen}}\enspace \themetag{reines Matt} \themetag{Mustermatt} \themetag{Idealmatt} \themetag{Switchback} \themetag{ein Stein zieht} \themetag{Kniest} \themetag{Pendel} \themetag{ohne Schach}\par\smallskip{\footnotesize\textit{Hinweis:} die Partei am Zug steht im Diagramm im Schach; die Lösung beginnt mit einem Schlag.}\par}

\clearpage
\phantomsection\label{hm:KQNvkq}
\begin{tcolorbox}[enhanced,colback=ink,colframe=ink,arc=0pt,boxrule=0pt,left=4mm,right=4mm,top=2.5mm,bottom=2.5mm]
\color{white}\sffamily{\LARGE\bfseries KQNvkq}\hfill{\Large h\#5.5}\\[0.5mm]
{\small 5 Steine\quad Weiß am Zug}\hfill{\small Nr.~194\qquad\href{https://helpman.komtera.lt/?fen=7N/7k/8/6Q1/8/8/1q6/K7&moves=5.5}{\color{white}\bfseries lösen\,\ensuremath{\nearrow}}}
\end{tcolorbox}
\vspace{2mm}
\noindent\begin{minipage}[t]{36mm}\vspace{0pt}\sffamily\small\raggedright
\hmfigure{längste eindeutige}{{\Large h\#5.5}}\hmfigure{eindeutig bei dieser Länge}{{\Large 6} Stellungen}\hmfigure{längstes Matt}{{\Large h\#7}}\hmfigure{in der Klasse}{242 Millionen Stellungen,\\ 130 Millionen mit Hilfsmatt}
\end{minipage}\hfill
\begin{minipage}[t]{80mm}\vspace{0pt}\begin{tikzpicture}\begin{axis}[width=72mm,height=42mm,ybar,bar width=1.6pt,xmin=-1,xmax=15,ymin=0,ymax=41.69,axis y line*=left,axis x line*=bottom,xlabel={\scriptsize Halbzüge bis Matt},ylabel={\scriptsize Stellungen (Mio.)},tick label style={font=\tiny},label style={font=\scriptsize},enlarge x limits=false,ylabel shift=-3pt,xlabel shift=-3pt]\addplot[fill=ink!55,draw=none] coordinates {(0,0.1087) (1,0.9668) (2,4.502) (3,17.25) (4,37.9) (5,37.5) (6,23.73) (7,6.913) (8,0.9678) (9,0.07556) (10,0.007429) (11,8.3e-05) (12,0.001037) (13,1.7e-05) (14,2.9e-05)};\draw[densely dotted] (axis cs:11,0) -- (axis cs:11,41.69);\end{axis}\begin{axis}[width=72mm,height=42mm,axis y line*=right,axis x line=none,ymode=log,xmin=-1,xmax=15,ylabel={\scriptsize eindeutig (log)},tick label style={font=\tiny},label style={font=\scriptsize},ylabel shift=-3pt]\addplot[mark=*,mark size=0.7pt,draw=black,thin] coordinates {(0,108666) (1,858790) (2,2643186) (3,3592653) (4,2880172) (5,715236) (6,70122) (7,5898) (8,240) (9,3) (10,18) (11,6)};\end{axis}\end{tikzpicture}\end{minipage}
\vspace{3mm}

\noindent\begin{minipage}[t]{60mm}\vspace{0pt}\sbox\hmboard{\chessboard[setfen=7N/7k/8/6Q1/8/8/1q6/K7 w - - 0 1]}\usebox\hmboard\\[1mm]
\makebox[\wd\hmboard]{\sffamily\hspace{3mm}\textbf{h\#5.5}\hfill(3+2)}\end{minipage}\hfill
\begin{minipage}[t]{52mm}\vspace{0pt}
\textcolor{ink}{\sffamily\textsc{Lösung}}\par\smallskip\small
\begin{tabular}{@{}r@{\enspace}l@{\enspace}l@{}}
\textbf{1.} & \ldots & \symking{}xb2\\
\textbf{2.} & \symking{}xh8 & \symking{}c3\\
\textbf{3.} & \symking{}h7 & \symking{}d4\\
\textbf{4.} & \symking{}h8 & \symking{}e5\\
\textbf{5.} & \symking{}h7 & \symking{}f6\\
\textbf{6.} & \symking{}h8 & \symqueen{}g7\#\\
\end{tabular}\par\vspace{3mm}
\qrcode[height=12mm]{https://helpman.komtera.lt/?fen=7N/7k/8/6Q1/8/8/1q6/K7&moves=5.5}\\{\sffamily\tiny\color{ink}scannen zum Lösen}
\end{minipage}
\vspace{2mm}

{\raggedright\noindent\textcolor{ink}{\sffamily\textsc{Themen}}\enspace \themetag{reines Matt} \themetag{Mustermatt} \themetag{Idealmatt} \themetag{Switchback} \themetag{ein Stein zieht} \themetag{Kniest} \themetag{Pendel} \themetag{ohne Schach}\par\smallskip{\footnotesize\textit{Hinweis:} die Partei am Zug steht im Diagramm im Schach; die Lösung beginnt mit einem Schlag.}\par}

\clearpage
\phantomsection\label{hm:KQvkbb}
\begin{tcolorbox}[enhanced,colback=ink,colframe=ink,arc=0pt,boxrule=0pt,left=4mm,right=4mm,top=2.5mm,bottom=2.5mm]
\color{white}\sffamily{\LARGE\bfseries KQvkbb}\hfill{\Large h\#5.5}\\[0.5mm]
{\small 5 Steine\quad Weiß am Zug}\hfill{\small Nr.~195\qquad\href{https://helpman.komtera.lt/?fen=8/8/7k/8/8/1bb5/8/KQ6&moves=5.5}{\color{white}\bfseries lösen\,\ensuremath{\nearrow}}}
\end{tcolorbox}
\vspace{2mm}
\noindent\begin{minipage}[t]{36mm}\vspace{0pt}\sffamily\small\raggedright
\hmfigure{längste eindeutige}{{\Large h\#5.5}}\hmfigure{eindeutig bei dieser Länge}{{\Large 4} Stellungen}\hmfigure{längstes Matt}{{\Large h\#7}}\hmfigure{in der Klasse}{242 Millionen Stellungen,\\ 72 Millionen mit Hilfsmatt}
\end{minipage}\hfill
\begin{minipage}[t]{80mm}\vspace{0pt}\begin{tikzpicture}\begin{axis}[width=72mm,height=42mm,ybar,bar width=1.6pt,xmin=-1,xmax=15,ymin=0,ymax=17.26,axis y line*=left,axis x line*=bottom,xlabel={\scriptsize Halbzüge bis Matt},ylabel={\scriptsize Stellungen (Mio.)},tick label style={font=\tiny},label style={font=\scriptsize},enlarge x limits=false,ylabel shift=-3pt,xlabel shift=-3pt]\addplot[fill=ink!55,draw=none] coordinates {(0,0.03366) (1,0.2708) (2,0.8795) (3,3.316) (4,7.474) (5,11.77) (6,15.69) (7,13.2) (8,11.42) (9,5.122) (10,2.289) (11,0.383) (12,0.1578) (13,0.00075) (14,7.6e-05)};\draw[densely dotted] (axis cs:11,0) -- (axis cs:11,17.26);\end{axis}\begin{axis}[width=72mm,height=42mm,axis y line*=right,axis x line=none,ymode=log,xmin=-1,xmax=15,ylabel={\scriptsize eindeutig (log)},tick label style={font=\tiny},label style={font=\scriptsize},ylabel shift=-3pt]\addplot[mark=*,mark size=0.7pt,draw=black,thin] coordinates {(0,33659) (1,213070) (2,402766) (3,684942) (4,784677) (5,576981) (6,379156) (7,198916) (8,76472) (9,9097) (10,1374) (11,4)};\end{axis}\end{tikzpicture}\end{minipage}
\vspace{3mm}

\noindent\begin{minipage}[t]{60mm}\vspace{0pt}\sbox\hmboard{\chessboard[setfen=8/8/7k/8/8/1bb5/8/KQ6 w - - 0 1]}\usebox\hmboard\\[1mm]
\makebox[\wd\hmboard]{\sffamily\hspace{3mm}\textbf{h\#5.5}\hfill(2+3)}\end{minipage}\hfill
\begin{minipage}[t]{52mm}\vspace{0pt}
\textcolor{ink}{\sffamily\textsc{Lösung}}\par\smallskip\small
\begin{tabular}{@{}r@{\enspace}l@{\enspace}l@{}}
\textbf{1.} & \ldots & \symqueen{}b2\\
\textbf{2.} & \symking{}g5 & \symqueen{}xc3\\
\textbf{3.} & \symking{}f4 & \symqueen{}f6+\\
\textbf{4.} & \symking{}e3 & \symking{}b2\\
\textbf{5.} & \symking{}e2 & \symking{}c3\\
\textbf{6.} & \symking{}d1 & \symqueen{}f1\#\\
\end{tabular}\par\vspace{3mm}
\qrcode[height=12mm]{https://helpman.komtera.lt/?fen=8/8/7k/8/8/1bb5/8/KQ6&moves=5.5}\\{\sffamily\tiny\color{ink}scannen zum Lösen}
\end{minipage}
\vspace{2mm}

{\raggedright\noindent\textcolor{ink}{\sffamily\textsc{Themen}}\enspace \themetag{reines Matt} \themetag{Mustermatt} \themetag{Spiegelmatt} \themetag{ein Stein zieht}\par\smallskip{\footnotesize\textit{Hinweis:} die Partei am Zug steht im Diagramm im Schach.}\par}

\clearpage
\phantomsection\label{hm:KQvkbp}
\begin{tcolorbox}[enhanced,colback=ink,colframe=ink,arc=0pt,boxrule=0pt,left=4mm,right=4mm,top=2.5mm,bottom=2.5mm]
\color{white}\sffamily{\LARGE\bfseries KQvkbp}\hfill{\Large h\#5.5}\\[0.5mm]
{\small 5 Steine\quad Weiß am Zug}\hfill{\small Nr.~196\qquad\href{https://helpman.komtera.lt/?fen=6Q1/8/7k/8/3b4/8/p7/K7&moves=5.5}{\color{white}\bfseries lösen\,\ensuremath{\nearrow}}}
\end{tcolorbox}
\vspace{2mm}
\noindent\begin{minipage}[t]{36mm}\vspace{0pt}\sffamily\small\raggedright
\hmfigure{längste eindeutige}{{\Large h\#5.5}}\hmfigure{eindeutig bei dieser Länge}{{\Large 2} Stellungen}\hmfigure{längstes Matt}{{\Large h\#7}}\hmfigure{in der Klasse}{710 Millionen Stellungen,\\ 462 Millionen mit Hilfsmatt}
\end{minipage}\hfill
\begin{minipage}[t]{80mm}\vspace{0pt}\begin{tikzpicture}\begin{axis}[width=72mm,height=42mm,ybar,bar width=1.6pt,xmin=-1,xmax=15,ymin=0,ymax=121.3,axis y line*=left,axis x line*=bottom,xlabel={\scriptsize Halbzüge bis Matt},ylabel={\scriptsize Stellungen (Mio.)},tick label style={font=\tiny},label style={font=\scriptsize},enlarge x limits=false,ylabel shift=-3pt,xlabel shift=-3pt]\addplot[fill=ink!55,draw=none] coordinates {(0,0.3087) (1,2.276) (2,6.553) (3,22.28) (4,53.93) (5,73.67) (6,110.2) (7,79.23) (8,72.94) (9,26.07) (10,12.2) (11,1.685) (12,0.6124) (13,0.000136) (14,1.2e-05)};\draw[densely dotted] (axis cs:11,0) -- (axis cs:11,121.3);\end{axis}\begin{axis}[width=72mm,height=42mm,axis y line*=right,axis x line=none,ymode=log,xmin=-1,xmax=15,ylabel={\scriptsize eindeutig (log)},tick label style={font=\tiny},label style={font=\scriptsize},ylabel shift=-3pt]\addplot[mark=*,mark size=0.7pt,draw=black,thin] coordinates {(0,308732) (1,1671796) (2,3026013) (3,4445249) (4,5820488) (5,3550395) (6,2462137) (7,1001387) (8,399258) (9,37542) (10,8168) (11,2)};\end{axis}\end{tikzpicture}\end{minipage}
\vspace{3mm}

\noindent\begin{minipage}[t]{60mm}\vspace{0pt}\sbox\hmboard{\chessboard[setfen=6Q1/8/7k/8/3b4/8/p7/K7 w - - 0 1]}\usebox\hmboard\\[1mm]
\makebox[\wd\hmboard]{\sffamily\hspace{3mm}\textbf{h\#5.5}\hfill(2+3)}\end{minipage}\hfill
\begin{minipage}[t]{52mm}\vspace{0pt}
\textcolor{ink}{\sffamily\textsc{Lösung}}\par\smallskip\small
\begin{tabular}{@{}r@{\enspace}l@{\enspace}l@{}}
\textbf{1.} & \ldots & \symking{}xa2\\
\textbf{2.} & \symking{}h5 & \symking{}b1\\
\textbf{3.} & \symking{}h4 & \symqueen{}a2\\
\textbf{4.} & \symking{}g3 & \symking{}c2\\
\textbf{5.} & \symking{}f2 & \symking{}d3+\\
\textbf{6.} & \symking{}e1 & \symqueen{}e2\#\\
\end{tabular}\par\vspace{3mm}
\qrcode[height=12mm]{https://helpman.komtera.lt/?fen=6Q1/8/7k/8/3b4/8/p7/K7&moves=5.5}\\{\sffamily\tiny\color{ink}scannen zum Lösen}
\end{minipage}
\vspace{2mm}

{\raggedright\noindent\textcolor{ink}{\sffamily\textsc{Themen}}\enspace \themetag{ein Stein zieht}\par\smallskip{\footnotesize\textit{Hinweis:} die Partei am Zug steht im Diagramm im Schach; die Lösung beginnt mit einem Schlag.}\par}

\clearpage
\phantomsection\label{hm:KQvkqb}
\begin{tcolorbox}[enhanced,colback=ink,colframe=ink,arc=0pt,boxrule=0pt,left=4mm,right=4mm,top=2.5mm,bottom=2.5mm]
\color{white}\sffamily{\LARGE\bfseries KQvkqb}\hfill{\Large h\#5.5}\\[0.5mm]
{\small 5 Steine\quad Weiß am Zug}\hfill{\small Nr.~197\qquad\href{https://helpman.komtera.lt/?fen=8/8/7k/8/1b6/q7/8/KQ6&moves=5.5}{\color{white}\bfseries lösen\,\ensuremath{\nearrow}}}
\end{tcolorbox}
\vspace{2mm}
\noindent\begin{minipage}[t]{36mm}\vspace{0pt}\sffamily\small\raggedright
\hmfigure{längste eindeutige}{{\Large h\#5.5}}\hmfigure{eindeutig bei dieser Länge}{{\Large 25} Stellungen}\hmfigure{längstes Matt}{{\Large h\#7.5}}\hmfigure{in der Klasse}{242 Millionen Stellungen,\\ 126 Millionen mit Hilfsmatt}
\end{minipage}\hfill
\begin{minipage}[t]{80mm}\vspace{0pt}\begin{tikzpicture}\begin{axis}[width=72mm,height=42mm,ybar,bar width=1.6pt,xmin=-1,xmax=16,ymin=0,ymax=30.04,axis y line*=left,axis x line*=bottom,xlabel={\scriptsize Halbzüge bis Matt},ylabel={\scriptsize Stellungen (Mio.)},tick label style={font=\tiny},label style={font=\scriptsize},enlarge x limits=false,ylabel shift=-3pt,xlabel shift=-3pt]\addplot[fill=ink!55,draw=none] coordinates {(0,0.03541) (1,0.322) (2,1.245) (3,5.52) (4,11.46) (5,23.06) (6,24.91) (7,27.31) (8,17.22) (9,11.07) (10,3.195) (11,0.6683) (12,0.06417) (13,0.005616) (14,4.9e-05) (15,3.3e-05)};\draw[densely dotted] (axis cs:11,0) -- (axis cs:11,30.04);\end{axis}\begin{axis}[width=72mm,height=42mm,axis y line*=right,axis x line=none,ymode=log,xmin=-1,xmax=16,ylabel={\scriptsize eindeutig (log)},tick label style={font=\tiny},label style={font=\scriptsize},ylabel shift=-3pt]\addplot[mark=*,mark size=0.7pt,draw=black,thin] coordinates {(0,35412) (1,277000) (2,595988) (3,1267470) (4,1203226) (5,1107985) (6,550246) (7,411419) (8,118512) (9,25651) (10,1707) (11,25)};\end{axis}\end{tikzpicture}\end{minipage}
\vspace{3mm}

\noindent\begin{minipage}[t]{60mm}\vspace{0pt}\sbox\hmboard{\chessboard[setfen=8/8/7k/8/1b6/q7/8/KQ6 w - - 0 1]}\usebox\hmboard\\[1mm]
\makebox[\wd\hmboard]{\sffamily\hspace{3mm}\textbf{h\#5.5}\hfill(2+3)}\end{minipage}\hfill
\begin{minipage}[t]{52mm}\vspace{0pt}
\textcolor{ink}{\sffamily\textsc{Lösung}}\par\smallskip\small
\begin{tabular}{@{}r@{\enspace}l@{\enspace}l@{}}
\textbf{1.} & \ldots & \symqueen{}a2\\
\textbf{2.} & \symking{}g5 & \symqueen{}xa3\\
\textbf{3.} & \symking{}f4 & \symqueen{}a2\\
\textbf{4.} & \symking{}e3 & \symking{}b2\\
\textbf{5.} & \symking{}d2 & \symking{}b3+\\
\textbf{6.} & \symking{}c1 & \symqueen{}c2\#\\
\end{tabular}\par\vspace{3mm}
\qrcode[height=12mm]{https://helpman.komtera.lt/?fen=8/8/7k/8/1b6/q7/8/KQ6&moves=5.5}\\{\sffamily\tiny\color{ink}scannen zum Lösen}
\end{minipage}
\vspace{2mm}

{\raggedright\noindent\textcolor{ink}{\sffamily\textsc{Themen}}\enspace \themetag{Switchback} \themetag{ein Stein zieht}\par\smallskip{\footnotesize\textit{Hinweis:} die Partei am Zug steht im Diagramm im Schach.}\par}

\clearpage
\phantomsection\label{hm:KQvkqq}
\begin{tcolorbox}[enhanced,colback=ink,colframe=ink,arc=0pt,boxrule=0pt,left=4mm,right=4mm,top=2.5mm,bottom=2.5mm]
\color{white}\sffamily{\LARGE\bfseries KQvkqq}\hfill{\Large h\#5.5}\\[0.5mm]
{\small 5 Steine\quad Weiß am Zug}\hfill{\small Nr.~198\qquad\href{https://helpman.komtera.lt/?fen=8/7k/8/8/8/1q6/Q7/K2q4&moves=5.5}{\color{white}\bfseries lösen\,\ensuremath{\nearrow}}}
\end{tcolorbox}
\vspace{2mm}
\noindent\begin{minipage}[t]{36mm}\vspace{0pt}\sffamily\small\raggedright
\hmfigure{längste eindeutige}{{\Large h\#5.5}}\hmfigure{eindeutig bei dieser Länge}{{\Large 58} Stellungen}\hmfigure{längstes Matt}{{\Large h\#7.5}}\hmfigure{in der Klasse}{242 Millionen Stellungen,\\ 54,9 Millionen mit Hilfsmatt}
\end{minipage}\hfill
\begin{minipage}[t]{80mm}\vspace{0pt}\begin{tikzpicture}\begin{axis}[width=72mm,height=42mm,ybar,bar width=1.6pt,xmin=-1,xmax=16,ymin=0,ymax=15.04,axis y line*=left,axis x line*=bottom,xlabel={\scriptsize Halbzüge bis Matt},ylabel={\scriptsize Stellungen (Mio.)},tick label style={font=\tiny},label style={font=\scriptsize},enlarge x limits=false,ylabel shift=-3pt,xlabel shift=-3pt]\addplot[fill=ink!55,draw=none] coordinates {(0,0.009051) (1,0.09151) (2,0.4147) (3,2.15) (4,4.237) (5,10.72) (6,9.434) (7,13.67) (8,6.394) (9,6.135) (10,1.162) (11,0.4734) (12,0.0116) (13,0.008039) (14,6e-06) (15,3.3e-05)};\draw[densely dotted] (axis cs:11,0) -- (axis cs:11,15.04);\end{axis}\begin{axis}[width=72mm,height=42mm,axis y line*=right,axis x line=none,ymode=log,xmin=-1,xmax=16,ylabel={\scriptsize eindeutig (log)},tick label style={font=\tiny},label style={font=\scriptsize},ylabel shift=-3pt]\addplot[mark=*,mark size=0.7pt,draw=black,thin] coordinates {(0,9051) (1,82940) (2,202663) (3,537915) (4,451953) (5,584793) (6,222970) (7,226870) (8,47673) (9,19090) (10,479) (11,58)};\end{axis}\end{tikzpicture}\end{minipage}
\vspace{3mm}

\noindent\begin{minipage}[t]{60mm}\vspace{0pt}\sbox\hmboard{\chessboard[setfen=8/7k/8/8/8/1q6/Q7/K2q4 w - - 0 1]}\usebox\hmboard\\[1mm]
\makebox[\wd\hmboard]{\sffamily\hspace{3mm}\textbf{h\#5.5}\hfill(2+3)}\end{minipage}\hfill
\begin{minipage}[t]{52mm}\vspace{0pt}
\textcolor{ink}{\sffamily\textsc{Lösung}}\par\smallskip\small
\begin{tabular}{@{}r@{\enspace}l@{\enspace}l@{}}
\textbf{1.} & \ldots & \symqueen{}b1+\\
\textbf{2.} & \symqueen{}dd3 & \symqueen{}e1\\
\textbf{3.} & \symqueen{}f7 & \symking{}b2\\
\textbf{4.} & \symqueen{}dg6 & \symking{}c3\\
\textbf{5.} & \symking{}g7 & \symking{}d4\\
\textbf{6.} & \symking{}f6 & \symqueen{}e5\#\\
\end{tabular}\par\vspace{3mm}
\qrcode[height=12mm]{https://helpman.komtera.lt/?fen=8/7k/8/8/8/1q6/Q7/K2q4&moves=5.5}\\{\sffamily\tiny\color{ink}scannen zum Lösen}
\end{minipage}
\vspace{2mm}

{\raggedright\noindent\textcolor{ink}{\sffamily\textsc{Themen}}\enspace \themetag{reines Matt} \themetag{Mustermatt} \themetag{Idealmatt} \themetag{Selbstblock} \themetag{ohne Schlag}\par\smallskip{\footnotesize\textit{Hinweis:} die Partei am Zug steht im Diagramm im Schach.}\par}

\clearpage
\phantomsection\label{hm:KQvkqr}
\begin{tcolorbox}[enhanced,colback=ink,colframe=ink,arc=0pt,boxrule=0pt,left=4mm,right=4mm,top=2.5mm,bottom=2.5mm]
\color{white}\sffamily{\LARGE\bfseries KQvkqr}\hfill{\Large h\#5.5}\\[0.5mm]
{\small 5 Steine\quad Weiß am Zug}\hfill{\small Nr.~199\qquad\href{https://helpman.komtera.lt/?fen=8/8/5k2/8/6Q1/q7/8/Kr6&moves=5.5}{\color{white}\bfseries lösen\,\ensuremath{\nearrow}}}
\end{tcolorbox}
\vspace{2mm}
\noindent\begin{minipage}[t]{36mm}\vspace{0pt}\sffamily\small\raggedright
\hmfigure{längste eindeutige}{{\Large h\#5.5}}\hmfigure{eindeutig bei dieser Länge}{{\Large 10} Stellungen}\hmfigure{längstes Matt}{{\Large h\#6.5}}\hmfigure{in der Klasse}{242 Millionen Stellungen,\\ 120 Millionen mit Hilfsmatt}
\end{minipage}\hfill
\begin{minipage}[t]{80mm}\vspace{0pt}\begin{tikzpicture}\begin{axis}[width=72mm,height=42mm,ybar,bar width=1.6pt,xmin=-1,xmax=14,ymin=0,ymax=31.63,axis y line*=left,axis x line*=bottom,xlabel={\scriptsize Halbzüge bis Matt},ylabel={\scriptsize Stellungen (Mio.)},tick label style={font=\tiny},label style={font=\scriptsize},enlarge x limits=false,ylabel shift=-3pt,xlabel shift=-3pt]\addplot[fill=ink!55,draw=none] coordinates {(0,0.02929) (1,0.2811) (2,1.313) (3,6.074) (4,15.16) (5,28.75) (6,28.34) (7,28.64) (8,7.531) (9,3.749) (10,0.1473) (11,0.04284) (13,0.003047)};\draw[densely dotted] (axis cs:11,0) -- (axis cs:11,31.63);\end{axis}\begin{axis}[width=72mm,height=42mm,axis y line*=right,axis x line=none,ymode=log,xmin=-1,xmax=14,ylabel={\scriptsize eindeutig (log)},tick label style={font=\tiny},label style={font=\scriptsize},ylabel shift=-3pt]\addplot[mark=*,mark size=0.7pt,draw=black,thin] coordinates {(0,29287) (1,246344) (2,702427) (3,1347619) (4,1196062) (5,970414) (6,257934) (7,133588) (8,6352) (9,1985) (11,10)};\end{axis}\end{tikzpicture}\end{minipage}
\vspace{3mm}

\noindent\begin{minipage}[t]{60mm}\vspace{0pt}\sbox\hmboard{\chessboard[setfen=8/8/5k2/8/6Q1/q7/8/Kr6 w - - 0 1]}\usebox\hmboard\\[1mm]
\makebox[\wd\hmboard]{\sffamily\hspace{3mm}\textbf{h\#5.5}\hfill(2+3)}\end{minipage}\hfill
\begin{minipage}[t]{52mm}\vspace{0pt}
\textcolor{ink}{\sffamily\textsc{Lösung}}\par\smallskip\small
\begin{tabular}{@{}r@{\enspace}l@{\enspace}l@{}}
\textbf{1.} & \ldots & \symking{}xb1\\
\textbf{2.} & \symqueen{}d3+ & \symking{}c1\\
\textbf{3.} & \symqueen{}g6 & \symking{}d2\\
\textbf{4.} & \symking{}g7 & \symking{}e3\\
\textbf{5.} & \symking{}h6 & \symking{}f4\\
\textbf{6.} & \symqueen{}h7 & \symqueen{}g5\#\\
\end{tabular}\par\vspace{3mm}
\qrcode[height=12mm]{https://helpman.komtera.lt/?fen=8/8/5k2/8/6Q1/q7/8/Kr6&moves=5.5}\\{\sffamily\tiny\color{ink}scannen zum Lösen}
\end{minipage}
\vspace{2mm}

{\raggedright\noindent\textcolor{ink}{\sffamily\textsc{Themen}}\enspace \themetag{reines Matt} \themetag{Mustermatt} \themetag{Idealmatt} \themetag{Selbstblock}\par\smallskip{\footnotesize\textit{Hinweis:} die Partei am Zug steht im Diagramm im Schach; die Lösung beginnt mit einem Schlag.}\par}

\clearpage
\phantomsection\label{hm:KQvkrn}
\begin{tcolorbox}[enhanced,colback=ink,colframe=ink,arc=0pt,boxrule=0pt,left=4mm,right=4mm,top=2.5mm,bottom=2.5mm]
\color{white}\sffamily{\LARGE\bfseries KQvkrn}\hfill{\Large h\#5.5}\\[0.5mm]
{\small 5 Steine\quad Weiß am Zug}\hfill{\small Nr.~200\qquad\href{https://helpman.komtera.lt/?fen=8/6k1/8/8/8/8/8/Krn4Q&moves=5.5}{\color{white}\bfseries lösen\,\ensuremath{\nearrow}}}
\end{tcolorbox}
\vspace{2mm}
\noindent\begin{minipage}[t]{36mm}\vspace{0pt}\sffamily\small\raggedright
\hmfigure{längste eindeutige}{{\Large h\#5.5}}\hmfigure{eindeutig bei dieser Länge}{{\Large 2} Stellungen}\hmfigure{längstes Matt}{{\Large h\#6.5}}\hmfigure{in der Klasse}{242 Millionen Stellungen,\\ 142 Millionen mit Hilfsmatt}
\end{minipage}\hfill
\begin{minipage}[t]{80mm}\vspace{0pt}\begin{tikzpicture}\begin{axis}[width=72mm,height=42mm,ybar,bar width=1.6pt,xmin=-1,xmax=14,ymin=0,ymax=36.42,axis y line*=left,axis x line*=bottom,xlabel={\scriptsize Halbzüge bis Matt},ylabel={\scriptsize Stellungen (Mio.)},tick label style={font=\tiny},label style={font=\scriptsize},enlarge x limits=false,ylabel shift=-3pt,xlabel shift=-3pt]\addplot[fill=ink!55,draw=none] coordinates {(0,0.0664) (1,0.5399) (2,2.006) (3,7.272) (4,18.63) (5,25.75) (6,33.11) (7,26.2) (8,18.64) (9,8.219) (10,1.56) (11,0.1305) (12,0.006259) (13,0.000177)};\draw[densely dotted] (axis cs:11,0) -- (axis cs:11,36.42);\end{axis}\begin{axis}[width=72mm,height=42mm,axis y line*=right,axis x line=none,ymode=log,xmin=-1,xmax=14,ylabel={\scriptsize eindeutig (log)},tick label style={font=\tiny},label style={font=\scriptsize},ylabel shift=-3pt]\addplot[mark=*,mark size=0.7pt,draw=black,thin] coordinates {(0,66397) (1,427061) (2,1021936) (3,1353502) (4,1486722) (5,963566) (6,574820) (7,268466) (8,85921) (9,10521) (10,642) (11,2)};\end{axis}\end{tikzpicture}\end{minipage}
\vspace{3mm}

\noindent\begin{minipage}[t]{60mm}\vspace{0pt}\sbox\hmboard{\chessboard[setfen=8/6k1/8/8/8/8/8/Krn4Q w - - 0 1]}\usebox\hmboard\\[1mm]
\makebox[\wd\hmboard]{\sffamily\hspace{3mm}\textbf{h\#5.5}\hfill(2+3)}\end{minipage}\hfill
\begin{minipage}[t]{52mm}\vspace{0pt}
\textcolor{ink}{\sffamily\textsc{Lösung}}\par\smallskip\small
\begin{tabular}{@{}r@{\enspace}l@{\enspace}l@{}}
\textbf{1.} & \ldots & \symking{}xb1\\
\textbf{2.} & \symknight{}e2 & \symking{}c2\\
\textbf{3.} & \symknight{}g1 & \symking{}d3\\
\textbf{4.} & \symknight{}h3 & \symking{}e4\\
\textbf{5.} & \symking{}h6 & \symking{}f5\\
\textbf{6.} & \symking{}h5 & \symqueen{}xh3\#\\
\end{tabular}\par\vspace{3mm}
\qrcode[height=12mm]{https://helpman.komtera.lt/?fen=8/6k1/8/8/8/8/8/Krn4Q&moves=5.5}\\{\sffamily\tiny\color{ink}scannen zum Lösen}
\end{minipage}
\vspace{2mm}

{\raggedright\noindent\textcolor{ink}{\sffamily\textsc{Themen}}\enspace \themetag{Spiegelmatt} \themetag{ohne Schach}\par\smallskip{\footnotesize\textit{Hinweis:} die Partei am Zug steht im Diagramm im Schach; die Lösung beginnt mit einem Schlag.}\par}

\clearpage
\phantomsection\label{hm:KRNvkq}
\begin{tcolorbox}[enhanced,colback=ink,colframe=ink,arc=0pt,boxrule=0pt,left=4mm,right=4mm,top=2.5mm,bottom=2.5mm]
\color{white}\sffamily{\LARGE\bfseries KRNvkq}\hfill{\Large h\#5.5}\\[0.5mm]
{\small 5 Steine\quad Weiß am Zug}\hfill{\small Nr.~201\qquad\href{https://helpman.komtera.lt/?fen=8/N7/8/7k/R7/8/8/Kq6&moves=5.5}{\color{white}\bfseries lösen\,\ensuremath{\nearrow}}}
\end{tcolorbox}
\vspace{2mm}
\noindent\begin{minipage}[t]{36mm}\vspace{0pt}\sffamily\small\raggedright
\hmfigure{längste eindeutige}{{\Large h\#5.5}}\hmfigure{eindeutig bei dieser Länge}{{\Large 2} Stellungen}\hmfigure{längstes Matt}{{\Large h\#7}}\hmfigure{in der Klasse}{242 Millionen Stellungen,\\ 143 Millionen mit Hilfsmatt}
\end{minipage}\hfill
\begin{minipage}[t]{80mm}\vspace{0pt}\begin{tikzpicture}\begin{axis}[width=72mm,height=42mm,ybar,bar width=1.6pt,xmin=-1,xmax=15,ymin=0,ymax=43.89,axis y line*=left,axis x line*=bottom,xlabel={\scriptsize Halbzüge bis Matt},ylabel={\scriptsize Stellungen (Mio.)},tick label style={font=\tiny},label style={font=\scriptsize},enlarge x limits=false,ylabel shift=-3pt,xlabel shift=-3pt]\addplot[fill=ink!55,draw=none] coordinates {(0,0.02839) (1,0.1795) (2,0.7831) (3,4.196) (4,10.25) (5,28.16) (6,39.9) (7,37.21) (8,16.63) (9,5.006) (10,0.4791) (11,0.09662) (12,0.000692) (13,0.00229) (14,2.6e-05)};\draw[densely dotted] (axis cs:11,0) -- (axis cs:11,43.89);\end{axis}\begin{axis}[width=72mm,height=42mm,axis y line*=right,axis x line=none,ymode=log,xmin=-1,xmax=15,ylabel={\scriptsize eindeutig (log)},tick label style={font=\tiny},label style={font=\scriptsize},ylabel shift=-3pt]\addplot[mark=*,mark size=0.7pt,draw=black,thin] coordinates {(0,28386) (1,177357) (2,400366) (3,1589319) (4,2514298) (5,2645008) (6,928745) (7,234296) (8,16439) (9,1755) (10,5) (11,2)};\end{axis}\end{tikzpicture}\end{minipage}
\vspace{3mm}

\noindent\begin{minipage}[t]{60mm}\vspace{0pt}\sbox\hmboard{\chessboard[setfen=8/N7/8/7k/R7/8/8/Kq6 w - - 0 1]}\usebox\hmboard\\[1mm]
\makebox[\wd\hmboard]{\sffamily\hspace{3mm}\textbf{h\#5.5}\hfill(3+2)}\end{minipage}\hfill
\begin{minipage}[t]{52mm}\vspace{0pt}
\textcolor{ink}{\sffamily\textsc{Lösung}}\par\smallskip\small
\begin{tabular}{@{}r@{\enspace}l@{\enspace}l@{}}
\textbf{1.} & \ldots & \symking{}xb1\\
\textbf{2.} & \symking{}g5 & \symrook{}a6\\
\textbf{3.} & \symking{}f4 & \symknight{}b5\\
\textbf{4.} & \symking{}e3 & \symrook{}e6+\\
\textbf{5.} & \symking{}d2 & \symrook{}e2+\\
\textbf{6.} & \symking{}d1 & \symknight{}c3\#\\
\end{tabular}\par\vspace{3mm}
\qrcode[height=12mm]{https://helpman.komtera.lt/?fen=8/N7/8/7k/R7/8/8/Kq6&moves=5.5}\\{\sffamily\tiny\color{ink}scannen zum Lösen}
\end{minipage}
\vspace{2mm}

{\raggedright\noindent\textcolor{ink}{\sffamily\textsc{Themen}}\enspace \themetag{ein Stein zieht}\par\smallskip{\footnotesize\textit{Hinweis:} die Partei am Zug steht im Diagramm im Schach; die Lösung beginnt mit einem Schlag.}\par}

\clearpage
\phantomsection\label{hm:KQBvkr}
\begin{tcolorbox}[enhanced,colback=ink,colframe=ink,arc=0pt,boxrule=0pt,left=4mm,right=4mm,top=2.5mm,bottom=2.5mm]
\color{white}\sffamily{\LARGE\bfseries KQBvkr}\hfill{\Large h\#5}\\[0.5mm]
{\small 5 Steine\quad Schwarz am Zug}\hfill{\small Nr.~202\qquad\href{https://helpman.komtera.lt/?fen=8/8/8/8/6kB/1r3Q2/8/K7&moves=5}{\color{white}\bfseries lösen\,\ensuremath{\nearrow}}}
\end{tcolorbox}
\vspace{2mm}
\noindent\begin{minipage}[t]{36mm}\vspace{0pt}\sffamily\small\raggedright
\hmfigure{längste eindeutige}{{\Large h\#5}}\hmfigure{eindeutig bei dieser Länge}{{\Large 15} Stellungen}\hmfigure{längstes Matt}{{\Large h\#6}}\hmfigure{in der Klasse}{242 Millionen Stellungen,\\ 139 Millionen mit Hilfsmatt}
\end{minipage}\hfill
\begin{minipage}[t]{80mm}\vspace{0pt}\begin{tikzpicture}\begin{axis}[width=72mm,height=42mm,ybar,bar width=1.6pt,xmin=-1,xmax=13,ymin=0,ymax=69.1,axis y line*=left,axis x line*=bottom,xlabel={\scriptsize Halbzüge bis Matt},ylabel={\scriptsize Stellungen (Mio.)},tick label style={font=\tiny},label style={font=\scriptsize},enlarge x limits=false,ylabel shift=-3pt,xlabel shift=-3pt]\addplot[fill=ink!55,draw=none] coordinates {(0,0.2726) (1,1.936) (2,9.106) (3,30) (4,62.82) (5,26.64) (6,7.949) (7,0.2417) (8,0.07113) (9,6e-06) (10,0.01699) (12,0.000173)};\draw[densely dotted] (axis cs:10,0) -- (axis cs:10,69.1);\end{axis}\begin{axis}[width=72mm,height=42mm,axis y line*=right,axis x line=none,ymode=log,xmin=-1,xmax=13,ylabel={\scriptsize eindeutig (log)},tick label style={font=\tiny},label style={font=\scriptsize},ylabel shift=-3pt]\addplot[mark=*,mark size=0.7pt,draw=black,thin] coordinates {(0,272570) (1,1667931) (2,5595744) (3,6081260) (4,2876673) (5,62743) (6,2352) (7,3) (8,593) (10,15)};\end{axis}\end{tikzpicture}\end{minipage}
\vspace{3mm}

\noindent\begin{minipage}[t]{60mm}\vspace{0pt}\sbox\hmboard{\chessboard[setfen=8/8/8/8/6kB/1r3Q2/8/K7 b - - 0 1]}\usebox\hmboard\\[1mm]
\makebox[\wd\hmboard]{\sffamily\hspace{3mm}\textbf{h\#5}\hfill(3+2)}\end{minipage}\hfill
\begin{minipage}[t]{52mm}\vspace{0pt}
\textcolor{ink}{\sffamily\textsc{Lösung}}\par\smallskip\small
\begin{tabular}{@{}r@{\enspace}l@{\enspace}l@{}}
\textbf{1.} & \symking{}xh4 & \symqueen{}xb3\\
\textbf{2.} & \symking{}g4 & \symqueen{}a2\\
\textbf{3.} & \symking{}f3 & \symking{}b2\\
\textbf{4.} & \symking{}e2 & \symking{}c3+\\
\textbf{5.} & \symking{}d1 & \symqueen{}d2\#\\
\end{tabular}\par\vspace{3mm}
\qrcode[height=12mm]{https://helpman.komtera.lt/?fen=8/8/8/8/6kB/1r3Q2/8/K7&moves=5}\\{\sffamily\tiny\color{ink}scannen zum Lösen}
\end{minipage}
\vspace{2mm}

{\raggedright\noindent\textcolor{ink}{\sffamily\textsc{Themen}}\enspace \themetag{Switchback} \themetag{ein Stein zieht}\par\smallskip{\footnotesize\textit{Hinweis:} die Partei am Zug steht im Diagramm im Schach; die Lösung beginnt mit einem Schlag.}\par}

\clearpage
\phantomsection\label{hm:KQQvkb}
\begin{tcolorbox}[enhanced,colback=ink,colframe=ink,arc=0pt,boxrule=0pt,left=4mm,right=4mm,top=2.5mm,bottom=2.5mm]
\color{white}\sffamily{\LARGE\bfseries KQQvkb}\hfill{\Large h\#5}\\[0.5mm]
{\small 5 Steine\quad Schwarz am Zug}\hfill{\small Nr.~203\qquad\href{https://helpman.komtera.lt/?fen=4k2b/2Q3Q1/8/8/8/8/8/1K6&moves=5}{\color{white}\bfseries lösen\,\ensuremath{\nearrow}}}
\end{tcolorbox}
\vspace{2mm}
\noindent\begin{minipage}[t]{36mm}\vspace{0pt}\sffamily\small\raggedright
\hmfigure{längste eindeutige}{{\Large h\#5}}\hmfigure{eindeutig bei dieser Länge}{{\Large 408} Stellungen}\hmfigure{längstes Matt}{{\Large h\#6}}\hmfigure{in der Klasse}{242 Millionen Stellungen,\\ 66,4 Millionen mit Hilfsmatt}
\end{minipage}\hfill
\begin{minipage}[t]{80mm}\vspace{0pt}\begin{tikzpicture}\begin{axis}[width=72mm,height=42mm,ybar,bar width=1.6pt,xmin=-1,xmax=13,ymin=0,ymax=25.05,axis y line*=left,axis x line*=bottom,xlabel={\scriptsize Halbzüge bis Matt},ylabel={\scriptsize Stellungen (Mio.)},tick label style={font=\tiny},label style={font=\scriptsize},enlarge x limits=false,ylabel shift=-3pt,xlabel shift=-3pt]\addplot[fill=ink!55,draw=none] coordinates {(0,1.339) (1,4.262) (2,18.83) (3,16.48) (4,22.78) (5,1.181) (6,0.6372) (7,1.2e-05) (8,0.6648) (9,1e-06) (10,0.2394) (12,0.02884)};\draw[densely dotted] (axis cs:10,0) -- (axis cs:10,25.05);\end{axis}\begin{axis}[width=72mm,height=42mm,axis y line*=right,axis x line=none,ymode=log,xmin=-1,xmax=13,ylabel={\scriptsize eindeutig (log)},tick label style={font=\tiny},label style={font=\scriptsize},ylabel shift=-3pt]\addplot[mark=*,mark size=0.7pt,draw=black,thin] coordinates {(0,1339030) (1,2353932) (2,6171397) (3,305137) (4,49252) (5,9) (6,28416) (8,9790) (10,408)};\end{axis}\end{tikzpicture}\end{minipage}
\vspace{3mm}

\noindent\begin{minipage}[t]{60mm}\vspace{0pt}\sbox\hmboard{\chessboard[setfen=4k2b/2Q3Q1/8/8/8/8/8/1K6 b - - 0 1]}\usebox\hmboard\\[1mm]
\makebox[\wd\hmboard]{\sffamily\hspace{3mm}\textbf{h\#5}\hfill(3+2)}\end{minipage}\hfill
\begin{minipage}[t]{52mm}\vspace{0pt}
\textcolor{ink}{\sffamily\textsc{Lösung}}\par\smallskip\small
\begin{tabular}{@{}r@{\enspace}l@{\enspace}l@{}}
\textbf{1.} & \symbishop{}xg7 & \symking{}c2\\
\textbf{2.} & \symking{}f8 & \symking{}d3\\
\textbf{3.} & \symking{}g8 & \symking{}e4\\
\textbf{4.} & \symking{}h7 & \symking{}f5\\
\textbf{5.} & \symking{}h6 & \symqueen{}h2\#\\
\end{tabular}\par\vspace{3mm}
\qrcode[height=12mm]{https://helpman.komtera.lt/?fen=4k2b/2Q3Q1/8/8/8/8/8/1K6&moves=5}\\{\sffamily\tiny\color{ink}scannen zum Lösen}
\end{minipage}
\vspace{2mm}

{\raggedright\noindent\textcolor{ink}{\sffamily\textsc{Themen}}\enspace \themetag{reines Matt} \themetag{Mustermatt} \themetag{Idealmatt} \themetag{Selbstblock} \themetag{ohne Schach}\par\smallskip{\footnotesize\textit{Hinweis:} die Lösung beginnt mit einem Schlag.}\par}

\clearpage
\phantomsection\label{hm:KQQvkq}
\begin{tcolorbox}[enhanced,colback=ink,colframe=ink,arc=0pt,boxrule=0pt,left=4mm,right=4mm,top=2.5mm,bottom=2.5mm]
\color{white}\sffamily{\LARGE\bfseries KQQvkq}\hfill{\Large h\#5}\\[0.5mm]
{\small 5 Steine\quad Schwarz am Zug}\hfill{\small Nr.~204\qquad\href{https://helpman.komtera.lt/?fen=8/8/8/4Q3/8/8/4q2k/K4Q2&moves=5}{\color{white}\bfseries lösen\,\ensuremath{\nearrow}}}
\end{tcolorbox}
\vspace{2mm}
\noindent\begin{minipage}[t]{36mm}\vspace{0pt}\sffamily\small\raggedright
\hmfigure{längste eindeutige}{{\Large h\#5}}\hmfigure{eindeutig bei dieser Länge}{{\Large 441} Stellungen}\hmfigure{längstes Matt}{{\Large h\#7}}\hmfigure{in der Klasse}{242 Millionen Stellungen,\\ 56 Millionen mit Hilfsmatt}
\end{minipage}\hfill
\begin{minipage}[t]{80mm}\vspace{0pt}\begin{tikzpicture}\begin{axis}[width=72mm,height=42mm,ybar,bar width=1.6pt,xmin=-1,xmax=15,ymin=0,ymax=20.15,axis y line*=left,axis x line*=bottom,xlabel={\scriptsize Halbzüge bis Matt},ylabel={\scriptsize Stellungen (Mio.)},tick label style={font=\tiny},label style={font=\scriptsize},enlarge x limits=false,ylabel shift=-3pt,xlabel shift=-3pt]\addplot[fill=ink!55,draw=none] coordinates {(0,0.815) (1,2.749) (2,13.65) (3,16.2) (4,18.32) (5,2.962) (6,0.495) (7,3e-06) (8,0.5417) (9,5e-06) (10,0.2259) (12,0.0261) (14,0.000478)};\draw[densely dotted] (axis cs:10,0) -- (axis cs:10,20.15);\end{axis}\begin{axis}[width=72mm,height=42mm,axis y line*=right,axis x line=none,ymode=log,xmin=-1,xmax=15,ylabel={\scriptsize eindeutig (log)},tick label style={font=\tiny},label style={font=\scriptsize},ylabel shift=-3pt]\addplot[mark=*,mark size=0.7pt,draw=black,thin] coordinates {(0,815022) (1,1711485) (2,4907898) (3,682320) (4,45414) (6,18497) (8,8600) (10,441)};\end{axis}\end{tikzpicture}\end{minipage}
\vspace{3mm}

\noindent\begin{minipage}[t]{60mm}\vspace{0pt}\sbox\hmboard{\chessboard[setfen=8/8/8/4Q3/8/8/4q2k/K4Q2 b - - 0 1]}\usebox\hmboard\\[1mm]
\makebox[\wd\hmboard]{\sffamily\hspace{3mm}\textbf{h\#5}\hfill(3+2)}\end{minipage}\hfill
\begin{minipage}[t]{52mm}\vspace{0pt}
\textcolor{ink}{\sffamily\textsc{Lösung}}\par\smallskip\small
\begin{tabular}{@{}r@{\enspace}l@{\enspace}l@{}}
\textbf{1.} & \symqueen{}xe5+ & \symking{}b1\\
\textbf{2.} & \symking{}g3 & \symqueen{}b5\\
\textbf{3.} & \symking{}f2 & \symking{}c2\\
\textbf{4.} & \symking{}e1 & \symking{}d3\\
\textbf{5.} & \symking{}d1 & \symqueen{}b1\#\\
\end{tabular}\par\vspace{3mm}
\qrcode[height=12mm]{https://helpman.komtera.lt/?fen=8/8/8/4Q3/8/8/4q2k/K4Q2&moves=5}\\{\sffamily\tiny\color{ink}scannen zum Lösen}
\end{minipage}
\vspace{2mm}

{\raggedright\noindent\textcolor{ink}{\sffamily\textsc{Themen}}\enspace \themetag{Spiegelmatt}\par\smallskip{\footnotesize\textit{Hinweis:} die Partei am Zug steht im Diagramm im Schach; die Lösung beginnt mit einem Schlag.}\par}

\clearpage
\phantomsection\label{hm:KQQvkr}
\begin{tcolorbox}[enhanced,colback=ink,colframe=ink,arc=0pt,boxrule=0pt,left=4mm,right=4mm,top=2.5mm,bottom=2.5mm]
\color{white}\sffamily{\LARGE\bfseries KQQvkr}\hfill{\Large h\#5}\\[0.5mm]
{\small 5 Steine\quad Schwarz am Zug}\hfill{\small Nr.~205\qquad\href{https://helpman.komtera.lt/?fen=8/8/2rQ4/7Q/6k1/8/8/K7&moves=5}{\color{white}\bfseries lösen\,\ensuremath{\nearrow}}}
\end{tcolorbox}
\vspace{2mm}
\noindent\begin{minipage}[t]{36mm}\vspace{0pt}\sffamily\small\raggedright
\hmfigure{längste eindeutige}{{\Large h\#5}}\hmfigure{eindeutig bei dieser Länge}{{\Large 46} Stellungen}\hmfigure{längstes Matt}{{\Large h\#6}}\hmfigure{in der Klasse}{242 Millionen Stellungen,\\ 62,7 Millionen mit Hilfsmatt}
\end{minipage}\hfill
\begin{minipage}[t]{80mm}\vspace{0pt}\begin{tikzpicture}\begin{axis}[width=72mm,height=42mm,ybar,bar width=1.6pt,xmin=-1,xmax=13,ymin=0,ymax=22.65,axis y line*=left,axis x line*=bottom,xlabel={\scriptsize Halbzüge bis Matt},ylabel={\scriptsize Stellungen (Mio.)},tick label style={font=\tiny},label style={font=\scriptsize},enlarge x limits=false,ylabel shift=-3pt,xlabel shift=-3pt]\addplot[fill=ink!55,draw=none] coordinates {(0,1.154) (1,3.796) (2,17.63) (3,16.37) (4,20.59) (5,1.753) (6,0.7086) (7,1e-05) (8,0.5712) (9,1e-06) (10,0.1236) (12,0.001156)};\draw[densely dotted] (axis cs:10,0) -- (axis cs:10,22.65);\end{axis}\begin{axis}[width=72mm,height=42mm,axis y line*=right,axis x line=none,ymode=log,xmin=-1,xmax=13,ylabel={\scriptsize eindeutig (log)},tick label style={font=\tiny},label style={font=\scriptsize},ylabel shift=-3pt]\addplot[mark=*,mark size=0.7pt,draw=black,thin] coordinates {(0,1153526) (1,2209926) (2,6038016) (3,392558) (4,45898) (6,17179) (8,4503) (10,46)};\end{axis}\end{tikzpicture}\end{minipage}
\vspace{3mm}

\noindent\begin{minipage}[t]{60mm}\vspace{0pt}\sbox\hmboard{\chessboard[setfen=8/8/2rQ4/7Q/6k1/8/8/K7 b - - 0 1]}\usebox\hmboard\\[1mm]
\makebox[\wd\hmboard]{\sffamily\hspace{3mm}\textbf{h\#5}\hfill(3+2)}\end{minipage}\hfill
\begin{minipage}[t]{52mm}\vspace{0pt}
\textcolor{ink}{\sffamily\textsc{Lösung}}\par\smallskip\small
\begin{tabular}{@{}r@{\enspace}l@{\enspace}l@{}}
\textbf{1.} & \symking{}xh5 & \symqueen{}xc6\\
\textbf{2.} & \symking{}g4 & \symqueen{}h6\\
\textbf{3.} & \symking{}f3 & \symking{}b2\\
\textbf{4.} & \symking{}e2 & \symking{}c3\\
\textbf{5.} & \symking{}d1 & \symqueen{}d2\#\\
\end{tabular}\par\vspace{3mm}
\qrcode[height=12mm]{https://helpman.komtera.lt/?fen=8/8/2rQ4/7Q/6k1/8/8/K7&moves=5}\\{\sffamily\tiny\color{ink}scannen zum Lösen}
\end{minipage}
\vspace{2mm}

{\raggedright\noindent\textcolor{ink}{\sffamily\textsc{Themen}}\enspace \themetag{Switchback} \themetag{ein Stein zieht} \themetag{ohne Schach}\par\smallskip{\footnotesize\textit{Hinweis:} die Partei am Zug steht im Diagramm im Schach; die Lösung beginnt mit einem Schlag.}\par}

\clearpage
\phantomsection\label{hm:KQRNvk}
\begin{tcolorbox}[enhanced,colback=ink,colframe=ink,arc=0pt,boxrule=0pt,left=4mm,right=4mm,top=2.5mm,bottom=2.5mm]
\color{white}\sffamily{\LARGE\bfseries KQRNvk}\hfill{\Large h\#5}\\[0.5mm]
{\small 5 Steine\quad Schwarz am Zug}\hfill{\small Nr.~206\qquad\href{https://helpman.komtera.lt/?fen=8/8/8/8/6Qk/N4R2/8/K7&moves=5}{\color{white}\bfseries lösen\,\ensuremath{\nearrow}}}
\end{tcolorbox}
\vspace{2mm}
\noindent\begin{minipage}[t]{36mm}\vspace{0pt}\sffamily\small\raggedright
\hmfigure{längste eindeutige}{{\Large h\#5}}\hmfigure{eindeutig bei dieser Länge}{{\Large 26} Stellungen}\hmfigure{längstes Matt}{{\Large h\#7}}\hmfigure{in der Klasse}{242 Millionen Stellungen,\\ 150 Millionen mit Hilfsmatt}
\end{minipage}\hfill
\begin{minipage}[t]{80mm}\vspace{0pt}\begin{tikzpicture}\begin{axis}[width=72mm,height=42mm,ybar,bar width=1.6pt,xmin=-1,xmax=15,ymin=0,ymax=78.39,axis y line*=left,axis x line*=bottom,xlabel={\scriptsize Halbzüge bis Matt},ylabel={\scriptsize Stellungen (Mio.)},tick label style={font=\tiny},label style={font=\scriptsize},enlarge x limits=false,ylabel shift=-3pt,xlabel shift=-3pt]\addplot[fill=ink!55,draw=none] coordinates {(0,2.013) (1,7.55) (2,23.81) (3,36.69) (4,71.27) (5,4.775) (6,2.833) (7,0.000103) (8,0.7986) (10,0.2655) (12,0.01044) (14,1.8e-05)};\draw[densely dotted] (axis cs:10,0) -- (axis cs:10,78.39);\end{axis}\begin{axis}[width=72mm,height=42mm,axis y line*=right,axis x line=none,ymode=log,xmin=-1,xmax=15,ylabel={\scriptsize eindeutig (log)},tick label style={font=\tiny},label style={font=\scriptsize},ylabel shift=-3pt]\addplot[mark=*,mark size=0.7pt,draw=black,thin] coordinates {(0,2013490) (1,4769034) (2,11794748) (3,3440450) (4,1295872) (5,76) (6,88113) (8,7022) (10,26)};\end{axis}\end{tikzpicture}\end{minipage}
\vspace{3mm}

\noindent\begin{minipage}[t]{60mm}\vspace{0pt}\sbox\hmboard{\chessboard[setfen=8/8/8/8/6Qk/N4R2/8/K7 b - - 0 1]}\usebox\hmboard\\[1mm]
\makebox[\wd\hmboard]{\sffamily\hspace{3mm}\textbf{h\#5}\hfill(4+1)}\end{minipage}\hfill
\begin{minipage}[t]{52mm}\vspace{0pt}
\textcolor{ink}{\sffamily\textsc{Lösung}}\par\smallskip\small
\begin{tabular}{@{}r@{\enspace}l@{\enspace}l@{}}
\textbf{1.} & \symking{}xg4 & \symrook{}b3\\
\textbf{2.} & \symking{}f4 & \symrook{}b2\\
\textbf{3.} & \symking{}e3 & \symknight{}c2+\\
\textbf{4.} & \symking{}d2 & \symknight{}e3+\\
\textbf{5.} & \symking{}c1 & \symrook{}c2\#\\
\end{tabular}\par\vspace{3mm}
\qrcode[height=12mm]{https://helpman.komtera.lt/?fen=8/8/8/8/6Qk/N4R2/8/K7&moves=5}\\{\sffamily\tiny\color{ink}scannen zum Lösen}
\end{minipage}
\vspace{2mm}

{\raggedright\noindent\textcolor{ink}{\sffamily\textsc{Themen}}\enspace \themetag{ein Stein zieht} \themetag{Umnov}\par\smallskip{\footnotesize\textit{Hinweis:} die Partei am Zug steht im Diagramm im Schach; die Lösung beginnt mit einem Schlag.}\par}

\clearpage
\phantomsection\label{hm:KQRvkb}
\begin{tcolorbox}[enhanced,colback=ink,colframe=ink,arc=0pt,boxrule=0pt,left=4mm,right=4mm,top=2.5mm,bottom=2.5mm]
\color{white}\sffamily{\LARGE\bfseries KQRvkb}\hfill{\Large h\#5}\\[0.5mm]
{\small 5 Steine\quad Schwarz am Zug}\hfill{\small Nr.~207\qquad\href{https://helpman.komtera.lt/?fen=7b/6Q1/8/8/8/8/6R1/K6k&moves=5}{\color{white}\bfseries lösen\,\ensuremath{\nearrow}}}
\end{tcolorbox}
\vspace{2mm}
\noindent\begin{minipage}[t]{36mm}\vspace{0pt}\sffamily\small\raggedright
\hmfigure{längste eindeutige}{{\Large h\#5}}\hmfigure{eindeutig bei dieser Länge}{{\Large 1\,113} Stellungen}\hmfigure{längstes Matt}{{\Large h\#7}}\hmfigure{in der Klasse}{242 Millionen Stellungen,\\ 142 Millionen mit Hilfsmatt}
\end{minipage}\hfill
\begin{minipage}[t]{80mm}\vspace{0pt}\begin{tikzpicture}\begin{axis}[width=72mm,height=42mm,ybar,bar width=1.6pt,xmin=-1,xmax=15,ymin=0,ymax=73.86,axis y line*=left,axis x line*=bottom,xlabel={\scriptsize Halbzüge bis Matt},ylabel={\scriptsize Stellungen (Mio.)},tick label style={font=\tiny},label style={font=\scriptsize},enlarge x limits=false,ylabel shift=-3pt,xlabel shift=-3pt]\addplot[fill=ink!55,draw=none] coordinates {(0,1.073) (1,4.364) (2,15.98) (3,35.72) (4,67.14) (5,13.06) (6,3.59) (7,0.009494) (8,0.6962) (9,8e-06) (10,0.4722) (11,1e-06) (12,0.08254) (14,0.001481)};\draw[densely dotted] (axis cs:10,0) -- (axis cs:10,73.86);\end{axis}\begin{axis}[width=72mm,height=42mm,axis y line*=right,axis x line=none,ymode=log,xmin=-1,xmax=15,ylabel={\scriptsize eindeutig (log)},tick label style={font=\tiny},label style={font=\scriptsize},ylabel shift=-3pt]\addplot[mark=*,mark size=0.7pt,draw=black,thin] coordinates {(0,1072517) (1,3064642) (2,6112721) (3,3645804) (4,1685498) (5,3410) (6,41415) (7,4) (8,7214) (10,1113)};\end{axis}\end{tikzpicture}\end{minipage}
\vspace{3mm}

\noindent\begin{minipage}[t]{60mm}\vspace{0pt}\sbox\hmboard{\chessboard[setfen=7b/6Q1/8/8/8/8/6R1/K6k b - - 0 1]}\usebox\hmboard\\[1mm]
\makebox[\wd\hmboard]{\sffamily\hspace{3mm}\textbf{h\#5}\hfill(3+2)}\end{minipage}\hfill
\begin{minipage}[t]{52mm}\vspace{0pt}
\textcolor{ink}{\sffamily\textsc{Lösung}}\par\smallskip\small
\begin{tabular}{@{}r@{\enspace}l@{\enspace}l@{}}
\textbf{1.} & \symbishop{}xg7+ & \symrook{}b2\\
\textbf{2.} & \symking{}g1 & \symking{}b1\\
\textbf{3.} & \symking{}f1 & \symking{}c2\\
\textbf{4.} & \symking{}e1 & \symking{}d3\\
\textbf{5.} & \symking{}d1 & \symrook{}b1\#\\
\end{tabular}\par\vspace{3mm}
\qrcode[height=12mm]{https://helpman.komtera.lt/?fen=7b/6Q1/8/8/8/8/6R1/K6k&moves=5}\\{\sffamily\tiny\color{ink}scannen zum Lösen}
\end{minipage}
\vspace{2mm}

{\raggedright\noindent\textcolor{ink}{\sffamily\textsc{Themen}}\enspace \themetag{reines Matt} \themetag{Mustermatt} \themetag{Spiegelmatt}\par\smallskip{\footnotesize\textit{Hinweis:} die Lösung beginnt mit einem Schlag.}\par}

\clearpage
\phantomsection\label{hm:KQvkrb}
\begin{tcolorbox}[enhanced,colback=ink,colframe=ink,arc=0pt,boxrule=0pt,left=4mm,right=4mm,top=2.5mm,bottom=2.5mm]
\color{white}\sffamily{\LARGE\bfseries KQvkrb}\hfill{\Large h\#5}\\[0.5mm]
{\small 5 Steine\quad Schwarz am Zug}\hfill{\small Nr.~208\qquad\href{https://helpman.komtera.lt/?fen=8/2r5/8/8/8/1b5k/6Q1/K7&moves=5}{\color{white}\bfseries lösen\,\ensuremath{\nearrow}}}
\end{tcolorbox}
\vspace{2mm}
\noindent\begin{minipage}[t]{36mm}\vspace{0pt}\sffamily\small\raggedright
\hmfigure{längste eindeutige}{{\Large h\#5}}\hmfigure{eindeutig bei dieser Länge}{{\Large 23} Stellungen}\hmfigure{längstes Matt}{{\Large h\#6.5}}\hmfigure{in der Klasse}{242 Millionen Stellungen,\\ 138 Millionen mit Hilfsmatt}
\end{minipage}\hfill
\begin{minipage}[t]{80mm}\vspace{0pt}\begin{tikzpicture}\begin{axis}[width=72mm,height=42mm,ybar,bar width=1.6pt,xmin=-1,xmax=14,ymin=0,ymax=39.46,axis y line*=left,axis x line*=bottom,xlabel={\scriptsize Halbzüge bis Matt},ylabel={\scriptsize Stellungen (Mio.)},tick label style={font=\tiny},label style={font=\scriptsize},enlarge x limits=false,ylabel shift=-3pt,xlabel shift=-3pt]\addplot[fill=ink!55,draw=none] coordinates {(0,0.05808) (1,0.4944) (2,1.974) (3,7.689) (4,20.81) (5,29.61) (6,35.87) (7,26.38) (8,11.21) (9,3.92) (10,0.2931) (11,0.008664) (12,0.000242) (13,0.000162)};\draw[densely dotted] (axis cs:10,0) -- (axis cs:10,39.46);\end{axis}\begin{axis}[width=72mm,height=42mm,axis y line*=right,axis x line=none,ymode=log,xmin=-1,xmax=14,ylabel={\scriptsize eindeutig (log)},tick label style={font=\tiny},label style={font=\scriptsize},ylabel shift=-3pt]\addplot[mark=*,mark size=0.7pt,draw=black,thin] coordinates {(0,58079) (1,405718) (2,1052812) (3,1457020) (4,1593699) (5,933297) (6,370760) (7,151188) (8,16695) (9,615) (10,23)};\end{axis}\end{tikzpicture}\end{minipage}
\vspace{3mm}

\noindent\begin{minipage}[t]{60mm}\vspace{0pt}\sbox\hmboard{\chessboard[setfen=8/2r5/8/8/8/1b5k/6Q1/K7 b - - 0 1]}\usebox\hmboard\\[1mm]
\makebox[\wd\hmboard]{\sffamily\hspace{3mm}\textbf{h\#5}\hfill(2+3)}\end{minipage}\hfill
\begin{minipage}[t]{52mm}\vspace{0pt}
\textcolor{ink}{\sffamily\textsc{Lösung}}\par\smallskip\small
\begin{tabular}{@{}r@{\enspace}l@{\enspace}l@{}}
\textbf{1.} & \symking{}h4 & \symqueen{}c6\\
\textbf{2.} & \symking{}g3 & \symking{}b2\\
\textbf{3.} & \symking{}f2 & \symking{}c3\\
\textbf{4.} & \symking{}e1 & \symking{}d3\\
\textbf{5.} & \symking{}d1 & \symqueen{}h1\#\\
\end{tabular}\par\vspace{3mm}
\qrcode[height=12mm]{https://helpman.komtera.lt/?fen=8/2r5/8/8/8/1b5k/6Q1/K7&moves=5}\\{\sffamily\tiny\color{ink}scannen zum Lösen}
\end{minipage}
\vspace{2mm}

{\raggedright\noindent\textcolor{ink}{\sffamily\textsc{Themen}}\enspace \themetag{reines Matt} \themetag{Mustermatt} \themetag{Spiegelmatt} \themetag{ein Stein zieht} \themetag{ohne Schlag} \themetag{ohne Schach}\par\smallskip{\footnotesize\textit{Hinweis:} die Partei am Zug steht im Diagramm im Schach.}\par}

\clearpage
\phantomsection\label{hm:KQvkrp}
\begin{tcolorbox}[enhanced,colback=ink,colframe=ink,arc=0pt,boxrule=0pt,left=4mm,right=4mm,top=2.5mm,bottom=2.5mm]
\color{white}\sffamily{\LARGE\bfseries KQvkrp}\hfill{\Large h\#5}\\[0.5mm]
{\small 5 Steine\quad Schwarz am Zug}\hfill{\small Nr.~209\qquad\href{https://helpman.komtera.lt/?fen=8/2r5/8/8/8/1p5k/6Q1/K7&moves=5}{\color{white}\bfseries lösen\,\ensuremath{\nearrow}}}
\end{tcolorbox}
\vspace{2mm}
\noindent\begin{minipage}[t]{36mm}\vspace{0pt}\sffamily\small\raggedright
\hmfigure{längste eindeutige}{{\Large h\#5}}\hmfigure{eindeutig bei dieser Länge}{{\Large 631} Stellungen}\hmfigure{längstes Matt}{{\Large h\#6.5}}\hmfigure{in der Klasse}{710 Millionen Stellungen,\\ 441 Millionen mit Hilfsmatt}
\end{minipage}\hfill
\begin{minipage}[t]{80mm}\vspace{0pt}\begin{tikzpicture}\begin{axis}[width=72mm,height=42mm,ybar,bar width=1.6pt,xmin=-1,xmax=14,ymin=0,ymax=121.9,axis y line*=left,axis x line*=bottom,xlabel={\scriptsize Halbzüge bis Matt},ylabel={\scriptsize Stellungen (Mio.)},tick label style={font=\tiny},label style={font=\scriptsize},enlarge x limits=false,ylabel shift=-3pt,xlabel shift=-3pt]\addplot[fill=ink!55,draw=none] coordinates {(0,0.2606) (1,2.063) (2,7.577) (3,26.44) (4,72.55) (5,88.2) (6,110.8) (7,73.37) (8,42.14) (9,14.97) (10,2.039) (11,0.104) (12,0.005391) (13,7e-06)};\draw[densely dotted] (axis cs:10,0) -- (axis cs:10,121.9);\end{axis}\begin{axis}[width=72mm,height=42mm,axis y line*=right,axis x line=none,ymode=log,xmin=-1,xmax=14,ylabel={\scriptsize eindeutig (log)},tick label style={font=\tiny},label style={font=\scriptsize},ylabel shift=-3pt]\addplot[mark=*,mark size=0.7pt,draw=black,thin] coordinates {(0,260601) (1,1612666) (2,4155028) (3,4657695) (4,5378548) (5,2696856) (6,1297462) (7,519624) (8,122966) (9,10364) (10,631)};\end{axis}\end{tikzpicture}\end{minipage}
\vspace{3mm}

\noindent\begin{minipage}[t]{60mm}\vspace{0pt}\sbox\hmboard{\chessboard[setfen=8/2r5/8/8/8/1p5k/6Q1/K7 b - - 0 1]}\usebox\hmboard\\[1mm]
\makebox[\wd\hmboard]{\sffamily\hspace{3mm}\textbf{h\#5}\hfill(2+3)}\end{minipage}\hfill
\begin{minipage}[t]{52mm}\vspace{0pt}
\textcolor{ink}{\sffamily\textsc{Lösung}}\par\smallskip\small
\begin{tabular}{@{}r@{\enspace}l@{\enspace}l@{}}
\textbf{1.} & \symking{}h4 & \symqueen{}c6\\
\textbf{2.} & \symking{}g3 & \symking{}b2\\
\textbf{3.} & \symking{}f2 & \symking{}c3\\
\textbf{4.} & \symking{}e1 & \symking{}d3\\
\textbf{5.} & \symking{}d1 & \symqueen{}h1\#\\
\end{tabular}\par\vspace{3mm}
\qrcode[height=12mm]{https://helpman.komtera.lt/?fen=8/2r5/8/8/8/1p5k/6Q1/K7&moves=5}\\{\sffamily\tiny\color{ink}scannen zum Lösen}
\end{minipage}
\vspace{2mm}

{\raggedright\noindent\textcolor{ink}{\sffamily\textsc{Themen}}\enspace \themetag{reines Matt} \themetag{Mustermatt} \themetag{Spiegelmatt} \themetag{ein Stein zieht} \themetag{ohne Schlag} \themetag{ohne Schach}\par\smallskip{\footnotesize\textit{Hinweis:} die Partei am Zug steht im Diagramm im Schach.}\par}

\clearpage
\phantomsection\label{hm:KRBvkr}
\begin{tcolorbox}[enhanced,colback=ink,colframe=ink,arc=0pt,boxrule=0pt,left=4mm,right=4mm,top=2.5mm,bottom=2.5mm]
\color{white}\sffamily{\LARGE\bfseries KRBvkr}\hfill{\Large h\#5}\\[0.5mm]
{\small 5 Steine\quad Schwarz am Zug}\hfill{\small Nr.~210\qquad\href{https://helpman.komtera.lt/?fen=6rk/6B1/8/8/8/8/7R/1K6&moves=5}{\color{white}\bfseries lösen\,\ensuremath{\nearrow}}}
\end{tcolorbox}
\vspace{2mm}
\noindent\begin{minipage}[t]{36mm}\vspace{0pt}\sffamily\small\raggedright
\hmfigure{längste eindeutige}{{\Large h\#5}}\hmfigure{eindeutig bei dieser Länge}{{\Large 3} Stellungen}\hmfigure{längstes Matt}{{\Large h\#7}}\hmfigure{in der Klasse}{242 Millionen Stellungen,\\ 153 Millionen mit Hilfsmatt}
\end{minipage}\hfill
\begin{minipage}[t]{80mm}\vspace{0pt}\begin{tikzpicture}\begin{axis}[width=72mm,height=42mm,ybar,bar width=1.6pt,xmin=-1,xmax=15,ymin=0,ymax=58.59,axis y line*=left,axis x line*=bottom,xlabel={\scriptsize Halbzüge bis Matt},ylabel={\scriptsize Stellungen (Mio.)},tick label style={font=\tiny},label style={font=\scriptsize},enlarge x limits=false,ylabel shift=-3pt,xlabel shift=-3pt]\addplot[fill=ink!55,draw=none] coordinates {(0,0.06159) (1,0.3329) (2,1.273) (3,5.965) (4,20.94) (5,40.21) (6,53.27) (7,23.8) (8,5.911) (9,0.7444) (10,0.004248) (11,0.001678) (12,0.000533) (13,0.000156) (14,1e-05)};\draw[densely dotted] (axis cs:10,0) -- (axis cs:10,58.59);\end{axis}\begin{axis}[width=72mm,height=42mm,axis y line*=right,axis x line=none,ymode=log,xmin=-1,xmax=15,ylabel={\scriptsize eindeutig (log)},tick label style={font=\tiny},label style={font=\scriptsize},ylabel shift=-3pt]\addplot[mark=*,mark size=0.7pt,draw=black,thin] coordinates {(0,61592) (1,327837) (2,741415) (3,2567071) (4,7172590) (5,3544777) (6,905078) (7,77683) (8,438) (9,19) (10,3)};\end{axis}\end{tikzpicture}\end{minipage}
\vspace{3mm}

\noindent\begin{minipage}[t]{60mm}\vspace{0pt}\sbox\hmboard{\chessboard[setfen=6rk/6B1/8/8/8/8/7R/1K6 b - - 0 1]}\usebox\hmboard\\[1mm]
\makebox[\wd\hmboard]{\sffamily\hspace{3mm}\textbf{h\#5}\hfill(3+2)}\end{minipage}\hfill
\begin{minipage}[t]{52mm}\vspace{0pt}
\textcolor{ink}{\sffamily\textsc{Lösung}}\par\smallskip\small
\begin{tabular}{@{}r@{\enspace}l@{\enspace}l@{}}
\textbf{1.} & \symking{}xg7 & \symking{}c2\\
\textbf{2.} & \symrook{}h8 & \symking{}d3\\
\textbf{3.} & \symrook{}h3+ & \symking{}e4\\
\textbf{4.} & \symking{}h6 & \symking{}f5\\
\textbf{5.} & \symking{}h5 & \symrook{}xh3\#\\
\end{tabular}\par\vspace{3mm}
\qrcode[height=12mm]{https://helpman.komtera.lt/?fen=6rk/6B1/8/8/8/8/7R/1K6&moves=5}\\{\sffamily\tiny\color{ink}scannen zum Lösen}
\end{minipage}
\vspace{2mm}

{\raggedright\noindent\textcolor{ink}{\sffamily\textsc{Themen}}\enspace \themetag{reines Matt} \themetag{Mustermatt} \themetag{Idealmatt} \themetag{Spiegelmatt}\par\smallskip{\footnotesize\textit{Hinweis:} die Partei am Zug steht im Diagramm im Schach; die Lösung beginnt mit einem Schlag.}\par}

\clearpage
\phantomsection\label{hm:KRRNvk}
\begin{tcolorbox}[enhanced,colback=ink,colframe=ink,arc=0pt,boxrule=0pt,left=4mm,right=4mm,top=2.5mm,bottom=2.5mm]
\color{white}\sffamily{\LARGE\bfseries KRRNvk}\hfill{\Large h\#5}\\[0.5mm]
{\small 5 Steine\quad Schwarz am Zug}\hfill{\small Nr.~211\qquad\href{https://helpman.komtera.lt/?fen=8/8/8/6R1/7k/N4R2/8/K7&moves=5}{\color{white}\bfseries lösen\,\ensuremath{\nearrow}}}
\end{tcolorbox}
\vspace{2mm}
\noindent\begin{minipage}[t]{36mm}\vspace{0pt}\sffamily\small\raggedright
\hmfigure{längste eindeutige}{{\Large h\#5}}\hmfigure{eindeutig bei dieser Länge}{{\Large 4} Stellungen}\hmfigure{längstes Matt}{{\Large h\#7}}\hmfigure{in der Klasse}{242 Millionen Stellungen,\\ 80,2 Millionen mit Hilfsmatt}
\end{minipage}\hfill
\begin{minipage}[t]{80mm}\vspace{0pt}\begin{tikzpicture}\begin{axis}[width=72mm,height=42mm,ybar,bar width=1.6pt,xmin=-1,xmax=15,ymin=0,ymax=38.29,axis y line*=left,axis x line*=bottom,xlabel={\scriptsize Halbzüge bis Matt},ylabel={\scriptsize Stellungen (Mio.)},tick label style={font=\tiny},label style={font=\scriptsize},enlarge x limits=false,ylabel shift=-3pt,xlabel shift=-3pt]\addplot[fill=ink!55,draw=none] coordinates {(0,0.555) (1,2.249) (2,6.255) (3,16.29) (4,34.81) (5,10.49) (6,8.848) (7,0.2365) (8,0.3462) (9,8e-06) (10,0.08825) (12,0.003725) (14,8e-06)};\draw[densely dotted] (axis cs:10,0) -- (axis cs:10,38.29);\end{axis}\begin{axis}[width=72mm,height=42mm,axis y line*=right,axis x line=none,ymode=log,xmin=-1,xmax=15,ylabel={\scriptsize eindeutig (log)},tick label style={font=\tiny},label style={font=\scriptsize},ylabel shift=-3pt]\addplot[mark=*,mark size=0.7pt,draw=black,thin] coordinates {(0,554955) (1,1901501) (2,3889156) (3,2334534) (4,2464124) (5,42610) (6,43362) (7,1) (8,2513) (10,4)};\end{axis}\end{tikzpicture}\end{minipage}
\vspace{3mm}

\noindent\begin{minipage}[t]{60mm}\vspace{0pt}\sbox\hmboard{\chessboard[setfen=8/8/8/6R1/7k/N4R2/8/K7 b - - 0 1]}\usebox\hmboard\\[1mm]
\makebox[\wd\hmboard]{\sffamily\hspace{3mm}\textbf{h\#5}\hfill(4+1)}\end{minipage}\hfill
\begin{minipage}[t]{52mm}\vspace{0pt}
\textcolor{ink}{\sffamily\textsc{Lösung}}\par\smallskip\small
\begin{tabular}{@{}r@{\enspace}l@{\enspace}l@{}}
\textbf{1.} & \symking{}xg5 & \symrook{}b3\\
\textbf{2.} & \symking{}f4 & \symrook{}b2\\
\textbf{3.} & \symking{}e3 & \symknight{}c2+\\
\textbf{4.} & \symking{}d2 & \symknight{}e3+\\
\textbf{5.} & \symking{}c1 & \symrook{}c2\#\\
\end{tabular}\par\vspace{3mm}
\qrcode[height=12mm]{https://helpman.komtera.lt/?fen=8/8/8/6R1/7k/N4R2/8/K7&moves=5}\\{\sffamily\tiny\color{ink}scannen zum Lösen}
\end{minipage}
\vspace{2mm}

{\raggedright\noindent\textcolor{ink}{\sffamily\textsc{Themen}}\enspace \themetag{ein Stein zieht} \themetag{Umnov}\par\smallskip{\footnotesize\textit{Hinweis:} die Lösung beginnt mit einem Schlag.}\par}

\clearpage
\phantomsection\label{hm:KRRvkb}
\begin{tcolorbox}[enhanced,colback=ink,colframe=ink,arc=0pt,boxrule=0pt,left=4mm,right=4mm,top=2.5mm,bottom=2.5mm]
\color{white}\sffamily{\LARGE\bfseries KRRvkb}\hfill{\Large h\#5}\\[0.5mm]
{\small 5 Steine\quad Schwarz am Zug}\hfill{\small Nr.~212\qquad\href{https://helpman.komtera.lt/?fen=7b/6R1/8/8/8/8/6R1/K6k&moves=5}{\color{white}\bfseries lösen\,\ensuremath{\nearrow}}}
\end{tcolorbox}
\vspace{2mm}
\noindent\begin{minipage}[t]{36mm}\vspace{0pt}\sffamily\small\raggedright
\hmfigure{längste eindeutige}{{\Large h\#5}}\hmfigure{eindeutig bei dieser Länge}{{\Large 290} Stellungen}\hmfigure{längstes Matt}{{\Large h\#7}}\hmfigure{in der Klasse}{242 Millionen Stellungen,\\ 76,3 Millionen mit Hilfsmatt}
\end{minipage}\hfill
\begin{minipage}[t]{80mm}\vspace{0pt}\begin{tikzpicture}\begin{axis}[width=72mm,height=42mm,ybar,bar width=1.6pt,xmin=-1,xmax=15,ymin=0,ymax=33.55,axis y line*=left,axis x line*=bottom,xlabel={\scriptsize Halbzüge bis Matt},ylabel={\scriptsize Stellungen (Mio.)},tick label style={font=\tiny},label style={font=\scriptsize},enlarge x limits=false,ylabel shift=-3pt,xlabel shift=-3pt]\addplot[fill=ink!55,draw=none] coordinates {(0,0.2904) (1,1.306) (2,4.792) (3,16.27) (4,30.5) (5,13.69) (6,8.648) (7,0.5546) (8,0.1868) (9,0.000259) (10,0.08734) (11,2e-06) (12,0.01926) (13,1e-06) (14,0.000231)};\draw[densely dotted] (axis cs:10,0) -- (axis cs:10,33.55);\end{axis}\begin{axis}[width=72mm,height=42mm,axis y line*=right,axis x line=none,ymode=log,xmin=-1,xmax=15,ylabel={\scriptsize eindeutig (log)},tick label style={font=\tiny},label style={font=\scriptsize},ylabel shift=-3pt]\addplot[mark=*,mark size=0.7pt,draw=black,thin] coordinates {(0,290378) (1,1173688) (2,2032365) (3,2694293) (4,1428496) (5,68214) (6,14364) (7,22) (8,1801) (10,290)};\end{axis}\end{tikzpicture}\end{minipage}
\vspace{3mm}

\noindent\begin{minipage}[t]{60mm}\vspace{0pt}\sbox\hmboard{\chessboard[setfen=7b/6R1/8/8/8/8/6R1/K6k b - - 0 1]}\usebox\hmboard\\[1mm]
\makebox[\wd\hmboard]{\sffamily\hspace{3mm}\textbf{h\#5}\hfill(3+2)}\end{minipage}\hfill
\begin{minipage}[t]{52mm}\vspace{0pt}
\textcolor{ink}{\sffamily\textsc{Lösung}}\par\smallskip\small
\begin{tabular}{@{}r@{\enspace}l@{\enspace}l@{}}
\textbf{1.} & \symbishop{}xg7+ & \symrook{}b2\\
\textbf{2.} & \symking{}g1 & \symking{}b1\\
\textbf{3.} & \symking{}f1 & \symking{}c2\\
\textbf{4.} & \symking{}e1 & \symking{}d3\\
\textbf{5.} & \symking{}d1 & \symrook{}b1\#\\
\end{tabular}\par\vspace{3mm}
\qrcode[height=12mm]{https://helpman.komtera.lt/?fen=7b/6R1/8/8/8/8/6R1/K6k&moves=5}\\{\sffamily\tiny\color{ink}scannen zum Lösen}
\end{minipage}
\vspace{2mm}

{\raggedright\noindent\textcolor{ink}{\sffamily\textsc{Themen}}\enspace \themetag{reines Matt} \themetag{Mustermatt} \themetag{Spiegelmatt}\par\smallskip{\footnotesize\textit{Hinweis:} die Lösung beginnt mit einem Schlag.}\par}

\clearpage
\phantomsection\label{hm:KRRvkn}
\begin{tcolorbox}[enhanced,colback=ink,colframe=ink,arc=0pt,boxrule=0pt,left=4mm,right=4mm,top=2.5mm,bottom=2.5mm]
\color{white}\sffamily{\LARGE\bfseries KRRvkn}\hfill{\Large h\#5}\\[0.5mm]
{\small 5 Steine\quad Schwarz am Zug}\hfill{\small Nr.~213\qquad\href{https://helpman.komtera.lt/?fen=7k/6R1/8/8/7n/8/8/1K5R&moves=5}{\color{white}\bfseries lösen\,\ensuremath{\nearrow}}}
\end{tcolorbox}
\vspace{2mm}
\noindent\begin{minipage}[t]{36mm}\vspace{0pt}\sffamily\small\raggedright
\hmfigure{längste eindeutige}{{\Large h\#5}}\hmfigure{eindeutig bei dieser Länge}{{\Large 229} Stellungen}\hmfigure{längstes Matt}{{\Large h\#7}}\hmfigure{in der Klasse}{242 Millionen Stellungen,\\ 78,8 Millionen mit Hilfsmatt}
\end{minipage}\hfill
\begin{minipage}[t]{80mm}\vspace{0pt}\begin{tikzpicture}\begin{axis}[width=72mm,height=42mm,ybar,bar width=1.6pt,xmin=-1,xmax=15,ymin=0,ymax=35.42,axis y line*=left,axis x line*=bottom,xlabel={\scriptsize Halbzüge bis Matt},ylabel={\scriptsize Stellungen (Mio.)},tick label style={font=\tiny},label style={font=\scriptsize},enlarge x limits=false,ylabel shift=-3pt,xlabel shift=-3pt]\addplot[fill=ink!55,draw=none] coordinates {(0,0.3836) (1,1.715) (2,5.623) (3,17.11) (4,32.2) (5,12.5) (6,8.493) (7,0.4931) (8,0.1981) (9,0.000254) (10,0.08858) (12,0.02022) (14,0.000361)};\draw[densely dotted] (axis cs:10,0) -- (axis cs:10,35.42);\end{axis}\begin{axis}[width=72mm,height=42mm,axis y line*=right,axis x line=none,ymode=log,xmin=-1,xmax=15,ylabel={\scriptsize eindeutig (log)},tick label style={font=\tiny},label style={font=\scriptsize},ylabel shift=-3pt]\addplot[mark=*,mark size=0.7pt,draw=black,thin] coordinates {(0,383640) (1,1501919) (2,2376669) (3,2203435) (4,1419106) (5,44755) (6,13557) (7,51) (8,1841) (10,229)};\end{axis}\end{tikzpicture}\end{minipage}
\vspace{3mm}

\noindent\begin{minipage}[t]{60mm}\vspace{0pt}\sbox\hmboard{\chessboard[setfen=7k/6R1/8/8/7n/8/8/1K5R b - - 0 1]}\usebox\hmboard\\[1mm]
\makebox[\wd\hmboard]{\sffamily\hspace{3mm}\textbf{h\#5}\hfill(3+2)}\end{minipage}\hfill
\begin{minipage}[t]{52mm}\vspace{0pt}
\textcolor{ink}{\sffamily\textsc{Lösung}}\par\smallskip\small
\begin{tabular}{@{}r@{\enspace}l@{\enspace}l@{}}
\textbf{1.} & \symking{}xg7 & \symking{}c2\\
\textbf{2.} & \symknight{}f3 & \symking{}d3\\
\textbf{3.} & \symknight{}h2 & \symking{}e4\\
\textbf{4.} & \symking{}h6 & \symking{}f5\\
\textbf{5.} & \symking{}h5 & \symrook{}xh2\#\\
\end{tabular}\par\vspace{3mm}
\qrcode[height=12mm]{https://helpman.komtera.lt/?fen=7k/6R1/8/8/7n/8/8/1K5R&moves=5}\\{\sffamily\tiny\color{ink}scannen zum Lösen}
\end{minipage}
\vspace{2mm}

{\raggedright\noindent\textcolor{ink}{\sffamily\textsc{Themen}}\enspace \themetag{reines Matt} \themetag{Mustermatt} \themetag{Idealmatt} \themetag{Spiegelmatt} \themetag{ohne Schach}\par\smallskip{\footnotesize\textit{Hinweis:} die Lösung beginnt mit einem Schlag.}\par}

\clearpage
\phantomsection\label{hm:KRRvkq}
\begin{tcolorbox}[enhanced,colback=ink,colframe=ink,arc=0pt,boxrule=0pt,left=4mm,right=4mm,top=2.5mm,bottom=2.5mm]
\color{white}\sffamily{\LARGE\bfseries KRRvkq}\hfill{\Large h\#5}\\[0.5mm]
{\small 5 Steine\quad Schwarz am Zug}\hfill{\small Nr.~214\qquad\href{https://helpman.komtera.lt/?fen=8/8/8/8/8/8/5R2/K4Rqk&moves=5}{\color{white}\bfseries lösen\,\ensuremath{\nearrow}}}
\end{tcolorbox}
\vspace{2mm}
\noindent\begin{minipage}[t]{36mm}\vspace{0pt}\sffamily\small\raggedright
\hmfigure{längste eindeutige}{{\Large h\#5}}\hmfigure{eindeutig bei dieser Länge}{{\Large 219} Stellungen}\hmfigure{längstes Matt}{{\Large h\#7}}\hmfigure{in der Klasse}{242 Millionen Stellungen,\\ 65,9 Millionen mit Hilfsmatt}
\end{minipage}\hfill
\begin{minipage}[t]{80mm}\vspace{0pt}\begin{tikzpicture}\begin{axis}[width=72mm,height=42mm,ybar,bar width=1.6pt,xmin=-1,xmax=15,ymin=0,ymax=25.34,axis y line*=left,axis x line*=bottom,xlabel={\scriptsize Halbzüge bis Matt},ylabel={\scriptsize Stellungen (Mio.)},tick label style={font=\tiny},label style={font=\scriptsize},enlarge x limits=false,ylabel shift=-3pt,xlabel shift=-3pt]\addplot[fill=ink!55,draw=none] coordinates {(0,0.1344) (1,0.6329) (2,3.096) (3,12.99) (4,23.04) (5,17.11) (6,7.562) (7,1.051) (8,0.1479) (9,0.001775) (10,0.0706) (11,0.000493) (12,0.01904) (13,9.8e-05) (14,0.000403)};\draw[densely dotted] (axis cs:10,0) -- (axis cs:10,25.34);\end{axis}\begin{axis}[width=72mm,height=42mm,axis y line*=right,axis x line=none,ymode=log,xmin=-1,xmax=15,ylabel={\scriptsize eindeutig (log)},tick label style={font=\tiny},label style={font=\scriptsize},ylabel shift=-3pt]\addplot[mark=*,mark size=0.7pt,draw=black,thin] coordinates {(0,134401) (1,599853) (2,1344506) (3,2746569) (4,1097392) (5,127888) (6,11239) (7,79) (8,1161) (10,219)};\end{axis}\end{tikzpicture}\end{minipage}
\vspace{3mm}

\noindent\begin{minipage}[t]{60mm}\vspace{0pt}\sbox\hmboard{\chessboard[setfen=8/8/8/8/8/8/5R2/K4Rqk b - - 0 1]}\usebox\hmboard\\[1mm]
\makebox[\wd\hmboard]{\sffamily\hspace{3mm}\textbf{h\#5}\hfill(3+2)}\end{minipage}\hfill
\begin{minipage}[t]{52mm}\vspace{0pt}
\textcolor{ink}{\sffamily\textsc{Lösung}}\par\smallskip\small
\begin{tabular}{@{}r@{\enspace}l@{\enspace}l@{}}
\textbf{1.} & \symqueen{}xf1+ & \symking{}b2\\
\textbf{2.} & \symking{}g1 & \symrook{}c2\\
\textbf{3.} & \symqueen{}f2 & \symking{}c3\\
\textbf{4.} & \symking{}f1 & \symking{}d3\\
\textbf{5.} & \symking{}e1 & \symrook{}c1\#\\
\end{tabular}\par\vspace{3mm}
\qrcode[height=12mm]{https://helpman.komtera.lt/?fen=8/8/8/8/8/8/5R2/K4Rqk&moves=5}\\{\sffamily\tiny\color{ink}scannen zum Lösen}
\end{minipage}
\vspace{2mm}

{\raggedright\noindent\textcolor{ink}{\sffamily\textsc{Themen}}\enspace \themetag{reines Matt} \themetag{Mustermatt} \themetag{Idealmatt} \themetag{Selbstblock} \themetag{Umnov}\par\smallskip{\footnotesize\textit{Hinweis:} die Lösung beginnt mit einem Schlag.}\par}

\clearpage
\phantomsection\label{hm:KQvkrr}
\begin{tcolorbox}[enhanced,colback=ink,colframe=ink,arc=0pt,boxrule=0pt,left=4mm,right=4mm,top=2.5mm,bottom=2.5mm]
\color{white}\sffamily{\LARGE\bfseries KQvkrr}\hfill{\Large h\#4.5}\\[0.5mm]
{\small 5 Steine\quad Weiß am Zug}\hfill{\small Nr.~215\qquad\href{https://helpman.komtera.lt/?fen=8/8/8/6k1/8/8/1r2r3/K4Q2&moves=4.5}{\color{white}\bfseries lösen\,\ensuremath{\nearrow}}}
\end{tcolorbox}
\vspace{2mm}
\noindent\begin{minipage}[t]{36mm}\vspace{0pt}\sffamily\small\raggedright
\hmfigure{längste eindeutige}{{\Large h\#4.5}}\hmfigure{eindeutig bei dieser Länge}{{\Large 460} Stellungen}\hmfigure{längstes Matt}{{\Large h\#6}}\hmfigure{in der Klasse}{242 Millionen Stellungen,\\ 65,4 Millionen mit Hilfsmatt}
\end{minipage}\hfill
\begin{minipage}[t]{80mm}\vspace{0pt}\begin{tikzpicture}\begin{axis}[width=72mm,height=42mm,ybar,bar width=1.6pt,xmin=-1,xmax=13,ymin=0,ymax=17.41,axis y line*=left,axis x line*=bottom,xlabel={\scriptsize Halbzüge bis Matt},ylabel={\scriptsize Stellungen (Mio.)},tick label style={font=\tiny},label style={font=\scriptsize},enlarge x limits=false,ylabel shift=-3pt,xlabel shift=-3pt]\addplot[fill=ink!55,draw=none] coordinates {(0,0.0236) (1,0.2142) (2,1.024) (3,4.137) (4,11.22) (5,15.82) (6,15.51) (7,12.29) (8,3.574) (9,1.464) (10,0.074) (11,0.006066) (12,1.1e-05)};\draw[densely dotted] (axis cs:9,0) -- (axis cs:9,17.41);\end{axis}\begin{axis}[width=72mm,height=42mm,axis y line*=right,axis x line=none,ymode=log,xmin=-1,xmax=13,ylabel={\scriptsize eindeutig (log)},tick label style={font=\tiny},label style={font=\scriptsize},ylabel shift=-3pt]\addplot[mark=*,mark size=0.7pt,draw=black,thin] coordinates {(0,23595) (1,181350) (2,586616) (3,745816) (4,655970) (5,391666) (6,112467) (7,45392) (8,2947) (9,460)};\end{axis}\end{tikzpicture}\end{minipage}
\vspace{3mm}

\noindent\begin{minipage}[t]{60mm}\vspace{0pt}\sbox\hmboard{\chessboard[setfen=8/8/8/6k1/8/8/1r2r3/K4Q2 w - - 0 1]}\usebox\hmboard\\[1mm]
\makebox[\wd\hmboard]{\sffamily\hspace{3mm}\textbf{h\#4.5}\hfill(2+3)}\end{minipage}\hfill
\begin{minipage}[t]{52mm}\vspace{0pt}
\textcolor{ink}{\sffamily\textsc{Lösung}}\par\smallskip\small
\begin{tabular}{@{}r@{\enspace}l@{\enspace}l@{}}
\textbf{1.} & \ldots & \symqueen{}b1\\
\textbf{2.} & \symking{}f4 & \symqueen{}xb2\\
\textbf{3.} & \symking{}e3 & \symqueen{}a3+\\
\textbf{4.} & \symking{}d2 & \symking{}b1\\
\textbf{5.} & \symking{}d1 & \symqueen{}c1\#\\
\end{tabular}\par\vspace{3mm}
\qrcode[height=12mm]{https://helpman.komtera.lt/?fen=8/8/8/6k1/8/8/1r2r3/K4Q2&moves=4.5}\\{\sffamily\tiny\color{ink}scannen zum Lösen}
\end{minipage}
\vspace{2mm}

{\raggedright\noindent\textcolor{ink}{\sffamily\textsc{Themen}}\enspace \themetag{ein Stein zieht}\par}

\clearpage
\phantomsection\label{hm:KQQNvk}
\begin{tcolorbox}[enhanced,colback=ink,colframe=ink,arc=0pt,boxrule=0pt,left=4mm,right=4mm,top=2.5mm,bottom=2.5mm]
\color{white}\sffamily{\LARGE\bfseries KQQNvk}\hfill{\Large h\#4}\\[0.5mm]
{\small 5 Steine\quad Schwarz am Zug}\hfill{\small Nr.~216\qquad\href{https://helpman.komtera.lt/?fen=N7/8/8/8/8/1Q6/7Q/K6k&moves=4}{\color{white}\bfseries lösen\,\ensuremath{\nearrow}}}
\end{tcolorbox}
\vspace{2mm}
\noindent\begin{minipage}[t]{36mm}\vspace{0pt}\sffamily\small\raggedright
\hmfigure{längste eindeutige}{{\Large h\#4}}\hmfigure{eindeutig bei dieser Länge}{{\Large 18} Stellungen}\hmfigure{längstes Matt}{{\Large h\#5}}\hmfigure{in der Klasse}{242 Millionen Stellungen,\\ 70,3 Millionen mit Hilfsmatt}
\end{minipage}\hfill
\begin{minipage}[t]{80mm}\vspace{0pt}\begin{tikzpicture}\begin{axis}[width=72mm,height=42mm,ybar,bar width=1.6pt,xmin=-1,xmax=11,ymin=0,ymax=26.08,axis y line*=left,axis x line*=bottom,xlabel={\scriptsize Halbzüge bis Matt},ylabel={\scriptsize Stellungen (Mio.)},tick label style={font=\tiny},label style={font=\scriptsize},enlarge x limits=false,ylabel shift=-3pt,xlabel shift=-3pt]\addplot[fill=ink!55,draw=none] coordinates {(0,2.197) (1,6.202) (2,22.85) (3,13.94) (4,23.71) (5,0.1226) (6,1.117) (8,0.1352) (10,0.000186)};\draw[densely dotted] (axis cs:8,0) -- (axis cs:8,26.08);\end{axis}\begin{axis}[width=72mm,height=42mm,axis y line*=right,axis x line=none,ymode=log,xmin=-1,xmax=11,ylabel={\scriptsize eindeutig (log)},tick label style={font=\tiny},label style={font=\scriptsize},ylabel shift=-3pt]\addplot[mark=*,mark size=0.7pt,draw=black,thin] coordinates {(0,2197249) (1,2985478) (2,8223461) (3,124021) (4,171895) (6,15936) (8,18)};\end{axis}\end{tikzpicture}\end{minipage}
\vspace{3mm}

\noindent\begin{minipage}[t]{60mm}\vspace{0pt}\sbox\hmboard{\chessboard[setfen=N7/8/8/8/8/1Q6/7Q/K6k b - - 0 1]}\usebox\hmboard\\[1mm]
\makebox[\wd\hmboard]{\sffamily\hspace{3mm}\textbf{h\#4}\hfill(4+1)}\end{minipage}\hfill
\begin{minipage}[t]{52mm}\vspace{0pt}
\textcolor{ink}{\sffamily\textsc{Lösung}}\par\smallskip\small
\begin{tabular}{@{}r@{\enspace}l@{\enspace}l@{}}
\textbf{1.} & \symking{}xh2 & \symqueen{}g8\\
\textbf{2.} & \symking{}h3 & \symknight{}c7\\
\textbf{3.} & \symking{}h4 & \symknight{}e6\\
\textbf{4.} & \symking{}h5 & \symqueen{}g5\#\\
\end{tabular}\par\vspace{3mm}
\qrcode[height=12mm]{https://helpman.komtera.lt/?fen=N7/8/8/8/8/1Q6/7Q/K6k&moves=4}\\{\sffamily\tiny\color{ink}scannen zum Lösen}
\end{minipage}
\vspace{2mm}

{\raggedright\noindent\textcolor{ink}{\sffamily\textsc{Themen}}\enspace \themetag{reines Matt} \themetag{Mustermatt} \themetag{ein Stein zieht} \themetag{ohne Schach}\par\smallskip{\footnotesize\textit{Hinweis:} die Partei am Zug steht im Diagramm im Schach; die Lösung beginnt mit einem Schlag.}\par}

\clearpage
\phantomsection\label{hm:KQQBvk}
\begin{tcolorbox}[enhanced,colback=ink,colframe=ink,arc=0pt,boxrule=0pt,left=4mm,right=4mm,top=2.5mm,bottom=2.5mm]
\color{white}\sffamily{\LARGE\bfseries KQQBvk}\hfill{\Large h\#3}\\[0.5mm]
{\small 5 Steine\quad Schwarz am Zug}\hfill{\small Nr.~217\qquad\href{https://helpman.komtera.lt/?fen=8/8/8/8/4Q3/B2k4/1Q6/K7&moves=3}{\color{white}\bfseries lösen\,\ensuremath{\nearrow}}}
\end{tcolorbox}
\vspace{2mm}
\noindent\begin{minipage}[t]{36mm}\vspace{0pt}\sffamily\small\raggedright
\hmfigure{längste eindeutige}{{\Large h\#3}}\hmfigure{eindeutig bei dieser Länge}{{\Large 293} Stellungen}\hmfigure{längstes Matt}{{\Large h\#4}}\hmfigure{in der Klasse}{242 Millionen Stellungen,\\ 68,6 Millionen mit Hilfsmatt}
\end{minipage}\hfill
\begin{minipage}[t]{80mm}\vspace{0pt}\begin{tikzpicture}\begin{axis}[width=72mm,height=42mm,ybar,bar width=1.6pt,xmin=-1,xmax=9,ymin=0,ymax=26.71,axis y line*=left,axis x line*=bottom,xlabel={\scriptsize Halbzüge bis Matt},ylabel={\scriptsize Stellungen (Mio.)},tick label style={font=\tiny},label style={font=\scriptsize},enlarge x limits=false,ylabel shift=-3pt,xlabel shift=-3pt]\addplot[fill=ink!55,draw=none] coordinates {(0,2.728) (1,6.464) (2,24.29) (3,12.23) (4,22.13) (5,0.0807) (6,0.7177) (8,0.001253)};\draw[densely dotted] (axis cs:6,0) -- (axis cs:6,26.71);\end{axis}\begin{axis}[width=72mm,height=42mm,axis y line*=right,axis x line=none,ymode=log,xmin=-1,xmax=9,ylabel={\scriptsize eindeutig (log)},tick label style={font=\tiny},label style={font=\scriptsize},ylabel shift=-3pt]\addplot[mark=*,mark size=0.7pt,draw=black,thin] coordinates {(0,2728303) (1,2808189) (2,7901392) (3,84995) (4,394301) (6,293)};\end{axis}\end{tikzpicture}\end{minipage}
\vspace{3mm}

\noindent\begin{minipage}[t]{60mm}\vspace{0pt}\sbox\hmboard{\chessboard[setfen=8/8/8/8/4Q3/B2k4/1Q6/K7 b - - 0 1]}\usebox\hmboard\\[1mm]
\makebox[\wd\hmboard]{\sffamily\hspace{3mm}\textbf{h\#3}\hfill(4+1)}\end{minipage}\hfill
\begin{minipage}[t]{52mm}\vspace{0pt}
\textcolor{ink}{\sffamily\textsc{Lösung}}\par\smallskip\small
\begin{tabular}{@{}r@{\enspace}l@{\enspace}l@{}}
\textbf{1.} & \symking{}xe4 & \symking{}b1\\
\textbf{2.} & \symking{}e3 & \symbishop{}c5+\\
\textbf{3.} & \symking{}d3 & \symqueen{}c2\#\\
\end{tabular}\par\vspace{3mm}
\qrcode[height=12mm]{https://helpman.komtera.lt/?fen=8/8/8/8/4Q3/B2k4/1Q6/K7&moves=3}\\{\sffamily\tiny\color{ink}scannen zum Lösen}
\end{minipage}
\vspace{2mm}

{\raggedright\noindent\textcolor{ink}{\sffamily\textsc{Themen}}\enspace \themetag{reines Matt} \themetag{Mustermatt} \themetag{Idealmatt} \themetag{Rundlauf} \themetag{ein Stein zieht}\par\smallskip{\footnotesize\textit{Hinweis:} die Partei am Zug steht im Diagramm im Schach; die Lösung beginnt mit einem Schlag.}\par}

\clearpage
\phantomsection\label{hm:KQQRvk}
\begin{tcolorbox}[enhanced,colback=ink,colframe=ink,arc=0pt,boxrule=0pt,left=4mm,right=4mm,top=2.5mm,bottom=2.5mm]
\color{white}\sffamily{\LARGE\bfseries KQQRvk}\hfill{\Large h\#3}\\[0.5mm]
{\small 5 Steine\quad Schwarz am Zug}\hfill{\small Nr.~218\qquad\href{https://helpman.komtera.lt/?fen=8/8/8/8/3Q4/3k4/2R5/K1Q5&moves=3}{\color{white}\bfseries lösen\,\ensuremath{\nearrow}}}
\end{tcolorbox}
\vspace{2mm}
\noindent\begin{minipage}[t]{36mm}\vspace{0pt}\sffamily\small\raggedright
\hmfigure{längste eindeutige}{{\Large h\#3}}\hmfigure{eindeutig bei dieser Länge}{{\Large 775} Stellungen}\hmfigure{längstes Matt}{{\Large h\#4}}\hmfigure{in der Klasse}{242 Millionen Stellungen,\\ 66,3 Millionen mit Hilfsmatt}
\end{minipage}\hfill
\begin{minipage}[t]{80mm}\vspace{0pt}\begin{tikzpicture}\begin{axis}[width=72mm,height=42mm,ybar,bar width=1.6pt,xmin=-1,xmax=9,ymin=0,ymax=28.84,axis y line*=left,axis x line*=bottom,xlabel={\scriptsize Halbzüge bis Matt},ylabel={\scriptsize Stellungen (Mio.)},tick label style={font=\tiny},label style={font=\scriptsize},enlarge x limits=false,ylabel shift=-3pt,xlabel shift=-3pt]\addplot[fill=ink!55,draw=none] coordinates {(0,4.315) (1,7.224) (2,26.22) (3,9.454) (4,18.82) (5,0.006955) (6,0.3037) (8,0.002831)};\draw[densely dotted] (axis cs:6,0) -- (axis cs:6,28.84);\end{axis}\begin{axis}[width=72mm,height=42mm,axis y line*=right,axis x line=none,ymode=log,xmin=-1,xmax=9,ylabel={\scriptsize eindeutig (log)},tick label style={font=\tiny},label style={font=\scriptsize},ylabel shift=-3pt]\addplot[mark=*,mark size=0.7pt,draw=black,thin] coordinates {(0,4314588) (1,2055180) (2,5469556) (3,11446) (4,132859) (6,775)};\end{axis}\end{tikzpicture}\end{minipage}
\vspace{3mm}

\noindent\begin{minipage}[t]{60mm}\vspace{0pt}\sbox\hmboard{\chessboard[setfen=8/8/8/8/3Q4/3k4/2R5/K1Q5 b - - 0 1]}\usebox\hmboard\\[1mm]
\makebox[\wd\hmboard]{\sffamily\hspace{3mm}\textbf{h\#3}\hfill(4+1)}\end{minipage}\hfill
\begin{minipage}[t]{52mm}\vspace{0pt}
\textcolor{ink}{\sffamily\textsc{Lösung}}\par\smallskip\small
\begin{tabular}{@{}r@{\enspace}l@{\enspace}l@{}}
\textbf{1.} & \symking{}xd4 & \symking{}b2\\
\textbf{2.} & \symking{}e4 & \symqueen{}f4+\\
\textbf{3.} & \symking{}d3 & \symrook{}d2\#\\
\end{tabular}\par\vspace{3mm}
\qrcode[height=12mm]{https://helpman.komtera.lt/?fen=8/8/8/8/3Q4/3k4/2R5/K1Q5&moves=3}\\{\sffamily\tiny\color{ink}scannen zum Lösen}
\end{minipage}
\vspace{2mm}

{\raggedright\noindent\textcolor{ink}{\sffamily\textsc{Themen}}\enspace \themetag{Rundlauf} \themetag{ein Stein zieht}\par\smallskip{\footnotesize\textit{Hinweis:} die Partei am Zug steht im Diagramm im Schach; die Lösung beginnt mit einem Schlag.}\par}

\clearpage
\phantomsection\label{hm:KQRRvk}
\begin{tcolorbox}[enhanced,colback=ink,colframe=ink,arc=0pt,boxrule=0pt,left=4mm,right=4mm,top=2.5mm,bottom=2.5mm]
\color{white}\sffamily{\LARGE\bfseries KQRRvk}\hfill{\Large h\#3}\\[0.5mm]
{\small 5 Steine\quad Schwarz am Zug}\hfill{\small Nr.~219\qquad\href{https://helpman.komtera.lt/?fen=8/8/8/8/4R3/3k4/2R5/K1Q5&moves=3}{\color{white}\bfseries lösen\,\ensuremath{\nearrow}}}
\end{tcolorbox}
\vspace{2mm}
\noindent\begin{minipage}[t]{36mm}\vspace{0pt}\sffamily\small\raggedright
\hmfigure{längste eindeutige}{{\Large h\#3}}\hmfigure{eindeutig bei dieser Länge}{{\Large 1\,116} Stellungen}\hmfigure{längstes Matt}{{\Large h\#5}}\hmfigure{in der Klasse}{242 Millionen Stellungen,\\ 70,5 Millionen mit Hilfsmatt}
\end{minipage}\hfill
\begin{minipage}[t]{80mm}\vspace{0pt}\begin{tikzpicture}\begin{axis}[width=72mm,height=42mm,ybar,bar width=1.6pt,xmin=-1,xmax=11,ymin=0,ymax=31.08,axis y line*=left,axis x line*=bottom,xlabel={\scriptsize Halbzüge bis Matt},ylabel={\scriptsize Stellungen (Mio.)},tick label style={font=\tiny},label style={font=\scriptsize},enlarge x limits=false,ylabel shift=-3pt,xlabel shift=-3pt]\addplot[fill=ink!55,draw=none] coordinates {(0,2.643) (1,5.926) (2,18.88) (3,14.16) (4,28.26) (5,0.1497) (6,0.4347) (8,0.009621) (10,2e-06)};\draw[densely dotted] (axis cs:6,0) -- (axis cs:6,31.08);\end{axis}\begin{axis}[width=72mm,height=42mm,axis y line*=right,axis x line=none,ymode=log,xmin=-1,xmax=11,ylabel={\scriptsize eindeutig (log)},tick label style={font=\tiny},label style={font=\scriptsize},ylabel shift=-3pt]\addplot[mark=*,mark size=0.7pt,draw=black,thin] coordinates {(0,2642860) (1,2427322) (2,5183635) (3,107074) (4,155326) (6,1116)};\end{axis}\end{tikzpicture}\end{minipage}
\vspace{3mm}

\noindent\begin{minipage}[t]{60mm}\vspace{0pt}\sbox\hmboard{\chessboard[setfen=8/8/8/8/4R3/3k4/2R5/K1Q5 b - - 0 1]}\usebox\hmboard\\[1mm]
\makebox[\wd\hmboard]{\sffamily\hspace{3mm}\textbf{h\#3}\hfill(4+1)}\end{minipage}\hfill
\begin{minipage}[t]{52mm}\vspace{0pt}
\textcolor{ink}{\sffamily\textsc{Lösung}}\par\smallskip\small
\begin{tabular}{@{}r@{\enspace}l@{\enspace}l@{}}
\textbf{1.} & \symking{}xe4 & \symking{}b2\\
\textbf{2.} & \symking{}d4 & \symqueen{}f4+\\
\textbf{3.} & \symking{}d3 & \symrook{}d2\#\\
\end{tabular}\par\vspace{3mm}
\qrcode[height=12mm]{https://helpman.komtera.lt/?fen=8/8/8/8/4R3/3k4/2R5/K1Q5&moves=3}\\{\sffamily\tiny\color{ink}scannen zum Lösen}
\end{minipage}
\vspace{2mm}

{\raggedright\noindent\textcolor{ink}{\sffamily\textsc{Themen}}\enspace \themetag{Rundlauf} \themetag{ein Stein zieht}\par\smallskip{\footnotesize\textit{Hinweis:} die Lösung beginnt mit einem Schlag.}\par}

\clearpage
\phantomsection\label{hm:KRRRvk}
\begin{tcolorbox}[enhanced,colback=ink,colframe=ink,arc=0pt,boxrule=0pt,left=4mm,right=4mm,top=2.5mm,bottom=2.5mm]
\color{white}\sffamily{\LARGE\bfseries KRRRvk}\hfill{\Large h\#3}\\[0.5mm]
{\small 5 Steine\quad Schwarz am Zug}\hfill{\small Nr.~220\qquad\href{https://helpman.komtera.lt/?fen=8/8/8/8/3R4/2k5/1R6/K3R3&moves=3}{\color{white}\bfseries lösen\,\ensuremath{\nearrow}}}
\end{tcolorbox}
\vspace{2mm}
\noindent\begin{minipage}[t]{36mm}\vspace{0pt}\sffamily\small\raggedright
\hmfigure{längste eindeutige}{{\Large h\#3}}\hmfigure{eindeutig bei dieser Länge}{{\Large 217} Stellungen}\hmfigure{längstes Matt}{{\Large h\#5}}\hmfigure{in der Klasse}{242 Millionen Stellungen,\\ 25 Millionen mit Hilfsmatt}
\end{minipage}\hfill
\begin{minipage}[t]{80mm}\vspace{0pt}\begin{tikzpicture}\begin{axis}[width=72mm,height=42mm,ybar,bar width=1.6pt,xmin=-1,xmax=11,ymin=0,ymax=11.79,axis y line*=left,axis x line*=bottom,xlabel={\scriptsize Halbzüge bis Matt},ylabel={\scriptsize Stellungen (Mio.)},tick label style={font=\tiny},label style={font=\scriptsize},enlarge x limits=false,ylabel shift=-3pt,xlabel shift=-3pt]\addplot[fill=ink!55,draw=none] coordinates {(0,0.6094) (1,1.769) (2,5.428) (3,6.093) (4,10.72) (5,0.2149) (6,0.1677) (8,0.003144) (10,2e-06)};\draw[densely dotted] (axis cs:6,0) -- (axis cs:6,11.79);\end{axis}\begin{axis}[width=72mm,height=42mm,axis y line*=right,axis x line=none,ymode=log,xmin=-1,xmax=11,ylabel={\scriptsize eindeutig (log)},tick label style={font=\tiny},label style={font=\scriptsize},ylabel shift=-3pt]\addplot[mark=*,mark size=0.7pt,draw=black,thin] coordinates {(0,609449) (1,1107748) (2,2277776) (3,78560) (4,71114) (6,217)};\end{axis}\end{tikzpicture}\end{minipage}
\vspace{3mm}

\noindent\begin{minipage}[t]{60mm}\vspace{0pt}\sbox\hmboard{\chessboard[setfen=8/8/8/8/3R4/2k5/1R6/K3R3 b - - 0 1]}\usebox\hmboard\\[1mm]
\makebox[\wd\hmboard]{\sffamily\hspace{3mm}\textbf{h\#3}\hfill(4+1)}\end{minipage}\hfill
\begin{minipage}[t]{52mm}\vspace{0pt}
\textcolor{ink}{\sffamily\textsc{Lösung}}\par\smallskip\small
\begin{tabular}{@{}r@{\enspace}l@{\enspace}l@{}}
\textbf{1.} & \symking{}xd4 & \symrook{}bb1\\
\textbf{2.} & \symking{}d3 & \symking{}b2\\
\textbf{3.} & \symking{}d2 & \symrook{}bd1\#\\
\end{tabular}\par\vspace{3mm}
\qrcode[height=12mm]{https://helpman.komtera.lt/?fen=8/8/8/8/3R4/2k5/1R6/K3R3&moves=3}\\{\sffamily\tiny\color{ink}scannen zum Lösen}
\end{minipage}
\vspace{2mm}

{\raggedright\noindent\textcolor{ink}{\sffamily\textsc{Themen}}\enspace \themetag{ein Stein zieht} \themetag{ohne Schach}\par\smallskip{\footnotesize\textit{Hinweis:} die Lösung beginnt mit einem Schlag.}\par}

\clearpage
\phantomsection\label{hm:KQQQvk}
\begin{tcolorbox}[enhanced,colback=ink,colframe=ink,arc=0pt,boxrule=0pt,left=4mm,right=4mm,top=2.5mm,bottom=2.5mm]
\color{white}\sffamily{\LARGE\bfseries KQQQvk}\hfill{\Large h\#2}\\[0.5mm]
{\small 5 Steine\quad Schwarz am Zug}\hfill{\small Nr.~221\qquad\href{https://helpman.komtera.lt/?fen=5Q2/4Q3/8/8/8/8/8/K1Qk4&moves=2}{\color{white}\bfseries lösen\,\ensuremath{\nearrow}}}
\end{tcolorbox}
\vspace{2mm}
\noindent\begin{minipage}[t]{36mm}\vspace{0pt}\sffamily\small\raggedright
\hmfigure{längste eindeutige}{{\Large h\#2}}\hmfigure{eindeutig bei dieser Länge}{{\Large 3\,150} Stellungen}\hmfigure{längstes Matt}{{\Large h\#3}}\hmfigure{in der Klasse}{242 Millionen Stellungen,\\ 20,9 Millionen mit Hilfsmatt}
\end{minipage}\hfill
\begin{minipage}[t]{80mm}\vspace{0pt}\begin{tikzpicture}\begin{axis}[width=72mm,height=42mm,ybar,bar width=1.6pt,xmin=-1,xmax=7,ymin=0,ymax=11.11,axis y line*=left,axis x line*=bottom,xlabel={\scriptsize Halbzüge bis Matt},ylabel={\scriptsize Stellungen (Mio.)},tick label style={font=\tiny},label style={font=\scriptsize},enlarge x limits=false,ylabel shift=-3pt,xlabel shift=-3pt]\addplot[fill=ink!55,draw=none] coordinates {(0,2.289) (1,2.613) (2,10.1) (3,1.9) (4,4.009) (5,0.00011) (6,0.001777)};\draw[densely dotted] (axis cs:4,0) -- (axis cs:4,11.11);\end{axis}\begin{axis}[width=72mm,height=42mm,axis y line*=right,axis x line=none,ymode=log,xmin=-1,xmax=7,ylabel={\scriptsize eindeutig (log)},tick label style={font=\tiny},label style={font=\scriptsize},ylabel shift=-3pt]\addplot[mark=*,mark size=0.7pt,draw=black,thin] coordinates {(0,2288900) (1,514652) (2,1363338) (3,271) (4,3150)};\end{axis}\end{tikzpicture}\end{minipage}
\vspace{3mm}

\noindent\begin{minipage}[t]{60mm}\vspace{0pt}\sbox\hmboard{\chessboard[setfen=5Q2/4Q3/8/8/8/8/8/K1Qk4 b - - 0 1]}\usebox\hmboard\\[1mm]
\makebox[\wd\hmboard]{\sffamily\hspace{3mm}\textbf{h\#2}\hfill(4+1)}\end{minipage}\hfill
\begin{minipage}[t]{52mm}\vspace{0pt}
\textcolor{ink}{\sffamily\textsc{Lösung}}\par\smallskip\small
\begin{tabular}{@{}r@{\enspace}l@{\enspace}l@{}}
\textbf{1.} & \symking{}xc1 & \symqueen{}f2\\
\textbf{2.} & \symking{}d1 & \symqueen{}ee1\#\\
\end{tabular}\par\vspace{3mm}
\qrcode[height=12mm]{https://helpman.komtera.lt/?fen=5Q2/4Q3/8/8/8/8/8/K1Qk4&moves=2}\\{\sffamily\tiny\color{ink}scannen zum Lösen}
\end{minipage}
\vspace{2mm}

{\raggedright\noindent\textcolor{ink}{\sffamily\textsc{Themen}}\enspace \themetag{Switchback} \themetag{ein Stein zieht} \themetag{ohne Schach}\par\smallskip{\footnotesize\textit{Hinweis:} die Partei am Zug steht im Diagramm im Schach; die Lösung beginnt mit einem Schlag.}\par}

\clearpage

\chapter{Sechs Steine}
\markboth{Sechs Steine}{}

640 der 645 Materialklassen mit 6 Steinen, in denen ein Matt möglich ist, längste eindeutige Aufgabe zuerst.
\clearpage

\phantomsection\label{hm:KBvkbbp}
\begin{tcolorbox}[enhanced,colback=ink,colframe=ink,arc=0pt,boxrule=0pt,left=4mm,right=4mm,top=2.5mm,bottom=2.5mm]
\color{white}\sffamily{\LARGE\bfseries KBvkbbp}\hfill{\Large h\#13}\\[0.5mm]
{\small 6 Steine\quad Schwarz am Zug}\hfill{\small Nr.~222\qquad\href{https://helpman.komtera.lt/?fen=3k3b/4p3/8/8/K7/8/1b6/B7&moves=13}{\color{white}\bfseries lösen\,\ensuremath{\nearrow}}}
\end{tcolorbox}
\vspace{2mm}
\noindent\begin{minipage}[t]{36mm}\vspace{0pt}\sffamily\small\raggedright
\hmfigure{längste eindeutige}{{\Large h\#13}}\hmfigure{eindeutig bei dieser Länge}{{\Large 2} Stellungen}\hmfigure{längstes Matt}{{\Large h\#16.5}}\hmfigure{in der Klasse}{45,4 Milliarden Stellungen,\\ 14,6 Milliarden mit Hilfsmatt}
\end{minipage}\hfill
\begin{minipage}[t]{80mm}\vspace{0pt}\begin{tikzpicture}\begin{axis}[width=72mm,height=42mm,ybar,bar width=1.6pt,xmin=-1,xmax=34,ymin=0,ymax=2.322,axis y line*=left,axis x line*=bottom,xlabel={\scriptsize Halbzüge bis Matt},ylabel={\scriptsize Stellungen (Mrd.)},tick label style={font=\tiny},label style={font=\scriptsize},enlarge x limits=false,ylabel shift=-3pt,xlabel shift=-3pt]\addplot[fill=ink!55,draw=none] coordinates {(0,0.0001401) (1,0.0005582) (2,0.004159) (3,0.01342) (4,0.05655) (5,0.1421) (6,0.3431) (7,0.6794) (8,1.041) (9,1.637) (10,1.637) (11,2.111) (12,1.554) (13,1.265) (14,0.7062) (15,0.6196) (16,0.3742) (17,0.3884) (18,0.2768) (19,0.3006) (20,0.2403) (21,0.2691) (22,0.1944) (23,0.2193) (24,0.1261) (25,0.1449) (26,0.0717) (27,0.08325) (28,0.02809) (29,0.03369) (30,0.009178) (31,0.01102) (32,2.574e-05) (33,1.777e-06)};\draw[densely dotted] (axis cs:26,0) -- (axis cs:26,2.322);\end{axis}\begin{axis}[width=72mm,height=42mm,axis y line*=right,axis x line=none,ymode=log,xmin=-1,xmax=34,ylabel={\scriptsize eindeutig (log)},tick label style={font=\tiny},label style={font=\scriptsize},ylabel shift=-3pt]\addplot[mark=*,mark size=0.7pt,draw=black,thin] coordinates {(0,140074) (1,553567) (2,3498364) (3,5104420) (4,12138719) (5,18218366) (6,12087413) (7,14668891) (8,9183075) (9,10253500) (10,5727142) (11,4205417) (12,1466143) (13,568149) (14,282206) (15,125025) (16,57101) (17,21654) (18,12060) (19,3452) (20,2049) (21,915) (22,851) (23,753) (24,112) (25,89) (26,2)};\end{axis}\end{tikzpicture}\end{minipage}
\vspace{3mm}

\noindent\begin{minipage}[t]{60mm}\vspace{0pt}\sbox\hmboard{\chessboard[setfen=3k3b/4p3/8/8/K7/8/1b6/B7 b - - 0 1]}\usebox\hmboard\\[1mm]
\makebox[\wd\hmboard]{\sffamily\hspace{3mm}\textbf{h\#13}\hfill(2+4)}\end{minipage}\hfill
\begin{minipage}[t]{52mm}\vspace{0pt}
\textcolor{ink}{\sffamily\textsc{Lösung}}\par\smallskip\small
\begin{tabular}{@{}r@{\enspace}l@{\enspace}l@{}}
\textbf{1.} & e5 & \symking{}b5\\
\textbf{2.} & \symking{}e7 & \symking{}c6\\
\textbf{3.} & \symking{}f6 & \symking{}d7\\
\textbf{4.} & e4 & \symking{}e8\\
\textbf{5.} & e3 & \symking{}f8\\
\textbf{6.} & e2 & \symking{}g8\\
\textbf{7.} & e1=\symknight{} & \symking{}xh8\\
\textbf{8.} & \symknight{}f3 & \symking{}g8\\
\textbf{9.} & \symknight{}g5 & \symking{}f8\\
\textbf{10.} & \symknight{}h7+ & \symking{}e8\\
\textbf{11.} & \symking{}g7 & \symking{}e7\\
\textbf{12.} & \symking{}h8 & \symking{}f7\\
\textbf{13.} & \symbishop{}g7 & \symbishop{}xg7\#\\
\end{tabular}\par\vspace{3mm}
\qrcode[height=12mm]{https://helpman.komtera.lt/?fen=3k3b/4p3/8/8/K7/8/1b6/B7&moves=13}\\{\sffamily\tiny\color{ink}scannen zum Lösen}
\end{minipage}
\vspace{2mm}

{\raggedright\noindent\textcolor{ink}{\sffamily\textsc{Themen}}\enspace \themetag{reines Matt} \themetag{Mustermatt} \themetag{Idealmatt} \themetag{Umwandlung} \themetag{Unterverwandlung} \themetag{Excelsior} \themetag{Switchback} \themetag{Selbstblock} \themetag{Kniest}\par}

\clearpage
\phantomsection\label{hm:KQBvkrp}
\begin{tcolorbox}[enhanced,colback=ink,colframe=ink,arc=0pt,boxrule=0pt,left=4mm,right=4mm,top=2.5mm,bottom=2.5mm]
\color{white}\sffamily{\LARGE\bfseries KQBvkrp}\hfill{\Large h\#13}\\[0.5mm]
{\small 6 Steine\quad Schwarz am Zug}\hfill{\small Nr.~223\qquad\href{https://helpman.komtera.lt/?fen=K7/8/8/3p4/2k5/2Q5/8/5Br1&moves=13}{\color{white}\bfseries lösen\,\ensuremath{\nearrow}}}
\end{tcolorbox}
\vspace{2mm}
\noindent\begin{minipage}[t]{36mm}\vspace{0pt}\sffamily\small\raggedright
\hmfigure{längste eindeutige}{{\Large h\#13}}\hmfigure{eindeutig bei dieser Länge}{{\Large 2} Stellungen}\hmfigure{längstes Matt}{{\Large h\#17}}\hmfigure{in der Klasse}{45,4 Milliarden Stellungen,\\ 24,8 Milliarden mit Hilfsmatt}
\end{minipage}\hfill
\begin{minipage}[t]{80mm}\vspace{0pt}\begin{tikzpicture}\begin{axis}[width=72mm,height=42mm,ybar,bar width=1.6pt,xmin=-1,xmax=35,ymin=0,ymax=12,axis y line*=left,axis x line*=bottom,xlabel={\scriptsize Halbzüge bis Matt},ylabel={\scriptsize Stellungen (Mrd.)},tick label style={font=\tiny},label style={font=\scriptsize},enlarge x limits=false,ylabel shift=-3pt,xlabel shift=-3pt]\addplot[fill=ink!55,draw=none] coordinates {(0,0.0492) (1,0.3383) (2,1.597) (3,5.324) (4,10.91) (5,4.937) (6,1.372) (7,0.05361) (8,0.01503) (9,3.28e-07) (10,0.01478) (11,1.79e-07) (12,0.02797) (13,4.21e-07) (14,0.0343) (15,8.6e-08) (16,0.03423) (17,1.8e-08) (18,0.03494) (19,6.3e-08) (20,0.02911) (21,1.44e-07) (22,0.01802) (23,2.14e-07) (24,0.01277) (25,2.1e-08) (26,0.01101) (27,2.9e-08) (28,0.007166) (29,1.3e-08) (30,0.004782) (31,3e-09) (32,0.001441) (33,5e-09) (34,1.29e-07)};\draw[densely dotted] (axis cs:26,0) -- (axis cs:26,12);\end{axis}\begin{axis}[width=72mm,height=42mm,axis y line*=right,axis x line=none,ymode=log,xmin=-1,xmax=35,ylabel={\scriptsize eindeutig (log)},tick label style={font=\tiny},label style={font=\scriptsize},ylabel shift=-3pt]\addplot[mark=*,mark size=0.7pt,draw=black,thin] coordinates {(0,49198392) (1,290828957) (2,969841761) (3,1066175856) (4,508419513) (5,12351218) (6,489945) (7,162) (8,212132) (9,1) (10,122715) (11,1) (12,65656) (14,18702) (16,7657) (18,1089) (20,577) (22,243) (24,381) (26,2)};\end{axis}\end{tikzpicture}\end{minipage}
\vspace{3mm}

\noindent\begin{minipage}[t]{60mm}\vspace{0pt}\sbox\hmboard{\chessboard[setfen=K7/8/8/3p4/2k5/2Q5/8/5Br1 b - - 0 1]}\usebox\hmboard\\[1mm]
\makebox[\wd\hmboard]{\sffamily\hspace{3mm}\textbf{h\#13}\hfill(3+3)}\end{minipage}\hfill
\begin{minipage}[t]{52mm}\vspace{0pt}
\textcolor{ink}{\sffamily\textsc{Lösung}}\par\smallskip\small
\begin{tabular}{@{}r@{\enspace}l@{\enspace}l@{}}
\textbf{1.} & \symking{}xc3 & \symking{}b7\\
\textbf{2.} & d4 & \symking{}c6\\
\textbf{3.} & d3 & \symking{}d5\\
\textbf{4.} & d2 & \symking{}e4\\
\textbf{5.} & d1=\symknight{} & \symking{}f3\\
\textbf{6.} & \symknight{}e3 & \symking{}f2\\
\textbf{7.} & \symking{}d2 & \symking{}xg1\\
\textbf{8.} & \symking{}e1 & \symking{}h2\\
\textbf{9.} & \symking{}f2 & \symking{}h3\\
\textbf{10.} & \symking{}g1 & \symking{}g3\\
\textbf{11.} & \symking{}h1 & \symbishop{}h3\\
\textbf{12.} & \symknight{}f1+ & \symking{}f2\\
\textbf{13.} & \symknight{}h2 & \symbishop{}g2\#\\
\end{tabular}\par\vspace{3mm}
\qrcode[height=12mm]{https://helpman.komtera.lt/?fen=K7/8/8/3p4/2k5/2Q5/8/5Br1&moves=13}\\{\sffamily\tiny\color{ink}scannen zum Lösen}
\end{minipage}
\vspace{2mm}

{\raggedright\noindent\textcolor{ink}{\sffamily\textsc{Themen}}\enspace \themetag{reines Matt} \themetag{Mustermatt} \themetag{Idealmatt} \themetag{Umwandlung} \themetag{Unterverwandlung} \themetag{Rundlauf} \themetag{Selbstblock} \themetag{Umnov}\par\smallskip{\footnotesize\textit{Hinweis:} das Diagramm hat keinen legalen letzten Zug, kann also in keiner Partie entstehen.}\par}

\clearpage
\phantomsection\label{hm:KBvkqbp}
\begin{tcolorbox}[enhanced,colback=ink,colframe=ink,arc=0pt,boxrule=0pt,left=4mm,right=4mm,top=2.5mm,bottom=2.5mm]
\color{white}\sffamily{\LARGE\bfseries KBvkqbp}\hfill{\Large h\#12.5}\\[0.5mm]
{\small 6 Steine\quad Weiß am Zug}\hfill{\small Nr.~224\qquad\href{https://helpman.komtera.lt/?fen=7b/8/4pk2/8/8/2q5/8/B2K4&moves=12.5}{\color{white}\bfseries lösen\,\ensuremath{\nearrow}}}
\end{tcolorbox}
\vspace{2mm}
\noindent\begin{minipage}[t]{36mm}\vspace{0pt}\sffamily\small\raggedright
\hmfigure{längste eindeutige}{{\Large h\#12.5}}\hmfigure{eindeutig bei dieser Länge}{{\Large 59} Stellungen}\hmfigure{längstes Matt}{{\Large h\#16.5}}\hmfigure{in der Klasse}{45,4 Milliarden Stellungen,\\ 26,1 Milliarden mit Hilfsmatt}
\end{minipage}\hfill
\begin{minipage}[t]{80mm}\vspace{0pt}\begin{tikzpicture}\begin{axis}[width=72mm,height=42mm,ybar,bar width=1.6pt,xmin=-1,xmax=34,ymin=0,ymax=4.842,axis y line*=left,axis x line*=bottom,xlabel={\scriptsize Halbzüge bis Matt},ylabel={\scriptsize Stellungen (Mrd.)},tick label style={font=\tiny},label style={font=\scriptsize},enlarge x limits=false,ylabel shift=-3pt,xlabel shift=-3pt]\addplot[fill=ink!55,draw=none] coordinates {(0,4.973e-05) (1,0.0002185) (2,0.001753) (3,0.008493) (4,0.03542) (5,0.1293) (6,0.3214) (7,0.9645) (8,1.478) (9,3.322) (10,2.819) (11,4.402) (12,2.303) (13,2.22) (14,1.221) (15,1.614) (16,0.9804) (17,1.359) (18,0.5969) (19,0.8484) (20,0.3109) (21,0.4643) (22,0.1811) (23,0.282) (24,0.08084) (25,0.1282) (26,0.00802) (27,0.01755) (28,2.207e-05) (29,0.002724) (30,2.9e-08) (31,0.000931) (32,9e-09) (33,1.209e-05)};\draw[densely dotted] (axis cs:25,0) -- (axis cs:25,4.842);\end{axis}\begin{axis}[width=72mm,height=42mm,axis y line*=right,axis x line=none,ymode=log,xmin=-1,xmax=34,ylabel={\scriptsize eindeutig (log)},tick label style={font=\tiny},label style={font=\scriptsize},ylabel shift=-3pt]\addplot[mark=*,mark size=0.7pt,draw=black,thin] coordinates {(0,49734) (1,218169) (2,1212291) (3,4180712) (4,9817838) (5,22371373) (6,22031607) (7,42853322) (8,29648730) (9,37974587) (10,12598536) (11,11673068) (12,3417412) (13,2936401) (14,1268115) (15,1026816) (16,303118) (17,203161) (18,48043) (19,26693) (20,6315) (21,2993) (22,657) (23,333) (24,18) (25,59)};\end{axis}\end{tikzpicture}\end{minipage}
\vspace{3mm}

\noindent\begin{minipage}[t]{60mm}\vspace{0pt}\sbox\hmboard{\chessboard[setfen=7b/8/4pk2/8/8/2q5/8/B2K4 w - - 0 1]}\usebox\hmboard\\[1mm]
\makebox[\wd\hmboard]{\sffamily\hspace{3mm}\textbf{h\#12.5}\hfill(2+4)}\end{minipage}\hfill
\begin{minipage}[t]{52mm}\vspace{0pt}
\textcolor{ink}{\sffamily\textsc{Lösung}}\par\smallskip\small
\begin{tabular}{@{}r@{\enspace}l@{\enspace}l@{}}
\textbf{1.} & \ldots & \symking{}e2\\
\textbf{2.} & \symqueen{}b2+ & \symking{}f3\\
\textbf{3.} & e5 & \symking{}g4\\
\textbf{4.} & e4 & \symking{}h5\\
\textbf{5.} & e3 & \symking{}h6\\
\textbf{6.} & e2 & \symking{}h7\\
\textbf{7.} & e1=\symknight{} & \symking{}xh8\\
\textbf{8.} & \symknight{}f3 & \symking{}g8\\
\textbf{9.} & \symknight{}g5 & \symking{}f8\\
\textbf{10.} & \symknight{}h7+ & \symking{}e8\\
\textbf{11.} & \symking{}g7 & \symking{}e7\\
\textbf{12.} & \symking{}h8 & \symking{}f7\\
\textbf{13.} & \symqueen{}g7+ & \symbishop{}xg7\#\\
\end{tabular}\par\vspace{3mm}
\qrcode[height=12mm]{https://helpman.komtera.lt/?fen=7b/8/4pk2/8/8/2q5/8/B2K4&moves=12.5}\\{\sffamily\tiny\color{ink}scannen zum Lösen}
\end{minipage}
\vspace{2mm}

{\raggedright\noindent\textcolor{ink}{\sffamily\textsc{Themen}}\enspace \themetag{reines Matt} \themetag{Mustermatt} \themetag{Idealmatt} \themetag{Umwandlung} \themetag{Unterverwandlung} \themetag{Selbstblock} \themetag{Kniest}\par}

\clearpage
\phantomsection\label{hm:KBvkqqp}
\begin{tcolorbox}[enhanced,colback=ink,colframe=ink,arc=0pt,boxrule=0pt,left=4mm,right=4mm,top=2.5mm,bottom=2.5mm]
\color{white}\sffamily{\LARGE\bfseries KBvkqqp}\hfill{\Large h\#12.5}\\[0.5mm]
{\small 6 Steine\quad Weiß am Zug}\hfill{\small Nr.~225\qquad\href{https://helpman.komtera.lt/?fen=k7/B7/1qp5/8/8/K7/3q4/8&moves=12.5}{\color{white}\bfseries lösen\,\ensuremath{\nearrow}}}
\end{tcolorbox}
\vspace{2mm}
\noindent\begin{minipage}[t]{36mm}\vspace{0pt}\sffamily\small\raggedright
\hmfigure{längste eindeutige}{{\Large h\#12.5}}\hmfigure{eindeutig bei dieser Länge}{{\Large 210} Stellungen}\hmfigure{längstes Matt}{{\Large h\#17.5}}\hmfigure{in der Klasse}{45,4 Milliarden Stellungen,\\ 11,6 Milliarden mit Hilfsmatt}
\end{minipage}\hfill
\begin{minipage}[t]{80mm}\vspace{0pt}\begin{tikzpicture}\begin{axis}[width=72mm,height=42mm,ybar,bar width=1.6pt,xmin=-1,xmax=36,ymin=0,ymax=1.756,axis y line*=left,axis x line*=bottom,xlabel={\scriptsize Halbzüge bis Matt},ylabel={\scriptsize Stellungen (Mrd.)},tick label style={font=\tiny},label style={font=\scriptsize},enlarge x limits=false,ylabel shift=-3pt,xlabel shift=-3pt]\addplot[fill=ink!55,draw=none] coordinates {(0,3.372e-06) (1,1.35e-05) (2,8.84e-05) (3,0.0004982) (4,0.001633) (5,0.007684) (6,0.01807) (7,0.0703) (8,0.1165) (9,0.362) (10,0.4119) (11,1.028) (12,0.8082) (13,1.589) (14,0.9224) (15,1.597) (16,0.764) (17,1.342) (18,0.4517) (19,0.823) (20,0.2275) (21,0.4338) (22,0.1335) (23,0.2656) (24,0.05417) (25,0.1155) (26,0.00321) (27,0.01719) (28,3.488e-06) (29,0.006044) (30,4.9e-08) (31,0.002585) (32,1.9e-08) (33,0.0002983) (34,3e-09) (35,3e-09)};\draw[densely dotted] (axis cs:25,0) -- (axis cs:25,1.756);\end{axis}\begin{axis}[width=72mm,height=42mm,axis y line*=right,axis x line=none,ymode=log,xmin=-1,xmax=36,ylabel={\scriptsize eindeutig (log)},tick label style={font=\tiny},label style={font=\scriptsize},ylabel shift=-3pt]\addplot[mark=*,mark size=0.7pt,draw=black,thin] coordinates {(0,3372) (1,13498) (2,45352) (3,193910) (4,331500) (5,1219031) (6,1366384) (7,3950075) (8,2997221) (9,6426550) (10,3561383) (11,6011622) (12,2476846) (13,3291605) (14,1140397) (15,1417985) (16,349546) (17,384729) (18,73583) (19,74611) (20,16459) (21,17656) (22,4439) (23,4326) (24,223) (25,210)};\end{axis}\end{tikzpicture}\end{minipage}
\vspace{3mm}

\noindent\begin{minipage}[t]{60mm}\vspace{0pt}\sbox\hmboard{\chessboard[setfen=k7/B7/1qp5/8/8/K7/3q4/8 w - - 0 1]}\usebox\hmboard\\[1mm]
\makebox[\wd\hmboard]{\sffamily\hspace{3mm}\textbf{h\#12.5}\hfill(2+4)}\end{minipage}\hfill
\begin{minipage}[t]{52mm}\vspace{0pt}
\textcolor{ink}{\sffamily\textsc{Lösung}}\par\smallskip\small
\begin{tabular}{@{}r@{\enspace}l@{\enspace}l@{}}
\textbf{1.} & \ldots & \symbishop{}b8\\
\textbf{2.} & c5 & \symbishop{}d6\\
\textbf{3.} & c4 & \symbishop{}b4\\
\textbf{4.} & c3 & \symking{}b3\\
\textbf{5.} & c2 & \symking{}c4\\
\textbf{6.} & c1=\symknight{} & \symbishop{}d6\\
\textbf{7.} & \symknight{}d3 & \symking{}d5\\
\textbf{8.} & \symqueen{}da5+ & \symking{}e6\\
\textbf{9.} & \symknight{}c5+ & \symking{}e7\\
\textbf{10.} & \symknight{}b7 & \symbishop{}b4\\
\textbf{11.} & \symking{}b8 & \symking{}d7\\
\textbf{12.} & \symqueen{}a8 & \symbishop{}a5\\
\textbf{13.} & \symqueen{}ba7 & \symbishop{}c7\#\\
\end{tabular}\par\vspace{3mm}
\qrcode[height=12mm]{https://helpman.komtera.lt/?fen=k7/B7/1qp5/8/8/K7/3q4/8&moves=12.5}\\{\sffamily\tiny\color{ink}scannen zum Lösen}
\end{minipage}
\vspace{2mm}

{\raggedright\noindent\textcolor{ink}{\sffamily\textsc{Themen}}\enspace \themetag{reines Matt} \themetag{Mustermatt} \themetag{Idealmatt} \themetag{Umwandlung} \themetag{Unterverwandlung} \themetag{Switchback} \themetag{Rundlauf} \themetag{Selbstblock} \themetag{ohne Schlag} \themetag{Umnov} \themetag{Klasinc}\par}

\clearpage
\phantomsection\label{hm:KBvkqrp}
\begin{tcolorbox}[enhanced,colback=ink,colframe=ink,arc=0pt,boxrule=0pt,left=4mm,right=4mm,top=2.5mm,bottom=2.5mm]
\color{white}\sffamily{\LARGE\bfseries KBvkqrp}\hfill{\Large h\#12.5}\\[0.5mm]
{\small 6 Steine\quad Weiß am Zug}\hfill{\small Nr.~226\qquad\href{https://helpman.komtera.lt/?fen=1k6/8/4p3/8/1q6/8/1r6/K1B5&moves=12.5}{\color{white}\bfseries lösen\,\ensuremath{\nearrow}}}
\end{tcolorbox}
\vspace{2mm}
\noindent\begin{minipage}[t]{36mm}\vspace{0pt}\sffamily\small\raggedright
\hmfigure{längste eindeutige}{{\Large h\#12.5}}\hmfigure{eindeutig bei dieser Länge}{{\Large 495} Stellungen}\hmfigure{längstes Matt}{{\Large h\#16.5}}\hmfigure{in der Klasse}{45,4 Milliarden Stellungen,\\ 25 Milliarden mit Hilfsmatt}
\end{minipage}\hfill
\begin{minipage}[t]{80mm}\vspace{0pt}\begin{tikzpicture}\begin{axis}[width=72mm,height=42mm,ybar,bar width=1.6pt,xmin=-1,xmax=34,ymin=0,ymax=3.52,axis y line*=left,axis x line*=bottom,xlabel={\scriptsize Halbzüge bis Matt},ylabel={\scriptsize Stellungen (Mrd.)},tick label style={font=\tiny},label style={font=\scriptsize},enlarge x limits=false,ylabel shift=-3pt,xlabel shift=-3pt]\addplot[fill=ink!55,draw=none] coordinates {(0,1.016e-05) (1,4.023e-05) (2,0.0002259) (3,0.001196) (4,0.004106) (5,0.01759) (6,0.04528) (7,0.1527) (8,0.2862) (9,0.7566) (10,0.9974) (11,2.114) (12,1.932) (13,3.199) (14,2.16) (15,3.2) (16,1.819) (17,2.734) (18,1.098) (19,1.691) (20,0.5643) (21,0.9005) (22,0.3328) (23,0.5517) (24,0.1549) (25,0.2627) (26,0.01628) (27,0.03734) (28,6.272e-05) (29,0.005903) (30,6.1e-08) (31,0.002562) (32,8e-09) (33,6.465e-05)};\draw[densely dotted] (axis cs:25,0) -- (axis cs:25,3.52);\end{axis}\begin{axis}[width=72mm,height=42mm,axis y line*=right,axis x line=none,ymode=log,xmin=-1,xmax=34,ylabel={\scriptsize eindeutig (log)},tick label style={font=\tiny},label style={font=\scriptsize},ylabel shift=-3pt]\addplot[mark=*,mark size=0.7pt,draw=black,thin] coordinates {(0,10162) (1,40232) (2,116003) (3,463310) (4,1007126) (5,3015008) (6,3652588) (7,8460137) (8,7781273) (9,13010384) (10,8700696) (11,11485218) (12,5554387) (13,5752886) (14,2404245) (15,2312013) (16,690927) (17,597769) (18,134948) (19,108382) (20,26888) (21,21881) (22,6512) (23,6125) (24,256) (25,495)};\end{axis}\end{tikzpicture}\end{minipage}
\vspace{3mm}

\noindent\begin{minipage}[t]{60mm}\vspace{0pt}\sbox\hmboard{\chessboard[setfen=1k6/8/4p3/8/1q6/8/1r6/K1B5 w - - 0 1]}\usebox\hmboard\\[1mm]
\makebox[\wd\hmboard]{\sffamily\hspace{3mm}\textbf{h\#12.5}\hfill(2+4)}\end{minipage}\hfill
\begin{minipage}[t]{52mm}\vspace{0pt}
\textcolor{ink}{\sffamily\textsc{Lösung}}\par\smallskip\small
\begin{tabular}{@{}r@{\enspace}l@{\enspace}l@{}}
\textbf{1.} & \ldots & \symbishop{}e3\\
\textbf{2.} & \symqueen{}b7 & \symbishop{}b6\\
\textbf{3.} & \symrook{}a2+ & \symking{}b1\\
\textbf{4.} & e5 & \symking{}c1\\
\textbf{5.} & e4 & \symking{}d1\\
\textbf{6.} & e3 & \symking{}e1\\
\textbf{7.} & e2 & \symking{}f2\\
\textbf{8.} & e1=\symknight{}+ & \symking{}e3\\
\textbf{9.} & \symknight{}d3 & \symking{}d4\\
\textbf{10.} & \symknight{}c5 & \symking{}e5\\
\textbf{11.} & \symqueen{}a8 & \symking{}d6\\
\textbf{12.} & \symknight{}b7+ & \symking{}d7\\
\textbf{13.} & \symrook{}a7 & \symbishop{}c7\#\\
\end{tabular}\par\vspace{3mm}
\qrcode[height=12mm]{https://helpman.komtera.lt/?fen=1k6/8/4p3/8/1q6/8/1r6/K1B5&moves=12.5}\\{\sffamily\tiny\color{ink}scannen zum Lösen}
\end{minipage}
\vspace{2mm}

{\raggedright\noindent\textcolor{ink}{\sffamily\textsc{Themen}}\enspace \themetag{reines Matt} \themetag{Mustermatt} \themetag{Idealmatt} \themetag{Umwandlung} \themetag{Unterverwandlung} \themetag{Selbstblock} \themetag{ohne Schlag} \themetag{Umnov}\par}

\clearpage
\phantomsection\label{hm:KBvkrpp}
\begin{tcolorbox}[enhanced,colback=ink,colframe=ink,arc=0pt,boxrule=0pt,left=4mm,right=4mm,top=2.5mm,bottom=2.5mm]
\color{white}\sffamily{\LARGE\bfseries KBvkrpp}\hfill{\Large h\#12.5}\\[0.5mm]
{\small 6 Steine\quad Weiß am Zug}\hfill{\small Nr.~227\qquad\href{https://helpman.komtera.lt/?fen=8/3p4/3p4/5k2/8/K7/7r/3B4&moves=12.5}{\color{white}\bfseries lösen\,\ensuremath{\nearrow}}}
\end{tcolorbox}
\vspace{2mm}
\noindent\begin{minipage}[t]{36mm}\vspace{0pt}\sffamily\small\raggedright
\hmfigure{längste eindeutige}{{\Large h\#12.5}}\hmfigure{eindeutig bei dieser Länge}{{\Large 14} Stellungen}\hmfigure{längstes Matt}{{\Large h\#16.5}}\hmfigure{in der Klasse}{34,1 Milliarden Stellungen,\\ 11,1 Milliarden mit Hilfsmatt}
\end{minipage}\hfill
\begin{minipage}[t]{80mm}\vspace{0pt}\begin{tikzpicture}\begin{axis}[width=72mm,height=42mm,ybar,bar width=1.6pt,xmin=-1,xmax=34,ymin=0,ymax=1.554,axis y line*=left,axis x line*=bottom,xlabel={\scriptsize Halbzüge bis Matt},ylabel={\scriptsize Stellungen (Mrd.)},tick label style={font=\tiny},label style={font=\scriptsize},enlarge x limits=false,ylabel shift=-3pt,xlabel shift=-3pt]\addplot[fill=ink!55,draw=none] coordinates {(0,2.318e-05) (1,9.946e-05) (2,0.0003783) (3,0.00162) (4,0.005734) (5,0.0186) (6,0.05337) (7,0.1275) (8,0.2814) (9,0.508) (10,0.8283) (11,1.189) (12,1.371) (13,1.413) (14,1.122) (15,1.096) (16,0.7629) (17,0.792) (18,0.4394) (19,0.4611) (20,0.1766) (21,0.1888) (22,0.06432) (23,0.07174) (24,0.03023) (25,0.03446) (26,0.01029) (27,0.01205) (28,0.001612) (29,0.001896) (30,0.0002538) (31,0.0003146) (32,2.1e-08) (33,8e-09)};\draw[densely dotted] (axis cs:25,0) -- (axis cs:25,1.554);\end{axis}\begin{axis}[width=72mm,height=42mm,axis y line*=right,axis x line=none,ymode=log,xmin=-1,xmax=34,ylabel={\scriptsize eindeutig (log)},tick label style={font=\tiny},label style={font=\scriptsize},ylabel shift=-3pt]\addplot[mark=*,mark size=0.7pt,draw=black,thin] coordinates {(0,23183) (1,99288) (2,207391) (3,629140) (4,2030487) (5,3260163) (6,4850829) (7,6157620) (8,8091001) (9,6392586) (10,6085595) (11,4225323) (12,2734452) (13,1223043) (14,718939) (15,298825) (16,114116) (17,41899) (18,14869) (19,8316) (20,2231) (21,2238) (22,954) (23,960) (24,149) (25,14)};\end{axis}\end{tikzpicture}\end{minipage}
\vspace{3mm}

\noindent\begin{minipage}[t]{60mm}\vspace{0pt}\sbox\hmboard{\chessboard[setfen=8/3p4/3p4/5k2/8/K7/7r/3B4 w - - 0 1]}\usebox\hmboard\\[1mm]
\makebox[\wd\hmboard]{\sffamily\hspace{3mm}\textbf{h\#12.5}\hfill(2+4)}\end{minipage}\hfill
\begin{minipage}[t]{52mm}\vspace{0pt}
\textcolor{ink}{\sffamily\textsc{Lösung}}\par\smallskip\small
\begin{tabular}{@{}r@{\enspace}l@{\enspace}l@{}}
\textbf{1.} & \ldots & \symbishop{}c2+\\
\textbf{2.} & \symking{}g4 & \symking{}b2\\
\textbf{3.} & d5 & \symking{}c1\\
\textbf{4.} & d4 & \symking{}d1\\
\textbf{5.} & d3 & \symking{}e1\\
\textbf{6.} & d2+ & \symking{}f1\\
\textbf{7.} & d1=\symknight{} & \symking{}g1\\
\textbf{8.} & \symknight{}e3 & \symking{}xh2\\
\textbf{9.} & \symknight{}f1+ & \symking{}g1\\
\textbf{10.} & \symking{}h3 & \symking{}f2\\
\textbf{11.} & \symking{}h2 & \symbishop{}f5\\
\textbf{12.} & \symking{}h1 & \symbishop{}h3\\
\textbf{13.} & \symknight{}h2 & \symbishop{}g2\#\\
\end{tabular}\par\vspace{3mm}
\qrcode[height=12mm]{https://helpman.komtera.lt/?fen=8/3p4/3p4/5k2/8/K7/7r/3B4&moves=12.5}\\{\sffamily\tiny\color{ink}scannen zum Lösen}
\end{minipage}
\vspace{2mm}

{\raggedright\noindent\textcolor{ink}{\sffamily\textsc{Themen}}\enspace \themetag{reines Matt} \themetag{Mustermatt} \themetag{Umwandlung} \themetag{Unterverwandlung} \themetag{Switchback} \themetag{Selbstblock}\par}

\clearpage
\phantomsection\label{hm:KBvkrrp}
\begin{tcolorbox}[enhanced,colback=ink,colframe=ink,arc=0pt,boxrule=0pt,left=4mm,right=4mm,top=2.5mm,bottom=2.5mm]
\color{white}\sffamily{\LARGE\bfseries KBvkrrp}\hfill{\Large h\#12.5}\\[0.5mm]
{\small 6 Steine\quad Weiß am Zug}\hfill{\small Nr.~228\qquad\href{https://helpman.komtera.lt/?fen=8/2k2p2/5r2/8/8/6r1/1K6/7B&moves=12.5}{\color{white}\bfseries lösen\,\ensuremath{\nearrow}}}
\end{tcolorbox}
\vspace{2mm}
\noindent\begin{minipage}[t]{36mm}\vspace{0pt}\sffamily\small\raggedright
\hmfigure{längste eindeutige}{{\Large h\#12.5}}\hmfigure{eindeutig bei dieser Länge}{{\Large 79} Stellungen}\hmfigure{längstes Matt}{{\Large h\#16.5}}\hmfigure{in der Klasse}{45,4 Milliarden Stellungen,\\ 13,5 Milliarden mit Hilfsmatt}
\end{minipage}\hfill
\begin{minipage}[t]{80mm}\vspace{0pt}\begin{tikzpicture}\begin{axis}[width=72mm,height=42mm,ybar,bar width=1.6pt,xmin=-1,xmax=34,ymin=0,ymax=1.757,axis y line*=left,axis x line*=bottom,xlabel={\scriptsize Halbzüge bis Matt},ylabel={\scriptsize Stellungen (Mrd.)},tick label style={font=\tiny},label style={font=\scriptsize},enlarge x limits=false,ylabel shift=-3pt,xlabel shift=-3pt]\addplot[fill=ink!55,draw=none] coordinates {(0,7.338e-06) (1,2.887e-05) (2,0.000137) (3,0.0006845) (4,0.002499) (5,0.00955) (6,0.02722) (7,0.07867) (8,0.1687) (9,0.3785) (10,0.5825) (11,1.05) (12,1.129) (13,1.588) (14,1.261) (15,1.597) (16,1.078) (17,1.377) (18,0.6662) (19,0.8636) (20,0.3508) (21,0.4688) (22,0.2052) (23,0.2844) (24,0.104) (25,0.1441) (26,0.01766) (27,0.0257) (28,0.0001022) (29,0.0009388) (30,5.7e-08) (31,0.0003916) (32,4e-09) (33,1.149e-06)};\draw[densely dotted] (axis cs:25,0) -- (axis cs:25,1.757);\end{axis}\begin{axis}[width=72mm,height=42mm,axis y line*=right,axis x line=none,ymode=log,xmin=-1,xmax=34,ylabel={\scriptsize eindeutig (log)},tick label style={font=\tiny},label style={font=\scriptsize},ylabel shift=-3pt]\addplot[mark=*,mark size=0.7pt,draw=black,thin] coordinates {(0,7338) (1,28868) (2,71262) (3,261282) (4,795424) (5,1747974) (6,2431860) (7,4447343) (8,5095740) (9,6329539) (10,5067837) (11,4955884) (12,2979294) (13,2328429) (14,1222226) (15,921227) (16,311750) (17,201458) (18,58988) (19,37265) (20,11673) (21,7741) (22,2413) (23,2097) (24,130) (25,79)};\end{axis}\end{tikzpicture}\end{minipage}
\vspace{3mm}

\noindent\begin{minipage}[t]{60mm}\vspace{0pt}\sbox\hmboard{\chessboard[setfen=8/2k2p2/5r2/8/8/6r1/1K6/7B w - - 0 1]}\usebox\hmboard\\[1mm]
\makebox[\wd\hmboard]{\sffamily\hspace{3mm}\textbf{h\#12.5}\hfill(2+4)}\end{minipage}\hfill
\begin{minipage}[t]{52mm}\vspace{0pt}
\textcolor{ink}{\sffamily\textsc{Lösung}}\par\smallskip\small
\begin{tabular}{@{}r@{\enspace}l@{\enspace}l@{}}
\textbf{1.} & \ldots & \symbishop{}e4\\
\textbf{2.} & \symking{}d6 & \symbishop{}d3\\
\textbf{3.} & \symking{}e5 & \symking{}c3\\
\textbf{4.} & \symking{}f4 & \symking{}d4\\
\textbf{5.} & \symking{}f3 & \symking{}e5\\
\textbf{6.} & \symking{}g2 & \symking{}xf6\\
\textbf{7.} & \symking{}h1 & \symking{}e5\\
\textbf{8.} & f5 & \symbishop{}f1\\
\textbf{9.} & f4 & \symbishop{}h3\\
\textbf{10.} & f3 & \symking{}f4\\
\textbf{11.} & f2 & \symking{}xg3\\
\textbf{12.} & f1=\symknight{}+ & \symking{}f2\\
\textbf{13.} & \symknight{}h2 & \symbishop{}g2\#\\
\end{tabular}\par\vspace{3mm}
\qrcode[height=12mm]{https://helpman.komtera.lt/?fen=8/2k2p2/5r2/8/8/6r1/1K6/7B&moves=12.5}\\{\sffamily\tiny\color{ink}scannen zum Lösen}
\end{minipage}
\vspace{2mm}

{\raggedright\noindent\textcolor{ink}{\sffamily\textsc{Themen}}\enspace \themetag{reines Matt} \themetag{Mustermatt} \themetag{Idealmatt} \themetag{Umwandlung} \themetag{Unterverwandlung} \themetag{Excelsior} \themetag{Switchback} \themetag{Selbstblock} \themetag{Umnov}\par}

\clearpage
\phantomsection\label{hm:KPvkbpp}
\begin{tcolorbox}[enhanced,colback=ink,colframe=ink,arc=0pt,boxrule=0pt,left=4mm,right=4mm,top=2.5mm,bottom=2.5mm]
\color{white}\sffamily{\LARGE\bfseries KPvkbpp}\hfill{\Large h\#12.5}\\[0.5mm]
{\small 6 Steine\quad Weiß am Zug}\hfill{\small Nr.~229\qquad\href{https://helpman.komtera.lt/?fen=8/8/7p/7p/k6P/7b/8/K7&moves=12.5}{\color{white}\bfseries lösen\,\ensuremath{\nearrow}}}
\end{tcolorbox}
\vspace{2mm}
\noindent\begin{minipage}[t]{36mm}\vspace{0pt}\sffamily\small\raggedright
\hmfigure{längste eindeutige}{{\Large h\#12.5}}\hmfigure{eindeutig bei dieser Länge}{{\Large 7} Stellungen}\hmfigure{längstes Matt}{{\Large h\#15.5}}\hmfigure{in der Klasse}{25,6 Milliarden Stellungen,\\ 9,03 Milliarden mit Hilfsmatt}
\end{minipage}\hfill
\begin{minipage}[t]{80mm}\vspace{0pt}\begin{tikzpicture}\begin{axis}[width=72mm,height=42mm,ybar,bar width=1.6pt,xmin=-1,xmax=32,ymin=0,ymax=1.548,axis y line*=left,axis x line*=bottom,xlabel={\scriptsize Halbzüge bis Matt},ylabel={\scriptsize Stellungen (Mrd.)},tick label style={font=\tiny},label style={font=\scriptsize},enlarge x limits=false,ylabel shift=-3pt,xlabel shift=-3pt]\addplot[fill=ink!55,draw=none] coordinates {(0,8.777e-06) (1,0.001471) (2,0.004066) (3,0.02896) (4,0.07138) (5,0.3007) (6,0.4876) (7,0.9229) (8,0.9417) (9,1.262) (10,1.06) (11,1.407) (12,1.161) (13,0.7895) (14,0.3461) (15,0.1172) (16,0.05049) (17,0.03242) (18,0.0139) (19,0.01286) (20,0.007904) (21,0.006027) (22,0.00376) (23,0.002264) (24,0.001185) (25,0.0004712) (26,0.0002018) (27,0.0001158) (28,8.336e-05) (29,2.883e-05) (30,2.321e-05) (31,7.74e-07)};\draw[densely dotted] (axis cs:25,0) -- (axis cs:25,1.548);\end{axis}\begin{axis}[width=72mm,height=42mm,axis y line*=right,axis x line=none,ymode=log,xmin=-1,xmax=32,ylabel={\scriptsize eindeutig (log)},tick label style={font=\tiny},label style={font=\scriptsize},ylabel shift=-3pt]\addplot[mark=*,mark size=0.7pt,draw=black,thin] coordinates {(0,8777) (1,392363) (2,887991) (3,7344835) (4,13624297) (5,28777524) (6,17918393) (7,17885214) (8,10263134) (9,12117903) (10,6801692) (11,7811155) (12,3347677) (13,811288) (14,283513) (15,63969) (16,15537) (17,4856) (18,1310) (19,659) (20,202) (21,126) (22,53) (23,29) (24,15) (25,7)};\end{axis}\end{tikzpicture}\end{minipage}
\vspace{3mm}

\noindent\begin{minipage}[t]{60mm}\vspace{0pt}\sbox\hmboard{\chessboard[setfen=8/8/7p/7p/k6P/7b/8/K7 w - - 0 1]}\usebox\hmboard\\[1mm]
\makebox[\wd\hmboard]{\sffamily\hspace{3mm}\textbf{h\#12.5}\hfill(2+4)}\end{minipage}\hfill
\begin{minipage}[t]{52mm}\vspace{0pt}
\textcolor{ink}{\sffamily\textsc{Lösung}}\par\smallskip\small
\begin{tabular}{@{}r@{\enspace}l@{\enspace}l@{}}
\textbf{1.} & \ldots & \symking{}b2\\
\textbf{2.} & \symking{}b5 & \symking{}c3\\
\textbf{3.} & \symking{}c6 & \symking{}d4\\
\textbf{4.} & \symking{}d7 & \symking{}e5\\
\textbf{5.} & \symking{}e8 & \symking{}f6\\
\textbf{6.} & \symking{}f8 & \symking{}g6\\
\textbf{7.} & \symking{}g8 & \symking{}xh5\\
\textbf{8.} & \symking{}h8 & \symking{}g6\\
\textbf{9.} & \symbishop{}f5+ & \symking{}f7\\
\textbf{10.} & \symking{}h7 & h5\\
\textbf{11.} & \symbishop{}g6+ & hxg6+\\
\textbf{12.} & \symking{}h8 & g7+\\
\textbf{13.} & \symking{}h7 & g8=\symqueen{}\#\\
\end{tabular}\par\vspace{3mm}
\qrcode[height=12mm]{https://helpman.komtera.lt/?fen=8/8/7p/7p/k6P/7b/8/K7&moves=12.5}\\{\sffamily\tiny\color{ink}scannen zum Lösen}
\end{minipage}
\vspace{2mm}

{\raggedright\noindent\textcolor{ink}{\sffamily\textsc{Themen}}\enspace \themetag{Umwandlung} \themetag{Switchback}\par}

\clearpage
\phantomsection\label{hm:KBBPvkp}
\begin{tcolorbox}[enhanced,colback=ink,colframe=ink,arc=0pt,boxrule=0pt,left=4mm,right=4mm,top=2.5mm,bottom=2.5mm]
\color{white}\sffamily{\LARGE\bfseries KBBPvkp}\hfill{\Large h\#12}\\[0.5mm]
{\small 6 Steine\quad Schwarz am Zug}\hfill{\small Nr.~230\qquad\href{https://helpman.komtera.lt/?fen=8/4B3/8/6Bp/7P/k7/8/K7&moves=12}{\color{white}\bfseries lösen\,\ensuremath{\nearrow}}}
\end{tcolorbox}
\vspace{2mm}
\noindent\begin{minipage}[t]{36mm}\vspace{0pt}\sffamily\small\raggedright
\hmfigure{längste eindeutige}{{\Large h\#12}}\hmfigure{eindeutig bei dieser Länge}{{\Large 4} Stellungen}\hmfigure{längstes Matt}{{\Large h\#15}}\hmfigure{in der Klasse}{34,1 Milliarden Stellungen,\\ 11,6 Milliarden mit Hilfsmatt}
\end{minipage}\hfill
\begin{minipage}[t]{80mm}\vspace{0pt}\begin{tikzpicture}\begin{axis}[width=72mm,height=42mm,ybar,bar width=1.6pt,xmin=-1,xmax=31,ymin=0,ymax=2.458,axis y line*=left,axis x line*=bottom,xlabel={\scriptsize Halbzüge bis Matt},ylabel={\scriptsize Stellungen (Mrd.)},tick label style={font=\tiny},label style={font=\scriptsize},enlarge x limits=false,ylabel shift=-3pt,xlabel shift=-3pt]\addplot[fill=ink!55,draw=none] coordinates {(0,0.002208) (1,0.01285) (2,0.03977) (3,0.2528) (4,0.7416) (5,1.447) (6,2.235) (7,1.449) (8,1.655) (9,0.8966) (10,0.9482) (11,0.6844) (12,0.7652) (13,0.2119) (14,0.1607) (15,0.0204) (16,0.01619) (17,0.005301) (18,0.006293) (19,0.004168) (20,0.005185) (21,0.002289) (22,0.002596) (23,0.0009534) (24,0.001112) (25,0.0002044) (26,0.0002567) (27,0.0001306) (28,0.0001855) (29,8.869e-06) (30,5.183e-06)};\draw[densely dotted] (axis cs:24,0) -- (axis cs:24,2.458);\end{axis}\begin{axis}[width=72mm,height=42mm,axis y line*=right,axis x line=none,ymode=log,xmin=-1,xmax=31,ylabel={\scriptsize eindeutig (log)},tick label style={font=\tiny},label style={font=\scriptsize},ylabel shift=-3pt]\addplot[mark=*,mark size=0.7pt,draw=black,thin] coordinates {(0,2207720) (1,8782466) (2,25998334) (3,128045471) (4,177353477) (5,97670198) (6,72687927) (7,26408241) (8,16958122) (9,9597841) (10,7272378) (11,4154034) (12,2429105) (13,150548) (14,112306) (15,13477) (16,10217) (17,4285) (18,2782) (19,559) (20,305) (21,44) (22,32) (23,9) (24,4)};\end{axis}\end{tikzpicture}\end{minipage}
\vspace{3mm}

\noindent\begin{minipage}[t]{60mm}\vspace{0pt}\sbox\hmboard{\chessboard[setfen=8/4B3/8/6Bp/7P/k7/8/K7 b - - 0 1]}\usebox\hmboard\\[1mm]
\makebox[\wd\hmboard]{\sffamily\hspace{3mm}\textbf{h\#12}\hfill(4+2)}\end{minipage}\hfill
\begin{minipage}[t]{52mm}\vspace{0pt}
\textcolor{ink}{\sffamily\textsc{Lösung}}\par\smallskip\small
\begin{tabular}{@{}r@{\enspace}l@{\enspace}l@{}}
\textbf{1.} & \symking{}b3 & \symking{}b1\\
\textbf{2.} & \symking{}c4 & \symking{}c2\\
\textbf{3.} & \symking{}d5 & \symking{}d3\\
\textbf{4.} & \symking{}e6 & \symking{}e4\\
\textbf{5.} & \symking{}f7 & \symking{}f5\\
\textbf{6.} & \symking{}g8 & \symking{}g6\\
\textbf{7.} & \symking{}h8 & \symking{}xh5\\
\textbf{8.} & \symking{}g8 & \symking{}g6\\
\textbf{9.} & \symking{}h8 & h5\\
\textbf{10.} & \symking{}g8 & h6\\
\textbf{11.} & \symking{}h8 & \symbishop{}gf6+\\
\textbf{12.} & \symking{}g8 & h7\#\\
\end{tabular}\par\vspace{3mm}
\qrcode[height=12mm]{https://helpman.komtera.lt/?fen=8/4B3/8/6Bp/7P/k7/8/K7&moves=12}\\{\sffamily\tiny\color{ink}scannen zum Lösen}
\end{minipage}
\vspace{2mm}

{\raggedright\noindent\textcolor{ink}{\sffamily\textsc{Themen}}\enspace \themetag{Switchback} \themetag{ein Stein zieht}\par\smallskip{\footnotesize\textit{Hinweis:} die Partei am Zug steht im Diagramm im Schach.}\par}

\clearpage
\phantomsection\label{hm:KBvkrbp}
\begin{tcolorbox}[enhanced,colback=ink,colframe=ink,arc=0pt,boxrule=0pt,left=4mm,right=4mm,top=2.5mm,bottom=2.5mm]
\color{white}\sffamily{\LARGE\bfseries KBvkrbp}\hfill{\Large h\#12}\\[0.5mm]
{\small 6 Steine\quad Schwarz am Zug}\hfill{\small Nr.~231\qquad\href{https://helpman.komtera.lt/?fen=5b2/1r2p3/4k3/8/8/8/1B6/K7&moves=12}{\color{white}\bfseries lösen\,\ensuremath{\nearrow}}}
\end{tcolorbox}
\vspace{2mm}
\noindent\begin{minipage}[t]{36mm}\vspace{0pt}\sffamily\small\raggedright
\hmfigure{längste eindeutige}{{\Large h\#12}}\hmfigure{eindeutig bei dieser Länge}{{\Large 6} Stellungen}\hmfigure{längstes Matt}{{\Large h\#16.5}}\hmfigure{in der Klasse}{45,4 Milliarden Stellungen,\\ 28,2 Milliarden mit Hilfsmatt}
\end{minipage}\hfill
\begin{minipage}[t]{80mm}\vspace{0pt}\begin{tikzpicture}\begin{axis}[width=72mm,height=42mm,ybar,bar width=1.6pt,xmin=-1,xmax=34,ymin=0,ymax=5.009,axis y line*=left,axis x line*=bottom,xlabel={\scriptsize Halbzüge bis Matt},ylabel={\scriptsize Stellungen (Mrd.)},tick label style={font=\tiny},label style={font=\scriptsize},enlarge x limits=false,ylabel shift=-3pt,xlabel shift=-3pt]\addplot[fill=ink!55,draw=none] coordinates {(0,7.522e-05) (1,0.0003261) (2,0.002305) (3,0.009981) (4,0.04431) (5,0.1377) (6,0.3686) (7,0.9224) (8,1.666) (9,3.206) (10,3.376) (11,4.554) (12,2.818) (13,2.24) (14,1.468) (15,1.612) (16,1.182) (17,1.363) (18,0.7332) (19,0.8662) (20,0.3944) (21,0.4829) (22,0.2222) (23,0.2791) (24,0.1078) (25,0.134) (26,0.0146) (27,0.01757) (28,7.231e-05) (29,0.000151) (30,9e-09) (31,4.392e-05) (33,1.27e-07)};\draw[densely dotted] (axis cs:24,0) -- (axis cs:24,5.009);\end{axis}\begin{axis}[width=72mm,height=42mm,axis y line*=right,axis x line=none,ymode=log,xmin=-1,xmax=34,ylabel={\scriptsize eindeutig (log)},tick label style={font=\tiny},label style={font=\scriptsize},ylabel shift=-3pt]\addplot[mark=*,mark size=0.7pt,draw=black,thin] coordinates {(0,75217) (1,325625) (2,1701665) (3,4486515) (4,12589577) (5,22723940) (6,26280285) (7,37106847) (8,42225587) (9,38667625) (10,18165885) (11,11521238) (12,4311339) (13,2516609) (14,1422332) (15,801994) (16,277932) (17,108734) (18,31204) (19,12343) (20,4347) (21,1645) (22,739) (23,489) (24,6)};\end{axis}\end{tikzpicture}\end{minipage}
\vspace{3mm}

\noindent\begin{minipage}[t]{60mm}\vspace{0pt}\sbox\hmboard{\chessboard[setfen=5b2/1r2p3/4k3/8/8/8/1B6/K7 b - - 0 1]}\usebox\hmboard\\[1mm]
\makebox[\wd\hmboard]{\sffamily\hspace{3mm}\textbf{h\#12}\hfill(2+4)}\end{minipage}\hfill
\begin{minipage}[t]{52mm}\vspace{0pt}
\textcolor{ink}{\sffamily\textsc{Lösung}}\par\smallskip\small
\begin{tabular}{@{}r@{\enspace}l@{\enspace}l@{}}
\textbf{1.} & \symking{}f7 & \symking{}a2\\
\textbf{2.} & \symking{}g8 & \symking{}a3\\
\textbf{3.} & e5+ & \symking{}a4\\
\textbf{4.} & e4 & \symking{}a5\\
\textbf{5.} & e3 & \symking{}a6\\
\textbf{6.} & e2 & \symking{}xb7\\
\textbf{7.} & e1=\symknight{} & \symking{}c6\\
\textbf{8.} & \symknight{}f3 & \symking{}d5\\
\textbf{9.} & \symknight{}g5 & \symking{}e5\\
\textbf{10.} & \symking{}h8 & \symking{}f6\\
\textbf{11.} & \symknight{}h7+ & \symking{}f7+\\
\textbf{12.} & \symbishop{}g7 & \symbishop{}xg7\#\\
\end{tabular}\par\vspace{3mm}
\qrcode[height=12mm]{https://helpman.komtera.lt/?fen=5b2/1r2p3/4k3/8/8/8/1B6/K7&moves=12}\\{\sffamily\tiny\color{ink}scannen zum Lösen}
\end{minipage}
\vspace{2mm}

{\raggedright\noindent\textcolor{ink}{\sffamily\textsc{Themen}}\enspace \themetag{reines Matt} \themetag{Mustermatt} \themetag{Idealmatt} \themetag{Umwandlung} \themetag{Unterverwandlung} \themetag{Excelsior} \themetag{Selbstblock}\par}

\clearpage
\phantomsection\label{hm:KQBvkqp}
\begin{tcolorbox}[enhanced,colback=ink,colframe=ink,arc=0pt,boxrule=0pt,left=4mm,right=4mm,top=2.5mm,bottom=2.5mm]
\color{white}\sffamily{\LARGE\bfseries KQBvkqp}\hfill{\Large h\#12}\\[0.5mm]
{\small 6 Steine\quad Schwarz am Zug}\hfill{\small Nr.~232\qquad\href{https://helpman.komtera.lt/?fen=K7/1B6/8/4q3/8/8/Qp6/k7&moves=12}{\color{white}\bfseries lösen\,\ensuremath{\nearrow}}}
\end{tcolorbox}
\vspace{2mm}
\noindent\begin{minipage}[t]{36mm}\vspace{0pt}\sffamily\small\raggedright
\hmfigure{längste eindeutige}{{\Large h\#12}}\hmfigure{eindeutig bei dieser Länge}{{\Large 1} Stellung}\hmfigure{längstes Matt}{{\Large h\#17}}\hmfigure{in der Klasse}{45,4 Milliarden Stellungen,\\ 22,6 Milliarden mit Hilfsmatt}
\end{minipage}\hfill
\begin{minipage}[t]{80mm}\vspace{0pt}\begin{tikzpicture}\begin{axis}[width=72mm,height=42mm,ybar,bar width=1.6pt,xmin=-1,xmax=35,ymin=0,ymax=10.01,axis y line*=left,axis x line*=bottom,xlabel={\scriptsize Halbzüge bis Matt},ylabel={\scriptsize Stellungen (Mrd.)},tick label style={font=\tiny},label style={font=\scriptsize},enlarge x limits=false,ylabel shift=-3pt,xlabel shift=-3pt]\addplot[fill=ink!55,draw=none] coordinates {(0,0.03389) (1,0.241) (2,1.202) (3,4.591) (4,9.099) (5,5.697) (6,1.377) (7,0.1169) (8,0.01353) (9,2.083e-05) (10,0.01521) (11,2.818e-06) (12,0.02513) (13,6.83e-07) (14,0.0283) (15,4.38e-07) (16,0.02834) (17,4.48e-07) (18,0.02934) (19,9.07e-07) (20,0.02309) (21,9.33e-07) (22,0.01408) (23,1.667e-06) (24,0.01089) (25,7.66e-07) (26,0.009477) (27,9.33e-07) (28,0.006722) (29,3.42e-07) (30,0.004301) (31,1.6e-07) (32,0.001714) (33,9.5e-08) (34,0.0001358)};\draw[densely dotted] (axis cs:24,0) -- (axis cs:24,10.01);\end{axis}\begin{axis}[width=72mm,height=42mm,axis y line*=right,axis x line=none,ymode=log,xmin=-1,xmax=35,ylabel={\scriptsize eindeutig (log)},tick label style={font=\tiny},label style={font=\scriptsize},ylabel shift=-3pt]\addplot[mark=*,mark size=0.7pt,draw=black,thin] coordinates {(0,33893960) (1,213136210) (2,719279599) (3,994833043) (4,457960143) (5,27271890) (6,559398) (7,538) (8,262695) (9,93) (10,122812) (11,53) (12,75229) (13,16) (14,17840) (16,4700) (18,609) (20,127) (22,15) (24,1)};\end{axis}\end{tikzpicture}\end{minipage}
\vspace{3mm}

\noindent\begin{minipage}[t]{60mm}\vspace{0pt}\sbox\hmboard{\chessboard[setfen=K7/1B6/8/4q3/8/8/Qp6/k7 b - - 0 1]}\usebox\hmboard\\[1mm]
\makebox[\wd\hmboard]{\sffamily\hspace{3mm}\textbf{h\#12}\hfill(3+3)}\end{minipage}\hfill
\begin{minipage}[t]{52mm}\vspace{0pt}
\textcolor{ink}{\sffamily\textsc{Lösung}}\par\smallskip\small
\begin{tabular}{@{}r@{\enspace}l@{\enspace}l@{}}
\textbf{1.} & \symking{}xa2 & \symking{}a7\\
\textbf{2.} & \symking{}b3 & \symking{}b6\\
\textbf{3.} & \symking{}c4 & \symbishop{}a6+\\
\textbf{4.} & \symking{}d5 & \symking{}b5\\
\textbf{5.} & \symking{}e4+ & \symking{}c4\\
\textbf{6.} & \symking{}f3 & \symking{}d3\\
\textbf{7.} & \symking{}g2 & \symking{}d2\\
\textbf{8.} & \symking{}h1 & \symbishop{}e2\\
\textbf{9.} & b1=\symknight{}+ & \symking{}e1\\
\textbf{10.} & \symknight{}d2 & \symking{}f2\\
\textbf{11.} & \symknight{}f3 & \symbishop{}f1\\
\textbf{12.} & \symknight{}h2 & \symbishop{}g2\#\\
\end{tabular}\par\vspace{3mm}
\qrcode[height=12mm]{https://helpman.komtera.lt/?fen=K7/1B6/8/4q3/8/8/Qp6/k7&moves=12}\\{\sffamily\tiny\color{ink}scannen zum Lösen}
\end{minipage}
\vspace{2mm}

{\raggedright\noindent\textcolor{ink}{\sffamily\textsc{Themen}}\enspace \themetag{reines Matt} \themetag{Mustermatt} \themetag{Umwandlung} \themetag{Unterverwandlung} \themetag{Selbstblock} \themetag{Umnov}\par\smallskip{\footnotesize\textit{Hinweis:} die Partei am Zug steht im Diagramm im Schach; die Lösung beginnt mit einem Schlag.}\par}

\clearpage
\phantomsection\label{hm:KRBvkrp}
\begin{tcolorbox}[enhanced,colback=ink,colframe=ink,arc=0pt,boxrule=0pt,left=4mm,right=4mm,top=2.5mm,bottom=2.5mm]
\color{white}\sffamily{\LARGE\bfseries KRBvkrp}\hfill{\Large h\#12}\\[0.5mm]
{\small 6 Steine\quad Schwarz am Zug}\hfill{\small Nr.~233\qquad\href{https://helpman.komtera.lt/?fen=8/8/K7/1B3p2/k7/1R6/8/6r1&moves=12}{\color{white}\bfseries lösen\,\ensuremath{\nearrow}}}
\end{tcolorbox}
\vspace{2mm}
\noindent\begin{minipage}[t]{36mm}\vspace{0pt}\sffamily\small\raggedright
\hmfigure{längste eindeutige}{{\Large h\#12}}\hmfigure{eindeutig bei dieser Länge}{{\Large 3} Stellungen}\hmfigure{längstes Matt}{{\Large h\#17}}\hmfigure{in der Klasse}{45,4 Milliarden Stellungen,\\ 26,9 Milliarden mit Hilfsmatt}
\end{minipage}\hfill
\begin{minipage}[t]{80mm}\vspace{0pt}\begin{tikzpicture}\begin{axis}[width=72mm,height=42mm,ybar,bar width=1.6pt,xmin=-1,xmax=35,ymin=0,ymax=10.21,axis y line*=left,axis x line*=bottom,xlabel={\scriptsize Halbzüge bis Matt},ylabel={\scriptsize Stellungen (Mrd.)},tick label style={font=\tiny},label style={font=\scriptsize},enlarge x limits=false,ylabel shift=-3pt,xlabel shift=-3pt]\addplot[fill=ink!55,draw=none] coordinates {(0,0.01173) (1,0.05949) (2,0.2301) (3,1.108) (4,3.685) (5,7.179) (6,9.28) (7,4.305) (8,0.9501) (9,0.1031) (10,0.001864) (11,0.0001964) (12,0.001944) (13,7.508e-06) (14,0.001979) (15,2.5e-07) (16,0.001999) (17,2.38e-07) (18,0.002046) (19,3.08e-07) (20,0.001826) (21,2.81e-07) (22,0.001147) (23,3.76e-07) (24,0.0008405) (25,1.93e-07) (26,0.0006295) (27,8.1e-08) (28,0.0006232) (29,1.56e-07) (30,0.0003869) (31,7e-09) (32,0.0002917) (33,1e-09) (34,8.3e-08)};\draw[densely dotted] (axis cs:24,0) -- (axis cs:24,10.21);\end{axis}\begin{axis}[width=72mm,height=42mm,axis y line*=right,axis x line=none,ymode=log,xmin=-1,xmax=35,ylabel={\scriptsize eindeutig (log)},tick label style={font=\tiny},label style={font=\scriptsize},ylabel shift=-3pt]\addplot[mark=*,mark size=0.7pt,draw=black,thin] coordinates {(0,11725148) (1,58519102) (2,134169958) (3,458926497) (4,1194036577) (5,638788191) (6,155514895) (7,11649554) (8,87764) (9,6589) (10,8789) (11,54) (12,3208) (13,7) (14,925) (16,524) (18,37) (20,9) (22,13) (24,3)};\end{axis}\end{tikzpicture}\end{minipage}
\vspace{3mm}

\noindent\begin{minipage}[t]{60mm}\vspace{0pt}\sbox\hmboard{\chessboard[setfen=8/8/K7/1B3p2/k7/1R6/8/6r1 b - - 0 1]}\usebox\hmboard\\[1mm]
\makebox[\wd\hmboard]{\sffamily\hspace{3mm}\textbf{h\#12}\hfill(3+3)}\end{minipage}\hfill
\begin{minipage}[t]{52mm}\vspace{0pt}
\textcolor{ink}{\sffamily\textsc{Lösung}}\par\smallskip\small
\begin{tabular}{@{}r@{\enspace}l@{\enspace}l@{}}
\textbf{1.} & \symking{}xb3 & \symking{}b6\\
\textbf{2.} & f4 & \symking{}c5\\
\textbf{3.} & f3 & \symking{}d4\\
\textbf{4.} & f2 & \symking{}e3\\
\textbf{5.} & f1=\symknight{}+ & \symking{}f2\\
\textbf{6.} & \symking{}c3 & \symking{}xg1\\
\textbf{7.} & \symking{}d4 & \symking{}f2\\
\textbf{8.} & \symking{}e4 & \symking{}e1\\
\textbf{9.} & \symking{}f3 & \symbishop{}d7\\
\textbf{10.} & \symking{}g2 & \symbishop{}h3+\\
\textbf{11.} & \symking{}h1 & \symking{}f2\\
\textbf{12.} & \symknight{}h2 & \symbishop{}g2\#\\
\end{tabular}\par\vspace{3mm}
\qrcode[height=12mm]{https://helpman.komtera.lt/?fen=8/8/K7/1B3p2/k7/1R6/8/6r1&moves=12}\\{\sffamily\tiny\color{ink}scannen zum Lösen}
\end{minipage}
\vspace{2mm}

{\raggedright\noindent\textcolor{ink}{\sffamily\textsc{Themen}}\enspace \themetag{reines Matt} \themetag{Mustermatt} \themetag{Idealmatt} \themetag{Umwandlung} \themetag{Unterverwandlung} \themetag{Switchback} \themetag{Selbstblock} \themetag{Umnov}\par\smallskip{\footnotesize\textit{Hinweis:} die Partei am Zug steht im Diagramm im Schach; die Lösung beginnt mit einem Schlag.}\par}

\clearpage
\phantomsection\label{hm:KBvkbpp}
\begin{tcolorbox}[enhanced,colback=ink,colframe=ink,arc=0pt,boxrule=0pt,left=4mm,right=4mm,top=2.5mm,bottom=2.5mm]
\color{white}\sffamily{\LARGE\bfseries KBvkbpp}\hfill{\Large h\#11.5}\\[0.5mm]
{\small 6 Steine\quad Weiß am Zug}\hfill{\small Nr.~234\qquad\href{https://helpman.komtera.lt/?fen=b7/8/2kp1p2/8/8/8/3K4/5B2&moves=11.5}{\color{white}\bfseries lösen\,\ensuremath{\nearrow}}}
\end{tcolorbox}
\vspace{2mm}
\noindent\begin{minipage}[t]{36mm}\vspace{0pt}\sffamily\small\raggedright
\hmfigure{längste eindeutige}{{\Large h\#11.5}}\hmfigure{eindeutig bei dieser Länge}{{\Large 34} Stellungen}\hmfigure{längstes Matt}{{\Large h\#16}}\hmfigure{in der Klasse}{34,1 Milliarden Stellungen,\\ 11,5 Milliarden mit Hilfsmatt}
\end{minipage}\hfill
\begin{minipage}[t]{80mm}\vspace{0pt}\begin{tikzpicture}\begin{axis}[width=72mm,height=42mm,ybar,bar width=1.6pt,xmin=-1,xmax=33,ymin=0,ymax=1.667,axis y line*=left,axis x line*=bottom,xlabel={\scriptsize Halbzüge bis Matt},ylabel={\scriptsize Stellungen (Mrd.)},tick label style={font=\tiny},label style={font=\scriptsize},enlarge x limits=false,ylabel shift=-3pt,xlabel shift=-3pt]\addplot[fill=ink!55,draw=none] coordinates {(0,0.0001085) (1,0.0004293) (2,0.002343) (3,0.007235) (4,0.03005) (5,0.07127) (6,0.1701) (7,0.3207) (8,0.5696) (9,0.8732) (10,1.207) (11,1.516) (12,1.443) (13,1.133) (14,0.8633) (15,0.722) (16,0.5582) (17,0.5257) (18,0.3623) (19,0.3585) (20,0.2133) (21,0.2198) (22,0.1044) (23,0.1098) (24,0.04136) (25,0.04462) (26,0.01369) (27,0.01508) (28,0.002443) (29,0.002693) (30,0.0004625) (31,0.0005121) (32,2.79e-07)};\draw[densely dotted] (axis cs:23,0) -- (axis cs:23,1.667);\end{axis}\begin{axis}[width=72mm,height=42mm,axis y line*=right,axis x line=none,ymode=log,xmin=-1,xmax=33,ylabel={\scriptsize eindeutig (log)},tick label style={font=\tiny},label style={font=\scriptsize},ylabel shift=-3pt]\addplot[mark=*,mark size=0.7pt,draw=black,thin] coordinates {(0,108537) (1,426253) (2,1964623) (3,2505749) (4,6294771) (5,9359536) (6,8059680) (7,7959709) (8,6227967) (9,6100899) (10,4287930) (11,2959765) (12,1584909) (13,578018) (14,348942) (15,101267) (16,48855) (17,10555) (18,5005) (19,1087) (20,536) (21,148) (22,50) (23,34)};\end{axis}\end{tikzpicture}\end{minipage}
\vspace{3mm}

\noindent\begin{minipage}[t]{60mm}\vspace{0pt}\sbox\hmboard{\chessboard[setfen=b7/8/2kp1p2/8/8/8/3K4/5B2 w - - 0 1]}\usebox\hmboard\\[1mm]
\makebox[\wd\hmboard]{\sffamily\hspace{3mm}\textbf{h\#11.5}\hfill(2+4)}\end{minipage}\hfill
\begin{minipage}[t]{52mm}\vspace{0pt}
\textcolor{ink}{\sffamily\textsc{Lösung}}\par\smallskip\small
\begin{tabular}{@{}r@{\enspace}l@{\enspace}l@{}}
\textbf{1.} & \ldots & \symking{}c3\\
\textbf{2.} & d5 & \symking{}b4\\
\textbf{3.} & d4 & \symking{}a5\\
\textbf{4.} & d3 & \symking{}a6\\
\textbf{5.} & d2 & \symking{}a7\\
\textbf{6.} & d1=\symknight{} & \symking{}xa8\\
\textbf{7.} & \symknight{}c3 & \symking{}b8\\
\textbf{8.} & \symknight{}b5 & \symking{}c8\\
\textbf{9.} & \symknight{}a7+ & \symking{}d8\\
\textbf{10.} & \symking{}b7 & \symbishop{}a6+\\
\textbf{11.} & \symking{}a8 & \symking{}c7\\
\textbf{12.} & f5 & \symbishop{}b7\#\\
\end{tabular}\par\vspace{3mm}
\qrcode[height=12mm]{https://helpman.komtera.lt/?fen=b7/8/2kp1p2/8/8/8/3K4/5B2&moves=11.5}\\{\sffamily\tiny\color{ink}scannen zum Lösen}
\end{minipage}
\vspace{2mm}

{\raggedright\noindent\textcolor{ink}{\sffamily\textsc{Themen}}\enspace \themetag{reines Matt} \themetag{Mustermatt} \themetag{Umwandlung} \themetag{Unterverwandlung} \themetag{Selbstblock} \themetag{Kniest}\par}

\clearpage
\phantomsection\label{hm:KBvkqnp}
\begin{tcolorbox}[enhanced,colback=ink,colframe=ink,arc=0pt,boxrule=0pt,left=4mm,right=4mm,top=2.5mm,bottom=2.5mm]
\color{white}\sffamily{\LARGE\bfseries KBvkqnp}\hfill{\Large h\#11.5}\\[0.5mm]
{\small 6 Steine\quad Weiß am Zug}\hfill{\small Nr.~235\qquad\href{https://helpman.komtera.lt/?fen=1k6/4n3/3K4/3q4/8/8/6p1/2B5&moves=11.5}{\color{white}\bfseries lösen\,\ensuremath{\nearrow}}}
\end{tcolorbox}
\vspace{2mm}
\noindent\begin{minipage}[t]{36mm}\vspace{0pt}\sffamily\small\raggedright
\hmfigure{längste eindeutige}{{\Large h\#11.5}}\hmfigure{eindeutig bei dieser Länge}{{\Large 2} Stellungen}\hmfigure{längstes Matt}{{\Large h\#16.5}}\hmfigure{in der Klasse}{45,4 Milliarden Stellungen,\\ 26,8 Milliarden mit Hilfsmatt}
\end{minipage}\hfill
\begin{minipage}[t]{80mm}\vspace{0pt}\begin{tikzpicture}\begin{axis}[width=72mm,height=42mm,ybar,bar width=1.6pt,xmin=-1,xmax=34,ymin=0,ymax=5.135,axis y line*=left,axis x line*=bottom,xlabel={\scriptsize Halbzüge bis Matt},ylabel={\scriptsize Stellungen (Mrd.)},tick label style={font=\tiny},label style={font=\scriptsize},enlarge x limits=false,ylabel shift=-3pt,xlabel shift=-3pt]\addplot[fill=ink!55,draw=none] coordinates {(0,2.083e-05) (1,7.674e-05) (2,0.0003642) (3,0.001872) (4,0.006527) (5,0.02813) (6,0.07446) (7,0.2386) (8,0.4605) (9,1.11) (10,1.542) (11,2.992) (12,3.097) (13,4.668) (14,3.409) (15,4.071) (16,1.8) (17,1.932) (18,0.4996) (19,0.6216) (20,0.09045) (21,0.1241) (22,0.006126) (23,0.009211) (24,2.776e-05) (25,0.000561) (26,3.6e-08) (27,0.000373) (28,9e-09) (29,0.0001859) (31,7.675e-05) (33,9.7e-06)};\draw[densely dotted] (axis cs:23,0) -- (axis cs:23,5.135);\end{axis}\begin{axis}[width=72mm,height=42mm,axis y line*=right,axis x line=none,ymode=log,xmin=-1,xmax=34,ylabel={\scriptsize eindeutig (log)},tick label style={font=\tiny},label style={font=\scriptsize},ylabel shift=-3pt]\addplot[mark=*,mark size=0.7pt,draw=black,thin] coordinates {(0,20830) (1,76744) (2,182532) (3,780627) (4,2041438) (5,5669320) (6,6500881) (7,12422199) (8,11566208) (9,17549668) (10,13282815) (11,15100126) (12,8005634) (13,7304716) (14,2033272) (15,1174496) (16,149624) (17,77568) (18,7695) (19,3050) (20,63) (21,34) (23,2)};\end{axis}\end{tikzpicture}\end{minipage}
\vspace{3mm}

\noindent\begin{minipage}[t]{60mm}\vspace{0pt}\sbox\hmboard{\chessboard[setfen=1k6/4n3/3K4/3q4/8/8/6p1/2B5 w - - 0 1]}\usebox\hmboard\\[1mm]
\makebox[\wd\hmboard]{\sffamily\hspace{3mm}\textbf{h\#11.5}\hfill(2+4)}\end{minipage}\hfill
\begin{minipage}[t]{52mm}\vspace{0pt}
\textcolor{ink}{\sffamily\textsc{Lösung}}\par\smallskip\small
\begin{tabular}{@{}r@{\enspace}l@{\enspace}l@{}}
\textbf{1.} & \ldots & \symking{}xe7\\
\textbf{2.} & \symking{}c7 & \symking{}f6\\
\textbf{3.} & \symking{}d6 & \symking{}g6\\
\textbf{4.} & \symking{}e5 & \symking{}g5\\
\textbf{5.} & \symking{}d4+ & \symking{}f4\\
\textbf{6.} & \symking{}c3 & \symbishop{}d2+\\
\textbf{7.} & \symking{}b2 & \symking{}e3\\
\textbf{8.} & \symking{}a1 & \symking{}e2\\
\textbf{9.} & g1=\symknight{}+ & \symking{}d1\\
\textbf{10.} & \symknight{}e2 & \symking{}c2\\
\textbf{11.} & \symknight{}c3 & \symbishop{}c1\\
\textbf{12.} & \symknight{}a2 & \symbishop{}b2\#\\
\end{tabular}\par\vspace{3mm}
\qrcode[height=12mm]{https://helpman.komtera.lt/?fen=1k6/4n3/3K4/3q4/8/8/6p1/2B5&moves=11.5}\\{\sffamily\tiny\color{ink}scannen zum Lösen}
\end{minipage}
\vspace{2mm}

{\raggedright\noindent\textcolor{ink}{\sffamily\textsc{Themen}}\enspace \themetag{reines Matt} \themetag{Mustermatt} \themetag{Umwandlung} \themetag{Unterverwandlung} \themetag{Switchback} \themetag{Selbstblock} \themetag{Phönix} \themetag{Umnov}\par\smallskip{\footnotesize\textit{Hinweis:} die Partei am Zug steht im Diagramm im Schach; die Lösung beginnt mit einem Schlag.}\par}

\clearpage
\phantomsection\label{hm:KBvkqpp}
\begin{tcolorbox}[enhanced,colback=ink,colframe=ink,arc=0pt,boxrule=0pt,left=4mm,right=4mm,top=2.5mm,bottom=2.5mm]
\color{white}\sffamily{\LARGE\bfseries KBvkqpp}\hfill{\Large h\#11.5}\\[0.5mm]
{\small 6 Steine\quad Weiß am Zug}\hfill{\small Nr.~236\qquad\href{https://helpman.komtera.lt/?fen=3K4/3p4/B7/1p1kq3/8/8/8/8&moves=11.5}{\color{white}\bfseries lösen\,\ensuremath{\nearrow}}}
\end{tcolorbox}
\vspace{2mm}
\noindent\begin{minipage}[t]{36mm}\vspace{0pt}\sffamily\small\raggedright
\hmfigure{längste eindeutige}{{\Large h\#11.5}}\hmfigure{eindeutig bei dieser Länge}{{\Large 3} Stellungen}\hmfigure{längstes Matt}{{\Large h\#16.5}}\hmfigure{in der Klasse}{34,1 Milliarden Stellungen,\\ 10,2 Milliarden mit Hilfsmatt}
\end{minipage}\hfill
\begin{minipage}[t]{80mm}\vspace{0pt}\begin{tikzpicture}\begin{axis}[width=72mm,height=42mm,ybar,bar width=1.6pt,xmin=-1,xmax=34,ymin=0,ymax=1.52,axis y line*=left,axis x line*=bottom,xlabel={\scriptsize Halbzüge bis Matt},ylabel={\scriptsize Stellungen (Mrd.)},tick label style={font=\tiny},label style={font=\scriptsize},enlarge x limits=false,ylabel shift=-3pt,xlabel shift=-3pt]\addplot[fill=ink!55,draw=none] coordinates {(0,1.555e-05) (1,6.739e-05) (2,0.0003249) (3,0.001505) (4,0.004948) (5,0.01855) (6,0.04714) (7,0.1336) (8,0.2509) (9,0.5337) (10,0.727) (11,1.216) (12,1.157) (13,1.382) (14,0.9182) (15,1.071) (16,0.6318) (17,0.7885) (18,0.3618) (19,0.4539) (20,0.1407) (21,0.1805) (22,0.05555) (23,0.07546) (24,0.02665) (25,0.03707) (26,0.009626) (27,0.0138) (28,0.00192) (29,0.0028) (30,0.0002277) (31,0.0003833) (32,8.403e-06) (33,1.607e-05)};\draw[densely dotted] (axis cs:23,0) -- (axis cs:23,1.52);\end{axis}\begin{axis}[width=72mm,height=42mm,axis y line*=right,axis x line=none,ymode=log,xmin=-1,xmax=34,ylabel={\scriptsize eindeutig (log)},tick label style={font=\tiny},label style={font=\scriptsize},ylabel shift=-3pt]\addplot[mark=*,mark size=0.7pt,draw=black,thin] coordinates {(0,15546) (1,67297) (2,168452) (3,626630) (4,1302802) (5,3009840) (6,3674035) (7,6167472) (8,6052671) (9,7101915) (10,5041303) (11,4792537) (12,2315994) (13,1465982) (14,609603) (15,358578) (16,108127) (17,51343) (18,11699) (19,4949) (20,711) (21,179) (22,16) (23,3)};\end{axis}\end{tikzpicture}\end{minipage}
\vspace{3mm}

\noindent\begin{minipage}[t]{60mm}\vspace{0pt}\sbox\hmboard{\chessboard[setfen=3K4/3p4/B7/1p1kq3/8/8/8/8 w - - 0 1]}\usebox\hmboard\\[1mm]
\makebox[\wd\hmboard]{\sffamily\hspace{3mm}\textbf{h\#11.5}\hfill(2+4)}\end{minipage}\hfill
\begin{minipage}[t]{52mm}\vspace{0pt}
\textcolor{ink}{\sffamily\textsc{Lösung}}\par\smallskip\small
\begin{tabular}{@{}r@{\enspace}l@{\enspace}l@{}}
\textbf{1.} & \ldots & \symking{}c8\\
\textbf{2.} & b4 & \symking{}b7\\
\textbf{3.} & b3 & \symking{}b6\\
\textbf{4.} & b2 & \symking{}b5\\
\textbf{5.} & \symking{}e4+ & \symking{}c4\\
\textbf{6.} & \symking{}f3 & \symking{}d3\\
\textbf{7.} & \symking{}g2 & \symking{}d2\\
\textbf{8.} & \symking{}h1 & \symbishop{}e2\\
\textbf{9.} & b1=\symknight{}+ & \symking{}e1\\
\textbf{10.} & \symknight{}d2 & \symking{}f2\\
\textbf{11.} & \symknight{}f3 & \symbishop{}f1\\
\textbf{12.} & \symknight{}h2 & \symbishop{}g2\#\\
\end{tabular}\par\vspace{3mm}
\qrcode[height=12mm]{https://helpman.komtera.lt/?fen=3K4/3p4/B7/1p1kq3/8/8/8/8&moves=11.5}\\{\sffamily\tiny\color{ink}scannen zum Lösen}
\end{minipage}
\vspace{2mm}

{\raggedright\noindent\textcolor{ink}{\sffamily\textsc{Themen}}\enspace \themetag{reines Matt} \themetag{Mustermatt} \themetag{Umwandlung} \themetag{Unterverwandlung} \themetag{Selbstblock} \themetag{ohne Schlag} \themetag{Umnov}\par}

\clearpage
\phantomsection\label{hm:KBvkrnp}
\begin{tcolorbox}[enhanced,colback=ink,colframe=ink,arc=0pt,boxrule=0pt,left=4mm,right=4mm,top=2.5mm,bottom=2.5mm]
\color{white}\sffamily{\LARGE\bfseries KBvkrnp}\hfill{\Large h\#11.5}\\[0.5mm]
{\small 6 Steine\quad Weiß am Zug}\hfill{\small Nr.~237\qquad\href{https://helpman.komtera.lt/?fen=8/8/3p4/8/6k1/1n6/B6r/1K6&moves=11.5}{\color{white}\bfseries lösen\,\ensuremath{\nearrow}}}
\end{tcolorbox}
\vspace{2mm}
\noindent\begin{minipage}[t]{36mm}\vspace{0pt}\sffamily\small\raggedright
\hmfigure{längste eindeutige}{{\Large h\#11.5}}\hmfigure{eindeutig bei dieser Länge}{{\Large 10} Stellungen}\hmfigure{längstes Matt}{{\Large h\#16.5}}\hmfigure{in der Klasse}{45,4 Milliarden Stellungen,\\ 28,9 Milliarden mit Hilfsmatt}
\end{minipage}\hfill
\begin{minipage}[t]{80mm}\vspace{0pt}\begin{tikzpicture}\begin{axis}[width=72mm,height=42mm,ybar,bar width=1.6pt,xmin=-1,xmax=34,ymin=0,ymax=5.317,axis y line*=left,axis x line*=bottom,xlabel={\scriptsize Halbzüge bis Matt},ylabel={\scriptsize Stellungen (Mrd.)},tick label style={font=\tiny},label style={font=\scriptsize},enlarge x limits=false,ylabel shift=-3pt,xlabel shift=-3pt]\addplot[fill=ink!55,draw=none] coordinates {(0,3.028e-05) (1,0.0001105) (2,0.0004211) (3,0.002051) (4,0.007554) (5,0.02988) (6,0.08614) (7,0.2451) (8,0.5376) (9,1.122) (10,1.824) (11,3.05) (12,3.737) (13,4.834) (14,4.098) (15,4.074) (16,2.094) (17,1.807) (18,0.5519) (19,0.56) (20,0.09031) (21,0.1016) (22,0.003711) (23,0.004147) (24,1.409e-05) (25,4.019e-05) (27,2.762e-05) (29,8.738e-06) (31,2.745e-06) (33,6.3e-08)};\draw[densely dotted] (axis cs:23,0) -- (axis cs:23,5.317);\end{axis}\begin{axis}[width=72mm,height=42mm,axis y line*=right,axis x line=none,ymode=log,xmin=-1,xmax=34,ylabel={\scriptsize eindeutig (log)},tick label style={font=\tiny},label style={font=\scriptsize},ylabel shift=-3pt]\addplot[mark=*,mark size=0.7pt,draw=black,thin] coordinates {(0,30284) (1,110492) (2,212545) (3,812852) (4,3053950) (5,6694410) (6,8727272) (7,12911783) (8,15393580) (9,17023045) (10,16164685) (11,13829251) (12,9401659) (13,6067343) (14,2146653) (15,882204) (16,152632) (17,59736) (18,7233) (19,3128) (20,122) (21,16) (23,10)};\end{axis}\end{tikzpicture}\end{minipage}
\vspace{3mm}

\noindent\begin{minipage}[t]{60mm}\vspace{0pt}\sbox\hmboard{\chessboard[setfen=8/8/3p4/8/6k1/1n6/B6r/1K6 w - - 0 1]}\usebox\hmboard\\[1mm]
\makebox[\wd\hmboard]{\sffamily\hspace{3mm}\textbf{h\#11.5}\hfill(2+4)}\end{minipage}\hfill
\begin{minipage}[t]{52mm}\vspace{0pt}
\textcolor{ink}{\sffamily\textsc{Lösung}}\par\smallskip\small
\begin{tabular}{@{}r@{\enspace}l@{\enspace}l@{}}
\textbf{1.} & \ldots & \symbishop{}xb3\\
\textbf{2.} & d5 & \symking{}c1\\
\textbf{3.} & d4 & \symking{}d1\\
\textbf{4.} & d3 & \symking{}e1\\
\textbf{5.} & d2+ & \symking{}f1\\
\textbf{6.} & d1=\symknight{} & \symking{}g1\\
\textbf{7.} & \symknight{}e3 & \symking{}xh2\\
\textbf{8.} & \symknight{}f1+ & \symking{}g1\\
\textbf{9.} & \symking{}h3 & \symking{}f2\\
\textbf{10.} & \symking{}h2 & \symbishop{}e6\\
\textbf{11.} & \symking{}h1 & \symbishop{}h3\\
\textbf{12.} & \symknight{}h2 & \symbishop{}g2\#\\
\end{tabular}\par\vspace{3mm}
\qrcode[height=12mm]{https://helpman.komtera.lt/?fen=8/8/3p4/8/6k1/1n6/B6r/1K6&moves=11.5}\\{\sffamily\tiny\color{ink}scannen zum Lösen}
\end{minipage}
\vspace{2mm}

{\raggedright\noindent\textcolor{ink}{\sffamily\textsc{Themen}}\enspace \themetag{reines Matt} \themetag{Mustermatt} \themetag{Idealmatt} \themetag{Umwandlung} \themetag{Unterverwandlung} \themetag{Switchback} \themetag{Selbstblock} \themetag{Phönix}\par\smallskip{\footnotesize\textit{Hinweis:} die Lösung beginnt mit einem Schlag.}\par}

\clearpage
\phantomsection\label{hm:KPvkppp}
\begin{tcolorbox}[enhanced,colback=ink,colframe=ink,arc=0pt,boxrule=0pt,left=4mm,right=4mm,top=2.5mm,bottom=2.5mm]
\color{white}\sffamily{\LARGE\bfseries KPvkppp}\hfill{\Large h\#11.5}\\[0.5mm]
{\small 6 Steine\quad Weiß am Zug}\hfill{\small Nr.~238\qquad\href{https://helpman.komtera.lt/?fen=8/8/4p3/1p3k2/8/1p6/1P6/K7&moves=11.5}{\color{white}\bfseries lösen\,\ensuremath{\nearrow}}}
\end{tcolorbox}
\vspace{2mm}
\noindent\begin{minipage}[t]{36mm}\vspace{0pt}\sffamily\small\raggedright
\hmfigure{längste eindeutige}{{\Large h\#11.5}}\hmfigure{eindeutig bei dieser Länge}{{\Large 4} Stellungen}\hmfigure{längstes Matt}{{\Large h\#16.5}}\hmfigure{in der Klasse}{19,2 Milliarden Stellungen,\\ 2,34 Milliarden mit Hilfsmatt}
\end{minipage}\hfill
\begin{minipage}[t]{80mm}\vspace{0pt}\begin{tikzpicture}\begin{axis}[width=72mm,height=42mm,ybar,bar width=1.6pt,xmin=-1,xmax=34,ymin=0,ymax=340.8,axis y line*=left,axis x line*=bottom,xlabel={\scriptsize Halbzüge bis Matt},ylabel={\scriptsize Stellungen (Mio.)},tick label style={font=\tiny},label style={font=\scriptsize},enlarge x limits=false,ylabel shift=-3pt,xlabel shift=-3pt]\addplot[fill=ink!55,draw=none] coordinates {(0,0.002201) (1,0.6556) (2,1.43) (3,8.727) (4,20.36) (5,79.29) (6,130.7) (7,222.9) (8,246.3) (9,294) (10,262.1) (11,309.8) (12,286.3) (13,186.2) (14,122.8) (15,52.38) (16,37.2) (17,28.21) (18,17.67) (19,13.34) (20,7.31) (21,4.659) (22,2.304) (23,1.222) (24,0.4144) (25,0.1885) (26,0.04957) (27,0.02617) (28,0.0186) (29,0.01519) (30,0.01539) (31,0.00645) (32,0.006277) (33,0.000128)};\draw[densely dotted] (axis cs:23,0) -- (axis cs:23,340.8);\end{axis}\begin{axis}[width=72mm,height=42mm,axis y line*=right,axis x line=none,ymode=log,xmin=-1,xmax=34,ylabel={\scriptsize eindeutig (log)},tick label style={font=\tiny},label style={font=\scriptsize},ylabel shift=-3pt]\addplot[mark=*,mark size=0.7pt,draw=black,thin] coordinates {(0,2201) (1,154685) (2,315085) (3,2138936) (4,3696199) (5,7337909) (6,5100053) (7,4119612) (8,2426249) (9,2577212) (10,1509152) (11,1652844) (12,794071) (13,125826) (14,67456) (15,17036) (16,9856) (17,2288) (18,1553) (19,295) (20,90) (21,43) (22,18) (23,4)};\end{axis}\end{tikzpicture}\end{minipage}
\vspace{3mm}

\noindent\begin{minipage}[t]{60mm}\vspace{0pt}\sbox\hmboard{\chessboard[setfen=8/8/4p3/1p3k2/8/1p6/1P6/K7 w - - 0 1]}\usebox\hmboard\\[1mm]
\makebox[\wd\hmboard]{\sffamily\hspace{3mm}\textbf{h\#11.5}\hfill(2+4)}\end{minipage}\hfill
\begin{minipage}[t]{52mm}\vspace{0pt}
\textcolor{ink}{\sffamily\textsc{Lösung}}\par\smallskip\small
\begin{tabular}{@{}r@{\enspace}l@{\enspace}l@{}}
\textbf{1.} & \ldots & \symking{}b1\\
\textbf{2.} & e5 & \symking{}c1\\
\textbf{3.} & e4 & \symking{}d2\\
\textbf{4.} & e3+ & \symking{}c3\\
\textbf{5.} & e2 & \symking{}xb3\\
\textbf{6.} & e1=\symbishop{} & \symking{}a3\\
\textbf{7.} & \symbishop{}a5 & b4\\
\textbf{8.} & \symking{}e4 & bxa5\\
\textbf{9.} & \symking{}d3 & a6\\
\textbf{10.} & \symking{}c2 & a7\\
\textbf{11.} & \symking{}b1 & a8=\symqueen{}\\
\textbf{12.} & \symking{}a1 & \symqueen{}h1\#\\
\end{tabular}\par\vspace{3mm}
\qrcode[height=12mm]{https://helpman.komtera.lt/?fen=8/8/4p3/1p3k2/8/1p6/1P6/K7&moves=11.5}\\{\sffamily\tiny\color{ink}scannen zum Lösen}
\end{minipage}
\vspace{2mm}

{\raggedright\noindent\textcolor{ink}{\sffamily\textsc{Themen}}\enspace \themetag{reines Matt} \themetag{Mustermatt} \themetag{Spiegelmatt} \themetag{Umwandlung} \themetag{Unterverwandlung} \themetag{Excelsior}\par}

\clearpage
\phantomsection\label{hm:KBBvkrp}
\begin{tcolorbox}[enhanced,colback=ink,colframe=ink,arc=0pt,boxrule=0pt,left=4mm,right=4mm,top=2.5mm,bottom=2.5mm]
\color{white}\sffamily{\LARGE\bfseries KBBvkrp}\hfill{\Large h\#11}\\[0.5mm]
{\small 6 Steine\quad Schwarz am Zug}\hfill{\small Nr.~239\qquad\href{https://helpman.komtera.lt/?fen=8/8/8/1K6/2B2p2/k7/1B6/6r1&moves=11}{\color{white}\bfseries lösen\,\ensuremath{\nearrow}}}
\end{tcolorbox}
\vspace{2mm}
\noindent\begin{minipage}[t]{36mm}\vspace{0pt}\sffamily\small\raggedright
\hmfigure{längste eindeutige}{{\Large h\#11}}\hmfigure{eindeutig bei dieser Länge}{{\Large 1} Stellung}\hmfigure{längstes Matt}{{\Large h\#17}}\hmfigure{in der Klasse}{45,4 Milliarden Stellungen,\\ 14 Milliarden mit Hilfsmatt}
\end{minipage}\hfill
\begin{minipage}[t]{80mm}\vspace{0pt}\begin{tikzpicture}\begin{axis}[width=72mm,height=42mm,ybar,bar width=1.6pt,xmin=-1,xmax=35,ymin=0,ymax=1.728,axis y line*=left,axis x line*=bottom,xlabel={\scriptsize Halbzüge bis Matt},ylabel={\scriptsize Stellungen (Mrd.)},tick label style={font=\tiny},label style={font=\scriptsize},enlarge x limits=false,ylabel shift=-3pt,xlabel shift=-3pt]\addplot[fill=ink!55,draw=none] coordinates {(0,0.000781) (1,0.00351) (2,0.01529) (3,0.1013) (4,0.3275) (5,0.8998) (6,1.565) (7,1.571) (8,1.535) (9,1.074) (10,0.8776) (11,0.7666) (12,0.9412) (13,0.8969) (14,0.9094) (15,0.7838) (16,0.5799) (17,0.4773) (18,0.2666) (19,0.2356) (20,0.06549) (21,0.05588) (22,0.002411) (23,0.0009723) (24,8.371e-05) (25,0.0001603) (26,6.421e-05) (27,1.168e-05) (28,3.564e-05) (29,3.711e-06) (30,2.449e-05) (31,4e-09) (32,1.139e-05) (33,9e-09) (34,1e-08)};\draw[densely dotted] (axis cs:22,0) -- (axis cs:22,1.728);\end{axis}\begin{axis}[width=72mm,height=42mm,axis y line*=right,axis x line=none,ymode=log,xmin=-1,xmax=35,ylabel={\scriptsize eindeutig (log)},tick label style={font=\tiny},label style={font=\scriptsize},ylabel shift=-3pt]\addplot[mark=*,mark size=0.7pt,draw=black,thin] coordinates {(0,781006) (1,3477045) (2,9465046) (3,45814081) (4,80215913) (5,75594853) (6,47809966) (7,17742245) (8,9421749) (9,4381874) (10,3447960) (11,2137281) (12,2116130) (13,965539) (14,676473) (15,190676) (16,71554) (17,12065) (18,2506) (19,352) (20,5) (21,1) (22,1)};\end{axis}\end{tikzpicture}\end{minipage}
\vspace{3mm}

\noindent\begin{minipage}[t]{60mm}\vspace{0pt}\sbox\hmboard{\chessboard[setfen=8/8/8/1K6/2B2p2/k7/1B6/6r1 b - - 0 1]}\usebox\hmboard\\[1mm]
\makebox[\wd\hmboard]{\sffamily\hspace{3mm}\textbf{h\#11}\hfill(3+3)}\end{minipage}\hfill
\begin{minipage}[t]{52mm}\vspace{0pt}
\textcolor{ink}{\sffamily\textsc{Lösung}}\par\smallskip\small
\begin{tabular}{@{}r@{\enspace}l@{\enspace}l@{}}
\textbf{1.} & \symking{}xb2 & \symking{}c5\\
\textbf{2.} & f3 & \symking{}d4\\
\textbf{3.} & f2 & \symking{}e3\\
\textbf{4.} & f1=\symknight{}+ & \symking{}f2\\
\textbf{5.} & \symking{}c3 & \symking{}xg1\\
\textbf{6.} & \symking{}d4 & \symking{}f2\\
\textbf{7.} & \symking{}e4 & \symking{}e1\\
\textbf{8.} & \symking{}f3 & \symbishop{}e6\\
\textbf{9.} & \symking{}g2 & \symbishop{}h3+\\
\textbf{10.} & \symking{}h1 & \symking{}f2\\
\textbf{11.} & \symknight{}h2 & \symbishop{}g2\#\\
\end{tabular}\par\vspace{3mm}
\qrcode[height=12mm]{https://helpman.komtera.lt/?fen=8/8/8/1K6/2B2p2/k7/1B6/6r1&moves=11}\\{\sffamily\tiny\color{ink}scannen zum Lösen}
\end{minipage}
\vspace{2mm}

{\raggedright\noindent\textcolor{ink}{\sffamily\textsc{Themen}}\enspace \themetag{reines Matt} \themetag{Mustermatt} \themetag{Idealmatt} \themetag{Umwandlung} \themetag{Unterverwandlung} \themetag{Switchback} \themetag{Selbstblock} \themetag{Umnov}\par\smallskip{\footnotesize\textit{Hinweis:} die Partei am Zug steht im Diagramm im Schach; die Lösung beginnt mit einem Schlag.}\par}

\clearpage
\phantomsection\label{hm:KBPPvkp}
\begin{tcolorbox}[enhanced,colback=ink,colframe=ink,arc=0pt,boxrule=0pt,left=4mm,right=4mm,top=2.5mm,bottom=2.5mm]
\color{white}\sffamily{\LARGE\bfseries KBPPvkp}\hfill{\Large h\#11}\\[0.5mm]
{\small 6 Steine\quad Schwarz am Zug}\hfill{\small Nr.~240\qquad\href{https://helpman.komtera.lt/?fen=5B2/8/8/7p/7P/kP6/8/1K6&moves=11}{\color{white}\bfseries lösen\,\ensuremath{\nearrow}}}
\end{tcolorbox}
\vspace{2mm}
\noindent\begin{minipage}[t]{36mm}\vspace{0pt}\sffamily\small\raggedright
\hmfigure{längste eindeutige}{{\Large h\#11}}\hmfigure{eindeutig bei dieser Länge}{{\Large 1} Stellung}\hmfigure{längstes Matt}{{\Large h\#14}}\hmfigure{in der Klasse}{25,6 Milliarden Stellungen,\\ 9,03 Milliarden mit Hilfsmatt}
\end{minipage}\hfill
\begin{minipage}[t]{80mm}\vspace{0pt}\begin{tikzpicture}\begin{axis}[width=72mm,height=42mm,ybar,bar width=1.6pt,xmin=-1,xmax=29,ymin=0,ymax=1.781,axis y line*=left,axis x line*=bottom,xlabel={\scriptsize Halbzüge bis Matt},ylabel={\scriptsize Stellungen (Mrd.)},tick label style={font=\tiny},label style={font=\scriptsize},enlarge x limits=false,ylabel shift=-3pt,xlabel shift=-3pt]\addplot[fill=ink!55,draw=none] coordinates {(0,0.0003445) (1,0.007462) (2,0.02077) (3,0.1533) (4,0.4196) (5,1.217) (6,1.619) (7,1.267) (8,1.358) (9,0.895) (10,0.9389) (11,0.5255) (12,0.4773) (13,0.08014) (14,0.04727) (15,0.002816) (16,0.002249) (17,0.0007431) (18,0.0007409) (19,0.0003536) (20,0.0003132) (21,9.603e-05) (22,7.959e-05) (23,7.52e-06) (24,7.966e-06) (25,2.6e-08) (26,2.85e-07) (28,2.24e-07)};\draw[densely dotted] (axis cs:22,0) -- (axis cs:22,1.781);\end{axis}\begin{axis}[width=72mm,height=42mm,axis y line*=right,axis x line=none,ymode=log,xmin=-1,xmax=29,ylabel={\scriptsize eindeutig (log)},tick label style={font=\tiny},label style={font=\scriptsize},ylabel shift=-3pt]\addplot[mark=*,mark size=0.7pt,draw=black,thin] coordinates {(0,344549) (1,2720919) (2,8212145) (3,66702110) (4,93383106) (5,90034473) (6,55336463) (7,32528795) (8,21062436) (9,10813406) (10,6497016) (11,2205641) (12,949067) (13,18612) (14,9456) (15,2171) (16,1231) (17,408) (18,286) (19,38) (20,8) (22,1)};\end{axis}\end{tikzpicture}\end{minipage}
\vspace{3mm}

\noindent\begin{minipage}[t]{60mm}\vspace{0pt}\sbox\hmboard{\chessboard[setfen=5B2/8/8/7p/7P/kP6/8/1K6 b - - 0 1]}\usebox\hmboard\\[1mm]
\makebox[\wd\hmboard]{\sffamily\hspace{3mm}\textbf{h\#11}\hfill(4+2)}\end{minipage}\hfill
\begin{minipage}[t]{52mm}\vspace{0pt}
\textcolor{ink}{\sffamily\textsc{Lösung}}\par\smallskip\small
\begin{tabular}{@{}r@{\enspace}l@{\enspace}l@{}}
\textbf{1.} & \symking{}xb3 & \symking{}c1\\
\textbf{2.} & \symking{}c4 & \symking{}d2\\
\textbf{3.} & \symking{}d5 & \symking{}e3\\
\textbf{4.} & \symking{}e6 & \symking{}f4\\
\textbf{5.} & \symking{}f7 & \symking{}g5\\
\textbf{6.} & \symking{}g8 & \symking{}xh5\\
\textbf{7.} & \symking{}h8 & \symking{}g6\\
\textbf{8.} & \symking{}g8 & h5\\
\textbf{9.} & \symking{}h8 & h6\\
\textbf{10.} & \symking{}g8 & h7+\\
\textbf{11.} & \symking{}h8 & \symbishop{}g7\#\\
\end{tabular}\par\vspace{3mm}
\qrcode[height=12mm]{https://helpman.komtera.lt/?fen=5B2/8/8/7p/7P/kP6/8/1K6&moves=11}\\{\sffamily\tiny\color{ink}scannen zum Lösen}
\end{minipage}
\vspace{2mm}

{\raggedright\noindent\textcolor{ink}{\sffamily\textsc{Themen}}\enspace \themetag{reines Matt} \themetag{Mustermatt} \themetag{Idealmatt} \themetag{Switchback} \themetag{ein Stein zieht}\par\smallskip{\footnotesize\textit{Hinweis:} die Partei am Zug steht im Diagramm im Schach; die Lösung beginnt mit einem Schlag.}\par}

\clearpage
\phantomsection\label{hm:KBPvkrp}
\begin{tcolorbox}[enhanced,colback=ink,colframe=ink,arc=0pt,boxrule=0pt,left=4mm,right=4mm,top=2.5mm,bottom=2.5mm]
\color{white}\sffamily{\LARGE\bfseries KBPvkrp}\hfill{\Large h\#11}\\[0.5mm]
{\small 6 Steine\quad Schwarz am Zug}\hfill{\small Nr.~241\qquad\href{https://helpman.komtera.lt/?fen=8/8/8/1K6/2B2p2/k7/1P6/6r1&moves=11}{\color{white}\bfseries lösen\,\ensuremath{\nearrow}}}
\end{tcolorbox}
\vspace{2mm}
\noindent\begin{minipage}[t]{36mm}\vspace{0pt}\sffamily\small\raggedright
\hmfigure{längste eindeutige}{{\Large h\#11}}\hmfigure{eindeutig bei dieser Länge}{{\Large 1} Stellung}\hmfigure{längstes Matt}{{\Large h\#16}}\hmfigure{in der Klasse}{34,1 Milliarden Stellungen,\\ 22,1 Milliarden mit Hilfsmatt}
\end{minipage}\hfill
\begin{minipage}[t]{80mm}\vspace{0pt}\begin{tikzpicture}\begin{axis}[width=72mm,height=42mm,ybar,bar width=1.6pt,xmin=-1,xmax=33,ymin=0,ymax=3.24,axis y line*=left,axis x line*=bottom,xlabel={\scriptsize Halbzüge bis Matt},ylabel={\scriptsize Stellungen (Mrd.)},tick label style={font=\tiny},label style={font=\scriptsize},enlarge x limits=false,ylabel shift=-3pt,xlabel shift=-3pt]\addplot[fill=ink!55,draw=none] coordinates {(0,0.0001857) (1,0.004697) (2,0.01984) (3,0.1798) (4,0.6189) (5,1.94) (6,2.453) (7,2.834) (8,2.553) (9,2.945) (10,2.79) (11,2.889) (12,2) (13,0.6536) (14,0.1463) (15,0.04382) (16,0.008043) (17,0.001657) (18,0.0001029) (19,3.306e-05) (20,6.447e-05) (21,2.998e-05) (22,4.552e-05) (23,1.198e-05) (24,3.222e-05) (25,6.236e-06) (26,2.622e-05) (27,5.53e-07) (28,1.13e-05) (29,1.054e-06) (30,1.056e-05) (31,1e-09) (32,1.513e-06)};\draw[densely dotted] (axis cs:22,0) -- (axis cs:22,3.24);\end{axis}\begin{axis}[width=72mm,height=42mm,axis y line*=right,axis x line=none,ymode=log,xmin=-1,xmax=33,ylabel={\scriptsize eindeutig (log)},tick label style={font=\tiny},label style={font=\scriptsize},ylabel shift=-3pt]\addplot[mark=*,mark size=0.7pt,draw=black,thin] coordinates {(0,185694) (1,2146387) (2,8414090) (3,71010956) (4,146112887) (5,164129651) (6,82235628) (7,65663384) (8,33582419) (9,29724637) (10,17389549) (11,6933077) (12,1589534) (13,304346) (14,79176) (15,7381) (16,237) (17,113) (18,4) (19,22) (20,1) (21,1) (22,1)};\end{axis}\end{tikzpicture}\end{minipage}
\vspace{3mm}

\noindent\begin{minipage}[t]{60mm}\vspace{0pt}\sbox\hmboard{\chessboard[setfen=8/8/8/1K6/2B2p2/k7/1P6/6r1 b - - 0 1]}\usebox\hmboard\\[1mm]
\makebox[\wd\hmboard]{\sffamily\hspace{3mm}\textbf{h\#11}\hfill(3+3)}\end{minipage}\hfill
\begin{minipage}[t]{52mm}\vspace{0pt}
\textcolor{ink}{\sffamily\textsc{Lösung}}\par\smallskip\small
\begin{tabular}{@{}r@{\enspace}l@{\enspace}l@{}}
\textbf{1.} & \symking{}xb2 & \symking{}c5\\
\textbf{2.} & f3 & \symking{}d4\\
\textbf{3.} & f2 & \symking{}e3\\
\textbf{4.} & f1=\symknight{}+ & \symking{}f2\\
\textbf{5.} & \symking{}c3 & \symking{}xg1\\
\textbf{6.} & \symking{}d4 & \symking{}f2\\
\textbf{7.} & \symking{}e4 & \symking{}e1\\
\textbf{8.} & \symking{}f3 & \symbishop{}e6\\
\textbf{9.} & \symking{}g2 & \symbishop{}h3+\\
\textbf{10.} & \symking{}h1 & \symking{}f2\\
\textbf{11.} & \symknight{}h2 & \symbishop{}g2\#\\
\end{tabular}\par\vspace{3mm}
\qrcode[height=12mm]{https://helpman.komtera.lt/?fen=8/8/8/1K6/2B2p2/k7/1P6/6r1&moves=11}\\{\sffamily\tiny\color{ink}scannen zum Lösen}
\end{minipage}
\vspace{2mm}

{\raggedright\noindent\textcolor{ink}{\sffamily\textsc{Themen}}\enspace \themetag{reines Matt} \themetag{Mustermatt} \themetag{Idealmatt} \themetag{Umwandlung} \themetag{Unterverwandlung} \themetag{Switchback} \themetag{Selbstblock} \themetag{Umnov}\par\smallskip{\footnotesize\textit{Hinweis:} das Diagramm hat keinen legalen letzten Zug, kann also in keiner Partie entstehen.}\par}

\clearpage
\phantomsection\label{hm:KBvkppp}
\begin{tcolorbox}[enhanced,colback=ink,colframe=ink,arc=0pt,boxrule=0pt,left=4mm,right=4mm,top=2.5mm,bottom=2.5mm]
\color{white}\sffamily{\LARGE\bfseries KBvkppp}\hfill{\Large h\#11}\\[0.5mm]
{\small 6 Steine\quad Schwarz am Zug}\hfill{\small Nr.~242\qquad\href{https://helpman.komtera.lt/?fen=1k1K4/4B3/8/6p1/4p1p1/8/8/8&moves=11}{\color{white}\bfseries lösen\,\ensuremath{\nearrow}}}
\end{tcolorbox}
\vspace{2mm}
\noindent\begin{minipage}[t]{36mm}\vspace{0pt}\sffamily\small\raggedright
\hmfigure{längste eindeutige}{{\Large h\#11}}\hmfigure{eindeutig bei dieser Länge}{{\Large 1} Stellung}\hmfigure{längstes Matt}{{\Large h\#15.5}}\hmfigure{in der Klasse}{25,6 Milliarden Stellungen,\\ 3 Milliarden mit Hilfsmatt}
\end{minipage}\hfill
\begin{minipage}[t]{80mm}\vspace{0pt}\begin{tikzpicture}\begin{axis}[width=72mm,height=42mm,ybar,bar width=1.6pt,xmin=-1,xmax=32,ymin=0,ymax=441,axis y line*=left,axis x line*=bottom,xlabel={\scriptsize Halbzüge bis Matt},ylabel={\scriptsize Stellungen (Mio.)},tick label style={font=\tiny},label style={font=\scriptsize},enlarge x limits=false,ylabel shift=-3pt,xlabel shift=-3pt]\addplot[fill=ink!55,draw=none] coordinates {(0,0.02281) (1,0.09036) (2,0.2456) (3,0.8394) (4,2.695) (5,6.751) (6,19.43) (7,35.62) (8,81.29) (9,118.8) (10,219.6) (11,272.9) (12,400.9) (13,350.2) (14,372) (15,293.7) (16,233.6) (17,196.4) (18,122.3) (19,109.7) (20,56.19) (21,53.48) (22,21.23) (23,20.62) (24,6.38) (25,6.313) (26,1.51) (27,1.524) (28,0.1568) (29,0.1603) (30,0.01537) (31,0.01527)};\draw[densely dotted] (axis cs:22,0) -- (axis cs:22,441);\end{axis}\begin{axis}[width=72mm,height=42mm,axis y line*=right,axis x line=none,ymode=log,xmin=-1,xmax=32,ylabel={\scriptsize eindeutig (log)},tick label style={font=\tiny},label style={font=\scriptsize},ylabel shift=-3pt]\addplot[mark=*,mark size=0.7pt,draw=black,thin] coordinates {(0,22815) (1,89719) (2,173561) (3,295444) (4,770336) (5,968077) (6,1467736) (7,1313635) (8,1388361) (9,1254666) (10,1225221) (11,925348) (12,661883) (13,172724) (14,108179) (15,19645) (16,10281) (17,1043) (18,328) (19,69) (20,4) (21,2) (22,1)};\end{axis}\end{tikzpicture}\end{minipage}
\vspace{3mm}

\noindent\begin{minipage}[t]{60mm}\vspace{0pt}\sbox\hmboard{\chessboard[setfen=1k1K4/4B3/8/6p1/4p1p1/8/8/8 b - - 0 1]}\usebox\hmboard\\[1mm]
\makebox[\wd\hmboard]{\sffamily\hspace{3mm}\textbf{h\#11}\hfill(2+4)}\end{minipage}\hfill
\begin{minipage}[t]{52mm}\vspace{0pt}
\textcolor{ink}{\sffamily\textsc{Lösung}}\par\smallskip\small
\begin{tabular}{@{}r@{\enspace}l@{\enspace}l@{}}
\textbf{1.} & e3 & \symking{}d7\\
\textbf{2.} & e2 & \symking{}c6\\
\textbf{3.} & e1=\symqueen{} & \symking{}b5\\
\textbf{4.} & \symking{}c7 & \symbishop{}d8+\\
\textbf{5.} & \symking{}d6 & \symking{}c4\\
\textbf{6.} & \symqueen{}e8 & \symking{}d3\\
\textbf{7.} & \symking{}e5 & \symking{}e2\\
\textbf{8.} & \symking{}f4+ & \symking{}f1\\
\textbf{9.} & \symking{}g3 & \symbishop{}c7+\\
\textbf{10.} & \symking{}h4 & \symking{}g2\\
\textbf{11.} & \symqueen{}h5 & \symbishop{}g3\#\\
\end{tabular}\par\vspace{3mm}
\qrcode[height=12mm]{https://helpman.komtera.lt/?fen=1k1K4/4B3/8/6p1/4p1p1/8/8/8&moves=11}\\{\sffamily\tiny\color{ink}scannen zum Lösen}
\end{minipage}
\vspace{2mm}

{\raggedright\noindent\textcolor{ink}{\sffamily\textsc{Themen}}\enspace \themetag{reines Matt} \themetag{Mustermatt} \themetag{Idealmatt} \themetag{Umwandlung} \themetag{Selbstblock} \themetag{ohne Schlag} \themetag{Inder}\par}

\clearpage
\phantomsection\label{hm:KQBPvkp}
\begin{tcolorbox}[enhanced,colback=ink,colframe=ink,arc=0pt,boxrule=0pt,left=4mm,right=4mm,top=2.5mm,bottom=2.5mm]
\color{white}\sffamily{\LARGE\bfseries KQBPvkp}\hfill{\Large h\#11}\\[0.5mm]
{\small 6 Steine\quad Schwarz am Zug}\hfill{\small Nr.~243\qquad\href{https://helpman.komtera.lt/?fen=5B2/8/8/7p/7P/kQ6/8/1K6&moves=11}{\color{white}\bfseries lösen\,\ensuremath{\nearrow}}}
\end{tcolorbox}
\vspace{2mm}
\noindent\begin{minipage}[t]{36mm}\vspace{0pt}\sffamily\small\raggedright
\hmfigure{längste eindeutige}{{\Large h\#11}}\hmfigure{eindeutig bei dieser Länge}{{\Large 1} Stellung}\hmfigure{längstes Matt}{{\Large h\#15}}\hmfigure{in der Klasse}{34,1 Milliarden Stellungen,\\ 20,9 Milliarden mit Hilfsmatt}
\end{minipage}\hfill
\begin{minipage}[t]{80mm}\vspace{0pt}\begin{tikzpicture}\begin{axis}[width=72mm,height=42mm,ybar,bar width=1.6pt,xmin=-1,xmax=31,ymin=0,ymax=10.78,axis y line*=left,axis x line*=bottom,xlabel={\scriptsize Halbzüge bis Matt},ylabel={\scriptsize Stellungen (Mrd.)},tick label style={font=\tiny},label style={font=\scriptsize},enlarge x limits=false,ylabel shift=-3pt,xlabel shift=-3pt]\addplot[fill=ink!55,draw=none] coordinates {(0,0.07363) (1,0.4385) (2,1.673) (3,4.604) (4,9.796) (5,2.695) (6,1.433) (7,0.01528) (8,0.05108) (9,1.239e-05) (10,0.04684) (11,2.44e-07) (12,0.05218) (13,1.7e-08) (14,0.02043) (16,0.003622) (18,0.00101) (20,0.0003677) (22,0.0002186) (24,8.352e-05) (26,1.624e-05) (28,1.057e-05) (30,2.33e-07)};\draw[densely dotted] (axis cs:22,0) -- (axis cs:22,10.78);\end{axis}\begin{axis}[width=72mm,height=42mm,axis y line*=right,axis x line=none,ymode=log,xmin=-1,xmax=31,ylabel={\scriptsize eindeutig (log)},tick label style={font=\tiny},label style={font=\scriptsize},ylabel shift=-3pt]\addplot[mark=*,mark size=0.7pt,draw=black,thin] coordinates {(0,73634181) (1,333656294) (2,1093758767) (3,1010834451) (4,488898360) (5,2442704) (6,2813640) (7,970) (8,1049864) (9,22) (10,562788) (12,301177) (14,13871) (16,6004) (18,292) (20,22) (22,1)};\end{axis}\end{tikzpicture}\end{minipage}
\vspace{3mm}

\noindent\begin{minipage}[t]{60mm}\vspace{0pt}\sbox\hmboard{\chessboard[setfen=5B2/8/8/7p/7P/kQ6/8/1K6 b - - 0 1]}\usebox\hmboard\\[1mm]
\makebox[\wd\hmboard]{\sffamily\hspace{3mm}\textbf{h\#11}\hfill(4+2)}\end{minipage}\hfill
\begin{minipage}[t]{52mm}\vspace{0pt}
\textcolor{ink}{\sffamily\textsc{Lösung}}\par\smallskip\small
\begin{tabular}{@{}r@{\enspace}l@{\enspace}l@{}}
\textbf{1.} & \symking{}xb3 & \symking{}c1\\
\textbf{2.} & \symking{}c4 & \symking{}d2\\
\textbf{3.} & \symking{}d5 & \symking{}e3\\
\textbf{4.} & \symking{}e6 & \symking{}f4\\
\textbf{5.} & \symking{}f7 & \symking{}g5\\
\textbf{6.} & \symking{}g8 & \symking{}xh5\\
\textbf{7.} & \symking{}h8 & \symking{}g6\\
\textbf{8.} & \symking{}g8 & h5\\
\textbf{9.} & \symking{}h8 & h6\\
\textbf{10.} & \symking{}g8 & h7+\\
\textbf{11.} & \symking{}h8 & \symbishop{}g7\#\\
\end{tabular}\par\vspace{3mm}
\qrcode[height=12mm]{https://helpman.komtera.lt/?fen=5B2/8/8/7p/7P/kQ6/8/1K6&moves=11}\\{\sffamily\tiny\color{ink}scannen zum Lösen}
\end{minipage}
\vspace{2mm}

{\raggedright\noindent\textcolor{ink}{\sffamily\textsc{Themen}}\enspace \themetag{reines Matt} \themetag{Mustermatt} \themetag{Idealmatt} \themetag{Switchback} \themetag{ein Stein zieht}\par\smallskip{\footnotesize\textit{Hinweis:} das Diagramm hat keinen legalen letzten Zug, kann also in keiner Partie entstehen.}\par}

\clearpage
\phantomsection\label{hm:KBBBvkp}
\begin{tcolorbox}[enhanced,colback=ink,colframe=ink,arc=0pt,boxrule=0pt,left=4mm,right=4mm,top=2.5mm,bottom=2.5mm]
\color{white}\sffamily{\LARGE\bfseries KBBBvkp}\hfill{\Large h\#10.5}\\[0.5mm]
{\small 6 Steine\quad Weiß am Zug}\hfill{\small Nr.~244\qquad\href{https://helpman.komtera.lt/?fen=8/K2p4/8/1B1k4/B7/3B4/8/8&moves=10.5}{\color{white}\bfseries lösen\,\ensuremath{\nearrow}}}
\end{tcolorbox}
\vspace{2mm}
\noindent\begin{minipage}[t]{36mm}\vspace{0pt}\sffamily\small\raggedright
\hmfigure{längste eindeutige}{{\Large h\#10.5}}\hmfigure{eindeutig bei dieser Länge}{{\Large 2} Stellungen}\hmfigure{längstes Matt}{{\Large h\#15}}\hmfigure{in der Klasse}{45,4 Milliarden Stellungen,\\ 4,91 Milliarden mit Hilfsmatt}
\end{minipage}\hfill
\begin{minipage}[t]{80mm}\vspace{0pt}\begin{tikzpicture}\begin{axis}[width=72mm,height=42mm,ybar,bar width=1.6pt,xmin=-1,xmax=31,ymin=0,ymax=1.309,axis y line*=left,axis x line*=bottom,xlabel={\scriptsize Halbzüge bis Matt},ylabel={\scriptsize Stellungen (Mrd.)},tick label style={font=\tiny},label style={font=\scriptsize},enlarge x limits=false,ylabel shift=-3pt,xlabel shift=-3pt]\addplot[fill=ink!55,draw=none] coordinates {(0,0.003634) (1,0.0121) (2,0.04091) (3,0.1736) (4,0.5646) (5,0.6886) (6,1.19) (7,0.4921) (8,0.3996) (9,0.1012) (10,0.07696) (11,0.04435) (12,0.08682) (13,0.06507) (14,0.1296) (15,0.08822) (16,0.1428) (17,0.09296) (18,0.1211) (19,0.08105) (20,0.08642) (21,0.06091) (22,0.05769) (23,0.04079) (24,0.02457) (25,0.01845) (26,0.0105) (27,0.007589) (28,0.002767) (29,0.002631) (30,4.213e-06)};\draw[densely dotted] (axis cs:21,0) -- (axis cs:21,1.309);\end{axis}\begin{axis}[width=72mm,height=42mm,axis y line*=right,axis x line=none,ymode=log,xmin=-1,xmax=31,ylabel={\scriptsize eindeutig (log)},tick label style={font=\tiny},label style={font=\scriptsize},ylabel shift=-3pt]\addplot[mark=*,mark size=0.7pt,draw=black,thin] coordinates {(0,3634181) (1,10654986) (2,33913386) (3,84629859) (4,128736714) (5,21903397) (6,15219119) (7,1905578) (8,838075) (9,219230) (10,322727) (11,131320) (12,221805) (13,47517) (14,74891) (15,17042) (16,22497) (17,3086) (18,3549) (19,123) (20,28) (21,2)};\end{axis}\end{tikzpicture}\end{minipage}
\vspace{3mm}

\noindent\begin{minipage}[t]{60mm}\vspace{0pt}\sbox\hmboard{\chessboard[setfen=8/K2p4/8/1B1k4/B7/3B4/8/8 w - - 0 1]}\usebox\hmboard\\[1mm]
\makebox[\wd\hmboard]{\sffamily\hspace{3mm}\textbf{h\#10.5}\hfill(4+2)}\end{minipage}\hfill
\begin{minipage}[t]{52mm}\vspace{0pt}
\textcolor{ink}{\sffamily\textsc{Lösung}}\par\smallskip\small
\begin{tabular}{@{}r@{\enspace}l@{\enspace}l@{}}
\textbf{1.} & \ldots & \symbishop{}dc4+\\
\textbf{2.} & \symking{}e4 & \symking{}b6\\
\textbf{3.} & \symking{}f3 & \symking{}c5\\
\textbf{4.} & \symking{}g2 & \symbishop{}f1+\\
\textbf{5.} & \symking{}h1 & \symbishop{}be2\\
\textbf{6.} & d5 & \symbishop{}d7\\
\textbf{7.} & d4 & \symbishop{}dh3\\
\textbf{8.} & d3 & \symking{}d4\\
\textbf{9.} & dxe2 & \symking{}e3\\
\textbf{10.} & exf1=\symknight{}+ & \symking{}f2\\
\textbf{11.} & \symknight{}h2 & \symbishop{}g2\#\\
\end{tabular}\par\vspace{3mm}
\qrcode[height=12mm]{https://helpman.komtera.lt/?fen=8/K2p4/8/1B1k4/B7/3B4/8/8&moves=10.5}\\{\sffamily\tiny\color{ink}scannen zum Lösen}
\end{minipage}
\vspace{2mm}

{\raggedright\noindent\textcolor{ink}{\sffamily\textsc{Themen}}\enspace \themetag{reines Matt} \themetag{Mustermatt} \themetag{Idealmatt} \themetag{Umwandlung} \themetag{Unterverwandlung} \themetag{Excelsior} \themetag{Selbstblock} \themetag{Umnov}\par}

\clearpage
\phantomsection\label{hm:KBPvkpp}
\begin{tcolorbox}[enhanced,colback=ink,colframe=ink,arc=0pt,boxrule=0pt,left=4mm,right=4mm,top=2.5mm,bottom=2.5mm]
\color{white}\sffamily{\LARGE\bfseries KBPvkpp}\hfill{\Large h\#10.5}\\[0.5mm]
{\small 6 Steine\quad Weiß am Zug}\hfill{\small Nr.~245\qquad\href{https://helpman.komtera.lt/?fen=8/Kp6/8/2k5/8/4p3/4P3/1B6&moves=10.5}{\color{white}\bfseries lösen\,\ensuremath{\nearrow}}}
\end{tcolorbox}
\vspace{2mm}
\noindent\begin{minipage}[t]{36mm}\vspace{0pt}\sffamily\small\raggedright
\hmfigure{längste eindeutige}{{\Large h\#10.5}}\hmfigure{eindeutig bei dieser Länge}{{\Large 28} Stellungen}\hmfigure{längstes Matt}{{\Large h\#16}}\hmfigure{in der Klasse}{25,6 Milliarden Stellungen,\\ 9,02 Milliarden mit Hilfsmatt}
\end{minipage}\hfill
\begin{minipage}[t]{80mm}\vspace{0pt}\begin{tikzpicture}\begin{axis}[width=72mm,height=42mm,ybar,bar width=1.6pt,xmin=-1,xmax=33,ymin=0,ymax=1.274,axis y line*=left,axis x line*=bottom,xlabel={\scriptsize Halbzüge bis Matt},ylabel={\scriptsize Stellungen (Mrd.)},tick label style={font=\tiny},label style={font=\scriptsize},enlarge x limits=false,ylabel shift=-3pt,xlabel shift=-3pt]\addplot[fill=ink!55,draw=none] coordinates {(0,0.0001729) (1,0.003578) (2,0.009533) (3,0.07319) (4,0.1983) (5,0.6969) (6,0.9546) (7,0.9473) (8,1.077) (9,1.002) (10,1.158) (11,1.075) (12,1.043) (13,0.3654) (14,0.2344) (15,0.08607) (16,0.05883) (17,0.0198) (18,0.01144) (19,0.004485) (20,0.002356) (21,0.0008626) (22,0.0003238) (23,0.0001159) (24,2.805e-05) (25,6.499e-06) (26,4.903e-06) (27,8.91e-07) (28,1.757e-06) (29,2.4e-08) (30,1.31e-07) (32,3e-09)};\draw[densely dotted] (axis cs:21,0) -- (axis cs:21,1.274);\end{axis}\begin{axis}[width=72mm,height=42mm,axis y line*=right,axis x line=none,ymode=log,xmin=-1,xmax=33,ylabel={\scriptsize eindeutig (log)},tick label style={font=\tiny},label style={font=\scriptsize},ylabel shift=-3pt]\addplot[mark=*,mark size=0.7pt,draw=black,thin] coordinates {(0,172900) (1,1387653) (2,3683555) (3,28120821) (4,42351664) (5,52714794) (6,32298423) (7,23918345) (8,16349440) (9,12383387) (10,8209132) (11,5182433) (12,2380679) (13,348518) (14,239156) (15,80483) (16,28597) (17,4912) (18,1758) (19,222) (20,40) (21,28)};\end{axis}\end{tikzpicture}\end{minipage}
\vspace{3mm}

\noindent\begin{minipage}[t]{60mm}\vspace{0pt}\sbox\hmboard{\chessboard[setfen=8/Kp6/8/2k5/8/4p3/4P3/1B6 w - - 0 1]}\usebox\hmboard\\[1mm]
\makebox[\wd\hmboard]{\sffamily\hspace{3mm}\textbf{h\#10.5}\hfill(3+3)}\end{minipage}\hfill
\begin{minipage}[t]{52mm}\vspace{0pt}
\textcolor{ink}{\sffamily\textsc{Lösung}}\par\smallskip\small
\begin{tabular}{@{}r@{\enspace}l@{\enspace}l@{}}
\textbf{1.} & \ldots & \symbishop{}f5\\
\textbf{2.} & \symking{}d6 & \symking{}b6\\
\textbf{3.} & \symking{}e7 & \symking{}c5\\
\textbf{4.} & b5 & \symking{}d4\\
\textbf{5.} & b4 & \symking{}xe3\\
\textbf{6.} & b3 & \symking{}f4\\
\textbf{7.} & b2 & e4\\
\textbf{8.} & b1=\symrook{} & e5\\
\textbf{9.} & \symrook{}b8 & e6\\
\textbf{10.} & \symking{}f6 & e7\\
\textbf{11.} & \symrook{}f8 & exf8=\symqueen{}\#\\
\end{tabular}\par\vspace{3mm}
\qrcode[height=12mm]{https://helpman.komtera.lt/?fen=8/Kp6/8/2k5/8/4p3/4P3/1B6&moves=10.5}\\{\sffamily\tiny\color{ink}scannen zum Lösen}
\end{minipage}
\vspace{2mm}

{\raggedright\noindent\textcolor{ink}{\sffamily\textsc{Themen}}\enspace \themetag{reines Matt} \themetag{Mustermatt} \themetag{Idealmatt} \themetag{Umwandlung} \themetag{Unterverwandlung} \themetag{Excelsior} \themetag{ohne Schach} \themetag{Umnov}\par}

\clearpage
\phantomsection\label{hm:KBPvkqp}
\begin{tcolorbox}[enhanced,colback=ink,colframe=ink,arc=0pt,boxrule=0pt,left=4mm,right=4mm,top=2.5mm,bottom=2.5mm]
\color{white}\sffamily{\LARGE\bfseries KBPvkqp}\hfill{\Large h\#10.5}\\[0.5mm]
{\small 6 Steine\quad Weiß am Zug}\hfill{\small Nr.~246\qquad\href{https://helpman.komtera.lt/?fen=5B2/8/8/7p/7P/1k6/8/1Kq5&moves=10.5}{\color{white}\bfseries lösen\,\ensuremath{\nearrow}}}
\end{tcolorbox}
\vspace{2mm}
\noindent\begin{minipage}[t]{36mm}\vspace{0pt}\sffamily\small\raggedright
\hmfigure{längste eindeutige}{{\Large h\#10.5}}\hmfigure{eindeutig bei dieser Länge}{{\Large 1} Stellung}\hmfigure{längstes Matt}{{\Large h\#17}}\hmfigure{in der Klasse}{34,1 Milliarden Stellungen,\\ 20,4 Milliarden mit Hilfsmatt}
\end{minipage}\hfill
\begin{minipage}[t]{80mm}\vspace{0pt}\begin{tikzpicture}\begin{axis}[width=72mm,height=42mm,ybar,bar width=1.6pt,xmin=-1,xmax=35,ymin=0,ymax=3.272,axis y line*=left,axis x line*=bottom,xlabel={\scriptsize Halbzüge bis Matt},ylabel={\scriptsize Stellungen (Mrd.)},tick label style={font=\tiny},label style={font=\scriptsize},enlarge x limits=false,ylabel shift=-3pt,xlabel shift=-3pt]\addplot[fill=ink!55,draw=none] coordinates {(0,0.0001289) (1,0.002859) (2,0.01387) (3,0.1412) (4,0.4543) (5,1.67) (6,1.986) (7,2.786) (8,2.182) (9,2.936) (10,2.346) (11,2.975) (12,1.8) (13,0.9109) (14,0.1386) (15,0.05401) (16,0.00541) (17,0.003355) (18,7.32e-05) (19,0.0005212) (20,4.828e-05) (21,0.0002627) (22,3.026e-05) (23,0.0001414) (24,2.373e-05) (25,3.514e-05) (26,2.138e-05) (27,1.713e-05) (28,1.142e-05) (29,4.86e-06) (30,9.945e-06) (31,4.2e-08) (32,1.816e-06) (33,1e-08) (34,1.792e-06)};\draw[densely dotted] (axis cs:21,0) -- (axis cs:21,3.272);\end{axis}\begin{axis}[width=72mm,height=42mm,axis y line*=right,axis x line=none,ymode=log,xmin=-1,xmax=35,ylabel={\scriptsize eindeutig (log)},tick label style={font=\tiny},label style={font=\scriptsize},ylabel shift=-3pt]\addplot[mark=*,mark size=0.7pt,draw=black,thin] coordinates {(0,128929) (1,1376408) (2,5594174) (3,54878815) (4,97650482) (5,135904397) (6,59688373) (7,62433985) (8,27493565) (9,29796058) (10,13873420) (11,9010972) (12,1788867) (13,509209) (14,42420) (15,15024) (16,88) (17,685) (18,6) (19,77) (20,1) (21,1)};\end{axis}\end{tikzpicture}\end{minipage}
\vspace{3mm}

\noindent\begin{minipage}[t]{60mm}\vspace{0pt}\sbox\hmboard{\chessboard[setfen=5B2/8/8/7p/7P/1k6/8/1Kq5 w - - 0 1]}\usebox\hmboard\\[1mm]
\makebox[\wd\hmboard]{\sffamily\hspace{3mm}\textbf{h\#10.5}\hfill(3+3)}\end{minipage}\hfill
\begin{minipage}[t]{52mm}\vspace{0pt}
\textcolor{ink}{\sffamily\textsc{Lösung}}\par\smallskip\small
\begin{tabular}{@{}r@{\enspace}l@{\enspace}l@{}}
\textbf{1.} & \ldots & \symking{}xc1\\
\textbf{2.} & \symking{}c4 & \symking{}d2\\
\textbf{3.} & \symking{}d5 & \symking{}e3\\
\textbf{4.} & \symking{}e6 & \symking{}f4\\
\textbf{5.} & \symking{}f7 & \symking{}g5\\
\textbf{6.} & \symking{}g8 & \symking{}xh5\\
\textbf{7.} & \symking{}h8 & \symking{}g6\\
\textbf{8.} & \symking{}g8 & h5\\
\textbf{9.} & \symking{}h8 & h6\\
\textbf{10.} & \symking{}g8 & h7+\\
\textbf{11.} & \symking{}h8 & \symbishop{}g7\#\\
\end{tabular}\par\vspace{3mm}
\qrcode[height=12mm]{https://helpman.komtera.lt/?fen=5B2/8/8/7p/7P/1k6/8/1Kq5&moves=10.5}\\{\sffamily\tiny\color{ink}scannen zum Lösen}
\end{minipage}
\vspace{2mm}

{\raggedright\noindent\textcolor{ink}{\sffamily\textsc{Themen}}\enspace \themetag{reines Matt} \themetag{Mustermatt} \themetag{Idealmatt} \themetag{Switchback} \themetag{ein Stein zieht}\par\smallskip{\footnotesize\textit{Hinweis:} die Partei am Zug steht im Diagramm im Schach; die Lösung beginnt mit einem Schlag.}\par}

\clearpage
\phantomsection\label{hm:KBvkbnp}
\begin{tcolorbox}[enhanced,colback=ink,colframe=ink,arc=0pt,boxrule=0pt,left=4mm,right=4mm,top=2.5mm,bottom=2.5mm]
\color{white}\sffamily{\LARGE\bfseries KBvkbnp}\hfill{\Large h\#10.5}\\[0.5mm]
{\small 6 Steine\quad Weiß am Zug}\hfill{\small Nr.~247\qquad\href{https://helpman.komtera.lt/?fen=8/2p5/8/1b6/8/KB6/n7/k7&moves=10.5}{\color{white}\bfseries lösen\,\ensuremath{\nearrow}}}
\end{tcolorbox}
\vspace{2mm}
\noindent\begin{minipage}[t]{36mm}\vspace{0pt}\sffamily\small\raggedright
\hmfigure{längste eindeutige}{{\Large h\#10.5}}\hmfigure{eindeutig bei dieser Länge}{{\Large 1} Stellung}\hmfigure{längstes Matt}{{\Large h\#16.5}}\hmfigure{in der Klasse}{45,4 Milliarden Stellungen,\\ 30 Milliarden mit Hilfsmatt}
\end{minipage}\hfill
\begin{minipage}[t]{80mm}\vspace{0pt}\begin{tikzpicture}\begin{axis}[width=72mm,height=42mm,ybar,bar width=1.6pt,xmin=-1,xmax=34,ymin=0,ymax=4.705,axis y line*=left,axis x line*=bottom,xlabel={\scriptsize Halbzüge bis Matt},ylabel={\scriptsize Stellungen (Mrd.)},tick label style={font=\tiny},label style={font=\scriptsize},enlarge x limits=false,ylabel shift=-3pt,xlabel shift=-3pt]\addplot[fill=ink!55,draw=none] coordinates {(0,0.0001569) (1,0.0005979) (2,0.003691) (3,0.0128) (4,0.05347) (5,0.1391) (6,0.3178) (7,0.6664) (8,1.127) (9,1.967) (10,2.769) (11,4.028) (12,4.277) (13,4.144) (14,3.351) (15,2.892) (16,1.644) (17,1.38) (18,0.5005) (19,0.4955) (20,0.1013) (21,0.1106) (22,0.004642) (23,0.005514) (24,5.96e-07) (25,1.131e-05) (27,5.846e-06) (29,2.099e-06) (31,3.72e-07) (33,1.8e-08)};\draw[densely dotted] (axis cs:21,0) -- (axis cs:21,4.705);\end{axis}\begin{axis}[width=72mm,height=42mm,axis y line*=right,axis x line=none,ymode=log,xmin=-1,xmax=34,ylabel={\scriptsize eindeutig (log)},tick label style={font=\tiny},label style={font=\scriptsize},ylabel shift=-3pt]\addplot[mark=*,mark size=0.7pt,draw=black,thin] coordinates {(0,156893) (1,595603) (2,2877519) (3,4744441) (4,12164710) (5,20061121) (6,15933492) (7,18415657) (8,15045271) (9,17973117) (10,14082842) (11,10648907) (12,5674553) (13,2837894) (14,1041780) (15,397686) (16,86135) (17,27855) (18,5661) (19,282) (20,30) (21,1)};\end{axis}\end{tikzpicture}\end{minipage}
\vspace{3mm}

\noindent\begin{minipage}[t]{60mm}\vspace{0pt}\sbox\hmboard{\chessboard[setfen=8/2p5/8/1b6/8/KB6/n7/k7 w - - 0 1]}\usebox\hmboard\\[1mm]
\makebox[\wd\hmboard]{\sffamily\hspace{3mm}\textbf{h\#10.5}\hfill(2+4)}\end{minipage}\hfill
\begin{minipage}[t]{52mm}\vspace{0pt}
\textcolor{ink}{\sffamily\textsc{Lösung}}\par\smallskip\small
\begin{tabular}{@{}r@{\enspace}l@{\enspace}l@{}}
\textbf{1.} & \ldots & \symbishop{}a4\\
\textbf{2.} & \symknight{}b4 & \symking{}b3\\
\textbf{3.} & \symknight{}c6 & \symking{}c3\\
\textbf{4.} & \symking{}a2 & \symbishop{}xb5\\
\textbf{5.} & \symking{}a3 & \symbishop{}a6\\
\textbf{6.} & \symking{}a4 & \symking{}c4\\
\textbf{7.} & \symking{}a5 & \symking{}d5\\
\textbf{8.} & \symking{}b6 & \symking{}e6\\
\textbf{9.} & \symking{}a7 & \symking{}d7\\
\textbf{10.} & \symking{}a8 & \symking{}xc7\\
\textbf{11.} & \symknight{}a7 & \symbishop{}b7\#\\
\end{tabular}\par\vspace{3mm}
\qrcode[height=12mm]{https://helpman.komtera.lt/?fen=8/2p5/8/1b6/8/KB6/n7/k7&moves=10.5}\\{\sffamily\tiny\color{ink}scannen zum Lösen}
\end{minipage}
\vspace{2mm}

{\raggedright\noindent\textcolor{ink}{\sffamily\textsc{Themen}}\enspace \themetag{reines Matt} \themetag{Mustermatt} \themetag{Idealmatt} \themetag{Selbstblock} \themetag{ohne Schach}\par}

\clearpage
\phantomsection\label{hm:KBvknpp}
\begin{tcolorbox}[enhanced,colback=ink,colframe=ink,arc=0pt,boxrule=0pt,left=4mm,right=4mm,top=2.5mm,bottom=2.5mm]
\color{white}\sffamily{\LARGE\bfseries KBvknpp}\hfill{\Large h\#10.5}\\[0.5mm]
{\small 6 Steine\quad Weiß am Zug}\hfill{\small Nr.~248\qquad\href{https://helpman.komtera.lt/?fen=8/2p5/8/1p6/8/KB6/n7/k7&moves=10.5}{\color{white}\bfseries lösen\,\ensuremath{\nearrow}}}
\end{tcolorbox}
\vspace{2mm}
\noindent\begin{minipage}[t]{36mm}\vspace{0pt}\sffamily\small\raggedright
\hmfigure{längste eindeutige}{{\Large h\#10.5}}\hmfigure{eindeutig bei dieser Länge}{{\Large 1} Stellung}\hmfigure{längstes Matt}{{\Large h\#14.5}}\hmfigure{in der Klasse}{34,1 Milliarden Stellungen,\\ 11,8 Milliarden mit Hilfsmatt}
\end{minipage}\hfill
\begin{minipage}[t]{80mm}\vspace{0pt}\begin{tikzpicture}\begin{axis}[width=72mm,height=42mm,ybar,bar width=1.6pt,xmin=-1,xmax=30,ymin=0,ymax=1.964,axis y line*=left,axis x line*=bottom,xlabel={\scriptsize Halbzüge bis Matt},ylabel={\scriptsize Stellungen (Mrd.)},tick label style={font=\tiny},label style={font=\scriptsize},enlarge x limits=false,ylabel shift=-3pt,xlabel shift=-3pt]\addplot[fill=ink!55,draw=none] coordinates {(0,4.784e-05) (1,0.00018) (2,0.0005057) (3,0.001916) (4,0.006061) (5,0.0182) (6,0.04999) (7,0.1137) (8,0.2457) (9,0.4274) (10,0.7415) (11,1.072) (12,1.561) (13,1.65) (14,1.785) (15,1.499) (16,1.081) (17,0.7625) (18,0.3324) (19,0.2992) (20,0.072) (21,0.07344) (22,0.003138) (23,0.00341) (24,1.319e-06) (25,6.65e-07) (27,1.63e-07) (29,3.5e-08)};\draw[densely dotted] (axis cs:21,0) -- (axis cs:21,1.964);\end{axis}\begin{axis}[width=72mm,height=42mm,axis y line*=right,axis x line=none,ymode=log,xmin=-1,xmax=30,ylabel={\scriptsize eindeutig (log)},tick label style={font=\tiny},label style={font=\scriptsize},ylabel shift=-3pt]\addplot[mark=*,mark size=0.7pt,draw=black,thin] coordinates {(0,47838) (1,179384) (2,275023) (3,630286) (4,1817102) (5,3011882) (6,3882435) (7,4551245) (8,4999676) (9,5552571) (10,5929169) (11,4409672) (12,3383337) (13,1300755) (14,526742) (15,175252) (16,48232) (17,6593) (18,1549) (19,256) (20,5) (21,1)};\end{axis}\end{tikzpicture}\end{minipage}
\vspace{3mm}

\noindent\begin{minipage}[t]{60mm}\vspace{0pt}\sbox\hmboard{\chessboard[setfen=8/2p5/8/1p6/8/KB6/n7/k7 w - - 0 1]}\usebox\hmboard\\[1mm]
\makebox[\wd\hmboard]{\sffamily\hspace{3mm}\textbf{h\#10.5}\hfill(2+4)}\end{minipage}\hfill
\begin{minipage}[t]{52mm}\vspace{0pt}
\textcolor{ink}{\sffamily\textsc{Lösung}}\par\smallskip\small
\begin{tabular}{@{}r@{\enspace}l@{\enspace}l@{}}
\textbf{1.} & \ldots & \symbishop{}a4\\
\textbf{2.} & \symknight{}b4 & \symking{}b3\\
\textbf{3.} & \symknight{}c6 & \symking{}c3\\
\textbf{4.} & \symking{}a2 & \symbishop{}xb5\\
\textbf{5.} & \symking{}a3 & \symbishop{}a6\\
\textbf{6.} & \symking{}a4 & \symking{}c4\\
\textbf{7.} & \symking{}a5 & \symking{}d5\\
\textbf{8.} & \symking{}b6 & \symking{}e6\\
\textbf{9.} & \symking{}a7 & \symking{}d7\\
\textbf{10.} & \symking{}a8 & \symking{}xc7\\
\textbf{11.} & \symknight{}a7 & \symbishop{}b7\#\\
\end{tabular}\par\vspace{3mm}
\qrcode[height=12mm]{https://helpman.komtera.lt/?fen=8/2p5/8/1p6/8/KB6/n7/k7&moves=10.5}\\{\sffamily\tiny\color{ink}scannen zum Lösen}
\end{minipage}
\vspace{2mm}

{\raggedright\noindent\textcolor{ink}{\sffamily\textsc{Themen}}\enspace \themetag{reines Matt} \themetag{Mustermatt} \themetag{Idealmatt} \themetag{Selbstblock} \themetag{ohne Schach}\par}

\clearpage
\phantomsection\label{hm:KNvkqqp}
\begin{tcolorbox}[enhanced,colback=ink,colframe=ink,arc=0pt,boxrule=0pt,left=4mm,right=4mm,top=2.5mm,bottom=2.5mm]
\color{white}\sffamily{\LARGE\bfseries KNvkqqp}\hfill{\Large h\#10.5}\\[0.5mm]
{\small 6 Steine\quad Weiß am Zug}\hfill{\small Nr.~249\qquad\href{https://helpman.komtera.lt/?fen=7N/8/3k1p2/8/8/8/2qq4/K7&moves=10.5}{\color{white}\bfseries lösen\,\ensuremath{\nearrow}}}
\end{tcolorbox}
\vspace{2mm}
\noindent\begin{minipage}[t]{36mm}\vspace{0pt}\sffamily\small\raggedright
\hmfigure{längste eindeutige}{{\Large h\#10.5}}\hmfigure{eindeutig bei dieser Länge}{{\Large 11} Stellungen}\hmfigure{längstes Matt}{{\Large h\#12.5}}\hmfigure{in der Klasse}{45,4 Milliarden Stellungen,\\ 12 Milliarden mit Hilfsmatt}
\end{minipage}\hfill
\begin{minipage}[t]{80mm}\vspace{0pt}\begin{tikzpicture}\begin{axis}[width=72mm,height=42mm,ybar,bar width=1.6pt,xmin=-1,xmax=26,ymin=0,ymax=2.471,axis y line*=left,axis x line*=bottom,xlabel={\scriptsize Halbzüge bis Matt},ylabel={\scriptsize Stellungen (Mrd.)},tick label style={font=\tiny},label style={font=\scriptsize},enlarge x limits=false,ylabel shift=-3pt,xlabel shift=-3pt]\addplot[fill=ink!55,draw=none] coordinates {(0,2.25e-06) (1,1.151e-05) (2,6.483e-05) (3,0.0004333) (4,0.001764) (5,0.008873) (6,0.02223) (7,0.08955) (8,0.1479) (9,0.4938) (10,0.5379) (11,1.457) (12,1.016) (13,2.247) (14,1.111) (15,2.059) (16,0.7853) (17,1.289) (18,0.2691) (19,0.3728) (20,0.02161) (21,0.02304) (22,2.312e-05) (23,8.772e-06) (24,4e-09) (25,1.2e-08)};\draw[densely dotted] (axis cs:21,0) -- (axis cs:21,2.471);\end{axis}\begin{axis}[width=72mm,height=42mm,axis y line*=right,axis x line=none,ymode=log,xmin=-1,xmax=26,ylabel={\scriptsize eindeutig (log)},tick label style={font=\tiny},label style={font=\scriptsize},ylabel shift=-3pt]\addplot[mark=*,mark size=0.7pt,draw=black,thin] coordinates {(0,2250) (1,11366) (2,42082) (3,214689) (4,505227) (5,1473018) (6,1534838) (7,4297766) (8,3010154) (9,7376597) (10,3457734) (11,6700725) (12,2228922) (13,3772489) (14,994320) (15,1468501) (16,261878) (17,311477) (18,16143) (19,11977) (20,24) (21,11)};\end{axis}\end{tikzpicture}\end{minipage}
\vspace{3mm}

\noindent\begin{minipage}[t]{60mm}\vspace{0pt}\sbox\hmboard{\chessboard[setfen=7N/8/3k1p2/8/8/8/2qq4/K7 w - - 0 1]}\usebox\hmboard\\[1mm]
\makebox[\wd\hmboard]{\sffamily\hspace{3mm}\textbf{h\#10.5}\hfill(2+4)}\end{minipage}\hfill
\begin{minipage}[t]{52mm}\vspace{0pt}
\textcolor{ink}{\sffamily\textsc{Lösung}}\par\smallskip\small
\begin{tabular}{@{}r@{\enspace}l@{\enspace}l@{}}
\textbf{1.} & \ldots & \symknight{}f7+\\
\textbf{2.} & \symking{}e7 & \symknight{}e5\\
\textbf{3.} & \symqueen{}d8 & \symknight{}d7\\
\textbf{4.} & \symqueen{}b2+ & \symking{}xb2\\
\textbf{5.} & f5 & \symking{}c3\\
\textbf{6.} & f4 & \symking{}d4\\
\textbf{7.} & f3 & \symking{}e5\\
\textbf{8.} & f2 & \symknight{}f6\\
\textbf{9.} & f1=\symrook{} & \symknight{}h5\\
\textbf{10.} & \symking{}e8 & \symking{}e6\\
\textbf{11.} & \symrook{}f8 & \symknight{}g7\#\\
\end{tabular}\par\vspace{3mm}
\qrcode[height=12mm]{https://helpman.komtera.lt/?fen=7N/8/3k1p2/8/8/8/2qq4/K7&moves=10.5}\\{\sffamily\tiny\color{ink}scannen zum Lösen}
\end{minipage}
\vspace{2mm}

{\raggedright\noindent\textcolor{ink}{\sffamily\textsc{Themen}}\enspace \themetag{reines Matt} \themetag{Mustermatt} \themetag{Idealmatt} \themetag{Umwandlung} \themetag{Unterverwandlung} \themetag{Selbstblock}\par}

\clearpage
\phantomsection\label{hm:KPPvkpp}
\begin{tcolorbox}[enhanced,colback=ink,colframe=ink,arc=0pt,boxrule=0pt,left=4mm,right=4mm,top=2.5mm,bottom=2.5mm]
\color{white}\sffamily{\LARGE\bfseries KPPvkpp}\hfill{\Large h\#10.5}\\[0.5mm]
{\small 6 Steine\quad Weiß am Zug}\hfill{\small Nr.~250\qquad\href{https://helpman.komtera.lt/?fen=8/8/8/8/1p2p3/1P6/1Pk5/K7&moves=10.5}{\color{white}\bfseries lösen\,\ensuremath{\nearrow}}}
\end{tcolorbox}
\vspace{2mm}
\noindent\begin{minipage}[t]{36mm}\vspace{0pt}\sffamily\small\raggedright
\hmfigure{längste eindeutige}{{\Large h\#10.5}}\hmfigure{eindeutig bei dieser Länge}{{\Large 1} Stellung}\hmfigure{längstes Matt}{{\Large h\#15.5}}\hmfigure{in der Klasse}{19,2 Milliarden Stellungen,\\ 3,5 Milliarden mit Hilfsmatt}
\end{minipage}\hfill
\begin{minipage}[t]{80mm}\vspace{0pt}\begin{tikzpicture}\begin{axis}[width=72mm,height=42mm,ybar,bar width=1.6pt,xmin=-1,xmax=32,ymin=0,ymax=608.7,axis y line*=left,axis x line*=bottom,xlabel={\scriptsize Halbzüge bis Matt},ylabel={\scriptsize Stellungen (Mio.)},tick label style={font=\tiny},label style={font=\scriptsize},enlarge x limits=false,ylabel shift=-3pt,xlabel shift=-3pt]\addplot[fill=ink!55,draw=none] coordinates {(0,0.01639) (1,2.021) (2,4.673) (3,30.82) (4,74.8) (5,271.2) (6,415.7) (7,553.3) (8,551.8) (9,486.3) (10,396.3) (11,303.9) (12,246.6) (13,90.22) (14,51.15) (15,6.943) (16,4.621) (17,3.797) (18,3.621) (19,2.607) (20,2.061) (21,0.7993) (22,0.5879) (23,0.1363) (24,0.08835) (25,0.02905) (26,0.0227) (27,0.009137) (28,0.008831) (29,0.00276) (30,0.002739) (31,3e-05)};\draw[densely dotted] (axis cs:21,0) -- (axis cs:21,608.7);\end{axis}\begin{axis}[width=72mm,height=42mm,axis y line*=right,axis x line=none,ymode=log,xmin=-1,xmax=32,ylabel={\scriptsize eindeutig (log)},tick label style={font=\tiny},label style={font=\scriptsize},ylabel shift=-3pt]\addplot[mark=*,mark size=0.7pt,draw=black,thin] coordinates {(0,16392) (1,493103) (2,1189059) (3,9613416) (4,14729222) (5,21478474) (6,13999486) (7,9508459) (8,5336579) (9,3757368) (10,2042487) (11,1367775) (12,650675) (13,31069) (14,5783) (15,1741) (16,736) (17,393) (18,138) (19,48) (20,15) (21,1)};\end{axis}\end{tikzpicture}\end{minipage}
\vspace{3mm}

\noindent\begin{minipage}[t]{60mm}\vspace{0pt}\sbox\hmboard{\chessboard[setfen=8/8/8/8/1p2p3/1P6/1Pk5/K7 w - - 0 1]}\usebox\hmboard\\[1mm]
\makebox[\wd\hmboard]{\sffamily\hspace{3mm}\textbf{h\#10.5}\hfill(3+3)}\end{minipage}\hfill
\begin{minipage}[t]{52mm}\vspace{0pt}
\textcolor{ink}{\sffamily\textsc{Lösung}}\par\smallskip\small
\begin{tabular}{@{}r@{\enspace}l@{\enspace}l@{}}
\textbf{1.} & \ldots & \symking{}a2\\
\textbf{2.} & e3 & \symking{}a1\\
\textbf{3.} & e2 & \symking{}a2\\
\textbf{4.} & e1=\symbishop{} & \symking{}a1\\
\textbf{5.} & \symbishop{}c3 & bxc3\\
\textbf{6.} & bxc3 & b4\\
\textbf{7.} & \symking{}d1 & b5\\
\textbf{8.} & c2 & b6\\
\textbf{9.} & c1=\symbishop{} & b7\\
\textbf{10.} & \symbishop{}d2 & b8=\symqueen{}\\
\textbf{11.} & \symking{}c1 & \symqueen{}b1\#\\
\end{tabular}\par\vspace{3mm}
\qrcode[height=12mm]{https://helpman.komtera.lt/?fen=8/8/8/8/1p2p3/1P6/1Pk5/K7&moves=10.5}\\{\sffamily\tiny\color{ink}scannen zum Lösen}
\end{minipage}
\vspace{2mm}

{\raggedright\noindent\textcolor{ink}{\sffamily\textsc{Themen}}\enspace \themetag{Satzspiel} \themetag{Umwandlung} \themetag{Unterverwandlung} \themetag{Switchback} \themetag{Selbstblock} \themetag{Phönix} \themetag{Pendel} \themetag{ohne Schach} \themetag{Umnov}\par}

\clearpage
\phantomsection\label{hm:KPvkbbp}
\begin{tcolorbox}[enhanced,colback=ink,colframe=ink,arc=0pt,boxrule=0pt,left=4mm,right=4mm,top=2.5mm,bottom=2.5mm]
\color{white}\sffamily{\LARGE\bfseries KPvkbbp}\hfill{\Large h\#10.5}\\[0.5mm]
{\small 6 Steine\quad Weiß am Zug}\hfill{\small Nr.~251\qquad\href{https://helpman.komtera.lt/?fen=8/8/4b3/8/6p1/3b2P1/8/K6k&moves=10.5}{\color{white}\bfseries lösen\,\ensuremath{\nearrow}}}
\end{tcolorbox}
\vspace{2mm}
\noindent\begin{minipage}[t]{36mm}\vspace{0pt}\sffamily\small\raggedright
\hmfigure{längste eindeutige}{{\Large h\#10.5}}\hmfigure{eindeutig bei dieser Länge}{{\Large 41} Stellungen}\hmfigure{längstes Matt}{{\Large h\#15}}\hmfigure{in der Klasse}{34,1 Milliarden Stellungen,\\ 11,6 Milliarden mit Hilfsmatt}
\end{minipage}\hfill
\begin{minipage}[t]{80mm}\vspace{0pt}\begin{tikzpicture}\begin{axis}[width=72mm,height=42mm,ybar,bar width=1.6pt,xmin=-1,xmax=31,ymin=0,ymax=2.195,axis y line*=left,axis x line*=bottom,xlabel={\scriptsize Halbzüge bis Matt},ylabel={\scriptsize Stellungen (Mrd.)},tick label style={font=\tiny},label style={font=\scriptsize},enlarge x limits=false,ylabel shift=-3pt,xlabel shift=-3pt]\addplot[fill=ink!55,draw=none] coordinates {(0,1.128e-05) (1,0.001095) (2,0.003565) (3,0.03001) (4,0.0781) (5,0.3589) (6,0.5678) (7,1.217) (8,1.159) (9,1.759) (10,1.354) (11,1.995) (12,1.459) (13,1.09) (14,0.3313) (15,0.1011) (16,0.02207) (17,0.0146) (18,0.005788) (19,0.006632) (20,0.004733) (21,0.00306) (22,0.001948) (23,0.0009238) (24,0.00053) (25,0.000413) (26,0.0003269) (27,1.657e-05) (28,2.8e-08) (29,1.415e-06) (30,1e-08)};\draw[densely dotted] (axis cs:21,0) -- (axis cs:21,2.195);\end{axis}\begin{axis}[width=72mm,height=42mm,axis y line*=right,axis x line=none,ymode=log,xmin=-1,xmax=31,ylabel={\scriptsize eindeutig (log)},tick label style={font=\tiny},label style={font=\scriptsize},ylabel shift=-3pt]\addplot[mark=*,mark size=0.7pt,draw=black,thin] coordinates {(0,11276) (1,315099) (2,757790) (3,8187959) (4,14541257) (5,34994993) (6,19978709) (7,25081504) (8,12404814) (9,16165020) (10,7540264) (11,9319450) (12,3179083) (13,1112126) (14,163226) (15,43812) (16,3607) (17,2589) (18,590) (19,406) (20,227) (21,41)};\end{axis}\end{tikzpicture}\end{minipage}
\vspace{3mm}

\noindent\begin{minipage}[t]{60mm}\vspace{0pt}\sbox\hmboard{\chessboard[setfen=8/8/4b3/8/6p1/3b2P1/8/K6k w - - 0 1]}\usebox\hmboard\\[1mm]
\makebox[\wd\hmboard]{\sffamily\hspace{3mm}\textbf{h\#10.5}\hfill(2+4)}\end{minipage}\hfill
\begin{minipage}[t]{52mm}\vspace{0pt}
\textcolor{ink}{\sffamily\textsc{Lösung}}\par\smallskip\small
\begin{tabular}{@{}r@{\enspace}l@{\enspace}l@{}}
\textbf{1.} & \ldots & \symking{}b2\\
\textbf{2.} & \symking{}g2 & \symking{}c1\\
\textbf{3.} & \symking{}f3 & \symking{}d1\\
\textbf{4.} & \symking{}e4 & \symking{}d2\\
\textbf{5.} & \symking{}f5 & \symking{}e3\\
\textbf{6.} & \symking{}g6 & \symking{}f4\\
\textbf{7.} & \symbishop{}f7 & \symking{}xg4\\
\textbf{8.} & \symking{}h7 & \symking{}g5\\
\textbf{9.} & \symbishop{}h5 & \symking{}f6\\
\textbf{10.} & \symking{}h6 & g4\\
\textbf{11.} & \symbishop{}h7 & g5\#\\
\end{tabular}\par\vspace{3mm}
\qrcode[height=12mm]{https://helpman.komtera.lt/?fen=8/8/4b3/8/6p1/3b2P1/8/K6k&moves=10.5}\\{\sffamily\tiny\color{ink}scannen zum Lösen}
\end{minipage}
\vspace{2mm}

{\raggedright\noindent\textcolor{ink}{\sffamily\textsc{Themen}}\enspace \themetag{Satzspiel} \themetag{reines Matt} \themetag{Mustermatt} \themetag{Idealmatt} \themetag{Selbstblock} \themetag{ohne Schach}\par}

\clearpage
\phantomsection\label{hm:KBBvkbp}
\begin{tcolorbox}[enhanced,colback=ink,colframe=ink,arc=0pt,boxrule=0pt,left=4mm,right=4mm,top=2.5mm,bottom=2.5mm]
\color{white}\sffamily{\LARGE\bfseries KBBvkbp}\hfill{\Large h\#10}\\[0.5mm]
{\small 6 Steine\quad Schwarz am Zug}\hfill{\small Nr.~252\qquad\href{https://helpman.komtera.lt/?fen=8/8/8/8/8/4p3/2b5/kBKB4&moves=10}{\color{white}\bfseries lösen\,\ensuremath{\nearrow}}}
\end{tcolorbox}
\vspace{2mm}
\noindent\begin{minipage}[t]{36mm}\vspace{0pt}\sffamily\small\raggedright
\hmfigure{längste eindeutige}{{\Large h\#10}}\hmfigure{eindeutig bei dieser Länge}{{\Large 51} Stellungen}\hmfigure{längstes Matt}{{\Large h\#16}}\hmfigure{in der Klasse}{45,4 Milliarden Stellungen,\\ 14,6 Milliarden mit Hilfsmatt}
\end{minipage}\hfill
\begin{minipage}[t]{80mm}\vspace{0pt}\begin{tikzpicture}\begin{axis}[width=72mm,height=42mm,ybar,bar width=1.6pt,xmin=-1,xmax=33,ymin=0,ymax=2.195,axis y line*=left,axis x line*=bottom,xlabel={\scriptsize Halbzüge bis Matt},ylabel={\scriptsize Stellungen (Mrd.)},tick label style={font=\tiny},label style={font=\scriptsize},enlarge x limits=false,ylabel shift=-3pt,xlabel shift=-3pt]\addplot[fill=ink!55,draw=none] coordinates {(0,0.001733) (1,0.006999) (2,0.03415) (3,0.1853) (4,0.6486) (5,1.234) (6,1.996) (7,1.601) (8,1.48) (9,0.9644) (10,0.8746) (11,0.8216) (12,0.8861) (13,0.5699) (14,0.5191) (15,0.4044) (16,0.4064) (17,0.334) (18,0.3157) (19,0.2662) (20,0.2378) (21,0.2131) (22,0.1655) (23,0.1465) (24,0.08118) (25,0.07748) (26,0.03605) (27,0.03375) (28,0.00973) (29,0.01108) (30,3.445e-05) (32,3.19e-07)};\draw[densely dotted] (axis cs:20,0) -- (axis cs:20,2.195);\end{axis}\begin{axis}[width=72mm,height=42mm,axis y line*=right,axis x line=none,ymode=log,xmin=-1,xmax=33,ylabel={\scriptsize eindeutig (log)},tick label style={font=\tiny},label style={font=\scriptsize},ylabel shift=-3pt]\addplot[mark=*,mark size=0.7pt,draw=black,thin] coordinates {(0,1732659) (1,6780358) (2,26249502) (3,95315989) (4,149645141) (5,70950199) (6,40373595) (7,12164359) (8,5817336) (9,3350340) (10,2943601) (11,1660903) (12,1205114) (13,358439) (14,335542) (15,59395) (16,40052) (17,5515) (18,2434) (19,161) (20,51)};\end{axis}\end{tikzpicture}\end{minipage}
\vspace{3mm}

\noindent\begin{minipage}[t]{60mm}\vspace{0pt}\sbox\hmboard{\chessboard[setfen=8/8/8/8/8/4p3/2b5/kBKB4 b - - 0 1]}\usebox\hmboard\\[1mm]
\makebox[\wd\hmboard]{\sffamily\hspace{3mm}\textbf{h\#10}\hfill(3+3)}\end{minipage}\hfill
\begin{minipage}[t]{52mm}\vspace{0pt}
\textcolor{ink}{\sffamily\textsc{Lösung}}\par\smallskip\small
\begin{tabular}{@{}r@{\enspace}l@{\enspace}l@{}}
\textbf{1.} & e2 & \symking{}d2\\
\textbf{2.} & \symking{}b2 & \symking{}e1\\
\textbf{3.} & \symking{}c3 & \symbishop{}dxc2\\
\textbf{4.} & \symking{}d4 & \symbishop{}f5\\
\textbf{5.} & \symking{}e3 & \symbishop{}h3\\
\textbf{6.} & \symking{}f3 & \symbishop{}f1\\
\textbf{7.} & exf1=\symknight{} & \symbishop{}f5\\
\textbf{8.} & \symking{}g2 & \symbishop{}h3+\\
\textbf{9.} & \symking{}h1 & \symking{}f2\\
\textbf{10.} & \symknight{}h2 & \symbishop{}g2\#\\
\end{tabular}\par\vspace{3mm}
\qrcode[height=12mm]{https://helpman.komtera.lt/?fen=8/8/8/8/8/4p3/2b5/kBKB4&moves=10}\\{\sffamily\tiny\color{ink}scannen zum Lösen}
\end{minipage}
\vspace{2mm}

{\raggedright\noindent\textcolor{ink}{\sffamily\textsc{Themen}}\enspace \themetag{reines Matt} \themetag{Mustermatt} \themetag{Idealmatt} \themetag{Umwandlung} \themetag{Unterverwandlung} \themetag{Selbstblock}\par}

\clearpage
\phantomsection\label{hm:KBBvkqp}
\begin{tcolorbox}[enhanced,colback=ink,colframe=ink,arc=0pt,boxrule=0pt,left=4mm,right=4mm,top=2.5mm,bottom=2.5mm]
\color{white}\sffamily{\LARGE\bfseries KBBvkqp}\hfill{\Large h\#10}\\[0.5mm]
{\small 6 Steine\quad Schwarz am Zug}\hfill{\small Nr.~253\qquad\href{https://helpman.komtera.lt/?fen=8/8/3p4/1K1k4/3q1B2/2B5/8/8&moves=10}{\color{white}\bfseries lösen\,\ensuremath{\nearrow}}}
\end{tcolorbox}
\vspace{2mm}
\noindent\begin{minipage}[t]{36mm}\vspace{0pt}\sffamily\small\raggedright
\hmfigure{längste eindeutige}{{\Large h\#10}}\hmfigure{eindeutig bei dieser Länge}{{\Large 7} Stellungen}\hmfigure{längstes Matt}{{\Large h\#17}}\hmfigure{in der Klasse}{45,4 Milliarden Stellungen,\\ 12,8 Milliarden mit Hilfsmatt}
\end{minipage}\hfill
\begin{minipage}[t]{80mm}\vspace{0pt}\begin{tikzpicture}\begin{axis}[width=72mm,height=42mm,ybar,bar width=1.6pt,xmin=-1,xmax=35,ymin=0,ymax=1.762,axis y line*=left,axis x line*=bottom,xlabel={\scriptsize Halbzüge bis Matt},ylabel={\scriptsize Stellungen (Mrd.)},tick label style={font=\tiny},label style={font=\scriptsize},enlarge x limits=false,ylabel shift=-3pt,xlabel shift=-3pt]\addplot[fill=ink!55,draw=none] coordinates {(0,0.0005011) (1,0.002321) (2,0.0124) (3,0.08778) (4,0.2807) (5,0.8614) (6,1.327) (7,1.602) (8,1.282) (9,1.089) (10,0.745) (11,0.7768) (12,0.8067) (13,0.9075) (14,0.7684) (15,0.7782) (16,0.4846) (17,0.4767) (18,0.2097) (19,0.2217) (20,0.04325) (21,0.04854) (22,0.001133) (23,0.004582) (24,6.559e-05) (25,0.002115) (26,5.807e-05) (27,0.0008656) (28,3.358e-05) (29,0.0002766) (30,2.802e-05) (31,1.74e-07) (32,1.261e-05) (33,2.9e-08) (34,5.383e-06)};\draw[densely dotted] (axis cs:20,0) -- (axis cs:20,1.762);\end{axis}\begin{axis}[width=72mm,height=42mm,axis y line*=right,axis x line=none,ymode=log,xmin=-1,xmax=35,ylabel={\scriptsize eindeutig (log)},tick label style={font=\tiny},label style={font=\scriptsize},ylabel shift=-3pt]\addplot[mark=*,mark size=0.7pt,draw=black,thin] coordinates {(0,501113) (1,2302209) (2,7081446) (3,37032027) (4,57873241) (5,75156845) (6,33600162) (7,18900091) (8,7152887) (9,4262390) (10,2774895) (11,2392805) (12,1733823) (13,1111599) (14,526412) (15,233679) (16,63659) (17,23456) (18,2440) (19,561) (20,7)};\end{axis}\end{tikzpicture}\end{minipage}
\vspace{3mm}

\noindent\begin{minipage}[t]{60mm}\vspace{0pt}\sbox\hmboard{\chessboard[setfen=8/8/3p4/1K1k4/3q1B2/2B5/8/8 b - - 0 1]}\usebox\hmboard\\[1mm]
\makebox[\wd\hmboard]{\sffamily\hspace{3mm}\textbf{h\#10}\hfill(3+3)}\end{minipage}\hfill
\begin{minipage}[t]{52mm}\vspace{0pt}
\textcolor{ink}{\sffamily\textsc{Lösung}}\par\smallskip\small
\begin{tabular}{@{}r@{\enspace}l@{\enspace}l@{}}
\textbf{1.} & \symqueen{}a4+ & \symking{}b6\\
\textbf{2.} & \symking{}c4 & \symking{}c7\\
\textbf{3.} & d5 & \symking{}d6\\
\textbf{4.} & d4 & \symking{}e5\\
\textbf{5.} & d3 & \symking{}e4\\
\textbf{6.} & d2 & \symking{}e3\\
\textbf{7.} & d1=\symbishop{} & \symking{}d2\\
\textbf{8.} & \symking{}b3 & \symking{}c1\\
\textbf{9.} & \symking{}a3 & \symking{}b1\\
\textbf{10.} & \symbishop{}b3 & \symbishop{}c1\#\\
\end{tabular}\par\vspace{3mm}
\qrcode[height=12mm]{https://helpman.komtera.lt/?fen=8/8/3p4/1K1k4/3q1B2/2B5/8/8&moves=10}\\{\sffamily\tiny\color{ink}scannen zum Lösen}
\end{minipage}
\vspace{2mm}

{\raggedright\noindent\textcolor{ink}{\sffamily\textsc{Themen}}\enspace \themetag{Umwandlung} \themetag{Unterverwandlung} \themetag{Selbstblock} \themetag{ohne Schlag} \themetag{Umnov}\par}

\clearpage
\phantomsection\label{hm:KBNPvkp}
\begin{tcolorbox}[enhanced,colback=ink,colframe=ink,arc=0pt,boxrule=0pt,left=4mm,right=4mm,top=2.5mm,bottom=2.5mm]
\color{white}\sffamily{\LARGE\bfseries KBNPvkp}\hfill{\Large h\#10}\\[0.5mm]
{\small 6 Steine\quad Schwarz am Zug}\hfill{\small Nr.~254\qquad\href{https://helpman.komtera.lt/?fen=5B2/8/8/7p/7P/1N6/k7/2K5&moves=10}{\color{white}\bfseries lösen\,\ensuremath{\nearrow}}}
\end{tcolorbox}
\vspace{2mm}
\noindent\begin{minipage}[t]{36mm}\vspace{0pt}\sffamily\small\raggedright
\hmfigure{längste eindeutige}{{\Large h\#10}}\hmfigure{eindeutig bei dieser Länge}{{\Large 1} Stellung}\hmfigure{längstes Matt}{{\Large h\#15}}\hmfigure{in der Klasse}{34,1 Milliarden Stellungen,\\ 23,7 Milliarden mit Hilfsmatt}
\end{minipage}\hfill
\begin{minipage}[t]{80mm}\vspace{0pt}\begin{tikzpicture}\begin{axis}[width=72mm,height=42mm,ybar,bar width=1.6pt,xmin=-1,xmax=31,ymin=0,ymax=5.721,axis y line*=left,axis x line*=bottom,xlabel={\scriptsize Halbzüge bis Matt},ylabel={\scriptsize Stellungen (Mrd.)},tick label style={font=\tiny},label style={font=\scriptsize},enlarge x limits=false,ylabel shift=-3pt,xlabel shift=-3pt]\addplot[fill=ink!55,draw=none] coordinates {(0,0.0012) (1,0.01322) (2,0.03821) (3,0.3022) (4,0.8661) (5,2.319) (6,3.621) (7,3.763) (8,5.201) (9,3.282) (10,3.004) (11,0.863) (12,0.4206) (13,0.03806) (14,0.008901) (15,0.000262) (16,0.0001256) (17,1.774e-05) (18,4.382e-05) (19,3.961e-06) (20,1.16e-05) (21,2.62e-07) (22,2.904e-06) (23,1e-09) (24,1.073e-06) (26,1.51e-07) (28,9.3e-08) (30,1e-09)};\draw[densely dotted] (axis cs:20,0) -- (axis cs:20,5.721);\end{axis}\begin{axis}[width=72mm,height=42mm,axis y line*=right,axis x line=none,ymode=log,xmin=-1,xmax=31,ylabel={\scriptsize eindeutig (log)},tick label style={font=\tiny},label style={font=\scriptsize},ylabel shift=-3pt]\addplot[mark=*,mark size=0.7pt,draw=black,thin] coordinates {(0,1199989) (1,6411736) (2,18938468) (3,138739535) (4,210923877) (5,207073833) (6,192086849) (7,85982059) (8,60147752) (9,12262025) (10,4943356) (11,744798) (12,146063) (13,1543) (14,499) (15,41) (16,86) (17,1) (18,5) (20,1)};\end{axis}\end{tikzpicture}\end{minipage}
\vspace{3mm}

\noindent\begin{minipage}[t]{60mm}\vspace{0pt}\sbox\hmboard{\chessboard[setfen=5B2/8/8/7p/7P/1N6/k7/2K5 b - - 0 1]}\usebox\hmboard\\[1mm]
\makebox[\wd\hmboard]{\sffamily\hspace{3mm}\textbf{h\#10}\hfill(4+2)}\end{minipage}\hfill
\begin{minipage}[t]{52mm}\vspace{0pt}
\textcolor{ink}{\sffamily\textsc{Lösung}}\par\smallskip\small
\begin{tabular}{@{}r@{\enspace}l@{\enspace}l@{}}
\textbf{1.} & \symking{}xb3 & \symking{}d2\\
\textbf{2.} & \symking{}c4 & \symking{}e3\\
\textbf{3.} & \symking{}d5 & \symking{}f4\\
\textbf{4.} & \symking{}e6 & \symking{}g5\\
\textbf{5.} & \symking{}f7 & \symking{}xh5\\
\textbf{6.} & \symking{}g8 & \symking{}g6\\
\textbf{7.} & \symking{}h8 & h5\\
\textbf{8.} & \symking{}g8 & h6\\
\textbf{9.} & \symking{}h8 & \symbishop{}g7+\\
\textbf{10.} & \symking{}g8 & h7\#\\
\end{tabular}\par\vspace{3mm}
\qrcode[height=12mm]{https://helpman.komtera.lt/?fen=5B2/8/8/7p/7P/1N6/k7/2K5&moves=10}\\{\sffamily\tiny\color{ink}scannen zum Lösen}
\end{minipage}
\vspace{2mm}

{\raggedright\noindent\textcolor{ink}{\sffamily\textsc{Themen}}\enspace \themetag{reines Matt} \themetag{Mustermatt} \themetag{Idealmatt} \themetag{Switchback} \themetag{ein Stein zieht}\par\smallskip{\footnotesize\textit{Hinweis:} die Lösung beginnt mit einem Schlag.}\par}

\clearpage
\phantomsection\label{hm:KBvkbbn}
\begin{tcolorbox}[enhanced,colback=ink,colframe=ink,arc=0pt,boxrule=0pt,left=4mm,right=4mm,top=2.5mm,bottom=2.5mm]
\color{white}\sffamily{\LARGE\bfseries KBvkbbn}\hfill{\Large h\#10}\\[0.5mm]
{\small 6 Steine\quad Schwarz am Zug}\hfill{\small Nr.~255\qquad\href{https://helpman.komtera.lt/?fen=8/8/8/8/n7/8/2b5/kBK4b&moves=10}{\color{white}\bfseries lösen\,\ensuremath{\nearrow}}}
\end{tcolorbox}
\vspace{2mm}
\noindent\begin{minipage}[t]{36mm}\vspace{0pt}\sffamily\small\raggedright
\hmfigure{längste eindeutige}{{\Large h\#10}}\hmfigure{eindeutig bei dieser Länge}{{\Large 10} Stellungen}\hmfigure{längstes Matt}{{\Large h\#12}}\hmfigure{in der Klasse}{15,5 Milliarden Stellungen,\\ 4,74 Milliarden mit Hilfsmatt}
\end{minipage}\hfill
\begin{minipage}[t]{80mm}\vspace{0pt}\begin{tikzpicture}\begin{axis}[width=72mm,height=42mm,ybar,bar width=1.6pt,xmin=-1,xmax=25,ymin=0,ymax=786.8,axis y line*=left,axis x line*=bottom,xlabel={\scriptsize Halbzüge bis Matt},ylabel={\scriptsize Stellungen (Mio.)},tick label style={font=\tiny},label style={font=\scriptsize},enlarge x limits=false,ylabel shift=-3pt,xlabel shift=-3pt]\addplot[fill=ink!55,draw=none] coordinates {(0,0.02618) (1,0.09999) (2,0.8647) (3,3.079) (4,13.4) (5,35.94) (6,87.05) (7,185) (8,280.8) (9,483.2) (10,491.1) (11,715.3) (12,596.1) (13,622.8) (14,390) (15,354) (16,161.6) (17,147.1) (18,60.88) (19,66.11) (20,18.05) (21,20.79) (22,1.04) (23,1.279) (24,0.000539)};\draw[densely dotted] (axis cs:20,0) -- (axis cs:20,786.8);\end{axis}\begin{axis}[width=72mm,height=42mm,axis y line*=right,axis x line=none,ymode=log,xmin=-1,xmax=25,ylabel={\scriptsize eindeutig (log)},tick label style={font=\tiny},label style={font=\scriptsize},ylabel shift=-3pt]\addplot[mark=*,mark size=0.7pt,draw=black,thin] coordinates {(0,26185) (1,99417) (2,710658) (3,1262942) (4,2963726) (5,4908477) (6,3171045) (7,4250656) (8,2739618) (9,3631247) (10,2387408) (11,2146993) (12,789007) (13,482205) (14,116224) (15,53749) (16,10550) (17,4918) (18,1352) (19,93) (20,10)};\end{axis}\end{tikzpicture}\end{minipage}
\vspace{3mm}

\noindent\begin{minipage}[t]{60mm}\vspace{0pt}\sbox\hmboard{\chessboard[setfen=8/8/8/8/n7/8/2b5/kBK4b b - - 0 1]}\usebox\hmboard\\[1mm]
\makebox[\wd\hmboard]{\sffamily\hspace{3mm}\textbf{h\#10}\hfill(2+4)}\end{minipage}\hfill
\begin{minipage}[t]{52mm}\vspace{0pt}
\textcolor{ink}{\sffamily\textsc{Lösung}}\par\smallskip\small
\begin{tabular}{@{}r@{\enspace}l@{\enspace}l@{}}
\textbf{1.} & \symknight{}c5 & \symking{}xc2\\
\textbf{2.} & \symknight{}a6 & \symking{}d3\\
\textbf{3.} & \symking{}b2 & \symking{}c4\\
\textbf{4.} & \symking{}a3 & \symbishop{}e4\\
\textbf{5.} & \symking{}a4 & \symbishop{}xh1\\
\textbf{6.} & \symking{}a5 & \symbishop{}b7\\
\textbf{7.} & \symking{}b6 & \symbishop{}c8\\
\textbf{8.} & \symking{}a7 & \symking{}b5\\
\textbf{9.} & \symking{}a8 & \symking{}b6\\
\textbf{10.} & \symknight{}b8 & \symbishop{}b7\#\\
\end{tabular}\par\vspace{3mm}
\qrcode[height=12mm]{https://helpman.komtera.lt/?fen=8/8/8/8/n7/8/2b5/kBK4b&moves=10}\\{\sffamily\tiny\color{ink}scannen zum Lösen}
\end{minipage}
\vspace{2mm}

{\raggedright\noindent\textcolor{ink}{\sffamily\textsc{Themen}}\enspace \themetag{reines Matt} \themetag{Mustermatt} \themetag{Idealmatt} \themetag{Switchback} \themetag{Selbstblock} \themetag{ohne Schach}\par}

\clearpage
\phantomsection\label{hm:KBvkqnn}
\begin{tcolorbox}[enhanced,colback=ink,colframe=ink,arc=0pt,boxrule=0pt,left=4mm,right=4mm,top=2.5mm,bottom=2.5mm]
\color{white}\sffamily{\LARGE\bfseries KBvkqnn}\hfill{\Large h\#10}\\[0.5mm]
{\small 6 Steine\quad Schwarz am Zug}\hfill{\small Nr.~256\qquad\href{https://helpman.komtera.lt/?fen=2n5/1n6/2q5/8/8/8/7B/1K5k&moves=10}{\color{white}\bfseries lösen\,\ensuremath{\nearrow}}}
\end{tcolorbox}
\vspace{2mm}
\noindent\begin{minipage}[t]{36mm}\vspace{0pt}\sffamily\small\raggedright
\hmfigure{längste eindeutige}{{\Large h\#10}}\hmfigure{eindeutig bei dieser Länge}{{\Large 5} Stellungen}\hmfigure{längstes Matt}{{\Large h\#13}}\hmfigure{in der Klasse}{15,5 Milliarden Stellungen,\\ 4,35 Milliarden mit Hilfsmatt}
\end{minipage}\hfill
\begin{minipage}[t]{80mm}\vspace{0pt}\begin{tikzpicture}\begin{axis}[width=72mm,height=42mm,ybar,bar width=1.6pt,xmin=-1,xmax=27,ymin=0,ymax=1005,axis y line*=left,axis x line*=bottom,xlabel={\scriptsize Halbzüge bis Matt},ylabel={\scriptsize Stellungen (Mio.)},tick label style={font=\tiny},label style={font=\scriptsize},enlarge x limits=false,ylabel shift=-3pt,xlabel shift=-3pt]\addplot[fill=ink!55,draw=none] coordinates {(0,0.001457) (1,0.003516) (2,0.01219) (3,0.1058) (4,0.4409) (5,2.472) (6,6.698) (7,24.68) (8,48.14) (9,131.3) (10,193.5) (11,433.9) (12,501.9) (13,914) (14,684.6) (15,831.5) (16,258.7) (17,239.7) (18,30.88) (19,41.77) (20,3.042) (21,5.125) (22,0.0547) (23,0.1163) (24,8e-06) (25,0.000539) (26,5e-06)};\draw[densely dotted] (axis cs:20,0) -- (axis cs:20,1005);\end{axis}\begin{axis}[width=72mm,height=42mm,axis y line*=right,axis x line=none,ymode=log,xmin=-1,xmax=27,ylabel={\scriptsize eindeutig (log)},tick label style={font=\tiny},label style={font=\scriptsize},ylabel shift=-3pt]\addplot[mark=*,mark size=0.7pt,draw=black,thin] coordinates {(0,1457) (1,3516) (2,7754) (3,55699) (4,227652) (5,746847) (6,684200) (7,1444098) (8,1308031) (9,2451555) (10,1948182) (11,2711935) (12,1600273) (13,1891677) (14,386523) (15,225653) (16,13572) (17,6230) (18,178) (19,110) (20,5)};\end{axis}\end{tikzpicture}\end{minipage}
\vspace{3mm}

\noindent\begin{minipage}[t]{60mm}\vspace{0pt}\sbox\hmboard{\chessboard[setfen=2n5/1n6/2q5/8/8/8/7B/1K5k b - - 0 1]}\usebox\hmboard\\[1mm]
\makebox[\wd\hmboard]{\sffamily\hspace{3mm}\textbf{h\#10}\hfill(2+4)}\end{minipage}\hfill
\begin{minipage}[t]{52mm}\vspace{0pt}
\textcolor{ink}{\sffamily\textsc{Lösung}}\par\smallskip\small
\begin{tabular}{@{}r@{\enspace}l@{\enspace}l@{}}
\textbf{1.} & \symking{}g2 & \symbishop{}e5\\
\textbf{2.} & \symking{}f3 & \symbishop{}c3\\
\textbf{3.} & \symking{}e4 & \symking{}c2\\
\textbf{4.} & \symking{}d5 & \symking{}d3\\
\textbf{5.} & \symking{}d6 & \symking{}d4\\
\textbf{6.} & \symking{}c7 & \symking{}e5\\
\textbf{7.} & \symqueen{}a4 & \symking{}e6\\
\textbf{8.} & \symqueen{}a8 & \symbishop{}a5+\\
\textbf{9.} & \symking{}b8 & \symking{}d7\\
\textbf{10.} & \symknight{}a7 & \symbishop{}c7\#\\
\end{tabular}\par\vspace{3mm}
\qrcode[height=12mm]{https://helpman.komtera.lt/?fen=2n5/1n6/2q5/8/8/8/7B/1K5k&moves=10}\\{\sffamily\tiny\color{ink}scannen zum Lösen}
\end{minipage}
\vspace{2mm}

{\raggedright\noindent\textcolor{ink}{\sffamily\textsc{Themen}}\enspace \themetag{reines Matt} \themetag{Mustermatt} \themetag{Idealmatt} \themetag{Selbstblock} \themetag{ohne Schlag}\par}

\clearpage
\phantomsection\label{hm:KNPvkpp}
\begin{tcolorbox}[enhanced,colback=ink,colframe=ink,arc=0pt,boxrule=0pt,left=4mm,right=4mm,top=2.5mm,bottom=2.5mm]
\color{white}\sffamily{\LARGE\bfseries KNPvkpp}\hfill{\Large h\#10}\\[0.5mm]
{\small 6 Steine\quad Schwarz am Zug}\hfill{\small Nr.~257\qquad\href{https://helpman.komtera.lt/?fen=6N1/7k/8/2p5/8/2p5/2P5/1K6&moves=10}{\color{white}\bfseries lösen\,\ensuremath{\nearrow}}}
\end{tcolorbox}
\vspace{2mm}
\noindent\begin{minipage}[t]{36mm}\vspace{0pt}\sffamily\small\raggedright
\hmfigure{längste eindeutige}{{\Large h\#10}}\hmfigure{eindeutig bei dieser Länge}{{\Large 3} Stellungen}\hmfigure{längstes Matt}{{\Large h\#15}}\hmfigure{in der Klasse}{25,6 Milliarden Stellungen,\\ 9,24 Milliarden mit Hilfsmatt}
\end{minipage}\hfill
\begin{minipage}[t]{80mm}\vspace{0pt}\begin{tikzpicture}\begin{axis}[width=72mm,height=42mm,ybar,bar width=1.6pt,xmin=-1,xmax=31,ymin=0,ymax=1.313,axis y line*=left,axis x line*=bottom,xlabel={\scriptsize Halbzüge bis Matt},ylabel={\scriptsize Stellungen (Mrd.)},tick label style={font=\tiny},label style={font=\scriptsize},enlarge x limits=false,ylabel shift=-3pt,xlabel shift=-3pt]\addplot[fill=ink!55,draw=none] coordinates {(0,4.29e-05) (1,0.002919) (2,0.007091) (3,0.05758) (4,0.1486) (5,0.5582) (6,0.826) (7,1.006) (8,1.098) (9,1.089) (10,1.194) (11,1.17) (12,1.104) (13,0.494) (14,0.3047) (15,0.09278) (16,0.05696) (17,0.01782) (18,0.008084) (19,0.001118) (20,0.0003794) (21,1.843e-05) (22,5.697e-06) (23,1.1e-08) (24,1.76e-07) (25,1e-09) (26,7.2e-08) (28,3e-08) (30,4e-09)};\draw[densely dotted] (axis cs:20,0) -- (axis cs:20,1.313);\end{axis}\begin{axis}[width=72mm,height=42mm,axis y line*=right,axis x line=none,ymode=log,xmin=-1,xmax=31,ylabel={\scriptsize eindeutig (log)},tick label style={font=\tiny},label style={font=\scriptsize},ylabel shift=-3pt]\addplot[mark=*,mark size=0.7pt,draw=black,thin] coordinates {(0,42899) (1,815057) (2,2067204) (3,20845779) (4,29968197) (5,40158584) (6,27474599) (7,21319371) (8,14409033) (9,11661608) (10,7449566) (11,5989611) (12,2626175) (13,298927) (14,173865) (15,55618) (16,14230) (17,1822) (18,416) (19,17) (20,3)};\end{axis}\end{tikzpicture}\end{minipage}
\vspace{3mm}

\noindent\begin{minipage}[t]{60mm}\vspace{0pt}\sbox\hmboard{\chessboard[setfen=6N1/7k/8/2p5/8/2p5/2P5/1K6 b - - 0 1]}\usebox\hmboard\\[1mm]
\makebox[\wd\hmboard]{\sffamily\hspace{3mm}\textbf{h\#10}\hfill(3+3)}\end{minipage}\hfill
\begin{minipage}[t]{52mm}\vspace{0pt}
\textcolor{ink}{\sffamily\textsc{Lösung}}\par\smallskip\small
\begin{tabular}{@{}r@{\enspace}l@{\enspace}l@{}}
\textbf{1.} & \symking{}g6 & \symknight{}f6\\
\textbf{2.} & \symking{}f5 & \symknight{}e4\\
\textbf{3.} & \symking{}e5 & \symknight{}d2\\
\textbf{4.} & cxd2 & c4\\
\textbf{5.} & d1=\symrook{}+ & \symking{}c2\\
\textbf{6.} & \symrook{}d5 & cxd5\\
\textbf{7.} & \symking{}d4 & d6\\
\textbf{8.} & \symking{}c4 & d7\\
\textbf{9.} & \symking{}b4 & d8=\symqueen{}\\
\textbf{10.} & \symking{}a3 & \symqueen{}a5\#\\
\end{tabular}\par\vspace{3mm}
\qrcode[height=12mm]{https://helpman.komtera.lt/?fen=6N1/7k/8/2p5/8/2p5/2P5/1K6&moves=10}\\{\sffamily\tiny\color{ink}scannen zum Lösen}
\end{minipage}
\vspace{2mm}

{\raggedright\noindent\textcolor{ink}{\sffamily\textsc{Themen}}\enspace \themetag{reines Matt} \themetag{Mustermatt} \themetag{Spiegelmatt} \themetag{Umwandlung} \themetag{Unterverwandlung} \themetag{Excelsior}\par}

\clearpage
\phantomsection\label{hm:KQBvkbp}
\begin{tcolorbox}[enhanced,colback=ink,colframe=ink,arc=0pt,boxrule=0pt,left=4mm,right=4mm,top=2.5mm,bottom=2.5mm]
\color{white}\sffamily{\LARGE\bfseries KQBvkbp}\hfill{\Large h\#10}\\[0.5mm]
{\small 6 Steine\quad Schwarz am Zug}\hfill{\small Nr.~258\qquad\href{https://helpman.komtera.lt/?fen=8/K7/8/kQ6/8/5p2/B7/7b&moves=10}{\color{white}\bfseries lösen\,\ensuremath{\nearrow}}}
\end{tcolorbox}
\vspace{2mm}
\noindent\begin{minipage}[t]{36mm}\vspace{0pt}\sffamily\small\raggedright
\hmfigure{längste eindeutige}{{\Large h\#10}}\hmfigure{eindeutig bei dieser Länge}{{\Large 5} Stellungen}\hmfigure{längstes Matt}{{\Large h\#16.5}}\hmfigure{in der Klasse}{45,4 Milliarden Stellungen,\\ 26 Milliarden mit Hilfsmatt}
\end{minipage}\hfill
\begin{minipage}[t]{80mm}\vspace{0pt}\begin{tikzpicture}\begin{axis}[width=72mm,height=42mm,ybar,bar width=1.6pt,xmin=-1,xmax=34,ymin=0,ymax=12.86,axis y line*=left,axis x line*=bottom,xlabel={\scriptsize Halbzüge bis Matt},ylabel={\scriptsize Stellungen (Mrd.)},tick label style={font=\tiny},label style={font=\scriptsize},enlarge x limits=false,ylabel shift=-3pt,xlabel shift=-3pt]\addplot[fill=ink!55,draw=none] coordinates {(0,0.06152) (1,0.393) (2,1.707) (3,5.607) (4,11.69) (5,4.61) (6,1.665) (7,0.04522) (8,0.02839) (9,8.24e-07) (10,0.03416) (11,8.13e-07) (12,0.04809) (13,1.668e-06) (14,0.03531) (15,1.154e-06) (16,0.02694) (17,2.528e-06) (18,0.02133) (19,1.92e-07) (20,0.01723) (21,1.75e-07) (22,0.017) (23,1.31e-07) (24,0.01344) (25,1.2e-08) (26,0.009352) (27,3.8e-08) (28,0.006204) (29,8e-09) (30,0.002798) (32,0.0006389) (33,2e-09)};\draw[densely dotted] (axis cs:20,0) -- (axis cs:20,12.86);\end{axis}\begin{axis}[width=72mm,height=42mm,axis y line*=right,axis x line=none,ymode=log,xmin=-1,xmax=34,ylabel={\scriptsize eindeutig (log)},tick label style={font=\tiny},label style={font=\scriptsize},ylabel shift=-3pt]\addplot[mark=*,mark size=0.7pt,draw=black,thin] coordinates {(0,61516468) (1,328917810) (2,1026709073) (3,1160570089) (4,585793293) (5,8648340) (6,894020) (7,99) (8,357706) (9,11) (10,142065) (12,58961) (14,8377) (16,1471) (18,146) (20,5)};\end{axis}\end{tikzpicture}\end{minipage}
\vspace{3mm}

\noindent\begin{minipage}[t]{60mm}\vspace{0pt}\sbox\hmboard{\chessboard[setfen=8/K7/8/kQ6/8/5p2/B7/7b b - - 0 1]}\usebox\hmboard\\[1mm]
\makebox[\wd\hmboard]{\sffamily\hspace{3mm}\textbf{h\#10}\hfill(3+3)}\end{minipage}\hfill
\begin{minipage}[t]{52mm}\vspace{0pt}
\textcolor{ink}{\sffamily\textsc{Lösung}}\par\smallskip\small
\begin{tabular}{@{}r@{\enspace}l@{\enspace}l@{}}
\textbf{1.} & \symking{}xb5 & \symbishop{}d5\\
\textbf{2.} & f2 & \symbishop{}xh1\\
\textbf{3.} & \symking{}c4 & \symking{}b6\\
\textbf{4.} & \symking{}d3 & \symking{}c5\\
\textbf{5.} & \symking{}e2 & \symking{}d4\\
\textbf{6.} & \symking{}f1 & \symbishop{}g2+\\
\textbf{7.} & \symking{}g1 & \symbishop{}h3\\
\textbf{8.} & \symking{}h1 & \symking{}e3\\
\textbf{9.} & f1=\symknight{}+ & \symking{}f2\\
\textbf{10.} & \symknight{}h2 & \symbishop{}g2\#\\
\end{tabular}\par\vspace{3mm}
\qrcode[height=12mm]{https://helpman.komtera.lt/?fen=8/K7/8/kQ6/8/5p2/B7/7b&moves=10}\\{\sffamily\tiny\color{ink}scannen zum Lösen}
\end{minipage}
\vspace{2mm}

{\raggedright\noindent\textcolor{ink}{\sffamily\textsc{Themen}}\enspace \themetag{reines Matt} \themetag{Mustermatt} \themetag{Idealmatt} \themetag{Umwandlung} \themetag{Unterverwandlung} \themetag{Switchback} \themetag{Selbstblock} \themetag{Kniest} \themetag{Umnov}\par\smallskip{\footnotesize\textit{Hinweis:} die Partei am Zug steht im Diagramm im Schach; die Lösung beginnt mit einem Schlag.}\par}

\clearpage
\phantomsection\label{hm:KQBvknp}
\begin{tcolorbox}[enhanced,colback=ink,colframe=ink,arc=0pt,boxrule=0pt,left=4mm,right=4mm,top=2.5mm,bottom=2.5mm]
\color{white}\sffamily{\LARGE\bfseries KQBvknp}\hfill{\Large h\#10}\\[0.5mm]
{\small 6 Steine\quad Schwarz am Zug}\hfill{\small Nr.~259\qquad\href{https://helpman.komtera.lt/?fen=Q7/k1B2p2/8/5n2/8/8/K7/8&moves=10}{\color{white}\bfseries lösen\,\ensuremath{\nearrow}}}
\end{tcolorbox}
\vspace{2mm}
\noindent\begin{minipage}[t]{36mm}\vspace{0pt}\sffamily\small\raggedright
\hmfigure{längste eindeutige}{{\Large h\#10}}\hmfigure{eindeutig bei dieser Länge}{{\Large 2} Stellungen}\hmfigure{längstes Matt}{{\Large h\#16}}\hmfigure{in der Klasse}{45,4 Milliarden Stellungen,\\ 26,9 Milliarden mit Hilfsmatt}
\end{minipage}\hfill
\begin{minipage}[t]{80mm}\vspace{0pt}\begin{tikzpicture}\begin{axis}[width=72mm,height=42mm,ybar,bar width=1.6pt,xmin=-1,xmax=33,ymin=0,ymax=13.12,axis y line*=left,axis x line*=bottom,xlabel={\scriptsize Halbzüge bis Matt},ylabel={\scriptsize Stellungen (Mrd.)},tick label style={font=\tiny},label style={font=\scriptsize},enlarge x limits=false,ylabel shift=-3pt,xlabel shift=-3pt]\addplot[fill=ink!55,draw=none] coordinates {(0,0.06571) (1,0.4073) (2,1.641) (3,5.504) (4,11.93) (5,4.685) (6,2.275) (7,0.05871) (8,0.01953) (9,4.19e-07) (10,0.02243) (11,2.1e-07) (12,0.04022) (13,1.63e-06) (14,0.05718) (15,1.113e-06) (16,0.06282) (17,1.333e-06) (18,0.04449) (19,2.681e-06) (20,0.02754) (21,2.249e-06) (22,0.007155) (23,1.663e-06) (24,0.0003115) (25,1.254e-06) (26,5.3e-08) (27,1.407e-06) (28,5.9e-08) (29,4.65e-07) (30,2.4e-08) (31,2.96e-07) (32,1.8e-08)};\draw[densely dotted] (axis cs:20,0) -- (axis cs:20,13.12);\end{axis}\begin{axis}[width=72mm,height=42mm,axis y line*=right,axis x line=none,ymode=log,xmin=-1,xmax=33,ylabel={\scriptsize eindeutig (log)},tick label style={font=\tiny},label style={font=\scriptsize},ylabel shift=-3pt]\addplot[mark=*,mark size=0.7pt,draw=black,thin] coordinates {(0,65714828) (1,334897508) (2,962905247) (3,1186784238) (4,647991945) (5,9285707) (6,805455) (7,23) (8,326790) (9,6) (10,236428) (12,123489) (13,242) (14,28400) (15,3) (16,5365) (18,132) (20,2)};\end{axis}\end{tikzpicture}\end{minipage}
\vspace{3mm}

\noindent\begin{minipage}[t]{60mm}\vspace{0pt}\sbox\hmboard{\chessboard[setfen=Q7/k1B2p2/8/5n2/8/8/K7/8 b - - 0 1]}\usebox\hmboard\\[1mm]
\makebox[\wd\hmboard]{\sffamily\hspace{3mm}\textbf{h\#10}\hfill(3+3)}\end{minipage}\hfill
\begin{minipage}[t]{52mm}\vspace{0pt}
\textcolor{ink}{\sffamily\textsc{Lösung}}\par\smallskip\small
\begin{tabular}{@{}r@{\enspace}l@{\enspace}l@{}}
\textbf{1.} & \symking{}xa8 & \symking{}b3\\
\textbf{2.} & \symking{}b7 & \symking{}c4\\
\textbf{3.} & \symking{}c8 & \symking{}d5\\
\textbf{4.} & \symking{}d7 & \symking{}e5\\
\textbf{5.} & \symking{}e8 & \symking{}f6\\
\textbf{6.} & \symking{}f8 & \symbishop{}d6+\\
\textbf{7.} & \symking{}g8 & \symbishop{}f8\\
\textbf{8.} & \symking{}h8 & \symking{}xf7\\
\textbf{9.} & \symknight{}h6+ & \symking{}g6\\
\textbf{10.} & \symknight{}g8 & \symbishop{}g7\#\\
\end{tabular}\par\vspace{3mm}
\qrcode[height=12mm]{https://helpman.komtera.lt/?fen=Q7/k1B2p2/8/5n2/8/8/K7/8&moves=10}\\{\sffamily\tiny\color{ink}scannen zum Lösen}
\end{minipage}
\vspace{2mm}

{\raggedright\noindent\textcolor{ink}{\sffamily\textsc{Themen}}\enspace \themetag{reines Matt} \themetag{Mustermatt} \themetag{Idealmatt} \themetag{Selbstblock} \themetag{Umnov}\par\smallskip{\footnotesize\textit{Hinweis:} die Partei am Zug steht im Diagramm im Schach; die Lösung beginnt mit einem Schlag.}\par}

\clearpage
\phantomsection\label{hm:KQBvkqn}
\begin{tcolorbox}[enhanced,colback=ink,colframe=ink,arc=0pt,boxrule=0pt,left=4mm,right=4mm,top=2.5mm,bottom=2.5mm]
\color{white}\sffamily{\LARGE\bfseries KQBvkqn}\hfill{\Large h\#10}\\[0.5mm]
{\small 6 Steine\quad Schwarz am Zug}\hfill{\small Nr.~260\qquad\href{https://helpman.komtera.lt/?fen=8/8/5q2/8/8/kQ6/8/1K1n1B2&moves=10}{\color{white}\bfseries lösen\,\ensuremath{\nearrow}}}
\end{tcolorbox}
\vspace{2mm}
\noindent\begin{minipage}[t]{36mm}\vspace{0pt}\sffamily\small\raggedright
\hmfigure{längste eindeutige}{{\Large h\#10}}\hmfigure{eindeutig bei dieser Länge}{{\Large 17} Stellungen}\hmfigure{längstes Matt}{{\Large h\#13}}\hmfigure{in der Klasse}{15,5 Milliarden Stellungen,\\ 7,28 Milliarden mit Hilfsmatt}
\end{minipage}\hfill
\begin{minipage}[t]{80mm}\vspace{0pt}\begin{tikzpicture}\begin{axis}[width=72mm,height=42mm,ybar,bar width=1.6pt,xmin=-1,xmax=27,ymin=0,ymax=3.096,axis y line*=left,axis x line*=bottom,xlabel={\scriptsize Halbzüge bis Matt},ylabel={\scriptsize Stellungen (Mrd.)},tick label style={font=\tiny},label style={font=\scriptsize},enlarge x limits=false,ylabel shift=-3pt,xlabel shift=-3pt]\addplot[fill=ink!55,draw=none] coordinates {(0,0.009229) (1,0.06587) (2,0.3372) (3,1.371) (4,2.815) (5,2.035) (6,0.5252) (7,0.05624) (8,0.003807) (9,7.599e-06) (10,0.004002) (11,6.26e-07) (12,0.007342) (13,2.324e-06) (14,0.01304) (15,6.68e-07) (16,0.01455) (17,5.43e-07) (18,0.01237) (19,8.47e-07) (20,0.008905) (21,2.587e-06) (22,0.003188) (23,9.9e-08) (24,0.0003484) (25,7e-08) (26,3.018e-06)};\draw[densely dotted] (axis cs:20,0) -- (axis cs:20,3.096);\end{axis}\begin{axis}[width=72mm,height=42mm,axis y line*=right,axis x line=none,ymode=log,xmin=-1,xmax=27,ylabel={\scriptsize eindeutig (log)},tick label style={font=\tiny},label style={font=\scriptsize},ylabel shift=-3pt]\addplot[mark=*,mark size=0.7pt,draw=black,thin] coordinates {(0,9229340) (1,58989984) (2,198272824) (3,296976209) (4,156546402) (5,11862052) (6,197365) (7,238) (8,75123) (9,15) (10,44512) (11,186) (12,31243) (14,16965) (16,4192) (18,662) (20,17)};\end{axis}\end{tikzpicture}\end{minipage}
\vspace{3mm}

\noindent\begin{minipage}[t]{60mm}\vspace{0pt}\sbox\hmboard{\chessboard[setfen=8/8/5q2/8/8/kQ6/8/1K1n1B2 b - - 0 1]}\usebox\hmboard\\[1mm]
\makebox[\wd\hmboard]{\sffamily\hspace{3mm}\textbf{h\#10}\hfill(3+3)}\end{minipage}\hfill
\begin{minipage}[t]{52mm}\vspace{0pt}
\textcolor{ink}{\sffamily\textsc{Lösung}}\par\smallskip\small
\begin{tabular}{@{}r@{\enspace}l@{\enspace}l@{}}
\textbf{1.} & \symking{}xb3 & \symking{}c1\\
\textbf{2.} & \symknight{}f2 & \symking{}d2\\
\textbf{3.} & \symknight{}h3 & \symking{}e3\\
\textbf{4.} & \symking{}c2 & \symbishop{}d3+\\
\textbf{5.} & \symking{}d1 & \symbishop{}f5\\
\textbf{6.} & \symking{}e1 & \symking{}f3\\
\textbf{7.} & \symking{}f1 & \symking{}g3\\
\textbf{8.} & \symking{}g1 & \symbishop{}d3\\
\textbf{9.} & \symking{}h1 & \symbishop{}f1\\
\textbf{10.} & \symknight{}g1 & \symbishop{}g2\#\\
\end{tabular}\par\vspace{3mm}
\qrcode[height=12mm]{https://helpman.komtera.lt/?fen=8/8/5q2/8/8/kQ6/8/1K1n1B2&moves=10}\\{\sffamily\tiny\color{ink}scannen zum Lösen}
\end{minipage}
\vspace{2mm}

{\raggedright\noindent\textcolor{ink}{\sffamily\textsc{Themen}}\enspace \themetag{reines Matt} \themetag{Mustermatt} \themetag{Switchback} \themetag{Selbstblock}\par\smallskip{\footnotesize\textit{Hinweis:} die Partei am Zug steht im Diagramm im Schach; die Lösung beginnt mit einem Schlag.}\par}

\clearpage
\phantomsection\label{hm:KQBvkrn}
\begin{tcolorbox}[enhanced,colback=ink,colframe=ink,arc=0pt,boxrule=0pt,left=4mm,right=4mm,top=2.5mm,bottom=2.5mm]
\color{white}\sffamily{\LARGE\bfseries KQBvkrn}\hfill{\Large h\#10}\\[0.5mm]
{\small 6 Steine\quad Schwarz am Zug}\hfill{\small Nr.~261\qquad\href{https://helpman.komtera.lt/?fen=8/8/5r2/8/8/kQ6/4B3/1K1n4&moves=10}{\color{white}\bfseries lösen\,\ensuremath{\nearrow}}}
\end{tcolorbox}
\vspace{2mm}
\noindent\begin{minipage}[t]{36mm}\vspace{0pt}\sffamily\small\raggedright
\hmfigure{längste eindeutige}{{\Large h\#10}}\hmfigure{eindeutig bei dieser Länge}{{\Large 16} Stellungen}\hmfigure{längstes Matt}{{\Large h\#13}}\hmfigure{in der Klasse}{15,5 Milliarden Stellungen,\\ 7,99 Milliarden mit Hilfsmatt}
\end{minipage}\hfill
\begin{minipage}[t]{80mm}\vspace{0pt}\begin{tikzpicture}\begin{axis}[width=72mm,height=42mm,ybar,bar width=1.6pt,xmin=-1,xmax=27,ymin=0,ymax=3.72,axis y line*=left,axis x line*=bottom,xlabel={\scriptsize Halbzüge bis Matt},ylabel={\scriptsize Stellungen (Mrd.)},tick label style={font=\tiny},label style={font=\scriptsize},enlarge x limits=false,ylabel shift=-3pt,xlabel shift=-3pt]\addplot[fill=ink!55,draw=none] coordinates {(0,0.01322) (1,0.09174) (2,0.4465) (3,1.599) (4,3.382) (5,1.813) (6,0.5351) (7,0.02872) (8,0.004684) (9,2.24e-07) (10,0.0044) (11,2.52e-07) (12,0.008959) (13,5.74e-07) (14,0.01559) (15,9.2e-08) (16,0.01763) (17,8.5e-08) (18,0.01459) (19,4.59e-07) (20,0.009892) (21,4e-08) (22,0.003027) (23,6e-09) (24,0.0002108) (25,2e-09) (26,3.25e-07)};\draw[densely dotted] (axis cs:20,0) -- (axis cs:20,3.72);\end{axis}\begin{axis}[width=72mm,height=42mm,axis y line*=right,axis x line=none,ymode=log,xmin=-1,xmax=27,ylabel={\scriptsize eindeutig (log)},tick label style={font=\tiny},label style={font=\scriptsize},ylabel shift=-3pt]\addplot[mark=*,mark size=0.7pt,draw=black,thin] coordinates {(0,13220481) (1,80035387) (2,264483873) (3,324367408) (4,176508354) (5,6375230) (6,197044) (7,26) (8,67082) (10,46921) (11,154) (12,30576) (14,14273) (16,3561) (18,584) (20,16)};\end{axis}\end{tikzpicture}\end{minipage}
\vspace{3mm}

\noindent\begin{minipage}[t]{60mm}\vspace{0pt}\sbox\hmboard{\chessboard[setfen=8/8/5r2/8/8/kQ6/4B3/1K1n4 b - - 0 1]}\usebox\hmboard\\[1mm]
\makebox[\wd\hmboard]{\sffamily\hspace{3mm}\textbf{h\#10}\hfill(3+3)}\end{minipage}\hfill
\begin{minipage}[t]{52mm}\vspace{0pt}
\textcolor{ink}{\sffamily\textsc{Lösung}}\par\smallskip\small
\begin{tabular}{@{}r@{\enspace}l@{\enspace}l@{}}
\textbf{1.} & \symking{}xb3 & \symking{}c1\\
\textbf{2.} & \symknight{}f2 & \symking{}d2\\
\textbf{3.} & \symknight{}h3 & \symking{}e3\\
\textbf{4.} & \symking{}c2 & \symbishop{}d3+\\
\textbf{5.} & \symking{}d1 & \symbishop{}f5\\
\textbf{6.} & \symking{}e1 & \symking{}f3\\
\textbf{7.} & \symking{}f1 & \symking{}g3\\
\textbf{8.} & \symking{}g1 & \symbishop{}d3\\
\textbf{9.} & \symking{}h1 & \symbishop{}f1\\
\textbf{10.} & \symknight{}g1 & \symbishop{}g2\#\\
\end{tabular}\par\vspace{3mm}
\qrcode[height=12mm]{https://helpman.komtera.lt/?fen=8/8/5r2/8/8/kQ6/4B3/1K1n4&moves=10}\\{\sffamily\tiny\color{ink}scannen zum Lösen}
\end{minipage}
\vspace{2mm}

{\raggedright\noindent\textcolor{ink}{\sffamily\textsc{Themen}}\enspace \themetag{reines Matt} \themetag{Mustermatt} \themetag{Switchback} \themetag{Selbstblock}\par\smallskip{\footnotesize\textit{Hinweis:} die Partei am Zug steht im Diagramm im Schach; die Lösung beginnt mit einem Schlag.}\par}

\clearpage
\phantomsection\label{hm:KQNvkpp}
\begin{tcolorbox}[enhanced,colback=ink,colframe=ink,arc=0pt,boxrule=0pt,left=4mm,right=4mm,top=2.5mm,bottom=2.5mm]
\color{white}\sffamily{\LARGE\bfseries KQNvkpp}\hfill{\Large h\#10}\\[0.5mm]
{\small 6 Steine\quad Schwarz am Zug}\hfill{\small Nr.~262\qquad\href{https://helpman.komtera.lt/?fen=8/6p1/1K1kQ1p1/8/8/8/8/2N5&moves=10}{\color{white}\bfseries lösen\,\ensuremath{\nearrow}}}
\end{tcolorbox}
\vspace{2mm}
\noindent\begin{minipage}[t]{36mm}\vspace{0pt}\sffamily\small\raggedright
\hmfigure{längste eindeutige}{{\Large h\#10}}\hmfigure{eindeutig bei dieser Länge}{{\Large 1} Stellung}\hmfigure{längstes Matt}{{\Large h\#11}}\hmfigure{in der Klasse}{34,1 Milliarden Stellungen,\\ 10,6 Milliarden mit Hilfsmatt}
\end{minipage}\hfill
\begin{minipage}[t]{80mm}\vspace{0pt}\begin{tikzpicture}\begin{axis}[width=72mm,height=42mm,ybar,bar width=1.6pt,xmin=-1,xmax=23,ymin=0,ymax=3.88,axis y line*=left,axis x line*=bottom,xlabel={\scriptsize Halbzüge bis Matt},ylabel={\scriptsize Stellungen (Mrd.)},tick label style={font=\tiny},label style={font=\scriptsize},enlarge x limits=false,ylabel shift=-3pt,xlabel shift=-3pt]\addplot[fill=ink!55,draw=none] coordinates {(0,0.01926) (1,0.1383) (2,0.4893) (3,1.511) (4,3.528) (5,2.289) (6,2.169) (7,0.2946) (8,0.1042) (9,0.001032) (10,0.0096) (11,6.5e-08) (12,0.01913) (13,1.9e-08) (14,0.02518) (15,4e-09) (16,0.02016) (18,0.008737) (19,4e-09) (20,0.002585) (22,0.0004202)};\draw[densely dotted] (axis cs:20,0) -- (axis cs:20,3.88);\end{axis}\begin{axis}[width=72mm,height=42mm,axis y line*=right,axis x line=none,ymode=log,xmin=-1,xmax=23,ylabel={\scriptsize eindeutig (log)},tick label style={font=\tiny},label style={font=\scriptsize},ylabel shift=-3pt]\addplot[mark=*,mark size=0.7pt,draw=black,thin] coordinates {(0,19255478) (1,109677812) (2,296403027) (3,284930232) (4,269031402) (5,26517713) (6,7560222) (7,65303) (8,72909) (9,1) (10,62372) (12,34309) (14,4537) (16,342) (18,9) (20,1)};\end{axis}\end{tikzpicture}\end{minipage}
\vspace{3mm}

\noindent\begin{minipage}[t]{60mm}\vspace{0pt}\sbox\hmboard{\chessboard[setfen=8/6p1/1K1kQ1p1/8/8/8/8/2N5 b - - 0 1]}\usebox\hmboard\\[1mm]
\makebox[\wd\hmboard]{\sffamily\hspace{3mm}\textbf{h\#10}\hfill(3+3)}\end{minipage}\hfill
\begin{minipage}[t]{52mm}\vspace{0pt}
\textcolor{ink}{\sffamily\textsc{Lösung}}\par\smallskip\small
\begin{tabular}{@{}r@{\enspace}l@{\enspace}l@{}}
\textbf{1.} & \symking{}xe6 & \symking{}c7\\
\textbf{2.} & g5 & \symking{}d8\\
\textbf{3.} & g4 & \symking{}e8\\
\textbf{4.} & g3 & \symking{}f8\\
\textbf{5.} & g2 & \symking{}xg7\\
\textbf{6.} & g1=\symrook{}+ & \symking{}h7\\
\textbf{7.} & \symking{}f7 & \symking{}h6\\
\textbf{8.} & \symking{}g8 & \symknight{}d3\\
\textbf{9.} & \symking{}h8 & \symknight{}e5\\
\textbf{10.} & \symrook{}g8 & \symknight{}f7\#\\
\end{tabular}\par\vspace{3mm}
\qrcode[height=12mm]{https://helpman.komtera.lt/?fen=8/6p1/1K1kQ1p1/8/8/8/8/2N5&moves=10}\\{\sffamily\tiny\color{ink}scannen zum Lösen}
\end{minipage}
\vspace{2mm}

{\raggedright\noindent\textcolor{ink}{\sffamily\textsc{Themen}}\enspace \themetag{reines Matt} \themetag{Mustermatt} \themetag{Idealmatt} \themetag{Umwandlung} \themetag{Unterverwandlung} \themetag{Selbstblock}\par\smallskip{\footnotesize\textit{Hinweis:} die Partei am Zug steht im Diagramm im Schach; die Lösung beginnt mit einem Schlag.}\par}

\clearpage
\phantomsection\label{hm:KQNvkqp}
\begin{tcolorbox}[enhanced,colback=ink,colframe=ink,arc=0pt,boxrule=0pt,left=4mm,right=4mm,top=2.5mm,bottom=2.5mm]
\color{white}\sffamily{\LARGE\bfseries KQNvkqp}\hfill{\Large h\#10}\\[0.5mm]
{\small 6 Steine\quad Schwarz am Zug}\hfill{\small Nr.~263\qquad\href{https://helpman.komtera.lt/?fen=8/8/4p3/4Q3/3k4/8/1Nq5/K7&moves=10}{\color{white}\bfseries lösen\,\ensuremath{\nearrow}}}
\end{tcolorbox}
\vspace{2mm}
\noindent\begin{minipage}[t]{36mm}\vspace{0pt}\sffamily\small\raggedright
\hmfigure{längste eindeutige}{{\Large h\#10}}\hmfigure{eindeutig bei dieser Länge}{{\Large 35} Stellungen}\hmfigure{längstes Matt}{{\Large h\#12}}\hmfigure{in der Klasse}{45,4 Milliarden Stellungen,\\ 23,2 Milliarden mit Hilfsmatt}
\end{minipage}\hfill
\begin{minipage}[t]{80mm}\vspace{0pt}\begin{tikzpicture}\begin{axis}[width=72mm,height=42mm,ybar,bar width=1.6pt,xmin=-1,xmax=25,ymin=0,ymax=7.576,axis y line*=left,axis x line*=bottom,xlabel={\scriptsize Halbzüge bis Matt},ylabel={\scriptsize Stellungen (Mrd.)},tick label style={font=\tiny},label style={font=\scriptsize},enlarge x limits=false,ylabel shift=-3pt,xlabel shift=-3pt]\addplot[fill=ink!55,draw=none] coordinates {(0,0.02076) (1,0.1748) (2,0.8173) (3,3.198) (4,6.887) (5,6.76) (6,3.894) (7,1.152) (8,0.1385) (9,0.01117) (10,0.01927) (11,2.894e-06) (12,0.03692) (13,9.23e-07) (14,0.04266) (15,1.464e-06) (16,0.03585) (17,2.08e-07) (18,0.02382) (19,1.59e-07) (20,0.005797) (21,3e-08) (22,0.0003007) (23,1.3e-08) (24,1.8e-08)};\draw[densely dotted] (axis cs:20,0) -- (axis cs:20,7.576);\end{axis}\begin{axis}[width=72mm,height=42mm,axis y line*=right,axis x line=none,ymode=log,xmin=-1,xmax=25,ylabel={\scriptsize eindeutig (log)},tick label style={font=\tiny},label style={font=\scriptsize},ylabel shift=-3pt]\addplot[mark=*,mark size=0.7pt,draw=black,thin] coordinates {(0,20764907) (1,154284351) (2,475588379) (3,651992059) (4,497874987) (5,111977553) (6,10290328) (7,740862) (8,208235) (9,89) (10,179132) (11,10) (12,114517) (14,51348) (16,12019) (18,1416) (20,35)};\end{axis}\end{tikzpicture}\end{minipage}
\vspace{3mm}

\noindent\begin{minipage}[t]{60mm}\vspace{0pt}\sbox\hmboard{\chessboard[setfen=8/8/4p3/4Q3/3k4/8/1Nq5/K7 b - - 0 1]}\usebox\hmboard\\[1mm]
\makebox[\wd\hmboard]{\sffamily\hspace{3mm}\textbf{h\#10}\hfill(3+3)}\end{minipage}\hfill
\begin{minipage}[t]{52mm}\vspace{0pt}
\textcolor{ink}{\sffamily\textsc{Lösung}}\par\smallskip\small
\begin{tabular}{@{}r@{\enspace}l@{\enspace}l@{}}
\textbf{1.} & \symking{}xe5 & \symknight{}c4+\\
\textbf{2.} & \symking{}f4 & \symknight{}d2\\
\textbf{3.} & \symqueen{}d1+ & \symking{}b2\\
\textbf{4.} & e5 & \symking{}c3\\
\textbf{5.} & e4 & \symking{}d4\\
\textbf{6.} & \symking{}g3 & \symking{}e5\\
\textbf{7.} & \symking{}f2 & \symking{}f4\\
\textbf{8.} & \symking{}e1 & \symking{}g3\\
\textbf{9.} & e3 & \symking{}g2\\
\textbf{10.} & e2 & \symknight{}f3\#\\
\end{tabular}\par\vspace{3mm}
\qrcode[height=12mm]{https://helpman.komtera.lt/?fen=8/8/4p3/4Q3/3k4/8/1Nq5/K7&moves=10}\\{\sffamily\tiny\color{ink}scannen zum Lösen}
\end{minipage}
\vspace{2mm}

{\raggedright\noindent\textcolor{ink}{\sffamily\textsc{Themen}}\enspace \themetag{reines Matt} \themetag{Mustermatt} \themetag{Idealmatt} \themetag{Selbstblock}\par\smallskip{\footnotesize\textit{Hinweis:} die Partei am Zug steht im Diagramm im Schach; die Lösung beginnt mit einem Schlag.}\par}

\clearpage
\phantomsection\label{hm:KQPvkbp}
\begin{tcolorbox}[enhanced,colback=ink,colframe=ink,arc=0pt,boxrule=0pt,left=4mm,right=4mm,top=2.5mm,bottom=2.5mm]
\color{white}\sffamily{\LARGE\bfseries KQPvkbp}\hfill{\Large h\#10}\\[0.5mm]
{\small 6 Steine\quad Schwarz am Zug}\hfill{\small Nr.~264\qquad\href{https://helpman.komtera.lt/?fen=2K5/8/8/8/6p1/6P1/Q1b5/k7&moves=10}{\color{white}\bfseries lösen\,\ensuremath{\nearrow}}}
\end{tcolorbox}
\vspace{2mm}
\noindent\begin{minipage}[t]{36mm}\vspace{0pt}\sffamily\small\raggedright
\hmfigure{längste eindeutige}{{\Large h\#10}}\hmfigure{eindeutig bei dieser Länge}{{\Large 2} Stellungen}\hmfigure{längstes Matt}{{\Large h\#15}}\hmfigure{in der Klasse}{34,1 Milliarden Stellungen,\\ 20,4 Milliarden mit Hilfsmatt}
\end{minipage}\hfill
\begin{minipage}[t]{80mm}\vspace{0pt}\begin{tikzpicture}\begin{axis}[width=72mm,height=42mm,ybar,bar width=1.6pt,xmin=-1,xmax=31,ymin=0,ymax=5.392,axis y line*=left,axis x line*=bottom,xlabel={\scriptsize Halbzüge bis Matt},ylabel={\scriptsize Stellungen (Mrd.)},tick label style={font=\tiny},label style={font=\scriptsize},enlarge x limits=false,ylabel shift=-3pt,xlabel shift=-3pt]\addplot[fill=ink!55,draw=none] coordinates {(0,0.01833) (1,0.1424) (2,0.4876) (3,1.863) (4,4.491) (5,4.409) (6,4.902) (7,2.2) (8,1.408) (9,0.3071) (10,0.1236) (11,0.004891) (12,0.04176) (13,3.746e-06) (14,0.02369) (15,1.35e-07) (16,0.002494) (17,7e-08) (18,0.0005785) (19,1.9e-08) (20,0.000293) (21,1.2e-08) (22,0.0001995) (23,7e-09) (24,7.984e-05) (26,2.43e-05) (28,1.524e-05) (30,2.51e-07)};\draw[densely dotted] (axis cs:20,0) -- (axis cs:20,5.392);\end{axis}\begin{axis}[width=72mm,height=42mm,axis y line*=right,axis x line=none,ymode=log,xmin=-1,xmax=31,ylabel={\scriptsize eindeutig (log)},tick label style={font=\tiny},label style={font=\scriptsize},ylabel shift=-3pt]\addplot[mark=*,mark size=0.7pt,draw=black,thin] coordinates {(0,18333952) (1,102869053) (2,237684124) (3,313815114) (4,383735743) (5,133017590) (6,67376862) (7,14670805) (8,4127806) (9,149643) (10,245756) (11,41) (12,157083) (14,17244) (16,1454) (18,140) (20,2)};\end{axis}\end{tikzpicture}\end{minipage}
\vspace{3mm}

\noindent\begin{minipage}[t]{60mm}\vspace{0pt}\sbox\hmboard{\chessboard[setfen=2K5/8/8/8/6p1/6P1/Q1b5/k7 b - - 0 1]}\usebox\hmboard\\[1mm]
\makebox[\wd\hmboard]{\sffamily\hspace{3mm}\textbf{h\#10}\hfill(3+3)}\end{minipage}\hfill
\begin{minipage}[t]{52mm}\vspace{0pt}
\textcolor{ink}{\sffamily\textsc{Lösung}}\par\smallskip\small
\begin{tabular}{@{}r@{\enspace}l@{\enspace}l@{}}
\textbf{1.} & \symking{}xa2 & \symking{}d7\\
\textbf{2.} & \symking{}b2 & \symking{}e6\\
\textbf{3.} & \symbishop{}b3+ & \symking{}f5\\
\textbf{4.} & \symking{}c3 & \symking{}xg4\\
\textbf{5.} & \symking{}d4 & \symking{}h5\\
\textbf{6.} & \symking{}e5 & g4\\
\textbf{7.} & \symking{}f6 & g5+\\
\textbf{8.} & \symking{}g7 & g6\\
\textbf{9.} & \symking{}h8 & \symking{}h6\\
\textbf{10.} & \symbishop{}g8 & g7\#\\
\end{tabular}\par\vspace{3mm}
\qrcode[height=12mm]{https://helpman.komtera.lt/?fen=2K5/8/8/8/6p1/6P1/Q1b5/k7&moves=10}\\{\sffamily\tiny\color{ink}scannen zum Lösen}
\end{minipage}
\vspace{2mm}

{\raggedright\noindent\textcolor{ink}{\sffamily\textsc{Themen}}\enspace \themetag{reines Matt} \themetag{Mustermatt} \themetag{Idealmatt} \themetag{Selbstblock}\par\smallskip{\footnotesize\textit{Hinweis:} die Partei am Zug steht im Diagramm im Schach; die Lösung beginnt mit einem Schlag.}\par}

\clearpage
\phantomsection\label{hm:KQPvkpp}
\begin{tcolorbox}[enhanced,colback=ink,colframe=ink,arc=0pt,boxrule=0pt,left=4mm,right=4mm,top=2.5mm,bottom=2.5mm]
\color{white}\sffamily{\LARGE\bfseries KQPvkpp}\hfill{\Large h\#10}\\[0.5mm]
{\small 6 Steine\quad Schwarz am Zug}\hfill{\small Nr.~265\qquad\href{https://helpman.komtera.lt/?fen=6Qk/8/1p6/8/8/p7/P7/K7&moves=10}{\color{white}\bfseries lösen\,\ensuremath{\nearrow}}}
\end{tcolorbox}
\vspace{2mm}
\noindent\begin{minipage}[t]{36mm}\vspace{0pt}\sffamily\small\raggedright
\hmfigure{längste eindeutige}{{\Large h\#10}}\hmfigure{eindeutig bei dieser Länge}{{\Large 4} Stellungen}\hmfigure{längstes Matt}{{\Large h\#16}}\hmfigure{in der Klasse}{25,6 Milliarden Stellungen,\\ 8,07 Milliarden mit Hilfsmatt}
\end{minipage}\hfill
\begin{minipage}[t]{80mm}\vspace{0pt}\begin{tikzpicture}\begin{axis}[width=72mm,height=42mm,ybar,bar width=1.6pt,xmin=-1,xmax=33,ymin=0,ymax=2.239,axis y line*=left,axis x line*=bottom,xlabel={\scriptsize Halbzüge bis Matt},ylabel={\scriptsize Stellungen (Mrd.)},tick label style={font=\tiny},label style={font=\scriptsize},enlarge x limits=false,ylabel shift=-3pt,xlabel shift=-3pt]\addplot[fill=ink!55,draw=none] coordinates {(0,0.01002) (1,0.07287) (2,0.215) (3,0.7508) (4,1.797) (5,1.583) (6,2.036) (7,0.8082) (8,0.6089) (9,0.1065) (10,0.05163) (11,0.001906) (12,0.01556) (13,2.596e-06) (14,0.009658) (15,1.4e-08) (16,0.001892) (17,3e-09) (18,0.001243) (19,1e-09) (20,0.0007972) (21,3e-09) (22,0.0003188) (24,0.0001052) (26,2.749e-05) (28,1.077e-05) (30,4.287e-06) (32,4e-08)};\draw[densely dotted] (axis cs:20,0) -- (axis cs:20,2.239);\end{axis}\begin{axis}[width=72mm,height=42mm,axis y line*=right,axis x line=none,ymode=log,xmin=-1,xmax=33,ylabel={\scriptsize eindeutig (log)},tick label style={font=\tiny},label style={font=\scriptsize},ylabel shift=-3pt]\addplot[mark=*,mark size=0.7pt,draw=black,thin] coordinates {(0,10023853) (1,48212137) (2,107060442) (3,126613184) (4,155977091) (5,48286554) (6,27873243) (7,4782247) (8,1653885) (9,36194) (10,86793) (11,61) (12,66853) (14,5631) (16,1314) (18,209) (20,4)};\end{axis}\end{tikzpicture}\end{minipage}
\vspace{3mm}

\noindent\begin{minipage}[t]{60mm}\vspace{0pt}\sbox\hmboard{\chessboard[setfen=6Qk/8/1p6/8/8/p7/P7/K7 b - - 0 1]}\usebox\hmboard\\[1mm]
\makebox[\wd\hmboard]{\sffamily\hspace{3mm}\textbf{h\#10}\hfill(3+3)}\end{minipage}\hfill
\begin{minipage}[t]{52mm}\vspace{0pt}
\textcolor{ink}{\sffamily\textsc{Lösung}}\par\smallskip\small
\begin{tabular}{@{}r@{\enspace}l@{\enspace}l@{}}
\textbf{1.} & \symking{}xg8 & \symking{}b1\\
\textbf{2.} & b5 & \symking{}c2\\
\textbf{3.} & b4 & \symking{}c1\\
\textbf{4.} & b3 & axb3\\
\textbf{5.} & \symking{}f7 & b4\\
\textbf{6.} & \symking{}e6 & b5\\
\textbf{7.} & \symking{}d5 & b6\\
\textbf{8.} & \symking{}c4 & b7\\
\textbf{9.} & \symking{}b3 & b8=\symqueen{}+\\
\textbf{10.} & \symking{}a2 & \symqueen{}b1\#\\
\end{tabular}\par\vspace{3mm}
\qrcode[height=12mm]{https://helpman.komtera.lt/?fen=6Qk/8/1p6/8/8/p7/P7/K7&moves=10}\\{\sffamily\tiny\color{ink}scannen zum Lösen}
\end{minipage}
\vspace{2mm}

{\raggedright\noindent\textcolor{ink}{\sffamily\textsc{Themen}}\enspace \themetag{Umwandlung} \themetag{Excelsior} \themetag{Phönix}\par\smallskip{\footnotesize\textit{Hinweis:} die Partei am Zug steht im Diagramm im Schach; die Lösung beginnt mit einem Schlag.}\par}

\clearpage
\phantomsection\label{hm:KRBPvkp}
\begin{tcolorbox}[enhanced,colback=ink,colframe=ink,arc=0pt,boxrule=0pt,left=4mm,right=4mm,top=2.5mm,bottom=2.5mm]
\color{white}\sffamily{\LARGE\bfseries KRBPvkp}\hfill{\Large h\#10}\\[0.5mm]
{\small 6 Steine\quad Schwarz am Zug}\hfill{\small Nr.~266\qquad\href{https://helpman.komtera.lt/?fen=8/8/7B/5p2/1R3P2/k7/8/K7&moves=10}{\color{white}\bfseries lösen\,\ensuremath{\nearrow}}}
\end{tcolorbox}
\vspace{2mm}
\noindent\begin{minipage}[t]{36mm}\vspace{0pt}\sffamily\small\raggedright
\hmfigure{längste eindeutige}{{\Large h\#10}}\hmfigure{eindeutig bei dieser Länge}{{\Large 3} Stellungen}\hmfigure{längstes Matt}{{\Large h\#15}}\hmfigure{in der Klasse}{34,1 Milliarden Stellungen,\\ 22,4 Milliarden mit Hilfsmatt}
\end{minipage}\hfill
\begin{minipage}[t]{80mm}\vspace{0pt}\begin{tikzpicture}\begin{axis}[width=72mm,height=42mm,ybar,bar width=1.6pt,xmin=-1,xmax=31,ymin=0,ymax=7.272,axis y line*=left,axis x line*=bottom,xlabel={\scriptsize Halbzüge bis Matt},ylabel={\scriptsize Stellungen (Mrd.)},tick label style={font=\tiny},label style={font=\scriptsize},enlarge x limits=false,ylabel shift=-3pt,xlabel shift=-3pt]\addplot[fill=ink!55,draw=none] coordinates {(0,0.01894) (1,0.1104) (2,0.2898) (3,1.231) (4,3.137) (5,4.292) (6,6.611) (7,3.047) (8,2.78) (9,0.5735) (10,0.3074) (11,0.02051) (12,0.01129) (13,8.504e-05) (14,0.001151) (15,5.802e-06) (16,0.0003428) (17,7.41e-07) (18,0.0001442) (19,1.88e-07) (20,8.674e-05) (21,3.6e-08) (22,4.817e-05) (23,2.2e-08) (24,1.854e-05) (25,8e-09) (26,3.211e-06) (28,2.14e-06) (30,3.9e-08)};\draw[densely dotted] (axis cs:20,0) -- (axis cs:20,7.272);\end{axis}\begin{axis}[width=72mm,height=42mm,axis y line*=right,axis x line=none,ymode=log,xmin=-1,xmax=31,ylabel={\scriptsize eindeutig (log)},tick label style={font=\tiny},label style={font=\scriptsize},ylabel shift=-3pt]\addplot[mark=*,mark size=0.7pt,draw=black,thin] coordinates {(0,18937905) (1,88276257) (2,173906009) (3,409574777) (4,751103160) (5,387868461) (6,247814015) (7,23458470) (8,10543783) (9,509922) (10,226632) (11,2003) (12,18340) (13,34) (14,1488) (15,2) (16,553) (18,71) (20,3)};\end{axis}\end{tikzpicture}\end{minipage}
\vspace{3mm}

\noindent\begin{minipage}[t]{60mm}\vspace{0pt}\sbox\hmboard{\chessboard[setfen=8/8/7B/5p2/1R3P2/k7/8/K7 b - - 0 1]}\usebox\hmboard\\[1mm]
\makebox[\wd\hmboard]{\sffamily\hspace{3mm}\textbf{h\#10}\hfill(4+2)}\end{minipage}\hfill
\begin{minipage}[t]{52mm}\vspace{0pt}
\textcolor{ink}{\sffamily\textsc{Lösung}}\par\smallskip\small
\begin{tabular}{@{}r@{\enspace}l@{\enspace}l@{}}
\textbf{1.} & \symking{}xb4 & \symking{}b2\\
\textbf{2.} & \symking{}c5 & \symking{}c3\\
\textbf{3.} & \symking{}d6 & \symking{}d4\\
\textbf{4.} & \symking{}e7 & \symking{}e5\\
\textbf{5.} & \symking{}f7 & \symking{}xf5\\
\textbf{6.} & \symking{}g8 & \symking{}g6\\
\textbf{7.} & \symking{}h8 & f5\\
\textbf{8.} & \symking{}g8 & f6\\
\textbf{9.} & \symking{}h8 & \symbishop{}g7+\\
\textbf{10.} & \symking{}g8 & f7\#\\
\end{tabular}\par\vspace{3mm}
\qrcode[height=12mm]{https://helpman.komtera.lt/?fen=8/8/7B/5p2/1R3P2/k7/8/K7&moves=10}\\{\sffamily\tiny\color{ink}scannen zum Lösen}
\end{minipage}
\vspace{2mm}

{\raggedright\noindent\textcolor{ink}{\sffamily\textsc{Themen}}\enspace \themetag{reines Matt} \themetag{Mustermatt} \themetag{Idealmatt} \themetag{Switchback} \themetag{ein Stein zieht}\par\smallskip{\footnotesize\textit{Hinweis:} die Lösung beginnt mit einem Schlag.}\par}

\clearpage
\phantomsection\label{hm:KRBvkbp}
\begin{tcolorbox}[enhanced,colback=ink,colframe=ink,arc=0pt,boxrule=0pt,left=4mm,right=4mm,top=2.5mm,bottom=2.5mm]
\color{white}\sffamily{\LARGE\bfseries KRBvkbp}\hfill{\Large h\#10}\\[0.5mm]
{\small 6 Steine\quad Schwarz am Zug}\hfill{\small Nr.~267\qquad\href{https://helpman.komtera.lt/?fen=8/K7/8/kR6/8/1B3p2/8/7b&moves=10}{\color{white}\bfseries lösen\,\ensuremath{\nearrow}}}
\end{tcolorbox}
\vspace{2mm}
\noindent\begin{minipage}[t]{36mm}\vspace{0pt}\sffamily\small\raggedright
\hmfigure{längste eindeutige}{{\Large h\#10}}\hmfigure{eindeutig bei dieser Länge}{{\Large 1} Stellung}\hmfigure{längstes Matt}{{\Large h\#16.5}}\hmfigure{in der Klasse}{45,4 Milliarden Stellungen,\\ 28,1 Milliarden mit Hilfsmatt}
\end{minipage}\hfill
\begin{minipage}[t]{80mm}\vspace{0pt}\begin{tikzpicture}\begin{axis}[width=72mm,height=42mm,ybar,bar width=1.6pt,xmin=-1,xmax=34,ymin=0,ymax=9.159,axis y line*=left,axis x line*=bottom,xlabel={\scriptsize Halbzüge bis Matt},ylabel={\scriptsize Stellungen (Mrd.)},tick label style={font=\tiny},label style={font=\scriptsize},enlarge x limits=false,ylabel shift=-3pt,xlabel shift=-3pt]\addplot[fill=ink!55,draw=none] coordinates {(0,0.01419) (1,0.06903) (2,0.2308) (3,1.05) (4,3.173) (5,5.708) (6,8.326) (7,4.75) (8,3.051) (9,1.093) (10,0.5488) (11,0.08807) (12,0.03069) (13,0.0006712) (14,0.002042) (15,1.5e-08) (16,0.001624) (17,9e-09) (18,0.001295) (19,1.4e-08) (20,0.0009738) (21,4e-09) (22,0.001035) (23,1.65e-07) (24,0.0009481) (25,8e-09) (26,0.0006574) (27,3e-08) (28,0.000545) (29,9e-09) (30,0.0002416) (31,1e-09) (32,0.0001352) (33,2e-09)};\draw[densely dotted] (axis cs:20,0) -- (axis cs:20,9.159);\end{axis}\begin{axis}[width=72mm,height=42mm,axis y line*=right,axis x line=none,ymode=log,xmin=-1,xmax=34,ylabel={\scriptsize eindeutig (log)},tick label style={font=\tiny},label style={font=\scriptsize},ylabel shift=-3pt]\addplot[mark=*,mark size=0.7pt,draw=black,thin] coordinates {(0,14193534) (1,67340746) (2,129475884) (3,399446236) (4,923055454) (5,500971289) (6,224823985) (7,35544049) (8,14724404) (9,1873084) (10,484349) (11,5952) (12,2983) (13,1) (14,264) (16,69) (18,22) (20,1)};\end{axis}\end{tikzpicture}\end{minipage}
\vspace{3mm}

\noindent\begin{minipage}[t]{60mm}\vspace{0pt}\sbox\hmboard{\chessboard[setfen=8/K7/8/kR6/8/1B3p2/8/7b b - - 0 1]}\usebox\hmboard\\[1mm]
\makebox[\wd\hmboard]{\sffamily\hspace{3mm}\textbf{h\#10}\hfill(3+3)}\end{minipage}\hfill
\begin{minipage}[t]{52mm}\vspace{0pt}
\textcolor{ink}{\sffamily\textsc{Lösung}}\par\smallskip\small
\begin{tabular}{@{}r@{\enspace}l@{\enspace}l@{}}
\textbf{1.} & \symking{}xb5 & \symbishop{}d5\\
\textbf{2.} & f2 & \symbishop{}xh1\\
\textbf{3.} & \symking{}c4 & \symking{}b6\\
\textbf{4.} & \symking{}d3 & \symking{}c5\\
\textbf{5.} & \symking{}e2 & \symking{}d4\\
\textbf{6.} & \symking{}f1 & \symbishop{}g2+\\
\textbf{7.} & \symking{}g1 & \symbishop{}h3\\
\textbf{8.} & \symking{}h1 & \symking{}e3\\
\textbf{9.} & f1=\symknight{}+ & \symking{}f2\\
\textbf{10.} & \symknight{}h2 & \symbishop{}g2\#\\
\end{tabular}\par\vspace{3mm}
\qrcode[height=12mm]{https://helpman.komtera.lt/?fen=8/K7/8/kR6/8/1B3p2/8/7b&moves=10}\\{\sffamily\tiny\color{ink}scannen zum Lösen}
\end{minipage}
\vspace{2mm}

{\raggedright\noindent\textcolor{ink}{\sffamily\textsc{Themen}}\enspace \themetag{reines Matt} \themetag{Mustermatt} \themetag{Idealmatt} \themetag{Umwandlung} \themetag{Unterverwandlung} \themetag{Switchback} \themetag{Selbstblock} \themetag{Kniest} \themetag{Umnov}\par\smallskip{\footnotesize\textit{Hinweis:} die Partei am Zug steht im Diagramm im Schach; die Lösung beginnt mit einem Schlag.}\par}

\clearpage
\phantomsection\label{hm:KRBvknp}
\begin{tcolorbox}[enhanced,colback=ink,colframe=ink,arc=0pt,boxrule=0pt,left=4mm,right=4mm,top=2.5mm,bottom=2.5mm]
\color{white}\sffamily{\LARGE\bfseries KRBvknp}\hfill{\Large h\#10}\\[0.5mm]
{\small 6 Steine\quad Schwarz am Zug}\hfill{\small Nr.~268\qquad\href{https://helpman.komtera.lt/?fen=k7/R1B2p2/8/5n2/8/8/K7/8&moves=10}{\color{white}\bfseries lösen\,\ensuremath{\nearrow}}}
\end{tcolorbox}
\vspace{2mm}
\noindent\begin{minipage}[t]{36mm}\vspace{0pt}\sffamily\small\raggedright
\hmfigure{längste eindeutige}{{\Large h\#10}}\hmfigure{eindeutig bei dieser Länge}{{\Large 1} Stellung}\hmfigure{längstes Matt}{{\Large h\#15.5}}\hmfigure{in der Klasse}{45,4 Milliarden Stellungen,\\ 29 Milliarden mit Hilfsmatt}
\end{minipage}\hfill
\begin{minipage}[t]{80mm}\vspace{0pt}\begin{tikzpicture}\begin{axis}[width=72mm,height=42mm,ybar,bar width=1.6pt,xmin=-1,xmax=32,ymin=0,ymax=8.825,axis y line*=left,axis x line*=bottom,xlabel={\scriptsize Halbzüge bis Matt},ylabel={\scriptsize Stellungen (Mrd.)},tick label style={font=\tiny},label style={font=\scriptsize},enlarge x limits=false,ylabel shift=-3pt,xlabel shift=-3pt]\addplot[fill=ink!55,draw=none] coordinates {(0,0.0168) (1,0.08442) (2,0.2313) (3,0.9416) (4,2.459) (5,4.595) (6,8.023) (7,5.525) (8,4.863) (9,1.546) (10,0.5811) (11,0.06694) (12,0.008939) (13,4.258e-05) (14,0.00224) (15,3.9e-08) (16,0.002494) (17,1.12e-07) (18,0.003075) (19,2.8e-07) (20,0.00258) (21,1.78e-07) (22,0.001208) (23,2.23e-07) (24,9.084e-05) (25,1.16e-07) (26,4e-09) (27,1.48e-07) (28,6e-09) (29,4.5e-08) (31,2.6e-08)};\draw[densely dotted] (axis cs:20,0) -- (axis cs:20,8.825);\end{axis}\begin{axis}[width=72mm,height=42mm,axis y line*=right,axis x line=none,ymode=log,xmin=-1,xmax=32,ylabel={\scriptsize eindeutig (log)},tick label style={font=\tiny},label style={font=\scriptsize},ylabel shift=-3pt]\addplot[mark=*,mark size=0.7pt,draw=black,thin] coordinates {(0,16796463) (1,82447397) (2,125842666) (3,342140398) (4,711019089) (5,527865523) (6,356419040) (7,53167243) (8,19259879) (9,1942823) (10,232675) (11,839) (12,7119) (14,899) (16,442) (18,17) (20,1)};\end{axis}\end{tikzpicture}\end{minipage}
\vspace{3mm}

\noindent\begin{minipage}[t]{60mm}\vspace{0pt}\sbox\hmboard{\chessboard[setfen=k7/R1B2p2/8/5n2/8/8/K7/8 b - - 0 1]}\usebox\hmboard\\[1mm]
\makebox[\wd\hmboard]{\sffamily\hspace{3mm}\textbf{h\#10}\hfill(3+3)}\end{minipage}\hfill
\begin{minipage}[t]{52mm}\vspace{0pt}
\textcolor{ink}{\sffamily\textsc{Lösung}}\par\smallskip\small
\begin{tabular}{@{}r@{\enspace}l@{\enspace}l@{}}
\textbf{1.} & \symking{}xa7 & \symking{}b3\\
\textbf{2.} & \symking{}b7 & \symking{}c4\\
\textbf{3.} & \symking{}c8 & \symking{}d5\\
\textbf{4.} & \symking{}d7 & \symking{}e5\\
\textbf{5.} & \symking{}e8 & \symking{}f6\\
\textbf{6.} & \symking{}f8 & \symbishop{}d6+\\
\textbf{7.} & \symking{}g8 & \symbishop{}f8\\
\textbf{8.} & \symking{}h8 & \symking{}xf7\\
\textbf{9.} & \symknight{}h6+ & \symking{}g6\\
\textbf{10.} & \symknight{}g8 & \symbishop{}g7\#\\
\end{tabular}\par\vspace{3mm}
\qrcode[height=12mm]{https://helpman.komtera.lt/?fen=k7/R1B2p2/8/5n2/8/8/K7/8&moves=10}\\{\sffamily\tiny\color{ink}scannen zum Lösen}
\end{minipage}
\vspace{2mm}

{\raggedright\noindent\textcolor{ink}{\sffamily\textsc{Themen}}\enspace \themetag{reines Matt} \themetag{Mustermatt} \themetag{Idealmatt} \themetag{Selbstblock} \themetag{Umnov}\par\smallskip{\footnotesize\textit{Hinweis:} die Partei am Zug steht im Diagramm im Schach; die Lösung beginnt mit einem Schlag.}\par}

\clearpage
\phantomsection\label{hm:KRBvkqn}
\begin{tcolorbox}[enhanced,colback=ink,colframe=ink,arc=0pt,boxrule=0pt,left=4mm,right=4mm,top=2.5mm,bottom=2.5mm]
\color{white}\sffamily{\LARGE\bfseries KRBvkqn}\hfill{\Large h\#10}\\[0.5mm]
{\small 6 Steine\quad Schwarz am Zug}\hfill{\small Nr.~269\qquad\href{https://helpman.komtera.lt/?fen=8/8/8/8/8/3q4/R4n2/k1KB4&moves=10}{\color{white}\bfseries lösen\,\ensuremath{\nearrow}}}
\end{tcolorbox}
\vspace{2mm}
\noindent\begin{minipage}[t]{36mm}\vspace{0pt}\sffamily\small\raggedright
\hmfigure{längste eindeutige}{{\Large h\#10}}\hmfigure{eindeutig bei dieser Länge}{{\Large 4} Stellungen}\hmfigure{längstes Matt}{{\Large h\#13}}\hmfigure{in der Klasse}{15,5 Milliarden Stellungen,\\ 7,99 Milliarden mit Hilfsmatt}
\end{minipage}\hfill
\begin{minipage}[t]{80mm}\vspace{0pt}\begin{tikzpicture}\begin{axis}[width=72mm,height=42mm,ybar,bar width=1.6pt,xmin=-1,xmax=27,ymin=0,ymax=2.826,axis y line*=left,axis x line*=bottom,xlabel={\scriptsize Halbzüge bis Matt},ylabel={\scriptsize Stellungen (Mrd.)},tick label style={font=\tiny},label style={font=\scriptsize},enlarge x limits=false,ylabel shift=-3pt,xlabel shift=-3pt]\addplot[fill=ink!55,draw=none] coordinates {(0,0.001977) (1,0.00935) (2,0.0441) (3,0.2544) (4,0.772) (5,1.964) (6,2.569) (7,1.914) (8,0.3625) (9,0.08716) (10,0.0005748) (11,0.003945) (12,0.0005122) (13,5.571e-05) (14,0.000537) (15,3.4e-08) (16,0.0005505) (17,1.1e-08) (18,0.000689) (19,4.1e-08) (20,0.0007723) (21,5e-08) (22,0.000489) (23,1e-08) (24,8.062e-05) (25,2e-09) (26,7.21e-07)};\draw[densely dotted] (axis cs:20,0) -- (axis cs:20,2.826);\end{axis}\begin{axis}[width=72mm,height=42mm,axis y line*=right,axis x line=none,ymode=log,xmin=-1,xmax=27,ylabel={\scriptsize eindeutig (log)},tick label style={font=\tiny},label style={font=\scriptsize},ylabel shift=-3pt]\addplot[mark=*,mark size=0.7pt,draw=black,thin] coordinates {(0,1977443) (1,9233923) (2,23582639) (3,101442208) (4,207929771) (5,167090615) (6,38905639) (7,7674969) (8,24440) (9,104093) (10,3654) (11,1561) (12,1568) (14,588) (16,88) (18,124) (20,4)};\end{axis}\end{tikzpicture}\end{minipage}
\vspace{3mm}

\noindent\begin{minipage}[t]{60mm}\vspace{0pt}\sbox\hmboard{\chessboard[setfen=8/8/8/8/8/3q4/R4n2/k1KB4 b - - 0 1]}\usebox\hmboard\\[1mm]
\makebox[\wd\hmboard]{\sffamily\hspace{3mm}\textbf{h\#10}\hfill(3+3)}\end{minipage}\hfill
\begin{minipage}[t]{52mm}\vspace{0pt}
\textcolor{ink}{\sffamily\textsc{Lösung}}\par\smallskip\small
\begin{tabular}{@{}r@{\enspace}l@{\enspace}l@{}}
\textbf{1.} & \symking{}xa2 & \symbishop{}e2\\
\textbf{2.} & \symking{}b3 & \symbishop{}xd3\\
\textbf{3.} & \symking{}c3 & \symking{}b1\\
\textbf{4.} & \symking{}d2 & \symbishop{}f1\\
\textbf{5.} & \symking{}e1 & \symking{}c2\\
\textbf{6.} & \symknight{}h3 & \symking{}d3\\
\textbf{7.} & \symking{}f2 & \symking{}e4\\
\textbf{8.} & \symking{}g1 & \symking{}f3\\
\textbf{9.} & \symking{}h1 & \symking{}g3\\
\textbf{10.} & \symknight{}g1 & \symbishop{}g2\#\\
\end{tabular}\par\vspace{3mm}
\qrcode[height=12mm]{https://helpman.komtera.lt/?fen=8/8/8/8/8/3q4/R4n2/k1KB4&moves=10}\\{\sffamily\tiny\color{ink}scannen zum Lösen}
\end{minipage}
\vspace{2mm}

{\raggedright\noindent\textcolor{ink}{\sffamily\textsc{Themen}}\enspace \themetag{reines Matt} \themetag{Mustermatt} \themetag{Idealmatt} \themetag{Selbstblock} \themetag{ohne Schach}\par\smallskip{\footnotesize\textit{Hinweis:} die Partei am Zug steht im Diagramm im Schach; die Lösung beginnt mit einem Schlag.}\par}

\clearpage
\phantomsection\label{hm:KRBvkqp}
\begin{tcolorbox}[enhanced,colback=ink,colframe=ink,arc=0pt,boxrule=0pt,left=4mm,right=4mm,top=2.5mm,bottom=2.5mm]
\color{white}\sffamily{\LARGE\bfseries KRBvkqp}\hfill{\Large h\#10}\\[0.5mm]
{\small 6 Steine\quad Schwarz am Zug}\hfill{\small Nr.~270\qquad\href{https://helpman.komtera.lt/?fen=8/8/3R4/1K1kqp2/8/8/8/1B6&moves=10}{\color{white}\bfseries lösen\,\ensuremath{\nearrow}}}
\end{tcolorbox}
\vspace{2mm}
\noindent\begin{minipage}[t]{36mm}\vspace{0pt}\sffamily\small\raggedright
\hmfigure{längste eindeutige}{{\Large h\#10}}\hmfigure{eindeutig bei dieser Länge}{{\Large 4} Stellungen}\hmfigure{längstes Matt}{{\Large h\#17}}\hmfigure{in der Klasse}{45,4 Milliarden Stellungen,\\ 24,7 Milliarden mit Hilfsmatt}
\end{minipage}\hfill
\begin{minipage}[t]{80mm}\vspace{0pt}\begin{tikzpicture}\begin{axis}[width=72mm,height=42mm,ybar,bar width=1.6pt,xmin=-1,xmax=35,ymin=0,ymax=8.793,axis y line*=left,axis x line*=bottom,xlabel={\scriptsize Halbzüge bis Matt},ylabel={\scriptsize Stellungen (Mrd.)},tick label style={font=\tiny},label style={font=\scriptsize},enlarge x limits=false,ylabel shift=-3pt,xlabel shift=-3pt]\addplot[fill=ink!55,draw=none] coordinates {(0,0.007535) (1,0.03636) (2,0.1632) (3,0.8851) (4,2.589) (5,6.116) (6,7.994) (7,5.436) (8,1.155) (9,0.26) (10,0.00191) (11,0.01027) (12,0.001711) (13,0.0005651) (14,0.001652) (15,3.151e-06) (16,0.001687) (17,2.908e-06) (18,0.00172) (19,3.31e-06) (20,0.00145) (21,3.361e-06) (22,0.0008548) (23,3.874e-06) (24,0.0007235) (25,4.529e-06) (26,0.000559) (27,1.734e-06) (28,0.0005326) (29,1.098e-06) (30,0.0003597) (31,6.44e-07) (32,0.0002726) (33,8.5e-08) (34,2.78e-05)};\draw[densely dotted] (axis cs:20,0) -- (axis cs:20,8.793);\end{axis}\begin{axis}[width=72mm,height=42mm,axis y line*=right,axis x line=none,ymode=log,xmin=-1,xmax=35,ylabel={\scriptsize eindeutig (log)},tick label style={font=\tiny},label style={font=\scriptsize},ylabel shift=-3pt]\addplot[mark=*,mark size=0.7pt,draw=black,thin] coordinates {(0,7534565) (1,35845455) (2,88784186) (3,356923413) (4,692196563) (5,528576770) (6,123773990) (7,21263387) (8,80865) (9,212087) (10,10987) (11,5377) (12,3076) (13,21) (14,845) (16,258) (18,32) (20,4)};\end{axis}\end{tikzpicture}\end{minipage}
\vspace{3mm}

\noindent\begin{minipage}[t]{60mm}\vspace{0pt}\sbox\hmboard{\chessboard[setfen=8/8/3R4/1K1kqp2/8/8/8/1B6 b - - 0 1]}\usebox\hmboard\\[1mm]
\makebox[\wd\hmboard]{\sffamily\hspace{3mm}\textbf{h\#10}\hfill(3+3)}\end{minipage}\hfill
\begin{minipage}[t]{52mm}\vspace{0pt}
\textcolor{ink}{\sffamily\textsc{Lösung}}\par\smallskip\small
\begin{tabular}{@{}r@{\enspace}l@{\enspace}l@{}}
\textbf{1.} & \symqueen{}xd6 & \symking{}a4\\
\textbf{2.} & f4 & \symking{}b3\\
\textbf{3.} & f3 & \symking{}c2\\
\textbf{4.} & f2 & \symking{}d1\\
\textbf{5.} & f1=\symknight{} & \symbishop{}a2+\\
\textbf{6.} & \symking{}e4+ & \symking{}e1\\
\textbf{7.} & \symking{}f3 & \symbishop{}e6\\
\textbf{8.} & \symking{}g2 & \symbishop{}h3+\\
\textbf{9.} & \symking{}h1 & \symking{}f2\\
\textbf{10.} & \symknight{}h2 & \symbishop{}g2\#\\
\end{tabular}\par\vspace{3mm}
\qrcode[height=12mm]{https://helpman.komtera.lt/?fen=8/8/3R4/1K1kqp2/8/8/8/1B6&moves=10}\\{\sffamily\tiny\color{ink}scannen zum Lösen}
\end{minipage}
\vspace{2mm}

{\raggedright\noindent\textcolor{ink}{\sffamily\textsc{Themen}}\enspace \themetag{reines Matt} \themetag{Mustermatt} \themetag{Umwandlung} \themetag{Unterverwandlung} \themetag{Selbstblock}\par\smallskip{\footnotesize\textit{Hinweis:} die Partei am Zug steht im Diagramm im Schach; die Lösung beginnt mit einem Schlag.}\par}

\clearpage
\phantomsection\label{hm:KRBvkrn}
\begin{tcolorbox}[enhanced,colback=ink,colframe=ink,arc=0pt,boxrule=0pt,left=4mm,right=4mm,top=2.5mm,bottom=2.5mm]
\color{white}\sffamily{\LARGE\bfseries KRBvkrn}\hfill{\Large h\#10}\\[0.5mm]
{\small 6 Steine\quad Schwarz am Zug}\hfill{\small Nr.~271\qquad\href{https://helpman.komtera.lt/?fen=8/8/2B5/8/8/kR6/6nr/1K6&moves=10}{\color{white}\bfseries lösen\,\ensuremath{\nearrow}}}
\end{tcolorbox}
\vspace{2mm}
\noindent\begin{minipage}[t]{36mm}\vspace{0pt}\sffamily\small\raggedright
\hmfigure{längste eindeutige}{{\Large h\#10}}\hmfigure{eindeutig bei dieser Länge}{{\Large 3} Stellungen}\hmfigure{längstes Matt}{{\Large h\#13}}\hmfigure{in der Klasse}{15,5 Milliarden Stellungen,\\ 8,69 Milliarden mit Hilfsmatt}
\end{minipage}\hfill
\begin{minipage}[t]{80mm}\vspace{0pt}\begin{tikzpicture}\begin{axis}[width=72mm,height=42mm,ybar,bar width=1.6pt,xmin=-1,xmax=27,ymin=0,ymax=3.256,axis y line*=left,axis x line*=bottom,xlabel={\scriptsize Halbzüge bis Matt},ylabel={\scriptsize Stellungen (Mrd.)},tick label style={font=\tiny},label style={font=\scriptsize},enlarge x limits=false,ylabel shift=-3pt,xlabel shift=-3pt]\addplot[fill=ink!55,draw=none] coordinates {(0,0.003009) (1,0.0152) (2,0.06239) (3,0.3203) (4,1.087) (5,2.28) (6,2.96) (7,1.575) (8,0.3387) (9,0.04792) (10,0.0006594) (11,0.0001999) (12,0.0005841) (13,5.056e-06) (14,0.0006408) (15,1.1e-08) (16,0.00067) (17,6e-09) (18,0.000866) (19,8e-09) (20,0.0009012) (21,6e-09) (22,0.0005151) (23,1e-08) (24,5.473e-05) (26,1.89e-07)};\draw[densely dotted] (axis cs:20,0) -- (axis cs:20,3.256);\end{axis}\begin{axis}[width=72mm,height=42mm,axis y line*=right,axis x line=none,ymode=log,xmin=-1,xmax=27,ylabel={\scriptsize eindeutig (log)},tick label style={font=\tiny},label style={font=\scriptsize},ylabel shift=-3pt]\addplot[mark=*,mark size=0.7pt,draw=black,thin] coordinates {(0,3009094) (1,14983049) (2,35000994) (3,129597235) (4,347303725) (5,203499414) (6,51464493) (7,5049803) (8,32492) (9,5847) (10,3905) (11,130) (12,1859) (14,484) (16,144) (18,62) (20,3)};\end{axis}\end{tikzpicture}\end{minipage}
\vspace{3mm}

\noindent\begin{minipage}[t]{60mm}\vspace{0pt}\sbox\hmboard{\chessboard[setfen=8/8/2B5/8/8/kR6/6nr/1K6 b - - 0 1]}\usebox\hmboard\\[1mm]
\makebox[\wd\hmboard]{\sffamily\hspace{3mm}\textbf{h\#10}\hfill(3+3)}\end{minipage}\hfill
\begin{minipage}[t]{52mm}\vspace{0pt}
\textcolor{ink}{\sffamily\textsc{Lösung}}\par\smallskip\small
\begin{tabular}{@{}r@{\enspace}l@{\enspace}l@{}}
\textbf{1.} & \symking{}xb3 & \symking{}c1\\
\textbf{2.} & \symking{}c4 & \symking{}d2\\
\textbf{3.} & \symknight{}f4+ & \symking{}e3\\
\textbf{4.} & \symknight{}h3 & \symking{}f3\\
\textbf{5.} & \symking{}d3 & \symking{}g3\\
\textbf{6.} & \symking{}e2 & \symking{}xh2\\
\textbf{7.} & \symking{}f1 & \symking{}g3\\
\textbf{8.} & \symking{}g1 & \symbishop{}b5\\
\textbf{9.} & \symking{}h1 & \symbishop{}f1\\
\textbf{10.} & \symknight{}g1 & \symbishop{}g2\#\\
\end{tabular}\par\vspace{3mm}
\qrcode[height=12mm]{https://helpman.komtera.lt/?fen=8/8/2B5/8/8/kR6/6nr/1K6&moves=10}\\{\sffamily\tiny\color{ink}scannen zum Lösen}
\end{minipage}
\vspace{2mm}

{\raggedright\noindent\textcolor{ink}{\sffamily\textsc{Themen}}\enspace \themetag{reines Matt} \themetag{Mustermatt} \themetag{Idealmatt} \themetag{Switchback} \themetag{Selbstblock}\par\smallskip{\footnotesize\textit{Hinweis:} die Partei am Zug steht im Diagramm im Schach; die Lösung beginnt mit einem Schlag.}\par}

\clearpage
\phantomsection\label{hm:KBBvknp}
\begin{tcolorbox}[enhanced,colback=ink,colframe=ink,arc=0pt,boxrule=0pt,left=4mm,right=4mm,top=2.5mm,bottom=2.5mm]
\color{white}\sffamily{\LARGE\bfseries KBBvknp}\hfill{\Large h\#9.5}\\[0.5mm]
{\small 6 Steine\quad Weiß am Zug}\hfill{\small Nr.~272\qquad\href{https://helpman.komtera.lt/?fen=8/6pn/8/8/8/7B/k3B3/3K4&moves=9.5}{\color{white}\bfseries lösen\,\ensuremath{\nearrow}}}
\end{tcolorbox}
\vspace{2mm}
\noindent\begin{minipage}[t]{36mm}\vspace{0pt}\sffamily\small\raggedright
\hmfigure{längste eindeutige}{{\Large h\#9.5}}\hmfigure{eindeutig bei dieser Länge}{{\Large 7} Stellungen}\hmfigure{längstes Matt}{{\Large h\#14.5}}\hmfigure{in der Klasse}{45,4 Milliarden Stellungen,\\ 15 Milliarden mit Hilfsmatt}
\end{minipage}\hfill
\begin{minipage}[t]{80mm}\vspace{0pt}\begin{tikzpicture}\begin{axis}[width=72mm,height=42mm,ybar,bar width=1.6pt,xmin=-1,xmax=30,ymin=0,ymax=1.895,axis y line*=left,axis x line*=bottom,xlabel={\scriptsize Halbzüge bis Matt},ylabel={\scriptsize Stellungen (Mrd.)},tick label style={font=\tiny},label style={font=\scriptsize},enlarge x limits=false,ylabel shift=-3pt,xlabel shift=-3pt]\addplot[fill=ink!55,draw=none] coordinates {(0,0.001464) (1,0.005936) (2,0.02161) (3,0.1319) (4,0.407) (5,0.9503) (6,1.722) (7,1.47) (8,1.682) (9,1.054) (10,1.03) (11,0.8632) (12,1.24) (13,1.11) (14,1.259) (15,0.8385) (16,0.5873) (17,0.3351) (18,0.1258) (19,0.09248) (20,0.01894) (21,0.0166) (22,0.0008695) (23,0.0008074) (24,5.564e-06) (25,2.341e-06) (27,1.61e-07) (28,9e-09) (29,3.3e-08)};\draw[densely dotted] (axis cs:19,0) -- (axis cs:19,1.895);\end{axis}\begin{axis}[width=72mm,height=42mm,axis y line*=right,axis x line=none,ymode=log,xmin=-1,xmax=30,ylabel={\scriptsize eindeutig (log)},tick label style={font=\tiny},label style={font=\scriptsize},ylabel shift=-3pt]\addplot[mark=*,mark size=0.7pt,draw=black,thin] coordinates {(0,1463724) (1,5844469) (2,14279701) (3,64964901) (4,94732925) (5,62542034) (6,45717825) (7,13923440) (8,9057199) (9,4156369) (10,4158223) (11,2328713) (12,2259743) (13,684152) (14,340952) (15,55709) (16,18449) (17,1408) (18,99) (19,7)};\end{axis}\end{tikzpicture}\end{minipage}
\vspace{3mm}

\noindent\begin{minipage}[t]{60mm}\vspace{0pt}\sbox\hmboard{\chessboard[setfen=8/6pn/8/8/8/7B/k3B3/3K4 w - - 0 1]}\usebox\hmboard\\[1mm]
\makebox[\wd\hmboard]{\sffamily\hspace{3mm}\textbf{h\#9.5}\hfill(3+3)}\end{minipage}\hfill
\begin{minipage}[t]{52mm}\vspace{0pt}
\textcolor{ink}{\sffamily\textsc{Lösung}}\par\smallskip\small
\begin{tabular}{@{}r@{\enspace}l@{\enspace}l@{}}
\textbf{1.} & \ldots & \symbishop{}ef1\\
\textbf{2.} & \symking{}b3 & \symking{}e2\\
\textbf{3.} & \symking{}c4 & \symking{}f3+\\
\textbf{4.} & \symking{}d5 & \symking{}g4\\
\textbf{5.} & \symking{}e6 & \symbishop{}d3\\
\textbf{6.} & \symking{}f7 & \symking{}f5\\
\textbf{7.} & \symking{}g8 & \symking{}e6\\
\textbf{8.} & \symknight{}g5+ & \symking{}e7\\
\textbf{9.} & \symknight{}f7 & \symbishop{}hf5\\
\textbf{10.} & \symknight{}h8 & \symbishop{}h7\#\\
\end{tabular}\par\vspace{3mm}
\qrcode[height=12mm]{https://helpman.komtera.lt/?fen=8/6pn/8/8/8/7B/k3B3/3K4&moves=9.5}\\{\sffamily\tiny\color{ink}scannen zum Lösen}
\end{minipage}
\vspace{2mm}

{\raggedright\noindent\textcolor{ink}{\sffamily\textsc{Themen}}\enspace \themetag{reines Matt} \themetag{Mustermatt} \themetag{Idealmatt} \themetag{Selbstblock} \themetag{ohne Schlag}\par}

\clearpage
\phantomsection\label{hm:KBBvkpp}
\begin{tcolorbox}[enhanced,colback=ink,colframe=ink,arc=0pt,boxrule=0pt,left=4mm,right=4mm,top=2.5mm,bottom=2.5mm]
\color{white}\sffamily{\LARGE\bfseries KBBvkpp}\hfill{\Large h\#9.5}\\[0.5mm]
{\small 6 Steine\quad Weiß am Zug}\hfill{\small Nr.~273\qquad\href{https://helpman.komtera.lt/?fen=B7/1p2p3/8/8/1K2k3/8/8/3B4&moves=9.5}{\color{white}\bfseries lösen\,\ensuremath{\nearrow}}}
\end{tcolorbox}
\vspace{2mm}
\noindent\begin{minipage}[t]{36mm}\vspace{0pt}\sffamily\small\raggedright
\hmfigure{längste eindeutige}{{\Large h\#9.5}}\hmfigure{eindeutig bei dieser Länge}{{\Large 4} Stellungen}\hmfigure{längstes Matt}{{\Large h\#16}}\hmfigure{in der Klasse}{34,1 Milliarden Stellungen,\\ 5,77 Milliarden mit Hilfsmatt}
\end{minipage}\hfill
\begin{minipage}[t]{80mm}\vspace{0pt}\begin{tikzpicture}\begin{axis}[width=72mm,height=42mm,ybar,bar width=1.6pt,xmin=-1,xmax=33,ymin=0,ymax=774.8,axis y line*=left,axis x line*=bottom,xlabel={\scriptsize Halbzüge bis Matt},ylabel={\scriptsize Stellungen (Mio.)},tick label style={font=\tiny},label style={font=\scriptsize},enlarge x limits=false,ylabel shift=-3pt,xlabel shift=-3pt]\addplot[fill=ink!55,draw=none] coordinates {(0,0.7977) (1,3.204) (2,10.55) (3,60.16) (4,183.1) (5,385) (6,704.3) (7,538.5) (8,630.9) (9,345.9) (10,334.3) (11,249.4) (12,385.7) (13,298.6) (14,403.4) (15,287.4) (16,288.9) (17,203.8) (18,154.8) (19,117.3) (20,67.2) (21,53.68) (22,25.14) (23,20.95) (24,7.107) (25,6.222) (26,1.401) (27,1.234) (28,0.1796) (29,0.1811) (30,0.000359) (32,1.2e-05)};\draw[densely dotted] (axis cs:19,0) -- (axis cs:19,774.8);\end{axis}\begin{axis}[width=72mm,height=42mm,axis y line*=right,axis x line=none,ymode=log,xmin=-1,xmax=33,ylabel={\scriptsize eindeutig (log)},tick label style={font=\tiny},label style={font=\scriptsize},ylabel shift=-3pt]\addplot[mark=*,mark size=0.7pt,draw=black,thin] coordinates {(0,797660) (1,3129297) (2,8270082) (3,34992384) (4,48448101) (5,24822866) (6,18865072) (7,4537762) (8,2809241) (9,1193166) (10,1369653) (11,766683) (12,866427) (13,202943) (14,196789) (15,17873) (16,9163) (17,1026) (18,151) (19,4)};\end{axis}\end{tikzpicture}\end{minipage}
\vspace{3mm}

\noindent\begin{minipage}[t]{60mm}\vspace{0pt}\sbox\hmboard{\chessboard[setfen=B7/1p2p3/8/8/1K2k3/8/8/3B4 w - - 0 1]}\usebox\hmboard\\[1mm]
\makebox[\wd\hmboard]{\sffamily\hspace{3mm}\textbf{h\#9.5}\hfill(3+3)}\end{minipage}\hfill
\begin{minipage}[t]{52mm}\vspace{0pt}
\textcolor{ink}{\sffamily\textsc{Lösung}}\par\smallskip\small
\begin{tabular}{@{}r@{\enspace}l@{\enspace}l@{}}
\textbf{1.} & \ldots & \symbishop{}e2\\
\textbf{2.} & e5 & \symbishop{}f1\\
\textbf{3.} & \symking{}f3 & \symbishop{}xb7+\\
\textbf{4.} & e4 & \symbishop{}c8\\
\textbf{5.} & e3 & \symking{}c3\\
\textbf{6.} & e2 & \symking{}d2\\
\textbf{7.} & exf1=\symknight{}+ & \symking{}e1\\
\textbf{8.} & \symking{}g2 & \symbishop{}h3+\\
\textbf{9.} & \symking{}h1 & \symking{}f2\\
\textbf{10.} & \symknight{}h2 & \symbishop{}g2\#\\
\end{tabular}\par\vspace{3mm}
\qrcode[height=12mm]{https://helpman.komtera.lt/?fen=B7/1p2p3/8/8/1K2k3/8/8/3B4&moves=9.5}\\{\sffamily\tiny\color{ink}scannen zum Lösen}
\end{minipage}
\vspace{2mm}

{\raggedright\noindent\textcolor{ink}{\sffamily\textsc{Themen}}\enspace \themetag{reines Matt} \themetag{Mustermatt} \themetag{Idealmatt} \themetag{Umwandlung} \themetag{Unterverwandlung} \themetag{Excelsior} \themetag{Selbstblock}\par}

\clearpage
\phantomsection\label{hm:KBBvkqn}
\begin{tcolorbox}[enhanced,colback=ink,colframe=ink,arc=0pt,boxrule=0pt,left=4mm,right=4mm,top=2.5mm,bottom=2.5mm]
\color{white}\sffamily{\LARGE\bfseries KBBvkqn}\hfill{\Large h\#9.5}\\[0.5mm]
{\small 6 Steine\quad Weiß am Zug}\hfill{\small Nr.~274\qquad\href{https://helpman.komtera.lt/?fen=8/8/4n3/8/3B4/8/1B4k1/2Kq4&moves=9.5}{\color{white}\bfseries lösen\,\ensuremath{\nearrow}}}
\end{tcolorbox}
\vspace{2mm}
\noindent\begin{minipage}[t]{36mm}\vspace{0pt}\sffamily\small\raggedright
\hmfigure{längste eindeutige}{{\Large h\#9.5}}\hmfigure{eindeutig bei dieser Länge}{{\Large 3} Stellungen}\hmfigure{längstes Matt}{{\Large h\#13}}\hmfigure{in der Klasse}{15,5 Milliarden Stellungen,\\ 4,15 Milliarden mit Hilfsmatt}
\end{minipage}\hfill
\begin{minipage}[t]{80mm}\vspace{0pt}\begin{tikzpicture}\begin{axis}[width=72mm,height=42mm,ybar,bar width=1.6pt,xmin=-1,xmax=27,ymin=0,ymax=633.9,axis y line*=left,axis x line*=bottom,xlabel={\scriptsize Halbzüge bis Matt},ylabel={\scriptsize Stellungen (Mio.)},tick label style={font=\tiny},label style={font=\scriptsize},enlarge x limits=false,ylabel shift=-3pt,xlabel shift=-3pt]\addplot[fill=ink!55,draw=none] coordinates {(0,0.1224) (1,0.5481) (2,3.129) (3,23.65) (4,80.58) (5,267.1) (6,449.7) (7,576.3) (8,428.3) (9,391.4) (10,286.7) (11,395.3) (12,369.8) (13,442.7) (14,227.8) (15,168) (16,30.41) (17,9.595) (18,0.3616) (19,1.75) (20,0.06049) (21,0.6056) (22,0.03724) (23,0.09243) (24,0.01025) (25,1e-05) (26,0.000468)};\draw[densely dotted] (axis cs:19,0) -- (axis cs:19,633.9);\end{axis}\begin{axis}[width=72mm,height=42mm,axis y line*=right,axis x line=none,ymode=log,xmin=-1,xmax=27,ylabel={\scriptsize eindeutig (log)},tick label style={font=\tiny},label style={font=\scriptsize},ylabel shift=-3pt]\addplot[mark=*,mark size=0.7pt,draw=black,thin] coordinates {(0,122357) (1,546233) (2,1747679) (3,9407690) (4,16262229) (5,23450816) (6,11470426) (7,6328000) (8,2108044) (9,1763285) (10,911185) (11,883150) (12,383056) (13,270486) (14,36165) (15,9540) (16,177) (17,172) (18,10) (19,3)};\end{axis}\end{tikzpicture}\end{minipage}
\vspace{3mm}

\noindent\begin{minipage}[t]{60mm}\vspace{0pt}\sbox\hmboard{\chessboard[setfen=8/8/4n3/8/3B4/8/1B4k1/2Kq4 w - - 0 1]}\usebox\hmboard\\[1mm]
\makebox[\wd\hmboard]{\sffamily\hspace{3mm}\textbf{h\#9.5}\hfill(3+3)}\end{minipage}\hfill
\begin{minipage}[t]{52mm}\vspace{0pt}
\textcolor{ink}{\sffamily\textsc{Lösung}}\par\smallskip\small
\begin{tabular}{@{}r@{\enspace}l@{\enspace}l@{}}
\textbf{1.} & \ldots & \symking{}xd1\\
\textbf{2.} & \symking{}f3 & \symbishop{}c1\\
\textbf{3.} & \symking{}e4 & \symking{}e2\\
\textbf{4.} & \symking{}xd4 & \symking{}f3\\
\textbf{5.} & \symking{}e5 & \symking{}g4\\
\textbf{6.} & \symking{}f6 & \symking{}h5\\
\textbf{7.} & \symking{}g7 & \symbishop{}h6+\\
\textbf{8.} & \symking{}h8 & \symking{}g6\\
\textbf{9.} & \symknight{}f8+ & \symking{}f7\\
\textbf{10.} & \symknight{}h7 & \symbishop{}g7\#\\
\end{tabular}\par\vspace{3mm}
\qrcode[height=12mm]{https://helpman.komtera.lt/?fen=8/8/4n3/8/3B4/8/1B4k1/2Kq4&moves=9.5}\\{\sffamily\tiny\color{ink}scannen zum Lösen}
\end{minipage}
\vspace{2mm}

{\raggedright\noindent\textcolor{ink}{\sffamily\textsc{Themen}}\enspace \themetag{reines Matt} \themetag{Mustermatt} \themetag{Idealmatt} \themetag{Selbstblock}\par\smallskip{\footnotesize\textit{Hinweis:} die Partei am Zug steht im Diagramm im Schach; die Lösung beginnt mit einem Schlag.}\par}

\clearpage
\phantomsection\label{hm:KBPvkbp}
\begin{tcolorbox}[enhanced,colback=ink,colframe=ink,arc=0pt,boxrule=0pt,left=4mm,right=4mm,top=2.5mm,bottom=2.5mm]
\color{white}\sffamily{\LARGE\bfseries KBPvkbp}\hfill{\Large h\#9.5}\\[0.5mm]
{\small 6 Steine\quad Weiß am Zug}\hfill{\small Nr.~275\qquad\href{https://helpman.komtera.lt/?fen=8/8/8/K1b5/8/6p1/6Pk/7B&moves=9.5}{\color{white}\bfseries lösen\,\ensuremath{\nearrow}}}
\end{tcolorbox}
\vspace{2mm}
\noindent\begin{minipage}[t]{36mm}\vspace{0pt}\sffamily\small\raggedright
\hmfigure{längste eindeutige}{{\Large h\#9.5}}\hmfigure{eindeutig bei dieser Länge}{{\Large 5} Stellungen}\hmfigure{längstes Matt}{{\Large h\#16}}\hmfigure{in der Klasse}{34,1 Milliarden Stellungen,\\ 23 Milliarden mit Hilfsmatt}
\end{minipage}\hfill
\begin{minipage}[t]{80mm}\vspace{0pt}\begin{tikzpicture}\begin{axis}[width=72mm,height=42mm,ybar,bar width=1.6pt,xmin=-1,xmax=33,ymin=0,ymax=3.379,axis y line*=left,axis x line*=bottom,xlabel={\scriptsize Halbzüge bis Matt},ylabel={\scriptsize Stellungen (Mrd.)},tick label style={font=\tiny},label style={font=\scriptsize},enlarge x limits=false,ylabel shift=-3pt,xlabel shift=-3pt]\addplot[fill=ink!55,draw=none] coordinates {(0,0.000431) (1,0.006086) (2,0.02257) (3,0.1858) (4,0.5805) (5,1.93) (6,2.64) (7,2.949) (8,3.031) (9,3.072) (10,2.979) (11,2.713) (12,2.03) (13,0.5939) (14,0.1961) (15,0.04419) (16,0.01806) (17,0.00458) (18,0.0009126) (19,0.0002555) (20,6.205e-05) (21,6.406e-06) (22,4.774e-05) (23,4.2e-08) (24,3.163e-05) (25,7.8e-08) (26,1.759e-05) (27,1e-08) (28,1.191e-05) (29,3e-09) (30,5.517e-06) (32,4e-08)};\draw[densely dotted] (axis cs:19,0) -- (axis cs:19,3.379);\end{axis}\begin{axis}[width=72mm,height=42mm,axis y line*=right,axis x line=none,ymode=log,xmin=-1,xmax=33,ylabel={\scriptsize eindeutig (log)},tick label style={font=\tiny},label style={font=\scriptsize},ylabel shift=-3pt]\addplot[mark=*,mark size=0.7pt,draw=black,thin] coordinates {(0,430970) (1,3006170) (2,10852343) (3,70902659) (4,129412422) (5,155534128) (6,99735901) (7,74071275) (8,42720911) (9,32484251) (10,16178083) (11,8295038) (12,2800007) (13,279967) (14,114795) (15,20819) (16,5383) (17,331) (18,65) (19,5)};\end{axis}\end{tikzpicture}\end{minipage}
\vspace{3mm}

\noindent\begin{minipage}[t]{60mm}\vspace{0pt}\sbox\hmboard{\chessboard[setfen=8/8/8/K1b5/8/6p1/6Pk/7B w - - 0 1]}\usebox\hmboard\\[1mm]
\makebox[\wd\hmboard]{\sffamily\hspace{3mm}\textbf{h\#9.5}\hfill(3+3)}\end{minipage}\hfill
\begin{minipage}[t]{52mm}\vspace{0pt}
\textcolor{ink}{\sffamily\textsc{Lösung}}\par\smallskip\small
\begin{tabular}{@{}r@{\enspace}l@{\enspace}l@{}}
\textbf{1.} & \ldots & \symking{}a4\\
\textbf{2.} & \symking{}g1 & \symking{}b3\\
\textbf{3.} & \symking{}f2 & \symking{}c2\\
\textbf{4.} & \symking{}e3 & \symking{}d1\\
\textbf{5.} & \symking{}f4 & \symking{}e2\\
\textbf{6.} & \symking{}g5 & \symking{}f3\\
\textbf{7.} & \symbishop{}g1 & \symking{}xg3\\
\textbf{8.} & \symbishop{}h2+ & \symking{}f3\\
\textbf{9.} & \symking{}h4 & g3+\\
\textbf{10.} & \symking{}h3 & \symbishop{}g2\#\\
\end{tabular}\par\vspace{3mm}
\qrcode[height=12mm]{https://helpman.komtera.lt/?fen=8/8/8/K1b5/8/6p1/6Pk/7B&moves=9.5}\\{\sffamily\tiny\color{ink}scannen zum Lösen}
\end{minipage}
\vspace{2mm}

{\raggedright\noindent\textcolor{ink}{\sffamily\textsc{Themen}}\enspace \themetag{reines Matt} \themetag{Mustermatt} \themetag{Idealmatt} \themetag{Switchback} \themetag{Selbstblock}\par}

\clearpage
\phantomsection\label{hm:KBPvknp}
\begin{tcolorbox}[enhanced,colback=ink,colframe=ink,arc=0pt,boxrule=0pt,left=4mm,right=4mm,top=2.5mm,bottom=2.5mm]
\color{white}\sffamily{\LARGE\bfseries KBPvknp}\hfill{\Large h\#9.5}\\[0.5mm]
{\small 6 Steine\quad Weiß am Zug}\hfill{\small Nr.~276\qquad\href{https://helpman.komtera.lt/?fen=4n3/8/8/8/8/4p3/2B1P3/k1K5&moves=9.5}{\color{white}\bfseries lösen\,\ensuremath{\nearrow}}}
\end{tcolorbox}
\vspace{2mm}
\noindent\begin{minipage}[t]{36mm}\vspace{0pt}\sffamily\small\raggedright
\hmfigure{längste eindeutige}{{\Large h\#9.5}}\hmfigure{eindeutig bei dieser Länge}{{\Large 7} Stellungen}\hmfigure{längstes Matt}{{\Large h\#14.5}}\hmfigure{in der Klasse}{34,1 Milliarden Stellungen,\\ 23,6 Milliarden mit Hilfsmatt}
\end{minipage}\hfill
\begin{minipage}[t]{80mm}\vspace{0pt}\begin{tikzpicture}\begin{axis}[width=72mm,height=42mm,ybar,bar width=1.6pt,xmin=-1,xmax=30,ymin=0,ymax=3.467,axis y line*=left,axis x line*=bottom,xlabel={\scriptsize Halbzüge bis Matt},ylabel={\scriptsize Stellungen (Mrd.)},tick label style={font=\tiny},label style={font=\scriptsize},enlarge x limits=false,ylabel shift=-3pt,xlabel shift=-3pt]\addplot[fill=ink!55,draw=none] coordinates {(0,0.0002988) (1,0.006628) (2,0.01929) (3,0.1642) (4,0.4501) (5,1.738) (6,2.393) (7,2.605) (8,2.828) (9,2.856) (10,3.151) (11,3.146) (12,2.873) (13,0.9277) (14,0.3602) (15,0.05128) (16,0.01625) (17,0.004512) (18,0.0007833) (19,0.0002747) (20,8.282e-05) (21,6.276e-06) (22,4.909e-06) (23,5.9e-08) (24,2.518e-06) (25,9.8e-08) (27,8e-09) (29,7e-09)};\draw[densely dotted] (axis cs:19,0) -- (axis cs:19,3.467);\end{axis}\begin{axis}[width=72mm,height=42mm,axis y line*=right,axis x line=none,ymode=log,xmin=-1,xmax=30,ylabel={\scriptsize eindeutig (log)},tick label style={font=\tiny},label style={font=\scriptsize},ylabel shift=-3pt]\addplot[mark=*,mark size=0.7pt,draw=black,thin] coordinates {(0,298821) (1,2570413) (2,6533143) (3,57810676) (4,94323453) (5,130595293) (6,82079957) (7,66555484) (8,48110374) (9,38421690) (10,23969823) (11,14931863) (12,5270920) (13,278180) (14,73690) (15,7815) (16,965) (17,93) (18,4) (19,7)};\end{axis}\end{tikzpicture}\end{minipage}
\vspace{3mm}

\noindent\begin{minipage}[t]{60mm}\vspace{0pt}\sbox\hmboard{\chessboard[setfen=4n3/8/8/8/8/4p3/2B1P3/k1K5 w - - 0 1]}\usebox\hmboard\\[1mm]
\makebox[\wd\hmboard]{\sffamily\hspace{3mm}\textbf{h\#9.5}\hfill(3+3)}\end{minipage}\hfill
\begin{minipage}[t]{52mm}\vspace{0pt}
\textcolor{ink}{\sffamily\textsc{Lösung}}\par\smallskip\small
\begin{tabular}{@{}r@{\enspace}l@{\enspace}l@{}}
\textbf{1.} & \ldots & \symbishop{}f5\\
\textbf{2.} & \symking{}a2 & \symking{}c2\\
\textbf{3.} & \symking{}a3 & \symking{}d3\\
\textbf{4.} & \symking{}b4 & \symking{}xe3\\
\textbf{5.} & \symking{}c5 & \symking{}f4\\
\textbf{6.} & \symking{}d6 & e4\\
\textbf{7.} & \symking{}e7 & e5\\
\textbf{8.} & \symknight{}d6 & e6\\
\textbf{9.} & \symknight{}f7 & exf7\\
\textbf{10.} & \symking{}f6 & f8=\symqueen{}\#\\
\end{tabular}\par\vspace{3mm}
\qrcode[height=12mm]{https://helpman.komtera.lt/?fen=4n3/8/8/8/8/4p3/2B1P3/k1K5&moves=9.5}\\{\sffamily\tiny\color{ink}scannen zum Lösen}
\end{minipage}
\vspace{2mm}

{\raggedright\noindent\textcolor{ink}{\sffamily\textsc{Themen}}\enspace \themetag{reines Matt} \themetag{Mustermatt} \themetag{Idealmatt} \themetag{Umwandlung} \themetag{Excelsior} \themetag{ohne Schach}\par}

\clearpage
\phantomsection\label{hm:KBvkbnn}
\begin{tcolorbox}[enhanced,colback=ink,colframe=ink,arc=0pt,boxrule=0pt,left=4mm,right=4mm,top=2.5mm,bottom=2.5mm]
\color{white}\sffamily{\LARGE\bfseries KBvkbnn}\hfill{\Large h\#9.5}\\[0.5mm]
{\small 6 Steine\quad Weiß am Zug}\hfill{\small Nr.~277\qquad\href{https://helpman.komtera.lt/?fen=8/8/8/8/8/1n6/7B/bK2n2k&moves=9.5}{\color{white}\bfseries lösen\,\ensuremath{\nearrow}}}
\end{tcolorbox}
\vspace{2mm}
\noindent\begin{minipage}[t]{36mm}\vspace{0pt}\sffamily\small\raggedright
\hmfigure{längste eindeutige}{{\Large h\#9.5}}\hmfigure{eindeutig bei dieser Länge}{{\Large 109} Stellungen}\hmfigure{längstes Matt}{{\Large h\#12.5}}\hmfigure{in der Klasse}{15,5 Milliarden Stellungen,\\ 4,86 Milliarden mit Hilfsmatt}
\end{minipage}\hfill
\begin{minipage}[t]{80mm}\vspace{0pt}\begin{tikzpicture}\begin{axis}[width=72mm,height=42mm,ybar,bar width=1.6pt,xmin=-1,xmax=26,ymin=0,ymax=859.8,axis y line*=left,axis x line*=bottom,xlabel={\scriptsize Halbzüge bis Matt},ylabel={\scriptsize Stellungen (Mio.)},tick label style={font=\tiny},label style={font=\scriptsize},enlarge x limits=false,ylabel shift=-3pt,xlabel shift=-3pt]\addplot[fill=ink!55,draw=none] coordinates {(0,0.01274) (1,0.04483) (2,0.3239) (3,1.348) (4,5.658) (5,16.62) (6,35.61) (7,83.98) (8,130.8) (9,260.9) (10,364.8) (11,610.7) (12,694.9) (13,781.7) (14,617.4) (15,536.4) (16,278.2) (17,240.9) (18,84.74) (19,87.88) (20,13.24) (21,14.76) (22,0.08916) (23,0.0954) (24,8.3e-05) (25,3e-06)};\draw[densely dotted] (axis cs:19,0) -- (axis cs:19,859.8);\end{axis}\begin{axis}[width=72mm,height=42mm,axis y line*=right,axis x line=none,ymode=log,xmin=-1,xmax=26,ylabel={\scriptsize eindeutig (log)},tick label style={font=\tiny},label style={font=\scriptsize},ylabel shift=-3pt]\addplot[mark=*,mark size=0.7pt,draw=black,thin] coordinates {(0,12738) (1,44831) (2,249878) (3,566293) (4,1449951) (5,2702165) (6,1960329) (7,2596692) (8,2215801) (9,3171437) (10,2583545) (11,2237319) (12,1120237) (13,666700) (14,169809) (15,71881) (16,18183) (17,7144) (18,2088) (19,109)};\end{axis}\end{tikzpicture}\end{minipage}
\vspace{3mm}

\noindent\begin{minipage}[t]{60mm}\vspace{0pt}\sbox\hmboard{\chessboard[setfen=8/8/8/8/8/1n6/7B/bK2n2k w - - 0 1]}\usebox\hmboard\\[1mm]
\makebox[\wd\hmboard]{\sffamily\hspace{3mm}\textbf{h\#9.5}\hfill(2+4)}\end{minipage}\hfill
\begin{minipage}[t]{52mm}\vspace{0pt}
\textcolor{ink}{\sffamily\textsc{Lösung}}\par\smallskip\small
\begin{tabular}{@{}r@{\enspace}l@{\enspace}l@{}}
\textbf{1.} & \ldots & \symbishop{}d6\\
\textbf{2.} & \symknight{}c1 & \symking{}xa1\\
\textbf{3.} & \symking{}g2 & \symking{}b2\\
\textbf{4.} & \symking{}f3 & \symking{}c3\\
\textbf{5.} & \symking{}e4 & \symking{}d2\\
\textbf{6.} & \symking{}d4 & \symking{}xe1\\
\textbf{7.} & \symking{}c3 & \symking{}d1\\
\textbf{8.} & \symking{}b2 & \symbishop{}a3+\\
\textbf{9.} & \symking{}a1 & \symking{}c2\\
\textbf{10.} & \symknight{}a2 & \symbishop{}b2\#\\
\end{tabular}\par\vspace{3mm}
\qrcode[height=12mm]{https://helpman.komtera.lt/?fen=8/8/8/8/8/1n6/7B/bK2n2k&moves=9.5}\\{\sffamily\tiny\color{ink}scannen zum Lösen}
\end{minipage}
\vspace{2mm}

{\raggedright\noindent\textcolor{ink}{\sffamily\textsc{Themen}}\enspace \themetag{reines Matt} \themetag{Mustermatt} \themetag{Idealmatt} \themetag{Selbstblock} \themetag{Kniest}\par}

\clearpage
\phantomsection\label{hm:KBvknnn}
\begin{tcolorbox}[enhanced,colback=ink,colframe=ink,arc=0pt,boxrule=0pt,left=4mm,right=4mm,top=2.5mm,bottom=2.5mm]
\color{white}\sffamily{\LARGE\bfseries KBvknnn}\hfill{\Large h\#9.5}\\[0.5mm]
{\small 6 Steine\quad Weiß am Zug}\hfill{\small Nr.~278\qquad\href{https://helpman.komtera.lt/?fen=8/8/8/8/8/1n6/7B/nK2n2k&moves=9.5}{\color{white}\bfseries lösen\,\ensuremath{\nearrow}}}
\end{tcolorbox}
\vspace{2mm}
\noindent\begin{minipage}[t]{36mm}\vspace{0pt}\sffamily\small\raggedright
\hmfigure{längste eindeutige}{{\Large h\#9.5}}\hmfigure{eindeutig bei dieser Länge}{{\Large 35} Stellungen}\hmfigure{längstes Matt}{{\Large h\#12}}\hmfigure{in der Klasse}{15,5 Milliarden Stellungen,\\ 1,66 Milliarden mit Hilfsmatt}
\end{minipage}\hfill
\begin{minipage}[t]{80mm}\vspace{0pt}\begin{tikzpicture}\begin{axis}[width=72mm,height=42mm,ybar,bar width=1.6pt,xmin=-1,xmax=25,ymin=0,ymax=302.1,axis y line*=left,axis x line*=bottom,xlabel={\scriptsize Halbzüge bis Matt},ylabel={\scriptsize Stellungen (Mio.)},tick label style={font=\tiny},label style={font=\scriptsize},enlarge x limits=false,ylabel shift=-3pt,xlabel shift=-3pt]\addplot[fill=ink!55,draw=none] coordinates {(0,0.001275) (1,0.002891) (2,0.008747) (3,0.07734) (4,0.3199) (5,1.672) (6,4.289) (7,13.82) (8,25.26) (9,54.96) (10,81.25) (11,144.9) (12,192.5) (13,268.8) (14,274.7) (15,237.7) (16,154.7) (17,113) (18,40.85) (19,39.28) (20,3.955) (21,4.235) (22,0.003615) (23,0.000961) (24,3e-06)};\draw[densely dotted] (axis cs:19,0) -- (axis cs:19,302.1);\end{axis}\begin{axis}[width=72mm,height=42mm,axis y line*=right,axis x line=none,ymode=log,xmin=-1,xmax=25,ylabel={\scriptsize eindeutig (log)},tick label style={font=\tiny},label style={font=\scriptsize},ylabel shift=-3pt]\addplot[mark=*,mark size=0.7pt,draw=black,thin] coordinates {(0,1275) (1,2891) (2,5371) (3,39954) (4,170225) (5,420241) (6,326512) (7,517316) (8,642224) (9,921039) (10,923121) (11,814854) (12,534341) (13,283297) (14,72916) (15,30794) (16,7293) (17,3430) (18,724) (19,35)};\end{axis}\end{tikzpicture}\end{minipage}
\vspace{3mm}

\noindent\begin{minipage}[t]{60mm}\vspace{0pt}\sbox\hmboard{\chessboard[setfen=8/8/8/8/8/1n6/7B/nK2n2k w - - 0 1]}\usebox\hmboard\\[1mm]
\makebox[\wd\hmboard]{\sffamily\hspace{3mm}\textbf{h\#9.5}\hfill(2+4)}\end{minipage}\hfill
\begin{minipage}[t]{52mm}\vspace{0pt}
\textcolor{ink}{\sffamily\textsc{Lösung}}\par\smallskip\small
\begin{tabular}{@{}r@{\enspace}l@{\enspace}l@{}}
\textbf{1.} & \ldots & \symbishop{}d6\\
\textbf{2.} & \symknight{}c1 & \symking{}xa1\\
\textbf{3.} & \symking{}g2 & \symking{}b2\\
\textbf{4.} & \symking{}f3 & \symking{}c3\\
\textbf{5.} & \symking{}e4 & \symking{}d2\\
\textbf{6.} & \symking{}d4 & \symking{}xe1\\
\textbf{7.} & \symking{}c3 & \symking{}d1\\
\textbf{8.} & \symking{}b2 & \symbishop{}a3+\\
\textbf{9.} & \symking{}a1 & \symking{}c2\\
\textbf{10.} & \symknight{}a2 & \symbishop{}b2\#\\
\end{tabular}\par\vspace{3mm}
\qrcode[height=12mm]{https://helpman.komtera.lt/?fen=8/8/8/8/8/1n6/7B/nK2n2k&moves=9.5}\\{\sffamily\tiny\color{ink}scannen zum Lösen}
\end{minipage}
\vspace{2mm}

{\raggedright\noindent\textcolor{ink}{\sffamily\textsc{Themen}}\enspace \themetag{reines Matt} \themetag{Mustermatt} \themetag{Idealmatt} \themetag{Selbstblock} \themetag{Kniest}\par}

\clearpage
\phantomsection\label{hm:KBvknnp}
\begin{tcolorbox}[enhanced,colback=ink,colframe=ink,arc=0pt,boxrule=0pt,left=4mm,right=4mm,top=2.5mm,bottom=2.5mm]
\color{white}\sffamily{\LARGE\bfseries KBvknnp}\hfill{\Large h\#9.5}\\[0.5mm]
{\small 6 Steine\quad Weiß am Zug}\hfill{\small Nr.~279\qquad\href{https://helpman.komtera.lt/?fen=3n4/2p4k/8/8/1n6/8/B7/K7&moves=9.5}{\color{white}\bfseries lösen\,\ensuremath{\nearrow}}}
\end{tcolorbox}
\vspace{2mm}
\noindent\begin{minipage}[t]{36mm}\vspace{0pt}\sffamily\small\raggedright
\hmfigure{längste eindeutige}{{\Large h\#9.5}}\hmfigure{eindeutig bei dieser Länge}{{\Large 209} Stellungen}\hmfigure{längstes Matt}{{\Large h\#12.5}}\hmfigure{in der Klasse}{45,4 Milliarden Stellungen,\\ 15,4 Milliarden mit Hilfsmatt}
\end{minipage}\hfill
\begin{minipage}[t]{80mm}\vspace{0pt}\begin{tikzpicture}\begin{axis}[width=72mm,height=42mm,ybar,bar width=1.6pt,xmin=-1,xmax=26,ymin=0,ymax=2.783,axis y line*=left,axis x line*=bottom,xlabel={\scriptsize Halbzüge bis Matt},ylabel={\scriptsize Stellungen (Mrd.)},tick label style={font=\tiny},label style={font=\scriptsize},enlarge x limits=false,ylabel shift=-3pt,xlabel shift=-3pt]\addplot[fill=ink!55,draw=none] coordinates {(0,2.975e-05) (1,9.943e-05) (2,0.0002846) (3,0.001333) (4,0.004463) (5,0.01757) (6,0.04616) (7,0.1295) (8,0.2557) (9,0.5099) (10,0.8162) (11,1.335) (12,1.904) (13,2.356) (14,2.53) (15,2.157) (16,1.451) (17,1.011) (18,0.3769) (19,0.3502) (20,0.05012) (21,0.0533) (22,0.0002881) (23,0.0003025) (24,1.92e-07) (25,1.3e-08)};\draw[densely dotted] (axis cs:19,0) -- (axis cs:19,2.783);\end{axis}\begin{axis}[width=72mm,height=42mm,axis y line*=right,axis x line=none,ymode=log,xmin=-1,xmax=26,ylabel={\scriptsize eindeutig (log)},tick label style={font=\tiny},label style={font=\scriptsize},ylabel shift=-3pt]\addplot[mark=*,mark size=0.7pt,draw=black,thin] coordinates {(0,29754) (1,99430) (2,127322) (3,482792) (4,1715176) (5,3660637) (6,3663815) (7,5113093) (8,6030389) (9,7859724) (10,8403543) (11,6827131) (12,5036735) (13,2428029) (14,739447) (15,242183) (16,60000) (17,17379) (18,4543) (19,209)};\end{axis}\end{tikzpicture}\end{minipage}
\vspace{3mm}

\noindent\begin{minipage}[t]{60mm}\vspace{0pt}\sbox\hmboard{\chessboard[setfen=3n4/2p4k/8/8/1n6/8/B7/K7 w - - 0 1]}\usebox\hmboard\\[1mm]
\makebox[\wd\hmboard]{\sffamily\hspace{3mm}\textbf{h\#9.5}\hfill(2+4)}\end{minipage}\hfill
\begin{minipage}[t]{52mm}\vspace{0pt}
\textcolor{ink}{\sffamily\textsc{Lösung}}\par\smallskip\small
\begin{tabular}{@{}r@{\enspace}l@{\enspace}l@{}}
\textbf{1.} & \ldots & \symking{}b2\\
\textbf{2.} & \symking{}g7 & \symking{}c3\\
\textbf{3.} & \symking{}f8 & \symking{}d4\\
\textbf{4.} & \symking{}e8 & \symking{}e5\\
\textbf{5.} & \symknight{}dc6+ & \symking{}e6\\
\textbf{6.} & \symking{}d8 & \symbishop{}c4\\
\textbf{7.} & \symking{}c8 & \symbishop{}a6+\\
\textbf{8.} & \symking{}b8 & \symking{}d7\\
\textbf{9.} & \symking{}a8 & \symking{}xc7\\
\textbf{10.} & \symknight{}a7 & \symbishop{}b7\#\\
\end{tabular}\par\vspace{3mm}
\qrcode[height=12mm]{https://helpman.komtera.lt/?fen=3n4/2p4k/8/8/1n6/8/B7/K7&moves=9.5}\\{\sffamily\tiny\color{ink}scannen zum Lösen}
\end{minipage}
\vspace{2mm}

{\raggedright\noindent\textcolor{ink}{\sffamily\textsc{Themen}}\enspace \themetag{reines Matt} \themetag{Mustermatt} \themetag{Selbstblock}\par}

\clearpage
\phantomsection\label{hm:KBvkqbn}
\begin{tcolorbox}[enhanced,colback=ink,colframe=ink,arc=0pt,boxrule=0pt,left=4mm,right=4mm,top=2.5mm,bottom=2.5mm]
\color{white}\sffamily{\LARGE\bfseries KBvkqbn}\hfill{\Large h\#9.5}\\[0.5mm]
{\small 6 Steine\quad Weiß am Zug}\hfill{\small Nr.~280\qquad\href{https://helpman.komtera.lt/?fen=7n/8/2B5/3b4/8/2k5/2q5/1K6&moves=9.5}{\color{white}\bfseries lösen\,\ensuremath{\nearrow}}}
\end{tcolorbox}
\vspace{2mm}
\noindent\begin{minipage}[t]{36mm}\vspace{0pt}\sffamily\small\raggedright
\hmfigure{längste eindeutige}{{\Large h\#9.5}}\hmfigure{eindeutig bei dieser Länge}{{\Large 90} Stellungen}\hmfigure{längstes Matt}{{\Large h\#12.5}}\hmfigure{in der Klasse}{15,5 Milliarden Stellungen,\\ 8,49 Milliarden mit Hilfsmatt}
\end{minipage}\hfill
\begin{minipage}[t]{80mm}\vspace{0pt}\begin{tikzpicture}\begin{axis}[width=72mm,height=42mm,ybar,bar width=1.6pt,xmin=-1,xmax=26,ymin=0,ymax=1.747,axis y line*=left,axis x line*=bottom,xlabel={\scriptsize Halbzüge bis Matt},ylabel={\scriptsize Stellungen (Mrd.)},tick label style={font=\tiny},label style={font=\scriptsize},enlarge x limits=false,ylabel shift=-3pt,xlabel shift=-3pt]\addplot[fill=ink!55,draw=none] coordinates {(0,9.149e-06) (1,3.511e-05) (2,0.000314) (3,0.0017) (4,0.007644) (5,0.0304) (6,0.07825) (7,0.2537) (8,0.4059) (9,0.9937) (10,0.9004) (11,1.588) (12,0.9373) (13,1.216) (14,0.6854) (15,0.9174) (16,0.2337) (17,0.215) (18,0.009861) (19,0.01746) (20,6.615e-05) (21,0.001805) (22,5.9e-08) (23,0.0001395) (24,1e-09) (25,5.51e-07)};\draw[densely dotted] (axis cs:19,0) -- (axis cs:19,1.747);\end{axis}\begin{axis}[width=72mm,height=42mm,axis y line*=right,axis x line=none,ymode=log,xmin=-1,xmax=26,ylabel={\scriptsize eindeutig (log)},tick label style={font=\tiny},label style={font=\scriptsize},ylabel shift=-3pt]\addplot[mark=*,mark size=0.7pt,draw=black,thin] coordinates {(0,9149) (1,35108) (2,227493) (3,904936) (4,2238176) (5,5791020) (6,5691357) (7,12041185) (8,8418059) (9,12374751) (10,4509985) (11,4978298) (12,1599006) (13,2062643) (14,449328) (15,311088) (16,7195) (17,3881) (18,15) (19,90)};\end{axis}\end{tikzpicture}\end{minipage}
\vspace{3mm}

\noindent\begin{minipage}[t]{60mm}\vspace{0pt}\sbox\hmboard{\chessboard[setfen=7n/8/2B5/3b4/8/2k5/2q5/1K6 w - - 0 1]}\usebox\hmboard\\[1mm]
\makebox[\wd\hmboard]{\sffamily\hspace{3mm}\textbf{h\#9.5}\hfill(2+4)}\end{minipage}\hfill
\begin{minipage}[t]{52mm}\vspace{0pt}
\textcolor{ink}{\sffamily\textsc{Lösung}}\par\smallskip\small
\begin{tabular}{@{}r@{\enspace}l@{\enspace}l@{}}
\textbf{1.} & \ldots & \symking{}a1\\
\textbf{2.} & \symking{}b4 & \symbishop{}b5\\
\textbf{3.} & \symqueen{}c5 & \symking{}b2\\
\textbf{4.} & \symbishop{}b3 & \symbishop{}c4\\
\textbf{5.} & \symking{}a4 & \symking{}c3\\
\textbf{6.} & \symqueen{}a3 & \symking{}d4\\
\textbf{7.} & \symknight{}g6 & \symking{}d5\\
\textbf{8.} & \symknight{}f4+ & \symking{}c6\\
\textbf{9.} & \symknight{}d3 & \symking{}b6\\
\textbf{10.} & \symknight{}b4 & \symbishop{}b5\#\\
\end{tabular}\par\vspace{3mm}
\qrcode[height=12mm]{https://helpman.komtera.lt/?fen=7n/8/2B5/3b4/8/2k5/2q5/1K6&moves=9.5}\\{\sffamily\tiny\color{ink}scannen zum Lösen}
\end{minipage}
\vspace{2mm}

{\raggedright\noindent\textcolor{ink}{\sffamily\textsc{Themen}}\enspace \themetag{reines Matt} \themetag{Mustermatt} \themetag{Idealmatt} \themetag{Switchback} \themetag{Selbstblock} \themetag{ohne Schlag}\par\smallskip{\footnotesize\textit{Hinweis:} die Partei am Zug steht im Diagramm im Schach.}\par}

\clearpage
\phantomsection\label{hm:KBvkqqn}
\begin{tcolorbox}[enhanced,colback=ink,colframe=ink,arc=0pt,boxrule=0pt,left=4mm,right=4mm,top=2.5mm,bottom=2.5mm]
\color{white}\sffamily{\LARGE\bfseries KBvkqqn}\hfill{\Large h\#9.5}\\[0.5mm]
{\small 6 Steine\quad Weiß am Zug}\hfill{\small Nr.~281\qquad\href{https://helpman.komtera.lt/?fen=8/1q6/1n6/8/8/8/q7/KBk5&moves=9.5}{\color{white}\bfseries lösen\,\ensuremath{\nearrow}}}
\end{tcolorbox}
\vspace{2mm}
\noindent\begin{minipage}[t]{36mm}\vspace{0pt}\sffamily\small\raggedright
\hmfigure{längste eindeutige}{{\Large h\#9.5}}\hmfigure{eindeutig bei dieser Länge}{{\Large 341} Stellungen}\hmfigure{längstes Matt}{{\Large h\#12.5}}\hmfigure{in der Klasse}{15,5 Milliarden Stellungen,\\ 3,77 Milliarden mit Hilfsmatt}
\end{minipage}\hfill
\begin{minipage}[t]{80mm}\vspace{0pt}\begin{tikzpicture}\begin{axis}[width=72mm,height=42mm,ybar,bar width=1.6pt,xmin=-1,xmax=26,ymin=0,ymax=998.2,axis y line*=left,axis x line*=bottom,xlabel={\scriptsize Halbzüge bis Matt},ylabel={\scriptsize Stellungen (Mio.)},tick label style={font=\tiny},label style={font=\scriptsize},enlarge x limits=false,ylabel shift=-3pt,xlabel shift=-3pt]\addplot[fill=ink!55,draw=none] coordinates {(0,0.000383) (1,0.001005) (2,0.004201) (3,0.03551) (4,0.1576) (5,0.9359) (6,2.945) (7,11.86) (8,27.6) (9,88.67) (10,138.1) (11,379.3) (12,376.5) (13,884.4) (14,502.6) (15,907.4) (16,180.2) (17,234.1) (18,5.172) (19,21.89) (20,0.002922) (21,5.427) (22,9.3e-05) (23,0.7184) (24,3e-06) (25,0.005678)};\draw[densely dotted] (axis cs:19,0) -- (axis cs:19,998.2);\end{axis}\begin{axis}[width=72mm,height=42mm,axis y line*=right,axis x line=none,ymode=log,xmin=-1,xmax=26,ylabel={\scriptsize eindeutig (log)},tick label style={font=\tiny},label style={font=\scriptsize},ylabel shift=-3pt]\addplot[mark=*,mark size=0.7pt,draw=black,thin] coordinates {(0,383) (1,1005) (2,2857) (3,19701) (4,79242) (5,307289) (6,374388) (7,884149) (8,796594) (9,1983038) (10,1324008) (11,2766309) (12,1226309) (13,2370830) (14,507149) (15,625483) (16,19348) (17,16823) (18,3) (19,341)};\end{axis}\end{tikzpicture}\end{minipage}
\vspace{3mm}

\noindent\begin{minipage}[t]{60mm}\vspace{0pt}\sbox\hmboard{\chessboard[setfen=8/1q6/1n6/8/8/8/q7/KBk5 w - - 0 1]}\usebox\hmboard\\[1mm]
\makebox[\wd\hmboard]{\sffamily\hspace{3mm}\textbf{h\#9.5}\hfill(2+4)}\end{minipage}\hfill
\begin{minipage}[t]{52mm}\vspace{0pt}
\textcolor{ink}{\sffamily\textsc{Lösung}}\par\smallskip\small
\begin{tabular}{@{}r@{\enspace}l@{\enspace}l@{}}
\textbf{1.} & \ldots & \symking{}xa2\\
\textbf{2.} & \symking{}d2 & \symbishop{}e4\\
\textbf{3.} & \symking{}c3 & \symbishop{}xb7\\
\textbf{4.} & \symking{}b4 & \symbishop{}a6\\
\textbf{5.} & \symking{}a5 & \symking{}b3\\
\textbf{6.} & \symknight{}c8 & \symking{}c4\\
\textbf{7.} & \symking{}b6 & \symking{}d5\\
\textbf{8.} & \symking{}a7 & \symking{}c6\\
\textbf{9.} & \symking{}a8 & \symking{}c7\\
\textbf{10.} & \symknight{}a7 & \symbishop{}b7\#\\
\end{tabular}\par\vspace{3mm}
\qrcode[height=12mm]{https://helpman.komtera.lt/?fen=8/1q6/1n6/8/8/8/q7/KBk5&moves=9.5}\\{\sffamily\tiny\color{ink}scannen zum Lösen}
\end{minipage}
\vspace{2mm}

{\raggedright\noindent\textcolor{ink}{\sffamily\textsc{Themen}}\enspace \themetag{reines Matt} \themetag{Mustermatt} \themetag{Idealmatt} \themetag{Switchback} \themetag{Selbstblock} \themetag{ohne Schach}\par\smallskip{\footnotesize\textit{Hinweis:} die Partei am Zug steht im Diagramm im Schach; die Lösung beginnt mit einem Schlag.}\par}

\clearpage
\phantomsection\label{hm:KBvkqrn}
\begin{tcolorbox}[enhanced,colback=ink,colframe=ink,arc=0pt,boxrule=0pt,left=4mm,right=4mm,top=2.5mm,bottom=2.5mm]
\color{white}\sffamily{\LARGE\bfseries KBvkqrn}\hfill{\Large h\#9.5}\\[0.5mm]
{\small 6 Steine\quad Weiß am Zug}\hfill{\small Nr.~282\qquad\href{https://helpman.komtera.lt/?fen=n7/8/4r3/8/8/5k2/2qB4/K7&moves=9.5}{\color{white}\bfseries lösen\,\ensuremath{\nearrow}}}
\end{tcolorbox}
\vspace{2mm}
\noindent\begin{minipage}[t]{36mm}\vspace{0pt}\sffamily\small\raggedright
\hmfigure{längste eindeutige}{{\Large h\#9.5}}\hmfigure{eindeutig bei dieser Länge}{{\Large 394} Stellungen}\hmfigure{längstes Matt}{{\Large h\#12.5}}\hmfigure{in der Klasse}{15,5 Milliarden Stellungen,\\ 8,15 Milliarden mit Hilfsmatt}
\end{minipage}\hfill
\begin{minipage}[t]{80mm}\vspace{0pt}\begin{tikzpicture}\begin{axis}[width=72mm,height=42mm,ybar,bar width=1.6pt,xmin=-1,xmax=26,ymin=0,ymax=2.083,axis y line*=left,axis x line*=bottom,xlabel={\scriptsize Halbzüge bis Matt},ylabel={\scriptsize Stellungen (Mrd.)},tick label style={font=\tiny},label style={font=\scriptsize},enlarge x limits=false,ylabel shift=-3pt,xlabel shift=-3pt]\addplot[fill=ink!55,draw=none] coordinates {(0,1.15e-06) (1,2.838e-06) (2,1.045e-05) (3,8.8e-05) (4,0.0003881) (5,0.002262) (6,0.006974) (7,0.02669) (8,0.06326) (9,0.183) (10,0.3084) (11,0.7456) (12,0.8591) (13,1.761) (14,1.214) (15,1.893) (16,0.4882) (17,0.5269) (18,0.0214) (19,0.04054) (20,3.572e-05) (21,0.005302) (22,7.8e-08) (23,0.0005055) (24,6e-09) (25,3.591e-06)};\draw[densely dotted] (axis cs:19,0) -- (axis cs:19,2.083);\end{axis}\begin{axis}[width=72mm,height=42mm,axis y line*=right,axis x line=none,ymode=log,xmin=-1,xmax=26,ylabel={\scriptsize eindeutig (log)},tick label style={font=\tiny},label style={font=\scriptsize},ylabel shift=-3pt]\addplot[mark=*,mark size=0.7pt,draw=black,thin] coordinates {(0,1150) (1,2838) (2,7018) (3,48284) (4,211044) (5,739609) (6,876685) (7,1853181) (8,1809842) (9,3712516) (10,2820379) (11,4650181) (12,2605788) (13,4006931) (14,1013803) (15,999901) (16,47076) (17,28633) (18,73) (19,394)};\end{axis}\end{tikzpicture}\end{minipage}
\vspace{3mm}

\noindent\begin{minipage}[t]{60mm}\vspace{0pt}\sbox\hmboard{\chessboard[setfen=n7/8/4r3/8/8/5k2/2qB4/K7 w - - 0 1]}\usebox\hmboard\\[1mm]
\makebox[\wd\hmboard]{\sffamily\hspace{3mm}\textbf{h\#9.5}\hfill(2+4)}\end{minipage}\hfill
\begin{minipage}[t]{52mm}\vspace{0pt}
\textcolor{ink}{\sffamily\textsc{Lösung}}\par\smallskip\small
\begin{tabular}{@{}r@{\enspace}l@{\enspace}l@{}}
\textbf{1.} & \ldots & \symbishop{}e3\\
\textbf{2.} & \symqueen{}d1+ & \symking{}b2\\
\textbf{3.} & \symknight{}c7 & \symking{}c3\\
\textbf{4.} & \symknight{}b5+ & \symking{}c4\\
\textbf{5.} & \symknight{}d4 & \symking{}d5\\
\textbf{6.} & \symrook{}d6+ & \symking{}e5\\
\textbf{7.} & \symking{}e2 & \symking{}f4\\
\textbf{8.} & \symking{}e1 & \symking{}g3\\
\textbf{9.} & \symknight{}e2+ & \symking{}g2\\
\textbf{10.} & \symrook{}d2 & \symbishop{}f2\#\\
\end{tabular}\par\vspace{3mm}
\qrcode[height=12mm]{https://helpman.komtera.lt/?fen=n7/8/4r3/8/8/5k2/2qB4/K7&moves=9.5}\\{\sffamily\tiny\color{ink}scannen zum Lösen}
\end{minipage}
\vspace{2mm}

{\raggedright\noindent\textcolor{ink}{\sffamily\textsc{Themen}}\enspace \themetag{Satzspiel} \themetag{reines Matt} \themetag{Mustermatt} \themetag{Idealmatt} \themetag{Selbstblock} \themetag{ohne Schlag}\par}

\clearpage
\phantomsection\label{hm:KBvkrbn}
\begin{tcolorbox}[enhanced,colback=ink,colframe=ink,arc=0pt,boxrule=0pt,left=4mm,right=4mm,top=2.5mm,bottom=2.5mm]
\color{white}\sffamily{\LARGE\bfseries KBvkrbn}\hfill{\Large h\#9.5}\\[0.5mm]
{\small 6 Steine\quad Weiß am Zug}\hfill{\small Nr.~283\qquad\href{https://helpman.komtera.lt/?fen=8/8/1n6/3b4/r7/8/8/K1kB4&moves=9.5}{\color{white}\bfseries lösen\,\ensuremath{\nearrow}}}
\end{tcolorbox}
\vspace{2mm}
\noindent\begin{minipage}[t]{36mm}\vspace{0pt}\sffamily\small\raggedright
\hmfigure{längste eindeutige}{{\Large h\#9.5}}\hmfigure{eindeutig bei dieser Länge}{{\Large 27} Stellungen}\hmfigure{längstes Matt}{{\Large h\#12.5}}\hmfigure{in der Klasse}{15,5 Milliarden Stellungen,\\ 9,17 Milliarden mit Hilfsmatt}
\end{minipage}\hfill
\begin{minipage}[t]{80mm}\vspace{0pt}\begin{tikzpicture}\begin{axis}[width=72mm,height=42mm,ybar,bar width=1.6pt,xmin=-1,xmax=26,ymin=0,ymax=1.926,axis y line*=left,axis x line*=bottom,xlabel={\scriptsize Halbzüge bis Matt},ylabel={\scriptsize Stellungen (Mrd.)},tick label style={font=\tiny},label style={font=\scriptsize},enlarge x limits=false,ylabel shift=-3pt,xlabel shift=-3pt]\addplot[fill=ink!55,draw=none] coordinates {(0,1.339e-05) (1,5.096e-05) (2,0.0004126) (3,0.002036) (4,0.009465) (5,0.03269) (6,0.08831) (7,0.2443) (8,0.458) (9,0.9819) (10,1.129) (11,1.751) (12,1.246) (13,1.318) (14,0.7735) (15,0.7615) (16,0.188) (17,0.1601) (18,0.009917) (19,0.01153) (20,0.000183) (21,0.000353) (22,1.8e-08) (23,5.469e-06) (25,1e-09)};\draw[densely dotted] (axis cs:19,0) -- (axis cs:19,1.926);\end{axis}\begin{axis}[width=72mm,height=42mm,axis y line*=right,axis x line=none,ymode=log,xmin=-1,xmax=26,ylabel={\scriptsize eindeutig (log)},tick label style={font=\tiny},label style={font=\scriptsize},ylabel shift=-3pt]\addplot[mark=*,mark size=0.7pt,draw=black,thin] coordinates {(0,13392) (1,50963) (2,313630) (3,1002588) (4,2856148) (5,5980128) (6,6817264) (7,10657046) (8,12371574) (9,13543211) (10,6967194) (11,5474465) (12,1962086) (13,1445523) (14,324039) (15,144207) (16,4016) (17,1649) (18,3) (19,27)};\end{axis}\end{tikzpicture}\end{minipage}
\vspace{3mm}

\noindent\begin{minipage}[t]{60mm}\vspace{0pt}\sbox\hmboard{\chessboard[setfen=8/8/1n6/3b4/r7/8/8/K1kB4 w - - 0 1]}\usebox\hmboard\\[1mm]
\makebox[\wd\hmboard]{\sffamily\hspace{3mm}\textbf{h\#9.5}\hfill(2+4)}\end{minipage}\hfill
\begin{minipage}[t]{52mm}\vspace{0pt}
\textcolor{ink}{\sffamily\textsc{Lösung}}\par\smallskip\small
\begin{tabular}{@{}r@{\enspace}l@{\enspace}l@{}}
\textbf{1.} & \ldots & \symbishop{}xa4\\
\textbf{2.} & \symking{}d2 & \symbishop{}b5\\
\textbf{3.} & \symking{}c3 & \symbishop{}a6\\
\textbf{4.} & \symking{}b4 & \symking{}b2\\
\textbf{5.} & \symking{}a5 & \symking{}c3\\
\textbf{6.} & \symknight{}c8 & \symking{}d4\\
\textbf{7.} & \symking{}b6 & \symking{}xd5\\
\textbf{8.} & \symking{}a7 & \symking{}c6\\
\textbf{9.} & \symking{}a8 & \symking{}c7\\
\textbf{10.} & \symknight{}a7 & \symbishop{}b7\#\\
\end{tabular}\par\vspace{3mm}
\qrcode[height=12mm]{https://helpman.komtera.lt/?fen=8/8/1n6/3b4/r7/8/8/K1kB4&moves=9.5}\\{\sffamily\tiny\color{ink}scannen zum Lösen}
\end{minipage}
\vspace{2mm}

{\raggedright\noindent\textcolor{ink}{\sffamily\textsc{Themen}}\enspace \themetag{reines Matt} \themetag{Mustermatt} \themetag{Idealmatt} \themetag{Selbstblock} \themetag{ohne Schach}\par\smallskip{\footnotesize\textit{Hinweis:} die Partei am Zug steht im Diagramm im Schach; die Lösung beginnt mit einem Schlag.}\par}

\clearpage
\phantomsection\label{hm:KBvkrnn}
\begin{tcolorbox}[enhanced,colback=ink,colframe=ink,arc=0pt,boxrule=0pt,left=4mm,right=4mm,top=2.5mm,bottom=2.5mm]
\color{white}\sffamily{\LARGE\bfseries KBvkrnn}\hfill{\Large h\#9.5}\\[0.5mm]
{\small 6 Steine\quad Weiß am Zug}\hfill{\small Nr.~284\qquad\href{https://helpman.komtera.lt/?fen=3n1n2/8/8/8/5Br1/8/8/K6k&moves=9.5}{\color{white}\bfseries lösen\,\ensuremath{\nearrow}}}
\end{tcolorbox}
\vspace{2mm}
\noindent\begin{minipage}[t]{36mm}\vspace{0pt}\sffamily\small\raggedright
\hmfigure{längste eindeutige}{{\Large h\#9.5}}\hmfigure{eindeutig bei dieser Länge}{{\Large 36} Stellungen}\hmfigure{längstes Matt}{{\Large h\#13}}\hmfigure{in der Klasse}{15,5 Milliarden Stellungen,\\ 4,68 Milliarden mit Hilfsmatt}
\end{minipage}\hfill
\begin{minipage}[t]{80mm}\vspace{0pt}\begin{tikzpicture}\begin{axis}[width=72mm,height=42mm,ybar,bar width=1.6pt,xmin=-1,xmax=27,ymin=0,ymax=1003,axis y line*=left,axis x line*=bottom,xlabel={\scriptsize Halbzüge bis Matt},ylabel={\scriptsize Stellungen (Mio.)},tick label style={font=\tiny},label style={font=\scriptsize},enlarge x limits=false,ylabel shift=-3pt,xlabel shift=-3pt]\addplot[fill=ink!55,draw=none] coordinates {(0,0.002094) (1,0.004814) (2,0.0151) (3,0.1303) (4,0.5536) (5,3.048) (6,8.365) (7,28.48) (8,58.47) (9,139.2) (10,226.4) (11,436.7) (12,582.7) (13,912) (14,803.9) (15,823.8) (16,320.4) (17,240.3) (18,41.05) (19,43.58) (20,4.366) (21,5.499) (22,0.06735) (23,0.0725) (24,0.00022) (25,0.000165) (26,3e-06)};\draw[densely dotted] (axis cs:19,0) -- (axis cs:19,1003);\end{axis}\begin{axis}[width=72mm,height=42mm,axis y line*=right,axis x line=none,ymode=log,xmin=-1,xmax=27,ylabel={\scriptsize eindeutig (log)},tick label style={font=\tiny},label style={font=\scriptsize},ylabel shift=-3pt]\addplot[mark=*,mark size=0.7pt,draw=black,thin] coordinates {(0,2094) (1,4814) (2,9652) (3,69190) (4,310718) (5,915056) (6,880892) (7,1483513) (8,1603278) (9,2409596) (10,2178420) (11,2364849) (12,1669547) (13,1446661) (14,378265) (15,155209) (16,13150) (17,4055) (18,177) (19,36)};\end{axis}\end{tikzpicture}\end{minipage}
\vspace{3mm}

\noindent\begin{minipage}[t]{60mm}\vspace{0pt}\sbox\hmboard{\chessboard[setfen=3n1n2/8/8/8/5Br1/8/8/K6k w - - 0 1]}\usebox\hmboard\\[1mm]
\makebox[\wd\hmboard]{\sffamily\hspace{3mm}\textbf{h\#9.5}\hfill(2+4)}\end{minipage}\hfill
\begin{minipage}[t]{52mm}\vspace{0pt}
\textcolor{ink}{\sffamily\textsc{Lösung}}\par\smallskip\small
\begin{tabular}{@{}r@{\enspace}l@{\enspace}l@{}}
\textbf{1.} & \ldots & \symking{}b2\\
\textbf{2.} & \symking{}g2 & \symking{}c3\\
\textbf{3.} & \symking{}h3 & \symking{}d4\\
\textbf{4.} & \symking{}h4 & \symking{}e5\\
\textbf{5.} & \symking{}h5 & \symbishop{}g5\\
\textbf{6.} & \symking{}g6 & \symbishop{}xd8\\
\textbf{7.} & \symking{}h7 & \symking{}f6\\
\textbf{8.} & \symking{}h8 & \symking{}f7\\
\textbf{9.} & \symknight{}h7 & \symbishop{}f6+\\
\textbf{10.} & \symrook{}g7+ & \symbishop{}xg7\#\\
\end{tabular}\par\vspace{3mm}
\qrcode[height=12mm]{https://helpman.komtera.lt/?fen=3n1n2/8/8/8/5Br1/8/8/K6k&moves=9.5}\\{\sffamily\tiny\color{ink}scannen zum Lösen}
\end{minipage}
\vspace{2mm}

{\raggedright\noindent\textcolor{ink}{\sffamily\textsc{Themen}}\enspace \themetag{reines Matt} \themetag{Mustermatt} \themetag{Idealmatt} \themetag{Selbstblock}\par}

\clearpage
\phantomsection\label{hm:KBvkrrn}
\begin{tcolorbox}[enhanced,colback=ink,colframe=ink,arc=0pt,boxrule=0pt,left=4mm,right=4mm,top=2.5mm,bottom=2.5mm]
\color{white}\sffamily{\LARGE\bfseries KBvkrrn}\hfill{\Large h\#9.5}\\[0.5mm]
{\small 6 Steine\quad Weiß am Zug}\hfill{\small Nr.~285\qquad\href{https://helpman.komtera.lt/?fen=1n5B/8/8/8/3rr3/8/8/K6k&moves=9.5}{\color{white}\bfseries lösen\,\ensuremath{\nearrow}}}
\end{tcolorbox}
\vspace{2mm}
\noindent\begin{minipage}[t]{36mm}\vspace{0pt}\sffamily\small\raggedright
\hmfigure{längste eindeutige}{{\Large h\#9.5}}\hmfigure{eindeutig bei dieser Länge}{{\Large 128} Stellungen}\hmfigure{längstes Matt}{{\Large h\#12.5}}\hmfigure{in der Klasse}{15,5 Milliarden Stellungen,\\ 4,37 Milliarden mit Hilfsmatt}
\end{minipage}\hfill
\begin{minipage}[t]{80mm}\vspace{0pt}\begin{tikzpicture}\begin{axis}[width=72mm,height=42mm,ybar,bar width=1.6pt,xmin=-1,xmax=26,ymin=0,ymax=1064,axis y line*=left,axis x line*=bottom,xlabel={\scriptsize Halbzüge bis Matt},ylabel={\scriptsize Stellungen (Mio.)},tick label style={font=\tiny},label style={font=\scriptsize},enlarge x limits=false,ylabel shift=-3pt,xlabel shift=-3pt]\addplot[fill=ink!55,draw=none] coordinates {(0,0.000825) (1,0.001937) (2,0.006409) (3,0.05351) (4,0.2352) (5,1.357) (6,4.145) (7,14.93) (8,36.27) (9,93.54) (10,170.4) (11,359.1) (12,479.9) (13,854.3) (14,717.2) (15,967.6) (16,327.3) (17,299.1) (18,21.98) (19,24.8) (20,0.1937) (21,1.09) (22,1.1e-05) (23,0.07524) (24,4e-06) (25,0.000358)};\draw[densely dotted] (axis cs:19,0) -- (axis cs:19,1064);\end{axis}\begin{axis}[width=72mm,height=42mm,axis y line*=right,axis x line=none,ymode=log,xmin=-1,xmax=26,ylabel={\scriptsize eindeutig (log)},tick label style={font=\tiny},label style={font=\scriptsize},ylabel shift=-3pt]\addplot[mark=*,mark size=0.7pt,draw=black,thin] coordinates {(0,825) (1,1937) (2,4285) (3,29191) (4,135039) (5,431414) (6,501229) (7,922120) (8,982118) (9,1624767) (10,1350111) (11,1740282) (12,1243778) (13,1509326) (14,503648) (15,433951) (16,25635) (17,9855) (18,63) (19,128)};\end{axis}\end{tikzpicture}\end{minipage}
\vspace{3mm}

\noindent\begin{minipage}[t]{60mm}\vspace{0pt}\sbox\hmboard{\chessboard[setfen=1n5B/8/8/8/3rr3/8/8/K6k w - - 0 1]}\usebox\hmboard\\[1mm]
\makebox[\wd\hmboard]{\sffamily\hspace{3mm}\textbf{h\#9.5}\hfill(2+4)}\end{minipage}\hfill
\begin{minipage}[t]{52mm}\vspace{0pt}
\textcolor{ink}{\sffamily\textsc{Lösung}}\par\smallskip\small
\begin{tabular}{@{}r@{\enspace}l@{\enspace}l@{}}
\textbf{1.} & \ldots & \symking{}b2\\
\textbf{2.} & \symrook{}d1 & \symbishop{}c3\\
\textbf{3.} & \symknight{}c6 & \symbishop{}d2\\
\textbf{4.} & \symknight{}d4 & \symking{}c3\\
\textbf{5.} & \symknight{}e2+ & \symking{}d3\\
\textbf{6.} & \symrook{}d4+ & \symking{}e3\\
\textbf{7.} & \symking{}g1 & \symking{}f3\\
\textbf{8.} & \symking{}f1 & \symbishop{}e3\\
\textbf{9.} & \symking{}e1 & \symking{}g2\\
\textbf{10.} & \symrook{}4d2 & \symbishop{}f2\#\\
\end{tabular}\par\vspace{3mm}
\qrcode[height=12mm]{https://helpman.komtera.lt/?fen=1n5B/8/8/8/3rr3/8/8/K6k&moves=9.5}\\{\sffamily\tiny\color{ink}scannen zum Lösen}
\end{minipage}
\vspace{2mm}

{\raggedright\noindent\textcolor{ink}{\sffamily\textsc{Themen}}\enspace \themetag{reines Matt} \themetag{Mustermatt} \themetag{Idealmatt} \themetag{Selbstblock} \themetag{ohne Schlag}\par}

\clearpage
\phantomsection\label{hm:KNPvkqp}
\begin{tcolorbox}[enhanced,colback=ink,colframe=ink,arc=0pt,boxrule=0pt,left=4mm,right=4mm,top=2.5mm,bottom=2.5mm]
\color{white}\sffamily{\LARGE\bfseries KNPvkqp}\hfill{\Large h\#9.5}\\[0.5mm]
{\small 6 Steine\quad Weiß am Zug}\hfill{\small Nr.~286\qquad\href{https://helpman.komtera.lt/?fen=8/7k/7N/8/8/1p6/1P6/Kq6&moves=9.5}{\color{white}\bfseries lösen\,\ensuremath{\nearrow}}}
\end{tcolorbox}
\vspace{2mm}
\noindent\begin{minipage}[t]{36mm}\vspace{0pt}\sffamily\small\raggedright
\hmfigure{längste eindeutige}{{\Large h\#9.5}}\hmfigure{eindeutig bei dieser Länge}{{\Large 4} Stellungen}\hmfigure{längstes Matt}{{\Large h\#14}}\hmfigure{in der Klasse}{34,1 Milliarden Stellungen,\\ 21 Milliarden mit Hilfsmatt}
\end{minipage}\hfill
\begin{minipage}[t]{80mm}\vspace{0pt}\begin{tikzpicture}\begin{axis}[width=72mm,height=42mm,ybar,bar width=1.6pt,xmin=-1,xmax=29,ymin=0,ymax=3.619,axis y line*=left,axis x line*=bottom,xlabel={\scriptsize Halbzüge bis Matt},ylabel={\scriptsize Stellungen (Mrd.)},tick label style={font=\tiny},label style={font=\scriptsize},enlarge x limits=false,ylabel shift=-3pt,xlabel shift=-3pt]\addplot[fill=ink!55,draw=none] coordinates {(0,4.963e-05) (1,0.002304) (2,0.01014) (3,0.1073) (4,0.3265) (5,1.347) (6,1.732) (7,2.851) (8,2.215) (9,3.141) (10,2.397) (11,3.29) (12,1.999) (13,1.248) (14,0.2404) (15,0.07927) (16,0.006027) (17,0.003956) (18,5.755e-05) (19,0.0007923) (20,1.08e-05) (21,7.943e-05) (22,2.28e-07) (23,2.624e-06) (24,5.2e-08) (25,1.3e-08) (26,4.3e-08) (27,4e-09) (28,1.8e-08)};\draw[densely dotted] (axis cs:19,0) -- (axis cs:19,3.619);\end{axis}\begin{axis}[width=72mm,height=42mm,axis y line*=right,axis x line=none,ymode=log,xmin=-1,xmax=29,ylabel={\scriptsize eindeutig (log)},tick label style={font=\tiny},label style={font=\scriptsize},ylabel shift=-3pt]\addplot[mark=*,mark size=0.7pt,draw=black,thin] coordinates {(0,49626) (1,854910) (2,3269066) (3,39228191) (4,68950420) (5,111731876) (6,49007998) (7,56174533) (8,24710368) (9,30052610) (10,12420883) (11,12199655) (12,2754231) (13,788707) (14,35208) (15,13621) (16,17) (17,847) (18,4) (19,4)};\end{axis}\end{tikzpicture}\end{minipage}
\vspace{3mm}

\noindent\begin{minipage}[t]{60mm}\vspace{0pt}\sbox\hmboard{\chessboard[setfen=8/7k/7N/8/8/1p6/1P6/Kq6 w - - 0 1]}\usebox\hmboard\\[1mm]
\makebox[\wd\hmboard]{\sffamily\hspace{3mm}\textbf{h\#9.5}\hfill(3+3)}\end{minipage}\hfill
\begin{minipage}[t]{52mm}\vspace{0pt}
\textcolor{ink}{\sffamily\textsc{Lösung}}\par\smallskip\small
\begin{tabular}{@{}r@{\enspace}l@{\enspace}l@{}}
\textbf{1.} & \ldots & \symking{}xb1\\
\textbf{2.} & \symking{}g6 & \symknight{}g4\\
\textbf{3.} & \symking{}f5 & \symking{}c1\\
\textbf{4.} & \symking{}e4 & \symking{}d1\\
\textbf{5.} & \symking{}d3 & \symking{}e1\\
\textbf{6.} & \symking{}c2 & \symknight{}e3+\\
\textbf{7.} & \symking{}xb2 & \symking{}d2\\
\textbf{8.} & \symking{}a1 & \symking{}c3\\
\textbf{9.} & b2 & \symking{}b3\\
\textbf{10.} & b1=\symknight{} & \symknight{}c2\#\\
\end{tabular}\par\vspace{3mm}
\qrcode[height=12mm]{https://helpman.komtera.lt/?fen=8/7k/7N/8/8/1p6/1P6/Kq6&moves=9.5}\\{\sffamily\tiny\color{ink}scannen zum Lösen}
\end{minipage}
\vspace{2mm}

{\raggedright\noindent\textcolor{ink}{\sffamily\textsc{Themen}}\enspace \themetag{reines Matt} \themetag{Mustermatt} \themetag{Idealmatt} \themetag{Umwandlung} \themetag{Unterverwandlung} \themetag{Selbstblock} \themetag{Umnov}\par\smallskip{\footnotesize\textit{Hinweis:} die Partei am Zug steht im Diagramm im Schach; die Lösung beginnt mit einem Schlag.}\par}

\clearpage
\phantomsection\label{hm:KNPvkrp}
\begin{tcolorbox}[enhanced,colback=ink,colframe=ink,arc=0pt,boxrule=0pt,left=4mm,right=4mm,top=2.5mm,bottom=2.5mm]
\color{white}\sffamily{\LARGE\bfseries KNPvkrp}\hfill{\Large h\#9.5}\\[0.5mm]
{\small 6 Steine\quad Weiß am Zug}\hfill{\small Nr.~287\qquad\href{https://helpman.komtera.lt/?fen=8/7k/7N/8/8/1p6/1P6/Kr6&moves=9.5}{\color{white}\bfseries lösen\,\ensuremath{\nearrow}}}
\end{tcolorbox}
\vspace{2mm}
\noindent\begin{minipage}[t]{36mm}\vspace{0pt}\sffamily\small\raggedright
\hmfigure{längste eindeutige}{{\Large h\#9.5}}\hmfigure{eindeutig bei dieser Länge}{{\Large 2} Stellungen}\hmfigure{längstes Matt}{{\Large h\#12.5}}\hmfigure{in der Klasse}{34,1 Milliarden Stellungen,\\ 22,7 Milliarden mit Hilfsmatt}
\end{minipage}\hfill
\begin{minipage}[t]{80mm}\vspace{0pt}\begin{tikzpicture}\begin{axis}[width=72mm,height=42mm,ybar,bar width=1.6pt,xmin=-1,xmax=26,ymin=0,ymax=3.878,axis y line*=left,axis x line*=bottom,xlabel={\scriptsize Halbzüge bis Matt},ylabel={\scriptsize Stellungen (Mrd.)},tick label style={font=\tiny},label style={font=\scriptsize},enlarge x limits=false,ylabel shift=-3pt,xlabel shift=-3pt]\addplot[fill=ink!55,draw=none] coordinates {(0,7.857e-05) (1,0.00394) (2,0.01586) (3,0.1466) (4,0.491) (5,1.718) (6,2.384) (7,3.265) (8,2.961) (9,3.525) (10,3.061) (11,2.852) (12,1.558) (13,0.5507) (14,0.1166) (15,0.02108) (16,0.002085) (17,0.0003648) (18,5.561e-06) (19,6.667e-05) (20,8e-09) (21,8.628e-06) (22,2.9e-08) (23,2.77e-07) (24,1e-09) (25,2e-09)};\draw[densely dotted] (axis cs:19,0) -- (axis cs:19,3.878);\end{axis}\begin{axis}[width=72mm,height=42mm,axis y line*=right,axis x line=none,ymode=log,xmin=-1,xmax=26,ylabel={\scriptsize eindeutig (log)},tick label style={font=\tiny},label style={font=\scriptsize},ylabel shift=-3pt]\addplot[mark=*,mark size=0.7pt,draw=black,thin] coordinates {(0,78566) (1,1419919) (2,5795834) (3,55039110) (4,113003828) (5,147498193) (6,76784824) (7,66463574) (8,31789847) (9,28740056) (10,12765927) (11,6157530) (12,1168348) (13,196583) (14,26719) (15,2736) (16,25) (17,115) (19,2)};\end{axis}\end{tikzpicture}\end{minipage}
\vspace{3mm}

\noindent\begin{minipage}[t]{60mm}\vspace{0pt}\sbox\hmboard{\chessboard[setfen=8/7k/7N/8/8/1p6/1P6/Kr6 w - - 0 1]}\usebox\hmboard\\[1mm]
\makebox[\wd\hmboard]{\sffamily\hspace{3mm}\textbf{h\#9.5}\hfill(3+3)}\end{minipage}\hfill
\begin{minipage}[t]{52mm}\vspace{0pt}
\textcolor{ink}{\sffamily\textsc{Lösung}}\par\smallskip\small
\begin{tabular}{@{}r@{\enspace}l@{\enspace}l@{}}
\textbf{1.} & \ldots & \symking{}xb1\\
\textbf{2.} & \symking{}g6 & \symknight{}g4\\
\textbf{3.} & \symking{}f5 & \symking{}c1\\
\textbf{4.} & \symking{}e4 & \symking{}d1\\
\textbf{5.} & \symking{}d3 & \symking{}e1\\
\textbf{6.} & \symking{}c2 & \symknight{}e3+\\
\textbf{7.} & \symking{}xb2 & \symking{}d2\\
\textbf{8.} & \symking{}a1 & \symking{}c3\\
\textbf{9.} & b2 & \symking{}b3\\
\textbf{10.} & b1=\symknight{} & \symknight{}c2\#\\
\end{tabular}\par\vspace{3mm}
\qrcode[height=12mm]{https://helpman.komtera.lt/?fen=8/7k/7N/8/8/1p6/1P6/Kr6&moves=9.5}\\{\sffamily\tiny\color{ink}scannen zum Lösen}
\end{minipage}
\vspace{2mm}

{\raggedright\noindent\textcolor{ink}{\sffamily\textsc{Themen}}\enspace \themetag{reines Matt} \themetag{Mustermatt} \themetag{Idealmatt} \themetag{Umwandlung} \themetag{Unterverwandlung} \themetag{Selbstblock} \themetag{Umnov}\par\smallskip{\footnotesize\textit{Hinweis:} die Partei am Zug steht im Diagramm im Schach; die Lösung beginnt mit einem Schlag.}\par}

\clearpage
\phantomsection\label{hm:KNvkppp}
\begin{tcolorbox}[enhanced,colback=ink,colframe=ink,arc=0pt,boxrule=0pt,left=4mm,right=4mm,top=2.5mm,bottom=2.5mm]
\color{white}\sffamily{\LARGE\bfseries KNvkppp}\hfill{\Large h\#9.5}\\[0.5mm]
{\small 6 Steine\quad Weiß am Zug}\hfill{\small Nr.~288\qquad\href{https://helpman.komtera.lt/?fen=8/6p1/1Kp1k1p1/8/8/8/8/2N5&moves=9.5}{\color{white}\bfseries lösen\,\ensuremath{\nearrow}}}
\end{tcolorbox}
\vspace{2mm}
\noindent\begin{minipage}[t]{36mm}\vspace{0pt}\sffamily\small\raggedright
\hmfigure{längste eindeutige}{{\Large h\#9.5}}\hmfigure{eindeutig bei dieser Länge}{{\Large 17} Stellungen}\hmfigure{längstes Matt}{{\Large h\#11.5}}\hmfigure{in der Klasse}{25,6 Milliarden Stellungen,\\ 3,08 Milliarden mit Hilfsmatt}
\end{minipage}\hfill
\begin{minipage}[t]{80mm}\vspace{0pt}\begin{tikzpicture}\begin{axis}[width=72mm,height=42mm,ybar,bar width=1.6pt,xmin=-1,xmax=24,ymin=0,ymax=556,axis y line*=left,axis x line*=bottom,xlabel={\scriptsize Halbzüge bis Matt},ylabel={\scriptsize Stellungen (Mio.)},tick label style={font=\tiny},label style={font=\scriptsize},enlarge x limits=false,ylabel shift=-3pt,xlabel shift=-3pt]\addplot[fill=ink!55,draw=none] coordinates {(0,0.004468) (1,0.02088) (2,0.05065) (3,0.3092) (4,1.165) (5,4.445) (6,15.17) (7,38.94) (8,93.82) (9,166.8) (10,290.2) (11,394.5) (12,497.1) (13,505.4) (14,419.6) (15,308.3) (16,180) (17,99.33) (18,37.51) (19,19.05) (20,2.953) (21,1.46) (22,0.01531) (23,4e-06)};\draw[densely dotted] (axis cs:19,0) -- (axis cs:19,556);\end{axis}\begin{axis}[width=72mm,height=42mm,axis y line*=right,axis x line=none,ymode=log,xmin=-1,xmax=24,ylabel={\scriptsize eindeutig (log)},tick label style={font=\tiny},label style={font=\scriptsize},ylabel shift=-3pt]\addplot[mark=*,mark size=0.7pt,draw=black,thin] coordinates {(0,4468) (1,19654) (2,30869) (3,130332) (4,464747) (5,618775) (6,1211323) (7,1161229) (8,1511371) (9,1363152) (10,1564231) (11,1013364) (12,686139) (13,147183) (14,62235) (15,7697) (16,2010) (17,210) (18,92) (19,17)};\end{axis}\end{tikzpicture}\end{minipage}
\vspace{3mm}

\noindent\begin{minipage}[t]{60mm}\vspace{0pt}\sbox\hmboard{\chessboard[setfen=8/6p1/1Kp1k1p1/8/8/8/8/2N5 w - - 0 1]}\usebox\hmboard\\[1mm]
\makebox[\wd\hmboard]{\sffamily\hspace{3mm}\textbf{h\#9.5}\hfill(2+4)}\end{minipage}\hfill
\begin{minipage}[t]{52mm}\vspace{0pt}
\textcolor{ink}{\sffamily\textsc{Lösung}}\par\smallskip\small
\begin{tabular}{@{}r@{\enspace}l@{\enspace}l@{}}
\textbf{1.} & \ldots & \symking{}c7\\
\textbf{2.} & g5 & \symking{}d8\\
\textbf{3.} & g4 & \symking{}e8\\
\textbf{4.} & g3 & \symking{}f8\\
\textbf{5.} & g2 & \symking{}xg7\\
\textbf{6.} & g1=\symrook{}+ & \symking{}h7\\
\textbf{7.} & \symking{}f7 & \symking{}h6\\
\textbf{8.} & \symking{}g8 & \symknight{}d3\\
\textbf{9.} & \symking{}h8 & \symknight{}e5\\
\textbf{10.} & \symrook{}g8 & \symknight{}f7\#\\
\end{tabular}\par\vspace{3mm}
\qrcode[height=12mm]{https://helpman.komtera.lt/?fen=8/6p1/1Kp1k1p1/8/8/8/8/2N5&moves=9.5}\\{\sffamily\tiny\color{ink}scannen zum Lösen}
\end{minipage}
\vspace{2mm}

{\raggedright\noindent\textcolor{ink}{\sffamily\textsc{Themen}}\enspace \themetag{reines Matt} \themetag{Mustermatt} \themetag{Umwandlung} \themetag{Unterverwandlung} \themetag{Selbstblock}\par}

\clearpage
\phantomsection\label{hm:KNvkqbb}
\begin{tcolorbox}[enhanced,colback=ink,colframe=ink,arc=0pt,boxrule=0pt,left=4mm,right=4mm,top=2.5mm,bottom=2.5mm]
\color{white}\sffamily{\LARGE\bfseries KNvkqbb}\hfill{\Large h\#9.5}\\[0.5mm]
{\small 6 Steine\quad Weiß am Zug}\hfill{\small Nr.~289\qquad\href{https://helpman.komtera.lt/?fen=7k/8/8/2b5/8/q7/2b5/K1N5&moves=9.5}{\color{white}\bfseries lösen\,\ensuremath{\nearrow}}}
\end{tcolorbox}
\vspace{2mm}
\noindent\begin{minipage}[t]{36mm}\vspace{0pt}\sffamily\small\raggedright
\hmfigure{längste eindeutige}{{\Large h\#9.5}}\hmfigure{eindeutig bei dieser Länge}{{\Large 23} Stellungen}\hmfigure{längstes Matt}{{\Large h\#11.5}}\hmfigure{in der Klasse}{15,5 Milliarden Stellungen,\\ 4,26 Milliarden mit Hilfsmatt}
\end{minipage}\hfill
\begin{minipage}[t]{80mm}\vspace{0pt}\begin{tikzpicture}\begin{axis}[width=72mm,height=42mm,ybar,bar width=1.6pt,xmin=-1,xmax=24,ymin=0,ymax=1.169,axis y line*=left,axis x line*=bottom,xlabel={\scriptsize Halbzüge bis Matt},ylabel={\scriptsize Stellungen (Mrd.)},tick label style={font=\tiny},label style={font=\scriptsize},enlarge x limits=false,ylabel shift=-3pt,xlabel shift=-3pt]\addplot[fill=ink!55,draw=none] coordinates {(0,1.327e-06) (1,6.813e-06) (2,5.3e-05) (3,0.000336) (4,0.001371) (5,0.006682) (6,0.01383) (7,0.05397) (8,0.08494) (9,0.2626) (10,0.333) (11,0.8014) (12,0.5889) (13,1.063) (14,0.3922) (15,0.4815) (16,0.09679) (17,0.07292) (18,0.007247) (19,0.003562) (20,3.164e-05) (21,1.282e-05) (23,4e-09)};\draw[densely dotted] (axis cs:19,0) -- (axis cs:19,1.169);\end{axis}\begin{axis}[width=72mm,height=42mm,axis y line*=right,axis x line=none,ymode=log,xmin=-1,xmax=24,ylabel={\scriptsize eindeutig (log)},tick label style={font=\tiny},label style={font=\scriptsize},ylabel shift=-3pt]\addplot[mark=*,mark size=0.7pt,draw=black,thin] coordinates {(0,1327) (1,6813) (2,42501) (3,205177) (4,511071) (5,1282514) (6,698374) (7,1798589) (8,1520037) (9,3109665) (10,2206435) (11,3429288) (12,801911) (13,854483) (14,120889) (15,87888) (16,12626) (17,6646) (18,34) (19,23)};\end{axis}\end{tikzpicture}\end{minipage}
\vspace{3mm}

\noindent\begin{minipage}[t]{60mm}\vspace{0pt}\sbox\hmboard{\chessboard[setfen=7k/8/8/2b5/8/q7/2b5/K1N5 w - - 0 1]}\usebox\hmboard\\[1mm]
\makebox[\wd\hmboard]{\sffamily\hspace{3mm}\textbf{h\#9.5}\hfill(2+4)}\end{minipage}\hfill
\begin{minipage}[t]{52mm}\vspace{0pt}
\textcolor{ink}{\sffamily\textsc{Lösung}}\par\smallskip\small
\begin{tabular}{@{}r@{\enspace}l@{\enspace}l@{}}
\textbf{1.} & \ldots & \symknight{}a2\\
\textbf{2.} & \symqueen{}b2+ & \symking{}xb2\\
\textbf{3.} & \symking{}g7 & \symking{}c3\\
\textbf{4.} & \symking{}f6 & \symking{}c4\\
\textbf{5.} & \symking{}e5 & \symking{}b5\\
\textbf{6.} & \symking{}d4 & \symknight{}c1\\
\textbf{7.} & \symking{}c3 & \symking{}xc5\\
\textbf{8.} & \symking{}b2 & \symking{}b4\\
\textbf{9.} & \symking{}a1 & \symking{}a3\\
\textbf{10.} & \symbishop{}b1 & \symknight{}b3\#\\
\end{tabular}\par\vspace{3mm}
\qrcode[height=12mm]{https://helpman.komtera.lt/?fen=7k/8/8/2b5/8/q7/2b5/K1N5&moves=9.5}\\{\sffamily\tiny\color{ink}scannen zum Lösen}
\end{minipage}
\vspace{2mm}

{\raggedright\noindent\textcolor{ink}{\sffamily\textsc{Themen}}\enspace \themetag{reines Matt} \themetag{Mustermatt} \themetag{Idealmatt} \themetag{Switchback} \themetag{Selbstblock}\par\smallskip{\footnotesize\textit{Hinweis:} die Partei am Zug steht im Diagramm im Schach.}\par}

\clearpage
\phantomsection\label{hm:KNvkqbn}
\begin{tcolorbox}[enhanced,colback=ink,colframe=ink,arc=0pt,boxrule=0pt,left=4mm,right=4mm,top=2.5mm,bottom=2.5mm]
\color{white}\sffamily{\LARGE\bfseries KNvkqbn}\hfill{\Large h\#9.5}\\[0.5mm]
{\small 6 Steine\quad Weiß am Zug}\hfill{\small Nr.~290\qquad\href{https://helpman.komtera.lt/?fen=7k/6N1/8/8/8/8/b2q4/nK6&moves=9.5}{\color{white}\bfseries lösen\,\ensuremath{\nearrow}}}
\end{tcolorbox}
\vspace{2mm}
\noindent\begin{minipage}[t]{36mm}\vspace{0pt}\sffamily\small\raggedright
\hmfigure{längste eindeutige}{{\Large h\#9.5}}\hmfigure{eindeutig bei dieser Länge}{{\Large 3} Stellungen}\hmfigure{längstes Matt}{{\Large h\#11.5}}\hmfigure{in der Klasse}{15,5 Milliarden Stellungen,\\ 8,77 Milliarden mit Hilfsmatt}
\end{minipage}\hfill
\begin{minipage}[t]{80mm}\vspace{0pt}\begin{tikzpicture}\begin{axis}[width=72mm,height=42mm,ybar,bar width=1.6pt,xmin=-1,xmax=24,ymin=0,ymax=2.461,axis y line*=left,axis x line*=bottom,xlabel={\scriptsize Halbzüge bis Matt},ylabel={\scriptsize Stellungen (Mrd.)},tick label style={font=\tiny},label style={font=\scriptsize},enlarge x limits=false,ylabel shift=-3pt,xlabel shift=-3pt]\addplot[fill=ink!55,draw=none] coordinates {(0,3.262e-06) (1,1.735e-05) (2,0.0001187) (3,0.0008348) (4,0.004318) (5,0.02189) (6,0.06811) (7,0.2491) (8,0.466) (9,1.201) (10,1.182) (11,2.237) (12,1.173) (13,1.564) (14,0.3475) (15,0.2304) (16,0.01771) (17,0.006037) (18,6.344e-05) (19,0.0004616) (20,4.3e-08) (21,9.84e-06) (22,2e-09) (23,9e-09)};\draw[densely dotted] (axis cs:19,0) -- (axis cs:19,2.461);\end{axis}\begin{axis}[width=72mm,height=42mm,axis y line*=right,axis x line=none,ymode=log,xmin=-1,xmax=24,ylabel={\scriptsize eindeutig (log)},tick label style={font=\tiny},label style={font=\scriptsize},ylabel shift=-3pt]\addplot[mark=*,mark size=0.7pt,draw=black,thin] coordinates {(0,3262) (1,17342) (2,93673) (3,499834) (4,1767353) (5,4528643) (6,4917858) (7,9907380) (8,6528374) (9,10730939) (10,4127540) (11,5433990) (12,782076) (13,628717) (14,26030) (15,10364) (16,187) (17,580) (19,3)};\end{axis}\end{tikzpicture}\end{minipage}
\vspace{3mm}

\noindent\begin{minipage}[t]{60mm}\vspace{0pt}\sbox\hmboard{\chessboard[setfen=7k/6N1/8/8/8/8/b2q4/nK6 w - - 0 1]}\usebox\hmboard\\[1mm]
\makebox[\wd\hmboard]{\sffamily\hspace{3mm}\textbf{h\#9.5}\hfill(2+4)}\end{minipage}\hfill
\begin{minipage}[t]{52mm}\vspace{0pt}
\textcolor{ink}{\sffamily\textsc{Lösung}}\par\smallskip\small
\begin{tabular}{@{}r@{\enspace}l@{\enspace}l@{}}
\textbf{1.} & \ldots & \symking{}xa1\\
\textbf{2.} & \symking{}g8 & \symknight{}f5\\
\textbf{3.} & \symking{}f7 & \symknight{}e3\\
\textbf{4.} & \symking{}e6 & \symknight{}c2\\
\textbf{5.} & \symking{}d5 & \symking{}b2\\
\textbf{6.} & \symking{}c4 & \symking{}a3\\
\textbf{7.} & \symking{}c3 & \symking{}a4\\
\textbf{8.} & \symking{}b2 & \symknight{}d4\\
\textbf{9.} & \symking{}a1 & \symking{}a3\\
\textbf{10.} & \symbishop{}b1 & \symknight{}b3\#\\
\end{tabular}\par\vspace{3mm}
\qrcode[height=12mm]{https://helpman.komtera.lt/?fen=7k/6N1/8/8/8/8/b2q4/nK6&moves=9.5}\\{\sffamily\tiny\color{ink}scannen zum Lösen}
\end{minipage}
\vspace{2mm}

{\raggedright\noindent\textcolor{ink}{\sffamily\textsc{Themen}}\enspace \themetag{reines Matt} \themetag{Mustermatt} \themetag{Switchback} \themetag{Selbstblock} \themetag{Kniest} \themetag{ohne Schach}\par\smallskip{\footnotesize\textit{Hinweis:} die Partei am Zug steht im Diagramm im Schach; die Lösung beginnt mit einem Schlag.}\par}

\clearpage
\phantomsection\label{hm:KNvkqbp}
\begin{tcolorbox}[enhanced,colback=ink,colframe=ink,arc=0pt,boxrule=0pt,left=4mm,right=4mm,top=2.5mm,bottom=2.5mm]
\color{white}\sffamily{\LARGE\bfseries KNvkqbp}\hfill{\Large h\#9.5}\\[0.5mm]
{\small 6 Steine\quad Weiß am Zug}\hfill{\small Nr.~291\qquad\href{https://helpman.komtera.lt/?fen=8/8/6p1/8/8/q1N5/8/Kbk5&moves=9.5}{\color{white}\bfseries lösen\,\ensuremath{\nearrow}}}
\end{tcolorbox}
\vspace{2mm}
\noindent\begin{minipage}[t]{36mm}\vspace{0pt}\sffamily\small\raggedright
\hmfigure{längste eindeutige}{{\Large h\#9.5}}\hmfigure{eindeutig bei dieser Länge}{{\Large 50} Stellungen}\hmfigure{längstes Matt}{{\Large h\#10.5}}\hmfigure{in der Klasse}{45,4 Milliarden Stellungen,\\ 26,9 Milliarden mit Hilfsmatt}
\end{minipage}\hfill
\begin{minipage}[t]{80mm}\vspace{0pt}\begin{tikzpicture}\begin{axis}[width=72mm,height=42mm,ybar,bar width=1.6pt,xmin=-1,xmax=22,ymin=0,ymax=6.188,axis y line*=left,axis x line*=bottom,xlabel={\scriptsize Halbzüge bis Matt},ylabel={\scriptsize Stellungen (Mrd.)},tick label style={font=\tiny},label style={font=\scriptsize},enlarge x limits=false,ylabel shift=-3pt,xlabel shift=-3pt]\addplot[fill=ink!55,draw=none] coordinates {(0,1.263e-05) (1,6.366e-05) (2,0.0003798) (3,0.002429) (4,0.01086) (5,0.05079) (6,0.1237) (7,0.4381) (8,0.7203) (9,1.936) (10,2.245) (11,4.655) (12,3.552) (13,5.626) (14,2.729) (15,3.169) (16,0.8981) (17,0.6757) (18,0.05736) (19,0.01995) (20,0.00012) (21,0.0001109)};\draw[densely dotted] (axis cs:19,0) -- (axis cs:19,6.188);\end{axis}\begin{axis}[width=72mm,height=42mm,axis y line*=right,axis x line=none,ymode=log,xmin=-1,xmax=22,ylabel={\scriptsize eindeutig (log)},tick label style={font=\tiny},label style={font=\scriptsize},ylabel shift=-3pt]\addplot[mark=*,mark size=0.7pt,draw=black,thin] coordinates {(0,12629) (1,62883) (2,284058) (3,1362064) (4,4223864) (5,9482208) (6,7514958) (7,16344357) (8,13026352) (9,22624916) (10,14838755) (11,19248909) (12,6641722) (13,5984581) (14,1588716) (15,1024307) (16,113167) (17,35935) (18,103) (19,50)};\end{axis}\end{tikzpicture}\end{minipage}
\vspace{3mm}

\noindent\begin{minipage}[t]{60mm}\vspace{0pt}\sbox\hmboard{\chessboard[setfen=8/8/6p1/8/8/q1N5/8/Kbk5 w - - 0 1]}\usebox\hmboard\\[1mm]
\makebox[\wd\hmboard]{\sffamily\hspace{3mm}\textbf{h\#9.5}\hfill(2+4)}\end{minipage}\hfill
\begin{minipage}[t]{52mm}\vspace{0pt}
\textcolor{ink}{\sffamily\textsc{Lösung}}\par\smallskip\small
\begin{tabular}{@{}r@{\enspace}l@{\enspace}l@{}}
\textbf{1.} & \ldots & \symknight{}a2+\\
\textbf{2.} & \symking{}d2 & \symking{}xb1\\
\textbf{3.} & g5 & \symknight{}c3\\
\textbf{4.} & \symqueen{}c1+ & \symking{}a2\\
\textbf{5.} & g4 & \symking{}b3\\
\textbf{6.} & g3 & \symking{}c4\\
\textbf{7.} & g2 & \symking{}d4\\
\textbf{8.} & g1=\symrook{} & \symknight{}e4+\\
\textbf{9.} & \symking{}d1 & \symking{}d3\\
\textbf{10.} & \symrook{}e1 & \symknight{}f2\#\\
\end{tabular}\par\vspace{3mm}
\qrcode[height=12mm]{https://helpman.komtera.lt/?fen=8/8/6p1/8/8/q1N5/8/Kbk5&moves=9.5}\\{\sffamily\tiny\color{ink}scannen zum Lösen}
\end{minipage}
\vspace{2mm}

{\raggedright\noindent\textcolor{ink}{\sffamily\textsc{Themen}}\enspace \themetag{reines Matt} \themetag{Mustermatt} \themetag{Idealmatt} \themetag{Umwandlung} \themetag{Unterverwandlung} \themetag{Switchback} \themetag{Selbstblock}\par\smallskip{\footnotesize\textit{Hinweis:} die Partei am Zug steht im Diagramm im Schach.}\par}

\clearpage
\phantomsection\label{hm:KNvkqnp}
\begin{tcolorbox}[enhanced,colback=ink,colframe=ink,arc=0pt,boxrule=0pt,left=4mm,right=4mm,top=2.5mm,bottom=2.5mm]
\color{white}\sffamily{\LARGE\bfseries KNvkqnp}\hfill{\Large h\#9.5}\\[0.5mm]
{\small 6 Steine\quad Weiß am Zug}\hfill{\small Nr.~292\qquad\href{https://helpman.komtera.lt/?fen=8/8/6p1/8/8/q1N5/8/Knk5&moves=9.5}{\color{white}\bfseries lösen\,\ensuremath{\nearrow}}}
\end{tcolorbox}
\vspace{2mm}
\noindent\begin{minipage}[t]{36mm}\vspace{0pt}\sffamily\small\raggedright
\hmfigure{längste eindeutige}{{\Large h\#9.5}}\hmfigure{eindeutig bei dieser Länge}{{\Large 50} Stellungen}\hmfigure{längstes Matt}{{\Large h\#11.5}}\hmfigure{in der Klasse}{45,4 Milliarden Stellungen,\\ 27,6 Milliarden mit Hilfsmatt}
\end{minipage}\hfill
\begin{minipage}[t]{80mm}\vspace{0pt}\begin{tikzpicture}\begin{axis}[width=72mm,height=42mm,ybar,bar width=1.6pt,xmin=-1,xmax=24,ymin=0,ymax=7.394,axis y line*=left,axis x line*=bottom,xlabel={\scriptsize Halbzüge bis Matt},ylabel={\scriptsize Stellungen (Mrd.)},tick label style={font=\tiny},label style={font=\scriptsize},enlarge x limits=false,ylabel shift=-3pt,xlabel shift=-3pt]\addplot[fill=ink!55,draw=none] coordinates {(0,1.398e-05) (1,7.229e-05) (2,0.000365) (3,0.002606) (4,0.01338) (5,0.06598) (6,0.2319) (7,0.7906) (8,1.631) (9,3.79) (10,4.048) (11,6.722) (12,3.844) (13,4.507) (14,1.151) (15,0.6956) (16,0.06709) (17,0.02823) (18,0.0002172) (19,0.001812) (20,3.1e-08) (21,0.0001016) (23,2.9e-08)};\draw[densely dotted] (axis cs:19,0) -- (axis cs:19,7.394);\end{axis}\begin{axis}[width=72mm,height=42mm,axis y line*=right,axis x line=none,ymode=log,xmin=-1,xmax=24,ylabel={\scriptsize eindeutig (log)},tick label style={font=\tiny},label style={font=\scriptsize},ylabel shift=-3pt]\addplot[mark=*,mark size=0.7pt,draw=black,thin] coordinates {(0,13983) (1,71355) (2,264699) (3,1431695) (4,4897451) (5,12631570) (6,18019229) (7,31538428) (8,21539345) (9,29471174) (10,10851546) (11,13383156) (12,2410900) (13,1892534) (14,119282) (15,27409) (16,201) (17,894) (18,4) (19,50)};\end{axis}\end{tikzpicture}\end{minipage}
\vspace{3mm}

\noindent\begin{minipage}[t]{60mm}\vspace{0pt}\sbox\hmboard{\chessboard[setfen=8/8/6p1/8/8/q1N5/8/Knk5 w - - 0 1]}\usebox\hmboard\\[1mm]
\makebox[\wd\hmboard]{\sffamily\hspace{3mm}\textbf{h\#9.5}\hfill(2+4)}\end{minipage}\hfill
\begin{minipage}[t]{52mm}\vspace{0pt}
\textcolor{ink}{\sffamily\textsc{Lösung}}\par\smallskip\small
\begin{tabular}{@{}r@{\enspace}l@{\enspace}l@{}}
\textbf{1.} & \ldots & \symknight{}a2+\\
\textbf{2.} & \symking{}d2 & \symking{}xb1\\
\textbf{3.} & g5 & \symknight{}c3\\
\textbf{4.} & \symqueen{}c1+ & \symking{}a2\\
\textbf{5.} & g4 & \symking{}b3\\
\textbf{6.} & g3 & \symking{}c4\\
\textbf{7.} & g2 & \symking{}d4\\
\textbf{8.} & g1=\symrook{} & \symknight{}e4+\\
\textbf{9.} & \symking{}d1 & \symking{}d3\\
\textbf{10.} & \symrook{}e1 & \symknight{}f2\#\\
\end{tabular}\par\vspace{3mm}
\qrcode[height=12mm]{https://helpman.komtera.lt/?fen=8/8/6p1/8/8/q1N5/8/Knk5&moves=9.5}\\{\sffamily\tiny\color{ink}scannen zum Lösen}
\end{minipage}
\vspace{2mm}

{\raggedright\noindent\textcolor{ink}{\sffamily\textsc{Themen}}\enspace \themetag{reines Matt} \themetag{Mustermatt} \themetag{Idealmatt} \themetag{Umwandlung} \themetag{Unterverwandlung} \themetag{Switchback} \themetag{Selbstblock}\par\smallskip{\footnotesize\textit{Hinweis:} die Partei am Zug steht im Diagramm im Schach.}\par}

\clearpage
\phantomsection\label{hm:KNvkqpp}
\begin{tcolorbox}[enhanced,colback=ink,colframe=ink,arc=0pt,boxrule=0pt,left=4mm,right=4mm,top=2.5mm,bottom=2.5mm]
\color{white}\sffamily{\LARGE\bfseries KNvkqpp}\hfill{\Large h\#9.5}\\[0.5mm]
{\small 6 Steine\quad Weiß am Zug}\hfill{\small Nr.~293\qquad\href{https://helpman.komtera.lt/?fen=8/3p4/4p3/4k3/8/8/2q5/K3N3&moves=9.5}{\color{white}\bfseries lösen\,\ensuremath{\nearrow}}}
\end{tcolorbox}
\vspace{2mm}
\noindent\begin{minipage}[t]{36mm}\vspace{0pt}\sffamily\small\raggedright
\hmfigure{längste eindeutige}{{\Large h\#9.5}}\hmfigure{eindeutig bei dieser Länge}{{\Large 993} Stellungen}\hmfigure{längstes Matt}{{\Large h\#11.5}}\hmfigure{in der Klasse}{34,1 Milliarden Stellungen,\\ 10,5 Milliarden mit Hilfsmatt}
\end{minipage}\hfill
\begin{minipage}[t]{80mm}\vspace{0pt}\begin{tikzpicture}\begin{axis}[width=72mm,height=42mm,ybar,bar width=1.6pt,xmin=-1,xmax=24,ymin=0,ymax=2.197,axis y line*=left,axis x line*=bottom,xlabel={\scriptsize Halbzüge bis Matt},ylabel={\scriptsize Stellungen (Mrd.)},tick label style={font=\tiny},label style={font=\scriptsize},enlarge x limits=false,ylabel shift=-3pt,xlabel shift=-3pt]\addplot[fill=ink!55,draw=none] coordinates {(0,6.626e-06) (1,3.256e-05) (2,0.0001369) (3,0.0008763) (4,0.003974) (5,0.01748) (6,0.05249) (7,0.1673) (8,0.3345) (9,0.7933) (10,1.042) (11,1.856) (12,1.499) (13,1.998) (14,1.003) (15,0.9986) (16,0.341) (17,0.3104) (18,0.05639) (19,0.05827) (20,0.002831) (21,0.003431) (22,7.06e-07) (23,1.77e-06)};\draw[densely dotted] (axis cs:19,0) -- (axis cs:19,2.197);\end{axis}\begin{axis}[width=72mm,height=42mm,axis y line*=right,axis x line=none,ymode=log,xmin=-1,xmax=24,ylabel={\scriptsize eindeutig (log)},tick label style={font=\tiny},label style={font=\scriptsize},ylabel shift=-3pt]\addplot[mark=*,mark size=0.7pt,draw=black,thin] coordinates {(0,6626) (1,31633) (2,94711) (3,438201) (4,1269508) (5,2747481) (6,3698617) (7,6896930) (8,6319144) (9,9366222) (10,5598886) (11,6729874) (12,2568351) (13,2249793) (14,557780) (15,365700) (16,54394) (17,35768) (18,1825) (19,993)};\end{axis}\end{tikzpicture}\end{minipage}
\vspace{3mm}

\noindent\begin{minipage}[t]{60mm}\vspace{0pt}\sbox\hmboard{\chessboard[setfen=8/3p4/4p3/4k3/8/8/2q5/K3N3 w - - 0 1]}\usebox\hmboard\\[1mm]
\makebox[\wd\hmboard]{\sffamily\hspace{3mm}\textbf{h\#9.5}\hfill(2+4)}\end{minipage}\hfill
\begin{minipage}[t]{52mm}\vspace{0pt}
\textcolor{ink}{\sffamily\textsc{Lösung}}\par\smallskip\small
\begin{tabular}{@{}r@{\enspace}l@{\enspace}l@{}}
\textbf{1.} & \ldots & \symknight{}f3+\\
\textbf{2.} & \symking{}f4 & \symknight{}d2\\
\textbf{3.} & \symqueen{}d1+ & \symking{}b2\\
\textbf{4.} & e5 & \symking{}c3\\
\textbf{5.} & e4 & \symking{}d4\\
\textbf{6.} & \symking{}g3 & \symking{}e5\\
\textbf{7.} & \symking{}f2 & \symking{}f4\\
\textbf{8.} & \symking{}e1 & \symking{}g3\\
\textbf{9.} & e3 & \symking{}g2\\
\textbf{10.} & e2 & \symknight{}f3\#\\
\end{tabular}\par\vspace{3mm}
\qrcode[height=12mm]{https://helpman.komtera.lt/?fen=8/3p4/4p3/4k3/8/8/2q5/K3N3&moves=9.5}\\{\sffamily\tiny\color{ink}scannen zum Lösen}
\end{minipage}
\vspace{2mm}

{\raggedright\noindent\textcolor{ink}{\sffamily\textsc{Themen}}\enspace \themetag{reines Matt} \themetag{Mustermatt} \themetag{Switchback} \themetag{Selbstblock} \themetag{ohne Schlag}\par}

\clearpage
\phantomsection\label{hm:KNvkqqb}
\begin{tcolorbox}[enhanced,colback=ink,colframe=ink,arc=0pt,boxrule=0pt,left=4mm,right=4mm,top=2.5mm,bottom=2.5mm]
\color{white}\sffamily{\LARGE\bfseries KNvkqqb}\hfill{\Large h\#9.5}\\[0.5mm]
{\small 6 Steine\quad Weiß am Zug}\hfill{\small Nr.~294\qquad\href{https://helpman.komtera.lt/?fen=7k/8/8/8/2q5/8/4q3/KN1b4&moves=9.5}{\color{white}\bfseries lösen\,\ensuremath{\nearrow}}}
\end{tcolorbox}
\vspace{2mm}
\noindent\begin{minipage}[t]{36mm}\vspace{0pt}\sffamily\small\raggedright
\hmfigure{längste eindeutige}{{\Large h\#9.5}}\hmfigure{eindeutig bei dieser Länge}{{\Large 140} Stellungen}\hmfigure{längstes Matt}{{\Large h\#11.5}}\hmfigure{in der Klasse}{15,5 Milliarden Stellungen,\\ 3,79 Milliarden mit Hilfsmatt}
\end{minipage}\hfill
\begin{minipage}[t]{80mm}\vspace{0pt}\begin{tikzpicture}\begin{axis}[width=72mm,height=42mm,ybar,bar width=1.6pt,xmin=-1,xmax=24,ymin=0,ymax=1.167,axis y line*=left,axis x line*=bottom,xlabel={\scriptsize Halbzüge bis Matt},ylabel={\scriptsize Stellungen (Mrd.)},tick label style={font=\tiny},label style={font=\scriptsize},enlarge x limits=false,ylabel shift=-3pt,xlabel shift=-3pt]\addplot[fill=ink!55,draw=none] coordinates {(0,4.4e-07) (1,2.285e-06) (2,1.781e-05) (3,0.0001149) (4,0.0005158) (5,0.002654) (6,0.005927) (7,0.02568) (8,0.04255) (9,0.1542) (10,0.2026) (11,0.6026) (12,0.4403) (13,1.061) (14,0.3429) (15,0.6335) (16,0.09977) (17,0.1416) (18,0.01321) (19,0.01662) (20,0.0002855) (21,0.0001758) (22,7.9e-08) (23,8e-08)};\draw[densely dotted] (axis cs:19,0) -- (axis cs:19,1.167);\end{axis}\begin{axis}[width=72mm,height=42mm,axis y line*=right,axis x line=none,ymode=log,xmin=-1,xmax=24,ylabel={\scriptsize eindeutig (log)},tick label style={font=\tiny},label style={font=\scriptsize},ylabel shift=-3pt]\addplot[mark=*,mark size=0.7pt,draw=black,thin] coordinates {(0,440) (1,2285) (2,13882) (3,68603) (4,205249) (5,548251) (6,293351) (7,889967) (8,811193) (9,1990449) (10,1338926) (11,2848138) (12,645286) (13,1160302) (14,97765) (15,136168) (16,12616) (17,14111) (18,319) (19,140)};\end{axis}\end{tikzpicture}\end{minipage}
\vspace{3mm}

\noindent\begin{minipage}[t]{60mm}\vspace{0pt}\sbox\hmboard{\chessboard[setfen=7k/8/8/8/2q5/8/4q3/KN1b4 w - - 0 1]}\usebox\hmboard\\[1mm]
\makebox[\wd\hmboard]{\sffamily\hspace{3mm}\textbf{h\#9.5}\hfill(2+4)}\end{minipage}\hfill
\begin{minipage}[t]{52mm}\vspace{0pt}
\textcolor{ink}{\sffamily\textsc{Lösung}}\par\smallskip\small
\begin{tabular}{@{}r@{\enspace}l@{\enspace}l@{}}
\textbf{1.} & \ldots & \symknight{}d2\\
\textbf{2.} & \symking{}g7 & \symking{}b2\\
\textbf{3.} & \symking{}f6 & \symking{}a3\\
\textbf{4.} & \symking{}e5 & \symknight{}xc4+\\
\textbf{5.} & \symking{}d4 & \symking{}b4\\
\textbf{6.} & \symbishop{}c2 & \symking{}a5\\
\textbf{7.} & \symking{}c3 & \symknight{}d2\\
\textbf{8.} & \symking{}b2 & \symking{}b4\\
\textbf{9.} & \symking{}a1 & \symking{}a3\\
\textbf{10.} & \symbishop{}b1 & \symknight{}b3\#\\
\end{tabular}\par\vspace{3mm}
\qrcode[height=12mm]{https://helpman.komtera.lt/?fen=7k/8/8/8/2q5/8/4q3/KN1b4&moves=9.5}\\{\sffamily\tiny\color{ink}scannen zum Lösen}
\end{minipage}
\vspace{2mm}

{\raggedright\noindent\textcolor{ink}{\sffamily\textsc{Themen}}\enspace \themetag{reines Matt} \themetag{Mustermatt} \themetag{Switchback} \themetag{Selbstblock}\par}

\clearpage
\phantomsection\label{hm:KNvkqrb}
\begin{tcolorbox}[enhanced,colback=ink,colframe=ink,arc=0pt,boxrule=0pt,left=4mm,right=4mm,top=2.5mm,bottom=2.5mm]
\color{white}\sffamily{\LARGE\bfseries KNvkqrb}\hfill{\Large h\#9.5}\\[0.5mm]
{\small 6 Steine\quad Weiß am Zug}\hfill{\small Nr.~295\qquad\href{https://helpman.komtera.lt/?fen=7k/6N1/8/8/8/8/b2q4/rK6&moves=9.5}{\color{white}\bfseries lösen\,\ensuremath{\nearrow}}}
\end{tcolorbox}
\vspace{2mm}
\noindent\begin{minipage}[t]{36mm}\vspace{0pt}\sffamily\small\raggedright
\hmfigure{längste eindeutige}{{\Large h\#9.5}}\hmfigure{eindeutig bei dieser Länge}{{\Large 1} Stellung}\hmfigure{längstes Matt}{{\Large h\#11.5}}\hmfigure{in der Klasse}{15,5 Milliarden Stellungen,\\ 8,23 Milliarden mit Hilfsmatt}
\end{minipage}\hfill
\begin{minipage}[t]{80mm}\vspace{0pt}\begin{tikzpicture}\begin{axis}[width=72mm,height=42mm,ybar,bar width=1.6pt,xmin=-1,xmax=24,ymin=0,ymax=2.478,axis y line*=left,axis x line*=bottom,xlabel={\scriptsize Halbzüge bis Matt},ylabel={\scriptsize Stellungen (Mrd.)},tick label style={font=\tiny},label style={font=\scriptsize},enlarge x limits=false,ylabel shift=-3pt,xlabel shift=-3pt]\addplot[fill=ink!55,draw=none] coordinates {(0,3.364e-06) (1,1.56e-05) (2,0.0001679) (3,0.00114) (4,0.006644) (5,0.03401) (6,0.108) (7,0.4241) (8,0.6794) (9,1.81) (10,1.199) (11,2.252) (12,0.7121) (13,0.8444) (14,0.08595) (15,0.05681) (16,0.001159) (17,0.007392) (18,6.95e-07) (19,0.002078) (20,6.3e-08) (21,2.93e-05) (23,4e-09)};\draw[densely dotted] (axis cs:19,0) -- (axis cs:19,2.478);\end{axis}\begin{axis}[width=72mm,height=42mm,axis y line*=right,axis x line=none,ymode=log,xmin=-1,xmax=24,ylabel={\scriptsize eindeutig (log)},tick label style={font=\tiny},label style={font=\scriptsize},ylabel shift=-3pt]\addplot[mark=*,mark size=0.7pt,draw=black,thin] coordinates {(0,3364) (1,15514) (2,142834) (3,776695) (4,2438339) (5,7679314) (6,8921951) (7,23946812) (8,13724758) (9,25430843) (10,6047787) (11,8181266) (12,620555) (13,447884) (14,4407) (15,10417) (16,1) (17,1993) (19,1)};\end{axis}\end{tikzpicture}\end{minipage}
\vspace{3mm}

\noindent\begin{minipage}[t]{60mm}\vspace{0pt}\sbox\hmboard{\chessboard[setfen=7k/6N1/8/8/8/8/b2q4/rK6 w - - 0 1]}\usebox\hmboard\\[1mm]
\makebox[\wd\hmboard]{\sffamily\hspace{3mm}\textbf{h\#9.5}\hfill(2+4)}\end{minipage}\hfill
\begin{minipage}[t]{52mm}\vspace{0pt}
\textcolor{ink}{\sffamily\textsc{Lösung}}\par\smallskip\small
\begin{tabular}{@{}r@{\enspace}l@{\enspace}l@{}}
\textbf{1.} & \ldots & \symking{}xa1\\
\textbf{2.} & \symking{}g8 & \symknight{}f5\\
\textbf{3.} & \symking{}f7 & \symknight{}e3\\
\textbf{4.} & \symking{}e6 & \symknight{}c2\\
\textbf{5.} & \symking{}d5 & \symking{}b2\\
\textbf{6.} & \symking{}c4 & \symking{}a3\\
\textbf{7.} & \symking{}c3 & \symking{}a4\\
\textbf{8.} & \symking{}b2 & \symknight{}d4\\
\textbf{9.} & \symking{}a1 & \symking{}a3\\
\textbf{10.} & \symbishop{}b1 & \symknight{}b3\#\\
\end{tabular}\par\vspace{3mm}
\qrcode[height=12mm]{https://helpman.komtera.lt/?fen=7k/6N1/8/8/8/8/b2q4/rK6&moves=9.5}\\{\sffamily\tiny\color{ink}scannen zum Lösen}
\end{minipage}
\vspace{2mm}

{\raggedright\noindent\textcolor{ink}{\sffamily\textsc{Themen}}\enspace \themetag{reines Matt} \themetag{Mustermatt} \themetag{Switchback} \themetag{Selbstblock} \themetag{Kniest} \themetag{ohne Schach}\par\smallskip{\footnotesize\textit{Hinweis:} das Diagramm hat keinen legalen letzten Zug, kann also in keiner Partie entstehen.}\par}

\clearpage
\phantomsection\label{hm:KNvkqrp}
\begin{tcolorbox}[enhanced,colback=ink,colframe=ink,arc=0pt,boxrule=0pt,left=4mm,right=4mm,top=2.5mm,bottom=2.5mm]
\color{white}\sffamily{\LARGE\bfseries KNvkqrp}\hfill{\Large h\#9.5}\\[0.5mm]
{\small 6 Steine\quad Weiß am Zug}\hfill{\small Nr.~296\qquad\href{https://helpman.komtera.lt/?fen=K1Nq4/1r6/5p2/2k5/8/8/8/8&moves=9.5}{\color{white}\bfseries lösen\,\ensuremath{\nearrow}}}
\end{tcolorbox}
\vspace{2mm}
\noindent\begin{minipage}[t]{36mm}\vspace{0pt}\sffamily\small\raggedright
\hmfigure{längste eindeutige}{{\Large h\#9.5}}\hmfigure{eindeutig bei dieser Länge}{{\Large 351} Stellungen}\hmfigure{längstes Matt}{{\Large h\#11.5}}\hmfigure{in der Klasse}{45,4 Milliarden Stellungen,\\ 25,8 Milliarden mit Hilfsmatt}
\end{minipage}\hfill
\begin{minipage}[t]{80mm}\vspace{0pt}\begin{tikzpicture}\begin{axis}[width=72mm,height=42mm,ybar,bar width=1.6pt,xmin=-1,xmax=24,ymin=0,ymax=7.33,axis y line*=left,axis x line*=bottom,xlabel={\scriptsize Halbzüge bis Matt},ylabel={\scriptsize Stellungen (Mrd.)},tick label style={font=\tiny},label style={font=\scriptsize},enlarge x limits=false,ylabel shift=-3pt,xlabel shift=-3pt]\addplot[fill=ink!55,draw=none] coordinates {(0,1.344e-05) (1,6.185e-05) (2,0.000532) (3,0.003608) (4,0.02053) (5,0.1024) (6,0.3317) (7,1.243) (8,2.186) (9,5.35) (10,3.938) (11,6.663) (12,2.517) (13,2.757) (14,0.4089) (15,0.2645) (16,0.00832) (17,0.02678) (18,1.944e-06) (19,0.006) (20,3.2e-07) (21,0.0006645) (22,5e-09) (23,2.3e-08)};\draw[densely dotted] (axis cs:19,0) -- (axis cs:19,7.33);\end{axis}\begin{axis}[width=72mm,height=42mm,axis y line*=right,axis x line=none,ymode=log,xmin=-1,xmax=24,ylabel={\scriptsize eindeutig (log)},tick label style={font=\tiny},label style={font=\scriptsize},ylabel shift=-3pt]\addplot[mark=*,mark size=0.7pt,draw=black,thin] coordinates {(0,13443) (1,61090) (2,438425) (3,2388188) (4,7296995) (5,22809565) (6,30133066) (7,73940829) (8,45365690) (9,72103630) (10,20347404) (11,26216369) (12,2875733) (13,2179754) (14,39990) (15,54878) (16,4) (17,5201) (19,351)};\end{axis}\end{tikzpicture}\end{minipage}
\vspace{3mm}

\noindent\begin{minipage}[t]{60mm}\vspace{0pt}\sbox\hmboard{\chessboard[setfen=K1Nq4/1r6/5p2/2k5/8/8/8/8 w - - 0 1]}\usebox\hmboard\\[1mm]
\makebox[\wd\hmboard]{\sffamily\hspace{3mm}\textbf{h\#9.5}\hfill(2+4)}\end{minipage}\hfill
\begin{minipage}[t]{52mm}\vspace{0pt}
\textcolor{ink}{\sffamily\textsc{Lösung}}\par\smallskip\small
\begin{tabular}{@{}r@{\enspace}l@{\enspace}l@{}}
\textbf{1.} & \ldots & \symking{}xb7\\
\textbf{2.} & f5 & \symknight{}b6\\
\textbf{3.} & \symking{}d6 & \symknight{}d7\\
\textbf{4.} & \symking{}e7 & \symking{}c6\\
\textbf{5.} & f4 & \symking{}d5\\
\textbf{6.} & f3 & \symking{}e5\\
\textbf{7.} & f2 & \symknight{}f6\\
\textbf{8.} & f1=\symrook{} & \symknight{}h5\\
\textbf{9.} & \symking{}e8 & \symking{}e6\\
\textbf{10.} & \symrook{}f8 & \symknight{}g7\#\\
\end{tabular}\par\vspace{3mm}
\qrcode[height=12mm]{https://helpman.komtera.lt/?fen=K1Nq4/1r6/5p2/2k5/8/8/8/8&moves=9.5}\\{\sffamily\tiny\color{ink}scannen zum Lösen}
\end{minipage}
\vspace{2mm}

{\raggedright\noindent\textcolor{ink}{\sffamily\textsc{Themen}}\enspace \themetag{reines Matt} \themetag{Mustermatt} \themetag{Idealmatt} \themetag{Umwandlung} \themetag{Unterverwandlung} \themetag{Selbstblock} \themetag{Phönix} \themetag{ohne Schach}\par\smallskip{\footnotesize\textit{Hinweis:} die Lösung beginnt mit einem Schlag.}\par}

\clearpage
\phantomsection\label{hm:KNvkrpp}
\begin{tcolorbox}[enhanced,colback=ink,colframe=ink,arc=0pt,boxrule=0pt,left=4mm,right=4mm,top=2.5mm,bottom=2.5mm]
\color{white}\sffamily{\LARGE\bfseries KNvkrpp}\hfill{\Large h\#9.5}\\[0.5mm]
{\small 6 Steine\quad Weiß am Zug}\hfill{\small Nr.~297\qquad\href{https://helpman.komtera.lt/?fen=2rK4/6p1/4k1p1/8/8/8/8/2N5&moves=9.5}{\color{white}\bfseries lösen\,\ensuremath{\nearrow}}}
\end{tcolorbox}
\vspace{2mm}
\noindent\begin{minipage}[t]{36mm}\vspace{0pt}\sffamily\small\raggedright
\hmfigure{längste eindeutige}{{\Large h\#9.5}}\hmfigure{eindeutig bei dieser Länge}{{\Large 1} Stellung}\hmfigure{längstes Matt}{{\Large h\#11.5}}\hmfigure{in der Klasse}{34,1 Milliarden Stellungen,\\ 11,4 Milliarden mit Hilfsmatt}
\end{minipage}\hfill
\begin{minipage}[t]{80mm}\vspace{0pt}\begin{tikzpicture}\begin{axis}[width=72mm,height=42mm,ybar,bar width=1.6pt,xmin=-1,xmax=24,ymin=0,ymax=2.655,axis y line*=left,axis x line*=bottom,xlabel={\scriptsize Halbzüge bis Matt},ylabel={\scriptsize Stellungen (Mrd.)},tick label style={font=\tiny},label style={font=\scriptsize},enlarge x limits=false,ylabel shift=-3pt,xlabel shift=-3pt]\addplot[fill=ink!55,draw=none] coordinates {(0,1.387e-05) (1,6.352e-05) (2,0.0003536) (3,0.002343) (4,0.01162) (5,0.05107) (6,0.1363) (7,0.3909) (8,0.7134) (9,1.364) (10,1.729) (11,2.414) (12,1.769) (13,1.596) (14,0.6905) (15,0.3694) (16,0.09651) (17,0.02396) (18,0.001808) (19,0.0003639) (20,8e-09) (21,2.69e-05) (23,4.6e-08)};\draw[densely dotted] (axis cs:19,0) -- (axis cs:19,2.655);\end{axis}\begin{axis}[width=72mm,height=42mm,axis y line*=right,axis x line=none,ymode=log,xmin=-1,xmax=24,ylabel={\scriptsize eindeutig (log)},tick label style={font=\tiny},label style={font=\scriptsize},ylabel shift=-3pt]\addplot[mark=*,mark size=0.7pt,draw=black,thin] coordinates {(0,13866) (1,61248) (2,281768) (3,1476094) (4,5004953) (5,10112702) (6,9258267) (7,12483859) (8,10981081) (9,12884520) (10,9758462) (11,6530512) (12,1781327) (13,808180) (14,207279) (15,44266) (16,3579) (17,434) (19,1)};\end{axis}\end{tikzpicture}\end{minipage}
\vspace{3mm}

\noindent\begin{minipage}[t]{60mm}\vspace{0pt}\sbox\hmboard{\chessboard[setfen=2rK4/6p1/4k1p1/8/8/8/8/2N5 w - - 0 1]}\usebox\hmboard\\[1mm]
\makebox[\wd\hmboard]{\sffamily\hspace{3mm}\textbf{h\#9.5}\hfill(2+4)}\end{minipage}\hfill
\begin{minipage}[t]{52mm}\vspace{0pt}
\textcolor{ink}{\sffamily\textsc{Lösung}}\par\smallskip\small
\begin{tabular}{@{}r@{\enspace}l@{\enspace}l@{}}
\textbf{1.} & \ldots & \symking{}xc8\\
\textbf{2.} & g5 & \symking{}d8\\
\textbf{3.} & g4 & \symking{}e8\\
\textbf{4.} & g3 & \symking{}f8\\
\textbf{5.} & g2 & \symking{}xg7\\
\textbf{6.} & g1=\symrook{}+ & \symking{}h7\\
\textbf{7.} & \symking{}f7 & \symking{}h6\\
\textbf{8.} & \symking{}g8 & \symknight{}d3\\
\textbf{9.} & \symking{}h8 & \symknight{}e5\\
\textbf{10.} & \symrook{}g8 & \symknight{}f7\#\\
\end{tabular}\par\vspace{3mm}
\qrcode[height=12mm]{https://helpman.komtera.lt/?fen=2rK4/6p1/4k1p1/8/8/8/8/2N5&moves=9.5}\\{\sffamily\tiny\color{ink}scannen zum Lösen}
\end{minipage}
\vspace{2mm}

{\raggedright\noindent\textcolor{ink}{\sffamily\textsc{Themen}}\enspace \themetag{reines Matt} \themetag{Mustermatt} \themetag{Idealmatt} \themetag{Umwandlung} \themetag{Unterverwandlung} \themetag{Switchback} \themetag{Selbstblock} \themetag{Phönix}\par\smallskip{\footnotesize\textit{Hinweis:} die Partei am Zug steht im Diagramm im Schach; die Lösung beginnt mit einem Schlag.}\par}

\clearpage
\phantomsection\label{hm:KPvkbnp}
\begin{tcolorbox}[enhanced,colback=ink,colframe=ink,arc=0pt,boxrule=0pt,left=4mm,right=4mm,top=2.5mm,bottom=2.5mm]
\color{white}\sffamily{\LARGE\bfseries KPvkbnp}\hfill{\Large h\#9.5}\\[0.5mm]
{\small 6 Steine\quad Weiß am Zug}\hfill{\small Nr.~298\qquad\href{https://helpman.komtera.lt/?fen=6k1/8/7n/8/8/p7/P7/K3b3&moves=9.5}{\color{white}\bfseries lösen\,\ensuremath{\nearrow}}}
\end{tcolorbox}
\vspace{2mm}
\noindent\begin{minipage}[t]{36mm}\vspace{0pt}\sffamily\small\raggedright
\hmfigure{längste eindeutige}{{\Large h\#9.5}}\hmfigure{eindeutig bei dieser Länge}{{\Large 14} Stellungen}\hmfigure{längstes Matt}{{\Large h\#14.5}}\hmfigure{in der Klasse}{34,1 Milliarden Stellungen,\\ 23,7 Milliarden mit Hilfsmatt}
\end{minipage}\hfill
\begin{minipage}[t]{80mm}\vspace{0pt}\begin{tikzpicture}\begin{axis}[width=72mm,height=42mm,ybar,bar width=1.6pt,xmin=-1,xmax=30,ymin=0,ymax=4.298,axis y line*=left,axis x line*=bottom,xlabel={\scriptsize Halbzüge bis Matt},ylabel={\scriptsize Stellungen (Mrd.)},tick label style={font=\tiny},label style={font=\scriptsize},enlarge x limits=false,ylabel shift=-3pt,xlabel shift=-3pt]\addplot[fill=ink!55,draw=none] coordinates {(0,1.845e-05) (1,0.002766) (2,0.008205) (3,0.06312) (4,0.1502) (5,0.6628) (6,1.039) (7,2.225) (8,2.279) (9,3.384) (10,2.793) (11,3.907) (12,3.142) (13,2.593) (14,1.101) (15,0.2846) (16,0.06237) (17,0.03118) (18,0.006305) (19,0.00481) (20,0.0008264) (21,0.000511) (22,1.838e-05) (23,1.393e-05) (24,5.9e-08) (25,1.627e-06) (27,6.59e-07) (29,4.22e-07)};\draw[densely dotted] (axis cs:19,0) -- (axis cs:19,4.298);\end{axis}\begin{axis}[width=72mm,height=42mm,axis y line*=right,axis x line=none,ymode=log,xmin=-1,xmax=30,ylabel={\scriptsize eindeutig (log)},tick label style={font=\tiny},label style={font=\scriptsize},ylabel shift=-3pt]\addplot[mark=*,mark size=0.7pt,draw=black,thin] coordinates {(0,18453) (1,686416) (2,1522759) (3,14589296) (4,25906853) (5,62570911) (6,37910299) (7,43653213) (8,23201764) (9,29556148) (10,16417009) (11,21385147) (12,9600520) (13,3127425) (14,407753) (15,98632) (16,8378) (17,2265) (18,64) (19,14)};\end{axis}\end{tikzpicture}\end{minipage}
\vspace{3mm}

\noindent\begin{minipage}[t]{60mm}\vspace{0pt}\sbox\hmboard{\chessboard[setfen=6k1/8/7n/8/8/p7/P7/K3b3 w - - 0 1]}\usebox\hmboard\\[1mm]
\makebox[\wd\hmboard]{\sffamily\hspace{3mm}\textbf{h\#9.5}\hfill(2+4)}\end{minipage}\hfill
\begin{minipage}[t]{52mm}\vspace{0pt}
\textcolor{ink}{\sffamily\textsc{Lösung}}\par\smallskip\small
\begin{tabular}{@{}r@{\enspace}l@{\enspace}l@{}}
\textbf{1.} & \ldots & \symking{}b1\\
\textbf{2.} & \symknight{}f5 & \symking{}c2\\
\textbf{3.} & \symknight{}d4+ & \symking{}c1\\
\textbf{4.} & \symknight{}b3+ & axb3\\
\textbf{5.} & \symking{}f7 & b4\\
\textbf{6.} & \symking{}e6 & b5\\
\textbf{7.} & \symking{}d5 & b6\\
\textbf{8.} & \symking{}c4 & b7\\
\textbf{9.} & \symking{}b3 & b8=\symqueen{}+\\
\textbf{10.} & \symking{}a2 & \symqueen{}b1\#\\
\end{tabular}\par\vspace{3mm}
\qrcode[height=12mm]{https://helpman.komtera.lt/?fen=6k1/8/7n/8/8/p7/P7/K3b3&moves=9.5}\\{\sffamily\tiny\color{ink}scannen zum Lösen}
\end{minipage}
\vspace{2mm}

{\raggedright\noindent\textcolor{ink}{\sffamily\textsc{Themen}}\enspace \themetag{Satzspiel} \themetag{Umwandlung} \themetag{Excelsior}\par}

\clearpage
\phantomsection\label{hm:KPvknnp}
\begin{tcolorbox}[enhanced,colback=ink,colframe=ink,arc=0pt,boxrule=0pt,left=4mm,right=4mm,top=2.5mm,bottom=2.5mm]
\color{white}\sffamily{\LARGE\bfseries KPvknnp}\hfill{\Large h\#9.5}\\[0.5mm]
{\small 6 Steine\quad Weiß am Zug}\hfill{\small Nr.~299\qquad\href{https://helpman.komtera.lt/?fen=6k1/8/7n/8/8/p7/P7/K2n4&moves=9.5}{\color{white}\bfseries lösen\,\ensuremath{\nearrow}}}
\end{tcolorbox}
\vspace{2mm}
\noindent\begin{minipage}[t]{36mm}\vspace{0pt}\sffamily\small\raggedright
\hmfigure{längste eindeutige}{{\Large h\#9.5}}\hmfigure{eindeutig bei dieser Länge}{{\Large 1} Stellung}\hmfigure{längstes Matt}{{\Large h\#11.5}}\hmfigure{in der Klasse}{34,1 Milliarden Stellungen,\\ 12,1 Milliarden mit Hilfsmatt}
\end{minipage}\hfill
\begin{minipage}[t]{80mm}\vspace{0pt}\begin{tikzpicture}\begin{axis}[width=72mm,height=42mm,ybar,bar width=1.6pt,xmin=-1,xmax=24,ymin=0,ymax=2.068,axis y line*=left,axis x line*=bottom,xlabel={\scriptsize Halbzüge bis Matt},ylabel={\scriptsize Stellungen (Mrd.)},tick label style={font=\tiny},label style={font=\scriptsize},enlarge x limits=false,ylabel shift=-3pt,xlabel shift=-3pt]\addplot[fill=ink!55,draw=none] coordinates {(0,7.174e-06) (1,0.001757) (2,0.004602) (3,0.03283) (4,0.07206) (5,0.3014) (6,0.4582) (7,1.01) (8,1.102) (9,1.626) (10,1.413) (11,1.88) (12,1.585) (13,1.469) (14,0.843) (15,0.2363) (16,0.07054) (17,0.01936) (18,0.004132) (19,0.002286) (20,0.0002052) (21,0.0001084) (22,2.055e-06) (23,2.8e-07)};\draw[densely dotted] (axis cs:19,0) -- (axis cs:19,2.068);\end{axis}\begin{axis}[width=72mm,height=42mm,axis y line*=right,axis x line=none,ymode=log,xmin=-1,xmax=24,ylabel={\scriptsize eindeutig (log)},tick label style={font=\tiny},label style={font=\scriptsize},ylabel shift=-3pt]\addplot[mark=*,mark size=0.7pt,draw=black,thin] coordinates {(0,7174) (1,366417) (2,742951) (3,6428158) (4,10565107) (5,26712798) (6,18173918) (7,19292217) (8,11147720) (9,13394624) (10,8116036) (11,10992310) (12,5850228) (13,2103370) (14,298308) (15,61740) (16,6360) (17,1835) (18,33) (19,1)};\end{axis}\end{tikzpicture}\end{minipage}
\vspace{3mm}

\noindent\begin{minipage}[t]{60mm}\vspace{0pt}\sbox\hmboard{\chessboard[setfen=6k1/8/7n/8/8/p7/P7/K2n4 w - - 0 1]}\usebox\hmboard\\[1mm]
\makebox[\wd\hmboard]{\sffamily\hspace{3mm}\textbf{h\#9.5}\hfill(2+4)}\end{minipage}\hfill
\begin{minipage}[t]{52mm}\vspace{0pt}
\textcolor{ink}{\sffamily\textsc{Lösung}}\par\smallskip\small
\begin{tabular}{@{}r@{\enspace}l@{\enspace}l@{}}
\textbf{1.} & \ldots & \symking{}b1\\
\textbf{2.} & \symknight{}f5 & \symking{}c2\\
\textbf{3.} & \symknight{}d4+ & \symking{}c1\\
\textbf{4.} & \symknight{}b3+ & axb3\\
\textbf{5.} & \symking{}f7 & b4\\
\textbf{6.} & \symking{}e6 & b5\\
\textbf{7.} & \symking{}d5 & b6\\
\textbf{8.} & \symking{}c4 & b7\\
\textbf{9.} & \symking{}b3 & b8=\symqueen{}+\\
\textbf{10.} & \symking{}a2 & \symqueen{}b1\#\\
\end{tabular}\par\vspace{3mm}
\qrcode[height=12mm]{https://helpman.komtera.lt/?fen=6k1/8/7n/8/8/p7/P7/K2n4&moves=9.5}\\{\sffamily\tiny\color{ink}scannen zum Lösen}
\end{minipage}
\vspace{2mm}

{\raggedright\noindent\textcolor{ink}{\sffamily\textsc{Themen}}\enspace \themetag{Satzspiel} \themetag{Umwandlung} \themetag{Excelsior}\par}

\clearpage
\phantomsection\label{hm:KPvknpp}
\begin{tcolorbox}[enhanced,colback=ink,colframe=ink,arc=0pt,boxrule=0pt,left=4mm,right=4mm,top=2.5mm,bottom=2.5mm]
\color{white}\sffamily{\LARGE\bfseries KPvknpp}\hfill{\Large h\#9.5}\\[0.5mm]
{\small 6 Steine\quad Weiß am Zug}\hfill{\small Nr.~300\qquad\href{https://helpman.komtera.lt/?fen=8/8/2p5/8/8/1p5n/1P6/K2k4&moves=9.5}{\color{white}\bfseries lösen\,\ensuremath{\nearrow}}}
\end{tcolorbox}
\vspace{2mm}
\noindent\begin{minipage}[t]{36mm}\vspace{0pt}\sffamily\small\raggedright
\hmfigure{längste eindeutige}{{\Large h\#9.5}}\hmfigure{eindeutig bei dieser Länge}{{\Large 112} Stellungen}\hmfigure{längstes Matt}{{\Large h\#15.5}}\hmfigure{in der Klasse}{25,6 Milliarden Stellungen,\\ 9,25 Milliarden mit Hilfsmatt}
\end{minipage}\hfill
\begin{minipage}[t]{80mm}\vspace{0pt}\begin{tikzpicture}\begin{axis}[width=72mm,height=42mm,ybar,bar width=1.6pt,xmin=-1,xmax=32,ymin=0,ymax=1.521,axis y line*=left,axis x line*=bottom,xlabel={\scriptsize Halbzüge bis Matt},ylabel={\scriptsize Stellungen (Mrd.)},tick label style={font=\tiny},label style={font=\scriptsize},enlarge x limits=false,ylabel shift=-3pt,xlabel shift=-3pt]\addplot[fill=ink!55,draw=none] coordinates {(0,6.85e-06) (1,0.001861) (2,0.004458) (3,0.02942) (4,0.06632) (5,0.2668) (6,0.4273) (7,0.828) (8,0.9254) (9,1.228) (10,1.098) (11,1.383) (12,1.232) (13,0.9446) (14,0.5352) (15,0.1621) (16,0.0651) (17,0.03367) (18,0.01021) (19,0.00668) (20,0.001154) (21,0.00065) (22,3.368e-05) (23,1.53e-05) (24,2.15e-07) (25,8.77e-07) (26,6.5e-08) (27,2.38e-07) (28,3e-09) (29,8.1e-08) (31,2.8e-08)};\draw[densely dotted] (axis cs:19,0) -- (axis cs:19,1.521);\end{axis}\begin{axis}[width=72mm,height=42mm,axis y line*=right,axis x line=none,ymode=log,xmin=-1,xmax=32,ylabel={\scriptsize eindeutig (log)},tick label style={font=\tiny},label style={font=\scriptsize},ylabel shift=-3pt]\addplot[mark=*,mark size=0.7pt,draw=black,thin] coordinates {(0,6850) (1,418200) (2,832615) (3,6261963) (4,10952295) (5,24054389) (6,17383457) (7,15796163) (8,10028730) (9,11157563) (10,7000318) (11,8504032) (12,4408498) (13,1003851) (14,282031) (15,58698) (16,13728) (17,3345) (18,615) (19,112)};\end{axis}\end{tikzpicture}\end{minipage}
\vspace{3mm}

\noindent\begin{minipage}[t]{60mm}\vspace{0pt}\sbox\hmboard{\chessboard[setfen=8/8/2p5/8/8/1p5n/1P6/K2k4 w - - 0 1]}\usebox\hmboard\\[1mm]
\makebox[\wd\hmboard]{\sffamily\hspace{3mm}\textbf{h\#9.5}\hfill(2+4)}\end{minipage}\hfill
\begin{minipage}[t]{52mm}\vspace{0pt}
\textcolor{ink}{\sffamily\textsc{Lösung}}\par\smallskip\small
\begin{tabular}{@{}r@{\enspace}l@{\enspace}l@{}}
\textbf{1.} & \ldots & \symking{}b1\\
\textbf{2.} & c5 & \symking{}a1\\
\textbf{3.} & c4 & \symking{}b1\\
\textbf{4.} & c3 & bxc3\\
\textbf{5.} & \symknight{}f4 & c4\\
\textbf{6.} & \symknight{}d5 & cxd5\\
\textbf{7.} & \symking{}d2 & d6\\
\textbf{8.} & \symking{}c3 & d7\\
\textbf{9.} & \symking{}b4 & d8=\symqueen{}\\
\textbf{10.} & \symking{}a3 & \symqueen{}a5\#\\
\end{tabular}\par\vspace{3mm}
\qrcode[height=12mm]{https://helpman.komtera.lt/?fen=8/8/2p5/8/8/1p5n/1P6/K2k4&moves=9.5}\\{\sffamily\tiny\color{ink}scannen zum Lösen}
\end{minipage}
\vspace{2mm}

{\raggedright\noindent\textcolor{ink}{\sffamily\textsc{Themen}}\enspace \themetag{Satzspiel} \themetag{Umwandlung} \themetag{Excelsior} \themetag{Switchback} \themetag{Pendel} \themetag{ohne Schach}\par}

\clearpage
\phantomsection\label{hm:KPvkqbp}
\begin{tcolorbox}[enhanced,colback=ink,colframe=ink,arc=0pt,boxrule=0pt,left=4mm,right=4mm,top=2.5mm,bottom=2.5mm]
\color{white}\sffamily{\LARGE\bfseries KPvkqbp}\hfill{\Large h\#9.5}\\[0.5mm]
{\small 6 Steine\quad Weiß am Zug}\hfill{\small Nr.~301\qquad\href{https://helpman.komtera.lt/?fen=8/8/6p1/5bq1/5P2/k7/8/K7&moves=9.5}{\color{white}\bfseries lösen\,\ensuremath{\nearrow}}}
\end{tcolorbox}
\vspace{2mm}
\noindent\begin{minipage}[t]{36mm}\vspace{0pt}\sffamily\small\raggedright
\hmfigure{längste eindeutige}{{\Large h\#9.5}}\hmfigure{eindeutig bei dieser Länge}{{\Large 9} Stellungen}\hmfigure{längstes Matt}{{\Large h\#14.5}}\hmfigure{in der Klasse}{34,1 Milliarden Stellungen,\\ 20,8 Milliarden mit Hilfsmatt}
\end{minipage}\hfill
\begin{minipage}[t]{80mm}\vspace{0pt}\begin{tikzpicture}\begin{axis}[width=72mm,height=42mm,ybar,bar width=1.6pt,xmin=-1,xmax=30,ymin=0,ymax=4.292,axis y line*=left,axis x line*=bottom,xlabel={\scriptsize Halbzüge bis Matt},ylabel={\scriptsize Stellungen (Mrd.)},tick label style={font=\tiny},label style={font=\scriptsize},enlarge x limits=false,ylabel shift=-3pt,xlabel shift=-3pt]\addplot[fill=ink!55,draw=none] coordinates {(0,8.603e-06) (1,0.001136) (2,0.004881) (3,0.04449) (4,0.1178) (5,0.6165) (6,0.8906) (7,2.335) (8,1.79) (9,3.426) (10,2.137) (11,3.901) (12,2.274) (13,2.49) (14,0.5155) (15,0.2472) (16,0.0221) (17,0.01519) (18,0.0002991) (19,0.0005313) (20,2.174e-06) (21,0.0004577) (22,8e-09) (23,0.0002947) (24,2e-09) (25,9.758e-05) (27,5.294e-05) (29,1.865e-05)};\draw[densely dotted] (axis cs:19,0) -- (axis cs:19,4.292);\end{axis}\begin{axis}[width=72mm,height=42mm,axis y line*=right,axis x line=none,ymode=log,xmin=-1,xmax=30,ylabel={\scriptsize eindeutig (log)},tick label style={font=\tiny},label style={font=\scriptsize},ylabel shift=-3pt]\addplot[mark=*,mark size=0.7pt,draw=black,thin] coordinates {(0,8603) (1,351603) (2,1054145) (3,12315025) (4,20698344) (5,63134594) (6,28583661) (7,49207314) (8,18047095) (9,30645632) (10,11277648) (11,17947410) (12,4798765) (13,2810000) (14,161101) (15,71281) (16,612) (17,897) (18,2) (19,9)};\end{axis}\end{tikzpicture}\end{minipage}
\vspace{3mm}

\noindent\begin{minipage}[t]{60mm}\vspace{0pt}\sbox\hmboard{\chessboard[setfen=8/8/6p1/5bq1/5P2/k7/8/K7 w - - 0 1]}\usebox\hmboard\\[1mm]
\makebox[\wd\hmboard]{\sffamily\hspace{3mm}\textbf{h\#9.5}\hfill(2+4)}\end{minipage}\hfill
\begin{minipage}[t]{52mm}\vspace{0pt}
\textcolor{ink}{\sffamily\textsc{Lösung}}\par\smallskip\small
\begin{tabular}{@{}r@{\enspace}l@{\enspace}l@{}}
\textbf{1.} & \ldots & fxg5\\
\textbf{2.} & \symking{}b4 & \symking{}b2\\
\textbf{3.} & \symking{}c5 & \symking{}c3\\
\textbf{4.} & \symking{}d6 & \symking{}d4\\
\textbf{5.} & \symking{}e7 & \symking{}e5\\
\textbf{6.} & \symking{}f8 & \symking{}f6\\
\textbf{7.} & \symbishop{}e6 & \symking{}xg6\\
\textbf{8.} & \symking{}g8 & \symking{}h6\\
\textbf{9.} & \symking{}h8 & g6\\
\textbf{10.} & \symbishop{}g8 & g7\#\\
\end{tabular}\par\vspace{3mm}
\qrcode[height=12mm]{https://helpman.komtera.lt/?fen=8/8/6p1/5bq1/5P2/k7/8/K7&moves=9.5}\\{\sffamily\tiny\color{ink}scannen zum Lösen}
\end{minipage}
\vspace{2mm}

{\raggedright\noindent\textcolor{ink}{\sffamily\textsc{Themen}}\enspace \themetag{reines Matt} \themetag{Mustermatt} \themetag{Idealmatt} \themetag{Selbstblock} \themetag{ohne Schach}\par\smallskip{\footnotesize\textit{Hinweis:} die Lösung beginnt mit einem Schlag.}\par}

\clearpage
\phantomsection\label{hm:KPvkqnp}
\begin{tcolorbox}[enhanced,colback=ink,colframe=ink,arc=0pt,boxrule=0pt,left=4mm,right=4mm,top=2.5mm,bottom=2.5mm]
\color{white}\sffamily{\LARGE\bfseries KPvkqnp}\hfill{\Large h\#9.5}\\[0.5mm]
{\small 6 Steine\quad Weiß am Zug}\hfill{\small Nr.~302\qquad\href{https://helpman.komtera.lt/?fen=8/7k/n7/8/7p/8/q6P/K7&moves=9.5}{\color{white}\bfseries lösen\,\ensuremath{\nearrow}}}
\end{tcolorbox}
\vspace{2mm}
\noindent\begin{minipage}[t]{36mm}\vspace{0pt}\sffamily\small\raggedright
\hmfigure{längste eindeutige}{{\Large h\#9.5}}\hmfigure{eindeutig bei dieser Länge}{{\Large 6} Stellungen}\hmfigure{längstes Matt}{{\Large h\#11.5}}\hmfigure{in der Klasse}{34,1 Milliarden Stellungen,\\ 21,3 Milliarden mit Hilfsmatt}
\end{minipage}\hfill
\begin{minipage}[t]{80mm}\vspace{0pt}\begin{tikzpicture}\begin{axis}[width=72mm,height=42mm,ybar,bar width=1.6pt,xmin=-1,xmax=24,ymin=0,ymax=4.261,axis y line*=left,axis x line*=bottom,xlabel={\scriptsize Halbzüge bis Matt},ylabel={\scriptsize Stellungen (Mrd.)},tick label style={font=\tiny},label style={font=\scriptsize},enlarge x limits=false,ylabel shift=-3pt,xlabel shift=-3pt]\addplot[fill=ink!55,draw=none] coordinates {(0,5.762e-06) (1,0.001395) (2,0.005777) (3,0.04858) (4,0.1177) (5,0.5793) (6,0.8293) (7,2.145) (8,1.799) (9,3.347) (10,2.203) (11,3.874) (12,2.478) (13,2.75) (14,0.7594) (15,0.3399) (16,0.03545) (17,0.02136) (18,0.0007312) (19,0.0004697) (20,1.087e-06) (21,3.793e-05) (23,9.12e-07)};\draw[densely dotted] (axis cs:19,0) -- (axis cs:19,4.261);\end{axis}\begin{axis}[width=72mm,height=42mm,axis y line*=right,axis x line=none,ymode=log,xmin=-1,xmax=24,ylabel={\scriptsize eindeutig (log)},tick label style={font=\tiny},label style={font=\scriptsize},ylabel shift=-3pt]\addplot[mark=*,mark size=0.7pt,draw=black,thin] coordinates {(0,5762) (1,380904) (2,1079172) (3,11400540) (4,17942561) (5,54387936) (6,25663135) (7,42252720) (8,16459547) (9,27373416) (10,11448228) (11,18986645) (12,6304161) (13,3788787) (14,211402) (15,96557) (16,979) (17,912) (18,1) (19,6)};\end{axis}\end{tikzpicture}\end{minipage}
\vspace{3mm}

\noindent\begin{minipage}[t]{60mm}\vspace{0pt}\sbox\hmboard{\chessboard[setfen=8/7k/n7/8/7p/8/q6P/K7 w - - 0 1]}\usebox\hmboard\\[1mm]
\makebox[\wd\hmboard]{\sffamily\hspace{3mm}\textbf{h\#9.5}\hfill(2+4)}\end{minipage}\hfill
\begin{minipage}[t]{52mm}\vspace{0pt}
\textcolor{ink}{\sffamily\textsc{Lösung}}\par\smallskip\small
\begin{tabular}{@{}r@{\enspace}l@{\enspace}l@{}}
\textbf{1.} & \ldots & \symking{}xa2\\
\textbf{2.} & \symknight{}c5 & \symking{}b2\\
\textbf{3.} & \symknight{}e4 & \symking{}a3\\
\textbf{4.} & \symknight{}g3 & hxg3\\
\textbf{5.} & \symking{}g6 & gxh4\\
\textbf{6.} & \symking{}f5 & h5\\
\textbf{7.} & \symking{}e4 & h6\\
\textbf{8.} & \symking{}d3 & h7\\
\textbf{9.} & \symking{}c2 & h8=\symqueen{}\\
\textbf{10.} & \symking{}b1 & \symqueen{}b2\#\\
\end{tabular}\par\vspace{3mm}
\qrcode[height=12mm]{https://helpman.komtera.lt/?fen=8/7k/n7/8/7p/8/q6P/K7&moves=9.5}\\{\sffamily\tiny\color{ink}scannen zum Lösen}
\end{minipage}
\vspace{2mm}

{\raggedright\noindent\textcolor{ink}{\sffamily\textsc{Themen}}\enspace \themetag{Umwandlung} \themetag{Excelsior} \themetag{ohne Schach}\par\smallskip{\footnotesize\textit{Hinweis:} die Partei am Zug steht im Diagramm im Schach; die Lösung beginnt mit einem Schlag.}\par}

\clearpage
\phantomsection\label{hm:KPvkqpp}
\begin{tcolorbox}[enhanced,colback=ink,colframe=ink,arc=0pt,boxrule=0pt,left=4mm,right=4mm,top=2.5mm,bottom=2.5mm]
\color{white}\sffamily{\LARGE\bfseries KPvkqpp}\hfill{\Large h\#9.5}\\[0.5mm]
{\small 6 Steine\quad Weiß am Zug}\hfill{\small Nr.~303\qquad\href{https://helpman.komtera.lt/?fen=6k1/8/1p6/8/8/p5q1/P7/K7&moves=9.5}{\color{white}\bfseries lösen\,\ensuremath{\nearrow}}}
\end{tcolorbox}
\vspace{2mm}
\noindent\begin{minipage}[t]{36mm}\vspace{0pt}\sffamily\small\raggedright
\hmfigure{längste eindeutige}{{\Large h\#9.5}}\hmfigure{eindeutig bei dieser Länge}{{\Large 236} Stellungen}\hmfigure{längstes Matt}{{\Large h\#15.5}}\hmfigure{in der Klasse}{25,6 Milliarden Stellungen,\\ 8,1 Milliarden mit Hilfsmatt}
\end{minipage}\hfill
\begin{minipage}[t]{80mm}\vspace{0pt}\begin{tikzpicture}\begin{axis}[width=72mm,height=42mm,ybar,bar width=1.6pt,xmin=-1,xmax=32,ymin=0,ymax=1.549,axis y line*=left,axis x line*=bottom,xlabel={\scriptsize Halbzüge bis Matt},ylabel={\scriptsize Stellungen (Mrd.)},tick label style={font=\tiny},label style={font=\scriptsize},enlarge x limits=false,ylabel shift=-3pt,xlabel shift=-3pt]\addplot[fill=ink!55,draw=none] coordinates {(0,2.892e-06) (1,0.0007487) (2,0.00295) (3,0.02295) (4,0.05692) (5,0.2673) (6,0.3939) (7,0.907) (8,0.7467) (9,1.258) (10,0.8425) (11,1.409) (12,0.9215) (13,0.8683) (14,0.2378) (15,0.1188) (16,0.02397) (17,0.0161) (18,0.000877) (19,0.001425) (20,2.892e-06) (21,0.0006175) (23,0.000231) (25,5.943e-05) (27,2.218e-05) (29,1.438e-05) (31,2.912e-06)};\draw[densely dotted] (axis cs:19,0) -- (axis cs:19,1.549);\end{axis}\begin{axis}[width=72mm,height=42mm,axis y line*=right,axis x line=none,ymode=log,xmin=-1,xmax=32,ylabel={\scriptsize eindeutig (log)},tick label style={font=\tiny},label style={font=\scriptsize},ylabel shift=-3pt]\addplot[mark=*,mark size=0.7pt,draw=black,thin] coordinates {(0,2892) (1,218333) (2,647088) (3,5792834) (4,9602488) (5,25442322) (6,12019644) (7,17024733) (8,6995608) (9,10655217) (10,4571197) (11,6669307) (12,2094612) (13,985671) (14,104171) (15,52673) (16,1455) (17,1130) (18,79) (19,236)};\end{axis}\end{tikzpicture}\end{minipage}
\vspace{3mm}

\noindent\begin{minipage}[t]{60mm}\vspace{0pt}\sbox\hmboard{\chessboard[setfen=6k1/8/1p6/8/8/p5q1/P7/K7 w - - 0 1]}\usebox\hmboard\\[1mm]
\makebox[\wd\hmboard]{\sffamily\hspace{3mm}\textbf{h\#9.5}\hfill(2+4)}\end{minipage}\hfill
\begin{minipage}[t]{52mm}\vspace{0pt}
\textcolor{ink}{\sffamily\textsc{Lösung}}\par\smallskip\small
\begin{tabular}{@{}r@{\enspace}l@{\enspace}l@{}}
\textbf{1.} & \ldots & \symking{}b1\\
\textbf{2.} & b5 & \symking{}c2\\
\textbf{3.} & b4 & \symking{}c1\\
\textbf{4.} & b3 & axb3\\
\textbf{5.} & \symking{}f7 & b4\\
\textbf{6.} & \symking{}e6 & b5\\
\textbf{7.} & \symking{}d5 & b6\\
\textbf{8.} & \symking{}c4 & b7\\
\textbf{9.} & \symking{}b3 & b8=\symqueen{}+\\
\textbf{10.} & \symking{}a2 & \symqueen{}b1\#\\
\end{tabular}\par\vspace{3mm}
\qrcode[height=12mm]{https://helpman.komtera.lt/?fen=6k1/8/1p6/8/8/p5q1/P7/K7&moves=9.5}\\{\sffamily\tiny\color{ink}scannen zum Lösen}
\end{minipage}
\vspace{2mm}

{\raggedright\noindent\textcolor{ink}{\sffamily\textsc{Themen}}\enspace \themetag{Satzspiel} \themetag{Umwandlung} \themetag{Excelsior}\par}

\clearpage
\phantomsection\label{hm:KPvkrbp}
\begin{tcolorbox}[enhanced,colback=ink,colframe=ink,arc=0pt,boxrule=0pt,left=4mm,right=4mm,top=2.5mm,bottom=2.5mm]
\color{white}\sffamily{\LARGE\bfseries KPvkrbp}\hfill{\Large h\#9.5}\\[0.5mm]
{\small 6 Steine\quad Weiß am Zug}\hfill{\small Nr.~304\qquad\href{https://helpman.komtera.lt/?fen=8/8/6p1/8/6P1/k7/1rb5/K7&moves=9.5}{\color{white}\bfseries lösen\,\ensuremath{\nearrow}}}
\end{tcolorbox}
\vspace{2mm}
\noindent\begin{minipage}[t]{36mm}\vspace{0pt}\sffamily\small\raggedright
\hmfigure{längste eindeutige}{{\Large h\#9.5}}\hmfigure{eindeutig bei dieser Länge}{{\Large 11} Stellungen}\hmfigure{längstes Matt}{{\Large h\#15}}\hmfigure{in der Klasse}{34,1 Milliarden Stellungen,\\ 22,4 Milliarden mit Hilfsmatt}
\end{minipage}\hfill
\begin{minipage}[t]{80mm}\vspace{0pt}\begin{tikzpicture}\begin{axis}[width=72mm,height=42mm,ybar,bar width=1.6pt,xmin=-1,xmax=31,ymin=0,ymax=4.184,axis y line*=left,axis x line*=bottom,xlabel={\scriptsize Halbzüge bis Matt},ylabel={\scriptsize Stellungen (Mrd.)},tick label style={font=\tiny},label style={font=\scriptsize},enlarge x limits=false,ylabel shift=-3pt,xlabel shift=-3pt]\addplot[fill=ink!55,draw=none] coordinates {(0,1.173e-05) (1,0.001977) (2,0.008001) (3,0.06428) (4,0.1994) (5,0.8901) (6,1.525) (7,3.05) (8,2.461) (9,3.739) (10,2.539) (11,3.804) (12,2.292) (13,1.466) (14,0.2286) (15,0.1083) (16,0.02135) (17,0.01453) (18,0.0004899) (19,0.0002443) (20,6.098e-06) (21,0.000119) (22,2.118e-06) (23,7.256e-05) (24,1.36e-06) (25,2.038e-05) (26,1.76e-07) (27,1.407e-05) (28,2.05e-07) (29,3.757e-06) (30,9e-09)};\draw[densely dotted] (axis cs:19,0) -- (axis cs:19,4.184);\end{axis}\begin{axis}[width=72mm,height=42mm,axis y line*=right,axis x line=none,ymode=log,xmin=-1,xmax=31,ylabel={\scriptsize eindeutig (log)},tick label style={font=\tiny},label style={font=\scriptsize},ylabel shift=-3pt]\addplot[mark=*,mark size=0.7pt,draw=black,thin] coordinates {(0,11732) (1,644638) (2,2481807) (3,20319177) (4,46658031) (5,95223144) (6,54260030) (7,64807641) (8,26954266) (9,35882072) (10,12425709) (11,10330802) (12,1688682) (13,747784) (14,153163) (15,55157) (16,3440) (17,1645) (18,18) (19,11)};\end{axis}\end{tikzpicture}\end{minipage}
\vspace{3mm}

\noindent\begin{minipage}[t]{60mm}\vspace{0pt}\sbox\hmboard{\chessboard[setfen=8/8/6p1/8/6P1/k7/1rb5/K7 w - - 0 1]}\usebox\hmboard\\[1mm]
\makebox[\wd\hmboard]{\sffamily\hspace{3mm}\textbf{h\#9.5}\hfill(2+4)}\end{minipage}\hfill
\begin{minipage}[t]{52mm}\vspace{0pt}
\textcolor{ink}{\sffamily\textsc{Lösung}}\par\smallskip\small
\begin{tabular}{@{}r@{\enspace}l@{\enspace}l@{}}
\textbf{1.} & \ldots & g5\\
\textbf{2.} & \symking{}b4 & \symking{}xb2\\
\textbf{3.} & \symking{}c5 & \symking{}c3\\
\textbf{4.} & \symking{}d6 & \symking{}d4\\
\textbf{5.} & \symking{}e7 & \symking{}e5\\
\textbf{6.} & \symking{}f8 & \symking{}f6\\
\textbf{7.} & \symbishop{}b3 & \symking{}xg6\\
\textbf{8.} & \symking{}g8 & \symking{}h6\\
\textbf{9.} & \symking{}h8 & g6\\
\textbf{10.} & \symbishop{}g8 & g7\#\\
\end{tabular}\par\vspace{3mm}
\qrcode[height=12mm]{https://helpman.komtera.lt/?fen=8/8/6p1/8/6P1/k7/1rb5/K7&moves=9.5}\\{\sffamily\tiny\color{ink}scannen zum Lösen}
\end{minipage}
\vspace{2mm}

{\raggedright\noindent\textcolor{ink}{\sffamily\textsc{Themen}}\enspace \themetag{reines Matt} \themetag{Mustermatt} \themetag{Idealmatt} \themetag{Selbstblock} \themetag{ohne Schach}\par}

\clearpage
\phantomsection\label{hm:KPvkrpp}
\begin{tcolorbox}[enhanced,colback=ink,colframe=ink,arc=0pt,boxrule=0pt,left=4mm,right=4mm,top=2.5mm,bottom=2.5mm]
\color{white}\sffamily{\LARGE\bfseries KPvkrpp}\hfill{\Large h\#9.5}\\[0.5mm]
{\small 6 Steine\quad Weiß am Zug}\hfill{\small Nr.~305\qquad\href{https://helpman.komtera.lt/?fen=6k1/8/1p6/8/8/p7/P1r5/K7&moves=9.5}{\color{white}\bfseries lösen\,\ensuremath{\nearrow}}}
\end{tcolorbox}
\vspace{2mm}
\noindent\begin{minipage}[t]{36mm}\vspace{0pt}\sffamily\small\raggedright
\hmfigure{längste eindeutige}{{\Large h\#9.5}}\hmfigure{eindeutig bei dieser Länge}{{\Large 105} Stellungen}\hmfigure{längstes Matt}{{\Large h\#15.5}}\hmfigure{in der Klasse}{25,6 Milliarden Stellungen,\\ 8,7 Milliarden mit Hilfsmatt}
\end{minipage}\hfill
\begin{minipage}[t]{80mm}\vspace{0pt}\begin{tikzpicture}\begin{axis}[width=72mm,height=42mm,ybar,bar width=1.6pt,xmin=-1,xmax=32,ymin=0,ymax=1.482,axis y line*=left,axis x line*=bottom,xlabel={\scriptsize Halbzüge bis Matt},ylabel={\scriptsize Stellungen (Mrd.)},tick label style={font=\tiny},label style={font=\scriptsize},enlarge x limits=false,ylabel shift=-3pt,xlabel shift=-3pt]\addplot[fill=ink!55,draw=none] coordinates {(0,3.91e-06) (1,0.001278) (2,0.004509) (3,0.03146) (4,0.09174) (5,0.3714) (6,0.6111) (7,1.098) (8,0.9269) (9,1.278) (10,0.9534) (11,1.347) (12,0.9951) (13,0.6334) (14,0.2085) (15,0.09934) (16,0.03214) (17,0.01779) (18,0.001504) (19,0.000582) (20,6.247e-06) (21,8.604e-05) (22,7.3e-07) (23,3.774e-05) (24,3.91e-07) (25,1.136e-05) (26,8.5e-08) (27,3.255e-06) (28,4.1e-08) (29,1.784e-06) (30,1.4e-08) (31,5.4e-07)};\draw[densely dotted] (axis cs:19,0) -- (axis cs:19,1.482);\end{axis}\begin{axis}[width=72mm,height=42mm,axis y line*=right,axis x line=none,ymode=log,xmin=-1,xmax=32,ylabel={\scriptsize eindeutig (log)},tick label style={font=\tiny},label style={font=\scriptsize},ylabel shift=-3pt]\addplot[mark=*,mark size=0.7pt,draw=black,thin] coordinates {(0,3910) (1,381095) (2,1334982) (3,9338572) (4,20993072) (5,38250769) (6,18984656) (7,20070790) (8,8312936) (9,10491758) (10,5269592) (11,4915059) (12,1607336) (13,615334) (14,212066) (15,74275) (16,9563) (17,3236) (18,121) (19,105)};\end{axis}\end{tikzpicture}\end{minipage}
\vspace{3mm}

\noindent\begin{minipage}[t]{60mm}\vspace{0pt}\sbox\hmboard{\chessboard[setfen=6k1/8/1p6/8/8/p7/P1r5/K7 w - - 0 1]}\usebox\hmboard\\[1mm]
\makebox[\wd\hmboard]{\sffamily\hspace{3mm}\textbf{h\#9.5}\hfill(2+4)}\end{minipage}\hfill
\begin{minipage}[t]{52mm}\vspace{0pt}
\textcolor{ink}{\sffamily\textsc{Lösung}}\par\smallskip\small
\begin{tabular}{@{}r@{\enspace}l@{\enspace}l@{}}
\textbf{1.} & \ldots & \symking{}b1\\
\textbf{2.} & b5 & \symking{}xc2\\
\textbf{3.} & b4 & \symking{}c1\\
\textbf{4.} & b3 & axb3\\
\textbf{5.} & \symking{}f7 & b4\\
\textbf{6.} & \symking{}e6 & b5\\
\textbf{7.} & \symking{}d5 & b6\\
\textbf{8.} & \symking{}c4 & b7\\
\textbf{9.} & \symking{}b3 & b8=\symqueen{}+\\
\textbf{10.} & \symking{}a2 & \symqueen{}b1\#\\
\end{tabular}\par\vspace{3mm}
\qrcode[height=12mm]{https://helpman.komtera.lt/?fen=6k1/8/1p6/8/8/p7/P1r5/K7&moves=9.5}\\{\sffamily\tiny\color{ink}scannen zum Lösen}
\end{minipage}
\vspace{2mm}

{\raggedright\noindent\textcolor{ink}{\sffamily\textsc{Themen}}\enspace \themetag{Umwandlung} \themetag{Excelsior}\par}

\clearpage
\phantomsection\label{hm:KBBBvkn}
\begin{tcolorbox}[enhanced,colback=ink,colframe=ink,arc=0pt,boxrule=0pt,left=4mm,right=4mm,top=2.5mm,bottom=2.5mm]
\color{white}\sffamily{\LARGE\bfseries KBBBvkn}\hfill{\Large h\#9}\\[0.5mm]
{\small 6 Steine\quad Schwarz am Zug}\hfill{\small Nr.~306\qquad\href{https://helpman.komtera.lt/?fen=8/8/5k2/8/3B4/8/1B6/BK5n&moves=9}{\color{white}\bfseries lösen\,\ensuremath{\nearrow}}}
\end{tcolorbox}
\vspace{2mm}
\noindent\begin{minipage}[t]{36mm}\vspace{0pt}\sffamily\small\raggedright
\hmfigure{längste eindeutige}{{\Large h\#9}}\hmfigure{eindeutig bei dieser Länge}{{\Large 22} Stellungen}\hmfigure{längstes Matt}{{\Large h\#12}}\hmfigure{in der Klasse}{15,5 Milliarden Stellungen,\\ 1,58 Milliarden mit Hilfsmatt}
\end{minipage}\hfill
\begin{minipage}[t]{80mm}\vspace{0pt}\begin{tikzpicture}\begin{axis}[width=72mm,height=42mm,ybar,bar width=1.6pt,xmin=-1,xmax=25,ymin=0,ymax=445.2,axis y line*=left,axis x line*=bottom,xlabel={\scriptsize Halbzüge bis Matt},ylabel={\scriptsize Stellungen (Mio.)},tick label style={font=\tiny},label style={font=\scriptsize},enlarge x limits=false,ylabel shift=-3pt,xlabel shift=-3pt]\addplot[fill=ink!55,draw=none] coordinates {(0,0.8829) (1,2.997) (2,11.84) (3,54.49) (4,183.7) (5,240.3) (6,404.7) (7,168.5) (8,102.6) (9,24.54) (10,18.46) (11,18.85) (12,38.34) (13,37.34) (14,66.98) (15,45.79) (16,59.45) (17,32.33) (18,26.8) (19,19.03) (20,9.405) (21,7.646) (22,1.26) (23,1.257) (24,0.02128)};\draw[densely dotted] (axis cs:18,0) -- (axis cs:18,445.2);\end{axis}\begin{axis}[width=72mm,height=42mm,axis y line*=right,axis x line=none,ymode=log,xmin=-1,xmax=25,ylabel={\scriptsize eindeutig (log)},tick label style={font=\tiny},label style={font=\scriptsize},ylabel shift=-3pt]\addplot[mark=*,mark size=0.7pt,draw=black,thin] coordinates {(0,882933) (1,2689871) (2,8018980) (3,21099673) (4,36900595) (5,6815315) (6,3612354) (7,384684) (8,118514) (9,59745) (10,95707) (11,43941) (12,68537) (13,20298) (14,14395) (15,3120) (16,2117) (17,59) (18,22)};\end{axis}\end{tikzpicture}\end{minipage}
\vspace{3mm}

\noindent\begin{minipage}[t]{60mm}\vspace{0pt}\sbox\hmboard{\chessboard[setfen=8/8/5k2/8/3B4/8/1B6/BK5n b - - 0 1]}\usebox\hmboard\\[1mm]
\makebox[\wd\hmboard]{\sffamily\hspace{3mm}\textbf{h\#9}\hfill(4+2)}\end{minipage}\hfill
\begin{minipage}[t]{52mm}\vspace{0pt}
\textcolor{ink}{\sffamily\textsc{Lösung}}\par\smallskip\small
\begin{tabular}{@{}r@{\enspace}l@{\enspace}l@{}}
\textbf{1.} & \symking{}f7 & \symking{}c2\\
\textbf{2.} & \symknight{}g3 & \symking{}d3\\
\textbf{3.} & \symknight{}e2 & \symking{}e4\\
\textbf{4.} & \symking{}g8 & \symking{}f5\\
\textbf{5.} & \symknight{}xd4+ & \symking{}g6\\
\textbf{6.} & \symknight{}f5 & \symbishop{}a3\\
\textbf{7.} & \symknight{}h6 & \symbishop{}h8\\
\textbf{8.} & \symking{}xh8 & \symbishop{}f8\\
\textbf{9.} & \symknight{}g8 & \symbishop{}g7\#\\
\end{tabular}\par\vspace{3mm}
\qrcode[height=12mm]{https://helpman.komtera.lt/?fen=8/8/5k2/8/3B4/8/1B6/BK5n&moves=9}\\{\sffamily\tiny\color{ink}scannen zum Lösen}
\end{minipage}
\vspace{2mm}

{\raggedright\noindent\textcolor{ink}{\sffamily\textsc{Themen}}\enspace \themetag{reines Matt} \themetag{Mustermatt} \themetag{Idealmatt} \themetag{Selbstblock} \themetag{Kniest} \themetag{Umnov}\par\smallskip{\footnotesize\textit{Hinweis:} die Partei am Zug steht im Diagramm im Schach.}\par}

\clearpage
\phantomsection\label{hm:KBBvkbn}
\begin{tcolorbox}[enhanced,colback=ink,colframe=ink,arc=0pt,boxrule=0pt,left=4mm,right=4mm,top=2.5mm,bottom=2.5mm]
\color{white}\sffamily{\LARGE\bfseries KBBvkbn}\hfill{\Large h\#9}\\[0.5mm]
{\small 6 Steine\quad Schwarz am Zug}\hfill{\small Nr.~307\qquad\href{https://helpman.komtera.lt/?fen=8/8/1n6/3b4/8/8/2B5/KBk5&moves=9}{\color{white}\bfseries lösen\,\ensuremath{\nearrow}}}
\end{tcolorbox}
\vspace{2mm}
\noindent\begin{minipage}[t]{36mm}\vspace{0pt}\sffamily\small\raggedright
\hmfigure{längste eindeutige}{{\Large h\#9}}\hmfigure{eindeutig bei dieser Länge}{{\Large 244} Stellungen}\hmfigure{längstes Matt}{{\Large h\#12}}\hmfigure{in der Klasse}{15,5 Milliarden Stellungen,\\ 4,7 Milliarden mit Hilfsmatt}
\end{minipage}\hfill
\begin{minipage}[t]{80mm}\vspace{0pt}\begin{tikzpicture}\begin{axis}[width=72mm,height=42mm,ybar,bar width=1.6pt,xmin=-1,xmax=25,ymin=0,ymax=698.7,axis y line*=left,axis x line*=bottom,xlabel={\scriptsize Halbzüge bis Matt},ylabel={\scriptsize Stellungen (Mio.)},tick label style={font=\tiny},label style={font=\scriptsize},enlarge x limits=false,ylabel shift=-3pt,xlabel shift=-3pt]\addplot[fill=ink!55,draw=none] coordinates {(0,0.3993) (1,1.624) (2,8.726) (3,50.74) (4,184.2) (5,382.1) (6,635.2) (7,563.3) (8,498.7) (9,360) (10,323.7) (11,325.6) (12,323.9) (13,255.7) (14,227.1) (15,163.8) (16,138.4) (17,102.4) (18,61.7) (19,55.95) (20,19.32) (21,19.3) (22,1.241) (23,1.419) (24,0.001486)};\draw[densely dotted] (axis cs:18,0) -- (axis cs:18,698.7);\end{axis}\begin{axis}[width=72mm,height=42mm,axis y line*=right,axis x line=none,ymode=log,xmin=-1,xmax=25,ylabel={\scriptsize eindeutig (log)},tick label style={font=\tiny},label style={font=\scriptsize},ylabel shift=-3pt]\addplot[mark=*,mark size=0.7pt,draw=black,thin] coordinates {(0,399312) (1,1588991) (2,6248108) (3,24658438) (4,40331058) (5,23462757) (6,12816707) (7,4413503) (8,2101794) (9,1320613) (10,987366) (11,512372) (12,305329) (13,127495) (14,60373) (15,16780) (16,5923) (17,629) (18,244)};\end{axis}\end{tikzpicture}\end{minipage}
\vspace{3mm}

\noindent\begin{minipage}[t]{60mm}\vspace{0pt}\sbox\hmboard{\chessboard[setfen=8/8/1n6/3b4/8/8/2B5/KBk5 b - - 0 1]}\usebox\hmboard\\[1mm]
\makebox[\wd\hmboard]{\sffamily\hspace{3mm}\textbf{h\#9}\hfill(3+3)}\end{minipage}\hfill
\begin{minipage}[t]{52mm}\vspace{0pt}
\textcolor{ink}{\sffamily\textsc{Lösung}}\par\smallskip\small
\begin{tabular}{@{}r@{\enspace}l@{\enspace}l@{}}
\textbf{1.} & \symking{}d2 & \symbishop{}d3\\
\textbf{2.} & \symking{}c3 & \symbishop{}a6\\
\textbf{3.} & \symking{}b4 & \symking{}b2\\
\textbf{4.} & \symking{}a5 & \symking{}c3\\
\textbf{5.} & \symknight{}c8 & \symking{}d4\\
\textbf{6.} & \symking{}b6 & \symking{}xd5\\
\textbf{7.} & \symking{}a7 & \symking{}c6\\
\textbf{8.} & \symking{}a8 & \symking{}c7\\
\textbf{9.} & \symknight{}a7 & \symbishop{}b7\#\\
\end{tabular}\par\vspace{3mm}
\qrcode[height=12mm]{https://helpman.komtera.lt/?fen=8/8/1n6/3b4/8/8/2B5/KBk5&moves=9}\\{\sffamily\tiny\color{ink}scannen zum Lösen}
\end{minipage}
\vspace{2mm}

{\raggedright\noindent\textcolor{ink}{\sffamily\textsc{Themen}}\enspace \themetag{reines Matt} \themetag{Selbstblock} \themetag{ohne Schach}\par}

\clearpage
\phantomsection\label{hm:KBBvkrn}
\begin{tcolorbox}[enhanced,colback=ink,colframe=ink,arc=0pt,boxrule=0pt,left=4mm,right=4mm,top=2.5mm,bottom=2.5mm]
\color{white}\sffamily{\LARGE\bfseries KBBvkrn}\hfill{\Large h\#9}\\[0.5mm]
{\small 6 Steine\quad Schwarz am Zug}\hfill{\small Nr.~308\qquad\href{https://helpman.komtera.lt/?fen=8/8/1n6/8/8/1B6/3B4/K1k3r1&moves=9}{\color{white}\bfseries lösen\,\ensuremath{\nearrow}}}
\end{tcolorbox}
\vspace{2mm}
\noindent\begin{minipage}[t]{36mm}\vspace{0pt}\sffamily\small\raggedright
\hmfigure{längste eindeutige}{{\Large h\#9}}\hmfigure{eindeutig bei dieser Länge}{{\Large 8} Stellungen}\hmfigure{längstes Matt}{{\Large h\#13}}\hmfigure{in der Klasse}{15,5 Milliarden Stellungen,\\ 4,51 Milliarden mit Hilfsmatt}
\end{minipage}\hfill
\begin{minipage}[t]{80mm}\vspace{0pt}\begin{tikzpicture}\begin{axis}[width=72mm,height=42mm,ybar,bar width=1.6pt,xmin=-1,xmax=27,ymin=0,ymax=622.4,axis y line*=left,axis x line*=bottom,xlabel={\scriptsize Halbzüge bis Matt},ylabel={\scriptsize Stellungen (Mio.)},tick label style={font=\tiny},label style={font=\scriptsize},enlarge x limits=false,ylabel shift=-3pt,xlabel shift=-3pt]\addplot[fill=ink!55,draw=none] coordinates {(0,0.1877) (1,0.8209) (2,3.865) (3,27.46) (4,92.56) (5,277.8) (6,514.1) (7,565.8) (8,542.1) (9,437.9) (10,390.7) (11,439.1) (12,453) (13,427.7) (14,220.3) (15,102.5) (16,10.92) (17,1.589) (18,0.07508) (19,0.1281) (20,0.06968) (21,0.02617) (22,0.03177) (23,0.003393) (24,0.00488) (25,1e-06) (26,2.7e-05)};\draw[densely dotted] (axis cs:18,0) -- (axis cs:18,622.4);\end{axis}\begin{axis}[width=72mm,height=42mm,axis y line*=right,axis x line=none,ymode=log,xmin=-1,xmax=27,ylabel={\scriptsize eindeutig (log)},tick label style={font=\tiny},label style={font=\scriptsize},ylabel shift=-3pt]\addplot[mark=*,mark size=0.7pt,draw=black,thin] coordinates {(0,187667) (1,817679) (2,2290786) (3,11584656) (4,21780133) (5,24049529) (6,17579525) (7,6850632) (8,3765618) (9,2392734) (10,1529500) (11,1090072) (12,573662) (13,289433) (14,30709) (15,3120) (16,28) (17,108) (18,8)};\end{axis}\end{tikzpicture}\end{minipage}
\vspace{3mm}

\noindent\begin{minipage}[t]{60mm}\vspace{0pt}\sbox\hmboard{\chessboard[setfen=8/8/1n6/8/8/1B6/3B4/K1k3r1 b - - 0 1]}\usebox\hmboard\\[1mm]
\makebox[\wd\hmboard]{\sffamily\hspace{3mm}\textbf{h\#9}\hfill(3+3)}\end{minipage}\hfill
\begin{minipage}[t]{52mm}\vspace{0pt}
\textcolor{ink}{\sffamily\textsc{Lösung}}\par\smallskip\small
\begin{tabular}{@{}r@{\enspace}l@{\enspace}l@{}}
\textbf{1.} & \symking{}xd2+ & \symking{}a2\\
\textbf{2.} & \symking{}c3 & \symbishop{}c4\\
\textbf{3.} & \symking{}b4 & \symbishop{}a6\\
\textbf{4.} & \symking{}a5 & \symking{}b3\\
\textbf{5.} & \symknight{}c8 & \symking{}c4\\
\textbf{6.} & \symking{}b6 & \symking{}d5\\
\textbf{7.} & \symking{}a7 & \symking{}c6\\
\textbf{8.} & \symking{}a8 & \symking{}c7\\
\textbf{9.} & \symknight{}a7 & \symbishop{}b7\#\\
\end{tabular}\par\vspace{3mm}
\qrcode[height=12mm]{https://helpman.komtera.lt/?fen=8/8/1n6/8/8/1B6/3B4/K1k3r1&moves=9}\\{\sffamily\tiny\color{ink}scannen zum Lösen}
\end{minipage}
\vspace{2mm}

{\raggedright\noindent\textcolor{ink}{\sffamily\textsc{Themen}}\enspace \themetag{reines Matt} \themetag{Mustermatt} \themetag{Selbstblock}\par\smallskip{\footnotesize\textit{Hinweis:} die Partei am Zug steht im Diagramm im Schach; die Lösung beginnt mit einem Schlag.}\par}

\clearpage
\phantomsection\label{hm:KBNvkqp}
\begin{tcolorbox}[enhanced,colback=ink,colframe=ink,arc=0pt,boxrule=0pt,left=4mm,right=4mm,top=2.5mm,bottom=2.5mm]
\color{white}\sffamily{\LARGE\bfseries KBNvkqp}\hfill{\Large h\#9}\\[0.5mm]
{\small 6 Steine\quad Schwarz am Zug}\hfill{\small Nr.~309\qquad\href{https://helpman.komtera.lt/?fen=4k3/4Bp2/1q3N2/8/8/8/8/K7&moves=9}{\color{white}\bfseries lösen\,\ensuremath{\nearrow}}}
\end{tcolorbox}
\vspace{2mm}
\noindent\begin{minipage}[t]{36mm}\vspace{0pt}\sffamily\small\raggedright
\hmfigure{längste eindeutige}{{\Large h\#9}}\hmfigure{eindeutig bei dieser Länge}{{\Large 16} Stellungen}\hmfigure{längstes Matt}{{\Large h\#17}}\hmfigure{in der Klasse}{45,4 Milliarden Stellungen,\\ 26,5 Milliarden mit Hilfsmatt}
\end{minipage}\hfill
\begin{minipage}[t]{80mm}\vspace{0pt}\begin{tikzpicture}\begin{axis}[width=72mm,height=42mm,ybar,bar width=1.6pt,xmin=-1,xmax=35,ymin=0,ymax=7.333,axis y line*=left,axis x line*=bottom,xlabel={\scriptsize Halbzüge bis Matt},ylabel={\scriptsize Stellungen (Mrd.)},tick label style={font=\tiny},label style={font=\scriptsize},enlarge x limits=false,ylabel shift=-3pt,xlabel shift=-3pt]\addplot[fill=ink!55,draw=none] coordinates {(0,0.0002687) (1,0.00113) (2,0.004913) (3,0.03648) (4,0.1245) (5,0.6258) (6,1.423) (7,3.807) (8,5.023) (9,6.666) (10,4.4) (11,3.095) (12,0.9152) (13,0.3108) (14,0.03119) (15,0.00517) (16,0.0002133) (17,3.158e-06) (18,0.0001187) (19,1.328e-06) (20,4.726e-05) (21,9.43e-07) (22,1.708e-05) (23,9.07e-07) (24,1.333e-05) (25,3.5e-07) (26,6.547e-06) (27,2.79e-07) (28,8.002e-06) (29,2.08e-07) (30,4.548e-06) (31,4.9e-08) (32,3.784e-06) (33,9e-09) (34,3.6e-08)};\draw[densely dotted] (axis cs:18,0) -- (axis cs:18,7.333);\end{axis}\begin{axis}[width=72mm,height=42mm,axis y line*=right,axis x line=none,ymode=log,xmin=-1,xmax=35,ylabel={\scriptsize eindeutig (log)},tick label style={font=\tiny},label style={font=\scriptsize},ylabel shift=-3pt]\addplot[mark=*,mark size=0.7pt,draw=black,thin] coordinates {(0,268698) (1,1078714) (2,2808113) (3,14502203) (4,32713976) (5,89151656) (6,90468547) (7,105877329) (8,47657062) (9,26666311) (10,5598382) (11,1666352) (12,156222) (13,26261) (14,503) (15,20) (16,35) (18,16)};\end{axis}\end{tikzpicture}\end{minipage}
\vspace{3mm}

\noindent\begin{minipage}[t]{60mm}\vspace{0pt}\sbox\hmboard{\chessboard[setfen=4k3/4Bp2/1q3N2/8/8/8/8/K7 b - - 0 1]}\usebox\hmboard\\[1mm]
\makebox[\wd\hmboard]{\sffamily\hspace{3mm}\textbf{h\#9}\hfill(3+3)}\end{minipage}\hfill
\begin{minipage}[t]{52mm}\vspace{0pt}
\textcolor{ink}{\sffamily\textsc{Lösung}}\par\smallskip\small
\begin{tabular}{@{}r@{\enspace}l@{\enspace}l@{}}
\textbf{1.} & \symking{}xe7 & \symknight{}d7\\
\textbf{2.} & \symqueen{}d8 & \symking{}b2\\
\textbf{3.} & f5 & \symking{}c3\\
\textbf{4.} & f4 & \symking{}d4\\
\textbf{5.} & f3 & \symking{}e5\\
\textbf{6.} & f2 & \symknight{}f6\\
\textbf{7.} & f1=\symrook{} & \symknight{}h5\\
\textbf{8.} & \symking{}e8 & \symking{}e6\\
\textbf{9.} & \symrook{}f8 & \symknight{}g7\#\\
\end{tabular}\par\vspace{3mm}
\qrcode[height=12mm]{https://helpman.komtera.lt/?fen=4k3/4Bp2/1q3N2/8/8/8/8/K7&moves=9}\\{\sffamily\tiny\color{ink}scannen zum Lösen}
\end{minipage}
\vspace{2mm}

{\raggedright\noindent\textcolor{ink}{\sffamily\textsc{Themen}}\enspace \themetag{reines Matt} \themetag{Mustermatt} \themetag{Idealmatt} \themetag{Umwandlung} \themetag{Unterverwandlung} \themetag{Excelsior} \themetag{Switchback} \themetag{Selbstblock} \themetag{ohne Schach}\par\smallskip{\footnotesize\textit{Hinweis:} die Partei am Zug steht im Diagramm im Schach; die Lösung beginnt mit einem Schlag.}\par}

\clearpage
\phantomsection\label{hm:KBPvkbn}
\begin{tcolorbox}[enhanced,colback=ink,colframe=ink,arc=0pt,boxrule=0pt,left=4mm,right=4mm,top=2.5mm,bottom=2.5mm]
\color{white}\sffamily{\LARGE\bfseries KBPvkbn}\hfill{\Large h\#9}\\[0.5mm]
{\small 6 Steine\quad Schwarz am Zug}\hfill{\small Nr.~310\qquad\href{https://helpman.komtera.lt/?fen=K7/5n2/kB6/1P2b3/8/8/8/8&moves=9}{\color{white}\bfseries lösen\,\ensuremath{\nearrow}}}
\end{tcolorbox}
\vspace{2mm}
\noindent\begin{minipage}[t]{36mm}\vspace{0pt}\sffamily\small\raggedright
\hmfigure{längste eindeutige}{{\Large h\#9}}\hmfigure{eindeutig bei dieser Länge}{{\Large 3} Stellungen}\hmfigure{längstes Matt}{{\Large h\#12}}\hmfigure{in der Klasse}{45,4 Milliarden Stellungen,\\ 29,9 Milliarden mit Hilfsmatt}
\end{minipage}\hfill
\begin{minipage}[t]{80mm}\vspace{0pt}\begin{tikzpicture}\begin{axis}[width=72mm,height=42mm,ybar,bar width=1.6pt,xmin=-1,xmax=25,ymin=0,ymax=4.612,axis y line*=left,axis x line*=bottom,xlabel={\scriptsize Halbzüge bis Matt},ylabel={\scriptsize Stellungen (Mrd.)},tick label style={font=\tiny},label style={font=\scriptsize},enlarge x limits=false,ylabel shift=-3pt,xlabel shift=-3pt]\addplot[fill=ink!55,draw=none] coordinates {(0,0.0003659) (1,0.005534) (2,0.022) (3,0.2053) (4,0.6515) (5,2.379) (6,3.262) (7,3.949) (8,3.851) (9,4.193) (10,3.904) (11,3.853) (12,2.731) (13,0.7954) (14,0.1446) (15,0.00384) (16,0.0001308) (17,1.2e-08) (18,7.679e-05) (20,6.201e-05) (22,3.349e-06) (24,2.23e-06)};\draw[densely dotted] (axis cs:18,0) -- (axis cs:18,4.612);\end{axis}\begin{axis}[width=72mm,height=42mm,axis y line*=right,axis x line=none,ymode=log,xmin=-1,xmax=25,ylabel={\scriptsize eindeutig (log)},tick label style={font=\tiny},label style={font=\scriptsize},ylabel shift=-3pt]\addplot[mark=*,mark size=0.7pt,draw=black,thin] coordinates {(0,365908) (1,2667986) (2,9679064) (3,76015807) (4,143398494) (5,189004781) (6,116408882) (7,94225437) (8,52699142) (9,43036056) (10,20160000) (11,9928029) (12,1894645) (13,50154) (14,229) (15,1) (16,14) (18,3)};\end{axis}\end{tikzpicture}\end{minipage}
\vspace{3mm}

\noindent\begin{minipage}[t]{60mm}\vspace{0pt}\sbox\hmboard{\chessboard[setfen=K7/5n2/kB6/1P2b3/8/8/8/8 b - - 0 1]}\usebox\hmboard\\[1mm]
\makebox[\wd\hmboard]{\sffamily\hspace{3mm}\textbf{h\#9}\hfill(3+3)}\end{minipage}\hfill
\begin{minipage}[t]{52mm}\vspace{0pt}
\textcolor{ink}{\sffamily\textsc{Lösung}}\par\smallskip\small
\begin{tabular}{@{}r@{\enspace}l@{\enspace}l@{}}
\textbf{1.} & \symking{}xb5 & \symbishop{}c5\\
\textbf{2.} & \symking{}c6 & \symbishop{}f8\\
\textbf{3.} & \symking{}d7 & \symking{}b7\\
\textbf{4.} & \symking{}e8 & \symking{}c6\\
\textbf{5.} & \symknight{}h6 & \symking{}d5\\
\textbf{6.} & \symking{}f7 & \symking{}xe5\\
\textbf{7.} & \symking{}g8 & \symking{}f6\\
\textbf{8.} & \symking{}h8 & \symking{}g6\\
\textbf{9.} & \symknight{}g8 & \symbishop{}g7\#\\
\end{tabular}\par\vspace{3mm}
\qrcode[height=12mm]{https://helpman.komtera.lt/?fen=K7/5n2/kB6/1P2b3/8/8/8/8&moves=9}\\{\sffamily\tiny\color{ink}scannen zum Lösen}
\end{minipage}
\vspace{2mm}

{\raggedright\noindent\textcolor{ink}{\sffamily\textsc{Themen}}\enspace \themetag{reines Matt} \themetag{Mustermatt} \themetag{Idealmatt} \themetag{Selbstblock} \themetag{ohne Schach}\par\smallskip{\footnotesize\textit{Hinweis:} die Partei am Zug steht im Diagramm im Schach; die Lösung beginnt mit einem Schlag.}\par}

\clearpage
\phantomsection\label{hm:KBPvkqn}
\begin{tcolorbox}[enhanced,colback=ink,colframe=ink,arc=0pt,boxrule=0pt,left=4mm,right=4mm,top=2.5mm,bottom=2.5mm]
\color{white}\sffamily{\LARGE\bfseries KBPvkqn}\hfill{\Large h\#9}\\[0.5mm]
{\small 6 Steine\quad Schwarz am Zug}\hfill{\small Nr.~311\qquad\href{https://helpman.komtera.lt/?fen=2k5/K2P4/8/B7/8/1n6/1q6/8&moves=9}{\color{white}\bfseries lösen\,\ensuremath{\nearrow}}}
\end{tcolorbox}
\vspace{2mm}
\noindent\begin{minipage}[t]{36mm}\vspace{0pt}\sffamily\small\raggedright
\hmfigure{längste eindeutige}{{\Large h\#9}}\hmfigure{eindeutig bei dieser Länge}{{\Large 20} Stellungen}\hmfigure{längstes Matt}{{\Large h\#13}}\hmfigure{in der Klasse}{45,4 Milliarden Stellungen,\\ 26,7 Milliarden mit Hilfsmatt}
\end{minipage}\hfill
\begin{minipage}[t]{80mm}\vspace{0pt}\begin{tikzpicture}\begin{axis}[width=72mm,height=42mm,ybar,bar width=1.6pt,xmin=-1,xmax=27,ymin=0,ymax=4.55,axis y line*=left,axis x line*=bottom,xlabel={\scriptsize Halbzüge bis Matt},ylabel={\scriptsize Stellungen (Mrd.)},tick label style={font=\tiny},label style={font=\scriptsize},enlarge x limits=false,ylabel shift=-3pt,xlabel shift=-3pt]\addplot[fill=ink!55,draw=none] coordinates {(0,0.0001217) (1,0.00266) (2,0.01372) (3,0.1553) (4,0.5119) (5,2.051) (6,2.475) (7,3.736) (8,2.796) (9,3.972) (10,3.04) (11,4.137) (12,2.363) (13,1.275) (14,0.1067) (15,0.01962) (16,0.0001116) (17,2.767e-06) (18,8.108e-05) (19,3.1e-08) (20,7.815e-05) (21,2e-09) (22,2.305e-05) (24,1.243e-05) (26,8.88e-07)};\draw[densely dotted] (axis cs:18,0) -- (axis cs:18,4.55);\end{axis}\begin{axis}[width=72mm,height=42mm,axis y line*=right,axis x line=none,ymode=log,xmin=-1,xmax=27,ylabel={\scriptsize eindeutig (log)},tick label style={font=\tiny},label style={font=\scriptsize},ylabel shift=-3pt]\addplot[mark=*,mark size=0.7pt,draw=black,thin] coordinates {(0,121650) (1,1275932) (2,5466332) (3,59027961) (4,106546068) (5,164738075) (6,72456504) (7,80388175) (8,33892640) (9,39613392) (10,17514316) (11,11707744) (12,1328954) (13,274627) (14,522) (15,18) (16,32) (18,20)};\end{axis}\end{tikzpicture}\end{minipage}
\vspace{3mm}

\noindent\begin{minipage}[t]{60mm}\vspace{0pt}\sbox\hmboard{\chessboard[setfen=2k5/K2P4/8/B7/8/1n6/1q6/8 b - - 0 1]}\usebox\hmboard\\[1mm]
\makebox[\wd\hmboard]{\sffamily\hspace{3mm}\textbf{h\#9}\hfill(3+3)}\end{minipage}\hfill
\begin{minipage}[t]{52mm}\vspace{0pt}
\textcolor{ink}{\sffamily\textsc{Lösung}}\par\smallskip\small
\begin{tabular}{@{}r@{\enspace}l@{\enspace}l@{}}
\textbf{1.} & \symking{}xd7 & \symbishop{}c3\\
\textbf{2.} & \symking{}c6 & \symbishop{}xb2\\
\textbf{3.} & \symking{}b5 & \symbishop{}a3\\
\textbf{4.} & \symking{}a4 & \symking{}b6\\
\textbf{5.} & \symknight{}c1 & \symking{}c5\\
\textbf{6.} & \symking{}b3 & \symking{}d4\\
\textbf{7.} & \symking{}a2 & \symking{}c3\\
\textbf{8.} & \symking{}a1 & \symking{}c2\\
\textbf{9.} & \symknight{}a2 & \symbishop{}b2\#\\
\end{tabular}\par\vspace{3mm}
\qrcode[height=12mm]{https://helpman.komtera.lt/?fen=2k5/K2P4/8/B7/8/1n6/1q6/8&moves=9}\\{\sffamily\tiny\color{ink}scannen zum Lösen}
\end{minipage}
\vspace{2mm}

{\raggedright\noindent\textcolor{ink}{\sffamily\textsc{Themen}}\enspace \themetag{reines Matt} \themetag{Mustermatt} \themetag{Idealmatt} \themetag{Switchback} \themetag{Selbstblock} \themetag{ohne Schach}\par\smallskip{\footnotesize\textit{Hinweis:} die Partei am Zug steht im Diagramm im Schach; die Lösung beginnt mit einem Schlag.}\par}

\clearpage
\phantomsection\label{hm:KBPvkrn}
\begin{tcolorbox}[enhanced,colback=ink,colframe=ink,arc=0pt,boxrule=0pt,left=4mm,right=4mm,top=2.5mm,bottom=2.5mm]
\color{white}\sffamily{\LARGE\bfseries KBPvkrn}\hfill{\Large h\#9}\\[0.5mm]
{\small 6 Steine\quad Schwarz am Zug}\hfill{\small Nr.~312\qquad\href{https://helpman.komtera.lt/?fen=2k5/K2P4/8/B7/8/1r6/8/2n5&moves=9}{\color{white}\bfseries lösen\,\ensuremath{\nearrow}}}
\end{tcolorbox}
\vspace{2mm}
\noindent\begin{minipage}[t]{36mm}\vspace{0pt}\sffamily\small\raggedright
\hmfigure{längste eindeutige}{{\Large h\#9}}\hmfigure{eindeutig bei dieser Länge}{{\Large 16} Stellungen}\hmfigure{längstes Matt}{{\Large h\#13}}\hmfigure{in der Klasse}{45,4 Milliarden Stellungen,\\ 28,8 Milliarden mit Hilfsmatt}
\end{minipage}\hfill
\begin{minipage}[t]{80mm}\vspace{0pt}\begin{tikzpicture}\begin{axis}[width=72mm,height=42mm,ybar,bar width=1.6pt,xmin=-1,xmax=27,ymin=0,ymax=4.451,axis y line*=left,axis x line*=bottom,xlabel={\scriptsize Halbzüge bis Matt},ylabel={\scriptsize Stellungen (Mrd.)},tick label style={font=\tiny},label style={font=\scriptsize},enlarge x limits=false,ylabel shift=-3pt,xlabel shift=-3pt]\addplot[fill=ink!55,draw=none] coordinates {(0,0.0001722) (1,0.004368) (2,0.01982) (3,0.1986) (4,0.6975) (5,2.393) (6,3.072) (7,3.85) (8,3.314) (9,4.047) (10,3.63) (11,4.029) (12,2.584) (13,0.848) (14,0.07503) (15,0.003811) (16,0.0001485) (17,1e-07) (18,9.496e-05) (19,1.7e-08) (20,9.377e-05) (22,1.668e-05) (24,5.73e-06) (26,2.4e-08)};\draw[densely dotted] (axis cs:18,0) -- (axis cs:18,4.451);\end{axis}\begin{axis}[width=72mm,height=42mm,axis y line*=right,axis x line=none,ymode=log,xmin=-1,xmax=27,ylabel={\scriptsize eindeutig (log)},tick label style={font=\tiny},label style={font=\scriptsize},ylabel shift=-3pt]\addplot[mark=*,mark size=0.7pt,draw=black,thin] coordinates {(0,172197) (1,1999455) (2,8216244) (3,76091149) (4,161186699) (5,200498779) (6,103432012) (7,87823146) (8,42758331) (9,42323947) (10,22185118) (11,8689470) (12,897032) (13,43095) (14,546) (15,1) (16,36) (18,16)};\end{axis}\end{tikzpicture}\end{minipage}
\vspace{3mm}

\noindent\begin{minipage}[t]{60mm}\vspace{0pt}\sbox\hmboard{\chessboard[setfen=2k5/K2P4/8/B7/8/1r6/8/2n5 b - - 0 1]}\usebox\hmboard\\[1mm]
\makebox[\wd\hmboard]{\sffamily\hspace{3mm}\textbf{h\#9}\hfill(3+3)}\end{minipage}\hfill
\begin{minipage}[t]{52mm}\vspace{0pt}
\textcolor{ink}{\sffamily\textsc{Lösung}}\par\smallskip\small
\begin{tabular}{@{}r@{\enspace}l@{\enspace}l@{}}
\textbf{1.} & \symking{}xd7 & \symbishop{}b4\\
\textbf{2.} & \symking{}c6 & \symbishop{}a3\\
\textbf{3.} & \symking{}b5 & \symking{}b7\\
\textbf{4.} & \symking{}a4+ & \symking{}c6\\
\textbf{5.} & \symrook{}c3+ & \symking{}d5\\
\textbf{6.} & \symking{}b3 & \symking{}d4\\
\textbf{7.} & \symking{}a2 & \symking{}xc3\\
\textbf{8.} & \symking{}a1 & \symking{}c2\\
\textbf{9.} & \symknight{}a2 & \symbishop{}b2\#\\
\end{tabular}\par\vspace{3mm}
\qrcode[height=12mm]{https://helpman.komtera.lt/?fen=2k5/K2P4/8/B7/8/1r6/8/2n5&moves=9}\\{\sffamily\tiny\color{ink}scannen zum Lösen}
\end{minipage}
\vspace{2mm}

{\raggedright\noindent\textcolor{ink}{\sffamily\textsc{Themen}}\enspace \themetag{reines Matt} \themetag{Mustermatt} \themetag{Idealmatt} \themetag{Selbstblock}\par\smallskip{\footnotesize\textit{Hinweis:} die Partei am Zug steht im Diagramm im Schach; die Lösung beginnt mit einem Schlag.}\par}

\clearpage
\phantomsection\label{hm:KNNvkqb}
\begin{tcolorbox}[enhanced,colback=ink,colframe=ink,arc=0pt,boxrule=0pt,left=4mm,right=4mm,top=2.5mm,bottom=2.5mm]
\color{white}\sffamily{\LARGE\bfseries KNNvkqb}\hfill{\Large h\#9}\\[0.5mm]
{\small 6 Steine\quad Schwarz am Zug}\hfill{\small Nr.~313\qquad\href{https://helpman.komtera.lt/?fen=8/N7/8/8/8/1N4b1/q7/k1K5&moves=9}{\color{white}\bfseries lösen\,\ensuremath{\nearrow}}}
\end{tcolorbox}
\vspace{2mm}
\noindent\begin{minipage}[t]{36mm}\vspace{0pt}\sffamily\small\raggedright
\hmfigure{längste eindeutige}{{\Large h\#9}}\hmfigure{eindeutig bei dieser Länge}{{\Large 2} Stellungen}\hmfigure{längstes Matt}{{\Large h\#11}}\hmfigure{in der Klasse}{15,5 Milliarden Stellungen,\\ 4,3 Milliarden mit Hilfsmatt}
\end{minipage}\hfill
\begin{minipage}[t]{80mm}\vspace{0pt}\begin{tikzpicture}\begin{axis}[width=72mm,height=42mm,ybar,bar width=1.6pt,xmin=-1,xmax=23,ymin=0,ymax=1.42,axis y line*=left,axis x line*=bottom,xlabel={\scriptsize Halbzüge bis Matt},ylabel={\scriptsize Stellungen (Mrd.)},tick label style={font=\tiny},label style={font=\scriptsize},enlarge x limits=false,ylabel shift=-3pt,xlabel shift=-3pt]\addplot[fill=ink!55,draw=none] coordinates {(0,2.676e-05) (1,0.0001341) (2,0.0006001) (3,0.005345) (4,0.02148) (5,0.1137) (6,0.2923) (7,0.8163) (8,0.9941) (9,1.291) (10,0.4325) (11,0.2957) (12,0.02555) (13,0.01305) (14,0.0001567) (15,0.0002804) (16,6.163e-06) (17,1.66e-07) (18,7.924e-06) (19,1.5e-08) (20,5.599e-06) (21,3e-09) (22,2.05e-07)};\draw[densely dotted] (axis cs:18,0) -- (axis cs:18,1.42);\end{axis}\begin{axis}[width=72mm,height=42mm,axis y line*=right,axis x line=none,ymode=log,xmin=-1,xmax=23,ylabel={\scriptsize eindeutig (log)},tick label style={font=\tiny},label style={font=\scriptsize},ylabel shift=-3pt]\addplot[mark=*,mark size=0.7pt,draw=black,thin] coordinates {(0,26758) (1,130169) (2,390904) (3,2464399) (4,6864964) (5,14415233) (6,14758936) (7,15703635) (8,3914737) (9,1973911) (10,79233) (11,35487) (12,228) (13,194) (14,1) (16,4) (18,2)};\end{axis}\end{tikzpicture}\end{minipage}
\vspace{3mm}

\noindent\begin{minipage}[t]{60mm}\vspace{0pt}\sbox\hmboard{\chessboard[setfen=8/N7/8/8/8/1N4b1/q7/k1K5 b - - 0 1]}\usebox\hmboard\\[1mm]
\makebox[\wd\hmboard]{\sffamily\hspace{3mm}\textbf{h\#9}\hfill(3+3)}\end{minipage}\hfill
\begin{minipage}[t]{52mm}\vspace{0pt}
\textcolor{ink}{\sffamily\textsc{Lösung}}\par\smallskip\small
\begin{tabular}{@{}r@{\enspace}l@{\enspace}l@{}}
\textbf{1.} & \symqueen{}xb3 & \symking{}d2\\
\textbf{2.} & \symking{}b2 & \symking{}e2\\
\textbf{3.} & \symking{}c3 & \symknight{}c8\\
\textbf{4.} & \symking{}d4 & \symking{}d2\\
\textbf{5.} & \symking{}e4 & \symknight{}d6+\\
\textbf{6.} & \symking{}f3 & \symknight{}e4\\
\textbf{7.} & \symking{}g2 & \symking{}e2\\
\textbf{8.} & \symking{}h1 & \symking{}f1\\
\textbf{9.} & \symbishop{}h2 & \symknight{}f2\#\\
\end{tabular}\par\vspace{3mm}
\qrcode[height=12mm]{https://helpman.komtera.lt/?fen=8/N7/8/8/8/1N4b1/q7/k1K5&moves=9}\\{\sffamily\tiny\color{ink}scannen zum Lösen}
\end{minipage}
\vspace{2mm}

{\raggedright\noindent\textcolor{ink}{\sffamily\textsc{Themen}}\enspace \themetag{reines Matt} \themetag{Mustermatt} \themetag{Switchback} \themetag{Selbstblock} \themetag{Umnov}\par\smallskip{\footnotesize\textit{Hinweis:} die Partei am Zug steht im Diagramm im Schach; die Lösung beginnt mit einem Schlag.}\par}

\clearpage
\phantomsection\label{hm:KNPPvkp}
\begin{tcolorbox}[enhanced,colback=ink,colframe=ink,arc=0pt,boxrule=0pt,left=4mm,right=4mm,top=2.5mm,bottom=2.5mm]
\color{white}\sffamily{\LARGE\bfseries KNPPvkp}\hfill{\Large h\#9}\\[0.5mm]
{\small 6 Steine\quad Schwarz am Zug}\hfill{\small Nr.~314\qquad\href{https://helpman.komtera.lt/?fen=8/8/1N6/8/4p3/4P3/4P3/3K1k2&moves=9}{\color{white}\bfseries lösen\,\ensuremath{\nearrow}}}
\end{tcolorbox}
\vspace{2mm}
\noindent\begin{minipage}[t]{36mm}\vspace{0pt}\sffamily\small\raggedright
\hmfigure{längste eindeutige}{{\Large h\#9}}\hmfigure{eindeutig bei dieser Länge}{{\Large 4} Stellungen}\hmfigure{längstes Matt}{{\Large h\#14}}\hmfigure{in der Klasse}{25,6 Milliarden Stellungen,\\ 9,25 Milliarden mit Hilfsmatt}
\end{minipage}\hfill
\begin{minipage}[t]{80mm}\vspace{0pt}\begin{tikzpicture}\begin{axis}[width=72mm,height=42mm,ybar,bar width=1.6pt,xmin=-1,xmax=29,ymin=0,ymax=1.624,axis y line*=left,axis x line*=bottom,xlabel={\scriptsize Halbzüge bis Matt},ylabel={\scriptsize Stellungen (Mrd.)},tick label style={font=\tiny},label style={font=\scriptsize},enlarge x limits=false,ylabel shift=-3pt,xlabel shift=-3pt]\addplot[fill=ink!55,draw=none] coordinates {(0,0.0001085) (1,0.006069) (2,0.0153) (3,0.124) (4,0.3272) (5,1.046) (6,1.477) (7,1.408) (8,1.426) (9,1.023) (10,0.9992) (11,0.6246) (12,0.558) (13,0.1294) (14,0.07878) (15,0.004338) (16,0.002806) (17,0.0005928) (18,0.0004643) (19,4.731e-05) (20,3.393e-05) (21,2.043e-06) (22,1.635e-06) (23,7e-09) (24,4.7e-08) (26,2.4e-08) (28,1e-08)};\draw[densely dotted] (axis cs:18,0) -- (axis cs:18,1.624);\end{axis}\begin{axis}[width=72mm,height=42mm,axis y line*=right,axis x line=none,ymode=log,xmin=-1,xmax=29,ylabel={\scriptsize eindeutig (log)},tick label style={font=\tiny},label style={font=\scriptsize},ylabel shift=-3pt]\addplot[mark=*,mark size=0.7pt,draw=black,thin] coordinates {(0,108472) (1,1679425) (2,4790142) (3,52883911) (4,70692766) (5,74647792) (6,48082974) (7,29935406) (8,18205820) (9,10215697) (10,5718941) (11,2755901) (12,1151907) (13,23272) (14,9712) (15,1907) (16,500) (17,15) (18,4)};\end{axis}\end{tikzpicture}\end{minipage}
\vspace{3mm}

\noindent\begin{minipage}[t]{60mm}\vspace{0pt}\sbox\hmboard{\chessboard[setfen=8/8/1N6/8/4p3/4P3/4P3/3K1k2 b - - 0 1]}\usebox\hmboard\\[1mm]
\makebox[\wd\hmboard]{\sffamily\hspace{3mm}\textbf{h\#9}\hfill(4+2)}\end{minipage}\hfill
\begin{minipage}[t]{52mm}\vspace{0pt}
\textcolor{ink}{\sffamily\textsc{Lösung}}\par\smallskip\small
\begin{tabular}{@{}r@{\enspace}l@{\enspace}l@{}}
\textbf{1.} & \symking{}f2 & \symknight{}d7\\
\textbf{2.} & \symking{}xe3 & \symknight{}e5\\
\textbf{3.} & \symking{}f2 & \symknight{}f3\\
\textbf{4.} & exf3 & e4\\
\textbf{5.} & \symking{}g1 & e5\\
\textbf{6.} & f2 & e6\\
\textbf{7.} & f1=\symbishop{} & e7\\
\textbf{8.} & \symbishop{}g2 & e8=\symqueen{}\\
\textbf{9.} & \symking{}f1 & \symqueen{}e1\#\\
\end{tabular}\par\vspace{3mm}
\qrcode[height=12mm]{https://helpman.komtera.lt/?fen=8/8/1N6/8/4p3/4P3/4P3/3K1k2&moves=9}\\{\sffamily\tiny\color{ink}scannen zum Lösen}
\end{minipage}
\vspace{2mm}

{\raggedright\noindent\textcolor{ink}{\sffamily\textsc{Themen}}\enspace \themetag{Umwandlung} \themetag{Unterverwandlung} \themetag{Excelsior} \themetag{Switchback} \themetag{Selbstblock} \themetag{ohne Schach} \themetag{Umnov}\par}

\clearpage
\phantomsection\label{hm:KNvkbpp}
\begin{tcolorbox}[enhanced,colback=ink,colframe=ink,arc=0pt,boxrule=0pt,left=4mm,right=4mm,top=2.5mm,bottom=2.5mm]
\color{white}\sffamily{\LARGE\bfseries KNvkbpp}\hfill{\Large h\#9}\\[0.5mm]
{\small 6 Steine\quad Schwarz am Zug}\hfill{\small Nr.~315\qquad\href{https://helpman.komtera.lt/?fen=N7/2p5/2p5/8/8/K1b5/8/7k&moves=9}{\color{white}\bfseries lösen\,\ensuremath{\nearrow}}}
\end{tcolorbox}
\vspace{2mm}
\noindent\begin{minipage}[t]{36mm}\vspace{0pt}\sffamily\small\raggedright
\hmfigure{längste eindeutige}{{\Large h\#9}}\hmfigure{eindeutig bei dieser Länge}{{\Large 2} Stellungen}\hmfigure{längstes Matt}{{\Large h\#10.5}}\hmfigure{in der Klasse}{34,1 Milliarden Stellungen,\\ 11,8 Milliarden mit Hilfsmatt}
\end{minipage}\hfill
\begin{minipage}[t]{80mm}\vspace{0pt}\begin{tikzpicture}\begin{axis}[width=72mm,height=42mm,ybar,bar width=1.6pt,xmin=-1,xmax=22,ymin=0,ymax=2.388,axis y line*=left,axis x line*=bottom,xlabel={\scriptsize Halbzüge bis Matt},ylabel={\scriptsize Stellungen (Mrd.)},tick label style={font=\tiny},label style={font=\scriptsize},enlarge x limits=false,ylabel shift=-3pt,xlabel shift=-3pt]\addplot[fill=ink!55,draw=none] coordinates {(0,1.384e-05) (1,6.71e-05) (2,0.0002887) (3,0.001787) (4,0.007829) (5,0.03275) (6,0.08905) (7,0.2547) (8,0.5186) (9,1.023) (10,1.533) (11,2.171) (12,2.067) (13,1.983) (14,1.138) (15,0.681) (16,0.2196) (17,0.06842) (18,0.008747) (19,0.00106) (20,7.695e-06) (21,2.5e-07)};\draw[densely dotted] (axis cs:18,0) -- (axis cs:18,2.388);\end{axis}\begin{axis}[width=72mm,height=42mm,axis y line*=right,axis x line=none,ymode=log,xmin=-1,xmax=22,ylabel={\scriptsize eindeutig (log)},tick label style={font=\tiny},label style={font=\scriptsize},ylabel shift=-3pt]\addplot[mark=*,mark size=0.7pt,draw=black,thin] coordinates {(0,13843) (1,64978) (2,215059) (3,970822) (4,3563355) (5,5764842) (6,6254467) (7,7825649) (8,9198926) (9,9403146) (10,8550323) (11,5293813) (12,2230038) (13,711391) (14,245762) (15,31082) (16,4504) (17,164) (18,2)};\end{axis}\end{tikzpicture}\end{minipage}
\vspace{3mm}

\noindent\begin{minipage}[t]{60mm}\vspace{0pt}\sbox\hmboard{\chessboard[setfen=N7/2p5/2p5/8/8/K1b5/8/7k b - - 0 1]}\usebox\hmboard\\[1mm]
\makebox[\wd\hmboard]{\sffamily\hspace{3mm}\textbf{h\#9}\hfill(2+4)}\end{minipage}\hfill
\begin{minipage}[t]{52mm}\vspace{0pt}
\textcolor{ink}{\sffamily\textsc{Lösung}}\par\smallskip\small
\begin{tabular}{@{}r@{\enspace}l@{\enspace}l@{}}
\textbf{1.} & \symbishop{}e5 & \symking{}b4\\
\textbf{2.} & \symking{}g2 & \symking{}c5\\
\textbf{3.} & \symking{}f3 & \symking{}xc6\\
\textbf{4.} & \symking{}e4 & \symking{}b5\\
\textbf{5.} & \symking{}d5 & \symking{}a5\\
\textbf{6.} & \symking{}c6 & \symknight{}xc7\\
\textbf{7.} & \symking{}b7 & \symknight{}d5\\
\textbf{8.} & \symking{}a8 & \symking{}a6\\
\textbf{9.} & \symbishop{}b8 & \symknight{}b6\#\\
\end{tabular}\par\vspace{3mm}
\qrcode[height=12mm]{https://helpman.komtera.lt/?fen=N7/2p5/2p5/8/8/K1b5/8/7k&moves=9}\\{\sffamily\tiny\color{ink}scannen zum Lösen}
\end{minipage}
\vspace{2mm}

{\raggedright\noindent\textcolor{ink}{\sffamily\textsc{Themen}}\enspace \themetag{reines Matt} \themetag{Mustermatt} \themetag{Idealmatt} \themetag{Selbstblock} \themetag{ohne Schach}\par}

\clearpage
\phantomsection\label{hm:KNvknpp}
\begin{tcolorbox}[enhanced,colback=ink,colframe=ink,arc=0pt,boxrule=0pt,left=4mm,right=4mm,top=2.5mm,bottom=2.5mm]
\color{white}\sffamily{\LARGE\bfseries KNvknpp}\hfill{\Large h\#9}\\[0.5mm]
{\small 6 Steine\quad Schwarz am Zug}\hfill{\small Nr.~316\qquad\href{https://helpman.komtera.lt/?fen=K7/2p5/p7/2n5/8/7k/8/N7&moves=9}{\color{white}\bfseries lösen\,\ensuremath{\nearrow}}}
\end{tcolorbox}
\vspace{2mm}
\noindent\begin{minipage}[t]{36mm}\vspace{0pt}\sffamily\small\raggedright
\hmfigure{längste eindeutige}{{\Large h\#9}}\hmfigure{eindeutig bei dieser Länge}{{\Large 2} Stellungen}\hmfigure{längstes Matt}{{\Large h\#10.5}}\hmfigure{in der Klasse}{34,1 Milliarden Stellungen,\\ 12,1 Milliarden mit Hilfsmatt}
\end{minipage}\hfill
\begin{minipage}[t]{80mm}\vspace{0pt}\begin{tikzpicture}\begin{axis}[width=72mm,height=42mm,ybar,bar width=1.6pt,xmin=-1,xmax=22,ymin=0,ymax=2.365,axis y line*=left,axis x line*=bottom,xlabel={\scriptsize Halbzüge bis Matt},ylabel={\scriptsize Stellungen (Mrd.)},tick label style={font=\tiny},label style={font=\scriptsize},enlarge x limits=false,ylabel shift=-3pt,xlabel shift=-3pt]\addplot[fill=ink!55,draw=none] coordinates {(0,1.505e-05) (1,7.352e-05) (2,0.0002131) (3,0.001428) (4,0.006013) (5,0.02562) (6,0.08258) (7,0.2225) (8,0.5001) (9,0.9173) (10,1.445) (11,1.996) (12,2.138) (13,2.15) (14,1.386) (15,0.8243) (16,0.3049) (17,0.07807) (18,0.01528) (19,0.001093) (20,9.393e-05) (21,1e-09)};\draw[densely dotted] (axis cs:18,0) -- (axis cs:18,2.365);\end{axis}\begin{axis}[width=72mm,height=42mm,axis y line*=right,axis x line=none,ymode=log,xmin=-1,xmax=22,ylabel={\scriptsize eindeutig (log)},tick label style={font=\tiny},label style={font=\scriptsize},ylabel shift=-3pt]\addplot[mark=*,mark size=0.7pt,draw=black,thin] coordinates {(0,15046) (1,70886) (2,137545) (3,674774) (4,2652260) (5,4414146) (6,6005659) (7,5635388) (8,6148364) (9,6014126) (10,5567326) (11,3880657) (12,1747039) (13,551893) (14,124949) (15,13124) (16,1914) (17,20) (18,2)};\end{axis}\end{tikzpicture}\end{minipage}
\vspace{3mm}

\noindent\begin{minipage}[t]{60mm}\vspace{0pt}\sbox\hmboard{\chessboard[setfen=K7/2p5/p7/2n5/8/7k/8/N7 b - - 0 1]}\usebox\hmboard\\[1mm]
\makebox[\wd\hmboard]{\sffamily\hspace{3mm}\textbf{h\#9}\hfill(2+4)}\end{minipage}\hfill
\begin{minipage}[t]{52mm}\vspace{0pt}
\textcolor{ink}{\sffamily\textsc{Lösung}}\par\smallskip\small
\begin{tabular}{@{}r@{\enspace}l@{\enspace}l@{}}
\textbf{1.} & \symking{}g4 & \symking{}b8\\
\textbf{2.} & \symking{}f5 & \symking{}xc7\\
\textbf{3.} & \symking{}e6 & \symking{}b6\\
\textbf{4.} & \symking{}d7 & \symknight{}c2\\
\textbf{5.} & \symking{}c8 & \symknight{}b4\\
\textbf{6.} & \symking{}b8 & \symknight{}xa6+\\
\textbf{7.} & \symking{}a8 & \symknight{}b4\\
\textbf{8.} & \symknight{}a6 & \symknight{}d5\\
\textbf{9.} & \symknight{}b8 & \symknight{}c7\#\\
\end{tabular}\par\vspace{3mm}
\qrcode[height=12mm]{https://helpman.komtera.lt/?fen=K7/2p5/p7/2n5/8/7k/8/N7&moves=9}\\{\sffamily\tiny\color{ink}scannen zum Lösen}
\end{minipage}
\vspace{2mm}

{\raggedright\noindent\textcolor{ink}{\sffamily\textsc{Themen}}\enspace \themetag{reines Matt} \themetag{Mustermatt} \themetag{Idealmatt} \themetag{Switchback} \themetag{Selbstblock} \themetag{Umnov}\par}

\clearpage
\phantomsection\label{hm:KQBBvkn}
\begin{tcolorbox}[enhanced,colback=ink,colframe=ink,arc=0pt,boxrule=0pt,left=4mm,right=4mm,top=2.5mm,bottom=2.5mm]
\color{white}\sffamily{\LARGE\bfseries KQBBvkn}\hfill{\Large h\#9}\\[0.5mm]
{\small 6 Steine\quad Schwarz am Zug}\hfill{\small Nr.~317\qquad\href{https://helpman.komtera.lt/?fen=8/8/1nB5/8/k7/Q7/8/1K5B&moves=9}{\color{white}\bfseries lösen\,\ensuremath{\nearrow}}}
\end{tcolorbox}
\vspace{2mm}
\noindent\begin{minipage}[t]{36mm}\vspace{0pt}\sffamily\small\raggedright
\hmfigure{längste eindeutige}{{\Large h\#9}}\hmfigure{eindeutig bei dieser Länge}{{\Large 8} Stellungen}\hmfigure{längstes Matt}{{\Large h\#12}}\hmfigure{in der Klasse}{15,5 Milliarden Stellungen,\\ 4,29 Milliarden mit Hilfsmatt}
\end{minipage}\hfill
\begin{minipage}[t]{80mm}\vspace{0pt}\begin{tikzpicture}\begin{axis}[width=72mm,height=42mm,ybar,bar width=1.6pt,xmin=-1,xmax=25,ymin=0,ymax=2.236,axis y line*=left,axis x line*=bottom,xlabel={\scriptsize Halbzüge bis Matt},ylabel={\scriptsize Stellungen (Mrd.)},tick label style={font=\tiny},label style={font=\scriptsize},enlarge x limits=false,ylabel shift=-3pt,xlabel shift=-3pt]\addplot[fill=ink!55,draw=none] coordinates {(0,0.02748) (1,0.1262) (2,0.5572) (3,1.046) (4,2.033) (5,0.3451) (6,0.1107) (7,0.0009905) (8,0.009063) (9,1.73e-07) (10,0.006344) (12,0.004699) (14,0.006612) (16,0.007491) (18,0.005771) (20,0.003526) (22,0.001298) (24,0.0001918)};\draw[densely dotted] (axis cs:18,0) -- (axis cs:18,2.236);\end{axis}\begin{axis}[width=72mm,height=42mm,axis y line*=right,axis x line=none,ymode=log,xmin=-1,xmax=25,ylabel={\scriptsize eindeutig (log)},tick label style={font=\tiny},label style={font=\scriptsize},ylabel shift=-3pt]\addplot[mark=*,mark size=0.7pt,draw=black,thin] coordinates {(0,27475621) (1,97598747) (2,294116751) (3,102395050) (4,35930057) (5,129088) (6,425277) (7,32) (8,65812) (10,18453) (12,8046) (14,2848) (16,291) (18,8)};\end{axis}\end{tikzpicture}\end{minipage}
\vspace{3mm}

\noindent\begin{minipage}[t]{60mm}\vspace{0pt}\sbox\hmboard{\chessboard[setfen=8/8/1nB5/8/k7/Q7/8/1K5B b - - 0 1]}\usebox\hmboard\\[1mm]
\makebox[\wd\hmboard]{\sffamily\hspace{3mm}\textbf{h\#9}\hfill(4+2)}\end{minipage}\hfill
\begin{minipage}[t]{52mm}\vspace{0pt}
\textcolor{ink}{\sffamily\textsc{Lösung}}\par\smallskip\small
\begin{tabular}{@{}r@{\enspace}l@{\enspace}l@{}}
\textbf{1.} & \symking{}xa3 & \symking{}c2\\
\textbf{2.} & \symking{}b4 & \symking{}d3\\
\textbf{3.} & \symking{}c5 & \symking{}e4\\
\textbf{4.} & \symking{}xc6 & \symking{}e5+\\
\textbf{5.} & \symking{}c7 & \symbishop{}b7\\
\textbf{6.} & \symking{}b8 & \symbishop{}a6\\
\textbf{7.} & \symking{}a8 & \symking{}d6\\
\textbf{8.} & \symknight{}c8+ & \symking{}c7\\
\textbf{9.} & \symknight{}a7 & \symbishop{}b7\#\\
\end{tabular}\par\vspace{3mm}
\qrcode[height=12mm]{https://helpman.komtera.lt/?fen=8/8/1nB5/8/k7/Q7/8/1K5B&moves=9}\\{\sffamily\tiny\color{ink}scannen zum Lösen}
\end{minipage}
\vspace{2mm}

{\raggedright\noindent\textcolor{ink}{\sffamily\textsc{Themen}}\enspace \themetag{reines Matt} \themetag{Mustermatt} \themetag{Idealmatt} \themetag{Switchback} \themetag{Selbstblock}\par\smallskip{\footnotesize\textit{Hinweis:} das Diagramm hat keinen legalen letzten Zug, kann also in keiner Partie entstehen.}\par}

\clearpage
\phantomsection\label{hm:KQBBvkp}
\begin{tcolorbox}[enhanced,colback=ink,colframe=ink,arc=0pt,boxrule=0pt,left=4mm,right=4mm,top=2.5mm,bottom=2.5mm]
\color{white}\sffamily{\LARGE\bfseries KQBBvkp}\hfill{\Large h\#9}\\[0.5mm]
{\small 6 Steine\quad Schwarz am Zug}\hfill{\small Nr.~318\qquad\href{https://helpman.komtera.lt/?fen=8/8/8/3B4/8/8/3Kp3/kQ1B4&moves=9}{\color{white}\bfseries lösen\,\ensuremath{\nearrow}}}
\end{tcolorbox}
\vspace{2mm}
\noindent\begin{minipage}[t]{36mm}\vspace{0pt}\sffamily\small\raggedright
\hmfigure{längste eindeutige}{{\Large h\#9}}\hmfigure{eindeutig bei dieser Länge}{{\Large 19} Stellungen}\hmfigure{längstes Matt}{{\Large h\#15}}\hmfigure{in der Klasse}{45,4 Milliarden Stellungen,\\ 13,4 Milliarden mit Hilfsmatt}
\end{minipage}\hfill
\begin{minipage}[t]{80mm}\vspace{0pt}\begin{tikzpicture}\begin{axis}[width=72mm,height=42mm,ybar,bar width=1.6pt,xmin=-1,xmax=31,ymin=0,ymax=7.01,axis y line*=left,axis x line*=bottom,xlabel={\scriptsize Halbzüge bis Matt},ylabel={\scriptsize Stellungen (Mrd.)},tick label style={font=\tiny},label style={font=\scriptsize},enlarge x limits=false,ylabel shift=-3pt,xlabel shift=-3pt]\addplot[fill=ink!55,draw=none] coordinates {(0,0.1009) (1,0.4574) (2,1.881) (3,3.264) (4,6.373) (5,0.8645) (6,0.3153) (7,0.001042) (8,0.02747) (9,1e-09) (10,0.01695) (12,0.009827) (14,0.01149) (15,2e-09) (16,0.01523) (17,2e-09) (18,0.01617) (20,0.01324) (21,6.7e-08) (22,0.01142) (24,0.007909) (25,2e-09) (26,0.004239) (28,0.002278) (30,0.0004146)};\draw[densely dotted] (axis cs:18,0) -- (axis cs:18,7.01);\end{axis}\begin{axis}[width=72mm,height=42mm,axis y line*=right,axis x line=none,ymode=log,xmin=-1,xmax=31,ylabel={\scriptsize eindeutig (log)},tick label style={font=\tiny},label style={font=\scriptsize},ylabel shift=-3pt]\addplot[mark=*,mark size=0.7pt,draw=black,thin] coordinates {(0,100926473) (1,343486278) (2,1135683986) (3,362146759) (4,124084048) (5,162637) (6,1596169) (7,5) (8,198143) (10,53146) (12,24145) (14,6042) (16,386) (18,19)};\end{axis}\end{tikzpicture}\end{minipage}
\vspace{3mm}

\noindent\begin{minipage}[t]{60mm}\vspace{0pt}\sbox\hmboard{\chessboard[setfen=8/8/8/3B4/8/8/3Kp3/kQ1B4 b - - 0 1]}\usebox\hmboard\\[1mm]
\makebox[\wd\hmboard]{\sffamily\hspace{3mm}\textbf{h\#9}\hfill(4+2)}\end{minipage}\hfill
\begin{minipage}[t]{52mm}\vspace{0pt}
\textcolor{ink}{\sffamily\textsc{Lösung}}\par\smallskip\small
\begin{tabular}{@{}r@{\enspace}l@{\enspace}l@{}}
\textbf{1.} & \symking{}xb1 & \symking{}e3\\
\textbf{2.} & \symking{}c1 & \symking{}f4\\
\textbf{3.} & \symking{}d2 & \symbishop{}g2\\
\textbf{4.} & \symking{}e1 & \symbishop{}f1\\
\textbf{5.} & exf1=\symknight{} & \symbishop{}g4\\
\textbf{6.} & \symking{}f2 & \symbishop{}h3\\
\textbf{7.} & \symking{}g1 & \symking{}f3\\
\textbf{8.} & \symking{}h1 & \symking{}f2\\
\textbf{9.} & \symknight{}h2 & \symbishop{}g2\#\\
\end{tabular}\par\vspace{3mm}
\qrcode[height=12mm]{https://helpman.komtera.lt/?fen=8/8/8/3B4/8/8/3Kp3/kQ1B4&moves=9}\\{\sffamily\tiny\color{ink}scannen zum Lösen}
\end{minipage}
\vspace{2mm}

{\raggedright\noindent\textcolor{ink}{\sffamily\textsc{Themen}}\enspace \themetag{reines Matt} \themetag{Mustermatt} \themetag{Idealmatt} \themetag{Umwandlung} \themetag{Unterverwandlung} \themetag{Selbstblock} \themetag{ohne Schach}\par\smallskip{\footnotesize\textit{Hinweis:} die Partei am Zug steht im Diagramm im Schach; die Lösung beginnt mit einem Schlag.}\par}

\clearpage
\phantomsection\label{hm:KQBvkbn}
\begin{tcolorbox}[enhanced,colback=ink,colframe=ink,arc=0pt,boxrule=0pt,left=4mm,right=4mm,top=2.5mm,bottom=2.5mm]
\color{white}\sffamily{\LARGE\bfseries KQBvkbn}\hfill{\Large h\#9}\\[0.5mm]
{\small 6 Steine\quad Schwarz am Zug}\hfill{\small Nr.~319\qquad\href{https://helpman.komtera.lt/?fen=8/8/1n6/8/8/1b6/2Q5/K1k4B&moves=9}{\color{white}\bfseries lösen\,\ensuremath{\nearrow}}}
\end{tcolorbox}
\vspace{2mm}
\noindent\begin{minipage}[t]{36mm}\vspace{0pt}\sffamily\small\raggedright
\hmfigure{längste eindeutige}{{\Large h\#9}}\hmfigure{eindeutig bei dieser Länge}{{\Large 160} Stellungen}\hmfigure{längstes Matt}{{\Large h\#12.5}}\hmfigure{in der Klasse}{15,5 Milliarden Stellungen,\\ 8,38 Milliarden mit Hilfsmatt}
\end{minipage}\hfill
\begin{minipage}[t]{80mm}\vspace{0pt}\begin{tikzpicture}\begin{axis}[width=72mm,height=42mm,ybar,bar width=1.6pt,xmin=-1,xmax=26,ymin=0,ymax=3.997,axis y line*=left,axis x line*=bottom,xlabel={\scriptsize Halbzüge bis Matt},ylabel={\scriptsize Stellungen (Mrd.)},tick label style={font=\tiny},label style={font=\scriptsize},enlarge x limits=false,ylabel shift=-3pt,xlabel shift=-3pt]\addplot[fill=ink!55,draw=none] coordinates {(0,0.01651) (1,0.1069) (2,0.4803) (3,1.693) (4,3.633) (5,1.708) (6,0.6304) (7,0.02526) (8,0.00816) (9,5.38e-07) (10,0.009853) (11,1.569e-06) (12,0.01604) (13,2.209e-06) (14,0.01666) (15,9.08e-07) (16,0.01546) (17,2.637e-06) (18,0.01138) (19,3.09e-07) (20,0.005428) (21,1.92e-07) (22,0.00132) (23,1e-09) (24,6.77e-05) (25,1e-09)};\draw[densely dotted] (axis cs:18,0) -- (axis cs:18,3.997);\end{axis}\begin{axis}[width=72mm,height=42mm,axis y line*=right,axis x line=none,ymode=log,xmin=-1,xmax=26,ylabel={\scriptsize eindeutig (log)},tick label style={font=\tiny},label style={font=\scriptsize},ylabel shift=-3pt]\addplot[mark=*,mark size=0.7pt,draw=black,thin] coordinates {(0,16505538) (1,91038195) (2,281268765) (3,351213758) (4,193246273) (5,4613085) (6,315297) (7,12) (8,128205) (9,13) (10,66882) (12,29867) (14,5936) (16,1353) (18,160)};\end{axis}\end{tikzpicture}\end{minipage}
\vspace{3mm}

\noindent\begin{minipage}[t]{60mm}\vspace{0pt}\sbox\hmboard{\chessboard[setfen=8/8/1n6/8/8/1b6/2Q5/K1k4B b - - 0 1]}\usebox\hmboard\\[1mm]
\makebox[\wd\hmboard]{\sffamily\hspace{3mm}\textbf{h\#9}\hfill(3+3)}\end{minipage}\hfill
\begin{minipage}[t]{52mm}\vspace{0pt}
\textcolor{ink}{\sffamily\textsc{Lösung}}\par\smallskip\small
\begin{tabular}{@{}r@{\enspace}l@{\enspace}l@{}}
\textbf{1.} & \symking{}xc2 & \symbishop{}b7\\
\textbf{2.} & \symking{}c3 & \symbishop{}a6\\
\textbf{3.} & \symking{}b4 & \symking{}b2\\
\textbf{4.} & \symking{}a5 & \symking{}xb3\\
\textbf{5.} & \symknight{}c8 & \symking{}c4\\
\textbf{6.} & \symking{}b6 & \symking{}d5\\
\textbf{7.} & \symking{}a7 & \symking{}c6\\
\textbf{8.} & \symking{}a8 & \symking{}c7\\
\textbf{9.} & \symknight{}a7 & \symbishop{}b7\#\\
\end{tabular}\par\vspace{3mm}
\qrcode[height=12mm]{https://helpman.komtera.lt/?fen=8/8/1n6/8/8/1b6/2Q5/K1k4B&moves=9}\\{\sffamily\tiny\color{ink}scannen zum Lösen}
\end{minipage}
\vspace{2mm}

{\raggedright\noindent\textcolor{ink}{\sffamily\textsc{Themen}}\enspace \themetag{reines Matt} \themetag{Mustermatt} \themetag{Idealmatt} \themetag{Switchback} \themetag{Selbstblock} \themetag{ohne Schach}\par\smallskip{\footnotesize\textit{Hinweis:} die Partei am Zug steht im Diagramm im Schach; die Lösung beginnt mit einem Schlag.}\par}

\clearpage
\phantomsection\label{hm:KQBvknn}
\begin{tcolorbox}[enhanced,colback=ink,colframe=ink,arc=0pt,boxrule=0pt,left=4mm,right=4mm,top=2.5mm,bottom=2.5mm]
\color{white}\sffamily{\LARGE\bfseries KQBvknn}\hfill{\Large h\#9}\\[0.5mm]
{\small 6 Steine\quad Schwarz am Zug}\hfill{\small Nr.~320\qquad\href{https://helpman.komtera.lt/?fen=8/8/1n6/8/8/n7/B1Q5/K1k5&moves=9}{\color{white}\bfseries lösen\,\ensuremath{\nearrow}}}
\end{tcolorbox}
\vspace{2mm}
\noindent\begin{minipage}[t]{36mm}\vspace{0pt}\sffamily\small\raggedright
\hmfigure{längste eindeutige}{{\Large h\#9}}\hmfigure{eindeutig bei dieser Länge}{{\Large 112} Stellungen}\hmfigure{längstes Matt}{{\Large h\#13}}\hmfigure{in der Klasse}{15,5 Milliarden Stellungen,\\ 4,31 Milliarden mit Hilfsmatt}
\end{minipage}\hfill
\begin{minipage}[t]{80mm}\vspace{0pt}\begin{tikzpicture}\begin{axis}[width=72mm,height=42mm,ybar,bar width=1.6pt,xmin=-1,xmax=27,ymin=0,ymax=2.024,axis y line*=left,axis x line*=bottom,xlabel={\scriptsize Halbzüge bis Matt},ylabel={\scriptsize Stellungen (Mrd.)},tick label style={font=\tiny},label style={font=\scriptsize},enlarge x limits=false,ylabel shift=-3pt,xlabel shift=-3pt]\addplot[fill=ink!55,draw=none] coordinates {(0,0.008747) (1,0.05508) (2,0.2301) (3,0.8282) (4,1.84) (5,0.8666) (6,0.4202) (7,0.01672) (8,0.003301) (9,1.78e-07) (10,0.004019) (11,6.48e-07) (12,0.006682) (13,1.09e-06) (14,0.01024) (15,6.16e-07) (16,0.01043) (17,7.97e-07) (18,0.006995) (19,1.655e-06) (20,0.003088) (21,7.52e-07) (22,0.0004053) (23,2e-09) (24,1.979e-06) (26,1e-09)};\draw[densely dotted] (axis cs:18,0) -- (axis cs:18,2.024);\end{axis}\begin{axis}[width=72mm,height=42mm,axis y line*=right,axis x line=none,ymode=log,xmin=-1,xmax=27,ylabel={\scriptsize eindeutig (log)},tick label style={font=\tiny},label style={font=\scriptsize},ylabel shift=-3pt]\addplot[mark=*,mark size=0.7pt,draw=black,thin] coordinates {(0,8747043) (1,46129770) (2,131172452) (3,175761226) (4,102249735) (5,2566913) (6,168775) (7,5) (8,51201) (9,10) (10,37986) (11,548) (12,22040) (13,18) (14,5140) (16,828) (18,112)};\end{axis}\end{tikzpicture}\end{minipage}
\vspace{3mm}

\noindent\begin{minipage}[t]{60mm}\vspace{0pt}\sbox\hmboard{\chessboard[setfen=8/8/1n6/8/8/n7/B1Q5/K1k5 b - - 0 1]}\usebox\hmboard\\[1mm]
\makebox[\wd\hmboard]{\sffamily\hspace{3mm}\textbf{h\#9}\hfill(3+3)}\end{minipage}\hfill
\begin{minipage}[t]{52mm}\vspace{0pt}
\textcolor{ink}{\sffamily\textsc{Lösung}}\par\smallskip\small
\begin{tabular}{@{}r@{\enspace}l@{\enspace}l@{}}
\textbf{1.} & \symking{}xc2 & \symbishop{}c4\\
\textbf{2.} & \symking{}c3 & \symbishop{}a6\\
\textbf{3.} & \symking{}b4 & \symking{}b2\\
\textbf{4.} & \symking{}a5 & \symking{}c3\\
\textbf{5.} & \symknight{}c8 & \symking{}d4\\
\textbf{6.} & \symking{}b6 & \symking{}d5\\
\textbf{7.} & \symking{}a7 & \symking{}c6\\
\textbf{8.} & \symking{}a8 & \symking{}c7\\
\textbf{9.} & \symknight{}a7 & \symbishop{}b7\#\\
\end{tabular}\par\vspace{3mm}
\qrcode[height=12mm]{https://helpman.komtera.lt/?fen=8/8/1n6/8/8/n7/B1Q5/K1k5&moves=9}\\{\sffamily\tiny\color{ink}scannen zum Lösen}
\end{minipage}
\vspace{2mm}

{\raggedright\noindent\textcolor{ink}{\sffamily\textsc{Themen}}\enspace \themetag{reines Matt} \themetag{Mustermatt} \themetag{Selbstblock} \themetag{ohne Schach}\par\smallskip{\footnotesize\textit{Hinweis:} die Partei am Zug steht im Diagramm im Schach; die Lösung beginnt mit einem Schlag.}\par}

\clearpage
\phantomsection\label{hm:KQBvkpp}
\begin{tcolorbox}[enhanced,colback=ink,colframe=ink,arc=0pt,boxrule=0pt,left=4mm,right=4mm,top=2.5mm,bottom=2.5mm]
\color{white}\sffamily{\LARGE\bfseries KQBvkpp}\hfill{\Large h\#9}\\[0.5mm]
{\small 6 Steine\quad Schwarz am Zug}\hfill{\small Nr.~321\qquad\href{https://helpman.komtera.lt/?fen=8/8/8/kQ6/8/5p2/5p2/KB6&moves=9}{\color{white}\bfseries lösen\,\ensuremath{\nearrow}}}
\end{tcolorbox}
\vspace{2mm}
\noindent\begin{minipage}[t]{36mm}\vspace{0pt}\sffamily\small\raggedright
\hmfigure{längste eindeutige}{{\Large h\#9}}\hmfigure{eindeutig bei dieser Länge}{{\Large 19} Stellungen}\hmfigure{längstes Matt}{{\Large h\#16}}\hmfigure{in der Klasse}{34,1 Milliarden Stellungen,\\ 10,4 Milliarden mit Hilfsmatt}
\end{minipage}\hfill
\begin{minipage}[t]{80mm}\vspace{0pt}\begin{tikzpicture}\begin{axis}[width=72mm,height=42mm,ybar,bar width=1.6pt,xmin=-1,xmax=33,ymin=0,ymax=5.248,axis y line*=left,axis x line*=bottom,xlabel={\scriptsize Halbzüge bis Matt},ylabel={\scriptsize Stellungen (Mrd.)},tick label style={font=\tiny},label style={font=\scriptsize},enlarge x limits=false,ylabel shift=-3pt,xlabel shift=-3pt]\addplot[fill=ink!55,draw=none] coordinates {(0,0.03089) (1,0.1876) (2,0.7238) (3,2.252) (4,4.771) (5,1.545) (6,0.7614) (7,0.01096) (8,0.007415) (9,2.59e-07) (10,0.007694) (11,1.01e-07) (12,0.01328) (13,1.7e-08) (14,0.01779) (15,1.2e-08) (16,0.01914) (17,1e-08) (18,0.01595) (19,1.5e-08) (20,0.01148) (21,2.2e-08) (22,0.008054) (23,1.1e-08) (24,0.004345) (25,2e-08) (26,0.002144) (27,5e-09) (28,0.0008433) (29,4e-09) (30,0.0001851) (32,3e-05)};\draw[densely dotted] (axis cs:18,0) -- (axis cs:18,5.248);\end{axis}\begin{axis}[width=72mm,height=42mm,axis y line*=right,axis x line=none,ymode=log,xmin=-1,xmax=33,ylabel={\scriptsize eindeutig (log)},tick label style={font=\tiny},label style={font=\scriptsize},ylabel shift=-3pt]\addplot[mark=*,mark size=0.7pt,draw=black,thin] coordinates {(0,30890392) (1,151149344) (2,456333503) (3,504134331) (4,252629591) (5,1768355) (6,241505) (7,9) (8,86982) (9,1) (10,52622) (12,26841) (14,4538) (16,397) (18,19)};\end{axis}\end{tikzpicture}\end{minipage}
\vspace{3mm}

\noindent\begin{minipage}[t]{60mm}\vspace{0pt}\sbox\hmboard{\chessboard[setfen=8/8/8/kQ6/8/5p2/5p2/KB6 b - - 0 1]}\usebox\hmboard\\[1mm]
\makebox[\wd\hmboard]{\sffamily\hspace{3mm}\textbf{h\#9}\hfill(3+3)}\end{minipage}\hfill
\begin{minipage}[t]{52mm}\vspace{0pt}
\textcolor{ink}{\sffamily\textsc{Lösung}}\par\smallskip\small
\begin{tabular}{@{}r@{\enspace}l@{\enspace}l@{}}
\textbf{1.} & \symking{}xb5 & \symbishop{}f5\\
\textbf{2.} & \symking{}c4 & \symbishop{}h3\\
\textbf{3.} & \symking{}d3 & \symking{}b2\\
\textbf{4.} & \symking{}e2 & \symking{}c3\\
\textbf{5.} & f1=\symknight{} & \symking{}d4\\
\textbf{6.} & \symking{}f2 & \symking{}e4\\
\textbf{7.} & \symking{}g1 & \symking{}xf3\\
\textbf{8.} & \symking{}h1 & \symking{}f2\\
\textbf{9.} & \symknight{}h2 & \symbishop{}g2\#\\
\end{tabular}\par\vspace{3mm}
\qrcode[height=12mm]{https://helpman.komtera.lt/?fen=8/8/8/kQ6/8/5p2/5p2/KB6&moves=9}\\{\sffamily\tiny\color{ink}scannen zum Lösen}
\end{minipage}
\vspace{2mm}

{\raggedright\noindent\textcolor{ink}{\sffamily\textsc{Themen}}\enspace \themetag{reines Matt} \themetag{Mustermatt} \themetag{Idealmatt} \themetag{Umwandlung} \themetag{Unterverwandlung} \themetag{Selbstblock} \themetag{ohne Schach}\par\smallskip{\footnotesize\textit{Hinweis:} die Partei am Zug steht im Diagramm im Schach; die Lösung beginnt mit einem Schlag.}\par}

\clearpage
\phantomsection\label{hm:KQNPvkp}
\begin{tcolorbox}[enhanced,colback=ink,colframe=ink,arc=0pt,boxrule=0pt,left=4mm,right=4mm,top=2.5mm,bottom=2.5mm]
\color{white}\sffamily{\LARGE\bfseries KQNPvkp}\hfill{\Large h\#9}\\[0.5mm]
{\small 6 Steine\quad Schwarz am Zug}\hfill{\small Nr.~322\qquad\href{https://helpman.komtera.lt/?fen=N7/8/8/8/Q3p3/k7/4P3/K7&moves=9}{\color{white}\bfseries lösen\,\ensuremath{\nearrow}}}
\end{tcolorbox}
\vspace{2mm}
\noindent\begin{minipage}[t]{36mm}\vspace{0pt}\sffamily\small\raggedright
\hmfigure{längste eindeutige}{{\Large h\#9}}\hmfigure{eindeutig bei dieser Länge}{{\Large 79} Stellungen}\hmfigure{längstes Matt}{{\Large h\#12}}\hmfigure{in der Klasse}{34,1 Milliarden Stellungen,\\ 21,4 Milliarden mit Hilfsmatt}
\end{minipage}\hfill
\begin{minipage}[t]{80mm}\vspace{0pt}\begin{tikzpicture}\begin{axis}[width=72mm,height=42mm,ybar,bar width=1.6pt,xmin=-1,xmax=25,ymin=0,ymax=8.797,axis y line*=left,axis x line*=bottom,xlabel={\scriptsize Halbzüge bis Matt},ylabel={\scriptsize Stellungen (Mrd.)},tick label style={font=\tiny},label style={font=\scriptsize},enlarge x limits=false,ylabel shift=-3pt,xlabel shift=-3pt]\addplot[fill=ink!55,draw=none] coordinates {(0,0.0465) (1,0.3361) (2,1.194) (3,3.478) (4,7.997) (5,4.044) (6,3.633) (7,0.3692) (8,0.1683) (9,0.000866) (10,0.04146) (11,1.27e-07) (12,0.04481) (13,1.4e-08) (14,0.0231) (16,0.003563) (18,0.001207) (20,0.0001327) (22,6.751e-06) (24,1.3e-08)};\draw[densely dotted] (axis cs:18,0) -- (axis cs:18,8.797);\end{axis}\begin{axis}[width=72mm,height=42mm,axis y line*=right,axis x line=none,ymode=log,xmin=-1,xmax=25,ylabel={\scriptsize eindeutig (log)},tick label style={font=\tiny},label style={font=\scriptsize},ylabel shift=-3pt]\addplot[mark=*,mark size=0.7pt,draw=black,thin] coordinates {(0,46502994) (1,251429773) (2,756276781) (3,632879559) (4,567731249) (5,38478263) (6,13437662) (7,60109) (8,717683) (9,15) (10,385022) (12,227917) (14,13254) (16,7077) (18,79)};\end{axis}\end{tikzpicture}\end{minipage}
\vspace{3mm}

\noindent\begin{minipage}[t]{60mm}\vspace{0pt}\sbox\hmboard{\chessboard[setfen=N7/8/8/8/Q3p3/k7/4P3/K7 b - - 0 1]}\usebox\hmboard\\[1mm]
\makebox[\wd\hmboard]{\sffamily\hspace{3mm}\textbf{h\#9}\hfill(4+2)}\end{minipage}\hfill
\begin{minipage}[t]{52mm}\vspace{0pt}
\textcolor{ink}{\sffamily\textsc{Lösung}}\par\smallskip\small
\begin{tabular}{@{}r@{\enspace}l@{\enspace}l@{}}
\textbf{1.} & \symking{}xa4 & \symknight{}b6+\\
\textbf{2.} & \symking{}b3 & \symknight{}c4\\
\textbf{3.} & e3 & \symknight{}d2+\\
\textbf{4.} & exd2 & e4\\
\textbf{5.} & \symking{}c2 & e5\\
\textbf{6.} & \symking{}c1 & e6\\
\textbf{7.} & d1=\symrook{} & e7\\
\textbf{8.} & \symrook{}d2 & e8=\symqueen{}\\
\textbf{9.} & \symrook{}c2 & \symqueen{}e1\#\\
\end{tabular}\par\vspace{3mm}
\qrcode[height=12mm]{https://helpman.komtera.lt/?fen=N7/8/8/8/Q3p3/k7/4P3/K7&moves=9}\\{\sffamily\tiny\color{ink}scannen zum Lösen}
\end{minipage}
\vspace{2mm}

{\raggedright\noindent\textcolor{ink}{\sffamily\textsc{Themen}}\enspace \themetag{Umwandlung} \themetag{Unterverwandlung} \themetag{Excelsior} \themetag{Switchback} \themetag{Selbstblock} \themetag{Phönix}\par\smallskip{\footnotesize\textit{Hinweis:} die Partei am Zug steht im Diagramm im Schach; die Lösung beginnt mit einem Schlag.}\par}

\clearpage
\phantomsection\label{hm:KQNvkbb}
\begin{tcolorbox}[enhanced,colback=ink,colframe=ink,arc=0pt,boxrule=0pt,left=4mm,right=4mm,top=2.5mm,bottom=2.5mm]
\color{white}\sffamily{\LARGE\bfseries KQNvkbb}\hfill{\Large h\#9}\\[0.5mm]
{\small 6 Steine\quad Schwarz am Zug}\hfill{\small Nr.~323\qquad\href{https://helpman.komtera.lt/?fen=7Q/7k/4b1b1/8/8/8/2N5/1K6&moves=9}{\color{white}\bfseries lösen\,\ensuremath{\nearrow}}}
\end{tcolorbox}
\vspace{2mm}
\noindent\begin{minipage}[t]{36mm}\vspace{0pt}\sffamily\small\raggedright
\hmfigure{längste eindeutige}{{\Large h\#9}}\hmfigure{eindeutig bei dieser Länge}{{\Large 9} Stellungen}\hmfigure{längstes Matt}{{\Large h\#10}}\hmfigure{in der Klasse}{15,5 Milliarden Stellungen,\\ 4,16 Milliarden mit Hilfsmatt}
\end{minipage}\hfill
\begin{minipage}[t]{80mm}\vspace{0pt}\begin{tikzpicture}\begin{axis}[width=72mm,height=42mm,ybar,bar width=1.6pt,xmin=-1,xmax=21,ymin=0,ymax=1.452,axis y line*=left,axis x line*=bottom,xlabel={\scriptsize Halbzüge bis Matt},ylabel={\scriptsize Stellungen (Mrd.)},tick label style={font=\tiny},label style={font=\scriptsize},enlarge x limits=false,ylabel shift=-3pt,xlabel shift=-3pt]\addplot[fill=ink!55,draw=none] coordinates {(0,0.004341) (1,0.03559) (2,0.1577) (3,0.5824) (4,1.32) (5,1.1) (6,0.7427) (7,0.158) (8,0.02766) (9,0.001136) (10,0.00439) (11,2.99e-07) (12,0.007498) (13,4.24e-07) (14,0.007727) (15,2.43e-07) (16,0.007502) (17,3.17e-07) (18,0.002674) (19,4e-09) (20,0.0001113)};\draw[densely dotted] (axis cs:18,0) -- (axis cs:18,1.452);\end{axis}\begin{axis}[width=72mm,height=42mm,axis y line*=right,axis x line=none,ymode=log,xmin=-1,xmax=21,ylabel={\scriptsize eindeutig (log)},tick label style={font=\tiny},label style={font=\scriptsize},ylabel shift=-3pt]\addplot[mark=*,mark size=0.7pt,draw=black,thin] coordinates {(0,4340897) (1,30788423) (2,90476683) (3,114662424) (4,104046613) (5,16905978) (6,2582446) (7,81424) (8,44599) (10,34357) (11,3) (12,17683) (14,1223) (16,1275) (18,9)};\end{axis}\end{tikzpicture}\end{minipage}
\vspace{3mm}

\noindent\begin{minipage}[t]{60mm}\vspace{0pt}\sbox\hmboard{\chessboard[setfen=7Q/7k/4b1b1/8/8/8/2N5/1K6 b - - 0 1]}\usebox\hmboard\\[1mm]
\makebox[\wd\hmboard]{\sffamily\hspace{3mm}\textbf{h\#9}\hfill(3+3)}\end{minipage}\hfill
\begin{minipage}[t]{52mm}\vspace{0pt}
\textcolor{ink}{\sffamily\textsc{Lösung}}\par\smallskip\small
\begin{tabular}{@{}r@{\enspace}l@{\enspace}l@{}}
\textbf{1.} & \symking{}xh8 & \symking{}b2\\
\textbf{2.} & \symking{}g7 & \symknight{}d4\\
\textbf{3.} & \symking{}f6 & \symknight{}xe6\\
\textbf{4.} & \symking{}e5 & \symknight{}c5\\
\textbf{5.} & \symking{}d4 & \symking{}a3\\
\textbf{6.} & \symking{}c3 & \symking{}a4\\
\textbf{7.} & \symking{}b2 & \symking{}b4\\
\textbf{8.} & \symking{}a1 & \symking{}a3\\
\textbf{9.} & \symbishop{}b1 & \symknight{}b3\#\\
\end{tabular}\par\vspace{3mm}
\qrcode[height=12mm]{https://helpman.komtera.lt/?fen=7Q/7k/4b1b1/8/8/8/2N5/1K6&moves=9}\\{\sffamily\tiny\color{ink}scannen zum Lösen}
\end{minipage}
\vspace{2mm}

{\raggedright\noindent\textcolor{ink}{\sffamily\textsc{Themen}}\enspace \themetag{reines Matt} \themetag{Mustermatt} \themetag{Idealmatt} \themetag{Rundlauf} \themetag{Selbstblock} \themetag{ohne Schach}\par\smallskip{\footnotesize\textit{Hinweis:} die Partei am Zug steht im Diagramm im Schach; die Lösung beginnt mit einem Schlag.}\par}

\clearpage
\phantomsection\label{hm:KQNvknp}
\begin{tcolorbox}[enhanced,colback=ink,colframe=ink,arc=0pt,boxrule=0pt,left=4mm,right=4mm,top=2.5mm,bottom=2.5mm]
\color{white}\sffamily{\LARGE\bfseries KQNvknp}\hfill{\Large h\#9}\\[0.5mm]
{\small 6 Steine\quad Schwarz am Zug}\hfill{\small Nr.~324\qquad\href{https://helpman.komtera.lt/?fen=N5Qk/8/8/8/8/1p6/8/Kn6&moves=9}{\color{white}\bfseries lösen\,\ensuremath{\nearrow}}}
\end{tcolorbox}
\vspace{2mm}
\noindent\begin{minipage}[t]{36mm}\vspace{0pt}\sffamily\small\raggedright
\hmfigure{längste eindeutige}{{\Large h\#9}}\hmfigure{eindeutig bei dieser Länge}{{\Large 1} Stellung}\hmfigure{längstes Matt}{{\Large h\#10.5}}\hmfigure{in der Klasse}{45,4 Milliarden Stellungen,\\ 27,5 Milliarden mit Hilfsmatt}
\end{minipage}\hfill
\begin{minipage}[t]{80mm}\vspace{0pt}\begin{tikzpicture}\begin{axis}[width=72mm,height=42mm,ybar,bar width=1.6pt,xmin=-1,xmax=22,ymin=0,ymax=9.5,axis y line*=left,axis x line*=bottom,xlabel={\scriptsize Halbzüge bis Matt},ylabel={\scriptsize Stellungen (Mrd.)},tick label style={font=\tiny},label style={font=\scriptsize},enlarge x limits=false,ylabel shift=-3pt,xlabel shift=-3pt]\addplot[fill=ink!55,draw=none] coordinates {(0,0.04114) (1,0.303) (2,1.116) (3,3.674) (4,8.636) (5,6.382) (6,5.885) (7,0.9411) (8,0.3005) (9,0.006537) (10,0.03819) (11,4.91e-07) (12,0.06797) (13,2.687e-06) (14,0.07001) (15,2.736e-06) (16,0.03404) (17,3.064e-06) (18,0.004529) (19,2.299e-06) (20,8.698e-05) (21,9.33e-07)};\draw[densely dotted] (axis cs:18,0) -- (axis cs:18,9.5);\end{axis}\begin{axis}[width=72mm,height=42mm,axis y line*=right,axis x line=none,ymode=log,xmin=-1,xmax=22,ylabel={\scriptsize eindeutig (log)},tick label style={font=\tiny},label style={font=\scriptsize},ylabel shift=-3pt]\addplot[mark=*,mark size=0.7pt,draw=black,thin] coordinates {(0,41139789) (1,245488686) (2,630694715) (3,698174319) (4,711012801) (5,84981406) (6,21286507) (7,351290) (8,321575) (9,10) (10,180607) (11,32) (12,85097) (13,14) (14,8739) (15,12) (16,912) (17,2) (18,1)};\end{axis}\end{tikzpicture}\end{minipage}
\vspace{3mm}

\noindent\begin{minipage}[t]{60mm}\vspace{0pt}\sbox\hmboard{\chessboard[setfen=N5Qk/8/8/8/8/1p6/8/Kn6 b - - 0 1]}\usebox\hmboard\\[1mm]
\makebox[\wd\hmboard]{\sffamily\hspace{3mm}\textbf{h\#9}\hfill(3+3)}\end{minipage}\hfill
\begin{minipage}[t]{52mm}\vspace{0pt}
\textcolor{ink}{\sffamily\textsc{Lösung}}\par\smallskip\small
\begin{tabular}{@{}r@{\enspace}l@{\enspace}l@{}}
\textbf{1.} & \symking{}xg8 & \symking{}b2\\
\textbf{2.} & \symknight{}c3 & \symking{}a3\\
\textbf{3.} & \symking{}f7 & \symking{}b4\\
\textbf{4.} & \symking{}e6 & \symking{}a5\\
\textbf{5.} & \symking{}d5 & \symking{}b6\\
\textbf{6.} & \symking{}c4 & \symking{}a6\\
\textbf{7.} & \symking{}b4 & \symknight{}b6\\
\textbf{8.} & \symking{}a3 & \symking{}a5\\
\textbf{9.} & \symknight{}a2 & \symknight{}c4\#\\
\end{tabular}\par\vspace{3mm}
\qrcode[height=12mm]{https://helpman.komtera.lt/?fen=N5Qk/8/8/8/8/1p6/8/Kn6&moves=9}\\{\sffamily\tiny\color{ink}scannen zum Lösen}
\end{minipage}
\vspace{2mm}

{\raggedright\noindent\textcolor{ink}{\sffamily\textsc{Themen}}\enspace \themetag{reines Matt} \themetag{Mustermatt} \themetag{Idealmatt} \themetag{Rundlauf} \themetag{Selbstblock} \themetag{ohne Schach}\par\smallskip{\footnotesize\textit{Hinweis:} die Partei am Zug steht im Diagramm im Schach; die Lösung beginnt mit einem Schlag.}\par}

\clearpage
\phantomsection\label{hm:KQNvkqb}
\begin{tcolorbox}[enhanced,colback=ink,colframe=ink,arc=0pt,boxrule=0pt,left=4mm,right=4mm,top=2.5mm,bottom=2.5mm]
\color{white}\sffamily{\LARGE\bfseries KQNvkqb}\hfill{\Large h\#9}\\[0.5mm]
{\small 6 Steine\quad Schwarz am Zug}\hfill{\small Nr.~325\qquad\href{https://helpman.komtera.lt/?fen=8/8/6Qk/7b/8/8/2q5/K6N&moves=9}{\color{white}\bfseries lösen\,\ensuremath{\nearrow}}}
\end{tcolorbox}
\vspace{2mm}
\noindent\begin{minipage}[t]{36mm}\vspace{0pt}\sffamily\small\raggedright
\hmfigure{längste eindeutige}{{\Large h\#9}}\hmfigure{eindeutig bei dieser Länge}{{\Large 111} Stellungen}\hmfigure{längstes Matt}{{\Large h\#12}}\hmfigure{in der Klasse}{15,5 Milliarden Stellungen,\\ 7,28 Milliarden mit Hilfsmatt}
\end{minipage}\hfill
\begin{minipage}[t]{80mm}\vspace{0pt}\begin{tikzpicture}\begin{axis}[width=72mm,height=42mm,ybar,bar width=1.6pt,xmin=-1,xmax=25,ymin=0,ymax=2.533,axis y line*=left,axis x line*=bottom,xlabel={\scriptsize Halbzüge bis Matt},ylabel={\scriptsize Stellungen (Mrd.)},tick label style={font=\tiny},label style={font=\scriptsize},enlarge x limits=false,ylabel shift=-3pt,xlabel shift=-3pt]\addplot[fill=ink!55,draw=none] coordinates {(0,0.00503) (1,0.04419) (2,0.2232) (3,0.9589) (4,2.035) (5,2.302) (6,1.18) (7,0.4341) (8,0.04109) (9,0.00506) (10,0.00452) (11,5.1e-07) (12,0.007779) (13,2.067e-06) (14,0.01123) (15,9.17e-07) (16,0.0133) (17,1.352e-06) (18,0.01048) (19,8.81e-07) (20,0.001299) (21,1.208e-06) (22,3.274e-06) (23,1e-08) (24,4e-09)};\draw[densely dotted] (axis cs:18,0) -- (axis cs:18,2.533);\end{axis}\begin{axis}[width=72mm,height=42mm,axis y line*=right,axis x line=none,ymode=log,xmin=-1,xmax=25,ylabel={\scriptsize eindeutig (log)},tick label style={font=\tiny},label style={font=\scriptsize},ylabel shift=-3pt]\addplot[mark=*,mark size=0.7pt,draw=black,thin] coordinates {(0,5030078) (1,40012444) (2,128297850) (3,202873316) (4,153847481) (5,42053624) (6,3154279) (7,327362) (8,41524) (9,6) (10,28054) (11,7) (12,25104) (14,5874) (16,3044) (18,111)};\end{axis}\end{tikzpicture}\end{minipage}
\vspace{3mm}

\noindent\begin{minipage}[t]{60mm}\vspace{0pt}\sbox\hmboard{\chessboard[setfen=8/8/6Qk/7b/8/8/2q5/K6N b - - 0 1]}\usebox\hmboard\\[1mm]
\makebox[\wd\hmboard]{\sffamily\hspace{3mm}\textbf{h\#9}\hfill(3+3)}\end{minipage}\hfill
\begin{minipage}[t]{52mm}\vspace{0pt}
\textcolor{ink}{\sffamily\textsc{Lösung}}\par\smallskip\small
\begin{tabular}{@{}r@{\enspace}l@{\enspace}l@{}}
\textbf{1.} & \symbishop{}xg6 & \symknight{}g3\\
\textbf{2.} & \symking{}g5 & \symknight{}f5\\
\textbf{3.} & \symking{}f4 & \symknight{}d4\\
\textbf{4.} & \symking{}e3 & \symknight{}xc2+\\
\textbf{5.} & \symking{}d2 & \symking{}a2\\
\textbf{6.} & \symking{}c1 & \symking{}b3\\
\textbf{7.} & \symking{}b1 & \symknight{}d4\\
\textbf{8.} & \symking{}a1 & \symking{}a3\\
\textbf{9.} & \symbishop{}b1 & \symknight{}b3\#\\
\end{tabular}\par\vspace{3mm}
\qrcode[height=12mm]{https://helpman.komtera.lt/?fen=8/8/6Qk/7b/8/8/2q5/K6N&moves=9}\\{\sffamily\tiny\color{ink}scannen zum Lösen}
\end{minipage}
\vspace{2mm}

{\raggedright\noindent\textcolor{ink}{\sffamily\textsc{Themen}}\enspace \themetag{reines Matt} \themetag{Mustermatt} \themetag{Idealmatt} \themetag{Switchback} \themetag{Selbstblock}\par\smallskip{\footnotesize\textit{Hinweis:} die Partei am Zug steht im Diagramm im Schach; die Lösung beginnt mit einem Schlag.}\par}

\clearpage
\phantomsection\label{hm:KQNvkrp}
\begin{tcolorbox}[enhanced,colback=ink,colframe=ink,arc=0pt,boxrule=0pt,left=4mm,right=4mm,top=2.5mm,bottom=2.5mm]
\color{white}\sffamily{\LARGE\bfseries KQNvkrp}\hfill{\Large h\#9}\\[0.5mm]
{\small 6 Steine\quad Schwarz am Zug}\hfill{\small Nr.~326\qquad\href{https://helpman.komtera.lt/?fen=6Qk/8/8/1p6/8/8/1r6/K6N&moves=9}{\color{white}\bfseries lösen\,\ensuremath{\nearrow}}}
\end{tcolorbox}
\vspace{2mm}
\noindent\begin{minipage}[t]{36mm}\vspace{0pt}\sffamily\small\raggedright
\hmfigure{längste eindeutige}{{\Large h\#9}}\hmfigure{eindeutig bei dieser Länge}{{\Large 7} Stellungen}\hmfigure{längstes Matt}{{\Large h\#10.5}}\hmfigure{in der Klasse}{45,4 Milliarden Stellungen,\\ 25,5 Milliarden mit Hilfsmatt}
\end{minipage}\hfill
\begin{minipage}[t]{80mm}\vspace{0pt}\begin{tikzpicture}\begin{axis}[width=72mm,height=42mm,ybar,bar width=1.6pt,xmin=-1,xmax=22,ymin=0,ymax=9.802,axis y line*=left,axis x line*=bottom,xlabel={\scriptsize Halbzüge bis Matt},ylabel={\scriptsize Stellungen (Mrd.)},tick label style={font=\tiny},label style={font=\scriptsize},enlarge x limits=false,ylabel shift=-3pt,xlabel shift=-3pt]\addplot[fill=ink!55,draw=none] coordinates {(0,0.03169) (1,0.2539) (2,1.134) (3,3.856) (4,8.911) (5,6.562) (6,3.847) (7,0.6304) (8,0.07218) (9,0.0007812) (10,0.04521) (11,2.5e-07) (12,0.06768) (13,1.69e-07) (14,0.04497) (15,2.2e-07) (16,0.01697) (17,8e-09) (18,0.002069) (19,1e-09) (20,6.309e-06) (21,3e-09)};\draw[densely dotted] (axis cs:18,0) -- (axis cs:18,9.802);\end{axis}\begin{axis}[width=72mm,height=42mm,axis y line*=right,axis x line=none,ymode=log,xmin=-1,xmax=22,ylabel={\scriptsize eindeutig (log)},tick label style={font=\tiny},label style={font=\scriptsize},ylabel shift=-3pt]\addplot[mark=*,mark size=0.7pt,draw=black,thin] coordinates {(0,31691803) (1,216921101) (2,679655169) (3,699922518) (4,594038347) (5,85244747) (6,8513251) (7,69967) (8,426278) (9,57) (10,362984) (12,150441) (14,20544) (16,1906) (18,7)};\end{axis}\end{tikzpicture}\end{minipage}
\vspace{3mm}

\noindent\begin{minipage}[t]{60mm}\vspace{0pt}\sbox\hmboard{\chessboard[setfen=6Qk/8/8/1p6/8/8/1r6/K6N b - - 0 1]}\usebox\hmboard\\[1mm]
\makebox[\wd\hmboard]{\sffamily\hspace{3mm}\textbf{h\#9}\hfill(3+3)}\end{minipage}\hfill
\begin{minipage}[t]{52mm}\vspace{0pt}
\textcolor{ink}{\sffamily\textsc{Lösung}}\par\smallskip\small
\begin{tabular}{@{}r@{\enspace}l@{\enspace}l@{}}
\textbf{1.} & \symking{}xg8 & \symknight{}g3\\
\textbf{2.} & \symrook{}h2 & \symking{}b1\\
\textbf{3.} & \symrook{}h8 & \symking{}c2\\
\textbf{4.} & b4 & \symking{}d3\\
\textbf{5.} & b3 & \symking{}e4\\
\textbf{6.} & b2 & \symking{}f5\\
\textbf{7.} & b1=\symrook{} & \symking{}g6\\
\textbf{8.} & \symrook{}b8 & \symknight{}f5\\
\textbf{9.} & \symrook{}f8 & \symknight{}e7\#\\
\end{tabular}\par\vspace{3mm}
\qrcode[height=12mm]{https://helpman.komtera.lt/?fen=6Qk/8/8/1p6/8/8/1r6/K6N&moves=9}\\{\sffamily\tiny\color{ink}scannen zum Lösen}
\end{minipage}
\vspace{2mm}

{\raggedright\noindent\textcolor{ink}{\sffamily\textsc{Themen}}\enspace \themetag{reines Matt} \themetag{Mustermatt} \themetag{Idealmatt} \themetag{Umwandlung} \themetag{Unterverwandlung} \themetag{Selbstblock} \themetag{Kniest} \themetag{ohne Schach}\par\smallskip{\footnotesize\textit{Hinweis:} die Partei am Zug steht im Diagramm im Schach; die Lösung beginnt mit einem Schlag.}\par}

\clearpage
\phantomsection\label{hm:KQPvknp}
\begin{tcolorbox}[enhanced,colback=ink,colframe=ink,arc=0pt,boxrule=0pt,left=4mm,right=4mm,top=2.5mm,bottom=2.5mm]
\color{white}\sffamily{\LARGE\bfseries KQPvknp}\hfill{\Large h\#9}\\[0.5mm]
{\small 6 Steine\quad Schwarz am Zug}\hfill{\small Nr.~327\qquad\href{https://helpman.komtera.lt/?fen=6kn/5Q2/8/8/2p5/8/2P5/K7&moves=9}{\color{white}\bfseries lösen\,\ensuremath{\nearrow}}}
\end{tcolorbox}
\vspace{2mm}
\noindent\begin{minipage}[t]{36mm}\vspace{0pt}\sffamily\small\raggedright
\hmfigure{längste eindeutige}{{\Large h\#9}}\hmfigure{eindeutig bei dieser Länge}{{\Large 91} Stellungen}\hmfigure{längstes Matt}{{\Large h\#15.5}}\hmfigure{in der Klasse}{34,1 Milliarden Stellungen,\\ 21 Milliarden mit Hilfsmatt}
\end{minipage}\hfill
\begin{minipage}[t]{80mm}\vspace{0pt}\begin{tikzpicture}\begin{axis}[width=72mm,height=42mm,ybar,bar width=1.6pt,xmin=-1,xmax=32,ymin=0,ymax=5.855,axis y line*=left,axis x line*=bottom,xlabel={\scriptsize Halbzüge bis Matt},ylabel={\scriptsize Stellungen (Mrd.)},tick label style={font=\tiny},label style={font=\scriptsize},enlarge x limits=false,ylabel shift=-3pt,xlabel shift=-3pt]\addplot[fill=ink!55,draw=none] coordinates {(0,0.02086) (1,0.155) (2,0.4782) (3,1.772) (4,4.15) (5,4.184) (6,5.323) (7,2.427) (8,1.865) (9,0.3812) (10,0.1751) (11,0.007939) (12,0.04331) (13,7.022e-06) (14,0.03081) (15,2.92e-06) (16,0.004512) (17,3.94e-07) (18,0.0006401) (19,2.9e-08) (20,0.000126) (21,9.2e-08) (22,1.359e-05) (23,1.76e-07) (24,4.18e-07) (25,8.6e-08) (26,2.3e-08) (27,5.7e-08) (28,1e-08) (29,3.8e-08) (30,4e-09) (31,6e-09)};\draw[densely dotted] (axis cs:18,0) -- (axis cs:18,5.855);\end{axis}\begin{axis}[width=72mm,height=42mm,axis y line*=right,axis x line=none,ymode=log,xmin=-1,xmax=32,ylabel={\scriptsize eindeutig (log)},tick label style={font=\tiny},label style={font=\scriptsize},ylabel shift=-3pt]\addplot[mark=*,mark size=0.7pt,draw=black,thin] coordinates {(0,20858903) (1,104675044) (2,214890424) (3,295122048) (4,358675780) (5,144885957) (6,90974454) (7,17700459) (8,6175062) (9,187933) (10,240666) (11,131) (12,180008) (13,10) (14,29439) (16,1260) (18,91)};\end{axis}\end{tikzpicture}\end{minipage}
\vspace{3mm}

\noindent\begin{minipage}[t]{60mm}\vspace{0pt}\sbox\hmboard{\chessboard[setfen=6kn/5Q2/8/8/2p5/8/2P5/K7 b - - 0 1]}\usebox\hmboard\\[1mm]
\makebox[\wd\hmboard]{\sffamily\hspace{3mm}\textbf{h\#9}\hfill(3+3)}\end{minipage}\hfill
\begin{minipage}[t]{52mm}\vspace{0pt}
\textcolor{ink}{\sffamily\textsc{Lösung}}\par\smallskip\small
\begin{tabular}{@{}r@{\enspace}l@{\enspace}l@{}}
\textbf{1.} & \symknight{}xf7 & \symking{}b2\\
\textbf{2.} & \symknight{}e5 & \symking{}c3\\
\textbf{3.} & \symknight{}d3 & cxd3\\
\textbf{4.} & \symking{}f7 & dxc4\\
\textbf{5.} & \symking{}e6 & c5\\
\textbf{6.} & \symking{}d5 & c6\\
\textbf{7.} & \symking{}c5 & c7\\
\textbf{8.} & \symking{}b5 & c8=\symqueen{}\\
\textbf{9.} & \symking{}a4 & \symqueen{}a6\#\\
\end{tabular}\par\vspace{3mm}
\qrcode[height=12mm]{https://helpman.komtera.lt/?fen=6kn/5Q2/8/8/2p5/8/2P5/K7&moves=9}\\{\sffamily\tiny\color{ink}scannen zum Lösen}
\end{minipage}
\vspace{2mm}

{\raggedright\noindent\textcolor{ink}{\sffamily\textsc{Themen}}\enspace \themetag{reines Matt} \themetag{Mustermatt} \themetag{Idealmatt} \themetag{Spiegelmatt} \themetag{Umwandlung} \themetag{Excelsior} \themetag{Phönix} \themetag{ohne Schach} \themetag{Umnov}\par\smallskip{\footnotesize\textit{Hinweis:} die Partei am Zug steht im Diagramm im Schach; die Lösung beginnt mit einem Schlag.}\par}

\clearpage
\phantomsection\label{hm:KQPvkrp}
\begin{tcolorbox}[enhanced,colback=ink,colframe=ink,arc=0pt,boxrule=0pt,left=4mm,right=4mm,top=2.5mm,bottom=2.5mm]
\color{white}\sffamily{\LARGE\bfseries KQPvkrp}\hfill{\Large h\#9}\\[0.5mm]
{\small 6 Steine\quad Schwarz am Zug}\hfill{\small Nr.~328\qquad\href{https://helpman.komtera.lt/?fen=8/8/8/8/6Qk/5p2/2r2P2/K7&moves=9}{\color{white}\bfseries lösen\,\ensuremath{\nearrow}}}
\end{tcolorbox}
\vspace{2mm}
\noindent\begin{minipage}[t]{36mm}\vspace{0pt}\sffamily\small\raggedright
\hmfigure{längste eindeutige}{{\Large h\#9}}\hmfigure{eindeutig bei dieser Länge}{{\Large 71} Stellungen}\hmfigure{längstes Matt}{{\Large h\#14}}\hmfigure{in der Klasse}{34,1 Milliarden Stellungen,\\ 19,5 Milliarden mit Hilfsmatt}
\end{minipage}\hfill
\begin{minipage}[t]{80mm}\vspace{0pt}\begin{tikzpicture}\begin{axis}[width=72mm,height=42mm,ybar,bar width=1.6pt,xmin=-1,xmax=29,ymin=0,ymax=5.444,axis y line*=left,axis x line*=bottom,xlabel={\scriptsize Halbzüge bis Matt},ylabel={\scriptsize Stellungen (Mrd.)},tick label style={font=\tiny},label style={font=\scriptsize},enlarge x limits=false,ylabel shift=-3pt,xlabel shift=-3pt]\addplot[fill=ink!55,draw=none] coordinates {(0,0.01577) (1,0.1286) (2,0.5169) (3,1.97) (4,4.949) (5,4.8) (6,4.326) (7,1.887) (8,0.6838) (9,0.1379) (10,0.04541) (11,0.0003252) (12,0.03591) (13,2.576e-06) (14,0.01426) (15,8.5e-08) (16,0.001648) (17,2e-09) (18,0.0003107) (20,6.159e-06) (21,2.3e-08) (22,7.08e-07) (23,3e-08) (24,4.28e-07) (25,1.6e-08) (26,1.45e-07) (28,7.8e-08)};\draw[densely dotted] (axis cs:18,0) -- (axis cs:18,5.444);\end{axis}\begin{axis}[width=72mm,height=42mm,axis y line*=right,axis x line=none,ymode=log,xmin=-1,xmax=29,ylabel={\scriptsize eindeutig (log)},tick label style={font=\tiny},label style={font=\scriptsize},ylabel shift=-3pt]\addplot[mark=*,mark size=0.7pt,draw=black,thin] coordinates {(0,15772826) (1,96979702) (2,275387931) (3,320790210) (4,358774980) (5,114585986) (6,35473773) (7,6538633) (8,1373480) (9,36759) (10,204282) (11,41) (12,95145) (14,8044) (16,1424) (18,71)};\end{axis}\end{tikzpicture}\end{minipage}
\vspace{3mm}

\noindent\begin{minipage}[t]{60mm}\vspace{0pt}\sbox\hmboard{\chessboard[setfen=8/8/8/8/6Qk/5p2/2r2P2/K7 b - - 0 1]}\usebox\hmboard\\[1mm]
\makebox[\wd\hmboard]{\sffamily\hspace{3mm}\textbf{h\#9}\hfill(3+3)}\end{minipage}\hfill
\begin{minipage}[t]{52mm}\vspace{0pt}
\textcolor{ink}{\sffamily\textsc{Lösung}}\par\smallskip\small
\begin{tabular}{@{}r@{\enspace}l@{\enspace}l@{}}
\textbf{1.} & \symking{}xg4 & \symking{}b1\\
\textbf{2.} & \symrook{}e2 & \symking{}c1\\
\textbf{3.} & \symrook{}e3 & fxe3\\
\textbf{4.} & \symking{}f5 & e4+\\
\textbf{5.} & \symking{}e6 & e5\\
\textbf{6.} & \symking{}d5 & e6\\
\textbf{7.} & \symking{}c4 & e7\\
\textbf{8.} & \symking{}b3 & e8=\symqueen{}\\
\textbf{9.} & \symking{}a2 & \symqueen{}a4\#\\
\end{tabular}\par\vspace{3mm}
\qrcode[height=12mm]{https://helpman.komtera.lt/?fen=8/8/8/8/6Qk/5p2/2r2P2/K7&moves=9}\\{\sffamily\tiny\color{ink}scannen zum Lösen}
\end{minipage}
\vspace{2mm}

{\raggedright\noindent\textcolor{ink}{\sffamily\textsc{Themen}}\enspace \themetag{reines Matt} \themetag{Mustermatt} \themetag{Spiegelmatt} \themetag{Umwandlung} \themetag{Excelsior} \themetag{Phönix} \themetag{Umnov}\par\smallskip{\footnotesize\textit{Hinweis:} die Partei am Zug steht im Diagramm im Schach; die Lösung beginnt mit einem Schlag.}\par}

\clearpage
\phantomsection\label{hm:KRBBvkp}
\begin{tcolorbox}[enhanced,colback=ink,colframe=ink,arc=0pt,boxrule=0pt,left=4mm,right=4mm,top=2.5mm,bottom=2.5mm]
\color{white}\sffamily{\LARGE\bfseries KRBBvkp}\hfill{\Large h\#9}\\[0.5mm]
{\small 6 Steine\quad Schwarz am Zug}\hfill{\small Nr.~329\qquad\href{https://helpman.komtera.lt/?fen=8/8/8/3B4/8/8/3Kp3/kR1B4&moves=9}{\color{white}\bfseries lösen\,\ensuremath{\nearrow}}}
\end{tcolorbox}
\vspace{2mm}
\noindent\begin{minipage}[t]{36mm}\vspace{0pt}\sffamily\small\raggedright
\hmfigure{längste eindeutige}{{\Large h\#9}}\hmfigure{eindeutig bei dieser Länge}{{\Large 4} Stellungen}\hmfigure{längstes Matt}{{\Large h\#16}}\hmfigure{in der Klasse}{45,4 Milliarden Stellungen,\\ 14,4 Milliarden mit Hilfsmatt}
\end{minipage}\hfill
\begin{minipage}[t]{80mm}\vspace{0pt}\begin{tikzpicture}\begin{axis}[width=72mm,height=42mm,ybar,bar width=1.6pt,xmin=-1,xmax=33,ymin=0,ymax=4.31,axis y line*=left,axis x line*=bottom,xlabel={\scriptsize Halbzüge bis Matt},ylabel={\scriptsize Stellungen (Mrd.)},tick label style={font=\tiny},label style={font=\scriptsize},enlarge x limits=false,ylabel shift=-3pt,xlabel shift=-3pt]\addplot[fill=ink!55,draw=none] coordinates {(0,0.02613) (1,0.1011) (2,0.3304) (3,1.254) (4,3.574) (5,3.009) (6,3.918) (7,1.046) (8,0.8875) (9,0.1326) (10,0.06825) (11,0.004801) (12,0.002836) (13,3.936e-05) (14,0.0005576) (15,1.504e-06) (16,0.0006646) (17,1.745e-06) (18,0.000836) (19,1.847e-06) (20,0.0007782) (21,1.967e-06) (22,0.0008132) (23,3.323e-06) (24,0.0006856) (25,1.237e-06) (26,0.000344) (27,4.71e-07) (28,0.0003062) (29,5.75e-07) (30,3.263e-05) (31,3e-09) (32,2e-09)};\draw[densely dotted] (axis cs:18,0) -- (axis cs:18,4.31);\end{axis}\begin{axis}[width=72mm,height=42mm,axis y line*=right,axis x line=none,ymode=log,xmin=-1,xmax=33,ylabel={\scriptsize eindeutig (log)},tick label style={font=\tiny},label style={font=\scriptsize},ylabel shift=-3pt]\addplot[mark=*,mark size=0.7pt,draw=black,thin] coordinates {(0,26128020) (1,92708667) (2,252493535) (3,574322270) (4,783052584) (5,170720338) (6,112697847) (7,6303077) (8,2653607) (9,99277) (10,39360) (11,188) (12,1367) (13,7) (14,171) (15,1) (16,5) (18,4)};\end{axis}\end{tikzpicture}\end{minipage}
\vspace{3mm}

\noindent\begin{minipage}[t]{60mm}\vspace{0pt}\sbox\hmboard{\chessboard[setfen=8/8/8/3B4/8/8/3Kp3/kR1B4 b - - 0 1]}\usebox\hmboard\\[1mm]
\makebox[\wd\hmboard]{\sffamily\hspace{3mm}\textbf{h\#9}\hfill(4+2)}\end{minipage}\hfill
\begin{minipage}[t]{52mm}\vspace{0pt}
\textcolor{ink}{\sffamily\textsc{Lösung}}\par\smallskip\small
\begin{tabular}{@{}r@{\enspace}l@{\enspace}l@{}}
\textbf{1.} & \symking{}xb1 & \symking{}e3\\
\textbf{2.} & \symking{}c1 & \symking{}f4\\
\textbf{3.} & \symking{}d2 & \symbishop{}g2\\
\textbf{4.} & \symking{}e1 & \symbishop{}f1\\
\textbf{5.} & exf1=\symknight{} & \symbishop{}g4\\
\textbf{6.} & \symking{}f2 & \symbishop{}h3\\
\textbf{7.} & \symking{}g1 & \symking{}f3\\
\textbf{8.} & \symking{}h1 & \symking{}f2\\
\textbf{9.} & \symknight{}h2 & \symbishop{}g2\#\\
\end{tabular}\par\vspace{3mm}
\qrcode[height=12mm]{https://helpman.komtera.lt/?fen=8/8/8/3B4/8/8/3Kp3/kR1B4&moves=9}\\{\sffamily\tiny\color{ink}scannen zum Lösen}
\end{minipage}
\vspace{2mm}

{\raggedright\noindent\textcolor{ink}{\sffamily\textsc{Themen}}\enspace \themetag{reines Matt} \themetag{Mustermatt} \themetag{Idealmatt} \themetag{Umwandlung} \themetag{Unterverwandlung} \themetag{Selbstblock} \themetag{ohne Schach}\par\smallskip{\footnotesize\textit{Hinweis:} die Partei am Zug steht im Diagramm im Schach; die Lösung beginnt mit einem Schlag.}\par}

\clearpage
\phantomsection\label{hm:KRBvkbn}
\begin{tcolorbox}[enhanced,colback=ink,colframe=ink,arc=0pt,boxrule=0pt,left=4mm,right=4mm,top=2.5mm,bottom=2.5mm]
\color{white}\sffamily{\LARGE\bfseries KRBvkbn}\hfill{\Large h\#9}\\[0.5mm]
{\small 6 Steine\quad Schwarz am Zug}\hfill{\small Nr.~330\qquad\href{https://helpman.komtera.lt/?fen=8/8/1n6/7B/8/1b6/2R5/K1k5&moves=9}{\color{white}\bfseries lösen\,\ensuremath{\nearrow}}}
\end{tcolorbox}
\vspace{2mm}
\noindent\begin{minipage}[t]{36mm}\vspace{0pt}\sffamily\small\raggedright
\hmfigure{längste eindeutige}{{\Large h\#9}}\hmfigure{eindeutig bei dieser Länge}{{\Large 37} Stellungen}\hmfigure{längstes Matt}{{\Large h\#12}}\hmfigure{in der Klasse}{15,5 Milliarden Stellungen,\\ 9,09 Milliarden mit Hilfsmatt}
\end{minipage}\hfill
\begin{minipage}[t]{80mm}\vspace{0pt}\begin{tikzpicture}\begin{axis}[width=72mm,height=42mm,ybar,bar width=1.6pt,xmin=-1,xmax=25,ymin=0,ymax=2.9,axis y line*=left,axis x line*=bottom,xlabel={\scriptsize Halbzüge bis Matt},ylabel={\scriptsize Stellungen (Mrd.)},tick label style={font=\tiny},label style={font=\scriptsize},enlarge x limits=false,ylabel shift=-3pt,xlabel shift=-3pt]\addplot[fill=ink!55,draw=none] coordinates {(0,0.003627) (1,0.01764) (2,0.0624) (3,0.2997) (4,0.9145) (5,1.808) (6,2.636) (7,1.698) (8,1.071) (9,0.3948) (10,0.151) (11,0.02259) (12,0.002985) (13,2.203e-05) (14,0.00066) (15,3.1e-08) (16,0.0008138) (17,6.3e-08) (18,0.0009227) (19,5e-09) (20,0.0005427) (22,0.0002281) (24,2.125e-05)};\draw[densely dotted] (axis cs:18,0) -- (axis cs:18,2.9);\end{axis}\begin{axis}[width=72mm,height=42mm,axis y line*=right,axis x line=none,ymode=log,xmin=-1,xmax=25,ylabel={\scriptsize eindeutig (log)},tick label style={font=\tiny},label style={font=\scriptsize},ylabel shift=-3pt]\addplot[mark=*,mark size=0.7pt,draw=black,thin] coordinates {(0,3627395) (1,17267926) (2,33005057) (3,109111998) (4,261741146) (5,159807716) (6,71898128) (7,13728842) (8,4934381) (9,666970) (10,71888) (11,367) (12,1709) (14,385) (16,37) (18,37)};\end{axis}\end{tikzpicture}\end{minipage}
\vspace{3mm}

\noindent\begin{minipage}[t]{60mm}\vspace{0pt}\sbox\hmboard{\chessboard[setfen=8/8/1n6/7B/8/1b6/2R5/K1k5 b - - 0 1]}\usebox\hmboard\\[1mm]
\makebox[\wd\hmboard]{\sffamily\hspace{3mm}\textbf{h\#9}\hfill(3+3)}\end{minipage}\hfill
\begin{minipage}[t]{52mm}\vspace{0pt}
\textcolor{ink}{\sffamily\textsc{Lösung}}\par\smallskip\small
\begin{tabular}{@{}r@{\enspace}l@{\enspace}l@{}}
\textbf{1.} & \symking{}xc2 & \symbishop{}e2\\
\textbf{2.} & \symking{}c3 & \symbishop{}a6\\
\textbf{3.} & \symking{}b4 & \symking{}b2\\
\textbf{4.} & \symking{}a5 & \symking{}xb3\\
\textbf{5.} & \symknight{}c8 & \symking{}c4\\
\textbf{6.} & \symking{}b6 & \symking{}d5\\
\textbf{7.} & \symking{}a7 & \symking{}c6\\
\textbf{8.} & \symking{}a8 & \symking{}c7\\
\textbf{9.} & \symknight{}a7 & \symbishop{}b7\#\\
\end{tabular}\par\vspace{3mm}
\qrcode[height=12mm]{https://helpman.komtera.lt/?fen=8/8/1n6/7B/8/1b6/2R5/K1k5&moves=9}\\{\sffamily\tiny\color{ink}scannen zum Lösen}
\end{minipage}
\vspace{2mm}

{\raggedright\noindent\textcolor{ink}{\sffamily\textsc{Themen}}\enspace \themetag{reines Matt} \themetag{Mustermatt} \themetag{Idealmatt} \themetag{Selbstblock} \themetag{ohne Schach}\par\smallskip{\footnotesize\textit{Hinweis:} die Partei am Zug steht im Diagramm im Schach; die Lösung beginnt mit einem Schlag.}\par}

\clearpage
\phantomsection\label{hm:KRBvknn}
\begin{tcolorbox}[enhanced,colback=ink,colframe=ink,arc=0pt,boxrule=0pt,left=4mm,right=4mm,top=2.5mm,bottom=2.5mm]
\color{white}\sffamily{\LARGE\bfseries KRBvknn}\hfill{\Large h\#9}\\[0.5mm]
{\small 6 Steine\quad Schwarz am Zug}\hfill{\small Nr.~331\qquad\href{https://helpman.komtera.lt/?fen=8/8/8/8/8/n2nB3/7R/K6k&moves=9}{\color{white}\bfseries lösen\,\ensuremath{\nearrow}}}
\end{tcolorbox}
\vspace{2mm}
\noindent\begin{minipage}[t]{36mm}\vspace{0pt}\sffamily\small\raggedright
\hmfigure{längste eindeutige}{{\Large h\#9}}\hmfigure{eindeutig bei dieser Länge}{{\Large 35} Stellungen}\hmfigure{längstes Matt}{{\Large h\#12}}\hmfigure{in der Klasse}{15,5 Milliarden Stellungen,\\ 4,66 Milliarden mit Hilfsmatt}
\end{minipage}\hfill
\begin{minipage}[t]{80mm}\vspace{0pt}\begin{tikzpicture}\begin{axis}[width=72mm,height=42mm,ybar,bar width=1.6pt,xmin=-1,xmax=25,ymin=0,ymax=1.329,axis y line*=left,axis x line*=bottom,xlabel={\scriptsize Halbzüge bis Matt},ylabel={\scriptsize Stellungen (Mrd.)},tick label style={font=\tiny},label style={font=\scriptsize},enlarge x limits=false,ylabel shift=-3pt,xlabel shift=-3pt]\addplot[fill=ink!55,draw=none] coordinates {(0,0.002128) (1,0.01076) (2,0.03136) (3,0.1334) (4,0.3393) (5,0.6993) (6,1.208) (7,0.9535) (8,0.8548) (9,0.3088) (10,0.1057) (11,0.01451) (12,0.001308) (13,7.116e-06) (14,0.0003703) (15,2e-09) (16,0.000403) (17,5.7e-08) (18,0.000534) (19,1.99e-07) (20,0.0003779) (21,9e-08) (22,8.105e-05) (23,2e-09) (24,9.9e-07)};\draw[densely dotted] (axis cs:18,0) -- (axis cs:18,1.329);\end{axis}\begin{axis}[width=72mm,height=42mm,axis y line*=right,axis x line=none,ymode=log,xmin=-1,xmax=25,ylabel={\scriptsize eindeutig (log)},tick label style={font=\tiny},label style={font=\scriptsize},ylabel shift=-3pt]\addplot[mark=*,mark size=0.7pt,draw=black,thin] coordinates {(0,2128470) (1,10544201) (2,15936016) (3,44999079) (4,92896733) (5,80309944) (6,54911535) (7,10580604) (8,3683489) (9,522077) (10,43241) (11,187) (12,1324) (14,232) (16,69) (18,35)};\end{axis}\end{tikzpicture}\end{minipage}
\vspace{3mm}

\noindent\begin{minipage}[t]{60mm}\vspace{0pt}\sbox\hmboard{\chessboard[setfen=8/8/8/8/8/n2nB3/7R/K6k b - - 0 1]}\usebox\hmboard\\[1mm]
\makebox[\wd\hmboard]{\sffamily\hspace{3mm}\textbf{h\#9}\hfill(3+3)}\end{minipage}\hfill
\begin{minipage}[t]{52mm}\vspace{0pt}
\textcolor{ink}{\sffamily\textsc{Lösung}}\par\smallskip\small
\begin{tabular}{@{}r@{\enspace}l@{\enspace}l@{}}
\textbf{1.} & \symking{}xh2 & \symking{}a2\\
\textbf{2.} & \symking{}g2 & \symking{}b3\\
\textbf{3.} & \symking{}f1 & \symking{}c3\\
\textbf{4.} & \symking{}e1 & \symking{}xd3\\
\textbf{5.} & \symking{}d1 & \symking{}d4\\
\textbf{6.} & \symking{}c2 & \symbishop{}c1\\
\textbf{7.} & \symking{}b1 & \symking{}c3\\
\textbf{8.} & \symking{}a1 & \symking{}b3\\
\textbf{9.} & \symknight{}b1 & \symbishop{}b2\#\\
\end{tabular}\par\vspace{3mm}
\qrcode[height=12mm]{https://helpman.komtera.lt/?fen=8/8/8/8/8/n2nB3/7R/K6k&moves=9}\\{\sffamily\tiny\color{ink}scannen zum Lösen}
\end{minipage}
\vspace{2mm}

{\raggedright\noindent\textcolor{ink}{\sffamily\textsc{Themen}}\enspace \themetag{reines Matt} \themetag{Mustermatt} \themetag{Idealmatt} \themetag{Rundlauf} \themetag{Selbstblock} \themetag{ohne Schach}\par\smallskip{\footnotesize\textit{Hinweis:} die Partei am Zug steht im Diagramm im Schach; die Lösung beginnt mit einem Schlag.}\par}

\clearpage
\phantomsection\label{hm:KRBvkpp}
\begin{tcolorbox}[enhanced,colback=ink,colframe=ink,arc=0pt,boxrule=0pt,left=4mm,right=4mm,top=2.5mm,bottom=2.5mm]
\color{white}\sffamily{\LARGE\bfseries KRBvkpp}\hfill{\Large h\#9}\\[0.5mm]
{\small 6 Steine\quad Schwarz am Zug}\hfill{\small Nr.~332\qquad\href{https://helpman.komtera.lt/?fen=8/8/p7/8/8/8/1B5p/2K3Rk&moves=9}{\color{white}\bfseries lösen\,\ensuremath{\nearrow}}}
\end{tcolorbox}
\vspace{2mm}
\noindent\begin{minipage}[t]{36mm}\vspace{0pt}\sffamily\small\raggedright
\hmfigure{längste eindeutige}{{\Large h\#9}}\hmfigure{eindeutig bei dieser Länge}{{\Large 4} Stellungen}\hmfigure{längstes Matt}{{\Large h\#16}}\hmfigure{in der Klasse}{34,1 Milliarden Stellungen,\\ 11,2 Milliarden mit Hilfsmatt}
\end{minipage}\hfill
\begin{minipage}[t]{80mm}\vspace{0pt}\begin{tikzpicture}\begin{axis}[width=72mm,height=42mm,ybar,bar width=1.6pt,xmin=-1,xmax=33,ymin=0,ymax=3.546,axis y line*=left,axis x line*=bottom,xlabel={\scriptsize Halbzüge bis Matt},ylabel={\scriptsize Stellungen (Mrd.)},tick label style={font=\tiny},label style={font=\scriptsize},enlarge x limits=false,ylabel shift=-3pt,xlabel shift=-3pt]\addplot[fill=ink!55,draw=none] coordinates {(0,0.00835) (1,0.04135) (2,0.1044) (3,0.4064) (4,1.082) (5,1.84) (6,3.223) (7,1.958) (8,1.736) (9,0.5018) (10,0.2282) (11,0.02569) (12,0.008067) (13,0.0001598) (14,0.0009478) (15,2.31e-07) (16,0.001007) (17,2.33e-07) (18,0.0009336) (19,2.64e-07) (20,0.0007307) (21,1.2e-07) (22,0.0005528) (23,2.27e-07) (24,0.0003467) (25,1.15e-07) (26,0.0001754) (27,6.6e-08) (28,8.578e-05) (29,4.5e-08) (30,1.928e-05) (31,2.8e-08) (32,6.351e-06)};\draw[densely dotted] (axis cs:18,0) -- (axis cs:18,3.546);\end{axis}\begin{axis}[width=72mm,height=42mm,axis y line*=right,axis x line=none,ymode=log,xmin=-1,xmax=33,ylabel={\scriptsize eindeutig (log)},tick label style={font=\tiny},label style={font=\scriptsize},ylabel shift=-3pt]\addplot[mark=*,mark size=0.7pt,draw=black,thin] coordinates {(0,8349571) (1,40217257) (2,65142960) (3,165271071) (4,334747468) (5,212131949) (6,139856122) (7,17303222) (8,6926174) (9,488036) (10,120222) (11,1594) (12,1120) (14,159) (16,14) (18,4)};\end{axis}\end{tikzpicture}\end{minipage}
\vspace{3mm}

\noindent\begin{minipage}[t]{60mm}\vspace{0pt}\sbox\hmboard{\chessboard[setfen=8/8/p7/8/8/8/1B5p/2K3Rk b - - 0 1]}\usebox\hmboard\\[1mm]
\makebox[\wd\hmboard]{\sffamily\hspace{3mm}\textbf{h\#9}\hfill(3+3)}\end{minipage}\hfill
\begin{minipage}[t]{52mm}\vspace{0pt}
\textcolor{ink}{\sffamily\textsc{Lösung}}\par\smallskip\small
\begin{tabular}{@{}r@{\enspace}l@{\enspace}l@{}}
\textbf{1.} & \symking{}xg1 & \symking{}d2\\
\textbf{2.} & h1=\symbishop{} & \symking{}e3\\
\textbf{3.} & \symbishop{}a8 & \symking{}f4\\
\textbf{4.} & \symking{}f2 & \symking{}f5\\
\textbf{5.} & \symking{}e3 & \symbishop{}c1+\\
\textbf{6.} & \symking{}d4 & \symking{}e6\\
\textbf{7.} & \symking{}c5 & \symking{}d7\\
\textbf{8.} & \symking{}b6 & \symking{}c8\\
\textbf{9.} & \symking{}a7 & \symbishop{}e3\#\\
\end{tabular}\par\vspace{3mm}
\qrcode[height=12mm]{https://helpman.komtera.lt/?fen=8/8/p7/8/8/8/1B5p/2K3Rk&moves=9}\\{\sffamily\tiny\color{ink}scannen zum Lösen}
\end{minipage}
\vspace{2mm}

{\raggedright\noindent\textcolor{ink}{\sffamily\textsc{Themen}}\enspace \themetag{reines Matt} \themetag{Mustermatt} \themetag{Idealmatt} \themetag{Umwandlung} \themetag{Unterverwandlung} \themetag{Selbstblock}\par\smallskip{\footnotesize\textit{Hinweis:} die Partei am Zug steht im Diagramm im Schach; die Lösung beginnt mit einem Schlag.}\par}

\clearpage
\phantomsection\label{hm:KRNPvkp}
\begin{tcolorbox}[enhanced,colback=ink,colframe=ink,arc=0pt,boxrule=0pt,left=4mm,right=4mm,top=2.5mm,bottom=2.5mm]
\color{white}\sffamily{\LARGE\bfseries KRNPvkp}\hfill{\Large h\#9}\\[0.5mm]
{\small 6 Steine\quad Schwarz am Zug}\hfill{\small Nr.~333\qquad\href{https://helpman.komtera.lt/?fen=8/3K4/8/8/6N1/2p5/2P5/6Rk&moves=9}{\color{white}\bfseries lösen\,\ensuremath{\nearrow}}}
\end{tcolorbox}
\vspace{2mm}
\noindent\begin{minipage}[t]{36mm}\vspace{0pt}\sffamily\small\raggedright
\hmfigure{längste eindeutige}{{\Large h\#9}}\hmfigure{eindeutig bei dieser Länge}{{\Large 1} Stellung}\hmfigure{längstes Matt}{{\Large h\#12}}\hmfigure{in der Klasse}{34,1 Milliarden Stellungen,\\ 22,9 Milliarden mit Hilfsmatt}
\end{minipage}\hfill
\begin{minipage}[t]{80mm}\vspace{0pt}\begin{tikzpicture}\begin{axis}[width=72mm,height=42mm,ybar,bar width=1.6pt,xmin=-1,xmax=25,ymin=0,ymax=7.538,axis y line*=left,axis x line*=bottom,xlabel={\scriptsize Halbzüge bis Matt},ylabel={\scriptsize Stellungen (Mrd.)},tick label style={font=\tiny},label style={font=\scriptsize},enlarge x limits=false,ylabel shift=-3pt,xlabel shift=-3pt]\addplot[fill=ink!55,draw=none] coordinates {(0,0.01596) (1,0.1037) (2,0.2629) (3,1.197) (4,2.951) (5,4.567) (6,6.853) (7,3.438) (8,2.897) (9,0.4295) (10,0.1738) (11,0.003595) (12,0.004152) (13,5.37e-06) (14,0.001156) (15,4.419e-06) (16,0.0003509) (17,7.33e-07) (18,0.0001366) (19,6.9e-08) (20,2.118e-05) (22,1.227e-06) (24,2e-09)};\draw[densely dotted] (axis cs:18,0) -- (axis cs:18,7.538);\end{axis}\begin{axis}[width=72mm,height=42mm,axis y line*=right,axis x line=none,ymode=log,xmin=-1,xmax=25,ylabel={\scriptsize eindeutig (log)},tick label style={font=\tiny},label style={font=\scriptsize},ylabel shift=-3pt]\addplot[mark=*,mark size=0.7pt,draw=black,thin] coordinates {(0,15964482) (1,82673757) (2,156695676) (3,391552980) (4,685755005) (5,365548710) (6,222236577) (7,18344855) (8,4975540) (9,61314) (10,44069) (11,37) (12,11437) (13,6) (14,1806) (15,6) (16,429) (18,1)};\end{axis}\end{tikzpicture}\end{minipage}
\vspace{3mm}

\noindent\begin{minipage}[t]{60mm}\vspace{0pt}\sbox\hmboard{\chessboard[setfen=8/3K4/8/8/6N1/2p5/2P5/6Rk b - - 0 1]}\usebox\hmboard\\[1mm]
\makebox[\wd\hmboard]{\sffamily\hspace{3mm}\textbf{h\#9}\hfill(4+2)}\end{minipage}\hfill
\begin{minipage}[t]{52mm}\vspace{0pt}
\textcolor{ink}{\sffamily\textsc{Lösung}}\par\smallskip\small
\begin{tabular}{@{}r@{\enspace}l@{\enspace}l@{}}
\textbf{1.} & \symking{}xg1 & \symknight{}f6\\
\textbf{2.} & \symking{}f2 & \symknight{}e4+\\
\textbf{3.} & \symking{}e3 & \symknight{}xc3\\
\textbf{4.} & \symking{}d4 & \symknight{}b5+\\
\textbf{5.} & \symking{}c5 & c4\\
\textbf{6.} & \symking{}b6 & c5+\\
\textbf{7.} & \symking{}b7 & c6+\\
\textbf{8.} & \symking{}b8 & c7+\\
\textbf{9.} & \symking{}a8 & c8=\symqueen{}\#\\
\end{tabular}\par\vspace{3mm}
\qrcode[height=12mm]{https://helpman.komtera.lt/?fen=8/3K4/8/8/6N1/2p5/2P5/6Rk&moves=9}\\{\sffamily\tiny\color{ink}scannen zum Lösen}
\end{minipage}
\vspace{2mm}

{\raggedright\noindent\textcolor{ink}{\sffamily\textsc{Themen}}\enspace \themetag{reines Matt} \themetag{Mustermatt} \themetag{Spiegelmatt} \themetag{Umwandlung} \themetag{Excelsior} \themetag{ein Stein zieht} \themetag{Umnov}\par\smallskip{\footnotesize\textit{Hinweis:} die Partei am Zug steht im Diagramm im Schach; die Lösung beginnt mit einem Schlag.}\par}

\clearpage
\phantomsection\label{hm:KRNvkqb}
\begin{tcolorbox}[enhanced,colback=ink,colframe=ink,arc=0pt,boxrule=0pt,left=4mm,right=4mm,top=2.5mm,bottom=2.5mm]
\color{white}\sffamily{\LARGE\bfseries KRNvkqb}\hfill{\Large h\#9}\\[0.5mm]
{\small 6 Steine\quad Schwarz am Zug}\hfill{\small Nr.~334\qquad\href{https://helpman.komtera.lt/?fen=7k/4N2R/8/8/8/1q6/8/Kb6&moves=9}{\color{white}\bfseries lösen\,\ensuremath{\nearrow}}}
\end{tcolorbox}
\vspace{2mm}
\noindent\begin{minipage}[t]{36mm}\vspace{0pt}\sffamily\small\raggedright
\hmfigure{längste eindeutige}{{\Large h\#9}}\hmfigure{eindeutig bei dieser Länge}{{\Large 18} Stellungen}\hmfigure{längstes Matt}{{\Large h\#11}}\hmfigure{in der Klasse}{15,5 Milliarden Stellungen,\\ 7,98 Milliarden mit Hilfsmatt}
\end{minipage}\hfill
\begin{minipage}[t]{80mm}\vspace{0pt}\begin{tikzpicture}\begin{axis}[width=72mm,height=42mm,ybar,bar width=1.6pt,xmin=-1,xmax=23,ymin=0,ymax=2.507,axis y line*=left,axis x line*=bottom,xlabel={\scriptsize Halbzüge bis Matt},ylabel={\scriptsize Stellungen (Mrd.)},tick label style={font=\tiny},label style={font=\scriptsize},enlarge x limits=false,ylabel shift=-3pt,xlabel shift=-3pt]\addplot[fill=ink!55,draw=none] coordinates {(0,0.001173) (1,0.00686) (2,0.03242) (3,0.2083) (4,0.5446) (5,1.689) (6,2.193) (7,2.28) (8,0.748) (9,0.2572) (10,0.01155) (11,0.001356) (12,0.0006809) (13,5.51e-07) (14,0.0006732) (15,1.22e-07) (16,0.0006427) (17,1.27e-07) (18,0.0007021) (19,1.2e-08) (20,0.0002128) (21,2.4e-08) (22,8.29e-07)};\draw[densely dotted] (axis cs:18,0) -- (axis cs:18,2.507);\end{axis}\begin{axis}[width=72mm,height=42mm,axis y line*=right,axis x line=none,ymode=log,xmin=-1,xmax=23,ylabel={\scriptsize eindeutig (log)},tick label style={font=\tiny},label style={font=\scriptsize},ylabel shift=-3pt]\addplot[mark=*,mark size=0.7pt,draw=black,thin] coordinates {(0,1172890) (1,6781191) (2,16158616) (3,74824737) (4,134947140) (5,155736580) (6,48839883) (7,13289563) (8,415817) (9,48259) (10,4077) (11,16) (12,1568) (14,331) (16,57) (18,18)};\end{axis}\end{tikzpicture}\end{minipage}
\vspace{3mm}

\noindent\begin{minipage}[t]{60mm}\vspace{0pt}\sbox\hmboard{\chessboard[setfen=7k/4N2R/8/8/8/1q6/8/Kb6 b - - 0 1]}\usebox\hmboard\\[1mm]
\makebox[\wd\hmboard]{\sffamily\hspace{3mm}\textbf{h\#9}\hfill(3+3)}\end{minipage}\hfill
\begin{minipage}[t]{52mm}\vspace{0pt}
\textcolor{ink}{\sffamily\textsc{Lösung}}\par\smallskip\small
\begin{tabular}{@{}r@{\enspace}l@{\enspace}l@{}}
\textbf{1.} & \symking{}xh7 & \symknight{}c6\\
\textbf{2.} & \symking{}g6 & \symknight{}a5\\
\textbf{3.} & \symking{}f5 & \symknight{}xb3\\
\textbf{4.} & \symking{}e4 & \symking{}b2\\
\textbf{5.} & \symking{}d3 & \symking{}c1\\
\textbf{6.} & \symking{}c3 & \symking{}d1\\
\textbf{7.} & \symking{}b2 & \symknight{}d4\\
\textbf{8.} & \symking{}a1 & \symking{}c1\\
\textbf{9.} & \symbishop{}a2 & \symknight{}c2\#\\
\end{tabular}\par\vspace{3mm}
\qrcode[height=12mm]{https://helpman.komtera.lt/?fen=7k/4N2R/8/8/8/1q6/8/Kb6&moves=9}\\{\sffamily\tiny\color{ink}scannen zum Lösen}
\end{minipage}
\vspace{2mm}

{\raggedright\noindent\textcolor{ink}{\sffamily\textsc{Themen}}\enspace \themetag{reines Matt} \themetag{Mustermatt} \themetag{Idealmatt} \themetag{Switchback} \themetag{Selbstblock} \themetag{ohne Schach}\par\smallskip{\footnotesize\textit{Hinweis:} die Partei am Zug steht im Diagramm im Schach; die Lösung beginnt mit einem Schlag.}\par}

\clearpage
\phantomsection\label{hm:KRNvkqp}
\begin{tcolorbox}[enhanced,colback=ink,colframe=ink,arc=0pt,boxrule=0pt,left=4mm,right=4mm,top=2.5mm,bottom=2.5mm]
\color{white}\sffamily{\LARGE\bfseries KRNvkqp}\hfill{\Large h\#9}\\[0.5mm]
{\small 6 Steine\quad Schwarz am Zug}\hfill{\small Nr.~335\qquad\href{https://helpman.komtera.lt/?fen=8/8/1p6/8/5NRk/8/2q5/K7&moves=9}{\color{white}\bfseries lösen\,\ensuremath{\nearrow}}}
\end{tcolorbox}
\vspace{2mm}
\noindent\begin{minipage}[t]{36mm}\vspace{0pt}\sffamily\small\raggedright
\hmfigure{längste eindeutige}{{\Large h\#9}}\hmfigure{eindeutig bei dieser Länge}{{\Large 81} Stellungen}\hmfigure{längstes Matt}{{\Large h\#12}}\hmfigure{in der Klasse}{45,4 Milliarden Stellungen,\\ 25,3 Milliarden mit Hilfsmatt}
\end{minipage}\hfill
\begin{minipage}[t]{80mm}\vspace{0pt}\begin{tikzpicture}\begin{axis}[width=72mm,height=42mm,ybar,bar width=1.6pt,xmin=-1,xmax=25,ymin=0,ymax=7.876,axis y line*=left,axis x line*=bottom,xlabel={\scriptsize Halbzüge bis Matt},ylabel={\scriptsize Stellungen (Mrd.)},tick label style={font=\tiny},label style={font=\scriptsize},enlarge x limits=false,ylabel shift=-3pt,xlabel shift=-3pt]\addplot[fill=ink!55,draw=none] coordinates {(0,0.005645) (1,0.03299) (2,0.1407) (3,0.7805) (4,1.914) (5,5.249) (6,7.16) (7,6.552) (8,2.633) (9,0.7538) (10,0.05895) (11,0.008455) (12,0.003137) (13,6.539e-05) (14,0.002847) (15,5.901e-06) (16,0.002319) (17,6.366e-06) (18,0.00136) (19,3.638e-06) (20,0.0002443) (21,1.162e-06) (22,6.126e-06) (23,6.54e-07) (24,1.1e-08)};\draw[densely dotted] (axis cs:18,0) -- (axis cs:18,7.876);\end{axis}\begin{axis}[width=72mm,height=42mm,axis y line*=right,axis x line=none,ymode=log,xmin=-1,xmax=25,ylabel={\scriptsize eindeutig (log)},tick label style={font=\tiny},label style={font=\scriptsize},ylabel shift=-3pt]\addplot[mark=*,mark size=0.7pt,draw=black,thin] coordinates {(0,5645105) (1,32562535) (2,71281593) (3,285240315) (4,457660830) (5,476945386) (6,164868892) (7,38290892) (8,2185723) (9,229095) (10,27068) (11,649) (12,10202) (13,19) (14,4415) (15,1) (16,575) (18,81)};\end{axis}\end{tikzpicture}\end{minipage}
\vspace{3mm}

\noindent\begin{minipage}[t]{60mm}\vspace{0pt}\sbox\hmboard{\chessboard[setfen=8/8/1p6/8/5NRk/8/2q5/K7 b - - 0 1]}\usebox\hmboard\\[1mm]
\makebox[\wd\hmboard]{\sffamily\hspace{3mm}\textbf{h\#9}\hfill(3+3)}\end{minipage}\hfill
\begin{minipage}[t]{52mm}\vspace{0pt}
\textcolor{ink}{\sffamily\textsc{Lösung}}\par\smallskip\small
\begin{tabular}{@{}r@{\enspace}l@{\enspace}l@{}}
\textbf{1.} & \symking{}xg4 & \symknight{}d5\\
\textbf{2.} & \symqueen{}a4+ & \symking{}b2\\
\textbf{3.} & \symking{}f3 & \symking{}c3\\
\textbf{4.} & \symqueen{}a5+ & \symking{}d4\\
\textbf{5.} & \symking{}e2 & \symking{}e5\\
\textbf{6.} & \symking{}d3 & \symking{}d6\\
\textbf{7.} & \symking{}c4 & \symking{}c7\\
\textbf{8.} & \symking{}b5 & \symking{}b8\\
\textbf{9.} & \symking{}a6 & \symknight{}c7\#\\
\end{tabular}\par\vspace{3mm}
\qrcode[height=12mm]{https://helpman.komtera.lt/?fen=8/8/1p6/8/5NRk/8/2q5/K7&moves=9}\\{\sffamily\tiny\color{ink}scannen zum Lösen}
\end{minipage}
\vspace{2mm}

{\raggedright\noindent\textcolor{ink}{\sffamily\textsc{Themen}}\enspace \themetag{reines Matt} \themetag{Mustermatt} \themetag{Idealmatt} \themetag{Selbstblock}\par\smallskip{\footnotesize\textit{Hinweis:} die Partei am Zug steht im Diagramm im Schach; die Lösung beginnt mit einem Schlag.}\par}

\clearpage
\phantomsection\label{hm:KRPvkbp}
\begin{tcolorbox}[enhanced,colback=ink,colframe=ink,arc=0pt,boxrule=0pt,left=4mm,right=4mm,top=2.5mm,bottom=2.5mm]
\color{white}\sffamily{\LARGE\bfseries KRPvkbp}\hfill{\Large h\#9}\\[0.5mm]
{\small 6 Steine\quad Schwarz am Zug}\hfill{\small Nr.~336\qquad\href{https://helpman.komtera.lt/?fen=8/8/6p1/6P1/R7/kb6/8/K7&moves=9}{\color{white}\bfseries lösen\,\ensuremath{\nearrow}}}
\end{tcolorbox}
\vspace{2mm}
\noindent\begin{minipage}[t]{36mm}\vspace{0pt}\sffamily\small\raggedright
\hmfigure{längste eindeutige}{{\Large h\#9}}\hmfigure{eindeutig bei dieser Länge}{{\Large 5} Stellungen}\hmfigure{längstes Matt}{{\Large h\#15}}\hmfigure{in der Klasse}{34,1 Milliarden Stellungen,\\ 22,1 Milliarden mit Hilfsmatt}
\end{minipage}\hfill
\begin{minipage}[t]{80mm}\vspace{0pt}\begin{tikzpicture}\begin{axis}[width=72mm,height=42mm,ybar,bar width=1.6pt,xmin=-1,xmax=31,ymin=0,ymax=5.038,axis y line*=left,axis x line*=bottom,xlabel={\scriptsize Halbzüge bis Matt},ylabel={\scriptsize Stellungen (Mrd.)},tick label style={font=\tiny},label style={font=\scriptsize},enlarge x limits=false,ylabel shift=-3pt,xlabel shift=-3pt]\addplot[fill=ink!55,draw=none] coordinates {(0,0.006994) (1,0.05238) (2,0.1589) (3,0.8546) (4,1.882) (5,3.624) (6,4.58) (7,3.879) (8,3.678) (9,1.956) (10,1.108) (11,0.2226) (12,0.08411) (13,0.0009719) (14,0.0005359) (15,2.73e-07) (16,6.565e-05) (17,1.7e-08) (18,3.392e-05) (19,3.5e-08) (20,2.234e-05) (21,8e-09) (22,1.243e-05) (23,3e-09) (24,4.289e-06) (26,8.98e-07) (28,3.88e-07) (30,4e-09)};\draw[densely dotted] (axis cs:18,0) -- (axis cs:18,5.038);\end{axis}\begin{axis}[width=72mm,height=42mm,axis y line*=right,axis x line=none,ymode=log,xmin=-1,xmax=31,ylabel={\scriptsize eindeutig (log)},tick label style={font=\tiny},label style={font=\scriptsize},ylabel shift=-3pt]\addplot[mark=*,mark size=0.7pt,draw=black,thin] coordinates {(0,6994146) (1,41186530) (2,60655756) (3,196192489) (4,308506313) (5,213374836) (6,126723969) (7,40897475) (8,20663699) (9,3676760) (10,1093253) (11,13117) (12,5092) (13,8) (14,212) (16,15) (18,5)};\end{axis}\end{tikzpicture}\end{minipage}
\vspace{3mm}

\noindent\begin{minipage}[t]{60mm}\vspace{0pt}\sbox\hmboard{\chessboard[setfen=8/8/6p1/6P1/R7/kb6/8/K7 b - - 0 1]}\usebox\hmboard\\[1mm]
\makebox[\wd\hmboard]{\sffamily\hspace{3mm}\textbf{h\#9}\hfill(3+3)}\end{minipage}\hfill
\begin{minipage}[t]{52mm}\vspace{0pt}
\textcolor{ink}{\sffamily\textsc{Lösung}}\par\smallskip\small
\begin{tabular}{@{}r@{\enspace}l@{\enspace}l@{}}
\textbf{1.} & \symking{}xa4 & \symking{}b2\\
\textbf{2.} & \symking{}b5 & \symking{}c3\\
\textbf{3.} & \symking{}c6 & \symking{}d4\\
\textbf{4.} & \symking{}d7 & \symking{}e5\\
\textbf{5.} & \symking{}e8 & \symking{}f6\\
\textbf{6.} & \symking{}f8 & \symking{}xg6\\
\textbf{7.} & \symking{}g8 & \symking{}h6\\
\textbf{8.} & \symking{}h8 & g6\\
\textbf{9.} & \symbishop{}g8 & g7\#\\
\end{tabular}\par\vspace{3mm}
\qrcode[height=12mm]{https://helpman.komtera.lt/?fen=8/8/6p1/6P1/R7/kb6/8/K7&moves=9}\\{\sffamily\tiny\color{ink}scannen zum Lösen}
\end{minipage}
\vspace{2mm}

{\raggedright\noindent\textcolor{ink}{\sffamily\textsc{Themen}}\enspace \themetag{reines Matt} \themetag{Mustermatt} \themetag{Idealmatt} \themetag{Selbstblock} \themetag{ohne Schach}\par\smallskip{\footnotesize\textit{Hinweis:} die Partei am Zug steht im Diagramm im Schach; die Lösung beginnt mit einem Schlag.}\par}

\clearpage
\phantomsection\label{hm:KRPvkpp}
\begin{tcolorbox}[enhanced,colback=ink,colframe=ink,arc=0pt,boxrule=0pt,left=4mm,right=4mm,top=2.5mm,bottom=2.5mm]
\color{white}\sffamily{\LARGE\bfseries KRPvkpp}\hfill{\Large h\#9}\\[0.5mm]
{\small 6 Steine\quad Schwarz am Zug}\hfill{\small Nr.~337\qquad\href{https://helpman.komtera.lt/?fen=8/8/4p3/8/R1p5/k7/2P5/K7&moves=9}{\color{white}\bfseries lösen\,\ensuremath{\nearrow}}}
\end{tcolorbox}
\vspace{2mm}
\noindent\begin{minipage}[t]{36mm}\vspace{0pt}\sffamily\small\raggedright
\hmfigure{längste eindeutige}{{\Large h\#9}}\hmfigure{eindeutig bei dieser Länge}{{\Large 24} Stellungen}\hmfigure{längstes Matt}{{\Large h\#15}}\hmfigure{in der Klasse}{25,6 Milliarden Stellungen,\\ 8,68 Milliarden mit Hilfsmatt}
\end{minipage}\hfill
\begin{minipage}[t]{80mm}\vspace{0pt}\begin{tikzpicture}\begin{axis}[width=72mm,height=42mm,ybar,bar width=1.6pt,xmin=-1,xmax=31,ymin=0,ymax=2.055,axis y line*=left,axis x line*=bottom,xlabel={\scriptsize Halbzüge bis Matt},ylabel={\scriptsize Stellungen (Mrd.)},tick label style={font=\tiny},label style={font=\scriptsize},enlarge x limits=false,ylabel shift=-3pt,xlabel shift=-3pt]\addplot[fill=ink!55,draw=none] coordinates {(0,0.0046) (1,0.03266) (2,0.07859) (3,0.3727) (4,0.8096) (5,1.362) (6,1.868) (7,1.386) (8,1.49) (9,0.6978) (10,0.4533) (11,0.08288) (12,0.04265) (13,0.0007644) (14,0.000514) (15,8.81e-07) (16,8e-05) (17,1.77e-07) (18,7.213e-05) (19,6.6e-08) (20,3.804e-05) (21,3e-08) (22,1.68e-05) (23,2.2e-08) (24,7.468e-06) (25,3e-09) (26,2.615e-06) (27,5e-09) (28,8.08e-07) (29,2e-09) (30,2.88e-07)};\draw[densely dotted] (axis cs:18,0) -- (axis cs:18,2.055);\end{axis}\begin{axis}[width=72mm,height=42mm,axis y line*=right,axis x line=none,ymode=log,xmin=-1,xmax=31,ylabel={\scriptsize eindeutig (log)},tick label style={font=\tiny},label style={font=\scriptsize},ylabel shift=-3pt]\addplot[mark=*,mark size=0.7pt,draw=black,thin] coordinates {(0,4600093) (1,25462892) (2,35340775) (3,94400576) (4,157069747) (5,80591506) (6,50355949) (7,12148332) (8,6366504) (9,944614) (10,358425) (11,11667) (12,8646) (13,2) (14,136) (15,1) (16,48) (18,24)};\end{axis}\end{tikzpicture}\end{minipage}
\vspace{3mm}

\noindent\begin{minipage}[t]{60mm}\vspace{0pt}\sbox\hmboard{\chessboard[setfen=8/8/4p3/8/R1p5/k7/2P5/K7 b - - 0 1]}\usebox\hmboard\\[1mm]
\makebox[\wd\hmboard]{\sffamily\hspace{3mm}\textbf{h\#9}\hfill(3+3)}\end{minipage}\hfill
\begin{minipage}[t]{52mm}\vspace{0pt}
\textcolor{ink}{\sffamily\textsc{Lösung}}\par\smallskip\small
\begin{tabular}{@{}r@{\enspace}l@{\enspace}l@{}}
\textbf{1.} & \symking{}xa4 & \symking{}b2\\
\textbf{2.} & c3+ & \symking{}xc3\\
\textbf{3.} & e5 & \symking{}b2\\
\textbf{4.} & e4 & c4\\
\textbf{5.} & e3 & c5\\
\textbf{6.} & e2 & c6\\
\textbf{7.} & e1=\symrook{} & c7\\
\textbf{8.} & \symrook{}e5 & c8=\symqueen{}\\
\textbf{9.} & \symrook{}a5 & \symqueen{}c4\#\\
\end{tabular}\par\vspace{3mm}
\qrcode[height=12mm]{https://helpman.komtera.lt/?fen=8/8/4p3/8/R1p5/k7/2P5/K7&moves=9}\\{\sffamily\tiny\color{ink}scannen zum Lösen}
\end{minipage}
\vspace{2mm}

{\raggedright\noindent\textcolor{ink}{\sffamily\textsc{Themen}}\enspace \themetag{Umwandlung} \themetag{Unterverwandlung} \themetag{Excelsior} \themetag{Switchback} \themetag{Rundlauf} \themetag{Selbstblock} \themetag{Kniest}\par\smallskip{\footnotesize\textit{Hinweis:} die Partei am Zug steht im Diagramm im Schach; die Lösung beginnt mit einem Schlag.}\par}

\clearpage
\phantomsection\label{hm:KBBvknn}
\begin{tcolorbox}[enhanced,colback=ink,colframe=ink,arc=0pt,boxrule=0pt,left=4mm,right=4mm,top=2.5mm,bottom=2.5mm]
\color{white}\sffamily{\LARGE\bfseries KBBvknn}\hfill{\Large h\#8.5}\\[0.5mm]
{\small 6 Steine\quad Weiß am Zug}\hfill{\small Nr.~338\qquad\href{https://helpman.komtera.lt/?fen=n7/8/5B2/8/7B/7k/2n5/1K6&moves=8.5}{\color{white}\bfseries lösen\,\ensuremath{\nearrow}}}
\end{tcolorbox}
\vspace{2mm}
\noindent\begin{minipage}[t]{36mm}\vspace{0pt}\sffamily\small\raggedright
\hmfigure{längste eindeutige}{{\Large h\#8.5}}\hmfigure{eindeutig bei dieser Länge}{{\Large 17} Stellungen}\hmfigure{längstes Matt}{{\Large h\#13}}\hmfigure{in der Klasse}{15,5 Milliarden Stellungen,\\ 2,41 Milliarden mit Hilfsmatt}
\end{minipage}\hfill
\begin{minipage}[t]{80mm}\vspace{0pt}\begin{tikzpicture}\begin{axis}[width=72mm,height=42mm,ybar,bar width=1.6pt,xmin=-1,xmax=27,ymin=0,ymax=305.3,axis y line*=left,axis x line*=bottom,xlabel={\scriptsize Halbzüge bis Matt},ylabel={\scriptsize Stellungen (Mio.)},tick label style={font=\tiny},label style={font=\scriptsize},enlarge x limits=false,ylabel shift=-3pt,xlabel shift=-3pt]\addplot[fill=ink!55,draw=none] coordinates {(0,0.1649) (1,0.6719) (2,2.642) (3,17.38) (4,52.93) (5,138.8) (6,242.7) (7,237.2) (8,277.6) (9,207.5) (10,213.7) (11,185.1) (12,240.3) (13,217.9) (14,192.7) (15,114.7) (16,46.36) (17,19.84) (18,2.916) (19,1.783) (20,0.1611) (21,0.1312) (22,0.009922) (23,0.0001) (24,7.8e-05) (25,7e-06) (26,1e-06)};\draw[densely dotted] (axis cs:17,0) -- (axis cs:17,305.3);\end{axis}\begin{axis}[width=72mm,height=42mm,axis y line*=right,axis x line=none,ymode=log,xmin=-1,xmax=27,ylabel={\scriptsize eindeutig (log)},tick label style={font=\tiny},label style={font=\scriptsize},ylabel shift=-3pt]\addplot[mark=*,mark size=0.7pt,draw=black,thin] coordinates {(0,164866) (1,666719) (2,1557248) (3,7531729) (4,10717638) (5,9416877) (6,6774951) (7,2605072) (8,1620862) (9,682278) (10,531481) (11,331578) (12,245347) (13,78243) (14,16845) (15,1829) (16,162) (17,17)};\end{axis}\end{tikzpicture}\end{minipage}
\vspace{3mm}

\noindent\begin{minipage}[t]{60mm}\vspace{0pt}\sbox\hmboard{\chessboard[setfen=n7/8/5B2/8/7B/7k/2n5/1K6 w - - 0 1]}\usebox\hmboard\\[1mm]
\makebox[\wd\hmboard]{\sffamily\hspace{3mm}\textbf{h\#8.5}\hfill(3+3)}\end{minipage}\hfill
\begin{minipage}[t]{52mm}\vspace{0pt}
\textcolor{ink}{\sffamily\textsc{Lösung}}\par\smallskip\small
\begin{tabular}{@{}r@{\enspace}l@{\enspace}l@{}}
\textbf{1.} & \ldots & \symbishop{}g3\\
\textbf{2.} & \symknight{}e3 & \symbishop{}b8\\
\textbf{3.} & \symknight{}g4 & \symking{}c2\\
\textbf{4.} & \symknight{}c7 & \symking{}d3\\
\textbf{5.} & \symking{}g3 & \symbishop{}e5+\\
\textbf{6.} & \symking{}h4 & \symking{}e4\\
\textbf{7.} & \symknight{}e6 & \symking{}f5\\
\textbf{8.} & \symknight{}g5 & \symking{}g6\\
\textbf{9.} & \symknight{}h3 & \symbishop{}g3\#\\
\end{tabular}\par\vspace{3mm}
\qrcode[height=12mm]{https://helpman.komtera.lt/?fen=n7/8/5B2/8/7B/7k/2n5/1K6&moves=8.5}\\{\sffamily\tiny\color{ink}scannen zum Lösen}
\end{minipage}
\vspace{2mm}

{\raggedright\noindent\textcolor{ink}{\sffamily\textsc{Themen}}\enspace \themetag{reines Matt} \themetag{Mustermatt} \themetag{Idealmatt} \themetag{Selbstblock} \themetag{ohne Schlag}\par}

\clearpage
\phantomsection\label{hm:KBvkbbb}
\begin{tcolorbox}[enhanced,colback=ink,colframe=ink,arc=0pt,boxrule=0pt,left=4mm,right=4mm,top=2.5mm,bottom=2.5mm]
\color{white}\sffamily{\LARGE\bfseries KBvkbbb}\hfill{\Large h\#8.5}\\[0.5mm]
{\small 6 Steine\quad Weiß am Zug}\hfill{\small Nr.~339\qquad\href{https://helpman.komtera.lt/?fen=8/8/8/8/b7/8/5b2/kBK4b&moves=8.5}{\color{white}\bfseries lösen\,\ensuremath{\nearrow}}}
\end{tcolorbox}
\vspace{2mm}
\noindent\begin{minipage}[t]{36mm}\vspace{0pt}\sffamily\small\raggedright
\hmfigure{längste eindeutige}{{\Large h\#8.5}}\hmfigure{eindeutig bei dieser Länge}{{\Large 446} Stellungen}\hmfigure{längstes Matt}{{\Large h\#10}}\hmfigure{in der Klasse}{15,5 Milliarden Stellungen,\\ 1,35 Milliarden mit Hilfsmatt}
\end{minipage}\hfill
\begin{minipage}[t]{80mm}\vspace{0pt}\begin{tikzpicture}\begin{axis}[width=72mm,height=42mm,ybar,bar width=1.6pt,xmin=-1,xmax=21,ymin=0,ymax=251.5,axis y line*=left,axis x line*=bottom,xlabel={\scriptsize Halbzüge bis Matt},ylabel={\scriptsize Stellungen (Mio.)},tick label style={font=\tiny},label style={font=\scriptsize},enlarge x limits=false,ylabel shift=-3pt,xlabel shift=-3pt]\addplot[fill=ink!55,draw=none] coordinates {(0,0.01371) (1,0.0551) (2,0.5019) (3,1.741) (4,7.093) (5,19.39) (6,45.36) (7,97.99) (8,133.1) (9,225.3) (10,172.3) (11,228.6) (12,124.5) (13,111) (14,59.19) (15,58.54) (16,23.43) (17,26.25) (18,5.784) (19,7.37) (20,0.002171)};\draw[densely dotted] (axis cs:17,0) -- (axis cs:17,251.5);\end{axis}\begin{axis}[width=72mm,height=42mm,axis y line*=right,axis x line=none,ymode=log,xmin=-1,xmax=21,ylabel={\scriptsize eindeutig (log)},tick label style={font=\tiny},label style={font=\scriptsize},ylabel shift=-3pt]\addplot[mark=*,mark size=0.7pt,draw=black,thin] coordinates {(0,13706) (1,54536) (2,418135) (3,715058) (4,1442622) (5,2517759) (6,1546970) (7,2119238) (8,1083374) (9,1339077) (10,567730) (11,465997) (12,129695) (13,56020) (14,14468) (15,6361) (16,1414) (17,446)};\end{axis}\end{tikzpicture}\end{minipage}
\vspace{3mm}

\noindent\begin{minipage}[t]{60mm}\vspace{0pt}\sbox\hmboard{\chessboard[setfen=8/8/8/8/b7/8/5b2/kBK4b w - - 0 1]}\usebox\hmboard\\[1mm]
\makebox[\wd\hmboard]{\sffamily\hspace{3mm}\textbf{h\#8.5}\hfill(2+4)}\end{minipage}\hfill
\begin{minipage}[t]{52mm}\vspace{0pt}
\textcolor{ink}{\sffamily\textsc{Lösung}}\par\smallskip\small
\begin{tabular}{@{}r@{\enspace}l@{\enspace}l@{}}
\textbf{1.} & \ldots & \symking{}d2\\
\textbf{2.} & \symking{}b2 & \symking{}d3\\
\textbf{3.} & \symking{}c1 & \symbishop{}a2\\
\textbf{4.} & \symking{}d1 & \symbishop{}d5\\
\textbf{5.} & \symking{}e1 & \symbishop{}xh1\\
\textbf{6.} & \symking{}f1 & \symbishop{}d5\\
\textbf{7.} & \symking{}g1 & \symking{}e4\\
\textbf{8.} & \symking{}h1 & \symking{}f3\\
\textbf{9.} & \symbishop{}g1 & \symking{}g3\#\\
\end{tabular}\par\vspace{3mm}
\qrcode[height=12mm]{https://helpman.komtera.lt/?fen=8/8/8/8/b7/8/5b2/kBK4b&moves=8.5}\\{\sffamily\tiny\color{ink}scannen zum Lösen}
\end{minipage}
\vspace{2mm}

{\raggedright\noindent\textcolor{ink}{\sffamily\textsc{Themen}}\enspace \themetag{Switchback} \themetag{Selbstblock} \themetag{Kniest} \themetag{ohne Schach} \themetag{Inder}\par}

\clearpage
\phantomsection\label{hm:KBvkqbb}
\begin{tcolorbox}[enhanced,colback=ink,colframe=ink,arc=0pt,boxrule=0pt,left=4mm,right=4mm,top=2.5mm,bottom=2.5mm]
\color{white}\sffamily{\LARGE\bfseries KBvkqbb}\hfill{\Large h\#8.5}\\[0.5mm]
{\small 6 Steine\quad Weiß am Zug}\hfill{\small Nr.~340\qquad\href{https://helpman.komtera.lt/?fen=5b2/8/8/8/7B/7b/7k/Kq6&moves=8.5}{\color{white}\bfseries lösen\,\ensuremath{\nearrow}}}
\end{tcolorbox}
\vspace{2mm}
\noindent\begin{minipage}[t]{36mm}\vspace{0pt}\sffamily\small\raggedright
\hmfigure{längste eindeutige}{{\Large h\#8.5}}\hmfigure{eindeutig bei dieser Länge}{{\Large 25} Stellungen}\hmfigure{längstes Matt}{{\Large h\#10.5}}\hmfigure{in der Klasse}{15,5 Milliarden Stellungen,\\ 3,12 Milliarden mit Hilfsmatt}
\end{minipage}\hfill
\begin{minipage}[t]{80mm}\vspace{0pt}\begin{tikzpicture}\begin{axis}[width=72mm,height=42mm,ybar,bar width=1.6pt,xmin=-1,xmax=22,ymin=0,ymax=908.8,axis y line*=left,axis x line*=bottom,xlabel={\scriptsize Halbzüge bis Matt},ylabel={\scriptsize Stellungen (Mio.)},tick label style={font=\tiny},label style={font=\scriptsize},enlarge x limits=false,ylabel shift=-3pt,xlabel shift=-3pt]\addplot[fill=ink!55,draw=none] coordinates {(0,0.008319) (1,0.03624) (2,0.3669) (3,1.871) (4,7.942) (5,29.86) (6,70.78) (7,219) (8,290.6) (9,696.4) (10,490.8) (11,826.2) (12,272.1) (13,194.5) (14,9.675) (15,6.358) (16,0.001179) (17,1.841) (18,5e-06) (19,0.4818) (21,1e-06)};\draw[densely dotted] (axis cs:17,0) -- (axis cs:17,908.8);\end{axis}\begin{axis}[width=72mm,height=42mm,axis y line*=right,axis x line=none,ymode=log,xmin=-1,xmax=22,ylabel={\scriptsize eindeutig (log)},tick label style={font=\tiny},label style={font=\scriptsize},ylabel shift=-3pt]\addplot[mark=*,mark size=0.7pt,draw=black,thin] coordinates {(0,8319) (1,36167) (2,268826) (3,963134) (4,1943881) (5,4972666) (6,4149274) (7,9265646) (8,5157723) (9,8346912) (10,2275795) (11,2250275) (12,134191) (13,41302) (14,39) (15,841) (17,25)};\end{axis}\end{tikzpicture}\end{minipage}
\vspace{3mm}

\noindent\begin{minipage}[t]{60mm}\vspace{0pt}\sbox\hmboard{\chessboard[setfen=5b2/8/8/8/7B/7b/7k/Kq6 w - - 0 1]}\usebox\hmboard\\[1mm]
\makebox[\wd\hmboard]{\sffamily\hspace{3mm}\textbf{h\#8.5}\hfill(2+4)}\end{minipage}\hfill
\begin{minipage}[t]{52mm}\vspace{0pt}
\textcolor{ink}{\sffamily\textsc{Lösung}}\par\smallskip\small
\begin{tabular}{@{}r@{\enspace}l@{\enspace}l@{}}
\textbf{1.} & \ldots & \symking{}xb1\\
\textbf{2.} & \symbishop{}f5+ & \symking{}b2\\
\textbf{3.} & \symking{}h3 & \symking{}c3\\
\textbf{4.} & \symking{}g4 & \symking{}d4\\
\textbf{5.} & \symking{}h5 & \symking{}e5\\
\textbf{6.} & \symking{}h6 & \symking{}f6\\
\textbf{7.} & \symking{}h7 & \symking{}f7\\
\textbf{8.} & \symking{}h8 & \symking{}xf8\\
\textbf{9.} & \symbishop{}h7 & \symbishop{}f6\#\\
\end{tabular}\par\vspace{3mm}
\qrcode[height=12mm]{https://helpman.komtera.lt/?fen=5b2/8/8/8/7B/7b/7k/Kq6&moves=8.5}\\{\sffamily\tiny\color{ink}scannen zum Lösen}
\end{minipage}
\vspace{2mm}

{\raggedright\noindent\textcolor{ink}{\sffamily\textsc{Themen}}\enspace \themetag{Selbstblock}\par\smallskip{\footnotesize\textit{Hinweis:} die Partei am Zug steht im Diagramm im Schach; die Lösung beginnt mit einem Schlag.}\par}

\clearpage
\phantomsection\label{hm:KBvkrbb}
\begin{tcolorbox}[enhanced,colback=ink,colframe=ink,arc=0pt,boxrule=0pt,left=4mm,right=4mm,top=2.5mm,bottom=2.5mm]
\color{white}\sffamily{\LARGE\bfseries KBvkrbb}\hfill{\Large h\#8.5}\\[0.5mm]
{\small 6 Steine\quad Weiß am Zug}\hfill{\small Nr.~341\qquad\href{https://helpman.komtera.lt/?fen=b7/8/8/8/8/8/2B2b2/k1Kr4&moves=8.5}{\color{white}\bfseries lösen\,\ensuremath{\nearrow}}}
\end{tcolorbox}
\vspace{2mm}
\noindent\begin{minipage}[t]{36mm}\vspace{0pt}\sffamily\small\raggedright
\hmfigure{längste eindeutige}{{\Large h\#8.5}}\hmfigure{eindeutig bei dieser Länge}{{\Large 2} Stellungen}\hmfigure{längstes Matt}{{\Large h\#9.5}}\hmfigure{in der Klasse}{15,5 Milliarden Stellungen,\\ 3,38 Milliarden mit Hilfsmatt}
\end{minipage}\hfill
\begin{minipage}[t]{80mm}\vspace{0pt}\begin{tikzpicture}\begin{axis}[width=72mm,height=42mm,ybar,bar width=1.6pt,xmin=-1,xmax=20,ymin=0,ymax=939.2,axis y line*=left,axis x line*=bottom,xlabel={\scriptsize Halbzüge bis Matt},ylabel={\scriptsize Stellungen (Mio.)},tick label style={font=\tiny},label style={font=\scriptsize},enlarge x limits=false,ylabel shift=-3pt,xlabel shift=-3pt]\addplot[fill=ink!55,draw=none] coordinates {(0,0.01277) (1,0.05474) (2,0.51) (3,2.307) (4,10.43) (5,33.23) (6,85.66) (7,220) (8,340.7) (9,689.6) (10,593) (11,853.8) (12,341.7) (13,193) (14,12.83) (15,2.214) (16,0.00415) (17,0.17) (19,0.03284)};\draw[densely dotted] (axis cs:17,0) -- (axis cs:17,939.2);\end{axis}\begin{axis}[width=72mm,height=42mm,axis y line*=right,axis x line=none,ymode=log,xmin=-1,xmax=20,ylabel={\scriptsize eindeutig (log)},tick label style={font=\tiny},label style={font=\scriptsize},ylabel shift=-3pt]\addplot[mark=*,mark size=0.7pt,draw=black,thin] coordinates {(0,12770) (1,54595) (2,396695) (3,1092093) (4,2584617) (5,5077837) (6,4841145) (7,8172882) (8,6899868) (9,8793122) (10,3489251) (11,2567956) (12,261894) (13,33161) (14,118) (15,145) (17,2)};\end{axis}\end{tikzpicture}\end{minipage}
\vspace{3mm}

\noindent\begin{minipage}[t]{60mm}\vspace{0pt}\sbox\hmboard{\chessboard[setfen=b7/8/8/8/8/8/2B2b2/k1Kr4 w - - 0 1]}\usebox\hmboard\\[1mm]
\makebox[\wd\hmboard]{\sffamily\hspace{3mm}\textbf{h\#8.5}\hfill(2+4)}\end{minipage}\hfill
\begin{minipage}[t]{52mm}\vspace{0pt}
\textcolor{ink}{\sffamily\textsc{Lösung}}\par\smallskip\small
\begin{tabular}{@{}r@{\enspace}l@{\enspace}l@{}}
\textbf{1.} & \ldots & \symking{}xd1\\
\textbf{2.} & \symking{}b2 & \symbishop{}e4\\
\textbf{3.} & \symking{}c3 & \symbishop{}xa8\\
\textbf{4.} & \symking{}d3 & \symking{}c1\\
\textbf{5.} & \symking{}e2 & \symking{}c2\\
\textbf{6.} & \symking{}f1 & \symking{}d3\\
\textbf{7.} & \symking{}g1 & \symking{}e4\\
\textbf{8.} & \symking{}h1 & \symking{}f3\\
\textbf{9.} & \symbishop{}g1 & \symking{}g3\#\\
\end{tabular}\par\vspace{3mm}
\qrcode[height=12mm]{https://helpman.komtera.lt/?fen=b7/8/8/8/8/8/2B2b2/k1Kr4&moves=8.5}\\{\sffamily\tiny\color{ink}scannen zum Lösen}
\end{minipage}
\vspace{2mm}

{\raggedright\noindent\textcolor{ink}{\sffamily\textsc{Themen}}\enspace \themetag{Switchback} \themetag{Selbstblock} \themetag{ohne Schach}\par\smallskip{\footnotesize\textit{Hinweis:} die Partei am Zug steht im Diagramm im Schach; die Lösung beginnt mit einem Schlag.}\par}

\clearpage
\phantomsection\label{hm:KNPvkbp}
\begin{tcolorbox}[enhanced,colback=ink,colframe=ink,arc=0pt,boxrule=0pt,left=4mm,right=4mm,top=2.5mm,bottom=2.5mm]
\color{white}\sffamily{\LARGE\bfseries KNPvkbp}\hfill{\Large h\#8.5}\\[0.5mm]
{\small 6 Steine\quad Weiß am Zug}\hfill{\small Nr.~342\qquad\href{https://helpman.komtera.lt/?fen=8/8/b7/8/p7/P7/8/K3N2k&moves=8.5}{\color{white}\bfseries lösen\,\ensuremath{\nearrow}}}
\end{tcolorbox}
\vspace{2mm}
\noindent\begin{minipage}[t]{36mm}\vspace{0pt}\sffamily\small\raggedright
\hmfigure{längste eindeutige}{{\Large h\#8.5}}\hmfigure{eindeutig bei dieser Länge}{{\Large 33} Stellungen}\hmfigure{längstes Matt}{{\Large h\#14}}\hmfigure{in der Klasse}{34,1 Milliarden Stellungen,\\ 23,6 Milliarden mit Hilfsmatt}
\end{minipage}\hfill
\begin{minipage}[t]{80mm}\vspace{0pt}\begin{tikzpicture}\begin{axis}[width=72mm,height=42mm,ybar,bar width=1.6pt,xmin=-1,xmax=29,ymin=0,ymax=3.63,axis y line*=left,axis x line*=bottom,xlabel={\scriptsize Halbzüge bis Matt},ylabel={\scriptsize Stellungen (Mrd.)},tick label style={font=\tiny},label style={font=\scriptsize},enlarge x limits=false,ylabel shift=-3pt,xlabel shift=-3pt]\addplot[fill=ink!55,draw=none] coordinates {(0,9.441e-05) (1,0.00443) (2,0.01386) (3,0.1345) (4,0.3868) (5,1.5) (6,2.174) (7,2.942) (8,2.96) (9,3.3) (10,3.225) (11,3.228) (12,2.411) (13,0.9288) (14,0.3129) (15,0.05095) (16,0.01422) (17,0.003175) (18,0.0006406) (19,6.94e-05) (20,1.44e-06) (21,5.5e-08) (22,9.6e-08) (23,1.4e-08) (24,4e-08) (25,8e-09) (26,1.3e-08) (28,2e-09)};\draw[densely dotted] (axis cs:17,0) -- (axis cs:17,3.63);\end{axis}\begin{axis}[width=72mm,height=42mm,axis y line*=right,axis x line=none,ymode=log,xmin=-1,xmax=29,ylabel={\scriptsize eindeutig (log)},tick label style={font=\tiny},label style={font=\scriptsize},ylabel shift=-3pt]\addplot[mark=*,mark size=0.7pt,draw=black,thin] coordinates {(0,94412) (1,1459818) (2,4242454) (3,49634853) (4,87588638) (5,122182055) (6,82835780) (7,69665403) (8,41808024) (9,35565086) (10,18099054) (11,13123683) (12,4145500) (13,412435) (14,86141) (15,7849) (16,453) (17,33)};\end{axis}\end{tikzpicture}\end{minipage}
\vspace{3mm}

\noindent\begin{minipage}[t]{60mm}\vspace{0pt}\sbox\hmboard{\chessboard[setfen=8/8/b7/8/p7/P7/8/K3N2k w - - 0 1]}\usebox\hmboard\\[1mm]
\makebox[\wd\hmboard]{\sffamily\hspace{3mm}\textbf{h\#8.5}\hfill(3+3)}\end{minipage}\hfill
\begin{minipage}[t]{52mm}\vspace{0pt}
\textcolor{ink}{\sffamily\textsc{Lösung}}\par\smallskip\small
\begin{tabular}{@{}r@{\enspace}l@{\enspace}l@{}}
\textbf{1.} & \ldots & \symking{}b2\\
\textbf{2.} & \symking{}g1 & \symking{}c3\\
\textbf{3.} & \symking{}f2 & \symking{}b4\\
\textbf{4.} & \symking{}e3 & \symking{}a5\\
\textbf{5.} & \symking{}d4 & \symking{}b6\\
\textbf{6.} & \symking{}c4 & \symking{}a7\\
\textbf{7.} & \symking{}b5 & \symknight{}d3\\
\textbf{8.} & \symking{}a5 & \symknight{}c5\\
\textbf{9.} & \symbishop{}b5 & \symknight{}b7\#\\
\end{tabular}\par\vspace{3mm}
\qrcode[height=12mm]{https://helpman.komtera.lt/?fen=8/8/b7/8/p7/P7/8/K3N2k&moves=8.5}\\{\sffamily\tiny\color{ink}scannen zum Lösen}
\end{minipage}
\vspace{2mm}

{\raggedright\noindent\textcolor{ink}{\sffamily\textsc{Themen}}\enspace \themetag{reines Matt} \themetag{Mustermatt} \themetag{Idealmatt} \themetag{Selbstblock} \themetag{ohne Schlag} \themetag{ohne Schach}\par}

\clearpage
\phantomsection\label{hm:KNPvknp}
\begin{tcolorbox}[enhanced,colback=ink,colframe=ink,arc=0pt,boxrule=0pt,left=4mm,right=4mm,top=2.5mm,bottom=2.5mm]
\color{white}\sffamily{\LARGE\bfseries KNPvknp}\hfill{\Large h\#8.5}\\[0.5mm]
{\small 6 Steine\quad Weiß am Zug}\hfill{\small Nr.~343\qquad\href{https://helpman.komtera.lt/?fen=N1K1n3/8/8/8/8/2p5/2P5/6k1&moves=8.5}{\color{white}\bfseries lösen\,\ensuremath{\nearrow}}}
\end{tcolorbox}
\vspace{2mm}
\noindent\begin{minipage}[t]{36mm}\vspace{0pt}\sffamily\small\raggedright
\hmfigure{längste eindeutige}{{\Large h\#8.5}}\hmfigure{eindeutig bei dieser Länge}{{\Large 1} Stellung}\hmfigure{längstes Matt}{{\Large h\#10.5}}\hmfigure{in der Klasse}{34,1 Milliarden Stellungen,\\ 24,2 Milliarden mit Hilfsmatt}
\end{minipage}\hfill
\begin{minipage}[t]{80mm}\vspace{0pt}\begin{tikzpicture}\begin{axis}[width=72mm,height=42mm,ybar,bar width=1.6pt,xmin=-1,xmax=22,ymin=0,ymax=3.791,axis y line*=left,axis x line*=bottom,xlabel={\scriptsize Halbzüge bis Matt},ylabel={\scriptsize Stellungen (Mrd.)},tick label style={font=\tiny},label style={font=\scriptsize},enlarge x limits=false,ylabel shift=-3pt,xlabel shift=-3pt]\addplot[fill=ink!55,draw=none] coordinates {(0,9.408e-05) (1,0.005495) (2,0.01452) (3,0.1308) (4,0.3399) (5,1.386) (6,2.082) (7,2.814) (8,3.075) (9,3.282) (10,3.446) (11,3.322) (12,2.69) (13,1.102) (14,0.4345) (15,0.04757) (16,0.009292) (17,0.00112) (18,0.000143) (19,6.043e-06) (20,2.2e-07) (21,4e-09)};\draw[densely dotted] (axis cs:17,0) -- (axis cs:17,3.791);\end{axis}\begin{axis}[width=72mm,height=42mm,axis y line*=right,axis x line=none,ymode=log,xmin=-1,xmax=22,ylabel={\scriptsize eindeutig (log)},tick label style={font=\tiny},label style={font=\scriptsize},ylabel shift=-3pt]\addplot[mark=*,mark size=0.7pt,draw=black,thin] coordinates {(0,94081) (1,1538648) (2,3696526) (3,44257233) (4,68790827) (5,107284955) (6,74722401) (7,61326239) (8,41431463) (9,33753177) (10,19200041) (11,15533447) (12,5170867) (13,370248) (14,55093) (15,3467) (16,273) (17,1)};\end{axis}\end{tikzpicture}\end{minipage}
\vspace{3mm}

\noindent\begin{minipage}[t]{60mm}\vspace{0pt}\sbox\hmboard{\chessboard[setfen=N1K1n3/8/8/8/8/2p5/2P5/6k1 w - - 0 1]}\usebox\hmboard\\[1mm]
\makebox[\wd\hmboard]{\sffamily\hspace{3mm}\textbf{h\#8.5}\hfill(3+3)}\end{minipage}\hfill
\begin{minipage}[t]{52mm}\vspace{0pt}
\textcolor{ink}{\sffamily\textsc{Lösung}}\par\smallskip\small
\begin{tabular}{@{}r@{\enspace}l@{\enspace}l@{}}
\textbf{1.} & \ldots & \symknight{}b6\\
\textbf{2.} & \symking{}f2 & \symknight{}a4\\
\textbf{3.} & \symking{}e3 & \symknight{}b2\\
\textbf{4.} & cxb2 & c4\\
\textbf{5.} & \symking{}d4 & c5\\
\textbf{6.} & \symknight{}d6+ & cxd6\\
\textbf{7.} & \symking{}c5 & d7\\
\textbf{8.} & \symking{}b6 & d8=\symqueen{}+\\
\textbf{9.} & \symking{}a7 & \symqueen{}a5\#\\
\end{tabular}\par\vspace{3mm}
\qrcode[height=12mm]{https://helpman.komtera.lt/?fen=N1K1n3/8/8/8/8/2p5/2P5/6k1&moves=8.5}\\{\sffamily\tiny\color{ink}scannen zum Lösen}
\end{minipage}
\vspace{2mm}

{\raggedright\noindent\textcolor{ink}{\sffamily\textsc{Themen}}\enspace \themetag{reines Matt} \themetag{Mustermatt} \themetag{Spiegelmatt} \themetag{Umwandlung} \themetag{Excelsior} \themetag{Umnov}\par}

\clearpage
\phantomsection\label{hm:KNvkbbb}
\begin{tcolorbox}[enhanced,colback=ink,colframe=ink,arc=0pt,boxrule=0pt,left=4mm,right=4mm,top=2.5mm,bottom=2.5mm]
\color{white}\sffamily{\LARGE\bfseries KNvkbbb}\hfill{\Large h\#8.5}\\[0.5mm]
{\small 6 Steine\quad Weiß am Zug}\hfill{\small Nr.~344\qquad\href{https://helpman.komtera.lt/?fen=8/8/8/8/b1b4k/8/4b3/K4N2&moves=8.5}{\color{white}\bfseries lösen\,\ensuremath{\nearrow}}}
\end{tcolorbox}
\vspace{2mm}
\noindent\begin{minipage}[t]{36mm}\vspace{0pt}\sffamily\small\raggedright
\hmfigure{längste eindeutige}{{\Large h\#8.5}}\hmfigure{eindeutig bei dieser Länge}{{\Large 1\,145} Stellungen}\hmfigure{längstes Matt}{{\Large h\#10.5}}\hmfigure{in der Klasse}{15,5 Milliarden Stellungen,\\ 1,58 Milliarden mit Hilfsmatt}
\end{minipage}\hfill
\begin{minipage}[t]{80mm}\vspace{0pt}\begin{tikzpicture}\begin{axis}[width=72mm,height=42mm,ybar,bar width=1.6pt,xmin=-1,xmax=22,ymin=0,ymax=338.3,axis y line*=left,axis x line*=bottom,xlabel={\scriptsize Halbzüge bis Matt},ylabel={\scriptsize Stellungen (Mio.)},tick label style={font=\tiny},label style={font=\scriptsize},enlarge x limits=false,ylabel shift=-3pt,xlabel shift=-3pt]\addplot[fill=ink!55,draw=none] coordinates {(0,0.000951) (1,0.004838) (2,0.0378) (3,0.2361) (4,0.879) (5,4.055) (6,7.906) (7,27.55) (8,42.31) (9,112.1) (10,144.6) (11,285.4) (12,227.6) (13,307.6) (14,164.5) (15,149.3) (16,59.81) (17,35.29) (18,5.421) (19,2.01) (20,0.00506) (21,0.000144)};\draw[densely dotted] (axis cs:17,0) -- (axis cs:17,338.3);\end{axis}\begin{axis}[width=72mm,height=42mm,axis y line*=right,axis x line=none,ymode=log,xmin=-1,xmax=22,ylabel={\scriptsize eindeutig (log)},tick label style={font=\tiny},label style={font=\scriptsize},ylabel shift=-3pt]\addplot[mark=*,mark size=0.7pt,draw=black,thin] coordinates {(0,951) (1,4838) (2,31099) (3,146991) (4,314860) (5,758065) (6,412486) (7,906267) (8,770257) (9,1407983) (10,1092873) (11,1187143) (12,355689) (13,188937) (14,65148) (15,18436) (16,7568) (17,1145)};\end{axis}\end{tikzpicture}\end{minipage}
\vspace{3mm}

\noindent\begin{minipage}[t]{60mm}\vspace{0pt}\sbox\hmboard{\chessboard[setfen=8/8/8/8/b1b4k/8/4b3/K4N2 w - - 0 1]}\usebox\hmboard\\[1mm]
\makebox[\wd\hmboard]{\sffamily\hspace{3mm}\textbf{h\#8.5}\hfill(2+4)}\end{minipage}\hfill
\begin{minipage}[t]{52mm}\vspace{0pt}
\textcolor{ink}{\sffamily\textsc{Lösung}}\par\smallskip\small
\begin{tabular}{@{}r@{\enspace}l@{\enspace}l@{}}
\textbf{1.} & \ldots & \symknight{}e3\\
\textbf{2.} & \symking{}g3 & \symking{}b2\\
\textbf{3.} & \symking{}f2 & \symking{}a3\\
\textbf{4.} & \symking{}e1 & \symking{}xa4\\
\textbf{5.} & \symking{}d2 & \symking{}b4\\
\textbf{6.} & \symking{}c1 & \symking{}c3\\
\textbf{7.} & \symking{}b1 & \symking{}d2\\
\textbf{8.} & \symking{}a1 & \symking{}c1\\
\textbf{9.} & \symbishop{}a2 & \symknight{}c2\#\\
\end{tabular}\par\vspace{3mm}
\qrcode[height=12mm]{https://helpman.komtera.lt/?fen=8/8/8/8/b1b4k/8/4b3/K4N2&moves=8.5}\\{\sffamily\tiny\color{ink}scannen zum Lösen}
\end{minipage}
\vspace{2mm}

{\raggedright\noindent\textcolor{ink}{\sffamily\textsc{Themen}}\enspace \themetag{reines Matt} \themetag{Mustermatt} \themetag{Selbstblock} \themetag{ohne Schach}\par}

\clearpage
\phantomsection\label{hm:KNvkbbn}
\begin{tcolorbox}[enhanced,colback=ink,colframe=ink,arc=0pt,boxrule=0pt,left=4mm,right=4mm,top=2.5mm,bottom=2.5mm]
\color{white}\sffamily{\LARGE\bfseries KNvkbbn}\hfill{\Large h\#8.5}\\[0.5mm]
{\small 6 Steine\quad Weiß am Zug}\hfill{\small Nr.~345\qquad\href{https://helpman.komtera.lt/?fen=6k1/8/8/8/8/2bb4/n1N5/1K6&moves=8.5}{\color{white}\bfseries lösen\,\ensuremath{\nearrow}}}
\end{tcolorbox}
\vspace{2mm}
\noindent\begin{minipage}[t]{36mm}\vspace{0pt}\sffamily\small\raggedright
\hmfigure{längste eindeutige}{{\Large h\#8.5}}\hmfigure{eindeutig bei dieser Länge}{{\Large 18} Stellungen}\hmfigure{längstes Matt}{{\Large h\#10.5}}\hmfigure{in der Klasse}{15,5 Milliarden Stellungen,\\ 4,87 Milliarden mit Hilfsmatt}
\end{minipage}\hfill
\begin{minipage}[t]{80mm}\vspace{0pt}\begin{tikzpicture}\begin{axis}[width=72mm,height=42mm,ybar,bar width=1.6pt,xmin=-1,xmax=22,ymin=0,ymax=1.249,axis y line*=left,axis x line*=bottom,xlabel={\scriptsize Halbzüge bis Matt},ylabel={\scriptsize Stellungen (Mrd.)},tick label style={font=\tiny},label style={font=\scriptsize},enlarge x limits=false,ylabel shift=-3pt,xlabel shift=-3pt]\addplot[fill=ink!55,draw=none] coordinates {(0,3.663e-06) (1,1.898e-05) (2,0.0001219) (3,0.0008033) (4,0.003704) (5,0.01697) (6,0.04617) (7,0.146) (8,0.2847) (9,0.6182) (10,0.743) (11,1.135) (12,0.7634) (13,0.7463) (14,0.2449) (15,0.107) (16,0.01484) (17,0.001872) (18,5.58e-05) (19,5.806e-06) (21,3e-09)};\draw[densely dotted] (axis cs:17,0) -- (axis cs:17,1.249);\end{axis}\begin{axis}[width=72mm,height=42mm,axis y line*=right,axis x line=none,ymode=log,xmin=-1,xmax=22,ylabel={\scriptsize eindeutig (log)},tick label style={font=\tiny},label style={font=\scriptsize},ylabel shift=-3pt]\addplot[mark=*,mark size=0.7pt,draw=black,thin] coordinates {(0,3663) (1,18972) (2,99364) (3,490315) (4,1598647) (5,3401909) (6,3189954) (7,5008194) (8,4547594) (9,5987792) (10,3359850) (11,2810444) (12,570762) (13,212436) (14,16829) (15,1284) (16,32) (17,18)};\end{axis}\end{tikzpicture}\end{minipage}
\vspace{3mm}

\noindent\begin{minipage}[t]{60mm}\vspace{0pt}\sbox\hmboard{\chessboard[setfen=6k1/8/8/8/8/2bb4/n1N5/1K6 w - - 0 1]}\usebox\hmboard\\[1mm]
\makebox[\wd\hmboard]{\sffamily\hspace{3mm}\textbf{h\#8.5}\hfill(2+4)}\end{minipage}\hfill
\begin{minipage}[t]{52mm}\vspace{0pt}
\textcolor{ink}{\sffamily\textsc{Lösung}}\par\smallskip\small
\begin{tabular}{@{}r@{\enspace}l@{\enspace}l@{}}
\textbf{1.} & \ldots & \symking{}xa2\\
\textbf{2.} & \symking{}f7 & \symking{}b3\\
\textbf{3.} & \symking{}e6 & \symking{}xc3\\
\textbf{4.} & \symking{}d5 & \symking{}b2\\
\textbf{5.} & \symking{}c4 & \symking{}a3\\
\textbf{6.} & \symking{}c3 & \symking{}a4\\
\textbf{7.} & \symking{}b2 & \symknight{}d4\\
\textbf{8.} & \symking{}a1 & \symking{}a3\\
\textbf{9.} & \symbishop{}b1 & \symknight{}b3\#\\
\end{tabular}\par\vspace{3mm}
\qrcode[height=12mm]{https://helpman.komtera.lt/?fen=6k1/8/8/8/8/2bb4/n1N5/1K6&moves=8.5}\\{\sffamily\tiny\color{ink}scannen zum Lösen}
\end{minipage}
\vspace{2mm}

{\raggedright\noindent\textcolor{ink}{\sffamily\textsc{Themen}}\enspace \themetag{reines Matt} \themetag{Mustermatt} \themetag{Idealmatt} \themetag{Switchback} \themetag{Selbstblock} \themetag{ohne Schach}\par\smallskip{\footnotesize\textit{Hinweis:} die Lösung beginnt mit einem Schlag.}\par}

\clearpage
\phantomsection\label{hm:KNvkbbp}
\begin{tcolorbox}[enhanced,colback=ink,colframe=ink,arc=0pt,boxrule=0pt,left=4mm,right=4mm,top=2.5mm,bottom=2.5mm]
\color{white}\sffamily{\LARGE\bfseries KNvkbbp}\hfill{\Large h\#8.5}\\[0.5mm]
{\small 6 Steine\quad Weiß am Zug}\hfill{\small Nr.~346\qquad\href{https://helpman.komtera.lt/?fen=N7/2p5/3b4/8/8/8/8/K2b3k&moves=8.5}{\color{white}\bfseries lösen\,\ensuremath{\nearrow}}}
\end{tcolorbox}
\vspace{2mm}
\noindent\begin{minipage}[t]{36mm}\vspace{0pt}\sffamily\small\raggedright
\hmfigure{längste eindeutige}{{\Large h\#8.5}}\hmfigure{eindeutig bei dieser Länge}{{\Large 1\,621} Stellungen}\hmfigure{längstes Matt}{{\Large h\#10.5}}\hmfigure{in der Klasse}{45,4 Milliarden Stellungen,\\ 15 Milliarden mit Hilfsmatt}
\end{minipage}\hfill
\begin{minipage}[t]{80mm}\vspace{0pt}\begin{tikzpicture}\begin{axis}[width=72mm,height=42mm,ybar,bar width=1.6pt,xmin=-1,xmax=22,ymin=0,ymax=3.093,axis y line*=left,axis x line*=bottom,xlabel={\scriptsize Halbzüge bis Matt},ylabel={\scriptsize Stellungen (Mrd.)},tick label style={font=\tiny},label style={font=\scriptsize},enlarge x limits=false,ylabel shift=-3pt,xlabel shift=-3pt]\addplot[fill=ink!55,draw=none] coordinates {(0,1.326e-05) (1,6.604e-05) (2,0.0003963) (3,0.002468) (4,0.01025) (5,0.04511) (6,0.1047) (7,0.3284) (8,0.5721) (9,1.285) (10,1.714) (11,2.812) (12,2.408) (13,2.648) (14,1.479) (15,1.065) (16,0.3609) (17,0.1371) (18,0.01393) (19,0.001944) (20,3.351e-06) (21,8.9e-08)};\draw[densely dotted] (axis cs:17,0) -- (axis cs:17,3.093);\end{axis}\begin{axis}[width=72mm,height=42mm,axis y line*=right,axis x line=none,ymode=log,xmin=-1,xmax=22,ylabel={\scriptsize eindeutig (log)},tick label style={font=\tiny},label style={font=\scriptsize},ylabel shift=-3pt]\addplot[mark=*,mark size=0.7pt,draw=black,thin] coordinates {(0,13260) (1,65137) (2,312330) (3,1448214) (4,4235436) (5,8444518) (6,6592839) (7,11062907) (8,10904768) (9,14694068) (10,12284351) (11,9829210) (12,3569916) (13,1465006) (14,577885) (15,126792) (16,21326) (17,1621)};\end{axis}\end{tikzpicture}\end{minipage}
\vspace{3mm}

\noindent\begin{minipage}[t]{60mm}\vspace{0pt}\sbox\hmboard{\chessboard[setfen=N7/2p5/3b4/8/8/8/8/K2b3k w - - 0 1]}\usebox\hmboard\\[1mm]
\makebox[\wd\hmboard]{\sffamily\hspace{3mm}\textbf{h\#8.5}\hfill(2+4)}\end{minipage}\hfill
\begin{minipage}[t]{52mm}\vspace{0pt}
\textcolor{ink}{\sffamily\textsc{Lösung}}\par\smallskip\small
\begin{tabular}{@{}r@{\enspace}l@{\enspace}l@{}}
\textbf{1.} & \ldots & \symking{}b2\\
\textbf{2.} & \symking{}g2 & \symking{}c3\\
\textbf{3.} & \symking{}f3 & \symking{}c4\\
\textbf{4.} & \symking{}e4 & \symking{}b5\\
\textbf{5.} & \symking{}d5 & \symking{}a5\\
\textbf{6.} & \symking{}c6 & \symknight{}xc7\\
\textbf{7.} & \symking{}b7 & \symknight{}d5\\
\textbf{8.} & \symking{}a8 & \symking{}a6\\
\textbf{9.} & \symbishop{}b8 & \symknight{}b6\#\\
\end{tabular}\par\vspace{3mm}
\qrcode[height=12mm]{https://helpman.komtera.lt/?fen=N7/2p5/3b4/8/8/8/8/K2b3k&moves=8.5}\\{\sffamily\tiny\color{ink}scannen zum Lösen}
\end{minipage}
\vspace{2mm}

{\raggedright\noindent\textcolor{ink}{\sffamily\textsc{Themen}}\enspace \themetag{reines Matt} \themetag{Mustermatt} \themetag{Selbstblock} \themetag{ohne Schach}\par}

\clearpage
\phantomsection\label{hm:KNvkbnn}
\begin{tcolorbox}[enhanced,colback=ink,colframe=ink,arc=0pt,boxrule=0pt,left=4mm,right=4mm,top=2.5mm,bottom=2.5mm]
\color{white}\sffamily{\LARGE\bfseries KNvkbnn}\hfill{\Large h\#8.5}\\[0.5mm]
{\small 6 Steine\quad Weiß am Zug}\hfill{\small Nr.~347\qquad\href{https://helpman.komtera.lt/?fen=7k/8/8/8/8/7b/8/KNn1n3&moves=8.5}{\color{white}\bfseries lösen\,\ensuremath{\nearrow}}}
\end{tcolorbox}
\vspace{2mm}
\noindent\begin{minipage}[t]{36mm}\vspace{0pt}\sffamily\small\raggedright
\hmfigure{längste eindeutige}{{\Large h\#8.5}}\hmfigure{eindeutig bei dieser Länge}{{\Large 9} Stellungen}\hmfigure{längstes Matt}{{\Large h\#9.5}}\hmfigure{in der Klasse}{15,5 Milliarden Stellungen,\\ 5 Milliarden mit Hilfsmatt}
\end{minipage}\hfill
\begin{minipage}[t]{80mm}\vspace{0pt}\begin{tikzpicture}\begin{axis}[width=72mm,height=42mm,ybar,bar width=1.6pt,xmin=-1,xmax=20,ymin=0,ymax=1.251,axis y line*=left,axis x line*=bottom,xlabel={\scriptsize Halbzüge bis Matt},ylabel={\scriptsize Stellungen (Mrd.)},tick label style={font=\tiny},label style={font=\scriptsize},enlarge x limits=false,ylabel shift=-3pt,xlabel shift=-3pt]\addplot[fill=ink!55,draw=none] coordinates {(0,4.33e-06) (1,2.287e-05) (2,0.0001201) (3,0.0008475) (4,0.004227) (5,0.01959) (6,0.05952) (7,0.1798) (8,0.3604) (9,0.7134) (10,0.8358) (11,1.137) (12,0.7448) (13,0.6447) (14,0.2098) (15,0.07645) (16,0.01113) (17,0.001051) (18,3.282e-05) (19,4.23e-07)};\draw[densely dotted] (axis cs:17,0) -- (axis cs:17,1.251);\end{axis}\begin{axis}[width=72mm,height=42mm,axis y line*=right,axis x line=none,ymode=log,xmin=-1,xmax=20,ylabel={\scriptsize eindeutig (log)},tick label style={font=\tiny},label style={font=\scriptsize},ylabel shift=-3pt]\addplot[mark=*,mark size=0.7pt,draw=black,thin] coordinates {(0,4330) (1,22813) (2,94478) (3,499039) (4,1940689) (5,3920886) (6,4351977) (7,5644435) (8,4910899) (9,5347494) (10,2732391) (11,1989668) (12,375664) (13,119260) (14,13121) (15,872) (16,47) (17,9)};\end{axis}\end{tikzpicture}\end{minipage}
\vspace{3mm}

\noindent\begin{minipage}[t]{60mm}\vspace{0pt}\sbox\hmboard{\chessboard[setfen=7k/8/8/8/8/7b/8/KNn1n3 w - - 0 1]}\usebox\hmboard\\[1mm]
\makebox[\wd\hmboard]{\sffamily\hspace{3mm}\textbf{h\#8.5}\hfill(2+4)}\end{minipage}\hfill
\begin{minipage}[t]{52mm}\vspace{0pt}
\textcolor{ink}{\sffamily\textsc{Lösung}}\par\smallskip\small
\begin{tabular}{@{}r@{\enspace}l@{\enspace}l@{}}
\textbf{1.} & \ldots & \symking{}b2\\
\textbf{2.} & \symking{}g7 & \symking{}c3\\
\textbf{3.} & \symking{}f6 & \symking{}d2\\
\textbf{4.} & \symking{}e5 & \symking{}d1\\
\textbf{5.} & \symking{}d4 & \symknight{}d2\\
\textbf{6.} & \symking{}c3 & \symking{}xe1\\
\textbf{7.} & \symking{}b2 & \symking{}d1\\
\textbf{8.} & \symking{}a1 & \symking{}c2\\
\textbf{9.} & \symknight{}a2 & \symknight{}b3\#\\
\end{tabular}\par\vspace{3mm}
\qrcode[height=12mm]{https://helpman.komtera.lt/?fen=7k/8/8/8/8/7b/8/KNn1n3&moves=8.5}\\{\sffamily\tiny\color{ink}scannen zum Lösen}
\end{minipage}
\vspace{2mm}

{\raggedright\noindent\textcolor{ink}{\sffamily\textsc{Themen}}\enspace \themetag{reines Matt} \themetag{Mustermatt} \themetag{Switchback} \themetag{Selbstblock} \themetag{ohne Schach}\par}

\clearpage
\phantomsection\label{hm:KNvkbnp}
\begin{tcolorbox}[enhanced,colback=ink,colframe=ink,arc=0pt,boxrule=0pt,left=4mm,right=4mm,top=2.5mm,bottom=2.5mm]
\color{white}\sffamily{\LARGE\bfseries KNvkbnp}\hfill{\Large h\#8.5}\\[0.5mm]
{\small 6 Steine\quad Weiß am Zug}\hfill{\small Nr.~348\qquad\href{https://helpman.komtera.lt/?fen=8/7k/7n/4p3/4b3/8/N7/K7&moves=8.5}{\color{white}\bfseries lösen\,\ensuremath{\nearrow}}}
\end{tcolorbox}
\vspace{2mm}
\noindent\begin{minipage}[t]{36mm}\vspace{0pt}\sffamily\small\raggedright
\hmfigure{längste eindeutige}{{\Large h\#8.5}}\hmfigure{eindeutig bei dieser Länge}{{\Large 2} Stellungen}\hmfigure{längstes Matt}{{\Large h\#10}}\hmfigure{in der Klasse}{45,4 Milliarden Stellungen,\\ 30,8 Milliarden mit Hilfsmatt}
\end{minipage}\hfill
\begin{minipage}[t]{80mm}\vspace{0pt}\begin{tikzpicture}\begin{axis}[width=72mm,height=42mm,ybar,bar width=1.6pt,xmin=-1,xmax=21,ymin=0,ymax=7.36,axis y line*=left,axis x line*=bottom,xlabel={\scriptsize Halbzüge bis Matt},ylabel={\scriptsize Stellungen (Mrd.)},tick label style={font=\tiny},label style={font=\scriptsize},enlarge x limits=false,ylabel shift=-3pt,xlabel shift=-3pt]\addplot[fill=ink!55,draw=none] coordinates {(0,3.099e-05) (1,0.0001565) (2,0.0007416) (3,0.004892) (4,0.02298) (5,0.1013) (6,0.2932) (7,0.859) (8,1.801) (9,3.56) (10,4.824) (11,6.691) (12,5.186) (13,4.643) (14,1.862) (15,0.7733) (16,0.1584) (17,0.01973) (18,0.001445) (19,2.428e-05) (20,3.4e-08)};\draw[densely dotted] (axis cs:17,0) -- (axis cs:17,7.36);\end{axis}\begin{axis}[width=72mm,height=42mm,axis y line*=right,axis x line=none,ymode=log,xmin=-1,xmax=21,ylabel={\scriptsize eindeutig (log)},tick label style={font=\tiny},label style={font=\scriptsize},ylabel shift=-3pt]\addplot[mark=*,mark size=0.7pt,draw=black,thin] coordinates {(0,30995) (1,154173) (2,564732) (3,2759018) (4,10752760) (5,19436558) (6,21253252) (7,26289298) (8,27732211) (9,29205950) (10,19508302) (11,13076424) (12,3515068) (13,1068214) (14,163447) (15,9659) (16,859) (17,2)};\end{axis}\end{tikzpicture}\end{minipage}
\vspace{3mm}

\noindent\begin{minipage}[t]{60mm}\vspace{0pt}\sbox\hmboard{\chessboard[setfen=8/7k/7n/4p3/4b3/8/N7/K7 w - - 0 1]}\usebox\hmboard\\[1mm]
\makebox[\wd\hmboard]{\sffamily\hspace{3mm}\textbf{h\#8.5}\hfill(2+4)}\end{minipage}\hfill
\begin{minipage}[t]{52mm}\vspace{0pt}
\textcolor{ink}{\sffamily\textsc{Lösung}}\par\smallskip\small
\begin{tabular}{@{}r@{\enspace}l@{\enspace}l@{}}
\textbf{1.} & \ldots & \symking{}b2\\
\textbf{2.} & \symbishop{}d5 & \symking{}c2\\
\textbf{3.} & e4 & \symking{}c3\\
\textbf{4.} & e3 & \symking{}d4\\
\textbf{5.} & e2 & \symking{}e5\\
\textbf{6.} & e1=\symrook{}+ & \symking{}f6\\
\textbf{7.} & \symrook{}e8 & \symknight{}c3\\
\textbf{8.} & \symrook{}h8 & \symknight{}e4\\
\textbf{9.} & \symbishop{}g8 & \symknight{}g5\#\\
\end{tabular}\par\vspace{3mm}
\qrcode[height=12mm]{https://helpman.komtera.lt/?fen=8/7k/7n/4p3/4b3/8/N7/K7&moves=8.5}\\{\sffamily\tiny\color{ink}scannen zum Lösen}
\end{minipage}
\vspace{2mm}

{\raggedright\noindent\textcolor{ink}{\sffamily\textsc{Themen}}\enspace \themetag{Satzspiel} \themetag{reines Matt} \themetag{Mustermatt} \themetag{Idealmatt} \themetag{Umwandlung} \themetag{Unterverwandlung} \themetag{Selbstblock} \themetag{ohne Schlag}\par}

\clearpage
\phantomsection\label{hm:KNvknnn}
\begin{tcolorbox}[enhanced,colback=ink,colframe=ink,arc=0pt,boxrule=0pt,left=4mm,right=4mm,top=2.5mm,bottom=2.5mm]
\color{white}\sffamily{\LARGE\bfseries KNvknnn}\hfill{\Large h\#8.5}\\[0.5mm]
{\small 6 Steine\quad Weiß am Zug}\hfill{\small Nr.~349\qquad\href{https://helpman.komtera.lt/?fen=7k/8/1n6/8/8/8/8/KNn1n3&moves=8.5}{\color{white}\bfseries lösen\,\ensuremath{\nearrow}}}
\end{tcolorbox}
\vspace{2mm}
\noindent\begin{minipage}[t]{36mm}\vspace{0pt}\sffamily\small\raggedright
\hmfigure{längste eindeutige}{{\Large h\#8.5}}\hmfigure{eindeutig bei dieser Länge}{{\Large 4} Stellungen}\hmfigure{längstes Matt}{{\Large h\#9.5}}\hmfigure{in der Klasse}{15,5 Milliarden Stellungen,\\ 1,7 Milliarden mit Hilfsmatt}
\end{minipage}\hfill
\begin{minipage}[t]{80mm}\vspace{0pt}\begin{tikzpicture}\begin{axis}[width=72mm,height=42mm,ybar,bar width=1.6pt,xmin=-1,xmax=20,ymin=0,ymax=394.6,axis y line*=left,axis x line*=bottom,xlabel={\scriptsize Halbzüge bis Matt},ylabel={\scriptsize Stellungen (Mio.)},tick label style={font=\tiny},label style={font=\scriptsize},enlarge x limits=false,ylabel shift=-3pt,xlabel shift=-3pt]\addplot[fill=ink!55,draw=none] coordinates {(0,0.001579) (1,0.008542) (2,0.03543) (3,0.2777) (4,1.339) (5,6.554) (6,21.23) (7,63.28) (8,124.4) (9,232.3) (10,275.9) (11,358.7) (12,261.1) (13,228.8) (14,87.61) (15,33.81) (16,6.036) (17,0.6139) (18,0.01361) (19,0.000167)};\draw[densely dotted] (axis cs:17,0) -- (axis cs:17,394.6);\end{axis}\begin{axis}[width=72mm,height=42mm,axis y line*=right,axis x line=none,ymode=log,xmin=-1,xmax=20,ylabel={\scriptsize eindeutig (log)},tick label style={font=\tiny},label style={font=\scriptsize},ylabel shift=-3pt]\addplot[mark=*,mark size=0.7pt,draw=black,thin] coordinates {(0,1579) (1,8513) (2,25887) (3,155906) (4,566863) (5,1275953) (6,1547552) (7,1737013) (8,1238068) (9,1170590) (10,580346) (11,476954) (12,97491) (13,35977) (14,5211) (15,334) (16,10) (17,4)};\end{axis}\end{tikzpicture}\end{minipage}
\vspace{3mm}

\noindent\begin{minipage}[t]{60mm}\vspace{0pt}\sbox\hmboard{\chessboard[setfen=7k/8/1n6/8/8/8/8/KNn1n3 w - - 0 1]}\usebox\hmboard\\[1mm]
\makebox[\wd\hmboard]{\sffamily\hspace{3mm}\textbf{h\#8.5}\hfill(2+4)}\end{minipage}\hfill
\begin{minipage}[t]{52mm}\vspace{0pt}
\textcolor{ink}{\sffamily\textsc{Lösung}}\par\smallskip\small
\begin{tabular}{@{}r@{\enspace}l@{\enspace}l@{}}
\textbf{1.} & \ldots & \symking{}b2\\
\textbf{2.} & \symking{}g7 & \symking{}c3\\
\textbf{3.} & \symking{}f6 & \symking{}d2\\
\textbf{4.} & \symking{}e5 & \symking{}d1\\
\textbf{5.} & \symking{}d4 & \symknight{}d2\\
\textbf{6.} & \symking{}c3 & \symking{}xe1\\
\textbf{7.} & \symking{}b2 & \symking{}d1\\
\textbf{8.} & \symking{}a1 & \symking{}c2\\
\textbf{9.} & \symknight{}a2 & \symknight{}b3\#\\
\end{tabular}\par\vspace{3mm}
\qrcode[height=12mm]{https://helpman.komtera.lt/?fen=7k/8/1n6/8/8/8/8/KNn1n3&moves=8.5}\\{\sffamily\tiny\color{ink}scannen zum Lösen}
\end{minipage}
\vspace{2mm}

{\raggedright\noindent\textcolor{ink}{\sffamily\textsc{Themen}}\enspace \themetag{reines Matt} \themetag{Mustermatt} \themetag{Switchback} \themetag{Selbstblock} \themetag{ohne Schach}\par}

\clearpage
\phantomsection\label{hm:KNvknnp}
\begin{tcolorbox}[enhanced,colback=ink,colframe=ink,arc=0pt,boxrule=0pt,left=4mm,right=4mm,top=2.5mm,bottom=2.5mm]
\color{white}\sffamily{\LARGE\bfseries KNvknnp}\hfill{\Large h\#8.5}\\[0.5mm]
{\small 6 Steine\quad Weiß am Zug}\hfill{\small Nr.~350\qquad\href{https://helpman.komtera.lt/?fen=7k/8/8/8/8/5p2/8/KNn1n3&moves=8.5}{\color{white}\bfseries lösen\,\ensuremath{\nearrow}}}
\end{tcolorbox}
\vspace{2mm}
\noindent\begin{minipage}[t]{36mm}\vspace{0pt}\sffamily\small\raggedright
\hmfigure{längste eindeutige}{{\Large h\#8.5}}\hmfigure{eindeutig bei dieser Länge}{{\Large 211} Stellungen}\hmfigure{längstes Matt}{{\Large h\#10}}\hmfigure{in der Klasse}{45,4 Milliarden Stellungen,\\ 15,8 Milliarden mit Hilfsmatt}
\end{minipage}\hfill
\begin{minipage}[t]{80mm}\vspace{0pt}\begin{tikzpicture}\begin{axis}[width=72mm,height=42mm,ybar,bar width=1.6pt,xmin=-1,xmax=21,ymin=0,ymax=3.386,axis y line*=left,axis x line*=bottom,xlabel={\scriptsize Halbzüge bis Matt},ylabel={\scriptsize Stellungen (Mrd.)},tick label style={font=\tiny},label style={font=\scriptsize},enlarge x limits=false,ylabel shift=-3pt,xlabel shift=-3pt]\addplot[fill=ink!55,draw=none] coordinates {(0,1.693e-05) (1,8.691e-05) (2,0.0003028) (3,0.002208) (4,0.01011) (5,0.04666) (6,0.1513) (7,0.4308) (8,0.912) (9,1.688) (10,2.32) (11,3.078) (12,2.699) (13,2.505) (14,1.213) (15,0.5566) (16,0.1278) (17,0.01843) (18,0.00108) (19,2.762e-05) (20,1e-09)};\draw[densely dotted] (axis cs:17,0) -- (axis cs:17,3.386);\end{axis}\begin{axis}[width=72mm,height=42mm,axis y line*=right,axis x line=none,ymode=log,xmin=-1,xmax=21,ylabel={\scriptsize eindeutig (log)},tick label style={font=\tiny},label style={font=\scriptsize},ylabel shift=-3pt]\addplot[mark=*,mark size=0.7pt,draw=black,thin] coordinates {(0,16934) (1,85399) (2,210352) (3,1158529) (4,4497896) (5,8824657) (6,11088217) (7,11226766) (8,9845643) (9,9288667) (10,6259800) (11,4745062) (12,1384952) (13,481931) (14,84318) (15,6313) (16,770) (17,211)};\end{axis}\end{tikzpicture}\end{minipage}
\vspace{3mm}

\noindent\begin{minipage}[t]{60mm}\vspace{0pt}\sbox\hmboard{\chessboard[setfen=7k/8/8/8/8/5p2/8/KNn1n3 w - - 0 1]}\usebox\hmboard\\[1mm]
\makebox[\wd\hmboard]{\sffamily\hspace{3mm}\textbf{h\#8.5}\hfill(2+4)}\end{minipage}\hfill
\begin{minipage}[t]{52mm}\vspace{0pt}
\textcolor{ink}{\sffamily\textsc{Lösung}}\par\smallskip\small
\begin{tabular}{@{}r@{\enspace}l@{\enspace}l@{}}
\textbf{1.} & \ldots & \symking{}b2\\
\textbf{2.} & \symking{}g7 & \symking{}c3\\
\textbf{3.} & \symking{}f6 & \symking{}d2\\
\textbf{4.} & \symking{}e5 & \symking{}d1\\
\textbf{5.} & \symking{}d4 & \symknight{}d2\\
\textbf{6.} & \symking{}c3 & \symking{}xe1\\
\textbf{7.} & \symking{}b2 & \symking{}d1\\
\textbf{8.} & \symking{}a1 & \symking{}c2\\
\textbf{9.} & \symknight{}a2 & \symknight{}b3\#\\
\end{tabular}\par\vspace{3mm}
\qrcode[height=12mm]{https://helpman.komtera.lt/?fen=7k/8/8/8/8/5p2/8/KNn1n3&moves=8.5}\\{\sffamily\tiny\color{ink}scannen zum Lösen}
\end{minipage}
\vspace{2mm}

{\raggedright\noindent\textcolor{ink}{\sffamily\textsc{Themen}}\enspace \themetag{reines Matt} \themetag{Mustermatt} \themetag{Switchback} \themetag{Selbstblock} \themetag{ohne Schach}\par}

\clearpage
\phantomsection\label{hm:KNvkqnn}
\begin{tcolorbox}[enhanced,colback=ink,colframe=ink,arc=0pt,boxrule=0pt,left=4mm,right=4mm,top=2.5mm,bottom=2.5mm]
\color{white}\sffamily{\LARGE\bfseries KNvkqnn}\hfill{\Large h\#8.5}\\[0.5mm]
{\small 6 Steine\quad Weiß am Zug}\hfill{\small Nr.~351\qquad\href{https://helpman.komtera.lt/?fen=6k1/8/8/8/8/nq6/1n6/1KN5&moves=8.5}{\color{white}\bfseries lösen\,\ensuremath{\nearrow}}}
\end{tcolorbox}
\vspace{2mm}
\noindent\begin{minipage}[t]{36mm}\vspace{0pt}\sffamily\small\raggedright
\hmfigure{längste eindeutige}{{\Large h\#8.5}}\hmfigure{eindeutig bei dieser Länge}{{\Large 7} Stellungen}\hmfigure{längstes Matt}{{\Large h\#9.5}}\hmfigure{in der Klasse}{15,5 Milliarden Stellungen,\\ 4,49 Milliarden mit Hilfsmatt}
\end{minipage}\hfill
\begin{minipage}[t]{80mm}\vspace{0pt}\begin{tikzpicture}\begin{axis}[width=72mm,height=42mm,ybar,bar width=1.6pt,xmin=-1,xmax=20,ymin=0,ymax=1.244,axis y line*=left,axis x line*=bottom,xlabel={\scriptsize Halbzüge bis Matt},ylabel={\scriptsize Stellungen (Mrd.)},tick label style={font=\tiny},label style={font=\scriptsize},enlarge x limits=false,ylabel shift=-3pt,xlabel shift=-3pt]\addplot[fill=ink!55,draw=none] coordinates {(0,1.801e-06) (1,9.938e-06) (2,6.009e-05) (3,0.000474) (4,0.002619) (5,0.01385) (6,0.05002) (7,0.1768) (8,0.3345) (9,0.7813) (10,0.6917) (11,1.131) (12,0.5219) (13,0.5902) (14,0.1215) (15,0.06615) (16,0.005723) (17,0.001406) (18,3.826e-06) (19,7.343e-06)};\draw[densely dotted] (axis cs:17,0) -- (axis cs:17,1.244);\end{axis}\begin{axis}[width=72mm,height=42mm,axis y line*=right,axis x line=none,ymode=log,xmin=-1,xmax=20,ylabel={\scriptsize eindeutig (log)},tick label style={font=\tiny},label style={font=\scriptsize},ylabel shift=-3pt]\addplot[mark=*,mark size=0.7pt,draw=black,thin] coordinates {(0,1801) (1,9938) (2,45367) (3,274586) (4,975438) (5,2790950) (6,3824666) (7,6863962) (8,3738471) (9,5115940) (10,1276722) (11,1613172) (12,156746) (13,119432) (14,5343) (15,1159) (16,10) (17,7)};\end{axis}\end{tikzpicture}\end{minipage}
\vspace{3mm}

\noindent\begin{minipage}[t]{60mm}\vspace{0pt}\sbox\hmboard{\chessboard[setfen=6k1/8/8/8/8/nq6/1n6/1KN5 w - - 0 1]}\usebox\hmboard\\[1mm]
\makebox[\wd\hmboard]{\sffamily\hspace{3mm}\textbf{h\#8.5}\hfill(2+4)}\end{minipage}\hfill
\begin{minipage}[t]{52mm}\vspace{0pt}
\textcolor{ink}{\sffamily\textsc{Lösung}}\par\smallskip\small
\begin{tabular}{@{}r@{\enspace}l@{\enspace}l@{}}
\textbf{1.} & \ldots & \symking{}a1\\
\textbf{2.} & \symqueen{}d1 & \symking{}a2\\
\textbf{3.} & \symking{}f7 & \symking{}xa3\\
\textbf{4.} & \symking{}e6 & \symking{}b4\\
\textbf{5.} & \symking{}d5 & \symking{}a5\\
\textbf{6.} & \symking{}c4 & \symknight{}e2\\
\textbf{7.} & \symking{}b3 & \symking{}b5\\
\textbf{8.} & \symking{}a2 & \symking{}b4\\
\textbf{9.} & \symqueen{}a1 & \symknight{}c3\#\\
\end{tabular}\par\vspace{3mm}
\qrcode[height=12mm]{https://helpman.komtera.lt/?fen=6k1/8/8/8/8/nq6/1n6/1KN5&moves=8.5}\\{\sffamily\tiny\color{ink}scannen zum Lösen}
\end{minipage}
\vspace{2mm}

{\raggedright\noindent\textcolor{ink}{\sffamily\textsc{Themen}}\enspace \themetag{reines Matt} \themetag{Mustermatt} \themetag{Idealmatt} \themetag{Rundlauf} \themetag{Selbstblock} \themetag{ohne Schach}\par\smallskip{\footnotesize\textit{Hinweis:} die Partei am Zug steht im Diagramm im Schach.}\par}

\clearpage
\phantomsection\label{hm:KNvkqqn}
\begin{tcolorbox}[enhanced,colback=ink,colframe=ink,arc=0pt,boxrule=0pt,left=4mm,right=4mm,top=2.5mm,bottom=2.5mm]
\color{white}\sffamily{\LARGE\bfseries KNvkqqn}\hfill{\Large h\#8.5}\\[0.5mm]
{\small 6 Steine\quad Weiß am Zug}\hfill{\small Nr.~352\qquad\href{https://helpman.komtera.lt/?fen=6k1/8/8/8/8/3n4/5q2/1KNq4&moves=8.5}{\color{white}\bfseries lösen\,\ensuremath{\nearrow}}}
\end{tcolorbox}
\vspace{2mm}
\noindent\begin{minipage}[t]{36mm}\vspace{0pt}\sffamily\small\raggedright
\hmfigure{längste eindeutige}{{\Large h\#8.5}}\hmfigure{eindeutig bei dieser Länge}{{\Large 34} Stellungen}\hmfigure{längstes Matt}{{\Large h\#9.5}}\hmfigure{in der Klasse}{15,5 Milliarden Stellungen,\\ 3,9 Milliarden mit Hilfsmatt}
\end{minipage}\hfill
\begin{minipage}[t]{80mm}\vspace{0pt}\begin{tikzpicture}\begin{axis}[width=72mm,height=42mm,ybar,bar width=1.6pt,xmin=-1,xmax=20,ymin=0,ymax=1.118,axis y line*=left,axis x line*=bottom,xlabel={\scriptsize Halbzüge bis Matt},ylabel={\scriptsize Stellungen (Mrd.)},tick label style={font=\tiny},label style={font=\scriptsize},enlarge x limits=false,ylabel shift=-3pt,xlabel shift=-3pt]\addplot[fill=ink!55,draw=none] coordinates {(0,4.77e-07) (1,2.7e-06) (2,2.057e-05) (3,0.0001669) (4,0.0008733) (5,0.005148) (6,0.01777) (7,0.07708) (8,0.1424) (9,0.4585) (10,0.4065) (11,1.016) (12,0.4636) (13,0.8882) (14,0.1848) (15,0.2102) (16,0.01705) (17,0.007104) (18,0.0001353) (19,3.799e-05)};\draw[densely dotted] (axis cs:17,0) -- (axis cs:17,1.118);\end{axis}\begin{axis}[width=72mm,height=42mm,axis y line*=right,axis x line=none,ymode=log,xmin=-1,xmax=20,ylabel={\scriptsize eindeutig (log)},tick label style={font=\tiny},label style={font=\scriptsize},ylabel shift=-3pt]\addplot[mark=*,mark size=0.7pt,draw=black,thin] coordinates {(0,477) (1,2700) (2,15076) (3,96798) (4,311912) (5,1095421) (6,1470386) (7,3675255) (8,1967852) (9,4173497) (10,1122542) (11,2128387) (12,245672) (13,386376) (14,16364) (15,10066) (16,43) (17,34)};\end{axis}\end{tikzpicture}\end{minipage}
\vspace{3mm}

\noindent\begin{minipage}[t]{60mm}\vspace{0pt}\sbox\hmboard{\chessboard[setfen=6k1/8/8/8/8/3n4/5q2/1KNq4 w - - 0 1]}\usebox\hmboard\\[1mm]
\makebox[\wd\hmboard]{\sffamily\hspace{3mm}\textbf{h\#8.5}\hfill(2+4)}\end{minipage}\hfill
\begin{minipage}[t]{52mm}\vspace{0pt}
\textcolor{ink}{\sffamily\textsc{Lösung}}\par\smallskip\small
\begin{tabular}{@{}r@{\enspace}l@{\enspace}l@{}}
\textbf{1.} & \ldots & \symking{}a1\\
\textbf{2.} & \symknight{}b2 & \symking{}a2\\
\textbf{3.} & \symking{}f7 & \symking{}a3\\
\textbf{4.} & \symking{}e6 & \symking{}b4\\
\textbf{5.} & \symking{}d5 & \symking{}a5\\
\textbf{6.} & \symking{}c4 & \symknight{}e2\\
\textbf{7.} & \symking{}b3 & \symking{}b5\\
\textbf{8.} & \symking{}a2 & \symking{}b4\\
\textbf{9.} & \symqueen{}a1 & \symknight{}c3\#\\
\end{tabular}\par\vspace{3mm}
\qrcode[height=12mm]{https://helpman.komtera.lt/?fen=6k1/8/8/8/8/3n4/5q2/1KNq4&moves=8.5}\\{\sffamily\tiny\color{ink}scannen zum Lösen}
\end{minipage}
\vspace{2mm}

{\raggedright\noindent\textcolor{ink}{\sffamily\textsc{Themen}}\enspace \themetag{Satzspiel} \themetag{reines Matt} \themetag{Mustermatt} \themetag{Rundlauf} \themetag{Selbstblock} \themetag{ohne Schlag} \themetag{ohne Schach}\par}

\clearpage
\phantomsection\label{hm:KNvkqqr}
\begin{tcolorbox}[enhanced,colback=ink,colframe=ink,arc=0pt,boxrule=0pt,left=4mm,right=4mm,top=2.5mm,bottom=2.5mm]
\color{white}\sffamily{\LARGE\bfseries KNvkqqr}\hfill{\Large h\#8.5}\\[0.5mm]
{\small 6 Steine\quad Weiß am Zug}\hfill{\small Nr.~353\qquad\href{https://helpman.komtera.lt/?fen=7k/6N1/8/8/8/8/1r1q3q/1K6&moves=8.5}{\color{white}\bfseries lösen\,\ensuremath{\nearrow}}}
\end{tcolorbox}
\vspace{2mm}
\noindent\begin{minipage}[t]{36mm}\vspace{0pt}\sffamily\small\raggedright
\hmfigure{längste eindeutige}{{\Large h\#8.5}}\hmfigure{eindeutig bei dieser Länge}{{\Large 2} Stellungen}\hmfigure{längstes Matt}{{\Large h\#9.5}}\hmfigure{in der Klasse}{15,5 Milliarden Stellungen,\\ 3,58 Milliarden mit Hilfsmatt}
\end{minipage}\hfill
\begin{minipage}[t]{80mm}\vspace{0pt}\begin{tikzpicture}\begin{axis}[width=72mm,height=42mm,ybar,bar width=1.6pt,xmin=-1,xmax=20,ymin=0,ymax=1.153,axis y line*=left,axis x line*=bottom,xlabel={\scriptsize Halbzüge bis Matt},ylabel={\scriptsize Stellungen (Mrd.)},tick label style={font=\tiny},label style={font=\scriptsize},enlarge x limits=false,ylabel shift=-3pt,xlabel shift=-3pt]\addplot[fill=ink!55,draw=none] coordinates {(0,5.72e-07) (1,2.459e-06) (2,3.462e-05) (3,0.0002413) (4,0.00153) (5,0.008423) (6,0.02696) (7,0.1266) (8,0.196) (9,0.6867) (10,0.4052) (11,1.048) (12,0.3162) (13,0.5924) (14,0.0785) (15,0.08739) (16,0.003394) (17,0.001905) (18,1.24e-06) (19,7.63e-07)};\draw[densely dotted] (axis cs:17,0) -- (axis cs:17,1.153);\end{axis}\begin{axis}[width=72mm,height=42mm,axis y line*=right,axis x line=none,ymode=log,xmin=-1,xmax=20,ylabel={\scriptsize eindeutig (log)},tick label style={font=\tiny},label style={font=\scriptsize},ylabel shift=-3pt]\addplot[mark=*,mark size=0.7pt,draw=black,thin] coordinates {(0,572) (1,2459) (2,29344) (3,169159) (4,492959) (5,1975625) (6,2427971) (7,8981707) (8,4460215) (9,11386841) (10,2261558) (11,4419676) (12,399766) (13,585472) (14,23024) (15,17900) (16,18) (17,2)};\end{axis}\end{tikzpicture}\end{minipage}
\vspace{3mm}

\noindent\begin{minipage}[t]{60mm}\vspace{0pt}\sbox\hmboard{\chessboard[setfen=7k/6N1/8/8/8/8/1r1q3q/1K6 w - - 0 1]}\usebox\hmboard\\[1mm]
\makebox[\wd\hmboard]{\sffamily\hspace{3mm}\textbf{h\#8.5}\hfill(2+4)}\end{minipage}\hfill
\begin{minipage}[t]{52mm}\vspace{0pt}
\textcolor{ink}{\sffamily\textsc{Lösung}}\par\smallskip\small
\begin{tabular}{@{}r@{\enspace}l@{\enspace}l@{}}
\textbf{1.} & \ldots & \symking{}a1\\
\textbf{2.} & \symrook{}b6 & \symknight{}e8\\
\textbf{3.} & \symrook{}h6 & \symking{}b1\\
\textbf{4.} & \symqueen{}c1+ & \symking{}xc1\\
\textbf{5.} & \symqueen{}h1+ & \symking{}d2\\
\textbf{6.} & \symking{}h7 & \symking{}e3\\
\textbf{7.} & \symking{}g6 & \symking{}f4\\
\textbf{8.} & \symking{}h5 & \symking{}f5\\
\textbf{9.} & \symqueen{}h4 & \symknight{}g7\#\\
\end{tabular}\par\vspace{3mm}
\qrcode[height=12mm]{https://helpman.komtera.lt/?fen=7k/6N1/8/8/8/8/1r1q3q/1K6&moves=8.5}\\{\sffamily\tiny\color{ink}scannen zum Lösen}
\end{minipage}
\vspace{2mm}

{\raggedright\noindent\textcolor{ink}{\sffamily\textsc{Themen}}\enspace \themetag{reines Matt} \themetag{Mustermatt} \themetag{Idealmatt} \themetag{Switchback} \themetag{Selbstblock}\par\smallskip{\footnotesize\textit{Hinweis:} die Partei am Zug steht im Diagramm im Schach.}\par}

\clearpage
\phantomsection\label{hm:KNvkqrn}
\begin{tcolorbox}[enhanced,colback=ink,colframe=ink,arc=0pt,boxrule=0pt,left=4mm,right=4mm,top=2.5mm,bottom=2.5mm]
\color{white}\sffamily{\LARGE\bfseries KNvkqrn}\hfill{\Large h\#8.5}\\[0.5mm]
{\small 6 Steine\quad Weiß am Zug}\hfill{\small Nr.~354\qquad\href{https://helpman.komtera.lt/?fen=8/7k/8/8/8/r7/1nq5/KN6&moves=8.5}{\color{white}\bfseries lösen\,\ensuremath{\nearrow}}}
\end{tcolorbox}
\vspace{2mm}
\noindent\begin{minipage}[t]{36mm}\vspace{0pt}\sffamily\small\raggedright
\hmfigure{längste eindeutige}{{\Large h\#8.5}}\hmfigure{eindeutig bei dieser Länge}{{\Large 13} Stellungen}\hmfigure{längstes Matt}{{\Large h\#9.5}}\hmfigure{in der Klasse}{15,5 Milliarden Stellungen,\\ 8,42 Milliarden mit Hilfsmatt}
\end{minipage}\hfill
\begin{minipage}[t]{80mm}\vspace{0pt}\begin{tikzpicture}\begin{axis}[width=72mm,height=42mm,ybar,bar width=1.6pt,xmin=-1,xmax=20,ymin=0,ymax=2.422,axis y line*=left,axis x line*=bottom,xlabel={\scriptsize Halbzüge bis Matt},ylabel={\scriptsize Stellungen (Mrd.)},tick label style={font=\tiny},label style={font=\scriptsize},enlarge x limits=false,ylabel shift=-3pt,xlabel shift=-3pt]\addplot[fill=ink!55,draw=none] coordinates {(0,3.444e-06) (1,1.647e-05) (2,0.0001662) (3,0.00119) (4,0.007142) (5,0.03673) (6,0.1223) (7,0.4606) (8,0.7483) (9,1.864) (10,1.255) (11,2.202) (12,0.737) (13,0.8329) (14,0.0908) (15,0.05561) (16,0.0008401) (17,0.001167) (19,3.873e-06)};\draw[densely dotted] (axis cs:17,0) -- (axis cs:17,2.422);\end{axis}\begin{axis}[width=72mm,height=42mm,axis y line*=right,axis x line=none,ymode=log,xmin=-1,xmax=20,ylabel={\scriptsize eindeutig (log)},tick label style={font=\tiny},label style={font=\scriptsize},ylabel shift=-3pt]\addplot[mark=*,mark size=0.7pt,draw=black,thin] coordinates {(0,3444) (1,16444) (2,139449) (3,800096) (4,2554099) (5,8233155) (6,10560866) (7,25719084) (8,13916580) (9,23179034) (10,5874168) (11,7762810) (12,739465) (13,482364) (14,3882) (15,2674) (17,13)};\end{axis}\end{tikzpicture}\end{minipage}
\vspace{3mm}

\noindent\begin{minipage}[t]{60mm}\vspace{0pt}\sbox\hmboard{\chessboard[setfen=8/7k/8/8/8/r7/1nq5/KN6 w - - 0 1]}\usebox\hmboard\\[1mm]
\makebox[\wd\hmboard]{\sffamily\hspace{3mm}\textbf{h\#8.5}\hfill(2+4)}\end{minipage}\hfill
\begin{minipage}[t]{52mm}\vspace{0pt}
\textcolor{ink}{\sffamily\textsc{Lösung}}\par\smallskip\small
\begin{tabular}{@{}r@{\enspace}l@{\enspace}l@{}}
\textbf{1.} & \ldots & \symknight{}xa3\\
\textbf{2.} & \symqueen{}a4 & \symking{}b1\\
\textbf{3.} & \symking{}g6 & \symking{}c1\\
\textbf{4.} & \symking{}f5 & \symking{}d2\\
\textbf{5.} & \symking{}e4 & \symking{}e1\\
\textbf{6.} & \symking{}d3 & \symknight{}b5\\
\textbf{7.} & \symking{}c2 & \symking{}e2\\
\textbf{8.} & \symking{}b1 & \symking{}d2\\
\textbf{9.} & \symqueen{}a1 & \symknight{}c3\#\\
\end{tabular}\par\vspace{3mm}
\qrcode[height=12mm]{https://helpman.komtera.lt/?fen=8/7k/8/8/8/r7/1nq5/KN6&moves=8.5}\\{\sffamily\tiny\color{ink}scannen zum Lösen}
\end{minipage}
\vspace{2mm}

{\raggedright\noindent\textcolor{ink}{\sffamily\textsc{Themen}}\enspace \themetag{reines Matt} \themetag{Mustermatt} \themetag{Idealmatt} \themetag{Rundlauf} \themetag{Selbstblock} \themetag{ohne Schach}\par\smallskip{\footnotesize\textit{Hinweis:} die Partei am Zug steht im Diagramm im Schach; die Lösung beginnt mit einem Schlag.}\par}

\clearpage
\phantomsection\label{hm:KNvkrbb}
\begin{tcolorbox}[enhanced,colback=ink,colframe=ink,arc=0pt,boxrule=0pt,left=4mm,right=4mm,top=2.5mm,bottom=2.5mm]
\color{white}\sffamily{\LARGE\bfseries KNvkrbb}\hfill{\Large h\#8.5}\\[0.5mm]
{\small 6 Steine\quad Weiß am Zug}\hfill{\small Nr.~355\qquad\href{https://helpman.komtera.lt/?fen=8/8/8/8/8/1b6/8/KrN2b1k&moves=8.5}{\color{white}\bfseries lösen\,\ensuremath{\nearrow}}}
\end{tcolorbox}
\vspace{2mm}
\noindent\begin{minipage}[t]{36mm}\vspace{0pt}\sffamily\small\raggedright
\hmfigure{längste eindeutige}{{\Large h\#8.5}}\hmfigure{eindeutig bei dieser Länge}{{\Large 130} Stellungen}\hmfigure{längstes Matt}{{\Large h\#10.5}}\hmfigure{in der Klasse}{15,5 Milliarden Stellungen,\\ 4,61 Milliarden mit Hilfsmatt}
\end{minipage}\hfill
\begin{minipage}[t]{80mm}\vspace{0pt}\begin{tikzpicture}\begin{axis}[width=72mm,height=42mm,ybar,bar width=1.6pt,xmin=-1,xmax=22,ymin=0,ymax=1.242,axis y line*=left,axis x line*=bottom,xlabel={\scriptsize Halbzüge bis Matt},ylabel={\scriptsize Stellungen (Mrd.)},tick label style={font=\tiny},label style={font=\scriptsize},enlarge x limits=false,ylabel shift=-3pt,xlabel shift=-3pt]\addplot[fill=ink!55,draw=none] coordinates {(0,3.676e-06) (1,1.734e-05) (2,0.0001663) (3,0.001104) (4,0.005765) (5,0.02727) (6,0.08255) (7,0.2721) (8,0.4608) (9,0.9881) (10,0.7925) (11,1.129) (12,0.4501) (13,0.3351) (14,0.04843) (15,0.01311) (16,0.0008797) (17,0.0005925) (18,1.24e-07) (19,4.268e-05) (21,1.4e-08)};\draw[densely dotted] (axis cs:17,0) -- (axis cs:17,1.242);\end{axis}\begin{axis}[width=72mm,height=42mm,axis y line*=right,axis x line=none,ymode=log,xmin=-1,xmax=22,ylabel={\scriptsize eindeutig (log)},tick label style={font=\tiny},label style={font=\scriptsize},ylabel shift=-3pt]\addplot[mark=*,mark size=0.7pt,draw=black,thin] coordinates {(0,3676) (1,17165) (2,141766) (3,739908) (4,2253535) (5,5784244) (6,6033401) (7,11332471) (8,8555814) (9,11663002) (10,3489158) (11,2863683) (12,208382) (13,59877) (14,530) (15,998) (17,130)};\end{axis}\end{tikzpicture}\end{minipage}
\vspace{3mm}

\noindent\begin{minipage}[t]{60mm}\vspace{0pt}\sbox\hmboard{\chessboard[setfen=8/8/8/8/8/1b6/8/KrN2b1k w - - 0 1]}\usebox\hmboard\\[1mm]
\makebox[\wd\hmboard]{\sffamily\hspace{3mm}\textbf{h\#8.5}\hfill(2+4)}\end{minipage}\hfill
\begin{minipage}[t]{52mm}\vspace{0pt}
\textcolor{ink}{\sffamily\textsc{Lösung}}\par\smallskip\small
\begin{tabular}{@{}r@{\enspace}l@{\enspace}l@{}}
\textbf{1.} & \ldots & \symking{}xb1\\
\textbf{2.} & \symking{}g2 & \symknight{}d3\\
\textbf{3.} & \symking{}f3 & \symking{}c1\\
\textbf{4.} & \symking{}e4 & \symking{}d2\\
\textbf{5.} & \symking{}d4 & \symking{}e1\\
\textbf{6.} & \symking{}c3 & \symknight{}b4\\
\textbf{7.} & \symking{}b2 & \symking{}d2\\
\textbf{8.} & \symking{}a1 & \symking{}c1\\
\textbf{9.} & \symbishop{}a2 & \symknight{}c2\#\\
\end{tabular}\par\vspace{3mm}
\qrcode[height=12mm]{https://helpman.komtera.lt/?fen=8/8/8/8/8/1b6/8/KrN2b1k&moves=8.5}\\{\sffamily\tiny\color{ink}scannen zum Lösen}
\end{minipage}
\vspace{2mm}

{\raggedright\noindent\textcolor{ink}{\sffamily\textsc{Themen}}\enspace \themetag{reines Matt} \themetag{Mustermatt} \themetag{Switchback} \themetag{Selbstblock} \themetag{ohne Schach}\par\smallskip{\footnotesize\textit{Hinweis:} die Partei am Zug steht im Diagramm im Schach; die Lösung beginnt mit einem Schlag.}\par}

\clearpage
\phantomsection\label{hm:KNvkrbn}
\begin{tcolorbox}[enhanced,colback=ink,colframe=ink,arc=0pt,boxrule=0pt,left=4mm,right=4mm,top=2.5mm,bottom=2.5mm]
\color{white}\sffamily{\LARGE\bfseries KNvkrbn}\hfill{\Large h\#8.5}\\[0.5mm]
{\small 6 Steine\quad Weiß am Zug}\hfill{\small Nr.~356\qquad\href{https://helpman.komtera.lt/?fen=7k/8/8/8/n7/1b6/2N5/K3r3&moves=8.5}{\color{white}\bfseries lösen\,\ensuremath{\nearrow}}}
\end{tcolorbox}
\vspace{2mm}
\noindent\begin{minipage}[t]{36mm}\vspace{0pt}\sffamily\small\raggedright
\hmfigure{längste eindeutige}{{\Large h\#8.5}}\hmfigure{eindeutig bei dieser Länge}{{\Large 5} Stellungen}\hmfigure{längstes Matt}{{\Large h\#9.5}}\hmfigure{in der Klasse}{15,5 Milliarden Stellungen,\\ 9,44 Milliarden mit Hilfsmatt}
\end{minipage}\hfill
\begin{minipage}[t]{80mm}\vspace{0pt}\begin{tikzpicture}\begin{axis}[width=72mm,height=42mm,ybar,bar width=1.6pt,xmin=-1,xmax=20,ymin=0,ymax=2.459,axis y line*=left,axis x line*=bottom,xlabel={\scriptsize Halbzüge bis Matt},ylabel={\scriptsize Stellungen (Mrd.)},tick label style={font=\tiny},label style={font=\scriptsize},enlarge x limits=false,ylabel shift=-3pt,xlabel shift=-3pt]\addplot[fill=ink!55,draw=none] coordinates {(0,8.099e-06) (1,3.915e-05) (2,0.0003276) (3,0.002239) (4,0.01224) (5,0.05675) (6,0.1753) (7,0.5498) (8,0.9673) (9,1.945) (10,1.659) (11,2.235) (12,0.9735) (13,0.7197) (14,0.1137) (15,0.02852) (16,0.001138) (17,0.0002133) (19,3.542e-06)};\draw[densely dotted] (axis cs:17,0) -- (axis cs:17,2.459);\end{axis}\begin{axis}[width=72mm,height=42mm,axis y line*=right,axis x line=none,ymode=log,xmin=-1,xmax=20,ylabel={\scriptsize eindeutig (log)},tick label style={font=\tiny},label style={font=\scriptsize},ylabel shift=-3pt]\addplot[mark=*,mark size=0.7pt,draw=black,thin] coordinates {(0,8099) (1,38796) (2,275749) (3,1476578) (4,5006759) (5,11890684) (6,12672417) (7,20249908) (8,15504137) (9,18609056) (10,6047849) (11,4541580) (12,393775) (13,109344) (14,1790) (15,422) (17,5)};\end{axis}\end{tikzpicture}\end{minipage}
\vspace{3mm}

\noindent\begin{minipage}[t]{60mm}\vspace{0pt}\sbox\hmboard{\chessboard[setfen=7k/8/8/8/n7/1b6/2N5/K3r3 w - - 0 1]}\usebox\hmboard\\[1mm]
\makebox[\wd\hmboard]{\sffamily\hspace{3mm}\textbf{h\#8.5}\hfill(2+4)}\end{minipage}\hfill
\begin{minipage}[t]{52mm}\vspace{0pt}
\textcolor{ink}{\sffamily\textsc{Lösung}}\par\smallskip\small
\begin{tabular}{@{}r@{\enspace}l@{\enspace}l@{}}
\textbf{1.} & \ldots & \symknight{}xe1\\
\textbf{2.} & \symking{}g7 & \symking{}b1\\
\textbf{3.} & \symking{}f6 & \symking{}c1\\
\textbf{4.} & \symking{}e5 & \symking{}d2\\
\textbf{5.} & \symking{}d4 & \symking{}e2\\
\textbf{6.} & \symking{}c3 & \symking{}e3\\
\textbf{7.} & \symking{}b2 & \symking{}d2\\
\textbf{8.} & \symking{}a1 & \symking{}c1\\
\textbf{9.} & \symbishop{}a2 & \symknight{}c2\#\\
\end{tabular}\par\vspace{3mm}
\qrcode[height=12mm]{https://helpman.komtera.lt/?fen=7k/8/8/8/n7/1b6/2N5/K3r3&moves=8.5}\\{\sffamily\tiny\color{ink}scannen zum Lösen}
\end{minipage}
\vspace{2mm}

{\raggedright\noindent\textcolor{ink}{\sffamily\textsc{Themen}}\enspace \themetag{reines Matt} \themetag{Mustermatt} \themetag{Switchback} \themetag{Rundlauf} \themetag{Selbstblock} \themetag{ohne Schach}\par\smallskip{\footnotesize\textit{Hinweis:} die Partei am Zug steht im Diagramm im Schach; die Lösung beginnt mit einem Schlag.}\par}

\clearpage
\phantomsection\label{hm:KNvkrbp}
\begin{tcolorbox}[enhanced,colback=ink,colframe=ink,arc=0pt,boxrule=0pt,left=4mm,right=4mm,top=2.5mm,bottom=2.5mm]
\color{white}\sffamily{\LARGE\bfseries KNvkrbp}\hfill{\Large h\#8.5}\\[0.5mm]
{\small 6 Steine\quad Weiß am Zug}\hfill{\small Nr.~357\qquad\href{https://helpman.komtera.lt/?fen=8/7k/8/8/3p4/1b6/N7/Kr6&moves=8.5}{\color{white}\bfseries lösen\,\ensuremath{\nearrow}}}
\end{tcolorbox}
\vspace{2mm}
\noindent\begin{minipage}[t]{36mm}\vspace{0pt}\sffamily\small\raggedright
\hmfigure{längste eindeutige}{{\Large h\#8.5}}\hmfigure{eindeutig bei dieser Länge}{{\Large 98} Stellungen}\hmfigure{längstes Matt}{{\Large h\#10.5}}\hmfigure{in der Klasse}{45,4 Milliarden Stellungen,\\ 29 Milliarden mit Hilfsmatt}
\end{minipage}\hfill
\begin{minipage}[t]{80mm}\vspace{0pt}\begin{tikzpicture}\begin{axis}[width=72mm,height=42mm,ybar,bar width=1.6pt,xmin=-1,xmax=22,ymin=0,ymax=7.576,axis y line*=left,axis x line*=bottom,xlabel={\scriptsize Halbzüge bis Matt},ylabel={\scriptsize Stellungen (Mrd.)},tick label style={font=\tiny},label style={font=\scriptsize},enlarge x limits=false,ylabel shift=-3pt,xlabel shift=-3pt]\addplot[fill=ink!55,draw=none] coordinates {(0,2.96e-05) (1,0.0001383) (2,0.001052) (3,0.006966) (4,0.03615) (5,0.1642) (6,0.4857) (7,1.494) (8,2.753) (9,5.5) (10,5.192) (11,6.887) (12,3.388) (13,2.444) (14,0.529) (15,0.1351) (16,0.01141) (17,0.001966) (18,3.976e-06) (19,5.415e-05) (21,3e-09)};\draw[densely dotted] (axis cs:17,0) -- (axis cs:17,7.576);\end{axis}\begin{axis}[width=72mm,height=42mm,axis y line*=right,axis x line=none,ymode=log,xmin=-1,xmax=22,ylabel={\scriptsize eindeutig (log)},tick label style={font=\tiny},label style={font=\scriptsize},ylabel shift=-3pt]\addplot[mark=*,mark size=0.7pt,draw=black,thin] coordinates {(0,29602) (1,135489) (2,873350) (3,4548525) (4,15075781) (5,33779090) (6,35294114) (7,54629169) (8,48898570) (9,56142497) (10,20868946) (11,14424463) (12,1560169) (13,405099) (14,13042) (15,3066) (16,4) (17,98)};\end{axis}\end{tikzpicture}\end{minipage}
\vspace{3mm}

\noindent\begin{minipage}[t]{60mm}\vspace{0pt}\sbox\hmboard{\chessboard[setfen=8/7k/8/8/3p4/1b6/N7/Kr6 w - - 0 1]}\usebox\hmboard\\[1mm]
\makebox[\wd\hmboard]{\sffamily\hspace{3mm}\textbf{h\#8.5}\hfill(2+4)}\end{minipage}\hfill
\begin{minipage}[t]{52mm}\vspace{0pt}
\textcolor{ink}{\sffamily\textsc{Lösung}}\par\smallskip\small
\begin{tabular}{@{}r@{\enspace}l@{\enspace}l@{}}
\textbf{1.} & \ldots & \symking{}xb1\\
\textbf{2.} & \symking{}g6 & \symking{}c1\\
\textbf{3.} & \symking{}f5 & \symking{}d2\\
\textbf{4.} & \symking{}e4 & \symking{}e1\\
\textbf{5.} & \symking{}d3 & \symknight{}b4+\\
\textbf{6.} & \symking{}c3 & \symking{}e2\\
\textbf{7.} & \symking{}b2 & \symking{}d2\\
\textbf{8.} & \symking{}a1 & \symking{}c1\\
\textbf{9.} & \symbishop{}a2 & \symknight{}c2\#\\
\end{tabular}\par\vspace{3mm}
\qrcode[height=12mm]{https://helpman.komtera.lt/?fen=8/7k/8/8/3p4/1b6/N7/Kr6&moves=8.5}\\{\sffamily\tiny\color{ink}scannen zum Lösen}
\end{minipage}
\vspace{2mm}

{\raggedright\noindent\textcolor{ink}{\sffamily\textsc{Themen}}\enspace \themetag{reines Matt} \themetag{Mustermatt} \themetag{Rundlauf} \themetag{Selbstblock}\par\smallskip{\footnotesize\textit{Hinweis:} die Partei am Zug steht im Diagramm im Schach; die Lösung beginnt mit einem Schlag.}\par}

\clearpage
\phantomsection\label{hm:KNvkrnn}
\begin{tcolorbox}[enhanced,colback=ink,colframe=ink,arc=0pt,boxrule=0pt,left=4mm,right=4mm,top=2.5mm,bottom=2.5mm]
\color{white}\sffamily{\LARGE\bfseries KNvkrnn}\hfill{\Large h\#8.5}\\[0.5mm]
{\small 6 Steine\quad Weiß am Zug}\hfill{\small Nr.~358\qquad\href{https://helpman.komtera.lt/?fen=6k1/8/8/8/8/n7/4r3/nKN5&moves=8.5}{\color{white}\bfseries lösen\,\ensuremath{\nearrow}}}
\end{tcolorbox}
\vspace{2mm}
\noindent\begin{minipage}[t]{36mm}\vspace{0pt}\sffamily\small\raggedright
\hmfigure{längste eindeutige}{{\Large h\#8.5}}\hmfigure{eindeutig bei dieser Länge}{{\Large 1} Stellung}\hmfigure{längstes Matt}{{\Large h\#9.5}}\hmfigure{in der Klasse}{15,5 Milliarden Stellungen,\\ 4,82 Milliarden mit Hilfsmatt}
\end{minipage}\hfill
\begin{minipage}[t]{80mm}\vspace{0pt}\begin{tikzpicture}\begin{axis}[width=72mm,height=42mm,ybar,bar width=1.6pt,xmin=-1,xmax=20,ymin=0,ymax=1.203,axis y line*=left,axis x line*=bottom,xlabel={\scriptsize Halbzüge bis Matt},ylabel={\scriptsize Stellungen (Mrd.)},tick label style={font=\tiny},label style={font=\scriptsize},enlarge x limits=false,ylabel shift=-3pt,xlabel shift=-3pt]\addplot[fill=ink!55,draw=none] coordinates {(0,4.128e-06) (1,2.056e-05) (2,0.0001482) (3,0.001063) (4,0.005741) (5,0.02677) (6,0.08057) (7,0.2468) (8,0.4404) (9,0.8583) (10,0.8185) (11,1.093) (12,0.5696) (13,0.4947) (14,0.126) (15,0.04819) (16,0.005468) (17,0.0004345) (18,2.6e-06) (19,1.02e-06)};\draw[densely dotted] (axis cs:17,0) -- (axis cs:17,1.203);\end{axis}\begin{axis}[width=72mm,height=42mm,axis y line*=right,axis x line=none,ymode=log,xmin=-1,xmax=20,ylabel={\scriptsize eindeutig (log)},tick label style={font=\tiny},label style={font=\scriptsize},ylabel shift=-3pt]\addplot[mark=*,mark size=0.7pt,draw=black,thin] coordinates {(0,4128) (1,20464) (2,121832) (3,686968) (4,2320281) (5,5583280) (6,5739385) (7,8281832) (8,5589316) (9,6187987) (10,2514843) (11,1884737) (12,255344) (13,131638) (14,10312) (15,535) (16,5) (17,1)};\end{axis}\end{tikzpicture}\end{minipage}
\vspace{3mm}

\noindent\begin{minipage}[t]{60mm}\vspace{0pt}\sbox\hmboard{\chessboard[setfen=6k1/8/8/8/8/n7/4r3/nKN5 w - - 0 1]}\usebox\hmboard\\[1mm]
\makebox[\wd\hmboard]{\sffamily\hspace{3mm}\textbf{h\#8.5}\hfill(2+4)}\end{minipage}\hfill
\begin{minipage}[t]{52mm}\vspace{0pt}
\textcolor{ink}{\sffamily\textsc{Lösung}}\par\smallskip\small
\begin{tabular}{@{}r@{\enspace}l@{\enspace}l@{}}
\textbf{1.} & \ldots & \symking{}xa1\\
\textbf{2.} & \symking{}f7 & \symknight{}xe2\\
\textbf{3.} & \symking{}e6 & \symking{}b2\\
\textbf{4.} & \symking{}d5 & \symking{}c1\\
\textbf{5.} & \symking{}c4 & \symking{}d2\\
\textbf{6.} & \symking{}b3 & \symknight{}d4+\\
\textbf{7.} & \symking{}a2 & \symking{}c3\\
\textbf{8.} & \symking{}a1 & \symking{}b3\\
\textbf{9.} & \symknight{}b1 & \symknight{}c2\#\\
\end{tabular}\par\vspace{3mm}
\qrcode[height=12mm]{https://helpman.komtera.lt/?fen=6k1/8/8/8/8/n7/4r3/nKN5&moves=8.5}\\{\sffamily\tiny\color{ink}scannen zum Lösen}
\end{minipage}
\vspace{2mm}

{\raggedright\noindent\textcolor{ink}{\sffamily\textsc{Themen}}\enspace \themetag{reines Matt} \themetag{Mustermatt} \themetag{Idealmatt} \themetag{Selbstblock} \themetag{Kniest}\par\smallskip{\footnotesize\textit{Hinweis:} die Partei am Zug steht im Diagramm im Schach; die Lösung beginnt mit einem Schlag.}\par}

\clearpage
\phantomsection\label{hm:KNvkrnp}
\begin{tcolorbox}[enhanced,colback=ink,colframe=ink,arc=0pt,boxrule=0pt,left=4mm,right=4mm,top=2.5mm,bottom=2.5mm]
\color{white}\sffamily{\LARGE\bfseries KNvkrnp}\hfill{\Large h\#8.5}\\[0.5mm]
{\small 6 Steine\quad Weiß am Zug}\hfill{\small Nr.~359\qquad\href{https://helpman.komtera.lt/?fen=6k1/8/8/2p5/8/n7/4r3/1KN5&moves=8.5}{\color{white}\bfseries lösen\,\ensuremath{\nearrow}}}
\end{tcolorbox}
\vspace{2mm}
\noindent\begin{minipage}[t]{36mm}\vspace{0pt}\sffamily\small\raggedright
\hmfigure{längste eindeutige}{{\Large h\#8.5}}\hmfigure{eindeutig bei dieser Länge}{{\Large 54} Stellungen}\hmfigure{längstes Matt}{{\Large h\#10.5}}\hmfigure{in der Klasse}{45,4 Milliarden Stellungen,\\ 29,7 Milliarden mit Hilfsmatt}
\end{minipage}\hfill
\begin{minipage}[t]{80mm}\vspace{0pt}\begin{tikzpicture}\begin{axis}[width=72mm,height=42mm,ybar,bar width=1.6pt,xmin=-1,xmax=22,ymin=0,ymax=7.376,axis y line*=left,axis x line*=bottom,xlabel={\scriptsize Halbzüge bis Matt},ylabel={\scriptsize Stellungen (Mrd.)},tick label style={font=\tiny},label style={font=\scriptsize},enlarge x limits=false,ylabel shift=-3pt,xlabel shift=-3pt]\addplot[fill=ink!55,draw=none] coordinates {(0,3.042e-05) (1,0.0001447) (2,0.0009193) (3,0.006322) (4,0.03291) (5,0.1491) (6,0.4329) (7,1.285) (8,2.403) (9,4.601) (10,5.031) (11,6.705) (12,4.036) (13,3.471) (14,1.035) (15,0.4132) (16,0.0569) (17,0.005503) (18,0.0001113) (19,4.603e-05) (21,5e-08)};\draw[densely dotted] (axis cs:17,0) -- (axis cs:17,7.376);\end{axis}\begin{axis}[width=72mm,height=42mm,axis y line*=right,axis x line=none,ymode=log,xmin=-1,xmax=22,ylabel={\scriptsize eindeutig (log)},tick label style={font=\tiny},label style={font=\scriptsize},ylabel shift=-3pt]\addplot[mark=*,mark size=0.7pt,draw=black,thin] coordinates {(0,30424) (1,141985) (2,742155) (3,4028333) (4,13892703) (5,30798466) (6,31179212) (7,42592116) (8,33032766) (9,36118451) (10,19764039) (11,13933031) (12,2733806) (13,1298398) (14,127967) (15,11138) (16,205) (17,54)};\end{axis}\end{tikzpicture}\end{minipage}
\vspace{3mm}

\noindent\begin{minipage}[t]{60mm}\vspace{0pt}\sbox\hmboard{\chessboard[setfen=6k1/8/8/2p5/8/n7/4r3/1KN5 w - - 0 1]}\usebox\hmboard\\[1mm]
\makebox[\wd\hmboard]{\sffamily\hspace{3mm}\textbf{h\#8.5}\hfill(2+4)}\end{minipage}\hfill
\begin{minipage}[t]{52mm}\vspace{0pt}
\textcolor{ink}{\sffamily\textsc{Lösung}}\par\smallskip\small
\begin{tabular}{@{}r@{\enspace}l@{\enspace}l@{}}
\textbf{1.} & \ldots & \symking{}a1\\
\textbf{2.} & \symking{}f7 & \symknight{}xe2\\
\textbf{3.} & \symking{}e6 & \symking{}b2\\
\textbf{4.} & \symking{}d5 & \symking{}c1\\
\textbf{5.} & \symking{}c4 & \symking{}d2\\
\textbf{6.} & \symking{}b3 & \symknight{}d4+\\
\textbf{7.} & \symking{}a2 & \symking{}c3\\
\textbf{8.} & \symking{}a1 & \symking{}b3\\
\textbf{9.} & \symknight{}b1 & \symknight{}c2\#\\
\end{tabular}\par\vspace{3mm}
\qrcode[height=12mm]{https://helpman.komtera.lt/?fen=6k1/8/8/2p5/8/n7/4r3/1KN5&moves=8.5}\\{\sffamily\tiny\color{ink}scannen zum Lösen}
\end{minipage}
\vspace{2mm}

{\raggedright\noindent\textcolor{ink}{\sffamily\textsc{Themen}}\enspace \themetag{reines Matt} \themetag{Mustermatt} \themetag{Selbstblock}\par\smallskip{\footnotesize\textit{Hinweis:} die Partei am Zug steht im Diagramm im Schach.}\par}

\clearpage
\phantomsection\label{hm:KNvkrrn}
\begin{tcolorbox}[enhanced,colback=ink,colframe=ink,arc=0pt,boxrule=0pt,left=4mm,right=4mm,top=2.5mm,bottom=2.5mm]
\color{white}\sffamily{\LARGE\bfseries KNvkrrn}\hfill{\Large h\#8.5}\\[0.5mm]
{\small 6 Steine\quad Weiß am Zug}\hfill{\small Nr.~360\qquad\href{https://helpman.komtera.lt/?fen=8/7k/8/1r6/8/r7/8/KNn5&moves=8.5}{\color{white}\bfseries lösen\,\ensuremath{\nearrow}}}
\end{tcolorbox}
\vspace{2mm}
\noindent\begin{minipage}[t]{36mm}\vspace{0pt}\sffamily\small\raggedright
\hmfigure{längste eindeutige}{{\Large h\#8.5}}\hmfigure{eindeutig bei dieser Länge}{{\Large 7} Stellungen}\hmfigure{längstes Matt}{{\Large h\#9.5}}\hmfigure{in der Klasse}{15,5 Milliarden Stellungen,\\ 4,51 Milliarden mit Hilfsmatt}
\end{minipage}\hfill
\begin{minipage}[t]{80mm}\vspace{0pt}\begin{tikzpicture}\begin{axis}[width=72mm,height=42mm,ybar,bar width=1.6pt,xmin=-1,xmax=20,ymin=0,ymax=1.181,axis y line*=left,axis x line*=bottom,xlabel={\scriptsize Halbzüge bis Matt},ylabel={\scriptsize Stellungen (Mrd.)},tick label style={font=\tiny},label style={font=\scriptsize},enlarge x limits=false,ylabel shift=-3pt,xlabel shift=-3pt]\addplot[fill=ink!55,draw=none] coordinates {(0,3.531e-06) (1,1.608e-05) (2,0.000159) (3,0.001102) (4,0.005698) (5,0.02769) (6,0.08569) (7,0.286) (8,0.4692) (9,0.9809) (10,0.7404) (11,1.074) (12,0.4061) (13,0.3597) (14,0.0499) (15,0.01989) (16,0.0005001) (17,0.0002498) (18,4e-09) (19,6.2e-07)};\draw[densely dotted] (axis cs:17,0) -- (axis cs:17,1.181);\end{axis}\begin{axis}[width=72mm,height=42mm,axis y line*=right,axis x line=none,ymode=log,xmin=-1,xmax=20,ylabel={\scriptsize eindeutig (log)},tick label style={font=\tiny},label style={font=\scriptsize},ylabel shift=-3pt]\addplot[mark=*,mark size=0.7pt,draw=black,thin] coordinates {(0,3531) (1,16019) (2,132641) (3,737001) (4,1875735) (5,5834731) (6,6975044) (7,12972351) (8,9080366) (9,11463982) (10,3568036) (11,3420823) (12,402634) (13,205676) (14,4314) (15,1296) (17,7)};\end{axis}\end{tikzpicture}\end{minipage}
\vspace{3mm}

\noindent\begin{minipage}[t]{60mm}\vspace{0pt}\sbox\hmboard{\chessboard[setfen=8/7k/8/1r6/8/r7/8/KNn5 w - - 0 1]}\usebox\hmboard\\[1mm]
\makebox[\wd\hmboard]{\sffamily\hspace{3mm}\textbf{h\#8.5}\hfill(2+4)}\end{minipage}\hfill
\begin{minipage}[t]{52mm}\vspace{0pt}
\textcolor{ink}{\sffamily\textsc{Lösung}}\par\smallskip\small
\begin{tabular}{@{}r@{\enspace}l@{\enspace}l@{}}
\textbf{1.} & \ldots & \symknight{}xa3\\
\textbf{2.} & \symking{}g6 & \symknight{}xb5\\
\textbf{3.} & \symking{}f5 & \symking{}b2\\
\textbf{4.} & \symking{}e4 & \symking{}a3\\
\textbf{5.} & \symking{}d3 & \symking{}b4\\
\textbf{6.} & \symking{}c2 & \symknight{}d4+\\
\textbf{7.} & \symking{}b1 & \symking{}c3\\
\textbf{8.} & \symking{}a1 & \symking{}c2\\
\textbf{9.} & \symknight{}a2 & \symknight{}b3\#\\
\end{tabular}\par\vspace{3mm}
\qrcode[height=12mm]{https://helpman.komtera.lt/?fen=8/7k/8/1r6/8/r7/8/KNn5&moves=8.5}\\{\sffamily\tiny\color{ink}scannen zum Lösen}
\end{minipage}
\vspace{2mm}

{\raggedright\noindent\textcolor{ink}{\sffamily\textsc{Themen}}\enspace \themetag{reines Matt} \themetag{Mustermatt} \themetag{Idealmatt} \themetag{Selbstblock}\par\smallskip{\footnotesize\textit{Hinweis:} die Partei am Zug steht im Diagramm im Schach; die Lösung beginnt mit einem Schlag.}\par}

\clearpage
\phantomsection\label{hm:KNvkrrp}
\begin{tcolorbox}[enhanced,colback=ink,colframe=ink,arc=0pt,boxrule=0pt,left=4mm,right=4mm,top=2.5mm,bottom=2.5mm]
\color{white}\sffamily{\LARGE\bfseries KNvkrrp}\hfill{\Large h\#8.5}\\[0.5mm]
{\small 6 Steine\quad Weiß am Zug}\hfill{\small Nr.~361\qquad\href{https://helpman.komtera.lt/?fen=8/7k/N7/2p5/8/8/4r3/Kr6&moves=8.5}{\color{white}\bfseries lösen\,\ensuremath{\nearrow}}}
\end{tcolorbox}
\vspace{2mm}
\noindent\begin{minipage}[t]{36mm}\vspace{0pt}\sffamily\small\raggedright
\hmfigure{längste eindeutige}{{\Large h\#8.5}}\hmfigure{eindeutig bei dieser Länge}{{\Large 167} Stellungen}\hmfigure{längstes Matt}{{\Large h\#9.5}}\hmfigure{in der Klasse}{45,4 Milliarden Stellungen,\\ 13,9 Milliarden mit Hilfsmatt}
\end{minipage}\hfill
\begin{minipage}[t]{80mm}\vspace{0pt}\begin{tikzpicture}\begin{axis}[width=72mm,height=42mm,ybar,bar width=1.6pt,xmin=-1,xmax=20,ymin=0,ymax=3.586,axis y line*=left,axis x line*=bottom,xlabel={\scriptsize Halbzüge bis Matt},ylabel={\scriptsize Stellungen (Mrd.)},tick label style={font=\tiny},label style={font=\scriptsize},enlarge x limits=false,ylabel shift=-3pt,xlabel shift=-3pt]\addplot[fill=ink!55,draw=none] coordinates {(0,1.328e-05) (1,5.871e-05) (2,0.0005119) (3,0.003452) (4,0.01708) (5,0.08147) (6,0.2463) (7,0.8072) (8,1.407) (9,2.851) (10,2.337) (11,3.26) (12,1.367) (13,1.169) (14,0.2145) (15,0.08558) (16,0.004432) (17,0.00284) (18,1.61e-07) (19,4.494e-05)};\draw[densely dotted] (axis cs:17,0) -- (axis cs:17,3.586);\end{axis}\begin{axis}[width=72mm,height=42mm,axis y line*=right,axis x line=none,ymode=log,xmin=-1,xmax=20,ylabel={\scriptsize eindeutig (log)},tick label style={font=\tiny},label style={font=\scriptsize},ylabel shift=-3pt]\addplot[mark=*,mark size=0.7pt,draw=black,thin] coordinates {(0,13278) (1,57711) (2,423202) (3,2293212) (4,5699057) (5,17001716) (6,20554997) (7,37403718) (8,29591336) (9,34774417) (10,12035831) (11,11115381) (12,1494606) (13,875823) (14,28739) (15,16894) (16,9) (17,167)};\end{axis}\end{tikzpicture}\end{minipage}
\vspace{3mm}

\noindent\begin{minipage}[t]{60mm}\vspace{0pt}\sbox\hmboard{\chessboard[setfen=8/7k/N7/2p5/8/8/4r3/Kr6 w - - 0 1]}\usebox\hmboard\\[1mm]
\makebox[\wd\hmboard]{\sffamily\hspace{3mm}\textbf{h\#8.5}\hfill(2+4)}\end{minipage}\hfill
\begin{minipage}[t]{52mm}\vspace{0pt}
\textcolor{ink}{\sffamily\textsc{Lösung}}\par\smallskip\small
\begin{tabular}{@{}r@{\enspace}l@{\enspace}l@{}}
\textbf{1.} & \ldots & \symking{}xb1\\
\textbf{2.} & \symking{}g6 & \symknight{}b4\\
\textbf{3.} & \symking{}f5 & \symknight{}c2\\
\textbf{4.} & \symking{}e4 & \symking{}b2\\
\textbf{5.} & \symking{}d3 & \symking{}c1\\
\textbf{6.} & \symking{}c3 & \symking{}d1\\
\textbf{7.} & \symking{}b2 & \symknight{}d4\\
\textbf{8.} & \symking{}a1 & \symking{}c1\\
\textbf{9.} & \symrook{}a2 & \symknight{}b3\#\\
\end{tabular}\par\vspace{3mm}
\qrcode[height=12mm]{https://helpman.komtera.lt/?fen=8/7k/N7/2p5/8/8/4r3/Kr6&moves=8.5}\\{\sffamily\tiny\color{ink}scannen zum Lösen}
\end{minipage}
\vspace{2mm}

{\raggedright\noindent\textcolor{ink}{\sffamily\textsc{Themen}}\enspace \themetag{reines Matt} \themetag{Mustermatt} \themetag{Switchback} \themetag{Selbstblock} \themetag{ohne Schach}\par\smallskip{\footnotesize\textit{Hinweis:} die Partei am Zug steht im Diagramm im Schach; die Lösung beginnt mit einem Schlag.}\par}

\clearpage
\phantomsection\label{hm:KPPPvkp}
\begin{tcolorbox}[enhanced,colback=ink,colframe=ink,arc=0pt,boxrule=0pt,left=4mm,right=4mm,top=2.5mm,bottom=2.5mm]
\color{white}\sffamily{\LARGE\bfseries KPPPvkp}\hfill{\Large h\#8.5}\\[0.5mm]
{\small 6 Steine\quad Weiß am Zug}\hfill{\small Nr.~362\qquad\href{https://helpman.komtera.lt/?fen=8/8/p7/8/P7/P7/P2K4/7k&moves=8.5}{\color{white}\bfseries lösen\,\ensuremath{\nearrow}}}
\end{tcolorbox}
\vspace{2mm}
\noindent\begin{minipage}[t]{36mm}\vspace{0pt}\sffamily\small\raggedright
\hmfigure{längste eindeutige}{{\Large h\#8.5}}\hmfigure{eindeutig bei dieser Länge}{{\Large 2} Stellungen}\hmfigure{längstes Matt}{{\Large h\#13.5}}\hmfigure{in der Klasse}{19,2 Milliarden Stellungen,\\ 2,34 Milliarden mit Hilfsmatt}
\end{minipage}\hfill
\begin{minipage}[t]{80mm}\vspace{0pt}\begin{tikzpicture}\begin{axis}[width=72mm,height=42mm,ybar,bar width=1.6pt,xmin=-1,xmax=28,ymin=0,ymax=474.3,axis y line*=left,axis x line*=bottom,xlabel={\scriptsize Halbzüge bis Matt},ylabel={\scriptsize Stellungen (Mio.)},tick label style={font=\tiny},label style={font=\scriptsize},enlarge x limits=false,ylabel shift=-3pt,xlabel shift=-3pt]\addplot[fill=ink!55,draw=none] coordinates {(0,0.02749) (1,2.113) (2,5.063) (3,34.99) (4,88.51) (5,282.4) (6,411.8) (7,431.1) (8,403.7) (9,266.9) (10,205.1) (11,102.1) (12,78.53) (13,14.95) (14,8.677) (15,0.2862) (16,0.09479) (17,0.0395) (18,0.04224) (19,0.03261) (20,0.03493) (21,0.02156) (22,0.02152) (23,0.005457) (24,0.005375) (25,0.001337) (26,0.001367) (27,1.5e-05)};\draw[densely dotted] (axis cs:17,0) -- (axis cs:17,474.3);\end{axis}\begin{axis}[width=72mm,height=42mm,axis y line*=right,axis x line=none,ymode=log,xmin=-1,xmax=28,ylabel={\scriptsize eindeutig (log)},tick label style={font=\tiny},label style={font=\scriptsize},ylabel shift=-3pt]\addplot[mark=*,mark size=0.7pt,draw=black,thin] coordinates {(0,27494) (1,530090) (2,1440042) (3,13587769) (4,19124609) (5,21109255) (6,13197405) (7,7148834) (8,3703557) (9,1758756) (10,874236) (11,344796) (12,147484) (13,2291) (14,228) (15,17) (16,5) (17,2)};\end{axis}\end{tikzpicture}\end{minipage}
\vspace{3mm}

\noindent\begin{minipage}[t]{60mm}\vspace{0pt}\sbox\hmboard{\chessboard[setfen=8/8/p7/8/P7/P7/P2K4/7k w - - 0 1]}\usebox\hmboard\\[1mm]
\makebox[\wd\hmboard]{\sffamily\hspace{3mm}\textbf{h\#8.5}\hfill(4+2)}\end{minipage}\hfill
\begin{minipage}[t]{52mm}\vspace{0pt}
\textcolor{ink}{\sffamily\textsc{Lösung}}\par\smallskip\small
\begin{tabular}{@{}r@{\enspace}l@{\enspace}l@{}}
\textbf{1.} & \ldots & \symking{}c3\\
\textbf{2.} & \symking{}g2 & \symking{}b4\\
\textbf{3.} & a5+ & \symking{}xa5\\
\textbf{4.} & \symking{}f3 & \symking{}b4\\
\textbf{5.} & \symking{}e4 & a5\\
\textbf{6.} & \symking{}d5 & a6\\
\textbf{7.} & \symking{}c6 & a7\\
\textbf{8.} & \symking{}b7 & \symking{}c5\\
\textbf{9.} & \symking{}a6 & a8=\symqueen{}\#\\
\end{tabular}\par\vspace{3mm}
\qrcode[height=12mm]{https://helpman.komtera.lt/?fen=8/8/p7/8/P7/P7/P2K4/7k&moves=8.5}\\{\sffamily\tiny\color{ink}scannen zum Lösen}
\end{minipage}
\vspace{2mm}

{\raggedright\noindent\textcolor{ink}{\sffamily\textsc{Themen}}\enspace \themetag{reines Matt} \themetag{Mustermatt} \themetag{Spiegelmatt} \themetag{Umwandlung} \themetag{Switchback}\par}

\clearpage
\phantomsection\label{hm:KPPvkbp}
\begin{tcolorbox}[enhanced,colback=ink,colframe=ink,arc=0pt,boxrule=0pt,left=4mm,right=4mm,top=2.5mm,bottom=2.5mm]
\color{white}\sffamily{\LARGE\bfseries KPPvkbp}\hfill{\Large h\#8.5}\\[0.5mm]
{\small 6 Steine\quad Weiß am Zug}\hfill{\small Nr.~363\qquad\href{https://helpman.komtera.lt/?fen=8/8/8/8/1p6/1P6/1Pk5/K1b5&moves=8.5}{\color{white}\bfseries lösen\,\ensuremath{\nearrow}}}
\end{tcolorbox}
\vspace{2mm}
\noindent\begin{minipage}[t]{36mm}\vspace{0pt}\sffamily\small\raggedright
\hmfigure{längste eindeutige}{{\Large h\#8.5}}\hmfigure{eindeutig bei dieser Länge}{{\Large 125} Stellungen}\hmfigure{längstes Matt}{{\Large h\#14.5}}\hmfigure{in der Klasse}{25,6 Milliarden Stellungen,\\ 9,02 Milliarden mit Hilfsmatt}
\end{minipage}\hfill
\begin{minipage}[t]{80mm}\vspace{0pt}\begin{tikzpicture}\begin{axis}[width=72mm,height=42mm,ybar,bar width=1.6pt,xmin=-1,xmax=30,ymin=0,ymax=1.698,axis y line*=left,axis x line*=bottom,xlabel={\scriptsize Halbzüge bis Matt},ylabel={\scriptsize Stellungen (Mrd.)},tick label style={font=\tiny},label style={font=\scriptsize},enlarge x limits=false,ylabel shift=-3pt,xlabel shift=-3pt]\addplot[fill=ink!55,draw=none] coordinates {(0,3.939e-05) (1,0.003033) (2,0.008866) (3,0.07093) (4,0.1861) (5,0.7052) (6,1.057) (7,1.543) (8,1.392) (9,1.379) (10,1.001) (11,0.8491) (12,0.568) (13,0.1916) (14,0.05994) (15,0.004417) (16,0.0007468) (17,0.0002444) (18,4.938e-05) (19,1.446e-05) (20,4.24e-06) (21,2.972e-06) (22,1.528e-06) (23,1.196e-06) (24,4.55e-07) (25,9.8e-08) (26,6e-09) (27,1.8e-08) (29,6e-09)};\draw[densely dotted] (axis cs:17,0) -- (axis cs:17,1.698);\end{axis}\begin{axis}[width=72mm,height=42mm,axis y line*=right,axis x line=none,ymode=log,xmin=-1,xmax=30,ylabel={\scriptsize eindeutig (log)},tick label style={font=\tiny},label style={font=\scriptsize},ylabel shift=-3pt]\addplot[mark=*,mark size=0.7pt,draw=black,thin] coordinates {(0,39392) (1,845593) (2,2221676) (3,22554892) (4,39453173) (5,59994980) (6,36999907) (7,29274301) (8,15545520) (9,11945567) (10,4998675) (11,3119792) (12,959050) (13,64741) (14,6800) (15,1213) (16,376) (17,125)};\end{axis}\end{tikzpicture}\end{minipage}
\vspace{3mm}

\noindent\begin{minipage}[t]{60mm}\vspace{0pt}\sbox\hmboard{\chessboard[setfen=8/8/8/8/1p6/1P6/1Pk5/K1b5 w - - 0 1]}\usebox\hmboard\\[1mm]
\makebox[\wd\hmboard]{\sffamily\hspace{3mm}\textbf{h\#8.5}\hfill(3+3)}\end{minipage}\hfill
\begin{minipage}[t]{52mm}\vspace{0pt}
\textcolor{ink}{\sffamily\textsc{Lösung}}\par\smallskip\small
\begin{tabular}{@{}r@{\enspace}l@{\enspace}l@{}}
\textbf{1.} & \ldots & \symking{}a2\\
\textbf{2.} & \symbishop{}d2 & \symking{}a1\\
\textbf{3.} & \symbishop{}c3 & bxc3\\
\textbf{4.} & bxc3 & b4\\
\textbf{5.} & \symking{}d1 & b5\\
\textbf{6.} & c2 & b6\\
\textbf{7.} & c1=\symbishop{} & b7\\
\textbf{8.} & \symbishop{}d2 & b8=\symqueen{}\\
\textbf{9.} & \symking{}c1 & \symqueen{}b1\#\\
\end{tabular}\par\vspace{3mm}
\qrcode[height=12mm]{https://helpman.komtera.lt/?fen=8/8/8/8/1p6/1P6/1Pk5/K1b5&moves=8.5}\\{\sffamily\tiny\color{ink}scannen zum Lösen}
\end{minipage}
\vspace{2mm}

{\raggedright\noindent\textcolor{ink}{\sffamily\textsc{Themen}}\enspace \themetag{Satzspiel} \themetag{Umwandlung} \themetag{Unterverwandlung} \themetag{Switchback} \themetag{Selbstblock} \themetag{Phönix} \themetag{ohne Schach} \themetag{Umnov}\par}

\clearpage
\phantomsection\label{hm:KPPvknp}
\begin{tcolorbox}[enhanced,colback=ink,colframe=ink,arc=0pt,boxrule=0pt,left=4mm,right=4mm,top=2.5mm,bottom=2.5mm]
\color{white}\sffamily{\LARGE\bfseries KPPvknp}\hfill{\Large h\#8.5}\\[0.5mm]
{\small 6 Steine\quad Weiß am Zug}\hfill{\small Nr.~364\qquad\href{https://helpman.komtera.lt/?fen=3n4/8/6k1/8/8/1p6/1P3P2/K7&moves=8.5}{\color{white}\bfseries lösen\,\ensuremath{\nearrow}}}
\end{tcolorbox}
\vspace{2mm}
\noindent\begin{minipage}[t]{36mm}\vspace{0pt}\sffamily\small\raggedright
\hmfigure{längste eindeutige}{{\Large h\#8.5}}\hmfigure{eindeutig bei dieser Länge}{{\Large 9} Stellungen}\hmfigure{längstes Matt}{{\Large h\#14.5}}\hmfigure{in der Klasse}{25,6 Milliarden Stellungen,\\ 9,24 Milliarden mit Hilfsmatt}
\end{minipage}\hfill
\begin{minipage}[t]{80mm}\vspace{0pt}\begin{tikzpicture}\begin{axis}[width=72mm,height=42mm,ybar,bar width=1.6pt,xmin=-1,xmax=30,ymin=0,ymax=1.6,axis y line*=left,axis x line*=bottom,xlabel={\scriptsize Halbzüge bis Matt},ylabel={\scriptsize Stellungen (Mrd.)},tick label style={font=\tiny},label style={font=\scriptsize},enlarge x limits=false,ylabel shift=-3pt,xlabel shift=-3pt]\addplot[fill=ink!55,draw=none] coordinates {(0,3.601e-05) (1,0.003798) (2,0.009586) (3,0.07035) (4,0.1683) (5,0.6414) (6,0.9712) (7,1.454) (8,1.441) (9,1.402) (10,1.113) (11,0.9065) (12,0.6718) (13,0.2588) (14,0.1136) (15,0.009479) (16,0.001565) (17,0.0002929) (18,4.874e-05) (19,1.642e-05) (20,5.49e-07) (21,1.48e-07) (22,4e-09) (23,1.01e-07) (25,2.6e-08) (27,1.2e-08) (29,4e-09)};\draw[densely dotted] (axis cs:17,0) -- (axis cs:17,1.6);\end{axis}\begin{axis}[width=72mm,height=42mm,axis y line*=right,axis x line=none,ymode=log,xmin=-1,xmax=30,ylabel={\scriptsize eindeutig (log)},tick label style={font=\tiny},label style={font=\scriptsize},ylabel shift=-3pt]\addplot[mark=*,mark size=0.7pt,draw=black,thin] coordinates {(0,36009) (1,890570) (2,2009533) (3,19633821) (4,31137903) (5,50637976) (6,34299460) (7,26125676) (8,15542707) (9,11568763) (10,5999860) (11,4239333) (12,1654424) (13,107941) (14,8606) (15,582) (16,88) (17,9)};\end{axis}\end{tikzpicture}\end{minipage}
\vspace{3mm}

\noindent\begin{minipage}[t]{60mm}\vspace{0pt}\sbox\hmboard{\chessboard[setfen=3n4/8/6k1/8/8/1p6/1P3P2/K7 w - - 0 1]}\usebox\hmboard\\[1mm]
\makebox[\wd\hmboard]{\sffamily\hspace{3mm}\textbf{h\#8.5}\hfill(3+3)}\end{minipage}\hfill
\begin{minipage}[t]{52mm}\vspace{0pt}
\textcolor{ink}{\sffamily\textsc{Lösung}}\par\smallskip\small
\begin{tabular}{@{}r@{\enspace}l@{\enspace}l@{}}
\textbf{1.} & \ldots & f4\\
\textbf{2.} & \symknight{}e6 & f5+\\
\textbf{3.} & \symking{}g5 & f6\\
\textbf{4.} & \symknight{}g7 & fxg7\\
\textbf{5.} & \symking{}f4 & g8=\symqueen{}\\
\textbf{6.} & \symking{}e3 & \symqueen{}xb3+\\
\textbf{7.} & \symking{}d2 & \symking{}a2\\
\textbf{8.} & \symking{}c1 & \symking{}a3\\
\textbf{9.} & \symking{}b1 & \symqueen{}d1\#\\
\end{tabular}\par\vspace{3mm}
\qrcode[height=12mm]{https://helpman.komtera.lt/?fen=3n4/8/6k1/8/8/1p6/1P3P2/K7&moves=8.5}\\{\sffamily\tiny\color{ink}scannen zum Lösen}
\end{minipage}
\vspace{2mm}

{\raggedright\noindent\textcolor{ink}{\sffamily\textsc{Themen}}\enspace \themetag{reines Matt} \themetag{Mustermatt} \themetag{Idealmatt} \themetag{Umwandlung} \themetag{Excelsior}\par}

\clearpage
\phantomsection\label{hm:KPPvkqp}
\begin{tcolorbox}[enhanced,colback=ink,colframe=ink,arc=0pt,boxrule=0pt,left=4mm,right=4mm,top=2.5mm,bottom=2.5mm]
\color{white}\sffamily{\LARGE\bfseries KPPvkqp}\hfill{\Large h\#8.5}\\[0.5mm]
{\small 6 Steine\quad Weiß am Zug}\hfill{\small Nr.~365\qquad\href{https://helpman.komtera.lt/?fen=8/8/p7/8/P7/1Kq5/P7/7k&moves=8.5}{\color{white}\bfseries lösen\,\ensuremath{\nearrow}}}
\end{tcolorbox}
\vspace{2mm}
\noindent\begin{minipage}[t]{36mm}\vspace{0pt}\sffamily\small\raggedright
\hmfigure{längste eindeutige}{{\Large h\#8.5}}\hmfigure{eindeutig bei dieser Länge}{{\Large 4} Stellungen}\hmfigure{längstes Matt}{{\Large h\#14.5}}\hmfigure{in der Klasse}{25,6 Milliarden Stellungen,\\ 8,06 Milliarden mit Hilfsmatt}
\end{minipage}\hfill
\begin{minipage}[t]{80mm}\vspace{0pt}\begin{tikzpicture}\begin{axis}[width=72mm,height=42mm,ybar,bar width=1.6pt,xmin=-1,xmax=30,ymin=0,ymax=1.698,axis y line*=left,axis x line*=bottom,xlabel={\scriptsize Halbzüge bis Matt},ylabel={\scriptsize Stellungen (Mrd.)},tick label style={font=\tiny},label style={font=\scriptsize},enlarge x limits=false,ylabel shift=-3pt,xlabel shift=-3pt]\addplot[fill=ink!55,draw=none] coordinates {(0,2.033e-05) (1,0.001585) (2,0.006519) (3,0.05695) (4,0.1546) (5,0.6331) (6,0.8515) (7,1.543) (8,1.086) (9,1.399) (10,0.7682) (11,0.8654) (12,0.4234) (13,0.231) (14,0.03354) (15,0.007026) (16,0.0001352) (17,8.289e-05) (18,7.46e-07) (19,4.129e-05) (20,1e-08) (21,4.277e-05) (22,8e-09) (23,2.039e-05) (24,2e-09) (25,9.441e-06) (27,4.408e-06) (29,7.2e-07)};\draw[densely dotted] (axis cs:17,0) -- (axis cs:17,1.698);\end{axis}\begin{axis}[width=72mm,height=42mm,axis y line*=right,axis x line=none,ymode=log,xmin=-1,xmax=30,ylabel={\scriptsize eindeutig (log)},tick label style={font=\tiny},label style={font=\scriptsize},ylabel shift=-3pt]\addplot[mark=*,mark size=0.7pt,draw=black,thin] coordinates {(0,20327) (1,487270) (2,1689488) (3,17796642) (4,30071557) (5,56067403) (6,24370576) (7,26917721) (8,9050977) (9,9562523) (10,2919691) (11,2444682) (12,480067) (13,110644) (14,1243) (15,292) (17,4)};\end{axis}\end{tikzpicture}\end{minipage}
\vspace{3mm}

\noindent\begin{minipage}[t]{60mm}\vspace{0pt}\sbox\hmboard{\chessboard[setfen=8/8/p7/8/P7/1Kq5/P7/7k w - - 0 1]}\usebox\hmboard\\[1mm]
\makebox[\wd\hmboard]{\sffamily\hspace{3mm}\textbf{h\#8.5}\hfill(3+3)}\end{minipage}\hfill
\begin{minipage}[t]{52mm}\vspace{0pt}
\textcolor{ink}{\sffamily\textsc{Lösung}}\par\smallskip\small
\begin{tabular}{@{}r@{\enspace}l@{\enspace}l@{}}
\textbf{1.} & \ldots & \symking{}xc3\\
\textbf{2.} & \symking{}g2 & \symking{}b4\\
\textbf{3.} & a5+ & \symking{}xa5\\
\textbf{4.} & \symking{}f3 & \symking{}b4\\
\textbf{5.} & \symking{}e4 & a5\\
\textbf{6.} & \symking{}d5 & a6\\
\textbf{7.} & \symking{}c6 & a7\\
\textbf{8.} & \symking{}b7 & \symking{}c5\\
\textbf{9.} & \symking{}a6 & a8=\symqueen{}\#\\
\end{tabular}\par\vspace{3mm}
\qrcode[height=12mm]{https://helpman.komtera.lt/?fen=8/8/p7/8/P7/1Kq5/P7/7k&moves=8.5}\\{\sffamily\tiny\color{ink}scannen zum Lösen}
\end{minipage}
\vspace{2mm}

{\raggedright\noindent\textcolor{ink}{\sffamily\textsc{Themen}}\enspace \themetag{reines Matt} \themetag{Mustermatt} \themetag{Spiegelmatt} \themetag{Umwandlung} \themetag{Switchback}\par\smallskip{\footnotesize\textit{Hinweis:} die Partei am Zug steht im Diagramm im Schach; die Lösung beginnt mit einem Schlag.}\par}

\clearpage
\phantomsection\label{hm:KPvkbbn}
\begin{tcolorbox}[enhanced,colback=ink,colframe=ink,arc=0pt,boxrule=0pt,left=4mm,right=4mm,top=2.5mm,bottom=2.5mm]
\color{white}\sffamily{\LARGE\bfseries KPvkbbn}\hfill{\Large h\#8.5}\\[0.5mm]
{\small 6 Steine\quad Weiß am Zug}\hfill{\small Nr.~366\qquad\href{https://helpman.komtera.lt/?fen=K7/8/b7/n7/8/8/P7/6bk&moves=8.5}{\color{white}\bfseries lösen\,\ensuremath{\nearrow}}}
\end{tcolorbox}
\vspace{2mm}
\noindent\begin{minipage}[t]{36mm}\vspace{0pt}\sffamily\small\raggedright
\hmfigure{längste eindeutige}{{\Large h\#8.5}}\hmfigure{eindeutig bei dieser Länge}{{\Large 11} Stellungen}\hmfigure{längstes Matt}{{\Large h\#9.5}}\hmfigure{in der Klasse}{45,4 Milliarden Stellungen,\\ 15,1 Milliarden mit Hilfsmatt}
\end{minipage}\hfill
\begin{minipage}[t]{80mm}\vspace{0pt}\begin{tikzpicture}\begin{axis}[width=72mm,height=42mm,ybar,bar width=1.6pt,xmin=-1,xmax=20,ymin=0,ymax=2.944,axis y line*=left,axis x line*=bottom,xlabel={\scriptsize Halbzüge bis Matt},ylabel={\scriptsize Stellungen (Mrd.)},tick label style={font=\tiny},label style={font=\scriptsize},enlarge x limits=false,ylabel shift=-3pt,xlabel shift=-3pt]\addplot[fill=ink!55,draw=none] coordinates {(0,1.193e-05) (1,0.001022) (2,0.003494) (3,0.03197) (4,0.07995) (5,0.3874) (6,0.5883) (7,1.424) (8,1.347) (9,2.278) (10,1.713) (11,2.676) (12,1.93) (13,1.811) (14,0.6473) (15,0.1976) (16,0.01768) (17,0.004068) (18,1.779e-05) (19,2.148e-06)};\draw[densely dotted] (axis cs:17,0) -- (axis cs:17,2.944);\end{axis}\begin{axis}[width=72mm,height=42mm,axis y line*=right,axis x line=none,ymode=log,xmin=-1,xmax=20,ylabel={\scriptsize eindeutig (log)},tick label style={font=\tiny},label style={font=\scriptsize},ylabel shift=-3pt]\addplot[mark=*,mark size=0.7pt,draw=black,thin] coordinates {(0,11932) (1,266502) (2,601323) (3,8109937) (4,13460283) (5,37731667) (6,19617400) (7,29518136) (8,12283006) (9,17834696) (10,7343117) (11,11425875) (12,3651708) (13,2377304) (14,124364) (15,30255) (16,20) (17,11)};\end{axis}\end{tikzpicture}\end{minipage}
\vspace{3mm}

\noindent\begin{minipage}[t]{60mm}\vspace{0pt}\sbox\hmboard{\chessboard[setfen=K7/8/b7/n7/8/8/P7/6bk w - - 0 1]}\usebox\hmboard\\[1mm]
\makebox[\wd\hmboard]{\sffamily\hspace{3mm}\textbf{h\#8.5}\hfill(2+4)}\end{minipage}\hfill
\begin{minipage}[t]{52mm}\vspace{0pt}
\textcolor{ink}{\sffamily\textsc{Lösung}}\par\smallskip\small
\begin{tabular}{@{}r@{\enspace}l@{\enspace}l@{}}
\textbf{1.} & \ldots & a4\\
\textbf{2.} & \symbishop{}b5 & axb5\\
\textbf{3.} & \symking{}g2 & b6\\
\textbf{4.} & \symking{}f3 & b7\\
\textbf{5.} & \symking{}e4 & b8=\symqueen{}\\
\textbf{6.} & \symking{}d5 & \symqueen{}e5+\\
\textbf{7.} & \symking{}c6 & \symking{}b8\\
\textbf{8.} & \symking{}b6 & \symking{}c8\\
\textbf{9.} & \symking{}a7 & \symqueen{}xa5\#\\
\end{tabular}\par\vspace{3mm}
\qrcode[height=12mm]{https://helpman.komtera.lt/?fen=K7/8/b7/n7/8/8/P7/6bk&moves=8.5}\\{\sffamily\tiny\color{ink}scannen zum Lösen}
\end{minipage}
\vspace{2mm}

{\raggedright\noindent\textcolor{ink}{\sffamily\textsc{Themen}}\enspace \themetag{reines Matt} \themetag{Mustermatt} \themetag{Spiegelmatt} \themetag{Umwandlung} \themetag{Excelsior}\par}

\clearpage
\phantomsection\label{hm:KPvkbnn}
\begin{tcolorbox}[enhanced,colback=ink,colframe=ink,arc=0pt,boxrule=0pt,left=4mm,right=4mm,top=2.5mm,bottom=2.5mm]
\color{white}\sffamily{\LARGE\bfseries KPvkbnn}\hfill{\Large h\#8.5}\\[0.5mm]
{\small 6 Steine\quad Weiß am Zug}\hfill{\small Nr.~367\qquad\href{https://helpman.komtera.lt/?fen=8/K7/8/2n5/2n5/8/2P3b1/7k&moves=8.5}{\color{white}\bfseries lösen\,\ensuremath{\nearrow}}}
\end{tcolorbox}
\vspace{2mm}
\noindent\begin{minipage}[t]{36mm}\vspace{0pt}\sffamily\small\raggedright
\hmfigure{längste eindeutige}{{\Large h\#8.5}}\hmfigure{eindeutig bei dieser Länge}{{\Large 20} Stellungen}\hmfigure{längstes Matt}{{\Large h\#9.5}}\hmfigure{in der Klasse}{45,4 Milliarden Stellungen,\\ 15,5 Milliarden mit Hilfsmatt}
\end{minipage}\hfill
\begin{minipage}[t]{80mm}\vspace{0pt}\begin{tikzpicture}\begin{axis}[width=72mm,height=42mm,ybar,bar width=1.6pt,xmin=-1,xmax=20,ymin=0,ymax=2.803,axis y line*=left,axis x line*=bottom,xlabel={\scriptsize Halbzüge bis Matt},ylabel={\scriptsize Stellungen (Mrd.)},tick label style={font=\tiny},label style={font=\scriptsize},enlarge x limits=false,ylabel shift=-3pt,xlabel shift=-3pt]\addplot[fill=ink!55,draw=none] coordinates {(0,9.766e-06) (1,0.001297) (2,0.004111) (3,0.03427) (4,0.07845) (5,0.3635) (6,0.5318) (7,1.295) (8,1.292) (9,2.157) (10,1.711) (11,2.548) (12,1.978) (13,2.091) (14,1.059) (15,0.3164) (16,0.04012) (17,0.007662) (18,5.099e-05) (19,5.814e-06)};\draw[densely dotted] (axis cs:17,0) -- (axis cs:17,2.803);\end{axis}\begin{axis}[width=72mm,height=42mm,axis y line*=right,axis x line=none,ymode=log,xmin=-1,xmax=20,ylabel={\scriptsize eindeutig (log)},tick label style={font=\tiny},label style={font=\scriptsize},ylabel shift=-3pt]\addplot[mark=*,mark size=0.7pt,draw=black,thin] coordinates {(0,9766) (1,291768) (2,632998) (3,7426828) (4,11985975) (5,34136635) (6,18888977) (7,26499330) (8,11394378) (9,16261357) (10,7849426) (11,12573657) (12,5397873) (13,3451833) (14,222439) (15,48253) (16,34) (17,20)};\end{axis}\end{tikzpicture}\end{minipage}
\vspace{3mm}

\noindent\begin{minipage}[t]{60mm}\vspace{0pt}\sbox\hmboard{\chessboard[setfen=8/K7/8/2n5/2n5/8/2P3b1/7k w - - 0 1]}\usebox\hmboard\\[1mm]
\makebox[\wd\hmboard]{\sffamily\hspace{3mm}\textbf{h\#8.5}\hfill(2+4)}\end{minipage}\hfill
\begin{minipage}[t]{52mm}\vspace{0pt}
\textcolor{ink}{\sffamily\textsc{Lösung}}\par\smallskip\small
\begin{tabular}{@{}r@{\enspace}l@{\enspace}l@{}}
\textbf{1.} & \ldots & \symking{}b8\\
\textbf{2.} & \symknight{}d3 & cxd3\\
\textbf{3.} & \symking{}h2 & d4\\
\textbf{4.} & \symking{}g3 & d5\\
\textbf{5.} & \symking{}f4 & d6\\
\textbf{6.} & \symking{}e5 & d7\\
\textbf{7.} & \symking{}d6 & \symking{}a7\\
\textbf{8.} & \symking{}c7 & \symking{}a6\\
\textbf{9.} & \symking{}b8 & d8=\symqueen{}\#\\
\end{tabular}\par\vspace{3mm}
\qrcode[height=12mm]{https://helpman.komtera.lt/?fen=8/K7/8/2n5/2n5/8/2P3b1/7k&moves=8.5}\\{\sffamily\tiny\color{ink}scannen zum Lösen}
\end{minipage}
\vspace{2mm}

{\raggedright\noindent\textcolor{ink}{\sffamily\textsc{Themen}}\enspace \themetag{reines Matt} \themetag{Mustermatt} \themetag{Spiegelmatt} \themetag{Umwandlung} \themetag{Excelsior} \themetag{Switchback} \themetag{ohne Schach} \themetag{Umnov}\par}

\clearpage
\phantomsection\label{hm:KPvknnn}
\begin{tcolorbox}[enhanced,colback=ink,colframe=ink,arc=0pt,boxrule=0pt,left=4mm,right=4mm,top=2.5mm,bottom=2.5mm]
\color{white}\sffamily{\LARGE\bfseries KPvknnn}\hfill{\Large h\#8.5}\\[0.5mm]
{\small 6 Steine\quad Weiß am Zug}\hfill{\small Nr.~368\qquad\href{https://helpman.komtera.lt/?fen=K7/8/3nn3/n7/8/8/P7/7k&moves=8.5}{\color{white}\bfseries lösen\,\ensuremath{\nearrow}}}
\end{tcolorbox}
\vspace{2mm}
\noindent\begin{minipage}[t]{36mm}\vspace{0pt}\sffamily\small\raggedright
\hmfigure{längste eindeutige}{{\Large h\#8.5}}\hmfigure{eindeutig bei dieser Länge}{{\Large 8} Stellungen}\hmfigure{längstes Matt}{{\Large h\#9.5}}\hmfigure{in der Klasse}{45,4 Milliarden Stellungen,\\ 5,28 Milliarden mit Hilfsmatt}
\end{minipage}\hfill
\begin{minipage}[t]{80mm}\vspace{0pt}\begin{tikzpicture}\begin{axis}[width=72mm,height=42mm,ybar,bar width=1.6pt,xmin=-1,xmax=20,ymin=0,ymax=898.7,axis y line*=left,axis x line*=bottom,xlabel={\scriptsize Halbzüge bis Matt},ylabel={\scriptsize Stellungen (Mio.)},tick label style={font=\tiny},label style={font=\scriptsize},enlarge x limits=false,ylabel shift=-3pt,xlabel shift=-3pt]\addplot[fill=ink!55,draw=none] coordinates {(0,0.00255) (1,0.5521) (2,1.578) (3,12.19) (4,26.36) (5,115.7) (6,165.6) (7,413.8) (8,428) (9,699.3) (10,577) (11,817) (12,643.3) (13,720.3) (14,447.1) (15,156.3) (16,47.75) (17,3.097) (18,0.0336) (19,0.004255)};\draw[densely dotted] (axis cs:17,0) -- (axis cs:17,898.7);\end{axis}\begin{axis}[width=72mm,height=42mm,axis y line*=right,axis x line=none,ymode=log,xmin=-1,xmax=20,ylabel={\scriptsize eindeutig (log)},tick label style={font=\tiny},label style={font=\scriptsize},ylabel shift=-3pt]\addplot[mark=*,mark size=0.7pt,draw=black,thin] coordinates {(0,2550) (1,104445) (2,220882) (3,2310724) (4,3442429) (5,10219007) (6,5977929) (7,7648551) (8,3555870) (9,4810940) (10,2480814) (11,3888334) (12,1791240) (13,1153646) (14,96120) (15,14307) (16,18) (17,8)};\end{axis}\end{tikzpicture}\end{minipage}
\vspace{3mm}

\noindent\begin{minipage}[t]{60mm}\vspace{0pt}\sbox\hmboard{\chessboard[setfen=K7/8/3nn3/n7/8/8/P7/7k w - - 0 1]}\usebox\hmboard\\[1mm]
\makebox[\wd\hmboard]{\sffamily\hspace{3mm}\textbf{h\#8.5}\hfill(2+4)}\end{minipage}\hfill
\begin{minipage}[t]{52mm}\vspace{0pt}
\textcolor{ink}{\sffamily\textsc{Lösung}}\par\smallskip\small
\begin{tabular}{@{}r@{\enspace}l@{\enspace}l@{}}
\textbf{1.} & \ldots & a4\\
\textbf{2.} & \symknight{}b5 & axb5\\
\textbf{3.} & \symking{}g2 & b6\\
\textbf{4.} & \symking{}f3 & b7\\
\textbf{5.} & \symking{}e4 & b8=\symqueen{}\\
\textbf{6.} & \symking{}d5 & \symqueen{}e5+\\
\textbf{7.} & \symking{}c6 & \symking{}b8\\
\textbf{8.} & \symking{}b6 & \symking{}c8\\
\textbf{9.} & \symking{}a7 & \symqueen{}xa5\#\\
\end{tabular}\par\vspace{3mm}
\qrcode[height=12mm]{https://helpman.komtera.lt/?fen=K7/8/3nn3/n7/8/8/P7/7k&moves=8.5}\\{\sffamily\tiny\color{ink}scannen zum Lösen}
\end{minipage}
\vspace{2mm}

{\raggedright\noindent\textcolor{ink}{\sffamily\textsc{Themen}}\enspace \themetag{reines Matt} \themetag{Mustermatt} \themetag{Spiegelmatt} \themetag{Umwandlung} \themetag{Excelsior}\par}

\clearpage
\phantomsection\label{hm:KPvkqbn}
\begin{tcolorbox}[enhanced,colback=ink,colframe=ink,arc=0pt,boxrule=0pt,left=4mm,right=4mm,top=2.5mm,bottom=2.5mm]
\color{white}\sffamily{\LARGE\bfseries KPvkqbn}\hfill{\Large h\#8.5}\\[0.5mm]
{\small 6 Steine\quad Weiß am Zug}\hfill{\small Nr.~369\qquad\href{https://helpman.komtera.lt/?fen=7k/8/8/8/7n/b6P/4q3/1K6&moves=8.5}{\color{white}\bfseries lösen\,\ensuremath{\nearrow}}}
\end{tcolorbox}
\vspace{2mm}
\noindent\begin{minipage}[t]{36mm}\vspace{0pt}\sffamily\small\raggedright
\hmfigure{längste eindeutige}{{\Large h\#8.5}}\hmfigure{eindeutig bei dieser Länge}{{\Large 2} Stellungen}\hmfigure{längstes Matt}{{\Large h\#11.5}}\hmfigure{in der Klasse}{45,4 Milliarden Stellungen,\\ 27,3 Milliarden mit Hilfsmatt}
\end{minipage}\hfill
\begin{minipage}[t]{80mm}\vspace{0pt}\begin{tikzpicture}\begin{axis}[width=72mm,height=42mm,ybar,bar width=1.6pt,xmin=-1,xmax=24,ymin=0,ymax=5.706,axis y line*=left,axis x line*=bottom,xlabel={\scriptsize Halbzüge bis Matt},ylabel={\scriptsize Stellungen (Mrd.)},tick label style={font=\tiny},label style={font=\scriptsize},enlarge x limits=false,ylabel shift=-3pt,xlabel shift=-3pt]\addplot[fill=ink!55,draw=none] coordinates {(0,8.803e-06) (1,0.001051) (2,0.004659) (3,0.04631) (4,0.1188) (5,0.6556) (6,0.9233) (7,2.714) (8,2.069) (9,4.386) (10,2.673) (11,5.187) (12,3.009) (13,3.992) (14,1.016) (15,0.4992) (16,0.01715) (17,0.005171) (18,9.022e-06) (19,1.846e-06) (20,1.02e-07) (21,1.18e-07) (23,2.1e-08)};\draw[densely dotted] (axis cs:17,0) -- (axis cs:17,5.706);\end{axis}\begin{axis}[width=72mm,height=42mm,axis y line*=right,axis x line=none,ymode=log,xmin=-1,xmax=24,ylabel={\scriptsize eindeutig (log)},tick label style={font=\tiny},label style={font=\scriptsize},ylabel shift=-3pt]\addplot[mark=*,mark size=0.7pt,draw=black,thin] coordinates {(0,8803) (1,296087) (2,812880) (3,12055989) (4,19049967) (5,67198887) (6,28518283) (7,58305755) (8,18303295) (9,35495239) (10,11642109) (11,22495763) (12,5615154) (13,5082210) (14,108549) (15,37654) (16,12) (17,2)};\end{axis}\end{tikzpicture}\end{minipage}
\vspace{3mm}

\noindent\begin{minipage}[t]{60mm}\vspace{0pt}\sbox\hmboard{\chessboard[setfen=7k/8/8/8/7n/b6P/4q3/1K6 w - - 0 1]}\usebox\hmboard\\[1mm]
\makebox[\wd\hmboard]{\sffamily\hspace{3mm}\textbf{h\#8.5}\hfill(2+4)}\end{minipage}\hfill
\begin{minipage}[t]{52mm}\vspace{0pt}
\textcolor{ink}{\sffamily\textsc{Lösung}}\par\smallskip\small
\begin{tabular}{@{}r@{\enspace}l@{\enspace}l@{}}
\textbf{1.} & \ldots & \symking{}a1\\
\textbf{2.} & \symknight{}g6 & h4\\
\textbf{3.} & \symking{}g7 & h5\\
\textbf{4.} & \symking{}f6 & hxg6\\
\textbf{5.} & \symking{}e5 & g7\\
\textbf{6.} & \symking{}d4 & g8=\symqueen{}\\
\textbf{7.} & \symqueen{}e5 & \symking{}a2\\
\textbf{8.} & \symbishop{}c1 & \symking{}b3\\
\textbf{9.} & \symbishop{}e3 & \symqueen{}c4\#\\
\end{tabular}\par\vspace{3mm}
\qrcode[height=12mm]{https://helpman.komtera.lt/?fen=7k/8/8/8/7n/b6P/4q3/1K6&moves=8.5}\\{\sffamily\tiny\color{ink}scannen zum Lösen}
\end{minipage}
\vspace{2mm}

{\raggedright\noindent\textcolor{ink}{\sffamily\textsc{Themen}}\enspace \themetag{Satzspiel} \themetag{Umwandlung} \themetag{Selbstblock} \themetag{ohne Schach} \themetag{Umnov}\par}

\clearpage
\phantomsection\label{hm:KPvkqnn}
\begin{tcolorbox}[enhanced,colback=ink,colframe=ink,arc=0pt,boxrule=0pt,left=4mm,right=4mm,top=2.5mm,bottom=2.5mm]
\color{white}\sffamily{\LARGE\bfseries KPvkqnn}\hfill{\Large h\#8.5}\\[0.5mm]
{\small 6 Steine\quad Weiß am Zug}\hfill{\small Nr.~370\qquad\href{https://helpman.komtera.lt/?fen=8/3nn3/K7/8/1q6/1P6/8/7k&moves=8.5}{\color{white}\bfseries lösen\,\ensuremath{\nearrow}}}
\end{tcolorbox}
\vspace{2mm}
\noindent\begin{minipage}[t]{36mm}\vspace{0pt}\sffamily\small\raggedright
\hmfigure{längste eindeutige}{{\Large h\#8.5}}\hmfigure{eindeutig bei dieser Länge}{{\Large 40} Stellungen}\hmfigure{längstes Matt}{{\Large h\#10.5}}\hmfigure{in der Klasse}{45,4 Milliarden Stellungen,\\ 14 Milliarden mit Hilfsmatt}
\end{minipage}\hfill
\begin{minipage}[t]{80mm}\vspace{0pt}\begin{tikzpicture}\begin{axis}[width=72mm,height=42mm,ybar,bar width=1.6pt,xmin=-1,xmax=22,ymin=0,ymax=2.778,axis y line*=left,axis x line*=bottom,xlabel={\scriptsize Halbzüge bis Matt},ylabel={\scriptsize Stellungen (Mrd.)},tick label style={font=\tiny},label style={font=\scriptsize},enlarge x limits=false,ylabel shift=-3pt,xlabel shift=-3pt]\addplot[fill=ink!55,draw=none] coordinates {(0,2.914e-06) (1,0.0006463) (2,0.002797) (3,0.02553) (4,0.0596) (5,0.307) (6,0.4091) (7,1.197) (8,0.9997) (9,2.106) (10,1.347) (11,2.525) (12,1.573) (13,2.176) (14,0.7838) (15,0.4145) (16,0.03137) (17,0.01331) (18,7.88e-05) (19,2.007e-05) (20,2e-09) (21,1e-09)};\draw[densely dotted] (axis cs:17,0) -- (axis cs:17,2.778);\end{axis}\begin{axis}[width=72mm,height=42mm,axis y line*=right,axis x line=none,ymode=log,xmin=-1,xmax=22,ylabel={\scriptsize eindeutig (log)},tick label style={font=\tiny},label style={font=\scriptsize},ylabel shift=-3pt]\addplot[mark=*,mark size=0.7pt,draw=black,thin] coordinates {(0,2914) (1,160936) (2,419121) (3,5604652) (4,8027163) (5,28538671) (6,12548416) (7,25292979) (8,8233103) (9,15795043) (10,5572455) (11,11637998) (12,3660883) (13,3944800) (14,181184) (15,91486) (16,85) (17,40)};\end{axis}\end{tikzpicture}\end{minipage}
\vspace{3mm}

\noindent\begin{minipage}[t]{60mm}\vspace{0pt}\sbox\hmboard{\chessboard[setfen=8/3nn3/K7/8/1q6/1P6/8/7k w - - 0 1]}\usebox\hmboard\\[1mm]
\makebox[\wd\hmboard]{\sffamily\hspace{3mm}\textbf{h\#8.5}\hfill(2+4)}\end{minipage}\hfill
\begin{minipage}[t]{52mm}\vspace{0pt}
\textcolor{ink}{\sffamily\textsc{Lösung}}\par\smallskip\small
\begin{tabular}{@{}r@{\enspace}l@{\enspace}l@{}}
\textbf{1.} & \ldots & \symking{}a7\\
\textbf{2.} & \symqueen{}c4 & bxc4\\
\textbf{3.} & \symking{}g2 & c5\\
\textbf{4.} & \symking{}f3 & c6\\
\textbf{5.} & \symking{}e4 & cxd7\\
\textbf{6.} & \symking{}d5 & d8=\symqueen{}+\\
\textbf{7.} & \symking{}c6 & \symqueen{}d5+\\
\textbf{8.} & \symking{}c7 & \symking{}a6\\
\textbf{9.} & \symking{}b8 & \symqueen{}b7\#\\
\end{tabular}\par\vspace{3mm}
\qrcode[height=12mm]{https://helpman.komtera.lt/?fen=8/3nn3/K7/8/1q6/1P6/8/7k&moves=8.5}\\{\sffamily\tiny\color{ink}scannen zum Lösen}
\end{minipage}
\vspace{2mm}

{\raggedright\noindent\textcolor{ink}{\sffamily\textsc{Themen}}\enspace \themetag{Satzspiel} \themetag{Umwandlung} \themetag{Switchback} \themetag{Umnov}\par}

\clearpage
\phantomsection\label{hm:KPvkqqn}
\begin{tcolorbox}[enhanced,colback=ink,colframe=ink,arc=0pt,boxrule=0pt,left=4mm,right=4mm,top=2.5mm,bottom=2.5mm]
\color{white}\sffamily{\LARGE\bfseries KPvkqqn}\hfill{\Large h\#8.5}\\[0.5mm]
{\small 6 Steine\quad Weiß am Zug}\hfill{\small Nr.~371\qquad\href{https://helpman.komtera.lt/?fen=8/q7/K7/8/8/2n5/2P5/2q4k&moves=8.5}{\color{white}\bfseries lösen\,\ensuremath{\nearrow}}}
\end{tcolorbox}
\vspace{2mm}
\noindent\begin{minipage}[t]{36mm}\vspace{0pt}\sffamily\small\raggedright
\hmfigure{längste eindeutige}{{\Large h\#8.5}}\hmfigure{eindeutig bei dieser Länge}{{\Large 7} Stellungen}\hmfigure{längstes Matt}{{\Large h\#9.5}}\hmfigure{in der Klasse}{45,4 Milliarden Stellungen,\\ 12,2 Milliarden mit Hilfsmatt}
\end{minipage}\hfill
\begin{minipage}[t]{80mm}\vspace{0pt}\begin{tikzpicture}\begin{axis}[width=72mm,height=42mm,ybar,bar width=1.6pt,xmin=-1,xmax=20,ymin=0,ymax=2.667,axis y line*=left,axis x line*=bottom,xlabel={\scriptsize Halbzüge bis Matt},ylabel={\scriptsize Stellungen (Mrd.)},tick label style={font=\tiny},label style={font=\scriptsize},enlarge x limits=false,ylabel shift=-3pt,xlabel shift=-3pt]\addplot[fill=ink!55,draw=none] coordinates {(0,7.66e-07) (1,0.0002542) (2,0.001381) (3,0.01546) (4,0.04312) (5,0.2664) (6,0.3574) (7,1.243) (8,0.7798) (9,2.047) (10,0.9841) (11,2.425) (12,1.16) (13,2.107) (14,0.3959) (15,0.3263) (16,0.005878) (17,0.003546) (18,1.299e-06) (19,2.242e-06)};\draw[densely dotted] (axis cs:17,0) -- (axis cs:17,2.667);\end{axis}\begin{axis}[width=72mm,height=42mm,axis y line*=right,axis x line=none,ymode=log,xmin=-1,xmax=20,ylabel={\scriptsize eindeutig (log)},tick label style={font=\tiny},label style={font=\scriptsize},ylabel shift=-3pt]\addplot[mark=*,mark size=0.7pt,draw=black,thin] coordinates {(0,766) (1,76404) (2,254375) (3,4286409) (4,7209220) (5,29974527) (6,11260064) (7,29633062) (8,6658834) (9,16658267) (10,4094864) (11,10750550) (12,2217116) (13,3004436) (14,51891) (15,33680) (16,4) (17,7)};\end{axis}\end{tikzpicture}\end{minipage}
\vspace{3mm}

\noindent\begin{minipage}[t]{60mm}\vspace{0pt}\sbox\hmboard{\chessboard[setfen=8/q7/K7/8/8/2n5/2P5/2q4k w - - 0 1]}\usebox\hmboard\\[1mm]
\makebox[\wd\hmboard]{\sffamily\hspace{3mm}\textbf{h\#8.5}\hfill(2+4)}\end{minipage}\hfill
\begin{minipage}[t]{52mm}\vspace{0pt}
\textcolor{ink}{\sffamily\textsc{Lösung}}\par\smallskip\small
\begin{tabular}{@{}r@{\enspace}l@{\enspace}l@{}}
\textbf{1.} & \ldots & \symking{}xa7\\
\textbf{2.} & \symknight{}d5 & c4\\
\textbf{3.} & \symking{}g2 & cxd5\\
\textbf{4.} & \symking{}f3 & d6\\
\textbf{5.} & \symking{}e4 & d7\\
\textbf{6.} & \symking{}d5 & d8=\symqueen{}+\\
\textbf{7.} & \symking{}c6 & \symqueen{}d5+\\
\textbf{8.} & \symking{}c7 & \symking{}a6\\
\textbf{9.} & \symking{}b8 & \symqueen{}b7\#\\
\end{tabular}\par\vspace{3mm}
\qrcode[height=12mm]{https://helpman.komtera.lt/?fen=8/q7/K7/8/8/2n5/2P5/2q4k&moves=8.5}\\{\sffamily\tiny\color{ink}scannen zum Lösen}
\end{minipage}
\vspace{2mm}

{\raggedright\noindent\textcolor{ink}{\sffamily\textsc{Themen}}\enspace \themetag{Umwandlung} \themetag{Excelsior} \themetag{Switchback} \themetag{Rundlauf} \themetag{Umnov}\par\smallskip{\footnotesize\textit{Hinweis:} die Partei am Zug steht im Diagramm im Schach; die Lösung beginnt mit einem Schlag.}\par}

\clearpage
\phantomsection\label{hm:KPvkqqp}
\begin{tcolorbox}[enhanced,colback=ink,colframe=ink,arc=0pt,boxrule=0pt,left=4mm,right=4mm,top=2.5mm,bottom=2.5mm]
\color{white}\sffamily{\LARGE\bfseries KPvkqqp}\hfill{\Large h\#8.5}\\[0.5mm]
{\small 6 Steine\quad Weiß am Zug}\hfill{\small Nr.~372\qquad\href{https://helpman.komtera.lt/?fen=8/7k/8/5p2/1q6/5P2/2q5/K7&moves=8.5}{\color{white}\bfseries lösen\,\ensuremath{\nearrow}}}
\end{tcolorbox}
\vspace{2mm}
\noindent\begin{minipage}[t]{36mm}\vspace{0pt}\sffamily\small\raggedright
\hmfigure{längste eindeutige}{{\Large h\#8.5}}\hmfigure{eindeutig bei dieser Länge}{{\Large 61} Stellungen}\hmfigure{längstes Matt}{{\Large h\#10.5}}\hmfigure{in der Klasse}{34,1 Milliarden Stellungen,\\ 9,27 Milliarden mit Hilfsmatt}
\end{minipage}\hfill
\begin{minipage}[t]{80mm}\vspace{0pt}\begin{tikzpicture}\begin{axis}[width=72mm,height=42mm,ybar,bar width=1.6pt,xmin=-1,xmax=22,ymin=0,ymax=2.029,axis y line*=left,axis x line*=bottom,xlabel={\scriptsize Halbzüge bis Matt},ylabel={\scriptsize Stellungen (Mrd.)},tick label style={font=\tiny},label style={font=\scriptsize},enlarge x limits=false,ylabel shift=-3pt,xlabel shift=-3pt]\addplot[fill=ink!55,draw=none] coordinates {(0,8.5e-07) (1,0.0002785) (2,0.001479) (3,0.01516) (4,0.04302) (5,0.253) (6,0.3399) (7,1.076) (8,0.6721) (9,1.614) (10,0.7921) (11,1.844) (12,0.8841) (13,1.366) (14,0.2006) (15,0.1553) (16,0.006758) (17,0.006602) (18,5.559e-05) (19,6.172e-05) (20,1.1e-08) (21,1.2e-08)};\draw[densely dotted] (axis cs:17,0) -- (axis cs:17,2.029);\end{axis}\begin{axis}[width=72mm,height=42mm,axis y line*=right,axis x line=none,ymode=log,xmin=-1,xmax=22,ylabel={\scriptsize eindeutig (log)},tick label style={font=\tiny},label style={font=\scriptsize},ylabel shift=-3pt]\addplot[mark=*,mark size=0.7pt,draw=black,thin] coordinates {(0,850) (1,91600) (2,332182) (3,4433420) (4,7789163) (5,28557685) (6,10833511) (7,24534987) (8,6305640) (9,13869814) (10,3853811) (11,8417921) (12,1744046) (13,1698492) (14,54954) (15,36384) (16,41) (17,61)};\end{axis}\end{tikzpicture}\end{minipage}
\vspace{3mm}

\noindent\begin{minipage}[t]{60mm}\vspace{0pt}\sbox\hmboard{\chessboard[setfen=8/7k/8/5p2/1q6/5P2/2q5/K7 w - - 0 1]}\usebox\hmboard\\[1mm]
\makebox[\wd\hmboard]{\sffamily\hspace{3mm}\textbf{h\#8.5}\hfill(2+4)}\end{minipage}\hfill
\begin{minipage}[t]{52mm}\vspace{0pt}
\textcolor{ink}{\sffamily\textsc{Lösung}}\par\smallskip\small
\begin{tabular}{@{}r@{\enspace}l@{\enspace}l@{}}
\textbf{1.} & \ldots & f4\\
\textbf{2.} & \symqueen{}c6 & \symking{}a2\\
\textbf{3.} & \symqueen{}e7 & \symking{}b3\\
\textbf{4.} & \symqueen{}g5 & fxg5\\
\textbf{5.} & \symking{}g7 & g6\\
\textbf{6.} & \symking{}f6 & g7\\
\textbf{7.} & \symking{}e5 & g8=\symqueen{}\\
\textbf{8.} & \symqueen{}f6 & \symking{}c4\\
\textbf{9.} & f4 & \symqueen{}d5\#\\
\end{tabular}\par\vspace{3mm}
\qrcode[height=12mm]{https://helpman.komtera.lt/?fen=8/7k/8/5p2/1q6/5P2/2q5/K7&moves=8.5}\\{\sffamily\tiny\color{ink}scannen zum Lösen}
\end{minipage}
\vspace{2mm}

{\raggedright\noindent\textcolor{ink}{\sffamily\textsc{Themen}}\enspace \themetag{Satzspiel} \themetag{Umwandlung} \themetag{Selbstblock} \themetag{ohne Schach} \themetag{Umnov}\par}

\clearpage
\phantomsection\label{hm:KPvkqrn}
\begin{tcolorbox}[enhanced,colback=ink,colframe=ink,arc=0pt,boxrule=0pt,left=4mm,right=4mm,top=2.5mm,bottom=2.5mm]
\color{white}\sffamily{\LARGE\bfseries KPvkqrn}\hfill{\Large h\#8.5}\\[0.5mm]
{\small 6 Steine\quad Weiß am Zug}\hfill{\small Nr.~373\qquad\href{https://helpman.komtera.lt/?fen=8/r7/K4q2/8/8/2n5/2P5/7k&moves=8.5}{\color{white}\bfseries lösen\,\ensuremath{\nearrow}}}
\end{tcolorbox}
\vspace{2mm}
\noindent\begin{minipage}[t]{36mm}\vspace{0pt}\sffamily\small\raggedright
\hmfigure{längste eindeutige}{{\Large h\#8.5}}\hmfigure{eindeutig bei dieser Länge}{{\Large 1} Stellung}\hmfigure{längstes Matt}{{\Large h\#9.5}}\hmfigure{in der Klasse}{45,4 Milliarden Stellungen,\\ 26,2 Milliarden mit Hilfsmatt}
\end{minipage}\hfill
\begin{minipage}[t]{80mm}\vspace{0pt}\begin{tikzpicture}\begin{axis}[width=72mm,height=42mm,ybar,bar width=1.6pt,xmin=-1,xmax=20,ymin=0,ymax=5.703,axis y line*=left,axis x line*=bottom,xlabel={\scriptsize Halbzüge bis Matt},ylabel={\scriptsize Stellungen (Mrd.)},tick label style={font=\tiny},label style={font=\scriptsize},enlarge x limits=false,ylabel shift=-3pt,xlabel shift=-3pt]\addplot[fill=ink!55,draw=none] coordinates {(0,2.3e-06) (1,0.0008939) (2,0.004818) (3,0.04694) (4,0.152) (5,0.8036) (6,1.362) (7,3.673) (8,2.522) (9,5.154) (10,2.553) (11,5.185) (12,2.268) (13,2.364) (14,0.07) (15,0.04814) (16,5.139e-05) (17,0.0005028) (18,1e-08) (19,9.41e-07)};\draw[densely dotted] (axis cs:17,0) -- (axis cs:17,5.703);\end{axis}\begin{axis}[width=72mm,height=42mm,axis y line*=right,axis x line=none,ymode=log,xmin=-1,xmax=20,ylabel={\scriptsize eindeutig (log)},tick label style={font=\tiny},label style={font=\scriptsize},ylabel shift=-3pt]\addplot[mark=*,mark size=0.7pt,draw=black,thin] coordinates {(0,2300) (1,292708) (2,1360372) (3,14431150) (4,33603282) (5,91683699) (6,49086026) (7,81724557) (8,20535398) (9,37278114) (10,9977943) (11,11970171) (12,466788) (13,290169) (14,510) (15,3295) (17,1)};\end{axis}\end{tikzpicture}\end{minipage}
\vspace{3mm}

\noindent\begin{minipage}[t]{60mm}\vspace{0pt}\sbox\hmboard{\chessboard[setfen=8/r7/K4q2/8/8/2n5/2P5/7k w - - 0 1]}\usebox\hmboard\\[1mm]
\makebox[\wd\hmboard]{\sffamily\hspace{3mm}\textbf{h\#8.5}\hfill(2+4)}\end{minipage}\hfill
\begin{minipage}[t]{52mm}\vspace{0pt}
\textcolor{ink}{\sffamily\textsc{Lösung}}\par\smallskip\small
\begin{tabular}{@{}r@{\enspace}l@{\enspace}l@{}}
\textbf{1.} & \ldots & \symking{}xa7\\
\textbf{2.} & \symknight{}d5 & c4\\
\textbf{3.} & \symking{}g2 & cxd5\\
\textbf{4.} & \symking{}f3 & d6\\
\textbf{5.} & \symking{}e4 & d7\\
\textbf{6.} & \symking{}d5 & d8=\symqueen{}+\\
\textbf{7.} & \symking{}c6 & \symqueen{}xf6+\\
\textbf{8.} & \symking{}c7 & \symking{}a6\\
\textbf{9.} & \symking{}b8 & \symqueen{}d8\#\\
\end{tabular}\par\vspace{3mm}
\qrcode[height=12mm]{https://helpman.komtera.lt/?fen=8/r7/K4q2/8/8/2n5/2P5/7k&moves=8.5}\\{\sffamily\tiny\color{ink}scannen zum Lösen}
\end{minipage}
\vspace{2mm}

{\raggedright\noindent\textcolor{ink}{\sffamily\textsc{Themen}}\enspace \themetag{reines Matt} \themetag{Mustermatt} \themetag{Idealmatt} \themetag{Spiegelmatt} \themetag{Umwandlung} \themetag{Excelsior} \themetag{Switchback}\par\smallskip{\footnotesize\textit{Hinweis:} das Diagramm hat keinen legalen letzten Zug, kann also in keiner Partie entstehen.}\par}

\clearpage
\phantomsection\label{hm:KPvkqrp}
\begin{tcolorbox}[enhanced,colback=ink,colframe=ink,arc=0pt,boxrule=0pt,left=4mm,right=4mm,top=2.5mm,bottom=2.5mm]
\color{white}\sffamily{\LARGE\bfseries KPvkqrp}\hfill{\Large h\#8.5}\\[0.5mm]
{\small 6 Steine\quad Weiß am Zug}\hfill{\small Nr.~374\qquad\href{https://helpman.komtera.lt/?fen=8/7k/8/8/8/p7/P2qr3/K7&moves=8.5}{\color{white}\bfseries lösen\,\ensuremath{\nearrow}}}
\end{tcolorbox}
\vspace{2mm}
\noindent\begin{minipage}[t]{36mm}\vspace{0pt}\sffamily\small\raggedright
\hmfigure{längste eindeutige}{{\Large h\#8.5}}\hmfigure{eindeutig bei dieser Länge}{{\Large 644} Stellungen}\hmfigure{längstes Matt}{{\Large h\#12.5}}\hmfigure{in der Klasse}{34,1 Milliarden Stellungen,\\ 20 Milliarden mit Hilfsmatt}
\end{minipage}\hfill
\begin{minipage}[t]{80mm}\vspace{0pt}\begin{tikzpicture}\begin{axis}[width=72mm,height=42mm,ybar,bar width=1.6pt,xmin=-1,xmax=26,ymin=0,ymax=4.065,axis y line*=left,axis x line*=bottom,xlabel={\scriptsize Halbzüge bis Matt},ylabel={\scriptsize Stellungen (Mrd.)},tick label style={font=\tiny},label style={font=\scriptsize},enlarge x limits=false,ylabel shift=-3pt,xlabel shift=-3pt]\addplot[fill=ink!55,draw=none] coordinates {(0,2.408e-06) (1,0.0009722) (2,0.005038) (3,0.04539) (4,0.148) (5,0.744) (6,1.205) (7,2.977) (8,1.927) (9,3.683) (10,1.918) (11,3.696) (12,1.716) (13,1.668) (14,0.1265) (15,0.1111) (16,0.009071) (17,0.009785) (18,9.55e-05) (19,0.0001115) (20,1e-08) (21,3.2e-08) (23,2e-09) (25,7e-09)};\draw[densely dotted] (axis cs:17,0) -- (axis cs:17,4.065);\end{axis}\begin{axis}[width=72mm,height=42mm,axis y line*=right,axis x line=none,ymode=log,xmin=-1,xmax=26,ylabel={\scriptsize eindeutig (log)},tick label style={font=\tiny},label style={font=\scriptsize},ylabel shift=-3pt]\addplot[mark=*,mark size=0.7pt,draw=black,thin] coordinates {(0,2408) (1,339673) (2,1552095) (3,14608152) (4,33942550) (5,84617072) (6,42047879) (7,63248673) (8,15163489) (9,26166981) (10,7641765) (11,8490726) (12,773970) (13,680258) (14,56611) (15,39876) (16,253) (17,644)};\end{axis}\end{tikzpicture}\end{minipage}
\vspace{3mm}

\noindent\begin{minipage}[t]{60mm}\vspace{0pt}\sbox\hmboard{\chessboard[setfen=8/7k/8/8/8/p7/P2qr3/K7 w - - 0 1]}\usebox\hmboard\\[1mm]
\makebox[\wd\hmboard]{\sffamily\hspace{3mm}\textbf{h\#8.5}\hfill(2+4)}\end{minipage}\hfill
\begin{minipage}[t]{52mm}\vspace{0pt}
\textcolor{ink}{\sffamily\textsc{Lösung}}\par\smallskip\small
\begin{tabular}{@{}r@{\enspace}l@{\enspace}l@{}}
\textbf{1.} & \ldots & \symking{}b1\\
\textbf{2.} & \symrook{}e3 & \symking{}a1\\
\textbf{3.} & \symrook{}b3 & axb3\\
\textbf{4.} & \symking{}g6 & b4\\
\textbf{5.} & \symking{}f5 & b5\\
\textbf{6.} & \symking{}e4 & b6\\
\textbf{7.} & \symking{}d3 & b7\\
\textbf{8.} & \symking{}c2 & b8=\symqueen{}\\
\textbf{9.} & \symking{}c1 & \symqueen{}b1\#\\
\end{tabular}\par\vspace{3mm}
\qrcode[height=12mm]{https://helpman.komtera.lt/?fen=8/7k/8/8/8/p7/P2qr3/K7&moves=8.5}\\{\sffamily\tiny\color{ink}scannen zum Lösen}
\end{minipage}
\vspace{2mm}

{\raggedright\noindent\textcolor{ink}{\sffamily\textsc{Themen}}\enspace \themetag{Umwandlung} \themetag{Excelsior} \themetag{Switchback} \themetag{ohne Schach}\par}

\clearpage
\phantomsection\label{hm:KPvkrnp}
\begin{tcolorbox}[enhanced,colback=ink,colframe=ink,arc=0pt,boxrule=0pt,left=4mm,right=4mm,top=2.5mm,bottom=2.5mm]
\color{white}\sffamily{\LARGE\bfseries KPvkrnp}\hfill{\Large h\#8.5}\\[0.5mm]
{\small 6 Steine\quad Weiß am Zug}\hfill{\small Nr.~375\qquad\href{https://helpman.komtera.lt/?fen=8/8/8/8/8/1p5n/1P5r/K3k3&moves=8.5}{\color{white}\bfseries lösen\,\ensuremath{\nearrow}}}
\end{tcolorbox}
\vspace{2mm}
\noindent\begin{minipage}[t]{36mm}\vspace{0pt}\sffamily\small\raggedright
\hmfigure{längste eindeutige}{{\Large h\#8.5}}\hmfigure{eindeutig bei dieser Länge}{{\Large 2\,170} Stellungen}\hmfigure{längstes Matt}{{\Large h\#12.5}}\hmfigure{in der Klasse}{34,1 Milliarden Stellungen,\\ 22,9 Milliarden mit Hilfsmatt}
\end{minipage}\hfill
\begin{minipage}[t]{80mm}\vspace{0pt}\begin{tikzpicture}\begin{axis}[width=72mm,height=42mm,ybar,bar width=1.6pt,xmin=-1,xmax=26,ymin=0,ymax=4.242,axis y line*=left,axis x line*=bottom,xlabel={\scriptsize Halbzüge bis Matt},ylabel={\scriptsize Stellungen (Mrd.)},tick label style={font=\tiny},label style={font=\scriptsize},enlarge x limits=false,ylabel shift=-3pt,xlabel shift=-3pt]\addplot[fill=ink!55,draw=none] coordinates {(0,8.002e-06) (1,0.002387) (2,0.008954) (3,0.06748) (4,0.1919) (5,0.8173) (6,1.367) (7,2.738) (8,2.443) (9,3.65) (10,2.601) (11,3.856) (12,2.698) (13,1.852) (14,0.3986) (15,0.1406) (16,0.02997) (17,0.01895) (18,0.0009727) (19,0.0003327) (20,1.947e-06) (21,1.217e-05) (22,5e-08) (23,2.77e-07) (25,7e-09)};\draw[densely dotted] (axis cs:17,0) -- (axis cs:17,4.242);\end{axis}\begin{axis}[width=72mm,height=42mm,axis y line*=right,axis x line=none,ymode=log,xmin=-1,xmax=26,ylabel={\scriptsize eindeutig (log)},tick label style={font=\tiny},label style={font=\scriptsize},ylabel shift=-3pt]\addplot[mark=*,mark size=0.7pt,draw=black,thin] coordinates {(0,8002) (1,680944) (2,2439292) (3,18666044) (4,41903055) (5,84172811) (6,46860660) (7,54530337) (8,24347088) (9,32873350) (10,14123601) (11,14521689) (12,3050185) (13,1108208) (14,186359) (15,80159) (16,4376) (17,2170)};\end{axis}\end{tikzpicture}\end{minipage}
\vspace{3mm}

\noindent\begin{minipage}[t]{60mm}\vspace{0pt}\sbox\hmboard{\chessboard[setfen=8/8/8/8/8/1p5n/1P5r/K3k3 w - - 0 1]}\usebox\hmboard\\[1mm]
\makebox[\wd\hmboard]{\sffamily\hspace{3mm}\textbf{h\#8.5}\hfill(2+4)}\end{minipage}\hfill
\begin{minipage}[t]{52mm}\vspace{0pt}
\textcolor{ink}{\sffamily\textsc{Lösung}}\par\smallskip\small
\begin{tabular}{@{}r@{\enspace}l@{\enspace}l@{}}
\textbf{1.} & \ldots & \symking{}b1\\
\textbf{2.} & \symrook{}c2 & \symking{}a1\\
\textbf{3.} & \symrook{}c3 & bxc3\\
\textbf{4.} & \symknight{}f4 & c4\\
\textbf{5.} & \symknight{}d5 & cxd5\\
\textbf{6.} & \symking{}d2 & d6\\
\textbf{7.} & \symking{}c3 & d7\\
\textbf{8.} & \symking{}b4 & d8=\symqueen{}\\
\textbf{9.} & \symking{}a3 & \symqueen{}a5\#\\
\end{tabular}\par\vspace{3mm}
\qrcode[height=12mm]{https://helpman.komtera.lt/?fen=8/8/8/8/8/1p5n/1P5r/K3k3&moves=8.5}\\{\sffamily\tiny\color{ink}scannen zum Lösen}
\end{minipage}
\vspace{2mm}

{\raggedright\noindent\textcolor{ink}{\sffamily\textsc{Themen}}\enspace \themetag{Satzspiel} \themetag{Umwandlung} \themetag{Excelsior} \themetag{Switchback} \themetag{ohne Schach}\par}

\clearpage
\phantomsection\label{hm:KPvkrrp}
\begin{tcolorbox}[enhanced,colback=ink,colframe=ink,arc=0pt,boxrule=0pt,left=4mm,right=4mm,top=2.5mm,bottom=2.5mm]
\color{white}\sffamily{\LARGE\bfseries KPvkrrp}\hfill{\Large h\#8.5}\\[0.5mm]
{\small 6 Steine\quad Weiß am Zug}\hfill{\small Nr.~376\qquad\href{https://helpman.komtera.lt/?fen=1r6/8/8/1r6/6k1/1p6/1P6/K7&moves=8.5}{\color{white}\bfseries lösen\,\ensuremath{\nearrow}}}
\end{tcolorbox}
\vspace{2mm}
\noindent\begin{minipage}[t]{36mm}\vspace{0pt}\sffamily\small\raggedright
\hmfigure{längste eindeutige}{{\Large h\#8.5}}\hmfigure{eindeutig bei dieser Länge}{{\Large 369} Stellungen}\hmfigure{längstes Matt}{{\Large h\#10.5}}\hmfigure{in der Klasse}{34,1 Milliarden Stellungen,\\ 10,7 Milliarden mit Hilfsmatt}
\end{minipage}\hfill
\begin{minipage}[t]{80mm}\vspace{0pt}\begin{tikzpicture}\begin{axis}[width=72mm,height=42mm,ybar,bar width=1.6pt,xmin=-1,xmax=22,ymin=0,ymax=2.004,axis y line*=left,axis x line*=bottom,xlabel={\scriptsize Halbzüge bis Matt},ylabel={\scriptsize Stellungen (Mrd.)},tick label style={font=\tiny},label style={font=\scriptsize},enlarge x limits=false,ylabel shift=-3pt,xlabel shift=-3pt]\addplot[fill=ink!55,draw=none] coordinates {(0,1.645e-06) (1,0.000821) (2,0.003969) (3,0.03236) (4,0.1125) (5,0.49) (6,0.8032) (7,1.621) (8,1.108) (9,1.822) (10,1.107) (11,1.807) (12,0.9717) (13,0.6898) (14,0.07178) (15,0.04733) (16,0.00608) (17,0.005315) (18,6.903e-05) (19,5.075e-05) (20,6e-09) (21,1.8e-08)};\draw[densely dotted] (axis cs:17,0) -- (axis cs:17,2.004);\end{axis}\begin{axis}[width=72mm,height=42mm,axis y line*=right,axis x line=none,ymode=log,xmin=-1,xmax=22,ylabel={\scriptsize eindeutig (log)},tick label style={font=\tiny},label style={font=\scriptsize},ylabel shift=-3pt]\addplot[mark=*,mark size=0.7pt,draw=black,thin] coordinates {(0,1645) (1,298097) (2,1468592) (3,11473361) (4,27321201) (5,52026114) (6,26040206) (7,31661183) (8,8883605) (9,13142973) (10,5323950) (11,3714360) (12,456917) (13,262450) (14,40492) (15,21533) (16,290) (17,369)};\end{axis}\end{tikzpicture}\end{minipage}
\vspace{3mm}

\noindent\begin{minipage}[t]{60mm}\vspace{0pt}\sbox\hmboard{\chessboard[setfen=1r6/8/8/1r6/6k1/1p6/1P6/K7 w - - 0 1]}\usebox\hmboard\\[1mm]
\makebox[\wd\hmboard]{\sffamily\hspace{3mm}\textbf{h\#8.5}\hfill(2+4)}\end{minipage}\hfill
\begin{minipage}[t]{52mm}\vspace{0pt}
\textcolor{ink}{\sffamily\textsc{Lösung}}\par\smallskip\small
\begin{tabular}{@{}r@{\enspace}l@{\enspace}l@{}}
\textbf{1.} & \ldots & \symking{}b1\\
\textbf{2.} & \symrook{}c5 & \symking{}a1\\
\textbf{3.} & \symrook{}c3 & bxc3\\
\textbf{4.} & \symking{}f5 & c4\\
\textbf{5.} & \symking{}e6 & c5\\
\textbf{6.} & \symking{}d7 & c6+\\
\textbf{7.} & \symking{}c8 & c7\\
\textbf{8.} & \symking{}b7 & c8=\symqueen{}+\\
\textbf{9.} & \symking{}a8 & \symqueen{}a6\#\\
\end{tabular}\par\vspace{3mm}
\qrcode[height=12mm]{https://helpman.komtera.lt/?fen=1r6/8/8/1r6/6k1/1p6/1P6/K7&moves=8.5}\\{\sffamily\tiny\color{ink}scannen zum Lösen}
\end{minipage}
\vspace{2mm}

{\raggedright\noindent\textcolor{ink}{\sffamily\textsc{Themen}}\enspace \themetag{Satzspiel} \themetag{reines Matt} \themetag{Mustermatt} \themetag{Umwandlung} \themetag{Excelsior} \themetag{Switchback} \themetag{Umnov}\par}

\clearpage
\phantomsection\label{hm:KBBBPvk}
\begin{tcolorbox}[enhanced,colback=ink,colframe=ink,arc=0pt,boxrule=0pt,left=4mm,right=4mm,top=2.5mm,bottom=2.5mm]
\color{white}\sffamily{\LARGE\bfseries KBBBPvk}\hfill{\Large h\#8}\\[0.5mm]
{\small 6 Steine\quad Schwarz am Zug}\hfill{\small Nr.~377\qquad\href{https://helpman.komtera.lt/?fen=8/8/7B/8/7B/8/7P/K5Bk&moves=8}{\color{white}\bfseries lösen\,\ensuremath{\nearrow}}}
\end{tcolorbox}
\vspace{2mm}
\noindent\begin{minipage}[t]{36mm}\vspace{0pt}\sffamily\small\raggedright
\hmfigure{längste eindeutige}{{\Large h\#8}}\hmfigure{eindeutig bei dieser Länge}{{\Large 11} Stellungen}\hmfigure{längstes Matt}{{\Large h\#10}}\hmfigure{in der Klasse}{45,4 Milliarden Stellungen,\\ 4,92 Milliarden mit Hilfsmatt}
\end{minipage}\hfill
\begin{minipage}[t]{80mm}\vspace{0pt}\begin{tikzpicture}\begin{axis}[width=72mm,height=42mm,ybar,bar width=1.6pt,xmin=-1,xmax=21,ymin=0,ymax=1.465,axis y line*=left,axis x line*=bottom,xlabel={\scriptsize Halbzüge bis Matt},ylabel={\scriptsize Stellungen (Mrd.)},tick label style={font=\tiny},label style={font=\scriptsize},enlarge x limits=false,ylabel shift=-3pt,xlabel shift=-3pt]\addplot[fill=ink!55,draw=none] coordinates {(0,0.004637) (1,0.01657) (2,0.045) (3,0.1842) (4,0.6864) (5,0.8512) (6,1.332) (7,0.4998) (8,0.5023) (9,0.1582) (10,0.1789) (11,0.1312) (12,0.1896) (13,0.06252) (14,0.0697) (15,0.003826) (16,0.003323) (17,4.552e-05) (18,3.234e-05) (19,3.9e-08) (20,8e-09)};\draw[densely dotted] (axis cs:16,0) -- (axis cs:16,1.465);\end{axis}\begin{axis}[width=72mm,height=42mm,axis y line*=right,axis x line=none,ymode=log,xmin=-1,xmax=21,ylabel={\scriptsize eindeutig (log)},tick label style={font=\tiny},label style={font=\scriptsize},ylabel shift=-3pt]\addplot[mark=*,mark size=0.7pt,draw=black,thin] coordinates {(0,4636904) (1,12500749) (2,33880469) (3,89973189) (4,189337444) (5,46831451) (6,37002061) (7,7969458) (8,5256272) (9,2724852) (10,2152356) (11,1752183) (12,1159864) (13,23780) (14,15690) (15,54) (16,11)};\end{axis}\end{tikzpicture}\end{minipage}
\vspace{3mm}

\noindent\begin{minipage}[t]{60mm}\vspace{0pt}\sbox\hmboard{\chessboard[setfen=8/8/7B/8/7B/8/7P/K5Bk b - - 0 1]}\usebox\hmboard\\[1mm]
\makebox[\wd\hmboard]{\sffamily\hspace{3mm}\textbf{h\#8}\hfill(5+1)}\end{minipage}\hfill
\begin{minipage}[t]{52mm}\vspace{0pt}
\textcolor{ink}{\sffamily\textsc{Lösung}}\par\smallskip\small
\begin{tabular}{@{}r@{\enspace}l@{\enspace}l@{}}
\textbf{1.} & \symking{}g2 & \symbishop{}d8\\
\textbf{2.} & \symking{}f3 & h4\\
\textbf{3.} & \symking{}e4 & h5\\
\textbf{4.} & \symking{}e5 & \symbishop{}f4+\\
\textbf{5.} & \symking{}d5 & h6\\
\textbf{6.} & \symking{}c6 & h7\\
\textbf{7.} & \symking{}b7 & h8=\symqueen{}\\
\textbf{8.} & \symking{}a8 & \symqueen{}h1\#\\
\end{tabular}\par\vspace{3mm}
\qrcode[height=12mm]{https://helpman.komtera.lt/?fen=8/8/7B/8/7B/8/7P/K5Bk&moves=8}\\{\sffamily\tiny\color{ink}scannen zum Lösen}
\end{minipage}
\vspace{2mm}

{\raggedright\noindent\textcolor{ink}{\sffamily\textsc{Themen}}\enspace \themetag{Satzspiel} \themetag{reines Matt} \themetag{Spiegelmatt} \themetag{Umwandlung} \themetag{Excelsior} \themetag{ein Stein zieht} \themetag{ohne Schlag}\par}

\clearpage
\phantomsection\label{hm:KBBPvkn}
\begin{tcolorbox}[enhanced,colback=ink,colframe=ink,arc=0pt,boxrule=0pt,left=4mm,right=4mm,top=2.5mm,bottom=2.5mm]
\color{white}\sffamily{\LARGE\bfseries KBBPvkn}\hfill{\Large h\#8}\\[0.5mm]
{\small 6 Steine\quad Schwarz am Zug}\hfill{\small Nr.~378\qquad\href{https://helpman.komtera.lt/?fen=5k2/5P2/8/8/1B6/n7/1B6/K7&moves=8}{\color{white}\bfseries lösen\,\ensuremath{\nearrow}}}
\end{tcolorbox}
\vspace{2mm}
\noindent\begin{minipage}[t]{36mm}\vspace{0pt}\sffamily\small\raggedright
\hmfigure{längste eindeutige}{{\Large h\#8}}\hmfigure{eindeutig bei dieser Länge}{{\Large 10} Stellungen}\hmfigure{längstes Matt}{{\Large h\#12}}\hmfigure{in der Klasse}{45,4 Milliarden Stellungen,\\ 15 Milliarden mit Hilfsmatt}
\end{minipage}\hfill
\begin{minipage}[t]{80mm}\vspace{0pt}\begin{tikzpicture}\begin{axis}[width=72mm,height=42mm,ybar,bar width=1.6pt,xmin=-1,xmax=25,ymin=0,ymax=3.153,axis y line*=left,axis x line*=bottom,xlabel={\scriptsize Halbzüge bis Matt},ylabel={\scriptsize Stellungen (Mrd.)},tick label style={font=\tiny},label style={font=\scriptsize},enlarge x limits=false,ylabel shift=-3pt,xlabel shift=-3pt]\addplot[fill=ink!55,draw=none] coordinates {(0,0.002061) (1,0.01221) (2,0.04341) (3,0.299) (4,0.8848) (5,1.861) (6,2.866) (7,2.026) (8,2.197) (9,1.3) (10,1.312) (11,0.966) (12,0.9439) (13,0.1961) (14,0.07659) (15,0.001289) (16,0.0003484) (17,2.22e-07) (18,0.0001807) (19,6e-09) (20,0.0001375) (21,4.6e-08) (22,3.037e-05) (24,2.108e-06)};\draw[densely dotted] (axis cs:16,0) -- (axis cs:16,3.153);\end{axis}\begin{axis}[width=72mm,height=42mm,axis y line*=right,axis x line=none,ymode=log,xmin=-1,xmax=25,ylabel={\scriptsize eindeutig (log)},tick label style={font=\tiny},label style={font=\scriptsize},ylabel shift=-3pt]\addplot[mark=*,mark size=0.7pt,draw=black,thin] coordinates {(0,2060584) (1,8521928) (2,22873454) (3,117208311) (4,175632721) (5,115234334) (6,85844651) (7,33714326) (8,21337974) (9,12428058) (10,7438591) (11,3694180) (12,1120983) (13,13089) (14,916) (15,2) (16,10)};\end{axis}\end{tikzpicture}\end{minipage}
\vspace{3mm}

\noindent\begin{minipage}[t]{60mm}\vspace{0pt}\sbox\hmboard{\chessboard[setfen=5k2/5P2/8/8/1B6/n7/1B6/K7 b - - 0 1]}\usebox\hmboard\\[1mm]
\makebox[\wd\hmboard]{\sffamily\hspace{3mm}\textbf{h\#8}\hfill(4+2)}\end{minipage}\hfill
\begin{minipage}[t]{52mm}\vspace{0pt}
\textcolor{ink}{\sffamily\textsc{Lösung}}\par\smallskip\small
\begin{tabular}{@{}r@{\enspace}l@{\enspace}l@{}}
\textbf{1.} & \symking{}xf7 & \symking{}a2\\
\textbf{2.} & \symking{}e6 & \symking{}b3\\
\textbf{3.} & \symking{}d5 & \symking{}a4\\
\textbf{4.} & \symking{}c4 & \symking{}a5\\
\textbf{5.} & \symking{}b3 & \symbishop{}c1\\
\textbf{6.} & \symking{}a2 & \symking{}a4\\
\textbf{7.} & \symking{}a1 & \symking{}b3\\
\textbf{8.} & \symknight{}b1 & \symbishop{}b2\#\\
\end{tabular}\par\vspace{3mm}
\qrcode[height=12mm]{https://helpman.komtera.lt/?fen=5k2/5P2/8/8/1B6/n7/1B6/K7&moves=8}\\{\sffamily\tiny\color{ink}scannen zum Lösen}
\end{minipage}
\vspace{2mm}

{\raggedright\noindent\textcolor{ink}{\sffamily\textsc{Themen}}\enspace \themetag{reines Matt} \themetag{Switchback} \themetag{Selbstblock} \themetag{ohne Schach}\par\smallskip{\footnotesize\textit{Hinweis:} die Partei am Zug steht im Diagramm im Schach; die Lösung beginnt mit einem Schlag.}\par}

\clearpage
\phantomsection\label{hm:KBBvkbb}
\begin{tcolorbox}[enhanced,colback=ink,colframe=ink,arc=0pt,boxrule=0pt,left=4mm,right=4mm,top=2.5mm,bottom=2.5mm]
\color{white}\sffamily{\LARGE\bfseries KBBvkbb}\hfill{\Large h\#8}\\[0.5mm]
{\small 6 Steine\quad Schwarz am Zug}\hfill{\small Nr.~379\qquad\href{https://helpman.komtera.lt/?fen=5b2/8/8/6B1/8/7b/7k/K3B3&moves=8}{\color{white}\bfseries lösen\,\ensuremath{\nearrow}}}
\end{tcolorbox}
\vspace{2mm}
\noindent\begin{minipage}[t]{36mm}\vspace{0pt}\sffamily\small\raggedright
\hmfigure{längste eindeutige}{{\Large h\#8}}\hmfigure{eindeutig bei dieser Länge}{{\Large 129} Stellungen}\hmfigure{längstes Matt}{{\Large h\#10}}\hmfigure{in der Klasse}{15,5 Milliarden Stellungen,\\ 2,01 Milliarden mit Hilfsmatt}
\end{minipage}\hfill
\begin{minipage}[t]{80mm}\vspace{0pt}\begin{tikzpicture}\begin{axis}[width=72mm,height=42mm,ybar,bar width=1.6pt,xmin=-1,xmax=21,ymin=0,ymax=368.5,axis y line*=left,axis x line*=bottom,xlabel={\scriptsize Halbzüge bis Matt},ylabel={\scriptsize Stellungen (Mio.)},tick label style={font=\tiny},label style={font=\scriptsize},enlarge x limits=false,ylabel shift=-3pt,xlabel shift=-3pt]\addplot[fill=ink!55,draw=none] coordinates {(0,0.2237) (1,0.9154) (2,5.721) (3,31.19) (4,115.3) (5,228.6) (6,335) (7,291.8) (8,214.9) (9,172.1) (10,135) (11,133.7) (12,119.1) (13,77.12) (14,53.37) (15,42.35) (16,20.79) (17,20.1) (18,5.665) (19,7.111) (20,0.01741)};\draw[densely dotted] (axis cs:16,0) -- (axis cs:16,368.5);\end{axis}\begin{axis}[width=72mm,height=42mm,axis y line*=right,axis x line=none,ymode=log,xmin=-1,xmax=21,ylabel={\scriptsize eindeutig (log)},tick label style={font=\tiny},label style={font=\scriptsize},ylabel shift=-3pt]\addplot[mark=*,mark size=0.7pt,draw=black,thin] coordinates {(0,223680) (1,881558) (2,4253535) (3,15397836) (4,24189108) (5,12568753) (6,5121042) (7,2046425) (8,834101) (9,612151) (10,371517) (11,210734) (12,93793) (13,24161) (14,9640) (15,1134) (16,129)};\end{axis}\end{tikzpicture}\end{minipage}
\vspace{3mm}

\noindent\begin{minipage}[t]{60mm}\vspace{0pt}\sbox\hmboard{\chessboard[setfen=5b2/8/8/6B1/8/7b/7k/K3B3 b - - 0 1]}\usebox\hmboard\\[1mm]
\makebox[\wd\hmboard]{\sffamily\hspace{3mm}\textbf{h\#8}\hfill(3+3)}\end{minipage}\hfill
\begin{minipage}[t]{52mm}\vspace{0pt}
\textcolor{ink}{\sffamily\textsc{Lösung}}\par\smallskip\small
\begin{tabular}{@{}r@{\enspace}l@{\enspace}l@{}}
\textbf{1.} & \symbishop{}f5 & \symking{}b2\\
\textbf{2.} & \symking{}h3 & \symking{}c3\\
\textbf{3.} & \symking{}g4 & \symking{}d4\\
\textbf{4.} & \symking{}xg5 & \symking{}e5\\
\textbf{5.} & \symking{}h6 & \symking{}f6\\
\textbf{6.} & \symking{}h7 & \symking{}f7\\
\textbf{7.} & \symking{}h8 & \symking{}xf8\\
\textbf{8.} & \symbishop{}h7 & \symbishop{}c3\#\\
\end{tabular}\par\vspace{3mm}
\qrcode[height=12mm]{https://helpman.komtera.lt/?fen=5b2/8/8/6B1/8/7b/7k/K3B3&moves=8}\\{\sffamily\tiny\color{ink}scannen zum Lösen}
\end{minipage}
\vspace{2mm}

{\raggedright\noindent\textcolor{ink}{\sffamily\textsc{Themen}}\enspace \themetag{Selbstblock} \themetag{ohne Schach}\par}

\clearpage
\phantomsection\label{hm:KBNPvkq}
\begin{tcolorbox}[enhanced,colback=ink,colframe=ink,arc=0pt,boxrule=0pt,left=4mm,right=4mm,top=2.5mm,bottom=2.5mm]
\color{white}\sffamily{\LARGE\bfseries KBNPvkq}\hfill{\Large h\#8}\\[0.5mm]
{\small 6 Steine\quad Schwarz am Zug}\hfill{\small Nr.~380\qquad\href{https://helpman.komtera.lt/?fen=8/8/8/8/8/7N/qB5P/k1K5&moves=8}{\color{white}\bfseries lösen\,\ensuremath{\nearrow}}}
\end{tcolorbox}
\vspace{2mm}
\noindent\begin{minipage}[t]{36mm}\vspace{0pt}\sffamily\small\raggedright
\hmfigure{längste eindeutige}{{\Large h\#8}}\hmfigure{eindeutig bei dieser Länge}{{\Large 1} Stellung}\hmfigure{längstes Matt}{{\Large h\#9}}\hmfigure{in der Klasse}{45,4 Milliarden Stellungen,\\ 26,4 Milliarden mit Hilfsmatt}
\end{minipage}\hfill
\begin{minipage}[t]{80mm}\vspace{0pt}\begin{tikzpicture}\begin{axis}[width=72mm,height=42mm,ybar,bar width=1.6pt,xmin=-1,xmax=19,ymin=0,ymax=5.727,axis y line*=left,axis x line*=bottom,xlabel={\scriptsize Halbzüge bis Matt},ylabel={\scriptsize Stellungen (Mrd.)},tick label style={font=\tiny},label style={font=\scriptsize},enlarge x limits=false,ylabel shift=-3pt,xlabel shift=-3pt]\addplot[fill=ink!55,draw=none] coordinates {(0,0.0005098) (1,0.005536) (2,0.02782) (3,0.2893) (4,0.9507) (5,2.784) (6,3.678) (7,5.206) (8,4.859) (9,4.57) (10,2.479) (11,1.227) (12,0.2781) (13,0.05363) (14,0.001802) (15,0.0005007) (16,3.012e-06) (17,4.07e-07) (18,8e-09)};\draw[densely dotted] (axis cs:16,0) -- (axis cs:16,5.727);\end{axis}\begin{axis}[width=72mm,height=42mm,axis y line*=right,axis x line=none,ymode=log,xmin=-1,xmax=19,ylabel={\scriptsize eindeutig (log)},tick label style={font=\tiny},label style={font=\scriptsize},ylabel shift=-3pt]\addplot[mark=*,mark size=0.7pt,draw=black,thin] coordinates {(0,509774) (1,3240392) (2,12460226) (3,109701690) (4,192876185) (5,232635452) (6,138921992) (7,117667350) (8,38794777) (9,17507355) (10,2791188) (11,685090) (12,28376) (13,6347) (14,29) (16,1)};\end{axis}\end{tikzpicture}\end{minipage}
\vspace{3mm}

\noindent\begin{minipage}[t]{60mm}\vspace{0pt}\sbox\hmboard{\chessboard[setfen=8/8/8/8/8/7N/qB5P/k1K5 b - - 0 1]}\usebox\hmboard\\[1mm]
\makebox[\wd\hmboard]{\sffamily\hspace{3mm}\textbf{h\#8}\hfill(4+2)}\end{minipage}\hfill
\begin{minipage}[t]{52mm}\vspace{0pt}
\textcolor{ink}{\sffamily\textsc{Lösung}}\par\smallskip\small
\begin{tabular}{@{}r@{\enspace}l@{\enspace}l@{}}
\textbf{1.} & \symqueen{}xb2+ & \symking{}d1\\
\textbf{2.} & \symqueen{}g7 & \symknight{}f4\\
\textbf{3.} & \symking{}b2 & h4\\
\textbf{4.} & \symking{}c3 & h5\\
\textbf{5.} & \symking{}d4 & h6\\
\textbf{6.} & \symking{}e3 & hxg7\\
\textbf{7.} & \symking{}f2 & g8=\symqueen{}\\
\textbf{8.} & \symking{}f1 & \symqueen{}g2\#\\
\end{tabular}\par\vspace{3mm}
\qrcode[height=12mm]{https://helpman.komtera.lt/?fen=8/8/8/8/8/7N/qB5P/k1K5&moves=8}\\{\sffamily\tiny\color{ink}scannen zum Lösen}
\end{minipage}
\vspace{2mm}

{\raggedright\noindent\textcolor{ink}{\sffamily\textsc{Themen}}\enspace \themetag{Umwandlung} \themetag{Excelsior}\par\smallskip{\footnotesize\textit{Hinweis:} die Partei am Zug steht im Diagramm im Schach; die Lösung beginnt mit einem Schlag.}\par}

\clearpage
\phantomsection\label{hm:KBNvkbp}
\begin{tcolorbox}[enhanced,colback=ink,colframe=ink,arc=0pt,boxrule=0pt,left=4mm,right=4mm,top=2.5mm,bottom=2.5mm]
\color{white}\sffamily{\LARGE\bfseries KBNvkbp}\hfill{\Large h\#8}\\[0.5mm]
{\small 6 Steine\quad Schwarz am Zug}\hfill{\small Nr.~381\qquad\href{https://helpman.komtera.lt/?fen=5kNb/6p1/6B1/8/8/8/8/K7&moves=8}{\color{white}\bfseries lösen\,\ensuremath{\nearrow}}}
\end{tcolorbox}
\vspace{2mm}
\noindent\begin{minipage}[t]{36mm}\vspace{0pt}\sffamily\small\raggedright
\hmfigure{längste eindeutige}{{\Large h\#8}}\hmfigure{eindeutig bei dieser Länge}{{\Large 4} Stellungen}\hmfigure{längstes Matt}{{\Large h\#16}}\hmfigure{in der Klasse}{45,4 Milliarden Stellungen,\\ 30 Milliarden mit Hilfsmatt}
\end{minipage}\hfill
\begin{minipage}[t]{80mm}\vspace{0pt}\begin{tikzpicture}\begin{axis}[width=72mm,height=42mm,ybar,bar width=1.6pt,xmin=-1,xmax=33,ymin=0,ymax=7.99,axis y line*=left,axis x line*=bottom,xlabel={\scriptsize Halbzüge bis Matt},ylabel={\scriptsize Stellungen (Mrd.)},tick label style={font=\tiny},label style={font=\scriptsize},enlarge x limits=false,ylabel shift=-3pt,xlabel shift=-3pt]\addplot[fill=ink!55,draw=none] coordinates {(0,0.0008407) (1,0.003315) (2,0.01635) (3,0.08977) (4,0.3729) (5,1.316) (6,3.096) (7,5.328) (8,7.264) (9,5.969) (10,4.046) (11,1.732) (12,0.5687) (13,0.1268) (14,0.01917) (15,0.00194) (16,0.0001626) (17,2.1e-08) (18,3.409e-05) (20,1.278e-05) (22,1.438e-05) (23,2e-09) (24,1.307e-05) (25,2e-09) (26,9.474e-06) (27,2e-09) (28,5.22e-06) (30,4.03e-06) (32,2.222e-06)};\draw[densely dotted] (axis cs:16,0) -- (axis cs:16,7.99);\end{axis}\begin{axis}[width=72mm,height=42mm,axis y line*=right,axis x line=none,ymode=log,xmin=-1,xmax=33,ylabel={\scriptsize eindeutig (log)},tick label style={font=\tiny},label style={font=\scriptsize},ylabel shift=-3pt]\addplot[mark=*,mark size=0.7pt,draw=black,thin] coordinates {(0,840711) (1,3163852) (2,13145113) (3,42686980) (4,113234549) (5,162750584) (6,151956106) (7,77337674) (8,43014117) (9,10073412) (10,2600234) (11,405017) (12,118058) (13,8198) (14,98) (16,4)};\end{axis}\end{tikzpicture}\end{minipage}
\vspace{3mm}

\noindent\begin{minipage}[t]{60mm}\vspace{0pt}\sbox\hmboard{\chessboard[setfen=5kNb/6p1/6B1/8/8/8/8/K7 b - - 0 1]}\usebox\hmboard\\[1mm]
\makebox[\wd\hmboard]{\sffamily\hspace{3mm}\textbf{h\#8}\hfill(3+3)}\end{minipage}\hfill
\begin{minipage}[t]{52mm}\vspace{0pt}
\textcolor{ink}{\sffamily\textsc{Lösung}}\par\smallskip\small
\begin{tabular}{@{}r@{\enspace}l@{\enspace}l@{}}
\textbf{1.} & \symking{}xg8 & \symbishop{}e8\\
\textbf{2.} & g5+ & \symking{}b1\\
\textbf{3.} & g4 & \symking{}c2\\
\textbf{4.} & g3 & \symking{}d3\\
\textbf{5.} & g2 & \symking{}e4\\
\textbf{6.} & g1=\symbishop{} & \symking{}f5\\
\textbf{7.} & \symbishop{}c5 & \symking{}g6\\
\textbf{8.} & \symbishop{}f8 & \symbishop{}f7\#\\
\end{tabular}\par\vspace{3mm}
\qrcode[height=12mm]{https://helpman.komtera.lt/?fen=5kNb/6p1/6B1/8/8/8/8/K7&moves=8}\\{\sffamily\tiny\color{ink}scannen zum Lösen}
\end{minipage}
\vspace{2mm}

{\raggedright\noindent\textcolor{ink}{\sffamily\textsc{Themen}}\enspace \themetag{reines Matt} \themetag{Mustermatt} \themetag{Idealmatt} \themetag{Umwandlung} \themetag{Unterverwandlung} \themetag{Excelsior} \themetag{Selbstblock} \themetag{Kniest}\par\smallskip{\footnotesize\textit{Hinweis:} die Lösung beginnt mit einem Schlag.}\par}

\clearpage
\phantomsection\label{hm:KBNvkqb}
\begin{tcolorbox}[enhanced,colback=ink,colframe=ink,arc=0pt,boxrule=0pt,left=4mm,right=4mm,top=2.5mm,bottom=2.5mm]
\color{white}\sffamily{\LARGE\bfseries KBNvkqb}\hfill{\Large h\#8}\\[0.5mm]
{\small 6 Steine\quad Schwarz am Zug}\hfill{\small Nr.~382\qquad\href{https://helpman.komtera.lt/?fen=8/8/8/8/8/4q3/2K1kB2/5bN1&moves=8}{\color{white}\bfseries lösen\,\ensuremath{\nearrow}}}
\end{tcolorbox}
\vspace{2mm}
\noindent\begin{minipage}[t]{36mm}\vspace{0pt}\sffamily\small\raggedright
\hmfigure{längste eindeutige}{{\Large h\#8}}\hmfigure{eindeutig bei dieser Länge}{{\Large 1} Stellung}\hmfigure{längstes Matt}{{\Large h\#11}}\hmfigure{in der Klasse}{15,5 Milliarden Stellungen,\\ 8,36 Milliarden mit Hilfsmatt}
\end{minipage}\hfill
\begin{minipage}[t]{80mm}\vspace{0pt}\begin{tikzpicture}\begin{axis}[width=72mm,height=42mm,ybar,bar width=1.6pt,xmin=-1,xmax=23,ymin=0,ymax=2.381,axis y line*=left,axis x line*=bottom,xlabel={\scriptsize Halbzüge bis Matt},ylabel={\scriptsize Stellungen (Mrd.)},tick label style={font=\tiny},label style={font=\scriptsize},enlarge x limits=false,ylabel shift=-3pt,xlabel shift=-3pt]\addplot[fill=ink!55,draw=none] coordinates {(0,7.453e-05) (1,0.0003101) (2,0.001939) (3,0.01439) (4,0.06173) (5,0.2901) (6,0.712) (7,1.757) (8,1.851) (9,2.165) (10,0.8002) (11,0.5623) (12,0.1051) (13,0.03902) (14,0.001549) (15,0.0005863) (16,4.934e-05) (17,1.122e-06) (18,4.239e-05) (19,1.47e-07) (20,1.359e-05) (21,9e-09) (22,4.05e-07)};\draw[densely dotted] (axis cs:16,0) -- (axis cs:16,2.381);\end{axis}\begin{axis}[width=72mm,height=42mm,axis y line*=right,axis x line=none,ymode=log,xmin=-1,xmax=23,ylabel={\scriptsize eindeutig (log)},tick label style={font=\tiny},label style={font=\scriptsize},ylabel shift=-3pt]\addplot[mark=*,mark size=0.7pt,draw=black,thin] coordinates {(0,74526) (1,298834) (2,1284057) (3,6477715) (4,17551791) (5,41317807) (6,36953095) (7,38342937) (8,8999777) (9,5116123) (10,636012) (11,222607) (12,6497) (13,1699) (14,21) (15,4) (16,1)};\end{axis}\end{tikzpicture}\end{minipage}
\vspace{3mm}

\noindent\begin{minipage}[t]{60mm}\vspace{0pt}\sbox\hmboard{\chessboard[setfen=8/8/8/8/8/4q3/2K1kB2/5bN1 b - - 0 1]}\usebox\hmboard\\[1mm]
\makebox[\wd\hmboard]{\sffamily\hspace{3mm}\textbf{h\#8}\hfill(3+3)}\end{minipage}\hfill
\begin{minipage}[t]{52mm}\vspace{0pt}
\textcolor{ink}{\sffamily\textsc{Lösung}}\par\smallskip\small
\begin{tabular}{@{}r@{\enspace}l@{\enspace}l@{}}
\textbf{1.} & \symking{}xf2 & \symknight{}f3\\
\textbf{2.} & \symbishop{}c4 & \symknight{}d2\\
\textbf{3.} & \symking{}e2 & \symking{}c1\\
\textbf{4.} & \symking{}d3 & \symking{}d1\\
\textbf{5.} & \symking{}c3 & \symknight{}f1\\
\textbf{6.} & \symking{}b2 & \symknight{}xe3\\
\textbf{7.} & \symking{}a1 & \symking{}c1\\
\textbf{8.} & \symbishop{}a2 & \symknight{}c2\#\\
\end{tabular}\par\vspace{3mm}
\qrcode[height=12mm]{https://helpman.komtera.lt/?fen=8/8/8/8/8/4q3/2K1kB2/5bN1&moves=8}\\{\sffamily\tiny\color{ink}scannen zum Lösen}
\end{minipage}
\vspace{2mm}

{\raggedright\noindent\textcolor{ink}{\sffamily\textsc{Themen}}\enspace \themetag{reines Matt} \themetag{Mustermatt} \themetag{Idealmatt} \themetag{Switchback} \themetag{Selbstblock} \themetag{ohne Schach} \themetag{Klasinc}\par\smallskip{\footnotesize\textit{Hinweis:} die Partei am Zug steht im Diagramm im Schach; die Lösung beginnt mit einem Schlag.}\par}

\clearpage
\phantomsection\label{hm:KBNvkrp}
\begin{tcolorbox}[enhanced,colback=ink,colframe=ink,arc=0pt,boxrule=0pt,left=4mm,right=4mm,top=2.5mm,bottom=2.5mm]
\color{white}\sffamily{\LARGE\bfseries KBNvkrp}\hfill{\Large h\#8}\\[0.5mm]
{\small 6 Steine\quad Schwarz am Zug}\hfill{\small Nr.~383\qquad\href{https://helpman.komtera.lt/?fen=8/7p/7r/8/8/8/NB6/k1K5&moves=8}{\color{white}\bfseries lösen\,\ensuremath{\nearrow}}}
\end{tcolorbox}
\vspace{2mm}
\noindent\begin{minipage}[t]{36mm}\vspace{0pt}\sffamily\small\raggedright
\hmfigure{längste eindeutige}{{\Large h\#8}}\hmfigure{eindeutig bei dieser Länge}{{\Large 7} Stellungen}\hmfigure{längstes Matt}{{\Large h\#16.5}}\hmfigure{in der Klasse}{45,4 Milliarden Stellungen,\\ 28,7 Milliarden mit Hilfsmatt}
\end{minipage}\hfill
\begin{minipage}[t]{80mm}\vspace{0pt}\begin{tikzpicture}\begin{axis}[width=72mm,height=42mm,ybar,bar width=1.6pt,xmin=-1,xmax=34,ymin=0,ymax=7.236,axis y line*=left,axis x line*=bottom,xlabel={\scriptsize Halbzüge bis Matt},ylabel={\scriptsize Stellungen (Mrd.)},tick label style={font=\tiny},label style={font=\scriptsize},enlarge x limits=false,ylabel shift=-3pt,xlabel shift=-3pt]\addplot[fill=ink!55,draw=none] coordinates {(0,0.0004215) (1,0.001776) (2,0.007415) (3,0.05396) (4,0.2162) (5,0.9357) (6,2.317) (7,4.728) (8,6.578) (9,6.359) (10,4.327) (11,2.301) (12,0.7077) (13,0.1784) (14,0.01899) (15,0.001662) (16,7.297e-05) (17,6.5e-08) (18,2.298e-05) (19,9.2e-08) (20,1.803e-05) (21,8e-08) (22,1.371e-05) (23,7e-08) (24,6.385e-06) (25,1.1e-08) (26,3.889e-06) (27,3e-08) (28,3.903e-06) (29,2.2e-08) (30,1.39e-06) (32,2.336e-06) (33,2e-09)};\draw[densely dotted] (axis cs:16,0) -- (axis cs:16,7.236);\end{axis}\begin{axis}[width=72mm,height=42mm,axis y line*=right,axis x line=none,ymode=log,xmin=-1,xmax=34,ylabel={\scriptsize eindeutig (log)},tick label style={font=\tiny},label style={font=\scriptsize},ylabel shift=-3pt]\addplot[mark=*,mark size=0.7pt,draw=black,thin] coordinates {(0,421453) (1,1675475) (2,4882578) (3,24541809) (4,74865881) (5,132627309) (6,150069665) (7,98449184) (8,49767973) (9,16650832) (10,4110054) (11,943299) (12,106634) (13,7653) (14,95) (16,7)};\end{axis}\end{tikzpicture}\end{minipage}
\vspace{3mm}

\noindent\begin{minipage}[t]{60mm}\vspace{0pt}\sbox\hmboard{\chessboard[setfen=8/7p/7r/8/8/8/NB6/k1K5 b - - 0 1]}\usebox\hmboard\\[1mm]
\makebox[\wd\hmboard]{\sffamily\hspace{3mm}\textbf{h\#8}\hfill(3+3)}\end{minipage}\hfill
\begin{minipage}[t]{52mm}\vspace{0pt}
\textcolor{ink}{\sffamily\textsc{Lösung}}\par\smallskip\small
\begin{tabular}{@{}r@{\enspace}l@{\enspace}l@{}}
\textbf{1.} & \symking{}xa2 & \symking{}d2\\
\textbf{2.} & \symking{}b3 & \symking{}e3\\
\textbf{3.} & \symking{}c4 & \symking{}f4\\
\textbf{4.} & \symking{}d5 & \symking{}g5\\
\textbf{5.} & \symking{}e6 & \symking{}xh6\\
\textbf{6.} & \symking{}f7 & \symking{}g5\\
\textbf{7.} & \symking{}g8 & \symking{}f6\\
\textbf{8.} & \symking{}h8 & \symking{}f7\#\\
\end{tabular}\par\vspace{3mm}
\qrcode[height=12mm]{https://helpman.komtera.lt/?fen=8/7p/7r/8/8/8/NB6/k1K5&moves=8}\\{\sffamily\tiny\color{ink}scannen zum Lösen}
\end{minipage}
\vspace{2mm}

{\raggedright\noindent\textcolor{ink}{\sffamily\textsc{Themen}}\enspace \themetag{Switchback} \themetag{ein Stein zieht} \themetag{ohne Schach}\par\smallskip{\footnotesize\textit{Hinweis:} die Partei am Zug steht im Diagramm im Schach; die Lösung beginnt mit einem Schlag.}\par}

\clearpage
\phantomsection\label{hm:KBPvknn}
\begin{tcolorbox}[enhanced,colback=ink,colframe=ink,arc=0pt,boxrule=0pt,left=4mm,right=4mm,top=2.5mm,bottom=2.5mm]
\color{white}\sffamily{\LARGE\bfseries KBPvknn}\hfill{\Large h\#8}\\[0.5mm]
{\small 6 Steine\quad Schwarz am Zug}\hfill{\small Nr.~384\qquad\href{https://helpman.komtera.lt/?fen=8/8/8/8/2K5/kB6/1P6/5n1n&moves=8}{\color{white}\bfseries lösen\,\ensuremath{\nearrow}}}
\end{tcolorbox}
\vspace{2mm}
\noindent\begin{minipage}[t]{36mm}\vspace{0pt}\sffamily\small\raggedright
\hmfigure{längste eindeutige}{{\Large h\#8}}\hmfigure{eindeutig bei dieser Länge}{{\Large 6} Stellungen}\hmfigure{längstes Matt}{{\Large h\#13}}\hmfigure{in der Klasse}{45,4 Milliarden Stellungen,\\ 15,3 Milliarden mit Hilfsmatt}
\end{minipage}\hfill
\begin{minipage}[t]{80mm}\vspace{0pt}\begin{tikzpicture}\begin{axis}[width=72mm,height=42mm,ybar,bar width=1.6pt,xmin=-1,xmax=27,ymin=0,ymax=2.356,axis y line*=left,axis x line*=bottom,xlabel={\scriptsize Halbzüge bis Matt},ylabel={\scriptsize Stellungen (Mrd.)},tick label style={font=\tiny},label style={font=\scriptsize},enlarge x limits=false,ylabel shift=-3pt,xlabel shift=-3pt]\addplot[fill=ink!55,draw=none] coordinates {(0,0.0001294) (1,0.003048) (2,0.009581) (3,0.09108) (4,0.2529) (5,1.072) (6,1.482) (7,1.763) (8,1.788) (9,1.933) (10,1.988) (11,2.142) (12,1.898) (13,0.6811) (14,0.2227) (15,0.006537) (16,0.0002374) (17,2.5e-08) (18,4.718e-05) (20,2.922e-05) (22,2.04e-06) (24,1.56e-07) (26,2e-09)};\draw[densely dotted] (axis cs:16,0) -- (axis cs:16,2.356);\end{axis}\begin{axis}[width=72mm,height=42mm,axis y line*=right,axis x line=none,ymode=log,xmin=-1,xmax=27,ylabel={\scriptsize eindeutig (log)},tick label style={font=\tiny},label style={font=\scriptsize},ylabel shift=-3pt]\addplot[mark=*,mark size=0.7pt,draw=black,thin] coordinates {(0,129430) (1,1161340) (2,3057637) (3,31047360) (4,51908333) (5,79857321) (6,47774068) (7,41523697) (8,28555326) (9,24736551) (10,14418345) (11,9932306) (12,3199402) (13,112551) (14,1815) (16,6)};\end{axis}\end{tikzpicture}\end{minipage}
\vspace{3mm}

\noindent\begin{minipage}[t]{60mm}\vspace{0pt}\sbox\hmboard{\chessboard[setfen=8/8/8/8/2K5/kB6/1P6/5n1n b - - 0 1]}\usebox\hmboard\\[1mm]
\makebox[\wd\hmboard]{\sffamily\hspace{3mm}\textbf{h\#8}\hfill(3+3)}\end{minipage}\hfill
\begin{minipage}[t]{52mm}\vspace{0pt}
\textcolor{ink}{\sffamily\textsc{Lösung}}\par\smallskip\small
\begin{tabular}{@{}r@{\enspace}l@{\enspace}l@{}}
\textbf{1.} & \symking{}xb2 & \symking{}d3\\
\textbf{2.} & \symking{}c1 & \symbishop{}d5\\
\textbf{3.} & \symking{}d1 & \symbishop{}xh1\\
\textbf{4.} & \symking{}e1 & \symbishop{}g2\\
\textbf{5.} & \symking{}f2 & \symbishop{}h3\\
\textbf{6.} & \symking{}g1 & \symking{}e2\\
\textbf{7.} & \symking{}h1 & \symking{}f2\\
\textbf{8.} & \symknight{}h2 & \symbishop{}g2\#\\
\end{tabular}\par\vspace{3mm}
\qrcode[height=12mm]{https://helpman.komtera.lt/?fen=8/8/8/8/2K5/kB6/1P6/5n1n&moves=8}\\{\sffamily\tiny\color{ink}scannen zum Lösen}
\end{minipage}
\vspace{2mm}

{\raggedright\noindent\textcolor{ink}{\sffamily\textsc{Themen}}\enspace \themetag{reines Matt} \themetag{Mustermatt} \themetag{Idealmatt} \themetag{Switchback} \themetag{Selbstblock} \themetag{Kniest} \themetag{ohne Schach}\par\smallskip{\footnotesize\textit{Hinweis:} das Diagramm hat keinen legalen letzten Zug, kann also in keiner Partie entstehen.}\par}

\clearpage
\phantomsection\label{hm:KNNPvkp}
\begin{tcolorbox}[enhanced,colback=ink,colframe=ink,arc=0pt,boxrule=0pt,left=4mm,right=4mm,top=2.5mm,bottom=2.5mm]
\color{white}\sffamily{\LARGE\bfseries KNNPvkp}\hfill{\Large h\#8}\\[0.5mm]
{\small 6 Steine\quad Schwarz am Zug}\hfill{\small Nr.~385\qquad\href{https://helpman.komtera.lt/?fen=8/8/8/8/8/1p3N2/1P2N3/1K1k4&moves=8}{\color{white}\bfseries lösen\,\ensuremath{\nearrow}}}
\end{tcolorbox}
\vspace{2mm}
\noindent\begin{minipage}[t]{36mm}\vspace{0pt}\sffamily\small\raggedright
\hmfigure{längste eindeutige}{{\Large h\#8}}\hmfigure{eindeutig bei dieser Länge}{{\Large 6} Stellungen}\hmfigure{längstes Matt}{{\Large h\#11}}\hmfigure{in der Klasse}{34,1 Milliarden Stellungen,\\ 12,1 Milliarden mit Hilfsmatt}
\end{minipage}\hfill
\begin{minipage}[t]{80mm}\vspace{0pt}\begin{tikzpicture}\begin{axis}[width=72mm,height=42mm,ybar,bar width=1.6pt,xmin=-1,xmax=23,ymin=0,ymax=2.858,axis y line*=left,axis x line*=bottom,xlabel={\scriptsize Halbzüge bis Matt},ylabel={\scriptsize Stellungen (Mrd.)},tick label style={font=\tiny},label style={font=\scriptsize},enlarge x limits=false,ylabel shift=-3pt,xlabel shift=-3pt]\addplot[fill=ink!55,draw=none] coordinates {(0,0.0002863) (1,0.005284) (2,0.0145) (3,0.127) (4,0.3565) (5,1.068) (6,1.691) (7,1.929) (8,2.598) (9,1.835) (10,1.654) (11,0.5541) (12,0.26) (13,0.03473) (14,0.006211) (15,0.0002527) (16,1.304e-05) (17,2.27e-07) (18,4.563e-06) (19,2e-08) (20,6.74e-07) (22,7e-09)};\draw[densely dotted] (axis cs:16,0) -- (axis cs:16,2.858);\end{axis}\begin{axis}[width=72mm,height=42mm,axis y line*=right,axis x line=none,ymode=log,xmin=-1,xmax=23,ylabel={\scriptsize eindeutig (log)},tick label style={font=\tiny},label style={font=\scriptsize},ylabel shift=-3pt]\addplot[mark=*,mark size=0.7pt,draw=black,thin] coordinates {(0,286254) (1,2121875) (2,6290604) (3,59432765) (4,84580993) (5,82290805) (6,73384347) (7,39298642) (8,28413089) (9,6488580) (10,2513762) (11,510029) (12,83584) (13,1578) (14,96) (16,6)};\end{axis}\end{tikzpicture}\end{minipage}
\vspace{3mm}

\noindent\begin{minipage}[t]{60mm}\vspace{0pt}\sbox\hmboard{\chessboard[setfen=8/8/8/8/8/1p3N2/1P2N3/1K1k4 b - - 0 1]}\usebox\hmboard\\[1mm]
\makebox[\wd\hmboard]{\sffamily\hspace{3mm}\textbf{h\#8}\hfill(4+2)}\end{minipage}\hfill
\begin{minipage}[t]{52mm}\vspace{0pt}
\textcolor{ink}{\sffamily\textsc{Lösung}}\par\smallskip\small
\begin{tabular}{@{}r@{\enspace}l@{\enspace}l@{}}
\textbf{1.} & \symking{}xe2 & \symking{}c1\\
\textbf{2.} & \symking{}e3 & \symking{}d1\\
\textbf{3.} & \symking{}d3 & \symking{}e1\\
\textbf{4.} & \symking{}c2 & \symknight{}d4+\\
\textbf{5.} & \symking{}xb2 & \symking{}d2\\
\textbf{6.} & \symking{}a1 & \symking{}c3\\
\textbf{7.} & b2 & \symking{}b3\\
\textbf{8.} & b1=\symknight{} & \symknight{}c2\#\\
\end{tabular}\par\vspace{3mm}
\qrcode[height=12mm]{https://helpman.komtera.lt/?fen=8/8/8/8/8/1p3N2/1P2N3/1K1k4&moves=8}\\{\sffamily\tiny\color{ink}scannen zum Lösen}
\end{minipage}
\vspace{2mm}

{\raggedright\noindent\textcolor{ink}{\sffamily\textsc{Themen}}\enspace \themetag{reines Matt} \themetag{Mustermatt} \themetag{Idealmatt} \themetag{Umwandlung} \themetag{Unterverwandlung} \themetag{Selbstblock} \themetag{Umnov}\par\smallskip{\footnotesize\textit{Hinweis:} die Lösung beginnt mit einem Schlag.}\par}

\clearpage
\phantomsection\label{hm:KNNvkpp}
\begin{tcolorbox}[enhanced,colback=ink,colframe=ink,arc=0pt,boxrule=0pt,left=4mm,right=4mm,top=2.5mm,bottom=2.5mm]
\color{white}\sffamily{\LARGE\bfseries KNNvkpp}\hfill{\Large h\#8}\\[0.5mm]
{\small 6 Steine\quad Schwarz am Zug}\hfill{\small Nr.~386\qquad\href{https://helpman.komtera.lt/?fen=7k/6p1/5Np1/6N1/8/8/8/K7&moves=8}{\color{white}\bfseries lösen\,\ensuremath{\nearrow}}}
\end{tcolorbox}
\vspace{2mm}
\noindent\begin{minipage}[t]{36mm}\vspace{0pt}\sffamily\small\raggedright
\hmfigure{längste eindeutige}{{\Large h\#8}}\hmfigure{eindeutig bei dieser Länge}{{\Large 1} Stellung}\hmfigure{längstes Matt}{{\Large h\#11}}\hmfigure{in der Klasse}{34,1 Milliarden Stellungen,\\ 6,05 Milliarden mit Hilfsmatt}
\end{minipage}\hfill
\begin{minipage}[t]{80mm}\vspace{0pt}\begin{tikzpicture}\begin{axis}[width=72mm,height=42mm,ybar,bar width=1.6pt,xmin=-1,xmax=23,ymin=0,ymax=1.427,axis y line*=left,axis x line*=bottom,xlabel={\scriptsize Halbzüge bis Matt},ylabel={\scriptsize Stellungen (Mrd.)},tick label style={font=\tiny},label style={font=\scriptsize},enlarge x limits=false,ylabel shift=-3pt,xlabel shift=-3pt]\addplot[fill=ink!55,draw=none] coordinates {(0,9.181e-05) (1,0.0004461) (2,0.001384) (3,0.01039) (4,0.0339) (5,0.1467) (6,0.3605) (7,0.7556) (8,1.297) (9,1.283) (10,1.232) (11,0.5929) (12,0.2635) (13,0.06329) (14,0.009393) (15,0.001453) (16,4.651e-06) (17,2.1e-08) (18,1.944e-06) (19,2e-09) (20,8.94e-07) (21,3e-09) (22,4.21e-07)};\draw[densely dotted] (axis cs:16,0) -- (axis cs:16,1.427);\end{axis}\begin{axis}[width=72mm,height=42mm,axis y line*=right,axis x line=none,ymode=log,xmin=-1,xmax=23,ylabel={\scriptsize eindeutig (log)},tick label style={font=\tiny},label style={font=\scriptsize},ylabel shift=-3pt]\addplot[mark=*,mark size=0.7pt,draw=black,thin] coordinates {(0,91811) (1,418556) (2,1021385) (3,5548976) (4,12001304) (5,15501902) (6,21294211) (7,12629054) (8,12273599) (9,2606237) (10,1124923) (11,165455) (12,36476) (13,6473) (14,12) (16,1)};\end{axis}\end{tikzpicture}\end{minipage}
\vspace{3mm}

\noindent\begin{minipage}[t]{60mm}\vspace{0pt}\sbox\hmboard{\chessboard[setfen=7k/6p1/5Np1/6N1/8/8/8/K7 b - - 0 1]}\usebox\hmboard\\[1mm]
\makebox[\wd\hmboard]{\sffamily\hspace{3mm}\textbf{h\#8}\hfill(3+3)}\end{minipage}\hfill
\begin{minipage}[t]{52mm}\vspace{0pt}
\textcolor{ink}{\sffamily\textsc{Lösung}}\par\smallskip\small
\begin{tabular}{@{}r@{\enspace}l@{\enspace}l@{}}
\textbf{1.} & gxf6 & \symking{}b2\\
\textbf{2.} & f5 & \symking{}c3\\
\textbf{3.} & f4 & \symking{}d4\\
\textbf{4.} & f3 & \symking{}e5\\
\textbf{5.} & f2 & \symking{}f6\\
\textbf{6.} & f1=\symrook{}+ & \symking{}xg6\\
\textbf{7.} & \symrook{}f8 & \symking{}h6\\
\textbf{8.} & \symrook{}g8 & \symknight{}f7\#\\
\end{tabular}\par\vspace{3mm}
\qrcode[height=12mm]{https://helpman.komtera.lt/?fen=7k/6p1/5Np1/6N1/8/8/8/K7&moves=8}\\{\sffamily\tiny\color{ink}scannen zum Lösen}
\end{minipage}
\vspace{2mm}

{\raggedright\noindent\textcolor{ink}{\sffamily\textsc{Themen}}\enspace \themetag{reines Matt} \themetag{Mustermatt} \themetag{Idealmatt} \themetag{Umwandlung} \themetag{Unterverwandlung} \themetag{Excelsior} \themetag{Selbstblock} \themetag{ein Stein zieht}\par\smallskip{\footnotesize\textit{Hinweis:} die Lösung beginnt mit einem Schlag.}\par}

\clearpage
\phantomsection\label{hm:KNNvkqp}
\begin{tcolorbox}[enhanced,colback=ink,colframe=ink,arc=0pt,boxrule=0pt,left=4mm,right=4mm,top=2.5mm,bottom=2.5mm]
\color{white}\sffamily{\LARGE\bfseries KNNvkqp}\hfill{\Large h\#8}\\[0.5mm]
{\small 6 Steine\quad Schwarz am Zug}\hfill{\small Nr.~387\qquad\href{https://helpman.komtera.lt/?fen=4qk2/4Np2/4N3/8/8/8/8/2K5&moves=8}{\color{white}\bfseries lösen\,\ensuremath{\nearrow}}}
\end{tcolorbox}
\vspace{2mm}
\noindent\begin{minipage}[t]{36mm}\vspace{0pt}\sffamily\small\raggedright
\hmfigure{längste eindeutige}{{\Large h\#8}}\hmfigure{eindeutig bei dieser Länge}{{\Large 1} Stellung}\hmfigure{längstes Matt}{{\Large h\#12}}\hmfigure{in der Klasse}{45,4 Milliarden Stellungen,\\ 13,6 Milliarden mit Hilfsmatt}
\end{minipage}\hfill
\begin{minipage}[t]{80mm}\vspace{0pt}\begin{tikzpicture}\begin{axis}[width=72mm,height=42mm,ybar,bar width=1.6pt,xmin=-1,xmax=25,ymin=0,ymax=3.906,axis y line*=left,axis x line*=bottom,xlabel={\scriptsize Halbzüge bis Matt},ylabel={\scriptsize Stellungen (Mrd.)},tick label style={font=\tiny},label style={font=\scriptsize},enlarge x limits=false,ylabel shift=-3pt,xlabel shift=-3pt]\addplot[fill=ink!55,draw=none] coordinates {(0,0.0001083) (1,0.0005414) (2,0.00204) (3,0.01726) (4,0.05931) (5,0.3022) (6,0.6921) (7,1.884) (8,2.551) (9,3.551) (10,2.227) (11,1.69) (12,0.4162) (13,0.1805) (14,0.01386) (15,0.004408) (16,2.698e-05) (17,1.209e-06) (18,1.205e-05) (19,2.02e-07) (20,7.15e-07) (21,7.1e-08) (22,1.23e-07) (23,3e-08) (24,3e-09)};\draw[densely dotted] (axis cs:16,0) -- (axis cs:16,3.906);\end{axis}\begin{axis}[width=72mm,height=42mm,axis y line*=right,axis x line=none,ymode=log,xmin=-1,xmax=25,ylabel={\scriptsize eindeutig (log)},tick label style={font=\tiny},label style={font=\scriptsize},ylabel shift=-3pt]\addplot[mark=*,mark size=0.7pt,draw=black,thin] coordinates {(0,108269) (1,522375) (2,1288777) (3,7710750) (4,16976557) (5,35805799) (6,36626802) (7,43031037) (8,22912579) (9,12900137) (10,2219217) (11,726127) (12,39149) (13,13403) (14,37) (15,1) (16,1)};\end{axis}\end{tikzpicture}\end{minipage}
\vspace{3mm}

\noindent\begin{minipage}[t]{60mm}\vspace{0pt}\sbox\hmboard{\chessboard[setfen=4qk2/4Np2/4N3/8/8/8/8/2K5 b - - 0 1]}\usebox\hmboard\\[1mm]
\makebox[\wd\hmboard]{\sffamily\hspace{3mm}\textbf{h\#8}\hfill(3+3)}\end{minipage}\hfill
\begin{minipage}[t]{52mm}\vspace{0pt}
\textcolor{ink}{\sffamily\textsc{Lösung}}\par\smallskip\small
\begin{tabular}{@{}r@{\enspace}l@{\enspace}l@{}}
\textbf{1.} & fxe6 & \symking{}d2\\
\textbf{2.} & e5 & \symking{}e3\\
\textbf{3.} & e4 & \symking{}f4\\
\textbf{4.} & e3 & \symking{}g5\\
\textbf{5.} & e2 & \symking{}h6\\
\textbf{6.} & e1=\symrook{} & \symking{}h7\\
\textbf{7.} & \symrook{}f1 & \symking{}h8\\
\textbf{8.} & \symrook{}f7 & \symknight{}g6\#\\
\end{tabular}\par\vspace{3mm}
\qrcode[height=12mm]{https://helpman.komtera.lt/?fen=4qk2/4Np2/4N3/8/8/8/8/2K5&moves=8}\\{\sffamily\tiny\color{ink}scannen zum Lösen}
\end{minipage}
\vspace{2mm}

{\raggedright\noindent\textcolor{ink}{\sffamily\textsc{Themen}}\enspace \themetag{reines Matt} \themetag{Mustermatt} \themetag{Idealmatt} \themetag{Umwandlung} \themetag{Unterverwandlung} \themetag{Excelsior} \themetag{Rundlauf} \themetag{Selbstblock} \themetag{ein Stein zieht} \themetag{ohne Schach} \themetag{Umnov}\par\smallskip{\footnotesize\textit{Hinweis:} die Partei am Zug steht im Diagramm im Schach; die Lösung beginnt mit einem Schlag.}\par}

\clearpage
\phantomsection\label{hm:KNPvkqb}
\begin{tcolorbox}[enhanced,colback=ink,colframe=ink,arc=0pt,boxrule=0pt,left=4mm,right=4mm,top=2.5mm,bottom=2.5mm]
\color{white}\sffamily{\LARGE\bfseries KNPvkqb}\hfill{\Large h\#8}\\[0.5mm]
{\small 6 Steine\quad Schwarz am Zug}\hfill{\small Nr.~388\qquad\href{https://helpman.komtera.lt/?fen=8/8/1K6/8/1kq5/bP6/N7/8&moves=8}{\color{white}\bfseries lösen\,\ensuremath{\nearrow}}}
\end{tcolorbox}
\vspace{2mm}
\noindent\begin{minipage}[t]{36mm}\vspace{0pt}\sffamily\small\raggedright
\hmfigure{längste eindeutige}{{\Large h\#8}}\hmfigure{eindeutig bei dieser Länge}{{\Large 2} Stellungen}\hmfigure{längstes Matt}{{\Large h\#11.5}}\hmfigure{in der Klasse}{45,4 Milliarden Stellungen,\\ 26,8 Milliarden mit Hilfsmatt}
\end{minipage}\hfill
\begin{minipage}[t]{80mm}\vspace{0pt}\begin{tikzpicture}\begin{axis}[width=72mm,height=42mm,ybar,bar width=1.6pt,xmin=-1,xmax=24,ymin=0,ymax=4.875,axis y line*=left,axis x line*=bottom,xlabel={\scriptsize Halbzüge bis Matt},ylabel={\scriptsize Stellungen (Mrd.)},tick label style={font=\tiny},label style={font=\scriptsize},enlarge x limits=false,ylabel shift=-3pt,xlabel shift=-3pt]\addplot[fill=ink!55,draw=none] coordinates {(0,5.29e-05) (1,0.001793) (2,0.00888) (3,0.1137) (4,0.3739) (5,1.675) (6,2.101) (7,3.883) (8,2.697) (9,4.25) (10,2.854) (11,4.432) (12,2.428) (13,1.728) (14,0.1914) (15,0.04566) (16,0.0002677) (17,2.354e-05) (18,6.048e-05) (19,2.11e-07) (20,9.038e-06) (21,2.7e-08) (22,2.62e-07) (23,5e-09)};\draw[densely dotted] (axis cs:16,0) -- (axis cs:16,4.875);\end{axis}\begin{axis}[width=72mm,height=42mm,axis y line*=right,axis x line=none,ymode=log,xmin=-1,xmax=24,ylabel={\scriptsize eindeutig (log)},tick label style={font=\tiny},label style={font=\scriptsize},ylabel shift=-3pt]\addplot[mark=*,mark size=0.7pt,draw=black,thin] coordinates {(0,52904) (1,741665) (2,3133901) (3,44098928) (4,81809075) (5,147597440) (6,62442409) (7,79312883) (8,28637781) (9,38431156) (10,13497054) (11,14464993) (12,2066794) (13,677917) (14,4006) (15,332) (16,2)};\end{axis}\end{tikzpicture}\end{minipage}
\vspace{3mm}

\noindent\begin{minipage}[t]{60mm}\vspace{0pt}\sbox\hmboard{\chessboard[setfen=8/8/1K6/8/1kq5/bP6/N7/8 b - - 0 1]}\usebox\hmboard\\[1mm]
\makebox[\wd\hmboard]{\sffamily\hspace{3mm}\textbf{h\#8}\hfill(3+3)}\end{minipage}\hfill
\begin{minipage}[t]{52mm}\vspace{0pt}
\textcolor{ink}{\sffamily\textsc{Lösung}}\par\smallskip\small
\begin{tabular}{@{}r@{\enspace}l@{\enspace}l@{}}
\textbf{1.} & \symking{}xb3 & \symknight{}c3\\
\textbf{2.} & \symbishop{}d6 & \symknight{}b5\\
\textbf{3.} & \symking{}b4 & \symking{}a6\\
\textbf{4.} & \symking{}c5 & \symking{}a5\\
\textbf{5.} & \symking{}c6 & \symknight{}a3\\
\textbf{6.} & \symking{}b7 & \symknight{}xc4\\
\textbf{7.} & \symking{}a8 & \symking{}a6\\
\textbf{8.} & \symbishop{}b8 & \symknight{}b6\#\\
\end{tabular}\par\vspace{3mm}
\qrcode[height=12mm]{https://helpman.komtera.lt/?fen=8/8/1K6/8/1kq5/bP6/N7/8&moves=8}\\{\sffamily\tiny\color{ink}scannen zum Lösen}
\end{minipage}
\vspace{2mm}

{\raggedright\noindent\textcolor{ink}{\sffamily\textsc{Themen}}\enspace \themetag{reines Matt} \themetag{Mustermatt} \themetag{Idealmatt} \themetag{Switchback} \themetag{Selbstblock} \themetag{ohne Schach} \themetag{Klasinc}\par\smallskip{\footnotesize\textit{Hinweis:} die Partei am Zug steht im Diagramm im Schach; die Lösung beginnt mit einem Schlag.}\par}

\clearpage
\phantomsection\label{hm:KNvkqrr}
\begin{tcolorbox}[enhanced,colback=ink,colframe=ink,arc=0pt,boxrule=0pt,left=4mm,right=4mm,top=2.5mm,bottom=2.5mm]
\color{white}\sffamily{\LARGE\bfseries KNvkqrr}\hfill{\Large h\#8}\\[0.5mm]
{\small 6 Steine\quad Schwarz am Zug}\hfill{\small Nr.~389\qquad\href{https://helpman.komtera.lt/?fen=7k/8/6N1/1rr5/8/1q6/8/K7&moves=8}{\color{white}\bfseries lösen\,\ensuremath{\nearrow}}}
\end{tcolorbox}
\vspace{2mm}
\noindent\begin{minipage}[t]{36mm}\vspace{0pt}\sffamily\small\raggedright
\hmfigure{längste eindeutige}{{\Large h\#8}}\hmfigure{eindeutig bei dieser Länge}{{\Large 1} Stellung}\hmfigure{längstes Matt}{{\Large h\#9.5}}\hmfigure{in der Klasse}{15,5 Milliarden Stellungen,\\ 3,89 Milliarden mit Hilfsmatt}
\end{minipage}\hfill
\begin{minipage}[t]{80mm}\vspace{0pt}\begin{tikzpicture}\begin{axis}[width=72mm,height=42mm,ybar,bar width=1.6pt,xmin=-1,xmax=20,ymin=0,ymax=1.166,axis y line*=left,axis x line*=bottom,xlabel={\scriptsize Halbzüge bis Matt},ylabel={\scriptsize Stellungen (Mrd.)},tick label style={font=\tiny},label style={font=\scriptsize},enlarge x limits=false,ylabel shift=-3pt,xlabel shift=-3pt]\addplot[fill=ink!55,draw=none] coordinates {(0,1.498e-06) (1,6.333e-06) (2,8.302e-05) (3,0.0005724) (4,0.003161) (5,0.01689) (6,0.04994) (7,0.2127) (8,0.2994) (9,0.8658) (10,0.5028) (11,1.06) (12,0.3197) (13,0.454) (14,0.06055) (15,0.04517) (16,0.001839) (17,0.0005723) (18,3.6e-08) (19,3.7e-08)};\draw[densely dotted] (axis cs:16,0) -- (axis cs:16,1.166);\end{axis}\begin{axis}[width=72mm,height=42mm,axis y line*=right,axis x line=none,ymode=log,xmin=-1,xmax=20,ylabel={\scriptsize eindeutig (log)},tick label style={font=\tiny},label style={font=\scriptsize},ylabel shift=-3pt]\addplot[mark=*,mark size=0.7pt,draw=black,thin] coordinates {(0,1498) (1,6333) (2,69228) (3,392467) (4,878358) (5,3610679) (6,4522051) (7,13699623) (8,6560711) (9,13313037) (10,2668253) (11,4233577) (12,333139) (13,351107) (14,7420) (15,3422) (16,1)};\end{axis}\end{tikzpicture}\end{minipage}
\vspace{3mm}

\noindent\begin{minipage}[t]{60mm}\vspace{0pt}\sbox\hmboard{\chessboard[setfen=7k/8/6N1/1rr5/8/1q6/8/K7 b - - 0 1]}\usebox\hmboard\\[1mm]
\makebox[\wd\hmboard]{\sffamily\hspace{3mm}\textbf{h\#8}\hfill(2+4)}\end{minipage}\hfill
\begin{minipage}[t]{52mm}\vspace{0pt}
\textcolor{ink}{\sffamily\textsc{Lösung}}\par\smallskip\small
\begin{tabular}{@{}r@{\enspace}l@{\enspace}l@{}}
\textbf{1.} & \symking{}g7 & \symknight{}e7\\
\textbf{2.} & \symqueen{}d1+ & \symking{}a2\\
\textbf{3.} & \symrook{}b2+ & \symking{}xb2\\
\textbf{4.} & \symrook{}h5 & \symking{}c3\\
\textbf{5.} & \symqueen{}h1 & \symking{}d4\\
\textbf{6.} & \symrook{}h7 & \symking{}e5\\
\textbf{7.} & \symking{}h6 & \symking{}f6\\
\textbf{8.} & \symqueen{}h5 & \symknight{}g8\#\\
\end{tabular}\par\vspace{3mm}
\qrcode[height=12mm]{https://helpman.komtera.lt/?fen=7k/8/6N1/1rr5/8/1q6/8/K7&moves=8}\\{\sffamily\tiny\color{ink}scannen zum Lösen}
\end{minipage}
\vspace{2mm}

{\raggedright\noindent\textcolor{ink}{\sffamily\textsc{Themen}}\enspace \themetag{reines Matt} \themetag{Mustermatt} \themetag{Idealmatt} \themetag{Selbstblock}\par\smallskip{\footnotesize\textit{Hinweis:} die Partei am Zug steht im Diagramm im Schach.}\par}

\clearpage
\phantomsection\label{hm:KPPvkrp}
\begin{tcolorbox}[enhanced,colback=ink,colframe=ink,arc=0pt,boxrule=0pt,left=4mm,right=4mm,top=2.5mm,bottom=2.5mm]
\color{white}\sffamily{\LARGE\bfseries KPPvkrp}\hfill{\Large h\#8}\\[0.5mm]
{\small 6 Steine\quad Schwarz am Zug}\hfill{\small Nr.~390\qquad\href{https://helpman.komtera.lt/?fen=8/8/8/8/5pk1/5P2/2r2P2/K7&moves=8}{\color{white}\bfseries lösen\,\ensuremath{\nearrow}}}
\end{tcolorbox}
\vspace{2mm}
\noindent\begin{minipage}[t]{36mm}\vspace{0pt}\sffamily\small\raggedright
\hmfigure{längste eindeutige}{{\Large h\#8}}\hmfigure{eindeutig bei dieser Länge}{{\Large 5} Stellungen}\hmfigure{längstes Matt}{{\Large h\#14.5}}\hmfigure{in der Klasse}{25,6 Milliarden Stellungen,\\ 8,68 Milliarden mit Hilfsmatt}
\end{minipage}\hfill
\begin{minipage}[t]{80mm}\vspace{0pt}\begin{tikzpicture}\begin{axis}[width=72mm,height=42mm,ybar,bar width=1.6pt,xmin=-1,xmax=30,ymin=0,ymax=1.895,axis y line*=left,axis x line*=bottom,xlabel={\scriptsize Halbzüge bis Matt},ylabel={\scriptsize Stellungen (Mrd.)},tick label style={font=\tiny},label style={font=\scriptsize},enlarge x limits=false,ylabel shift=-3pt,xlabel shift=-3pt]\addplot[fill=ink!55,draw=none] coordinates {(0,2.68e-05) (1,0.002703) (2,0.01011) (3,0.0754) (4,0.2281) (5,0.812) (6,1.202) (7,1.722) (8,1.256) (9,1.283) (10,0.8024) (11,0.7083) (12,0.4115) (13,0.1371) (14,0.02348) (15,0.002332) (16,0.000184) (17,4.179e-05) (18,1.568e-06) (19,4.269e-06) (20,8.5e-08) (21,5.302e-06) (22,7.8e-08) (23,2.855e-06) (24,4.1e-08) (25,8.76e-07) (26,1e-08) (27,4.9e-07) (28,3e-09) (29,1.32e-07)};\draw[densely dotted] (axis cs:16,0) -- (axis cs:16,1.895);\end{axis}\begin{axis}[width=72mm,height=42mm,axis y line*=right,axis x line=none,ymode=log,xmin=-1,xmax=30,ylabel={\scriptsize eindeutig (log)},tick label style={font=\tiny},label style={font=\scriptsize},ylabel shift=-3pt]\addplot[mark=*,mark size=0.7pt,draw=black,thin] coordinates {(0,26800) (1,826488) (2,3174447) (3,25263735) (4,52245129) (5,75057892) (6,36126025) (7,29933079) (8,10291363) (9,8569865) (10,3140968) (11,1565625) (12,323833) (13,31844) (14,1998) (15,366) (16,5)};\end{axis}\end{tikzpicture}\end{minipage}
\vspace{3mm}

\noindent\begin{minipage}[t]{60mm}\vspace{0pt}\sbox\hmboard{\chessboard[setfen=8/8/8/8/5pk1/5P2/2r2P2/K7 b - - 0 1]}\usebox\hmboard\\[1mm]
\makebox[\wd\hmboard]{\sffamily\hspace{3mm}\textbf{h\#8}\hfill(3+3)}\end{minipage}\hfill
\begin{minipage}[t]{52mm}\vspace{0pt}
\textcolor{ink}{\sffamily\textsc{Lösung}}\par\smallskip\small
\begin{tabular}{@{}r@{\enspace}l@{\enspace}l@{}}
\textbf{1.} & \symking{}f5 & \symking{}b1\\
\textbf{2.} & \symrook{}e2 & \symking{}c1\\
\textbf{3.} & \symrook{}e4 & fxe4+\\
\textbf{4.} & \symking{}e6 & e5\\
\textbf{5.} & \symking{}d5 & e6\\
\textbf{6.} & \symking{}c4 & e7\\
\textbf{7.} & \symking{}b3 & e8=\symqueen{}\\
\textbf{8.} & \symking{}a2 & \symqueen{}a4\#\\
\end{tabular}\par\vspace{3mm}
\qrcode[height=12mm]{https://helpman.komtera.lt/?fen=8/8/8/8/5pk1/5P2/2r2P2/K7&moves=8}\\{\sffamily\tiny\color{ink}scannen zum Lösen}
\end{minipage}
\vspace{2mm}

{\raggedright\noindent\textcolor{ink}{\sffamily\textsc{Themen}}\enspace \themetag{reines Matt} \themetag{Mustermatt} \themetag{Spiegelmatt} \themetag{Umwandlung} \themetag{Umnov}\par\smallskip{\footnotesize\textit{Hinweis:} die Partei am Zug steht im Diagramm im Schach.}\par}

\clearpage
\phantomsection\label{hm:KPvkqbb}
\begin{tcolorbox}[enhanced,colback=ink,colframe=ink,arc=0pt,boxrule=0pt,left=4mm,right=4mm,top=2.5mm,bottom=2.5mm]
\color{white}\sffamily{\LARGE\bfseries KPvkqbb}\hfill{\Large h\#8}\\[0.5mm]
{\small 6 Steine\quad Schwarz am Zug}\hfill{\small Nr.~391\qquad\href{https://helpman.komtera.lt/?fen=K7/8/8/8/6b1/1q2b2k/6P1/8&moves=8}{\color{white}\bfseries lösen\,\ensuremath{\nearrow}}}
\end{tcolorbox}
\vspace{2mm}
\noindent\begin{minipage}[t]{36mm}\vspace{0pt}\sffamily\small\raggedright
\hmfigure{längste eindeutige}{{\Large h\#8}}\hmfigure{eindeutig bei dieser Länge}{{\Large 3} Stellungen}\hmfigure{längstes Matt}{{\Large h\#11.5}}\hmfigure{in der Klasse}{45,4 Milliarden Stellungen,\\ 13,3 Milliarden mit Hilfsmatt}
\end{minipage}\hfill
\begin{minipage}[t]{80mm}\vspace{0pt}\begin{tikzpicture}\begin{axis}[width=72mm,height=42mm,ybar,bar width=1.6pt,xmin=-1,xmax=24,ymin=0,ymax=2.9,axis y line*=left,axis x line*=bottom,xlabel={\scriptsize Halbzüge bis Matt},ylabel={\scriptsize Stellungen (Mrd.)},tick label style={font=\tiny},label style={font=\scriptsize},enlarge x limits=false,ylabel shift=-3pt,xlabel shift=-3pt]\addplot[fill=ink!55,draw=none] coordinates {(0,5.822e-06) (1,0.000423) (2,0.001921) (3,0.02101) (4,0.05895) (5,0.3429) (6,0.4844) (7,1.451) (8,1.034) (9,2.266) (10,1.286) (11,2.636) (12,1.403) (13,1.831) (14,0.3159) (15,0.1561) (16,0.003124) (17,0.0009403) (18,2.83e-07) (19,1.95e-07) (20,4.6e-08) (21,4e-08) (23,4e-09)};\draw[densely dotted] (axis cs:16,0) -- (axis cs:16,2.9);\end{axis}\begin{axis}[width=72mm,height=42mm,axis y line*=right,axis x line=none,ymode=log,xmin=-1,xmax=24,ylabel={\scriptsize eindeutig (log)},tick label style={font=\tiny},label style={font=\scriptsize},ylabel shift=-3pt]\addplot[mark=*,mark size=0.7pt,draw=black,thin] coordinates {(0,5822) (1,133199) (2,386833) (3,6347481) (4,10620581) (5,37742633) (6,15271213) (7,33107278) (8,9828258) (9,19448651) (10,5587102) (11,10657199) (12,2016638) (13,1696531) (14,26703) (15,8361) (16,3)};\end{axis}\end{tikzpicture}\end{minipage}
\vspace{3mm}

\noindent\begin{minipage}[t]{60mm}\vspace{0pt}\sbox\hmboard{\chessboard[setfen=K7/8/8/8/6b1/1q2b2k/6P1/8 b - - 0 1]}\usebox\hmboard\\[1mm]
\makebox[\wd\hmboard]{\sffamily\hspace{3mm}\textbf{h\#8}\hfill(2+4)}\end{minipage}\hfill
\begin{minipage}[t]{52mm}\vspace{0pt}
\textcolor{ink}{\sffamily\textsc{Lösung}}\par\smallskip\small
\begin{tabular}{@{}r@{\enspace}l@{\enspace}l@{}}
\textbf{1.} & \symking{}h2 & g3\\
\textbf{2.} & \symbishop{}d7 & g4\\
\textbf{3.} & \symking{}g3 & g5\\
\textbf{4.} & \symking{}f4 & g6\\
\textbf{5.} & \symking{}e5 & g7\\
\textbf{6.} & \symking{}d6 & g8=\symqueen{}\\
\textbf{7.} & \symking{}c7 & \symqueen{}xb3\\
\textbf{8.} & \symking{}c8 & \symqueen{}b8\#\\
\end{tabular}\par\vspace{3mm}
\qrcode[height=12mm]{https://helpman.komtera.lt/?fen=K7/8/8/8/6b1/1q2b2k/6P1/8&moves=8}\\{\sffamily\tiny\color{ink}scannen zum Lösen}
\end{minipage}
\vspace{2mm}

{\raggedright\noindent\textcolor{ink}{\sffamily\textsc{Themen}}\enspace \themetag{Umwandlung} \themetag{Excelsior} \themetag{Selbstblock} \themetag{ein Stein zieht} \themetag{ohne Schach} \themetag{Umnov}\par\smallskip{\footnotesize\textit{Hinweis:} das Diagramm hat keinen legalen letzten Zug, kann also in keiner Partie entstehen.}\par}

\clearpage
\phantomsection\label{hm:KPvkqqb}
\begin{tcolorbox}[enhanced,colback=ink,colframe=ink,arc=0pt,boxrule=0pt,left=4mm,right=4mm,top=2.5mm,bottom=2.5mm]
\color{white}\sffamily{\LARGE\bfseries KPvkqqb}\hfill{\Large h\#8}\\[0.5mm]
{\small 6 Steine\quad Schwarz am Zug}\hfill{\small Nr.~392\qquad\href{https://helpman.komtera.lt/?fen=K7/8/8/8/6b1/1q5k/6P1/6q1&moves=8}{\color{white}\bfseries lösen\,\ensuremath{\nearrow}}}
\end{tcolorbox}
\vspace{2mm}
\noindent\begin{minipage}[t]{36mm}\vspace{0pt}\sffamily\small\raggedright
\hmfigure{längste eindeutige}{{\Large h\#8}}\hmfigure{eindeutig bei dieser Länge}{{\Large 15} Stellungen}\hmfigure{längstes Matt}{{\Large h\#11.5}}\hmfigure{in der Klasse}{45,4 Milliarden Stellungen,\\ 11,8 Milliarden mit Hilfsmatt}
\end{minipage}\hfill
\begin{minipage}[t]{80mm}\vspace{0pt}\begin{tikzpicture}\begin{axis}[width=72mm,height=42mm,ybar,bar width=1.6pt,xmin=-1,xmax=24,ymin=0,ymax=2.711,axis y line*=left,axis x line*=bottom,xlabel={\scriptsize Halbzüge bis Matt},ylabel={\scriptsize Stellungen (Mrd.)},tick label style={font=\tiny},label style={font=\scriptsize},enlarge x limits=false,ylabel shift=-3pt,xlabel shift=-3pt]\addplot[fill=ink!55,draw=none] coordinates {(0,1.709e-06) (1,0.0002096) (2,0.00116) (3,0.01407) (4,0.04236) (5,0.2733) (6,0.3615) (7,1.292) (8,0.7603) (9,2.088) (10,0.9587) (11,2.464) (12,1.074) (13,1.981) (14,0.2754) (15,0.2352) (16,0.002305) (17,0.001143) (18,2.75e-07) (19,2.75e-07) (20,6.6e-08) (21,5.5e-08) (23,2e-09)};\draw[densely dotted] (axis cs:16,0) -- (axis cs:16,2.711);\end{axis}\begin{axis}[width=72mm,height=42mm,axis y line*=right,axis x line=none,ymode=log,xmin=-1,xmax=24,ylabel={\scriptsize eindeutig (log)},tick label style={font=\tiny},label style={font=\scriptsize},ylabel shift=-3pt]\addplot[mark=*,mark size=0.7pt,draw=black,thin] coordinates {(0,1709) (1,69411) (2,246768) (3,4437604) (4,7920819) (5,33349981) (6,11915688) (7,32964292) (8,7541174) (9,19102964) (10,4409203) (11,10908438) (12,1817210) (13,2318218) (14,25265) (15,13971) (16,15)};\end{axis}\end{tikzpicture}\end{minipage}
\vspace{3mm}

\noindent\begin{minipage}[t]{60mm}\vspace{0pt}\sbox\hmboard{\chessboard[setfen=K7/8/8/8/6b1/1q5k/6P1/6q1 b - - 0 1]}\usebox\hmboard\\[1mm]
\makebox[\wd\hmboard]{\sffamily\hspace{3mm}\textbf{h\#8}\hfill(2+4)}\end{minipage}\hfill
\begin{minipage}[t]{52mm}\vspace{0pt}
\textcolor{ink}{\sffamily\textsc{Lösung}}\par\smallskip\small
\begin{tabular}{@{}r@{\enspace}l@{\enspace}l@{}}
\textbf{1.} & \symking{}h2 & g3\\
\textbf{2.} & \symbishop{}d7 & g4\\
\textbf{3.} & \symking{}g3 & g5\\
\textbf{4.} & \symking{}f4 & g6\\
\textbf{5.} & \symking{}e5 & g7\\
\textbf{6.} & \symking{}d6 & g8=\symqueen{}\\
\textbf{7.} & \symking{}c7 & \symqueen{}xb3\\
\textbf{8.} & \symking{}c8 & \symqueen{}b8\#\\
\end{tabular}\par\vspace{3mm}
\qrcode[height=12mm]{https://helpman.komtera.lt/?fen=K7/8/8/8/6b1/1q5k/6P1/6q1&moves=8}\\{\sffamily\tiny\color{ink}scannen zum Lösen}
\end{minipage}
\vspace{2mm}

{\raggedright\noindent\textcolor{ink}{\sffamily\textsc{Themen}}\enspace \themetag{Umwandlung} \themetag{Excelsior} \themetag{Selbstblock} \themetag{ein Stein zieht} \themetag{ohne Schach} \themetag{Umnov}\par\smallskip{\footnotesize\textit{Hinweis:} das Diagramm hat keinen legalen letzten Zug, kann also in keiner Partie entstehen.}\par}

\clearpage
\phantomsection\label{hm:KPvkqqq}
\begin{tcolorbox}[enhanced,colback=ink,colframe=ink,arc=0pt,boxrule=0pt,left=4mm,right=4mm,top=2.5mm,bottom=2.5mm]
\color{white}\sffamily{\LARGE\bfseries KPvkqqq}\hfill{\Large h\#8}\\[0.5mm]
{\small 6 Steine\quad Schwarz am Zug}\hfill{\small Nr.~393\qquad\href{https://helpman.komtera.lt/?fen=K7/7q/8/8/8/1q6/P7/1q5k&moves=8}{\color{white}\bfseries lösen\,\ensuremath{\nearrow}}}
\end{tcolorbox}
\vspace{2mm}
\noindent\begin{minipage}[t]{36mm}\vspace{0pt}\sffamily\small\raggedright
\hmfigure{längste eindeutige}{{\Large h\#8}}\hmfigure{eindeutig bei dieser Länge}{{\Large 11} Stellungen}\hmfigure{längstes Matt}{{\Large h\#9.5}}\hmfigure{in der Klasse}{45,4 Milliarden Stellungen,\\ 3,41 Milliarden mit Hilfsmatt}
\end{minipage}\hfill
\begin{minipage}[t]{80mm}\vspace{0pt}\begin{tikzpicture}\begin{axis}[width=72mm,height=42mm,ybar,bar width=1.6pt,xmin=-1,xmax=20,ymin=0,ymax=804.3,axis y line*=left,axis x line*=bottom,xlabel={\scriptsize Halbzüge bis Matt},ylabel={\scriptsize Stellungen (Mio.)},tick label style={font=\tiny},label style={font=\scriptsize},enlarge x limits=false,ylabel shift=-3pt,xlabel shift=-3pt]\addplot[fill=ink!55,draw=none] coordinates {(1,0.03173) (2,0.2111) (3,2.932) (4,9.706) (5,68.62) (6,86.95) (7,365.9) (8,182.3) (9,616.7) (10,227.3) (11,731.2) (12,269.1) (13,670.4) (14,75.37) (15,103) (16,0.732) (17,0.5454) (18,0.000176) (19,1.9e-05)};\draw[densely dotted] (axis cs:16,0) -- (axis cs:16,804.3);\end{axis}\begin{axis}[width=72mm,height=42mm,axis y line*=right,axis x line=none,ymode=log,xmin=-1,xmax=20,ylabel={\scriptsize eindeutig (log)},tick label style={font=\tiny},label style={font=\scriptsize},ylabel shift=-3pt]\addplot[mark=*,mark size=0.7pt,draw=black,thin] coordinates {(1,11028) (2,48072) (3,970475) (4,1891399) (5,9179471) (6,3003988) (7,10186150) (8,1811897) (9,5761784) (10,1014604) (11,3286633) (12,473115) (13,910275) (14,10419) (15,9213) (16,11)};\end{axis}\end{tikzpicture}\end{minipage}
\vspace{3mm}

\noindent\begin{minipage}[t]{60mm}\vspace{0pt}\sbox\hmboard{\chessboard[setfen=K7/7q/8/8/8/1q6/P7/1q5k b - - 0 1]}\usebox\hmboard\\[1mm]
\makebox[\wd\hmboard]{\sffamily\hspace{3mm}\textbf{h\#8}\hfill(2+4)}\end{minipage}\hfill
\begin{minipage}[t]{52mm}\vspace{0pt}
\textcolor{ink}{\sffamily\textsc{Lösung}}\par\smallskip\small
\begin{tabular}{@{}r@{\enspace}l@{\enspace}l@{}}
\textbf{1.} & \symking{}g2 & a4\\
\textbf{2.} & \symking{}f3 & a5\\
\textbf{3.} & \symking{}e4 & a6\\
\textbf{4.} & \symqueen{}hb7+ & axb7\\
\textbf{5.} & \symking{}d5 & \symking{}b8\\
\textbf{6.} & \symking{}c6 & \symking{}c8\\
\textbf{7.} & \symqueen{}d5 & b8=\symqueen{}\\
\textbf{8.} & \symqueen{}bb5 & \symqueen{}c7\#\\
\end{tabular}\par\vspace{3mm}
\qrcode[height=12mm]{https://helpman.komtera.lt/?fen=K7/7q/8/8/8/1q6/P7/1q5k&moves=8}\\{\sffamily\tiny\color{ink}scannen zum Lösen}
\end{minipage}
\vspace{2mm}

{\raggedright\noindent\textcolor{ink}{\sffamily\textsc{Themen}}\enspace \themetag{Umwandlung} \themetag{Excelsior} \themetag{Selbstblock}\par}

\clearpage
\phantomsection\label{hm:KQBvkbb}
\begin{tcolorbox}[enhanced,colback=ink,colframe=ink,arc=0pt,boxrule=0pt,left=4mm,right=4mm,top=2.5mm,bottom=2.5mm]
\color{white}\sffamily{\LARGE\bfseries KQBvkbb}\hfill{\Large h\#8}\\[0.5mm]
{\small 6 Steine\quad Schwarz am Zug}\hfill{\small Nr.~394\qquad\href{https://helpman.komtera.lt/?fen=5b2/8/8/5b2/7B/8/7Q/K6k&moves=8}{\color{white}\bfseries lösen\,\ensuremath{\nearrow}}}
\end{tcolorbox}
\vspace{2mm}
\noindent\begin{minipage}[t]{36mm}\vspace{0pt}\sffamily\small\raggedright
\hmfigure{längste eindeutige}{{\Large h\#8}}\hmfigure{eindeutig bei dieser Länge}{{\Large 22} Stellungen}\hmfigure{längstes Matt}{{\Large h\#10.5}}\hmfigure{in der Klasse}{15,5 Milliarden Stellungen,\\ 4,04 Milliarden mit Hilfsmatt}
\end{minipage}\hfill
\begin{minipage}[t]{80mm}\vspace{0pt}\begin{tikzpicture}\begin{axis}[width=72mm,height=42mm,ybar,bar width=1.6pt,xmin=-1,xmax=22,ymin=0,ymax=1.921,axis y line*=left,axis x line*=bottom,xlabel={\scriptsize Halbzüge bis Matt},ylabel={\scriptsize Stellungen (Mrd.)},tick label style={font=\tiny},label style={font=\scriptsize},enlarge x limits=false,ylabel shift=-3pt,xlabel shift=-3pt]\addplot[fill=ink!55,draw=none] coordinates {(0,0.007707) (1,0.05122) (2,0.2446) (3,0.8554) (4,1.746) (5,0.8495) (6,0.2452) (7,0.01033) (8,0.005658) (9,4.35e-07) (10,0.007242) (11,9.67e-07) (12,0.007044) (13,1.8e-07) (14,0.003976) (15,1.5e-07) (16,0.002925) (17,1.113e-06) (18,0.001988) (19,2.61e-07) (20,0.0002862) (21,6.6e-08)};\draw[densely dotted] (axis cs:16,0) -- (axis cs:16,1.921);\end{axis}\begin{axis}[width=72mm,height=42mm,axis y line*=right,axis x line=none,ymode=log,xmin=-1,xmax=22,ylabel={\scriptsize eindeutig (log)},tick label style={font=\tiny},label style={font=\scriptsize},ylabel shift=-3pt]\addplot[mark=*,mark size=0.7pt,draw=black,thin] coordinates {(0,7707170) (1,44173174) (2,145243437) (3,174792321) (4,83745745) (5,2217251) (6,195226) (7,34) (8,77084) (9,9) (10,25000) (12,8691) (14,1069) (16,22)};\end{axis}\end{tikzpicture}\end{minipage}
\vspace{3mm}

\noindent\begin{minipage}[t]{60mm}\vspace{0pt}\sbox\hmboard{\chessboard[setfen=5b2/8/8/5b2/7B/8/7Q/K6k b - - 0 1]}\usebox\hmboard\\[1mm]
\makebox[\wd\hmboard]{\sffamily\hspace{3mm}\textbf{h\#8}\hfill(3+3)}\end{minipage}\hfill
\begin{minipage}[t]{52mm}\vspace{0pt}
\textcolor{ink}{\sffamily\textsc{Lösung}}\par\smallskip\small
\begin{tabular}{@{}r@{\enspace}l@{\enspace}l@{}}
\textbf{1.} & \symking{}xh2 & \symking{}b2\\
\textbf{2.} & \symking{}h3 & \symking{}c3\\
\textbf{3.} & \symking{}g4 & \symking{}d4\\
\textbf{4.} & \symking{}h5 & \symking{}e5\\
\textbf{5.} & \symking{}h6 & \symking{}f6\\
\textbf{6.} & \symking{}h7 & \symking{}f7\\
\textbf{7.} & \symking{}h8 & \symking{}xf8\\
\textbf{8.} & \symbishop{}h7 & \symbishop{}f6\#\\
\end{tabular}\par\vspace{3mm}
\qrcode[height=12mm]{https://helpman.komtera.lt/?fen=5b2/8/8/5b2/7B/8/7Q/K6k&moves=8}\\{\sffamily\tiny\color{ink}scannen zum Lösen}
\end{minipage}
\vspace{2mm}

{\raggedright\noindent\textcolor{ink}{\sffamily\textsc{Themen}}\enspace \themetag{Selbstblock} \themetag{ohne Schach}\par\smallskip{\footnotesize\textit{Hinweis:} die Partei am Zug steht im Diagramm im Schach; die Lösung beginnt mit einem Schlag.}\par}

\clearpage
\phantomsection\label{hm:KQNPvkq}
\begin{tcolorbox}[enhanced,colback=ink,colframe=ink,arc=0pt,boxrule=0pt,left=4mm,right=4mm,top=2.5mm,bottom=2.5mm]
\color{white}\sffamily{\LARGE\bfseries KQNPvkq}\hfill{\Large h\#8}\\[0.5mm]
{\small 6 Steine\quad Schwarz am Zug}\hfill{\small Nr.~395\qquad\href{https://helpman.komtera.lt/?fen=8/8/8/8/8/7N/qQ5P/k1K5&moves=8}{\color{white}\bfseries lösen\,\ensuremath{\nearrow}}}
\end{tcolorbox}
\vspace{2mm}
\noindent\begin{minipage}[t]{36mm}\vspace{0pt}\sffamily\small\raggedright
\hmfigure{längste eindeutige}{{\Large h\#8}}\hmfigure{eindeutig bei dieser Länge}{{\Large 16} Stellungen}\hmfigure{längstes Matt}{{\Large h\#9}}\hmfigure{in der Klasse}{45,4 Milliarden Stellungen,\\ 23,3 Milliarden mit Hilfsmatt}
\end{minipage}\hfill
\begin{minipage}[t]{80mm}\vspace{0pt}\begin{tikzpicture}\begin{axis}[width=72mm,height=42mm,ybar,bar width=1.6pt,xmin=-1,xmax=19,ymin=0,ymax=8.664,axis y line*=left,axis x line*=bottom,xlabel={\scriptsize Halbzüge bis Matt},ylabel={\scriptsize Stellungen (Mrd.)},tick label style={font=\tiny},label style={font=\scriptsize},enlarge x limits=false,ylabel shift=-3pt,xlabel shift=-3pt]\addplot[fill=ink!55,draw=none] coordinates {(0,0.026) (1,0.2127) (2,1.017) (3,3.739) (4,7.876) (5,6.242) (6,3.134) (7,0.7771) (8,0.118) (9,0.006533) (10,0.04033) (11,5.507e-06) (12,0.04767) (13,4.41e-07) (14,0.0166) (15,1.58e-07) (16,0.0001778) (17,1.8e-08) (18,1.7e-08)};\draw[densely dotted] (axis cs:16,0) -- (axis cs:16,8.664);\end{axis}\begin{axis}[width=72mm,height=42mm,axis y line*=right,axis x line=none,ymode=log,xmin=-1,xmax=19,ylabel={\scriptsize eindeutig (log)},tick label style={font=\tiny},label style={font=\scriptsize},ylabel shift=-3pt]\addplot[mark=*,mark size=0.7pt,draw=black,thin] coordinates {(0,26003700) (1,181082571) (2,575700676) (3,674303584) (4,471536091) (5,87540468) (6,8400004) (7,549342) (8,708951) (9,242) (10,339153) (11,92) (12,174283) (14,1899) (16,16)};\end{axis}\end{tikzpicture}\end{minipage}
\vspace{3mm}

\noindent\begin{minipage}[t]{60mm}\vspace{0pt}\sbox\hmboard{\chessboard[setfen=8/8/8/8/8/7N/qQ5P/k1K5 b - - 0 1]}\usebox\hmboard\\[1mm]
\makebox[\wd\hmboard]{\sffamily\hspace{3mm}\textbf{h\#8}\hfill(4+2)}\end{minipage}\hfill
\begin{minipage}[t]{52mm}\vspace{0pt}
\textcolor{ink}{\sffamily\textsc{Lösung}}\par\smallskip\small
\begin{tabular}{@{}r@{\enspace}l@{\enspace}l@{}}
\textbf{1.} & \symqueen{}xb2+ & \symking{}d1\\
\textbf{2.} & \symqueen{}g7 & \symknight{}f4\\
\textbf{3.} & \symking{}b2 & h4\\
\textbf{4.} & \symking{}c3 & h5\\
\textbf{5.} & \symking{}d4 & h6\\
\textbf{6.} & \symking{}e3 & hxg7\\
\textbf{7.} & \symking{}f2 & g8=\symqueen{}\\
\textbf{8.} & \symking{}f1 & \symqueen{}g2\#\\
\end{tabular}\par\vspace{3mm}
\qrcode[height=12mm]{https://helpman.komtera.lt/?fen=8/8/8/8/8/7N/qQ5P/k1K5&moves=8}\\{\sffamily\tiny\color{ink}scannen zum Lösen}
\end{minipage}
\vspace{2mm}

{\raggedright\noindent\textcolor{ink}{\sffamily\textsc{Themen}}\enspace \themetag{Umwandlung} \themetag{Excelsior} \themetag{Phönix}\par\smallskip{\footnotesize\textit{Hinweis:} die Partei am Zug steht im Diagramm im Schach; die Lösung beginnt mit einem Schlag.}\par}

\clearpage
\phantomsection\label{hm:KQNvkbn}
\begin{tcolorbox}[enhanced,colback=ink,colframe=ink,arc=0pt,boxrule=0pt,left=4mm,right=4mm,top=2.5mm,bottom=2.5mm]
\color{white}\sffamily{\LARGE\bfseries KQNvkbn}\hfill{\Large h\#8}\\[0.5mm]
{\small 6 Steine\quad Schwarz am Zug}\hfill{\small Nr.~396\qquad\href{https://helpman.komtera.lt/?fen=2n5/8/8/8/1N4Qk/8/4b3/K7&moves=8}{\color{white}\bfseries lösen\,\ensuremath{\nearrow}}}
\end{tcolorbox}
\vspace{2mm}
\noindent\begin{minipage}[t]{36mm}\vspace{0pt}\sffamily\small\raggedright
\hmfigure{längste eindeutige}{{\Large h\#8}}\hmfigure{eindeutig bei dieser Länge}{{\Large 21} Stellungen}\hmfigure{längstes Matt}{{\Large h\#10}}\hmfigure{in der Klasse}{15,5 Milliarden Stellungen,\\ 8,6 Milliarden mit Hilfsmatt}
\end{minipage}\hfill
\begin{minipage}[t]{80mm}\vspace{0pt}\begin{tikzpicture}\begin{axis}[width=72mm,height=42mm,ybar,bar width=1.6pt,xmin=-1,xmax=21,ymin=0,ymax=2.948,axis y line*=left,axis x line*=bottom,xlabel={\scriptsize Halbzüge bis Matt},ylabel={\scriptsize Stellungen (Mrd.)},tick label style={font=\tiny},label style={font=\scriptsize},enlarge x limits=false,ylabel shift=-3pt,xlabel shift=-3pt]\addplot[fill=ink!55,draw=none] coordinates {(0,0.009813) (1,0.07726) (2,0.3189) (3,1.145) (4,2.68) (5,2.196) (6,1.698) (7,0.3333) (8,0.07598) (9,0.002784) (10,0.01647) (11,2.665e-06) (12,0.0253) (13,5.45e-06) (14,0.01673) (15,9.67e-07) (16,0.003439) (17,1.064e-06) (18,0.0001255) (19,4.25e-07) (20,2.38e-07)};\draw[densely dotted] (axis cs:16,0) -- (axis cs:16,2.948);\end{axis}\begin{axis}[width=72mm,height=42mm,axis y line*=right,axis x line=none,ymode=log,xmin=-1,xmax=21,ylabel={\scriptsize eindeutig (log)},tick label style={font=\tiny},label style={font=\scriptsize},ylabel shift=-3pt]\addplot[mark=*,mark size=0.7pt,draw=black,thin] coordinates {(0,9813458) (1,65557884) (2,180802522) (3,224285070) (4,227795390) (5,34583265) (6,6628008) (7,171576) (8,135431) (9,4) (10,91142) (11,40) (12,26320) (13,4) (14,1917) (16,21)};\end{axis}\end{tikzpicture}\end{minipage}
\vspace{3mm}

\noindent\begin{minipage}[t]{60mm}\vspace{0pt}\sbox\hmboard{\chessboard[setfen=2n5/8/8/8/1N4Qk/8/4b3/K7 b - - 0 1]}\usebox\hmboard\\[1mm]
\makebox[\wd\hmboard]{\sffamily\hspace{3mm}\textbf{h\#8}\hfill(3+3)}\end{minipage}\hfill
\begin{minipage}[t]{52mm}\vspace{0pt}
\textcolor{ink}{\sffamily\textsc{Lösung}}\par\smallskip\small
\begin{tabular}{@{}r@{\enspace}l@{\enspace}l@{}}
\textbf{1.} & \symking{}xg4 & \symking{}b2\\
\textbf{2.} & \symking{}f5 & \symking{}c3\\
\textbf{3.} & \symking{}e6 & \symking{}d4\\
\textbf{4.} & \symking{}d7 & \symking{}c5\\
\textbf{5.} & \symking{}c7 & \symknight{}d5+\\
\textbf{6.} & \symking{}b8 & \symking{}c6\\
\textbf{7.} & \symking{}a8 & \symking{}c7\\
\textbf{8.} & \symknight{}a7 & \symknight{}b6\#\\
\end{tabular}\par\vspace{3mm}
\qrcode[height=12mm]{https://helpman.komtera.lt/?fen=2n5/8/8/8/1N4Qk/8/4b3/K7&moves=8}\\{\sffamily\tiny\color{ink}scannen zum Lösen}
\end{minipage}
\vspace{2mm}

{\raggedright\noindent\textcolor{ink}{\sffamily\textsc{Themen}}\enspace \themetag{reines Matt} \themetag{Mustermatt} \themetag{Selbstblock}\par\smallskip{\footnotesize\textit{Hinweis:} die Partei am Zug steht im Diagramm im Schach; die Lösung beginnt mit einem Schlag.}\par}

\clearpage
\phantomsection\label{hm:KQNvkbp}
\begin{tcolorbox}[enhanced,colback=ink,colframe=ink,arc=0pt,boxrule=0pt,left=4mm,right=4mm,top=2.5mm,bottom=2.5mm]
\color{white}\sffamily{\LARGE\bfseries KQNvkbp}\hfill{\Large h\#8}\\[0.5mm]
{\small 6 Steine\quad Schwarz am Zug}\hfill{\small Nr.~397\qquad\href{https://helpman.komtera.lt/?fen=8/6N1/8/6p1/8/8/2Q5/K1k3b1&moves=8}{\color{white}\bfseries lösen\,\ensuremath{\nearrow}}}
\end{tcolorbox}
\vspace{2mm}
\noindent\begin{minipage}[t]{36mm}\vspace{0pt}\sffamily\small\raggedright
\hmfigure{längste eindeutige}{{\Large h\#8}}\hmfigure{eindeutig bei dieser Länge}{{\Large 1\,700} Stellungen}\hmfigure{längstes Matt}{{\Large h\#10}}\hmfigure{in der Klasse}{45,4 Milliarden Stellungen,\\ 26,7 Milliarden mit Hilfsmatt}
\end{minipage}\hfill
\begin{minipage}[t]{80mm}\vspace{0pt}\begin{tikzpicture}\begin{axis}[width=72mm,height=42mm,ybar,bar width=1.6pt,xmin=-1,xmax=21,ymin=0,ymax=9.707,axis y line*=left,axis x line*=bottom,xlabel={\scriptsize Halbzüge bis Matt},ylabel={\scriptsize Stellungen (Mrd.)},tick label style={font=\tiny},label style={font=\scriptsize},enlarge x limits=false,ylabel shift=-3pt,xlabel shift=-3pt]\addplot[fill=ink!55,draw=none] coordinates {(0,0.03637) (1,0.2813) (2,1.127) (3,3.805) (4,8.825) (5,6.386) (6,4.994) (7,0.8298) (8,0.2014) (9,0.004454) (10,0.03535) (11,3.15e-07) (12,0.06034) (13,7.52e-07) (14,0.05965) (15,6.05e-07) (16,0.03755) (17,9.17e-07) (18,0.007334) (19,6.8e-08) (20,7.401e-05)};\draw[densely dotted] (axis cs:16,0) -- (axis cs:16,9.707);\end{axis}\begin{axis}[width=72mm,height=42mm,axis y line*=right,axis x line=none,ymode=log,xmin=-1,xmax=21,ylabel={\scriptsize eindeutig (log)},tick label style={font=\tiny},label style={font=\scriptsize},ylabel shift=-3pt]\addplot[mark=*,mark size=0.7pt,draw=black,thin] coordinates {(0,36368760) (1,234437758) (2,654340415) (3,722283387) (4,699131485) (5,88635014) (6,18917751) (7,318885) (8,289157) (9,62) (10,226254) (12,100799) (14,14840) (16,1700)};\end{axis}\end{tikzpicture}\end{minipage}
\vspace{3mm}

\noindent\begin{minipage}[t]{60mm}\vspace{0pt}\sbox\hmboard{\chessboard[setfen=8/6N1/8/6p1/8/8/2Q5/K1k3b1 b - - 0 1]}\usebox\hmboard\\[1mm]
\makebox[\wd\hmboard]{\sffamily\hspace{3mm}\textbf{h\#8}\hfill(3+3)}\end{minipage}\hfill
\begin{minipage}[t]{52mm}\vspace{0pt}
\textcolor{ink}{\sffamily\textsc{Lösung}}\par\smallskip\small
\begin{tabular}{@{}r@{\enspace}l@{\enspace}l@{}}
\textbf{1.} & \symking{}xc2 & \symknight{}e8\\
\textbf{2.} & \symking{}d1 & \symking{}b2\\
\textbf{3.} & \symbishop{}f2 & \symking{}c3\\
\textbf{4.} & \symbishop{}h4 & \symking{}d4\\
\textbf{5.} & \symking{}e2 & \symking{}e5\\
\textbf{6.} & \symking{}f3 & \symking{}f6\\
\textbf{7.} & \symking{}g4 & \symking{}g7\\
\textbf{8.} & \symking{}h5 & \symknight{}f6\#\\
\end{tabular}\par\vspace{3mm}
\qrcode[height=12mm]{https://helpman.komtera.lt/?fen=8/6N1/8/6p1/8/8/2Q5/K1k3b1&moves=8}\\{\sffamily\tiny\color{ink}scannen zum Lösen}
\end{minipage}
\vspace{2mm}

{\raggedright\noindent\textcolor{ink}{\sffamily\textsc{Themen}}\enspace \themetag{reines Matt} \themetag{Mustermatt} \themetag{Idealmatt} \themetag{Selbstblock} \themetag{ohne Schach}\par\smallskip{\footnotesize\textit{Hinweis:} die Partei am Zug steht im Diagramm im Schach; die Lösung beginnt mit einem Schlag.}\par}

\clearpage
\phantomsection\label{hm:KQNvknn}
\begin{tcolorbox}[enhanced,colback=ink,colframe=ink,arc=0pt,boxrule=0pt,left=4mm,right=4mm,top=2.5mm,bottom=2.5mm]
\color{white}\sffamily{\LARGE\bfseries KQNvknn}\hfill{\Large h\#8}\\[0.5mm]
{\small 6 Steine\quad Schwarz am Zug}\hfill{\small Nr.~398\qquad\href{https://helpman.komtera.lt/?fen=3n4/N1n5/8/8/6Qk/8/8/K7&moves=8}{\color{white}\bfseries lösen\,\ensuremath{\nearrow}}}
\end{tcolorbox}
\vspace{2mm}
\noindent\begin{minipage}[t]{36mm}\vspace{0pt}\sffamily\small\raggedright
\hmfigure{längste eindeutige}{{\Large h\#8}}\hmfigure{eindeutig bei dieser Länge}{{\Large 18} Stellungen}\hmfigure{längstes Matt}{{\Large h\#10}}\hmfigure{in der Klasse}{15,5 Milliarden Stellungen,\\ 4,42 Milliarden mit Hilfsmatt}
\end{minipage}\hfill
\begin{minipage}[t]{80mm}\vspace{0pt}\begin{tikzpicture}\begin{axis}[width=72mm,height=42mm,ybar,bar width=1.6pt,xmin=-1,xmax=21,ymin=0,ymax=1.45,axis y line*=left,axis x line*=bottom,xlabel={\scriptsize Halbzüge bis Matt},ylabel={\scriptsize Stellungen (Mrd.)},tick label style={font=\tiny},label style={font=\scriptsize},enlarge x limits=false,ylabel shift=-3pt,xlabel shift=-3pt]\addplot[fill=ink!55,draw=none] coordinates {(0,0.005514) (1,0.04148) (2,0.1581) (3,0.554) (4,1.318) (5,1.097) (6,0.9771) (7,0.1827) (8,0.05285) (9,0.002029) (10,0.00876) (11,7.01e-07) (12,0.01219) (13,1.994e-06) (14,0.009282) (15,1.842e-06) (16,0.002518) (17,4.23e-07) (18,8.593e-05) (20,3.8e-08)};\draw[densely dotted] (axis cs:16,0) -- (axis cs:16,1.45);\end{axis}\begin{axis}[width=72mm,height=42mm,axis y line*=right,axis x line=none,ymode=log,xmin=-1,xmax=21,ylabel={\scriptsize eindeutig (log)},tick label style={font=\tiny},label style={font=\scriptsize},ylabel shift=-3pt]\addplot[mark=*,mark size=0.7pt,draw=black,thin] coordinates {(0,5514013) (1,34258507) (2,87430052) (3,108591334) (4,115969015) (5,16321665) (6,3487872) (7,99271) (8,74561) (9,5) (10,25964) (11,35) (12,9755) (13,59) (14,1046) (16,18)};\end{axis}\end{tikzpicture}\end{minipage}
\vspace{3mm}

\noindent\begin{minipage}[t]{60mm}\vspace{0pt}\sbox\hmboard{\chessboard[setfen=3n4/N1n5/8/8/6Qk/8/8/K7 b - - 0 1]}\usebox\hmboard\\[1mm]
\makebox[\wd\hmboard]{\sffamily\hspace{3mm}\textbf{h\#8}\hfill(3+3)}\end{minipage}\hfill
\begin{minipage}[t]{52mm}\vspace{0pt}
\textcolor{ink}{\sffamily\textsc{Lösung}}\par\smallskip\small
\begin{tabular}{@{}r@{\enspace}l@{\enspace}l@{}}
\textbf{1.} & \symking{}xg4 & \symking{}b2\\
\textbf{2.} & \symking{}f5 & \symking{}c3\\
\textbf{3.} & \symking{}e6 & \symking{}d4\\
\textbf{4.} & \symking{}d7 & \symking{}e5\\
\textbf{5.} & \symknight{}a8 & \symking{}f6\\
\textbf{6.} & \symking{}c7 & \symking{}e7\\
\textbf{7.} & \symking{}b8 & \symking{}d7\\
\textbf{8.} & \symknight{}b7 & \symknight{}c6\#\\
\end{tabular}\par\vspace{3mm}
\qrcode[height=12mm]{https://helpman.komtera.lt/?fen=3n4/N1n5/8/8/6Qk/8/8/K7&moves=8}\\{\sffamily\tiny\color{ink}scannen zum Lösen}
\end{minipage}
\vspace{2mm}

{\raggedright\noindent\textcolor{ink}{\sffamily\textsc{Themen}}\enspace \themetag{reines Matt} \themetag{Mustermatt} \themetag{Idealmatt} \themetag{Selbstblock} \themetag{ohne Schach}\par\smallskip{\footnotesize\textit{Hinweis:} die Partei am Zug steht im Diagramm im Schach; die Lösung beginnt mit einem Schlag.}\par}

\clearpage
\phantomsection\label{hm:KQNvkqn}
\begin{tcolorbox}[enhanced,colback=ink,colframe=ink,arc=0pt,boxrule=0pt,left=4mm,right=4mm,top=2.5mm,bottom=2.5mm]
\color{white}\sffamily{\LARGE\bfseries KQNvkqn}\hfill{\Large h\#8}\\[0.5mm]
{\small 6 Steine\quad Schwarz am Zug}\hfill{\small Nr.~399\qquad\href{https://helpman.komtera.lt/?fen=nN6/8/8/8/8/1q4Qk/8/K7&moves=8}{\color{white}\bfseries lösen\,\ensuremath{\nearrow}}}
\end{tcolorbox}
\vspace{2mm}
\noindent\begin{minipage}[t]{36mm}\vspace{0pt}\sffamily\small\raggedright
\hmfigure{längste eindeutige}{{\Large h\#8}}\hmfigure{eindeutig bei dieser Länge}{{\Large 23} Stellungen}\hmfigure{längstes Matt}{{\Large h\#9}}\hmfigure{in der Klasse}{15,5 Milliarden Stellungen,\\ 7,5 Milliarden mit Hilfsmatt}
\end{minipage}\hfill
\begin{minipage}[t]{80mm}\vspace{0pt}\begin{tikzpicture}\begin{axis}[width=72mm,height=42mm,ybar,bar width=1.6pt,xmin=-1,xmax=19,ymin=0,ymax=2.513,axis y line*=left,axis x line*=bottom,xlabel={\scriptsize Halbzüge bis Matt},ylabel={\scriptsize Stellungen (Mrd.)},tick label style={font=\tiny},label style={font=\scriptsize},enlarge x limits=false,ylabel shift=-3pt,xlabel shift=-3pt]\addplot[fill=ink!55,draw=none] coordinates {(0,0.005623) (1,0.04797) (2,0.2297) (3,0.9502) (4,2.071) (5,2.284) (6,1.342) (7,0.4601) (8,0.05588) (9,0.005755) (10,0.015) (11,1.455e-06) (12,0.02086) (13,3.613e-06) (14,0.01158) (15,4.76e-07) (16,0.001526) (17,1.6e-08) (18,1.69e-05)};\draw[densely dotted] (axis cs:16,0) -- (axis cs:16,2.513);\end{axis}\begin{axis}[width=72mm,height=42mm,axis y line*=right,axis x line=none,ymode=log,xmin=-1,xmax=19,ylabel={\scriptsize eindeutig (log)},tick label style={font=\tiny},label style={font=\scriptsize},ylabel shift=-3pt]\addplot[mark=*,mark size=0.7pt,draw=black,thin] coordinates {(0,5622715) (1,42901888) (2,131332638) (3,198252141) (4,165513282) (5,43508799) (6,3957142) (7,341728) (8,140993) (9,16) (10,69411) (11,73) (12,20562) (13,30) (14,2353) (16,23)};\end{axis}\end{tikzpicture}\end{minipage}
\vspace{3mm}

\noindent\begin{minipage}[t]{60mm}\vspace{0pt}\sbox\hmboard{\chessboard[setfen=nN6/8/8/8/8/1q4Qk/8/K7 b - - 0 1]}\usebox\hmboard\\[1mm]
\makebox[\wd\hmboard]{\sffamily\hspace{3mm}\textbf{h\#8}\hfill(3+3)}\end{minipage}\hfill
\begin{minipage}[t]{52mm}\vspace{0pt}
\textcolor{ink}{\sffamily\textsc{Lösung}}\par\smallskip\small
\begin{tabular}{@{}r@{\enspace}l@{\enspace}l@{}}
\textbf{1.} & \symqueen{}xg3 & \symking{}b2\\
\textbf{2.} & \symqueen{}g5 & \symking{}c3\\
\textbf{3.} & \symknight{}c7 & \symking{}d4\\
\textbf{4.} & \symqueen{}h4+ & \symking{}e5\\
\textbf{5.} & \symknight{}e6 & \symknight{}d7\\
\textbf{6.} & \symknight{}g5 & \symking{}f6\\
\textbf{7.} & \symking{}g4 & \symking{}g7\\
\textbf{8.} & \symking{}h5 & \symknight{}f6\#\\
\end{tabular}\par\vspace{3mm}
\qrcode[height=12mm]{https://helpman.komtera.lt/?fen=nN6/8/8/8/8/1q4Qk/8/K7&moves=8}\\{\sffamily\tiny\color{ink}scannen zum Lösen}
\end{minipage}
\vspace{2mm}

{\raggedright\noindent\textcolor{ink}{\sffamily\textsc{Themen}}\enspace \themetag{reines Matt} \themetag{Mustermatt} \themetag{Idealmatt} \themetag{Selbstblock}\par\smallskip{\footnotesize\textit{Hinweis:} die Partei am Zug steht im Diagramm im Schach; die Lösung beginnt mit einem Schlag.}\par}

\clearpage
\phantomsection\label{hm:KQNvkrn}
\begin{tcolorbox}[enhanced,colback=ink,colframe=ink,arc=0pt,boxrule=0pt,left=4mm,right=4mm,top=2.5mm,bottom=2.5mm]
\color{white}\sffamily{\LARGE\bfseries KQNvkrn}\hfill{\Large h\#8}\\[0.5mm]
{\small 6 Steine\quad Schwarz am Zug}\hfill{\small Nr.~400\qquad\href{https://helpman.komtera.lt/?fen=8/8/6Qk/8/n7/N7/6r1/K7&moves=8}{\color{white}\bfseries lösen\,\ensuremath{\nearrow}}}
\end{tcolorbox}
\vspace{2mm}
\noindent\begin{minipage}[t]{36mm}\vspace{0pt}\sffamily\small\raggedright
\hmfigure{längste eindeutige}{{\Large h\#8}}\hmfigure{eindeutig bei dieser Länge}{{\Large 9} Stellungen}\hmfigure{längstes Matt}{{\Large h\#9}}\hmfigure{in der Klasse}{15,5 Milliarden Stellungen,\\ 8,21 Milliarden mit Hilfsmatt}
\end{minipage}\hfill
\begin{minipage}[t]{80mm}\vspace{0pt}\begin{tikzpicture}\begin{axis}[width=72mm,height=42mm,ybar,bar width=1.6pt,xmin=-1,xmax=19,ymin=0,ymax=2.955,axis y line*=left,axis x line*=bottom,xlabel={\scriptsize Halbzüge bis Matt},ylabel={\scriptsize Stellungen (Mrd.)},tick label style={font=\tiny},label style={font=\scriptsize},enlarge x limits=false,ylabel shift=-3pt,xlabel shift=-3pt]\addplot[fill=ink!55,draw=none] coordinates {(0,0.008485) (1,0.06923) (2,0.3172) (3,1.148) (4,2.686) (5,2.265) (6,1.362) (7,0.2699) (8,0.02857) (9,0.0004719) (10,0.01912) (11,2.369e-06) (12,0.02209) (13,3.239e-06) (14,0.01082) (15,6.18e-07) (16,0.001552) (18,1.804e-05)};\draw[densely dotted] (axis cs:16,0) -- (axis cs:16,2.955);\end{axis}\begin{axis}[width=72mm,height=42mm,axis y line*=right,axis x line=none,ymode=log,xmin=-1,xmax=19,ylabel={\scriptsize eindeutig (log)},tick label style={font=\tiny},label style={font=\scriptsize},ylabel shift=-3pt]\addplot[mark=*,mark size=0.7pt,draw=black,thin] coordinates {(0,8485468) (1,60044453) (2,184995264) (3,216900323) (4,200942681) (5,34236632) (6,3257781) (7,39153) (8,194263) (9,8) (10,94206) (11,100) (12,29152) (13,3) (14,2924) (16,9)};\end{axis}\end{tikzpicture}\end{minipage}
\vspace{3mm}

\noindent\begin{minipage}[t]{60mm}\vspace{0pt}\sbox\hmboard{\chessboard[setfen=8/8/6Qk/8/n7/N7/6r1/K7 b - - 0 1]}\usebox\hmboard\\[1mm]
\makebox[\wd\hmboard]{\sffamily\hspace{3mm}\textbf{h\#8}\hfill(3+3)}\end{minipage}\hfill
\begin{minipage}[t]{52mm}\vspace{0pt}
\textcolor{ink}{\sffamily\textsc{Lösung}}\par\smallskip\small
\begin{tabular}{@{}r@{\enspace}l@{\enspace}l@{}}
\textbf{1.} & \symking{}xg6 & \symknight{}c2\\
\textbf{2.} & \symking{}f5 & \symking{}a2\\
\textbf{3.} & \symking{}e4 & \symking{}b1\\
\textbf{4.} & \symking{}d3 & \symking{}c1\\
\textbf{5.} & \symking{}c3 & \symking{}d1\\
\textbf{6.} & \symking{}b2 & \symknight{}d4\\
\textbf{7.} & \symking{}a1 & \symking{}c1\\
\textbf{8.} & \symrook{}a2 & \symknight{}b3\#\\
\end{tabular}\par\vspace{3mm}
\qrcode[height=12mm]{https://helpman.komtera.lt/?fen=8/8/6Qk/8/n7/N7/6r1/K7&moves=8}\\{\sffamily\tiny\color{ink}scannen zum Lösen}
\end{minipage}
\vspace{2mm}

{\raggedright\noindent\textcolor{ink}{\sffamily\textsc{Themen}}\enspace \themetag{reines Matt} \themetag{Mustermatt} \themetag{Switchback} \themetag{Selbstblock} \themetag{ohne Schach}\par\smallskip{\footnotesize\textit{Hinweis:} die Partei am Zug steht im Diagramm im Schach; die Lösung beginnt mit einem Schlag.}\par}

\clearpage
\phantomsection\label{hm:KQNvkrr}
\begin{tcolorbox}[enhanced,colback=ink,colframe=ink,arc=0pt,boxrule=0pt,left=4mm,right=4mm,top=2.5mm,bottom=2.5mm]
\color{white}\sffamily{\LARGE\bfseries KQNvkrr}\hfill{\Large h\#8}\\[0.5mm]
{\small 6 Steine\quad Schwarz am Zug}\hfill{\small Nr.~401\qquad\href{https://helpman.komtera.lt/?fen=8/6Qk/8/8/6N1/8/1r2r3/K7&moves=8}{\color{white}\bfseries lösen\,\ensuremath{\nearrow}}}
\end{tcolorbox}
\vspace{2mm}
\noindent\begin{minipage}[t]{36mm}\vspace{0pt}\sffamily\small\raggedright
\hmfigure{längste eindeutige}{{\Large h\#8}}\hmfigure{eindeutig bei dieser Länge}{{\Large 4} Stellungen}\hmfigure{längstes Matt}{{\Large h\#9.5}}\hmfigure{in der Klasse}{15,5 Milliarden Stellungen,\\ 3,78 Milliarden mit Hilfsmatt}
\end{minipage}\hfill
\begin{minipage}[t]{80mm}\vspace{0pt}\begin{tikzpicture}\begin{axis}[width=72mm,height=42mm,ybar,bar width=1.6pt,xmin=-1,xmax=20,ymin=0,ymax=1.377,axis y line*=left,axis x line*=bottom,xlabel={\scriptsize Halbzüge bis Matt},ylabel={\scriptsize Stellungen (Mrd.)},tick label style={font=\tiny},label style={font=\scriptsize},enlarge x limits=false,ylabel shift=-3pt,xlabel shift=-3pt]\addplot[fill=ink!55,draw=none] coordinates {(0,0.00321) (1,0.02799) (2,0.1483) (3,0.5728) (4,1.252) (5,1.148) (6,0.4763) (7,0.1213) (8,0.009419) (9,0.0001571) (10,0.01081) (11,2.727e-06) (12,0.008198) (13,4.459e-06) (14,0.002474) (15,2.371e-06) (16,0.0001651) (17,8.47e-07) (18,3.26e-07) (19,1.71e-07)};\draw[densely dotted] (axis cs:16,0) -- (axis cs:16,1.377);\end{axis}\begin{axis}[width=72mm,height=42mm,axis y line*=right,axis x line=none,ymode=log,xmin=-1,xmax=20,ylabel={\scriptsize eindeutig (log)},tick label style={font=\tiny},label style={font=\scriptsize},ylabel shift=-3pt]\addplot[mark=*,mark size=0.7pt,draw=black,thin] coordinates {(0,3210433) (1,25037641) (2,88500296) (3,107964091) (4,73035204) (5,14122565) (6,745025) (7,15458) (8,138241) (10,98388) (11,26) (12,25876) (14,2896) (15,23) (16,4)};\end{axis}\end{tikzpicture}\end{minipage}
\vspace{3mm}

\noindent\begin{minipage}[t]{60mm}\vspace{0pt}\sbox\hmboard{\chessboard[setfen=8/6Qk/8/8/6N1/8/1r2r3/K7 b - - 0 1]}\usebox\hmboard\\[1mm]
\makebox[\wd\hmboard]{\sffamily\hspace{3mm}\textbf{h\#8}\hfill(3+3)}\end{minipage}\hfill
\begin{minipage}[t]{52mm}\vspace{0pt}
\textcolor{ink}{\sffamily\textsc{Lösung}}\par\smallskip\small
\begin{tabular}{@{}r@{\enspace}l@{\enspace}l@{}}
\textbf{1.} & \symking{}xg7 & \symknight{}h2\\
\textbf{2.} & \symrook{}e6 & \symknight{}f1\\
\textbf{3.} & \symrook{}h2 & \symking{}b1\\
\textbf{4.} & \symrook{}h1 & \symking{}c2\\
\textbf{5.} & \symrook{}eh6 & \symking{}d3\\
\textbf{6.} & \symking{}g6 & \symking{}e4\\
\textbf{7.} & \symking{}h5 & \symking{}f5\\
\textbf{8.} & \symrook{}h4 & \symknight{}g3\#\\
\end{tabular}\par\vspace{3mm}
\qrcode[height=12mm]{https://helpman.komtera.lt/?fen=8/6Qk/8/8/6N1/8/1r2r3/K7&moves=8}\\{\sffamily\tiny\color{ink}scannen zum Lösen}
\end{minipage}
\vspace{2mm}

{\raggedright\noindent\textcolor{ink}{\sffamily\textsc{Themen}}\enspace \themetag{reines Matt} \themetag{Mustermatt} \themetag{Idealmatt} \themetag{Selbstblock} \themetag{ohne Schach} \themetag{Umnov}\par\smallskip{\footnotesize\textit{Hinweis:} die Partei am Zug steht im Diagramm im Schach; die Lösung beginnt mit einem Schlag.}\par}

\clearpage
\phantomsection\label{hm:KQPPvkp}
\begin{tcolorbox}[enhanced,colback=ink,colframe=ink,arc=0pt,boxrule=0pt,left=4mm,right=4mm,top=2.5mm,bottom=2.5mm]
\color{white}\sffamily{\LARGE\bfseries KQPPvkp}\hfill{\Large h\#8}\\[0.5mm]
{\small 6 Steine\quad Schwarz am Zug}\hfill{\small Nr.~402\qquad\href{https://helpman.komtera.lt/?fen=8/8/p7/8/P7/K7/P5Q1/6k1&moves=8}{\color{white}\bfseries lösen\,\ensuremath{\nearrow}}}
\end{tcolorbox}
\vspace{2mm}
\noindent\begin{minipage}[t]{36mm}\vspace{0pt}\sffamily\small\raggedright
\hmfigure{längste eindeutige}{{\Large h\#8}}\hmfigure{eindeutig bei dieser Länge}{{\Large 48} Stellungen}\hmfigure{längstes Matt}{{\Large h\#15}}\hmfigure{in der Klasse}{25,6 Milliarden Stellungen,\\ 8,11 Milliarden mit Hilfsmatt}
\end{minipage}\hfill
\begin{minipage}[t]{80mm}\vspace{0pt}\begin{tikzpicture}\begin{axis}[width=72mm,height=42mm,ybar,bar width=1.6pt,xmin=-1,xmax=31,ymin=0,ymax=2.557,axis y line*=left,axis x line*=bottom,xlabel={\scriptsize Halbzüge bis Matt},ylabel={\scriptsize Stellungen (Mrd.)},tick label style={font=\tiny},label style={font=\scriptsize},enlarge x limits=false,ylabel shift=-3pt,xlabel shift=-3pt]\addplot[fill=ink!55,draw=none] coordinates {(0,0.0126) (1,0.09441) (2,0.2941) (3,0.9845) (4,2.324) (5,1.604) (6,1.855) (7,0.5042) (8,0.3532) (9,0.03959) (10,0.02912) (11,0.0003118) (12,0.01036) (13,1.54e-07) (14,0.003549) (15,2e-09) (16,0.0001579) (18,4.358e-05) (20,3.937e-05) (22,2.272e-05) (24,7.869e-06) (26,3.822e-06) (28,1.076e-06) (30,4e-09)};\draw[densely dotted] (axis cs:16,0) -- (axis cs:16,2.557);\end{axis}\begin{axis}[width=72mm,height=42mm,axis y line*=right,axis x line=none,ymode=log,xmin=-1,xmax=31,ylabel={\scriptsize eindeutig (log)},tick label style={font=\tiny},label style={font=\scriptsize},ylabel shift=-3pt]\addplot[mark=*,mark size=0.7pt,draw=black,thin] coordinates {(0,12600928) (1,59747549) (2,158405819) (3,164748117) (4,191604778) (5,38244816) (6,20291271) (7,1953893) (8,783449) (9,5796) (10,88784) (11,33) (12,35329) (14,1030) (16,48)};\end{axis}\end{tikzpicture}\end{minipage}
\vspace{3mm}

\noindent\begin{minipage}[t]{60mm}\vspace{0pt}\sbox\hmboard{\chessboard[setfen=8/8/p7/8/P7/K7/P5Q1/6k1 b - - 0 1]}\usebox\hmboard\\[1mm]
\makebox[\wd\hmboard]{\sffamily\hspace{3mm}\textbf{h\#8}\hfill(4+2)}\end{minipage}\hfill
\begin{minipage}[t]{52mm}\vspace{0pt}
\textcolor{ink}{\sffamily\textsc{Lösung}}\par\smallskip\small
\begin{tabular}{@{}r@{\enspace}l@{\enspace}l@{}}
\textbf{1.} & \symking{}xg2 & \symking{}b4\\
\textbf{2.} & a5+ & \symking{}xa5\\
\textbf{3.} & \symking{}f3 & \symking{}b4\\
\textbf{4.} & \symking{}e4 & a5\\
\textbf{5.} & \symking{}d5 & a6\\
\textbf{6.} & \symking{}c6 & a7\\
\textbf{7.} & \symking{}b7 & \symking{}c5\\
\textbf{8.} & \symking{}a6 & a8=\symqueen{}\#\\
\end{tabular}\par\vspace{3mm}
\qrcode[height=12mm]{https://helpman.komtera.lt/?fen=8/8/p7/8/P7/K7/P5Q1/6k1&moves=8}\\{\sffamily\tiny\color{ink}scannen zum Lösen}
\end{minipage}
\vspace{2mm}

{\raggedright\noindent\textcolor{ink}{\sffamily\textsc{Themen}}\enspace \themetag{reines Matt} \themetag{Mustermatt} \themetag{Spiegelmatt} \themetag{Umwandlung} \themetag{Switchback} \themetag{Phönix}\par\smallskip{\footnotesize\textit{Hinweis:} die Partei am Zug steht im Diagramm im Schach; die Lösung beginnt mit einem Schlag.}\par}

\clearpage
\phantomsection\label{hm:KQPvkbb}
\begin{tcolorbox}[enhanced,colback=ink,colframe=ink,arc=0pt,boxrule=0pt,left=4mm,right=4mm,top=2.5mm,bottom=2.5mm]
\color{white}\sffamily{\LARGE\bfseries KQPvkbb}\hfill{\Large h\#8}\\[0.5mm]
{\small 6 Steine\quad Schwarz am Zug}\hfill{\small Nr.~403\qquad\href{https://helpman.komtera.lt/?fen=7k/7Q/8/3b4/8/bP6/8/K7&moves=8}{\color{white}\bfseries lösen\,\ensuremath{\nearrow}}}
\end{tcolorbox}
\vspace{2mm}
\noindent\begin{minipage}[t]{36mm}\vspace{0pt}\sffamily\small\raggedright
\hmfigure{längste eindeutige}{{\Large h\#8}}\hmfigure{eindeutig bei dieser Länge}{{\Large 108} Stellungen}\hmfigure{längstes Matt}{{\Large h\#9.5}}\hmfigure{in der Klasse}{45,4 Milliarden Stellungen,\\ 12,8 Milliarden mit Hilfsmatt}
\end{minipage}\hfill
\begin{minipage}[t]{80mm}\vspace{0pt}\begin{tikzpicture}\begin{axis}[width=72mm,height=42mm,ybar,bar width=1.6pt,xmin=-1,xmax=20,ymin=0,ymax=3.266,axis y line*=left,axis x line*=bottom,xlabel={\scriptsize Halbzüge bis Matt},ylabel={\scriptsize Stellungen (Mrd.)},tick label style={font=\tiny},label style={font=\scriptsize},enlarge x limits=false,ylabel shift=-3pt,xlabel shift=-3pt]\addplot[fill=ink!55,draw=none] coordinates {(0,0.008538) (1,0.06985) (2,0.2686) (3,1.12) (4,2.652) (5,2.969) (6,2.931) (7,1.554) (8,0.8837) (9,0.2438) (10,0.08401) (11,0.00428) (12,0.02604) (13,4.019e-06) (14,0.01362) (15,3.15e-07) (16,0.00132) (17,1.04e-07) (18,6.559e-06) (19,1e-09)};\draw[densely dotted] (axis cs:16,0) -- (axis cs:16,3.266);\end{axis}\begin{axis}[width=72mm,height=42mm,axis y line*=right,axis x line=none,ymode=log,xmin=-1,xmax=20,ylabel={\scriptsize eindeutig (log)},tick label style={font=\tiny},label style={font=\scriptsize},ylabel shift=-3pt]\addplot[mark=*,mark size=0.7pt,draw=black,thin] coordinates {(0,8537826) (1,54026838) (2,132297311) (3,196223605) (4,222847554) (5,91748874) (6,40826269) (7,11424463) (8,2832417) (9,141305) (10,136807) (11,149) (12,62698) (13,5) (14,14357) (16,108)};\end{axis}\end{tikzpicture}\end{minipage}
\vspace{3mm}

\noindent\begin{minipage}[t]{60mm}\vspace{0pt}\sbox\hmboard{\chessboard[setfen=7k/7Q/8/3b4/8/bP6/8/K7 b - - 0 1]}\usebox\hmboard\\[1mm]
\makebox[\wd\hmboard]{\sffamily\hspace{3mm}\textbf{h\#8}\hfill(3+3)}\end{minipage}\hfill
\begin{minipage}[t]{52mm}\vspace{0pt}
\textcolor{ink}{\sffamily\textsc{Lösung}}\par\smallskip\small
\begin{tabular}{@{}r@{\enspace}l@{\enspace}l@{}}
\textbf{1.} & \symking{}xh7 & \symking{}a2\\
\textbf{2.} & \symking{}g6 & \symking{}xa3\\
\textbf{3.} & \symking{}f5 & b4\\
\textbf{4.} & \symking{}e4 & b5\\
\textbf{5.} & \symking{}d3 & b6\\
\textbf{6.} & \symking{}c2 & b7\\
\textbf{7.} & \symking{}b1 & b8=\symqueen{}+\\
\textbf{8.} & \symking{}a1 & \symqueen{}b2\#\\
\end{tabular}\par\vspace{3mm}
\qrcode[height=12mm]{https://helpman.komtera.lt/?fen=7k/7Q/8/3b4/8/bP6/8/K7&moves=8}\\{\sffamily\tiny\color{ink}scannen zum Lösen}
\end{minipage}
\vspace{2mm}

{\raggedright\noindent\textcolor{ink}{\sffamily\textsc{Themen}}\enspace \themetag{Umwandlung} \themetag{ein Stein zieht} \themetag{Phönix}\par\smallskip{\footnotesize\textit{Hinweis:} die Partei am Zug steht im Diagramm im Schach; die Lösung beginnt mit einem Schlag.}\par}

\clearpage
\phantomsection\label{hm:KQPvkbn}
\begin{tcolorbox}[enhanced,colback=ink,colframe=ink,arc=0pt,boxrule=0pt,left=4mm,right=4mm,top=2.5mm,bottom=2.5mm]
\color{white}\sffamily{\LARGE\bfseries KQPvkbn}\hfill{\Large h\#8}\\[0.5mm]
{\small 6 Steine\quad Schwarz am Zug}\hfill{\small Nr.~404\qquad\href{https://helpman.komtera.lt/?fen=8/6Qk/8/8/2n5/b7/P7/K7&moves=8}{\color{white}\bfseries lösen\,\ensuremath{\nearrow}}}
\end{tcolorbox}
\vspace{2mm}
\noindent\begin{minipage}[t]{36mm}\vspace{0pt}\sffamily\small\raggedright
\hmfigure{längste eindeutige}{{\Large h\#8}}\hmfigure{eindeutig bei dieser Länge}{{\Large 552} Stellungen}\hmfigure{längstes Matt}{{\Large h\#9}}\hmfigure{in der Klasse}{45,4 Milliarden Stellungen,\\ 26,5 Milliarden mit Hilfsmatt}
\end{minipage}\hfill
\begin{minipage}[t]{80mm}\vspace{0pt}\begin{tikzpicture}\begin{axis}[width=72mm,height=42mm,ybar,bar width=1.6pt,xmin=-1,xmax=19,ymin=0,ymax=6.976,axis y line*=left,axis x line*=bottom,xlabel={\scriptsize Halbzüge bis Matt},ylabel={\scriptsize Stellungen (Mrd.)},tick label style={font=\tiny},label style={font=\scriptsize},enlarge x limits=false,ylabel shift=-3pt,xlabel shift=-3pt]\addplot[fill=ink!55,draw=none] coordinates {(0,0.01924) (1,0.152) (2,0.536) (3,2.174) (4,5.144) (5,5.72) (6,6.342) (7,3.29) (8,2.198) (9,0.573) (10,0.2304) (11,0.01311) (12,0.05273) (13,1.041e-05) (14,0.03907) (15,8.619e-06) (16,0.004917) (17,1.57e-07) (18,3.008e-05)};\draw[densely dotted] (axis cs:16,0) -- (axis cs:16,6.976);\end{axis}\begin{axis}[width=72mm,height=42mm,axis y line*=right,axis x line=none,ymode=log,xmin=-1,xmax=19,ylabel={\scriptsize eindeutig (log)},tick label style={font=\tiny},label style={font=\scriptsize},ylabel shift=-3pt]\addplot[mark=*,mark size=0.7pt,draw=black,thin] coordinates {(0,19244106) (1,111864112) (2,249793988) (3,370982029) (4,448906915) (5,197118333) (6,104843749) (7,26110564) (8,7899120) (9,358134) (10,274629) (11,111) (12,155224) (13,2) (14,48619) (16,552)};\end{axis}\end{tikzpicture}\end{minipage}
\vspace{3mm}

\noindent\begin{minipage}[t]{60mm}\vspace{0pt}\sbox\hmboard{\chessboard[setfen=8/6Qk/8/8/2n5/b7/P7/K7 b - - 0 1]}\usebox\hmboard\\[1mm]
\makebox[\wd\hmboard]{\sffamily\hspace{3mm}\textbf{h\#8}\hfill(3+3)}\end{minipage}\hfill
\begin{minipage}[t]{52mm}\vspace{0pt}
\textcolor{ink}{\sffamily\textsc{Lösung}}\par\smallskip\small
\begin{tabular}{@{}r@{\enspace}l@{\enspace}l@{}}
\textbf{1.} & \symking{}xg7 & \symking{}b1\\
\textbf{2.} & \symbishop{}b4 & a4\\
\textbf{3.} & \symking{}f6 & a5\\
\textbf{4.} & \symking{}e5 & a6\\
\textbf{5.} & \symking{}d4 & a7\\
\textbf{6.} & \symking{}c3 & a8=\symqueen{}\\
\textbf{7.} & \symking{}b3 & \symqueen{}g2\\
\textbf{8.} & \symking{}a3 & \symqueen{}a2\#\\
\end{tabular}\par\vspace{3mm}
\qrcode[height=12mm]{https://helpman.komtera.lt/?fen=8/6Qk/8/8/2n5/b7/P7/K7&moves=8}\\{\sffamily\tiny\color{ink}scannen zum Lösen}
\end{minipage}
\vspace{2mm}

{\raggedright\noindent\textcolor{ink}{\sffamily\textsc{Themen}}\enspace \themetag{Umwandlung} \themetag{Excelsior} \themetag{Rundlauf} \themetag{Selbstblock} \themetag{Phönix} \themetag{ohne Schach}\par\smallskip{\footnotesize\textit{Hinweis:} die Partei am Zug steht im Diagramm im Schach; die Lösung beginnt mit einem Schlag.}\par}

\clearpage
\phantomsection\label{hm:KQPvknn}
\begin{tcolorbox}[enhanced,colback=ink,colframe=ink,arc=0pt,boxrule=0pt,left=4mm,right=4mm,top=2.5mm,bottom=2.5mm]
\color{white}\sffamily{\LARGE\bfseries KQPvknn}\hfill{\Large h\#8}\\[0.5mm]
{\small 6 Steine\quad Schwarz am Zug}\hfill{\small Nr.~405\qquad\href{https://helpman.komtera.lt/?fen=6Qk/8/8/8/3n4/8/3P4/1K1n4&moves=8}{\color{white}\bfseries lösen\,\ensuremath{\nearrow}}}
\end{tcolorbox}
\vspace{2mm}
\noindent\begin{minipage}[t]{36mm}\vspace{0pt}\sffamily\small\raggedright
\hmfigure{längste eindeutige}{{\Large h\#8}}\hmfigure{eindeutig bei dieser Länge}{{\Large 345} Stellungen}\hmfigure{längstes Matt}{{\Large h\#10}}\hmfigure{in der Klasse}{45,4 Milliarden Stellungen,\\ 13,6 Milliarden mit Hilfsmatt}
\end{minipage}\hfill
\begin{minipage}[t]{80mm}\vspace{0pt}\begin{tikzpicture}\begin{axis}[width=72mm,height=42mm,ybar,bar width=1.6pt,xmin=-1,xmax=21,ymin=0,ymax=3.75,axis y line*=left,axis x line*=bottom,xlabel={\scriptsize Halbzüge bis Matt},ylabel={\scriptsize Stellungen (Mrd.)},tick label style={font=\tiny},label style={font=\scriptsize},enlarge x limits=false,ylabel shift=-3pt,xlabel shift=-3pt]\addplot[fill=ink!55,draw=none] coordinates {(0,0.01089) (1,0.08251) (2,0.265) (3,1.042) (4,2.401) (5,2.736) (6,3.409) (7,1.764) (8,1.353) (9,0.3269) (10,0.1479) (11,0.008561) (12,0.02732) (13,5.173e-06) (14,0.02214) (15,5.044e-06) (16,0.004251) (17,1.857e-06) (18,2.042e-05) (20,2e-09)};\draw[densely dotted] (axis cs:16,0) -- (axis cs:16,3.75);\end{axis}\begin{axis}[width=72mm,height=42mm,axis y line*=right,axis x line=none,ymode=log,xmin=-1,xmax=21,ylabel={\scriptsize eindeutig (log)},tick label style={font=\tiny},label style={font=\scriptsize},ylabel shift=-3pt]\addplot[mark=*,mark size=0.7pt,draw=black,thin] coordinates {(0,10885088) (1,56835227) (2,114559677) (3,176051906) (4,208194234) (5,103146986) (6,66280303) (7,14615138) (8,5087629) (9,188640) (10,134495) (11,40) (12,77335) (13,4) (14,29302) (16,345)};\end{axis}\end{tikzpicture}\end{minipage}
\vspace{3mm}

\noindent\begin{minipage}[t]{60mm}\vspace{0pt}\sbox\hmboard{\chessboard[setfen=6Qk/8/8/8/3n4/8/3P4/1K1n4 b - - 0 1]}\usebox\hmboard\\[1mm]
\makebox[\wd\hmboard]{\sffamily\hspace{3mm}\textbf{h\#8}\hfill(3+3)}\end{minipage}\hfill
\begin{minipage}[t]{52mm}\vspace{0pt}
\textcolor{ink}{\sffamily\textsc{Lösung}}\par\smallskip\small
\begin{tabular}{@{}r@{\enspace}l@{\enspace}l@{}}
\textbf{1.} & \symking{}xg8 & \symking{}c1\\
\textbf{2.} & \symknight{}e3 & dxe3\\
\textbf{3.} & \symking{}f7 & e4\\
\textbf{4.} & \symking{}e6 & e5\\
\textbf{5.} & \symking{}d5 & e6\\
\textbf{6.} & \symking{}c4 & e7\\
\textbf{7.} & \symking{}b3 & e8=\symqueen{}\\
\textbf{8.} & \symking{}a2 & \symqueen{}a4\#\\
\end{tabular}\par\vspace{3mm}
\qrcode[height=12mm]{https://helpman.komtera.lt/?fen=6Qk/8/8/8/3n4/8/3P4/1K1n4&moves=8}\\{\sffamily\tiny\color{ink}scannen zum Lösen}
\end{minipage}
\vspace{2mm}

{\raggedright\noindent\textcolor{ink}{\sffamily\textsc{Themen}}\enspace \themetag{reines Matt} \themetag{Mustermatt} \themetag{Spiegelmatt} \themetag{Umwandlung} \themetag{Excelsior} \themetag{Phönix} \themetag{ohne Schach} \themetag{Umnov}\par\smallskip{\footnotesize\textit{Hinweis:} die Partei am Zug steht im Diagramm im Schach; die Lösung beginnt mit einem Schlag.}\par}

\clearpage
\phantomsection\label{hm:KQPvkqb}
\begin{tcolorbox}[enhanced,colback=ink,colframe=ink,arc=0pt,boxrule=0pt,left=4mm,right=4mm,top=2.5mm,bottom=2.5mm]
\color{white}\sffamily{\LARGE\bfseries KQPvkqb}\hfill{\Large h\#8}\\[0.5mm]
{\small 6 Steine\quad Schwarz am Zug}\hfill{\small Nr.~406\qquad\href{https://helpman.komtera.lt/?fen=8/6Qk/8/8/2b5/q7/P7/K7&moves=8}{\color{white}\bfseries lösen\,\ensuremath{\nearrow}}}
\end{tcolorbox}
\vspace{2mm}
\noindent\begin{minipage}[t]{36mm}\vspace{0pt}\sffamily\small\raggedright
\hmfigure{längste eindeutige}{{\Large h\#8}}\hmfigure{eindeutig bei dieser Länge}{{\Large 33} Stellungen}\hmfigure{längstes Matt}{{\Large h\#12}}\hmfigure{in der Klasse}{45,4 Milliarden Stellungen,\\ 22,5 Milliarden mit Hilfsmatt}
\end{minipage}\hfill
\begin{minipage}[t]{80mm}\vspace{0pt}\begin{tikzpicture}\begin{axis}[width=72mm,height=42mm,ybar,bar width=1.6pt,xmin=-1,xmax=25,ymin=0,ymax=6.577,axis y line*=left,axis x line*=bottom,xlabel={\scriptsize Halbzüge bis Matt},ylabel={\scriptsize Stellungen (Mrd.)},tick label style={font=\tiny},label style={font=\scriptsize},enlarge x limits=false,ylabel shift=-3pt,xlabel shift=-3pt]\addplot[fill=ink!55,draw=none] coordinates {(0,0.009732) (1,0.08602) (2,0.385) (3,1.843) (4,4.132) (5,5.979) (6,4.585) (7,3.372) (8,1.347) (9,0.5956) (10,0.1319) (11,0.01278) (12,0.04005) (13,3.225e-05) (14,0.02282) (15,5.855e-06) (16,0.001627) (17,1.823e-06) (18,2.898e-06) (19,4e-09) (20,3.1e-08) (22,3.2e-08) (24,1.8e-08)};\draw[densely dotted] (axis cs:16,0) -- (axis cs:16,6.577);\end{axis}\begin{axis}[width=72mm,height=42mm,axis y line*=right,axis x line=none,ymode=log,xmin=-1,xmax=25,ylabel={\scriptsize eindeutig (log)},tick label style={font=\tiny},label style={font=\scriptsize},ylabel shift=-3pt]\addplot[mark=*,mark size=0.7pt,draw=black,thin] coordinates {(0,9731653) (1,71634996) (2,194158447) (3,359177463) (4,347358637) (5,197724234) (6,58135584) (7,26423706) (8,4568774) (9,549143) (10,227217) (11,409) (12,110233) (13,3) (14,16170) (15,6) (16,33)};\end{axis}\end{tikzpicture}\end{minipage}
\vspace{3mm}

\noindent\begin{minipage}[t]{60mm}\vspace{0pt}\sbox\hmboard{\chessboard[setfen=8/6Qk/8/8/2b5/q7/P7/K7 b - - 0 1]}\usebox\hmboard\\[1mm]
\makebox[\wd\hmboard]{\sffamily\hspace{3mm}\textbf{h\#8}\hfill(3+3)}\end{minipage}\hfill
\begin{minipage}[t]{52mm}\vspace{0pt}
\textcolor{ink}{\sffamily\textsc{Lösung}}\par\smallskip\small
\begin{tabular}{@{}r@{\enspace}l@{\enspace}l@{}}
\textbf{1.} & \symking{}xg7 & \symking{}b1\\
\textbf{2.} & \symqueen{}c5 & a4\\
\textbf{3.} & \symking{}f6 & a5\\
\textbf{4.} & \symking{}e5 & a6\\
\textbf{5.} & \symking{}d4 & a7\\
\textbf{6.} & \symking{}c3 & a8=\symqueen{}\\
\textbf{7.} & \symking{}b4 & \symking{}b2\\
\textbf{8.} & \symqueen{}b5 & \symqueen{}a3\#\\
\end{tabular}\par\vspace{3mm}
\qrcode[height=12mm]{https://helpman.komtera.lt/?fen=8/6Qk/8/8/2b5/q7/P7/K7&moves=8}\\{\sffamily\tiny\color{ink}scannen zum Lösen}
\end{minipage}
\vspace{2mm}

{\raggedright\noindent\textcolor{ink}{\sffamily\textsc{Themen}}\enspace \themetag{Umwandlung} \themetag{Excelsior} \themetag{Selbstblock} \themetag{Phönix} \themetag{ohne Schach}\par\smallskip{\footnotesize\textit{Hinweis:} die Partei am Zug steht im Diagramm im Schach; die Lösung beginnt mit einem Schlag.}\par}

\clearpage
\phantomsection\label{hm:KQPvkqn}
\begin{tcolorbox}[enhanced,colback=ink,colframe=ink,arc=0pt,boxrule=0pt,left=4mm,right=4mm,top=2.5mm,bottom=2.5mm]
\color{white}\sffamily{\LARGE\bfseries KQPvkqn}\hfill{\Large h\#8}\\[0.5mm]
{\small 6 Steine\quad Schwarz am Zug}\hfill{\small Nr.~407\qquad\href{https://helpman.komtera.lt/?fen=8/7Q/7k/8/7P/3n4/2q5/K7&moves=8}{\color{white}\bfseries lösen\,\ensuremath{\nearrow}}}
\end{tcolorbox}
\vspace{2mm}
\noindent\begin{minipage}[t]{36mm}\vspace{0pt}\sffamily\small\raggedright
\hmfigure{längste eindeutige}{{\Large h\#8}}\hmfigure{eindeutig bei dieser Länge}{{\Large 422} Stellungen}\hmfigure{längstes Matt}{{\Large h\#10}}\hmfigure{in der Klasse}{45,4 Milliarden Stellungen,\\ 23,2 Milliarden mit Hilfsmatt}
\end{minipage}\hfill
\begin{minipage}[t]{80mm}\vspace{0pt}\begin{tikzpicture}\begin{axis}[width=72mm,height=42mm,ybar,bar width=1.6pt,xmin=-1,xmax=21,ymin=0,ymax=6.277,axis y line*=left,axis x line*=bottom,xlabel={\scriptsize Halbzüge bis Matt},ylabel={\scriptsize Stellungen (Mrd.)},tick label style={font=\tiny},label style={font=\scriptsize},enlarge x limits=false,ylabel shift=-3pt,xlabel shift=-3pt]\addplot[fill=ink!55,draw=none] coordinates {(0,0.01078) (1,0.09301) (2,0.3925) (3,1.811) (4,4.027) (5,5.706) (6,4.928) (7,3.57) (8,1.685) (9,0.6993) (10,0.1874) (11,0.02191) (12,0.04244) (13,7.319e-05) (14,0.02932) (15,1.48e-05) (16,0.004413) (17,1.259e-06) (18,4.348e-05) (20,9e-09)};\draw[densely dotted] (axis cs:16,0) -- (axis cs:16,6.277);\end{axis}\begin{axis}[width=72mm,height=42mm,axis y line*=right,axis x line=none,ymode=log,xmin=-1,xmax=21,ylabel={\scriptsize eindeutig (log)},tick label style={font=\tiny},label style={font=\scriptsize},ylabel shift=-3pt]\addplot[mark=*,mark size=0.7pt,draw=black,thin] coordinates {(0,10777265) (1,74926952) (2,189472125) (3,339774214) (4,346269979) (5,205923143) (6,74568862) (7,29845262) (8,6228967) (9,772596) (10,218812) (11,1105) (12,119285) (13,3) (14,34567) (16,422)};\end{axis}\end{tikzpicture}\end{minipage}
\vspace{3mm}

\noindent\begin{minipage}[t]{60mm}\vspace{0pt}\sbox\hmboard{\chessboard[setfen=8/7Q/7k/8/7P/3n4/2q5/K7 b - - 0 1]}\usebox\hmboard\\[1mm]
\makebox[\wd\hmboard]{\sffamily\hspace{3mm}\textbf{h\#8}\hfill(3+3)}\end{minipage}\hfill
\begin{minipage}[t]{52mm}\vspace{0pt}
\textcolor{ink}{\sffamily\textsc{Lösung}}\par\smallskip\small
\begin{tabular}{@{}r@{\enspace}l@{\enspace}l@{}}
\textbf{1.} & \symking{}xh7 & h5\\
\textbf{2.} & \symqueen{}a2+ & \symking{}xa2\\
\textbf{3.} & \symking{}h6 & \symking{}b3\\
\textbf{4.} & \symking{}g5 & h6\\
\textbf{5.} & \symking{}f4 & h7\\
\textbf{6.} & \symking{}e3 & h8=\symqueen{}\\
\textbf{7.} & \symking{}d2 & \symqueen{}h2+\\
\textbf{8.} & \symking{}c1 & \symqueen{}c2\#\\
\end{tabular}\par\vspace{3mm}
\qrcode[height=12mm]{https://helpman.komtera.lt/?fen=8/7Q/7k/8/7P/3n4/2q5/K7&moves=8}\\{\sffamily\tiny\color{ink}scannen zum Lösen}
\end{minipage}
\vspace{2mm}

{\raggedright\noindent\textcolor{ink}{\sffamily\textsc{Themen}}\enspace \themetag{Umwandlung} \themetag{Switchback} \themetag{Phönix} \themetag{Umnov}\par\smallskip{\footnotesize\textit{Hinweis:} die Partei am Zug steht im Diagramm im Schach; die Lösung beginnt mit einem Schlag.}\par}

\clearpage
\phantomsection\label{hm:KQPvkqp}
\begin{tcolorbox}[enhanced,colback=ink,colframe=ink,arc=0pt,boxrule=0pt,left=4mm,right=4mm,top=2.5mm,bottom=2.5mm]
\color{white}\sffamily{\LARGE\bfseries KQPvkqp}\hfill{\Large h\#8}\\[0.5mm]
{\small 6 Steine\quad Schwarz am Zug}\hfill{\small Nr.~408\qquad\href{https://helpman.komtera.lt/?fen=8/8/8/8/2p5/2P5/2q1Q3/K3k3&moves=8}{\color{white}\bfseries lösen\,\ensuremath{\nearrow}}}
\end{tcolorbox}
\vspace{2mm}
\noindent\begin{minipage}[t]{36mm}\vspace{0pt}\sffamily\small\raggedright
\hmfigure{längste eindeutige}{{\Large h\#8}}\hmfigure{eindeutig bei dieser Länge}{{\Large 945} Stellungen}\hmfigure{längstes Matt}{{\Large h\#13}}\hmfigure{in der Klasse}{34,1 Milliarden Stellungen,\\ 17,8 Milliarden mit Hilfsmatt}
\end{minipage}\hfill
\begin{minipage}[t]{80mm}\vspace{0pt}\begin{tikzpicture}\begin{axis}[width=72mm,height=42mm,ybar,bar width=1.6pt,xmin=-1,xmax=27,ymin=0,ymax=4.952,axis y line*=left,axis x line*=bottom,xlabel={\scriptsize Halbzüge bis Matt},ylabel={\scriptsize Stellungen (Mrd.)},tick label style={font=\tiny},label style={font=\scriptsize},enlarge x limits=false,ylabel shift=-3pt,xlabel shift=-3pt]\addplot[fill=ink!55,draw=none] coordinates {(0,0.01026) (1,0.08711) (2,0.3606) (3,1.578) (4,3.566) (5,4.502) (6,3.798) (7,2.378) (8,1.043) (9,0.3648) (10,0.09554) (11,0.007475) (12,0.03361) (13,2.718e-05) (14,0.01771) (15,9.73e-07) (16,0.001852) (17,2.9e-07) (18,0.0001823) (19,8e-09) (20,3.187e-06) (21,1.6e-08) (22,3e-09) (23,3e-08) (24,4.2e-08) (25,2.7e-08) (26,1e-09)};\draw[densely dotted] (axis cs:16,0) -- (axis cs:16,4.952);\end{axis}\begin{axis}[width=72mm,height=42mm,axis y line*=right,axis x line=none,ymode=log,xmin=-1,xmax=27,ylabel={\scriptsize eindeutig (log)},tick label style={font=\tiny},label style={font=\scriptsize},ylabel shift=-3pt]\addplot[mark=*,mark size=0.7pt,draw=black,thin] coordinates {(0,10258133) (1,69188856) (2,180372600) (3,294086895) (4,295074312) (5,142682495) (6,45588117) (7,15805085) (8,2951571) (9,286499) (10,175819) (11,796) (12,117059) (13,10) (14,12938) (16,945)};\end{axis}\end{tikzpicture}\end{minipage}
\vspace{3mm}

\noindent\begin{minipage}[t]{60mm}\vspace{0pt}\sbox\hmboard{\chessboard[setfen=8/8/8/8/2p5/2P5/2q1Q3/K3k3 b - - 0 1]}\usebox\hmboard\\[1mm]
\makebox[\wd\hmboard]{\sffamily\hspace{3mm}\textbf{h\#8}\hfill(3+3)}\end{minipage}\hfill
\begin{minipage}[t]{52mm}\vspace{0pt}
\textcolor{ink}{\sffamily\textsc{Lösung}}\par\smallskip\small
\begin{tabular}{@{}r@{\enspace}l@{\enspace}l@{}}
\textbf{1.} & \symqueen{}xe2 & \symking{}b1\\
\textbf{2.} & \symqueen{}e7 & \symking{}c1\\
\textbf{3.} & \symqueen{}b4 & cxb4\\
\textbf{4.} & \symking{}e2 & b5\\
\textbf{5.} & \symking{}d3 & b6\\
\textbf{6.} & \symking{}c3 & b7\\
\textbf{7.} & \symking{}b3 & b8=\symqueen{}+\\
\textbf{8.} & \symking{}a2 & \symqueen{}b2\#\\
\end{tabular}\par\vspace{3mm}
\qrcode[height=12mm]{https://helpman.komtera.lt/?fen=8/8/8/8/2p5/2P5/2q1Q3/K3k3&moves=8}\\{\sffamily\tiny\color{ink}scannen zum Lösen}
\end{minipage}
\vspace{2mm}

{\raggedright\noindent\textcolor{ink}{\sffamily\textsc{Themen}}\enspace \themetag{Umwandlung} \themetag{Phönix}\par\smallskip{\footnotesize\textit{Hinweis:} die Partei am Zug steht im Diagramm im Schach; die Lösung beginnt mit einem Schlag.}\par}

\clearpage
\phantomsection\label{hm:KQPvkqq}
\begin{tcolorbox}[enhanced,colback=ink,colframe=ink,arc=0pt,boxrule=0pt,left=4mm,right=4mm,top=2.5mm,bottom=2.5mm]
\color{white}\sffamily{\LARGE\bfseries KQPvkqq}\hfill{\Large h\#8}\\[0.5mm]
{\small 6 Steine\quad Schwarz am Zug}\hfill{\small Nr.~409\qquad\href{https://helpman.komtera.lt/?fen=8/1q6/8/7Q/7k/6q1/6P1/K7&moves=8}{\color{white}\bfseries lösen\,\ensuremath{\nearrow}}}
\end{tcolorbox}
\vspace{2mm}
\noindent\begin{minipage}[t]{36mm}\vspace{0pt}\sffamily\small\raggedright
\hmfigure{längste eindeutige}{{\Large h\#8}}\hmfigure{eindeutig bei dieser Länge}{{\Large 14} Stellungen}\hmfigure{längstes Matt}{{\Large h\#9.5}}\hmfigure{in der Klasse}{45,4 Milliarden Stellungen,\\ 9,85 Milliarden mit Hilfsmatt}
\end{minipage}\hfill
\begin{minipage}[t]{80mm}\vspace{0pt}\begin{tikzpicture}\begin{axis}[width=72mm,height=42mm,ybar,bar width=1.6pt,xmin=-1,xmax=20,ymin=0,ymax=3.187,axis y line*=left,axis x line*=bottom,xlabel={\scriptsize Halbzüge bis Matt},ylabel={\scriptsize Stellungen (Mrd.)},tick label style={font=\tiny},label style={font=\scriptsize},enlarge x limits=false,ylabel shift=-3pt,xlabel shift=-3pt]\addplot[fill=ink!55,draw=none] coordinates {(0,0.002692) (1,0.02557) (2,0.1309) (3,0.7226) (4,1.551) (5,2.897) (6,1.766) (7,1.807) (8,0.507) (9,0.3557) (10,0.05063) (11,0.01097) (12,0.01531) (13,5.809e-05) (14,0.008813) (15,2.542e-05) (16,0.0005536) (17,3.83e-06) (18,6.49e-07) (19,3e-09)};\draw[densely dotted] (axis cs:16,0) -- (axis cs:16,3.187);\end{axis}\begin{axis}[width=72mm,height=42mm,axis y line*=right,axis x line=none,ymode=log,xmin=-1,xmax=20,ylabel={\scriptsize eindeutig (log)},tick label style={font=\tiny},label style={font=\scriptsize},ylabel shift=-3pt]\addplot[mark=*,mark size=0.7pt,draw=black,thin] coordinates {(0,2692048) (1,22369418) (2,66816346) (3,155524352) (4,132748456) (5,108951116) (6,22158143) (7,14944862) (8,1800703) (9,478266) (10,81932) (11,820) (12,50678) (13,50) (14,6690) (15,6) (16,14)};\end{axis}\end{tikzpicture}\end{minipage}
\vspace{3mm}

\noindent\begin{minipage}[t]{60mm}\vspace{0pt}\sbox\hmboard{\chessboard[setfen=8/1q6/8/7Q/7k/6q1/6P1/K7 b - - 0 1]}\usebox\hmboard\\[1mm]
\makebox[\wd\hmboard]{\sffamily\hspace{3mm}\textbf{h\#8}\hfill(3+3)}\end{minipage}\hfill
\begin{minipage}[t]{52mm}\vspace{0pt}
\textcolor{ink}{\sffamily\textsc{Lösung}}\par\smallskip\small
\begin{tabular}{@{}r@{\enspace}l@{\enspace}l@{}}
\textbf{1.} & \symking{}xh5 & \symking{}a2\\
\textbf{2.} & \symqueen{}e3 & g4+\\
\textbf{3.} & \symking{}g6 & g5\\
\textbf{4.} & \symking{}f5 & g6\\
\textbf{5.} & \symking{}e4 & g7\\
\textbf{6.} & \symking{}d3 & g8=\symqueen{}\\
\textbf{7.} & \symqueen{}g2+ & \symking{}b3\\
\textbf{8.} & \symqueen{}gd2 & \symqueen{}c4\#\\
\end{tabular}\par\vspace{3mm}
\qrcode[height=12mm]{https://helpman.komtera.lt/?fen=8/1q6/8/7Q/7k/6q1/6P1/K7&moves=8}\\{\sffamily\tiny\color{ink}scannen zum Lösen}
\end{minipage}
\vspace{2mm}

{\raggedright\noindent\textcolor{ink}{\sffamily\textsc{Themen}}\enspace \themetag{Umwandlung} \themetag{Excelsior} \themetag{Selbstblock} \themetag{Phönix} \themetag{Umnov}\par\smallskip{\footnotesize\textit{Hinweis:} die Partei am Zug steht im Diagramm im Schach; die Lösung beginnt mit einem Schlag.}\par}

\clearpage
\phantomsection\label{hm:KQPvkqr}
\begin{tcolorbox}[enhanced,colback=ink,colframe=ink,arc=0pt,boxrule=0pt,left=4mm,right=4mm,top=2.5mm,bottom=2.5mm]
\color{white}\sffamily{\LARGE\bfseries KQPvkqr}\hfill{\Large h\#8}\\[0.5mm]
{\small 6 Steine\quad Schwarz am Zug}\hfill{\small Nr.~410\qquad\href{https://helpman.komtera.lt/?fen=K7/1r6/8/8/2q5/2P5/7Q/7k&moves=8}{\color{white}\bfseries lösen\,\ensuremath{\nearrow}}}
\end{tcolorbox}
\vspace{2mm}
\noindent\begin{minipage}[t]{36mm}\vspace{0pt}\sffamily\small\raggedright
\hmfigure{längste eindeutige}{{\Large h\#8}}\hmfigure{eindeutig bei dieser Länge}{{\Large 1} Stellung}\hmfigure{längstes Matt}{{\Large h\#9}}\hmfigure{in der Klasse}{45,4 Milliarden Stellungen,\\ 21,5 Milliarden mit Hilfsmatt}
\end{minipage}\hfill
\begin{minipage}[t]{80mm}\vspace{0pt}\begin{tikzpicture}\begin{axis}[width=72mm,height=42mm,ybar,bar width=1.6pt,xmin=-1,xmax=19,ymin=0,ymax=7.259,axis y line*=left,axis x line*=bottom,xlabel={\scriptsize Halbzüge bis Matt},ylabel={\scriptsize Stellungen (Mrd.)},tick label style={font=\tiny},label style={font=\scriptsize},enlarge x limits=false,ylabel shift=-3pt,xlabel shift=-3pt]\addplot[fill=ink!55,draw=none] coordinates {(0,0.008154) (1,0.0753) (2,0.3809) (3,1.867) (4,4.408) (5,6.599) (6,4.303) (7,3.114) (8,0.5055) (9,0.1601) (10,0.0396) (11,0.0009315) (12,0.03478) (13,3.557e-05) (14,0.005224) (15,1.998e-05) (16,5.134e-05) (17,5.2e-07) (18,1.74e-07)};\draw[densely dotted] (axis cs:16,0) -- (axis cs:16,7.259);\end{axis}\begin{axis}[width=72mm,height=42mm,axis y line*=right,axis x line=none,ymode=log,xmin=-1,xmax=19,ylabel={\scriptsize eindeutig (log)},tick label style={font=\tiny},label style={font=\scriptsize},ylabel shift=-3pt]\addplot[mark=*,mark size=0.7pt,draw=black,thin] coordinates {(0,8153658) (1,63763813) (2,202923738) (3,368114109) (4,334688017) (5,185066305) (6,27081261) (7,7760898) (8,781858) (9,46842) (10,198095) (11,299) (12,31386) (13,26) (14,136) (16,1)};\end{axis}\end{tikzpicture}\end{minipage}
\vspace{3mm}

\noindent\begin{minipage}[t]{60mm}\vspace{0pt}\sbox\hmboard{\chessboard[setfen=K7/1r6/8/8/2q5/2P5/7Q/7k b - - 0 1]}\usebox\hmboard\\[1mm]
\makebox[\wd\hmboard]{\sffamily\hspace{3mm}\textbf{h\#8}\hfill(3+3)}\end{minipage}\hfill
\begin{minipage}[t]{52mm}\vspace{0pt}
\textcolor{ink}{\sffamily\textsc{Lösung}}\par\smallskip\small
\begin{tabular}{@{}r@{\enspace}l@{\enspace}l@{}}
\textbf{1.} & \symking{}xh2 & \symking{}xb7\\
\textbf{2.} & \symqueen{}d4 & cxd4\\
\textbf{3.} & \symking{}g3 & d5\\
\textbf{4.} & \symking{}f4 & d6\\
\textbf{5.} & \symking{}e5 & d7\\
\textbf{6.} & \symking{}d6 & \symking{}a7\\
\textbf{7.} & \symking{}c7 & \symking{}a6\\
\textbf{8.} & \symking{}b8 & d8=\symqueen{}\#\\
\end{tabular}\par\vspace{3mm}
\qrcode[height=12mm]{https://helpman.komtera.lt/?fen=K7/1r6/8/8/2q5/2P5/7Q/7k&moves=8}\\{\sffamily\tiny\color{ink}scannen zum Lösen}
\end{minipage}
\vspace{2mm}

{\raggedright\noindent\textcolor{ink}{\sffamily\textsc{Themen}}\enspace \themetag{reines Matt} \themetag{Mustermatt} \themetag{Idealmatt} \themetag{Spiegelmatt} \themetag{Umwandlung} \themetag{Phönix} \themetag{ohne Schach} \themetag{Umnov}\par\smallskip{\footnotesize\textit{Hinweis:} die Partei am Zug steht im Diagramm im Schach; die Lösung beginnt mit einem Schlag.}\par}

\clearpage
\phantomsection\label{hm:KQPvkrb}
\begin{tcolorbox}[enhanced,colback=ink,colframe=ink,arc=0pt,boxrule=0pt,left=4mm,right=4mm,top=2.5mm,bottom=2.5mm]
\color{white}\sffamily{\LARGE\bfseries KQPvkrb}\hfill{\Large h\#8}\\[0.5mm]
{\small 6 Steine\quad Schwarz am Zug}\hfill{\small Nr.~411\qquad\href{https://helpman.komtera.lt/?fen=1r6/K7/8/8/2b5/2P5/7Q/7k&moves=8}{\color{white}\bfseries lösen\,\ensuremath{\nearrow}}}
\end{tcolorbox}
\vspace{2mm}
\noindent\begin{minipage}[t]{36mm}\vspace{0pt}\sffamily\small\raggedright
\hmfigure{längste eindeutige}{{\Large h\#8}}\hmfigure{eindeutig bei dieser Länge}{{\Large 2} Stellungen}\hmfigure{längstes Matt}{{\Large h\#9}}\hmfigure{in der Klasse}{45,4 Milliarden Stellungen,\\ 24,7 Milliarden mit Hilfsmatt}
\end{minipage}\hfill
\begin{minipage}[t]{80mm}\vspace{0pt}\begin{tikzpicture}\begin{axis}[width=72mm,height=42mm,ybar,bar width=1.6pt,xmin=-1,xmax=19,ymin=0,ymax=7.217,axis y line*=left,axis x line*=bottom,xlabel={\scriptsize Halbzüge bis Matt},ylabel={\scriptsize Stellungen (Mrd.)},tick label style={font=\tiny},label style={font=\scriptsize},enlarge x limits=false,ylabel shift=-3pt,xlabel shift=-3pt]\addplot[fill=ink!55,draw=none] coordinates {(0,0.01489) (1,0.127) (2,0.5584) (3,2.344) (4,5.884) (5,6.561) (6,5.45) (7,2.725) (8,0.7439) (9,0.1566) (10,0.05356) (11,0.0001779) (12,0.0459) (13,3.098e-06) (14,0.009267) (15,6.8e-07) (16,9.47e-05) (17,3.4e-08) (18,1.02e-07)};\draw[densely dotted] (axis cs:16,0) -- (axis cs:16,7.217);\end{axis}\begin{axis}[width=72mm,height=42mm,axis y line*=right,axis x line=none,ymode=log,xmin=-1,xmax=19,ylabel={\scriptsize eindeutig (log)},tick label style={font=\tiny},label style={font=\scriptsize},ylabel shift=-3pt]\addplot[mark=*,mark size=0.7pt,draw=black,thin] coordinates {(0,14885567) (1,101190240) (2,295708850) (3,403540736) (4,434420486) (5,159993589) (6,39230653) (7,7752962) (8,1286861) (9,10832) (10,394206) (11,47) (12,71203) (14,1286) (16,2)};\end{axis}\end{tikzpicture}\end{minipage}
\vspace{3mm}

\noindent\begin{minipage}[t]{60mm}\vspace{0pt}\sbox\hmboard{\chessboard[setfen=1r6/K7/8/8/2b5/2P5/7Q/7k b - - 0 1]}\usebox\hmboard\\[1mm]
\makebox[\wd\hmboard]{\sffamily\hspace{3mm}\textbf{h\#8}\hfill(3+3)}\end{minipage}\hfill
\begin{minipage}[t]{52mm}\vspace{0pt}
\textcolor{ink}{\sffamily\textsc{Lösung}}\par\smallskip\small
\begin{tabular}{@{}r@{\enspace}l@{\enspace}l@{}}
\textbf{1.} & \symking{}xh2 & \symking{}xb8\\
\textbf{2.} & \symbishop{}d5 & c4\\
\textbf{3.} & \symking{}g3 & cxd5\\
\textbf{4.} & \symking{}f4 & d6\\
\textbf{5.} & \symking{}e5 & d7\\
\textbf{6.} & \symking{}d6 & \symking{}a7\\
\textbf{7.} & \symking{}c7 & \symking{}a6\\
\textbf{8.} & \symking{}b8 & d8=\symqueen{}\#\\
\end{tabular}\par\vspace{3mm}
\qrcode[height=12mm]{https://helpman.komtera.lt/?fen=1r6/K7/8/8/2b5/2P5/7Q/7k&moves=8}\\{\sffamily\tiny\color{ink}scannen zum Lösen}
\end{minipage}
\vspace{2mm}

{\raggedright\noindent\textcolor{ink}{\sffamily\textsc{Themen}}\enspace \themetag{reines Matt} \themetag{Mustermatt} \themetag{Idealmatt} \themetag{Spiegelmatt} \themetag{Umwandlung} \themetag{Switchback} \themetag{Kniest} \themetag{Phönix} \themetag{ohne Schach} \themetag{Umnov}\par\smallskip{\footnotesize\textit{Hinweis:} die Partei am Zug steht im Diagramm im Schach; die Lösung beginnt mit einem Schlag.}\par}

\clearpage
\phantomsection\label{hm:KQPvkrn}
\begin{tcolorbox}[enhanced,colback=ink,colframe=ink,arc=0pt,boxrule=0pt,left=4mm,right=4mm,top=2.5mm,bottom=2.5mm]
\color{white}\sffamily{\LARGE\bfseries KQPvkrn}\hfill{\Large h\#8}\\[0.5mm]
{\small 6 Steine\quad Schwarz am Zug}\hfill{\small Nr.~412\qquad\href{https://helpman.komtera.lt/?fen=7k/7Q/8/8/8/6rn/K6P/8&moves=8}{\color{white}\bfseries lösen\,\ensuremath{\nearrow}}}
\end{tcolorbox}
\vspace{2mm}
\noindent\begin{minipage}[t]{36mm}\vspace{0pt}\sffamily\small\raggedright
\hmfigure{längste eindeutige}{{\Large h\#8}}\hmfigure{eindeutig bei dieser Länge}{{\Large 20} Stellungen}\hmfigure{längstes Matt}{{\Large h\#9}}\hmfigure{in der Klasse}{45,4 Milliarden Stellungen,\\ 25,3 Milliarden mit Hilfsmatt}
\end{minipage}\hfill
\begin{minipage}[t]{80mm}\vspace{0pt}\begin{tikzpicture}\begin{axis}[width=72mm,height=42mm,ybar,bar width=1.6pt,xmin=-1,xmax=19,ymin=0,ymax=6.843,axis y line*=left,axis x line*=bottom,xlabel={\scriptsize Halbzüge bis Matt},ylabel={\scriptsize Stellungen (Mrd.)},tick label style={font=\tiny},label style={font=\scriptsize},enlarge x limits=false,ylabel shift=-3pt,xlabel shift=-3pt]\addplot[fill=ink!55,draw=none] coordinates {(0,0.0164) (1,0.1363) (2,0.561) (3,2.267) (4,5.621) (5,6.221) (6,5.838) (7,2.996) (8,1.214) (9,0.2947) (10,0.07671) (11,0.001311) (12,0.0481) (13,8.796e-06) (14,0.01825) (15,4.415e-06) (16,0.0004066) (17,4.1e-07) (18,9.74e-07)};\draw[densely dotted] (axis cs:16,0) -- (axis cs:16,6.843);\end{axis}\begin{axis}[width=72mm,height=42mm,axis y line*=right,axis x line=none,ymode=log,xmin=-1,xmax=19,ylabel={\scriptsize eindeutig (log)},tick label style={font=\tiny},label style={font=\scriptsize},ylabel shift=-3pt]\addplot[mark=*,mark size=0.7pt,draw=black,thin] coordinates {(0,16400258) (1,104367706) (2,284836190) (3,376894412) (4,433742007) (5,177879718) (6,63860582) (7,13739594) (8,2785648) (9,124560) (10,361502) (11,70) (12,134626) (13,14) (14,4975) (16,20)};\end{axis}\end{tikzpicture}\end{minipage}
\vspace{3mm}

\noindent\begin{minipage}[t]{60mm}\vspace{0pt}\sbox\hmboard{\chessboard[setfen=7k/7Q/8/8/8/6rn/K6P/8 b - - 0 1]}\usebox\hmboard\\[1mm]
\makebox[\wd\hmboard]{\sffamily\hspace{3mm}\textbf{h\#8}\hfill(3+3)}\end{minipage}\hfill
\begin{minipage}[t]{52mm}\vspace{0pt}
\textcolor{ink}{\sffamily\textsc{Lösung}}\par\smallskip\small
\begin{tabular}{@{}r@{\enspace}l@{\enspace}l@{}}
\textbf{1.} & \symking{}xh7 & hxg3\\
\textbf{2.} & \symking{}g6 & \symking{}a3\\
\textbf{3.} & \symking{}f5 & g4+\\
\textbf{4.} & \symking{}e4 & g5\\
\textbf{5.} & \symking{}d3 & g6\\
\textbf{6.} & \symking{}c2 & g7\\
\textbf{7.} & \symking{}b1 & g8=\symqueen{}\\
\textbf{8.} & \symking{}a1 & \symqueen{}a2\#\\
\end{tabular}\par\vspace{3mm}
\qrcode[height=12mm]{https://helpman.komtera.lt/?fen=7k/7Q/8/8/8/6rn/K6P/8&moves=8}\\{\sffamily\tiny\color{ink}scannen zum Lösen}
\end{minipage}
\vspace{2mm}

{\raggedright\noindent\textcolor{ink}{\sffamily\textsc{Themen}}\enspace \themetag{Umwandlung} \themetag{Excelsior} \themetag{ein Stein zieht} \themetag{Phönix}\par\smallskip{\footnotesize\textit{Hinweis:} die Partei am Zug steht im Diagramm im Schach; die Lösung beginnt mit einem Schlag.}\par}

\clearpage
\phantomsection\label{hm:KRBBvkn}
\begin{tcolorbox}[enhanced,colback=ink,colframe=ink,arc=0pt,boxrule=0pt,left=4mm,right=4mm,top=2.5mm,bottom=2.5mm]
\color{white}\sffamily{\LARGE\bfseries KRBBvkn}\hfill{\Large h\#8}\\[0.5mm]
{\small 6 Steine\quad Schwarz am Zug}\hfill{\small Nr.~413\qquad\href{https://helpman.komtera.lt/?fen=6n1/7k/7R/8/3B4/8/8/B1K5&moves=8}{\color{white}\bfseries lösen\,\ensuremath{\nearrow}}}
\end{tcolorbox}
\vspace{2mm}
\noindent\begin{minipage}[t]{36mm}\vspace{0pt}\sffamily\small\raggedright
\hmfigure{längste eindeutige}{{\Large h\#8}}\hmfigure{eindeutig bei dieser Länge}{{\Large 8} Stellungen}\hmfigure{längstes Matt}{{\Large h\#12}}\hmfigure{in der Klasse}{15,5 Milliarden Stellungen,\\ 4,61 Milliarden mit Hilfsmatt}
\end{minipage}\hfill
\begin{minipage}[t]{80mm}\vspace{0pt}\begin{tikzpicture}\begin{axis}[width=72mm,height=42mm,ybar,bar width=1.6pt,xmin=-1,xmax=25,ymin=0,ymax=1.399,axis y line*=left,axis x line*=bottom,xlabel={\scriptsize Halbzüge bis Matt},ylabel={\scriptsize Stellungen (Mrd.)},tick label style={font=\tiny},label style={font=\scriptsize},enlarge x limits=false,ylabel shift=-3pt,xlabel shift=-3pt]\addplot[fill=ink!55,draw=none] coordinates {(0,0.006809) (1,0.02594) (2,0.09437) (3,0.3836) (4,1.106) (5,1.005) (6,1.272) (7,0.3904) (8,0.2818) (9,0.03548) (10,0.01061) (11,0.0005138) (12,0.0003767) (13,3.3e-08) (14,0.0001979) (16,0.0002038) (18,0.0002712) (20,0.0002887) (22,0.0001696) (24,2.337e-05)};\draw[densely dotted] (axis cs:16,0) -- (axis cs:16,1.399);\end{axis}\begin{axis}[width=72mm,height=42mm,axis y line*=right,axis x line=none,ymode=log,xmin=-1,xmax=25,ylabel={\scriptsize eindeutig (log)},tick label style={font=\tiny},label style={font=\scriptsize},ylabel shift=-3pt]\addplot[mark=*,mark size=0.7pt,draw=black,thin] coordinates {(0,6809246) (1,24043569) (2,59520632) (3,139827840) (4,196301133) (5,52620718) (6,34518549) (7,1891108) (8,381536) (9,10401) (10,4434) (11,1) (12,336) (14,38) (16,8)};\end{axis}\end{tikzpicture}\end{minipage}
\vspace{3mm}

\noindent\begin{minipage}[t]{60mm}\vspace{0pt}\sbox\hmboard{\chessboard[setfen=6n1/7k/7R/8/3B4/8/8/B1K5 b - - 0 1]}\usebox\hmboard\\[1mm]
\makebox[\wd\hmboard]{\sffamily\hspace{3mm}\textbf{h\#8}\hfill(4+2)}\end{minipage}\hfill
\begin{minipage}[t]{52mm}\vspace{0pt}
\textcolor{ink}{\sffamily\textsc{Lösung}}\par\smallskip\small
\begin{tabular}{@{}r@{\enspace}l@{\enspace}l@{}}
\textbf{1.} & \symknight{}xh6 & \symking{}d2\\
\textbf{2.} & \symking{}g8 & \symking{}e3\\
\textbf{3.} & \symknight{}f5+ & \symking{}f4\\
\textbf{4.} & \symknight{}xd4 & \symking{}g5\\
\textbf{5.} & \symknight{}f5 & \symbishop{}g7\\
\textbf{6.} & \symknight{}h6 & \symbishop{}f8\\
\textbf{7.} & \symking{}h8 & \symking{}g6\\
\textbf{8.} & \symknight{}g8 & \symbishop{}g7\#\\
\end{tabular}\par\vspace{3mm}
\qrcode[height=12mm]{https://helpman.komtera.lt/?fen=6n1/7k/7R/8/3B4/8/8/B1K5&moves=8}\\{\sffamily\tiny\color{ink}scannen zum Lösen}
\end{minipage}
\vspace{2mm}

{\raggedright\noindent\textcolor{ink}{\sffamily\textsc{Themen}}\enspace \themetag{reines Matt} \themetag{Mustermatt} \themetag{Idealmatt} \themetag{Switchback} \themetag{Selbstblock}\par\smallskip{\footnotesize\textit{Hinweis:} die Partei am Zug steht im Diagramm im Schach; die Lösung beginnt mit einem Schlag.}\par}

\clearpage
\phantomsection\label{hm:KRBvkbb}
\begin{tcolorbox}[enhanced,colback=ink,colframe=ink,arc=0pt,boxrule=0pt,left=4mm,right=4mm,top=2.5mm,bottom=2.5mm]
\color{white}\sffamily{\LARGE\bfseries KRBvkbb}\hfill{\Large h\#8}\\[0.5mm]
{\small 6 Steine\quad Schwarz am Zug}\hfill{\small Nr.~414\qquad\href{https://helpman.komtera.lt/?fen=8/8/8/4b3/4b2B/8/8/3K1kR1&moves=8}{\color{white}\bfseries lösen\,\ensuremath{\nearrow}}}
\end{tcolorbox}
\vspace{2mm}
\noindent\begin{minipage}[t]{36mm}\vspace{0pt}\sffamily\small\raggedright
\hmfigure{längste eindeutige}{{\Large h\#8}}\hmfigure{eindeutig bei dieser Länge}{{\Large 5} Stellungen}\hmfigure{längstes Matt}{{\Large h\#10}}\hmfigure{in der Klasse}{15,5 Milliarden Stellungen,\\ 4,4 Milliarden mit Hilfsmatt}
\end{minipage}\hfill
\begin{minipage}[t]{80mm}\vspace{0pt}\begin{tikzpicture}\begin{axis}[width=72mm,height=42mm,ybar,bar width=1.6pt,xmin=-1,xmax=21,ymin=0,ymax=1.433,axis y line*=left,axis x line*=bottom,xlabel={\scriptsize Halbzüge bis Matt},ylabel={\scriptsize Stellungen (Mrd.)},tick label style={font=\tiny},label style={font=\scriptsize},enlarge x limits=false,ylabel shift=-3pt,xlabel shift=-3pt]\addplot[fill=ink!55,draw=none] coordinates {(0,0.001576) (1,0.007327) (2,0.03057) (3,0.1614) (4,0.5247) (5,1.021) (6,1.303) (7,0.7706) (8,0.3345) (9,0.142) (10,0.07856) (11,0.01774) (12,0.007756) (13,0.000246) (14,0.0002495) (15,8.8e-08) (16,0.0002217) (18,0.0002148) (20,5.313e-05)};\draw[densely dotted] (axis cs:16,0) -- (axis cs:16,1.433);\end{axis}\begin{axis}[width=72mm,height=42mm,axis y line*=right,axis x line=none,ymode=log,xmin=-1,xmax=21,ylabel={\scriptsize eindeutig (log)},tick label style={font=\tiny},label style={font=\scriptsize},ylabel shift=-3pt]\addplot[mark=*,mark size=0.7pt,draw=black,thin] coordinates {(0,1575700) (1,7155219) (2,17078028) (3,61710722) (4,141195420) (5,76782912) (6,22841183) (7,4220708) (8,1659165) (9,314649) (10,81631) (11,855) (12,635) (14,305) (16,5)};\end{axis}\end{tikzpicture}\end{minipage}
\vspace{3mm}

\noindent\begin{minipage}[t]{60mm}\vspace{0pt}\sbox\hmboard{\chessboard[setfen=8/8/8/4b3/4b2B/8/8/3K1kR1 b - - 0 1]}\usebox\hmboard\\[1mm]
\makebox[\wd\hmboard]{\sffamily\hspace{3mm}\textbf{h\#8}\hfill(3+3)}\end{minipage}\hfill
\begin{minipage}[t]{52mm}\vspace{0pt}
\textcolor{ink}{\sffamily\textsc{Lösung}}\par\smallskip\small
\begin{tabular}{@{}r@{\enspace}l@{\enspace}l@{}}
\textbf{1.} & \symking{}xg1 & \symbishop{}f6\\
\textbf{2.} & \symking{}f2 & \symbishop{}xe5\\
\textbf{3.} & \symking{}e3 & \symking{}e1\\
\textbf{4.} & \symking{}d3 & \symking{}f2\\
\textbf{5.} & \symking{}c2 & \symking{}e3\\
\textbf{6.} & \symking{}b1 & \symking{}d4\\
\textbf{7.} & \symking{}a1 & \symking{}c3\\
\textbf{8.} & \symbishop{}b1 & \symking{}b3\#\\
\end{tabular}\par\vspace{3mm}
\qrcode[height=12mm]{https://helpman.komtera.lt/?fen=8/8/8/4b3/4b2B/8/8/3K1kR1&moves=8}\\{\sffamily\tiny\color{ink}scannen zum Lösen}
\end{minipage}
\vspace{2mm}

{\raggedright\noindent\textcolor{ink}{\sffamily\textsc{Themen}}\enspace \themetag{Selbstblock} \themetag{ohne Schach}\par\smallskip{\footnotesize\textit{Hinweis:} die Partei am Zug steht im Diagramm im Schach; die Lösung beginnt mit einem Schlag.}\par}

\clearpage
\phantomsection\label{hm:KRNvkbb}
\begin{tcolorbox}[enhanced,colback=ink,colframe=ink,arc=0pt,boxrule=0pt,left=4mm,right=4mm,top=2.5mm,bottom=2.5mm]
\color{white}\sffamily{\LARGE\bfseries KRNvkbb}\hfill{\Large h\#8}\\[0.5mm]
{\small 6 Steine\quad Schwarz am Zug}\hfill{\small Nr.~415\qquad\href{https://helpman.komtera.lt/?fen=8/8/Nb6/8/8/2b5/R7/k1K5&moves=8}{\color{white}\bfseries lösen\,\ensuremath{\nearrow}}}
\end{tcolorbox}
\vspace{2mm}
\noindent\begin{minipage}[t]{36mm}\vspace{0pt}\sffamily\small\raggedright
\hmfigure{längste eindeutige}{{\Large h\#8}}\hmfigure{eindeutig bei dieser Länge}{{\Large 20} Stellungen}\hmfigure{längstes Matt}{{\Large h\#10}}\hmfigure{in der Klasse}{15,5 Milliarden Stellungen,\\ 4,51 Milliarden mit Hilfsmatt}
\end{minipage}\hfill
\begin{minipage}[t]{80mm}\vspace{0pt}\begin{tikzpicture}\begin{axis}[width=72mm,height=42mm,ybar,bar width=1.6pt,xmin=-1,xmax=21,ymin=0,ymax=1.533,axis y line*=left,axis x line*=bottom,xlabel={\scriptsize Halbzüge bis Matt},ylabel={\scriptsize Stellungen (Mrd.)},tick label style={font=\tiny},label style={font=\scriptsize},enlarge x limits=false,ylabel shift=-3pt,xlabel shift=-3pt]\addplot[fill=ink!55,draw=none] coordinates {(0,0.00112) (1,0.006415) (2,0.02424) (3,0.1352) (4,0.3803) (5,0.9653) (6,1.394) (7,1.034) (8,0.472) (9,0.08677) (10,0.009186) (11,0.0002598) (12,0.0005819) (14,0.0003573) (16,0.0002888) (18,0.0001765) (20,2.112e-05)};\draw[densely dotted] (axis cs:16,0) -- (axis cs:16,1.533);\end{axis}\begin{axis}[width=72mm,height=42mm,axis y line*=right,axis x line=none,ymode=log,xmin=-1,xmax=21,ylabel={\scriptsize eindeutig (log)},tick label style={font=\tiny},label style={font=\scriptsize},ylabel shift=-3pt]\addplot[mark=*,mark size=0.7pt,draw=black,thin] coordinates {(0,1119589) (1,6322748) (2,12402830) (3,47142013) (4,102161652) (5,87339653) (6,33912665) (7,4368258) (8,327178) (9,7485) (10,3269) (12,1503) (14,63) (16,20)};\end{axis}\end{tikzpicture}\end{minipage}
\vspace{3mm}

\noindent\begin{minipage}[t]{60mm}\vspace{0pt}\sbox\hmboard{\chessboard[setfen=8/8/Nb6/8/8/2b5/R7/k1K5 b - - 0 1]}\usebox\hmboard\\[1mm]
\makebox[\wd\hmboard]{\sffamily\hspace{3mm}\textbf{h\#8}\hfill(3+3)}\end{minipage}\hfill
\begin{minipage}[t]{52mm}\vspace{0pt}
\textcolor{ink}{\sffamily\textsc{Lösung}}\par\smallskip\small
\begin{tabular}{@{}r@{\enspace}l@{\enspace}l@{}}
\textbf{1.} & \symking{}xa2 & \symking{}c2\\
\textbf{2.} & \symking{}a3 & \symking{}d3\\
\textbf{3.} & \symking{}a4 & \symking{}e4\\
\textbf{4.} & \symking{}b5 & \symking{}f5\\
\textbf{5.} & \symking{}c6 & \symking{}e6\\
\textbf{6.} & \symking{}b7 & \symking{}d7\\
\textbf{7.} & \symking{}a8 & \symking{}c8\\
\textbf{8.} & \symbishop{}a7 & \symknight{}c7\#\\
\end{tabular}\par\vspace{3mm}
\qrcode[height=12mm]{https://helpman.komtera.lt/?fen=8/8/Nb6/8/8/2b5/R7/k1K5&moves=8}\\{\sffamily\tiny\color{ink}scannen zum Lösen}
\end{minipage}
\vspace{2mm}

{\raggedright\noindent\textcolor{ink}{\sffamily\textsc{Themen}}\enspace \themetag{reines Matt} \themetag{Mustermatt} \themetag{Selbstblock} \themetag{ohne Schach}\par\smallskip{\footnotesize\textit{Hinweis:} die Partei am Zug steht im Diagramm im Schach; die Lösung beginnt mit einem Schlag.}\par}

\clearpage
\phantomsection\label{hm:KRNvkbp}
\begin{tcolorbox}[enhanced,colback=ink,colframe=ink,arc=0pt,boxrule=0pt,left=4mm,right=4mm,top=2.5mm,bottom=2.5mm]
\color{white}\sffamily{\LARGE\bfseries KRNvkbp}\hfill{\Large h\#8}\\[0.5mm]
{\small 6 Steine\quad Schwarz am Zug}\hfill{\small Nr.~416\qquad\href{https://helpman.komtera.lt/?fen=7R/6pk/8/8/7N/b7/8/K7&moves=8}{\color{white}\bfseries lösen\,\ensuremath{\nearrow}}}
\end{tcolorbox}
\vspace{2mm}
\noindent\begin{minipage}[t]{36mm}\vspace{0pt}\sffamily\small\raggedright
\hmfigure{längste eindeutige}{{\Large h\#8}}\hmfigure{eindeutig bei dieser Länge}{{\Large 36} Stellungen}\hmfigure{längstes Matt}{{\Large h\#10}}\hmfigure{in der Klasse}{45,4 Milliarden Stellungen,\\ 28,8 Milliarden mit Hilfsmatt}
\end{minipage}\hfill
\begin{minipage}[t]{80mm}\vspace{0pt}\begin{tikzpicture}\begin{axis}[width=72mm,height=42mm,ybar,bar width=1.6pt,xmin=-1,xmax=21,ymin=0,ymax=10.02,axis y line*=left,axis x line*=bottom,xlabel={\scriptsize Halbzüge bis Matt},ylabel={\scriptsize Stellungen (Mrd.)},tick label style={font=\tiny},label style={font=\scriptsize},enlarge x limits=false,ylabel shift=-3pt,xlabel shift=-3pt]\addplot[fill=ink!55,draw=none] coordinates {(0,0.0112) (1,0.06324) (2,0.2002) (3,0.9622) (4,2.605) (5,5.821) (6,9.112) (7,6.045) (8,3.38) (9,0.5006) (10,0.0671) (11,0.0009803) (12,0.00424) (13,1.92e-07) (14,0.003) (15,7.9e-08) (16,0.001575) (17,5e-09) (18,0.0002515) (19,2e-09) (20,3.27e-07)};\draw[densely dotted] (axis cs:16,0) -- (axis cs:16,10.02);\end{axis}\begin{axis}[width=72mm,height=42mm,axis y line*=right,axis x line=none,ymode=log,xmin=-1,xmax=21,ylabel={\scriptsize eindeutig (log)},tick label style={font=\tiny},label style={font=\scriptsize},ylabel shift=-3pt]\addplot[mark=*,mark size=0.7pt,draw=black,thin] coordinates {(0,11203269) (1,62231168) (2,107543443) (3,344057596) (4,731487254) (5,551333747) (6,273243117) (7,27259850) (8,2708109) (9,31474) (10,21112) (11,1) (12,6950) (14,866) (16,36)};\end{axis}\end{tikzpicture}\end{minipage}
\vspace{3mm}

\noindent\begin{minipage}[t]{60mm}\vspace{0pt}\sbox\hmboard{\chessboard[setfen=7R/6pk/8/8/7N/b7/8/K7 b - - 0 1]}\usebox\hmboard\\[1mm]
\makebox[\wd\hmboard]{\sffamily\hspace{3mm}\textbf{h\#8}\hfill(3+3)}\end{minipage}\hfill
\begin{minipage}[t]{52mm}\vspace{0pt}
\textcolor{ink}{\sffamily\textsc{Lösung}}\par\smallskip\small
\begin{tabular}{@{}r@{\enspace}l@{\enspace}l@{}}
\textbf{1.} & \symking{}xh8 & \symking{}a2\\
\textbf{2.} & g5 & \symking{}xa3\\
\textbf{3.} & g4 & \symking{}b4\\
\textbf{4.} & g3 & \symking{}c5\\
\textbf{5.} & g2 & \symking{}d6\\
\textbf{6.} & g1=\symrook{} & \symking{}e7\\
\textbf{7.} & \symrook{}g7+ & \symking{}f8\\
\textbf{8.} & \symrook{}h7 & \symknight{}g6\#\\
\end{tabular}\par\vspace{3mm}
\qrcode[height=12mm]{https://helpman.komtera.lt/?fen=7R/6pk/8/8/7N/b7/8/K7&moves=8}\\{\sffamily\tiny\color{ink}scannen zum Lösen}
\end{minipage}
\vspace{2mm}

{\raggedright\noindent\textcolor{ink}{\sffamily\textsc{Themen}}\enspace \themetag{reines Matt} \themetag{Mustermatt} \themetag{Idealmatt} \themetag{Umwandlung} \themetag{Unterverwandlung} \themetag{Excelsior} \themetag{Rundlauf} \themetag{Selbstblock} \themetag{Kniest}\par\smallskip{\footnotesize\textit{Hinweis:} die Partei am Zug steht im Diagramm im Schach; die Lösung beginnt mit einem Schlag.}\par}

\clearpage
\phantomsection\label{hm:KRNvknp}
\begin{tcolorbox}[enhanced,colback=ink,colframe=ink,arc=0pt,boxrule=0pt,left=4mm,right=4mm,top=2.5mm,bottom=2.5mm]
\color{white}\sffamily{\LARGE\bfseries KRNvknp}\hfill{\Large h\#8}\\[0.5mm]
{\small 6 Steine\quad Schwarz am Zug}\hfill{\small Nr.~417\qquad\href{https://helpman.komtera.lt/?fen=8/8/5NRk/8/8/2p5/n7/K7&moves=8}{\color{white}\bfseries lösen\,\ensuremath{\nearrow}}}
\end{tcolorbox}
\vspace{2mm}
\noindent\begin{minipage}[t]{36mm}\vspace{0pt}\sffamily\small\raggedright
\hmfigure{längste eindeutige}{{\Large h\#8}}\hmfigure{eindeutig bei dieser Länge}{{\Large 55} Stellungen}\hmfigure{längstes Matt}{{\Large h\#10.5}}\hmfigure{in der Klasse}{45,4 Milliarden Stellungen,\\ 29,6 Milliarden mit Hilfsmatt}
\end{minipage}\hfill
\begin{minipage}[t]{80mm}\vspace{0pt}\begin{tikzpicture}\begin{axis}[width=72mm,height=42mm,ybar,bar width=1.6pt,xmin=-1,xmax=22,ymin=0,ymax=9.821,axis y line*=left,axis x line*=bottom,xlabel={\scriptsize Halbzüge bis Matt},ylabel={\scriptsize Stellungen (Mrd.)},tick label style={font=\tiny},label style={font=\scriptsize},enlarge x limits=false,ylabel shift=-3pt,xlabel shift=-3pt]\addplot[fill=ink!55,draw=none] coordinates {(0,0.01444) (1,0.08201) (2,0.2237) (3,0.9775) (4,2.534) (5,5.299) (6,8.928) (7,6.251) (8,4.344) (9,0.781) (10,0.1404) (11,0.002925) (12,0.005046) (13,8.9e-08) (14,0.003281) (15,1.74e-07) (16,0.001025) (17,1.65e-07) (18,9.695e-05) (19,8.7e-08) (20,1.119e-06) (21,3.3e-08)};\draw[densely dotted] (axis cs:16,0) -- (axis cs:16,9.821);\end{axis}\begin{axis}[width=72mm,height=42mm,axis y line*=right,axis x line=none,ymode=log,xmin=-1,xmax=22,ylabel={\scriptsize eindeutig (log)},tick label style={font=\tiny},label style={font=\scriptsize},ylabel shift=-3pt]\addplot[mark=*,mark size=0.7pt,draw=black,thin] coordinates {(0,14440384) (1,80496101) (2,123733729) (3,359184106) (4,716529390) (5,529667632) (6,311158506) (7,36233967) (8,4885175) (9,83101) (10,14310) (11,4) (12,4933) (14,185) (16,55)};\end{axis}\end{tikzpicture}\end{minipage}
\vspace{3mm}

\noindent\begin{minipage}[t]{60mm}\vspace{0pt}\sbox\hmboard{\chessboard[setfen=8/8/5NRk/8/8/2p5/n7/K7 b - - 0 1]}\usebox\hmboard\\[1mm]
\makebox[\wd\hmboard]{\sffamily\hspace{3mm}\textbf{h\#8}\hfill(3+3)}\end{minipage}\hfill
\begin{minipage}[t]{52mm}\vspace{0pt}
\textcolor{ink}{\sffamily\textsc{Lösung}}\par\smallskip\small
\begin{tabular}{@{}r@{\enspace}l@{\enspace}l@{}}
\textbf{1.} & \symking{}xg6 & \symknight{}d5\\
\textbf{2.} & c2 & \symking{}b2\\
\textbf{3.} & c1=\symbishop{}+ & \symking{}c2\\
\textbf{4.} & \symknight{}c3 & \symking{}d3\\
\textbf{5.} & \symknight{}e4 & \symking{}e2\\
\textbf{6.} & \symbishop{}h6 & \symking{}f3\\
\textbf{7.} & \symknight{}g5+ & \symking{}g3\\
\textbf{8.} & \symking{}h5 & \symknight{}f4\#\\
\end{tabular}\par\vspace{3mm}
\qrcode[height=12mm]{https://helpman.komtera.lt/?fen=8/8/5NRk/8/8/2p5/n7/K7&moves=8}\\{\sffamily\tiny\color{ink}scannen zum Lösen}
\end{minipage}
\vspace{2mm}

{\raggedright\noindent\textcolor{ink}{\sffamily\textsc{Themen}}\enspace \themetag{reines Matt} \themetag{Mustermatt} \themetag{Idealmatt} \themetag{Umwandlung} \themetag{Unterverwandlung} \themetag{Selbstblock} \themetag{Umnov}\par\smallskip{\footnotesize\textit{Hinweis:} die Partei am Zug steht im Diagramm im Schach; die Lösung beginnt mit einem Schlag.}\par}

\clearpage
\phantomsection\label{hm:KRNvkpp}
\begin{tcolorbox}[enhanced,colback=ink,colframe=ink,arc=0pt,boxrule=0pt,left=4mm,right=4mm,top=2.5mm,bottom=2.5mm]
\color{white}\sffamily{\LARGE\bfseries KRNvkpp}\hfill{\Large h\#8}\\[0.5mm]
{\small 6 Steine\quad Schwarz am Zug}\hfill{\small Nr.~418\qquad\href{https://helpman.komtera.lt/?fen=7k/6R1/8/8/6N1/p7/p7/K7&moves=8}{\color{white}\bfseries lösen\,\ensuremath{\nearrow}}}
\end{tcolorbox}
\vspace{2mm}
\noindent\begin{minipage}[t]{36mm}\vspace{0pt}\sffamily\small\raggedright
\hmfigure{längste eindeutige}{{\Large h\#8}}\hmfigure{eindeutig bei dieser Länge}{{\Large 15} Stellungen}\hmfigure{längstes Matt}{{\Large h\#11.5}}\hmfigure{in der Klasse}{34,1 Milliarden Stellungen,\\ 11,4 Milliarden mit Hilfsmatt}
\end{minipage}\hfill
\begin{minipage}[t]{80mm}\vspace{0pt}\begin{tikzpicture}\begin{axis}[width=72mm,height=42mm,ybar,bar width=1.6pt,xmin=-1,xmax=24,ymin=0,ymax=3.69,axis y line*=left,axis x line*=bottom,xlabel={\scriptsize Halbzüge bis Matt},ylabel={\scriptsize Stellungen (Mrd.)},tick label style={font=\tiny},label style={font=\scriptsize},enlarge x limits=false,ylabel shift=-3pt,xlabel shift=-3pt]\addplot[fill=ink!55,draw=none] coordinates {(0,0.007149) (1,0.03912) (2,0.09683) (3,0.4025) (4,1.027) (5,1.994) (6,3.355) (7,2.228) (8,1.795) (9,0.3385) (10,0.1097) (11,0.003131) (12,0.002211) (13,4.63e-07) (14,0.001453) (15,3.39e-07) (16,0.000893) (17,2.26e-07) (18,0.0003998) (19,3.3e-08) (20,0.0001432) (21,1.2e-08) (22,3.121e-05) (23,8e-09)};\draw[densely dotted] (axis cs:16,0) -- (axis cs:16,3.69);\end{axis}\begin{axis}[width=72mm,height=42mm,axis y line*=right,axis x line=none,ymode=log,xmin=-1,xmax=24,ylabel={\scriptsize eindeutig (log)},tick label style={font=\tiny},label style={font=\scriptsize},ylabel shift=-3pt]\addplot[mark=*,mark size=0.7pt,draw=black,thin] coordinates {(0,7148602) (1,38393218) (2,60298674) (3,159397828) (4,299082916) (5,195905673) (6,125172069) (7,13993396) (8,3206548) (9,56863) (10,14702) (11,24) (12,2104) (13,5) (14,247) (16,15)};\end{axis}\end{tikzpicture}\end{minipage}
\vspace{3mm}

\noindent\begin{minipage}[t]{60mm}\vspace{0pt}\sbox\hmboard{\chessboard[setfen=7k/6R1/8/8/6N1/p7/p7/K7 b - - 0 1]}\usebox\hmboard\\[1mm]
\makebox[\wd\hmboard]{\sffamily\hspace{3mm}\textbf{h\#8}\hfill(3+3)}\end{minipage}\hfill
\begin{minipage}[t]{52mm}\vspace{0pt}
\textcolor{ink}{\sffamily\textsc{Lösung}}\par\smallskip\small
\begin{tabular}{@{}r@{\enspace}l@{\enspace}l@{}}
\textbf{1.} & \symking{}xg7 & \symknight{}e3\\
\textbf{2.} & \symking{}f6 & \symking{}xa2\\
\textbf{3.} & \symking{}e5 & \symking{}b1\\
\textbf{4.} & \symking{}d4 & \symking{}c1\\
\textbf{5.} & \symking{}c3 & \symking{}d1\\
\textbf{6.} & \symking{}b2 & \symking{}d2\\
\textbf{7.} & \symking{}a1 & \symking{}c1\\
\textbf{8.} & a2 & \symknight{}c2\#\\
\end{tabular}\par\vspace{3mm}
\qrcode[height=12mm]{https://helpman.komtera.lt/?fen=7k/6R1/8/8/6N1/p7/p7/K7&moves=8}\\{\sffamily\tiny\color{ink}scannen zum Lösen}
\end{minipage}
\vspace{2mm}

{\raggedright\noindent\textcolor{ink}{\sffamily\textsc{Themen}}\enspace \themetag{reines Matt} \themetag{Mustermatt} \themetag{Idealmatt} \themetag{Rundlauf} \themetag{Selbstblock} \themetag{ohne Schach}\par\smallskip{\footnotesize\textit{Hinweis:} die Lösung beginnt mit einem Schlag.}\par}

\clearpage
\phantomsection\label{hm:KRNvkrp}
\begin{tcolorbox}[enhanced,colback=ink,colframe=ink,arc=0pt,boxrule=0pt,left=4mm,right=4mm,top=2.5mm,bottom=2.5mm]
\color{white}\sffamily{\LARGE\bfseries KRNvkrp}\hfill{\Large h\#8}\\[0.5mm]
{\small 6 Steine\quad Schwarz am Zug}\hfill{\small Nr.~419\qquad\href{https://helpman.komtera.lt/?fen=8/7R/4N1pk/8/1r6/8/8/K7&moves=8}{\color{white}\bfseries lösen\,\ensuremath{\nearrow}}}
\end{tcolorbox}
\vspace{2mm}
\noindent\begin{minipage}[t]{36mm}\vspace{0pt}\sffamily\small\raggedright
\hmfigure{längste eindeutige}{{\Large h\#8}}\hmfigure{eindeutig bei dieser Länge}{{\Large 39} Stellungen}\hmfigure{längstes Matt}{{\Large h\#11.5}}\hmfigure{in der Klasse}{45,4 Milliarden Stellungen,\\ 27,6 Milliarden mit Hilfsmatt}
\end{minipage}\hfill
\begin{minipage}[t]{80mm}\vspace{0pt}\begin{tikzpicture}\begin{axis}[width=72mm,height=42mm,ybar,bar width=1.6pt,xmin=-1,xmax=24,ymin=0,ymax=9.991,axis y line*=left,axis x line*=bottom,xlabel={\scriptsize Halbzüge bis Matt},ylabel={\scriptsize Stellungen (Mrd.)},tick label style={font=\tiny},label style={font=\scriptsize},enlarge x limits=false,ylabel shift=-3pt,xlabel shift=-3pt]\addplot[fill=ink!55,draw=none] coordinates {(0,0.009638) (1,0.05627) (2,0.2153) (3,1.067) (4,3.177) (5,6.678) (6,9.083) (7,5.396) (8,1.673) (9,0.1896) (10,0.008755) (11,0.0001879) (12,0.004046) (13,8.43e-07) (14,0.001274) (15,2.97e-07) (16,0.0003389) (17,1.93e-07) (18,2.965e-05) (19,1.92e-07) (20,4e-08) (21,2.09e-07) (23,1.61e-07)};\draw[densely dotted] (axis cs:16,0) -- (axis cs:16,9.991);\end{axis}\begin{axis}[width=72mm,height=42mm,axis y line*=right,axis x line=none,ymode=log,xmin=-1,xmax=24,ylabel={\scriptsize eindeutig (log)},tick label style={font=\tiny},label style={font=\scriptsize},ylabel shift=-3pt]\addplot[mark=*,mark size=0.7pt,draw=black,thin] coordinates {(0,9637591) (1,55433519) (2,125293476) (3,406022662) (4,865744492) (5,582927859) (6,182647411) (7,15733612) (8,380777) (9,8509) (10,38412) (11,12) (12,6423) (13,2) (14,771) (15,1) (16,39)};\end{axis}\end{tikzpicture}\end{minipage}
\vspace{3mm}

\noindent\begin{minipage}[t]{60mm}\vspace{0pt}\sbox\hmboard{\chessboard[setfen=8/7R/4N1pk/8/1r6/8/8/K7 b - - 0 1]}\usebox\hmboard\\[1mm]
\makebox[\wd\hmboard]{\sffamily\hspace{3mm}\textbf{h\#8}\hfill(3+3)}\end{minipage}\hfill
\begin{minipage}[t]{52mm}\vspace{0pt}
\textcolor{ink}{\sffamily\textsc{Lösung}}\par\smallskip\small
\begin{tabular}{@{}r@{\enspace}l@{\enspace}l@{}}
\textbf{1.} & \symking{}xh7 & \symknight{}g7\\
\textbf{2.} & \symrook{}h4 & \symking{}b2\\
\textbf{3.} & \symking{}h6 & \symking{}c3\\
\textbf{4.} & \symrook{}h3+ & \symking{}d4\\
\textbf{5.} & \symking{}g5 & \symking{}e5\\
\textbf{6.} & \symking{}h4 & \symking{}f6\\
\textbf{7.} & g5 & \symking{}g6\\
\textbf{8.} & g4 & \symknight{}f5\#\\
\end{tabular}\par\vspace{3mm}
\qrcode[height=12mm]{https://helpman.komtera.lt/?fen=8/7R/4N1pk/8/1r6/8/8/K7&moves=8}\\{\sffamily\tiny\color{ink}scannen zum Lösen}
\end{minipage}
\vspace{2mm}

{\raggedright\noindent\textcolor{ink}{\sffamily\textsc{Themen}}\enspace \themetag{reines Matt} \themetag{Mustermatt} \themetag{Idealmatt} \themetag{Switchback} \themetag{Selbstblock} \themetag{Umnov}\par\smallskip{\footnotesize\textit{Hinweis:} die Partei am Zug steht im Diagramm im Schach; die Lösung beginnt mit einem Schlag.}\par}

\clearpage
\phantomsection\label{hm:KRPvkbb}
\begin{tcolorbox}[enhanced,colback=ink,colframe=ink,arc=0pt,boxrule=0pt,left=4mm,right=4mm,top=2.5mm,bottom=2.5mm]
\color{white}\sffamily{\LARGE\bfseries KRPvkbb}\hfill{\Large h\#8}\\[0.5mm]
{\small 6 Steine\quad Schwarz am Zug}\hfill{\small Nr.~420\qquad\href{https://helpman.komtera.lt/?fen=6bk/7R/8/8/8/bP6/8/K7&moves=8}{\color{white}\bfseries lösen\,\ensuremath{\nearrow}}}
\end{tcolorbox}
\vspace{2mm}
\noindent\begin{minipage}[t]{36mm}\vspace{0pt}\sffamily\small\raggedright
\hmfigure{längste eindeutige}{{\Large h\#8}}\hmfigure{eindeutig bei dieser Länge}{{\Large 4} Stellungen}\hmfigure{längstes Matt}{{\Large h\#9}}\hmfigure{in der Klasse}{45,4 Milliarden Stellungen,\\ 14 Milliarden mit Hilfsmatt}
\end{minipage}\hfill
\begin{minipage}[t]{80mm}\vspace{0pt}\begin{tikzpicture}\begin{axis}[width=72mm,height=42mm,ybar,bar width=1.6pt,xmin=-1,xmax=19,ymin=0,ymax=3.007,axis y line*=left,axis x line*=bottom,xlabel={\scriptsize Halbzüge bis Matt},ylabel={\scriptsize Stellungen (Mrd.)},tick label style={font=\tiny},label style={font=\scriptsize},enlarge x limits=false,ylabel shift=-3pt,xlabel shift=-3pt]\addplot[fill=ink!55,draw=none] coordinates {(0,0.002677) (1,0.02092) (2,0.07485) (3,0.4636) (4,1.049) (5,2.33) (6,2.734) (7,2.702) (8,2.254) (9,1.398) (10,0.7086) (11,0.1702) (12,0.04673) (13,0.0007095) (14,0.000233) (16,1.857e-05) (18,7.7e-08)};\draw[densely dotted] (axis cs:16,0) -- (axis cs:16,3.007);\end{axis}\begin{axis}[width=72mm,height=42mm,axis y line*=right,axis x line=none,ymode=log,xmin=-1,xmax=19,ylabel={\scriptsize eindeutig (log)},tick label style={font=\tiny},label style={font=\scriptsize},ylabel shift=-3pt]\addplot[mark=*,mark size=0.7pt,draw=black,thin] coordinates {(0,2677384) (1,16580358) (2,26706054) (3,103823112) (4,152659144) (5,133884916) (6,72789959) (7,31036559) (8,13926468) (9,3316030) (10,658139) (11,8058) (12,1163) (14,248) (16,4)};\end{axis}\end{tikzpicture}\end{minipage}
\vspace{3mm}

\noindent\begin{minipage}[t]{60mm}\vspace{0pt}\sbox\hmboard{\chessboard[setfen=6bk/7R/8/8/8/bP6/8/K7 b - - 0 1]}\usebox\hmboard\\[1mm]
\makebox[\wd\hmboard]{\sffamily\hspace{3mm}\textbf{h\#8}\hfill(3+3)}\end{minipage}\hfill
\begin{minipage}[t]{52mm}\vspace{0pt}
\textcolor{ink}{\sffamily\textsc{Lösung}}\par\smallskip\small
\begin{tabular}{@{}r@{\enspace}l@{\enspace}l@{}}
\textbf{1.} & \symking{}xh7 & \symking{}a2\\
\textbf{2.} & \symking{}g6 & \symking{}xa3\\
\textbf{3.} & \symking{}f5 & b4\\
\textbf{4.} & \symking{}e4 & b5\\
\textbf{5.} & \symking{}d3 & b6\\
\textbf{6.} & \symking{}c2 & b7\\
\textbf{7.} & \symking{}b1 & b8=\symqueen{}+\\
\textbf{8.} & \symking{}a1 & \symqueen{}b2\#\\
\end{tabular}\par\vspace{3mm}
\qrcode[height=12mm]{https://helpman.komtera.lt/?fen=6bk/7R/8/8/8/bP6/8/K7&moves=8}\\{\sffamily\tiny\color{ink}scannen zum Lösen}
\end{minipage}
\vspace{2mm}

{\raggedright\noindent\textcolor{ink}{\sffamily\textsc{Themen}}\enspace \themetag{Umwandlung} \themetag{ein Stein zieht}\par\smallskip{\footnotesize\textit{Hinweis:} die Partei am Zug steht im Diagramm im Schach; die Lösung beginnt mit einem Schlag.}\par}

\clearpage
\phantomsection\label{hm:KRPvkbn}
\begin{tcolorbox}[enhanced,colback=ink,colframe=ink,arc=0pt,boxrule=0pt,left=4mm,right=4mm,top=2.5mm,bottom=2.5mm]
\color{white}\sffamily{\LARGE\bfseries KRPvkbn}\hfill{\Large h\#8}\\[0.5mm]
{\small 6 Steine\quad Schwarz am Zug}\hfill{\small Nr.~421\qquad\href{https://helpman.komtera.lt/?fen=6Rk/7b/8/8/8/3n4/3P4/K7&moves=8}{\color{white}\bfseries lösen\,\ensuremath{\nearrow}}}
\end{tcolorbox}
\vspace{2mm}
\noindent\begin{minipage}[t]{36mm}\vspace{0pt}\sffamily\small\raggedright
\hmfigure{längste eindeutige}{{\Large h\#8}}\hmfigure{eindeutig bei dieser Länge}{{\Large 17} Stellungen}\hmfigure{längstes Matt}{{\Large h\#9}}\hmfigure{in der Klasse}{45,4 Milliarden Stellungen,\\ 28,7 Milliarden mit Hilfsmatt}
\end{minipage}\hfill
\begin{minipage}[t]{80mm}\vspace{0pt}\begin{tikzpicture}\begin{axis}[width=72mm,height=42mm,ybar,bar width=1.6pt,xmin=-1,xmax=19,ymin=0,ymax=6.403,axis y line*=left,axis x line*=bottom,xlabel={\scriptsize Halbzüge bis Matt},ylabel={\scriptsize Stellungen (Mrd.)},tick label style={font=\tiny},label style={font=\scriptsize},enlarge x limits=false,ylabel shift=-3pt,xlabel shift=-3pt]\addplot[fill=ink!55,draw=none] coordinates {(0,0.007032) (1,0.05503) (2,0.1798) (3,1.019) (4,2.262) (5,4.717) (6,5.821) (7,5.353) (8,4.783) (9,2.726) (10,1.423) (11,0.3009) (12,0.08903) (13,0.0009065) (14,0.0006753) (15,8.9e-08) (16,7.997e-05) (18,4.2e-07)};\draw[densely dotted] (axis cs:16,0) -- (axis cs:16,6.403);\end{axis}\begin{axis}[width=72mm,height=42mm,axis y line*=right,axis x line=none,ymode=log,xmin=-1,xmax=19,ylabel={\scriptsize eindeutig (log)},tick label style={font=\tiny},label style={font=\scriptsize},ylabel shift=-3pt]\addplot[mark=*,mark size=0.7pt,draw=black,thin] coordinates {(0,7031930) (1,43466633) (2,65588401) (3,227423752) (4,348321396) (5,277420555) (6,161958210) (7,58219678) (8,27084313) (9,5388986) (10,1265761) (11,8126) (12,2358) (14,946) (16,17)};\end{axis}\end{tikzpicture}\end{minipage}
\vspace{3mm}

\noindent\begin{minipage}[t]{60mm}\vspace{0pt}\sbox\hmboard{\chessboard[setfen=6Rk/7b/8/8/8/3n4/3P4/K7 b - - 0 1]}\usebox\hmboard\\[1mm]
\makebox[\wd\hmboard]{\sffamily\hspace{3mm}\textbf{h\#8}\hfill(3+3)}\end{minipage}\hfill
\begin{minipage}[t]{52mm}\vspace{0pt}
\textcolor{ink}{\sffamily\textsc{Lösung}}\par\smallskip\small
\begin{tabular}{@{}r@{\enspace}l@{\enspace}l@{}}
\textbf{1.} & \symking{}xg8 & \symking{}b1\\
\textbf{2.} & \symknight{}e5+ & \symking{}c1\\
\textbf{3.} & \symking{}f7 & d4\\
\textbf{4.} & \symking{}e6 & dxe5\\
\textbf{5.} & \symking{}d5 & e6\\
\textbf{6.} & \symking{}c4 & e7\\
\textbf{7.} & \symking{}b3 & e8=\symqueen{}\\
\textbf{8.} & \symking{}a2 & \symqueen{}a4\#\\
\end{tabular}\par\vspace{3mm}
\qrcode[height=12mm]{https://helpman.komtera.lt/?fen=6Rk/7b/8/8/8/3n4/3P4/K7&moves=8}\\{\sffamily\tiny\color{ink}scannen zum Lösen}
\end{minipage}
\vspace{2mm}

{\raggedright\noindent\textcolor{ink}{\sffamily\textsc{Themen}}\enspace \themetag{reines Matt} \themetag{Mustermatt} \themetag{Spiegelmatt} \themetag{Umwandlung} \themetag{Excelsior} \themetag{Umnov}\par\smallskip{\footnotesize\textit{Hinweis:} die Partei am Zug steht im Diagramm im Schach; die Lösung beginnt mit einem Schlag.}\par}

\clearpage
\phantomsection\label{hm:KRPvknn}
\begin{tcolorbox}[enhanced,colback=ink,colframe=ink,arc=0pt,boxrule=0pt,left=4mm,right=4mm,top=2.5mm,bottom=2.5mm]
\color{white}\sffamily{\LARGE\bfseries KRPvknn}\hfill{\Large h\#8}\\[0.5mm]
{\small 6 Steine\quad Schwarz am Zug}\hfill{\small Nr.~422\qquad\href{https://helpman.komtera.lt/?fen=K7/8/2n5/8/8/8/2P4R/6nk&moves=8}{\color{white}\bfseries lösen\,\ensuremath{\nearrow}}}
\end{tcolorbox}
\vspace{2mm}
\noindent\begin{minipage}[t]{36mm}\vspace{0pt}\sffamily\small\raggedright
\hmfigure{längste eindeutige}{{\Large h\#8}}\hmfigure{eindeutig bei dieser Länge}{{\Large 7} Stellungen}\hmfigure{längstes Matt}{{\Large h\#9}}\hmfigure{in der Klasse}{45,4 Milliarden Stellungen,\\ 14,7 Milliarden mit Hilfsmatt}
\end{minipage}\hfill
\begin{minipage}[t]{80mm}\vspace{0pt}\begin{tikzpicture}\begin{axis}[width=72mm,height=42mm,ybar,bar width=1.6pt,xmin=-1,xmax=19,ymin=0,ymax=3.279,axis y line*=left,axis x line*=bottom,xlabel={\scriptsize Halbzüge bis Matt},ylabel={\scriptsize Stellungen (Mrd.)},tick label style={font=\tiny},label style={font=\scriptsize},enlarge x limits=false,ylabel shift=-3pt,xlabel shift=-3pt]\addplot[fill=ink!55,draw=none] coordinates {(0,0.004605) (1,0.03586) (2,0.1038) (3,0.5436) (4,1.177) (5,2.336) (6,2.98) (7,2.612) (8,2.503) (9,1.388) (10,0.813) (11,0.1698) (12,0.05955) (13,0.0006653) (14,0.0004278) (15,8e-08) (16,6.097e-05) (17,2.3e-08) (18,2.61e-07)};\draw[densely dotted] (axis cs:16,0) -- (axis cs:16,3.279);\end{axis}\begin{axis}[width=72mm,height=42mm,axis y line*=right,axis x line=none,ymode=log,xmin=-1,xmax=19,ylabel={\scriptsize eindeutig (log)},tick label style={font=\tiny},label style={font=\scriptsize},ylabel shift=-3pt]\addplot[mark=*,mark size=0.7pt,draw=black,thin] coordinates {(0,4604812) (1,28198621) (2,39428863) (3,124204709) (4,189383839) (5,137664417) (6,83300601) (7,26503526) (8,13772494) (9,2553872) (10,716304) (11,3542) (12,1275) (14,634) (16,7)};\end{axis}\end{tikzpicture}\end{minipage}
\vspace{3mm}

\noindent\begin{minipage}[t]{60mm}\vspace{0pt}\sbox\hmboard{\chessboard[setfen=K7/8/2n5/8/8/8/2P4R/6nk b - - 0 1]}\usebox\hmboard\\[1mm]
\makebox[\wd\hmboard]{\sffamily\hspace{3mm}\textbf{h\#8}\hfill(3+3)}\end{minipage}\hfill
\begin{minipage}[t]{52mm}\vspace{0pt}
\textcolor{ink}{\sffamily\textsc{Lösung}}\par\smallskip\small
\begin{tabular}{@{}r@{\enspace}l@{\enspace}l@{}}
\textbf{1.} & \symking{}xh2 & c3\\
\textbf{2.} & \symknight{}d4 & cxd4\\
\textbf{3.} & \symking{}g3 & d5\\
\textbf{4.} & \symking{}f4 & d6\\
\textbf{5.} & \symking{}e5 & d7\\
\textbf{6.} & \symking{}d6 & \symking{}a7\\
\textbf{7.} & \symking{}c7 & \symking{}a6\\
\textbf{8.} & \symking{}b8 & d8=\symqueen{}\#\\
\end{tabular}\par\vspace{3mm}
\qrcode[height=12mm]{https://helpman.komtera.lt/?fen=K7/8/2n5/8/8/8/2P4R/6nk&moves=8}\\{\sffamily\tiny\color{ink}scannen zum Lösen}
\end{minipage}
\vspace{2mm}

{\raggedright\noindent\textcolor{ink}{\sffamily\textsc{Themen}}\enspace \themetag{reines Matt} \themetag{Mustermatt} \themetag{Spiegelmatt} \themetag{Umwandlung} \themetag{Excelsior} \themetag{ohne Schach} \themetag{Umnov}\par\smallskip{\footnotesize\textit{Hinweis:} die Partei am Zug steht im Diagramm im Schach; die Lösung beginnt mit einem Schlag.}\par}

\clearpage
\phantomsection\label{hm:KRPvknp}
\begin{tcolorbox}[enhanced,colback=ink,colframe=ink,arc=0pt,boxrule=0pt,left=4mm,right=4mm,top=2.5mm,bottom=2.5mm]
\color{white}\sffamily{\LARGE\bfseries KRPvknp}\hfill{\Large h\#8}\\[0.5mm]
{\small 6 Steine\quad Schwarz am Zug}\hfill{\small Nr.~423\qquad\href{https://helpman.komtera.lt/?fen=8/8/8/2R5/2kp4/K2n4/3P4/8&moves=8}{\color{white}\bfseries lösen\,\ensuremath{\nearrow}}}
\end{tcolorbox}
\vspace{2mm}
\noindent\begin{minipage}[t]{36mm}\vspace{0pt}\sffamily\small\raggedright
\hmfigure{längste eindeutige}{{\Large h\#8}}\hmfigure{eindeutig bei dieser Länge}{{\Large 4} Stellungen}\hmfigure{längstes Matt}{{\Large h\#12.5}}\hmfigure{in der Klasse}{34,1 Milliarden Stellungen,\\ 22,7 Milliarden mit Hilfsmatt}
\end{minipage}\hfill
\begin{minipage}[t]{80mm}\vspace{0pt}\begin{tikzpicture}\begin{axis}[width=72mm,height=42mm,ybar,bar width=1.6pt,xmin=-1,xmax=26,ymin=0,ymax=5.26,axis y line*=left,axis x line*=bottom,xlabel={\scriptsize Halbzüge bis Matt},ylabel={\scriptsize Stellungen (Mrd.)},tick label style={font=\tiny},label style={font=\scriptsize},enlarge x limits=false,ylabel shift=-3pt,xlabel shift=-3pt]\addplot[fill=ink!55,draw=none] coordinates {(0,0.00921) (1,0.06847) (2,0.1822) (3,0.9091) (4,1.972) (5,3.601) (6,4.782) (7,3.839) (8,3.88) (9,1.948) (10,1.176) (11,0.2229) (12,0.08861) (13,0.001085) (14,0.0007176) (15,1.61e-07) (16,9.433e-05) (17,8e-09) (18,2.443e-05) (19,1e-09) (20,5.567e-06) (21,2.3e-08) (22,6.91e-07) (23,7e-09) (24,2.7e-08) (25,1e-09)};\draw[densely dotted] (axis cs:16,0) -- (axis cs:16,5.26);\end{axis}\begin{axis}[width=72mm,height=42mm,axis y line*=right,axis x line=none,ymode=log,xmin=-1,xmax=26,ylabel={\scriptsize eindeutig (log)},tick label style={font=\tiny},label style={font=\scriptsize},ylabel shift=-3pt]\addplot[mark=*,mark size=0.7pt,draw=black,thin] coordinates {(0,9210260) (1,53632816) (2,72806523) (3,216156160) (4,349952821) (5,213238577) (6,132442457) (7,36424239) (8,18972896) (9,3076401) (10,994833) (11,10256) (12,6003) (14,591) (16,4)};\end{axis}\end{tikzpicture}\end{minipage}
\vspace{3mm}

\noindent\begin{minipage}[t]{60mm}\vspace{0pt}\sbox\hmboard{\chessboard[setfen=8/8/8/2R5/2kp4/K2n4/3P4/8 b - - 0 1]}\usebox\hmboard\\[1mm]
\makebox[\wd\hmboard]{\sffamily\hspace{3mm}\textbf{h\#8}\hfill(3+3)}\end{minipage}\hfill
\begin{minipage}[t]{52mm}\vspace{0pt}
\textcolor{ink}{\sffamily\textsc{Lösung}}\par\smallskip\small
\begin{tabular}{@{}r@{\enspace}l@{\enspace}l@{}}
\textbf{1.} & \symking{}xc5 & \symking{}b3\\
\textbf{2.} & \symknight{}f2 & d3\\
\textbf{3.} & \symknight{}e4 & dxe4\\
\textbf{4.} & d3 & e5\\
\textbf{5.} & \symking{}d4 & e6\\
\textbf{6.} & \symking{}e3 & e7\\
\textbf{7.} & \symking{}d2 & e8=\symqueen{}\\
\textbf{8.} & \symking{}c1 & \symqueen{}e1\#\\
\end{tabular}\par\vspace{3mm}
\qrcode[height=12mm]{https://helpman.komtera.lt/?fen=8/8/8/2R5/2kp4/K2n4/3P4/8&moves=8}\\{\sffamily\tiny\color{ink}scannen zum Lösen}
\end{minipage}
\vspace{2mm}

{\raggedright\noindent\textcolor{ink}{\sffamily\textsc{Themen}}\enspace \themetag{reines Matt} \themetag{Mustermatt} \themetag{Spiegelmatt} \themetag{Umwandlung} \themetag{Excelsior} \themetag{ohne Schach} \themetag{Umnov}\par\smallskip{\footnotesize\textit{Hinweis:} die Partei am Zug steht im Diagramm im Schach; die Lösung beginnt mit einem Schlag.}\par}

\clearpage
\phantomsection\label{hm:KRPvkqb}
\begin{tcolorbox}[enhanced,colback=ink,colframe=ink,arc=0pt,boxrule=0pt,left=4mm,right=4mm,top=2.5mm,bottom=2.5mm]
\color{white}\sffamily{\LARGE\bfseries KRPvkqb}\hfill{\Large h\#8}\\[0.5mm]
{\small 6 Steine\quad Schwarz am Zug}\hfill{\small Nr.~424\qquad\href{https://helpman.komtera.lt/?fen=6qk/7R/8/8/8/bP6/8/K7&moves=8}{\color{white}\bfseries lösen\,\ensuremath{\nearrow}}}
\end{tcolorbox}
\vspace{2mm}
\noindent\begin{minipage}[t]{36mm}\vspace{0pt}\sffamily\small\raggedright
\hmfigure{längste eindeutige}{{\Large h\#8}}\hmfigure{eindeutig bei dieser Länge}{{\Large 1} Stellung}\hmfigure{längstes Matt}{{\Large h\#12}}\hmfigure{in der Klasse}{45,4 Milliarden Stellungen,\\ 24,8 Milliarden mit Hilfsmatt}
\end{minipage}\hfill
\begin{minipage}[t]{80mm}\vspace{0pt}\begin{tikzpicture}\begin{axis}[width=72mm,height=42mm,ybar,bar width=1.6pt,xmin=-1,xmax=25,ymin=0,ymax=6.108,axis y line*=left,axis x line*=bottom,xlabel={\scriptsize Halbzüge bis Matt},ylabel={\scriptsize Stellungen (Mrd.)},tick label style={font=\tiny},label style={font=\scriptsize},enlarge x limits=false,ylabel shift=-3pt,xlabel shift=-3pt]\addplot[fill=ink!55,draw=none] coordinates {(0,0.002681) (1,0.02192) (2,0.1008) (3,0.7089) (4,1.559) (5,4.198) (6,4.181) (7,5.552) (8,3.557) (9,3.148) (10,1.182) (11,0.4887) (12,0.07258) (13,0.005432) (14,0.0004744) (15,1.782e-06) (16,3.147e-05) (17,1.2e-07) (18,2.5e-08) (20,7e-09) (22,1.9e-08) (24,9e-09)};\draw[densely dotted] (axis cs:16,0) -- (axis cs:16,6.108);\end{axis}\begin{axis}[width=72mm,height=42mm,axis y line*=right,axis x line=none,ymode=log,xmin=-1,xmax=25,ylabel={\scriptsize eindeutig (log)},tick label style={font=\tiny},label style={font=\scriptsize},ylabel shift=-3pt]\addplot[mark=*,mark size=0.7pt,draw=black,thin] coordinates {(0,2681168) (1,17713104) (2,36248847) (3,173313205) (4,201642163) (5,247982162) (6,109328784) (7,76875568) (8,23944875) (9,9713161) (10,1172773) (11,80403) (12,1701) (13,6) (14,404) (16,1)};\end{axis}\end{tikzpicture}\end{minipage}
\vspace{3mm}

\noindent\begin{minipage}[t]{60mm}\vspace{0pt}\sbox\hmboard{\chessboard[setfen=6qk/7R/8/8/8/bP6/8/K7 b - - 0 1]}\usebox\hmboard\\[1mm]
\makebox[\wd\hmboard]{\sffamily\hspace{3mm}\textbf{h\#8}\hfill(3+3)}\end{minipage}\hfill
\begin{minipage}[t]{52mm}\vspace{0pt}
\textcolor{ink}{\sffamily\textsc{Lösung}}\par\smallskip\small
\begin{tabular}{@{}r@{\enspace}l@{\enspace}l@{}}
\textbf{1.} & \symking{}xh7 & \symking{}a2\\
\textbf{2.} & \symking{}g6 & \symking{}xa3\\
\textbf{3.} & \symking{}f5 & b4\\
\textbf{4.} & \symking{}e4 & b5\\
\textbf{5.} & \symking{}d3 & b6\\
\textbf{6.} & \symking{}c2 & b7\\
\textbf{7.} & \symking{}b1 & b8=\symqueen{}+\\
\textbf{8.} & \symking{}a1 & \symqueen{}b2\#\\
\end{tabular}\par\vspace{3mm}
\qrcode[height=12mm]{https://helpman.komtera.lt/?fen=6qk/7R/8/8/8/bP6/8/K7&moves=8}\\{\sffamily\tiny\color{ink}scannen zum Lösen}
\end{minipage}
\vspace{2mm}

{\raggedright\noindent\textcolor{ink}{\sffamily\textsc{Themen}}\enspace \themetag{Umwandlung} \themetag{ein Stein zieht}\par\smallskip{\footnotesize\textit{Hinweis:} die Partei am Zug steht im Diagramm im Schach; die Lösung beginnt mit einem Schlag.}\par}

\clearpage
\phantomsection\label{hm:KRPvkqn}
\begin{tcolorbox}[enhanced,colback=ink,colframe=ink,arc=0pt,boxrule=0pt,left=4mm,right=4mm,top=2.5mm,bottom=2.5mm]
\color{white}\sffamily{\LARGE\bfseries KRPvkqn}\hfill{\Large h\#8}\\[0.5mm]
{\small 6 Steine\quad Schwarz am Zug}\hfill{\small Nr.~425\qquad\href{https://helpman.komtera.lt/?fen=6Rk/7q/8/8/8/3n4/3P4/K7&moves=8}{\color{white}\bfseries lösen\,\ensuremath{\nearrow}}}
\end{tcolorbox}
\vspace{2mm}
\noindent\begin{minipage}[t]{36mm}\vspace{0pt}\sffamily\small\raggedright
\hmfigure{längste eindeutige}{{\Large h\#8}}\hmfigure{eindeutig bei dieser Länge}{{\Large 11} Stellungen}\hmfigure{längstes Matt}{{\Large h\#9}}\hmfigure{in der Klasse}{45,4 Milliarden Stellungen,\\ 25,4 Milliarden mit Hilfsmatt}
\end{minipage}\hfill
\begin{minipage}[t]{80mm}\vspace{0pt}\begin{tikzpicture}\begin{axis}[width=72mm,height=42mm,ybar,bar width=1.6pt,xmin=-1,xmax=19,ymin=0,ymax=6.121,axis y line*=left,axis x line*=bottom,xlabel={\scriptsize Halbzüge bis Matt},ylabel={\scriptsize Stellungen (Mrd.)},tick label style={font=\tiny},label style={font=\scriptsize},enlarge x limits=false,ylabel shift=-3pt,xlabel shift=-3pt]\addplot[fill=ink!55,draw=none] coordinates {(0,0.003383) (1,0.02806) (2,0.1243) (3,0.7977) (4,1.703) (5,4.314) (6,4.524) (7,5.564) (8,3.761) (9,3.02) (10,1.125) (11,0.4126) (12,0.06504) (13,0.004174) (14,0.0005575) (15,1.975e-06) (16,7.232e-05) (17,3.5e-08) (18,5.69e-07)};\draw[densely dotted] (axis cs:16,0) -- (axis cs:16,6.121);\end{axis}\begin{axis}[width=72mm,height=42mm,axis y line*=right,axis x line=none,ymode=log,xmin=-1,xmax=19,ylabel={\scriptsize eindeutig (log)},tick label style={font=\tiny},label style={font=\scriptsize},ylabel shift=-3pt]\addplot[mark=*,mark size=0.7pt,draw=black,thin] coordinates {(0,3382622) (1,22572642) (2,45295234) (3,194002287) (4,223542117) (5,259837344) (6,123466614) (7,75898729) (8,23283560) (9,7867546) (10,1021223) (11,40154) (12,1516) (14,756) (16,11)};\end{axis}\end{tikzpicture}\end{minipage}
\vspace{3mm}

\noindent\begin{minipage}[t]{60mm}\vspace{0pt}\sbox\hmboard{\chessboard[setfen=6Rk/7q/8/8/8/3n4/3P4/K7 b - - 0 1]}\usebox\hmboard\\[1mm]
\makebox[\wd\hmboard]{\sffamily\hspace{3mm}\textbf{h\#8}\hfill(3+3)}\end{minipage}\hfill
\begin{minipage}[t]{52mm}\vspace{0pt}
\textcolor{ink}{\sffamily\textsc{Lösung}}\par\smallskip\small
\begin{tabular}{@{}r@{\enspace}l@{\enspace}l@{}}
\textbf{1.} & \symking{}xg8 & \symking{}b1\\
\textbf{2.} & \symknight{}e5+ & \symking{}c1\\
\textbf{3.} & \symking{}f7 & d4\\
\textbf{4.} & \symking{}e6 & dxe5\\
\textbf{5.} & \symking{}d5 & e6\\
\textbf{6.} & \symking{}c4 & e7\\
\textbf{7.} & \symking{}b3 & e8=\symqueen{}\\
\textbf{8.} & \symking{}a2 & \symqueen{}a4\#\\
\end{tabular}\par\vspace{3mm}
\qrcode[height=12mm]{https://helpman.komtera.lt/?fen=6Rk/7q/8/8/8/3n4/3P4/K7&moves=8}\\{\sffamily\tiny\color{ink}scannen zum Lösen}
\end{minipage}
\vspace{2mm}

{\raggedright\noindent\textcolor{ink}{\sffamily\textsc{Themen}}\enspace \themetag{reines Matt} \themetag{Mustermatt} \themetag{Spiegelmatt} \themetag{Umwandlung} \themetag{Excelsior} \themetag{Umnov}\par\smallskip{\footnotesize\textit{Hinweis:} die Partei am Zug steht im Diagramm im Schach; die Lösung beginnt mit einem Schlag.}\par}

\clearpage
\phantomsection\label{hm:KRPvkqp}
\begin{tcolorbox}[enhanced,colback=ink,colframe=ink,arc=0pt,boxrule=0pt,left=4mm,right=4mm,top=2.5mm,bottom=2.5mm]
\color{white}\sffamily{\LARGE\bfseries KRPvkqp}\hfill{\Large h\#8}\\[0.5mm]
{\small 6 Steine\quad Schwarz am Zug}\hfill{\small Nr.~426\qquad\href{https://helpman.komtera.lt/?fen=8/8/8/8/5p2/8/1R3P2/qk1K4&moves=8}{\color{white}\bfseries lösen\,\ensuremath{\nearrow}}}
\end{tcolorbox}
\vspace{2mm}
\noindent\begin{minipage}[t]{36mm}\vspace{0pt}\sffamily\small\raggedright
\hmfigure{längste eindeutige}{{\Large h\#8}}\hmfigure{eindeutig bei dieser Länge}{{\Large 13} Stellungen}\hmfigure{längstes Matt}{{\Large h\#14.5}}\hmfigure{in der Klasse}{34,1 Milliarden Stellungen,\\ 19,5 Milliarden mit Hilfsmatt}
\end{minipage}\hfill
\begin{minipage}[t]{80mm}\vspace{0pt}\begin{tikzpicture}\begin{axis}[width=72mm,height=42mm,ybar,bar width=1.6pt,xmin=-1,xmax=30,ymin=0,ymax=4.424,axis y line*=left,axis x line*=bottom,xlabel={\scriptsize Halbzüge bis Matt},ylabel={\scriptsize Stellungen (Mrd.)},tick label style={font=\tiny},label style={font=\scriptsize},enlarge x limits=false,ylabel shift=-3pt,xlabel shift=-3pt]\addplot[fill=ink!55,draw=none] coordinates {(0,0.003401) (1,0.02687) (2,0.1124) (3,0.6827) (4,1.427) (5,3.343) (6,3.534) (7,4.022) (8,2.878) (9,2.181) (10,0.9093) (11,0.3149) (12,0.06306) (13,0.003379) (14,0.0004061) (15,7.246e-06) (16,4.1e-05) (17,4.71e-07) (18,6.202e-06) (19,3.98e-07) (20,4.4e-08) (21,7.09e-07) (22,3e-09) (23,4.79e-07) (24,3e-09) (25,2.66e-07) (27,1.47e-07) (29,2.8e-08)};\draw[densely dotted] (axis cs:16,0) -- (axis cs:16,4.424);\end{axis}\begin{axis}[width=72mm,height=42mm,axis y line*=right,axis x line=none,ymode=log,xmin=-1,xmax=30,ylabel={\scriptsize eindeutig (log)},tick label style={font=\tiny},label style={font=\scriptsize},ylabel shift=-3pt]\addplot[mark=*,mark size=0.7pt,draw=black,thin] coordinates {(0,3400830) (1,21552176) (2,41782307) (3,168680571) (4,190651729) (5,199470019) (6,90870704) (7,51236346) (8,18360158) (9,5727482) (10,844175) (11,33189) (12,2157) (13,94) (14,193) (15,1) (16,13)};\end{axis}\end{tikzpicture}\end{minipage}
\vspace{3mm}

\noindent\begin{minipage}[t]{60mm}\vspace{0pt}\sbox\hmboard{\chessboard[setfen=8/8/8/8/5p2/8/1R3P2/qk1K4 b - - 0 1]}\usebox\hmboard\\[1mm]
\makebox[\wd\hmboard]{\sffamily\hspace{3mm}\textbf{h\#8}\hfill(3+3)}\end{minipage}\hfill
\begin{minipage}[t]{52mm}\vspace{0pt}
\textcolor{ink}{\sffamily\textsc{Lösung}}\par\smallskip\small
\begin{tabular}{@{}r@{\enspace}l@{\enspace}l@{}}
\textbf{1.} & \symking{}xb2+ & \symking{}d2\\
\textbf{2.} & \symqueen{}g1 & f3\\
\textbf{3.} & \symqueen{}g4 & fxg4\\
\textbf{4.} & f3 & g5\\
\textbf{5.} & f2 & g6\\
\textbf{6.} & f1=\symrook{} & g7\\
\textbf{7.} & \symrook{}a1 & g8=\symqueen{}\\
\textbf{8.} & \symking{}b1 & \symqueen{}b3\#\\
\end{tabular}\par\vspace{3mm}
\qrcode[height=12mm]{https://helpman.komtera.lt/?fen=8/8/8/8/5p2/8/1R3P2/qk1K4&moves=8}\\{\sffamily\tiny\color{ink}scannen zum Lösen}
\end{minipage}
\vspace{2mm}

{\raggedright\noindent\textcolor{ink}{\sffamily\textsc{Themen}}\enspace \themetag{Umwandlung} \themetag{Unterverwandlung} \themetag{Excelsior} \themetag{Switchback} \themetag{Selbstblock} \themetag{Umnov} \themetag{Klasinc}\par\smallskip{\footnotesize\textit{Hinweis:} die Partei am Zug steht im Diagramm im Schach; die Lösung beginnt mit einem Schlag.}\par}

\clearpage
\phantomsection\label{hm:KRPvkrp}
\begin{tcolorbox}[enhanced,colback=ink,colframe=ink,arc=0pt,boxrule=0pt,left=4mm,right=4mm,top=2.5mm,bottom=2.5mm]
\color{white}\sffamily{\LARGE\bfseries KRPvkrp}\hfill{\Large h\#8}\\[0.5mm]
{\small 6 Steine\quad Schwarz am Zug}\hfill{\small Nr.~427\qquad\href{https://helpman.komtera.lt/?fen=8/8/8/8/4p3/8/2r1P3/K1kR4&moves=8}{\color{white}\bfseries lösen\,\ensuremath{\nearrow}}}
\end{tcolorbox}
\vspace{2mm}
\noindent\begin{minipage}[t]{36mm}\vspace{0pt}\sffamily\small\raggedright
\hmfigure{längste eindeutige}{{\Large h\#8}}\hmfigure{eindeutig bei dieser Länge}{{\Large 12} Stellungen}\hmfigure{längstes Matt}{{\Large h\#12.5}}\hmfigure{in der Klasse}{34,1 Milliarden Stellungen,\\ 21,2 Milliarden mit Hilfsmatt}
\end{minipage}\hfill
\begin{minipage}[t]{80mm}\vspace{0pt}\begin{tikzpicture}\begin{axis}[width=72mm,height=42mm,ybar,bar width=1.6pt,xmin=-1,xmax=26,ymin=0,ymax=4.766,axis y line*=left,axis x line*=bottom,xlabel={\scriptsize Halbzüge bis Matt},ylabel={\scriptsize Stellungen (Mrd.)},tick label style={font=\tiny},label style={font=\scriptsize},enlarge x limits=false,ylabel shift=-3pt,xlabel shift=-3pt]\addplot[fill=ink!55,draw=none] coordinates {(0,0.005838) (1,0.04539) (2,0.1605) (3,0.8496) (4,1.915) (5,3.621) (6,4.333) (7,3.978) (8,3.268) (9,1.885) (10,0.8605) (11,0.2045) (12,0.04641) (13,0.0008854) (14,0.0003187) (15,1.77e-06) (16,4.775e-05) (17,1.8e-08) (18,1.035e-05) (19,1.7e-08) (20,1.13e-07) (21,7.6e-08) (22,3e-08) (23,1.3e-08) (24,9e-09) (25,2e-08)};\draw[densely dotted] (axis cs:16,0) -- (axis cs:16,4.766);\end{axis}\begin{axis}[width=72mm,height=42mm,axis y line*=right,axis x line=none,ymode=log,xmin=-1,xmax=26,ylabel={\scriptsize eindeutig (log)},tick label style={font=\tiny},label style={font=\scriptsize},ylabel shift=-3pt]\addplot[mark=*,mark size=0.7pt,draw=black,thin] coordinates {(0,5838029) (1,36037226) (2,66590402) (3,203957768) (4,302162794) (5,215837178) (6,115093974) (7,43368004) (8,17494492) (9,3455412) (10,779293) (11,13853) (12,1540) (13,1) (14,195) (16,12)};\end{axis}\end{tikzpicture}\end{minipage}
\vspace{3mm}

\noindent\begin{minipage}[t]{60mm}\vspace{0pt}\sbox\hmboard{\chessboard[setfen=8/8/8/8/4p3/8/2r1P3/K1kR4 b - - 0 1]}\usebox\hmboard\\[1mm]
\makebox[\wd\hmboard]{\sffamily\hspace{3mm}\textbf{h\#8}\hfill(3+3)}\end{minipage}\hfill
\begin{minipage}[t]{52mm}\vspace{0pt}
\textcolor{ink}{\sffamily\textsc{Lösung}}\par\smallskip\small
\begin{tabular}{@{}r@{\enspace}l@{\enspace}l@{}}
\textbf{1.} & \symking{}xd1 & e3\\
\textbf{2.} & \symrook{}f2 & \symking{}b1\\
\textbf{3.} & \symrook{}f4 & exf4\\
\textbf{4.} & e3 & f5\\
\textbf{5.} & e2 & f6\\
\textbf{6.} & e1=\symrook{} & f7\\
\textbf{7.} & \symrook{}e2 & f8=\symqueen{}\\
\textbf{8.} & \symrook{}d2 & \symqueen{}f1\#\\
\end{tabular}\par\vspace{3mm}
\qrcode[height=12mm]{https://helpman.komtera.lt/?fen=8/8/8/8/4p3/8/2r1P3/K1kR4&moves=8}\\{\sffamily\tiny\color{ink}scannen zum Lösen}
\end{minipage}
\vspace{2mm}

{\raggedright\noindent\textcolor{ink}{\sffamily\textsc{Themen}}\enspace \themetag{Umwandlung} \themetag{Unterverwandlung} \themetag{Excelsior} \themetag{Switchback} \themetag{Selbstblock} \themetag{Kniest} \themetag{Phönix} \themetag{ohne Schach} \themetag{Umnov}\par\smallskip{\footnotesize\textit{Hinweis:} die Partei am Zug steht im Diagramm im Schach; die Lösung beginnt mit einem Schlag.}\par}

\clearpage
\phantomsection\label{hm:KBBNvkq}
\begin{tcolorbox}[enhanced,colback=ink,colframe=ink,arc=0pt,boxrule=0pt,left=4mm,right=4mm,top=2.5mm,bottom=2.5mm]
\color{white}\sffamily{\LARGE\bfseries KBBNvkq}\hfill{\Large h\#7.5}\\[0.5mm]
{\small 6 Steine\quad Weiß am Zug}\hfill{\small Nr.~428\qquad\href{https://helpman.komtera.lt/?fen=6kN/8/7B/8/8/8/2B5/qK6&moves=7.5}{\color{white}\bfseries lösen\,\ensuremath{\nearrow}}}
\end{tcolorbox}
\vspace{2mm}
\noindent\begin{minipage}[t]{36mm}\vspace{0pt}\sffamily\small\raggedright
\hmfigure{längste eindeutige}{{\Large h\#7.5}}\hmfigure{eindeutig bei dieser Länge}{{\Large 1} Stellung}\hmfigure{längstes Matt}{{\Large h\#8.5}}\hmfigure{in der Klasse}{15,5 Milliarden Stellungen,\\ 4,14 Milliarden mit Hilfsmatt}
\end{minipage}\hfill
\begin{minipage}[t]{80mm}\vspace{0pt}\begin{tikzpicture}\begin{axis}[width=72mm,height=42mm,ybar,bar width=1.6pt,xmin=-1,xmax=18,ymin=0,ymax=1087,axis y line*=left,axis x line*=bottom,xlabel={\scriptsize Halbzüge bis Matt},ylabel={\scriptsize Stellungen (Mio.)},tick label style={font=\tiny},label style={font=\scriptsize},enlarge x limits=false,ylabel shift=-3pt,xlabel shift=-3pt]\addplot[fill=ink!55,draw=none] coordinates {(0,0.3537) (1,1.425) (2,7.631) (3,52.21) (4,195.5) (5,572.1) (6,961.5) (7,987.9) (8,706.4) (9,432.3) (10,162.8) (11,51.94) (12,3.971) (13,1.244) (14,0.03356) (15,0.07914) (16,0.000519) (17,6e-06)};\draw[densely dotted] (axis cs:15,0) -- (axis cs:15,1087);\end{axis}\begin{axis}[width=72mm,height=42mm,axis y line*=right,axis x line=none,ymode=log,xmin=-1,xmax=18,ylabel={\scriptsize eindeutig (log)},tick label style={font=\tiny},label style={font=\scriptsize},ylabel shift=-3pt]\addplot[mark=*,mark size=0.7pt,draw=black,thin] coordinates {(0,353674) (1,1343501) (2,4547369) (3,21418432) (4,45741048) (5,47283309) (6,23700539) (7,9901719) (8,2366201) (9,637450) (10,35906) (11,5071) (12,97) (13,28) (15,1)};\end{axis}\end{tikzpicture}\end{minipage}
\vspace{3mm}

\noindent\begin{minipage}[t]{60mm}\vspace{0pt}\sbox\hmboard{\chessboard[setfen=6kN/8/7B/8/8/8/2B5/qK6 w - - 0 1]}\usebox\hmboard\\[1mm]
\makebox[\wd\hmboard]{\sffamily\hspace{3mm}\textbf{h\#7.5}\hfill(4+2)}\end{minipage}\hfill
\begin{minipage}[t]{52mm}\vspace{0pt}
\textcolor{ink}{\sffamily\textsc{Lösung}}\par\smallskip\small
\begin{tabular}{@{}r@{\enspace}l@{\enspace}l@{}}
\textbf{1.} & \ldots & \symking{}xa1\\
\textbf{2.} & \symking{}xh8 & \symking{}b2\\
\textbf{3.} & \symking{}g8 & \symking{}c3\\
\textbf{4.} & \symking{}f7 & \symking{}d3\\
\textbf{5.} & \symking{}g6 & \symking{}e4\\
\textbf{6.} & \symking{}h5 & \symking{}f5\\
\textbf{7.} & \symking{}h4 & \symbishop{}g5+\\
\textbf{8.} & \symking{}h5 & \symbishop{}d1\#\\
\end{tabular}\par\vspace{3mm}
\qrcode[height=12mm]{https://helpman.komtera.lt/?fen=6kN/8/7B/8/8/8/2B5/qK6&moves=7.5}\\{\sffamily\tiny\color{ink}scannen zum Lösen}
\end{minipage}
\vspace{2mm}

{\raggedright\noindent\textcolor{ink}{\sffamily\textsc{Themen}}\enspace \themetag{Switchback} \themetag{ein Stein zieht}\par\smallskip{\footnotesize\textit{Hinweis:} die Partei am Zug steht im Diagramm im Schach; die Lösung beginnt mit einem Schlag.}\par}

\clearpage
\phantomsection\label{hm:KBBPvkq}
\begin{tcolorbox}[enhanced,colback=ink,colframe=ink,arc=0pt,boxrule=0pt,left=4mm,right=4mm,top=2.5mm,bottom=2.5mm]
\color{white}\sffamily{\LARGE\bfseries KBBPvkq}\hfill{\Large h\#7.5}\\[0.5mm]
{\small 6 Steine\quad Weiß am Zug}\hfill{\small Nr.~429\qquad\href{https://helpman.komtera.lt/?fen=8/6B1/8/6B1/8/6P1/8/Kq1k4&moves=7.5}{\color{white}\bfseries lösen\,\ensuremath{\nearrow}}}
\end{tcolorbox}
\vspace{2mm}
\noindent\begin{minipage}[t]{36mm}\vspace{0pt}\sffamily\small\raggedright
\hmfigure{längste eindeutige}{{\Large h\#7.5}}\hmfigure{eindeutig bei dieser Länge}{{\Large 195} Stellungen}\hmfigure{längstes Matt}{{\Large h\#9.5}}\hmfigure{in der Klasse}{45,4 Milliarden Stellungen,\\ 12,8 Milliarden mit Hilfsmatt}
\end{minipage}\hfill
\begin{minipage}[t]{80mm}\vspace{0pt}\begin{tikzpicture}\begin{axis}[width=72mm,height=42mm,ybar,bar width=1.6pt,xmin=-1,xmax=20,ymin=0,ymax=2.486,axis y line*=left,axis x line*=bottom,xlabel={\scriptsize Halbzüge bis Matt},ylabel={\scriptsize Stellungen (Mrd.)},tick label style={font=\tiny},label style={font=\scriptsize},enlarge x limits=false,ylabel shift=-3pt,xlabel shift=-3pt]\addplot[fill=ink!55,draw=none] coordinates {(0,0.000763) (1,0.005213) (2,0.02832) (3,0.2346) (4,0.7515) (5,1.8) (6,2.26) (7,2.217) (8,1.603) (9,1.326) (10,0.9111) (11,0.8476) (12,0.5524) (13,0.2201) (14,0.02966) (15,0.006625) (16,1.788e-05) (17,6.287e-05) (19,2.1e-08)};\draw[densely dotted] (axis cs:15,0) -- (axis cs:15,2.486);\end{axis}\begin{axis}[width=72mm,height=42mm,axis y line*=right,axis x line=none,ymode=log,xmin=-1,xmax=20,ylabel={\scriptsize eindeutig (log)},tick label style={font=\tiny},label style={font=\scriptsize},ylabel shift=-3pt]\addplot[mark=*,mark size=0.7pt,draw=black,thin] coordinates {(0,763045) (1,3905653) (2,14434642) (3,87334377) (4,139728127) (5,131502057) (6,57567938) (7,33503712) (8,11924345) (9,8800057) (10,4465667) (11,2212343) (12,415587) (13,82917) (14,384) (15,195)};\end{axis}\end{tikzpicture}\end{minipage}
\vspace{3mm}

\noindent\begin{minipage}[t]{60mm}\vspace{0pt}\sbox\hmboard{\chessboard[setfen=8/6B1/8/6B1/8/6P1/8/Kq1k4 w - - 0 1]}\usebox\hmboard\\[1mm]
\makebox[\wd\hmboard]{\sffamily\hspace{3mm}\textbf{h\#7.5}\hfill(4+2)}\end{minipage}\hfill
\begin{minipage}[t]{52mm}\vspace{0pt}
\textcolor{ink}{\sffamily\textsc{Lösung}}\par\smallskip\small
\begin{tabular}{@{}r@{\enspace}l@{\enspace}l@{}}
\textbf{1.} & \ldots & \symking{}xb1\\
\textbf{2.} & \symking{}e2 & \symking{}c2\\
\textbf{3.} & \symking{}f3 & \symking{}d3\\
\textbf{4.} & \symking{}g4 & \symking{}e4\\
\textbf{5.} & \symking{}xg5 & \symbishop{}f6+\\
\textbf{6.} & \symking{}h5 & \symking{}f5\\
\textbf{7.} & \symking{}h6 & \symbishop{}g5+\\
\textbf{8.} & \symking{}h5 & g4\#\\
\end{tabular}\par\vspace{3mm}
\qrcode[height=12mm]{https://helpman.komtera.lt/?fen=8/6B1/8/6B1/8/6P1/8/Kq1k4&moves=7.5}\\{\sffamily\tiny\color{ink}scannen zum Lösen}
\end{minipage}
\vspace{2mm}

{\raggedright\noindent\textcolor{ink}{\sffamily\textsc{Themen}}\enspace \themetag{reines Matt} \themetag{Mustermatt} \themetag{Idealmatt} \themetag{Switchback} \themetag{ein Stein zieht}\par\smallskip{\footnotesize\textit{Hinweis:} die Partei am Zug steht im Diagramm im Schach; die Lösung beginnt mit einem Schlag.}\par}

\clearpage
\phantomsection\label{hm:KBBPvkr}
\begin{tcolorbox}[enhanced,colback=ink,colframe=ink,arc=0pt,boxrule=0pt,left=4mm,right=4mm,top=2.5mm,bottom=2.5mm]
\color{white}\sffamily{\LARGE\bfseries KBBPvkr}\hfill{\Large h\#7.5}\\[0.5mm]
{\small 6 Steine\quad Weiß am Zug}\hfill{\small Nr.~430\qquad\href{https://helpman.komtera.lt/?fen=6B1/8/8/8/6B1/2k5/6P1/rK6&moves=7.5}{\color{white}\bfseries lösen\,\ensuremath{\nearrow}}}
\end{tcolorbox}
\vspace{2mm}
\noindent\begin{minipage}[t]{36mm}\vspace{0pt}\sffamily\small\raggedright
\hmfigure{längste eindeutige}{{\Large h\#7.5}}\hmfigure{eindeutig bei dieser Länge}{{\Large 12} Stellungen}\hmfigure{längstes Matt}{{\Large h\#9.5}}\hmfigure{in der Klasse}{45,4 Milliarden Stellungen,\\ 13,9 Milliarden mit Hilfsmatt}
\end{minipage}\hfill
\begin{minipage}[t]{80mm}\vspace{0pt}\begin{tikzpicture}\begin{axis}[width=72mm,height=42mm,ybar,bar width=1.6pt,xmin=-1,xmax=20,ymin=0,ymax=2.97,axis y line*=left,axis x line*=bottom,xlabel={\scriptsize Halbzüge bis Matt},ylabel={\scriptsize Stellungen (Mrd.)},tick label style={font=\tiny},label style={font=\scriptsize},enlarge x limits=false,ylabel shift=-3pt,xlabel shift=-3pt]\addplot[fill=ink!55,draw=none] coordinates {(0,0.00116) (1,0.008035) (2,0.03673) (3,0.2787) (4,0.9209) (5,1.941) (6,2.7) (7,2.166) (8,1.883) (9,1.294) (10,1.092) (11,0.826) (12,0.6253) (13,0.1458) (14,0.02709) (15,0.001468) (16,1.891e-05) (17,6.206e-06) (19,9e-09)};\draw[densely dotted] (axis cs:15,0) -- (axis cs:15,2.97);\end{axis}\begin{axis}[width=72mm,height=42mm,axis y line*=right,axis x line=none,ymode=log,xmin=-1,xmax=20,ylabel={\scriptsize eindeutig (log)},tick label style={font=\tiny},label style={font=\scriptsize},ylabel shift=-3pt]\addplot[mark=*,mark size=0.7pt,draw=black,thin] coordinates {(0,1160243) (1,5795539) (2,19108681) (3,105783080) (4,190727073) (5,145233360) (6,79829818) (7,32092287) (8,14979253) (9,8845706) (10,5729913) (11,1700945) (12,352823) (13,17677) (14,537) (15,12)};\end{axis}\end{tikzpicture}\end{minipage}
\vspace{3mm}

\noindent\begin{minipage}[t]{60mm}\vspace{0pt}\sbox\hmboard{\chessboard[setfen=6B1/8/8/8/6B1/2k5/6P1/rK6 w - - 0 1]}\usebox\hmboard\\[1mm]
\makebox[\wd\hmboard]{\sffamily\hspace{3mm}\textbf{h\#7.5}\hfill(4+2)}\end{minipage}\hfill
\begin{minipage}[t]{52mm}\vspace{0pt}
\textcolor{ink}{\sffamily\textsc{Lösung}}\par\smallskip\small
\begin{tabular}{@{}r@{\enspace}l@{\enspace}l@{}}
\textbf{1.} & \ldots & \symking{}xa1\\
\textbf{2.} & \symking{}d4 & \symbishop{}h5\\
\textbf{3.} & \symking{}e5 & g4\\
\textbf{4.} & \symking{}f6 & g5+\\
\textbf{5.} & \symking{}g7 & g6\\
\textbf{6.} & \symking{}xg8 & g7\\
\textbf{7.} & \symking{}h7 & g8=\symqueen{}+\\
\textbf{8.} & \symking{}h6 & \symqueen{}g6\#\\
\end{tabular}\par\vspace{3mm}
\qrcode[height=12mm]{https://helpman.komtera.lt/?fen=6B1/8/8/8/6B1/2k5/6P1/rK6&moves=7.5}\\{\sffamily\tiny\color{ink}scannen zum Lösen}
\end{minipage}
\vspace{2mm}

{\raggedright\noindent\textcolor{ink}{\sffamily\textsc{Themen}}\enspace \themetag{reines Matt} \themetag{Mustermatt} \themetag{Umwandlung} \themetag{Excelsior} \themetag{Rundlauf} \themetag{ein Stein zieht} \themetag{Umnov}\par\smallskip{\footnotesize\textit{Hinweis:} die Partei am Zug steht im Diagramm im Schach; die Lösung beginnt mit einem Schlag.}\par}

\clearpage
\phantomsection\label{hm:KBBvkqb}
\begin{tcolorbox}[enhanced,colback=ink,colframe=ink,arc=0pt,boxrule=0pt,left=4mm,right=4mm,top=2.5mm,bottom=2.5mm]
\color{white}\sffamily{\LARGE\bfseries KBBvkqb}\hfill{\Large h\#7.5}\\[0.5mm]
{\small 6 Steine\quad Weiß am Zug}\hfill{\small Nr.~431\qquad\href{https://helpman.komtera.lt/?fen=8/8/8/8/2Bb4/7B/8/1k1Kq3&moves=7.5}{\color{white}\bfseries lösen\,\ensuremath{\nearrow}}}
\end{tcolorbox}
\vspace{2mm}
\noindent\begin{minipage}[t]{36mm}\vspace{0pt}\sffamily\small\raggedright
\hmfigure{längste eindeutige}{{\Large h\#7.5}}\hmfigure{eindeutig bei dieser Länge}{{\Large 48} Stellungen}\hmfigure{längstes Matt}{{\Large h\#10.5}}\hmfigure{in der Klasse}{15,5 Milliarden Stellungen,\\ 3,05 Milliarden mit Hilfsmatt}
\end{minipage}\hfill
\begin{minipage}[t]{80mm}\vspace{0pt}\begin{tikzpicture}\begin{axis}[width=72mm,height=42mm,ybar,bar width=1.6pt,xmin=-1,xmax=22,ymin=0,ymax=734.6,axis y line*=left,axis x line*=bottom,xlabel={\scriptsize Halbzüge bis Matt},ylabel={\scriptsize Stellungen (Mio.)},tick label style={font=\tiny},label style={font=\scriptsize},enlarge x limits=false,ylabel shift=-3pt,xlabel shift=-3pt]\addplot[fill=ink!55,draw=none] coordinates {(0,0.1315) (1,0.6117) (2,4.191) (3,30.19) (4,111.3) (5,334.8) (6,515.6) (7,667.8) (8,421.7) (9,416.5) (10,208.4) (11,227.3) (12,78.72) (13,30.97) (14,1.331) (15,1.614) (16,0.000896) (17,0.6401) (19,0.218) (21,5.9e-05)};\draw[densely dotted] (axis cs:15,0) -- (axis cs:15,734.6);\end{axis}\begin{axis}[width=72mm,height=42mm,axis y line*=right,axis x line=none,ymode=log,xmin=-1,xmax=22,ylabel={\scriptsize eindeutig (log)},tick label style={font=\tiny},label style={font=\scriptsize},ylabel shift=-3pt]\addplot[mark=*,mark size=0.7pt,draw=black,thin] coordinates {(0,131524) (1,604575) (2,2602770) (3,13244160) (4,24287627) (5,28790998) (6,11002495) (7,6420468) (8,1794245) (9,1537042) (10,533316) (11,288492) (12,17071) (13,3501) (14,5) (15,48)};\end{axis}\end{tikzpicture}\end{minipage}
\vspace{3mm}

\noindent\begin{minipage}[t]{60mm}\vspace{0pt}\sbox\hmboard{\chessboard[setfen=8/8/8/8/2Bb4/7B/8/1k1Kq3 w - - 0 1]}\usebox\hmboard\\[1mm]
\makebox[\wd\hmboard]{\sffamily\hspace{3mm}\textbf{h\#7.5}\hfill(3+3)}\end{minipage}\hfill
\begin{minipage}[t]{52mm}\vspace{0pt}
\textcolor{ink}{\sffamily\textsc{Lösung}}\par\smallskip\small
\begin{tabular}{@{}r@{\enspace}l@{\enspace}l@{}}
\textbf{1.} & \ldots & \symking{}xe1\\
\textbf{2.} & \symking{}c2 & \symbishop{}d5\\
\textbf{3.} & \symking{}d3 & \symking{}d1\\
\textbf{4.} & \symking{}e3 & \symking{}c2\\
\textbf{5.} & \symking{}f2 & \symking{}d3\\
\textbf{6.} & \symking{}g1 & \symking{}e4\\
\textbf{7.} & \symking{}h1 & \symking{}f3\\
\textbf{8.} & \symbishop{}g1 & \symking{}g3\#\\
\end{tabular}\par\vspace{3mm}
\qrcode[height=12mm]{https://helpman.komtera.lt/?fen=8/8/8/8/2Bb4/7B/8/1k1Kq3&moves=7.5}\\{\sffamily\tiny\color{ink}scannen zum Lösen}
\end{minipage}
\vspace{2mm}

{\raggedright\noindent\textcolor{ink}{\sffamily\textsc{Themen}}\enspace \themetag{Switchback} \themetag{Selbstblock} \themetag{ohne Schach}\par\smallskip{\footnotesize\textit{Hinweis:} die Partei am Zug steht im Diagramm im Schach; die Lösung beginnt mit einem Schlag.}\par}

\clearpage
\phantomsection\label{hm:KBNvkqn}
\begin{tcolorbox}[enhanced,colback=ink,colframe=ink,arc=0pt,boxrule=0pt,left=4mm,right=4mm,top=2.5mm,bottom=2.5mm]
\color{white}\sffamily{\LARGE\bfseries KBNvkqn}\hfill{\Large h\#7.5}\\[0.5mm]
{\small 6 Steine\quad Weiß am Zug}\hfill{\small Nr.~432\qquad\href{https://helpman.komtera.lt/?fen=7N/7k/3B4/8/8/8/2q5/nK6&moves=7.5}{\color{white}\bfseries lösen\,\ensuremath{\nearrow}}}
\end{tcolorbox}
\vspace{2mm}
\noindent\begin{minipage}[t]{36mm}\vspace{0pt}\sffamily\small\raggedright
\hmfigure{längste eindeutige}{{\Large h\#7.5}}\hmfigure{eindeutig bei dieser Länge}{{\Large 8} Stellungen}\hmfigure{längstes Matt}{{\Large h\#12.5}}\hmfigure{in der Klasse}{15,5 Milliarden Stellungen,\\ 8,59 Milliarden mit Hilfsmatt}
\end{minipage}\hfill
\begin{minipage}[t]{80mm}\vspace{0pt}\begin{tikzpicture}\begin{axis}[width=72mm,height=42mm,ybar,bar width=1.6pt,xmin=-1,xmax=26,ymin=0,ymax=2.44,axis y line*=left,axis x line*=bottom,xlabel={\scriptsize Halbzüge bis Matt},ylabel={\scriptsize Stellungen (Mrd.)},tick label style={font=\tiny},label style={font=\scriptsize},enlarge x limits=false,ylabel shift=-3pt,xlabel shift=-3pt]\addplot[fill=ink!55,draw=none] coordinates {(0,6.415e-05) (1,0.0002642) (2,0.001198) (3,0.01008) (4,0.03753) (5,0.1998) (6,0.4911) (7,1.338) (8,1.685) (9,2.218) (10,1.289) (11,0.9763) (12,0.2445) (13,0.09173) (14,0.005385) (15,0.001328) (16,8.773e-06) (17,1.87e-06) (18,5.417e-06) (19,2.2e-08) (20,7.491e-06) (21,3e-08) (22,7.501e-06) (23,1.25e-07) (24,1.497e-06) (25,3e-08)};\draw[densely dotted] (axis cs:15,0) -- (axis cs:15,2.44);\end{axis}\begin{axis}[width=72mm,height=42mm,axis y line*=right,axis x line=none,ymode=log,xmin=-1,xmax=26,ylabel={\scriptsize eindeutig (log)},tick label style={font=\tiny},label style={font=\scriptsize},ylabel shift=-3pt]\addplot[mark=*,mark size=0.7pt,draw=black,thin] coordinates {(0,64150) (1,252305) (2,695362) (3,3973843) (4,9763082) (5,26578480) (6,27595253) (7,33234798) (8,12039920) (9,7382696) (10,1397100) (11,476240) (12,28169) (13,4689) (14,18) (15,8)};\end{axis}\end{tikzpicture}\end{minipage}
\vspace{3mm}

\noindent\begin{minipage}[t]{60mm}\vspace{0pt}\sbox\hmboard{\chessboard[setfen=7N/7k/3B4/8/8/8/2q5/nK6 w - - 0 1]}\usebox\hmboard\\[1mm]
\makebox[\wd\hmboard]{\sffamily\hspace{3mm}\textbf{h\#7.5}\hfill(3+3)}\end{minipage}\hfill
\begin{minipage}[t]{52mm}\vspace{0pt}
\textcolor{ink}{\sffamily\textsc{Lösung}}\par\smallskip\small
\begin{tabular}{@{}r@{\enspace}l@{\enspace}l@{}}
\textbf{1.} & \ldots & \symking{}xa1\\
\textbf{2.} & \symking{}h6 & \symknight{}g6\\
\textbf{3.} & \symking{}g5 & \symknight{}e5\\
\textbf{4.} & \symking{}f4 & \symknight{}d3+\\
\textbf{5.} & \symking{}e3 & \symknight{}b2\\
\textbf{6.} & \symking{}d2 & \symking{}a2\\
\textbf{7.} & \symqueen{}b1+ & \symking{}b3\\
\textbf{8.} & \symking{}c1 & \symbishop{}f4\#\\
\end{tabular}\par\vspace{3mm}
\qrcode[height=12mm]{https://helpman.komtera.lt/?fen=7N/7k/3B4/8/8/8/2q5/nK6&moves=7.5}\\{\sffamily\tiny\color{ink}scannen zum Lösen}
\end{minipage}
\vspace{2mm}

{\raggedright\noindent\textcolor{ink}{\sffamily\textsc{Themen}}\enspace \themetag{reines Matt} \themetag{Mustermatt} \themetag{Idealmatt} \themetag{Selbstblock}\par\smallskip{\footnotesize\textit{Hinweis:} die Partei am Zug steht im Diagramm im Schach; die Lösung beginnt mit einem Schlag.}\par}

\clearpage
\phantomsection\label{hm:KBNvkqq}
\begin{tcolorbox}[enhanced,colback=ink,colframe=ink,arc=0pt,boxrule=0pt,left=4mm,right=4mm,top=2.5mm,bottom=2.5mm]
\color{white}\sffamily{\LARGE\bfseries KBNvkqq}\hfill{\Large h\#7.5}\\[0.5mm]
{\small 6 Steine\quad Weiß am Zug}\hfill{\small Nr.~433\qquad\href{https://helpman.komtera.lt/?fen=7N/6k1/8/1q6/8/q7/8/K1B5&moves=7.5}{\color{white}\bfseries lösen\,\ensuremath{\nearrow}}}
\end{tcolorbox}
\vspace{2mm}
\noindent\begin{minipage}[t]{36mm}\vspace{0pt}\sffamily\small\raggedright
\hmfigure{längste eindeutige}{{\Large h\#7.5}}\hmfigure{eindeutig bei dieser Länge}{{\Large 42} Stellungen}\hmfigure{längstes Matt}{{\Large h\#8.5}}\hmfigure{in der Klasse}{15,5 Milliarden Stellungen,\\ 3,69 Milliarden mit Hilfsmatt}
\end{minipage}\hfill
\begin{minipage}[t]{80mm}\vspace{0pt}\begin{tikzpicture}\begin{axis}[width=72mm,height=42mm,ybar,bar width=1.6pt,xmin=-1,xmax=18,ymin=0,ymax=1.238,axis y line*=left,axis x line*=bottom,xlabel={\scriptsize Halbzüge bis Matt},ylabel={\scriptsize Stellungen (Mrd.)},tick label style={font=\tiny},label style={font=\scriptsize},enlarge x limits=false,ylabel shift=-3pt,xlabel shift=-3pt]\addplot[fill=ink!55,draw=none] coordinates {(0,1.478e-05) (1,5.785e-05) (2,0.0002665) (3,0.002466) (4,0.007453) (5,0.04991) (6,0.105) (7,0.4352) (8,0.5302) (9,1.126) (10,0.5662) (11,0.6571) (12,0.1215) (13,0.08488) (14,0.003386) (15,0.002577) (16,5.32e-07) (17,1.326e-05)};\draw[densely dotted] (axis cs:15,0) -- (axis cs:15,1.238);\end{axis}\begin{axis}[width=72mm,height=42mm,axis y line*=right,axis x line=none,ymode=log,xmin=-1,xmax=18,ylabel={\scriptsize eindeutig (log)},tick label style={font=\tiny},label style={font=\scriptsize},ylabel shift=-3pt]\addplot[mark=*,mark size=0.7pt,draw=black,thin] coordinates {(0,14781) (1,55852) (2,128047) (3,763878) (4,1401203) (5,6592616) (6,6895938) (7,15461065) (8,5065323) (9,5478589) (10,706730) (11,552651) (12,21446) (13,11668) (14,15) (15,42)};\end{axis}\end{tikzpicture}\end{minipage}
\vspace{3mm}

\noindent\begin{minipage}[t]{60mm}\vspace{0pt}\sbox\hmboard{\chessboard[setfen=7N/6k1/8/1q6/8/q7/8/K1B5 w - - 0 1]}\usebox\hmboard\\[1mm]
\makebox[\wd\hmboard]{\sffamily\hspace{3mm}\textbf{h\#7.5}\hfill(3+3)}\end{minipage}\hfill
\begin{minipage}[t]{52mm}\vspace{0pt}
\textcolor{ink}{\sffamily\textsc{Lösung}}\par\smallskip\small
\begin{tabular}{@{}r@{\enspace}l@{\enspace}l@{}}
\textbf{1.} & \ldots & \symbishop{}xa3\\
\textbf{2.} & \symking{}f6 & \symknight{}f7\\
\textbf{3.} & \symking{}e6 & \symknight{}d6\\
\textbf{4.} & \symking{}d5 & \symknight{}xb5\\
\textbf{5.} & \symking{}c4 & \symking{}b1\\
\textbf{6.} & \symking{}b3 & \symking{}c1\\
\textbf{7.} & \symking{}a2 & \symknight{}c3+\\
\textbf{8.} & \symking{}a1 & \symbishop{}b2\#\\
\end{tabular}\par\vspace{3mm}
\qrcode[height=12mm]{https://helpman.komtera.lt/?fen=7N/6k1/8/1q6/8/q7/8/K1B5&moves=7.5}\\{\sffamily\tiny\color{ink}scannen zum Lösen}
\end{minipage}
\vspace{2mm}

{\raggedright\noindent\textcolor{ink}{\sffamily\textsc{Themen}}\enspace \themetag{ein Stein zieht}\par\smallskip{\footnotesize\textit{Hinweis:} die Partei am Zug steht im Diagramm im Schach; die Lösung beginnt mit einem Schlag.}\par}

\clearpage
\phantomsection\label{hm:KBNvkqr}
\begin{tcolorbox}[enhanced,colback=ink,colframe=ink,arc=0pt,boxrule=0pt,left=4mm,right=4mm,top=2.5mm,bottom=2.5mm]
\color{white}\sffamily{\LARGE\bfseries KBNvkqr}\hfill{\Large h\#7.5}\\[0.5mm]
{\small 6 Steine\quad Weiß am Zug}\hfill{\small Nr.~434\qquad\href{https://helpman.komtera.lt/?fen=7N/6k1/8/1q6/8/r7/8/K1B5&moves=7.5}{\color{white}\bfseries lösen\,\ensuremath{\nearrow}}}
\end{tcolorbox}
\vspace{2mm}
\noindent\begin{minipage}[t]{36mm}\vspace{0pt}\sffamily\small\raggedright
\hmfigure{längste eindeutige}{{\Large h\#7.5}}\hmfigure{eindeutig bei dieser Länge}{{\Large 28} Stellungen}\hmfigure{längstes Matt}{{\Large h\#8.5}}\hmfigure{in der Klasse}{15,5 Milliarden Stellungen,\\ 8,01 Milliarden mit Hilfsmatt}
\end{minipage}\hfill
\begin{minipage}[t]{80mm}\vspace{0pt}\begin{tikzpicture}\begin{axis}[width=72mm,height=42mm,ybar,bar width=1.6pt,xmin=-1,xmax=18,ymin=0,ymax=2.52,axis y line*=left,axis x line*=bottom,xlabel={\scriptsize Halbzüge bis Matt},ylabel={\scriptsize Stellungen (Mrd.)},tick label style={font=\tiny},label style={font=\scriptsize},enlarge x limits=false,ylabel shift=-3pt,xlabel shift=-3pt]\addplot[fill=ink!55,draw=none] coordinates {(0,4.782e-05) (1,0.0001939) (2,0.0009566) (3,0.008448) (4,0.03281) (5,0.1822) (6,0.4323) (7,1.303) (8,1.524) (9,2.291) (10,1.06) (11,0.9315) (12,0.1593) (13,0.07725) (14,0.002619) (15,0.001597) (16,3.5e-07) (17,7.026e-06)};\draw[densely dotted] (axis cs:15,0) -- (axis cs:15,2.52);\end{axis}\begin{axis}[width=72mm,height=42mm,axis y line*=right,axis x line=none,ymode=log,xmin=-1,xmax=18,ylabel={\scriptsize eindeutig (log)},tick label style={font=\tiny},label style={font=\scriptsize},ylabel shift=-3pt]\addplot[mark=*,mark size=0.7pt,draw=black,thin] coordinates {(0,47821) (1,185632) (2,570149) (3,3402558) (4,9106940) (5,27556871) (6,28514570) (7,41118598) (8,13290420) (9,10414410) (10,1581982) (11,843525) (12,40164) (13,12115) (14,13) (15,28)};\end{axis}\end{tikzpicture}\end{minipage}
\vspace{3mm}

\noindent\begin{minipage}[t]{60mm}\vspace{0pt}\sbox\hmboard{\chessboard[setfen=7N/6k1/8/1q6/8/r7/8/K1B5 w - - 0 1]}\usebox\hmboard\\[1mm]
\makebox[\wd\hmboard]{\sffamily\hspace{3mm}\textbf{h\#7.5}\hfill(3+3)}\end{minipage}\hfill
\begin{minipage}[t]{52mm}\vspace{0pt}
\textcolor{ink}{\sffamily\textsc{Lösung}}\par\smallskip\small
\begin{tabular}{@{}r@{\enspace}l@{\enspace}l@{}}
\textbf{1.} & \ldots & \symbishop{}xa3\\
\textbf{2.} & \symking{}f6 & \symknight{}f7\\
\textbf{3.} & \symking{}e6 & \symknight{}d6\\
\textbf{4.} & \symking{}d5 & \symknight{}xb5\\
\textbf{5.} & \symking{}c4 & \symking{}b1\\
\textbf{6.} & \symking{}b3 & \symking{}c1\\
\textbf{7.} & \symking{}a2 & \symknight{}c3+\\
\textbf{8.} & \symking{}a1 & \symbishop{}b2\#\\
\end{tabular}\par\vspace{3mm}
\qrcode[height=12mm]{https://helpman.komtera.lt/?fen=7N/6k1/8/1q6/8/r7/8/K1B5&moves=7.5}\\{\sffamily\tiny\color{ink}scannen zum Lösen}
\end{minipage}
\vspace{2mm}

{\raggedright\noindent\textcolor{ink}{\sffamily\textsc{Themen}}\enspace \themetag{ein Stein zieht}\par\smallskip{\footnotesize\textit{Hinweis:} die Partei am Zug steht im Diagramm im Schach; die Lösung beginnt mit einem Schlag.}\par}

\clearpage
\phantomsection\label{hm:KBNvkrr}
\begin{tcolorbox}[enhanced,colback=ink,colframe=ink,arc=0pt,boxrule=0pt,left=4mm,right=4mm,top=2.5mm,bottom=2.5mm]
\color{white}\sffamily{\LARGE\bfseries KBNvkrr}\hfill{\Large h\#7.5}\\[0.5mm]
{\small 6 Steine\quad Weiß am Zug}\hfill{\small Nr.~435\qquad\href{https://helpman.komtera.lt/?fen=7k/8/B7/7N/8/8/6r1/rK6&moves=7.5}{\color{white}\bfseries lösen\,\ensuremath{\nearrow}}}
\end{tcolorbox}
\vspace{2mm}
\noindent\begin{minipage}[t]{36mm}\vspace{0pt}\sffamily\small\raggedright
\hmfigure{längste eindeutige}{{\Large h\#7.5}}\hmfigure{eindeutig bei dieser Länge}{{\Large 1} Stellung}\hmfigure{längstes Matt}{{\Large h\#8.5}}\hmfigure{in der Klasse}{15,5 Milliarden Stellungen,\\ 4,32 Milliarden mit Hilfsmatt}
\end{minipage}\hfill
\begin{minipage}[t]{80mm}\vspace{0pt}\begin{tikzpicture}\begin{axis}[width=72mm,height=42mm,ybar,bar width=1.6pt,xmin=-1,xmax=18,ymin=0,ymax=1.209,axis y line*=left,axis x line*=bottom,xlabel={\scriptsize Halbzüge bis Matt},ylabel={\scriptsize Stellungen (Mrd.)},tick label style={font=\tiny},label style={font=\scriptsize},enlarge x limits=false,ylabel shift=-3pt,xlabel shift=-3pt]\addplot[fill=ink!55,draw=none] coordinates {(0,3.718e-05) (1,0.0001535) (2,0.0007606) (3,0.006455) (4,0.02539) (5,0.1252) (6,0.2996) (7,0.7454) (8,0.9135) (9,1.099) (10,0.587) (11,0.4057) (12,0.0828) (13,0.02716) (14,0.001423) (15,0.0003387) (16,2.34e-07) (17,7.56e-07)};\draw[densely dotted] (axis cs:15,0) -- (axis cs:15,1.209);\end{axis}\begin{axis}[width=72mm,height=42mm,axis y line*=right,axis x line=none,ymode=log,xmin=-1,xmax=18,ylabel={\scriptsize eindeutig (log)},tick label style={font=\tiny},label style={font=\scriptsize},ylabel shift=-3pt]\addplot[mark=*,mark size=0.7pt,draw=black,thin] coordinates {(0,37179) (1,145290) (2,493535) (3,2831923) (4,7300654) (5,18131018) (6,18621840) (7,18914146) (8,7091590) (9,4202696) (10,791863) (11,304190) (12,19582) (13,3502) (14,19) (15,1)};\end{axis}\end{tikzpicture}\end{minipage}
\vspace{3mm}

\noindent\begin{minipage}[t]{60mm}\vspace{0pt}\sbox\hmboard{\chessboard[setfen=7k/8/B7/7N/8/8/6r1/rK6 w - - 0 1]}\usebox\hmboard\\[1mm]
\makebox[\wd\hmboard]{\sffamily\hspace{3mm}\textbf{h\#7.5}\hfill(3+3)}\end{minipage}\hfill
\begin{minipage}[t]{52mm}\vspace{0pt}
\textcolor{ink}{\sffamily\textsc{Lösung}}\par\smallskip\small
\begin{tabular}{@{}r@{\enspace}l@{\enspace}l@{}}
\textbf{1.} & \ldots & \symking{}xa1\\
\textbf{2.} & \symking{}h7 & \symbishop{}f1\\
\textbf{3.} & \symking{}g6 & \symknight{}f4+\\
\textbf{4.} & \symking{}f5 & \symknight{}d3\\
\textbf{5.} & \symking{}e4 & \symbishop{}xg2+\\
\textbf{6.} & \symking{}e3 & \symking{}b2\\
\textbf{7.} & \symking{}e2 & \symking{}c3\\
\textbf{8.} & \symking{}d1 & \symbishop{}f3\#\\
\end{tabular}\par\vspace{3mm}
\qrcode[height=12mm]{https://helpman.komtera.lt/?fen=7k/8/B7/7N/8/8/6r1/rK6&moves=7.5}\\{\sffamily\tiny\color{ink}scannen zum Lösen}
\end{minipage}
\vspace{2mm}

{\raggedright\noindent\textcolor{ink}{\sffamily\textsc{Themen}}\enspace \themetag{reines Matt} \themetag{Mustermatt} \themetag{Idealmatt} \themetag{Spiegelmatt} \themetag{ein Stein zieht}\par\smallskip{\footnotesize\textit{Hinweis:} die Partei am Zug steht im Diagramm im Schach; die Lösung beginnt mit einem Schlag.}\par}

\clearpage
\phantomsection\label{hm:KBPvkqb}
\begin{tcolorbox}[enhanced,colback=ink,colframe=ink,arc=0pt,boxrule=0pt,left=4mm,right=4mm,top=2.5mm,bottom=2.5mm]
\color{white}\sffamily{\LARGE\bfseries KBPvkqb}\hfill{\Large h\#7.5}\\[0.5mm]
{\small 6 Steine\quad Weiß am Zug}\hfill{\small Nr.~436\qquad\href{https://helpman.komtera.lt/?fen=8/8/8/1B1b4/8/K1q5/P6k/8&moves=7.5}{\color{white}\bfseries lösen\,\ensuremath{\nearrow}}}
\end{tcolorbox}
\vspace{2mm}
\noindent\begin{minipage}[t]{36mm}\vspace{0pt}\sffamily\small\raggedright
\hmfigure{längste eindeutige}{{\Large h\#7.5}}\hmfigure{eindeutig bei dieser Länge}{{\Large 5} Stellungen}\hmfigure{längstes Matt}{{\Large h\#11}}\hmfigure{in der Klasse}{45,4 Milliarden Stellungen,\\ 26 Milliarden mit Hilfsmatt}
\end{minipage}\hfill
\begin{minipage}[t]{80mm}\vspace{0pt}\begin{tikzpicture}\begin{axis}[width=72mm,height=42mm,ybar,bar width=1.6pt,xmin=-1,xmax=23,ymin=0,ymax=5.133,axis y line*=left,axis x line*=bottom,xlabel={\scriptsize Halbzüge bis Matt},ylabel={\scriptsize Stellungen (Mrd.)},tick label style={font=\tiny},label style={font=\scriptsize},enlarge x limits=false,ylabel shift=-3pt,xlabel shift=-3pt]\addplot[fill=ink!55,draw=none] coordinates {(0,0.0001521) (1,0.002434) (2,0.01395) (3,0.1621) (4,0.568) (5,2.195) (6,2.634) (7,4.276) (8,3.197) (9,4.666) (10,2.849) (11,3.301) (12,1.349) (13,0.7102) (14,0.04385) (15,0.01044) (16,1.066e-05) (17,1.422e-06) (18,1e-09) (19,1.35e-07) (20,2.4e-08) (21,2e-08) (22,1e-09)};\draw[densely dotted] (axis cs:15,0) -- (axis cs:15,5.133);\end{axis}\begin{axis}[width=72mm,height=42mm,axis y line*=right,axis x line=none,ymode=log,xmin=-1,xmax=23,ylabel={\scriptsize eindeutig (log)},tick label style={font=\tiny},label style={font=\scriptsize},ylabel shift=-3pt]\addplot[mark=*,mark size=0.7pt,draw=black,thin] coordinates {(0,152122) (1,1358084) (2,6483478) (3,65738873) (4,122838342) (5,186868766) (6,85848816) (7,99569755) (8,38156052) (9,39465888) (10,11087860) (11,6281530) (12,486742) (13,139595) (14,521) (15,5)};\end{axis}\end{tikzpicture}\end{minipage}
\vspace{3mm}

\noindent\begin{minipage}[t]{60mm}\vspace{0pt}\sbox\hmboard{\chessboard[setfen=8/8/8/1B1b4/8/K1q5/P6k/8 w - - 0 1]}\usebox\hmboard\\[1mm]
\makebox[\wd\hmboard]{\sffamily\hspace{3mm}\textbf{h\#7.5}\hfill(3+3)}\end{minipage}\hfill
\begin{minipage}[t]{52mm}\vspace{0pt}
\textcolor{ink}{\sffamily\textsc{Lösung}}\par\smallskip\small
\begin{tabular}{@{}r@{\enspace}l@{\enspace}l@{}}
\textbf{1.} & \ldots & \symking{}a4\\
\textbf{2.} & \symking{}g3 & \symbishop{}d7\\
\textbf{3.} & \symking{}f4 & \symking{}b5\\
\textbf{4.} & \symking{}e5 & a4\\
\textbf{5.} & \symking{}d6 & a5\\
\textbf{6.} & \symking{}c7 & a6\\
\textbf{7.} & \symking{}b8 & \symking{}b6\\
\textbf{8.} & \symbishop{}a8 & a7\#\\
\end{tabular}\par\vspace{3mm}
\qrcode[height=12mm]{https://helpman.komtera.lt/?fen=8/8/8/1B1b4/8/K1q5/P6k/8&moves=7.5}\\{\sffamily\tiny\color{ink}scannen zum Lösen}
\end{minipage}
\vspace{2mm}

{\raggedright\noindent\textcolor{ink}{\sffamily\textsc{Themen}}\enspace \themetag{reines Matt} \themetag{Mustermatt} \themetag{Selbstblock} \themetag{ohne Schlag} \themetag{ohne Schach}\par\smallskip{\footnotesize\textit{Hinweis:} die Partei am Zug steht im Diagramm im Schach.}\par}

\clearpage
\phantomsection\label{hm:KBPvkqq}
\begin{tcolorbox}[enhanced,colback=ink,colframe=ink,arc=0pt,boxrule=0pt,left=4mm,right=4mm,top=2.5mm,bottom=2.5mm]
\color{white}\sffamily{\LARGE\bfseries KBPvkqq}\hfill{\Large h\#7.5}\\[0.5mm]
{\small 6 Steine\quad Weiß am Zug}\hfill{\small Nr.~437\qquad\href{https://helpman.komtera.lt/?fen=8/8/8/8/B7/8/PK6/qq3k2&moves=7.5}{\color{white}\bfseries lösen\,\ensuremath{\nearrow}}}
\end{tcolorbox}
\vspace{2mm}
\noindent\begin{minipage}[t]{36mm}\vspace{0pt}\sffamily\small\raggedright
\hmfigure{längste eindeutige}{{\Large h\#7.5}}\hmfigure{eindeutig bei dieser Länge}{{\Large 91} Stellungen}\hmfigure{längstes Matt}{{\Large h\#9}}\hmfigure{in der Klasse}{45,4 Milliarden Stellungen,\\ 11,5 Milliarden mit Hilfsmatt}
\end{minipage}\hfill
\begin{minipage}[t]{80mm}\vspace{0pt}\begin{tikzpicture}\begin{axis}[width=72mm,height=42mm,ybar,bar width=1.6pt,xmin=-1,xmax=19,ymin=0,ymax=2.159,axis y line*=left,axis x line*=bottom,xlabel={\scriptsize Halbzüge bis Matt},ylabel={\scriptsize Stellungen (Mrd.)},tick label style={font=\tiny},label style={font=\scriptsize},enlarge x limits=false,ylabel shift=-3pt,xlabel shift=-3pt]\addplot[fill=ink!55,draw=none] coordinates {(0,2.821e-05) (1,0.00059) (2,0.004016) (3,0.05658) (4,0.2025) (5,0.9027) (6,0.9178) (7,1.824) (8,0.949) (9,1.818) (10,1.021) (11,1.963) (12,0.9009) (13,0.8847) (14,0.03684) (15,0.01633) (16,2.24e-05) (17,2.492e-06) (18,4e-09)};\draw[densely dotted] (axis cs:15,0) -- (axis cs:15,2.159);\end{axis}\begin{axis}[width=72mm,height=42mm,axis y line*=right,axis x line=none,ymode=log,xmin=-1,xmax=19,ylabel={\scriptsize eindeutig (log)},tick label style={font=\tiny},label style={font=\scriptsize},ylabel shift=-3pt]\addplot[mark=*,mark size=0.7pt,draw=black,thin] coordinates {(0,28207) (1,330366) (2,1812162) (3,23270673) (4,43522359) (5,80508116) (6,25798774) (7,36385370) (8,10086861) (9,17159206) (10,6081311) (11,6314731) (12,466634) (13,229757) (14,611) (15,91)};\end{axis}\end{tikzpicture}\end{minipage}
\vspace{3mm}

\noindent\begin{minipage}[t]{60mm}\vspace{0pt}\sbox\hmboard{\chessboard[setfen=8/8/8/8/B7/8/PK6/qq3k2 w - - 0 1]}\usebox\hmboard\\[1mm]
\makebox[\wd\hmboard]{\sffamily\hspace{3mm}\textbf{h\#7.5}\hfill(3+3)}\end{minipage}\hfill
\begin{minipage}[t]{52mm}\vspace{0pt}
\textcolor{ink}{\sffamily\textsc{Lösung}}\par\smallskip\small
\begin{tabular}{@{}r@{\enspace}l@{\enspace}l@{}}
\textbf{1.} & \ldots & \symking{}a3\\
\textbf{2.} & \symqueen{}e1 & \symking{}b3\\
\textbf{3.} & \symking{}e2 & \symking{}c4\\
\textbf{4.} & \symking{}d2 & \symking{}d5\\
\textbf{5.} & \symking{}c3 & \symbishop{}b5\\
\textbf{6.} & \symking{}b4 & a3+\\
\textbf{7.} & \symking{}a5 & \symking{}c5\\
\textbf{8.} & \symqueen{}b4+ & axb4\#\\
\end{tabular}\par\vspace{3mm}
\qrcode[height=12mm]{https://helpman.komtera.lt/?fen=8/8/8/8/B7/8/PK6/qq3k2&moves=7.5}\\{\sffamily\tiny\color{ink}scannen zum Lösen}
\end{minipage}
\vspace{2mm}

{\raggedright\noindent\textcolor{ink}{\sffamily\textsc{Themen}}\enspace \themetag{reines Matt} \themetag{Mustermatt}\par\smallskip{\footnotesize\textit{Hinweis:} das Diagramm hat keinen legalen letzten Zug, kann also in keiner Partie entstehen.}\par}

\clearpage
\phantomsection\label{hm:KBPvkqr}
\begin{tcolorbox}[enhanced,colback=ink,colframe=ink,arc=0pt,boxrule=0pt,left=4mm,right=4mm,top=2.5mm,bottom=2.5mm]
\color{white}\sffamily{\LARGE\bfseries KBPvkqr}\hfill{\Large h\#7.5}\\[0.5mm]
{\small 6 Steine\quad Weiß am Zug}\hfill{\small Nr.~438\qquad\href{https://helpman.komtera.lt/?fen=8/8/B7/2r5/8/qK6/P7/7k&moves=7.5}{\color{white}\bfseries lösen\,\ensuremath{\nearrow}}}
\end{tcolorbox}
\vspace{2mm}
\noindent\begin{minipage}[t]{36mm}\vspace{0pt}\sffamily\small\raggedright
\hmfigure{längste eindeutige}{{\Large h\#7.5}}\hmfigure{eindeutig bei dieser Länge}{{\Large 33} Stellungen}\hmfigure{längstes Matt}{{\Large h\#9}}\hmfigure{in der Klasse}{45,4 Milliarden Stellungen,\\ 24,9 Milliarden mit Hilfsmatt}
\end{minipage}\hfill
\begin{minipage}[t]{80mm}\vspace{0pt}\begin{tikzpicture}\begin{axis}[width=72mm,height=42mm,ybar,bar width=1.6pt,xmin=-1,xmax=19,ymin=0,ymax=4.31,axis y line*=left,axis x line*=bottom,xlabel={\scriptsize Halbzüge bis Matt},ylabel={\scriptsize Stellungen (Mrd.)},tick label style={font=\tiny},label style={font=\scriptsize},enlarge x limits=false,ylabel shift=-3pt,xlabel shift=-3pt]\addplot[fill=ink!55,draw=none] coordinates {(0,8.164e-05) (1,0.001945) (2,0.0123) (3,0.1504) (4,0.5542) (5,2.13) (6,2.311) (7,3.846) (8,2.328) (9,3.845) (10,2.556) (11,3.918) (12,1.873) (13,1.298) (14,0.049) (15,0.01239) (16,3.109e-05) (17,1.725e-06) (18,1e-09)};\draw[densely dotted] (axis cs:15,0) -- (axis cs:15,4.31);\end{axis}\begin{axis}[width=72mm,height=42mm,axis y line*=right,axis x line=none,ymode=log,xmin=-1,xmax=19,ylabel={\scriptsize eindeutig (log)},tick label style={font=\tiny},label style={font=\scriptsize},ylabel shift=-3pt]\addplot[mark=*,mark size=0.7pt,draw=black,thin] coordinates {(0,81641) (1,1054337) (2,5707843) (3,61769634) (4,123829144) (5,186292929) (6,71061545) (7,80797837) (8,25126806) (9,35186791) (10,15674103) (11,9942226) (12,561765) (13,136813) (14,463) (15,33)};\end{axis}\end{tikzpicture}\end{minipage}
\vspace{3mm}

\noindent\begin{minipage}[t]{60mm}\vspace{0pt}\sbox\hmboard{\chessboard[setfen=8/8/B7/2r5/8/qK6/P7/7k w - - 0 1]}\usebox\hmboard\\[1mm]
\makebox[\wd\hmboard]{\sffamily\hspace{3mm}\textbf{h\#7.5}\hfill(3+3)}\end{minipage}\hfill
\begin{minipage}[t]{52mm}\vspace{0pt}
\textcolor{ink}{\sffamily\textsc{Lösung}}\par\smallskip\small
\begin{tabular}{@{}r@{\enspace}l@{\enspace}l@{}}
\textbf{1.} & \ldots & \symking{}xa3\\
\textbf{2.} & \symking{}g2 & \symking{}b4\\
\textbf{3.} & \symking{}f3 & a4\\
\textbf{4.} & \symrook{}b5+ & axb5\\
\textbf{5.} & \symking{}e4 & b6\\
\textbf{6.} & \symking{}d5 & b7\\
\textbf{7.} & \symking{}c6 & \symbishop{}b5+\\
\textbf{8.} & \symking{}b6 & b8=\symqueen{}\#\\
\end{tabular}\par\vspace{3mm}
\qrcode[height=12mm]{https://helpman.komtera.lt/?fen=8/8/B7/2r5/8/qK6/P7/7k&moves=7.5}\\{\sffamily\tiny\color{ink}scannen zum Lösen}
\end{minipage}
\vspace{2mm}

{\raggedright\noindent\textcolor{ink}{\sffamily\textsc{Themen}}\enspace \themetag{reines Matt} \themetag{Mustermatt} \themetag{Idealmatt} \themetag{Umwandlung} \themetag{Excelsior}\par\smallskip{\footnotesize\textit{Hinweis:} die Partei am Zug steht im Diagramm im Schach; die Lösung beginnt mit einem Schlag.}\par}

\clearpage
\phantomsection\label{hm:KBvkqqb}
\begin{tcolorbox}[enhanced,colback=ink,colframe=ink,arc=0pt,boxrule=0pt,left=4mm,right=4mm,top=2.5mm,bottom=2.5mm]
\color{white}\sffamily{\LARGE\bfseries KBvkqqb}\hfill{\Large h\#7.5}\\[0.5mm]
{\small 6 Steine\quad Weiß am Zug}\hfill{\small Nr.~439\qquad\href{https://helpman.komtera.lt/?fen=7b/8/8/4k3/8/8/B1q3q1/1K6&moves=7.5}{\color{white}\bfseries lösen\,\ensuremath{\nearrow}}}
\end{tcolorbox}
\vspace{2mm}
\noindent\begin{minipage}[t]{36mm}\vspace{0pt}\sffamily\small\raggedright
\hmfigure{längste eindeutige}{{\Large h\#7.5}}\hmfigure{eindeutig bei dieser Länge}{{\Large 8} Stellungen}\hmfigure{längstes Matt}{{\Large h\#8.5}}\hmfigure{in der Klasse}{15,5 Milliarden Stellungen,\\ 1,85 Milliarden mit Hilfsmatt}
\end{minipage}\hfill
\begin{minipage}[t]{80mm}\vspace{0pt}\begin{tikzpicture}\begin{axis}[width=72mm,height=42mm,ybar,bar width=1.6pt,xmin=-1,xmax=18,ymin=0,ymax=636,axis y line*=left,axis x line*=bottom,xlabel={\scriptsize Halbzüge bis Matt},ylabel={\scriptsize Stellungen (Mio.)},tick label style={font=\tiny},label style={font=\scriptsize},enlarge x limits=false,ylabel shift=-3pt,xlabel shift=-3pt]\addplot[fill=ink!55,draw=none] coordinates {(0,0.00169) (1,0.00663) (2,0.06767) (3,0.431) (4,1.902) (5,9.486) (6,22.29) (7,94.44) (8,123) (9,401.3) (10,257) (11,578.2) (12,167) (13,180.6) (14,8.783) (15,3.441) (16,0.007011) (17,0.000502)};\draw[densely dotted] (axis cs:15,0) -- (axis cs:15,636);\end{axis}\begin{axis}[width=72mm,height=42mm,axis y line*=right,axis x line=none,ymode=log,xmin=-1,xmax=18,ylabel={\scriptsize eindeutig (log)},tick label style={font=\tiny},label style={font=\scriptsize},ylabel shift=-3pt]\addplot[mark=*,mark size=0.7pt,draw=black,thin] coordinates {(0,1690) (1,6630) (2,46525) (3,236996) (4,525777) (5,1946384) (6,1625387) (7,5399410) (8,2786398) (9,6006057) (10,1463046) (11,1900840) (12,110386) (13,49357) (14,35) (15,8)};\end{axis}\end{tikzpicture}\end{minipage}
\vspace{3mm}

\noindent\begin{minipage}[t]{60mm}\vspace{0pt}\sbox\hmboard{\chessboard[setfen=7b/8/8/4k3/8/8/B1q3q1/1K6 w - - 0 1]}\usebox\hmboard\\[1mm]
\makebox[\wd\hmboard]{\sffamily\hspace{3mm}\textbf{h\#7.5}\hfill(2+4)}\end{minipage}\hfill
\begin{minipage}[t]{52mm}\vspace{0pt}
\textcolor{ink}{\sffamily\textsc{Lösung}}\par\smallskip\small
\begin{tabular}{@{}r@{\enspace}l@{\enspace}l@{}}
\textbf{1.} & \ldots & \symking{}a1\\
\textbf{2.} & \symking{}d4 & \symbishop{}d5\\
\textbf{3.} & \symking{}c3 & \symbishop{}xg2\\
\textbf{4.} & \symqueen{}c1+ & \symking{}a2\\
\textbf{5.} & \symking{}d2 & \symking{}b3\\
\textbf{6.} & \symbishop{}c3 & \symking{}c4\\
\textbf{7.} & \symking{}d1 & \symking{}d3\\
\textbf{8.} & \symbishop{}e1 & \symbishop{}f3\#\\
\end{tabular}\par\vspace{3mm}
\qrcode[height=12mm]{https://helpman.komtera.lt/?fen=7b/8/8/4k3/8/8/B1q3q1/1K6&moves=7.5}\\{\sffamily\tiny\color{ink}scannen zum Lösen}
\end{minipage}
\vspace{2mm}

{\raggedright\noindent\textcolor{ink}{\sffamily\textsc{Themen}}\enspace \themetag{Selbstblock}\par\smallskip{\footnotesize\textit{Hinweis:} die Partei am Zug steht im Diagramm im Schach.}\par}

\clearpage
\phantomsection\label{hm:KBvkqrb}
\begin{tcolorbox}[enhanced,colback=ink,colframe=ink,arc=0pt,boxrule=0pt,left=4mm,right=4mm,top=2.5mm,bottom=2.5mm]
\color{white}\sffamily{\LARGE\bfseries KBvkqrb}\hfill{\Large h\#7.5}\\[0.5mm]
{\small 6 Steine\quad Weiß am Zug}\hfill{\small Nr.~440\qquad\href{https://helpman.komtera.lt/?fen=8/7k/7b/8/8/8/B2r4/K1q5&moves=7.5}{\color{white}\bfseries lösen\,\ensuremath{\nearrow}}}
\end{tcolorbox}
\vspace{2mm}
\noindent\begin{minipage}[t]{36mm}\vspace{0pt}\sffamily\small\raggedright
\hmfigure{längste eindeutige}{{\Large h\#7.5}}\hmfigure{eindeutig bei dieser Länge}{{\Large 2} Stellungen}\hmfigure{längstes Matt}{{\Large h\#8.5}}\hmfigure{in der Klasse}{15,5 Milliarden Stellungen,\\ 4,02 Milliarden mit Hilfsmatt}
\end{minipage}\hfill
\begin{minipage}[t]{80mm}\vspace{0pt}\begin{tikzpicture}\begin{axis}[width=72mm,height=42mm,ybar,bar width=1.6pt,xmin=-1,xmax=18,ymin=0,ymax=1.303,axis y line*=left,axis x line*=bottom,xlabel={\scriptsize Halbzüge bis Matt},ylabel={\scriptsize Stellungen (Mrd.)},tick label style={font=\tiny},label style={font=\scriptsize},enlarge x limits=false,ylabel shift=-3pt,xlabel shift=-3pt]\addplot[fill=ink!55,draw=none] coordinates {(0,5.095e-06) (1,1.988e-05) (2,0.0001907) (3,0.001134) (4,0.005242) (5,0.02307) (6,0.05839) (7,0.2124) (8,0.3088) (9,0.8581) (10,0.6252) (11,1.185) (12,0.3945) (13,0.3227) (14,0.01785) (15,0.004341) (16,9.903e-06) (17,2.86e-07)};\draw[densely dotted] (axis cs:15,0) -- (axis cs:15,1.303);\end{axis}\begin{axis}[width=72mm,height=42mm,axis y line*=right,axis x line=none,ymode=log,xmin=-1,xmax=18,ylabel={\scriptsize eindeutig (log)},tick label style={font=\tiny},label style={font=\scriptsize},ylabel shift=-3pt]\addplot[mark=*,mark size=0.7pt,draw=black,thin] coordinates {(0,5095) (1,19878) (2,136843) (3,616264) (4,1524602) (5,4550461) (6,4246598) (7,11123538) (8,7142162) (9,12439232) (10,3410192) (11,3429244) (12,240292) (13,65581) (14,91) (15,2)};\end{axis}\end{tikzpicture}\end{minipage}
\vspace{3mm}

\noindent\begin{minipage}[t]{60mm}\vspace{0pt}\sbox\hmboard{\chessboard[setfen=8/7k/7b/8/8/8/B2r4/K1q5 w - - 0 1]}\usebox\hmboard\\[1mm]
\makebox[\wd\hmboard]{\sffamily\hspace{3mm}\textbf{h\#7.5}\hfill(2+4)}\end{minipage}\hfill
\begin{minipage}[t]{52mm}\vspace{0pt}
\textcolor{ink}{\sffamily\textsc{Lösung}}\par\smallskip\small
\begin{tabular}{@{}r@{\enspace}l@{\enspace}l@{}}
\textbf{1.} & \ldots & \symbishop{}b1+\\
\textbf{2.} & \symqueen{}c2 & \symbishop{}a2\\
\textbf{3.} & \symqueen{}b1+ & \symking{}xb1\\
\textbf{4.} & \symrook{}d1+ & \symking{}c2\\
\textbf{5.} & \symrook{}h1 & \symking{}d3\\
\textbf{6.} & \symking{}g6 & \symking{}e4\\
\textbf{7.} & \symking{}h5 & \symking{}f5\\
\textbf{8.} & \symrook{}h4 & \symbishop{}f7\#\\
\end{tabular}\par\vspace{3mm}
\qrcode[height=12mm]{https://helpman.komtera.lt/?fen=8/7k/7b/8/8/8/B2r4/K1q5&moves=7.5}\\{\sffamily\tiny\color{ink}scannen zum Lösen}
\end{minipage}
\vspace{2mm}

{\raggedright\noindent\textcolor{ink}{\sffamily\textsc{Themen}}\enspace \themetag{Switchback} \themetag{Selbstblock} \themetag{Umnov}\par\smallskip{\footnotesize\textit{Hinweis:} die Partei am Zug steht im Diagramm im Schach.}\par}

\clearpage
\phantomsection\label{hm:KBvkrrb}
\begin{tcolorbox}[enhanced,colback=ink,colframe=ink,arc=0pt,boxrule=0pt,left=4mm,right=4mm,top=2.5mm,bottom=2.5mm]
\color{white}\sffamily{\LARGE\bfseries KBvkrrb}\hfill{\Large h\#7.5}\\[0.5mm]
{\small 6 Steine\quad Weiß am Zug}\hfill{\small Nr.~441\qquad\href{https://helpman.komtera.lt/?fen=8/8/6k1/4b3/8/8/Br6/K1r5&moves=7.5}{\color{white}\bfseries lösen\,\ensuremath{\nearrow}}}
\end{tcolorbox}
\vspace{2mm}
\noindent\begin{minipage}[t]{36mm}\vspace{0pt}\sffamily\small\raggedright
\hmfigure{längste eindeutige}{{\Large h\#7.5}}\hmfigure{eindeutig bei dieser Länge}{{\Large 1} Stellung}\hmfigure{längstes Matt}{{\Large h\#8.5}}\hmfigure{in der Klasse}{15,5 Milliarden Stellungen,\\ 2,16 Milliarden mit Hilfsmatt}
\end{minipage}\hfill
\begin{minipage}[t]{80mm}\vspace{0pt}\begin{tikzpicture}\begin{axis}[width=72mm,height=42mm,ybar,bar width=1.6pt,xmin=-1,xmax=18,ymin=0,ymax=663.6,axis y line*=left,axis x line*=bottom,xlabel={\scriptsize Halbzüge bis Matt},ylabel={\scriptsize Stellungen (Mio.)},tick label style={font=\tiny},label style={font=\scriptsize},enlarge x limits=false,ylabel shift=-3pt,xlabel shift=-3pt]\addplot[fill=ink!55,draw=none] coordinates {(0,0.003679) (1,0.01431) (2,0.128) (3,0.7113) (4,3.444) (5,13.26) (6,36.17) (7,111.7) (8,185) (9,438.8) (10,375) (11,603.2) (12,235) (13,150.2) (14,10.7) (15,1.524) (16,0.003984) (17,4.5e-05)};\draw[densely dotted] (axis cs:15,0) -- (axis cs:15,663.6);\end{axis}\begin{axis}[width=72mm,height=42mm,axis y line*=right,axis x line=none,ymode=log,xmin=-1,xmax=18,ylabel={\scriptsize eindeutig (log)},tick label style={font=\tiny},label style={font=\scriptsize},ylabel shift=-3pt]\addplot[mark=*,mark size=0.7pt,draw=black,thin] coordinates {(0,3679) (1,14315) (2,95963) (3,375601) (4,1050725) (5,2460192) (6,2671822) (7,5391819) (8,4965934) (9,6697830) (10,2216856) (11,1652408) (12,146349) (13,20269) (14,115) (15,1)};\end{axis}\end{tikzpicture}\end{minipage}
\vspace{3mm}

\noindent\begin{minipage}[t]{60mm}\vspace{0pt}\sbox\hmboard{\chessboard[setfen=8/8/6k1/4b3/8/8/Br6/K1r5 w - - 0 1]}\usebox\hmboard\\[1mm]
\makebox[\wd\hmboard]{\sffamily\hspace{3mm}\textbf{h\#7.5}\hfill(2+4)}\end{minipage}\hfill
\begin{minipage}[t]{52mm}\vspace{0pt}
\textcolor{ink}{\sffamily\textsc{Lösung}}\par\smallskip\small
\begin{tabular}{@{}r@{\enspace}l@{\enspace}l@{}}
\textbf{1.} & \ldots & \symbishop{}b1+\\
\textbf{2.} & \symrook{}cc2 & \symbishop{}a2\\
\textbf{3.} & \symrook{}c3 & \symking{}xb2\\
\textbf{4.} & \symrook{}h3+ & \symking{}c2\\
\textbf{5.} & \symrook{}h4 & \symking{}d3\\
\textbf{6.} & \symbishop{}f4 & \symking{}e4\\
\textbf{7.} & \symking{}h5 & \symking{}f5\\
\textbf{8.} & \symbishop{}h6 & \symbishop{}f7\#\\
\end{tabular}\par\vspace{3mm}
\qrcode[height=12mm]{https://helpman.komtera.lt/?fen=8/8/6k1/4b3/8/8/Br6/K1r5&moves=7.5}\\{\sffamily\tiny\color{ink}scannen zum Lösen}
\end{minipage}
\vspace{2mm}

{\raggedright\noindent\textcolor{ink}{\sffamily\textsc{Themen}}\enspace \themetag{Switchback} \themetag{Selbstblock}\par\smallskip{\footnotesize\textit{Hinweis:} die Partei am Zug steht im Diagramm im Schach.}\par}

\clearpage
\phantomsection\label{hm:KNNPvkq}
\begin{tcolorbox}[enhanced,colback=ink,colframe=ink,arc=0pt,boxrule=0pt,left=4mm,right=4mm,top=2.5mm,bottom=2.5mm]
\color{white}\sffamily{\LARGE\bfseries KNNPvkq}\hfill{\Large h\#7.5}\\[0.5mm]
{\small 6 Steine\quad Weiß am Zug}\hfill{\small Nr.~442\qquad\href{https://helpman.komtera.lt/?fen=8/8/8/8/8/5N2/q4PN1/K6k&moves=7.5}{\color{white}\bfseries lösen\,\ensuremath{\nearrow}}}
\end{tcolorbox}
\vspace{2mm}
\noindent\begin{minipage}[t]{36mm}\vspace{0pt}\sffamily\small\raggedright
\hmfigure{längste eindeutige}{{\Large h\#7.5}}\hmfigure{eindeutig bei dieser Länge}{{\Large 6} Stellungen}\hmfigure{längstes Matt}{{\Large h\#9}}\hmfigure{in der Klasse}{45,4 Milliarden Stellungen,\\ 13,6 Milliarden mit Hilfsmatt}
\end{minipage}\hfill
\begin{minipage}[t]{80mm}\vspace{0pt}\begin{tikzpicture}\begin{axis}[width=72mm,height=42mm,ybar,bar width=1.6pt,xmin=-1,xmax=19,ymin=0,ymax=2.936,axis y line*=left,axis x line*=bottom,xlabel={\scriptsize Halbzüge bis Matt},ylabel={\scriptsize Stellungen (Mrd.)},tick label style={font=\tiny},label style={font=\scriptsize},enlarge x limits=false,ylabel shift=-3pt,xlabel shift=-3pt]\addplot[fill=ink!55,draw=none] coordinates {(0,0.0001719) (1,0.002333) (2,0.01073) (3,0.1208) (4,0.391) (5,1.284) (6,1.736) (7,2.669) (8,2.496) (9,2.561) (10,1.353) (11,0.7483) (12,0.1499) (13,0.03928) (14,0.001422) (15,0.0004949) (16,7.47e-07) (17,3.4e-08) (18,1e-09)};\draw[densely dotted] (axis cs:15,0) -- (axis cs:15,2.936);\end{axis}\begin{axis}[width=72mm,height=42mm,axis y line*=right,axis x line=none,ymode=log,xmin=-1,xmax=19,ylabel={\scriptsize eindeutig (log)},tick label style={font=\tiny},label style={font=\scriptsize},ylabel shift=-3pt]\addplot[mark=*,mark size=0.7pt,draw=black,thin] coordinates {(0,171917) (1,1251486) (2,4482557) (3,47019088) (4,83188788) (5,104210855) (6,60546792) (7,55331979) (8,20993490) (9,9884518) (10,1383502) (11,427160) (12,15726) (13,5924) (14,57) (15,6)};\end{axis}\end{tikzpicture}\end{minipage}
\vspace{3mm}

\noindent\begin{minipage}[t]{60mm}\vspace{0pt}\sbox\hmboard{\chessboard[setfen=8/8/8/8/8/5N2/q4PN1/K6k w - - 0 1]}\usebox\hmboard\\[1mm]
\makebox[\wd\hmboard]{\sffamily\hspace{3mm}\textbf{h\#7.5}\hfill(4+2)}\end{minipage}\hfill
\begin{minipage}[t]{52mm}\vspace{0pt}
\textcolor{ink}{\sffamily\textsc{Lösung}}\par\smallskip\small
\begin{tabular}{@{}r@{\enspace}l@{\enspace}l@{}}
\textbf{1.} & \ldots & \symking{}xa2\\
\textbf{2.} & \symking{}xg2 & \symknight{}d4\\
\textbf{3.} & \symking{}f1 & f4\\
\textbf{4.} & \symking{}e1 & f5\\
\textbf{5.} & \symking{}d2 & f6\\
\textbf{6.} & \symking{}c3 & f7\\
\textbf{7.} & \symking{}b4 & f8=\symqueen{}+\\
\textbf{8.} & \symking{}a4 & \symqueen{}a3\#\\
\end{tabular}\par\vspace{3mm}
\qrcode[height=12mm]{https://helpman.komtera.lt/?fen=8/8/8/8/8/5N2/q4PN1/K6k&moves=7.5}\\{\sffamily\tiny\color{ink}scannen zum Lösen}
\end{minipage}
\vspace{2mm}

{\raggedright\noindent\textcolor{ink}{\sffamily\textsc{Themen}}\enspace \themetag{Umwandlung} \themetag{Excelsior} \themetag{ein Stein zieht}\par\smallskip{\footnotesize\textit{Hinweis:} die Partei am Zug steht im Diagramm im Schach; die Lösung beginnt mit einem Schlag.}\par}

\clearpage
\phantomsection\label{hm:KNPPvkq}
\begin{tcolorbox}[enhanced,colback=ink,colframe=ink,arc=0pt,boxrule=0pt,left=4mm,right=4mm,top=2.5mm,bottom=2.5mm]
\color{white}\sffamily{\LARGE\bfseries KNPPvkq}\hfill{\Large h\#7.5}\\[0.5mm]
{\small 6 Steine\quad Weiß am Zug}\hfill{\small Nr.~443\qquad\href{https://helpman.komtera.lt/?fen=8/8/8/8/8/5N2/q4PP1/K6k&moves=7.5}{\color{white}\bfseries lösen\,\ensuremath{\nearrow}}}
\end{tcolorbox}
\vspace{2mm}
\noindent\begin{minipage}[t]{36mm}\vspace{0pt}\sffamily\small\raggedright
\hmfigure{längste eindeutige}{{\Large h\#7.5}}\hmfigure{eindeutig bei dieser Länge}{{\Large 5} Stellungen}\hmfigure{längstes Matt}{{\Large h\#8.5}}\hmfigure{in der Klasse}{34,1 Milliarden Stellungen,\\ 10,5 Milliarden mit Hilfsmatt}
\end{minipage}\hfill
\begin{minipage}[t]{80mm}\vspace{0pt}\begin{tikzpicture}\begin{axis}[width=72mm,height=42mm,ybar,bar width=1.6pt,xmin=-1,xmax=18,ymin=0,ymax=2.236,axis y line*=left,axis x line*=bottom,xlabel={\scriptsize Halbzüge bis Matt},ylabel={\scriptsize Stellungen (Mrd.)},tick label style={font=\tiny},label style={font=\scriptsize},enlarge x limits=false,ylabel shift=-3pt,xlabel shift=-3pt]\addplot[fill=ink!55,draw=none] coordinates {(0,6.857e-05) (1,0.002472) (2,0.01125) (3,0.1196) (4,0.3706) (5,1.278) (6,1.528) (7,2.032) (8,1.372) (9,1.458) (10,0.8876) (11,0.817) (12,0.4112) (13,0.1598) (14,0.01438) (15,0.002939) (16,2.92e-06) (17,1.732e-06)};\draw[densely dotted] (axis cs:15,0) -- (axis cs:15,2.236);\end{axis}\begin{axis}[width=72mm,height=42mm,axis y line*=right,axis x line=none,ymode=log,xmin=-1,xmax=18,ylabel={\scriptsize eindeutig (log)},tick label style={font=\tiny},label style={font=\scriptsize},ylabel shift=-3pt]\addplot[mark=*,mark size=0.7pt,draw=black,thin] coordinates {(0,68571) (1,915623) (2,3791977) (3,45235588) (4,77854192) (5,99688155) (6,40660032) (7,36978149) (8,12807589) (9,10897410) (10,3074058) (11,1781223) (12,214969) (13,43691) (14,50) (15,5)};\end{axis}\end{tikzpicture}\end{minipage}
\vspace{3mm}

\noindent\begin{minipage}[t]{60mm}\vspace{0pt}\sbox\hmboard{\chessboard[setfen=8/8/8/8/8/5N2/q4PP1/K6k w - - 0 1]}\usebox\hmboard\\[1mm]
\makebox[\wd\hmboard]{\sffamily\hspace{3mm}\textbf{h\#7.5}\hfill(4+2)}\end{minipage}\hfill
\begin{minipage}[t]{52mm}\vspace{0pt}
\textcolor{ink}{\sffamily\textsc{Lösung}}\par\smallskip\small
\begin{tabular}{@{}r@{\enspace}l@{\enspace}l@{}}
\textbf{1.} & \ldots & \symking{}xa2\\
\textbf{2.} & \symking{}xg2 & \symknight{}d4\\
\textbf{3.} & \symking{}f1 & f4\\
\textbf{4.} & \symking{}e1 & f5\\
\textbf{5.} & \symking{}d2 & f6\\
\textbf{6.} & \symking{}c3 & f7\\
\textbf{7.} & \symking{}b4 & f8=\symqueen{}+\\
\textbf{8.} & \symking{}a4 & \symqueen{}a3\#\\
\end{tabular}\par\vspace{3mm}
\qrcode[height=12mm]{https://helpman.komtera.lt/?fen=8/8/8/8/8/5N2/q4PP1/K6k&moves=7.5}\\{\sffamily\tiny\color{ink}scannen zum Lösen}
\end{minipage}
\vspace{2mm}

{\raggedright\noindent\textcolor{ink}{\sffamily\textsc{Themen}}\enspace \themetag{Umwandlung} \themetag{Excelsior} \themetag{ein Stein zieht}\par\smallskip{\footnotesize\textit{Hinweis:} die Partei am Zug steht im Diagramm im Schach; die Lösung beginnt mit einem Schlag.}\par}

\clearpage
\phantomsection\label{hm:KNPvkbb}
\begin{tcolorbox}[enhanced,colback=ink,colframe=ink,arc=0pt,boxrule=0pt,left=4mm,right=4mm,top=2.5mm,bottom=2.5mm]
\color{white}\sffamily{\LARGE\bfseries KNPvkbb}\hfill{\Large h\#7.5}\\[0.5mm]
{\small 6 Steine\quad Weiß am Zug}\hfill{\small Nr.~444\qquad\href{https://helpman.komtera.lt/?fen=8/7b/7k/N3b3/8/P7/8/K7&moves=7.5}{\color{white}\bfseries lösen\,\ensuremath{\nearrow}}}
\end{tcolorbox}
\vspace{2mm}
\noindent\begin{minipage}[t]{36mm}\vspace{0pt}\sffamily\small\raggedright
\hmfigure{längste eindeutige}{{\Large h\#7.5}}\hmfigure{eindeutig bei dieser Länge}{{\Large 20} Stellungen}\hmfigure{längstes Matt}{{\Large h\#10}}\hmfigure{in der Klasse}{45,4 Milliarden Stellungen,\\ 15 Milliarden mit Hilfsmatt}
\end{minipage}\hfill
\begin{minipage}[t]{80mm}\vspace{0pt}\begin{tikzpicture}\begin{axis}[width=72mm,height=42mm,ybar,bar width=1.6pt,xmin=-1,xmax=21,ymin=0,ymax=2.456,axis y line*=left,axis x line*=bottom,xlabel={\scriptsize Halbzüge bis Matt},ylabel={\scriptsize Stellungen (Mrd.)},tick label style={font=\tiny},label style={font=\scriptsize},enlarge x limits=false,ylabel shift=-3pt,xlabel shift=-3pt]\addplot[fill=ink!55,draw=none] coordinates {(0,5.031e-05) (1,0.001684) (2,0.006376) (3,0.075) (4,0.2328) (5,0.954) (6,1.34) (7,2.022) (8,1.8) (9,2.233) (10,1.895) (11,2.194) (12,1.475) (13,0.6075) (14,0.121) (15,0.008734) (16,0.000192) (17,1.043e-06) (18,1.317e-05) (20,6.12e-07)};\draw[densely dotted] (axis cs:15,0) -- (axis cs:15,2.456);\end{axis}\begin{axis}[width=72mm,height=42mm,axis y line*=right,axis x line=none,ymode=log,xmin=-1,xmax=21,ylabel={\scriptsize eindeutig (log)},tick label style={font=\tiny},label style={font=\scriptsize},ylabel shift=-3pt]\addplot[mark=*,mark size=0.7pt,draw=black,thin] coordinates {(0,50311) (1,629160) (2,2129823) (3,29033674) (4,52648515) (5,80212700) (6,47838534) (7,46865369) (8,22657769) (9,22372149) (10,9351718) (11,6944421) (12,1312181) (13,149740) (14,3965) (15,20)};\end{axis}\end{tikzpicture}\end{minipage}
\vspace{3mm}

\noindent\begin{minipage}[t]{60mm}\vspace{0pt}\sbox\hmboard{\chessboard[setfen=8/7b/7k/N3b3/8/P7/8/K7 w - - 0 1]}\usebox\hmboard\\[1mm]
\makebox[\wd\hmboard]{\sffamily\hspace{3mm}\textbf{h\#7.5}\hfill(3+3)}\end{minipage}\hfill
\begin{minipage}[t]{52mm}\vspace{0pt}
\textcolor{ink}{\sffamily\textsc{Lösung}}\par\smallskip\small
\begin{tabular}{@{}r@{\enspace}l@{\enspace}l@{}}
\textbf{1.} & \ldots & \symking{}a2\\
\textbf{2.} & \symbishop{}d3 & a4\\
\textbf{3.} & \symbishop{}b5 & axb5\\
\textbf{4.} & \symking{}g5 & b6\\
\textbf{5.} & \symking{}f4 & b7\\
\textbf{6.} & \symking{}e3 & b8=\symqueen{}\\
\textbf{7.} & \symking{}d2 & \symknight{}c4+\\
\textbf{8.} & \symking{}c1 & \symqueen{}b1\#\\
\end{tabular}\par\vspace{3mm}
\qrcode[height=12mm]{https://helpman.komtera.lt/?fen=8/7b/7k/N3b3/8/P7/8/K7&moves=7.5}\\{\sffamily\tiny\color{ink}scannen zum Lösen}
\end{minipage}
\vspace{2mm}

{\raggedright\noindent\textcolor{ink}{\sffamily\textsc{Themen}}\enspace \themetag{Umwandlung}\par\smallskip{\footnotesize\textit{Hinweis:} die Partei am Zug steht im Diagramm im Schach.}\par}

\clearpage
\phantomsection\label{hm:KNPvkbn}
\begin{tcolorbox}[enhanced,colback=ink,colframe=ink,arc=0pt,boxrule=0pt,left=4mm,right=4mm,top=2.5mm,bottom=2.5mm]
\color{white}\sffamily{\LARGE\bfseries KNPvkbn}\hfill{\Large h\#7.5}\\[0.5mm]
{\small 6 Steine\quad Weiß am Zug}\hfill{\small Nr.~445\qquad\href{https://helpman.komtera.lt/?fen=8/7k/8/8/1b6/1n6/1P6/KN6&moves=7.5}{\color{white}\bfseries lösen\,\ensuremath{\nearrow}}}
\end{tcolorbox}
\vspace{2mm}
\noindent\begin{minipage}[t]{36mm}\vspace{0pt}\sffamily\small\raggedright
\hmfigure{längste eindeutige}{{\Large h\#7.5}}\hmfigure{eindeutig bei dieser Länge}{{\Large 14} Stellungen}\hmfigure{längstes Matt}{{\Large h\#9}}\hmfigure{in der Klasse}{45,4 Milliarden Stellungen,\\ 30,8 Milliarden mit Hilfsmatt}
\end{minipage}\hfill
\begin{minipage}[t]{80mm}\vspace{0pt}\begin{tikzpicture}\begin{axis}[width=72mm,height=42mm,ybar,bar width=1.6pt,xmin=-1,xmax=19,ymin=0,ymax=5.378,axis y line*=left,axis x line*=bottom,xlabel={\scriptsize Halbzüge bis Matt},ylabel={\scriptsize Stellungen (Mrd.)},tick label style={font=\tiny},label style={font=\scriptsize},enlarge x limits=false,ylabel shift=-3pt,xlabel shift=-3pt]\addplot[fill=ink!55,draw=none] coordinates {(0,0.0001027) (1,0.004158) (2,0.01392) (3,0.1507) (4,0.4399) (5,1.846) (6,2.726) (7,4.062) (8,4.139) (9,4.889) (10,4.501) (11,4.286) (12,2.543) (13,0.9443) (14,0.2029) (15,0.01213) (16,0.0003135) (17,1.244e-06) (18,1.579e-06)};\draw[densely dotted] (axis cs:15,0) -- (axis cs:15,5.378);\end{axis}\begin{axis}[width=72mm,height=42mm,axis y line*=right,axis x line=none,ymode=log,xmin=-1,xmax=19,ylabel={\scriptsize eindeutig (log)},tick label style={font=\tiny},label style={font=\scriptsize},ylabel shift=-3pt]\addplot[mark=*,mark size=0.7pt,draw=black,thin] coordinates {(0,102738) (1,1358971) (2,4037646) (3,54609892) (4,99248702) (5,155199158) (6,101126461) (7,90983676) (8,51571972) (9,43493019) (10,17460074) (11,12095844) (12,2205922) (13,218936) (14,4983) (15,14)};\end{axis}\end{tikzpicture}\end{minipage}
\vspace{3mm}

\noindent\begin{minipage}[t]{60mm}\vspace{0pt}\sbox\hmboard{\chessboard[setfen=8/7k/8/8/1b6/1n6/1P6/KN6 w - - 0 1]}\usebox\hmboard\\[1mm]
\makebox[\wd\hmboard]{\sffamily\hspace{3mm}\textbf{h\#7.5}\hfill(3+3)}\end{minipage}\hfill
\begin{minipage}[t]{52mm}\vspace{0pt}
\textcolor{ink}{\sffamily\textsc{Lösung}}\par\smallskip\small
\begin{tabular}{@{}r@{\enspace}l@{\enspace}l@{}}
\textbf{1.} & \ldots & \symking{}a2\\
\textbf{2.} & \symbishop{}c3 & bxc3\\
\textbf{3.} & \symknight{}d4 & cxd4\\
\textbf{4.} & \symking{}g6 & d5\\
\textbf{5.} & \symking{}f5 & d6\\
\textbf{6.} & \symking{}e4 & d7\\
\textbf{7.} & \symking{}d3 & d8=\symqueen{}+\\
\textbf{8.} & \symking{}c2 & \symqueen{}d2\#\\
\end{tabular}\par\vspace{3mm}
\qrcode[height=12mm]{https://helpman.komtera.lt/?fen=8/7k/8/8/1b6/1n6/1P6/KN6&moves=7.5}\\{\sffamily\tiny\color{ink}scannen zum Lösen}
\end{minipage}
\vspace{2mm}

{\raggedright\noindent\textcolor{ink}{\sffamily\textsc{Themen}}\enspace \themetag{Umwandlung} \themetag{Excelsior}\par\smallskip{\footnotesize\textit{Hinweis:} die Partei am Zug steht im Diagramm im Schach.}\par}

\clearpage
\phantomsection\label{hm:KNPvknn}
\begin{tcolorbox}[enhanced,colback=ink,colframe=ink,arc=0pt,boxrule=0pt,left=4mm,right=4mm,top=2.5mm,bottom=2.5mm]
\color{white}\sffamily{\LARGE\bfseries KNPvknn}\hfill{\Large h\#7.5}\\[0.5mm]
{\small 6 Steine\quad Weiß am Zug}\hfill{\small Nr.~446\qquad\href{https://helpman.komtera.lt/?fen=7k/8/8/8/8/n7/Pn6/NK6&moves=7.5}{\color{white}\bfseries lösen\,\ensuremath{\nearrow}}}
\end{tcolorbox}
\vspace{2mm}
\noindent\begin{minipage}[t]{36mm}\vspace{0pt}\sffamily\small\raggedright
\hmfigure{längste eindeutige}{{\Large h\#7.5}}\hmfigure{eindeutig bei dieser Länge}{{\Large 13} Stellungen}\hmfigure{längstes Matt}{{\Large h\#9}}\hmfigure{in der Klasse}{45,4 Milliarden Stellungen,\\ 15,7 Milliarden mit Hilfsmatt}
\end{minipage}\hfill
\begin{minipage}[t]{80mm}\vspace{0pt}\begin{tikzpicture}\begin{axis}[width=72mm,height=42mm,ybar,bar width=1.6pt,xmin=-1,xmax=19,ymin=0,ymax=2.669,axis y line*=left,axis x line*=bottom,xlabel={\scriptsize Halbzüge bis Matt},ylabel={\scriptsize Stellungen (Mrd.)},tick label style={font=\tiny},label style={font=\scriptsize},enlarge x limits=false,ylabel shift=-3pt,xlabel shift=-3pt]\addplot[fill=ink!55,draw=none] coordinates {(0,5.143e-05) (1,0.002577) (2,0.007333) (3,0.07379) (4,0.1948) (5,0.8625) (6,1.318) (7,1.97) (8,2.141) (9,2.426) (10,2.362) (11,2.169) (12,1.449) (13,0.5822) (14,0.1686) (15,0.01077) (16,0.0004109) (17,2.878e-06) (18,8.6e-08)};\draw[densely dotted] (axis cs:15,0) -- (axis cs:15,2.669);\end{axis}\begin{axis}[width=72mm,height=42mm,axis y line*=right,axis x line=none,ymode=log,xmin=-1,xmax=19,ylabel={\scriptsize eindeutig (log)},tick label style={font=\tiny},label style={font=\scriptsize},ylabel shift=-3pt]\addplot[mark=*,mark size=0.7pt,draw=black,thin] coordinates {(0,51433) (1,712317) (2,1745637) (3,24690292) (4,39323622) (5,70438432) (6,46761116) (7,41063318) (8,25234974) (9,20404495) (10,8719054) (11,7104281) (12,1513441) (13,158750) (14,3970) (15,13)};\end{axis}\end{tikzpicture}\end{minipage}
\vspace{3mm}

\noindent\begin{minipage}[t]{60mm}\vspace{0pt}\sbox\hmboard{\chessboard[setfen=7k/8/8/8/8/n7/Pn6/NK6 w - - 0 1]}\usebox\hmboard\\[1mm]
\makebox[\wd\hmboard]{\sffamily\hspace{3mm}\textbf{h\#7.5}\hfill(3+3)}\end{minipage}\hfill
\begin{minipage}[t]{52mm}\vspace{0pt}
\textcolor{ink}{\sffamily\textsc{Lösung}}\par\smallskip\small
\begin{tabular}{@{}r@{\enspace}l@{\enspace}l@{}}
\textbf{1.} & \ldots & \symking{}c1\\
\textbf{2.} & \symknight{}ac4 & a4\\
\textbf{3.} & \symking{}g7 & a5\\
\textbf{4.} & \symking{}f6 & a6\\
\textbf{5.} & \symking{}e5 & a7\\
\textbf{6.} & \symking{}d4 & a8=\symqueen{}\\
\textbf{7.} & \symknight{}e5 & \symqueen{}f3\\
\textbf{8.} & \symknight{}bc4 & \symknight{}b3\#\\
\end{tabular}\par\vspace{3mm}
\qrcode[height=12mm]{https://helpman.komtera.lt/?fen=7k/8/8/8/8/n7/Pn6/NK6&moves=7.5}\\{\sffamily\tiny\color{ink}scannen zum Lösen}
\end{minipage}
\vspace{2mm}

{\raggedright\noindent\textcolor{ink}{\sffamily\textsc{Themen}}\enspace \themetag{reines Matt} \themetag{Mustermatt} \themetag{Umwandlung} \themetag{Excelsior} \themetag{Selbstblock} \themetag{ohne Schlag} \themetag{ohne Schach}\par\smallskip{\footnotesize\textit{Hinweis:} die Partei am Zug steht im Diagramm im Schach.}\par}

\clearpage
\phantomsection\label{hm:KNPvkqn}
\begin{tcolorbox}[enhanced,colback=ink,colframe=ink,arc=0pt,boxrule=0pt,left=4mm,right=4mm,top=2.5mm,bottom=2.5mm]
\color{white}\sffamily{\LARGE\bfseries KNPvkqn}\hfill{\Large h\#7.5}\\[0.5mm]
{\small 6 Steine\quad Weiß am Zug}\hfill{\small Nr.~447\qquad\href{https://helpman.komtera.lt/?fen=8/8/8/5k2/N7/P7/q6n/K7&moves=7.5}{\color{white}\bfseries lösen\,\ensuremath{\nearrow}}}
\end{tcolorbox}
\vspace{2mm}
\noindent\begin{minipage}[t]{36mm}\vspace{0pt}\sffamily\small\raggedright
\hmfigure{längste eindeutige}{{\Large h\#7.5}}\hmfigure{eindeutig bei dieser Länge}{{\Large 254} Stellungen}\hmfigure{längstes Matt}{{\Large h\#8.5}}\hmfigure{in der Klasse}{45,4 Milliarden Stellungen,\\ 27,5 Milliarden mit Hilfsmatt}
\end{minipage}\hfill
\begin{minipage}[t]{80mm}\vspace{0pt}\begin{tikzpicture}\begin{axis}[width=72mm,height=42mm,ybar,bar width=1.6pt,xmin=-1,xmax=18,ymin=0,ymax=5.508,axis y line*=left,axis x line*=bottom,xlabel={\scriptsize Halbzüge bis Matt},ylabel={\scriptsize Stellungen (Mrd.)},tick label style={font=\tiny},label style={font=\scriptsize},enlarge x limits=false,ylabel shift=-3pt,xlabel shift=-3pt]\addplot[fill=ink!55,draw=none] coordinates {(0,5.232e-05) (1,0.002153) (2,0.01002) (3,0.1182) (4,0.3681) (5,1.643) (6,2.2) (7,4.024) (8,3.342) (9,5.008) (10,3.523) (11,4.313) (12,1.767) (13,1.018) (14,0.09611) (15,0.0273) (16,7.339e-05) (17,4.471e-05)};\draw[densely dotted] (axis cs:15,0) -- (axis cs:15,5.508);\end{axis}\begin{axis}[width=72mm,height=42mm,axis y line*=right,axis x line=none,ymode=log,xmin=-1,xmax=18,ylabel={\scriptsize eindeutig (log)},tick label style={font=\tiny},label style={font=\scriptsize},ylabel shift=-3pt]\addplot[mark=*,mark size=0.7pt,draw=black,thin] coordinates {(0,52315) (1,795671) (2,3147156) (3,42997015) (4,76913795) (5,140756606) (6,65925975) (7,80232888) (8,32533131) (9,37581217) (10,11078799) (11,10320585) (12,1086292) (13,372615) (14,586) (15,254)};\end{axis}\end{tikzpicture}\end{minipage}
\vspace{3mm}

\noindent\begin{minipage}[t]{60mm}\vspace{0pt}\sbox\hmboard{\chessboard[setfen=8/8/8/5k2/N7/P7/q6n/K7 w - - 0 1]}\usebox\hmboard\\[1mm]
\makebox[\wd\hmboard]{\sffamily\hspace{3mm}\textbf{h\#7.5}\hfill(3+3)}\end{minipage}\hfill
\begin{minipage}[t]{52mm}\vspace{0pt}
\textcolor{ink}{\sffamily\textsc{Lösung}}\par\smallskip\small
\begin{tabular}{@{}r@{\enspace}l@{\enspace}l@{}}
\textbf{1.} & \ldots & \symking{}xa2\\
\textbf{2.} & \symknight{}f3 & \symknight{}c3\\
\textbf{3.} & \symknight{}d4 & a4\\
\textbf{4.} & \symknight{}b5 & axb5\\
\textbf{5.} & \symking{}f4 & b6\\
\textbf{6.} & \symking{}e3 & b7\\
\textbf{7.} & \symking{}d2 & b8=\symqueen{}\\
\textbf{8.} & \symking{}c1 & \symqueen{}b2\#\\
\end{tabular}\par\vspace{3mm}
\qrcode[height=12mm]{https://helpman.komtera.lt/?fen=8/8/8/5k2/N7/P7/q6n/K7&moves=7.5}\\{\sffamily\tiny\color{ink}scannen zum Lösen}
\end{minipage}
\vspace{2mm}

{\raggedright\noindent\textcolor{ink}{\sffamily\textsc{Themen}}\enspace \themetag{Umwandlung} \themetag{ohne Schach}\par\smallskip{\footnotesize\textit{Hinweis:} die Partei am Zug steht im Diagramm im Schach; die Lösung beginnt mit einem Schlag.}\par}

\clearpage
\phantomsection\label{hm:KNPvkqq}
\begin{tcolorbox}[enhanced,colback=ink,colframe=ink,arc=0pt,boxrule=0pt,left=4mm,right=4mm,top=2.5mm,bottom=2.5mm]
\color{white}\sffamily{\LARGE\bfseries KNPvkqq}\hfill{\Large h\#7.5}\\[0.5mm]
{\small 6 Steine\quad Weiß am Zug}\hfill{\small Nr.~448\qquad\href{https://helpman.komtera.lt/?fen=8/8/8/8/8/7N/1q5P/k1Kq4&moves=7.5}{\color{white}\bfseries lösen\,\ensuremath{\nearrow}}}
\end{tcolorbox}
\vspace{2mm}
\noindent\begin{minipage}[t]{36mm}\vspace{0pt}\sffamily\small\raggedright
\hmfigure{längste eindeutige}{{\Large h\#7.5}}\hmfigure{eindeutig bei dieser Länge}{{\Large 366} Stellungen}\hmfigure{längstes Matt}{{\Large h\#8.5}}\hmfigure{in der Klasse}{45,4 Milliarden Stellungen,\\ 11,9 Milliarden mit Hilfsmatt}
\end{minipage}\hfill
\begin{minipage}[t]{80mm}\vspace{0pt}\begin{tikzpicture}\begin{axis}[width=72mm,height=42mm,ybar,bar width=1.6pt,xmin=-1,xmax=18,ymin=0,ymax=2.341,axis y line*=left,axis x line*=bottom,xlabel={\scriptsize Halbzüge bis Matt},ylabel={\scriptsize Stellungen (Mrd.)},tick label style={font=\tiny},label style={font=\scriptsize},enlarge x limits=false,ylabel shift=-3pt,xlabel shift=-3pt]\addplot[fill=ink!55,draw=none] coordinates {(0,1.256e-05) (1,0.0004553) (2,0.002785) (3,0.04035) (4,0.1404) (5,0.7077) (6,0.7914) (7,1.801) (8,0.9511) (9,1.898) (10,0.9904) (11,2.128) (12,1.069) (13,1.227) (14,0.08725) (15,0.0464) (16,8.586e-05) (17,2.702e-05)};\draw[densely dotted] (axis cs:15,0) -- (axis cs:15,2.341);\end{axis}\begin{axis}[width=72mm,height=42mm,axis y line*=right,axis x line=none,ymode=log,xmin=-1,xmax=18,ylabel={\scriptsize eindeutig (log)},tick label style={font=\tiny},label style={font=\scriptsize},ylabel shift=-3pt]\addplot[mark=*,mark size=0.7pt,draw=black,thin] coordinates {(0,12563) (1,204930) (2,1036376) (3,15872597) (4,30788396) (5,67527986) (6,21624892) (7,34040905) (8,8624441) (9,16204511) (10,4946611) (11,7777341) (12,929886) (13,598547) (14,2413) (15,366)};\end{axis}\end{tikzpicture}\end{minipage}
\vspace{3mm}

\noindent\begin{minipage}[t]{60mm}\vspace{0pt}\sbox\hmboard{\chessboard[setfen=8/8/8/8/8/7N/1q5P/k1Kq4 w - - 0 1]}\usebox\hmboard\\[1mm]
\makebox[\wd\hmboard]{\sffamily\hspace{3mm}\textbf{h\#7.5}\hfill(3+3)}\end{minipage}\hfill
\begin{minipage}[t]{52mm}\vspace{0pt}
\textcolor{ink}{\sffamily\textsc{Lösung}}\par\smallskip\small
\begin{tabular}{@{}r@{\enspace}l@{\enspace}l@{}}
\textbf{1.} & \ldots & \symking{}xd1\\
\textbf{2.} & \symqueen{}g7 & \symknight{}f4\\
\textbf{3.} & \symking{}b2 & h4\\
\textbf{4.} & \symking{}c3 & h5\\
\textbf{5.} & \symking{}d4 & h6\\
\textbf{6.} & \symking{}e3 & hxg7\\
\textbf{7.} & \symking{}f2 & g8=\symqueen{}\\
\textbf{8.} & \symking{}f1 & \symqueen{}g2\#\\
\end{tabular}\par\vspace{3mm}
\qrcode[height=12mm]{https://helpman.komtera.lt/?fen=8/8/8/8/8/7N/1q5P/k1Kq4&moves=7.5}\\{\sffamily\tiny\color{ink}scannen zum Lösen}
\end{minipage}
\vspace{2mm}

{\raggedright\noindent\textcolor{ink}{\sffamily\textsc{Themen}}\enspace \themetag{Umwandlung} \themetag{Excelsior} \themetag{ohne Schach}\par\smallskip{\footnotesize\textit{Hinweis:} das Diagramm hat keinen legalen letzten Zug, kann also in keiner Partie entstehen.}\par}

\clearpage
\phantomsection\label{hm:KNPvkqr}
\begin{tcolorbox}[enhanced,colback=ink,colframe=ink,arc=0pt,boxrule=0pt,left=4mm,right=4mm,top=2.5mm,bottom=2.5mm]
\color{white}\sffamily{\LARGE\bfseries KNPvkqr}\hfill{\Large h\#7.5}\\[0.5mm]
{\small 6 Steine\quad Weiß am Zug}\hfill{\small Nr.~449\qquad\href{https://helpman.komtera.lt/?fen=8/8/8/8/8/7N/1q5P/k1Kr4&moves=7.5}{\color{white}\bfseries lösen\,\ensuremath{\nearrow}}}
\end{tcolorbox}
\vspace{2mm}
\noindent\begin{minipage}[t]{36mm}\vspace{0pt}\sffamily\small\raggedright
\hmfigure{längste eindeutige}{{\Large h\#7.5}}\hmfigure{eindeutig bei dieser Länge}{{\Large 115} Stellungen}\hmfigure{längstes Matt}{{\Large h\#8.5}}\hmfigure{in der Klasse}{45,4 Milliarden Stellungen,\\ 25,7 Milliarden mit Hilfsmatt}
\end{minipage}\hfill
\begin{minipage}[t]{80mm}\vspace{0pt}\begin{tikzpicture}\begin{axis}[width=72mm,height=42mm,ybar,bar width=1.6pt,xmin=-1,xmax=18,ymin=0,ymax=6.166,axis y line*=left,axis x line*=bottom,xlabel={\scriptsize Halbzüge bis Matt},ylabel={\scriptsize Stellungen (Mrd.)},tick label style={font=\tiny},label style={font=\scriptsize},enlarge x limits=false,ylabel shift=-3pt,xlabel shift=-3pt]\addplot[fill=ink!55,draw=none] coordinates {(0,4.305e-05) (1,0.001587) (2,0.009421) (3,0.1154) (4,0.4252) (5,1.842) (6,2.373) (7,4.676) (8,3.365) (9,5.606) (10,2.79) (11,3.31) (12,0.7151) (13,0.42) (14,0.006299) (15,0.009859) (16,2.519e-06) (17,1.294e-05)};\draw[densely dotted] (axis cs:15,0) -- (axis cs:15,6.166);\end{axis}\begin{axis}[width=72mm,height=42mm,axis y line*=right,axis x line=none,ymode=log,xmin=-1,xmax=18,ylabel={\scriptsize eindeutig (log)},tick label style={font=\tiny},label style={font=\scriptsize},ylabel shift=-3pt]\addplot[mark=*,mark size=0.7pt,draw=black,thin] coordinates {(0,43046) (1,715033) (2,3904514) (3,46013123) (4,97895911) (5,176208445) (6,82982011) (7,110815223) (8,35917436) (9,43191717) (10,8432841) (11,5173945) (12,79870) (13,97925) (14,65) (15,115)};\end{axis}\end{tikzpicture}\end{minipage}
\vspace{3mm}

\noindent\begin{minipage}[t]{60mm}\vspace{0pt}\sbox\hmboard{\chessboard[setfen=8/8/8/8/8/7N/1q5P/k1Kr4 w - - 0 1]}\usebox\hmboard\\[1mm]
\makebox[\wd\hmboard]{\sffamily\hspace{3mm}\textbf{h\#7.5}\hfill(3+3)}\end{minipage}\hfill
\begin{minipage}[t]{52mm}\vspace{0pt}
\textcolor{ink}{\sffamily\textsc{Lösung}}\par\smallskip\small
\begin{tabular}{@{}r@{\enspace}l@{\enspace}l@{}}
\textbf{1.} & \ldots & \symking{}xd1\\
\textbf{2.} & \symqueen{}g7 & \symknight{}f4\\
\textbf{3.} & \symking{}b2 & h4\\
\textbf{4.} & \symking{}c3 & h5\\
\textbf{5.} & \symking{}d4 & h6\\
\textbf{6.} & \symking{}e3 & hxg7\\
\textbf{7.} & \symking{}f2 & g8=\symqueen{}\\
\textbf{8.} & \symking{}f1 & \symqueen{}g2\#\\
\end{tabular}\par\vspace{3mm}
\qrcode[height=12mm]{https://helpman.komtera.lt/?fen=8/8/8/8/8/7N/1q5P/k1Kr4&moves=7.5}\\{\sffamily\tiny\color{ink}scannen zum Lösen}
\end{minipage}
\vspace{2mm}

{\raggedright\noindent\textcolor{ink}{\sffamily\textsc{Themen}}\enspace \themetag{Umwandlung} \themetag{Excelsior} \themetag{ohne Schach}\par\smallskip{\footnotesize\textit{Hinweis:} das Diagramm hat keinen legalen letzten Zug, kann also in keiner Partie entstehen.}\par}

\clearpage
\phantomsection\label{hm:KNPvkrb}
\begin{tcolorbox}[enhanced,colback=ink,colframe=ink,arc=0pt,boxrule=0pt,left=4mm,right=4mm,top=2.5mm,bottom=2.5mm]
\color{white}\sffamily{\LARGE\bfseries KNPvkrb}\hfill{\Large h\#7.5}\\[0.5mm]
{\small 6 Steine\quad Weiß am Zug}\hfill{\small Nr.~450\qquad\href{https://helpman.komtera.lt/?fen=8/8/7k/N7/8/r7/1Pb5/K7&moves=7.5}{\color{white}\bfseries lösen\,\ensuremath{\nearrow}}}
\end{tcolorbox}
\vspace{2mm}
\noindent\begin{minipage}[t]{36mm}\vspace{0pt}\sffamily\small\raggedright
\hmfigure{längste eindeutige}{{\Large h\#7.5}}\hmfigure{eindeutig bei dieser Länge}{{\Large 28} Stellungen}\hmfigure{längstes Matt}{{\Large h\#9.5}}\hmfigure{in der Klasse}{45,4 Milliarden Stellungen,\\ 28,9 Milliarden mit Hilfsmatt}
\end{minipage}\hfill
\begin{minipage}[t]{80mm}\vspace{0pt}\begin{tikzpicture}\begin{axis}[width=72mm,height=42mm,ybar,bar width=1.6pt,xmin=-1,xmax=20,ymin=0,ymax=6.116,axis y line*=left,axis x line*=bottom,xlabel={\scriptsize Halbzüge bis Matt},ylabel={\scriptsize Stellungen (Mrd.)},tick label style={font=\tiny},label style={font=\scriptsize},enlarge x limits=false,ylabel shift=-3pt,xlabel shift=-3pt]\addplot[fill=ink!55,draw=none] coordinates {(0,8.687e-05) (1,0.003127) (2,0.01474) (3,0.1603) (4,0.5711) (5,2.171) (6,3.091) (7,4.813) (8,4.321) (9,5.56) (10,3.73) (11,3.19) (12,1.016) (13,0.2801) (14,0.01485) (15,0.001661) (16,8.238e-06) (17,1.467e-06) (19,2e-09)};\draw[densely dotted] (axis cs:15,0) -- (axis cs:15,6.116);\end{axis}\begin{axis}[width=72mm,height=42mm,axis y line*=right,axis x line=none,ymode=log,xmin=-1,xmax=20,ylabel={\scriptsize eindeutig (log)},tick label style={font=\tiny},label style={font=\scriptsize},ylabel shift=-3pt]\addplot[mark=*,mark size=0.7pt,draw=black,thin] coordinates {(0,86865) (1,1283734) (2,5914032) (3,63374553) (4,135317664) (5,195396432) (6,110778850) (7,106147856) (8,48273985) (9,41961149) (10,9765510) (11,3453513) (12,154925) (13,20780) (14,47) (15,28)};\end{axis}\end{tikzpicture}\end{minipage}
\vspace{3mm}

\noindent\begin{minipage}[t]{60mm}\vspace{0pt}\sbox\hmboard{\chessboard[setfen=8/8/7k/N7/8/r7/1Pb5/K7 w - - 0 1]}\usebox\hmboard\\[1mm]
\makebox[\wd\hmboard]{\sffamily\hspace{3mm}\textbf{h\#7.5}\hfill(3+3)}\end{minipage}\hfill
\begin{minipage}[t]{52mm}\vspace{0pt}
\textcolor{ink}{\sffamily\textsc{Lösung}}\par\smallskip\small
\begin{tabular}{@{}r@{\enspace}l@{\enspace}l@{}}
\textbf{1.} & \ldots & bxa3\\
\textbf{2.} & \symbishop{}d3 & a4\\
\textbf{3.} & \symbishop{}b5 & axb5\\
\textbf{4.} & \symking{}g5 & b6\\
\textbf{5.} & \symking{}f4 & b7\\
\textbf{6.} & \symking{}e3 & b8=\symqueen{}\\
\textbf{7.} & \symking{}d2 & \symknight{}c4+\\
\textbf{8.} & \symking{}c1 & \symqueen{}b1\#\\
\end{tabular}\par\vspace{3mm}
\qrcode[height=12mm]{https://helpman.komtera.lt/?fen=8/8/7k/N7/8/r7/1Pb5/K7&moves=7.5}\\{\sffamily\tiny\color{ink}scannen zum Lösen}
\end{minipage}
\vspace{2mm}

{\raggedright\noindent\textcolor{ink}{\sffamily\textsc{Themen}}\enspace \themetag{Umwandlung} \themetag{Excelsior}\par\smallskip{\footnotesize\textit{Hinweis:} die Partei am Zug steht im Diagramm im Schach; die Lösung beginnt mit einem Schlag.}\par}

\clearpage
\phantomsection\label{hm:KNPvkrn}
\begin{tcolorbox}[enhanced,colback=ink,colframe=ink,arc=0pt,boxrule=0pt,left=4mm,right=4mm,top=2.5mm,bottom=2.5mm]
\color{white}\sffamily{\LARGE\bfseries KNPvkrn}\hfill{\Large h\#7.5}\\[0.5mm]
{\small 6 Steine\quad Weiß am Zug}\hfill{\small Nr.~451\qquad\href{https://helpman.komtera.lt/?fen=8/8/8/N6k/8/r7/1P1n4/K7&moves=7.5}{\color{white}\bfseries lösen\,\ensuremath{\nearrow}}}
\end{tcolorbox}
\vspace{2mm}
\noindent\begin{minipage}[t]{36mm}\vspace{0pt}\sffamily\small\raggedright
\hmfigure{längste eindeutige}{{\Large h\#7.5}}\hmfigure{eindeutig bei dieser Länge}{{\Large 38} Stellungen}\hmfigure{längstes Matt}{{\Large h\#9}}\hmfigure{in der Klasse}{45,4 Milliarden Stellungen,\\ 29,6 Milliarden mit Hilfsmatt}
\end{minipage}\hfill
\begin{minipage}[t]{80mm}\vspace{0pt}\begin{tikzpicture}\begin{axis}[width=72mm,height=42mm,ybar,bar width=1.6pt,xmin=-1,xmax=19,ymin=0,ymax=5.747,axis y line*=left,axis x line*=bottom,xlabel={\scriptsize Halbzüge bis Matt},ylabel={\scriptsize Stellungen (Mrd.)},tick label style={font=\tiny},label style={font=\scriptsize},enlarge x limits=false,ylabel shift=-3pt,xlabel shift=-3pt]\addplot[fill=ink!55,draw=none] coordinates {(0,8.484e-05) (1,0.003694) (2,0.01592) (3,0.1624) (4,0.5518) (5,2.097) (6,3.028) (7,4.579) (8,4.232) (9,5.225) (10,4.01) (11,3.607) (12,1.514) (13,0.5023) (14,0.04172) (15,0.003792) (16,3.156e-05) (17,3.756e-06) (18,1e-09)};\draw[densely dotted] (axis cs:15,0) -- (axis cs:15,5.747);\end{axis}\begin{axis}[width=72mm,height=42mm,axis y line*=right,axis x line=none,ymode=log,xmin=-1,xmax=19,ylabel={\scriptsize eindeutig (log)},tick label style={font=\tiny},label style={font=\scriptsize},ylabel shift=-3pt]\addplot[mark=*,mark size=0.7pt,draw=black,thin] coordinates {(0,84837) (1,1339228) (2,5670533) (3,59976425) (4,126207791) (5,185342016) (6,103427898) (7,93913674) (8,42839180) (9,37542896) (10,12246198) (11,5873144) (12,431733) (13,51204) (14,496) (15,38)};\end{axis}\end{tikzpicture}\end{minipage}
\vspace{3mm}

\noindent\begin{minipage}[t]{60mm}\vspace{0pt}\sbox\hmboard{\chessboard[setfen=8/8/8/N6k/8/r7/1P1n4/K7 w - - 0 1]}\usebox\hmboard\\[1mm]
\makebox[\wd\hmboard]{\sffamily\hspace{3mm}\textbf{h\#7.5}\hfill(3+3)}\end{minipage}\hfill
\begin{minipage}[t]{52mm}\vspace{0pt}
\textcolor{ink}{\sffamily\textsc{Lösung}}\par\smallskip\small
\begin{tabular}{@{}r@{\enspace}l@{\enspace}l@{}}
\textbf{1.} & \ldots & bxa3\\
\textbf{2.} & \symknight{}c4 & a4\\
\textbf{3.} & \symknight{}b6 & \symknight{}c4\\
\textbf{4.} & \symking{}g4 & a5\\
\textbf{5.} & \symking{}f3 & axb6\\
\textbf{6.} & \symking{}e2 & b7\\
\textbf{7.} & \symking{}d1 & b8=\symqueen{}\\
\textbf{8.} & \symking{}c1 & \symqueen{}b1\#\\
\end{tabular}\par\vspace{3mm}
\qrcode[height=12mm]{https://helpman.komtera.lt/?fen=8/8/8/N6k/8/r7/1P1n4/K7&moves=7.5}\\{\sffamily\tiny\color{ink}scannen zum Lösen}
\end{minipage}
\vspace{2mm}

{\raggedright\noindent\textcolor{ink}{\sffamily\textsc{Themen}}\enspace \themetag{Umwandlung} \themetag{Excelsior} \themetag{ohne Schach} \themetag{Umnov}\par\smallskip{\footnotesize\textit{Hinweis:} die Partei am Zug steht im Diagramm im Schach; die Lösung beginnt mit einem Schlag.}\par}

\clearpage
\phantomsection\label{hm:KNPvkrr}
\begin{tcolorbox}[enhanced,colback=ink,colframe=ink,arc=0pt,boxrule=0pt,left=4mm,right=4mm,top=2.5mm,bottom=2.5mm]
\color{white}\sffamily{\LARGE\bfseries KNPvkrr}\hfill{\Large h\#7.5}\\[0.5mm]
{\small 6 Steine\quad Weiß am Zug}\hfill{\small Nr.~452\qquad\href{https://helpman.komtera.lt/?fen=8/8/1r6/3N4/r7/K7/P7/7k&moves=7.5}{\color{white}\bfseries lösen\,\ensuremath{\nearrow}}}
\end{tcolorbox}
\vspace{2mm}
\noindent\begin{minipage}[t]{36mm}\vspace{0pt}\sffamily\small\raggedright
\hmfigure{längste eindeutige}{{\Large h\#7.5}}\hmfigure{eindeutig bei dieser Länge}{{\Large 18} Stellungen}\hmfigure{längstes Matt}{{\Large h\#8.5}}\hmfigure{in der Klasse}{45,4 Milliarden Stellungen,\\ 13,8 Milliarden mit Hilfsmatt}
\end{minipage}\hfill
\begin{minipage}[t]{80mm}\vspace{0pt}\begin{tikzpicture}\begin{axis}[width=72mm,height=42mm,ybar,bar width=1.6pt,xmin=-1,xmax=18,ymin=0,ymax=3.053,axis y line*=left,axis x line*=bottom,xlabel={\scriptsize Halbzüge bis Matt},ylabel={\scriptsize Stellungen (Mrd.)},tick label style={font=\tiny},label style={font=\scriptsize},enlarge x limits=false,ylabel shift=-3pt,xlabel shift=-3pt]\addplot[fill=ink!55,draw=none] coordinates {(0,3.411e-05) (1,0.001333) (2,0.007395) (3,0.07958) (4,0.2955) (5,1.106) (6,1.487) (7,2.474) (8,2.008) (9,2.775) (10,1.584) (11,1.465) (12,0.3641) (13,0.1328) (14,0.003186) (15,0.001361) (16,1.654e-06) (17,2.149e-06)};\draw[densely dotted] (axis cs:15,0) -- (axis cs:15,3.053);\end{axis}\begin{axis}[width=72mm,height=42mm,axis y line*=right,axis x line=none,ymode=log,xmin=-1,xmax=18,ylabel={\scriptsize eindeutig (log)},tick label style={font=\tiny},label style={font=\scriptsize},ylabel shift=-3pt]\addplot[mark=*,mark size=0.7pt,draw=black,thin] coordinates {(0,34106) (1,592029) (2,3229198) (3,32051911) (4,66788794) (5,99520608) (6,51808545) (7,55050519) (8,20982744) (9,19283673) (10,4645387) (11,1699917) (12,30234) (13,13294) (14,62) (15,18)};\end{axis}\end{tikzpicture}\end{minipage}
\vspace{3mm}

\noindent\begin{minipage}[t]{60mm}\vspace{0pt}\sbox\hmboard{\chessboard[setfen=8/8/1r6/3N4/r7/K7/P7/7k w - - 0 1]}\usebox\hmboard\\[1mm]
\makebox[\wd\hmboard]{\sffamily\hspace{3mm}\textbf{h\#7.5}\hfill(3+3)}\end{minipage}\hfill
\begin{minipage}[t]{52mm}\vspace{0pt}
\textcolor{ink}{\sffamily\textsc{Lösung}}\par\smallskip\small
\begin{tabular}{@{}r@{\enspace}l@{\enspace}l@{}}
\textbf{1.} & \ldots & \symking{}xa4\\
\textbf{2.} & \symrook{}e6 & \symking{}b5\\
\textbf{3.} & \symking{}g2 & a4\\
\textbf{4.} & \symking{}f3 & a5\\
\textbf{5.} & \symking{}e4 & a6\\
\textbf{6.} & \symking{}e5 & a7\\
\textbf{7.} & \symking{}d6 & a8=\symqueen{}\\
\textbf{8.} & \symrook{}e5 & \symqueen{}c6\#\\
\end{tabular}\par\vspace{3mm}
\qrcode[height=12mm]{https://helpman.komtera.lt/?fen=8/8/1r6/3N4/r7/K7/P7/7k&moves=7.5}\\{\sffamily\tiny\color{ink}scannen zum Lösen}
\end{minipage}
\vspace{2mm}

{\raggedright\noindent\textcolor{ink}{\sffamily\textsc{Themen}}\enspace \themetag{Umwandlung} \themetag{Excelsior} \themetag{Selbstblock} \themetag{ohne Schach}\par\smallskip{\footnotesize\textit{Hinweis:} die Partei am Zug steht im Diagramm im Schach; die Lösung beginnt mit einem Schlag.}\par}

\clearpage
\phantomsection\label{hm:KNvkrrb}
\begin{tcolorbox}[enhanced,colback=ink,colframe=ink,arc=0pt,boxrule=0pt,left=4mm,right=4mm,top=2.5mm,bottom=2.5mm]
\color{white}\sffamily{\LARGE\bfseries KNvkrrb}\hfill{\Large h\#7.5}\\[0.5mm]
{\small 6 Steine\quad Weiß am Zug}\hfill{\small Nr.~453\qquad\href{https://helpman.komtera.lt/?fen=8/1r5k/8/8/8/8/Nr2b3/K7&moves=7.5}{\color{white}\bfseries lösen\,\ensuremath{\nearrow}}}
\end{tcolorbox}
\vspace{2mm}
\noindent\begin{minipage}[t]{36mm}\vspace{0pt}\sffamily\small\raggedright
\hmfigure{längste eindeutige}{{\Large h\#7.5}}\hmfigure{eindeutig bei dieser Länge}{{\Large 215} Stellungen}\hmfigure{längstes Matt}{{\Large h\#8.5}}\hmfigure{in der Klasse}{15,5 Milliarden Stellungen,\\ 4,42 Milliarden mit Hilfsmatt}
\end{minipage}\hfill
\begin{minipage}[t]{80mm}\vspace{0pt}\begin{tikzpicture}\begin{axis}[width=72mm,height=42mm,ybar,bar width=1.6pt,xmin=-1,xmax=18,ymin=0,ymax=1.18,axis y line*=left,axis x line*=bottom,xlabel={\scriptsize Halbzüge bis Matt},ylabel={\scriptsize Stellungen (Mrd.)},tick label style={font=\tiny},label style={font=\scriptsize},enlarge x limits=false,ylabel shift=-3pt,xlabel shift=-3pt]\addplot[fill=ink!55,draw=none] coordinates {(0,3.624e-06) (1,1.623e-05) (2,0.0001738) (3,0.001178) (4,0.006122) (5,0.02982) (6,0.08909) (7,0.3041) (8,0.4669) (9,1.019) (10,0.7131) (11,1.073) (12,0.3639) (13,0.3085) (14,0.03587) (15,0.0118) (16,0.0003548) (17,6.714e-05)};\draw[densely dotted] (axis cs:15,0) -- (axis cs:15,1.18);\end{axis}\begin{axis}[width=72mm,height=42mm,axis y line*=right,axis x line=none,ymode=log,xmin=-1,xmax=18,ylabel={\scriptsize eindeutig (log)},tick label style={font=\tiny},label style={font=\scriptsize},ylabel shift=-3pt]\addplot[mark=*,mark size=0.7pt,draw=black,thin] coordinates {(0,3624) (1,16048) (2,145976) (3,791012) (4,2066638) (5,6111633) (6,6758386) (7,13726941) (8,8841029) (9,12556161) (10,3122662) (11,2943172) (12,210419) (13,80974) (14,882) (15,215)};\end{axis}\end{tikzpicture}\end{minipage}
\vspace{3mm}

\noindent\begin{minipage}[t]{60mm}\vspace{0pt}\sbox\hmboard{\chessboard[setfen=8/1r5k/8/8/8/8/Nr2b3/K7 w - - 0 1]}\usebox\hmboard\\[1mm]
\makebox[\wd\hmboard]{\sffamily\hspace{3mm}\textbf{h\#7.5}\hfill(2+4)}\end{minipage}\hfill
\begin{minipage}[t]{52mm}\vspace{0pt}
\textcolor{ink}{\sffamily\textsc{Lösung}}\par\smallskip\small
\begin{tabular}{@{}r@{\enspace}l@{\enspace}l@{}}
\textbf{1.} & \ldots & \symknight{}c3\\
\textbf{2.} & \symrook{}2b6 & \symknight{}e4\\
\textbf{3.} & \symbishop{}b5 & \symking{}b2\\
\textbf{4.} & \symking{}g6 & \symking{}c3\\
\textbf{5.} & \symrook{}h7 & \symking{}d4\\
\textbf{6.} & \symrook{}h4 & \symking{}e5\\
\textbf{7.} & \symking{}h5 & \symking{}f5\\
\textbf{8.} & \symrook{}h6 & \symknight{}g3\#\\
\end{tabular}\par\vspace{3mm}
\qrcode[height=12mm]{https://helpman.komtera.lt/?fen=8/1r5k/8/8/8/8/Nr2b3/K7&moves=7.5}\\{\sffamily\tiny\color{ink}scannen zum Lösen}
\end{minipage}
\vspace{2mm}

{\raggedright\noindent\textcolor{ink}{\sffamily\textsc{Themen}}\enspace \themetag{reines Matt} \themetag{Mustermatt} \themetag{Selbstblock} \themetag{ohne Schlag} \themetag{ohne Schach}\par}

\clearpage
\phantomsection\label{hm:KNvkrrr}
\begin{tcolorbox}[enhanced,colback=ink,colframe=ink,arc=0pt,boxrule=0pt,left=4mm,right=4mm,top=2.5mm,bottom=2.5mm]
\color{white}\sffamily{\LARGE\bfseries KNvkrrr}\hfill{\Large h\#7.5}\\[0.5mm]
{\small 6 Steine\quad Weiß am Zug}\hfill{\small Nr.~454\qquad\href{https://helpman.komtera.lt/?fen=8/rr6/1r6/6k1/N7/8/8/K7&moves=7.5}{\color{white}\bfseries lösen\,\ensuremath{\nearrow}}}
\end{tcolorbox}
\vspace{2mm}
\noindent\begin{minipage}[t]{36mm}\vspace{0pt}\sffamily\small\raggedright
\hmfigure{längste eindeutige}{{\Large h\#7.5}}\hmfigure{eindeutig bei dieser Länge}{{\Large 651} Stellungen}\hmfigure{längstes Matt}{{\Large h\#9}}\hmfigure{in der Klasse}{15,5 Milliarden Stellungen,\\ 1,39 Milliarden mit Hilfsmatt}
\end{minipage}\hfill
\begin{minipage}[t]{80mm}\vspace{0pt}\begin{tikzpicture}\begin{axis}[width=72mm,height=42mm,ybar,bar width=1.6pt,xmin=-1,xmax=19,ymin=0,ymax=382.9,axis y line*=left,axis x line*=bottom,xlabel={\scriptsize Halbzüge bis Matt},ylabel={\scriptsize Stellungen (Mio.)},tick label style={font=\tiny},label style={font=\scriptsize},enlarge x limits=false,ylabel shift=-3pt,xlabel shift=-3pt]\addplot[fill=ink!55,draw=none] coordinates {(0,0.00095) (1,0.003965) (2,0.04815) (3,0.3286) (4,1.59) (5,8.2) (6,23.37) (7,88.91) (8,128.1) (9,310.5) (10,199.8) (11,348.1) (12,116.6) (13,131.4) (14,20.78) (15,11.21) (16,0.5885) (17,0.1107) (18,2e-06)};\draw[densely dotted] (axis cs:15,0) -- (axis cs:15,382.9);\end{axis}\begin{axis}[width=72mm,height=42mm,axis y line*=right,axis x line=none,ymode=log,xmin=-1,xmax=19,ylabel={\scriptsize eindeutig (log)},tick label style={font=\tiny},label style={font=\scriptsize},ylabel shift=-3pt]\addplot[mark=*,mark size=0.7pt,draw=black,thin] coordinates {(0,950) (1,3965) (2,39367) (3,218821) (4,424331) (5,1667175) (6,2101373) (7,5012579) (8,2871316) (9,4532339) (10,1041795) (11,1273042) (12,110062) (13,82714) (14,2086) (15,651)};\end{axis}\end{tikzpicture}\end{minipage}
\vspace{3mm}

\noindent\begin{minipage}[t]{60mm}\vspace{0pt}\sbox\hmboard{\chessboard[setfen=8/rr6/1r6/6k1/N7/8/8/K7 w - - 0 1]}\usebox\hmboard\\[1mm]
\makebox[\wd\hmboard]{\sffamily\hspace{3mm}\textbf{h\#7.5}\hfill(2+4)}\end{minipage}\hfill
\begin{minipage}[t]{52mm}\vspace{0pt}
\textcolor{ink}{\sffamily\textsc{Lösung}}\par\smallskip\small
\begin{tabular}{@{}r@{\enspace}l@{\enspace}l@{}}
\textbf{1.} & \ldots & \symking{}a2\\
\textbf{2.} & \symrook{}b8 & \symking{}a3\\
\textbf{3.} & \symrook{}d7 & \symknight{}xb6\\
\textbf{4.} & \symking{}f6 & \symking{}b4\\
\textbf{5.} & \symking{}e7 & \symking{}c5\\
\textbf{6.} & \symking{}d8 & \symknight{}d5\\
\textbf{7.} & \symking{}c8 & \symking{}c6\\
\textbf{8.} & \symrook{}d8 & \symknight{}e7\#\\
\end{tabular}\par\vspace{3mm}
\qrcode[height=12mm]{https://helpman.komtera.lt/?fen=8/rr6/1r6/6k1/N7/8/8/K7&moves=7.5}\\{\sffamily\tiny\color{ink}scannen zum Lösen}
\end{minipage}
\vspace{2mm}

{\raggedright\noindent\textcolor{ink}{\sffamily\textsc{Themen}}\enspace \themetag{reines Matt} \themetag{Mustermatt} \themetag{Idealmatt} \themetag{Selbstblock} \themetag{ohne Schach}\par}

\clearpage
\phantomsection\label{hm:KPPvkbb}
\begin{tcolorbox}[enhanced,colback=ink,colframe=ink,arc=0pt,boxrule=0pt,left=4mm,right=4mm,top=2.5mm,bottom=2.5mm]
\color{white}\sffamily{\LARGE\bfseries KPPvkbb}\hfill{\Large h\#7.5}\\[0.5mm]
{\small 6 Steine\quad Weiß am Zug}\hfill{\small Nr.~455\qquad\href{https://helpman.komtera.lt/?fen=8/K7/8/b7/8/P7/P7/3b3k&moves=7.5}{\color{white}\bfseries lösen\,\ensuremath{\nearrow}}}
\end{tcolorbox}
\vspace{2mm}
\noindent\begin{minipage}[t]{36mm}\vspace{0pt}\sffamily\small\raggedright
\hmfigure{längste eindeutige}{{\Large h\#7.5}}\hmfigure{eindeutig bei dieser Länge}{{\Large 21} Stellungen}\hmfigure{längstes Matt}{{\Large h\#8.5}}\hmfigure{in der Klasse}{34,1 Milliarden Stellungen,\\ 5,77 Milliarden mit Hilfsmatt}
\end{minipage}\hfill
\begin{minipage}[t]{80mm}\vspace{0pt}\begin{tikzpicture}\begin{axis}[width=72mm,height=42mm,ybar,bar width=1.6pt,xmin=-1,xmax=18,ymin=0,ymax=1.153,axis y line*=left,axis x line*=bottom,xlabel={\scriptsize Halbzüge bis Matt},ylabel={\scriptsize Stellungen (Mrd.)},tick label style={font=\tiny},label style={font=\scriptsize},enlarge x limits=false,ylabel shift=-3pt,xlabel shift=-3pt]\addplot[fill=ink!55,draw=none] coordinates {(0,2.311e-05) (1,0.001133) (2,0.003914) (3,0.03861) (4,0.1086) (5,0.4372) (6,0.6388) (7,1.049) (8,0.8621) (9,0.969) (10,0.6142) (11,0.5824) (12,0.324) (13,0.1186) (14,0.01982) (15,0.001522) (16,1.77e-05) (17,3.95e-07)};\draw[densely dotted] (axis cs:15,0) -- (axis cs:15,1.153);\end{axis}\begin{axis}[width=72mm,height=42mm,axis y line*=right,axis x line=none,ymode=log,xmin=-1,xmax=18,ylabel={\scriptsize eindeutig (log)},tick label style={font=\tiny},label style={font=\scriptsize},ylabel shift=-3pt]\addplot[mark=*,mark size=0.7pt,draw=black,thin] coordinates {(0,23113) (1,346360) (2,1011575) (3,13024765) (4,22815481) (5,38320877) (6,21642299) (7,20491886) (8,8426210) (9,7161346) (10,2185003) (11,1553694) (12,236800) (13,28326) (14,487) (15,21)};\end{axis}\end{tikzpicture}\end{minipage}
\vspace{3mm}

\noindent\begin{minipage}[t]{60mm}\vspace{0pt}\sbox\hmboard{\chessboard[setfen=8/K7/8/b7/8/P7/P7/3b3k w - - 0 1]}\usebox\hmboard\\[1mm]
\makebox[\wd\hmboard]{\sffamily\hspace{3mm}\textbf{h\#7.5}\hfill(3+3)}\end{minipage}\hfill
\begin{minipage}[t]{52mm}\vspace{0pt}
\textcolor{ink}{\sffamily\textsc{Lösung}}\par\smallskip\small
\begin{tabular}{@{}r@{\enspace}l@{\enspace}l@{}}
\textbf{1.} & \ldots & \symking{}b8\\
\textbf{2.} & \symbishop{}b4 & axb4\\
\textbf{3.} & \symking{}g2 & a4\\
\textbf{4.} & \symking{}f3 & a5\\
\textbf{5.} & \symking{}e4 & a6\\
\textbf{6.} & \symking{}d5 & a7\\
\textbf{7.} & \symking{}c6 & a8=\symqueen{}+\\
\textbf{8.} & \symking{}b6 & \symqueen{}b7\#\\
\end{tabular}\par\vspace{3mm}
\qrcode[height=12mm]{https://helpman.komtera.lt/?fen=8/K7/8/b7/8/P7/P7/3b3k&moves=7.5}\\{\sffamily\tiny\color{ink}scannen zum Lösen}
\end{minipage}
\vspace{2mm}

{\raggedright\noindent\textcolor{ink}{\sffamily\textsc{Themen}}\enspace \themetag{Umwandlung} \themetag{Excelsior}\par}

\clearpage
\phantomsection\label{hm:KPPvkbn}
\begin{tcolorbox}[enhanced,colback=ink,colframe=ink,arc=0pt,boxrule=0pt,left=4mm,right=4mm,top=2.5mm,bottom=2.5mm]
\color{white}\sffamily{\LARGE\bfseries KPPvkbn}\hfill{\Large h\#7.5}\\[0.5mm]
{\small 6 Steine\quad Weiß am Zug}\hfill{\small Nr.~456\qquad\href{https://helpman.komtera.lt/?fen=8/7k/1b6/8/n7/P7/P7/K7&moves=7.5}{\color{white}\bfseries lösen\,\ensuremath{\nearrow}}}
\end{tcolorbox}
\vspace{2mm}
\noindent\begin{minipage}[t]{36mm}\vspace{0pt}\sffamily\small\raggedright
\hmfigure{längste eindeutige}{{\Large h\#7.5}}\hmfigure{eindeutig bei dieser Länge}{{\Large 73} Stellungen}\hmfigure{längstes Matt}{{\Large h\#8.5}}\hmfigure{in der Klasse}{34,1 Milliarden Stellungen,\\ 11,8 Milliarden mit Hilfsmatt}
\end{minipage}\hfill
\begin{minipage}[t]{80mm}\vspace{0pt}\begin{tikzpicture}\begin{axis}[width=72mm,height=42mm,ybar,bar width=1.6pt,xmin=-1,xmax=18,ymin=0,ymax=2.182,axis y line*=left,axis x line*=bottom,xlabel={\scriptsize Halbzüge bis Matt},ylabel={\scriptsize Stellungen (Mrd.)},tick label style={font=\tiny},label style={font=\scriptsize},enlarge x limits=false,ylabel shift=-3pt,xlabel shift=-3pt]\addplot[fill=ink!55,draw=none] coordinates {(0,4.316e-05) (1,0.002828) (2,0.008825) (3,0.0789) (4,0.2044) (5,0.8222) (6,1.213) (7,1.984) (8,1.766) (9,1.932) (10,1.358) (11,1.247) (12,0.801) (13,0.322) (14,0.0938) (15,0.006441) (16,0.0001319) (17,2.56e-06)};\draw[densely dotted] (axis cs:15,0) -- (axis cs:15,2.182);\end{axis}\begin{axis}[width=72mm,height=42mm,axis y line*=right,axis x line=none,ymode=log,xmin=-1,xmax=18,ylabel={\scriptsize eindeutig (log)},tick label style={font=\tiny},label style={font=\scriptsize},ylabel shift=-3pt]\addplot[mark=*,mark size=0.7pt,draw=black,thin] coordinates {(0,43155) (1,742992) (2,1957652) (3,23744640) (4,40871812) (5,70071526) (6,41904661) (7,37172090) (8,17796965) (9,14817776) (10,5621142) (11,4492902) (12,1092468) (13,113472) (14,2500) (15,73)};\end{axis}\end{tikzpicture}\end{minipage}
\vspace{3mm}

\noindent\begin{minipage}[t]{60mm}\vspace{0pt}\sbox\hmboard{\chessboard[setfen=8/7k/1b6/8/n7/P7/P7/K7 w - - 0 1]}\usebox\hmboard\\[1mm]
\makebox[\wd\hmboard]{\sffamily\hspace{3mm}\textbf{h\#7.5}\hfill(3+3)}\end{minipage}\hfill
\begin{minipage}[t]{52mm}\vspace{0pt}
\textcolor{ink}{\sffamily\textsc{Lösung}}\par\smallskip\small
\begin{tabular}{@{}r@{\enspace}l@{\enspace}l@{}}
\textbf{1.} & \ldots & \symking{}b1\\
\textbf{2.} & \symbishop{}d4 & \symking{}c1\\
\textbf{3.} & \symknight{}b6 & a4\\
\textbf{4.} & \symking{}g6 & a5\\
\textbf{5.} & \symking{}f5 & axb6\\
\textbf{6.} & \symking{}e4 & b7\\
\textbf{7.} & \symking{}d3 & b8=\symqueen{}\\
\textbf{8.} & \symking{}c3 & \symqueen{}b3\#\\
\end{tabular}\par\vspace{3mm}
\qrcode[height=12mm]{https://helpman.komtera.lt/?fen=8/7k/1b6/8/n7/P7/P7/K7&moves=7.5}\\{\sffamily\tiny\color{ink}scannen zum Lösen}
\end{minipage}
\vspace{2mm}

{\raggedright\noindent\textcolor{ink}{\sffamily\textsc{Themen}}\enspace \themetag{Umwandlung} \themetag{Selbstblock} \themetag{ohne Schach} \themetag{Umnov}\par}

\clearpage
\phantomsection\label{hm:KPPvknn}
\begin{tcolorbox}[enhanced,colback=ink,colframe=ink,arc=0pt,boxrule=0pt,left=4mm,right=4mm,top=2.5mm,bottom=2.5mm]
\color{white}\sffamily{\LARGE\bfseries KPPvknn}\hfill{\Large h\#7.5}\\[0.5mm]
{\small 6 Steine\quad Weiß am Zug}\hfill{\small Nr.~457\qquad\href{https://helpman.komtera.lt/?fen=7k/8/8/8/1n6/1n6/1PP5/K7&moves=7.5}{\color{white}\bfseries lösen\,\ensuremath{\nearrow}}}
\end{tcolorbox}
\vspace{2mm}
\noindent\begin{minipage}[t]{36mm}\vspace{0pt}\sffamily\small\raggedright
\hmfigure{längste eindeutige}{{\Large h\#7.5}}\hmfigure{eindeutig bei dieser Länge}{{\Large 71} Stellungen}\hmfigure{längstes Matt}{{\Large h\#9}}\hmfigure{in der Klasse}{34,1 Milliarden Stellungen,\\ 6,05 Milliarden mit Hilfsmatt}
\end{minipage}\hfill
\begin{minipage}[t]{80mm}\vspace{0pt}\begin{tikzpicture}\begin{axis}[width=72mm,height=42mm,ybar,bar width=1.6pt,xmin=-1,xmax=19,ymin=0,ymax=1067,axis y line*=left,axis x line*=bottom,xlabel={\scriptsize Halbzüge bis Matt},ylabel={\scriptsize Stellungen (Mio.)},tick label style={font=\tiny},label style={font=\scriptsize},enlarge x limits=false,ylabel shift=-3pt,xlabel shift=-3pt]\addplot[fill=ink!55,draw=none] coordinates {(0,0.01976) (1,1.776) (2,4.842) (3,39.77) (4,94.68) (5,380.2) (6,563) (7,942.3) (8,904.3) (9,970.4) (10,739.8) (11,650.4) (12,456.2) (13,204.5) (14,89.89) (15,8.233) (16,0.7523) (17,0.006248) (18,1e-06)};\draw[densely dotted] (axis cs:15,0) -- (axis cs:15,1067);\end{axis}\begin{axis}[width=72mm,height=42mm,axis y line*=right,axis x line=none,ymode=log,xmin=-1,xmax=19,ylabel={\scriptsize eindeutig (log)},tick label style={font=\tiny},label style={font=\scriptsize},ylabel shift=-3pt]\addplot[mark=*,mark size=0.7pt,draw=black,thin] coordinates {(0,19764) (1,389910) (2,894760) (3,10575438) (4,16420735) (5,30139194) (6,19458586) (7,16805138) (8,8974735) (9,7357634) (10,3228049) (11,2758399) (12,926105) (13,99534) (14,2322) (15,71)};\end{axis}\end{tikzpicture}\end{minipage}
\vspace{3mm}

\noindent\begin{minipage}[t]{60mm}\vspace{0pt}\sbox\hmboard{\chessboard[setfen=7k/8/8/8/1n6/1n6/1PP5/K7 w - - 0 1]}\usebox\hmboard\\[1mm]
\makebox[\wd\hmboard]{\sffamily\hspace{3mm}\textbf{h\#7.5}\hfill(3+3)}\end{minipage}\hfill
\begin{minipage}[t]{52mm}\vspace{0pt}
\textcolor{ink}{\sffamily\textsc{Lösung}}\par\smallskip\small
\begin{tabular}{@{}r@{\enspace}l@{\enspace}l@{}}
\textbf{1.} & \ldots & \symking{}b1\\
\textbf{2.} & \symknight{}c5 & \symking{}c1\\
\textbf{3.} & \symknight{}c6 & b4\\
\textbf{4.} & \symking{}g7 & b5\\
\textbf{5.} & \symking{}f6 & bxc6\\
\textbf{6.} & \symking{}e5 & c7\\
\textbf{7.} & \symking{}d4 & c8=\symqueen{}\\
\textbf{8.} & \symking{}c3 & \symqueen{}xc5\#\\
\end{tabular}\par\vspace{3mm}
\qrcode[height=12mm]{https://helpman.komtera.lt/?fen=7k/8/8/8/1n6/1n6/1PP5/K7&moves=7.5}\\{\sffamily\tiny\color{ink}scannen zum Lösen}
\end{minipage}
\vspace{2mm}

{\raggedright\noindent\textcolor{ink}{\sffamily\textsc{Themen}}\enspace \themetag{reines Matt} \themetag{Mustermatt} \themetag{Idealmatt} \themetag{Umwandlung} \themetag{Excelsior} \themetag{ohne Schach} \themetag{Umnov}\par\smallskip{\footnotesize\textit{Hinweis:} die Partei am Zug steht im Diagramm im Schach.}\par}

\clearpage
\phantomsection\label{hm:KPPvkqb}
\begin{tcolorbox}[enhanced,colback=ink,colframe=ink,arc=0pt,boxrule=0pt,left=4mm,right=4mm,top=2.5mm,bottom=2.5mm]
\color{white}\sffamily{\LARGE\bfseries KPPvkqb}\hfill{\Large h\#7.5}\\[0.5mm]
{\small 6 Steine\quad Weiß am Zug}\hfill{\small Nr.~458\qquad\href{https://helpman.komtera.lt/?fen=8/7k/8/5q2/1b6/1P6/1P6/K7&moves=7.5}{\color{white}\bfseries lösen\,\ensuremath{\nearrow}}}
\end{tcolorbox}
\vspace{2mm}
\noindent\begin{minipage}[t]{36mm}\vspace{0pt}\sffamily\small\raggedright
\hmfigure{längste eindeutige}{{\Large h\#7.5}}\hmfigure{eindeutig bei dieser Länge}{{\Large 24} Stellungen}\hmfigure{längstes Matt}{{\Large h\#11}}\hmfigure{in der Klasse}{34,1 Milliarden Stellungen,\\ 10,4 Milliarden mit Hilfsmatt}
\end{minipage}\hfill
\begin{minipage}[t]{80mm}\vspace{0pt}\begin{tikzpicture}\begin{axis}[width=72mm,height=42mm,ybar,bar width=1.6pt,xmin=-1,xmax=23,ymin=0,ymax=2.264,axis y line*=left,axis x line*=bottom,xlabel={\scriptsize Halbzüge bis Matt},ylabel={\scriptsize Stellungen (Mrd.)},tick label style={font=\tiny},label style={font=\scriptsize},enlarge x limits=false,ylabel shift=-3pt,xlabel shift=-3pt]\addplot[fill=ink!55,draw=none] coordinates {(0,2.408e-05) (1,0.00121) (2,0.005493) (3,0.05858) (4,0.1715) (5,0.7652) (6,1.009) (7,2.059) (8,1.326) (9,1.946) (10,0.9576) (11,1.211) (12,0.499) (13,0.3173) (14,0.02577) (15,0.006357) (16,2.017e-05) (17,2.791e-06) (18,2e-09) (20,2.4e-08) (22,1e-09)};\draw[densely dotted] (axis cs:15,0) -- (axis cs:15,2.264);\end{axis}\begin{axis}[width=72mm,height=42mm,axis y line*=right,axis x line=none,ymode=log,xmin=-1,xmax=23,ylabel={\scriptsize eindeutig (log)},tick label style={font=\tiny},label style={font=\scriptsize},ylabel shift=-3pt]\addplot[mark=*,mark size=0.7pt,draw=black,thin] coordinates {(0,24078) (1,402591) (2,1483437) (3,19757066) (4,34575151) (5,71873277) (6,30143464) (7,38665110) (8,11155390) (9,13228860) (10,3055372) (11,3017888) (12,288668) (13,89348) (14,462) (15,24)};\end{axis}\end{tikzpicture}\end{minipage}
\vspace{3mm}

\noindent\begin{minipage}[t]{60mm}\vspace{0pt}\sbox\hmboard{\chessboard[setfen=8/7k/8/5q2/1b6/1P6/1P6/K7 w - - 0 1]}\usebox\hmboard\\[1mm]
\makebox[\wd\hmboard]{\sffamily\hspace{3mm}\textbf{h\#7.5}\hfill(3+3)}\end{minipage}\hfill
\begin{minipage}[t]{52mm}\vspace{0pt}
\textcolor{ink}{\sffamily\textsc{Lösung}}\par\smallskip\small
\begin{tabular}{@{}r@{\enspace}l@{\enspace}l@{}}
\textbf{1.} & \ldots & \symking{}a2\\
\textbf{2.} & \symbishop{}d2 & b4\\
\textbf{3.} & \symqueen{}a5+ & bxa5\\
\textbf{4.} & \symking{}g6 & a6\\
\textbf{5.} & \symking{}f5 & a7\\
\textbf{6.} & \symking{}e4 & a8=\symqueen{}+\\
\textbf{7.} & \symking{}d3 & \symqueen{}h1\\
\textbf{8.} & \symking{}c2 & \symqueen{}b1\#\\
\end{tabular}\par\vspace{3mm}
\qrcode[height=12mm]{https://helpman.komtera.lt/?fen=8/7k/8/5q2/1b6/1P6/1P6/K7&moves=7.5}\\{\sffamily\tiny\color{ink}scannen zum Lösen}
\end{minipage}
\vspace{2mm}

{\raggedright\noindent\textcolor{ink}{\sffamily\textsc{Themen}}\enspace \themetag{reines Matt} \themetag{Mustermatt} \themetag{Idealmatt} \themetag{Umwandlung} \themetag{Selbstblock} \themetag{Umnov}\par}

\clearpage
\phantomsection\label{hm:KPPvkqn}
\begin{tcolorbox}[enhanced,colback=ink,colframe=ink,arc=0pt,boxrule=0pt,left=4mm,right=4mm,top=2.5mm,bottom=2.5mm]
\color{white}\sffamily{\LARGE\bfseries KPPvkqn}\hfill{\Large h\#7.5}\\[0.5mm]
{\small 6 Steine\quad Weiß am Zug}\hfill{\small Nr.~459\qquad\href{https://helpman.komtera.lt/?fen=8/7k/8/8/q1n5/P7/P7/K7&moves=7.5}{\color{white}\bfseries lösen\,\ensuremath{\nearrow}}}
\end{tcolorbox}
\vspace{2mm}
\noindent\begin{minipage}[t]{36mm}\vspace{0pt}\sffamily\small\raggedright
\hmfigure{längste eindeutige}{{\Large h\#7.5}}\hmfigure{eindeutig bei dieser Länge}{{\Large 42} Stellungen}\hmfigure{längstes Matt}{{\Large h\#9}}\hmfigure{in der Klasse}{34,1 Milliarden Stellungen,\\ 10,6 Milliarden mit Hilfsmatt}
\end{minipage}\hfill
\begin{minipage}[t]{80mm}\vspace{0pt}\begin{tikzpicture}\begin{axis}[width=72mm,height=42mm,ybar,bar width=1.6pt,xmin=-1,xmax=19,ymin=0,ymax=2.15,axis y line*=left,axis x line*=bottom,xlabel={\scriptsize Halbzüge bis Matt},ylabel={\scriptsize Stellungen (Mrd.)},tick label style={font=\tiny},label style={font=\scriptsize},enlarge x limits=false,ylabel shift=-3pt,xlabel shift=-3pt]\addplot[fill=ink!55,draw=none] coordinates {(0,2.196e-05) (1,0.00147) (2,0.006362) (3,0.06177) (4,0.1669) (5,0.7254) (6,0.9664) (7,1.955) (8,1.377) (9,1.953) (10,1.047) (11,1.277) (12,0.6118) (13,0.39) (14,0.05843) (15,0.01219) (16,7.969e-05) (17,2.081e-05) (18,1e-09)};\draw[densely dotted] (axis cs:15,0) -- (axis cs:15,2.15);\end{axis}\begin{axis}[width=72mm,height=42mm,axis y line*=right,axis x line=none,ymode=log,xmin=-1,xmax=19,ylabel={\scriptsize eindeutig (log)},tick label style={font=\tiny},label style={font=\scriptsize},ylabel shift=-3pt]\addplot[mark=*,mark size=0.7pt,draw=black,thin] coordinates {(0,21960) (1,430499) (2,1489584) (3,18518345) (4,30767065) (5,64573865) (6,28528684) (7,35449991) (8,11201068) (9,12923183) (10,3541393) (11,3808619) (12,569141) (13,173951) (14,1423) (15,42)};\end{axis}\end{tikzpicture}\end{minipage}
\vspace{3mm}

\noindent\begin{minipage}[t]{60mm}\vspace{0pt}\sbox\hmboard{\chessboard[setfen=8/7k/8/8/q1n5/P7/P7/K7 w - - 0 1]}\usebox\hmboard\\[1mm]
\makebox[\wd\hmboard]{\sffamily\hspace{3mm}\textbf{h\#7.5}\hfill(3+3)}\end{minipage}\hfill
\begin{minipage}[t]{52mm}\vspace{0pt}
\textcolor{ink}{\sffamily\textsc{Lösung}}\par\smallskip\small
\begin{tabular}{@{}r@{\enspace}l@{\enspace}l@{}}
\textbf{1.} & \ldots & \symking{}b1\\
\textbf{2.} & \symknight{}b6 & \symking{}c1\\
\textbf{3.} & \symqueen{}d4 & a4\\
\textbf{4.} & \symking{}g6 & a5\\
\textbf{5.} & \symking{}f5 & axb6\\
\textbf{6.} & \symking{}e4 & b7\\
\textbf{7.} & \symking{}d3 & b8=\symqueen{}\\
\textbf{8.} & \symking{}c3 & \symqueen{}b3\#\\
\end{tabular}\par\vspace{3mm}
\qrcode[height=12mm]{https://helpman.komtera.lt/?fen=8/7k/8/8/q1n5/P7/P7/K7&moves=7.5}\\{\sffamily\tiny\color{ink}scannen zum Lösen}
\end{minipage}
\vspace{2mm}

{\raggedright\noindent\textcolor{ink}{\sffamily\textsc{Themen}}\enspace \themetag{Umwandlung} \themetag{Selbstblock} \themetag{ohne Schach} \themetag{Umnov}\par}

\clearpage
\phantomsection\label{hm:KPPvkqq}
\begin{tcolorbox}[enhanced,colback=ink,colframe=ink,arc=0pt,boxrule=0pt,left=4mm,right=4mm,top=2.5mm,bottom=2.5mm]
\color{white}\sffamily{\LARGE\bfseries KPPvkqq}\hfill{\Large h\#7.5}\\[0.5mm]
{\small 6 Steine\quad Weiß am Zug}\hfill{\small Nr.~460\qquad\href{https://helpman.komtera.lt/?fen=8/8/7k/8/8/1P2q3/KP6/q7&moves=7.5}{\color{white}\bfseries lösen\,\ensuremath{\nearrow}}}
\end{tcolorbox}
\vspace{2mm}
\noindent\begin{minipage}[t]{36mm}\vspace{0pt}\sffamily\small\raggedright
\hmfigure{längste eindeutige}{{\Large h\#7.5}}\hmfigure{eindeutig bei dieser Länge}{{\Large 56} Stellungen}\hmfigure{längstes Matt}{{\Large h\#9}}\hmfigure{in der Klasse}{34,1 Milliarden Stellungen,\\ 4,6 Milliarden mit Hilfsmatt}
\end{minipage}\hfill
\begin{minipage}[t]{80mm}\vspace{0pt}\begin{tikzpicture}\begin{axis}[width=72mm,height=42mm,ybar,bar width=1.6pt,xmin=-1,xmax=19,ymin=0,ymax=1080,axis y line*=left,axis x line*=bottom,xlabel={\scriptsize Halbzüge bis Matt},ylabel={\scriptsize Stellungen (Mio.)},tick label style={font=\tiny},label style={font=\scriptsize},enlarge x limits=false,ylabel shift=-3pt,xlabel shift=-3pt]\addplot[fill=ink!55,draw=none] coordinates {(0,0.005742) (1,0.3059) (2,1.726) (3,20.67) (4,64.71) (5,323.2) (6,390.8) (7,982.2) (8,501.5) (9,953.6) (10,354.4) (11,605.3) (12,193.3) (13,197.7) (14,8.561) (15,4.943) (16,0.006272) (17,0.002495) (18,1e-06)};\draw[densely dotted] (axis cs:15,0) -- (axis cs:15,1080);\end{axis}\begin{axis}[width=72mm,height=42mm,axis y line*=right,axis x line=none,ymode=log,xmin=-1,xmax=19,ylabel={\scriptsize eindeutig (log)},tick label style={font=\tiny},label style={font=\scriptsize},ylabel shift=-3pt]\addplot[mark=*,mark size=0.7pt,draw=black,thin] coordinates {(0,5742) (1,110129) (2,495194) (3,7172628) (4,13240826) (5,33410865) (6,11201748) (7,18772932) (8,3732677) (9,6018801) (10,979813) (11,1432528) (12,90316) (13,59766) (14,125) (15,56)};\end{axis}\end{tikzpicture}\end{minipage}
\vspace{3mm}

\noindent\begin{minipage}[t]{60mm}\vspace{0pt}\sbox\hmboard{\chessboard[setfen=8/8/7k/8/8/1P2q3/KP6/q7 w - - 0 1]}\usebox\hmboard\\[1mm]
\makebox[\wd\hmboard]{\sffamily\hspace{3mm}\textbf{h\#7.5}\hfill(3+3)}\end{minipage}\hfill
\begin{minipage}[t]{52mm}\vspace{0pt}
\textcolor{ink}{\sffamily\textsc{Lösung}}\par\smallskip\small
\begin{tabular}{@{}r@{\enspace}l@{\enspace}l@{}}
\textbf{1.} & \ldots & \symking{}xa1\\
\textbf{2.} & \symking{}g5 & b4\\
\textbf{3.} & \symking{}f4 & b3\\
\textbf{4.} & \symqueen{}xb3 & b5\\
\textbf{5.} & \symking{}e3 & b6\\
\textbf{6.} & \symking{}d2 & b7\\
\textbf{7.} & \symking{}c1 & b8=\symqueen{}\\
\textbf{8.} & \symqueen{}d1 & \symqueen{}b2\#\\
\end{tabular}\par\vspace{3mm}
\qrcode[height=12mm]{https://helpman.komtera.lt/?fen=8/8/7k/8/8/1P2q3/KP6/q7&moves=7.5}\\{\sffamily\tiny\color{ink}scannen zum Lösen}
\end{minipage}
\vspace{2mm}

{\raggedright\noindent\textcolor{ink}{\sffamily\textsc{Themen}}\enspace \themetag{Umwandlung} \themetag{Selbstblock} \themetag{ohne Schach}\par\smallskip{\footnotesize\textit{Hinweis:} die Partei am Zug steht im Diagramm im Schach; die Lösung beginnt mit einem Schlag.}\par}

\clearpage
\phantomsection\label{hm:KPPvkqr}
\begin{tcolorbox}[enhanced,colback=ink,colframe=ink,arc=0pt,boxrule=0pt,left=4mm,right=4mm,top=2.5mm,bottom=2.5mm]
\color{white}\sffamily{\LARGE\bfseries KPPvkqr}\hfill{\Large h\#7.5}\\[0.5mm]
{\small 6 Steine\quad Weiß am Zug}\hfill{\small Nr.~461\qquad\href{https://helpman.komtera.lt/?fen=8/8/7k/8/8/1P2q3/KP6/r7&moves=7.5}{\color{white}\bfseries lösen\,\ensuremath{\nearrow}}}
\end{tcolorbox}
\vspace{2mm}
\noindent\begin{minipage}[t]{36mm}\vspace{0pt}\sffamily\small\raggedright
\hmfigure{längste eindeutige}{{\Large h\#7.5}}\hmfigure{eindeutig bei dieser Länge}{{\Large 7} Stellungen}\hmfigure{längstes Matt}{{\Large h\#8.5}}\hmfigure{in der Klasse}{34,1 Milliarden Stellungen,\\ 9,93 Milliarden mit Hilfsmatt}
\end{minipage}\hfill
\begin{minipage}[t]{80mm}\vspace{0pt}\begin{tikzpicture}\begin{axis}[width=72mm,height=42mm,ybar,bar width=1.6pt,xmin=-1,xmax=18,ymin=0,ymax=2.659,axis y line*=left,axis x line*=bottom,xlabel={\scriptsize Halbzüge bis Matt},ylabel={\scriptsize Stellungen (Mrd.)},tick label style={font=\tiny},label style={font=\scriptsize},enlarge x limits=false,ylabel shift=-3pt,xlabel shift=-3pt]\addplot[fill=ink!55,draw=none] coordinates {(0,1.562e-05) (1,0.001062) (2,0.005805) (3,0.05849) (4,0.1965) (5,0.8581) (6,1.195) (7,2.417) (8,1.257) (9,1.87) (10,0.7073) (11,0.9123) (12,0.2672) (13,0.1782) (14,0.0009578) (15,0.001761) (16,4.03e-07) (17,6.81e-07)};\draw[densely dotted] (axis cs:15,0) -- (axis cs:15,2.659);\end{axis}\begin{axis}[width=72mm,height=42mm,axis y line*=right,axis x line=none,ymode=log,xmin=-1,xmax=18,ylabel={\scriptsize eindeutig (log)},tick label style={font=\tiny},label style={font=\scriptsize},ylabel shift=-3pt]\addplot[mark=*,mark size=0.7pt,draw=black,thin] coordinates {(0,15618) (1,389678) (2,1966874) (3,20794026) (4,44638493) (5,87212328) (6,36882504) (7,44029855) (8,7823474) (9,9003586) (10,1586029) (11,978136) (12,6984) (13,16355) (15,7)};\end{axis}\end{tikzpicture}\end{minipage}
\vspace{3mm}

\noindent\begin{minipage}[t]{60mm}\vspace{0pt}\sbox\hmboard{\chessboard[setfen=8/8/7k/8/8/1P2q3/KP6/r7 w - - 0 1]}\usebox\hmboard\\[1mm]
\makebox[\wd\hmboard]{\sffamily\hspace{3mm}\textbf{h\#7.5}\hfill(3+3)}\end{minipage}\hfill
\begin{minipage}[t]{52mm}\vspace{0pt}
\textcolor{ink}{\sffamily\textsc{Lösung}}\par\smallskip\small
\begin{tabular}{@{}r@{\enspace}l@{\enspace}l@{}}
\textbf{1.} & \ldots & \symking{}xa1\\
\textbf{2.} & \symking{}g5 & b4\\
\textbf{3.} & \symking{}f4 & b3\\
\textbf{4.} & \symqueen{}xb3 & b5\\
\textbf{5.} & \symking{}e3 & b6\\
\textbf{6.} & \symking{}d2 & b7\\
\textbf{7.} & \symking{}c1 & b8=\symqueen{}\\
\textbf{8.} & \symqueen{}d1 & \symqueen{}b2\#\\
\end{tabular}\par\vspace{3mm}
\qrcode[height=12mm]{https://helpman.komtera.lt/?fen=8/8/7k/8/8/1P2q3/KP6/r7&moves=7.5}\\{\sffamily\tiny\color{ink}scannen zum Lösen}
\end{minipage}
\vspace{2mm}

{\raggedright\noindent\textcolor{ink}{\sffamily\textsc{Themen}}\enspace \themetag{Umwandlung} \themetag{Selbstblock} \themetag{ohne Schach}\par\smallskip{\footnotesize\textit{Hinweis:} die Partei am Zug steht im Diagramm im Schach; die Lösung beginnt mit einem Schlag.}\par}

\clearpage
\phantomsection\label{hm:KPPvkrb}
\begin{tcolorbox}[enhanced,colback=ink,colframe=ink,arc=0pt,boxrule=0pt,left=4mm,right=4mm,top=2.5mm,bottom=2.5mm]
\color{white}\sffamily{\LARGE\bfseries KPPvkrb}\hfill{\Large h\#7.5}\\[0.5mm]
{\small 6 Steine\quad Weiß am Zug}\hfill{\small Nr.~462\qquad\href{https://helpman.komtera.lt/?fen=8/8/7k/8/8/1P6/KP1b4/r7&moves=7.5}{\color{white}\bfseries lösen\,\ensuremath{\nearrow}}}
\end{tcolorbox}
\vspace{2mm}
\noindent\begin{minipage}[t]{36mm}\vspace{0pt}\sffamily\small\raggedright
\hmfigure{längste eindeutige}{{\Large h\#7.5}}\hmfigure{eindeutig bei dieser Länge}{{\Large 1} Stellung}\hmfigure{längstes Matt}{{\Large h\#8.5}}\hmfigure{in der Klasse}{34,1 Milliarden Stellungen,\\ 11,2 Milliarden mit Hilfsmatt}
\end{minipage}\hfill
\begin{minipage}[t]{80mm}\vspace{0pt}\begin{tikzpicture}\begin{axis}[width=72mm,height=42mm,ybar,bar width=1.6pt,xmin=-1,xmax=18,ymin=0,ymax=2.663,axis y line*=left,axis x line*=bottom,xlabel={\scriptsize Halbzüge bis Matt},ylabel={\scriptsize Stellungen (Mrd.)},tick label style={font=\tiny},label style={font=\scriptsize},enlarge x limits=false,ylabel shift=-3pt,xlabel shift=-3pt]\addplot[fill=ink!55,draw=none] coordinates {(0,3.243e-05) (1,0.002102) (2,0.00905) (3,0.0807) (4,0.2591) (5,1.004) (6,1.514) (7,2.421) (8,1.635) (9,1.852) (10,0.9801) (11,0.9017) (12,0.3723) (13,0.1258) (14,0.00391) (15,0.0003894) (16,5.26e-07) (17,3.81e-07)};\draw[densely dotted] (axis cs:15,0) -- (axis cs:15,2.663);\end{axis}\begin{axis}[width=72mm,height=42mm,axis y line*=right,axis x line=none,ymode=log,xmin=-1,xmax=18,ylabel={\scriptsize eindeutig (log)},tick label style={font=\tiny},label style={font=\scriptsize},ylabel shift=-3pt]\addplot[mark=*,mark size=0.7pt,draw=black,thin] coordinates {(0,32435) (1,710391) (2,3001171) (3,28405545) (4,60004764) (5,95278503) (6,49628141) (7,46093521) (8,15495744) (9,13273699) (10,2949149) (11,1118373) (12,51203) (13,5191) (14,62) (15,1)};\end{axis}\end{tikzpicture}\end{minipage}
\vspace{3mm}

\noindent\begin{minipage}[t]{60mm}\vspace{0pt}\sbox\hmboard{\chessboard[setfen=8/8/7k/8/8/1P6/KP1b4/r7 w - - 0 1]}\usebox\hmboard\\[1mm]
\makebox[\wd\hmboard]{\sffamily\hspace{3mm}\textbf{h\#7.5}\hfill(3+3)}\end{minipage}\hfill
\begin{minipage}[t]{52mm}\vspace{0pt}
\textcolor{ink}{\sffamily\textsc{Lösung}}\par\smallskip\small
\begin{tabular}{@{}r@{\enspace}l@{\enspace}l@{}}
\textbf{1.} & \ldots & \symking{}xa1\\
\textbf{2.} & \symking{}g5 & b4\\
\textbf{3.} & \symking{}f4 & b5\\
\textbf{4.} & \symking{}e3 & b4\\
\textbf{5.} & \symbishop{}xb4 & b6\\
\textbf{6.} & \symking{}d2 & b7\\
\textbf{7.} & \symking{}c1 & b8=\symqueen{}\\
\textbf{8.} & \symbishop{}d2 & \symqueen{}b1\#\\
\end{tabular}\par\vspace{3mm}
\qrcode[height=12mm]{https://helpman.komtera.lt/?fen=8/8/7k/8/8/1P6/KP1b4/r7&moves=7.5}\\{\sffamily\tiny\color{ink}scannen zum Lösen}
\end{minipage}
\vspace{2mm}

{\raggedright\noindent\textcolor{ink}{\sffamily\textsc{Themen}}\enspace \themetag{Umwandlung} \themetag{Switchback} \themetag{Selbstblock} \themetag{ohne Schach}\par\smallskip{\footnotesize\textit{Hinweis:} die Partei am Zug steht im Diagramm im Schach; die Lösung beginnt mit einem Schlag.}\par}

\clearpage
\phantomsection\label{hm:KPPvkrn}
\begin{tcolorbox}[enhanced,colback=ink,colframe=ink,arc=0pt,boxrule=0pt,left=4mm,right=4mm,top=2.5mm,bottom=2.5mm]
\color{white}\sffamily{\LARGE\bfseries KPPvkrn}\hfill{\Large h\#7.5}\\[0.5mm]
{\small 6 Steine\quad Weiß am Zug}\hfill{\small Nr.~463\qquad\href{https://helpman.komtera.lt/?fen=1r6/3n4/8/8/K7/P6k/P7/8&moves=7.5}{\color{white}\bfseries lösen\,\ensuremath{\nearrow}}}
\end{tcolorbox}
\vspace{2mm}
\noindent\begin{minipage}[t]{36mm}\vspace{0pt}\sffamily\small\raggedright
\hmfigure{längste eindeutige}{{\Large h\#7.5}}\hmfigure{eindeutig bei dieser Länge}{{\Large 1} Stellung}\hmfigure{längstes Matt}{{\Large h\#8.5}}\hmfigure{in der Klasse}{34,1 Milliarden Stellungen,\\ 11,4 Milliarden mit Hilfsmatt}
\end{minipage}\hfill
\begin{minipage}[t]{80mm}\vspace{0pt}\begin{tikzpicture}\begin{axis}[width=72mm,height=42mm,ybar,bar width=1.6pt,xmin=-1,xmax=18,ymin=0,ymax=2.488,axis y line*=left,axis x line*=bottom,xlabel={\scriptsize Halbzüge bis Matt},ylabel={\scriptsize Stellungen (Mrd.)},tick label style={font=\tiny},label style={font=\scriptsize},enlarge x limits=false,ylabel shift=-3pt,xlabel shift=-3pt]\addplot[fill=ink!55,draw=none] coordinates {(0,2.919e-05) (1,0.002517) (2,0.01002) (3,0.08262) (4,0.2468) (5,0.9366) (6,1.41) (7,2.262) (8,1.696) (9,1.889) (10,1.099) (11,1.027) (12,0.5285) (13,0.188) (14,0.01506) (15,0.0009021) (16,1.487e-06) (17,3.818e-06)};\draw[densely dotted] (axis cs:15,0) -- (axis cs:15,2.488);\end{axis}\begin{axis}[width=72mm,height=42mm,axis y line*=right,axis x line=none,ymode=log,xmin=-1,xmax=18,ylabel={\scriptsize eindeutig (log)},tick label style={font=\tiny},label style={font=\scriptsize},ylabel shift=-3pt]\addplot[mark=*,mark size=0.7pt,draw=black,thin] coordinates {(0,29186) (1,742503) (2,2906763) (3,26195590) (4,54092902) (5,86814848) (6,46109791) (7,41842024) (8,15625891) (9,13917184) (10,4308565) (11,2285014) (12,213115) (13,13652) (14,241) (15,1)};\end{axis}\end{tikzpicture}\end{minipage}
\vspace{3mm}

\noindent\begin{minipage}[t]{60mm}\vspace{0pt}\sbox\hmboard{\chessboard[setfen=1r6/3n4/8/8/K7/P6k/P7/8 w - - 0 1]}\usebox\hmboard\\[1mm]
\makebox[\wd\hmboard]{\sffamily\hspace{3mm}\textbf{h\#7.5}\hfill(3+3)}\end{minipage}\hfill
\begin{minipage}[t]{52mm}\vspace{0pt}
\textcolor{ink}{\sffamily\textsc{Lösung}}\par\smallskip\small
\begin{tabular}{@{}r@{\enspace}l@{\enspace}l@{}}
\textbf{1.} & \ldots & \symking{}a5\\
\textbf{2.} & \symrook{}d8 & a4\\
\textbf{3.} & \symknight{}b8 & \symking{}b6\\
\textbf{4.} & \symking{}g4 & a5\\
\textbf{5.} & \symking{}f5 & a6\\
\textbf{6.} & \symking{}e6 & a7\\
\textbf{7.} & \symking{}d7 & a8=\symqueen{}\\
\textbf{8.} & \symking{}c8 & \symqueen{}b7\#\\
\end{tabular}\par\vspace{3mm}
\qrcode[height=12mm]{https://helpman.komtera.lt/?fen=1r6/3n4/8/8/K7/P6k/P7/8&moves=7.5}\\{\sffamily\tiny\color{ink}scannen zum Lösen}
\end{minipage}
\vspace{2mm}

{\raggedright\noindent\textcolor{ink}{\sffamily\textsc{Themen}}\enspace \themetag{Satzspiel} \themetag{Umwandlung} \themetag{Selbstblock} \themetag{ohne Schlag} \themetag{ohne Schach}\par}

\clearpage
\phantomsection\label{hm:KPvkbbb}
\begin{tcolorbox}[enhanced,colback=ink,colframe=ink,arc=0pt,boxrule=0pt,left=4mm,right=4mm,top=2.5mm,bottom=2.5mm]
\color{white}\sffamily{\LARGE\bfseries KPvkbbb}\hfill{\Large h\#7.5}\\[0.5mm]
{\small 6 Steine\quad Weiß am Zug}\hfill{\small Nr.~464\qquad\href{https://helpman.komtera.lt/?fen=8/8/7k/8/8/3b4/3P4/K1b3b1&moves=7.5}{\color{white}\bfseries lösen\,\ensuremath{\nearrow}}}
\end{tcolorbox}
\vspace{2mm}
\noindent\begin{minipage}[t]{36mm}\vspace{0pt}\sffamily\small\raggedright
\hmfigure{längste eindeutige}{{\Large h\#7.5}}\hmfigure{eindeutig bei dieser Länge}{{\Large 907} Stellungen}\hmfigure{längstes Matt}{{\Large h\#9.5}}\hmfigure{in der Klasse}{45,4 Milliarden Stellungen,\\ 4,9 Milliarden mit Hilfsmatt}
\end{minipage}\hfill
\begin{minipage}[t]{80mm}\vspace{0pt}\begin{tikzpicture}\begin{axis}[width=72mm,height=42mm,ybar,bar width=1.6pt,xmin=-1,xmax=20,ymin=0,ymax=1006,axis y line*=left,axis x line*=bottom,xlabel={\scriptsize Halbzüge bis Matt},ylabel={\scriptsize Stellungen (Mio.)},tick label style={font=\tiny},label style={font=\scriptsize},enlarge x limits=false,ylabel shift=-3pt,xlabel shift=-3pt]\addplot[fill=ink!55,draw=none] coordinates {(0,0.004645) (1,0.2684) (2,0.9724) (3,9.896) (4,26.95) (5,137) (6,210.7) (7,519) (8,454.9) (9,791) (10,551.7) (11,914.4) (12,603.6) (13,532.7) (14,119.4) (15,29.98) (16,1.295) (17,0.2006) (18,0.000268) (19,8e-06)};\draw[densely dotted] (axis cs:15,0) -- (axis cs:15,1006);\end{axis}\begin{axis}[width=72mm,height=42mm,axis y line*=right,axis x line=none,ymode=log,xmin=-1,xmax=20,ylabel={\scriptsize eindeutig (log)},tick label style={font=\tiny},label style={font=\scriptsize},ylabel shift=-3pt]\addplot[mark=*,mark size=0.7pt,draw=black,thin] coordinates {(0,4645) (1,79403) (2,188342) (3,2921480) (4,4833138) (5,13800987) (6,6609934) (7,10752950) (8,4124713) (9,6243686) (10,2231582) (11,3349208) (12,747368) (13,365854) (14,10217) (15,907)};\end{axis}\end{tikzpicture}\end{minipage}
\vspace{3mm}

\noindent\begin{minipage}[t]{60mm}\vspace{0pt}\sbox\hmboard{\chessboard[setfen=8/8/7k/8/8/3b4/3P4/K1b3b1 w - - 0 1]}\usebox\hmboard\\[1mm]
\makebox[\wd\hmboard]{\sffamily\hspace{3mm}\textbf{h\#7.5}\hfill(2+4)}\end{minipage}\hfill
\begin{minipage}[t]{52mm}\vspace{0pt}
\textcolor{ink}{\sffamily\textsc{Lösung}}\par\smallskip\small
\begin{tabular}{@{}r@{\enspace}l@{\enspace}l@{}}
\textbf{1.} & \ldots & \symking{}a2\\
\textbf{2.} & \symbishop{}e2 & \symking{}b1\\
\textbf{3.} & \symbishop{}c5 & d4\\
\textbf{4.} & \symking{}g5 & dxc5\\
\textbf{5.} & \symking{}f4 & c6\\
\textbf{6.} & \symking{}e3 & c7\\
\textbf{7.} & \symking{}d2 & c8=\symqueen{}\\
\textbf{8.} & \symking{}d1 & \symqueen{}xc1\#\\
\end{tabular}\par\vspace{3mm}
\qrcode[height=12mm]{https://helpman.komtera.lt/?fen=8/8/7k/8/8/3b4/3P4/K1b3b1&moves=7.5}\\{\sffamily\tiny\color{ink}scannen zum Lösen}
\end{minipage}
\vspace{2mm}

{\raggedright\noindent\textcolor{ink}{\sffamily\textsc{Themen}}\enspace \themetag{Satzspiel} \themetag{Umwandlung} \themetag{Excelsior} \themetag{Selbstblock} \themetag{ohne Schach}\par}

\clearpage
\phantomsection\label{hm:KPvkqqr}
\begin{tcolorbox}[enhanced,colback=ink,colframe=ink,arc=0pt,boxrule=0pt,left=4mm,right=4mm,top=2.5mm,bottom=2.5mm]
\color{white}\sffamily{\LARGE\bfseries KPvkqqr}\hfill{\Large h\#7.5}\\[0.5mm]
{\small 6 Steine\quad Weiß am Zug}\hfill{\small Nr.~465\qquad\href{https://helpman.komtera.lt/?fen=3q4/q7/8/1K6/2rk4/8/2P5/8&moves=7.5}{\color{white}\bfseries lösen\,\ensuremath{\nearrow}}}
\end{tcolorbox}
\vspace{2mm}
\noindent\begin{minipage}[t]{36mm}\vspace{0pt}\sffamily\small\raggedright
\hmfigure{längste eindeutige}{{\Large h\#7.5}}\hmfigure{eindeutig bei dieser Länge}{{\Large 768} Stellungen}\hmfigure{längstes Matt}{{\Large h\#9.5}}\hmfigure{in der Klasse}{45,4 Milliarden Stellungen,\\ 11,2 Milliarden mit Hilfsmatt}
\end{minipage}\hfill
\begin{minipage}[t]{80mm}\vspace{0pt}\begin{tikzpicture}\begin{axis}[width=72mm,height=42mm,ybar,bar width=1.6pt,xmin=-1,xmax=20,ymin=0,ymax=2.649,axis y line*=left,axis x line*=bottom,xlabel={\scriptsize Halbzüge bis Matt},ylabel={\scriptsize Stellungen (Mrd.)},tick label style={font=\tiny},label style={font=\scriptsize},enlarge x limits=false,ylabel shift=-3pt,xlabel shift=-3pt]\addplot[fill=ink!55,draw=none] coordinates {(1,0.0001749) (2,0.001182) (3,0.01397) (4,0.0517) (5,0.3224) (6,0.5024) (7,1.729) (8,0.8977) (9,2.408) (10,0.8818) (11,2.356) (12,0.7599) (13,1.246) (14,0.01711) (15,0.02095) (16,2.722e-05) (17,7.717e-05) (18,1.3e-08) (19,2e-09)};\draw[densely dotted] (axis cs:15,0) -- (axis cs:15,2.649);\end{axis}\begin{axis}[width=72mm,height=42mm,axis y line*=right,axis x line=none,ymode=log,xmin=-1,xmax=20,ylabel={\scriptsize eindeutig (log)},tick label style={font=\tiny},label style={font=\scriptsize},ylabel shift=-3pt]\addplot[mark=*,mark size=0.7pt,draw=black,thin] coordinates {(1,66517) (2,373668) (3,4949104) (4,12242216) (5,43378170) (6,19553179) (7,44380747) (8,7840488) (9,17809022) (10,3396949) (11,5652386) (12,93927) (13,109349) (14,309) (15,768)};\end{axis}\end{tikzpicture}\end{minipage}
\vspace{3mm}

\noindent\begin{minipage}[t]{60mm}\vspace{0pt}\sbox\hmboard{\chessboard[setfen=3q4/q7/8/1K6/2rk4/8/2P5/8 w - - 0 1]}\usebox\hmboard\\[1mm]
\makebox[\wd\hmboard]{\sffamily\hspace{3mm}\textbf{h\#7.5}\hfill(2+4)}\end{minipage}\hfill
\begin{minipage}[t]{52mm}\vspace{0pt}
\textcolor{ink}{\sffamily\textsc{Lösung}}\par\smallskip\small
\begin{tabular}{@{}r@{\enspace}l@{\enspace}l@{}}
\textbf{1.} & \ldots & c3+\\
\textbf{2.} & \symking{}e5 & \symking{}xc4\\
\textbf{3.} & \symqueen{}ad4+ & cxd4+\\
\textbf{4.} & \symking{}d6 & d5\\
\textbf{5.} & \symking{}c7 & d6+\\
\textbf{6.} & \symking{}b6 & d7\\
\textbf{7.} & \symqueen{}e8 & dxe8=\symqueen{}\\
\textbf{8.} & \symking{}a5 & \symqueen{}b5\#\\
\end{tabular}\par\vspace{3mm}
\qrcode[height=12mm]{https://helpman.komtera.lt/?fen=3q4/q7/8/1K6/2rk4/8/2P5/8&moves=7.5}\\{\sffamily\tiny\color{ink}scannen zum Lösen}
\end{minipage}
\vspace{2mm}

{\raggedright\noindent\textcolor{ink}{\sffamily\textsc{Themen}}\enspace \themetag{Umwandlung} \themetag{Excelsior} \themetag{Umnov}\par}

\clearpage
\phantomsection\label{hm:KPvkqrb}
\begin{tcolorbox}[enhanced,colback=ink,colframe=ink,arc=0pt,boxrule=0pt,left=4mm,right=4mm,top=2.5mm,bottom=2.5mm]
\color{white}\sffamily{\LARGE\bfseries KPvkqrb}\hfill{\Large h\#7.5}\\[0.5mm]
{\small 6 Steine\quad Weiß am Zug}\hfill{\small Nr.~466\qquad\href{https://helpman.komtera.lt/?fen=7b/8/8/8/4rk2/8/3qP3/1K6&moves=7.5}{\color{white}\bfseries lösen\,\ensuremath{\nearrow}}}
\end{tcolorbox}
\vspace{2mm}
\noindent\begin{minipage}[t]{36mm}\vspace{0pt}\sffamily\small\raggedright
\hmfigure{längste eindeutige}{{\Large h\#7.5}}\hmfigure{eindeutig bei dieser Länge}{{\Large 854} Stellungen}\hmfigure{längstes Matt}{{\Large h\#11.5}}\hmfigure{in der Klasse}{45,4 Milliarden Stellungen,\\ 25,6 Milliarden mit Hilfsmatt}
\end{minipage}\hfill
\begin{minipage}[t]{80mm}\vspace{0pt}\begin{tikzpicture}\begin{axis}[width=72mm,height=42mm,ybar,bar width=1.6pt,xmin=-1,xmax=24,ymin=0,ymax=5.678,axis y line*=left,axis x line*=bottom,xlabel={\scriptsize Halbzüge bis Matt},ylabel={\scriptsize Stellungen (Mrd.)},tick label style={font=\tiny},label style={font=\scriptsize},enlarge x limits=false,ylabel shift=-3pt,xlabel shift=-3pt]\addplot[fill=ink!55,draw=none] coordinates {(0,4.84e-06) (1,0.0007582) (2,0.004234) (3,0.04387) (4,0.1522) (5,0.8414) (6,1.396) (7,3.877) (8,2.428) (9,5.161) (10,2.442) (11,5.072) (12,1.968) (13,2.174) (14,0.04187) (15,0.02824) (16,2.919e-05) (17,0.0001399) (18,8.2e-08) (19,1.9e-07) (20,5.8e-08) (21,3.1e-08) (23,2e-09)};\draw[densely dotted] (axis cs:15,0) -- (axis cs:15,5.678);\end{axis}\begin{axis}[width=72mm,height=42mm,axis y line*=right,axis x line=none,ymode=log,xmin=-1,xmax=24,ylabel={\scriptsize eindeutig (log)},tick label style={font=\tiny},label style={font=\scriptsize},ylabel shift=-3pt]\addplot[mark=*,mark size=0.7pt,draw=black,thin] coordinates {(0,4840) (1,273680) (2,1332266) (3,15046553) (4,35801059) (5,102007647) (6,52142616) (7,91091434) (8,22507958) (9,41388605) (10,9990894) (11,11395745) (12,282394) (13,177417) (14,402) (15,854)};\end{axis}\end{tikzpicture}\end{minipage}
\vspace{3mm}

\noindent\begin{minipage}[t]{60mm}\vspace{0pt}\sbox\hmboard{\chessboard[setfen=7b/8/8/8/4rk2/8/3qP3/1K6 w - - 0 1]}\usebox\hmboard\\[1mm]
\makebox[\wd\hmboard]{\sffamily\hspace{3mm}\textbf{h\#7.5}\hfill(2+4)}\end{minipage}\hfill
\begin{minipage}[t]{52mm}\vspace{0pt}
\textcolor{ink}{\sffamily\textsc{Lösung}}\par\smallskip\small
\begin{tabular}{@{}r@{\enspace}l@{\enspace}l@{}}
\textbf{1.} & \ldots & e3+\\
\textbf{2.} & \symking{}e5 & \symking{}a1\\
\textbf{3.} & \symrook{}g4 & e4\\
\textbf{4.} & \symking{}f6 & e5+\\
\textbf{5.} & \symking{}g7 & e6\\
\textbf{6.} & \symrook{}g6 & e7\\
\textbf{7.} & \symqueen{}g5 & e8=\symqueen{}\\
\textbf{8.} & \symking{}h6+ & \symqueen{}xh8\#\\
\end{tabular}\par\vspace{3mm}
\qrcode[height=12mm]{https://helpman.komtera.lt/?fen=7b/8/8/8/4rk2/8/3qP3/1K6&moves=7.5}\\{\sffamily\tiny\color{ink}scannen zum Lösen}
\end{minipage}
\vspace{2mm}

{\raggedright\noindent\textcolor{ink}{\sffamily\textsc{Themen}}\enspace \themetag{reines Matt} \themetag{Mustermatt} \themetag{Umwandlung} \themetag{Excelsior} \themetag{Selbstblock} \themetag{Umnov}\par}

\clearpage
\phantomsection\label{hm:KPvkqrr}
\begin{tcolorbox}[enhanced,colback=ink,colframe=ink,arc=0pt,boxrule=0pt,left=4mm,right=4mm,top=2.5mm,bottom=2.5mm]
\color{white}\sffamily{\LARGE\bfseries KPvkqrr}\hfill{\Large h\#7.5}\\[0.5mm]
{\small 6 Steine\quad Weiß am Zug}\hfill{\small Nr.~467\qquad\href{https://helpman.komtera.lt/?fen=8/8/2r5/K7/1rk5/8/1P6/6q1&moves=7.5}{\color{white}\bfseries lösen\,\ensuremath{\nearrow}}}
\end{tcolorbox}
\vspace{2mm}
\noindent\begin{minipage}[t]{36mm}\vspace{0pt}\sffamily\small\raggedright
\hmfigure{längste eindeutige}{{\Large h\#7.5}}\hmfigure{eindeutig bei dieser Länge}{{\Large 22} Stellungen}\hmfigure{längstes Matt}{{\Large h\#8.5}}\hmfigure{in der Klasse}{45,4 Milliarden Stellungen,\\ 12,1 Milliarden mit Hilfsmatt}
\end{minipage}\hfill
\begin{minipage}[t]{80mm}\vspace{0pt}\begin{tikzpicture}\begin{axis}[width=72mm,height=42mm,ybar,bar width=1.6pt,xmin=-1,xmax=18,ymin=0,ymax=2.671,axis y line*=left,axis x line*=bottom,xlabel={\scriptsize Halbzüge bis Matt},ylabel={\scriptsize Stellungen (Mrd.)},tick label style={font=\tiny},label style={font=\scriptsize},enlarge x limits=false,ylabel shift=-3pt,xlabel shift=-3pt]\addplot[fill=ink!55,draw=none] coordinates {(1,0.0003054) (2,0.002017) (3,0.02113) (4,0.08256) (5,0.4512) (6,0.7068) (7,2.021) (8,1.05) (9,2.428) (10,1.022) (11,2.352) (12,0.8578) (13,1.095) (14,0.01657) (15,0.008987) (16,2.024e-05) (17,2.032e-06)};\draw[densely dotted] (axis cs:15,0) -- (axis cs:15,2.671);\end{axis}\begin{axis}[width=72mm,height=42mm,axis y line*=right,axis x line=none,ymode=log,xmin=-1,xmax=18,ylabel={\scriptsize eindeutig (log)},tick label style={font=\tiny},label style={font=\scriptsize},ylabel shift=-3pt]\addplot[mark=*,mark size=0.7pt,draw=black,thin] coordinates {(1,124474) (2,769205) (3,7950282) (4,19995741) (5,54894364) (6,24644308) (7,44125484) (8,8171756) (9,16057124) (10,4309414) (11,5024418) (12,89754) (13,53706) (14,375) (15,22)};\end{axis}\end{tikzpicture}\end{minipage}
\vspace{3mm}

\noindent\begin{minipage}[t]{60mm}\vspace{0pt}\sbox\hmboard{\chessboard[setfen=8/8/2r5/K7/1rk5/8/1P6/6q1 w - - 0 1]}\usebox\hmboard\\[1mm]
\makebox[\wd\hmboard]{\sffamily\hspace{3mm}\textbf{h\#7.5}\hfill(2+4)}\end{minipage}\hfill
\begin{minipage}[t]{52mm}\vspace{0pt}
\textcolor{ink}{\sffamily\textsc{Lösung}}\par\smallskip\small
\begin{tabular}{@{}r@{\enspace}l@{\enspace}l@{}}
\textbf{1.} & \ldots & b3+\\
\textbf{2.} & \symking{}d5 & \symking{}xb4\\
\textbf{3.} & \symrook{}a6 & \symking{}c3\\
\textbf{4.} & \symking{}c6 & b4\\
\textbf{5.} & \symqueen{}c5+ & bxc5\\
\textbf{6.} & \symking{}b7 & c6+\\
\textbf{7.} & \symking{}a8 & c7\\
\textbf{8.} & \symrook{}a7 & c8=\symqueen{}\#\\
\end{tabular}\par\vspace{3mm}
\qrcode[height=12mm]{https://helpman.komtera.lt/?fen=8/8/2r5/K7/1rk5/8/1P6/6q1&moves=7.5}\\{\sffamily\tiny\color{ink}scannen zum Lösen}
\end{minipage}
\vspace{2mm}

{\raggedright\noindent\textcolor{ink}{\sffamily\textsc{Themen}}\enspace \themetag{reines Matt} \themetag{Mustermatt} \themetag{Umwandlung} \themetag{Excelsior} \themetag{Selbstblock} \themetag{Umnov}\par}

\clearpage
\phantomsection\label{hm:KPvkrbb}
\begin{tcolorbox}[enhanced,colback=ink,colframe=ink,arc=0pt,boxrule=0pt,left=4mm,right=4mm,top=2.5mm,bottom=2.5mm]
\color{white}\sffamily{\LARGE\bfseries KPvkrbb}\hfill{\Large h\#7.5}\\[0.5mm]
{\small 6 Steine\quad Weiß am Zug}\hfill{\small Nr.~468\qquad\href{https://helpman.komtera.lt/?fen=6k1/8/8/8/b4b2/5P2/1r6/K7&moves=7.5}{\color{white}\bfseries lösen\,\ensuremath{\nearrow}}}
\end{tcolorbox}
\vspace{2mm}
\noindent\begin{minipage}[t]{36mm}\vspace{0pt}\sffamily\small\raggedright
\hmfigure{längste eindeutige}{{\Large h\#7.5}}\hmfigure{eindeutig bei dieser Länge}{{\Large 323} Stellungen}\hmfigure{längstes Matt}{{\Large h\#9.5}}\hmfigure{in der Klasse}{45,4 Milliarden Stellungen,\\ 14,3 Milliarden mit Hilfsmatt}
\end{minipage}\hfill
\begin{minipage}[t]{80mm}\vspace{0pt}\begin{tikzpicture}\begin{axis}[width=72mm,height=42mm,ybar,bar width=1.6pt,xmin=-1,xmax=20,ymin=0,ymax=2.913,axis y line*=left,axis x line*=bottom,xlabel={\scriptsize Halbzüge bis Matt},ylabel={\scriptsize Stellungen (Mrd.)},tick label style={font=\tiny},label style={font=\scriptsize},enlarge x limits=false,ylabel shift=-3pt,xlabel shift=-3pt]\addplot[fill=ink!55,draw=none] coordinates {(0,8.096e-06) (1,0.0007589) (2,0.003413) (3,0.03189) (4,0.1051) (5,0.5171) (6,0.9092) (7,2.043) (8,1.57) (9,2.648) (10,1.618) (11,2.635) (12,1.309) (13,0.9043) (14,0.033) (15,0.005363) (16,9.98e-06) (17,5.02e-05) (19,1.3e-08)};\draw[densely dotted] (axis cs:15,0) -- (axis cs:15,2.913);\end{axis}\begin{axis}[width=72mm,height=42mm,axis y line*=right,axis x line=none,ymode=log,xmin=-1,xmax=20,ylabel={\scriptsize eindeutig (log)},tick label style={font=\tiny},label style={font=\scriptsize},ylabel shift=-3pt]\addplot[mark=*,mark size=0.7pt,draw=black,thin] coordinates {(0,8096) (1,259121) (2,1087743) (3,10764695) (4,24887074) (5,57514887) (6,33887958) (7,46407775) (8,16468084) (9,24074868) (10,7220782) (11,5893964) (12,278276) (13,42462) (14,205) (15,323)};\end{axis}\end{tikzpicture}\end{minipage}
\vspace{3mm}

\noindent\begin{minipage}[t]{60mm}\vspace{0pt}\sbox\hmboard{\chessboard[setfen=6k1/8/8/8/b4b2/5P2/1r6/K7 w - - 0 1]}\usebox\hmboard\\[1mm]
\makebox[\wd\hmboard]{\sffamily\hspace{3mm}\textbf{h\#7.5}\hfill(2+4)}\end{minipage}\hfill
\begin{minipage}[t]{52mm}\vspace{0pt}
\textcolor{ink}{\sffamily\textsc{Lösung}}\par\smallskip\small
\begin{tabular}{@{}r@{\enspace}l@{\enspace}l@{}}
\textbf{1.} & \ldots & \symking{}xb2\\
\textbf{2.} & \symbishop{}e5+ & \symking{}c1\\
\textbf{3.} & \symking{}f7 & f4\\
\textbf{4.} & \symking{}e6 & fxe5\\
\textbf{5.} & \symking{}d5 & e6\\
\textbf{6.} & \symking{}c4 & e7\\
\textbf{7.} & \symking{}b3 & e8=\symqueen{}\\
\textbf{8.} & \symking{}a2 & \symqueen{}xa4\#\\
\end{tabular}\par\vspace{3mm}
\qrcode[height=12mm]{https://helpman.komtera.lt/?fen=6k1/8/8/8/b4b2/5P2/1r6/K7&moves=7.5}\\{\sffamily\tiny\color{ink}scannen zum Lösen}
\end{minipage}
\vspace{2mm}

{\raggedright\noindent\textcolor{ink}{\sffamily\textsc{Themen}}\enspace \themetag{reines Matt} \themetag{Mustermatt} \themetag{Idealmatt} \themetag{Spiegelmatt} \themetag{Umwandlung} \themetag{Umnov}\par\smallskip{\footnotesize\textit{Hinweis:} die Lösung beginnt mit einem Schlag.}\par}

\clearpage
\phantomsection\label{hm:KPvkrbn}
\begin{tcolorbox}[enhanced,colback=ink,colframe=ink,arc=0pt,boxrule=0pt,left=4mm,right=4mm,top=2.5mm,bottom=2.5mm]
\color{white}\sffamily{\LARGE\bfseries KPvkrbn}\hfill{\Large h\#7.5}\\[0.5mm]
{\small 6 Steine\quad Weiß am Zug}\hfill{\small Nr.~469\qquad\href{https://helpman.komtera.lt/?fen=8/n6k/8/8/b7/r7/P7/2K5&moves=7.5}{\color{white}\bfseries lösen\,\ensuremath{\nearrow}}}
\end{tcolorbox}
\vspace{2mm}
\noindent\begin{minipage}[t]{36mm}\vspace{0pt}\sffamily\small\raggedright
\hmfigure{längste eindeutige}{{\Large h\#7.5}}\hmfigure{eindeutig bei dieser Länge}{{\Large 744} Stellungen}\hmfigure{längstes Matt}{{\Large h\#9.5}}\hmfigure{in der Klasse}{45,4 Milliarden Stellungen,\\ 29,3 Milliarden mit Hilfsmatt}
\end{minipage}\hfill
\begin{minipage}[t]{80mm}\vspace{0pt}\begin{tikzpicture}\begin{axis}[width=72mm,height=42mm,ybar,bar width=1.6pt,xmin=-1,xmax=20,ymin=0,ymax=5.91,axis y line*=left,axis x line*=bottom,xlabel={\scriptsize Halbzüge bis Matt},ylabel={\scriptsize Stellungen (Mrd.)},tick label style={font=\tiny},label style={font=\scriptsize},enlarge x limits=false,ylabel shift=-3pt,xlabel shift=-3pt]\addplot[fill=ink!55,draw=none] coordinates {(0,1.226e-05) (1,0.001839) (2,0.007798) (3,0.06781) (4,0.2055) (5,0.9687) (6,1.707) (7,3.784) (8,3.211) (9,5.244) (10,3.394) (11,5.373) (12,3.063) (13,2.128) (14,0.1559) (15,0.0247) (16,3.91e-05) (17,0.0001329) (19,8.1e-08)};\draw[densely dotted] (axis cs:15,0) -- (axis cs:15,5.91);\end{axis}\begin{axis}[width=72mm,height=42mm,axis y line*=right,axis x line=none,ymode=log,xmin=-1,xmax=20,ylabel={\scriptsize eindeutig (log)},tick label style={font=\tiny},label style={font=\scriptsize},ylabel shift=-3pt]\addplot[mark=*,mark size=0.7pt,draw=black,thin] coordinates {(0,12260) (1,558901) (2,2211372) (3,20309035) (4,46136745) (5,104311432) (6,62730713) (7,82862660) (8,33599507) (9,48136842) (10,16152664) (11,14943931) (12,1314104) (13,300477) (14,815) (15,744)};\end{axis}\end{tikzpicture}\end{minipage}
\vspace{3mm}

\noindent\begin{minipage}[t]{60mm}\vspace{0pt}\sbox\hmboard{\chessboard[setfen=8/n6k/8/8/b7/r7/P7/2K5 w - - 0 1]}\usebox\hmboard\\[1mm]
\makebox[\wd\hmboard]{\sffamily\hspace{3mm}\textbf{h\#7.5}\hfill(2+4)}\end{minipage}\hfill
\begin{minipage}[t]{52mm}\vspace{0pt}
\textcolor{ink}{\sffamily\textsc{Lösung}}\par\smallskip\small
\begin{tabular}{@{}r@{\enspace}l@{\enspace}l@{}}
\textbf{1.} & \ldots & \symking{}d2\\
\textbf{2.} & \symbishop{}b3 & \symking{}e3\\
\textbf{3.} & \symrook{}a6 & a4\\
\textbf{4.} & \symknight{}b5 & axb5\\
\textbf{5.} & \symking{}g6 & b6\\
\textbf{6.} & \symking{}f5 & b7\\
\textbf{7.} & \symrook{}g6 & b8=\symqueen{}\\
\textbf{8.} & \symbishop{}e6 & \symqueen{}f4\#\\
\end{tabular}\par\vspace{3mm}
\qrcode[height=12mm]{https://helpman.komtera.lt/?fen=8/n6k/8/8/b7/r7/P7/2K5&moves=7.5}\\{\sffamily\tiny\color{ink}scannen zum Lösen}
\end{minipage}
\vspace{2mm}

{\raggedright\noindent\textcolor{ink}{\sffamily\textsc{Themen}}\enspace \themetag{Umwandlung} \themetag{Excelsior} \themetag{Selbstblock} \themetag{ohne Schach}\par}

\clearpage
\phantomsection\label{hm:KPvkrnn}
\begin{tcolorbox}[enhanced,colback=ink,colframe=ink,arc=0pt,boxrule=0pt,left=4mm,right=4mm,top=2.5mm,bottom=2.5mm]
\color{white}\sffamily{\LARGE\bfseries KPvkrnn}\hfill{\Large h\#7.5}\\[0.5mm]
{\small 6 Steine\quad Weiß am Zug}\hfill{\small Nr.~470\qquad\href{https://helpman.komtera.lt/?fen=8/8/6k1/1r6/n6n/8/7P/K7&moves=7.5}{\color{white}\bfseries lösen\,\ensuremath{\nearrow}}}
\end{tcolorbox}
\vspace{2mm}
\noindent\begin{minipage}[t]{36mm}\vspace{0pt}\sffamily\small\raggedright
\hmfigure{längste eindeutige}{{\Large h\#7.5}}\hmfigure{eindeutig bei dieser Länge}{{\Large 680} Stellungen}\hmfigure{längstes Matt}{{\Large h\#9.5}}\hmfigure{in der Klasse}{45,4 Milliarden Stellungen,\\ 15 Milliarden mit Hilfsmatt}
\end{minipage}\hfill
\begin{minipage}[t]{80mm}\vspace{0pt}\begin{tikzpicture}\begin{axis}[width=72mm,height=42mm,ybar,bar width=1.6pt,xmin=-1,xmax=20,ymin=0,ymax=2.963,axis y line*=left,axis x line*=bottom,xlabel={\scriptsize Halbzüge bis Matt},ylabel={\scriptsize Stellungen (Mrd.)},tick label style={font=\tiny},label style={font=\scriptsize},enlarge x limits=false,ylabel shift=-3pt,xlabel shift=-3pt]\addplot[fill=ink!55,draw=none] coordinates {(0,4.188e-06) (1,0.001108) (2,0.004399) (3,0.03591) (4,0.09853) (5,0.4418) (6,0.7285) (7,1.628) (8,1.53) (9,2.509) (10,1.731) (11,2.693) (12,1.801) (13,1.445) (14,0.2592) (15,0.04554) (16,0.0002261) (17,0.0002668) (18,1.2e-08) (19,1.62e-07)};\draw[densely dotted] (axis cs:15,0) -- (axis cs:15,2.963);\end{axis}\begin{axis}[width=72mm,height=42mm,axis y line*=right,axis x line=none,ymode=log,xmin=-1,xmax=20,ylabel={\scriptsize eindeutig (log)},tick label style={font=\tiny},label style={font=\scriptsize},ylabel shift=-3pt]\addplot[mark=*,mark size=0.7pt,draw=black,thin] coordinates {(0,4188) (1,295143) (2,1089937) (3,9377738) (4,20159729) (5,45532711) (6,26836662) (7,35484490) (8,15828094) (9,22972548) (10,8747417) (11,10834851) (12,1930665) (13,693590) (14,2322) (15,680)};\end{axis}\end{tikzpicture}\end{minipage}
\vspace{3mm}

\noindent\begin{minipage}[t]{60mm}\vspace{0pt}\sbox\hmboard{\chessboard[setfen=8/8/6k1/1r6/n6n/8/7P/K7 w - - 0 1]}\usebox\hmboard\\[1mm]
\makebox[\wd\hmboard]{\sffamily\hspace{3mm}\textbf{h\#7.5}\hfill(2+4)}\end{minipage}\hfill
\begin{minipage}[t]{52mm}\vspace{0pt}
\textcolor{ink}{\sffamily\textsc{Lösung}}\par\smallskip\small
\begin{tabular}{@{}r@{\enspace}l@{\enspace}l@{}}
\textbf{1.} & \ldots & \symking{}a2\\
\textbf{2.} & \symking{}f5 & \symking{}a3\\
\textbf{3.} & \symknight{}g6 & h4\\
\textbf{4.} & \symking{}e4 & h5\\
\textbf{5.} & \symking{}d3 & hxg6\\
\textbf{6.} & \symking{}c2 & g7\\
\textbf{7.} & \symking{}b1 & g8=\symqueen{}\\
\textbf{8.} & \symking{}a1 & \symqueen{}a2\#\\
\end{tabular}\par\vspace{3mm}
\qrcode[height=12mm]{https://helpman.komtera.lt/?fen=8/8/6k1/1r6/n6n/8/7P/K7&moves=7.5}\\{\sffamily\tiny\color{ink}scannen zum Lösen}
\end{minipage}
\vspace{2mm}

{\raggedright\noindent\textcolor{ink}{\sffamily\textsc{Themen}}\enspace \themetag{Umwandlung} \themetag{Excelsior} \themetag{ohne Schach} \themetag{Umnov}\par}

\clearpage
\phantomsection\label{hm:KPvkrrb}
\begin{tcolorbox}[enhanced,colback=ink,colframe=ink,arc=0pt,boxrule=0pt,left=4mm,right=4mm,top=2.5mm,bottom=2.5mm]
\color{white}\sffamily{\LARGE\bfseries KPvkrrb}\hfill{\Large h\#7.5}\\[0.5mm]
{\small 6 Steine\quad Weiß am Zug}\hfill{\small Nr.~471\qquad\href{https://helpman.komtera.lt/?fen=8/7k/8/1r6/K7/b7/Pr6/8&moves=7.5}{\color{white}\bfseries lösen\,\ensuremath{\nearrow}}}
\end{tcolorbox}
\vspace{2mm}
\noindent\begin{minipage}[t]{36mm}\vspace{0pt}\sffamily\small\raggedright
\hmfigure{längste eindeutige}{{\Large h\#7.5}}\hmfigure{eindeutig bei dieser Länge}{{\Large 21} Stellungen}\hmfigure{längstes Matt}{{\Large h\#8.5}}\hmfigure{in der Klasse}{45,4 Milliarden Stellungen,\\ 13,8 Milliarden mit Hilfsmatt}
\end{minipage}\hfill
\begin{minipage}[t]{80mm}\vspace{0pt}\begin{tikzpicture}\begin{axis}[width=72mm,height=42mm,ybar,bar width=1.6pt,xmin=-1,xmax=18,ymin=0,ymax=2.822,axis y line*=left,axis x line*=bottom,xlabel={\scriptsize Halbzüge bis Matt},ylabel={\scriptsize Stellungen (Mrd.)},tick label style={font=\tiny},label style={font=\scriptsize},enlarge x limits=false,ylabel shift=-3pt,xlabel shift=-3pt]\addplot[fill=ink!55,draw=none] coordinates {(0,3.305e-06) (1,0.0006564) (2,0.003522) (3,0.03244) (4,0.1199) (5,0.5723) (6,0.9632) (7,2.174) (8,1.414) (9,2.566) (10,1.41) (11,2.47) (12,1.116) (13,0.8915) (14,0.02163) (15,0.005284) (16,1.203e-05) (17,1.57e-06)};\draw[densely dotted] (axis cs:15,0) -- (axis cs:15,2.822);\end{axis}\begin{axis}[width=72mm,height=42mm,axis y line*=right,axis x line=none,ymode=log,xmin=-1,xmax=18,ylabel={\scriptsize eindeutig (log)},tick label style={font=\tiny},label style={font=\scriptsize},ylabel shift=-3pt]\addplot[mark=*,mark size=0.7pt,draw=black,thin] coordinates {(0,3305) (1,252442) (2,1365728) (3,12037238) (4,29249137) (5,63360551) (6,33200443) (7,45991884) (8,12600035) (9,19847282) (10,6377614) (11,4818760) (12,172722) (13,37578) (14,198) (15,21)};\end{axis}\end{tikzpicture}\end{minipage}
\vspace{3mm}

\noindent\begin{minipage}[t]{60mm}\vspace{0pt}\sbox\hmboard{\chessboard[setfen=8/7k/8/1r6/K7/b7/Pr6/8 w - - 0 1]}\usebox\hmboard\\[1mm]
\makebox[\wd\hmboard]{\sffamily\hspace{3mm}\textbf{h\#7.5}\hfill(2+4)}\end{minipage}\hfill
\begin{minipage}[t]{52mm}\vspace{0pt}
\textcolor{ink}{\sffamily\textsc{Lösung}}\par\smallskip\small
\begin{tabular}{@{}r@{\enspace}l@{\enspace}l@{}}
\textbf{1.} & \ldots & \symking{}xa3\\
\textbf{2.} & \symrook{}5b3+ & axb3\\
\textbf{3.} & \symking{}g6 & b4\\
\textbf{4.} & \symking{}f5 & b5\\
\textbf{5.} & \symking{}e4 & b6\\
\textbf{6.} & \symking{}d3 & b7\\
\textbf{7.} & \symking{}c2 & b8=\symqueen{}\\
\textbf{8.} & \symking{}b1 & \symqueen{}xb2\#\\
\end{tabular}\par\vspace{3mm}
\qrcode[height=12mm]{https://helpman.komtera.lt/?fen=8/7k/8/1r6/K7/b7/Pr6/8&moves=7.5}\\{\sffamily\tiny\color{ink}scannen zum Lösen}
\end{minipage}
\vspace{2mm}

{\raggedright\noindent\textcolor{ink}{\sffamily\textsc{Themen}}\enspace \themetag{Satzspiel} \themetag{Umwandlung} \themetag{Excelsior}\par\smallskip{\footnotesize\textit{Hinweis:} die Lösung beginnt mit einem Schlag.}\par}

\clearpage
\phantomsection\label{hm:KPvkrrn}
\begin{tcolorbox}[enhanced,colback=ink,colframe=ink,arc=0pt,boxrule=0pt,left=4mm,right=4mm,top=2.5mm,bottom=2.5mm]
\color{white}\sffamily{\LARGE\bfseries KPvkrrn}\hfill{\Large h\#7.5}\\[0.5mm]
{\small 6 Steine\quad Weiß am Zug}\hfill{\small Nr.~472\qquad\href{https://helpman.komtera.lt/?fen=8/8/2r5/K7/1rk1n3/8/1P6/8&moves=7.5}{\color{white}\bfseries lösen\,\ensuremath{\nearrow}}}
\end{tcolorbox}
\vspace{2mm}
\noindent\begin{minipage}[t]{36mm}\vspace{0pt}\sffamily\small\raggedright
\hmfigure{längste eindeutige}{{\Large h\#7.5}}\hmfigure{eindeutig bei dieser Länge}{{\Large 60} Stellungen}\hmfigure{längstes Matt}{{\Large h\#9.5}}\hmfigure{in der Klasse}{45,4 Milliarden Stellungen,\\ 14 Milliarden mit Hilfsmatt}
\end{minipage}\hfill
\begin{minipage}[t]{80mm}\vspace{0pt}\begin{tikzpicture}\begin{axis}[width=72mm,height=42mm,ybar,bar width=1.6pt,xmin=-1,xmax=20,ymin=0,ymax=2.845,axis y line*=left,axis x line*=bottom,xlabel={\scriptsize Halbzüge bis Matt},ylabel={\scriptsize Stellungen (Mrd.)},tick label style={font=\tiny},label style={font=\scriptsize},enlarge x limits=false,ylabel shift=-3pt,xlabel shift=-3pt]\addplot[fill=ink!55,draw=none] coordinates {(0,1.65e-06) (1,0.0007579) (2,0.003849) (3,0.0338) (4,0.117) (5,0.5361) (6,0.9239) (7,2.047) (8,1.471) (9,2.586) (10,1.469) (11,2.547) (12,1.273) (13,0.9608) (14,0.03288) (15,0.009582) (16,2.238e-05) (17,1.463e-05) (19,2e-09)};\draw[densely dotted] (axis cs:15,0) -- (axis cs:15,2.845);\end{axis}\begin{axis}[width=72mm,height=42mm,axis y line*=right,axis x line=none,ymode=log,xmin=-1,xmax=20,ylabel={\scriptsize eindeutig (log)},tick label style={font=\tiny},label style={font=\scriptsize},ylabel shift=-3pt]\addplot[mark=*,mark size=0.7pt,draw=black,thin] coordinates {(0,1650) (1,263489) (2,1357902) (3,11450617) (4,27685021) (5,57372940) (6,31033417) (7,41342766) (8,11952621) (9,18486252) (10,6541771) (11,5008128) (12,223396) (13,61556) (14,161) (15,60)};\end{axis}\end{tikzpicture}\end{minipage}
\vspace{3mm}

\noindent\begin{minipage}[t]{60mm}\vspace{0pt}\sbox\hmboard{\chessboard[setfen=8/8/2r5/K7/1rk1n3/8/1P6/8 w - - 0 1]}\usebox\hmboard\\[1mm]
\makebox[\wd\hmboard]{\sffamily\hspace{3mm}\textbf{h\#7.5}\hfill(2+4)}\end{minipage}\hfill
\begin{minipage}[t]{52mm}\vspace{0pt}
\textcolor{ink}{\sffamily\textsc{Lösung}}\par\smallskip\small
\begin{tabular}{@{}r@{\enspace}l@{\enspace}l@{}}
\textbf{1.} & \ldots & b3+\\
\textbf{2.} & \symking{}d5 & \symking{}xb4\\
\textbf{3.} & \symknight{}c5 & \symking{}c3\\
\textbf{4.} & \symrook{}a6 & b4\\
\textbf{5.} & \symking{}c6 & bxc5\\
\textbf{6.} & \symking{}b7 & c6+\\
\textbf{7.} & \symking{}a8 & c7\\
\textbf{8.} & \symrook{}a7 & c8=\symqueen{}\#\\
\end{tabular}\par\vspace{3mm}
\qrcode[height=12mm]{https://helpman.komtera.lt/?fen=8/8/2r5/K7/1rk1n3/8/1P6/8&moves=7.5}\\{\sffamily\tiny\color{ink}scannen zum Lösen}
\end{minipage}
\vspace{2mm}

{\raggedright\noindent\textcolor{ink}{\sffamily\textsc{Themen}}\enspace \themetag{reines Matt} \themetag{Mustermatt} \themetag{Umwandlung} \themetag{Excelsior} \themetag{Selbstblock} \themetag{Umnov}\par}

\clearpage
\phantomsection\label{hm:KPvkrrr}
\begin{tcolorbox}[enhanced,colback=ink,colframe=ink,arc=0pt,boxrule=0pt,left=4mm,right=4mm,top=2.5mm,bottom=2.5mm]
\color{white}\sffamily{\LARGE\bfseries KPvkrrr}\hfill{\Large h\#7.5}\\[0.5mm]
{\small 6 Steine\quad Weiß am Zug}\hfill{\small Nr.~473\qquad\href{https://helpman.komtera.lt/?fen=8/7k/8/1r6/K7/r7/Pr6/8&moves=7.5}{\color{white}\bfseries lösen\,\ensuremath{\nearrow}}}
\end{tcolorbox}
\vspace{2mm}
\noindent\begin{minipage}[t]{36mm}\vspace{0pt}\sffamily\small\raggedright
\hmfigure{längste eindeutige}{{\Large h\#7.5}}\hmfigure{eindeutig bei dieser Länge}{{\Large 6} Stellungen}\hmfigure{längstes Matt}{{\Large h\#8.5}}\hmfigure{in der Klasse}{45,4 Milliarden Stellungen,\\ 4,32 Milliarden mit Hilfsmatt}
\end{minipage}\hfill
\begin{minipage}[t]{80mm}\vspace{0pt}\begin{tikzpicture}\begin{axis}[width=72mm,height=42mm,ybar,bar width=1.6pt,xmin=-1,xmax=18,ymin=0,ymax=879.3,axis y line*=left,axis x line*=bottom,xlabel={\scriptsize Halbzüge bis Matt},ylabel={\scriptsize Stellungen (Mio.)},tick label style={font=\tiny},label style={font=\scriptsize},enlarge x limits=false,ylabel shift=-3pt,xlabel shift=-3pt]\addplot[fill=ink!55,draw=none] coordinates {(1,0.1709) (2,1.069) (3,10.23) (4,40.97) (5,196.5) (6,304.2) (7,727.6) (8,399.4) (9,799.4) (10,395.6) (11,781.8) (12,332.1) (13,323.2) (14,6.367) (15,2.217) (16,0.006618) (17,0.000386)};\draw[densely dotted] (axis cs:15,0) -- (axis cs:15,879.3);\end{axis}\begin{axis}[width=72mm,height=42mm,axis y line*=right,axis x line=none,ymode=log,xmin=-1,xmax=18,ylabel={\scriptsize eindeutig (log)},tick label style={font=\tiny},label style={font=\scriptsize},ylabel shift=-3pt]\addplot[mark=*,mark size=0.7pt,draw=black,thin] coordinates {(1,73285) (2,462364) (3,4089027) (4,9960187) (5,21363798) (6,9726768) (7,14040420) (8,2979944) (9,5199493) (10,1886233) (11,1518806) (12,35744) (13,13136) (14,136) (15,6)};\end{axis}\end{tikzpicture}\end{minipage}
\vspace{3mm}

\noindent\begin{minipage}[t]{60mm}\vspace{0pt}\sbox\hmboard{\chessboard[setfen=8/7k/8/1r6/K7/r7/Pr6/8 w - - 0 1]}\usebox\hmboard\\[1mm]
\makebox[\wd\hmboard]{\sffamily\hspace{3mm}\textbf{h\#7.5}\hfill(2+4)}\end{minipage}\hfill
\begin{minipage}[t]{52mm}\vspace{0pt}
\textcolor{ink}{\sffamily\textsc{Lösung}}\par\smallskip\small
\begin{tabular}{@{}r@{\enspace}l@{\enspace}l@{}}
\textbf{1.} & \ldots & \symking{}xa3\\
\textbf{2.} & \symrook{}5b3+ & axb3\\
\textbf{3.} & \symking{}g6 & b4\\
\textbf{4.} & \symking{}f5 & b5\\
\textbf{5.} & \symking{}e4 & b6\\
\textbf{6.} & \symking{}d3 & b7\\
\textbf{7.} & \symking{}c2 & b8=\symqueen{}\\
\textbf{8.} & \symking{}b1 & \symqueen{}xb2\#\\
\end{tabular}\par\vspace{3mm}
\qrcode[height=12mm]{https://helpman.komtera.lt/?fen=8/7k/8/1r6/K7/r7/Pr6/8&moves=7.5}\\{\sffamily\tiny\color{ink}scannen zum Lösen}
\end{minipage}
\vspace{2mm}

{\raggedright\noindent\textcolor{ink}{\sffamily\textsc{Themen}}\enspace \themetag{Umwandlung} \themetag{Excelsior}\par\smallskip{\footnotesize\textit{Hinweis:} die Partei am Zug steht im Diagramm im Schach; die Lösung beginnt mit einem Schlag.}\par}

\clearpage
\phantomsection\label{hm:KQNPvkn}
\begin{tcolorbox}[enhanced,colback=ink,colframe=ink,arc=0pt,boxrule=0pt,left=4mm,right=4mm,top=2.5mm,bottom=2.5mm]
\color{white}\sffamily{\LARGE\bfseries KQNPvkn}\hfill{\Large h\#7.5}\\[0.5mm]
{\small 6 Steine\quad Weiß am Zug}\hfill{\small Nr.~474\qquad\href{https://helpman.komtera.lt/?fen=K1k5/Q1n5/8/8/8/2N5/2P5/8&moves=7.5}{\color{white}\bfseries lösen\,\ensuremath{\nearrow}}}
\end{tcolorbox}
\vspace{2mm}
\noindent\begin{minipage}[t]{36mm}\vspace{0pt}\sffamily\small\raggedright
\hmfigure{längste eindeutige}{{\Large h\#7.5}}\hmfigure{eindeutig bei dieser Länge}{{\Large 1} Stellung}\hmfigure{längstes Matt}{{\Large h\#9}}\hmfigure{in der Klasse}{45,4 Milliarden Stellungen,\\ 27,6 Milliarden mit Hilfsmatt}
\end{minipage}\hfill
\begin{minipage}[t]{80mm}\vspace{0pt}\begin{tikzpicture}\begin{axis}[width=72mm,height=42mm,ybar,bar width=1.6pt,xmin=-1,xmax=19,ymin=0,ymax=10.94,axis y line*=left,axis x line*=bottom,xlabel={\scriptsize Halbzüge bis Matt},ylabel={\scriptsize Stellungen (Mrd.)},tick label style={font=\tiny},label style={font=\scriptsize},enlarge x limits=false,ylabel shift=-3pt,xlabel shift=-3pt]\addplot[fill=ink!55,draw=none] coordinates {(0,0.05003) (1,0.3683) (2,1.394) (3,4.325) (4,9.948) (5,5.686) (6,4.886) (7,0.6037) (8,0.2301) (9,0.003429) (10,0.06024) (11,1.903e-06) (12,0.06112) (13,2.409e-06) (14,0.02302) (15,3.16e-06) (16,0.0003117) (17,1.94e-07) (18,3.6e-08)};\draw[densely dotted] (axis cs:15,0) -- (axis cs:15,10.94);\end{axis}\begin{axis}[width=72mm,height=42mm,axis y line*=right,axis x line=none,ymode=log,xmin=-1,xmax=19,ylabel={\scriptsize eindeutig (log)},tick label style={font=\tiny},label style={font=\scriptsize},ylabel shift=-3pt]\addplot[mark=*,mark size=0.7pt,draw=black,thin] coordinates {(0,50031037) (1,283189391) (2,760390550) (3,706217185) (4,672510503) (5,58363816) (6,15971382) (7,191581) (8,1070707) (9,55) (10,564911) (11,9) (12,273517) (13,17) (14,4917) (15,1)};\end{axis}\end{tikzpicture}\end{minipage}
\vspace{3mm}

\noindent\begin{minipage}[t]{60mm}\vspace{0pt}\sbox\hmboard{\chessboard[setfen=K1k5/Q1n5/8/8/8/2N5/2P5/8 w - - 0 1]}\usebox\hmboard\\[1mm]
\makebox[\wd\hmboard]{\sffamily\hspace{3mm}\textbf{h\#7.5}\hfill(4+2)}\end{minipage}\hfill
\begin{minipage}[t]{52mm}\vspace{0pt}
\textcolor{ink}{\sffamily\textsc{Lösung}}\par\smallskip\small
\begin{tabular}{@{}r@{\enspace}l@{\enspace}l@{}}
\textbf{1.} & \ldots & \symqueen{}xc7+\\
\textbf{2.} & \symking{}xc7 & \symknight{}e2\\
\textbf{3.} & \symking{}d6 & c4\\
\textbf{4.} & \symking{}e5 & c5\\
\textbf{5.} & \symking{}e4 & c6\\
\textbf{6.} & \symking{}f3 & c7\\
\textbf{7.} & \symking{}g2 & c8=\symqueen{}\\
\textbf{8.} & \symking{}h1 & \symqueen{}h3\#\\
\end{tabular}\par\vspace{3mm}
\qrcode[height=12mm]{https://helpman.komtera.lt/?fen=K1k5/Q1n5/8/8/8/2N5/2P5/8&moves=7.5}\\{\sffamily\tiny\color{ink}scannen zum Lösen}
\end{minipage}
\vspace{2mm}

{\raggedright\noindent\textcolor{ink}{\sffamily\textsc{Themen}}\enspace \themetag{reines Matt} \themetag{Mustermatt} \themetag{Spiegelmatt} \themetag{Umwandlung} \themetag{Excelsior} \themetag{ein Stein zieht} \themetag{Phönix}\par\smallskip{\footnotesize\textit{Hinweis:} die Partei am Zug steht im Diagramm im Schach; die Lösung beginnt mit einem Schlag.}\par}

\clearpage
\phantomsection\label{hm:KQNvkqr}
\begin{tcolorbox}[enhanced,colback=ink,colframe=ink,arc=0pt,boxrule=0pt,left=4mm,right=4mm,top=2.5mm,bottom=2.5mm]
\color{white}\sffamily{\LARGE\bfseries KQNvkqr}\hfill{\Large h\#7.5}\\[0.5mm]
{\small 6 Steine\quad Weiß am Zug}\hfill{\small Nr.~475\qquad\href{https://helpman.komtera.lt/?fen=8/8/N7/8/8/7Q/7r/K4qk1&moves=7.5}{\color{white}\bfseries lösen\,\ensuremath{\nearrow}}}
\end{tcolorbox}
\vspace{2mm}
\noindent\begin{minipage}[t]{36mm}\vspace{0pt}\sffamily\small\raggedright
\hmfigure{längste eindeutige}{{\Large h\#7.5}}\hmfigure{eindeutig bei dieser Länge}{{\Large 23} Stellungen}\hmfigure{längstes Matt}{{\Large h\#9.5}}\hmfigure{in der Klasse}{15,5 Milliarden Stellungen,\\ 6,94 Milliarden mit Hilfsmatt}
\end{minipage}\hfill
\begin{minipage}[t]{80mm}\vspace{0pt}\begin{tikzpicture}\begin{axis}[width=72mm,height=42mm,ybar,bar width=1.6pt,xmin=-1,xmax=20,ymin=0,ymax=2.62,axis y line*=left,axis x line*=bottom,xlabel={\scriptsize Halbzüge bis Matt},ylabel={\scriptsize Stellungen (Mrd.)},tick label style={font=\tiny},label style={font=\scriptsize},enlarge x limits=false,ylabel shift=-3pt,xlabel shift=-3pt]\addplot[fill=ink!55,draw=none] coordinates {(0,0.00428) (1,0.0386) (2,0.2123) (3,0.9427) (4,1.997) (5,2.382) (6,0.9413) (7,0.3602) (8,0.01924) (9,0.0007754) (10,0.01833) (11,4.418e-06) (12,0.01432) (13,5.822e-06) (14,0.004541) (15,2.958e-06) (16,0.0002793) (17,8.53e-07) (18,2.87e-07) (19,1.76e-07)};\draw[densely dotted] (axis cs:15,0) -- (axis cs:15,2.62);\end{axis}\begin{axis}[width=72mm,height=42mm,axis y line*=right,axis x line=none,ymode=log,xmin=-1,xmax=20,ylabel={\scriptsize eindeutig (log)},tick label style={font=\tiny},label style={font=\scriptsize},ylabel shift=-3pt]\addplot[mark=*,mark size=0.7pt,draw=black,thin] coordinates {(0,4280175) (1,35296770) (2,124006994) (3,198932533) (4,135334425) (5,40146426) (6,1458341) (7,64657) (8,358153) (9,46) (10,179847) (11,67) (12,51707) (14,4447) (15,23)};\end{axis}\end{tikzpicture}\end{minipage}
\vspace{3mm}

\noindent\begin{minipage}[t]{60mm}\vspace{0pt}\sbox\hmboard{\chessboard[setfen=8/8/N7/8/8/7Q/7r/K4qk1 w - - 0 1]}\usebox\hmboard\\[1mm]
\makebox[\wd\hmboard]{\sffamily\hspace{3mm}\textbf{h\#7.5}\hfill(3+3)}\end{minipage}\hfill
\begin{minipage}[t]{52mm}\vspace{0pt}
\textcolor{ink}{\sffamily\textsc{Lösung}}\par\smallskip\small
\begin{tabular}{@{}r@{\enspace}l@{\enspace}l@{}}
\textbf{1.} & \ldots & \symqueen{}xf1+\\
\textbf{2.} & \symking{}xf1 & \symking{}b1\\
\textbf{3.} & \symking{}e2 & \symking{}b2\\
\textbf{4.} & \symking{}d3+ & \symking{}c1\\
\textbf{5.} & \symking{}c3 & \symking{}d1\\
\textbf{6.} & \symking{}b2 & \symknight{}c5\\
\textbf{7.} & \symking{}a1 & \symking{}c1\\
\textbf{8.} & \symrook{}a2 & \symknight{}b3\#\\
\end{tabular}\par\vspace{3mm}
\qrcode[height=12mm]{https://helpman.komtera.lt/?fen=8/8/N7/8/8/7Q/7r/K4qk1&moves=7.5}\\{\sffamily\tiny\color{ink}scannen zum Lösen}
\end{minipage}
\vspace{2mm}

{\raggedright\noindent\textcolor{ink}{\sffamily\textsc{Themen}}\enspace \themetag{reines Matt} \themetag{Mustermatt} \themetag{Idealmatt} \themetag{Switchback} \themetag{Selbstblock}\par\smallskip{\footnotesize\textit{Hinweis:} die Partei am Zug steht im Diagramm im Schach; die Lösung beginnt mit einem Schlag.}\par}

\clearpage
\phantomsection\label{hm:KRNPvkq}
\begin{tcolorbox}[enhanced,colback=ink,colframe=ink,arc=0pt,boxrule=0pt,left=4mm,right=4mm,top=2.5mm,bottom=2.5mm]
\color{white}\sffamily{\LARGE\bfseries KRNPvkq}\hfill{\Large h\#7.5}\\[0.5mm]
{\small 6 Steine\quad Weiß am Zug}\hfill{\small Nr.~476\qquad\href{https://helpman.komtera.lt/?fen=8/8/8/8/8/5N2/q4PR1/K6k&moves=7.5}{\color{white}\bfseries lösen\,\ensuremath{\nearrow}}}
\end{tcolorbox}
\vspace{2mm}
\noindent\begin{minipage}[t]{36mm}\vspace{0pt}\sffamily\small\raggedright
\hmfigure{längste eindeutige}{{\Large h\#7.5}}\hmfigure{eindeutig bei dieser Länge}{{\Large 12} Stellungen}\hmfigure{längstes Matt}{{\Large h\#8.5}}\hmfigure{in der Klasse}{45,4 Milliarden Stellungen,\\ 25,3 Milliarden mit Hilfsmatt}
\end{minipage}\hfill
\begin{minipage}[t]{80mm}\vspace{0pt}\begin{tikzpicture}\begin{axis}[width=72mm,height=42mm,ybar,bar width=1.6pt,xmin=-1,xmax=18,ymin=0,ymax=7.896,axis y line*=left,axis x line*=bottom,xlabel={\scriptsize Halbzüge bis Matt},ylabel={\scriptsize Stellungen (Mrd.)},tick label style={font=\tiny},label style={font=\scriptsize},enlarge x limits=false,ylabel shift=-3pt,xlabel shift=-3pt]\addplot[fill=ink!55,draw=none] coordinates {(0,0.006561) (1,0.04439) (2,0.1931) (3,1.153) (4,2.913) (5,6.146) (6,7.178) (7,5.186) (8,1.943) (9,0.4818) (10,0.0376) (11,0.006142) (12,0.00309) (13,0.0001086) (14,0.0005429) (15,5.158e-05) (16,5.77e-06) (17,1.404e-06)};\draw[densely dotted] (axis cs:15,0) -- (axis cs:15,7.896);\end{axis}\begin{axis}[width=72mm,height=42mm,axis y line*=right,axis x line=none,ymode=log,xmin=-1,xmax=18,ylabel={\scriptsize eindeutig (log)},tick label style={font=\tiny},label style={font=\scriptsize},ylabel shift=-3pt]\addplot[mark=*,mark size=0.7pt,draw=black,thin] coordinates {(0,6560553) (1,36827261) (2,85817785) (3,340821150) (4,523254543) (5,436421318) (6,130053367) (7,25642655) (8,1302164) (9,146540) (10,16475) (11,1518) (12,4451) (13,183) (14,40) (15,12)};\end{axis}\end{tikzpicture}\end{minipage}
\vspace{3mm}

\noindent\begin{minipage}[t]{60mm}\vspace{0pt}\sbox\hmboard{\chessboard[setfen=8/8/8/8/8/5N2/q4PR1/K6k w - - 0 1]}\usebox\hmboard\\[1mm]
\makebox[\wd\hmboard]{\sffamily\hspace{3mm}\textbf{h\#7.5}\hfill(4+2)}\end{minipage}\hfill
\begin{minipage}[t]{52mm}\vspace{0pt}
\textcolor{ink}{\sffamily\textsc{Lösung}}\par\smallskip\small
\begin{tabular}{@{}r@{\enspace}l@{\enspace}l@{}}
\textbf{1.} & \ldots & \symking{}xa2\\
\textbf{2.} & \symking{}xg2 & \symknight{}d4\\
\textbf{3.} & \symking{}f1 & f4\\
\textbf{4.} & \symking{}e1 & f5\\
\textbf{5.} & \symking{}d2 & f6\\
\textbf{6.} & \symking{}c3 & f7\\
\textbf{7.} & \symking{}b4 & f8=\symqueen{}+\\
\textbf{8.} & \symking{}a4 & \symqueen{}a3\#\\
\end{tabular}\par\vspace{3mm}
\qrcode[height=12mm]{https://helpman.komtera.lt/?fen=8/8/8/8/8/5N2/q4PR1/K6k&moves=7.5}\\{\sffamily\tiny\color{ink}scannen zum Lösen}
\end{minipage}
\vspace{2mm}

{\raggedright\noindent\textcolor{ink}{\sffamily\textsc{Themen}}\enspace \themetag{Umwandlung} \themetag{Excelsior} \themetag{ein Stein zieht}\par\smallskip{\footnotesize\textit{Hinweis:} die Partei am Zug steht im Diagramm im Schach; die Lösung beginnt mit einem Schlag.}\par}

\clearpage
\phantomsection\label{hm:KRNPvkr}
\begin{tcolorbox}[enhanced,colback=ink,colframe=ink,arc=0pt,boxrule=0pt,left=4mm,right=4mm,top=2.5mm,bottom=2.5mm]
\color{white}\sffamily{\LARGE\bfseries KRNPvkr}\hfill{\Large h\#7.5}\\[0.5mm]
{\small 6 Steine\quad Weiß am Zug}\hfill{\small Nr.~477\qquad\href{https://helpman.komtera.lt/?fen=KN6/r7/8/8/8/8/3P2R1/7k&moves=7.5}{\color{white}\bfseries lösen\,\ensuremath{\nearrow}}}
\end{tcolorbox}
\vspace{2mm}
\noindent\begin{minipage}[t]{36mm}\vspace{0pt}\sffamily\small\raggedright
\hmfigure{längste eindeutige}{{\Large h\#7.5}}\hmfigure{eindeutig bei dieser Länge}{{\Large 2} Stellungen}\hmfigure{längstes Matt}{{\Large h\#8.5}}\hmfigure{in der Klasse}{45,4 Milliarden Stellungen,\\ 27,6 Milliarden mit Hilfsmatt}
\end{minipage}\hfill
\begin{minipage}[t]{80mm}\vspace{0pt}\begin{tikzpicture}\begin{axis}[width=72mm,height=42mm,ybar,bar width=1.6pt,xmin=-1,xmax=18,ymin=0,ymax=9.596,axis y line*=left,axis x line*=bottom,xlabel={\scriptsize Halbzüge bis Matt},ylabel={\scriptsize Stellungen (Mrd.)},tick label style={font=\tiny},label style={font=\scriptsize},enlarge x limits=false,ylabel shift=-3pt,xlabel shift=-3pt]\addplot[fill=ink!55,draw=none] coordinates {(0,0.01097) (1,0.07503) (2,0.2911) (3,1.525) (4,4.381) (5,7.274) (6,8.724) (7,4.049) (8,1.159) (9,0.1021) (10,0.005693) (11,0.0001405) (12,0.001757) (13,4.836e-06) (14,9.354e-05) (15,2.875e-06) (16,6.92e-07) (17,1.03e-07)};\draw[densely dotted] (axis cs:15,0) -- (axis cs:15,9.596);\end{axis}\begin{axis}[width=72mm,height=42mm,axis y line*=right,axis x line=none,ymode=log,xmin=-1,xmax=18,ylabel={\scriptsize eindeutig (log)},tick label style={font=\tiny},label style={font=\scriptsize},ylabel shift=-3pt]\addplot[mark=*,mark size=0.7pt,draw=black,thin] coordinates {(0,10972882) (1,61534624) (2,147119471) (3,461979773) (4,952332935) (5,510981632) (6,144289478) (7,9321474) (8,263730) (9,4455) (10,20610) (11,43) (12,1457) (13,6) (14,4) (15,2)};\end{axis}\end{tikzpicture}\end{minipage}
\vspace{3mm}

\noindent\begin{minipage}[t]{60mm}\vspace{0pt}\sbox\hmboard{\chessboard[setfen=KN6/r7/8/8/8/8/3P2R1/7k w - - 0 1]}\usebox\hmboard\\[1mm]
\makebox[\wd\hmboard]{\sffamily\hspace{3mm}\textbf{h\#7.5}\hfill(4+2)}\end{minipage}\hfill
\begin{minipage}[t]{52mm}\vspace{0pt}
\textcolor{ink}{\sffamily\textsc{Lösung}}\par\smallskip\small
\begin{tabular}{@{}r@{\enspace}l@{\enspace}l@{}}
\textbf{1.} & \ldots & \symking{}xa7\\
\textbf{2.} & \symking{}xg2 & d4\\
\textbf{3.} & \symking{}f3 & d5\\
\textbf{4.} & \symking{}e4 & d6\\
\textbf{5.} & \symking{}d5 & \symknight{}c6\\
\textbf{6.} & \symking{}xc6 & d7\\
\textbf{7.} & \symking{}c7 & \symking{}a6\\
\textbf{8.} & \symking{}b8 & d8=\symqueen{}\#\\
\end{tabular}\par\vspace{3mm}
\qrcode[height=12mm]{https://helpman.komtera.lt/?fen=KN6/r7/8/8/8/8/3P2R1/7k&moves=7.5}\\{\sffamily\tiny\color{ink}scannen zum Lösen}
\end{minipage}
\vspace{2mm}

{\raggedright\noindent\textcolor{ink}{\sffamily\textsc{Themen}}\enspace \themetag{reines Matt} \themetag{Mustermatt} \themetag{Idealmatt} \themetag{Spiegelmatt} \themetag{Umwandlung} \themetag{Excelsior} \themetag{ein Stein zieht} \themetag{ohne Schach} \themetag{Umnov}\par\smallskip{\footnotesize\textit{Hinweis:} die Partei am Zug steht im Diagramm im Schach; die Lösung beginnt mit einem Schlag.}\par}

\clearpage
\phantomsection\label{hm:KRPPvkq}
\begin{tcolorbox}[enhanced,colback=ink,colframe=ink,arc=0pt,boxrule=0pt,left=4mm,right=4mm,top=2.5mm,bottom=2.5mm]
\color{white}\sffamily{\LARGE\bfseries KRPPvkq}\hfill{\Large h\#7.5}\\[0.5mm]
{\small 6 Steine\quad Weiß am Zug}\hfill{\small Nr.~478\qquad\href{https://helpman.komtera.lt/?fen=8/8/6R1/7k/8/6P1/q5P1/K7&moves=7.5}{\color{white}\bfseries lösen\,\ensuremath{\nearrow}}}
\end{tcolorbox}
\vspace{2mm}
\noindent\begin{minipage}[t]{36mm}\vspace{0pt}\sffamily\small\raggedright
\hmfigure{längste eindeutige}{{\Large h\#7.5}}\hmfigure{eindeutig bei dieser Länge}{{\Large 1} Stellung}\hmfigure{längstes Matt}{{\Large h\#8.5}}\hmfigure{in der Klasse}{34,1 Milliarden Stellungen,\\ 9,74 Milliarden mit Hilfsmatt}
\end{minipage}\hfill
\begin{minipage}[t]{80mm}\vspace{0pt}\begin{tikzpicture}\begin{axis}[width=72mm,height=42mm,ybar,bar width=1.6pt,xmin=-1,xmax=18,ymin=0,ymax=2.341,axis y line*=left,axis x line*=bottom,xlabel={\scriptsize Halbzüge bis Matt},ylabel={\scriptsize Stellungen (Mrd.)},tick label style={font=\tiny},label style={font=\scriptsize},enlarge x limits=false,ylabel shift=-3pt,xlabel shift=-3pt]\addplot[fill=ink!55,draw=none] coordinates {(0,0.001988) (1,0.01758) (2,0.07455) (3,0.4678) (4,1.056) (5,2.128) (6,2.111) (7,1.855) (8,1.127) (9,0.6209) (10,0.2169) (11,0.05217) (12,0.007943) (13,0.0004579) (14,3.091e-05) (15,7.191e-06) (16,2.71e-07) (17,4e-08)};\draw[densely dotted] (axis cs:15,0) -- (axis cs:15,2.341);\end{axis}\begin{axis}[width=72mm,height=42mm,axis y line*=right,axis x line=none,ymode=log,xmin=-1,xmax=18,ylabel={\scriptsize eindeutig (log)},tick label style={font=\tiny},label style={font=\scriptsize},ylabel shift=-3pt]\addplot[mark=*,mark size=0.7pt,draw=black,thin] coordinates {(0,1988228) (1,12150255) (2,25276166) (3,104006654) (4,125302758) (5,107845480) (6,43013814) (7,19324598) (8,4845861) (9,1020622) (10,110228) (11,4683) (12,125) (13,29) (14,1) (15,1)};\end{axis}\end{tikzpicture}\end{minipage}
\vspace{3mm}

\noindent\begin{minipage}[t]{60mm}\vspace{0pt}\sbox\hmboard{\chessboard[setfen=8/8/6R1/7k/8/6P1/q5P1/K7 w - - 0 1]}\usebox\hmboard\\[1mm]
\makebox[\wd\hmboard]{\sffamily\hspace{3mm}\textbf{h\#7.5}\hfill(4+2)}\end{minipage}\hfill
\begin{minipage}[t]{52mm}\vspace{0pt}
\textcolor{ink}{\sffamily\textsc{Lösung}}\par\smallskip\small
\begin{tabular}{@{}r@{\enspace}l@{\enspace}l@{}}
\textbf{1.} & \ldots & \symking{}xa2\\
\textbf{2.} & \symking{}xg6 & \symking{}a3\\
\textbf{3.} & \symking{}f5 & g4+\\
\textbf{4.} & \symking{}e4 & g5\\
\textbf{5.} & \symking{}d3 & g6\\
\textbf{6.} & \symking{}c2 & g7\\
\textbf{7.} & \symking{}b1 & g8=\symqueen{}\\
\textbf{8.} & \symking{}a1 & \symqueen{}a2\#\\
\end{tabular}\par\vspace{3mm}
\qrcode[height=12mm]{https://helpman.komtera.lt/?fen=8/8/6R1/7k/8/6P1/q5P1/K7&moves=7.5}\\{\sffamily\tiny\color{ink}scannen zum Lösen}
\end{minipage}
\vspace{2mm}

{\raggedright\noindent\textcolor{ink}{\sffamily\textsc{Themen}}\enspace \themetag{Umwandlung} \themetag{ein Stein zieht}\par\smallskip{\footnotesize\textit{Hinweis:} die Partei am Zug steht im Diagramm im Schach; die Lösung beginnt mit einem Schlag.}\par}

\clearpage
\phantomsection\label{hm:KBBBBvk}
\begin{tcolorbox}[enhanced,colback=ink,colframe=ink,arc=0pt,boxrule=0pt,left=4mm,right=4mm,top=2.5mm,bottom=2.5mm]
\color{white}\sffamily{\LARGE\bfseries KBBBBvk}\hfill{\Large h\#7}\\[0.5mm]
{\small 6 Steine\quad Schwarz am Zug}\hfill{\small Nr.~479\qquad\href{https://helpman.komtera.lt/?fen=5B2/5B1k/7B/8/8/8/7B/K7&moves=7}{\color{white}\bfseries lösen\,\ensuremath{\nearrow}}}
\end{tcolorbox}
\vspace{2mm}
\noindent\begin{minipage}[t]{36mm}\vspace{0pt}\sffamily\small\raggedright
\hmfigure{längste eindeutige}{{\Large h\#7}}\hmfigure{eindeutig bei dieser Länge}{{\Large 10} Stellungen}\hmfigure{längstes Matt}{{\Large h\#8}}\hmfigure{in der Klasse}{15,5 Milliarden Stellungen,\\ 341 Millionen mit Hilfsmatt}
\end{minipage}\hfill
\begin{minipage}[t]{80mm}\vspace{0pt}\begin{tikzpicture}\begin{axis}[width=72mm,height=42mm,ybar,bar width=1.6pt,xmin=-1,xmax=17,ymin=0,ymax=97.48,axis y line*=left,axis x line*=bottom,xlabel={\scriptsize Halbzüge bis Matt},ylabel={\scriptsize Stellungen (Mio.)},tick label style={font=\tiny},label style={font=\scriptsize},enlarge x limits=false,ylabel shift=-3pt,xlabel shift=-3pt]\addplot[fill=ink!55,draw=none] coordinates {(0,1.228) (1,2.896) (2,8.967) (3,23.3) (4,88.47) (5,62.57) (6,88.62) (7,23.55) (8,26.85) (9,6.254) (10,6.063) (11,1.357) (12,1.159) (13,0.1212) (14,0.08748) (15,0.000272) (16,1.3e-05)};\draw[densely dotted] (axis cs:14,0) -- (axis cs:14,97.48);\end{axis}\begin{axis}[width=72mm,height=42mm,axis y line*=right,axis x line=none,ymode=log,xmin=-1,xmax=17,ylabel={\scriptsize eindeutig (log)},tick label style={font=\tiny},label style={font=\scriptsize},ylabel shift=-3pt]\addplot[mark=*,mark size=0.7pt,draw=black,thin] coordinates {(0,1228121) (1,2165905) (2,6834897) (3,7221464) (4,14140046) (5,897342) (6,883437) (7,140927) (8,76824) (9,16698) (10,6448) (11,2740) (12,945) (13,8) (14,10)};\end{axis}\end{tikzpicture}\end{minipage}
\vspace{3mm}

\noindent\begin{minipage}[t]{60mm}\vspace{0pt}\sbox\hmboard{\chessboard[setfen=5B2/5B1k/7B/8/8/8/7B/K7 b - - 0 1]}\usebox\hmboard\\[1mm]
\makebox[\wd\hmboard]{\sffamily\hspace{3mm}\textbf{h\#7}\hfill(5+1)}\end{minipage}\hfill
\begin{minipage}[t]{52mm}\vspace{0pt}
\textcolor{ink}{\sffamily\textsc{Lösung}}\par\smallskip\small
\begin{tabular}{@{}r@{\enspace}l@{\enspace}l@{}}
\textbf{1.} & \symking{}h8 & \symking{}b2\\
\textbf{2.} & \symking{}h7 & \symking{}c3\\
\textbf{3.} & \symking{}h8 & \symking{}d4\\
\textbf{4.} & \symking{}h7 & \symking{}e5\\
\textbf{5.} & \symking{}h8 & \symking{}f6\\
\textbf{6.} & \symking{}h7 & \symbishop{}e5\\
\textbf{7.} & \symking{}h8 & \symking{}g6\#\\
\end{tabular}\par\vspace{3mm}
\qrcode[height=12mm]{https://helpman.komtera.lt/?fen=5B2/5B1k/7B/8/8/8/7B/K7&moves=7}\\{\sffamily\tiny\color{ink}scannen zum Lösen}
\end{minipage}
\vspace{2mm}

{\raggedright\noindent\textcolor{ink}{\sffamily\textsc{Themen}}\enspace \themetag{Satzspiel} \themetag{Spiegelmatt} \themetag{Switchback} \themetag{ein Stein zieht} \themetag{Pendel} \themetag{ohne Schlag} \themetag{ohne Schach}\par}

\clearpage
\phantomsection\label{hm:KBBBNvk}
\begin{tcolorbox}[enhanced,colback=ink,colframe=ink,arc=0pt,boxrule=0pt,left=4mm,right=4mm,top=2.5mm,bottom=2.5mm]
\color{white}\sffamily{\LARGE\bfseries KBBBNvk}\hfill{\Large h\#7}\\[0.5mm]
{\small 6 Steine\quad Schwarz am Zug}\hfill{\small Nr.~480\qquad\href{https://helpman.komtera.lt/?fen=5B1N/5B1k/8/8/8/8/7B/K7&moves=7}{\color{white}\bfseries lösen\,\ensuremath{\nearrow}}}
\end{tcolorbox}
\vspace{2mm}
\noindent\begin{minipage}[t]{36mm}\vspace{0pt}\sffamily\small\raggedright
\hmfigure{längste eindeutige}{{\Large h\#7}}\hmfigure{eindeutig bei dieser Länge}{{\Large 3} Stellungen}\hmfigure{längstes Matt}{{\Large h\#8}}\hmfigure{in der Klasse}{15,5 Milliarden Stellungen,\\ 1,6 Milliarden mit Hilfsmatt}
\end{minipage}\hfill
\begin{minipage}[t]{80mm}\vspace{0pt}\begin{tikzpicture}\begin{axis}[width=72mm,height=42mm,ybar,bar width=1.6pt,xmin=-1,xmax=17,ymin=0,ymax=526.9,axis y line*=left,axis x line*=bottom,xlabel={\scriptsize Halbzüge bis Matt},ylabel={\scriptsize Stellungen (Mio.)},tick label style={font=\tiny},label style={font=\scriptsize},enlarge x limits=false,ylabel shift=-3pt,xlabel shift=-3pt]\addplot[fill=ink!55,draw=none] coordinates {(0,2.114) (1,6.208) (2,18.11) (3,61.83) (4,252.1) (5,310) (6,479) (7,137.7) (8,160.8) (9,65.43) (10,76.46) (11,18.41) (12,13.11) (13,2.477) (14,1) (15,0.1407) (16,0.00021)};\draw[densely dotted] (axis cs:14,0) -- (axis cs:14,526.9);\end{axis}\begin{axis}[width=72mm,height=42mm,axis y line*=right,axis x line=none,ymode=log,xmin=-1,xmax=17,ylabel={\scriptsize eindeutig (log)},tick label style={font=\tiny},label style={font=\scriptsize},ylabel shift=-3pt]\addplot[mark=*,mark size=0.7pt,draw=black,thin] coordinates {(0,2113750) (1,5137598) (2,14995180) (3,30481300) (4,73882842) (5,13839598) (6,9817451) (7,990662) (8,911161) (9,68077) (10,34221) (11,3327) (12,698) (13,70) (14,3)};\end{axis}\end{tikzpicture}\end{minipage}
\vspace{3mm}

\noindent\begin{minipage}[t]{60mm}\vspace{0pt}\sbox\hmboard{\chessboard[setfen=5B1N/5B1k/8/8/8/8/7B/K7 b - - 0 1]}\usebox\hmboard\\[1mm]
\makebox[\wd\hmboard]{\sffamily\hspace{3mm}\textbf{h\#7}\hfill(5+1)}\end{minipage}\hfill
\begin{minipage}[t]{52mm}\vspace{0pt}
\textcolor{ink}{\sffamily\textsc{Lösung}}\par\smallskip\small
\begin{tabular}{@{}r@{\enspace}l@{\enspace}l@{}}
\textbf{1.} & \symking{}xh8 & \symking{}b2\\
\textbf{2.} & \symking{}h7 & \symking{}c3\\
\textbf{3.} & \symking{}h8 & \symking{}d4\\
\textbf{4.} & \symking{}h7 & \symking{}e5\\
\textbf{5.} & \symking{}h8 & \symking{}f6\\
\textbf{6.} & \symking{}h7 & \symbishop{}e5\\
\textbf{7.} & \symking{}h8 & \symking{}g6\#\\
\end{tabular}\par\vspace{3mm}
\qrcode[height=12mm]{https://helpman.komtera.lt/?fen=5B1N/5B1k/8/8/8/8/7B/K7&moves=7}\\{\sffamily\tiny\color{ink}scannen zum Lösen}
\end{minipage}
\vspace{2mm}

{\raggedright\noindent\textcolor{ink}{\sffamily\textsc{Themen}}\enspace \themetag{Spiegelmatt} \themetag{Switchback} \themetag{ein Stein zieht} \themetag{Kniest} \themetag{Pendel} \themetag{ohne Schach}\par\smallskip{\footnotesize\textit{Hinweis:} die Lösung beginnt mit einem Schlag.}\par}

\clearpage
\phantomsection\label{hm:KBBBvkb}
\begin{tcolorbox}[enhanced,colback=ink,colframe=ink,arc=0pt,boxrule=0pt,left=4mm,right=4mm,top=2.5mm,bottom=2.5mm]
\color{white}\sffamily{\LARGE\bfseries KBBBvkb}\hfill{\Large h\#7}\\[0.5mm]
{\small 6 Steine\quad Schwarz am Zug}\hfill{\small Nr.~481\qquad\href{https://helpman.komtera.lt/?fen=8/8/6b1/8/8/8/1B5k/K1B1B3&moves=7}{\color{white}\bfseries lösen\,\ensuremath{\nearrow}}}
\end{tcolorbox}
\vspace{2mm}
\noindent\begin{minipage}[t]{36mm}\vspace{0pt}\sffamily\small\raggedright
\hmfigure{längste eindeutige}{{\Large h\#7}}\hmfigure{eindeutig bei dieser Länge}{{\Large 948} Stellungen}\hmfigure{längstes Matt}{{\Large h\#11}}\hmfigure{in der Klasse}{15,5 Milliarden Stellungen,\\ 1,35 Milliarden mit Hilfsmatt}
\end{minipage}\hfill
\begin{minipage}[t]{80mm}\vspace{0pt}\begin{tikzpicture}\begin{axis}[width=72mm,height=42mm,ybar,bar width=1.6pt,xmin=-1,xmax=23,ymin=0,ymax=359.4,axis y line*=left,axis x line*=bottom,xlabel={\scriptsize Halbzüge bis Matt},ylabel={\scriptsize Stellungen (Mio.)},tick label style={font=\tiny},label style={font=\scriptsize},enlarge x limits=false,ylabel shift=-3pt,xlabel shift=-3pt]\addplot[fill=ink!55,draw=none] coordinates {(0,0.9065) (1,3.061) (2,14.87) (3,67.01) (4,241.8) (5,271.1) (6,326.7) (7,122.7) (8,84.41) (9,36.96) (10,39.02) (11,30.23) (12,41.07) (13,20.24) (14,18.7) (15,11.66) (16,7.366) (17,5.877) (18,2.314) (19,2.371) (20,0.03784) (22,1e-06)};\draw[densely dotted] (axis cs:14,0) -- (axis cs:14,359.4);\end{axis}\begin{axis}[width=72mm,height=42mm,axis y line*=right,axis x line=none,ymode=log,xmin=-1,xmax=23,ylabel={\scriptsize eindeutig (log)},tick label style={font=\tiny},label style={font=\scriptsize},ylabel shift=-3pt]\addplot[mark=*,mark size=0.7pt,draw=black,thin] coordinates {(0,906512) (1,2673656) (2,10313490) (3,24104474) (4,42665414) (5,5297474) (6,2476277) (7,547517) (8,270014) (9,158627) (10,155042) (11,42433) (12,22548) (13,2695) (14,948)};\end{axis}\end{tikzpicture}\end{minipage}
\vspace{3mm}

\noindent\begin{minipage}[t]{60mm}\vspace{0pt}\sbox\hmboard{\chessboard[setfen=8/8/6b1/8/8/8/1B5k/K1B1B3 b - - 0 1]}\usebox\hmboard\\[1mm]
\makebox[\wd\hmboard]{\sffamily\hspace{3mm}\textbf{h\#7}\hfill(4+2)}\end{minipage}\hfill
\begin{minipage}[t]{52mm}\vspace{0pt}
\textcolor{ink}{\sffamily\textsc{Lösung}}\par\smallskip\small
\begin{tabular}{@{}r@{\enspace}l@{\enspace}l@{}}
\textbf{1.} & \symking{}h3 & \symking{}a2\\
\textbf{2.} & \symking{}g4 & \symking{}b3\\
\textbf{3.} & \symking{}f5 & \symking{}c3\\
\textbf{4.} & \symking{}f6 & \symking{}d4\\
\textbf{5.} & \symking{}g7 & \symking{}e5\\
\textbf{6.} & \symking{}h8 & \symking{}f6\\
\textbf{7.} & \symbishop{}h7 & \symking{}f7\#\\
\end{tabular}\par\vspace{3mm}
\qrcode[height=12mm]{https://helpman.komtera.lt/?fen=8/8/6b1/8/8/8/1B5k/K1B1B3&moves=7}\\{\sffamily\tiny\color{ink}scannen zum Lösen}
\end{minipage}
\vspace{2mm}

{\raggedright\noindent\textcolor{ink}{\sffamily\textsc{Themen}}\enspace \themetag{Selbstblock} \themetag{ein Stein zieht} \themetag{ohne Schlag} \themetag{ohne Schach}\par}

\clearpage
\phantomsection\label{hm:KBBNPvk}
\begin{tcolorbox}[enhanced,colback=ink,colframe=ink,arc=0pt,boxrule=0pt,left=4mm,right=4mm,top=2.5mm,bottom=2.5mm]
\color{white}\sffamily{\LARGE\bfseries KBBNPvk}\hfill{\Large h\#7}\\[0.5mm]
{\small 6 Steine\quad Schwarz am Zug}\hfill{\small Nr.~482\qquad\href{https://helpman.komtera.lt/?fen=8/8/8/7k/5B1N/4B2P/8/1K6&moves=7}{\color{white}\bfseries lösen\,\ensuremath{\nearrow}}}
\end{tcolorbox}
\vspace{2mm}
\noindent\begin{minipage}[t]{36mm}\vspace{0pt}\sffamily\small\raggedright
\hmfigure{längste eindeutige}{{\Large h\#7}}\hmfigure{eindeutig bei dieser Länge}{{\Large 65} Stellungen}\hmfigure{längstes Matt}{{\Large h\#9}}\hmfigure{in der Klasse}{45,4 Milliarden Stellungen,\\ 15,2 Milliarden mit Hilfsmatt}
\end{minipage}\hfill
\begin{minipage}[t]{80mm}\vspace{0pt}\begin{tikzpicture}\begin{axis}[width=72mm,height=42mm,ybar,bar width=1.6pt,xmin=-1,xmax=19,ymin=0,ymax=4.348,axis y line*=left,axis x line*=bottom,xlabel={\scriptsize Halbzüge bis Matt},ylabel={\scriptsize Stellungen (Mrd.)},tick label style={font=\tiny},label style={font=\scriptsize},enlarge x limits=false,ylabel shift=-3pt,xlabel shift=-3pt]\addplot[fill=ink!55,draw=none] coordinates {(0,0.005389) (1,0.02511) (2,0.06684) (3,0.3599) (4,1.283) (5,2.343) (6,3.953) (7,2.428) (8,2.988) (9,0.9078) (10,0.7192) (11,0.07497) (12,0.0309) (13,0.0005947) (14,0.0005077) (15,2.92e-07) (16,4.774e-06) (17,4e-09) (18,2.3e-08)};\draw[densely dotted] (axis cs:14,0) -- (axis cs:14,4.348);\end{axis}\begin{axis}[width=72mm,height=42mm,axis y line*=right,axis x line=none,ymode=log,xmin=-1,xmax=19,ylabel={\scriptsize eindeutig (log)},tick label style={font=\tiny},label style={font=\scriptsize},ylabel shift=-3pt]\addplot[mark=*,mark size=0.7pt,draw=black,thin] coordinates {(0,5388800) (1,17631619) (2,46857749) (3,194808258) (4,403722801) (5,208396535) (6,194694508) (7,35865983) (8,21365099) (9,1432285) (10,512000) (11,15273) (12,9717) (14,65)};\end{axis}\end{tikzpicture}\end{minipage}
\vspace{3mm}

\noindent\begin{minipage}[t]{60mm}\vspace{0pt}\sbox\hmboard{\chessboard[setfen=8/8/8/7k/5B1N/4B2P/8/1K6 b - - 0 1]}\usebox\hmboard\\[1mm]
\makebox[\wd\hmboard]{\sffamily\hspace{3mm}\textbf{h\#7}\hfill(5+1)}\end{minipage}\hfill
\begin{minipage}[t]{52mm}\vspace{0pt}
\textcolor{ink}{\sffamily\textsc{Lösung}}\par\smallskip\small
\begin{tabular}{@{}r@{\enspace}l@{\enspace}l@{}}
\textbf{1.} & \symking{}xh4 & \symbishop{}d2\\
\textbf{2.} & \symking{}h5 & h4\\
\textbf{3.} & \symking{}g4 & h5\\
\textbf{4.} & \symking{}f3 & h6\\
\textbf{5.} & \symking{}f2 & h7\\
\textbf{6.} & \symking{}e2 & h8=\symqueen{}\\
\textbf{7.} & \symking{}d1 & \symqueen{}h5\#\\
\end{tabular}\par\vspace{3mm}
\qrcode[height=12mm]{https://helpman.komtera.lt/?fen=8/8/8/7k/5B1N/4B2P/8/1K6&moves=7}\\{\sffamily\tiny\color{ink}scannen zum Lösen}
\end{minipage}
\vspace{2mm}

{\raggedright\noindent\textcolor{ink}{\sffamily\textsc{Themen}}\enspace \themetag{Umwandlung} \themetag{Switchback} \themetag{Rundlauf} \themetag{ein Stein zieht} \themetag{ohne Schach} \themetag{Umnov}\par\smallskip{\footnotesize\textit{Hinweis:} die Lösung beginnt mit einem Schlag.}\par}

\clearpage
\phantomsection\label{hm:KBBNvkn}
\begin{tcolorbox}[enhanced,colback=ink,colframe=ink,arc=0pt,boxrule=0pt,left=4mm,right=4mm,top=2.5mm,bottom=2.5mm]
\color{white}\sffamily{\LARGE\bfseries KBBNvkn}\hfill{\Large h\#7}\\[0.5mm]
{\small 6 Steine\quad Schwarz am Zug}\hfill{\small Nr.~483\qquad\href{https://helpman.komtera.lt/?fen=8/6Bk/5N2/7B/3n4/8/8/K7&moves=7}{\color{white}\bfseries lösen\,\ensuremath{\nearrow}}}
\end{tcolorbox}
\vspace{2mm}
\noindent\begin{minipage}[t]{36mm}\vspace{0pt}\sffamily\small\raggedright
\hmfigure{längste eindeutige}{{\Large h\#7}}\hmfigure{eindeutig bei dieser Länge}{{\Large 6} Stellungen}\hmfigure{längstes Matt}{{\Large h\#12}}\hmfigure{in der Klasse}{15,5 Milliarden Stellungen,\\ 4,87 Milliarden mit Hilfsmatt}
\end{minipage}\hfill
\begin{minipage}[t]{80mm}\vspace{0pt}\begin{tikzpicture}\begin{axis}[width=72mm,height=42mm,ybar,bar width=1.6pt,xmin=-1,xmax=25,ymin=0,ymax=1.402,axis y line*=left,axis x line*=bottom,xlabel={\scriptsize Halbzüge bis Matt},ylabel={\scriptsize Stellungen (Mrd.)},tick label style={font=\tiny},label style={font=\scriptsize},enlarge x limits=false,ylabel shift=-3pt,xlabel shift=-3pt]\addplot[fill=ink!55,draw=none] coordinates {(0,0.0009442) (1,0.003469) (2,0.01333) (3,0.07667) (4,0.2733) (5,0.6177) (6,1.275) (7,1.006) (8,1.063) (9,0.3748) (10,0.146) (11,0.02214) (12,0.00199) (13,0.0001302) (14,8.841e-05) (15,1e-09) (16,1.075e-06) (18,1.315e-06) (20,4e-08) (22,1.289e-06) (24,5.3e-08)};\draw[densely dotted] (axis cs:14,0) -- (axis cs:14,1.402);\end{axis}\begin{axis}[width=72mm,height=42mm,axis y line*=right,axis x line=none,ymode=log,xmin=-1,xmax=25,ylabel={\scriptsize eindeutig (log)},tick label style={font=\tiny},label style={font=\scriptsize},ylabel shift=-3pt]\addplot[mark=*,mark size=0.7pt,draw=black,thin] coordinates {(0,944236) (1,3214929) (2,9170828) (3,35290910) (4,72393225) (5,37842291) (6,34383238) (7,5420148) (8,1587242) (9,95271) (10,6610) (11,260) (12,200) (14,6)};\end{axis}\end{tikzpicture}\end{minipage}
\vspace{3mm}

\noindent\begin{minipage}[t]{60mm}\vspace{0pt}\sbox\hmboard{\chessboard[setfen=8/6Bk/5N2/7B/3n4/8/8/K7 b - - 0 1]}\usebox\hmboard\\[1mm]
\makebox[\wd\hmboard]{\sffamily\hspace{3mm}\textbf{h\#7}\hfill(4+2)}\end{minipage}\hfill
\begin{minipage}[t]{52mm}\vspace{0pt}
\textcolor{ink}{\sffamily\textsc{Lösung}}\par\smallskip\small
\begin{tabular}{@{}r@{\enspace}l@{\enspace}l@{}}
\textbf{1.} & \symking{}xg7 & \symking{}b2\\
\textbf{2.} & \symking{}f8 & \symking{}c3\\
\textbf{3.} & \symking{}e7 & \symking{}xd4\\
\textbf{4.} & \symking{}d8 & \symking{}c5\\
\textbf{5.} & \symking{}c8 & \symking{}b6\\
\textbf{6.} & \symking{}b8 & \symknight{}d7+\\
\textbf{7.} & \symking{}a8 & \symbishop{}f3\#\\
\end{tabular}\par\vspace{3mm}
\qrcode[height=12mm]{https://helpman.komtera.lt/?fen=8/6Bk/5N2/7B/3n4/8/8/K7&moves=7}\\{\sffamily\tiny\color{ink}scannen zum Lösen}
\end{minipage}
\vspace{2mm}

{\raggedright\noindent\textcolor{ink}{\sffamily\textsc{Themen}}\enspace \themetag{Spiegelmatt} \themetag{ein Stein zieht}\par\smallskip{\footnotesize\textit{Hinweis:} die Partei am Zug steht im Diagramm im Schach; die Lösung beginnt mit einem Schlag.}\par}

\clearpage
\phantomsection\label{hm:KBBNvkp}
\begin{tcolorbox}[enhanced,colback=ink,colframe=ink,arc=0pt,boxrule=0pt,left=4mm,right=4mm,top=2.5mm,bottom=2.5mm]
\color{white}\sffamily{\LARGE\bfseries KBBNvkp}\hfill{\Large h\#7}\\[0.5mm]
{\small 6 Steine\quad Schwarz am Zug}\hfill{\small Nr.~484\qquad\href{https://helpman.komtera.lt/?fen=7B/8/8/8/8/4B3/2p5/K1kN4&moves=7}{\color{white}\bfseries lösen\,\ensuremath{\nearrow}}}
\end{tcolorbox}
\vspace{2mm}
\noindent\begin{minipage}[t]{36mm}\vspace{0pt}\sffamily\small\raggedright
\hmfigure{längste eindeutige}{{\Large h\#7}}\hmfigure{eindeutig bei dieser Länge}{{\Large 17} Stellungen}\hmfigure{längstes Matt}{{\Large h\#15}}\hmfigure{in der Klasse}{45,4 Milliarden Stellungen,\\ 15,1 Milliarden mit Hilfsmatt}
\end{minipage}\hfill
\begin{minipage}[t]{80mm}\vspace{0pt}\begin{tikzpicture}\begin{axis}[width=72mm,height=42mm,ybar,bar width=1.6pt,xmin=-1,xmax=31,ymin=0,ymax=3.804,axis y line*=left,axis x line*=bottom,xlabel={\scriptsize Halbzüge bis Matt},ylabel={\scriptsize Stellungen (Mrd.)},tick label style={font=\tiny},label style={font=\scriptsize},enlarge x limits=false,ylabel shift=-3pt,xlabel shift=-3pt]\addplot[fill=ink!55,draw=none] coordinates {(0,0.003991) (1,0.01447) (2,0.04679) (3,0.2444) (4,0.8124) (5,1.702) (6,3.439) (7,2.782) (8,3.458) (9,1.385) (10,0.9904) (11,0.192) (12,0.06257) (13,0.006777) (14,0.001104) (15,5.17e-05) (16,3.45e-05) (17,6e-09) (18,3.763e-05) (19,2.8e-07) (20,2.872e-05) (21,1.91e-07) (22,2.983e-05) (23,3.34e-07) (24,2.81e-05) (25,3.73e-07) (26,9.026e-06) (27,1.25e-07) (28,3.201e-06) (29,1.1e-08) (30,3.25e-07)};\draw[densely dotted] (axis cs:14,0) -- (axis cs:14,3.804);\end{axis}\begin{axis}[width=72mm,height=42mm,axis y line*=right,axis x line=none,ymode=log,xmin=-1,xmax=31,ylabel={\scriptsize eindeutig (log)},tick label style={font=\tiny},label style={font=\scriptsize},ylabel shift=-3pt]\addplot[mark=*,mark size=0.7pt,draw=black,thin] coordinates {(0,3990926) (1,13264612) (2,39938251) (3,143923879) (4,246347628) (5,124373880) (6,131124650) (7,22635273) (8,14005283) (9,922209) (10,262580) (11,14752) (12,4363) (13,197) (14,17)};\end{axis}\end{tikzpicture}\end{minipage}
\vspace{3mm}

\noindent\begin{minipage}[t]{60mm}\vspace{0pt}\sbox\hmboard{\chessboard[setfen=7B/8/8/8/8/4B3/2p5/K1kN4 b - - 0 1]}\usebox\hmboard\\[1mm]
\makebox[\wd\hmboard]{\sffamily\hspace{3mm}\textbf{h\#7}\hfill(4+2)}\end{minipage}\hfill
\begin{minipage}[t]{52mm}\vspace{0pt}
\textcolor{ink}{\sffamily\textsc{Lösung}}\par\smallskip\small
\begin{tabular}{@{}r@{\enspace}l@{\enspace}l@{}}
\textbf{1.} & \symking{}xd1 & \symbishop{}c5\\
\textbf{2.} & \symking{}d2 & \symking{}a2\\
\textbf{3.} & \symking{}c1 & \symking{}b3\\
\textbf{4.} & \symking{}b1 & \symbishop{}a1\\
\textbf{5.} & \symking{}xa1 & \symbishop{}a3\\
\textbf{6.} & c1=\symknight{}+ & \symking{}c2\\
\textbf{7.} & \symknight{}a2 & \symbishop{}b2\#\\
\end{tabular}\par\vspace{3mm}
\qrcode[height=12mm]{https://helpman.komtera.lt/?fen=7B/8/8/8/8/4B3/2p5/K1kN4&moves=7}\\{\sffamily\tiny\color{ink}scannen zum Lösen}
\end{minipage}
\vspace{2mm}

{\raggedright\noindent\textcolor{ink}{\sffamily\textsc{Themen}}\enspace \themetag{reines Matt} \themetag{Mustermatt} \themetag{Idealmatt} \themetag{Umwandlung} \themetag{Unterverwandlung} \themetag{Rundlauf} \themetag{Selbstblock} \themetag{Kniest} \themetag{Umnov}\par\smallskip{\footnotesize\textit{Hinweis:} die Partei am Zug steht im Diagramm im Schach; die Lösung beginnt mit einem Schlag.}\par}

\clearpage
\phantomsection\label{hm:KBBPPvk}
\begin{tcolorbox}[enhanced,colback=ink,colframe=ink,arc=0pt,boxrule=0pt,left=4mm,right=4mm,top=2.5mm,bottom=2.5mm]
\color{white}\sffamily{\LARGE\bfseries KBBPPvk}\hfill{\Large h\#7}\\[0.5mm]
{\small 6 Steine\quad Schwarz am Zug}\hfill{\small Nr.~485\qquad\href{https://helpman.komtera.lt/?fen=6B1/8/8/8/6B1/2k3P1/6P1/K7&moves=7}{\color{white}\bfseries lösen\,\ensuremath{\nearrow}}}
\end{tcolorbox}
\vspace{2mm}
\noindent\begin{minipage}[t]{36mm}\vspace{0pt}\sffamily\small\raggedright
\hmfigure{längste eindeutige}{{\Large h\#7}}\hmfigure{eindeutig bei dieser Länge}{{\Large 432} Stellungen}\hmfigure{längstes Matt}{{\Large h\#9}}\hmfigure{in der Klasse}{34,1 Milliarden Stellungen,\\ 5,8 Milliarden mit Hilfsmatt}
\end{minipage}\hfill
\begin{minipage}[t]{80mm}\vspace{0pt}\begin{tikzpicture}\begin{axis}[width=72mm,height=42mm,ybar,bar width=1.6pt,xmin=-1,xmax=19,ymin=0,ymax=1.543,axis y line*=left,axis x line*=bottom,xlabel={\scriptsize Halbzüge bis Matt},ylabel={\scriptsize Stellungen (Mrd.)},tick label style={font=\tiny},label style={font=\scriptsize},enlarge x limits=false,ylabel shift=-3pt,xlabel shift=-3pt]\addplot[fill=ink!55,draw=none] coordinates {(0,0.001568) (1,0.01031) (2,0.02718) (3,0.1608) (4,0.5272) (5,0.9659) (6,1.403) (7,0.7531) (8,0.8609) (9,0.3531) (10,0.3672) (11,0.1559) (12,0.1753) (13,0.01949) (14,0.01529) (15,9.236e-05) (16,6.271e-05) (17,1.84e-07) (18,5.2e-08)};\draw[densely dotted] (axis cs:14,0) -- (axis cs:14,1.543);\end{axis}\begin{axis}[width=72mm,height=42mm,axis y line*=right,axis x line=none,ymode=log,xmin=-1,xmax=19,ylabel={\scriptsize eindeutig (log)},tick label style={font=\tiny},label style={font=\scriptsize},ylabel shift=-3pt]\addplot[mark=*,mark size=0.7pt,draw=black,thin] coordinates {(0,1568193) (1,5915704) (2,15849138) (3,82460445) (4,148759099) (5,78702266) (6,61090732) (7,19788351) (8,11783972) (9,4512548) (10,2936381) (11,920187) (12,450966) (13,730) (14,432)};\end{axis}\end{tikzpicture}\end{minipage}
\vspace{3mm}

\noindent\begin{minipage}[t]{60mm}\vspace{0pt}\sbox\hmboard{\chessboard[setfen=6B1/8/8/8/6B1/2k3P1/6P1/K7 b - - 0 1]}\usebox\hmboard\\[1mm]
\makebox[\wd\hmboard]{\sffamily\hspace{3mm}\textbf{h\#7}\hfill(5+1)}\end{minipage}\hfill
\begin{minipage}[t]{52mm}\vspace{0pt}
\textcolor{ink}{\sffamily\textsc{Lösung}}\par\smallskip\small
\begin{tabular}{@{}r@{\enspace}l@{\enspace}l@{}}
\textbf{1.} & \symking{}d4 & \symbishop{}h5\\
\textbf{2.} & \symking{}e5 & g4\\
\textbf{3.} & \symking{}f6 & g5+\\
\textbf{4.} & \symking{}g7 & g6\\
\textbf{5.} & \symking{}xg8 & g7\\
\textbf{6.} & \symking{}h7 & g8=\symqueen{}+\\
\textbf{7.} & \symking{}h6 & \symqueen{}g6\#\\
\end{tabular}\par\vspace{3mm}
\qrcode[height=12mm]{https://helpman.komtera.lt/?fen=6B1/8/8/8/6B1/2k3P1/6P1/K7&moves=7}\\{\sffamily\tiny\color{ink}scannen zum Lösen}
\end{minipage}
\vspace{2mm}

{\raggedright\noindent\textcolor{ink}{\sffamily\textsc{Themen}}\enspace \themetag{reines Matt} \themetag{Mustermatt} \themetag{Umwandlung} \themetag{Rundlauf} \themetag{ein Stein zieht} \themetag{Umnov}\par}

\clearpage
\phantomsection\label{hm:KBBPvkb}
\begin{tcolorbox}[enhanced,colback=ink,colframe=ink,arc=0pt,boxrule=0pt,left=4mm,right=4mm,top=2.5mm,bottom=2.5mm]
\color{white}\sffamily{\LARGE\bfseries KBBPvkb}\hfill{\Large h\#7}\\[0.5mm]
{\small 6 Steine\quad Schwarz am Zug}\hfill{\small Nr.~486\qquad\href{https://helpman.komtera.lt/?fen=1B5b/6k1/3B4/8/8/1P6/8/K7&moves=7}{\color{white}\bfseries lösen\,\ensuremath{\nearrow}}}
\end{tcolorbox}
\vspace{2mm}
\noindent\begin{minipage}[t]{36mm}\vspace{0pt}\sffamily\small\raggedright
\hmfigure{längste eindeutige}{{\Large h\#7}}\hmfigure{eindeutig bei dieser Länge}{{\Large 906} Stellungen}\hmfigure{längstes Matt}{{\Large h\#10}}\hmfigure{in der Klasse}{45,4 Milliarden Stellungen,\\ 14,6 Milliarden mit Hilfsmatt}
\end{minipage}\hfill
\begin{minipage}[t]{80mm}\vspace{0pt}\begin{tikzpicture}\begin{axis}[width=72mm,height=42mm,ybar,bar width=1.6pt,xmin=-1,xmax=21,ymin=0,ymax=3.314,axis y line*=left,axis x line*=bottom,xlabel={\scriptsize Halbzüge bis Matt},ylabel={\scriptsize Stellungen (Mrd.)},tick label style={font=\tiny},label style={font=\scriptsize},enlarge x limits=false,ylabel shift=-3pt,xlabel shift=-3pt]\addplot[fill=ink!55,draw=none] coordinates {(0,0.00234) (1,0.0125) (2,0.05501) (3,0.361) (4,1.19) (5,2.142) (6,3.012) (7,2.057) (8,1.847) (9,1.148) (10,1.04) (11,0.7849) (12,0.7141) (13,0.1562) (14,0.05921) (15,0.00154) (16,0.0001349) (17,1.07e-07) (18,7.833e-05) (20,7.504e-06)};\draw[densely dotted] (axis cs:14,0) -- (axis cs:14,3.314);\end{axis}\begin{axis}[width=72mm,height=42mm,axis y line*=right,axis x line=none,ymode=log,xmin=-1,xmax=21,ylabel={\scriptsize eindeutig (log)},tick label style={font=\tiny},label style={font=\scriptsize},ylabel shift=-3pt]\addplot[mark=*,mark size=0.7pt,draw=black,thin] coordinates {(0,2340353) (1,9459463) (2,34905118) (3,152614949) (4,246689692) (5,125422022) (6,79250743) (7,31746685) (8,17889849) (9,10742792) (10,5984127) (11,2643471) (12,973595) (13,19787) (14,906)};\end{axis}\end{tikzpicture}\end{minipage}
\vspace{3mm}

\noindent\begin{minipage}[t]{60mm}\vspace{0pt}\sbox\hmboard{\chessboard[setfen=1B5b/6k1/3B4/8/8/1P6/8/K7 b - - 0 1]}\usebox\hmboard\\[1mm]
\makebox[\wd\hmboard]{\sffamily\hspace{3mm}\textbf{h\#7}\hfill(4+2)}\end{minipage}\hfill
\begin{minipage}[t]{52mm}\vspace{0pt}
\textcolor{ink}{\sffamily\textsc{Lösung}}\par\smallskip\small
\begin{tabular}{@{}r@{\enspace}l@{\enspace}l@{}}
\textbf{1.} & \symking{}f7+ & \symking{}b1\\
\textbf{2.} & \symbishop{}d4 & b4\\
\textbf{3.} & \symbishop{}c5 & bxc5\\
\textbf{4.} & \symking{}e6 & c6\\
\textbf{5.} & \symking{}d5 & c7\\
\textbf{6.} & \symking{}c4 & c8=\symqueen{}+\\
\textbf{7.} & \symking{}b3 & \symqueen{}c2\#\\
\end{tabular}\par\vspace{3mm}
\qrcode[height=12mm]{https://helpman.komtera.lt/?fen=1B5b/6k1/3B4/8/8/1P6/8/K7&moves=7}\\{\sffamily\tiny\color{ink}scannen zum Lösen}
\end{minipage}
\vspace{2mm}

{\raggedright\noindent\textcolor{ink}{\sffamily\textsc{Themen}}\enspace \themetag{Umwandlung}\par}

\clearpage
\phantomsection\label{hm:KBBvkrb}
\begin{tcolorbox}[enhanced,colback=ink,colframe=ink,arc=0pt,boxrule=0pt,left=4mm,right=4mm,top=2.5mm,bottom=2.5mm]
\color{white}\sffamily{\LARGE\bfseries KBBvkrb}\hfill{\Large h\#7}\\[0.5mm]
{\small 6 Steine\quad Schwarz am Zug}\hfill{\small Nr.~487\qquad\href{https://helpman.komtera.lt/?fen=7k/5bB1/7B/8/8/8/1r6/K7&moves=7}{\color{white}\bfseries lösen\,\ensuremath{\nearrow}}}
\end{tcolorbox}
\vspace{2mm}
\noindent\begin{minipage}[t]{36mm}\vspace{0pt}\sffamily\small\raggedright
\hmfigure{längste eindeutige}{{\Large h\#7}}\hmfigure{eindeutig bei dieser Länge}{{\Large 3} Stellungen}\hmfigure{längstes Matt}{{\Large h\#10.5}}\hmfigure{in der Klasse}{15,5 Milliarden Stellungen,\\ 3,32 Milliarden mit Hilfsmatt}
\end{minipage}\hfill
\begin{minipage}[t]{80mm}\vspace{0pt}\begin{tikzpicture}\begin{axis}[width=72mm,height=42mm,ybar,bar width=1.6pt,xmin=-1,xmax=22,ymin=0,ymax=704.7,axis y line*=left,axis x line*=bottom,xlabel={\scriptsize Halbzüge bis Matt},ylabel={\scriptsize Stellungen (Mio.)},tick label style={font=\tiny},label style={font=\scriptsize},enlarge x limits=false,ylabel shift=-3pt,xlabel shift=-3pt]\addplot[fill=ink!55,draw=none] coordinates {(0,0.2101) (1,0.9463) (2,5.57) (3,36.84) (4,137) (5,353.7) (6,598) (7,640.6) (8,511.7) (9,418.2) (10,258.4) (11,239.1) (12,96.6) (13,26.05) (14,1.481) (15,0.1827) (16,0.001697) (17,0.02206) (19,0.005165) (21,2e-06)};\draw[densely dotted] (axis cs:14,0) -- (axis cs:14,704.7);\end{axis}\begin{axis}[width=72mm,height=42mm,axis y line*=right,axis x line=none,ymode=log,xmin=-1,xmax=22,ylabel={\scriptsize eindeutig (log)},tick label style={font=\tiny},label style={font=\scriptsize},ylabel shift=-3pt]\addplot[mark=*,mark size=0.7pt,draw=black,thin] coordinates {(0,210054) (1,933414) (2,3718400) (3,16878555) (4,32800045) (5,28228813) (6,14399123) (7,5943882) (8,2460758) (9,1638944) (10,784767) (11,318215) (12,35814) (13,1693) (14,3)};\end{axis}\end{tikzpicture}\end{minipage}
\vspace{3mm}

\noindent\begin{minipage}[t]{60mm}\vspace{0pt}\sbox\hmboard{\chessboard[setfen=7k/5bB1/7B/8/8/8/1r6/K7 b - - 0 1]}\usebox\hmboard\\[1mm]
\makebox[\wd\hmboard]{\sffamily\hspace{3mm}\textbf{h\#7}\hfill(3+3)}\end{minipage}\hfill
\begin{minipage}[t]{52mm}\vspace{0pt}
\textcolor{ink}{\sffamily\textsc{Lösung}}\par\smallskip\small
\begin{tabular}{@{}r@{\enspace}l@{\enspace}l@{}}
\textbf{1.} & \symking{}h7 & \symbishop{}e3\\
\textbf{2.} & \symrook{}h2 & \symking{}b1\\
\textbf{3.} & \symrook{}h3 & \symking{}c2\\
\textbf{4.} & \symking{}g6 & \symking{}d3\\
\textbf{5.} & \symking{}h5 & \symking{}e4\\
\textbf{6.} & \symking{}h4 & \symking{}f4\\
\textbf{7.} & \symbishop{}h5 & \symbishop{}f6\#\\
\end{tabular}\par\vspace{3mm}
\qrcode[height=12mm]{https://helpman.komtera.lt/?fen=7k/5bB1/7B/8/8/8/1r6/K7&moves=7}\\{\sffamily\tiny\color{ink}scannen zum Lösen}
\end{minipage}
\vspace{2mm}

{\raggedright\noindent\textcolor{ink}{\sffamily\textsc{Themen}}\enspace \themetag{Selbstblock} \themetag{ohne Schlag} \themetag{ohne Schach}\par\smallskip{\footnotesize\textit{Hinweis:} die Partei am Zug steht im Diagramm im Schach.}\par}

\clearpage
\phantomsection\label{hm:KBNNPvk}
\begin{tcolorbox}[enhanced,colback=ink,colframe=ink,arc=0pt,boxrule=0pt,left=4mm,right=4mm,top=2.5mm,bottom=2.5mm]
\color{white}\sffamily{\LARGE\bfseries KBNNPvk}\hfill{\Large h\#7}\\[0.5mm]
{\small 6 Steine\quad Schwarz am Zug}\hfill{\small Nr.~488\qquad\href{https://helpman.komtera.lt/?fen=8/K7/8/8/6N1/4N3/4P3/6kB&moves=7}{\color{white}\bfseries lösen\,\ensuremath{\nearrow}}}
\end{tcolorbox}
\vspace{2mm}
\noindent\begin{minipage}[t]{36mm}\vspace{0pt}\sffamily\small\raggedright
\hmfigure{längste eindeutige}{{\Large h\#7}}\hmfigure{eindeutig bei dieser Länge}{{\Large 2} Stellungen}\hmfigure{längstes Matt}{{\Large h\#8}}\hmfigure{in der Klasse}{45,4 Milliarden Stellungen,\\ 15,6 Milliarden mit Hilfsmatt}
\end{minipage}\hfill
\begin{minipage}[t]{80mm}\vspace{0pt}\begin{tikzpicture}\begin{axis}[width=72mm,height=42mm,ybar,bar width=1.6pt,xmin=-1,xmax=17,ymin=0,ymax=4.255,axis y line*=left,axis x line*=bottom,xlabel={\scriptsize Halbzüge bis Matt},ylabel={\scriptsize Stellungen (Mrd.)},tick label style={font=\tiny},label style={font=\scriptsize},enlarge x limits=false,ylabel shift=-3pt,xlabel shift=-3pt]\addplot[fill=ink!55,draw=none] coordinates {(0,0.002093) (1,0.01415) (2,0.03787) (3,0.2482) (4,0.8526) (5,2.016) (6,3.539) (7,2.993) (8,3.868) (9,1.172) (10,0.7356) (11,0.06266) (12,0.017) (13,0.0002245) (14,0.0002435) (15,1.48e-07) (16,6.5e-08)};\draw[densely dotted] (axis cs:14,0) -- (axis cs:14,4.255);\end{axis}\begin{axis}[width=72mm,height=42mm,axis y line*=right,axis x line=none,ymode=log,xmin=-1,xmax=17,ylabel={\scriptsize eindeutig (log)},tick label style={font=\tiny},label style={font=\scriptsize},ylabel shift=-3pt]\addplot[mark=*,mark size=0.7pt,draw=black,thin] coordinates {(0,2093129) (1,8360750) (2,23313265) (3,133656461) (4,290962171) (5,237528061) (6,233753038) (7,43598473) (8,22492030) (9,1001587) (10,266625) (11,4026) (12,3427) (14,2)};\end{axis}\end{tikzpicture}\end{minipage}
\vspace{3mm}

\noindent\begin{minipage}[t]{60mm}\vspace{0pt}\sbox\hmboard{\chessboard[setfen=8/K7/8/8/6N1/4N3/4P3/6kB b - - 0 1]}\usebox\hmboard\\[1mm]
\makebox[\wd\hmboard]{\sffamily\hspace{3mm}\textbf{h\#7}\hfill(5+1)}\end{minipage}\hfill
\begin{minipage}[t]{52mm}\vspace{0pt}
\textcolor{ink}{\sffamily\textsc{Lösung}}\par\smallskip\small
\begin{tabular}{@{}r@{\enspace}l@{\enspace}l@{}}
\textbf{1.} & \symking{}xh1 & \symknight{}f5\\
\textbf{2.} & \symking{}g2 & e4\\
\textbf{3.} & \symking{}f3 & e5\\
\textbf{4.} & \symking{}f4 & e6\\
\textbf{5.} & \symking{}g5 & e7\\
\textbf{6.} & \symking{}g6 & e8=\symqueen{}+\\
\textbf{7.} & \symking{}h7 & \symknight{}f6\#\\
\end{tabular}\par\vspace{3mm}
\qrcode[height=12mm]{https://helpman.komtera.lt/?fen=8/K7/8/8/6N1/4N3/4P3/6kB&moves=7}\\{\sffamily\tiny\color{ink}scannen zum Lösen}
\end{minipage}
\vspace{2mm}

{\raggedright\noindent\textcolor{ink}{\sffamily\textsc{Themen}}\enspace \themetag{Spiegelmatt} \themetag{Umwandlung} \themetag{Excelsior} \themetag{ein Stein zieht}\par\smallskip{\footnotesize\textit{Hinweis:} die Lösung beginnt mit einem Schlag.}\par}

\clearpage
\phantomsection\label{hm:KBNNvkp}
\begin{tcolorbox}[enhanced,colback=ink,colframe=ink,arc=0pt,boxrule=0pt,left=4mm,right=4mm,top=2.5mm,bottom=2.5mm]
\color{white}\sffamily{\LARGE\bfseries KBNNvkp}\hfill{\Large h\#7}\\[0.5mm]
{\small 6 Steine\quad Schwarz am Zug}\hfill{\small Nr.~489\qquad\href{https://helpman.komtera.lt/?fen=5N2/8/7k/7N/2p2B2/8/8/K7&moves=7}{\color{white}\bfseries lösen\,\ensuremath{\nearrow}}}
\end{tcolorbox}
\vspace{2mm}
\noindent\begin{minipage}[t]{36mm}\vspace{0pt}\sffamily\small\raggedright
\hmfigure{längste eindeutige}{{\Large h\#7}}\hmfigure{eindeutig bei dieser Länge}{{\Large 12} Stellungen}\hmfigure{längstes Matt}{{\Large h\#9}}\hmfigure{in der Klasse}{45,4 Milliarden Stellungen,\\ 15,5 Milliarden mit Hilfsmatt}
\end{minipage}\hfill
\begin{minipage}[t]{80mm}\vspace{0pt}\begin{tikzpicture}\begin{axis}[width=72mm,height=42mm,ybar,bar width=1.6pt,xmin=-1,xmax=19,ymin=0,ymax=4.939,axis y line*=left,axis x line*=bottom,xlabel={\scriptsize Halbzüge bis Matt},ylabel={\scriptsize Stellungen (Mrd.)},tick label style={font=\tiny},label style={font=\scriptsize},enlarge x limits=false,ylabel shift=-3pt,xlabel shift=-3pt]\addplot[fill=ink!55,draw=none] coordinates {(0,0.001438) (1,0.005554) (2,0.01756) (3,0.1034) (4,0.3613) (5,1.196) (6,2.828) (7,3.35) (8,4.49) (9,1.857) (10,1.076) (11,0.1788) (12,0.04167) (13,0.003661) (14,0.0003243) (15,5.75e-06) (16,3.65e-07) (18,2e-09)};\draw[densely dotted] (axis cs:14,0) -- (axis cs:14,4.939);\end{axis}\begin{axis}[width=72mm,height=42mm,axis y line*=right,axis x line=none,ymode=log,xmin=-1,xmax=19,ylabel={\scriptsize eindeutig (log)},tick label style={font=\tiny},label style={font=\scriptsize},ylabel shift=-3pt]\addplot[mark=*,mark size=0.7pt,draw=black,thin] coordinates {(0,1437826) (1,5048568) (2,14859482) (3,61417540) (4,139077687) (5,139444987) (6,165406672) (7,30801952) (8,16393221) (9,913223) (10,190112) (11,6040) (12,895) (14,12)};\end{axis}\end{tikzpicture}\end{minipage}
\vspace{3mm}

\noindent\begin{minipage}[t]{60mm}\vspace{0pt}\sbox\hmboard{\chessboard[setfen=5N2/8/7k/7N/2p2B2/8/8/K7 b - - 0 1]}\usebox\hmboard\\[1mm]
\makebox[\wd\hmboard]{\sffamily\hspace{3mm}\textbf{h\#7}\hfill(4+2)}\end{minipage}\hfill
\begin{minipage}[t]{52mm}\vspace{0pt}
\textcolor{ink}{\sffamily\textsc{Lösung}}\par\smallskip\small
\begin{tabular}{@{}r@{\enspace}l@{\enspace}l@{}}
\textbf{1.} & \symking{}xh5 & \symking{}b2\\
\textbf{2.} & \symking{}g4 & \symking{}c3\\
\textbf{3.} & \symking{}f5 & \symking{}d4\\
\textbf{4.} & \symking{}f6 & \symking{}d5\\
\textbf{5.} & \symking{}g7 & \symking{}e6\\
\textbf{6.} & \symking{}h8 & \symking{}f7\\
\textbf{7.} & c3 & \symbishop{}e5\#\\
\end{tabular}\par\vspace{3mm}
\qrcode[height=12mm]{https://helpman.komtera.lt/?fen=5N2/8/7k/7N/2p2B2/8/8/K7&moves=7}\\{\sffamily\tiny\color{ink}scannen zum Lösen}
\end{minipage}
\vspace{2mm}

{\raggedright\noindent\textcolor{ink}{\sffamily\textsc{Themen}}\enspace \themetag{Spiegelmatt} \themetag{ohne Schach}\par\smallskip{\footnotesize\textit{Hinweis:} die Partei am Zug steht im Diagramm im Schach; die Lösung beginnt mit einem Schlag.}\par}

\clearpage
\phantomsection\label{hm:KBNNvkq}
\begin{tcolorbox}[enhanced,colback=ink,colframe=ink,arc=0pt,boxrule=0pt,left=4mm,right=4mm,top=2.5mm,bottom=2.5mm]
\color{white}\sffamily{\LARGE\bfseries KBNNvkq}\hfill{\Large h\#7}\\[0.5mm]
{\small 6 Steine\quad Schwarz am Zug}\hfill{\small Nr.~490\qquad\href{https://helpman.komtera.lt/?fen=8/6NN/8/8/8/8/qB6/k1K5&moves=7}{\color{white}\bfseries lösen\,\ensuremath{\nearrow}}}
\end{tcolorbox}
\vspace{2mm}
\noindent\begin{minipage}[t]{36mm}\vspace{0pt}\sffamily\small\raggedright
\hmfigure{längste eindeutige}{{\Large h\#7}}\hmfigure{eindeutig bei dieser Länge}{{\Large 1} Stellung}\hmfigure{längstes Matt}{{\Large h\#8.5}}\hmfigure{in der Klasse}{15,5 Milliarden Stellungen,\\ 4,26 Milliarden mit Hilfsmatt}
\end{minipage}\hfill
\begin{minipage}[t]{80mm}\vspace{0pt}\begin{tikzpicture}\begin{axis}[width=72mm,height=42mm,ybar,bar width=1.6pt,xmin=-1,xmax=18,ymin=0,ymax=1.292,axis y line*=left,axis x line*=bottom,xlabel={\scriptsize Halbzüge bis Matt},ylabel={\scriptsize Stellungen (Mrd.)},tick label style={font=\tiny},label style={font=\scriptsize},enlarge x limits=false,ylabel shift=-3pt,xlabel shift=-3pt]\addplot[fill=ink!55,draw=none] coordinates {(0,0.0001637) (1,0.0006557) (2,0.002926) (3,0.02219) (4,0.08457) (5,0.3483) (6,0.7546) (7,1.174) (8,1.023) (9,0.6303) (10,0.171) (11,0.04698) (12,0.001578) (13,0.001139) (14,1.112e-05) (15,8.122e-05) (16,1.32e-07) (17,5e-09)};\draw[densely dotted] (axis cs:14,0) -- (axis cs:14,1.292);\end{axis}\begin{axis}[width=72mm,height=42mm,axis y line*=right,axis x line=none,ymode=log,xmin=-1,xmax=18,ylabel={\scriptsize eindeutig (log)},tick label style={font=\tiny},label style={font=\scriptsize},ylabel shift=-3pt]\addplot[mark=*,mark size=0.7pt,draw=black,thin] coordinates {(0,163662) (1,609176) (2,1785360) (3,8909758) (4,22677330) (5,39685248) (6,29939711) (7,15162257) (8,2340501) (9,497780) (10,11587) (11,3348) (12,54) (13,52) (14,1)};\end{axis}\end{tikzpicture}\end{minipage}
\vspace{3mm}

\noindent\begin{minipage}[t]{60mm}\vspace{0pt}\sbox\hmboard{\chessboard[setfen=8/6NN/8/8/8/8/qB6/k1K5 b - - 0 1]}\usebox\hmboard\\[1mm]
\makebox[\wd\hmboard]{\sffamily\hspace{3mm}\textbf{h\#7}\hfill(4+2)}\end{minipage}\hfill
\begin{minipage}[t]{52mm}\vspace{0pt}
\textcolor{ink}{\sffamily\textsc{Lösung}}\par\smallskip\small
\begin{tabular}{@{}r@{\enspace}l@{\enspace}l@{}}
\textbf{1.} & \symqueen{}xb2+ & \symking{}d1\\
\textbf{2.} & \symking{}b1 & \symking{}e1\\
\textbf{3.} & \symking{}c2 & \symknight{}g5\\
\textbf{4.} & \symking{}d3 & \symking{}d1\\
\textbf{5.} & \symking{}e3 & \symknight{}f5+\\
\textbf{6.} & \symking{}f2 & \symknight{}h3+\\
\textbf{7.} & \symking{}f1 & \symknight{}e3\#\\
\end{tabular}\par\vspace{3mm}
\qrcode[height=12mm]{https://helpman.komtera.lt/?fen=8/6NN/8/8/8/8/qB6/k1K5&moves=7}\\{\sffamily\tiny\color{ink}scannen zum Lösen}
\end{minipage}
\vspace{2mm}

{\raggedright\noindent\textcolor{ink}{\sffamily\textsc{Themen}}\enspace \themetag{reines Matt} \themetag{Mustermatt} \themetag{Spiegelmatt} \themetag{Switchback}\par\smallskip{\footnotesize\textit{Hinweis:} die Partei am Zug steht im Diagramm im Schach; die Lösung beginnt mit einem Schlag.}\par}

\clearpage
\phantomsection\label{hm:KBNPPvk}
\begin{tcolorbox}[enhanced,colback=ink,colframe=ink,arc=0pt,boxrule=0pt,left=4mm,right=4mm,top=2.5mm,bottom=2.5mm]
\color{white}\sffamily{\LARGE\bfseries KBNPPvk}\hfill{\Large h\#7}\\[0.5mm]
{\small 6 Steine\quad Schwarz am Zug}\hfill{\small Nr.~491\qquad\href{https://helpman.komtera.lt/?fen=8/8/8/8/4N3/K3P3/4P3/kB6&moves=7}{\color{white}\bfseries lösen\,\ensuremath{\nearrow}}}
\end{tcolorbox}
\vspace{2mm}
\noindent\begin{minipage}[t]{36mm}\vspace{0pt}\sffamily\small\raggedright
\hmfigure{längste eindeutige}{{\Large h\#7}}\hmfigure{eindeutig bei dieser Länge}{{\Large 3} Stellungen}\hmfigure{längstes Matt}{{\Large h\#8}}\hmfigure{in der Klasse}{34,1 Milliarden Stellungen,\\ 11,9 Milliarden mit Hilfsmatt}
\end{minipage}\hfill
\begin{minipage}[t]{80mm}\vspace{0pt}\begin{tikzpicture}\begin{axis}[width=72mm,height=42mm,ybar,bar width=1.6pt,xmin=-1,xmax=17,ymin=0,ymax=2.752,axis y line*=left,axis x line*=bottom,xlabel={\scriptsize Halbzüge bis Matt},ylabel={\scriptsize Stellungen (Mrd.)},tick label style={font=\tiny},label style={font=\scriptsize},enlarge x limits=false,ylabel shift=-3pt,xlabel shift=-3pt]\addplot[fill=ink!55,draw=none] coordinates {(0,0.0009999) (1,0.01278) (2,0.03293) (3,0.2406) (4,0.7581) (5,1.76) (6,2.502) (7,1.852) (8,2.346) (9,1.09) (10,1.017) (11,0.1789) (12,0.1026) (13,0.003466) (14,0.001091) (15,9.62e-07) (16,2.03e-07)};\draw[densely dotted] (axis cs:14,0) -- (axis cs:14,2.752);\end{axis}\begin{axis}[width=72mm,height=42mm,axis y line*=right,axis x line=none,ymode=log,xmin=-1,xmax=17,ylabel={\scriptsize eindeutig (log)},tick label style={font=\tiny},label style={font=\scriptsize},ylabel shift=-3pt]\addplot[mark=*,mark size=0.7pt,draw=black,thin] coordinates {(0,999886) (1,5466135) (2,15578812) (3,122540544) (4,224966747) (5,173469004) (6,138836690) (7,53340258) (8,34282735) (9,4902107) (10,2073089) (11,139346) (12,26247) (13,1) (14,3)};\end{axis}\end{tikzpicture}\end{minipage}
\vspace{3mm}

\noindent\begin{minipage}[t]{60mm}\vspace{0pt}\sbox\hmboard{\chessboard[setfen=8/8/8/8/4N3/K3P3/4P3/kB6 b - - 0 1]}\usebox\hmboard\\[1mm]
\makebox[\wd\hmboard]{\sffamily\hspace{3mm}\textbf{h\#7}\hfill(5+1)}\end{minipage}\hfill
\begin{minipage}[t]{52mm}\vspace{0pt}
\textcolor{ink}{\sffamily\textsc{Lösung}}\par\smallskip\small
\begin{tabular}{@{}r@{\enspace}l@{\enspace}l@{}}
\textbf{1.} & \symking{}xb1 & \symknight{}d6\\
\textbf{2.} & \symking{}c2 & e4\\
\textbf{3.} & \symking{}c3 & e5\\
\textbf{4.} & \symking{}d4 & e6\\
\textbf{5.} & \symking{}c5 & e7\\
\textbf{6.} & \symking{}b6 & e8=\symqueen{}\\
\textbf{7.} & \symking{}a5 & \symqueen{}b5\#\\
\end{tabular}\par\vspace{3mm}
\qrcode[height=12mm]{https://helpman.komtera.lt/?fen=8/8/8/8/4N3/K3P3/4P3/kB6&moves=7}\\{\sffamily\tiny\color{ink}scannen zum Lösen}
\end{minipage}
\vspace{2mm}

{\raggedright\noindent\textcolor{ink}{\sffamily\textsc{Themen}}\enspace \themetag{Umwandlung} \themetag{ein Stein zieht} \themetag{ohne Schach}\par\smallskip{\footnotesize\textit{Hinweis:} die Lösung beginnt mit einem Schlag.}\par}

\clearpage
\phantomsection\label{hm:KBNPvkn}
\begin{tcolorbox}[enhanced,colback=ink,colframe=ink,arc=0pt,boxrule=0pt,left=4mm,right=4mm,top=2.5mm,bottom=2.5mm]
\color{white}\sffamily{\LARGE\bfseries KBNPvkn}\hfill{\Large h\#7}\\[0.5mm]
{\small 6 Steine\quad Schwarz am Zug}\hfill{\small Nr.~492\qquad\href{https://helpman.komtera.lt/?fen=6k1/6P1/5N2/7B/3n4/8/8/K7&moves=7}{\color{white}\bfseries lösen\,\ensuremath{\nearrow}}}
\end{tcolorbox}
\vspace{2mm}
\noindent\begin{minipage}[t]{36mm}\vspace{0pt}\sffamily\small\raggedright
\hmfigure{längste eindeutige}{{\Large h\#7}}\hmfigure{eindeutig bei dieser Länge}{{\Large 7} Stellungen}\hmfigure{längstes Matt}{{\Large h\#8}}\hmfigure{in der Klasse}{45,4 Milliarden Stellungen,\\ 30,8 Milliarden mit Hilfsmatt}
\end{minipage}\hfill
\begin{minipage}[t]{80mm}\vspace{0pt}\begin{tikzpicture}\begin{axis}[width=72mm,height=42mm,ybar,bar width=1.6pt,xmin=-1,xmax=17,ymin=0,ymax=7.821,axis y line*=left,axis x line*=bottom,xlabel={\scriptsize Halbzüge bis Matt},ylabel={\scriptsize Stellungen (Mrd.)},tick label style={font=\tiny},label style={font=\scriptsize},enlarge x limits=false,ylabel shift=-3pt,xlabel shift=-3pt]\addplot[fill=ink!55,draw=none] coordinates {(0,0.001113) (1,0.01245) (2,0.04125) (3,0.3608) (4,1.067) (5,3.043) (6,4.93) (7,5.412) (8,7.11) (9,4.362) (10,3.214) (11,0.9329) (12,0.2881) (13,0.02769) (14,0.002985) (15,5.408e-05) (16,5.9e-07)};\draw[densely dotted] (axis cs:14,0) -- (axis cs:14,7.821);\end{axis}\begin{axis}[width=72mm,height=42mm,axis y line*=right,axis x line=none,ymode=log,xmin=-1,xmax=17,ylabel={\scriptsize eindeutig (log)},tick label style={font=\tiny},label style={font=\scriptsize},ylabel shift=-3pt]\addplot[mark=*,mark size=0.7pt,draw=black,thin] coordinates {(0,1112945) (1,6188373) (2,17364413) (3,134366891) (4,230329276) (5,234083490) (6,214384249) (7,98389363) (8,54794181) (9,11762422) (10,2730065) (11,471744) (12,38347) (13,515) (14,7)};\end{axis}\end{tikzpicture}\end{minipage}
\vspace{3mm}

\noindent\begin{minipage}[t]{60mm}\vspace{0pt}\sbox\hmboard{\chessboard[setfen=6k1/6P1/5N2/7B/3n4/8/8/K7 b - - 0 1]}\usebox\hmboard\\[1mm]
\makebox[\wd\hmboard]{\sffamily\hspace{3mm}\textbf{h\#7}\hfill(4+2)}\end{minipage}\hfill
\begin{minipage}[t]{52mm}\vspace{0pt}
\textcolor{ink}{\sffamily\textsc{Lösung}}\par\smallskip\small
\begin{tabular}{@{}r@{\enspace}l@{\enspace}l@{}}
\textbf{1.} & \symking{}xg7 & \symking{}b2\\
\textbf{2.} & \symking{}f8 & \symking{}c3\\
\textbf{3.} & \symking{}e7 & \symking{}xd4\\
\textbf{4.} & \symking{}d8 & \symking{}c5\\
\textbf{5.} & \symking{}c8 & \symking{}b6\\
\textbf{6.} & \symking{}b8 & \symknight{}d7+\\
\textbf{7.} & \symking{}a8 & \symbishop{}f3\#\\
\end{tabular}\par\vspace{3mm}
\qrcode[height=12mm]{https://helpman.komtera.lt/?fen=6k1/6P1/5N2/7B/3n4/8/8/K7&moves=7}\\{\sffamily\tiny\color{ink}scannen zum Lösen}
\end{minipage}
\vspace{2mm}

{\raggedright\noindent\textcolor{ink}{\sffamily\textsc{Themen}}\enspace \themetag{Spiegelmatt} \themetag{ein Stein zieht}\par\smallskip{\footnotesize\textit{Hinweis:} die Partei am Zug steht im Diagramm im Schach; die Lösung beginnt mit einem Schlag.}\par}

\clearpage
\phantomsection\label{hm:KBNPvkr}
\begin{tcolorbox}[enhanced,colback=ink,colframe=ink,arc=0pt,boxrule=0pt,left=4mm,right=4mm,top=2.5mm,bottom=2.5mm]
\color{white}\sffamily{\LARGE\bfseries KBNPvkr}\hfill{\Large h\#7}\\[0.5mm]
{\small 6 Steine\quad Schwarz am Zug}\hfill{\small Nr.~493\qquad\href{https://helpman.komtera.lt/?fen=8/8/3B4/8/8/8/r1N4P/k1K5&moves=7}{\color{white}\bfseries lösen\,\ensuremath{\nearrow}}}
\end{tcolorbox}
\vspace{2mm}
\noindent\begin{minipage}[t]{36mm}\vspace{0pt}\sffamily\small\raggedright
\hmfigure{längste eindeutige}{{\Large h\#7}}\hmfigure{eindeutig bei dieser Länge}{{\Large 5} Stellungen}\hmfigure{längstes Matt}{{\Large h\#9}}\hmfigure{in der Klasse}{45,4 Milliarden Stellungen,\\ 28,7 Milliarden mit Hilfsmatt}
\end{minipage}\hfill
\begin{minipage}[t]{80mm}\vspace{0pt}\begin{tikzpicture}\begin{axis}[width=72mm,height=42mm,ybar,bar width=1.6pt,xmin=-1,xmax=19,ymin=0,ymax=6.473,axis y line*=left,axis x line*=bottom,xlabel={\scriptsize Halbzüge bis Matt},ylabel={\scriptsize Stellungen (Mrd.)},tick label style={font=\tiny},label style={font=\scriptsize},enlarge x limits=false,ylabel shift=-3pt,xlabel shift=-3pt]\addplot[fill=ink!55,draw=none] coordinates {(0,0.0007616) (1,0.008963) (2,0.04002) (3,0.3707) (4,1.298) (5,3.308) (6,4.902) (7,5.677) (8,5.885) (9,3.994) (10,2.272) (11,0.7679) (12,0.1751) (13,0.01735) (14,0.0008393) (15,3.509e-05) (16,4.78e-07) (17,2.8e-08) (18,5e-09)};\draw[densely dotted] (axis cs:14,0) -- (axis cs:14,6.473);\end{axis}\begin{axis}[width=72mm,height=42mm,axis y line*=right,axis x line=none,ymode=log,xmin=-1,xmax=19,ylabel={\scriptsize eindeutig (log)},tick label style={font=\tiny},label style={font=\scriptsize},ylabel shift=-3pt]\addplot[mark=*,mark size=0.7pt,draw=black,thin] coordinates {(0,761553) (1,4945105) (2,18484096) (3,141771551) (4,294119872) (5,292395012) (6,210874545) (7,103500492) (8,37839361) (9,9476541) (10,1642822) (11,207737) (12,8109) (13,391) (14,5)};\end{axis}\end{tikzpicture}\end{minipage}
\vspace{3mm}

\noindent\begin{minipage}[t]{60mm}\vspace{0pt}\sbox\hmboard{\chessboard[setfen=8/8/3B4/8/8/8/r1N4P/k1K5 b - - 0 1]}\usebox\hmboard\\[1mm]
\makebox[\wd\hmboard]{\sffamily\hspace{3mm}\textbf{h\#7}\hfill(4+2)}\end{minipage}\hfill
\begin{minipage}[t]{52mm}\vspace{0pt}
\textcolor{ink}{\sffamily\textsc{Lösung}}\par\smallskip\small
\begin{tabular}{@{}r@{\enspace}l@{\enspace}l@{}}
\textbf{1.} & \symrook{}xc2+ & \symking{}xc2\\
\textbf{2.} & \symking{}a2 & h4\\
\textbf{3.} & \symking{}a1 & h5\\
\textbf{4.} & \symking{}a2 & h6\\
\textbf{5.} & \symking{}a1 & h7\\
\textbf{6.} & \symking{}a2 & h8=\symrook{}\\
\textbf{7.} & \symking{}a1 & \symrook{}a8\#\\
\end{tabular}\par\vspace{3mm}
\qrcode[height=12mm]{https://helpman.komtera.lt/?fen=8/8/3B4/8/8/8/r1N4P/k1K5&moves=7}\\{\sffamily\tiny\color{ink}scannen zum Lösen}
\end{minipage}
\vspace{2mm}

{\raggedright\noindent\textcolor{ink}{\sffamily\textsc{Themen}}\enspace \themetag{reines Matt} \themetag{Spiegelmatt} \themetag{Umwandlung} \themetag{Unterverwandlung} \themetag{Excelsior} \themetag{Switchback} \themetag{Pendel}\par\smallskip{\footnotesize\textit{Hinweis:} die Partei am Zug steht im Diagramm im Schach; die Lösung beginnt mit einem Schlag.}\par}

\clearpage
\phantomsection\label{hm:KBNvkbb}
\begin{tcolorbox}[enhanced,colback=ink,colframe=ink,arc=0pt,boxrule=0pt,left=4mm,right=4mm,top=2.5mm,bottom=2.5mm]
\color{white}\sffamily{\LARGE\bfseries KBNvkbb}\hfill{\Large h\#7}\\[0.5mm]
{\small 6 Steine\quad Schwarz am Zug}\hfill{\small Nr.~494\qquad\href{https://helpman.komtera.lt/?fen=7N/b6k/8/8/3B4/b7/8/K7&moves=7}{\color{white}\bfseries lösen\,\ensuremath{\nearrow}}}
\end{tcolorbox}
\vspace{2mm}
\noindent\begin{minipage}[t]{36mm}\vspace{0pt}\sffamily\small\raggedright
\hmfigure{längste eindeutige}{{\Large h\#7}}\hmfigure{eindeutig bei dieser Länge}{{\Large 161} Stellungen}\hmfigure{längstes Matt}{{\Large h\#10}}\hmfigure{in der Klasse}{15,5 Milliarden Stellungen,\\ 4,7 Milliarden mit Hilfsmatt}
\end{minipage}\hfill
\begin{minipage}[t]{80mm}\vspace{0pt}\begin{tikzpicture}\begin{axis}[width=72mm,height=42mm,ybar,bar width=1.6pt,xmin=-1,xmax=21,ymin=0,ymax=1.215,axis y line*=left,axis x line*=bottom,xlabel={\scriptsize Halbzüge bis Matt},ylabel={\scriptsize Stellungen (Mrd.)},tick label style={font=\tiny},label style={font=\scriptsize},enlarge x limits=false,ylabel shift=-3pt,xlabel shift=-3pt]\addplot[fill=ink!55,draw=none] coordinates {(0,0.0001161) (1,0.000458) (2,0.003025) (3,0.01641) (4,0.07047) (5,0.2488) (6,0.5818) (7,1.008) (8,1.104) (9,0.9161) (10,0.4365) (11,0.2107) (12,0.08025) (13,0.02127) (14,0.005528) (15,0.0006778) (16,4.845e-05) (17,1.8e-08) (18,1.203e-05) (20,3.211e-06)};\draw[densely dotted] (axis cs:14,0) -- (axis cs:14,1.215);\end{axis}\begin{axis}[width=72mm,height=42mm,axis y line*=right,axis x line=none,ymode=log,xmin=-1,xmax=21,ylabel={\scriptsize eindeutig (log)},tick label style={font=\tiny},label style={font=\scriptsize},ylabel shift=-3pt]\addplot[mark=*,mark size=0.7pt,draw=black,thin] coordinates {(0,116052) (1,440140) (2,2398618) (3,7653601) (4,19326849) (5,29857293) (6,23608639) (7,12680071) (8,4557405) (9,1447159) (10,363642) (11,54201) (12,16204) (13,851) (14,161)};\end{axis}\end{tikzpicture}\end{minipage}
\vspace{3mm}

\noindent\begin{minipage}[t]{60mm}\vspace{0pt}\sbox\hmboard{\chessboard[setfen=7N/b6k/8/8/3B4/b7/8/K7 b - - 0 1]}\usebox\hmboard\\[1mm]
\makebox[\wd\hmboard]{\sffamily\hspace{3mm}\textbf{h\#7}\hfill(3+3)}\end{minipage}\hfill
\begin{minipage}[t]{52mm}\vspace{0pt}
\textcolor{ink}{\sffamily\textsc{Lösung}}\par\smallskip\small
\begin{tabular}{@{}r@{\enspace}l@{\enspace}l@{}}
\textbf{1.} & \symbishop{}b8 & \symknight{}f7\\
\textbf{2.} & \symking{}g6 & \symking{}a2\\
\textbf{3.} & \symking{}f5 & \symking{}b3\\
\textbf{4.} & \symking{}e4 & \symknight{}d6+\\
\textbf{5.} & \symking{}d3 & \symknight{}b5\\
\textbf{6.} & \symking{}d2 & \symknight{}c3\\
\textbf{7.} & \symking{}c1 & \symbishop{}e3\#\\
\end{tabular}\par\vspace{3mm}
\qrcode[height=12mm]{https://helpman.komtera.lt/?fen=7N/b6k/8/8/3B4/b7/8/K7&moves=7}\\{\sffamily\tiny\color{ink}scannen zum Lösen}
\end{minipage}
\vspace{2mm}

{\raggedright\noindent\textcolor{ink}{\sffamily\textsc{Themen}}\enspace \themetag{reines Matt} \themetag{Mustermatt} \themetag{Spiegelmatt} \themetag{ohne Schlag}\par}

\clearpage
\phantomsection\label{hm:KBNvkbn}
\begin{tcolorbox}[enhanced,colback=ink,colframe=ink,arc=0pt,boxrule=0pt,left=4mm,right=4mm,top=2.5mm,bottom=2.5mm]
\color{white}\sffamily{\LARGE\bfseries KBNvkbn}\hfill{\Large h\#7}\\[0.5mm]
{\small 6 Steine\quad Schwarz am Zug}\hfill{\small Nr.~495\qquad\href{https://helpman.komtera.lt/?fen=8/8/8/2n2Nb1/6Bk/8/8/K7&moves=7}{\color{white}\bfseries lösen\,\ensuremath{\nearrow}}}
\end{tcolorbox}
\vspace{2mm}
\noindent\begin{minipage}[t]{36mm}\vspace{0pt}\sffamily\small\raggedright
\hmfigure{längste eindeutige}{{\Large h\#7}}\hmfigure{eindeutig bei dieser Länge}{{\Large 18} Stellungen}\hmfigure{längstes Matt}{{\Large h\#12}}\hmfigure{in der Klasse}{15,5 Milliarden Stellungen,\\ 9,69 Milliarden mit Hilfsmatt}
\end{minipage}\hfill
\begin{minipage}[t]{80mm}\vspace{0pt}\begin{tikzpicture}\begin{axis}[width=72mm,height=42mm,ybar,bar width=1.6pt,xmin=-1,xmax=25,ymin=0,ymax=2.549,axis y line*=left,axis x line*=bottom,xlabel={\scriptsize Halbzüge bis Matt},ylabel={\scriptsize Stellungen (Mrd.)},tick label style={font=\tiny},label style={font=\scriptsize},enlarge x limits=false,ylabel shift=-3pt,xlabel shift=-3pt]\addplot[fill=ink!55,draw=none] coordinates {(0,0.0001828) (1,0.0007258) (2,0.004037) (3,0.02529) (4,0.1105) (5,0.4116) (6,1.008) (7,1.793) (8,2.318) (9,1.968) (10,1.187) (11,0.5915) (12,0.2087) (13,0.05376) (14,0.009017) (15,0.0007913) (16,1.503e-05) (17,1e-09) (18,5.532e-06) (19,1.6e-08) (20,5.451e-06) (22,4.159e-06) (23,1e-09) (24,7.02e-07)};\draw[densely dotted] (axis cs:14,0) -- (axis cs:14,2.549);\end{axis}\begin{axis}[width=72mm,height=42mm,axis y line*=right,axis x line=none,ymode=log,xmin=-1,xmax=25,ylabel={\scriptsize eindeutig (log)},tick label style={font=\tiny},label style={font=\scriptsize},ylabel shift=-3pt]\addplot[mark=*,mark size=0.7pt,draw=black,thin] coordinates {(0,182822) (1,693133) (2,3025171) (3,11660303) (4,32042468) (5,46453899) (6,41213523) (7,23703916) (8,10847867) (9,3386949) (10,924122) (11,159023) (12,31387) (13,942) (14,18)};\end{axis}\end{tikzpicture}\end{minipage}
\vspace{3mm}

\noindent\begin{minipage}[t]{60mm}\vspace{0pt}\sbox\hmboard{\chessboard[setfen=8/8/8/2n2Nb1/6Bk/8/8/K7 b - - 0 1]}\usebox\hmboard\\[1mm]
\makebox[\wd\hmboard]{\sffamily\hspace{3mm}\textbf{h\#7}\hfill(3+3)}\end{minipage}\hfill
\begin{minipage}[t]{52mm}\vspace{0pt}
\textcolor{ink}{\sffamily\textsc{Lösung}}\par\smallskip\small
\begin{tabular}{@{}r@{\enspace}l@{\enspace}l@{}}
\textbf{1.} & \symking{}xg4 & \symking{}b2\\
\textbf{2.} & \symking{}f3 & \symking{}c3\\
\textbf{3.} & \symking{}e2 & \symking{}d4\\
\textbf{4.} & \symknight{}b3+ & \symking{}e4\\
\textbf{5.} & \symking{}d1 & \symking{}f3\\
\textbf{6.} & \symbishop{}c1 & \symking{}f2\\
\textbf{7.} & \symknight{}d2 & \symknight{}e3\#\\
\end{tabular}\par\vspace{3mm}
\qrcode[height=12mm]{https://helpman.komtera.lt/?fen=8/8/8/2n2Nb1/6Bk/8/8/K7&moves=7}\\{\sffamily\tiny\color{ink}scannen zum Lösen}
\end{minipage}
\vspace{2mm}

{\raggedright\noindent\textcolor{ink}{\sffamily\textsc{Themen}}\enspace \themetag{reines Matt} \themetag{Mustermatt} \themetag{Idealmatt} \themetag{Selbstblock}\par\smallskip{\footnotesize\textit{Hinweis:} die Partei am Zug steht im Diagramm im Schach; die Lösung beginnt mit einem Schlag.}\par}

\clearpage
\phantomsection\label{hm:KBNvknn}
\begin{tcolorbox}[enhanced,colback=ink,colframe=ink,arc=0pt,boxrule=0pt,left=4mm,right=4mm,top=2.5mm,bottom=2.5mm]
\color{white}\sffamily{\LARGE\bfseries KBNvknn}\hfill{\Large h\#7}\\[0.5mm]
{\small 6 Steine\quad Schwarz am Zug}\hfill{\small Nr.~496\qquad\href{https://helpman.komtera.lt/?fen=4N3/7k/8/n7/8/8/n7/KB6&moves=7}{\color{white}\bfseries lösen\,\ensuremath{\nearrow}}}
\end{tcolorbox}
\vspace{2mm}
\noindent\begin{minipage}[t]{36mm}\vspace{0pt}\sffamily\small\raggedright
\hmfigure{längste eindeutige}{{\Large h\#7}}\hmfigure{eindeutig bei dieser Länge}{{\Large 33} Stellungen}\hmfigure{längstes Matt}{{\Large h\#11}}\hmfigure{in der Klasse}{15,5 Milliarden Stellungen,\\ 4,97 Milliarden mit Hilfsmatt}
\end{minipage}\hfill
\begin{minipage}[t]{80mm}\vspace{0pt}\begin{tikzpicture}\begin{axis}[width=72mm,height=42mm,ybar,bar width=1.6pt,xmin=-1,xmax=23,ymin=0,ymax=1.295,axis y line*=left,axis x line*=bottom,xlabel={\scriptsize Halbzüge bis Matt},ylabel={\scriptsize Stellungen (Mrd.)},tick label style={font=\tiny},label style={font=\scriptsize},enlarge x limits=false,ylabel shift=-3pt,xlabel shift=-3pt]\addplot[fill=ink!55,draw=none] coordinates {(0,6.917e-05) (1,0.0002832) (2,0.001117) (3,0.008625) (4,0.03296) (5,0.1438) (6,0.3721) (7,0.7459) (8,1.177) (9,1.09) (10,0.8326) (11,0.3986) (12,0.1244) (13,0.03474) (14,0.004131) (15,0.0005459) (16,1.085e-05) (18,1.142e-06) (20,8.84e-07) (22,1.142e-06)};\draw[densely dotted] (axis cs:14,0) -- (axis cs:14,1.295);\end{axis}\begin{axis}[width=72mm,height=42mm,axis y line*=right,axis x line=none,ymode=log,xmin=-1,xmax=23,ylabel={\scriptsize eindeutig (log)},tick label style={font=\tiny},label style={font=\scriptsize},ylabel shift=-3pt]\addplot[mark=*,mark size=0.7pt,draw=black,thin] coordinates {(0,69174) (1,267220) (2,705977) (3,3706475) (4,9503033) (5,15233243) (6,17035357) (7,10541668) (8,6378742) (9,1896966) (10,480154) (11,83865) (12,13796) (13,2232) (14,33)};\end{axis}\end{tikzpicture}\end{minipage}
\vspace{3mm}

\noindent\begin{minipage}[t]{60mm}\vspace{0pt}\sbox\hmboard{\chessboard[setfen=4N3/7k/8/n7/8/8/n7/KB6 b - - 0 1]}\usebox\hmboard\\[1mm]
\makebox[\wd\hmboard]{\sffamily\hspace{3mm}\textbf{h\#7}\hfill(3+3)}\end{minipage}\hfill
\begin{minipage}[t]{52mm}\vspace{0pt}
\textcolor{ink}{\sffamily\textsc{Lösung}}\par\smallskip\small
\begin{tabular}{@{}r@{\enspace}l@{\enspace}l@{}}
\textbf{1.} & \symking{}g8 & \symking{}b2\\
\textbf{2.} & \symking{}f7 & \symking{}a3\\
\textbf{3.} & \symking{}e6 & \symking{}a4\\
\textbf{4.} & \symking{}d5 & \symking{}xa5\\
\textbf{5.} & \symking{}c4 & \symknight{}d6+\\
\textbf{6.} & \symking{}b3 & \symbishop{}c2+\\
\textbf{7.} & \symking{}a3 & \symknight{}c4\#\\
\end{tabular}\par\vspace{3mm}
\qrcode[height=12mm]{https://helpman.komtera.lt/?fen=4N3/7k/8/n7/8/8/n7/KB6&moves=7}\\{\sffamily\tiny\color{ink}scannen zum Lösen}
\end{minipage}
\vspace{2mm}

{\raggedright\noindent\textcolor{ink}{\sffamily\textsc{Themen}}\enspace \themetag{ein Stein zieht}\par\smallskip{\footnotesize\textit{Hinweis:} das Diagramm hat keinen legalen letzten Zug, kann also in keiner Partie entstehen.}\par}

\clearpage
\phantomsection\label{hm:KBNvknp}
\begin{tcolorbox}[enhanced,colback=ink,colframe=ink,arc=0pt,boxrule=0pt,left=4mm,right=4mm,top=2.5mm,bottom=2.5mm]
\color{white}\sffamily{\LARGE\bfseries KBNvknp}\hfill{\Large h\#7}\\[0.5mm]
{\small 6 Steine\quad Schwarz am Zug}\hfill{\small Nr.~497\qquad\href{https://helpman.komtera.lt/?fen=n6N/7k/8/8/2pB4/8/8/K7&moves=7}{\color{white}\bfseries lösen\,\ensuremath{\nearrow}}}
\end{tcolorbox}
\vspace{2mm}
\noindent\begin{minipage}[t]{36mm}\vspace{0pt}\sffamily\small\raggedright
\hmfigure{längste eindeutige}{{\Large h\#7}}\hmfigure{eindeutig bei dieser Länge}{{\Large 220} Stellungen}\hmfigure{längstes Matt}{{\Large h\#12}}\hmfigure{in der Klasse}{45,4 Milliarden Stellungen,\\ 30,8 Milliarden mit Hilfsmatt}
\end{minipage}\hfill
\begin{minipage}[t]{80mm}\vspace{0pt}\begin{tikzpicture}\begin{axis}[width=72mm,height=42mm,ybar,bar width=1.6pt,xmin=-1,xmax=25,ymin=0,ymax=7.998,axis y line*=left,axis x line*=bottom,xlabel={\scriptsize Halbzüge bis Matt},ylabel={\scriptsize Stellungen (Mrd.)},tick label style={font=\tiny},label style={font=\scriptsize},enlarge x limits=false,ylabel shift=-3pt,xlabel shift=-3pt]\addplot[fill=ink!55,draw=none] coordinates {(0,0.000629) (1,0.002524) (2,0.008737) (3,0.05946) (4,0.2113) (5,0.8924) (6,2.278) (7,4.478) (8,7.271) (9,6.459) (10,5.45) (11,2.453) (12,0.9411) (13,0.218) (14,0.03423) (15,0.003621) (16,0.0001253) (17,5.9e-08) (18,1.377e-05) (20,4.376e-06) (22,9.664e-06) (23,2e-09) (24,6.25e-07)};\draw[densely dotted] (axis cs:14,0) -- (axis cs:14,7.998);\end{axis}\begin{axis}[width=72mm,height=42mm,axis y line*=right,axis x line=none,ymode=log,xmin=-1,xmax=25,ylabel={\scriptsize eindeutig (log)},tick label style={font=\tiny},label style={font=\scriptsize},ylabel shift=-3pt]\addplot[mark=*,mark size=0.7pt,draw=black,thin] coordinates {(0,629021) (1,2384480) (2,5735237) (3,27134755) (4,66429131) (5,105874158) (6,129050409) (7,69343993) (8,48977797) (9,11948851) (10,4007319) (11,620226) (12,144448) (13,11674) (14,220)};\end{axis}\end{tikzpicture}\end{minipage}
\vspace{3mm}

\noindent\begin{minipage}[t]{60mm}\vspace{0pt}\sbox\hmboard{\chessboard[setfen=n6N/7k/8/8/2pB4/8/8/K7 b - - 0 1]}\usebox\hmboard\\[1mm]
\makebox[\wd\hmboard]{\sffamily\hspace{3mm}\textbf{h\#7}\hfill(3+3)}\end{minipage}\hfill
\begin{minipage}[t]{52mm}\vspace{0pt}
\textcolor{ink}{\sffamily\textsc{Lösung}}\par\smallskip\small
\begin{tabular}{@{}r@{\enspace}l@{\enspace}l@{}}
\textbf{1.} & c3 & \symknight{}f7\\
\textbf{2.} & \symking{}g6 & \symking{}a2\\
\textbf{3.} & \symking{}f5 & \symking{}b3\\
\textbf{4.} & \symking{}e4 & \symknight{}d6+\\
\textbf{5.} & \symking{}d3 & \symknight{}b5\\
\textbf{6.} & \symking{}d2 & \symknight{}xc3\\
\textbf{7.} & \symking{}c1 & \symbishop{}e3\#\\
\end{tabular}\par\vspace{3mm}
\qrcode[height=12mm]{https://helpman.komtera.lt/?fen=n6N/7k/8/8/2pB4/8/8/K7&moves=7}\\{\sffamily\tiny\color{ink}scannen zum Lösen}
\end{minipage}
\vspace{2mm}

{\raggedright\noindent\textcolor{ink}{\sffamily\textsc{Themen}}\enspace \themetag{reines Matt} \themetag{Mustermatt} \themetag{Spiegelmatt}\par}

\clearpage
\phantomsection\label{hm:KBNvkpp}
\begin{tcolorbox}[enhanced,colback=ink,colframe=ink,arc=0pt,boxrule=0pt,left=4mm,right=4mm,top=2.5mm,bottom=2.5mm]
\color{white}\sffamily{\LARGE\bfseries KBNvkpp}\hfill{\Large h\#7}\\[0.5mm]
{\small 6 Steine\quad Schwarz am Zug}\hfill{\small Nr.~498\qquad\href{https://helpman.komtera.lt/?fen=8/8/pp6/kB6/8/3N4/8/K7&moves=7}{\color{white}\bfseries lösen\,\ensuremath{\nearrow}}}
\end{tcolorbox}
\vspace{2mm}
\noindent\begin{minipage}[t]{36mm}\vspace{0pt}\sffamily\small\raggedright
\hmfigure{längste eindeutige}{{\Large h\#7}}\hmfigure{eindeutig bei dieser Länge}{{\Large 238} Stellungen}\hmfigure{längstes Matt}{{\Large h\#16}}\hmfigure{in der Klasse}{34,1 Milliarden Stellungen,\\ 11,8 Milliarden mit Hilfsmatt}
\end{minipage}\hfill
\begin{minipage}[t]{80mm}\vspace{0pt}\begin{tikzpicture}\begin{axis}[width=72mm,height=42mm,ybar,bar width=1.6pt,xmin=-1,xmax=33,ymin=0,ymax=2.933,axis y line*=left,axis x line*=bottom,xlabel={\scriptsize Halbzüge bis Matt},ylabel={\scriptsize Stellungen (Mrd.)},tick label style={font=\tiny},label style={font=\scriptsize},enlarge x limits=false,ylabel shift=-3pt,xlabel shift=-3pt]\addplot[fill=ink!55,draw=none] coordinates {(0,0.0003572) (1,0.001412) (2,0.004237) (3,0.02538) (4,0.0824) (5,0.334) (6,0.8114) (7,1.585) (8,2.666) (9,2.391) (10,2.313) (11,1.003) (12,0.4936) (13,0.1032) (14,0.02416) (15,0.002297) (16,0.0001067) (17,1.53e-07) (18,3.459e-05) (19,8.6e-08) (20,1.341e-05) (21,3.1e-08) (22,5.542e-06) (23,6.6e-08) (24,1.627e-06) (25,3e-09) (26,1.529e-06) (27,3.5e-08) (28,9.2e-08) (30,2.1e-08) (32,8e-09)};\draw[densely dotted] (axis cs:14,0) -- (axis cs:14,2.933);\end{axis}\begin{axis}[width=72mm,height=42mm,axis y line*=right,axis x line=none,ymode=log,xmin=-1,xmax=33,ylabel={\scriptsize eindeutig (log)},tick label style={font=\tiny},label style={font=\scriptsize},ylabel shift=-3pt]\addplot[mark=*,mark size=0.7pt,draw=black,thin] coordinates {(0,357158) (1,1334139) (2,3266506) (3,13295225) (4,29262533) (5,46893889) (6,60103511) (7,29503938) (8,25822477) (9,5492436) (10,2636037) (11,349384) (12,122515) (13,9553) (14,238)};\end{axis}\end{tikzpicture}\end{minipage}
\vspace{3mm}

\noindent\begin{minipage}[t]{60mm}\vspace{0pt}\sbox\hmboard{\chessboard[setfen=8/8/pp6/kB6/8/3N4/8/K7 b - - 0 1]}\usebox\hmboard\\[1mm]
\makebox[\wd\hmboard]{\sffamily\hspace{3mm}\textbf{h\#7}\hfill(3+3)}\end{minipage}\hfill
\begin{minipage}[t]{52mm}\vspace{0pt}
\textcolor{ink}{\sffamily\textsc{Lösung}}\par\smallskip\small
\begin{tabular}{@{}r@{\enspace}l@{\enspace}l@{}}
\textbf{1.} & \symking{}xb5 & \symking{}b2\\
\textbf{2.} & a5 & \symking{}c3\\
\textbf{3.} & a4 & \symking{}d4\\
\textbf{4.} & a3 & \symking{}d5\\
\textbf{5.} & \symking{}a4 & \symking{}c6\\
\textbf{6.} & b5 & \symking{}b6\\
\textbf{7.} & b4 & \symknight{}c5\#\\
\end{tabular}\par\vspace{3mm}
\qrcode[height=12mm]{https://helpman.komtera.lt/?fen=8/8/pp6/kB6/8/3N4/8/K7&moves=7}\\{\sffamily\tiny\color{ink}scannen zum Lösen}
\end{minipage}
\vspace{2mm}

{\raggedright\noindent\textcolor{ink}{\sffamily\textsc{Themen}}\enspace \themetag{reines Matt} \themetag{Mustermatt} \themetag{Idealmatt} \themetag{Selbstblock} \themetag{ohne Schach} \themetag{Umnov}\par\smallskip{\footnotesize\textit{Hinweis:} die Lösung beginnt mit einem Schlag.}\par}

\clearpage
\phantomsection\label{hm:KBNvkrn}
\begin{tcolorbox}[enhanced,colback=ink,colframe=ink,arc=0pt,boxrule=0pt,left=4mm,right=4mm,top=2.5mm,bottom=2.5mm]
\color{white}\sffamily{\LARGE\bfseries KBNvkrn}\hfill{\Large h\#7}\\[0.5mm]
{\small 6 Steine\quad Schwarz am Zug}\hfill{\small Nr.~499\qquad\href{https://helpman.komtera.lt/?fen=6rk/6BN/8/8/2n5/8/8/K7&moves=7}{\color{white}\bfseries lösen\,\ensuremath{\nearrow}}}
\end{tcolorbox}
\vspace{2mm}
\noindent\begin{minipage}[t]{36mm}\vspace{0pt}\sffamily\small\raggedright
\hmfigure{längste eindeutige}{{\Large h\#7}}\hmfigure{eindeutig bei dieser Länge}{{\Large 5} Stellungen}\hmfigure{längstes Matt}{{\Large h\#12}}\hmfigure{in der Klasse}{15,5 Milliarden Stellungen,\\ 9,3 Milliarden mit Hilfsmatt}
\end{minipage}\hfill
\begin{minipage}[t]{80mm}\vspace{0pt}\begin{tikzpicture}\begin{axis}[width=72mm,height=42mm,ybar,bar width=1.6pt,xmin=-1,xmax=25,ymin=0,ymax=2.613,axis y line*=left,axis x line*=bottom,xlabel={\scriptsize Halbzüge bis Matt},ylabel={\scriptsize Stellungen (Mrd.)},tick label style={font=\tiny},label style={font=\scriptsize},enlarge x limits=false,ylabel shift=-3pt,xlabel shift=-3pt]\addplot[fill=ink!55,draw=none] coordinates {(0,0.0001013) (1,0.0004215) (2,0.001941) (3,0.01571) (4,0.06777) (5,0.3067) (6,0.831) (7,1.736) (8,2.375) (9,2.204) (10,1.103) (11,0.5584) (12,0.07616) (13,0.02116) (14,0.0004004) (15,5.935e-05) (16,6.1e-06) (17,7.4e-08) (18,4.656e-06) (19,4e-09) (20,4.517e-06) (22,5.527e-06) (23,2e-08) (24,8.18e-07)};\draw[densely dotted] (axis cs:14,0) -- (axis cs:14,2.613);\end{axis}\begin{axis}[width=72mm,height=42mm,axis y line*=right,axis x line=none,ymode=log,xmin=-1,xmax=25,ylabel={\scriptsize eindeutig (log)},tick label style={font=\tiny},label style={font=\scriptsize},ylabel shift=-3pt]\addplot[mark=*,mark size=0.7pt,draw=black,thin] coordinates {(0,101272) (1,397910) (2,1276010) (3,6987689) (4,22522448) (5,41267592) (6,48799648) (7,32431179) (8,11481303) (9,3855293) (10,444417) (11,112510) (12,2259) (13,207) (14,5)};\end{axis}\end{tikzpicture}\end{minipage}
\vspace{3mm}

\noindent\begin{minipage}[t]{60mm}\vspace{0pt}\sbox\hmboard{\chessboard[setfen=6rk/6BN/8/8/2n5/8/8/K7 b - - 0 1]}\usebox\hmboard\\[1mm]
\makebox[\wd\hmboard]{\sffamily\hspace{3mm}\textbf{h\#7}\hfill(3+3)}\end{minipage}\hfill
\begin{minipage}[t]{52mm}\vspace{0pt}
\textcolor{ink}{\sffamily\textsc{Lösung}}\par\smallskip\small
\begin{tabular}{@{}r@{\enspace}l@{\enspace}l@{}}
\textbf{1.} & \symking{}xg7 & \symking{}a2\\
\textbf{2.} & \symknight{}d6 & \symking{}b3\\
\textbf{3.} & \symknight{}e8 & \symking{}c4\\
\textbf{4.} & \symrook{}f8 & \symking{}d5\\
\textbf{5.} & \symrook{}h8 & \symking{}e6\\
\textbf{6.} & \symking{}g8 & \symking{}e7\\
\textbf{7.} & \symknight{}g7 & \symknight{}f6\#\\
\end{tabular}\par\vspace{3mm}
\qrcode[height=12mm]{https://helpman.komtera.lt/?fen=6rk/6BN/8/8/2n5/8/8/K7&moves=7}\\{\sffamily\tiny\color{ink}scannen zum Lösen}
\end{minipage}
\vspace{2mm}

{\raggedright\noindent\textcolor{ink}{\sffamily\textsc{Themen}}\enspace \themetag{reines Matt} \themetag{Mustermatt} \themetag{Idealmatt} \themetag{Selbstblock} \themetag{ohne Schach}\par\smallskip{\footnotesize\textit{Hinweis:} die Partei am Zug steht im Diagramm im Schach; die Lösung beginnt mit einem Schlag.}\par}

\clearpage
\phantomsection\label{hm:KBPPPvk}
\begin{tcolorbox}[enhanced,colback=ink,colframe=ink,arc=0pt,boxrule=0pt,left=4mm,right=4mm,top=2.5mm,bottom=2.5mm]
\color{white}\sffamily{\LARGE\bfseries KBPPPvk}\hfill{\Large h\#7}\\[0.5mm]
{\small 6 Steine\quad Schwarz am Zug}\hfill{\small Nr.~500\qquad\href{https://helpman.komtera.lt/?fen=k1K5/P7/8/8/2B5/2P5/2P5/8&moves=7}{\color{white}\bfseries lösen\,\ensuremath{\nearrow}}}
\end{tcolorbox}
\vspace{2mm}
\noindent\begin{minipage}[t]{36mm}\vspace{0pt}\sffamily\small\raggedright
\hmfigure{längste eindeutige}{{\Large h\#7}}\hmfigure{eindeutig bei dieser Länge}{{\Large 1} Stellung}\hmfigure{längstes Matt}{{\Large h\#8}}\hmfigure{in der Klasse}{25,6 Milliarden Stellungen,\\ 3,02 Milliarden mit Hilfsmatt}
\end{minipage}\hfill
\begin{minipage}[t]{80mm}\vspace{0pt}\begin{tikzpicture}\begin{axis}[width=72mm,height=42mm,ybar,bar width=1.6pt,xmin=-1,xmax=17,ymin=0,ymax=772.8,axis y line*=left,axis x line*=bottom,xlabel={\scriptsize Halbzüge bis Matt},ylabel={\scriptsize Stellungen (Mio.)},tick label style={font=\tiny},label style={font=\scriptsize},enlarge x limits=false,ylabel shift=-3pt,xlabel shift=-3pt]\addplot[fill=ink!55,draw=none] coordinates {(0,0.2098) (1,4.015) (2,10.1) (3,69.96) (4,210.9) (5,534.1) (6,702.5) (7,441) (8,457.9) (9,215.9) (10,219.3) (11,73.43) (12,69.46) (13,4.26) (14,2.594) (15,0.000372) (16,4.2e-05)};\draw[densely dotted] (axis cs:14,0) -- (axis cs:14,772.8);\end{axis}\begin{axis}[width=72mm,height=42mm,axis y line*=right,axis x line=none,ymode=log,xmin=-1,xmax=17,ylabel={\scriptsize eindeutig (log)},tick label style={font=\tiny},label style={font=\scriptsize},ylabel shift=-3pt]\addplot[mark=*,mark size=0.7pt,draw=black,thin] coordinates {(0,209795) (1,1442342) (2,4243008) (3,34953721) (4,61047578) (5,47328421) (6,28051549) (7,12729806) (8,7357245) (9,2518335) (10,1450705) (11,229453) (12,83023) (13,13) (14,1)};\end{axis}\end{tikzpicture}\end{minipage}
\vspace{3mm}

\noindent\begin{minipage}[t]{60mm}\vspace{0pt}\sbox\hmboard{\chessboard[setfen=k1K5/P7/8/8/2B5/2P5/2P5/8 b - - 0 1]}\usebox\hmboard\\[1mm]
\makebox[\wd\hmboard]{\sffamily\hspace{3mm}\textbf{h\#7}\hfill(5+1)}\end{minipage}\hfill
\begin{minipage}[t]{52mm}\vspace{0pt}
\textcolor{ink}{\sffamily\textsc{Lösung}}\par\smallskip\small
\begin{tabular}{@{}r@{\enspace}l@{\enspace}l@{}}
\textbf{1.} & \symking{}xa7 & \symbishop{}d5\\
\textbf{2.} & \symking{}b6 & c4\\
\textbf{3.} & \symking{}c5 & \symking{}b7\\
\textbf{4.} & \symking{}d6 & \symking{}a6\\
\textbf{5.} & \symking{}c7 & c5\\
\textbf{6.} & \symking{}b8 & c6\\
\textbf{7.} & \symking{}a8 & c7\#\\
\end{tabular}\par\vspace{3mm}
\qrcode[height=12mm]{https://helpman.komtera.lt/?fen=k1K5/P7/8/8/2B5/2P5/2P5/8&moves=7}\\{\sffamily\tiny\color{ink}scannen zum Lösen}
\end{minipage}
\vspace{2mm}

{\raggedright\noindent\textcolor{ink}{\sffamily\textsc{Themen}}\enspace \themetag{Spiegelmatt} \themetag{Rundlauf} \themetag{ein Stein zieht} \themetag{ohne Schach}\par\smallskip{\footnotesize\textit{Hinweis:} die Lösung beginnt mit einem Schlag.}\par}

\clearpage
\phantomsection\label{hm:KBPPvkb}
\begin{tcolorbox}[enhanced,colback=ink,colframe=ink,arc=0pt,boxrule=0pt,left=4mm,right=4mm,top=2.5mm,bottom=2.5mm]
\color{white}\sffamily{\LARGE\bfseries KBPPvkb}\hfill{\Large h\#7}\\[0.5mm]
{\small 6 Steine\quad Schwarz am Zug}\hfill{\small Nr.~501\qquad\href{https://helpman.komtera.lt/?fen=8/8/7k/8/B7/PK6/P7/5b2&moves=7}{\color{white}\bfseries lösen\,\ensuremath{\nearrow}}}
\end{tcolorbox}
\vspace{2mm}
\noindent\begin{minipage}[t]{36mm}\vspace{0pt}\sffamily\small\raggedright
\hmfigure{längste eindeutige}{{\Large h\#7}}\hmfigure{eindeutig bei dieser Länge}{{\Large 36} Stellungen}\hmfigure{längstes Matt}{{\Large h\#8}}\hmfigure{in der Klasse}{34,1 Milliarden Stellungen,\\ 11,5 Milliarden mit Hilfsmatt}
\end{minipage}\hfill
\begin{minipage}[t]{80mm}\vspace{0pt}\begin{tikzpicture}\begin{axis}[width=72mm,height=42mm,ybar,bar width=1.6pt,xmin=-1,xmax=17,ymin=0,ymax=2.377,axis y line*=left,axis x line*=bottom,xlabel={\scriptsize Halbzüge bis Matt},ylabel={\scriptsize Stellungen (Mrd.)},tick label style={font=\tiny},label style={font=\scriptsize},enlarge x limits=false,ylabel shift=-3pt,xlabel shift=-3pt]\addplot[fill=ink!55,draw=none] coordinates {(0,0.0003922) (1,0.006193) (2,0.02305) (3,0.1944) (4,0.5999) (5,1.667) (6,2.16) (7,1.888) (8,1.75) (9,1.243) (10,1.03) (11,0.5422) (12,0.3492) (13,0.04116) (14,0.006921) (15,6.631e-05) (16,5.79e-07)};\draw[densely dotted] (axis cs:14,0) -- (axis cs:14,2.377);\end{axis}\begin{axis}[width=72mm,height=42mm,axis y line*=right,axis x line=none,ymode=log,xmin=-1,xmax=17,ylabel={\scriptsize eindeutig (log)},tick label style={font=\tiny},label style={font=\scriptsize},ylabel shift=-3pt]\addplot[mark=*,mark size=0.7pt,draw=black,thin] coordinates {(0,392161) (1,2817940) (2,9884157) (3,74176596) (4,131788782) (5,123050795) (6,78777337) (7,45785138) (8,24174574) (9,11736740) (10,4180357) (11,873072) (12,178009) (13,1558) (14,36)};\end{axis}\end{tikzpicture}\end{minipage}
\vspace{3mm}

\noindent\begin{minipage}[t]{60mm}\vspace{0pt}\sbox\hmboard{\chessboard[setfen=8/8/7k/8/B7/PK6/P7/5b2 b - - 0 1]}\usebox\hmboard\\[1mm]
\makebox[\wd\hmboard]{\sffamily\hspace{3mm}\textbf{h\#7}\hfill(4+2)}\end{minipage}\hfill
\begin{minipage}[t]{52mm}\vspace{0pt}
\textcolor{ink}{\sffamily\textsc{Lösung}}\par\smallskip\small
\begin{tabular}{@{}r@{\enspace}l@{\enspace}l@{}}
\textbf{1.} & \symbishop{}b5 & \symbishop{}xb5\\
\textbf{2.} & \symking{}g5 & a4\\
\textbf{3.} & \symking{}f4 & a5\\
\textbf{4.} & \symking{}e3 & a6\\
\textbf{5.} & \symking{}d2 & a7\\
\textbf{6.} & \symking{}c1 & a8=\symqueen{}\\
\textbf{7.} & \symking{}b1 & \symqueen{}h1\#\\
\end{tabular}\par\vspace{3mm}
\qrcode[height=12mm]{https://helpman.komtera.lt/?fen=8/8/7k/8/B7/PK6/P7/5b2&moves=7}\\{\sffamily\tiny\color{ink}scannen zum Lösen}
\end{minipage}
\vspace{2mm}

{\raggedright\noindent\textcolor{ink}{\sffamily\textsc{Themen}}\enspace \themetag{reines Matt} \themetag{Umwandlung} \themetag{ohne Schach}\par}

\clearpage
\phantomsection\label{hm:KBPPvkn}
\begin{tcolorbox}[enhanced,colback=ink,colframe=ink,arc=0pt,boxrule=0pt,left=4mm,right=4mm,top=2.5mm,bottom=2.5mm]
\color{white}\sffamily{\LARGE\bfseries KBPPvkn}\hfill{\Large h\#7}\\[0.5mm]
{\small 6 Steine\quad Schwarz am Zug}\hfill{\small Nr.~502\qquad\href{https://helpman.komtera.lt/?fen=7k/6B1/8/8/n7/P7/P7/K7&moves=7}{\color{white}\bfseries lösen\,\ensuremath{\nearrow}}}
\end{tcolorbox}
\vspace{2mm}
\noindent\begin{minipage}[t]{36mm}\vspace{0pt}\sffamily\small\raggedright
\hmfigure{längste eindeutige}{{\Large h\#7}}\hmfigure{eindeutig bei dieser Länge}{{\Large 57} Stellungen}\hmfigure{längstes Matt}{{\Large h\#8}}\hmfigure{in der Klasse}{34,1 Milliarden Stellungen,\\ 11,8 Milliarden mit Hilfsmatt}
\end{minipage}\hfill
\begin{minipage}[t]{80mm}\vspace{0pt}\begin{tikzpicture}\begin{axis}[width=72mm,height=42mm,ybar,bar width=1.6pt,xmin=-1,xmax=17,ymin=0,ymax=2.268,axis y line*=left,axis x line*=bottom,xlabel={\scriptsize Halbzüge bis Matt},ylabel={\scriptsize Stellungen (Mrd.)},tick label style={font=\tiny},label style={font=\scriptsize},enlarge x limits=false,ylabel shift=-3pt,xlabel shift=-3pt]\addplot[fill=ink!55,draw=none] coordinates {(0,0.0003237) (1,0.007036) (2,0.02199) (3,0.179) (4,0.4928) (5,1.545) (6,2.061) (7,1.774) (8,1.813) (9,1.292) (10,1.26) (11,0.7129) (12,0.548) (13,0.07189) (14,0.01874) (15,0.0001253) (16,2.376e-06)};\draw[densely dotted] (axis cs:14,0) -- (axis cs:14,2.268);\end{axis}\begin{axis}[width=72mm,height=42mm,axis y line*=right,axis x line=none,ymode=log,xmin=-1,xmax=17,ylabel={\scriptsize eindeutig (log)},tick label style={font=\tiny},label style={font=\scriptsize},ylabel shift=-3pt]\addplot[mark=*,mark size=0.7pt,draw=black,thin] coordinates {(0,323748) (1,2627246) (2,7346424) (3,64350899) (4,99873903) (5,106411842) (6,67817590) (7,42660503) (8,28413212) (9,15145128) (10,7341417) (11,2119937) (12,478447) (13,3146) (14,57)};\end{axis}\end{tikzpicture}\end{minipage}
\vspace{3mm}

\noindent\begin{minipage}[t]{60mm}\vspace{0pt}\sbox\hmboard{\chessboard[setfen=7k/6B1/8/8/n7/P7/P7/K7 b - - 0 1]}\usebox\hmboard\\[1mm]
\makebox[\wd\hmboard]{\sffamily\hspace{3mm}\textbf{h\#7}\hfill(4+2)}\end{minipage}\hfill
\begin{minipage}[t]{52mm}\vspace{0pt}
\textcolor{ink}{\sffamily\textsc{Lösung}}\par\smallskip\small
\begin{tabular}{@{}r@{\enspace}l@{\enspace}l@{}}
\textbf{1.} & \symking{}h7 & \symbishop{}h6\\
\textbf{2.} & \symknight{}b6 & a4\\
\textbf{3.} & \symking{}g6 & a5\\
\textbf{4.} & \symking{}f5 & axb6\\
\textbf{5.} & \symking{}e4 & b7\\
\textbf{6.} & \symking{}d3 & b8=\symqueen{}\\
\textbf{7.} & \symking{}c2 & \symqueen{}b3\#\\
\end{tabular}\par\vspace{3mm}
\qrcode[height=12mm]{https://helpman.komtera.lt/?fen=7k/6B1/8/8/n7/P7/P7/K7&moves=7}\\{\sffamily\tiny\color{ink}scannen zum Lösen}
\end{minipage}
\vspace{2mm}

{\raggedright\noindent\textcolor{ink}{\sffamily\textsc{Themen}}\enspace \themetag{Umwandlung} \themetag{ohne Schach} \themetag{Umnov}\par\smallskip{\footnotesize\textit{Hinweis:} die Partei am Zug steht im Diagramm im Schach.}\par}

\clearpage
\phantomsection\label{hm:KBPPvkq}
\begin{tcolorbox}[enhanced,colback=ink,colframe=ink,arc=0pt,boxrule=0pt,left=4mm,right=4mm,top=2.5mm,bottom=2.5mm]
\color{white}\sffamily{\LARGE\bfseries KBPPvkq}\hfill{\Large h\#7}\\[0.5mm]
{\small 6 Steine\quad Schwarz am Zug}\hfill{\small Nr.~503\qquad\href{https://helpman.komtera.lt/?fen=8/8/B7/8/K7/P7/P4q2/6k1&moves=7}{\color{white}\bfseries lösen\,\ensuremath{\nearrow}}}
\end{tcolorbox}
\vspace{2mm}
\noindent\begin{minipage}[t]{36mm}\vspace{0pt}\sffamily\small\raggedright
\hmfigure{längste eindeutige}{{\Large h\#7}}\hmfigure{eindeutig bei dieser Länge}{{\Large 89} Stellungen}\hmfigure{längstes Matt}{{\Large h\#8.5}}\hmfigure{in der Klasse}{34,1 Milliarden Stellungen,\\ 10,2 Milliarden mit Hilfsmatt}
\end{minipage}\hfill
\begin{minipage}[t]{80mm}\vspace{0pt}\begin{tikzpicture}\begin{axis}[width=72mm,height=42mm,ybar,bar width=1.6pt,xmin=-1,xmax=18,ymin=0,ymax=2.097,axis y line*=left,axis x line*=bottom,xlabel={\scriptsize Halbzüge bis Matt},ylabel={\scriptsize Stellungen (Mrd.)},tick label style={font=\tiny},label style={font=\scriptsize},enlarge x limits=false,ylabel shift=-3pt,xlabel shift=-3pt]\addplot[fill=ink!55,draw=none] coordinates {(0,0.0001522) (1,0.003128) (2,0.01565) (3,0.1531) (4,0.4894) (5,1.492) (6,1.674) (7,1.906) (8,1.305) (9,1.284) (10,0.8175) (11,0.6447) (12,0.2889) (13,0.09118) (14,0.004066) (15,0.001219) (16,4.8e-07) (17,5.66e-07)};\draw[densely dotted] (axis cs:14,0) -- (axis cs:14,2.097);\end{axis}\begin{axis}[width=72mm,height=42mm,axis y line*=right,axis x line=none,ymode=log,xmin=-1,xmax=18,ylabel={\scriptsize eindeutig (log)},tick label style={font=\tiny},label style={font=\scriptsize},ylabel shift=-3pt]\addplot[mark=*,mark size=0.7pt,draw=black,thin] coordinates {(0,152173) (1,1465025) (2,6322393) (3,58858676) (4,101263055) (5,112372293) (6,45674304) (7,37474819) (8,12856782) (9,9528546) (10,2986036) (11,1084360) (12,85172) (13,23948) (14,89)};\end{axis}\end{tikzpicture}\end{minipage}
\vspace{3mm}

\noindent\begin{minipage}[t]{60mm}\vspace{0pt}\sbox\hmboard{\chessboard[setfen=8/8/B7/8/K7/P7/P4q2/6k1 b - - 0 1]}\usebox\hmboard\\[1mm]
\makebox[\wd\hmboard]{\sffamily\hspace{3mm}\textbf{h\#7}\hfill(4+2)}\end{minipage}\hfill
\begin{minipage}[t]{52mm}\vspace{0pt}
\textcolor{ink}{\sffamily\textsc{Lösung}}\par\smallskip\small
\begin{tabular}{@{}r@{\enspace}l@{\enspace}l@{}}
\textbf{1.} & \symking{}g2 & \symking{}b4\\
\textbf{2.} & \symking{}f3 & a4\\
\textbf{3.} & \symking{}e4 & a5\\
\textbf{4.} & \symqueen{}b6+ & axb6\\
\textbf{5.} & \symking{}d5 & b7\\
\textbf{6.} & \symking{}c6 & \symbishop{}b5+\\
\textbf{7.} & \symking{}b6 & b8=\symqueen{}\#\\
\end{tabular}\par\vspace{3mm}
\qrcode[height=12mm]{https://helpman.komtera.lt/?fen=8/8/B7/8/K7/P7/P4q2/6k1&moves=7}\\{\sffamily\tiny\color{ink}scannen zum Lösen}
\end{minipage}
\vspace{2mm}

{\raggedright\noindent\textcolor{ink}{\sffamily\textsc{Themen}}\enspace \themetag{reines Matt} \themetag{Mustermatt} \themetag{Umwandlung} \themetag{Kniest}\par}

\clearpage
\phantomsection\label{hm:KBPPvkr}
\begin{tcolorbox}[enhanced,colback=ink,colframe=ink,arc=0pt,boxrule=0pt,left=4mm,right=4mm,top=2.5mm,bottom=2.5mm]
\color{white}\sffamily{\LARGE\bfseries KBPPvkr}\hfill{\Large h\#7}\\[0.5mm]
{\small 6 Steine\quad Schwarz am Zug}\hfill{\small Nr.~504\qquad\href{https://helpman.komtera.lt/?fen=8/8/8/8/4B3/4P3/4P3/rk1K4&moves=7}{\color{white}\bfseries lösen\,\ensuremath{\nearrow}}}
\end{tcolorbox}
\vspace{2mm}
\noindent\begin{minipage}[t]{36mm}\vspace{0pt}\sffamily\small\raggedright
\hmfigure{längste eindeutige}{{\Large h\#7}}\hmfigure{eindeutig bei dieser Länge}{{\Large 56} Stellungen}\hmfigure{längstes Matt}{{\Large h\#8.5}}\hmfigure{in der Klasse}{34,1 Milliarden Stellungen,\\ 11 Milliarden mit Hilfsmatt}
\end{minipage}\hfill
\begin{minipage}[t]{80mm}\vspace{0pt}\begin{tikzpicture}\begin{axis}[width=72mm,height=42mm,ybar,bar width=1.6pt,xmin=-1,xmax=18,ymin=0,ymax=2.247,axis y line*=left,axis x line*=bottom,xlabel={\scriptsize Halbzüge bis Matt},ylabel={\scriptsize Stellungen (Mrd.)},tick label style={font=\tiny},label style={font=\scriptsize},enlarge x limits=false,ylabel shift=-3pt,xlabel shift=-3pt]\addplot[fill=ink!55,draw=none] coordinates {(0,0.0002154) (1,0.005109) (2,0.02233) (3,0.1917) (4,0.6498) (5,1.702) (6,2.043) (7,1.877) (8,1.496) (9,1.203) (10,0.9323) (11,0.5514) (12,0.2985) (13,0.04744) (14,0.002695) (15,0.0001181) (16,1.7e-07) (17,7.3e-08)};\draw[densely dotted] (axis cs:14,0) -- (axis cs:14,2.247);\end{axis}\begin{axis}[width=72mm,height=42mm,axis y line*=right,axis x line=none,ymode=log,xmin=-1,xmax=18,ylabel={\scriptsize eindeutig (log)},tick label style={font=\tiny},label style={font=\scriptsize},ylabel shift=-3pt]\addplot[mark=*,mark size=0.7pt,draw=black,thin] coordinates {(0,215436) (1,2248445) (2,9130081) (3,74116068) (4,146238849) (5,132776504) (6,62131881) (7,37667220) (8,15024039) (9,8395466) (10,3477566) (11,583178) (12,45334) (13,2219) (14,56)};\end{axis}\end{tikzpicture}\end{minipage}
\vspace{3mm}

\noindent\begin{minipage}[t]{60mm}\vspace{0pt}\sbox\hmboard{\chessboard[setfen=8/8/8/8/4B3/4P3/4P3/rk1K4 b - - 0 1]}\usebox\hmboard\\[1mm]
\makebox[\wd\hmboard]{\sffamily\hspace{3mm}\textbf{h\#7}\hfill(4+2)}\end{minipage}\hfill
\begin{minipage}[t]{52mm}\vspace{0pt}
\textcolor{ink}{\sffamily\textsc{Lösung}}\par\smallskip\small
\begin{tabular}{@{}r@{\enspace}l@{\enspace}l@{}}
\textbf{1.} & \symking{}b2+ & \symbishop{}b1\\
\textbf{2.} & \symking{}c3 & e4\\
\textbf{3.} & \symking{}d4 & e5\\
\textbf{4.} & \symking{}e3 & e6\\
\textbf{5.} & \symking{}f4 & e7\\
\textbf{6.} & \symking{}g5 & e8=\symqueen{}\\
\textbf{7.} & \symking{}h6 & \symqueen{}g6\#\\
\end{tabular}\par\vspace{3mm}
\qrcode[height=12mm]{https://helpman.komtera.lt/?fen=8/8/8/8/4B3/4P3/4P3/rk1K4&moves=7}\\{\sffamily\tiny\color{ink}scannen zum Lösen}
\end{minipage}
\vspace{2mm}

{\raggedright\noindent\textcolor{ink}{\sffamily\textsc{Themen}}\enspace \themetag{reines Matt} \themetag{Mustermatt} \themetag{Umwandlung} \themetag{ein Stein zieht} \themetag{ohne Schlag} \themetag{Umnov}\par\smallskip{\footnotesize\textit{Hinweis:} die Partei am Zug steht im Diagramm im Schach.}\par}

\clearpage
\phantomsection\label{hm:KBPvkbb}
\begin{tcolorbox}[enhanced,colback=ink,colframe=ink,arc=0pt,boxrule=0pt,left=4mm,right=4mm,top=2.5mm,bottom=2.5mm]
\color{white}\sffamily{\LARGE\bfseries KBPvkbb}\hfill{\Large h\#7}\\[0.5mm]
{\small 6 Steine\quad Schwarz am Zug}\hfill{\small Nr.~505\qquad\href{https://helpman.komtera.lt/?fen=7k/8/5B2/4b3/8/2b5/1P6/K7&moves=7}{\color{white}\bfseries lösen\,\ensuremath{\nearrow}}}
\end{tcolorbox}
\vspace{2mm}
\noindent\begin{minipage}[t]{36mm}\vspace{0pt}\sffamily\small\raggedright
\hmfigure{längste eindeutige}{{\Large h\#7}}\hmfigure{eindeutig bei dieser Länge}{{\Large 388} Stellungen}\hmfigure{längstes Matt}{{\Large h\#10}}\hmfigure{in der Klasse}{45,4 Milliarden Stellungen,\\ 14,6 Milliarden mit Hilfsmatt}
\end{minipage}\hfill
\begin{minipage}[t]{80mm}\vspace{0pt}\begin{tikzpicture}\begin{axis}[width=72mm,height=42mm,ybar,bar width=1.6pt,xmin=-1,xmax=21,ymin=0,ymax=2.466,axis y line*=left,axis x line*=bottom,xlabel={\scriptsize Halbzüge bis Matt},ylabel={\scriptsize Stellungen (Mrd.)},tick label style={font=\tiny},label style={font=\scriptsize},enlarge x limits=false,ylabel shift=-3pt,xlabel shift=-3pt]\addplot[fill=ink!55,draw=none] coordinates {(0,0.0002487) (1,0.002638) (2,0.01279) (3,0.1129) (4,0.3837) (5,1.285) (6,1.755) (7,2.216) (8,2.054) (9,2.242) (10,1.718) (11,1.568) (12,0.9064) (13,0.262) (14,0.04038) (15,0.001529) (16,3.51e-05) (17,1.6e-08) (18,1.657e-05) (20,5.798e-06)};\draw[densely dotted] (axis cs:14,0) -- (axis cs:14,2.466);\end{axis}\begin{axis}[width=72mm,height=42mm,axis y line*=right,axis x line=none,ymode=log,xmin=-1,xmax=21,ylabel={\scriptsize eindeutig (log)},tick label style={font=\tiny},label style={font=\scriptsize},ylabel shift=-3pt]\addplot[mark=*,mark size=0.7pt,draw=black,thin] coordinates {(0,248737) (1,1543321) (2,7057839) (3,44327490) (4,83437801) (5,102901250) (6,59317806) (7,49682304) (8,22736974) (9,17954203) (10,6895901) (11,2859597) (12,504657) (13,22478) (14,388)};\end{axis}\end{tikzpicture}\end{minipage}
\vspace{3mm}

\noindent\begin{minipage}[t]{60mm}\vspace{0pt}\sbox\hmboard{\chessboard[setfen=7k/8/5B2/4b3/8/2b5/1P6/K7 b - - 0 1]}\usebox\hmboard\\[1mm]
\makebox[\wd\hmboard]{\sffamily\hspace{3mm}\textbf{h\#7}\hfill(3+3)}\end{minipage}\hfill
\begin{minipage}[t]{52mm}\vspace{0pt}
\textcolor{ink}{\sffamily\textsc{Lösung}}\par\smallskip\small
\begin{tabular}{@{}r@{\enspace}l@{\enspace}l@{}}
\textbf{1.} & \symking{}h7 & \symbishop{}xe5\\
\textbf{2.} & \symbishop{}d2 & b4\\
\textbf{3.} & \symking{}g6 & b5\\
\textbf{4.} & \symking{}f5 & b6\\
\textbf{5.} & \symking{}e4 & b7\\
\textbf{6.} & \symking{}d3 & b8=\symqueen{}\\
\textbf{7.} & \symking{}c2 & \symqueen{}b1\#\\
\end{tabular}\par\vspace{3mm}
\qrcode[height=12mm]{https://helpman.komtera.lt/?fen=7k/8/5B2/4b3/8/2b5/1P6/K7&moves=7}\\{\sffamily\tiny\color{ink}scannen zum Lösen}
\end{minipage}
\vspace{2mm}

{\raggedright\noindent\textcolor{ink}{\sffamily\textsc{Themen}}\enspace \themetag{Umwandlung} \themetag{Excelsior} \themetag{Selbstblock} \themetag{ohne Schach}\par\smallskip{\footnotesize\textit{Hinweis:} die Partei am Zug steht im Diagramm im Schach.}\par}

\clearpage
\phantomsection\label{hm:KBPvkrb}
\begin{tcolorbox}[enhanced,colback=ink,colframe=ink,arc=0pt,boxrule=0pt,left=4mm,right=4mm,top=2.5mm,bottom=2.5mm]
\color{white}\sffamily{\LARGE\bfseries KBPvkrb}\hfill{\Large h\#7}\\[0.5mm]
{\small 6 Steine\quad Schwarz am Zug}\hfill{\small Nr.~506\qquad\href{https://helpman.komtera.lt/?fen=8/8/8/1b6/8/K6B/P7/1r3k2&moves=7}{\color{white}\bfseries lösen\,\ensuremath{\nearrow}}}
\end{tcolorbox}
\vspace{2mm}
\noindent\begin{minipage}[t]{36mm}\vspace{0pt}\sffamily\small\raggedright
\hmfigure{längste eindeutige}{{\Large h\#7}}\hmfigure{eindeutig bei dieser Länge}{{\Large 162} Stellungen}\hmfigure{längstes Matt}{{\Large h\#9.5}}\hmfigure{in der Klasse}{45,4 Milliarden Stellungen,\\ 28,1 Milliarden mit Hilfsmatt}
\end{minipage}\hfill
\begin{minipage}[t]{80mm}\vspace{0pt}\begin{tikzpicture}\begin{axis}[width=72mm,height=42mm,ybar,bar width=1.6pt,xmin=-1,xmax=20,ymin=0,ymax=5.106,axis y line*=left,axis x line*=bottom,xlabel={\scriptsize Halbzüge bis Matt},ylabel={\scriptsize Stellungen (Mrd.)},tick label style={font=\tiny},label style={font=\scriptsize},enlarge x limits=false,ylabel shift=-3pt,xlabel shift=-3pt]\addplot[fill=ink!55,draw=none] coordinates {(0,0.0002236) (1,0.004017) (2,0.02064) (3,0.2103) (4,0.7703) (5,2.543) (6,3.254) (7,4.326) (8,3.754) (9,4.642) (10,3.459) (11,3.196) (12,1.473) (13,0.4478) (14,0.02885) (15,0.001771) (16,6.094e-06) (17,8.1e-08) (18,1e-09) (19,4e-09)};\draw[densely dotted] (axis cs:14,0) -- (axis cs:14,5.106);\end{axis}\begin{axis}[width=72mm,height=42mm,axis y line*=right,axis x line=none,ymode=log,xmin=-1,xmax=20,ylabel={\scriptsize eindeutig (log)},tick label style={font=\tiny},label style={font=\scriptsize},ylabel shift=-3pt]\addplot[mark=*,mark size=0.7pt,draw=black,thin] coordinates {(0,223640) (1,2159122) (2,9954728) (3,85151147) (4,178189394) (5,217392509) (6,116316613) (7,102249328) (8,51226850) (9,42069777) (10,14512991) (11,4446485) (12,334051) (13,16732) (14,162)};\end{axis}\end{tikzpicture}\end{minipage}
\vspace{3mm}

\noindent\begin{minipage}[t]{60mm}\vspace{0pt}\sbox\hmboard{\chessboard[setfen=8/8/8/1b6/8/K6B/P7/1r3k2 b - - 0 1]}\usebox\hmboard\\[1mm]
\makebox[\wd\hmboard]{\sffamily\hspace{3mm}\textbf{h\#7}\hfill(3+3)}\end{minipage}\hfill
\begin{minipage}[t]{52mm}\vspace{0pt}
\textcolor{ink}{\sffamily\textsc{Lösung}}\par\smallskip\small
\begin{tabular}{@{}r@{\enspace}l@{\enspace}l@{}}
\textbf{1.} & \symking{}e2 & \symbishop{}c8\\
\textbf{2.} & \symrook{}a1 & \symking{}b2\\
\textbf{3.} & \symking{}d3 & a4\\
\textbf{4.} & \symking{}c4 & axb5\\
\textbf{5.} & \symrook{}a5 & b6\\
\textbf{6.} & \symking{}b5 & b7\\
\textbf{7.} & \symking{}a6 & b8=\symqueen{}\#\\
\end{tabular}\par\vspace{3mm}
\qrcode[height=12mm]{https://helpman.komtera.lt/?fen=8/8/8/1b6/8/K6B/P7/1r3k2&moves=7}\\{\sffamily\tiny\color{ink}scannen zum Lösen}
\end{minipage}
\vspace{2mm}

{\raggedright\noindent\textcolor{ink}{\sffamily\textsc{Themen}}\enspace \themetag{Umwandlung} \themetag{Excelsior} \themetag{Selbstblock} \themetag{ohne Schach} \themetag{Umnov}\par\smallskip{\footnotesize\textit{Hinweis:} die Partei am Zug steht im Diagramm im Schach.}\par}

\clearpage
\phantomsection\label{hm:KBPvkrr}
\begin{tcolorbox}[enhanced,colback=ink,colframe=ink,arc=0pt,boxrule=0pt,left=4mm,right=4mm,top=2.5mm,bottom=2.5mm]
\color{white}\sffamily{\LARGE\bfseries KBPvkrr}\hfill{\Large h\#7}\\[0.5mm]
{\small 6 Steine\quad Schwarz am Zug}\hfill{\small Nr.~507\qquad\href{https://helpman.komtera.lt/?fen=8/8/8/1r6/8/K7/P2rB3/5k2&moves=7}{\color{white}\bfseries lösen\,\ensuremath{\nearrow}}}
\end{tcolorbox}
\vspace{2mm}
\noindent\begin{minipage}[t]{36mm}\vspace{0pt}\sffamily\small\raggedright
\hmfigure{längste eindeutige}{{\Large h\#7}}\hmfigure{eindeutig bei dieser Länge}{{\Large 336} Stellungen}\hmfigure{längstes Matt}{{\Large h\#8.5}}\hmfigure{in der Klasse}{45,4 Milliarden Stellungen,\\ 13,4 Milliarden mit Hilfsmatt}
\end{minipage}\hfill
\begin{minipage}[t]{80mm}\vspace{0pt}\begin{tikzpicture}\begin{axis}[width=72mm,height=42mm,ybar,bar width=1.6pt,xmin=-1,xmax=18,ymin=0,ymax=2.134,axis y line*=left,axis x line*=bottom,xlabel={\scriptsize Halbzüge bis Matt},ylabel={\scriptsize Stellungen (Mrd.)},tick label style={font=\tiny},label style={font=\scriptsize},enlarge x limits=false,ylabel shift=-3pt,xlabel shift=-3pt]\addplot[fill=ink!55,draw=none] coordinates {(0,5.736e-05) (1,0.001568) (2,0.009) (3,0.09775) (4,0.3636) (5,1.21) (6,1.385) (7,1.937) (8,1.365) (9,1.94) (10,1.516) (11,1.934) (12,1.083) (13,0.5213) (14,0.0288) (15,0.003264) (16,1.962e-05) (17,2.12e-07)};\draw[densely dotted] (axis cs:14,0) -- (axis cs:14,2.134);\end{axis}\begin{axis}[width=72mm,height=42mm,axis y line*=right,axis x line=none,ymode=log,xmin=-1,xmax=18,ylabel={\scriptsize eindeutig (log)},tick label style={font=\tiny},label style={font=\scriptsize},ylabel shift=-3pt]\addplot[mark=*,mark size=0.7pt,draw=black,thin] coordinates {(0,57363) (1,823969) (2,4292179) (3,40569996) (4,81456482) (5,101180031) (6,42714749) (7,41556577) (8,15326628) (9,18447886) (10,9975415) (11,4189528) (12,302183) (13,29313) (14,336)};\end{axis}\end{tikzpicture}\end{minipage}
\vspace{3mm}

\noindent\begin{minipage}[t]{60mm}\vspace{0pt}\sbox\hmboard{\chessboard[setfen=8/8/8/1r6/8/K7/P2rB3/5k2 b - - 0 1]}\usebox\hmboard\\[1mm]
\makebox[\wd\hmboard]{\sffamily\hspace{3mm}\textbf{h\#7}\hfill(3+3)}\end{minipage}\hfill
\begin{minipage}[t]{52mm}\vspace{0pt}
\textcolor{ink}{\sffamily\textsc{Lösung}}\par\smallskip\small
\begin{tabular}{@{}r@{\enspace}l@{\enspace}l@{}}
\textbf{1.} & \symking{}e1 & \symking{}a4\\
\textbf{2.} & \symrook{}d4+ & \symking{}xb5\\
\textbf{3.} & \symking{}d2 & \symking{}c6\\
\textbf{4.} & \symking{}c3 & \symbishop{}b5\\
\textbf{5.} & \symking{}b4 & a3+\\
\textbf{6.} & \symking{}a5 & \symking{}c5\\
\textbf{7.} & \symrook{}b4 & axb4\#\\
\end{tabular}\par\vspace{3mm}
\qrcode[height=12mm]{https://helpman.komtera.lt/?fen=8/8/8/1r6/8/K7/P2rB3/5k2&moves=7}\\{\sffamily\tiny\color{ink}scannen zum Lösen}
\end{minipage}
\vspace{2mm}

{\raggedright\noindent\textcolor{ink}{\sffamily\textsc{Themen}}\enspace \themetag{reines Matt} \themetag{Mustermatt} \themetag{Idealmatt}\par\smallskip{\footnotesize\textit{Hinweis:} die Partei am Zug steht im Diagramm im Schach.}\par}

\clearpage
\phantomsection\label{hm:KNNNPvk}
\begin{tcolorbox}[enhanced,colback=ink,colframe=ink,arc=0pt,boxrule=0pt,left=4mm,right=4mm,top=2.5mm,bottom=2.5mm]
\color{white}\sffamily{\LARGE\bfseries KNNNPvk}\hfill{\Large h\#7}\\[0.5mm]
{\small 6 Steine\quad Schwarz am Zug}\hfill{\small Nr.~508\qquad\href{https://helpman.komtera.lt/?fen=5NkN/8/4N3/8/8/K7/P7/8&moves=7}{\color{white}\bfseries lösen\,\ensuremath{\nearrow}}}
\end{tcolorbox}
\vspace{2mm}
\noindent\begin{minipage}[t]{36mm}\vspace{0pt}\sffamily\small\raggedright
\hmfigure{längste eindeutige}{{\Large h\#7}}\hmfigure{eindeutig bei dieser Länge}{{\Large 4} Stellungen}\hmfigure{längstes Matt}{{\Large h\#8}}\hmfigure{in der Klasse}{45,4 Milliarden Stellungen,\\ 5,29 Milliarden mit Hilfsmatt}
\end{minipage}\hfill
\begin{minipage}[t]{80mm}\vspace{0pt}\begin{tikzpicture}\begin{axis}[width=72mm,height=42mm,ybar,bar width=1.6pt,xmin=-1,xmax=17,ymin=0,ymax=1.466,axis y line*=left,axis x line*=bottom,xlabel={\scriptsize Halbzüge bis Matt},ylabel={\scriptsize Stellungen (Mrd.)},tick label style={font=\tiny},label style={font=\scriptsize},enlarge x limits=false,ylabel shift=-3pt,xlabel shift=-3pt]\addplot[fill=ink!55,draw=none] coordinates {(0,0.0004393) (1,0.003818) (2,0.01007) (3,0.07137) (4,0.2447) (5,0.6292) (6,1.092) (7,1.033) (8,1.332) (9,0.5009) (10,0.3277) (11,0.03619) (12,0.01135) (13,0.0003987) (14,0.0001056) (15,1.29e-07) (16,2.1e-08)};\draw[densely dotted] (axis cs:14,0) -- (axis cs:14,1.466);\end{axis}\begin{axis}[width=72mm,height=42mm,axis y line*=right,axis x line=none,ymode=log,xmin=-1,xmax=17,ylabel={\scriptsize eindeutig (log)},tick label style={font=\tiny},label style={font=\scriptsize},ylabel shift=-3pt]\addplot[mark=*,mark size=0.7pt,draw=black,thin] coordinates {(0,439270) (1,2040437) (2,5848764) (3,38252417) (4,82893478) (5,64199563) (6,61546551) (7,15199406) (8,8373554) (9,424580) (10,136177) (11,5512) (12,980) (13,2) (14,4)};\end{axis}\end{tikzpicture}\end{minipage}
\vspace{3mm}

\noindent\begin{minipage}[t]{60mm}\vspace{0pt}\sbox\hmboard{\chessboard[setfen=5NkN/8/4N3/8/8/K7/P7/8 b - - 0 1]}\usebox\hmboard\\[1mm]
\makebox[\wd\hmboard]{\sffamily\hspace{3mm}\textbf{h\#7}\hfill(5+1)}\end{minipage}\hfill
\begin{minipage}[t]{52mm}\vspace{0pt}
\textcolor{ink}{\sffamily\textsc{Lösung}}\par\smallskip\small
\begin{tabular}{@{}r@{\enspace}l@{\enspace}l@{}}
\textbf{1.} & \symking{}xh8 & \symking{}b4\\
\textbf{2.} & \symking{}g8 & a4\\
\textbf{3.} & \symking{}f7 & a5\\
\textbf{4.} & \symking{}e7 & a6\\
\textbf{5.} & \symking{}d6 & a7\\
\textbf{6.} & \symking{}c6 & a8=\symqueen{}+\\
\textbf{7.} & \symking{}b6 & \symknight{}d7\#\\
\end{tabular}\par\vspace{3mm}
\qrcode[height=12mm]{https://helpman.komtera.lt/?fen=5NkN/8/4N3/8/8/K7/P7/8&moves=7}\\{\sffamily\tiny\color{ink}scannen zum Lösen}
\end{minipage}
\vspace{2mm}

{\raggedright\noindent\textcolor{ink}{\sffamily\textsc{Themen}}\enspace \themetag{Spiegelmatt} \themetag{Umwandlung} \themetag{Excelsior} \themetag{Switchback} \themetag{ein Stein zieht}\par\smallskip{\footnotesize\textit{Hinweis:} die Lösung beginnt mit einem Schlag.}\par}

\clearpage
\phantomsection\label{hm:KNNNvkp}
\begin{tcolorbox}[enhanced,colback=ink,colframe=ink,arc=0pt,boxrule=0pt,left=4mm,right=4mm,top=2.5mm,bottom=2.5mm]
\color{white}\sffamily{\LARGE\bfseries KNNNvkp}\hfill{\Large h\#7}\\[0.5mm]
{\small 6 Steine\quad Schwarz am Zug}\hfill{\small Nr.~509\qquad\href{https://helpman.komtera.lt/?fen=7N/4N2k/7p/7N/8/8/8/K7&moves=7}{\color{white}\bfseries lösen\,\ensuremath{\nearrow}}}
\end{tcolorbox}
\vspace{2mm}
\noindent\begin{minipage}[t]{36mm}\vspace{0pt}\sffamily\small\raggedright
\hmfigure{längste eindeutige}{{\Large h\#7}}\hmfigure{eindeutig bei dieser Länge}{{\Large 1} Stellung}\hmfigure{längstes Matt}{{\Large h\#8}}\hmfigure{in der Klasse}{45,4 Milliarden Stellungen,\\ 5,28 Milliarden mit Hilfsmatt}
\end{minipage}\hfill
\begin{minipage}[t]{80mm}\vspace{0pt}\begin{tikzpicture}\begin{axis}[width=72mm,height=42mm,ybar,bar width=1.6pt,xmin=-1,xmax=17,ymin=0,ymax=1.667,axis y line*=left,axis x line*=bottom,xlabel={\scriptsize Halbzüge bis Matt},ylabel={\scriptsize Stellungen (Mrd.)},tick label style={font=\tiny},label style={font=\scriptsize},enlarge x limits=false,ylabel shift=-3pt,xlabel shift=-3pt]\addplot[fill=ink!55,draw=none] coordinates {(0,0.0003204) (1,0.001427) (2,0.004689) (3,0.02892) (4,0.09984) (5,0.3391) (6,0.81) (7,1.085) (8,1.515) (9,0.7774) (10,0.4826) (11,0.1016) (12,0.02572) (13,0.003308) (14,0.0003084) (15,1.886e-05) (16,4.24e-07)};\draw[densely dotted] (axis cs:14,0) -- (axis cs:14,1.667);\end{axis}\begin{axis}[width=72mm,height=42mm,axis y line*=right,axis x line=none,ymode=log,xmin=-1,xmax=17,ylabel={\scriptsize eindeutig (log)},tick label style={font=\tiny},label style={font=\scriptsize},ylabel shift=-3pt]\addplot[mark=*,mark size=0.7pt,draw=black,thin] coordinates {(0,320430) (1,1263678) (2,3833275) (3,16054635) (4,34338105) (5,32032189) (6,39445696) (7,10325897) (8,6757453) (9,447824) (10,118919) (11,3898) (12,576) (13,18) (14,1)};\end{axis}\end{tikzpicture}\end{minipage}
\vspace{3mm}

\noindent\begin{minipage}[t]{60mm}\vspace{0pt}\sbox\hmboard{\chessboard[setfen=7N/4N2k/7p/7N/8/8/8/K7 b - - 0 1]}\usebox\hmboard\\[1mm]
\makebox[\wd\hmboard]{\sffamily\hspace{3mm}\textbf{h\#7}\hfill(4+2)}\end{minipage}\hfill
\begin{minipage}[t]{52mm}\vspace{0pt}
\textcolor{ink}{\sffamily\textsc{Lösung}}\par\smallskip\small
\begin{tabular}{@{}r@{\enspace}l@{\enspace}l@{}}
\textbf{1.} & \symking{}xh8 & \symknight{}f4\\
\textbf{2.} & h5 & \symking{}b2\\
\textbf{3.} & h4 & \symking{}c3\\
\textbf{4.} & h3 & \symking{}d4\\
\textbf{5.} & h2 & \symking{}e5\\
\textbf{6.} & h1=\symrook{} & \symking{}f6\\
\textbf{7.} & \symrook{}h7 & \symknight{}fg6\#\\
\end{tabular}\par\vspace{3mm}
\qrcode[height=12mm]{https://helpman.komtera.lt/?fen=7N/4N2k/7p/7N/8/8/8/K7&moves=7}\\{\sffamily\tiny\color{ink}scannen zum Lösen}
\end{minipage}
\vspace{2mm}

{\raggedright\noindent\textcolor{ink}{\sffamily\textsc{Themen}}\enspace \themetag{reines Matt} \themetag{Mustermatt} \themetag{Idealmatt} \themetag{Umwandlung} \themetag{Unterverwandlung} \themetag{Selbstblock} \themetag{Kniest} \themetag{ohne Schach} \themetag{Umnov}\par\smallskip{\footnotesize\textit{Hinweis:} die Lösung beginnt mit einem Schlag.}\par}

\clearpage
\phantomsection\label{hm:KNNPPvk}
\begin{tcolorbox}[enhanced,colback=ink,colframe=ink,arc=0pt,boxrule=0pt,left=4mm,right=4mm,top=2.5mm,bottom=2.5mm]
\color{white}\sffamily{\LARGE\bfseries KNNPPvk}\hfill{\Large h\#7}\\[0.5mm]
{\small 6 Steine\quad Schwarz am Zug}\hfill{\small Nr.~510\qquad\href{https://helpman.komtera.lt/?fen=8/8/8/8/4N3/K3P3/4P3/kN6&moves=7}{\color{white}\bfseries lösen\,\ensuremath{\nearrow}}}
\end{tcolorbox}
\vspace{2mm}
\noindent\begin{minipage}[t]{36mm}\vspace{0pt}\sffamily\small\raggedright
\hmfigure{längste eindeutige}{{\Large h\#7}}\hmfigure{eindeutig bei dieser Länge}{{\Large 2} Stellungen}\hmfigure{längstes Matt}{{\Large h\#8}}\hmfigure{in der Klasse}{34,1 Milliarden Stellungen,\\ 6,08 Milliarden mit Hilfsmatt}
\end{minipage}\hfill
\begin{minipage}[t]{80mm}\vspace{0pt}\begin{tikzpicture}\begin{axis}[width=72mm,height=42mm,ybar,bar width=1.6pt,xmin=-1,xmax=17,ymin=0,ymax=1.335,axis y line*=left,axis x line*=bottom,xlabel={\scriptsize Halbzüge bis Matt},ylabel={\scriptsize Stellungen (Mrd.)},tick label style={font=\tiny},label style={font=\scriptsize},enlarge x limits=false,ylabel shift=-3pt,xlabel shift=-3pt]\addplot[fill=ink!55,draw=none] coordinates {(0,0.0002382) (1,0.005137) (2,0.01257) (3,0.1035) (4,0.3215) (5,0.8411) (6,1.213) (7,1.001) (8,1.2) (9,0.6246) (10,0.5639) (11,0.1215) (12,0.06928) (13,0.00405) (14,0.001115) (15,3.54e-06) (16,3e-08)};\draw[densely dotted] (axis cs:14,0) -- (axis cs:14,1.335);\end{axis}\begin{axis}[width=72mm,height=42mm,axis y line*=right,axis x line=none,ymode=log,xmin=-1,xmax=17,ylabel={\scriptsize eindeutig (log)},tick label style={font=\tiny},label style={font=\scriptsize},ylabel shift=-3pt]\addplot[mark=*,mark size=0.7pt,draw=black,thin] coordinates {(0,238170) (1,1813279) (2,5325229) (3,53388835) (4,98835600) (5,75772201) (6,55262669) (7,24142526) (8,14807580) (9,2462634) (10,1012190) (11,115739) (12,17546) (13,41) (14,2)};\end{axis}\end{tikzpicture}\end{minipage}
\vspace{3mm}

\noindent\begin{minipage}[t]{60mm}\vspace{0pt}\sbox\hmboard{\chessboard[setfen=8/8/8/8/4N3/K3P3/4P3/kN6 b - - 0 1]}\usebox\hmboard\\[1mm]
\makebox[\wd\hmboard]{\sffamily\hspace{3mm}\textbf{h\#7}\hfill(5+1)}\end{minipage}\hfill
\begin{minipage}[t]{52mm}\vspace{0pt}
\textcolor{ink}{\sffamily\textsc{Lösung}}\par\smallskip\small
\begin{tabular}{@{}r@{\enspace}l@{\enspace}l@{}}
\textbf{1.} & \symking{}xb1 & \symknight{}d6\\
\textbf{2.} & \symking{}c2 & e4\\
\textbf{3.} & \symking{}c3 & e5\\
\textbf{4.} & \symking{}d4 & e6\\
\textbf{5.} & \symking{}c5 & e7\\
\textbf{6.} & \symking{}b6 & e8=\symqueen{}\\
\textbf{7.} & \symking{}a5 & \symqueen{}b5\#\\
\end{tabular}\par\vspace{3mm}
\qrcode[height=12mm]{https://helpman.komtera.lt/?fen=8/8/8/8/4N3/K3P3/4P3/kN6&moves=7}\\{\sffamily\tiny\color{ink}scannen zum Lösen}
\end{minipage}
\vspace{2mm}

{\raggedright\noindent\textcolor{ink}{\sffamily\textsc{Themen}}\enspace \themetag{Umwandlung} \themetag{ein Stein zieht} \themetag{ohne Schach}\par\smallskip{\footnotesize\textit{Hinweis:} die Lösung beginnt mit einem Schlag.}\par}

\clearpage
\phantomsection\label{hm:KNNPvkb}
\begin{tcolorbox}[enhanced,colback=ink,colframe=ink,arc=0pt,boxrule=0pt,left=4mm,right=4mm,top=2.5mm,bottom=2.5mm]
\color{white}\sffamily{\LARGE\bfseries KNNPvkb}\hfill{\Large h\#7}\\[0.5mm]
{\small 6 Steine\quad Schwarz am Zug}\hfill{\small Nr.~511\qquad\href{https://helpman.komtera.lt/?fen=8/8/8/8/1K6/2NN4/k2P4/b7&moves=7}{\color{white}\bfseries lösen\,\ensuremath{\nearrow}}}
\end{tcolorbox}
\vspace{2mm}
\noindent\begin{minipage}[t]{36mm}\vspace{0pt}\sffamily\small\raggedright
\hmfigure{längste eindeutige}{{\Large h\#7}}\hmfigure{eindeutig bei dieser Länge}{{\Large 1} Stellung}\hmfigure{längstes Matt}{{\Large h\#8}}\hmfigure{in der Klasse}{45,4 Milliarden Stellungen,\\ 15,4 Milliarden mit Hilfsmatt}
\end{minipage}\hfill
\begin{minipage}[t]{80mm}\vspace{0pt}\begin{tikzpicture}\begin{axis}[width=72mm,height=42mm,ybar,bar width=1.6pt,xmin=-1,xmax=17,ymin=0,ymax=3.766,axis y line*=left,axis x line*=bottom,xlabel={\scriptsize Halbzüge bis Matt},ylabel={\scriptsize Stellungen (Mrd.)},tick label style={font=\tiny},label style={font=\scriptsize},enlarge x limits=false,ylabel shift=-3pt,xlabel shift=-3pt]\addplot[fill=ink!55,draw=none] coordinates {(0,0.0003021) (1,0.00431) (2,0.01535) (3,0.1568) (4,0.4977) (5,1.514) (6,2.43) (7,2.993) (8,3.424) (9,2.33) (10,1.447) (11,0.4253) (12,0.1069) (13,0.01143) (14,0.0008661) (15,2.755e-05) (16,3.8e-08)};\draw[densely dotted] (axis cs:14,0) -- (axis cs:14,3.766);\end{axis}\begin{axis}[width=72mm,height=42mm,axis y line*=right,axis x line=none,ymode=log,xmin=-1,xmax=17,ylabel={\scriptsize eindeutig (log)},tick label style={font=\tiny},label style={font=\scriptsize},ylabel shift=-3pt]\addplot[mark=*,mark size=0.7pt,draw=black,thin] coordinates {(0,302091) (1,2073181) (2,6517009) (3,62967319) (4,122294926) (5,123434365) (6,108313181) (7,53635948) (8,26802020) (9,5286212) (10,937109) (11,138149) (12,7134) (13,258) (14,1)};\end{axis}\end{tikzpicture}\end{minipage}
\vspace{3mm}

\noindent\begin{minipage}[t]{60mm}\vspace{0pt}\sbox\hmboard{\chessboard[setfen=8/8/8/8/1K6/2NN4/k2P4/b7 b - - 0 1]}\usebox\hmboard\\[1mm]
\makebox[\wd\hmboard]{\sffamily\hspace{3mm}\textbf{h\#7}\hfill(4+2)}\end{minipage}\hfill
\begin{minipage}[t]{52mm}\vspace{0pt}
\textcolor{ink}{\sffamily\textsc{Lösung}}\par\smallskip\small
\begin{tabular}{@{}r@{\enspace}l@{\enspace}l@{}}
\textbf{1.} & \symbishop{}xc3+ & \symking{}b5\\
\textbf{2.} & \symbishop{}e5 & \symknight{}c5\\
\textbf{3.} & \symking{}b2 & d4\\
\textbf{4.} & \symking{}c3 & d5\\
\textbf{5.} & \symking{}d4 & d6\\
\textbf{6.} & \symking{}d5 & d7\\
\textbf{7.} & \symking{}d6 & d8=\symqueen{}\#\\
\end{tabular}\par\vspace{3mm}
\qrcode[height=12mm]{https://helpman.komtera.lt/?fen=8/8/8/8/1K6/2NN4/k2P4/b7&moves=7}\\{\sffamily\tiny\color{ink}scannen zum Lösen}
\end{minipage}
\vspace{2mm}

{\raggedright\noindent\textcolor{ink}{\sffamily\textsc{Themen}}\enspace \themetag{Umwandlung} \themetag{Excelsior} \themetag{Selbstblock} \themetag{Umnov}\par\smallskip{\footnotesize\textit{Hinweis:} die Partei am Zug steht im Diagramm im Schach; die Lösung beginnt mit einem Schlag.}\par}

\clearpage
\phantomsection\label{hm:KNNvkbp}
\begin{tcolorbox}[enhanced,colback=ink,colframe=ink,arc=0pt,boxrule=0pt,left=4mm,right=4mm,top=2.5mm,bottom=2.5mm]
\color{white}\sffamily{\LARGE\bfseries KNNvkbp}\hfill{\Large h\#7}\\[0.5mm]
{\small 6 Steine\quad Schwarz am Zug}\hfill{\small Nr.~512\qquad\href{https://helpman.komtera.lt/?fen=2k5/1bp5/1NN5/8/2K5/8/8/8&moves=7}{\color{white}\bfseries lösen\,\ensuremath{\nearrow}}}
\end{tcolorbox}
\vspace{2mm}
\noindent\begin{minipage}[t]{36mm}\vspace{0pt}\sffamily\small\raggedright
\hmfigure{längste eindeutige}{{\Large h\#7}}\hmfigure{eindeutig bei dieser Länge}{{\Large 4} Stellungen}\hmfigure{längstes Matt}{{\Large h\#9}}\hmfigure{in der Klasse}{45,4 Milliarden Stellungen,\\ 15,3 Milliarden mit Hilfsmatt}
\end{minipage}\hfill
\begin{minipage}[t]{80mm}\vspace{0pt}\begin{tikzpicture}\begin{axis}[width=72mm,height=42mm,ybar,bar width=1.6pt,xmin=-1,xmax=19,ymin=0,ymax=4.015,axis y line*=left,axis x line*=bottom,xlabel={\scriptsize Halbzüge bis Matt},ylabel={\scriptsize Stellungen (Mrd.)},tick label style={font=\tiny},label style={font=\scriptsize},enlarge x limits=false,ylabel shift=-3pt,xlabel shift=-3pt]\addplot[fill=ink!55,draw=none] coordinates {(0,0.0001887) (1,0.0009234) (2,0.00358) (3,0.02826) (4,0.1093) (5,0.4791) (6,1.17) (7,2.462) (8,3.65) (9,3.522) (10,2.446) (11,1.079) (12,0.3063) (13,0.06899) (14,0.007223) (15,0.001213) (16,8.018e-06) (17,4e-09) (18,1.397e-06)};\draw[densely dotted] (axis cs:14,0) -- (axis cs:14,4.015);\end{axis}\begin{axis}[width=72mm,height=42mm,axis y line*=right,axis x line=none,ymode=log,xmin=-1,xmax=19,ylabel={\scriptsize eindeutig (log)},tick label style={font=\tiny},label style={font=\scriptsize},ylabel shift=-3pt]\addplot[mark=*,mark size=0.7pt,draw=black,thin] coordinates {(0,188659) (1,876319) (2,2535411) (3,14146999) (4,40851516) (5,56117325) (6,63229375) (7,38499887) (8,25778874) (9,5681220) (10,1219321) (11,123708) (12,15122) (13,1283) (14,4)};\end{axis}\end{tikzpicture}\end{minipage}
\vspace{3mm}

\noindent\begin{minipage}[t]{60mm}\vspace{0pt}\sbox\hmboard{\chessboard[setfen=2k5/1bp5/1NN5/8/2K5/8/8/8 b - - 0 1]}\usebox\hmboard\\[1mm]
\makebox[\wd\hmboard]{\sffamily\hspace{3mm}\textbf{h\#7}\hfill(3+3)}\end{minipage}\hfill
\begin{minipage}[t]{52mm}\vspace{0pt}
\textcolor{ink}{\sffamily\textsc{Lösung}}\par\smallskip\small
\begin{tabular}{@{}r@{\enspace}l@{\enspace}l@{}}
\textbf{1.} & cxb6 & \symknight{}e5\\
\textbf{2.} & \symbishop{}c6 & \symking{}d4\\
\textbf{3.} & \symbishop{}a4 & \symking{}d5\\
\textbf{4.} & \symking{}b7 & \symking{}d6\\
\textbf{5.} & \symking{}a6 & \symking{}c7\\
\textbf{6.} & \symking{}a5 & \symking{}b7\\
\textbf{7.} & b5 & \symknight{}c6\#\\
\end{tabular}\par\vspace{3mm}
\qrcode[height=12mm]{https://helpman.komtera.lt/?fen=2k5/1bp5/1NN5/8/2K5/8/8/8&moves=7}\\{\sffamily\tiny\color{ink}scannen zum Lösen}
\end{minipage}
\vspace{2mm}

{\raggedright\noindent\textcolor{ink}{\sffamily\textsc{Themen}}\enspace \themetag{reines Matt} \themetag{Mustermatt} \themetag{Idealmatt} \themetag{Switchback} \themetag{Selbstblock} \themetag{ohne Schach} \themetag{Umnov}\par\smallskip{\footnotesize\textit{Hinweis:} die Partei am Zug steht im Diagramm im Schach; die Lösung beginnt mit einem Schlag.}\par}

\clearpage
\phantomsection\label{hm:KNNvkqr}
\begin{tcolorbox}[enhanced,colback=ink,colframe=ink,arc=0pt,boxrule=0pt,left=4mm,right=4mm,top=2.5mm,bottom=2.5mm]
\color{white}\sffamily{\LARGE\bfseries KNNvkqr}\hfill{\Large h\#7}\\[0.5mm]
{\small 6 Steine\quad Schwarz am Zug}\hfill{\small Nr.~513\qquad\href{https://helpman.komtera.lt/?fen=5N1N/7k/8/8/8/8/1r1q4/K7&moves=7}{\color{white}\bfseries lösen\,\ensuremath{\nearrow}}}
\end{tcolorbox}
\vspace{2mm}
\noindent\begin{minipage}[t]{36mm}\vspace{0pt}\sffamily\small\raggedright
\hmfigure{längste eindeutige}{{\Large h\#7}}\hmfigure{eindeutig bei dieser Länge}{{\Large 3} Stellungen}\hmfigure{längstes Matt}{{\Large h\#8.5}}\hmfigure{in der Klasse}{15,5 Milliarden Stellungen,\\ 4,12 Milliarden mit Hilfsmatt}
\end{minipage}\hfill
\begin{minipage}[t]{80mm}\vspace{0pt}\begin{tikzpicture}\begin{axis}[width=72mm,height=42mm,ybar,bar width=1.6pt,xmin=-1,xmax=18,ymin=0,ymax=1.267,axis y line*=left,axis x line*=bottom,xlabel={\scriptsize Halbzüge bis Matt},ylabel={\scriptsize Stellungen (Mrd.)},tick label style={font=\tiny},label style={font=\scriptsize},enlarge x limits=false,ylabel shift=-3pt,xlabel shift=-3pt]\addplot[fill=ink!55,draw=none] coordinates {(0,2.418e-05) (1,0.0001206) (2,0.0005928) (3,0.005291) (4,0.02226) (5,0.1166) (6,0.276) (7,0.763) (8,0.7806) (9,1.152) (10,0.4555) (11,0.436) (12,0.06977) (13,0.04116) (14,0.001051) (15,0.0007371) (16,6e-08) (17,3.65e-07)};\draw[densely dotted] (axis cs:14,0) -- (axis cs:14,1.267);\end{axis}\begin{axis}[width=72mm,height=42mm,axis y line*=right,axis x line=none,ymode=log,xmin=-1,xmax=18,ylabel={\scriptsize eindeutig (log)},tick label style={font=\tiny},label style={font=\scriptsize},ylabel shift=-3pt]\addplot[mark=*,mark size=0.7pt,draw=black,thin] coordinates {(0,24175) (1,117626) (2,396930) (3,2509952) (4,7269434) (5,18030104) (6,16982794) (7,22830445) (8,6893837) (9,5860878) (10,861545) (11,516683) (12,19468) (13,5290) (14,3)};\end{axis}\end{tikzpicture}\end{minipage}
\vspace{3mm}

\noindent\begin{minipage}[t]{60mm}\vspace{0pt}\sbox\hmboard{\chessboard[setfen=5N1N/7k/8/8/8/8/1r1q4/K7 b - - 0 1]}\usebox\hmboard\\[1mm]
\makebox[\wd\hmboard]{\sffamily\hspace{3mm}\textbf{h\#7}\hfill(3+3)}\end{minipage}\hfill
\begin{minipage}[t]{52mm}\vspace{0pt}
\textcolor{ink}{\sffamily\textsc{Lösung}}\par\smallskip\small
\begin{tabular}{@{}r@{\enspace}l@{\enspace}l@{}}
\textbf{1.} & \symking{}h6 & \symknight{}hg6\\
\textbf{2.} & \symrook{}b4 & \symknight{}h4\\
\textbf{3.} & \symrook{}g4 & \symking{}b1\\
\textbf{4.} & \symqueen{}e1+ & \symking{}c2\\
\textbf{5.} & \symking{}g5 & \symking{}d3\\
\textbf{6.} & \symking{}f4 & \symking{}d4\\
\textbf{7.} & \symqueen{}g3 & \symknight{}e6\#\\
\end{tabular}\par\vspace{3mm}
\qrcode[height=12mm]{https://helpman.komtera.lt/?fen=5N1N/7k/8/8/8/8/1r1q4/K7&moves=7}\\{\sffamily\tiny\color{ink}scannen zum Lösen}
\end{minipage}
\vspace{2mm}

{\raggedright\noindent\textcolor{ink}{\sffamily\textsc{Themen}}\enspace \themetag{reines Matt} \themetag{Mustermatt} \themetag{Idealmatt} \themetag{Selbstblock} \themetag{ohne Schlag}\par\smallskip{\footnotesize\textit{Hinweis:} die Partei am Zug steht im Diagramm im Schach.}\par}

\clearpage
\phantomsection\label{hm:KNNvkrp}
\begin{tcolorbox}[enhanced,colback=ink,colframe=ink,arc=0pt,boxrule=0pt,left=4mm,right=4mm,top=2.5mm,bottom=2.5mm]
\color{white}\sffamily{\LARGE\bfseries KNNvkrp}\hfill{\Large h\#7}\\[0.5mm]
{\small 6 Steine\quad Schwarz am Zug}\hfill{\small Nr.~514\qquad\href{https://helpman.komtera.lt/?fen=5k2/5pr1/K4NN1/8/8/8/8/8&moves=7}{\color{white}\bfseries lösen\,\ensuremath{\nearrow}}}
\end{tcolorbox}
\vspace{2mm}
\noindent\begin{minipage}[t]{36mm}\vspace{0pt}\sffamily\small\raggedright
\hmfigure{längste eindeutige}{{\Large h\#7}}\hmfigure{eindeutig bei dieser Länge}{{\Large 3} Stellungen}\hmfigure{längstes Matt}{{\Large h\#12}}\hmfigure{in der Klasse}{45,4 Milliarden Stellungen,\\ 14,7 Milliarden mit Hilfsmatt}
\end{minipage}\hfill
\begin{minipage}[t]{80mm}\vspace{0pt}\begin{tikzpicture}\begin{axis}[width=72mm,height=42mm,ybar,bar width=1.6pt,xmin=-1,xmax=25,ymin=0,ymax=3.667,axis y line*=left,axis x line*=bottom,xlabel={\scriptsize Halbzüge bis Matt},ylabel={\scriptsize Stellungen (Mrd.)},tick label style={font=\tiny},label style={font=\scriptsize},enlarge x limits=false,ylabel shift=-3pt,xlabel shift=-3pt]\addplot[fill=ink!55,draw=none] coordinates {(0,0.0001756) (1,0.0008567) (2,0.003597) (3,0.02881) (4,0.1247) (5,0.5377) (6,1.303) (7,2.563) (8,3.334) (9,3.327) (10,2.037) (11,1.094) (12,0.2769) (13,0.08354) (14,0.006866) (15,0.001415) (16,2.055e-06) (17,3.4e-08) (18,2.06e-07) (19,4e-09) (20,3.8e-08) (21,1.8e-08) (22,2e-09) (23,5e-09) (24,1e-09)};\draw[densely dotted] (axis cs:14,0) -- (axis cs:14,3.667);\end{axis}\begin{axis}[width=72mm,height=42mm,axis y line*=right,axis x line=none,ymode=log,xmin=-1,xmax=25,ylabel={\scriptsize eindeutig (log)},tick label style={font=\tiny},label style={font=\scriptsize},ylabel shift=-3pt]\addplot[mark=*,mark size=0.7pt,draw=black,thin] coordinates {(0,175572) (1,814932) (2,2618645) (3,14734601) (4,47834627) (5,68197077) (6,66830666) (7,42848635) (8,21411174) (9,7106970) (10,1335088) (11,256856) (12,16964) (13,1909) (14,3)};\end{axis}\end{tikzpicture}\end{minipage}
\vspace{3mm}

\noindent\begin{minipage}[t]{60mm}\vspace{0pt}\sbox\hmboard{\chessboard[setfen=5k2/5pr1/K4NN1/8/8/8/8/8 b - - 0 1]}\usebox\hmboard\\[1mm]
\makebox[\wd\hmboard]{\sffamily\hspace{3mm}\textbf{h\#7}\hfill(3+3)}\end{minipage}\hfill
\begin{minipage}[t]{52mm}\vspace{0pt}
\textcolor{ink}{\sffamily\textsc{Lösung}}\par\smallskip\small
\begin{tabular}{@{}r@{\enspace}l@{\enspace}l@{}}
\textbf{1.} & fxg6 & \symking{}b6\\
\textbf{2.} & \symking{}e7 & \symking{}b7\\
\textbf{3.} & \symking{}d6+ & \symking{}c8\\
\textbf{4.} & \symking{}c6 & \symking{}d8\\
\textbf{5.} & \symking{}b7 & \symknight{}d5\\
\textbf{6.} & \symking{}a8 & \symking{}c8\\
\textbf{7.} & \symrook{}a7 & \symknight{}b6\#\\
\end{tabular}\par\vspace{3mm}
\qrcode[height=12mm]{https://helpman.komtera.lt/?fen=5k2/5pr1/K4NN1/8/8/8/8/8&moves=7}\\{\sffamily\tiny\color{ink}scannen zum Lösen}
\end{minipage}
\vspace{2mm}

{\raggedright\noindent\textcolor{ink}{\sffamily\textsc{Themen}}\enspace \themetag{reines Matt} \themetag{Mustermatt} \themetag{Switchback} \themetag{Selbstblock}\par\smallskip{\footnotesize\textit{Hinweis:} die Partei am Zug steht im Diagramm im Schach; die Lösung beginnt mit einem Schlag.}\par}

\clearpage
\phantomsection\label{hm:KNPPvkb}
\begin{tcolorbox}[enhanced,colback=ink,colframe=ink,arc=0pt,boxrule=0pt,left=4mm,right=4mm,top=2.5mm,bottom=2.5mm]
\color{white}\sffamily{\LARGE\bfseries KNPPvkb}\hfill{\Large h\#7}\\[0.5mm]
{\small 6 Steine\quad Schwarz am Zug}\hfill{\small Nr.~515\qquad\href{https://helpman.komtera.lt/?fen=7b/6N1/8/8/2k5/3P4/3P4/K7&moves=7}{\color{white}\bfseries lösen\,\ensuremath{\nearrow}}}
\end{tcolorbox}
\vspace{2mm}
\noindent\begin{minipage}[t]{36mm}\vspace{0pt}\sffamily\small\raggedright
\hmfigure{längste eindeutige}{{\Large h\#7}}\hmfigure{eindeutig bei dieser Länge}{{\Large 100} Stellungen}\hmfigure{längstes Matt}{{\Large h\#8.5}}\hmfigure{in der Klasse}{34,1 Milliarden Stellungen,\\ 11,8 Milliarden mit Hilfsmatt}
\end{minipage}\hfill
\begin{minipage}[t]{80mm}\vspace{0pt}\begin{tikzpicture}\begin{axis}[width=72mm,height=42mm,ybar,bar width=1.6pt,xmin=-1,xmax=18,ymin=0,ymax=2.22,axis y line*=left,axis x line*=bottom,xlabel={\scriptsize Halbzüge bis Matt},ylabel={\scriptsize Stellungen (Mrd.)},tick label style={font=\tiny},label style={font=\scriptsize},enlarge x limits=false,ylabel shift=-3pt,xlabel shift=-3pt]\addplot[fill=ink!55,draw=none] coordinates {(0,0.000121) (1,0.004633) (2,0.01515) (3,0.1486) (4,0.4346) (5,1.406) (6,1.918) (7,2.018) (8,1.825) (9,1.456) (10,1.198) (11,0.7551) (12,0.5076) (13,0.09) (14,0.02118) (15,0.0003622) (16,5.511e-06) (17,1e-08)};\draw[densely dotted] (axis cs:14,0) -- (axis cs:14,2.22);\end{axis}\begin{axis}[width=72mm,height=42mm,axis y line*=right,axis x line=none,ymode=log,xmin=-1,xmax=18,ylabel={\scriptsize eindeutig (log)},tick label style={font=\tiny},label style={font=\scriptsize},ylabel shift=-3pt]\addplot[mark=*,mark size=0.7pt,draw=black,thin] coordinates {(0,121013) (1,1515266) (2,4727416) (3,56713537) (4,98064838) (5,106877104) (6,71297382) (7,47337484) (8,25925192) (9,14645077) (10,5272557) (11,1930128) (12,376945) (13,7575) (14,100)};\end{axis}\end{tikzpicture}\end{minipage}
\vspace{3mm}

\noindent\begin{minipage}[t]{60mm}\vspace{0pt}\sbox\hmboard{\chessboard[setfen=7b/6N1/8/8/2k5/3P4/3P4/K7 b - - 0 1]}\usebox\hmboard\\[1mm]
\makebox[\wd\hmboard]{\sffamily\hspace{3mm}\textbf{h\#7}\hfill(4+2)}\end{minipage}\hfill
\begin{minipage}[t]{52mm}\vspace{0pt}
\textcolor{ink}{\sffamily\textsc{Lösung}}\par\smallskip\small
\begin{tabular}{@{}r@{\enspace}l@{\enspace}l@{}}
\textbf{1.} & \symking{}d4 & \symknight{}e6+\\
\textbf{2.} & \symking{}e5 & d4+\\
\textbf{3.} & \symking{}f6 & d5\\
\textbf{4.} & \symbishop{}g7 & d6\\
\textbf{5.} & \symbishop{}f8 & d7\\
\textbf{6.} & \symking{}g6 & d8=\symqueen{}\\
\textbf{7.} & \symking{}h5 & \symqueen{}g5\#\\
\end{tabular}\par\vspace{3mm}
\qrcode[height=12mm]{https://helpman.komtera.lt/?fen=7b/6N1/8/8/2k5/3P4/3P4/K7&moves=7}\\{\sffamily\tiny\color{ink}scannen zum Lösen}
\end{minipage}
\vspace{2mm}

{\raggedright\noindent\textcolor{ink}{\sffamily\textsc{Themen}}\enspace \themetag{reines Matt} \themetag{Mustermatt} \themetag{Umwandlung} \themetag{ohne Schlag} \themetag{Umnov}\par\smallskip{\footnotesize\textit{Hinweis:} die Partei am Zug steht im Diagramm im Schach.}\par}

\clearpage
\phantomsection\label{hm:KNPPvkn}
\begin{tcolorbox}[enhanced,colback=ink,colframe=ink,arc=0pt,boxrule=0pt,left=4mm,right=4mm,top=2.5mm,bottom=2.5mm]
\color{white}\sffamily{\LARGE\bfseries KNPPvkn}\hfill{\Large h\#7}\\[0.5mm]
{\small 6 Steine\quad Schwarz am Zug}\hfill{\small Nr.~516\qquad\href{https://helpman.komtera.lt/?fen=8/N7/8/6n1/8/7k/P5P1/K7&moves=7}{\color{white}\bfseries lösen\,\ensuremath{\nearrow}}}
\end{tcolorbox}
\vspace{2mm}
\noindent\begin{minipage}[t]{36mm}\vspace{0pt}\sffamily\small\raggedright
\hmfigure{längste eindeutige}{{\Large h\#7}}\hmfigure{eindeutig bei dieser Länge}{{\Large 204} Stellungen}\hmfigure{längstes Matt}{{\Large h\#8.5}}\hmfigure{in der Klasse}{34,1 Milliarden Stellungen,\\ 12,1 Milliarden mit Hilfsmatt}
\end{minipage}\hfill
\begin{minipage}[t]{80mm}\vspace{0pt}\begin{tikzpicture}\begin{axis}[width=72mm,height=42mm,ybar,bar width=1.6pt,xmin=-1,xmax=18,ymin=0,ymax=2.159,axis y line*=left,axis x line*=bottom,xlabel={\scriptsize Halbzüge bis Matt},ylabel={\scriptsize Stellungen (Mrd.)},tick label style={font=\tiny},label style={font=\scriptsize},enlarge x limits=false,ylabel shift=-3pt,xlabel shift=-3pt]\addplot[fill=ink!55,draw=none] coordinates {(0,0.0001203) (1,0.005698) (2,0.01594) (3,0.1446) (4,0.3818) (5,1.313) (6,1.866) (7,1.963) (8,1.95) (9,1.508) (10,1.375) (11,0.8315) (12,0.5918) (13,0.1111) (14,0.03492) (15,0.0005775) (16,1.392e-05) (17,1.1e-08)};\draw[densely dotted] (axis cs:14,0) -- (axis cs:14,2.159);\end{axis}\begin{axis}[width=72mm,height=42mm,axis y line*=right,axis x line=none,ymode=log,xmin=-1,xmax=18,ylabel={\scriptsize eindeutig (log)},tick label style={font=\tiny},label style={font=\scriptsize},ylabel shift=-3pt]\addplot[mark=*,mark size=0.7pt,draw=black,thin] coordinates {(0,120345) (1,1579260) (2,4164411) (3,51807639) (4,76596704) (5,94446430) (6,63149379) (7,41921453) (8,26212484) (9,14618214) (10,6537069) (11,2805293) (12,564507) (13,11150) (14,204)};\end{axis}\end{tikzpicture}\end{minipage}
\vspace{3mm}

\noindent\begin{minipage}[t]{60mm}\vspace{0pt}\sbox\hmboard{\chessboard[setfen=8/N7/8/6n1/8/7k/P5P1/K7 b - - 0 1]}\usebox\hmboard\\[1mm]
\makebox[\wd\hmboard]{\sffamily\hspace{3mm}\textbf{h\#7}\hfill(4+2)}\end{minipage}\hfill
\begin{minipage}[t]{52mm}\vspace{0pt}
\textcolor{ink}{\sffamily\textsc{Lösung}}\par\smallskip\small
\begin{tabular}{@{}r@{\enspace}l@{\enspace}l@{}}
\textbf{1.} & \symking{}h4 & g4\\
\textbf{2.} & \symknight{}e6 & g5\\
\textbf{3.} & \symknight{}c7 & g6\\
\textbf{4.} & \symking{}g5 & g7\\
\textbf{5.} & \symking{}f6 & g8=\symqueen{}\\
\textbf{6.} & \symking{}e7 & \symqueen{}e6+\\
\textbf{7.} & \symking{}d8 & \symknight{}c6\#\\
\end{tabular}\par\vspace{3mm}
\qrcode[height=12mm]{https://helpman.komtera.lt/?fen=8/N7/8/6n1/8/7k/P5P1/K7&moves=7}\\{\sffamily\tiny\color{ink}scannen zum Lösen}
\end{minipage}
\vspace{2mm}

{\raggedright\noindent\textcolor{ink}{\sffamily\textsc{Themen}}\enspace \themetag{Umwandlung} \themetag{Excelsior} \themetag{Selbstblock} \themetag{ohne Schlag} \themetag{Umnov}\par\smallskip{\footnotesize\textit{Hinweis:} das Diagramm hat keinen legalen letzten Zug, kann also in keiner Partie entstehen.}\par}

\clearpage
\phantomsection\label{hm:KNPPvkr}
\begin{tcolorbox}[enhanced,colback=ink,colframe=ink,arc=0pt,boxrule=0pt,left=4mm,right=4mm,top=2.5mm,bottom=2.5mm]
\color{white}\sffamily{\LARGE\bfseries KNPPvkr}\hfill{\Large h\#7}\\[0.5mm]
{\small 6 Steine\quad Schwarz am Zug}\hfill{\small Nr.~517\qquad\href{https://helpman.komtera.lt/?fen=8/8/8/N7/K7/P1r5/P7/6k1&moves=7}{\color{white}\bfseries lösen\,\ensuremath{\nearrow}}}
\end{tcolorbox}
\vspace{2mm}
\noindent\begin{minipage}[t]{36mm}\vspace{0pt}\sffamily\small\raggedright
\hmfigure{längste eindeutige}{{\Large h\#7}}\hmfigure{eindeutig bei dieser Länge}{{\Large 29} Stellungen}\hmfigure{längstes Matt}{{\Large h\#8.5}}\hmfigure{in der Klasse}{34,1 Milliarden Stellungen,\\ 11,3 Milliarden mit Hilfsmatt}
\end{minipage}\hfill
\begin{minipage}[t]{80mm}\vspace{0pt}\begin{tikzpicture}\begin{axis}[width=72mm,height=42mm,ybar,bar width=1.6pt,xmin=-1,xmax=18,ymin=0,ymax=2.343,axis y line*=left,axis x line*=bottom,xlabel={\scriptsize Halbzüge bis Matt},ylabel={\scriptsize Stellungen (Mrd.)},tick label style={font=\tiny},label style={font=\scriptsize},enlarge x limits=false,ylabel shift=-3pt,xlabel shift=-3pt]\addplot[fill=ink!55,draw=none] coordinates {(0,9.814e-05) (1,0.004194) (2,0.01727) (3,0.1575) (4,0.5258) (5,1.554) (6,1.993) (7,2.13) (8,1.641) (9,1.389) (10,0.987) (11,0.5872) (12,0.274) (13,0.05181) (14,0.005775) (15,0.0003203) (16,2.195e-06) (17,1.68e-07)};\draw[densely dotted] (axis cs:14,0) -- (axis cs:14,2.343);\end{axis}\begin{axis}[width=72mm,height=42mm,axis y line*=right,axis x line=none,ymode=log,xmin=-1,xmax=18,ylabel={\scriptsize eindeutig (log)},tick label style={font=\tiny},label style={font=\scriptsize},ylabel shift=-3pt]\addplot[mark=*,mark size=0.7pt,draw=black,thin] coordinates {(0,98136) (1,1474997) (2,6147571) (3,59378660) (4,116542938) (5,121785761) (6,59240453) (7,39360543) (8,14758644) (9,8616004) (10,2737974) (11,754076) (12,80035) (13,4368) (14,29)};\end{axis}\end{tikzpicture}\end{minipage}
\vspace{3mm}

\noindent\begin{minipage}[t]{60mm}\vspace{0pt}\sbox\hmboard{\chessboard[setfen=8/8/8/N7/K7/P1r5/P7/6k1 b - - 0 1]}\usebox\hmboard\\[1mm]
\makebox[\wd\hmboard]{\sffamily\hspace{3mm}\textbf{h\#7}\hfill(4+2)}\end{minipage}\hfill
\begin{minipage}[t]{52mm}\vspace{0pt}
\textcolor{ink}{\sffamily\textsc{Lösung}}\par\smallskip\small
\begin{tabular}{@{}r@{\enspace}l@{\enspace}l@{}}
\textbf{1.} & \symrook{}b3 & axb3\\
\textbf{2.} & \symking{}f2 & b4\\
\textbf{3.} & \symking{}e3 & b5\\
\textbf{4.} & \symking{}d4 & b6\\
\textbf{5.} & \symking{}c5 & b7\\
\textbf{6.} & \symking{}b6 & b8=\symqueen{}+\\
\textbf{7.} & \symking{}a6 & \symqueen{}b7\#\\
\end{tabular}\par\vspace{3mm}
\qrcode[height=12mm]{https://helpman.komtera.lt/?fen=8/8/8/N7/K7/P1r5/P7/6k1&moves=7}\\{\sffamily\tiny\color{ink}scannen zum Lösen}
\end{minipage}
\vspace{2mm}

{\raggedright\noindent\textcolor{ink}{\sffamily\textsc{Themen}}\enspace \themetag{Umwandlung} \themetag{Excelsior} \themetag{Switchback} \themetag{ein Stein zieht} \themetag{Umnov}\par}

\clearpage
\phantomsection\label{hm:KPPPvkb}
\begin{tcolorbox}[enhanced,colback=ink,colframe=ink,arc=0pt,boxrule=0pt,left=4mm,right=4mm,top=2.5mm,bottom=2.5mm]
\color{white}\sffamily{\LARGE\bfseries KPPPvkb}\hfill{\Large h\#7}\\[0.5mm]
{\small 6 Steine\quad Schwarz am Zug}\hfill{\small Nr.~518\qquad\href{https://helpman.komtera.lt/?fen=8/8/8/8/K7/P7/P1P5/3b3k&moves=7}{\color{white}\bfseries lösen\,\ensuremath{\nearrow}}}
\end{tcolorbox}
\vspace{2mm}
\noindent\begin{minipage}[t]{36mm}\vspace{0pt}\sffamily\small\raggedright
\hmfigure{längste eindeutige}{{\Large h\#7}}\hmfigure{eindeutig bei dieser Länge}{{\Large 8} Stellungen}\hmfigure{längstes Matt}{{\Large h\#8}}\hmfigure{in der Klasse}{25,6 Milliarden Stellungen,\\ 3 Milliarden mit Hilfsmatt}
\end{minipage}\hfill
\begin{minipage}[t]{80mm}\vspace{0pt}\begin{tikzpicture}\begin{axis}[width=72mm,height=42mm,ybar,bar width=1.6pt,xmin=-1,xmax=17,ymin=0,ymax=663,axis y line*=left,axis x line*=bottom,xlabel={\scriptsize Halbzüge bis Matt},ylabel={\scriptsize Stellungen (Mio.)},tick label style={font=\tiny},label style={font=\scriptsize},enlarge x limits=false,ylabel shift=-3pt,xlabel shift=-3pt]\addplot[fill=ink!55,draw=none] coordinates {(0,0.03207) (1,1.598) (2,4.921) (3,41.54) (4,114.1) (5,372.6) (6,525.2) (7,602.7) (8,499.8) (9,374.4) (10,242.1) (11,132.1) (12,77.88) (13,12.54) (14,2.942) (15,0.04342) (16,5.3e-05)};\draw[densely dotted] (axis cs:14,0) -- (axis cs:14,663);\end{axis}\begin{axis}[width=72mm,height=42mm,axis y line*=right,axis x line=none,ymode=log,xmin=-1,xmax=17,ylabel={\scriptsize eindeutig (log)},tick label style={font=\tiny},label style={font=\scriptsize},ylabel shift=-3pt]\addplot[mark=*,mark size=0.7pt,draw=black,thin] coordinates {(0,32074) (1,461294) (2,1363149) (3,14759443) (4,25398789) (5,29295873) (6,18269407) (7,11152688) (8,5664835) (9,2902643) (10,886131) (11,299925) (12,54267) (13,987) (14,8)};\end{axis}\end{tikzpicture}\end{minipage}
\vspace{3mm}

\noindent\begin{minipage}[t]{60mm}\vspace{0pt}\sbox\hmboard{\chessboard[setfen=8/8/8/8/K7/P7/P1P5/3b3k b - - 0 1]}\usebox\hmboard\\[1mm]
\makebox[\wd\hmboard]{\sffamily\hspace{3mm}\textbf{h\#7}\hfill(4+2)}\end{minipage}\hfill
\begin{minipage}[t]{52mm}\vspace{0pt}
\textcolor{ink}{\sffamily\textsc{Lösung}}\par\smallskip\small
\begin{tabular}{@{}r@{\enspace}l@{\enspace}l@{}}
\textbf{1.} & \symking{}g2 & \symking{}b4\\
\textbf{2.} & \symking{}f3 & a4\\
\textbf{3.} & \symking{}e4 & a5\\
\textbf{4.} & \symking{}d5 & a6\\
\textbf{5.} & \symking{}c6 & a7\\
\textbf{6.} & \symking{}b7 & \symking{}c5\\
\textbf{7.} & \symking{}a6 & a8=\symqueen{}\#\\
\end{tabular}\par\vspace{3mm}
\qrcode[height=12mm]{https://helpman.komtera.lt/?fen=8/8/8/8/K7/P7/P1P5/3b3k&moves=7}\\{\sffamily\tiny\color{ink}scannen zum Lösen}
\end{minipage}
\vspace{2mm}

{\raggedright\noindent\textcolor{ink}{\sffamily\textsc{Themen}}\enspace \themetag{reines Matt} \themetag{Mustermatt} \themetag{Spiegelmatt} \themetag{Umwandlung} \themetag{ein Stein zieht} \themetag{ohne Schlag} \themetag{ohne Schach}\par}

\clearpage
\phantomsection\label{hm:KPPPvkn}
\begin{tcolorbox}[enhanced,colback=ink,colframe=ink,arc=0pt,boxrule=0pt,left=4mm,right=4mm,top=2.5mm,bottom=2.5mm]
\color{white}\sffamily{\LARGE\bfseries KPPPvkn}\hfill{\Large h\#7}\\[0.5mm]
{\small 6 Steine\quad Schwarz am Zug}\hfill{\small Nr.~519\qquad\href{https://helpman.komtera.lt/?fen=8/8/8/8/K7/P7/1P4P1/3n3k&moves=7}{\color{white}\bfseries lösen\,\ensuremath{\nearrow}}}
\end{tcolorbox}
\vspace{2mm}
\noindent\begin{minipage}[t]{36mm}\vspace{0pt}\sffamily\small\raggedright
\hmfigure{längste eindeutige}{{\Large h\#7}}\hmfigure{eindeutig bei dieser Länge}{{\Large 33} Stellungen}\hmfigure{längstes Matt}{{\Large h\#8.5}}\hmfigure{in der Klasse}{25,6 Milliarden Stellungen,\\ 3,08 Milliarden mit Hilfsmatt}
\end{minipage}\hfill
\begin{minipage}[t]{80mm}\vspace{0pt}\begin{tikzpicture}\begin{axis}[width=72mm,height=42mm,ybar,bar width=1.6pt,xmin=-1,xmax=18,ymin=0,ymax=643.8,axis y line*=left,axis x line*=bottom,xlabel={\scriptsize Halbzüge bis Matt},ylabel={\scriptsize Stellungen (Mio.)},tick label style={font=\tiny},label style={font=\scriptsize},enlarge x limits=false,ylabel shift=-3pt,xlabel shift=-3pt]\addplot[fill=ink!55,draw=none] coordinates {(0,0.03065) (1,1.981) (2,5.278) (3,40.97) (4,102.8) (5,345.8) (6,501) (7,585.3) (8,536.8) (9,393.7) (10,286.5) (11,151.2) (12,99.29) (13,18.52) (14,7.014) (15,0.1567) (16,0.001212) (17,4e-06)};\draw[densely dotted] (axis cs:14,0) -- (axis cs:14,643.8);\end{axis}\begin{axis}[width=72mm,height=42mm,axis y line*=right,axis x line=none,ymode=log,xmin=-1,xmax=18,ylabel={\scriptsize eindeutig (log)},tick label style={font=\tiny},label style={font=\scriptsize},ylabel shift=-3pt]\addplot[mark=*,mark size=0.7pt,draw=black,thin] coordinates {(0,30650) (1,481417) (2,1221233) (3,13310916) (4,20021493) (5,25000447) (6,16583786) (7,10011820) (8,5563119) (9,2818120) (10,1193156) (11,479901) (12,124474) (13,2757) (14,33)};\end{axis}\end{tikzpicture}\end{minipage}
\vspace{3mm}

\noindent\begin{minipage}[t]{60mm}\vspace{0pt}\sbox\hmboard{\chessboard[setfen=8/8/8/8/K7/P7/1P4P1/3n3k b - - 0 1]}\usebox\hmboard\\[1mm]
\makebox[\wd\hmboard]{\sffamily\hspace{3mm}\textbf{h\#7}\hfill(4+2)}\end{minipage}\hfill
\begin{minipage}[t]{52mm}\vspace{0pt}
\textcolor{ink}{\sffamily\textsc{Lösung}}\par\smallskip\small
\begin{tabular}{@{}r@{\enspace}l@{\enspace}l@{}}
\textbf{1.} & \symking{}xg2 & \symking{}b4\\
\textbf{2.} & \symking{}f3 & a4\\
\textbf{3.} & \symking{}e4 & a5\\
\textbf{4.} & \symking{}d5 & a6\\
\textbf{5.} & \symking{}c6 & a7\\
\textbf{6.} & \symking{}b7 & \symking{}c5\\
\textbf{7.} & \symking{}a6 & a8=\symqueen{}\#\\
\end{tabular}\par\vspace{3mm}
\qrcode[height=12mm]{https://helpman.komtera.lt/?fen=8/8/8/8/K7/P7/1P4P1/3n3k&moves=7}\\{\sffamily\tiny\color{ink}scannen zum Lösen}
\end{minipage}
\vspace{2mm}

{\raggedright\noindent\textcolor{ink}{\sffamily\textsc{Themen}}\enspace \themetag{reines Matt} \themetag{Mustermatt} \themetag{Spiegelmatt} \themetag{Umwandlung} \themetag{ein Stein zieht} \themetag{ohne Schach}\par\smallskip{\footnotesize\textit{Hinweis:} die Lösung beginnt mit einem Schlag.}\par}

\clearpage
\phantomsection\label{hm:KQBBBvk}
\begin{tcolorbox}[enhanced,colback=ink,colframe=ink,arc=0pt,boxrule=0pt,left=4mm,right=4mm,top=2.5mm,bottom=2.5mm]
\color{white}\sffamily{\LARGE\bfseries KQBBBvk}\hfill{\Large h\#7}\\[0.5mm]
{\small 6 Steine\quad Schwarz am Zug}\hfill{\small Nr.~520\qquad\href{https://helpman.komtera.lt/?fen=5B1Q/5B1k/8/8/8/8/7B/K7&moves=7}{\color{white}\bfseries lösen\,\ensuremath{\nearrow}}}
\end{tcolorbox}
\vspace{2mm}
\noindent\begin{minipage}[t]{36mm}\vspace{0pt}\sffamily\small\raggedright
\hmfigure{längste eindeutige}{{\Large h\#7}}\hmfigure{eindeutig bei dieser Länge}{{\Large 3} Stellungen}\hmfigure{längstes Matt}{{\Large h\#8}}\hmfigure{in der Klasse}{15,5 Milliarden Stellungen,\\ 1,41 Milliarden mit Hilfsmatt}
\end{minipage}\hfill
\begin{minipage}[t]{80mm}\vspace{0pt}\begin{tikzpicture}\begin{axis}[width=72mm,height=42mm,ybar,bar width=1.6pt,xmin=-1,xmax=17,ymin=0,ymax=703.8,axis y line*=left,axis x line*=bottom,xlabel={\scriptsize Halbzüge bis Matt},ylabel={\scriptsize Stellungen (Mio.)},tick label style={font=\tiny},label style={font=\scriptsize},enlarge x limits=false,ylabel shift=-3pt,xlabel shift=-3pt]\addplot[fill=ink!55,draw=none] coordinates {(0,24.81) (1,76.27) (2,287.6) (3,299.4) (4,639.8) (5,54.94) (6,24.24) (7,0.003516) (8,4.778) (10,1.334) (12,0.295) (14,0.03486) (16,2.8e-05)};\draw[densely dotted] (axis cs:14,0) -- (axis cs:14,703.8);\end{axis}\begin{axis}[width=72mm,height=42mm,axis y line*=right,axis x line=none,ymode=log,xmin=-1,xmax=17,ylabel={\scriptsize eindeutig (log)},tick label style={font=\tiny},label style={font=\scriptsize},ylabel shift=-3pt]\addplot[mark=*,mark size=0.7pt,draw=black,thin] coordinates {(0,24806519) (1,50211710) (2,157151489) (3,22583347) (4,9415228) (5,1128) (6,251991) (8,26666) (10,3136) (12,433) (14,3)};\end{axis}\end{tikzpicture}\end{minipage}
\vspace{3mm}

\noindent\begin{minipage}[t]{60mm}\vspace{0pt}\sbox\hmboard{\chessboard[setfen=5B1Q/5B1k/8/8/8/8/7B/K7 b - - 0 1]}\usebox\hmboard\\[1mm]
\makebox[\wd\hmboard]{\sffamily\hspace{3mm}\textbf{h\#7}\hfill(5+1)}\end{minipage}\hfill
\begin{minipage}[t]{52mm}\vspace{0pt}
\textcolor{ink}{\sffamily\textsc{Lösung}}\par\smallskip\small
\begin{tabular}{@{}r@{\enspace}l@{\enspace}l@{}}
\textbf{1.} & \symking{}xh8 & \symking{}b2\\
\textbf{2.} & \symking{}h7 & \symking{}c3\\
\textbf{3.} & \symking{}h8 & \symking{}d4\\
\textbf{4.} & \symking{}h7 & \symking{}e5\\
\textbf{5.} & \symking{}h8 & \symking{}f6\\
\textbf{6.} & \symking{}h7 & \symbishop{}e5\\
\textbf{7.} & \symking{}h8 & \symking{}g6\#\\
\end{tabular}\par\vspace{3mm}
\qrcode[height=12mm]{https://helpman.komtera.lt/?fen=5B1Q/5B1k/8/8/8/8/7B/K7&moves=7}\\{\sffamily\tiny\color{ink}scannen zum Lösen}
\end{minipage}
\vspace{2mm}

{\raggedright\noindent\textcolor{ink}{\sffamily\textsc{Themen}}\enspace \themetag{Spiegelmatt} \themetag{Switchback} \themetag{ein Stein zieht} \themetag{Kniest} \themetag{Pendel} \themetag{ohne Schach}\par\smallskip{\footnotesize\textit{Hinweis:} die Partei am Zug steht im Diagramm im Schach; die Lösung beginnt mit einem Schlag.}\par}

\clearpage
\phantomsection\label{hm:KQBBNvk}
\begin{tcolorbox}[enhanced,colback=ink,colframe=ink,arc=0pt,boxrule=0pt,left=4mm,right=4mm,top=2.5mm,bottom=2.5mm]
\color{white}\sffamily{\LARGE\bfseries KQBBNvk}\hfill{\Large h\#7}\\[0.5mm]
{\small 6 Steine\quad Schwarz am Zug}\hfill{\small Nr.~521\qquad\href{https://helpman.komtera.lt/?fen=7Q/5B1k/8/5B2/8/8/5N2/K7&moves=7}{\color{white}\bfseries lösen\,\ensuremath{\nearrow}}}
\end{tcolorbox}
\vspace{2mm}
\noindent\begin{minipage}[t]{36mm}\vspace{0pt}\sffamily\small\raggedright
\hmfigure{längste eindeutige}{{\Large h\#7}}\hmfigure{eindeutig bei dieser Länge}{{\Large 1} Stellung}\hmfigure{längstes Matt}{{\Large h\#8}}\hmfigure{in der Klasse}{15,5 Milliarden Stellungen,\\ 4,38 Milliarden mit Hilfsmatt}
\end{minipage}\hfill
\begin{minipage}[t]{80mm}\vspace{0pt}\begin{tikzpicture}\begin{axis}[width=72mm,height=42mm,ybar,bar width=1.6pt,xmin=-1,xmax=17,ymin=0,ymax=2.282,axis y line*=left,axis x line*=bottom,xlabel={\scriptsize Halbzüge bis Matt},ylabel={\scriptsize Stellungen (Mrd.)},tick label style={font=\tiny},label style={font=\scriptsize},enlarge x limits=false,ylabel shift=-3pt,xlabel shift=-3pt]\addplot[fill=ink!55,draw=none] coordinates {(0,0.0494) (1,0.1979) (2,0.7323) (3,0.9898) (4,2.075) (5,0.2151) (6,0.07549) (7,1.065e-06) (8,0.02259) (9,1e-09) (10,0.01586) (12,0.004153) (14,0.0006609) (16,1.456e-05)};\draw[densely dotted] (axis cs:14,0) -- (axis cs:14,2.282);\end{axis}\begin{axis}[width=72mm,height=42mm,axis y line*=right,axis x line=none,ymode=log,xmin=-1,xmax=17,ylabel={\scriptsize eindeutig (log)},tick label style={font=\tiny},label style={font=\scriptsize},ylabel shift=-3pt]\addplot[mark=*,mark size=0.7pt,draw=black,thin] coordinates {(0,49400898) (1,141701881) (2,442840680) (3,110730870) (4,40002556) (5,370) (6,1314472) (8,251818) (9,1) (10,19016) (12,484) (14,1)};\end{axis}\end{tikzpicture}\end{minipage}
\vspace{3mm}

\noindent\begin{minipage}[t]{60mm}\vspace{0pt}\sbox\hmboard{\chessboard[setfen=7Q/5B1k/8/5B2/8/8/5N2/K7 b - - 0 1]}\usebox\hmboard\\[1mm]
\makebox[\wd\hmboard]{\sffamily\hspace{3mm}\textbf{h\#7}\hfill(5+1)}\end{minipage}\hfill
\begin{minipage}[t]{52mm}\vspace{0pt}
\textcolor{ink}{\sffamily\textsc{Lösung}}\par\smallskip\small
\begin{tabular}{@{}r@{\enspace}l@{\enspace}l@{}}
\textbf{1.} & \symking{}xh8 & \symking{}b2\\
\textbf{2.} & \symking{}g7 & \symking{}c3\\
\textbf{3.} & \symking{}xf7 & \symking{}d4\\
\textbf{4.} & \symking{}g7 & \symking{}e5\\
\textbf{5.} & \symking{}h6 & \symking{}f6\\
\textbf{6.} & \symking{}h5 & \symbishop{}g6+\\
\textbf{7.} & \symking{}h6 & \symknight{}g4\#\\
\end{tabular}\par\vspace{3mm}
\qrcode[height=12mm]{https://helpman.komtera.lt/?fen=7Q/5B1k/8/5B2/8/8/5N2/K7&moves=7}\\{\sffamily\tiny\color{ink}scannen zum Lösen}
\end{minipage}
\vspace{2mm}

{\raggedright\noindent\textcolor{ink}{\sffamily\textsc{Themen}}\enspace \themetag{reines Matt} \themetag{Mustermatt} \themetag{Idealmatt} \themetag{Switchback} \themetag{ein Stein zieht}\par\smallskip{\footnotesize\textit{Hinweis:} das Diagramm hat keinen legalen letzten Zug, kann also in keiner Partie entstehen.}\par}

\clearpage
\phantomsection\label{hm:KQBBPvk}
\begin{tcolorbox}[enhanced,colback=ink,colframe=ink,arc=0pt,boxrule=0pt,left=4mm,right=4mm,top=2.5mm,bottom=2.5mm]
\color{white}\sffamily{\LARGE\bfseries KQBBPvk}\hfill{\Large h\#7}\\[0.5mm]
{\small 6 Steine\quad Schwarz am Zug}\hfill{\small Nr.~522\qquad\href{https://helpman.komtera.lt/?fen=1B6/8/8/8/7B/8/6QP/K4k2&moves=7}{\color{white}\bfseries lösen\,\ensuremath{\nearrow}}}
\end{tcolorbox}
\vspace{2mm}
\noindent\begin{minipage}[t]{36mm}\vspace{0pt}\sffamily\small\raggedright
\hmfigure{längste eindeutige}{{\Large h\#7}}\hmfigure{eindeutig bei dieser Länge}{{\Large 2\,289} Stellungen}\hmfigure{längstes Matt}{{\Large h\#9}}\hmfigure{in der Klasse}{45,4 Milliarden Stellungen,\\ 13,4 Milliarden mit Hilfsmatt}
\end{minipage}\hfill
\begin{minipage}[t]{80mm}\vspace{0pt}\begin{tikzpicture}\begin{axis}[width=72mm,height=42mm,ybar,bar width=1.6pt,xmin=-1,xmax=19,ymin=0,ymax=7.157,axis y line*=left,axis x line*=bottom,xlabel={\scriptsize Halbzüge bis Matt},ylabel={\scriptsize Stellungen (Mrd.)},tick label style={font=\tiny},label style={font=\scriptsize},enlarge x limits=false,ylabel shift=-3pt,xlabel shift=-3pt]\addplot[fill=ink!55,draw=none] coordinates {(0,0.1172) (1,0.5017) (2,1.809) (3,2.978) (4,6.507) (5,0.9642) (6,0.4049) (7,9.547e-05) (8,0.04535) (10,0.02873) (12,0.02613) (14,0.0118) (16,0.0003734) (18,1.298e-06)};\draw[densely dotted] (axis cs:14,0) -- (axis cs:14,7.157);\end{axis}\begin{axis}[width=72mm,height=42mm,axis y line*=right,axis x line=none,ymode=log,xmin=-1,xmax=19,ylabel={\scriptsize eindeutig (log)},tick label style={font=\tiny},label style={font=\scriptsize},ylabel shift=-3pt]\addplot[mark=*,mark size=0.7pt,draw=black,thin] coordinates {(0,117201069) (1,359910275) (2,1137393671) (3,417338696) (4,192223681) (5,38457) (6,4284241) (8,1087593) (10,483596) (12,256207) (14,2289)};\end{axis}\end{tikzpicture}\end{minipage}
\vspace{3mm}

\noindent\begin{minipage}[t]{60mm}\vspace{0pt}\sbox\hmboard{\chessboard[setfen=1B6/8/8/8/7B/8/6QP/K4k2 b - - 0 1]}\usebox\hmboard\\[1mm]
\makebox[\wd\hmboard]{\sffamily\hspace{3mm}\textbf{h\#7}\hfill(5+1)}\end{minipage}\hfill
\begin{minipage}[t]{52mm}\vspace{0pt}
\textcolor{ink}{\sffamily\textsc{Lösung}}\par\smallskip\small
\begin{tabular}{@{}r@{\enspace}l@{\enspace}l@{}}
\textbf{1.} & \symking{}xg2 & \symbishop{}hg3\\
\textbf{2.} & \symking{}f3 & h4\\
\textbf{3.} & \symking{}e4 & h5\\
\textbf{4.} & \symking{}d5 & h6\\
\textbf{5.} & \symking{}c6 & h7\\
\textbf{6.} & \symking{}b7 & h8=\symqueen{}\\
\textbf{7.} & \symking{}a8 & \symqueen{}h1\#\\
\end{tabular}\par\vspace{3mm}
\qrcode[height=12mm]{https://helpman.komtera.lt/?fen=1B6/8/8/8/7B/8/6QP/K4k2&moves=7}\\{\sffamily\tiny\color{ink}scannen zum Lösen}
\end{minipage}
\vspace{2mm}

{\raggedright\noindent\textcolor{ink}{\sffamily\textsc{Themen}}\enspace \themetag{reines Matt} \themetag{Mustermatt} \themetag{Umwandlung} \themetag{Excelsior} \themetag{ein Stein zieht} \themetag{Phönix} \themetag{ohne Schach}\par\smallskip{\footnotesize\textit{Hinweis:} die Partei am Zug steht im Diagramm im Schach; die Lösung beginnt mit einem Schlag.}\par}

\clearpage
\phantomsection\label{hm:KQBBvkb}
\begin{tcolorbox}[enhanced,colback=ink,colframe=ink,arc=0pt,boxrule=0pt,left=4mm,right=4mm,top=2.5mm,bottom=2.5mm]
\color{white}\sffamily{\LARGE\bfseries KQBBvkb}\hfill{\Large h\#7}\\[0.5mm]
{\small 6 Steine\quad Schwarz am Zug}\hfill{\small Nr.~523\qquad\href{https://helpman.komtera.lt/?fen=6b1/8/8/8/8/4B3/5B1Q/K6k&moves=7}{\color{white}\bfseries lösen\,\ensuremath{\nearrow}}}
\end{tcolorbox}
\vspace{2mm}
\noindent\begin{minipage}[t]{36mm}\vspace{0pt}\sffamily\small\raggedright
\hmfigure{längste eindeutige}{{\Large h\#7}}\hmfigure{eindeutig bei dieser Länge}{{\Large 289} Stellungen}\hmfigure{längstes Matt}{{\Large h\#11}}\hmfigure{in der Klasse}{15,5 Milliarden Stellungen,\\ 4,14 Milliarden mit Hilfsmatt}
\end{minipage}\hfill
\begin{minipage}[t]{80mm}\vspace{0pt}\begin{tikzpicture}\begin{axis}[width=72mm,height=42mm,ybar,bar width=1.6pt,xmin=-1,xmax=23,ymin=0,ymax=2.115,axis y line*=left,axis x line*=bottom,xlabel={\scriptsize Halbzüge bis Matt},ylabel={\scriptsize Stellungen (Mrd.)},tick label style={font=\tiny},label style={font=\scriptsize},enlarge x limits=false,ylabel shift=-3pt,xlabel shift=-3pt]\addplot[fill=ink!55,draw=none] coordinates {(0,0.02606) (1,0.1223) (2,0.5683) (3,1.051) (4,1.923) (5,0.3443) (6,0.07833) (7,0.0006754) (8,0.009594) (9,2.53e-07) (10,0.005793) (11,2.2e-08) (12,0.004695) (13,3e-08) (14,0.002973) (15,2.6e-08) (16,0.001913) (17,9.91e-07) (18,0.001231) (20,0.0003042) (22,1e-09)};\draw[densely dotted] (axis cs:14,0) -- (axis cs:14,2.115);\end{axis}\begin{axis}[width=72mm,height=42mm,axis y line*=right,axis x line=none,ymode=log,xmin=-1,xmax=23,ylabel={\scriptsize eindeutig (log)},tick label style={font=\tiny},label style={font=\scriptsize},ylabel shift=-3pt]\addplot[mark=*,mark size=0.7pt,draw=black,thin] coordinates {(0,26057164) (1,95433099) (2,301940120) (3,98198485) (4,28466307) (5,113999) (6,382011) (7,7) (8,64348) (10,20001) (12,5387) (14,289)};\end{axis}\end{tikzpicture}\end{minipage}
\vspace{3mm}

\noindent\begin{minipage}[t]{60mm}\vspace{0pt}\sbox\hmboard{\chessboard[setfen=6b1/8/8/8/8/4B3/5B1Q/K6k b - - 0 1]}\usebox\hmboard\\[1mm]
\makebox[\wd\hmboard]{\sffamily\hspace{3mm}\textbf{h\#7}\hfill(4+2)}\end{minipage}\hfill
\begin{minipage}[t]{52mm}\vspace{0pt}
\textcolor{ink}{\sffamily\textsc{Lösung}}\par\smallskip\small
\begin{tabular}{@{}r@{\enspace}l@{\enspace}l@{}}
\textbf{1.} & \symking{}xh2 & \symking{}b2\\
\textbf{2.} & \symking{}h3 & \symking{}c3\\
\textbf{3.} & \symking{}g4 & \symking{}d4\\
\textbf{4.} & \symking{}h5 & \symking{}e5\\
\textbf{5.} & \symking{}g6 & \symbishop{}d4\\
\textbf{6.} & \symking{}h7 & \symking{}f6\\
\textbf{7.} & \symking{}h8 & \symking{}g6\#\\
\end{tabular}\par\vspace{3mm}
\qrcode[height=12mm]{https://helpman.komtera.lt/?fen=6b1/8/8/8/8/4B3/5B1Q/K6k&moves=7}\\{\sffamily\tiny\color{ink}scannen zum Lösen}
\end{minipage}
\vspace{2mm}

{\raggedright\noindent\textcolor{ink}{\sffamily\textsc{Themen}}\enspace \themetag{ein Stein zieht} \themetag{ohne Schach}\par\smallskip{\footnotesize\textit{Hinweis:} die Partei am Zug steht im Diagramm im Schach; die Lösung beginnt mit einem Schlag.}\par}

\clearpage
\phantomsection\label{hm:KQBNNvk}
\begin{tcolorbox}[enhanced,colback=ink,colframe=ink,arc=0pt,boxrule=0pt,left=4mm,right=4mm,top=2.5mm,bottom=2.5mm]
\color{white}\sffamily{\LARGE\bfseries KQBNNvk}\hfill{\Large h\#7}\\[0.5mm]
{\small 6 Steine\quad Schwarz am Zug}\hfill{\small Nr.~524\qquad\href{https://helpman.komtera.lt/?fen=4N2Q/6k1/8/7B/8/5N2/8/K7&moves=7}{\color{white}\bfseries lösen\,\ensuremath{\nearrow}}}
\end{tcolorbox}
\vspace{2mm}
\noindent\begin{minipage}[t]{36mm}\vspace{0pt}\sffamily\small\raggedright
\hmfigure{längste eindeutige}{{\Large h\#7}}\hmfigure{eindeutig bei dieser Länge}{{\Large 6} Stellungen}\hmfigure{längstes Matt}{{\Large h\#8}}\hmfigure{in der Klasse}{15,5 Milliarden Stellungen,\\ 4,48 Milliarden mit Hilfsmatt}
\end{minipage}\hfill
\begin{minipage}[t]{80mm}\vspace{0pt}\begin{tikzpicture}\begin{axis}[width=72mm,height=42mm,ybar,bar width=1.6pt,xmin=-1,xmax=17,ymin=0,ymax=2.387,axis y line*=left,axis x line*=bottom,xlabel={\scriptsize Halbzüge bis Matt},ylabel={\scriptsize Stellungen (Mrd.)},tick label style={font=\tiny},label style={font=\scriptsize},enlarge x limits=false,ylabel shift=-3pt,xlabel shift=-3pt]\addplot[fill=ink!55,draw=none] coordinates {(0,0.03403) (1,0.1687) (2,0.6163) (3,1.027) (4,2.17) (5,0.3066) (6,0.1159) (7,7.392e-05) (8,0.02347) (10,0.01727) (12,0.003795) (14,0.000615) (16,6.374e-06)};\draw[densely dotted] (axis cs:14,0) -- (axis cs:14,2.387);\end{axis}\begin{axis}[width=72mm,height=42mm,axis y line*=right,axis x line=none,ymode=log,xmin=-1,xmax=17,ylabel={\scriptsize eindeutig (log)},tick label style={font=\tiny},label style={font=\scriptsize},ylabel shift=-3pt]\addplot[mark=*,mark size=0.7pt,draw=black,thin] coordinates {(0,34028759) (1,127414607) (2,400676576) (3,159756319) (4,64516219) (5,9534) (6,1403186) (8,307003) (10,21453) (12,660) (14,6)};\end{axis}\end{tikzpicture}\end{minipage}
\vspace{3mm}

\noindent\begin{minipage}[t]{60mm}\vspace{0pt}\sbox\hmboard{\chessboard[setfen=4N2Q/6k1/8/7B/8/5N2/8/K7 b - - 0 1]}\usebox\hmboard\\[1mm]
\makebox[\wd\hmboard]{\sffamily\hspace{3mm}\textbf{h\#7}\hfill(5+1)}\end{minipage}\hfill
\begin{minipage}[t]{52mm}\vspace{0pt}
\textcolor{ink}{\sffamily\textsc{Lösung}}\par\smallskip\small
\begin{tabular}{@{}r@{\enspace}l@{\enspace}l@{}}
\textbf{1.} & \symking{}xh8 & \symking{}b2\\
\textbf{2.} & \symking{}g8 & \symking{}c3\\
\textbf{3.} & \symking{}h7 & \symking{}d4\\
\textbf{4.} & \symking{}h6 & \symking{}e5\\
\textbf{5.} & \symking{}xh5 & \symking{}f6\\
\textbf{6.} & \symking{}g4 & \symking{}g7\\
\textbf{7.} & \symking{}h5 & \symknight{}f6\#\\
\end{tabular}\par\vspace{3mm}
\qrcode[height=12mm]{https://helpman.komtera.lt/?fen=4N2Q/6k1/8/7B/8/5N2/8/K7&moves=7}\\{\sffamily\tiny\color{ink}scannen zum Lösen}
\end{minipage}
\vspace{2mm}

{\raggedright\noindent\textcolor{ink}{\sffamily\textsc{Themen}}\enspace \themetag{reines Matt} \themetag{Mustermatt} \themetag{Idealmatt} \themetag{Spiegelmatt} \themetag{Switchback} \themetag{ein Stein zieht} \themetag{Kniest} \themetag{ohne Schach}\par\smallskip{\footnotesize\textit{Hinweis:} das Diagramm hat keinen legalen letzten Zug, kann also in keiner Partie entstehen.}\par}

\clearpage
\phantomsection\label{hm:KQBNPvk}
\begin{tcolorbox}[enhanced,colback=ink,colframe=ink,arc=0pt,boxrule=0pt,left=4mm,right=4mm,top=2.5mm,bottom=2.5mm]
\color{white}\sffamily{\LARGE\bfseries KQBNPvk}\hfill{\Large h\#7}\\[0.5mm]
{\small 6 Steine\quad Schwarz am Zug}\hfill{\small Nr.~525\qquad\href{https://helpman.komtera.lt/?fen=4N3/8/8/6kQ/8/4P3/B7/K7&moves=7}{\color{white}\bfseries lösen\,\ensuremath{\nearrow}}}
\end{tcolorbox}
\vspace{2mm}
\noindent\begin{minipage}[t]{36mm}\vspace{0pt}\sffamily\small\raggedright
\hmfigure{längste eindeutige}{{\Large h\#7}}\hmfigure{eindeutig bei dieser Länge}{{\Large 34} Stellungen}\hmfigure{längstes Matt}{{\Large h\#8}}\hmfigure{in der Klasse}{45,4 Milliarden Stellungen,\\ 27,5 Milliarden mit Hilfsmatt}
\end{minipage}\hfill
\begin{minipage}[t]{80mm}\vspace{0pt}\begin{tikzpicture}\begin{axis}[width=72mm,height=42mm,ybar,bar width=1.6pt,xmin=-1,xmax=17,ymin=0,ymax=14.75,axis y line*=left,axis x line*=bottom,xlabel={\scriptsize Halbzüge bis Matt},ylabel={\scriptsize Stellungen (Mrd.)},tick label style={font=\tiny},label style={font=\scriptsize},enlarge x limits=false,ylabel shift=-3pt,xlabel shift=-3pt]\addplot[fill=ink!55,draw=none] coordinates {(0,0.1529) (1,0.8289) (2,2.919) (3,6.049) (4,13.41) (5,2.66) (6,1.225) (7,0.002782) (8,0.1019) (10,0.09382) (11,2e-09) (12,0.03324) (14,0.002931) (16,3.954e-06)};\draw[densely dotted] (axis cs:14,0) -- (axis cs:14,14.75);\end{axis}\begin{axis}[width=72mm,height=42mm,axis y line*=right,axis x line=none,ymode=log,xmin=-1,xmax=17,ylabel={\scriptsize eindeutig (log)},tick label style={font=\tiny},label style={font=\scriptsize},ylabel shift=-3pt]\addplot[mark=*,mark size=0.7pt,draw=black,thin] coordinates {(0,152916998) (1,623906855) (2,1969872667) (3,1177051291) (4,608507143) (5,378915) (6,8145497) (8,2994805) (10,601778) (12,56543) (14,34)};\end{axis}\end{tikzpicture}\end{minipage}
\vspace{3mm}

\noindent\begin{minipage}[t]{60mm}\vspace{0pt}\sbox\hmboard{\chessboard[setfen=4N3/8/8/6kQ/8/4P3/B7/K7 b - - 0 1]}\usebox\hmboard\\[1mm]
\makebox[\wd\hmboard]{\sffamily\hspace{3mm}\textbf{h\#7}\hfill(5+1)}\end{minipage}\hfill
\begin{minipage}[t]{52mm}\vspace{0pt}
\textcolor{ink}{\sffamily\textsc{Lösung}}\par\smallskip\small
\begin{tabular}{@{}r@{\enspace}l@{\enspace}l@{}}
\textbf{1.} & \symking{}xh5 & e4\\
\textbf{2.} & \symking{}g6 & \symbishop{}b1\\
\textbf{3.} & \symking{}f7 & e5\\
\textbf{4.} & \symking{}xe8 & e6\\
\textbf{5.} & \symking{}f8 & e7+\\
\textbf{6.} & \symking{}g7 & e8=\symqueen{}\\
\textbf{7.} & \symking{}h6 & \symqueen{}g6\#\\
\end{tabular}\par\vspace{3mm}
\qrcode[height=12mm]{https://helpman.komtera.lt/?fen=4N3/8/8/6kQ/8/4P3/B7/K7&moves=7}\\{\sffamily\tiny\color{ink}scannen zum Lösen}
\end{minipage}
\vspace{2mm}

{\raggedright\noindent\textcolor{ink}{\sffamily\textsc{Themen}}\enspace \themetag{reines Matt} \themetag{Mustermatt} \themetag{Umwandlung} \themetag{ein Stein zieht} \themetag{Phönix}\par\smallskip{\footnotesize\textit{Hinweis:} die Partei am Zug steht im Diagramm im Schach; die Lösung beginnt mit einem Schlag.}\par}

\clearpage
\phantomsection\label{hm:KQBNvkb}
\begin{tcolorbox}[enhanced,colback=ink,colframe=ink,arc=0pt,boxrule=0pt,left=4mm,right=4mm,top=2.5mm,bottom=2.5mm]
\color{white}\sffamily{\LARGE\bfseries KQBNvkb}\hfill{\Large h\#7}\\[0.5mm]
{\small 6 Steine\quad Schwarz am Zug}\hfill{\small Nr.~526\qquad\href{https://helpman.komtera.lt/?fen=6N1/6Qk/8/8/2B5/1b6/8/K7&moves=7}{\color{white}\bfseries lösen\,\ensuremath{\nearrow}}}
\end{tcolorbox}
\vspace{2mm}
\noindent\begin{minipage}[t]{36mm}\vspace{0pt}\sffamily\small\raggedright
\hmfigure{längste eindeutige}{{\Large h\#7}}\hmfigure{eindeutig bei dieser Länge}{{\Large 39} Stellungen}\hmfigure{längstes Matt}{{\Large h\#9}}\hmfigure{in der Klasse}{15,5 Milliarden Stellungen,\\ 8,53 Milliarden mit Hilfsmatt}
\end{minipage}\hfill
\begin{minipage}[t]{80mm}\vspace{0pt}\begin{tikzpicture}\begin{axis}[width=72mm,height=42mm,ybar,bar width=1.6pt,xmin=-1,xmax=19,ymin=0,ymax=4.427,axis y line*=left,axis x line*=bottom,xlabel={\scriptsize Halbzüge bis Matt},ylabel={\scriptsize Stellungen (Mrd.)},tick label style={font=\tiny},label style={font=\scriptsize},enlarge x limits=false,ylabel shift=-3pt,xlabel shift=-3pt]\addplot[fill=ink!55,draw=none] coordinates {(0,0.03307) (1,0.1955) (2,0.8854) (3,2.108) (4,4.025) (5,0.9512) (6,0.2465) (7,0.003895) (8,0.03385) (9,8.24e-07) (10,0.03332) (11,5.55e-07) (12,0.01234) (13,3.13e-07) (14,0.001926) (15,4e-09) (16,3.06e-05) (18,8e-09)};\draw[densely dotted] (axis cs:14,0) -- (axis cs:14,4.427);\end{axis}\begin{axis}[width=72mm,height=42mm,axis y line*=right,axis x line=none,ymode=log,xmin=-1,xmax=19,ylabel={\scriptsize eindeutig (log)},tick label style={font=\tiny},label style={font=\scriptsize},ylabel shift=-3pt]\addplot[mark=*,mark size=0.7pt,draw=black,thin] coordinates {(0,33070435) (1,160246430) (2,501569241) (3,301184564) (4,103481016) (5,659370) (6,993635) (7,201) (8,348436) (10,59019) (12,1802) (14,39)};\end{axis}\end{tikzpicture}\end{minipage}
\vspace{3mm}

\noindent\begin{minipage}[t]{60mm}\vspace{0pt}\sbox\hmboard{\chessboard[setfen=6N1/6Qk/8/8/2B5/1b6/8/K7 b - - 0 1]}\usebox\hmboard\\[1mm]
\makebox[\wd\hmboard]{\sffamily\hspace{3mm}\textbf{h\#7}\hfill(4+2)}\end{minipage}\hfill
\begin{minipage}[t]{52mm}\vspace{0pt}
\textcolor{ink}{\sffamily\textsc{Lösung}}\par\smallskip\small
\begin{tabular}{@{}r@{\enspace}l@{\enspace}l@{}}
\textbf{1.} & \symking{}xg7 & \symknight{}e7\\
\textbf{2.} & \symking{}f6 & \symking{}b2\\
\textbf{3.} & \symking{}e5 & \symking{}a3\\
\textbf{4.} & \symking{}d4 & \symknight{}c6+\\
\textbf{5.} & \symking{}c3 & \symknight{}a5\\
\textbf{6.} & \symking{}c2 & \symknight{}xb3\\
\textbf{7.} & \symking{}b1 & \symbishop{}d3\#\\
\end{tabular}\par\vspace{3mm}
\qrcode[height=12mm]{https://helpman.komtera.lt/?fen=6N1/6Qk/8/8/2B5/1b6/8/K7&moves=7}\\{\sffamily\tiny\color{ink}scannen zum Lösen}
\end{minipage}
\vspace{2mm}

{\raggedright\noindent\textcolor{ink}{\sffamily\textsc{Themen}}\enspace \themetag{reines Matt} \themetag{Mustermatt} \themetag{Idealmatt} \themetag{Spiegelmatt} \themetag{ein Stein zieht}\par\smallskip{\footnotesize\textit{Hinweis:} die Partei am Zug steht im Diagramm im Schach; die Lösung beginnt mit einem Schlag.}\par}

\clearpage
\phantomsection\label{hm:KQBNvkn}
\begin{tcolorbox}[enhanced,colback=ink,colframe=ink,arc=0pt,boxrule=0pt,left=4mm,right=4mm,top=2.5mm,bottom=2.5mm]
\color{white}\sffamily{\LARGE\bfseries KQBNvkn}\hfill{\Large h\#7}\\[0.5mm]
{\small 6 Steine\quad Schwarz am Zug}\hfill{\small Nr.~527\qquad\href{https://helpman.komtera.lt/?fen=8/6QN/7k/8/2n5/8/4B3/K7&moves=7}{\color{white}\bfseries lösen\,\ensuremath{\nearrow}}}
\end{tcolorbox}
\vspace{2mm}
\noindent\begin{minipage}[t]{36mm}\vspace{0pt}\sffamily\small\raggedright
\hmfigure{längste eindeutige}{{\Large h\#7}}\hmfigure{eindeutig bei dieser Länge}{{\Large 101} Stellungen}\hmfigure{längstes Matt}{{\Large h\#9}}\hmfigure{in der Klasse}{15,5 Milliarden Stellungen,\\ 8,81 Milliarden mit Hilfsmatt}
\end{minipage}\hfill
\begin{minipage}[t]{80mm}\vspace{0pt}\begin{tikzpicture}\begin{axis}[width=72mm,height=42mm,ybar,bar width=1.6pt,xmin=-1,xmax=19,ymin=0,ymax=4.615,axis y line*=left,axis x line*=bottom,xlabel={\scriptsize Halbzüge bis Matt},ylabel={\scriptsize Stellungen (Mrd.)},tick label style={font=\tiny},label style={font=\scriptsize},enlarge x limits=false,ylabel shift=-3pt,xlabel shift=-3pt]\addplot[fill=ink!55,draw=none] coordinates {(0,0.03549) (1,0.2047) (2,0.8755) (3,2.095) (4,4.196) (5,0.9538) (6,0.3496) (7,0.005356) (8,0.02888) (9,1.239e-06) (10,0.03994) (11,1.839e-06) (12,0.01801) (13,1.439e-06) (14,0.003153) (15,1.35e-07) (16,5.074e-05) (18,3e-09)};\draw[densely dotted] (axis cs:14,0) -- (axis cs:14,4.615);\end{axis}\begin{axis}[width=72mm,height=42mm,axis y line*=right,axis x line=none,ymode=log,xmin=-1,xmax=19,ylabel={\scriptsize eindeutig (log)},tick label style={font=\tiny},label style={font=\scriptsize},ylabel shift=-3pt]\addplot[mark=*,mark size=0.7pt,draw=black,thin] coordinates {(0,35491109) (1,165500855) (2,490618257) (3,314808630) (4,121116944) (5,690809) (6,744807) (7,100) (8,338430) (9,30) (10,74202) (11,18) (12,3123) (13,1) (14,101)};\end{axis}\end{tikzpicture}\end{minipage}
\vspace{3mm}

\noindent\begin{minipage}[t]{60mm}\vspace{0pt}\sbox\hmboard{\chessboard[setfen=8/6QN/7k/8/2n5/8/4B3/K7 b - - 0 1]}\usebox\hmboard\\[1mm]
\makebox[\wd\hmboard]{\sffamily\hspace{3mm}\textbf{h\#7}\hfill(4+2)}\end{minipage}\hfill
\begin{minipage}[t]{52mm}\vspace{0pt}
\textcolor{ink}{\sffamily\textsc{Lösung}}\par\smallskip\small
\begin{tabular}{@{}r@{\enspace}l@{\enspace}l@{}}
\textbf{1.} & \symking{}xg7 & \symknight{}g5\\
\textbf{2.} & \symking{}f6 & \symking{}b1\\
\textbf{3.} & \symking{}e5 & \symking{}c1\\
\textbf{4.} & \symking{}d4 & \symknight{}f3+\\
\textbf{5.} & \symking{}c3 & \symknight{}e1\\
\textbf{6.} & \symking{}b3 & \symknight{}c2\\
\textbf{7.} & \symking{}a2 & \symbishop{}xc4\#\\
\end{tabular}\par\vspace{3mm}
\qrcode[height=12mm]{https://helpman.komtera.lt/?fen=8/6QN/7k/8/2n5/8/4B3/K7&moves=7}\\{\sffamily\tiny\color{ink}scannen zum Lösen}
\end{minipage}
\vspace{2mm}

{\raggedright\noindent\textcolor{ink}{\sffamily\textsc{Themen}}\enspace \themetag{reines Matt} \themetag{Mustermatt} \themetag{Idealmatt} \themetag{Spiegelmatt} \themetag{ein Stein zieht}\par\smallskip{\footnotesize\textit{Hinweis:} die Partei am Zug steht im Diagramm im Schach; die Lösung beginnt mit einem Schlag.}\par}

\clearpage
\phantomsection\label{hm:KQBNvkp}
\begin{tcolorbox}[enhanced,colback=ink,colframe=ink,arc=0pt,boxrule=0pt,left=4mm,right=4mm,top=2.5mm,bottom=2.5mm]
\color{white}\sffamily{\LARGE\bfseries KQBNvkp}\hfill{\Large h\#7}\\[0.5mm]
{\small 6 Steine\quad Schwarz am Zug}\hfill{\small Nr.~528\qquad\href{https://helpman.komtera.lt/?fen=5N2/8/8/8/2p5/4B2Q/6k1/K7&moves=7}{\color{white}\bfseries lösen\,\ensuremath{\nearrow}}}
\end{tcolorbox}
\vspace{2mm}
\noindent\begin{minipage}[t]{36mm}\vspace{0pt}\sffamily\small\raggedright
\hmfigure{längste eindeutige}{{\Large h\#7}}\hmfigure{eindeutig bei dieser Länge}{{\Large 412} Stellungen}\hmfigure{längstes Matt}{{\Large h\#12}}\hmfigure{in der Klasse}{45,4 Milliarden Stellungen,\\ 27,4 Milliarden mit Hilfsmatt}
\end{minipage}\hfill
\begin{minipage}[t]{80mm}\vspace{0pt}\begin{tikzpicture}\begin{axis}[width=72mm,height=42mm,ybar,bar width=1.6pt,xmin=-1,xmax=25,ymin=0,ymax=14.53,axis y line*=left,axis x line*=bottom,xlabel={\scriptsize Halbzüge bis Matt},ylabel={\scriptsize Stellungen (Mrd.)},tick label style={font=\tiny},label style={font=\scriptsize},enlarge x limits=false,ylabel shift=-3pt,xlabel shift=-3pt]\addplot[fill=ink!55,draw=none] coordinates {(0,0.1306) (1,0.7415) (2,2.973) (3,6.607) (4,13.21) (5,2.476) (6,1.025) (7,0.005962) (8,0.08462) (9,3.75e-07) (10,0.1257) (11,2.1e-08) (12,0.05958) (14,0.01019) (16,0.0001314) (18,1.1e-08) (20,1.8e-08) (22,1.2e-08) (24,1.4e-08)};\draw[densely dotted] (axis cs:14,0) -- (axis cs:14,14.53);\end{axis}\begin{axis}[width=72mm,height=42mm,axis y line*=right,axis x line=none,ymode=log,xmin=-1,xmax=25,ylabel={\scriptsize eindeutig (log)},tick label style={font=\tiny},label style={font=\scriptsize},ylabel shift=-3pt]\addplot[mark=*,mark size=0.7pt,draw=black,thin] coordinates {(0,130641188) (1,585081101) (2,1920175341) (3,1154932311) (4,418683998) (5,787741) (6,3405708) (7,19) (8,1551235) (10,323064) (12,19259) (14,412)};\end{axis}\end{tikzpicture}\end{minipage}
\vspace{3mm}

\noindent\begin{minipage}[t]{60mm}\vspace{0pt}\sbox\hmboard{\chessboard[setfen=5N2/8/8/8/2p5/4B2Q/6k1/K7 b - - 0 1]}\usebox\hmboard\\[1mm]
\makebox[\wd\hmboard]{\sffamily\hspace{3mm}\textbf{h\#7}\hfill(4+2)}\end{minipage}\hfill
\begin{minipage}[t]{52mm}\vspace{0pt}
\textcolor{ink}{\sffamily\textsc{Lösung}}\par\smallskip\small
\begin{tabular}{@{}r@{\enspace}l@{\enspace}l@{}}
\textbf{1.} & \symking{}xh3 & \symking{}b2\\
\textbf{2.} & \symking{}g4 & \symking{}c3\\
\textbf{3.} & \symking{}f5 & \symking{}d4\\
\textbf{4.} & \symking{}f6 & \symking{}d5\\
\textbf{5.} & \symking{}g7 & \symking{}e6\\
\textbf{6.} & \symking{}h8 & \symking{}f7\\
\textbf{7.} & c3 & \symbishop{}d4\#\\
\end{tabular}\par\vspace{3mm}
\qrcode[height=12mm]{https://helpman.komtera.lt/?fen=5N2/8/8/8/2p5/4B2Q/6k1/K7&moves=7}\\{\sffamily\tiny\color{ink}scannen zum Lösen}
\end{minipage}
\vspace{2mm}

{\raggedright\noindent\textcolor{ink}{\sffamily\textsc{Themen}}\enspace \themetag{Spiegelmatt} \themetag{ohne Schach}\par\smallskip{\footnotesize\textit{Hinweis:} die Partei am Zug steht im Diagramm im Schach; die Lösung beginnt mit einem Schlag.}\par}

\clearpage
\phantomsection\label{hm:KQBNvkq}
\begin{tcolorbox}[enhanced,colback=ink,colframe=ink,arc=0pt,boxrule=0pt,left=4mm,right=4mm,top=2.5mm,bottom=2.5mm]
\color{white}\sffamily{\LARGE\bfseries KQBNvkq}\hfill{\Large h\#7}\\[0.5mm]
{\small 6 Steine\quad Schwarz am Zug}\hfill{\small Nr.~529\qquad\href{https://helpman.komtera.lt/?fen=5N2/8/8/8/8/8/1B2q2Q/K6k&moves=7}{\color{white}\bfseries lösen\,\ensuremath{\nearrow}}}
\end{tcolorbox}
\vspace{2mm}
\noindent\begin{minipage}[t]{36mm}\vspace{0pt}\sffamily\small\raggedright
\hmfigure{längste eindeutige}{{\Large h\#7}}\hmfigure{eindeutig bei dieser Länge}{{\Large 168} Stellungen}\hmfigure{längstes Matt}{{\Large h\#9}}\hmfigure{in der Klasse}{15,5 Milliarden Stellungen,\\ 7,33 Milliarden mit Hilfsmatt}
\end{minipage}\hfill
\begin{minipage}[t]{80mm}\vspace{0pt}\begin{tikzpicture}\begin{axis}[width=72mm,height=42mm,ybar,bar width=1.6pt,xmin=-1,xmax=19,ymin=0,ymax=3.484,axis y line*=left,axis x line*=bottom,xlabel={\scriptsize Halbzüge bis Matt},ylabel={\scriptsize Stellungen (Mrd.)},tick label style={font=\tiny},label style={font=\scriptsize},enlarge x limits=false,ylabel shift=-3pt,xlabel shift=-3pt]\addplot[fill=ink!55,draw=none] coordinates {(0,0.01884) (1,0.1211) (2,0.6246) (3,1.795) (4,3.168) (5,1.329) (6,0.1954) (7,0.0118) (8,0.01941) (9,3.229e-06) (10,0.03006) (11,4.23e-07) (12,0.01725) (13,1.77e-07) (14,0.003044) (15,1.5e-08) (16,8.256e-05) (18,3.1e-08)};\draw[densely dotted] (axis cs:14,0) -- (axis cs:14,3.484);\end{axis}\begin{axis}[width=72mm,height=42mm,axis y line*=right,axis x line=none,ymode=log,xmin=-1,xmax=19,ylabel={\scriptsize eindeutig (log)},tick label style={font=\tiny},label style={font=\scriptsize},ylabel shift=-3pt]\addplot[mark=*,mark size=0.7pt,draw=black,thin] coordinates {(0,18838109) (1,105288341) (2,358140042) (3,290143614) (4,82425565) (5,2977255) (6,507492) (7,464) (8,439413) (9,3) (10,108338) (11,1) (12,6356) (14,168)};\end{axis}\end{tikzpicture}\end{minipage}
\vspace{3mm}

\noindent\begin{minipage}[t]{60mm}\vspace{0pt}\sbox\hmboard{\chessboard[setfen=5N2/8/8/8/8/8/1B2q2Q/K6k b - - 0 1]}\usebox\hmboard\\[1mm]
\makebox[\wd\hmboard]{\sffamily\hspace{3mm}\textbf{h\#7}\hfill(4+2)}\end{minipage}\hfill
\begin{minipage}[t]{52mm}\vspace{0pt}
\textcolor{ink}{\sffamily\textsc{Lösung}}\par\smallskip\small
\begin{tabular}{@{}r@{\enspace}l@{\enspace}l@{}}
\textbf{1.} & \symking{}xh2 & \symbishop{}h8\\
\textbf{2.} & \symqueen{}g4 & \symking{}b2\\
\textbf{3.} & \symking{}g3 & \symking{}c3\\
\textbf{4.} & \symking{}f4 & \symking{}d4\\
\textbf{5.} & \symking{}g5+ & \symking{}e5\\
\textbf{6.} & \symking{}h6 & \symking{}f6\\
\textbf{7.} & \symqueen{}h5 & \symbishop{}g7\#\\
\end{tabular}\par\vspace{3mm}
\qrcode[height=12mm]{https://helpman.komtera.lt/?fen=5N2/8/8/8/8/8/1B2q2Q/K6k&moves=7}\\{\sffamily\tiny\color{ink}scannen zum Lösen}
\end{minipage}
\vspace{2mm}

{\raggedright\noindent\textcolor{ink}{\sffamily\textsc{Themen}}\enspace \themetag{Selbstblock}\par\smallskip{\footnotesize\textit{Hinweis:} die Partei am Zug steht im Diagramm im Schach; die Lösung beginnt mit einem Schlag.}\par}

\clearpage
\phantomsection\label{hm:KQBNvkr}
\begin{tcolorbox}[enhanced,colback=ink,colframe=ink,arc=0pt,boxrule=0pt,left=4mm,right=4mm,top=2.5mm,bottom=2.5mm]
\color{white}\sffamily{\LARGE\bfseries KQBNvkr}\hfill{\Large h\#7}\\[0.5mm]
{\small 6 Steine\quad Schwarz am Zug}\hfill{\small Nr.~530\qquad\href{https://helpman.komtera.lt/?fen=8/7Q/6k1/8/7B/3r4/8/KN6&moves=7}{\color{white}\bfseries lösen\,\ensuremath{\nearrow}}}
\end{tcolorbox}
\vspace{2mm}
\noindent\begin{minipage}[t]{36mm}\vspace{0pt}\sffamily\small\raggedright
\hmfigure{längste eindeutige}{{\Large h\#7}}\hmfigure{eindeutig bei dieser Länge}{{\Large 49} Stellungen}\hmfigure{längstes Matt}{{\Large h\#9}}\hmfigure{in der Klasse}{15,5 Milliarden Stellungen,\\ 8,11 Milliarden mit Hilfsmatt}
\end{minipage}\hfill
\begin{minipage}[t]{80mm}\vspace{0pt}\begin{tikzpicture}\begin{axis}[width=72mm,height=42mm,ybar,bar width=1.6pt,xmin=-1,xmax=19,ymin=0,ymax=4.112,axis y line*=left,axis x line*=bottom,xlabel={\scriptsize Halbzüge bis Matt},ylabel={\scriptsize Stellungen (Mrd.)},tick label style={font=\tiny},label style={font=\scriptsize},enlarge x limits=false,ylabel shift=-3pt,xlabel shift=-3pt]\addplot[fill=ink!55,draw=none] coordinates {(0,0.02695) (1,0.1691) (2,0.8184) (3,2.019) (4,3.738) (5,1.066) (6,0.1872) (7,0.003836) (8,0.02789) (9,2.28e-07) (10,0.03326) (11,3.82e-07) (12,0.01523) (13,5.4e-08) (14,0.002236) (16,3.617e-05) (18,1e-09)};\draw[densely dotted] (axis cs:14,0) -- (axis cs:14,4.112);\end{axis}\begin{axis}[width=72mm,height=42mm,axis y line*=right,axis x line=none,ymode=log,xmin=-1,xmax=19,ylabel={\scriptsize eindeutig (log)},tick label style={font=\tiny},label style={font=\scriptsize},ylabel shift=-3pt]\addplot[mark=*,mark size=0.7pt,draw=black,thin] coordinates {(0,26951898) (1,142382368) (2,469247637) (3,288625022) (4,88493259) (5,861911) (6,861603) (7,17) (8,419345) (9,3) (10,91785) (12,5226) (14,49)};\end{axis}\end{tikzpicture}\end{minipage}
\vspace{3mm}

\noindent\begin{minipage}[t]{60mm}\vspace{0pt}\sbox\hmboard{\chessboard[setfen=8/7Q/6k1/8/7B/3r4/8/KN6 b - - 0 1]}\usebox\hmboard\\[1mm]
\makebox[\wd\hmboard]{\sffamily\hspace{3mm}\textbf{h\#7}\hfill(4+2)}\end{minipage}\hfill
\begin{minipage}[t]{52mm}\vspace{0pt}
\textcolor{ink}{\sffamily\textsc{Lösung}}\par\smallskip\small
\begin{tabular}{@{}r@{\enspace}l@{\enspace}l@{}}
\textbf{1.} & \symking{}xh7 & \symking{}b2\\
\textbf{2.} & \symking{}g6 & \symking{}c2\\
\textbf{3.} & \symking{}f5 & \symking{}xd3\\
\textbf{4.} & \symking{}f4 & \symking{}c4\\
\textbf{5.} & \symking{}e3 & \symknight{}c3\\
\textbf{6.} & \symking{}d2 & \symking{}b3\\
\textbf{7.} & \symking{}c1 & \symbishop{}g5\#\\
\end{tabular}\par\vspace{3mm}
\qrcode[height=12mm]{https://helpman.komtera.lt/?fen=8/7Q/6k1/8/7B/3r4/8/KN6&moves=7}\\{\sffamily\tiny\color{ink}scannen zum Lösen}
\end{minipage}
\vspace{2mm}

{\raggedright\noindent\textcolor{ink}{\sffamily\textsc{Themen}}\enspace \themetag{reines Matt} \themetag{Mustermatt} \themetag{Idealmatt} \themetag{Spiegelmatt} \themetag{Switchback} \themetag{ein Stein zieht} \themetag{ohne Schach}\par\smallskip{\footnotesize\textit{Hinweis:} die Partei am Zug steht im Diagramm im Schach; die Lösung beginnt mit einem Schlag.}\par}

\clearpage
\phantomsection\label{hm:KQBPPvk}
\begin{tcolorbox}[enhanced,colback=ink,colframe=ink,arc=0pt,boxrule=0pt,left=4mm,right=4mm,top=2.5mm,bottom=2.5mm]
\color{white}\sffamily{\LARGE\bfseries KQBPPvk}\hfill{\Large h\#7}\\[0.5mm]
{\small 6 Steine\quad Schwarz am Zug}\hfill{\small Nr.~531\qquad\href{https://helpman.komtera.lt/?fen=k7/Q7/8/8/8/5P2/6P1/K6B&moves=7}{\color{white}\bfseries lösen\,\ensuremath{\nearrow}}}
\end{tcolorbox}
\vspace{2mm}
\noindent\begin{minipage}[t]{36mm}\vspace{0pt}\sffamily\small\raggedright
\hmfigure{längste eindeutige}{{\Large h\#7}}\hmfigure{eindeutig bei dieser Länge}{{\Large 49} Stellungen}\hmfigure{längstes Matt}{{\Large h\#8}}\hmfigure{in der Klasse}{34,1 Milliarden Stellungen,\\ 10,5 Milliarden mit Hilfsmatt}
\end{minipage}\hfill
\begin{minipage}[t]{80mm}\vspace{0pt}\begin{tikzpicture}\begin{axis}[width=72mm,height=42mm,ybar,bar width=1.6pt,xmin=-1,xmax=17,ymin=0,ymax=5.469,axis y line*=left,axis x line*=bottom,xlabel={\scriptsize Halbzüge bis Matt},ylabel={\scriptsize Stellungen (Mrd.)},tick label style={font=\tiny},label style={font=\scriptsize},enlarge x limits=false,ylabel shift=-3pt,xlabel shift=-3pt]\addplot[fill=ink!55,draw=none] coordinates {(0,0.04396) (1,0.2516) (2,0.845) (3,2.132) (4,4.972) (5,1.371) (6,0.7723) (7,0.005405) (8,0.03476) (9,1.429e-06) (10,0.02192) (11,1e-09) (12,0.01434) (13,2e-09) (14,0.002644) (16,4.293e-06)};\draw[densely dotted] (axis cs:14,0) -- (axis cs:14,5.469);\end{axis}\begin{axis}[width=72mm,height=42mm,axis y line*=right,axis x line=none,ymode=log,xmin=-1,xmax=17,ylabel={\scriptsize eindeutig (log)},tick label style={font=\tiny},label style={font=\scriptsize},ylabel shift=-3pt]\addplot[mark=*,mark size=0.7pt,draw=black,thin] coordinates {(0,43957461) (1,181933061) (2,573890142) (3,501028753) (4,330044701) (5,924116) (6,3169038) (7,308) (8,915212) (10,320844) (12,73463) (14,49)};\end{axis}\end{tikzpicture}\end{minipage}
\vspace{3mm}

\noindent\begin{minipage}[t]{60mm}\vspace{0pt}\sbox\hmboard{\chessboard[setfen=k7/Q7/8/8/8/5P2/6P1/K6B b - - 0 1]}\usebox\hmboard\\[1mm]
\makebox[\wd\hmboard]{\sffamily\hspace{3mm}\textbf{h\#7}\hfill(5+1)}\end{minipage}\hfill
\begin{minipage}[t]{52mm}\vspace{0pt}
\textcolor{ink}{\sffamily\textsc{Lösung}}\par\smallskip\small
\begin{tabular}{@{}r@{\enspace}l@{\enspace}l@{}}
\textbf{1.} & \symking{}xa7 & g4\\
\textbf{2.} & \symking{}b6 & g5\\
\textbf{3.} & \symking{}c5 & g6\\
\textbf{4.} & \symking{}d4 & g7\\
\textbf{5.} & \symking{}e3 & g8=\symqueen{}\\
\textbf{6.} & \symking{}f2 & \symqueen{}a2+\\
\textbf{7.} & \symking{}g1 & \symqueen{}g2\#\\
\end{tabular}\par\vspace{3mm}
\qrcode[height=12mm]{https://helpman.komtera.lt/?fen=k7/Q7/8/8/8/5P2/6P1/K6B&moves=7}\\{\sffamily\tiny\color{ink}scannen zum Lösen}
\end{minipage}
\vspace{2mm}

{\raggedright\noindent\textcolor{ink}{\sffamily\textsc{Themen}}\enspace \themetag{reines Matt} \themetag{Mustermatt} \themetag{Umwandlung} \themetag{Excelsior} \themetag{Rundlauf} \themetag{ein Stein zieht} \themetag{Phönix}\par\smallskip{\footnotesize\textit{Hinweis:} die Partei am Zug steht im Diagramm im Schach; die Lösung beginnt mit einem Schlag.}\par}

\clearpage
\phantomsection\label{hm:KQBPvkb}
\begin{tcolorbox}[enhanced,colback=ink,colframe=ink,arc=0pt,boxrule=0pt,left=4mm,right=4mm,top=2.5mm,bottom=2.5mm]
\color{white}\sffamily{\LARGE\bfseries KQBPvkb}\hfill{\Large h\#7}\\[0.5mm]
{\small 6 Steine\quad Schwarz am Zug}\hfill{\small Nr.~532\qquad\href{https://helpman.komtera.lt/?fen=7b/7Q/7k/8/8/8/1P1B4/K7&moves=7}{\color{white}\bfseries lösen\,\ensuremath{\nearrow}}}
\end{tcolorbox}
\vspace{2mm}
\noindent\begin{minipage}[t]{36mm}\vspace{0pt}\sffamily\small\raggedright
\hmfigure{längste eindeutige}{{\Large h\#7}}\hmfigure{eindeutig bei dieser Länge}{{\Large 958} Stellungen}\hmfigure{längstes Matt}{{\Large h\#9}}\hmfigure{in der Klasse}{45,4 Milliarden Stellungen,\\ 26,2 Milliarden mit Hilfsmatt}
\end{minipage}\hfill
\begin{minipage}[t]{80mm}\vspace{0pt}\begin{tikzpicture}\begin{axis}[width=72mm,height=42mm,ybar,bar width=1.6pt,xmin=-1,xmax=19,ymin=0,ymax=13.28,axis y line*=left,axis x line*=bottom,xlabel={\scriptsize Halbzüge bis Matt},ylabel={\scriptsize Stellungen (Mrd.)},tick label style={font=\tiny},label style={font=\scriptsize},enlarge x limits=false,ylabel shift=-3pt,xlabel shift=-3pt]\addplot[fill=ink!55,draw=none] coordinates {(0,0.0742) (1,0.4594) (2,2.006) (3,5.847) (4,12.08) (5,3.994) (6,1.478) (7,0.03706) (8,0.07342) (9,4.277e-06) (10,0.06317) (11,1.386e-06) (12,0.05948) (13,1.134e-06) (14,0.01252) (15,1.213e-06) (16,3.467e-05) (17,7.8e-08) (18,2.4e-08)};\draw[densely dotted] (axis cs:14,0) -- (axis cs:14,13.28);\end{axis}\begin{axis}[width=72mm,height=42mm,axis y line*=right,axis x line=none,ymode=log,xmin=-1,xmax=19,ylabel={\scriptsize eindeutig (log)},tick label style={font=\tiny},label style={font=\scriptsize},ylabel shift=-3pt]\addplot[mark=*,mark size=0.7pt,draw=black,thin] coordinates {(0,74200511) (1,367790347) (2,1153015076) (3,1040038935) (4,505548840) (5,6366036) (6,3577139) (7,397) (8,1650421) (9,49) (10,660218) (12,206882) (14,958)};\end{axis}\end{tikzpicture}\end{minipage}
\vspace{3mm}

\noindent\begin{minipage}[t]{60mm}\vspace{0pt}\sbox\hmboard{\chessboard[setfen=7b/7Q/7k/8/8/8/1P1B4/K7 b - - 0 1]}\usebox\hmboard\\[1mm]
\makebox[\wd\hmboard]{\sffamily\hspace{3mm}\textbf{h\#7}\hfill(4+2)}\end{minipage}\hfill
\begin{minipage}[t]{52mm}\vspace{0pt}
\textcolor{ink}{\sffamily\textsc{Lösung}}\par\smallskip\small
\begin{tabular}{@{}r@{\enspace}l@{\enspace}l@{}}
\textbf{1.} & \symking{}xh7 & \symbishop{}c3\\
\textbf{2.} & \symking{}g6 & b4\\
\textbf{3.} & \symking{}f5 & b5\\
\textbf{4.} & \symking{}e4 & b6\\
\textbf{5.} & \symking{}d3 & b7\\
\textbf{6.} & \symking{}c2 & b8=\symqueen{}\\
\textbf{7.} & \symking{}c1 & \symqueen{}b1\#\\
\end{tabular}\par\vspace{3mm}
\qrcode[height=12mm]{https://helpman.komtera.lt/?fen=7b/7Q/7k/8/8/8/1P1B4/K7&moves=7}\\{\sffamily\tiny\color{ink}scannen zum Lösen}
\end{minipage}
\vspace{2mm}

{\raggedright\noindent\textcolor{ink}{\sffamily\textsc{Themen}}\enspace \themetag{Umwandlung} \themetag{Excelsior} \themetag{ein Stein zieht} \themetag{Phönix} \themetag{ohne Schach}\par\smallskip{\footnotesize\textit{Hinweis:} das Diagramm hat keinen legalen letzten Zug, kann also in keiner Partie entstehen.}\par}

\clearpage
\phantomsection\label{hm:KQBPvkn}
\begin{tcolorbox}[enhanced,colback=ink,colframe=ink,arc=0pt,boxrule=0pt,left=4mm,right=4mm,top=2.5mm,bottom=2.5mm]
\color{white}\sffamily{\LARGE\bfseries KQBPvkn}\hfill{\Large h\#7}\\[0.5mm]
{\small 6 Steine\quad Schwarz am Zug}\hfill{\small Nr.~533\qquad\href{https://helpman.komtera.lt/?fen=8/8/8/8/7B/8/3Q3P/K1k4n&moves=7}{\color{white}\bfseries lösen\,\ensuremath{\nearrow}}}
\end{tcolorbox}
\vspace{2mm}
\noindent\begin{minipage}[t]{36mm}\vspace{0pt}\sffamily\small\raggedright
\hmfigure{längste eindeutige}{{\Large h\#7}}\hmfigure{eindeutig bei dieser Länge}{{\Large 1\,271} Stellungen}\hmfigure{längstes Matt}{{\Large h\#9}}\hmfigure{in der Klasse}{45,4 Milliarden Stellungen,\\ 27 Milliarden mit Hilfsmatt}
\end{minipage}\hfill
\begin{minipage}[t]{80mm}\vspace{0pt}\begin{tikzpicture}\begin{axis}[width=72mm,height=42mm,ybar,bar width=1.6pt,xmin=-1,xmax=19,ymin=0,ymax=13.57,axis y line*=left,axis x line*=bottom,xlabel={\scriptsize Halbzüge bis Matt},ylabel={\scriptsize Stellungen (Mrd.)},tick label style={font=\tiny},label style={font=\scriptsize},enlarge x limits=false,ylabel shift=-3pt,xlabel shift=-3pt]\addplot[fill=ink!55,draw=none] coordinates {(0,0.07884) (1,0.4777) (2,1.95) (3,5.755) (4,12.34) (5,4.059) (6,2.054) (7,0.04499) (8,0.06808) (9,1.412e-05) (10,0.06539) (11,2.606e-06) (12,0.07315) (13,3.512e-06) (14,0.02252) (15,3.087e-06) (16,7.872e-05) (17,9.2e-08) (18,2e-08)};\draw[densely dotted] (axis cs:14,0) -- (axis cs:14,13.57);\end{axis}\begin{axis}[width=72mm,height=42mm,axis y line*=right,axis x line=none,ymode=log,xmin=-1,xmax=19,ylabel={\scriptsize eindeutig (log)},tick label style={font=\tiny},label style={font=\scriptsize},ylabel shift=-3pt]\addplot[mark=*,mark size=0.7pt,draw=black,thin] coordinates {(0,78836683) (1,372984285) (2,1097804711) (3,1071841675) (4,558144804) (5,6611277) (6,3225669) (7,776) (8,1466779) (9,179) (10,863126) (11,59) (12,387196) (13,66) (14,1271)};\end{axis}\end{tikzpicture}\end{minipage}
\vspace{3mm}

\noindent\begin{minipage}[t]{60mm}\vspace{0pt}\sbox\hmboard{\chessboard[setfen=8/8/8/8/7B/8/3Q3P/K1k4n b - - 0 1]}\usebox\hmboard\\[1mm]
\makebox[\wd\hmboard]{\sffamily\hspace{3mm}\textbf{h\#7}\hfill(4+2)}\end{minipage}\hfill
\begin{minipage}[t]{52mm}\vspace{0pt}
\textcolor{ink}{\sffamily\textsc{Lösung}}\par\smallskip\small
\begin{tabular}{@{}r@{\enspace}l@{\enspace}l@{}}
\textbf{1.} & \symking{}xd2 & \symbishop{}e1+\\
\textbf{2.} & \symking{}c2 & h4\\
\textbf{3.} & \symknight{}g3 & h5\\
\textbf{4.} & \symknight{}f5 & h6\\
\textbf{5.} & \symknight{}g7 & hxg7\\
\textbf{6.} & \symking{}b3 & g8=\symqueen{}+\\
\textbf{7.} & \symking{}a3 & \symqueen{}a2\#\\
\end{tabular}\par\vspace{3mm}
\qrcode[height=12mm]{https://helpman.komtera.lt/?fen=8/8/8/8/7B/8/3Q3P/K1k4n&moves=7}\\{\sffamily\tiny\color{ink}scannen zum Lösen}
\end{minipage}
\vspace{2mm}

{\raggedright\noindent\textcolor{ink}{\sffamily\textsc{Themen}}\enspace \themetag{Umwandlung} \themetag{Excelsior} \themetag{Phönix}\par\smallskip{\footnotesize\textit{Hinweis:} die Partei am Zug steht im Diagramm im Schach; die Lösung beginnt mit einem Schlag.}\par}

\clearpage
\phantomsection\label{hm:KQBPvkq}
\begin{tcolorbox}[enhanced,colback=ink,colframe=ink,arc=0pt,boxrule=0pt,left=4mm,right=4mm,top=2.5mm,bottom=2.5mm]
\color{white}\sffamily{\LARGE\bfseries KQBPvkq}\hfill{\Large h\#7}\\[0.5mm]
{\small 6 Steine\quad Schwarz am Zug}\hfill{\small Nr.~534\qquad\href{https://helpman.komtera.lt/?fen=8/8/8/8/8/P7/B1Q5/K1k3q1&moves=7}{\color{white}\bfseries lösen\,\ensuremath{\nearrow}}}
\end{tcolorbox}
\vspace{2mm}
\noindent\begin{minipage}[t]{36mm}\vspace{0pt}\sffamily\small\raggedright
\hmfigure{längste eindeutige}{{\Large h\#7}}\hmfigure{eindeutig bei dieser Länge}{{\Large 1\,425} Stellungen}\hmfigure{längstes Matt}{{\Large h\#8}}\hmfigure{in der Klasse}{45,4 Milliarden Stellungen,\\ 22,6 Milliarden mit Hilfsmatt}
\end{minipage}\hfill
\begin{minipage}[t]{80mm}\vspace{0pt}\begin{tikzpicture}\begin{axis}[width=72mm,height=42mm,ybar,bar width=1.6pt,xmin=-1,xmax=17,ymin=0,ymax=10.4,axis y line*=left,axis x line*=bottom,xlabel={\scriptsize Halbzüge bis Matt},ylabel={\scriptsize Stellungen (Mrd.)},tick label style={font=\tiny},label style={font=\scriptsize},enlarge x limits=false,ylabel shift=-3pt,xlabel shift=-3pt]\addplot[fill=ink!55,draw=none] coordinates {(0,0.04136) (1,0.2825) (2,1.423) (3,4.852) (4,9.454) (5,5.094) (6,1.189) (7,0.1002) (8,0.05424) (9,5.814e-05) (10,0.04851) (11,8.438e-06) (12,0.0534) (13,9.4e-07) (14,0.01102) (15,3.06e-07) (16,5.761e-05)};\draw[densely dotted] (axis cs:14,0) -- (axis cs:14,10.4);\end{axis}\begin{axis}[width=72mm,height=42mm,axis y line*=right,axis x line=none,ymode=log,xmin=-1,xmax=17,ylabel={\scriptsize eindeutig (log)},tick label style={font=\tiny},label style={font=\scriptsize},ylabel shift=-3pt]\addplot[mark=*,mark size=0.7pt,draw=black,thin] coordinates {(0,41361219) (1,241001469) (2,822184355) (3,924719713) (4,396799530) (5,23176012) (6,2846409) (7,4457) (8,1091532) (9,301) (10,431544) (11,94) (12,154459) (13,5) (14,1425)};\end{axis}\end{tikzpicture}\end{minipage}
\vspace{3mm}

\noindent\begin{minipage}[t]{60mm}\vspace{0pt}\sbox\hmboard{\chessboard[setfen=8/8/8/8/8/P7/B1Q5/K1k3q1 b - - 0 1]}\usebox\hmboard\\[1mm]
\makebox[\wd\hmboard]{\sffamily\hspace{3mm}\textbf{h\#7}\hfill(4+2)}\end{minipage}\hfill
\begin{minipage}[t]{52mm}\vspace{0pt}
\textcolor{ink}{\sffamily\textsc{Lösung}}\par\smallskip\small
\begin{tabular}{@{}r@{\enspace}l@{\enspace}l@{}}
\textbf{1.} & \symking{}xc2+ & \symbishop{}b1+\\
\textbf{2.} & \symking{}d2 & a4\\
\textbf{3.} & \symking{}e3 & a5\\
\textbf{4.} & \symking{}f4 & a6\\
\textbf{5.} & \symking{}g5 & a7\\
\textbf{6.} & \symking{}h6 & a8=\symqueen{}\\
\textbf{7.} & \symqueen{}g5 & \symqueen{}h8\#\\
\end{tabular}\par\vspace{3mm}
\qrcode[height=12mm]{https://helpman.komtera.lt/?fen=8/8/8/8/8/P7/B1Q5/K1k3q1&moves=7}\\{\sffamily\tiny\color{ink}scannen zum Lösen}
\end{minipage}
\vspace{2mm}

{\raggedright\noindent\textcolor{ink}{\sffamily\textsc{Themen}}\enspace \themetag{Umwandlung} \themetag{Selbstblock} \themetag{Phönix}\par\smallskip{\footnotesize\textit{Hinweis:} die Partei am Zug steht im Diagramm im Schach; die Lösung beginnt mit einem Schlag.}\par}

\clearpage
\phantomsection\label{hm:KQBPvkr}
\begin{tcolorbox}[enhanced,colback=ink,colframe=ink,arc=0pt,boxrule=0pt,left=4mm,right=4mm,top=2.5mm,bottom=2.5mm]
\color{white}\sffamily{\LARGE\bfseries KQBPvkr}\hfill{\Large h\#7}\\[0.5mm]
{\small 6 Steine\quad Schwarz am Zug}\hfill{\small Nr.~535\qquad\href{https://helpman.komtera.lt/?fen=8/8/8/8/8/P7/B1Q5/K1k3r1&moves=7}{\color{white}\bfseries lösen\,\ensuremath{\nearrow}}}
\end{tcolorbox}
\vspace{2mm}
\noindent\begin{minipage}[t]{36mm}\vspace{0pt}\sffamily\small\raggedright
\hmfigure{längste eindeutige}{{\Large h\#7}}\hmfigure{eindeutig bei dieser Länge}{{\Large 673} Stellungen}\hmfigure{längstes Matt}{{\Large h\#8}}\hmfigure{in der Klasse}{45,4 Milliarden Stellungen,\\ 24,9 Milliarden mit Hilfsmatt}
\end{minipage}\hfill
\begin{minipage}[t]{80mm}\vspace{0pt}\begin{tikzpicture}\begin{axis}[width=72mm,height=42mm,ybar,bar width=1.6pt,xmin=-1,xmax=17,ymin=0,ymax=12.39,axis y line*=left,axis x line*=bottom,xlabel={\scriptsize Halbzüge bis Matt},ylabel={\scriptsize Stellungen (Mrd.)},tick label style={font=\tiny},label style={font=\scriptsize},enlarge x limits=false,ylabel shift=-3pt,xlabel shift=-3pt]\addplot[fill=ink!55,draw=none] coordinates {(0,0.05924) (1,0.3945) (2,1.866) (3,5.565) (4,11.26) (5,4.333) (6,1.193) (7,0.04197) (8,0.06502) (9,2.013e-06) (10,0.05761) (11,6.25e-07) (12,0.06084) (13,4.98e-07) (14,0.008479) (15,1.48e-07) (16,2.535e-05)};\draw[densely dotted] (axis cs:14,0) -- (axis cs:14,12.39);\end{axis}\begin{axis}[width=72mm,height=42mm,axis y line*=right,axis x line=none,ymode=log,xmin=-1,xmax=17,ylabel={\scriptsize eindeutig (log)},tick label style={font=\tiny},label style={font=\scriptsize},ylabel shift=-3pt]\addplot[mark=*,mark size=0.7pt,draw=black,thin] coordinates {(0,59239147) (1,324962874) (2,1086392605) (3,971319275) (4,440940310) (5,9019657) (6,3581371) (7,460) (8,1366021) (9,4) (10,514473) (11,4) (12,119553) (14,673)};\end{axis}\end{tikzpicture}\end{minipage}
\vspace{3mm}

\noindent\begin{minipage}[t]{60mm}\vspace{0pt}\sbox\hmboard{\chessboard[setfen=8/8/8/8/8/P7/B1Q5/K1k3r1 b - - 0 1]}\usebox\hmboard\\[1mm]
\makebox[\wd\hmboard]{\sffamily\hspace{3mm}\textbf{h\#7}\hfill(4+2)}\end{minipage}\hfill
\begin{minipage}[t]{52mm}\vspace{0pt}
\textcolor{ink}{\sffamily\textsc{Lösung}}\par\smallskip\small
\begin{tabular}{@{}r@{\enspace}l@{\enspace}l@{}}
\textbf{1.} & \symking{}xc2+ & \symbishop{}b1+\\
\textbf{2.} & \symking{}d2 & a4\\
\textbf{3.} & \symking{}e3 & a5\\
\textbf{4.} & \symking{}f4 & a6\\
\textbf{5.} & \symking{}g5 & a7\\
\textbf{6.} & \symking{}h6 & a8=\symqueen{}\\
\textbf{7.} & \symrook{}g5 & \symqueen{}h8\#\\
\end{tabular}\par\vspace{3mm}
\qrcode[height=12mm]{https://helpman.komtera.lt/?fen=8/8/8/8/8/P7/B1Q5/K1k3r1&moves=7}\\{\sffamily\tiny\color{ink}scannen zum Lösen}
\end{minipage}
\vspace{2mm}

{\raggedright\noindent\textcolor{ink}{\sffamily\textsc{Themen}}\enspace \themetag{Umwandlung} \themetag{Selbstblock} \themetag{Phönix}\par\smallskip{\footnotesize\textit{Hinweis:} die Partei am Zug steht im Diagramm im Schach; die Lösung beginnt mit einem Schlag.}\par}

\clearpage
\phantomsection\label{hm:KQBvkqb}
\begin{tcolorbox}[enhanced,colback=ink,colframe=ink,arc=0pt,boxrule=0pt,left=4mm,right=4mm,top=2.5mm,bottom=2.5mm]
\color{white}\sffamily{\LARGE\bfseries KQBvkqb}\hfill{\Large h\#7}\\[0.5mm]
{\small 6 Steine\quad Schwarz am Zug}\hfill{\small Nr.~536\qquad\href{https://helpman.komtera.lt/?fen=8/6q1/5k2/5QB1/8/8/2b5/K7&moves=7}{\color{white}\bfseries lösen\,\ensuremath{\nearrow}}}
\end{tcolorbox}
\vspace{2mm}
\noindent\begin{minipage}[t]{36mm}\vspace{0pt}\sffamily\small\raggedright
\hmfigure{längste eindeutige}{{\Large h\#7}}\hmfigure{eindeutig bei dieser Länge}{{\Large 5} Stellungen}\hmfigure{längstes Matt}{{\Large h\#10.5}}\hmfigure{in der Klasse}{15,5 Milliarden Stellungen,\\ 7,03 Milliarden mit Hilfsmatt}
\end{minipage}\hfill
\begin{minipage}[t]{80mm}\vspace{0pt}\begin{tikzpicture}\begin{axis}[width=72mm,height=42mm,ybar,bar width=1.6pt,xmin=-1,xmax=22,ymin=0,ymax=2.976,axis y line*=left,axis x line*=bottom,xlabel={\scriptsize Halbzüge bis Matt},ylabel={\scriptsize Stellungen (Mrd.)},tick label style={font=\tiny},label style={font=\scriptsize},enlarge x limits=false,ylabel shift=-3pt,xlabel shift=-3pt]\addplot[fill=ink!55,draw=none] coordinates {(0,0.008638) (1,0.06299) (2,0.338) (3,1.386) (4,2.706) (5,2.029) (6,0.4191) (7,0.04678) (8,0.007943) (9,8.971e-06) (10,0.01299) (11,5.079e-06) (12,0.01049) (13,2.774e-06) (14,0.001186) (15,7.08e-07) (16,1.62e-06) (17,1.82e-06) (18,2.09e-07) (19,2.04e-07) (20,6.84e-07) (21,1.72e-07)};\draw[densely dotted] (axis cs:14,0) -- (axis cs:14,2.976);\end{axis}\begin{axis}[width=72mm,height=42mm,axis y line*=right,axis x line=none,ymode=log,xmin=-1,xmax=22,ylabel={\scriptsize eindeutig (log)},tick label style={font=\tiny},label style={font=\scriptsize},ylabel shift=-3pt]\addplot[mark=*,mark size=0.7pt,draw=black,thin] coordinates {(0,8638472) (1,56914763) (2,200098142) (3,298733252) (4,137175218) (5,11244247) (6,289580) (7,382) (8,234465) (9,189) (10,89952) (11,35) (12,8509) (13,5) (14,5)};\end{axis}\end{tikzpicture}\end{minipage}
\vspace{3mm}

\noindent\begin{minipage}[t]{60mm}\vspace{0pt}\sbox\hmboard{\chessboard[setfen=8/6q1/5k2/5QB1/8/8/2b5/K7 b - - 0 1]}\usebox\hmboard\\[1mm]
\makebox[\wd\hmboard]{\sffamily\hspace{3mm}\textbf{h\#7}\hfill(3+3)}\end{minipage}\hfill
\begin{minipage}[t]{52mm}\vspace{0pt}
\textcolor{ink}{\sffamily\textsc{Lösung}}\par\smallskip\small
\begin{tabular}{@{}r@{\enspace}l@{\enspace}l@{}}
\textbf{1.} & \symking{}xf5+ & \symking{}a2\\
\textbf{2.} & \symbishop{}b1+ & \symking{}b3\\
\textbf{3.} & \symking{}g6 & \symking{}c4\\
\textbf{4.} & \symking{}h7 & \symking{}d5\\
\textbf{5.} & \symking{}h8 & \symking{}e6\\
\textbf{6.} & \symqueen{}f7+ & \symking{}xf7\\
\textbf{7.} & \symbishop{}h7 & \symbishop{}f6\#\\
\end{tabular}\par\vspace{3mm}
\qrcode[height=12mm]{https://helpman.komtera.lt/?fen=8/6q1/5k2/5QB1/8/8/2b5/K7&moves=7}\\{\sffamily\tiny\color{ink}scannen zum Lösen}
\end{minipage}
\vspace{2mm}

{\raggedright\noindent\textcolor{ink}{\sffamily\textsc{Themen}}\enspace \themetag{Selbstblock}\par\smallskip{\footnotesize\textit{Hinweis:} das Diagramm hat keinen legalen letzten Zug, kann also in keiner Partie entstehen.}\par}

\clearpage
\phantomsection\label{hm:KQBvkrb}
\begin{tcolorbox}[enhanced,colback=ink,colframe=ink,arc=0pt,boxrule=0pt,left=4mm,right=4mm,top=2.5mm,bottom=2.5mm]
\color{white}\sffamily{\LARGE\bfseries KQBvkrb}\hfill{\Large h\#7}\\[0.5mm]
{\small 6 Steine\quad Schwarz am Zug}\hfill{\small Nr.~537\qquad\href{https://helpman.komtera.lt/?fen=7Q/7k/8/8/8/1r6/5b2/KB6&moves=7}{\color{white}\bfseries lösen\,\ensuremath{\nearrow}}}
\end{tcolorbox}
\vspace{2mm}
\noindent\begin{minipage}[t]{36mm}\vspace{0pt}\sffamily\small\raggedright
\hmfigure{längste eindeutige}{{\Large h\#7}}\hmfigure{eindeutig bei dieser Länge}{{\Large 20} Stellungen}\hmfigure{längstes Matt}{{\Large h\#10}}\hmfigure{in der Klasse}{15,5 Milliarden Stellungen,\\ 7,74 Milliarden mit Hilfsmatt}
\end{minipage}\hfill
\begin{minipage}[t]{80mm}\vspace{0pt}\begin{tikzpicture}\begin{axis}[width=72mm,height=42mm,ybar,bar width=1.6pt,xmin=-1,xmax=21,ymin=0,ymax=3.61,axis y line*=left,axis x line*=bottom,xlabel={\scriptsize Halbzüge bis Matt},ylabel={\scriptsize Stellungen (Mrd.)},tick label style={font=\tiny},label style={font=\scriptsize},enlarge x limits=false,ylabel shift=-3pt,xlabel shift=-3pt]\addplot[fill=ink!55,draw=none] coordinates {(0,0.0125) (1,0.08862) (2,0.4538) (3,1.627) (4,3.282) (5,1.792) (6,0.4196) (7,0.02333) (8,0.008795) (9,4.45e-07) (10,0.014) (11,2.962e-06) (12,0.01369) (13,1.323e-06) (14,0.001714) (15,1.5e-08) (16,3.915e-06) (17,1.2e-08) (18,3.4e-08) (20,3.5e-08)};\draw[densely dotted] (axis cs:14,0) -- (axis cs:14,3.61);\end{axis}\begin{axis}[width=72mm,height=42mm,axis y line*=right,axis x line=none,ymode=log,xmin=-1,xmax=21,ylabel={\scriptsize eindeutig (log)},tick label style={font=\tiny},label style={font=\scriptsize},ylabel shift=-3pt]\addplot[mark=*,mark size=0.7pt,draw=black,thin] coordinates {(0,12503013) (1,78133534) (2,271526098) (3,324940444) (4,152864143) (5,5591785) (6,264003) (7,28) (8,205012) (9,22) (10,114478) (11,203) (12,16115) (13,10) (14,20)};\end{axis}\end{tikzpicture}\end{minipage}
\vspace{3mm}

\noindent\begin{minipage}[t]{60mm}\vspace{0pt}\sbox\hmboard{\chessboard[setfen=7Q/7k/8/8/8/1r6/5b2/KB6 b - - 0 1]}\usebox\hmboard\\[1mm]
\makebox[\wd\hmboard]{\sffamily\hspace{3mm}\textbf{h\#7}\hfill(3+3)}\end{minipage}\hfill
\begin{minipage}[t]{52mm}\vspace{0pt}
\textcolor{ink}{\sffamily\textsc{Lösung}}\par\smallskip\small
\begin{tabular}{@{}r@{\enspace}l@{\enspace}l@{}}
\textbf{1.} & \symking{}xh8 & \symbishop{}c2\\
\textbf{2.} & \symrook{}h3 & \symking{}b2\\
\textbf{3.} & \symrook{}h6 & \symking{}c3\\
\textbf{4.} & \symking{}g7 & \symking{}d3\\
\textbf{5.} & \symking{}g6 & \symking{}e4\\
\textbf{6.} & \symking{}h5 & \symking{}f5\\
\textbf{7.} & \symbishop{}h4 & \symbishop{}d1\#\\
\end{tabular}\par\vspace{3mm}
\qrcode[height=12mm]{https://helpman.komtera.lt/?fen=7Q/7k/8/8/8/1r6/5b2/KB6&moves=7}\\{\sffamily\tiny\color{ink}scannen zum Lösen}
\end{minipage}
\vspace{2mm}

{\raggedright\noindent\textcolor{ink}{\sffamily\textsc{Themen}}\enspace \themetag{Selbstblock} \themetag{ohne Schach}\par\smallskip{\footnotesize\textit{Hinweis:} das Diagramm hat keinen legalen letzten Zug, kann also in keiner Partie entstehen.}\par}

\clearpage
\phantomsection\label{hm:KQNNPvk}
\begin{tcolorbox}[enhanced,colback=ink,colframe=ink,arc=0pt,boxrule=0pt,left=4mm,right=4mm,top=2.5mm,bottom=2.5mm]
\color{white}\sffamily{\LARGE\bfseries KQNNPvk}\hfill{\Large h\#7}\\[0.5mm]
{\small 6 Steine\quad Schwarz am Zug}\hfill{\small Nr.~538\qquad\href{https://helpman.komtera.lt/?fen=7Q/4N1k1/3N4/8/8/8/4P3/K7&moves=7}{\color{white}\bfseries lösen\,\ensuremath{\nearrow}}}
\end{tcolorbox}
\vspace{2mm}
\noindent\begin{minipage}[t]{36mm}\vspace{0pt}\sffamily\small\raggedright
\hmfigure{längste eindeutige}{{\Large h\#7}}\hmfigure{eindeutig bei dieser Länge}{{\Large 45} Stellungen}\hmfigure{längstes Matt}{{\Large h\#8}}\hmfigure{in der Klasse}{45,4 Milliarden Stellungen,\\ 14 Milliarden mit Hilfsmatt}
\end{minipage}\hfill
\begin{minipage}[t]{80mm}\vspace{0pt}\begin{tikzpicture}\begin{axis}[width=72mm,height=42mm,ybar,bar width=1.6pt,xmin=-1,xmax=17,ymin=0,ymax=6.832,axis y line*=left,axis x line*=bottom,xlabel={\scriptsize Halbzüge bis Matt},ylabel={\scriptsize Stellungen (Mrd.)},tick label style={font=\tiny},label style={font=\scriptsize},enlarge x limits=false,ylabel shift=-3pt,xlabel shift=-3pt]\addplot[fill=ink!55,draw=none] coordinates {(0,0.05233) (1,0.3403) (2,1.165) (3,2.666) (4,6.211) (5,2.003) (6,1.438) (7,0.05151) (8,0.05927) (9,5.15e-07) (10,0.0404) (12,0.0169) (14,0.003077) (16,6.028e-06)};\draw[densely dotted] (axis cs:14,0) -- (axis cs:14,6.832);\end{axis}\begin{axis}[width=72mm,height=42mm,axis y line*=right,axis x line=none,ymode=log,xmin=-1,xmax=17,ylabel={\scriptsize eindeutig (log)},tick label style={font=\tiny},label style={font=\scriptsize},ylabel shift=-3pt]\addplot[mark=*,mark size=0.7pt,draw=black,thin] coordinates {(0,52334510) (1,259281136) (2,806025837) (3,465513441) (4,382695692) (5,6025108) (6,4956219) (7,5) (8,943848) (10,218543) (12,35504) (14,45)};\end{axis}\end{tikzpicture}\end{minipage}
\vspace{3mm}

\noindent\begin{minipage}[t]{60mm}\vspace{0pt}\sbox\hmboard{\chessboard[setfen=7Q/4N1k1/3N4/8/8/8/4P3/K7 b - - 0 1]}\usebox\hmboard\\[1mm]
\makebox[\wd\hmboard]{\sffamily\hspace{3mm}\textbf{h\#7}\hfill(5+1)}\end{minipage}\hfill
\begin{minipage}[t]{52mm}\vspace{0pt}
\textcolor{ink}{\sffamily\textsc{Lösung}}\par\smallskip\small
\begin{tabular}{@{}r@{\enspace}l@{\enspace}l@{}}
\textbf{1.} & \symking{}xh8 & e4\\
\textbf{2.} & \symking{}h7 & e5\\
\textbf{3.} & \symking{}g7 & e6\\
\textbf{4.} & \symking{}f6 & \symknight{}ef5\\
\textbf{5.} & \symking{}e5 & e7\\
\textbf{6.} & \symking{}d5 & e8=\symqueen{}\\
\textbf{7.} & \symking{}c5 & \symqueen{}b5\#\\
\end{tabular}\par\vspace{3mm}
\qrcode[height=12mm]{https://helpman.komtera.lt/?fen=7Q/4N1k1/3N4/8/8/8/4P3/K7&moves=7}\\{\sffamily\tiny\color{ink}scannen zum Lösen}
\end{minipage}
\vspace{2mm}

{\raggedright\noindent\textcolor{ink}{\sffamily\textsc{Themen}}\enspace \themetag{Umwandlung} \themetag{Excelsior} \themetag{Rundlauf} \themetag{ein Stein zieht} \themetag{Phönix} \themetag{ohne Schach}\par\smallskip{\footnotesize\textit{Hinweis:} die Partei am Zug steht im Diagramm im Schach; die Lösung beginnt mit einem Schlag.}\par}

\clearpage
\phantomsection\label{hm:KQNNvkb}
\begin{tcolorbox}[enhanced,colback=ink,colframe=ink,arc=0pt,boxrule=0pt,left=4mm,right=4mm,top=2.5mm,bottom=2.5mm]
\color{white}\sffamily{\LARGE\bfseries KQNNvkb}\hfill{\Large h\#7}\\[0.5mm]
{\small 6 Steine\quad Schwarz am Zug}\hfill{\small Nr.~539\qquad\href{https://helpman.komtera.lt/?fen=kQ5N/8/8/8/8/8/7N/b2K4&moves=7}{\color{white}\bfseries lösen\,\ensuremath{\nearrow}}}
\end{tcolorbox}
\vspace{2mm}
\noindent\begin{minipage}[t]{36mm}\vspace{0pt}\sffamily\small\raggedright
\hmfigure{längste eindeutige}{{\Large h\#7}}\hmfigure{eindeutig bei dieser Länge}{{\Large 17} Stellungen}\hmfigure{längstes Matt}{{\Large h\#9}}\hmfigure{in der Klasse}{15,5 Milliarden Stellungen,\\ 4,36 Milliarden mit Hilfsmatt}
\end{minipage}\hfill
\begin{minipage}[t]{80mm}\vspace{0pt}\begin{tikzpicture}\begin{axis}[width=72mm,height=42mm,ybar,bar width=1.6pt,xmin=-1,xmax=19,ymin=0,ymax=2.039,axis y line*=left,axis x line*=bottom,xlabel={\scriptsize Halbzüge bis Matt},ylabel={\scriptsize Stellungen (Mrd.)},tick label style={font=\tiny},label style={font=\scriptsize},enlarge x limits=false,ylabel shift=-3pt,xlabel shift=-3pt]\addplot[fill=ink!55,draw=none] coordinates {(0,0.01097) (1,0.07741) (2,0.3391) (3,0.889) (4,1.854) (5,0.733) (6,0.3928) (7,0.02855) (8,0.0165) (9,2.862e-05) (10,0.01596) (11,1.01e-07) (12,0.005598) (13,2.3e-08) (14,0.001015) (16,8.354e-06) (18,2e-09)};\draw[densely dotted] (axis cs:14,0) -- (axis cs:14,2.039);\end{axis}\begin{axis}[width=72mm,height=42mm,axis y line*=right,axis x line=none,ymode=log,xmin=-1,xmax=19,ylabel={\scriptsize eindeutig (log)},tick label style={font=\tiny},label style={font=\scriptsize},ylabel shift=-3pt]\addplot[mark=*,mark size=0.7pt,draw=black,thin] coordinates {(0,10974816) (1,64505314) (2,194363070) (3,128639713) (4,86437325) (5,3552203) (6,608844) (7,2193) (8,147052) (10,23016) (12,298) (14,17)};\end{axis}\end{tikzpicture}\end{minipage}
\vspace{3mm}

\noindent\begin{minipage}[t]{60mm}\vspace{0pt}\sbox\hmboard{\chessboard[setfen=kQ5N/8/8/8/8/8/7N/b2K4 b - - 0 1]}\usebox\hmboard\\[1mm]
\makebox[\wd\hmboard]{\sffamily\hspace{3mm}\textbf{h\#7}\hfill(4+2)}\end{minipage}\hfill
\begin{minipage}[t]{52mm}\vspace{0pt}
\textcolor{ink}{\sffamily\textsc{Lösung}}\par\smallskip\small
\begin{tabular}{@{}r@{\enspace}l@{\enspace}l@{}}
\textbf{1.} & \symking{}xb8 & \symking{}e2\\
\textbf{2.} & \symking{}c7 & \symking{}f3\\
\textbf{3.} & \symking{}d6 & \symking{}g4\\
\textbf{4.} & \symking{}e5 & \symking{}h5\\
\textbf{5.} & \symking{}f4 & \symknight{}g6+\\
\textbf{6.} & \symking{}g3 & \symknight{}f1+\\
\textbf{7.} & \symking{}h3 & \symknight{}f4\#\\
\end{tabular}\par\vspace{3mm}
\qrcode[height=12mm]{https://helpman.komtera.lt/?fen=kQ5N/8/8/8/8/8/7N/b2K4&moves=7}\\{\sffamily\tiny\color{ink}scannen zum Lösen}
\end{minipage}
\vspace{2mm}

{\raggedright\noindent\textcolor{ink}{\sffamily\textsc{Themen}}\enspace \themetag{reines Matt} \themetag{Mustermatt} \themetag{Spiegelmatt} \themetag{ein Stein zieht}\par\smallskip{\footnotesize\textit{Hinweis:} die Partei am Zug steht im Diagramm im Schach; die Lösung beginnt mit einem Schlag.}\par}

\clearpage
\phantomsection\label{hm:KQNNvkp}
\begin{tcolorbox}[enhanced,colback=ink,colframe=ink,arc=0pt,boxrule=0pt,left=4mm,right=4mm,top=2.5mm,bottom=2.5mm]
\color{white}\sffamily{\LARGE\bfseries KQNNvkp}\hfill{\Large h\#7}\\[0.5mm]
{\small 6 Steine\quad Schwarz am Zug}\hfill{\small Nr.~540\qquad\href{https://helpman.komtera.lt/?fen=7Q/6Nk/6p1/8/8/8/8/K5N1&moves=7}{\color{white}\bfseries lösen\,\ensuremath{\nearrow}}}
\end{tcolorbox}
\vspace{2mm}
\noindent\begin{minipage}[t]{36mm}\vspace{0pt}\sffamily\small\raggedright
\hmfigure{längste eindeutige}{{\Large h\#7}}\hmfigure{eindeutig bei dieser Länge}{{\Large 923} Stellungen}\hmfigure{längstes Matt}{{\Large h\#9}}\hmfigure{in der Klasse}{45,4 Milliarden Stellungen,\\ 14 Milliarden mit Hilfsmatt}
\end{minipage}\hfill
\begin{minipage}[t]{80mm}\vspace{0pt}\begin{tikzpicture}\begin{axis}[width=72mm,height=42mm,ybar,bar width=1.6pt,xmin=-1,xmax=19,ymin=0,ymax=6.578,axis y line*=left,axis x line*=bottom,xlabel={\scriptsize Halbzüge bis Matt},ylabel={\scriptsize Stellungen (Mrd.)},tick label style={font=\tiny},label style={font=\scriptsize},enlarge x limits=false,ylabel shift=-3pt,xlabel shift=-3pt]\addplot[fill=ink!55,draw=none] coordinates {(0,0.04458) (1,0.2985) (2,1.157) (3,2.778) (4,5.98) (5,2.054) (6,1.49) (7,0.0756) (8,0.05078) (9,4.279e-05) (10,0.05275) (11,2e-09) (12,0.02767) (14,0.007776) (16,6.881e-05) (18,6e-09)};\draw[densely dotted] (axis cs:14,0) -- (axis cs:14,6.578);\end{axis}\begin{axis}[width=72mm,height=42mm,axis y line*=right,axis x line=none,ymode=log,xmin=-1,xmax=19,ylabel={\scriptsize eindeutig (log)},tick label style={font=\tiny},label style={font=\scriptsize},ylabel shift=-3pt]\addplot[mark=*,mark size=0.7pt,draw=black,thin] coordinates {(0,44582912) (1,238300611) (2,755175902) (3,451645098) (4,301076031) (5,7078563) (6,2124652) (7,2506) (8,518781) (10,137388) (12,13043) (14,923)};\end{axis}\end{tikzpicture}\end{minipage}
\vspace{3mm}

\noindent\begin{minipage}[t]{60mm}\vspace{0pt}\sbox\hmboard{\chessboard[setfen=7Q/6Nk/6p1/8/8/8/8/K5N1 b - - 0 1]}\usebox\hmboard\\[1mm]
\makebox[\wd\hmboard]{\sffamily\hspace{3mm}\textbf{h\#7}\hfill(4+2)}\end{minipage}\hfill
\begin{minipage}[t]{52mm}\vspace{0pt}
\textcolor{ink}{\sffamily\textsc{Lösung}}\par\smallskip\small
\begin{tabular}{@{}r@{\enspace}l@{\enspace}l@{}}
\textbf{1.} & \symking{}xh8 & \symking{}b2\\
\textbf{2.} & \symking{}h7 & \symking{}c3\\
\textbf{3.} & \symking{}h6 & \symking{}d4\\
\textbf{4.} & \symking{}g5 & \symking{}e5\\
\textbf{5.} & \symking{}h4 & \symking{}f6\\
\textbf{6.} & g5 & \symking{}g6\\
\textbf{7.} & g4 & \symknight{}f5\#\\
\end{tabular}\par\vspace{3mm}
\qrcode[height=12mm]{https://helpman.komtera.lt/?fen=7Q/6Nk/6p1/8/8/8/8/K5N1&moves=7}\\{\sffamily\tiny\color{ink}scannen zum Lösen}
\end{minipage}
\vspace{2mm}

{\raggedright\noindent\textcolor{ink}{\sffamily\textsc{Themen}}\enspace \themetag{reines Matt} \themetag{Mustermatt} \themetag{Idealmatt} \themetag{Switchback} \themetag{Selbstblock} \themetag{ohne Schach} \themetag{Umnov}\par\smallskip{\footnotesize\textit{Hinweis:} die Partei am Zug steht im Diagramm im Schach; die Lösung beginnt mit einem Schlag.}\par}

\clearpage
\phantomsection\label{hm:KQNNvkq}
\begin{tcolorbox}[enhanced,colback=ink,colframe=ink,arc=0pt,boxrule=0pt,left=4mm,right=4mm,top=2.5mm,bottom=2.5mm]
\color{white}\sffamily{\LARGE\bfseries KQNNvkq}\hfill{\Large h\#7}\\[0.5mm]
{\small 6 Steine\quad Schwarz am Zug}\hfill{\small Nr.~541\qquad\href{https://helpman.komtera.lt/?fen=8/7Q/7k/8/1q6/7N/8/KN6&moves=7}{\color{white}\bfseries lösen\,\ensuremath{\nearrow}}}
\end{tcolorbox}
\vspace{2mm}
\noindent\begin{minipage}[t]{36mm}\vspace{0pt}\sffamily\small\raggedright
\hmfigure{längste eindeutige}{{\Large h\#7}}\hmfigure{eindeutig bei dieser Länge}{{\Large 43} Stellungen}\hmfigure{längstes Matt}{{\Large h\#9}}\hmfigure{in der Klasse}{15,5 Milliarden Stellungen,\\ 3,77 Milliarden mit Hilfsmatt}
\end{minipage}\hfill
\begin{minipage}[t]{80mm}\vspace{0pt}\begin{tikzpicture}\begin{axis}[width=72mm,height=42mm,ybar,bar width=1.6pt,xmin=-1,xmax=19,ymin=0,ymax=1.597,axis y line*=left,axis x line*=bottom,xlabel={\scriptsize Halbzüge bis Matt},ylabel={\scriptsize Stellungen (Mrd.)},tick label style={font=\tiny},label style={font=\scriptsize},enlarge x limits=false,ylabel shift=-3pt,xlabel shift=-3pt]\addplot[fill=ink!55,draw=none] coordinates {(0,0.006484) (1,0.04849) (2,0.2436) (3,0.7602) (4,1.451) (5,0.8661) (6,0.3048) (7,0.05186) (8,0.01033) (9,0.0001075) (10,0.01284) (11,3.4e-08) (12,0.00684) (14,0.001888) (16,3.072e-05) (18,2e-09)};\draw[densely dotted] (axis cs:14,0) -- (axis cs:14,1.597);\end{axis}\begin{axis}[width=72mm,height=42mm,axis y line*=right,axis x line=none,ymode=log,xmin=-1,xmax=19,ylabel={\scriptsize eindeutig (log)},tick label style={font=\tiny},label style={font=\scriptsize},ylabel shift=-3pt]\addplot[mark=*,mark size=0.7pt,draw=black,thin] coordinates {(0,6483629) (1,42585153) (2,140978385) (3,123883526) (4,60373351) (5,6640386) (6,290319) (7,9937) (8,140532) (10,38962) (12,2309) (14,43)};\end{axis}\end{tikzpicture}\end{minipage}
\vspace{3mm}

\noindent\begin{minipage}[t]{60mm}\vspace{0pt}\sbox\hmboard{\chessboard[setfen=8/7Q/7k/8/1q6/7N/8/KN6 b - - 0 1]}\usebox\hmboard\\[1mm]
\makebox[\wd\hmboard]{\sffamily\hspace{3mm}\textbf{h\#7}\hfill(4+2)}\end{minipage}\hfill
\begin{minipage}[t]{52mm}\vspace{0pt}
\textcolor{ink}{\sffamily\textsc{Lösung}}\par\smallskip\small
\begin{tabular}{@{}r@{\enspace}l@{\enspace}l@{}}
\textbf{1.} & \symking{}xh7 & \symknight{}f2\\
\textbf{2.} & \symking{}g6 & \symknight{}d3\\
\textbf{3.} & \symking{}f5 & \symknight{}xb4\\
\textbf{4.} & \symking{}e4 & \symknight{}d5\\
\textbf{5.} & \symking{}d3 & \symknight{}d2\\
\textbf{6.} & \symking{}c2 & \symknight{}e3+\\
\textbf{7.} & \symking{}c1 & \symknight{}b3\#\\
\end{tabular}\par\vspace{3mm}
\qrcode[height=12mm]{https://helpman.komtera.lt/?fen=8/7Q/7k/8/1q6/7N/8/KN6&moves=7}\\{\sffamily\tiny\color{ink}scannen zum Lösen}
\end{minipage}
\vspace{2mm}

{\raggedright\noindent\textcolor{ink}{\sffamily\textsc{Themen}}\enspace \themetag{reines Matt} \themetag{Mustermatt} \themetag{Idealmatt} \themetag{Spiegelmatt} \themetag{ein Stein zieht}\par\smallskip{\footnotesize\textit{Hinweis:} die Partei am Zug steht im Diagramm im Schach; die Lösung beginnt mit einem Schlag.}\par}

\clearpage
\phantomsection\label{hm:KQNNvkr}
\begin{tcolorbox}[enhanced,colback=ink,colframe=ink,arc=0pt,boxrule=0pt,left=4mm,right=4mm,top=2.5mm,bottom=2.5mm]
\color{white}\sffamily{\LARGE\bfseries KQNNvkr}\hfill{\Large h\#7}\\[0.5mm]
{\small 6 Steine\quad Schwarz am Zug}\hfill{\small Nr.~542\qquad\href{https://helpman.komtera.lt/?fen=8/7Q/7k/8/1r6/7N/8/KN6&moves=7}{\color{white}\bfseries lösen\,\ensuremath{\nearrow}}}
\end{tcolorbox}
\vspace{2mm}
\noindent\begin{minipage}[t]{36mm}\vspace{0pt}\sffamily\small\raggedright
\hmfigure{längste eindeutige}{{\Large h\#7}}\hmfigure{eindeutig bei dieser Länge}{{\Large 6} Stellungen}\hmfigure{längstes Matt}{{\Large h\#8}}\hmfigure{in der Klasse}{15,5 Milliarden Stellungen,\\ 4,15 Milliarden mit Hilfsmatt}
\end{minipage}\hfill
\begin{minipage}[t]{80mm}\vspace{0pt}\begin{tikzpicture}\begin{axis}[width=72mm,height=42mm,ybar,bar width=1.6pt,xmin=-1,xmax=17,ymin=0,ymax=1.952,axis y line*=left,axis x line*=bottom,xlabel={\scriptsize Halbzüge bis Matt},ylabel={\scriptsize Stellungen (Mrd.)},tick label style={font=\tiny},label style={font=\scriptsize},enlarge x limits=false,ylabel shift=-3pt,xlabel shift=-3pt]\addplot[fill=ink!55,draw=none] coordinates {(0,0.009563) (1,0.06917) (2,0.3266) (3,0.8762) (4,1.775) (5,0.7614) (6,0.2804) (7,0.02094) (8,0.01387) (9,3.604e-06) (10,0.01335) (11,2.7e-08) (12,0.005401) (14,0.001175) (16,9.104e-06)};\draw[densely dotted] (axis cs:14,0) -- (axis cs:14,1.952);\end{axis}\begin{axis}[width=72mm,height=42mm,axis y line*=right,axis x line=none,ymode=log,xmin=-1,xmax=17,ylabel={\scriptsize eindeutig (log)},tick label style={font=\tiny},label style={font=\scriptsize},ylabel shift=-3pt]\addplot[mark=*,mark size=0.7pt,draw=black,thin] coordinates {(0,9562588) (1,58760283) (2,190336779) (3,125059787) (4,69397512) (5,3374151) (6,477501) (7,461) (8,149405) (10,25575) (12,1096) (14,6)};\end{axis}\end{tikzpicture}\end{minipage}
\vspace{3mm}

\noindent\begin{minipage}[t]{60mm}\vspace{0pt}\sbox\hmboard{\chessboard[setfen=8/7Q/7k/8/1r6/7N/8/KN6 b - - 0 1]}\usebox\hmboard\\[1mm]
\makebox[\wd\hmboard]{\sffamily\hspace{3mm}\textbf{h\#7}\hfill(4+2)}\end{minipage}\hfill
\begin{minipage}[t]{52mm}\vspace{0pt}
\textcolor{ink}{\sffamily\textsc{Lösung}}\par\smallskip\small
\begin{tabular}{@{}r@{\enspace}l@{\enspace}l@{}}
\textbf{1.} & \symking{}xh7 & \symknight{}f2\\
\textbf{2.} & \symking{}g6 & \symknight{}d3\\
\textbf{3.} & \symking{}f5 & \symknight{}xb4\\
\textbf{4.} & \symking{}e4 & \symknight{}d5\\
\textbf{5.} & \symking{}d3 & \symknight{}d2\\
\textbf{6.} & \symking{}c2 & \symknight{}e3+\\
\textbf{7.} & \symking{}c1 & \symknight{}b3\#\\
\end{tabular}\par\vspace{3mm}
\qrcode[height=12mm]{https://helpman.komtera.lt/?fen=8/7Q/7k/8/1r6/7N/8/KN6&moves=7}\\{\sffamily\tiny\color{ink}scannen zum Lösen}
\end{minipage}
\vspace{2mm}

{\raggedright\noindent\textcolor{ink}{\sffamily\textsc{Themen}}\enspace \themetag{reines Matt} \themetag{Mustermatt} \themetag{Idealmatt} \themetag{Spiegelmatt} \themetag{ein Stein zieht}\par\smallskip{\footnotesize\textit{Hinweis:} die Partei am Zug steht im Diagramm im Schach; die Lösung beginnt mit einem Schlag.}\par}

\clearpage
\phantomsection\label{hm:KQNPPvk}
\begin{tcolorbox}[enhanced,colback=ink,colframe=ink,arc=0pt,boxrule=0pt,left=4mm,right=4mm,top=2.5mm,bottom=2.5mm]
\color{white}\sffamily{\LARGE\bfseries KQNPPvk}\hfill{\Large h\#7}\\[0.5mm]
{\small 6 Steine\quad Schwarz am Zug}\hfill{\small Nr.~543\qquad\href{https://helpman.komtera.lt/?fen=8/8/8/2N5/8/2P5/2P5/K5Qk&moves=7}{\color{white}\bfseries lösen\,\ensuremath{\nearrow}}}
\end{tcolorbox}
\vspace{2mm}
\noindent\begin{minipage}[t]{36mm}\vspace{0pt}\sffamily\small\raggedright
\hmfigure{längste eindeutige}{{\Large h\#7}}\hmfigure{eindeutig bei dieser Länge}{{\Large 134} Stellungen}\hmfigure{längstes Matt}{{\Large h\#8}}\hmfigure{in der Klasse}{34,1 Milliarden Stellungen,\\ 10,7 Milliarden mit Hilfsmatt}
\end{minipage}\hfill
\begin{minipage}[t]{80mm}\vspace{0pt}\begin{tikzpicture}\begin{axis}[width=72mm,height=42mm,ybar,bar width=1.6pt,xmin=-1,xmax=17,ymin=0,ymax=4.695,axis y line*=left,axis x line*=bottom,xlabel={\scriptsize Halbzüge bis Matt},ylabel={\scriptsize Stellungen (Mrd.)},tick label style={font=\tiny},label style={font=\scriptsize},enlarge x limits=false,ylabel shift=-3pt,xlabel shift=-3pt]\addplot[fill=ink!55,draw=none] coordinates {(0,0.02809) (1,0.1987) (2,0.6312) (3,1.752) (4,4.268) (5,1.917) (6,1.673) (7,0.1291) (8,0.08149) (9,0.0001965) (10,0.02054) (12,0.01328) (14,0.003585) (16,1.302e-05)};\draw[densely dotted] (axis cs:14,0) -- (axis cs:14,4.695);\end{axis}\begin{axis}[width=72mm,height=42mm,axis y line*=right,axis x line=none,ymode=log,xmin=-1,xmax=17,ylabel={\scriptsize eindeutig (log)},tick label style={font=\tiny},label style={font=\scriptsize},ylabel shift=-3pt]\addplot[mark=*,mark size=0.7pt,draw=black,thin] coordinates {(0,28089672) (1,141523734) (2,425827059) (3,334339926) (4,351344117) (5,16419677) (6,9011820) (7,14479) (8,652330) (10,205409) (12,58168) (14,134)};\end{axis}\end{tikzpicture}\end{minipage}
\vspace{3mm}

\noindent\begin{minipage}[t]{60mm}\vspace{0pt}\sbox\hmboard{\chessboard[setfen=8/8/8/2N5/8/2P5/2P5/K5Qk b - - 0 1]}\usebox\hmboard\\[1mm]
\makebox[\wd\hmboard]{\sffamily\hspace{3mm}\textbf{h\#7}\hfill(5+1)}\end{minipage}\hfill
\begin{minipage}[t]{52mm}\vspace{0pt}
\textcolor{ink}{\sffamily\textsc{Lösung}}\par\smallskip\small
\begin{tabular}{@{}r@{\enspace}l@{\enspace}l@{}}
\textbf{1.} & \symking{}xg1 & c4\\
\textbf{2.} & \symking{}f2 & \symknight{}d3+\\
\textbf{3.} & \symking{}e3 & c5\\
\textbf{4.} & \symking{}d4 & c6\\
\textbf{5.} & \symking{}c4 & c7\\
\textbf{6.} & \symking{}b5 & c8=\symqueen{}\\
\textbf{7.} & \symking{}a4 & \symqueen{}a6\#\\
\end{tabular}\par\vspace{3mm}
\qrcode[height=12mm]{https://helpman.komtera.lt/?fen=8/8/8/2N5/8/2P5/2P5/K5Qk&moves=7}\\{\sffamily\tiny\color{ink}scannen zum Lösen}
\end{minipage}
\vspace{2mm}

{\raggedright\noindent\textcolor{ink}{\sffamily\textsc{Themen}}\enspace \themetag{reines Matt} \themetag{Mustermatt} \themetag{Spiegelmatt} \themetag{Umwandlung} \themetag{ein Stein zieht} \themetag{Phönix}\par\smallskip{\footnotesize\textit{Hinweis:} die Partei am Zug steht im Diagramm im Schach; die Lösung beginnt mit einem Schlag.}\par}

\clearpage
\phantomsection\label{hm:KQNPvkb}
\begin{tcolorbox}[enhanced,colback=ink,colframe=ink,arc=0pt,boxrule=0pt,left=4mm,right=4mm,top=2.5mm,bottom=2.5mm]
\color{white}\sffamily{\LARGE\bfseries KQNPvkb}\hfill{\Large h\#7}\\[0.5mm]
{\small 6 Steine\quad Schwarz am Zug}\hfill{\small Nr.~544\qquad\href{https://helpman.komtera.lt/?fen=N7/8/8/8/8/8/P2b4/K3kQ2&moves=7}{\color{white}\bfseries lösen\,\ensuremath{\nearrow}}}
\end{tcolorbox}
\vspace{2mm}
\noindent\begin{minipage}[t]{36mm}\vspace{0pt}\sffamily\small\raggedright
\hmfigure{längste eindeutige}{{\Large h\#7}}\hmfigure{eindeutig bei dieser Länge}{{\Large 3\,278} Stellungen}\hmfigure{längstes Matt}{{\Large h\#9}}\hmfigure{in der Klasse}{45,4 Milliarden Stellungen,\\ 26,8 Milliarden mit Hilfsmatt}
\end{minipage}\hfill
\begin{minipage}[t]{80mm}\vspace{0pt}\begin{tikzpicture}\begin{axis}[width=72mm,height=42mm,ybar,bar width=1.6pt,xmin=-1,xmax=19,ymin=0,ymax=11.06,axis y line*=left,axis x line*=bottom,xlabel={\scriptsize Halbzüge bis Matt},ylabel={\scriptsize Stellungen (Mrd.)},tick label style={font=\tiny},label style={font=\scriptsize},enlarge x limits=false,ylabel shift=-3pt,xlabel shift=-3pt]\addplot[fill=ink!55,draw=none] coordinates {(0,0.04494) (1,0.3425) (2,1.399) (3,4.432) (4,10.05) (5,5.682) (6,4.051) (7,0.5271) (8,0.1661) (9,0.002314) (10,0.05502) (11,4.76e-07) (12,0.05565) (13,5.07e-07) (14,0.0188) (15,5.17e-07) (16,0.0001727) (17,2.9e-08) (18,8e-09)};\draw[densely dotted] (axis cs:14,0) -- (axis cs:14,11.06);\end{axis}\begin{axis}[width=72mm,height=42mm,axis y line*=right,axis x line=none,ymode=log,xmin=-1,xmax=19,ylabel={\scriptsize eindeutig (log)},tick label style={font=\tiny},label style={font=\scriptsize},ylabel shift=-3pt]\addplot[mark=*,mark size=0.7pt,draw=black,thin] coordinates {(0,44935242) (1,272992708) (2,783382763) (3,726866789) (4,659316438) (5,61753547) (6,14372728) (7,184665) (8,1218574) (9,8) (10,557923) (11,10) (12,242379) (14,3278)};\end{axis}\end{tikzpicture}\end{minipage}
\vspace{3mm}

\noindent\begin{minipage}[t]{60mm}\vspace{0pt}\sbox\hmboard{\chessboard[setfen=N7/8/8/8/8/8/P2b4/K3kQ2 b - - 0 1]}\usebox\hmboard\\[1mm]
\makebox[\wd\hmboard]{\sffamily\hspace{3mm}\textbf{h\#7}\hfill(4+2)}\end{minipage}\hfill
\begin{minipage}[t]{52mm}\vspace{0pt}
\textcolor{ink}{\sffamily\textsc{Lösung}}\par\smallskip\small
\begin{tabular}{@{}r@{\enspace}l@{\enspace}l@{}}
\textbf{1.} & \symking{}xf1 & a3\\
\textbf{2.} & \symbishop{}b4 & axb4\\
\textbf{3.} & \symking{}e2 & b5\\
\textbf{4.} & \symking{}d3 & b6\\
\textbf{5.} & \symking{}c4 & b7\\
\textbf{6.} & \symking{}b5 & b8=\symqueen{}+\\
\textbf{7.} & \symking{}a6 & \symqueen{}b6\#\\
\end{tabular}\par\vspace{3mm}
\qrcode[height=12mm]{https://helpman.komtera.lt/?fen=N7/8/8/8/8/8/P2b4/K3kQ2&moves=7}\\{\sffamily\tiny\color{ink}scannen zum Lösen}
\end{minipage}
\vspace{2mm}

{\raggedright\noindent\textcolor{ink}{\sffamily\textsc{Themen}}\enspace \themetag{reines Matt} \themetag{Mustermatt} \themetag{Umwandlung} \themetag{Excelsior} \themetag{Rundlauf} \themetag{ein Stein zieht} \themetag{Phönix}\par\smallskip{\footnotesize\textit{Hinweis:} die Partei am Zug steht im Diagramm im Schach; die Lösung beginnt mit einem Schlag.}\par}

\clearpage
\phantomsection\label{hm:KQNPvkr}
\begin{tcolorbox}[enhanced,colback=ink,colframe=ink,arc=0pt,boxrule=0pt,left=4mm,right=4mm,top=2.5mm,bottom=2.5mm]
\color{white}\sffamily{\LARGE\bfseries KQNPvkr}\hfill{\Large h\#7}\\[0.5mm]
{\small 6 Steine\quad Schwarz am Zug}\hfill{\small Nr.~545\qquad\href{https://helpman.komtera.lt/?fen=8/8/Nr6/8/8/6Qk/6P1/K7&moves=7}{\color{white}\bfseries lösen\,\ensuremath{\nearrow}}}
\end{tcolorbox}
\vspace{2mm}
\noindent\begin{minipage}[t]{36mm}\vspace{0pt}\sffamily\small\raggedright
\hmfigure{längste eindeutige}{{\Large h\#7}}\hmfigure{eindeutig bei dieser Länge}{{\Large 844} Stellungen}\hmfigure{längstes Matt}{{\Large h\#8.5}}\hmfigure{in der Klasse}{45,4 Milliarden Stellungen,\\ 25,6 Milliarden mit Hilfsmatt}
\end{minipage}\hfill
\begin{minipage}[t]{80mm}\vspace{0pt}\begin{tikzpicture}\begin{axis}[width=72mm,height=42mm,ybar,bar width=1.6pt,xmin=-1,xmax=18,ymin=0,ymax=10.86,axis y line*=left,axis x line*=bottom,xlabel={\scriptsize Halbzüge bis Matt},ylabel={\scriptsize Stellungen (Mrd.)},tick label style={font=\tiny},label style={font=\scriptsize},enlarge x limits=false,ylabel shift=-3pt,xlabel shift=-3pt]\addplot[fill=ink!55,draw=none] coordinates {(0,0.03878) (1,0.3062) (2,1.378) (3,4.408) (4,9.877) (5,5.873) (6,3.088) (7,0.3955) (8,0.08725) (9,0.0004141) (10,0.05432) (11,6.3e-08) (12,0.04308) (13,1.54e-07) (14,0.00757) (15,6e-08) (16,5.505e-05) (17,1.3e-08)};\draw[densely dotted] (axis cs:14,0) -- (axis cs:14,10.86);\end{axis}\begin{axis}[width=72mm,height=42mm,axis y line*=right,axis x line=none,ymode=log,xmin=-1,xmax=18,ylabel={\scriptsize eindeutig (log)},tick label style={font=\tiny},label style={font=\scriptsize},ylabel shift=-3pt]\addplot[mark=*,mark size=0.7pt,draw=black,thin] coordinates {(0,38784939) (1,250429515) (2,793082981) (3,706202823) (4,558082174) (5,60915954) (6,7395295) (7,47008) (8,1085312) (9,17) (10,355383) (12,84497) (14,844)};\end{axis}\end{tikzpicture}\end{minipage}
\vspace{3mm}

\noindent\begin{minipage}[t]{60mm}\vspace{0pt}\sbox\hmboard{\chessboard[setfen=8/8/Nr6/8/8/6Qk/6P1/K7 b - - 0 1]}\usebox\hmboard\\[1mm]
\makebox[\wd\hmboard]{\sffamily\hspace{3mm}\textbf{h\#7}\hfill(4+2)}\end{minipage}\hfill
\begin{minipage}[t]{52mm}\vspace{0pt}
\textcolor{ink}{\sffamily\textsc{Lösung}}\par\smallskip\small
\begin{tabular}{@{}r@{\enspace}l@{\enspace}l@{}}
\textbf{1.} & \symking{}xg3 & \symknight{}c7\\
\textbf{2.} & \symking{}f4 & g4\\
\textbf{3.} & \symking{}e5 & g5\\
\textbf{4.} & \symking{}d6 & g6\\
\textbf{5.} & \symking{}c6 & g7\\
\textbf{6.} & \symking{}b7 & g8=\symqueen{}\\
\textbf{7.} & \symking{}a7 & \symqueen{}a8\#\\
\end{tabular}\par\vspace{3mm}
\qrcode[height=12mm]{https://helpman.komtera.lt/?fen=8/8/Nr6/8/8/6Qk/6P1/K7&moves=7}\\{\sffamily\tiny\color{ink}scannen zum Lösen}
\end{minipage}
\vspace{2mm}

{\raggedright\noindent\textcolor{ink}{\sffamily\textsc{Themen}}\enspace \themetag{Umwandlung} \themetag{Excelsior} \themetag{ein Stein zieht} \themetag{Phönix} \themetag{ohne Schach}\par\smallskip{\footnotesize\textit{Hinweis:} das Diagramm hat keinen legalen letzten Zug, kann also in keiner Partie entstehen.}\par}

\clearpage
\phantomsection\label{hm:KQNvkrb}
\begin{tcolorbox}[enhanced,colback=ink,colframe=ink,arc=0pt,boxrule=0pt,left=4mm,right=4mm,top=2.5mm,bottom=2.5mm]
\color{white}\sffamily{\LARGE\bfseries KQNvkrb}\hfill{\Large h\#7}\\[0.5mm]
{\small 6 Steine\quad Schwarz am Zug}\hfill{\small Nr.~546\qquad\href{https://helpman.komtera.lt/?fen=8/5b2/8/1r6/4Q3/5k2/8/K6N&moves=7}{\color{white}\bfseries lösen\,\ensuremath{\nearrow}}}
\end{tcolorbox}
\vspace{2mm}
\noindent\begin{minipage}[t]{36mm}\vspace{0pt}\sffamily\small\raggedright
\hmfigure{längste eindeutige}{{\Large h\#7}}\hmfigure{eindeutig bei dieser Länge}{{\Large 316} Stellungen}\hmfigure{längstes Matt}{{\Large h\#10}}\hmfigure{in der Klasse}{15,5 Milliarden Stellungen,\\ 7,99 Milliarden mit Hilfsmatt}
\end{minipage}\hfill
\begin{minipage}[t]{80mm}\vspace{0pt}\begin{tikzpicture}\begin{axis}[width=72mm,height=42mm,ybar,bar width=1.6pt,xmin=-1,xmax=21,ymin=0,ymax=2.941,axis y line*=left,axis x line*=bottom,xlabel={\scriptsize Halbzüge bis Matt},ylabel={\scriptsize Stellungen (Mrd.)},tick label style={font=\tiny},label style={font=\scriptsize},enlarge x limits=false,ylabel shift=-3pt,xlabel shift=-3pt]\addplot[fill=ink!55,draw=none] coordinates {(0,0.007673) (1,0.0646) (2,0.3136) (3,1.174) (4,2.673) (5,2.269) (6,1.174) (7,0.2446) (8,0.02316) (9,0.0002832) (10,0.02189) (11,1.358e-06) (12,0.02048) (13,4.39e-07) (14,0.006057) (15,4e-08) (16,0.0003112) (18,1.34e-06) (20,5e-09)};\draw[densely dotted] (axis cs:14,0) -- (axis cs:14,2.941);\end{axis}\begin{axis}[width=72mm,height=42mm,axis y line*=right,axis x line=none,ymode=log,xmin=-1,xmax=21,ylabel={\scriptsize eindeutig (log)},tick label style={font=\tiny},label style={font=\scriptsize},ylabel shift=-3pt]\addplot[mark=*,mark size=0.7pt,draw=black,thin] coordinates {(0,7672859) (1,56928584) (2,184859815) (3,223975435) (4,183208247) (5,30453316) (6,2263554) (7,25005) (8,225179) (9,137) (10,159616) (11,7) (12,27714) (14,316)};\end{axis}\end{tikzpicture}\end{minipage}
\vspace{3mm}

\noindent\begin{minipage}[t]{60mm}\vspace{0pt}\sbox\hmboard{\chessboard[setfen=8/5b2/8/1r6/4Q3/5k2/8/K6N b - - 0 1]}\usebox\hmboard\\[1mm]
\makebox[\wd\hmboard]{\sffamily\hspace{3mm}\textbf{h\#7}\hfill(3+3)}\end{minipage}\hfill
\begin{minipage}[t]{52mm}\vspace{0pt}
\textcolor{ink}{\sffamily\textsc{Lösung}}\par\smallskip\small
\begin{tabular}{@{}r@{\enspace}l@{\enspace}l@{}}
\textbf{1.} & \symking{}xe4 & \symknight{}f2+\\
\textbf{2.} & \symking{}e3 & \symknight{}d3\\
\textbf{3.} & \symbishop{}h5 & \symking{}a2\\
\textbf{4.} & \symrook{}b1 & \symknight{}b2\\
\textbf{5.} & \symking{}d2 & \symking{}b3\\
\textbf{6.} & \symking{}c1 & \symking{}c3\\
\textbf{7.} & \symbishop{}d1 & \symknight{}d3\#\\
\end{tabular}\par\vspace{3mm}
\qrcode[height=12mm]{https://helpman.komtera.lt/?fen=8/5b2/8/1r6/4Q3/5k2/8/K6N&moves=7}\\{\sffamily\tiny\color{ink}scannen zum Lösen}
\end{minipage}
\vspace{2mm}

{\raggedright\noindent\textcolor{ink}{\sffamily\textsc{Themen}}\enspace \themetag{Switchback} \themetag{Selbstblock}\par\smallskip{\footnotesize\textit{Hinweis:} die Partei am Zug steht im Diagramm im Schach; die Lösung beginnt mit einem Schlag.}\par}

\clearpage
\phantomsection\label{hm:KQPPPvk}
\begin{tcolorbox}[enhanced,colback=ink,colframe=ink,arc=0pt,boxrule=0pt,left=4mm,right=4mm,top=2.5mm,bottom=2.5mm]
\color{white}\sffamily{\LARGE\bfseries KQPPPvk}\hfill{\Large h\#7}\\[0.5mm]
{\small 6 Steine\quad Schwarz am Zug}\hfill{\small Nr.~547\qquad\href{https://helpman.komtera.lt/?fen=7k/7Q/8/8/6P1/P7/3P4/K7&moves=7}{\color{white}\bfseries lösen\,\ensuremath{\nearrow}}}
\end{tcolorbox}
\vspace{2mm}
\noindent\begin{minipage}[t]{36mm}\vspace{0pt}\sffamily\small\raggedright
\hmfigure{längste eindeutige}{{\Large h\#7}}\hmfigure{eindeutig bei dieser Länge}{{\Large 44} Stellungen}\hmfigure{längstes Matt}{{\Large h\#8}}\hmfigure{in der Klasse}{25,6 Milliarden Stellungen,\\ 2,71 Milliarden mit Hilfsmatt}
\end{minipage}\hfill
\begin{minipage}[t]{80mm}\vspace{0pt}\begin{tikzpicture}\begin{axis}[width=72mm,height=42mm,ybar,bar width=1.6pt,xmin=-1,xmax=17,ymin=0,ymax=976.4,axis y line*=left,axis x line*=bottom,xlabel={\scriptsize Halbzüge bis Matt},ylabel={\scriptsize Stellungen (Mio.)},tick label style={font=\tiny},label style={font=\scriptsize},enlarge x limits=false,ylabel shift=-3pt,xlabel shift=-3pt]\addplot[fill=ink!55,draw=none] coordinates {(0,5.257) (1,38.95) (2,112.1) (3,363.5) (4,887.7) (5,524.7) (6,573) (7,111.4) (8,78.61) (9,5.73) (10,6.57) (11,0.01613) (12,1.836) (13,4e-06) (14,0.3003) (16,0.001056)};\draw[densely dotted] (axis cs:14,0) -- (axis cs:14,976.4);\end{axis}\begin{axis}[width=72mm,height=42mm,axis y line*=right,axis x line=none,ymode=log,xmin=-1,xmax=17,ylabel={\scriptsize eindeutig (log)},tick label style={font=\tiny},label style={font=\scriptsize},ylabel shift=-3pt]\addplot[mark=*,mark size=0.7pt,draw=black,thin] coordinates {(0,5256739) (1,23772388) (2,66657837) (3,67067422) (4,83032144) (5,10873499) (6,6110345) (7,316109) (8,213615) (9,424) (10,23190) (12,4267) (14,44)};\end{axis}\end{tikzpicture}\end{minipage}
\vspace{3mm}

\noindent\begin{minipage}[t]{60mm}\vspace{0pt}\sbox\hmboard{\chessboard[setfen=7k/7Q/8/8/6P1/P7/3P4/K7 b - - 0 1]}\usebox\hmboard\\[1mm]
\makebox[\wd\hmboard]{\sffamily\hspace{3mm}\textbf{h\#7}\hfill(5+1)}\end{minipage}\hfill
\begin{minipage}[t]{52mm}\vspace{0pt}
\textcolor{ink}{\sffamily\textsc{Lösung}}\par\smallskip\small
\begin{tabular}{@{}r@{\enspace}l@{\enspace}l@{}}
\textbf{1.} & \symking{}xh7 & g5\\
\textbf{2.} & \symking{}g7 & g6\\
\textbf{3.} & \symking{}f6 & g7\\
\textbf{4.} & \symking{}e5 & g8=\symqueen{}\\
\textbf{5.} & \symking{}d4 & \symqueen{}f8\\
\textbf{6.} & \symking{}c4 & \symking{}b1\\
\textbf{7.} & \symking{}b3 & \symqueen{}b4\#\\
\end{tabular}\par\vspace{3mm}
\qrcode[height=12mm]{https://helpman.komtera.lt/?fen=7k/7Q/8/8/6P1/P7/3P4/K7&moves=7}\\{\sffamily\tiny\color{ink}scannen zum Lösen}
\end{minipage}
\vspace{2mm}

{\raggedright\noindent\textcolor{ink}{\sffamily\textsc{Themen}}\enspace \themetag{Umwandlung} \themetag{ein Stein zieht} \themetag{Phönix} \themetag{ohne Schach} \themetag{Umnov}\par\smallskip{\footnotesize\textit{Hinweis:} die Partei am Zug steht im Diagramm im Schach; die Lösung beginnt mit einem Schlag.}\par}

\clearpage
\phantomsection\label{hm:KQPPvkb}
\begin{tcolorbox}[enhanced,colback=ink,colframe=ink,arc=0pt,boxrule=0pt,left=4mm,right=4mm,top=2.5mm,bottom=2.5mm]
\color{white}\sffamily{\LARGE\bfseries KQPPvkb}\hfill{\Large h\#7}\\[0.5mm]
{\small 6 Steine\quad Schwarz am Zug}\hfill{\small Nr.~548\qquad\href{https://helpman.komtera.lt/?fen=8/8/8/6Qk/8/1P6/1P1b4/K7&moves=7}{\color{white}\bfseries lösen\,\ensuremath{\nearrow}}}
\end{tcolorbox}
\vspace{2mm}
\noindent\begin{minipage}[t]{36mm}\vspace{0pt}\sffamily\small\raggedright
\hmfigure{längste eindeutige}{{\Large h\#7}}\hmfigure{eindeutig bei dieser Länge}{{\Large 768} Stellungen}\hmfigure{längstes Matt}{{\Large h\#9}}\hmfigure{in der Klasse}{34,1 Milliarden Stellungen,\\ 10,3 Milliarden mit Hilfsmatt}
\end{minipage}\hfill
\begin{minipage}[t]{80mm}\vspace{0pt}\begin{tikzpicture}\begin{axis}[width=72mm,height=42mm,ybar,bar width=1.6pt,xmin=-1,xmax=19,ymin=0,ymax=3.225,axis y line*=left,axis x line*=bottom,xlabel={\scriptsize Halbzüge bis Matt},ylabel={\scriptsize Stellungen (Mrd.)},tick label style={font=\tiny},label style={font=\scriptsize},enlarge x limits=false,ylabel shift=-3pt,xlabel shift=-3pt]\addplot[fill=ink!55,draw=none] coordinates {(0,0.01192) (1,0.09317) (2,0.3403) (3,1.236) (4,2.932) (5,2.249) (6,2.191) (7,0.6973) (8,0.3967) (9,0.05822) (10,0.03309) (11,0.0003482) (12,0.01283) (13,1.31e-07) (14,0.002964) (16,2.578e-05) (18,3e-09)};\draw[densely dotted] (axis cs:14,0) -- (axis cs:14,3.225);\end{axis}\begin{axis}[width=72mm,height=42mm,axis y line*=right,axis x line=none,ymode=log,xmin=-1,xmax=19,ylabel={\scriptsize eindeutig (log)},tick label style={font=\tiny},label style={font=\scriptsize},ylabel shift=-3pt]\addplot[mark=*,mark size=0.7pt,draw=black,thin] coordinates {(0,11916468) (1,65155005) (2,165459153) (3,188432862) (4,217034007) (5,49923194) (6,22819921) (7,3054905) (8,962613) (9,11754) (10,155233) (12,38343) (14,768)};\end{axis}\end{tikzpicture}\end{minipage}
\vspace{3mm}

\noindent\begin{minipage}[t]{60mm}\vspace{0pt}\sbox\hmboard{\chessboard[setfen=8/8/8/6Qk/8/1P6/1P1b4/K7 b - - 0 1]}\usebox\hmboard\\[1mm]
\makebox[\wd\hmboard]{\sffamily\hspace{3mm}\textbf{h\#7}\hfill(4+2)}\end{minipage}\hfill
\begin{minipage}[t]{52mm}\vspace{0pt}
\textcolor{ink}{\sffamily\textsc{Lösung}}\par\smallskip\small
\begin{tabular}{@{}r@{\enspace}l@{\enspace}l@{}}
\textbf{1.} & \symking{}xg5 & b4\\
\textbf{2.} & \symking{}f4 & b5\\
\textbf{3.} & \symking{}e3 & b4\\
\textbf{4.} & \symbishop{}xb4 & b6\\
\textbf{5.} & \symking{}d2 & b7\\
\textbf{6.} & \symking{}c1 & b8=\symqueen{}\\
\textbf{7.} & \symbishop{}d2 & \symqueen{}b1\#\\
\end{tabular}\par\vspace{3mm}
\qrcode[height=12mm]{https://helpman.komtera.lt/?fen=8/8/8/6Qk/8/1P6/1P1b4/K7&moves=7}\\{\sffamily\tiny\color{ink}scannen zum Lösen}
\end{minipage}
\vspace{2mm}

{\raggedright\noindent\textcolor{ink}{\sffamily\textsc{Themen}}\enspace \themetag{Umwandlung} \themetag{Switchback} \themetag{Selbstblock} \themetag{Phönix} \themetag{ohne Schach}\par\smallskip{\footnotesize\textit{Hinweis:} die Partei am Zug steht im Diagramm im Schach; die Lösung beginnt mit einem Schlag.}\par}

\clearpage
\phantomsection\label{hm:KQPPvkn}
\begin{tcolorbox}[enhanced,colback=ink,colframe=ink,arc=0pt,boxrule=0pt,left=4mm,right=4mm,top=2.5mm,bottom=2.5mm]
\color{white}\sffamily{\LARGE\bfseries KQPPvkn}\hfill{\Large h\#7}\\[0.5mm]
{\small 6 Steine\quad Schwarz am Zug}\hfill{\small Nr.~549\qquad\href{https://helpman.komtera.lt/?fen=8/7n/8/8/8/6P1/4P3/K1kQ4&moves=7}{\color{white}\bfseries lösen\,\ensuremath{\nearrow}}}
\end{tcolorbox}
\vspace{2mm}
\noindent\begin{minipage}[t]{36mm}\vspace{0pt}\sffamily\small\raggedright
\hmfigure{längste eindeutige}{{\Large h\#7}}\hmfigure{eindeutig bei dieser Länge}{{\Large 1\,832} Stellungen}\hmfigure{längstes Matt}{{\Large h\#9}}\hmfigure{in der Klasse}{34,1 Milliarden Stellungen,\\ 10,6 Milliarden mit Hilfsmatt}
\end{minipage}\hfill
\begin{minipage}[t]{80mm}\vspace{0pt}\begin{tikzpicture}\begin{axis}[width=72mm,height=42mm,ybar,bar width=1.6pt,xmin=-1,xmax=19,ymin=0,ymax=3.087,axis y line*=left,axis x line*=bottom,xlabel={\scriptsize Halbzüge bis Matt},ylabel={\scriptsize Stellungen (Mrd.)},tick label style={font=\tiny},label style={font=\scriptsize},enlarge x limits=false,ylabel shift=-3pt,xlabel shift=-3pt]\addplot[fill=ink!55,draw=none] coordinates {(0,0.01329) (1,0.1012) (2,0.3383) (3,1.202) (4,2.807) (5,2.187) (6,2.473) (7,0.7727) (8,0.5246) (9,0.07054) (10,0.04188) (11,0.0005446) (12,0.01423) (13,9.32e-07) (14,0.004578) (15,5.32e-07) (16,0.0001065) (17,3e-08) (18,2e-08)};\draw[densely dotted] (axis cs:14,0) -- (axis cs:14,3.087);\end{axis}\begin{axis}[width=72mm,height=42mm,axis y line*=right,axis x line=none,ymode=log,xmin=-1,xmax=19,ylabel={\scriptsize eindeutig (log)},tick label style={font=\tiny},label style={font=\scriptsize},ylabel shift=-3pt]\addplot[mark=*,mark size=0.7pt,draw=black,thin] coordinates {(0,13291507) (1,65938582) (2,154340342) (3,180302082) (4,204872820) (5,54325848) (6,29926664) (7,3518009) (8,1217686) (9,13352) (10,145386) (11,4) (12,53454) (14,1832)};\end{axis}\end{tikzpicture}\end{minipage}
\vspace{3mm}

\noindent\begin{minipage}[t]{60mm}\vspace{0pt}\sbox\hmboard{\chessboard[setfen=8/7n/8/8/8/6P1/4P3/K1kQ4 b - - 0 1]}\usebox\hmboard\\[1mm]
\makebox[\wd\hmboard]{\sffamily\hspace{3mm}\textbf{h\#7}\hfill(4+2)}\end{minipage}\hfill
\begin{minipage}[t]{52mm}\vspace{0pt}
\textcolor{ink}{\sffamily\textsc{Lösung}}\par\smallskip\small
\begin{tabular}{@{}r@{\enspace}l@{\enspace}l@{}}
\textbf{1.} & \symking{}xd1 & g4\\
\textbf{2.} & \symking{}e1 & g5\\
\textbf{3.} & \symking{}f2 & g6\\
\textbf{4.} & \symknight{}g5 & g7\\
\textbf{5.} & \symknight{}f3 & exf3\\
\textbf{6.} & \symking{}g3 & g8=\symqueen{}+\\
\textbf{7.} & \symking{}h4 & \symqueen{}g4\#\\
\end{tabular}\par\vspace{3mm}
\qrcode[height=12mm]{https://helpman.komtera.lt/?fen=8/7n/8/8/8/6P1/4P3/K1kQ4&moves=7}\\{\sffamily\tiny\color{ink}scannen zum Lösen}
\end{minipage}
\vspace{2mm}

{\raggedright\noindent\textcolor{ink}{\sffamily\textsc{Themen}}\enspace \themetag{reines Matt} \themetag{Mustermatt} \themetag{Umwandlung} \themetag{Rundlauf} \themetag{Phönix} \themetag{Umnov}\par\smallskip{\footnotesize\textit{Hinweis:} die Partei am Zug steht im Diagramm im Schach; die Lösung beginnt mit einem Schlag.}\par}

\clearpage
\phantomsection\label{hm:KQPPvkq}
\begin{tcolorbox}[enhanced,colback=ink,colframe=ink,arc=0pt,boxrule=0pt,left=4mm,right=4mm,top=2.5mm,bottom=2.5mm]
\color{white}\sffamily{\LARGE\bfseries KQPPvkq}\hfill{\Large h\#7}\\[0.5mm]
{\small 6 Steine\quad Schwarz am Zug}\hfill{\small Nr.~550\qquad\href{https://helpman.komtera.lt/?fen=8/8/8/6Qk/8/1P2q3/1P6/K7&moves=7}{\color{white}\bfseries lösen\,\ensuremath{\nearrow}}}
\end{tcolorbox}
\vspace{2mm}
\noindent\begin{minipage}[t]{36mm}\vspace{0pt}\sffamily\small\raggedright
\hmfigure{längste eindeutige}{{\Large h\#7}}\hmfigure{eindeutig bei dieser Länge}{{\Large 422} Stellungen}\hmfigure{längstes Matt}{{\Large h\#9}}\hmfigure{in der Klasse}{34,1 Milliarden Stellungen,\\ 8,93 Milliarden mit Hilfsmatt}
\end{minipage}\hfill
\begin{minipage}[t]{80mm}\vspace{0pt}\begin{tikzpicture}\begin{axis}[width=72mm,height=42mm,ybar,bar width=1.6pt,xmin=-1,xmax=19,ymin=0,ymax=2.591,axis y line*=left,axis x line*=bottom,xlabel={\scriptsize Halbzüge bis Matt},ylabel={\scriptsize Stellungen (Mrd.)},tick label style={font=\tiny},label style={font=\scriptsize},enlarge x limits=false,ylabel shift=-3pt,xlabel shift=-3pt]\addplot[fill=ink!55,draw=none] coordinates {(0,0.006799) (1,0.05703) (2,0.2502) (3,1.04) (4,2.319) (5,2.356) (6,1.69) (7,0.8003) (8,0.2918) (9,0.0759) (10,0.02559) (11,0.000781) (12,0.01007) (13,1.512e-06) (14,0.002087) (15,4.6e-08) (16,2.138e-05) (18,5e-09)};\draw[densely dotted] (axis cs:14,0) -- (axis cs:14,2.591);\end{axis}\begin{axis}[width=72mm,height=42mm,axis y line*=right,axis x line=none,ymode=log,xmin=-1,xmax=19,ylabel={\scriptsize eindeutig (log)},tick label style={font=\tiny},label style={font=\scriptsize},ylabel shift=-3pt]\addplot[mark=*,mark size=0.7pt,draw=black,thin] coordinates {(0,6799248) (1,43993130) (2,125187694) (3,178091119) (4,167779087) (5,59141459) (6,15126608) (7,3656851) (8,674730) (9,33699) (10,85816) (11,237) (12,22432) (14,422)};\end{axis}\end{tikzpicture}\end{minipage}
\vspace{3mm}

\noindent\begin{minipage}[t]{60mm}\vspace{0pt}\sbox\hmboard{\chessboard[setfen=8/8/8/6Qk/8/1P2q3/1P6/K7 b - - 0 1]}\usebox\hmboard\\[1mm]
\makebox[\wd\hmboard]{\sffamily\hspace{3mm}\textbf{h\#7}\hfill(4+2)}\end{minipage}\hfill
\begin{minipage}[t]{52mm}\vspace{0pt}
\textcolor{ink}{\sffamily\textsc{Lösung}}\par\smallskip\small
\begin{tabular}{@{}r@{\enspace}l@{\enspace}l@{}}
\textbf{1.} & \symking{}xg5 & b4\\
\textbf{2.} & \symking{}f4 & b3\\
\textbf{3.} & \symqueen{}xb3 & b5\\
\textbf{4.} & \symking{}e3 & b6\\
\textbf{5.} & \symking{}d2 & b7\\
\textbf{6.} & \symking{}c1 & b8=\symqueen{}\\
\textbf{7.} & \symqueen{}d1 & \symqueen{}b2\#\\
\end{tabular}\par\vspace{3mm}
\qrcode[height=12mm]{https://helpman.komtera.lt/?fen=8/8/8/6Qk/8/1P2q3/1P6/K7&moves=7}\\{\sffamily\tiny\color{ink}scannen zum Lösen}
\end{minipage}
\vspace{2mm}

{\raggedright\noindent\textcolor{ink}{\sffamily\textsc{Themen}}\enspace \themetag{Umwandlung} \themetag{Selbstblock} \themetag{Phönix} \themetag{ohne Schach}\par\smallskip{\footnotesize\textit{Hinweis:} die Partei am Zug steht im Diagramm im Schach; die Lösung beginnt mit einem Schlag.}\par}

\clearpage
\phantomsection\label{hm:KQPPvkr}
\begin{tcolorbox}[enhanced,colback=ink,colframe=ink,arc=0pt,boxrule=0pt,left=4mm,right=4mm,top=2.5mm,bottom=2.5mm]
\color{white}\sffamily{\LARGE\bfseries KQPPvkr}\hfill{\Large h\#7}\\[0.5mm]
{\small 6 Steine\quad Schwarz am Zug}\hfill{\small Nr.~551\qquad\href{https://helpman.komtera.lt/?fen=8/8/6r1/6Qk/6P1/6P1/8/K7&moves=7}{\color{white}\bfseries lösen\,\ensuremath{\nearrow}}}
\end{tcolorbox}
\vspace{2mm}
\noindent\begin{minipage}[t]{36mm}\vspace{0pt}\sffamily\small\raggedright
\hmfigure{längste eindeutige}{{\Large h\#7}}\hmfigure{eindeutig bei dieser Länge}{{\Large 69} Stellungen}\hmfigure{längstes Matt}{{\Large h\#8}}\hmfigure{in der Klasse}{34,1 Milliarden Stellungen,\\ 9,78 Milliarden mit Hilfsmatt}
\end{minipage}\hfill
\begin{minipage}[t]{80mm}\vspace{0pt}\begin{tikzpicture}\begin{axis}[width=72mm,height=42mm,ybar,bar width=1.6pt,xmin=-1,xmax=17,ymin=0,ymax=3.326,axis y line*=left,axis x line*=bottom,xlabel={\scriptsize Halbzüge bis Matt},ylabel={\scriptsize Stellungen (Mrd.)},tick label style={font=\tiny},label style={font=\scriptsize},enlarge x limits=false,ylabel shift=-3pt,xlabel shift=-3pt]\addplot[fill=ink!55,draw=none] coordinates {(0,0.01016) (1,0.08311) (2,0.3448) (3,1.254) (4,3.024) (5,2.387) (6,1.851) (7,0.5837) (8,0.1864) (9,0.02472) (10,0.01778) (11,3.016e-05) (12,0.009021) (13,1.37e-07) (14,0.001428) (15,2.7e-08) (16,2.535e-06)};\draw[densely dotted] (axis cs:14,0) -- (axis cs:14,3.326);\end{axis}\begin{axis}[width=72mm,height=42mm,axis y line*=right,axis x line=none,ymode=log,xmin=-1,xmax=17,ylabel={\scriptsize eindeutig (log)},tick label style={font=\tiny},label style={font=\scriptsize},ylabel shift=-3pt]\addplot[mark=*,mark size=0.7pt,draw=black,thin] coordinates {(0,10155116) (1,60460892) (2,178178839) (3,189486358) (4,199602017) (5,45279066) (6,11783037) (7,1297156) (8,519740) (9,3686) (10,87609) (11,6) (12,16245) (14,69)};\end{axis}\end{tikzpicture}\end{minipage}
\vspace{3mm}

\noindent\begin{minipage}[t]{60mm}\vspace{0pt}\sbox\hmboard{\chessboard[setfen=8/8/6r1/6Qk/6P1/6P1/8/K7 b - - 0 1]}\usebox\hmboard\\[1mm]
\makebox[\wd\hmboard]{\sffamily\hspace{3mm}\textbf{h\#7}\hfill(4+2)}\end{minipage}\hfill
\begin{minipage}[t]{52mm}\vspace{0pt}
\textcolor{ink}{\sffamily\textsc{Lösung}}\par\smallskip\small
\begin{tabular}{@{}r@{\enspace}l@{\enspace}l@{}}
\textbf{1.} & \symking{}xg5 & \symking{}b2\\
\textbf{2.} & \symrook{}f6 & \symking{}c3\\
\textbf{3.} & \symrook{}f5 & gxf5\\
\textbf{4.} & \symking{}g4 & f6\\
\textbf{5.} & \symking{}f3 & f7\\
\textbf{6.} & \symking{}e2 & f8=\symqueen{}\\
\textbf{7.} & \symking{}d1 & \symqueen{}f1\#\\
\end{tabular}\par\vspace{3mm}
\qrcode[height=12mm]{https://helpman.komtera.lt/?fen=8/8/6r1/6Qk/6P1/6P1/8/K7&moves=7}\\{\sffamily\tiny\color{ink}scannen zum Lösen}
\end{minipage}
\vspace{2mm}

{\raggedright\noindent\textcolor{ink}{\sffamily\textsc{Themen}}\enspace \themetag{reines Matt} \themetag{Mustermatt} \themetag{Spiegelmatt} \themetag{Umwandlung} \themetag{Phönix} \themetag{ohne Schach} \themetag{Umnov}\par\smallskip{\footnotesize\textit{Hinweis:} das Diagramm hat keinen legalen letzten Zug, kann also in keiner Partie entstehen.}\par}

\clearpage
\phantomsection\label{hm:KQPvkrr}
\begin{tcolorbox}[enhanced,colback=ink,colframe=ink,arc=0pt,boxrule=0pt,left=4mm,right=4mm,top=2.5mm,bottom=2.5mm]
\color{white}\sffamily{\LARGE\bfseries KQPvkrr}\hfill{\Large h\#7}\\[0.5mm]
{\small 6 Steine\quad Schwarz am Zug}\hfill{\small Nr.~552\qquad\href{https://helpman.komtera.lt/?fen=8/8/8/8/1r6/5r2/5P2/K1kQ4&moves=7}{\color{white}\bfseries lösen\,\ensuremath{\nearrow}}}
\end{tcolorbox}
\vspace{2mm}
\noindent\begin{minipage}[t]{36mm}\vspace{0pt}\sffamily\small\raggedright
\hmfigure{längste eindeutige}{{\Large h\#7}}\hmfigure{eindeutig bei dieser Länge}{{\Large 111} Stellungen}\hmfigure{längstes Matt}{{\Large h\#9}}\hmfigure{in der Klasse}{45,4 Milliarden Stellungen,\\ 11,7 Milliarden mit Hilfsmatt}
\end{minipage}\hfill
\begin{minipage}[t]{80mm}\vspace{0pt}\begin{tikzpicture}\begin{axis}[width=72mm,height=42mm,ybar,bar width=1.6pt,xmin=-1,xmax=19,ymin=0,ymax=3.729,axis y line*=left,axis x line*=bottom,xlabel={\scriptsize Halbzüge bis Matt},ylabel={\scriptsize Stellungen (Mrd.)},tick label style={font=\tiny},label style={font=\scriptsize},enlarge x limits=false,ylabel shift=-3pt,xlabel shift=-3pt]\addplot[fill=ink!55,draw=none] coordinates {(0,0.006155) (1,0.05517) (2,0.2748) (3,1.179) (4,2.901) (5,3.39) (6,2.292) (7,1.259) (8,0.2308) (9,0.05426) (10,0.02245) (11,9.371e-05) (12,0.01989) (13,9.7e-06) (14,0.002512) (15,6.913e-06) (16,2.193e-05) (17,6e-08) (18,6.6e-08)};\draw[densely dotted] (axis cs:14,0) -- (axis cs:14,3.729);\end{axis}\begin{axis}[width=72mm,height=42mm,axis y line*=right,axis x line=none,ymode=log,xmin=-1,xmax=19,ylabel={\scriptsize eindeutig (log)},tick label style={font=\tiny},label style={font=\scriptsize},ylabel shift=-3pt]\addplot[mark=*,mark size=0.7pt,draw=black,thin] coordinates {(0,6155089) (1,45135705) (2,152173001) (3,202441299) (4,184641029) (5,73431873) (6,12365258) (7,2484947) (8,399986) (9,7973) (10,107817) (11,15) (12,13796) (13,3) (14,111)};\end{axis}\end{tikzpicture}\end{minipage}
\vspace{3mm}

\noindent\begin{minipage}[t]{60mm}\vspace{0pt}\sbox\hmboard{\chessboard[setfen=8/8/8/8/1r6/5r2/5P2/K1kQ4 b - - 0 1]}\usebox\hmboard\\[1mm]
\makebox[\wd\hmboard]{\sffamily\hspace{3mm}\textbf{h\#7}\hfill(3+3)}\end{minipage}\hfill
\begin{minipage}[t]{52mm}\vspace{0pt}
\textcolor{ink}{\sffamily\textsc{Lösung}}\par\smallskip\small
\begin{tabular}{@{}r@{\enspace}l@{\enspace}l@{}}
\textbf{1.} & \symking{}xd1 & \symking{}a2\\
\textbf{2.} & \symrook{}g3 & f4\\
\textbf{3.} & \symking{}e2 & f5\\
\textbf{4.} & \symking{}f3 & f6\\
\textbf{5.} & \symking{}g4 & f7\\
\textbf{6.} & \symking{}h4 & f8=\symqueen{}\\
\textbf{7.} & \symrook{}bg4 & \symqueen{}h6\#\\
\end{tabular}\par\vspace{3mm}
\qrcode[height=12mm]{https://helpman.komtera.lt/?fen=8/8/8/8/1r6/5r2/5P2/K1kQ4&moves=7}\\{\sffamily\tiny\color{ink}scannen zum Lösen}
\end{minipage}
\vspace{2mm}

{\raggedright\noindent\textcolor{ink}{\sffamily\textsc{Themen}}\enspace \themetag{reines Matt} \themetag{Mustermatt} \themetag{Umwandlung} \themetag{Excelsior} \themetag{Selbstblock} \themetag{Phönix} \themetag{ohne Schach}\par\smallskip{\footnotesize\textit{Hinweis:} die Partei am Zug steht im Diagramm im Schach; die Lösung beginnt mit einem Schlag.}\par}

\clearpage
\phantomsection\label{hm:KRBBBvk}
\begin{tcolorbox}[enhanced,colback=ink,colframe=ink,arc=0pt,boxrule=0pt,left=4mm,right=4mm,top=2.5mm,bottom=2.5mm]
\color{white}\sffamily{\LARGE\bfseries KRBBBvk}\hfill{\Large h\#7}\\[0.5mm]
{\small 6 Steine\quad Schwarz am Zug}\hfill{\small Nr.~553\qquad\href{https://helpman.komtera.lt/?fen=5B1R/5B1k/8/8/8/8/7B/K7&moves=7}{\color{white}\bfseries lösen\,\ensuremath{\nearrow}}}
\end{tcolorbox}
\vspace{2mm}
\noindent\begin{minipage}[t]{36mm}\vspace{0pt}\sffamily\small\raggedright
\hmfigure{längste eindeutige}{{\Large h\#7}}\hmfigure{eindeutig bei dieser Länge}{{\Large 3} Stellungen}\hmfigure{längstes Matt}{{\Large h\#8}}\hmfigure{in der Klasse}{15,5 Milliarden Stellungen,\\ 1,53 Milliarden mit Hilfsmatt}
\end{minipage}\hfill
\begin{minipage}[t]{80mm}\vspace{0pt}\begin{tikzpicture}\begin{axis}[width=72mm,height=42mm,ybar,bar width=1.6pt,xmin=-1,xmax=17,ymin=0,ymax=624.6,axis y line*=left,axis x line*=bottom,xlabel={\scriptsize Halbzüge bis Matt},ylabel={\scriptsize Stellungen (Mio.)},tick label style={font=\tiny},label style={font=\scriptsize},enlarge x limits=false,ylabel shift=-3pt,xlabel shift=-3pt]\addplot[fill=ink!55,draw=none] coordinates {(0,7.221) (1,21.31) (2,70.45) (3,183.4) (4,567.8) (5,269.8) (6,294.1) (7,49.91) (8,53.43) (9,4.132) (10,3.414) (11,0.03471) (12,0.1551) (13,0.000239) (14,0.01365) (16,1.4e-05)};\draw[densely dotted] (axis cs:14,0) -- (axis cs:14,624.6);\end{axis}\begin{axis}[width=72mm,height=42mm,axis y line*=right,axis x line=none,ymode=log,xmin=-1,xmax=17,ylabel={\scriptsize eindeutig (log)},tick label style={font=\tiny},label style={font=\scriptsize},ylabel shift=-3pt]\addplot[mark=*,mark size=0.7pt,draw=black,thin] coordinates {(0,7221464) (1,17734601) (2,56018552) (3,71438562) (4,103436979) (5,10605658) (6,9630364) (7,440656) (8,190418) (9,1426) (10,2562) (11,2) (12,99) (14,3)};\end{axis}\end{tikzpicture}\end{minipage}
\vspace{3mm}

\noindent\begin{minipage}[t]{60mm}\vspace{0pt}\sbox\hmboard{\chessboard[setfen=5B1R/5B1k/8/8/8/8/7B/K7 b - - 0 1]}\usebox\hmboard\\[1mm]
\makebox[\wd\hmboard]{\sffamily\hspace{3mm}\textbf{h\#7}\hfill(5+1)}\end{minipage}\hfill
\begin{minipage}[t]{52mm}\vspace{0pt}
\textcolor{ink}{\sffamily\textsc{Lösung}}\par\smallskip\small
\begin{tabular}{@{}r@{\enspace}l@{\enspace}l@{}}
\textbf{1.} & \symking{}xh8 & \symking{}b2\\
\textbf{2.} & \symking{}h7 & \symking{}c3\\
\textbf{3.} & \symking{}h8 & \symking{}d4\\
\textbf{4.} & \symking{}h7 & \symking{}e5\\
\textbf{5.} & \symking{}h8 & \symking{}f6\\
\textbf{6.} & \symking{}h7 & \symbishop{}e5\\
\textbf{7.} & \symking{}h8 & \symking{}g6\#\\
\end{tabular}\par\vspace{3mm}
\qrcode[height=12mm]{https://helpman.komtera.lt/?fen=5B1R/5B1k/8/8/8/8/7B/K7&moves=7}\\{\sffamily\tiny\color{ink}scannen zum Lösen}
\end{minipage}
\vspace{2mm}

{\raggedright\noindent\textcolor{ink}{\sffamily\textsc{Themen}}\enspace \themetag{Spiegelmatt} \themetag{Switchback} \themetag{ein Stein zieht} \themetag{Kniest} \themetag{Pendel} \themetag{ohne Schach}\par\smallskip{\footnotesize\textit{Hinweis:} die Partei am Zug steht im Diagramm im Schach; die Lösung beginnt mit einem Schlag.}\par}

\clearpage
\phantomsection\label{hm:KRBBPvk}
\begin{tcolorbox}[enhanced,colback=ink,colframe=ink,arc=0pt,boxrule=0pt,left=4mm,right=4mm,top=2.5mm,bottom=2.5mm]
\color{white}\sffamily{\LARGE\bfseries KRBBPvk}\hfill{\Large h\#7}\\[0.5mm]
{\small 6 Steine\quad Schwarz am Zug}\hfill{\small Nr.~554\qquad\href{https://helpman.komtera.lt/?fen=1B6/8/8/8/7B/8/6RP/K4k2&moves=7}{\color{white}\bfseries lösen\,\ensuremath{\nearrow}}}
\end{tcolorbox}
\vspace{2mm}
\noindent\begin{minipage}[t]{36mm}\vspace{0pt}\sffamily\small\raggedright
\hmfigure{längste eindeutige}{{\Large h\#7}}\hmfigure{eindeutig bei dieser Länge}{{\Large 500} Stellungen}\hmfigure{längstes Matt}{{\Large h\#9}}\hmfigure{in der Klasse}{45,4 Milliarden Stellungen,\\ 14,4 Milliarden mit Hilfsmatt}
\end{minipage}\hfill
\begin{minipage}[t]{80mm}\vspace{0pt}\begin{tikzpicture}\begin{axis}[width=72mm,height=42mm,ybar,bar width=1.6pt,xmin=-1,xmax=19,ymin=0,ymax=4.359,axis y line*=left,axis x line*=bottom,xlabel={\scriptsize Halbzüge bis Matt},ylabel={\scriptsize Stellungen (Mrd.)},tick label style={font=\tiny},label style={font=\scriptsize},enlarge x limits=false,ylabel shift=-3pt,xlabel shift=-3pt]\addplot[fill=ink!55,draw=none] coordinates {(0,0.03009) (1,0.1241) (2,0.3532) (3,1.253) (4,3.781) (5,3.079) (6,3.963) (7,0.8388) (8,0.8187) (9,0.08693) (10,0.06195) (11,0.002121) (12,0.007134) (13,7.604e-06) (14,0.00258) (15,8.3e-08) (16,7.745e-05) (18,2.39e-07)};\draw[densely dotted] (axis cs:14,0) -- (axis cs:14,4.359);\end{axis}\begin{axis}[width=72mm,height=42mm,axis y line*=right,axis x line=none,ymode=log,xmin=-1,xmax=19,ylabel={\scriptsize eindeutig (log)},tick label style={font=\tiny},label style={font=\scriptsize},ylabel shift=-3pt]\addplot[mark=*,mark size=0.7pt,draw=black,thin] coordinates {(0,30090336) (1,100792721) (2,272317386) (3,579021123) (4,994411971) (5,190158734) (6,141765424) (7,5930655) (8,3040793) (9,82952) (10,171451) (11,767) (12,50148) (13,2) (14,500)};\end{axis}\end{tikzpicture}\end{minipage}
\vspace{3mm}

\noindent\begin{minipage}[t]{60mm}\vspace{0pt}\sbox\hmboard{\chessboard[setfen=1B6/8/8/8/7B/8/6RP/K4k2 b - - 0 1]}\usebox\hmboard\\[1mm]
\makebox[\wd\hmboard]{\sffamily\hspace{3mm}\textbf{h\#7}\hfill(5+1)}\end{minipage}\hfill
\begin{minipage}[t]{52mm}\vspace{0pt}
\textcolor{ink}{\sffamily\textsc{Lösung}}\par\smallskip\small
\begin{tabular}{@{}r@{\enspace}l@{\enspace}l@{}}
\textbf{1.} & \symking{}xg2 & \symbishop{}hg3\\
\textbf{2.} & \symking{}f3 & h4\\
\textbf{3.} & \symking{}e4 & h5\\
\textbf{4.} & \symking{}d5 & h6\\
\textbf{5.} & \symking{}c6 & h7\\
\textbf{6.} & \symking{}b7 & h8=\symqueen{}\\
\textbf{7.} & \symking{}a8 & \symqueen{}h1\#\\
\end{tabular}\par\vspace{3mm}
\qrcode[height=12mm]{https://helpman.komtera.lt/?fen=1B6/8/8/8/7B/8/6RP/K4k2&moves=7}\\{\sffamily\tiny\color{ink}scannen zum Lösen}
\end{minipage}
\vspace{2mm}

{\raggedright\noindent\textcolor{ink}{\sffamily\textsc{Themen}}\enspace \themetag{reines Matt} \themetag{Mustermatt} \themetag{Umwandlung} \themetag{Excelsior} \themetag{ein Stein zieht} \themetag{ohne Schach}\par\smallskip{\footnotesize\textit{Hinweis:} die Lösung beginnt mit einem Schlag.}\par}

\clearpage
\phantomsection\label{hm:KRBBvkb}
\begin{tcolorbox}[enhanced,colback=ink,colframe=ink,arc=0pt,boxrule=0pt,left=4mm,right=4mm,top=2.5mm,bottom=2.5mm]
\color{white}\sffamily{\LARGE\bfseries KRBBvkb}\hfill{\Large h\#7}\\[0.5mm]
{\small 6 Steine\quad Schwarz am Zug}\hfill{\small Nr.~555\qquad\href{https://helpman.komtera.lt/?fen=6b1/8/8/8/8/4B3/5B1R/K6k&moves=7}{\color{white}\bfseries lösen\,\ensuremath{\nearrow}}}
\end{tcolorbox}
\vspace{2mm}
\noindent\begin{minipage}[t]{36mm}\vspace{0pt}\sffamily\small\raggedright
\hmfigure{längste eindeutige}{{\Large h\#7}}\hmfigure{eindeutig bei dieser Länge}{{\Large 5} Stellungen}\hmfigure{längstes Matt}{{\Large h\#10}}\hmfigure{in der Klasse}{15,5 Milliarden Stellungen,\\ 4,48 Milliarden mit Hilfsmatt}
\end{minipage}\hfill
\begin{minipage}[t]{80mm}\vspace{0pt}\begin{tikzpicture}\begin{axis}[width=72mm,height=42mm,ybar,bar width=1.6pt,xmin=-1,xmax=21,ymin=0,ymax=1.384,axis y line*=left,axis x line*=bottom,xlabel={\scriptsize Halbzüge bis Matt},ylabel={\scriptsize Stellungen (Mrd.)},tick label style={font=\tiny},label style={font=\scriptsize},enlarge x limits=false,ylabel shift=-3pt,xlabel shift=-3pt]\addplot[fill=ink!55,draw=none] coordinates {(0,0.006266) (1,0.02315) (2,0.09821) (3,0.4157) (4,1.258) (5,1.075) (6,1.065) (7,0.2855) (8,0.1839) (9,0.03982) (10,0.02181) (11,0.002327) (12,0.001141) (13,2.059e-05) (14,0.0001066) (15,1e-09) (16,0.0001649) (17,4.5e-08) (18,0.0001522) (20,4.024e-05)};\draw[densely dotted] (axis cs:14,0) -- (axis cs:14,1.384);\end{axis}\begin{axis}[width=72mm,height=42mm,axis y line*=right,axis x line=none,ymode=log,xmin=-1,xmax=21,ylabel={\scriptsize eindeutig (log)},tick label style={font=\tiny},label style={font=\scriptsize},ylabel shift=-3pt]\addplot[mark=*,mark size=0.7pt,draw=black,thin] coordinates {(0,6265936) (1,21350883) (2,62917255) (3,147076789) (4,212800821) (5,39328365) (6,17844290) (7,1369927) (8,525592) (9,34861) (10,13681) (11,247) (12,148) (14,5)};\end{axis}\end{tikzpicture}\end{minipage}
\vspace{3mm}

\noindent\begin{minipage}[t]{60mm}\vspace{0pt}\sbox\hmboard{\chessboard[setfen=6b1/8/8/8/8/4B3/5B1R/K6k b - - 0 1]}\usebox\hmboard\\[1mm]
\makebox[\wd\hmboard]{\sffamily\hspace{3mm}\textbf{h\#7}\hfill(4+2)}\end{minipage}\hfill
\begin{minipage}[t]{52mm}\vspace{0pt}
\textcolor{ink}{\sffamily\textsc{Lösung}}\par\smallskip\small
\begin{tabular}{@{}r@{\enspace}l@{\enspace}l@{}}
\textbf{1.} & \symking{}xh2 & \symking{}b2\\
\textbf{2.} & \symking{}h3 & \symking{}c3\\
\textbf{3.} & \symking{}g4 & \symking{}d4\\
\textbf{4.} & \symking{}h5 & \symking{}e5\\
\textbf{5.} & \symking{}g6 & \symbishop{}d4\\
\textbf{6.} & \symking{}h7 & \symking{}f6\\
\textbf{7.} & \symking{}h8 & \symking{}g6\#\\
\end{tabular}\par\vspace{3mm}
\qrcode[height=12mm]{https://helpman.komtera.lt/?fen=6b1/8/8/8/8/4B3/5B1R/K6k&moves=7}\\{\sffamily\tiny\color{ink}scannen zum Lösen}
\end{minipage}
\vspace{2mm}

{\raggedright\noindent\textcolor{ink}{\sffamily\textsc{Themen}}\enspace \themetag{ein Stein zieht} \themetag{ohne Schach}\par\smallskip{\footnotesize\textit{Hinweis:} die Partei am Zug steht im Diagramm im Schach; die Lösung beginnt mit einem Schlag.}\par}

\clearpage
\phantomsection\label{hm:KRBNNvk}
\begin{tcolorbox}[enhanced,colback=ink,colframe=ink,arc=0pt,boxrule=0pt,left=4mm,right=4mm,top=2.5mm,bottom=2.5mm]
\color{white}\sffamily{\LARGE\bfseries KRBNNvk}\hfill{\Large h\#7}\\[0.5mm]
{\small 6 Steine\quad Schwarz am Zug}\hfill{\small Nr.~556\qquad\href{https://helpman.komtera.lt/?fen=4N2R/6k1/8/7B/8/5N2/8/K7&moves=7}{\color{white}\bfseries lösen\,\ensuremath{\nearrow}}}
\end{tcolorbox}
\vspace{2mm}
\noindent\begin{minipage}[t]{36mm}\vspace{0pt}\sffamily\small\raggedright
\hmfigure{längste eindeutige}{{\Large h\#7}}\hmfigure{eindeutig bei dieser Länge}{{\Large 3} Stellungen}\hmfigure{längstes Matt}{{\Large h\#8}}\hmfigure{in der Klasse}{15,5 Milliarden Stellungen,\\ 4,8 Milliarden mit Hilfsmatt}
\end{minipage}\hfill
\begin{minipage}[t]{80mm}\vspace{0pt}\begin{tikzpicture}\begin{axis}[width=72mm,height=42mm,ybar,bar width=1.6pt,xmin=-1,xmax=17,ymin=0,ymax=1.721,axis y line*=left,axis x line*=bottom,xlabel={\scriptsize Halbzüge bis Matt},ylabel={\scriptsize Stellungen (Mrd.)},tick label style={font=\tiny},label style={font=\scriptsize},enlarge x limits=false,ylabel shift=-3pt,xlabel shift=-3pt]\addplot[fill=ink!55,draw=none] coordinates {(0,0.009234) (1,0.03897) (2,0.1107) (3,0.3799) (4,1.23) (5,1.233) (6,1.565) (7,0.1418) (8,0.08315) (9,0.0007454) (10,0.004188) (11,2.26e-07) (12,0.0008399) (13,5e-09) (14,0.0001869) (15,1e-09) (16,1.025e-06)};\draw[densely dotted] (axis cs:14,0) -- (axis cs:14,1.721);\end{axis}\begin{axis}[width=72mm,height=42mm,axis y line*=right,axis x line=none,ymode=log,xmin=-1,xmax=17,ylabel={\scriptsize eindeutig (log)},tick label style={font=\tiny},label style={font=\scriptsize},ylabel shift=-3pt]\addplot[mark=*,mark size=0.7pt,draw=black,thin] coordinates {(0,9234368) (1,35068207) (2,90919502) (3,195913020) (4,399615614) (5,43302897) (6,14240391) (7,52354) (8,89041) (9,2) (10,3999) (12,165) (14,3)};\end{axis}\end{tikzpicture}\end{minipage}
\vspace{3mm}

\noindent\begin{minipage}[t]{60mm}\vspace{0pt}\sbox\hmboard{\chessboard[setfen=4N2R/6k1/8/7B/8/5N2/8/K7 b - - 0 1]}\usebox\hmboard\\[1mm]
\makebox[\wd\hmboard]{\sffamily\hspace{3mm}\textbf{h\#7}\hfill(5+1)}\end{minipage}\hfill
\begin{minipage}[t]{52mm}\vspace{0pt}
\textcolor{ink}{\sffamily\textsc{Lösung}}\par\smallskip\small
\begin{tabular}{@{}r@{\enspace}l@{\enspace}l@{}}
\textbf{1.} & \symking{}xh8 & \symking{}b2\\
\textbf{2.} & \symking{}g8 & \symking{}c3\\
\textbf{3.} & \symking{}h7 & \symking{}d4\\
\textbf{4.} & \symking{}h6 & \symking{}e5\\
\textbf{5.} & \symking{}xh5 & \symking{}f6\\
\textbf{6.} & \symking{}g4 & \symking{}g7\\
\textbf{7.} & \symking{}h5 & \symknight{}f6\#\\
\end{tabular}\par\vspace{3mm}
\qrcode[height=12mm]{https://helpman.komtera.lt/?fen=4N2R/6k1/8/7B/8/5N2/8/K7&moves=7}\\{\sffamily\tiny\color{ink}scannen zum Lösen}
\end{minipage}
\vspace{2mm}

{\raggedright\noindent\textcolor{ink}{\sffamily\textsc{Themen}}\enspace \themetag{reines Matt} \themetag{Mustermatt} \themetag{Idealmatt} \themetag{Spiegelmatt} \themetag{Switchback} \themetag{ein Stein zieht} \themetag{Kniest} \themetag{ohne Schach}\par\smallskip{\footnotesize\textit{Hinweis:} die Partei am Zug steht im Diagramm im Schach; die Lösung beginnt mit einem Schlag.}\par}

\clearpage
\phantomsection\label{hm:KRBNPvk}
\begin{tcolorbox}[enhanced,colback=ink,colframe=ink,arc=0pt,boxrule=0pt,left=4mm,right=4mm,top=2.5mm,bottom=2.5mm]
\color{white}\sffamily{\LARGE\bfseries KRBNPvk}\hfill{\Large h\#7}\\[0.5mm]
{\small 6 Steine\quad Schwarz am Zug}\hfill{\small Nr.~557\qquad\href{https://helpman.komtera.lt/?fen=4N3/8/6k1/7R/8/8/B3P3/K7&moves=7}{\color{white}\bfseries lösen\,\ensuremath{\nearrow}}}
\end{tcolorbox}
\vspace{2mm}
\noindent\begin{minipage}[t]{36mm}\vspace{0pt}\sffamily\small\raggedright
\hmfigure{längste eindeutige}{{\Large h\#7}}\hmfigure{eindeutig bei dieser Länge}{{\Large 5} Stellungen}\hmfigure{längstes Matt}{{\Large h\#8}}\hmfigure{in der Klasse}{45,4 Milliarden Stellungen,\\ 29,5 Milliarden mit Hilfsmatt}
\end{minipage}\hfill
\begin{minipage}[t]{80mm}\vspace{0pt}\begin{tikzpicture}\begin{axis}[width=72mm,height=42mm,ybar,bar width=1.6pt,xmin=-1,xmax=17,ymin=0,ymax=11.12,axis y line*=left,axis x line*=bottom,xlabel={\scriptsize Halbzüge bis Matt},ylabel={\scriptsize Stellungen (Mrd.)},tick label style={font=\tiny},label style={font=\scriptsize},enlarge x limits=false,ylabel shift=-3pt,xlabel shift=-3pt]\addplot[fill=ink!55,draw=none] coordinates {(0,0.03905) (1,0.2012) (2,0.5181) (3,2.044) (4,6.081) (5,7.294) (6,10.11) (7,1.837) (8,1.269) (9,0.03822) (10,0.0367) (11,2.422e-05) (12,0.007151) (13,2.31e-07) (14,0.00068) (15,4e-09) (16,2.58e-07)};\draw[densely dotted] (axis cs:14,0) -- (axis cs:14,11.12);\end{axis}\begin{axis}[width=72mm,height=42mm,axis y line*=right,axis x line=none,ymode=log,xmin=-1,xmax=17,ylabel={\scriptsize eindeutig (log)},tick label style={font=\tiny},label style={font=\scriptsize},ylabel shift=-3pt]\addplot[mark=*,mark size=0.7pt,draw=black,thin] coordinates {(0,39047328) (1,164241069) (2,387526861) (3,941586621) (4,1924622131) (5,468514947) (6,205939814) (7,2464299) (8,1570640) (9,940) (10,141750) (11,17) (12,11511) (14,5)};\end{axis}\end{tikzpicture}\end{minipage}
\vspace{3mm}

\noindent\begin{minipage}[t]{60mm}\vspace{0pt}\sbox\hmboard{\chessboard[setfen=4N3/8/6k1/7R/8/8/B3P3/K7 b - - 0 1]}\usebox\hmboard\\[1mm]
\makebox[\wd\hmboard]{\sffamily\hspace{3mm}\textbf{h\#7}\hfill(5+1)}\end{minipage}\hfill
\begin{minipage}[t]{52mm}\vspace{0pt}
\textcolor{ink}{\sffamily\textsc{Lösung}}\par\smallskip\small
\begin{tabular}{@{}r@{\enspace}l@{\enspace}l@{}}
\textbf{1.} & \symking{}xh5 & e4\\
\textbf{2.} & \symking{}g6 & \symbishop{}b1\\
\textbf{3.} & \symking{}f7 & e5\\
\textbf{4.} & \symking{}xe8 & e6\\
\textbf{5.} & \symking{}f8 & e7+\\
\textbf{6.} & \symking{}g7 & e8=\symqueen{}\\
\textbf{7.} & \symking{}h6 & \symqueen{}g6\#\\
\end{tabular}\par\vspace{3mm}
\qrcode[height=12mm]{https://helpman.komtera.lt/?fen=4N3/8/6k1/7R/8/8/B3P3/K7&moves=7}\\{\sffamily\tiny\color{ink}scannen zum Lösen}
\end{minipage}
\vspace{2mm}

{\raggedright\noindent\textcolor{ink}{\sffamily\textsc{Themen}}\enspace \themetag{reines Matt} \themetag{Mustermatt} \themetag{Umwandlung} \themetag{Excelsior} \themetag{Switchback} \themetag{ein Stein zieht}\par\smallskip{\footnotesize\textit{Hinweis:} die Lösung beginnt mit einem Schlag.}\par}

\clearpage
\phantomsection\label{hm:KRBNvkb}
\begin{tcolorbox}[enhanced,colback=ink,colframe=ink,arc=0pt,boxrule=0pt,left=4mm,right=4mm,top=2.5mm,bottom=2.5mm]
\color{white}\sffamily{\LARGE\bfseries KRBNvkb}\hfill{\Large h\#7}\\[0.5mm]
{\small 6 Steine\quad Schwarz am Zug}\hfill{\small Nr.~558\qquad\href{https://helpman.komtera.lt/?fen=7R/7k/8/7N/7b/8/2B5/K7&moves=7}{\color{white}\bfseries lösen\,\ensuremath{\nearrow}}}
\end{tcolorbox}
\vspace{2mm}
\noindent\begin{minipage}[t]{36mm}\vspace{0pt}\sffamily\small\raggedright
\hmfigure{längste eindeutige}{{\Large h\#7}}\hmfigure{eindeutig bei dieser Länge}{{\Large 8} Stellungen}\hmfigure{längstes Matt}{{\Large h\#8}}\hmfigure{in der Klasse}{15,5 Milliarden Stellungen,\\ 9,17 Milliarden mit Hilfsmatt}
\end{minipage}\hfill
\begin{minipage}[t]{80mm}\vspace{0pt}\begin{tikzpicture}\begin{axis}[width=72mm,height=42mm,ybar,bar width=1.6pt,xmin=-1,xmax=17,ymin=0,ymax=3.265,axis y line*=left,axis x line*=bottom,xlabel={\scriptsize Halbzüge bis Matt},ylabel={\scriptsize Stellungen (Mrd.)},tick label style={font=\tiny},label style={font=\scriptsize},enlarge x limits=false,ylabel shift=-3pt,xlabel shift=-3pt]\addplot[fill=ink!55,draw=none] coordinates {(0,0.007604) (1,0.03471) (2,0.1321) (3,0.6025) (4,1.915) (5,2.565) (6,2.968) (7,0.6876) (8,0.2445) (9,0.01245) (10,0.003225) (11,9.974e-06) (12,0.0005951) (13,5e-09) (14,0.0001066) (15,3e-09) (16,2.997e-06)};\draw[densely dotted] (axis cs:14,0) -- (axis cs:14,3.265);\end{axis}\begin{axis}[width=72mm,height=42mm,axis y line*=right,axis x line=none,ymode=log,xmin=-1,xmax=17,ylabel={\scriptsize eindeutig (log)},tick label style={font=\tiny},label style={font=\scriptsize},ylabel shift=-3pt]\addplot[mark=*,mark size=0.7pt,draw=black,thin] coordinates {(0,7603848) (1,32944710) (2,80212441) (3,221488830) (4,452661865) (5,118979797) (6,27241789) (7,673166) (8,71583) (9,248) (10,4630) (12,132) (14,8)};\end{axis}\end{tikzpicture}\end{minipage}
\vspace{3mm}

\noindent\begin{minipage}[t]{60mm}\vspace{0pt}\sbox\hmboard{\chessboard[setfen=7R/7k/8/7N/7b/8/2B5/K7 b - - 0 1]}\usebox\hmboard\\[1mm]
\makebox[\wd\hmboard]{\sffamily\hspace{3mm}\textbf{h\#7}\hfill(4+2)}\end{minipage}\hfill
\begin{minipage}[t]{52mm}\vspace{0pt}
\textcolor{ink}{\sffamily\textsc{Lösung}}\par\smallskip\small
\begin{tabular}{@{}r@{\enspace}l@{\enspace}l@{}}
\textbf{1.} & \symking{}xh8 & \symking{}b2\\
\textbf{2.} & \symking{}g8 & \symking{}c3\\
\textbf{3.} & \symking{}f7 & \symking{}d3\\
\textbf{4.} & \symking{}g6 & \symking{}e2+\\
\textbf{5.} & \symking{}g5 & \symking{}f1\\
\textbf{6.} & \symking{}g4 & \symking{}g1\\
\textbf{7.} & \symking{}h3 & \symbishop{}f5\#\\
\end{tabular}\par\vspace{3mm}
\qrcode[height=12mm]{https://helpman.komtera.lt/?fen=7R/7k/8/7N/7b/8/2B5/K7&moves=7}\\{\sffamily\tiny\color{ink}scannen zum Lösen}
\end{minipage}
\vspace{2mm}

{\raggedright\noindent\textcolor{ink}{\sffamily\textsc{Themen}}\enspace \themetag{reines Matt} \themetag{Mustermatt} \themetag{Idealmatt} \themetag{ein Stein zieht}\par\smallskip{\footnotesize\textit{Hinweis:} das Diagramm hat keinen legalen letzten Zug, kann also in keiner Partie entstehen.}\par}

\clearpage
\phantomsection\label{hm:KRBNvkn}
\begin{tcolorbox}[enhanced,colback=ink,colframe=ink,arc=0pt,boxrule=0pt,left=4mm,right=4mm,top=2.5mm,bottom=2.5mm]
\color{white}\sffamily{\LARGE\bfseries KRBNvkn}\hfill{\Large h\#7}\\[0.5mm]
{\small 6 Steine\quad Schwarz am Zug}\hfill{\small Nr.~559\qquad\href{https://helpman.komtera.lt/?fen=8/7R/6nk/8/5B2/5N2/8/K7&moves=7}{\color{white}\bfseries lösen\,\ensuremath{\nearrow}}}
\end{tcolorbox}
\vspace{2mm}
\noindent\begin{minipage}[t]{36mm}\vspace{0pt}\sffamily\small\raggedright
\hmfigure{längste eindeutige}{{\Large h\#7}}\hmfigure{eindeutig bei dieser Länge}{{\Large 4} Stellungen}\hmfigure{längstes Matt}{{\Large h\#8}}\hmfigure{in der Klasse}{15,5 Milliarden Stellungen,\\ 9,45 Milliarden mit Hilfsmatt}
\end{minipage}\hfill
\begin{minipage}[t]{80mm}\vspace{0pt}\begin{tikzpicture}\begin{axis}[width=72mm,height=42mm,ybar,bar width=1.6pt,xmin=-1,xmax=17,ymin=0,ymax=3.677,axis y line*=left,axis x line*=bottom,xlabel={\scriptsize Halbzüge bis Matt},ylabel={\scriptsize Stellungen (Mrd.)},tick label style={font=\tiny},label style={font=\scriptsize},enlarge x limits=false,ylabel shift=-3pt,xlabel shift=-3pt]\addplot[fill=ink!55,draw=none] coordinates {(0,0.008764) (1,0.04141) (2,0.135) (3,0.5753) (4,1.706) (5,2.466) (6,3.342) (7,0.8033) (8,0.3499) (9,0.01595) (10,0.003483) (11,1.178e-05) (12,0.0008488) (13,1.42e-07) (14,0.0001525) (15,1.4e-08) (16,3.897e-06)};\draw[densely dotted] (axis cs:14,0) -- (axis cs:14,3.677);\end{axis}\begin{axis}[width=72mm,height=42mm,axis y line*=right,axis x line=none,ymode=log,xmin=-1,xmax=17,ylabel={\scriptsize eindeutig (log)},tick label style={font=\tiny},label style={font=\scriptsize},ylabel shift=-3pt]\addplot[mark=*,mark size=0.7pt,draw=black,thin] coordinates {(0,8763679) (1,39290068) (2,80834610) (3,211668441) (4,411253676) (5,132181256) (6,39558561) (7,787660) (8,74288) (9,321) (10,5127) (12,270) (14,4)};\end{axis}\end{tikzpicture}\end{minipage}
\vspace{3mm}

\noindent\begin{minipage}[t]{60mm}\vspace{0pt}\sbox\hmboard{\chessboard[setfen=8/7R/6nk/8/5B2/5N2/8/K7 b - - 0 1]}\usebox\hmboard\\[1mm]
\makebox[\wd\hmboard]{\sffamily\hspace{3mm}\textbf{h\#7}\hfill(4+2)}\end{minipage}\hfill
\begin{minipage}[t]{52mm}\vspace{0pt}
\textcolor{ink}{\sffamily\textsc{Lösung}}\par\smallskip\small
\begin{tabular}{@{}r@{\enspace}l@{\enspace}l@{}}
\textbf{1.} & \symking{}xh7 & \symking{}b2\\
\textbf{2.} & \symknight{}h4 & \symking{}c3\\
\textbf{3.} & \symking{}g6 & \symking{}d4\\
\textbf{4.} & \symking{}h5 & \symking{}e5\\
\textbf{5.} & \symking{}g4 & \symking{}f6\\
\textbf{6.} & \symking{}h3 & \symking{}g5\\
\textbf{7.} & \symknight{}g2 & \symknight{}g1\#\\
\end{tabular}\par\vspace{3mm}
\qrcode[height=12mm]{https://helpman.komtera.lt/?fen=8/7R/6nk/8/5B2/5N2/8/K7&moves=7}\\{\sffamily\tiny\color{ink}scannen zum Lösen}
\end{minipage}
\vspace{2mm}

{\raggedright\noindent\textcolor{ink}{\sffamily\textsc{Themen}}\enspace \themetag{reines Matt} \themetag{Mustermatt} \themetag{Idealmatt} \themetag{Selbstblock} \themetag{ohne Schach}\par\smallskip{\footnotesize\textit{Hinweis:} das Diagramm hat keinen legalen letzten Zug, kann also in keiner Partie entstehen.}\par}

\clearpage
\phantomsection\label{hm:KRBNvkp}
\begin{tcolorbox}[enhanced,colback=ink,colframe=ink,arc=0pt,boxrule=0pt,left=4mm,right=4mm,top=2.5mm,bottom=2.5mm]
\color{white}\sffamily{\LARGE\bfseries KRBNvkp}\hfill{\Large h\#7}\\[0.5mm]
{\small 6 Steine\quad Schwarz am Zug}\hfill{\small Nr.~560\qquad\href{https://helpman.komtera.lt/?fen=5N2/8/8/7k/2p4R/8/8/K1B5&moves=7}{\color{white}\bfseries lösen\,\ensuremath{\nearrow}}}
\end{tcolorbox}
\vspace{2mm}
\noindent\begin{minipage}[t]{36mm}\vspace{0pt}\sffamily\small\raggedright
\hmfigure{längste eindeutige}{{\Large h\#7}}\hmfigure{eindeutig bei dieser Länge}{{\Large 59} Stellungen}\hmfigure{längstes Matt}{{\Large h\#13}}\hmfigure{in der Klasse}{45,4 Milliarden Stellungen,\\ 29,4 Milliarden mit Hilfsmatt}
\end{minipage}\hfill
\begin{minipage}[t]{80mm}\vspace{0pt}\begin{tikzpicture}\begin{axis}[width=72mm,height=42mm,ybar,bar width=1.6pt,xmin=-1,xmax=27,ymin=0,ymax=11.38,axis y line*=left,axis x line*=bottom,xlabel={\scriptsize Halbzüge bis Matt},ylabel={\scriptsize Stellungen (Mrd.)},tick label style={font=\tiny},label style={font=\scriptsize},enlarge x limits=false,ylabel shift=-3pt,xlabel shift=-3pt]\addplot[fill=ink!55,draw=none] coordinates {(0,0.03413) (1,0.1601) (2,0.4684) (3,1.871) (4,5.431) (5,7.354) (6,10.35) (7,2.311) (8,1.322) (9,0.04732) (10,0.02037) (11,2.706e-05) (12,0.003475) (13,2.432e-06) (14,0.0006225) (15,4.29e-07) (16,1.019e-05) (18,1e-09) (20,1e-09) (22,1.9e-08) (24,1e-09) (26,2e-09)};\draw[densely dotted] (axis cs:14,0) -- (axis cs:14,11.38);\end{axis}\begin{axis}[width=72mm,height=42mm,axis y line*=right,axis x line=none,ymode=log,xmin=-1,xmax=27,ylabel={\scriptsize eindeutig (log)},tick label style={font=\tiny},label style={font=\scriptsize},ylabel shift=-3pt]\addplot[mark=*,mark size=0.7pt,draw=black,thin] coordinates {(0,34125748) (1,151237851) (2,349607028) (3,864938711) (4,1543783584) (5,431087222) (6,162087830) (7,2362796) (8,592925) (9,640) (10,26001) (11,35) (12,1530) (14,59)};\end{axis}\end{tikzpicture}\end{minipage}
\vspace{3mm}

\noindent\begin{minipage}[t]{60mm}\vspace{0pt}\sbox\hmboard{\chessboard[setfen=5N2/8/8/7k/2p4R/8/8/K1B5 b - - 0 1]}\usebox\hmboard\\[1mm]
\makebox[\wd\hmboard]{\sffamily\hspace{3mm}\textbf{h\#7}\hfill(4+2)}\end{minipage}\hfill
\begin{minipage}[t]{52mm}\vspace{0pt}
\textcolor{ink}{\sffamily\textsc{Lösung}}\par\smallskip\small
\begin{tabular}{@{}r@{\enspace}l@{\enspace}l@{}}
\textbf{1.} & \symking{}xh4 & \symking{}b2\\
\textbf{2.} & c3+ & \symking{}xc3\\
\textbf{3.} & \symking{}h5 & \symbishop{}b2\\
\textbf{4.} & \symking{}h6 & \symking{}d4\\
\textbf{5.} & \symking{}g7 & \symking{}e5\\
\textbf{6.} & \symking{}g8 & \symking{}f6\\
\textbf{7.} & \symking{}h8 & \symking{}f7\#\\
\end{tabular}\par\vspace{3mm}
\qrcode[height=12mm]{https://helpman.komtera.lt/?fen=5N2/8/8/7k/2p4R/8/8/K1B5&moves=7}\\{\sffamily\tiny\color{ink}scannen zum Lösen}
\end{minipage}
\vspace{2mm}

{\raggedright\noindent\textcolor{ink}{\sffamily\textsc{Themen}}\enspace \themetag{Spiegelmatt} \themetag{Switchback}\par\smallskip{\footnotesize\textit{Hinweis:} die Partei am Zug steht im Diagramm im Schach; die Lösung beginnt mit einem Schlag.}\par}

\clearpage
\phantomsection\label{hm:KRBNvkq}
\begin{tcolorbox}[enhanced,colback=ink,colframe=ink,arc=0pt,boxrule=0pt,left=4mm,right=4mm,top=2.5mm,bottom=2.5mm]
\color{white}\sffamily{\LARGE\bfseries KRBNvkq}\hfill{\Large h\#7}\\[0.5mm]
{\small 6 Steine\quad Schwarz am Zug}\hfill{\small Nr.~561\qquad\href{https://helpman.komtera.lt/?fen=7k/7R/8/8/8/8/BN1q4/K7&moves=7}{\color{white}\bfseries lösen\,\ensuremath{\nearrow}}}
\end{tcolorbox}
\vspace{2mm}
\noindent\begin{minipage}[t]{36mm}\vspace{0pt}\sffamily\small\raggedright
\hmfigure{längste eindeutige}{{\Large h\#7}}\hmfigure{eindeutig bei dieser Länge}{{\Large 25} Stellungen}\hmfigure{längstes Matt}{{\Large h\#8.5}}\hmfigure{in der Klasse}{15,5 Milliarden Stellungen,\\ 7,98 Milliarden mit Hilfsmatt}
\end{minipage}\hfill
\begin{minipage}[t]{80mm}\vspace{0pt}\begin{tikzpicture}\begin{axis}[width=72mm,height=42mm,ybar,bar width=1.6pt,xmin=-1,xmax=18,ymin=0,ymax=2.761,axis y line*=left,axis x line*=bottom,xlabel={\scriptsize Halbzüge bis Matt},ylabel={\scriptsize Stellungen (Mrd.)},tick label style={font=\tiny},label style={font=\scriptsize},enlarge x limits=false,ylabel shift=-3pt,xlabel shift=-3pt]\addplot[fill=ink!55,draw=none] coordinates {(0,0.004115) (1,0.01871) (2,0.08962) (3,0.4796) (4,1.463) (5,2.51) (6,2.436) (7,0.8813) (8,0.08072) (9,0.009444) (10,0.001722) (11,0.0001334) (12,0.0007475) (13,3.787e-05) (14,0.0001417) (15,1.154e-05) (16,6.08e-06) (17,1.5e-08)};\draw[densely dotted] (axis cs:14,0) -- (axis cs:14,2.761);\end{axis}\begin{axis}[width=72mm,height=42mm,axis y line*=right,axis x line=none,ymode=log,xmin=-1,xmax=18,ylabel={\scriptsize eindeutig (log)},tick label style={font=\tiny},label style={font=\scriptsize},ylabel shift=-3pt]\addplot[mark=*,mark size=0.7pt,draw=black,thin] coordinates {(0,4115119) (1,18037896) (2,52343881) (3,177577077) (4,312365081) (5,143885440) (6,13470398) (7,1060909) (8,37890) (9,3572) (10,6043) (11,355) (12,403) (13,22) (14,25)};\end{axis}\end{tikzpicture}\end{minipage}
\vspace{3mm}

\noindent\begin{minipage}[t]{60mm}\vspace{0pt}\sbox\hmboard{\chessboard[setfen=7k/7R/8/8/8/8/BN1q4/K7 b - - 0 1]}\usebox\hmboard\\[1mm]
\makebox[\wd\hmboard]{\sffamily\hspace{3mm}\textbf{h\#7}\hfill(4+2)}\end{minipage}\hfill
\begin{minipage}[t]{52mm}\vspace{0pt}
\textcolor{ink}{\sffamily\textsc{Lösung}}\par\smallskip\small
\begin{tabular}{@{}r@{\enspace}l@{\enspace}l@{}}
\textbf{1.} & \symking{}xh7 & \symbishop{}g8+\\
\textbf{2.} & \symking{}g6 & \symking{}a2\\
\textbf{3.} & \symking{}f5 & \symking{}a3\\
\textbf{4.} & \symking{}e4 & \symknight{}c4\\
\textbf{5.} & \symking{}d3 & \symknight{}xd2\\
\textbf{6.} & \symking{}c2 & \symknight{}b3\\
\textbf{7.} & \symking{}b1 & \symbishop{}h7\#\\
\end{tabular}\par\vspace{3mm}
\qrcode[height=12mm]{https://helpman.komtera.lt/?fen=7k/7R/8/8/8/8/BN1q4/K7&moves=7}\\{\sffamily\tiny\color{ink}scannen zum Lösen}
\end{minipage}
\vspace{2mm}

{\raggedright\noindent\textcolor{ink}{\sffamily\textsc{Themen}}\enspace \themetag{reines Matt} \themetag{Mustermatt} \themetag{Idealmatt} \themetag{Spiegelmatt} \themetag{ein Stein zieht}\par\smallskip{\footnotesize\textit{Hinweis:} die Partei am Zug steht im Diagramm im Schach; die Lösung beginnt mit einem Schlag.}\par}

\clearpage
\phantomsection\label{hm:KRBNvkr}
\begin{tcolorbox}[enhanced,colback=ink,colframe=ink,arc=0pt,boxrule=0pt,left=4mm,right=4mm,top=2.5mm,bottom=2.5mm]
\color{white}\sffamily{\LARGE\bfseries KRBNvkr}\hfill{\Large h\#7}\\[0.5mm]
{\small 6 Steine\quad Schwarz am Zug}\hfill{\small Nr.~562\qquad\href{https://helpman.komtera.lt/?fen=7k/7R/2r5/1N6/8/8/B7/K7&moves=7}{\color{white}\bfseries lösen\,\ensuremath{\nearrow}}}
\end{tcolorbox}
\vspace{2mm}
\noindent\begin{minipage}[t]{36mm}\vspace{0pt}\sffamily\small\raggedright
\hmfigure{längste eindeutige}{{\Large h\#7}}\hmfigure{eindeutig bei dieser Länge}{{\Large 8} Stellungen}\hmfigure{längstes Matt}{{\Large h\#8}}\hmfigure{in der Klasse}{15,5 Milliarden Stellungen,\\ 8,75 Milliarden mit Hilfsmatt}
\end{minipage}\hfill
\begin{minipage}[t]{80mm}\vspace{0pt}\begin{tikzpicture}\begin{axis}[width=72mm,height=42mm,ybar,bar width=1.6pt,xmin=-1,xmax=17,ymin=0,ymax=2.997,axis y line*=left,axis x line*=bottom,xlabel={\scriptsize Halbzüge bis Matt},ylabel={\scriptsize Stellungen (Mrd.)},tick label style={font=\tiny},label style={font=\scriptsize},enlarge x limits=false,ylabel shift=-3pt,xlabel shift=-3pt]\addplot[fill=ink!55,draw=none] coordinates {(0,0.006322) (1,0.03026) (2,0.1305) (3,0.6192) (4,2.071) (5,2.724) (6,2.597) (7,0.5266) (8,0.04277) (9,0.0006257) (10,0.001604) (11,7.839e-06) (12,0.0006972) (13,1.475e-06) (14,0.0001133) (15,1.8e-07) (16,2.491e-06)};\draw[densely dotted] (axis cs:14,0) -- (axis cs:14,2.997);\end{axis}\begin{axis}[width=72mm,height=42mm,axis y line*=right,axis x line=none,ymode=log,xmin=-1,xmax=17,ylabel={\scriptsize eindeutig (log)},tick label style={font=\tiny},label style={font=\scriptsize},ylabel shift=-3pt]\addplot[mark=*,mark size=0.7pt,draw=black,thin] coordinates {(0,6322165) (1,29022712) (2,81806059) (3,235481415) (4,504877172) (5,122847197) (6,11362451) (7,155585) (8,29742) (9,201) (10,5609) (11,33) (12,355) (13,1) (14,8)};\end{axis}\end{tikzpicture}\end{minipage}
\vspace{3mm}

\noindent\begin{minipage}[t]{60mm}\vspace{0pt}\sbox\hmboard{\chessboard[setfen=7k/7R/2r5/1N6/8/8/B7/K7 b - - 0 1]}\usebox\hmboard\\[1mm]
\makebox[\wd\hmboard]{\sffamily\hspace{3mm}\textbf{h\#7}\hfill(4+2)}\end{minipage}\hfill
\begin{minipage}[t]{52mm}\vspace{0pt}
\textcolor{ink}{\sffamily\textsc{Lösung}}\par\smallskip\small
\begin{tabular}{@{}r@{\enspace}l@{\enspace}l@{}}
\textbf{1.} & \symking{}xh7 & \symking{}b2\\
\textbf{2.} & \symking{}g6 & \symking{}a3\\
\textbf{3.} & \symking{}f5 & \symknight{}d4+\\
\textbf{4.} & \symking{}e4 & \symbishop{}d5+\\
\textbf{5.} & \symking{}d3 & \symknight{}b3\\
\textbf{6.} & \symking{}c2 & \symbishop{}xc6\\
\textbf{7.} & \symking{}b1 & \symbishop{}e4\#\\
\end{tabular}\par\vspace{3mm}
\qrcode[height=12mm]{https://helpman.komtera.lt/?fen=7k/7R/2r5/1N6/8/8/B7/K7&moves=7}\\{\sffamily\tiny\color{ink}scannen zum Lösen}
\end{minipage}
\vspace{2mm}

{\raggedright\noindent\textcolor{ink}{\sffamily\textsc{Themen}}\enspace \themetag{reines Matt} \themetag{Mustermatt} \themetag{Idealmatt} \themetag{Spiegelmatt} \themetag{ein Stein zieht}\par\smallskip{\footnotesize\textit{Hinweis:} die Partei am Zug steht im Diagramm im Schach; die Lösung beginnt mit einem Schlag.}\par}

\clearpage
\phantomsection\label{hm:KRBPPvk}
\begin{tcolorbox}[enhanced,colback=ink,colframe=ink,arc=0pt,boxrule=0pt,left=4mm,right=4mm,top=2.5mm,bottom=2.5mm]
\color{white}\sffamily{\LARGE\bfseries KRBPPvk}\hfill{\Large h\#7}\\[0.5mm]
{\small 6 Steine\quad Schwarz am Zug}\hfill{\small Nr.~563\qquad\href{https://helpman.komtera.lt/?fen=8/8/8/1B6/K7/1P6/P6R/6k1&moves=7}{\color{white}\bfseries lösen\,\ensuremath{\nearrow}}}
\end{tcolorbox}
\vspace{2mm}
\noindent\begin{minipage}[t]{36mm}\vspace{0pt}\sffamily\small\raggedright
\hmfigure{längste eindeutige}{{\Large h\#7}}\hmfigure{eindeutig bei dieser Länge}{{\Large 10} Stellungen}\hmfigure{längstes Matt}{{\Large h\#8}}\hmfigure{in der Klasse}{34,1 Milliarden Stellungen,\\ 11,3 Milliarden mit Hilfsmatt}
\end{minipage}\hfill
\begin{minipage}[t]{80mm}\vspace{0pt}\begin{tikzpicture}\begin{axis}[width=72mm,height=42mm,ybar,bar width=1.6pt,xmin=-1,xmax=17,ymin=0,ymax=3.695,axis y line*=left,axis x line*=bottom,xlabel={\scriptsize Halbzüge bis Matt},ylabel={\scriptsize Stellungen (Mrd.)},tick label style={font=\tiny},label style={font=\scriptsize},enlarge x limits=false,ylabel shift=-3pt,xlabel shift=-3pt]\addplot[fill=ink!55,draw=none] coordinates {(0,0.01095) (1,0.07022) (2,0.1666) (3,0.7315) (4,1.928) (5,2.329) (6,3.359) (7,1.196) (8,1.167) (9,0.1734) (10,0.1115) (11,0.003879) (12,0.004919) (13,3.969e-06) (14,0.0005661) (15,1e-09) (16,8.35e-07)};\draw[densely dotted] (axis cs:14,0) -- (axis cs:14,3.695);\end{axis}\begin{axis}[width=72mm,height=42mm,axis y line*=right,axis x line=none,ymode=log,xmin=-1,xmax=17,ylabel={\scriptsize eindeutig (log)},tick label style={font=\tiny},label style={font=\scriptsize},ylabel shift=-3pt]\addplot[mark=*,mark size=0.7pt,draw=black,thin] coordinates {(0,10950130) (1,48644182) (2,102010497) (3,259456142) (4,476048589) (5,197192624) (6,131529037) (7,9735817) (8,5234503) (9,164042) (10,171128) (11,535) (12,14728) (14,10)};\end{axis}\end{tikzpicture}\end{minipage}
\vspace{3mm}

\noindent\begin{minipage}[t]{60mm}\vspace{0pt}\sbox\hmboard{\chessboard[setfen=8/8/8/1B6/K7/1P6/P6R/6k1 b - - 0 1]}\usebox\hmboard\\[1mm]
\makebox[\wd\hmboard]{\sffamily\hspace{3mm}\textbf{h\#7}\hfill(5+1)}\end{minipage}\hfill
\begin{minipage}[t]{52mm}\vspace{0pt}
\textcolor{ink}{\sffamily\textsc{Lösung}}\par\smallskip\small
\begin{tabular}{@{}r@{\enspace}l@{\enspace}l@{}}
\textbf{1.} & \symking{}xh2 & \symking{}b4\\
\textbf{2.} & \symking{}g3 & a4\\
\textbf{3.} & \symking{}f4 & a5\\
\textbf{4.} & \symking{}e5 & a6\\
\textbf{5.} & \symking{}d6 & a7\\
\textbf{6.} & \symking{}c7 & a8=\symqueen{}\\
\textbf{7.} & \symking{}b6 & \symqueen{}b8\#\\
\end{tabular}\par\vspace{3mm}
\qrcode[height=12mm]{https://helpman.komtera.lt/?fen=8/8/8/1B6/K7/1P6/P6R/6k1&moves=7}\\{\sffamily\tiny\color{ink}scannen zum Lösen}
\end{minipage}
\vspace{2mm}

{\raggedright\noindent\textcolor{ink}{\sffamily\textsc{Themen}}\enspace \themetag{reines Matt} \themetag{Mustermatt} \themetag{Umwandlung} \themetag{Excelsior} \themetag{ein Stein zieht} \themetag{ohne Schach}\par\smallskip{\footnotesize\textit{Hinweis:} die Lösung beginnt mit einem Schlag.}\par}

\clearpage
\phantomsection\label{hm:KRBPvkb}
\begin{tcolorbox}[enhanced,colback=ink,colframe=ink,arc=0pt,boxrule=0pt,left=4mm,right=4mm,top=2.5mm,bottom=2.5mm]
\color{white}\sffamily{\LARGE\bfseries KRBPvkb}\hfill{\Large h\#7}\\[0.5mm]
{\small 6 Steine\quad Schwarz am Zug}\hfill{\small Nr.~564\qquad\href{https://helpman.komtera.lt/?fen=7b/6R1/7k/8/8/K7/P3B3/8&moves=7}{\color{white}\bfseries lösen\,\ensuremath{\nearrow}}}
\end{tcolorbox}
\vspace{2mm}
\noindent\begin{minipage}[t]{36mm}\vspace{0pt}\sffamily\small\raggedright
\hmfigure{längste eindeutige}{{\Large h\#7}}\hmfigure{eindeutig bei dieser Länge}{{\Large 11} Stellungen}\hmfigure{längstes Matt}{{\Large h\#8}}\hmfigure{in der Klasse}{45,4 Milliarden Stellungen,\\ 28,2 Milliarden mit Hilfsmatt}
\end{minipage}\hfill
\begin{minipage}[t]{80mm}\vspace{0pt}\begin{tikzpicture}\begin{axis}[width=72mm,height=42mm,ybar,bar width=1.6pt,xmin=-1,xmax=17,ymin=0,ymax=9.24,axis y line*=left,axis x line*=bottom,xlabel={\scriptsize Halbzüge bis Matt},ylabel={\scriptsize Stellungen (Mrd.)},tick label style={font=\tiny},label style={font=\scriptsize},enlarge x limits=false,ylabel shift=-3pt,xlabel shift=-3pt]\addplot[fill=ink!55,draw=none] coordinates {(0,0.01632) (1,0.0927) (2,0.3195) (3,1.554) (4,4.435) (5,6.426) (6,8.4) (7,3.708) (8,2.369) (9,0.5902) (10,0.2938) (11,0.02416) (12,0.009597) (13,4.06e-05) (14,0.0004077) (15,2.3e-08) (16,1.933e-06)};\draw[densely dotted] (axis cs:14,0) -- (axis cs:14,9.24);\end{axis}\begin{axis}[width=72mm,height=42mm,axis y line*=right,axis x line=none,ymode=log,xmin=-1,xmax=17,ylabel={\scriptsize eindeutig (log)},tick label style={font=\tiny},label style={font=\scriptsize},ylabel shift=-3pt]\addplot[mark=*,mark size=0.7pt,draw=black,thin] coordinates {(0,16317323) (1,75647192) (2,155686654) (3,455279662) (4,968775433) (5,442177326) (6,186904242) (7,24447949) (8,10519115) (9,708911) (10,211117) (11,1252) (12,6400) (14,11)};\end{axis}\end{tikzpicture}\end{minipage}
\vspace{3mm}

\noindent\begin{minipage}[t]{60mm}\vspace{0pt}\sbox\hmboard{\chessboard[setfen=7b/6R1/7k/8/8/K7/P3B3/8 b - - 0 1]}\usebox\hmboard\\[1mm]
\makebox[\wd\hmboard]{\sffamily\hspace{3mm}\textbf{h\#7}\hfill(4+2)}\end{minipage}\hfill
\begin{minipage}[t]{52mm}\vspace{0pt}
\textcolor{ink}{\sffamily\textsc{Lösung}}\par\smallskip\small
\begin{tabular}{@{}r@{\enspace}l@{\enspace}l@{}}
\textbf{1.} & \symbishop{}xg7 & \symking{}b3\\
\textbf{2.} & \symking{}g5 & a4\\
\textbf{3.} & \symking{}f4 & a5\\
\textbf{4.} & \symking{}e3 & a6\\
\textbf{5.} & \symking{}d2 & a7\\
\textbf{6.} & \symking{}c1 & a8=\symqueen{}\\
\textbf{7.} & \symking{}b1 & \symqueen{}h1\#\\
\end{tabular}\par\vspace{3mm}
\qrcode[height=12mm]{https://helpman.komtera.lt/?fen=7b/6R1/7k/8/8/K7/P3B3/8&moves=7}\\{\sffamily\tiny\color{ink}scannen zum Lösen}
\end{minipage}
\vspace{2mm}

{\raggedright\noindent\textcolor{ink}{\sffamily\textsc{Themen}}\enspace \themetag{reines Matt} \themetag{Spiegelmatt} \themetag{Umwandlung} \themetag{Excelsior} \themetag{ohne Schach}\par\smallskip{\footnotesize\textit{Hinweis:} die Lösung beginnt mit einem Schlag.}\par}

\clearpage
\phantomsection\label{hm:KRBPvkn}
\begin{tcolorbox}[enhanced,colback=ink,colframe=ink,arc=0pt,boxrule=0pt,left=4mm,right=4mm,top=2.5mm,bottom=2.5mm]
\color{white}\sffamily{\LARGE\bfseries KRBPvkn}\hfill{\Large h\#7}\\[0.5mm]
{\small 6 Steine\quad Schwarz am Zug}\hfill{\small Nr.~565\qquad\href{https://helpman.komtera.lt/?fen=8/8/8/8/4R3/4k3/3nP3/K5B1&moves=7}{\color{white}\bfseries lösen\,\ensuremath{\nearrow}}}
\end{tcolorbox}
\vspace{2mm}
\noindent\begin{minipage}[t]{36mm}\vspace{0pt}\sffamily\small\raggedright
\hmfigure{längste eindeutige}{{\Large h\#7}}\hmfigure{eindeutig bei dieser Länge}{{\Large 69} Stellungen}\hmfigure{längstes Matt}{{\Large h\#8.5}}\hmfigure{in der Klasse}{45,4 Milliarden Stellungen,\\ 29 Milliarden mit Hilfsmatt}
\end{minipage}\hfill
\begin{minipage}[t]{80mm}\vspace{0pt}\begin{tikzpicture}\begin{axis}[width=72mm,height=42mm,ybar,bar width=1.6pt,xmin=-1,xmax=18,ymin=0,ymax=9.281,axis y line*=left,axis x line*=bottom,xlabel={\scriptsize Halbzüge bis Matt},ylabel={\scriptsize Stellungen (Mrd.)},tick label style={font=\tiny},label style={font=\scriptsize},enlarge x limits=false,ylabel shift=-3pt,xlabel shift=-3pt]\addplot[fill=ink!55,draw=none] coordinates {(0,0.01908) (1,0.113) (2,0.3295) (3,1.474) (4,3.743) (5,5.562) (6,8.438) (7,4.393) (8,3.801) (9,0.8341) (10,0.3163) (11,0.01897) (12,0.005473) (13,4.393e-06) (14,0.000656) (15,2.07e-07) (16,4.581e-06) (17,2e-09)};\draw[densely dotted] (axis cs:14,0) -- (axis cs:14,9.281);\end{axis}\begin{axis}[width=72mm,height=42mm,axis y line*=right,axis x line=none,ymode=log,xmin=-1,xmax=18,ylabel={\scriptsize eindeutig (log)},tick label style={font=\tiny},label style={font=\scriptsize},ylabel shift=-3pt]\addplot[mark=*,mark size=0.7pt,draw=black,thin] coordinates {(0,19076850) (1,90878360) (2,153545518) (3,393852593) (4,739176466) (5,463264545) (6,302456302) (7,36472313) (8,13060929) (9,655030) (10,115584) (11,93) (12,11688) (14,69)};\end{axis}\end{tikzpicture}\end{minipage}
\vspace{3mm}

\noindent\begin{minipage}[t]{60mm}\vspace{0pt}\sbox\hmboard{\chessboard[setfen=8/8/8/8/4R3/4k3/3nP3/K5B1 b - - 0 1]}\usebox\hmboard\\[1mm]
\makebox[\wd\hmboard]{\sffamily\hspace{3mm}\textbf{h\#7}\hfill(4+2)}\end{minipage}\hfill
\begin{minipage}[t]{52mm}\vspace{0pt}
\textcolor{ink}{\sffamily\textsc{Lösung}}\par\smallskip\small
\begin{tabular}{@{}r@{\enspace}l@{\enspace}l@{}}
\textbf{1.} & \symking{}xe4 & e3\\
\textbf{2.} & \symking{}d3 & e4\\
\textbf{3.} & \symknight{}c4 & e5\\
\textbf{4.} & \symknight{}d6 & exd6\\
\textbf{5.} & \symking{}c4 & d7\\
\textbf{6.} & \symking{}b5 & d8=\symqueen{}\\
\textbf{7.} & \symking{}a6 & \symqueen{}b6\#\\
\end{tabular}\par\vspace{3mm}
\qrcode[height=12mm]{https://helpman.komtera.lt/?fen=8/8/8/8/4R3/4k3/3nP3/K5B1&moves=7}\\{\sffamily\tiny\color{ink}scannen zum Lösen}
\end{minipage}
\vspace{2mm}

{\raggedright\noindent\textcolor{ink}{\sffamily\textsc{Themen}}\enspace \themetag{reines Matt} \themetag{Mustermatt} \themetag{Umwandlung} \themetag{Excelsior} \themetag{ein Stein zieht} \themetag{ohne Schach} \themetag{Umnov}\par\smallskip{\footnotesize\textit{Hinweis:} das Diagramm hat keinen legalen letzten Zug, kann also in keiner Partie entstehen.}\par}

\clearpage
\phantomsection\label{hm:KRBPvkq}
\begin{tcolorbox}[enhanced,colback=ink,colframe=ink,arc=0pt,boxrule=0pt,left=4mm,right=4mm,top=2.5mm,bottom=2.5mm]
\color{white}\sffamily{\LARGE\bfseries KRBPvkq}\hfill{\Large h\#7}\\[0.5mm]
{\small 6 Steine\quad Schwarz am Zug}\hfill{\small Nr.~566\qquad\href{https://helpman.komtera.lt/?fen=8/5B2/8/8/8/5P2/q7/kRK5&moves=7}{\color{white}\bfseries lösen\,\ensuremath{\nearrow}}}
\end{tcolorbox}
\vspace{2mm}
\noindent\begin{minipage}[t]{36mm}\vspace{0pt}\sffamily\small\raggedright
\hmfigure{längste eindeutige}{{\Large h\#7}}\hmfigure{eindeutig bei dieser Länge}{{\Large 39} Stellungen}\hmfigure{längstes Matt}{{\Large h\#8.5}}\hmfigure{in der Klasse}{45,4 Milliarden Stellungen,\\ 24,7 Milliarden mit Hilfsmatt}
\end{minipage}\hfill
\begin{minipage}[t]{80mm}\vspace{0pt}\begin{tikzpicture}\begin{axis}[width=72mm,height=42mm,ybar,bar width=1.6pt,xmin=-1,xmax=18,ymin=0,ymax=8.376,axis y line*=left,axis x line*=bottom,xlabel={\scriptsize Halbzüge bis Matt},ylabel={\scriptsize Stellungen (Mrd.)},tick label style={font=\tiny},label style={font=\scriptsize},enlarge x limits=false,ylabel shift=-3pt,xlabel shift=-3pt]\addplot[fill=ink!55,draw=none] coordinates {(0,0.008731) (1,0.04903) (2,0.2252) (3,1.283) (4,3.607) (5,6.707) (6,7.614) (7,4.183) (8,0.8129) (9,0.1552) (10,0.00353) (11,0.007009) (12,0.002771) (13,0.0003787) (14,0.0004209) (15,5.965e-05) (16,3.257e-06) (17,9.33e-07)};\draw[densely dotted] (axis cs:14,0) -- (axis cs:14,8.376);\end{axis}\begin{axis}[width=72mm,height=42mm,axis y line*=right,axis x line=none,ymode=log,xmin=-1,xmax=18,ylabel={\scriptsize eindeutig (log)},tick label style={font=\tiny},label style={font=\scriptsize},ylabel shift=-3pt]\addplot[mark=*,mark size=0.7pt,draw=black,thin] coordinates {(0,8730628) (1,41280027) (2,108734251) (3,408197446) (4,728428686) (5,471569673) (6,95912019) (7,14050454) (8,96256) (9,184819) (10,16377) (11,10408) (12,4596) (13,801) (14,39)};\end{axis}\end{tikzpicture}\end{minipage}
\vspace{3mm}

\noindent\begin{minipage}[t]{60mm}\vspace{0pt}\sbox\hmboard{\chessboard[setfen=8/5B2/8/8/8/5P2/q7/kRK5 b - - 0 1]}\usebox\hmboard\\[1mm]
\makebox[\wd\hmboard]{\sffamily\hspace{3mm}\textbf{h\#7}\hfill(4+2)}\end{minipage}\hfill
\begin{minipage}[t]{52mm}\vspace{0pt}
\textcolor{ink}{\sffamily\textsc{Lösung}}\par\smallskip\small
\begin{tabular}{@{}r@{\enspace}l@{\enspace}l@{}}
\textbf{1.} & \symqueen{}xb1+ & \symking{}d2\\
\textbf{2.} & \symqueen{}e4 & fxe4\\
\textbf{3.} & \symking{}b2 & e5\\
\textbf{4.} & \symking{}a3 & e6\\
\textbf{5.} & \symking{}b4 & e7\\
\textbf{6.} & \symking{}c5 & e8=\symqueen{}\\
\textbf{7.} & \symking{}d4 & \symqueen{}e3\#\\
\end{tabular}\par\vspace{3mm}
\qrcode[height=12mm]{https://helpman.komtera.lt/?fen=8/5B2/8/8/8/5P2/q7/kRK5&moves=7}\\{\sffamily\tiny\color{ink}scannen zum Lösen}
\end{minipage}
\vspace{2mm}

{\raggedright\noindent\textcolor{ink}{\sffamily\textsc{Themen}}\enspace \themetag{Umwandlung}\par\smallskip{\footnotesize\textit{Hinweis:} die Partei am Zug steht im Diagramm im Schach; die Lösung beginnt mit einem Schlag.}\par}

\clearpage
\phantomsection\label{hm:KRBPvkr}
\begin{tcolorbox}[enhanced,colback=ink,colframe=ink,arc=0pt,boxrule=0pt,left=4mm,right=4mm,top=2.5mm,bottom=2.5mm]
\color{white}\sffamily{\LARGE\bfseries KRBPvkr}\hfill{\Large h\#7}\\[0.5mm]
{\small 6 Steine\quad Schwarz am Zug}\hfill{\small Nr.~567\qquad\href{https://helpman.komtera.lt/?fen=8/8/3B4/8/8/8/2RP4/K1kr4&moves=7}{\color{white}\bfseries lösen\,\ensuremath{\nearrow}}}
\end{tcolorbox}
\vspace{2mm}
\noindent\begin{minipage}[t]{36mm}\vspace{0pt}\sffamily\small\raggedright
\hmfigure{längste eindeutige}{{\Large h\#7}}\hmfigure{eindeutig bei dieser Länge}{{\Large 17} Stellungen}\hmfigure{längstes Matt}{{\Large h\#8.5}}\hmfigure{in der Klasse}{45,4 Milliarden Stellungen,\\ 27 Milliarden mit Hilfsmatt}
\end{minipage}\hfill
\begin{minipage}[t]{80mm}\vspace{0pt}\begin{tikzpicture}\begin{axis}[width=72mm,height=42mm,ybar,bar width=1.6pt,xmin=-1,xmax=18,ymin=0,ymax=9.584,axis y line*=left,axis x line*=bottom,xlabel={\scriptsize Halbzüge bis Matt},ylabel={\scriptsize Stellungen (Mrd.)},tick label style={font=\tiny},label style={font=\scriptsize},enlarge x limits=false,ylabel shift=-3pt,xlabel shift=-3pt]\addplot[fill=ink!55,draw=none] coordinates {(0,0.01342) (1,0.07985) (2,0.3179) (3,1.593) (4,4.902) (5,7.508) (6,8.713) (7,3.17) (8,0.6207) (9,0.04046) (10,0.003906) (11,0.0001466) (12,0.003168) (13,1.405e-05) (14,0.0003243) (15,3.768e-06) (16,1.493e-06) (17,8.1e-08)};\draw[densely dotted] (axis cs:14,0) -- (axis cs:14,9.584);\end{axis}\begin{axis}[width=72mm,height=42mm,axis y line*=right,axis x line=none,ymode=log,xmin=-1,xmax=18,ylabel={\scriptsize eindeutig (log)},tick label style={font=\tiny},label style={font=\scriptsize},ylabel shift=-3pt]\addplot[mark=*,mark size=0.7pt,draw=black,thin] coordinates {(0,13424481) (1,66170263) (2,163497455) (3,510976817) (4,1205176635) (5,537990451) (6,116635771) (7,5649835) (8,109600) (9,4343) (10,25532) (11,296) (12,5308) (13,64) (14,17)};\end{axis}\end{tikzpicture}\end{minipage}
\vspace{3mm}

\noindent\begin{minipage}[t]{60mm}\vspace{0pt}\sbox\hmboard{\chessboard[setfen=8/8/3B4/8/8/8/2RP4/K1kr4 b - - 0 1]}\usebox\hmboard\\[1mm]
\makebox[\wd\hmboard]{\sffamily\hspace{3mm}\textbf{h\#7}\hfill(4+2)}\end{minipage}\hfill
\begin{minipage}[t]{52mm}\vspace{0pt}
\textcolor{ink}{\sffamily\textsc{Lösung}}\par\smallskip\small
\begin{tabular}{@{}r@{\enspace}l@{\enspace}l@{}}
\textbf{1.} & \symking{}xc2+ & \symking{}a2\\
\textbf{2.} & \symrook{}c1 & d4\\
\textbf{3.} & \symking{}d3 & d5\\
\textbf{4.} & \symrook{}c6 & dxc6\\
\textbf{5.} & \symking{}c4 & c7\\
\textbf{6.} & \symking{}b5 & c8=\symqueen{}\\
\textbf{7.} & \symking{}a4 & \symqueen{}a6\#\\
\end{tabular}\par\vspace{3mm}
\qrcode[height=12mm]{https://helpman.komtera.lt/?fen=8/8/3B4/8/8/8/2RP4/K1kr4&moves=7}\\{\sffamily\tiny\color{ink}scannen zum Lösen}
\end{minipage}
\vspace{2mm}

{\raggedright\noindent\textcolor{ink}{\sffamily\textsc{Themen}}\enspace \themetag{Spiegelmatt} \themetag{Umwandlung} \themetag{Excelsior}\par\smallskip{\footnotesize\textit{Hinweis:} die Partei am Zug steht im Diagramm im Schach; die Lösung beginnt mit einem Schlag.}\par}

\clearpage
\phantomsection\label{hm:KRBvkrb}
\begin{tcolorbox}[enhanced,colback=ink,colframe=ink,arc=0pt,boxrule=0pt,left=4mm,right=4mm,top=2.5mm,bottom=2.5mm]
\color{white}\sffamily{\LARGE\bfseries KRBvkrb}\hfill{\Large h\#7}\\[0.5mm]
{\small 6 Steine\quad Schwarz am Zug}\hfill{\small Nr.~568\qquad\href{https://helpman.komtera.lt/?fen=6kR/6b1/8/8/8/8/B4r2/1K6&moves=7}{\color{white}\bfseries lösen\,\ensuremath{\nearrow}}}
\end{tcolorbox}
\vspace{2mm}
\noindent\begin{minipage}[t]{36mm}\vspace{0pt}\sffamily\small\raggedright
\hmfigure{längste eindeutige}{{\Large h\#7}}\hmfigure{eindeutig bei dieser Länge}{{\Large 1} Stellung}\hmfigure{längstes Matt}{{\Large h\#8.5}}\hmfigure{in der Klasse}{15,5 Milliarden Stellungen,\\ 8,48 Milliarden mit Hilfsmatt}
\end{minipage}\hfill
\begin{minipage}[t]{80mm}\vspace{0pt}\begin{tikzpicture}\begin{axis}[width=72mm,height=42mm,ybar,bar width=1.6pt,xmin=-1,xmax=18,ymin=0,ymax=3.044,axis y line*=left,axis x line*=bottom,xlabel={\scriptsize Halbzüge bis Matt},ylabel={\scriptsize Stellungen (Mrd.)},tick label style={font=\tiny},label style={font=\scriptsize},enlarge x limits=false,ylabel shift=-3pt,xlabel shift=-3pt]\addplot[fill=ink!55,draw=none] coordinates {(0,0.002638) (1,0.01281) (2,0.05861) (3,0.3296) (4,1.144) (5,2.385) (6,2.767) (7,1.474) (8,0.2648) (9,0.03695) (10,0.001181) (11,0.0001638) (12,0.0005265) (13,1.996e-05) (14,7.473e-05) (15,2.5e-08) (16,6.6e-08) (17,1e-08)};\draw[densely dotted] (axis cs:14,0) -- (axis cs:14,3.044);\end{axis}\begin{axis}[width=72mm,height=42mm,axis y line*=right,axis x line=none,ymode=log,xmin=-1,xmax=18,ylabel={\scriptsize eindeutig (log)},tick label style={font=\tiny},label style={font=\scriptsize},ylabel shift=-3pt]\addplot[mark=*,mark size=0.7pt,draw=black,thin] coordinates {(0,2637854) (1,12628338) (2,33380077) (3,135331861) (4,341741392) (5,191806618) (6,40265625) (7,3947045) (8,37754) (9,2862) (10,4016) (11,202) (12,379) (14,1)};\end{axis}\end{tikzpicture}\end{minipage}
\vspace{3mm}

\noindent\begin{minipage}[t]{60mm}\vspace{0pt}\sbox\hmboard{\chessboard[setfen=6kR/6b1/8/8/8/8/B4r2/1K6 b - - 0 1]}\usebox\hmboard\\[1mm]
\makebox[\wd\hmboard]{\sffamily\hspace{3mm}\textbf{h\#7}\hfill(3+3)}\end{minipage}\hfill
\begin{minipage}[t]{52mm}\vspace{0pt}
\textcolor{ink}{\sffamily\textsc{Lösung}}\par\smallskip\small
\begin{tabular}{@{}r@{\enspace}l@{\enspace}l@{}}
\textbf{1.} & \symking{}xh8 & \symbishop{}b3\\
\textbf{2.} & \symrook{}f8 & \symking{}c2\\
\textbf{3.} & \symbishop{}f6 & \symking{}d3\\
\textbf{4.} & \symking{}g7 & \symking{}c4\\
\textbf{5.} & \symking{}f7 & \symking{}d5\\
\textbf{6.} & \symking{}e8 & \symking{}e6\\
\textbf{7.} & \symbishop{}d8 & \symbishop{}a4\#\\
\end{tabular}\par\vspace{3mm}
\qrcode[height=12mm]{https://helpman.komtera.lt/?fen=6kR/6b1/8/8/8/8/B4r2/1K6&moves=7}\\{\sffamily\tiny\color{ink}scannen zum Lösen}
\end{minipage}
\vspace{2mm}

{\raggedright\noindent\textcolor{ink}{\sffamily\textsc{Themen}}\enspace \themetag{Selbstblock} \themetag{ohne Schach}\par\smallskip{\footnotesize\textit{Hinweis:} das Diagramm hat keinen legalen letzten Zug, kann also in keiner Partie entstehen.}\par}

\clearpage
\phantomsection\label{hm:KRNNPvk}
\begin{tcolorbox}[enhanced,colback=ink,colframe=ink,arc=0pt,boxrule=0pt,left=4mm,right=4mm,top=2.5mm,bottom=2.5mm]
\color{white}\sffamily{\LARGE\bfseries KRNNPvk}\hfill{\Large h\#7}\\[0.5mm]
{\small 6 Steine\quad Schwarz am Zug}\hfill{\small Nr.~569\qquad\href{https://helpman.komtera.lt/?fen=5N1R/7k/4N3/8/8/K7/P7/8&moves=7}{\color{white}\bfseries lösen\,\ensuremath{\nearrow}}}
\end{tcolorbox}
\vspace{2mm}
\noindent\begin{minipage}[t]{36mm}\vspace{0pt}\sffamily\small\raggedright
\hmfigure{längste eindeutige}{{\Large h\#7}}\hmfigure{eindeutig bei dieser Länge}{{\Large 7} Stellungen}\hmfigure{längstes Matt}{{\Large h\#8}}\hmfigure{in der Klasse}{45,4 Milliarden Stellungen,\\ 15 Milliarden mit Hilfsmatt}
\end{minipage}\hfill
\begin{minipage}[t]{80mm}\vspace{0pt}\begin{tikzpicture}\begin{axis}[width=72mm,height=42mm,ybar,bar width=1.6pt,xmin=-1,xmax=17,ymin=0,ymax=5.67,axis y line*=left,axis x line*=bottom,xlabel={\scriptsize Halbzüge bis Matt},ylabel={\scriptsize Stellungen (Mrd.)},tick label style={font=\tiny},label style={font=\scriptsize},enlarge x limits=false,ylabel shift=-3pt,xlabel shift=-3pt]\addplot[fill=ink!55,draw=none] coordinates {(0,0.01592) (1,0.09021) (2,0.2196) (3,0.9371) (4,2.672) (5,3.638) (6,5.154) (7,1.274) (8,0.9438) (9,0.04972) (10,0.02714) (11,2.817e-05) (12,0.003358) (13,1e-08) (14,0.0006669) (16,1.96e-07)};\draw[densely dotted] (axis cs:14,0) -- (axis cs:14,5.67);\end{axis}\begin{axis}[width=72mm,height=42mm,axis y line*=right,axis x line=none,ymode=log,xmin=-1,xmax=17,ylabel={\scriptsize eindeutig (log)},tick label style={font=\tiny},label style={font=\scriptsize},ylabel shift=-3pt]\addplot[mark=*,mark size=0.7pt,draw=black,thin] coordinates {(0,15916529) (1,72491738) (2,160644624) (3,405028869) (4,792108925) (5,249441607) (6,112334681) (7,1844921) (8,797193) (9,676) (10,51826) (12,7061) (14,7)};\end{axis}\end{tikzpicture}\end{minipage}
\vspace{3mm}

\noindent\begin{minipage}[t]{60mm}\vspace{0pt}\sbox\hmboard{\chessboard[setfen=5N1R/7k/4N3/8/8/K7/P7/8 b - - 0 1]}\usebox\hmboard\\[1mm]
\makebox[\wd\hmboard]{\sffamily\hspace{3mm}\textbf{h\#7}\hfill(5+1)}\end{minipage}\hfill
\begin{minipage}[t]{52mm}\vspace{0pt}
\textcolor{ink}{\sffamily\textsc{Lösung}}\par\smallskip\small
\begin{tabular}{@{}r@{\enspace}l@{\enspace}l@{}}
\textbf{1.} & \symking{}xh8 & \symking{}b4\\
\textbf{2.} & \symking{}g8 & a4\\
\textbf{3.} & \symking{}f7 & a5\\
\textbf{4.} & \symking{}e7 & a6\\
\textbf{5.} & \symking{}d6 & a7\\
\textbf{6.} & \symking{}c6 & a8=\symqueen{}+\\
\textbf{7.} & \symking{}b6 & \symknight{}d7\#\\
\end{tabular}\par\vspace{3mm}
\qrcode[height=12mm]{https://helpman.komtera.lt/?fen=5N1R/7k/4N3/8/8/K7/P7/8&moves=7}\\{\sffamily\tiny\color{ink}scannen zum Lösen}
\end{minipage}
\vspace{2mm}

{\raggedright\noindent\textcolor{ink}{\sffamily\textsc{Themen}}\enspace \themetag{Spiegelmatt} \themetag{Umwandlung} \themetag{Excelsior} \themetag{ein Stein zieht}\par\smallskip{\footnotesize\textit{Hinweis:} das Diagramm hat keinen legalen letzten Zug, kann also in keiner Partie entstehen.}\par}

\clearpage
\phantomsection\label{hm:KRNNvkb}
\begin{tcolorbox}[enhanced,colback=ink,colframe=ink,arc=0pt,boxrule=0pt,left=4mm,right=4mm,top=2.5mm,bottom=2.5mm]
\color{white}\sffamily{\LARGE\bfseries KRNNvkb}\hfill{\Large h\#7}\\[0.5mm]
{\small 6 Steine\quad Schwarz am Zug}\hfill{\small Nr.~570\qquad\href{https://helpman.komtera.lt/?fen=kR5N/b7/8/8/8/8/7N/3K4&moves=7}{\color{white}\bfseries lösen\,\ensuremath{\nearrow}}}
\end{tcolorbox}
\vspace{2mm}
\noindent\begin{minipage}[t]{36mm}\vspace{0pt}\sffamily\small\raggedright
\hmfigure{längste eindeutige}{{\Large h\#7}}\hmfigure{eindeutig bei dieser Länge}{{\Large 1} Stellung}\hmfigure{längstes Matt}{{\Large h\#8}}\hmfigure{in der Klasse}{15,5 Milliarden Stellungen,\\ 4,68 Milliarden mit Hilfsmatt}
\end{minipage}\hfill
\begin{minipage}[t]{80mm}\vspace{0pt}\begin{tikzpicture}\begin{axis}[width=72mm,height=42mm,ybar,bar width=1.6pt,xmin=-1,xmax=17,ymin=0,ymax=1.804,axis y line*=left,axis x line*=bottom,xlabel={\scriptsize Halbzüge bis Matt},ylabel={\scriptsize Stellungen (Mrd.)},tick label style={font=\tiny},label style={font=\scriptsize},enlarge x limits=false,ylabel shift=-3pt,xlabel shift=-3pt]\addplot[fill=ink!55,draw=none] coordinates {(0,0.002959) (1,0.01497) (2,0.05329) (3,0.2508) (4,0.7518) (5,1.254) (6,1.64) (7,0.5155) (8,0.1861) (9,0.0113) (10,0.001567) (11,2.597e-06) (12,0.0002424) (14,3.907e-05) (15,2e-09) (16,4.73e-07)};\draw[densely dotted] (axis cs:14,0) -- (axis cs:14,1.804);\end{axis}\begin{axis}[width=72mm,height=42mm,axis y line*=right,axis x line=none,ymode=log,xmin=-1,xmax=17,ylabel={\scriptsize eindeutig (log)},tick label style={font=\tiny},label style={font=\scriptsize},ylabel shift=-3pt]\addplot[mark=*,mark size=0.7pt,draw=black,thin] coordinates {(0,2958674) (1,14156560) (2,30365493) (3,86412627) (4,182426083) (5,74651588) (6,21865804) (7,594950) (8,33418) (9,80) (10,1477) (12,19) (14,1)};\end{axis}\end{tikzpicture}\end{minipage}
\vspace{3mm}

\noindent\begin{minipage}[t]{60mm}\vspace{0pt}\sbox\hmboard{\chessboard[setfen=kR5N/b7/8/8/8/8/7N/3K4 b - - 0 1]}\usebox\hmboard\\[1mm]
\makebox[\wd\hmboard]{\sffamily\hspace{3mm}\textbf{h\#7}\hfill(4+2)}\end{minipage}\hfill
\begin{minipage}[t]{52mm}\vspace{0pt}
\textcolor{ink}{\sffamily\textsc{Lösung}}\par\smallskip\small
\begin{tabular}{@{}r@{\enspace}l@{\enspace}l@{}}
\textbf{1.} & \symking{}xb8 & \symking{}e2\\
\textbf{2.} & \symking{}c7 & \symking{}f3\\
\textbf{3.} & \symking{}d6 & \symking{}g4\\
\textbf{4.} & \symking{}e5 & \symking{}h5\\
\textbf{5.} & \symking{}f4 & \symknight{}g6+\\
\textbf{6.} & \symking{}g3 & \symknight{}f1+\\
\textbf{7.} & \symking{}h3 & \symknight{}f4\#\\
\end{tabular}\par\vspace{3mm}
\qrcode[height=12mm]{https://helpman.komtera.lt/?fen=kR5N/b7/8/8/8/8/7N/3K4&moves=7}\\{\sffamily\tiny\color{ink}scannen zum Lösen}
\end{minipage}
\vspace{2mm}

{\raggedright\noindent\textcolor{ink}{\sffamily\textsc{Themen}}\enspace \themetag{reines Matt} \themetag{Mustermatt} \themetag{Spiegelmatt} \themetag{ein Stein zieht}\par\smallskip{\footnotesize\textit{Hinweis:} die Partei am Zug steht im Diagramm im Schach; die Lösung beginnt mit einem Schlag.}\par}

\clearpage
\phantomsection\label{hm:KRNNvkp}
\begin{tcolorbox}[enhanced,colback=ink,colframe=ink,arc=0pt,boxrule=0pt,left=4mm,right=4mm,top=2.5mm,bottom=2.5mm]
\color{white}\sffamily{\LARGE\bfseries KRNNvkp}\hfill{\Large h\#7}\\[0.5mm]
{\small 6 Steine\quad Schwarz am Zug}\hfill{\small Nr.~571\qquad\href{https://helpman.komtera.lt/?fen=7R/7k/8/4pN2/7N/8/8/K7&moves=7}{\color{white}\bfseries lösen\,\ensuremath{\nearrow}}}
\end{tcolorbox}
\vspace{2mm}
\noindent\begin{minipage}[t]{36mm}\vspace{0pt}\sffamily\small\raggedright
\hmfigure{längste eindeutige}{{\Large h\#7}}\hmfigure{eindeutig bei dieser Länge}{{\Large 47} Stellungen}\hmfigure{längstes Matt}{{\Large h\#11}}\hmfigure{in der Klasse}{45,4 Milliarden Stellungen,\\ 15 Milliarden mit Hilfsmatt}
\end{minipage}\hfill
\begin{minipage}[t]{80mm}\vspace{0pt}\begin{tikzpicture}\begin{axis}[width=72mm,height=42mm,ybar,bar width=1.6pt,xmin=-1,xmax=23,ymin=0,ymax=5.757,axis y line*=left,axis x line*=bottom,xlabel={\scriptsize Halbzüge bis Matt},ylabel={\scriptsize Stellungen (Mrd.)},tick label style={font=\tiny},label style={font=\scriptsize},enlarge x limits=false,ylabel shift=-3pt,xlabel shift=-3pt]\addplot[fill=ink!55,draw=none] coordinates {(0,0.01413) (1,0.07153) (2,0.1994) (3,0.8193) (4,2.288) (5,3.569) (6,5.233) (7,1.623) (8,1.054) (9,0.07007) (10,0.0236) (11,6.446e-05) (12,0.001698) (13,1.66e-06) (14,0.0003822) (15,7.11e-07) (16,3.178e-06) (18,1.5e-08) (20,3e-09) (22,2e-09)};\draw[densely dotted] (axis cs:14,0) -- (axis cs:14,5.757);\end{axis}\begin{axis}[width=72mm,height=42mm,axis y line*=right,axis x line=none,ymode=log,xmin=-1,xmax=23,ylabel={\scriptsize eindeutig (log)},tick label style={font=\tiny},label style={font=\scriptsize},ylabel shift=-3pt]\addplot[mark=*,mark size=0.7pt,draw=black,thin] coordinates {(0,14134841) (1,67422118) (2,145288945) (3,352983600) (4,622805619) (5,221717654) (6,93869459) (7,2392672) (8,590841) (9,1471) (10,10860) (11,2) (12,784) (13,3) (14,47)};\end{axis}\end{tikzpicture}\end{minipage}
\vspace{3mm}

\noindent\begin{minipage}[t]{60mm}\vspace{0pt}\sbox\hmboard{\chessboard[setfen=7R/7k/8/4pN2/7N/8/8/K7 b - - 0 1]}\usebox\hmboard\\[1mm]
\makebox[\wd\hmboard]{\sffamily\hspace{3mm}\textbf{h\#7}\hfill(4+2)}\end{minipage}\hfill
\begin{minipage}[t]{52mm}\vspace{0pt}
\textcolor{ink}{\sffamily\textsc{Lösung}}\par\smallskip\small
\begin{tabular}{@{}r@{\enspace}l@{\enspace}l@{}}
\textbf{1.} & \symking{}xh8 & \symking{}b2\\
\textbf{2.} & e4 & \symking{}c3\\
\textbf{3.} & e3 & \symking{}d4\\
\textbf{4.} & e2 & \symking{}e5\\
\textbf{5.} & e1=\symrook{}+ & \symking{}f6\\
\textbf{6.} & \symrook{}e7 & \symknight{}h6\\
\textbf{7.} & \symrook{}h7 & \symknight{}g6\#\\
\end{tabular}\par\vspace{3mm}
\qrcode[height=12mm]{https://helpman.komtera.lt/?fen=7R/7k/8/4pN2/7N/8/8/K7&moves=7}\\{\sffamily\tiny\color{ink}scannen zum Lösen}
\end{minipage}
\vspace{2mm}

{\raggedright\noindent\textcolor{ink}{\sffamily\textsc{Themen}}\enspace \themetag{reines Matt} \themetag{Mustermatt} \themetag{Idealmatt} \themetag{Umwandlung} \themetag{Unterverwandlung} \themetag{Selbstblock} \themetag{Kniest}\par\smallskip{\footnotesize\textit{Hinweis:} die Partei am Zug steht im Diagramm im Schach; die Lösung beginnt mit einem Schlag.}\par}

\clearpage
\phantomsection\label{hm:KRNNvkq}
\begin{tcolorbox}[enhanced,colback=ink,colframe=ink,arc=0pt,boxrule=0pt,left=4mm,right=4mm,top=2.5mm,bottom=2.5mm]
\color{white}\sffamily{\LARGE\bfseries KRNNvkq}\hfill{\Large h\#7}\\[0.5mm]
{\small 6 Steine\quad Schwarz am Zug}\hfill{\small Nr.~572\qquad\href{https://helpman.komtera.lt/?fen=kR5N/q7/8/8/8/8/7N/3K4&moves=7}{\color{white}\bfseries lösen\,\ensuremath{\nearrow}}}
\end{tcolorbox}
\vspace{2mm}
\noindent\begin{minipage}[t]{36mm}\vspace{0pt}\sffamily\small\raggedright
\hmfigure{längste eindeutige}{{\Large h\#7}}\hmfigure{eindeutig bei dieser Länge}{{\Large 1} Stellung}\hmfigure{längstes Matt}{{\Large h\#8}}\hmfigure{in der Klasse}{15,5 Milliarden Stellungen,\\ 4,08 Milliarden mit Hilfsmatt}
\end{minipage}\hfill
\begin{minipage}[t]{80mm}\vspace{0pt}\begin{tikzpicture}\begin{axis}[width=72mm,height=42mm,ybar,bar width=1.6pt,xmin=-1,xmax=17,ymin=0,ymax=1.402,axis y line*=left,axis x line*=bottom,xlabel={\scriptsize Halbzüge bis Matt},ylabel={\scriptsize Stellungen (Mrd.)},tick label style={font=\tiny},label style={font=\scriptsize},enlarge x limits=false,ylabel shift=-3pt,xlabel shift=-3pt]\addplot[fill=ink!55,draw=none] coordinates {(0,0.001597) (1,0.008153) (2,0.0373) (3,0.2039) (4,0.573) (5,1.155) (6,1.275) (7,0.6537) (8,0.1498) (9,0.0241) (10,0.001483) (11,0.0001556) (12,0.0003281) (13,7.221e-06) (14,7.354e-05) (15,8.983e-06) (16,1.31e-06)};\draw[densely dotted] (axis cs:14,0) -- (axis cs:14,1.402);\end{axis}\begin{axis}[width=72mm,height=42mm,axis y line*=right,axis x line=none,ymode=log,xmin=-1,xmax=17,ylabel={\scriptsize eindeutig (log)},tick label style={font=\tiny},label style={font=\scriptsize},ylabel shift=-3pt]\addplot[mark=*,mark size=0.7pt,draw=black,thin] coordinates {(0,1597438) (1,7768251) (2,20575041) (3,71855429) (4,117764246) (5,72706471) (6,12953110) (7,1457666) (8,32743) (9,3456) (10,2898) (11,80) (12,217) (13,6) (14,1)};\end{axis}\end{tikzpicture}\end{minipage}
\vspace{3mm}

\noindent\begin{minipage}[t]{60mm}\vspace{0pt}\sbox\hmboard{\chessboard[setfen=kR5N/q7/8/8/8/8/7N/3K4 b - - 0 1]}\usebox\hmboard\\[1mm]
\makebox[\wd\hmboard]{\sffamily\hspace{3mm}\textbf{h\#7}\hfill(4+2)}\end{minipage}\hfill
\begin{minipage}[t]{52mm}\vspace{0pt}
\textcolor{ink}{\sffamily\textsc{Lösung}}\par\smallskip\small
\begin{tabular}{@{}r@{\enspace}l@{\enspace}l@{}}
\textbf{1.} & \symking{}xb8 & \symking{}e2\\
\textbf{2.} & \symking{}c7 & \symking{}f3\\
\textbf{3.} & \symking{}d6 & \symking{}g4\\
\textbf{4.} & \symking{}e5 & \symking{}h5\\
\textbf{5.} & \symking{}f4 & \symknight{}g6+\\
\textbf{6.} & \symking{}g3 & \symknight{}f1+\\
\textbf{7.} & \symking{}h3 & \symknight{}f4\#\\
\end{tabular}\par\vspace{3mm}
\qrcode[height=12mm]{https://helpman.komtera.lt/?fen=kR5N/q7/8/8/8/8/7N/3K4&moves=7}\\{\sffamily\tiny\color{ink}scannen zum Lösen}
\end{minipage}
\vspace{2mm}

{\raggedright\noindent\textcolor{ink}{\sffamily\textsc{Themen}}\enspace \themetag{reines Matt} \themetag{Mustermatt} \themetag{Spiegelmatt} \themetag{ein Stein zieht}\par\smallskip{\footnotesize\textit{Hinweis:} die Partei am Zug steht im Diagramm im Schach; die Lösung beginnt mit einem Schlag.}\par}

\clearpage
\phantomsection\label{hm:KRNPPvk}
\begin{tcolorbox}[enhanced,colback=ink,colframe=ink,arc=0pt,boxrule=0pt,left=4mm,right=4mm,top=2.5mm,bottom=2.5mm]
\color{white}\sffamily{\LARGE\bfseries KRNPPvk}\hfill{\Large h\#7}\\[0.5mm]
{\small 6 Steine\quad Schwarz am Zug}\hfill{\small Nr.~573\qquad\href{https://helpman.komtera.lt/?fen=8/8/8/8/4N3/K3P3/4P3/kR6&moves=7}{\color{white}\bfseries lösen\,\ensuremath{\nearrow}}}
\end{tcolorbox}
\vspace{2mm}
\noindent\begin{minipage}[t]{36mm}\vspace{0pt}\sffamily\small\raggedright
\hmfigure{längste eindeutige}{{\Large h\#7}}\hmfigure{eindeutig bei dieser Länge}{{\Large 11} Stellungen}\hmfigure{längstes Matt}{{\Large h\#8}}\hmfigure{in der Klasse}{34,1 Milliarden Stellungen,\\ 11,5 Milliarden mit Hilfsmatt}
\end{minipage}\hfill
\begin{minipage}[t]{80mm}\vspace{0pt}\begin{tikzpicture}\begin{axis}[width=72mm,height=42mm,ybar,bar width=1.6pt,xmin=-1,xmax=17,ymin=0,ymax=3.834,axis y line*=left,axis x line*=bottom,xlabel={\scriptsize Halbzüge bis Matt},ylabel={\scriptsize Stellungen (Mrd.)},tick label style={font=\tiny},label style={font=\scriptsize},enlarge x limits=false,ylabel shift=-3pt,xlabel shift=-3pt]\addplot[fill=ink!55,draw=none] coordinates {(0,0.009046) (1,0.06537) (2,0.1475) (3,0.7133) (4,1.82) (5,2.468) (6,3.485) (7,1.348) (8,1.217) (9,0.1405) (10,0.07232) (11,0.0008472) (12,0.002957) (13,2.05e-07) (14,0.0006675) (15,3e-09) (16,1.667e-06)};\draw[densely dotted] (axis cs:14,0) -- (axis cs:14,3.834);\end{axis}\begin{axis}[width=72mm,height=42mm,axis y line*=right,axis x line=none,ymode=log,xmin=-1,xmax=17,ylabel={\scriptsize eindeutig (log)},tick label style={font=\tiny},label style={font=\scriptsize},ylabel shift=-3pt]\addplot[mark=*,mark size=0.7pt,draw=black,thin] coordinates {(0,9046312) (1,44719439) (2,89052207) (3,247385670) (4,437852211) (5,200870361) (6,123577607) (7,6970582) (8,2470427) (9,20081) (10,53129) (12,7844) (14,11)};\end{axis}\end{tikzpicture}\end{minipage}
\vspace{3mm}

\noindent\begin{minipage}[t]{60mm}\vspace{0pt}\sbox\hmboard{\chessboard[setfen=8/8/8/8/4N3/K3P3/4P3/kR6 b - - 0 1]}\usebox\hmboard\\[1mm]
\makebox[\wd\hmboard]{\sffamily\hspace{3mm}\textbf{h\#7}\hfill(5+1)}\end{minipage}\hfill
\begin{minipage}[t]{52mm}\vspace{0pt}
\textcolor{ink}{\sffamily\textsc{Lösung}}\par\smallskip\small
\begin{tabular}{@{}r@{\enspace}l@{\enspace}l@{}}
\textbf{1.} & \symking{}xb1 & \symknight{}d6\\
\textbf{2.} & \symking{}c2 & e4\\
\textbf{3.} & \symking{}c3 & e5\\
\textbf{4.} & \symking{}d4 & e6\\
\textbf{5.} & \symking{}c5 & e7\\
\textbf{6.} & \symking{}b6 & e8=\symqueen{}\\
\textbf{7.} & \symking{}a5 & \symqueen{}b5\#\\
\end{tabular}\par\vspace{3mm}
\qrcode[height=12mm]{https://helpman.komtera.lt/?fen=8/8/8/8/4N3/K3P3/4P3/kR6&moves=7}\\{\sffamily\tiny\color{ink}scannen zum Lösen}
\end{minipage}
\vspace{2mm}

{\raggedright\noindent\textcolor{ink}{\sffamily\textsc{Themen}}\enspace \themetag{Umwandlung} \themetag{ein Stein zieht} \themetag{ohne Schach}\par\smallskip{\footnotesize\textit{Hinweis:} die Partei am Zug steht im Diagramm im Schach; die Lösung beginnt mit einem Schlag.}\par}

\clearpage
\phantomsection\label{hm:KRNPvkb}
\begin{tcolorbox}[enhanced,colback=ink,colframe=ink,arc=0pt,boxrule=0pt,left=4mm,right=4mm,top=2.5mm,bottom=2.5mm]
\color{white}\sffamily{\LARGE\bfseries KRNPvkb}\hfill{\Large h\#7}\\[0.5mm]
{\small 6 Steine\quad Schwarz am Zug}\hfill{\small Nr.~574\qquad\href{https://helpman.komtera.lt/?fen=6bk/7R/8/8/1N6/8/3P4/K7&moves=7}{\color{white}\bfseries lösen\,\ensuremath{\nearrow}}}
\end{tcolorbox}
\vspace{2mm}
\noindent\begin{minipage}[t]{36mm}\vspace{0pt}\sffamily\small\raggedright
\hmfigure{längste eindeutige}{{\Large h\#7}}\hmfigure{eindeutig bei dieser Länge}{{\Large 8} Stellungen}\hmfigure{längstes Matt}{{\Large h\#8}}\hmfigure{in der Klasse}{45,4 Milliarden Stellungen,\\ 28,9 Milliarden mit Hilfsmatt}
\end{minipage}\hfill
\begin{minipage}[t]{80mm}\vspace{0pt}\begin{tikzpicture}\begin{axis}[width=72mm,height=42mm,ybar,bar width=1.6pt,xmin=-1,xmax=17,ymin=0,ymax=10.05,axis y line*=left,axis x line*=bottom,xlabel={\scriptsize Halbzüge bis Matt},ylabel={\scriptsize Stellungen (Mrd.)},tick label style={font=\tiny},label style={font=\scriptsize},enlarge x limits=false,ylabel shift=-3pt,xlabel shift=-3pt]\addplot[fill=ink!55,draw=none] coordinates {(0,0.01274) (1,0.08441) (2,0.2745) (3,1.427) (4,3.763) (5,6.569) (6,9.139) (7,4.641) (8,2.603) (9,0.3078) (10,0.04752) (11,0.0005906) (12,0.002998) (13,8e-09) (14,0.0004196) (15,8e-09) (16,2.144e-06)};\draw[densely dotted] (axis cs:14,0) -- (axis cs:14,10.05);\end{axis}\begin{axis}[width=72mm,height=42mm,axis y line*=right,axis x line=none,ymode=log,xmin=-1,xmax=17,ylabel={\scriptsize eindeutig (log)},tick label style={font=\tiny},label style={font=\scriptsize},ylabel shift=-3pt]\addplot[mark=*,mark size=0.7pt,draw=black,thin] coordinates {(0,12742921) (1,68301897) (2,124774038) (3,393186363) (4,790974774) (5,488647095) (6,229666259) (7,17191284) (8,1851276) (9,16428) (10,36830) (12,8084) (14,8)};\end{axis}\end{tikzpicture}\end{minipage}
\vspace{3mm}

\noindent\begin{minipage}[t]{60mm}\vspace{0pt}\sbox\hmboard{\chessboard[setfen=6bk/7R/8/8/1N6/8/3P4/K7 b - - 0 1]}\usebox\hmboard\\[1mm]
\makebox[\wd\hmboard]{\sffamily\hspace{3mm}\textbf{h\#7}\hfill(4+2)}\end{minipage}\hfill
\begin{minipage}[t]{52mm}\vspace{0pt}
\textcolor{ink}{\sffamily\textsc{Lösung}}\par\smallskip\small
\begin{tabular}{@{}r@{\enspace}l@{\enspace}l@{}}
\textbf{1.} & \symbishop{}xh7 & d4\\
\textbf{2.} & \symking{}g7 & d5\\
\textbf{3.} & \symking{}f6 & d6\\
\textbf{4.} & \symking{}e5 & d7\\
\textbf{5.} & \symking{}d4 & d8=\symqueen{}+\\
\textbf{6.} & \symking{}c3 & \symqueen{}d5\\
\textbf{7.} & \symbishop{}c2 & \symknight{}a2\#\\
\end{tabular}\par\vspace{3mm}
\qrcode[height=12mm]{https://helpman.komtera.lt/?fen=6bk/7R/8/8/1N6/8/3P4/K7&moves=7}\\{\sffamily\tiny\color{ink}scannen zum Lösen}
\end{minipage}
\vspace{2mm}

{\raggedright\noindent\textcolor{ink}{\sffamily\textsc{Themen}}\enspace \themetag{reines Matt} \themetag{Mustermatt} \themetag{Idealmatt} \themetag{Umwandlung} \themetag{Excelsior} \themetag{Rundlauf} \themetag{Selbstblock}\par\smallskip{\footnotesize\textit{Hinweis:} die Partei am Zug steht im Diagramm im Schach; die Lösung beginnt mit einem Schlag.}\par}

\clearpage
\phantomsection\label{hm:KRNPvkn}
\begin{tcolorbox}[enhanced,colback=ink,colframe=ink,arc=0pt,boxrule=0pt,left=4mm,right=4mm,top=2.5mm,bottom=2.5mm]
\color{white}\sffamily{\LARGE\bfseries KRNPvkn}\hfill{\Large h\#7}\\[0.5mm]
{\small 6 Steine\quad Schwarz am Zug}\hfill{\small Nr.~575\qquad\href{https://helpman.komtera.lt/?fen=8/2N5/8/8/8/K7/P6R/6nk&moves=7}{\color{white}\bfseries lösen\,\ensuremath{\nearrow}}}
\end{tcolorbox}
\vspace{2mm}
\noindent\begin{minipage}[t]{36mm}\vspace{0pt}\sffamily\small\raggedright
\hmfigure{längste eindeutige}{{\Large h\#7}}\hmfigure{eindeutig bei dieser Länge}{{\Large 15} Stellungen}\hmfigure{längstes Matt}{{\Large h\#9}}\hmfigure{in der Klasse}{45,4 Milliarden Stellungen,\\ 29,7 Milliarden mit Hilfsmatt}
\end{minipage}\hfill
\begin{minipage}[t]{80mm}\vspace{0pt}\begin{tikzpicture}\begin{axis}[width=72mm,height=42mm,ybar,bar width=1.6pt,xmin=-1,xmax=19,ymin=0,ymax=10.01,axis y line*=left,axis x line*=bottom,xlabel={\scriptsize Halbzüge bis Matt},ylabel={\scriptsize Stellungen (Mrd.)},tick label style={font=\tiny},label style={font=\scriptsize},enlarge x limits=false,ylabel shift=-3pt,xlabel shift=-3pt]\addplot[fill=ink!55,draw=none] coordinates {(0,0.01615) (1,0.108) (2,0.3057) (3,1.463) (4,3.657) (5,6.105) (6,9.098) (7,4.86) (8,3.468) (9,0.4925) (10,0.1045) (11,0.001948) (12,0.003279) (13,1.53e-07) (14,0.0004627) (15,1.28e-07) (16,3.272e-06) (18,3e-09)};\draw[densely dotted] (axis cs:14,0) -- (axis cs:14,10.01);\end{axis}\begin{axis}[width=72mm,height=42mm,axis y line*=right,axis x line=none,ymode=log,xmin=-1,xmax=19,ylabel={\scriptsize eindeutig (log)},tick label style={font=\tiny},label style={font=\scriptsize},ylabel shift=-3pt]\addplot[mark=*,mark size=0.7pt,draw=black,thin] coordinates {(0,16153289) (1,86604304) (2,141525008) (3,398338459) (4,746488588) (5,461202771) (6,265915534) (7,23082856) (8,3415752) (9,50121) (10,31790) (12,9294) (13,2) (14,15)};\end{axis}\end{tikzpicture}\end{minipage}
\vspace{3mm}

\noindent\begin{minipage}[t]{60mm}\vspace{0pt}\sbox\hmboard{\chessboard[setfen=8/2N5/8/8/8/K7/P6R/6nk b - - 0 1]}\usebox\hmboard\\[1mm]
\makebox[\wd\hmboard]{\sffamily\hspace{3mm}\textbf{h\#7}\hfill(4+2)}\end{minipage}\hfill
\begin{minipage}[t]{52mm}\vspace{0pt}
\textcolor{ink}{\sffamily\textsc{Lösung}}\par\smallskip\small
\begin{tabular}{@{}r@{\enspace}l@{\enspace}l@{}}
\textbf{1.} & \symking{}xh2 & \symking{}b4\\
\textbf{2.} & \symking{}g3 & a4\\
\textbf{3.} & \symking{}f4 & a5\\
\textbf{4.} & \symking{}e5 & a6\\
\textbf{5.} & \symking{}d6 & a7\\
\textbf{6.} & \symking{}c6 & a8=\symqueen{}+\\
\textbf{7.} & \symking{}b6 & \symknight{}d5\#\\
\end{tabular}\par\vspace{3mm}
\qrcode[height=12mm]{https://helpman.komtera.lt/?fen=8/2N5/8/8/8/K7/P6R/6nk&moves=7}\\{\sffamily\tiny\color{ink}scannen zum Lösen}
\end{minipage}
\vspace{2mm}

{\raggedright\noindent\textcolor{ink}{\sffamily\textsc{Themen}}\enspace \themetag{Spiegelmatt} \themetag{Umwandlung} \themetag{Excelsior} \themetag{ein Stein zieht}\par\smallskip{\footnotesize\textit{Hinweis:} die Partei am Zug steht im Diagramm im Schach; die Lösung beginnt mit einem Schlag.}\par}

\clearpage
\phantomsection\label{hm:KRNvkbn}
\begin{tcolorbox}[enhanced,colback=ink,colframe=ink,arc=0pt,boxrule=0pt,left=4mm,right=4mm,top=2.5mm,bottom=2.5mm]
\color{white}\sffamily{\LARGE\bfseries KRNvkbn}\hfill{\Large h\#7}\\[0.5mm]
{\small 6 Steine\quad Schwarz am Zug}\hfill{\small Nr.~576\qquad\href{https://helpman.komtera.lt/?fen=8/8/8/8/5NRk/8/n7/Kb6&moves=7}{\color{white}\bfseries lösen\,\ensuremath{\nearrow}}}
\end{tcolorbox}
\vspace{2mm}
\noindent\begin{minipage}[t]{36mm}\vspace{0pt}\sffamily\small\raggedright
\hmfigure{längste eindeutige}{{\Large h\#7}}\hmfigure{eindeutig bei dieser Länge}{{\Large 39} Stellungen}\hmfigure{längstes Matt}{{\Large h\#9.5}}\hmfigure{in der Klasse}{15,5 Milliarden Stellungen,\\ 9,3 Milliarden mit Hilfsmatt}
\end{minipage}\hfill
\begin{minipage}[t]{80mm}\vspace{0pt}\begin{tikzpicture}\begin{axis}[width=72mm,height=42mm,ybar,bar width=1.6pt,xmin=-1,xmax=20,ymin=0,ymax=3.233,axis y line*=left,axis x line*=bottom,xlabel={\scriptsize Halbzüge bis Matt},ylabel={\scriptsize Stellungen (Mrd.)},tick label style={font=\tiny},label style={font=\scriptsize},enlarge x limits=false,ylabel shift=-3pt,xlabel shift=-3pt]\addplot[fill=ink!55,draw=none] coordinates {(0,0.002868) (1,0.01672) (2,0.05684) (3,0.2877) (4,0.8007) (5,1.9) (6,2.939) (7,2.083) (8,1.028) (9,0.1684) (10,0.01561) (11,0.0003554) (12,0.001527) (13,1.94e-07) (14,0.0005133) (15,2.5e-08) (16,5.443e-05) (17,3.8e-08) (18,2.731e-06) (19,7e-09)};\draw[densely dotted] (axis cs:14,0) -- (axis cs:14,3.233);\end{axis}\begin{axis}[width=72mm,height=42mm,axis y line*=right,axis x line=none,ymode=log,xmin=-1,xmax=20,ylabel={\scriptsize eindeutig (log)},tick label style={font=\tiny},label style={font=\scriptsize},ylabel shift=-3pt]\addplot[mark=*,mark size=0.7pt,draw=black,thin] coordinates {(0,2868260) (1,16444750) (2,29306816) (3,100825090) (4,222194712) (5,179135329) (6,80233737) (7,9125020) (8,581237) (9,10000) (10,6072) (12,1043) (14,39)};\end{axis}\end{tikzpicture}\end{minipage}
\vspace{3mm}

\noindent\begin{minipage}[t]{60mm}\vspace{0pt}\sbox\hmboard{\chessboard[setfen=8/8/8/8/5NRk/8/n7/Kb6 b - - 0 1]}\usebox\hmboard\\[1mm]
\makebox[\wd\hmboard]{\sffamily\hspace{3mm}\textbf{h\#7}\hfill(3+3)}\end{minipage}\hfill
\begin{minipage}[t]{52mm}\vspace{0pt}
\textcolor{ink}{\sffamily\textsc{Lösung}}\par\smallskip\small
\begin{tabular}{@{}r@{\enspace}l@{\enspace}l@{}}
\textbf{1.} & \symking{}xg4 & \symking{}b2\\
\textbf{2.} & \symking{}f3 & \symking{}b3\\
\textbf{3.} & \symking{}e3 & \symking{}c4\\
\textbf{4.} & \symking{}d2 & \symking{}d4\\
\textbf{5.} & \symking{}c1 & \symking{}e3\\
\textbf{6.} & \symknight{}b4 & \symking{}e2\\
\textbf{7.} & \symknight{}c2 & \symknight{}d3\#\\
\end{tabular}\par\vspace{3mm}
\qrcode[height=12mm]{https://helpman.komtera.lt/?fen=8/8/8/8/5NRk/8/n7/Kb6&moves=7}\\{\sffamily\tiny\color{ink}scannen zum Lösen}
\end{minipage}
\vspace{2mm}

{\raggedright\noindent\textcolor{ink}{\sffamily\textsc{Themen}}\enspace \themetag{reines Matt} \themetag{Mustermatt} \themetag{Idealmatt} \themetag{Selbstblock} \themetag{ohne Schach}\par\smallskip{\footnotesize\textit{Hinweis:} die Partei am Zug steht im Diagramm im Schach; die Lösung beginnt mit einem Schlag.}\par}

\clearpage
\phantomsection\label{hm:KRNvknn}
\begin{tcolorbox}[enhanced,colback=ink,colframe=ink,arc=0pt,boxrule=0pt,left=4mm,right=4mm,top=2.5mm,bottom=2.5mm]
\color{white}\sffamily{\LARGE\bfseries KRNvknn}\hfill{\Large h\#7}\\[0.5mm]
{\small 6 Steine\quad Schwarz am Zug}\hfill{\small Nr.~577\qquad\href{https://helpman.komtera.lt/?fen=8/8/7R/3n3k/4N1n1/8/8/K7&moves=7}{\color{white}\bfseries lösen\,\ensuremath{\nearrow}}}
\end{tcolorbox}
\vspace{2mm}
\noindent\begin{minipage}[t]{36mm}\vspace{0pt}\sffamily\small\raggedright
\hmfigure{längste eindeutige}{{\Large h\#7}}\hmfigure{eindeutig bei dieser Länge}{{\Large 13} Stellungen}\hmfigure{längstes Matt}{{\Large h\#9}}\hmfigure{in der Klasse}{15,5 Milliarden Stellungen,\\ 4,77 Milliarden mit Hilfsmatt}
\end{minipage}\hfill
\begin{minipage}[t]{80mm}\vspace{0pt}\begin{tikzpicture}\begin{axis}[width=72mm,height=42mm,ybar,bar width=1.6pt,xmin=-1,xmax=19,ymin=0,ymax=1.581,axis y line*=left,axis x line*=bottom,xlabel={\scriptsize Halbzüge bis Matt},ylabel={\scriptsize Stellungen (Mrd.)},tick label style={font=\tiny},label style={font=\scriptsize},enlarge x limits=false,ylabel shift=-3pt,xlabel shift=-3pt]\addplot[fill=ink!55,draw=none] coordinates {(0,0.00183) (1,0.01077) (2,0.03176) (3,0.1455) (4,0.3822) (5,0.8629) (6,1.437) (7,1.081) (8,0.6724) (9,0.1272) (10,0.01756) (11,0.0004783) (12,0.0007374) (13,1e-07) (14,0.0003389) (15,5e-08) (16,4.722e-05) (17,1.6e-08) (18,2.11e-07)};\draw[densely dotted] (axis cs:14,0) -- (axis cs:14,1.581);\end{axis}\begin{axis}[width=72mm,height=42mm,axis y line*=right,axis x line=none,ymode=log,xmin=-1,xmax=19,ylabel={\scriptsize eindeutig (log)},tick label style={font=\tiny},label style={font=\scriptsize},ylabel shift=-3pt]\addplot[mark=*,mark size=0.7pt,draw=black,thin] coordinates {(0,1830184) (1,10562115) (2,16664839) (3,51143703) (4,103589262) (5,85237541) (6,45733444) (7,6067231) (8,588470) (9,12515) (10,1371) (11,2) (12,269) (14,13)};\end{axis}\end{tikzpicture}\end{minipage}
\vspace{3mm}

\noindent\begin{minipage}[t]{60mm}\vspace{0pt}\sbox\hmboard{\chessboard[setfen=8/8/7R/3n3k/4N1n1/8/8/K7 b - - 0 1]}\usebox\hmboard\\[1mm]
\makebox[\wd\hmboard]{\sffamily\hspace{3mm}\textbf{h\#7}\hfill(3+3)}\end{minipage}\hfill
\begin{minipage}[t]{52mm}\vspace{0pt}
\textcolor{ink}{\sffamily\textsc{Lösung}}\par\smallskip\small
\begin{tabular}{@{}r@{\enspace}l@{\enspace}l@{}}
\textbf{1.} & \symknight{}xh6 & \symking{}b2\\
\textbf{2.} & \symknight{}e7 & \symking{}c3\\
\textbf{3.} & \symking{}g6 & \symking{}d4\\
\textbf{4.} & \symking{}f7 & \symking{}e5\\
\textbf{5.} & \symking{}e8 & \symking{}f6\\
\textbf{6.} & \symknight{}f7 & \symking{}g7\\
\textbf{7.} & \symknight{}d8 & \symknight{}f6\#\\
\end{tabular}\par\vspace{3mm}
\qrcode[height=12mm]{https://helpman.komtera.lt/?fen=8/8/7R/3n3k/4N1n1/8/8/K7&moves=7}\\{\sffamily\tiny\color{ink}scannen zum Lösen}
\end{minipage}
\vspace{2mm}

{\raggedright\noindent\textcolor{ink}{\sffamily\textsc{Themen}}\enspace \themetag{reines Matt} \themetag{Mustermatt} \themetag{Idealmatt} \themetag{Selbstblock} \themetag{ohne Schach}\par\smallskip{\footnotesize\textit{Hinweis:} die Partei am Zug steht im Diagramm im Schach; die Lösung beginnt mit einem Schlag.}\par}

\clearpage
\phantomsection\label{hm:KRNvkqn}
\begin{tcolorbox}[enhanced,colback=ink,colframe=ink,arc=0pt,boxrule=0pt,left=4mm,right=4mm,top=2.5mm,bottom=2.5mm]
\color{white}\sffamily{\LARGE\bfseries KRNvkqn}\hfill{\Large h\#7}\\[0.5mm]
{\small 6 Steine\quad Schwarz am Zug}\hfill{\small Nr.~578\qquad\href{https://helpman.komtera.lt/?fen=1n6/8/8/8/5NRk/8/4q3/K7&moves=7}{\color{white}\bfseries lösen\,\ensuremath{\nearrow}}}
\end{tcolorbox}
\vspace{2mm}
\noindent\begin{minipage}[t]{36mm}\vspace{0pt}\sffamily\small\raggedright
\hmfigure{längste eindeutige}{{\Large h\#7}}\hmfigure{eindeutig bei dieser Länge}{{\Large 63} Stellungen}\hmfigure{längstes Matt}{{\Large h\#9}}\hmfigure{in der Klasse}{15,5 Milliarden Stellungen,\\ 8,2 Milliarden mit Hilfsmatt}
\end{minipage}\hfill
\begin{minipage}[t]{80mm}\vspace{0pt}\begin{tikzpicture}\begin{axis}[width=72mm,height=42mm,ybar,bar width=1.6pt,xmin=-1,xmax=19,ymin=0,ymax=2.552,axis y line*=left,axis x line*=bottom,xlabel={\scriptsize Halbzüge bis Matt},ylabel={\scriptsize Stellungen (Mrd.)},tick label style={font=\tiny},label style={font=\scriptsize},enlarge x limits=false,ylabel shift=-3pt,xlabel shift=-3pt]\addplot[fill=ink!55,draw=none] coordinates {(0,0.001443) (1,0.008695) (2,0.03917) (3,0.2295) (4,0.5922) (5,1.714) (6,2.32) (7,2.245) (8,0.7878) (9,0.249) (10,0.0121) (11,0.001576) (12,0.001296) (13,1.69e-07) (14,0.0003405) (16,2.832e-05) (18,1.55e-07)};\draw[densely dotted] (axis cs:14,0) -- (axis cs:14,2.552);\end{axis}\begin{axis}[width=72mm,height=42mm,axis y line*=right,axis x line=none,ymode=log,xmin=-1,xmax=19,ylabel={\scriptsize eindeutig (log)},tick label style={font=\tiny},label style={font=\scriptsize},ylabel shift=-3pt]\addplot[mark=*,mark size=0.7pt,draw=black,thin] coordinates {(0,1443091) (1,8579189) (2,19655215) (3,84063115) (4,147483536) (5,155945123) (6,50964316) (7,12961091) (8,424987) (9,54553) (10,4883) (12,590) (14,63)};\end{axis}\end{tikzpicture}\end{minipage}
\vspace{3mm}

\noindent\begin{minipage}[t]{60mm}\vspace{0pt}\sbox\hmboard{\chessboard[setfen=1n6/8/8/8/5NRk/8/4q3/K7 b - - 0 1]}\usebox\hmboard\\[1mm]
\makebox[\wd\hmboard]{\sffamily\hspace{3mm}\textbf{h\#7}\hfill(3+3)}\end{minipage}\hfill
\begin{minipage}[t]{52mm}\vspace{0pt}
\textcolor{ink}{\sffamily\textsc{Lösung}}\par\smallskip\small
\begin{tabular}{@{}r@{\enspace}l@{\enspace}l@{}}
\textbf{1.} & \symqueen{}xg4 & \symking{}b2\\
\textbf{2.} & \symking{}g5 & \symking{}c3\\
\textbf{3.} & \symking{}f6 & \symking{}d4\\
\textbf{4.} & \symking{}e7 & \symking{}e5\\
\textbf{5.} & \symking{}d8 & \symking{}f6\\
\textbf{6.} & \symqueen{}c8 & \symking{}f7\\
\textbf{7.} & \symknight{}d7 & \symknight{}e6\#\\
\end{tabular}\par\vspace{3mm}
\qrcode[height=12mm]{https://helpman.komtera.lt/?fen=1n6/8/8/8/5NRk/8/4q3/K7&moves=7}\\{\sffamily\tiny\color{ink}scannen zum Lösen}
\end{minipage}
\vspace{2mm}

{\raggedright\noindent\textcolor{ink}{\sffamily\textsc{Themen}}\enspace \themetag{reines Matt} \themetag{Mustermatt} \themetag{Idealmatt} \themetag{Selbstblock} \themetag{ohne Schach}\par\smallskip{\footnotesize\textit{Hinweis:} die Partei am Zug steht im Diagramm im Schach; die Lösung beginnt mit einem Schlag.}\par}

\clearpage
\phantomsection\label{hm:KRNvkqr}
\begin{tcolorbox}[enhanced,colback=ink,colframe=ink,arc=0pt,boxrule=0pt,left=4mm,right=4mm,top=2.5mm,bottom=2.5mm]
\color{white}\sffamily{\LARGE\bfseries KRNvkqr}\hfill{\Large h\#7}\\[0.5mm]
{\small 6 Steine\quad Schwarz am Zug}\hfill{\small Nr.~579\qquad\href{https://helpman.komtera.lt/?fen=8/7q/8/5NRk/8/8/1r6/K7&moves=7}{\color{white}\bfseries lösen\,\ensuremath{\nearrow}}}
\end{tcolorbox}
\vspace{2mm}
\noindent\begin{minipage}[t]{36mm}\vspace{0pt}\sffamily\small\raggedright
\hmfigure{längste eindeutige}{{\Large h\#7}}\hmfigure{eindeutig bei dieser Länge}{{\Large 97} Stellungen}\hmfigure{längstes Matt}{{\Large h\#9.5}}\hmfigure{in der Klasse}{15,5 Milliarden Stellungen,\\ 7,62 Milliarden mit Hilfsmatt}
\end{minipage}\hfill
\begin{minipage}[t]{80mm}\vspace{0pt}\begin{tikzpicture}\begin{axis}[width=72mm,height=42mm,ybar,bar width=1.6pt,xmin=-1,xmax=20,ymin=0,ymax=2.428,axis y line*=left,axis x line*=bottom,xlabel={\scriptsize Halbzüge bis Matt},ylabel={\scriptsize Stellungen (Mrd.)},tick label style={font=\tiny},label style={font=\scriptsize},enlarge x limits=false,ylabel shift=-3pt,xlabel shift=-3pt]\addplot[fill=ink!55,draw=none] coordinates {(0,0.0009899) (1,0.006018) (2,0.03179) (3,0.2156) (4,0.6119) (5,1.839) (6,2.117) (7,2.207) (8,0.4468) (9,0.143) (10,0.002727) (11,0.0002317) (12,0.0004709) (13,5.14e-07) (14,0.000123) (15,3.07e-07) (16,8.339e-06) (17,2.6e-07) (18,1.6e-08) (19,6e-09)};\draw[densely dotted] (axis cs:14,0) -- (axis cs:14,2.428);\end{axis}\begin{axis}[width=72mm,height=42mm,axis y line*=right,axis x line=none,ymode=log,xmin=-1,xmax=20,ylabel={\scriptsize eindeutig (log)},tick label style={font=\tiny},label style={font=\scriptsize},ylabel shift=-3pt]\addplot[mark=*,mark size=0.7pt,draw=black,thin] coordinates {(0,989906) (1,5955855) (2,16711827) (3,81022343) (4,154671243) (5,174689376) (6,41478812) (7,10407828) (8,111144) (9,9502) (10,9195) (11,1) (12,1739) (14,97)};\end{axis}\end{tikzpicture}\end{minipage}
\vspace{3mm}

\noindent\begin{minipage}[t]{60mm}\vspace{0pt}\sbox\hmboard{\chessboard[setfen=8/7q/8/5NRk/8/8/1r6/K7 b - - 0 1]}\usebox\hmboard\\[1mm]
\makebox[\wd\hmboard]{\sffamily\hspace{3mm}\textbf{h\#7}\hfill(3+3)}\end{minipage}\hfill
\begin{minipage}[t]{52mm}\vspace{0pt}
\textcolor{ink}{\sffamily\textsc{Lösung}}\par\smallskip\small
\begin{tabular}{@{}r@{\enspace}l@{\enspace}l@{}}
\textbf{1.} & \symking{}xg5 & \symknight{}e7\\
\textbf{2.} & \symrook{}c2 & \symking{}b1\\
\textbf{3.} & \symrook{}h2+ & \symking{}c1\\
\textbf{4.} & \symrook{}h5 & \symking{}d2\\
\textbf{5.} & \symking{}h4 & \symking{}e3\\
\textbf{6.} & \symqueen{}d3+ & \symking{}f4\\
\textbf{7.} & \symqueen{}h3 & \symknight{}g6\#\\
\end{tabular}\par\vspace{3mm}
\qrcode[height=12mm]{https://helpman.komtera.lt/?fen=8/7q/8/5NRk/8/8/1r6/K7&moves=7}\\{\sffamily\tiny\color{ink}scannen zum Lösen}
\end{minipage}
\vspace{2mm}

{\raggedright\noindent\textcolor{ink}{\sffamily\textsc{Themen}}\enspace \themetag{reines Matt} \themetag{Mustermatt} \themetag{Idealmatt} \themetag{Selbstblock}\par\smallskip{\footnotesize\textit{Hinweis:} die Partei am Zug steht im Diagramm im Schach; die Lösung beginnt mit einem Schlag.}\par}

\clearpage
\phantomsection\label{hm:KRNvkrb}
\begin{tcolorbox}[enhanced,colback=ink,colframe=ink,arc=0pt,boxrule=0pt,left=4mm,right=4mm,top=2.5mm,bottom=2.5mm]
\color{white}\sffamily{\LARGE\bfseries KRNvkrb}\hfill{\Large h\#7}\\[0.5mm]
{\small 6 Steine\quad Schwarz am Zug}\hfill{\small Nr.~580\qquad\href{https://helpman.komtera.lt/?fen=8/6b1/7k/4N2R/8/1r6/8/K7&moves=7}{\color{white}\bfseries lösen\,\ensuremath{\nearrow}}}
\end{tcolorbox}
\vspace{2mm}
\noindent\begin{minipage}[t]{36mm}\vspace{0pt}\sffamily\small\raggedright
\hmfigure{längste eindeutige}{{\Large h\#7}}\hmfigure{eindeutig bei dieser Länge}{{\Large 4} Stellungen}\hmfigure{längstes Matt}{{\Large h\#10}}\hmfigure{in der Klasse}{15,5 Milliarden Stellungen,\\ 8,69 Milliarden mit Hilfsmatt}
\end{minipage}\hfill
\begin{minipage}[t]{80mm}\vspace{0pt}\begin{tikzpicture}\begin{axis}[width=72mm,height=42mm,ybar,bar width=1.6pt,xmin=-1,xmax=21,ymin=0,ymax=3.064,axis y line*=left,axis x line*=bottom,xlabel={\scriptsize Halbzüge bis Matt},ylabel={\scriptsize Stellungen (Mrd.)},tick label style={font=\tiny},label style={font=\scriptsize},enlarge x limits=false,ylabel shift=-3pt,xlabel shift=-3pt]\addplot[fill=ink!55,draw=none] coordinates {(0,0.001984) (1,0.01177) (2,0.05156) (3,0.2955) (4,0.9132) (5,2.163) (6,2.785) (7,1.91) (8,0.485) (9,0.07302) (10,0.002614) (11,4.286e-05) (12,0.0006819) (13,7.6e-08) (14,0.0001547) (16,8.644e-06) (17,3e-09) (18,6.5e-08) (20,5e-09)};\draw[densely dotted] (axis cs:14,0) -- (axis cs:14,3.064);\end{axis}\begin{axis}[width=72mm,height=42mm,axis y line*=right,axis x line=none,ymode=log,xmin=-1,xmax=21,ylabel={\scriptsize eindeutig (log)},tick label style={font=\tiny},label style={font=\scriptsize},ylabel shift=-3pt]\addplot[mark=*,mark size=0.7pt,draw=black,thin] coordinates {(0,1983858) (1,11610950) (2,28481668) (3,109959650) (4,245722740) (5,184491385) (6,50013001) (7,5626481) (8,74441) (9,2113) (10,7525) (11,2) (12,870) (14,4)};\end{axis}\end{tikzpicture}\end{minipage}
\vspace{3mm}

\noindent\begin{minipage}[t]{60mm}\vspace{0pt}\sbox\hmboard{\chessboard[setfen=8/6b1/7k/4N2R/8/1r6/8/K7 b - - 0 1]}\usebox\hmboard\\[1mm]
\makebox[\wd\hmboard]{\sffamily\hspace{3mm}\textbf{h\#7}\hfill(3+3)}\end{minipage}\hfill
\begin{minipage}[t]{52mm}\vspace{0pt}
\textcolor{ink}{\sffamily\textsc{Lösung}}\par\smallskip\small
\begin{tabular}{@{}r@{\enspace}l@{\enspace}l@{}}
\textbf{1.} & \symking{}xh5 & \symking{}a2\\
\textbf{2.} & \symrook{}b8 & \symknight{}d7\\
\textbf{3.} & \symrook{}d8 & \symking{}b3\\
\textbf{4.} & \symking{}g6 & \symking{}c4\\
\textbf{5.} & \symking{}f7 & \symking{}d5\\
\textbf{6.} & \symking{}e8 & \symking{}e6\\
\textbf{7.} & \symbishop{}f8 & \symknight{}f6\#\\
\end{tabular}\par\vspace{3mm}
\qrcode[height=12mm]{https://helpman.komtera.lt/?fen=8/6b1/7k/4N2R/8/1r6/8/K7&moves=7}\\{\sffamily\tiny\color{ink}scannen zum Lösen}
\end{minipage}
\vspace{2mm}

{\raggedright\noindent\textcolor{ink}{\sffamily\textsc{Themen}}\enspace \themetag{Selbstblock} \themetag{ohne Schach}\par\smallskip{\footnotesize\textit{Hinweis:} die Partei am Zug steht im Diagramm im Schach; die Lösung beginnt mit einem Schlag.}\par}

\clearpage
\phantomsection\label{hm:KRNvkrn}
\begin{tcolorbox}[enhanced,colback=ink,colframe=ink,arc=0pt,boxrule=0pt,left=4mm,right=4mm,top=2.5mm,bottom=2.5mm]
\color{white}\sffamily{\LARGE\bfseries KRNvkrn}\hfill{\Large h\#7}\\[0.5mm]
{\small 6 Steine\quad Schwarz am Zug}\hfill{\small Nr.~581\qquad\href{https://helpman.komtera.lt/?fen=n7/7N/8/8/6rk/7R/8/K7&moves=7}{\color{white}\bfseries lösen\,\ensuremath{\nearrow}}}
\end{tcolorbox}
\vspace{2mm}
\noindent\begin{minipage}[t]{36mm}\vspace{0pt}\sffamily\small\raggedright
\hmfigure{längste eindeutige}{{\Large h\#7}}\hmfigure{eindeutig bei dieser Länge}{{\Large 20} Stellungen}\hmfigure{längstes Matt}{{\Large h\#9}}\hmfigure{in der Klasse}{15,5 Milliarden Stellungen,\\ 8,91 Milliarden mit Hilfsmatt}
\end{minipage}\hfill
\begin{minipage}[t]{80mm}\vspace{0pt}\begin{tikzpicture}\begin{axis}[width=72mm,height=42mm,ybar,bar width=1.6pt,xmin=-1,xmax=19,ymin=0,ymax=3.187,axis y line*=left,axis x line*=bottom,xlabel={\scriptsize Halbzüge bis Matt},ylabel={\scriptsize Stellungen (Mrd.)},tick label style={font=\tiny},label style={font=\scriptsize},enlarge x limits=false,ylabel shift=-3pt,xlabel shift=-3pt]\addplot[fill=ink!55,draw=none] coordinates {(0,0.002431) (1,0.0147) (2,0.05989) (3,0.3149) (4,0.9638) (5,2.151) (6,2.897) (7,1.898) (8,0.5284) (9,0.0742) (10,0.002672) (11,7.651e-05) (12,0.0009474) (13,1.19e-07) (14,0.0002471) (15,2e-08) (16,2.443e-05) (18,1.57e-07)};\draw[densely dotted] (axis cs:14,0) -- (axis cs:14,3.187);\end{axis}\begin{axis}[width=72mm,height=42mm,axis y line*=right,axis x line=none,ymode=log,xmin=-1,xmax=19,ylabel={\scriptsize eindeutig (log)},tick label style={font=\tiny},label style={font=\scriptsize},ylabel shift=-3pt]\addplot[mark=*,mark size=0.7pt,draw=black,thin] coordinates {(0,2431423) (1,14480602) (2,33513231) (3,118420012) (4,264174911) (5,187619117) (6,55670638) (7,6031686) (8,92846) (9,3308) (10,6361) (12,1046) (14,20)};\end{axis}\end{tikzpicture}\end{minipage}
\vspace{3mm}

\noindent\begin{minipage}[t]{60mm}\vspace{0pt}\sbox\hmboard{\chessboard[setfen=n7/7N/8/8/6rk/7R/8/K7 b - - 0 1]}\usebox\hmboard\\[1mm]
\makebox[\wd\hmboard]{\sffamily\hspace{3mm}\textbf{h\#7}\hfill(3+3)}\end{minipage}\hfill
\begin{minipage}[t]{52mm}\vspace{0pt}
\textcolor{ink}{\sffamily\textsc{Lösung}}\par\smallskip\small
\begin{tabular}{@{}r@{\enspace}l@{\enspace}l@{}}
\textbf{1.} & \symking{}xh3 & \symking{}b2\\
\textbf{2.} & \symrook{}h4 & \symking{}c3\\
\textbf{3.} & \symking{}g4 & \symking{}d4\\
\textbf{4.} & \symknight{}c7 & \symking{}e5\\
\textbf{5.} & \symknight{}e6 & \symking{}f6\\
\textbf{6.} & \symknight{}g5 & \symking{}g7\\
\textbf{7.} & \symking{}h5 & \symknight{}f6\#\\
\end{tabular}\par\vspace{3mm}
\qrcode[height=12mm]{https://helpman.komtera.lt/?fen=n7/7N/8/8/6rk/7R/8/K7&moves=7}\\{\sffamily\tiny\color{ink}scannen zum Lösen}
\end{minipage}
\vspace{2mm}

{\raggedright\noindent\textcolor{ink}{\sffamily\textsc{Themen}}\enspace \themetag{reines Matt} \themetag{Mustermatt} \themetag{Idealmatt} \themetag{Selbstblock} \themetag{ohne Schach}\par\smallskip{\footnotesize\textit{Hinweis:} die Partei am Zug steht im Diagramm im Schach; die Lösung beginnt mit einem Schlag.}\par}

\clearpage
\phantomsection\label{hm:KRNvkrr}
\begin{tcolorbox}[enhanced,colback=ink,colframe=ink,arc=0pt,boxrule=0pt,left=4mm,right=4mm,top=2.5mm,bottom=2.5mm]
\color{white}\sffamily{\LARGE\bfseries KRNvkrr}\hfill{\Large h\#7}\\[0.5mm]
{\small 6 Steine\quad Schwarz am Zug}\hfill{\small Nr.~582\qquad\href{https://helpman.komtera.lt/?fen=8/8/8/5NRk/8/1r6/2r5/K7&moves=7}{\color{white}\bfseries lösen\,\ensuremath{\nearrow}}}
\end{tcolorbox}
\vspace{2mm}
\noindent\begin{minipage}[t]{36mm}\vspace{0pt}\sffamily\small\raggedright
\hmfigure{längste eindeutige}{{\Large h\#7}}\hmfigure{eindeutig bei dieser Länge}{{\Large 102} Stellungen}\hmfigure{längstes Matt}{{\Large h\#9.5}}\hmfigure{in der Klasse}{15,5 Milliarden Stellungen,\\ 4,13 Milliarden mit Hilfsmatt}
\end{minipage}\hfill
\begin{minipage}[t]{80mm}\vspace{0pt}\begin{tikzpicture}\begin{axis}[width=72mm,height=42mm,ybar,bar width=1.6pt,xmin=-1,xmax=20,ymin=0,ymax=1.367,axis y line*=left,axis x line*=bottom,xlabel={\scriptsize Halbzüge bis Matt},ylabel={\scriptsize Stellungen (Mrd.)},tick label style={font=\tiny},label style={font=\scriptsize},enlarge x limits=false,ylabel shift=-3pt,xlabel shift=-3pt]\addplot[fill=ink!55,draw=none] coordinates {(0,0.0008082) (1,0.004993) (2,0.02469) (3,0.1489) (4,0.4462) (5,1.076) (6,1.243) (7,0.9496) (8,0.1944) (9,0.03781) (10,0.001135) (11,1.499e-05) (12,0.000262) (13,4.55e-07) (14,6.84e-05) (15,2.92e-07) (16,5.029e-06) (17,2.45e-07) (18,1.3e-08) (19,6e-09)};\draw[densely dotted] (axis cs:14,0) -- (axis cs:14,1.367);\end{axis}\begin{axis}[width=72mm,height=42mm,axis y line*=right,axis x line=none,ymode=log,xmin=-1,xmax=20,ylabel={\scriptsize eindeutig (log)},tick label style={font=\tiny},label style={font=\scriptsize},ylabel shift=-3pt]\addplot[mark=*,mark size=0.7pt,draw=black,thin] coordinates {(0,808197) (1,4933150) (2,13880347) (3,56531747) (4,105603064) (5,87629607) (6,19737747) (7,2734177) (8,35545) (9,880) (10,5024) (11,1) (12,820) (14,102)};\end{axis}\end{tikzpicture}\end{minipage}
\vspace{3mm}

\noindent\begin{minipage}[t]{60mm}\vspace{0pt}\sbox\hmboard{\chessboard[setfen=8/8/8/5NRk/8/1r6/2r5/K7 b - - 0 1]}\usebox\hmboard\\[1mm]
\makebox[\wd\hmboard]{\sffamily\hspace{3mm}\textbf{h\#7}\hfill(3+3)}\end{minipage}\hfill
\begin{minipage}[t]{52mm}\vspace{0pt}
\textcolor{ink}{\sffamily\textsc{Lösung}}\par\smallskip\small
\begin{tabular}{@{}r@{\enspace}l@{\enspace}l@{}}
\textbf{1.} & \symking{}xg5 & \symknight{}e7\\
\textbf{2.} & \symrook{}h3 & \symking{}b1\\
\textbf{3.} & \symrook{}ch2 & \symking{}c1\\
\textbf{4.} & \symrook{}h1+ & \symking{}d2\\
\textbf{5.} & \symrook{}h5 & \symking{}e3\\
\textbf{6.} & \symking{}h4 & \symking{}f4\\
\textbf{7.} & \symrook{}h3 & \symknight{}g6\#\\
\end{tabular}\par\vspace{3mm}
\qrcode[height=12mm]{https://helpman.komtera.lt/?fen=8/8/8/5NRk/8/1r6/2r5/K7&moves=7}\\{\sffamily\tiny\color{ink}scannen zum Lösen}
\end{minipage}
\vspace{2mm}

{\raggedright\noindent\textcolor{ink}{\sffamily\textsc{Themen}}\enspace \themetag{reines Matt} \themetag{Mustermatt} \themetag{Idealmatt} \themetag{Selbstblock}\par\smallskip{\footnotesize\textit{Hinweis:} die Partei am Zug steht im Diagramm im Schach; die Lösung beginnt mit einem Schlag.}\par}

\clearpage
\phantomsection\label{hm:KRPPPvk}
\begin{tcolorbox}[enhanced,colback=ink,colframe=ink,arc=0pt,boxrule=0pt,left=4mm,right=4mm,top=2.5mm,bottom=2.5mm]
\color{white}\sffamily{\LARGE\bfseries KRPPPvk}\hfill{\Large h\#7}\\[0.5mm]
{\small 6 Steine\quad Schwarz am Zug}\hfill{\small Nr.~583\qquad\href{https://helpman.komtera.lt/?fen=8/8/8/8/K7/PP4P1/6R1/7k&moves=7}{\color{white}\bfseries lösen\,\ensuremath{\nearrow}}}
\end{tcolorbox}
\vspace{2mm}
\noindent\begin{minipage}[t]{36mm}\vspace{0pt}\sffamily\small\raggedright
\hmfigure{längste eindeutige}{{\Large h\#7}}\hmfigure{eindeutig bei dieser Länge}{{\Large 1} Stellung}\hmfigure{längstes Matt}{{\Large h\#8}}\hmfigure{in der Klasse}{25,6 Milliarden Stellungen,\\ 2,91 Milliarden mit Hilfsmatt}
\end{minipage}\hfill
\begin{minipage}[t]{80mm}\vspace{0pt}\begin{tikzpicture}\begin{axis}[width=72mm,height=42mm,ybar,bar width=1.6pt,xmin=-1,xmax=17,ymin=0,ymax=870.5,axis y line*=left,axis x line*=bottom,xlabel={\scriptsize Halbzüge bis Matt},ylabel={\scriptsize Stellungen (Mio.)},tick label style={font=\tiny},label style={font=\scriptsize},enlarge x limits=false,ylabel shift=-3pt,xlabel shift=-3pt]\addplot[fill=ink!55,draw=none] coordinates {(0,1.948) (1,17.37) (2,38.07) (3,199.5) (4,466.8) (5,617.3) (6,791.4) (7,330.7) (8,324.7) (9,71.24) (10,47.79) (11,3.209) (12,2.106) (13,0.007481) (14,0.05373) (16,2.6e-05)};\draw[densely dotted] (axis cs:14,0) -- (axis cs:14,870.5);\end{axis}\begin{axis}[width=72mm,height=42mm,axis y line*=right,axis x line=none,ymode=log,xmin=-1,xmax=17,ylabel={\scriptsize eindeutig (log)},tick label style={font=\tiny},label style={font=\scriptsize},ylabel shift=-3pt]\addplot[mark=*,mark size=0.7pt,draw=black,thin] coordinates {(0,1947888) (1,9964700) (2,18181275) (3,54861470) (4,83009930) (5,33599689) (6,21024233) (7,2640060) (8,1118179) (9,61069) (10,32585) (11,599) (12,682) (14,1)};\end{axis}\end{tikzpicture}\end{minipage}
\vspace{3mm}

\noindent\begin{minipage}[t]{60mm}\vspace{0pt}\sbox\hmboard{\chessboard[setfen=8/8/8/8/K7/PP4P1/6R1/7k b - - 0 1]}\usebox\hmboard\\[1mm]
\makebox[\wd\hmboard]{\sffamily\hspace{3mm}\textbf{h\#7}\hfill(5+1)}\end{minipage}\hfill
\begin{minipage}[t]{52mm}\vspace{0pt}
\textcolor{ink}{\sffamily\textsc{Lösung}}\par\smallskip\small
\begin{tabular}{@{}r@{\enspace}l@{\enspace}l@{}}
\textbf{1.} & \symking{}xg2 & \symking{}b4\\
\textbf{2.} & \symking{}f3 & a4\\
\textbf{3.} & \symking{}e4 & a5\\
\textbf{4.} & \symking{}d5 & a6\\
\textbf{5.} & \symking{}c6 & a7\\
\textbf{6.} & \symking{}b7 & \symking{}c5\\
\textbf{7.} & \symking{}a6 & a8=\symqueen{}\#\\
\end{tabular}\par\vspace{3mm}
\qrcode[height=12mm]{https://helpman.komtera.lt/?fen=8/8/8/8/K7/PP4P1/6R1/7k&moves=7}\\{\sffamily\tiny\color{ink}scannen zum Lösen}
\end{minipage}
\vspace{2mm}

{\raggedright\noindent\textcolor{ink}{\sffamily\textsc{Themen}}\enspace \themetag{reines Matt} \themetag{Mustermatt} \themetag{Spiegelmatt} \themetag{Umwandlung} \themetag{ein Stein zieht} \themetag{ohne Schach}\par\smallskip{\footnotesize\textit{Hinweis:} die Lösung beginnt mit einem Schlag.}\par}

\clearpage
\phantomsection\label{hm:KRPPvkb}
\begin{tcolorbox}[enhanced,colback=ink,colframe=ink,arc=0pt,boxrule=0pt,left=4mm,right=4mm,top=2.5mm,bottom=2.5mm]
\color{white}\sffamily{\LARGE\bfseries KRPPvkb}\hfill{\Large h\#7}\\[0.5mm]
{\small 6 Steine\quad Schwarz am Zug}\hfill{\small Nr.~584\qquad\href{https://helpman.komtera.lt/?fen=6Rk/7b/8/8/8/1P6/1P6/2K5&moves=7}{\color{white}\bfseries lösen\,\ensuremath{\nearrow}}}
\end{tcolorbox}
\vspace{2mm}
\noindent\begin{minipage}[t]{36mm}\vspace{0pt}\sffamily\small\raggedright
\hmfigure{längste eindeutige}{{\Large h\#7}}\hmfigure{eindeutig bei dieser Länge}{{\Large 1} Stellung}\hmfigure{längstes Matt}{{\Large h\#8}}\hmfigure{in der Klasse}{34,1 Milliarden Stellungen,\\ 11,1 Milliarden mit Hilfsmatt}
\end{minipage}\hfill
\begin{minipage}[t]{80mm}\vspace{0pt}\begin{tikzpicture}\begin{axis}[width=72mm,height=42mm,ybar,bar width=1.6pt,xmin=-1,xmax=17,ymin=0,ymax=2.971,axis y line*=left,axis x line*=bottom,xlabel={\scriptsize Halbzüge bis Matt},ylabel={\scriptsize Stellungen (Mrd.)},tick label style={font=\tiny},label style={font=\scriptsize},enlarge x limits=false,ylabel shift=-3pt,xlabel shift=-3pt]\addplot[fill=ink!55,draw=none] coordinates {(0,0.003983) (1,0.03395) (2,0.1055) (3,0.5871) (4,1.352) (5,2.248) (6,2.701) (7,1.708) (8,1.466) (9,0.5358) (10,0.281) (11,0.03472) (12,0.01178) (13,7.121e-05) (14,3.376e-05) (16,2.85e-07)};\draw[densely dotted] (axis cs:14,0) -- (axis cs:14,2.971);\end{axis}\begin{axis}[width=72mm,height=42mm,axis y line*=right,axis x line=none,ymode=log,xmin=-1,xmax=17,ylabel={\scriptsize eindeutig (log)},tick label style={font=\tiny},label style={font=\scriptsize},ylabel shift=-3pt]\addplot[mark=*,mark size=0.7pt,draw=black,thin] coordinates {(0,3982538) (1,22599442) (2,35724548) (3,117747077) (4,187019769) (5,108378060) (6,63269807) (7,15382712) (8,6057674) (9,664853) (10,162737) (11,908) (12,850) (14,1)};\end{axis}\end{tikzpicture}\end{minipage}
\vspace{3mm}

\noindent\begin{minipage}[t]{60mm}\vspace{0pt}\sbox\hmboard{\chessboard[setfen=6Rk/7b/8/8/8/1P6/1P6/2K5 b - - 0 1]}\usebox\hmboard\\[1mm]
\makebox[\wd\hmboard]{\sffamily\hspace{3mm}\textbf{h\#7}\hfill(4+2)}\end{minipage}\hfill
\begin{minipage}[t]{52mm}\vspace{0pt}
\textcolor{ink}{\sffamily\textsc{Lösung}}\par\smallskip\small
\begin{tabular}{@{}r@{\enspace}l@{\enspace}l@{}}
\textbf{1.} & \symking{}xg8 & b4\\
\textbf{2.} & \symking{}f7 & b5\\
\textbf{3.} & \symking{}e6 & b6\\
\textbf{4.} & \symking{}d5 & b7\\
\textbf{5.} & \symking{}c4 & b3+\\
\textbf{6.} & \symking{}xb3 & b8=\symqueen{}+\\
\textbf{7.} & \symking{}a2 & \symqueen{}b2\#\\
\end{tabular}\par\vspace{3mm}
\qrcode[height=12mm]{https://helpman.komtera.lt/?fen=6Rk/7b/8/8/8/1P6/1P6/2K5&moves=7}\\{\sffamily\tiny\color{ink}scannen zum Lösen}
\end{minipage}
\vspace{2mm}

{\raggedright\noindent\textcolor{ink}{\sffamily\textsc{Themen}}\enspace \themetag{Umwandlung} \themetag{ein Stein zieht}\par\smallskip{\footnotesize\textit{Hinweis:} die Partei am Zug steht im Diagramm im Schach; die Lösung beginnt mit einem Schlag.}\par}

\clearpage
\phantomsection\label{hm:KRPPvkn}
\begin{tcolorbox}[enhanced,colback=ink,colframe=ink,arc=0pt,boxrule=0pt,left=4mm,right=4mm,top=2.5mm,bottom=2.5mm]
\color{white}\sffamily{\LARGE\bfseries KRPPvkn}\hfill{\Large h\#7}\\[0.5mm]
{\small 6 Steine\quad Schwarz am Zug}\hfill{\small Nr.~585\qquad\href{https://helpman.komtera.lt/?fen=6Rk/7n/8/8/8/1P6/1P6/2K5&moves=7}{\color{white}\bfseries lösen\,\ensuremath{\nearrow}}}
\end{tcolorbox}
\vspace{2mm}
\noindent\begin{minipage}[t]{36mm}\vspace{0pt}\sffamily\small\raggedright
\hmfigure{längste eindeutige}{{\Large h\#7}}\hmfigure{eindeutig bei dieser Länge}{{\Large 25} Stellungen}\hmfigure{längstes Matt}{{\Large h\#8.5}}\hmfigure{in der Klasse}{34,1 Milliarden Stellungen,\\ 11,4 Milliarden mit Hilfsmatt}
\end{minipage}\hfill
\begin{minipage}[t]{80mm}\vspace{0pt}\begin{tikzpicture}\begin{axis}[width=72mm,height=42mm,ybar,bar width=1.6pt,xmin=-1,xmax=18,ymin=0,ymax=3.08,axis y line*=left,axis x line*=bottom,xlabel={\scriptsize Halbzüge bis Matt},ylabel={\scriptsize Stellungen (Mrd.)},tick label style={font=\tiny},label style={font=\scriptsize},enlarge x limits=false,ylabel shift=-3pt,xlabel shift=-3pt]\addplot[fill=ink!55,draw=none] coordinates {(0,0.005154) (1,0.04386) (2,0.12) (3,0.621) (4,1.386) (5,2.212) (6,2.8) (7,1.688) (8,1.578) (9,0.5461) (10,0.3145) (11,0.03645) (12,0.01347) (13,6.804e-05) (14,5.415e-05) (15,1.6e-08) (16,1.182e-06) (17,1e-09)};\draw[densely dotted] (axis cs:14,0) -- (axis cs:14,3.08);\end{axis}\begin{axis}[width=72mm,height=42mm,axis y line*=right,axis x line=none,ymode=log,xmin=-1,xmax=18,ylabel={\scriptsize eindeutig (log)},tick label style={font=\tiny},label style={font=\scriptsize},ylabel shift=-3pt]\addplot[mark=*,mark size=0.7pt,draw=black,thin] coordinates {(0,5154289) (1,28868592) (2,41867337) (3,123187716) (4,194105186) (5,102771065) (6,63691019) (7,13568089) (8,5933301) (9,599889) (10,173815) (11,487) (12,950) (14,25)};\end{axis}\end{tikzpicture}\end{minipage}
\vspace{3mm}

\noindent\begin{minipage}[t]{60mm}\vspace{0pt}\sbox\hmboard{\chessboard[setfen=6Rk/7n/8/8/8/1P6/1P6/2K5 b - - 0 1]}\usebox\hmboard\\[1mm]
\makebox[\wd\hmboard]{\sffamily\hspace{3mm}\textbf{h\#7}\hfill(4+2)}\end{minipage}\hfill
\begin{minipage}[t]{52mm}\vspace{0pt}
\textcolor{ink}{\sffamily\textsc{Lösung}}\par\smallskip\small
\begin{tabular}{@{}r@{\enspace}l@{\enspace}l@{}}
\textbf{1.} & \symking{}xg8 & b4\\
\textbf{2.} & \symking{}f7 & b5\\
\textbf{3.} & \symking{}e6 & b6\\
\textbf{4.} & \symking{}d5 & b7\\
\textbf{5.} & \symking{}c4 & b3+\\
\textbf{6.} & \symking{}xb3 & b8=\symqueen{}+\\
\textbf{7.} & \symking{}a2 & \symqueen{}b2\#\\
\end{tabular}\par\vspace{3mm}
\qrcode[height=12mm]{https://helpman.komtera.lt/?fen=6Rk/7n/8/8/8/1P6/1P6/2K5&moves=7}\\{\sffamily\tiny\color{ink}scannen zum Lösen}
\end{minipage}
\vspace{2mm}

{\raggedright\noindent\textcolor{ink}{\sffamily\textsc{Themen}}\enspace \themetag{Umwandlung} \themetag{ein Stein zieht}\par\smallskip{\footnotesize\textit{Hinweis:} die Partei am Zug steht im Diagramm im Schach; die Lösung beginnt mit einem Schlag.}\par}

\clearpage
\phantomsection\label{hm:KRPPvkp}
\begin{tcolorbox}[enhanced,colback=ink,colframe=ink,arc=0pt,boxrule=0pt,left=4mm,right=4mm,top=2.5mm,bottom=2.5mm]
\color{white}\sffamily{\LARGE\bfseries KRPPvkp}\hfill{\Large h\#7}\\[0.5mm]
{\small 6 Steine\quad Schwarz am Zug}\hfill{\small Nr.~586\qquad\href{https://helpman.komtera.lt/?fen=8/3p4/3P4/8/8/8/5PR1/K6k&moves=7}{\color{white}\bfseries lösen\,\ensuremath{\nearrow}}}
\end{tcolorbox}
\vspace{2mm}
\noindent\begin{minipage}[t]{36mm}\vspace{0pt}\sffamily\small\raggedright
\hmfigure{längste eindeutige}{{\Large h\#7}}\hmfigure{eindeutig bei dieser Länge}{{\Large 76} Stellungen}\hmfigure{längstes Matt}{{\Large h\#14}}\hmfigure{in der Klasse}{25,6 Milliarden Stellungen,\\ 8,71 Milliarden mit Hilfsmatt}
\end{minipage}\hfill
\begin{minipage}[t]{80mm}\vspace{0pt}\begin{tikzpicture}\begin{axis}[width=72mm,height=42mm,ybar,bar width=1.6pt,xmin=-1,xmax=29,ymin=0,ymax=2.385,axis y line*=left,axis x line*=bottom,xlabel={\scriptsize Halbzüge bis Matt},ylabel={\scriptsize Stellungen (Mrd.)},tick label style={font=\tiny},label style={font=\scriptsize},enlarge x limits=false,ylabel shift=-3pt,xlabel shift=-3pt]\addplot[fill=ink!55,draw=none] coordinates {(0,0.005144) (1,0.04204) (2,0.1029) (3,0.5076) (4,1.127) (5,1.659) (6,2.169) (7,1.197) (8,1.212) (9,0.3904) (10,0.2526) (11,0.02915) (12,0.01557) (13,0.0001494) (14,0.0001942) (15,5.59e-07) (16,1.24e-05) (17,1.3e-08) (18,6.396e-06) (19,3e-09) (20,5.438e-06) (21,1e-09) (22,3.611e-06) (24,9.46e-07) (26,3.6e-07) (28,9.6e-08)};\draw[densely dotted] (axis cs:14,0) -- (axis cs:14,2.385);\end{axis}\begin{axis}[width=72mm,height=42mm,axis y line*=right,axis x line=none,ymode=log,xmin=-1,xmax=29,ylabel={\scriptsize eindeutig (log)},tick label style={font=\tiny},label style={font=\scriptsize},ylabel shift=-3pt]\addplot[mark=*,mark size=0.7pt,draw=black,thin] coordinates {(0,5143949) (1,27479190) (2,47063816) (3,117730685) (4,188638508) (5,85363563) (6,52165714) (7,9552893) (8,4420070) (9,411631) (10,163426) (11,3898) (12,4094) (13,1) (14,76)};\end{axis}\end{tikzpicture}\end{minipage}
\vspace{3mm}

\noindent\begin{minipage}[t]{60mm}\vspace{0pt}\sbox\hmboard{\chessboard[setfen=8/3p4/3P4/8/8/8/5PR1/K6k b - - 0 1]}\usebox\hmboard\\[1mm]
\makebox[\wd\hmboard]{\sffamily\hspace{3mm}\textbf{h\#7}\hfill(4+2)}\end{minipage}\hfill
\begin{minipage}[t]{52mm}\vspace{0pt}
\textcolor{ink}{\sffamily\textsc{Lösung}}\par\smallskip\small
\begin{tabular}{@{}r@{\enspace}l@{\enspace}l@{}}
\textbf{1.} & \symking{}xg2 & f4\\
\textbf{2.} & \symking{}f3 & f5\\
\textbf{3.} & \symking{}e4 & f6\\
\textbf{4.} & \symking{}d5 & f7\\
\textbf{5.} & \symking{}c6 & f8=\symqueen{}\\
\textbf{6.} & \symking{}b7 & \symqueen{}f3+\\
\textbf{7.} & \symking{}c8 & \symqueen{}a8\#\\
\end{tabular}\par\vspace{3mm}
\qrcode[height=12mm]{https://helpman.komtera.lt/?fen=8/3p4/3P4/8/8/8/5PR1/K6k&moves=7}\\{\sffamily\tiny\color{ink}scannen zum Lösen}
\end{minipage}
\vspace{2mm}

{\raggedright\noindent\textcolor{ink}{\sffamily\textsc{Themen}}\enspace \themetag{reines Matt} \themetag{Mustermatt} \themetag{Umwandlung} \themetag{Excelsior} \themetag{ein Stein zieht}\par\smallskip{\footnotesize\textit{Hinweis:} die Lösung beginnt mit einem Schlag.}\par}

\clearpage
\phantomsection\label{hm:KRPPvkr}
\begin{tcolorbox}[enhanced,colback=ink,colframe=ink,arc=0pt,boxrule=0pt,left=4mm,right=4mm,top=2.5mm,bottom=2.5mm]
\color{white}\sffamily{\LARGE\bfseries KRPPvkr}\hfill{\Large h\#7}\\[0.5mm]
{\small 6 Steine\quad Schwarz am Zug}\hfill{\small Nr.~587\qquad\href{https://helpman.komtera.lt/?fen=8/8/8/8/K7/8/kR2P1P1/r7&moves=7}{\color{white}\bfseries lösen\,\ensuremath{\nearrow}}}
\end{tcolorbox}
\vspace{2mm}
\noindent\begin{minipage}[t]{36mm}\vspace{0pt}\sffamily\small\raggedright
\hmfigure{längste eindeutige}{{\Large h\#7}}\hmfigure{eindeutig bei dieser Länge}{{\Large 5} Stellungen}\hmfigure{längstes Matt}{{\Large h\#8.5}}\hmfigure{in der Klasse}{34,1 Milliarden Stellungen,\\ 10,6 Milliarden mit Hilfsmatt}
\end{minipage}\hfill
\begin{minipage}[t]{80mm}\vspace{0pt}\begin{tikzpicture}\begin{axis}[width=72mm,height=42mm,ybar,bar width=1.6pt,xmin=-1,xmax=18,ymin=0,ymax=2.79,axis y line*=left,axis x line*=bottom,xlabel={\scriptsize Halbzüge bis Matt},ylabel={\scriptsize Stellungen (Mrd.)},tick label style={font=\tiny},label style={font=\scriptsize},enlarge x limits=false,ylabel shift=-3pt,xlabel shift=-3pt]\addplot[fill=ink!55,draw=none] coordinates {(0,0.003339) (1,0.02957) (2,0.1072) (3,0.5845) (4,1.393) (5,2.275) (6,2.537) (7,1.749) (8,1.216) (9,0.4797) (10,0.1831) (11,0.02769) (12,0.004912) (13,6.24e-05) (14,1.701e-05) (15,3.18e-07) (16,5.5e-08) (17,1.2e-08)};\draw[densely dotted] (axis cs:14,0) -- (axis cs:14,2.79);\end{axis}\begin{axis}[width=72mm,height=42mm,axis y line*=right,axis x line=none,ymode=log,xmin=-1,xmax=18,ylabel={\scriptsize eindeutig (log)},tick label style={font=\tiny},label style={font=\scriptsize},ylabel shift=-3pt]\addplot[mark=*,mark size=0.7pt,draw=black,thin] coordinates {(0,3339206) (1,20008782) (2,40211048) (3,123868413) (4,194701073) (5,117848754) (6,56869137) (7,15557992) (8,4682704) (9,590403) (10,96032) (11,889) (12,66) (13,2) (14,5)};\end{axis}\end{tikzpicture}\end{minipage}
\vspace{3mm}

\noindent\begin{minipage}[t]{60mm}\vspace{0pt}\sbox\hmboard{\chessboard[setfen=8/8/8/8/K7/8/kR2P1P1/r7 b - - 0 1]}\usebox\hmboard\\[1mm]
\makebox[\wd\hmboard]{\sffamily\hspace{3mm}\textbf{h\#7}\hfill(4+2)}\end{minipage}\hfill
\begin{minipage}[t]{52mm}\vspace{0pt}
\textcolor{ink}{\sffamily\textsc{Lösung}}\par\smallskip\small
\begin{tabular}{@{}r@{\enspace}l@{\enspace}l@{}}
\textbf{1.} & \symking{}xb2+ & \symking{}b5\\
\textbf{2.} & \symking{}c3 & e4\\
\textbf{3.} & \symking{}d4 & e5\\
\textbf{4.} & \symking{}d5 & e6\\
\textbf{5.} & \symking{}d6 & e7\\
\textbf{6.} & \symking{}d7 & \symking{}b6\\
\textbf{7.} & \symking{}c8 & e8=\symqueen{}\#\\
\end{tabular}\par\vspace{3mm}
\qrcode[height=12mm]{https://helpman.komtera.lt/?fen=8/8/8/8/K7/8/kR2P1P1/r7&moves=7}\\{\sffamily\tiny\color{ink}scannen zum Lösen}
\end{minipage}
\vspace{2mm}

{\raggedright\noindent\textcolor{ink}{\sffamily\textsc{Themen}}\enspace \themetag{reines Matt} \themetag{Mustermatt} \themetag{Spiegelmatt} \themetag{Umwandlung} \themetag{Excelsior} \themetag{ein Stein zieht}\par\smallskip{\footnotesize\textit{Hinweis:} die Partei am Zug steht im Diagramm im Schach; die Lösung beginnt mit einem Schlag.}\par}

\clearpage
\phantomsection\label{hm:KRPvkqq}
\begin{tcolorbox}[enhanced,colback=ink,colframe=ink,arc=0pt,boxrule=0pt,left=4mm,right=4mm,top=2.5mm,bottom=2.5mm]
\color{white}\sffamily{\LARGE\bfseries KRPvkqq}\hfill{\Large h\#7}\\[0.5mm]
{\small 6 Steine\quad Schwarz am Zug}\hfill{\small Nr.~588\qquad\href{https://helpman.komtera.lt/?fen=8/8/8/7R/6qk/8/2q4P/K7&moves=7}{\color{white}\bfseries lösen\,\ensuremath{\nearrow}}}
\end{tcolorbox}
\vspace{2mm}
\noindent\begin{minipage}[t]{36mm}\vspace{0pt}\sffamily\small\raggedright
\hmfigure{längste eindeutige}{{\Large h\#7}}\hmfigure{eindeutig bei dieser Länge}{{\Large 157} Stellungen}\hmfigure{längstes Matt}{{\Large h\#9}}\hmfigure{in der Klasse}{45,4 Milliarden Stellungen,\\ 10,9 Milliarden mit Hilfsmatt}
\end{minipage}\hfill
\begin{minipage}[t]{80mm}\vspace{0pt}\begin{tikzpicture}\begin{axis}[width=72mm,height=42mm,ybar,bar width=1.6pt,xmin=-1,xmax=19,ymin=0,ymax=2.987,axis y line*=left,axis x line*=bottom,xlabel={\scriptsize Halbzüge bis Matt},ylabel={\scriptsize Stellungen (Mrd.)},tick label style={font=\tiny},label style={font=\scriptsize},enlarge x limits=false,ylabel shift=-3pt,xlabel shift=-3pt]\addplot[fill=ink!55,draw=none] coordinates {(0,0.000646) (1,0.005523) (2,0.03001) (3,0.2475) (4,0.5468) (5,1.807) (6,1.537) (7,2.715) (8,1.357) (9,1.723) (10,0.5208) (11,0.3808) (12,0.03974) (13,0.009333) (14,0.0002258) (15,4.779e-06) (16,1.255e-05) (17,1.66e-07) (18,7e-09)};\draw[densely dotted] (axis cs:14,0) -- (axis cs:14,2.987);\end{axis}\begin{axis}[width=72mm,height=42mm,axis y line*=right,axis x line=none,ymode=log,xmin=-1,xmax=19,ylabel={\scriptsize eindeutig (log)},tick label style={font=\tiny},label style={font=\scriptsize},ylabel shift=-3pt]\addplot[mark=*,mark size=0.7pt,draw=black,thin] coordinates {(0,646047) (1,4551960) (2,10653377) (3,63320035) (4,67038162) (5,108821134) (6,38313774) (7,40419350) (8,9373399) (9,7220852) (10,715357) (11,170814) (12,851) (13,17) (14,157)};\end{axis}\end{tikzpicture}\end{minipage}
\vspace{3mm}

\noindent\begin{minipage}[t]{60mm}\vspace{0pt}\sbox\hmboard{\chessboard[setfen=8/8/8/7R/6qk/8/2q4P/K7 b - - 0 1]}\usebox\hmboard\\[1mm]
\makebox[\wd\hmboard]{\sffamily\hspace{3mm}\textbf{h\#7}\hfill(3+3)}\end{minipage}\hfill
\begin{minipage}[t]{52mm}\vspace{0pt}
\textcolor{ink}{\sffamily\textsc{Lösung}}\par\smallskip\small
\begin{tabular}{@{}r@{\enspace}l@{\enspace}l@{}}
\textbf{1.} & \symking{}xh5 & h3\\
\textbf{2.} & \symking{}g5 & h4+\\
\textbf{3.} & \symking{}f4 & h5\\
\textbf{4.} & \symking{}e3 & h6\\
\textbf{5.} & \symking{}d2 & h7\\
\textbf{6.} & \symking{}c1 & h8=\symqueen{}\\
\textbf{7.} & \symqueen{}cd1 & \symqueen{}b2\#\\
\end{tabular}\par\vspace{3mm}
\qrcode[height=12mm]{https://helpman.komtera.lt/?fen=8/8/8/7R/6qk/8/2q4P/K7&moves=7}\\{\sffamily\tiny\color{ink}scannen zum Lösen}
\end{minipage}
\vspace{2mm}

{\raggedright\noindent\textcolor{ink}{\sffamily\textsc{Themen}}\enspace \themetag{Umwandlung} \themetag{Excelsior} \themetag{Selbstblock} \themetag{ein Stein zieht}\par\smallskip{\footnotesize\textit{Hinweis:} die Partei am Zug steht im Diagramm im Schach; die Lösung beginnt mit einem Schlag.}\par}

\clearpage
\phantomsection\label{hm:KRPvkqr}
\begin{tcolorbox}[enhanced,colback=ink,colframe=ink,arc=0pt,boxrule=0pt,left=4mm,right=4mm,top=2.5mm,bottom=2.5mm]
\color{white}\sffamily{\LARGE\bfseries KRPvkqr}\hfill{\Large h\#7}\\[0.5mm]
{\small 6 Steine\quad Schwarz am Zug}\hfill{\small Nr.~589\qquad\href{https://helpman.komtera.lt/?fen=4r3/8/8/8/8/4q3/R1K1P3/k7&moves=7}{\color{white}\bfseries lösen\,\ensuremath{\nearrow}}}
\end{tcolorbox}
\vspace{2mm}
\noindent\begin{minipage}[t]{36mm}\vspace{0pt}\sffamily\small\raggedright
\hmfigure{längste eindeutige}{{\Large h\#7}}\hmfigure{eindeutig bei dieser Länge}{{\Large 8} Stellungen}\hmfigure{längstes Matt}{{\Large h\#8.5}}\hmfigure{in der Klasse}{45,4 Milliarden Stellungen,\\ 23,7 Milliarden mit Hilfsmatt}
\end{minipage}\hfill
\begin{minipage}[t]{80mm}\vspace{0pt}\begin{tikzpicture}\begin{axis}[width=72mm,height=42mm,ybar,bar width=1.6pt,xmin=-1,xmax=18,ymin=0,ymax=6.261,axis y line*=left,axis x line*=bottom,xlabel={\scriptsize Halbzüge bis Matt},ylabel={\scriptsize Stellungen (Mrd.)},tick label style={font=\tiny},label style={font=\scriptsize},enlarge x limits=false,ylabel shift=-3pt,xlabel shift=-3pt]\addplot[fill=ink!55,draw=none] coordinates {(0,0.002172) (1,0.01864) (2,0.09361) (3,0.6721) (4,1.527) (5,4.153) (6,3.982) (7,5.692) (8,3.188) (9,3.041) (10,0.8462) (11,0.4274) (12,0.04464) (13,0.00627) (14,0.0002515) (15,4.002e-06) (16,2.05e-06) (17,1.11e-07)};\draw[densely dotted] (axis cs:14,0) -- (axis cs:14,6.261);\end{axis}\begin{axis}[width=72mm,height=42mm,axis y line*=right,axis x line=none,ymode=log,xmin=-1,xmax=18,ylabel={\scriptsize eindeutig (log)},tick label style={font=\tiny},label style={font=\scriptsize},ylabel shift=-3pt]\addplot[mark=*,mark size=0.7pt,draw=black,thin] coordinates {(0,2172031) (1,15208060) (2,35574620) (3,170735221) (4,208894178) (5,253368705) (6,107686509) (7,82913055) (8,19143755) (9,8592620) (10,864258) (11,108112) (12,753) (13,7) (14,8)};\end{axis}\end{tikzpicture}\end{minipage}
\vspace{3mm}

\noindent\begin{minipage}[t]{60mm}\vspace{0pt}\sbox\hmboard{\chessboard[setfen=4r3/8/8/8/8/4q3/R1K1P3/k7 b - - 0 1]}\usebox\hmboard\\[1mm]
\makebox[\wd\hmboard]{\sffamily\hspace{3mm}\textbf{h\#7}\hfill(3+3)}\end{minipage}\hfill
\begin{minipage}[t]{52mm}\vspace{0pt}
\textcolor{ink}{\sffamily\textsc{Lösung}}\par\smallskip\small
\begin{tabular}{@{}r@{\enspace}l@{\enspace}l@{}}
\textbf{1.} & \symking{}xa2 & \symking{}d1\\
\textbf{2.} & \symqueen{}c3 & e4\\
\textbf{3.} & \symking{}b3 & e5\\
\textbf{4.} & \symking{}c4 & e6\\
\textbf{5.} & \symking{}d3 & e7\\
\textbf{6.} & \symrook{}d8 & e8=\symqueen{}\\
\textbf{7.} & \symrook{}d4 & \symqueen{}e2\#\\
\end{tabular}\par\vspace{3mm}
\qrcode[height=12mm]{https://helpman.komtera.lt/?fen=4r3/8/8/8/8/4q3/R1K1P3/k7&moves=7}\\{\sffamily\tiny\color{ink}scannen zum Lösen}
\end{minipage}
\vspace{2mm}

{\raggedright\noindent\textcolor{ink}{\sffamily\textsc{Themen}}\enspace \themetag{Umwandlung} \themetag{Excelsior} \themetag{Rundlauf} \themetag{Selbstblock} \themetag{ohne Schach} \themetag{Umnov}\par\smallskip{\footnotesize\textit{Hinweis:} die Partei am Zug steht im Diagramm im Schach; die Lösung beginnt mit einem Schlag.}\par}

\clearpage
\phantomsection\label{hm:KRPvkrb}
\begin{tcolorbox}[enhanced,colback=ink,colframe=ink,arc=0pt,boxrule=0pt,left=4mm,right=4mm,top=2.5mm,bottom=2.5mm]
\color{white}\sffamily{\LARGE\bfseries KRPvkrb}\hfill{\Large h\#7}\\[0.5mm]
{\small 6 Steine\quad Schwarz am Zug}\hfill{\small Nr.~590\qquad\href{https://helpman.komtera.lt/?fen=6Rk/7b/8/8/8/8/KP1r4/8&moves=7}{\color{white}\bfseries lösen\,\ensuremath{\nearrow}}}
\end{tcolorbox}
\vspace{2mm}
\noindent\begin{minipage}[t]{36mm}\vspace{0pt}\sffamily\small\raggedright
\hmfigure{längste eindeutige}{{\Large h\#7}}\hmfigure{eindeutig bei dieser Länge}{{\Large 7} Stellungen}\hmfigure{längstes Matt}{{\Large h\#9}}\hmfigure{in der Klasse}{45,4 Milliarden Stellungen,\\ 26,9 Milliarden mit Hilfsmatt}
\end{minipage}\hfill
\begin{minipage}[t]{80mm}\vspace{0pt}\begin{tikzpicture}\begin{axis}[width=72mm,height=42mm,ybar,bar width=1.6pt,xmin=-1,xmax=19,ymin=0,ymax=6.137,axis y line*=left,axis x line*=bottom,xlabel={\scriptsize Halbzüge bis Matt},ylabel={\scriptsize Stellungen (Mrd.)},tick label style={font=\tiny},label style={font=\scriptsize},enlarge x limits=false,ylabel shift=-3pt,xlabel shift=-3pt]\addplot[fill=ink!55,draw=none] coordinates {(0,0.004598) (1,0.0374) (2,0.1506) (3,0.9231) (4,2.142) (5,4.66) (6,5.275) (7,5.579) (8,4.104) (9,2.679) (10,1.033) (11,0.2814) (12,0.04938) (13,0.001186) (14,0.0002745) (15,1.3e-08) (16,2.476e-06) (18,3e-09)};\draw[densely dotted] (axis cs:14,0) -- (axis cs:14,6.137);\end{axis}\begin{axis}[width=72mm,height=42mm,axis y line*=right,axis x line=none,ymode=log,xmin=-1,xmax=19,ylabel={\scriptsize eindeutig (log)},tick label style={font=\tiny},label style={font=\scriptsize},ylabel shift=-3pt]\addplot[mark=*,mark size=0.7pt,draw=black,thin] coordinates {(0,4597982) (1,29904339) (2,58409147) (3,218724360) (4,318809014) (5,272847022) (6,144514440) (7,69943067) (8,24071139) (9,5716868) (10,942947) (11,21944) (12,1466) (14,7)};\end{axis}\end{tikzpicture}\end{minipage}
\vspace{3mm}

\noindent\begin{minipage}[t]{60mm}\vspace{0pt}\sbox\hmboard{\chessboard[setfen=6Rk/7b/8/8/8/8/KP1r4/8 b - - 0 1]}\usebox\hmboard\\[1mm]
\makebox[\wd\hmboard]{\sffamily\hspace{3mm}\textbf{h\#7}\hfill(3+3)}\end{minipage}\hfill
\begin{minipage}[t]{52mm}\vspace{0pt}
\textcolor{ink}{\sffamily\textsc{Lösung}}\par\smallskip\small
\begin{tabular}{@{}r@{\enspace}l@{\enspace}l@{}}
\textbf{1.} & \symking{}xg8 & \symking{}a3\\
\textbf{2.} & \symking{}f7 & b4\\
\textbf{3.} & \symking{}e6 & b5\\
\textbf{4.} & \symking{}d5 & b6\\
\textbf{5.} & \symking{}c4 & b7\\
\textbf{6.} & \symrook{}d5 & b8=\symqueen{}\\
\textbf{7.} & \symbishop{}d3 & \symqueen{}b4\#\\
\end{tabular}\par\vspace{3mm}
\qrcode[height=12mm]{https://helpman.komtera.lt/?fen=6Rk/7b/8/8/8/8/KP1r4/8&moves=7}\\{\sffamily\tiny\color{ink}scannen zum Lösen}
\end{minipage}
\vspace{2mm}

{\raggedright\noindent\textcolor{ink}{\sffamily\textsc{Themen}}\enspace \themetag{Umwandlung} \themetag{Excelsior} \themetag{Rundlauf} \themetag{Selbstblock} \themetag{ohne Schach}\par\smallskip{\footnotesize\textit{Hinweis:} die Partei am Zug steht im Diagramm im Schach; die Lösung beginnt mit einem Schlag.}\par}

\clearpage
\phantomsection\label{hm:KRPvkrn}
\begin{tcolorbox}[enhanced,colback=ink,colframe=ink,arc=0pt,boxrule=0pt,left=4mm,right=4mm,top=2.5mm,bottom=2.5mm]
\color{white}\sffamily{\LARGE\bfseries KRPvkrn}\hfill{\Large h\#7}\\[0.5mm]
{\small 6 Steine\quad Schwarz am Zug}\hfill{\small Nr.~591\qquad\href{https://helpman.komtera.lt/?fen=8/8/8/8/8/4r3/4P2R/K5nk&moves=7}{\color{white}\bfseries lösen\,\ensuremath{\nearrow}}}
\end{tcolorbox}
\vspace{2mm}
\noindent\begin{minipage}[t]{36mm}\vspace{0pt}\sffamily\small\raggedright
\hmfigure{längste eindeutige}{{\Large h\#7}}\hmfigure{eindeutig bei dieser Länge}{{\Large 100} Stellungen}\hmfigure{längstes Matt}{{\Large h\#9}}\hmfigure{in der Klasse}{45,4 Milliarden Stellungen,\\ 27,6 Milliarden mit Hilfsmatt}
\end{minipage}\hfill
\begin{minipage}[t]{80mm}\vspace{0pt}\begin{tikzpicture}\begin{axis}[width=72mm,height=42mm,ybar,bar width=1.6pt,xmin=-1,xmax=19,ymin=0,ymax=6.158,axis y line*=left,axis x line*=bottom,xlabel={\scriptsize Halbzüge bis Matt},ylabel={\scriptsize Stellungen (Mrd.)},tick label style={font=\tiny},label style={font=\scriptsize},enlarge x limits=false,ylabel shift=-3pt,xlabel shift=-3pt]\addplot[fill=ink!55,draw=none] coordinates {(0,0.005769) (1,0.04722) (2,0.1783) (3,1.005) (4,2.301) (5,4.72) (6,5.598) (7,5.538) (8,4.271) (9,2.596) (10,0.9996) (11,0.2564) (12,0.03869) (13,0.0007862) (14,0.000342) (15,9.4e-08) (16,8.496e-06) (17,1.5e-08) (18,1.1e-08)};\draw[densely dotted] (axis cs:14,0) -- (axis cs:14,6.158);\end{axis}\begin{axis}[width=72mm,height=42mm,axis y line*=right,axis x line=none,ymode=log,xmin=-1,xmax=19,ylabel={\scriptsize eindeutig (log)},tick label style={font=\tiny},label style={font=\scriptsize},ylabel shift=-3pt]\addplot[mark=*,mark size=0.7pt,draw=black,thin] coordinates {(0,5769323) (1,37587803) (2,70786105) (3,238616214) (4,362654591) (5,282423299) (6,156091100) (7,65300252) (8,22992262) (9,4898291) (10,802263) (11,14835) (12,1683) (14,100)};\end{axis}\end{tikzpicture}\end{minipage}
\vspace{3mm}

\noindent\begin{minipage}[t]{60mm}\vspace{0pt}\sbox\hmboard{\chessboard[setfen=8/8/8/8/8/4r3/4P2R/K5nk b - - 0 1]}\usebox\hmboard\\[1mm]
\makebox[\wd\hmboard]{\sffamily\hspace{3mm}\textbf{h\#7}\hfill(3+3)}\end{minipage}\hfill
\begin{minipage}[t]{52mm}\vspace{0pt}
\textcolor{ink}{\sffamily\textsc{Lösung}}\par\smallskip\small
\begin{tabular}{@{}r@{\enspace}l@{\enspace}l@{}}
\textbf{1.} & \symking{}xh2 & \symking{}b1\\
\textbf{2.} & \symrook{}h3 & e4\\
\textbf{3.} & \symknight{}f3 & e5\\
\textbf{4.} & \symknight{}g5 & e6\\
\textbf{5.} & \symknight{}f7 & exf7\\
\textbf{6.} & \symking{}h1 & f8=\symqueen{}\\
\textbf{7.} & \symrook{}h2 & \symqueen{}f1\#\\
\end{tabular}\par\vspace{3mm}
\qrcode[height=12mm]{https://helpman.komtera.lt/?fen=8/8/8/8/8/4r3/4P2R/K5nk&moves=7}\\{\sffamily\tiny\color{ink}scannen zum Lösen}
\end{minipage}
\vspace{2mm}

{\raggedright\noindent\textcolor{ink}{\sffamily\textsc{Themen}}\enspace \themetag{reines Matt} \themetag{Mustermatt} \themetag{Umwandlung} \themetag{Excelsior} \themetag{Switchback} \themetag{Selbstblock} \themetag{ohne Schach}\par\smallskip{\footnotesize\textit{Hinweis:} die Partei am Zug steht im Diagramm im Schach; die Lösung beginnt mit einem Schlag.}\par}

\clearpage
\phantomsection\label{hm:KRPvkrr}
\begin{tcolorbox}[enhanced,colback=ink,colframe=ink,arc=0pt,boxrule=0pt,left=4mm,right=4mm,top=2.5mm,bottom=2.5mm]
\color{white}\sffamily{\LARGE\bfseries KRPvkrr}\hfill{\Large h\#7}\\[0.5mm]
{\small 6 Steine\quad Schwarz am Zug}\hfill{\small Nr.~592\qquad\href{https://helpman.komtera.lt/?fen=8/8/8/8/K7/8/kR2P3/r1r5&moves=7}{\color{white}\bfseries lösen\,\ensuremath{\nearrow}}}
\end{tcolorbox}
\vspace{2mm}
\noindent\begin{minipage}[t]{36mm}\vspace{0pt}\sffamily\small\raggedright
\hmfigure{längste eindeutige}{{\Large h\#7}}\hmfigure{eindeutig bei dieser Länge}{{\Large 27} Stellungen}\hmfigure{längstes Matt}{{\Large h\#8.5}}\hmfigure{in der Klasse}{45,4 Milliarden Stellungen,\\ 12,8 Milliarden mit Hilfsmatt}
\end{minipage}\hfill
\begin{minipage}[t]{80mm}\vspace{0pt}\begin{tikzpicture}\begin{axis}[width=72mm,height=42mm,ybar,bar width=1.6pt,xmin=-1,xmax=18,ymin=0,ymax=3.089,axis y line*=left,axis x line*=bottom,xlabel={\scriptsize Halbzüge bis Matt},ylabel={\scriptsize Stellungen (Mrd.)},tick label style={font=\tiny},label style={font=\scriptsize},enlarge x limits=false,ylabel shift=-3pt,xlabel shift=-3pt]\addplot[fill=ink!55,draw=none] coordinates {(0,0.001802) (1,0.01538) (2,0.06887) (3,0.4318) (4,1.005) (5,2.259) (6,2.398) (7,2.808) (8,1.808) (9,1.364) (10,0.4449) (11,0.1651) (12,0.02192) (13,0.001553) (14,0.0001159) (15,1.092e-06) (16,9.71e-07) (17,6e-09)};\draw[densely dotted] (axis cs:14,0) -- (axis cs:14,3.089);\end{axis}\begin{axis}[width=72mm,height=42mm,axis y line*=right,axis x line=none,ymode=log,xmin=-1,xmax=18,ylabel={\scriptsize eindeutig (log)},tick label style={font=\tiny},label style={font=\scriptsize},ylabel shift=-3pt]\addplot[mark=*,mark size=0.7pt,draw=black,thin] coordinates {(0,1802057) (1,12428188) (2,27459769) (3,105166325) (4,139633802) (5,126209867) (6,60162267) (7,35269689) (8,9407167) (9,3017409) (10,397896) (11,24436) (12,308) (14,27)};\end{axis}\end{tikzpicture}\end{minipage}
\vspace{3mm}

\noindent\begin{minipage}[t]{60mm}\vspace{0pt}\sbox\hmboard{\chessboard[setfen=8/8/8/8/K7/8/kR2P3/r1r5 b - - 0 1]}\usebox\hmboard\\[1mm]
\makebox[\wd\hmboard]{\sffamily\hspace{3mm}\textbf{h\#7}\hfill(3+3)}\end{minipage}\hfill
\begin{minipage}[t]{52mm}\vspace{0pt}
\textcolor{ink}{\sffamily\textsc{Lösung}}\par\smallskip\small
\begin{tabular}{@{}r@{\enspace}l@{\enspace}l@{}}
\textbf{1.} & \symking{}xb2+ & \symking{}b5\\
\textbf{2.} & \symking{}c3 & e4\\
\textbf{3.} & \symking{}d4 & e5\\
\textbf{4.} & \symking{}d5 & e6\\
\textbf{5.} & \symking{}d6 & e7\\
\textbf{6.} & \symking{}d7 & \symking{}b6\\
\textbf{7.} & \symking{}c8 & e8=\symqueen{}\#\\
\end{tabular}\par\vspace{3mm}
\qrcode[height=12mm]{https://helpman.komtera.lt/?fen=8/8/8/8/K7/8/kR2P3/r1r5&moves=7}\\{\sffamily\tiny\color{ink}scannen zum Lösen}
\end{minipage}
\vspace{2mm}

{\raggedright\noindent\textcolor{ink}{\sffamily\textsc{Themen}}\enspace \themetag{reines Matt} \themetag{Mustermatt} \themetag{Spiegelmatt} \themetag{Umwandlung} \themetag{Excelsior} \themetag{ein Stein zieht}\par\smallskip{\footnotesize\textit{Hinweis:} die Partei am Zug steht im Diagramm im Schach; die Lösung beginnt mit einem Schlag.}\par}

\clearpage
\phantomsection\label{hm:KRvkqbn}
\begin{tcolorbox}[enhanced,colback=ink,colframe=ink,arc=0pt,boxrule=0pt,left=4mm,right=4mm,top=2.5mm,bottom=2.5mm]
\color{white}\sffamily{\LARGE\bfseries KRvkqbn}\hfill{\Large h\#7}\\[0.5mm]
{\small 6 Steine\quad Schwarz am Zug}\hfill{\small Nr.~593\qquad\href{https://helpman.komtera.lt/?fen=8/6Rk/8/8/8/b1n5/2q5/K7&moves=7}{\color{white}\bfseries lösen\,\ensuremath{\nearrow}}}
\end{tcolorbox}
\vspace{2mm}
\noindent\begin{minipage}[t]{36mm}\vspace{0pt}\sffamily\small\raggedright
\hmfigure{längste eindeutige}{{\Large h\#7}}\hmfigure{eindeutig bei dieser Länge}{{\Large 1} Stellung}\hmfigure{längstes Matt}{{\Large h\#8}}\hmfigure{in der Klasse}{15,5 Milliarden Stellungen,\\ 8,08 Milliarden mit Hilfsmatt}
\end{minipage}\hfill
\begin{minipage}[t]{80mm}\vspace{0pt}\begin{tikzpicture}\begin{axis}[width=72mm,height=42mm,ybar,bar width=1.6pt,xmin=-1,xmax=17,ymin=0,ymax=1.852,axis y line*=left,axis x line*=bottom,xlabel={\scriptsize Halbzüge bis Matt},ylabel={\scriptsize Stellungen (Mrd.)},tick label style={font=\tiny},label style={font=\scriptsize},enlarge x limits=false,ylabel shift=-3pt,xlabel shift=-3pt]\addplot[fill=ink!55,draw=none] coordinates {(0,0.0005665) (1,0.004115) (2,0.01951) (3,0.1259) (4,0.2284) (5,0.7789) (6,0.8476) (7,1.682) (8,1.278) (9,1.683) (10,0.7729) (11,0.5319) (12,0.1081) (13,0.01646) (14,0.001165) (15,4.411e-05) (16,2.7e-08)};\draw[densely dotted] (axis cs:14,0) -- (axis cs:14,1.852);\end{axis}\begin{axis}[width=72mm,height=42mm,axis y line*=right,axis x line=none,ymode=log,xmin=-1,xmax=17,ylabel={\scriptsize eindeutig (log)},tick label style={font=\tiny},label style={font=\scriptsize},ylabel shift=-3pt]\addplot[mark=*,mark size=0.7pt,draw=black,thin] coordinates {(0,566454) (1,4114532) (2,8379550) (3,40058111) (4,38830077) (5,72197751) (6,34373960) (7,33598368) (8,12960275) (9,8824764) (10,1635377) (11,214752) (12,5049) (13,213) (14,1)};\end{axis}\end{tikzpicture}\end{minipage}
\vspace{3mm}

\noindent\begin{minipage}[t]{60mm}\vspace{0pt}\sbox\hmboard{\chessboard[setfen=8/6Rk/8/8/8/b1n5/2q5/K7 b - - 0 1]}\usebox\hmboard\\[1mm]
\makebox[\wd\hmboard]{\sffamily\hspace{3mm}\textbf{h\#7}\hfill(2+4)}\end{minipage}\hfill
\begin{minipage}[t]{52mm}\vspace{0pt}
\textcolor{ink}{\sffamily\textsc{Lösung}}\par\smallskip\small
\begin{tabular}{@{}r@{\enspace}l@{\enspace}l@{}}
\textbf{1.} & \symking{}h6 & \symrook{}g2\\
\textbf{2.} & \symking{}h5 & \symrook{}xc2\\
\textbf{3.} & \symking{}g4 & \symrook{}a2\\
\textbf{4.} & \symking{}f3 & \symrook{}xa3\\
\textbf{5.} & \symking{}e2 & \symking{}b2\\
\textbf{6.} & \symking{}d1 & \symking{}xc3\\
\textbf{7.} & \symking{}c1 & \symrook{}a1\#\\
\end{tabular}\par\vspace{3mm}
\qrcode[height=12mm]{https://helpman.komtera.lt/?fen=8/6Rk/8/8/8/b1n5/2q5/K7&moves=7}\\{\sffamily\tiny\color{ink}scannen zum Lösen}
\end{minipage}
\vspace{2mm}

{\raggedright\noindent\textcolor{ink}{\sffamily\textsc{Themen}}\enspace \themetag{reines Matt} \themetag{Mustermatt} \themetag{Idealmatt} \themetag{Spiegelmatt} \themetag{ein Stein zieht} \themetag{ohne Schach}\par\smallskip{\footnotesize\textit{Hinweis:} die Partei am Zug steht im Diagramm im Schach.}\par}

\clearpage
\phantomsection\label{hm:KRvkqnn}
\begin{tcolorbox}[enhanced,colback=ink,colframe=ink,arc=0pt,boxrule=0pt,left=4mm,right=4mm,top=2.5mm,bottom=2.5mm]
\color{white}\sffamily{\LARGE\bfseries KRvkqnn}\hfill{\Large h\#7}\\[0.5mm]
{\small 6 Steine\quad Schwarz am Zug}\hfill{\small Nr.~594\qquad\href{https://helpman.komtera.lt/?fen=8/6Rk/8/8/n7/2n5/2q5/K7&moves=7}{\color{white}\bfseries lösen\,\ensuremath{\nearrow}}}
\end{tcolorbox}
\vspace{2mm}
\noindent\begin{minipage}[t]{36mm}\vspace{0pt}\sffamily\small\raggedright
\hmfigure{längste eindeutige}{{\Large h\#7}}\hmfigure{eindeutig bei dieser Länge}{{\Large 1} Stellung}\hmfigure{längstes Matt}{{\Large h\#8}}\hmfigure{in der Klasse}{15,5 Milliarden Stellungen,\\ 4,14 Milliarden mit Hilfsmatt}
\end{minipage}\hfill
\begin{minipage}[t]{80mm}\vspace{0pt}\begin{tikzpicture}\begin{axis}[width=72mm,height=42mm,ybar,bar width=1.6pt,xmin=-1,xmax=17,ymin=0,ymax=932.2,axis y line*=left,axis x line*=bottom,xlabel={\scriptsize Halbzüge bis Matt},ylabel={\scriptsize Stellungen (Mio.)},tick label style={font=\tiny},label style={font=\scriptsize},enlarge x limits=false,ylabel shift=-3pt,xlabel shift=-3pt]\addplot[fill=ink!55,draw=none] coordinates {(0,0.3628) (1,2.657) (2,12.25) (3,72.34) (4,129.8) (5,410.4) (6,463.3) (7,847.4) (8,671.7) (9,826.3) (10,393.9) (11,249.1) (12,54.79) (13,7.152) (14,0.5417) (15,0.01552) (16,7.6e-05)};\draw[densely dotted] (axis cs:14,0) -- (axis cs:14,932.2);\end{axis}\begin{axis}[width=72mm,height=42mm,axis y line*=right,axis x line=none,ymode=log,xmin=-1,xmax=17,ylabel={\scriptsize eindeutig (log)},tick label style={font=\tiny},label style={font=\scriptsize},ylabel shift=-3pt]\addplot[mark=*,mark size=0.7pt,draw=black,thin] coordinates {(0,362806) (1,2657442) (2,5401324) (3,23825192) (4,23565518) (5,39944492) (6,20471782) (7,17392480) (8,6933935) (9,4074823) (10,804487) (11,75702) (12,4027) (13,105) (14,1)};\end{axis}\end{tikzpicture}\end{minipage}
\vspace{3mm}

\noindent\begin{minipage}[t]{60mm}\vspace{0pt}\sbox\hmboard{\chessboard[setfen=8/6Rk/8/8/n7/2n5/2q5/K7 b - - 0 1]}\usebox\hmboard\\[1mm]
\makebox[\wd\hmboard]{\sffamily\hspace{3mm}\textbf{h\#7}\hfill(2+4)}\end{minipage}\hfill
\begin{minipage}[t]{52mm}\vspace{0pt}
\textcolor{ink}{\sffamily\textsc{Lösung}}\par\smallskip\small
\begin{tabular}{@{}r@{\enspace}l@{\enspace}l@{}}
\textbf{1.} & \symking{}h6 & \symrook{}g2\\
\textbf{2.} & \symking{}h5 & \symrook{}xc2\\
\textbf{3.} & \symking{}g4 & \symrook{}a2\\
\textbf{4.} & \symking{}f3 & \symrook{}xa4\\
\textbf{5.} & \symking{}e2 & \symking{}b2\\
\textbf{6.} & \symking{}d1 & \symking{}xc3\\
\textbf{7.} & \symking{}c1 & \symrook{}a1\#\\
\end{tabular}\par\vspace{3mm}
\qrcode[height=12mm]{https://helpman.komtera.lt/?fen=8/6Rk/8/8/n7/2n5/2q5/K7&moves=7}\\{\sffamily\tiny\color{ink}scannen zum Lösen}
\end{minipage}
\vspace{2mm}

{\raggedright\noindent\textcolor{ink}{\sffamily\textsc{Themen}}\enspace \themetag{reines Matt} \themetag{Mustermatt} \themetag{Idealmatt} \themetag{Spiegelmatt} \themetag{ein Stein zieht} \themetag{ohne Schach}\par\smallskip{\footnotesize\textit{Hinweis:} die Partei am Zug steht im Diagramm im Schach.}\par}

\clearpage
\phantomsection\label{hm:KRvkqnp}
\begin{tcolorbox}[enhanced,colback=ink,colframe=ink,arc=0pt,boxrule=0pt,left=4mm,right=4mm,top=2.5mm,bottom=2.5mm]
\color{white}\sffamily{\LARGE\bfseries KRvkqnp}\hfill{\Large h\#7}\\[0.5mm]
{\small 6 Steine\quad Schwarz am Zug}\hfill{\small Nr.~595\qquad\href{https://helpman.komtera.lt/?fen=8/6Rk/8/8/8/p1n5/2q5/K7&moves=7}{\color{white}\bfseries lösen\,\ensuremath{\nearrow}}}
\end{tcolorbox}
\vspace{2mm}
\noindent\begin{minipage}[t]{36mm}\vspace{0pt}\sffamily\small\raggedright
\hmfigure{längste eindeutige}{{\Large h\#7}}\hmfigure{eindeutig bei dieser Länge}{{\Large 1} Stellung}\hmfigure{längstes Matt}{{\Large h\#8}}\hmfigure{in der Klasse}{45,4 Milliarden Stellungen,\\ 25,5 Milliarden mit Hilfsmatt}
\end{minipage}\hfill
\begin{minipage}[t]{80mm}\vspace{0pt}\begin{tikzpicture}\begin{axis}[width=72mm,height=42mm,ybar,bar width=1.6pt,xmin=-1,xmax=17,ymin=0,ymax=5.746,axis y line*=left,axis x line*=bottom,xlabel={\scriptsize Halbzüge bis Matt},ylabel={\scriptsize Stellungen (Mrd.)},tick label style={font=\tiny},label style={font=\scriptsize},enlarge x limits=false,ylabel shift=-3pt,xlabel shift=-3pt]\addplot[fill=ink!55,draw=none] coordinates {(0,0.002951) (1,0.02028) (2,0.08852) (3,0.5015) (4,0.9017) (5,2.651) (6,3.098) (7,5.223) (8,4.324) (9,4.828) (10,2.28) (11,1.278) (12,0.2803) (13,0.03189) (14,0.002241) (15,4.006e-05) (16,3.4e-08)};\draw[densely dotted] (axis cs:14,0) -- (axis cs:14,5.746);\end{axis}\begin{axis}[width=72mm,height=42mm,axis y line*=right,axis x line=none,ymode=log,xmin=-1,xmax=17,ylabel={\scriptsize eindeutig (log)},tick label style={font=\tiny},label style={font=\scriptsize},ylabel shift=-3pt]\addplot[mark=*,mark size=0.7pt,draw=black,thin] coordinates {(0,2950945) (1,20280823) (2,39211390) (3,163265689) (4,166563931) (5,245052825) (6,126435622) (7,95457486) (8,39639941) (9,19929379) (10,3768079) (11,299287) (12,16854) (13,1069) (14,1)};\end{axis}\end{tikzpicture}\end{minipage}
\vspace{3mm}

\noindent\begin{minipage}[t]{60mm}\vspace{0pt}\sbox\hmboard{\chessboard[setfen=8/6Rk/8/8/8/p1n5/2q5/K7 b - - 0 1]}\usebox\hmboard\\[1mm]
\makebox[\wd\hmboard]{\sffamily\hspace{3mm}\textbf{h\#7}\hfill(2+4)}\end{minipage}\hfill
\begin{minipage}[t]{52mm}\vspace{0pt}
\textcolor{ink}{\sffamily\textsc{Lösung}}\par\smallskip\small
\begin{tabular}{@{}r@{\enspace}l@{\enspace}l@{}}
\textbf{1.} & \symking{}h6 & \symrook{}g2\\
\textbf{2.} & \symking{}h5 & \symrook{}xc2\\
\textbf{3.} & \symking{}g4 & \symrook{}a2\\
\textbf{4.} & \symking{}f3 & \symrook{}xa3\\
\textbf{5.} & \symking{}e2 & \symking{}b2\\
\textbf{6.} & \symking{}d1 & \symking{}xc3\\
\textbf{7.} & \symking{}c1 & \symrook{}a1\#\\
\end{tabular}\par\vspace{3mm}
\qrcode[height=12mm]{https://helpman.komtera.lt/?fen=8/6Rk/8/8/8/p1n5/2q5/K7&moves=7}\\{\sffamily\tiny\color{ink}scannen zum Lösen}
\end{minipage}
\vspace{2mm}

{\raggedright\noindent\textcolor{ink}{\sffamily\textsc{Themen}}\enspace \themetag{reines Matt} \themetag{Mustermatt} \themetag{Idealmatt} \themetag{Spiegelmatt} \themetag{ein Stein zieht} \themetag{ohne Schach}\par\smallskip{\footnotesize\textit{Hinweis:} die Partei am Zug steht im Diagramm im Schach.}\par}

\clearpage
\phantomsection\label{hm:KRvkqqb}
\begin{tcolorbox}[enhanced,colback=ink,colframe=ink,arc=0pt,boxrule=0pt,left=4mm,right=4mm,top=2.5mm,bottom=2.5mm]
\color{white}\sffamily{\LARGE\bfseries KRvkqqb}\hfill{\Large h\#7}\\[0.5mm]
{\small 6 Steine\quad Schwarz am Zug}\hfill{\small Nr.~596\qquad\href{https://helpman.komtera.lt/?fen=8/6Rk/8/8/2b5/1q6/2q5/K7&moves=7}{\color{white}\bfseries lösen\,\ensuremath{\nearrow}}}
\end{tcolorbox}
\vspace{2mm}
\noindent\begin{minipage}[t]{36mm}\vspace{0pt}\sffamily\small\raggedright
\hmfigure{längste eindeutige}{{\Large h\#7}}\hmfigure{eindeutig bei dieser Länge}{{\Large 4} Stellungen}\hmfigure{längstes Matt}{{\Large h\#9}}\hmfigure{in der Klasse}{15,5 Milliarden Stellungen,\\ 3,47 Milliarden mit Hilfsmatt}
\end{minipage}\hfill
\begin{minipage}[t]{80mm}\vspace{0pt}\begin{tikzpicture}\begin{axis}[width=72mm,height=42mm,ybar,bar width=1.6pt,xmin=-1,xmax=19,ymin=0,ymax=952.7,axis y line*=left,axis x line*=bottom,xlabel={\scriptsize Halbzüge bis Matt},ylabel={\scriptsize Stellungen (Mio.)},tick label style={font=\tiny},label style={font=\scriptsize},enlarge x limits=false,ylabel shift=-3pt,xlabel shift=-3pt]\addplot[fill=ink!55,draw=none] coordinates {(0,0.101) (1,0.7847) (2,4.413) (3,35.5) (4,61.59) (5,276.3) (6,256.6) (7,732) (8,447.7) (9,866.1) (10,318.9) (11,387.7) (12,56.92) (13,25.95) (14,0.9674) (15,0.2001) (16,0.001242) (17,0.000347) (18,1e-06)};\draw[densely dotted] (axis cs:14,0) -- (axis cs:14,952.7);\end{axis}\begin{axis}[width=72mm,height=42mm,axis y line*=right,axis x line=none,ymode=log,xmin=-1,xmax=19,ylabel={\scriptsize eindeutig (log)},tick label style={font=\tiny},label style={font=\scriptsize},ylabel shift=-3pt]\addplot[mark=*,mark size=0.7pt,draw=black,thin] coordinates {(0,100997) (1,784660) (2,1779837) (3,10935389) (4,8726653) (5,25081557) (6,9232039) (7,15683117) (8,4387617) (9,5652690) (10,846906) (11,433809) (12,6049) (13,1182) (14,4)};\end{axis}\end{tikzpicture}\end{minipage}
\vspace{3mm}

\noindent\begin{minipage}[t]{60mm}\vspace{0pt}\sbox\hmboard{\chessboard[setfen=8/6Rk/8/8/2b5/1q6/2q5/K7 b - - 0 1]}\usebox\hmboard\\[1mm]
\makebox[\wd\hmboard]{\sffamily\hspace{3mm}\textbf{h\#7}\hfill(2+4)}\end{minipage}\hfill
\begin{minipage}[t]{52mm}\vspace{0pt}
\textcolor{ink}{\sffamily\textsc{Lösung}}\par\smallskip\small
\begin{tabular}{@{}r@{\enspace}l@{\enspace}l@{}}
\textbf{1.} & \symking{}h6 & \symrook{}g2\\
\textbf{2.} & \symking{}h5 & \symrook{}xc2\\
\textbf{3.} & \symking{}g4 & \symrook{}b2\\
\textbf{4.} & \symking{}f3 & \symrook{}xb3+\\
\textbf{5.} & \symking{}e2 & \symking{}b2\\
\textbf{6.} & \symking{}d1 & \symking{}c3\\
\textbf{7.} & \symbishop{}e2 & \symrook{}b1\#\\
\end{tabular}\par\vspace{3mm}
\qrcode[height=12mm]{https://helpman.komtera.lt/?fen=8/6Rk/8/8/2b5/1q6/2q5/K7&moves=7}\\{\sffamily\tiny\color{ink}scannen zum Lösen}
\end{minipage}
\vspace{2mm}

{\raggedright\noindent\textcolor{ink}{\sffamily\textsc{Themen}}\enspace \themetag{reines Matt} \themetag{Mustermatt} \themetag{Idealmatt} \themetag{Selbstblock}\par\smallskip{\footnotesize\textit{Hinweis:} die Partei am Zug steht im Diagramm im Schach.}\par}

\clearpage
\phantomsection\label{hm:KRvkqqn}
\begin{tcolorbox}[enhanced,colback=ink,colframe=ink,arc=0pt,boxrule=0pt,left=4mm,right=4mm,top=2.5mm,bottom=2.5mm]
\color{white}\sffamily{\LARGE\bfseries KRvkqqn}\hfill{\Large h\#7}\\[0.5mm]
{\small 6 Steine\quad Schwarz am Zug}\hfill{\small Nr.~597\qquad\href{https://helpman.komtera.lt/?fen=8/7k/7R/8/7n/8/2q1q3/K7&moves=7}{\color{white}\bfseries lösen\,\ensuremath{\nearrow}}}
\end{tcolorbox}
\vspace{2mm}
\noindent\begin{minipage}[t]{36mm}\vspace{0pt}\sffamily\small\raggedright
\hmfigure{längste eindeutige}{{\Large h\#7}}\hmfigure{eindeutig bei dieser Länge}{{\Large 5} Stellungen}\hmfigure{längstes Matt}{{\Large h\#8}}\hmfigure{in der Klasse}{15,5 Milliarden Stellungen,\\ 3,57 Milliarden mit Hilfsmatt}
\end{minipage}\hfill
\begin{minipage}[t]{80mm}\vspace{0pt}\begin{tikzpicture}\begin{axis}[width=72mm,height=42mm,ybar,bar width=1.6pt,xmin=-1,xmax=17,ymin=0,ymax=937.7,axis y line*=left,axis x line*=bottom,xlabel={\scriptsize Halbzüge bis Matt},ylabel={\scriptsize Stellungen (Mio.)},tick label style={font=\tiny},label style={font=\scriptsize},enlarge x limits=false,ylabel shift=-3pt,xlabel shift=-3pt]\addplot[fill=ink!55,draw=none] coordinates {(0,0.129) (1,1.018) (2,5.659) (3,42.78) (4,74.16) (5,308.2) (6,293.9) (7,760.3) (8,477.8) (9,852.4) (10,324.9) (11,355.8) (12,54.81) (13,21.56) (14,0.9459) (15,0.1532) (16,0.001033)};\draw[densely dotted] (axis cs:14,0) -- (axis cs:14,937.7);\end{axis}\begin{axis}[width=72mm,height=42mm,axis y line*=right,axis x line=none,ymode=log,xmin=-1,xmax=17,ylabel={\scriptsize eindeutig (log)},tick label style={font=\tiny},label style={font=\scriptsize},ylabel shift=-3pt]\addplot[mark=*,mark size=0.7pt,draw=black,thin] coordinates {(0,128955) (1,1017898) (2,2332209) (3,13492652) (4,11139151) (5,29195056) (6,11800810) (7,17148314) (8,4925996) (9,5515820) (10,872534) (11,362435) (12,7103) (13,890) (14,5)};\end{axis}\end{tikzpicture}\end{minipage}
\vspace{3mm}

\noindent\begin{minipage}[t]{60mm}\vspace{0pt}\sbox\hmboard{\chessboard[setfen=8/7k/7R/8/7n/8/2q1q3/K7 b - - 0 1]}\usebox\hmboard\\[1mm]
\makebox[\wd\hmboard]{\sffamily\hspace{3mm}\textbf{h\#7}\hfill(2+4)}\end{minipage}\hfill
\begin{minipage}[t]{52mm}\vspace{0pt}
\textcolor{ink}{\sffamily\textsc{Lösung}}\par\smallskip\small
\begin{tabular}{@{}r@{\enspace}l@{\enspace}l@{}}
\textbf{1.} & \symking{}g7 & \symrook{}h8\\
\textbf{2.} & \symqueen{}b1+ & \symking{}xb1\\
\textbf{3.} & \symqueen{}h5 & \symking{}c2\\
\textbf{4.} & \symknight{}g6 & \symking{}d3\\
\textbf{5.} & \symqueen{}h7 & \symking{}e4\\
\textbf{6.} & \symking{}h6 & \symking{}f5\\
\textbf{7.} & \symking{}h5 & \symrook{}xh7\#\\
\end{tabular}\par\vspace{3mm}
\qrcode[height=12mm]{https://helpman.komtera.lt/?fen=8/7k/7R/8/7n/8/2q1q3/K7&moves=7}\\{\sffamily\tiny\color{ink}scannen zum Lösen}
\end{minipage}
\vspace{2mm}

{\raggedright\noindent\textcolor{ink}{\sffamily\textsc{Themen}}\enspace \themetag{Maslar}\par\smallskip{\footnotesize\textit{Hinweis:} die Partei am Zug steht im Diagramm im Schach.}\par}

\clearpage
\phantomsection\label{hm:KRvkqqp}
\begin{tcolorbox}[enhanced,colback=ink,colframe=ink,arc=0pt,boxrule=0pt,left=4mm,right=4mm,top=2.5mm,bottom=2.5mm]
\color{white}\sffamily{\LARGE\bfseries KRvkqqp}\hfill{\Large h\#7}\\[0.5mm]
{\small 6 Steine\quad Schwarz am Zug}\hfill{\small Nr.~598\qquad\href{https://helpman.komtera.lt/?fen=7k/7R/8/3p4/8/1q6/6q1/K7&moves=7}{\color{white}\bfseries lösen\,\ensuremath{\nearrow}}}
\end{tcolorbox}
\vspace{2mm}
\noindent\begin{minipage}[t]{36mm}\vspace{0pt}\sffamily\small\raggedright
\hmfigure{längste eindeutige}{{\Large h\#7}}\hmfigure{eindeutig bei dieser Länge}{{\Large 6} Stellungen}\hmfigure{längstes Matt}{{\Large h\#8}}\hmfigure{in der Klasse}{45,4 Milliarden Stellungen,\\ 11 Milliarden mit Hilfsmatt}
\end{minipage}\hfill
\begin{minipage}[t]{80mm}\vspace{0pt}\begin{tikzpicture}\begin{axis}[width=72mm,height=42mm,ybar,bar width=1.6pt,xmin=-1,xmax=17,ymin=0,ymax=2.784,axis y line*=left,axis x line*=bottom,xlabel={\scriptsize Halbzüge bis Matt},ylabel={\scriptsize Stellungen (Mrd.)},tick label style={font=\tiny},label style={font=\scriptsize},enlarge x limits=false,ylabel shift=-3pt,xlabel shift=-3pt]\addplot[fill=ink!55,draw=none] coordinates {(0,0.0005294) (1,0.003888) (2,0.0207) (3,0.1502) (4,0.2558) (5,1.004) (6,0.9601) (7,2.338) (8,1.525) (9,2.531) (10,0.9859) (11,0.9876) (12,0.159) (13,0.05648) (14,0.002702) (15,0.0003908) (16,2.457e-06)};\draw[densely dotted] (axis cs:14,0) -- (axis cs:14,2.784);\end{axis}\begin{axis}[width=72mm,height=42mm,axis y line*=right,axis x line=none,ymode=log,xmin=-1,xmax=17,ylabel={\scriptsize eindeutig (log)},tick label style={font=\tiny},label style={font=\scriptsize},ylabel shift=-3pt]\addplot[mark=*,mark size=0.7pt,draw=black,thin] coordinates {(0,529423) (1,3888205) (2,8449936) (3,46258016) (4,36827933) (5,87145744) (6,34529681) (7,46973414) (8,14338770) (9,14706046) (10,2364862) (11,891795) (12,17836) (13,2116) (14,6)};\end{axis}\end{tikzpicture}\end{minipage}
\vspace{3mm}

\noindent\begin{minipage}[t]{60mm}\vspace{0pt}\sbox\hmboard{\chessboard[setfen=7k/7R/8/3p4/8/1q6/6q1/K7 b - - 0 1]}\usebox\hmboard\\[1mm]
\makebox[\wd\hmboard]{\sffamily\hspace{3mm}\textbf{h\#7}\hfill(2+4)}\end{minipage}\hfill
\begin{minipage}[t]{52mm}\vspace{0pt}
\textcolor{ink}{\sffamily\textsc{Lösung}}\par\smallskip\small
\begin{tabular}{@{}r@{\enspace}l@{\enspace}l@{}}
\textbf{1.} & \symking{}g8 & \symrook{}h1\\
\textbf{2.} & \symqueen{}bh3 & \symking{}b1\\
\textbf{3.} & \symqueen{}g7 & \symking{}c2\\
\textbf{4.} & \symqueen{}h2+ & \symking{}d3\\
\textbf{5.} & d4 & \symking{}e4\\
\textbf{6.} & \symking{}h7 & \symking{}f5\\
\textbf{7.} & \symking{}h6 & \symrook{}xh2\#\\
\end{tabular}\par\vspace{3mm}
\qrcode[height=12mm]{https://helpman.komtera.lt/?fen=7k/7R/8/3p4/8/1q6/6q1/K7&moves=7}\\{\sffamily\tiny\color{ink}scannen zum Lösen}
\end{minipage}
\vspace{2mm}

{\raggedright\noindent\textcolor{ink}{\sffamily\textsc{Themen}}\enspace \themetag{reines Matt} \themetag{Mustermatt} \themetag{Selbstblock} \themetag{Maslar}\par\smallskip{\footnotesize\textit{Hinweis:} die Partei am Zug steht im Diagramm im Schach.}\par}

\clearpage
\phantomsection\label{hm:KRvkqqq}
\begin{tcolorbox}[enhanced,colback=ink,colframe=ink,arc=0pt,boxrule=0pt,left=4mm,right=4mm,top=2.5mm,bottom=2.5mm]
\color{white}\sffamily{\LARGE\bfseries KRvkqqq}\hfill{\Large h\#7}\\[0.5mm]
{\small 6 Steine\quad Schwarz am Zug}\hfill{\small Nr.~599\qquad\href{https://helpman.komtera.lt/?fen=8/8/6k1/6R1/1qq5/1q6/8/K7&moves=7}{\color{white}\bfseries lösen\,\ensuremath{\nearrow}}}
\end{tcolorbox}
\vspace{2mm}
\noindent\begin{minipage}[t]{36mm}\vspace{0pt}\sffamily\small\raggedright
\hmfigure{längste eindeutige}{{\Large h\#7}}\hmfigure{eindeutig bei dieser Länge}{{\Large 22} Stellungen}\hmfigure{längstes Matt}{{\Large h\#8.5}}\hmfigure{in der Klasse}{15,5 Milliarden Stellungen,\\ 1 Milliarden mit Hilfsmatt}
\end{minipage}\hfill
\begin{minipage}[t]{80mm}\vspace{0pt}\begin{tikzpicture}\begin{axis}[width=72mm,height=42mm,ybar,bar width=1.6pt,xmin=-1,xmax=18,ymin=0,ymax=299.3,axis y line*=left,axis x line*=bottom,xlabel={\scriptsize Halbzüge bis Matt},ylabel={\scriptsize Stellungen (Mio.)},tick label style={font=\tiny},label style={font=\scriptsize},enlarge x limits=false,ylabel shift=-3pt,xlabel shift=-3pt]\addplot[fill=ink!55,draw=none] coordinates {(0,0.01442) (1,0.1239) (2,0.783) (3,7.239) (4,12.26) (5,65.98) (6,54.78) (7,201.4) (8,106.1) (9,272.1) (10,86.11) (11,157.5) (12,19.75) (13,18.15) (14,0.5577) (15,0.2366) (16,0.001532) (17,0.000183)};\draw[densely dotted] (axis cs:14,0) -- (axis cs:14,299.3);\end{axis}\begin{axis}[width=72mm,height=42mm,axis y line*=right,axis x line=none,ymode=log,xmin=-1,xmax=18,ylabel={\scriptsize eindeutig (log)},tick label style={font=\tiny},label style={font=\scriptsize},ylabel shift=-3pt]\addplot[mark=*,mark size=0.7pt,draw=black,thin] coordinates {(0,14415) (1,123889) (2,300403) (3,2158991) (4,1540005) (5,5890310) (6,1798072) (7,4401981) (8,933500) (9,1818714) (10,238883) (11,255890) (12,6743) (13,3935) (14,22)};\end{axis}\end{tikzpicture}\end{minipage}
\vspace{3mm}

\noindent\begin{minipage}[t]{60mm}\vspace{0pt}\sbox\hmboard{\chessboard[setfen=8/8/6k1/6R1/1qq5/1q6/8/K7 b - - 0 1]}\usebox\hmboard\\[1mm]
\makebox[\wd\hmboard]{\sffamily\hspace{3mm}\textbf{h\#7}\hfill(2+4)}\end{minipage}\hfill
\begin{minipage}[t]{52mm}\vspace{0pt}
\textcolor{ink}{\sffamily\textsc{Lösung}}\par\smallskip\small
\begin{tabular}{@{}r@{\enspace}l@{\enspace}l@{}}
\textbf{1.} & \symking{}f6 & \symrook{}g3\\
\textbf{2.} & \symking{}e5 & \symrook{}xb3\\
\textbf{3.} & \symking{}d4 & \symrook{}xb4\\
\textbf{4.} & \symking{}c3 & \symrook{}xc4+\\
\textbf{5.} & \symking{}b3 & \symking{}b1\\
\textbf{6.} & \symking{}a3 & \symking{}c2\\
\textbf{7.} & \symking{}a2 & \symrook{}a4\#\\
\end{tabular}\par\vspace{3mm}
\qrcode[height=12mm]{https://helpman.komtera.lt/?fen=8/8/6k1/6R1/1qq5/1q6/8/K7&moves=7}\\{\sffamily\tiny\color{ink}scannen zum Lösen}
\end{minipage}
\vspace{2mm}

{\raggedright\noindent\textcolor{ink}{\sffamily\textsc{Themen}}\enspace \themetag{reines Matt} \themetag{Mustermatt} \themetag{Idealmatt} \themetag{Spiegelmatt} \themetag{ein Stein zieht}\par\smallskip{\footnotesize\textit{Hinweis:} die Partei am Zug steht im Diagramm im Schach.}\par}

\clearpage
\phantomsection\label{hm:KRvkqqr}
\begin{tcolorbox}[enhanced,colback=ink,colframe=ink,arc=0pt,boxrule=0pt,left=4mm,right=4mm,top=2.5mm,bottom=2.5mm]
\color{white}\sffamily{\LARGE\bfseries KRvkqqr}\hfill{\Large h\#7}\\[0.5mm]
{\small 6 Steine\quad Schwarz am Zug}\hfill{\small Nr.~600\qquad\href{https://helpman.komtera.lt/?fen=8/8/7k/7R/1r6/1q6/3q4/K7&moves=7}{\color{white}\bfseries lösen\,\ensuremath{\nearrow}}}
\end{tcolorbox}
\vspace{2mm}
\noindent\begin{minipage}[t]{36mm}\vspace{0pt}\sffamily\small\raggedright
\hmfigure{längste eindeutige}{{\Large h\#7}}\hmfigure{eindeutig bei dieser Länge}{{\Large 49} Stellungen}\hmfigure{längstes Matt}{{\Large h\#8.5}}\hmfigure{in der Klasse}{15,5 Milliarden Stellungen,\\ 3,3 Milliarden mit Hilfsmatt}
\end{minipage}\hfill
\begin{minipage}[t]{80mm}\vspace{0pt}\begin{tikzpicture}\begin{axis}[width=72mm,height=42mm,ybar,bar width=1.6pt,xmin=-1,xmax=18,ymin=0,ymax=916.8,axis y line*=left,axis x line*=bottom,xlabel={\scriptsize Halbzüge bis Matt},ylabel={\scriptsize Stellungen (Mio.)},tick label style={font=\tiny},label style={font=\scriptsize},enlarge x limits=false,ylabel shift=-3pt,xlabel shift=-3pt]\addplot[fill=ink!55,draw=none] coordinates {(0,0.07767) (1,0.648) (2,3.897) (3,32.26) (4,56.47) (5,255.8) (6,234.4) (7,690.4) (8,388.1) (9,833.5) (10,285.9) (11,416.9) (12,56.94) (13,38.21) (14,1.422) (15,0.3844) (16,0.002435) (17,3e-06)};\draw[densely dotted] (axis cs:14,0) -- (axis cs:14,916.8);\end{axis}\begin{axis}[width=72mm,height=42mm,axis y line*=right,axis x line=none,ymode=log,xmin=-1,xmax=18,ylabel={\scriptsize eindeutig (log)},tick label style={font=\tiny},label style={font=\scriptsize},ylabel shift=-3pt]\addplot[mark=*,mark size=0.7pt,draw=black,thin] coordinates {(0,77671) (1,647995) (2,1591965) (3,9971096) (4,8115345) (5,23238036) (6,8388222) (7,15081459) (8,3656389) (9,5411224) (10,805669) (11,588609) (12,15675) (13,4869) (14,49)};\end{axis}\end{tikzpicture}\end{minipage}
\vspace{3mm}

\noindent\begin{minipage}[t]{60mm}\vspace{0pt}\sbox\hmboard{\chessboard[setfen=8/8/7k/7R/1r6/1q6/3q4/K7 b - - 0 1]}\usebox\hmboard\\[1mm]
\makebox[\wd\hmboard]{\sffamily\hspace{3mm}\textbf{h\#7}\hfill(2+4)}\end{minipage}\hfill
\begin{minipage}[t]{52mm}\vspace{0pt}
\textcolor{ink}{\sffamily\textsc{Lösung}}\par\smallskip\small
\begin{tabular}{@{}r@{\enspace}l@{\enspace}l@{}}
\textbf{1.} & \symking{}g6 & \symrook{}h3\\
\textbf{2.} & \symking{}f5 & \symrook{}xb3\\
\textbf{3.} & \symking{}e4 & \symrook{}xb4+\\
\textbf{4.} & \symking{}d3 & \symrook{}a4\\
\textbf{5.} & \symking{}c2 & \symking{}a2\\
\textbf{6.} & \symking{}d1+ & \symking{}b3\\
\textbf{7.} & \symking{}c1 & \symrook{}a1\#\\
\end{tabular}\par\vspace{3mm}
\qrcode[height=12mm]{https://helpman.komtera.lt/?fen=8/8/7k/7R/1r6/1q6/3q4/K7&moves=7}\\{\sffamily\tiny\color{ink}scannen zum Lösen}
\end{minipage}
\vspace{2mm}

{\raggedright\noindent\textcolor{ink}{\sffamily\textsc{Themen}}\enspace \themetag{reines Matt} \themetag{Mustermatt} \themetag{Idealmatt} \themetag{ein Stein zieht}\par\smallskip{\footnotesize\textit{Hinweis:} die Partei am Zug steht im Diagramm im Schach.}\par}

\clearpage
\phantomsection\label{hm:KRvkqrb}
\begin{tcolorbox}[enhanced,colback=ink,colframe=ink,arc=0pt,boxrule=0pt,left=4mm,right=4mm,top=2.5mm,bottom=2.5mm]
\color{white}\sffamily{\LARGE\bfseries KRvkqrb}\hfill{\Large h\#7}\\[0.5mm]
{\small 6 Steine\quad Schwarz am Zug}\hfill{\small Nr.~601\qquad\href{https://helpman.komtera.lt/?fen=8/6Rk/8/4b3/8/8/1rq5/K7&moves=7}{\color{white}\bfseries lösen\,\ensuremath{\nearrow}}}
\end{tcolorbox}
\vspace{2mm}
\noindent\begin{minipage}[t]{36mm}\vspace{0pt}\sffamily\small\raggedright
\hmfigure{längste eindeutige}{{\Large h\#7}}\hmfigure{eindeutig bei dieser Länge}{{\Large 5} Stellungen}\hmfigure{längstes Matt}{{\Large h\#8}}\hmfigure{in der Klasse}{15,5 Milliarden Stellungen,\\ 7,55 Milliarden mit Hilfsmatt}
\end{minipage}\hfill
\begin{minipage}[t]{80mm}\vspace{0pt}\begin{tikzpicture}\begin{axis}[width=72mm,height=42mm,ybar,bar width=1.6pt,xmin=-1,xmax=17,ymin=0,ymax=1.88,axis y line*=left,axis x line*=bottom,xlabel={\scriptsize Halbzüge bis Matt},ylabel={\scriptsize Stellungen (Mrd.)},tick label style={font=\tiny},label style={font=\scriptsize},enlarge x limits=false,ylabel shift=-3pt,xlabel shift=-3pt]\addplot[fill=ink!55,draw=none] coordinates {(0,0.0003597) (1,0.002725) (2,0.01417) (3,0.1008) (4,0.1821) (5,0.6788) (6,0.7007) (7,1.592) (8,1.081) (9,1.709) (10,0.6993) (11,0.6477) (12,0.1108) (13,0.0319) (14,0.001536) (15,0.0001388) (16,8.63e-07)};\draw[densely dotted] (axis cs:14,0) -- (axis cs:14,1.88);\end{axis}\begin{axis}[width=72mm,height=42mm,axis y line*=right,axis x line=none,ymode=log,xmin=-1,xmax=17,ylabel={\scriptsize eindeutig (log)},tick label style={font=\tiny},label style={font=\scriptsize},ylabel shift=-3pt]\addplot[mark=*,mark size=0.7pt,draw=black,thin] coordinates {(0,359719) (1,2725354) (2,6031517) (3,31439062) (4,28809732) (5,60111873) (6,25812754) (7,32068931) (8,10554547) (9,10001808) (10,1753466) (11,519855) (12,12575) (13,1014) (14,5)};\end{axis}\end{tikzpicture}\end{minipage}
\vspace{3mm}

\noindent\begin{minipage}[t]{60mm}\vspace{0pt}\sbox\hmboard{\chessboard[setfen=8/6Rk/8/4b3/8/8/1rq5/K7 b - - 0 1]}\usebox\hmboard\\[1mm]
\makebox[\wd\hmboard]{\sffamily\hspace{3mm}\textbf{h\#7}\hfill(2+4)}\end{minipage}\hfill
\begin{minipage}[t]{52mm}\vspace{0pt}
\textcolor{ink}{\sffamily\textsc{Lösung}}\par\smallskip\small
\begin{tabular}{@{}r@{\enspace}l@{\enspace}l@{}}
\textbf{1.} & \symking{}h6 & \symrook{}c7\\
\textbf{2.} & \symking{}g5 & \symrook{}xc2\\
\textbf{3.} & \symking{}f4 & \symrook{}c5\\
\textbf{4.} & \symking{}e3 & \symrook{}xe5+\\
\textbf{5.} & \symking{}d2 & \symking{}xb2\\
\textbf{6.} & \symking{}d1 & \symking{}c3\\
\textbf{7.} & \symking{}c1 & \symrook{}e1\#\\
\end{tabular}\par\vspace{3mm}
\qrcode[height=12mm]{https://helpman.komtera.lt/?fen=8/6Rk/8/4b3/8/8/1rq5/K7&moves=7}\\{\sffamily\tiny\color{ink}scannen zum Lösen}
\end{minipage}
\vspace{2mm}

{\raggedright\noindent\textcolor{ink}{\sffamily\textsc{Themen}}\enspace \themetag{reines Matt} \themetag{Mustermatt} \themetag{Idealmatt} \themetag{Spiegelmatt} \themetag{ein Stein zieht}\par\smallskip{\footnotesize\textit{Hinweis:} die Partei am Zug steht im Diagramm im Schach.}\par}

\clearpage
\phantomsection\label{hm:KRvkqrp}
\begin{tcolorbox}[enhanced,colback=ink,colframe=ink,arc=0pt,boxrule=0pt,left=4mm,right=4mm,top=2.5mm,bottom=2.5mm]
\color{white}\sffamily{\LARGE\bfseries KRvkqrp}\hfill{\Large h\#7}\\[0.5mm]
{\small 6 Steine\quad Schwarz am Zug}\hfill{\small Nr.~602\qquad\href{https://helpman.komtera.lt/?fen=8/7k/7R/8/3p4/8/2qr4/K7&moves=7}{\color{white}\bfseries lösen\,\ensuremath{\nearrow}}}
\end{tcolorbox}
\vspace{2mm}
\noindent\begin{minipage}[t]{36mm}\vspace{0pt}\sffamily\small\raggedright
\hmfigure{längste eindeutige}{{\Large h\#7}}\hmfigure{eindeutig bei dieser Länge}{{\Large 9} Stellungen}\hmfigure{längstes Matt}{{\Large h\#8}}\hmfigure{in der Klasse}{45,4 Milliarden Stellungen,\\ 23,8 Milliarden mit Hilfsmatt}
\end{minipage}\hfill
\begin{minipage}[t]{80mm}\vspace{0pt}\begin{tikzpicture}\begin{axis}[width=72mm,height=42mm,ybar,bar width=1.6pt,xmin=-1,xmax=17,ymin=0,ymax=5.508,axis y line*=left,axis x line*=bottom,xlabel={\scriptsize Halbzüge bis Matt},ylabel={\scriptsize Stellungen (Mrd.)},tick label style={font=\tiny},label style={font=\scriptsize},enlarge x limits=false,ylabel shift=-3pt,xlabel shift=-3pt]\addplot[fill=ink!55,draw=none] coordinates {(0,0.00184) (1,0.01319) (2,0.06408) (3,0.4067) (4,0.7296) (5,2.363) (6,2.568) (7,5.007) (8,3.693) (9,4.972) (10,2.066) (11,1.572) (12,0.2775) (13,0.06243) (14,0.003301) (15,0.0001676) (16,8.31e-07)};\draw[densely dotted] (axis cs:14,0) -- (axis cs:14,5.508);\end{axis}\begin{axis}[width=72mm,height=42mm,axis y line*=right,axis x line=none,ymode=log,xmin=-1,xmax=17,ylabel={\scriptsize eindeutig (log)},tick label style={font=\tiny},label style={font=\scriptsize},ylabel shift=-3pt]\addplot[mark=*,mark size=0.7pt,draw=black,thin] coordinates {(0,1839679) (1,13186338) (2,27877683) (3,126227919) (4,117646039) (5,197185273) (6,91304603) (7,90788154) (8,32289080) (9,24104858) (10,4317910) (11,884687) (12,32771) (13,1498) (14,9)};\end{axis}\end{tikzpicture}\end{minipage}
\vspace{3mm}

\noindent\begin{minipage}[t]{60mm}\vspace{0pt}\sbox\hmboard{\chessboard[setfen=8/7k/7R/8/3p4/8/2qr4/K7 b - - 0 1]}\usebox\hmboard\\[1mm]
\makebox[\wd\hmboard]{\sffamily\hspace{3mm}\textbf{h\#7}\hfill(2+4)}\end{minipage}\hfill
\begin{minipage}[t]{52mm}\vspace{0pt}
\textcolor{ink}{\sffamily\textsc{Lösung}}\par\smallskip\small
\begin{tabular}{@{}r@{\enspace}l@{\enspace}l@{}}
\textbf{1.} & \symking{}g7 & \symrook{}h8\\
\textbf{2.} & \symqueen{}b1+ & \symking{}xb1\\
\textbf{3.} & \symrook{}d1+ & \symking{}c2\\
\textbf{4.} & \symrook{}h1 & \symking{}d3\\
\textbf{5.} & \symrook{}h7 & \symking{}e4\\
\textbf{6.} & \symking{}h6 & \symking{}f5\\
\textbf{7.} & \symking{}h5 & \symrook{}xh7\#\\
\end{tabular}\par\vspace{3mm}
\qrcode[height=12mm]{https://helpman.komtera.lt/?fen=8/7k/7R/8/3p4/8/2qr4/K7&moves=7}\\{\sffamily\tiny\color{ink}scannen zum Lösen}
\end{minipage}
\vspace{2mm}

{\raggedright\noindent\textcolor{ink}{\sffamily\textsc{Themen}}\enspace \themetag{reines Matt} \themetag{Mustermatt} \themetag{Spiegelmatt} \themetag{Maslar}\par\smallskip{\footnotesize\textit{Hinweis:} die Partei am Zug steht im Diagramm im Schach.}\par}

\clearpage
\phantomsection\label{hm:KRvkqrr}
\begin{tcolorbox}[enhanced,colback=ink,colframe=ink,arc=0pt,boxrule=0pt,left=4mm,right=4mm,top=2.5mm,bottom=2.5mm]
\color{white}\sffamily{\LARGE\bfseries KRvkqrr}\hfill{\Large h\#7}\\[0.5mm]
{\small 6 Steine\quad Schwarz am Zug}\hfill{\small Nr.~603\qquad\href{https://helpman.komtera.lt/?fen=8/8/7k/7R/1r6/1q6/3r4/K7&moves=7}{\color{white}\bfseries lösen\,\ensuremath{\nearrow}}}
\end{tcolorbox}
\vspace{2mm}
\noindent\begin{minipage}[t]{36mm}\vspace{0pt}\sffamily\small\raggedright
\hmfigure{längste eindeutige}{{\Large h\#7}}\hmfigure{eindeutig bei dieser Länge}{{\Large 22} Stellungen}\hmfigure{längstes Matt}{{\Large h\#8}}\hmfigure{in der Klasse}{15,5 Milliarden Stellungen,\\ 3,57 Milliarden mit Hilfsmatt}
\end{minipage}\hfill
\begin{minipage}[t]{80mm}\vspace{0pt}\begin{tikzpicture}\begin{axis}[width=72mm,height=42mm,ybar,bar width=1.6pt,xmin=-1,xmax=17,ymin=0,ymax=922.4,axis y line*=left,axis x line*=bottom,xlabel={\scriptsize Halbzüge bis Matt},ylabel={\scriptsize Stellungen (Mio.)},tick label style={font=\tiny},label style={font=\scriptsize},enlarge x limits=false,ylabel shift=-3pt,xlabel shift=-3pt]\addplot[fill=ink!55,draw=none] coordinates {(0,0.1332) (1,1.083) (2,6.052) (3,44.57) (4,79.93) (5,306.8) (6,305.8) (7,746.7) (8,467.6) (9,838.5) (10,317.6) (11,369.7) (12,58.08) (13,27.41) (14,1.239) (15,0.1749) (16,0.001025)};\draw[densely dotted] (axis cs:14,0) -- (axis cs:14,922.4);\end{axis}\begin{axis}[width=72mm,height=42mm,axis y line*=right,axis x line=none,ymode=log,xmin=-1,xmax=17,ylabel={\scriptsize eindeutig (log)},tick label style={font=\tiny},label style={font=\scriptsize},ylabel shift=-3pt]\addplot[mark=*,mark size=0.7pt,draw=black,thin] coordinates {(0,133179) (1,1082910) (2,2547014) (3,13638439) (4,11852882) (5,26259354) (6,10613988) (7,14766777) (8,4081756) (9,4833465) (10,834417) (11,417814) (12,13966) (13,2169) (14,22)};\end{axis}\end{tikzpicture}\end{minipage}
\vspace{3mm}

\noindent\begin{minipage}[t]{60mm}\vspace{0pt}\sbox\hmboard{\chessboard[setfen=8/8/7k/7R/1r6/1q6/3r4/K7 b - - 0 1]}\usebox\hmboard\\[1mm]
\makebox[\wd\hmboard]{\sffamily\hspace{3mm}\textbf{h\#7}\hfill(2+4)}\end{minipage}\hfill
\begin{minipage}[t]{52mm}\vspace{0pt}
\textcolor{ink}{\sffamily\textsc{Lösung}}\par\smallskip\small
\begin{tabular}{@{}r@{\enspace}l@{\enspace}l@{}}
\textbf{1.} & \symking{}g6 & \symrook{}h3\\
\textbf{2.} & \symking{}f5 & \symrook{}xb3\\
\textbf{3.} & \symking{}e4 & \symrook{}xb4+\\
\textbf{4.} & \symking{}d3 & \symrook{}a4\\
\textbf{5.} & \symking{}c2 & \symking{}a2\\
\textbf{6.} & \symking{}d1+ & \symking{}b3\\
\textbf{7.} & \symking{}c1 & \symrook{}a1\#\\
\end{tabular}\par\vspace{3mm}
\qrcode[height=12mm]{https://helpman.komtera.lt/?fen=8/8/7k/7R/1r6/1q6/3r4/K7&moves=7}\\{\sffamily\tiny\color{ink}scannen zum Lösen}
\end{minipage}
\vspace{2mm}

{\raggedright\noindent\textcolor{ink}{\sffamily\textsc{Themen}}\enspace \themetag{reines Matt} \themetag{Mustermatt} \themetag{Idealmatt} \themetag{ein Stein zieht}\par\smallskip{\footnotesize\textit{Hinweis:} die Partei am Zug steht im Diagramm im Schach.}\par}

\clearpage
\phantomsection\label{hm:KRvkrpp}
\begin{tcolorbox}[enhanced,colback=ink,colframe=ink,arc=0pt,boxrule=0pt,left=4mm,right=4mm,top=2.5mm,bottom=2.5mm]
\color{white}\sffamily{\LARGE\bfseries KRvkrpp}\hfill{\Large h\#7}\\[0.5mm]
{\small 6 Steine\quad Schwarz am Zug}\hfill{\small Nr.~604\qquad\href{https://helpman.komtera.lt/?fen=6Rk/8/8/8/8/2p5/p1r5/K7&moves=7}{\color{white}\bfseries lösen\,\ensuremath{\nearrow}}}
\end{tcolorbox}
\vspace{2mm}
\noindent\begin{minipage}[t]{36mm}\vspace{0pt}\sffamily\small\raggedright
\hmfigure{längste eindeutige}{{\Large h\#7}}\hmfigure{eindeutig bei dieser Länge}{{\Large 1} Stellung}\hmfigure{längstes Matt}{{\Large h\#7.5}}\hmfigure{in der Klasse}{34,1 Milliarden Stellungen,\\ 10,6 Milliarden mit Hilfsmatt}
\end{minipage}\hfill
\begin{minipage}[t]{80mm}\vspace{0pt}\begin{tikzpicture}\begin{axis}[width=72mm,height=42mm,ybar,bar width=1.6pt,xmin=-1,xmax=16,ymin=0,ymax=2.176,axis y line*=left,axis x line*=bottom,xlabel={\scriptsize Halbzüge bis Matt},ylabel={\scriptsize Stellungen (Mrd.)},tick label style={font=\tiny},label style={font=\scriptsize},enlarge x limits=false,ylabel shift=-3pt,xlabel shift=-3pt]\addplot[fill=ink!55,draw=none] coordinates {(0,0.002606) (1,0.01634) (2,0.05498) (3,0.2568) (4,0.5057) (5,1.121) (6,1.534) (7,1.971) (8,1.978) (9,1.692) (10,0.9511) (11,0.3785) (12,0.1156) (13,0.006131) (14,0.0008136) (15,4.12e-07)};\draw[densely dotted] (axis cs:14,0) -- (axis cs:14,2.176);\end{axis}\begin{axis}[width=72mm,height=42mm,axis y line*=right,axis x line=none,ymode=log,xmin=-1,xmax=16,ylabel={\scriptsize eindeutig (log)},tick label style={font=\tiny},label style={font=\scriptsize},ylabel shift=-3pt]\addplot[mark=*,mark size=0.7pt,draw=black,thin] coordinates {(0,2605747) (1,16340840) (2,27144881) (3,81591183) (4,106995178) (5,90700560) (6,54868993) (7,26853492) (8,14880774) (9,4998231) (10,1505980) (11,71472) (12,5379) (13,7) (14,1)};\end{axis}\end{tikzpicture}\end{minipage}
\vspace{3mm}

\noindent\begin{minipage}[t]{60mm}\vspace{0pt}\sbox\hmboard{\chessboard[setfen=6Rk/8/8/8/8/2p5/p1r5/K7 b - - 0 1]}\usebox\hmboard\\[1mm]
\makebox[\wd\hmboard]{\sffamily\hspace{3mm}\textbf{h\#7}\hfill(2+4)}\end{minipage}\hfill
\begin{minipage}[t]{52mm}\vspace{0pt}
\textcolor{ink}{\sffamily\textsc{Lösung}}\par\smallskip\small
\begin{tabular}{@{}r@{\enspace}l@{\enspace}l@{}}
\textbf{1.} & \symking{}h7 & \symrook{}a8\\
\textbf{2.} & \symrook{}c1+ & \symking{}xa2\\
\textbf{3.} & \symking{}g7 & \symking{}b3\\
\textbf{4.} & \symrook{}b1+ & \symking{}c4\\
\textbf{5.} & \symrook{}b8 & \symking{}d5\\
\textbf{6.} & \symking{}f8 & \symking{}e6\\
\textbf{7.} & \symking{}e8 & \symrook{}xb8\#\\
\end{tabular}\par\vspace{3mm}
\qrcode[height=12mm]{https://helpman.komtera.lt/?fen=6Rk/8/8/8/8/2p5/p1r5/K7&moves=7}\\{\sffamily\tiny\color{ink}scannen zum Lösen}
\end{minipage}
\vspace{2mm}

{\raggedright\noindent\textcolor{ink}{\sffamily\textsc{Themen}}\enspace \themetag{reines Matt} \themetag{Mustermatt} \themetag{Spiegelmatt} \themetag{Maslar}\par\smallskip{\footnotesize\textit{Hinweis:} die Partei am Zug steht im Diagramm im Schach.}\par}

\clearpage
\phantomsection\label{hm:KRvkrrp}
\begin{tcolorbox}[enhanced,colback=ink,colframe=ink,arc=0pt,boxrule=0pt,left=4mm,right=4mm,top=2.5mm,bottom=2.5mm]
\color{white}\sffamily{\LARGE\bfseries KRvkrrp}\hfill{\Large h\#7}\\[0.5mm]
{\small 6 Steine\quad Schwarz am Zug}\hfill{\small Nr.~605\qquad\href{https://helpman.komtera.lt/?fen=7k/7R/8/8/8/7p/3r3r/K7&moves=7}{\color{white}\bfseries lösen\,\ensuremath{\nearrow}}}
\end{tcolorbox}
\vspace{2mm}
\noindent\begin{minipage}[t]{36mm}\vspace{0pt}\sffamily\small\raggedright
\hmfigure{längste eindeutige}{{\Large h\#7}}\hmfigure{eindeutig bei dieser Länge}{{\Large 3} Stellungen}\hmfigure{längstes Matt}{{\Large h\#8}}\hmfigure{in der Klasse}{45,4 Milliarden Stellungen,\\ 12,8 Milliarden mit Hilfsmatt}
\end{minipage}\hfill
\begin{minipage}[t]{80mm}\vspace{0pt}\begin{tikzpicture}\begin{axis}[width=72mm,height=42mm,ybar,bar width=1.6pt,xmin=-1,xmax=17,ymin=0,ymax=2.817,axis y line*=left,axis x line*=bottom,xlabel={\scriptsize Halbzüge bis Matt},ylabel={\scriptsize Stellungen (Mrd.)},tick label style={font=\tiny},label style={font=\scriptsize},enlarge x limits=false,ylabel shift=-3pt,xlabel shift=-3pt]\addplot[fill=ink!55,draw=none] coordinates {(0,0.001566) (1,0.01094) (2,0.04638) (3,0.2563) (4,0.4853) (5,1.297) (6,1.592) (7,2.561) (8,2.174) (9,2.419) (10,1.142) (11,0.6775) (12,0.1465) (13,0.01974) (14,0.001437) (15,1.103e-05) (16,1.63e-07)};\draw[densely dotted] (axis cs:14,0) -- (axis cs:14,2.817);\end{axis}\begin{axis}[width=72mm,height=42mm,axis y line*=right,axis x line=none,ymode=log,xmin=-1,xmax=17,ylabel={\scriptsize eindeutig (log)},tick label style={font=\tiny},label style={font=\scriptsize},ylabel shift=-3pt]\addplot[mark=*,mark size=0.7pt,draw=black,thin] coordinates {(0,1566191) (1,10934082) (2,21007316) (3,76752488) (4,82930306) (5,96651591) (6,51634334) (7,39716769) (8,16990944) (9,9637324) (10,2195521) (11,274837) (12,13626) (13,30) (14,3)};\end{axis}\end{tikzpicture}\end{minipage}
\vspace{3mm}

\noindent\begin{minipage}[t]{60mm}\vspace{0pt}\sbox\hmboard{\chessboard[setfen=7k/7R/8/8/8/7p/3r3r/K7 b - - 0 1]}\usebox\hmboard\\[1mm]
\makebox[\wd\hmboard]{\sffamily\hspace{3mm}\textbf{h\#7}\hfill(2+4)}\end{minipage}\hfill
\begin{minipage}[t]{52mm}\vspace{0pt}
\textcolor{ink}{\sffamily\textsc{Lösung}}\par\smallskip\small
\begin{tabular}{@{}r@{\enspace}l@{\enspace}l@{}}
\textbf{1.} & \symking{}g8 & \symking{}b1\\
\textbf{2.} & \symrook{}hg2 & \symking{}c1\\
\textbf{3.} & \symrook{}g7 & \symking{}xd2\\
\textbf{4.} & \symking{}f7 & \symking{}e3\\
\textbf{5.} & \symking{}g6 & \symrook{}xg7+\\
\textbf{6.} & \symking{}h5 & \symking{}f4\\
\textbf{7.} & \symking{}h4 & \symrook{}h7\#\\
\end{tabular}\par\vspace{3mm}
\qrcode[height=12mm]{https://helpman.komtera.lt/?fen=7k/7R/8/8/8/7p/3r3r/K7&moves=7}\\{\sffamily\tiny\color{ink}scannen zum Lösen}
\end{minipage}
\vspace{2mm}

{\raggedright\noindent\textcolor{ink}{\sffamily\textsc{Themen}}\enspace \themetag{Switchback}\par\smallskip{\footnotesize\textit{Hinweis:} die Partei am Zug steht im Diagramm im Schach.}\par}

\clearpage
\phantomsection\label{hm:KBBBvkq}
\begin{tcolorbox}[enhanced,colback=ink,colframe=ink,arc=0pt,boxrule=0pt,left=4mm,right=4mm,top=2.5mm,bottom=2.5mm]
\color{white}\sffamily{\LARGE\bfseries KBBBvkq}\hfill{\Large h\#6.5}\\[0.5mm]
{\small 6 Steine\quad Weiß am Zug}\hfill{\small Nr.~606\qquad\href{https://helpman.komtera.lt/?fen=4B3/4B3/7k/8/8/8/q7/K5B1&moves=6.5}{\color{white}\bfseries lösen\,\ensuremath{\nearrow}}}
\end{tcolorbox}
\vspace{2mm}
\noindent\begin{minipage}[t]{36mm}\vspace{0pt}\sffamily\small\raggedright
\hmfigure{längste eindeutige}{{\Large h\#6.5}}\hmfigure{eindeutig bei dieser Länge}{{\Large 29} Stellungen}\hmfigure{längstes Matt}{{\Large h\#8}}\hmfigure{in der Klasse}{15,5 Milliarden Stellungen,\\ 1 Milliarden mit Hilfsmatt}
\end{minipage}\hfill
\begin{minipage}[t]{80mm}\vspace{0pt}\begin{tikzpicture}\begin{axis}[width=72mm,height=42mm,ybar,bar width=1.6pt,xmin=-1,xmax=17,ymin=0,ymax=374.4,axis y line*=left,axis x line*=bottom,xlabel={\scriptsize Halbzüge bis Matt},ylabel={\scriptsize Stellungen (Mio.)},tick label style={font=\tiny},label style={font=\scriptsize},enlarge x limits=false,ylabel shift=-3pt,xlabel shift=-3pt]\addplot[fill=ink!55,draw=none] coordinates {(0,0.2923) (1,1.136) (2,6.609) (3,38.53) (4,138.7) (5,270.9) (6,340.3) (7,167.9) (8,31.12) (9,4.457) (10,0.1424) (11,0.2743) (12,0.02935) (13,0.08609) (14,0.002915) (15,0.001605) (16,3e-06)};\draw[densely dotted] (axis cs:13,0) -- (axis cs:13,374.4);\end{axis}\begin{axis}[width=72mm,height=42mm,axis y line*=right,axis x line=none,ymode=log,xmin=-1,xmax=17,ylabel={\scriptsize eindeutig (log)},tick label style={font=\tiny},label style={font=\scriptsize},ylabel shift=-3pt]\addplot[mark=*,mark size=0.7pt,draw=black,thin] coordinates {(0,292316) (1,1047646) (2,3850815) (3,13705280) (4,26783998) (5,12888068) (6,2412486) (7,173198) (8,3015) (9,712) (10,537) (11,719) (12,6) (13,29)};\end{axis}\end{tikzpicture}\end{minipage}
\vspace{3mm}

\noindent\begin{minipage}[t]{60mm}\vspace{0pt}\sbox\hmboard{\chessboard[setfen=4B3/4B3/7k/8/8/8/q7/K5B1 w - - 0 1]}\usebox\hmboard\\[1mm]
\makebox[\wd\hmboard]{\sffamily\hspace{3mm}\textbf{h\#6.5}\hfill(4+2)}\end{minipage}\hfill
\begin{minipage}[t]{52mm}\vspace{0pt}
\textcolor{ink}{\sffamily\textsc{Lösung}}\par\smallskip\small
\begin{tabular}{@{}r@{\enspace}l@{\enspace}l@{}}
\textbf{1.} & \ldots & \symking{}xa2\\
\textbf{2.} & \symking{}g7 & \symbishop{}d7\\
\textbf{3.} & \symking{}f7 & \symking{}b3\\
\textbf{4.} & \symking{}xe7 & \symking{}c4\\
\textbf{5.} & \symking{}d8 & \symking{}d5\\
\textbf{6.} & \symking{}c7 & \symking{}e6\\
\textbf{7.} & \symking{}d8 & \symbishop{}b6\#\\
\end{tabular}\par\vspace{3mm}
\qrcode[height=12mm]{https://helpman.komtera.lt/?fen=4B3/4B3/7k/8/8/8/q7/K5B1&moves=6.5}\\{\sffamily\tiny\color{ink}scannen zum Lösen}
\end{minipage}
\vspace{2mm}

{\raggedright\noindent\textcolor{ink}{\sffamily\textsc{Themen}}\enspace \themetag{reines Matt} \themetag{Mustermatt} \themetag{Idealmatt} \themetag{Switchback} \themetag{ein Stein zieht} \themetag{ohne Schach}\par\smallskip{\footnotesize\textit{Hinweis:} die Partei am Zug steht im Diagramm im Schach; die Lösung beginnt mit einem Schlag.}\par}

\clearpage
\phantomsection\label{hm:KBBNNvk}
\begin{tcolorbox}[enhanced,colback=ink,colframe=ink,arc=0pt,boxrule=0pt,left=4mm,right=4mm,top=2.5mm,bottom=2.5mm]
\color{white}\sffamily{\LARGE\bfseries KBBNNvk}\hfill{\Large h\#6.5}\\[0.5mm]
{\small 6 Steine\quad Weiß am Zug}\hfill{\small Nr.~607\qquad\href{https://helpman.komtera.lt/?fen=7k/6NB/1N4B1/8/8/8/8/K7&moves=6.5}{\color{white}\bfseries lösen\,\ensuremath{\nearrow}}}
\end{tcolorbox}
\vspace{2mm}
\noindent\begin{minipage}[t]{36mm}\vspace{0pt}\sffamily\small\raggedright
\hmfigure{längste eindeutige}{{\Large h\#6.5}}\hmfigure{eindeutig bei dieser Länge}{{\Large 1} Stellung}\hmfigure{längstes Matt}{{\Large h\#8}}\hmfigure{in der Klasse}{15,5 Milliarden Stellungen,\\ 2,47 Milliarden mit Hilfsmatt}
\end{minipage}\hfill
\begin{minipage}[t]{80mm}\vspace{0pt}\begin{tikzpicture}\begin{axis}[width=72mm,height=42mm,ybar,bar width=1.6pt,xmin=-1,xmax=17,ymin=0,ymax=909.8,axis y line*=left,axis x line*=bottom,xlabel={\scriptsize Halbzüge bis Matt},ylabel={\scriptsize Stellungen (Mio.)},tick label style={font=\tiny},label style={font=\scriptsize},enlarge x limits=false,ylabel shift=-3pt,xlabel shift=-3pt]\addplot[fill=ink!55,draw=none] coordinates {(0,1.397) (1,4.609) (2,13.44) (3,54.15) (4,226.3) (5,444.7) (6,827.1) (7,409.1) (8,417.5) (9,53.96) (10,19.57) (11,0.2268) (12,0.3098) (13,0.001402) (14,0.05905) (15,0.000102) (16,0.000124)};\draw[densely dotted] (axis cs:13,0) -- (axis cs:13,909.8);\end{axis}\begin{axis}[width=72mm,height=42mm,axis y line*=right,axis x line=none,ymode=log,xmin=-1,xmax=17,ylabel={\scriptsize eindeutig (log)},tick label style={font=\tiny},label style={font=\scriptsize},ylabel shift=-3pt]\addplot[mark=*,mark size=0.7pt,draw=black,thin] coordinates {(0,1397459) (1,3997292) (2,11615554) (3,32261279) (4,85429335) (5,38111605) (6,34906958) (7,1759166) (8,597792) (9,2066) (10,2202) (11,3) (12,106) (13,1)};\end{axis}\end{tikzpicture}\end{minipage}
\vspace{3mm}

\noindent\begin{minipage}[t]{60mm}\vspace{0pt}\sbox\hmboard{\chessboard[setfen=7k/6NB/1N4B1/8/8/8/8/K7 w - - 0 1]}\usebox\hmboard\\[1mm]
\makebox[\wd\hmboard]{\sffamily\hspace{3mm}\textbf{h\#6.5}\hfill(5+1)}\end{minipage}\hfill
\begin{minipage}[t]{52mm}\vspace{0pt}
\textcolor{ink}{\sffamily\textsc{Lösung}}\par\smallskip\small
\begin{tabular}{@{}r@{\enspace}l@{\enspace}l@{}}
\textbf{1.} & \ldots & \symbishop{}g8\\
\textbf{2.} & \symking{}xg8 & \symking{}b2\\
\textbf{3.} & \symking{}f8 & \symking{}c3\\
\textbf{4.} & \symking{}e7 & \symking{}d4\\
\textbf{5.} & \symking{}d8 & \symking{}e5\\
\textbf{6.} & \symking{}c7 & \symking{}f6\\
\textbf{7.} & \symking{}d8 & \symknight{}e6\#\\
\end{tabular}\par\vspace{3mm}
\qrcode[height=12mm]{https://helpman.komtera.lt/?fen=7k/6NB/1N4B1/8/8/8/8/K7&moves=6.5}\\{\sffamily\tiny\color{ink}scannen zum Lösen}
\end{minipage}
\vspace{2mm}

{\raggedright\noindent\textcolor{ink}{\sffamily\textsc{Themen}}\enspace \themetag{reines Matt} \themetag{Mustermatt} \themetag{Idealmatt} \themetag{Spiegelmatt} \themetag{Switchback} \themetag{ein Stein zieht} \themetag{ohne Schach}\par}

\clearpage
\phantomsection\label{hm:KBBNvkb}
\begin{tcolorbox}[enhanced,colback=ink,colframe=ink,arc=0pt,boxrule=0pt,left=4mm,right=4mm,top=2.5mm,bottom=2.5mm]
\color{white}\sffamily{\LARGE\bfseries KBBNvkb}\hfill{\Large h\#6.5}\\[0.5mm]
{\small 6 Steine\quad Weiß am Zug}\hfill{\small Nr.~608\qquad\href{https://helpman.komtera.lt/?fen=N1B5/5b1k/8/8/B7/8/8/K7&moves=6.5}{\color{white}\bfseries lösen\,\ensuremath{\nearrow}}}
\end{tcolorbox}
\vspace{2mm}
\noindent\begin{minipage}[t]{36mm}\vspace{0pt}\sffamily\small\raggedright
\hmfigure{längste eindeutige}{{\Large h\#6.5}}\hmfigure{eindeutig bei dieser Länge}{{\Large 7} Stellungen}\hmfigure{längstes Matt}{{\Large h\#10}}\hmfigure{in der Klasse}{15,5 Milliarden Stellungen,\\ 4,74 Milliarden mit Hilfsmatt}
\end{minipage}\hfill
\begin{minipage}[t]{80mm}\vspace{0pt}\begin{tikzpicture}\begin{axis}[width=72mm,height=42mm,ybar,bar width=1.6pt,xmin=-1,xmax=21,ymin=0,ymax=1.447,axis y line*=left,axis x line*=bottom,xlabel={\scriptsize Halbzüge bis Matt},ylabel={\scriptsize Stellungen (Mrd.)},tick label style={font=\tiny},label style={font=\scriptsize},enlarge x limits=false,ylabel shift=-3pt,xlabel shift=-3pt]\addplot[fill=ink!55,draw=none] coordinates {(0,0.001041) (1,0.003823) (2,0.01817) (3,0.09922) (4,0.3907) (5,0.7683) (6,1.316) (7,0.8673) (8,0.6648) (9,0.2848) (10,0.2091) (11,0.07004) (12,0.03429) (13,0.007057) (14,0.001846) (15,0.0002176) (16,2.609e-06) (18,5.32e-07) (20,7.21e-07)};\draw[densely dotted] (axis cs:13,0) -- (axis cs:13,1.447);\end{axis}\begin{axis}[width=72mm,height=42mm,axis y line*=right,axis x line=none,ymode=log,xmin=-1,xmax=21,ylabel={\scriptsize eindeutig (log)},tick label style={font=\tiny},label style={font=\scriptsize},ylabel shift=-3pt]\addplot[mark=*,mark size=0.7pt,draw=black,thin] coordinates {(0,1040776) (1,3509585) (2,13507686) (3,46811215) (4,100246162) (5,44503587) (6,27724209) (7,4271537) (8,2290551) (9,300208) (10,87671) (11,4927) (12,1515) (13,7)};\end{axis}\end{tikzpicture}\end{minipage}
\vspace{3mm}

\noindent\begin{minipage}[t]{60mm}\vspace{0pt}\sbox\hmboard{\chessboard[setfen=N1B5/5b1k/8/8/B7/8/8/K7 w - - 0 1]}\usebox\hmboard\\[1mm]
\makebox[\wd\hmboard]{\sffamily\hspace{3mm}\textbf{h\#6.5}\hfill(4+2)}\end{minipage}\hfill
\begin{minipage}[t]{52mm}\vspace{0pt}
\textcolor{ink}{\sffamily\textsc{Lösung}}\par\smallskip\small
\begin{tabular}{@{}r@{\enspace}l@{\enspace}l@{}}
\textbf{1.} & \ldots & \symknight{}c7\\
\textbf{2.} & \symking{}g6 & \symknight{}e6\\
\textbf{3.} & \symking{}f5 & \symknight{}d4+\\
\textbf{4.} & \symking{}e4 & \symknight{}b3\\
\textbf{5.} & \symking{}d3 & \symking{}a2\\
\textbf{6.} & \symking{}c2 & \symking{}a3\\
\textbf{7.} & \symking{}b1 & \symbishop{}f5\#\\
\end{tabular}\par\vspace{3mm}
\qrcode[height=12mm]{https://helpman.komtera.lt/?fen=N1B5/5b1k/8/8/B7/8/8/K7&moves=6.5}\\{\sffamily\tiny\color{ink}scannen zum Lösen}
\end{minipage}
\vspace{2mm}

{\raggedright\noindent\textcolor{ink}{\sffamily\textsc{Themen}}\enspace \themetag{reines Matt} \themetag{Spiegelmatt} \themetag{ein Stein zieht} \themetag{ohne Schlag}\par}

\clearpage
\phantomsection\label{hm:KBBvkqq}
\begin{tcolorbox}[enhanced,colback=ink,colframe=ink,arc=0pt,boxrule=0pt,left=4mm,right=4mm,top=2.5mm,bottom=2.5mm]
\color{white}\sffamily{\LARGE\bfseries KBBvkqq}\hfill{\Large h\#6.5}\\[0.5mm]
{\small 6 Steine\quad Weiß am Zug}\hfill{\small Nr.~609\qquad\href{https://helpman.komtera.lt/?fen=7k/8/7B/8/1q6/8/3q4/KB6&moves=6.5}{\color{white}\bfseries lösen\,\ensuremath{\nearrow}}}
\end{tcolorbox}
\vspace{2mm}
\noindent\begin{minipage}[t]{36mm}\vspace{0pt}\sffamily\small\raggedright
\hmfigure{längste eindeutige}{{\Large h\#6.5}}\hmfigure{eindeutig bei dieser Länge}{{\Large 273} Stellungen}\hmfigure{längstes Matt}{{\Large h\#7.5}}\hmfigure{in der Klasse}{15,5 Milliarden Stellungen,\\ 888 Millionen mit Hilfsmatt}
\end{minipage}\hfill
\begin{minipage}[t]{80mm}\vspace{0pt}\begin{tikzpicture}\begin{axis}[width=72mm,height=42mm,ybar,bar width=1.6pt,xmin=-1,xmax=16,ymin=0,ymax=244.4,axis y line*=left,axis x line*=bottom,xlabel={\scriptsize Halbzüge bis Matt},ylabel={\scriptsize Stellungen (Mio.)},tick label style={font=\tiny},label style={font=\scriptsize},enlarge x limits=false,ylabel shift=-3pt,xlabel shift=-3pt]\addplot[fill=ink!55,draw=none] coordinates {(0,0.02257) (1,0.1027) (2,0.6838) (3,6.059) (4,18.45) (5,83.11) (6,109.2) (7,222.2) (8,138.2) (9,184.8) (10,63.76) (11,49.63) (12,8.263) (13,3.615) (14,0.2962) (15,0.03611)};\draw[densely dotted] (axis cs:13,0) -- (axis cs:13,244.4);\end{axis}\begin{axis}[width=72mm,height=42mm,axis y line*=right,axis x line=none,ymode=log,xmin=-1,xmax=16,ylabel={\scriptsize eindeutig (log)},tick label style={font=\tiny},label style={font=\scriptsize},ylabel shift=-3pt]\addplot[mark=*,mark size=0.7pt,draw=black,thin] coordinates {(0,22574) (1,102523) (2,330671) (3,2001671) (4,2702266) (5,7304193) (6,2777607) (7,3365179) (8,860697) (9,669226) (10,59663) (11,28314) (12,2948) (13,273)};\end{axis}\end{tikzpicture}\end{minipage}
\vspace{3mm}

\noindent\begin{minipage}[t]{60mm}\vspace{0pt}\sbox\hmboard{\chessboard[setfen=7k/8/7B/8/1q6/8/3q4/KB6 w - - 0 1]}\usebox\hmboard\\[1mm]
\makebox[\wd\hmboard]{\sffamily\hspace{3mm}\textbf{h\#6.5}\hfill(3+3)}\end{minipage}\hfill
\begin{minipage}[t]{52mm}\vspace{0pt}
\textcolor{ink}{\sffamily\textsc{Lösung}}\par\smallskip\small
\begin{tabular}{@{}r@{\enspace}l@{\enspace}l@{}}
\textbf{1.} & \ldots & \symbishop{}f5\\
\textbf{2.} & \symqueen{}g5 & \symking{}a2\\
\textbf{3.} & \symqueen{}b3+ & \symking{}xb3\\
\textbf{4.} & \symking{}g8 & \symking{}c4\\
\textbf{5.} & \symking{}f7 & \symking{}d5\\
\textbf{6.} & \symking{}e8 & \symking{}e6\\
\textbf{7.} & \symqueen{}d8 & \symbishop{}g6\#\\
\end{tabular}\par\vspace{3mm}
\qrcode[height=12mm]{https://helpman.komtera.lt/?fen=7k/8/7B/8/1q6/8/3q4/KB6&moves=6.5}\\{\sffamily\tiny\color{ink}scannen zum Lösen}
\end{minipage}
\vspace{2mm}

{\raggedright\noindent\textcolor{ink}{\sffamily\textsc{Themen}}\enspace \themetag{Selbstblock}\par}

\clearpage
\phantomsection\label{hm:KBBvkqr}
\begin{tcolorbox}[enhanced,colback=ink,colframe=ink,arc=0pt,boxrule=0pt,left=4mm,right=4mm,top=2.5mm,bottom=2.5mm]
\color{white}\sffamily{\LARGE\bfseries KBBvkqr}\hfill{\Large h\#6.5}\\[0.5mm]
{\small 6 Steine\quad Weiß am Zug}\hfill{\small Nr.~610\qquad\href{https://helpman.komtera.lt/?fen=7k/1r6/7B/8/8/8/3q4/KB6&moves=6.5}{\color{white}\bfseries lösen\,\ensuremath{\nearrow}}}
\end{tcolorbox}
\vspace{2mm}
\noindent\begin{minipage}[t]{36mm}\vspace{0pt}\sffamily\small\raggedright
\hmfigure{längste eindeutige}{{\Large h\#6.5}}\hmfigure{eindeutig bei dieser Länge}{{\Large 214} Stellungen}\hmfigure{längstes Matt}{{\Large h\#7.5}}\hmfigure{in der Klasse}{15,5 Milliarden Stellungen,\\ 1,93 Milliarden mit Hilfsmatt}
\end{minipage}\hfill
\begin{minipage}[t]{80mm}\vspace{0pt}\begin{tikzpicture}\begin{axis}[width=72mm,height=42mm,ybar,bar width=1.6pt,xmin=-1,xmax=16,ymin=0,ymax=507.6,axis y line*=left,axis x line*=bottom,xlabel={\scriptsize Halbzüge bis Matt},ylabel={\scriptsize Stellungen (Mio.)},tick label style={font=\tiny},label style={font=\scriptsize},enlarge x limits=false,ylabel shift=-3pt,xlabel shift=-3pt]\addplot[fill=ink!55,draw=none] coordinates {(0,0.07141) (1,0.3199) (2,1.89) (3,15.58) (4,48.21) (5,192.7) (6,274.4) (7,461.4) (8,330.4) (9,355.6) (10,141.6) (11,85.57) (12,17.98) (13,5.769) (14,0.629) (15,0.03623)};\draw[densely dotted] (axis cs:13,0) -- (axis cs:13,507.6);\end{axis}\begin{axis}[width=72mm,height=42mm,axis y line*=right,axis x line=none,ymode=log,xmin=-1,xmax=16,ylabel={\scriptsize eindeutig (log)},tick label style={font=\tiny},label style={font=\scriptsize},ylabel shift=-3pt]\addplot[mark=*,mark size=0.7pt,draw=black,thin] coordinates {(0,71410) (1,319383) (2,983163) (3,5547539) (4,8103025) (5,17555372) (6,7514460) (7,6859595) (8,2017248) (9,1173731) (10,121901) (11,39106) (12,3071) (13,214)};\end{axis}\end{tikzpicture}\end{minipage}
\vspace{3mm}

\noindent\begin{minipage}[t]{60mm}\vspace{0pt}\sbox\hmboard{\chessboard[setfen=7k/1r6/7B/8/8/8/3q4/KB6 w - - 0 1]}\usebox\hmboard\\[1mm]
\makebox[\wd\hmboard]{\sffamily\hspace{3mm}\textbf{h\#6.5}\hfill(3+3)}\end{minipage}\hfill
\begin{minipage}[t]{52mm}\vspace{0pt}
\textcolor{ink}{\sffamily\textsc{Lösung}}\par\smallskip\small
\begin{tabular}{@{}r@{\enspace}l@{\enspace}l@{}}
\textbf{1.} & \ldots & \symbishop{}f5\\
\textbf{2.} & \symqueen{}g5 & \symking{}a2\\
\textbf{3.} & \symrook{}b3 & \symking{}xb3\\
\textbf{4.} & \symking{}g8 & \symking{}c4\\
\textbf{5.} & \symking{}f7 & \symking{}d5\\
\textbf{6.} & \symking{}e8 & \symking{}e6\\
\textbf{7.} & \symqueen{}d8 & \symbishop{}g6\#\\
\end{tabular}\par\vspace{3mm}
\qrcode[height=12mm]{https://helpman.komtera.lt/?fen=7k/1r6/7B/8/8/8/3q4/KB6&moves=6.5}\\{\sffamily\tiny\color{ink}scannen zum Lösen}
\end{minipage}
\vspace{2mm}

{\raggedright\noindent\textcolor{ink}{\sffamily\textsc{Themen}}\enspace \themetag{Selbstblock} \themetag{ohne Schach}\par}

\clearpage
\phantomsection\label{hm:KBBvkrr}
\begin{tcolorbox}[enhanced,colback=ink,colframe=ink,arc=0pt,boxrule=0pt,left=4mm,right=4mm,top=2.5mm,bottom=2.5mm]
\color{white}\sffamily{\LARGE\bfseries KBBvkrr}\hfill{\Large h\#6.5}\\[0.5mm]
{\small 6 Steine\quad Weiß am Zug}\hfill{\small Nr.~611\qquad\href{https://helpman.komtera.lt/?fen=7k/8/7B/8/8/8/6r1/KB4r1&moves=6.5}{\color{white}\bfseries lösen\,\ensuremath{\nearrow}}}
\end{tcolorbox}
\vspace{2mm}
\noindent\begin{minipage}[t]{36mm}\vspace{0pt}\sffamily\small\raggedright
\hmfigure{längste eindeutige}{{\Large h\#6.5}}\hmfigure{eindeutig bei dieser Länge}{{\Large 52} Stellungen}\hmfigure{längstes Matt}{{\Large h\#7.5}}\hmfigure{in der Klasse}{15,5 Milliarden Stellungen,\\ 1,05 Milliarden mit Hilfsmatt}
\end{minipage}\hfill
\begin{minipage}[t]{80mm}\vspace{0pt}\begin{tikzpicture}\begin{axis}[width=72mm,height=42mm,ybar,bar width=1.6pt,xmin=-1,xmax=16,ymin=0,ymax=257,axis y line*=left,axis x line*=bottom,xlabel={\scriptsize Halbzüge bis Matt},ylabel={\scriptsize Stellungen (Mio.)},tick label style={font=\tiny},label style={font=\scriptsize},enlarge x limits=false,ylabel shift=-3pt,xlabel shift=-3pt]\addplot[fill=ink!55,draw=none] coordinates {(0,0.05437) (1,0.2403) (2,1.235) (3,9.519) (4,30.23) (5,107.8) (6,168.5) (7,233.6) (8,195.6) (9,169.7) (10,78.76) (11,37.93) (12,9.92) (13,2.443) (14,0.3471) (15,0.008987)};\draw[densely dotted] (axis cs:13,0) -- (axis cs:13,257);\end{axis}\begin{axis}[width=72mm,height=42mm,axis y line*=right,axis x line=none,ymode=log,xmin=-1,xmax=16,ylabel={\scriptsize eindeutig (log)},tick label style={font=\tiny},label style={font=\scriptsize},ylabel shift=-3pt]\addplot[mark=*,mark size=0.7pt,draw=black,thin] coordinates {(0,54368) (1,239763) (2,690893) (3,3644136) (4,5968236) (5,9840388) (6,5123357) (7,3325784) (8,1193351) (9,507151) (10,62533) (11,14007) (12,2326) (13,52)};\end{axis}\end{tikzpicture}\end{minipage}
\vspace{3mm}

\noindent\begin{minipage}[t]{60mm}\vspace{0pt}\sbox\hmboard{\chessboard[setfen=7k/8/7B/8/8/8/6r1/KB4r1 w - - 0 1]}\usebox\hmboard\\[1mm]
\makebox[\wd\hmboard]{\sffamily\hspace{3mm}\textbf{h\#6.5}\hfill(3+3)}\end{minipage}\hfill
\begin{minipage}[t]{52mm}\vspace{0pt}
\textcolor{ink}{\sffamily\textsc{Lösung}}\par\smallskip\small
\begin{tabular}{@{}r@{\enspace}l@{\enspace}l@{}}
\textbf{1.} & \ldots & \symbishop{}f8\\
\textbf{2.} & \symrook{}g7 & \symking{}b2\\
\textbf{3.} & \symrook{}1g6 & \symking{}c3\\
\textbf{4.} & \symking{}h7 & \symking{}d4\\
\textbf{5.} & \symking{}h6 & \symking{}e5\\
\textbf{6.} & \symrook{}g5+ & \symking{}f6\\
\textbf{7.} & \symrook{}h5 & \symbishop{}xg7\#\\
\end{tabular}\par\vspace{3mm}
\qrcode[height=12mm]{https://helpman.komtera.lt/?fen=7k/8/7B/8/8/8/6r1/KB4r1&moves=6.5}\\{\sffamily\tiny\color{ink}scannen zum Lösen}
\end{minipage}
\vspace{2mm}

{\raggedright\noindent\textcolor{ink}{\sffamily\textsc{Themen}}\enspace \themetag{Satzspiel} \themetag{Selbstblock} \themetag{Maslar}\par}

\clearpage
\phantomsection\label{hm:KBNPvkb}
\begin{tcolorbox}[enhanced,colback=ink,colframe=ink,arc=0pt,boxrule=0pt,left=4mm,right=4mm,top=2.5mm,bottom=2.5mm]
\color{white}\sffamily{\LARGE\bfseries KBNPvkb}\hfill{\Large h\#6.5}\\[0.5mm]
{\small 6 Steine\quad Weiß am Zug}\hfill{\small Nr.~612\qquad\href{https://helpman.komtera.lt/?fen=8/8/8/8/8/2b5/B6P/KN3k2&moves=6.5}{\color{white}\bfseries lösen\,\ensuremath{\nearrow}}}
\end{tcolorbox}
\vspace{2mm}
\noindent\begin{minipage}[t]{36mm}\vspace{0pt}\sffamily\small\raggedright
\hmfigure{längste eindeutige}{{\Large h\#6.5}}\hmfigure{eindeutig bei dieser Länge}{{\Large 196} Stellungen}\hmfigure{längstes Matt}{{\Large h\#8}}\hmfigure{in der Klasse}{45,4 Milliarden Stellungen,\\ 30 Milliarden mit Hilfsmatt}
\end{minipage}\hfill
\begin{minipage}[t]{80mm}\vspace{0pt}\begin{tikzpicture}\begin{axis}[width=72mm,height=42mm,ybar,bar width=1.6pt,xmin=-1,xmax=17,ymin=0,ymax=7.142,axis y line*=left,axis x line*=bottom,xlabel={\scriptsize Halbzüge bis Matt},ylabel={\scriptsize Stellungen (Mrd.)},tick label style={font=\tiny},label style={font=\scriptsize},enlarge x limits=false,ylabel shift=-3pt,xlabel shift=-3pt]\addplot[fill=ink!55,draw=none] coordinates {(0,0.00134) (1,0.0117) (2,0.04728) (3,0.4018) (4,1.345) (5,3.514) (6,5.602) (7,5.92) (8,6.493) (9,3.707) (10,2.205) (11,0.584) (12,0.148) (13,0.01158) (14,0.0008681) (15,1.448e-05) (16,3.72e-07)};\draw[densely dotted] (axis cs:13,0) -- (axis cs:13,7.142);\end{axis}\begin{axis}[width=72mm,height=42mm,axis y line*=right,axis x line=none,ymode=log,xmin=-1,xmax=17,ylabel={\scriptsize eindeutig (log)},tick label style={font=\tiny},label style={font=\scriptsize},ylabel shift=-3pt]\addplot[mark=*,mark size=0.7pt,draw=black,thin] coordinates {(0,1339884) (1,6846358) (2,25300288) (3,158432716) (4,315237594) (5,299053780) (6,236019340) (7,97564521) (8,43251865) (9,8302908) (10,1529668) (11,188665) (12,12841) (13,196)};\end{axis}\end{tikzpicture}\end{minipage}
\vspace{3mm}

\noindent\begin{minipage}[t]{60mm}\vspace{0pt}\sbox\hmboard{\chessboard[setfen=8/8/8/8/8/2b5/B6P/KN3k2 w - - 0 1]}\usebox\hmboard\\[1mm]
\makebox[\wd\hmboard]{\sffamily\hspace{3mm}\textbf{h\#6.5}\hfill(4+2)}\end{minipage}\hfill
\begin{minipage}[t]{52mm}\vspace{0pt}
\textcolor{ink}{\sffamily\textsc{Lösung}}\par\smallskip\small
\begin{tabular}{@{}r@{\enspace}l@{\enspace}l@{}}
\textbf{1.} & \ldots & \symknight{}xc3\\
\textbf{2.} & \symking{}f2 & h4\\
\textbf{3.} & \symking{}e3 & h5\\
\textbf{4.} & \symking{}d4 & h6\\
\textbf{5.} & \symking{}c5 & h7\\
\textbf{6.} & \symking{}b4 & h8=\symqueen{}\\
\textbf{7.} & \symking{}a3 & \symqueen{}f8\#\\
\end{tabular}\par\vspace{3mm}
\qrcode[height=12mm]{https://helpman.komtera.lt/?fen=8/8/8/8/8/2b5/B6P/KN3k2&moves=6.5}\\{\sffamily\tiny\color{ink}scannen zum Lösen}
\end{minipage}
\vspace{2mm}

{\raggedright\noindent\textcolor{ink}{\sffamily\textsc{Themen}}\enspace \themetag{reines Matt} \themetag{Mustermatt} \themetag{Idealmatt} \themetag{Umwandlung} \themetag{Excelsior} \themetag{ein Stein zieht} \themetag{ohne Schach}\par\smallskip{\footnotesize\textit{Hinweis:} die Partei am Zug steht im Diagramm im Schach; die Lösung beginnt mit einem Schlag.}\par}

\clearpage
\phantomsection\label{hm:KBNvkrb}
\begin{tcolorbox}[enhanced,colback=ink,colframe=ink,arc=0pt,boxrule=0pt,left=4mm,right=4mm,top=2.5mm,bottom=2.5mm]
\color{white}\sffamily{\LARGE\bfseries KBNvkrb}\hfill{\Large h\#6.5}\\[0.5mm]
{\small 6 Steine\quad Weiß am Zug}\hfill{\small Nr.~613\qquad\href{https://helpman.komtera.lt/?fen=1r1N4/8/5b2/8/8/8/8/K1B4k&moves=6.5}{\color{white}\bfseries lösen\,\ensuremath{\nearrow}}}
\end{tcolorbox}
\vspace{2mm}
\noindent\begin{minipage}[t]{36mm}\vspace{0pt}\sffamily\small\raggedright
\hmfigure{längste eindeutige}{{\Large h\#6.5}}\hmfigure{eindeutig bei dieser Länge}{{\Large 440} Stellungen}\hmfigure{längstes Matt}{{\Large h\#10}}\hmfigure{in der Klasse}{15,5 Milliarden Stellungen,\\ 9,08 Milliarden mit Hilfsmatt}
\end{minipage}\hfill
\begin{minipage}[t]{80mm}\vspace{0pt}\begin{tikzpicture}\begin{axis}[width=72mm,height=42mm,ybar,bar width=1.6pt,xmin=-1,xmax=21,ymin=0,ymax=2.436,axis y line*=left,axis x line*=bottom,xlabel={\scriptsize Halbzüge bis Matt},ylabel={\scriptsize Stellungen (Mrd.)},tick label style={font=\tiny},label style={font=\scriptsize},enlarge x limits=false,ylabel shift=-3pt,xlabel shift=-3pt]\addplot[fill=ink!55,draw=none] coordinates {(0,0.0001189) (1,0.0004978) (2,0.002999) (3,0.02137) (4,0.1008) (5,0.4089) (6,1.058) (7,2.039) (8,2.214) (9,1.931) (10,0.7567) (11,0.4156) (12,0.1054) (13,0.02474) (14,0.001623) (15,0.0001141) (16,1.65e-06) (17,2.2e-08) (18,1.8e-08) (19,1e-09) (20,1.1e-08)};\draw[densely dotted] (axis cs:13,0) -- (axis cs:13,2.436);\end{axis}\begin{axis}[width=72mm,height=42mm,axis y line*=right,axis x line=none,ymode=log,xmin=-1,xmax=21,ylabel={\scriptsize eindeutig (log)},tick label style={font=\tiny},label style={font=\scriptsize},ylabel shift=-3pt]\addplot[mark=*,mark size=0.7pt,draw=black,thin] coordinates {(0,118924) (1,474745) (2,2165362) (3,10138785) (4,32243776) (5,54683976) (6,53586853) (7,34558504) (8,9820679) (9,4015698) (10,776183) (11,182462) (12,9610) (13,440)};\end{axis}\end{tikzpicture}\end{minipage}
\vspace{3mm}

\noindent\begin{minipage}[t]{60mm}\vspace{0pt}\sbox\hmboard{\chessboard[setfen=1r1N4/8/5b2/8/8/8/8/K1B4k w - - 0 1]}\usebox\hmboard\\[1mm]
\makebox[\wd\hmboard]{\sffamily\hspace{3mm}\textbf{h\#6.5}\hfill(3+3)}\end{minipage}\hfill
\begin{minipage}[t]{52mm}\vspace{0pt}
\textcolor{ink}{\sffamily\textsc{Lösung}}\par\smallskip\small
\begin{tabular}{@{}r@{\enspace}l@{\enspace}l@{}}
\textbf{1.} & \ldots & \symking{}a2\\
\textbf{2.} & \symking{}g2 & \symknight{}f7\\
\textbf{3.} & \symrook{}d8 & \symking{}b3\\
\textbf{4.} & \symrook{}d3+ & \symking{}c4\\
\textbf{5.} & \symbishop{}d4 & \symking{}d5\\
\textbf{6.} & \symking{}f3 & \symking{}e6\\
\textbf{7.} & \symking{}e4 & \symknight{}g5\#\\
\end{tabular}\par\vspace{3mm}
\qrcode[height=12mm]{https://helpman.komtera.lt/?fen=1r1N4/8/5b2/8/8/8/8/K1B4k&moves=6.5}\\{\sffamily\tiny\color{ink}scannen zum Lösen}
\end{minipage}
\vspace{2mm}

{\raggedright\noindent\textcolor{ink}{\sffamily\textsc{Themen}}\enspace \themetag{reines Matt} \themetag{Mustermatt} \themetag{Idealmatt} \themetag{Selbstblock} \themetag{ohne Schlag} \themetag{Umnov}\par\smallskip{\footnotesize\textit{Hinweis:} die Partei am Zug steht im Diagramm im Schach.}\par}

\clearpage
\phantomsection\label{hm:KNNPvkn}
\begin{tcolorbox}[enhanced,colback=ink,colframe=ink,arc=0pt,boxrule=0pt,left=4mm,right=4mm,top=2.5mm,bottom=2.5mm]
\color{white}\sffamily{\LARGE\bfseries KNNPvkn}\hfill{\Large h\#6.5}\\[0.5mm]
{\small 6 Steine\quad Weiß am Zug}\hfill{\small Nr.~614\qquad\href{https://helpman.komtera.lt/?fen=7k/8/8/8/8/P7/1Nn5/K1N5&moves=6.5}{\color{white}\bfseries lösen\,\ensuremath{\nearrow}}}
\end{tcolorbox}
\vspace{2mm}
\noindent\begin{minipage}[t]{36mm}\vspace{0pt}\sffamily\small\raggedright
\hmfigure{längste eindeutige}{{\Large h\#6.5}}\hmfigure{eindeutig bei dieser Länge}{{\Large 229} Stellungen}\hmfigure{längstes Matt}{{\Large h\#8}}\hmfigure{in der Klasse}{45,4 Milliarden Stellungen,\\ 15,8 Milliarden mit Hilfsmatt}
\end{minipage}\hfill
\begin{minipage}[t]{80mm}\vspace{0pt}\begin{tikzpicture}\begin{axis}[width=72mm,height=42mm,ybar,bar width=1.6pt,xmin=-1,xmax=17,ymin=0,ymax=4.042,axis y line*=left,axis x line*=bottom,xlabel={\scriptsize Halbzüge bis Matt},ylabel={\scriptsize Stellungen (Mrd.)},tick label style={font=\tiny},label style={font=\scriptsize},enlarge x limits=false,ylabel shift=-3pt,xlabel shift=-3pt]\addplot[fill=ink!55,draw=none] coordinates {(0,0.0003255) (1,0.005197) (2,0.01617) (3,0.1554) (4,0.4548) (5,1.437) (6,2.405) (7,2.873) (8,3.675) (9,2.462) (10,1.666) (11,0.4919) (12,0.1082) (13,0.00893) (14,0.0006167) (15,1.517e-05) (16,8e-09)};\draw[densely dotted] (axis cs:13,0) -- (axis cs:13,4.042);\end{axis}\begin{axis}[width=72mm,height=42mm,axis y line*=right,axis x line=none,ymode=log,xmin=-1,xmax=17,ylabel={\scriptsize eindeutig (log)},tick label style={font=\tiny},label style={font=\scriptsize},ylabel shift=-3pt]\addplot[mark=*,mark size=0.7pt,draw=black,thin] coordinates {(0,325474) (1,2261236) (2,6386516) (3,60632610) (4,103357148) (5,103839718) (6,92510877) (7,47547213) (8,27233296) (9,5478350) (10,879485) (11,103489) (12,3695) (13,229)};\end{axis}\end{tikzpicture}\end{minipage}
\vspace{3mm}

\noindent\begin{minipage}[t]{60mm}\vspace{0pt}\sbox\hmboard{\chessboard[setfen=7k/8/8/8/8/P7/1Nn5/K1N5 w - - 0 1]}\usebox\hmboard\\[1mm]
\makebox[\wd\hmboard]{\sffamily\hspace{3mm}\textbf{h\#6.5}\hfill(4+2)}\end{minipage}\hfill
\begin{minipage}[t]{52mm}\vspace{0pt}
\textcolor{ink}{\sffamily\textsc{Lösung}}\par\smallskip\small
\begin{tabular}{@{}r@{\enspace}l@{\enspace}l@{}}
\textbf{1.} & \ldots & \symking{}b1\\
\textbf{2.} & \symking{}g7 & a4\\
\textbf{3.} & \symking{}f6 & a5\\
\textbf{4.} & \symking{}e5 & a6\\
\textbf{5.} & \symking{}d4 & a7\\
\textbf{6.} & \symking{}c3 & a8=\symqueen{}\\
\textbf{7.} & \symknight{}d4 & \symqueen{}a5\#\\
\end{tabular}\par\vspace{3mm}
\qrcode[height=12mm]{https://helpman.komtera.lt/?fen=7k/8/8/8/8/P7/1Nn5/K1N5&moves=6.5}\\{\sffamily\tiny\color{ink}scannen zum Lösen}
\end{minipage}
\vspace{2mm}

{\raggedright\noindent\textcolor{ink}{\sffamily\textsc{Themen}}\enspace \themetag{Umwandlung} \themetag{Rundlauf} \themetag{Selbstblock} \themetag{ohne Schlag} \themetag{ohne Schach}\par\smallskip{\footnotesize\textit{Hinweis:} die Partei am Zug steht im Diagramm im Schach.}\par}

\clearpage
\phantomsection\label{hm:KNNPvkr}
\begin{tcolorbox}[enhanced,colback=ink,colframe=ink,arc=0pt,boxrule=0pt,left=4mm,right=4mm,top=2.5mm,bottom=2.5mm]
\color{white}\sffamily{\LARGE\bfseries KNNPvkr}\hfill{\Large h\#6.5}\\[0.5mm]
{\small 6 Steine\quad Weiß am Zug}\hfill{\small Nr.~615\qquad\href{https://helpman.komtera.lt/?fen=3N1N2/8/8/8/8/8/P6k/Kr6&moves=6.5}{\color{white}\bfseries lösen\,\ensuremath{\nearrow}}}
\end{tcolorbox}
\vspace{2mm}
\noindent\begin{minipage}[t]{36mm}\vspace{0pt}\sffamily\small\raggedright
\hmfigure{längste eindeutige}{{\Large h\#6.5}}\hmfigure{eindeutig bei dieser Länge}{{\Large 467} Stellungen}\hmfigure{längstes Matt}{{\Large h\#9}}\hmfigure{in der Klasse}{45,4 Milliarden Stellungen,\\ 14,7 Milliarden mit Hilfsmatt}
\end{minipage}\hfill
\begin{minipage}[t]{80mm}\vspace{0pt}\begin{tikzpicture}\begin{axis}[width=72mm,height=42mm,ybar,bar width=1.6pt,xmin=-1,xmax=19,ymin=0,ymax=3.348,axis y line*=left,axis x line*=bottom,xlabel={\scriptsize Halbzüge bis Matt},ylabel={\scriptsize Stellungen (Mrd.)},tick label style={font=\tiny},label style={font=\scriptsize},enlarge x limits=false,ylabel shift=-3pt,xlabel shift=-3pt]\addplot[fill=ink!55,draw=none] coordinates {(0,0.0002619) (1,0.003857) (2,0.01655) (3,0.163) (4,0.5794) (5,1.65) (6,2.507) (7,3.044) (8,2.97) (9,2.146) (10,1.129) (11,0.4112) (12,0.08413) (13,0.01206) (14,0.0005662) (15,4.654e-05) (16,1.4e-08) (17,2e-09) (18,1e-09)};\draw[densely dotted] (axis cs:13,0) -- (axis cs:13,3.348);\end{axis}\begin{axis}[width=72mm,height=42mm,axis y line*=right,axis x line=none,ymode=log,xmin=-1,xmax=19,ylabel={\scriptsize eindeutig (log)},tick label style={font=\tiny},label style={font=\scriptsize},ylabel shift=-3pt]\addplot[mark=*,mark size=0.7pt,draw=black,thin] coordinates {(0,261852) (1,1937844) (2,7369714) (3,64167704) (4,138666382) (5,139311687) (6,94141806) (7,48351092) (8,17523705) (9,4672125) (10,695437) (11,118368) (12,3908) (13,467)};\end{axis}\end{tikzpicture}\end{minipage}
\vspace{3mm}

\noindent\begin{minipage}[t]{60mm}\vspace{0pt}\sbox\hmboard{\chessboard[setfen=3N1N2/8/8/8/8/8/P6k/Kr6 w - - 0 1]}\usebox\hmboard\\[1mm]
\makebox[\wd\hmboard]{\sffamily\hspace{3mm}\textbf{h\#6.5}\hfill(4+2)}\end{minipage}\hfill
\begin{minipage}[t]{52mm}\vspace{0pt}
\textcolor{ink}{\sffamily\textsc{Lösung}}\par\smallskip\small
\begin{tabular}{@{}r@{\enspace}l@{\enspace}l@{}}
\textbf{1.} & \ldots & \symking{}xb1\\
\textbf{2.} & \symking{}g3 & a4\\
\textbf{3.} & \symking{}f4 & a5\\
\textbf{4.} & \symking{}e5 & a6\\
\textbf{5.} & \symking{}d6 & a7\\
\textbf{6.} & \symking{}c7 & \symknight{}de6+\\
\textbf{7.} & \symking{}c8 & a8=\symqueen{}\#\\
\end{tabular}\par\vspace{3mm}
\qrcode[height=12mm]{https://helpman.komtera.lt/?fen=3N1N2/8/8/8/8/8/P6k/Kr6&moves=6.5}\\{\sffamily\tiny\color{ink}scannen zum Lösen}
\end{minipage}
\vspace{2mm}

{\raggedright\noindent\textcolor{ink}{\sffamily\textsc{Themen}}\enspace \themetag{Spiegelmatt} \themetag{Umwandlung} \themetag{Excelsior} \themetag{ein Stein zieht}\par\smallskip{\footnotesize\textit{Hinweis:} die Partei am Zug steht im Diagramm im Schach; die Lösung beginnt mit einem Schlag.}\par}

\clearpage
\phantomsection\label{hm:KNNvkbb}
\begin{tcolorbox}[enhanced,colback=ink,colframe=ink,arc=0pt,boxrule=0pt,left=4mm,right=4mm,top=2.5mm,bottom=2.5mm]
\color{white}\sffamily{\LARGE\bfseries KNNvkbb}\hfill{\Large h\#6.5}\\[0.5mm]
{\small 6 Steine\quad Weiß am Zug}\hfill{\small Nr.~616\qquad\href{https://helpman.komtera.lt/?fen=2N3N1/8/8/2b5/3b4/8/6k1/K7&moves=6.5}{\color{white}\bfseries lösen\,\ensuremath{\nearrow}}}
\end{tcolorbox}
\vspace{2mm}
\noindent\begin{minipage}[t]{36mm}\vspace{0pt}\sffamily\small\raggedright
\hmfigure{längste eindeutige}{{\Large h\#6.5}}\hmfigure{eindeutig bei dieser Länge}{{\Large 215} Stellungen}\hmfigure{längstes Matt}{{\Large h\#10}}\hmfigure{in der Klasse}{15,5 Milliarden Stellungen,\\ 2,41 Milliarden mit Hilfsmatt}
\end{minipage}\hfill
\begin{minipage}[t]{80mm}\vspace{0pt}\begin{tikzpicture}\begin{axis}[width=72mm,height=42mm,ybar,bar width=1.6pt,xmin=-1,xmax=21,ymin=0,ymax=618.2,axis y line*=left,axis x line*=bottom,xlabel={\scriptsize Halbzüge bis Matt},ylabel={\scriptsize Stellungen (Mio.)},tick label style={font=\tiny},label style={font=\scriptsize},enlarge x limits=false,ylabel shift=-3pt,xlabel shift=-3pt]\addplot[fill=ink!55,draw=none] coordinates {(0,0.02336) (1,0.1148) (2,0.5307) (3,4.399) (4,18.27) (5,83.59) (6,204.4) (7,447.1) (8,561.7) (9,562) (10,302.8) (11,158.5) (12,50.44) (13,15.51) (14,2.956) (15,0.7386) (16,0.002445) (18,0.001492) (20,2.4e-05)};\draw[densely dotted] (axis cs:13,0) -- (axis cs:13,618.2);\end{axis}\begin{axis}[width=72mm,height=42mm,axis y line*=right,axis x line=none,ymode=log,xmin=-1,xmax=21,ylabel={\scriptsize eindeutig (log)},tick label style={font=\tiny},label style={font=\scriptsize},ylabel shift=-3pt]\addplot[mark=*,mark size=0.7pt,draw=black,thin] coordinates {(0,23356) (1,109968) (2,373303) (3,2180455) (4,6184448) (5,9670260) (6,8891578) (7,6234383) (8,2882733) (9,926104) (10,190162) (11,22512) (12,3119) (13,215)};\end{axis}\end{tikzpicture}\end{minipage}
\vspace{3mm}

\noindent\begin{minipage}[t]{60mm}\vspace{0pt}\sbox\hmboard{\chessboard[setfen=2N3N1/8/8/2b5/3b4/8/6k1/K7 w - - 0 1]}\usebox\hmboard\\[1mm]
\makebox[\wd\hmboard]{\sffamily\hspace{3mm}\textbf{h\#6.5}\hfill(3+3)}\end{minipage}\hfill
\begin{minipage}[t]{52mm}\vspace{0pt}
\textcolor{ink}{\sffamily\textsc{Lösung}}\par\smallskip\small
\begin{tabular}{@{}r@{\enspace}l@{\enspace}l@{}}
\textbf{1.} & \ldots & \symking{}a2\\
\textbf{2.} & \symking{}f3 & \symking{}b3\\
\textbf{3.} & \symking{}e4 & \symknight{}d6+\\
\textbf{4.} & \symking{}d5 & \symknight{}f6+\\
\textbf{5.} & \symking{}c6 & \symknight{}fe8\\
\textbf{6.} & \symking{}b6 & \symknight{}c7\\
\textbf{7.} & \symking{}a5 & \symknight{}c4\#\\
\end{tabular}\par\vspace{3mm}
\qrcode[height=12mm]{https://helpman.komtera.lt/?fen=2N3N1/8/8/2b5/3b4/8/6k1/K7&moves=6.5}\\{\sffamily\tiny\color{ink}scannen zum Lösen}
\end{minipage}
\vspace{2mm}

{\raggedright\noindent\textcolor{ink}{\sffamily\textsc{Themen}}\enspace \themetag{reines Matt} \themetag{Mustermatt} \themetag{Spiegelmatt} \themetag{ein Stein zieht} \themetag{ohne Schlag}\par\smallskip{\footnotesize\textit{Hinweis:} die Partei am Zug steht im Diagramm im Schach.}\par}

\clearpage
\phantomsection\label{hm:KNNvkbn}
\begin{tcolorbox}[enhanced,colback=ink,colframe=ink,arc=0pt,boxrule=0pt,left=4mm,right=4mm,top=2.5mm,bottom=2.5mm]
\color{white}\sffamily{\LARGE\bfseries KNNvkbn}\hfill{\Large h\#6.5}\\[0.5mm]
{\small 6 Steine\quad Weiß am Zug}\hfill{\small Nr.~617\qquad\href{https://helpman.komtera.lt/?fen=6k1/5b2/8/8/8/n7/8/NKN5&moves=6.5}{\color{white}\bfseries lösen\,\ensuremath{\nearrow}}}
\end{tcolorbox}
\vspace{2mm}
\noindent\begin{minipage}[t]{36mm}\vspace{0pt}\sffamily\small\raggedright
\hmfigure{längste eindeutige}{{\Large h\#6.5}}\hmfigure{eindeutig bei dieser Länge}{{\Large 12} Stellungen}\hmfigure{längstes Matt}{{\Large h\#9}}\hmfigure{in der Klasse}{15,5 Milliarden Stellungen,\\ 4,97 Milliarden mit Hilfsmatt}
\end{minipage}\hfill
\begin{minipage}[t]{80mm}\vspace{0pt}\begin{tikzpicture}\begin{axis}[width=72mm,height=42mm,ybar,bar width=1.6pt,xmin=-1,xmax=19,ymin=0,ymax=1.418,axis y line*=left,axis x line*=bottom,xlabel={\scriptsize Halbzüge bis Matt},ylabel={\scriptsize Stellungen (Mrd.)},tick label style={font=\tiny},label style={font=\scriptsize},enlarge x limits=false,ylabel shift=-3pt,xlabel shift=-3pt]\addplot[fill=ink!55,draw=none] coordinates {(0,5.439e-05) (1,0.0002752) (2,0.001202) (3,0.009843) (4,0.04222) (5,0.1803) (6,0.4679) (7,0.9383) (8,1.289) (9,1.167) (10,0.5871) (11,0.2421) (12,0.03458) (13,0.006016) (14,0.0001841) (15,1.782e-05) (16,1.03e-07) (18,1e-08)};\draw[densely dotted] (axis cs:13,0) -- (axis cs:13,1.418);\end{axis}\begin{axis}[width=72mm,height=42mm,axis y line*=right,axis x line=none,ymode=log,xmin=-1,xmax=19,ylabel={\scriptsize eindeutig (log)},tick label style={font=\tiny},label style={font=\scriptsize},ylabel shift=-3pt]\addplot[mark=*,mark size=0.7pt,draw=black,thin] coordinates {(0,54385) (1,264298) (2,856349) (3,4898290) (4,14947975) (5,19270166) (6,20067937) (7,11884555) (8,5601266) (9,1267820) (10,116980) (11,12521) (12,137) (13,12)};\end{axis}\end{tikzpicture}\end{minipage}
\vspace{3mm}

\noindent\begin{minipage}[t]{60mm}\vspace{0pt}\sbox\hmboard{\chessboard[setfen=6k1/5b2/8/8/8/n7/8/NKN5 w - - 0 1]}\usebox\hmboard\\[1mm]
\makebox[\wd\hmboard]{\sffamily\hspace{3mm}\textbf{h\#6.5}\hfill(3+3)}\end{minipage}\hfill
\begin{minipage}[t]{52mm}\vspace{0pt}
\textcolor{ink}{\sffamily\textsc{Lösung}}\par\smallskip\small
\begin{tabular}{@{}r@{\enspace}l@{\enspace}l@{}}
\textbf{1.} & \ldots & \symking{}b2\\
\textbf{2.} & \symking{}f8 & \symking{}xa3\\
\textbf{3.} & \symking{}e7 & \symking{}b2\\
\textbf{4.} & \symking{}d6 & \symking{}b1\\
\textbf{5.} & \symking{}c5 & \symknight{}cb3+\\
\textbf{6.} & \symking{}b4 & \symknight{}c5\\
\textbf{7.} & \symking{}a3 & \symknight{}c2\#\\
\end{tabular}\par\vspace{3mm}
\qrcode[height=12mm]{https://helpman.komtera.lt/?fen=6k1/5b2/8/8/8/n7/8/NKN5&moves=6.5}\\{\sffamily\tiny\color{ink}scannen zum Lösen}
\end{minipage}
\vspace{2mm}

{\raggedright\noindent\textcolor{ink}{\sffamily\textsc{Themen}}\enspace \themetag{reines Matt} \themetag{Mustermatt} \themetag{Spiegelmatt} \themetag{Switchback} \themetag{ein Stein zieht} \themetag{Kniest} \themetag{Umnov}\par\smallskip{\footnotesize\textit{Hinweis:} die Partei am Zug steht im Diagramm im Schach.}\par}

\clearpage
\phantomsection\label{hm:KNNvknn}
\begin{tcolorbox}[enhanced,colback=ink,colframe=ink,arc=0pt,boxrule=0pt,left=4mm,right=4mm,top=2.5mm,bottom=2.5mm]
\color{white}\sffamily{\LARGE\bfseries KNNvknn}\hfill{\Large h\#6.5}\\[0.5mm]
{\small 6 Steine\quad Weiß am Zug}\hfill{\small Nr.~618\qquad\href{https://helpman.komtera.lt/?fen=1k5n/8/8/8/8/6Nn/8/3K3N&moves=6.5}{\color{white}\bfseries lösen\,\ensuremath{\nearrow}}}
\end{tcolorbox}
\vspace{2mm}
\noindent\begin{minipage}[t]{36mm}\vspace{0pt}\sffamily\small\raggedright
\hmfigure{längste eindeutige}{{\Large h\#6.5}}\hmfigure{eindeutig bei dieser Länge}{{\Large 31} Stellungen}\hmfigure{längstes Matt}{{\Large h\#8}}\hmfigure{in der Klasse}{15,5 Milliarden Stellungen,\\ 2,54 Milliarden mit Hilfsmatt}
\end{minipage}\hfill
\begin{minipage}[t]{80mm}\vspace{0pt}\begin{tikzpicture}\begin{axis}[width=72mm,height=42mm,ybar,bar width=1.6pt,xmin=-1,xmax=17,ymin=0,ymax=704.2,axis y line*=left,axis x line*=bottom,xlabel={\scriptsize Halbzüge bis Matt},ylabel={\scriptsize Stellungen (Mio.)},tick label style={font=\tiny},label style={font=\scriptsize},enlarge x limits=false,ylabel shift=-3pt,xlabel shift=-3pt]\addplot[fill=ink!55,draw=none] coordinates {(0,0.0305) (1,0.1591) (2,0.6532) (3,5.31) (4,21.93) (5,88.81) (6,233.9) (7,441.9) (8,640.2) (9,578.5) (10,348) (11,152.6) (12,27.15) (13,4.777) (14,0.2449) (15,0.03131) (16,1e-05)};\draw[densely dotted] (axis cs:13,0) -- (axis cs:13,704.2);\end{axis}\begin{axis}[width=72mm,height=42mm,axis y line*=right,axis x line=none,ymode=log,xmin=-1,xmax=17,ylabel={\scriptsize eindeutig (log)},tick label style={font=\tiny},label style={font=\scriptsize},ylabel shift=-3pt]\addplot[mark=*,mark size=0.7pt,draw=black,thin] coordinates {(0,30501) (1,153490) (2,466489) (3,2641404) (4,7018048) (5,7862676) (6,8319502) (7,4972975) (8,3003768) (9,723217) (10,86537) (11,8953) (12,181) (13,31)};\end{axis}\end{tikzpicture}\end{minipage}
\vspace{3mm}

\noindent\begin{minipage}[t]{60mm}\vspace{0pt}\sbox\hmboard{\chessboard[setfen=1k5n/8/8/8/8/6Nn/8/3K3N w - - 0 1]}\usebox\hmboard\\[1mm]
\makebox[\wd\hmboard]{\sffamily\hspace{3mm}\textbf{h\#6.5}\hfill(3+3)}\end{minipage}\hfill
\begin{minipage}[t]{52mm}\vspace{0pt}
\textcolor{ink}{\sffamily\textsc{Lösung}}\par\smallskip\small
\begin{tabular}{@{}r@{\enspace}l@{\enspace}l@{}}
\textbf{1.} & \ldots & \symking{}e2\\
\textbf{2.} & \symking{}c7 & \symking{}f3\\
\textbf{3.} & \symking{}d6 & \symking{}g4\\
\textbf{4.} & \symking{}e5 & \symking{}h5\\
\textbf{5.} & \symking{}f4 & \symking{}h6\\
\textbf{6.} & \symking{}g4 & \symknight{}f2+\\
\textbf{7.} & \symking{}h4 & \symknight{}f5\#\\
\end{tabular}\par\vspace{3mm}
\qrcode[height=12mm]{https://helpman.komtera.lt/?fen=1k5n/8/8/8/8/6Nn/8/3K3N&moves=6.5}\\{\sffamily\tiny\color{ink}scannen zum Lösen}
\end{minipage}
\vspace{2mm}

{\raggedright\noindent\textcolor{ink}{\sffamily\textsc{Themen}}\enspace \themetag{ein Stein zieht} \themetag{ohne Schlag}\par}

\clearpage
\phantomsection\label{hm:KNNvknp}
\begin{tcolorbox}[enhanced,colback=ink,colframe=ink,arc=0pt,boxrule=0pt,left=4mm,right=4mm,top=2.5mm,bottom=2.5mm]
\color{white}\sffamily{\LARGE\bfseries KNNvknp}\hfill{\Large h\#6.5}\\[0.5mm]
{\small 6 Steine\quad Weiß am Zug}\hfill{\small Nr.~619\qquad\href{https://helpman.komtera.lt/?fen=7k/8/8/8/8/pn6/1N6/K1N5&moves=6.5}{\color{white}\bfseries lösen\,\ensuremath{\nearrow}}}
\end{tcolorbox}
\vspace{2mm}
\noindent\begin{minipage}[t]{36mm}\vspace{0pt}\sffamily\small\raggedright
\hmfigure{längste eindeutige}{{\Large h\#6.5}}\hmfigure{eindeutig bei dieser Länge}{{\Large 120} Stellungen}\hmfigure{längstes Matt}{{\Large h\#9}}\hmfigure{in der Klasse}{45,4 Milliarden Stellungen,\\ 15,7 Milliarden mit Hilfsmatt}
\end{minipage}\hfill
\begin{minipage}[t]{80mm}\vspace{0pt}\begin{tikzpicture}\begin{axis}[width=72mm,height=42mm,ybar,bar width=1.6pt,xmin=-1,xmax=19,ymin=0,ymax=4.172,axis y line*=left,axis x line*=bottom,xlabel={\scriptsize Halbzüge bis Matt},ylabel={\scriptsize Stellungen (Mrd.)},tick label style={font=\tiny},label style={font=\scriptsize},enlarge x limits=false,ylabel shift=-3pt,xlabel shift=-3pt]\addplot[fill=ink!55,draw=none] coordinates {(0,0.0002128) (1,0.001071) (2,0.003843) (3,0.03013) (4,0.1128) (5,0.4669) (6,1.209) (7,2.374) (8,3.793) (9,3.547) (10,2.684) (11,1.179) (12,0.2916) (13,0.0418) (14,0.003698) (15,0.0002541) (16,1.139e-06) (18,1.28e-07)};\draw[densely dotted] (axis cs:13,0) -- (axis cs:13,4.172);\end{axis}\begin{axis}[width=72mm,height=42mm,axis y line*=right,axis x line=none,ymode=log,xmin=-1,xmax=19,ylabel={\scriptsize eindeutig (log)},tick label style={font=\tiny},label style={font=\scriptsize},ylabel shift=-3pt]\addplot[mark=*,mark size=0.7pt,draw=black,thin] coordinates {(0,212793) (1,1020002) (2,2739885) (3,15154420) (4,38622767) (5,45246269) (6,55915366) (7,32645568) (8,25008007) (9,5418388) (10,1068039) (11,75201) (12,2751) (13,120)};\end{axis}\end{tikzpicture}\end{minipage}
\vspace{3mm}

\noindent\begin{minipage}[t]{60mm}\vspace{0pt}\sbox\hmboard{\chessboard[setfen=7k/8/8/8/8/pn6/1N6/K1N5 w - - 0 1]}\usebox\hmboard\\[1mm]
\makebox[\wd\hmboard]{\sffamily\hspace{3mm}\textbf{h\#6.5}\hfill(3+3)}\end{minipage}\hfill
\begin{minipage}[t]{52mm}\vspace{0pt}
\textcolor{ink}{\sffamily\textsc{Lösung}}\par\smallskip\small
\begin{tabular}{@{}r@{\enspace}l@{\enspace}l@{}}
\textbf{1.} & \ldots & \symking{}b1\\
\textbf{2.} & \symking{}g7 & \symking{}c2\\
\textbf{3.} & \symking{}f6 & \symking{}c3\\
\textbf{4.} & \symking{}e5 & \symking{}b4\\
\textbf{5.} & \symking{}d4 & \symking{}b5\\
\textbf{6.} & \symking{}c3 & \symking{}b6\\
\textbf{7.} & \symking{}b4 & \symknight{}a2\#\\
\end{tabular}\par\vspace{3mm}
\qrcode[height=12mm]{https://helpman.komtera.lt/?fen=7k/8/8/8/8/pn6/1N6/K1N5&moves=6.5}\\{\sffamily\tiny\color{ink}scannen zum Lösen}
\end{minipage}
\vspace{2mm}

{\raggedright\noindent\textcolor{ink}{\sffamily\textsc{Themen}}\enspace \themetag{reines Matt} \themetag{Mustermatt} \themetag{Idealmatt} \themetag{ein Stein zieht} \themetag{ohne Schlag} \themetag{ohne Schach}\par\smallskip{\footnotesize\textit{Hinweis:} die Partei am Zug steht im Diagramm im Schach.}\par}

\clearpage
\phantomsection\label{hm:KNNvkqn}
\begin{tcolorbox}[enhanced,colback=ink,colframe=ink,arc=0pt,boxrule=0pt,left=4mm,right=4mm,top=2.5mm,bottom=2.5mm]
\color{white}\sffamily{\LARGE\bfseries KNNvkqn}\hfill{\Large h\#6.5}\\[0.5mm]
{\small 6 Steine\quad Weiß am Zug}\hfill{\small Nr.~620\qquad\href{https://helpman.komtera.lt/?fen=1N3N2/8/8/8/n6k/8/q7/K7&moves=6.5}{\color{white}\bfseries lösen\,\ensuremath{\nearrow}}}
\end{tcolorbox}
\vspace{2mm}
\noindent\begin{minipage}[t]{36mm}\vspace{0pt}\sffamily\small\raggedright
\hmfigure{längste eindeutige}{{\Large h\#6.5}}\hmfigure{eindeutig bei dieser Länge}{{\Large 229} Stellungen}\hmfigure{längstes Matt}{{\Large h\#9}}\hmfigure{in der Klasse}{15,5 Milliarden Stellungen,\\ 4,42 Milliarden mit Hilfsmatt}
\end{minipage}\hfill
\begin{minipage}[t]{80mm}\vspace{0pt}\begin{tikzpicture}\begin{axis}[width=72mm,height=42mm,ybar,bar width=1.6pt,xmin=-1,xmax=19,ymin=0,ymax=1.387,axis y line*=left,axis x line*=bottom,xlabel={\scriptsize Halbzüge bis Matt},ylabel={\scriptsize Stellungen (Mrd.)},tick label style={font=\tiny},label style={font=\scriptsize},enlarge x limits=false,ylabel shift=-3pt,xlabel shift=-3pt]\addplot[fill=ink!55,draw=none] coordinates {(0,3.03e-05) (1,0.000157) (2,0.0006651) (3,0.005886) (4,0.02331) (5,0.1165) (6,0.304) (7,0.786) (8,0.9991) (9,1.261) (10,0.5154) (11,0.3554) (12,0.03419) (13,0.01384) (14,0.0002609) (15,0.0001243) (16,2.52e-07) (17,1.54e-07) (18,1.2e-08)};\draw[densely dotted] (axis cs:13,0) -- (axis cs:13,1.387);\end{axis}\begin{axis}[width=72mm,height=42mm,axis y line*=right,axis x line=none,ymode=log,xmin=-1,xmax=19,ylabel={\scriptsize eindeutig (log)},tick label style={font=\tiny},label style={font=\scriptsize},ylabel shift=-3pt]\addplot[mark=*,mark size=0.7pt,draw=black,thin] coordinates {(0,30295) (1,153063) (2,435193) (3,2715646) (4,7096276) (5,13262233) (6,13472664) (7,14261970) (8,4764669) (9,2541036) (10,145741) (11,50531) (12,383) (13,229)};\end{axis}\end{tikzpicture}\end{minipage}
\vspace{3mm}

\noindent\begin{minipage}[t]{60mm}\vspace{0pt}\sbox\hmboard{\chessboard[setfen=1N3N2/8/8/8/n6k/8/q7/K7 w - - 0 1]}\usebox\hmboard\\[1mm]
\makebox[\wd\hmboard]{\sffamily\hspace{3mm}\textbf{h\#6.5}\hfill(3+3)}\end{minipage}\hfill
\begin{minipage}[t]{52mm}\vspace{0pt}
\textcolor{ink}{\sffamily\textsc{Lösung}}\par\smallskip\small
\begin{tabular}{@{}r@{\enspace}l@{\enspace}l@{}}
\textbf{1.} & \ldots & \symking{}xa2\\
\textbf{2.} & \symking{}g5 & \symking{}b3\\
\textbf{3.} & \symking{}f6 & \symking{}c4\\
\textbf{4.} & \symking{}e7 & \symking{}d5\\
\textbf{5.} & \symking{}d8 & \symking{}e6\\
\textbf{6.} & \symknight{}b6 & \symking{}f7\\
\textbf{7.} & \symknight{}c8 & \symknight{}e6\#\\
\end{tabular}\par\vspace{3mm}
\qrcode[height=12mm]{https://helpman.komtera.lt/?fen=1N3N2/8/8/8/n6k/8/q7/K7&moves=6.5}\\{\sffamily\tiny\color{ink}scannen zum Lösen}
\end{minipage}
\vspace{2mm}

{\raggedright\noindent\textcolor{ink}{\sffamily\textsc{Themen}}\enspace \themetag{reines Matt} \themetag{Mustermatt} \themetag{Idealmatt} \themetag{Selbstblock} \themetag{ohne Schach}\par\smallskip{\footnotesize\textit{Hinweis:} die Partei am Zug steht im Diagramm im Schach; die Lösung beginnt mit einem Schlag.}\par}

\clearpage
\phantomsection\label{hm:KNNvkqq}
\begin{tcolorbox}[enhanced,colback=ink,colframe=ink,arc=0pt,boxrule=0pt,left=4mm,right=4mm,top=2.5mm,bottom=2.5mm]
\color{white}\sffamily{\LARGE\bfseries KNNvkqq}\hfill{\Large h\#6.5}\\[0.5mm]
{\small 6 Steine\quad Weiß am Zug}\hfill{\small Nr.~621\qquad\href{https://helpman.komtera.lt/?fen=1q1N4/8/8/N7/8/2q5/8/K6k&moves=6.5}{\color{white}\bfseries lösen\,\ensuremath{\nearrow}}}
\end{tcolorbox}
\vspace{2mm}
\noindent\begin{minipage}[t]{36mm}\vspace{0pt}\sffamily\small\raggedright
\hmfigure{längste eindeutige}{{\Large h\#6.5}}\hmfigure{eindeutig bei dieser Länge}{{\Large 3\,470} Stellungen}\hmfigure{längstes Matt}{{\Large h\#8.5}}\hmfigure{in der Klasse}{15,5 Milliarden Stellungen,\\ 1,9 Milliarden mit Hilfsmatt}
\end{minipage}\hfill
\begin{minipage}[t]{80mm}\vspace{0pt}\begin{tikzpicture}\begin{axis}[width=72mm,height=42mm,ybar,bar width=1.6pt,xmin=-1,xmax=18,ymin=0,ymax=646.8,axis y line*=left,axis x line*=bottom,xlabel={\scriptsize Halbzüge bis Matt},ylabel={\scriptsize Stellungen (Mio.)},tick label style={font=\tiny},label style={font=\scriptsize},enlarge x limits=false,ylabel shift=-3pt,xlabel shift=-3pt]\addplot[fill=ink!55,draw=none] coordinates {(0,0.007067) (1,0.03611) (2,0.1494) (3,1.459) (4,4.62) (5,29.54) (6,67.81) (7,256.9) (8,297.4) (9,588) (10,250.3) (11,319.4) (12,46.03) (13,40.38) (14,0.7848) (15,1.065) (16,1.5e-05) (17,0.000789)};\draw[densely dotted] (axis cs:13,0) -- (axis cs:13,646.8);\end{axis}\begin{axis}[width=72mm,height=42mm,axis y line*=right,axis x line=none,ymode=log,xmin=-1,xmax=18,ylabel={\scriptsize eindeutig (log)},tick label style={font=\tiny},label style={font=\scriptsize},ylabel shift=-3pt]\addplot[mark=*,mark size=0.7pt,draw=black,thin] coordinates {(0,7067) (1,35536) (2,83002) (3,575827) (4,1115742) (5,3703740) (6,4289478) (7,8505679) (8,2629972) (9,2998224) (10,349498) (11,272118) (12,7062) (13,3470)};\end{axis}\end{tikzpicture}\end{minipage}
\vspace{3mm}

\noindent\begin{minipage}[t]{60mm}\vspace{0pt}\sbox\hmboard{\chessboard[setfen=1q1N4/8/8/N7/8/2q5/8/K6k w - - 0 1]}\usebox\hmboard\\[1mm]
\makebox[\wd\hmboard]{\sffamily\hspace{3mm}\textbf{h\#6.5}\hfill(3+3)}\end{minipage}\hfill
\begin{minipage}[t]{52mm}\vspace{0pt}
\textcolor{ink}{\sffamily\textsc{Lösung}}\par\smallskip\small
\begin{tabular}{@{}r@{\enspace}l@{\enspace}l@{}}
\textbf{1.} & \ldots & \symking{}a2\\
\textbf{2.} & \symqueen{}cg3 & \symknight{}c4\\
\textbf{3.} & \symqueen{}a8+ & \symking{}b2\\
\textbf{4.} & \symking{}g2 & \symknight{}e3+\\
\textbf{5.} & \symking{}f3 & \symking{}c3\\
\textbf{6.} & \symking{}f4 & \symking{}d4\\
\textbf{7.} & \symqueen{}af3 & \symknight{}e6\#\\
\end{tabular}\par\vspace{3mm}
\qrcode[height=12mm]{https://helpman.komtera.lt/?fen=1q1N4/8/8/N7/8/2q5/8/K6k&moves=6.5}\\{\sffamily\tiny\color{ink}scannen zum Lösen}
\end{minipage}
\vspace{2mm}

{\raggedright\noindent\textcolor{ink}{\sffamily\textsc{Themen}}\enspace \themetag{reines Matt} \themetag{Mustermatt} \themetag{Idealmatt} \themetag{Selbstblock} \themetag{ohne Schlag}\par\smallskip{\footnotesize\textit{Hinweis:} die Partei am Zug steht im Diagramm im Schach.}\par}

\clearpage
\phantomsection\label{hm:KNNvkrb}
\begin{tcolorbox}[enhanced,colback=ink,colframe=ink,arc=0pt,boxrule=0pt,left=4mm,right=4mm,top=2.5mm,bottom=2.5mm]
\color{white}\sffamily{\LARGE\bfseries KNNvkrb}\hfill{\Large h\#6.5}\\[0.5mm]
{\small 6 Steine\quad Weiß am Zug}\hfill{\small Nr.~622\qquad\href{https://helpman.komtera.lt/?fen=7N/7k/8/N2b4/8/8/8/Kr6&moves=6.5}{\color{white}\bfseries lösen\,\ensuremath{\nearrow}}}
\end{tcolorbox}
\vspace{2mm}
\noindent\begin{minipage}[t]{36mm}\vspace{0pt}\sffamily\small\raggedright
\hmfigure{längste eindeutige}{{\Large h\#6.5}}\hmfigure{eindeutig bei dieser Länge}{{\Large 1} Stellung}\hmfigure{längstes Matt}{{\Large h\#10}}\hmfigure{in der Klasse}{15,5 Milliarden Stellungen,\\ 4,66 Milliarden mit Hilfsmatt}
\end{minipage}\hfill
\begin{minipage}[t]{80mm}\vspace{0pt}\begin{tikzpicture}\begin{axis}[width=72mm,height=42mm,ybar,bar width=1.6pt,xmin=-1,xmax=21,ymin=0,ymax=1.358,axis y line*=left,axis x line*=bottom,xlabel={\scriptsize Halbzüge bis Matt},ylabel={\scriptsize Stellungen (Mrd.)},tick label style={font=\tiny},label style={font=\scriptsize},enlarge x limits=false,ylabel shift=-3pt,xlabel shift=-3pt]\addplot[fill=ink!55,draw=none] coordinates {(0,4.537e-05) (1,0.000222) (2,0.001128) (3,0.009367) (4,0.04514) (5,0.1991) (6,0.5322) (7,1.095) (8,1.235) (9,1.121) (10,0.3001) (11,0.1162) (12,0.006923) (13,0.001264) (14,2.383e-06) (15,1.512e-05) (16,1.64e-07) (17,6e-09) (18,1.4e-08) (19,4e-09) (20,1.7e-08)};\draw[densely dotted] (axis cs:13,0) -- (axis cs:13,1.358);\end{axis}\begin{axis}[width=72mm,height=42mm,axis y line*=right,axis x line=none,ymode=log,xmin=-1,xmax=21,ylabel={\scriptsize eindeutig (log)},tick label style={font=\tiny},label style={font=\scriptsize},ylabel shift=-3pt]\addplot[mark=*,mark size=0.7pt,draw=black,thin] coordinates {(0,45372) (1,212721) (2,821925) (3,4785110) (4,17201179) (5,26356244) (6,27267773) (7,16963274) (8,3788407) (9,1026812) (10,25795) (11,5231) (12,2) (13,1)};\end{axis}\end{tikzpicture}\end{minipage}
\vspace{3mm}

\noindent\begin{minipage}[t]{60mm}\vspace{0pt}\sbox\hmboard{\chessboard[setfen=7N/7k/8/N2b4/8/8/8/Kr6 w - - 0 1]}\usebox\hmboard\\[1mm]
\makebox[\wd\hmboard]{\sffamily\hspace{3mm}\textbf{h\#6.5}\hfill(3+3)}\end{minipage}\hfill
\begin{minipage}[t]{52mm}\vspace{0pt}
\textcolor{ink}{\sffamily\textsc{Lösung}}\par\smallskip\small
\begin{tabular}{@{}r@{\enspace}l@{\enspace}l@{}}
\textbf{1.} & \ldots & \symking{}xb1\\
\textbf{2.} & \symking{}h6 & \symknight{}g6\\
\textbf{3.} & \symking{}g5 & \symknight{}e7\\
\textbf{4.} & \symking{}f4 & \symking{}a1\\
\textbf{5.} & \symking{}e3 & \symknight{}xd5+\\
\textbf{6.} & \symking{}d2 & \symknight{}e3\\
\textbf{7.} & \symking{}c1 & \symknight{}b3\#\\
\end{tabular}\par\vspace{3mm}
\qrcode[height=12mm]{https://helpman.komtera.lt/?fen=7N/7k/8/N2b4/8/8/8/Kr6&moves=6.5}\\{\sffamily\tiny\color{ink}scannen zum Lösen}
\end{minipage}
\vspace{2mm}

{\raggedright\noindent\textcolor{ink}{\sffamily\textsc{Themen}}\enspace \themetag{reines Matt} \themetag{Mustermatt} \themetag{Idealmatt} \themetag{Spiegelmatt} \themetag{Switchback} \themetag{ein Stein zieht} \themetag{Umnov}\par\smallskip{\footnotesize\textit{Hinweis:} die Partei am Zug steht im Diagramm im Schach; die Lösung beginnt mit einem Schlag.}\par}

\clearpage
\phantomsection\label{hm:KNNvkrn}
\begin{tcolorbox}[enhanced,colback=ink,colframe=ink,arc=0pt,boxrule=0pt,left=4mm,right=4mm,top=2.5mm,bottom=2.5mm]
\color{white}\sffamily{\LARGE\bfseries KNNvkrn}\hfill{\Large h\#6.5}\\[0.5mm]
{\small 6 Steine\quad Weiß am Zug}\hfill{\small Nr.~623\qquad\href{https://helpman.komtera.lt/?fen=N5k1/8/8/8/8/8/N2r4/n2K4&moves=6.5}{\color{white}\bfseries lösen\,\ensuremath{\nearrow}}}
\end{tcolorbox}
\vspace{2mm}
\noindent\begin{minipage}[t]{36mm}\vspace{0pt}\sffamily\small\raggedright
\hmfigure{längste eindeutige}{{\Large h\#6.5}}\hmfigure{eindeutig bei dieser Länge}{{\Large 1} Stellung}\hmfigure{längstes Matt}{{\Large h\#8}}\hmfigure{in der Klasse}{15,5 Milliarden Stellungen,\\ 4,77 Milliarden mit Hilfsmatt}
\end{minipage}\hfill
\begin{minipage}[t]{80mm}\vspace{0pt}\begin{tikzpicture}\begin{axis}[width=72mm,height=42mm,ybar,bar width=1.6pt,xmin=-1,xmax=17,ymin=0,ymax=1.335,axis y line*=left,axis x line*=bottom,xlabel={\scriptsize Halbzüge bis Matt},ylabel={\scriptsize Stellungen (Mrd.)},tick label style={font=\tiny},label style={font=\scriptsize},enlarge x limits=false,ylabel shift=-3pt,xlabel shift=-3pt]\addplot[fill=ink!55,draw=none] coordinates {(0,4.935e-05) (1,0.0002491) (2,0.001175) (3,0.009729) (4,0.04509) (5,0.1941) (6,0.51) (7,1.005) (8,1.214) (9,1.123) (10,0.4408) (11,0.2075) (12,0.01718) (13,0.003725) (14,1.85e-05) (15,4.987e-06) (16,1.49e-07)};\draw[densely dotted] (axis cs:13,0) -- (axis cs:13,1.335);\end{axis}\begin{axis}[width=72mm,height=42mm,axis y line*=right,axis x line=none,ymode=log,xmin=-1,xmax=17,ylabel={\scriptsize eindeutig (log)},tick label style={font=\tiny},label style={font=\scriptsize},ylabel shift=-3pt]\addplot[mark=*,mark size=0.7pt,draw=black,thin] coordinates {(0,49349) (1,240190) (2,847977) (3,4922207) (4,16563300) (5,23521918) (6,22725701) (7,14188430) (8,4344928) (9,1327834) (10,73849) (11,14202) (12,35) (13,1)};\end{axis}\end{tikzpicture}\end{minipage}
\vspace{3mm}

\noindent\begin{minipage}[t]{60mm}\vspace{0pt}\sbox\hmboard{\chessboard[setfen=N5k1/8/8/8/8/8/N2r4/n2K4 w - - 0 1]}\usebox\hmboard\\[1mm]
\makebox[\wd\hmboard]{\sffamily\hspace{3mm}\textbf{h\#6.5}\hfill(3+3)}\end{minipage}\hfill
\begin{minipage}[t]{52mm}\vspace{0pt}
\textcolor{ink}{\sffamily\textsc{Lösung}}\par\smallskip\small
\begin{tabular}{@{}r@{\enspace}l@{\enspace}l@{}}
\textbf{1.} & \ldots & \symking{}xd2\\
\textbf{2.} & \symking{}f7 & \symking{}c3\\
\textbf{3.} & \symking{}e6 & \symking{}b4\\
\textbf{4.} & \symking{}d5 & \symking{}a5\\
\textbf{5.} & \symking{}c4 & \symknight{}b6+\\
\textbf{6.} & \symking{}b3 & \symknight{}c1+\\
\textbf{7.} & \symking{}a3 & \symknight{}c4\#\\
\end{tabular}\par\vspace{3mm}
\qrcode[height=12mm]{https://helpman.komtera.lt/?fen=N5k1/8/8/8/8/8/N2r4/n2K4&moves=6.5}\\{\sffamily\tiny\color{ink}scannen zum Lösen}
\end{minipage}
\vspace{2mm}

{\raggedright\noindent\textcolor{ink}{\sffamily\textsc{Themen}}\enspace \themetag{reines Matt} \themetag{Mustermatt} \themetag{Spiegelmatt} \themetag{ein Stein zieht}\par\smallskip{\footnotesize\textit{Hinweis:} die Partei am Zug steht im Diagramm im Schach; die Lösung beginnt mit einem Schlag.}\par}

\clearpage
\phantomsection\label{hm:KNNvkrr}
\begin{tcolorbox}[enhanced,colback=ink,colframe=ink,arc=0pt,boxrule=0pt,left=4mm,right=4mm,top=2.5mm,bottom=2.5mm]
\color{white}\sffamily{\LARGE\bfseries KNNvkrr}\hfill{\Large h\#6.5}\\[0.5mm]
{\small 6 Steine\quad Weiß am Zug}\hfill{\small Nr.~624\qquad\href{https://helpman.komtera.lt/?fen=8/7N/7k/8/8/rr6/8/K1N5&moves=6.5}{\color{white}\bfseries lösen\,\ensuremath{\nearrow}}}
\end{tcolorbox}
\vspace{2mm}
\noindent\begin{minipage}[t]{36mm}\vspace{0pt}\sffamily\small\raggedright
\hmfigure{längste eindeutige}{{\Large h\#6.5}}\hmfigure{eindeutig bei dieser Länge}{{\Large 2\,263} Stellungen}\hmfigure{längstes Matt}{{\Large h\#8.5}}\hmfigure{in der Klasse}{15,5 Milliarden Stellungen,\\ 2,22 Milliarden mit Hilfsmatt}
\end{minipage}\hfill
\begin{minipage}[t]{80mm}\vspace{0pt}\begin{tikzpicture}\begin{axis}[width=72mm,height=42mm,ybar,bar width=1.6pt,xmin=-1,xmax=18,ymin=0,ymax=604.5,axis y line*=left,axis x line*=bottom,xlabel={\scriptsize Halbzüge bis Matt},ylabel={\scriptsize Stellungen (Mio.)},tick label style={font=\tiny},label style={font=\scriptsize},enlarge x limits=false,ylabel shift=-3pt,xlabel shift=-3pt]\addplot[fill=ink!55,draw=none] coordinates {(0,0.0192) (1,0.09388) (2,0.4893) (3,4.143) (4,17.55) (5,81.84) (6,187.5) (7,425.1) (8,459) (9,549.5) (10,252.5) (11,187.1) (12,37.59) (13,15.95) (14,0.6925) (15,0.1885) (16,1.7e-05) (17,6e-06)};\draw[densely dotted] (axis cs:13,0) -- (axis cs:13,604.5);\end{axis}\begin{axis}[width=72mm,height=42mm,axis y line*=right,axis x line=none,ymode=log,xmin=-1,xmax=18,ylabel={\scriptsize eindeutig (log)},tick label style={font=\tiny},label style={font=\scriptsize},ylabel shift=-3pt]\addplot[mark=*,mark size=0.7pt,draw=black,thin] coordinates {(0,19196) (1,90611) (2,347498) (3,2074789) (4,5628536) (5,11435710) (6,10091179) (7,9743724) (8,3640741) (9,2328093) (10,424917) (11,191887) (12,10672) (13,2263)};\end{axis}\end{tikzpicture}\end{minipage}
\vspace{3mm}

\noindent\begin{minipage}[t]{60mm}\vspace{0pt}\sbox\hmboard{\chessboard[setfen=8/7N/7k/8/8/rr6/8/K1N5 w - - 0 1]}\usebox\hmboard\\[1mm]
\makebox[\wd\hmboard]{\sffamily\hspace{3mm}\textbf{h\#6.5}\hfill(3+3)}\end{minipage}\hfill
\begin{minipage}[t]{52mm}\vspace{0pt}
\textcolor{ink}{\sffamily\textsc{Lösung}}\par\smallskip\small
\begin{tabular}{@{}r@{\enspace}l@{\enspace}l@{}}
\textbf{1.} & \ldots & \symknight{}a2\\
\textbf{2.} & \symrook{}b2 & \symking{}xb2\\
\textbf{3.} & \symking{}h5 & \symking{}xa3\\
\textbf{4.} & \symking{}g4 & \symking{}b2\\
\textbf{5.} & \symking{}f3 & \symknight{}g5+\\
\textbf{6.} & \symking{}e2 & \symknight{}f3\\
\textbf{7.} & \symking{}d1 & \symknight{}c3\#\\
\end{tabular}\par\vspace{3mm}
\qrcode[height=12mm]{https://helpman.komtera.lt/?fen=8/7N/7k/8/8/rr6/8/K1N5&moves=6.5}\\{\sffamily\tiny\color{ink}scannen zum Lösen}
\end{minipage}
\vspace{2mm}

{\raggedright\noindent\textcolor{ink}{\sffamily\textsc{Themen}}\enspace \themetag{reines Matt} \themetag{Mustermatt} \themetag{Idealmatt} \themetag{Spiegelmatt} \themetag{Switchback} \themetag{Umnov}\par\smallskip{\footnotesize\textit{Hinweis:} die Partei am Zug steht im Diagramm im Schach.}\par}

\clearpage
\phantomsection\label{hm:KNPPPvk}
\begin{tcolorbox}[enhanced,colback=ink,colframe=ink,arc=0pt,boxrule=0pt,left=4mm,right=4mm,top=2.5mm,bottom=2.5mm]
\color{white}\sffamily{\LARGE\bfseries KNPPPvk}\hfill{\Large h\#6.5}\\[0.5mm]
{\small 6 Steine\quad Weiß am Zug}\hfill{\small Nr.~625\qquad\href{https://helpman.komtera.lt/?fen=8/8/8/8/8/1N1P4/2PP4/K6k&moves=6.5}{\color{white}\bfseries lösen\,\ensuremath{\nearrow}}}
\end{tcolorbox}
\vspace{2mm}
\noindent\begin{minipage}[t]{36mm}\vspace{0pt}\sffamily\small\raggedright
\hmfigure{längste eindeutige}{{\Large h\#6.5}}\hmfigure{eindeutig bei dieser Länge}{{\Large 103} Stellungen}\hmfigure{längstes Matt}{{\Large h\#8}}\hmfigure{in der Klasse}{25,6 Milliarden Stellungen,\\ 3,09 Milliarden mit Hilfsmatt}
\end{minipage}\hfill
\begin{minipage}[t]{80mm}\vspace{0pt}\begin{tikzpicture}\begin{axis}[width=72mm,height=42mm,ybar,bar width=1.6pt,xmin=-1,xmax=17,ymin=0,ymax=733,axis y line*=left,axis x line*=bottom,xlabel={\scriptsize Halbzüge bis Matt},ylabel={\scriptsize Stellungen (Mio.)},tick label style={font=\tiny},label style={font=\scriptsize},enlarge x limits=false,ylabel shift=-3pt,xlabel shift=-3pt]\addplot[fill=ink!55,draw=none] coordinates {(0,0.07687) (1,3.216) (2,7.422) (3,58.44) (4,171.1) (5,483.4) (6,666.4) (7,507.9) (8,495.6) (9,261.5) (10,243.3) (11,91.98) (12,84.64) (13,7.93) (14,5.202) (15,0.005842) (16,0.000106)};\draw[densely dotted] (axis cs:13,0) -- (axis cs:13,733);\end{axis}\begin{axis}[width=72mm,height=42mm,axis y line*=right,axis x line=none,ymode=log,xmin=-1,xmax=17,ylabel={\scriptsize eindeutig (log)},tick label style={font=\tiny},label style={font=\scriptsize},ylabel shift=-3pt]\addplot[mark=*,mark size=0.7pt,draw=black,thin] coordinates {(0,76865) (1,897902) (2,2657034) (3,28881431) (4,51160212) (5,41527601) (6,24148907) (7,11814752) (8,6280457) (9,2330273) (10,1155289) (11,294035) (12,109313) (13,103)};\end{axis}\end{tikzpicture}\end{minipage}
\vspace{3mm}

\noindent\begin{minipage}[t]{60mm}\vspace{0pt}\sbox\hmboard{\chessboard[setfen=8/8/8/8/8/1N1P4/2PP4/K6k w - - 0 1]}\usebox\hmboard\\[1mm]
\makebox[\wd\hmboard]{\sffamily\hspace{3mm}\textbf{h\#6.5}\hfill(5+1)}\end{minipage}\hfill
\begin{minipage}[t]{52mm}\vspace{0pt}
\textcolor{ink}{\sffamily\textsc{Lösung}}\par\smallskip\small
\begin{tabular}{@{}r@{\enspace}l@{\enspace}l@{}}
\textbf{1.} & \ldots & d4\\
\textbf{2.} & \symking{}g2 & d5\\
\textbf{3.} & \symking{}f3 & d6\\
\textbf{4.} & \symking{}e4 & d7\\
\textbf{5.} & \symking{}d5 & d8=\symqueen{}+\\
\textbf{6.} & \symking{}c4 & d3+\\
\textbf{7.} & \symking{}b4 & \symqueen{}a5\#\\
\end{tabular}\par\vspace{3mm}
\qrcode[height=12mm]{https://helpman.komtera.lt/?fen=8/8/8/8/8/1N1P4/2PP4/K6k&moves=6.5}\\{\sffamily\tiny\color{ink}scannen zum Lösen}
\end{minipage}
\vspace{2mm}

{\raggedright\noindent\textcolor{ink}{\sffamily\textsc{Themen}}\enspace \themetag{Umwandlung} \themetag{ein Stein zieht} \themetag{ohne Schlag}\par}

\clearpage
\phantomsection\label{hm:KPPPPvk}
\begin{tcolorbox}[enhanced,colback=ink,colframe=ink,arc=0pt,boxrule=0pt,left=4mm,right=4mm,top=2.5mm,bottom=2.5mm]
\color{white}\sffamily{\LARGE\bfseries KPPPPvk}\hfill{\Large h\#6.5}\\[0.5mm]
{\small 6 Steine\quad Weiß am Zug}\hfill{\small Nr.~626\qquad\href{https://helpman.komtera.lt/?fen=8/8/7k/8/6P1/P7/3P2P1/K7&moves=6.5}{\color{white}\bfseries lösen\,\ensuremath{\nearrow}}}
\end{tcolorbox}
\vspace{2mm}
\noindent\begin{minipage}[t]{36mm}\vspace{0pt}\sffamily\small\raggedright
\hmfigure{längste eindeutige}{{\Large h\#6.5}}\hmfigure{eindeutig bei dieser Länge}{{\Large 33} Stellungen}\hmfigure{längstes Matt}{{\Large h\#7.5}}\hmfigure{in der Klasse}{19,2 Milliarden Stellungen,\\ 585 Millionen mit Hilfsmatt}
\end{minipage}\hfill
\begin{minipage}[t]{80mm}\vspace{0pt}\begin{tikzpicture}\begin{axis}[width=72mm,height=42mm,ybar,bar width=1.6pt,xmin=-1,xmax=16,ymin=0,ymax=145.1,axis y line*=left,axis x line*=bottom,xlabel={\scriptsize Halbzüge bis Matt},ylabel={\scriptsize Stellungen (Mio.)},tick label style={font=\tiny},label style={font=\scriptsize},enlarge x limits=false,ylabel shift=-3pt,xlabel shift=-3pt]\addplot[fill=ink!55,draw=none] coordinates {(0,0.01418) (1,0.7469) (2,1.643) (3,11.67) (4,32.52) (5,92.95) (6,131.9) (7,109.8) (8,98.06) (9,48.69) (10,36.04) (11,11.1) (12,8.351) (13,0.6325) (14,0.3648) (15,0.001657)};\draw[densely dotted] (axis cs:13,0) -- (axis cs:13,145.1);\end{axis}\begin{axis}[width=72mm,height=42mm,axis y line*=right,axis x line=none,ymode=log,xmin=-1,xmax=16,ylabel={\scriptsize eindeutig (log)},tick label style={font=\tiny},label style={font=\scriptsize},ylabel shift=-3pt]\addplot[mark=*,mark size=0.7pt,draw=black,thin] coordinates {(0,14181) (1,192038) (2,560134) (3,5433110) (4,9673624) (5,7686343) (6,4516759) (7,1875764) (8,819233) (9,248192) (10,107639) (11,22895) (12,7617) (13,33)};\end{axis}\end{tikzpicture}\end{minipage}
\vspace{3mm}

\noindent\begin{minipage}[t]{60mm}\vspace{0pt}\sbox\hmboard{\chessboard[setfen=8/8/7k/8/6P1/P7/3P2P1/K7 w - - 0 1]}\usebox\hmboard\\[1mm]
\makebox[\wd\hmboard]{\sffamily\hspace{3mm}\textbf{h\#6.5}\hfill(5+1)}\end{minipage}\hfill
\begin{minipage}[t]{52mm}\vspace{0pt}
\textcolor{ink}{\sffamily\textsc{Lösung}}\par\smallskip\small
\begin{tabular}{@{}r@{\enspace}l@{\enspace}l@{}}
\textbf{1.} & \ldots & g5+\\
\textbf{2.} & \symking{}g7 & g6\\
\textbf{3.} & \symking{}f6 & g7\\
\textbf{4.} & \symking{}e5 & g8=\symqueen{}\\
\textbf{5.} & \symking{}d4 & \symqueen{}f8\\
\textbf{6.} & \symking{}c4 & \symking{}b1\\
\textbf{7.} & \symking{}b3 & \symqueen{}b4\#\\
\end{tabular}\par\vspace{3mm}
\qrcode[height=12mm]{https://helpman.komtera.lt/?fen=8/8/7k/8/6P1/P7/3P2P1/K7&moves=6.5}\\{\sffamily\tiny\color{ink}scannen zum Lösen}
\end{minipage}
\vspace{2mm}

{\raggedright\noindent\textcolor{ink}{\sffamily\textsc{Themen}}\enspace \themetag{Umwandlung} \themetag{ein Stein zieht} \themetag{ohne Schlag} \themetag{Umnov}\par}

\clearpage
\phantomsection\label{hm:KPPPvkq}
\begin{tcolorbox}[enhanced,colback=ink,colframe=ink,arc=0pt,boxrule=0pt,left=4mm,right=4mm,top=2.5mm,bottom=2.5mm]
\color{white}\sffamily{\LARGE\bfseries KPPPvkq}\hfill{\Large h\#6.5}\\[0.5mm]
{\small 6 Steine\quad Weiß am Zug}\hfill{\small Nr.~627\qquad\href{https://helpman.komtera.lt/?fen=8/8/8/8/8/q7/1PP3P1/K1k5&moves=6.5}{\color{white}\bfseries lösen\,\ensuremath{\nearrow}}}
\end{tcolorbox}
\vspace{2mm}
\noindent\begin{minipage}[t]{36mm}\vspace{0pt}\sffamily\small\raggedright
\hmfigure{längste eindeutige}{{\Large h\#6.5}}\hmfigure{eindeutig bei dieser Länge}{{\Large 3\,694} Stellungen}\hmfigure{längstes Matt}{{\Large h\#8.5}}\hmfigure{in der Klasse}{25,6 Milliarden Stellungen,\\ 2,68 Milliarden mit Hilfsmatt}
\end{minipage}\hfill
\begin{minipage}[t]{80mm}\vspace{0pt}\begin{tikzpicture}\begin{axis}[width=72mm,height=42mm,ybar,bar width=1.6pt,xmin=-1,xmax=18,ymin=0,ymax=675.1,axis y line*=left,axis x line*=bottom,xlabel={\scriptsize Halbzüge bis Matt},ylabel={\scriptsize Stellungen (Mio.)},tick label style={font=\tiny},label style={font=\scriptsize},enlarge x limits=false,ylabel shift=-3pt,xlabel shift=-3pt]\addplot[fill=ink!55,draw=none] coordinates {(0,0.01837) (1,0.8577) (2,3.666) (3,33.47) (4,96.1) (5,335.9) (6,418.8) (7,613.7) (8,383.1) (9,388.6) (10,180.3) (11,142.1) (12,56.83) (13,20.23) (14,1.654) (15,0.2519) (16,9.3e-05) (17,3e-05)};\draw[densely dotted] (axis cs:13,0) -- (axis cs:13,675.1);\end{axis}\begin{axis}[width=72mm,height=42mm,axis y line*=right,axis x line=none,ymode=log,xmin=-1,xmax=18,ylabel={\scriptsize eindeutig (log)},tick label style={font=\tiny},label style={font=\scriptsize},ylabel shift=-3pt]\addplot[mark=*,mark size=0.7pt,draw=black,thin] coordinates {(0,18371) (1,273973) (2,1064449) (3,11581992) (4,19443736) (5,27710083) (6,11354390) (7,9858709) (8,2761836) (9,2096134) (10,445061) (11,251423) (12,23991) (13,3694)};\end{axis}\end{tikzpicture}\end{minipage}
\vspace{3mm}

\noindent\begin{minipage}[t]{60mm}\vspace{0pt}\sbox\hmboard{\chessboard[setfen=8/8/8/8/8/q7/1PP3P1/K1k5 w - - 0 1]}\usebox\hmboard\\[1mm]
\makebox[\wd\hmboard]{\sffamily\hspace{3mm}\textbf{h\#6.5}\hfill(4+2)}\end{minipage}\hfill
\begin{minipage}[t]{52mm}\vspace{0pt}
\textcolor{ink}{\sffamily\textsc{Lösung}}\par\smallskip\small
\begin{tabular}{@{}r@{\enspace}l@{\enspace}l@{}}
\textbf{1.} & \ldots & bxa3\\
\textbf{2.} & \symking{}d2 & g4\\
\textbf{3.} & \symking{}c3 & g5\\
\textbf{4.} & \symking{}c4 & g6\\
\textbf{5.} & \symking{}b5 & g7\\
\textbf{6.} & \symking{}a4 & g8=\symqueen{}\\
\textbf{7.} & \symking{}xa3 & \symqueen{}b3\#\\
\end{tabular}\par\vspace{3mm}
\qrcode[height=12mm]{https://helpman.komtera.lt/?fen=8/8/8/8/8/q7/1PP3P1/K1k5&moves=6.5}\\{\sffamily\tiny\color{ink}scannen zum Lösen}
\end{minipage}
\vspace{2mm}

{\raggedright\noindent\textcolor{ink}{\sffamily\textsc{Themen}}\enspace \themetag{Umwandlung} \themetag{Excelsior} \themetag{ein Stein zieht} \themetag{Kniest} \themetag{Zajic} \themetag{ohne Schach}\par\smallskip{\footnotesize\textit{Hinweis:} die Partei am Zug steht im Diagramm im Schach; die Lösung beginnt mit einem Schlag.}\par}

\clearpage
\phantomsection\label{hm:KPPPvkr}
\begin{tcolorbox}[enhanced,colback=ink,colframe=ink,arc=0pt,boxrule=0pt,left=4mm,right=4mm,top=2.5mm,bottom=2.5mm]
\color{white}\sffamily{\LARGE\bfseries KPPPvkr}\hfill{\Large h\#6.5}\\[0.5mm]
{\small 6 Steine\quad Weiß am Zug}\hfill{\small Nr.~628\qquad\href{https://helpman.komtera.lt/?fen=8/8/8/8/8/r7/1PP3P1/K1k5&moves=6.5}{\color{white}\bfseries lösen\,\ensuremath{\nearrow}}}
\end{tcolorbox}
\vspace{2mm}
\noindent\begin{minipage}[t]{36mm}\vspace{0pt}\sffamily\small\raggedright
\hmfigure{längste eindeutige}{{\Large h\#6.5}}\hmfigure{eindeutig bei dieser Länge}{{\Large 554} Stellungen}\hmfigure{längstes Matt}{{\Large h\#8.5}}\hmfigure{in der Klasse}{25,6 Milliarden Stellungen,\\ 2,88 Milliarden mit Hilfsmatt}
\end{minipage}\hfill
\begin{minipage}[t]{80mm}\vspace{0pt}\begin{tikzpicture}\begin{axis}[width=72mm,height=42mm,ybar,bar width=1.6pt,xmin=-1,xmax=18,ymin=0,ymax=712.8,axis y line*=left,axis x line*=bottom,xlabel={\scriptsize Halbzüge bis Matt},ylabel={\scriptsize Stellungen (Mio.)},tick label style={font=\tiny},label style={font=\scriptsize},enlarge x limits=false,ylabel shift=-3pt,xlabel shift=-3pt]\addplot[fill=ink!55,draw=none] coordinates {(0,0.02408) (1,1.455) (2,5.645) (3,43.26) (4,134.3) (5,411.6) (6,562.1) (7,648) (8,420.3) (9,322.5) (10,173.9) (11,99.56) (12,50.99) (13,9.827) (14,1.003) (15,0.03294) (16,1e-06) (17,2.2e-05)};\draw[densely dotted] (axis cs:13,0) -- (axis cs:13,712.8);\end{axis}\begin{axis}[width=72mm,height=42mm,axis y line*=right,axis x line=none,ymode=log,xmin=-1,xmax=18,ylabel={\scriptsize eindeutig (log)},tick label style={font=\tiny},label style={font=\scriptsize},ylabel shift=-3pt]\addplot[mark=*,mark size=0.7pt,draw=black,thin] coordinates {(0,24077) (1,452757) (2,1821998) (3,15346831) (4,30341606) (5,34734945) (6,16204765) (7,10385070) (8,3106066) (9,1745729) (10,465115) (11,145614) (12,17886) (13,554)};\end{axis}\end{tikzpicture}\end{minipage}
\vspace{3mm}

\noindent\begin{minipage}[t]{60mm}\vspace{0pt}\sbox\hmboard{\chessboard[setfen=8/8/8/8/8/r7/1PP3P1/K1k5 w - - 0 1]}\usebox\hmboard\\[1mm]
\makebox[\wd\hmboard]{\sffamily\hspace{3mm}\textbf{h\#6.5}\hfill(4+2)}\end{minipage}\hfill
\begin{minipage}[t]{52mm}\vspace{0pt}
\textcolor{ink}{\sffamily\textsc{Lösung}}\par\smallskip\small
\begin{tabular}{@{}r@{\enspace}l@{\enspace}l@{}}
\textbf{1.} & \ldots & bxa3\\
\textbf{2.} & \symking{}d2 & g4\\
\textbf{3.} & \symking{}c3 & g5\\
\textbf{4.} & \symking{}c4 & g6\\
\textbf{5.} & \symking{}b5 & g7\\
\textbf{6.} & \symking{}a4 & g8=\symqueen{}\\
\textbf{7.} & \symking{}xa3 & \symqueen{}b3\#\\
\end{tabular}\par\vspace{3mm}
\qrcode[height=12mm]{https://helpman.komtera.lt/?fen=8/8/8/8/8/r7/1PP3P1/K1k5&moves=6.5}\\{\sffamily\tiny\color{ink}scannen zum Lösen}
\end{minipage}
\vspace{2mm}

{\raggedright\noindent\textcolor{ink}{\sffamily\textsc{Themen}}\enspace \themetag{Umwandlung} \themetag{Excelsior} \themetag{ein Stein zieht} \themetag{Kniest} \themetag{Zajic} \themetag{ohne Schach}\par\smallskip{\footnotesize\textit{Hinweis:} die Partei am Zug steht im Diagramm im Schach; die Lösung beginnt mit einem Schlag.}\par}

\clearpage
\phantomsection\label{hm:KPPvkrr}
\begin{tcolorbox}[enhanced,colback=ink,colframe=ink,arc=0pt,boxrule=0pt,left=4mm,right=4mm,top=2.5mm,bottom=2.5mm]
\color{white}\sffamily{\LARGE\bfseries KPPvkrr}\hfill{\Large h\#6.5}\\[0.5mm]
{\small 6 Steine\quad Weiß am Zug}\hfill{\small Nr.~629\qquad\href{https://helpman.komtera.lt/?fen=8/8/8/8/1r6/3P4/3P4/Krk5&moves=6.5}{\color{white}\bfseries lösen\,\ensuremath{\nearrow}}}
\end{tcolorbox}
\vspace{2mm}
\noindent\begin{minipage}[t]{36mm}\vspace{0pt}\sffamily\small\raggedright
\hmfigure{längste eindeutige}{{\Large h\#6.5}}\hmfigure{eindeutig bei dieser Länge}{{\Large 3\,010} Stellungen}\hmfigure{längstes Matt}{{\Large h\#8.5}}\hmfigure{in der Klasse}{34,1 Milliarden Stellungen,\\ 5,32 Milliarden mit Hilfsmatt}
\end{minipage}\hfill
\begin{minipage}[t]{80mm}\vspace{0pt}\begin{tikzpicture}\begin{axis}[width=72mm,height=42mm,ybar,bar width=1.6pt,xmin=-1,xmax=18,ymin=0,ymax=1.376,axis y line*=left,axis x line*=bottom,xlabel={\scriptsize Halbzüge bis Matt},ylabel={\scriptsize Stellungen (Mrd.)},tick label style={font=\tiny},label style={font=\scriptsize},enlarge x limits=false,ylabel shift=-3pt,xlabel shift=-3pt]\addplot[fill=ink!55,draw=none] coordinates {(0,1.033e-05) (1,0.0008935) (2,0.004561) (3,0.03988) (4,0.1388) (5,0.5293) (6,0.7563) (7,1.251) (8,0.7069) (9,0.8738) (10,0.3959) (11,0.4083) (12,0.1479) (13,0.0651) (14,0.0003977) (15,0.0002815) (16,6.7e-08) (17,4.1e-08)};\draw[densely dotted] (axis cs:13,0) -- (axis cs:13,1.376);\end{axis}\begin{axis}[width=72mm,height=42mm,axis y line*=right,axis x line=none,ymode=log,xmin=-1,xmax=18,ylabel={\scriptsize eindeutig (log)},tick label style={font=\tiny},label style={font=\scriptsize},ylabel shift=-3pt]\addplot[mark=*,mark size=0.7pt,draw=black,thin] coordinates {(0,10330) (1,330888) (2,1711953) (3,14753792) (4,32267833) (5,50102562) (6,22040779) (7,21222065) (8,4336330) (9,4015824) (10,1145654) (11,376025) (12,2884) (13,3010)};\end{axis}\end{tikzpicture}\end{minipage}
\vspace{3mm}

\noindent\begin{minipage}[t]{60mm}\vspace{0pt}\sbox\hmboard{\chessboard[setfen=8/8/8/8/1r6/3P4/3P4/Krk5 w - - 0 1]}\usebox\hmboard\\[1mm]
\makebox[\wd\hmboard]{\sffamily\hspace{3mm}\textbf{h\#6.5}\hfill(3+3)}\end{minipage}\hfill
\begin{minipage}[t]{52mm}\vspace{0pt}
\textcolor{ink}{\sffamily\textsc{Lösung}}\par\smallskip\small
\begin{tabular}{@{}r@{\enspace}l@{\enspace}l@{}}
\textbf{1.} & \ldots & \symking{}a2\\
\textbf{2.} & \symrook{}c4 & dxc4\\
\textbf{3.} & \symking{}xd2 & c5\\
\textbf{4.} & \symking{}c3 & c6\\
\textbf{5.} & \symking{}b4 & c7\\
\textbf{6.} & \symking{}a4 & c8=\symqueen{}\\
\textbf{7.} & \symrook{}b4 & \symqueen{}a6\#\\
\end{tabular}\par\vspace{3mm}
\qrcode[height=12mm]{https://helpman.komtera.lt/?fen=8/8/8/8/1r6/3P4/3P4/Krk5&moves=6.5}\\{\sffamily\tiny\color{ink}scannen zum Lösen}
\end{minipage}
\vspace{2mm}

{\raggedright\noindent\textcolor{ink}{\sffamily\textsc{Themen}}\enspace \themetag{Umwandlung} \themetag{Selbstblock} \themetag{ohne Schach}\par\smallskip{\footnotesize\textit{Hinweis:} die Partei am Zug steht im Diagramm im Schach.}\par}

\clearpage
\phantomsection\label{hm:KQvkqnn}
\begin{tcolorbox}[enhanced,colback=ink,colframe=ink,arc=0pt,boxrule=0pt,left=4mm,right=4mm,top=2.5mm,bottom=2.5mm]
\color{white}\sffamily{\LARGE\bfseries KQvkqnn}\hfill{\Large h\#6.5}\\[0.5mm]
{\small 6 Steine\quad Weiß am Zug}\hfill{\small Nr.~630\qquad\href{https://helpman.komtera.lt/?fen=6k1/8/8/8/8/2n5/3q3Q/1K1n4&moves=6.5}{\color{white}\bfseries lösen\,\ensuremath{\nearrow}}}
\end{tcolorbox}
\vspace{2mm}
\noindent\begin{minipage}[t]{36mm}\vspace{0pt}\sffamily\small\raggedright
\hmfigure{längste eindeutige}{{\Large h\#6.5}}\hmfigure{eindeutig bei dieser Länge}{{\Large 17} Stellungen}\hmfigure{längstes Matt}{{\Large h\#8}}\hmfigure{in der Klasse}{15,5 Milliarden Stellungen,\\ 3,74 Milliarden mit Hilfsmatt}
\end{minipage}\hfill
\begin{minipage}[t]{80mm}\vspace{0pt}\begin{tikzpicture}\begin{axis}[width=72mm,height=42mm,ybar,bar width=1.6pt,xmin=-1,xmax=17,ymin=0,ymax=874.2,axis y line*=left,axis x line*=bottom,xlabel={\scriptsize Halbzüge bis Matt},ylabel={\scriptsize Stellungen (Mio.)},tick label style={font=\tiny},label style={font=\scriptsize},enlarge x limits=false,ylabel shift=-3pt,xlabel shift=-3pt]\addplot[fill=ink!55,draw=none] coordinates {(0,1.005) (1,8.562) (2,31.82) (3,140.5) (4,266.3) (5,569.4) (6,616.9) (7,794.7) (8,592) (9,457.7) (10,180.1) (11,59.83) (12,18.07) (13,1.09) (14,0.121) (15,0.000432) (16,5e-06)};\draw[densely dotted] (axis cs:13,0) -- (axis cs:13,874.2);\end{axis}\begin{axis}[width=72mm,height=42mm,axis y line*=right,axis x line=none,ymode=log,xmin=-1,xmax=17,ylabel={\scriptsize eindeutig (log)},tick label style={font=\tiny},label style={font=\scriptsize},ylabel shift=-3pt]\addplot[mark=*,mark size=0.7pt,draw=black,thin] coordinates {(0,1005354) (1,7216205) (2,14398637) (3,31746745) (4,28210714) (5,32433004) (6,18985709) (7,14773153) (8,5347678) (9,1603401) (10,200327) (11,8987) (12,560) (13,17)};\end{axis}\end{tikzpicture}\end{minipage}
\vspace{3mm}

\noindent\begin{minipage}[t]{60mm}\vspace{0pt}\sbox\hmboard{\chessboard[setfen=6k1/8/8/8/8/2n5/3q3Q/1K1n4 w - - 0 1]}\usebox\hmboard\\[1mm]
\makebox[\wd\hmboard]{\sffamily\hspace{3mm}\textbf{h\#6.5}\hfill(2+4)}\end{minipage}\hfill
\begin{minipage}[t]{52mm}\vspace{0pt}
\textcolor{ink}{\sffamily\textsc{Lösung}}\par\smallskip\small
\begin{tabular}{@{}r@{\enspace}l@{\enspace}l@{}}
\textbf{1.} & \ldots & \symking{}a1\\
\textbf{2.} & \symking{}f7 & \symqueen{}xd2\\
\textbf{3.} & \symking{}e6 & \symqueen{}xc3\\
\textbf{4.} & \symking{}d5 & \symqueen{}e5+\\
\textbf{5.} & \symking{}c4 & \symking{}b1\\
\textbf{6.} & \symking{}b4 & \symking{}c2\\
\textbf{7.} & \symking{}a3 & \symqueen{}a5\#\\
\end{tabular}\par\vspace{3mm}
\qrcode[height=12mm]{https://helpman.komtera.lt/?fen=6k1/8/8/8/8/2n5/3q3Q/1K1n4&moves=6.5}\\{\sffamily\tiny\color{ink}scannen zum Lösen}
\end{minipage}
\vspace{2mm}

{\raggedright\noindent\textcolor{ink}{\sffamily\textsc{Themen}}\enspace \themetag{reines Matt} \themetag{Mustermatt} \themetag{Spiegelmatt} \themetag{Switchback} \themetag{ein Stein zieht}\par\smallskip{\footnotesize\textit{Hinweis:} die Partei am Zug steht im Diagramm im Schach.}\par}

\clearpage
\phantomsection\label{hm:KRBBvkq}
\begin{tcolorbox}[enhanced,colback=ink,colframe=ink,arc=0pt,boxrule=0pt,left=4mm,right=4mm,top=2.5mm,bottom=2.5mm]
\color{white}\sffamily{\LARGE\bfseries KRBBvkq}\hfill{\Large h\#6.5}\\[0.5mm]
{\small 6 Steine\quad Weiß am Zug}\hfill{\small Nr.~631\qquad\href{https://helpman.komtera.lt/?fen=5B2/7k/3B2R1/8/8/8/q7/K7&moves=6.5}{\color{white}\bfseries lösen\,\ensuremath{\nearrow}}}
\end{tcolorbox}
\vspace{2mm}
\noindent\begin{minipage}[t]{36mm}\vspace{0pt}\sffamily\small\raggedright
\hmfigure{längste eindeutige}{{\Large h\#6.5}}\hmfigure{eindeutig bei dieser Länge}{{\Large 7} Stellungen}\hmfigure{längstes Matt}{{\Large h\#7.5}}\hmfigure{in der Klasse}{15,5 Milliarden Stellungen,\\ 3,88 Milliarden mit Hilfsmatt}
\end{minipage}\hfill
\begin{minipage}[t]{80mm}\vspace{0pt}\begin{tikzpicture}\begin{axis}[width=72mm,height=42mm,ybar,bar width=1.6pt,xmin=-1,xmax=16,ymin=0,ymax=1.313,axis y line*=left,axis x line*=bottom,xlabel={\scriptsize Halbzüge bis Matt},ylabel={\scriptsize Stellungen (Mrd.)},tick label style={font=\tiny},label style={font=\scriptsize},enlarge x limits=false,ylabel shift=-3pt,xlabel shift=-3pt]\addplot[fill=ink!55,draw=none] coordinates {(0,0.00334) (1,0.01206) (2,0.06324) (3,0.3191) (4,0.9727) (5,1.193) (6,0.9657) (7,0.3098) (8,0.03166) (9,0.005033) (10,0.0003444) (11,0.0004264) (12,6.602e-05) (13,3.46e-05) (14,6.097e-06) (15,4.52e-07)};\draw[densely dotted] (axis cs:13,0) -- (axis cs:13,1.313);\end{axis}\begin{axis}[width=72mm,height=42mm,axis y line*=right,axis x line=none,ymode=log,xmin=-1,xmax=16,ylabel={\scriptsize eindeutig (log)},tick label style={font=\tiny},label style={font=\scriptsize},ylabel shift=-3pt]\addplot[mark=*,mark size=0.7pt,draw=black,thin] coordinates {(0,3339763) (1,11461287) (2,38564987) (3,112537363) (4,163723527) (5,46694781) (6,5103278) (7,670207) (8,7143) (9,9289) (10,1115) (11,554) (12,75) (13,7)};\end{axis}\end{tikzpicture}\end{minipage}
\vspace{3mm}

\noindent\begin{minipage}[t]{60mm}\vspace{0pt}\sbox\hmboard{\chessboard[setfen=5B2/7k/3B2R1/8/8/8/q7/K7 w - - 0 1]}\usebox\hmboard\\[1mm]
\makebox[\wd\hmboard]{\sffamily\hspace{3mm}\textbf{h\#6.5}\hfill(4+2)}\end{minipage}\hfill
\begin{minipage}[t]{52mm}\vspace{0pt}
\textcolor{ink}{\sffamily\textsc{Lösung}}\par\smallskip\small
\begin{tabular}{@{}r@{\enspace}l@{\enspace}l@{}}
\textbf{1.} & \ldots & \symking{}xa2\\
\textbf{2.} & \symking{}h8 & \symking{}b3\\
\textbf{3.} & \symking{}h7 & \symking{}c4\\
\textbf{4.} & \symking{}h8 & \symking{}d5\\
\textbf{5.} & \symking{}h7 & \symking{}e6\\
\textbf{6.} & \symking{}h8 & \symking{}f7\\
\textbf{7.} & \symking{}h7 & \symrook{}h6\#\\
\end{tabular}\par\vspace{3mm}
\qrcode[height=12mm]{https://helpman.komtera.lt/?fen=5B2/7k/3B2R1/8/8/8/q7/K7&moves=6.5}\\{\sffamily\tiny\color{ink}scannen zum Lösen}
\end{minipage}
\vspace{2mm}

{\raggedright\noindent\textcolor{ink}{\sffamily\textsc{Themen}}\enspace \themetag{Switchback} \themetag{ein Stein zieht} \themetag{Pendel} \themetag{ohne Schach}\par\smallskip{\footnotesize\textit{Hinweis:} die Partei am Zug steht im Diagramm im Schach; die Lösung beginnt mit einem Schlag.}\par}

\clearpage
\phantomsection\label{hm:KRvkbbp}
\begin{tcolorbox}[enhanced,colback=ink,colframe=ink,arc=0pt,boxrule=0pt,left=4mm,right=4mm,top=2.5mm,bottom=2.5mm]
\color{white}\sffamily{\LARGE\bfseries KRvkbbp}\hfill{\Large h\#6.5}\\[0.5mm]
{\small 6 Steine\quad Weiß am Zug}\hfill{\small Nr.~632\qquad\href{https://helpman.komtera.lt/?fen=7k/8/6R1/8/8/bp6/2b5/1K6&moves=6.5}{\color{white}\bfseries lösen\,\ensuremath{\nearrow}}}
\end{tcolorbox}
\vspace{2mm}
\noindent\begin{minipage}[t]{36mm}\vspace{0pt}\sffamily\small\raggedright
\hmfigure{längste eindeutige}{{\Large h\#6.5}}\hmfigure{eindeutig bei dieser Länge}{{\Large 3} Stellungen}\hmfigure{längstes Matt}{{\Large h\#7.5}}\hmfigure{in der Klasse}{45,4 Milliarden Stellungen,\\ 13,9 Milliarden mit Hilfsmatt}
\end{minipage}\hfill
\begin{minipage}[t]{80mm}\vspace{0pt}\begin{tikzpicture}\begin{axis}[width=72mm,height=42mm,ybar,bar width=1.6pt,xmin=-1,xmax=16,ymin=0,ymax=2.927,axis y line*=left,axis x line*=bottom,xlabel={\scriptsize Halbzüge bis Matt},ylabel={\scriptsize Stellungen (Mrd.)},tick label style={font=\tiny},label style={font=\scriptsize},enlarge x limits=false,ylabel shift=-3pt,xlabel shift=-3pt]\addplot[fill=ink!55,draw=none] coordinates {(0,0.002346) (1,0.01501) (2,0.05196) (3,0.2826) (4,0.5602) (5,1.467) (6,1.888) (7,2.661) (8,2.621) (9,2.33) (10,1.355) (11,0.5227) (12,0.1779) (13,0.006674) (14,0.001553) (15,1.099e-06)};\draw[densely dotted] (axis cs:13,0) -- (axis cs:13,2.927);\end{axis}\begin{axis}[width=72mm,height=42mm,axis y line*=right,axis x line=none,ymode=log,xmin=-1,xmax=16,ylabel={\scriptsize eindeutig (log)},tick label style={font=\tiny},label style={font=\scriptsize},ylabel shift=-3pt]\addplot[mark=*,mark size=0.7pt,draw=black,thin] coordinates {(0,2345965) (1,15011891) (2,22954536) (3,85781443) (4,115067423) (5,126210637) (6,70027184) (7,37305647) (8,21555986) (9,7811221) (10,2193536) (11,69443) (12,11590) (13,3)};\end{axis}\end{tikzpicture}\end{minipage}
\vspace{3mm}

\noindent\begin{minipage}[t]{60mm}\vspace{0pt}\sbox\hmboard{\chessboard[setfen=7k/8/6R1/8/8/bp6/2b5/1K6 w - - 0 1]}\usebox\hmboard\\[1mm]
\makebox[\wd\hmboard]{\sffamily\hspace{3mm}\textbf{h\#6.5}\hfill(2+4)}\end{minipage}\hfill
\begin{minipage}[t]{52mm}\vspace{0pt}
\textcolor{ink}{\sffamily\textsc{Lösung}}\par\smallskip\small
\begin{tabular}{@{}r@{\enspace}l@{\enspace}l@{}}
\textbf{1.} & \ldots & \symking{}a1\\
\textbf{2.} & \symbishop{}b2+ & \symking{}xb2\\
\textbf{3.} & \symbishop{}d1 & \symking{}c3\\
\textbf{4.} & b2 & \symking{}d4\\
\textbf{5.} & b1=\symrook{} & \symking{}e5\\
\textbf{6.} & \symrook{}b8 & \symking{}f6\\
\textbf{7.} & \symrook{}g8 & \symrook{}h6\#\\
\end{tabular}\par\vspace{3mm}
\qrcode[height=12mm]{https://helpman.komtera.lt/?fen=7k/8/6R1/8/8/bp6/2b5/1K6&moves=6.5}\\{\sffamily\tiny\color{ink}scannen zum Lösen}
\end{minipage}
\vspace{2mm}

{\raggedright\noindent\textcolor{ink}{\sffamily\textsc{Themen}}\enspace \themetag{reines Matt} \themetag{Mustermatt} \themetag{Umwandlung} \themetag{Unterverwandlung} \themetag{Selbstblock} \themetag{Umnov}\par\smallskip{\footnotesize\textit{Hinweis:} die Partei am Zug steht im Diagramm im Schach.}\par}

\clearpage
\phantomsection\label{hm:KRvkbnp}
\begin{tcolorbox}[enhanced,colback=ink,colframe=ink,arc=0pt,boxrule=0pt,left=4mm,right=4mm,top=2.5mm,bottom=2.5mm]
\color{white}\sffamily{\LARGE\bfseries KRvkbnp}\hfill{\Large h\#6.5}\\[0.5mm]
{\small 6 Steine\quad Weiß am Zug}\hfill{\small Nr.~633\qquad\href{https://helpman.komtera.lt/?fen=8/7k/6R1/8/8/bp6/3n4/1K6&moves=6.5}{\color{white}\bfseries lösen\,\ensuremath{\nearrow}}}
\end{tcolorbox}
\vspace{2mm}
\noindent\begin{minipage}[t]{36mm}\vspace{0pt}\sffamily\small\raggedright
\hmfigure{längste eindeutige}{{\Large h\#6.5}}\hmfigure{eindeutig bei dieser Länge}{{\Large 10} Stellungen}\hmfigure{längstes Matt}{{\Large h\#7.5}}\hmfigure{in der Klasse}{45,4 Milliarden Stellungen,\\ 28,7 Milliarden mit Hilfsmatt}
\end{minipage}\hfill
\begin{minipage}[t]{80mm}\vspace{0pt}\begin{tikzpicture}\begin{axis}[width=72mm,height=42mm,ybar,bar width=1.6pt,xmin=-1,xmax=16,ymin=0,ymax=6.082,axis y line*=left,axis x line*=bottom,xlabel={\scriptsize Halbzüge bis Matt},ylabel={\scriptsize Stellungen (Mrd.)},tick label style={font=\tiny},label style={font=\scriptsize},enlarge x limits=false,ylabel shift=-3pt,xlabel shift=-3pt]\addplot[fill=ink!55,draw=none] coordinates {(0,0.006289) (1,0.04022) (2,0.1267) (3,0.6306) (4,1.273) (5,3.042) (6,4.141) (7,5.351) (8,5.529) (9,4.548) (10,2.722) (11,0.9501) (12,0.3406) (13,0.01103) (14,0.002568) (15,9.96e-07)};\draw[densely dotted] (axis cs:13,0) -- (axis cs:13,6.082);\end{axis}\begin{axis}[width=72mm,height=42mm,axis y line*=right,axis x line=none,ymode=log,xmin=-1,xmax=16,ylabel={\scriptsize eindeutig (log)},tick label style={font=\tiny},label style={font=\scriptsize},ylabel shift=-3pt]\addplot[mark=*,mark size=0.7pt,draw=black,thin] coordinates {(0,6289496) (1,40215878) (2,58107893) (3,202147318) (4,293434356) (5,277343073) (6,161255409) (7,70280997) (8,41734048) (9,13345175) (10,4360992) (11,93017) (12,25712) (13,10)};\end{axis}\end{tikzpicture}\end{minipage}
\vspace{3mm}

\noindent\begin{minipage}[t]{60mm}\vspace{0pt}\sbox\hmboard{\chessboard[setfen=8/7k/6R1/8/8/bp6/3n4/1K6 w - - 0 1]}\usebox\hmboard\\[1mm]
\makebox[\wd\hmboard]{\sffamily\hspace{3mm}\textbf{h\#6.5}\hfill(2+4)}\end{minipage}\hfill
\begin{minipage}[t]{52mm}\vspace{0pt}
\textcolor{ink}{\sffamily\textsc{Lösung}}\par\smallskip\small
\begin{tabular}{@{}r@{\enspace}l@{\enspace}l@{}}
\textbf{1.} & \ldots & \symking{}a1\\
\textbf{2.} & \symbishop{}b2+ & \symking{}xb2\\
\textbf{3.} & \symking{}h8 & \symking{}c3\\
\textbf{4.} & b2 & \symking{}d4\\
\textbf{5.} & b1=\symrook{} & \symking{}e5\\
\textbf{6.} & \symrook{}b8 & \symking{}f6\\
\textbf{7.} & \symrook{}g8 & \symrook{}h6\#\\
\end{tabular}\par\vspace{3mm}
\qrcode[height=12mm]{https://helpman.komtera.lt/?fen=8/7k/6R1/8/8/bp6/3n4/1K6&moves=6.5}\\{\sffamily\tiny\color{ink}scannen zum Lösen}
\end{minipage}
\vspace{2mm}

{\raggedright\noindent\textcolor{ink}{\sffamily\textsc{Themen}}\enspace \themetag{reines Matt} \themetag{Mustermatt} \themetag{Umwandlung} \themetag{Unterverwandlung} \themetag{Selbstblock} \themetag{Umnov}\par\smallskip{\footnotesize\textit{Hinweis:} die Partei am Zug steht im Diagramm im Schach.}\par}

\clearpage
\phantomsection\label{hm:KRvkbpp}
\begin{tcolorbox}[enhanced,colback=ink,colframe=ink,arc=0pt,boxrule=0pt,left=4mm,right=4mm,top=2.5mm,bottom=2.5mm]
\color{white}\sffamily{\LARGE\bfseries KRvkbpp}\hfill{\Large h\#6.5}\\[0.5mm]
{\small 6 Steine\quad Weiß am Zug}\hfill{\small Nr.~634\qquad\href{https://helpman.komtera.lt/?fen=4R1bk/8/8/8/8/p7/p7/K7&moves=6.5}{\color{white}\bfseries lösen\,\ensuremath{\nearrow}}}
\end{tcolorbox}
\vspace{2mm}
\noindent\begin{minipage}[t]{36mm}\vspace{0pt}\sffamily\small\raggedright
\hmfigure{längste eindeutige}{{\Large h\#6.5}}\hmfigure{eindeutig bei dieser Länge}{{\Large 11} Stellungen}\hmfigure{längstes Matt}{{\Large h\#8}}\hmfigure{in der Klasse}{34,1 Milliarden Stellungen,\\ 11 Milliarden mit Hilfsmatt}
\end{minipage}\hfill
\begin{minipage}[t]{80mm}\vspace{0pt}\begin{tikzpicture}\begin{axis}[width=72mm,height=42mm,ybar,bar width=1.6pt,xmin=-1,xmax=17,ymin=0,ymax=2.376,axis y line*=left,axis x line*=bottom,xlabel={\scriptsize Halbzüge bis Matt},ylabel={\scriptsize Stellungen (Mrd.)},tick label style={font=\tiny},label style={font=\scriptsize},enlarge x limits=false,ylabel shift=-3pt,xlabel shift=-3pt]\addplot[fill=ink!55,draw=none] coordinates {(0,0.003131) (1,0.01895) (2,0.05503) (3,0.2613) (4,0.5303) (5,1.181) (6,1.661) (7,1.99) (8,2.16) (9,1.658) (10,1.028) (11,0.3308) (12,0.1394) (13,0.004164) (14,0.001197) (15,6.26e-07) (16,1e-09)};\draw[densely dotted] (axis cs:13,0) -- (axis cs:13,2.376);\end{axis}\begin{axis}[width=72mm,height=42mm,axis y line*=right,axis x line=none,ymode=log,xmin=-1,xmax=17,ylabel={\scriptsize eindeutig (log)},tick label style={font=\tiny},label style={font=\scriptsize},ylabel shift=-3pt]\addplot[mark=*,mark size=0.7pt,draw=black,thin] coordinates {(0,3131047) (1,18953319) (2,26097105) (3,84136564) (4,128122459) (5,104365013) (6,61709985) (7,23573903) (8,14802282) (9,4295920) (10,1636757) (11,49799) (12,16865) (13,11)};\end{axis}\end{tikzpicture}\end{minipage}
\vspace{3mm}

\noindent\begin{minipage}[t]{60mm}\vspace{0pt}\sbox\hmboard{\chessboard[setfen=4R1bk/8/8/8/8/p7/p7/K7 w - - 0 1]}\usebox\hmboard\\[1mm]
\makebox[\wd\hmboard]{\sffamily\hspace{3mm}\textbf{h\#6.5}\hfill(2+4)}\end{minipage}\hfill
\begin{minipage}[t]{52mm}\vspace{0pt}
\textcolor{ink}{\sffamily\textsc{Lösung}}\par\smallskip\small
\begin{tabular}{@{}r@{\enspace}l@{\enspace}l@{}}
\textbf{1.} & \ldots & \symrook{}e7\\
\textbf{2.} & \symbishop{}h7 & \symking{}xa2\\
\textbf{3.} & \symking{}g8 & \symking{}b3\\
\textbf{4.} & a2 & \symking{}c4\\
\textbf{5.} & a1=\symrook{} & \symking{}d5\\
\textbf{6.} & \symrook{}g1 & \symking{}e6\\
\textbf{7.} & \symrook{}g7 & \symrook{}e8\#\\
\end{tabular}\par\vspace{3mm}
\qrcode[height=12mm]{https://helpman.komtera.lt/?fen=4R1bk/8/8/8/8/p7/p7/K7&moves=6.5}\\{\sffamily\tiny\color{ink}scannen zum Lösen}
\end{minipage}
\vspace{2mm}

{\raggedright\noindent\textcolor{ink}{\sffamily\textsc{Themen}}\enspace \themetag{reines Matt} \themetag{Mustermatt} \themetag{Idealmatt} \themetag{Umwandlung} \themetag{Unterverwandlung} \themetag{Switchback} \themetag{Selbstblock} \themetag{ohne Schach} \themetag{Umnov}\par}

\clearpage
\phantomsection\label{hm:KRvknnn}
\begin{tcolorbox}[enhanced,colback=ink,colframe=ink,arc=0pt,boxrule=0pt,left=4mm,right=4mm,top=2.5mm,bottom=2.5mm]
\color{white}\sffamily{\LARGE\bfseries KRvknnn}\hfill{\Large h\#6.5}\\[0.5mm]
{\small 6 Steine\quad Weiß am Zug}\hfill{\small Nr.~635\qquad\href{https://helpman.komtera.lt/?fen=7k/5R2/8/8/3n4/2nn4/8/1K6&moves=6.5}{\color{white}\bfseries lösen\,\ensuremath{\nearrow}}}
\end{tcolorbox}
\vspace{2mm}
\noindent\begin{minipage}[t]{36mm}\vspace{0pt}\sffamily\small\raggedright
\hmfigure{längste eindeutige}{{\Large h\#6.5}}\hmfigure{eindeutig bei dieser Länge}{{\Large 1} Stellung}\hmfigure{längstes Matt}{{\Large h\#7.5}}\hmfigure{in der Klasse}{15,5 Milliarden Stellungen,\\ 1,59 Milliarden mit Hilfsmatt}
\end{minipage}\hfill
\begin{minipage}[t]{80mm}\vspace{0pt}\begin{tikzpicture}\begin{axis}[width=72mm,height=42mm,ybar,bar width=1.6pt,xmin=-1,xmax=16,ymin=0,ymax=330.6,axis y line*=left,axis x line*=bottom,xlabel={\scriptsize Halbzüge bis Matt},ylabel={\scriptsize Stellungen (Mio.)},tick label style={font=\tiny},label style={font=\scriptsize},enlarge x limits=false,ylabel shift=-3pt,xlabel shift=-3pt]\addplot[fill=ink!55,draw=none] coordinates {(0,0.3479) (1,2.338) (2,7.015) (3,33.27) (4,65.87) (5,159.1) (6,214.9) (7,290.9) (8,300.6) (9,259.3) (10,165.1) (11,62.47) (12,23.25) (13,0.9471) (14,0.2319) (15,0.000194)};\draw[densely dotted] (axis cs:13,0) -- (axis cs:13,330.6);\end{axis}\begin{axis}[width=72mm,height=42mm,axis y line*=right,axis x line=none,ymode=log,xmin=-1,xmax=16,ylabel={\scriptsize eindeutig (log)},tick label style={font=\tiny},label style={font=\scriptsize},ylabel shift=-3pt]\addplot[mark=*,mark size=0.7pt,draw=black,thin] coordinates {(0,347945) (1,2337625) (2,3295002) (3,11393297) (4,15913139) (5,15604507) (6,9341836) (7,4276755) (8,2647026) (9,858999) (10,291212) (11,4883) (12,1057) (13,1)};\end{axis}\end{tikzpicture}\end{minipage}
\vspace{3mm}

\noindent\begin{minipage}[t]{60mm}\vspace{0pt}\sbox\hmboard{\chessboard[setfen=7k/5R2/8/8/3n4/2nn4/8/1K6 w - - 0 1]}\usebox\hmboard\\[1mm]
\makebox[\wd\hmboard]{\sffamily\hspace{3mm}\textbf{h\#6.5}\hfill(2+4)}\end{minipage}\hfill
\begin{minipage}[t]{52mm}\vspace{0pt}
\textcolor{ink}{\sffamily\textsc{Lösung}}\par\smallskip\small
\begin{tabular}{@{}r@{\enspace}l@{\enspace}l@{}}
\textbf{1.} & \ldots & \symking{}a1\\
\textbf{2.} & \symknight{}f4 & \symking{}b2\\
\textbf{3.} & \symknight{}de6 & \symking{}xc3\\
\textbf{4.} & \symknight{}f8 & \symking{}d4\\
\textbf{5.} & \symking{}g8 & \symking{}e5\\
\textbf{6.} & \symknight{}4g6+ & \symking{}f6\\
\textbf{7.} & \symknight{}h8 & \symrook{}g7\#\\
\end{tabular}\par\vspace{3mm}
\qrcode[height=12mm]{https://helpman.komtera.lt/?fen=7k/5R2/8/8/3n4/2nn4/8/1K6&moves=6.5}\\{\sffamily\tiny\color{ink}scannen zum Lösen}
\end{minipage}
\vspace{2mm}

{\raggedright\noindent\textcolor{ink}{\sffamily\textsc{Themen}}\enspace \themetag{Selbstblock}\par\smallskip{\footnotesize\textit{Hinweis:} die Partei am Zug steht im Diagramm im Schach.}\par}

\clearpage
\phantomsection\label{hm:KRvknpp}
\begin{tcolorbox}[enhanced,colback=ink,colframe=ink,arc=0pt,boxrule=0pt,left=4mm,right=4mm,top=2.5mm,bottom=2.5mm]
\color{white}\sffamily{\LARGE\bfseries KRvknpp}\hfill{\Large h\#6.5}\\[0.5mm]
{\small 6 Steine\quad Weiß am Zug}\hfill{\small Nr.~636\qquad\href{https://helpman.komtera.lt/?fen=8/4R3/7k/8/8/2p5/2p5/K1n5&moves=6.5}{\color{white}\bfseries lösen\,\ensuremath{\nearrow}}}
\end{tcolorbox}
\vspace{2mm}
\noindent\begin{minipage}[t]{36mm}\vspace{0pt}\sffamily\small\raggedright
\hmfigure{längste eindeutige}{{\Large h\#6.5}}\hmfigure{eindeutig bei dieser Länge}{{\Large 2} Stellungen}\hmfigure{längstes Matt}{{\Large h\#7.5}}\hmfigure{in der Klasse}{34,1 Milliarden Stellungen,\\ 11,3 Milliarden mit Hilfsmatt}
\end{minipage}\hfill
\begin{minipage}[t]{80mm}\vspace{0pt}\begin{tikzpicture}\begin{axis}[width=72mm,height=42mm,ybar,bar width=1.6pt,xmin=-1,xmax=16,ymin=0,ymax=2.479,axis y line*=left,axis x line*=bottom,xlabel={\scriptsize Halbzüge bis Matt},ylabel={\scriptsize Stellungen (Mrd.)},tick label style={font=\tiny},label style={font=\scriptsize},enlarge x limits=false,ylabel shift=-3pt,xlabel shift=-3pt]\addplot[fill=ink!55,draw=none] coordinates {(0,0.004186) (1,0.02518) (2,0.06357) (3,0.2787) (4,0.5818) (5,1.195) (6,1.772) (7,1.986) (8,2.253) (9,1.633) (10,1.057) (11,0.3229) (12,0.1405) (13,0.004314) (14,0.001217) (15,4.4e-08)};\draw[densely dotted] (axis cs:13,0) -- (axis cs:13,2.479);\end{axis}\begin{axis}[width=72mm,height=42mm,axis y line*=right,axis x line=none,ymode=log,xmin=-1,xmax=16,ylabel={\scriptsize eindeutig (log)},tick label style={font=\tiny},label style={font=\scriptsize},ylabel shift=-3pt]\addplot[mark=*,mark size=0.7pt,draw=black,thin] coordinates {(0,4186382) (1,25179446) (2,32161855) (3,99782360) (4,161689315) (5,109896376) (6,66662434) (7,21079876) (8,13054529) (9,3513092) (10,1294298) (11,49513) (12,24250) (13,2)};\end{axis}\end{tikzpicture}\end{minipage}
\vspace{3mm}

\noindent\begin{minipage}[t]{60mm}\vspace{0pt}\sbox\hmboard{\chessboard[setfen=8/4R3/7k/8/8/2p5/2p5/K1n5 w - - 0 1]}\usebox\hmboard\\[1mm]
\makebox[\wd\hmboard]{\sffamily\hspace{3mm}\textbf{h\#6.5}\hfill(2+4)}\end{minipage}\hfill
\begin{minipage}[t]{52mm}\vspace{0pt}
\textcolor{ink}{\sffamily\textsc{Lösung}}\par\smallskip\small
\begin{tabular}{@{}r@{\enspace}l@{\enspace}l@{}}
\textbf{1.} & \ldots & \symrook{}c7\\
\textbf{2.} & \symknight{}d3 & \symking{}a2\\
\textbf{3.} & \symknight{}e5 & \symking{}b3\\
\textbf{4.} & \symknight{}d7 & \symking{}c4\\
\textbf{5.} & \symking{}g7 & \symking{}d5\\
\textbf{6.} & \symking{}f8 & \symking{}e6\\
\textbf{7.} & \symking{}e8 & \symrook{}c8\#\\
\end{tabular}\par\vspace{3mm}
\qrcode[height=12mm]{https://helpman.komtera.lt/?fen=8/4R3/7k/8/8/2p5/2p5/K1n5&moves=6.5}\\{\sffamily\tiny\color{ink}scannen zum Lösen}
\end{minipage}
\vspace{2mm}

{\raggedright\noindent\textcolor{ink}{\sffamily\textsc{Themen}}\enspace \themetag{ohne Schlag} \themetag{ohne Schach}\par}

\clearpage
\phantomsection\label{hm:KRvkqbb}
\begin{tcolorbox}[enhanced,colback=ink,colframe=ink,arc=0pt,boxrule=0pt,left=4mm,right=4mm,top=2.5mm,bottom=2.5mm]
\color{white}\sffamily{\LARGE\bfseries KRvkqbb}\hfill{\Large h\#6.5}\\[0.5mm]
{\small 6 Steine\quad Weiß am Zug}\hfill{\small Nr.~637\qquad\href{https://helpman.komtera.lt/?fen=4R3/8/7k/8/4b3/1q6/8/K1b5&moves=6.5}{\color{white}\bfseries lösen\,\ensuremath{\nearrow}}}
\end{tcolorbox}
\vspace{2mm}
\noindent\begin{minipage}[t]{36mm}\vspace{0pt}\sffamily\small\raggedright
\hmfigure{längste eindeutige}{{\Large h\#6.5}}\hmfigure{eindeutig bei dieser Länge}{{\Large 160} Stellungen}\hmfigure{längstes Matt}{{\Large h\#8}}\hmfigure{in der Klasse}{15,5 Milliarden Stellungen,\\ 3,92 Milliarden mit Hilfsmatt}
\end{minipage}\hfill
\begin{minipage}[t]{80mm}\vspace{0pt}\begin{tikzpicture}\begin{axis}[width=72mm,height=42mm,ybar,bar width=1.6pt,xmin=-1,xmax=17,ymin=0,ymax=948.8,axis y line*=left,axis x line*=bottom,xlabel={\scriptsize Halbzüge bis Matt},ylabel={\scriptsize Stellungen (Mio.)},tick label style={font=\tiny},label style={font=\scriptsize},enlarge x limits=false,ylabel shift=-3pt,xlabel shift=-3pt]\addplot[fill=ink!55,draw=none] coordinates {(0,0.2162) (1,1.554) (2,7.546) (3,52.84) (4,96) (5,357.5) (6,374.9) (7,824.9) (8,601.4) (9,862.5) (10,380.9) (11,294.7) (12,55.52) (13,10.41) (14,0.6164) (15,0.03818) (16,2e-06)};\draw[densely dotted] (axis cs:13,0) -- (axis cs:13,948.8);\end{axis}\begin{axis}[width=72mm,height=42mm,axis y line*=right,axis x line=none,ymode=log,xmin=-1,xmax=17,ylabel={\scriptsize eindeutig (log)},tick label style={font=\tiny},label style={font=\scriptsize},ylabel shift=-3pt]\addplot[mark=*,mark size=0.7pt,draw=black,thin] coordinates {(0,216161) (1,1554268) (2,3157507) (3,16375468) (4,15306552) (5,31872150) (6,13844836) (7,16077044) (8,5920368) (9,4745395) (10,850136) (11,154587) (12,1555) (13,160)};\end{axis}\end{tikzpicture}\end{minipage}
\vspace{3mm}

\noindent\begin{minipage}[t]{60mm}\vspace{0pt}\sbox\hmboard{\chessboard[setfen=4R3/8/7k/8/4b3/1q6/8/K1b5 w - - 0 1]}\usebox\hmboard\\[1mm]
\makebox[\wd\hmboard]{\sffamily\hspace{3mm}\textbf{h\#6.5}\hfill(2+4)}\end{minipage}\hfill
\begin{minipage}[t]{52mm}\vspace{0pt}
\textcolor{ink}{\sffamily\textsc{Lösung}}\par\smallskip\small
\begin{tabular}{@{}r@{\enspace}l@{\enspace}l@{}}
\textbf{1.} & \ldots & \symrook{}a8\\
\textbf{2.} & \symqueen{}b8 & \symking{}a2\\
\textbf{3.} & \symbishop{}b7 & \symking{}b3\\
\textbf{4.} & \symking{}g7 & \symking{}c4\\
\textbf{5.} & \symbishop{}d5+ & \symking{}xd5\\
\textbf{6.} & \symking{}f8 & \symking{}e6\\
\textbf{7.} & \symking{}e8 & \symrook{}xb8\#\\
\end{tabular}\par\vspace{3mm}
\qrcode[height=12mm]{https://helpman.komtera.lt/?fen=4R3/8/7k/8/4b3/1q6/8/K1b5&moves=6.5}\\{\sffamily\tiny\color{ink}scannen zum Lösen}
\end{minipage}
\vspace{2mm}

{\raggedright\noindent\textcolor{ink}{\sffamily\textsc{Themen}}\enspace \themetag{reines Matt} \themetag{Mustermatt} \themetag{Spiegelmatt} \themetag{Maslar}\par}

\clearpage
\phantomsection\label{hm:KRvkqbp}
\begin{tcolorbox}[enhanced,colback=ink,colframe=ink,arc=0pt,boxrule=0pt,left=4mm,right=4mm,top=2.5mm,bottom=2.5mm]
\color{white}\sffamily{\LARGE\bfseries KRvkqbp}\hfill{\Large h\#6.5}\\[0.5mm]
{\small 6 Steine\quad Weiß am Zug}\hfill{\small Nr.~638\qquad\href{https://helpman.komtera.lt/?fen=5k2/8/8/3b3R/8/p7/2q5/K7&moves=6.5}{\color{white}\bfseries lösen\,\ensuremath{\nearrow}}}
\end{tcolorbox}
\vspace{2mm}
\noindent\begin{minipage}[t]{36mm}\vspace{0pt}\sffamily\small\raggedright
\hmfigure{längste eindeutige}{{\Large h\#6.5}}\hmfigure{eindeutig bei dieser Länge}{{\Large 855} Stellungen}\hmfigure{längstes Matt}{{\Large h\#8}}\hmfigure{in der Klasse}{45,4 Milliarden Stellungen,\\ 24,8 Milliarden mit Hilfsmatt}
\end{minipage}\hfill
\begin{minipage}[t]{80mm}\vspace{0pt}\begin{tikzpicture}\begin{axis}[width=72mm,height=42mm,ybar,bar width=1.6pt,xmin=-1,xmax=17,ymin=0,ymax=5.632,axis y line*=left,axis x line*=bottom,xlabel={\scriptsize Halbzüge bis Matt},ylabel={\scriptsize Stellungen (Mrd.)},tick label style={font=\tiny},label style={font=\scriptsize},enlarge x limits=false,ylabel shift=-3pt,xlabel shift=-3pt]\addplot[fill=ink!55,draw=none] coordinates {(0,0.002286) (1,0.01559) (2,0.07029) (3,0.435) (4,0.7807) (5,2.502) (6,2.764) (7,5.12) (8,4.063) (9,4.946) (10,2.328) (11,1.453) (12,0.3172) (13,0.04271) (14,0.003283) (15,8.42e-05) (16,1e-09)};\draw[densely dotted] (axis cs:13,0) -- (axis cs:13,5.632);\end{axis}\begin{axis}[width=72mm,height=42mm,axis y line*=right,axis x line=none,ymode=log,xmin=-1,xmax=17,ylabel={\scriptsize eindeutig (log)},tick label style={font=\tiny},label style={font=\scriptsize},ylabel shift=-3pt]\addplot[mark=*,mark size=0.7pt,draw=black,thin] coordinates {(0,2285922) (1,15589754) (2,30302473) (3,136344162) (4,132330005) (5,218482978) (6,103946151) (7,92380544) (8,38629219) (9,23161361) (10,4373760) (11,499877) (12,13288) (13,855)};\end{axis}\end{tikzpicture}\end{minipage}
\vspace{3mm}

\noindent\begin{minipage}[t]{60mm}\vspace{0pt}\sbox\hmboard{\chessboard[setfen=5k2/8/8/3b3R/8/p7/2q5/K7 w - - 0 1]}\usebox\hmboard\\[1mm]
\makebox[\wd\hmboard]{\sffamily\hspace{3mm}\textbf{h\#6.5}\hfill(2+4)}\end{minipage}\hfill
\begin{minipage}[t]{52mm}\vspace{0pt}
\textcolor{ink}{\sffamily\textsc{Lösung}}\par\smallskip\small
\begin{tabular}{@{}r@{\enspace}l@{\enspace}l@{}}
\textbf{1.} & \ldots & \symrook{}h1\\
\textbf{2.} & \symqueen{}h2 & \symking{}b1\\
\textbf{3.} & \symbishop{}g2 & \symking{}c2\\
\textbf{4.} & \symking{}g7 & \symking{}d3\\
\textbf{5.} & \symbishop{}e4+ & \symking{}xe4\\
\textbf{6.} & \symking{}h6 & \symking{}f5\\
\textbf{7.} & \symking{}h5 & \symrook{}xh2\#\\
\end{tabular}\par\vspace{3mm}
\qrcode[height=12mm]{https://helpman.komtera.lt/?fen=5k2/8/8/3b3R/8/p7/2q5/K7&moves=6.5}\\{\sffamily\tiny\color{ink}scannen zum Lösen}
\end{minipage}
\vspace{2mm}

{\raggedright\noindent\textcolor{ink}{\sffamily\textsc{Themen}}\enspace \themetag{reines Matt} \themetag{Mustermatt} \themetag{Spiegelmatt} \themetag{Maslar}\par}

\clearpage
\phantomsection\label{hm:KRvkqpp}
\begin{tcolorbox}[enhanced,colback=ink,colframe=ink,arc=0pt,boxrule=0pt,left=4mm,right=4mm,top=2.5mm,bottom=2.5mm]
\color{white}\sffamily{\LARGE\bfseries KRvkqpp}\hfill{\Large h\#6.5}\\[0.5mm]
{\small 6 Steine\quad Weiß am Zug}\hfill{\small Nr.~639\qquad\href{https://helpman.komtera.lt/?fen=6k1/5R2/8/8/8/1p6/1p4q1/K7&moves=6.5}{\color{white}\bfseries lösen\,\ensuremath{\nearrow}}}
\end{tcolorbox}
\vspace{2mm}
\noindent\begin{minipage}[t]{36mm}\vspace{0pt}\sffamily\small\raggedright
\hmfigure{längste eindeutige}{{\Large h\#6.5}}\hmfigure{eindeutig bei dieser Länge}{{\Large 346} Stellungen}\hmfigure{längstes Matt}{{\Large h\#7.5}}\hmfigure{in der Klasse}{34,1 Milliarden Stellungen,\\ 9,76 Milliarden mit Hilfsmatt}
\end{minipage}\hfill
\begin{minipage}[t]{80mm}\vspace{0pt}\begin{tikzpicture}\begin{axis}[width=72mm,height=42mm,ybar,bar width=1.6pt,xmin=-1,xmax=16,ymin=0,ymax=2.142,axis y line*=left,axis x line*=bottom,xlabel={\scriptsize Halbzüge bis Matt},ylabel={\scriptsize Stellungen (Mrd.)},tick label style={font=\tiny},label style={font=\scriptsize},enlarge x limits=false,ylabel shift=-3pt,xlabel shift=-3pt]\addplot[fill=ink!55,draw=none] coordinates {(0,0.001492) (1,0.009646) (2,0.03955) (3,0.2133) (4,0.3782) (5,1.04) (6,1.239) (7,1.947) (8,1.685) (9,1.759) (10,0.8727) (11,0.4529) (12,0.1143) (13,0.01166) (14,0.001008) (15,8.469e-06)};\draw[densely dotted] (axis cs:13,0) -- (axis cs:13,2.142);\end{axis}\begin{axis}[width=72mm,height=42mm,axis y line*=right,axis x line=none,ymode=log,xmin=-1,xmax=16,ylabel={\scriptsize eindeutig (log)},tick label style={font=\tiny},label style={font=\scriptsize},ylabel shift=-3pt]\addplot[mark=*,mark size=0.7pt,draw=black,thin] coordinates {(0,1491590) (1,9645465) (2,17619712) (3,68242724) (4,70052857) (5,91146745) (6,46470655) (7,31393609) (8,14572518) (9,6475672) (10,1314718) (11,86989) (12,5227) (13,346)};\end{axis}\end{tikzpicture}\end{minipage}
\vspace{3mm}

\noindent\begin{minipage}[t]{60mm}\vspace{0pt}\sbox\hmboard{\chessboard[setfen=6k1/5R2/8/8/8/1p6/1p4q1/K7 w - - 0 1]}\usebox\hmboard\\[1mm]
\makebox[\wd\hmboard]{\sffamily\hspace{3mm}\textbf{h\#6.5}\hfill(2+4)}\end{minipage}\hfill
\begin{minipage}[t]{52mm}\vspace{0pt}
\textcolor{ink}{\sffamily\textsc{Lösung}}\par\smallskip\small
\begin{tabular}{@{}r@{\enspace}l@{\enspace}l@{}}
\textbf{1.} & \ldots & \symking{}b1\\
\textbf{2.} & \symqueen{}b7 & \symking{}xb2\\
\textbf{3.} & \symking{}h8 & \symking{}c3\\
\textbf{4.} & b2 & \symking{}d4\\
\textbf{5.} & b1=\symrook{} & \symking{}e5\\
\textbf{6.} & \symrook{}h1 & \symking{}f6\\
\textbf{7.} & \symrook{}h7 & \symrook{}f8\#\\
\end{tabular}\par\vspace{3mm}
\qrcode[height=12mm]{https://helpman.komtera.lt/?fen=6k1/5R2/8/8/8/1p6/1p4q1/K7&moves=6.5}\\{\sffamily\tiny\color{ink}scannen zum Lösen}
\end{minipage}
\vspace{2mm}

{\raggedright\noindent\textcolor{ink}{\sffamily\textsc{Themen}}\enspace \themetag{reines Matt} \themetag{Mustermatt} \themetag{Umwandlung} \themetag{Unterverwandlung} \themetag{Selbstblock} \themetag{ohne Schach} \themetag{Umnov}\par\smallskip{\footnotesize\textit{Hinweis:} die Partei am Zug steht im Diagramm im Schach.}\par}

\clearpage
\phantomsection\label{hm:KRvkqrn}
\begin{tcolorbox}[enhanced,colback=ink,colframe=ink,arc=0pt,boxrule=0pt,left=4mm,right=4mm,top=2.5mm,bottom=2.5mm]
\color{white}\sffamily{\LARGE\bfseries KRvkqrn}\hfill{\Large h\#6.5}\\[0.5mm]
{\small 6 Steine\quad Weiß am Zug}\hfill{\small Nr.~640\qquad\href{https://helpman.komtera.lt/?fen=8/8/7k/6R1/n7/8/1rq5/K7&moves=6.5}{\color{white}\bfseries lösen\,\ensuremath{\nearrow}}}
\end{tcolorbox}
\vspace{2mm}
\noindent\begin{minipage}[t]{36mm}\vspace{0pt}\sffamily\small\raggedright
\hmfigure{längste eindeutige}{{\Large h\#6.5}}\hmfigure{eindeutig bei dieser Länge}{{\Large 188} Stellungen}\hmfigure{längstes Matt}{{\Large h\#8}}\hmfigure{in der Klasse}{15,5 Milliarden Stellungen,\\ 7,74 Milliarden mit Hilfsmatt}
\end{minipage}\hfill
\begin{minipage}[t]{80mm}\vspace{0pt}\begin{tikzpicture}\begin{axis}[width=72mm,height=42mm,ybar,bar width=1.6pt,xmin=-1,xmax=17,ymin=0,ymax=1.886,axis y line*=left,axis x line*=bottom,xlabel={\scriptsize Halbzüge bis Matt},ylabel={\scriptsize Stellungen (Mrd.)},tick label style={font=\tiny},label style={font=\scriptsize},enlarge x limits=false,ylabel shift=-3pt,xlabel shift=-3pt]\addplot[fill=ink!55,draw=none] coordinates {(0,0.0004478) (1,0.003443) (2,0.01757) (3,0.1172) (4,0.2177) (5,0.7457) (6,0.8251) (7,1.697) (8,1.236) (9,1.714) (10,0.6095) (11,0.4866) (12,0.05295) (13,0.01422) (14,0.0003169) (15,2.388e-05) (16,1.66e-07)};\draw[densely dotted] (axis cs:13,0) -- (axis cs:13,1.886);\end{axis}\begin{axis}[width=72mm,height=42mm,axis y line*=right,axis x line=none,ymode=log,xmin=-1,xmax=17,ylabel={\scriptsize eindeutig (log)},tick label style={font=\tiny},label style={font=\scriptsize},ylabel shift=-3pt]\addplot[mark=*,mark size=0.7pt,draw=black,thin] coordinates {(0,447785) (1,3443450) (2,7647605) (3,37560780) (4,38322377) (5,68366911) (6,31746166) (7,34415841) (8,10771464) (9,8467990) (10,1032937) (11,269179) (12,4830) (13,188)};\end{axis}\end{tikzpicture}\end{minipage}
\vspace{3mm}

\noindent\begin{minipage}[t]{60mm}\vspace{0pt}\sbox\hmboard{\chessboard[setfen=8/8/7k/6R1/n7/8/1rq5/K7 w - - 0 1]}\usebox\hmboard\\[1mm]
\makebox[\wd\hmboard]{\sffamily\hspace{3mm}\textbf{h\#6.5}\hfill(2+4)}\end{minipage}\hfill
\begin{minipage}[t]{52mm}\vspace{0pt}
\textcolor{ink}{\sffamily\textsc{Lösung}}\par\smallskip\small
\begin{tabular}{@{}r@{\enspace}l@{\enspace}l@{}}
\textbf{1.} & \ldots & \symrook{}g2\\
\textbf{2.} & \symking{}h5 & \symrook{}xc2\\
\textbf{3.} & \symking{}g4 & \symrook{}c4+\\
\textbf{4.} & \symking{}f3 & \symrook{}xa4\\
\textbf{5.} & \symking{}e2 & \symking{}xb2\\
\textbf{6.} & \symking{}d1 & \symking{}c3\\
\textbf{7.} & \symking{}c1 & \symrook{}a1\#\\
\end{tabular}\par\vspace{3mm}
\qrcode[height=12mm]{https://helpman.komtera.lt/?fen=8/8/7k/6R1/n7/8/1rq5/K7&moves=6.5}\\{\sffamily\tiny\color{ink}scannen zum Lösen}
\end{minipage}
\vspace{2mm}

{\raggedright\noindent\textcolor{ink}{\sffamily\textsc{Themen}}\enspace \themetag{reines Matt} \themetag{Mustermatt} \themetag{Idealmatt} \themetag{Spiegelmatt} \themetag{ein Stein zieht}\par}

\clearpage
\phantomsection\label{hm:KRvkrbb}
\begin{tcolorbox}[enhanced,colback=ink,colframe=ink,arc=0pt,boxrule=0pt,left=4mm,right=4mm,top=2.5mm,bottom=2.5mm]
\color{white}\sffamily{\LARGE\bfseries KRvkrbb}\hfill{\Large h\#6.5}\\[0.5mm]
{\small 6 Steine\quad Weiß am Zug}\hfill{\small Nr.~641\qquad\href{https://helpman.komtera.lt/?fen=8/8/6k1/5b2/4R3/b7/6r1/1K6&moves=6.5}{\color{white}\bfseries lösen\,\ensuremath{\nearrow}}}
\end{tcolorbox}
\vspace{2mm}
\noindent\begin{minipage}[t]{36mm}\vspace{0pt}\sffamily\small\raggedright
\hmfigure{längste eindeutige}{{\Large h\#6.5}}\hmfigure{eindeutig bei dieser Länge}{{\Large 19} Stellungen}\hmfigure{längstes Matt}{{\Large h\#7.5}}\hmfigure{in der Klasse}{15,5 Milliarden Stellungen,\\ 4,26 Milliarden mit Hilfsmatt}
\end{minipage}\hfill
\begin{minipage}[t]{80mm}\vspace{0pt}\begin{tikzpicture}\begin{axis}[width=72mm,height=42mm,ybar,bar width=1.6pt,xmin=-1,xmax=16,ymin=0,ymax=946.3,axis y line*=left,axis x line*=bottom,xlabel={\scriptsize Halbzüge bis Matt},ylabel={\scriptsize Stellungen (Mio.)},tick label style={font=\tiny},label style={font=\scriptsize},enlarge x limits=false,ylabel shift=-3pt,xlabel shift=-3pt]\addplot[fill=ink!55,draw=none] coordinates {(0,0.3867) (1,2.712) (2,11.55) (3,71.33) (4,137.5) (5,418.3) (6,495.1) (7,860.3) (8,721.1) (9,826.4) (10,417.2) (11,235.9) (12,56.4) (13,5.036) (14,0.4634) (15,0.005666)};\draw[densely dotted] (axis cs:13,0) -- (axis cs:13,946.3);\end{axis}\begin{axis}[width=72mm,height=42mm,axis y line*=right,axis x line=none,ymode=log,xmin=-1,xmax=16,ylabel={\scriptsize eindeutig (log)},tick label style={font=\tiny},label style={font=\scriptsize},ylabel shift=-3pt]\addplot[mark=*,mark size=0.7pt,draw=black,thin] coordinates {(0,386677) (1,2711768) (2,5015906) (3,21619665) (4,24549579) (5,35052884) (6,17901897) (7,14541533) (8,6561469) (9,3723754) (10,898006) (11,76963) (12,2785) (13,19)};\end{axis}\end{tikzpicture}\end{minipage}
\vspace{3mm}

\noindent\begin{minipage}[t]{60mm}\vspace{0pt}\sbox\hmboard{\chessboard[setfen=8/8/6k1/5b2/4R3/b7/6r1/1K6 w - - 0 1]}\usebox\hmboard\\[1mm]
\makebox[\wd\hmboard]{\sffamily\hspace{3mm}\textbf{h\#6.5}\hfill(2+4)}\end{minipage}\hfill
\begin{minipage}[t]{52mm}\vspace{0pt}
\textcolor{ink}{\sffamily\textsc{Lösung}}\par\smallskip\small
\begin{tabular}{@{}r@{\enspace}l@{\enspace}l@{}}
\textbf{1.} & \ldots & \symking{}a1\\
\textbf{2.} & \symking{}g5 & \symrook{}e5\\
\textbf{3.} & \symking{}f4 & \symrook{}xf5+\\
\textbf{4.} & \symking{}e3 & \symking{}b1\\
\textbf{5.} & \symking{}e2 & \symking{}c2\\
\textbf{6.} & \symking{}e1+ & \symking{}d3\\
\textbf{7.} & \symking{}d1 & \symrook{}f1\#\\
\end{tabular}\par\vspace{3mm}
\qrcode[height=12mm]{https://helpman.komtera.lt/?fen=8/8/6k1/5b2/4R3/b7/6r1/1K6&moves=6.5}\\{\sffamily\tiny\color{ink}scannen zum Lösen}
\end{minipage}
\vspace{2mm}

{\raggedright\noindent\textcolor{ink}{\sffamily\textsc{Themen}}\enspace \themetag{Satzspiel} \themetag{reines Matt} \themetag{Mustermatt} \themetag{Spiegelmatt} \themetag{Switchback} \themetag{ein Stein zieht}\par}

\clearpage
\phantomsection\label{hm:KRvkrbn}
\begin{tcolorbox}[enhanced,colback=ink,colframe=ink,arc=0pt,boxrule=0pt,left=4mm,right=4mm,top=2.5mm,bottom=2.5mm]
\color{white}\sffamily{\LARGE\bfseries KRvkrbn}\hfill{\Large h\#6.5}\\[0.5mm]
{\small 6 Steine\quad Weiß am Zug}\hfill{\small Nr.~642\qquad\href{https://helpman.komtera.lt/?fen=8/8/5b2/8/4k3/n7/2r5/K2R4&moves=6.5}{\color{white}\bfseries lösen\,\ensuremath{\nearrow}}}
\end{tcolorbox}
\vspace{2mm}
\noindent\begin{minipage}[t]{36mm}\vspace{0pt}\sffamily\small\raggedright
\hmfigure{längste eindeutige}{{\Large h\#6.5}}\hmfigure{eindeutig bei dieser Länge}{{\Large 14} Stellungen}\hmfigure{längstes Matt}{{\Large h\#8}}\hmfigure{in der Klasse}{15,5 Milliarden Stellungen,\\ 8,74 Milliarden mit Hilfsmatt}
\end{minipage}\hfill
\begin{minipage}[t]{80mm}\vspace{0pt}\begin{tikzpicture}\begin{axis}[width=72mm,height=42mm,ybar,bar width=1.6pt,xmin=-1,xmax=17,ymin=0,ymax=1.965,axis y line*=left,axis x line*=bottom,xlabel={\scriptsize Halbzüge bis Matt},ylabel={\scriptsize Stellungen (Mrd.)},tick label style={font=\tiny},label style={font=\scriptsize},enlarge x limits=false,ylabel shift=-3pt,xlabel shift=-3pt]\addplot[fill=ink!55,draw=none] coordinates {(0,0.0009941) (1,0.007022) (2,0.02857) (3,0.1631) (4,0.324) (5,0.8953) (6,1.136) (7,1.787) (8,1.59) (9,1.616) (10,0.7546) (11,0.3707) (12,0.06441) (13,0.004806) (14,0.000222) (15,2.279e-06) (16,4e-09)};\draw[densely dotted] (axis cs:13,0) -- (axis cs:13,1.965);\end{axis}\begin{axis}[width=72mm,height=42mm,axis y line*=right,axis x line=none,ymode=log,xmin=-1,xmax=17,ylabel={\scriptsize eindeutig (log)},tick label style={font=\tiny},label style={font=\scriptsize},ylabel shift=-3pt]\addplot[mark=*,mark size=0.7pt,draw=black,thin] coordinates {(0,994119) (1,7022389) (2,12794431) (3,51207431) (4,64788944) (5,77370031) (6,42795864) (7,29481136) (8,13065005) (9,6117077) (10,1182088) (11,79083) (12,2064) (13,14)};\end{axis}\end{tikzpicture}\end{minipage}
\vspace{3mm}

\noindent\begin{minipage}[t]{60mm}\vspace{0pt}\sbox\hmboard{\chessboard[setfen=8/8/5b2/8/4k3/n7/2r5/K2R4 w - - 0 1]}\usebox\hmboard\\[1mm]
\makebox[\wd\hmboard]{\sffamily\hspace{3mm}\textbf{h\#6.5}\hfill(2+4)}\end{minipage}\hfill
\begin{minipage}[t]{52mm}\vspace{0pt}
\textcolor{ink}{\sffamily\textsc{Lösung}}\par\smallskip\small
\begin{tabular}{@{}r@{\enspace}l@{\enspace}l@{}}
\textbf{1.} & \ldots & \symrook{}d4+\\
\textbf{2.} & \symking{}e5 & \symrook{}c4\\
\textbf{3.} & \symknight{}b5 & \symking{}b1\\
\textbf{4.} & \symrook{}a2 & \symking{}c1\\
\textbf{5.} & \symrook{}a6 & \symking{}d2\\
\textbf{6.} & \symrook{}e6 & \symking{}e3\\
\textbf{7.} & \symknight{}d6 & \symrook{}c5\#\\
\end{tabular}\par\vspace{3mm}
\qrcode[height=12mm]{https://helpman.komtera.lt/?fen=8/8/5b2/8/4k3/n7/2r5/K2R4&moves=6.5}\\{\sffamily\tiny\color{ink}scannen zum Lösen}
\end{minipage}
\vspace{2mm}

{\raggedright\noindent\textcolor{ink}{\sffamily\textsc{Themen}}\enspace \themetag{reines Matt} \themetag{Mustermatt} \themetag{Idealmatt} \themetag{Selbstblock} \themetag{ohne Schlag}\par\smallskip{\footnotesize\textit{Hinweis:} die Partei am Zug steht im Diagramm im Schach.}\par}

\clearpage
\phantomsection\label{hm:KRvkrbp}
\begin{tcolorbox}[enhanced,colback=ink,colframe=ink,arc=0pt,boxrule=0pt,left=4mm,right=4mm,top=2.5mm,bottom=2.5mm]
\color{white}\sffamily{\LARGE\bfseries KRvkrbp}\hfill{\Large h\#6.5}\\[0.5mm]
{\small 6 Steine\quad Weiß am Zug}\hfill{\small Nr.~643\qquad\href{https://helpman.komtera.lt/?fen=8/3R4/5b2/8/4k3/8/p2r4/K7&moves=6.5}{\color{white}\bfseries lösen\,\ensuremath{\nearrow}}}
\end{tcolorbox}
\vspace{2mm}
\noindent\begin{minipage}[t]{36mm}\vspace{0pt}\sffamily\small\raggedright
\hmfigure{längste eindeutige}{{\Large h\#6.5}}\hmfigure{eindeutig bei dieser Länge}{{\Large 37} Stellungen}\hmfigure{längstes Matt}{{\Large h\#7.5}}\hmfigure{in der Klasse}{45,4 Milliarden Stellungen,\\ 26,9 Milliarden mit Hilfsmatt}
\end{minipage}\hfill
\begin{minipage}[t]{80mm}\vspace{0pt}\begin{tikzpicture}\begin{axis}[width=72mm,height=42mm,ybar,bar width=1.6pt,xmin=-1,xmax=16,ymin=0,ymax=5.792,axis y line*=left,axis x line*=bottom,xlabel={\scriptsize Halbzüge bis Matt},ylabel={\scriptsize Stellungen (Mrd.)},tick label style={font=\tiny},label style={font=\scriptsize},enlarge x limits=false,ylabel shift=-3pt,xlabel shift=-3pt]\addplot[fill=ink!55,draw=none] coordinates {(0,0.004018) (1,0.02668) (2,0.1026) (3,0.5555) (4,1.079) (5,2.804) (6,3.555) (7,5.266) (8,4.845) (9,4.736) (10,2.493) (11,1.152) (12,0.3048) (13,0.01987) (14,0.002187) (15,2.649e-06)};\draw[densely dotted] (axis cs:13,0) -- (axis cs:13,5.792);\end{axis}\begin{axis}[width=72mm,height=42mm,axis y line*=right,axis x line=none,ymode=log,xmin=-1,xmax=16,ylabel={\scriptsize eindeutig (log)},tick label style={font=\tiny},label style={font=\scriptsize},ylabel shift=-3pt]\addplot[mark=*,mark size=0.7pt,draw=black,thin] coordinates {(0,4017571) (1,26678398) (2,47002976) (3,170714386) (4,207036631) (5,227703123) (6,127393400) (7,79030561) (8,39614372) (9,17306700) (10,4509982) (11,250968) (12,14497) (13,37)};\end{axis}\end{tikzpicture}\end{minipage}
\vspace{3mm}

\noindent\begin{minipage}[t]{60mm}\vspace{0pt}\sbox\hmboard{\chessboard[setfen=8/3R4/5b2/8/4k3/8/p2r4/K7 w - - 0 1]}\usebox\hmboard\\[1mm]
\makebox[\wd\hmboard]{\sffamily\hspace{3mm}\textbf{h\#6.5}\hfill(2+4)}\end{minipage}\hfill
\begin{minipage}[t]{52mm}\vspace{0pt}
\textcolor{ink}{\sffamily\textsc{Lösung}}\par\smallskip\small
\begin{tabular}{@{}r@{\enspace}l@{\enspace}l@{}}
\textbf{1.} & \ldots & \symrook{}d4+\\
\textbf{2.} & \symking{}e5 & \symrook{}d6\\
\textbf{3.} & \symking{}e4+ & \symrook{}xf6\\
\textbf{4.} & \symking{}d3 & \symrook{}a6\\
\textbf{5.} & \symking{}c2 & \symking{}xa2\\
\textbf{6.} & \symking{}d1+ & \symking{}b3\\
\textbf{7.} & \symking{}c1 & \symrook{}a1\#\\
\end{tabular}\par\vspace{3mm}
\qrcode[height=12mm]{https://helpman.komtera.lt/?fen=8/3R4/5b2/8/4k3/8/p2r4/K7&moves=6.5}\\{\sffamily\tiny\color{ink}scannen zum Lösen}
\end{minipage}
\vspace{2mm}

{\raggedright\noindent\textcolor{ink}{\sffamily\textsc{Themen}}\enspace \themetag{reines Matt} \themetag{Mustermatt} \themetag{Idealmatt} \themetag{Switchback} \themetag{ein Stein zieht}\par\smallskip{\footnotesize\textit{Hinweis:} die Partei am Zug steht im Diagramm im Schach.}\par}

\clearpage
\phantomsection\label{hm:KRvkrnp}
\begin{tcolorbox}[enhanced,colback=ink,colframe=ink,arc=0pt,boxrule=0pt,left=4mm,right=4mm,top=2.5mm,bottom=2.5mm]
\color{white}\sffamily{\LARGE\bfseries KRvkrnp}\hfill{\Large h\#6.5}\\[0.5mm]
{\small 6 Steine\quad Weiß am Zug}\hfill{\small Nr.~644\qquad\href{https://helpman.komtera.lt/?fen=8/7k/5R2/7p/8/n7/2r5/1K6&moves=6.5}{\color{white}\bfseries lösen\,\ensuremath{\nearrow}}}
\end{tcolorbox}
\vspace{2mm}
\noindent\begin{minipage}[t]{36mm}\vspace{0pt}\sffamily\small\raggedright
\hmfigure{längste eindeutige}{{\Large h\#6.5}}\hmfigure{eindeutig bei dieser Länge}{{\Large 5} Stellungen}\hmfigure{längstes Matt}{{\Large h\#8}}\hmfigure{in der Klasse}{45,4 Milliarden Stellungen,\\ 27,6 Milliarden mit Hilfsmatt}
\end{minipage}\hfill
\begin{minipage}[t]{80mm}\vspace{0pt}\begin{tikzpicture}\begin{axis}[width=72mm,height=42mm,ybar,bar width=1.6pt,xmin=-1,xmax=17,ymin=0,ymax=5.93,axis y line*=left,axis x line*=bottom,xlabel={\scriptsize Halbzüge bis Matt},ylabel={\scriptsize Stellungen (Mrd.)},tick label style={font=\tiny},label style={font=\scriptsize},enlarge x limits=false,ylabel shift=-3pt,xlabel shift=-3pt]\addplot[fill=ink!55,draw=none] coordinates {(0,0.005135) (1,0.0342) (2,0.1235) (3,0.6152) (4,1.232) (5,2.921) (6,3.96) (7,5.391) (8,5.154) (9,4.602) (10,2.333) (11,0.9867) (12,0.2114) (13,0.01191) (14,0.0006819) (15,2.413e-06) (16,4e-09)};\draw[densely dotted] (axis cs:13,0) -- (axis cs:13,5.93);\end{axis}\begin{axis}[width=72mm,height=42mm,axis y line*=right,axis x line=none,ymode=log,xmin=-1,xmax=17,ylabel={\scriptsize eindeutig (log)},tick label style={font=\tiny},label style={font=\scriptsize},ylabel shift=-3pt]\addplot[mark=*,mark size=0.7pt,draw=black,thin] coordinates {(0,5134712) (1,34194598) (2,58804364) (3,199979923) (4,269155851) (5,249929702) (6,149171639) (7,80305859) (8,39579885) (9,14630511) (10,3467145) (11,170445) (12,6452) (13,5)};\end{axis}\end{tikzpicture}\end{minipage}
\vspace{3mm}

\noindent\begin{minipage}[t]{60mm}\vspace{0pt}\sbox\hmboard{\chessboard[setfen=8/7k/5R2/7p/8/n7/2r5/1K6 w - - 0 1]}\usebox\hmboard\\[1mm]
\makebox[\wd\hmboard]{\sffamily\hspace{3mm}\textbf{h\#6.5}\hfill(2+4)}\end{minipage}\hfill
\begin{minipage}[t]{52mm}\vspace{0pt}
\textcolor{ink}{\sffamily\textsc{Lösung}}\par\smallskip\small
\begin{tabular}{@{}r@{\enspace}l@{\enspace}l@{}}
\textbf{1.} & \ldots & \symking{}a1\\
\textbf{2.} & \symrook{}c7 & \symking{}b2\\
\textbf{3.} & \symknight{}c4+ & \symking{}c3\\
\textbf{4.} & \symknight{}e5+ & \symking{}d4\\
\textbf{5.} & \symknight{}g6 & \symking{}e4\\
\textbf{6.} & \symking{}h6 & \symking{}f5\\
\textbf{7.} & \symrook{}h7 & \symrook{}xg6\#\\
\end{tabular}\par\vspace{3mm}
\qrcode[height=12mm]{https://helpman.komtera.lt/?fen=8/7k/5R2/7p/8/n7/2r5/1K6&moves=6.5}\\{\sffamily\tiny\color{ink}scannen zum Lösen}
\end{minipage}
\vspace{2mm}

{\raggedright\noindent\textcolor{ink}{\sffamily\textsc{Themen}}\enspace \themetag{Selbstblock} \themetag{Inder}\par\smallskip{\footnotesize\textit{Hinweis:} die Partei am Zug steht im Diagramm im Schach.}\par}

\clearpage
\phantomsection\label{hm:KRvkrrb}
\begin{tcolorbox}[enhanced,colback=ink,colframe=ink,arc=0pt,boxrule=0pt,left=4mm,right=4mm,top=2.5mm,bottom=2.5mm]
\color{white}\sffamily{\LARGE\bfseries KRvkrrb}\hfill{\Large h\#6.5}\\[0.5mm]
{\small 6 Steine\quad Weiß am Zug}\hfill{\small Nr.~645\qquad\href{https://helpman.komtera.lt/?fen=3b1R2/1r5k/1r6/8/8/8/8/K7&moves=6.5}{\color{white}\bfseries lösen\,\ensuremath{\nearrow}}}
\end{tcolorbox}
\vspace{2mm}
\noindent\begin{minipage}[t]{36mm}\vspace{0pt}\sffamily\small\raggedright
\hmfigure{längste eindeutige}{{\Large h\#6.5}}\hmfigure{eindeutig bei dieser Länge}{{\Large 84} Stellungen}\hmfigure{längstes Matt}{{\Large h\#8}}\hmfigure{in der Klasse}{15,5 Milliarden Stellungen,\\ 4,08 Milliarden mit Hilfsmatt}
\end{minipage}\hfill
\begin{minipage}[t]{80mm}\vspace{0pt}\begin{tikzpicture}\begin{axis}[width=72mm,height=42mm,ybar,bar width=1.6pt,xmin=-1,xmax=17,ymin=0,ymax=918.1,axis y line*=left,axis x line*=bottom,xlabel={\scriptsize Halbzüge bis Matt},ylabel={\scriptsize Stellungen (Mio.)},tick label style={font=\tiny},label style={font=\scriptsize},enlarge x limits=false,ylabel shift=-3pt,xlabel shift=-3pt]\addplot[fill=ink!55,draw=none] coordinates {(0,0.3094) (1,2.287) (2,10.58) (3,66.44) (4,125.3) (5,388.3) (6,446.2) (7,828.8) (8,645.1) (9,834.6) (10,388.1) (11,275.8) (12,57.9) (13,9.814) (14,0.6674) (15,0.01196) (16,6.4e-05)};\draw[densely dotted] (axis cs:13,0) -- (axis cs:13,918.1);\end{axis}\begin{axis}[width=72mm,height=42mm,axis y line*=right,axis x line=none,ymode=log,xmin=-1,xmax=17,ylabel={\scriptsize eindeutig (log)},tick label style={font=\tiny},label style={font=\scriptsize},ylabel shift=-3pt]\addplot[mark=*,mark size=0.7pt,draw=black,thin] coordinates {(0,309365) (1,2287303) (2,4596833) (3,19989867) (4,20744526) (5,31099063) (6,15141772) (7,14248692) (8,5611441) (9,4032625) (10,879572) (11,155359) (12,5553) (13,84)};\end{axis}\end{tikzpicture}\end{minipage}
\vspace{3mm}

\noindent\begin{minipage}[t]{60mm}\vspace{0pt}\sbox\hmboard{\chessboard[setfen=3b1R2/1r5k/1r6/8/8/8/8/K7 w - - 0 1]}\usebox\hmboard\\[1mm]
\makebox[\wd\hmboard]{\sffamily\hspace{3mm}\textbf{h\#6.5}\hfill(2+4)}\end{minipage}\hfill
\begin{minipage}[t]{52mm}\vspace{0pt}
\textcolor{ink}{\sffamily\textsc{Lösung}}\par\smallskip\small
\begin{tabular}{@{}r@{\enspace}l@{\enspace}l@{}}
\textbf{1.} & \ldots & \symking{}a2\\
\textbf{2.} & \symbishop{}e7 & \symrook{}a8\\
\textbf{3.} & \symbishop{}b4 & \symking{}b3\\
\textbf{4.} & \symrook{}g7 & \symking{}c4\\
\textbf{5.} & \symrook{}b8 & \symking{}d5\\
\textbf{6.} & \symking{}g8 & \symking{}e6\\
\textbf{7.} & \symking{}f8 & \symrook{}xb8\#\\
\end{tabular}\par\vspace{3mm}
\qrcode[height=12mm]{https://helpman.komtera.lt/?fen=3b1R2/1r5k/1r6/8/8/8/8/K7&moves=6.5}\\{\sffamily\tiny\color{ink}scannen zum Lösen}
\end{minipage}
\vspace{2mm}

{\raggedright\noindent\textcolor{ink}{\sffamily\textsc{Themen}}\enspace \themetag{reines Matt} \themetag{Mustermatt} \themetag{Selbstblock} \themetag{ohne Schach} \themetag{Maslar}\par}

\clearpage
\phantomsection\label{hm:KRvkrrn}
\begin{tcolorbox}[enhanced,colback=ink,colframe=ink,arc=0pt,boxrule=0pt,left=4mm,right=4mm,top=2.5mm,bottom=2.5mm]
\color{white}\sffamily{\LARGE\bfseries KRvkrrn}\hfill{\Large h\#6.5}\\[0.5mm]
{\small 6 Steine\quad Weiß am Zug}\hfill{\small Nr.~646\qquad\href{https://helpman.komtera.lt/?fen=8/6R1/7k/8/8/1r6/1r6/K2n4&moves=6.5}{\color{white}\bfseries lösen\,\ensuremath{\nearrow}}}
\end{tcolorbox}
\vspace{2mm}
\noindent\begin{minipage}[t]{36mm}\vspace{0pt}\sffamily\small\raggedright
\hmfigure{längste eindeutige}{{\Large h\#6.5}}\hmfigure{eindeutig bei dieser Länge}{{\Large 9} Stellungen}\hmfigure{längstes Matt}{{\Large h\#8}}\hmfigure{in der Klasse}{15,5 Milliarden Stellungen,\\ 4,16 Milliarden mit Hilfsmatt}
\end{minipage}\hfill
\begin{minipage}[t]{80mm}\vspace{0pt}\begin{tikzpicture}\begin{axis}[width=72mm,height=42mm,ybar,bar width=1.6pt,xmin=-1,xmax=17,ymin=0,ymax=971.7,axis y line*=left,axis x line*=bottom,xlabel={\scriptsize Halbzüge bis Matt},ylabel={\scriptsize Stellungen (Mio.)},tick label style={font=\tiny},label style={font=\scriptsize},enlarge x limits=false,ylabel shift=-3pt,xlabel shift=-3pt]\addplot[fill=ink!55,draw=none] coordinates {(0,0.3793) (1,2.838) (2,12.72) (3,74.9) (4,148.4) (5,418.9) (6,523.1) (7,883.3) (8,723.2) (9,823.6) (10,322.9) (11,200.6) (12,25.35) (13,4.174) (14,0.1203) (15,0.003379) (16,6.1e-05)};\draw[densely dotted] (axis cs:13,0) -- (axis cs:13,971.7);\end{axis}\begin{axis}[width=72mm,height=42mm,axis y line*=right,axis x line=none,ymode=log,xmin=-1,xmax=17,ylabel={\scriptsize eindeutig (log)},tick label style={font=\tiny},label style={font=\scriptsize},ylabel shift=-3pt]\addplot[mark=*,mark size=0.7pt,draw=black,thin] coordinates {(0,379316) (1,2838015) (2,5663288) (3,23415577) (4,28008319) (5,34288055) (6,18150249) (7,15157730) (8,5390019) (9,3241233) (10,486075) (11,74933) (12,1600) (13,9)};\end{axis}\end{tikzpicture}\end{minipage}
\vspace{3mm}

\noindent\begin{minipage}[t]{60mm}\vspace{0pt}\sbox\hmboard{\chessboard[setfen=8/6R1/7k/8/8/1r6/1r6/K2n4 w - - 0 1]}\usebox\hmboard\\[1mm]
\makebox[\wd\hmboard]{\sffamily\hspace{3mm}\textbf{h\#6.5}\hfill(2+4)}\end{minipage}\hfill
\begin{minipage}[t]{52mm}\vspace{0pt}
\textcolor{ink}{\sffamily\textsc{Lösung}}\par\smallskip\small
\begin{tabular}{@{}r@{\enspace}l@{\enspace}l@{}}
\textbf{1.} & \ldots & \symrook{}a7\\
\textbf{2.} & \symrook{}a2+ & \symking{}xa2\\
\textbf{3.} & \symrook{}b7 & \symking{}a3\\
\textbf{4.} & \symrook{}f7 & \symking{}b4\\
\textbf{5.} & \symking{}g7 & \symking{}c5\\
\textbf{6.} & \symking{}f8 & \symking{}d6\\
\textbf{7.} & \symking{}e8 & \symrook{}a8\#\\
\end{tabular}\par\vspace{3mm}
\qrcode[height=12mm]{https://helpman.komtera.lt/?fen=8/6R1/7k/8/8/1r6/1r6/K2n4&moves=6.5}\\{\sffamily\tiny\color{ink}scannen zum Lösen}
\end{minipage}
\vspace{2mm}

{\raggedright\noindent\textcolor{ink}{\sffamily\textsc{Themen}}\enspace \themetag{reines Matt} \themetag{Mustermatt} \themetag{Selbstblock}\par}

\clearpage
\phantomsection\label{hm:KRvkrrr}
\begin{tcolorbox}[enhanced,colback=ink,colframe=ink,arc=0pt,boxrule=0pt,left=4mm,right=4mm,top=2.5mm,bottom=2.5mm]
\color{white}\sffamily{\LARGE\bfseries KRvkrrr}\hfill{\Large h\#6.5}\\[0.5mm]
{\small 6 Steine\quad Weiß am Zug}\hfill{\small Nr.~647\qquad\href{https://helpman.komtera.lt/?fen=5R2/8/1r4k1/1r6/8/8/3r4/K7&moves=6.5}{\color{white}\bfseries lösen\,\ensuremath{\nearrow}}}
\end{tcolorbox}
\vspace{2mm}
\noindent\begin{minipage}[t]{36mm}\vspace{0pt}\sffamily\small\raggedright
\hmfigure{längste eindeutige}{{\Large h\#6.5}}\hmfigure{eindeutig bei dieser Länge}{{\Large 152} Stellungen}\hmfigure{längstes Matt}{{\Large h\#8}}\hmfigure{in der Klasse}{15,5 Milliarden Stellungen,\\ 1,28 Milliarden mit Hilfsmatt}
\end{minipage}\hfill
\begin{minipage}[t]{80mm}\vspace{0pt}\begin{tikzpicture}\begin{axis}[width=72mm,height=42mm,ybar,bar width=1.6pt,xmin=-1,xmax=17,ymin=0,ymax=305.7,axis y line*=left,axis x line*=bottom,xlabel={\scriptsize Halbzüge bis Matt},ylabel={\scriptsize Stellungen (Mio.)},tick label style={font=\tiny},label style={font=\scriptsize},enlarge x limits=false,ylabel shift=-3pt,xlabel shift=-3pt]\addplot[fill=ink!55,draw=none] coordinates {(0,0.07386) (1,0.5859) (2,2.954) (3,19.25) (4,35.63) (5,115.4) (6,127.1) (7,259.5) (8,185) (9,277.9) (10,118.7) (11,110.6) (12,20.59) (13,6.616) (14,0.3723) (15,0.0195) (16,0.000208)};\draw[densely dotted] (axis cs:13,0) -- (axis cs:13,305.7);\end{axis}\begin{axis}[width=72mm,height=42mm,axis y line*=right,axis x line=none,ymode=log,xmin=-1,xmax=17,ylabel={\scriptsize eindeutig (log)},tick label style={font=\tiny},label style={font=\scriptsize},ylabel shift=-3pt]\addplot[mark=*,mark size=0.7pt,draw=black,thin] coordinates {(0,73855) (1,585914) (2,1241947) (3,5607111) (4,5344073) (5,8784291) (6,3941550) (7,4372962) (8,1434672) (9,1362981) (10,282754) (11,100094) (12,4511) (13,152)};\end{axis}\end{tikzpicture}\end{minipage}
\vspace{3mm}

\noindent\begin{minipage}[t]{60mm}\vspace{0pt}\sbox\hmboard{\chessboard[setfen=5R2/8/1r4k1/1r6/8/8/3r4/K7 w - - 0 1]}\usebox\hmboard\\[1mm]
\makebox[\wd\hmboard]{\sffamily\hspace{3mm}\textbf{h\#6.5}\hfill(2+4)}\end{minipage}\hfill
\begin{minipage}[t]{52mm}\vspace{0pt}
\textcolor{ink}{\sffamily\textsc{Lösung}}\par\smallskip\small
\begin{tabular}{@{}r@{\enspace}l@{\enspace}l@{}}
\textbf{1.} & \ldots & \symrook{}b8\\
\textbf{2.} & \symking{}f5 & \symrook{}xb6\\
\textbf{3.} & \symking{}e4 & \symrook{}xb5\\
\textbf{4.} & \symking{}d3 & \symrook{}a5\\
\textbf{5.} & \symking{}c2 & \symking{}a2\\
\textbf{6.} & \symking{}d1+ & \symking{}b3\\
\textbf{7.} & \symking{}c1 & \symrook{}a1\#\\
\end{tabular}\par\vspace{3mm}
\qrcode[height=12mm]{https://helpman.komtera.lt/?fen=5R2/8/1r4k1/1r6/8/8/3r4/K7&moves=6.5}\\{\sffamily\tiny\color{ink}scannen zum Lösen}
\end{minipage}
\vspace{2mm}

{\raggedright\noindent\textcolor{ink}{\sffamily\textsc{Themen}}\enspace \themetag{reines Matt} \themetag{Mustermatt} \themetag{Idealmatt} \themetag{ein Stein zieht}\par}

\clearpage
\phantomsection\label{hm:KBBBvkr}
\begin{tcolorbox}[enhanced,colback=ink,colframe=ink,arc=0pt,boxrule=0pt,left=4mm,right=4mm,top=2.5mm,bottom=2.5mm]
\color{white}\sffamily{\LARGE\bfseries KBBBvkr}\hfill{\Large h\#6}\\[0.5mm]
{\small 6 Steine\quad Schwarz am Zug}\hfill{\small Nr.~648\qquad\href{https://helpman.komtera.lt/?fen=6B1/6Bk/8/8/1r6/8/8/KB6&moves=6}{\color{white}\bfseries lösen\,\ensuremath{\nearrow}}}
\end{tcolorbox}
\vspace{2mm}
\noindent\begin{minipage}[t]{36mm}\vspace{0pt}\sffamily\small\raggedright
\hmfigure{längste eindeutige}{{\Large h\#6}}\hmfigure{eindeutig bei dieser Länge}{{\Large 28} Stellungen}\hmfigure{längstes Matt}{{\Large h\#8}}\hmfigure{in der Klasse}{15,5 Milliarden Stellungen,\\ 1,1 Milliarden mit Hilfsmatt}
\end{minipage}\hfill
\begin{minipage}[t]{80mm}\vspace{0pt}\begin{tikzpicture}\begin{axis}[width=72mm,height=42mm,ybar,bar width=1.6pt,xmin=-1,xmax=17,ymin=0,ymax=449.7,axis y line*=left,axis x line*=bottom,xlabel={\scriptsize Halbzüge bis Matt},ylabel={\scriptsize Stellungen (Mio.)},tick label style={font=\tiny},label style={font=\scriptsize},enlarge x limits=false,ylabel shift=-3pt,xlabel shift=-3pt]\addplot[fill=ink!55,draw=none] coordinates {(0,0.4532) (1,1.717) (2,8.089) (3,42.22) (4,153.1) (5,269.6) (6,408.8) (7,166.8) (8,44.36) (9,3.167) (10,0.1666) (11,0.002655) (12,0.03631) (13,0.000861) (14,0.003718) (15,1e-05) (16,3e-06)};\draw[densely dotted] (axis cs:12,0) -- (axis cs:12,449.7);\end{axis}\begin{axis}[width=72mm,height=42mm,axis y line*=right,axis x line=none,ymode=log,xmin=-1,xmax=17,ylabel={\scriptsize eindeutig (log)},tick label style={font=\tiny},label style={font=\scriptsize},ylabel shift=-3pt]\addplot[mark=*,mark size=0.7pt,draw=black,thin] coordinates {(0,453227) (1,1566478) (2,4939867) (3,15185185) (4,34934440) (5,14356990) (6,5445380) (7,202489) (8,4748) (9,5) (10,419) (12,28)};\end{axis}\end{tikzpicture}\end{minipage}
\vspace{3mm}

\noindent\begin{minipage}[t]{60mm}\vspace{0pt}\sbox\hmboard{\chessboard[setfen=6B1/6Bk/8/8/1r6/8/8/KB6 b - - 0 1]}\usebox\hmboard\\[1mm]
\makebox[\wd\hmboard]{\sffamily\hspace{3mm}\textbf{h\#6}\hfill(4+2)}\end{minipage}\hfill
\begin{minipage}[t]{52mm}\vspace{0pt}
\textcolor{ink}{\sffamily\textsc{Lösung}}\par\smallskip\small
\begin{tabular}{@{}r@{\enspace}l@{\enspace}l@{}}
\textbf{1.} & \symking{}xg8 & \symking{}a2\\
\textbf{2.} & \symrook{}a4+ & \symking{}b3\\
\textbf{3.} & \symrook{}a8 & \symking{}c4\\
\textbf{4.} & \symking{}f7 & \symking{}d5\\
\textbf{5.} & \symking{}e8 & \symking{}e6\\
\textbf{6.} & \symrook{}d8 & \symbishop{}g6\#\\
\end{tabular}\par\vspace{3mm}
\qrcode[height=12mm]{https://helpman.komtera.lt/?fen=6B1/6Bk/8/8/1r6/8/8/KB6&moves=6}\\{\sffamily\tiny\color{ink}scannen zum Lösen}
\end{minipage}
\vspace{2mm}

{\raggedright\noindent\textcolor{ink}{\sffamily\textsc{Themen}}\enspace \themetag{Selbstblock}\par\smallskip{\footnotesize\textit{Hinweis:} das Diagramm hat keinen legalen letzten Zug, kann also in keiner Partie entstehen.}\par}

\clearpage
\phantomsection\label{hm:KBBNvkr}
\begin{tcolorbox}[enhanced,colback=ink,colframe=ink,arc=0pt,boxrule=0pt,left=4mm,right=4mm,top=2.5mm,bottom=2.5mm]
\color{white}\sffamily{\LARGE\bfseries KBBNvkr}\hfill{\Large h\#6}\\[0.5mm]
{\small 6 Steine\quad Schwarz am Zug}\hfill{\small Nr.~649\qquad\href{https://helpman.komtera.lt/?fen=8/6N1/7k/8/4B3/8/8/K1B4r&moves=6}{\color{white}\bfseries lösen\,\ensuremath{\nearrow}}}
\end{tcolorbox}
\vspace{2mm}
\noindent\begin{minipage}[t]{36mm}\vspace{0pt}\sffamily\small\raggedright
\hmfigure{längste eindeutige}{{\Large h\#6}}\hmfigure{eindeutig bei dieser Länge}{{\Large 92} Stellungen}\hmfigure{längstes Matt}{{\Large h\#8.5}}\hmfigure{in der Klasse}{15,5 Milliarden Stellungen,\\ 4,53 Milliarden mit Hilfsmatt}
\end{minipage}\hfill
\begin{minipage}[t]{80mm}\vspace{0pt}\begin{tikzpicture}\begin{axis}[width=72mm,height=42mm,ybar,bar width=1.6pt,xmin=-1,xmax=18,ymin=0,ymax=1.297,axis y line*=left,axis x line*=bottom,xlabel={\scriptsize Halbzüge bis Matt},ylabel={\scriptsize Stellungen (Mrd.)},tick label style={font=\tiny},label style={font=\scriptsize},enlarge x limits=false,ylabel shift=-3pt,xlabel shift=-3pt]\addplot[fill=ink!55,draw=none] coordinates {(0,0.0005427) (1,0.002153) (2,0.009781) (3,0.06124) (4,0.2432) (5,0.6317) (6,1.179) (7,0.9717) (8,0.7982) (9,0.3891) (10,0.189) (11,0.04415) (12,0.005288) (13,0.0002848) (14,2.823e-05) (15,2.112e-06) (16,1.81e-07) (17,1e-09)};\draw[densely dotted] (axis cs:12,0) -- (axis cs:12,1.297);\end{axis}\begin{axis}[width=72mm,height=42mm,axis y line*=right,axis x line=none,ymode=log,xmin=-1,xmax=18,ylabel={\scriptsize eindeutig (log)},tick label style={font=\tiny},label style={font=\scriptsize},ylabel shift=-3pt]\addplot[mark=*,mark size=0.7pt,draw=black,thin] coordinates {(0,542701) (1,2006077) (2,6236892) (3,25966028) (4,69844736) (5,52790151) (6,36606598) (7,8257021) (8,3002923) (9,592281) (10,58237) (11,3863) (12,92)};\end{axis}\end{tikzpicture}\end{minipage}
\vspace{3mm}

\noindent\begin{minipage}[t]{60mm}\vspace{0pt}\sbox\hmboard{\chessboard[setfen=8/6N1/7k/8/4B3/8/8/K1B4r b - - 0 1]}\usebox\hmboard\\[1mm]
\makebox[\wd\hmboard]{\sffamily\hspace{3mm}\textbf{h\#6}\hfill(4+2)}\end{minipage}\hfill
\begin{minipage}[t]{52mm}\vspace{0pt}
\textcolor{ink}{\sffamily\textsc{Lösung}}\par\smallskip\small
\begin{tabular}{@{}r@{\enspace}l@{\enspace}l@{}}
\textbf{1.} & \symrook{}xc1+ & \symking{}b2\\
\textbf{2.} & \symking{}g5 & \symknight{}f5\\
\textbf{3.} & \symking{}f4 & \symknight{}d4\\
\textbf{4.} & \symking{}e3 & \symknight{}c2+\\
\textbf{5.} & \symking{}e2 & \symking{}c3\\
\textbf{6.} & \symking{}d1 & \symbishop{}f3\#\\
\end{tabular}\par\vspace{3mm}
\qrcode[height=12mm]{https://helpman.komtera.lt/?fen=8/6N1/7k/8/4B3/8/8/K1B4r&moves=6}\\{\sffamily\tiny\color{ink}scannen zum Lösen}
\end{minipage}
\vspace{2mm}

{\raggedright\noindent\textcolor{ink}{\sffamily\textsc{Themen}}\enspace \themetag{reines Matt} \themetag{Mustermatt} \themetag{Idealmatt} \themetag{Selbstblock}\par\smallskip{\footnotesize\textit{Hinweis:} die Partei am Zug steht im Diagramm im Schach; die Lösung beginnt mit einem Schlag.}\par}

\clearpage
\phantomsection\label{hm:KBNNNvk}
\begin{tcolorbox}[enhanced,colback=ink,colframe=ink,arc=0pt,boxrule=0pt,left=4mm,right=4mm,top=2.5mm,bottom=2.5mm]
\color{white}\sffamily{\LARGE\bfseries KBNNNvk}\hfill{\Large h\#6}\\[0.5mm]
{\small 6 Steine\quad Schwarz am Zug}\hfill{\small Nr.~650\qquad\href{https://helpman.komtera.lt/?fen=8/6NN/7k/8/5N2/8/8/K2B4&moves=6}{\color{white}\bfseries lösen\,\ensuremath{\nearrow}}}
\end{tcolorbox}
\vspace{2mm}
\noindent\begin{minipage}[t]{36mm}\vspace{0pt}\sffamily\small\raggedright
\hmfigure{längste eindeutige}{{\Large h\#6}}\hmfigure{eindeutig bei dieser Länge}{{\Large 40} Stellungen}\hmfigure{längstes Matt}{{\Large h\#8}}\hmfigure{in der Klasse}{15,5 Milliarden Stellungen,\\ 1,69 Milliarden mit Hilfsmatt}
\end{minipage}\hfill
\begin{minipage}[t]{80mm}\vspace{0pt}\begin{tikzpicture}\begin{axis}[width=72mm,height=42mm,ybar,bar width=1.6pt,xmin=-1,xmax=17,ymin=0,ymax=585.1,axis y line*=left,axis x line*=bottom,xlabel={\scriptsize Halbzüge bis Matt},ylabel={\scriptsize Stellungen (Mio.)},tick label style={font=\tiny},label style={font=\scriptsize},enlarge x limits=false,ylabel shift=-3pt,xlabel shift=-3pt]\addplot[fill=ink!55,draw=none] coordinates {(0,0.4638) (1,1.673) (2,5.057) (3,21.11) (4,91.79) (5,253.9) (6,531.9) (7,347.5) (8,358.3) (9,57.04) (10,16.81) (11,0.3124) (12,0.1473) (13,0.000141) (14,0.02082) (15,4e-06) (16,4.9e-05)};\draw[densely dotted] (axis cs:12,0) -- (axis cs:12,585.1);\end{axis}\begin{axis}[width=72mm,height=42mm,axis y line*=right,axis x line=none,ymode=log,xmin=-1,xmax=17,ylabel={\scriptsize eindeutig (log)},tick label style={font=\tiny},label style={font=\scriptsize},ylabel shift=-3pt]\addplot[mark=*,mark size=0.7pt,draw=black,thin] coordinates {(0,463783) (1,1436494) (2,4314292) (3,13102909) (4,41212173) (5,28305933) (6,27492969) (7,1418133) (8,484509) (9,2597) (10,1205) (12,40)};\end{axis}\end{tikzpicture}\end{minipage}
\vspace{3mm}

\noindent\begin{minipage}[t]{60mm}\vspace{0pt}\sbox\hmboard{\chessboard[setfen=8/6NN/7k/8/5N2/8/8/K2B4 b - - 0 1]}\usebox\hmboard\\[1mm]
\makebox[\wd\hmboard]{\sffamily\hspace{3mm}\textbf{h\#6}\hfill(5+1)}\end{minipage}\hfill
\begin{minipage}[t]{52mm}\vspace{0pt}
\textcolor{ink}{\sffamily\textsc{Lösung}}\par\smallskip\small
\begin{tabular}{@{}r@{\enspace}l@{\enspace}l@{}}
\textbf{1.} & \symking{}xh7 & \symking{}b2\\
\textbf{2.} & \symking{}h6 & \symking{}c3\\
\textbf{3.} & \symking{}g5 & \symking{}d4\\
\textbf{4.} & \symking{}h4 & \symking{}e5\\
\textbf{5.} & \symking{}g3 & \symking{}f6\\
\textbf{6.} & \symking{}h4 & \symknight{}f5\#\\
\end{tabular}\par\vspace{3mm}
\qrcode[height=12mm]{https://helpman.komtera.lt/?fen=8/6NN/7k/8/5N2/8/8/K2B4&moves=6}\\{\sffamily\tiny\color{ink}scannen zum Lösen}
\end{minipage}
\vspace{2mm}

{\raggedright\noindent\textcolor{ink}{\sffamily\textsc{Themen}}\enspace \themetag{Spiegelmatt} \themetag{Switchback} \themetag{ein Stein zieht} \themetag{ohne Schach}\par\smallskip{\footnotesize\textit{Hinweis:} die Lösung beginnt mit einem Schlag.}\par}

\clearpage
\phantomsection\label{hm:KBNNvkb}
\begin{tcolorbox}[enhanced,colback=ink,colframe=ink,arc=0pt,boxrule=0pt,left=4mm,right=4mm,top=2.5mm,bottom=2.5mm]
\color{white}\sffamily{\LARGE\bfseries KBNNvkb}\hfill{\Large h\#6}\\[0.5mm]
{\small 6 Steine\quad Schwarz am Zug}\hfill{\small Nr.~651\qquad\href{https://helpman.komtera.lt/?fen=8/5NN1/7k/8/4B3/8/8/K4b2&moves=6}{\color{white}\bfseries lösen\,\ensuremath{\nearrow}}}
\end{tcolorbox}
\vspace{2mm}
\noindent\begin{minipage}[t]{36mm}\vspace{0pt}\sffamily\small\raggedright
\hmfigure{längste eindeutige}{{\Large h\#6}}\hmfigure{eindeutig bei dieser Länge}{{\Large 25} Stellungen}\hmfigure{längstes Matt}{{\Large h\#8}}\hmfigure{in der Klasse}{15,5 Milliarden Stellungen,\\ 4,86 Milliarden mit Hilfsmatt}
\end{minipage}\hfill
\begin{minipage}[t]{80mm}\vspace{0pt}\begin{tikzpicture}\begin{axis}[width=72mm,height=42mm,ybar,bar width=1.6pt,xmin=-1,xmax=17,ymin=0,ymax=1.38,axis y line*=left,axis x line*=bottom,xlabel={\scriptsize Halbzüge bis Matt},ylabel={\scriptsize Stellungen (Mrd.)},tick label style={font=\tiny},label style={font=\scriptsize},enlarge x limits=false,ylabel shift=-3pt,xlabel shift=-3pt]\addplot[fill=ink!55,draw=none] coordinates {(0,0.0003884) (1,0.001502) (2,0.007) (3,0.0437) (4,0.193) (5,0.5811) (6,1.255) (7,1.222) (8,1.061) (9,0.3624) (10,0.118) (11,0.01502) (12,0.0009851) (13,1.469e-05) (14,4.709e-06) (15,1.24e-07) (16,2.7e-08)};\draw[densely dotted] (axis cs:12,0) -- (axis cs:12,1.38);\end{axis}\begin{axis}[width=72mm,height=42mm,axis y line*=right,axis x line=none,ymode=log,xmin=-1,xmax=17,ylabel={\scriptsize eindeutig (log)},tick label style={font=\tiny},label style={font=\scriptsize},ylabel shift=-3pt]\addplot[mark=*,mark size=0.7pt,draw=black,thin] coordinates {(0,388409) (1,1383796) (2,5370127) (3,20681053) (4,62216860) (5,53870349) (6,42696048) (7,6774328) (8,1497046) (9,72263) (10,4236) (11,53) (12,25)};\end{axis}\end{tikzpicture}\end{minipage}
\vspace{3mm}

\noindent\begin{minipage}[t]{60mm}\vspace{0pt}\sbox\hmboard{\chessboard[setfen=8/5NN1/7k/8/4B3/8/8/K4b2 b - - 0 1]}\usebox\hmboard\\[1mm]
\makebox[\wd\hmboard]{\sffamily\hspace{3mm}\textbf{h\#6}\hfill(4+2)}\end{minipage}\hfill
\begin{minipage}[t]{52mm}\vspace{0pt}
\textcolor{ink}{\sffamily\textsc{Lösung}}\par\smallskip\small
\begin{tabular}{@{}r@{\enspace}l@{\enspace}l@{}}
\textbf{1.} & \symking{}xg7 & \symking{}b2\\
\textbf{2.} & \symbishop{}h3 & \symking{}c3\\
\textbf{3.} & \symbishop{}f5 & \symking{}d4\\
\textbf{4.} & \symbishop{}h7 & \symking{}e5\\
\textbf{5.} & \symbishop{}g8 & \symking{}f5\\
\textbf{6.} & \symking{}h7 & \symking{}f6\#\\
\end{tabular}\par\vspace{3mm}
\qrcode[height=12mm]{https://helpman.komtera.lt/?fen=8/5NN1/7k/8/4B3/8/8/K4b2&moves=6}\\{\sffamily\tiny\color{ink}scannen zum Lösen}
\end{minipage}
\vspace{2mm}

{\raggedright\noindent\textcolor{ink}{\sffamily\textsc{Themen}}\enspace \themetag{Selbstblock} \themetag{ein Stein zieht} \themetag{ohne Schach}\par\smallskip{\footnotesize\textit{Hinweis:} die Partei am Zug steht im Diagramm im Schach; die Lösung beginnt mit einem Schlag.}\par}

\clearpage
\phantomsection\label{hm:KBNNvkn}
\begin{tcolorbox}[enhanced,colback=ink,colframe=ink,arc=0pt,boxrule=0pt,left=4mm,right=4mm,top=2.5mm,bottom=2.5mm]
\color{white}\sffamily{\LARGE\bfseries KBNNvkn}\hfill{\Large h\#6}\\[0.5mm]
{\small 6 Steine\quad Schwarz am Zug}\hfill{\small Nr.~652\qquad\href{https://helpman.komtera.lt/?fen=8/8/8/n4N2/6k1/7N/8/K2B4&moves=6}{\color{white}\bfseries lösen\,\ensuremath{\nearrow}}}
\end{tcolorbox}
\vspace{2mm}
\noindent\begin{minipage}[t]{36mm}\vspace{0pt}\sffamily\small\raggedright
\hmfigure{längste eindeutige}{{\Large h\#6}}\hmfigure{eindeutig bei dieser Länge}{{\Large 63} Stellungen}\hmfigure{längstes Matt}{{\Large h\#8}}\hmfigure{in der Klasse}{15,5 Milliarden Stellungen,\\ 5 Milliarden mit Hilfsmatt}
\end{minipage}\hfill
\begin{minipage}[t]{80mm}\vspace{0pt}\begin{tikzpicture}\begin{axis}[width=72mm,height=42mm,ybar,bar width=1.6pt,xmin=-1,xmax=17,ymin=0,ymax=1.427,axis y line*=left,axis x line*=bottom,xlabel={\scriptsize Halbzüge bis Matt},ylabel={\scriptsize Stellungen (Mrd.)},tick label style={font=\tiny},label style={font=\scriptsize},enlarge x limits=false,ylabel shift=-3pt,xlabel shift=-3pt]\addplot[fill=ink!55,draw=none] coordinates {(0,0.0003539) (1,0.001404) (2,0.005473) (3,0.03656) (4,0.1457) (5,0.4779) (6,1.17) (7,1.232) (8,1.298) (9,0.456) (10,0.1526) (11,0.02148) (12,0.001666) (13,6.816e-05) (14,4.19e-06) (15,1e-09) (16,3e-09)};\draw[densely dotted] (axis cs:12,0) -- (axis cs:12,1.427);\end{axis}\begin{axis}[width=72mm,height=42mm,axis y line*=right,axis x line=none,ymode=log,xmin=-1,xmax=17,ylabel={\scriptsize eindeutig (log)},tick label style={font=\tiny},label style={font=\scriptsize},ylabel shift=-3pt]\addplot[mark=*,mark size=0.7pt,draw=black,thin] coordinates {(0,353875) (1,1283856) (2,3872331) (3,16974345) (4,46779353) (5,40441846) (6,39986688) (7,6270439) (8,1572306) (9,79245) (10,5320) (11,165) (12,63)};\end{axis}\end{tikzpicture}\end{minipage}
\vspace{3mm}

\noindent\begin{minipage}[t]{60mm}\vspace{0pt}\sbox\hmboard{\chessboard[setfen=8/8/8/n4N2/6k1/7N/8/K2B4 b - - 0 1]}\usebox\hmboard\\[1mm]
\makebox[\wd\hmboard]{\sffamily\hspace{3mm}\textbf{h\#6}\hfill(4+2)}\end{minipage}\hfill
\begin{minipage}[t]{52mm}\vspace{0pt}
\textcolor{ink}{\sffamily\textsc{Lösung}}\par\smallskip\small
\begin{tabular}{@{}r@{\enspace}l@{\enspace}l@{}}
\textbf{1.} & \symking{}xh3 & \symking{}b2\\
\textbf{2.} & \symknight{}b3 & \symking{}c3\\
\textbf{3.} & \symknight{}c1 & \symking{}d4\\
\textbf{4.} & \symknight{}e2+ & \symking{}e5\\
\textbf{5.} & \symking{}g4 & \symking{}f6\\
\textbf{6.} & \symking{}h5 & \symbishop{}xe2\#\\
\end{tabular}\par\vspace{3mm}
\qrcode[height=12mm]{https://helpman.komtera.lt/?fen=8/8/8/n4N2/6k1/7N/8/K2B4&moves=6}\\{\sffamily\tiny\color{ink}scannen zum Lösen}
\end{minipage}
\vspace{2mm}

{\raggedright\noindent\textcolor{ink}{\sffamily\textsc{Themen}}\enspace \themetag{reines Matt} \themetag{Mustermatt} \themetag{Idealmatt} \themetag{Spiegelmatt} \themetag{Switchback}\par\smallskip{\footnotesize\textit{Hinweis:} die Partei am Zug steht im Diagramm im Schach; die Lösung beginnt mit einem Schlag.}\par}

\clearpage
\phantomsection\label{hm:KBNNvkr}
\begin{tcolorbox}[enhanced,colback=ink,colframe=ink,arc=0pt,boxrule=0pt,left=4mm,right=4mm,top=2.5mm,bottom=2.5mm]
\color{white}\sffamily{\LARGE\bfseries KBNNvkr}\hfill{\Large h\#6}\\[0.5mm]
{\small 6 Steine\quad Schwarz am Zug}\hfill{\small Nr.~653\qquad\href{https://helpman.komtera.lt/?fen=8/6N1/6N1/7k/6r1/8/8/K1B5&moves=6}{\color{white}\bfseries lösen\,\ensuremath{\nearrow}}}
\end{tcolorbox}
\vspace{2mm}
\noindent\begin{minipage}[t]{36mm}\vspace{0pt}\sffamily\small\raggedright
\hmfigure{längste eindeutige}{{\Large h\#6}}\hmfigure{eindeutig bei dieser Länge}{{\Large 43} Stellungen}\hmfigure{längstes Matt}{{\Large h\#8}}\hmfigure{in der Klasse}{15,5 Milliarden Stellungen,\\ 4,65 Milliarden mit Hilfsmatt}
\end{minipage}\hfill
\begin{minipage}[t]{80mm}\vspace{0pt}\begin{tikzpicture}\begin{axis}[width=72mm,height=42mm,ybar,bar width=1.6pt,xmin=-1,xmax=17,ymin=0,ymax=1.416,axis y line*=left,axis x line*=bottom,xlabel={\scriptsize Halbzüge bis Matt},ylabel={\scriptsize Stellungen (Mrd.)},tick label style={font=\tiny},label style={font=\scriptsize},enlarge x limits=false,ylabel shift=-3pt,xlabel shift=-3pt]\addplot[fill=ink!55,draw=none] coordinates {(0,0.0002497) (1,0.001006) (2,0.004342) (3,0.03087) (4,0.1407) (5,0.4995) (6,1.157) (7,1.287) (8,1.03) (9,0.3928) (10,0.09182) (11,0.01395) (12,0.0004629) (13,3.619e-05) (14,7.17e-06) (15,2.933e-06) (16,2.1e-08)};\draw[densely dotted] (axis cs:12,0) -- (axis cs:12,1.416);\end{axis}\begin{axis}[width=72mm,height=42mm,axis y line*=right,axis x line=none,ymode=log,xmin=-1,xmax=17,ylabel={\scriptsize eindeutig (log)},tick label style={font=\tiny},label style={font=\scriptsize},ylabel shift=-3pt]\addplot[mark=*,mark size=0.7pt,draw=black,thin] coordinates {(0,249661) (1,917605) (2,2932361) (3,13391027) (4,47693661) (5,54210900) (6,46632021) (7,9140751) (8,1168978) (9,97568) (10,2590) (11,183) (12,43)};\end{axis}\end{tikzpicture}\end{minipage}
\vspace{3mm}

\noindent\begin{minipage}[t]{60mm}\vspace{0pt}\sbox\hmboard{\chessboard[setfen=8/6N1/6N1/7k/6r1/8/8/K1B5 b - - 0 1]}\usebox\hmboard\\[1mm]
\makebox[\wd\hmboard]{\sffamily\hspace{3mm}\textbf{h\#6}\hfill(4+2)}\end{minipage}\hfill
\begin{minipage}[t]{52mm}\vspace{0pt}
\textcolor{ink}{\sffamily\textsc{Lösung}}\par\smallskip\small
\begin{tabular}{@{}r@{\enspace}l@{\enspace}l@{}}
\textbf{1.} & \symking{}xg6 & \symking{}b2\\
\textbf{2.} & \symrook{}h4 & \symking{}c3\\
\textbf{3.} & \symrook{}h8 & \symking{}d4\\
\textbf{4.} & \symking{}f7 & \symking{}e5\\
\textbf{5.} & \symking{}f8 & \symking{}f6\\
\textbf{6.} & \symrook{}g8 & \symbishop{}a3\#\\
\end{tabular}\par\vspace{3mm}
\qrcode[height=12mm]{https://helpman.komtera.lt/?fen=8/6N1/6N1/7k/6r1/8/8/K1B5&moves=6}\\{\sffamily\tiny\color{ink}scannen zum Lösen}
\end{minipage}
\vspace{2mm}

{\raggedright\noindent\textcolor{ink}{\sffamily\textsc{Themen}}\enspace \themetag{Selbstblock} \themetag{ohne Schach}\par\smallskip{\footnotesize\textit{Hinweis:} die Partei am Zug steht im Diagramm im Schach; die Lösung beginnt mit einem Schlag.}\par}

\clearpage
\phantomsection\label{hm:KNNNvkb}
\begin{tcolorbox}[enhanced,colback=ink,colframe=ink,arc=0pt,boxrule=0pt,left=4mm,right=4mm,top=2.5mm,bottom=2.5mm]
\color{white}\sffamily{\LARGE\bfseries KNNNvkb}\hfill{\Large h\#6}\\[0.5mm]
{\small 6 Steine\quad Schwarz am Zug}\hfill{\small Nr.~654\qquad\href{https://helpman.komtera.lt/?fen=8/7N/7k/5N2/5N2/8/8/Kb6&moves=6}{\color{white}\bfseries lösen\,\ensuremath{\nearrow}}}
\end{tcolorbox}
\vspace{2mm}
\noindent\begin{minipage}[t]{36mm}\vspace{0pt}\sffamily\small\raggedright
\hmfigure{längste eindeutige}{{\Large h\#6}}\hmfigure{eindeutig bei dieser Länge}{{\Large 11} Stellungen}\hmfigure{längstes Matt}{{\Large h\#8}}\hmfigure{in der Klasse}{15,5 Milliarden Stellungen,\\ 1,66 Milliarden mit Hilfsmatt}
\end{minipage}\hfill
\begin{minipage}[t]{80mm}\vspace{0pt}\begin{tikzpicture}\begin{axis}[width=72mm,height=42mm,ybar,bar width=1.6pt,xmin=-1,xmax=17,ymin=0,ymax=466.2,axis y line*=left,axis x line*=bottom,xlabel={\scriptsize Halbzüge bis Matt},ylabel={\scriptsize Stellungen (Mio.)},tick label style={font=\tiny},label style={font=\scriptsize},enlarge x limits=false,ylabel shift=-3pt,xlabel shift=-3pt]\addplot[fill=ink!55,draw=none] coordinates {(0,0.08245) (1,0.368) (2,1.513) (3,10.62) (4,44.65) (5,152.2) (6,355.1) (7,423.8) (8,414.5) (9,179.8) (10,61.97) (11,10.9) (12,0.7399) (13,0.03093) (14,0.000545) (15,5.7e-05) (16,3.2e-05)};\draw[densely dotted] (axis cs:12,0) -- (axis cs:12,466.2);\end{axis}\begin{axis}[width=72mm,height=42mm,axis y line*=right,axis x line=none,ymode=log,xmin=-1,xmax=17,ylabel={\scriptsize eindeutig (log)},tick label style={font=\tiny},label style={font=\scriptsize},ylabel shift=-3pt]\addplot[mark=*,mark size=0.7pt,draw=black,thin] coordinates {(0,82449) (1,330835) (2,1058316) (3,4650683) (4,14835087) (5,13300384) (6,12846251) (7,2553845) (8,771521) (9,49033) (10,2336) (11,67) (12,11)};\end{axis}\end{tikzpicture}\end{minipage}
\vspace{3mm}

\noindent\begin{minipage}[t]{60mm}\vspace{0pt}\sbox\hmboard{\chessboard[setfen=8/7N/7k/5N2/5N2/8/8/Kb6 b - - 0 1]}\usebox\hmboard\\[1mm]
\makebox[\wd\hmboard]{\sffamily\hspace{3mm}\textbf{h\#6}\hfill(4+2)}\end{minipage}\hfill
\begin{minipage}[t]{52mm}\vspace{0pt}
\textcolor{ink}{\sffamily\textsc{Lösung}}\par\smallskip\small
\begin{tabular}{@{}r@{\enspace}l@{\enspace}l@{}}
\textbf{1.} & \symbishop{}xf5 & \symknight{}d5\\
\textbf{2.} & \symking{}h5 & \symking{}a2\\
\textbf{3.} & \symking{}g4 & \symking{}b2\\
\textbf{4.} & \symking{}f3 & \symknight{}g5+\\
\textbf{5.} & \symking{}e2 & \symknight{}f3\\
\textbf{6.} & \symking{}d1 & \symknight{}c3\#\\
\end{tabular}\par\vspace{3mm}
\qrcode[height=12mm]{https://helpman.komtera.lt/?fen=8/7N/7k/5N2/5N2/8/8/Kb6&moves=6}\\{\sffamily\tiny\color{ink}scannen zum Lösen}
\end{minipage}
\vspace{2mm}

{\raggedright\noindent\textcolor{ink}{\sffamily\textsc{Themen}}\enspace \themetag{reines Matt} \themetag{Mustermatt} \themetag{Spiegelmatt} \themetag{Umnov}\par\smallskip{\footnotesize\textit{Hinweis:} die Partei am Zug steht im Diagramm im Schach; die Lösung beginnt mit einem Schlag.}\par}

\clearpage
\phantomsection\label{hm:KNNNvkq}
\begin{tcolorbox}[enhanced,colback=ink,colframe=ink,arc=0pt,boxrule=0pt,left=4mm,right=4mm,top=2.5mm,bottom=2.5mm]
\color{white}\sffamily{\LARGE\bfseries KNNNvkq}\hfill{\Large h\#6}\\[0.5mm]
{\small 6 Steine\quad Schwarz am Zug}\hfill{\small Nr.~655\qquad\href{https://helpman.komtera.lt/?fen=N5N1/8/8/7N/8/8/4q3/K6k&moves=6}{\color{white}\bfseries lösen\,\ensuremath{\nearrow}}}
\end{tcolorbox}
\vspace{2mm}
\noindent\begin{minipage}[t]{36mm}\vspace{0pt}\sffamily\small\raggedright
\hmfigure{längste eindeutige}{{\Large h\#6}}\hmfigure{eindeutig bei dieser Länge}{{\Large 76} Stellungen}\hmfigure{längstes Matt}{{\Large h\#8.5}}\hmfigure{in der Klasse}{15,5 Milliarden Stellungen,\\ 1,46 Milliarden mit Hilfsmatt}
\end{minipage}\hfill
\begin{minipage}[t]{80mm}\vspace{0pt}\begin{tikzpicture}\begin{axis}[width=72mm,height=42mm,ybar,bar width=1.6pt,xmin=-1,xmax=18,ymin=0,ymax=413.7,axis y line*=left,axis x line*=bottom,xlabel={\scriptsize Halbzüge bis Matt},ylabel={\scriptsize Stellungen (Mio.)},tick label style={font=\tiny},label style={font=\scriptsize},enlarge x limits=false,ylabel shift=-3pt,xlabel shift=-3pt]\addplot[fill=ink!55,draw=none] coordinates {(0,0.04755) (1,0.2143) (2,0.867) (3,6.62) (4,24.17) (5,100.8) (6,220.7) (7,376.1) (8,350.2) (9,264.8) (10,82.15) (11,28.02) (12,1.315) (13,0.5394) (14,0.001721) (15,0.03816) (16,5.1e-05) (17,2e-06)};\draw[densely dotted] (axis cs:12,0) -- (axis cs:12,413.7);\end{axis}\begin{axis}[width=72mm,height=42mm,axis y line*=right,axis x line=none,ymode=log,xmin=-1,xmax=18,ylabel={\scriptsize eindeutig (log)},tick label style={font=\tiny},label style={font=\scriptsize},ylabel shift=-3pt]\addplot[mark=*,mark size=0.7pt,draw=black,thin] coordinates {(0,47554) (1,196830) (2,544512) (3,2643336) (4,6495309) (5,10388517) (6,8594878) (7,5175512) (8,1167351) (9,315850) (10,12951) (11,2726) (12,76)};\end{axis}\end{tikzpicture}\end{minipage}
\vspace{3mm}

\noindent\begin{minipage}[t]{60mm}\vspace{0pt}\sbox\hmboard{\chessboard[setfen=N5N1/8/8/7N/8/8/4q3/K6k b - - 0 1]}\usebox\hmboard\\[1mm]
\makebox[\wd\hmboard]{\sffamily\hspace{3mm}\textbf{h\#6}\hfill(4+2)}\end{minipage}\hfill
\begin{minipage}[t]{52mm}\vspace{0pt}
\textcolor{ink}{\sffamily\textsc{Lösung}}\par\smallskip\small
\begin{tabular}{@{}r@{\enspace}l@{\enspace}l@{}}
\textbf{1.} & \symking{}g2 & \symking{}b1\\
\textbf{2.} & \symqueen{}b5+ & \symking{}c2\\
\textbf{3.} & \symking{}f3 & \symking{}c3\\
\textbf{4.} & \symking{}e4 & \symknight{}g3+\\
\textbf{5.} & \symking{}d5 & \symknight{}e7+\\
\textbf{6.} & \symking{}c5 & \symknight{}e4\#\\
\end{tabular}\par\vspace{3mm}
\qrcode[height=12mm]{https://helpman.komtera.lt/?fen=N5N1/8/8/7N/8/8/4q3/K6k&moves=6}\\{\sffamily\tiny\color{ink}scannen zum Lösen}
\end{minipage}
\vspace{2mm}

{\raggedright\noindent\textcolor{ink}{\sffamily\textsc{Themen}}\enspace \themetag{reines Matt} \themetag{Mustermatt} \themetag{Idealmatt} \themetag{Selbstblock} \themetag{ohne Schlag}\par}

\clearpage
\phantomsection\label{hm:KNNNvkr}
\begin{tcolorbox}[enhanced,colback=ink,colframe=ink,arc=0pt,boxrule=0pt,left=4mm,right=4mm,top=2.5mm,bottom=2.5mm]
\color{white}\sffamily{\LARGE\bfseries KNNNvkr}\hfill{\Large h\#6}\\[0.5mm]
{\small 6 Steine\quad Schwarz am Zug}\hfill{\small Nr.~656\qquad\href{https://helpman.komtera.lt/?fen=7k/5Nr1/5N2/8/8/7N/8/1K6&moves=6}{\color{white}\bfseries lösen\,\ensuremath{\nearrow}}}
\end{tcolorbox}
\vspace{2mm}
\noindent\begin{minipage}[t]{36mm}\vspace{0pt}\sffamily\small\raggedright
\hmfigure{längste eindeutige}{{\Large h\#6}}\hmfigure{eindeutig bei dieser Länge}{{\Large 1} Stellung}\hmfigure{längstes Matt}{{\Large h\#8}}\hmfigure{in der Klasse}{15,5 Milliarden Stellungen,\\ 1,59 Milliarden mit Hilfsmatt}
\end{minipage}\hfill
\begin{minipage}[t]{80mm}\vspace{0pt}\begin{tikzpicture}\begin{axis}[width=72mm,height=42mm,ybar,bar width=1.6pt,xmin=-1,xmax=17,ymin=0,ymax=475.2,axis y line*=left,axis x line*=bottom,xlabel={\scriptsize Halbzüge bis Matt},ylabel={\scriptsize Stellungen (Mio.)},tick label style={font=\tiny},label style={font=\scriptsize},enlarge x limits=false,ylabel shift=-3pt,xlabel shift=-3pt]\addplot[fill=ink!55,draw=none] coordinates {(0,0.07388) (1,0.332) (2,1.426) (3,10.26) (4,48.6) (5,167.6) (6,375) (7,432) (8,341.5) (9,159.1) (10,41.44) (11,8.122) (12,0.3142) (13,0.02368) (14,0.000607) (15,0.001323) (16,1e-05)};\draw[densely dotted] (axis cs:12,0) -- (axis cs:12,475.2);\end{axis}\begin{axis}[width=72mm,height=42mm,axis y line*=right,axis x line=none,ymode=log,xmin=-1,xmax=17,ylabel={\scriptsize eindeutig (log)},tick label style={font=\tiny},label style={font=\scriptsize},ylabel shift=-3pt]\addplot[mark=*,mark size=0.7pt,draw=black,thin] coordinates {(0,73879) (1,299688) (2,1004692) (3,4477915) (4,16537730) (5,15967142) (6,12127672) (7,2734696) (8,476639) (9,60467) (10,1879) (11,178) (12,1)};\end{axis}\end{tikzpicture}\end{minipage}
\vspace{3mm}

\noindent\begin{minipage}[t]{60mm}\vspace{0pt}\sbox\hmboard{\chessboard[setfen=7k/5Nr1/5N2/8/8/7N/8/1K6 b - - 0 1]}\usebox\hmboard\\[1mm]
\makebox[\wd\hmboard]{\sffamily\hspace{3mm}\textbf{h\#6}\hfill(4+2)}\end{minipage}\hfill
\begin{minipage}[t]{52mm}\vspace{0pt}
\textcolor{ink}{\sffamily\textsc{Lösung}}\par\smallskip\small
\begin{tabular}{@{}r@{\enspace}l@{\enspace}l@{}}
\textbf{1.} & \symrook{}xf7 & \symking{}c2\\
\textbf{2.} & \symrook{}h7 & \symking{}d3\\
\textbf{3.} & \symking{}g7 & \symking{}e4\\
\textbf{4.} & \symrook{}h4+ & \symking{}f5\\
\textbf{5.} & \symking{}h6 & \symknight{}g4+\\
\textbf{6.} & \symking{}h5 & \symknight{}f4\#\\
\end{tabular}\par\vspace{3mm}
\qrcode[height=12mm]{https://helpman.komtera.lt/?fen=7k/5Nr1/5N2/8/8/7N/8/1K6&moves=6}\\{\sffamily\tiny\color{ink}scannen zum Lösen}
\end{minipage}
\vspace{2mm}

{\raggedright\noindent\textcolor{ink}{\sffamily\textsc{Themen}}\enspace \themetag{Selbstblock}\par\smallskip{\footnotesize\textit{Hinweis:} die Partei am Zug steht im Diagramm im Schach; die Lösung beginnt mit einem Schlag.}\par}

\clearpage
\phantomsection\label{hm:KQBBvkq}
\begin{tcolorbox}[enhanced,colback=ink,colframe=ink,arc=0pt,boxrule=0pt,left=4mm,right=4mm,top=2.5mm,bottom=2.5mm]
\color{white}\sffamily{\LARGE\bfseries KQBBvkq}\hfill{\Large h\#6}\\[0.5mm]
{\small 6 Steine\quad Schwarz am Zug}\hfill{\small Nr.~657\qquad\href{https://helpman.komtera.lt/?fen=2B5/6Q1/5k2/8/8/4B3/3q4/K7&moves=6}{\color{white}\bfseries lösen\,\ensuremath{\nearrow}}}
\end{tcolorbox}
\vspace{2mm}
\noindent\begin{minipage}[t]{36mm}\vspace{0pt}\sffamily\small\raggedright
\hmfigure{längste eindeutige}{{\Large h\#6}}\hmfigure{eindeutig bei dieser Länge}{{\Large 256} Stellungen}\hmfigure{längstes Matt}{{\Large h\#7}}\hmfigure{in der Klasse}{15,5 Milliarden Stellungen,\\ 3,53 Milliarden mit Hilfsmatt}
\end{minipage}\hfill
\begin{minipage}[t]{80mm}\vspace{0pt}\begin{tikzpicture}\begin{axis}[width=72mm,height=42mm,ybar,bar width=1.6pt,xmin=-1,xmax=15,ymin=0,ymax=1.681,axis y line*=left,axis x line*=bottom,xlabel={\scriptsize Halbzüge bis Matt},ylabel={\scriptsize Stellungen (Mrd.)},tick label style={font=\tiny},label style={font=\scriptsize},enlarge x limits=false,ylabel shift=-3pt,xlabel shift=-3pt]\addplot[fill=ink!55,draw=none] coordinates {(0,0.0145) (1,0.07539) (2,0.3978) (3,0.9221) (4,1.528) (5,0.5175) (6,0.06294) (7,0.002196) (8,0.007259) (9,1.71e-07) (10,0.004229) (11,7e-09) (12,0.0007552) (14,5.666e-05)};\draw[densely dotted] (axis cs:12,0) -- (axis cs:12,1.681);\end{axis}\begin{axis}[width=72mm,height=42mm,axis y line*=right,axis x line=none,ymode=log,xmin=-1,xmax=15,ylabel={\scriptsize eindeutig (log)},tick label style={font=\tiny},label style={font=\scriptsize},ylabel shift=-3pt]\addplot[mark=*,mark size=0.7pt,draw=black,thin] coordinates {(0,14499370) (1,63499659) (2,217822433) (3,114977019) (4,25864903) (5,585629) (6,338897) (7,6) (8,67661) (9,4) (10,8642) (12,256)};\end{axis}\end{tikzpicture}\end{minipage}
\vspace{3mm}

\noindent\begin{minipage}[t]{60mm}\vspace{0pt}\sbox\hmboard{\chessboard[setfen=2B5/6Q1/5k2/8/8/4B3/3q4/K7 b - - 0 1]}\usebox\hmboard\\[1mm]
\makebox[\wd\hmboard]{\sffamily\hspace{3mm}\textbf{h\#6}\hfill(4+2)}\end{minipage}\hfill
\begin{minipage}[t]{52mm}\vspace{0pt}
\textcolor{ink}{\sffamily\textsc{Lösung}}\par\smallskip\small
\begin{tabular}{@{}r@{\enspace}l@{\enspace}l@{}}
\textbf{1.} & \symking{}xg7 & \symking{}b1\\
\textbf{2.} & \symqueen{}a5 & \symking{}c2\\
\textbf{3.} & \symking{}g6 & \symking{}d3\\
\textbf{4.} & \symking{}h5 & \symking{}e4\\
\textbf{5.} & \symking{}h4 & \symking{}f4\\
\textbf{6.} & \symqueen{}h5 & \symbishop{}f2\#\\
\end{tabular}\par\vspace{3mm}
\qrcode[height=12mm]{https://helpman.komtera.lt/?fen=2B5/6Q1/5k2/8/8/4B3/3q4/K7&moves=6}\\{\sffamily\tiny\color{ink}scannen zum Lösen}
\end{minipage}
\vspace{2mm}

{\raggedright\noindent\textcolor{ink}{\sffamily\textsc{Themen}}\enspace \themetag{Selbstblock} \themetag{ohne Schach}\par\smallskip{\footnotesize\textit{Hinweis:} die Partei am Zug steht im Diagramm im Schach; die Lösung beginnt mit einem Schlag.}\par}

\clearpage
\phantomsection\label{hm:KQBBvkr}
\begin{tcolorbox}[enhanced,colback=ink,colframe=ink,arc=0pt,boxrule=0pt,left=4mm,right=4mm,top=2.5mm,bottom=2.5mm]
\color{white}\sffamily{\LARGE\bfseries KQBBvkr}\hfill{\Large h\#6}\\[0.5mm]
{\small 6 Steine\quad Schwarz am Zug}\hfill{\small Nr.~658\qquad\href{https://helpman.komtera.lt/?fen=8/7Q/6k1/3B4/7B/6r1/8/K7&moves=6}{\color{white}\bfseries lösen\,\ensuremath{\nearrow}}}
\end{tcolorbox}
\vspace{2mm}
\noindent\begin{minipage}[t]{36mm}\vspace{0pt}\sffamily\small\raggedright
\hmfigure{längste eindeutige}{{\Large h\#6}}\hmfigure{eindeutig bei dieser Länge}{{\Large 291} Stellungen}\hmfigure{längstes Matt}{{\Large h\#7}}\hmfigure{in der Klasse}{15,5 Milliarden Stellungen,\\ 3,92 Milliarden mit Hilfsmatt}
\end{minipage}\hfill
\begin{minipage}[t]{80mm}\vspace{0pt}\begin{tikzpicture}\begin{axis}[width=72mm,height=42mm,ybar,bar width=1.6pt,xmin=-1,xmax=15,ymin=0,ymax=1.97,axis y line*=left,axis x line*=bottom,xlabel={\scriptsize Halbzüge bis Matt},ylabel={\scriptsize Stellungen (Mrd.)},tick label style={font=\tiny},label style={font=\scriptsize},enlarge x limits=false,ylabel shift=-3pt,xlabel shift=-3pt]\addplot[fill=ink!55,draw=none] coordinates {(0,0.02021) (1,0.103) (2,0.5095) (3,1.013) (4,1.791) (5,0.4019) (6,0.06399) (7,0.0006062) (8,0.008285) (9,9.8e-08) (10,0.004811) (11,7e-09) (12,0.0008877) (14,6.838e-05)};\draw[densely dotted] (axis cs:12,0) -- (axis cs:12,1.97);\end{axis}\begin{axis}[width=72mm,height=42mm,axis y line*=right,axis x line=none,ymode=log,xmin=-1,xmax=15,ylabel={\scriptsize eindeutig (log)},tick label style={font=\tiny},label style={font=\scriptsize},ylabel shift=-3pt]\addplot[mark=*,mark size=0.7pt,draw=black,thin] coordinates {(0,20211757) (1,83865060) (2,277372508) (3,105786119) (4,27667941) (5,152927) (6,397445) (7,1) (8,83976) (9,4) (10,11455) (12,291)};\end{axis}\end{tikzpicture}\end{minipage}
\vspace{3mm}

\noindent\begin{minipage}[t]{60mm}\vspace{0pt}\sbox\hmboard{\chessboard[setfen=8/7Q/6k1/3B4/7B/6r1/8/K7 b - - 0 1]}\usebox\hmboard\\[1mm]
\makebox[\wd\hmboard]{\sffamily\hspace{3mm}\textbf{h\#6}\hfill(4+2)}\end{minipage}\hfill
\begin{minipage}[t]{52mm}\vspace{0pt}
\textcolor{ink}{\sffamily\textsc{Lösung}}\par\smallskip\small
\begin{tabular}{@{}r@{\enspace}l@{\enspace}l@{}}
\textbf{1.} & \symking{}xh7 & \symking{}b2\\
\textbf{2.} & \symrook{}g8 & \symking{}c3\\
\textbf{3.} & \symking{}g7 & \symking{}d4\\
\textbf{4.} & \symking{}f8 & \symking{}e5\\
\textbf{5.} & \symking{}e8 & \symking{}e6\\
\textbf{6.} & \symrook{}f8 & \symbishop{}c6\#\\
\end{tabular}\par\vspace{3mm}
\qrcode[height=12mm]{https://helpman.komtera.lt/?fen=8/7Q/6k1/3B4/7B/6r1/8/K7&moves=6}\\{\sffamily\tiny\color{ink}scannen zum Lösen}
\end{minipage}
\vspace{2mm}

{\raggedright\noindent\textcolor{ink}{\sffamily\textsc{Themen}}\enspace \themetag{Selbstblock} \themetag{ohne Schach}\par\smallskip{\footnotesize\textit{Hinweis:} die Partei am Zug steht im Diagramm im Schach; die Lösung beginnt mit einem Schlag.}\par}

\clearpage
\phantomsection\label{hm:KQNNNvk}
\begin{tcolorbox}[enhanced,colback=ink,colframe=ink,arc=0pt,boxrule=0pt,left=4mm,right=4mm,top=2.5mm,bottom=2.5mm]
\color{white}\sffamily{\LARGE\bfseries KQNNNvk}\hfill{\Large h\#6}\\[0.5mm]
{\small 6 Steine\quad Schwarz am Zug}\hfill{\small Nr.~659\qquad\href{https://helpman.komtera.lt/?fen=8/5N1Q/7k/8/8/8/8/1K3NN1&moves=6}{\color{white}\bfseries lösen\,\ensuremath{\nearrow}}}
\end{tcolorbox}
\vspace{2mm}
\noindent\begin{minipage}[t]{36mm}\vspace{0pt}\sffamily\small\raggedright
\hmfigure{längste eindeutige}{{\Large h\#6}}\hmfigure{eindeutig bei dieser Länge}{{\Large 81} Stellungen}\hmfigure{längstes Matt}{{\Large h\#8}}\hmfigure{in der Klasse}{15,5 Milliarden Stellungen,\\ 1,53 Milliarden mit Hilfsmatt}
\end{minipage}\hfill
\begin{minipage}[t]{80mm}\vspace{0pt}\begin{tikzpicture}\begin{axis}[width=72mm,height=42mm,ybar,bar width=1.6pt,xmin=-1,xmax=17,ymin=0,ymax=772.6,axis y line*=left,axis x line*=bottom,xlabel={\scriptsize Halbzüge bis Matt},ylabel={\scriptsize Stellungen (Mio.)},tick label style={font=\tiny},label style={font=\scriptsize},enlarge x limits=false,ylabel shift=-3pt,xlabel shift=-3pt]\addplot[fill=ink!55,draw=none] coordinates {(0,8.27) (1,47.86) (2,172.6) (3,317.7) (4,702.4) (5,162.8) (6,98.24) (7,1.816) (8,6.947) (10,5.42) (12,1.491) (14,0.3339) (16,0.000655)};\draw[densely dotted] (axis cs:12,0) -- (axis cs:12,772.6);\end{axis}\begin{axis}[width=72mm,height=42mm,axis y line*=right,axis x line=none,ymode=log,xmin=-1,xmax=17,ylabel={\scriptsize eindeutig (log)},tick label style={font=\tiny},label style={font=\scriptsize},ylabel shift=-3pt]\addplot[mark=*,mark size=0.7pt,draw=black,thin] coordinates {(0,8269867) (1,37025985) (2,116406033) (3,47579543) (4,30694548) (5,172060) (6,302220) (8,67155) (10,6971) (12,81)};\end{axis}\end{tikzpicture}\end{minipage}
\vspace{3mm}

\noindent\begin{minipage}[t]{60mm}\vspace{0pt}\sbox\hmboard{\chessboard[setfen=8/5N1Q/7k/8/8/8/8/1K3NN1 b - - 0 1]}\usebox\hmboard\\[1mm]
\makebox[\wd\hmboard]{\sffamily\hspace{3mm}\textbf{h\#6}\hfill(5+1)}\end{minipage}\hfill
\begin{minipage}[t]{52mm}\vspace{0pt}
\textcolor{ink}{\sffamily\textsc{Lösung}}\par\smallskip\small
\begin{tabular}{@{}r@{\enspace}l@{\enspace}l@{}}
\textbf{1.} & \symking{}xh7 & \symking{}c2\\
\textbf{2.} & \symking{}g7 & \symking{}d3\\
\textbf{3.} & \symking{}g6 & \symking{}e4\\
\textbf{4.} & \symking{}h5 & \symking{}f5\\
\textbf{5.} & \symking{}h4 & \symknight{}f3+\\
\textbf{6.} & \symking{}h5 & \symknight{}g3\#\\
\end{tabular}\par\vspace{3mm}
\qrcode[height=12mm]{https://helpman.komtera.lt/?fen=8/5N1Q/7k/8/8/8/8/1K3NN1&moves=6}\\{\sffamily\tiny\color{ink}scannen zum Lösen}
\end{minipage}
\vspace{2mm}

{\raggedright\noindent\textcolor{ink}{\sffamily\textsc{Themen}}\enspace \themetag{Spiegelmatt} \themetag{Switchback} \themetag{ein Stein zieht}\par\smallskip{\footnotesize\textit{Hinweis:} das Diagramm hat keinen legalen letzten Zug, kann also in keiner Partie entstehen.}\par}

\clearpage
\phantomsection\label{hm:KQNNvkn}
\begin{tcolorbox}[enhanced,colback=ink,colframe=ink,arc=0pt,boxrule=0pt,left=4mm,right=4mm,top=2.5mm,bottom=2.5mm]
\color{white}\sffamily{\LARGE\bfseries KQNNvkn}\hfill{\Large h\#6}\\[0.5mm]
{\small 6 Steine\quad Schwarz am Zug}\hfill{\small Nr.~660\qquad\href{https://helpman.komtera.lt/?fen=4N3/8/6Qk/n7/8/8/8/K2N4&moves=6}{\color{white}\bfseries lösen\,\ensuremath{\nearrow}}}
\end{tcolorbox}
\vspace{2mm}
\noindent\begin{minipage}[t]{36mm}\vspace{0pt}\sffamily\small\raggedright
\hmfigure{längste eindeutige}{{\Large h\#6}}\hmfigure{eindeutig bei dieser Länge}{{\Large 248} Stellungen}\hmfigure{längstes Matt}{{\Large h\#8}}\hmfigure{in der Klasse}{15,5 Milliarden Stellungen,\\ 4,5 Milliarden mit Hilfsmatt}
\end{minipage}\hfill
\begin{minipage}[t]{80mm}\vspace{0pt}\begin{tikzpicture}\begin{axis}[width=72mm,height=42mm,ybar,bar width=1.6pt,xmin=-1,xmax=17,ymin=0,ymax=2.079,axis y line*=left,axis x line*=bottom,xlabel={\scriptsize Halbzüge bis Matt},ylabel={\scriptsize Stellungen (Mrd.)},tick label style={font=\tiny},label style={font=\scriptsize},enlarge x limits=false,ylabel shift=-3pt,xlabel shift=-3pt]\addplot[fill=ink!55,draw=none] coordinates {(0,0.01221) (1,0.08292) (2,0.3412) (3,0.8787) (4,1.89) (5,0.7325) (6,0.4863) (7,0.03383) (8,0.01835) (9,7.919e-05) (10,0.01814) (11,5.41e-07) (12,0.006904) (13,1.43e-07) (14,0.0004033) (15,1.3e-08) (16,1.47e-07)};\draw[densely dotted] (axis cs:12,0) -- (axis cs:12,2.079);\end{axis}\begin{axis}[width=72mm,height=42mm,axis y line*=right,axis x line=none,ymode=log,xmin=-1,xmax=17,ylabel={\scriptsize eindeutig (log)},tick label style={font=\tiny},label style={font=\scriptsize},ylabel shift=-3pt]\addplot[mark=*,mark size=0.7pt,draw=black,thin] coordinates {(0,12214984) (1,67722546) (2,193229663) (3,128452394) (4,92668274) (5,3119632) (6,562010) (7,3058) (8,141418) (9,32) (10,27372) (12,248)};\end{axis}\end{tikzpicture}\end{minipage}
\vspace{3mm}

\noindent\begin{minipage}[t]{60mm}\vspace{0pt}\sbox\hmboard{\chessboard[setfen=4N3/8/6Qk/n7/8/8/8/K2N4 b - - 0 1]}\usebox\hmboard\\[1mm]
\makebox[\wd\hmboard]{\sffamily\hspace{3mm}\textbf{h\#6}\hfill(4+2)}\end{minipage}\hfill
\begin{minipage}[t]{52mm}\vspace{0pt}
\textcolor{ink}{\sffamily\textsc{Lösung}}\par\smallskip\small
\begin{tabular}{@{}r@{\enspace}l@{\enspace}l@{}}
\textbf{1.} & \symking{}xg6 & \symking{}a2\\
\textbf{2.} & \symking{}f5 & \symknight{}e3+\\
\textbf{3.} & \symking{}e4 & \symknight{}d6+\\
\textbf{4.} & \symking{}d3 & \symknight{}b7\\
\textbf{5.} & \symking{}d2 & \symknight{}xa5\\
\textbf{6.} & \symking{}c1 & \symknight{}b3\#\\
\end{tabular}\par\vspace{3mm}
\qrcode[height=12mm]{https://helpman.komtera.lt/?fen=4N3/8/6Qk/n7/8/8/8/K2N4&moves=6}\\{\sffamily\tiny\color{ink}scannen zum Lösen}
\end{minipage}
\vspace{2mm}

{\raggedright\noindent\textcolor{ink}{\sffamily\textsc{Themen}}\enspace \themetag{reines Matt} \themetag{Mustermatt} \themetag{Idealmatt} \themetag{Spiegelmatt} \themetag{ein Stein zieht}\par\smallskip{\footnotesize\textit{Hinweis:} die Partei am Zug steht im Diagramm im Schach; die Lösung beginnt mit einem Schlag.}\par}

\clearpage
\phantomsection\label{hm:KQQPPvk}
\begin{tcolorbox}[enhanced,colback=ink,colframe=ink,arc=0pt,boxrule=0pt,left=4mm,right=4mm,top=2.5mm,bottom=2.5mm]
\color{white}\sffamily{\LARGE\bfseries KQQPPvk}\hfill{\Large h\#6}\\[0.5mm]
{\small 6 Steine\quad Schwarz am Zug}\hfill{\small Nr.~661\qquad\href{https://helpman.komtera.lt/?fen=K7/8/8/8/8/P7/P3Q3/6kQ&moves=6}{\color{white}\bfseries lösen\,\ensuremath{\nearrow}}}
\end{tcolorbox}
\vspace{2mm}
\noindent\begin{minipage}[t]{36mm}\vspace{0pt}\sffamily\small\raggedright
\hmfigure{längste eindeutige}{{\Large h\#6}}\hmfigure{eindeutig bei dieser Länge}{{\Large 3} Stellungen}\hmfigure{längstes Matt}{{\Large h\#7}}\hmfigure{in der Klasse}{34,1 Milliarden Stellungen,\\ 4,74 Milliarden mit Hilfsmatt}
\end{minipage}\hfill
\begin{minipage}[t]{80mm}\vspace{0pt}\begin{tikzpicture}\begin{axis}[width=72mm,height=42mm,ybar,bar width=1.6pt,xmin=-1,xmax=15,ymin=0,ymax=1.936,axis y line*=left,axis x line*=bottom,xlabel={\scriptsize Halbzüge bis Matt},ylabel={\scriptsize Stellungen (Mrd.)},tick label style={font=\tiny},label style={font=\scriptsize},enlarge x limits=false,ylabel shift=-3pt,xlabel shift=-3pt]\addplot[fill=ink!55,draw=none] coordinates {(0,0.1304) (1,0.3611) (2,1.337) (3,1.029) (4,1.76) (5,0.03583) (6,0.06776) (7,2.67e-07) (8,0.01837) (10,0.001383) (12,1.499e-05) (14,1.7e-08)};\draw[densely dotted] (axis cs:12,0) -- (axis cs:12,1.936);\end{axis}\begin{axis}[width=72mm,height=42mm,axis y line*=right,axis x line=none,ymode=log,xmin=-1,xmax=15,ylabel={\scriptsize eindeutig (log)},tick label style={font=\tiny},label style={font=\scriptsize},ylabel shift=-3pt]\addplot[mark=*,mark size=0.7pt,draw=black,thin] coordinates {(0,130375432) (1,174008605) (2,501927421) (3,24248486) (4,9323287) (5,30) (6,1665065) (8,74698) (10,280) (12,3)};\end{axis}\end{tikzpicture}\end{minipage}
\vspace{3mm}

\noindent\begin{minipage}[t]{60mm}\vspace{0pt}\sbox\hmboard{\chessboard[setfen=K7/8/8/8/8/P7/P3Q3/6kQ b - - 0 1]}\usebox\hmboard\\[1mm]
\makebox[\wd\hmboard]{\sffamily\hspace{3mm}\textbf{h\#6}\hfill(5+1)}\end{minipage}\hfill
\begin{minipage}[t]{52mm}\vspace{0pt}
\textcolor{ink}{\sffamily\textsc{Lösung}}\par\smallskip\small
\begin{tabular}{@{}r@{\enspace}l@{\enspace}l@{}}
\textbf{1.} & \symking{}xh1 & \symking{}b7\\
\textbf{2.} & \symking{}g1 & \symking{}c6\\
\textbf{3.} & \symking{}h1 & \symking{}d5\\
\textbf{4.} & \symking{}g1 & \symking{}e4\\
\textbf{5.} & \symking{}h1 & \symking{}f3\\
\textbf{6.} & \symking{}g1 & \symqueen{}g2\#\\
\end{tabular}\par\vspace{3mm}
\qrcode[height=12mm]{https://helpman.komtera.lt/?fen=K7/8/8/8/8/P7/P3Q3/6kQ&moves=6}\\{\sffamily\tiny\color{ink}scannen zum Lösen}
\end{minipage}
\vspace{2mm}

{\raggedright\noindent\textcolor{ink}{\sffamily\textsc{Themen}}\enspace \themetag{Switchback} \themetag{ein Stein zieht} \themetag{Pendel} \themetag{ohne Schach}\par\smallskip{\footnotesize\textit{Hinweis:} die Partei am Zug steht im Diagramm im Schach; die Lösung beginnt mit einem Schlag.}\par}

\clearpage
\phantomsection\label{hm:KQQPvkp}
\begin{tcolorbox}[enhanced,colback=ink,colframe=ink,arc=0pt,boxrule=0pt,left=4mm,right=4mm,top=2.5mm,bottom=2.5mm]
\color{white}\sffamily{\LARGE\bfseries KQQPvkp}\hfill{\Large h\#6}\\[0.5mm]
{\small 6 Steine\quad Schwarz am Zug}\hfill{\small Nr.~662\qquad\href{https://helpman.komtera.lt/?fen=7Q/7k/8/6Q1/8/3p4/3P4/K7&moves=6}{\color{white}\bfseries lösen\,\ensuremath{\nearrow}}}
\end{tcolorbox}
\vspace{2mm}
\noindent\begin{minipage}[t]{36mm}\vspace{0pt}\sffamily\small\raggedright
\hmfigure{längste eindeutige}{{\Large h\#6}}\hmfigure{eindeutig bei dieser Länge}{{\Large 34} Stellungen}\hmfigure{längstes Matt}{{\Large h\#7}}\hmfigure{in der Klasse}{34,1 Milliarden Stellungen,\\ 9,51 Milliarden mit Hilfsmatt}
\end{minipage}\hfill
\begin{minipage}[t]{80mm}\vspace{0pt}\begin{tikzpicture}\begin{axis}[width=72mm,height=42mm,ybar,bar width=1.6pt,xmin=-1,xmax=15,ymin=0,ymax=3.709,axis y line*=left,axis x line*=bottom,xlabel={\scriptsize Halbzüge bis Matt},ylabel={\scriptsize Stellungen (Mrd.)},tick label style={font=\tiny},label style={font=\scriptsize},enlarge x limits=false,ylabel shift=-3pt,xlabel shift=-3pt]\addplot[fill=ink!55,draw=none] coordinates {(0,0.2374) (1,0.688) (2,2.77) (3,2.17) (4,3.372) (5,0.08195) (6,0.1233) (7,4.222e-06) (8,0.05808) (9,4.3e-08) (10,0.00786) (11,3e-09) (12,0.0002269) (14,4.4e-07)};\draw[densely dotted] (axis cs:12,0) -- (axis cs:12,3.709);\end{axis}\begin{axis}[width=72mm,height=42mm,axis y line*=right,axis x line=none,ymode=log,xmin=-1,xmax=15,ylabel={\scriptsize eindeutig (log)},tick label style={font=\tiny},label style={font=\scriptsize},ylabel shift=-3pt]\addplot[mark=*,mark size=0.7pt,draw=black,thin] coordinates {(0,237367889) (1,344925466) (2,943138338) (3,36759823) (4,12908892) (5,498) (6,3648275) (7,5) (8,346385) (10,2398) (12,34)};\end{axis}\end{tikzpicture}\end{minipage}
\vspace{3mm}

\noindent\begin{minipage}[t]{60mm}\vspace{0pt}\sbox\hmboard{\chessboard[setfen=7Q/7k/8/6Q1/8/3p4/3P4/K7 b - - 0 1]}\usebox\hmboard\\[1mm]
\makebox[\wd\hmboard]{\sffamily\hspace{3mm}\textbf{h\#6}\hfill(4+2)}\end{minipage}\hfill
\begin{minipage}[t]{52mm}\vspace{0pt}
\textcolor{ink}{\sffamily\textsc{Lösung}}\par\smallskip\small
\begin{tabular}{@{}r@{\enspace}l@{\enspace}l@{}}
\textbf{1.} & \symking{}xh8 & \symking{}b2\\
\textbf{2.} & \symking{}h7 & \symking{}c3\\
\textbf{3.} & \symking{}h8 & \symking{}d4\\
\textbf{4.} & \symking{}h7 & \symking{}e5\\
\textbf{5.} & \symking{}h8 & \symking{}f6\\
\textbf{6.} & \symking{}h7 & \symqueen{}g7\#\\
\end{tabular}\par\vspace{3mm}
\qrcode[height=12mm]{https://helpman.komtera.lt/?fen=7Q/7k/8/6Q1/8/3p4/3P4/K7&moves=6}\\{\sffamily\tiny\color{ink}scannen zum Lösen}
\end{minipage}
\vspace{2mm}

{\raggedright\noindent\textcolor{ink}{\sffamily\textsc{Themen}}\enspace \themetag{Switchback} \themetag{ein Stein zieht} \themetag{Pendel} \themetag{ohne Schach}\par\smallskip{\footnotesize\textit{Hinweis:} die Partei am Zug steht im Diagramm im Schach; die Lösung beginnt mit einem Schlag.}\par}

\clearpage
\phantomsection\label{hm:KQQPvkq}
\begin{tcolorbox}[enhanced,colback=ink,colframe=ink,arc=0pt,boxrule=0pt,left=4mm,right=4mm,top=2.5mm,bottom=2.5mm]
\color{white}\sffamily{\LARGE\bfseries KQQPvkq}\hfill{\Large h\#6}\\[0.5mm]
{\small 6 Steine\quad Schwarz am Zug}\hfill{\small Nr.~663\qquad\href{https://helpman.komtera.lt/?fen=K7/8/8/Qq6/8/P7/7Q/7k&moves=6}{\color{white}\bfseries lösen\,\ensuremath{\nearrow}}}
\end{tcolorbox}
\vspace{2mm}
\noindent\begin{minipage}[t]{36mm}\vspace{0pt}\sffamily\small\raggedright
\hmfigure{längste eindeutige}{{\Large h\#6}}\hmfigure{eindeutig bei dieser Länge}{{\Large 3} Stellungen}\hmfigure{längstes Matt}{{\Large h\#7}}\hmfigure{in der Klasse}{45,4 Milliarden Stellungen,\\ 10 Milliarden mit Hilfsmatt}
\end{minipage}\hfill
\begin{minipage}[t]{80mm}\vspace{0pt}\begin{tikzpicture}\begin{axis}[width=72mm,height=42mm,ybar,bar width=1.6pt,xmin=-1,xmax=15,ymin=0,ymax=3.68,axis y line*=left,axis x line*=bottom,xlabel={\scriptsize Halbzüge bis Matt},ylabel={\scriptsize Stellungen (Mrd.)},tick label style={font=\tiny},label style={font=\scriptsize},enlarge x limits=false,ylabel shift=-3pt,xlabel shift=-3pt]\addplot[fill=ink!55,draw=none] coordinates {(0,0.1541) (1,0.4902) (2,2.44) (3,2.854) (4,3.346) (5,0.5659) (6,0.127) (7,4.371e-05) (8,0.05961) (9,8.13e-07) (10,0.009291) (11,5e-09) (12,0.0002391) (14,1.39e-07)};\draw[densely dotted] (axis cs:12,0) -- (axis cs:12,3.68);\end{axis}\begin{axis}[width=72mm,height=42mm,axis y line*=right,axis x line=none,ymode=log,xmin=-1,xmax=15,ylabel={\scriptsize eindeutig (log)},tick label style={font=\tiny},label style={font=\scriptsize},ylabel shift=-3pt]\addplot[mark=*,mark size=0.7pt,draw=black,thin] coordinates {(0,154070320) (1,298402564) (2,861209705) (3,126838524) (4,12989233) (5,8569) (6,3350084) (7,38) (8,437557) (10,5775) (12,3)};\end{axis}\end{tikzpicture}\end{minipage}
\vspace{3mm}

\noindent\begin{minipage}[t]{60mm}\vspace{0pt}\sbox\hmboard{\chessboard[setfen=K7/8/8/Qq6/8/P7/7Q/7k b - - 0 1]}\usebox\hmboard\\[1mm]
\makebox[\wd\hmboard]{\sffamily\hspace{3mm}\textbf{h\#6}\hfill(4+2)}\end{minipage}\hfill
\begin{minipage}[t]{52mm}\vspace{0pt}
\textcolor{ink}{\sffamily\textsc{Lösung}}\par\smallskip\small
\begin{tabular}{@{}r@{\enspace}l@{\enspace}l@{}}
\textbf{1.} & \symking{}xh2 & a4\\
\textbf{2.} & \symking{}g3 & axb5\\
\textbf{3.} & \symking{}f4 & \symqueen{}a4+\\
\textbf{4.} & \symking{}e5 & b6\\
\textbf{5.} & \symking{}d6 & b7\\
\textbf{6.} & \symking{}c7 & b8=\symqueen{}\#\\
\end{tabular}\par\vspace{3mm}
\qrcode[height=12mm]{https://helpman.komtera.lt/?fen=K7/8/8/Qq6/8/P7/7Q/7k&moves=6}\\{\sffamily\tiny\color{ink}scannen zum Lösen}
\end{minipage}
\vspace{2mm}

{\raggedright\noindent\textcolor{ink}{\sffamily\textsc{Themen}}\enspace \themetag{Umwandlung} \themetag{ein Stein zieht} \themetag{Phönix}\par\smallskip{\footnotesize\textit{Hinweis:} die Partei am Zug steht im Diagramm im Schach; die Lösung beginnt mit einem Schlag.}\par}

\clearpage
\phantomsection\label{hm:KQQvkbn}
\begin{tcolorbox}[enhanced,colback=ink,colframe=ink,arc=0pt,boxrule=0pt,left=4mm,right=4mm,top=2.5mm,bottom=2.5mm]
\color{white}\sffamily{\LARGE\bfseries KQQvkbn}\hfill{\Large h\#6}\\[0.5mm]
{\small 6 Steine\quad Schwarz am Zug}\hfill{\small Nr.~664\qquad\href{https://helpman.komtera.lt/?fen=8/6Q1/7k/8/8/b7/7Q/K1n5&moves=6}{\color{white}\bfseries lösen\,\ensuremath{\nearrow}}}
\end{tcolorbox}
\vspace{2mm}
\noindent\begin{minipage}[t]{36mm}\vspace{0pt}\sffamily\small\raggedright
\hmfigure{längste eindeutige}{{\Large h\#6}}\hmfigure{eindeutig bei dieser Länge}{{\Large 9} Stellungen}\hmfigure{längstes Matt}{{\Large h\#7}}\hmfigure{in der Klasse}{15,5 Milliarden Stellungen,\\ 3,76 Milliarden mit Hilfsmatt}
\end{minipage}\hfill
\begin{minipage}[t]{80mm}\vspace{0pt}\begin{tikzpicture}\begin{axis}[width=72mm,height=42mm,ybar,bar width=1.6pt,xmin=-1,xmax=15,ymin=0,ymax=1.468,axis y line*=left,axis x line*=bottom,xlabel={\scriptsize Halbzüge bis Matt},ylabel={\scriptsize Stellungen (Mrd.)},tick label style={font=\tiny},label style={font=\scriptsize},enlarge x limits=false,ylabel shift=-3pt,xlabel shift=-3pt]\addplot[fill=ink!55,draw=none] coordinates {(0,0.0647) (1,0.2004) (2,0.933) (3,1.005) (4,1.335) (5,0.1292) (6,0.03502) (7,3.988e-06) (8,0.038) (9,2.703e-06) (10,0.01537) (11,1.2e-08) (12,0.001779) (14,2.381e-06)};\draw[densely dotted] (axis cs:12,0) -- (axis cs:12,1.468);\end{axis}\begin{axis}[width=72mm,height=42mm,axis y line*=right,axis x line=none,ymode=log,xmin=-1,xmax=15,ylabel={\scriptsize eindeutig (log)},tick label style={font=\tiny},label style={font=\scriptsize},ylabel shift=-3pt]\addplot[mark=*,mark size=0.7pt,draw=black,thin] coordinates {(0,64700753) (1,116613130) (2,315817905) (3,26227048) (4,2577118) (5,275) (6,1619843) (7,486) (8,599417) (9,1) (10,29305) (12,9)};\end{axis}\end{tikzpicture}\end{minipage}
\vspace{3mm}

\noindent\begin{minipage}[t]{60mm}\vspace{0pt}\sbox\hmboard{\chessboard[setfen=8/6Q1/7k/8/8/b7/7Q/K1n5 b - - 0 1]}\usebox\hmboard\\[1mm]
\makebox[\wd\hmboard]{\sffamily\hspace{3mm}\textbf{h\#6}\hfill(3+3)}\end{minipage}\hfill
\begin{minipage}[t]{52mm}\vspace{0pt}
\textcolor{ink}{\sffamily\textsc{Lösung}}\par\smallskip\small
\begin{tabular}{@{}r@{\enspace}l@{\enspace}l@{}}
\textbf{1.} & \symking{}xg7 & \symking{}b1\\
\textbf{2.} & \symknight{}e2 & \symking{}c2\\
\textbf{3.} & \symknight{}g1 & \symking{}d3\\
\textbf{4.} & \symknight{}h3 & \symking{}e4\\
\textbf{5.} & \symking{}h6 & \symking{}f5\\
\textbf{6.} & \symking{}h5 & \symqueen{}xh3\#\\
\end{tabular}\par\vspace{3mm}
\qrcode[height=12mm]{https://helpman.komtera.lt/?fen=8/6Q1/7k/8/8/b7/7Q/K1n5&moves=6}\\{\sffamily\tiny\color{ink}scannen zum Lösen}
\end{minipage}
\vspace{2mm}

{\raggedright\noindent\textcolor{ink}{\sffamily\textsc{Themen}}\enspace \themetag{Spiegelmatt} \themetag{Switchback} \themetag{ohne Schach}\par\smallskip{\footnotesize\textit{Hinweis:} das Diagramm hat keinen legalen letzten Zug, kann also in keiner Partie entstehen.}\par}

\clearpage
\phantomsection\label{hm:KQQvknn}
\begin{tcolorbox}[enhanced,colback=ink,colframe=ink,arc=0pt,boxrule=0pt,left=4mm,right=4mm,top=2.5mm,bottom=2.5mm]
\color{white}\sffamily{\LARGE\bfseries KQQvknn}\hfill{\Large h\#6}\\[0.5mm]
{\small 6 Steine\quad Schwarz am Zug}\hfill{\small Nr.~665\qquad\href{https://helpman.komtera.lt/?fen=8/6Q1/7k/8/8/8/8/K1nn3Q&moves=6}{\color{white}\bfseries lösen\,\ensuremath{\nearrow}}}
\end{tcolorbox}
\vspace{2mm}
\noindent\begin{minipage}[t]{36mm}\vspace{0pt}\sffamily\small\raggedright
\hmfigure{längste eindeutige}{{\Large h\#6}}\hmfigure{eindeutig bei dieser Länge}{{\Large 40} Stellungen}\hmfigure{längstes Matt}{{\Large h\#7}}\hmfigure{in der Klasse}{15,5 Milliarden Stellungen,\\ 1,94 Milliarden mit Hilfsmatt}
\end{minipage}\hfill
\begin{minipage}[t]{80mm}\vspace{0pt}\begin{tikzpicture}\begin{axis}[width=72mm,height=42mm,ybar,bar width=1.6pt,xmin=-1,xmax=15,ymin=0,ymax=786.7,axis y line*=left,axis x line*=bottom,xlabel={\scriptsize Halbzüge bis Matt},ylabel={\scriptsize Stellungen (Mio.)},tick label style={font=\tiny},label style={font=\scriptsize},enlarge x limits=false,ylabel shift=-3pt,xlabel shift=-3pt]\addplot[fill=ink!55,draw=none] coordinates {(0,35.37) (1,105.8) (2,473.4) (3,502.2) (4,715.1) (5,59.14) (6,18.35) (7,0.003094) (8,20.5) (9,0.000966) (10,8.408) (12,0.9949) (14,0.002262)};\draw[densely dotted] (axis cs:12,0) -- (axis cs:12,786.7);\end{axis}\begin{axis}[width=72mm,height=42mm,axis y line*=right,axis x line=none,ymode=log,xmin=-1,xmax=15,ylabel={\scriptsize eindeutig (log)},tick label style={font=\tiny},label style={font=\scriptsize},ylabel shift=-3pt]\addplot[mark=*,mark size=0.7pt,draw=black,thin] coordinates {(0,35368746) (1,60158093) (2,155827732) (3,11571574) (4,1371015) (5,111) (6,996485) (7,430) (8,328665) (10,15962) (12,40)};\end{axis}\end{tikzpicture}\end{minipage}
\vspace{3mm}

\noindent\begin{minipage}[t]{60mm}\vspace{0pt}\sbox\hmboard{\chessboard[setfen=8/6Q1/7k/8/8/8/8/K1nn3Q b - - 0 1]}\usebox\hmboard\\[1mm]
\makebox[\wd\hmboard]{\sffamily\hspace{3mm}\textbf{h\#6}\hfill(3+3)}\end{minipage}\hfill
\begin{minipage}[t]{52mm}\vspace{0pt}
\textcolor{ink}{\sffamily\textsc{Lösung}}\par\smallskip\small
\begin{tabular}{@{}r@{\enspace}l@{\enspace}l@{}}
\textbf{1.} & \symking{}xg7 & \symking{}b1\\
\textbf{2.} & \symknight{}e2 & \symking{}c2\\
\textbf{3.} & \symknight{}g1 & \symking{}d3\\
\textbf{4.} & \symknight{}h3 & \symking{}e4\\
\textbf{5.} & \symking{}h6 & \symking{}f5\\
\textbf{6.} & \symking{}h5 & \symqueen{}xh3\#\\
\end{tabular}\par\vspace{3mm}
\qrcode[height=12mm]{https://helpman.komtera.lt/?fen=8/6Q1/7k/8/8/8/8/K1nn3Q&moves=6}\\{\sffamily\tiny\color{ink}scannen zum Lösen}
\end{minipage}
\vspace{2mm}

{\raggedright\noindent\textcolor{ink}{\sffamily\textsc{Themen}}\enspace \themetag{Spiegelmatt} \themetag{Switchback} \themetag{ohne Schach}\par\smallskip{\footnotesize\textit{Hinweis:} das Diagramm hat keinen legalen letzten Zug, kann also in keiner Partie entstehen.}\par}

\clearpage
\phantomsection\label{hm:KQQvknp}
\begin{tcolorbox}[enhanced,colback=ink,colframe=ink,arc=0pt,boxrule=0pt,left=4mm,right=4mm,top=2.5mm,bottom=2.5mm]
\color{white}\sffamily{\LARGE\bfseries KQQvknp}\hfill{\Large h\#6}\\[0.5mm]
{\small 6 Steine\quad Schwarz am Zug}\hfill{\small Nr.~666\qquad\href{https://helpman.komtera.lt/?fen=8/7Q/5Q1k/8/2p5/8/1n6/K7&moves=6}{\color{white}\bfseries lösen\,\ensuremath{\nearrow}}}
\end{tcolorbox}
\vspace{2mm}
\noindent\begin{minipage}[t]{36mm}\vspace{0pt}\sffamily\small\raggedright
\hmfigure{längste eindeutige}{{\Large h\#6}}\hmfigure{eindeutig bei dieser Länge}{{\Large 50} Stellungen}\hmfigure{längstes Matt}{{\Large h\#7}}\hmfigure{in der Klasse}{45,4 Milliarden Stellungen,\\ 12,1 Milliarden mit Hilfsmatt}
\end{minipage}\hfill
\begin{minipage}[t]{80mm}\vspace{0pt}\begin{tikzpicture}\begin{axis}[width=72mm,height=42mm,ybar,bar width=1.6pt,xmin=-1,xmax=15,ymin=0,ymax=4.804,axis y line*=left,axis x line*=bottom,xlabel={\scriptsize Halbzüge bis Matt},ylabel={\scriptsize Stellungen (Mrd.)},tick label style={font=\tiny},label style={font=\scriptsize},enlarge x limits=false,ylabel shift=-3pt,xlabel shift=-3pt]\addplot[fill=ink!55,draw=none] coordinates {(0,0.2486) (1,0.7454) (2,3.189) (3,3.043) (4,4.367) (5,0.2439) (6,0.1192) (7,1.356e-05) (8,0.1239) (9,1.846e-06) (10,0.04384) (12,0.004677) (14,2.017e-06)};\draw[densely dotted] (axis cs:12,0) -- (axis cs:12,4.804);\end{axis}\begin{axis}[width=72mm,height=42mm,axis y line*=right,axis x line=none,ymode=log,xmin=-1,xmax=15,ylabel={\scriptsize eindeutig (log)},tick label style={font=\tiny},label style={font=\scriptsize},ylabel shift=-3pt]\addplot[mark=*,mark size=0.7pt,draw=black,thin] coordinates {(0,248633838) (1,406506411) (2,1035607364) (3,61179350) (4,9248065) (5,909) (6,5736283) (7,1367) (8,1699986) (10,77903) (12,50)};\end{axis}\end{tikzpicture}\end{minipage}
\vspace{3mm}

\noindent\begin{minipage}[t]{60mm}\vspace{0pt}\sbox\hmboard{\chessboard[setfen=8/7Q/5Q1k/8/2p5/8/1n6/K7 b - - 0 1]}\usebox\hmboard\\[1mm]
\makebox[\wd\hmboard]{\sffamily\hspace{3mm}\textbf{h\#6}\hfill(3+3)}\end{minipage}\hfill
\begin{minipage}[t]{52mm}\vspace{0pt}
\textcolor{ink}{\sffamily\textsc{Lösung}}\par\smallskip\small
\begin{tabular}{@{}r@{\enspace}l@{\enspace}l@{}}
\textbf{1.} & \symking{}xh7 & \symking{}b1\\
\textbf{2.} & c3 & \symking{}c2\\
\textbf{3.} & \symknight{}c4 & \symking{}d3\\
\textbf{4.} & \symknight{}e5+ & \symking{}e4\\
\textbf{5.} & \symknight{}g6 & \symking{}f5\\
\textbf{6.} & \symking{}h6 & \symqueen{}xg6\#\\
\end{tabular}\par\vspace{3mm}
\qrcode[height=12mm]{https://helpman.komtera.lt/?fen=8/7Q/5Q1k/8/2p5/8/1n6/K7&moves=6}\\{\sffamily\tiny\color{ink}scannen zum Lösen}
\end{minipage}
\vspace{2mm}

{\raggedright\noindent\textcolor{ink}{\sffamily\textsc{Themen}}\enspace \themetag{Switchback}\par\smallskip{\footnotesize\textit{Hinweis:} das Diagramm hat keinen legalen letzten Zug, kann also in keiner Partie entstehen.}\par}

\clearpage
\phantomsection\label{hm:KQQvkpp}
\begin{tcolorbox}[enhanced,colback=ink,colframe=ink,arc=0pt,boxrule=0pt,left=4mm,right=4mm,top=2.5mm,bottom=2.5mm]
\color{white}\sffamily{\LARGE\bfseries KQQvkpp}\hfill{\Large h\#6}\\[0.5mm]
{\small 6 Steine\quad Schwarz am Zug}\hfill{\small Nr.~667\qquad\href{https://helpman.komtera.lt/?fen=8/7Q/5Q1k/3p4/p7/8/8/K7&moves=6}{\color{white}\bfseries lösen\,\ensuremath{\nearrow}}}
\end{tcolorbox}
\vspace{2mm}
\noindent\begin{minipage}[t]{36mm}\vspace{0pt}\sffamily\small\raggedright
\hmfigure{längste eindeutige}{{\Large h\#6}}\hmfigure{eindeutig bei dieser Länge}{{\Large 23} Stellungen}\hmfigure{längstes Matt}{{\Large h\#7}}\hmfigure{in der Klasse}{34,1 Milliarden Stellungen,\\ 4,71 Milliarden mit Hilfsmatt}
\end{minipage}\hfill
\begin{minipage}[t]{80mm}\vspace{0pt}\begin{tikzpicture}\begin{axis}[width=72mm,height=42mm,ybar,bar width=1.6pt,xmin=-1,xmax=15,ymin=0,ymax=1.818,axis y line*=left,axis x line*=bottom,xlabel={\scriptsize Halbzüge bis Matt},ylabel={\scriptsize Stellungen (Mrd.)},tick label style={font=\tiny},label style={font=\scriptsize},enlarge x limits=false,ylabel shift=-3pt,xlabel shift=-3pt]\addplot[fill=ink!55,draw=none] coordinates {(0,0.1095) (1,0.3262) (2,1.327) (3,1.136) (4,1.653) (5,0.05309) (6,0.04872) (7,1.518e-06) (8,0.0452) (9,3.4e-08) (10,0.01286) (12,0.001131) (14,1.239e-06)};\draw[densely dotted] (axis cs:12,0) -- (axis cs:12,1.818);\end{axis}\begin{axis}[width=72mm,height=42mm,axis y line*=right,axis x line=none,ymode=log,xmin=-1,xmax=15,ylabel={\scriptsize eindeutig (log)},tick label style={font=\tiny},label style={font=\scriptsize},ylabel shift=-3pt]\addplot[mark=*,mark size=0.7pt,draw=black,thin] coordinates {(0,109516942) (1,169739277) (2,435059155) (3,18540305) (4,3897613) (5,228) (6,1977234) (7,38) (8,468088) (10,8405) (12,23)};\end{axis}\end{tikzpicture}\end{minipage}
\vspace{3mm}

\noindent\begin{minipage}[t]{60mm}\vspace{0pt}\sbox\hmboard{\chessboard[setfen=8/7Q/5Q1k/3p4/p7/8/8/K7 b - - 0 1]}\usebox\hmboard\\[1mm]
\makebox[\wd\hmboard]{\sffamily\hspace{3mm}\textbf{h\#6}\hfill(3+3)}\end{minipage}\hfill
\begin{minipage}[t]{52mm}\vspace{0pt}
\textcolor{ink}{\sffamily\textsc{Lösung}}\par\smallskip\small
\begin{tabular}{@{}r@{\enspace}l@{\enspace}l@{}}
\textbf{1.} & \symking{}xh7 & \symking{}b2\\
\textbf{2.} & a3+ & \symking{}c3\\
\textbf{3.} & a2 & \symking{}d4\\
\textbf{4.} & a1=\symrook{} & \symking{}e5\\
\textbf{5.} & \symrook{}a8 & \symking{}f5\\
\textbf{6.} & \symrook{}h8 & \symqueen{}g6\#\\
\end{tabular}\par\vspace{3mm}
\qrcode[height=12mm]{https://helpman.komtera.lt/?fen=8/7Q/5Q1k/3p4/p7/8/8/K7&moves=6}\\{\sffamily\tiny\color{ink}scannen zum Lösen}
\end{minipage}
\vspace{2mm}

{\raggedright\noindent\textcolor{ink}{\sffamily\textsc{Themen}}\enspace \themetag{reines Matt} \themetag{Mustermatt} \themetag{Umwandlung} \themetag{Unterverwandlung} \themetag{Selbstblock} \themetag{Kniest}\par\smallskip{\footnotesize\textit{Hinweis:} das Diagramm hat keinen legalen letzten Zug, kann also in keiner Partie entstehen.}\par}

\clearpage
\phantomsection\label{hm:KQQvkqn}
\begin{tcolorbox}[enhanced,colback=ink,colframe=ink,arc=0pt,boxrule=0pt,left=4mm,right=4mm,top=2.5mm,bottom=2.5mm]
\color{white}\sffamily{\LARGE\bfseries KQQvkqn}\hfill{\Large h\#6}\\[0.5mm]
{\small 6 Steine\quad Schwarz am Zug}\hfill{\small Nr.~668\qquad\href{https://helpman.komtera.lt/?fen=8/8/7Q/7k/1q6/4n3/6Q1/K7&moves=6}{\color{white}\bfseries lösen\,\ensuremath{\nearrow}}}
\end{tcolorbox}
\vspace{2mm}
\noindent\begin{minipage}[t]{36mm}\vspace{0pt}\sffamily\small\raggedright
\hmfigure{längste eindeutige}{{\Large h\#6}}\hmfigure{eindeutig bei dieser Länge}{{\Large 75} Stellungen}\hmfigure{längstes Matt}{{\Large h\#7}}\hmfigure{in der Klasse}{15,5 Milliarden Stellungen,\\ 3,21 Milliarden mit Hilfsmatt}
\end{minipage}\hfill
\begin{minipage}[t]{80mm}\vspace{0pt}\begin{tikzpicture}\begin{axis}[width=72mm,height=42mm,ybar,bar width=1.6pt,xmin=-1,xmax=15,ymin=0,ymax=1.192,axis y line*=left,axis x line*=bottom,xlabel={\scriptsize Halbzüge bis Matt},ylabel={\scriptsize Stellungen (Mrd.)},tick label style={font=\tiny},label style={font=\scriptsize},enlarge x limits=false,ylabel shift=-3pt,xlabel shift=-3pt]\addplot[fill=ink!55,draw=none] coordinates {(0,0.04135) (1,0.1316) (2,0.6784) (3,0.9611) (4,1.083) (5,0.2401) (6,0.02734) (7,1.413e-05) (8,0.03058) (9,7.225e-06) (10,0.01409) (11,9e-08) (12,0.001719) (13,3e-09) (14,3.271e-05)};\draw[densely dotted] (axis cs:12,0) -- (axis cs:12,1.192);\end{axis}\begin{axis}[width=72mm,height=42mm,axis y line*=right,axis x line=none,ymode=log,xmin=-1,xmax=15,ylabel={\scriptsize eindeutig (log)},tick label style={font=\tiny},label style={font=\scriptsize},ylabel shift=-3pt]\addplot[mark=*,mark size=0.7pt,draw=black,thin] coordinates {(0,41347674) (1,84085811) (2,248415321) (3,43384361) (4,2258247) (5,1124) (6,1080271) (7,794) (8,517129) (9,2) (10,31344) (12,75)};\end{axis}\end{tikzpicture}\end{minipage}
\vspace{3mm}

\noindent\begin{minipage}[t]{60mm}\vspace{0pt}\sbox\hmboard{\chessboard[setfen=8/8/7Q/7k/1q6/4n3/6Q1/K7 b - - 0 1]}\usebox\hmboard\\[1mm]
\makebox[\wd\hmboard]{\sffamily\hspace{3mm}\textbf{h\#6}\hfill(3+3)}\end{minipage}\hfill
\begin{minipage}[t]{52mm}\vspace{0pt}
\textcolor{ink}{\sffamily\textsc{Lösung}}\par\smallskip\small
\begin{tabular}{@{}r@{\enspace}l@{\enspace}l@{}}
\textbf{1.} & \symking{}xh6 & \symking{}a2\\
\textbf{2.} & \symqueen{}b3+ & \symking{}xb3\\
\textbf{3.} & \symknight{}g4 & \symking{}c4\\
\textbf{4.} & \symking{}g7 & \symking{}d5\\
\textbf{5.} & \symking{}f8 & \symking{}e6\\
\textbf{6.} & \symking{}e8 & \symqueen{}a8\#\\
\end{tabular}\par\vspace{3mm}
\qrcode[height=12mm]{https://helpman.komtera.lt/?fen=8/8/7Q/7k/1q6/4n3/6Q1/K7&moves=6}\\{\sffamily\tiny\color{ink}scannen zum Lösen}
\end{minipage}
\vspace{2mm}

{\raggedright\noindent\textcolor{ink}{\sffamily\textsc{Themen}}\enspace \themetag{reines Matt} \themetag{Mustermatt} \themetag{Spiegelmatt}\par\smallskip{\footnotesize\textit{Hinweis:} die Partei am Zug steht im Diagramm im Schach; die Lösung beginnt mit einem Schlag.}\par}

\clearpage
\phantomsection\label{hm:KQQvkqp}
\begin{tcolorbox}[enhanced,colback=ink,colframe=ink,arc=0pt,boxrule=0pt,left=4mm,right=4mm,top=2.5mm,bottom=2.5mm]
\color{white}\sffamily{\LARGE\bfseries KQQvkqp}\hfill{\Large h\#6}\\[0.5mm]
{\small 6 Steine\quad Schwarz am Zug}\hfill{\small Nr.~669\qquad\href{https://helpman.komtera.lt/?fen=8/KQ4q1/5p2/8/8/6Q1/8/7k&moves=6}{\color{white}\bfseries lösen\,\ensuremath{\nearrow}}}
\end{tcolorbox}
\vspace{2mm}
\noindent\begin{minipage}[t]{36mm}\vspace{0pt}\sffamily\small\raggedright
\hmfigure{längste eindeutige}{{\Large h\#6}}\hmfigure{eindeutig bei dieser Länge}{{\Large 47} Stellungen}\hmfigure{längstes Matt}{{\Large h\#7}}\hmfigure{in der Klasse}{45,4 Milliarden Stellungen,\\ 9,99 Milliarden mit Hilfsmatt}
\end{minipage}\hfill
\begin{minipage}[t]{80mm}\vspace{0pt}\begin{tikzpicture}\begin{axis}[width=72mm,height=42mm,ybar,bar width=1.6pt,xmin=-1,xmax=15,ymin=0,ymax=3.641,axis y line*=left,axis x line*=bottom,xlabel={\scriptsize Halbzüge bis Matt},ylabel={\scriptsize Stellungen (Mrd.)},tick label style={font=\tiny},label style={font=\scriptsize},enlarge x limits=false,ylabel shift=-3pt,xlabel shift=-3pt]\addplot[fill=ink!55,draw=none] coordinates {(0,0.1419) (1,0.4605) (2,2.286) (3,2.957) (4,3.31) (5,0.6121) (6,0.09789) (7,2.06e-05) (8,0.0922) (9,3.772e-06) (10,0.03019) (11,1.048e-06) (12,0.002111) (13,3.2e-08) (14,1.497e-05)};\draw[densely dotted] (axis cs:12,0) -- (axis cs:12,3.641);\end{axis}\begin{axis}[width=72mm,height=42mm,axis y line*=right,axis x line=none,ymode=log,xmin=-1,xmax=15,ylabel={\scriptsize eindeutig (log)},tick label style={font=\tiny},label style={font=\scriptsize},ylabel shift=-3pt]\addplot[mark=*,mark size=0.7pt,draw=black,thin] coordinates {(0,141919657) (1,285779047) (2,820949355) (3,131252889) (4,8216545) (5,4638) (6,3147310) (7,473) (8,1153046) (9,93) (10,48651) (12,47)};\end{axis}\end{tikzpicture}\end{minipage}
\vspace{3mm}

\noindent\begin{minipage}[t]{60mm}\vspace{0pt}\sbox\hmboard{\chessboard[setfen=8/KQ4q1/5p2/8/8/6Q1/8/7k b - - 0 1]}\usebox\hmboard\\[1mm]
\makebox[\wd\hmboard]{\sffamily\hspace{3mm}\textbf{h\#6}\hfill(3+3)}\end{minipage}\hfill
\begin{minipage}[t]{52mm}\vspace{0pt}
\textcolor{ink}{\sffamily\textsc{Lösung}}\par\smallskip\small
\begin{tabular}{@{}r@{\enspace}l@{\enspace}l@{}}
\textbf{1.} & \symqueen{}xb7+ & \symking{}xb7\\
\textbf{2.} & f5 & \symking{}c6\\
\textbf{3.} & f4 & \symking{}d5\\
\textbf{4.} & f3 & \symking{}e4\\
\textbf{5.} & f2 & \symking{}f3\\
\textbf{6.} & f1=\symknight{} & \symqueen{}g2\#\\
\end{tabular}\par\vspace{3mm}
\qrcode[height=12mm]{https://helpman.komtera.lt/?fen=8/KQ4q1/5p2/8/8/6Q1/8/7k&moves=6}\\{\sffamily\tiny\color{ink}scannen zum Lösen}
\end{minipage}
\vspace{2mm}

{\raggedright\noindent\textcolor{ink}{\sffamily\textsc{Themen}}\enspace \themetag{reines Matt} \themetag{Mustermatt} \themetag{Umwandlung} \themetag{Unterverwandlung} \themetag{Umnov}\par\smallskip{\footnotesize\textit{Hinweis:} die Partei am Zug steht im Diagramm im Schach; die Lösung beginnt mit einem Schlag.}\par}

\clearpage
\phantomsection\label{hm:KQQvkqq}
\begin{tcolorbox}[enhanced,colback=ink,colframe=ink,arc=0pt,boxrule=0pt,left=4mm,right=4mm,top=2.5mm,bottom=2.5mm]
\color{white}\sffamily{\LARGE\bfseries KQQvkqq}\hfill{\Large h\#6}\\[0.5mm]
{\small 6 Steine\quad Schwarz am Zug}\hfill{\small Nr.~670\qquad\href{https://helpman.komtera.lt/?fen=8/7Q/7k/8/8/1q6/8/K1Qq4&moves=6}{\color{white}\bfseries lösen\,\ensuremath{\nearrow}}}
\end{tcolorbox}
\vspace{2mm}
\noindent\begin{minipage}[t]{36mm}\vspace{0pt}\sffamily\small\raggedright
\hmfigure{längste eindeutige}{{\Large h\#6}}\hmfigure{eindeutig bei dieser Länge}{{\Large 2} Stellungen}\hmfigure{längstes Matt}{{\Large h\#8}}\hmfigure{in der Klasse}{15,5 Milliarden Stellungen,\\ 1,32 Milliarden mit Hilfsmatt}
\end{minipage}\hfill
\begin{minipage}[t]{80mm}\vspace{0pt}\begin{tikzpicture}\begin{axis}[width=72mm,height=42mm,ybar,bar width=1.6pt,xmin=-1,xmax=17,ymin=0,ymax=485.3,axis y line*=left,axis x line*=bottom,xlabel={\scriptsize Halbzüge bis Matt},ylabel={\scriptsize Stellungen (Mio.)},tick label style={font=\tiny},label style={font=\scriptsize},enlarge x limits=false,ylabel shift=-3pt,xlabel shift=-3pt]\addplot[fill=ink!55,draw=none] coordinates {(0,12.18) (1,39.12) (2,223.7) (3,441.1) (4,404.5) (5,174.6) (6,11.8) (7,0.04071) (8,10.53) (9,0.01015) (10,4.138) (11,0.002692) (12,0.121) (13,0.000162) (14,0.000216) (15,2e-06) (16,2.2e-05)};\draw[densely dotted] (axis cs:12,0) -- (axis cs:12,485.3);\end{axis}\begin{axis}[width=72mm,height=42mm,axis y line*=right,axis x line=none,ymode=log,xmin=-1,xmax=17,ylabel={\scriptsize eindeutig (log)},tick label style={font=\tiny},label style={font=\scriptsize},ylabel shift=-3pt]\addplot[mark=*,mark size=0.7pt,draw=black,thin] coordinates {(0,12175059) (1,26893437) (2,86727648) (3,31235349) (4,985514) (5,8057) (6,397993) (7,693) (8,167616) (9,8) (10,5623) (12,2)};\end{axis}\end{tikzpicture}\end{minipage}
\vspace{3mm}

\noindent\begin{minipage}[t]{60mm}\vspace{0pt}\sbox\hmboard{\chessboard[setfen=8/7Q/7k/8/8/1q6/8/K1Qq4 b - - 0 1]}\usebox\hmboard\\[1mm]
\makebox[\wd\hmboard]{\sffamily\hspace{3mm}\textbf{h\#6}\hfill(3+3)}\end{minipage}\hfill
\begin{minipage}[t]{52mm}\vspace{0pt}
\textcolor{ink}{\sffamily\textsc{Lösung}}\par\smallskip\small
\begin{tabular}{@{}r@{\enspace}l@{\enspace}l@{}}
\textbf{1.} & \symking{}xh7 & \symqueen{}b1+\\
\textbf{2.} & \symqueen{}dd3 & \symqueen{}e1\\
\textbf{3.} & \symqueen{}f7 & \symking{}b2\\
\textbf{4.} & \symqueen{}dg6 & \symking{}c3\\
\textbf{5.} & \symking{}g7 & \symking{}d4\\
\textbf{6.} & \symking{}f6 & \symqueen{}e5\#\\
\end{tabular}\par\vspace{3mm}
\qrcode[height=12mm]{https://helpman.komtera.lt/?fen=8/7Q/7k/8/8/1q6/8/K1Qq4&moves=6}\\{\sffamily\tiny\color{ink}scannen zum Lösen}
\end{minipage}
\vspace{2mm}

{\raggedright\noindent\textcolor{ink}{\sffamily\textsc{Themen}}\enspace \themetag{reines Matt} \themetag{Mustermatt} \themetag{Idealmatt} \themetag{Selbstblock}\par\smallskip{\footnotesize\textit{Hinweis:} das Diagramm hat keinen legalen letzten Zug, kann also in keiner Partie entstehen.}\par}

\clearpage
\phantomsection\label{hm:KQRBBvk}
\begin{tcolorbox}[enhanced,colback=ink,colframe=ink,arc=0pt,boxrule=0pt,left=4mm,right=4mm,top=2.5mm,bottom=2.5mm]
\color{white}\sffamily{\LARGE\bfseries KQRBBvk}\hfill{\Large h\#6}\\[0.5mm]
{\small 6 Steine\quad Schwarz am Zug}\hfill{\small Nr.~671\qquad\href{https://helpman.komtera.lt/?fen=8/7Q/6k1/8/7B/6R1/7B/K7&moves=6}{\color{white}\bfseries lösen\,\ensuremath{\nearrow}}}
\end{tcolorbox}
\vspace{2mm}
\noindent\begin{minipage}[t]{36mm}\vspace{0pt}\sffamily\small\raggedright
\hmfigure{längste eindeutige}{{\Large h\#6}}\hmfigure{eindeutig bei dieser Länge}{{\Large 41} Stellungen}\hmfigure{längstes Matt}{{\Large h\#8}}\hmfigure{in der Klasse}{15,5 Milliarden Stellungen,\\ 4,16 Milliarden mit Hilfsmatt}
\end{minipage}\hfill
\begin{minipage}[t]{80mm}\vspace{0pt}\begin{tikzpicture}\begin{axis}[width=72mm,height=42mm,ybar,bar width=1.6pt,xmin=-1,xmax=17,ymin=0,ymax=1.884,axis y line*=left,axis x line*=bottom,xlabel={\scriptsize Halbzüge bis Matt},ylabel={\scriptsize Stellungen (Mrd.)},tick label style={font=\tiny},label style={font=\scriptsize},enlarge x limits=false,ylabel shift=-3pt,xlabel shift=-3pt]\addplot[fill=ink!55,draw=none] coordinates {(0,0.1127) (1,0.2919) (2,1.067) (3,0.8684) (4,1.713) (5,0.04417) (6,0.04461) (7,1e-09) (8,0.01372) (10,0.002645) (12,0.0001843) (14,3.149e-06) (16,3e-09)};\draw[densely dotted] (axis cs:12,0) -- (axis cs:12,1.884);\end{axis}\begin{axis}[width=72mm,height=42mm,axis y line*=right,axis x line=none,ymode=log,xmin=-1,xmax=17,ylabel={\scriptsize eindeutig (log)},tick label style={font=\tiny},label style={font=\scriptsize},ylabel shift=-3pt]\addplot[mark=*,mark size=0.7pt,draw=black,thin] coordinates {(0,112718393) (1,162281766) (2,465686042) (3,39589159) (4,16449252) (5,1) (6,2928736) (8,109722) (10,2936) (12,41)};\end{axis}\end{tikzpicture}\end{minipage}
\vspace{3mm}

\noindent\begin{minipage}[t]{60mm}\vspace{0pt}\sbox\hmboard{\chessboard[setfen=8/7Q/6k1/8/7B/6R1/7B/K7 b - - 0 1]}\usebox\hmboard\\[1mm]
\makebox[\wd\hmboard]{\sffamily\hspace{3mm}\textbf{h\#6}\hfill(5+1)}\end{minipage}\hfill
\begin{minipage}[t]{52mm}\vspace{0pt}
\textcolor{ink}{\sffamily\textsc{Lösung}}\par\smallskip\small
\begin{tabular}{@{}r@{\enspace}l@{\enspace}l@{}}
\textbf{1.} & \symking{}xh7 & \symking{}b2\\
\textbf{2.} & \symking{}h6 & \symking{}c3\\
\textbf{3.} & \symking{}h5 & \symking{}d4\\
\textbf{4.} & \symking{}xh4 & \symking{}e5\\
\textbf{5.} & \symking{}h5 & \symking{}f6\\
\textbf{6.} & \symking{}h6 & \symrook{}h3\#\\
\end{tabular}\par\vspace{3mm}
\qrcode[height=12mm]{https://helpman.komtera.lt/?fen=8/7Q/6k1/8/7B/6R1/7B/K7&moves=6}\\{\sffamily\tiny\color{ink}scannen zum Lösen}
\end{minipage}
\vspace{2mm}

{\raggedright\noindent\textcolor{ink}{\sffamily\textsc{Themen}}\enspace \themetag{reines Matt} \themetag{Spiegelmatt} \themetag{Switchback} \themetag{ein Stein zieht} \themetag{ohne Schach}\par\smallskip{\footnotesize\textit{Hinweis:} das Diagramm hat keinen legalen letzten Zug, kann also in keiner Partie entstehen.}\par}

\clearpage
\phantomsection\label{hm:KQRBPvk}
\begin{tcolorbox}[enhanced,colback=ink,colframe=ink,arc=0pt,boxrule=0pt,left=4mm,right=4mm,top=2.5mm,bottom=2.5mm]
\color{white}\sffamily{\LARGE\bfseries KQRBPvk}\hfill{\Large h\#6}\\[0.5mm]
{\small 6 Steine\quad Schwarz am Zug}\hfill{\small Nr.~672\qquad\href{https://helpman.komtera.lt/?fen=8/8/6kQ/6RB/8/8/1P6/K7&moves=6}{\color{white}\bfseries lösen\,\ensuremath{\nearrow}}}
\end{tcolorbox}
\vspace{2mm}
\noindent\begin{minipage}[t]{36mm}\vspace{0pt}\sffamily\small\raggedright
\hmfigure{längste eindeutige}{{\Large h\#6}}\hmfigure{eindeutig bei dieser Länge}{{\Large 321} Stellungen}\hmfigure{längstes Matt}{{\Large h\#7}}\hmfigure{in der Klasse}{45,4 Milliarden Stellungen,\\ 26 Milliarden mit Hilfsmatt}
\end{minipage}\hfill
\begin{minipage}[t]{80mm}\vspace{0pt}\begin{tikzpicture}\begin{axis}[width=72mm,height=42mm,ybar,bar width=1.6pt,xmin=-1,xmax=15,ymin=0,ymax=13.09,axis y line*=left,axis x line*=bottom,xlabel={\scriptsize Halbzüge bis Matt},ylabel={\scriptsize Stellungen (Mrd.)},tick label style={font=\tiny},label style={font=\scriptsize},enlarge x limits=false,ylabel shift=-3pt,xlabel shift=-3pt]\addplot[fill=ink!55,draw=none] coordinates {(0,0.4501) (1,1.453) (2,4.984) (3,6.102) (4,11.9) (5,0.6298) (6,0.3619) (7,9.39e-06) (8,0.1125) (10,0.02972) (12,0.001674) (14,6.939e-06)};\draw[densely dotted] (axis cs:12,0) -- (axis cs:12,13.09);\end{axis}\begin{axis}[width=72mm,height=42mm,axis y line*=right,axis x line=none,ymode=log,xmin=-1,xmax=15,ylabel={\scriptsize eindeutig (log)},tick label style={font=\tiny},label style={font=\scriptsize},ylabel shift=-3pt]\addplot[mark=*,mark size=0.7pt,draw=black,thin] coordinates {(0,450062636) (1,872986528) (2,2419882935) (3,499008508) (4,158248550) (5,2485) (6,15740690) (8,1084391) (10,43595) (12,321)};\end{axis}\end{tikzpicture}\end{minipage}
\vspace{3mm}

\noindent\begin{minipage}[t]{60mm}\vspace{0pt}\sbox\hmboard{\chessboard[setfen=8/8/6kQ/6RB/8/8/1P6/K7 b - - 0 1]}\usebox\hmboard\\[1mm]
\makebox[\wd\hmboard]{\sffamily\hspace{3mm}\textbf{h\#6}\hfill(5+1)}\end{minipage}\hfill
\begin{minipage}[t]{52mm}\vspace{0pt}
\textcolor{ink}{\sffamily\textsc{Lösung}}\par\smallskip\small
\begin{tabular}{@{}r@{\enspace}l@{\enspace}l@{}}
\textbf{1.} & \symking{}xh6 & b4\\
\textbf{2.} & \symking{}xg5 & b5\\
\textbf{3.} & \symking{}f4 & b6\\
\textbf{4.} & \symking{}e3 & b7\\
\textbf{5.} & \symking{}d2 & b8=\symqueen{}\\
\textbf{6.} & \symking{}c1 & \symqueen{}b2\#\\
\end{tabular}\par\vspace{3mm}
\qrcode[height=12mm]{https://helpman.komtera.lt/?fen=8/8/6kQ/6RB/8/8/1P6/K7&moves=6}\\{\sffamily\tiny\color{ink}scannen zum Lösen}
\end{minipage}
\vspace{2mm}

{\raggedright\noindent\textcolor{ink}{\sffamily\textsc{Themen}}\enspace \themetag{Umwandlung} \themetag{Excelsior} \themetag{Rundlauf} \themetag{ein Stein zieht} \themetag{Phönix} \themetag{ohne Schach}\par\smallskip{\footnotesize\textit{Hinweis:} das Diagramm hat keinen legalen letzten Zug, kann also in keiner Partie entstehen.}\par}

\clearpage
\phantomsection\label{hm:KQRBvkb}
\begin{tcolorbox}[enhanced,colback=ink,colframe=ink,arc=0pt,boxrule=0pt,left=4mm,right=4mm,top=2.5mm,bottom=2.5mm]
\color{white}\sffamily{\LARGE\bfseries KQRBvkb}\hfill{\Large h\#6}\\[0.5mm]
{\small 6 Steine\quad Schwarz am Zug}\hfill{\small Nr.~673\qquad\href{https://helpman.komtera.lt/?fen=8/4B1b1/6Q1/4R2k/8/8/8/K7&moves=6}{\color{white}\bfseries lösen\,\ensuremath{\nearrow}}}
\end{tcolorbox}
\vspace{2mm}
\noindent\begin{minipage}[t]{36mm}\vspace{0pt}\sffamily\small\raggedright
\hmfigure{längste eindeutige}{{\Large h\#6}}\hmfigure{eindeutig bei dieser Länge}{{\Large 149} Stellungen}\hmfigure{längstes Matt}{{\Large h\#7}}\hmfigure{in der Klasse}{15,5 Milliarden Stellungen,\\ 8,06 Milliarden mit Hilfsmatt}
\end{minipage}\hfill
\begin{minipage}[t]{80mm}\vspace{0pt}\begin{tikzpicture}\begin{axis}[width=72mm,height=42mm,ybar,bar width=1.6pt,xmin=-1,xmax=15,ymin=0,ymax=3.888,axis y line*=left,axis x line*=bottom,xlabel={\scriptsize Halbzüge bis Matt},ylabel={\scriptsize Stellungen (Mrd.)},tick label style={font=\tiny},label style={font=\scriptsize},enlarge x limits=false,ylabel shift=-3pt,xlabel shift=-3pt]\addplot[fill=ink!55,draw=none] coordinates {(0,0.09797) (1,0.3495) (2,1.512) (3,2.122) (4,3.535) (5,0.3212) (6,0.0747) (7,3.247e-05) (8,0.03488) (9,5.03e-07) (10,0.01198) (11,4.5e-08) (12,0.00168) (13,1e-08) (14,3.394e-05)};\draw[densely dotted] (axis cs:12,0) -- (axis cs:12,3.888);\end{axis}\begin{axis}[width=72mm,height=42mm,axis y line*=right,axis x line=none,ymode=log,xmin=-1,xmax=15,ylabel={\scriptsize eindeutig (log)},tick label style={font=\tiny},label style={font=\scriptsize},ylabel shift=-3pt]\addplot[mark=*,mark size=0.7pt,draw=black,thin] coordinates {(0,97969709) (1,236402228) (2,628456614) (3,120888339) (4,19383219) (5,4310) (6,3821107) (7,294) (8,241528) (10,24636) (12,149)};\end{axis}\end{tikzpicture}\end{minipage}
\vspace{3mm}

\noindent\begin{minipage}[t]{60mm}\vspace{0pt}\sbox\hmboard{\chessboard[setfen=8/4B1b1/6Q1/4R2k/8/8/8/K7 b - - 0 1]}\usebox\hmboard\\[1mm]
\makebox[\wd\hmboard]{\sffamily\hspace{3mm}\textbf{h\#6}\hfill(4+2)}\end{minipage}\hfill
\begin{minipage}[t]{52mm}\vspace{0pt}
\textcolor{ink}{\sffamily\textsc{Lösung}}\par\smallskip\small
\begin{tabular}{@{}r@{\enspace}l@{\enspace}l@{}}
\textbf{1.} & \symking{}xg6 & \symbishop{}g5\\
\textbf{2.} & \symking{}h5 & \symking{}b1\\
\textbf{3.} & \symking{}g4 & \symrook{}e8\\
\textbf{4.} & \symking{}f3 & \symbishop{}e7\\
\textbf{5.} & \symking{}e2 & \symbishop{}b4+\\
\textbf{6.} & \symking{}d1 & \symrook{}e1\#\\
\end{tabular}\par\vspace{3mm}
\qrcode[height=12mm]{https://helpman.komtera.lt/?fen=8/4B1b1/6Q1/4R2k/8/8/8/K7&moves=6}\\{\sffamily\tiny\color{ink}scannen zum Lösen}
\end{minipage}
\vspace{2mm}

{\raggedright\noindent\textcolor{ink}{\sffamily\textsc{Themen}}\enspace \themetag{Switchback} \themetag{ein Stein zieht} \themetag{Klasinc} \themetag{Inder}\par\smallskip{\footnotesize\textit{Hinweis:} das Diagramm hat keinen legalen letzten Zug, kann also in keiner Partie entstehen.}\par}

\clearpage
\phantomsection\label{hm:KQRBvkn}
\begin{tcolorbox}[enhanced,colback=ink,colframe=ink,arc=0pt,boxrule=0pt,left=4mm,right=4mm,top=2.5mm,bottom=2.5mm]
\color{white}\sffamily{\LARGE\bfseries KQRBvkn}\hfill{\Large h\#6}\\[0.5mm]
{\small 6 Steine\quad Schwarz am Zug}\hfill{\small Nr.~674\qquad\href{https://helpman.komtera.lt/?fen=8/2B5/6Q1/5k2/4R3/3n4/8/K7&moves=6}{\color{white}\bfseries lösen\,\ensuremath{\nearrow}}}
\end{tcolorbox}
\vspace{2mm}
\noindent\begin{minipage}[t]{36mm}\vspace{0pt}\sffamily\small\raggedright
\hmfigure{längste eindeutige}{{\Large h\#6}}\hmfigure{eindeutig bei dieser Länge}{{\Large 28} Stellungen}\hmfigure{längstes Matt}{{\Large h\#7}}\hmfigure{in der Klasse}{15,5 Milliarden Stellungen,\\ 8,34 Milliarden mit Hilfsmatt}
\end{minipage}\hfill
\begin{minipage}[t]{80mm}\vspace{0pt}\begin{tikzpicture}\begin{axis}[width=72mm,height=42mm,ybar,bar width=1.6pt,xmin=-1,xmax=15,ymin=0,ymax=4.121,axis y line*=left,axis x line*=bottom,xlabel={\scriptsize Halbzüge bis Matt},ylabel={\scriptsize Stellungen (Mrd.)},tick label style={font=\tiny},label style={font=\scriptsize},enlarge x limits=false,ylabel shift=-3pt,xlabel shift=-3pt]\addplot[fill=ink!55,draw=none] coordinates {(0,0.1105) (1,0.3807) (2,1.546) (3,2.128) (4,3.746) (5,0.2844) (6,0.07508) (7,8.441e-05) (8,0.04573) (9,9.3e-07) (10,0.01855) (11,2.8e-08) (12,0.001357) (14,2.785e-06)};\draw[densely dotted] (axis cs:12,0) -- (axis cs:12,4.121);\end{axis}\begin{axis}[width=72mm,height=42mm,axis y line*=right,axis x line=none,ymode=log,xmin=-1,xmax=15,ylabel={\scriptsize eindeutig (log)},tick label style={font=\tiny},label style={font=\scriptsize},ylabel shift=-3pt]\addplot[mark=*,mark size=0.7pt,draw=black,thin] coordinates {(0,110539046) (1,249120520) (2,618587078) (3,118959018) (4,22505753) (5,9276) (6,4057717) (7,2380) (8,432033) (10,32941) (12,28)};\end{axis}\end{tikzpicture}\end{minipage}
\vspace{3mm}

\noindent\begin{minipage}[t]{60mm}\vspace{0pt}\sbox\hmboard{\chessboard[setfen=8/2B5/6Q1/5k2/4R3/3n4/8/K7 b - - 0 1]}\usebox\hmboard\\[1mm]
\makebox[\wd\hmboard]{\sffamily\hspace{3mm}\textbf{h\#6}\hfill(4+2)}\end{minipage}\hfill
\begin{minipage}[t]{52mm}\vspace{0pt}
\textcolor{ink}{\sffamily\textsc{Lösung}}\par\smallskip\small
\begin{tabular}{@{}r@{\enspace}l@{\enspace}l@{}}
\textbf{1.} & \symking{}xg6 & \symrook{}a4\\
\textbf{2.} & \symking{}f5 & \symrook{}a3\\
\textbf{3.} & \symking{}e4 & \symking{}a2\\
\textbf{4.} & \symking{}e3 & \symking{}b3\\
\textbf{5.} & \symking{}d2 & \symbishop{}a5+\\
\textbf{6.} & \symking{}c1 & \symrook{}a1\#\\
\end{tabular}\par\vspace{3mm}
\qrcode[height=12mm]{https://helpman.komtera.lt/?fen=8/2B5/6Q1/5k2/4R3/3n4/8/K7&moves=6}\\{\sffamily\tiny\color{ink}scannen zum Lösen}
\end{minipage}
\vspace{2mm}

{\raggedright\noindent\textcolor{ink}{\sffamily\textsc{Themen}}\enspace \themetag{reines Matt} \themetag{Mustermatt} \themetag{Spiegelmatt} \themetag{Switchback} \themetag{ein Stein zieht}\par\smallskip{\footnotesize\textit{Hinweis:} die Partei am Zug steht im Diagramm im Schach; die Lösung beginnt mit einem Schlag.}\par}

\clearpage
\phantomsection\label{hm:KQRBvkp}
\begin{tcolorbox}[enhanced,colback=ink,colframe=ink,arc=0pt,boxrule=0pt,left=4mm,right=4mm,top=2.5mm,bottom=2.5mm]
\color{white}\sffamily{\LARGE\bfseries KQRBvkp}\hfill{\Large h\#6}\\[0.5mm]
{\small 6 Steine\quad Schwarz am Zug}\hfill{\small Nr.~675\qquad\href{https://helpman.komtera.lt/?fen=6B1/8/6kQ/6Rp/8/8/8/K7&moves=6}{\color{white}\bfseries lösen\,\ensuremath{\nearrow}}}
\end{tcolorbox}
\vspace{2mm}
\noindent\begin{minipage}[t]{36mm}\vspace{0pt}\sffamily\small\raggedright
\hmfigure{längste eindeutige}{{\Large h\#6}}\hmfigure{eindeutig bei dieser Länge}{{\Large 600} Stellungen}\hmfigure{längstes Matt}{{\Large h\#14}}\hmfigure{in der Klasse}{45,4 Milliarden Stellungen,\\ 26 Milliarden mit Hilfsmatt}
\end{minipage}\hfill
\begin{minipage}[t]{80mm}\vspace{0pt}\begin{tikzpicture}\begin{axis}[width=72mm,height=42mm,ybar,bar width=1.6pt,xmin=-1,xmax=29,ymin=0,ymax=12.76,axis y line*=left,axis x line*=bottom,xlabel={\scriptsize Halbzüge bis Matt},ylabel={\scriptsize Stellungen (Mrd.)},tick label style={font=\tiny},label style={font=\scriptsize},enlarge x limits=false,ylabel shift=-3pt,xlabel shift=-3pt]\addplot[fill=ink!55,draw=none] coordinates {(0,0.4027) (1,1.364) (2,5.16) (3,6.503) (4,11.6) (5,0.565) (6,0.2469) (7,8.971e-05) (8,0.1335) (9,6e-08) (10,0.04988) (11,6e-09) (12,0.00471) (14,7.363e-05) (16,5e-09) (18,3e-09) (20,3e-09) (22,1.9e-08) (24,1.2e-08) (26,1e-09) (28,4e-09)};\draw[densely dotted] (axis cs:12,0) -- (axis cs:12,12.76);\end{axis}\begin{axis}[width=72mm,height=42mm,axis y line*=right,axis x line=none,ymode=log,xmin=-1,xmax=29,ylabel={\scriptsize eindeutig (log)},tick label style={font=\tiny},label style={font=\scriptsize},ylabel shift=-3pt]\addplot[mark=*,mark size=0.7pt,draw=black,thin] coordinates {(0,402729203) (1,856585401) (2,2279535172) (3,376086788) (4,83796860) (5,12484) (6,14070599) (7,16) (8,1224883) (9,1) (10,54854) (12,600)};\end{axis}\end{tikzpicture}\end{minipage}
\vspace{3mm}

\noindent\begin{minipage}[t]{60mm}\vspace{0pt}\sbox\hmboard{\chessboard[setfen=6B1/8/6kQ/6Rp/8/8/8/K7 b - - 0 1]}\usebox\hmboard\\[1mm]
\makebox[\wd\hmboard]{\sffamily\hspace{3mm}\textbf{h\#6}\hfill(4+2)}\end{minipage}\hfill
\begin{minipage}[t]{52mm}\vspace{0pt}
\textcolor{ink}{\sffamily\textsc{Lösung}}\par\smallskip\small
\begin{tabular}{@{}r@{\enspace}l@{\enspace}l@{}}
\textbf{1.} & \symking{}xh6 & \symking{}b2\\
\textbf{2.} & h4 & \symking{}c3\\
\textbf{3.} & h3 & \symking{}d4\\
\textbf{4.} & h2 & \symking{}e5\\
\textbf{5.} & h1=\symrook{} & \symking{}f6\\
\textbf{6.} & \symrook{}h5 & \symrook{}g6\#\\
\end{tabular}\par\vspace{3mm}
\qrcode[height=12mm]{https://helpman.komtera.lt/?fen=6B1/8/6kQ/6Rp/8/8/8/K7&moves=6}\\{\sffamily\tiny\color{ink}scannen zum Lösen}
\end{minipage}
\vspace{2mm}

{\raggedright\noindent\textcolor{ink}{\sffamily\textsc{Themen}}\enspace \themetag{Umwandlung} \themetag{Unterverwandlung} \themetag{Rundlauf} \themetag{Selbstblock} \themetag{Kniest} \themetag{ohne Schach}\par\smallskip{\footnotesize\textit{Hinweis:} das Diagramm hat keinen legalen letzten Zug, kann also in keiner Partie entstehen.}\par}

\clearpage
\phantomsection\label{hm:KQRPPvk}
\begin{tcolorbox}[enhanced,colback=ink,colframe=ink,arc=0pt,boxrule=0pt,left=4mm,right=4mm,top=2.5mm,bottom=2.5mm]
\color{white}\sffamily{\LARGE\bfseries KQRPPvk}\hfill{\Large h\#6}\\[0.5mm]
{\small 6 Steine\quad Schwarz am Zug}\hfill{\small Nr.~676\qquad\href{https://helpman.komtera.lt/?fen=8/7Q/6Rk/8/8/8/P4P2/K7&moves=6}{\color{white}\bfseries lösen\,\ensuremath{\nearrow}}}
\end{tcolorbox}
\vspace{2mm}
\noindent\begin{minipage}[t]{36mm}\vspace{0pt}\sffamily\small\raggedright
\hmfigure{längste eindeutige}{{\Large h\#6}}\hmfigure{eindeutig bei dieser Länge}{{\Large 113} Stellungen}\hmfigure{längstes Matt}{{\Large h\#7}}\hmfigure{in der Klasse}{34,1 Milliarden Stellungen,\\ 10,1 Milliarden mit Hilfsmatt}
\end{minipage}\hfill
\begin{minipage}[t]{80mm}\vspace{0pt}\begin{tikzpicture}\begin{axis}[width=72mm,height=42mm,ybar,bar width=1.6pt,xmin=-1,xmax=15,ymin=0,ymax=5.301,axis y line*=left,axis x line*=bottom,xlabel={\scriptsize Halbzüge bis Matt},ylabel={\scriptsize Stellungen (Mrd.)},tick label style={font=\tiny},label style={font=\scriptsize},enlarge x limits=false,ylabel shift=-3pt,xlabel shift=-3pt]\addplot[fill=ink!55,draw=none] coordinates {(0,0.1172) (1,0.4353) (2,1.322) (3,2.378) (4,4.819) (5,0.613) (6,0.3837) (7,0.0007797) (8,0.03919) (9,3.2e-08) (10,0.01091) (12,0.0008621) (14,5.975e-06)};\draw[densely dotted] (axis cs:12,0) -- (axis cs:12,5.301);\end{axis}\begin{axis}[width=72mm,height=42mm,axis y line*=right,axis x line=none,ymode=log,xmin=-1,xmax=15,ylabel={\scriptsize eindeutig (log)},tick label style={font=\tiny},label style={font=\scriptsize},ylabel shift=-3pt]\addplot[mark=*,mark size=0.7pt,draw=black,thin] coordinates {(0,117169176) (1,264988649) (2,629031668) (3,286641085) (4,145851971) (5,175243) (6,3723710) (7,1) (8,213356) (10,8669) (12,113)};\end{axis}\end{tikzpicture}\end{minipage}
\vspace{3mm}

\noindent\begin{minipage}[t]{60mm}\vspace{0pt}\sbox\hmboard{\chessboard[setfen=8/7Q/6Rk/8/8/8/P4P2/K7 b - - 0 1]}\usebox\hmboard\\[1mm]
\makebox[\wd\hmboard]{\sffamily\hspace{3mm}\textbf{h\#6}\hfill(5+1)}\end{minipage}\hfill
\begin{minipage}[t]{52mm}\vspace{0pt}
\textcolor{ink}{\sffamily\textsc{Lösung}}\par\smallskip\small
\begin{tabular}{@{}r@{\enspace}l@{\enspace}l@{}}
\textbf{1.} & \symking{}xh7 & f4\\
\textbf{2.} & \symking{}h8 & f5\\
\textbf{3.} & \symking{}h7 & f6\\
\textbf{4.} & \symking{}h8 & f7\\
\textbf{5.} & \symking{}h7 & \symrook{}g7+\\
\textbf{6.} & \symking{}h8 & f8=\symqueen{}\#\\
\end{tabular}\par\vspace{3mm}
\qrcode[height=12mm]{https://helpman.komtera.lt/?fen=8/7Q/6Rk/8/8/8/P4P2/K7&moves=6}\\{\sffamily\tiny\color{ink}scannen zum Lösen}
\end{minipage}
\vspace{2mm}

{\raggedright\noindent\textcolor{ink}{\sffamily\textsc{Themen}}\enspace \themetag{Umwandlung} \themetag{Excelsior} \themetag{Switchback} \themetag{ein Stein zieht} \themetag{Phönix}\par\smallskip{\footnotesize\textit{Hinweis:} das Diagramm hat keinen legalen letzten Zug, kann also in keiner Partie entstehen.}\par}

\clearpage
\phantomsection\label{hm:KQRPvkb}
\begin{tcolorbox}[enhanced,colback=ink,colframe=ink,arc=0pt,boxrule=0pt,left=4mm,right=4mm,top=2.5mm,bottom=2.5mm]
\color{white}\sffamily{\LARGE\bfseries KQRPvkb}\hfill{\Large h\#6}\\[0.5mm]
{\small 6 Steine\quad Schwarz am Zug}\hfill{\small Nr.~677\qquad\href{https://helpman.komtera.lt/?fen=8/7Q/6Rk/8/8/b7/6P1/K7&moves=6}{\color{white}\bfseries lösen\,\ensuremath{\nearrow}}}
\end{tcolorbox}
\vspace{2mm}
\noindent\begin{minipage}[t]{36mm}\vspace{0pt}\sffamily\small\raggedright
\hmfigure{längste eindeutige}{{\Large h\#6}}\hmfigure{eindeutig bei dieser Länge}{{\Large 207} Stellungen}\hmfigure{längstes Matt}{{\Large h\#7}}\hmfigure{in der Klasse}{45,4 Milliarden Stellungen,\\ 25,3 Milliarden mit Hilfsmatt}
\end{minipage}\hfill
\begin{minipage}[t]{80mm}\vspace{0pt}\begin{tikzpicture}\begin{axis}[width=72mm,height=42mm,ybar,bar width=1.6pt,xmin=-1,xmax=15,ymin=0,ymax=13.02,axis y line*=left,axis x line*=bottom,xlabel={\scriptsize Halbzüge bis Matt},ylabel={\scriptsize Stellungen (Mrd.)},tick label style={font=\tiny},label style={font=\scriptsize},enlarge x limits=false,ylabel shift=-3pt,xlabel shift=-3pt]\addplot[fill=ink!55,draw=none] coordinates {(0,0.1993) (1,0.7889) (2,3.01) (3,6.374) (4,11.84) (5,2.244) (6,0.6349) (7,0.004016) (8,0.1078) (9,1.313e-06) (10,0.04725) (11,1.67e-07) (12,0.005211) (13,1e-09) (14,2.314e-05)};\draw[densely dotted] (axis cs:12,0) -- (axis cs:12,13.02);\end{axis}\begin{axis}[width=72mm,height=42mm,axis y line*=right,axis x line=none,ymode=log,xmin=-1,xmax=15,ylabel={\scriptsize eindeutig (log)},tick label style={font=\tiny},label style={font=\scriptsize},ylabel shift=-3pt]\addplot[mark=*,mark size=0.7pt,draw=black,thin] coordinates {(0,199325806) (1,543966431) (2,1193040404) (3,644428951) (4,253872837) (5,825060) (6,6571181) (7,438) (8,886825) (9,14) (10,68221) (12,207)};\end{axis}\end{tikzpicture}\end{minipage}
\vspace{3mm}

\noindent\begin{minipage}[t]{60mm}\vspace{0pt}\sbox\hmboard{\chessboard[setfen=8/7Q/6Rk/8/8/b7/6P1/K7 b - - 0 1]}\usebox\hmboard\\[1mm]
\makebox[\wd\hmboard]{\sffamily\hspace{3mm}\textbf{h\#6}\hfill(4+2)}\end{minipage}\hfill
\begin{minipage}[t]{52mm}\vspace{0pt}
\textcolor{ink}{\sffamily\textsc{Lösung}}\par\smallskip\small
\begin{tabular}{@{}r@{\enspace}l@{\enspace}l@{}}
\textbf{1.} & \symking{}xh7 & g4\\
\textbf{2.} & \symbishop{}e7 & g5\\
\textbf{3.} & \symbishop{}f6+ & gxf6\\
\textbf{4.} & \symking{}h8 & f7\\
\textbf{5.} & \symking{}h7 & \symrook{}g7+\\
\textbf{6.} & \symking{}h8 & f8=\symqueen{}\#\\
\end{tabular}\par\vspace{3mm}
\qrcode[height=12mm]{https://helpman.komtera.lt/?fen=8/7Q/6Rk/8/8/b7/6P1/K7&moves=6}\\{\sffamily\tiny\color{ink}scannen zum Lösen}
\end{minipage}
\vspace{2mm}

{\raggedright\noindent\textcolor{ink}{\sffamily\textsc{Themen}}\enspace \themetag{Umwandlung} \themetag{Excelsior} \themetag{Switchback} \themetag{Phönix}\par\smallskip{\footnotesize\textit{Hinweis:} das Diagramm hat keinen legalen letzten Zug, kann also in keiner Partie entstehen.}\par}

\clearpage
\phantomsection\label{hm:KQRPvkn}
\begin{tcolorbox}[enhanced,colback=ink,colframe=ink,arc=0pt,boxrule=0pt,left=4mm,right=4mm,top=2.5mm,bottom=2.5mm]
\color{white}\sffamily{\LARGE\bfseries KQRPvkn}\hfill{\Large h\#6}\\[0.5mm]
{\small 6 Steine\quad Schwarz am Zug}\hfill{\small Nr.~678\qquad\href{https://helpman.komtera.lt/?fen=7R/6Q1/7k/8/1n6/8/1P6/K7&moves=6}{\color{white}\bfseries lösen\,\ensuremath{\nearrow}}}
\end{tcolorbox}
\vspace{2mm}
\noindent\begin{minipage}[t]{36mm}\vspace{0pt}\sffamily\small\raggedright
\hmfigure{längste eindeutige}{{\Large h\#6}}\hmfigure{eindeutig bei dieser Länge}{{\Large 30} Stellungen}\hmfigure{längstes Matt}{{\Large h\#8}}\hmfigure{in der Klasse}{45,4 Milliarden Stellungen,\\ 26,1 Milliarden mit Hilfsmatt}
\end{minipage}\hfill
\begin{minipage}[t]{80mm}\vspace{0pt}\begin{tikzpicture}\begin{axis}[width=72mm,height=42mm,ybar,bar width=1.6pt,xmin=-1,xmax=17,ymin=0,ymax=13.5,axis y line*=left,axis x line*=bottom,xlabel={\scriptsize Halbzüge bis Matt},ylabel={\scriptsize Stellungen (Mrd.)},tick label style={font=\tiny},label style={font=\scriptsize},enlarge x limits=false,ylabel shift=-3pt,xlabel shift=-3pt]\addplot[fill=ink!55,draw=none] coordinates {(0,0.2331) (1,0.8889) (2,3.124) (3,6.371) (4,12.27) (5,2.145) (6,0.8509) (7,0.00625) (8,0.1133) (9,4.687e-06) (10,0.05021) (11,5.2e-08) (12,0.005776) (14,3.823e-05) (16,2.5e-08)};\draw[densely dotted] (axis cs:12,0) -- (axis cs:12,13.5);\end{axis}\begin{axis}[width=72mm,height=42mm,axis y line*=right,axis x line=none,ymode=log,xmin=-1,xmax=17,ylabel={\scriptsize eindeutig (log)},tick label style={font=\tiny},label style={font=\scriptsize},ylabel shift=-3pt]\addplot[mark=*,mark size=0.7pt,draw=black,thin] coordinates {(0,233070833) (1,585795381) (2,1155817312) (3,620027031) (4,308158365) (5,996107) (6,7499690) (7,7327) (8,861977) (9,13) (10,65313) (12,30)};\end{axis}\end{tikzpicture}\end{minipage}
\vspace{3mm}

\noindent\begin{minipage}[t]{60mm}\vspace{0pt}\sbox\hmboard{\chessboard[setfen=7R/6Q1/7k/8/1n6/8/1P6/K7 b - - 0 1]}\usebox\hmboard\\[1mm]
\makebox[\wd\hmboard]{\sffamily\hspace{3mm}\textbf{h\#6}\hfill(4+2)}\end{minipage}\hfill
\begin{minipage}[t]{52mm}\vspace{0pt}
\textcolor{ink}{\sffamily\textsc{Lösung}}\par\smallskip\small
\begin{tabular}{@{}r@{\enspace}l@{\enspace}l@{}}
\textbf{1.} & \symking{}xg7 & \symking{}b1\\
\textbf{2.} & \symknight{}d5 & \symking{}c2\\
\textbf{3.} & \symknight{}f6 & \symking{}d3\\
\textbf{4.} & \symknight{}h7 & \symking{}e4\\
\textbf{5.} & \symking{}h6 & \symking{}f5\\
\textbf{6.} & \symking{}h5 & \symrook{}xh7\#\\
\end{tabular}\par\vspace{3mm}
\qrcode[height=12mm]{https://helpman.komtera.lt/?fen=7R/6Q1/7k/8/1n6/8/1P6/K7&moves=6}\\{\sffamily\tiny\color{ink}scannen zum Lösen}
\end{minipage}
\vspace{2mm}

{\raggedright\noindent\textcolor{ink}{\sffamily\textsc{Themen}}\enspace \themetag{reines Matt} \themetag{Mustermatt} \themetag{Spiegelmatt} \themetag{Switchback} \themetag{ohne Schach}\par\smallskip{\footnotesize\textit{Hinweis:} das Diagramm hat keinen legalen letzten Zug, kann also in keiner Partie entstehen.}\par}

\clearpage
\phantomsection\label{hm:KQRPvkp}
\begin{tcolorbox}[enhanced,colback=ink,colframe=ink,arc=0pt,boxrule=0pt,left=4mm,right=4mm,top=2.5mm,bottom=2.5mm]
\color{white}\sffamily{\LARGE\bfseries KQRPvkp}\hfill{\Large h\#6}\\[0.5mm]
{\small 6 Steine\quad Schwarz am Zug}\hfill{\small Nr.~679\qquad\href{https://helpman.komtera.lt/?fen=8/8/7Q/5pk1/8/8/6RP/K7&moves=6}{\color{white}\bfseries lösen\,\ensuremath{\nearrow}}}
\end{tcolorbox}
\vspace{2mm}
\noindent\begin{minipage}[t]{36mm}\vspace{0pt}\sffamily\small\raggedright
\hmfigure{längste eindeutige}{{\Large h\#6}}\hmfigure{eindeutig bei dieser Länge}{{\Large 720} Stellungen}\hmfigure{längstes Matt}{{\Large h\#10}}\hmfigure{in der Klasse}{34,1 Milliarden Stellungen,\\ 20,2 Milliarden mit Hilfsmatt}
\end{minipage}\hfill
\begin{minipage}[t]{80mm}\vspace{0pt}\begin{tikzpicture}\begin{axis}[width=72mm,height=42mm,ybar,bar width=1.6pt,xmin=-1,xmax=21,ymin=0,ymax=10.47,axis y line*=left,axis x line*=bottom,xlabel={\scriptsize Halbzüge bis Matt},ylabel={\scriptsize Stellungen (Mrd.)},tick label style={font=\tiny},label style={font=\scriptsize},enlarge x limits=false,ylabel shift=-3pt,xlabel shift=-3pt]\addplot[fill=ink!55,draw=none] coordinates {(0,0.2135) (1,0.8022) (2,2.618) (3,4.939) (4,9.516) (5,1.312) (6,0.6718) (7,0.002159) (8,0.08602) (9,4.033e-06) (10,0.03721) (11,3.29e-07) (12,0.004515) (13,2e-09) (14,5.87e-05) (16,2.2e-08) (18,1.1e-08) (20,1e-09)};\draw[densely dotted] (axis cs:12,0) -- (axis cs:12,10.47);\end{axis}\begin{axis}[width=72mm,height=42mm,axis y line*=right,axis x line=none,ymode=log,xmin=-1,xmax=21,ylabel={\scriptsize eindeutig (log)},tick label style={font=\tiny},label style={font=\scriptsize},ylabel shift=-3pt]\addplot[mark=*,mark size=0.7pt,draw=black,thin] coordinates {(0,213529039) (1,510951113) (2,1083634269) (3,476920990) (4,231644637) (5,525170) (6,6112243) (7,401) (8,540840) (9,31) (10,34108) (12,720)};\end{axis}\end{tikzpicture}\end{minipage}
\vspace{3mm}

\noindent\begin{minipage}[t]{60mm}\vspace{0pt}\sbox\hmboard{\chessboard[setfen=8/8/7Q/5pk1/8/8/6RP/K7 b - - 0 1]}\usebox\hmboard\\[1mm]
\makebox[\wd\hmboard]{\sffamily\hspace{3mm}\textbf{h\#6}\hfill(4+2)}\end{minipage}\hfill
\begin{minipage}[t]{52mm}\vspace{0pt}
\textcolor{ink}{\sffamily\textsc{Lösung}}\par\smallskip\small
\begin{tabular}{@{}r@{\enspace}l@{\enspace}l@{}}
\textbf{1.} & \symking{}xh6 & \symking{}b2\\
\textbf{2.} & f4 & \symking{}c3\\
\textbf{3.} & f3 & \symking{}d4\\
\textbf{4.} & f2 & \symking{}e5\\
\textbf{5.} & f1=\symknight{} & \symking{}f6\\
\textbf{6.} & \symknight{}xh2 & \symrook{}xh2\#\\
\end{tabular}\par\vspace{3mm}
\qrcode[height=12mm]{https://helpman.komtera.lt/?fen=8/8/7Q/5pk1/8/8/6RP/K7&moves=6}\\{\sffamily\tiny\color{ink}scannen zum Lösen}
\end{minipage}
\vspace{2mm}

{\raggedright\noindent\textcolor{ink}{\sffamily\textsc{Themen}}\enspace \themetag{reines Matt} \themetag{Mustermatt} \themetag{Idealmatt} \themetag{Spiegelmatt} \themetag{Umwandlung} \themetag{Unterverwandlung} \themetag{Kniest} \themetag{ohne Schach}\par\smallskip{\footnotesize\textit{Hinweis:} das Diagramm hat keinen legalen letzten Zug, kann also in keiner Partie entstehen.}\par}

\clearpage
\phantomsection\label{hm:KQRPvkq}
\begin{tcolorbox}[enhanced,colback=ink,colframe=ink,arc=0pt,boxrule=0pt,left=4mm,right=4mm,top=2.5mm,bottom=2.5mm]
\color{white}\sffamily{\LARGE\bfseries KQRPvkq}\hfill{\Large h\#6}\\[0.5mm]
{\small 6 Steine\quad Schwarz am Zug}\hfill{\small Nr.~680\qquad\href{https://helpman.komtera.lt/?fen=8/5Q2/4k3/3R4/8/2q5/1P6/K7&moves=6}{\color{white}\bfseries lösen\,\ensuremath{\nearrow}}}
\end{tcolorbox}
\vspace{2mm}
\noindent\begin{minipage}[t]{36mm}\vspace{0pt}\sffamily\small\raggedright
\hmfigure{längste eindeutige}{{\Large h\#6}}\hmfigure{eindeutig bei dieser Länge}{{\Large 99} Stellungen}\hmfigure{längstes Matt}{{\Large h\#8}}\hmfigure{in der Klasse}{45,4 Milliarden Stellungen,\\ 21,7 Milliarden mit Hilfsmatt}
\end{minipage}\hfill
\begin{minipage}[t]{80mm}\vspace{0pt}\begin{tikzpicture}\begin{axis}[width=72mm,height=42mm,ybar,bar width=1.6pt,xmin=-1,xmax=17,ymin=0,ymax=10.3,axis y line*=left,axis x line*=bottom,xlabel={\scriptsize Halbzüge bis Matt},ylabel={\scriptsize Stellungen (Mrd.)},tick label style={font=\tiny},label style={font=\scriptsize},enlarge x limits=false,ylabel shift=-3pt,xlabel shift=-3pt]\addplot[fill=ink!55,draw=none] coordinates {(0,0.117) (1,0.4837) (2,2.216) (3,5.622) (4,9.36) (5,3.284) (6,0.4468) (7,0.01325) (8,0.08943) (9,0.0001361) (10,0.03985) (11,2.171e-05) (12,0.004787) (13,6.92e-07) (14,2.269e-05) (16,2.4e-08)};\draw[densely dotted] (axis cs:12,0) -- (axis cs:12,10.3);\end{axis}\begin{axis}[width=72mm,height=42mm,axis y line*=right,axis x line=none,ymode=log,xmin=-1,xmax=17,ylabel={\scriptsize eindeutig (log)},tick label style={font=\tiny},label style={font=\scriptsize},ylabel shift=-3pt]\addplot[mark=*,mark size=0.7pt,draw=black,thin] coordinates {(0,116973840) (1,364923477) (2,952501756) (3,670291858) (4,181409888) (5,3926852) (6,4634430) (7,9906) (8,825678) (9,459) (10,84900) (11,10) (12,99)};\end{axis}\end{tikzpicture}\end{minipage}
\vspace{3mm}

\noindent\begin{minipage}[t]{60mm}\vspace{0pt}\sbox\hmboard{\chessboard[setfen=8/5Q2/4k3/3R4/8/2q5/1P6/K7 b - - 0 1]}\usebox\hmboard\\[1mm]
\makebox[\wd\hmboard]{\sffamily\hspace{3mm}\textbf{h\#6}\hfill(4+2)}\end{minipage}\hfill
\begin{minipage}[t]{52mm}\vspace{0pt}
\textcolor{ink}{\sffamily\textsc{Lösung}}\par\smallskip\small
\begin{tabular}{@{}r@{\enspace}l@{\enspace}l@{}}
\textbf{1.} & \symking{}xf7 & \symrook{}d2\\
\textbf{2.} & \symking{}e6 & \symrook{}c2\\
\textbf{3.} & \symking{}d5 & \symking{}b1\\
\textbf{4.} & \symking{}c4 & \symking{}c1\\
\textbf{5.} & \symking{}b3 & \symrook{}xc3+\\
\textbf{6.} & \symking{}a2 & \symrook{}a3\#\\
\end{tabular}\par\vspace{3mm}
\qrcode[height=12mm]{https://helpman.komtera.lt/?fen=8/5Q2/4k3/3R4/8/2q5/1P6/K7&moves=6}\\{\sffamily\tiny\color{ink}scannen zum Lösen}
\end{minipage}
\vspace{2mm}

{\raggedright\noindent\textcolor{ink}{\sffamily\textsc{Themen}}\enspace \themetag{reines Matt} \themetag{Mustermatt} \themetag{Idealmatt} \themetag{Switchback} \themetag{ein Stein zieht}\par\smallskip{\footnotesize\textit{Hinweis:} die Partei am Zug steht im Diagramm im Schach; die Lösung beginnt mit einem Schlag.}\par}

\clearpage
\phantomsection\label{hm:KQRPvkr}
\begin{tcolorbox}[enhanced,colback=ink,colframe=ink,arc=0pt,boxrule=0pt,left=4mm,right=4mm,top=2.5mm,bottom=2.5mm]
\color{white}\sffamily{\LARGE\bfseries KQRPvkr}\hfill{\Large h\#6}\\[0.5mm]
{\small 6 Steine\quad Schwarz am Zug}\hfill{\small Nr.~681\qquad\href{https://helpman.komtera.lt/?fen=8/K7/8/8/8/1r6/1P5Q/1R4k1&moves=6}{\color{white}\bfseries lösen\,\ensuremath{\nearrow}}}
\end{tcolorbox}
\vspace{2mm}
\noindent\begin{minipage}[t]{36mm}\vspace{0pt}\sffamily\small\raggedright
\hmfigure{längste eindeutige}{{\Large h\#6}}\hmfigure{eindeutig bei dieser Länge}{{\Large 106} Stellungen}\hmfigure{längstes Matt}{{\Large h\#7}}\hmfigure{in der Klasse}{45,4 Milliarden Stellungen,\\ 24 Milliarden mit Hilfsmatt}
\end{minipage}\hfill
\begin{minipage}[t]{80mm}\vspace{0pt}\begin{tikzpicture}\begin{axis}[width=72mm,height=42mm,ybar,bar width=1.6pt,xmin=-1,xmax=15,ymin=0,ymax=11.94,axis y line*=left,axis x line*=bottom,xlabel={\scriptsize Halbzüge bis Matt},ylabel={\scriptsize Stellungen (Mrd.)},tick label style={font=\tiny},label style={font=\scriptsize},enlarge x limits=false,ylabel shift=-3pt,xlabel shift=-3pt]\addplot[fill=ink!55,draw=none] coordinates {(0,0.1765) (1,0.7165) (2,2.944) (3,6.23) (4,10.86) (5,2.458) (6,0.4533) (7,0.004338) (8,0.1004) (9,6.412e-06) (10,0.0383) (11,5.11e-07) (12,0.003823) (13,4.2e-08) (14,1.604e-05)};\draw[densely dotted] (axis cs:12,0) -- (axis cs:12,11.94);\end{axis}\begin{axis}[width=72mm,height=42mm,axis y line*=right,axis x line=none,ymode=log,xmin=-1,xmax=15,ylabel={\scriptsize eindeutig (log)},tick label style={font=\tiny},label style={font=\scriptsize},ylabel shift=-3pt]\addplot[mark=*,mark size=0.7pt,draw=black,thin] coordinates {(0,176497057) (1,506889642) (2,1209970336) (3,601818839) (4,188203577) (5,1105717) (6,5934014) (7,1008) (8,752535) (9,10) (10,67520) (12,106)};\end{axis}\end{tikzpicture}\end{minipage}
\vspace{3mm}

\noindent\begin{minipage}[t]{60mm}\vspace{0pt}\sbox\hmboard{\chessboard[setfen=8/K7/8/8/8/1r6/1P5Q/1R4k1 b - - 0 1]}\usebox\hmboard\\[1mm]
\makebox[\wd\hmboard]{\sffamily\hspace{3mm}\textbf{h\#6}\hfill(4+2)}\end{minipage}\hfill
\begin{minipage}[t]{52mm}\vspace{0pt}
\textcolor{ink}{\sffamily\textsc{Lösung}}\par\smallskip\small
\begin{tabular}{@{}r@{\enspace}l@{\enspace}l@{}}
\textbf{1.} & \symking{}xh2 & \symking{}a6\\
\textbf{2.} & \symrook{}xb2 & \symrook{}xb2+\\
\textbf{3.} & \symking{}g1 & \symking{}b5\\
\textbf{4.} & \symking{}f1 & \symking{}c4\\
\textbf{5.} & \symking{}e1 & \symking{}d3\\
\textbf{6.} & \symking{}d1 & \symrook{}b1\#\\
\end{tabular}\par\vspace{3mm}
\qrcode[height=12mm]{https://helpman.komtera.lt/?fen=8/K7/8/8/8/1r6/1P5Q/1R4k1&moves=6}\\{\sffamily\tiny\color{ink}scannen zum Lösen}
\end{minipage}
\vspace{2mm}

{\raggedright\noindent\textcolor{ink}{\sffamily\textsc{Themen}}\enspace \themetag{reines Matt} \themetag{Mustermatt} \themetag{Idealmatt} \themetag{Spiegelmatt} \themetag{Switchback}\par\smallskip{\footnotesize\textit{Hinweis:} das Diagramm hat keinen legalen letzten Zug, kann also in keiner Partie entstehen.}\par}

\clearpage
\phantomsection\label{hm:KQRvkbb}
\begin{tcolorbox}[enhanced,colback=ink,colframe=ink,arc=0pt,boxrule=0pt,left=4mm,right=4mm,top=2.5mm,bottom=2.5mm]
\color{white}\sffamily{\LARGE\bfseries KQRvkbb}\hfill{\Large h\#6}\\[0.5mm]
{\small 6 Steine\quad Schwarz am Zug}\hfill{\small Nr.~682\qquad\href{https://helpman.komtera.lt/?fen=8/8/8/5bQ1/5k2/4R3/8/K1b5&moves=6}{\color{white}\bfseries lösen\,\ensuremath{\nearrow}}}
\end{tcolorbox}
\vspace{2mm}
\noindent\begin{minipage}[t]{36mm}\vspace{0pt}\sffamily\small\raggedright
\hmfigure{längste eindeutige}{{\Large h\#6}}\hmfigure{eindeutig bei dieser Länge}{{\Large 92} Stellungen}\hmfigure{längstes Matt}{{\Large h\#7}}\hmfigure{in der Klasse}{15,5 Milliarden Stellungen,\\ 3,89 Milliarden mit Hilfsmatt}
\end{minipage}\hfill
\begin{minipage}[t]{80mm}\vspace{0pt}\begin{tikzpicture}\begin{axis}[width=72mm,height=42mm,ybar,bar width=1.6pt,xmin=-1,xmax=15,ymin=0,ymax=1.939,axis y line*=left,axis x line*=bottom,xlabel={\scriptsize Halbzüge bis Matt},ylabel={\scriptsize Stellungen (Mrd.)},tick label style={font=\tiny},label style={font=\scriptsize},enlarge x limits=false,ylabel shift=-3pt,xlabel shift=-3pt]\addplot[fill=ink!55,draw=none] coordinates {(0,0.02148) (1,0.08734) (2,0.3714) (3,0.9836) (4,1.763) (5,0.5341) (6,0.09211) (7,0.0006946) (8,0.01849) (9,1.456e-06) (10,0.01339) (11,1.39e-07) (12,0.002594) (13,3e-09) (14,3.713e-05)};\draw[densely dotted] (axis cs:12,0) -- (axis cs:12,1.939);\end{axis}\begin{axis}[width=72mm,height=42mm,axis y line*=right,axis x line=none,ymode=log,xmin=-1,xmax=15,ylabel={\scriptsize eindeutig (log)},tick label style={font=\tiny},label style={font=\scriptsize},ylabel shift=-3pt]\addplot[mark=*,mark size=0.7pt,draw=black,thin] coordinates {(0,21477373) (1,65783299) (2,154139690) (3,120205629) (4,42419336) (5,244307) (6,922545) (7,126) (8,229410) (9,81) (10,38163) (12,92)};\end{axis}\end{tikzpicture}\end{minipage}
\vspace{3mm}

\noindent\begin{minipage}[t]{60mm}\vspace{0pt}\sbox\hmboard{\chessboard[setfen=8/8/8/5bQ1/5k2/4R3/8/K1b5 b - - 0 1]}\usebox\hmboard\\[1mm]
\makebox[\wd\hmboard]{\sffamily\hspace{3mm}\textbf{h\#6}\hfill(3+3)}\end{minipage}\hfill
\begin{minipage}[t]{52mm}\vspace{0pt}
\textcolor{ink}{\sffamily\textsc{Lösung}}\par\smallskip\small
\begin{tabular}{@{}r@{\enspace}l@{\enspace}l@{}}
\textbf{1.} & \symking{}xg5 & \symrook{}e5\\
\textbf{2.} & \symking{}f4 & \symrook{}xf5+\\
\textbf{3.} & \symking{}e3 & \symking{}b1\\
\textbf{4.} & \symking{}e2 & \symking{}c2\\
\textbf{5.} & \symking{}e1 & \symking{}d3\\
\textbf{6.} & \symking{}d1 & \symrook{}f1\#\\
\end{tabular}\par\vspace{3mm}
\qrcode[height=12mm]{https://helpman.komtera.lt/?fen=8/8/8/5bQ1/5k2/4R3/8/K1b5&moves=6}\\{\sffamily\tiny\color{ink}scannen zum Lösen}
\end{minipage}
\vspace{2mm}

{\raggedright\noindent\textcolor{ink}{\sffamily\textsc{Themen}}\enspace \themetag{Switchback} \themetag{ein Stein zieht}\par\smallskip{\footnotesize\textit{Hinweis:} die Partei am Zug steht im Diagramm im Schach; die Lösung beginnt mit einem Schlag.}\par}

\clearpage
\phantomsection\label{hm:KQRvkbn}
\begin{tcolorbox}[enhanced,colback=ink,colframe=ink,arc=0pt,boxrule=0pt,left=4mm,right=4mm,top=2.5mm,bottom=2.5mm]
\color{white}\sffamily{\LARGE\bfseries KQRvkbn}\hfill{\Large h\#6}\\[0.5mm]
{\small 6 Steine\quad Schwarz am Zug}\hfill{\small Nr.~683\qquad\href{https://helpman.komtera.lt/?fen=8/8/8/6Q1/2b4k/8/7R/K2n4&moves=6}{\color{white}\bfseries lösen\,\ensuremath{\nearrow}}}
\end{tcolorbox}
\vspace{2mm}
\noindent\begin{minipage}[t]{36mm}\vspace{0pt}\sffamily\small\raggedright
\hmfigure{längste eindeutige}{{\Large h\#6}}\hmfigure{eindeutig bei dieser Länge}{{\Large 141} Stellungen}\hmfigure{längstes Matt}{{\Large h\#7}}\hmfigure{in der Klasse}{15,5 Milliarden Stellungen,\\ 8,06 Milliarden mit Hilfsmatt}
\end{minipage}\hfill
\begin{minipage}[t]{80mm}\vspace{0pt}\begin{tikzpicture}\begin{axis}[width=72mm,height=42mm,ybar,bar width=1.6pt,xmin=-1,xmax=15,ymin=0,ymax=4.095,axis y line*=left,axis x line*=bottom,xlabel={\scriptsize Halbzüge bis Matt},ylabel={\scriptsize Stellungen (Mrd.)},tick label style={font=\tiny},label style={font=\scriptsize},enlarge x limits=false,ylabel shift=-3pt,xlabel shift=-3pt]\addplot[fill=ink!55,draw=none] coordinates {(0,0.05024) (1,0.1986) (2,0.7856) (3,1.994) (4,3.723) (5,1.018) (6,0.2133) (7,0.001397) (8,0.03931) (9,7.563e-06) (10,0.02812) (11,1.02e-07) (12,0.0055) (14,9.405e-05)};\draw[densely dotted] (axis cs:12,0) -- (axis cs:12,4.095);\end{axis}\begin{axis}[width=72mm,height=42mm,axis y line*=right,axis x line=none,ymode=log,xmin=-1,xmax=15,ylabel={\scriptsize eindeutig (log)},tick label style={font=\tiny},label style={font=\scriptsize},ylabel shift=-3pt]\addplot[mark=*,mark size=0.7pt,draw=black,thin] coordinates {(0,50238444) (1,145166391) (2,308119896) (3,239988126) (4,105880632) (5,422119) (6,2311061) (7,961) (8,488831) (9,123) (10,80277) (12,141)};\end{axis}\end{tikzpicture}\end{minipage}
\vspace{3mm}

\noindent\begin{minipage}[t]{60mm}\vspace{0pt}\sbox\hmboard{\chessboard[setfen=8/8/8/6Q1/2b4k/8/7R/K2n4 b - - 0 1]}\usebox\hmboard\\[1mm]
\makebox[\wd\hmboard]{\sffamily\hspace{3mm}\textbf{h\#6}\hfill(3+3)}\end{minipage}\hfill
\begin{minipage}[t]{52mm}\vspace{0pt}
\textcolor{ink}{\sffamily\textsc{Lösung}}\par\smallskip\small
\begin{tabular}{@{}r@{\enspace}l@{\enspace}l@{}}
\textbf{1.} & \symking{}xg5 & \symrook{}c2\\
\textbf{2.} & \symking{}f4 & \symrook{}xc4+\\
\textbf{3.} & \symking{}e3 & \symking{}a2\\
\textbf{4.} & \symking{}d2 & \symrook{}d4+\\
\textbf{5.} & \symking{}c1 & \symking{}b3\\
\textbf{6.} & \symking{}b1 & \symrook{}xd1\#\\
\end{tabular}\par\vspace{3mm}
\qrcode[height=12mm]{https://helpman.komtera.lt/?fen=8/8/8/6Q1/2b4k/8/7R/K2n4&moves=6}\\{\sffamily\tiny\color{ink}scannen zum Lösen}
\end{minipage}
\vspace{2mm}

{\raggedright\noindent\textcolor{ink}{\sffamily\textsc{Themen}}\enspace \themetag{reines Matt} \themetag{Mustermatt} \themetag{Idealmatt} \themetag{Spiegelmatt} \themetag{ein Stein zieht}\par\smallskip{\footnotesize\textit{Hinweis:} das Diagramm hat keinen legalen letzten Zug, kann also in keiner Partie entstehen.}\par}

\clearpage
\phantomsection\label{hm:KQRvkbp}
\begin{tcolorbox}[enhanced,colback=ink,colframe=ink,arc=0pt,boxrule=0pt,left=4mm,right=4mm,top=2.5mm,bottom=2.5mm]
\color{white}\sffamily{\LARGE\bfseries KQRvkbp}\hfill{\Large h\#6}\\[0.5mm]
{\small 6 Steine\quad Schwarz am Zug}\hfill{\small Nr.~684\qquad\href{https://helpman.komtera.lt/?fen=8/8/8/7Q/7k/p7/2b3R1/K7&moves=6}{\color{white}\bfseries lösen\,\ensuremath{\nearrow}}}
\end{tcolorbox}
\vspace{2mm}
\noindent\begin{minipage}[t]{36mm}\vspace{0pt}\sffamily\small\raggedright
\hmfigure{längste eindeutige}{{\Large h\#6}}\hmfigure{eindeutig bei dieser Länge}{{\Large 1\,435} Stellungen}\hmfigure{längstes Matt}{{\Large h\#7}}\hmfigure{in der Klasse}{45,4 Milliarden Stellungen,\\ 25,1 Milliarden mit Hilfsmatt}
\end{minipage}\hfill
\begin{minipage}[t]{80mm}\vspace{0pt}\begin{tikzpicture}\begin{axis}[width=72mm,height=42mm,ybar,bar width=1.6pt,xmin=-1,xmax=15,ymin=0,ymax=12.78,axis y line*=left,axis x line*=bottom,xlabel={\scriptsize Halbzüge bis Matt},ylabel={\scriptsize Stellungen (Mrd.)},tick label style={font=\tiny},label style={font=\scriptsize},enlarge x limits=false,ylabel shift=-3pt,xlabel shift=-3pt]\addplot[fill=ink!55,draw=none] coordinates {(0,0.1819) (1,0.7169) (2,2.705) (3,6.332) (4,11.62) (5,2.632) (6,0.6564) (7,0.003673) (8,0.1218) (9,1.704e-06) (10,0.08274) (11,1.45e-07) (12,0.01525) (13,2.1e-08) (14,0.0002578)};\draw[densely dotted] (axis cs:12,0) -- (axis cs:12,12.78);\end{axis}\begin{axis}[width=72mm,height=42mm,axis y line*=right,axis x line=none,ymode=log,xmin=-1,xmax=15,ylabel={\scriptsize eindeutig (log)},tick label style={font=\tiny},label style={font=\scriptsize},ylabel shift=-3pt]\addplot[mark=*,mark size=0.7pt,draw=black,thin] coordinates {(0,181856514) (1,512188952) (2,1078064933) (3,693918276) (4,290539707) (5,1000249) (6,7079453) (7,173) (8,1262864) (9,7) (10,190031) (11,1) (12,1435)};\end{axis}\end{tikzpicture}\end{minipage}
\vspace{3mm}

\noindent\begin{minipage}[t]{60mm}\vspace{0pt}\sbox\hmboard{\chessboard[setfen=8/8/8/7Q/7k/p7/2b3R1/K7 b - - 0 1]}\usebox\hmboard\\[1mm]
\makebox[\wd\hmboard]{\sffamily\hspace{3mm}\textbf{h\#6}\hfill(3+3)}\end{minipage}\hfill
\begin{minipage}[t]{52mm}\vspace{0pt}
\textcolor{ink}{\sffamily\textsc{Lösung}}\par\smallskip\small
\begin{tabular}{@{}r@{\enspace}l@{\enspace}l@{}}
\textbf{1.} & \symking{}xh5 & \symrook{}xc2\\
\textbf{2.} & \symking{}g4 & \symrook{}a2\\
\textbf{3.} & \symking{}f3 & \symrook{}xa3+\\
\textbf{4.} & \symking{}e2 & \symking{}b2\\
\textbf{5.} & \symking{}d1 & \symking{}c3\\
\textbf{6.} & \symking{}c1 & \symrook{}a1\#\\
\end{tabular}\par\vspace{3mm}
\qrcode[height=12mm]{https://helpman.komtera.lt/?fen=8/8/8/7Q/7k/p7/2b3R1/K7&moves=6}\\{\sffamily\tiny\color{ink}scannen zum Lösen}
\end{minipage}
\vspace{2mm}

{\raggedright\noindent\textcolor{ink}{\sffamily\textsc{Themen}}\enspace \themetag{reines Matt} \themetag{Mustermatt} \themetag{Idealmatt} \themetag{Spiegelmatt} \themetag{ein Stein zieht}\par\smallskip{\footnotesize\textit{Hinweis:} die Partei am Zug steht im Diagramm im Schach; die Lösung beginnt mit einem Schlag.}\par}

\clearpage
\phantomsection\label{hm:KQRvknn}
\begin{tcolorbox}[enhanced,colback=ink,colframe=ink,arc=0pt,boxrule=0pt,left=4mm,right=4mm,top=2.5mm,bottom=2.5mm]
\color{white}\sffamily{\LARGE\bfseries KQRvknn}\hfill{\Large h\#6}\\[0.5mm]
{\small 6 Steine\quad Schwarz am Zug}\hfill{\small Nr.~685\qquad\href{https://helpman.komtera.lt/?fen=8/8/8/6Q1/7k/2nn4/8/K6R&moves=6}{\color{white}\bfseries lösen\,\ensuremath{\nearrow}}}
\end{tcolorbox}
\vspace{2mm}
\noindent\begin{minipage}[t]{36mm}\vspace{0pt}\sffamily\small\raggedright
\hmfigure{längste eindeutige}{{\Large h\#6}}\hmfigure{eindeutig bei dieser Länge}{{\Large 294} Stellungen}\hmfigure{längstes Matt}{{\Large h\#7}}\hmfigure{in der Klasse}{15,5 Milliarden Stellungen,\\ 4,15 Milliarden mit Hilfsmatt}
\end{minipage}\hfill
\begin{minipage}[t]{80mm}\vspace{0pt}\begin{tikzpicture}\begin{axis}[width=72mm,height=42mm,ybar,bar width=1.6pt,xmin=-1,xmax=15,ymin=0,ymax=2.12,axis y line*=left,axis x line*=bottom,xlabel={\scriptsize Halbzüge bis Matt},ylabel={\scriptsize Stellungen (Mrd.)},tick label style={font=\tiny},label style={font=\scriptsize},enlarge x limits=false,ylabel shift=-3pt,xlabel shift=-3pt]\addplot[fill=ink!55,draw=none] coordinates {(0,0.02931) (1,0.1122) (2,0.4093) (3,0.9966) (4,1.927) (5,0.4958) (6,0.1399) (7,0.001178) (8,0.02052) (9,2.583e-06) (10,0.01468) (11,1.02e-07) (12,0.003011) (14,6.777e-05)};\draw[densely dotted] (axis cs:12,0) -- (axis cs:12,2.12);\end{axis}\begin{axis}[width=72mm,height=42mm,axis y line*=right,axis x line=none,ymode=log,xmin=-1,xmax=15,ylabel={\scriptsize eindeutig (log)},tick label style={font=\tiny},label style={font=\scriptsize},ylabel shift=-3pt]\addplot[mark=*,mark size=0.7pt,draw=black,thin] coordinates {(0,29305553) (1,79022863) (2,148534555) (3,116226374) (4,63929012) (5,258031) (6,1388431) (7,1262) (8,255238) (9,13) (10,36010) (12,294)};\end{axis}\end{tikzpicture}\end{minipage}
\vspace{3mm}

\noindent\begin{minipage}[t]{60mm}\vspace{0pt}\sbox\hmboard{\chessboard[setfen=8/8/8/6Q1/7k/2nn4/8/K6R b - - 0 1]}\usebox\hmboard\\[1mm]
\makebox[\wd\hmboard]{\sffamily\hspace{3mm}\textbf{h\#6}\hfill(3+3)}\end{minipage}\hfill
\begin{minipage}[t]{52mm}\vspace{0pt}
\textcolor{ink}{\sffamily\textsc{Lösung}}\par\smallskip\small
\begin{tabular}{@{}r@{\enspace}l@{\enspace}l@{}}
\textbf{1.} & \symking{}xg5 & \symrook{}c1\\
\textbf{2.} & \symking{}f4 & \symrook{}xc3\\
\textbf{3.} & \symking{}e3 & \symking{}a2\\
\textbf{4.} & \symking{}d2 & \symrook{}xd3+\\
\textbf{5.} & \symking{}c1 & \symking{}b3\\
\textbf{6.} & \symking{}b1 & \symrook{}d1\#\\
\end{tabular}\par\vspace{3mm}
\qrcode[height=12mm]{https://helpman.komtera.lt/?fen=8/8/8/6Q1/7k/2nn4/8/K6R&moves=6}\\{\sffamily\tiny\color{ink}scannen zum Lösen}
\end{minipage}
\vspace{2mm}

{\raggedright\noindent\textcolor{ink}{\sffamily\textsc{Themen}}\enspace \themetag{reines Matt} \themetag{Mustermatt} \themetag{Idealmatt} \themetag{Spiegelmatt} \themetag{ein Stein zieht}\par\smallskip{\footnotesize\textit{Hinweis:} das Diagramm hat keinen legalen letzten Zug, kann also in keiner Partie entstehen.}\par}

\clearpage
\phantomsection\label{hm:KQRvknp}
\begin{tcolorbox}[enhanced,colback=ink,colframe=ink,arc=0pt,boxrule=0pt,left=4mm,right=4mm,top=2.5mm,bottom=2.5mm]
\color{white}\sffamily{\LARGE\bfseries KQRvknp}\hfill{\Large h\#6}\\[0.5mm]
{\small 6 Steine\quad Schwarz am Zug}\hfill{\small Nr.~686\qquad\href{https://helpman.komtera.lt/?fen=8/8/8/7Q/7k/p1n5/8/K5R1&moves=6}{\color{white}\bfseries lösen\,\ensuremath{\nearrow}}}
\end{tcolorbox}
\vspace{2mm}
\noindent\begin{minipage}[t]{36mm}\vspace{0pt}\sffamily\small\raggedright
\hmfigure{längste eindeutige}{{\Large h\#6}}\hmfigure{eindeutig bei dieser Länge}{{\Large 2\,648} Stellungen}\hmfigure{längstes Matt}{{\Large h\#7}}\hmfigure{in der Klasse}{45,4 Milliarden Stellungen,\\ 25,9 Milliarden mit Hilfsmatt}
\end{minipage}\hfill
\begin{minipage}[t]{80mm}\vspace{0pt}\begin{tikzpicture}\begin{axis}[width=72mm,height=42mm,ybar,bar width=1.6pt,xmin=-1,xmax=15,ymin=0,ymax=13.27,axis y line*=left,axis x line*=bottom,xlabel={\scriptsize Halbzüge bis Matt},ylabel={\scriptsize Stellungen (Mrd.)},tick label style={font=\tiny},label style={font=\scriptsize},enlarge x limits=false,ylabel shift=-3pt,xlabel shift=-3pt]\addplot[fill=ink!55,draw=none] coordinates {(0,0.2143) (1,0.8132) (2,2.815) (3,6.325) (4,12.06) (5,2.54) (6,0.8715) (7,0.006196) (8,0.1279) (9,7.568e-06) (10,0.08683) (11,7.6e-08) (12,0.01662) (14,0.000338)};\draw[densely dotted] (axis cs:12,0) -- (axis cs:12,13.27);\end{axis}\begin{axis}[width=72mm,height=42mm,axis y line*=right,axis x line=none,ymode=log,xmin=-1,xmax=15,ylabel={\scriptsize eindeutig (log)},tick label style={font=\tiny},label style={font=\scriptsize},ylabel shift=-3pt]\addplot[mark=*,mark size=0.7pt,draw=black,thin] coordinates {(0,214313222) (1,558349582) (2,1035248205) (3,666344341) (4,357442777) (5,1272202) (6,8550707) (7,4784) (8,1270259) (9,1) (10,163607) (12,2648)};\end{axis}\end{tikzpicture}\end{minipage}
\vspace{3mm}

\noindent\begin{minipage}[t]{60mm}\vspace{0pt}\sbox\hmboard{\chessboard[setfen=8/8/8/7Q/7k/p1n5/8/K5R1 b - - 0 1]}\usebox\hmboard\\[1mm]
\makebox[\wd\hmboard]{\sffamily\hspace{3mm}\textbf{h\#6}\hfill(3+3)}\end{minipage}\hfill
\begin{minipage}[t]{52mm}\vspace{0pt}
\textcolor{ink}{\sffamily\textsc{Lösung}}\par\smallskip\small
\begin{tabular}{@{}r@{\enspace}l@{\enspace}l@{}}
\textbf{1.} & \symking{}xh5 & \symrook{}b1\\
\textbf{2.} & \symking{}g4 & \symrook{}b3\\
\textbf{3.} & \symking{}f3 & \symrook{}xa3\\
\textbf{4.} & \symking{}e2 & \symking{}b2\\
\textbf{5.} & \symking{}d1 & \symking{}xc3\\
\textbf{6.} & \symking{}c1 & \symrook{}a1\#\\
\end{tabular}\par\vspace{3mm}
\qrcode[height=12mm]{https://helpman.komtera.lt/?fen=8/8/8/7Q/7k/p1n5/8/K5R1&moves=6}\\{\sffamily\tiny\color{ink}scannen zum Lösen}
\end{minipage}
\vspace{2mm}

{\raggedright\noindent\textcolor{ink}{\sffamily\textsc{Themen}}\enspace \themetag{reines Matt} \themetag{Mustermatt} \themetag{Idealmatt} \themetag{Spiegelmatt} \themetag{ein Stein zieht} \themetag{ohne Schach}\par\smallskip{\footnotesize\textit{Hinweis:} die Partei am Zug steht im Diagramm im Schach; die Lösung beginnt mit einem Schlag.}\par}

\clearpage
\phantomsection\label{hm:KQRvkpp}
\begin{tcolorbox}[enhanced,colback=ink,colframe=ink,arc=0pt,boxrule=0pt,left=4mm,right=4mm,top=2.5mm,bottom=2.5mm]
\color{white}\sffamily{\LARGE\bfseries KQRvkpp}\hfill{\Large h\#6}\\[0.5mm]
{\small 6 Steine\quad Schwarz am Zug}\hfill{\small Nr.~687\qquad\href{https://helpman.komtera.lt/?fen=8/8/7Q/6k1/5R1p/2p5/8/K7&moves=6}{\color{white}\bfseries lösen\,\ensuremath{\nearrow}}}
\end{tcolorbox}
\vspace{2mm}
\noindent\begin{minipage}[t]{36mm}\vspace{0pt}\sffamily\small\raggedright
\hmfigure{längste eindeutige}{{\Large h\#6}}\hmfigure{eindeutig bei dieser Länge}{{\Large 1\,431} Stellungen}\hmfigure{längstes Matt}{{\Large h\#7}}\hmfigure{in der Klasse}{34,1 Milliarden Stellungen,\\ 10 Milliarden mit Hilfsmatt}
\end{minipage}\hfill
\begin{minipage}[t]{80mm}\vspace{0pt}\begin{tikzpicture}\begin{axis}[width=72mm,height=42mm,ybar,bar width=1.6pt,xmin=-1,xmax=15,ymin=0,ymax=5.137,axis y line*=left,axis x line*=bottom,xlabel={\scriptsize Halbzüge bis Matt},ylabel={\scriptsize Stellungen (Mrd.)},tick label style={font=\tiny},label style={font=\scriptsize},enlarge x limits=false,ylabel shift=-3pt,xlabel shift=-3pt]\addplot[fill=ink!55,draw=none] coordinates {(0,0.09828) (1,0.3677) (2,1.198) (3,2.475) (4,4.67) (5,0.7854) (6,0.3419) (7,0.001761) (8,0.04896) (9,5.84e-07) (10,0.03188) (11,1.78e-07) (12,0.005915) (13,3e-09) (14,0.0001294)};\draw[densely dotted] (axis cs:12,0) -- (axis cs:12,5.137);\end{axis}\begin{axis}[width=72mm,height=42mm,axis y line*=right,axis x line=none,ymode=log,xmin=-1,xmax=15,ylabel={\scriptsize eindeutig (log)},tick label style={font=\tiny},label style={font=\scriptsize},ylabel shift=-3pt]\addplot[mark=*,mark size=0.7pt,draw=black,thin] coordinates {(0,98279397) (1,245407891) (2,463337813) (3,232158005) (4,117720355) (5,477373) (6,3162102) (7,83) (8,366686) (9,3) (10,38024) (12,1431)};\end{axis}\end{tikzpicture}\end{minipage}
\vspace{3mm}

\noindent\begin{minipage}[t]{60mm}\vspace{0pt}\sbox\hmboard{\chessboard[setfen=8/8/7Q/6k1/5R1p/2p5/8/K7 b - - 0 1]}\usebox\hmboard\\[1mm]
\makebox[\wd\hmboard]{\sffamily\hspace{3mm}\textbf{h\#6}\hfill(3+3)}\end{minipage}\hfill
\begin{minipage}[t]{52mm}\vspace{0pt}
\textcolor{ink}{\sffamily\textsc{Lösung}}\par\smallskip\small
\begin{tabular}{@{}r@{\enspace}l@{\enspace}l@{}}
\textbf{1.} & \symking{}xh6 & \symking{}b1\\
\textbf{2.} & h3 & \symking{}c2\\
\textbf{3.} & h2 & \symking{}d3\\
\textbf{4.} & h1=\symrook{} & \symking{}e4\\
\textbf{5.} & \symrook{}g1 & \symking{}f5\\
\textbf{6.} & \symrook{}g7 & \symrook{}h4\#\\
\end{tabular}\par\vspace{3mm}
\qrcode[height=12mm]{https://helpman.komtera.lt/?fen=8/8/7Q/6k1/5R1p/2p5/8/K7&moves=6}\\{\sffamily\tiny\color{ink}scannen zum Lösen}
\end{minipage}
\vspace{2mm}

{\raggedright\noindent\textcolor{ink}{\sffamily\textsc{Themen}}\enspace \themetag{reines Matt} \themetag{Mustermatt} \themetag{Umwandlung} \themetag{Unterverwandlung} \themetag{Selbstblock} \themetag{Kniest} \themetag{ohne Schach}\par\smallskip{\footnotesize\textit{Hinweis:} die Partei am Zug steht im Diagramm im Schach; die Lösung beginnt mit einem Schlag.}\par}

\clearpage
\phantomsection\label{hm:KQRvkqb}
\begin{tcolorbox}[enhanced,colback=ink,colframe=ink,arc=0pt,boxrule=0pt,left=4mm,right=4mm,top=2.5mm,bottom=2.5mm]
\color{white}\sffamily{\LARGE\bfseries KQRvkqb}\hfill{\Large h\#6}\\[0.5mm]
{\small 6 Steine\quad Schwarz am Zug}\hfill{\small Nr.~688\qquad\href{https://helpman.komtera.lt/?fen=8/8/8/8/7Q/b2q2Rk/8/K7&moves=6}{\color{white}\bfseries lösen\,\ensuremath{\nearrow}}}
\end{tcolorbox}
\vspace{2mm}
\noindent\begin{minipage}[t]{36mm}\vspace{0pt}\sffamily\small\raggedright
\hmfigure{längste eindeutige}{{\Large h\#6}}\hmfigure{eindeutig bei dieser Länge}{{\Large 454} Stellungen}\hmfigure{längstes Matt}{{\Large h\#8.5}}\hmfigure{in der Klasse}{15,5 Milliarden Stellungen,\\ 6,74 Milliarden mit Hilfsmatt}
\end{minipage}\hfill
\begin{minipage}[t]{80mm}\vspace{0pt}\begin{tikzpicture}\begin{axis}[width=72mm,height=42mm,ybar,bar width=1.6pt,xmin=-1,xmax=18,ymin=0,ymax=3.055,axis y line*=left,axis x line*=bottom,xlabel={\scriptsize Halbzüge bis Matt},ylabel={\scriptsize Stellungen (Mrd.)},tick label style={font=\tiny},label style={font=\scriptsize},enlarge x limits=false,ylabel shift=-3pt,xlabel shift=-3pt]\addplot[fill=ink!55,draw=none] coordinates {(0,0.02609) (1,0.1079) (2,0.534) (3,1.692) (4,2.777) (5,1.401) (6,0.1387) (7,0.004094) (8,0.02956) (9,1.517e-05) (10,0.02237) (11,3.507e-06) (12,0.005023) (13,5.39e-07) (14,0.0001405) (15,9.2e-08) (16,6.4e-08) (17,1e-09)};\draw[densely dotted] (axis cs:12,0) -- (axis cs:12,3.055);\end{axis}\begin{axis}[width=72mm,height=42mm,axis y line*=right,axis x line=none,ymode=log,xmin=-1,xmax=18,ylabel={\scriptsize eindeutig (log)},tick label style={font=\tiny},label style={font=\scriptsize},ylabel shift=-3pt]\addplot[mark=*,mark size=0.7pt,draw=black,thin] coordinates {(0,26087521) (1,86752764) (2,237920739) (3,235687736) (4,66528289) (5,1558851) (6,1217275) (7,877) (8,426894) (9,239) (10,87168) (11,1) (12,454)};\end{axis}\end{tikzpicture}\end{minipage}
\vspace{3mm}

\noindent\begin{minipage}[t]{60mm}\vspace{0pt}\sbox\hmboard{\chessboard[setfen=8/8/8/8/7Q/b2q2Rk/8/K7 b - - 0 1]}\usebox\hmboard\\[1mm]
\makebox[\wd\hmboard]{\sffamily\hspace{3mm}\textbf{h\#6}\hfill(3+3)}\end{minipage}\hfill
\begin{minipage}[t]{52mm}\vspace{0pt}
\textcolor{ink}{\sffamily\textsc{Lösung}}\par\smallskip\small
\begin{tabular}{@{}r@{\enspace}l@{\enspace}l@{}}
\textbf{1.} & \symking{}xh4 & \symrook{}g7\\
\textbf{2.} & \symqueen{}h7 & \symrook{}xh7+\\
\textbf{3.} & \symking{}g3 & \symking{}b1\\
\textbf{4.} & \symking{}f2 & \symking{}c2\\
\textbf{5.} & \symking{}e1 & \symking{}d3\\
\textbf{6.} & \symking{}d1 & \symrook{}h1\#\\
\end{tabular}\par\vspace{3mm}
\qrcode[height=12mm]{https://helpman.komtera.lt/?fen=8/8/8/8/7Q/b2q2Rk/8/K7&moves=6}\\{\sffamily\tiny\color{ink}scannen zum Lösen}
\end{minipage}
\vspace{2mm}

{\raggedright\noindent\textcolor{ink}{\sffamily\textsc{Themen}}\enspace \themetag{reines Matt} \themetag{Mustermatt} \themetag{Spiegelmatt}\par\smallskip{\footnotesize\textit{Hinweis:} das Diagramm hat keinen legalen letzten Zug, kann also in keiner Partie entstehen.}\par}

\clearpage
\phantomsection\label{hm:KQRvkqn}
\begin{tcolorbox}[enhanced,colback=ink,colframe=ink,arc=0pt,boxrule=0pt,left=4mm,right=4mm,top=2.5mm,bottom=2.5mm]
\color{white}\sffamily{\LARGE\bfseries KQRvkqn}\hfill{\Large h\#6}\\[0.5mm]
{\small 6 Steine\quad Schwarz am Zug}\hfill{\small Nr.~689\qquad\href{https://helpman.komtera.lt/?fen=8/8/8/8/7Q/6nk/2q5/K5R1&moves=6}{\color{white}\bfseries lösen\,\ensuremath{\nearrow}}}
\end{tcolorbox}
\vspace{2mm}
\noindent\begin{minipage}[t]{36mm}\vspace{0pt}\sffamily\small\raggedright
\hmfigure{längste eindeutige}{{\Large h\#6}}\hmfigure{eindeutig bei dieser Länge}{{\Large 656} Stellungen}\hmfigure{längstes Matt}{{\Large h\#7}}\hmfigure{in der Klasse}{15,5 Milliarden Stellungen,\\ 6,96 Milliarden mit Hilfsmatt}
\end{minipage}\hfill
\begin{minipage}[t]{80mm}\vspace{0pt}\begin{tikzpicture}\begin{axis}[width=72mm,height=42mm,ybar,bar width=1.6pt,xmin=-1,xmax=15,ymin=0,ymax=3.228,axis y line*=left,axis x line*=bottom,xlabel={\scriptsize Halbzüge bis Matt},ylabel={\scriptsize Stellungen (Mrd.)},tick label style={font=\tiny},label style={font=\scriptsize},enlarge x limits=false,ylabel shift=-3pt,xlabel shift=-3pt]\addplot[fill=ink!55,draw=none] coordinates {(0,0.02974) (1,0.1212) (2,0.5708) (3,1.725) (4,2.935) (5,1.357) (6,0.1577) (7,0.004373) (8,0.03164) (9,2.215e-05) (10,0.02369) (11,1.932e-06) (12,0.005321) (13,2.3e-08) (14,0.0001526)};\draw[densely dotted] (axis cs:12,0) -- (axis cs:12,3.228);\end{axis}\begin{axis}[width=72mm,height=42mm,axis y line*=right,axis x line=none,ymode=log,xmin=-1,xmax=15,ylabel={\scriptsize eindeutig (log)},tick label style={font=\tiny},label style={font=\scriptsize},ylabel shift=-3pt]\addplot[mark=*,mark size=0.7pt,draw=black,thin] coordinates {(0,29739107) (1,95645444) (2,244359594) (3,236709313) (4,80529969) (5,1551315) (6,1501007) (7,1838) (8,472461) (9,420) (10,89584) (11,1) (12,656)};\end{axis}\end{tikzpicture}\end{minipage}
\vspace{3mm}

\noindent\begin{minipage}[t]{60mm}\vspace{0pt}\sbox\hmboard{\chessboard[setfen=8/8/8/8/7Q/6nk/2q5/K5R1 b - - 0 1]}\usebox\hmboard\\[1mm]
\makebox[\wd\hmboard]{\sffamily\hspace{3mm}\textbf{h\#6}\hfill(3+3)}\end{minipage}\hfill
\begin{minipage}[t]{52mm}\vspace{0pt}
\textcolor{ink}{\sffamily\textsc{Lösung}}\par\smallskip\small
\begin{tabular}{@{}r@{\enspace}l@{\enspace}l@{}}
\textbf{1.} & \symking{}xh4 & \symrook{}c1\\
\textbf{2.} & \symqueen{}f2 & \symrook{}c2\\
\textbf{3.} & \symking{}h3 & \symking{}b2\\
\textbf{4.} & \symking{}g2 & \symking{}c3\\
\textbf{5.} & \symking{}f1 & \symking{}d3\\
\textbf{6.} & \symking{}e1 & \symrook{}c1\#\\
\end{tabular}\par\vspace{3mm}
\qrcode[height=12mm]{https://helpman.komtera.lt/?fen=8/8/8/8/7Q/6nk/2q5/K5R1&moves=6}\\{\sffamily\tiny\color{ink}scannen zum Lösen}
\end{minipage}
\vspace{2mm}

{\raggedright\noindent\textcolor{ink}{\sffamily\textsc{Themen}}\enspace \themetag{reines Matt} \themetag{Mustermatt} \themetag{Switchback} \themetag{Selbstblock} \themetag{ohne Schach} \themetag{Umnov}\par\smallskip{\footnotesize\textit{Hinweis:} die Partei am Zug steht im Diagramm im Schach; die Lösung beginnt mit einem Schlag.}\par}

\clearpage
\phantomsection\label{hm:KQRvkqp}
\begin{tcolorbox}[enhanced,colback=ink,colframe=ink,arc=0pt,boxrule=0pt,left=4mm,right=4mm,top=2.5mm,bottom=2.5mm]
\color{white}\sffamily{\LARGE\bfseries KQRvkqp}\hfill{\Large h\#6}\\[0.5mm]
{\small 6 Steine\quad Schwarz am Zug}\hfill{\small Nr.~690\qquad\href{https://helpman.komtera.lt/?fen=8/8/8/5Q2/2p2k2/7R/3q4/K7&moves=6}{\color{white}\bfseries lösen\,\ensuremath{\nearrow}}}
\end{tcolorbox}
\vspace{2mm}
\noindent\begin{minipage}[t]{36mm}\vspace{0pt}\sffamily\small\raggedright
\hmfigure{längste eindeutige}{{\Large h\#6}}\hmfigure{eindeutig bei dieser Länge}{{\Large 1\,359} Stellungen}\hmfigure{längstes Matt}{{\Large h\#7}}\hmfigure{in der Klasse}{45,4 Milliarden Stellungen,\\ 21,6 Milliarden mit Hilfsmatt}
\end{minipage}\hfill
\begin{minipage}[t]{80mm}\vspace{0pt}\begin{tikzpicture}\begin{axis}[width=72mm,height=42mm,ybar,bar width=1.6pt,xmin=-1,xmax=15,ymin=0,ymax=10.1,axis y line*=left,axis x line*=bottom,xlabel={\scriptsize Halbzüge bis Matt},ylabel={\scriptsize Stellungen (Mrd.)},tick label style={font=\tiny},label style={font=\scriptsize},enlarge x limits=false,ylabel shift=-3pt,xlabel shift=-3pt]\addplot[fill=ink!55,draw=none] coordinates {(0,0.1063) (1,0.4374) (2,1.984) (3,5.553) (4,9.182) (5,3.674) (6,0.4652) (7,0.01263) (8,0.09928) (9,3.159e-05) (10,0.07002) (11,6.065e-06) (12,0.01481) (13,4.54e-07) (14,0.0004147)};\draw[densely dotted] (axis cs:12,0) -- (axis cs:12,10.1);\end{axis}\begin{axis}[width=72mm,height=42mm,axis y line*=right,axis x line=none,ymode=log,xmin=-1,xmax=15,ylabel={\scriptsize eindeutig (log)},tick label style={font=\tiny},label style={font=\scriptsize},ylabel shift=-3pt]\addplot[mark=*,mark size=0.7pt,draw=black,thin] coordinates {(0,106306358) (1,339118772) (2,854763680) (3,713311860) (4,213037823) (5,4002838) (6,4441382) (7,2276) (8,1235105) (9,127) (10,234785) (11,1) (12,1359)};\end{axis}\end{tikzpicture}\end{minipage}
\vspace{3mm}

\noindent\begin{minipage}[t]{60mm}\vspace{0pt}\sbox\hmboard{\chessboard[setfen=8/8/8/5Q2/2p2k2/7R/3q4/K7 b - - 0 1]}\usebox\hmboard\\[1mm]
\makebox[\wd\hmboard]{\sffamily\hspace{3mm}\textbf{h\#6}\hfill(3+3)}\end{minipage}\hfill
\begin{minipage}[t]{52mm}\vspace{0pt}
\textcolor{ink}{\sffamily\textsc{Lösung}}\par\smallskip\small
\begin{tabular}{@{}r@{\enspace}l@{\enspace}l@{}}
\textbf{1.} & \symking{}xf5 & \symrook{}c3\\
\textbf{2.} & \symking{}e4 & \symrook{}xc4+\\
\textbf{3.} & \symking{}d3 & \symrook{}a4\\
\textbf{4.} & \symking{}c2 & \symking{}a2\\
\textbf{5.} & \symking{}d1+ & \symking{}b3\\
\textbf{6.} & \symking{}c1 & \symrook{}a1\#\\
\end{tabular}\par\vspace{3mm}
\qrcode[height=12mm]{https://helpman.komtera.lt/?fen=8/8/8/5Q2/2p2k2/7R/3q4/K7&moves=6}\\{\sffamily\tiny\color{ink}scannen zum Lösen}
\end{minipage}
\vspace{2mm}

{\raggedright\noindent\textcolor{ink}{\sffamily\textsc{Themen}}\enspace \themetag{reines Matt} \themetag{Mustermatt} \themetag{Idealmatt} \themetag{ein Stein zieht}\par\smallskip{\footnotesize\textit{Hinweis:} die Partei am Zug steht im Diagramm im Schach; die Lösung beginnt mit einem Schlag.}\par}

\clearpage
\phantomsection\label{hm:KQRvkqq}
\begin{tcolorbox}[enhanced,colback=ink,colframe=ink,arc=0pt,boxrule=0pt,left=4mm,right=4mm,top=2.5mm,bottom=2.5mm]
\color{white}\sffamily{\LARGE\bfseries KQRvkqq}\hfill{\Large h\#6}\\[0.5mm]
{\small 6 Steine\quad Schwarz am Zug}\hfill{\small Nr.~691\qquad\href{https://helpman.komtera.lt/?fen=8/8/8/8/7Q/4q1Rk/2q5/K7&moves=6}{\color{white}\bfseries lösen\,\ensuremath{\nearrow}}}
\end{tcolorbox}
\vspace{2mm}
\noindent\begin{minipage}[t]{36mm}\vspace{0pt}\sffamily\small\raggedright
\hmfigure{längste eindeutige}{{\Large h\#6}}\hmfigure{eindeutig bei dieser Länge}{{\Large 1\,440} Stellungen}\hmfigure{längstes Matt}{{\Large h\#8.5}}\hmfigure{in der Klasse}{15,5 Milliarden Stellungen,\\ 2,91 Milliarden mit Hilfsmatt}
\end{minipage}\hfill
\begin{minipage}[t]{80mm}\vspace{0pt}\begin{tikzpicture}\begin{axis}[width=72mm,height=42mm,ybar,bar width=1.6pt,xmin=-1,xmax=18,ymin=0,ymax=1.174,axis y line*=left,axis x line*=bottom,xlabel={\scriptsize Halbzüge bis Matt},ylabel={\scriptsize Stellungen (Mrd.)},tick label style={font=\tiny},label style={font=\scriptsize},enlarge x limits=false,ylabel shift=-3pt,xlabel shift=-3pt]\addplot[fill=ink!55,draw=none] coordinates {(0,0.007926) (1,0.03239) (2,0.1812) (3,0.7002) (4,1.067) (5,0.8357) (6,0.05542) (7,0.0047) (8,0.01116) (9,2.631e-05) (10,0.00899) (11,7.411e-06) (12,0.002458) (13,9.35e-07) (14,0.0001116) (15,9.3e-08) (16,2.96e-07) (17,1e-09)};\draw[densely dotted] (axis cs:12,0) -- (axis cs:12,1.174);\end{axis}\begin{axis}[width=72mm,height=42mm,axis y line*=right,axis x line=none,ymode=log,xmin=-1,xmax=18,ylabel={\scriptsize eindeutig (log)},tick label style={font=\tiny},label style={font=\scriptsize},ylabel shift=-3pt]\addplot[mark=*,mark size=0.7pt,draw=black,thin] coordinates {(0,7925770) (1,27277649) (2,85028324) (3,111058673) (4,26741708) (5,1794855) (6,397928) (7,1822) (8,167922) (9,986) (10,34722) (11,7) (12,1440)};\end{axis}\end{tikzpicture}\end{minipage}
\vspace{3mm}

\noindent\begin{minipage}[t]{60mm}\vspace{0pt}\sbox\hmboard{\chessboard[setfen=8/8/8/8/7Q/4q1Rk/2q5/K7 b - - 0 1]}\usebox\hmboard\\[1mm]
\makebox[\wd\hmboard]{\sffamily\hspace{3mm}\textbf{h\#6}\hfill(3+3)}\end{minipage}\hfill
\begin{minipage}[t]{52mm}\vspace{0pt}
\textcolor{ink}{\sffamily\textsc{Lösung}}\par\smallskip\small
\begin{tabular}{@{}r@{\enspace}l@{\enspace}l@{}}
\textbf{1.} & \symking{}xh4 & \symrook{}g2\\
\textbf{2.} & \symking{}h3 & \symrook{}xc2\\
\textbf{3.} & \symqueen{}f2 & \symking{}b2\\
\textbf{4.} & \symking{}g2 & \symking{}c3\\
\textbf{5.} & \symking{}f1 & \symking{}d3\\
\textbf{6.} & \symking{}e1 & \symrook{}c1\#\\
\end{tabular}\par\vspace{3mm}
\qrcode[height=12mm]{https://helpman.komtera.lt/?fen=8/8/8/8/7Q/4q1Rk/2q5/K7&moves=6}\\{\sffamily\tiny\color{ink}scannen zum Lösen}
\end{minipage}
\vspace{2mm}

{\raggedright\noindent\textcolor{ink}{\sffamily\textsc{Themen}}\enspace \themetag{reines Matt} \themetag{Mustermatt} \themetag{Idealmatt} \themetag{Switchback} \themetag{Selbstblock} \themetag{ohne Schach}\par\smallskip{\footnotesize\textit{Hinweis:} das Diagramm hat keinen legalen letzten Zug, kann also in keiner Partie entstehen.}\par}

\clearpage
\phantomsection\label{hm:KQRvkqr}
\begin{tcolorbox}[enhanced,colback=ink,colframe=ink,arc=0pt,boxrule=0pt,left=4mm,right=4mm,top=2.5mm,bottom=2.5mm]
\color{white}\sffamily{\LARGE\bfseries KQRvkqr}\hfill{\Large h\#6}\\[0.5mm]
{\small 6 Steine\quad Schwarz am Zug}\hfill{\small Nr.~692\qquad\href{https://helpman.komtera.lt/?fen=8/8/8/4Q3/1qrk4/8/3R4/K7&moves=6}{\color{white}\bfseries lösen\,\ensuremath{\nearrow}}}
\end{tcolorbox}
\vspace{2mm}
\noindent\begin{minipage}[t]{36mm}\vspace{0pt}\sffamily\small\raggedright
\hmfigure{längste eindeutige}{{\Large h\#6}}\hmfigure{eindeutig bei dieser Länge}{{\Large 2\,025} Stellungen}\hmfigure{längstes Matt}{{\Large h\#8}}\hmfigure{in der Klasse}{15,5 Milliarden Stellungen,\\ 6,4 Milliarden mit Hilfsmatt}
\end{minipage}\hfill
\begin{minipage}[t]{80mm}\vspace{0pt}\begin{tikzpicture}\begin{axis}[width=72mm,height=42mm,ybar,bar width=1.6pt,xmin=-1,xmax=17,ymin=0,ymax=2.778,axis y line*=left,axis x line*=bottom,xlabel={\scriptsize Halbzüge bis Matt},ylabel={\scriptsize Stellungen (Mrd.)},tick label style={font=\tiny},label style={font=\scriptsize},enlarge x limits=false,ylabel shift=-3pt,xlabel shift=-3pt]\addplot[fill=ink!55,draw=none] coordinates {(0,0.0229) (1,0.09602) (2,0.5039) (3,1.639) (4,2.526) (5,1.449) (6,0.1091) (7,0.004818) (8,0.02668) (9,2.882e-05) (10,0.01811) (11,4.487e-06) (12,0.004989) (13,1.23e-07) (14,0.0001959) (16,1.6e-08)};\draw[densely dotted] (axis cs:12,0) -- (axis cs:12,2.778);\end{axis}\begin{axis}[width=72mm,height=42mm,axis y line*=right,axis x line=none,ymode=log,xmin=-1,xmax=17,ylabel={\scriptsize eindeutig (log)},tick label style={font=\tiny},label style={font=\scriptsize},ylabel shift=-3pt]\addplot[mark=*,mark size=0.7pt,draw=black,thin] coordinates {(0,22900158) (1,78184628) (2,232247453) (3,230590823) (4,55568525) (5,1914044) (6,971745) (7,1377) (8,289122) (9,836) (10,69164) (11,9) (12,2025)};\end{axis}\end{tikzpicture}\end{minipage}
\vspace{3mm}

\noindent\begin{minipage}[t]{60mm}\vspace{0pt}\sbox\hmboard{\chessboard[setfen=8/8/8/4Q3/1qrk4/8/3R4/K7 b - - 0 1]}\usebox\hmboard\\[1mm]
\makebox[\wd\hmboard]{\sffamily\hspace{3mm}\textbf{h\#6}\hfill(3+3)}\end{minipage}\hfill
\begin{minipage}[t]{52mm}\vspace{0pt}
\textcolor{ink}{\sffamily\textsc{Lösung}}\par\smallskip\small
\begin{tabular}{@{}r@{\enspace}l@{\enspace}l@{}}
\textbf{1.} & \symking{}xe5 & \symrook{}b2\\
\textbf{2.} & \symking{}d4 & \symrook{}xb4\\
\textbf{3.} & \symking{}c3 & \symrook{}xc4+\\
\textbf{4.} & \symking{}b3 & \symking{}b1\\
\textbf{5.} & \symking{}a3 & \symking{}c2\\
\textbf{6.} & \symking{}a2 & \symrook{}a4\#\\
\end{tabular}\par\vspace{3mm}
\qrcode[height=12mm]{https://helpman.komtera.lt/?fen=8/8/8/4Q3/1qrk4/8/3R4/K7&moves=6}\\{\sffamily\tiny\color{ink}scannen zum Lösen}
\end{minipage}
\vspace{2mm}

{\raggedright\noindent\textcolor{ink}{\sffamily\textsc{Themen}}\enspace \themetag{reines Matt} \themetag{Mustermatt} \themetag{Idealmatt} \themetag{Spiegelmatt} \themetag{Switchback} \themetag{ein Stein zieht}\par\smallskip{\footnotesize\textit{Hinweis:} das Diagramm hat keinen legalen letzten Zug, kann also in keiner Partie entstehen.}\par}

\clearpage
\phantomsection\label{hm:KQRvkrb}
\begin{tcolorbox}[enhanced,colback=ink,colframe=ink,arc=0pt,boxrule=0pt,left=4mm,right=4mm,top=2.5mm,bottom=2.5mm]
\color{white}\sffamily{\LARGE\bfseries KQRvkrb}\hfill{\Large h\#6}\\[0.5mm]
{\small 6 Steine\quad Schwarz am Zug}\hfill{\small Nr.~693\qquad\href{https://helpman.komtera.lt/?fen=8/8/8/6Q1/7k/8/b1r5/K6R&moves=6}{\color{white}\bfseries lösen\,\ensuremath{\nearrow}}}
\end{tcolorbox}
\vspace{2mm}
\noindent\begin{minipage}[t]{36mm}\vspace{0pt}\sffamily\small\raggedright
\hmfigure{längste eindeutige}{{\Large h\#6}}\hmfigure{eindeutig bei dieser Länge}{{\Large 850} Stellungen}\hmfigure{längstes Matt}{{\Large h\#7}}\hmfigure{in der Klasse}{15,5 Milliarden Stellungen,\\ 7,45 Milliarden mit Hilfsmatt}
\end{minipage}\hfill
\begin{minipage}[t]{80mm}\vspace{0pt}\begin{tikzpicture}\begin{axis}[width=72mm,height=42mm,ybar,bar width=1.6pt,xmin=-1,xmax=15,ymin=0,ymax=3.598,axis y line*=left,axis x line*=bottom,xlabel={\scriptsize Halbzüge bis Matt},ylabel={\scriptsize Stellungen (Mrd.)},tick label style={font=\tiny},label style={font=\scriptsize},enlarge x limits=false,ylabel shift=-3pt,xlabel shift=-3pt]\addplot[fill=ink!55,draw=none] coordinates {(0,0.03903) (1,0.1615) (2,0.7329) (3,1.934) (4,3.271) (5,1.113) (6,0.1359) (7,0.001482) (8,0.03437) (9,1.691e-06) (10,0.02273) (11,1.5e-07) (12,0.005109) (13,5e-09) (14,0.0001159)};\draw[densely dotted] (axis cs:12,0) -- (axis cs:12,3.598);\end{axis}\begin{axis}[width=72mm,height=42mm,axis y line*=right,axis x line=none,ymode=log,xmin=-1,xmax=15,ylabel={\scriptsize eindeutig (log)},tick label style={font=\tiny},label style={font=\scriptsize},ylabel shift=-3pt]\addplot[mark=*,mark size=0.7pt,draw=black,thin] coordinates {(0,39034046) (1,123855970) (2,317820951) (3,232085835) (4,69351873) (5,542643) (6,1500858) (7,362) (8,334603) (9,41) (10,76454) (12,850)};\end{axis}\end{tikzpicture}\end{minipage}
\vspace{3mm}

\noindent\begin{minipage}[t]{60mm}\vspace{0pt}\sbox\hmboard{\chessboard[setfen=8/8/8/6Q1/7k/8/b1r5/K6R b - - 0 1]}\usebox\hmboard\\[1mm]
\makebox[\wd\hmboard]{\sffamily\hspace{3mm}\textbf{h\#6}\hfill(3+3)}\end{minipage}\hfill
\begin{minipage}[t]{52mm}\vspace{0pt}
\textcolor{ink}{\sffamily\textsc{Lösung}}\par\smallskip\small
\begin{tabular}{@{}r@{\enspace}l@{\enspace}l@{}}
\textbf{1.} & \symking{}xg5 & \symrook{}b1\\
\textbf{2.} & \symking{}f4 & \symrook{}b2\\
\textbf{3.} & \symking{}e3 & \symking{}xa2\\
\textbf{4.} & \symking{}d2 & \symking{}b3\\
\textbf{5.} & \symking{}c1 & \symrook{}a2\\
\textbf{6.} & \symrook{}d2 & \symrook{}a1\#\\
\end{tabular}\par\vspace{3mm}
\qrcode[height=12mm]{https://helpman.komtera.lt/?fen=8/8/8/6Q1/7k/8/b1r5/K6R&moves=6}\\{\sffamily\tiny\color{ink}scannen zum Lösen}
\end{minipage}
\vspace{2mm}

{\raggedright\noindent\textcolor{ink}{\sffamily\textsc{Themen}}\enspace \themetag{reines Matt} \themetag{Mustermatt} \themetag{Idealmatt} \themetag{Selbstblock} \themetag{ohne Schach}\par\smallskip{\footnotesize\textit{Hinweis:} das Diagramm hat keinen legalen letzten Zug, kann also in keiner Partie entstehen.}\par}

\clearpage
\phantomsection\label{hm:KQRvkrn}
\begin{tcolorbox}[enhanced,colback=ink,colframe=ink,arc=0pt,boxrule=0pt,left=4mm,right=4mm,top=2.5mm,bottom=2.5mm]
\color{white}\sffamily{\LARGE\bfseries KQRvkrn}\hfill{\Large h\#6}\\[0.5mm]
{\small 6 Steine\quad Schwarz am Zug}\hfill{\small Nr.~694\qquad\href{https://helpman.komtera.lt/?fen=8/8/8/7Q/1r5k/8/8/K3n1R1&moves=6}{\color{white}\bfseries lösen\,\ensuremath{\nearrow}}}
\end{tcolorbox}
\vspace{2mm}
\noindent\begin{minipage}[t]{36mm}\vspace{0pt}\sffamily\small\raggedright
\hmfigure{längste eindeutige}{{\Large h\#6}}\hmfigure{eindeutig bei dieser Länge}{{\Large 562} Stellungen}\hmfigure{längstes Matt}{{\Large h\#7}}\hmfigure{in der Klasse}{15,5 Milliarden Stellungen,\\ 7,67 Milliarden mit Hilfsmatt}
\end{minipage}\hfill
\begin{minipage}[t]{80mm}\vspace{0pt}\begin{tikzpicture}\begin{axis}[width=72mm,height=42mm,ybar,bar width=1.6pt,xmin=-1,xmax=15,ymin=0,ymax=3.758,axis y line*=left,axis x line*=bottom,xlabel={\scriptsize Halbzüge bis Matt},ylabel={\scriptsize Stellungen (Mrd.)},tick label style={font=\tiny},label style={font=\scriptsize},enlarge x limits=false,ylabel shift=-3pt,xlabel shift=-3pt]\addplot[fill=ink!55,draw=none] coordinates {(0,0.04447) (1,0.1803) (2,0.7677) (3,1.944) (4,3.417) (5,1.085) (6,0.1616) (7,0.001657) (8,0.03693) (9,5.123e-06) (10,0.02391) (11,1.75e-07) (12,0.004629) (14,7.512e-05)};\draw[densely dotted] (axis cs:12,0) -- (axis cs:12,3.758);\end{axis}\begin{axis}[width=72mm,height=42mm,axis y line*=right,axis x line=none,ymode=log,xmin=-1,xmax=15,ylabel={\scriptsize eindeutig (log)},tick label style={font=\tiny},label style={font=\scriptsize},ylabel shift=-3pt]\addplot[mark=*,mark size=0.7pt,draw=black,thin] coordinates {(0,44472328) (1,134959224) (2,315845806) (3,228441682) (4,85288553) (5,562886) (6,1868871) (7,1365) (8,393618) (9,46) (10,71072) (12,562)};\end{axis}\end{tikzpicture}\end{minipage}
\vspace{3mm}

\noindent\begin{minipage}[t]{60mm}\vspace{0pt}\sbox\hmboard{\chessboard[setfen=8/8/8/7Q/1r5k/8/8/K3n1R1 b - - 0 1]}\usebox\hmboard\\[1mm]
\makebox[\wd\hmboard]{\sffamily\hspace{3mm}\textbf{h\#6}\hfill(3+3)}\end{minipage}\hfill
\begin{minipage}[t]{52mm}\vspace{0pt}
\textcolor{ink}{\sffamily\textsc{Lösung}}\par\smallskip\small
\begin{tabular}{@{}r@{\enspace}l@{\enspace}l@{}}
\textbf{1.} & \symking{}xh5 & \symrook{}h1+\\
\textbf{2.} & \symking{}g4 & \symking{}a2\\
\textbf{3.} & \symking{}f3 & \symking{}a3\\
\textbf{4.} & \symking{}e2 & \symking{}xb4\\
\textbf{5.} & \symking{}d1 & \symking{}c3\\
\textbf{6.} & \symking{}c1 & \symrook{}xe1\#\\
\end{tabular}\par\vspace{3mm}
\qrcode[height=12mm]{https://helpman.komtera.lt/?fen=8/8/8/7Q/1r5k/8/8/K3n1R1&moves=6}\\{\sffamily\tiny\color{ink}scannen zum Lösen}
\end{minipage}
\vspace{2mm}

{\raggedright\noindent\textcolor{ink}{\sffamily\textsc{Themen}}\enspace \themetag{reines Matt} \themetag{Mustermatt} \themetag{Idealmatt} \themetag{Spiegelmatt} \themetag{ein Stein zieht}\par\smallskip{\footnotesize\textit{Hinweis:} die Partei am Zug steht im Diagramm im Schach; die Lösung beginnt mit einem Schlag.}\par}

\clearpage
\phantomsection\label{hm:KQRvkrp}
\begin{tcolorbox}[enhanced,colback=ink,colframe=ink,arc=0pt,boxrule=0pt,left=4mm,right=4mm,top=2.5mm,bottom=2.5mm]
\color{white}\sffamily{\LARGE\bfseries KQRvkrp}\hfill{\Large h\#6}\\[0.5mm]
{\small 6 Steine\quad Schwarz am Zug}\hfill{\small Nr.~695\qquad\href{https://helpman.komtera.lt/?fen=8/8/8/6Q1/7k/8/p1r5/K6R&moves=6}{\color{white}\bfseries lösen\,\ensuremath{\nearrow}}}
\end{tcolorbox}
\vspace{2mm}
\noindent\begin{minipage}[t]{36mm}\vspace{0pt}\sffamily\small\raggedright
\hmfigure{längste eindeutige}{{\Large h\#6}}\hmfigure{eindeutig bei dieser Länge}{{\Large 2\,337} Stellungen}\hmfigure{längstes Matt}{{\Large h\#7}}\hmfigure{in der Klasse}{45,4 Milliarden Stellungen,\\ 23,9 Milliarden mit Hilfsmatt}
\end{minipage}\hfill
\begin{minipage}[t]{80mm}\vspace{0pt}\begin{tikzpicture}\begin{axis}[width=72mm,height=42mm,ybar,bar width=1.6pt,xmin=-1,xmax=15,ymin=0,ymax=11.74,axis y line*=left,axis x line*=bottom,xlabel={\scriptsize Halbzüge bis Matt},ylabel={\scriptsize Stellungen (Mrd.)},tick label style={font=\tiny},label style={font=\scriptsize},enlarge x limits=false,ylabel shift=-3pt,xlabel shift=-3pt]\addplot[fill=ink!55,draw=none] coordinates {(0,0.1618) (1,0.6549) (2,2.671) (3,6.209) (4,10.68) (5,2.815) (6,0.4635) (7,0.00404) (8,0.1122) (9,2.167e-06) (10,0.07265) (11,1.1e-07) (12,0.01495) (13,6e-09) (14,0.0003397)};\draw[densely dotted] (axis cs:12,0) -- (axis cs:12,11.74);\end{axis}\begin{axis}[width=72mm,height=42mm,axis y line*=right,axis x line=none,ymode=log,xmin=-1,xmax=15,ylabel={\scriptsize eindeutig (log)},tick label style={font=\tiny},label style={font=\scriptsize},ylabel shift=-3pt]\addplot[mark=*,mark size=0.7pt,draw=black,thin] coordinates {(0,161762682) (1,479320911) (2,1112043078) (3,655616423) (4,218206215) (5,1263170) (6,5587496) (7,289) (8,1042136) (9,49) (10,215042) (12,2337)};\end{axis}\end{tikzpicture}\end{minipage}
\vspace{3mm}

\noindent\begin{minipage}[t]{60mm}\vspace{0pt}\sbox\hmboard{\chessboard[setfen=8/8/8/6Q1/7k/8/p1r5/K6R b - - 0 1]}\usebox\hmboard\\[1mm]
\makebox[\wd\hmboard]{\sffamily\hspace{3mm}\textbf{h\#6}\hfill(3+3)}\end{minipage}\hfill
\begin{minipage}[t]{52mm}\vspace{0pt}
\textcolor{ink}{\sffamily\textsc{Lösung}}\par\smallskip\small
\begin{tabular}{@{}r@{\enspace}l@{\enspace}l@{}}
\textbf{1.} & \symking{}xg5 & \symrook{}b1\\
\textbf{2.} & \symking{}f4 & \symrook{}b2\\
\textbf{3.} & \symking{}e3 & \symking{}xa2\\
\textbf{4.} & \symking{}d2 & \symking{}b3\\
\textbf{5.} & \symking{}c1 & \symrook{}a2\\
\textbf{6.} & \symrook{}d2 & \symrook{}a1\#\\
\end{tabular}\par\vspace{3mm}
\qrcode[height=12mm]{https://helpman.komtera.lt/?fen=8/8/8/6Q1/7k/8/p1r5/K6R&moves=6}\\{\sffamily\tiny\color{ink}scannen zum Lösen}
\end{minipage}
\vspace{2mm}

{\raggedright\noindent\textcolor{ink}{\sffamily\textsc{Themen}}\enspace \themetag{reines Matt} \themetag{Mustermatt} \themetag{Idealmatt} \themetag{Selbstblock} \themetag{ohne Schach}\par\smallskip{\footnotesize\textit{Hinweis:} das Diagramm hat keinen legalen letzten Zug, kann also in keiner Partie entstehen.}\par}

\clearpage
\phantomsection\label{hm:KQRvkrr}
\begin{tcolorbox}[enhanced,colback=ink,colframe=ink,arc=0pt,boxrule=0pt,left=4mm,right=4mm,top=2.5mm,bottom=2.5mm]
\color{white}\sffamily{\LARGE\bfseries KQRvkrr}\hfill{\Large h\#6}\\[0.5mm]
{\small 6 Steine\quad Schwarz am Zug}\hfill{\small Nr.~696\qquad\href{https://helpman.komtera.lt/?fen=8/8/8/4Q3/1rrk4/8/8/K2R4&moves=6}{\color{white}\bfseries lösen\,\ensuremath{\nearrow}}}
\end{tcolorbox}
\vspace{2mm}
\noindent\begin{minipage}[t]{36mm}\vspace{0pt}\sffamily\small\raggedright
\hmfigure{längste eindeutige}{{\Large h\#6}}\hmfigure{eindeutig bei dieser Länge}{{\Large 1\,093} Stellungen}\hmfigure{längstes Matt}{{\Large h\#8}}\hmfigure{in der Klasse}{15,5 Milliarden Stellungen,\\ 3,51 Milliarden mit Hilfsmatt}
\end{minipage}\hfill
\begin{minipage}[t]{80mm}\vspace{0pt}\begin{tikzpicture}\begin{axis}[width=72mm,height=42mm,ybar,bar width=1.6pt,xmin=-1,xmax=17,ymin=0,ymax=1.615,axis y line*=left,axis x line*=bottom,xlabel={\scriptsize Halbzüge bis Matt},ylabel={\scriptsize Stellungen (Mrd.)},tick label style={font=\tiny},label style={font=\scriptsize},enlarge x limits=false,ylabel shift=-3pt,xlabel shift=-3pt]\addplot[fill=ink!55,draw=none] coordinates {(0,0.01673) (1,0.07081) (2,0.3426) (3,0.93) (4,1.468) (5,0.5993) (6,0.05499) (7,0.001387) (8,0.01491) (9,8.572e-06) (10,0.01026) (11,7.79e-07) (12,0.002732) (13,2.5e-08) (14,8.573e-05) (16,3e-09)};\draw[densely dotted] (axis cs:12,0) -- (axis cs:12,1.615);\end{axis}\begin{axis}[width=72mm,height=42mm,axis y line*=right,axis x line=none,ymode=log,xmin=-1,xmax=17,ylabel={\scriptsize eindeutig (log)},tick label style={font=\tiny},label style={font=\scriptsize},ylabel shift=-3pt]\addplot[mark=*,mark size=0.7pt,draw=black,thin] coordinates {(0,16727495) (1,55332885) (2,152628511) (3,110953229) (4,27523242) (5,466403) (6,556912) (7,416) (8,143491) (9,170) (10,35678) (11,2) (12,1093)};\end{axis}\end{tikzpicture}\end{minipage}
\vspace{3mm}

\noindent\begin{minipage}[t]{60mm}\vspace{0pt}\sbox\hmboard{\chessboard[setfen=8/8/8/4Q3/1rrk4/8/8/K2R4 b - - 0 1]}\usebox\hmboard\\[1mm]
\makebox[\wd\hmboard]{\sffamily\hspace{3mm}\textbf{h\#6}\hfill(3+3)}\end{minipage}\hfill
\begin{minipage}[t]{52mm}\vspace{0pt}
\textcolor{ink}{\sffamily\textsc{Lösung}}\par\smallskip\small
\begin{tabular}{@{}r@{\enspace}l@{\enspace}l@{}}
\textbf{1.} & \symking{}xe5 & \symrook{}b1\\
\textbf{2.} & \symking{}d4 & \symrook{}xb4\\
\textbf{3.} & \symking{}c3 & \symrook{}xc4+\\
\textbf{4.} & \symking{}b3 & \symking{}b1\\
\textbf{5.} & \symking{}a3 & \symking{}c2\\
\textbf{6.} & \symking{}a2 & \symrook{}a4\#\\
\end{tabular}\par\vspace{3mm}
\qrcode[height=12mm]{https://helpman.komtera.lt/?fen=8/8/8/4Q3/1rrk4/8/8/K2R4&moves=6}\\{\sffamily\tiny\color{ink}scannen zum Lösen}
\end{minipage}
\vspace{2mm}

{\raggedright\noindent\textcolor{ink}{\sffamily\textsc{Themen}}\enspace \themetag{reines Matt} \themetag{Mustermatt} \themetag{Idealmatt} \themetag{Spiegelmatt} \themetag{Switchback} \themetag{ein Stein zieht}\par\smallskip{\footnotesize\textit{Hinweis:} das Diagramm hat keinen legalen letzten Zug, kann also in keiner Partie entstehen.}\par}

\clearpage
\phantomsection\label{hm:KQvkbbn}
\begin{tcolorbox}[enhanced,colback=ink,colframe=ink,arc=0pt,boxrule=0pt,left=4mm,right=4mm,top=2.5mm,bottom=2.5mm]
\color{white}\sffamily{\LARGE\bfseries KQvkbbn}\hfill{\Large h\#6}\\[0.5mm]
{\small 6 Steine\quad Schwarz am Zug}\hfill{\small Nr.~697\qquad\href{https://helpman.komtera.lt/?fen=8/8/8/7k/1n4Q1/b7/2b5/K7&moves=6}{\color{white}\bfseries lösen\,\ensuremath{\nearrow}}}
\end{tcolorbox}
\vspace{2mm}
\noindent\begin{minipage}[t]{36mm}\vspace{0pt}\sffamily\small\raggedright
\hmfigure{längste eindeutige}{{\Large h\#6}}\hmfigure{eindeutig bei dieser Länge}{{\Large 352} Stellungen}\hmfigure{längstes Matt}{{\Large h\#7.5}}\hmfigure{in der Klasse}{15,5 Milliarden Stellungen,\\ 4,11 Milliarden mit Hilfsmatt}
\end{minipage}\hfill
\begin{minipage}[t]{80mm}\vspace{0pt}\begin{tikzpicture}\begin{axis}[width=72mm,height=42mm,ybar,bar width=1.6pt,xmin=-1,xmax=16,ymin=0,ymax=921.4,axis y line*=left,axis x line*=bottom,xlabel={\scriptsize Halbzüge bis Matt},ylabel={\scriptsize Stellungen (Mio.)},tick label style={font=\tiny},label style={font=\scriptsize},enlarge x limits=false,ylabel shift=-3pt,xlabel shift=-3pt]\addplot[fill=ink!55,draw=none] coordinates {(0,1.605) (1,12.82) (2,43.38) (3,177) (4,382.6) (5,675) (6,837.7) (7,800.1) (8,651.4) (9,344.1) (10,147.8) (11,28.72) (12,10.18) (13,0.08909) (14,0.0116) (15,3.8e-05)};\draw[densely dotted] (axis cs:12,0) -- (axis cs:12,921.4);\end{axis}\begin{axis}[width=72mm,height=42mm,axis y line*=right,axis x line=none,ymode=log,xmin=-1,xmax=16,ylabel={\scriptsize eindeutig (log)},tick label style={font=\tiny},label style={font=\scriptsize},ylabel shift=-3pt]\addplot[mark=*,mark size=0.7pt,draw=black,thin] coordinates {(0,1604786) (1,10364405) (2,19955114) (3,37534714) (4,39948380) (5,34457995) (6,21364736) (7,12244873) (8,4838575) (9,714065) (10,103423) (11,1230) (12,352)};\end{axis}\end{tikzpicture}\end{minipage}
\vspace{3mm}

\noindent\begin{minipage}[t]{60mm}\vspace{0pt}\sbox\hmboard{\chessboard[setfen=8/8/8/7k/1n4Q1/b7/2b5/K7 b - - 0 1]}\usebox\hmboard\\[1mm]
\makebox[\wd\hmboard]{\sffamily\hspace{3mm}\textbf{h\#6}\hfill(2+4)}\end{minipage}\hfill
\begin{minipage}[t]{52mm}\vspace{0pt}
\textcolor{ink}{\sffamily\textsc{Lösung}}\par\smallskip\small
\begin{tabular}{@{}r@{\enspace}l@{\enspace}l@{}}
\textbf{1.} & \symking{}h6 & \symqueen{}xb4\\
\textbf{2.} & \symking{}g5 & \symqueen{}xa3\\
\textbf{3.} & \symking{}f4 & \symqueen{}f8+\\
\textbf{4.} & \symking{}e3 & \symking{}b2\\
\textbf{5.} & \symking{}e2 & \symking{}c3\\
\textbf{6.} & \symking{}d1 & \symqueen{}f1\#\\
\end{tabular}\par\vspace{3mm}
\qrcode[height=12mm]{https://helpman.komtera.lt/?fen=8/8/8/7k/1n4Q1/b7/2b5/K7&moves=6}\\{\sffamily\tiny\color{ink}scannen zum Lösen}
\end{minipage}
\vspace{2mm}

{\raggedright\noindent\textcolor{ink}{\sffamily\textsc{Themen}}\enspace \themetag{ein Stein zieht}\par\smallskip{\footnotesize\textit{Hinweis:} die Partei am Zug steht im Diagramm im Schach.}\par}

\clearpage
\phantomsection\label{hm:KQvkbbp}
\begin{tcolorbox}[enhanced,colback=ink,colframe=ink,arc=0pt,boxrule=0pt,left=4mm,right=4mm,top=2.5mm,bottom=2.5mm]
\color{white}\sffamily{\LARGE\bfseries KQvkbbp}\hfill{\Large h\#6}\\[0.5mm]
{\small 6 Steine\quad Schwarz am Zug}\hfill{\small Nr.~698\qquad\href{https://helpman.komtera.lt/?fen=8/8/7k/6Q1/p7/8/b7/K1b5&moves=6}{\color{white}\bfseries lösen\,\ensuremath{\nearrow}}}
\end{tcolorbox}
\vspace{2mm}
\noindent\begin{minipage}[t]{36mm}\vspace{0pt}\sffamily\small\raggedright
\hmfigure{längste eindeutige}{{\Large h\#6}}\hmfigure{eindeutig bei dieser Länge}{{\Large 131} Stellungen}\hmfigure{längstes Matt}{{\Large h\#7}}\hmfigure{in der Klasse}{45,4 Milliarden Stellungen,\\ 12,7 Milliarden mit Hilfsmatt}
\end{minipage}\hfill
\begin{minipage}[t]{80mm}\vspace{0pt}\begin{tikzpicture}\begin{axis}[width=72mm,height=42mm,ybar,bar width=1.6pt,xmin=-1,xmax=15,ymin=0,ymax=3.196,axis y line*=left,axis x line*=bottom,xlabel={\scriptsize Halbzüge bis Matt},ylabel={\scriptsize Stellungen (Mrd.)},tick label style={font=\tiny},label style={font=\scriptsize},enlarge x limits=false,ylabel shift=-3pt,xlabel shift=-3pt]\addplot[fill=ink!55,draw=none] coordinates {(0,0.006176) (1,0.04846) (2,0.1641) (3,0.6461) (4,1.516) (5,2.34) (6,2.905) (7,2.329) (8,1.725) (9,0.7374) (10,0.2551) (11,0.03599) (12,0.008545) (13,3.337e-05) (14,3.052e-06)};\draw[densely dotted] (axis cs:12,0) -- (axis cs:12,3.196);\end{axis}\begin{axis}[width=72mm,height=42mm,axis y line*=right,axis x line=none,ymode=log,xmin=-1,xmax=15,ylabel={\scriptsize eindeutig (log)},tick label style={font=\tiny},label style={font=\scriptsize},ylabel shift=-3pt]\addplot[mark=*,mark size=0.7pt,draw=black,thin] coordinates {(0,6175533) (1,38702757) (2,79567712) (3,131559820) (4,150211537) (5,101149540) (6,56879414) (7,27485707) (8,8953886) (9,1041186) (10,136906) (11,632) (12,131)};\end{axis}\end{tikzpicture}\end{minipage}
\vspace{3mm}

\noindent\begin{minipage}[t]{60mm}\vspace{0pt}\sbox\hmboard{\chessboard[setfen=8/8/7k/6Q1/p7/8/b7/K1b5 b - - 0 1]}\usebox\hmboard\\[1mm]
\makebox[\wd\hmboard]{\sffamily\hspace{3mm}\textbf{h\#6}\hfill(2+4)}\end{minipage}\hfill
\begin{minipage}[t]{52mm}\vspace{0pt}
\textcolor{ink}{\sffamily\textsc{Lösung}}\par\smallskip\small
\begin{tabular}{@{}r@{\enspace}l@{\enspace}l@{}}
\textbf{1.} & \symking{}h7 & \symking{}xa2\\
\textbf{2.} & a3 & \symking{}b3\\
\textbf{3.} & \symbishop{}b2 & \symking{}c4\\
\textbf{4.} & \symbishop{}g7 & \symking{}d5\\
\textbf{5.} & \symking{}g8 & \symking{}e6\\
\textbf{6.} & \symking{}f8 & \symqueen{}d8\#\\
\end{tabular}\par\vspace{3mm}
\qrcode[height=12mm]{https://helpman.komtera.lt/?fen=8/8/7k/6Q1/p7/8/b7/K1b5&moves=6}\\{\sffamily\tiny\color{ink}scannen zum Lösen}
\end{minipage}
\vspace{2mm}

{\raggedright\noindent\textcolor{ink}{\sffamily\textsc{Themen}}\enspace \themetag{Selbstblock} \themetag{ohne Schach}\par\smallskip{\footnotesize\textit{Hinweis:} die Partei am Zug steht im Diagramm im Schach.}\par}

\clearpage
\phantomsection\label{hm:KQvkbnn}
\begin{tcolorbox}[enhanced,colback=ink,colframe=ink,arc=0pt,boxrule=0pt,left=4mm,right=4mm,top=2.5mm,bottom=2.5mm]
\color{white}\sffamily{\LARGE\bfseries KQvkbnn}\hfill{\Large h\#6}\\[0.5mm]
{\small 6 Steine\quad Schwarz am Zug}\hfill{\small Nr.~699\qquad\href{https://helpman.komtera.lt/?fen=8/8/8/7k/6Q1/3n4/2b5/K2n4&moves=6}{\color{white}\bfseries lösen\,\ensuremath{\nearrow}}}
\end{tcolorbox}
\vspace{2mm}
\noindent\begin{minipage}[t]{36mm}\vspace{0pt}\sffamily\small\raggedright
\hmfigure{längste eindeutige}{{\Large h\#6}}\hmfigure{eindeutig bei dieser Länge}{{\Large 310} Stellungen}\hmfigure{längstes Matt}{{\Large h\#7.5}}\hmfigure{in der Klasse}{15,5 Milliarden Stellungen,\\ 4,24 Milliarden mit Hilfsmatt}
\end{minipage}\hfill
\begin{minipage}[t]{80mm}\vspace{0pt}\begin{tikzpicture}\begin{axis}[width=72mm,height=42mm,ybar,bar width=1.6pt,xmin=-1,xmax=16,ymin=0,ymax=908.2,axis y line*=left,axis x line*=bottom,xlabel={\scriptsize Halbzüge bis Matt},ylabel={\scriptsize Stellungen (Mio.)},tick label style={font=\tiny},label style={font=\scriptsize},enlarge x limits=false,ylabel shift=-3pt,xlabel shift=-3pt]\addplot[fill=ink!55,draw=none] coordinates {(0,1.853) (1,14.02) (2,43.29) (3,167.1) (4,348.4) (5,617.9) (6,825.7) (7,805.7) (8,753.8) (9,391.6) (10,205.3) (11,41.56) (12,19.71) (13,0.1406) (14,0.03229) (15,4.3e-05)};\draw[densely dotted] (axis cs:12,0) -- (axis cs:12,908.2);\end{axis}\begin{axis}[width=72mm,height=42mm,axis y line*=right,axis x line=none,ymode=log,xmin=-1,xmax=16,ylabel={\scriptsize eindeutig (log)},tick label style={font=\tiny},label style={font=\scriptsize},ylabel shift=-3pt]\addplot[mark=*,mark size=0.7pt,draw=black,thin] coordinates {(0,1852777) (1,10839985) (2,18629408) (3,35648058) (4,39748140) (5,36634350) (6,26913085) (7,13879448) (8,6545115) (9,865576) (10,177835) (11,1289) (12,310)};\end{axis}\end{tikzpicture}\end{minipage}
\vspace{3mm}

\noindent\begin{minipage}[t]{60mm}\vspace{0pt}\sbox\hmboard{\chessboard[setfen=8/8/8/7k/6Q1/3n4/2b5/K2n4 b - - 0 1]}\usebox\hmboard\\[1mm]
\makebox[\wd\hmboard]{\sffamily\hspace{3mm}\textbf{h\#6}\hfill(2+4)}\end{minipage}\hfill
\begin{minipage}[t]{52mm}\vspace{0pt}
\textcolor{ink}{\sffamily\textsc{Lösung}}\par\smallskip\small
\begin{tabular}{@{}r@{\enspace}l@{\enspace}l@{}}
\textbf{1.} & \symking{}h6 & \symqueen{}xd1\\
\textbf{2.} & \symking{}g5 & \symqueen{}xd3\\
\textbf{3.} & \symking{}f4 & \symqueen{}f5+\\
\textbf{4.} & \symking{}e3 & \symking{}b2\\
\textbf{5.} & \symking{}e2 & \symking{}c3\\
\textbf{6.} & \symking{}d1 & \symqueen{}f1\#\\
\end{tabular}\par\vspace{3mm}
\qrcode[height=12mm]{https://helpman.komtera.lt/?fen=8/8/8/7k/6Q1/3n4/2b5/K2n4&moves=6}\\{\sffamily\tiny\color{ink}scannen zum Lösen}
\end{minipage}
\vspace{2mm}

{\raggedright\noindent\textcolor{ink}{\sffamily\textsc{Themen}}\enspace \themetag{ein Stein zieht} \themetag{Kniest}\par\smallskip{\footnotesize\textit{Hinweis:} die Partei am Zug steht im Diagramm im Schach.}\par}

\clearpage
\phantomsection\label{hm:KQvkbnp}
\begin{tcolorbox}[enhanced,colback=ink,colframe=ink,arc=0pt,boxrule=0pt,left=4mm,right=4mm,top=2.5mm,bottom=2.5mm]
\color{white}\sffamily{\LARGE\bfseries KQvkbnp}\hfill{\Large h\#6}\\[0.5mm]
{\small 6 Steine\quad Schwarz am Zug}\hfill{\small Nr.~700\qquad\href{https://helpman.komtera.lt/?fen=8/8/8/7k/3p2Q1/3n4/2b5/K7&moves=6}{\color{white}\bfseries lösen\,\ensuremath{\nearrow}}}
\end{tcolorbox}
\vspace{2mm}
\noindent\begin{minipage}[t]{36mm}\vspace{0pt}\sffamily\small\raggedright
\hmfigure{längste eindeutige}{{\Large h\#6}}\hmfigure{eindeutig bei dieser Länge}{{\Large 171} Stellungen}\hmfigure{längstes Matt}{{\Large h\#7.5}}\hmfigure{in der Klasse}{45,4 Milliarden Stellungen,\\ 26,2 Milliarden mit Hilfsmatt}
\end{minipage}\hfill
\begin{minipage}[t]{80mm}\vspace{0pt}\begin{tikzpicture}\begin{axis}[width=72mm,height=42mm,ybar,bar width=1.6pt,xmin=-1,xmax=16,ymin=0,ymax=6.466,axis y line*=left,axis x line*=bottom,xlabel={\scriptsize Halbzüge bis Matt},ylabel={\scriptsize Stellungen (Mrd.)},tick label style={font=\tiny},label style={font=\scriptsize},enlarge x limits=false,ylabel shift=-3pt,xlabel shift=-3pt]\addplot[fill=ink!55,draw=none] coordinates {(0,0.01432) (1,0.106) (2,0.323) (3,1.201) (4,2.746) (5,4.252) (6,5.878) (7,4.826) (8,4.185) (9,1.769) (10,0.7869) (11,0.1209) (12,0.03957) (13,8.582e-05) (14,1.635e-05) (15,2.6e-08)};\draw[densely dotted] (axis cs:12,0) -- (axis cs:12,6.466);\end{axis}\begin{axis}[width=72mm,height=42mm,axis y line*=right,axis x line=none,ymode=log,xmin=-1,xmax=16,ylabel={\scriptsize eindeutig (log)},tick label style={font=\tiny},label style={font=\scriptsize},ylabel shift=-3pt]\addplot[mark=*,mark size=0.7pt,draw=black,thin] coordinates {(0,14317535) (1,80946000) (2,146681195) (3,243480063) (4,294672092) (5,215105870) (6,143595334) (7,63781826) (8,25972199) (9,2878666) (10,531924) (11,1019) (12,171)};\end{axis}\end{tikzpicture}\end{minipage}
\vspace{3mm}

\noindent\begin{minipage}[t]{60mm}\vspace{0pt}\sbox\hmboard{\chessboard[setfen=8/8/8/7k/3p2Q1/3n4/2b5/K7 b - - 0 1]}\usebox\hmboard\\[1mm]
\makebox[\wd\hmboard]{\sffamily\hspace{3mm}\textbf{h\#6}\hfill(2+4)}\end{minipage}\hfill
\begin{minipage}[t]{52mm}\vspace{0pt}
\textcolor{ink}{\sffamily\textsc{Lösung}}\par\smallskip\small
\begin{tabular}{@{}r@{\enspace}l@{\enspace}l@{}}
\textbf{1.} & \symking{}h6 & \symqueen{}xd4\\
\textbf{2.} & \symking{}g5 & \symqueen{}xd3\\
\textbf{3.} & \symking{}f4 & \symqueen{}f5+\\
\textbf{4.} & \symking{}e3 & \symking{}b2\\
\textbf{5.} & \symking{}e2 & \symking{}c3\\
\textbf{6.} & \symking{}d1 & \symqueen{}f1\#\\
\end{tabular}\par\vspace{3mm}
\qrcode[height=12mm]{https://helpman.komtera.lt/?fen=8/8/8/7k/3p2Q1/3n4/2b5/K7&moves=6}\\{\sffamily\tiny\color{ink}scannen zum Lösen}
\end{minipage}
\vspace{2mm}

{\raggedright\noindent\textcolor{ink}{\sffamily\textsc{Themen}}\enspace \themetag{ein Stein zieht}\par\smallskip{\footnotesize\textit{Hinweis:} die Partei am Zug steht im Diagramm im Schach.}\par}

\clearpage
\phantomsection\label{hm:KQvkbpp}
\begin{tcolorbox}[enhanced,colback=ink,colframe=ink,arc=0pt,boxrule=0pt,left=4mm,right=4mm,top=2.5mm,bottom=2.5mm]
\color{white}\sffamily{\LARGE\bfseries KQvkbpp}\hfill{\Large h\#6}\\[0.5mm]
{\small 6 Steine\quad Schwarz am Zug}\hfill{\small Nr.~701\qquad\href{https://helpman.komtera.lt/?fen=3Q3k/8/8/1p6/8/2p5/2b5/K7&moves=6}{\color{white}\bfseries lösen\,\ensuremath{\nearrow}}}
\end{tcolorbox}
\vspace{2mm}
\noindent\begin{minipage}[t]{36mm}\vspace{0pt}\sffamily\small\raggedright
\hmfigure{längste eindeutige}{{\Large h\#6}}\hmfigure{eindeutig bei dieser Länge}{{\Large 3} Stellungen}\hmfigure{längstes Matt}{{\Large h\#7}}\hmfigure{in der Klasse}{34,1 Milliarden Stellungen,\\ 10,1 Milliarden mit Hilfsmatt}
\end{minipage}\hfill
\begin{minipage}[t]{80mm}\vspace{0pt}\begin{tikzpicture}\begin{axis}[width=72mm,height=42mm,ybar,bar width=1.6pt,xmin=-1,xmax=15,ymin=0,ymax=2.733,axis y line*=left,axis x line*=bottom,xlabel={\scriptsize Halbzüge bis Matt},ylabel={\scriptsize Stellungen (Mrd.)},tick label style={font=\tiny},label style={font=\scriptsize},enlarge x limits=false,ylabel shift=-3pt,xlabel shift=-3pt]\addplot[fill=ink!55,draw=none] coordinates {(0,0.006942) (1,0.05026) (2,0.1516) (3,0.5425) (4,1.336) (5,1.794) (6,2.484) (7,1.717) (8,1.358) (9,0.4711) (10,0.1704) (11,0.01976) (12,0.004286) (13,2.473e-06) (14,1.28e-07)};\draw[densely dotted] (axis cs:12,0) -- (axis cs:12,2.733);\end{axis}\begin{axis}[width=72mm,height=42mm,axis y line*=right,axis x line=none,ymode=log,xmin=-1,xmax=15,ylabel={\scriptsize eindeutig (log)},tick label style={font=\tiny},label style={font=\scriptsize},ylabel shift=-3pt]\addplot[mark=*,mark size=0.7pt,draw=black,thin] coordinates {(0,6942430) (1,37977927) (2,73634398) (3,104782893) (4,135306273) (5,75475436) (6,45625086) (7,17591603) (8,5924551) (9,546118) (10,96019) (11,60) (12,3)};\end{axis}\end{tikzpicture}\end{minipage}
\vspace{3mm}

\noindent\begin{minipage}[t]{60mm}\vspace{0pt}\sbox\hmboard{\chessboard[setfen=3Q3k/8/8/1p6/8/2p5/2b5/K7 b - - 0 1]}\usebox\hmboard\\[1mm]
\makebox[\wd\hmboard]{\sffamily\hspace{3mm}\textbf{h\#6}\hfill(2+4)}\end{minipage}\hfill
\begin{minipage}[t]{52mm}\vspace{0pt}
\textcolor{ink}{\sffamily\textsc{Lösung}}\par\smallskip\small
\begin{tabular}{@{}r@{\enspace}l@{\enspace}l@{}}
\textbf{1.} & \symking{}g7 & \symking{}a2\\
\textbf{2.} & \symbishop{}h7 & \symking{}b3\\
\textbf{3.} & c2 & \symking{}c3\\
\textbf{4.} & c1=\symqueen{}+ & \symking{}d4\\
\textbf{5.} & \symqueen{}c4+ & \symking{}e5\\
\textbf{6.} & \symqueen{}g8 & \symqueen{}f6\#\\
\end{tabular}\par\vspace{3mm}
\qrcode[height=12mm]{https://helpman.komtera.lt/?fen=3Q3k/8/8/1p6/8/2p5/2b5/K7&moves=6}\\{\sffamily\tiny\color{ink}scannen zum Lösen}
\end{minipage}
\vspace{2mm}

{\raggedright\noindent\textcolor{ink}{\sffamily\textsc{Themen}}\enspace \themetag{reines Matt} \themetag{Mustermatt} \themetag{Umwandlung} \themetag{Selbstblock} \themetag{ohne Schlag} \themetag{Umnov}\par\smallskip{\footnotesize\textit{Hinweis:} die Partei am Zug steht im Diagramm im Schach.}\par}

\clearpage
\phantomsection\label{hm:KQvknnn}
\begin{tcolorbox}[enhanced,colback=ink,colframe=ink,arc=0pt,boxrule=0pt,left=4mm,right=4mm,top=2.5mm,bottom=2.5mm]
\color{white}\sffamily{\LARGE\bfseries KQvknnn}\hfill{\Large h\#6}\\[0.5mm]
{\small 6 Steine\quad Schwarz am Zug}\hfill{\small Nr.~702\qquad\href{https://helpman.komtera.lt/?fen=8/8/8/7k/2n3Q1/n2n4/8/K7&moves=6}{\color{white}\bfseries lösen\,\ensuremath{\nearrow}}}
\end{tcolorbox}
\vspace{2mm}
\noindent\begin{minipage}[t]{36mm}\vspace{0pt}\sffamily\small\raggedright
\hmfigure{längste eindeutige}{{\Large h\#6}}\hmfigure{eindeutig bei dieser Länge}{{\Large 292} Stellungen}\hmfigure{längstes Matt}{{\Large h\#7.5}}\hmfigure{in der Klasse}{15,5 Milliarden Stellungen,\\ 1,45 Milliarden mit Hilfsmatt}
\end{minipage}\hfill
\begin{minipage}[t]{80mm}\vspace{0pt}\begin{tikzpicture}\begin{axis}[width=72mm,height=42mm,ybar,bar width=1.6pt,xmin=-1,xmax=16,ymin=0,ymax=316.9,axis y line*=left,axis x line*=bottom,xlabel={\scriptsize Halbzüge bis Matt},ylabel={\scriptsize Stellungen (Mio.)},tick label style={font=\tiny},label style={font=\scriptsize},enlarge x limits=false,ylabel shift=-3pt,xlabel shift=-3pt]\addplot[fill=ink!55,draw=none] coordinates {(0,0.7171) (1,5.094) (2,14.42) (3,52.35) (4,103.3) (5,198.2) (6,288.1) (7,274.7) (8,277.8) (9,134.4) (10,76.35) (11,14.57) (12,7.362) (13,0.06875) (14,0.0196) (15,3e-06)};\draw[densely dotted] (axis cs:12,0) -- (axis cs:12,316.9);\end{axis}\begin{axis}[width=72mm,height=42mm,axis y line*=right,axis x line=none,ymode=log,xmin=-1,xmax=16,ylabel={\scriptsize eindeutig (log)},tick label style={font=\tiny},label style={font=\scriptsize},ylabel shift=-3pt]\addplot[mark=*,mark size=0.7pt,draw=black,thin] coordinates {(0,717052) (1,3704099) (2,5702807) (3,11921383) (4,13488391) (5,13037969) (6,10819013) (7,4736252) (8,2470975) (9,263395) (10,58952) (11,730) (12,292)};\end{axis}\end{tikzpicture}\end{minipage}
\vspace{3mm}

\noindent\begin{minipage}[t]{60mm}\vspace{0pt}\sbox\hmboard{\chessboard[setfen=8/8/8/7k/2n3Q1/n2n4/8/K7 b - - 0 1]}\usebox\hmboard\\[1mm]
\makebox[\wd\hmboard]{\sffamily\hspace{3mm}\textbf{h\#6}\hfill(2+4)}\end{minipage}\hfill
\begin{minipage}[t]{52mm}\vspace{0pt}
\textcolor{ink}{\sffamily\textsc{Lösung}}\par\smallskip\small
\begin{tabular}{@{}r@{\enspace}l@{\enspace}l@{}}
\textbf{1.} & \symking{}h6 & \symqueen{}xc4\\
\textbf{2.} & \symking{}g5 & \symqueen{}xd3\\
\textbf{3.} & \symking{}f4 & \symqueen{}f5+\\
\textbf{4.} & \symking{}e3 & \symking{}b2\\
\textbf{5.} & \symking{}e2 & \symking{}c3\\
\textbf{6.} & \symking{}d1 & \symqueen{}f1\#\\
\end{tabular}\par\vspace{3mm}
\qrcode[height=12mm]{https://helpman.komtera.lt/?fen=8/8/8/7k/2n3Q1/n2n4/8/K7&moves=6}\\{\sffamily\tiny\color{ink}scannen zum Lösen}
\end{minipage}
\vspace{2mm}

{\raggedright\noindent\textcolor{ink}{\sffamily\textsc{Themen}}\enspace \themetag{reines Matt} \themetag{Mustermatt} \themetag{Spiegelmatt} \themetag{ein Stein zieht}\par\smallskip{\footnotesize\textit{Hinweis:} die Partei am Zug steht im Diagramm im Schach.}\par}

\clearpage
\phantomsection\label{hm:KQvknnp}
\begin{tcolorbox}[enhanced,colback=ink,colframe=ink,arc=0pt,boxrule=0pt,left=4mm,right=4mm,top=2.5mm,bottom=2.5mm]
\color{white}\sffamily{\LARGE\bfseries KQvknnp}\hfill{\Large h\#6}\\[0.5mm]
{\small 6 Steine\quad Schwarz am Zug}\hfill{\small Nr.~703\qquad\href{https://helpman.komtera.lt/?fen=K7/2p5/3n4/n7/6Q1/8/7k/8&moves=6}{\color{white}\bfseries lösen\,\ensuremath{\nearrow}}}
\end{tcolorbox}
\vspace{2mm}
\noindent\begin{minipage}[t]{36mm}\vspace{0pt}\sffamily\small\raggedright
\hmfigure{längste eindeutige}{{\Large h\#6}}\hmfigure{eindeutig bei dieser Länge}{{\Large 1\,102} Stellungen}\hmfigure{längstes Matt}{{\Large h\#7.5}}\hmfigure{in der Klasse}{45,4 Milliarden Stellungen,\\ 13,5 Milliarden mit Hilfsmatt}
\end{minipage}\hfill
\begin{minipage}[t]{80mm}\vspace{0pt}\begin{tikzpicture}\begin{axis}[width=72mm,height=42mm,ybar,bar width=1.6pt,xmin=-1,xmax=16,ymin=0,ymax=3.19,axis y line*=left,axis x line*=bottom,xlabel={\scriptsize Halbzüge bis Matt},ylabel={\scriptsize Stellungen (Mrd.)},tick label style={font=\tiny},label style={font=\scriptsize},enlarge x limits=false,ylabel shift=-3pt,xlabel shift=-3pt]\addplot[fill=ink!55,draw=none] coordinates {(0,0.008343) (1,0.05775) (2,0.1577) (3,0.5483) (4,1.163) (5,1.919) (6,2.9) (7,2.47) (8,2.502) (9,1.048) (10,0.5691) (11,0.09432) (12,0.04029) (13,0.0001432) (14,4.532e-05) (15,5.6e-08)};\draw[densely dotted] (axis cs:12,0) -- (axis cs:12,3.19);\end{axis}\begin{axis}[width=72mm,height=42mm,axis y line*=right,axis x line=none,ymode=log,xmin=-1,xmax=16,ylabel={\scriptsize eindeutig (log)},tick label style={font=\tiny},label style={font=\scriptsize},ylabel shift=-3pt]\addplot[mark=*,mark size=0.7pt,draw=black,thin] coordinates {(0,8343014) (1,41492049) (2,65200186) (3,117397716) (4,140077718) (5,114804132) (6,94540250) (7,38284985) (8,19318073) (9,1858294) (10,336267) (11,1948) (12,1102)};\end{axis}\end{tikzpicture}\end{minipage}
\vspace{3mm}

\noindent\begin{minipage}[t]{60mm}\vspace{0pt}\sbox\hmboard{\chessboard[setfen=K7/2p5/3n4/n7/6Q1/8/7k/8 b - - 0 1]}\usebox\hmboard\\[1mm]
\makebox[\wd\hmboard]{\sffamily\hspace{3mm}\textbf{h\#6}\hfill(2+4)}\end{minipage}\hfill
\begin{minipage}[t]{52mm}\vspace{0pt}
\textcolor{ink}{\sffamily\textsc{Lösung}}\par\smallskip\small
\begin{tabular}{@{}r@{\enspace}l@{\enspace}l@{}}
\textbf{1.} & c5 & \symking{}b8\\
\textbf{2.} & c4 & \symking{}c7\\
\textbf{3.} & c3 & \symking{}xd6\\
\textbf{4.} & c2 & \symking{}e5\\
\textbf{5.} & c1=\symrook{} & \symking{}f4\\
\textbf{6.} & \symrook{}h1 & \symqueen{}g3\#\\
\end{tabular}\par\vspace{3mm}
\qrcode[height=12mm]{https://helpman.komtera.lt/?fen=K7/2p5/3n4/n7/6Q1/8/7k/8&moves=6}\\{\sffamily\tiny\color{ink}scannen zum Lösen}
\end{minipage}
\vspace{2mm}

{\raggedright\noindent\textcolor{ink}{\sffamily\textsc{Themen}}\enspace \themetag{reines Matt} \themetag{Mustermatt} \themetag{Umwandlung} \themetag{Unterverwandlung} \themetag{Excelsior} \themetag{Selbstblock} \themetag{ein Stein zieht} \themetag{ohne Schach}\par}

\clearpage
\phantomsection\label{hm:KQvknpp}
\begin{tcolorbox}[enhanced,colback=ink,colframe=ink,arc=0pt,boxrule=0pt,left=4mm,right=4mm,top=2.5mm,bottom=2.5mm]
\color{white}\sffamily{\LARGE\bfseries KQvknpp}\hfill{\Large h\#6}\\[0.5mm]
{\small 6 Steine\quad Schwarz am Zug}\hfill{\small Nr.~704\qquad\href{https://helpman.komtera.lt/?fen=8/7k/4pQ1p/8/8/8/n7/K7&moves=6}{\color{white}\bfseries lösen\,\ensuremath{\nearrow}}}
\end{tcolorbox}
\vspace{2mm}
\noindent\begin{minipage}[t]{36mm}\vspace{0pt}\sffamily\small\raggedright
\hmfigure{längste eindeutige}{{\Large h\#6}}\hmfigure{eindeutig bei dieser Länge}{{\Large 163} Stellungen}\hmfigure{längstes Matt}{{\Large h\#7}}\hmfigure{in der Klasse}{34,1 Milliarden Stellungen,\\ 10,4 Milliarden mit Hilfsmatt}
\end{minipage}\hfill
\begin{minipage}[t]{80mm}\vspace{0pt}\begin{tikzpicture}\begin{axis}[width=72mm,height=42mm,ybar,bar width=1.6pt,xmin=-1,xmax=15,ymin=0,ymax=2.7,axis y line*=left,axis x line*=bottom,xlabel={\scriptsize Halbzüge bis Matt},ylabel={\scriptsize Stellungen (Mrd.)},tick label style={font=\tiny},label style={font=\scriptsize},enlarge x limits=false,ylabel shift=-3pt,xlabel shift=-3pt]\addplot[fill=ink!55,draw=none] coordinates {(0,0.008122) (1,0.05482) (2,0.1457) (3,0.4869) (4,1.12) (5,1.578) (6,2.455) (7,1.811) (8,1.758) (9,0.6218) (10,0.3039) (11,0.0426) (12,0.01379) (13,1.319e-05) (14,4.656e-06)};\draw[densely dotted] (axis cs:12,0) -- (axis cs:12,2.7);\end{axis}\begin{axis}[width=72mm,height=42mm,axis y line*=right,axis x line=none,ymode=log,xmin=-1,xmax=15,ylabel={\scriptsize eindeutig (log)},tick label style={font=\tiny},label style={font=\scriptsize},ylabel shift=-3pt]\addplot[mark=*,mark size=0.7pt,draw=black,thin] coordinates {(0,8121956) (1,39000876) (2,64427408) (3,97486243) (4,123989373) (5,81352827) (6,62337803) (7,22887736) (8,10405498) (9,886961) (10,139268) (11,306) (12,163)};\end{axis}\end{tikzpicture}\end{minipage}
\vspace{3mm}

\noindent\begin{minipage}[t]{60mm}\vspace{0pt}\sbox\hmboard{\chessboard[setfen=8/7k/4pQ1p/8/8/8/n7/K7 b - - 0 1]}\usebox\hmboard\\[1mm]
\makebox[\wd\hmboard]{\sffamily\hspace{3mm}\textbf{h\#6}\hfill(2+4)}\end{minipage}\hfill
\begin{minipage}[t]{52mm}\vspace{0pt}
\textcolor{ink}{\sffamily\textsc{Lösung}}\par\smallskip\small
\begin{tabular}{@{}r@{\enspace}l@{\enspace}l@{}}
\textbf{1.} & \symknight{}b4 & \symking{}b2\\
\textbf{2.} & \symknight{}c6 & \symking{}c3\\
\textbf{3.} & \symknight{}d8 & \symking{}d4\\
\textbf{4.} & \symking{}g8 & \symking{}e5\\
\textbf{5.} & \symknight{}f7+ & \symking{}xe6\\
\textbf{6.} & \symking{}f8 & \symqueen{}xf7\#\\
\end{tabular}\par\vspace{3mm}
\qrcode[height=12mm]{https://helpman.komtera.lt/?fen=8/7k/4pQ1p/8/8/8/n7/K7&moves=6}\\{\sffamily\tiny\color{ink}scannen zum Lösen}
\end{minipage}
\vspace{2mm}

{\raggedright\noindent\textcolor{ink}{\sffamily\textsc{Themen}}\enspace \textit{keine erkannt}\par}

\clearpage
\phantomsection\label{hm:KQvkppp}
\begin{tcolorbox}[enhanced,colback=ink,colframe=ink,arc=0pt,boxrule=0pt,left=4mm,right=4mm,top=2.5mm,bottom=2.5mm]
\color{white}\sffamily{\LARGE\bfseries KQvkppp}\hfill{\Large h\#6}\\[0.5mm]
{\small 6 Steine\quad Schwarz am Zug}\hfill{\small Nr.~705\qquad\href{https://helpman.komtera.lt/?fen=8/7k/2p2Q2/p4p2/8/8/8/K7&moves=6}{\color{white}\bfseries lösen\,\ensuremath{\nearrow}}}
\end{tcolorbox}
\vspace{2mm}
\noindent\begin{minipage}[t]{36mm}\vspace{0pt}\sffamily\small\raggedright
\hmfigure{längste eindeutige}{{\Large h\#6}}\hmfigure{eindeutig bei dieser Länge}{{\Large 35} Stellungen}\hmfigure{längstes Matt}{{\Large h\#7}}\hmfigure{in der Klasse}{25,6 Milliarden Stellungen,\\ 2,66 Milliarden mit Hilfsmatt}
\end{minipage}\hfill
\begin{minipage}[t]{80mm}\vspace{0pt}\begin{tikzpicture}\begin{axis}[width=72mm,height=42mm,ybar,bar width=1.6pt,xmin=-1,xmax=15,ymin=0,ymax=747.2,axis y line*=left,axis x line*=bottom,xlabel={\scriptsize Halbzüge bis Matt},ylabel={\scriptsize Stellungen (Mio.)},tick label style={font=\tiny},label style={font=\scriptsize},enlarge x limits=false,ylabel shift=-3pt,xlabel shift=-3pt]\addplot[fill=ink!55,draw=none] coordinates {(0,2.643) (1,17.39) (2,45.08) (3,145.2) (4,360.4) (5,432.8) (6,679.3) (7,426.9) (8,382.9) (9,112.7) (10,46.78) (11,5.436) (12,1.303) (13,0.001466) (14,0.00077)};\draw[densely dotted] (axis cs:12,0) -- (axis cs:12,747.2);\end{axis}\begin{axis}[width=72mm,height=42mm,axis y line*=right,axis x line=none,ymode=log,xmin=-1,xmax=15,ylabel={\scriptsize eindeutig (log)},tick label style={font=\tiny},label style={font=\scriptsize},ylabel shift=-3pt]\addplot[mark=*,mark size=0.7pt,draw=black,thin] coordinates {(0,2642923) (1,12260867) (2,22227473) (3,27538699) (4,37308213) (5,18449430) (6,12471300) (7,4104315) (8,1553203) (9,98410) (10,14981) (11,56) (12,35)};\end{axis}\end{tikzpicture}\end{minipage}
\vspace{3mm}

\noindent\begin{minipage}[t]{60mm}\vspace{0pt}\sbox\hmboard{\chessboard[setfen=8/7k/2p2Q2/p4p2/8/8/8/K7 b - - 0 1]}\usebox\hmboard\\[1mm]
\makebox[\wd\hmboard]{\sffamily\hspace{3mm}\textbf{h\#6}\hfill(2+4)}\end{minipage}\hfill
\begin{minipage}[t]{52mm}\vspace{0pt}
\textcolor{ink}{\sffamily\textsc{Lösung}}\par\smallskip\small
\begin{tabular}{@{}r@{\enspace}l@{\enspace}l@{}}
\textbf{1.} & a4 & \symking{}b2\\
\textbf{2.} & a3+ & \symking{}c3\\
\textbf{3.} & a2 & \symking{}d4\\
\textbf{4.} & a1=\symrook{} & \symking{}e5\\
\textbf{5.} & \symrook{}a8 & \symking{}xf5\\
\textbf{6.} & \symrook{}h8 & \symqueen{}g6\#\\
\end{tabular}\par\vspace{3mm}
\qrcode[height=12mm]{https://helpman.komtera.lt/?fen=8/7k/2p2Q2/p4p2/8/8/8/K7&moves=6}\\{\sffamily\tiny\color{ink}scannen zum Lösen}
\end{minipage}
\vspace{2mm}

{\raggedright\noindent\textcolor{ink}{\sffamily\textsc{Themen}}\enspace \themetag{reines Matt} \themetag{Mustermatt} \themetag{Umwandlung} \themetag{Unterverwandlung} \themetag{Selbstblock} \themetag{ein Stein zieht}\par}

\clearpage
\phantomsection\label{hm:KQvkqbb}
\begin{tcolorbox}[enhanced,colback=ink,colframe=ink,arc=0pt,boxrule=0pt,left=4mm,right=4mm,top=2.5mm,bottom=2.5mm]
\color{white}\sffamily{\LARGE\bfseries KQvkqbb}\hfill{\Large h\#6}\\[0.5mm]
{\small 6 Steine\quad Schwarz am Zug}\hfill{\small Nr.~706\qquad\href{https://helpman.komtera.lt/?fen=8/7k/6Q1/8/8/8/8/K1b1qb2&moves=6}{\color{white}\bfseries lösen\,\ensuremath{\nearrow}}}
\end{tcolorbox}
\vspace{2mm}
\noindent\begin{minipage}[t]{36mm}\vspace{0pt}\sffamily\small\raggedright
\hmfigure{längste eindeutige}{{\Large h\#6}}\hmfigure{eindeutig bei dieser Länge}{{\Large 12} Stellungen}\hmfigure{längstes Matt}{{\Large h\#7}}\hmfigure{in der Klasse}{15,5 Milliarden Stellungen,\\ 3,52 Milliarden mit Hilfsmatt}
\end{minipage}\hfill
\begin{minipage}[t]{80mm}\vspace{0pt}\begin{tikzpicture}\begin{axis}[width=72mm,height=42mm,ybar,bar width=1.6pt,xmin=-1,xmax=15,ymin=0,ymax=890.2,axis y line*=left,axis x line*=bottom,xlabel={\scriptsize Halbzüge bis Matt},ylabel={\scriptsize Stellungen (Mio.)},tick label style={font=\tiny},label style={font=\scriptsize},enlarge x limits=false,ylabel shift=-3pt,xlabel shift=-3pt]\addplot[fill=ink!55,draw=none] coordinates {(0,0.7824) (1,7.183) (2,30.22) (3,150.1) (4,314) (5,709) (6,650.8) (7,809.3) (8,426.8) (9,328) (10,75.65) (11,20.03) (12,1.251) (13,0.1162) (14,0.000975)};\draw[densely dotted] (axis cs:12,0) -- (axis cs:12,890.2);\end{axis}\begin{axis}[width=72mm,height=42mm,axis y line*=right,axis x line=none,ymode=log,xmin=-1,xmax=15,ylabel={\scriptsize eindeutig (log)},tick label style={font=\tiny},label style={font=\scriptsize},ylabel shift=-3pt]\addplot[mark=*,mark size=0.7pt,draw=black,thin] coordinates {(0,782381) (1,6365794) (2,15000987) (3,36246528) (4,31862981) (5,34023130) (6,13890156) (7,12200118) (8,3071661) (9,906456) (10,55340) (11,2272) (12,12)};\end{axis}\end{tikzpicture}\end{minipage}
\vspace{3mm}

\noindent\begin{minipage}[t]{60mm}\vspace{0pt}\sbox\hmboard{\chessboard[setfen=8/7k/6Q1/8/8/8/8/K1b1qb2 b - - 0 1]}\usebox\hmboard\\[1mm]
\makebox[\wd\hmboard]{\sffamily\hspace{3mm}\textbf{h\#6}\hfill(2+4)}\end{minipage}\hfill
\begin{minipage}[t]{52mm}\vspace{0pt}
\textcolor{ink}{\sffamily\textsc{Lösung}}\par\smallskip\small
\begin{tabular}{@{}r@{\enspace}l@{\enspace}l@{}}
\textbf{1.} & \symking{}h8 & \symqueen{}g4\\
\textbf{2.} & \symbishop{}g5+ & \symking{}b2\\
\textbf{3.} & \symqueen{}e7 & \symking{}c3\\
\textbf{4.} & \symking{}g7 & \symking{}d4\\
\textbf{5.} & \symking{}f6 & \symking{}d5\\
\textbf{6.} & \symqueen{}g7 & \symqueen{}e6\#\\
\end{tabular}\par\vspace{3mm}
\qrcode[height=12mm]{https://helpman.komtera.lt/?fen=8/7k/6Q1/8/8/8/8/K1b1qb2&moves=6}\\{\sffamily\tiny\color{ink}scannen zum Lösen}
\end{minipage}
\vspace{2mm}

{\raggedright\noindent\textcolor{ink}{\sffamily\textsc{Themen}}\enspace \themetag{Selbstblock} \themetag{ohne Schlag}\par\smallskip{\footnotesize\textit{Hinweis:} die Partei am Zug steht im Diagramm im Schach.}\par}

\clearpage
\phantomsection\label{hm:KQvkqbn}
\begin{tcolorbox}[enhanced,colback=ink,colframe=ink,arc=0pt,boxrule=0pt,left=4mm,right=4mm,top=2.5mm,bottom=2.5mm]
\color{white}\sffamily{\LARGE\bfseries KQvkqbn}\hfill{\Large h\#6}\\[0.5mm]
{\small 6 Steine\quad Schwarz am Zug}\hfill{\small Nr.~707\qquad\href{https://helpman.komtera.lt/?fen=8/8/8/3n3k/b5Q1/1q6/8/K7&moves=6}{\color{white}\bfseries lösen\,\ensuremath{\nearrow}}}
\end{tcolorbox}
\vspace{2mm}
\noindent\begin{minipage}[t]{36mm}\vspace{0pt}\sffamily\small\raggedright
\hmfigure{längste eindeutige}{{\Large h\#6}}\hmfigure{eindeutig bei dieser Länge}{{\Large 218} Stellungen}\hmfigure{längstes Matt}{{\Large h\#7.5}}\hmfigure{in der Klasse}{15,5 Milliarden Stellungen,\\ 7,28 Milliarden mit Hilfsmatt}
\end{minipage}\hfill
\begin{minipage}[t]{80mm}\vspace{0pt}\begin{tikzpicture}\begin{axis}[width=72mm,height=42mm,ybar,bar width=1.6pt,xmin=-1,xmax=16,ymin=0,ymax=1.812,axis y line*=left,axis x line*=bottom,xlabel={\scriptsize Halbzüge bis Matt},ylabel={\scriptsize Stellungen (Mrd.)},tick label style={font=\tiny},label style={font=\scriptsize},enlarge x limits=false,ylabel shift=-3pt,xlabel shift=-3pt]\addplot[fill=ink!55,draw=none] coordinates {(0,0.001784) (1,0.01576) (2,0.06224) (3,0.2925) (4,0.5926) (5,1.301) (6,1.339) (7,1.647) (8,0.9981) (9,0.742) (10,0.218) (11,0.05893) (12,0.006158) (13,0.0005255) (14,9.581e-06) (15,1.917e-06)};\draw[densely dotted] (axis cs:12,0) -- (axis cs:12,1.812);\end{axis}\begin{axis}[width=72mm,height=42mm,axis y line*=right,axis x line=none,ymode=log,xmin=-1,xmax=16,ylabel={\scriptsize eindeutig (log)},tick label style={font=\tiny},label style={font=\scriptsize},ylabel shift=-3pt]\addplot[mark=*,mark size=0.7pt,draw=black,thin] coordinates {(0,1783594) (1,13658627) (2,29687809) (3,68168205) (4,62192160) (5,67244348) (6,31977665) (7,25866736) (8,7595600) (9,2185193) (10,166984) (11,6526) (12,218)};\end{axis}\end{tikzpicture}\end{minipage}
\vspace{3mm}

\noindent\begin{minipage}[t]{60mm}\vspace{0pt}\sbox\hmboard{\chessboard[setfen=8/8/8/3n3k/b5Q1/1q6/8/K7 b - - 0 1]}\usebox\hmboard\\[1mm]
\makebox[\wd\hmboard]{\sffamily\hspace{3mm}\textbf{h\#6}\hfill(2+4)}\end{minipage}\hfill
\begin{minipage}[t]{52mm}\vspace{0pt}
\textcolor{ink}{\sffamily\textsc{Lösung}}\par\smallskip\small
\begin{tabular}{@{}r@{\enspace}l@{\enspace}l@{}}
\textbf{1.} & \symking{}h6 & \symqueen{}xa4\\
\textbf{2.} & \symking{}g5 & \symqueen{}xb3\\
\textbf{3.} & \symking{}f4 & \symqueen{}a2\\
\textbf{4.} & \symking{}e3 & \symking{}b2\\
\textbf{5.} & \symking{}d2 & \symking{}b3+\\
\textbf{6.} & \symking{}c1 & \symqueen{}c2\#\\
\end{tabular}\par\vspace{3mm}
\qrcode[height=12mm]{https://helpman.komtera.lt/?fen=8/8/8/3n3k/b5Q1/1q6/8/K7&moves=6}\\{\sffamily\tiny\color{ink}scannen zum Lösen}
\end{minipage}
\vspace{2mm}

{\raggedright\noindent\textcolor{ink}{\sffamily\textsc{Themen}}\enspace \themetag{ein Stein zieht}\par\smallskip{\footnotesize\textit{Hinweis:} die Partei am Zug steht im Diagramm im Schach.}\par}

\clearpage
\phantomsection\label{hm:KQvkqbp}
\begin{tcolorbox}[enhanced,colback=ink,colframe=ink,arc=0pt,boxrule=0pt,left=4mm,right=4mm,top=2.5mm,bottom=2.5mm]
\color{white}\sffamily{\LARGE\bfseries KQvkqbp}\hfill{\Large h\#6}\\[0.5mm]
{\small 6 Steine\quad Schwarz am Zug}\hfill{\small Nr.~708\qquad\href{https://helpman.komtera.lt/?fen=8/8/7k/6Q1/1qp5/b7/8/K7&moves=6}{\color{white}\bfseries lösen\,\ensuremath{\nearrow}}}
\end{tcolorbox}
\vspace{2mm}
\noindent\begin{minipage}[t]{36mm}\vspace{0pt}\sffamily\small\raggedright
\hmfigure{längste eindeutige}{{\Large h\#6}}\hmfigure{eindeutig bei dieser Länge}{{\Large 28} Stellungen}\hmfigure{längstes Matt}{{\Large h\#7.5}}\hmfigure{in der Klasse}{45,4 Milliarden Stellungen,\\ 22,4 Milliarden mit Hilfsmatt}
\end{minipage}\hfill
\begin{minipage}[t]{80mm}\vspace{0pt}\begin{tikzpicture}\begin{axis}[width=72mm,height=42mm,ybar,bar width=1.6pt,xmin=-1,xmax=16,ymin=0,ymax=5.276,axis y line*=left,axis x line*=bottom,xlabel={\scriptsize Halbzüge bis Matt},ylabel={\scriptsize Stellungen (Mrd.)},tick label style={font=\tiny},label style={font=\scriptsize},enlarge x limits=false,ylabel shift=-3pt,xlabel shift=-3pt]\addplot[fill=ink!55,draw=none] coordinates {(0,0.006877) (1,0.05965) (2,0.2368) (3,1.085) (4,2.365) (5,4.6) (6,4.582) (7,4.796) (8,2.636) (9,1.617) (10,0.3753) (11,0.07444) (12,0.005537) (13,0.0003009) (14,2.979e-06) (15,1.042e-06)};\draw[densely dotted] (axis cs:12,0) -- (axis cs:12,5.276);\end{axis}\begin{axis}[width=72mm,height=42mm,axis y line*=right,axis x line=none,ymode=log,xmin=-1,xmax=16,ylabel={\scriptsize eindeutig (log)},tick label style={font=\tiny},label style={font=\scriptsize},ylabel shift=-3pt]\addplot[mark=*,mark size=0.7pt,draw=black,thin] coordinates {(0,6877110) (1,51227344) (2,117786744) (3,246285048) (4,235717505) (5,203885988) (6,86158632) (7,59703357) (8,14468088) (9,3073324) (10,190641) (11,4080) (12,28)};\end{axis}\end{tikzpicture}\end{minipage}
\vspace{3mm}

\noindent\begin{minipage}[t]{60mm}\vspace{0pt}\sbox\hmboard{\chessboard[setfen=8/8/7k/6Q1/1qp5/b7/8/K7 b - - 0 1]}\usebox\hmboard\\[1mm]
\makebox[\wd\hmboard]{\sffamily\hspace{3mm}\textbf{h\#6}\hfill(2+4)}\end{minipage}\hfill
\begin{minipage}[t]{52mm}\vspace{0pt}
\textcolor{ink}{\sffamily\textsc{Lösung}}\par\smallskip\small
\begin{tabular}{@{}r@{\enspace}l@{\enspace}l@{}}
\textbf{1.} & \symking{}h7 & \symking{}a2\\
\textbf{2.} & \symqueen{}b7 & \symking{}xa3\\
\textbf{3.} & \symqueen{}g7 & \symking{}b4\\
\textbf{4.} & \symking{}g8 & \symking{}c5\\
\textbf{5.} & \symking{}f8 & \symking{}d6\\
\textbf{6.} & \symqueen{}g8 & \symqueen{}e7\#\\
\end{tabular}\par\vspace{3mm}
\qrcode[height=12mm]{https://helpman.komtera.lt/?fen=8/8/7k/6Q1/1qp5/b7/8/K7&moves=6}\\{\sffamily\tiny\color{ink}scannen zum Lösen}
\end{minipage}
\vspace{2mm}

{\raggedright\noindent\textcolor{ink}{\sffamily\textsc{Themen}}\enspace \themetag{reines Matt} \themetag{Mustermatt} \themetag{Selbstblock} \themetag{ohne Schach}\par\smallskip{\footnotesize\textit{Hinweis:} die Partei am Zug steht im Diagramm im Schach.}\par}

\clearpage
\phantomsection\label{hm:KQvkqnp}
\begin{tcolorbox}[enhanced,colback=ink,colframe=ink,arc=0pt,boxrule=0pt,left=4mm,right=4mm,top=2.5mm,bottom=2.5mm]
\color{white}\sffamily{\LARGE\bfseries KQvkqnp}\hfill{\Large h\#6}\\[0.5mm]
{\small 6 Steine\quad Schwarz am Zug}\hfill{\small Nr.~709\qquad\href{https://helpman.komtera.lt/?fen=8/8/8/7k/1p4Q1/2n5/7q/K7&moves=6}{\color{white}\bfseries lösen\,\ensuremath{\nearrow}}}
\end{tcolorbox}
\vspace{2mm}
\noindent\begin{minipage}[t]{36mm}\vspace{0pt}\sffamily\small\raggedright
\hmfigure{längste eindeutige}{{\Large h\#6}}\hmfigure{eindeutig bei dieser Länge}{{\Large 605} Stellungen}\hmfigure{längstes Matt}{{\Large h\#7.5}}\hmfigure{in der Klasse}{45,4 Milliarden Stellungen,\\ 23,1 Milliarden mit Hilfsmatt}
\end{minipage}\hfill
\begin{minipage}[t]{80mm}\vspace{0pt}\begin{tikzpicture}\begin{axis}[width=72mm,height=42mm,ybar,bar width=1.6pt,xmin=-1,xmax=16,ymin=0,ymax=5.374,axis y line*=left,axis x line*=bottom,xlabel={\scriptsize Halbzüge bis Matt},ylabel={\scriptsize Stellungen (Mrd.)},tick label style={font=\tiny},label style={font=\scriptsize},enlarge x limits=false,ylabel shift=-3pt,xlabel shift=-3pt]\addplot[fill=ink!55,draw=none] coordinates {(0,0.007802) (1,0.06493) (2,0.2409) (3,1.035) (4,2.182) (5,4.131) (6,4.567) (7,4.886) (8,3.203) (9,1.976) (10,0.6169) (11,0.1548) (12,0.02997) (13,0.00134) (14,0.0001065) (15,1.22e-07)};\draw[densely dotted] (axis cs:12,0) -- (axis cs:12,5.374);\end{axis}\begin{axis}[width=72mm,height=42mm,axis y line*=right,axis x line=none,ymode=log,xmin=-1,xmax=16,ylabel={\scriptsize eindeutig (log)},tick label style={font=\tiny},label style={font=\scriptsize},ylabel shift=-3pt]\addplot[mark=*,mark size=0.7pt,draw=black,thin] coordinates {(0,7802074) (1,54189420) (2,114345882) (3,226463754) (4,219944248) (5,199921734) (6,100612995) (7,66048720) (8,19592453) (9,4735688) (10,456477) (11,12071) (12,605)};\end{axis}\end{tikzpicture}\end{minipage}
\vspace{3mm}

\noindent\begin{minipage}[t]{60mm}\vspace{0pt}\sbox\hmboard{\chessboard[setfen=8/8/8/7k/1p4Q1/2n5/7q/K7 b - - 0 1]}\usebox\hmboard\\[1mm]
\makebox[\wd\hmboard]{\sffamily\hspace{3mm}\textbf{h\#6}\hfill(2+4)}\end{minipage}\hfill
\begin{minipage}[t]{52mm}\vspace{0pt}
\textcolor{ink}{\sffamily\textsc{Lösung}}\par\smallskip\small
\begin{tabular}{@{}r@{\enspace}l@{\enspace}l@{}}
\textbf{1.} & \symking{}h6 & \symqueen{}e2\\
\textbf{2.} & \symking{}g5 & \symking{}b2\\
\textbf{3.} & \symking{}f4 & \symqueen{}xh2+\\
\textbf{4.} & \symking{}e3 & \symqueen{}h7\\
\textbf{5.} & \symking{}d2 & \symking{}b3\\
\textbf{6.} & \symking{}c1 & \symqueen{}c2\#\\
\end{tabular}\par\vspace{3mm}
\qrcode[height=12mm]{https://helpman.komtera.lt/?fen=8/8/8/7k/1p4Q1/2n5/7q/K7&moves=6}\\{\sffamily\tiny\color{ink}scannen zum Lösen}
\end{minipage}
\vspace{2mm}

{\raggedright\noindent\textcolor{ink}{\sffamily\textsc{Themen}}\enspace \themetag{ein Stein zieht}\par\smallskip{\footnotesize\textit{Hinweis:} die Partei am Zug steht im Diagramm im Schach.}\par}

\clearpage
\phantomsection\label{hm:KQvkqpp}
\begin{tcolorbox}[enhanced,colback=ink,colframe=ink,arc=0pt,boxrule=0pt,left=4mm,right=4mm,top=2.5mm,bottom=2.5mm]
\color{white}\sffamily{\LARGE\bfseries KQvkqpp}\hfill{\Large h\#6}\\[0.5mm]
{\small 6 Steine\quad Schwarz am Zug}\hfill{\small Nr.~710\qquad\href{https://helpman.komtera.lt/?fen=8/8/8/7k/p5Q1/1qp5/8/K7&moves=6}{\color{white}\bfseries lösen\,\ensuremath{\nearrow}}}
\end{tcolorbox}
\vspace{2mm}
\noindent\begin{minipage}[t]{36mm}\vspace{0pt}\sffamily\small\raggedright
\hmfigure{längste eindeutige}{{\Large h\#6}}\hmfigure{eindeutig bei dieser Länge}{{\Large 60} Stellungen}\hmfigure{längstes Matt}{{\Large h\#7.5}}\hmfigure{in der Klasse}{34,1 Milliarden Stellungen,\\ 8,87 Milliarden mit Hilfsmatt}
\end{minipage}\hfill
\begin{minipage}[t]{80mm}\vspace{0pt}\begin{tikzpicture}\begin{axis}[width=72mm,height=42mm,ybar,bar width=1.6pt,xmin=-1,xmax=16,ymin=0,ymax=2.139,axis y line*=left,axis x line*=bottom,xlabel={\scriptsize Halbzüge bis Matt},ylabel={\scriptsize Stellungen (Mrd.)},tick label style={font=\tiny},label style={font=\scriptsize},enlarge x limits=false,ylabel shift=-3pt,xlabel shift=-3pt]\addplot[fill=ink!55,draw=none] coordinates {(0,0.00379) (1,0.03081) (2,0.114) (3,0.4752) (4,1.076) (5,1.796) (6,1.944) (7,1.75) (8,1.014) (9,0.5123) (10,0.1246) (11,0.02367) (12,0.00301) (13,0.0001097) (14,4.503e-06) (15,1e-09)};\draw[densely dotted] (axis cs:12,0) -- (axis cs:12,2.139);\end{axis}\begin{axis}[width=72mm,height=42mm,axis y line*=right,axis x line=none,ymode=log,xmin=-1,xmax=16,ylabel={\scriptsize eindeutig (log)},tick label style={font=\tiny},label style={font=\scriptsize},ylabel shift=-3pt]\addplot[mark=*,mark size=0.7pt,draw=black,thin] coordinates {(0,3789917) (1,25470026) (2,56924011) (3,101066604) (4,103533088) (5,73185824) (6,31822772) (7,17858614) (8,4338930) (9,813747) (10,62146) (11,1188) (12,60)};\end{axis}\end{tikzpicture}\end{minipage}
\vspace{3mm}

\noindent\begin{minipage}[t]{60mm}\vspace{0pt}\sbox\hmboard{\chessboard[setfen=8/8/8/7k/p5Q1/1qp5/8/K7 b - - 0 1]}\usebox\hmboard\\[1mm]
\makebox[\wd\hmboard]{\sffamily\hspace{3mm}\textbf{h\#6}\hfill(2+4)}\end{minipage}\hfill
\begin{minipage}[t]{52mm}\vspace{0pt}
\textcolor{ink}{\sffamily\textsc{Lösung}}\par\smallskip\small
\begin{tabular}{@{}r@{\enspace}l@{\enspace}l@{}}
\textbf{1.} & \symking{}h6 & \symqueen{}xa4\\
\textbf{2.} & \symking{}g5 & \symqueen{}xb3\\
\textbf{3.} & \symking{}f4 & \symking{}a2\\
\textbf{4.} & \symking{}e3 & \symqueen{}e6+\\
\textbf{5.} & \symking{}d2 & \symking{}b3\\
\textbf{6.} & \symking{}c1 & \symqueen{}e1\#\\
\end{tabular}\par\vspace{3mm}
\qrcode[height=12mm]{https://helpman.komtera.lt/?fen=8/8/8/7k/p5Q1/1qp5/8/K7&moves=6}\\{\sffamily\tiny\color{ink}scannen zum Lösen}
\end{minipage}
\vspace{2mm}

{\raggedright\noindent\textcolor{ink}{\sffamily\textsc{Themen}}\enspace \themetag{reines Matt} \themetag{Mustermatt} \themetag{Spiegelmatt} \themetag{ein Stein zieht}\par\smallskip{\footnotesize\textit{Hinweis:} die Partei am Zug steht im Diagramm im Schach.}\par}

\clearpage
\phantomsection\label{hm:KQvkqqb}
\begin{tcolorbox}[enhanced,colback=ink,colframe=ink,arc=0pt,boxrule=0pt,left=4mm,right=4mm,top=2.5mm,bottom=2.5mm]
\color{white}\sffamily{\LARGE\bfseries KQvkqqb}\hfill{\Large h\#6}\\[0.5mm]
{\small 6 Steine\quad Schwarz am Zug}\hfill{\small Nr.~711\qquad\href{https://helpman.komtera.lt/?fen=8/7k/6Q1/8/8/3q4/2q5/K1b5&moves=6}{\color{white}\bfseries lösen\,\ensuremath{\nearrow}}}
\end{tcolorbox}
\vspace{2mm}
\noindent\begin{minipage}[t]{36mm}\vspace{0pt}\sffamily\small\raggedright
\hmfigure{längste eindeutige}{{\Large h\#6}}\hmfigure{eindeutig bei dieser Länge}{{\Large 8} Stellungen}\hmfigure{längstes Matt}{{\Large h\#7}}\hmfigure{in der Klasse}{15,5 Milliarden Stellungen,\\ 3,1 Milliarden mit Hilfsmatt}
\end{minipage}\hfill
\begin{minipage}[t]{80mm}\vspace{0pt}\begin{tikzpicture}\begin{axis}[width=72mm,height=42mm,ybar,bar width=1.6pt,xmin=-1,xmax=15,ymin=0,ymax=878.8,axis y line*=left,axis x line*=bottom,xlabel={\scriptsize Halbzüge bis Matt},ylabel={\scriptsize Stellungen (Mio.)},tick label style={font=\tiny},label style={font=\scriptsize},enlarge x limits=false,ylabel shift=-3pt,xlabel shift=-3pt]\addplot[fill=ink!55,draw=none] coordinates {(0,0.4265) (1,4.271) (2,20.48) (3,115.3) (4,230.1) (5,638.1) (6,483.8) (7,798.9) (8,331.5) (9,376.6) (10,66.65) (11,31.55) (12,0.8265) (13,0.1236) (14,0.000473)};\draw[densely dotted] (axis cs:12,0) -- (axis cs:12,878.8);\end{axis}\begin{axis}[width=72mm,height=42mm,axis y line*=right,axis x line=none,ymode=log,xmin=-1,xmax=15,ylabel={\scriptsize eindeutig (log)},tick label style={font=\tiny},label style={font=\scriptsize},ylabel shift=-3pt]\addplot[mark=*,mark size=0.7pt,draw=black,thin] coordinates {(0,426484) (1,3928810) (2,10189932) (3,29852049) (4,23882364) (5,35245628) (6,11192832) (7,13737419) (8,2728645) (9,1392305) (10,37407) (11,3132) (12,8)};\end{axis}\end{tikzpicture}\end{minipage}
\vspace{3mm}

\noindent\begin{minipage}[t]{60mm}\vspace{0pt}\sbox\hmboard{\chessboard[setfen=8/7k/6Q1/8/8/3q4/2q5/K1b5 b - - 0 1]}\usebox\hmboard\\[1mm]
\makebox[\wd\hmboard]{\sffamily\hspace{3mm}\textbf{h\#6}\hfill(2+4)}\end{minipage}\hfill
\begin{minipage}[t]{52mm}\vspace{0pt}
\textcolor{ink}{\sffamily\textsc{Lösung}}\par\smallskip\small
\begin{tabular}{@{}r@{\enspace}l@{\enspace}l@{}}
\textbf{1.} & \symking{}h8 & \symqueen{}g3\\
\textbf{2.} & \symqueen{}c7 & \symking{}a2\\
\textbf{3.} & \symqueen{}g6 & \symking{}b3\\
\textbf{4.} & \symqueen{}cf7+ & \symking{}c3\\
\textbf{5.} & \symking{}g7 & \symking{}d4\\
\textbf{6.} & \symking{}f6 & \symqueen{}e5\#\\
\end{tabular}\par\vspace{3mm}
\qrcode[height=12mm]{https://helpman.komtera.lt/?fen=8/7k/6Q1/8/8/3q4/2q5/K1b5&moves=6}\\{\sffamily\tiny\color{ink}scannen zum Lösen}
\end{minipage}
\vspace{2mm}

{\raggedright\noindent\textcolor{ink}{\sffamily\textsc{Themen}}\enspace \themetag{reines Matt} \themetag{Mustermatt} \themetag{Selbstblock} \themetag{ohne Schlag}\par\smallskip{\footnotesize\textit{Hinweis:} die Partei am Zug steht im Diagramm im Schach.}\par}

\clearpage
\phantomsection\label{hm:KQvkqqn}
\begin{tcolorbox}[enhanced,colback=ink,colframe=ink,arc=0pt,boxrule=0pt,left=4mm,right=4mm,top=2.5mm,bottom=2.5mm]
\color{white}\sffamily{\LARGE\bfseries KQvkqqn}\hfill{\Large h\#6}\\[0.5mm]
{\small 6 Steine\quad Schwarz am Zug}\hfill{\small Nr.~712\qquad\href{https://helpman.komtera.lt/?fen=8/8/4n2k/6Q1/8/8/2q4q/K7&moves=6}{\color{white}\bfseries lösen\,\ensuremath{\nearrow}}}
\end{tcolorbox}
\vspace{2mm}
\noindent\begin{minipage}[t]{36mm}\vspace{0pt}\sffamily\small\raggedright
\hmfigure{längste eindeutige}{{\Large h\#6}}\hmfigure{eindeutig bei dieser Länge}{{\Large 36} Stellungen}\hmfigure{längstes Matt}{{\Large h\#7.5}}\hmfigure{in der Klasse}{15,5 Milliarden Stellungen,\\ 3,2 Milliarden mit Hilfsmatt}
\end{minipage}\hfill
\begin{minipage}[t]{80mm}\vspace{0pt}\begin{tikzpicture}\begin{axis}[width=72mm,height=42mm,ybar,bar width=1.6pt,xmin=-1,xmax=16,ymin=0,ymax=901.8,axis y line*=left,axis x line*=bottom,xlabel={\scriptsize Halbzüge bis Matt},ylabel={\scriptsize Stellungen (Mio.)},tick label style={font=\tiny},label style={font=\scriptsize},enlarge x limits=false,ylabel shift=-3pt,xlabel shift=-3pt]\addplot[fill=ink!55,draw=none] coordinates {(0,0.4787) (1,4.65) (2,21.24) (3,114.3) (4,222.3) (5,597.5) (6,515.3) (7,819.8) (8,376.8) (9,402.3) (10,80.4) (11,39.43) (12,1.261) (13,0.6231) (14,0.000857) (15,0.001933)};\draw[densely dotted] (axis cs:12,0) -- (axis cs:12,901.8);\end{axis}\begin{axis}[width=72mm,height=42mm,axis y line*=right,axis x line=none,ymode=log,xmin=-1,xmax=16,ylabel={\scriptsize eindeutig (log)},tick label style={font=\tiny},label style={font=\scriptsize},ylabel shift=-3pt]\addplot[mark=*,mark size=0.7pt,draw=black,thin] coordinates {(0,478719) (1,4215347) (2,10340528) (3,28671288) (4,23707833) (5,34576136) (6,13054256) (7,14510460) (8,3147655) (9,1570232) (10,49830) (11,7780) (12,36)};\end{axis}\end{tikzpicture}\end{minipage}
\vspace{3mm}

\noindent\begin{minipage}[t]{60mm}\vspace{0pt}\sbox\hmboard{\chessboard[setfen=8/8/4n2k/6Q1/8/8/2q4q/K7 b - - 0 1]}\usebox\hmboard\\[1mm]
\makebox[\wd\hmboard]{\sffamily\hspace{3mm}\textbf{h\#6}\hfill(2+4)}\end{minipage}\hfill
\begin{minipage}[t]{52mm}\vspace{0pt}
\textcolor{ink}{\sffamily\textsc{Lösung}}\par\smallskip\small
\begin{tabular}{@{}r@{\enspace}l@{\enspace}l@{}}
\textbf{1.} & \symking{}h7 & \symqueen{}d2\\
\textbf{2.} & \symqueen{}g6 & \symking{}b2\\
\textbf{3.} & \symknight{}f4 & \symking{}c3\\
\textbf{4.} & \symking{}h6 & \symking{}d4\\
\textbf{5.} & \symking{}g5 & \symking{}e5\\
\textbf{6.} & \symqueen{}hh5 & \symqueen{}xf4\#\\
\end{tabular}\par\vspace{3mm}
\qrcode[height=12mm]{https://helpman.komtera.lt/?fen=8/8/4n2k/6Q1/8/8/2q4q/K7&moves=6}\\{\sffamily\tiny\color{ink}scannen zum Lösen}
\end{minipage}
\vspace{2mm}

{\raggedright\noindent\textcolor{ink}{\sffamily\textsc{Themen}}\enspace \themetag{Switchback} \themetag{Selbstblock} \themetag{ohne Schach} \themetag{Maslar}\par\smallskip{\footnotesize\textit{Hinweis:} die Partei am Zug steht im Diagramm im Schach.}\par}

\clearpage
\phantomsection\label{hm:KQvkqqp}
\begin{tcolorbox}[enhanced,colback=ink,colframe=ink,arc=0pt,boxrule=0pt,left=4mm,right=4mm,top=2.5mm,bottom=2.5mm]
\color{white}\sffamily{\LARGE\bfseries KQvkqqp}\hfill{\Large h\#6}\\[0.5mm]
{\small 6 Steine\quad Schwarz am Zug}\hfill{\small Nr.~713\qquad\href{https://helpman.komtera.lt/?fen=8/7k/6Q1/2p5/8/7q/7q/K7&moves=6}{\color{white}\bfseries lösen\,\ensuremath{\nearrow}}}
\end{tcolorbox}
\vspace{2mm}
\noindent\begin{minipage}[t]{36mm}\vspace{0pt}\sffamily\small\raggedright
\hmfigure{längste eindeutige}{{\Large h\#6}}\hmfigure{eindeutig bei dieser Länge}{{\Large 6} Stellungen}\hmfigure{längstes Matt}{{\Large h\#7.5}}\hmfigure{in der Klasse}{45,4 Milliarden Stellungen,\\ 9,85 Milliarden mit Hilfsmatt}
\end{minipage}\hfill
\begin{minipage}[t]{80mm}\vspace{0pt}\begin{tikzpicture}\begin{axis}[width=72mm,height=42mm,ybar,bar width=1.6pt,xmin=-1,xmax=16,ymin=0,ymax=2.67,axis y line*=left,axis x line*=bottom,xlabel={\scriptsize Halbzüge bis Matt},ylabel={\scriptsize Stellungen (Mrd.)},tick label style={font=\tiny},label style={font=\scriptsize},enlarge x limits=false,ylabel shift=-3pt,xlabel shift=-3pt]\addplot[fill=ink!55,draw=none] coordinates {(0,0.00186) (1,0.01771) (2,0.08116) (3,0.4286) (4,0.8809) (5,2.146) (6,1.746) (7,2.428) (8,1.007) (9,0.9095) (10,0.145) (11,0.05233) (12,0.001338) (13,0.0003348) (14,3.9e-07) (15,1.046e-06)};\draw[densely dotted] (axis cs:12,0) -- (axis cs:12,2.67);\end{axis}\begin{axis}[width=72mm,height=42mm,axis y line*=right,axis x line=none,ymode=log,xmin=-1,xmax=16,ylabel={\scriptsize eindeutig (log)},tick label style={font=\tiny},label style={font=\scriptsize},ylabel shift=-3pt]\addplot[mark=*,mark size=0.7pt,draw=black,thin] coordinates {(0,1859570) (1,15930722) (2,40710996) (3,106257840) (4,89612465) (5,108290253) (6,35095045) (7,33385472) (8,6082116) (9,2227324) (10,59828) (11,5348) (12,6)};\end{axis}\end{tikzpicture}\end{minipage}
\vspace{3mm}

\noindent\begin{minipage}[t]{60mm}\vspace{0pt}\sbox\hmboard{\chessboard[setfen=8/7k/6Q1/2p5/8/7q/7q/K7 b - - 0 1]}\usebox\hmboard\\[1mm]
\makebox[\wd\hmboard]{\sffamily\hspace{3mm}\textbf{h\#6}\hfill(2+4)}\end{minipage}\hfill
\begin{minipage}[t]{52mm}\vspace{0pt}
\textcolor{ink}{\sffamily\textsc{Lösung}}\par\smallskip\small
\begin{tabular}{@{}r@{\enspace}l@{\enspace}l@{}}
\textbf{1.} & \symking{}h8 & \symqueen{}b6\\
\textbf{2.} & \symqueen{}b8 & \symking{}b2\\
\textbf{3.} & \symqueen{}h7 & \symking{}c3\\
\textbf{4.} & c4 & \symking{}d4\\
\textbf{5.} & \symqueen{}bg8 & \symking{}e5\\
\textbf{6.} & \symking{}g7 & \symqueen{}f6\#\\
\end{tabular}\par\vspace{3mm}
\qrcode[height=12mm]{https://helpman.komtera.lt/?fen=8/7k/6Q1/2p5/8/7q/7q/K7&moves=6}\\{\sffamily\tiny\color{ink}scannen zum Lösen}
\end{minipage}
\vspace{2mm}

{\raggedright\noindent\textcolor{ink}{\sffamily\textsc{Themen}}\enspace \themetag{reines Matt} \themetag{Mustermatt} \themetag{Selbstblock} \themetag{ohne Schlag} \themetag{ohne Schach}\par\smallskip{\footnotesize\textit{Hinweis:} die Partei am Zug steht im Diagramm im Schach.}\par}

\clearpage
\phantomsection\label{hm:KQvkqqq}
\begin{tcolorbox}[enhanced,colback=ink,colframe=ink,arc=0pt,boxrule=0pt,left=4mm,right=4mm,top=2.5mm,bottom=2.5mm]
\color{white}\sffamily{\LARGE\bfseries KQvkqqq}\hfill{\Large h\#6}\\[0.5mm]
{\small 6 Steine\quad Schwarz am Zug}\hfill{\small Nr.~714\qquad\href{https://helpman.komtera.lt/?fen=8/7k/6Q1/8/8/q3q3/3q4/1K6&moves=6}{\color{white}\bfseries lösen\,\ensuremath{\nearrow}}}
\end{tcolorbox}
\vspace{2mm}
\noindent\begin{minipage}[t]{36mm}\vspace{0pt}\sffamily\small\raggedright
\hmfigure{längste eindeutige}{{\Large h\#6}}\hmfigure{eindeutig bei dieser Länge}{{\Large 2} Stellungen}\hmfigure{längstes Matt}{{\Large h\#7}}\hmfigure{in der Klasse}{15,5 Milliarden Stellungen,\\ 893 Millionen mit Hilfsmatt}
\end{minipage}\hfill
\begin{minipage}[t]{80mm}\vspace{0pt}\begin{tikzpicture}\begin{axis}[width=72mm,height=42mm,ybar,bar width=1.6pt,xmin=-1,xmax=15,ymin=0,ymax=278.4,axis y line*=left,axis x line*=bottom,xlabel={\scriptsize Halbzüge bis Matt},ylabel={\scriptsize Stellungen (Mio.)},tick label style={font=\tiny},label style={font=\scriptsize},enlarge x limits=false,ylabel shift=-3pt,xlabel shift=-3pt]\addplot[fill=ink!55,draw=none] coordinates {(0,0.0728) (1,0.8035) (2,4.384) (3,27.76) (4,53.67) (5,180.9) (6,115.8) (7,253.1) (8,82.83) (9,136.7) (10,20.08) (11,16.96) (12,0.294) (13,0.06056) (14,0.000168)};\draw[densely dotted] (axis cs:12,0) -- (axis cs:12,278.4);\end{axis}\begin{axis}[width=72mm,height=42mm,axis y line*=right,axis x line=none,ymode=log,xmin=-1,xmax=15,ylabel={\scriptsize eindeutig (log)},tick label style={font=\tiny},label style={font=\scriptsize},ylabel shift=-3pt]\addplot[mark=*,mark size=0.7pt,draw=black,thin] coordinates {(0,72796) (1,759752) (2,2149064) (3,7465544) (4,5580009) (5,11494508) (6,2723765) (7,4747592) (8,722522) (9,656580) (10,12344) (11,1547) (12,2)};\end{axis}\end{tikzpicture}\end{minipage}
\vspace{3mm}

\noindent\begin{minipage}[t]{60mm}\vspace{0pt}\sbox\hmboard{\chessboard[setfen=8/7k/6Q1/8/8/q3q3/3q4/1K6 b - - 0 1]}\usebox\hmboard\\[1mm]
\makebox[\wd\hmboard]{\sffamily\hspace{3mm}\textbf{h\#6}\hfill(2+4)}\end{minipage}\hfill
\begin{minipage}[t]{52mm}\vspace{0pt}
\textcolor{ink}{\sffamily\textsc{Lösung}}\par\smallskip\small
\begin{tabular}{@{}r@{\enspace}l@{\enspace}l@{}}
\textbf{1.} & \symking{}h8 & \symqueen{}d3\\
\textbf{2.} & \symqueen{}c2+ & \symking{}xc2\\
\textbf{3.} & \symqueen{}d2+ & \symking{}xd2\\
\textbf{4.} & \symking{}g7 & \symking{}e3\\
\textbf{5.} & \symking{}h6 & \symking{}f4\\
\textbf{6.} & \symking{}h5 & \symqueen{}h7\#\\
\end{tabular}\par\vspace{3mm}
\qrcode[height=12mm]{https://helpman.komtera.lt/?fen=8/7k/6Q1/8/8/q3q3/3q4/1K6&moves=6}\\{\sffamily\tiny\color{ink}scannen zum Lösen}
\end{minipage}
\vspace{2mm}

{\raggedright\noindent\textcolor{ink}{\sffamily\textsc{Themen}}\enspace \themetag{reines Matt} \themetag{Mustermatt} \themetag{Spiegelmatt}\par\smallskip{\footnotesize\textit{Hinweis:} die Partei am Zug steht im Diagramm im Schach.}\par}

\clearpage
\phantomsection\label{hm:KRBBNvk}
\begin{tcolorbox}[enhanced,colback=ink,colframe=ink,arc=0pt,boxrule=0pt,left=4mm,right=4mm,top=2.5mm,bottom=2.5mm]
\color{white}\sffamily{\LARGE\bfseries KRBBNvk}\hfill{\Large h\#6}\\[0.5mm]
{\small 6 Steine\quad Schwarz am Zug}\hfill{\small Nr.~715\qquad\href{https://helpman.komtera.lt/?fen=4BN2/3B2R1/7k/8/8/8/8/K7&moves=6}{\color{white}\bfseries lösen\,\ensuremath{\nearrow}}}
\end{tcolorbox}
\vspace{2mm}
\noindent\begin{minipage}[t]{36mm}\vspace{0pt}\sffamily\small\raggedright
\hmfigure{längste eindeutige}{{\Large h\#6}}\hmfigure{eindeutig bei dieser Länge}{{\Large 94} Stellungen}\hmfigure{längstes Matt}{{\Large h\#8}}\hmfigure{in der Klasse}{15,5 Milliarden Stellungen,\\ 4,7 Milliarden mit Hilfsmatt}
\end{minipage}\hfill
\begin{minipage}[t]{80mm}\vspace{0pt}\begin{tikzpicture}\begin{axis}[width=72mm,height=42mm,ybar,bar width=1.6pt,xmin=-1,xmax=17,ymin=0,ymax=1.624,axis y line*=left,axis x line*=bottom,xlabel={\scriptsize Halbzüge bis Matt},ylabel={\scriptsize Stellungen (Mrd.)},tick label style={font=\tiny},label style={font=\scriptsize},enlarge x limits=false,ylabel shift=-3pt,xlabel shift=-3pt]\addplot[fill=ink!55,draw=none] coordinates {(0,0.01349) (1,0.04873) (2,0.1496) (3,0.463) (4,1.476) (5,1.069) (6,1.273) (7,0.1151) (8,0.08076) (9,0.001171) (10,0.004243) (11,1.192e-06) (12,0.0009061) (13,1.17e-07) (14,0.0001986) (15,3.3e-08) (16,7.19e-07)};\draw[densely dotted] (axis cs:12,0) -- (axis cs:12,1.624);\end{axis}\begin{axis}[width=72mm,height=42mm,axis y line*=right,axis x line=none,ymode=log,xmin=-1,xmax=17,ylabel={\scriptsize eindeutig (log)},tick label style={font=\tiny},label style={font=\scriptsize},ylabel shift=-3pt]\addplot[mark=*,mark size=0.7pt,draw=black,thin] coordinates {(0,13488456) (1,42921259) (2,123288038) (3,223031243) (4,389235857) (5,35600361) (6,16950085) (7,110057) (8,98517) (9,147) (10,4628) (12,94)};\end{axis}\end{tikzpicture}\end{minipage}
\vspace{3mm}

\noindent\begin{minipage}[t]{60mm}\vspace{0pt}\sbox\hmboard{\chessboard[setfen=4BN2/3B2R1/7k/8/8/8/8/K7 b - - 0 1]}\usebox\hmboard\\[1mm]
\makebox[\wd\hmboard]{\sffamily\hspace{3mm}\textbf{h\#6}\hfill(5+1)}\end{minipage}\hfill
\begin{minipage}[t]{52mm}\vspace{0pt}
\textcolor{ink}{\sffamily\textsc{Lösung}}\par\smallskip\small
\begin{tabular}{@{}r@{\enspace}l@{\enspace}l@{}}
\textbf{1.} & \symking{}xg7 & \symking{}b2\\
\textbf{2.} & \symking{}f6 & \symking{}c3\\
\textbf{3.} & \symking{}e7 & \symking{}d4\\
\textbf{4.} & \symking{}d8 & \symking{}e5\\
\textbf{5.} & \symking{}c7 & \symking{}f6\\
\textbf{6.} & \symking{}d8 & \symknight{}e6\#\\
\end{tabular}\par\vspace{3mm}
\qrcode[height=12mm]{https://helpman.komtera.lt/?fen=4BN2/3B2R1/7k/8/8/8/8/K7&moves=6}\\{\sffamily\tiny\color{ink}scannen zum Lösen}
\end{minipage}
\vspace{2mm}

{\raggedright\noindent\textcolor{ink}{\sffamily\textsc{Themen}}\enspace \themetag{reines Matt} \themetag{Mustermatt} \themetag{Idealmatt} \themetag{Switchback} \themetag{ein Stein zieht} \themetag{ohne Schach}\par\smallskip{\footnotesize\textit{Hinweis:} die Lösung beginnt mit einem Schlag.}\par}

\clearpage
\phantomsection\label{hm:KRBBvkr}
\begin{tcolorbox}[enhanced,colback=ink,colframe=ink,arc=0pt,boxrule=0pt,left=4mm,right=4mm,top=2.5mm,bottom=2.5mm]
\color{white}\sffamily{\LARGE\bfseries KRBBvkr}\hfill{\Large h\#6}\\[0.5mm]
{\small 6 Steine\quad Schwarz am Zug}\hfill{\small Nr.~716\qquad\href{https://helpman.komtera.lt/?fen=8/5B1R/7k/4B1r1/8/8/8/K7&moves=6}{\color{white}\bfseries lösen\,\ensuremath{\nearrow}}}
\end{tcolorbox}
\vspace{2mm}
\noindent\begin{minipage}[t]{36mm}\vspace{0pt}\sffamily\small\raggedright
\hmfigure{längste eindeutige}{{\Large h\#6}}\hmfigure{eindeutig bei dieser Länge}{{\Large 66} Stellungen}\hmfigure{längstes Matt}{{\Large h\#7.5}}\hmfigure{in der Klasse}{15,5 Milliarden Stellungen,\\ 4,26 Milliarden mit Hilfsmatt}
\end{minipage}\hfill
\begin{minipage}[t]{80mm}\vspace{0pt}\begin{tikzpicture}\begin{axis}[width=72mm,height=42mm,ybar,bar width=1.6pt,xmin=-1,xmax=16,ymin=0,ymax=1.368,axis y line*=left,axis x line*=bottom,xlabel={\scriptsize Halbzüge bis Matt},ylabel={\scriptsize Stellungen (Mrd.)},tick label style={font=\tiny},label style={font=\scriptsize},enlarge x limits=false,ylabel shift=-3pt,xlabel shift=-3pt]\addplot[fill=ink!55,draw=none] coordinates {(0,0.004792) (1,0.01855) (2,0.08561) (3,0.3841) (4,1.244) (5,1.22) (6,1.057) (7,0.2173) (8,0.03113) (9,0.001245) (10,0.0004073) (11,6.381e-06) (12,8.071e-05) (13,7e-08) (14,8.032e-06) (15,1.3e-08)};\draw[densely dotted] (axis cs:12,0) -- (axis cs:12,1.368);\end{axis}\begin{axis}[width=72mm,height=42mm,axis y line*=right,axis x line=none,ymode=log,xmin=-1,xmax=16,ylabel={\scriptsize eindeutig (log)},tick label style={font=\tiny},label style={font=\scriptsize},ylabel shift=-3pt]\addplot[mark=*,mark size=0.7pt,draw=black,thin] coordinates {(0,4792064) (1,17510885) (2,54426805) (3,138162846) (4,240111492) (5,40028533) (6,6450383) (7,206151) (8,7890) (9,197) (10,1460) (12,66)};\end{axis}\end{tikzpicture}\end{minipage}
\vspace{3mm}

\noindent\begin{minipage}[t]{60mm}\vspace{0pt}\sbox\hmboard{\chessboard[setfen=8/5B1R/7k/4B1r1/8/8/8/K7 b - - 0 1]}\usebox\hmboard\\[1mm]
\makebox[\wd\hmboard]{\sffamily\hspace{3mm}\textbf{h\#6}\hfill(4+2)}\end{minipage}\hfill
\begin{minipage}[t]{52mm}\vspace{0pt}
\textcolor{ink}{\sffamily\textsc{Lösung}}\par\smallskip\small
\begin{tabular}{@{}r@{\enspace}l@{\enspace}l@{}}
\textbf{1.} & \symking{}xh7 & \symking{}b2\\
\textbf{2.} & \symking{}h6 & \symking{}c3\\
\textbf{3.} & \symrook{}g3+ & \symking{}d4\\
\textbf{4.} & \symking{}g5 & \symking{}e4\\
\textbf{5.} & \symking{}h4 & \symking{}f4\\
\textbf{6.} & \symrook{}h3 & \symbishop{}f6\#\\
\end{tabular}\par\vspace{3mm}
\qrcode[height=12mm]{https://helpman.komtera.lt/?fen=8/5B1R/7k/4B1r1/8/8/8/K7&moves=6}\\{\sffamily\tiny\color{ink}scannen zum Lösen}
\end{minipage}
\vspace{2mm}

{\raggedright\noindent\textcolor{ink}{\sffamily\textsc{Themen}}\enspace \themetag{Switchback} \themetag{Selbstblock}\par\smallskip{\footnotesize\textit{Hinweis:} die Partei am Zug steht im Diagramm im Schach; die Lösung beginnt mit einem Schlag.}\par}

\clearpage
\phantomsection\label{hm:KRBvkqb}
\begin{tcolorbox}[enhanced,colback=ink,colframe=ink,arc=0pt,boxrule=0pt,left=4mm,right=4mm,top=2.5mm,bottom=2.5mm]
\color{white}\sffamily{\LARGE\bfseries KRBvkqb}\hfill{\Large h\#6}\\[0.5mm]
{\small 6 Steine\quad Schwarz am Zug}\hfill{\small Nr.~717\qquad\href{https://helpman.komtera.lt/?fen=8/8/8/8/8/Bb6/2R5/K1kq4&moves=6}{\color{white}\bfseries lösen\,\ensuremath{\nearrow}}}
\end{tcolorbox}
\vspace{2mm}
\noindent\begin{minipage}[t]{36mm}\vspace{0pt}\sffamily\small\raggedright
\hmfigure{längste eindeutige}{{\Large h\#6}}\hmfigure{eindeutig bei dieser Länge}{{\Large 273} Stellungen}\hmfigure{längstes Matt}{{\Large h\#10.5}}\hmfigure{in der Klasse}{15,5 Milliarden Stellungen,\\ 7,76 Milliarden mit Hilfsmatt}
\end{minipage}\hfill
\begin{minipage}[t]{80mm}\vspace{0pt}\begin{tikzpicture}\begin{axis}[width=72mm,height=42mm,ybar,bar width=1.6pt,xmin=-1,xmax=22,ymin=0,ymax=2.597,axis y line*=left,axis x line*=bottom,xlabel={\scriptsize Halbzüge bis Matt},ylabel={\scriptsize Stellungen (Mrd.)},tick label style={font=\tiny},label style={font=\scriptsize},enlarge x limits=false,ylabel shift=-3pt,xlabel shift=-3pt]\addplot[fill=ink!55,draw=none] coordinates {(0,0.001747) (1,0.007898) (2,0.04069) (3,0.2573) (4,0.813) (5,2.047) (6,2.361) (7,1.821) (8,0.3139) (9,0.09041) (10,0.001082) (11,0.004399) (12,0.0003844) (13,0.0003536) (14,5.191e-05) (15,1.863e-06) (16,4.3e-08) (17,2.1e-07) (18,1.08e-07) (19,3.5e-08) (20,4.55e-07) (21,7e-08)};\draw[densely dotted] (axis cs:12,0) -- (axis cs:12,2.597);\end{axis}\begin{axis}[width=72mm,height=42mm,axis y line*=right,axis x line=none,ymode=log,xmin=-1,xmax=22,ylabel={\scriptsize eindeutig (log)},tick label style={font=\tiny},label style={font=\scriptsize},ylabel shift=-3pt]\addplot[mark=*,mark size=0.7pt,draw=black,thin] coordinates {(0,1747169) (1,7798738) (2,22458156) (3,104470632) (4,216682851) (5,164933417) (6,31333524) (7,6561691) (8,33237) (9,75005) (10,3555) (11,2674) (12,273)};\end{axis}\end{tikzpicture}\end{minipage}
\vspace{3mm}

\noindent\begin{minipage}[t]{60mm}\vspace{0pt}\sbox\hmboard{\chessboard[setfen=8/8/8/8/8/Bb6/2R5/K1kq4 b - - 0 1]}\usebox\hmboard\\[1mm]
\makebox[\wd\hmboard]{\sffamily\hspace{3mm}\textbf{h\#6}\hfill(3+3)}\end{minipage}\hfill
\begin{minipage}[t]{52mm}\vspace{0pt}
\textcolor{ink}{\sffamily\textsc{Lösung}}\par\smallskip\small
\begin{tabular}{@{}r@{\enspace}l@{\enspace}l@{}}
\textbf{1.} & \symking{}xc2+ & \symbishop{}c1\\
\textbf{2.} & \symking{}d3 & \symking{}b2\\
\textbf{3.} & \symking{}e2 & \symking{}c3\\
\textbf{4.} & \symqueen{}f1 & \symking{}d4\\
\textbf{5.} & \symking{}e1 & \symking{}e3\\
\textbf{6.} & \symbishop{}d1 & \symbishop{}d2\#\\
\end{tabular}\par\vspace{3mm}
\qrcode[height=12mm]{https://helpman.komtera.lt/?fen=8/8/8/8/8/Bb6/2R5/K1kq4&moves=6}\\{\sffamily\tiny\color{ink}scannen zum Lösen}
\end{minipage}
\vspace{2mm}

{\raggedright\noindent\textcolor{ink}{\sffamily\textsc{Themen}}\enspace \themetag{reines Matt} \themetag{Mustermatt} \themetag{Idealmatt} \themetag{Selbstblock} \themetag{Umnov}\par\smallskip{\footnotesize\textit{Hinweis:} die Partei am Zug steht im Diagramm im Schach; die Lösung beginnt mit einem Schlag.}\par}

\clearpage
\phantomsection\label{hm:KRBvkqq}
\begin{tcolorbox}[enhanced,colback=ink,colframe=ink,arc=0pt,boxrule=0pt,left=4mm,right=4mm,top=2.5mm,bottom=2.5mm]
\color{white}\sffamily{\LARGE\bfseries KRBvkqq}\hfill{\Large h\#6}\\[0.5mm]
{\small 6 Steine\quad Schwarz am Zug}\hfill{\small Nr.~718\qquad\href{https://helpman.komtera.lt/?fen=6qk/6B1/8/8/8/7R/7q/K7&moves=6}{\color{white}\bfseries lösen\,\ensuremath{\nearrow}}}
\end{tcolorbox}
\vspace{2mm}
\noindent\begin{minipage}[t]{36mm}\vspace{0pt}\sffamily\small\raggedright
\hmfigure{längste eindeutige}{{\Large h\#6}}\hmfigure{eindeutig bei dieser Länge}{{\Large 21} Stellungen}\hmfigure{längstes Matt}{{\Large h\#8}}\hmfigure{in der Klasse}{15,5 Milliarden Stellungen,\\ 3,4 Milliarden mit Hilfsmatt}
\end{minipage}\hfill
\begin{minipage}[t]{80mm}\vspace{0pt}\begin{tikzpicture}\begin{axis}[width=72mm,height=42mm,ybar,bar width=1.6pt,xmin=-1,xmax=17,ymin=0,ymax=1.118,axis y line*=left,axis x line*=bottom,xlabel={\scriptsize Halbzüge bis Matt},ylabel={\scriptsize Stellungen (Mrd.)},tick label style={font=\tiny},label style={font=\scriptsize},enlarge x limits=false,ylabel shift=-3pt,xlabel shift=-3pt]\addplot[fill=ink!55,draw=none] coordinates {(0,0.000506) (1,0.002141) (2,0.01273) (3,0.09662) (4,0.2842) (5,0.911) (6,0.9244) (7,1.016) (8,0.1108) (9,0.04221) (10,6.711e-05) (11,4.7e-06) (12,1.526e-05) (13,1.1e-08) (14,7.44e-07) (16,1e-09)};\draw[densely dotted] (axis cs:12,0) -- (axis cs:12,1.118);\end{axis}\begin{axis}[width=72mm,height=42mm,axis y line*=right,axis x line=none,ymode=log,xmin=-1,xmax=17,ylabel={\scriptsize eindeutig (log)},tick label style={font=\tiny},label style={font=\scriptsize},ylabel shift=-3pt]\addplot[mark=*,mark size=0.7pt,draw=black,thin] coordinates {(0,505960) (1,2125043) (2,6912364) (3,39354314) (4,76108562) (5,85203187) (6,13755141) (7,4644492) (8,4833) (9,579) (10,220) (12,21)};\end{axis}\end{tikzpicture}\end{minipage}
\vspace{3mm}

\noindent\begin{minipage}[t]{60mm}\vspace{0pt}\sbox\hmboard{\chessboard[setfen=6qk/6B1/8/8/8/7R/7q/K7 b - - 0 1]}\usebox\hmboard\\[1mm]
\makebox[\wd\hmboard]{\sffamily\hspace{3mm}\textbf{h\#6}\hfill(3+3)}\end{minipage}\hfill
\begin{minipage}[t]{52mm}\vspace{0pt}
\textcolor{ink}{\sffamily\textsc{Lösung}}\par\smallskip\small
\begin{tabular}{@{}r@{\enspace}l@{\enspace}l@{}}
\textbf{1.} & \symking{}xg7 & \symking{}b1\\
\textbf{2.} & \symqueen{}e5 & \symking{}c2\\
\textbf{3.} & \symqueen{}h5 & \symking{}d3\\
\textbf{4.} & \symking{}h6 & \symking{}e4\\
\textbf{5.} & \symqueen{}h4+ & \symking{}f5\\
\textbf{6.} & \symqueen{}g7 & \symrook{}xh4\#\\
\end{tabular}\par\vspace{3mm}
\qrcode[height=12mm]{https://helpman.komtera.lt/?fen=6qk/6B1/8/8/8/7R/7q/K7&moves=6}\\{\sffamily\tiny\color{ink}scannen zum Lösen}
\end{minipage}
\vspace{2mm}

{\raggedright\noindent\textcolor{ink}{\sffamily\textsc{Themen}}\enspace \themetag{reines Matt} \themetag{Mustermatt} \themetag{Idealmatt} \themetag{Selbstblock}\par\smallskip{\footnotesize\textit{Hinweis:} die Partei am Zug steht im Diagramm im Schach; die Lösung beginnt mit einem Schlag.}\par}

\clearpage
\phantomsection\label{hm:KRBvkqr}
\begin{tcolorbox}[enhanced,colback=ink,colframe=ink,arc=0pt,boxrule=0pt,left=4mm,right=4mm,top=2.5mm,bottom=2.5mm]
\color{white}\sffamily{\LARGE\bfseries KRBvkqr}\hfill{\Large h\#6}\\[0.5mm]
{\small 6 Steine\quad Schwarz am Zug}\hfill{\small Nr.~719\qquad\href{https://helpman.komtera.lt/?fen=8/8/2q3r1/6Bk/7R/8/8/K7&moves=6}{\color{white}\bfseries lösen\,\ensuremath{\nearrow}}}
\end{tcolorbox}
\vspace{2mm}
\noindent\begin{minipage}[t]{36mm}\vspace{0pt}\sffamily\small\raggedright
\hmfigure{längste eindeutige}{{\Large h\#6}}\hmfigure{eindeutig bei dieser Länge}{{\Large 47} Stellungen}\hmfigure{längstes Matt}{{\Large h\#8}}\hmfigure{in der Klasse}{15,5 Milliarden Stellungen,\\ 7,41 Milliarden mit Hilfsmatt}
\end{minipage}\hfill
\begin{minipage}[t]{80mm}\vspace{0pt}\begin{tikzpicture}\begin{axis}[width=72mm,height=42mm,ybar,bar width=1.6pt,xmin=-1,xmax=17,ymin=0,ymax=2.399,axis y line*=left,axis x line*=bottom,xlabel={\scriptsize Halbzüge bis Matt},ylabel={\scriptsize Stellungen (Mrd.)},tick label style={font=\tiny},label style={font=\scriptsize},enlarge x limits=false,ylabel shift=-3pt,xlabel shift=-3pt]\addplot[fill=ink!55,draw=none] coordinates {(0,0.001462) (1,0.006781) (2,0.03774) (3,0.2546) (4,0.8123) (5,2.181) (6,2.176) (7,1.72) (8,0.1804) (9,0.03956) (10,0.0001051) (11,3.167e-06) (12,3.257e-05) (13,1.3e-08) (14,1.275e-06) (16,2e-09)};\draw[densely dotted] (axis cs:12,0) -- (axis cs:12,2.399);\end{axis}\begin{axis}[width=72mm,height=42mm,axis y line*=right,axis x line=none,ymode=log,xmin=-1,xmax=17,ylabel={\scriptsize eindeutig (log)},tick label style={font=\tiny},label style={font=\scriptsize},ylabel shift=-3pt]\addplot[mark=*,mark size=0.7pt,draw=black,thin] coordinates {(0,1462436) (1,6723759) (2,21154579) (3,106615552) (4,230140464) (5,203368623) (6,29900238) (7,5804088) (8,6263) (9,391) (10,366) (12,47)};\end{axis}\end{tikzpicture}\end{minipage}
\vspace{3mm}

\noindent\begin{minipage}[t]{60mm}\vspace{0pt}\sbox\hmboard{\chessboard[setfen=8/8/2q3r1/6Bk/7R/8/8/K7 b - - 0 1]}\usebox\hmboard\\[1mm]
\makebox[\wd\hmboard]{\sffamily\hspace{3mm}\textbf{h\#6}\hfill(3+3)}\end{minipage}\hfill
\begin{minipage}[t]{52mm}\vspace{0pt}
\textcolor{ink}{\sffamily\textsc{Lösung}}\par\smallskip\small
\begin{tabular}{@{}r@{\enspace}l@{\enspace}l@{}}
\textbf{1.} & \symking{}xg5 & \symking{}b2\\
\textbf{2.} & \symqueen{}e8 & \symking{}c3\\
\textbf{3.} & \symrook{}g8 & \symking{}d4\\
\textbf{4.} & \symqueen{}h5 & \symking{}e5\\
\textbf{5.} & \symking{}h6+ & \symking{}f6\\
\textbf{6.} & \symking{}h7 & \symrook{}xh5\#\\
\end{tabular}\par\vspace{3mm}
\qrcode[height=12mm]{https://helpman.komtera.lt/?fen=8/8/2q3r1/6Bk/7R/8/8/K7&moves=6}\\{\sffamily\tiny\color{ink}scannen zum Lösen}
\end{minipage}
\vspace{2mm}

{\raggedright\noindent\textcolor{ink}{\sffamily\textsc{Themen}}\enspace \themetag{reines Matt} \themetag{Mustermatt} \themetag{Idealmatt} \themetag{Selbstblock}\par\smallskip{\footnotesize\textit{Hinweis:} die Partei am Zug steht im Diagramm im Schach; die Lösung beginnt mit einem Schlag.}\par}

\clearpage
\phantomsection\label{hm:KRBvkrr}
\begin{tcolorbox}[enhanced,colback=ink,colframe=ink,arc=0pt,boxrule=0pt,left=4mm,right=4mm,top=2.5mm,bottom=2.5mm]
\color{white}\sffamily{\LARGE\bfseries KRBvkrr}\hfill{\Large h\#6}\\[0.5mm]
{\small 6 Steine\quad Schwarz am Zug}\hfill{\small Nr.~720\qquad\href{https://helpman.komtera.lt/?fen=6rk/6B1/8/8/8/7R/3r4/K7&moves=6}{\color{white}\bfseries lösen\,\ensuremath{\nearrow}}}
\end{tcolorbox}
\vspace{2mm}
\noindent\begin{minipage}[t]{36mm}\vspace{0pt}\sffamily\small\raggedright
\hmfigure{längste eindeutige}{{\Large h\#6}}\hmfigure{eindeutig bei dieser Länge}{{\Large 21} Stellungen}\hmfigure{längstes Matt}{{\Large h\#7}}\hmfigure{in der Klasse}{15,5 Milliarden Stellungen,\\ 4,02 Milliarden mit Hilfsmatt}
\end{minipage}\hfill
\begin{minipage}[t]{80mm}\vspace{0pt}\begin{tikzpicture}\begin{axis}[width=72mm,height=42mm,ybar,bar width=1.6pt,xmin=-1,xmax=15,ymin=0,ymax=1.376,axis y line*=left,axis x line*=bottom,xlabel={\scriptsize Halbzüge bis Matt},ylabel={\scriptsize Stellungen (Mrd.)},tick label style={font=\tiny},label style={font=\scriptsize},enlarge x limits=false,ylabel shift=-3pt,xlabel shift=-3pt]\addplot[fill=ink!55,draw=none] coordinates {(0,0.00107) (1,0.005346) (2,0.02701) (3,0.1621) (4,0.5344) (5,1.187) (6,1.251) (7,0.7432) (8,0.09453) (9,0.01431) (10,5.9e-05) (11,8.02e-07) (12,1.799e-05) (14,5.89e-07)};\draw[densely dotted] (axis cs:12,0) -- (axis cs:12,1.376);\end{axis}\begin{axis}[width=72mm,height=42mm,axis y line*=right,axis x line=none,ymode=log,xmin=-1,xmax=15,ylabel={\scriptsize eindeutig (log)},tick label style={font=\tiny},label style={font=\scriptsize},ylabel shift=-3pt]\addplot[mark=*,mark size=0.7pt,draw=black,thin] coordinates {(0,1070218) (1,5293788) (2,15474928) (3,68353449) (4,145781031) (5,97007293) (6,17010913) (7,1949622) (8,4162) (9,98) (10,181) (12,21)};\end{axis}\end{tikzpicture}\end{minipage}
\vspace{3mm}

\noindent\begin{minipage}[t]{60mm}\vspace{0pt}\sbox\hmboard{\chessboard[setfen=6rk/6B1/8/8/8/7R/3r4/K7 b - - 0 1]}\usebox\hmboard\\[1mm]
\makebox[\wd\hmboard]{\sffamily\hspace{3mm}\textbf{h\#6}\hfill(3+3)}\end{minipage}\hfill
\begin{minipage}[t]{52mm}\vspace{0pt}
\textcolor{ink}{\sffamily\textsc{Lösung}}\par\smallskip\small
\begin{tabular}{@{}r@{\enspace}l@{\enspace}l@{}}
\textbf{1.} & \symking{}xg7 & \symking{}b1\\
\textbf{2.} & \symrook{}d5 & \symking{}c2\\
\textbf{3.} & \symrook{}h5 & \symking{}d3\\
\textbf{4.} & \symking{}h6 & \symking{}e4\\
\textbf{5.} & \symrook{}h4+ & \symking{}f5\\
\textbf{6.} & \symrook{}g7 & \symrook{}xh4\#\\
\end{tabular}\par\vspace{3mm}
\qrcode[height=12mm]{https://helpman.komtera.lt/?fen=6rk/6B1/8/8/8/7R/3r4/K7&moves=6}\\{\sffamily\tiny\color{ink}scannen zum Lösen}
\end{minipage}
\vspace{2mm}

{\raggedright\noindent\textcolor{ink}{\sffamily\textsc{Themen}}\enspace \themetag{reines Matt} \themetag{Mustermatt} \themetag{Idealmatt} \themetag{Selbstblock}\par\smallskip{\footnotesize\textit{Hinweis:} die Partei am Zug steht im Diagramm im Schach; die Lösung beginnt mit einem Schlag.}\par}

\clearpage
\phantomsection\label{hm:KRNNNvk}
\begin{tcolorbox}[enhanced,colback=ink,colframe=ink,arc=0pt,boxrule=0pt,left=4mm,right=4mm,top=2.5mm,bottom=2.5mm]
\color{white}\sffamily{\LARGE\bfseries KRNNNvk}\hfill{\Large h\#6}\\[0.5mm]
{\small 6 Steine\quad Schwarz am Zug}\hfill{\small Nr.~721\qquad\href{https://helpman.komtera.lt/?fen=8/7N/8/7N/8/8/k1K5/R1N5&moves=6}{\color{white}\bfseries lösen\,\ensuremath{\nearrow}}}
\end{tcolorbox}
\vspace{2mm}
\noindent\begin{minipage}[t]{36mm}\vspace{0pt}\sffamily\small\raggedright
\hmfigure{längste eindeutige}{{\Large h\#6}}\hmfigure{eindeutig bei dieser Länge}{{\Large 12} Stellungen}\hmfigure{längstes Matt}{{\Large h\#7}}\hmfigure{in der Klasse}{15,5 Milliarden Stellungen,\\ 1,63 Milliarden mit Hilfsmatt}
\end{minipage}\hfill
\begin{minipage}[t]{80mm}\vspace{0pt}\begin{tikzpicture}\begin{axis}[width=72mm,height=42mm,ybar,bar width=1.6pt,xmin=-1,xmax=15,ymin=0,ymax=619.3,axis y line*=left,axis x line*=bottom,xlabel={\scriptsize Halbzüge bis Matt},ylabel={\scriptsize Stellungen (Mio.)},tick label style={font=\tiny},label style={font=\scriptsize},enlarge x limits=false,ylabel shift=-3pt,xlabel shift=-3pt]\addplot[fill=ink!55,draw=none] coordinates {(0,2.534) (1,11.37) (2,30.74) (3,111.5) (4,351.3) (5,420.2) (6,563) (7,81.64) (8,52.39) (9,1.58) (10,1.714) (11,5.9e-05) (12,0.2982) (14,0.07585)};\draw[densely dotted] (axis cs:12,0) -- (axis cs:12,619.3);\end{axis}\begin{axis}[width=72mm,height=42mm,axis y line*=right,axis x line=none,ymode=log,xmin=-1,xmax=15,ylabel={\scriptsize eindeutig (log)},tick label style={font=\tiny},label style={font=\scriptsize},ylabel shift=-3pt]\addplot[mark=*,mark size=0.7pt,draw=black,thin] coordinates {(0,2534273) (1,10125174) (2,24788554) (3,54683745) (4,110421953) (5,18149794) (6,6231184) (7,39820) (8,28409) (9,6) (10,1809) (12,12)};\end{axis}\end{tikzpicture}\end{minipage}
\vspace{3mm}

\noindent\begin{minipage}[t]{60mm}\vspace{0pt}\sbox\hmboard{\chessboard[setfen=8/7N/8/7N/8/8/k1K5/R1N5 b - - 0 1]}\usebox\hmboard\\[1mm]
\makebox[\wd\hmboard]{\sffamily\hspace{3mm}\textbf{h\#6}\hfill(5+1)}\end{minipage}\hfill
\begin{minipage}[t]{52mm}\vspace{0pt}
\textcolor{ink}{\sffamily\textsc{Lösung}}\par\smallskip\small
\begin{tabular}{@{}r@{\enspace}l@{\enspace}l@{}}
\textbf{1.} & \symking{}xa1 & \symking{}d1\\
\textbf{2.} & \symking{}b2 & \symknight{}b3\\
\textbf{3.} & \symking{}c3 & \symking{}c1\\
\textbf{4.} & \symking{}d3 & \symknight{}c5+\\
\textbf{5.} & \symking{}e2 & \symknight{}g3+\\
\textbf{6.} & \symking{}e1 & \symknight{}d3\#\\
\end{tabular}\par\vspace{3mm}
\qrcode[height=12mm]{https://helpman.komtera.lt/?fen=8/7N/8/7N/8/8/k1K5/R1N5&moves=6}\\{\sffamily\tiny\color{ink}scannen zum Lösen}
\end{minipage}
\vspace{2mm}

{\raggedright\noindent\textcolor{ink}{\sffamily\textsc{Themen}}\enspace \themetag{reines Matt} \themetag{Spiegelmatt} \themetag{ein Stein zieht}\par\smallskip{\footnotesize\textit{Hinweis:} das Diagramm hat keinen legalen letzten Zug, kann also in keiner Partie entstehen.}\par}

\clearpage
\phantomsection\label{hm:KRNNvkn}
\begin{tcolorbox}[enhanced,colback=ink,colframe=ink,arc=0pt,boxrule=0pt,left=4mm,right=4mm,top=2.5mm,bottom=2.5mm]
\color{white}\sffamily{\LARGE\bfseries KRNNvkn}\hfill{\Large h\#6}\\[0.5mm]
{\small 6 Steine\quad Schwarz am Zug}\hfill{\small Nr.~722\qquad\href{https://helpman.komtera.lt/?fen=3Rk1Nn/8/7N/8/8/8/8/2K5&moves=6}{\color{white}\bfseries lösen\,\ensuremath{\nearrow}}}
\end{tcolorbox}
\vspace{2mm}
\noindent\begin{minipage}[t]{36mm}\vspace{0pt}\sffamily\small\raggedright
\hmfigure{längste eindeutige}{{\Large h\#6}}\hmfigure{eindeutig bei dieser Länge}{{\Large 6} Stellungen}\hmfigure{längstes Matt}{{\Large h\#7}}\hmfigure{in der Klasse}{15,5 Milliarden Stellungen,\\ 4,82 Milliarden mit Hilfsmatt}
\end{minipage}\hfill
\begin{minipage}[t]{80mm}\vspace{0pt}\begin{tikzpicture}\begin{axis}[width=72mm,height=42mm,ybar,bar width=1.6pt,xmin=-1,xmax=15,ymin=0,ymax=1.872,axis y line*=left,axis x line*=bottom,xlabel={\scriptsize Halbzüge bis Matt},ylabel={\scriptsize Stellungen (Mrd.)},tick label style={font=\tiny},label style={font=\scriptsize},enlarge x limits=false,ylabel shift=-3pt,xlabel shift=-3pt]\addplot[fill=ink!55,draw=none] coordinates {(0,0.003663) (1,0.01885) (2,0.05877) (3,0.2563) (4,0.737) (5,1.204) (6,1.702) (7,0.547) (8,0.2689) (9,0.02007) (10,0.003164) (11,2.379e-05) (12,0.0002954) (13,8e-09) (14,6.837e-06)};\draw[densely dotted] (axis cs:12,0) -- (axis cs:12,1.872);\end{axis}\begin{axis}[width=72mm,height=42mm,axis y line*=right,axis x line=none,ymode=log,xmin=-1,xmax=15,ylabel={\scriptsize eindeutig (log)},tick label style={font=\tiny},label style={font=\scriptsize},ylabel shift=-3pt]\addplot[mark=*,mark size=0.7pt,draw=black,thin] coordinates {(0,3663338) (1,17732502) (2,34006608) (3,90383525) (4,177215754) (5,70548415) (6,23986772) (7,734726) (8,78703) (9,764) (10,1714) (12,6)};\end{axis}\end{tikzpicture}\end{minipage}
\vspace{3mm}

\noindent\begin{minipage}[t]{60mm}\vspace{0pt}\sbox\hmboard{\chessboard[setfen=3Rk1Nn/8/7N/8/8/8/8/2K5 b - - 0 1]}\usebox\hmboard\\[1mm]
\makebox[\wd\hmboard]{\sffamily\hspace{3mm}\textbf{h\#6}\hfill(4+2)}\end{minipage}\hfill
\begin{minipage}[t]{52mm}\vspace{0pt}
\textcolor{ink}{\sffamily\textsc{Lösung}}\par\smallskip\small
\begin{tabular}{@{}r@{\enspace}l@{\enspace}l@{}}
\textbf{1.} & \symking{}xd8 & \symking{}d2\\
\textbf{2.} & \symking{}d7 & \symking{}e3\\
\textbf{3.} & \symking{}e8 & \symking{}f4\\
\textbf{4.} & \symking{}f8 & \symking{}g5\\
\textbf{5.} & \symking{}g7 & \symknight{}f5+\\
\textbf{6.} & \symking{}h7 & \symknight{}f6\#\\
\end{tabular}\par\vspace{3mm}
\qrcode[height=12mm]{https://helpman.komtera.lt/?fen=3Rk1Nn/8/7N/8/8/8/8/2K5&moves=6}\\{\sffamily\tiny\color{ink}scannen zum Lösen}
\end{minipage}
\vspace{2mm}

{\raggedright\noindent\textcolor{ink}{\sffamily\textsc{Themen}}\enspace \themetag{Rundlauf} \themetag{ein Stein zieht}\par\smallskip{\footnotesize\textit{Hinweis:} die Partei am Zug steht im Diagramm im Schach; die Lösung beginnt mit einem Schlag.}\par}

\clearpage
\phantomsection\label{hm:KRNNvkr}
\begin{tcolorbox}[enhanced,colback=ink,colframe=ink,arc=0pt,boxrule=0pt,left=4mm,right=4mm,top=2.5mm,bottom=2.5mm]
\color{white}\sffamily{\LARGE\bfseries KRNNvkr}\hfill{\Large h\#6}\\[0.5mm]
{\small 6 Steine\quad Schwarz am Zug}\hfill{\small Nr.~723\qquad\href{https://helpman.komtera.lt/?fen=8/8/1r3N1R/7k/8/5N2/8/K7&moves=6}{\color{white}\bfseries lösen\,\ensuremath{\nearrow}}}
\end{tcolorbox}
\vspace{2mm}
\noindent\begin{minipage}[t]{36mm}\vspace{0pt}\sffamily\small\raggedright
\hmfigure{längste eindeutige}{{\Large h\#6}}\hmfigure{eindeutig bei dieser Länge}{{\Large 122} Stellungen}\hmfigure{längstes Matt}{{\Large h\#8}}\hmfigure{in der Klasse}{15,5 Milliarden Stellungen,\\ 4,47 Milliarden mit Hilfsmatt}
\end{minipage}\hfill
\begin{minipage}[t]{80mm}\vspace{0pt}\begin{tikzpicture}\begin{axis}[width=72mm,height=42mm,ybar,bar width=1.6pt,xmin=-1,xmax=17,ymin=0,ymax=1.544,axis y line*=left,axis x line*=bottom,xlabel={\scriptsize Halbzüge bis Matt},ylabel={\scriptsize Stellungen (Mrd.)},tick label style={font=\tiny},label style={font=\scriptsize},enlarge x limits=false,ylabel shift=-3pt,xlabel shift=-3pt]\addplot[fill=ink!55,draw=none] coordinates {(0,0.002614) (1,0.01366) (2,0.05758) (3,0.279) (4,0.9003) (5,1.34) (6,1.403) (7,0.4108) (8,0.06008) (9,0.002226) (10,0.0008095) (11,8.11e-06) (12,0.0002282) (13,8.4e-08) (14,4.345e-05) (15,7.4e-08) (16,3.52e-07)};\draw[densely dotted] (axis cs:12,0) -- (axis cs:12,1.544);\end{axis}\begin{axis}[width=72mm,height=42mm,axis y line*=right,axis x line=none,ymode=log,xmin=-1,xmax=17,ylabel={\scriptsize eindeutig (log)},tick label style={font=\tiny},label style={font=\scriptsize},ylabel shift=-3pt]\addplot[mark=*,mark size=0.7pt,draw=black,thin] coordinates {(0,2613765) (1,12941909) (2,35348828) (3,101685832) (4,211598445) (5,68837480) (6,10299444) (7,284380) (8,12348) (9,287) (10,1559) (12,122)};\end{axis}\end{tikzpicture}\end{minipage}
\vspace{3mm}

\noindent\begin{minipage}[t]{60mm}\vspace{0pt}\sbox\hmboard{\chessboard[setfen=8/8/1r3N1R/7k/8/5N2/8/K7 b - - 0 1]}\usebox\hmboard\\[1mm]
\makebox[\wd\hmboard]{\sffamily\hspace{3mm}\textbf{h\#6}\hfill(4+2)}\end{minipage}\hfill
\begin{minipage}[t]{52mm}\vspace{0pt}
\textcolor{ink}{\sffamily\textsc{Lösung}}\par\smallskip\small
\begin{tabular}{@{}r@{\enspace}l@{\enspace}l@{}}
\textbf{1.} & \symking{}xh6 & \symknight{}d4\\
\textbf{2.} & \symking{}g5 & \symknight{}d7\\
\textbf{3.} & \symking{}f4 & \symknight{}xb6\\
\textbf{4.} & \symking{}e3 & \symknight{}d5+\\
\textbf{5.} & \symking{}d2 & \symknight{}e3\\
\textbf{6.} & \symking{}c1 & \symknight{}b3\#\\
\end{tabular}\par\vspace{3mm}
\qrcode[height=12mm]{https://helpman.komtera.lt/?fen=8/8/1r3N1R/7k/8/5N2/8/K7&moves=6}\\{\sffamily\tiny\color{ink}scannen zum Lösen}
\end{minipage}
\vspace{2mm}

{\raggedright\noindent\textcolor{ink}{\sffamily\textsc{Themen}}\enspace \themetag{reines Matt} \themetag{Mustermatt} \themetag{Idealmatt} \themetag{Spiegelmatt} \themetag{ein Stein zieht} \themetag{Umnov}\par\smallskip{\footnotesize\textit{Hinweis:} das Diagramm hat keinen legalen letzten Zug, kann also in keiner Partie entstehen.}\par}

\clearpage
\phantomsection\label{hm:KRNvkqq}
\begin{tcolorbox}[enhanced,colback=ink,colframe=ink,arc=0pt,boxrule=0pt,left=4mm,right=4mm,top=2.5mm,bottom=2.5mm]
\color{white}\sffamily{\LARGE\bfseries KRNvkqq}\hfill{\Large h\#6}\\[0.5mm]
{\small 6 Steine\quad Schwarz am Zug}\hfill{\small Nr.~724\qquad\href{https://helpman.komtera.lt/?fen=8/8/8/7q/7k/5NR1/2q5/K7&moves=6}{\color{white}\bfseries lösen\,\ensuremath{\nearrow}}}
\end{tcolorbox}
\vspace{2mm}
\noindent\begin{minipage}[t]{36mm}\vspace{0pt}\sffamily\small\raggedright
\hmfigure{längste eindeutige}{{\Large h\#6}}\hmfigure{eindeutig bei dieser Länge}{{\Large 22} Stellungen}\hmfigure{längstes Matt}{{\Large h\#8}}\hmfigure{in der Klasse}{15,5 Milliarden Stellungen,\\ 3,5 Milliarden mit Hilfsmatt}
\end{minipage}\hfill
\begin{minipage}[t]{80mm}\vspace{0pt}\begin{tikzpicture}\begin{axis}[width=72mm,height=42mm,ybar,bar width=1.6pt,xmin=-1,xmax=17,ymin=0,ymax=1.318,axis y line*=left,axis x line*=bottom,xlabel={\scriptsize Halbzüge bis Matt},ylabel={\scriptsize Stellungen (Mrd.)},tick label style={font=\tiny},label style={font=\scriptsize},enlarge x limits=false,ylabel shift=-3pt,xlabel shift=-3pt]\addplot[fill=ink!55,draw=none] coordinates {(0,0.0003024) (1,0.001779) (2,0.009711) (3,0.07398) (4,0.179) (5,0.6847) (6,0.7965) (7,1.199) (8,0.3385) (9,0.2099) (10,0.00871) (11,0.002488) (12,2.591e-05) (13,5e-09) (14,1.792e-06) (16,1e-09)};\draw[densely dotted] (axis cs:12,0) -- (axis cs:12,1.318);\end{axis}\begin{axis}[width=72mm,height=42mm,axis y line*=right,axis x line=none,ymode=log,xmin=-1,xmax=17,ylabel={\scriptsize eindeutig (log)},tick label style={font=\tiny},label style={font=\scriptsize},ylabel shift=-3pt]\addplot[mark=*,mark size=0.7pt,draw=black,thin] coordinates {(0,302410) (1,1763736) (2,4642049) (3,26154603) (4,41004203) (5,66619092) (6,18786056) (7,9499508) (8,295808) (9,75303) (10,276) (12,22)};\end{axis}\end{tikzpicture}\end{minipage}
\vspace{3mm}

\noindent\begin{minipage}[t]{60mm}\vspace{0pt}\sbox\hmboard{\chessboard[setfen=8/8/8/7q/7k/5NR1/2q5/K7 b - - 0 1]}\usebox\hmboard\\[1mm]
\makebox[\wd\hmboard]{\sffamily\hspace{3mm}\textbf{h\#6}\hfill(3+3)}\end{minipage}\hfill
\begin{minipage}[t]{52mm}\vspace{0pt}
\textcolor{ink}{\sffamily\textsc{Lösung}}\par\smallskip\small
\begin{tabular}{@{}r@{\enspace}l@{\enspace}l@{}}
\textbf{1.} & \symqueen{}xf3 & \symrook{}g2\\
\textbf{2.} & \symking{}h3 & \symrook{}xc2\\
\textbf{3.} & \symqueen{}f2 & \symking{}b2\\
\textbf{4.} & \symking{}g2 & \symking{}c3\\
\textbf{5.} & \symking{}f1 & \symking{}d3\\
\textbf{6.} & \symking{}e1 & \symrook{}c1\#\\
\end{tabular}\par\vspace{3mm}
\qrcode[height=12mm]{https://helpman.komtera.lt/?fen=8/8/8/7q/7k/5NR1/2q5/K7&moves=6}\\{\sffamily\tiny\color{ink}scannen zum Lösen}
\end{minipage}
\vspace{2mm}

{\raggedright\noindent\textcolor{ink}{\sffamily\textsc{Themen}}\enspace \themetag{reines Matt} \themetag{Mustermatt} \themetag{Idealmatt} \themetag{Selbstblock} \themetag{ohne Schach}\par\smallskip{\footnotesize\textit{Hinweis:} die Partei am Zug steht im Diagramm im Schach; die Lösung beginnt mit einem Schlag.}\par}

\clearpage
\phantomsection\label{hm:KRRBBvk}
\begin{tcolorbox}[enhanced,colback=ink,colframe=ink,arc=0pt,boxrule=0pt,left=4mm,right=4mm,top=2.5mm,bottom=2.5mm]
\color{white}\sffamily{\LARGE\bfseries KRRBBvk}\hfill{\Large h\#6}\\[0.5mm]
{\small 6 Steine\quad Schwarz am Zug}\hfill{\small Nr.~725\qquad\href{https://helpman.komtera.lt/?fen=8/7R/7k/8/7B/6R1/7B/K7&moves=6}{\color{white}\bfseries lösen\,\ensuremath{\nearrow}}}
\end{tcolorbox}
\vspace{2mm}
\noindent\begin{minipage}[t]{36mm}\vspace{0pt}\sffamily\small\raggedright
\hmfigure{längste eindeutige}{{\Large h\#6}}\hmfigure{eindeutig bei dieser Länge}{{\Large 29} Stellungen}\hmfigure{längstes Matt}{{\Large h\#8}}\hmfigure{in der Klasse}{15,5 Milliarden Stellungen,\\ 2,22 Milliarden mit Hilfsmatt}
\end{minipage}\hfill
\begin{minipage}[t]{80mm}\vspace{0pt}\begin{tikzpicture}\begin{axis}[width=72mm,height=42mm,ybar,bar width=1.6pt,xmin=-1,xmax=17,ymin=0,ymax=1.157,axis y line*=left,axis x line*=bottom,xlabel={\scriptsize Halbzüge bis Matt},ylabel={\scriptsize Stellungen (Mrd.)},tick label style={font=\tiny},label style={font=\scriptsize},enlarge x limits=false,ylabel shift=-3pt,xlabel shift=-3pt]\addplot[fill=ink!55,draw=none] coordinates {(0,0.02505) (1,0.07304) (2,0.2318) (3,0.4431) (4,1.052) (5,0.2131) (6,0.1768) (7,0.0006606) (8,0.005038) (9,1.143e-06) (10,0.0009129) (11,1.4e-08) (12,7.589e-05) (14,1.487e-06) (16,1e-09)};\draw[densely dotted] (axis cs:12,0) -- (axis cs:12,1.157);\end{axis}\begin{axis}[width=72mm,height=42mm,axis y line*=right,axis x line=none,ymode=log,xmin=-1,xmax=17,ylabel={\scriptsize eindeutig (log)},tick label style={font=\tiny},label style={font=\scriptsize},ylabel shift=-3pt]\addplot[mark=*,mark size=0.7pt,draw=black,thin] coordinates {(0,25052822) (1,55476948) (2,138462957) (3,81615503) (4,75913236) (5,407125) (6,1170989) (7,96) (8,34335) (10,1332) (12,29)};\end{axis}\end{tikzpicture}\end{minipage}
\vspace{3mm}

\noindent\begin{minipage}[t]{60mm}\vspace{0pt}\sbox\hmboard{\chessboard[setfen=8/7R/7k/8/7B/6R1/7B/K7 b - - 0 1]}\usebox\hmboard\\[1mm]
\makebox[\wd\hmboard]{\sffamily\hspace{3mm}\textbf{h\#6}\hfill(5+1)}\end{minipage}\hfill
\begin{minipage}[t]{52mm}\vspace{0pt}
\textcolor{ink}{\sffamily\textsc{Lösung}}\par\smallskip\small
\begin{tabular}{@{}r@{\enspace}l@{\enspace}l@{}}
\textbf{1.} & \symking{}xh7 & \symking{}b2\\
\textbf{2.} & \symking{}h6 & \symking{}c3\\
\textbf{3.} & \symking{}h5 & \symking{}d4\\
\textbf{4.} & \symking{}xh4 & \symking{}e5\\
\textbf{5.} & \symking{}h5 & \symking{}f6\\
\textbf{6.} & \symking{}h6 & \symrook{}h3\#\\
\end{tabular}\par\vspace{3mm}
\qrcode[height=12mm]{https://helpman.komtera.lt/?fen=8/7R/7k/8/7B/6R1/7B/K7&moves=6}\\{\sffamily\tiny\color{ink}scannen zum Lösen}
\end{minipage}
\vspace{2mm}

{\raggedright\noindent\textcolor{ink}{\sffamily\textsc{Themen}}\enspace \themetag{reines Matt} \themetag{Spiegelmatt} \themetag{Switchback} \themetag{ein Stein zieht} \themetag{ohne Schach}\par\smallskip{\footnotesize\textit{Hinweis:} die Partei am Zug steht im Diagramm im Schach; die Lösung beginnt mit einem Schlag.}\par}

\clearpage
\phantomsection\label{hm:KRRBPvk}
\begin{tcolorbox}[enhanced,colback=ink,colframe=ink,arc=0pt,boxrule=0pt,left=4mm,right=4mm,top=2.5mm,bottom=2.5mm]
\color{white}\sffamily{\LARGE\bfseries KRRBPvk}\hfill{\Large h\#6}\\[0.5mm]
{\small 6 Steine\quad Schwarz am Zug}\hfill{\small Nr.~726\qquad\href{https://helpman.komtera.lt/?fen=5B2/7R/7k/8/8/P7/8/K5R1&moves=6}{\color{white}\bfseries lösen\,\ensuremath{\nearrow}}}
\end{tcolorbox}
\vspace{2mm}
\noindent\begin{minipage}[t]{36mm}\vspace{0pt}\sffamily\small\raggedright
\hmfigure{längste eindeutige}{{\Large h\#6}}\hmfigure{eindeutig bei dieser Länge}{{\Large 182} Stellungen}\hmfigure{längstes Matt}{{\Large h\#7}}\hmfigure{in der Klasse}{45,4 Milliarden Stellungen,\\ 13,9 Milliarden mit Hilfsmatt}
\end{minipage}\hfill
\begin{minipage}[t]{80mm}\vspace{0pt}\begin{tikzpicture}\begin{axis}[width=72mm,height=42mm,ybar,bar width=1.6pt,xmin=-1,xmax=15,ymin=0,ymax=6.702,axis y line*=left,axis x line*=bottom,xlabel={\scriptsize Halbzüge bis Matt},ylabel={\scriptsize Stellungen (Mrd.)},tick label style={font=\tiny},label style={font=\scriptsize},enlarge x limits=false,ylabel shift=-3pt,xlabel shift=-3pt]\addplot[fill=ink!55,draw=none] coordinates {(0,0.1058) (1,0.3818) (2,1.1) (3,2.663) (4,6.092) (5,1.84) (6,1.642) (7,0.02129) (8,0.04689) (9,6.451e-06) (10,0.01071) (11,1.96e-07) (12,0.0006362) (13,1e-09) (14,2.584e-06)};\draw[densely dotted] (axis cs:12,0) -- (axis cs:12,6.702);\end{axis}\begin{axis}[width=72mm,height=42mm,axis y line*=right,axis x line=none,ymode=log,xmin=-1,xmax=15,ylabel={\scriptsize eindeutig (log)},tick label style={font=\tiny},label style={font=\scriptsize},ylabel shift=-3pt]\addplot[mark=*,mark size=0.7pt,draw=black,thin] coordinates {(0,105833917) (1,298918695) (2,647239709) (3,481391965) (4,498834478) (5,7376522) (6,7360538) (7,519) (8,454906) (9,8) (10,18909) (12,182)};\end{axis}\end{tikzpicture}\end{minipage}
\vspace{3mm}

\noindent\begin{minipage}[t]{60mm}\vspace{0pt}\sbox\hmboard{\chessboard[setfen=5B2/7R/7k/8/8/P7/8/K5R1 b - - 0 1]}\usebox\hmboard\\[1mm]
\makebox[\wd\hmboard]{\sffamily\hspace{3mm}\textbf{h\#6}\hfill(5+1)}\end{minipage}\hfill
\begin{minipage}[t]{52mm}\vspace{0pt}
\textcolor{ink}{\sffamily\textsc{Lösung}}\par\smallskip\small
\begin{tabular}{@{}r@{\enspace}l@{\enspace}l@{}}
\textbf{1.} & \symking{}xh7 & a4\\
\textbf{2.} & \symking{}h8 & a5\\
\textbf{3.} & \symking{}h7 & a6\\
\textbf{4.} & \symking{}h8 & a7\\
\textbf{5.} & \symking{}h7 & a8=\symqueen{}\\
\textbf{6.} & \symking{}h8 & \symqueen{}h1\#\\
\end{tabular}\par\vspace{3mm}
\qrcode[height=12mm]{https://helpman.komtera.lt/?fen=5B2/7R/7k/8/8/P7/8/K5R1&moves=6}\\{\sffamily\tiny\color{ink}scannen zum Lösen}
\end{minipage}
\vspace{2mm}

{\raggedright\noindent\textcolor{ink}{\sffamily\textsc{Themen}}\enspace \themetag{Spiegelmatt} \themetag{Umwandlung} \themetag{Switchback} \themetag{ein Stein zieht} \themetag{ohne Schach}\par\smallskip{\footnotesize\textit{Hinweis:} die Partei am Zug steht im Diagramm im Schach; die Lösung beginnt mit einem Schlag.}\par}

\clearpage
\phantomsection\label{hm:KRRBvkb}
\begin{tcolorbox}[enhanced,colback=ink,colframe=ink,arc=0pt,boxrule=0pt,left=4mm,right=4mm,top=2.5mm,bottom=2.5mm]
\color{white}\sffamily{\LARGE\bfseries KRRBvkb}\hfill{\Large h\#6}\\[0.5mm]
{\small 6 Steine\quad Schwarz am Zug}\hfill{\small Nr.~727\qquad\href{https://helpman.komtera.lt/?fen=8/4B3/7R/6Rk/7b/8/8/K7&moves=6}{\color{white}\bfseries lösen\,\ensuremath{\nearrow}}}
\end{tcolorbox}
\vspace{2mm}
\noindent\begin{minipage}[t]{36mm}\vspace{0pt}\sffamily\small\raggedright
\hmfigure{längste eindeutige}{{\Large h\#6}}\hmfigure{eindeutig bei dieser Länge}{{\Large 26} Stellungen}\hmfigure{längstes Matt}{{\Large h\#8}}\hmfigure{in der Klasse}{15,5 Milliarden Stellungen,\\ 4,31 Milliarden mit Hilfsmatt}
\end{minipage}\hfill
\begin{minipage}[t]{80mm}\vspace{0pt}\begin{tikzpicture}\begin{axis}[width=72mm,height=42mm,ybar,bar width=1.6pt,xmin=-1,xmax=17,ymin=0,ymax=2.077,axis y line*=left,axis x line*=bottom,xlabel={\scriptsize Halbzüge bis Matt},ylabel={\scriptsize Stellungen (Mrd.)},tick label style={font=\tiny},label style={font=\scriptsize},enlarge x limits=false,ylabel shift=-3pt,xlabel shift=-3pt]\addplot[fill=ink!55,draw=none] coordinates {(0,0.02124) (1,0.07665) (2,0.2939) (3,0.8717) (4,1.888) (5,0.7138) (6,0.421) (7,0.01333) (8,0.007872) (9,2.374e-06) (10,0.002393) (11,1.48e-07) (12,0.0004081) (13,4.8e-08) (14,6.215e-06) (15,5e-09) (16,5e-09)};\draw[densely dotted] (axis cs:12,0) -- (axis cs:12,2.077);\end{axis}\begin{axis}[width=72mm,height=42mm,axis y line*=right,axis x line=none,ymode=log,xmin=-1,xmax=17,ylabel={\scriptsize eindeutig (log)},tick label style={font=\tiny},label style={font=\scriptsize},ylabel shift=-3pt]\addplot[mark=*,mark size=0.7pt,draw=black,thin] coordinates {(0,21243067) (1,66709452) (2,131838961) (3,153755299) (4,122403996) (5,2677637) (6,804601) (7,530) (8,62556) (9,2) (10,9315) (12,26)};\end{axis}\end{tikzpicture}\end{minipage}
\vspace{3mm}

\noindent\begin{minipage}[t]{60mm}\vspace{0pt}\sbox\hmboard{\chessboard[setfen=8/4B3/7R/6Rk/7b/8/8/K7 b - - 0 1]}\usebox\hmboard\\[1mm]
\makebox[\wd\hmboard]{\sffamily\hspace{3mm}\textbf{h\#6}\hfill(4+2)}\end{minipage}\hfill
\begin{minipage}[t]{52mm}\vspace{0pt}
\textcolor{ink}{\sffamily\textsc{Lösung}}\par\smallskip\small
\begin{tabular}{@{}r@{\enspace}l@{\enspace}l@{}}
\textbf{1.} & \symking{}xh6 & \symrook{}g8\\
\textbf{2.} & \symking{}h5 & \symrook{}e8\\
\textbf{3.} & \symking{}g4 & \symking{}b2\\
\textbf{4.} & \symking{}f3 & \symking{}c3\\
\textbf{5.} & \symking{}e2 & \symbishop{}xh4+\\
\textbf{6.} & \symking{}d1 & \symrook{}e1\#\\
\end{tabular}\par\vspace{3mm}
\qrcode[height=12mm]{https://helpman.komtera.lt/?fen=8/4B3/7R/6Rk/7b/8/8/K7&moves=6}\\{\sffamily\tiny\color{ink}scannen zum Lösen}
\end{minipage}
\vspace{2mm}

{\raggedright\noindent\textcolor{ink}{\sffamily\textsc{Themen}}\enspace \themetag{reines Matt} \themetag{Mustermatt} \themetag{Idealmatt} \themetag{Switchback} \themetag{ein Stein zieht}\par\smallskip{\footnotesize\textit{Hinweis:} das Diagramm hat keinen legalen letzten Zug, kann also in keiner Partie entstehen.}\par}

\clearpage
\phantomsection\label{hm:KRRBvkn}
\begin{tcolorbox}[enhanced,colback=ink,colframe=ink,arc=0pt,boxrule=0pt,left=4mm,right=4mm,top=2.5mm,bottom=2.5mm]
\color{white}\sffamily{\LARGE\bfseries KRRBvkn}\hfill{\Large h\#6}\\[0.5mm]
{\small 6 Steine\quad Schwarz am Zug}\hfill{\small Nr.~728\qquad\href{https://helpman.komtera.lt/?fen=8/7R/7k/8/5B2/8/1n6/K5R1&moves=6}{\color{white}\bfseries lösen\,\ensuremath{\nearrow}}}
\end{tcolorbox}
\vspace{2mm}
\noindent\begin{minipage}[t]{36mm}\vspace{0pt}\sffamily\small\raggedright
\hmfigure{längste eindeutige}{{\Large h\#6}}\hmfigure{eindeutig bei dieser Länge}{{\Large 10} Stellungen}\hmfigure{längstes Matt}{{\Large h\#7}}\hmfigure{in der Klasse}{15,5 Milliarden Stellungen,\\ 4,45 Milliarden mit Hilfsmatt}
\end{minipage}\hfill
\begin{minipage}[t]{80mm}\vspace{0pt}\begin{tikzpicture}\begin{axis}[width=72mm,height=42mm,ybar,bar width=1.6pt,xmin=-1,xmax=15,ymin=0,ymax=2.128,axis y line*=left,axis x line*=bottom,xlabel={\scriptsize Halbzüge bis Matt},ylabel={\scriptsize Stellungen (Mrd.)},tick label style={font=\tiny},label style={font=\scriptsize},enlarge x limits=false,ylabel shift=-3pt,xlabel shift=-3pt]\addplot[fill=ink!55,draw=none] coordinates {(0,0.02588) (1,0.09468) (2,0.3266) (3,0.8909) (4,1.935) (5,0.679) (6,0.4721) (7,0.0109) (8,0.009441) (9,1.499e-06) (10,0.003603) (11,4e-08) (12,0.0003269) (14,5.47e-07)};\draw[densely dotted] (axis cs:12,0) -- (axis cs:12,2.128);\end{axis}\begin{axis}[width=72mm,height=42mm,axis y line*=right,axis x line=none,ymode=log,xmin=-1,xmax=15,ylabel={\scriptsize eindeutig (log)},tick label style={font=\tiny},label style={font=\scriptsize},ylabel shift=-3pt]\addplot[mark=*,mark size=0.7pt,draw=black,thin] coordinates {(0,25882080) (1,80292193) (2,142323375) (3,138552394) (4,123345574) (5,2255370) (6,876784) (7,1211) (8,107638) (9,6) (10,9368) (12,10)};\end{axis}\end{tikzpicture}\end{minipage}
\vspace{3mm}

\noindent\begin{minipage}[t]{60mm}\vspace{0pt}\sbox\hmboard{\chessboard[setfen=8/7R/7k/8/5B2/8/1n6/K5R1 b - - 0 1]}\usebox\hmboard\\[1mm]
\makebox[\wd\hmboard]{\sffamily\hspace{3mm}\textbf{h\#6}\hfill(4+2)}\end{minipage}\hfill
\begin{minipage}[t]{52mm}\vspace{0pt}
\textcolor{ink}{\sffamily\textsc{Lösung}}\par\smallskip\small
\begin{tabular}{@{}r@{\enspace}l@{\enspace}l@{}}
\textbf{1.} & \symking{}xh7 & \symking{}b1\\
\textbf{2.} & \symknight{}d3 & \symking{}c2\\
\textbf{3.} & \symknight{}xf4 & \symking{}d2\\
\textbf{4.} & \symking{}h6 & \symking{}e3\\
\textbf{5.} & \symking{}h5 & \symking{}xf4\\
\textbf{6.} & \symking{}h4 & \symrook{}h1\#\\
\end{tabular}\par\vspace{3mm}
\qrcode[height=12mm]{https://helpman.komtera.lt/?fen=8/7R/7k/8/5B2/8/1n6/K5R1&moves=6}\\{\sffamily\tiny\color{ink}scannen zum Lösen}
\end{minipage}
\vspace{2mm}

{\raggedright\noindent\textcolor{ink}{\sffamily\textsc{Themen}}\enspace \themetag{reines Matt} \themetag{Mustermatt} \themetag{Idealmatt} \themetag{Spiegelmatt} \themetag{Switchback} \themetag{ohne Schach}\par\smallskip{\footnotesize\textit{Hinweis:} das Diagramm hat keinen legalen letzten Zug, kann also in keiner Partie entstehen.}\par}

\clearpage
\phantomsection\label{hm:KRRBvkp}
\begin{tcolorbox}[enhanced,colback=ink,colframe=ink,arc=0pt,boxrule=0pt,left=4mm,right=4mm,top=2.5mm,bottom=2.5mm]
\color{white}\sffamily{\LARGE\bfseries KRRBvkp}\hfill{\Large h\#6}\\[0.5mm]
{\small 6 Steine\quad Schwarz am Zug}\hfill{\small Nr.~729\qquad\href{https://helpman.komtera.lt/?fen=8/2p4R/7k/8/8/8/3B4/K5R1&moves=6}{\color{white}\bfseries lösen\,\ensuremath{\nearrow}}}
\end{tcolorbox}
\vspace{2mm}
\noindent\begin{minipage}[t]{36mm}\vspace{0pt}\sffamily\small\raggedright
\hmfigure{längste eindeutige}{{\Large h\#6}}\hmfigure{eindeutig bei dieser Länge}{{\Large 255} Stellungen}\hmfigure{längstes Matt}{{\Large h\#8}}\hmfigure{in der Klasse}{45,4 Milliarden Stellungen,\\ 13,9 Milliarden mit Hilfsmatt}
\end{minipage}\hfill
\begin{minipage}[t]{80mm}\vspace{0pt}\begin{tikzpicture}\begin{axis}[width=72mm,height=42mm,ybar,bar width=1.6pt,xmin=-1,xmax=17,ymin=0,ymax=6.666,axis y line*=left,axis x line*=bottom,xlabel={\scriptsize Halbzüge bis Matt},ylabel={\scriptsize Stellungen (Mrd.)},tick label style={font=\tiny},label style={font=\scriptsize},enlarge x limits=false,ylabel shift=-3pt,xlabel shift=-3pt]\addplot[fill=ink!55,draw=none] coordinates {(0,0.09728) (1,0.3503) (2,1.103) (3,2.811) (4,6.06) (5,1.859) (6,1.503) (7,0.02779) (8,0.03401) (9,1.022e-05) (10,0.01125) (11,1.049e-06) (12,0.001261) (13,8e-09) (14,1.813e-05) (16,1e-09)};\draw[densely dotted] (axis cs:12,0) -- (axis cs:12,6.666);\end{axis}\begin{axis}[width=72mm,height=42mm,axis y line*=right,axis x line=none,ymode=log,xmin=-1,xmax=17,ylabel={\scriptsize eindeutig (log)},tick label style={font=\tiny},label style={font=\scriptsize},ylabel shift=-3pt]\addplot[mark=*,mark size=0.7pt,draw=black,thin] coordinates {(0,97275632) (1,291749154) (2,554298515) (3,470538755) (4,413752751) (5,7372829) (6,3865690) (7,1658) (8,356633) (9,53) (10,18556) (12,255)};\end{axis}\end{tikzpicture}\end{minipage}
\vspace{3mm}

\noindent\begin{minipage}[t]{60mm}\vspace{0pt}\sbox\hmboard{\chessboard[setfen=8/2p4R/7k/8/8/8/3B4/K5R1 b - - 0 1]}\usebox\hmboard\\[1mm]
\makebox[\wd\hmboard]{\sffamily\hspace{3mm}\textbf{h\#6}\hfill(4+2)}\end{minipage}\hfill
\begin{minipage}[t]{52mm}\vspace{0pt}
\textcolor{ink}{\sffamily\textsc{Lösung}}\par\smallskip\small
\begin{tabular}{@{}r@{\enspace}l@{\enspace}l@{}}
\textbf{1.} & \symking{}xh7 & \symking{}b2\\
\textbf{2.} & c5 & \symking{}c3\\
\textbf{3.} & c4 & \symking{}d4\\
\textbf{4.} & c3 & \symking{}e5\\
\textbf{5.} & cxd2 & \symking{}f6\\
\textbf{6.} & \symking{}h6 & \symrook{}h1\#\\
\end{tabular}\par\vspace{3mm}
\qrcode[height=12mm]{https://helpman.komtera.lt/?fen=8/2p4R/7k/8/8/8/3B4/K5R1&moves=6}\\{\sffamily\tiny\color{ink}scannen zum Lösen}
\end{minipage}
\vspace{2mm}

{\raggedright\noindent\textcolor{ink}{\sffamily\textsc{Themen}}\enspace \themetag{reines Matt} \themetag{Mustermatt} \themetag{Spiegelmatt} \themetag{Switchback} \themetag{ohne Schach} \themetag{Umnov}\par\smallskip{\footnotesize\textit{Hinweis:} das Diagramm hat keinen legalen letzten Zug, kann also in keiner Partie entstehen.}\par}

\clearpage
\phantomsection\label{hm:KRRPPvk}
\begin{tcolorbox}[enhanced,colback=ink,colframe=ink,arc=0pt,boxrule=0pt,left=4mm,right=4mm,top=2.5mm,bottom=2.5mm]
\color{white}\sffamily{\LARGE\bfseries KRRPPvk}\hfill{\Large h\#6}\\[0.5mm]
{\small 6 Steine\quad Schwarz am Zug}\hfill{\small Nr.~730\qquad\href{https://helpman.komtera.lt/?fen=7R/7k/8/8/6R1/8/1P5P/K7&moves=6}{\color{white}\bfseries lösen\,\ensuremath{\nearrow}}}
\end{tcolorbox}
\vspace{2mm}
\noindent\begin{minipage}[t]{36mm}\vspace{0pt}\sffamily\small\raggedright
\hmfigure{längste eindeutige}{{\Large h\#6}}\hmfigure{eindeutig bei dieser Länge}{{\Large 41} Stellungen}\hmfigure{längstes Matt}{{\Large h\#7}}\hmfigure{in der Klasse}{34,1 Milliarden Stellungen,\\ 5,41 Milliarden mit Hilfsmatt}
\end{minipage}\hfill
\begin{minipage}[t]{80mm}\vspace{0pt}\begin{tikzpicture}\begin{axis}[width=72mm,height=42mm,ybar,bar width=1.6pt,xmin=-1,xmax=15,ymin=0,ymax=2.478,axis y line*=left,axis x line*=bottom,xlabel={\scriptsize Halbzüge bis Matt},ylabel={\scriptsize Stellungen (Mrd.)},tick label style={font=\tiny},label style={font=\scriptsize},enlarge x limits=false,ylabel shift=-3pt,xlabel shift=-3pt]\addplot[fill=ink!55,draw=none] coordinates {(0,0.0329) (1,0.1388) (2,0.385) (3,1.072) (4,2.252) (5,0.8051) (6,0.682) (7,0.01989) (8,0.0159) (9,1.073e-05) (10,0.003366) (11,3.42e-07) (12,0.0002872) (13,4.1e-08) (14,1.944e-06)};\draw[densely dotted] (axis cs:12,0) -- (axis cs:12,2.478);\end{axis}\begin{axis}[width=72mm,height=42mm,axis y line*=right,axis x line=none,ymode=log,xmin=-1,xmax=15,ylabel={\scriptsize eindeutig (log)},tick label style={font=\tiny},label style={font=\scriptsize},ylabel shift=-3pt]\addplot[mark=*,mark size=0.7pt,draw=black,thin] coordinates {(0,32897724) (1,105995327) (2,216454493) (3,160616202) (4,146659296) (5,3158299) (6,1719645) (7,435) (8,75179) (9,5) (10,4010) (12,41)};\end{axis}\end{tikzpicture}\end{minipage}
\vspace{3mm}

\noindent\begin{minipage}[t]{60mm}\vspace{0pt}\sbox\hmboard{\chessboard[setfen=7R/7k/8/8/6R1/8/1P5P/K7 b - - 0 1]}\usebox\hmboard\\[1mm]
\makebox[\wd\hmboard]{\sffamily\hspace{3mm}\textbf{h\#6}\hfill(5+1)}\end{minipage}\hfill
\begin{minipage}[t]{52mm}\vspace{0pt}
\textcolor{ink}{\sffamily\textsc{Lösung}}\par\smallskip\small
\begin{tabular}{@{}r@{\enspace}l@{\enspace}l@{}}
\textbf{1.} & \symking{}xh8 & b4\\
\textbf{2.} & \symking{}h7 & b5\\
\textbf{3.} & \symking{}h6 & b6\\
\textbf{4.} & \symking{}h5 & b7\\
\textbf{5.} & \symking{}xg4 & b8=\symqueen{}\\
\textbf{6.} & \symking{}h3 & \symqueen{}g3\#\\
\end{tabular}\par\vspace{3mm}
\qrcode[height=12mm]{https://helpman.komtera.lt/?fen=7R/7k/8/8/6R1/8/1P5P/K7&moves=6}\\{\sffamily\tiny\color{ink}scannen zum Lösen}
\end{minipage}
\vspace{2mm}

{\raggedright\noindent\textcolor{ink}{\sffamily\textsc{Themen}}\enspace \themetag{reines Matt} \themetag{Mustermatt} \themetag{Umwandlung} \themetag{Excelsior} \themetag{Switchback} \themetag{ein Stein zieht} \themetag{ohne Schach}\par\smallskip{\footnotesize\textit{Hinweis:} die Partei am Zug steht im Diagramm im Schach; die Lösung beginnt mit einem Schlag.}\par}

\clearpage
\phantomsection\label{hm:KRRPvkb}
\begin{tcolorbox}[enhanced,colback=ink,colframe=ink,arc=0pt,boxrule=0pt,left=4mm,right=4mm,top=2.5mm,bottom=2.5mm]
\color{white}\sffamily{\LARGE\bfseries KRRPvkb}\hfill{\Large h\#6}\\[0.5mm]
{\small 6 Steine\quad Schwarz am Zug}\hfill{\small Nr.~731\qquad\href{https://helpman.komtera.lt/?fen=7k/7R/6R1/8/8/b7/6P1/K7&moves=6}{\color{white}\bfseries lösen\,\ensuremath{\nearrow}}}
\end{tcolorbox}
\vspace{2mm}
\noindent\begin{minipage}[t]{36mm}\vspace{0pt}\sffamily\small\raggedright
\hmfigure{längste eindeutige}{{\Large h\#6}}\hmfigure{eindeutig bei dieser Länge}{{\Large 49} Stellungen}\hmfigure{längstes Matt}{{\Large h\#8}}\hmfigure{in der Klasse}{45,4 Milliarden Stellungen,\\ 13,5 Milliarden mit Hilfsmatt}
\end{minipage}\hfill
\begin{minipage}[t]{80mm}\vspace{0pt}\begin{tikzpicture}\begin{axis}[width=72mm,height=42mm,ybar,bar width=1.6pt,xmin=-1,xmax=17,ymin=0,ymax=5.874,axis y line*=left,axis x line*=bottom,xlabel={\scriptsize Halbzüge bis Matt},ylabel={\scriptsize Stellungen (Mrd.)},tick label style={font=\tiny},label style={font=\scriptsize},enlarge x limits=false,ylabel shift=-3pt,xlabel shift=-3pt]\addplot[fill=ink!55,draw=none] coordinates {(0,0.05101) (1,0.2183) (2,0.8005) (3,2.7) (4,5.34) (5,2.579) (6,1.693) (7,0.09663) (8,0.02625) (9,2.411e-05) (10,0.008755) (11,4.14e-07) (12,0.001068) (13,8e-08) (14,1.586e-06) (15,2e-09) (16,2e-09)};\draw[densely dotted] (axis cs:12,0) -- (axis cs:12,5.874);\end{axis}\begin{axis}[width=72mm,height=42mm,axis y line*=right,axis x line=none,ymode=log,xmin=-1,xmax=17,ylabel={\scriptsize eindeutig (log)},tick label style={font=\tiny},label style={font=\scriptsize},ylabel shift=-3pt]\addplot[mark=*,mark size=0.7pt,draw=black,thin] coordinates {(0,51010011) (1,184679002) (2,331269355) (3,449468862) (4,286768178) (5,13812961) (6,2270407) (7,2869) (8,226286) (9,22) (10,19373) (11,12) (12,49)};\end{axis}\end{tikzpicture}\end{minipage}
\vspace{3mm}

\noindent\begin{minipage}[t]{60mm}\vspace{0pt}\sbox\hmboard{\chessboard[setfen=7k/7R/6R1/8/8/b7/6P1/K7 b - - 0 1]}\usebox\hmboard\\[1mm]
\makebox[\wd\hmboard]{\sffamily\hspace{3mm}\textbf{h\#6}\hfill(4+2)}\end{minipage}\hfill
\begin{minipage}[t]{52mm}\vspace{0pt}
\textcolor{ink}{\sffamily\textsc{Lösung}}\par\smallskip\small
\begin{tabular}{@{}r@{\enspace}l@{\enspace}l@{}}
\textbf{1.} & \symking{}xh7 & g4\\
\textbf{2.} & \symbishop{}e7 & g5\\
\textbf{3.} & \symbishop{}f6+ & gxf6\\
\textbf{4.} & \symking{}h8 & f7\\
\textbf{5.} & \symking{}h7 & \symrook{}g7+\\
\textbf{6.} & \symking{}h8 & f8=\symqueen{}\#\\
\end{tabular}\par\vspace{3mm}
\qrcode[height=12mm]{https://helpman.komtera.lt/?fen=7k/7R/6R1/8/8/b7/6P1/K7&moves=6}\\{\sffamily\tiny\color{ink}scannen zum Lösen}
\end{minipage}
\vspace{2mm}

{\raggedright\noindent\textcolor{ink}{\sffamily\textsc{Themen}}\enspace \themetag{Umwandlung} \themetag{Excelsior} \themetag{Switchback} \themetag{Pendel}\par\smallskip{\footnotesize\textit{Hinweis:} die Partei am Zug steht im Diagramm im Schach; die Lösung beginnt mit einem Schlag.}\par}

\clearpage
\phantomsection\label{hm:KRRPvkn}
\begin{tcolorbox}[enhanced,colback=ink,colframe=ink,arc=0pt,boxrule=0pt,left=4mm,right=4mm,top=2.5mm,bottom=2.5mm]
\color{white}\sffamily{\LARGE\bfseries KRRPvkn}\hfill{\Large h\#6}\\[0.5mm]
{\small 6 Steine\quad Schwarz am Zug}\hfill{\small Nr.~732\qquad\href{https://helpman.komtera.lt/?fen=7R/6R1/7k/8/1n6/8/1P6/K7&moves=6}{\color{white}\bfseries lösen\,\ensuremath{\nearrow}}}
\end{tcolorbox}
\vspace{2mm}
\noindent\begin{minipage}[t]{36mm}\vspace{0pt}\sffamily\small\raggedright
\hmfigure{längste eindeutige}{{\Large h\#6}}\hmfigure{eindeutig bei dieser Länge}{{\Large 8} Stellungen}\hmfigure{längstes Matt}{{\Large h\#8}}\hmfigure{in der Klasse}{45,4 Milliarden Stellungen,\\ 13,9 Milliarden mit Hilfsmatt}
\end{minipage}\hfill
\begin{minipage}[t]{80mm}\vspace{0pt}\begin{tikzpicture}\begin{axis}[width=72mm,height=42mm,ybar,bar width=1.6pt,xmin=-1,xmax=17,ymin=0,ymax=6.162,axis y line*=left,axis x line*=bottom,xlabel={\scriptsize Halbzüge bis Matt},ylabel={\scriptsize Stellungen (Mrd.)},tick label style={font=\tiny},label style={font=\scriptsize},enlarge x limits=false,ylabel shift=-3pt,xlabel shift=-3pt]\addplot[fill=ink!55,draw=none] coordinates {(0,0.06515) (1,0.2787) (2,0.9236) (3,2.833) (4,5.602) (5,2.399) (6,1.697) (7,0.08397) (8,0.02785) (9,2.762e-05) (10,0.009144) (11,1.55e-07) (12,0.00115) (13,1.1e-08) (14,1.891e-06) (16,5e-09)};\draw[densely dotted] (axis cs:12,0) -- (axis cs:12,6.162);\end{axis}\begin{axis}[width=72mm,height=42mm,axis y line*=right,axis x line=none,ymode=log,xmin=-1,xmax=17,ylabel={\scriptsize eindeutig (log)},tick label style={font=\tiny},label style={font=\scriptsize},ylabel shift=-3pt]\addplot[mark=*,mark size=0.7pt,draw=black,thin] coordinates {(0,65150906) (1,229736896) (2,377611302) (3,391104654) (4,281188194) (5,9804403) (6,2203518) (7,6252) (8,215004) (9,6) (10,21181) (12,8)};\end{axis}\end{tikzpicture}\end{minipage}
\vspace{3mm}

\noindent\begin{minipage}[t]{60mm}\vspace{0pt}\sbox\hmboard{\chessboard[setfen=7R/6R1/7k/8/1n6/8/1P6/K7 b - - 0 1]}\usebox\hmboard\\[1mm]
\makebox[\wd\hmboard]{\sffamily\hspace{3mm}\textbf{h\#6}\hfill(4+2)}\end{minipage}\hfill
\begin{minipage}[t]{52mm}\vspace{0pt}
\textcolor{ink}{\sffamily\textsc{Lösung}}\par\smallskip\small
\begin{tabular}{@{}r@{\enspace}l@{\enspace}l@{}}
\textbf{1.} & \symking{}xg7 & \symking{}b1\\
\textbf{2.} & \symknight{}d5 & \symking{}c2\\
\textbf{3.} & \symknight{}f6 & \symking{}d3\\
\textbf{4.} & \symknight{}h7 & \symking{}e4\\
\textbf{5.} & \symking{}h6 & \symking{}f5\\
\textbf{6.} & \symking{}h5 & \symrook{}xh7\#\\
\end{tabular}\par\vspace{3mm}
\qrcode[height=12mm]{https://helpman.komtera.lt/?fen=7R/6R1/7k/8/1n6/8/1P6/K7&moves=6}\\{\sffamily\tiny\color{ink}scannen zum Lösen}
\end{minipage}
\vspace{2mm}

{\raggedright\noindent\textcolor{ink}{\sffamily\textsc{Themen}}\enspace \themetag{reines Matt} \themetag{Mustermatt} \themetag{Spiegelmatt} \themetag{Switchback} \themetag{ohne Schach}\par\smallskip{\footnotesize\textit{Hinweis:} die Partei am Zug steht im Diagramm im Schach; die Lösung beginnt mit einem Schlag.}\par}

\clearpage
\phantomsection\label{hm:KRRPvkp}
\begin{tcolorbox}[enhanced,colback=ink,colframe=ink,arc=0pt,boxrule=0pt,left=4mm,right=4mm,top=2.5mm,bottom=2.5mm]
\color{white}\sffamily{\LARGE\bfseries KRRPvkp}\hfill{\Large h\#6}\\[0.5mm]
{\small 6 Steine\quad Schwarz am Zug}\hfill{\small Nr.~733\qquad\href{https://helpman.komtera.lt/?fen=8/8/7R/7k/4p3/8/7P/K5R1&moves=6}{\color{white}\bfseries lösen\,\ensuremath{\nearrow}}}
\end{tcolorbox}
\vspace{2mm}
\noindent\begin{minipage}[t]{36mm}\vspace{0pt}\sffamily\small\raggedright
\hmfigure{längste eindeutige}{{\Large h\#6}}\hmfigure{eindeutig bei dieser Länge}{{\Large 224} Stellungen}\hmfigure{längstes Matt}{{\Large h\#10}}\hmfigure{in der Klasse}{34,1 Milliarden Stellungen,\\ 10,8 Milliarden mit Hilfsmatt}
\end{minipage}\hfill
\begin{minipage}[t]{80mm}\vspace{0pt}\begin{tikzpicture}\begin{axis}[width=72mm,height=42mm,ybar,bar width=1.6pt,xmin=-1,xmax=21,ymin=0,ymax=4.839,axis y line*=left,axis x line*=bottom,xlabel={\scriptsize Halbzüge bis Matt},ylabel={\scriptsize Stellungen (Mrd.)},tick label style={font=\tiny},label style={font=\scriptsize},enlarge x limits=false,ylabel shift=-3pt,xlabel shift=-3pt]\addplot[fill=ink!55,draw=none] coordinates {(0,0.06156) (1,0.2554) (2,0.7701) (3,2.23) (4,4.399) (5,1.651) (6,1.313) (7,0.04856) (8,0.02517) (9,2.707e-05) (10,0.007686) (11,7.367e-06) (12,0.001059) (13,5.97e-07) (14,1.174e-05) (16,1.3e-08) (18,1e-09) (20,1e-09)};\draw[densely dotted] (axis cs:12,0) -- (axis cs:12,4.839);\end{axis}\begin{axis}[width=72mm,height=42mm,axis y line*=right,axis x line=none,ymode=log,xmin=-1,xmax=21,ylabel={\scriptsize eindeutig (log)},tick label style={font=\tiny},label style={font=\scriptsize},ylabel shift=-3pt]\addplot[mark=*,mark size=0.7pt,draw=black,thin] coordinates {(0,61563315) (1,207856783) (2,364983733) (3,324031617) (4,244582372) (5,6304458) (6,2020212) (7,1466) (8,160914) (9,4) (10,12134) (11,5) (12,224)};\end{axis}\end{tikzpicture}\end{minipage}
\vspace{3mm}

\noindent\begin{minipage}[t]{60mm}\vspace{0pt}\sbox\hmboard{\chessboard[setfen=8/8/7R/7k/4p3/8/7P/K5R1 b - - 0 1]}\usebox\hmboard\\[1mm]
\makebox[\wd\hmboard]{\sffamily\hspace{3mm}\textbf{h\#6}\hfill(4+2)}\end{minipage}\hfill
\begin{minipage}[t]{52mm}\vspace{0pt}
\textcolor{ink}{\sffamily\textsc{Lösung}}\par\smallskip\small
\begin{tabular}{@{}r@{\enspace}l@{\enspace}l@{}}
\textbf{1.} & \symking{}xh6 & \symking{}b2\\
\textbf{2.} & \symking{}h5 & \symking{}c3\\
\textbf{3.} & \symking{}h4 & \symking{}d4\\
\textbf{4.} & \symking{}h3 & \symking{}xe4\\
\textbf{5.} & \symking{}xh2 & \symking{}f3\\
\textbf{6.} & \symking{}h3 & \symrook{}h1\#\\
\end{tabular}\par\vspace{3mm}
\qrcode[height=12mm]{https://helpman.komtera.lt/?fen=8/8/7R/7k/4p3/8/7P/K5R1&moves=6}\\{\sffamily\tiny\color{ink}scannen zum Lösen}
\end{minipage}
\vspace{2mm}

{\raggedright\noindent\textcolor{ink}{\sffamily\textsc{Themen}}\enspace \themetag{reines Matt} \themetag{Mustermatt} \themetag{Idealmatt} \themetag{Spiegelmatt} \themetag{Switchback} \themetag{ein Stein zieht} \themetag{ohne Schach}\par\smallskip{\footnotesize\textit{Hinweis:} die Partei am Zug steht im Diagramm im Schach; die Lösung beginnt mit einem Schlag.}\par}

\clearpage
\phantomsection\label{hm:KRRPvkq}
\begin{tcolorbox}[enhanced,colback=ink,colframe=ink,arc=0pt,boxrule=0pt,left=4mm,right=4mm,top=2.5mm,bottom=2.5mm]
\color{white}\sffamily{\LARGE\bfseries KRRPvkq}\hfill{\Large h\#6}\\[0.5mm]
{\small 6 Steine\quad Schwarz am Zug}\hfill{\small Nr.~734\qquad\href{https://helpman.komtera.lt/?fen=8/K7/8/8/8/8/3P2R1/1q4Rk&moves=6}{\color{white}\bfseries lösen\,\ensuremath{\nearrow}}}
\end{tcolorbox}
\vspace{2mm}
\noindent\begin{minipage}[t]{36mm}\vspace{0pt}\sffamily\small\raggedright
\hmfigure{längste eindeutige}{{\Large h\#6}}\hmfigure{eindeutig bei dieser Länge}{{\Large 16} Stellungen}\hmfigure{längstes Matt}{{\Large h\#8}}\hmfigure{in der Klasse}{45,4 Milliarden Stellungen,\\ 11,7 Milliarden mit Hilfsmatt}
\end{minipage}\hfill
\begin{minipage}[t]{80mm}\vspace{0pt}\begin{tikzpicture}\begin{axis}[width=72mm,height=42mm,ybar,bar width=1.6pt,xmin=-1,xmax=17,ymin=0,ymax=4.552,axis y line*=left,axis x line*=bottom,xlabel={\scriptsize Halbzüge bis Matt},ylabel={\scriptsize Stellungen (Mrd.)},tick label style={font=\tiny},label style={font=\scriptsize},enlarge x limits=false,ylabel shift=-3pt,xlabel shift=-3pt]\addplot[fill=ink!55,draw=none] coordinates {(0,0.0258) (1,0.1137) (2,0.5383) (3,2.197) (4,4.138) (5,3.087) (6,1.408) (7,0.1907) (8,0.01959) (9,0.0002383) (10,0.007133) (11,6.228e-05) (12,0.0009513) (13,8.37e-06) (14,2.315e-06) (15,3e-09) (16,2.4e-08)};\draw[densely dotted] (axis cs:12,0) -- (axis cs:12,4.552);\end{axis}\begin{axis}[width=72mm,height=42mm,axis y line*=right,axis x line=none,ymode=log,xmin=-1,xmax=17,ylabel={\scriptsize eindeutig (log)},tick label style={font=\tiny},label style={font=\scriptsize},ylabel shift=-3pt]\addplot[mark=*,mark size=0.7pt,draw=black,thin] coordinates {(0,25799844) (1,101092428) (2,225694284) (3,439866469) (4,224766703) (5,25343996) (6,1492060) (7,18082) (8,150513) (9,184) (10,17067) (11,21) (12,16)};\end{axis}\end{tikzpicture}\end{minipage}
\vspace{3mm}

\noindent\begin{minipage}[t]{60mm}\vspace{0pt}\sbox\hmboard{\chessboard[setfen=8/K7/8/8/8/8/3P2R1/1q4Rk b - - 0 1]}\usebox\hmboard\\[1mm]
\makebox[\wd\hmboard]{\sffamily\hspace{3mm}\textbf{h\#6}\hfill(4+2)}\end{minipage}\hfill
\begin{minipage}[t]{52mm}\vspace{0pt}
\textcolor{ink}{\sffamily\textsc{Lösung}}\par\smallskip\small
\begin{tabular}{@{}r@{\enspace}l@{\enspace}l@{}}
\textbf{1.} & \symqueen{}xg1+ & \symking{}a6\\
\textbf{2.} & \symqueen{}h2 & \symrook{}xh2+\\
\textbf{3.} & \symking{}g1 & \symking{}b5\\
\textbf{4.} & \symking{}f1 & \symking{}c4\\
\textbf{5.} & \symking{}e1 & \symking{}d3\\
\textbf{6.} & \symking{}d1 & \symrook{}h1\#\\
\end{tabular}\par\vspace{3mm}
\qrcode[height=12mm]{https://helpman.komtera.lt/?fen=8/K7/8/8/8/8/3P2R1/1q4Rk&moves=6}\\{\sffamily\tiny\color{ink}scannen zum Lösen}
\end{minipage}
\vspace{2mm}

{\raggedright\noindent\textcolor{ink}{\sffamily\textsc{Themen}}\enspace \themetag{reines Matt} \themetag{Mustermatt} \themetag{Idealmatt}\par\smallskip{\footnotesize\textit{Hinweis:} die Partei am Zug steht im Diagramm im Schach; die Lösung beginnt mit einem Schlag.}\par}

\clearpage
\phantomsection\label{hm:KRRPvkr}
\begin{tcolorbox}[enhanced,colback=ink,colframe=ink,arc=0pt,boxrule=0pt,left=4mm,right=4mm,top=2.5mm,bottom=2.5mm]
\color{white}\sffamily{\LARGE\bfseries KRRPvkr}\hfill{\Large h\#6}\\[0.5mm]
{\small 6 Steine\quad Schwarz am Zug}\hfill{\small Nr.~735\qquad\href{https://helpman.komtera.lt/?fen=8/K7/8/8/8/1r6/1P5R/1R4k1&moves=6}{\color{white}\bfseries lösen\,\ensuremath{\nearrow}}}
\end{tcolorbox}
\vspace{2mm}
\noindent\begin{minipage}[t]{36mm}\vspace{0pt}\sffamily\small\raggedright
\hmfigure{längste eindeutige}{{\Large h\#6}}\hmfigure{eindeutig bei dieser Länge}{{\Large 11} Stellungen}\hmfigure{längstes Matt}{{\Large h\#7}}\hmfigure{in der Klasse}{45,4 Milliarden Stellungen,\\ 12,9 Milliarden mit Hilfsmatt}
\end{minipage}\hfill
\begin{minipage}[t]{80mm}\vspace{0pt}\begin{tikzpicture}\begin{axis}[width=72mm,height=42mm,ybar,bar width=1.6pt,xmin=-1,xmax=15,ymin=0,ymax=5.551,axis y line*=left,axis x line*=bottom,xlabel={\scriptsize Halbzüge bis Matt},ylabel={\scriptsize Stellungen (Mrd.)},tick label style={font=\tiny},label style={font=\scriptsize},enlarge x limits=false,ylabel shift=-3pt,xlabel shift=-3pt]\addplot[fill=ink!55,draw=none] coordinates {(0,0.04417) (1,0.1956) (2,0.7732) (3,2.634) (4,5.047) (5,2.651) (6,1.392) (7,0.1113) (8,0.02163) (9,4.475e-05) (10,0.007929) (11,2.76e-06) (12,0.0008062) (13,9.5e-08) (14,1.371e-06)};\draw[densely dotted] (axis cs:12,0) -- (axis cs:12,5.551);\end{axis}\begin{axis}[width=72mm,height=42mm,axis y line*=right,axis x line=none,ymode=log,xmin=-1,xmax=15,ylabel={\scriptsize eindeutig (log)},tick label style={font=\tiny},label style={font=\scriptsize},ylabel shift=-3pt]\addplot[mark=*,mark size=0.7pt,draw=black,thin] coordinates {(0,44170246) (1,170131366) (2,332434485) (3,449737646) (4,271746912) (5,14657584) (6,1710305) (7,4415) (8,203368) (9,10) (10,19609) (12,11)};\end{axis}\end{tikzpicture}\end{minipage}
\vspace{3mm}

\noindent\begin{minipage}[t]{60mm}\vspace{0pt}\sbox\hmboard{\chessboard[setfen=8/K7/8/8/8/1r6/1P5R/1R4k1 b - - 0 1]}\usebox\hmboard\\[1mm]
\makebox[\wd\hmboard]{\sffamily\hspace{3mm}\textbf{h\#6}\hfill(4+2)}\end{minipage}\hfill
\begin{minipage}[t]{52mm}\vspace{0pt}
\textcolor{ink}{\sffamily\textsc{Lösung}}\par\smallskip\small
\begin{tabular}{@{}r@{\enspace}l@{\enspace}l@{}}
\textbf{1.} & \symking{}xh2 & \symking{}a6\\
\textbf{2.} & \symrook{}xb2 & \symrook{}xb2+\\
\textbf{3.} & \symking{}g1 & \symking{}b5\\
\textbf{4.} & \symking{}f1 & \symking{}c4\\
\textbf{5.} & \symking{}e1 & \symking{}d3\\
\textbf{6.} & \symking{}d1 & \symrook{}b1\#\\
\end{tabular}\par\vspace{3mm}
\qrcode[height=12mm]{https://helpman.komtera.lt/?fen=8/K7/8/8/8/1r6/1P5R/1R4k1&moves=6}\\{\sffamily\tiny\color{ink}scannen zum Lösen}
\end{minipage}
\vspace{2mm}

{\raggedright\noindent\textcolor{ink}{\sffamily\textsc{Themen}}\enspace \themetag{reines Matt} \themetag{Mustermatt} \themetag{Idealmatt} \themetag{Spiegelmatt} \themetag{Switchback}\par\smallskip{\footnotesize\textit{Hinweis:} die Partei am Zug steht im Diagramm im Schach; die Lösung beginnt mit einem Schlag.}\par}

\clearpage
\phantomsection\label{hm:KRRvkbb}
\begin{tcolorbox}[enhanced,colback=ink,colframe=ink,arc=0pt,boxrule=0pt,left=4mm,right=4mm,top=2.5mm,bottom=2.5mm]
\color{white}\sffamily{\LARGE\bfseries KRRvkbb}\hfill{\Large h\#6}\\[0.5mm]
{\small 6 Steine\quad Schwarz am Zug}\hfill{\small Nr.~736\qquad\href{https://helpman.komtera.lt/?fen=8/8/8/7R/7k/b7/2b3R1/K7&moves=6}{\color{white}\bfseries lösen\,\ensuremath{\nearrow}}}
\end{tcolorbox}
\vspace{2mm}
\noindent\begin{minipage}[t]{36mm}\vspace{0pt}\sffamily\small\raggedright
\hmfigure{längste eindeutige}{{\Large h\#6}}\hmfigure{eindeutig bei dieser Länge}{{\Large 72} Stellungen}\hmfigure{längstes Matt}{{\Large h\#7}}\hmfigure{in der Klasse}{15,5 Milliarden Stellungen,\\ 2,1 Milliarden mit Hilfsmatt}
\end{minipage}\hfill
\begin{minipage}[t]{80mm}\vspace{0pt}\begin{tikzpicture}\begin{axis}[width=72mm,height=42mm,ybar,bar width=1.6pt,xmin=-1,xmax=15,ymin=0,ymax=816.8,axis y line*=left,axis x line*=bottom,xlabel={\scriptsize Halbzüge bis Matt},ylabel={\scriptsize Stellungen (Mio.)},tick label style={font=\tiny},label style={font=\scriptsize},enlarge x limits=false,ylabel shift=-3pt,xlabel shift=-3pt]\addplot[fill=ink!55,draw=none] coordinates {(0,4.93) (1,21.09) (2,90.24) (3,377.9) (4,742.5) (5,523.4) (6,295) (7,32.78) (8,5.473) (9,0.01918) (10,2.277) (11,2e-06) (12,0.6082) (14,0.006158)};\draw[densely dotted] (axis cs:12,0) -- (axis cs:12,816.8);\end{axis}\begin{axis}[width=72mm,height=42mm,axis y line*=right,axis x line=none,ymode=log,xmin=-1,xmax=15,ylabel={\scriptsize eindeutig (log)},tick label style={font=\tiny},label style={font=\scriptsize},ylabel shift=-3pt]\addplot[mark=*,mark size=0.7pt,draw=black,thin] coordinates {(0,4930283) (1,19451468) (2,38161931) (3,75054177) (4,37876714) (5,3761697) (6,491886) (7,2208) (8,43260) (9,50) (10,10906) (12,72)};\end{axis}\end{tikzpicture}\end{minipage}
\vspace{3mm}

\noindent\begin{minipage}[t]{60mm}\vspace{0pt}\sbox\hmboard{\chessboard[setfen=8/8/8/7R/7k/b7/2b3R1/K7 b - - 0 1]}\usebox\hmboard\\[1mm]
\makebox[\wd\hmboard]{\sffamily\hspace{3mm}\textbf{h\#6}\hfill(3+3)}\end{minipage}\hfill
\begin{minipage}[t]{52mm}\vspace{0pt}
\textcolor{ink}{\sffamily\textsc{Lösung}}\par\smallskip\small
\begin{tabular}{@{}r@{\enspace}l@{\enspace}l@{}}
\textbf{1.} & \symking{}xh5 & \symrook{}xc2\\
\textbf{2.} & \symking{}g4 & \symrook{}a2\\
\textbf{3.} & \symking{}f3 & \symrook{}xa3+\\
\textbf{4.} & \symking{}e2 & \symking{}b2\\
\textbf{5.} & \symking{}d1 & \symking{}c3\\
\textbf{6.} & \symking{}c1 & \symrook{}a1\#\\
\end{tabular}\par\vspace{3mm}
\qrcode[height=12mm]{https://helpman.komtera.lt/?fen=8/8/8/7R/7k/b7/2b3R1/K7&moves=6}\\{\sffamily\tiny\color{ink}scannen zum Lösen}
\end{minipage}
\vspace{2mm}

{\raggedright\noindent\textcolor{ink}{\sffamily\textsc{Themen}}\enspace \themetag{reines Matt} \themetag{Mustermatt} \themetag{Idealmatt} \themetag{Spiegelmatt} \themetag{ein Stein zieht}\par\smallskip{\footnotesize\textit{Hinweis:} die Partei am Zug steht im Diagramm im Schach; die Lösung beginnt mit einem Schlag.}\par}

\clearpage
\phantomsection\label{hm:KRRvkbn}
\begin{tcolorbox}[enhanced,colback=ink,colframe=ink,arc=0pt,boxrule=0pt,left=4mm,right=4mm,top=2.5mm,bottom=2.5mm]
\color{white}\sffamily{\LARGE\bfseries KRRvkbn}\hfill{\Large h\#6}\\[0.5mm]
{\small 6 Steine\quad Schwarz am Zug}\hfill{\small Nr.~737\qquad\href{https://helpman.komtera.lt/?fen=8/8/8/7R/7k/b1n5/8/K5R1&moves=6}{\color{white}\bfseries lösen\,\ensuremath{\nearrow}}}
\end{tcolorbox}
\vspace{2mm}
\noindent\begin{minipage}[t]{36mm}\vspace{0pt}\sffamily\small\raggedright
\hmfigure{längste eindeutige}{{\Large h\#6}}\hmfigure{eindeutig bei dieser Länge}{{\Large 74} Stellungen}\hmfigure{längstes Matt}{{\Large h\#7}}\hmfigure{in der Klasse}{15,5 Milliarden Stellungen,\\ 4,33 Milliarden mit Hilfsmatt}
\end{minipage}\hfill
\begin{minipage}[t]{80mm}\vspace{0pt}\begin{tikzpicture}\begin{axis}[width=72mm,height=42mm,ybar,bar width=1.6pt,xmin=-1,xmax=15,ymin=0,ymax=1.767,axis y line*=left,axis x line*=bottom,xlabel={\scriptsize Halbzüge bis Matt},ylabel={\scriptsize Stellungen (Mrd.)},tick label style={font=\tiny},label style={font=\scriptsize},enlarge x limits=false,ylabel shift=-3pt,xlabel shift=-3pt]\addplot[fill=ink!55,draw=none] coordinates {(0,0.0127) (1,0.05493) (2,0.2174) (3,0.8236) (4,1.607) (5,0.9816) (6,0.5706) (7,0.05029) (8,0.009199) (9,1.415e-05) (10,0.004751) (11,3.6e-08) (12,0.001225) (14,1.368e-05)};\draw[densely dotted] (axis cs:12,0) -- (axis cs:12,1.767);\end{axis}\begin{axis}[width=72mm,height=42mm,axis y line*=right,axis x line=none,ymode=log,xmin=-1,xmax=15,ylabel={\scriptsize eindeutig (log)},tick label style={font=\tiny},label style={font=\scriptsize},ylabel shift=-3pt]\addplot[mark=*,mark size=0.7pt,draw=black,thin] coordinates {(0,12699369) (1,49770766) (2,90138823) (3,148542205) (4,81826611) (5,5887331) (6,778688) (7,1344) (8,99554) (9,61) (10,20183) (12,74)};\end{axis}\end{tikzpicture}\end{minipage}
\vspace{3mm}

\noindent\begin{minipage}[t]{60mm}\vspace{0pt}\sbox\hmboard{\chessboard[setfen=8/8/8/7R/7k/b1n5/8/K5R1 b - - 0 1]}\usebox\hmboard\\[1mm]
\makebox[\wd\hmboard]{\sffamily\hspace{3mm}\textbf{h\#6}\hfill(3+3)}\end{minipage}\hfill
\begin{minipage}[t]{52mm}\vspace{0pt}
\textcolor{ink}{\sffamily\textsc{Lösung}}\par\smallskip\small
\begin{tabular}{@{}r@{\enspace}l@{\enspace}l@{}}
\textbf{1.} & \symking{}xh5 & \symrook{}b1\\
\textbf{2.} & \symking{}g4 & \symrook{}b3\\
\textbf{3.} & \symking{}f3 & \symrook{}xa3\\
\textbf{4.} & \symking{}e2 & \symking{}b2\\
\textbf{5.} & \symking{}d1 & \symking{}xc3\\
\textbf{6.} & \symking{}c1 & \symrook{}a1\#\\
\end{tabular}\par\vspace{3mm}
\qrcode[height=12mm]{https://helpman.komtera.lt/?fen=8/8/8/7R/7k/b1n5/8/K5R1&moves=6}\\{\sffamily\tiny\color{ink}scannen zum Lösen}
\end{minipage}
\vspace{2mm}

{\raggedright\noindent\textcolor{ink}{\sffamily\textsc{Themen}}\enspace \themetag{reines Matt} \themetag{Mustermatt} \themetag{Idealmatt} \themetag{Spiegelmatt} \themetag{ein Stein zieht} \themetag{ohne Schach}\par\smallskip{\footnotesize\textit{Hinweis:} die Partei am Zug steht im Diagramm im Schach; die Lösung beginnt mit einem Schlag.}\par}

\clearpage
\phantomsection\label{hm:KRRvkbp}
\begin{tcolorbox}[enhanced,colback=ink,colframe=ink,arc=0pt,boxrule=0pt,left=4mm,right=4mm,top=2.5mm,bottom=2.5mm]
\color{white}\sffamily{\LARGE\bfseries KRRvkbp}\hfill{\Large h\#6}\\[0.5mm]
{\small 6 Steine\quad Schwarz am Zug}\hfill{\small Nr.~738\qquad\href{https://helpman.komtera.lt/?fen=K7/8/8/4R3/7p/8/6Rk/7b&moves=6}{\color{white}\bfseries lösen\,\ensuremath{\nearrow}}}
\end{tcolorbox}
\vspace{2mm}
\noindent\begin{minipage}[t]{36mm}\vspace{0pt}\sffamily\small\raggedright
\hmfigure{längste eindeutige}{{\Large h\#6}}\hmfigure{eindeutig bei dieser Länge}{{\Large 723} Stellungen}\hmfigure{längstes Matt}{{\Large h\#7}}\hmfigure{in der Klasse}{45,4 Milliarden Stellungen,\\ 13,4 Milliarden mit Hilfsmatt}
\end{minipage}\hfill
\begin{minipage}[t]{80mm}\vspace{0pt}\begin{tikzpicture}\begin{axis}[width=72mm,height=42mm,ybar,bar width=1.6pt,xmin=-1,xmax=15,ymin=0,ymax=5.61,axis y line*=left,axis x line*=bottom,xlabel={\scriptsize Halbzüge bis Matt},ylabel={\scriptsize Stellungen (Mrd.)},tick label style={font=\tiny},label style={font=\scriptsize},enlarge x limits=false,ylabel shift=-3pt,xlabel shift=-3pt]\addplot[fill=ink!55,draw=none] coordinates {(0,0.04764) (1,0.1999) (2,0.7397) (3,2.653) (4,5.1) (5,2.769) (6,1.753) (7,0.1269) (8,0.03247) (9,4.616e-05) (10,0.01513) (11,2.03e-07) (12,0.003595) (13,2.8e-08) (14,3.816e-05)};\draw[densely dotted] (axis cs:12,0) -- (axis cs:12,5.61);\end{axis}\begin{axis}[width=72mm,height=42mm,axis y line*=right,axis x line=none,ymode=log,xmin=-1,xmax=15,ylabel={\scriptsize eindeutig (log)},tick label style={font=\tiny},label style={font=\scriptsize},ylabel shift=-3pt]\addplot[mark=*,mark size=0.7pt,draw=black,thin] coordinates {(0,47642419) (1,179736637) (2,316691986) (3,484473569) (4,274368462) (5,16430594) (6,2814920) (7,3595) (8,317855) (9,2) (10,55210) (11,1) (12,723)};\end{axis}\end{tikzpicture}\end{minipage}
\vspace{3mm}

\noindent\begin{minipage}[t]{60mm}\vspace{0pt}\sbox\hmboard{\chessboard[setfen=K7/8/8/4R3/7p/8/6Rk/7b b - - 0 1]}\usebox\hmboard\\[1mm]
\makebox[\wd\hmboard]{\sffamily\hspace{3mm}\textbf{h\#6}\hfill(3+3)}\end{minipage}\hfill
\begin{minipage}[t]{52mm}\vspace{0pt}
\textcolor{ink}{\sffamily\textsc{Lösung}}\par\smallskip\small
\begin{tabular}{@{}r@{\enspace}l@{\enspace}l@{}}
\textbf{1.} & \symking{}h3 & \symking{}b7\\
\textbf{2.} & \symking{}xg2 & \symking{}c6\\
\textbf{3.} & h3 & \symking{}d5\\
\textbf{4.} & h2 & \symking{}e4\\
\textbf{5.} & \symking{}h3+ & \symking{}f4\\
\textbf{6.} & \symbishop{}g2 & \symrook{}h5\#\\
\end{tabular}\par\vspace{3mm}
\qrcode[height=12mm]{https://helpman.komtera.lt/?fen=K7/8/8/4R3/7p/8/6Rk/7b&moves=6}\\{\sffamily\tiny\color{ink}scannen zum Lösen}
\end{minipage}
\vspace{2mm}

{\raggedright\noindent\textcolor{ink}{\sffamily\textsc{Themen}}\enspace \themetag{Switchback} \themetag{Selbstblock}\par\smallskip{\footnotesize\textit{Hinweis:} die Partei am Zug steht im Diagramm im Schach.}\par}

\clearpage
\phantomsection\label{hm:KRRvknn}
\begin{tcolorbox}[enhanced,colback=ink,colframe=ink,arc=0pt,boxrule=0pt,left=4mm,right=4mm,top=2.5mm,bottom=2.5mm]
\color{white}\sffamily{\LARGE\bfseries KRRvknn}\hfill{\Large h\#6}\\[0.5mm]
{\small 6 Steine\quad Schwarz am Zug}\hfill{\small Nr.~739\qquad\href{https://helpman.komtera.lt/?fen=8/8/8/6R1/7k/2nn4/8/K6R&moves=6}{\color{white}\bfseries lösen\,\ensuremath{\nearrow}}}
\end{tcolorbox}
\vspace{2mm}
\noindent\begin{minipage}[t]{36mm}\vspace{0pt}\sffamily\small\raggedright
\hmfigure{längste eindeutige}{{\Large h\#6}}\hmfigure{eindeutig bei dieser Länge}{{\Large 175} Stellungen}\hmfigure{längstes Matt}{{\Large h\#7}}\hmfigure{in der Klasse}{15,5 Milliarden Stellungen,\\ 2,23 Milliarden mit Hilfsmatt}
\end{minipage}\hfill
\begin{minipage}[t]{80mm}\vspace{0pt}\begin{tikzpicture}\begin{axis}[width=72mm,height=42mm,ybar,bar width=1.6pt,xmin=-1,xmax=15,ymin=0,ymax=928.3,axis y line*=left,axis x line*=bottom,xlabel={\scriptsize Halbzüge bis Matt},ylabel={\scriptsize Stellungen (Mio.)},tick label style={font=\tiny},label style={font=\scriptsize},enlarge x limits=false,ylabel shift=-3pt,xlabel shift=-3pt]\addplot[fill=ink!55,draw=none] coordinates {(0,8.111) (1,35.14) (2,126.8) (3,436.6) (4,843.9) (5,459.2) (6,284.5) (7,24.29) (8,5.674) (9,0.0131) (10,2.449) (11,2e-06) (12,0.6394) (14,0.008156)};\draw[densely dotted] (axis cs:12,0) -- (axis cs:12,928.3);\end{axis}\begin{axis}[width=72mm,height=42mm,axis y line*=right,axis x line=none,ymode=log,xmin=-1,xmax=15,ylabel={\scriptsize eindeutig (log)},tick label style={font=\tiny},label style={font=\scriptsize},ylabel shift=-3pt]\addplot[mark=*,mark size=0.7pt,draw=black,thin] coordinates {(0,8110554) (1,31139741) (2,51436282) (3,67188429) (4,39394832) (5,2263196) (6,472012) (7,1090) (8,53362) (10,9460) (12,175)};\end{axis}\end{tikzpicture}\end{minipage}
\vspace{3mm}

\noindent\begin{minipage}[t]{60mm}\vspace{0pt}\sbox\hmboard{\chessboard[setfen=8/8/8/6R1/7k/2nn4/8/K6R b - - 0 1]}\usebox\hmboard\\[1mm]
\makebox[\wd\hmboard]{\sffamily\hspace{3mm}\textbf{h\#6}\hfill(3+3)}\end{minipage}\hfill
\begin{minipage}[t]{52mm}\vspace{0pt}
\textcolor{ink}{\sffamily\textsc{Lösung}}\par\smallskip\small
\begin{tabular}{@{}r@{\enspace}l@{\enspace}l@{}}
\textbf{1.} & \symking{}xg5 & \symrook{}c1\\
\textbf{2.} & \symking{}f4 & \symrook{}xc3\\
\textbf{3.} & \symking{}e3 & \symking{}a2\\
\textbf{4.} & \symking{}d2 & \symrook{}xd3+\\
\textbf{5.} & \symking{}c1 & \symking{}b3\\
\textbf{6.} & \symking{}b1 & \symrook{}d1\#\\
\end{tabular}\par\vspace{3mm}
\qrcode[height=12mm]{https://helpman.komtera.lt/?fen=8/8/8/6R1/7k/2nn4/8/K6R&moves=6}\\{\sffamily\tiny\color{ink}scannen zum Lösen}
\end{minipage}
\vspace{2mm}

{\raggedright\noindent\textcolor{ink}{\sffamily\textsc{Themen}}\enspace \themetag{reines Matt} \themetag{Mustermatt} \themetag{Idealmatt} \themetag{Spiegelmatt} \themetag{ein Stein zieht}\par\smallskip{\footnotesize\textit{Hinweis:} die Partei am Zug steht im Diagramm im Schach; die Lösung beginnt mit einem Schlag.}\par}

\clearpage
\phantomsection\label{hm:KRRvknp}
\begin{tcolorbox}[enhanced,colback=ink,colframe=ink,arc=0pt,boxrule=0pt,left=4mm,right=4mm,top=2.5mm,bottom=2.5mm]
\color{white}\sffamily{\LARGE\bfseries KRRvknp}\hfill{\Large h\#6}\\[0.5mm]
{\small 6 Steine\quad Schwarz am Zug}\hfill{\small Nr.~740\qquad\href{https://helpman.komtera.lt/?fen=8/8/8/7R/7k/p1n5/8/K5R1&moves=6}{\color{white}\bfseries lösen\,\ensuremath{\nearrow}}}
\end{tcolorbox}
\vspace{2mm}
\noindent\begin{minipage}[t]{36mm}\vspace{0pt}\sffamily\small\raggedright
\hmfigure{längste eindeutige}{{\Large h\#6}}\hmfigure{eindeutig bei dieser Länge}{{\Large 1\,311} Stellungen}\hmfigure{längstes Matt}{{\Large h\#7}}\hmfigure{in der Klasse}{45,4 Milliarden Stellungen,\\ 13,8 Milliarden mit Hilfsmatt}
\end{minipage}\hfill
\begin{minipage}[t]{80mm}\vspace{0pt}\begin{tikzpicture}\begin{axis}[width=72mm,height=42mm,ybar,bar width=1.6pt,xmin=-1,xmax=15,ymin=0,ymax=5.937,axis y line*=left,axis x line*=bottom,xlabel={\scriptsize Halbzüge bis Matt},ylabel={\scriptsize Stellungen (Mrd.)},tick label style={font=\tiny},label style={font=\scriptsize},enlarge x limits=false,ylabel shift=-3pt,xlabel shift=-3pt]\addplot[fill=ink!55,draw=none] coordinates {(0,0.06127) (1,0.2567) (2,0.8608) (3,2.812) (4,5.398) (5,2.573) (6,1.726) (7,0.1077) (8,0.03284) (9,4.442e-05) (10,0.01536) (11,3.6e-08) (12,0.003636) (14,4.933e-05)};\draw[densely dotted] (axis cs:12,0) -- (axis cs:12,5.937);\end{axis}\begin{axis}[width=72mm,height=42mm,axis y line*=right,axis x line=none,ymode=log,xmin=-1,xmax=15,ylabel={\scriptsize eindeutig (log)},tick label style={font=\tiny},label style={font=\scriptsize},ylabel shift=-3pt]\addplot[mark=*,mark size=0.7pt,draw=black,thin] coordinates {(0,61274213) (1,225698514) (2,362223164) (3,435431833) (4,279668525) (5,11528341) (6,2599405) (7,3187) (8,308963) (10,46951) (12,1311)};\end{axis}\end{tikzpicture}\end{minipage}
\vspace{3mm}

\noindent\begin{minipage}[t]{60mm}\vspace{0pt}\sbox\hmboard{\chessboard[setfen=8/8/8/7R/7k/p1n5/8/K5R1 b - - 0 1]}\usebox\hmboard\\[1mm]
\makebox[\wd\hmboard]{\sffamily\hspace{3mm}\textbf{h\#6}\hfill(3+3)}\end{minipage}\hfill
\begin{minipage}[t]{52mm}\vspace{0pt}
\textcolor{ink}{\sffamily\textsc{Lösung}}\par\smallskip\small
\begin{tabular}{@{}r@{\enspace}l@{\enspace}l@{}}
\textbf{1.} & \symking{}xh5 & \symrook{}b1\\
\textbf{2.} & \symking{}g4 & \symrook{}b3\\
\textbf{3.} & \symking{}f3 & \symrook{}xa3\\
\textbf{4.} & \symking{}e2 & \symking{}b2\\
\textbf{5.} & \symking{}d1 & \symking{}xc3\\
\textbf{6.} & \symking{}c1 & \symrook{}a1\#\\
\end{tabular}\par\vspace{3mm}
\qrcode[height=12mm]{https://helpman.komtera.lt/?fen=8/8/8/7R/7k/p1n5/8/K5R1&moves=6}\\{\sffamily\tiny\color{ink}scannen zum Lösen}
\end{minipage}
\vspace{2mm}

{\raggedright\noindent\textcolor{ink}{\sffamily\textsc{Themen}}\enspace \themetag{reines Matt} \themetag{Mustermatt} \themetag{Idealmatt} \themetag{Spiegelmatt} \themetag{ein Stein zieht} \themetag{ohne Schach}\par\smallskip{\footnotesize\textit{Hinweis:} die Partei am Zug steht im Diagramm im Schach; die Lösung beginnt mit einem Schlag.}\par}

\clearpage
\phantomsection\label{hm:KRRvkpp}
\begin{tcolorbox}[enhanced,colback=ink,colframe=ink,arc=0pt,boxrule=0pt,left=4mm,right=4mm,top=2.5mm,bottom=2.5mm]
\color{white}\sffamily{\LARGE\bfseries KRRvkpp}\hfill{\Large h\#6}\\[0.5mm]
{\small 6 Steine\quad Schwarz am Zug}\hfill{\small Nr.~741\qquad\href{https://helpman.komtera.lt/?fen=K7/8/8/7p/5R1p/7k/6R1/8&moves=6}{\color{white}\bfseries lösen\,\ensuremath{\nearrow}}}
\end{tcolorbox}
\vspace{2mm}
\noindent\begin{minipage}[t]{36mm}\vspace{0pt}\sffamily\small\raggedright
\hmfigure{längste eindeutige}{{\Large h\#6}}\hmfigure{eindeutig bei dieser Länge}{{\Large 912} Stellungen}\hmfigure{längstes Matt}{{\Large h\#7}}\hmfigure{in der Klasse}{34,1 Milliarden Stellungen,\\ 5,35 Milliarden mit Hilfsmatt}
\end{minipage}\hfill
\begin{minipage}[t]{80mm}\vspace{0pt}\begin{tikzpicture}\begin{axis}[width=72mm,height=42mm,ybar,bar width=1.6pt,xmin=-1,xmax=15,ymin=0,ymax=2.339,axis y line*=left,axis x line*=bottom,xlabel={\scriptsize Halbzüge bis Matt},ylabel={\scriptsize Stellungen (Mrd.)},tick label style={font=\tiny},label style={font=\scriptsize},enlarge x limits=false,ylabel shift=-3pt,xlabel shift=-3pt]\addplot[fill=ink!55,draw=none] coordinates {(0,0.02894) (1,0.1171) (2,0.3601) (3,1.114) (4,2.126) (5,0.8878) (6,0.6612) (7,0.03147) (8,0.0137) (9,1.006e-05) (10,0.006071) (11,7.49e-07) (12,0.00139) (13,1.4e-08) (14,2.635e-05)};\draw[densely dotted] (axis cs:12,0) -- (axis cs:12,2.339);\end{axis}\begin{axis}[width=72mm,height=42mm,axis y line*=right,axis x line=none,ymode=log,xmin=-1,xmax=15,ylabel={\scriptsize eindeutig (log)},tick label style={font=\tiny},label style={font=\scriptsize},ylabel shift=-3pt]\addplot[mark=*,mark size=0.7pt,draw=black,thin] coordinates {(0,28943248) (1,102146661) (2,162010650) (3,173700238) (4,117868827) (5,3643916) (6,972489) (7,314) (8,109001) (9,19) (10,13321) (12,912)};\end{axis}\end{tikzpicture}\end{minipage}
\vspace{3mm}

\noindent\begin{minipage}[t]{60mm}\vspace{0pt}\sbox\hmboard{\chessboard[setfen=K7/8/8/7p/5R1p/7k/6R1/8 b - - 0 1]}\usebox\hmboard\\[1mm]
\makebox[\wd\hmboard]{\sffamily\hspace{3mm}\textbf{h\#6}\hfill(3+3)}\end{minipage}\hfill
\begin{minipage}[t]{52mm}\vspace{0pt}
\textcolor{ink}{\sffamily\textsc{Lösung}}\par\smallskip\small
\begin{tabular}{@{}r@{\enspace}l@{\enspace}l@{}}
\textbf{1.} & \symking{}xg2 & \symking{}b7\\
\textbf{2.} & h3 & \symking{}c6\\
\textbf{3.} & h2 & \symking{}d5\\
\textbf{4.} & h1=\symrook{} & \symking{}e4\\
\textbf{5.} & \symking{}h2 & \symking{}f3\\
\textbf{6.} & \symrook{}g1 & \symrook{}h4\#\\
\end{tabular}\par\vspace{3mm}
\qrcode[height=12mm]{https://helpman.komtera.lt/?fen=K7/8/8/7p/5R1p/7k/6R1/8&moves=6}\\{\sffamily\tiny\color{ink}scannen zum Lösen}
\end{minipage}
\vspace{2mm}

{\raggedright\noindent\textcolor{ink}{\sffamily\textsc{Themen}}\enspace \themetag{reines Matt} \themetag{Mustermatt} \themetag{Umwandlung} \themetag{Unterverwandlung} \themetag{Selbstblock} \themetag{ohne Schach}\par\smallskip{\footnotesize\textit{Hinweis:} die Lösung beginnt mit einem Schlag.}\par}

\clearpage
\phantomsection\label{hm:KRRvkqb}
\begin{tcolorbox}[enhanced,colback=ink,colframe=ink,arc=0pt,boxrule=0pt,left=4mm,right=4mm,top=2.5mm,bottom=2.5mm]
\color{white}\sffamily{\LARGE\bfseries KRRvkqb}\hfill{\Large h\#6}\\[0.5mm]
{\small 6 Steine\quad Schwarz am Zug}\hfill{\small Nr.~742\qquad\href{https://helpman.komtera.lt/?fen=8/8/8/7R/1q5k/8/8/K4bR1&moves=6}{\color{white}\bfseries lösen\,\ensuremath{\nearrow}}}
\end{tcolorbox}
\vspace{2mm}
\noindent\begin{minipage}[t]{36mm}\vspace{0pt}\sffamily\small\raggedright
\hmfigure{längste eindeutige}{{\Large h\#6}}\hmfigure{eindeutig bei dieser Länge}{{\Large 140} Stellungen}\hmfigure{längstes Matt}{{\Large h\#8}}\hmfigure{in der Klasse}{15,5 Milliarden Stellungen,\\ 3,67 Milliarden mit Hilfsmatt}
\end{minipage}\hfill
\begin{minipage}[t]{80mm}\vspace{0pt}\begin{tikzpicture}\begin{axis}[width=72mm,height=42mm,ybar,bar width=1.6pt,xmin=-1,xmax=17,ymin=0,ymax=1.3,axis y line*=left,axis x line*=bottom,xlabel={\scriptsize Halbzüge bis Matt},ylabel={\scriptsize Stellungen (Mrd.)},tick label style={font=\tiny},label style={font=\scriptsize},enlarge x limits=false,ylabel shift=-3pt,xlabel shift=-3pt]\addplot[fill=ink!55,draw=none] coordinates {(0,0.005117) (1,0.02203) (2,0.116) (3,0.5872) (4,1.119) (5,1.181) (6,0.5127) (7,0.1146) (8,0.009282) (9,0.0001439) (10,0.003618) (11,1.511e-06) (12,0.001162) (13,1.02e-07) (14,2.72e-05) (16,5e-09)};\draw[densely dotted] (axis cs:12,0) -- (axis cs:12,1.3);\end{axis}\begin{axis}[width=72mm,height=42mm,axis y line*=right,axis x line=none,ymode=log,xmin=-1,xmax=17,ylabel={\scriptsize eindeutig (log)},tick label style={font=\tiny},label style={font=\scriptsize},ylabel shift=-3pt]\addplot[mark=*,mark size=0.7pt,draw=black,thin] coordinates {(0,5116594) (1,21011676) (2,49926258) (3,127977913) (4,58312003) (5,11492325) (6,858961) (7,16689) (8,60082) (9,147) (10,17468) (11,4) (12,140)};\end{axis}\end{tikzpicture}\end{minipage}
\vspace{3mm}

\noindent\begin{minipage}[t]{60mm}\vspace{0pt}\sbox\hmboard{\chessboard[setfen=8/8/8/7R/1q5k/8/8/K4bR1 b - - 0 1]}\usebox\hmboard\\[1mm]
\makebox[\wd\hmboard]{\sffamily\hspace{3mm}\textbf{h\#6}\hfill(3+3)}\end{minipage}\hfill
\begin{minipage}[t]{52mm}\vspace{0pt}
\textcolor{ink}{\sffamily\textsc{Lösung}}\par\smallskip\small
\begin{tabular}{@{}r@{\enspace}l@{\enspace}l@{}}
\textbf{1.} & \symking{}xh5 & \symrook{}h1+\\
\textbf{2.} & \symking{}g4 & \symrook{}h4+\\
\textbf{3.} & \symking{}f3 & \symrook{}xb4\\
\textbf{4.} & \symking{}e2 & \symking{}b2\\
\textbf{5.} & \symking{}d1 & \symking{}c3\\
\textbf{6.} & \symbishop{}e2 & \symrook{}b1\#\\
\end{tabular}\par\vspace{3mm}
\qrcode[height=12mm]{https://helpman.komtera.lt/?fen=8/8/8/7R/1q5k/8/8/K4bR1&moves=6}\\{\sffamily\tiny\color{ink}scannen zum Lösen}
\end{minipage}
\vspace{2mm}

{\raggedright\noindent\textcolor{ink}{\sffamily\textsc{Themen}}\enspace \themetag{reines Matt} \themetag{Mustermatt} \themetag{Idealmatt} \themetag{Selbstblock}\par\smallskip{\footnotesize\textit{Hinweis:} die Partei am Zug steht im Diagramm im Schach; die Lösung beginnt mit einem Schlag.}\par}

\clearpage
\phantomsection\label{hm:KRRvkqn}
\begin{tcolorbox}[enhanced,colback=ink,colframe=ink,arc=0pt,boxrule=0pt,left=4mm,right=4mm,top=2.5mm,bottom=2.5mm]
\color{white}\sffamily{\LARGE\bfseries KRRvkqn}\hfill{\Large h\#6}\\[0.5mm]
{\small 6 Steine\quad Schwarz am Zug}\hfill{\small Nr.~743\qquad\href{https://helpman.komtera.lt/?fen=8/8/8/8/7R/6nk/2q5/K5R1&moves=6}{\color{white}\bfseries lösen\,\ensuremath{\nearrow}}}
\end{tcolorbox}
\vspace{2mm}
\noindent\begin{minipage}[t]{36mm}\vspace{0pt}\sffamily\small\raggedright
\hmfigure{längste eindeutige}{{\Large h\#6}}\hmfigure{eindeutig bei dieser Länge}{{\Large 149} Stellungen}\hmfigure{längstes Matt}{{\Large h\#7}}\hmfigure{in der Klasse}{15,5 Milliarden Stellungen,\\ 3,78 Milliarden mit Hilfsmatt}
\end{minipage}\hfill
\begin{minipage}[t]{80mm}\vspace{0pt}\begin{tikzpicture}\begin{axis}[width=72mm,height=42mm,ybar,bar width=1.6pt,xmin=-1,xmax=15,ymin=0,ymax=1.347,axis y line*=left,axis x line*=bottom,xlabel={\scriptsize Halbzüge bis Matt},ylabel={\scriptsize Stellungen (Mrd.)},tick label style={font=\tiny},label style={font=\scriptsize},enlarge x limits=false,ylabel shift=-3pt,xlabel shift=-3pt]\addplot[fill=ink!55,draw=none] coordinates {(0,0.006374) (1,0.02825) (2,0.1422) (3,0.6513) (4,1.225) (5,1.141) (6,0.4918) (7,0.08701) (8,0.006761) (9,5.838e-05) (10,0.00383) (11,1.061e-06) (12,0.001183) (13,3.4e-08) (14,2.702e-05)};\draw[densely dotted] (axis cs:12,0) -- (axis cs:12,1.347);\end{axis}\begin{axis}[width=72mm,height=42mm,axis y line*=right,axis x line=none,ymode=log,xmin=-1,xmax=15,ylabel={\scriptsize eindeutig (log)},tick label style={font=\tiny},label style={font=\scriptsize},ylabel shift=-3pt]\addplot[mark=*,mark size=0.7pt,draw=black,thin] coordinates {(0,6374086) (1,26714556) (2,60393762) (3,137002506) (4,66480908) (5,9520521) (6,526160) (7,5282) (8,64854) (9,102) (10,16342) (11,1) (12,149)};\end{axis}\end{tikzpicture}\end{minipage}
\vspace{3mm}

\noindent\begin{minipage}[t]{60mm}\vspace{0pt}\sbox\hmboard{\chessboard[setfen=8/8/8/8/7R/6nk/2q5/K5R1 b - - 0 1]}\usebox\hmboard\\[1mm]
\makebox[\wd\hmboard]{\sffamily\hspace{3mm}\textbf{h\#6}\hfill(3+3)}\end{minipage}\hfill
\begin{minipage}[t]{52mm}\vspace{0pt}
\textcolor{ink}{\sffamily\textsc{Lösung}}\par\smallskip\small
\begin{tabular}{@{}r@{\enspace}l@{\enspace}l@{}}
\textbf{1.} & \symking{}xh4 & \symrook{}c1\\
\textbf{2.} & \symqueen{}f2 & \symrook{}c2\\
\textbf{3.} & \symking{}h3 & \symking{}b2\\
\textbf{4.} & \symking{}g2 & \symking{}c3\\
\textbf{5.} & \symking{}f1 & \symking{}d3\\
\textbf{6.} & \symking{}e1 & \symrook{}c1\#\\
\end{tabular}\par\vspace{3mm}
\qrcode[height=12mm]{https://helpman.komtera.lt/?fen=8/8/8/8/7R/6nk/2q5/K5R1&moves=6}\\{\sffamily\tiny\color{ink}scannen zum Lösen}
\end{minipage}
\vspace{2mm}

{\raggedright\noindent\textcolor{ink}{\sffamily\textsc{Themen}}\enspace \themetag{reines Matt} \themetag{Mustermatt} \themetag{Switchback} \themetag{Selbstblock} \themetag{ohne Schach} \themetag{Umnov}\par\smallskip{\footnotesize\textit{Hinweis:} die Partei am Zug steht im Diagramm im Schach; die Lösung beginnt mit einem Schlag.}\par}

\clearpage
\phantomsection\label{hm:KRRvkqp}
\begin{tcolorbox}[enhanced,colback=ink,colframe=ink,arc=0pt,boxrule=0pt,left=4mm,right=4mm,top=2.5mm,bottom=2.5mm]
\color{white}\sffamily{\LARGE\bfseries KRRvkqp}\hfill{\Large h\#6}\\[0.5mm]
{\small 6 Steine\quad Schwarz am Zug}\hfill{\small Nr.~744\qquad\href{https://helpman.komtera.lt/?fen=8/8/8/7R/7k/p7/2q3R1/K7&moves=6}{\color{white}\bfseries lösen\,\ensuremath{\nearrow}}}
\end{tcolorbox}
\vspace{2mm}
\noindent\begin{minipage}[t]{36mm}\vspace{0pt}\sffamily\small\raggedright
\hmfigure{längste eindeutige}{{\Large h\#6}}\hmfigure{eindeutig bei dieser Länge}{{\Large 453} Stellungen}\hmfigure{längstes Matt}{{\Large h\#7.5}}\hmfigure{in der Klasse}{45,4 Milliarden Stellungen,\\ 11,7 Milliarden mit Hilfsmatt}
\end{minipage}\hfill
\begin{minipage}[t]{80mm}\vspace{0pt}\begin{tikzpicture}\begin{axis}[width=72mm,height=42mm,ybar,bar width=1.6pt,xmin=-1,xmax=16,ymin=0,ymax=4.288,axis y line*=left,axis x line*=bottom,xlabel={\scriptsize Halbzüge bis Matt},ylabel={\scriptsize Stellungen (Mrd.)},tick label style={font=\tiny},label style={font=\scriptsize},enlarge x limits=false,ylabel shift=-3pt,xlabel shift=-3pt]\addplot[fill=ink!55,draw=none] coordinates {(0,0.02386) (1,0.1032) (2,0.4924) (3,2.127) (4,3.898) (5,3.277) (6,1.505) (7,0.2358) (8,0.02548) (9,0.000224) (10,0.01238) (11,1.162e-05) (12,0.003506) (13,1.443e-06) (14,7.57e-05) (15,3e-09)};\draw[densely dotted] (axis cs:12,0) -- (axis cs:12,4.288);\end{axis}\begin{axis}[width=72mm,height=42mm,axis y line*=right,axis x line=none,ymode=log,xmin=-1,xmax=16,ylabel={\scriptsize eindeutig (log)},tick label style={font=\tiny},label style={font=\scriptsize},ylabel shift=-3pt]\addplot[mark=*,mark size=0.7pt,draw=black,thin] coordinates {(0,23858389) (1,97009518) (2,212062360) (3,456314005) (4,210284457) (5,28215265) (6,2101396) (7,17017) (8,192864) (9,88) (10,47622) (11,15) (12,453)};\end{axis}\end{tikzpicture}\end{minipage}
\vspace{3mm}

\noindent\begin{minipage}[t]{60mm}\vspace{0pt}\sbox\hmboard{\chessboard[setfen=8/8/8/7R/7k/p7/2q3R1/K7 b - - 0 1]}\usebox\hmboard\\[1mm]
\makebox[\wd\hmboard]{\sffamily\hspace{3mm}\textbf{h\#6}\hfill(3+3)}\end{minipage}\hfill
\begin{minipage}[t]{52mm}\vspace{0pt}
\textcolor{ink}{\sffamily\textsc{Lösung}}\par\smallskip\small
\begin{tabular}{@{}r@{\enspace}l@{\enspace}l@{}}
\textbf{1.} & \symking{}xh5 & \symrook{}xc2\\
\textbf{2.} & \symking{}g4 & \symrook{}a2\\
\textbf{3.} & \symking{}f3 & \symrook{}xa3+\\
\textbf{4.} & \symking{}e2 & \symking{}b2\\
\textbf{5.} & \symking{}d1 & \symking{}c3\\
\textbf{6.} & \symking{}c1 & \symrook{}a1\#\\
\end{tabular}\par\vspace{3mm}
\qrcode[height=12mm]{https://helpman.komtera.lt/?fen=8/8/8/7R/7k/p7/2q3R1/K7&moves=6}\\{\sffamily\tiny\color{ink}scannen zum Lösen}
\end{minipage}
\vspace{2mm}

{\raggedright\noindent\textcolor{ink}{\sffamily\textsc{Themen}}\enspace \themetag{reines Matt} \themetag{Mustermatt} \themetag{Idealmatt} \themetag{Spiegelmatt} \themetag{ein Stein zieht}\par\smallskip{\footnotesize\textit{Hinweis:} die Partei am Zug steht im Diagramm im Schach; die Lösung beginnt mit einem Schlag.}\par}

\clearpage
\phantomsection\label{hm:KRRvkqq}
\begin{tcolorbox}[enhanced,colback=ink,colframe=ink,arc=0pt,boxrule=0pt,left=4mm,right=4mm,top=2.5mm,bottom=2.5mm]
\color{white}\sffamily{\LARGE\bfseries KRRvkqq}\hfill{\Large h\#6}\\[0.5mm]
{\small 6 Steine\quad Schwarz am Zug}\hfill{\small Nr.~745\qquad\href{https://helpman.komtera.lt/?fen=8/8/8/7R/7k/1q2q3/8/K5R1&moves=6}{\color{white}\bfseries lösen\,\ensuremath{\nearrow}}}
\end{tcolorbox}
\vspace{2mm}
\noindent\begin{minipage}[t]{36mm}\vspace{0pt}\sffamily\small\raggedright
\hmfigure{längste eindeutige}{{\Large h\#6}}\hmfigure{eindeutig bei dieser Länge}{{\Large 498} Stellungen}\hmfigure{längstes Matt}{{\Large h\#8}}\hmfigure{in der Klasse}{15,5 Milliarden Stellungen,\\ 1,6 Milliarden mit Hilfsmatt}
\end{minipage}\hfill
\begin{minipage}[t]{80mm}\vspace{0pt}\begin{tikzpicture}\begin{axis}[width=72mm,height=42mm,ybar,bar width=1.6pt,xmin=-1,xmax=17,ymin=0,ymax=686.6,axis y line*=left,axis x line*=bottom,xlabel={\scriptsize Halbzüge bis Matt},ylabel={\scriptsize Stellungen (Mio.)},tick label style={font=\tiny},label style={font=\scriptsize},enlarge x limits=false,ylabel shift=-3pt,xlabel shift=-3pt]\addplot[fill=ink!55,draw=none] coordinates {(0,1.316) (1,5.493) (2,33.39) (3,212.2) (4,405) (5,624.2) (6,221.5) (7,90.96) (8,4.02) (9,0.1818) (10,1.36) (11,0.001191) (12,0.5375) (13,0.000181) (14,0.02387) (16,3.7e-05)};\draw[densely dotted] (axis cs:12,0) -- (axis cs:12,686.6);\end{axis}\begin{axis}[width=72mm,height=42mm,axis y line*=right,axis x line=none,ymode=log,xmin=-1,xmax=17,ylabel={\scriptsize eindeutig (log)},tick label style={font=\tiny},label style={font=\scriptsize},ylabel shift=-3pt]\addplot[mark=*,mark size=0.7pt,draw=black,thin] coordinates {(0,1316340) (1,5339250) (2,14284291) (3,48406987) (4,21756427) (5,8009973) (6,406567) (7,23765) (8,21333) (9,513) (10,5756) (11,7) (12,498)};\end{axis}\end{tikzpicture}\end{minipage}
\vspace{3mm}

\noindent\begin{minipage}[t]{60mm}\vspace{0pt}\sbox\hmboard{\chessboard[setfen=8/8/8/7R/7k/1q2q3/8/K5R1 b - - 0 1]}\usebox\hmboard\\[1mm]
\makebox[\wd\hmboard]{\sffamily\hspace{3mm}\textbf{h\#6}\hfill(3+3)}\end{minipage}\hfill
\begin{minipage}[t]{52mm}\vspace{0pt}
\textcolor{ink}{\sffamily\textsc{Lösung}}\par\smallskip\small
\begin{tabular}{@{}r@{\enspace}l@{\enspace}l@{}}
\textbf{1.} & \symking{}xh5 & \symrook{}b1\\
\textbf{2.} & \symking{}g4 & \symrook{}xb3\\
\textbf{3.} & \symking{}f3 & \symking{}b2\\
\textbf{4.} & \symking{}e2 & \symrook{}xe3+\\
\textbf{5.} & \symking{}d1 & \symking{}c3\\
\textbf{6.} & \symking{}c1 & \symrook{}e1\#\\
\end{tabular}\par\vspace{3mm}
\qrcode[height=12mm]{https://helpman.komtera.lt/?fen=8/8/8/7R/7k/1q2q3/8/K5R1&moves=6}\\{\sffamily\tiny\color{ink}scannen zum Lösen}
\end{minipage}
\vspace{2mm}

{\raggedright\noindent\textcolor{ink}{\sffamily\textsc{Themen}}\enspace \themetag{reines Matt} \themetag{Mustermatt} \themetag{Idealmatt} \themetag{Spiegelmatt} \themetag{ein Stein zieht}\par\smallskip{\footnotesize\textit{Hinweis:} die Partei am Zug steht im Diagramm im Schach; die Lösung beginnt mit einem Schlag.}\par}

\clearpage
\phantomsection\label{hm:KRRvkqr}
\begin{tcolorbox}[enhanced,colback=ink,colframe=ink,arc=0pt,boxrule=0pt,left=4mm,right=4mm,top=2.5mm,bottom=2.5mm]
\color{white}\sffamily{\LARGE\bfseries KRRvkqr}\hfill{\Large h\#6}\\[0.5mm]
{\small 6 Steine\quad Schwarz am Zug}\hfill{\small Nr.~746\qquad\href{https://helpman.komtera.lt/?fen=8/8/8/4R3/R4k2/5rq1/8/K7&moves=6}{\color{white}\bfseries lösen\,\ensuremath{\nearrow}}}
\end{tcolorbox}
\vspace{2mm}
\noindent\begin{minipage}[t]{36mm}\vspace{0pt}\sffamily\small\raggedright
\hmfigure{längste eindeutige}{{\Large h\#6}}\hmfigure{eindeutig bei dieser Länge}{{\Large 797} Stellungen}\hmfigure{längstes Matt}{{\Large h\#8}}\hmfigure{in der Klasse}{15,5 Milliarden Stellungen,\\ 3,5 Milliarden mit Hilfsmatt}
\end{minipage}\hfill
\begin{minipage}[t]{80mm}\vspace{0pt}\begin{tikzpicture}\begin{axis}[width=72mm,height=42mm,ybar,bar width=1.6pt,xmin=-1,xmax=17,ymin=0,ymax=1.326,axis y line*=left,axis x line*=bottom,xlabel={\scriptsize Halbzüge bis Matt},ylabel={\scriptsize Stellungen (Mrd.)},tick label style={font=\tiny},label style={font=\scriptsize},enlarge x limits=false,ylabel shift=-3pt,xlabel shift=-3pt]\addplot[fill=ink!55,draw=none] coordinates {(0,0.004365) (1,0.01911) (2,0.1058) (3,0.5559) (4,1.05) (5,1.205) (6,0.4358) (7,0.1134) (8,0.005586) (9,8.135e-05) (10,0.003327) (11,1.086e-06) (12,0.001188) (13,2.51e-07) (14,4e-05) (16,1.4e-08)};\draw[densely dotted] (axis cs:12,0) -- (axis cs:12,1.326);\end{axis}\begin{axis}[width=72mm,height=42mm,axis y line*=right,axis x line=none,ymode=log,xmin=-1,xmax=17,ylabel={\scriptsize eindeutig (log)},tick label style={font=\tiny},label style={font=\scriptsize},ylabel shift=-3pt]\addplot[mark=*,mark size=0.7pt,draw=black,thin] coordinates {(0,4365335) (1,18404243) (2,47245187) (3,124756041) (4,60404410) (5,11985856) (6,413580) (7,10665) (8,48767) (9,258) (10,14767) (11,7) (12,797)};\end{axis}\end{tikzpicture}\end{minipage}
\vspace{3mm}

\noindent\begin{minipage}[t]{60mm}\vspace{0pt}\sbox\hmboard{\chessboard[setfen=8/8/8/4R3/R4k2/5rq1/8/K7 b - - 0 1]}\usebox\hmboard\\[1mm]
\makebox[\wd\hmboard]{\sffamily\hspace{3mm}\textbf{h\#6}\hfill(3+3)}\end{minipage}\hfill
\begin{minipage}[t]{52mm}\vspace{0pt}
\textcolor{ink}{\sffamily\textsc{Lösung}}\par\smallskip\small
\begin{tabular}{@{}r@{\enspace}l@{\enspace}l@{}}
\textbf{1.} & \symking{}xe5 & \symrook{}g4\\
\textbf{2.} & \symking{}d5 & \symrook{}xg3\\
\textbf{3.} & \symking{}c4 & \symking{}b1\\
\textbf{4.} & \symking{}b3 & \symking{}c1\\
\textbf{5.} & \symking{}a2 & \symrook{}xf3\\
\textbf{6.} & \symking{}a1 & \symrook{}a3\#\\
\end{tabular}\par\vspace{3mm}
\qrcode[height=12mm]{https://helpman.komtera.lt/?fen=8/8/8/4R3/R4k2/5rq1/8/K7&moves=6}\\{\sffamily\tiny\color{ink}scannen zum Lösen}
\end{minipage}
\vspace{2mm}

{\raggedright\noindent\textcolor{ink}{\sffamily\textsc{Themen}}\enspace \themetag{reines Matt} \themetag{Mustermatt} \themetag{Idealmatt} \themetag{Spiegelmatt} \themetag{ein Stein zieht} \themetag{ohne Schach}\par\smallskip{\footnotesize\textit{Hinweis:} die Partei am Zug steht im Diagramm im Schach; die Lösung beginnt mit einem Schlag.}\par}

\clearpage
\phantomsection\label{hm:KRRvkrb}
\begin{tcolorbox}[enhanced,colback=ink,colframe=ink,arc=0pt,boxrule=0pt,left=4mm,right=4mm,top=2.5mm,bottom=2.5mm]
\color{white}\sffamily{\LARGE\bfseries KRRvkrb}\hfill{\Large h\#6}\\[0.5mm]
{\small 6 Steine\quad Schwarz am Zug}\hfill{\small Nr.~747\qquad\href{https://helpman.komtera.lt/?fen=8/8/1r6/7R/7k/8/8/K4bR1&moves=6}{\color{white}\bfseries lösen\,\ensuremath{\nearrow}}}
\end{tcolorbox}
\vspace{2mm}
\noindent\begin{minipage}[t]{36mm}\vspace{0pt}\sffamily\small\raggedright
\hmfigure{längste eindeutige}{{\Large h\#6}}\hmfigure{eindeutig bei dieser Länge}{{\Large 388} Stellungen}\hmfigure{längstes Matt}{{\Large h\#7}}\hmfigure{in der Klasse}{15,5 Milliarden Stellungen,\\ 4,03 Milliarden mit Hilfsmatt}
\end{minipage}\hfill
\begin{minipage}[t]{80mm}\vspace{0pt}\begin{tikzpicture}\begin{axis}[width=72mm,height=42mm,ybar,bar width=1.6pt,xmin=-1,xmax=15,ymin=0,ymax=1.555,axis y line*=left,axis x line*=bottom,xlabel={\scriptsize Halbzüge bis Matt},ylabel={\scriptsize Stellungen (Mrd.)},tick label style={font=\tiny},label style={font=\scriptsize},enlarge x limits=false,ylabel shift=-3pt,xlabel shift=-3pt]\addplot[fill=ink!55,draw=none] coordinates {(0,0.00881) (1,0.03869) (2,0.1757) (3,0.7378) (4,1.413) (5,1.065) (6,0.5091) (7,0.06781) (8,0.007673) (9,2.153e-05) (10,0.004429) (11,7.1e-08) (12,0.001265) (13,1e-09) (14,2.093e-05)};\draw[densely dotted] (axis cs:12,0) -- (axis cs:12,1.555);\end{axis}\begin{axis}[width=72mm,height=42mm,axis y line*=right,axis x line=none,ymode=log,xmin=-1,xmax=15,ylabel={\scriptsize eindeutig (log)},tick label style={font=\tiny},label style={font=\scriptsize},ylabel shift=-3pt]\addplot[mark=*,mark size=0.7pt,draw=black,thin] coordinates {(0,8809794) (1,36323638) (2,78074864) (3,150311246) (4,76065460) (5,7732707) (6,643736) (7,2969) (8,67905) (9,35) (10,21880) (12,388)};\end{axis}\end{tikzpicture}\end{minipage}
\vspace{3mm}

\noindent\begin{minipage}[t]{60mm}\vspace{0pt}\sbox\hmboard{\chessboard[setfen=8/8/1r6/7R/7k/8/8/K4bR1 b - - 0 1]}\usebox\hmboard\\[1mm]
\makebox[\wd\hmboard]{\sffamily\hspace{3mm}\textbf{h\#6}\hfill(3+3)}\end{minipage}\hfill
\begin{minipage}[t]{52mm}\vspace{0pt}
\textcolor{ink}{\sffamily\textsc{Lösung}}\par\smallskip\small
\begin{tabular}{@{}r@{\enspace}l@{\enspace}l@{}}
\textbf{1.} & \symking{}xh5 & \symrook{}h1+\\
\textbf{2.} & \symking{}g4 & \symrook{}h6\\
\textbf{3.} & \symking{}f3 & \symrook{}xb6\\
\textbf{4.} & \symking{}e2 & \symking{}b2\\
\textbf{5.} & \symking{}d1 & \symking{}c3\\
\textbf{6.} & \symbishop{}e2 & \symrook{}b1\#\\
\end{tabular}\par\vspace{3mm}
\qrcode[height=12mm]{https://helpman.komtera.lt/?fen=8/8/1r6/7R/7k/8/8/K4bR1&moves=6}\\{\sffamily\tiny\color{ink}scannen zum Lösen}
\end{minipage}
\vspace{2mm}

{\raggedright\noindent\textcolor{ink}{\sffamily\textsc{Themen}}\enspace \themetag{reines Matt} \themetag{Mustermatt} \themetag{Idealmatt} \themetag{Selbstblock}\par\smallskip{\footnotesize\textit{Hinweis:} die Partei am Zug steht im Diagramm im Schach; die Lösung beginnt mit einem Schlag.}\par}

\clearpage
\phantomsection\label{hm:KRRvkrn}
\begin{tcolorbox}[enhanced,colback=ink,colframe=ink,arc=0pt,boxrule=0pt,left=4mm,right=4mm,top=2.5mm,bottom=2.5mm]
\color{white}\sffamily{\LARGE\bfseries KRRvkrn}\hfill{\Large h\#6}\\[0.5mm]
{\small 6 Steine\quad Schwarz am Zug}\hfill{\small Nr.~748\qquad\href{https://helpman.komtera.lt/?fen=8/8/8/7R/1r5k/8/8/K3n1R1&moves=6}{\color{white}\bfseries lösen\,\ensuremath{\nearrow}}}
\end{tcolorbox}
\vspace{2mm}
\noindent\begin{minipage}[t]{36mm}\vspace{0pt}\sffamily\small\raggedright
\hmfigure{längste eindeutige}{{\Large h\#6}}\hmfigure{eindeutig bei dieser Länge}{{\Large 166} Stellungen}\hmfigure{längstes Matt}{{\Large h\#7}}\hmfigure{in der Klasse}{15,5 Milliarden Stellungen,\\ 4,14 Milliarden mit Hilfsmatt}
\end{minipage}\hfill
\begin{minipage}[t]{80mm}\vspace{0pt}\begin{tikzpicture}\begin{axis}[width=72mm,height=42mm,ybar,bar width=1.6pt,xmin=-1,xmax=15,ymin=0,ymax=1.649,axis y line*=left,axis x line*=bottom,xlabel={\scriptsize Halbzüge bis Matt},ylabel={\scriptsize Stellungen (Mrd.)},tick label style={font=\tiny},label style={font=\scriptsize},enlarge x limits=false,ylabel shift=-3pt,xlabel shift=-3pt]\addplot[fill=ink!55,draw=none] coordinates {(0,0.01091) (1,0.04876) (2,0.2063) (3,0.7927) (4,1.5) (5,1.01) (6,0.4975) (7,0.05811) (8,0.008093) (9,3.026e-05) (10,0.004521) (11,5.5e-08) (12,0.001067) (13,1e-09) (14,1.303e-05)};\draw[densely dotted] (axis cs:12,0) -- (axis cs:12,1.649);\end{axis}\begin{axis}[width=72mm,height=42mm,axis y line*=right,axis x line=none,ymode=log,xmin=-1,xmax=15,ylabel={\scriptsize eindeutig (log)},tick label style={font=\tiny},label style={font=\scriptsize},ylabel shift=-3pt]\addplot[mark=*,mark size=0.7pt,draw=black,thin] coordinates {(0,10911489) (1,45267799) (2,89408868) (3,148701919) (4,79421932) (5,6300999) (6,613135) (7,3091) (8,77784) (9,28) (10,18560) (12,166)};\end{axis}\end{tikzpicture}\end{minipage}
\vspace{3mm}

\noindent\begin{minipage}[t]{60mm}\vspace{0pt}\sbox\hmboard{\chessboard[setfen=8/8/8/7R/1r5k/8/8/K3n1R1 b - - 0 1]}\usebox\hmboard\\[1mm]
\makebox[\wd\hmboard]{\sffamily\hspace{3mm}\textbf{h\#6}\hfill(3+3)}\end{minipage}\hfill
\begin{minipage}[t]{52mm}\vspace{0pt}
\textcolor{ink}{\sffamily\textsc{Lösung}}\par\smallskip\small
\begin{tabular}{@{}r@{\enspace}l@{\enspace}l@{}}
\textbf{1.} & \symking{}xh5 & \symrook{}h1+\\
\textbf{2.} & \symking{}g4 & \symking{}a2\\
\textbf{3.} & \symking{}f3 & \symking{}a3\\
\textbf{4.} & \symking{}e2 & \symking{}xb4\\
\textbf{5.} & \symking{}d1 & \symking{}c3\\
\textbf{6.} & \symking{}c1 & \symrook{}xe1\#\\
\end{tabular}\par\vspace{3mm}
\qrcode[height=12mm]{https://helpman.komtera.lt/?fen=8/8/8/7R/1r5k/8/8/K3n1R1&moves=6}\\{\sffamily\tiny\color{ink}scannen zum Lösen}
\end{minipage}
\vspace{2mm}

{\raggedright\noindent\textcolor{ink}{\sffamily\textsc{Themen}}\enspace \themetag{reines Matt} \themetag{Mustermatt} \themetag{Idealmatt} \themetag{Spiegelmatt} \themetag{ein Stein zieht}\par\smallskip{\footnotesize\textit{Hinweis:} die Partei am Zug steht im Diagramm im Schach; die Lösung beginnt mit einem Schlag.}\par}

\clearpage
\phantomsection\label{hm:KRRvkrp}
\begin{tcolorbox}[enhanced,colback=ink,colframe=ink,arc=0pt,boxrule=0pt,left=4mm,right=4mm,top=2.5mm,bottom=2.5mm]
\color{white}\sffamily{\LARGE\bfseries KRRvkrp}\hfill{\Large h\#6}\\[0.5mm]
{\small 6 Steine\quad Schwarz am Zug}\hfill{\small Nr.~749\qquad\href{https://helpman.komtera.lt/?fen=8/8/8/7R/7k/1r6/5p2/K5R1&moves=6}{\color{white}\bfseries lösen\,\ensuremath{\nearrow}}}
\end{tcolorbox}
\vspace{2mm}
\noindent\begin{minipage}[t]{36mm}\vspace{0pt}\sffamily\small\raggedright
\hmfigure{längste eindeutige}{{\Large h\#6}}\hmfigure{eindeutig bei dieser Länge}{{\Large 797} Stellungen}\hmfigure{längstes Matt}{{\Large h\#7.5}}\hmfigure{in der Klasse}{45,4 Milliarden Stellungen,\\ 12,8 Milliarden mit Hilfsmatt}
\end{minipage}\hfill
\begin{minipage}[t]{80mm}\vspace{0pt}\begin{tikzpicture}\begin{axis}[width=72mm,height=42mm,ybar,bar width=1.6pt,xmin=-1,xmax=16,ymin=0,ymax=5.263,axis y line*=left,axis x line*=bottom,xlabel={\scriptsize Halbzüge bis Matt},ylabel={\scriptsize Stellungen (Mrd.)},tick label style={font=\tiny},label style={font=\scriptsize},enlarge x limits=false,ylabel shift=-3pt,xlabel shift=-3pt]\addplot[fill=ink!55,draw=none] coordinates {(0,0.04127) (1,0.179) (2,0.7105) (3,2.574) (4,4.785) (5,2.848) (6,1.503) (7,0.1472) (8,0.0273) (9,7.608e-05) (10,0.01505) (11,4.56e-07) (12,0.003719) (13,6e-09) (14,6.144e-05) (15,1e-09)};\draw[densely dotted] (axis cs:12,0) -- (axis cs:12,5.263);\end{axis}\begin{axis}[width=72mm,height=42mm,axis y line*=right,axis x line=none,ymode=log,xmin=-1,xmax=16,ylabel={\scriptsize eindeutig (log)},tick label style={font=\tiny},label style={font=\scriptsize},ylabel shift=-3pt]\addplot[mark=*,mark size=0.7pt,draw=black,thin] coordinates {(0,41270205) (1,165029301) (2,315946295) (3,486063823) (4,259038461) (5,17620814) (6,2207903) (7,7577) (8,253365) (9,28) (10,66761) (12,797)};\end{axis}\end{tikzpicture}\end{minipage}
\vspace{3mm}

\noindent\begin{minipage}[t]{60mm}\vspace{0pt}\sbox\hmboard{\chessboard[setfen=8/8/8/7R/7k/1r6/5p2/K5R1 b - - 0 1]}\usebox\hmboard\\[1mm]
\makebox[\wd\hmboard]{\sffamily\hspace{3mm}\textbf{h\#6}\hfill(3+3)}\end{minipage}\hfill
\begin{minipage}[t]{52mm}\vspace{0pt}
\textcolor{ink}{\sffamily\textsc{Lösung}}\par\smallskip\small
\begin{tabular}{@{}r@{\enspace}l@{\enspace}l@{}}
\textbf{1.} & \symking{}xh5 & \symrook{}f1\\
\textbf{2.} & \symking{}g4 & \symking{}a2\\
\textbf{3.} & \symking{}f3 & \symking{}xb3\\
\textbf{4.} & \symking{}e2 & \symrook{}xf2+\\
\textbf{5.} & \symking{}d1 & \symking{}c3\\
\textbf{6.} & \symking{}c1 & \symrook{}f1\#\\
\end{tabular}\par\vspace{3mm}
\qrcode[height=12mm]{https://helpman.komtera.lt/?fen=8/8/8/7R/7k/1r6/5p2/K5R1&moves=6}\\{\sffamily\tiny\color{ink}scannen zum Lösen}
\end{minipage}
\vspace{2mm}

{\raggedright\noindent\textcolor{ink}{\sffamily\textsc{Themen}}\enspace \themetag{reines Matt} \themetag{Mustermatt} \themetag{Idealmatt} \themetag{Spiegelmatt} \themetag{Switchback} \themetag{ein Stein zieht}\par\smallskip{\footnotesize\textit{Hinweis:} die Partei am Zug steht im Diagramm im Schach; die Lösung beginnt mit einem Schlag.}\par}

\clearpage
\phantomsection\label{hm:KRRvkrr}
\begin{tcolorbox}[enhanced,colback=ink,colframe=ink,arc=0pt,boxrule=0pt,left=4mm,right=4mm,top=2.5mm,bottom=2.5mm]
\color{white}\sffamily{\LARGE\bfseries KRRvkrr}\hfill{\Large h\#6}\\[0.5mm]
{\small 6 Steine\quad Schwarz am Zug}\hfill{\small Nr.~750\qquad\href{https://helpman.komtera.lt/?fen=8/8/8/4R3/R4k2/5rr1/8/K7&moves=6}{\color{white}\bfseries lösen\,\ensuremath{\nearrow}}}
\end{tcolorbox}
\vspace{2mm}
\noindent\begin{minipage}[t]{36mm}\vspace{0pt}\sffamily\small\raggedright
\hmfigure{längste eindeutige}{{\Large h\#6}}\hmfigure{eindeutig bei dieser Länge}{{\Large 424} Stellungen}\hmfigure{längstes Matt}{{\Large h\#8}}\hmfigure{in der Klasse}{15,5 Milliarden Stellungen,\\ 1,91 Milliarden mit Hilfsmatt}
\end{minipage}\hfill
\begin{minipage}[t]{80mm}\vspace{0pt}\begin{tikzpicture}\begin{axis}[width=72mm,height=42mm,ybar,bar width=1.6pt,xmin=-1,xmax=17,ymin=0,ymax=713.9,axis y line*=left,axis x line*=bottom,xlabel={\scriptsize Halbzüge bis Matt},ylabel={\scriptsize Stellungen (Mio.)},tick label style={font=\tiny},label style={font=\scriptsize},enlarge x limits=false,ylabel shift=-3pt,xlabel shift=-3pt]\addplot[fill=ink!55,draw=none] coordinates {(0,3.628) (1,16.32) (2,78.84) (3,345.3) (4,649) (5,551.8) (6,217.9) (7,38.15) (8,3.094) (9,0.03069) (10,2.054) (11,0.000244) (12,0.682) (13,7.5e-05) (14,0.01819) (16,3e-06)};\draw[densely dotted] (axis cs:12,0) -- (axis cs:12,713.9);\end{axis}\begin{axis}[width=72mm,height=42mm,axis y line*=right,axis x line=none,ymode=log,xmin=-1,xmax=17,ylabel={\scriptsize eindeutig (log)},tick label style={font=\tiny},label style={font=\scriptsize},ylabel shift=-3pt]\addplot[mark=*,mark size=0.7pt,draw=black,thin] coordinates {(0,3627896) (1,15535438) (2,35819865) (3,72358807) (4,36326764) (5,3937339) (6,185472) (7,3213) (8,26712) (9,10) (10,9291) (12,424)};\end{axis}\end{tikzpicture}\end{minipage}
\vspace{3mm}

\noindent\begin{minipage}[t]{60mm}\vspace{0pt}\sbox\hmboard{\chessboard[setfen=8/8/8/4R3/R4k2/5rr1/8/K7 b - - 0 1]}\usebox\hmboard\\[1mm]
\makebox[\wd\hmboard]{\sffamily\hspace{3mm}\textbf{h\#6}\hfill(3+3)}\end{minipage}\hfill
\begin{minipage}[t]{52mm}\vspace{0pt}
\textcolor{ink}{\sffamily\textsc{Lösung}}\par\smallskip\small
\begin{tabular}{@{}r@{\enspace}l@{\enspace}l@{}}
\textbf{1.} & \symking{}xe5 & \symrook{}g4\\
\textbf{2.} & \symking{}d5 & \symrook{}xg3\\
\textbf{3.} & \symking{}c4 & \symking{}b1\\
\textbf{4.} & \symking{}b3 & \symking{}c1\\
\textbf{5.} & \symking{}a2 & \symrook{}xf3\\
\textbf{6.} & \symking{}a1 & \symrook{}a3\#\\
\end{tabular}\par\vspace{3mm}
\qrcode[height=12mm]{https://helpman.komtera.lt/?fen=8/8/8/4R3/R4k2/5rr1/8/K7&moves=6}\\{\sffamily\tiny\color{ink}scannen zum Lösen}
\end{minipage}
\vspace{2mm}

{\raggedright\noindent\textcolor{ink}{\sffamily\textsc{Themen}}\enspace \themetag{reines Matt} \themetag{Mustermatt} \themetag{Idealmatt} \themetag{Spiegelmatt} \themetag{ein Stein zieht} \themetag{ohne Schach}\par\smallskip{\footnotesize\textit{Hinweis:} die Partei am Zug steht im Diagramm im Schach; die Lösung beginnt mit einem Schlag.}\par}

\clearpage
\phantomsection\label{hm:KRvkbbb}
\begin{tcolorbox}[enhanced,colback=ink,colframe=ink,arc=0pt,boxrule=0pt,left=4mm,right=4mm,top=2.5mm,bottom=2.5mm]
\color{white}\sffamily{\LARGE\bfseries KRvkbbb}\hfill{\Large h\#6}\\[0.5mm]
{\small 6 Steine\quad Schwarz am Zug}\hfill{\small Nr.~751\qquad\href{https://helpman.komtera.lt/?fen=8/8/7k/5b2/4R3/b7/8/K1b5&moves=6}{\color{white}\bfseries lösen\,\ensuremath{\nearrow}}}
\end{tcolorbox}
\vspace{2mm}
\noindent\begin{minipage}[t]{36mm}\vspace{0pt}\sffamily\small\raggedright
\hmfigure{längste eindeutige}{{\Large h\#6}}\hmfigure{eindeutig bei dieser Länge}{{\Large 408} Stellungen}\hmfigure{längstes Matt}{{\Large h\#7.5}}\hmfigure{in der Klasse}{15,5 Milliarden Stellungen,\\ 1,46 Milliarden mit Hilfsmatt}
\end{minipage}\hfill
\begin{minipage}[t]{80mm}\vspace{0pt}\begin{tikzpicture}\begin{axis}[width=72mm,height=42mm,ybar,bar width=1.6pt,xmin=-1,xmax=16,ymin=0,ymax=321.6,axis y line*=left,axis x line*=bottom,xlabel={\scriptsize Halbzüge bis Matt},ylabel={\scriptsize Stellungen (Mio.)},tick label style={font=\tiny},label style={font=\scriptsize},enlarge x limits=false,ylabel shift=-3pt,xlabel shift=-3pt]\addplot[fill=ink!55,draw=none] coordinates {(0,0.1454) (1,0.9809) (2,3.858) (3,23.92) (4,46.93) (5,144.7) (6,173.3) (7,292.4) (8,260.4) (9,273.3) (10,148.2) (11,71.45) (12,19.51) (13,1.096) (14,0.1749) (15,0.001525)};\draw[densely dotted] (axis cs:12,0) -- (axis cs:12,321.6);\end{axis}\begin{axis}[width=72mm,height=42mm,axis y line*=right,axis x line=none,ymode=log,xmin=-1,xmax=16,ylabel={\scriptsize eindeutig (log)},tick label style={font=\tiny},label style={font=\scriptsize},ylabel shift=-3pt]\addplot[mark=*,mark size=0.7pt,draw=black,thin] coordinates {(0,145403) (1,980894) (2,1623961) (3,7085406) (4,8519912) (5,12306541) (6,6204965) (7,4580950) (8,2290507) (9,1096856) (10,264785) (11,14075) (12,408)};\end{axis}\end{tikzpicture}\end{minipage}
\vspace{3mm}

\noindent\begin{minipage}[t]{60mm}\vspace{0pt}\sbox\hmboard{\chessboard[setfen=8/8/7k/5b2/4R3/b7/8/K1b5 b - - 0 1]}\usebox\hmboard\\[1mm]
\makebox[\wd\hmboard]{\sffamily\hspace{3mm}\textbf{h\#6}\hfill(2+4)}\end{minipage}\hfill
\begin{minipage}[t]{52mm}\vspace{0pt}
\textcolor{ink}{\sffamily\textsc{Lösung}}\par\smallskip\small
\begin{tabular}{@{}r@{\enspace}l@{\enspace}l@{}}
\textbf{1.} & \symking{}g5 & \symrook{}e5\\
\textbf{2.} & \symking{}f4 & \symrook{}xf5+\\
\textbf{3.} & \symking{}e3 & \symking{}b1\\
\textbf{4.} & \symking{}e2 & \symking{}c2\\
\textbf{5.} & \symking{}e1 & \symking{}d3\\
\textbf{6.} & \symking{}d1 & \symrook{}f1\#\\
\end{tabular}\par\vspace{3mm}
\qrcode[height=12mm]{https://helpman.komtera.lt/?fen=8/8/7k/5b2/4R3/b7/8/K1b5&moves=6}\\{\sffamily\tiny\color{ink}scannen zum Lösen}
\end{minipage}
\vspace{2mm}

{\raggedright\noindent\textcolor{ink}{\sffamily\textsc{Themen}}\enspace \themetag{ein Stein zieht}\par}

\clearpage
\phantomsection\label{hm:KRvkbbn}
\begin{tcolorbox}[enhanced,colback=ink,colframe=ink,arc=0pt,boxrule=0pt,left=4mm,right=4mm,top=2.5mm,bottom=2.5mm]
\color{white}\sffamily{\LARGE\bfseries KRvkbbn}\hfill{\Large h\#6}\\[0.5mm]
{\small 6 Steine\quad Schwarz am Zug}\hfill{\small Nr.~752\qquad\href{https://helpman.komtera.lt/?fen=8/8/6k1/5bn1/3R4/8/2b5/K7&moves=6}{\color{white}\bfseries lösen\,\ensuremath{\nearrow}}}
\end{tcolorbox}
\vspace{2mm}
\noindent\begin{minipage}[t]{36mm}\vspace{0pt}\sffamily\small\raggedright
\hmfigure{längste eindeutige}{{\Large h\#6}}\hmfigure{eindeutig bei dieser Länge}{{\Large 1\,232} Stellungen}\hmfigure{längstes Matt}{{\Large h\#7.5}}\hmfigure{in der Klasse}{15,5 Milliarden Stellungen,\\ 4,52 Milliarden mit Hilfsmatt}
\end{minipage}\hfill
\begin{minipage}[t]{80mm}\vspace{0pt}\begin{tikzpicture}\begin{axis}[width=72mm,height=42mm,ybar,bar width=1.6pt,xmin=-1,xmax=16,ymin=0,ymax=970.2,axis y line*=left,axis x line*=bottom,xlabel={\scriptsize Halbzüge bis Matt},ylabel={\scriptsize Stellungen (Mio.)},tick label style={font=\tiny},label style={font=\scriptsize},enlarge x limits=false,ylabel shift=-3pt,xlabel shift=-3pt]\addplot[fill=ink!55,draw=none] coordinates {(0,0.5901) (1,3.995) (2,14.72) (3,83.36) (4,165.1) (5,462.3) (6,579.7) (7,882) (8,826.9) (9,797) (10,452.8) (11,193.1) (12,58.95) (13,2.474) (14,0.5257) (15,0.001217)};\draw[densely dotted] (axis cs:12,0) -- (axis cs:12,970.2);\end{axis}\begin{axis}[width=72mm,height=42mm,axis y line*=right,axis x line=none,ymode=log,xmin=-1,xmax=16,ylabel={\scriptsize eindeutig (log)},tick label style={font=\tiny},label style={font=\scriptsize},ylabel shift=-3pt]\addplot[mark=*,mark size=0.7pt,draw=black,thin] coordinates {(0,590140) (1,3994818) (2,6330322) (3,25311878) (4,32623729) (5,41051273) (6,22257296) (7,13441795) (8,7224756) (9,2955895) (10,792142) (11,22931) (12,1232)};\end{axis}\end{tikzpicture}\end{minipage}
\vspace{3mm}

\noindent\begin{minipage}[t]{60mm}\vspace{0pt}\sbox\hmboard{\chessboard[setfen=8/8/6k1/5bn1/3R4/8/2b5/K7 b - - 0 1]}\usebox\hmboard\\[1mm]
\makebox[\wd\hmboard]{\sffamily\hspace{3mm}\textbf{h\#6}\hfill(2+4)}\end{minipage}\hfill
\begin{minipage}[t]{52mm}\vspace{0pt}
\textcolor{ink}{\sffamily\textsc{Lösung}}\par\smallskip\small
\begin{tabular}{@{}r@{\enspace}l@{\enspace}l@{}}
\textbf{1.} & \symking{}h5 & \symrook{}d5\\
\textbf{2.} & \symking{}g4 & \symking{}b2\\
\textbf{3.} & \symking{}f3 & \symrook{}xf5+\\
\textbf{4.} & \symking{}e2 & \symking{}xc2\\
\textbf{5.} & \symking{}e1 & \symking{}d3\\
\textbf{6.} & \symking{}d1 & \symrook{}f1\#\\
\end{tabular}\par\vspace{3mm}
\qrcode[height=12mm]{https://helpman.komtera.lt/?fen=8/8/6k1/5bn1/3R4/8/2b5/K7&moves=6}\\{\sffamily\tiny\color{ink}scannen zum Lösen}
\end{minipage}
\vspace{2mm}

{\raggedright\noindent\textcolor{ink}{\sffamily\textsc{Themen}}\enspace \themetag{reines Matt} \themetag{Mustermatt} \themetag{Spiegelmatt} \themetag{ein Stein zieht}\par}

\clearpage
\phantomsection\label{hm:KRvkbnn}
\begin{tcolorbox}[enhanced,colback=ink,colframe=ink,arc=0pt,boxrule=0pt,left=4mm,right=4mm,top=2.5mm,bottom=2.5mm]
\color{white}\sffamily{\LARGE\bfseries KRvkbnn}\hfill{\Large h\#6}\\[0.5mm]
{\small 6 Steine\quad Schwarz am Zug}\hfill{\small Nr.~753\qquad\href{https://helpman.komtera.lt/?fen=8/8/8/7k/n2R4/8/2b5/K2n4&moves=6}{\color{white}\bfseries lösen\,\ensuremath{\nearrow}}}
\end{tcolorbox}
\vspace{2mm}
\noindent\begin{minipage}[t]{36mm}\vspace{0pt}\sffamily\small\raggedright
\hmfigure{längste eindeutige}{{\Large h\#6}}\hmfigure{eindeutig bei dieser Länge}{{\Large 1\,834} Stellungen}\hmfigure{längstes Matt}{{\Large h\#7.5}}\hmfigure{in der Klasse}{15,5 Milliarden Stellungen,\\ 4,65 Milliarden mit Hilfsmatt}
\end{minipage}\hfill
\begin{minipage}[t]{80mm}\vspace{0pt}\begin{tikzpicture}\begin{axis}[width=72mm,height=42mm,ybar,bar width=1.6pt,xmin=-1,xmax=16,ymin=0,ymax=968.2,axis y line*=left,axis x line*=bottom,xlabel={\scriptsize Halbzüge bis Matt},ylabel={\scriptsize Stellungen (Mio.)},tick label style={font=\tiny},label style={font=\scriptsize},enlarge x limits=false,ylabel shift=-3pt,xlabel shift=-3pt]\addplot[fill=ink!55,draw=none] coordinates {(0,0.787) (1,5.323) (2,17.98) (3,93.02) (4,184.8) (5,477.1) (6,623.1) (7,880.1) (8,866.6) (9,783) (10,467.9) (11,183.7) (12,62.39) (13,2.321) (14,0.553) (15,0.000692)};\draw[densely dotted] (axis cs:12,0) -- (axis cs:12,968.2);\end{axis}\begin{axis}[width=72mm,height=42mm,axis y line*=right,axis x line=none,ymode=log,xmin=-1,xmax=16,ylabel={\scriptsize eindeutig (log)},tick label style={font=\tiny},label style={font=\scriptsize},ylabel shift=-3pt]\addplot[mark=*,mark size=0.7pt,draw=black,thin] coordinates {(0,786962) (1,5323329) (2,8001781) (3,29562505) (4,40248156) (5,44863512) (6,25579555) (7,13006005) (8,7501522) (9,2712515) (10,840406) (11,15838) (12,1834)};\end{axis}\end{tikzpicture}\end{minipage}
\vspace{3mm}

\noindent\begin{minipage}[t]{60mm}\vspace{0pt}\sbox\hmboard{\chessboard[setfen=8/8/8/7k/n2R4/8/2b5/K2n4 b - - 0 1]}\usebox\hmboard\\[1mm]
\makebox[\wd\hmboard]{\sffamily\hspace{3mm}\textbf{h\#6}\hfill(2+4)}\end{minipage}\hfill
\begin{minipage}[t]{52mm}\vspace{0pt}
\textcolor{ink}{\sffamily\textsc{Lösung}}\par\smallskip\small
\begin{tabular}{@{}r@{\enspace}l@{\enspace}l@{}}
\textbf{1.} & \symbishop{}b3 & \symrook{}xd1\\
\textbf{2.} & \symking{}g4 & \symrook{}d4+\\
\textbf{3.} & \symking{}f3 & \symrook{}xa4\\
\textbf{4.} & \symking{}e2 & \symking{}b2\\
\textbf{5.} & \symking{}d1 & \symking{}c3\\
\textbf{6.} & \symking{}c1 & \symrook{}a1\#\\
\end{tabular}\par\vspace{3mm}
\qrcode[height=12mm]{https://helpman.komtera.lt/?fen=8/8/8/7k/n2R4/8/2b5/K2n4&moves=6}\\{\sffamily\tiny\color{ink}scannen zum Lösen}
\end{minipage}
\vspace{2mm}

{\raggedright\noindent\textcolor{ink}{\sffamily\textsc{Themen}}\enspace \themetag{reines Matt} \themetag{Mustermatt} \themetag{Spiegelmatt} \themetag{Switchback}\par}

\clearpage
\phantomsection\label{hm:KRvknnp}
\begin{tcolorbox}[enhanced,colback=ink,colframe=ink,arc=0pt,boxrule=0pt,left=4mm,right=4mm,top=2.5mm,bottom=2.5mm]
\color{white}\sffamily{\LARGE\bfseries KRvknnp}\hfill{\Large h\#6}\\[0.5mm]
{\small 6 Steine\quad Schwarz am Zug}\hfill{\small Nr.~754\qquad\href{https://helpman.komtera.lt/?fen=8/8/8/6pk/6R1/8/n7/K1n5&moves=6}{\color{white}\bfseries lösen\,\ensuremath{\nearrow}}}
\end{tcolorbox}
\vspace{2mm}
\noindent\begin{minipage}[t]{36mm}\vspace{0pt}\sffamily\small\raggedright
\hmfigure{längste eindeutige}{{\Large h\#6}}\hmfigure{eindeutig bei dieser Länge}{{\Large 23\,475} Stellungen}\hmfigure{längstes Matt}{{\Large h\#7.5}}\hmfigure{in der Klasse}{45,4 Milliarden Stellungen,\\ 14,7 Milliarden mit Hilfsmatt}
\end{minipage}\hfill
\begin{minipage}[t]{80mm}\vspace{0pt}\begin{tikzpicture}\begin{axis}[width=72mm,height=42mm,ybar,bar width=1.6pt,xmin=-1,xmax=16,ymin=0,ymax=3.151,axis y line*=left,axis x line*=bottom,xlabel={\scriptsize Halbzüge bis Matt},ylabel={\scriptsize Stellungen (Mrd.)},tick label style={font=\tiny},label style={font=\scriptsize},enlarge x limits=false,ylabel shift=-3pt,xlabel shift=-3pt]\addplot[fill=ink!55,draw=none] coordinates {(0,0.004189) (1,0.02661) (2,0.07374) (3,0.3374) (4,0.6865) (5,1.526) (6,2.157) (7,2.643) (8,2.865) (9,2.259) (10,1.447) (11,0.4877) (12,0.1931) (13,0.006622) (14,0.001621) (15,1.01e-07)};\draw[densely dotted] (axis cs:12,0) -- (axis cs:12,3.151);\end{axis}\begin{axis}[width=72mm,height=42mm,axis y line*=right,axis x line=none,ymode=log,xmin=-1,xmax=16,ylabel={\scriptsize eindeutig (log)},tick label style={font=\tiny},label style={font=\scriptsize},ylabel shift=-3pt]\addplot[mark=*,mark size=0.7pt,draw=black,thin] coordinates {(0,4188766) (1,26609027) (2,35856807) (3,117918868) (4,177961936) (5,145248425) (6,87508133) (7,33098370) (8,20827792) (9,6079261) (10,2162447) (11,53534) (12,23475)};\end{axis}\end{tikzpicture}\end{minipage}
\vspace{3mm}

\noindent\begin{minipage}[t]{60mm}\vspace{0pt}\sbox\hmboard{\chessboard[setfen=8/8/8/6pk/6R1/8/n7/K1n5 b - - 0 1]}\usebox\hmboard\\[1mm]
\makebox[\wd\hmboard]{\sffamily\hspace{3mm}\textbf{h\#6}\hfill(2+4)}\end{minipage}\hfill
\begin{minipage}[t]{52mm}\vspace{0pt}
\textcolor{ink}{\sffamily\textsc{Lösung}}\par\smallskip\small
\begin{tabular}{@{}r@{\enspace}l@{\enspace}l@{}}
\textbf{1.} & \symknight{}d3 & \symrook{}xg5+\\
\textbf{2.} & \symking{}h4 & \symrook{}h5+\\
\textbf{3.} & \symking{}g3 & \symking{}b1\\
\textbf{4.} & \symking{}f2 & \symking{}c2\\
\textbf{5.} & \symking{}e1 & \symking{}xd3\\
\textbf{6.} & \symking{}d1 & \symrook{}h1\#\\
\end{tabular}\par\vspace{3mm}
\qrcode[height=12mm]{https://helpman.komtera.lt/?fen=8/8/8/6pk/6R1/8/n7/K1n5&moves=6}\\{\sffamily\tiny\color{ink}scannen zum Lösen}
\end{minipage}
\vspace{2mm}

{\raggedright\noindent\textcolor{ink}{\sffamily\textsc{Themen}}\enspace \themetag{reines Matt} \themetag{Mustermatt} \themetag{Spiegelmatt} \themetag{Umnov}\par}

\clearpage
\phantomsection\label{hm:KRvkppp}
\begin{tcolorbox}[enhanced,colback=ink,colframe=ink,arc=0pt,boxrule=0pt,left=4mm,right=4mm,top=2.5mm,bottom=2.5mm]
\color{white}\sffamily{\LARGE\bfseries KRvkppp}\hfill{\Large h\#6}\\[0.5mm]
{\small 6 Steine\quad Schwarz am Zug}\hfill{\small Nr.~755\qquad\href{https://helpman.komtera.lt/?fen=8/8/8/8/p7/p6k/p7/K5R1&moves=6}{\color{white}\bfseries lösen\,\ensuremath{\nearrow}}}
\end{tcolorbox}
\vspace{2mm}
\noindent\begin{minipage}[t]{36mm}\vspace{0pt}\sffamily\small\raggedright
\hmfigure{längste eindeutige}{{\Large h\#6}}\hmfigure{eindeutig bei dieser Länge}{{\Large 8\,358} Stellungen}\hmfigure{längstes Matt}{{\Large h\#7.5}}\hmfigure{in der Klasse}{25,6 Milliarden Stellungen,\\ 2,88 Milliarden mit Hilfsmatt}
\end{minipage}\hfill
\begin{minipage}[t]{80mm}\vspace{0pt}\begin{tikzpicture}\begin{axis}[width=72mm,height=42mm,ybar,bar width=1.6pt,xmin=-1,xmax=16,ymin=0,ymax=641.9,axis y line*=left,axis x line*=bottom,xlabel={\scriptsize Halbzüge bis Matt},ylabel={\scriptsize Stellungen (Mio.)},tick label style={font=\tiny},label style={font=\scriptsize},enlarge x limits=false,ylabel shift=-3pt,xlabel shift=-3pt]\addplot[fill=ink!55,draw=none] coordinates {(0,1.392) (1,7.937) (2,17.91) (3,74.56) (4,159.2) (5,303.8) (6,468.8) (7,489.2) (8,583.5) (9,396) (10,265.6) (11,75.69) (12,39.58) (13,1.204) (14,0.519) (15,2.8e-05)};\draw[densely dotted] (axis cs:12,0) -- (axis cs:12,641.9);\end{axis}\begin{axis}[width=72mm,height=42mm,axis y line*=right,axis x line=none,ymode=log,xmin=-1,xmax=16,ylabel={\scriptsize eindeutig (log)},tick label style={font=\tiny},label style={font=\scriptsize},ylabel shift=-3pt]\addplot[mark=*,mark size=0.7pt,draw=black,thin] coordinates {(0,1391859) (1,7936045) (2,9677135) (3,27405074) (4,47000116) (5,26792758) (6,16382288) (7,4474565) (8,2738509) (9,692880) (10,246367) (11,16661) (12,8358)};\end{axis}\end{tikzpicture}\end{minipage}
\vspace{3mm}

\noindent\begin{minipage}[t]{60mm}\vspace{0pt}\sbox\hmboard{\chessboard[setfen=8/8/8/8/p7/p6k/p7/K5R1 b - - 0 1]}\usebox\hmboard\\[1mm]
\makebox[\wd\hmboard]{\sffamily\hspace{3mm}\textbf{h\#6}\hfill(2+4)}\end{minipage}\hfill
\begin{minipage}[t]{52mm}\vspace{0pt}
\textcolor{ink}{\sffamily\textsc{Lösung}}\par\smallskip\small
\begin{tabular}{@{}r@{\enspace}l@{\enspace}l@{}}
\textbf{1.} & \symking{}h2 & \symrook{}g2+\\
\textbf{2.} & \symking{}h1 & \symrook{}xa2\\
\textbf{3.} & \symking{}g1 & \symking{}b1\\
\textbf{4.} & \symking{}f1 & \symking{}c2\\
\textbf{5.} & \symking{}e1 & \symking{}d3\\
\textbf{6.} & \symking{}d1 & \symrook{}a1\#\\
\end{tabular}\par\vspace{3mm}
\qrcode[height=12mm]{https://helpman.komtera.lt/?fen=8/8/8/8/p7/p6k/p7/K5R1&moves=6}\\{\sffamily\tiny\color{ink}scannen zum Lösen}
\end{minipage}
\vspace{2mm}

{\raggedright\noindent\textcolor{ink}{\sffamily\textsc{Themen}}\enspace \themetag{reines Matt} \themetag{Mustermatt} \themetag{Spiegelmatt} \themetag{ein Stein zieht}\par}

\clearpage
\phantomsection\label{hm:KRvkrnn}
\begin{tcolorbox}[enhanced,colback=ink,colframe=ink,arc=0pt,boxrule=0pt,left=4mm,right=4mm,top=2.5mm,bottom=2.5mm]
\color{white}\sffamily{\LARGE\bfseries KRvkrnn}\hfill{\Large h\#6}\\[0.5mm]
{\small 6 Steine\quad Schwarz am Zug}\hfill{\small Nr.~756\qquad\href{https://helpman.komtera.lt/?fen=8/8/8/7k/nr2R3/8/n7/K7&moves=6}{\color{white}\bfseries lösen\,\ensuremath{\nearrow}}}
\end{tcolorbox}
\vspace{2mm}
\noindent\begin{minipage}[t]{36mm}\vspace{0pt}\sffamily\small\raggedright
\hmfigure{längste eindeutige}{{\Large h\#6}}\hmfigure{eindeutig bei dieser Länge}{{\Large 778} Stellungen}\hmfigure{längstes Matt}{{\Large h\#7.5}}\hmfigure{in der Klasse}{15,5 Milliarden Stellungen,\\ 4,47 Milliarden mit Hilfsmatt}
\end{minipage}\hfill
\begin{minipage}[t]{80mm}\vspace{0pt}\begin{tikzpicture}\begin{axis}[width=72mm,height=42mm,ybar,bar width=1.6pt,xmin=-1,xmax=16,ymin=0,ymax=981.3,axis y line*=left,axis x line*=bottom,xlabel={\scriptsize Halbzüge bis Matt},ylabel={\scriptsize Stellungen (Mio.)},tick label style={font=\tiny},label style={font=\scriptsize},enlarge x limits=false,ylabel shift=-3pt,xlabel shift=-3pt]\addplot[fill=ink!55,draw=none] coordinates {(0,0.6297) (1,4.464) (2,17.13) (3,90.1) (4,181) (5,461.9) (6,608.2) (7,892.1) (8,814.5) (9,787.8) (10,388.8) (11,183.7) (12,34.4) (13,2.233) (14,0.08164) (15,0.001052)};\draw[densely dotted] (axis cs:12,0) -- (axis cs:12,981.3);\end{axis}\begin{axis}[width=72mm,height=42mm,axis y line*=right,axis x line=none,ymode=log,xmin=-1,xmax=16,ylabel={\scriptsize eindeutig (log)},tick label style={font=\tiny},label style={font=\scriptsize},ylabel shift=-3pt]\addplot[mark=*,mark size=0.7pt,draw=black,thin] coordinates {(0,629683) (1,4464141) (2,7929507) (3,29536453) (4,39749267) (5,41838146) (6,24329424) (7,14773524) (8,7020979) (9,2926514) (10,663503) (11,36811) (12,778)};\end{axis}\end{tikzpicture}\end{minipage}
\vspace{3mm}

\noindent\begin{minipage}[t]{60mm}\vspace{0pt}\sbox\hmboard{\chessboard[setfen=8/8/8/7k/nr2R3/8/n7/K7 b - - 0 1]}\usebox\hmboard\\[1mm]
\makebox[\wd\hmboard]{\sffamily\hspace{3mm}\textbf{h\#6}\hfill(2+4)}\end{minipage}\hfill
\begin{minipage}[t]{52mm}\vspace{0pt}
\textcolor{ink}{\sffamily\textsc{Lösung}}\par\smallskip\small
\begin{tabular}{@{}r@{\enspace}l@{\enspace}l@{}}
\textbf{1.} & \symrook{}d4 & \symrook{}e2\\
\textbf{2.} & \symking{}g4 & \symrook{}xa2\\
\textbf{3.} & \symking{}f3 & \symrook{}xa4\\
\textbf{4.} & \symking{}e2 & \symking{}b2\\
\textbf{5.} & \symking{}d1 & \symking{}c3\\
\textbf{6.} & \symking{}c1 & \symrook{}a1\#\\
\end{tabular}\par\vspace{3mm}
\qrcode[height=12mm]{https://helpman.komtera.lt/?fen=8/8/8/7k/nr2R3/8/n7/K7&moves=6}\\{\sffamily\tiny\color{ink}scannen zum Lösen}
\end{minipage}
\vspace{2mm}

{\raggedright\noindent\textcolor{ink}{\sffamily\textsc{Themen}}\enspace \themetag{reines Matt} \themetag{Mustermatt} \themetag{Spiegelmatt} \themetag{ohne Schach}\par}

\clearpage
\phantomsection\label{hm:KNNNvkn}
\begin{tcolorbox}[enhanced,colback=ink,colframe=ink,arc=0pt,boxrule=0pt,left=4mm,right=4mm,top=2.5mm,bottom=2.5mm]
\color{white}\sffamily{\LARGE\bfseries KNNNvkn}\hfill{\Large h\#5.5}\\[0.5mm]
{\small 6 Steine\quad Weiß am Zug}\hfill{\small Nr.~757\qquad\href{https://helpman.komtera.lt/?fen=8/7k/8/8/8/8/1n6/KN1N3N&moves=5.5}{\color{white}\bfseries lösen\,\ensuremath{\nearrow}}}
\end{tcolorbox}
\vspace{2mm}
\noindent\begin{minipage}[t]{36mm}\vspace{0pt}\sffamily\small\raggedright
\hmfigure{längste eindeutige}{{\Large h\#5.5}}\hmfigure{eindeutig bei dieser Länge}{{\Large 147} Stellungen}\hmfigure{längstes Matt}{{\Large h\#7}}\hmfigure{in der Klasse}{15,5 Milliarden Stellungen,\\ 1,7 Milliarden mit Hilfsmatt}
\end{minipage}\hfill
\begin{minipage}[t]{80mm}\vspace{0pt}\begin{tikzpicture}\begin{axis}[width=72mm,height=42mm,ybar,bar width=1.6pt,xmin=-1,xmax=15,ymin=0,ymax=492.5,axis y line*=left,axis x line*=bottom,xlabel={\scriptsize Halbzüge bis Matt},ylabel={\scriptsize Stellungen (Mio.)},tick label style={font=\tiny},label style={font=\scriptsize},enlarge x limits=false,ylabel shift=-3pt,xlabel shift=-3pt]\addplot[fill=ink!55,draw=none] coordinates {(0,0.09342) (1,0.4295) (2,1.708) (3,11.61) (4,46.43) (5,146.1) (6,356) (7,409.5) (8,447.7) (9,197.2) (10,71.64) (11,12.86) (12,0.9049) (13,0.05035) (14,0.000172)};\draw[densely dotted] (axis cs:11,0) -- (axis cs:11,492.5);\end{axis}\begin{axis}[width=72mm,height=42mm,axis y line*=right,axis x line=none,ymode=log,xmin=-1,xmax=15,ylabel={\scriptsize eindeutig (log)},tick label style={font=\tiny},label style={font=\scriptsize},ylabel shift=-3pt]\addplot[mark=*,mark size=0.7pt,draw=black,thin] coordinates {(0,93424) (1,385264) (2,1214043) (3,5063045) (4,13700064) (5,9736213) (6,9931671) (7,2165697) (8,807282) (9,60436) (10,3728) (11,147)};\end{axis}\end{tikzpicture}\end{minipage}
\vspace{3mm}

\noindent\begin{minipage}[t]{60mm}\vspace{0pt}\sbox\hmboard{\chessboard[setfen=8/7k/8/8/8/8/1n6/KN1N3N w - - 0 1]}\usebox\hmboard\\[1mm]
\makebox[\wd\hmboard]{\sffamily\hspace{3mm}\textbf{h\#5.5}\hfill(4+2)}\end{minipage}\hfill
\begin{minipage}[t]{52mm}\vspace{0pt}
\textcolor{ink}{\sffamily\textsc{Lösung}}\par\smallskip\small
\begin{tabular}{@{}r@{\enspace}l@{\enspace}l@{}}
\textbf{1.} & \ldots & \symknight{}hf2\\
\textbf{2.} & \symknight{}d3 & \symknight{}g4\\
\textbf{3.} & \symknight{}f4 & \symking{}b2\\
\textbf{4.} & \symking{}g6 & \symking{}c3\\
\textbf{5.} & \symking{}f5 & \symknight{}de3+\\
\textbf{6.} & \symking{}e4 & \symknight{}d2\#\\
\end{tabular}\par\vspace{3mm}
\qrcode[height=12mm]{https://helpman.komtera.lt/?fen=8/7k/8/8/8/8/1n6/KN1N3N&moves=5.5}\\{\sffamily\tiny\color{ink}scannen zum Lösen}
\end{minipage}
\vspace{2mm}

{\raggedright\noindent\textcolor{ink}{\sffamily\textsc{Themen}}\enspace \themetag{reines Matt} \themetag{Mustermatt} \themetag{Idealmatt} \themetag{Selbstblock} \themetag{ohne Schlag}\par}

\clearpage
\phantomsection\label{hm:KQQvkqb}
\begin{tcolorbox}[enhanced,colback=ink,colframe=ink,arc=0pt,boxrule=0pt,left=4mm,right=4mm,top=2.5mm,bottom=2.5mm]
\color{white}\sffamily{\LARGE\bfseries KQQvkqb}\hfill{\Large h\#5.5}\\[0.5mm]
{\small 6 Steine\quad Weiß am Zug}\hfill{\small Nr.~758\qquad\href{https://helpman.komtera.lt/?fen=7b/6Q1/8/7k/8/6Q1/q7/K7&moves=5.5}{\color{white}\bfseries lösen\,\ensuremath{\nearrow}}}
\end{tcolorbox}
\vspace{2mm}
\noindent\begin{minipage}[t]{36mm}\vspace{0pt}\sffamily\small\raggedright
\hmfigure{längste eindeutige}{{\Large h\#5.5}}\hmfigure{eindeutig bei dieser Länge}{{\Large 1} Stellung}\hmfigure{längstes Matt}{{\Large h\#8}}\hmfigure{in der Klasse}{15,5 Milliarden Stellungen,\\ 3,1 Milliarden mit Hilfsmatt}
\end{minipage}\hfill
\begin{minipage}[t]{80mm}\vspace{0pt}\begin{tikzpicture}\begin{axis}[width=72mm,height=42mm,ybar,bar width=1.6pt,xmin=-1,xmax=17,ymin=0,ymax=1.112,axis y line*=left,axis x line*=bottom,xlabel={\scriptsize Halbzüge bis Matt},ylabel={\scriptsize Stellungen (Mrd.)},tick label style={font=\tiny},label style={font=\scriptsize},enlarge x limits=false,ylabel shift=-3pt,xlabel shift=-3pt]\addplot[fill=ink!55,draw=none] coordinates {(0,0.03788) (1,0.1233) (2,0.6502) (3,0.9595) (4,1.011) (5,0.249) (6,0.02941) (7,1.827e-05) (8,0.02763) (9,5.395e-06) (10,0.01031) (11,1.251e-06) (12,0.0004553) (13,1.79e-07) (14,5.66e-07) (15,2e-09) (16,8e-09)};\draw[densely dotted] (axis cs:11,0) -- (axis cs:11,1.112);\end{axis}\begin{axis}[width=72mm,height=42mm,axis y line*=right,axis x line=none,ymode=log,xmin=-1,xmax=17,ylabel={\scriptsize eindeutig (log)},tick label style={font=\tiny},label style={font=\scriptsize},ylabel shift=-3pt]\addplot[mark=*,mark size=0.7pt,draw=black,thin] coordinates {(0,37876100) (1,79690217) (2,239782148) (3,47653988) (4,2364258) (5,2259) (6,992095) (7,846) (8,424318) (9,11) (10,14400) (11,1)};\end{axis}\end{tikzpicture}\end{minipage}
\vspace{3mm}

\noindent\begin{minipage}[t]{60mm}\vspace{0pt}\sbox\hmboard{\chessboard[setfen=7b/6Q1/8/7k/8/6Q1/q7/K7 w - - 0 1]}\usebox\hmboard\\[1mm]
\makebox[\wd\hmboard]{\sffamily\hspace{3mm}\textbf{h\#5.5}\hfill(3+3)}\end{minipage}\hfill
\begin{minipage}[t]{52mm}\vspace{0pt}
\textcolor{ink}{\sffamily\textsc{Lösung}}\par\smallskip\small
\begin{tabular}{@{}r@{\enspace}l@{\enspace}l@{}}
\textbf{1.} & \ldots & \symking{}xa2\\
\textbf{2.} & \symbishop{}xg7 & \symking{}b3\\
\textbf{3.} & \symking{}h6 & \symking{}c4\\
\textbf{4.} & \symking{}h7 & \symking{}d5\\
\textbf{5.} & \symking{}g8 & \symking{}e6\\
\textbf{6.} & \symking{}f8 & \symqueen{}b8\#\\
\end{tabular}\par\vspace{3mm}
\qrcode[height=12mm]{https://helpman.komtera.lt/?fen=7b/6Q1/8/7k/8/6Q1/q7/K7&moves=5.5}\\{\sffamily\tiny\color{ink}scannen zum Lösen}
\end{minipage}
\vspace{2mm}

{\raggedright\noindent\textcolor{ink}{\sffamily\textsc{Themen}}\enspace \themetag{reines Matt} \themetag{Mustermatt} \themetag{Idealmatt} \themetag{Selbstblock} \themetag{ohne Schach}\par\smallskip{\footnotesize\textit{Hinweis:} die Partei am Zug steht im Diagramm im Schach; die Lösung beginnt mit einem Schlag.}\par}

\clearpage
\phantomsection\label{hm:KQvkbbb}
\begin{tcolorbox}[enhanced,colback=ink,colframe=ink,arc=0pt,boxrule=0pt,left=4mm,right=4mm,top=2.5mm,bottom=2.5mm]
\color{white}\sffamily{\LARGE\bfseries KQvkbbb}\hfill{\Large h\#5.5}\\[0.5mm]
{\small 6 Steine\quad Weiß am Zug}\hfill{\small Nr.~759\qquad\href{https://helpman.komtera.lt/?fen=8/7k/3b4/8/4b3/b7/Q7/K7&moves=5.5}{\color{white}\bfseries lösen\,\ensuremath{\nearrow}}}
\end{tcolorbox}
\vspace{2mm}
\noindent\begin{minipage}[t]{36mm}\vspace{0pt}\sffamily\small\raggedright
\hmfigure{längste eindeutige}{{\Large h\#5.5}}\hmfigure{eindeutig bei dieser Länge}{{\Large 66} Stellungen}\hmfigure{längstes Matt}{{\Large h\#7}}\hmfigure{in der Klasse}{15,5 Milliarden Stellungen,\\ 1,32 Milliarden mit Hilfsmatt}
\end{minipage}\hfill
\begin{minipage}[t]{80mm}\vspace{0pt}\begin{tikzpicture}\begin{axis}[width=72mm,height=42mm,ybar,bar width=1.6pt,xmin=-1,xmax=15,ymin=0,ymax=307.4,axis y line*=left,axis x line*=bottom,xlabel={\scriptsize Halbzüge bis Matt},ylabel={\scriptsize Stellungen (Mio.)},tick label style={font=\tiny},label style={font=\scriptsize},enlarge x limits=false,ylabel shift=-3pt,xlabel shift=-3pt]\addplot[fill=ink!55,draw=none] coordinates {(0,0.4625) (1,3.885) (2,14.39) (3,61.76) (4,136.9) (5,247.2) (6,279.5) (7,264.9) (8,181.9) (9,96.59) (10,31.12) (11,4.764) (12,0.7254) (13,0.002457) (14,0.00014)};\draw[densely dotted] (axis cs:11,0) -- (axis cs:11,307.4);\end{axis}\begin{axis}[width=72mm,height=42mm,axis y line*=right,axis x line=none,ymode=log,xmin=-1,xmax=15,ylabel={\scriptsize eindeutig (log)},tick label style={font=\tiny},label style={font=\scriptsize},ylabel shift=-3pt]\addplot[mark=*,mark size=0.7pt,draw=black,thin] coordinates {(0,462452) (1,3243822) (2,6972920) (3,13440287) (4,13469083) (5,11017158) (6,5587418) (7,3420547) (8,1126655) (9,164936) (10,17445) (11,66)};\end{axis}\end{tikzpicture}\end{minipage}
\vspace{3mm}

\noindent\begin{minipage}[t]{60mm}\vspace{0pt}\sbox\hmboard{\chessboard[setfen=8/7k/3b4/8/4b3/b7/Q7/K7 w - - 0 1]}\usebox\hmboard\\[1mm]
\makebox[\wd\hmboard]{\sffamily\hspace{3mm}\textbf{h\#5.5}\hfill(2+4)}\end{minipage}\hfill
\begin{minipage}[t]{52mm}\vspace{0pt}
\textcolor{ink}{\sffamily\textsc{Lösung}}\par\smallskip\small
\begin{tabular}{@{}r@{\enspace}l@{\enspace}l@{}}
\textbf{1.} & \ldots & \symqueen{}c2\\
\textbf{2.} & \symbishop{}b2+ & \symking{}b1\\
\textbf{3.} & \symbishop{}f6 & \symking{}c1\\
\textbf{4.} & \symking{}g6 & \symking{}d2\\
\textbf{5.} & \symking{}f5 & \symking{}e3\\
\textbf{6.} & \symking{}e5 & \symqueen{}xe4\#\\
\end{tabular}\par\vspace{3mm}
\qrcode[height=12mm]{https://helpman.komtera.lt/?fen=8/7k/3b4/8/4b3/b7/Q7/K7&moves=5.5}\\{\sffamily\tiny\color{ink}scannen zum Lösen}
\end{minipage}
\vspace{2mm}

{\raggedright\noindent\textcolor{ink}{\sffamily\textsc{Themen}}\enspace \themetag{Selbstblock}\par}

\clearpage
\phantomsection\label{hm:KQvkqqr}
\begin{tcolorbox}[enhanced,colback=ink,colframe=ink,arc=0pt,boxrule=0pt,left=4mm,right=4mm,top=2.5mm,bottom=2.5mm]
\color{white}\sffamily{\LARGE\bfseries KQvkqqr}\hfill{\Large h\#5.5}\\[0.5mm]
{\small 6 Steine\quad Weiß am Zug}\hfill{\small Nr.~760\qquad\href{https://helpman.komtera.lt/?fen=8/7k/8/8/8/1q6/8/KrQq4&moves=5.5}{\color{white}\bfseries lösen\,\ensuremath{\nearrow}}}
\end{tcolorbox}
\vspace{2mm}
\noindent\begin{minipage}[t]{36mm}\vspace{0pt}\sffamily\small\raggedright
\hmfigure{längste eindeutige}{{\Large h\#5.5}}\hmfigure{eindeutig bei dieser Länge}{{\Large 84} Stellungen}\hmfigure{längstes Matt}{{\Large h\#6.5}}\hmfigure{in der Klasse}{15,5 Milliarden Stellungen,\\ 2,94 Milliarden mit Hilfsmatt}
\end{minipage}\hfill
\begin{minipage}[t]{80mm}\vspace{0pt}\begin{tikzpicture}\begin{axis}[width=72mm,height=42mm,ybar,bar width=1.6pt,xmin=-1,xmax=14,ymin=0,ymax=1005,axis y line*=left,axis x line*=bottom,xlabel={\scriptsize Halbzüge bis Matt},ylabel={\scriptsize Stellungen (Mio.)},tick label style={font=\tiny},label style={font=\scriptsize},enlarge x limits=false,ylabel shift=-3pt,xlabel shift=-3pt]\addplot[fill=ink!55,draw=none] coordinates {(0,0.348) (1,3.671) (2,20.39) (3,120) (4,280.2) (5,786.5) (6,572.1) (7,914) (8,139.9) (9,101.1) (10,2.159) (11,1.138) (12,3.6e-05) (13,0.003418)};\draw[densely dotted] (axis cs:11,0) -- (axis cs:11,1005);\end{axis}\begin{axis}[width=72mm,height=42mm,axis y line*=right,axis x line=none,ymode=log,xmin=-1,xmax=14,ylabel={\scriptsize eindeutig (log)},tick label style={font=\tiny},label style={font=\scriptsize},ylabel shift=-3pt]\addplot[mark=*,mark size=0.7pt,draw=black,thin] coordinates {(0,347956) (1,3398092) (2,10733560) (3,31298615) (4,26157042) (5,36368509) (6,5661654) (7,4587758) (8,108564) (9,49691) (10,1) (11,84)};\end{axis}\end{tikzpicture}\end{minipage}
\vspace{3mm}

\noindent\begin{minipage}[t]{60mm}\vspace{0pt}\sbox\hmboard{\chessboard[setfen=8/7k/8/8/8/1q6/8/KrQq4 w - - 0 1]}\usebox\hmboard\\[1mm]
\makebox[\wd\hmboard]{\sffamily\hspace{3mm}\textbf{h\#5.5}\hfill(2+4)}\end{minipage}\hfill
\begin{minipage}[t]{52mm}\vspace{0pt}
\textcolor{ink}{\sffamily\textsc{Lösung}}\par\smallskip\small
\begin{tabular}{@{}r@{\enspace}l@{\enspace}l@{}}
\textbf{1.} & \ldots & \symqueen{}xb1+\\
\textbf{2.} & \symqueen{}dd3 & \symqueen{}e1\\
\textbf{3.} & \symqueen{}f7 & \symking{}b2\\
\textbf{4.} & \symqueen{}dg6 & \symking{}c3\\
\textbf{5.} & \symking{}g7 & \symking{}d4\\
\textbf{6.} & \symking{}f6 & \symqueen{}e5\#\\
\end{tabular}\par\vspace{3mm}
\qrcode[height=12mm]{https://helpman.komtera.lt/?fen=8/7k/8/8/8/1q6/8/KrQq4&moves=5.5}\\{\sffamily\tiny\color{ink}scannen zum Lösen}
\end{minipage}
\vspace{2mm}

{\raggedright\noindent\textcolor{ink}{\sffamily\textsc{Themen}}\enspace \themetag{reines Matt} \themetag{Mustermatt} \themetag{Idealmatt} \themetag{Selbstblock}\par\smallskip{\footnotesize\textit{Hinweis:} die Partei am Zug steht im Diagramm im Schach; die Lösung beginnt mit einem Schlag.}\par}

\clearpage
\phantomsection\label{hm:KQvkqrb}
\begin{tcolorbox}[enhanced,colback=ink,colframe=ink,arc=0pt,boxrule=0pt,left=4mm,right=4mm,top=2.5mm,bottom=2.5mm]
\color{white}\sffamily{\LARGE\bfseries KQvkqrb}\hfill{\Large h\#5.5}\\[0.5mm]
{\small 6 Steine\quad Weiß am Zug}\hfill{\small Nr.~761\qquad\href{https://helpman.komtera.lt/?fen=8/8/8/7k/1b6/1q6/8/K2rQ3&moves=5.5}{\color{white}\bfseries lösen\,\ensuremath{\nearrow}}}
\end{tcolorbox}
\vspace{2mm}
\noindent\begin{minipage}[t]{36mm}\vspace{0pt}\sffamily\small\raggedright
\hmfigure{längste eindeutige}{{\Large h\#5.5}}\hmfigure{eindeutig bei dieser Länge}{{\Large 359} Stellungen}\hmfigure{längstes Matt}{{\Large h\#6.5}}\hmfigure{in der Klasse}{15,5 Milliarden Stellungen,\\ 6,77 Milliarden mit Hilfsmatt}
\end{minipage}\hfill
\begin{minipage}[t]{80mm}\vspace{0pt}\begin{tikzpicture}\begin{axis}[width=72mm,height=42mm,ybar,bar width=1.6pt,xmin=-1,xmax=14,ymin=0,ymax=1.941,axis y line*=left,axis x line*=bottom,xlabel={\scriptsize Halbzüge bis Matt},ylabel={\scriptsize Stellungen (Mrd.)},tick label style={font=\tiny},label style={font=\scriptsize},enlarge x limits=false,ylabel shift=-3pt,xlabel shift=-3pt]\addplot[fill=ink!55,draw=none] coordinates {(0,0.001355) (1,0.01296) (2,0.06406) (3,0.3285) (4,0.8151) (5,1.765) (6,1.513) (7,1.729) (8,0.3586) (9,0.1794) (10,0.005846) (11,0.001414) (12,7e-08) (13,1.911e-05)};\draw[densely dotted] (axis cs:11,0) -- (axis cs:11,1.941);\end{axis}\begin{axis}[width=72mm,height=42mm,axis y line*=right,axis x line=none,ymode=log,xmin=-1,xmax=14,ylabel={\scriptsize eindeutig (log)},tick label style={font=\tiny},label style={font=\scriptsize},ylabel shift=-3pt]\addplot[mark=*,mark size=0.7pt,draw=black,thin] coordinates {(0,1354681) (1,11604132) (2,34465464) (3,77905916) (4,71017226) (5,65946430) (6,13820547) (7,7461320) (8,288025) (9,60548) (10,3) (11,359)};\end{axis}\end{tikzpicture}\end{minipage}
\vspace{3mm}

\noindent\begin{minipage}[t]{60mm}\vspace{0pt}\sbox\hmboard{\chessboard[setfen=8/8/8/7k/1b6/1q6/8/K2rQ3 w - - 0 1]}\usebox\hmboard\\[1mm]
\makebox[\wd\hmboard]{\sffamily\hspace{3mm}\textbf{h\#5.5}\hfill(2+4)}\end{minipage}\hfill
\begin{minipage}[t]{52mm}\vspace{0pt}
\textcolor{ink}{\sffamily\textsc{Lösung}}\par\smallskip\small
\begin{tabular}{@{}r@{\enspace}l@{\enspace}l@{}}
\textbf{1.} & \ldots & \symqueen{}xd1+\\
\textbf{2.} & \symking{}g5 & \symqueen{}xb3\\
\textbf{3.} & \symking{}f4 & \symqueen{}a2\\
\textbf{4.} & \symking{}e3 & \symking{}b2\\
\textbf{5.} & \symking{}d2 & \symking{}b3+\\
\textbf{6.} & \symking{}c1 & \symqueen{}c2\#\\
\end{tabular}\par\vspace{3mm}
\qrcode[height=12mm]{https://helpman.komtera.lt/?fen=8/8/8/7k/1b6/1q6/8/K2rQ3&moves=5.5}\\{\sffamily\tiny\color{ink}scannen zum Lösen}
\end{minipage}
\vspace{2mm}

{\raggedright\noindent\textcolor{ink}{\sffamily\textsc{Themen}}\enspace \themetag{ein Stein zieht}\par\smallskip{\footnotesize\textit{Hinweis:} die Partei am Zug steht im Diagramm im Schach; die Lösung beginnt mit einem Schlag.}\par}

\clearpage
\phantomsection\label{hm:KQvkqrn}
\begin{tcolorbox}[enhanced,colback=ink,colframe=ink,arc=0pt,boxrule=0pt,left=4mm,right=4mm,top=2.5mm,bottom=2.5mm]
\color{white}\sffamily{\LARGE\bfseries KQvkqrn}\hfill{\Large h\#5.5}\\[0.5mm]
{\small 6 Steine\quad Weiß am Zug}\hfill{\small Nr.~762\qquad\href{https://helpman.komtera.lt/?fen=8/8/6k1/8/n7/rq6/Q7/K7&moves=5.5}{\color{white}\bfseries lösen\,\ensuremath{\nearrow}}}
\end{tcolorbox}
\vspace{2mm}
\noindent\begin{minipage}[t]{36mm}\vspace{0pt}\sffamily\small\raggedright
\hmfigure{längste eindeutige}{{\Large h\#5.5}}\hmfigure{eindeutig bei dieser Länge}{{\Large 1\,825} Stellungen}\hmfigure{längstes Matt}{{\Large h\#7.5}}\hmfigure{in der Klasse}{15,5 Milliarden Stellungen,\\ 6,95 Milliarden mit Hilfsmatt}
\end{minipage}\hfill
\begin{minipage}[t]{80mm}\vspace{0pt}\begin{tikzpicture}\begin{axis}[width=72mm,height=42mm,ybar,bar width=1.6pt,xmin=-1,xmax=16,ymin=0,ymax=1.966,axis y line*=left,axis x line*=bottom,xlabel={\scriptsize Halbzüge bis Matt},ylabel={\scriptsize Stellungen (Mrd.)},tick label style={font=\tiny},label style={font=\scriptsize},enlarge x limits=false,ylabel shift=-3pt,xlabel shift=-3pt]\addplot[fill=ink!55,draw=none] coordinates {(0,0.001498) (1,0.01393) (2,0.06565) (3,0.3203) (4,0.7814) (5,1.619) (6,1.569) (7,1.787) (8,0.4953) (9,0.285) (10,0.0117) (11,0.003456) (12,1.354e-06) (13,0.0002051) (15,2.2e-08)};\draw[densely dotted] (axis cs:11,0) -- (axis cs:11,1.966);\end{axis}\begin{axis}[width=72mm,height=42mm,axis y line*=right,axis x line=none,ymode=log,xmin=-1,xmax=16,ylabel={\scriptsize eindeutig (log)},tick label style={font=\tiny},label style={font=\scriptsize},ylabel shift=-3pt]\addplot[mark=*,mark size=0.7pt,draw=black,thin] coordinates {(0,1498199) (1,12259349) (2,34713242) (3,72810236) (4,66752901) (5,62652117) (6,17637811) (7,10172402) (8,512393) (9,156789) (10,88) (11,1825)};\end{axis}\end{tikzpicture}\end{minipage}
\vspace{3mm}

\noindent\begin{minipage}[t]{60mm}\vspace{0pt}\sbox\hmboard{\chessboard[setfen=8/8/6k1/8/n7/rq6/Q7/K7 w - - 0 1]}\usebox\hmboard\\[1mm]
\makebox[\wd\hmboard]{\sffamily\hspace{3mm}\textbf{h\#5.5}\hfill(2+4)}\end{minipage}\hfill
\begin{minipage}[t]{52mm}\vspace{0pt}
\textcolor{ink}{\sffamily\textsc{Lösung}}\par\smallskip\small
\begin{tabular}{@{}r@{\enspace}l@{\enspace}l@{}}
\textbf{1.} & \ldots & \symqueen{}xa3\\
\textbf{2.} & \symking{}f5 & \symqueen{}xb3\\
\textbf{3.} & \symking{}e4 & \symqueen{}e6+\\
\textbf{4.} & \symking{}d3 & \symking{}a2\\
\textbf{5.} & \symking{}d2 & \symking{}b3\\
\textbf{6.} & \symking{}c1 & \symqueen{}e1\#\\
\end{tabular}\par\vspace{3mm}
\qrcode[height=12mm]{https://helpman.komtera.lt/?fen=8/8/6k1/8/n7/rq6/Q7/K7&moves=5.5}\\{\sffamily\tiny\color{ink}scannen zum Lösen}
\end{minipage}
\vspace{2mm}

{\raggedright\noindent\textcolor{ink}{\sffamily\textsc{Themen}}\enspace \themetag{reines Matt} \themetag{Mustermatt} \themetag{Spiegelmatt} \themetag{ein Stein zieht}\par\smallskip{\footnotesize\textit{Hinweis:} die Lösung beginnt mit einem Schlag.}\par}

\clearpage
\phantomsection\label{hm:KQvkqrp}
\begin{tcolorbox}[enhanced,colback=ink,colframe=ink,arc=0pt,boxrule=0pt,left=4mm,right=4mm,top=2.5mm,bottom=2.5mm]
\color{white}\sffamily{\LARGE\bfseries KQvkqrp}\hfill{\Large h\#5.5}\\[0.5mm]
{\small 6 Steine\quad Weiß am Zug}\hfill{\small Nr.~763\qquad\href{https://helpman.komtera.lt/?fen=8/8/7k/8/8/3p4/2q5/KQr5&moves=5.5}{\color{white}\bfseries lösen\,\ensuremath{\nearrow}}}
\end{tcolorbox}
\vspace{2mm}
\noindent\begin{minipage}[t]{36mm}\vspace{0pt}\sffamily\small\raggedright
\hmfigure{längste eindeutige}{{\Large h\#5.5}}\hmfigure{eindeutig bei dieser Länge}{{\Large 1\,434} Stellungen}\hmfigure{längstes Matt}{{\Large h\#7.5}}\hmfigure{in der Klasse}{45,4 Milliarden Stellungen,\\ 21,4 Milliarden mit Hilfsmatt}
\end{minipage}\hfill
\begin{minipage}[t]{80mm}\vspace{0pt}\begin{tikzpicture}\begin{axis}[width=72mm,height=42mm,ybar,bar width=1.6pt,xmin=-1,xmax=16,ymin=0,ymax=6.135,axis y line*=left,axis x line*=bottom,xlabel={\scriptsize Halbzüge bis Matt},ylabel={\scriptsize Stellungen (Mrd.)},tick label style={font=\tiny},label style={font=\scriptsize},enlarge x limits=false,ylabel shift=-3pt,xlabel shift=-3pt]\addplot[fill=ink!55,draw=none] coordinates {(0,0.00586) (1,0.0533) (2,0.2514) (3,1.194) (4,2.983) (5,5.577) (6,4.921) (7,4.82) (8,1.111) (9,0.5093) (10,0.01889) (11,0.004126) (12,1.5e-08) (13,0.0001094) (15,4e-09)};\draw[densely dotted] (axis cs:11,0) -- (axis cs:11,6.135);\end{axis}\begin{axis}[width=72mm,height=42mm,axis y line*=right,axis x line=none,ymode=log,xmin=-1,xmax=16,ylabel={\scriptsize eindeutig (log)},tick label style={font=\tiny},label style={font=\scriptsize},ylabel shift=-3pt]\addplot[mark=*,mark size=0.7pt,draw=black,thin] coordinates {(0,5860208) (1,46586219) (2,137215472) (3,265404447) (4,244750928) (5,185020371) (6,41431560) (7,19047090) (8,845818) (9,196595) (11,1434)};\end{axis}\end{tikzpicture}\end{minipage}
\vspace{3mm}

\noindent\begin{minipage}[t]{60mm}\vspace{0pt}\sbox\hmboard{\chessboard[setfen=8/8/7k/8/8/3p4/2q5/KQr5 w - - 0 1]}\usebox\hmboard\\[1mm]
\makebox[\wd\hmboard]{\sffamily\hspace{3mm}\textbf{h\#5.5}\hfill(2+4)}\end{minipage}\hfill
\begin{minipage}[t]{52mm}\vspace{0pt}
\textcolor{ink}{\sffamily\textsc{Lösung}}\par\smallskip\small
\begin{tabular}{@{}r@{\enspace}l@{\enspace}l@{}}
\textbf{1.} & \ldots & \symqueen{}xc1+\\
\textbf{2.} & \symking{}h5 & \symqueen{}xc2\\
\textbf{3.} & \symking{}g4 & \symqueen{}a2\\
\textbf{4.} & \symking{}f3 & \symking{}b2\\
\textbf{5.} & \symking{}e2 & \symking{}c3+\\
\textbf{6.} & \symking{}d1 & \symqueen{}d2\#\\
\end{tabular}\par\vspace{3mm}
\qrcode[height=12mm]{https://helpman.komtera.lt/?fen=8/8/7k/8/8/3p4/2q5/KQr5&moves=5.5}\\{\sffamily\tiny\color{ink}scannen zum Lösen}
\end{minipage}
\vspace{2mm}

{\raggedright\noindent\textcolor{ink}{\sffamily\textsc{Themen}}\enspace \themetag{ein Stein zieht} \themetag{Inder}\par\smallskip{\footnotesize\textit{Hinweis:} die Lösung beginnt mit einem Schlag.}\par}

\clearpage
\phantomsection\label{hm:KQvkrbb}
\begin{tcolorbox}[enhanced,colback=ink,colframe=ink,arc=0pt,boxrule=0pt,left=4mm,right=4mm,top=2.5mm,bottom=2.5mm]
\color{white}\sffamily{\LARGE\bfseries KQvkrbb}\hfill{\Large h\#5.5}\\[0.5mm]
{\small 6 Steine\quad Weiß am Zug}\hfill{\small Nr.~764\qquad\href{https://helpman.komtera.lt/?fen=8/8/5k2/8/6Q1/bb6/8/Kr6&moves=5.5}{\color{white}\bfseries lösen\,\ensuremath{\nearrow}}}
\end{tcolorbox}
\vspace{2mm}
\noindent\begin{minipage}[t]{36mm}\vspace{0pt}\sffamily\small\raggedright
\hmfigure{längste eindeutige}{{\Large h\#5.5}}\hmfigure{eindeutig bei dieser Länge}{{\Large 93} Stellungen}\hmfigure{längstes Matt}{{\Large h\#6.5}}\hmfigure{in der Klasse}{15,5 Milliarden Stellungen,\\ 3,85 Milliarden mit Hilfsmatt}
\end{minipage}\hfill
\begin{minipage}[t]{80mm}\vspace{0pt}\begin{tikzpicture}\begin{axis}[width=72mm,height=42mm,ybar,bar width=1.6pt,xmin=-1,xmax=14,ymin=0,ymax=1054,axis y line*=left,axis x line*=bottom,xlabel={\scriptsize Halbzüge bis Matt},ylabel={\scriptsize Stellungen (Mio.)},tick label style={font=\tiny},label style={font=\scriptsize},enlarge x limits=false,ylabel shift=-3pt,xlabel shift=-3pt]\addplot[fill=ink!55,draw=none] coordinates {(0,1.249) (1,10.89) (2,47.57) (3,211.3) (4,562.3) (5,929.7) (6,958.2) (7,795.3) (8,244.1) (9,87.52) (10,4.657) (11,0.218) (12,0.000789) (13,0.007827)};\draw[densely dotted] (axis cs:11,0) -- (axis cs:11,1054);\end{axis}\begin{axis}[width=72mm,height=42mm,axis y line*=right,axis x line=none,ymode=log,xmin=-1,xmax=14,ylabel={\scriptsize eindeutig (log)},tick label style={font=\tiny},label style={font=\scriptsize},ylabel shift=-3pt]\addplot[mark=*,mark size=0.7pt,draw=black,thin] coordinates {(0,1248543) (1,9313908) (2,25824841) (3,43861176) (4,45905419) (5,30037876) (6,9168581) (7,3651207) (8,254200) (9,9821) (10,98) (11,93)};\end{axis}\end{tikzpicture}\end{minipage}
\vspace{3mm}

\noindent\begin{minipage}[t]{60mm}\vspace{0pt}\sbox\hmboard{\chessboard[setfen=8/8/5k2/8/6Q1/bb6/8/Kr6 w - - 0 1]}\usebox\hmboard\\[1mm]
\makebox[\wd\hmboard]{\sffamily\hspace{3mm}\textbf{h\#5.5}\hfill(2+4)}\end{minipage}\hfill
\begin{minipage}[t]{52mm}\vspace{0pt}
\textcolor{ink}{\sffamily\textsc{Lösung}}\par\smallskip\small
\begin{tabular}{@{}r@{\enspace}l@{\enspace}l@{}}
\textbf{1.} & \ldots & \symking{}xb1\\
\textbf{2.} & \symbishop{}f7 & \symking{}c2\\
\textbf{3.} & \symbishop{}g6+ & \symking{}d2\\
\textbf{4.} & \symking{}g7 & \symking{}e3\\
\textbf{5.} & \symking{}h6 & \symking{}f4\\
\textbf{6.} & \symbishop{}h7 & \symqueen{}g5\#\\
\end{tabular}\par\vspace{3mm}
\qrcode[height=12mm]{https://helpman.komtera.lt/?fen=8/8/5k2/8/6Q1/bb6/8/Kr6&moves=5.5}\\{\sffamily\tiny\color{ink}scannen zum Lösen}
\end{minipage}
\vspace{2mm}

{\raggedright\noindent\textcolor{ink}{\sffamily\textsc{Themen}}\enspace \themetag{reines Matt} \themetag{Mustermatt} \themetag{Selbstblock}\par\smallskip{\footnotesize\textit{Hinweis:} die Partei am Zug steht im Diagramm im Schach; die Lösung beginnt mit einem Schlag.}\par}

\clearpage
\phantomsection\label{hm:KQvkrbn}
\begin{tcolorbox}[enhanced,colback=ink,colframe=ink,arc=0pt,boxrule=0pt,left=4mm,right=4mm,top=2.5mm,bottom=2.5mm]
\color{white}\sffamily{\LARGE\bfseries KQvkrbn}\hfill{\Large h\#5.5}\\[0.5mm]
{\small 6 Steine\quad Weiß am Zug}\hfill{\small Nr.~765\qquad\href{https://helpman.komtera.lt/?fen=8/8/7k/8/1n6/b7/8/KQr5&moves=5.5}{\color{white}\bfseries lösen\,\ensuremath{\nearrow}}}
\end{tcolorbox}
\vspace{2mm}
\noindent\begin{minipage}[t]{36mm}\vspace{0pt}\sffamily\small\raggedright
\hmfigure{längste eindeutige}{{\Large h\#5.5}}\hmfigure{eindeutig bei dieser Länge}{{\Large 304} Stellungen}\hmfigure{längstes Matt}{{\Large h\#7.5}}\hmfigure{in der Klasse}{15,5 Milliarden Stellungen,\\ 7,93 Milliarden mit Hilfsmatt}
\end{minipage}\hfill
\begin{minipage}[t]{80mm}\vspace{0pt}\begin{tikzpicture}\begin{axis}[width=72mm,height=42mm,ybar,bar width=1.6pt,xmin=-1,xmax=16,ymin=0,ymax=2.159,axis y line*=left,axis x line*=bottom,xlabel={\scriptsize Halbzüge bis Matt},ylabel={\scriptsize Stellungen (Mrd.)},tick label style={font=\tiny},label style={font=\scriptsize},enlarge x limits=false,ylabel shift=-3pt,xlabel shift=-3pt]\addplot[fill=ink!55,draw=none] coordinates {(0,0.002815) (1,0.0236) (2,0.09651) (3,0.4051) (4,1.056) (5,1.692) (6,1.963) (7,1.66) (8,0.712) (9,0.2908) (10,0.02315) (11,0.001238) (12,2.832e-05) (13,2.949e-05) (15,1.7e-08)};\draw[densely dotted] (axis cs:11,0) -- (axis cs:11,2.159);\end{axis}\begin{axis}[width=72mm,height=42mm,axis y line*=right,axis x line=none,ymode=log,xmin=-1,xmax=16,ylabel={\scriptsize eindeutig (log)},tick label style={font=\tiny},label style={font=\scriptsize},ylabel shift=-3pt]\addplot[mark=*,mark size=0.7pt,draw=black,thin] coordinates {(0,2814908) (1,19666504) (2,50804916) (3,80755146) (4,88205625) (5,60652924) (6,24483787) (7,10667581) (8,1244235) (9,78810) (10,2726) (11,304)};\end{axis}\end{tikzpicture}\end{minipage}
\vspace{3mm}

\noindent\begin{minipage}[t]{60mm}\vspace{0pt}\sbox\hmboard{\chessboard[setfen=8/8/7k/8/1n6/b7/8/KQr5 w - - 0 1]}\usebox\hmboard\\[1mm]
\makebox[\wd\hmboard]{\sffamily\hspace{3mm}\textbf{h\#5.5}\hfill(2+4)}\end{minipage}\hfill
\begin{minipage}[t]{52mm}\vspace{0pt}
\textcolor{ink}{\sffamily\textsc{Lösung}}\par\smallskip\small
\begin{tabular}{@{}r@{\enspace}l@{\enspace}l@{}}
\textbf{1.} & \ldots & \symqueen{}xc1+\\
\textbf{2.} & \symking{}h5 & \symqueen{}xa3\\
\textbf{3.} & \symking{}g4 & \symqueen{}a2\\
\textbf{4.} & \symking{}f3 & \symking{}b2\\
\textbf{5.} & \symking{}e2 & \symking{}c3+\\
\textbf{6.} & \symking{}d1 & \symqueen{}d2\#\\
\end{tabular}\par\vspace{3mm}
\qrcode[height=12mm]{https://helpman.komtera.lt/?fen=8/8/7k/8/1n6/b7/8/KQr5&moves=5.5}\\{\sffamily\tiny\color{ink}scannen zum Lösen}
\end{minipage}
\vspace{2mm}

{\raggedright\noindent\textcolor{ink}{\sffamily\textsc{Themen}}\enspace \themetag{ein Stein zieht}\par\smallskip{\footnotesize\textit{Hinweis:} die Lösung beginnt mit einem Schlag.}\par}

\clearpage
\phantomsection\label{hm:KQvkrbp}
\begin{tcolorbox}[enhanced,colback=ink,colframe=ink,arc=0pt,boxrule=0pt,left=4mm,right=4mm,top=2.5mm,bottom=2.5mm]
\color{white}\sffamily{\LARGE\bfseries KQvkrbp}\hfill{\Large h\#5.5}\\[0.5mm]
{\small 6 Steine\quad Weiß am Zug}\hfill{\small Nr.~766\qquad\href{https://helpman.komtera.lt/?fen=8/8/7k/8/8/r1p5/Q1b5/K7&moves=5.5}{\color{white}\bfseries lösen\,\ensuremath{\nearrow}}}
\end{tcolorbox}
\vspace{2mm}
\noindent\begin{minipage}[t]{36mm}\vspace{0pt}\sffamily\small\raggedright
\hmfigure{längste eindeutige}{{\Large h\#5.5}}\hmfigure{eindeutig bei dieser Länge}{{\Large 369} Stellungen}\hmfigure{längstes Matt}{{\Large h\#6.5}}\hmfigure{in der Klasse}{45,4 Milliarden Stellungen,\\ 24,5 Milliarden mit Hilfsmatt}
\end{minipage}\hfill
\begin{minipage}[t]{80mm}\vspace{0pt}\begin{tikzpicture}\begin{axis}[width=72mm,height=42mm,ybar,bar width=1.6pt,xmin=-1,xmax=14,ymin=0,ymax=6.806,axis y line*=left,axis x line*=bottom,xlabel={\scriptsize Halbzüge bis Matt},ylabel={\scriptsize Stellungen (Mrd.)},tick label style={font=\tiny},label style={font=\scriptsize},enlarge x limits=false,ylabel shift=-3pt,xlabel shift=-3pt]\addplot[fill=ink!55,draw=none] coordinates {(0,0.01095) (1,0.08992) (2,0.368) (3,1.492) (4,4.037) (5,5.721) (6,6.187) (7,4.431) (8,1.6) (9,0.5311) (10,0.0359) (11,0.001089) (12,2.497e-05) (13,1.097e-05)};\draw[densely dotted] (axis cs:11,0) -- (axis cs:11,6.806);\end{axis}\begin{axis}[width=72mm,height=42mm,axis y line*=right,axis x line=none,ymode=log,xmin=-1,xmax=14,ylabel={\scriptsize eindeutig (log)},tick label style={font=\tiny},label style={font=\scriptsize},ylabel shift=-3pt]\addplot[mark=*,mark size=0.7pt,draw=black,thin] coordinates {(0,10946675) (1,74314168) (2,201939190) (3,285559324) (4,318647323) (5,170809944) (6,58249863) (7,20891488) (8,1991987) (9,75611) (10,2207) (11,369)};\end{axis}\end{tikzpicture}\end{minipage}
\vspace{3mm}

\noindent\begin{minipage}[t]{60mm}\vspace{0pt}\sbox\hmboard{\chessboard[setfen=8/8/7k/8/8/r1p5/Q1b5/K7 w - - 0 1]}\usebox\hmboard\\[1mm]
\makebox[\wd\hmboard]{\sffamily\hspace{3mm}\textbf{h\#5.5}\hfill(2+4)}\end{minipage}\hfill
\begin{minipage}[t]{52mm}\vspace{0pt}
\textcolor{ink}{\sffamily\textsc{Lösung}}\par\smallskip\small
\begin{tabular}{@{}r@{\enspace}l@{\enspace}l@{}}
\textbf{1.} & \ldots & \symqueen{}xa3\\
\textbf{2.} & \symking{}g5 & \symqueen{}xc3\\
\textbf{3.} & \symking{}f4 & \symqueen{}f6+\\
\textbf{4.} & \symking{}e3 & \symking{}b2\\
\textbf{5.} & \symking{}e2 & \symking{}c3\\
\textbf{6.} & \symking{}d1 & \symqueen{}f1\#\\
\end{tabular}\par\vspace{3mm}
\qrcode[height=12mm]{https://helpman.komtera.lt/?fen=8/8/7k/8/8/r1p5/Q1b5/K7&moves=5.5}\\{\sffamily\tiny\color{ink}scannen zum Lösen}
\end{minipage}
\vspace{2mm}

{\raggedright\noindent\textcolor{ink}{\sffamily\textsc{Themen}}\enspace \themetag{ein Stein zieht}\par\smallskip{\footnotesize\textit{Hinweis:} die Lösung beginnt mit einem Schlag.}\par}

\clearpage
\phantomsection\label{hm:KQvkrnn}
\begin{tcolorbox}[enhanced,colback=ink,colframe=ink,arc=0pt,boxrule=0pt,left=4mm,right=4mm,top=2.5mm,bottom=2.5mm]
\color{white}\sffamily{\LARGE\bfseries KQvkrnn}\hfill{\Large h\#5.5}\\[0.5mm]
{\small 6 Steine\quad Weiß am Zug}\hfill{\small Nr.~767\qquad\href{https://helpman.komtera.lt/?fen=8/8/6k1/8/8/r1n5/Q7/K2n4&moves=5.5}{\color{white}\bfseries lösen\,\ensuremath{\nearrow}}}
\end{tcolorbox}
\vspace{2mm}
\noindent\begin{minipage}[t]{36mm}\vspace{0pt}\sffamily\small\raggedright
\hmfigure{längste eindeutige}{{\Large h\#5.5}}\hmfigure{eindeutig bei dieser Länge}{{\Large 329} Stellungen}\hmfigure{längstes Matt}{{\Large h\#7.5}}\hmfigure{in der Klasse}{15,5 Milliarden Stellungen,\\ 4,06 Milliarden mit Hilfsmatt}
\end{minipage}\hfill
\begin{minipage}[t]{80mm}\vspace{0pt}\begin{tikzpicture}\begin{axis}[width=72mm,height=42mm,ybar,bar width=1.6pt,xmin=-1,xmax=16,ymin=0,ymax=988.5,axis y line*=left,axis x line*=bottom,xlabel={\scriptsize Halbzüge bis Matt},ylabel={\scriptsize Stellungen (Mio.)},tick label style={font=\tiny},label style={font=\scriptsize},enlarge x limits=false,ylabel shift=-3pt,xlabel shift=-3pt]\addplot[fill=ink!55,draw=none] coordinates {(0,1.577) (1,12.69) (2,48.47) (3,191.2) (4,471.4) (5,741.2) (6,898.6) (7,809.5) (8,545.6) (9,275.3) (10,54.49) (11,6.721) (12,0.3012) (13,0.01671) (15,1.3e-05)};\draw[densely dotted] (axis cs:11,0) -- (axis cs:11,988.5);\end{axis}\begin{axis}[width=72mm,height=42mm,axis y line*=right,axis x line=none,ymode=log,xmin=-1,xmax=16,ylabel={\scriptsize eindeutig (log)},tick label style={font=\tiny},label style={font=\scriptsize},ylabel shift=-3pt]\addplot[mark=*,mark size=0.7pt,draw=black,thin] coordinates {(0,1577012) (1,10233169) (2,24390723) (3,37427019) (4,41187341) (5,30587820) (6,17758916) (7,8959923) (8,2826706) (9,457509) (10,24422) (11,329)};\end{axis}\end{tikzpicture}\end{minipage}
\vspace{3mm}

\noindent\begin{minipage}[t]{60mm}\vspace{0pt}\sbox\hmboard{\chessboard[setfen=8/8/6k1/8/8/r1n5/Q7/K2n4 w - - 0 1]}\usebox\hmboard\\[1mm]
\makebox[\wd\hmboard]{\sffamily\hspace{3mm}\textbf{h\#5.5}\hfill(2+4)}\end{minipage}\hfill
\begin{minipage}[t]{52mm}\vspace{0pt}
\textcolor{ink}{\sffamily\textsc{Lösung}}\par\smallskip\small
\begin{tabular}{@{}r@{\enspace}l@{\enspace}l@{}}
\textbf{1.} & \ldots & \symqueen{}xa3\\
\textbf{2.} & \symking{}f5 & \symqueen{}xc3\\
\textbf{3.} & \symking{}e4 & \symqueen{}e1+\\
\textbf{4.} & \symking{}d3 & \symking{}a2\\
\textbf{5.} & \symking{}c2 & \symking{}a3\\
\textbf{6.} & \symking{}b1 & \symqueen{}xd1\#\\
\end{tabular}\par\vspace{3mm}
\qrcode[height=12mm]{https://helpman.komtera.lt/?fen=8/8/6k1/8/8/r1n5/Q7/K2n4&moves=5.5}\\{\sffamily\tiny\color{ink}scannen zum Lösen}
\end{minipage}
\vspace{2mm}

{\raggedright\noindent\textcolor{ink}{\sffamily\textsc{Themen}}\enspace \themetag{reines Matt} \themetag{Mustermatt} \themetag{Idealmatt} \themetag{Spiegelmatt} \themetag{ein Stein zieht}\par\smallskip{\footnotesize\textit{Hinweis:} die Lösung beginnt mit einem Schlag.}\par}

\clearpage
\phantomsection\label{hm:KQvkrnp}
\begin{tcolorbox}[enhanced,colback=ink,colframe=ink,arc=0pt,boxrule=0pt,left=4mm,right=4mm,top=2.5mm,bottom=2.5mm]
\color{white}\sffamily{\LARGE\bfseries KQvkrnp}\hfill{\Large h\#5.5}\\[0.5mm]
{\small 6 Steine\quad Weiß am Zug}\hfill{\small Nr.~768\qquad\href{https://helpman.komtera.lt/?fen=8/8/7k/8/1n6/p7/8/KQr5&moves=5.5}{\color{white}\bfseries lösen\,\ensuremath{\nearrow}}}
\end{tcolorbox}
\vspace{2mm}
\noindent\begin{minipage}[t]{36mm}\vspace{0pt}\sffamily\small\raggedright
\hmfigure{längste eindeutige}{{\Large h\#5.5}}\hmfigure{eindeutig bei dieser Länge}{{\Large 390} Stellungen}\hmfigure{längstes Matt}{{\Large h\#7.5}}\hmfigure{in der Klasse}{45,4 Milliarden Stellungen,\\ 25,1 Milliarden mit Hilfsmatt}
\end{minipage}\hfill
\begin{minipage}[t]{80mm}\vspace{0pt}\begin{tikzpicture}\begin{axis}[width=72mm,height=42mm,ybar,bar width=1.6pt,xmin=-1,xmax=16,ymin=0,ymax=6.7,axis y line*=left,axis x line*=bottom,xlabel={\scriptsize Halbzüge bis Matt},ylabel={\scriptsize Stellungen (Mrd.)},tick label style={font=\tiny},label style={font=\scriptsize},enlarge x limits=false,ylabel shift=-3pt,xlabel shift=-3pt]\addplot[fill=ink!55,draw=none] coordinates {(0,0.01235) (1,0.097) (2,0.3685) (3,1.4) (4,3.679) (5,5.105) (6,6.091) (7,4.619) (8,2.564) (9,1.036) (10,0.1496) (11,0.0121) (12,0.0005418) (13,2.17e-05) (15,2e-09)};\draw[densely dotted] (axis cs:11,0) -- (axis cs:11,6.7);\end{axis}\begin{axis}[width=72mm,height=42mm,axis y line*=right,axis x line=none,ymode=log,xmin=-1,xmax=16,ylabel={\scriptsize eindeutig (log)},tick label style={font=\tiny},label style={font=\scriptsize},ylabel shift=-3pt]\addplot[mark=*,mark size=0.7pt,draw=black,thin] coordinates {(0,12350835) (1,77509410) (2,194174269) (3,259120271) (4,300013592) (5,177018420) (6,86369915) (7,35581654) (8,8443313) (9,986635) (10,52527) (11,390)};\end{axis}\end{tikzpicture}\end{minipage}
\vspace{3mm}

\noindent\begin{minipage}[t]{60mm}\vspace{0pt}\sbox\hmboard{\chessboard[setfen=8/8/7k/8/1n6/p7/8/KQr5 w - - 0 1]}\usebox\hmboard\\[1mm]
\makebox[\wd\hmboard]{\sffamily\hspace{3mm}\textbf{h\#5.5}\hfill(2+4)}\end{minipage}\hfill
\begin{minipage}[t]{52mm}\vspace{0pt}
\textcolor{ink}{\sffamily\textsc{Lösung}}\par\smallskip\small
\begin{tabular}{@{}r@{\enspace}l@{\enspace}l@{}}
\textbf{1.} & \ldots & \symqueen{}xc1+\\
\textbf{2.} & \symking{}h5 & \symqueen{}xa3\\
\textbf{3.} & \symking{}g4 & \symqueen{}a2\\
\textbf{4.} & \symking{}f3 & \symking{}b2\\
\textbf{5.} & \symking{}e2 & \symking{}c3+\\
\textbf{6.} & \symking{}d1 & \symqueen{}d2\#\\
\end{tabular}\par\vspace{3mm}
\qrcode[height=12mm]{https://helpman.komtera.lt/?fen=8/8/7k/8/1n6/p7/8/KQr5&moves=5.5}\\{\sffamily\tiny\color{ink}scannen zum Lösen}
\end{minipage}
\vspace{2mm}

{\raggedright\noindent\textcolor{ink}{\sffamily\textsc{Themen}}\enspace \themetag{ein Stein zieht}\par\smallskip{\footnotesize\textit{Hinweis:} die Lösung beginnt mit einem Schlag.}\par}

\clearpage
\phantomsection\label{hm:KQvkrpp}
\begin{tcolorbox}[enhanced,colback=ink,colframe=ink,arc=0pt,boxrule=0pt,left=4mm,right=4mm,top=2.5mm,bottom=2.5mm]
\color{white}\sffamily{\LARGE\bfseries KQvkrpp}\hfill{\Large h\#5.5}\\[0.5mm]
{\small 6 Steine\quad Weiß am Zug}\hfill{\small Nr.~769\qquad\href{https://helpman.komtera.lt/?fen=8/8/6k1/8/8/pp6/8/K2rQ3&moves=5.5}{\color{white}\bfseries lösen\,\ensuremath{\nearrow}}}
\end{tcolorbox}
\vspace{2mm}
\noindent\begin{minipage}[t]{36mm}\vspace{0pt}\sffamily\small\raggedright
\hmfigure{längste eindeutige}{{\Large h\#5.5}}\hmfigure{eindeutig bei dieser Länge}{{\Large 75} Stellungen}\hmfigure{längstes Matt}{{\Large h\#6.5}}\hmfigure{in der Klasse}{34,1 Milliarden Stellungen,\\ 9,68 Milliarden mit Hilfsmatt}
\end{minipage}\hfill
\begin{minipage}[t]{80mm}\vspace{0pt}\begin{tikzpicture}\begin{axis}[width=72mm,height=42mm,ybar,bar width=1.6pt,xmin=-1,xmax=14,ymin=0,ymax=2.661,axis y line*=left,axis x line*=bottom,xlabel={\scriptsize Halbzüge bis Matt},ylabel={\scriptsize Stellungen (Mrd.)},tick label style={font=\tiny},label style={font=\scriptsize},enlarge x limits=false,ylabel shift=-3pt,xlabel shift=-3pt]\addplot[fill=ink!55,draw=none] coordinates {(0,0.006053) (1,0.04636) (2,0.1747) (3,0.6387) (4,1.748) (5,2.13) (6,2.419) (7,1.543) (8,0.709) (9,0.2331) (10,0.02588) (11,0.001275) (12,6.139e-05) (13,9.05e-07)};\draw[densely dotted] (axis cs:11,0) -- (axis cs:11,2.661);\end{axis}\begin{axis}[width=72mm,height=42mm,axis y line*=right,axis x line=none,ymode=log,xmin=-1,xmax=14,ylabel={\scriptsize eindeutig (log)},tick label style={font=\tiny},label style={font=\scriptsize},ylabel shift=-3pt]\addplot[mark=*,mark size=0.7pt,draw=black,thin] coordinates {(0,6053170) (1,36716501) (2,97185153) (3,112555197) (4,133331996) (5,60124911) (6,24167482) (7,8511727) (8,1563817) (9,124996) (10,6685) (11,75)};\end{axis}\end{tikzpicture}\end{minipage}
\vspace{3mm}

\noindent\begin{minipage}[t]{60mm}\vspace{0pt}\sbox\hmboard{\chessboard[setfen=8/8/6k1/8/8/pp6/8/K2rQ3 w - - 0 1]}\usebox\hmboard\\[1mm]
\makebox[\wd\hmboard]{\sffamily\hspace{3mm}\textbf{h\#5.5}\hfill(2+4)}\end{minipage}\hfill
\begin{minipage}[t]{52mm}\vspace{0pt}
\textcolor{ink}{\sffamily\textsc{Lösung}}\par\smallskip\small
\begin{tabular}{@{}r@{\enspace}l@{\enspace}l@{}}
\textbf{1.} & \ldots & \symqueen{}xd1\\
\textbf{2.} & \symking{}f5 & \symqueen{}xb3\\
\textbf{3.} & \symking{}e4 & \symqueen{}e6+\\
\textbf{4.} & \symking{}d3 & \symking{}a2\\
\textbf{5.} & \symking{}d2 & \symking{}b3\\
\textbf{6.} & \symking{}c1 & \symqueen{}e1\#\\
\end{tabular}\par\vspace{3mm}
\qrcode[height=12mm]{https://helpman.komtera.lt/?fen=8/8/6k1/8/8/pp6/8/K2rQ3&moves=5.5}\\{\sffamily\tiny\color{ink}scannen zum Lösen}
\end{minipage}
\vspace{2mm}

{\raggedright\noindent\textcolor{ink}{\sffamily\textsc{Themen}}\enspace \themetag{reines Matt} \themetag{Mustermatt} \themetag{Spiegelmatt} \themetag{Rundlauf} \themetag{ein Stein zieht}\par\smallskip{\footnotesize\textit{Hinweis:} die Partei am Zug steht im Diagramm im Schach; die Lösung beginnt mit einem Schlag.}\par}

\clearpage
\phantomsection\label{hm:KRRBvkq}
\begin{tcolorbox}[enhanced,colback=ink,colframe=ink,arc=0pt,boxrule=0pt,left=4mm,right=4mm,top=2.5mm,bottom=2.5mm]
\color{white}\sffamily{\LARGE\bfseries KRRBvkq}\hfill{\Large h\#5.5}\\[0.5mm]
{\small 6 Steine\quad Weiß am Zug}\hfill{\small Nr.~770\qquad\href{https://helpman.komtera.lt/?fen=1B6/8/8/7R/6k1/1R6/q7/K7&moves=5.5}{\color{white}\bfseries lösen\,\ensuremath{\nearrow}}}
\end{tcolorbox}
\vspace{2mm}
\noindent\begin{minipage}[t]{36mm}\vspace{0pt}\sffamily\small\raggedright
\hmfigure{längste eindeutige}{{\Large h\#5.5}}\hmfigure{eindeutig bei dieser Länge}{{\Large 90} Stellungen}\hmfigure{längstes Matt}{{\Large h\#7.5}}\hmfigure{in der Klasse}{15,5 Milliarden Stellungen,\\ 3,71 Milliarden mit Hilfsmatt}
\end{minipage}\hfill
\begin{minipage}[t]{80mm}\vspace{0pt}\begin{tikzpicture}\begin{axis}[width=72mm,height=42mm,ybar,bar width=1.6pt,xmin=-1,xmax=16,ymin=0,ymax=1.645,axis y line*=left,axis x line*=bottom,xlabel={\scriptsize Halbzüge bis Matt},ylabel={\scriptsize Stellungen (Mrd.)},tick label style={font=\tiny},label style={font=\scriptsize},enlarge x limits=false,ylabel shift=-3pt,xlabel shift=-3pt]\addplot[fill=ink!55,draw=none] coordinates {(0,0.01156) (1,0.04106) (2,0.2005) (3,0.7276) (4,1.496) (5,0.8835) (6,0.3249) (7,0.02197) (8,0.005069) (9,7.981e-05) (10,0.0001352) (11,3.834e-05) (12,5.89e-07) (13,6.578e-06) (15,1.7e-08)};\draw[densely dotted] (axis cs:11,0) -- (axis cs:11,1.645);\end{axis}\begin{axis}[width=72mm,height=42mm,axis y line*=right,axis x line=none,ymode=log,xmin=-1,xmax=16,ylabel={\scriptsize eindeutig (log)},tick label style={font=\tiny},label style={font=\scriptsize},ylabel shift=-3pt]\addplot[mark=*,mark size=0.7pt,draw=black,thin] coordinates {(0,11556958) (1,37669320) (2,93188785) (3,153257554) (4,105367544) (5,4804610) (6,492990) (7,10435) (8,10150) (9,719) (10,79) (11,90)};\end{axis}\end{tikzpicture}\end{minipage}
\vspace{3mm}

\noindent\begin{minipage}[t]{60mm}\vspace{0pt}\sbox\hmboard{\chessboard[setfen=1B6/8/8/7R/6k1/1R6/q7/K7 w - - 0 1]}\usebox\hmboard\\[1mm]
\makebox[\wd\hmboard]{\sffamily\hspace{3mm}\textbf{h\#5.5}\hfill(4+2)}\end{minipage}\hfill
\begin{minipage}[t]{52mm}\vspace{0pt}
\textcolor{ink}{\sffamily\textsc{Lösung}}\par\smallskip\small
\begin{tabular}{@{}r@{\enspace}l@{\enspace}l@{}}
\textbf{1.} & \ldots & \symking{}xa2\\
\textbf{2.} & \symking{}xh5 & \symking{}b2\\
\textbf{3.} & \symking{}g4 & \symking{}c3\\
\textbf{4.} & \symking{}f3 & \symking{}d3\\
\textbf{5.} & \symking{}f2 & \symbishop{}a7+\\
\textbf{6.} & \symking{}e1 & \symrook{}b1\#\\
\end{tabular}\par\vspace{3mm}
\qrcode[height=12mm]{https://helpman.komtera.lt/?fen=1B6/8/8/7R/6k1/1R6/q7/K7&moves=5.5}\\{\sffamily\tiny\color{ink}scannen zum Lösen}
\end{minipage}
\vspace{2mm}

{\raggedright\noindent\textcolor{ink}{\sffamily\textsc{Themen}}\enspace \themetag{reines Matt} \themetag{Mustermatt} \themetag{Idealmatt} \themetag{Spiegelmatt} \themetag{Switchback} \themetag{ein Stein zieht}\par\smallskip{\footnotesize\textit{Hinweis:} die Partei am Zug steht im Diagramm im Schach; die Lösung beginnt mit einem Schlag.}\par}

\clearpage
\phantomsection\label{hm:KRRNvkq}
\begin{tcolorbox}[enhanced,colback=ink,colframe=ink,arc=0pt,boxrule=0pt,left=4mm,right=4mm,top=2.5mm,bottom=2.5mm]
\color{white}\sffamily{\LARGE\bfseries KRRNvkq}\hfill{\Large h\#5.5}\\[0.5mm]
{\small 6 Steine\quad Weiß am Zug}\hfill{\small Nr.~771\qquad\href{https://helpman.komtera.lt/?fen=3N4/8/8/6R1/7k/2R5/q7/K7&moves=5.5}{\color{white}\bfseries lösen\,\ensuremath{\nearrow}}}
\end{tcolorbox}
\vspace{2mm}
\noindent\begin{minipage}[t]{36mm}\vspace{0pt}\sffamily\small\raggedright
\hmfigure{längste eindeutige}{{\Large h\#5.5}}\hmfigure{eindeutig bei dieser Länge}{{\Large 1} Stellung}\hmfigure{längstes Matt}{{\Large h\#6.5}}\hmfigure{in der Klasse}{15,5 Milliarden Stellungen,\\ 3,79 Milliarden mit Hilfsmatt}
\end{minipage}\hfill
\begin{minipage}[t]{80mm}\vspace{0pt}\begin{tikzpicture}\begin{axis}[width=72mm,height=42mm,ybar,bar width=1.6pt,xmin=-1,xmax=14,ymin=0,ymax=1.597,axis y line*=left,axis x line*=bottom,xlabel={\scriptsize Halbzüge bis Matt},ylabel={\scriptsize Stellungen (Mrd.)},tick label style={font=\tiny},label style={font=\scriptsize},enlarge x limits=false,ylabel shift=-3pt,xlabel shift=-3pt]\addplot[fill=ink!55,draw=none] coordinates {(0,0.009927) (1,0.03864) (2,0.1868) (3,0.7285) (4,1.452) (5,0.9574) (6,0.383) (7,0.03167) (8,0.006307) (9,7.433e-05) (10,0.0005267) (11,2.094e-05) (12,8.96e-07) (13,9.61e-07)};\draw[densely dotted] (axis cs:11,0) -- (axis cs:11,1.597);\end{axis}\begin{axis}[width=72mm,height=42mm,axis y line*=right,axis x line=none,ymode=log,xmin=-1,xmax=14,ylabel={\scriptsize eindeutig (log)},tick label style={font=\tiny},label style={font=\scriptsize},ylabel shift=-3pt]\addplot[mark=*,mark size=0.7pt,draw=black,thin] coordinates {(0,9926595) (1,35635402) (2,85416518) (3,158002012) (4,100738670) (5,5784083) (6,488029) (7,8374) (8,21541) (9,617) (10,99) (11,1)};\end{axis}\end{tikzpicture}\end{minipage}
\vspace{3mm}

\noindent\begin{minipage}[t]{60mm}\vspace{0pt}\sbox\hmboard{\chessboard[setfen=3N4/8/8/6R1/7k/2R5/q7/K7 w - - 0 1]}\usebox\hmboard\\[1mm]
\makebox[\wd\hmboard]{\sffamily\hspace{3mm}\textbf{h\#5.5}\hfill(4+2)}\end{minipage}\hfill
\begin{minipage}[t]{52mm}\vspace{0pt}
\textcolor{ink}{\sffamily\textsc{Lösung}}\par\smallskip\small
\begin{tabular}{@{}r@{\enspace}l@{\enspace}l@{}}
\textbf{1.} & \ldots & \symking{}xa2\\
\textbf{2.} & \symking{}xg5 & \symking{}b3\\
\textbf{3.} & \symking{}f6 & \symking{}c4\\
\textbf{4.} & \symking{}e7 & \symking{}d5\\
\textbf{5.} & \symking{}xd8 & \symking{}e6\\
\textbf{6.} & \symking{}e8 & \symrook{}c8\#\\
\end{tabular}\par\vspace{3mm}
\qrcode[height=12mm]{https://helpman.komtera.lt/?fen=3N4/8/8/6R1/7k/2R5/q7/K7&moves=5.5}\\{\sffamily\tiny\color{ink}scannen zum Lösen}
\end{minipage}
\vspace{2mm}

{\raggedright\noindent\textcolor{ink}{\sffamily\textsc{Themen}}\enspace \themetag{reines Matt} \themetag{Mustermatt} \themetag{Idealmatt} \themetag{Spiegelmatt} \themetag{ein Stein zieht} \themetag{ohne Schach}\par\smallskip{\footnotesize\textit{Hinweis:} die Partei am Zug steht im Diagramm im Schach; die Lösung beginnt mit einem Schlag.}\par}

\clearpage
\phantomsection\label{hm:KNNNNvk}
\begin{tcolorbox}[enhanced,colback=ink,colframe=ink,arc=0pt,boxrule=0pt,left=4mm,right=4mm,top=2.5mm,bottom=2.5mm]
\color{white}\sffamily{\LARGE\bfseries KNNNNvk}\hfill{\Large h\#5}\\[0.5mm]
{\small 6 Steine\quad Schwarz am Zug}\hfill{\small Nr.~772\qquad\href{https://helpman.komtera.lt/?fen=N5N1/8/8/7N/8/N7/8/K6k&moves=5}{\color{white}\bfseries lösen\,\ensuremath{\nearrow}}}
\end{tcolorbox}
\vspace{2mm}
\noindent\begin{minipage}[t]{36mm}\vspace{0pt}\sffamily\small\raggedright
\hmfigure{längste eindeutige}{{\Large h\#5}}\hmfigure{eindeutig bei dieser Länge}{{\Large 768} Stellungen}\hmfigure{längstes Matt}{{\Large h\#7}}\hmfigure{in der Klasse}{15,5 Milliarden Stellungen,\\ 430 Millionen mit Hilfsmatt}
\end{minipage}\hfill
\begin{minipage}[t]{80mm}\vspace{0pt}\begin{tikzpicture}\begin{axis}[width=72mm,height=42mm,ybar,bar width=1.6pt,xmin=-1,xmax=15,ymin=0,ymax=128.6,axis y line*=left,axis x line*=bottom,xlabel={\scriptsize Halbzüge bis Matt},ylabel={\scriptsize Stellungen (Mio.)},tick label style={font=\tiny},label style={font=\scriptsize},enlarge x limits=false,ylabel shift=-3pt,xlabel shift=-3pt]\addplot[fill=ink!55,draw=none] coordinates {(0,0.08828) (1,0.3496) (2,1.074) (3,4.269) (4,18.72) (5,54.08) (6,116.9) (7,95.53) (8,106.1) (9,23.77) (10,8.333) (11,0.2839) (12,0.04028) (13,5.3e-05) (14,0.000396)};\draw[densely dotted] (axis cs:10,0) -- (axis cs:10,128.6);\end{axis}\begin{axis}[width=72mm,height=42mm,axis y line*=right,axis x line=none,ymode=log,xmin=-1,xmax=15,ylabel={\scriptsize eindeutig (log)},tick label style={font=\tiny},label style={font=\scriptsize},ylabel shift=-3pt]\addplot[mark=*,mark size=0.7pt,draw=black,thin] coordinates {(0,88275) (1,287545) (2,872383) (3,2401668) (4,7696671) (5,5150465) (6,5463919) (7,444849) (8,213827) (9,2070) (10,768)};\end{axis}\end{tikzpicture}\end{minipage}
\vspace{3mm}

\noindent\begin{minipage}[t]{60mm}\vspace{0pt}\sbox\hmboard{\chessboard[setfen=N5N1/8/8/7N/8/N7/8/K6k b - - 0 1]}\usebox\hmboard\\[1mm]
\makebox[\wd\hmboard]{\sffamily\hspace{3mm}\textbf{h\#5}\hfill(5+1)}\end{minipage}\hfill
\begin{minipage}[t]{52mm}\vspace{0pt}
\textcolor{ink}{\sffamily\textsc{Lösung}}\par\smallskip\small
\begin{tabular}{@{}r@{\enspace}l@{\enspace}l@{}}
\textbf{1.} & \symking{}g2 & \symking{}b2\\
\textbf{2.} & \symking{}f3 & \symking{}c3\\
\textbf{3.} & \symking{}e4 & \symknight{}g3+\\
\textbf{4.} & \symking{}d5 & \symknight{}e7+\\
\textbf{5.} & \symking{}c5 & \symknight{}e4\#\\
\end{tabular}\par\vspace{3mm}
\qrcode[height=12mm]{https://helpman.komtera.lt/?fen=N5N1/8/8/7N/8/N7/8/K6k&moves=5}\\{\sffamily\tiny\color{ink}scannen zum Lösen}
\end{minipage}
\vspace{2mm}

{\raggedright\noindent\textcolor{ink}{\sffamily\textsc{Themen}}\enspace \themetag{Spiegelmatt} \themetag{ein Stein zieht} \themetag{ohne Schlag}\par}

\clearpage
\phantomsection\label{hm:KQBvkqq}
\begin{tcolorbox}[enhanced,colback=ink,colframe=ink,arc=0pt,boxrule=0pt,left=4mm,right=4mm,top=2.5mm,bottom=2.5mm]
\color{white}\sffamily{\LARGE\bfseries KQBvkqq}\hfill{\Large h\#5}\\[0.5mm]
{\small 6 Steine\quad Schwarz am Zug}\hfill{\small Nr.~773\qquad\href{https://helpman.komtera.lt/?fen=8/8/8/8/1q6/3q4/7B/K4Qk1&moves=5}{\color{white}\bfseries lösen\,\ensuremath{\nearrow}}}
\end{tcolorbox}
\vspace{2mm}
\noindent\begin{minipage}[t]{36mm}\vspace{0pt}\sffamily\small\raggedright
\hmfigure{längste eindeutige}{{\Large h\#5}}\hmfigure{eindeutig bei dieser Länge}{{\Large 673} Stellungen}\hmfigure{längstes Matt}{{\Large h\#8}}\hmfigure{in der Klasse}{15,5 Milliarden Stellungen,\\ 3,04 Milliarden mit Hilfsmatt}
\end{minipage}\hfill
\begin{minipage}[t]{80mm}\vspace{0pt}\begin{tikzpicture}\begin{axis}[width=72mm,height=42mm,ybar,bar width=1.6pt,xmin=-1,xmax=17,ymin=0,ymax=1.235,axis y line*=left,axis x line*=bottom,xlabel={\scriptsize Halbzüge bis Matt},ylabel={\scriptsize Stellungen (Mrd.)},tick label style={font=\tiny},label style={font=\scriptsize},enlarge x limits=false,ylabel shift=-3pt,xlabel shift=-3pt]\addplot[fill=ink!55,draw=none] coordinates {(0,0.002444) (1,0.0189) (2,0.1124) (3,0.5419) (4,1.009) (5,1.122) (6,0.1866) (7,0.04556) (8,0.00118) (9,1.286e-05) (10,0.0003796) (11,7.76e-07) (12,1.523e-05) (13,1.08e-07) (14,8e-09) (15,7.5e-08) (16,2.2e-08)};\draw[densely dotted] (axis cs:10,0) -- (axis cs:10,1.235);\end{axis}\begin{axis}[width=72mm,height=42mm,axis y line*=right,axis x line=none,ymode=log,xmin=-1,xmax=17,ylabel={\scriptsize eindeutig (log)},tick label style={font=\tiny},label style={font=\scriptsize},ylabel shift=-3pt]\addplot[mark=*,mark size=0.7pt,draw=black,thin] coordinates {(0,2444393) (1,17601302) (2,65718146) (3,124438502) (4,53305347) (5,10511847) (6,88349) (7,791) (8,15198) (9,78) (10,673)};\end{axis}\end{tikzpicture}\end{minipage}
\vspace{3mm}

\noindent\begin{minipage}[t]{60mm}\vspace{0pt}\sbox\hmboard{\chessboard[setfen=8/8/8/8/1q6/3q4/7B/K4Qk1 b - - 0 1]}\usebox\hmboard\\[1mm]
\makebox[\wd\hmboard]{\sffamily\hspace{3mm}\textbf{h\#5}\hfill(3+3)}\end{minipage}\hfill
\begin{minipage}[t]{52mm}\vspace{0pt}
\textcolor{ink}{\sffamily\textsc{Lösung}}\par\smallskip\small
\begin{tabular}{@{}r@{\enspace}l@{\enspace}l@{}}
\textbf{1.} & \symking{}xh2 & \symqueen{}e1\\
\textbf{2.} & \symqueen{}f4 & \symking{}b2\\
\textbf{3.} & \symqueen{}dg3 & \symking{}c2\\
\textbf{4.} & \symking{}g2 & \symking{}d1\\
\textbf{5.} & \symking{}f3 & \symqueen{}e2\#\\
\end{tabular}\par\vspace{3mm}
\qrcode[height=12mm]{https://helpman.komtera.lt/?fen=8/8/8/8/1q6/3q4/7B/K4Qk1&moves=5}\\{\sffamily\tiny\color{ink}scannen zum Lösen}
\end{minipage}
\vspace{2mm}

{\raggedright\noindent\textcolor{ink}{\sffamily\textsc{Themen}}\enspace \themetag{reines Matt} \themetag{Mustermatt} \themetag{Idealmatt} \themetag{Selbstblock} \themetag{ohne Schach}\par\smallskip{\footnotesize\textit{Hinweis:} das Diagramm hat keinen legalen letzten Zug, kann also in keiner Partie entstehen.}\par}

\clearpage
\phantomsection\label{hm:KQBvkrr}
\begin{tcolorbox}[enhanced,colback=ink,colframe=ink,arc=0pt,boxrule=0pt,left=4mm,right=4mm,top=2.5mm,bottom=2.5mm]
\color{white}\sffamily{\LARGE\bfseries KQBvkrr}\hfill{\Large h\#5}\\[0.5mm]
{\small 6 Steine\quad Schwarz am Zug}\hfill{\small Nr.~774\qquad\href{https://helpman.komtera.lt/?fen=8/8/8/5B2/5k2/1rr1Q3/8/K7&moves=5}{\color{white}\bfseries lösen\,\ensuremath{\nearrow}}}
\end{tcolorbox}
\vspace{2mm}
\noindent\begin{minipage}[t]{36mm}\vspace{0pt}\sffamily\small\raggedright
\hmfigure{längste eindeutige}{{\Large h\#5}}\hmfigure{eindeutig bei dieser Länge}{{\Large 7} Stellungen}\hmfigure{längstes Matt}{{\Large h\#6}}\hmfigure{in der Klasse}{15,5 Milliarden Stellungen,\\ 3,64 Milliarden mit Hilfsmatt}
\end{minipage}\hfill
\begin{minipage}[t]{80mm}\vspace{0pt}\begin{tikzpicture}\begin{axis}[width=72mm,height=42mm,ybar,bar width=1.6pt,xmin=-1,xmax=13,ymin=0,ymax=1.626,axis y line*=left,axis x line*=bottom,xlabel={\scriptsize Halbzüge bis Matt},ylabel={\scriptsize Stellungen (Mrd.)},tick label style={font=\tiny},label style={font=\scriptsize},enlarge x limits=false,ylabel shift=-3pt,xlabel shift=-3pt]\addplot[fill=ink!55,draw=none] coordinates {(0,0.004956) (1,0.03711) (2,0.2041) (3,0.7616) (4,1.478) (5,0.9461) (6,0.1892) (7,0.01666) (8,0.00151) (9,3.91e-07) (10,0.0001596) (11,5e-09) (12,7.2e-08)};\draw[densely dotted] (axis cs:10,0) -- (axis cs:10,1.626);\end{axis}\begin{axis}[width=72mm,height=42mm,axis y line*=right,axis x line=none,ymode=log,xmin=-1,xmax=13,ylabel={\scriptsize eindeutig (log)},tick label style={font=\tiny},label style={font=\scriptsize},ylabel shift=-3pt]\addplot[mark=*,mark size=0.7pt,draw=black,thin] coordinates {(0,4956073) (1,33277413) (2,122523652) (3,150643451) (4,63915144) (5,3897080) (6,60100) (7,43) (8,5362) (10,7)};\end{axis}\end{tikzpicture}\end{minipage}
\vspace{3mm}

\noindent\begin{minipage}[t]{60mm}\vspace{0pt}\sbox\hmboard{\chessboard[setfen=8/8/8/5B2/5k2/1rr1Q3/8/K7 b - - 0 1]}\usebox\hmboard\\[1mm]
\makebox[\wd\hmboard]{\sffamily\hspace{3mm}\textbf{h\#5}\hfill(3+3)}\end{minipage}\hfill
\begin{minipage}[t]{52mm}\vspace{0pt}
\textcolor{ink}{\sffamily\textsc{Lösung}}\par\smallskip\small
\begin{tabular}{@{}r@{\enspace}l@{\enspace}l@{}}
\textbf{1.} & \symking{}xf5 & \symqueen{}xc3\\
\textbf{2.} & \symking{}e4 & \symqueen{}e5+\\
\textbf{3.} & \symking{}d3 & \symking{}a2\\
\textbf{4.} & \symking{}d2 & \symking{}xb3\\
\textbf{5.} & \symking{}c1 & \symqueen{}e1\#\\
\end{tabular}\par\vspace{3mm}
\qrcode[height=12mm]{https://helpman.komtera.lt/?fen=8/8/8/5B2/5k2/1rr1Q3/8/K7&moves=5}\\{\sffamily\tiny\color{ink}scannen zum Lösen}
\end{minipage}
\vspace{2mm}

{\raggedright\noindent\textcolor{ink}{\sffamily\textsc{Themen}}\enspace \themetag{reines Matt} \themetag{Mustermatt} \themetag{Idealmatt} \themetag{Spiegelmatt} \themetag{ein Stein zieht}\par\smallskip{\footnotesize\textit{Hinweis:} die Partei am Zug steht im Diagramm im Schach; die Lösung beginnt mit einem Schlag.}\par}

\clearpage
\phantomsection\label{hm:KQNvkqq}
\begin{tcolorbox}[enhanced,colback=ink,colframe=ink,arc=0pt,boxrule=0pt,left=4mm,right=4mm,top=2.5mm,bottom=2.5mm]
\color{white}\sffamily{\LARGE\bfseries KQNvkqq}\hfill{\Large h\#5}\\[0.5mm]
{\small 6 Steine\quad Schwarz am Zug}\hfill{\small Nr.~775\qquad\href{https://helpman.komtera.lt/?fen=8/8/8/8/1q6/3q4/7N/K4Qk1&moves=5}{\color{white}\bfseries lösen\,\ensuremath{\nearrow}}}
\end{tcolorbox}
\vspace{2mm}
\noindent\begin{minipage}[t]{36mm}\vspace{0pt}\sffamily\small\raggedright
\hmfigure{längste eindeutige}{{\Large h\#5}}\hmfigure{eindeutig bei dieser Länge}{{\Large 326} Stellungen}\hmfigure{längstes Matt}{{\Large h\#7}}\hmfigure{in der Klasse}{15,5 Milliarden Stellungen,\\ 3,15 Milliarden mit Hilfsmatt}
\end{minipage}\hfill
\begin{minipage}[t]{80mm}\vspace{0pt}\begin{tikzpicture}\begin{axis}[width=72mm,height=42mm,ybar,bar width=1.6pt,xmin=-1,xmax=15,ymin=0,ymax=1.271,axis y line*=left,axis x line*=bottom,xlabel={\scriptsize Halbzüge bis Matt},ylabel={\scriptsize Stellungen (Mrd.)},tick label style={font=\tiny},label style={font=\scriptsize},enlarge x limits=false,ylabel shift=-3pt,xlabel shift=-3pt]\addplot[fill=ink!55,draw=none] coordinates {(0,0.001418) (1,0.01317) (2,0.07507) (3,0.3792) (4,0.758) (5,1.156) (6,0.4641) (7,0.283) (8,0.01639) (9,0.004523) (10,0.0001576) (11,1.73e-07) (12,6.917e-06) (13,1e-09) (14,9e-09)};\draw[densely dotted] (axis cs:10,0) -- (axis cs:10,1.271);\end{axis}\begin{axis}[width=72mm,height=42mm,axis y line*=right,axis x line=none,ymode=log,xmin=-1,xmax=15,ylabel={\scriptsize eindeutig (log)},tick label style={font=\tiny},label style={font=\scriptsize},ylabel shift=-3pt]\addplot[mark=*,mark size=0.7pt,draw=black,thin] coordinates {(0,1418175) (1,12283438) (2,42593968) (3,86223698) (4,55596783) (5,25226030) (6,1072130) (7,228609) (8,5399) (9,9) (10,326)};\end{axis}\end{tikzpicture}\end{minipage}
\vspace{3mm}

\noindent\begin{minipage}[t]{60mm}\vspace{0pt}\sbox\hmboard{\chessboard[setfen=8/8/8/8/1q6/3q4/7N/K4Qk1 b - - 0 1]}\usebox\hmboard\\[1mm]
\makebox[\wd\hmboard]{\sffamily\hspace{3mm}\textbf{h\#5}\hfill(3+3)}\end{minipage}\hfill
\begin{minipage}[t]{52mm}\vspace{0pt}
\textcolor{ink}{\sffamily\textsc{Lösung}}\par\smallskip\small
\begin{tabular}{@{}r@{\enspace}l@{\enspace}l@{}}
\textbf{1.} & \symking{}xh2 & \symqueen{}e1\\
\textbf{2.} & \symqueen{}f4 & \symking{}b2\\
\textbf{3.} & \symqueen{}dg3 & \symking{}c2\\
\textbf{4.} & \symking{}g2 & \symking{}d1\\
\textbf{5.} & \symking{}f3 & \symqueen{}e2\#\\
\end{tabular}\par\vspace{3mm}
\qrcode[height=12mm]{https://helpman.komtera.lt/?fen=8/8/8/8/1q6/3q4/7N/K4Qk1&moves=5}\\{\sffamily\tiny\color{ink}scannen zum Lösen}
\end{minipage}
\vspace{2mm}

{\raggedright\noindent\textcolor{ink}{\sffamily\textsc{Themen}}\enspace \themetag{reines Matt} \themetag{Mustermatt} \themetag{Idealmatt} \themetag{Selbstblock} \themetag{ohne Schach}\par\smallskip{\footnotesize\textit{Hinweis:} die Partei am Zug steht im Diagramm im Schach; die Lösung beginnt mit einem Schlag.}\par}

\clearpage
\phantomsection\label{hm:KQQPvkb}
\begin{tcolorbox}[enhanced,colback=ink,colframe=ink,arc=0pt,boxrule=0pt,left=4mm,right=4mm,top=2.5mm,bottom=2.5mm]
\color{white}\sffamily{\LARGE\bfseries KQQPvkb}\hfill{\Large h\#5}\\[0.5mm]
{\small 6 Steine\quad Schwarz am Zug}\hfill{\small Nr.~776\qquad\href{https://helpman.komtera.lt/?fen=8/8/8/8/8/8/1Pb4Q/K2Q2k1&moves=5}{\color{white}\bfseries lösen\,\ensuremath{\nearrow}}}
\end{tcolorbox}
\vspace{2mm}
\noindent\begin{minipage}[t]{36mm}\vspace{0pt}\sffamily\small\raggedright
\hmfigure{längste eindeutige}{{\Large h\#5}}\hmfigure{eindeutig bei dieser Länge}{{\Large 4\,745} Stellungen}\hmfigure{längstes Matt}{{\Large h\#7}}\hmfigure{in der Klasse}{45,4 Milliarden Stellungen,\\ 11,8 Milliarden mit Hilfsmatt}
\end{minipage}\hfill
\begin{minipage}[t]{80mm}\vspace{0pt}\begin{tikzpicture}\begin{axis}[width=72mm,height=42mm,ybar,bar width=1.6pt,xmin=-1,xmax=15,ymin=0,ymax=4.529,axis y line*=left,axis x line*=bottom,xlabel={\scriptsize Halbzüge bis Matt},ylabel={\scriptsize Stellungen (Mrd.)},tick label style={font=\tiny},label style={font=\scriptsize},enlarge x limits=false,ylabel shift=-3pt,xlabel shift=-3pt]\addplot[fill=ink!55,draw=none] coordinates {(0,0.2451) (1,0.7445) (2,3.32) (3,2.922) (4,4.117) (5,0.2466) (6,0.1534) (7,1.42e-05) (8,0.07238) (9,1.35e-07) (10,0.01048) (11,8e-09) (12,0.0002539) (14,1e-08)};\draw[densely dotted] (axis cs:10,0) -- (axis cs:10,4.529);\end{axis}\begin{axis}[width=72mm,height=42mm,axis y line*=right,axis x line=none,ymode=log,xmin=-1,xmax=15,ylabel={\scriptsize eindeutig (log)},tick label style={font=\tiny},label style={font=\scriptsize},ylabel shift=-3pt]\addplot[mark=*,mark size=0.7pt,draw=black,thin] coordinates {(0,245089531) (1,405176381) (2,1075893518) (3,65649127) (4,14258553) (5,2602) (6,4139069) (7,27) (8,531813) (10,4745)};\end{axis}\end{tikzpicture}\end{minipage}
\vspace{3mm}

\noindent\begin{minipage}[t]{60mm}\vspace{0pt}\sbox\hmboard{\chessboard[setfen=8/8/8/8/8/8/1Pb4Q/K2Q2k1 b - - 0 1]}\usebox\hmboard\\[1mm]
\makebox[\wd\hmboard]{\sffamily\hspace{3mm}\textbf{h\#5}\hfill(4+2)}\end{minipage}\hfill
\begin{minipage}[t]{52mm}\vspace{0pt}
\textcolor{ink}{\sffamily\textsc{Lösung}}\par\smallskip\small
\begin{tabular}{@{}r@{\enspace}l@{\enspace}l@{}}
\textbf{1.} & \symking{}xh2 & \symqueen{}d4\\
\textbf{2.} & \symbishop{}b1 & \symking{}xb1\\
\textbf{3.} & \symking{}g2 & \symking{}c2\\
\textbf{4.} & \symking{}f1 & \symking{}d3\\
\textbf{5.} & \symking{}e1 & \symqueen{}g1\#\\
\end{tabular}\par\vspace{3mm}
\qrcode[height=12mm]{https://helpman.komtera.lt/?fen=8/8/8/8/8/8/1Pb4Q/K2Q2k1&moves=5}\\{\sffamily\tiny\color{ink}scannen zum Lösen}
\end{minipage}
\vspace{2mm}

{\raggedright\noindent\textcolor{ink}{\sffamily\textsc{Themen}}\enspace \themetag{reines Matt} \themetag{Mustermatt} \themetag{Spiegelmatt} \themetag{ohne Schach}\par\smallskip{\footnotesize\textit{Hinweis:} das Diagramm hat keinen legalen letzten Zug, kann also in keiner Partie entstehen.}\par}

\clearpage
\phantomsection\label{hm:KQQPvkn}
\begin{tcolorbox}[enhanced,colback=ink,colframe=ink,arc=0pt,boxrule=0pt,left=4mm,right=4mm,top=2.5mm,bottom=2.5mm]
\color{white}\sffamily{\LARGE\bfseries KQQPvkn}\hfill{\Large h\#5}\\[0.5mm]
{\small 6 Steine\quad Schwarz am Zug}\hfill{\small Nr.~777\qquad\href{https://helpman.komtera.lt/?fen=8/1P6/8/8/8/8/Q1Q5/K1k1n3&moves=5}{\color{white}\bfseries lösen\,\ensuremath{\nearrow}}}
\end{tcolorbox}
\vspace{2mm}
\noindent\begin{minipage}[t]{36mm}\vspace{0pt}\sffamily\small\raggedright
\hmfigure{längste eindeutige}{{\Large h\#5}}\hmfigure{eindeutig bei dieser Länge}{{\Large 7\,801} Stellungen}\hmfigure{längstes Matt}{{\Large h\#9}}\hmfigure{in der Klasse}{45,4 Milliarden Stellungen,\\ 12,2 Milliarden mit Hilfsmatt}
\end{minipage}\hfill
\begin{minipage}[t]{80mm}\vspace{0pt}\begin{tikzpicture}\begin{axis}[width=72mm,height=42mm,ybar,bar width=1.6pt,xmin=-1,xmax=19,ymin=0,ymax=4.862,axis y line*=left,axis x line*=bottom,xlabel={\scriptsize Halbzüge bis Matt},ylabel={\scriptsize Stellungen (Mrd.)},tick label style={font=\tiny},label style={font=\scriptsize},enlarge x limits=false,ylabel shift=-3pt,xlabel shift=-3pt]\addplot[fill=ink!55,draw=none] coordinates {(0,0.2693) (1,0.788) (2,3.377) (3,2.917) (4,4.42) (5,0.2077) (6,0.161) (7,3.294e-05) (8,0.08251) (9,1.077e-06) (10,0.01304) (12,0.0003729) (14,1.7e-07) (16,3.11e-07) (18,1.33e-07)};\draw[densely dotted] (axis cs:10,0) -- (axis cs:10,4.862);\end{axis}\begin{axis}[width=72mm,height=42mm,axis y line*=right,axis x line=none,ymode=log,xmin=-1,xmax=19,ylabel={\scriptsize eindeutig (log)},tick label style={font=\tiny},label style={font=\scriptsize},ylabel shift=-3pt]\addplot[mark=*,mark size=0.7pt,draw=black,thin] coordinates {(0,269300377) (1,415980061) (2,1071962373) (3,56213428) (4,14445750) (5,3266) (6,5045542) (7,1012) (8,630304) (10,7801)};\end{axis}\end{tikzpicture}\end{minipage}
\vspace{3mm}

\noindent\begin{minipage}[t]{60mm}\vspace{0pt}\sbox\hmboard{\chessboard[setfen=8/1P6/8/8/8/8/Q1Q5/K1k1n3 b - - 0 1]}\usebox\hmboard\\[1mm]
\makebox[\wd\hmboard]{\sffamily\hspace{3mm}\textbf{h\#5}\hfill(4+2)}\end{minipage}\hfill
\begin{minipage}[t]{52mm}\vspace{0pt}
\textcolor{ink}{\sffamily\textsc{Lösung}}\par\smallskip\small
\begin{tabular}{@{}r@{\enspace}l@{\enspace}l@{}}
\textbf{1.} & \symknight{}xc2+ & \symqueen{}xc2+\\
\textbf{2.} & \symking{}xc2 & \symking{}a2\\
\textbf{3.} & \symking{}c1 & \symking{}a3\\
\textbf{4.} & \symking{}b1 & b8=\symqueen{}+\\
\textbf{5.} & \symking{}a1 & \symqueen{}b2\#\\
\end{tabular}\par\vspace{3mm}
\qrcode[height=12mm]{https://helpman.komtera.lt/?fen=8/1P6/8/8/8/8/Q1Q5/K1k1n3&moves=5}\\{\sffamily\tiny\color{ink}scannen zum Lösen}
\end{minipage}
\vspace{2mm}

{\raggedright\noindent\textcolor{ink}{\sffamily\textsc{Themen}}\enspace \themetag{Umwandlung} \themetag{Switchback} \themetag{Phönix}\par\smallskip{\footnotesize\textit{Hinweis:} die Partei am Zug steht im Diagramm im Schach; die Lösung beginnt mit einem Schlag.}\par}

\clearpage
\phantomsection\label{hm:KQQPvkr}
\begin{tcolorbox}[enhanced,colback=ink,colframe=ink,arc=0pt,boxrule=0pt,left=4mm,right=4mm,top=2.5mm,bottom=2.5mm]
\color{white}\sffamily{\LARGE\bfseries KQQPvkr}\hfill{\Large h\#5}\\[0.5mm]
{\small 6 Steine\quad Schwarz am Zug}\hfill{\small Nr.~778\qquad\href{https://helpman.komtera.lt/?fen=8/8/8/1r6/7Q/3P2k1/8/K4Q2&moves=5}{\color{white}\bfseries lösen\,\ensuremath{\nearrow}}}
\end{tcolorbox}
\vspace{2mm}
\noindent\begin{minipage}[t]{36mm}\vspace{0pt}\sffamily\small\raggedright
\hmfigure{längste eindeutige}{{\Large h\#5}}\hmfigure{eindeutig bei dieser Länge}{{\Large 487} Stellungen}\hmfigure{längstes Matt}{{\Large h\#6}}\hmfigure{in der Klasse}{45,4 Milliarden Stellungen,\\ 11,2 Milliarden mit Hilfsmatt}
\end{minipage}\hfill
\begin{minipage}[t]{80mm}\vspace{0pt}\begin{tikzpicture}\begin{axis}[width=72mm,height=42mm,ybar,bar width=1.6pt,xmin=-1,xmax=13,ymin=0,ymax=4.128,axis y line*=left,axis x line*=bottom,xlabel={\scriptsize Halbzüge bis Matt},ylabel={\scriptsize Stellungen (Mrd.)},tick label style={font=\tiny},label style={font=\scriptsize},enlarge x limits=false,ylabel shift=-3pt,xlabel shift=-3pt]\addplot[fill=ink!55,draw=none] coordinates {(0,0.2135) (1,0.6661) (2,3.108) (3,2.899) (4,3.753) (5,0.3468) (6,0.1519) (7,1.134e-05) (8,0.05543) (9,3.67e-07) (10,0.003556) (11,4e-09) (12,6.531e-06)};\draw[densely dotted] (axis cs:10,0) -- (axis cs:10,4.128);\end{axis}\begin{axis}[width=72mm,height=42mm,axis y line*=right,axis x line=none,ymode=log,xmin=-1,xmax=13,ylabel={\scriptsize eindeutig (log)},tick label style={font=\tiny},label style={font=\scriptsize},ylabel shift=-3pt]\addplot[mark=*,mark size=0.7pt,draw=black,thin] coordinates {(0,213540070) (1,379557870) (2,1046749856) (3,79192019) (4,13308003) (5,1823) (6,2980486) (7,24) (8,190135) (10,487)};\end{axis}\end{tikzpicture}\end{minipage}
\vspace{3mm}

\noindent\begin{minipage}[t]{60mm}\vspace{0pt}\sbox\hmboard{\chessboard[setfen=8/8/8/1r6/7Q/3P2k1/8/K4Q2 b - - 0 1]}\usebox\hmboard\\[1mm]
\makebox[\wd\hmboard]{\sffamily\hspace{3mm}\textbf{h\#5}\hfill(4+2)}\end{minipage}\hfill
\begin{minipage}[t]{52mm}\vspace{0pt}
\textcolor{ink}{\sffamily\textsc{Lösung}}\par\smallskip\small
\begin{tabular}{@{}r@{\enspace}l@{\enspace}l@{}}
\textbf{1.} & \symking{}xh4 & d4\\
\textbf{2.} & \symking{}g3 & \symqueen{}xb5\\
\textbf{3.} & \symking{}f2 & \symking{}b2\\
\textbf{4.} & \symking{}e1 & \symking{}c3\\
\textbf{5.} & \symking{}d1 & \symqueen{}f1\#\\
\end{tabular}\par\vspace{3mm}
\qrcode[height=12mm]{https://helpman.komtera.lt/?fen=8/8/8/1r6/7Q/3P2k1/8/K4Q2&moves=5}\\{\sffamily\tiny\color{ink}scannen zum Lösen}
\end{minipage}
\vspace{2mm}

{\raggedright\noindent\textcolor{ink}{\sffamily\textsc{Themen}}\enspace \themetag{reines Matt} \themetag{Mustermatt} \themetag{Spiegelmatt} \themetag{Switchback} \themetag{ein Stein zieht} \themetag{ohne Schach}\par\smallskip{\footnotesize\textit{Hinweis:} die Partei am Zug steht im Diagramm im Schach; die Lösung beginnt mit einem Schlag.}\par}

\clearpage
\phantomsection\label{hm:KQQvkbb}
\begin{tcolorbox}[enhanced,colback=ink,colframe=ink,arc=0pt,boxrule=0pt,left=4mm,right=4mm,top=2.5mm,bottom=2.5mm]
\color{white}\sffamily{\LARGE\bfseries KQQvkbb}\hfill{\Large h\#5}\\[0.5mm]
{\small 6 Steine\quad Schwarz am Zug}\hfill{\small Nr.~779\qquad\href{https://helpman.komtera.lt/?fen=8/1Q6/8/8/7Q/6k1/2b5/K1b5&moves=5}{\color{white}\bfseries lösen\,\ensuremath{\nearrow}}}
\end{tcolorbox}
\vspace{2mm}
\noindent\begin{minipage}[t]{36mm}\vspace{0pt}\sffamily\small\raggedright
\hmfigure{längste eindeutige}{{\Large h\#5}}\hmfigure{eindeutig bei dieser Länge}{{\Large 8\,228} Stellungen}\hmfigure{längstes Matt}{{\Large h\#7}}\hmfigure{in der Klasse}{15,5 Milliarden Stellungen,\\ 1,81 Milliarden mit Hilfsmatt}
\end{minipage}\hfill
\begin{minipage}[t]{80mm}\vspace{0pt}\begin{tikzpicture}\begin{axis}[width=72mm,height=42mm,ybar,bar width=1.6pt,xmin=-1,xmax=15,ymin=0,ymax=680.6,axis y line*=left,axis x line*=bottom,xlabel={\scriptsize Halbzüge bis Matt},ylabel={\scriptsize Stellungen (Mio.)},tick label style={font=\tiny},label style={font=\scriptsize},enlarge x limits=false,ylabel shift=-3pt,xlabel shift=-3pt]\addplot[fill=ink!55,draw=none] coordinates {(0,29.34) (1,93.89) (2,451.4) (3,502) (4,618.7) (5,71.2) (6,17.68) (7,0.001456) (8,17.13) (9,0.000863) (10,6.281) (12,0.514) (14,0.000835)};\draw[densely dotted] (axis cs:10,0) -- (axis cs:10,680.6);\end{axis}\begin{axis}[width=72mm,height=42mm,axis y line*=right,axis x line=none,ymode=log,xmin=-1,xmax=15,ylabel={\scriptsize eindeutig (log)},tick label style={font=\tiny},label style={font=\scriptsize},ylabel shift=-3pt]\addplot[mark=*,mark size=0.7pt,draw=black,thin] coordinates {(0,29343248) (1,55699674) (2,155146129) (3,15335009) (4,1300104) (5,53) (6,656096) (7,357) (8,246300) (9,3) (10,8228)};\end{axis}\end{tikzpicture}\end{minipage}
\vspace{3mm}

\noindent\begin{minipage}[t]{60mm}\vspace{0pt}\sbox\hmboard{\chessboard[setfen=8/1Q6/8/8/7Q/6k1/2b5/K1b5 b - - 0 1]}\usebox\hmboard\\[1mm]
\makebox[\wd\hmboard]{\sffamily\hspace{3mm}\textbf{h\#5}\hfill(3+3)}\end{minipage}\hfill
\begin{minipage}[t]{52mm}\vspace{0pt}
\textcolor{ink}{\sffamily\textsc{Lösung}}\par\smallskip\small
\begin{tabular}{@{}r@{\enspace}l@{\enspace}l@{}}
\textbf{1.} & \symking{}xh4 & \symqueen{}b2\\
\textbf{2.} & \symbishop{}b1 & \symking{}xb1\\
\textbf{3.} & \symking{}g3 & \symking{}c2\\
\textbf{4.} & \symking{}f2 & \symking{}d3+\\
\textbf{5.} & \symking{}e1 & \symqueen{}e2\#\\
\end{tabular}\par\vspace{3mm}
\qrcode[height=12mm]{https://helpman.komtera.lt/?fen=8/1Q6/8/8/7Q/6k1/2b5/K1b5&moves=5}\\{\sffamily\tiny\color{ink}scannen zum Lösen}
\end{minipage}
\vspace{2mm}

{\raggedright\noindent\textcolor{ink}{\sffamily\textsc{Themen}}\enspace \themetag{Switchback}\par\smallskip{\footnotesize\textit{Hinweis:} die Partei am Zug steht im Diagramm im Schach; die Lösung beginnt mit einem Schlag.}\par}

\clearpage
\phantomsection\label{hm:KQQvkbp}
\begin{tcolorbox}[enhanced,colback=ink,colframe=ink,arc=0pt,boxrule=0pt,left=4mm,right=4mm,top=2.5mm,bottom=2.5mm]
\color{white}\sffamily{\LARGE\bfseries KQQvkbp}\hfill{\Large h\#5}\\[0.5mm]
{\small 6 Steine\quad Schwarz am Zug}\hfill{\small Nr.~780\qquad\href{https://helpman.komtera.lt/?fen=8/5Q2/8/8/8/2p5/2b4Q/K5k1&moves=5}{\color{white}\bfseries lösen\,\ensuremath{\nearrow}}}
\end{tcolorbox}
\vspace{2mm}
\noindent\begin{minipage}[t]{36mm}\vspace{0pt}\sffamily\small\raggedright
\hmfigure{längste eindeutige}{{\Large h\#5}}\hmfigure{eindeutig bei dieser Länge}{{\Large 46\,878} Stellungen}\hmfigure{längstes Matt}{{\Large h\#7}}\hmfigure{in der Klasse}{45,4 Milliarden Stellungen,\\ 11,7 Milliarden mit Hilfsmatt}
\end{minipage}\hfill
\begin{minipage}[t]{80mm}\vspace{0pt}\begin{tikzpicture}\begin{axis}[width=72mm,height=42mm,ybar,bar width=1.6pt,xmin=-1,xmax=15,ymin=0,ymax=4.472,axis y line*=left,axis x line*=bottom,xlabel={\scriptsize Halbzüge bis Matt},ylabel={\scriptsize Stellungen (Mrd.)},tick label style={font=\tiny},label style={font=\scriptsize},enlarge x limits=false,ylabel shift=-3pt,xlabel shift=-3pt]\addplot[fill=ink!55,draw=none] coordinates {(0,0.2254) (1,0.7021) (2,3.132) (3,3.045) (4,4.065) (5,0.2851) (6,0.1197) (7,3.811e-06) (8,0.113) (9,1.87e-07) (10,0.03473) (11,1e-08) (12,0.002634) (14,1.28e-07)};\draw[densely dotted] (axis cs:10,0) -- (axis cs:10,4.472);\end{axis}\begin{axis}[width=72mm,height=42mm,axis y line*=right,axis x line=none,ymode=log,xmin=-1,xmax=15,ylabel={\scriptsize eindeutig (log)},tick label style={font=\tiny},label style={font=\scriptsize},ylabel shift=-3pt]\addplot[mark=*,mark size=0.7pt,draw=black,thin] coordinates {(0,225350678) (1,393099590) (2,1043827201) (3,71076164) (4,8998991) (5,386) (6,4611299) (7,141) (8,1346793) (9,2) (10,46878)};\end{axis}\end{tikzpicture}\end{minipage}
\vspace{3mm}

\noindent\begin{minipage}[t]{60mm}\vspace{0pt}\sbox\hmboard{\chessboard[setfen=8/5Q2/8/8/8/2p5/2b4Q/K5k1 b - - 0 1]}\usebox\hmboard\\[1mm]
\makebox[\wd\hmboard]{\sffamily\hspace{3mm}\textbf{h\#5}\hfill(3+3)}\end{minipage}\hfill
\begin{minipage}[t]{52mm}\vspace{0pt}
\textcolor{ink}{\sffamily\textsc{Lösung}}\par\smallskip\small
\begin{tabular}{@{}r@{\enspace}l@{\enspace}l@{}}
\textbf{1.} & \symking{}xh2 & \symqueen{}a7\\
\textbf{2.} & \symbishop{}b1 & \symking{}xb1\\
\textbf{3.} & \symking{}g2 & \symking{}c2\\
\textbf{4.} & \symking{}f1 & \symking{}d3\\
\textbf{5.} & \symking{}e1 & \symqueen{}g1\#\\
\end{tabular}\par\vspace{3mm}
\qrcode[height=12mm]{https://helpman.komtera.lt/?fen=8/5Q2/8/8/8/2p5/2b4Q/K5k1&moves=5}\\{\sffamily\tiny\color{ink}scannen zum Lösen}
\end{minipage}
\vspace{2mm}

{\raggedright\noindent\textcolor{ink}{\sffamily\textsc{Themen}}\enspace \themetag{reines Matt} \themetag{Mustermatt} \themetag{Spiegelmatt} \themetag{ohne Schach}\par\smallskip{\footnotesize\textit{Hinweis:} die Partei am Zug steht im Diagramm im Schach; die Lösung beginnt mit einem Schlag.}\par}

\clearpage
\phantomsection\label{hm:KQQvkqr}
\begin{tcolorbox}[enhanced,colback=ink,colframe=ink,arc=0pt,boxrule=0pt,left=4mm,right=4mm,top=2.5mm,bottom=2.5mm]
\color{white}\sffamily{\LARGE\bfseries KQQvkqr}\hfill{\Large h\#5}\\[0.5mm]
{\small 6 Steine\quad Schwarz am Zug}\hfill{\small Nr.~781\qquad\href{https://helpman.komtera.lt/?fen=8/1q6/8/8/8/1r6/Q7/KQ5k&moves=5}{\color{white}\bfseries lösen\,\ensuremath{\nearrow}}}
\end{tcolorbox}
\vspace{2mm}
\noindent\begin{minipage}[t]{36mm}\vspace{0pt}\sffamily\small\raggedright
\hmfigure{längste eindeutige}{{\Large h\#5}}\hmfigure{eindeutig bei dieser Länge}{{\Large 10} Stellungen}\hmfigure{längstes Matt}{{\Large h\#7}}\hmfigure{in der Klasse}{15,5 Milliarden Stellungen,\\ 2,93 Milliarden mit Hilfsmatt}
\end{minipage}\hfill
\begin{minipage}[t]{80mm}\vspace{0pt}\begin{tikzpicture}\begin{axis}[width=72mm,height=42mm,ybar,bar width=1.6pt,xmin=-1,xmax=15,ymin=0,ymax=1038,axis y line*=left,axis x line*=bottom,xlabel={\scriptsize Halbzüge bis Matt},ylabel={\scriptsize Stellungen (Mio.)},tick label style={font=\tiny},label style={font=\scriptsize},enlarge x limits=false,ylabel shift=-3pt,xlabel shift=-3pt]\addplot[fill=ink!55,draw=none] coordinates {(0,32.85) (1,107.7) (2,590.2) (3,943.7) (4,923.7) (5,274.9) (6,34.1) (7,0.03113) (8,23.01) (9,0.005256) (10,1.87) (11,0.000272) (12,0.000225) (13,1.8e-05) (14,3.9e-05)};\draw[densely dotted] (axis cs:10,0) -- (axis cs:10,1038);\end{axis}\begin{axis}[width=72mm,height=42mm,axis y line*=right,axis x line=none,ymode=log,xmin=-1,xmax=15,ylabel={\scriptsize eindeutig (log)},tick label style={font=\tiny},label style={font=\scriptsize},ylabel shift=-3pt]\addplot[mark=*,mark size=0.7pt,draw=black,thin] coordinates {(0,32850986) (1,71311109) (2,224593378) (3,53189408) (4,2247991) (5,8003) (6,649917) (7,915) (8,63998) (9,1) (10,10)};\end{axis}\end{tikzpicture}\end{minipage}
\vspace{3mm}

\noindent\begin{minipage}[t]{60mm}\vspace{0pt}\sbox\hmboard{\chessboard[setfen=8/1q6/8/8/8/1r6/Q7/KQ5k b - - 0 1]}\usebox\hmboard\\[1mm]
\makebox[\wd\hmboard]{\sffamily\hspace{3mm}\textbf{h\#5}\hfill(3+3)}\end{minipage}\hfill
\begin{minipage}[t]{52mm}\vspace{0pt}
\textcolor{ink}{\sffamily\textsc{Lösung}}\par\smallskip\small
\begin{tabular}{@{}r@{\enspace}l@{\enspace}l@{}}
\textbf{1.} & \symrook{}xb1+ & \symqueen{}xb1+\\
\textbf{2.} & \symking{}g2 & \symqueen{}b5\\
\textbf{3.} & \symking{}f2 & \symking{}b2\\
\textbf{4.} & \symking{}e1 & \symking{}c3\\
\textbf{5.} & \symking{}d1 & \symqueen{}f1\#\\
\end{tabular}\par\vspace{3mm}
\qrcode[height=12mm]{https://helpman.komtera.lt/?fen=8/1q6/8/8/8/1r6/Q7/KQ5k&moves=5}\\{\sffamily\tiny\color{ink}scannen zum Lösen}
\end{minipage}
\vspace{2mm}

{\raggedright\noindent\textcolor{ink}{\sffamily\textsc{Themen}}\enspace \themetag{reines Matt} \themetag{Mustermatt} \themetag{Spiegelmatt}\par\smallskip{\footnotesize\textit{Hinweis:} die Partei am Zug steht im Diagramm im Schach; die Lösung beginnt mit einem Schlag.}\par}

\clearpage
\phantomsection\label{hm:KQQvkrb}
\begin{tcolorbox}[enhanced,colback=ink,colframe=ink,arc=0pt,boxrule=0pt,left=4mm,right=4mm,top=2.5mm,bottom=2.5mm]
\color{white}\sffamily{\LARGE\bfseries KQQvkrb}\hfill{\Large h\#5}\\[0.5mm]
{\small 6 Steine\quad Schwarz am Zug}\hfill{\small Nr.~782\qquad\href{https://helpman.komtera.lt/?fen=8/1r6/8/8/2b4Q/6k1/8/K4Q2&moves=5}{\color{white}\bfseries lösen\,\ensuremath{\nearrow}}}
\end{tcolorbox}
\vspace{2mm}
\noindent\begin{minipage}[t]{36mm}\vspace{0pt}\sffamily\small\raggedright
\hmfigure{längste eindeutige}{{\Large h\#5}}\hmfigure{eindeutig bei dieser Länge}{{\Large 292} Stellungen}\hmfigure{längstes Matt}{{\Large h\#6.5}}\hmfigure{in der Klasse}{15,5 Milliarden Stellungen,\\ 3,45 Milliarden mit Hilfsmatt}
\end{minipage}\hfill
\begin{minipage}[t]{80mm}\vspace{0pt}\begin{tikzpicture}\begin{axis}[width=72mm,height=42mm,ybar,bar width=1.6pt,xmin=-1,xmax=14,ymin=0,ymax=1.26,axis y line*=left,axis x line*=bottom,xlabel={\scriptsize Halbzüge bis Matt},ylabel={\scriptsize Stellungen (Mrd.)},tick label style={font=\tiny},label style={font=\scriptsize},enlarge x limits=false,ylabel shift=-3pt,xlabel shift=-3pt]\addplot[fill=ink!55,draw=none] coordinates {(0,0.05219) (1,0.1695) (2,0.8484) (3,0.9937) (4,1.146) (5,0.1706) (6,0.04144) (7,2.841e-06) (8,0.02916) (9,2.43e-07) (10,0.003305) (11,2.8e-08) (12,3.51e-06) (13,1e-09)};\draw[densely dotted] (axis cs:10,0) -- (axis cs:10,1.26);\end{axis}\begin{axis}[width=72mm,height=42mm,axis y line*=right,axis x line=none,ymode=log,xmin=-1,xmax=14,ylabel={\scriptsize eindeutig (log)},tick label style={font=\tiny},label style={font=\scriptsize},ylabel shift=-3pt]\addplot[mark=*,mark size=0.7pt,draw=black,thin] coordinates {(0,52189020) (1,104150655) (2,304029076) (3,33797861) (4,2563400) (5,394) (6,890225) (7,24) (8,140910) (10,292)};\end{axis}\end{tikzpicture}\end{minipage}
\vspace{3mm}

\noindent\begin{minipage}[t]{60mm}\vspace{0pt}\sbox\hmboard{\chessboard[setfen=8/1r6/8/8/2b4Q/6k1/8/K4Q2 b - - 0 1]}\usebox\hmboard\\[1mm]
\makebox[\wd\hmboard]{\sffamily\hspace{3mm}\textbf{h\#5}\hfill(3+3)}\end{minipage}\hfill
\begin{minipage}[t]{52mm}\vspace{0pt}
\textcolor{ink}{\sffamily\textsc{Lösung}}\par\smallskip\small
\begin{tabular}{@{}r@{\enspace}l@{\enspace}l@{}}
\textbf{1.} & \symking{}xh4 & \symqueen{}xc4+\\
\textbf{2.} & \symking{}g3 & \symqueen{}b5\\
\textbf{3.} & \symking{}f2 & \symking{}b2\\
\textbf{4.} & \symking{}e1 & \symking{}c3\\
\textbf{5.} & \symking{}d1 & \symqueen{}f1\#\\
\end{tabular}\par\vspace{3mm}
\qrcode[height=12mm]{https://helpman.komtera.lt/?fen=8/1r6/8/8/2b4Q/6k1/8/K4Q2&moves=5}\\{\sffamily\tiny\color{ink}scannen zum Lösen}
\end{minipage}
\vspace{2mm}

{\raggedright\noindent\textcolor{ink}{\sffamily\textsc{Themen}}\enspace \themetag{reines Matt} \themetag{Mustermatt} \themetag{Spiegelmatt} \themetag{Switchback} \themetag{Rundlauf} \themetag{ein Stein zieht}\par\smallskip{\footnotesize\textit{Hinweis:} die Partei am Zug steht im Diagramm im Schach; die Lösung beginnt mit einem Schlag.}\par}

\clearpage
\phantomsection\label{hm:KQQvkrn}
\begin{tcolorbox}[enhanced,colback=ink,colframe=ink,arc=0pt,boxrule=0pt,left=4mm,right=4mm,top=2.5mm,bottom=2.5mm]
\color{white}\sffamily{\LARGE\bfseries KQQvkrn}\hfill{\Large h\#5}\\[0.5mm]
{\small 6 Steine\quad Schwarz am Zug}\hfill{\small Nr.~783\qquad\href{https://helpman.komtera.lt/?fen=8/8/8/8/7Q/3r2k1/7n/K4Q2&moves=5}{\color{white}\bfseries lösen\,\ensuremath{\nearrow}}}
\end{tcolorbox}
\vspace{2mm}
\noindent\begin{minipage}[t]{36mm}\vspace{0pt}\sffamily\small\raggedright
\hmfigure{längste eindeutige}{{\Large h\#5}}\hmfigure{eindeutig bei dieser Länge}{{\Large 6\,726} Stellungen}\hmfigure{längstes Matt}{{\Large h\#6}}\hmfigure{in der Klasse}{15,5 Milliarden Stellungen,\\ 3,56 Milliarden mit Hilfsmatt}
\end{minipage}\hfill
\begin{minipage}[t]{80mm}\vspace{0pt}\begin{tikzpicture}\begin{axis}[width=72mm,height=42mm,ybar,bar width=1.6pt,xmin=-1,xmax=13,ymin=0,ymax=1.342,axis y line*=left,axis x line*=bottom,xlabel={\scriptsize Halbzüge bis Matt},ylabel={\scriptsize Stellungen (Mrd.)},tick label style={font=\tiny},label style={font=\scriptsize},enlarge x limits=false,ylabel shift=-3pt,xlabel shift=-3pt]\addplot[fill=ink!55,draw=none] coordinates {(0,0.05658) (1,0.1792) (2,0.871) (3,0.9919) (4,1.22) (5,0.1628) (6,0.03857) (7,3.713e-06) (8,0.03379) (9,1.009e-06) (10,0.007954) (11,8e-09) (12,0.0001097)};\draw[densely dotted] (axis cs:10,0) -- (axis cs:10,1.342);\end{axis}\begin{axis}[width=72mm,height=42mm,axis y line*=right,axis x line=none,ymode=log,xmin=-1,xmax=13,ylabel={\scriptsize eindeutig (log)},tick label style={font=\tiny},label style={font=\scriptsize},ylabel shift=-3pt]\addplot[mark=*,mark size=0.7pt,draw=black,thin] coordinates {(0,56579851) (1,108554185) (2,307438986) (3,30455910) (4,2419435) (5,268) (6,1079549) (7,315) (8,286940) (10,6726)};\end{axis}\end{tikzpicture}\end{minipage}
\vspace{3mm}

\noindent\begin{minipage}[t]{60mm}\vspace{0pt}\sbox\hmboard{\chessboard[setfen=8/8/8/8/7Q/3r2k1/7n/K4Q2 b - - 0 1]}\usebox\hmboard\\[1mm]
\makebox[\wd\hmboard]{\sffamily\hspace{3mm}\textbf{h\#5}\hfill(3+3)}\end{minipage}\hfill
\begin{minipage}[t]{52mm}\vspace{0pt}
\textcolor{ink}{\sffamily\textsc{Lösung}}\par\smallskip\small
\begin{tabular}{@{}r@{\enspace}l@{\enspace}l@{}}
\textbf{1.} & \symking{}xh4 & \symking{}b2\\
\textbf{2.} & \symking{}g3 & \symqueen{}xd3+\\
\textbf{3.} & \symking{}f2 & \symking{}c3\\
\textbf{4.} & \symking{}e1 & \symqueen{}e3+\\
\textbf{5.} & \symking{}d1 & \symqueen{}d2\#\\
\end{tabular}\par\vspace{3mm}
\qrcode[height=12mm]{https://helpman.komtera.lt/?fen=8/8/8/8/7Q/3r2k1/7n/K4Q2&moves=5}\\{\sffamily\tiny\color{ink}scannen zum Lösen}
\end{minipage}
\vspace{2mm}

{\raggedright\noindent\textcolor{ink}{\sffamily\textsc{Themen}}\enspace \themetag{Switchback} \themetag{ein Stein zieht}\par\smallskip{\footnotesize\textit{Hinweis:} die Partei am Zug steht im Diagramm im Schach; die Lösung beginnt mit einem Schlag.}\par}

\clearpage
\phantomsection\label{hm:KQQvkrp}
\begin{tcolorbox}[enhanced,colback=ink,colframe=ink,arc=0pt,boxrule=0pt,left=4mm,right=4mm,top=2.5mm,bottom=2.5mm]
\color{white}\sffamily{\LARGE\bfseries KQQvkrp}\hfill{\Large h\#5}\\[0.5mm]
{\small 6 Steine\quad Schwarz am Zug}\hfill{\small Nr.~784\qquad\href{https://helpman.komtera.lt/?fen=8/8/8/1r6/2p4Q/6k1/8/K4Q2&moves=5}{\color{white}\bfseries lösen\,\ensuremath{\nearrow}}}
\end{tcolorbox}
\vspace{2mm}
\noindent\begin{minipage}[t]{36mm}\vspace{0pt}\sffamily\small\raggedright
\hmfigure{längste eindeutige}{{\Large h\#5}}\hmfigure{eindeutig bei dieser Länge}{{\Large 5\,256} Stellungen}\hmfigure{längstes Matt}{{\Large h\#6}}\hmfigure{in der Klasse}{45,4 Milliarden Stellungen,\\ 11,1 Milliarden mit Hilfsmatt}
\end{minipage}\hfill
\begin{minipage}[t]{80mm}\vspace{0pt}\begin{tikzpicture}\begin{axis}[width=72mm,height=42mm,ybar,bar width=1.6pt,xmin=-1,xmax=13,ymin=0,ymax=4.081,axis y line*=left,axis x line*=bottom,xlabel={\scriptsize Halbzüge bis Matt},ylabel={\scriptsize Stellungen (Mrd.)},tick label style={font=\tiny},label style={font=\scriptsize},enlarge x limits=false,ylabel shift=-3pt,xlabel shift=-3pt]\addplot[fill=ink!55,draw=none] coordinates {(0,0.1974) (1,0.6305) (2,2.94) (3,3.016) (4,3.71) (5,0.3852) (6,0.131) (7,2.78e-06) (8,0.09382) (9,3.97e-07) (10,0.01419) (11,2.4e-08) (12,8.291e-05)};\draw[densely dotted] (axis cs:10,0) -- (axis cs:10,4.081);\end{axis}\begin{axis}[width=72mm,height=42mm,axis y line*=right,axis x line=none,ymode=log,xmin=-1,xmax=13,ylabel={\scriptsize eindeutig (log)},tick label style={font=\tiny},label style={font=\scriptsize},ylabel shift=-3pt]\addplot[mark=*,mark size=0.7pt,draw=black,thin] coordinates {(0,197362580) (1,368564703) (2,1017081845) (3,84319818) (4,8551612) (5,666) (6,2889975) (7,45) (8,551911) (10,5256)};\end{axis}\end{tikzpicture}\end{minipage}
\vspace{3mm}

\noindent\begin{minipage}[t]{60mm}\vspace{0pt}\sbox\hmboard{\chessboard[setfen=8/8/8/1r6/2p4Q/6k1/8/K4Q2 b - - 0 1]}\usebox\hmboard\\[1mm]
\makebox[\wd\hmboard]{\sffamily\hspace{3mm}\textbf{h\#5}\hfill(3+3)}\end{minipage}\hfill
\begin{minipage}[t]{52mm}\vspace{0pt}
\textcolor{ink}{\sffamily\textsc{Lösung}}\par\smallskip\small
\begin{tabular}{@{}r@{\enspace}l@{\enspace}l@{}}
\textbf{1.} & \symking{}xh4 & \symqueen{}xc4+\\
\textbf{2.} & \symking{}g3 & \symqueen{}xb5\\
\textbf{3.} & \symking{}f2 & \symking{}b2\\
\textbf{4.} & \symking{}e1 & \symking{}c3\\
\textbf{5.} & \symking{}d1 & \symqueen{}f1\#\\
\end{tabular}\par\vspace{3mm}
\qrcode[height=12mm]{https://helpman.komtera.lt/?fen=8/8/8/1r6/2p4Q/6k1/8/K4Q2&moves=5}\\{\sffamily\tiny\color{ink}scannen zum Lösen}
\end{minipage}
\vspace{2mm}

{\raggedright\noindent\textcolor{ink}{\sffamily\textsc{Themen}}\enspace \themetag{reines Matt} \themetag{Mustermatt} \themetag{Idealmatt} \themetag{Spiegelmatt} \themetag{Switchback} \themetag{Rundlauf} \themetag{ein Stein zieht}\par\smallskip{\footnotesize\textit{Hinweis:} die Partei am Zug steht im Diagramm im Schach; die Lösung beginnt mit einem Schlag.}\par}

\clearpage
\phantomsection\label{hm:KQQvkrr}
\begin{tcolorbox}[enhanced,colback=ink,colframe=ink,arc=0pt,boxrule=0pt,left=4mm,right=4mm,top=2.5mm,bottom=2.5mm]
\color{white}\sffamily{\LARGE\bfseries KQQvkrr}\hfill{\Large h\#5}\\[0.5mm]
{\small 6 Steine\quad Schwarz am Zug}\hfill{\small Nr.~785\qquad\href{https://helpman.komtera.lt/?fen=8/8/8/5Q2/4k3/1rr5/3Q4/K7&moves=5}{\color{white}\bfseries lösen\,\ensuremath{\nearrow}}}
\end{tcolorbox}
\vspace{2mm}
\noindent\begin{minipage}[t]{36mm}\vspace{0pt}\sffamily\small\raggedright
\hmfigure{längste eindeutige}{{\Large h\#5}}\hmfigure{eindeutig bei dieser Länge}{{\Large 42} Stellungen}\hmfigure{längstes Matt}{{\Large h\#6}}\hmfigure{in der Klasse}{15,5 Milliarden Stellungen,\\ 1,62 Milliarden mit Hilfsmatt}
\end{minipage}\hfill
\begin{minipage}[t]{80mm}\vspace{0pt}\begin{tikzpicture}\begin{axis}[width=72mm,height=42mm,ybar,bar width=1.6pt,xmin=-1,xmax=13,ymin=0,ymax=570,axis y line*=left,axis x line*=bottom,xlabel={\scriptsize Halbzüge bis Matt},ylabel={\scriptsize Stellungen (Mio.)},tick label style={font=\tiny},label style={font=\scriptsize},enlarge x limits=false,ylabel shift=-3pt,xlabel shift=-3pt]\addplot[fill=ink!55,draw=none] coordinates {(0,22.25) (1,73.56) (2,382.8) (3,491.1) (4,518.2) (5,101.2) (6,20.41) (7,0.01154) (8,10.53) (9,0.000291) (10,1.005) (11,1.6e-05) (12,0.000525)};\draw[densely dotted] (axis cs:10,0) -- (axis cs:10,570);\end{axis}\begin{axis}[width=72mm,height=42mm,axis y line*=right,axis x line=none,ymode=log,xmin=-1,xmax=13,ylabel={\scriptsize eindeutig (log)},tick label style={font=\tiny},label style={font=\scriptsize},ylabel shift=-3pt]\addplot[mark=*,mark size=0.7pt,draw=black,thin] coordinates {(0,22248838) (1,46644081) (2,140988938) (3,20147348) (4,1169644) (5,2110) (6,285960) (7,7) (8,29916) (10,42)};\end{axis}\end{tikzpicture}\end{minipage}
\vspace{3mm}

\noindent\begin{minipage}[t]{60mm}\vspace{0pt}\sbox\hmboard{\chessboard[setfen=8/8/8/5Q2/4k3/1rr5/3Q4/K7 b - - 0 1]}\usebox\hmboard\\[1mm]
\makebox[\wd\hmboard]{\sffamily\hspace{3mm}\textbf{h\#5}\hfill(3+3)}\end{minipage}\hfill
\begin{minipage}[t]{52mm}\vspace{0pt}
\textcolor{ink}{\sffamily\textsc{Lösung}}\par\smallskip\small
\begin{tabular}{@{}r@{\enspace}l@{\enspace}l@{}}
\textbf{1.} & \symking{}xf5 & \symqueen{}xc3\\
\textbf{2.} & \symking{}e4 & \symqueen{}e5+\\
\textbf{3.} & \symking{}d3 & \symking{}a2\\
\textbf{4.} & \symking{}d2 & \symking{}xb3\\
\textbf{5.} & \symking{}c1 & \symqueen{}e1\#\\
\end{tabular}\par\vspace{3mm}
\qrcode[height=12mm]{https://helpman.komtera.lt/?fen=8/8/8/5Q2/4k3/1rr5/3Q4/K7&moves=5}\\{\sffamily\tiny\color{ink}scannen zum Lösen}
\end{minipage}
\vspace{2mm}

{\raggedright\noindent\textcolor{ink}{\sffamily\textsc{Themen}}\enspace \themetag{reines Matt} \themetag{Mustermatt} \themetag{Idealmatt} \themetag{Spiegelmatt} \themetag{Switchback} \themetag{ein Stein zieht}\par\smallskip{\footnotesize\textit{Hinweis:} die Partei am Zug steht im Diagramm im Schach; die Lösung beginnt mit einem Schlag.}\par}

\clearpage
\phantomsection\label{hm:KQRBNvk}
\begin{tcolorbox}[enhanced,colback=ink,colframe=ink,arc=0pt,boxrule=0pt,left=4mm,right=4mm,top=2.5mm,bottom=2.5mm]
\color{white}\sffamily{\LARGE\bfseries KQRBNvk}\hfill{\Large h\#5}\\[0.5mm]
{\small 6 Steine\quad Schwarz am Zug}\hfill{\small Nr.~786\qquad\href{https://helpman.komtera.lt/?fen=Q1B5/kR6/8/8/8/8/8/1K5N&moves=5}{\color{white}\bfseries lösen\,\ensuremath{\nearrow}}}
\end{tcolorbox}
\vspace{2mm}
\noindent\begin{minipage}[t]{36mm}\vspace{0pt}\sffamily\small\raggedright
\hmfigure{längste eindeutige}{{\Large h\#5}}\hmfigure{eindeutig bei dieser Länge}{{\Large 16} Stellungen}\hmfigure{längstes Matt}{{\Large h\#7}}\hmfigure{in der Klasse}{15,5 Milliarden Stellungen,\\ 8,5 Milliarden mit Hilfsmatt}
\end{minipage}\hfill
\begin{minipage}[t]{80mm}\vspace{0pt}\begin{tikzpicture}\begin{axis}[width=72mm,height=42mm,ybar,bar width=1.6pt,xmin=-1,xmax=15,ymin=0,ymax=4.058,axis y line*=left,axis x line*=bottom,xlabel={\scriptsize Halbzüge bis Matt},ylabel={\scriptsize Stellungen (Mrd.)},tick label style={font=\tiny},label style={font=\scriptsize},enlarge x limits=false,ylabel shift=-3pt,xlabel shift=-3pt]\addplot[fill=ink!55,draw=none] coordinates {(0,0.178) (1,0.5427) (2,1.92) (3,1.914) (4,3.69) (5,0.1228) (6,0.1116) (7,1.25e-07) (8,0.02376) (10,0.0005973) (12,4.04e-07) (14,3e-09)};\draw[densely dotted] (axis cs:10,0) -- (axis cs:10,4.058);\end{axis}\begin{axis}[width=72mm,height=42mm,axis y line*=right,axis x line=none,ymode=log,xmin=-1,xmax=15,ylabel={\scriptsize eindeutig (log)},tick label style={font=\tiny},label style={font=\scriptsize},ylabel shift=-3pt]\addplot[mark=*,mark size=0.7pt,draw=black,thin] coordinates {(0,178018989) (1,321916014) (2,898569433) (3,112014059) (4,38901506) (5,10) (6,5606325) (8,23810) (10,16)};\end{axis}\end{tikzpicture}\end{minipage}
\vspace{3mm}

\noindent\begin{minipage}[t]{60mm}\vspace{0pt}\sbox\hmboard{\chessboard[setfen=Q1B5/kR6/8/8/8/8/8/1K5N b - - 0 1]}\usebox\hmboard\\[1mm]
\makebox[\wd\hmboard]{\sffamily\hspace{3mm}\textbf{h\#5}\hfill(5+1)}\end{minipage}\hfill
\begin{minipage}[t]{52mm}\vspace{0pt}
\textcolor{ink}{\sffamily\textsc{Lösung}}\par\smallskip\small
\begin{tabular}{@{}r@{\enspace}l@{\enspace}l@{}}
\textbf{1.} & \symking{}xa8 & \symrook{}a7+\\
\textbf{2.} & \symking{}b8 & \symrook{}a8+\\
\textbf{3.} & \symking{}c7 & \symknight{}g3\\
\textbf{4.} & \symking{}d8 & \symknight{}f5\\
\textbf{5.} & \symking{}e8 & \symbishop{}e6\#\\
\end{tabular}\par\vspace{3mm}
\qrcode[height=12mm]{https://helpman.komtera.lt/?fen=Q1B5/kR6/8/8/8/8/8/1K5N&moves=5}\\{\sffamily\tiny\color{ink}scannen zum Lösen}
\end{minipage}
\vspace{2mm}

{\raggedright\noindent\textcolor{ink}{\sffamily\textsc{Themen}}\enspace \themetag{reines Matt} \themetag{Mustermatt} \themetag{Spiegelmatt} \themetag{ein Stein zieht} \themetag{Umnov}\par\smallskip{\footnotesize\textit{Hinweis:} das Diagramm hat keinen legalen letzten Zug, kann also in keiner Partie entstehen.}\par}

\clearpage
\phantomsection\label{hm:KQRBvkq}
\begin{tcolorbox}[enhanced,colback=ink,colframe=ink,arc=0pt,boxrule=0pt,left=4mm,right=4mm,top=2.5mm,bottom=2.5mm]
\color{white}\sffamily{\LARGE\bfseries KQRBvkq}\hfill{\Large h\#5}\\[0.5mm]
{\small 6 Steine\quad Schwarz am Zug}\hfill{\small Nr.~787\qquad\href{https://helpman.komtera.lt/?fen=8/8/8/8/8/4B3/1QqR4/K1k5&moves=5}{\color{white}\bfseries lösen\,\ensuremath{\nearrow}}}
\end{tcolorbox}
\vspace{2mm}
\noindent\begin{minipage}[t]{36mm}\vspace{0pt}\sffamily\small\raggedright
\hmfigure{längste eindeutige}{{\Large h\#5}}\hmfigure{eindeutig bei dieser Länge}{{\Large 214} Stellungen}\hmfigure{längstes Matt}{{\Large h\#7}}\hmfigure{in der Klasse}{15,5 Milliarden Stellungen,\\ 6,87 Milliarden mit Hilfsmatt}
\end{minipage}\hfill
\begin{minipage}[t]{80mm}\vspace{0pt}\begin{tikzpicture}\begin{axis}[width=72mm,height=42mm,ybar,bar width=1.6pt,xmin=-1,xmax=15,ymin=0,ymax=3.106,axis y line*=left,axis x line*=bottom,xlabel={\scriptsize Halbzüge bis Matt},ylabel={\scriptsize Stellungen (Mrd.)},tick label style={font=\tiny},label style={font=\scriptsize},enlarge x limits=false,ylabel shift=-3pt,xlabel shift=-3pt]\addplot[fill=ink!55,draw=none] coordinates {(0,0.05817) (1,0.2191) (2,1.093) (3,1.938) (4,2.824) (5,0.6331) (6,0.07076) (7,0.0001888) (8,0.0281) (9,8.66e-07) (10,0.0007158) (11,2e-08) (12,8.57e-07) (14,4e-09)};\draw[densely dotted] (axis cs:10,0) -- (axis cs:10,3.106);\end{axis}\begin{axis}[width=72mm,height=42mm,axis y line*=right,axis x line=none,ymode=log,xmin=-1,xmax=15,ylabel={\scriptsize eindeutig (log)},tick label style={font=\tiny},label style={font=\scriptsize},ylabel shift=-3pt]\addplot[mark=*,mark size=0.7pt,draw=black,thin] coordinates {(0,58172548) (1,163619913) (2,488289268) (3,154689747) (4,17395526) (5,101240) (6,3032376) (7,123) (8,48395) (10,214)};\end{axis}\end{tikzpicture}\end{minipage}
\vspace{3mm}

\noindent\begin{minipage}[t]{60mm}\vspace{0pt}\sbox\hmboard{\chessboard[setfen=8/8/8/8/8/4B3/1QqR4/K1k5 b - - 0 1]}\usebox\hmboard\\[1mm]
\makebox[\wd\hmboard]{\sffamily\hspace{3mm}\textbf{h\#5}\hfill(4+2)}\end{minipage}\hfill
\begin{minipage}[t]{52mm}\vspace{0pt}
\textcolor{ink}{\sffamily\textsc{Lösung}}\par\smallskip\small
\begin{tabular}{@{}r@{\enspace}l@{\enspace}l@{}}
\textbf{1.} & \symqueen{}xb2+ & \symrook{}xb2+\\
\textbf{2.} & \symking{}d1 & \symking{}b1\\
\textbf{3.} & \symking{}e1 & \symking{}c2\\
\textbf{4.} & \symking{}f1 & \symking{}d3\\
\textbf{5.} & \symking{}e1 & \symrook{}b1\#\\
\end{tabular}\par\vspace{3mm}
\qrcode[height=12mm]{https://helpman.komtera.lt/?fen=8/8/8/8/8/4B3/1QqR4/K1k5&moves=5}\\{\sffamily\tiny\color{ink}scannen zum Lösen}
\end{minipage}
\vspace{2mm}

{\raggedright\noindent\textcolor{ink}{\sffamily\textsc{Themen}}\enspace \themetag{Spiegelmatt} \themetag{Switchback}\par\smallskip{\footnotesize\textit{Hinweis:} die Partei am Zug steht im Diagramm im Schach; die Lösung beginnt mit einem Schlag.}\par}

\clearpage
\phantomsection\label{hm:KQRNNvk}
\begin{tcolorbox}[enhanced,colback=ink,colframe=ink,arc=0pt,boxrule=0pt,left=4mm,right=4mm,top=2.5mm,bottom=2.5mm]
\color{white}\sffamily{\LARGE\bfseries KQRNNvk}\hfill{\Large h\#5}\\[0.5mm]
{\small 6 Steine\quad Schwarz am Zug}\hfill{\small Nr.~788\qquad\href{https://helpman.komtera.lt/?fen=8/7Q/6k1/5R2/8/8/1K6/N3N3&moves=5}{\color{white}\bfseries lösen\,\ensuremath{\nearrow}}}
\end{tcolorbox}
\vspace{2mm}
\noindent\begin{minipage}[t]{36mm}\vspace{0pt}\sffamily\small\raggedright
\hmfigure{längste eindeutige}{{\Large h\#5}}\hmfigure{eindeutig bei dieser Länge}{{\Large 11} Stellungen}\hmfigure{längstes Matt}{{\Large h\#6}}\hmfigure{in der Klasse}{15,5 Milliarden Stellungen,\\ 4,33 Milliarden mit Hilfsmatt}
\end{minipage}\hfill
\begin{minipage}[t]{80mm}\vspace{0pt}\begin{tikzpicture}\begin{axis}[width=72mm,height=42mm,ybar,bar width=1.6pt,xmin=-1,xmax=13,ymin=0,ymax=2.151,axis y line*=left,axis x line*=bottom,xlabel={\scriptsize Halbzüge bis Matt},ylabel={\scriptsize Stellungen (Mrd.)},tick label style={font=\tiny},label style={font=\scriptsize},enlarge x limits=false,ylabel shift=-3pt,xlabel shift=-3pt]\addplot[fill=ink!55,draw=none] coordinates {(0,0.07346) (1,0.2505) (2,0.8544) (3,1.019) (4,1.956) (5,0.09535) (6,0.06957) (7,1.755e-06) (8,0.01488) (10,0.0008062) (12,1.029e-06)};\draw[densely dotted] (axis cs:10,0) -- (axis cs:10,2.151);\end{axis}\begin{axis}[width=72mm,height=42mm,axis y line*=right,axis x line=none,ymode=log,xmin=-1,xmax=13,ylabel={\scriptsize eindeutig (log)},tick label style={font=\tiny},label style={font=\scriptsize},ylabel shift=-3pt]\addplot[mark=*,mark size=0.7pt,draw=black,thin] coordinates {(0,73460939) (1,155050422) (2,414791342) (3,73741011) (4,25038245) (5,639) (6,2650532) (8,20228) (10,11)};\end{axis}\end{tikzpicture}\end{minipage}
\vspace{3mm}

\noindent\begin{minipage}[t]{60mm}\vspace{0pt}\sbox\hmboard{\chessboard[setfen=8/7Q/6k1/5R2/8/8/1K6/N3N3 b - - 0 1]}\usebox\hmboard\\[1mm]
\makebox[\wd\hmboard]{\sffamily\hspace{3mm}\textbf{h\#5}\hfill(5+1)}\end{minipage}\hfill
\begin{minipage}[t]{52mm}\vspace{0pt}
\textcolor{ink}{\sffamily\textsc{Lösung}}\par\smallskip\small
\begin{tabular}{@{}r@{\enspace}l@{\enspace}l@{}}
\textbf{1.} & \symking{}xh7 & \symrook{}f1\\
\textbf{2.} & \symking{}g6 & \symknight{}f3\\
\textbf{3.} & \symking{}f5 & \symknight{}b3\\
\textbf{4.} & \symking{}e4 & \symrook{}e1+\\
\textbf{5.} & \symking{}d3 & \symknight{}e5\#\\
\end{tabular}\par\vspace{3mm}
\qrcode[height=12mm]{https://helpman.komtera.lt/?fen=8/7Q/6k1/5R2/8/8/1K6/N3N3&moves=5}\\{\sffamily\tiny\color{ink}scannen zum Lösen}
\end{minipage}
\vspace{2mm}

{\raggedright\noindent\textcolor{ink}{\sffamily\textsc{Themen}}\enspace \themetag{reines Matt} \themetag{Mustermatt} \themetag{Idealmatt} \themetag{Spiegelmatt} \themetag{Switchback} \themetag{ein Stein zieht}\par\smallskip{\footnotesize\textit{Hinweis:} die Partei am Zug steht im Diagramm im Schach; die Lösung beginnt mit einem Schlag.}\par}

\clearpage
\phantomsection\label{hm:KQRNPvk}
\begin{tcolorbox}[enhanced,colback=ink,colframe=ink,arc=0pt,boxrule=0pt,left=4mm,right=4mm,top=2.5mm,bottom=2.5mm]
\color{white}\sffamily{\LARGE\bfseries KQRNPvk}\hfill{\Large h\#5}\\[0.5mm]
{\small 6 Steine\quad Schwarz am Zug}\hfill{\small Nr.~789\qquad\href{https://helpman.komtera.lt/?fen=8/8/8/8/8/8/1P4R1/K2N2kQ&moves=5}{\color{white}\bfseries lösen\,\ensuremath{\nearrow}}}
\end{tcolorbox}
\vspace{2mm}
\noindent\begin{minipage}[t]{36mm}\vspace{0pt}\sffamily\small\raggedright
\hmfigure{längste eindeutige}{{\Large h\#5}}\hmfigure{eindeutig bei dieser Länge}{{\Large 4\,322} Stellungen}\hmfigure{längstes Matt}{{\Large h\#7}}\hmfigure{in der Klasse}{45,4 Milliarden Stellungen,\\ 26,6 Milliarden mit Hilfsmatt}
\end{minipage}\hfill
\begin{minipage}[t]{80mm}\vspace{0pt}\begin{tikzpicture}\begin{axis}[width=72mm,height=42mm,ybar,bar width=1.6pt,xmin=-1,xmax=15,ymin=0,ymax=13.74,axis y line*=left,axis x line*=bottom,xlabel={\scriptsize Halbzüge bis Matt},ylabel={\scriptsize Stellungen (Mrd.)},tick label style={font=\tiny},label style={font=\scriptsize},enlarge x limits=false,ylabel shift=-3pt,xlabel shift=-3pt]\addplot[fill=ink!55,draw=none] coordinates {(0,0.3667) (1,1.324) (2,4.305) (3,6.372) (4,12.49) (5,0.9867) (6,0.5858) (7,0.0002543) (8,0.1132) (10,0.0188) (12,0.000273) (14,1.2e-08)};\draw[densely dotted] (axis cs:10,0) -- (axis cs:10,13.74);\end{axis}\begin{axis}[width=72mm,height=42mm,axis y line*=right,axis x line=none,ymode=log,xmin=-1,xmax=15,ylabel={\scriptsize eindeutig (log)},tick label style={font=\tiny},label style={font=\scriptsize},ylabel shift=-3pt]\addplot[mark=*,mark size=0.7pt,draw=black,thin] coordinates {(0,366699926) (1,824726783) (2,2129289862) (3,623372073) (4,240057620) (5,64798) (6,15214565) (8,630675) (10,4322)};\end{axis}\end{tikzpicture}\end{minipage}
\vspace{3mm}

\noindent\begin{minipage}[t]{60mm}\vspace{0pt}\sbox\hmboard{\chessboard[setfen=8/8/8/8/8/8/1P4R1/K2N2kQ b - - 0 1]}\usebox\hmboard\\[1mm]
\makebox[\wd\hmboard]{\sffamily\hspace{3mm}\textbf{h\#5}\hfill(5+1)}\end{minipage}\hfill
\begin{minipage}[t]{52mm}\vspace{0pt}
\textcolor{ink}{\sffamily\textsc{Lösung}}\par\smallskip\small
\begin{tabular}{@{}r@{\enspace}l@{\enspace}l@{}}
\textbf{1.} & \symking{}xh1 & \symrook{}h2+\\
\textbf{2.} & \symking{}g1 & \symking{}b1\\
\textbf{3.} & \symking{}f1 & \symking{}c2\\
\textbf{4.} & \symking{}e1 & \symking{}d3\\
\textbf{5.} & \symking{}xd1 & \symrook{}h1\#\\
\end{tabular}\par\vspace{3mm}
\qrcode[height=12mm]{https://helpman.komtera.lt/?fen=8/8/8/8/8/8/1P4R1/K2N2kQ&moves=5}\\{\sffamily\tiny\color{ink}scannen zum Lösen}
\end{minipage}
\vspace{2mm}

{\raggedright\noindent\textcolor{ink}{\sffamily\textsc{Themen}}\enspace \themetag{reines Matt} \themetag{Mustermatt} \themetag{Spiegelmatt} \themetag{Switchback} \themetag{ein Stein zieht} \themetag{Kniest}\par\smallskip{\footnotesize\textit{Hinweis:} das Diagramm hat keinen legalen letzten Zug, kann also in keiner Partie entstehen.}\par}

\clearpage
\phantomsection\label{hm:KQRNvkb}
\begin{tcolorbox}[enhanced,colback=ink,colframe=ink,arc=0pt,boxrule=0pt,left=4mm,right=4mm,top=2.5mm,bottom=2.5mm]
\color{white}\sffamily{\LARGE\bfseries KQRNvkb}\hfill{\Large h\#5}\\[0.5mm]
{\small 6 Steine\quad Schwarz am Zug}\hfill{\small Nr.~790\qquad\href{https://helpman.komtera.lt/?fen=8/8/8/6Qk/8/N5b1/8/K4R2&moves=5}{\color{white}\bfseries lösen\,\ensuremath{\nearrow}}}
\end{tcolorbox}
\vspace{2mm}
\noindent\begin{minipage}[t]{36mm}\vspace{0pt}\sffamily\small\raggedright
\hmfigure{längste eindeutige}{{\Large h\#5}}\hmfigure{eindeutig bei dieser Länge}{{\Large 7} Stellungen}\hmfigure{längstes Matt}{{\Large h\#6}}\hmfigure{in der Klasse}{15,5 Milliarden Stellungen,\\ 8,23 Milliarden mit Hilfsmatt}
\end{minipage}\hfill
\begin{minipage}[t]{80mm}\vspace{0pt}\begin{tikzpicture}\begin{axis}[width=72mm,height=42mm,ybar,bar width=1.6pt,xmin=-1,xmax=13,ymin=0,ymax=4.139,axis y line*=left,axis x line*=bottom,xlabel={\scriptsize Halbzüge bis Matt},ylabel={\scriptsize Stellungen (Mrd.)},tick label style={font=\tiny},label style={font=\scriptsize},enlarge x limits=false,ylabel shift=-3pt,xlabel shift=-3pt]\addplot[fill=ink!55,draw=none] coordinates {(0,0.0785) (1,0.3085) (2,1.265) (3,2.176) (4,3.763) (5,0.4768) (6,0.1195) (7,0.0001913) (8,0.04168) (9,2.54e-07) (10,0.002413) (12,9.23e-07)};\draw[densely dotted] (axis cs:10,0) -- (axis cs:10,4.139);\end{axis}\begin{axis}[width=72mm,height=42mm,axis y line*=right,axis x line=none,ymode=log,xmin=-1,xmax=13,ylabel={\scriptsize eindeutig (log)},tick label style={font=\tiny},label style={font=\scriptsize},ylabel shift=-3pt]\addplot[mark=*,mark size=0.7pt,draw=black,thin] coordinates {(0,78500663) (1,216482163) (2,529510546) (3,164300351) (4,40311448) (5,34077) (6,4080913) (7,186) (8,120898) (10,7)};\end{axis}\end{tikzpicture}\end{minipage}
\vspace{3mm}

\noindent\begin{minipage}[t]{60mm}\vspace{0pt}\sbox\hmboard{\chessboard[setfen=8/8/8/6Qk/8/N5b1/8/K4R2 b - - 0 1]}\usebox\hmboard\\[1mm]
\makebox[\wd\hmboard]{\sffamily\hspace{3mm}\textbf{h\#5}\hfill(4+2)}\end{minipage}\hfill
\begin{minipage}[t]{52mm}\vspace{0pt}
\textcolor{ink}{\sffamily\textsc{Lösung}}\par\smallskip\small
\begin{tabular}{@{}r@{\enspace}l@{\enspace}l@{}}
\textbf{1.} & \symking{}xg5 & \symrook{}b1\\
\textbf{2.} & \symking{}f4 & \symrook{}b2\\
\textbf{3.} & \symking{}e3 & \symknight{}c2+\\
\textbf{4.} & \symking{}d2 & \symknight{}e3+\\
\textbf{5.} & \symking{}c1 & \symrook{}c2\#\\
\end{tabular}\par\vspace{3mm}
\qrcode[height=12mm]{https://helpman.komtera.lt/?fen=8/8/8/6Qk/8/N5b1/8/K4R2&moves=5}\\{\sffamily\tiny\color{ink}scannen zum Lösen}
\end{minipage}
\vspace{2mm}

{\raggedright\noindent\textcolor{ink}{\sffamily\textsc{Themen}}\enspace \themetag{ein Stein zieht} \themetag{Umnov}\par\smallskip{\footnotesize\textit{Hinweis:} die Partei am Zug steht im Diagramm im Schach; die Lösung beginnt mit einem Schlag.}\par}

\clearpage
\phantomsection\label{hm:KQRNvkn}
\begin{tcolorbox}[enhanced,colback=ink,colframe=ink,arc=0pt,boxrule=0pt,left=4mm,right=4mm,top=2.5mm,bottom=2.5mm]
\color{white}\sffamily{\LARGE\bfseries KQRNvkn}\hfill{\Large h\#5}\\[0.5mm]
{\small 6 Steine\quad Schwarz am Zug}\hfill{\small Nr.~791\qquad\href{https://helpman.komtera.lt/?fen=8/8/8/8/N7/7Q/n6k/K3R3&moves=5}{\color{white}\bfseries lösen\,\ensuremath{\nearrow}}}
\end{tcolorbox}
\vspace{2mm}
\noindent\begin{minipage}[t]{36mm}\vspace{0pt}\sffamily\small\raggedright
\hmfigure{längste eindeutige}{{\Large h\#5}}\hmfigure{eindeutig bei dieser Länge}{{\Large 754} Stellungen}\hmfigure{längstes Matt}{{\Large h\#6}}\hmfigure{in der Klasse}{15,5 Milliarden Stellungen,\\ 8,51 Milliarden mit Hilfsmatt}
\end{minipage}\hfill
\begin{minipage}[t]{80mm}\vspace{0pt}\begin{tikzpicture}\begin{axis}[width=72mm,height=42mm,ybar,bar width=1.6pt,xmin=-1,xmax=13,ymin=0,ymax=4.345,axis y line*=left,axis x line*=bottom,xlabel={\scriptsize Halbzüge bis Matt},ylabel={\scriptsize Stellungen (Mrd.)},tick label style={font=\tiny},label style={font=\scriptsize},enlarge x limits=false,ylabel shift=-3pt,xlabel shift=-3pt]\addplot[fill=ink!55,draw=none] coordinates {(0,0.09108) (1,0.3441) (2,1.311) (3,2.182) (4,3.95) (5,0.4355) (6,0.1412) (7,0.0004199) (8,0.04564) (9,1.017e-06) (10,0.006869) (11,9e-09) (12,3.259e-05)};\draw[densely dotted] (axis cs:10,0) -- (axis cs:10,4.345);\end{axis}\begin{axis}[width=72mm,height=42mm,axis y line*=right,axis x line=none,ymode=log,xmin=-1,xmax=13,ylabel={\scriptsize eindeutig (log)},tick label style={font=\tiny},label style={font=\scriptsize},ylabel shift=-3pt]\addplot[mark=*,mark size=0.7pt,draw=black,thin] coordinates {(0,91075086) (1,231971407) (2,521878167) (3,158609295) (4,48374226) (5,52038) (6,4149519) (7,2006) (8,254565) (9,1) (10,754)};\end{axis}\end{tikzpicture}\end{minipage}
\vspace{3mm}

\noindent\begin{minipage}[t]{60mm}\vspace{0pt}\sbox\hmboard{\chessboard[setfen=8/8/8/8/N7/7Q/n6k/K3R3 b - - 0 1]}\usebox\hmboard\\[1mm]
\makebox[\wd\hmboard]{\sffamily\hspace{3mm}\textbf{h\#5}\hfill(4+2)}\end{minipage}\hfill
\begin{minipage}[t]{52mm}\vspace{0pt}
\textcolor{ink}{\sffamily\textsc{Lösung}}\par\smallskip\small
\begin{tabular}{@{}r@{\enspace}l@{\enspace}l@{}}
\textbf{1.} & \symking{}xh3 & \symking{}b1\\
\textbf{2.} & \symking{}g2 & \symrook{}e2+\\
\textbf{3.} & \symking{}f1 & \symrook{}xa2\\
\textbf{4.} & \symking{}e1 & \symrook{}e2+\\
\textbf{5.} & \symking{}d1 & \symknight{}c3\#\\
\end{tabular}\par\vspace{3mm}
\qrcode[height=12mm]{https://helpman.komtera.lt/?fen=8/8/8/8/N7/7Q/n6k/K3R3&moves=5}\\{\sffamily\tiny\color{ink}scannen zum Lösen}
\end{minipage}
\vspace{2mm}

{\raggedright\noindent\textcolor{ink}{\sffamily\textsc{Themen}}\enspace \themetag{Switchback} \themetag{ein Stein zieht}\par\smallskip{\footnotesize\textit{Hinweis:} die Partei am Zug steht im Diagramm im Schach; die Lösung beginnt mit einem Schlag.}\par}

\clearpage
\phantomsection\label{hm:KQRNvkp}
\begin{tcolorbox}[enhanced,colback=ink,colframe=ink,arc=0pt,boxrule=0pt,left=4mm,right=4mm,top=2.5mm,bottom=2.5mm]
\color{white}\sffamily{\LARGE\bfseries KQRNvkp}\hfill{\Large h\#5}\\[0.5mm]
{\small 6 Steine\quad Schwarz am Zug}\hfill{\small Nr.~792\qquad\href{https://helpman.komtera.lt/?fen=8/8/8/8/6R1/p7/NQ6/K6k&moves=5}{\color{white}\bfseries lösen\,\ensuremath{\nearrow}}}
\end{tcolorbox}
\vspace{2mm}
\noindent\begin{minipage}[t]{36mm}\vspace{0pt}\sffamily\small\raggedright
\hmfigure{längste eindeutige}{{\Large h\#5}}\hmfigure{eindeutig bei dieser Länge}{{\Large 4\,367} Stellungen}\hmfigure{längstes Matt}{{\Large h\#8}}\hmfigure{in der Klasse}{45,4 Milliarden Stellungen,\\ 26,5 Milliarden mit Hilfsmatt}
\end{minipage}\hfill
\begin{minipage}[t]{80mm}\vspace{0pt}\begin{tikzpicture}\begin{axis}[width=72mm,height=42mm,ybar,bar width=1.6pt,xmin=-1,xmax=17,ymin=0,ymax=13.5,axis y line*=left,axis x line*=bottom,xlabel={\scriptsize Halbzüge bis Matt},ylabel={\scriptsize Stellungen (Mrd.)},tick label style={font=\tiny},label style={font=\scriptsize},enlarge x limits=false,ylabel shift=-3pt,xlabel shift=-3pt]\addplot[fill=ink!55,draw=none] coordinates {(0,0.3324) (1,1.235) (2,4.388) (3,6.719) (4,12.27) (5,0.9765) (6,0.4551) (7,0.0004611) (8,0.135) (9,1.687e-06) (10,0.02851) (12,0.00055) (14,8.8e-08) (16,2e-09)};\draw[densely dotted] (axis cs:10,0) -- (axis cs:10,13.5);\end{axis}\begin{axis}[width=72mm,height=42mm,axis y line*=right,axis x line=none,ymode=log,xmin=-1,xmax=17,ylabel={\scriptsize eindeutig (log)},tick label style={font=\tiny},label style={font=\scriptsize},ylabel shift=-3pt]\addplot[mark=*,mark size=0.7pt,draw=black,thin] coordinates {(0,332356821) (1,803895304) (2,1943438033) (3,490707901) (4,159378667) (5,90501) (6,12543642) (7,109) (8,862502) (10,4367)};\end{axis}\end{tikzpicture}\end{minipage}
\vspace{3mm}

\noindent\begin{minipage}[t]{60mm}\vspace{0pt}\sbox\hmboard{\chessboard[setfen=8/8/8/8/6R1/p7/NQ6/K6k b - - 0 1]}\usebox\hmboard\\[1mm]
\makebox[\wd\hmboard]{\sffamily\hspace{3mm}\textbf{h\#5}\hfill(4+2)}\end{minipage}\hfill
\begin{minipage}[t]{52mm}\vspace{0pt}
\textcolor{ink}{\sffamily\textsc{Lösung}}\par\smallskip\small
\begin{tabular}{@{}r@{\enspace}l@{\enspace}l@{}}
\textbf{1.} & axb2+ & \symking{}b1\\
\textbf{2.} & \symking{}h2 & \symking{}c2\\
\textbf{3.} & b1=\symrook{} & \symknight{}c3\\
\textbf{4.} & \symrook{}g1 & \symknight{}e2\\
\textbf{5.} & \symrook{}g2 & \symrook{}h4\#\\
\end{tabular}\par\vspace{3mm}
\qrcode[height=12mm]{https://helpman.komtera.lt/?fen=8/8/8/8/6R1/p7/NQ6/K6k&moves=5}\\{\sffamily\tiny\color{ink}scannen zum Lösen}
\end{minipage}
\vspace{2mm}

{\raggedright\noindent\textcolor{ink}{\sffamily\textsc{Themen}}\enspace \themetag{reines Matt} \themetag{Mustermatt} \themetag{Umwandlung} \themetag{Unterverwandlung} \themetag{Selbstblock} \themetag{Umnov}\par\smallskip{\footnotesize\textit{Hinweis:} die Lösung beginnt mit einem Schlag.}\par}

\clearpage
\phantomsection\label{hm:KQRNvkq}
\begin{tcolorbox}[enhanced,colback=ink,colframe=ink,arc=0pt,boxrule=0pt,left=4mm,right=4mm,top=2.5mm,bottom=2.5mm]
\color{white}\sffamily{\LARGE\bfseries KQRNvkq}\hfill{\Large h\#5}\\[0.5mm]
{\small 6 Steine\quad Schwarz am Zug}\hfill{\small Nr.~793\qquad\href{https://helpman.komtera.lt/?fen=8/8/N7/8/8/8/3q1R2/K3Q2k&moves=5}{\color{white}\bfseries lösen\,\ensuremath{\nearrow}}}
\end{tcolorbox}
\vspace{2mm}
\noindent\begin{minipage}[t]{36mm}\vspace{0pt}\sffamily\small\raggedright
\hmfigure{längste eindeutige}{{\Large h\#5}}\hmfigure{eindeutig bei dieser Länge}{{\Large 119} Stellungen}\hmfigure{längstes Matt}{{\Large h\#6}}\hmfigure{in der Klasse}{15,5 Milliarden Stellungen,\\ 7,03 Milliarden mit Hilfsmatt}
\end{minipage}\hfill
\begin{minipage}[t]{80mm}\vspace{0pt}\begin{tikzpicture}\begin{axis}[width=72mm,height=42mm,ybar,bar width=1.6pt,xmin=-1,xmax=13,ymin=0,ymax=3.277,axis y line*=left,axis x line*=bottom,xlabel={\scriptsize Halbzüge bis Matt},ylabel={\scriptsize Stellungen (Mrd.)},tick label style={font=\tiny},label style={font=\scriptsize},enlarge x limits=false,ylabel shift=-3pt,xlabel shift=-3pt]\addplot[fill=ink!55,draw=none] coordinates {(0,0.04662) (1,0.1914) (2,0.9223) (3,1.952) (4,2.979) (5,0.8148) (6,0.09116) (7,0.0008355) (8,0.03456) (9,1.218e-05) (10,0.002077) (11,2.2e-08) (12,1.413e-06)};\draw[densely dotted] (axis cs:10,0) -- (axis cs:10,3.277);\end{axis}\begin{axis}[width=72mm,height=42mm,axis y line*=right,axis x line=none,ymode=log,xmin=-1,xmax=13,ylabel={\scriptsize eindeutig (log)},tick label style={font=\tiny},label style={font=\scriptsize},ylabel shift=-3pt]\addplot[mark=*,mark size=0.7pt,draw=black,thin] coordinates {(0,46617314) (1,146502311) (2,415315736) (3,187804568) (4,30971507) (5,320463) (6,2713219) (7,1748) (8,94539) (9,18) (10,119)};\end{axis}\end{tikzpicture}\end{minipage}
\vspace{3mm}

\noindent\begin{minipage}[t]{60mm}\vspace{0pt}\sbox\hmboard{\chessboard[setfen=8/8/N7/8/8/8/3q1R2/K3Q2k b - - 0 1]}\usebox\hmboard\\[1mm]
\makebox[\wd\hmboard]{\sffamily\hspace{3mm}\textbf{h\#5}\hfill(4+2)}\end{minipage}\hfill
\begin{minipage}[t]{52mm}\vspace{0pt}
\textcolor{ink}{\sffamily\textsc{Lösung}}\par\smallskip\small
\begin{tabular}{@{}r@{\enspace}l@{\enspace}l@{}}
\textbf{1.} & \symqueen{}xe1+ & \symking{}b2\\
\textbf{2.} & \symking{}g1 & \symrook{}c2\\
\textbf{3.} & \symqueen{}f2 & \symking{}c3\\
\textbf{4.} & \symking{}f1 & \symking{}d3\\
\textbf{5.} & \symking{}e1 & \symrook{}c1\#\\
\end{tabular}\par\vspace{3mm}
\qrcode[height=12mm]{https://helpman.komtera.lt/?fen=8/8/N7/8/8/8/3q1R2/K3Q2k&moves=5}\\{\sffamily\tiny\color{ink}scannen zum Lösen}
\end{minipage}
\vspace{2mm}

{\raggedright\noindent\textcolor{ink}{\sffamily\textsc{Themen}}\enspace \themetag{reines Matt} \themetag{Selbstblock} \themetag{Kniest} \themetag{Umnov}\par\smallskip{\footnotesize\textit{Hinweis:} die Partei am Zug steht im Diagramm im Schach; die Lösung beginnt mit einem Schlag.}\par}

\clearpage
\phantomsection\label{hm:KQRNvkr}
\begin{tcolorbox}[enhanced,colback=ink,colframe=ink,arc=0pt,boxrule=0pt,left=4mm,right=4mm,top=2.5mm,bottom=2.5mm]
\color{white}\sffamily{\LARGE\bfseries KQRNvkr}\hfill{\Large h\#5}\\[0.5mm]
{\small 6 Steine\quad Schwarz am Zug}\hfill{\small Nr.~794\qquad\href{https://helpman.komtera.lt/?fen=8/8/N7/8/8/8/5R2/K4Qrk&moves=5}{\color{white}\bfseries lösen\,\ensuremath{\nearrow}}}
\end{tcolorbox}
\vspace{2mm}
\noindent\begin{minipage}[t]{36mm}\vspace{0pt}\sffamily\small\raggedright
\hmfigure{längste eindeutige}{{\Large h\#5}}\hmfigure{eindeutig bei dieser Länge}{{\Large 4} Stellungen}\hmfigure{längstes Matt}{{\Large h\#6}}\hmfigure{in der Klasse}{15,5 Milliarden Stellungen,\\ 7,81 Milliarden mit Hilfsmatt}
\end{minipage}\hfill
\begin{minipage}[t]{80mm}\vspace{0pt}\begin{tikzpicture}\begin{axis}[width=72mm,height=42mm,ybar,bar width=1.6pt,xmin=-1,xmax=13,ymin=0,ymax=3.766,axis y line*=left,axis x line*=bottom,xlabel={\scriptsize Halbzüge bis Matt},ylabel={\scriptsize Stellungen (Mrd.)},tick label style={font=\tiny},label style={font=\scriptsize},enlarge x limits=false,ylabel shift=-3pt,xlabel shift=-3pt]\addplot[fill=ink!55,draw=none] coordinates {(0,0.06967) (1,0.28) (2,1.223) (3,2.12) (4,3.423) (5,0.5613) (6,0.1) (7,0.000182) (8,0.03197) (9,6.11e-07) (10,0.0003956) (11,1e-09) (12,4e-08)};\draw[densely dotted] (axis cs:10,0) -- (axis cs:10,3.766);\end{axis}\begin{axis}[width=72mm,height=42mm,axis y line*=right,axis x line=none,ymode=log,xmin=-1,xmax=13,ylabel={\scriptsize eindeutig (log)},tick label style={font=\tiny},label style={font=\scriptsize},ylabel shift=-3pt]\addplot[mark=*,mark size=0.7pt,draw=black,thin] coordinates {(0,69671355) (1,200705208) (2,523963719) (3,153479680) (4,31226969) (5,51325) (6,3818845) (7,107) (8,42881) (9,5) (10,4)};\end{axis}\end{tikzpicture}\end{minipage}
\vspace{3mm}

\noindent\begin{minipage}[t]{60mm}\vspace{0pt}\sbox\hmboard{\chessboard[setfen=8/8/N7/8/8/8/5R2/K4Qrk b - - 0 1]}\usebox\hmboard\\[1mm]
\makebox[\wd\hmboard]{\sffamily\hspace{3mm}\textbf{h\#5}\hfill(4+2)}\end{minipage}\hfill
\begin{minipage}[t]{52mm}\vspace{0pt}
\textcolor{ink}{\sffamily\textsc{Lösung}}\par\smallskip\small
\begin{tabular}{@{}r@{\enspace}l@{\enspace}l@{}}
\textbf{1.} & \symrook{}xf1+ & \symking{}b2\\
\textbf{2.} & \symking{}g1 & \symrook{}c2\\
\textbf{3.} & \symrook{}f2 & \symking{}c3\\
\textbf{4.} & \symking{}f1 & \symking{}d3\\
\textbf{5.} & \symking{}e1 & \symrook{}c1\#\\
\end{tabular}\par\vspace{3mm}
\qrcode[height=12mm]{https://helpman.komtera.lt/?fen=8/8/N7/8/8/8/5R2/K4Qrk&moves=5}\\{\sffamily\tiny\color{ink}scannen zum Lösen}
\end{minipage}
\vspace{2mm}

{\raggedright\noindent\textcolor{ink}{\sffamily\textsc{Themen}}\enspace \themetag{reines Matt} \themetag{Selbstblock} \themetag{Umnov}\par\smallskip{\footnotesize\textit{Hinweis:} die Lösung beginnt mit einem Schlag.}\par}

\clearpage
\phantomsection\label{hm:KQvkqrr}
\begin{tcolorbox}[enhanced,colback=ink,colframe=ink,arc=0pt,boxrule=0pt,left=4mm,right=4mm,top=2.5mm,bottom=2.5mm]
\color{white}\sffamily{\LARGE\bfseries KQvkqrr}\hfill{\Large h\#5}\\[0.5mm]
{\small 6 Steine\quad Schwarz am Zug}\hfill{\small Nr.~795\qquad\href{https://helpman.komtera.lt/?fen=8/8/8/8/8/2r1Q3/3K1kq1/2r5&moves=5}{\color{white}\bfseries lösen\,\ensuremath{\nearrow}}}
\end{tcolorbox}
\vspace{2mm}
\noindent\begin{minipage}[t]{36mm}\vspace{0pt}\sffamily\small\raggedright
\hmfigure{längste eindeutige}{{\Large h\#5}}\hmfigure{eindeutig bei dieser Länge}{{\Large 1} Stellung}\hmfigure{längstes Matt}{{\Large h\#6}}\hmfigure{in der Klasse}{15,5 Milliarden Stellungen,\\ 3,2 Milliarden mit Hilfsmatt}
\end{minipage}\hfill
\begin{minipage}[t]{80mm}\vspace{0pt}\begin{tikzpicture}\begin{axis}[width=72mm,height=42mm,ybar,bar width=1.6pt,xmin=-1,xmax=13,ymin=0,ymax=1039,axis y line*=left,axis x line*=bottom,xlabel={\scriptsize Halbzüge bis Matt},ylabel={\scriptsize Stellungen (Mio.)},tick label style={font=\tiny},label style={font=\scriptsize},enlarge x limits=false,ylabel shift=-3pt,xlabel shift=-3pt]\addplot[fill=ink!55,draw=none] coordinates {(0,0.5454) (1,5.518) (2,31.34) (3,167.6) (4,423.2) (5,944.5) (6,651.7) (7,796.2) (8,113.5) (9,63.99) (10,1.543) (11,0.02496) (12,1.6e-05)};\draw[densely dotted] (axis cs:10,0) -- (axis cs:10,1039);\end{axis}\begin{axis}[width=72mm,height=42mm,axis y line*=right,axis x line=none,ymode=log,xmin=-1,xmax=13,ylabel={\scriptsize eindeutig (log)},tick label style={font=\tiny},label style={font=\scriptsize},ylabel shift=-3pt]\addplot[mark=*,mark size=0.7pt,draw=black,thin] coordinates {(0,545378) (1,4997979) (2,17548021) (3,39785537) (4,31549681) (5,30892570) (6,4407990) (7,2534411) (8,63715) (9,2185) (10,1)};\end{axis}\end{tikzpicture}\end{minipage}
\vspace{3mm}

\noindent\begin{minipage}[t]{60mm}\vspace{0pt}\sbox\hmboard{\chessboard[setfen=8/8/8/8/8/2r1Q3/3K1kq1/2r5 b - - 0 1]}\usebox\hmboard\\[1mm]
\makebox[\wd\hmboard]{\sffamily\hspace{3mm}\textbf{h\#5}\hfill(2+4)}\end{minipage}\hfill
\begin{minipage}[t]{52mm}\vspace{0pt}
\textcolor{ink}{\sffamily\textsc{Lösung}}\par\smallskip\small
\begin{tabular}{@{}r@{\enspace}l@{\enspace}l@{}}
\textbf{1.} & \symking{}f1+ & \symqueen{}e2+\\
\textbf{2.} & \symking{}g1 & \symqueen{}f2+\\
\textbf{3.} & \symking{}h2 & \symqueen{}e2\\
\textbf{4.} & \symrook{}g3 & \symking{}xc1\\
\textbf{5.} & \symking{}h3 & \symqueen{}h5\#\\
\end{tabular}\par\vspace{3mm}
\qrcode[height=12mm]{https://helpman.komtera.lt/?fen=8/8/8/8/8/2r1Q3/3K1kq1/2r5&moves=5}\\{\sffamily\tiny\color{ink}scannen zum Lösen}
\end{minipage}
\vspace{2mm}

{\raggedright\noindent\textcolor{ink}{\sffamily\textsc{Themen}}\enspace \themetag{reines Matt} \themetag{Mustermatt} \themetag{Switchback} \themetag{Selbstblock}\par\smallskip{\footnotesize\textit{Hinweis:} die Partei am Zug steht im Diagramm im Schach.}\par}

\clearpage
\phantomsection\label{hm:KQvkrrb}
\begin{tcolorbox}[enhanced,colback=ink,colframe=ink,arc=0pt,boxrule=0pt,left=4mm,right=4mm,top=2.5mm,bottom=2.5mm]
\color{white}\sffamily{\LARGE\bfseries KQvkrrb}\hfill{\Large h\#5}\\[0.5mm]
{\small 6 Steine\quad Schwarz am Zug}\hfill{\small Nr.~796\qquad\href{https://helpman.komtera.lt/?fen=2r5/1r6/8/8/8/7k/b6Q/K7&moves=5}{\color{white}\bfseries lösen\,\ensuremath{\nearrow}}}
\end{tcolorbox}
\vspace{2mm}
\noindent\begin{minipage}[t]{36mm}\vspace{0pt}\sffamily\small\raggedright
\hmfigure{längste eindeutige}{{\Large h\#5}}\hmfigure{eindeutig bei dieser Länge}{{\Large 2} Stellungen}\hmfigure{längstes Matt}{{\Large h\#6}}\hmfigure{in der Klasse}{15,5 Milliarden Stellungen,\\ 3,68 Milliarden mit Hilfsmatt}
\end{minipage}\hfill
\begin{minipage}[t]{80mm}\vspace{0pt}\begin{tikzpicture}\begin{axis}[width=72mm,height=42mm,ybar,bar width=1.6pt,xmin=-1,xmax=13,ymin=0,ymax=1.103,axis y line*=left,axis x line*=bottom,xlabel={\scriptsize Halbzüge bis Matt},ylabel={\scriptsize Stellungen (Mrd.)},tick label style={font=\tiny},label style={font=\scriptsize},enlarge x limits=false,ylabel shift=-3pt,xlabel shift=-3pt]\addplot[fill=ink!55,draw=none] coordinates {(0,0.001065) (1,0.009742) (2,0.04957) (3,0.2262) (4,0.6087) (5,1.003) (6,0.84) (7,0.7281) (8,0.1527) (9,0.05976) (10,0.002358) (11,4.121e-05) (12,1.11e-07)};\draw[densely dotted] (axis cs:10,0) -- (axis cs:10,1.103);\end{axis}\begin{axis}[width=72mm,height=42mm,axis y line*=right,axis x line=none,ymode=log,xmin=-1,xmax=13,ylabel={\scriptsize eindeutig (log)},tick label style={font=\tiny},label style={font=\scriptsize},ylabel shift=-3pt]\addplot[mark=*,mark size=0.7pt,draw=black,thin] coordinates {(0,1064712) (1,8487924) (2,28411883) (3,45627558) (4,41807583) (5,27053185) (6,5795830) (7,2264789) (8,108784) (9,3251) (10,2)};\end{axis}\end{tikzpicture}\end{minipage}
\vspace{3mm}

\noindent\begin{minipage}[t]{60mm}\vspace{0pt}\sbox\hmboard{\chessboard[setfen=2r5/1r6/8/8/8/7k/b6Q/K7 b - - 0 1]}\usebox\hmboard\\[1mm]
\makebox[\wd\hmboard]{\sffamily\hspace{3mm}\textbf{h\#5}\hfill(2+4)}\end{minipage}\hfill
\begin{minipage}[t]{52mm}\vspace{0pt}
\textcolor{ink}{\sffamily\textsc{Lösung}}\par\smallskip\small
\begin{tabular}{@{}r@{\enspace}l@{\enspace}l@{}}
\textbf{1.} & \symking{}g4 & \symqueen{}b2\\
\textbf{2.} & \symbishop{}c4 & \symking{}b1\\
\textbf{3.} & \symking{}f3 & \symking{}c2\\
\textbf{4.} & \symking{}e2 & \symking{}c3+\\
\textbf{5.} & \symking{}d1 & \symqueen{}d2\#\\
\end{tabular}\par\vspace{3mm}
\qrcode[height=12mm]{https://helpman.komtera.lt/?fen=2r5/1r6/8/8/8/7k/b6Q/K7&moves=5}\\{\sffamily\tiny\color{ink}scannen zum Lösen}
\end{minipage}
\vspace{2mm}

{\raggedright\noindent\textcolor{ink}{\sffamily\textsc{Themen}}\enspace \themetag{ohne Schlag} \themetag{Inder}\par\smallskip{\footnotesize\textit{Hinweis:} die Partei am Zug steht im Diagramm im Schach.}\par}

\clearpage
\phantomsection\label{hm:KQvkrrn}
\begin{tcolorbox}[enhanced,colback=ink,colframe=ink,arc=0pt,boxrule=0pt,left=4mm,right=4mm,top=2.5mm,bottom=2.5mm]
\color{white}\sffamily{\LARGE\bfseries KQvkrrn}\hfill{\Large h\#5}\\[0.5mm]
{\small 6 Steine\quad Schwarz am Zug}\hfill{\small Nr.~797\qquad\href{https://helpman.komtera.lt/?fen=1r6/8/8/8/2r5/4nk2/6Q1/K7&moves=5}{\color{white}\bfseries lösen\,\ensuremath{\nearrow}}}
\end{tcolorbox}
\vspace{2mm}
\noindent\begin{minipage}[t]{36mm}\vspace{0pt}\sffamily\small\raggedright
\hmfigure{längste eindeutige}{{\Large h\#5}}\hmfigure{eindeutig bei dieser Länge}{{\Large 110} Stellungen}\hmfigure{längstes Matt}{{\Large h\#6}}\hmfigure{in der Klasse}{15,5 Milliarden Stellungen,\\ 3,76 Milliarden mit Hilfsmatt}
\end{minipage}\hfill
\begin{minipage}[t]{80mm}\vspace{0pt}\begin{tikzpicture}\begin{axis}[width=72mm,height=42mm,ybar,bar width=1.6pt,xmin=-1,xmax=13,ymin=0,ymax=997.1,axis y line*=left,axis x line*=bottom,xlabel={\scriptsize Halbzüge bis Matt},ylabel={\scriptsize Stellungen (Mio.)},tick label style={font=\tiny},label style={font=\scriptsize},enlarge x limits=false,ylabel shift=-3pt,xlabel shift=-3pt]\addplot[fill=ink!55,draw=none] coordinates {(0,1.166) (1,10.36) (2,50.14) (3,217) (4,572.4) (5,906.5) (6,871.7) (7,785.5) (8,234.4) (9,108.4) (10,5.478) (11,0.5431) (12,0.001719)};\draw[densely dotted] (axis cs:10,0) -- (axis cs:10,997.1);\end{axis}\begin{axis}[width=72mm,height=42mm,axis y line*=right,axis x line=none,ymode=log,xmin=-1,xmax=13,ylabel={\scriptsize eindeutig (log)},tick label style={font=\tiny},label style={font=\scriptsize},ylabel shift=-3pt]\addplot[mark=*,mark size=0.7pt,draw=black,thin] coordinates {(0,1165987) (1,8859934) (2,28257863) (3,41685928) (4,38589681) (5,26613278) (6,8083547) (7,3494992) (8,219795) (9,33477) (10,110)};\end{axis}\end{tikzpicture}\end{minipage}
\vspace{3mm}

\noindent\begin{minipage}[t]{60mm}\vspace{0pt}\sbox\hmboard{\chessboard[setfen=1r6/8/8/8/2r5/4nk2/6Q1/K7 b - - 0 1]}\usebox\hmboard\\[1mm]
\makebox[\wd\hmboard]{\sffamily\hspace{3mm}\textbf{h\#5}\hfill(2+4)}\end{minipage}\hfill
\begin{minipage}[t]{52mm}\vspace{0pt}
\textcolor{ink}{\sffamily\textsc{Lösung}}\par\smallskip\small
\begin{tabular}{@{}r@{\enspace}l@{\enspace}l@{}}
\textbf{1.} & \symking{}f4 & \symqueen{}b7\\
\textbf{2.} & \symrook{}e4 & \symking{}b2\\
\textbf{3.} & \symking{}f3 & \symking{}c3\\
\textbf{4.} & \symking{}e2 & \symqueen{}b2+\\
\textbf{5.} & \symking{}d1 & \symqueen{}d2\#\\
\end{tabular}\par\vspace{3mm}
\qrcode[height=12mm]{https://helpman.komtera.lt/?fen=1r6/8/8/8/2r5/4nk2/6Q1/K7&moves=5}\\{\sffamily\tiny\color{ink}scannen zum Lösen}
\end{minipage}
\vspace{2mm}

{\raggedright\noindent\textcolor{ink}{\sffamily\textsc{Themen}}\enspace \themetag{Switchback} \themetag{ohne Schlag} \themetag{Klasinc}\par\smallskip{\footnotesize\textit{Hinweis:} die Partei am Zug steht im Diagramm im Schach.}\par}

\clearpage
\phantomsection\label{hm:KQvkrrp}
\begin{tcolorbox}[enhanced,colback=ink,colframe=ink,arc=0pt,boxrule=0pt,left=4mm,right=4mm,top=2.5mm,bottom=2.5mm]
\color{white}\sffamily{\LARGE\bfseries KQvkrrp}\hfill{\Large h\#5}\\[0.5mm]
{\small 6 Steine\quad Schwarz am Zug}\hfill{\small Nr.~798\qquad\href{https://helpman.komtera.lt/?fen=8/8/8/8/7k/2p4Q/1r2r3/K7&moves=5}{\color{white}\bfseries lösen\,\ensuremath{\nearrow}}}
\end{tcolorbox}
\vspace{2mm}
\noindent\begin{minipage}[t]{36mm}\vspace{0pt}\sffamily\small\raggedright
\hmfigure{längste eindeutige}{{\Large h\#5}}\hmfigure{eindeutig bei dieser Länge}{{\Large 17} Stellungen}\hmfigure{längstes Matt}{{\Large h\#6}}\hmfigure{in der Klasse}{45,4 Milliarden Stellungen,\\ 11,6 Milliarden mit Hilfsmatt}
\end{minipage}\hfill
\begin{minipage}[t]{80mm}\vspace{0pt}\begin{tikzpicture}\begin{axis}[width=72mm,height=42mm,ybar,bar width=1.6pt,xmin=-1,xmax=13,ymin=0,ymax=3.354,axis y line*=left,axis x line*=bottom,xlabel={\scriptsize Halbzüge bis Matt},ylabel={\scriptsize Stellungen (Mrd.)},tick label style={font=\tiny},label style={font=\scriptsize},enlarge x limits=false,ylabel shift=-3pt,xlabel shift=-3pt]\addplot[fill=ink!55,draw=none] coordinates {(0,0.004597) (1,0.03991) (2,0.1925) (3,0.8041) (4,2.154) (5,3.049) (6,2.664) (7,2.035) (8,0.4983) (9,0.1837) (10,0.008344) (11,0.0005052) (12,5.81e-07)};\draw[densely dotted] (axis cs:10,0) -- (axis cs:10,3.354);\end{axis}\begin{axis}[width=72mm,height=42mm,axis y line*=right,axis x line=none,ymode=log,xmin=-1,xmax=13,ylabel={\scriptsize eindeutig (log)},tick label style={font=\tiny},label style={font=\scriptsize},ylabel shift=-3pt]\addplot[mark=*,mark size=0.7pt,draw=black,thin] coordinates {(0,4596614) (1,33887445) (2,112322301) (3,148479121) (4,137190846) (5,74626177) (6,17930685) (7,6311678) (8,364827) (9,40660) (10,17)};\end{axis}\end{tikzpicture}\end{minipage}
\vspace{3mm}

\noindent\begin{minipage}[t]{60mm}\vspace{0pt}\sbox\hmboard{\chessboard[setfen=8/8/8/8/7k/2p4Q/1r2r3/K7 b - - 0 1]}\usebox\hmboard\\[1mm]
\makebox[\wd\hmboard]{\sffamily\hspace{3mm}\textbf{h\#5}\hfill(2+4)}\end{minipage}\hfill
\begin{minipage}[t]{52mm}\vspace{0pt}
\textcolor{ink}{\sffamily\textsc{Lösung}}\par\smallskip\small
\begin{tabular}{@{}r@{\enspace}l@{\enspace}l@{}}
\textbf{1.} & \symking{}g5 & \symqueen{}xc3\\
\textbf{2.} & \symking{}f4 & \symqueen{}xb2\\
\textbf{3.} & \symking{}e3 & \symqueen{}a3+\\
\textbf{4.} & \symking{}d2 & \symking{}b1\\
\textbf{5.} & \symking{}d1 & \symqueen{}c1\#\\
\end{tabular}\par\vspace{3mm}
\qrcode[height=12mm]{https://helpman.komtera.lt/?fen=8/8/8/8/7k/2p4Q/1r2r3/K7&moves=5}\\{\sffamily\tiny\color{ink}scannen zum Lösen}
\end{minipage}
\vspace{2mm}

{\raggedright\noindent\textcolor{ink}{\sffamily\textsc{Themen}}\enspace \themetag{ein Stein zieht}\par\smallskip{\footnotesize\textit{Hinweis:} die Partei am Zug steht im Diagramm im Schach.}\par}

\clearpage
\phantomsection\label{hm:KRRBNvk}
\begin{tcolorbox}[enhanced,colback=ink,colframe=ink,arc=0pt,boxrule=0pt,left=4mm,right=4mm,top=2.5mm,bottom=2.5mm]
\color{white}\sffamily{\LARGE\bfseries KRRBNvk}\hfill{\Large h\#5}\\[0.5mm]
{\small 6 Steine\quad Schwarz am Zug}\hfill{\small Nr.~799\qquad\href{https://helpman.komtera.lt/?fen=R1B5/kR6/8/8/8/8/8/1K5N&moves=5}{\color{white}\bfseries lösen\,\ensuremath{\nearrow}}}
\end{tcolorbox}
\vspace{2mm}
\noindent\begin{minipage}[t]{36mm}\vspace{0pt}\sffamily\small\raggedright
\hmfigure{längste eindeutige}{{\Large h\#5}}\hmfigure{eindeutig bei dieser Länge}{{\Large 16} Stellungen}\hmfigure{längstes Matt}{{\Large h\#6}}\hmfigure{in der Klasse}{15,5 Milliarden Stellungen,\\ 4,53 Milliarden mit Hilfsmatt}
\end{minipage}\hfill
\begin{minipage}[t]{80mm}\vspace{0pt}\begin{tikzpicture}\begin{axis}[width=72mm,height=42mm,ybar,bar width=1.6pt,xmin=-1,xmax=13,ymin=0,ymax=2.339,axis y line*=left,axis x line*=bottom,xlabel={\scriptsize Halbzüge bis Matt},ylabel={\scriptsize Stellungen (Mrd.)},tick label style={font=\tiny},label style={font=\scriptsize},enlarge x limits=false,ylabel shift=-3pt,xlabel shift=-3pt]\addplot[fill=ink!55,draw=none] coordinates {(0,0.04123) (1,0.1359) (2,0.4108) (3,0.9012) (4,2.127) (5,0.5061) (6,0.401) (7,0.0005262) (8,0.00899) (9,1e-08) (10,0.0002023) (12,3.3e-08)};\draw[densely dotted] (axis cs:10,0) -- (axis cs:10,2.339);\end{axis}\begin{axis}[width=72mm,height=42mm,axis y line*=right,axis x line=none,ymode=log,xmin=-1,xmax=13,ylabel={\scriptsize eindeutig (log)},tick label style={font=\tiny},label style={font=\scriptsize},ylabel shift=-3pt]\addplot[mark=*,mark size=0.7pt,draw=black,thin] coordinates {(0,41227308) (1,107168645) (2,247511425) (3,175157597) (4,175562549) (5,249572) (6,2235917) (7,3) (8,9685) (10,16)};\end{axis}\end{tikzpicture}\end{minipage}
\vspace{3mm}

\noindent\begin{minipage}[t]{60mm}\vspace{0pt}\sbox\hmboard{\chessboard[setfen=R1B5/kR6/8/8/8/8/8/1K5N b - - 0 1]}\usebox\hmboard\\[1mm]
\makebox[\wd\hmboard]{\sffamily\hspace{3mm}\textbf{h\#5}\hfill(5+1)}\end{minipage}\hfill
\begin{minipage}[t]{52mm}\vspace{0pt}
\textcolor{ink}{\sffamily\textsc{Lösung}}\par\smallskip\small
\begin{tabular}{@{}r@{\enspace}l@{\enspace}l@{}}
\textbf{1.} & \symking{}xa8 & \symrook{}a7+\\
\textbf{2.} & \symking{}b8 & \symrook{}a8+\\
\textbf{3.} & \symking{}c7 & \symknight{}g3\\
\textbf{4.} & \symking{}d8 & \symknight{}f5\\
\textbf{5.} & \symking{}e8 & \symbishop{}e6\#\\
\end{tabular}\par\vspace{3mm}
\qrcode[height=12mm]{https://helpman.komtera.lt/?fen=R1B5/kR6/8/8/8/8/8/1K5N&moves=5}\\{\sffamily\tiny\color{ink}scannen zum Lösen}
\end{minipage}
\vspace{2mm}

{\raggedright\noindent\textcolor{ink}{\sffamily\textsc{Themen}}\enspace \themetag{reines Matt} \themetag{Mustermatt} \themetag{Spiegelmatt} \themetag{ein Stein zieht} \themetag{Umnov}\par\smallskip{\footnotesize\textit{Hinweis:} das Diagramm hat keinen legalen letzten Zug, kann also in keiner Partie entstehen.}\par}

\clearpage
\phantomsection\label{hm:KRRNNvk}
\begin{tcolorbox}[enhanced,colback=ink,colframe=ink,arc=0pt,boxrule=0pt,left=4mm,right=4mm,top=2.5mm,bottom=2.5mm]
\color{white}\sffamily{\LARGE\bfseries KRRNNvk}\hfill{\Large h\#5}\\[0.5mm]
{\small 6 Steine\quad Schwarz am Zug}\hfill{\small Nr.~800\qquad\href{https://helpman.komtera.lt/?fen=6R1/7k/4R3/N7/8/8/1K6/N7&moves=5}{\color{white}\bfseries lösen\,\ensuremath{\nearrow}}}
\end{tcolorbox}
\vspace{2mm}
\noindent\begin{minipage}[t]{36mm}\vspace{0pt}\sffamily\small\raggedright
\hmfigure{längste eindeutige}{{\Large h\#5}}\hmfigure{eindeutig bei dieser Länge}{{\Large 3} Stellungen}\hmfigure{längstes Matt}{{\Large h\#6}}\hmfigure{in der Klasse}{15,5 Milliarden Stellungen,\\ 2,31 Milliarden mit Hilfsmatt}
\end{minipage}\hfill
\begin{minipage}[t]{80mm}\vspace{0pt}\begin{tikzpicture}\begin{axis}[width=72mm,height=42mm,ybar,bar width=1.6pt,xmin=-1,xmax=13,ymin=0,ymax=1.155,axis y line*=left,axis x line*=bottom,xlabel={\scriptsize Halbzüge bis Matt},ylabel={\scriptsize Stellungen (Mrd.)},tick label style={font=\tiny},label style={font=\scriptsize},enlarge x limits=false,ylabel shift=-3pt,xlabel shift=-3pt]\addplot[fill=ink!55,draw=none] coordinates {(0,0.01853) (1,0.06516) (2,0.1931) (3,0.4572) (4,1.05) (5,0.2852) (6,0.2296) (7,0.000685) (8,0.005465) (10,0.0002408) (12,2.81e-07)};\draw[densely dotted] (axis cs:10,0) -- (axis cs:10,1.155);\end{axis}\begin{axis}[width=72mm,height=42mm,axis y line*=right,axis x line=none,ymode=log,xmin=-1,xmax=13,ylabel={\scriptsize eindeutig (log)},tick label style={font=\tiny},label style={font=\scriptsize},ylabel shift=-3pt]\addplot[mark=*,mark size=0.7pt,draw=black,thin] coordinates {(0,18527922) (1,52479978) (2,117455069) (3,83212318) (4,85303391) (5,218664) (6,1121737) (8,7044) (10,3)};\end{axis}\end{tikzpicture}\end{minipage}
\vspace{3mm}

\noindent\begin{minipage}[t]{60mm}\vspace{0pt}\sbox\hmboard{\chessboard[setfen=6R1/7k/4R3/N7/8/8/1K6/N7 b - - 0 1]}\usebox\hmboard\\[1mm]
\makebox[\wd\hmboard]{\sffamily\hspace{3mm}\textbf{h\#5}\hfill(5+1)}\end{minipage}\hfill
\begin{minipage}[t]{52mm}\vspace{0pt}
\textcolor{ink}{\sffamily\textsc{Lösung}}\par\smallskip\small
\begin{tabular}{@{}r@{\enspace}l@{\enspace}l@{}}
\textbf{1.} & \symking{}xg8 & \symrook{}a6\\
\textbf{2.} & \symking{}f7 & \symknight{}c6\\
\textbf{3.} & \symking{}e6 & \symknight{}c2\\
\textbf{4.} & \symking{}d5 & \symrook{}a5+\\
\textbf{5.} & \symking{}c4 & \symknight{}e5\#\\
\end{tabular}\par\vspace{3mm}
\qrcode[height=12mm]{https://helpman.komtera.lt/?fen=6R1/7k/4R3/N7/8/8/1K6/N7&moves=5}\\{\sffamily\tiny\color{ink}scannen zum Lösen}
\end{minipage}
\vspace{2mm}

{\raggedright\noindent\textcolor{ink}{\sffamily\textsc{Themen}}\enspace \themetag{reines Matt} \themetag{Mustermatt} \themetag{Idealmatt} \themetag{Spiegelmatt} \themetag{ein Stein zieht}\par\smallskip{\footnotesize\textit{Hinweis:} die Lösung beginnt mit einem Schlag.}\par}

\clearpage
\phantomsection\label{hm:KRRNPvk}
\begin{tcolorbox}[enhanced,colback=ink,colframe=ink,arc=0pt,boxrule=0pt,left=4mm,right=4mm,top=2.5mm,bottom=2.5mm]
\color{white}\sffamily{\LARGE\bfseries KRRNPvk}\hfill{\Large h\#5}\\[0.5mm]
{\small 6 Steine\quad Schwarz am Zug}\hfill{\small Nr.~801\qquad\href{https://helpman.komtera.lt/?fen=8/8/8/8/5N2/8/1P5R/KR4k1&moves=5}{\color{white}\bfseries lösen\,\ensuremath{\nearrow}}}
\end{tcolorbox}
\vspace{2mm}
\noindent\begin{minipage}[t]{36mm}\vspace{0pt}\sffamily\small\raggedright
\hmfigure{längste eindeutige}{{\Large h\#5}}\hmfigure{eindeutig bei dieser Länge}{{\Large 1\,670} Stellungen}\hmfigure{längstes Matt}{{\Large h\#7}}\hmfigure{in der Klasse}{45,4 Milliarden Stellungen,\\ 14,2 Milliarden mit Hilfsmatt}
\end{minipage}\hfill
\begin{minipage}[t]{80mm}\vspace{0pt}\begin{tikzpicture}\begin{axis}[width=72mm,height=42mm,ybar,bar width=1.6pt,xmin=-1,xmax=15,ymin=0,ymax=6.7,axis y line*=left,axis x line*=bottom,xlabel={\scriptsize Halbzüge bis Matt},ylabel={\scriptsize Stellungen (Mrd.)},tick label style={font=\tiny},label style={font=\scriptsize},enlarge x limits=false,ylabel shift=-3pt,xlabel shift=-3pt]\addplot[fill=ink!55,draw=none] coordinates {(0,0.09622) (1,0.3729) (2,1.06) (3,2.751) (4,6.091) (5,2.001) (6,1.713) (7,0.02508) (8,0.04511) (9,3.961e-06) (10,0.006113) (11,2.4e-08) (12,8.505e-05) (14,1e-08)};\draw[densely dotted] (axis cs:10,0) -- (axis cs:10,6.7);\end{axis}\begin{axis}[width=72mm,height=42mm,axis y line*=right,axis x line=none,ymode=log,xmin=-1,xmax=15,ylabel={\scriptsize eindeutig (log)},tick label style={font=\tiny},label style={font=\scriptsize},ylabel shift=-3pt]\addplot[mark=*,mark size=0.7pt,draw=black,thin] coordinates {(0,96224079) (1,296508511) (2,630625035) (3,466745286) (4,483111759) (5,6793427) (6,7045396) (7,268) (8,228313) (9,1) (10,1670)};\end{axis}\end{tikzpicture}\end{minipage}
\vspace{3mm}

\noindent\begin{minipage}[t]{60mm}\vspace{0pt}\sbox\hmboard{\chessboard[setfen=8/8/8/8/5N2/8/1P5R/KR4k1 b - - 0 1]}\usebox\hmboard\\[1mm]
\makebox[\wd\hmboard]{\sffamily\hspace{3mm}\textbf{h\#5}\hfill(5+1)}\end{minipage}\hfill
\begin{minipage}[t]{52mm}\vspace{0pt}
\textcolor{ink}{\sffamily\textsc{Lösung}}\par\smallskip\small
\begin{tabular}{@{}r@{\enspace}l@{\enspace}l@{}}
\textbf{1.} & \symking{}xh2 & b4\\
\textbf{2.} & \symking{}g3 & \symrook{}b3+\\
\textbf{3.} & \symking{}f2 & \symking{}b2\\
\textbf{4.} & \symking{}e1 & \symking{}c3\\
\textbf{5.} & \symking{}d1 & \symrook{}b1\#\\
\end{tabular}\par\vspace{3mm}
\qrcode[height=12mm]{https://helpman.komtera.lt/?fen=8/8/8/8/5N2/8/1P5R/KR4k1&moves=5}\\{\sffamily\tiny\color{ink}scannen zum Lösen}
\end{minipage}
\vspace{2mm}

{\raggedright\noindent\textcolor{ink}{\sffamily\textsc{Themen}}\enspace \themetag{reines Matt} \themetag{Mustermatt} \themetag{Spiegelmatt} \themetag{Switchback} \themetag{ein Stein zieht}\par\smallskip{\footnotesize\textit{Hinweis:} das Diagramm hat keinen legalen letzten Zug, kann also in keiner Partie entstehen.}\par}

\clearpage
\phantomsection\label{hm:KRRNvkb}
\begin{tcolorbox}[enhanced,colback=ink,colframe=ink,arc=0pt,boxrule=0pt,left=4mm,right=4mm,top=2.5mm,bottom=2.5mm]
\color{white}\sffamily{\LARGE\bfseries KRRNvkb}\hfill{\Large h\#5}\\[0.5mm]
{\small 6 Steine\quad Schwarz am Zug}\hfill{\small Nr.~802\qquad\href{https://helpman.komtera.lt/?fen=7b/6R1/7k/3R4/8/8/1K6/N7&moves=5}{\color{white}\bfseries lösen\,\ensuremath{\nearrow}}}
\end{tcolorbox}
\vspace{2mm}
\noindent\begin{minipage}[t]{36mm}\vspace{0pt}\sffamily\small\raggedright
\hmfigure{längste eindeutige}{{\Large h\#5}}\hmfigure{eindeutig bei dieser Länge}{{\Large 1} Stellung}\hmfigure{längstes Matt}{{\Large h\#7}}\hmfigure{in der Klasse}{15,5 Milliarden Stellungen,\\ 4,39 Milliarden mit Hilfsmatt}
\end{minipage}\hfill
\begin{minipage}[t]{80mm}\vspace{0pt}\begin{tikzpicture}\begin{axis}[width=72mm,height=42mm,ybar,bar width=1.6pt,xmin=-1,xmax=15,ymin=0,ymax=2.038,axis y line*=left,axis x line*=bottom,xlabel={\scriptsize Halbzüge bis Matt},ylabel={\scriptsize Stellungen (Mrd.)},tick label style={font=\tiny},label style={font=\scriptsize},enlarge x limits=false,ylabel shift=-3pt,xlabel shift=-3pt]\addplot[fill=ink!55,draw=none] coordinates {(0,0.01893) (1,0.07349) (2,0.2778) (3,0.8863) (4,1.853) (5,0.7847) (6,0.4781) (7,0.01311) (8,0.007721) (9,6.11e-07) (10,0.0003625) (11,2.2e-08) (12,8.6e-08) (13,8e-09) (14,8e-09)};\draw[densely dotted] (axis cs:10,0) -- (axis cs:10,2.038);\end{axis}\begin{axis}[width=72mm,height=42mm,axis y line*=right,axis x line=none,ymode=log,xmin=-1,xmax=15,ylabel={\scriptsize eindeutig (log)},tick label style={font=\tiny},label style={font=\scriptsize},ylabel shift=-3pt]\addplot[mark=*,mark size=0.7pt,draw=black,thin] coordinates {(0,18933897) (1,64630260) (2,124113773) (3,161034566) (4,124761427) (5,2796823) (6,830715) (7,370) (8,21034) (10,1)};\end{axis}\end{tikzpicture}\end{minipage}
\vspace{3mm}

\noindent\begin{minipage}[t]{60mm}\vspace{0pt}\sbox\hmboard{\chessboard[setfen=7b/6R1/7k/3R4/8/8/1K6/N7 b - - 0 1]}\usebox\hmboard\\[1mm]
\makebox[\wd\hmboard]{\sffamily\hspace{3mm}\textbf{h\#5}\hfill(4+2)}\end{minipage}\hfill
\begin{minipage}[t]{52mm}\vspace{0pt}
\textcolor{ink}{\sffamily\textsc{Lösung}}\par\smallskip\small
\begin{tabular}{@{}r@{\enspace}l@{\enspace}l@{}}
\textbf{1.} & \symking{}xg7 & \symking{}c3\\
\textbf{2.} & \symking{}f6 & \symking{}d4\\
\textbf{3.} & \symking{}g7 & \symking{}e5\\
\textbf{4.} & \symking{}f8+ & \symking{}e6\\
\textbf{5.} & \symbishop{}g7 & \symrook{}d8\#\\
\end{tabular}\par\vspace{3mm}
\qrcode[height=12mm]{https://helpman.komtera.lt/?fen=7b/6R1/7k/3R4/8/8/1K6/N7&moves=5}\\{\sffamily\tiny\color{ink}scannen zum Lösen}
\end{minipage}
\vspace{2mm}

{\raggedright\noindent\textcolor{ink}{\sffamily\textsc{Themen}}\enspace \themetag{reines Matt} \themetag{Switchback} \themetag{Selbstblock}\par\smallskip{\footnotesize\textit{Hinweis:} die Lösung beginnt mit einem Schlag.}\par}

\clearpage
\phantomsection\label{hm:KRRNvkn}
\begin{tcolorbox}[enhanced,colback=ink,colframe=ink,arc=0pt,boxrule=0pt,left=4mm,right=4mm,top=2.5mm,bottom=2.5mm]
\color{white}\sffamily{\LARGE\bfseries KRRNvkn}\hfill{\Large h\#5}\\[0.5mm]
{\small 6 Steine\quad Schwarz am Zug}\hfill{\small Nr.~803\qquad\href{https://helpman.komtera.lt/?fen=8/8/6R1/7k/8/n7/8/K3N2R&moves=5}{\color{white}\bfseries lösen\,\ensuremath{\nearrow}}}
\end{tcolorbox}
\vspace{2mm}
\noindent\begin{minipage}[t]{36mm}\vspace{0pt}\sffamily\small\raggedright
\hmfigure{längste eindeutige}{{\Large h\#5}}\hmfigure{eindeutig bei dieser Länge}{{\Large 93} Stellungen}\hmfigure{längstes Matt}{{\Large h\#6}}\hmfigure{in der Klasse}{15,5 Milliarden Stellungen,\\ 4,53 Milliarden mit Hilfsmatt}
\end{minipage}\hfill
\begin{minipage}[t]{80mm}\vspace{0pt}\begin{tikzpicture}\begin{axis}[width=72mm,height=42mm,ybar,bar width=1.6pt,xmin=-1,xmax=13,ymin=0,ymax=2.12,axis y line*=left,axis x line*=bottom,xlabel={\scriptsize Halbzüge bis Matt},ylabel={\scriptsize Stellungen (Mrd.)},tick label style={font=\tiny},label style={font=\scriptsize},enlarge x limits=false,ylabel shift=-3pt,xlabel shift=-3pt]\addplot[fill=ink!55,draw=none] coordinates {(0,0.02386) (1,0.09333) (2,0.3186) (3,0.9245) (4,1.928) (5,0.7268) (6,0.4926) (7,0.01308) (8,0.009453) (9,1.019e-06) (10,0.0011) (11,3e-09) (12,5.734e-06)};\draw[densely dotted] (axis cs:10,0) -- (axis cs:10,2.12);\end{axis}\begin{axis}[width=72mm,height=42mm,axis y line*=right,axis x line=none,ymode=log,xmin=-1,xmax=13,ylabel={\scriptsize eindeutig (log)},tick label style={font=\tiny},label style={font=\scriptsize},ylabel shift=-3pt]\addplot[mark=*,mark size=0.7pt,draw=black,thin] coordinates {(0,23860001) (1,80014571) (2,140230294) (3,142468638) (4,117031429) (5,2184327) (6,869866) (7,1396) (8,47252) (10,93)};\end{axis}\end{tikzpicture}\end{minipage}
\vspace{3mm}

\noindent\begin{minipage}[t]{60mm}\vspace{0pt}\sbox\hmboard{\chessboard[setfen=8/8/6R1/7k/8/n7/8/K3N2R b - - 0 1]}\usebox\hmboard\\[1mm]
\makebox[\wd\hmboard]{\sffamily\hspace{3mm}\textbf{h\#5}\hfill(4+2)}\end{minipage}\hfill
\begin{minipage}[t]{52mm}\vspace{0pt}
\textcolor{ink}{\sffamily\textsc{Lösung}}\par\smallskip\small
\begin{tabular}{@{}r@{\enspace}l@{\enspace}l@{}}
\textbf{1.} & \symking{}xg6 & \symking{}b2\\
\textbf{2.} & \symknight{}c4+ & \symking{}c2\\
\textbf{3.} & \symking{}f5 & \symrook{}h5+\\
\textbf{4.} & \symking{}e4 & \symrook{}e5+\\
\textbf{5.} & \symking{}d4 & \symknight{}f3\#\\
\end{tabular}\par\vspace{3mm}
\qrcode[height=12mm]{https://helpman.komtera.lt/?fen=8/8/6R1/7k/8/n7/8/K3N2R&moves=5}\\{\sffamily\tiny\color{ink}scannen zum Lösen}
\end{minipage}
\vspace{2mm}

{\raggedright\noindent\textcolor{ink}{\sffamily\textsc{Themen}}\enspace \themetag{reines Matt} \themetag{Mustermatt} \themetag{Idealmatt} \themetag{Selbstblock}\par\smallskip{\footnotesize\textit{Hinweis:} die Partei am Zug steht im Diagramm im Schach; die Lösung beginnt mit einem Schlag.}\par}

\clearpage
\phantomsection\label{hm:KRRNvkp}
\begin{tcolorbox}[enhanced,colback=ink,colframe=ink,arc=0pt,boxrule=0pt,left=4mm,right=4mm,top=2.5mm,bottom=2.5mm]
\color{white}\sffamily{\LARGE\bfseries KRRNvkp}\hfill{\Large h\#5}\\[0.5mm]
{\small 6 Steine\quad Schwarz am Zug}\hfill{\small Nr.~804\qquad\href{https://helpman.komtera.lt/?fen=8/8/8/8/6R1/p7/NR6/K6k&moves=5}{\color{white}\bfseries lösen\,\ensuremath{\nearrow}}}
\end{tcolorbox}
\vspace{2mm}
\noindent\begin{minipage}[t]{36mm}\vspace{0pt}\sffamily\small\raggedright
\hmfigure{längste eindeutige}{{\Large h\#5}}\hmfigure{eindeutig bei dieser Länge}{{\Large 907} Stellungen}\hmfigure{längstes Matt}{{\Large h\#8}}\hmfigure{in der Klasse}{45,4 Milliarden Stellungen,\\ 14,1 Milliarden mit Hilfsmatt}
\end{minipage}\hfill
\begin{minipage}[t]{80mm}\vspace{0pt}\begin{tikzpicture}\begin{axis}[width=72mm,height=42mm,ybar,bar width=1.6pt,xmin=-1,xmax=17,ymin=0,ymax=6.613,axis y line*=left,axis x line*=bottom,xlabel={\scriptsize Halbzüge bis Matt},ylabel={\scriptsize Stellungen (Mrd.)},tick label style={font=\tiny},label style={font=\scriptsize},enlarge x limits=false,ylabel shift=-3pt,xlabel shift=-3pt]\addplot[fill=ink!55,draw=none] coordinates {(0,0.08982) (1,0.3433) (2,1.067) (3,2.908) (4,6.012) (5,2.003) (6,1.6) (7,0.03553) (8,0.03584) (9,1.17e-05) (10,0.005646) (11,7.36e-07) (12,0.0001087) (13,2e-09) (14,4.4e-08) (16,1e-09)};\draw[densely dotted] (axis cs:10,0) -- (axis cs:10,6.613);\end{axis}\begin{axis}[width=72mm,height=42mm,axis y line*=right,axis x line=none,ymode=log,xmin=-1,xmax=17,ylabel={\scriptsize eindeutig (log)},tick label style={font=\tiny},label style={font=\scriptsize},ylabel shift=-3pt]\addplot[mark=*,mark size=0.7pt,draw=black,thin] coordinates {(0,89817884) (1,290960938) (2,544545319) (3,478130827) (4,390820269) (5,7177203) (6,3810969) (7,2491) (8,200390) (9,13) (10,907)};\end{axis}\end{tikzpicture}\end{minipage}
\vspace{3mm}

\noindent\begin{minipage}[t]{60mm}\vspace{0pt}\sbox\hmboard{\chessboard[setfen=8/8/8/8/6R1/p7/NR6/K6k b - - 0 1]}\usebox\hmboard\\[1mm]
\makebox[\wd\hmboard]{\sffamily\hspace{3mm}\textbf{h\#5}\hfill(4+2)}\end{minipage}\hfill
\begin{minipage}[t]{52mm}\vspace{0pt}
\textcolor{ink}{\sffamily\textsc{Lösung}}\par\smallskip\small
\begin{tabular}{@{}r@{\enspace}l@{\enspace}l@{}}
\textbf{1.} & axb2+ & \symking{}b1\\
\textbf{2.} & \symking{}h2 & \symking{}c2\\
\textbf{3.} & b1=\symrook{} & \symknight{}c3\\
\textbf{4.} & \symrook{}g1 & \symknight{}e2\\
\textbf{5.} & \symrook{}g2 & \symrook{}h4\#\\
\end{tabular}\par\vspace{3mm}
\qrcode[height=12mm]{https://helpman.komtera.lt/?fen=8/8/8/8/6R1/p7/NR6/K6k&moves=5}\\{\sffamily\tiny\color{ink}scannen zum Lösen}
\end{minipage}
\vspace{2mm}

{\raggedright\noindent\textcolor{ink}{\sffamily\textsc{Themen}}\enspace \themetag{reines Matt} \themetag{Mustermatt} \themetag{Umwandlung} \themetag{Unterverwandlung} \themetag{Selbstblock} \themetag{Umnov}\par\smallskip{\footnotesize\textit{Hinweis:} die Lösung beginnt mit einem Schlag.}\par}

\clearpage
\phantomsection\label{hm:KRRNvkr}
\begin{tcolorbox}[enhanced,colback=ink,colframe=ink,arc=0pt,boxrule=0pt,left=4mm,right=4mm,top=2.5mm,bottom=2.5mm]
\color{white}\sffamily{\LARGE\bfseries KRRNvkr}\hfill{\Large h\#5}\\[0.5mm]
{\small 6 Steine\quad Schwarz am Zug}\hfill{\small Nr.~805\qquad\href{https://helpman.komtera.lt/?fen=8/8/N7/8/8/8/5R2/K4Rrk&moves=5}{\color{white}\bfseries lösen\,\ensuremath{\nearrow}}}
\end{tcolorbox}
\vspace{2mm}
\noindent\begin{minipage}[t]{36mm}\vspace{0pt}\sffamily\small\raggedright
\hmfigure{längste eindeutige}{{\Large h\#5}}\hmfigure{eindeutig bei dieser Länge}{{\Large 4} Stellungen}\hmfigure{längstes Matt}{{\Large h\#6.5}}\hmfigure{in der Klasse}{15,5 Milliarden Stellungen,\\ 4,18 Milliarden mit Hilfsmatt}
\end{minipage}\hfill
\begin{minipage}[t]{80mm}\vspace{0pt}\begin{tikzpicture}\begin{axis}[width=72mm,height=42mm,ybar,bar width=1.6pt,xmin=-1,xmax=14,ymin=0,ymax=1.97,axis y line*=left,axis x line*=bottom,xlabel={\scriptsize Halbzüge bis Matt},ylabel={\scriptsize Stellungen (Mrd.)},tick label style={font=\tiny},label style={font=\scriptsize},enlarge x limits=false,ylabel shift=-3pt,xlabel shift=-3pt]\addplot[fill=ink!55,draw=none] coordinates {(0,0.01659) (1,0.06647) (2,0.2741) (3,0.8819) (4,1.791) (5,0.7956) (6,0.3372) (7,0.01324) (8,0.005864) (9,4.097e-06) (10,7.565e-05) (11,4.56e-07) (12,7e-09) (13,2e-09)};\draw[densely dotted] (axis cs:10,0) -- (axis cs:10,1.97);\end{axis}\begin{axis}[width=72mm,height=42mm,axis y line*=right,axis x line=none,ymode=log,xmin=-1,xmax=14,ylabel={\scriptsize eindeutig (log)},tick label style={font=\tiny},label style={font=\scriptsize},ylabel shift=-3pt]\addplot[mark=*,mark size=0.7pt,draw=black,thin] coordinates {(0,16594570) (1,59686694) (2,128133455) (3,163216088) (4,114164541) (5,2436897) (6,812764) (7,591) (8,7035) (9,50) (10,4)};\end{axis}\end{tikzpicture}\end{minipage}
\vspace{3mm}

\noindent\begin{minipage}[t]{60mm}\vspace{0pt}\sbox\hmboard{\chessboard[setfen=8/8/N7/8/8/8/5R2/K4Rrk b - - 0 1]}\usebox\hmboard\\[1mm]
\makebox[\wd\hmboard]{\sffamily\hspace{3mm}\textbf{h\#5}\hfill(4+2)}\end{minipage}\hfill
\begin{minipage}[t]{52mm}\vspace{0pt}
\textcolor{ink}{\sffamily\textsc{Lösung}}\par\smallskip\small
\begin{tabular}{@{}r@{\enspace}l@{\enspace}l@{}}
\textbf{1.} & \symrook{}xf1+ & \symking{}b2\\
\textbf{2.} & \symking{}g1 & \symrook{}c2\\
\textbf{3.} & \symrook{}f2 & \symking{}c3\\
\textbf{4.} & \symking{}f1 & \symking{}d3\\
\textbf{5.} & \symking{}e1 & \symrook{}c1\#\\
\end{tabular}\par\vspace{3mm}
\qrcode[height=12mm]{https://helpman.komtera.lt/?fen=8/8/N7/8/8/8/5R2/K4Rrk&moves=5}\\{\sffamily\tiny\color{ink}scannen zum Lösen}
\end{minipage}
\vspace{2mm}

{\raggedright\noindent\textcolor{ink}{\sffamily\textsc{Themen}}\enspace \themetag{reines Matt} \themetag{Selbstblock} \themetag{Umnov}\par\smallskip{\footnotesize\textit{Hinweis:} die Lösung beginnt mit einem Schlag.}\par}

\clearpage
\phantomsection\label{hm:KQBvkqr}
\begin{tcolorbox}[enhanced,colback=ink,colframe=ink,arc=0pt,boxrule=0pt,left=4mm,right=4mm,top=2.5mm,bottom=2.5mm]
\color{white}\sffamily{\LARGE\bfseries KQBvkqr}\hfill{\Large h\#4.5}\\[0.5mm]
{\small 6 Steine\quad Weiß am Zug}\hfill{\small Nr.~806\qquad\href{https://helpman.komtera.lt/?fen=8/5q2/8/8/1Q6/2B5/2k5/K2r4&moves=4.5}{\color{white}\bfseries lösen\,\ensuremath{\nearrow}}}
\end{tcolorbox}
\vspace{2mm}
\noindent\begin{minipage}[t]{36mm}\vspace{0pt}\sffamily\small\raggedright
\hmfigure{längste eindeutige}{{\Large h\#4.5}}\hmfigure{eindeutig bei dieser Länge}{{\Large 51} Stellungen}\hmfigure{längstes Matt}{{\Large h\#6.5}}\hmfigure{in der Klasse}{15,5 Milliarden Stellungen,\\ 6,66 Milliarden mit Hilfsmatt}
\end{minipage}\hfill
\begin{minipage}[t]{80mm}\vspace{0pt}\begin{tikzpicture}\begin{axis}[width=72mm,height=42mm,ybar,bar width=1.6pt,xmin=-1,xmax=14,ymin=0,ymax=2.719,axis y line*=left,axis x line*=bottom,xlabel={\scriptsize Halbzüge bis Matt},ylabel={\scriptsize Stellungen (Mrd.)},tick label style={font=\tiny},label style={font=\scriptsize},enlarge x limits=false,ylabel shift=-3pt,xlabel shift=-3pt]\addplot[fill=ink!55,draw=none] coordinates {(0,0.006967) (1,0.0531) (2,0.3042) (3,1.298) (4,2.472) (5,2.102) (6,0.3716) (7,0.05393) (8,0.00285) (9,7.094e-06) (10,0.0002967) (11,1.83e-07) (12,1.2e-08) (13,1.6e-08)};\draw[densely dotted] (axis cs:9,0) -- (axis cs:9,2.719);\end{axis}\begin{axis}[width=72mm,height=42mm,axis y line*=right,axis x line=none,ymode=log,xmin=-1,xmax=14,ylabel={\scriptsize eindeutig (log)},tick label style={font=\tiny},label style={font=\scriptsize},ylabel shift=-3pt]\addplot[mark=*,mark size=0.7pt,draw=black,thin] coordinates {(0,6966924) (1,48564497) (2,180604484) (3,279417330) (4,119776353) (5,13602601) (6,149460) (7,259) (8,11489) (9,51)};\end{axis}\end{tikzpicture}\end{minipage}
\vspace{3mm}

\noindent\begin{minipage}[t]{60mm}\vspace{0pt}\sbox\hmboard{\chessboard[setfen=8/5q2/8/8/1Q6/2B5/2k5/K2r4 w - - 0 1]}\usebox\hmboard\\[1mm]
\makebox[\wd\hmboard]{\sffamily\hspace{3mm}\textbf{h\#4.5}\hfill(3+3)}\end{minipage}\hfill
\begin{minipage}[t]{52mm}\vspace{0pt}
\textcolor{ink}{\sffamily\textsc{Lösung}}\par\smallskip\small
\begin{tabular}{@{}r@{\enspace}l@{\enspace}l@{}}
\textbf{1.} & \ldots & \symqueen{}b1+\\
\textbf{2.} & \symking{}xc3 & \symqueen{}xd1\\
\textbf{3.} & \symking{}b4 & \symking{}b1\\
\textbf{4.} & \symking{}a3 & \symking{}c1\\
\textbf{5.} & \symking{}a2 & \symqueen{}a4\#\\
\end{tabular}\par\vspace{3mm}
\qrcode[height=12mm]{https://helpman.komtera.lt/?fen=8/5q2/8/8/1Q6/2B5/2k5/K2r4&moves=4.5}\\{\sffamily\tiny\color{ink}scannen zum Lösen}
\end{minipage}
\vspace{2mm}

{\raggedright\noindent\textcolor{ink}{\sffamily\textsc{Themen}}\enspace \themetag{reines Matt} \themetag{Mustermatt} \themetag{Spiegelmatt} \themetag{ein Stein zieht}\par\smallskip{\footnotesize\textit{Hinweis:} die Partei am Zug steht im Diagramm im Schach.}\par}

\clearpage
\phantomsection\label{hm:KQvkrrr}
\begin{tcolorbox}[enhanced,colback=ink,colframe=ink,arc=0pt,boxrule=0pt,left=4mm,right=4mm,top=2.5mm,bottom=2.5mm]
\color{white}\sffamily{\LARGE\bfseries KQvkrrr}\hfill{\Large h\#4.5}\\[0.5mm]
{\small 6 Steine\quad Weiß am Zug}\hfill{\small Nr.~807\qquad\href{https://helpman.komtera.lt/?fen=8/8/4k3/8/5Q2/r7/rr6/2K5&moves=4.5}{\color{white}\bfseries lösen\,\ensuremath{\nearrow}}}
\end{tcolorbox}
\vspace{2mm}
\noindent\begin{minipage}[t]{36mm}\vspace{0pt}\sffamily\small\raggedright
\hmfigure{längste eindeutige}{{\Large h\#4.5}}\hmfigure{eindeutig bei dieser Länge}{{\Large 577} Stellungen}\hmfigure{längstes Matt}{{\Large h\#6}}\hmfigure{in der Klasse}{15,5 Milliarden Stellungen,\\ 1,15 Milliarden mit Hilfsmatt}
\end{minipage}\hfill
\begin{minipage}[t]{80mm}\vspace{0pt}\begin{tikzpicture}\begin{axis}[width=72mm,height=42mm,ybar,bar width=1.6pt,xmin=-1,xmax=13,ymin=0,ymax=381.5,axis y line*=left,axis x line*=bottom,xlabel={\scriptsize Halbzüge bis Matt},ylabel={\scriptsize Stellungen (Mio.)},tick label style={font=\tiny},label style={font=\scriptsize},enlarge x limits=false,ylabel shift=-3pt,xlabel shift=-3pt]\addplot[fill=ink!55,draw=none] coordinates {(0,0.2825) (1,2.745) (2,15.91) (3,75.72) (4,198.5) (5,346.8) (6,234.1) (7,226.2) (8,35.29) (9,16.76) (10,0.4925) (11,0.01233) (12,2.6e-05)};\draw[densely dotted] (axis cs:9,0) -- (axis cs:9,381.5);\end{axis}\begin{axis}[width=72mm,height=42mm,axis y line*=right,axis x line=none,ymode=log,xmin=-1,xmax=13,ylabel={\scriptsize eindeutig (log)},tick label style={font=\tiny},label style={font=\scriptsize},ylabel shift=-3pt]\addplot[mark=*,mark size=0.7pt,draw=black,thin] coordinates {(0,282465) (1,2434768) (2,9340361) (3,15236354) (4,11669091) (5,8274616) (6,1258396) (7,555259) (8,17514) (9,577)};\end{axis}\end{tikzpicture}\end{minipage}
\vspace{3mm}

\noindent\begin{minipage}[t]{60mm}\vspace{0pt}\sbox\hmboard{\chessboard[setfen=8/8/4k3/8/5Q2/r7/rr6/2K5 w - - 0 1]}\usebox\hmboard\\[1mm]
\makebox[\wd\hmboard]{\sffamily\hspace{3mm}\textbf{h\#4.5}\hfill(2+4)}\end{minipage}\hfill
\begin{minipage}[t]{52mm}\vspace{0pt}
\textcolor{ink}{\sffamily\textsc{Lösung}}\par\smallskip\small
\begin{tabular}{@{}r@{\enspace}l@{\enspace}l@{}}
\textbf{1.} & \ldots & \symking{}d1\\
\textbf{2.} & \symking{}d5 & \symqueen{}d2+\\
\textbf{3.} & \symking{}c4 & \symking{}e2\\
\textbf{4.} & \symking{}b3 & \symking{}d3\\
\textbf{5.} & \symrook{}a4 & \symqueen{}c3\#\\
\end{tabular}\par\vspace{3mm}
\qrcode[height=12mm]{https://helpman.komtera.lt/?fen=8/8/4k3/8/5Q2/r7/rr6/2K5&moves=4.5}\\{\sffamily\tiny\color{ink}scannen zum Lösen}
\end{minipage}
\vspace{2mm}

{\raggedright\noindent\textcolor{ink}{\sffamily\textsc{Themen}}\enspace \themetag{Selbstblock} \themetag{ohne Schlag}\par}

\clearpage
\phantomsection\label{hm:KRRBvkr}
\begin{tcolorbox}[enhanced,colback=ink,colframe=ink,arc=0pt,boxrule=0pt,left=4mm,right=4mm,top=2.5mm,bottom=2.5mm]
\color{white}\sffamily{\LARGE\bfseries KRRBvkr}\hfill{\Large h\#4.5}\\[0.5mm]
{\small 6 Steine\quad Weiß am Zug}\hfill{\small Nr.~808\qquad\href{https://helpman.komtera.lt/?fen=8/8/8/8/8/R2R4/rBk5/K7&moves=4.5}{\color{white}\bfseries lösen\,\ensuremath{\nearrow}}}
\end{tcolorbox}
\vspace{2mm}
\noindent\begin{minipage}[t]{36mm}\vspace{0pt}\sffamily\small\raggedright
\hmfigure{längste eindeutige}{{\Large h\#4.5}}\hmfigure{eindeutig bei dieser Länge}{{\Large 43} Stellungen}\hmfigure{längstes Matt}{{\Large h\#7.5}}\hmfigure{in der Klasse}{15,5 Milliarden Stellungen,\\ 4,1 Milliarden mit Hilfsmatt}
\end{minipage}\hfill
\begin{minipage}[t]{80mm}\vspace{0pt}\begin{tikzpicture}\begin{axis}[width=72mm,height=42mm,ybar,bar width=1.6pt,xmin=-1,xmax=16,ymin=0,ymax=2.013,axis y line*=left,axis x line*=bottom,xlabel={\scriptsize Halbzüge bis Matt},ylabel={\scriptsize Stellungen (Mrd.)},tick label style={font=\tiny},label style={font=\scriptsize},enlarge x limits=false,ylabel shift=-3pt,xlabel shift=-3pt]\addplot[fill=ink!55,draw=none] coordinates {(0,0.01845) (1,0.06888) (2,0.2866) (3,0.8668) (4,1.83) (5,0.7319) (6,0.285) (7,0.007591) (8,0.004299) (9,4.939e-06) (10,1.288e-05) (11,4.83e-07) (12,6e-09) (13,7.5e-08) (15,1e-09)};\draw[densely dotted] (axis cs:9,0) -- (axis cs:9,2.013);\end{axis}\begin{axis}[width=72mm,height=42mm,axis y line*=right,axis x line=none,ymode=log,xmin=-1,xmax=16,ylabel={\scriptsize eindeutig (log)},tick label style={font=\tiny},label style={font=\scriptsize},ylabel shift=-3pt]\addplot[mark=*,mark size=0.7pt,draw=black,thin] coordinates {(0,18453462) (1,61481732) (2,134203405) (3,159125675) (4,121509962) (5,1762716) (6,705439) (7,958) (8,2484) (9,43)};\end{axis}\end{tikzpicture}\end{minipage}
\vspace{3mm}

\noindent\begin{minipage}[t]{60mm}\vspace{0pt}\sbox\hmboard{\chessboard[setfen=8/8/8/8/8/R2R4/rBk5/K7 w - - 0 1]}\usebox\hmboard\\[1mm]
\makebox[\wd\hmboard]{\sffamily\hspace{3mm}\textbf{h\#4.5}\hfill(4+2)}\end{minipage}\hfill
\begin{minipage}[t]{52mm}\vspace{0pt}
\textcolor{ink}{\sffamily\textsc{Lösung}}\par\smallskip\small
\begin{tabular}{@{}r@{\enspace}l@{\enspace}l@{}}
\textbf{1.} & \ldots & \symrook{}xa2\\
\textbf{2.} & \symking{}xd3 & \symking{}b1\\
\textbf{3.} & \symking{}e2 & \symking{}c2\\
\textbf{4.} & \symking{}e1 & \symking{}d3\\
\textbf{5.} & \symking{}d1 & \symrook{}a1\#\\
\end{tabular}\par\vspace{3mm}
\qrcode[height=12mm]{https://helpman.komtera.lt/?fen=8/8/8/8/8/R2R4/rBk5/K7&moves=4.5}\\{\sffamily\tiny\color{ink}scannen zum Lösen}
\end{minipage}
\vspace{2mm}

{\raggedright\noindent\textcolor{ink}{\sffamily\textsc{Themen}}\enspace \themetag{Spiegelmatt} \themetag{ein Stein zieht} \themetag{ohne Schach}\par\smallskip{\footnotesize\textit{Hinweis:} das Diagramm hat keinen legalen letzten Zug, kann also in keiner Partie entstehen.}\par}

\clearpage
\phantomsection\label{hm:KQQBvkp}
\begin{tcolorbox}[enhanced,colback=ink,colframe=ink,arc=0pt,boxrule=0pt,left=4mm,right=4mm,top=2.5mm,bottom=2.5mm]
\color{white}\sffamily{\LARGE\bfseries KQQBvkp}\hfill{\Large h\#4}\\[0.5mm]
{\small 6 Steine\quad Schwarz am Zug}\hfill{\small Nr.~809\qquad\href{https://helpman.komtera.lt/?fen=8/8/8/8/8/2k1p3/3QB3/1QK5&moves=4}{\color{white}\bfseries lösen\,\ensuremath{\nearrow}}}
\end{tcolorbox}
\vspace{2mm}
\noindent\begin{minipage}[t]{36mm}\vspace{0pt}\sffamily\small\raggedright
\hmfigure{längste eindeutige}{{\Large h\#4}}\hmfigure{eindeutig bei dieser Länge}{{\Large 6} Stellungen}\hmfigure{längstes Matt}{{\Large h\#5}}\hmfigure{in der Klasse}{45,4 Milliarden Stellungen,\\ 12,3 Milliarden mit Hilfsmatt}
\end{minipage}\hfill
\begin{minipage}[t]{80mm}\vspace{0pt}\begin{tikzpicture}\begin{axis}[width=72mm,height=42mm,ybar,bar width=1.6pt,xmin=-1,xmax=11,ymin=0,ymax=4.883,axis y line*=left,axis x line*=bottom,xlabel={\scriptsize Halbzüge bis Matt},ylabel={\scriptsize Stellungen (Mrd.)},tick label style={font=\tiny},label style={font=\scriptsize},enlarge x limits=false,ylabel shift=-3pt,xlabel shift=-3pt]\addplot[fill=ink!55,draw=none] coordinates {(0,0.4363) (1,1.072) (2,4.439) (3,2.357) (4,3.799) (5,0.04764) (6,0.1118) (7,1.39e-07) (8,0.000211) (10,1.4e-08)};\draw[densely dotted] (axis cs:8,0) -- (axis cs:8,4.883);\end{axis}\begin{axis}[width=72mm,height=42mm,axis y line*=right,axis x line=none,ymode=log,xmin=-1,xmax=11,ylabel={\scriptsize eindeutig (log)},tick label style={font=\tiny},label style={font=\scriptsize},ylabel shift=-3pt]\addplot[mark=*,mark size=0.7pt,draw=black,thin] coordinates {(0,436340760) (1,485614567) (2,1312911724) (3,19178116) (4,60267808) (5,54) (6,45119) (8,6)};\end{axis}\end{tikzpicture}\end{minipage}
\vspace{3mm}

\noindent\begin{minipage}[t]{60mm}\vspace{0pt}\sbox\hmboard{\chessboard[setfen=8/8/8/8/8/2k1p3/3QB3/1QK5 b - - 0 1]}\usebox\hmboard\\[1mm]
\makebox[\wd\hmboard]{\sffamily\hspace{3mm}\textbf{h\#4}\hfill(4+2)}\end{minipage}\hfill
\begin{minipage}[t]{52mm}\vspace{0pt}
\textcolor{ink}{\sffamily\textsc{Lösung}}\par\smallskip\small
\begin{tabular}{@{}r@{\enspace}l@{\enspace}l@{}}
\textbf{1.} & exd2+ & \symking{}d1\\
\textbf{2.} & \symking{}d4 & \symking{}c2\\
\textbf{3.} & d1=\symbishop{}+ & \symking{}c1\\
\textbf{4.} & \symking{}c3 & \symqueen{}b2\#\\
\end{tabular}\par\vspace{3mm}
\qrcode[height=12mm]{https://helpman.komtera.lt/?fen=8/8/8/8/8/2k1p3/3QB3/1QK5&moves=4}\\{\sffamily\tiny\color{ink}scannen zum Lösen}
\end{minipage}
\vspace{2mm}

{\raggedright\noindent\textcolor{ink}{\sffamily\textsc{Themen}}\enspace \themetag{Umwandlung} \themetag{Unterverwandlung} \themetag{Switchback} \themetag{Rundlauf} \themetag{Umnov}\par\smallskip{\footnotesize\textit{Hinweis:} die Partei am Zug steht im Diagramm im Schach; die Lösung beginnt mit einem Schlag.}\par}

\clearpage
\phantomsection\label{hm:KQQBvkq}
\begin{tcolorbox}[enhanced,colback=ink,colframe=ink,arc=0pt,boxrule=0pt,left=4mm,right=4mm,top=2.5mm,bottom=2.5mm]
\color{white}\sffamily{\LARGE\bfseries KQQBvkq}\hfill{\Large h\#4}\\[0.5mm]
{\small 6 Steine\quad Schwarz am Zug}\hfill{\small Nr.~810\qquad\href{https://helpman.komtera.lt/?fen=8/8/8/8/2Q5/1q6/8/KQ1kB3&moves=4}{\color{white}\bfseries lösen\,\ensuremath{\nearrow}}}
\end{tcolorbox}
\vspace{2mm}
\noindent\begin{minipage}[t]{36mm}\vspace{0pt}\sffamily\small\raggedright
\hmfigure{längste eindeutige}{{\Large h\#4}}\hmfigure{eindeutig bei dieser Länge}{{\Large 10} Stellungen}\hmfigure{längstes Matt}{{\Large h\#5}}\hmfigure{in der Klasse}{15,5 Milliarden Stellungen,\\ 3,18 Milliarden mit Hilfsmatt}
\end{minipage}\hfill
\begin{minipage}[t]{80mm}\vspace{0pt}\begin{tikzpicture}\begin{axis}[width=72mm,height=42mm,ybar,bar width=1.6pt,xmin=-1,xmax=11,ymin=0,ymax=1079,axis y line*=left,axis x line*=bottom,xlabel={\scriptsize Halbzüge bis Matt},ylabel={\scriptsize Stellungen (Mio.)},tick label style={font=\tiny},label style={font=\scriptsize},enlarge x limits=false,ylabel shift=-3pt,xlabel shift=-3pt]\addplot[fill=ink!55,draw=none] coordinates {(0,70.48) (1,196) (2,981) (3,839.4) (4,951.3) (5,111.7) (6,32.55) (7,0.000766) (8,0.06691) (9,1.9e-05) (10,0.000274)};\draw[densely dotted] (axis cs:8,0) -- (axis cs:8,1079);\end{axis}\begin{axis}[width=72mm,height=42mm,axis y line*=right,axis x line=none,ymode=log,xmin=-1,xmax=11,ylabel={\scriptsize eindeutig (log)},tick label style={font=\tiny},label style={font=\scriptsize},ylabel shift=-3pt]\addplot[mark=*,mark size=0.7pt,draw=black,thin] coordinates {(0,70484096) (1,111070549) (2,308517297) (3,29432819) (4,11786145) (5,223) (6,24943) (7,23) (8,10)};\end{axis}\end{tikzpicture}\end{minipage}
\vspace{3mm}

\noindent\begin{minipage}[t]{60mm}\vspace{0pt}\sbox\hmboard{\chessboard[setfen=8/8/8/8/2Q5/1q6/8/KQ1kB3 b - - 0 1]}\usebox\hmboard\\[1mm]
\makebox[\wd\hmboard]{\sffamily\hspace{3mm}\textbf{h\#4}\hfill(4+2)}\end{minipage}\hfill
\begin{minipage}[t]{52mm}\vspace{0pt}
\textcolor{ink}{\sffamily\textsc{Lösung}}\par\smallskip\small
\begin{tabular}{@{}r@{\enspace}l@{\enspace}l@{}}
\textbf{1.} & \symqueen{}xb1+ & \symking{}xb1\\
\textbf{2.} & \symking{}xe1 & \symking{}b2\\
\textbf{3.} & \symking{}d1 & \symking{}c3\\
\textbf{4.} & \symking{}c1 & \symqueen{}f1\#\\
\end{tabular}\par\vspace{3mm}
\qrcode[height=12mm]{https://helpman.komtera.lt/?fen=8/8/8/8/2Q5/1q6/8/KQ1kB3&moves=4}\\{\sffamily\tiny\color{ink}scannen zum Lösen}
\end{minipage}
\vspace{2mm}

{\raggedright\noindent\textcolor{ink}{\sffamily\textsc{Themen}}\enspace \themetag{reines Matt} \themetag{Mustermatt} \themetag{Idealmatt} \themetag{Spiegelmatt} \themetag{Switchback}\par\smallskip{\footnotesize\textit{Hinweis:} die Partei am Zug steht im Diagramm im Schach; die Lösung beginnt mit einem Schlag.}\par}

\clearpage
\phantomsection\label{hm:KQQNPvk}
\begin{tcolorbox}[enhanced,colback=ink,colframe=ink,arc=0pt,boxrule=0pt,left=4mm,right=4mm,top=2.5mm,bottom=2.5mm]
\color{white}\sffamily{\LARGE\bfseries KQQNPvk}\hfill{\Large h\#4}\\[0.5mm]
{\small 6 Steine\quad Schwarz am Zug}\hfill{\small Nr.~811\qquad\href{https://helpman.komtera.lt/?fen=N7/8/8/8/3Q4/8/1P6/K5kQ&moves=4}{\color{white}\bfseries lösen\,\ensuremath{\nearrow}}}
\end{tcolorbox}
\vspace{2mm}
\noindent\begin{minipage}[t]{36mm}\vspace{0pt}\sffamily\small\raggedright
\hmfigure{längste eindeutige}{{\Large h\#4}}\hmfigure{eindeutig bei dieser Länge}{{\Large 1\,607} Stellungen}\hmfigure{längstes Matt}{{\Large h\#5}}\hmfigure{in der Klasse}{45,4 Milliarden Stellungen,\\ 12,5 Milliarden mit Hilfsmatt}
\end{minipage}\hfill
\begin{minipage}[t]{80mm}\vspace{0pt}\begin{tikzpicture}\begin{axis}[width=72mm,height=42mm,ybar,bar width=1.6pt,xmin=-1,xmax=11,ymin=0,ymax=4.731,axis y line*=left,axis x line*=bottom,xlabel={\scriptsize Halbzüge bis Matt},ylabel={\scriptsize Stellungen (Mrd.)},tick label style={font=\tiny},label style={font=\scriptsize},enlarge x limits=false,ylabel shift=-3pt,xlabel shift=-3pt]\addplot[fill=ink!55,draw=none] coordinates {(0,0.3981) (1,1.067) (2,3.965) (3,2.499) (4,4.301) (5,0.0519) (6,0.178) (7,1.9e-08) (8,0.01347) (10,1.515e-05)};\draw[densely dotted] (axis cs:8,0) -- (axis cs:8,4.731);\end{axis}\begin{axis}[width=72mm,height=42mm,axis y line*=right,axis x line=none,ymode=log,xmin=-1,xmax=11,ylabel={\scriptsize eindeutig (log)},tick label style={font=\tiny},label style={font=\scriptsize},ylabel shift=-3pt]\addplot[mark=*,mark size=0.7pt,draw=black,thin] coordinates {(0,398113791) (1,508218001) (2,1412592019) (3,38601225) (4,32639022) (5,8) (6,1826033) (8,1607)};\end{axis}\end{tikzpicture}\end{minipage}
\vspace{3mm}

\noindent\begin{minipage}[t]{60mm}\vspace{0pt}\sbox\hmboard{\chessboard[setfen=N7/8/8/8/3Q4/8/1P6/K5kQ b - - 0 1]}\usebox\hmboard\\[1mm]
\makebox[\wd\hmboard]{\sffamily\hspace{3mm}\textbf{h\#4}\hfill(5+1)}\end{minipage}\hfill
\begin{minipage}[t]{52mm}\vspace{0pt}
\textcolor{ink}{\sffamily\textsc{Lösung}}\par\smallskip\small
\begin{tabular}{@{}r@{\enspace}l@{\enspace}l@{}}
\textbf{1.} & \symking{}xh1 & \symking{}b1\\
\textbf{2.} & \symking{}g2 & \symking{}c2\\
\textbf{3.} & \symking{}f1 & \symking{}d3\\
\textbf{4.} & \symking{}e1 & \symqueen{}g1\#\\
\end{tabular}\par\vspace{3mm}
\qrcode[height=12mm]{https://helpman.komtera.lt/?fen=N7/8/8/8/3Q4/8/1P6/K5kQ&moves=4}\\{\sffamily\tiny\color{ink}scannen zum Lösen}
\end{minipage}
\vspace{2mm}

{\raggedright\noindent\textcolor{ink}{\sffamily\textsc{Themen}}\enspace \themetag{reines Matt} \themetag{Spiegelmatt} \themetag{ein Stein zieht} \themetag{ohne Schach}\par\smallskip{\footnotesize\textit{Hinweis:} das Diagramm hat keinen legalen letzten Zug, kann also in keiner Partie entstehen.}\par}

\clearpage
\phantomsection\label{hm:KQQNvkb}
\begin{tcolorbox}[enhanced,colback=ink,colframe=ink,arc=0pt,boxrule=0pt,left=4mm,right=4mm,top=2.5mm,bottom=2.5mm]
\color{white}\sffamily{\LARGE\bfseries KQQNvkb}\hfill{\Large h\#4}\\[0.5mm]
{\small 6 Steine\quad Schwarz am Zug}\hfill{\small Nr.~812\qquad\href{https://helpman.komtera.lt/?fen=N7/2b5/8/8/8/4Q3/8/K5kQ&moves=4}{\color{white}\bfseries lösen\,\ensuremath{\nearrow}}}
\end{tcolorbox}
\vspace{2mm}
\noindent\begin{minipage}[t]{36mm}\vspace{0pt}\sffamily\small\raggedright
\hmfigure{längste eindeutige}{{\Large h\#4}}\hmfigure{eindeutig bei dieser Länge}{{\Large 552} Stellungen}\hmfigure{längstes Matt}{{\Large h\#5}}\hmfigure{in der Klasse}{15,5 Milliarden Stellungen,\\ 3,87 Milliarden mit Hilfsmatt}
\end{minipage}\hfill
\begin{minipage}[t]{80mm}\vspace{0pt}\begin{tikzpicture}\begin{axis}[width=72mm,height=42mm,ybar,bar width=1.6pt,xmin=-1,xmax=11,ymin=0,ymax=1.364,axis y line*=left,axis x line*=bottom,xlabel={\scriptsize Halbzüge bis Matt},ylabel={\scriptsize Stellungen (Mrd.)},tick label style={font=\tiny},label style={font=\scriptsize},enlarge x limits=false,ylabel shift=-3pt,xlabel shift=-3pt]\addplot[fill=ink!55,draw=none] coordinates {(0,0.09193) (1,0.2779) (2,1.24) (3,0.8979) (4,1.238) (5,0.05752) (6,0.05728) (7,1.003e-06) (8,0.005464) (9,9e-09) (10,4.832e-06)};\draw[densely dotted] (axis cs:8,0) -- (axis cs:8,1.364);\end{axis}\begin{axis}[width=72mm,height=42mm,axis y line*=right,axis x line=none,ymode=log,xmin=-1,xmax=11,ylabel={\scriptsize eindeutig (log)},tick label style={font=\tiny},label style={font=\scriptsize},ylabel shift=-3pt]\addplot[mark=*,mark size=0.7pt,draw=black,thin] coordinates {(0,91929402) (1,149404113) (2,379063599) (3,15503209) (4,8638400) (5,297) (6,634535) (7,1) (8,552)};\end{axis}\end{tikzpicture}\end{minipage}
\vspace{3mm}

\noindent\begin{minipage}[t]{60mm}\vspace{0pt}\sbox\hmboard{\chessboard[setfen=N7/2b5/8/8/8/4Q3/8/K5kQ b - - 0 1]}\usebox\hmboard\\[1mm]
\makebox[\wd\hmboard]{\sffamily\hspace{3mm}\textbf{h\#4}\hfill(4+2)}\end{minipage}\hfill
\begin{minipage}[t]{52mm}\vspace{0pt}
\textcolor{ink}{\sffamily\textsc{Lösung}}\par\smallskip\small
\begin{tabular}{@{}r@{\enspace}l@{\enspace}l@{}}
\textbf{1.} & \symking{}xh1 & \symknight{}b6\\
\textbf{2.} & \symbishop{}h2 & \symknight{}d5\\
\textbf{3.} & \symbishop{}g1 & \symknight{}f4\\
\textbf{4.} & \symking{}h2 & \symqueen{}h3\#\\
\end{tabular}\par\vspace{3mm}
\qrcode[height=12mm]{https://helpman.komtera.lt/?fen=N7/2b5/8/8/8/4Q3/8/K5kQ&moves=4}\\{\sffamily\tiny\color{ink}scannen zum Lösen}
\end{minipage}
\vspace{2mm}

{\raggedright\noindent\textcolor{ink}{\sffamily\textsc{Themen}}\enspace \themetag{Selbstblock} \themetag{ohne Schach}\par\smallskip{\footnotesize\textit{Hinweis:} das Diagramm hat keinen legalen letzten Zug, kann also in keiner Partie entstehen.}\par}

\clearpage
\phantomsection\label{hm:KQQNvkn}
\begin{tcolorbox}[enhanced,colback=ink,colframe=ink,arc=0pt,boxrule=0pt,left=4mm,right=4mm,top=2.5mm,bottom=2.5mm]
\color{white}\sffamily{\LARGE\bfseries KQQNvkn}\hfill{\Large h\#4}\\[0.5mm]
{\small 6 Steine\quad Schwarz am Zug}\hfill{\small Nr.~813\qquad\href{https://helpman.komtera.lt/?fen=N7/4Q3/8/8/8/8/8/K2n1kQ1&moves=4}{\color{white}\bfseries lösen\,\ensuremath{\nearrow}}}
\end{tcolorbox}
\vspace{2mm}
\noindent\begin{minipage}[t]{36mm}\vspace{0pt}\sffamily\small\raggedright
\hmfigure{längste eindeutige}{{\Large h\#4}}\hmfigure{eindeutig bei dieser Länge}{{\Large 891} Stellungen}\hmfigure{längstes Matt}{{\Large h\#5}}\hmfigure{in der Klasse}{15,5 Milliarden Stellungen,\\ 4 Milliarden mit Hilfsmatt}
\end{minipage}\hfill
\begin{minipage}[t]{80mm}\vspace{0pt}\begin{tikzpicture}\begin{axis}[width=72mm,height=42mm,ybar,bar width=1.6pt,xmin=-1,xmax=11,ymin=0,ymax=1.464,axis y line*=left,axis x line*=bottom,xlabel={\scriptsize Halbzüge bis Matt},ylabel={\scriptsize Stellungen (Mrd.)},tick label style={font=\tiny},label style={font=\scriptsize},enlarge x limits=false,ylabel shift=-3pt,xlabel shift=-3pt]\addplot[fill=ink!55,draw=none] coordinates {(0,0.1013) (1,0.2946) (2,1.269) (3,0.8908) (4,1.331) (5,0.04781) (6,0.06232) (7,3.747e-06) (8,0.006923) (9,1.4e-08) (10,9.118e-06)};\draw[densely dotted] (axis cs:8,0) -- (axis cs:8,1.464);\end{axis}\begin{axis}[width=72mm,height=42mm,axis y line*=right,axis x line=none,ymode=log,xmin=-1,xmax=11,ylabel={\scriptsize eindeutig (log)},tick label style={font=\tiny},label style={font=\scriptsize},ylabel shift=-3pt]\addplot[mark=*,mark size=0.7pt,draw=black,thin] coordinates {(0,101282424) (1,154069016) (2,381454168) (3,12231925) (4,8801038) (5,1002) (6,667140) (7,20) (8,891)};\end{axis}\end{tikzpicture}\end{minipage}
\vspace{3mm}

\noindent\begin{minipage}[t]{60mm}\vspace{0pt}\sbox\hmboard{\chessboard[setfen=N7/4Q3/8/8/8/8/8/K2n1kQ1 b - - 0 1]}\usebox\hmboard\\[1mm]
\makebox[\wd\hmboard]{\sffamily\hspace{3mm}\textbf{h\#4}\hfill(4+2)}\end{minipage}\hfill
\begin{minipage}[t]{52mm}\vspace{0pt}
\textcolor{ink}{\sffamily\textsc{Lösung}}\par\smallskip\small
\begin{tabular}{@{}r@{\enspace}l@{\enspace}l@{}}
\textbf{1.} & \symking{}xg1 & \symknight{}b6\\
\textbf{2.} & \symknight{}e3 & \symknight{}c4\\
\textbf{3.} & \symknight{}g2 & \symqueen{}e3+\\
\textbf{4.} & \symking{}f1 & \symknight{}d2\#\\
\end{tabular}\par\vspace{3mm}
\qrcode[height=12mm]{https://helpman.komtera.lt/?fen=N7/4Q3/8/8/8/8/8/K2n1kQ1&moves=4}\\{\sffamily\tiny\color{ink}scannen zum Lösen}
\end{minipage}
\vspace{2mm}

{\raggedright\noindent\textcolor{ink}{\sffamily\textsc{Themen}}\enspace \themetag{reines Matt} \themetag{Mustermatt} \themetag{Switchback} \themetag{Selbstblock} \themetag{Umnov}\par\smallskip{\footnotesize\textit{Hinweis:} die Partei am Zug steht im Diagramm im Schach; die Lösung beginnt mit einem Schlag.}\par}

\clearpage
\phantomsection\label{hm:KQQNvkp}
\begin{tcolorbox}[enhanced,colback=ink,colframe=ink,arc=0pt,boxrule=0pt,left=4mm,right=4mm,top=2.5mm,bottom=2.5mm]
\color{white}\sffamily{\LARGE\bfseries KQQNvkp}\hfill{\Large h\#4}\\[0.5mm]
{\small 6 Steine\quad Schwarz am Zug}\hfill{\small Nr.~814\qquad\href{https://helpman.komtera.lt/?fen=N7/8/8/5Q1p/8/8/7Q/K5k1&moves=4}{\color{white}\bfseries lösen\,\ensuremath{\nearrow}}}
\end{tcolorbox}
\vspace{2mm}
\noindent\begin{minipage}[t]{36mm}\vspace{0pt}\sffamily\small\raggedright
\hmfigure{längste eindeutige}{{\Large h\#4}}\hmfigure{eindeutig bei dieser Länge}{{\Large 1\,680} Stellungen}\hmfigure{längstes Matt}{{\Large h\#5}}\hmfigure{in der Klasse}{45,4 Milliarden Stellungen,\\ 12,5 Milliarden mit Hilfsmatt}
\end{minipage}\hfill
\begin{minipage}[t]{80mm}\vspace{0pt}\begin{tikzpicture}\begin{axis}[width=72mm,height=42mm,ybar,bar width=1.6pt,xmin=-1,xmax=11,ymin=0,ymax=4.572,axis y line*=left,axis x line*=bottom,xlabel={\scriptsize Halbzüge bis Matt},ylabel={\scriptsize Stellungen (Mrd.)},tick label style={font=\tiny},label style={font=\scriptsize},enlarge x limits=false,ylabel shift=-3pt,xlabel shift=-3pt]\addplot[fill=ink!55,draw=none] coordinates {(0,0.3577) (1,1.025) (2,4.156) (3,2.637) (4,4.074) (5,0.06559) (6,0.184) (7,1.173e-06) (8,0.01996) (9,4e-09) (10,2.143e-05)};\draw[densely dotted] (axis cs:8,0) -- (axis cs:8,4.572);\end{axis}\begin{axis}[width=72mm,height=42mm,axis y line*=right,axis x line=none,ymode=log,xmin=-1,xmax=11,ylabel={\scriptsize eindeutig (log)},tick label style={font=\tiny},label style={font=\scriptsize},ylabel shift=-3pt]\addplot[mark=*,mark size=0.7pt,draw=black,thin] coordinates {(0,357739667) (1,507369661) (2,1339423513) (3,27873718) (4,29743107) (5,117) (6,1985904) (8,1680)};\end{axis}\end{tikzpicture}\end{minipage}
\vspace{3mm}

\noindent\begin{minipage}[t]{60mm}\vspace{0pt}\sbox\hmboard{\chessboard[setfen=N7/8/8/5Q1p/8/8/7Q/K5k1 b - - 0 1]}\usebox\hmboard\\[1mm]
\makebox[\wd\hmboard]{\sffamily\hspace{3mm}\textbf{h\#4}\hfill(4+2)}\end{minipage}\hfill
\begin{minipage}[t]{52mm}\vspace{0pt}
\textcolor{ink}{\sffamily\textsc{Lösung}}\par\smallskip\small
\begin{tabular}{@{}r@{\enspace}l@{\enspace}l@{}}
\textbf{1.} & \symking{}xh2 & \symknight{}c7\\
\textbf{2.} & \symking{}g2 & \symknight{}e6\\
\textbf{3.} & \symking{}g3 & \symknight{}g5\\
\textbf{4.} & \symking{}h4 & \symqueen{}f4\#\\
\end{tabular}\par\vspace{3mm}
\qrcode[height=12mm]{https://helpman.komtera.lt/?fen=N7/8/8/5Q1p/8/8/7Q/K5k1&moves=4}\\{\sffamily\tiny\color{ink}scannen zum Lösen}
\end{minipage}
\vspace{2mm}

{\raggedright\noindent\textcolor{ink}{\sffamily\textsc{Themen}}\enspace \themetag{reines Matt} \themetag{Mustermatt} \themetag{ein Stein zieht} \themetag{ohne Schach}\par\smallskip{\footnotesize\textit{Hinweis:} die Partei am Zug steht im Diagramm im Schach; die Lösung beginnt mit einem Schlag.}\par}

\clearpage
\phantomsection\label{hm:KQQNvkq}
\begin{tcolorbox}[enhanced,colback=ink,colframe=ink,arc=0pt,boxrule=0pt,left=4mm,right=4mm,top=2.5mm,bottom=2.5mm]
\color{white}\sffamily{\LARGE\bfseries KQQNvkq}\hfill{\Large h\#4}\\[0.5mm]
{\small 6 Steine\quad Schwarz am Zug}\hfill{\small Nr.~815\qquad\href{https://helpman.komtera.lt/?fen=N7/8/8/8/8/1q6/Q7/KQk5&moves=4}{\color{white}\bfseries lösen\,\ensuremath{\nearrow}}}
\end{tcolorbox}
\vspace{2mm}
\noindent\begin{minipage}[t]{36mm}\vspace{0pt}\sffamily\small\raggedright
\hmfigure{längste eindeutige}{{\Large h\#4}}\hmfigure{eindeutig bei dieser Länge}{{\Large 1\,019} Stellungen}\hmfigure{längstes Matt}{{\Large h\#5}}\hmfigure{in der Klasse}{15,5 Milliarden Stellungen,\\ 3,27 Milliarden mit Hilfsmatt}
\end{minipage}\hfill
\begin{minipage}[t]{80mm}\vspace{0pt}\begin{tikzpicture}\begin{axis}[width=72mm,height=42mm,ybar,bar width=1.6pt,xmin=-1,xmax=11,ymin=0,ymax=1.119,axis y line*=left,axis x line*=bottom,xlabel={\scriptsize Halbzüge bis Matt},ylabel={\scriptsize Stellungen (Mrd.)},tick label style={font=\tiny},label style={font=\scriptsize},enlarge x limits=false,ylabel shift=-3pt,xlabel shift=-3pt]\addplot[fill=ink!55,draw=none] coordinates {(0,0.0577) (1,0.1834) (2,0.9085) (3,0.9074) (4,1.017) (5,0.1414) (6,0.04826) (7,1.589e-06) (8,0.004578) (9,5.4e-08) (10,9.513e-06)};\draw[densely dotted] (axis cs:8,0) -- (axis cs:8,1.119);\end{axis}\begin{axis}[width=72mm,height=42mm,axis y line*=right,axis x line=none,ymode=log,xmin=-1,xmax=11,ylabel={\scriptsize eindeutig (log)},tick label style={font=\tiny},label style={font=\scriptsize},ylabel shift=-3pt]\addplot[mark=*,mark size=0.7pt,draw=black,thin] coordinates {(0,57700487) (1,110646216) (2,306056287) (3,35096316) (4,7025053) (5,291) (6,462184) (7,13) (8,1019)};\end{axis}\end{tikzpicture}\end{minipage}
\vspace{3mm}

\noindent\begin{minipage}[t]{60mm}\vspace{0pt}\sbox\hmboard{\chessboard[setfen=N7/8/8/8/8/1q6/Q7/KQk5 b - - 0 1]}\usebox\hmboard\\[1mm]
\makebox[\wd\hmboard]{\sffamily\hspace{3mm}\textbf{h\#4}\hfill(4+2)}\end{minipage}\hfill
\begin{minipage}[t]{52mm}\vspace{0pt}
\textcolor{ink}{\sffamily\textsc{Lösung}}\par\smallskip\small
\begin{tabular}{@{}r@{\enspace}l@{\enspace}l@{}}
\textbf{1.} & \symqueen{}xb1+ & \symqueen{}xb1+\\
\textbf{2.} & \symking{}d2 & \symknight{}b6\\
\textbf{3.} & \symking{}c3 & \symqueen{}c1+\\
\textbf{4.} & \symking{}b3 & \symqueen{}b2\#\\
\end{tabular}\par\vspace{3mm}
\qrcode[height=12mm]{https://helpman.komtera.lt/?fen=N7/8/8/8/8/1q6/Q7/KQk5&moves=4}\\{\sffamily\tiny\color{ink}scannen zum Lösen}
\end{minipage}
\vspace{2mm}

{\raggedright\noindent\textcolor{ink}{\sffamily\textsc{Themen}}\enspace \textit{keine erkannt}\par\smallskip{\footnotesize\textit{Hinweis:} die Partei am Zug steht im Diagramm im Schach; die Lösung beginnt mit einem Schlag.}\par}

\clearpage
\phantomsection\label{hm:KQQNvkr}
\begin{tcolorbox}[enhanced,colback=ink,colframe=ink,arc=0pt,boxrule=0pt,left=4mm,right=4mm,top=2.5mm,bottom=2.5mm]
\color{white}\sffamily{\LARGE\bfseries KQQNvkr}\hfill{\Large h\#4}\\[0.5mm]
{\small 6 Steine\quad Schwarz am Zug}\hfill{\small Nr.~816\qquad\href{https://helpman.komtera.lt/?fen=N7/8/8/8/8/1r6/Q7/KQk5&moves=4}{\color{white}\bfseries lösen\,\ensuremath{\nearrow}}}
\end{tcolorbox}
\vspace{2mm}
\noindent\begin{minipage}[t]{36mm}\vspace{0pt}\sffamily\small\raggedright
\hmfigure{längste eindeutige}{{\Large h\#4}}\hmfigure{eindeutig bei dieser Länge}{{\Large 29} Stellungen}\hmfigure{längstes Matt}{{\Large h\#5}}\hmfigure{in der Klasse}{15,5 Milliarden Stellungen,\\ 3,66 Milliarden mit Hilfsmatt}
\end{minipage}\hfill
\begin{minipage}[t]{80mm}\vspace{0pt}\begin{tikzpicture}\begin{axis}[width=72mm,height=42mm,ybar,bar width=1.6pt,xmin=-1,xmax=11,ymin=0,ymax=1.271,axis y line*=left,axis x line*=bottom,xlabel={\scriptsize Halbzüge bis Matt},ylabel={\scriptsize Stellungen (Mrd.)},tick label style={font=\tiny},label style={font=\scriptsize},enlarge x limits=false,ylabel shift=-3pt,xlabel shift=-3pt]\addplot[fill=ink!55,draw=none] coordinates {(0,0.08023) (1,0.2485) (2,1.156) (3,0.9015) (4,1.131) (5,0.0831) (6,0.05242) (7,5.58e-07) (8,0.003282) (9,5.1e-08) (10,1.17e-07)};\draw[densely dotted] (axis cs:8,0) -- (axis cs:8,1.271);\end{axis}\begin{axis}[width=72mm,height=42mm,axis y line*=right,axis x line=none,ymode=log,xmin=-1,xmax=11,ylabel={\scriptsize eindeutig (log)},tick label style={font=\tiny},label style={font=\scriptsize},ylabel shift=-3pt]\addplot[mark=*,mark size=0.7pt,draw=black,thin] coordinates {(0,80226189) (1,139486975) (2,367102620) (3,20124174) (4,7716419) (5,63) (6,477799) (7,2) (8,29)};\end{axis}\end{tikzpicture}\end{minipage}
\vspace{3mm}

\noindent\begin{minipage}[t]{60mm}\vspace{0pt}\sbox\hmboard{\chessboard[setfen=N7/8/8/8/8/1r6/Q7/KQk5 b - - 0 1]}\usebox\hmboard\\[1mm]
\makebox[\wd\hmboard]{\sffamily\hspace{3mm}\textbf{h\#4}\hfill(4+2)}\end{minipage}\hfill
\begin{minipage}[t]{52mm}\vspace{0pt}
\textcolor{ink}{\sffamily\textsc{Lösung}}\par\smallskip\small
\begin{tabular}{@{}r@{\enspace}l@{\enspace}l@{}}
\textbf{1.} & \symrook{}xb1+ & \symqueen{}xb1+\\
\textbf{2.} & \symking{}d2 & \symknight{}b6\\
\textbf{3.} & \symking{}c3 & \symqueen{}c1+\\
\textbf{4.} & \symking{}b3 & \symqueen{}b2\#\\
\end{tabular}\par\vspace{3mm}
\qrcode[height=12mm]{https://helpman.komtera.lt/?fen=N7/8/8/8/8/1r6/Q7/KQk5&moves=4}\\{\sffamily\tiny\color{ink}scannen zum Lösen}
\end{minipage}
\vspace{2mm}

{\raggedright\noindent\textcolor{ink}{\sffamily\textsc{Themen}}\enspace \textit{keine erkannt}\par\smallskip{\footnotesize\textit{Hinweis:} die Partei am Zug steht im Diagramm im Schach; die Lösung beginnt mit einem Schlag.}\par}

\clearpage
\phantomsection\label{hm:KQQQvkn}
\begin{tcolorbox}[enhanced,colback=ink,colframe=ink,arc=0pt,boxrule=0pt,left=4mm,right=4mm,top=2.5mm,bottom=2.5mm]
\color{white}\sffamily{\LARGE\bfseries KQQQvkn}\hfill{\Large h\#4}\\[0.5mm]
{\small 6 Steine\quad Schwarz am Zug}\hfill{\small Nr.~817\qquad\href{https://helpman.komtera.lt/?fen=8/8/8/1Q6/8/8/Q1Q5/K1k1n3&moves=4}{\color{white}\bfseries lösen\,\ensuremath{\nearrow}}}
\end{tcolorbox}
\vspace{2mm}
\noindent\begin{minipage}[t]{36mm}\vspace{0pt}\sffamily\small\raggedright
\hmfigure{längste eindeutige}{{\Large h\#4}}\hmfigure{eindeutig bei dieser Länge}{{\Large 45} Stellungen}\hmfigure{längstes Matt}{{\Large h\#5}}\hmfigure{in der Klasse}{15,5 Milliarden Stellungen,\\ 1,2 Milliarden mit Hilfsmatt}
\end{minipage}\hfill
\begin{minipage}[t]{80mm}\vspace{0pt}\begin{tikzpicture}\begin{axis}[width=72mm,height=42mm,ybar,bar width=1.6pt,xmin=-1,xmax=11,ymin=0,ymax=638.5,axis y line*=left,axis x line*=bottom,xlabel={\scriptsize Halbzüge bis Matt},ylabel={\scriptsize Stellungen (Mio.)},tick label style={font=\tiny},label style={font=\scriptsize},enlarge x limits=false,ylabel shift=-3pt,xlabel shift=-3pt]\addplot[fill=ink!55,draw=none] coordinates {(0,100.2) (1,130.9) (2,580.4) (3,146.3) (4,236.4) (5,3.441) (6,0.2528) (8,0.000424) (10,1.3e-05)};\draw[densely dotted] (axis cs:8,0) -- (axis cs:8,638.5);\end{axis}\begin{axis}[width=72mm,height=42mm,axis y line*=right,axis x line=none,ymode=log,xmin=-1,xmax=11,ylabel={\scriptsize eindeutig (log)},tick label style={font=\tiny},label style={font=\scriptsize},ylabel shift=-3pt]\addplot[mark=*,mark size=0.7pt,draw=black,thin] coordinates {(0,100179096) (1,32945612) (2,75274603) (3,714976) (4,213190) (6,10) (8,45)};\end{axis}\end{tikzpicture}\end{minipage}
\vspace{3mm}

\noindent\begin{minipage}[t]{60mm}\vspace{0pt}\sbox\hmboard{\chessboard[setfen=8/8/8/1Q6/8/8/Q1Q5/K1k1n3 b - - 0 1]}\usebox\hmboard\\[1mm]
\makebox[\wd\hmboard]{\sffamily\hspace{3mm}\textbf{h\#4}\hfill(4+2)}\end{minipage}\hfill
\begin{minipage}[t]{52mm}\vspace{0pt}
\textcolor{ink}{\sffamily\textsc{Lösung}}\par\smallskip\small
\begin{tabular}{@{}r@{\enspace}l@{\enspace}l@{}}
\textbf{1.} & \symknight{}xc2+ & \symqueen{}xc2+\\
\textbf{2.} & \symking{}xc2 & \symking{}a2\\
\textbf{3.} & \symking{}c1 & \symking{}b3\\
\textbf{4.} & \symking{}b1 & \symqueen{}f1\#\\
\end{tabular}\par\vspace{3mm}
\qrcode[height=12mm]{https://helpman.komtera.lt/?fen=8/8/8/1Q6/8/8/Q1Q5/K1k1n3&moves=4}\\{\sffamily\tiny\color{ink}scannen zum Lösen}
\end{minipage}
\vspace{2mm}

{\raggedright\noindent\textcolor{ink}{\sffamily\textsc{Themen}}\enspace \themetag{reines Matt} \themetag{Mustermatt} \themetag{Idealmatt} \themetag{Spiegelmatt} \themetag{Switchback}\par\smallskip{\footnotesize\textit{Hinweis:} die Partei am Zug steht im Diagramm im Schach; die Lösung beginnt mit einem Schlag.}\par}

\clearpage
\phantomsection\label{hm:KQQRvkb}
\begin{tcolorbox}[enhanced,colback=ink,colframe=ink,arc=0pt,boxrule=0pt,left=4mm,right=4mm,top=2.5mm,bottom=2.5mm]
\color{white}\sffamily{\LARGE\bfseries KQQRvkb}\hfill{\Large h\#4}\\[0.5mm]
{\small 6 Steine\quad Schwarz am Zug}\hfill{\small Nr.~818\qquad\href{https://helpman.komtera.lt/?fen=6bk/7Q/8/8/8/3K4/6R1/7Q&moves=4}{\color{white}\bfseries lösen\,\ensuremath{\nearrow}}}
\end{tcolorbox}
\vspace{2mm}
\noindent\begin{minipage}[t]{36mm}\vspace{0pt}\sffamily\small\raggedright
\hmfigure{längste eindeutige}{{\Large h\#4}}\hmfigure{eindeutig bei dieser Länge}{{\Large 4} Stellungen}\hmfigure{längstes Matt}{{\Large h\#6}}\hmfigure{in der Klasse}{15,5 Milliarden Stellungen,\\ 3,66 Milliarden mit Hilfsmatt}
\end{minipage}\hfill
\begin{minipage}[t]{80mm}\vspace{0pt}\begin{tikzpicture}\begin{axis}[width=72mm,height=42mm,ybar,bar width=1.6pt,xmin=-1,xmax=13,ymin=0,ymax=1.567,axis y line*=left,axis x line*=bottom,xlabel={\scriptsize Halbzüge bis Matt},ylabel={\scriptsize Stellungen (Mrd.)},tick label style={font=\tiny},label style={font=\scriptsize},enlarge x limits=false,ylabel shift=-3pt,xlabel shift=-3pt]\addplot[fill=ink!55,draw=none] coordinates {(0,0.1686) (1,0.3251) (2,1.425) (3,0.6718) (4,1.017) (5,0.03112) (6,0.01714) (7,3.95e-07) (8,2.957e-06) (10,6e-09) (12,1e-09)};\draw[densely dotted] (axis cs:8,0) -- (axis cs:8,1.567);\end{axis}\begin{axis}[width=72mm,height=42mm,axis y line*=right,axis x line=none,ymode=log,xmin=-1,xmax=13,ylabel={\scriptsize eindeutig (log)},tick label style={font=\tiny},label style={font=\scriptsize},ylabel shift=-3pt]\addplot[mark=*,mark size=0.7pt,draw=black,thin] coordinates {(0,168597397) (1,125500091) (2,297575737) (3,6515849) (4,8106041) (5,73) (6,1322) (8,4)};\end{axis}\end{tikzpicture}\end{minipage}
\vspace{3mm}

\noindent\begin{minipage}[t]{60mm}\vspace{0pt}\sbox\hmboard{\chessboard[setfen=6bk/7Q/8/8/8/3K4/6R1/7Q b - - 0 1]}\usebox\hmboard\\[1mm]
\makebox[\wd\hmboard]{\sffamily\hspace{3mm}\textbf{h\#4}\hfill(4+2)}\end{minipage}\hfill
\begin{minipage}[t]{52mm}\vspace{0pt}
\textcolor{ink}{\sffamily\textsc{Lösung}}\par\smallskip\small
\begin{tabular}{@{}r@{\enspace}l@{\enspace}l@{}}
\textbf{1.} & \symbishop{}xh7+ & \symqueen{}xh7+\\
\textbf{2.} & \symking{}xh7 & \symking{}e4\\
\textbf{3.} & \symking{}h6 & \symking{}f5\\
\textbf{4.} & \symking{}h5 & \symrook{}h2\#\\
\end{tabular}\par\vspace{3mm}
\qrcode[height=12mm]{https://helpman.komtera.lt/?fen=6bk/7Q/8/8/8/3K4/6R1/7Q&moves=4}\\{\sffamily\tiny\color{ink}scannen zum Lösen}
\end{minipage}
\vspace{2mm}

{\raggedright\noindent\textcolor{ink}{\sffamily\textsc{Themen}}\enspace \themetag{reines Matt} \themetag{Mustermatt} \themetag{Idealmatt} \themetag{Spiegelmatt}\par\smallskip{\footnotesize\textit{Hinweis:} die Partei am Zug steht im Diagramm im Schach; die Lösung beginnt mit einem Schlag.}\par}

\clearpage
\phantomsection\label{hm:KQQRvkn}
\begin{tcolorbox}[enhanced,colback=ink,colframe=ink,arc=0pt,boxrule=0pt,left=4mm,right=4mm,top=2.5mm,bottom=2.5mm]
\color{white}\sffamily{\LARGE\bfseries KQQRvkn}\hfill{\Large h\#4}\\[0.5mm]
{\small 6 Steine\quad Schwarz am Zug}\hfill{\small Nr.~819\qquad\href{https://helpman.komtera.lt/?fen=8/8/8/8/8/3R4/Q1Q5/K1k1n3&moves=4}{\color{white}\bfseries lösen\,\ensuremath{\nearrow}}}
\end{tcolorbox}
\vspace{2mm}
\noindent\begin{minipage}[t]{36mm}\vspace{0pt}\sffamily\small\raggedright
\hmfigure{längste eindeutige}{{\Large h\#4}}\hmfigure{eindeutig bei dieser Länge}{{\Large 322} Stellungen}\hmfigure{längstes Matt}{{\Large h\#5}}\hmfigure{in der Klasse}{15,5 Milliarden Stellungen,\\ 3,79 Milliarden mit Hilfsmatt}
\end{minipage}\hfill
\begin{minipage}[t]{80mm}\vspace{0pt}\begin{tikzpicture}\begin{axis}[width=72mm,height=42mm,ybar,bar width=1.6pt,xmin=-1,xmax=11,ymin=0,ymax=1.61,axis y line*=left,axis x line*=bottom,xlabel={\scriptsize Halbzüge bis Matt},ylabel={\scriptsize Stellungen (Mrd.)},tick label style={font=\tiny},label style={font=\scriptsize},enlarge x limits=false,ylabel shift=-3pt,xlabel shift=-3pt]\addplot[fill=ink!55,draw=none] coordinates {(0,0.1897) (1,0.3465) (2,1.463) (3,0.655) (4,1.092) (5,0.02655) (6,0.01945) (7,2.01e-07) (8,1.755e-05) (10,8.2e-08)};\draw[densely dotted] (axis cs:8,0) -- (axis cs:8,1.61);\end{axis}\begin{axis}[width=72mm,height=42mm,axis y line*=right,axis x line=none,ymode=log,xmin=-1,xmax=11,ylabel={\scriptsize eindeutig (log)},tick label style={font=\tiny},label style={font=\scriptsize},ylabel shift=-3pt]\addplot[mark=*,mark size=0.7pt,draw=black,thin] coordinates {(0,189685242) (1,122784854) (2,286073574) (3,4593419) (4,8663843) (5,229) (6,6009) (8,322)};\end{axis}\end{tikzpicture}\end{minipage}
\vspace{3mm}

\noindent\begin{minipage}[t]{60mm}\vspace{0pt}\sbox\hmboard{\chessboard[setfen=8/8/8/8/8/3R4/Q1Q5/K1k1n3 b - - 0 1]}\usebox\hmboard\\[1mm]
\makebox[\wd\hmboard]{\sffamily\hspace{3mm}\textbf{h\#4}\hfill(4+2)}\end{minipage}\hfill
\begin{minipage}[t]{52mm}\vspace{0pt}
\textcolor{ink}{\sffamily\textsc{Lösung}}\par\smallskip\small
\begin{tabular}{@{}r@{\enspace}l@{\enspace}l@{}}
\textbf{1.} & \symknight{}xc2+ & \symqueen{}xc2+\\
\textbf{2.} & \symking{}xc2 & \symking{}a2\\
\textbf{3.} & \symking{}c1 & \symking{}b3\\
\textbf{4.} & \symking{}b1 & \symrook{}d1\#\\
\end{tabular}\par\vspace{3mm}
\qrcode[height=12mm]{https://helpman.komtera.lt/?fen=8/8/8/8/8/3R4/Q1Q5/K1k1n3&moves=4}\\{\sffamily\tiny\color{ink}scannen zum Lösen}
\end{minipage}
\vspace{2mm}

{\raggedright\noindent\textcolor{ink}{\sffamily\textsc{Themen}}\enspace \themetag{reines Matt} \themetag{Mustermatt} \themetag{Idealmatt} \themetag{Spiegelmatt} \themetag{Switchback}\par\smallskip{\footnotesize\textit{Hinweis:} die Partei am Zug steht im Diagramm im Schach; die Lösung beginnt mit einem Schlag.}\par}

\clearpage
\phantomsection\label{hm:KQQRvkq}
\begin{tcolorbox}[enhanced,colback=ink,colframe=ink,arc=0pt,boxrule=0pt,left=4mm,right=4mm,top=2.5mm,bottom=2.5mm]
\color{white}\sffamily{\LARGE\bfseries KQQRvkq}\hfill{\Large h\#4}\\[0.5mm]
{\small 6 Steine\quad Schwarz am Zug}\hfill{\small Nr.~820\qquad\href{https://helpman.komtera.lt/?fen=8/8/8/8/q2k4/5Q2/2R5/QK6&moves=4}{\color{white}\bfseries lösen\,\ensuremath{\nearrow}}}
\end{tcolorbox}
\vspace{2mm}
\noindent\begin{minipage}[t]{36mm}\vspace{0pt}\sffamily\small\raggedright
\hmfigure{längste eindeutige}{{\Large h\#4}}\hmfigure{eindeutig bei dieser Länge}{{\Large 4} Stellungen}\hmfigure{längstes Matt}{{\Large h\#7}}\hmfigure{in der Klasse}{15,5 Milliarden Stellungen,\\ 3,06 Milliarden mit Hilfsmatt}
\end{minipage}\hfill
\begin{minipage}[t]{80mm}\vspace{0pt}\begin{tikzpicture}\begin{axis}[width=72mm,height=42mm,ybar,bar width=1.6pt,xmin=-1,xmax=15,ymin=0,ymax=1.177,axis y line*=left,axis x line*=bottom,xlabel={\scriptsize Halbzüge bis Matt},ylabel={\scriptsize Stellungen (Mrd.)},tick label style={font=\tiny},label style={font=\scriptsize},enlarge x limits=false,ylabel shift=-3pt,xlabel shift=-3pt]\addplot[fill=ink!55,draw=none] coordinates {(0,0.1062) (1,0.2216) (2,1.07) (3,0.7215) (4,0.8414) (5,0.08395) (6,0.015) (7,2.08e-07) (8,5.142e-06) (9,8.9e-08) (10,4.29e-07) (11,1e-08) (12,1.36e-07) (14,3e-09)};\draw[densely dotted] (axis cs:8,0) -- (axis cs:8,1.177);\end{axis}\begin{axis}[width=72mm,height=42mm,axis y line*=right,axis x line=none,ymode=log,xmin=-1,xmax=15,ylabel={\scriptsize eindeutig (log)},tick label style={font=\tiny},label style={font=\scriptsize},ylabel shift=-3pt]\addplot[mark=*,mark size=0.7pt,draw=black,thin] coordinates {(0,106226258) (1,104704230) (2,265520128) (3,17615231) (4,7872392) (5,308) (6,3336) (7,26) (8,4)};\end{axis}\end{tikzpicture}\end{minipage}
\vspace{3mm}

\noindent\begin{minipage}[t]{60mm}\vspace{0pt}\sbox\hmboard{\chessboard[setfen=8/8/8/8/q2k4/5Q2/2R5/QK6 b - - 0 1]}\usebox\hmboard\\[1mm]
\makebox[\wd\hmboard]{\sffamily\hspace{3mm}\textbf{h\#4}\hfill(4+2)}\end{minipage}\hfill
\begin{minipage}[t]{52mm}\vspace{0pt}
\textcolor{ink}{\sffamily\textsc{Lösung}}\par\smallskip\small
\begin{tabular}{@{}r@{\enspace}l@{\enspace}l@{}}
\textbf{1.} & \symqueen{}xa1+ & \symking{}xa1\\
\textbf{2.} & \symking{}e5 & \symking{}b2\\
\textbf{3.} & \symking{}d4 & \symqueen{}f4+\\
\textbf{4.} & \symking{}d3 & \symrook{}d2\#\\
\end{tabular}\par\vspace{3mm}
\qrcode[height=12mm]{https://helpman.komtera.lt/?fen=8/8/8/8/q2k4/5Q2/2R5/QK6&moves=4}\\{\sffamily\tiny\color{ink}scannen zum Lösen}
\end{minipage}
\vspace{2mm}

{\raggedright\noindent\textcolor{ink}{\sffamily\textsc{Themen}}\enspace \themetag{Switchback}\par\smallskip{\footnotesize\textit{Hinweis:} die Partei am Zug steht im Diagramm im Schach; die Lösung beginnt mit einem Schlag.}\par}

\clearpage
\phantomsection\label{hm:KQRBvkr}
\begin{tcolorbox}[enhanced,colback=ink,colframe=ink,arc=0pt,boxrule=0pt,left=4mm,right=4mm,top=2.5mm,bottom=2.5mm]
\color{white}\sffamily{\LARGE\bfseries KQRBvkr}\hfill{\Large h\#4}\\[0.5mm]
{\small 6 Steine\quad Schwarz am Zug}\hfill{\small Nr.~821\qquad\href{https://helpman.komtera.lt/?fen=8/8/8/8/8/8/Br2R3/KQk5&moves=4}{\color{white}\bfseries lösen\,\ensuremath{\nearrow}}}
\end{tcolorbox}
\vspace{2mm}
\noindent\begin{minipage}[t]{36mm}\vspace{0pt}\sffamily\small\raggedright
\hmfigure{längste eindeutige}{{\Large h\#4}}\hmfigure{eindeutig bei dieser Länge}{{\Large 47\,252} Stellungen}\hmfigure{längstes Matt}{{\Large h\#6}}\hmfigure{in der Klasse}{15,5 Milliarden Stellungen,\\ 7,64 Milliarden mit Hilfsmatt}
\end{minipage}\hfill
\begin{minipage}[t]{80mm}\vspace{0pt}\begin{tikzpicture}\begin{axis}[width=72mm,height=42mm,ybar,bar width=1.6pt,xmin=-1,xmax=13,ymin=0,ymax=3.555,axis y line*=left,axis x line*=bottom,xlabel={\scriptsize Halbzüge bis Matt},ylabel={\scriptsize Stellungen (Mrd.)},tick label style={font=\tiny},label style={font=\scriptsize},enlarge x limits=false,ylabel shift=-3pt,xlabel shift=-3pt]\addplot[fill=ink!55,draw=none] coordinates {(0,0.08443) (1,0.3118) (2,1.423) (3,2.066) (4,3.232) (5,0.4141) (6,0.07945) (7,2.96e-05) (8,0.02861) (9,3.9e-08) (10,0.0004222) (11,6e-09) (12,1.2e-08)};\draw[densely dotted] (axis cs:8,0) -- (axis cs:8,3.555);\end{axis}\begin{axis}[width=72mm,height=42mm,axis y line*=right,axis x line=none,ymode=log,xmin=-1,xmax=13,ylabel={\scriptsize eindeutig (log)},tick label style={font=\tiny},label style={font=\scriptsize},ylabel shift=-3pt]\addplot[mark=*,mark size=0.7pt,draw=black,thin] coordinates {(0,84425455) (1,217791588) (2,603771136) (3,119268052) (4,18000904) (5,8793) (6,4321614) (7,18) (8,47252)};\end{axis}\end{tikzpicture}\end{minipage}
\vspace{3mm}

\noindent\begin{minipage}[t]{60mm}\vspace{0pt}\sbox\hmboard{\chessboard[setfen=8/8/8/8/8/8/Br2R3/KQk5 b - - 0 1]}\usebox\hmboard\\[1mm]
\makebox[\wd\hmboard]{\sffamily\hspace{3mm}\textbf{h\#4}\hfill(4+2)}\end{minipage}\hfill
\begin{minipage}[t]{52mm}\vspace{0pt}
\textcolor{ink}{\sffamily\textsc{Lösung}}\par\smallskip\small
\begin{tabular}{@{}r@{\enspace}l@{\enspace}l@{}}
\textbf{1.} & \symrook{}xb1+ & \symbishop{}xb1\\
\textbf{2.} & \symking{}d1 & \symking{}a2\\
\textbf{3.} & \symking{}c1 & \symking{}b3\\
\textbf{4.} & \symking{}xb1 & \symrook{}e1\#\\
\end{tabular}\par\vspace{3mm}
\qrcode[height=12mm]{https://helpman.komtera.lt/?fen=8/8/8/8/8/8/Br2R3/KQk5&moves=4}\\{\sffamily\tiny\color{ink}scannen zum Lösen}
\end{minipage}
\vspace{2mm}

{\raggedright\noindent\textcolor{ink}{\sffamily\textsc{Themen}}\enspace \themetag{reines Matt} \themetag{Mustermatt} \themetag{Idealmatt} \themetag{Spiegelmatt} \themetag{Switchback} \themetag{Kniest} \themetag{Zajic}\par\smallskip{\footnotesize\textit{Hinweis:} die Partei am Zug steht im Diagramm im Schach; die Lösung beginnt mit einem Schlag.}\par}

\clearpage
\phantomsection\label{hm:KQRRBvk}
\begin{tcolorbox}[enhanced,colback=ink,colframe=ink,arc=0pt,boxrule=0pt,left=4mm,right=4mm,top=2.5mm,bottom=2.5mm]
\color{white}\sffamily{\LARGE\bfseries KQRRBvk}\hfill{\Large h\#4}\\[0.5mm]
{\small 6 Steine\quad Schwarz am Zug}\hfill{\small Nr.~822\qquad\href{https://helpman.komtera.lt/?fen=8/8/8/8/8/6R1/6BQ/1K4Rk&moves=4}{\color{white}\bfseries lösen\,\ensuremath{\nearrow}}}
\end{tcolorbox}
\vspace{2mm}
\noindent\begin{minipage}[t]{36mm}\vspace{0pt}\sffamily\small\raggedright
\hmfigure{längste eindeutige}{{\Large h\#4}}\hmfigure{eindeutig bei dieser Länge}{{\Large 6} Stellungen}\hmfigure{längstes Matt}{{\Large h\#6}}\hmfigure{in der Klasse}{15,5 Milliarden Stellungen,\\ 4,02 Milliarden mit Hilfsmatt}
\end{minipage}\hfill
\begin{minipage}[t]{80mm}\vspace{0pt}\begin{tikzpicture}\begin{axis}[width=72mm,height=42mm,ybar,bar width=1.6pt,xmin=-1,xmax=13,ymin=0,ymax=1.602,axis y line*=left,axis x line*=bottom,xlabel={\scriptsize Halbzüge bis Matt},ylabel={\scriptsize Stellungen (Mrd.)},tick label style={font=\tiny},label style={font=\scriptsize},enlarge x limits=false,ylabel shift=-3pt,xlabel shift=-3pt]\addplot[fill=ink!55,draw=none] coordinates {(0,0.192) (1,0.3612) (2,1.271) (3,0.7111) (4,1.456) (5,0.007793) (6,0.02296) (8,0.0004278) (10,9.5e-08) (12,8e-09)};\draw[densely dotted] (axis cs:8,0) -- (axis cs:8,1.602);\end{axis}\begin{axis}[width=72mm,height=42mm,axis y line*=right,axis x line=none,ymode=log,xmin=-1,xmax=13,ylabel={\scriptsize eindeutig (log)},tick label style={font=\tiny},label style={font=\scriptsize},ylabel shift=-3pt]\addplot[mark=*,mark size=0.7pt,draw=black,thin] coordinates {(0,191976016) (1,146219790) (2,378398885) (3,7516418) (4,9736423) (6,103937) (8,6)};\end{axis}\end{tikzpicture}\end{minipage}
\vspace{3mm}

\noindent\begin{minipage}[t]{60mm}\vspace{0pt}\sbox\hmboard{\chessboard[setfen=8/8/8/8/8/6R1/6BQ/1K4Rk b - - 0 1]}\usebox\hmboard\\[1mm]
\makebox[\wd\hmboard]{\sffamily\hspace{3mm}\textbf{h\#4}\hfill(5+1)}\end{minipage}\hfill
\begin{minipage}[t]{52mm}\vspace{0pt}
\textcolor{ink}{\sffamily\textsc{Lösung}}\par\smallskip\small
\begin{tabular}{@{}r@{\enspace}l@{\enspace}l@{}}
\textbf{1.} & \symking{}xh2 & \symking{}c2\\
\textbf{2.} & \symking{}xg1 & \symking{}d3\\
\textbf{3.} & \symking{}f2 & \symrook{}f3+\\
\textbf{4.} & \symking{}e1 & \symrook{}f1\#\\
\end{tabular}\par\vspace{3mm}
\qrcode[height=12mm]{https://helpman.komtera.lt/?fen=8/8/8/8/8/6R1/6BQ/1K4Rk&moves=4}\\{\sffamily\tiny\color{ink}scannen zum Lösen}
\end{minipage}
\vspace{2mm}

{\raggedright\noindent\textcolor{ink}{\sffamily\textsc{Themen}}\enspace \themetag{reines Matt} \themetag{Mustermatt} \themetag{Idealmatt} \themetag{ein Stein zieht}\par\smallskip{\footnotesize\textit{Hinweis:} das Diagramm hat keinen legalen letzten Zug, kann also in keiner Partie entstehen.}\par}

\clearpage
\phantomsection\label{hm:KQRRNvk}
\begin{tcolorbox}[enhanced,colback=ink,colframe=ink,arc=0pt,boxrule=0pt,left=4mm,right=4mm,top=2.5mm,bottom=2.5mm]
\color{white}\sffamily{\LARGE\bfseries KQRRNvk}\hfill{\Large h\#4}\\[0.5mm]
{\small 6 Steine\quad Schwarz am Zug}\hfill{\small Nr.~823\qquad\href{https://helpman.komtera.lt/?fen=8/8/8/6N1/3Q4/2Rk4/2R5/K7&moves=4}{\color{white}\bfseries lösen\,\ensuremath{\nearrow}}}
\end{tcolorbox}
\vspace{2mm}
\noindent\begin{minipage}[t]{36mm}\vspace{0pt}\sffamily\small\raggedright
\hmfigure{längste eindeutige}{{\Large h\#4}}\hmfigure{eindeutig bei dieser Länge}{{\Large 3} Stellungen}\hmfigure{längstes Matt}{{\Large h\#5}}\hmfigure{in der Klasse}{15,5 Milliarden Stellungen,\\ 4,09 Milliarden mit Hilfsmatt}
\end{minipage}\hfill
\begin{minipage}[t]{80mm}\vspace{0pt}\begin{tikzpicture}\begin{axis}[width=72mm,height=42mm,ybar,bar width=1.6pt,xmin=-1,xmax=11,ymin=0,ymax=1.722,axis y line*=left,axis x line*=bottom,xlabel={\scriptsize Halbzüge bis Matt},ylabel={\scriptsize Stellungen (Mrd.)},tick label style={font=\tiny},label style={font=\scriptsize},enlarge x limits=false,ylabel shift=-3pt,xlabel shift=-3pt]\addplot[fill=ink!55,draw=none] coordinates {(0,0.169) (1,0.3521) (2,1.191) (3,0.7789) (4,1.565) (5,0.01179) (6,0.02416) (8,0.0004327) (10,1.7e-08)};\draw[densely dotted] (axis cs:8,0) -- (axis cs:8,1.722);\end{axis}\begin{axis}[width=72mm,height=42mm,axis y line*=right,axis x line=none,ymode=log,xmin=-1,xmax=11,ylabel={\scriptsize eindeutig (log)},tick label style={font=\tiny},label style={font=\scriptsize},ylabel shift=-3pt]\addplot[mark=*,mark size=0.7pt,draw=black,thin] coordinates {(0,169032503) (1,148238827) (2,355420068) (3,9569870) (4,10431716) (5,2) (6,73063) (8,3)};\end{axis}\end{tikzpicture}\end{minipage}
\vspace{3mm}

\noindent\begin{minipage}[t]{60mm}\vspace{0pt}\sbox\hmboard{\chessboard[setfen=8/8/8/6N1/3Q4/2Rk4/2R5/K7 b - - 0 1]}\usebox\hmboard\\[1mm]
\makebox[\wd\hmboard]{\sffamily\hspace{3mm}\textbf{h\#4}\hfill(5+1)}\end{minipage}\hfill
\begin{minipage}[t]{52mm}\vspace{0pt}
\textcolor{ink}{\sffamily\textsc{Lösung}}\par\smallskip\small
\begin{tabular}{@{}r@{\enspace}l@{\enspace}l@{}}
\textbf{1.} & \symking{}xd4 & \symking{}b2\\
\textbf{2.} & \symking{}e5 & \symrook{}e2+\\
\textbf{3.} & \symking{}d5 & \symrook{}e5+\\
\textbf{4.} & \symking{}d4 & \symknight{}f3\#\\
\end{tabular}\par\vspace{3mm}
\qrcode[height=12mm]{https://helpman.komtera.lt/?fen=8/8/8/6N1/3Q4/2Rk4/2R5/K7&moves=4}\\{\sffamily\tiny\color{ink}scannen zum Lösen}
\end{minipage}
\vspace{2mm}

{\raggedright\noindent\textcolor{ink}{\sffamily\textsc{Themen}}\enspace \themetag{Rundlauf} \themetag{ein Stein zieht} \themetag{Kniest} \themetag{Umnov}\par\smallskip{\footnotesize\textit{Hinweis:} das Diagramm hat keinen legalen letzten Zug, kann also in keiner Partie entstehen.}\par}

\clearpage
\phantomsection\label{hm:KQRRPvk}
\begin{tcolorbox}[enhanced,colback=ink,colframe=ink,arc=0pt,boxrule=0pt,left=4mm,right=4mm,top=2.5mm,bottom=2.5mm]
\color{white}\sffamily{\LARGE\bfseries KQRRPvk}\hfill{\Large h\#4}\\[0.5mm]
{\small 6 Steine\quad Schwarz am Zug}\hfill{\small Nr.~824\qquad\href{https://helpman.komtera.lt/?fen=2k5/2Q5/8/8/8/P7/R7/RK6&moves=4}{\color{white}\bfseries lösen\,\ensuremath{\nearrow}}}
\end{tcolorbox}
\vspace{2mm}
\noindent\begin{minipage}[t]{36mm}\vspace{0pt}\sffamily\small\raggedright
\hmfigure{längste eindeutige}{{\Large h\#4}}\hmfigure{eindeutig bei dieser Länge}{{\Large 3} Stellungen}\hmfigure{längstes Matt}{{\Large h\#7}}\hmfigure{in der Klasse}{45,4 Milliarden Stellungen,\\ 12,5 Milliarden mit Hilfsmatt}
\end{minipage}\hfill
\begin{minipage}[t]{80mm}\vspace{0pt}\begin{tikzpicture}\begin{axis}[width=72mm,height=42mm,ybar,bar width=1.6pt,xmin=-1,xmax=15,ymin=0,ymax=5.627,axis y line*=left,axis x line*=bottom,xlabel={\scriptsize Halbzüge bis Matt},ylabel={\scriptsize Stellungen (Mrd.)},tick label style={font=\tiny},label style={font=\scriptsize},enlarge x limits=false,ylabel shift=-3pt,xlabel shift=-3pt]\addplot[fill=ink!55,draw=none] coordinates {(0,0.451) (1,0.9911) (2,3.236) (3,2.574) (4,5.116) (5,0.06926) (6,0.08892) (7,5.55e-07) (8,0.001327) (10,2.54e-07) (12,1.3e-08) (14,2e-09)};\draw[densely dotted] (axis cs:8,0) -- (axis cs:8,5.627);\end{axis}\begin{axis}[width=72mm,height=42mm,axis y line*=right,axis x line=none,ymode=log,xmin=-1,xmax=15,ylabel={\scriptsize eindeutig (log)},tick label style={font=\tiny},label style={font=\scriptsize},ylabel shift=-3pt]\addplot[mark=*,mark size=0.7pt,draw=black,thin] coordinates {(0,451001682) (1,424352711) (2,966141863) (3,50265546) (4,36180989) (5,229) (6,161018) (8,3)};\end{axis}\end{tikzpicture}\end{minipage}
\vspace{3mm}

\noindent\begin{minipage}[t]{60mm}\vspace{0pt}\sbox\hmboard{\chessboard[setfen=2k5/2Q5/8/8/8/P7/R7/RK6 b - - 0 1]}\usebox\hmboard\\[1mm]
\makebox[\wd\hmboard]{\sffamily\hspace{3mm}\textbf{h\#4}\hfill(5+1)}\end{minipage}\hfill
\begin{minipage}[t]{52mm}\vspace{0pt}
\textcolor{ink}{\sffamily\textsc{Lösung}}\par\smallskip\small
\begin{tabular}{@{}r@{\enspace}l@{\enspace}l@{}}
\textbf{1.} & \symking{}xc7 & a4\\
\textbf{2.} & \symking{}b6 & \symrook{}b2+\\
\textbf{3.} & \symking{}a5 & \symking{}a2\\
\textbf{4.} & \symking{}xa4 & \symking{}b1\#\\
\end{tabular}\par\vspace{3mm}
\qrcode[height=12mm]{https://helpman.komtera.lt/?fen=2k5/2Q5/8/8/8/P7/R7/RK6&moves=4}\\{\sffamily\tiny\color{ink}scannen zum Lösen}
\end{minipage}
\vspace{2mm}

{\raggedright\noindent\textcolor{ink}{\sffamily\textsc{Themen}}\enspace \themetag{reines Matt} \themetag{Mustermatt} \themetag{Spiegelmatt} \themetag{Switchback} \themetag{ein Stein zieht} \themetag{Kniest}\par\smallskip{\footnotesize\textit{Hinweis:} die Partei am Zug steht im Diagramm im Schach; die Lösung beginnt mit einem Schlag.}\par}

\clearpage
\phantomsection\label{hm:KQRRvkb}
\begin{tcolorbox}[enhanced,colback=ink,colframe=ink,arc=0pt,boxrule=0pt,left=4mm,right=4mm,top=2.5mm,bottom=2.5mm]
\color{white}\sffamily{\LARGE\bfseries KQRRvkb}\hfill{\Large h\#4}\\[0.5mm]
{\small 6 Steine\quad Schwarz am Zug}\hfill{\small Nr.~825\qquad\href{https://helpman.komtera.lt/?fen=7b/8/8/4k1R1/5Q2/4R3/8/K7&moves=4}{\color{white}\bfseries lösen\,\ensuremath{\nearrow}}}
\end{tcolorbox}
\vspace{2mm}
\noindent\begin{minipage}[t]{36mm}\vspace{0pt}\sffamily\small\raggedright
\hmfigure{längste eindeutige}{{\Large h\#4}}\hmfigure{eindeutig bei dieser Länge}{{\Large 33} Stellungen}\hmfigure{längstes Matt}{{\Large h\#7}}\hmfigure{in der Klasse}{15,5 Milliarden Stellungen,\\ 3,87 Milliarden mit Hilfsmatt}
\end{minipage}\hfill
\begin{minipage}[t]{80mm}\vspace{0pt}\begin{tikzpicture}\begin{axis}[width=72mm,height=42mm,ybar,bar width=1.6pt,xmin=-1,xmax=15,ymin=0,ymax=1.706,axis y line*=left,axis x line*=bottom,xlabel={\scriptsize Halbzüge bis Matt},ylabel={\scriptsize Stellungen (Mrd.)},tick label style={font=\tiny},label style={font=\scriptsize},enlarge x limits=false,ylabel shift=-3pt,xlabel shift=-3pt]\addplot[fill=ink!55,draw=none] coordinates {(0,0.09715) (1,0.237) (2,0.9573) (3,0.9155) (4,1.551) (5,0.08501) (6,0.02595) (7,8.83e-07) (8,0.0005854) (10,7.3e-08) (12,3e-09) (14,1e-09)};\draw[densely dotted] (axis cs:8,0) -- (axis cs:8,1.706);\end{axis}\begin{axis}[width=72mm,height=42mm,axis y line*=right,axis x line=none,ymode=log,xmin=-1,xmax=15,ylabel={\scriptsize eindeutig (log)},tick label style={font=\tiny},label style={font=\scriptsize},ylabel shift=-3pt]\addplot[mark=*,mark size=0.7pt,draw=black,thin] coordinates {(0,97150374) (1,125194636) (2,266014156) (3,17112628) (4,8898599) (5,204) (6,69166) (7,2) (8,33)};\end{axis}\end{tikzpicture}\end{minipage}
\vspace{3mm}

\noindent\begin{minipage}[t]{60mm}\vspace{0pt}\sbox\hmboard{\chessboard[setfen=7b/8/8/4k1R1/5Q2/4R3/8/K7 b - - 0 1]}\usebox\hmboard\\[1mm]
\makebox[\wd\hmboard]{\sffamily\hspace{3mm}\textbf{h\#4}\hfill(4+2)}\end{minipage}\hfill
\begin{minipage}[t]{52mm}\vspace{0pt}
\textcolor{ink}{\sffamily\textsc{Lösung}}\par\smallskip\small
\begin{tabular}{@{}r@{\enspace}l@{\enspace}l@{}}
\textbf{1.} & \symking{}xf4+ & \symrook{}ge5\\
\textbf{2.} & \symking{}g4 & \symrook{}g3+\\
\textbf{3.} & \symking{}h4 & \symrook{}g7\\
\textbf{4.} & \symking{}h3 & \symrook{}h5\#\\
\end{tabular}\par\vspace{3mm}
\qrcode[height=12mm]{https://helpman.komtera.lt/?fen=7b/8/8/4k1R1/5Q2/4R3/8/K7&moves=4}\\{\sffamily\tiny\color{ink}scannen zum Lösen}
\end{minipage}
\vspace{2mm}

{\raggedright\noindent\textcolor{ink}{\sffamily\textsc{Themen}}\enspace \themetag{reines Matt} \themetag{Mustermatt} \themetag{Spiegelmatt} \themetag{ein Stein zieht} \themetag{Umnov}\par\smallskip{\footnotesize\textit{Hinweis:} das Diagramm hat keinen legalen letzten Zug, kann also in keiner Partie entstehen.}\par}

\clearpage
\phantomsection\label{hm:KQRRvkn}
\begin{tcolorbox}[enhanced,colback=ink,colframe=ink,arc=0pt,boxrule=0pt,left=4mm,right=4mm,top=2.5mm,bottom=2.5mm]
\color{white}\sffamily{\LARGE\bfseries KQRRvkn}\hfill{\Large h\#4}\\[0.5mm]
{\small 6 Steine\quad Schwarz am Zug}\hfill{\small Nr.~826\qquad\href{https://helpman.komtera.lt/?fen=8/8/8/8/8/3R4/Q1R5/K1k1n3&moves=4}{\color{white}\bfseries lösen\,\ensuremath{\nearrow}}}
\end{tcolorbox}
\vspace{2mm}
\noindent\begin{minipage}[t]{36mm}\vspace{0pt}\sffamily\small\raggedright
\hmfigure{längste eindeutige}{{\Large h\#4}}\hmfigure{eindeutig bei dieser Länge}{{\Large 390} Stellungen}\hmfigure{längstes Matt}{{\Large h\#5}}\hmfigure{in der Klasse}{15,5 Milliarden Stellungen,\\ 4,01 Milliarden mit Hilfsmatt}
\end{minipage}\hfill
\begin{minipage}[t]{80mm}\vspace{0pt}\begin{tikzpicture}\begin{axis}[width=72mm,height=42mm,ybar,bar width=1.6pt,xmin=-1,xmax=11,ymin=0,ymax=1.785,axis y line*=left,axis x line*=bottom,xlabel={\scriptsize Halbzüge bis Matt},ylabel={\scriptsize Stellungen (Mrd.)},tick label style={font=\tiny},label style={font=\scriptsize},enlarge x limits=false,ylabel shift=-3pt,xlabel shift=-3pt]\addplot[fill=ink!55,draw=none] coordinates {(0,0.1142) (1,0.2663) (2,1.007) (3,0.9031) (4,1.623) (5,0.06805) (6,0.02588) (7,4.47e-07) (8,0.0006005) (10,1.1e-07)};\draw[densely dotted] (axis cs:8,0) -- (axis cs:8,1.785);\end{axis}\begin{axis}[width=72mm,height=42mm,axis y line*=right,axis x line=none,ymode=log,xmin=-1,xmax=11,ylabel={\scriptsize eindeutig (log)},tick label style={font=\tiny},label style={font=\scriptsize},ylabel shift=-3pt]\addplot[mark=*,mark size=0.7pt,draw=black,thin] coordinates {(0,114207147) (1,127989724) (2,248875764) (3,13042418) (4,8731276) (5,322) (6,56618) (8,390)};\end{axis}\end{tikzpicture}\end{minipage}
\vspace{3mm}

\noindent\begin{minipage}[t]{60mm}\vspace{0pt}\sbox\hmboard{\chessboard[setfen=8/8/8/8/8/3R4/Q1R5/K1k1n3 b - - 0 1]}\usebox\hmboard\\[1mm]
\makebox[\wd\hmboard]{\sffamily\hspace{3mm}\textbf{h\#4}\hfill(4+2)}\end{minipage}\hfill
\begin{minipage}[t]{52mm}\vspace{0pt}
\textcolor{ink}{\sffamily\textsc{Lösung}}\par\smallskip\small
\begin{tabular}{@{}r@{\enspace}l@{\enspace}l@{}}
\textbf{1.} & \symknight{}xc2+ & \symqueen{}xc2+\\
\textbf{2.} & \symking{}xc2 & \symking{}a2\\
\textbf{3.} & \symking{}c1 & \symking{}b3\\
\textbf{4.} & \symking{}b1 & \symrook{}d1\#\\
\end{tabular}\par\vspace{3mm}
\qrcode[height=12mm]{https://helpman.komtera.lt/?fen=8/8/8/8/8/3R4/Q1R5/K1k1n3&moves=4}\\{\sffamily\tiny\color{ink}scannen zum Lösen}
\end{minipage}
\vspace{2mm}

{\raggedright\noindent\textcolor{ink}{\sffamily\textsc{Themen}}\enspace \themetag{reines Matt} \themetag{Mustermatt} \themetag{Idealmatt} \themetag{Spiegelmatt} \themetag{Switchback}\par\smallskip{\footnotesize\textit{Hinweis:} die Partei am Zug steht im Diagramm im Schach; die Lösung beginnt mit einem Schlag.}\par}

\clearpage
\phantomsection\label{hm:KQRRvkp}
\begin{tcolorbox}[enhanced,colback=ink,colframe=ink,arc=0pt,boxrule=0pt,left=4mm,right=4mm,top=2.5mm,bottom=2.5mm]
\color{white}\sffamily{\LARGE\bfseries KQRRvkp}\hfill{\Large h\#4}\\[0.5mm]
{\small 6 Steine\quad Schwarz am Zug}\hfill{\small Nr.~827\qquad\href{https://helpman.komtera.lt/?fen=8/8/8/1p6/QR6/k7/8/KR6&moves=4}{\color{white}\bfseries lösen\,\ensuremath{\nearrow}}}
\end{tcolorbox}
\vspace{2mm}
\noindent\begin{minipage}[t]{36mm}\vspace{0pt}\sffamily\small\raggedright
\hmfigure{längste eindeutige}{{\Large h\#4}}\hmfigure{eindeutig bei dieser Länge}{{\Large 3} Stellungen}\hmfigure{längstes Matt}{{\Large h\#7}}\hmfigure{in der Klasse}{45,4 Milliarden Stellungen,\\ 12,5 Milliarden mit Hilfsmatt}
\end{minipage}\hfill
\begin{minipage}[t]{80mm}\vspace{0pt}\begin{tikzpicture}\begin{axis}[width=72mm,height=42mm,ybar,bar width=1.6pt,xmin=-1,xmax=15,ymin=0,ymax=5.479,axis y line*=left,axis x line*=bottom,xlabel={\scriptsize Halbzüge bis Matt},ylabel={\scriptsize Stellungen (Mrd.)},tick label style={font=\tiny},label style={font=\scriptsize},enlarge x limits=false,ylabel shift=-3pt,xlabel shift=-3pt]\addplot[fill=ink!55,draw=none] coordinates {(0,0.4151) (1,0.9523) (2,3.321) (3,2.691) (4,4.981) (5,0.09501) (6,0.07497) (7,1.6e-06) (8,0.001562) (9,2e-09) (10,5.43e-07) (12,7.6e-08) (14,1.4e-08)};\draw[densely dotted] (axis cs:8,0) -- (axis cs:8,5.479);\end{axis}\begin{axis}[width=72mm,height=42mm,axis y line*=right,axis x line=none,ymode=log,xmin=-1,xmax=15,ylabel={\scriptsize eindeutig (log)},tick label style={font=\tiny},label style={font=\scriptsize},ylabel shift=-3pt]\addplot[mark=*,mark size=0.7pt,draw=black,thin] coordinates {(0,415080159) (1,426801756) (2,838601724) (3,30437868) (4,25825678) (5,970) (6,145507) (8,3)};\end{axis}\end{tikzpicture}\end{minipage}
\vspace{3mm}

\noindent\begin{minipage}[t]{60mm}\vspace{0pt}\sbox\hmboard{\chessboard[setfen=8/8/8/1p6/QR6/k7/8/KR6 b - - 0 1]}\usebox\hmboard\\[1mm]
\makebox[\wd\hmboard]{\sffamily\hspace{3mm}\textbf{h\#4}\hfill(4+2)}\end{minipage}\hfill
\begin{minipage}[t]{52mm}\vspace{0pt}
\textcolor{ink}{\sffamily\textsc{Lösung}}\par\smallskip\small
\begin{tabular}{@{}r@{\enspace}l@{\enspace}l@{}}
\textbf{1.} & bxa4 & \symrook{}1b3+\\
\textbf{2.} & axb3 & \symking{}b1\\
\textbf{3.} & b2 & \symking{}c2\\
\textbf{4.} & \symking{}a2 & \symrook{}a4\#\\
\end{tabular}\par\vspace{3mm}
\qrcode[height=12mm]{https://helpman.komtera.lt/?fen=8/8/8/1p6/QR6/k7/8/KR6&moves=4}\\{\sffamily\tiny\color{ink}scannen zum Lösen}
\end{minipage}
\vspace{2mm}

{\raggedright\noindent\textcolor{ink}{\sffamily\textsc{Themen}}\enspace \textit{keine erkannt}\par\smallskip{\footnotesize\textit{Hinweis:} die Partei am Zug steht im Diagramm im Schach; die Lösung beginnt mit einem Schlag.}\par}

\clearpage
\phantomsection\label{hm:KQRRvkq}
\begin{tcolorbox}[enhanced,colback=ink,colframe=ink,arc=0pt,boxrule=0pt,left=4mm,right=4mm,top=2.5mm,bottom=2.5mm]
\color{white}\sffamily{\LARGE\bfseries KQRRvkq}\hfill{\Large h\#4}\\[0.5mm]
{\small 6 Steine\quad Schwarz am Zug}\hfill{\small Nr.~828\qquad\href{https://helpman.komtera.lt/?fen=8/6R1/5k2/4Q3/4R3/8/7q/K7&moves=4}{\color{white}\bfseries lösen\,\ensuremath{\nearrow}}}
\end{tcolorbox}
\vspace{2mm}
\noindent\begin{minipage}[t]{36mm}\vspace{0pt}\sffamily\small\raggedright
\hmfigure{längste eindeutige}{{\Large h\#4}}\hmfigure{eindeutig bei dieser Länge}{{\Large 30} Stellungen}\hmfigure{längstes Matt}{{\Large h\#6}}\hmfigure{in der Klasse}{15,5 Milliarden Stellungen,\\ 3,27 Milliarden mit Hilfsmatt}
\end{minipage}\hfill
\begin{minipage}[t]{80mm}\vspace{0pt}\begin{tikzpicture}\begin{axis}[width=72mm,height=42mm,ybar,bar width=1.6pt,xmin=-1,xmax=13,ymin=0,ymax=1.369,axis y line*=left,axis x line*=bottom,xlabel={\scriptsize Halbzüge bis Matt},ylabel={\scriptsize Stellungen (Mrd.)},tick label style={font=\tiny},label style={font=\scriptsize},enlarge x limits=false,ylabel shift=-3pt,xlabel shift=-3pt]\addplot[fill=ink!55,draw=none] coordinates {(0,0.05845) (1,0.1507) (2,0.7092) (3,0.8805) (4,1.244) (5,0.2052) (6,0.02378) (7,2.421e-05) (8,0.0006069) (9,1.42e-07) (10,4.58e-07) (11,1e-09) (12,4.52e-07)};\draw[densely dotted] (axis cs:8,0) -- (axis cs:8,1.369);\end{axis}\begin{axis}[width=72mm,height=42mm,axis y line*=right,axis x line=none,ymode=log,xmin=-1,xmax=13,ylabel={\scriptsize eindeutig (log)},tick label style={font=\tiny},label style={font=\scriptsize},ylabel shift=-3pt]\addplot[mark=*,mark size=0.7pt,draw=black,thin] coordinates {(0,58453986) (1,94415561) (2,230313196) (3,29114256) (4,8146495) (5,12848) (6,81376) (7,64) (8,30)};\end{axis}\end{tikzpicture}\end{minipage}
\vspace{3mm}

\noindent\begin{minipage}[t]{60mm}\vspace{0pt}\sbox\hmboard{\chessboard[setfen=8/6R1/5k2/4Q3/4R3/8/7q/K7 b - - 0 1]}\usebox\hmboard\\[1mm]
\makebox[\wd\hmboard]{\sffamily\hspace{3mm}\textbf{h\#4}\hfill(4+2)}\end{minipage}\hfill
\begin{minipage}[t]{52mm}\vspace{0pt}
\textcolor{ink}{\sffamily\textsc{Lösung}}\par\smallskip\small
\begin{tabular}{@{}r@{\enspace}l@{\enspace}l@{}}
\textbf{1.} & \symqueen{}xe5+ & \symking{}b1\\
\textbf{2.} & \symking{}e6 & \symking{}c2\\
\textbf{3.} & \symking{}d5 & \symrook{}xe5+\\
\textbf{4.} & \symking{}c4 & \symrook{}g4\#\\
\end{tabular}\par\vspace{3mm}
\qrcode[height=12mm]{https://helpman.komtera.lt/?fen=8/6R1/5k2/4Q3/4R3/8/7q/K7&moves=4}\\{\sffamily\tiny\color{ink}scannen zum Lösen}
\end{minipage}
\vspace{2mm}

{\raggedright\noindent\textcolor{ink}{\sffamily\textsc{Themen}}\enspace \themetag{reines Matt} \themetag{Mustermatt} \themetag{Idealmatt} \themetag{Spiegelmatt}\par\smallskip{\footnotesize\textit{Hinweis:} die Partei am Zug steht im Diagramm im Schach; die Lösung beginnt mit einem Schlag.}\par}

\clearpage
\phantomsection\label{hm:KRRRNvk}
\begin{tcolorbox}[enhanced,colback=ink,colframe=ink,arc=0pt,boxrule=0pt,left=4mm,right=4mm,top=2.5mm,bottom=2.5mm]
\color{white}\sffamily{\LARGE\bfseries KRRRNvk}\hfill{\Large h\#4}\\[0.5mm]
{\small 6 Steine\quad Schwarz am Zug}\hfill{\small Nr.~829\qquad\href{https://helpman.komtera.lt/?fen=8/8/8/6N1/3R4/2Rk4/2R5/K7&moves=4}{\color{white}\bfseries lösen\,\ensuremath{\nearrow}}}
\end{tcolorbox}
\vspace{2mm}
\noindent\begin{minipage}[t]{36mm}\vspace{0pt}\sffamily\small\raggedright
\hmfigure{längste eindeutige}{{\Large h\#4}}\hmfigure{eindeutig bei dieser Länge}{{\Large 2} Stellungen}\hmfigure{längstes Matt}{{\Large h\#5}}\hmfigure{in der Klasse}{15,5 Milliarden Stellungen,\\ 1,45 Milliarden mit Hilfsmatt}
\end{minipage}\hfill
\begin{minipage}[t]{80mm}\vspace{0pt}\begin{tikzpicture}\begin{axis}[width=72mm,height=42mm,ybar,bar width=1.6pt,xmin=-1,xmax=11,ymin=0,ymax=712.5,axis y line*=left,axis x line*=bottom,xlabel={\scriptsize Halbzüge bis Matt},ylabel={\scriptsize Stellungen (Mio.)},tick label style={font=\tiny},label style={font=\scriptsize},enlarge x limits=false,ylabel shift=-3pt,xlabel shift=-3pt]\addplot[fill=ink!55,draw=none] coordinates {(0,36.33) (1,92.87) (2,294.3) (3,334.3) (4,647.7) (5,25.06) (6,14.85) (7,0.000263) (8,0.175) (10,1.6e-05)};\draw[densely dotted] (axis cs:8,0) -- (axis cs:8,712.5);\end{axis}\begin{axis}[width=72mm,height=42mm,axis y line*=right,axis x line=none,ymode=log,xmin=-1,xmax=11,ylabel={\scriptsize eindeutig (log)},tick label style={font=\tiny},label style={font=\scriptsize},ylabel shift=-3pt]\addplot[mark=*,mark size=0.7pt,draw=black,thin] coordinates {(0,36332705) (1,57113164) (2,124852195) (3,11965760) (4,7554082) (5,87) (6,31157) (8,2)};\end{axis}\end{tikzpicture}\end{minipage}
\vspace{3mm}

\noindent\begin{minipage}[t]{60mm}\vspace{0pt}\sbox\hmboard{\chessboard[setfen=8/8/8/6N1/3R4/2Rk4/2R5/K7 b - - 0 1]}\usebox\hmboard\\[1mm]
\makebox[\wd\hmboard]{\sffamily\hspace{3mm}\textbf{h\#4}\hfill(5+1)}\end{minipage}\hfill
\begin{minipage}[t]{52mm}\vspace{0pt}
\textcolor{ink}{\sffamily\textsc{Lösung}}\par\smallskip\small
\begin{tabular}{@{}r@{\enspace}l@{\enspace}l@{}}
\textbf{1.} & \symking{}xd4 & \symking{}b2\\
\textbf{2.} & \symking{}e5 & \symrook{}e2+\\
\textbf{3.} & \symking{}d5 & \symrook{}e5+\\
\textbf{4.} & \symking{}d4 & \symknight{}f3\#\\
\end{tabular}\par\vspace{3mm}
\qrcode[height=12mm]{https://helpman.komtera.lt/?fen=8/8/8/6N1/3R4/2Rk4/2R5/K7&moves=4}\\{\sffamily\tiny\color{ink}scannen zum Lösen}
\end{minipage}
\vspace{2mm}

{\raggedright\noindent\textcolor{ink}{\sffamily\textsc{Themen}}\enspace \themetag{Rundlauf} \themetag{ein Stein zieht} \themetag{Kniest} \themetag{Umnov}\par\smallskip{\footnotesize\textit{Hinweis:} das Diagramm hat keinen legalen letzten Zug, kann also in keiner Partie entstehen.}\par}

\clearpage
\phantomsection\label{hm:KRRRvkb}
\begin{tcolorbox}[enhanced,colback=ink,colframe=ink,arc=0pt,boxrule=0pt,left=4mm,right=4mm,top=2.5mm,bottom=2.5mm]
\color{white}\sffamily{\LARGE\bfseries KRRRvkb}\hfill{\Large h\#4}\\[0.5mm]
{\small 6 Steine\quad Schwarz am Zug}\hfill{\small Nr.~830\qquad\href{https://helpman.komtera.lt/?fen=7b/8/6R1/4kR2/6R1/8/8/K7&moves=4}{\color{white}\bfseries lösen\,\ensuremath{\nearrow}}}
\end{tcolorbox}
\vspace{2mm}
\noindent\begin{minipage}[t]{36mm}\vspace{0pt}\sffamily\small\raggedright
\hmfigure{längste eindeutige}{{\Large h\#4}}\hmfigure{eindeutig bei dieser Länge}{{\Large 10} Stellungen}\hmfigure{längstes Matt}{{\Large h\#7}}\hmfigure{in der Klasse}{15,5 Milliarden Stellungen,\\ 1,37 Milliarden mit Hilfsmatt}
\end{minipage}\hfill
\begin{minipage}[t]{80mm}\vspace{0pt}\begin{tikzpicture}\begin{axis}[width=72mm,height=42mm,ybar,bar width=1.6pt,xmin=-1,xmax=15,ymin=0,ymax=674.2,axis y line*=left,axis x line*=bottom,xlabel={\scriptsize Halbzüge bis Matt},ylabel={\scriptsize Stellungen (Mio.)},tick label style={font=\tiny},label style={font=\scriptsize},enlarge x limits=false,ylabel shift=-3pt,xlabel shift=-3pt]\addplot[fill=ink!55,draw=none] coordinates {(0,20.21) (1,58.28) (2,232.5) (3,373.1) (4,612.9) (5,59.42) (6,12.35) (7,0.01448) (8,0.1598) (9,1e-06) (10,5.5e-05) (12,2e-06) (14,1e-06)};\draw[densely dotted] (axis cs:8,0) -- (axis cs:8,674.2);\end{axis}\begin{axis}[width=72mm,height=42mm,axis y line*=right,axis x line=none,ymode=log,xmin=-1,xmax=15,ylabel={\scriptsize eindeutig (log)},tick label style={font=\tiny},label style={font=\scriptsize},ylabel shift=-3pt]\addplot[mark=*,mark size=0.7pt,draw=black,thin] coordinates {(0,20214005) (1,42762070) (2,85137964) (3,17539097) (4,5243545) (5,3485) (6,18578) (8,10)};\end{axis}\end{tikzpicture}\end{minipage}
\vspace{3mm}

\noindent\begin{minipage}[t]{60mm}\vspace{0pt}\sbox\hmboard{\chessboard[setfen=7b/8/6R1/4kR2/6R1/8/8/K7 b - - 0 1]}\usebox\hmboard\\[1mm]
\makebox[\wd\hmboard]{\sffamily\hspace{3mm}\textbf{h\#4}\hfill(4+2)}\end{minipage}\hfill
\begin{minipage}[t]{52mm}\vspace{0pt}
\textcolor{ink}{\sffamily\textsc{Lösung}}\par\smallskip\small
\begin{tabular}{@{}r@{\enspace}l@{\enspace}l@{}}
\textbf{1.} & \symking{}xf5+ & \symrook{}f6+\\
\textbf{2.} & \symking{}e5 & \symking{}b2\\
\textbf{3.} & \symking{}d5 & \symking{}c3\\
\textbf{4.} & \symking{}c5 & \symrook{}g5\#\\
\end{tabular}\par\vspace{3mm}
\qrcode[height=12mm]{https://helpman.komtera.lt/?fen=7b/8/6R1/4kR2/6R1/8/8/K7&moves=4}\\{\sffamily\tiny\color{ink}scannen zum Lösen}
\end{minipage}
\vspace{2mm}

{\raggedright\noindent\textcolor{ink}{\sffamily\textsc{Themen}}\enspace \themetag{reines Matt} \themetag{Mustermatt} \themetag{Spiegelmatt} \themetag{Switchback} \themetag{ein Stein zieht}\par\smallskip{\footnotesize\textit{Hinweis:} die Partei am Zug steht im Diagramm im Schach; die Lösung beginnt mit einem Schlag.}\par}

\clearpage
\phantomsection\label{hm:KRRRvkn}
\begin{tcolorbox}[enhanced,colback=ink,colframe=ink,arc=0pt,boxrule=0pt,left=4mm,right=4mm,top=2.5mm,bottom=2.5mm]
\color{white}\sffamily{\LARGE\bfseries KRRRvkn}\hfill{\Large h\#4}\\[0.5mm]
{\small 6 Steine\quad Schwarz am Zug}\hfill{\small Nr.~831\qquad\href{https://helpman.komtera.lt/?fen=8/8/8/8/8/3R4/R1R5/K1k1n3&moves=4}{\color{white}\bfseries lösen\,\ensuremath{\nearrow}}}
\end{tcolorbox}
\vspace{2mm}
\noindent\begin{minipage}[t]{36mm}\vspace{0pt}\sffamily\small\raggedright
\hmfigure{längste eindeutige}{{\Large h\#4}}\hmfigure{eindeutig bei dieser Länge}{{\Large 175} Stellungen}\hmfigure{längstes Matt}{{\Large h\#5}}\hmfigure{in der Klasse}{15,5 Milliarden Stellungen,\\ 1,41 Milliarden mit Hilfsmatt}
\end{minipage}\hfill
\begin{minipage}[t]{80mm}\vspace{0pt}\begin{tikzpicture}\begin{axis}[width=72mm,height=42mm,ybar,bar width=1.6pt,xmin=-1,xmax=11,ymin=0,ymax=685.8,axis y line*=left,axis x line*=bottom,xlabel={\scriptsize Halbzüge bis Matt},ylabel={\scriptsize Stellungen (Mio.)},tick label style={font=\tiny},label style={font=\scriptsize},enlarge x limits=false,ylabel shift=-3pt,xlabel shift=-3pt]\addplot[fill=ink!55,draw=none] coordinates {(0,25.26) (1,72.38) (2,264.2) (3,371.2) (4,623.4) (5,47.2) (6,11.07) (7,0.01415) (8,0.1582) (10,9.3e-05)};\draw[densely dotted] (axis cs:8,0) -- (axis cs:8,685.8);\end{axis}\begin{axis}[width=72mm,height=42mm,axis y line*=right,axis x line=none,ymode=log,xmin=-1,xmax=11,ylabel={\scriptsize eindeutig (log)},tick label style={font=\tiny},label style={font=\scriptsize},ylabel shift=-3pt]\addplot[mark=*,mark size=0.7pt,draw=black,thin] coordinates {(0,25256515) (1,48957659) (2,89104314) (3,12712596) (4,4727375) (5,2371) (6,12816) (8,175)};\end{axis}\end{tikzpicture}\end{minipage}
\vspace{3mm}

\noindent\begin{minipage}[t]{60mm}\vspace{0pt}\sbox\hmboard{\chessboard[setfen=8/8/8/8/8/3R4/R1R5/K1k1n3 b - - 0 1]}\usebox\hmboard\\[1mm]
\makebox[\wd\hmboard]{\sffamily\hspace{3mm}\textbf{h\#4}\hfill(4+2)}\end{minipage}\hfill
\begin{minipage}[t]{52mm}\vspace{0pt}
\textcolor{ink}{\sffamily\textsc{Lösung}}\par\smallskip\small
\begin{tabular}{@{}r@{\enspace}l@{\enspace}l@{}}
\textbf{1.} & \symknight{}xc2+ & \symrook{}xc2+\\
\textbf{2.} & \symking{}xc2 & \symking{}a2\\
\textbf{3.} & \symking{}c1 & \symking{}b3\\
\textbf{4.} & \symking{}b1 & \symrook{}d1\#\\
\end{tabular}\par\vspace{3mm}
\qrcode[height=12mm]{https://helpman.komtera.lt/?fen=8/8/8/8/8/3R4/R1R5/K1k1n3&moves=4}\\{\sffamily\tiny\color{ink}scannen zum Lösen}
\end{minipage}
\vspace{2mm}

{\raggedright\noindent\textcolor{ink}{\sffamily\textsc{Themen}}\enspace \themetag{reines Matt} \themetag{Mustermatt} \themetag{Idealmatt} \themetag{Spiegelmatt} \themetag{Switchback}\par\smallskip{\footnotesize\textit{Hinweis:} die Partei am Zug steht im Diagramm im Schach; die Lösung beginnt mit einem Schlag.}\par}

\clearpage
\phantomsection\label{hm:KRRRvkp}
\begin{tcolorbox}[enhanced,colback=ink,colframe=ink,arc=0pt,boxrule=0pt,left=4mm,right=4mm,top=2.5mm,bottom=2.5mm]
\color{white}\sffamily{\LARGE\bfseries KRRRvkp}\hfill{\Large h\#4}\\[0.5mm]
{\small 6 Steine\quad Schwarz am Zug}\hfill{\small Nr.~832\qquad\href{https://helpman.komtera.lt/?fen=8/8/8/1p6/RR6/k7/8/KR6&moves=4}{\color{white}\bfseries lösen\,\ensuremath{\nearrow}}}
\end{tcolorbox}
\vspace{2mm}
\noindent\begin{minipage}[t]{36mm}\vspace{0pt}\sffamily\small\raggedright
\hmfigure{längste eindeutige}{{\Large h\#4}}\hmfigure{eindeutig bei dieser Länge}{{\Large 3} Stellungen}\hmfigure{längstes Matt}{{\Large h\#7}}\hmfigure{in der Klasse}{45,4 Milliarden Stellungen,\\ 4,42 Milliarden mit Hilfsmatt}
\end{minipage}\hfill
\begin{minipage}[t]{80mm}\vspace{0pt}\begin{tikzpicture}\begin{axis}[width=72mm,height=42mm,ybar,bar width=1.6pt,xmin=-1,xmax=15,ymin=0,ymax=2.118,axis y line*=left,axis x line*=bottom,xlabel={\scriptsize Halbzüge bis Matt},ylabel={\scriptsize Stellungen (Mrd.)},tick label style={font=\tiny},label style={font=\scriptsize},enlarge x limits=false,ylabel shift=-3pt,xlabel shift=-3pt]\addplot[fill=ink!55,draw=none] coordinates {(0,0.09413) (1,0.2654) (2,0.8804) (3,1.125) (4,1.925) (5,0.08927) (6,0.03514) (7,1.009e-05) (8,0.0004598) (9,5e-09) (10,5.26e-07) (12,7.5e-08) (14,1.4e-08)};\draw[densely dotted] (axis cs:8,0) -- (axis cs:8,2.118);\end{axis}\begin{axis}[width=72mm,height=42mm,axis y line*=right,axis x line=none,ymode=log,xmin=-1,xmax=15,ylabel={\scriptsize eindeutig (log)},tick label style={font=\tiny},label style={font=\scriptsize},ylabel shift=-3pt]\addplot[mark=*,mark size=0.7pt,draw=black,thin] coordinates {(0,94131161) (1,174222181) (2,322210889) (3,35466972) (4,16629065) (5,3100) (6,35227) (7,1) (8,3)};\end{axis}\end{tikzpicture}\end{minipage}
\vspace{3mm}

\noindent\begin{minipage}[t]{60mm}\vspace{0pt}\sbox\hmboard{\chessboard[setfen=8/8/8/1p6/RR6/k7/8/KR6 b - - 0 1]}\usebox\hmboard\\[1mm]
\makebox[\wd\hmboard]{\sffamily\hspace{3mm}\textbf{h\#4}\hfill(4+2)}\end{minipage}\hfill
\begin{minipage}[t]{52mm}\vspace{0pt}
\textcolor{ink}{\sffamily\textsc{Lösung}}\par\smallskip\small
\begin{tabular}{@{}r@{\enspace}l@{\enspace}l@{}}
\textbf{1.} & bxa4 & \symrook{}1b3+\\
\textbf{2.} & axb3 & \symking{}b1\\
\textbf{3.} & b2 & \symking{}c2\\
\textbf{4.} & \symking{}a2 & \symrook{}a4\#\\
\end{tabular}\par\vspace{3mm}
\qrcode[height=12mm]{https://helpman.komtera.lt/?fen=8/8/8/1p6/RR6/k7/8/KR6&moves=4}\\{\sffamily\tiny\color{ink}scannen zum Lösen}
\end{minipage}
\vspace{2mm}

{\raggedright\noindent\textcolor{ink}{\sffamily\textsc{Themen}}\enspace \textit{keine erkannt}\par\smallskip{\footnotesize\textit{Hinweis:} die Partei am Zug steht im Diagramm im Schach; die Lösung beginnt mit einem Schlag.}\par}

\clearpage
\phantomsection\label{hm:KQQBvkr}
\begin{tcolorbox}[enhanced,colback=ink,colframe=ink,arc=0pt,boxrule=0pt,left=4mm,right=4mm,top=2.5mm,bottom=2.5mm]
\color{white}\sffamily{\LARGE\bfseries KQQBvkr}\hfill{\Large h\#3.5}\\[0.5mm]
{\small 6 Steine\quad Weiß am Zug}\hfill{\small Nr.~833\qquad\href{https://helpman.komtera.lt/?fen=8/8/BQ6/8/8/2k5/4Q3/3K1r2&moves=3.5}{\color{white}\bfseries lösen\,\ensuremath{\nearrow}}}
\end{tcolorbox}
\vspace{2mm}
\noindent\begin{minipage}[t]{36mm}\vspace{0pt}\sffamily\small\raggedright
\hmfigure{längste eindeutige}{{\Large h\#3.5}}\hmfigure{eindeutig bei dieser Länge}{{\Large 1} Stellung}\hmfigure{längstes Matt}{{\Large h\#5}}\hmfigure{in der Klasse}{15,5 Milliarden Stellungen,\\ 3,57 Milliarden mit Hilfsmatt}
\end{minipage}\hfill
\begin{minipage}[t]{80mm}\vspace{0pt}\begin{tikzpicture}\begin{axis}[width=72mm,height=42mm,ybar,bar width=1.6pt,xmin=-1,xmax=11,ymin=0,ymax=1.356,axis y line*=left,axis x line*=bottom,xlabel={\scriptsize Halbzüge bis Matt},ylabel={\scriptsize Stellungen (Mrd.)},tick label style={font=\tiny},label style={font=\scriptsize},enlarge x limits=false,ylabel shift=-3pt,xlabel shift=-3pt]\addplot[fill=ink!55,draw=none] coordinates {(0,0.09673) (1,0.261) (2,1.233) (3,0.8229) (4,1.058) (5,0.06407) (6,0.0337) (7,2.04e-07) (8,4.984e-05) (9,1.3e-08) (10,6e-09)};\draw[densely dotted] (axis cs:7,0) -- (axis cs:7,1.356);\end{axis}\begin{axis}[width=72mm,height=42mm,axis y line*=right,axis x line=none,ymode=log,xmin=-1,xmax=11,ylabel={\scriptsize eindeutig (log)},tick label style={font=\tiny},label style={font=\scriptsize},ylabel shift=-3pt]\addplot[mark=*,mark size=0.7pt,draw=black,thin] coordinates {(0,96726236) (1,137018279) (2,363826358) (3,16603412) (4,13513802) (5,31) (6,16125) (7,1)};\end{axis}\end{tikzpicture}\end{minipage}
\vspace{3mm}

\noindent\begin{minipage}[t]{60mm}\vspace{0pt}\sbox\hmboard{\chessboard[setfen=8/8/BQ6/8/8/2k5/4Q3/3K1r2 w - - 0 1]}\usebox\hmboard\\[1mm]
\makebox[\wd\hmboard]{\sffamily\hspace{3mm}\textbf{h\#3.5}\hfill(4+2)}\end{minipage}\hfill
\begin{minipage}[t]{52mm}\vspace{0pt}
\textcolor{ink}{\sffamily\textsc{Lösung}}\par\smallskip\small
\begin{tabular}{@{}r@{\enspace}l@{\enspace}l@{}}
\textbf{1.} & \ldots & \symqueen{}e1+\\
\textbf{2.} & \symrook{}xe1+ & \symking{}xe1\\
\textbf{3.} & \symking{}c2 & \symbishop{}d3+\\
\textbf{4.} & \symking{}c1 & \symqueen{}b1\#\\
\end{tabular}\par\vspace{3mm}
\qrcode[height=12mm]{https://helpman.komtera.lt/?fen=8/8/BQ6/8/8/2k5/4Q3/3K1r2&moves=3.5}\\{\sffamily\tiny\color{ink}scannen zum Lösen}
\end{minipage}
\vspace{2mm}

{\raggedright\noindent\textcolor{ink}{\sffamily\textsc{Themen}}\enspace \textit{keine erkannt}\par\smallskip{\footnotesize\textit{Hinweis:} die Partei am Zug steht im Diagramm im Schach.}\par}

\clearpage
\phantomsection\label{hm:KQRRvkr}
\begin{tcolorbox}[enhanced,colback=ink,colframe=ink,arc=0pt,boxrule=0pt,left=4mm,right=4mm,top=2.5mm,bottom=2.5mm]
\color{white}\sffamily{\LARGE\bfseries KQRRvkr}\hfill{\Large h\#3.5}\\[0.5mm]
{\small 6 Steine\quad Weiß am Zug}\hfill{\small Nr.~834\qquad\href{https://helpman.komtera.lt/?fen=8/8/8/3Q4/8/r3R3/2k5/KR6&moves=3.5}{\color{white}\bfseries lösen\,\ensuremath{\nearrow}}}
\end{tcolorbox}
\vspace{2mm}
\noindent\begin{minipage}[t]{36mm}\vspace{0pt}\sffamily\small\raggedright
\hmfigure{längste eindeutige}{{\Large h\#3.5}}\hmfigure{eindeutig bei dieser Länge}{{\Large 8} Stellungen}\hmfigure{längstes Matt}{{\Large h\#5}}\hmfigure{in der Klasse}{15,5 Milliarden Stellungen,\\ 3,66 Milliarden mit Hilfsmatt}
\end{minipage}\hfill
\begin{minipage}[t]{80mm}\vspace{0pt}\begin{tikzpicture}\begin{axis}[width=72mm,height=42mm,ybar,bar width=1.6pt,xmin=-1,xmax=11,ymin=0,ymax=1.528,axis y line*=left,axis x line*=bottom,xlabel={\scriptsize Halbzüge bis Matt},ylabel={\scriptsize Stellungen (Mrd.)},tick label style={font=\tiny},label style={font=\scriptsize},enlarge x limits=false,ylabel shift=-3pt,xlabel shift=-3pt]\addplot[fill=ink!55,draw=none] coordinates {(0,0.08776) (1,0.221) (2,0.9236) (3,0.8905) (4,1.389) (5,0.1259) (6,0.0222) (7,1.23e-06) (8,0.0005053) (9,6e-09) (10,1.9e-08)};\draw[densely dotted] (axis cs:7,0) -- (axis cs:7,1.528);\end{axis}\begin{axis}[width=72mm,height=42mm,axis y line*=right,axis x line=none,ymode=log,xmin=-1,xmax=11,ylabel={\scriptsize eindeutig (log)},tick label style={font=\tiny},label style={font=\scriptsize},ylabel shift=-3pt]\addplot[mark=*,mark size=0.7pt,draw=black,thin] coordinates {(0,87760679) (1,120302868) (2,253449057) (3,14412437) (4,8019245) (5,884) (6,51484) (7,8)};\end{axis}\end{tikzpicture}\end{minipage}
\vspace{3mm}

\noindent\begin{minipage}[t]{60mm}\vspace{0pt}\sbox\hmboard{\chessboard[setfen=8/8/8/3Q4/8/r3R3/2k5/KR6 w - - 0 1]}\usebox\hmboard\\[1mm]
\makebox[\wd\hmboard]{\sffamily\hspace{3mm}\textbf{h\#3.5}\hfill(4+2)}\end{minipage}\hfill
\begin{minipage}[t]{52mm}\vspace{0pt}
\textcolor{ink}{\sffamily\textsc{Lösung}}\par\smallskip\small
\begin{tabular}{@{}r@{\enspace}l@{\enspace}l@{}}
\textbf{1.} & \ldots & \symqueen{}a2+\\
\textbf{2.} & \symrook{}xa2+ & \symking{}xa2\\
\textbf{3.} & \symking{}d2 & \symrook{}b2+\\
\textbf{4.} & \symking{}c1 & \symrook{}e1\#\\
\end{tabular}\par\vspace{3mm}
\qrcode[height=12mm]{https://helpman.komtera.lt/?fen=8/8/8/3Q4/8/r3R3/2k5/KR6&moves=3.5}\\{\sffamily\tiny\color{ink}scannen zum Lösen}
\end{minipage}
\vspace{2mm}

{\raggedright\noindent\textcolor{ink}{\sffamily\textsc{Themen}}\enspace \textit{keine erkannt}\par\smallskip{\footnotesize\textit{Hinweis:} die Partei am Zug steht im Diagramm im Schach.}\par}

\clearpage
\phantomsection\label{hm:KRRRvkq}
\begin{tcolorbox}[enhanced,colback=ink,colframe=ink,arc=0pt,boxrule=0pt,left=4mm,right=4mm,top=2.5mm,bottom=2.5mm]
\color{white}\sffamily{\LARGE\bfseries KRRRvkq}\hfill{\Large h\#3.5}\\[0.5mm]
{\small 6 Steine\quad Weiß am Zug}\hfill{\small Nr.~835\qquad\href{https://helpman.komtera.lt/?fen=8/8/8/8/3R4/2k5/1q6/KR2R3&moves=3.5}{\color{white}\bfseries lösen\,\ensuremath{\nearrow}}}
\end{tcolorbox}
\vspace{2mm}
\noindent\begin{minipage}[t]{36mm}\vspace{0pt}\sffamily\small\raggedright
\hmfigure{längste eindeutige}{{\Large h\#3.5}}\hmfigure{eindeutig bei dieser Länge}{{\Large 56} Stellungen}\hmfigure{längstes Matt}{{\Large h\#6}}\hmfigure{in der Klasse}{15,5 Milliarden Stellungen,\\ 1,17 Milliarden mit Hilfsmatt}
\end{minipage}\hfill
\begin{minipage}[t]{80mm}\vspace{0pt}\begin{tikzpicture}\begin{axis}[width=72mm,height=42mm,ybar,bar width=1.6pt,xmin=-1,xmax=13,ymin=0,ymax=547.6,axis y line*=left,axis x line*=bottom,xlabel={\scriptsize Halbzüge bis Matt},ylabel={\scriptsize Stellungen (Mio.)},tick label style={font=\tiny},label style={font=\scriptsize},enlarge x limits=false,ylabel shift=-3pt,xlabel shift=-3pt]\addplot[fill=ink!55,draw=none] coordinates {(0,11.01) (1,32.28) (2,158.7) (3,338.6) (4,497.8) (5,119.5) (6,11.54) (7,0.1076) (8,0.1709) (9,0.00056) (10,7e-05) (12,4.4e-05)};\draw[densely dotted] (axis cs:7,0) -- (axis cs:7,547.6);\end{axis}\begin{axis}[width=72mm,height=42mm,axis y line*=right,axis x line=none,ymode=log,xmin=-1,xmax=13,ylabel={\scriptsize eindeutig (log)},tick label style={font=\tiny},label style={font=\scriptsize},ylabel shift=-3pt]\addplot[mark=*,mark size=0.7pt,draw=black,thin] coordinates {(0,11010902) (1,26654426) (2,62083203) (3,23548375) (4,4755725) (5,23861) (6,22748) (7,56)};\end{axis}\end{tikzpicture}\end{minipage}
\vspace{3mm}

\noindent\begin{minipage}[t]{60mm}\vspace{0pt}\sbox\hmboard{\chessboard[setfen=8/8/8/8/3R4/2k5/1q6/KR2R3 w - - 0 1]}\usebox\hmboard\\[1mm]
\makebox[\wd\hmboard]{\sffamily\hspace{3mm}\textbf{h\#3.5}\hfill(4+2)}\end{minipage}\hfill
\begin{minipage}[t]{52mm}\vspace{0pt}
\textcolor{ink}{\sffamily\textsc{Lösung}}\par\smallskip\small
\begin{tabular}{@{}r@{\enspace}l@{\enspace}l@{}}
\textbf{1.} & \ldots & \symrook{}xb2\\
\textbf{2.} & \symking{}xd4 & \symrook{}bb1\\
\textbf{3.} & \symking{}d3 & \symking{}b2\\
\textbf{4.} & \symking{}d2 & \symrook{}bd1\#\\
\end{tabular}\par\vspace{3mm}
\qrcode[height=12mm]{https://helpman.komtera.lt/?fen=8/8/8/8/3R4/2k5/1q6/KR2R3&moves=3.5}\\{\sffamily\tiny\color{ink}scannen zum Lösen}
\end{minipage}
\vspace{2mm}

{\raggedright\noindent\textcolor{ink}{\sffamily\textsc{Themen}}\enspace \themetag{Switchback} \themetag{ein Stein zieht} \themetag{ohne Schach}\par\smallskip{\footnotesize\textit{Hinweis:} die Partei am Zug steht im Diagramm im Schach; die Lösung beginnt mit einem Schlag.}\par}

\clearpage
\phantomsection\label{hm:KRRRvkr}
\begin{tcolorbox}[enhanced,colback=ink,colframe=ink,arc=0pt,boxrule=0pt,left=4mm,right=4mm,top=2.5mm,bottom=2.5mm]
\color{white}\sffamily{\LARGE\bfseries KRRRvkr}\hfill{\Large h\#3.5}\\[0.5mm]
{\small 6 Steine\quad Weiß am Zug}\hfill{\small Nr.~836\qquad\href{https://helpman.komtera.lt/?fen=8/8/8/8/3R4/2k5/1R6/rK2R3&moves=3.5}{\color{white}\bfseries lösen\,\ensuremath{\nearrow}}}
\end{tcolorbox}
\vspace{2mm}
\noindent\begin{minipage}[t]{36mm}\vspace{0pt}\sffamily\small\raggedright
\hmfigure{längste eindeutige}{{\Large h\#3.5}}\hmfigure{eindeutig bei dieser Länge}{{\Large 2} Stellungen}\hmfigure{längstes Matt}{{\Large h\#4.5}}\hmfigure{in der Klasse}{15,5 Milliarden Stellungen,\\ 1,3 Milliarden mit Hilfsmatt}
\end{minipage}\hfill
\begin{minipage}[t]{80mm}\vspace{0pt}\begin{tikzpicture}\begin{axis}[width=72mm,height=42mm,ybar,bar width=1.6pt,xmin=-1,xmax=10,ymin=0,ymax=610.9,axis y line*=left,axis x line*=bottom,xlabel={\scriptsize Halbzüge bis Matt},ylabel={\scriptsize Stellungen (Mio.)},tick label style={font=\tiny},label style={font=\scriptsize},enlarge x limits=false,ylabel shift=-3pt,xlabel shift=-3pt]\addplot[fill=ink!55,draw=none] coordinates {(0,18.16) (1,54.93) (2,224.9) (3,361.4) (4,555.4) (5,74.42) (6,9.527) (7,0.03154) (8,0.1549) (9,7.1e-05)};\draw[densely dotted] (axis cs:7,0) -- (axis cs:7,610.9);\end{axis}\begin{axis}[width=72mm,height=42mm,axis y line*=right,axis x line=none,ymode=log,xmin=-1,xmax=10,ylabel={\scriptsize eindeutig (log)},tick label style={font=\tiny},label style={font=\scriptsize},ylabel shift=-3pt]\addplot[mark=*,mark size=0.7pt,draw=black,thin] coordinates {(0,18155060) (1,41806903) (2,81044901) (3,14223865) (4,4291208) (5,7233) (6,15216) (7,2)};\end{axis}\end{tikzpicture}\end{minipage}
\vspace{3mm}

\noindent\begin{minipage}[t]{60mm}\vspace{0pt}\sbox\hmboard{\chessboard[setfen=8/8/8/8/3R4/2k5/1R6/rK2R3 w - - 0 1]}\usebox\hmboard\\[1mm]
\makebox[\wd\hmboard]{\sffamily\hspace{3mm}\textbf{h\#3.5}\hfill(4+2)}\end{minipage}\hfill
\begin{minipage}[t]{52mm}\vspace{0pt}
\textcolor{ink}{\sffamily\textsc{Lösung}}\par\smallskip\small
\begin{tabular}{@{}r@{\enspace}l@{\enspace}l@{}}
\textbf{1.} & \ldots & \symking{}xa1\\
\textbf{2.} & \symking{}xd4 & \symrook{}bb1\\
\textbf{3.} & \symking{}d3 & \symking{}b2\\
\textbf{4.} & \symking{}d2 & \symrook{}bd1\#\\
\end{tabular}\par\vspace{3mm}
\qrcode[height=12mm]{https://helpman.komtera.lt/?fen=8/8/8/8/3R4/2k5/1R6/rK2R3&moves=3.5}\\{\sffamily\tiny\color{ink}scannen zum Lösen}
\end{minipage}
\vspace{2mm}

{\raggedright\noindent\textcolor{ink}{\sffamily\textsc{Themen}}\enspace \themetag{ein Stein zieht} \themetag{ohne Schach}\par\smallskip{\footnotesize\textit{Hinweis:} die Partei am Zug steht im Diagramm im Schach; die Lösung beginnt mit einem Schlag.}\par}

\clearpage
\phantomsection\label{hm:KQQBBvk}
\begin{tcolorbox}[enhanced,colback=ink,colframe=ink,arc=0pt,boxrule=0pt,left=4mm,right=4mm,top=2.5mm,bottom=2.5mm]
\color{white}\sffamily{\LARGE\bfseries KQQBBvk}\hfill{\Large h\#3}\\[0.5mm]
{\small 6 Steine\quad Schwarz am Zug}\hfill{\small Nr.~837\qquad\href{https://helpman.komtera.lt/?fen=8/8/8/8/4Q3/3k4/8/K1B1BQ2&moves=3}{\color{white}\bfseries lösen\,\ensuremath{\nearrow}}}
\end{tcolorbox}
\vspace{2mm}
\noindent\begin{minipage}[t]{36mm}\vspace{0pt}\sffamily\small\raggedright
\hmfigure{längste eindeutige}{{\Large h\#3}}\hmfigure{eindeutig bei dieser Länge}{{\Large 631} Stellungen}\hmfigure{längstes Matt}{{\Large h\#4}}\hmfigure{in der Klasse}{15,5 Milliarden Stellungen,\\ 1,95 Milliarden mit Hilfsmatt}
\end{minipage}\hfill
\begin{minipage}[t]{80mm}\vspace{0pt}\begin{tikzpicture}\begin{axis}[width=72mm,height=42mm,ybar,bar width=1.6pt,xmin=-1,xmax=9,ymin=0,ymax=864.5,axis y line*=left,axis x line*=bottom,xlabel={\scriptsize Halbzüge bis Matt},ylabel={\scriptsize Stellungen (Mio.)},tick label style={font=\tiny},label style={font=\scriptsize},enlarge x limits=false,ylabel shift=-3pt,xlabel shift=-3pt]\addplot[fill=ink!55,draw=none] coordinates {(0,116.4) (1,203.9) (2,785.9) (3,282.9) (4,549.2) (5,2.255) (6,11.61) (8,0.001267)};\draw[densely dotted] (axis cs:6,0) -- (axis cs:6,864.5);\end{axis}\begin{axis}[width=72mm,height=42mm,axis y line*=right,axis x line=none,ymode=log,xmin=-1,xmax=9,ylabel={\scriptsize eindeutig (log)},tick label style={font=\tiny},label style={font=\scriptsize},ylabel shift=-3pt]\addplot[mark=*,mark size=0.7pt,draw=black,thin] coordinates {(0,116351848) (1,77893052) (2,214476112) (3,1860078) (4,5858912) (6,631)};\end{axis}\end{tikzpicture}\end{minipage}
\vspace{3mm}

\noindent\begin{minipage}[t]{60mm}\vspace{0pt}\sbox\hmboard{\chessboard[setfen=8/8/8/8/4Q3/3k4/8/K1B1BQ2 b - - 0 1]}\usebox\hmboard\\[1mm]
\makebox[\wd\hmboard]{\sffamily\hspace{3mm}\textbf{h\#3}\hfill(5+1)}\end{minipage}\hfill
\begin{minipage}[t]{52mm}\vspace{0pt}
\textcolor{ink}{\sffamily\textsc{Lösung}}\par\smallskip\small
\begin{tabular}{@{}r@{\enspace}l@{\enspace}l@{}}
\textbf{1.} & \symking{}xe4 & \symqueen{}f4+\\
\textbf{2.} & \symking{}d3 & \symking{}a2\\
\textbf{3.} & \symking{}c2 & \symqueen{}d2\#\\
\end{tabular}\par\vspace{3mm}
\qrcode[height=12mm]{https://helpman.komtera.lt/?fen=8/8/8/8/4Q3/3k4/8/K1B1BQ2&moves=3}\\{\sffamily\tiny\color{ink}scannen zum Lösen}
\end{minipage}
\vspace{2mm}

{\raggedright\noindent\textcolor{ink}{\sffamily\textsc{Themen}}\enspace \themetag{Switchback} \themetag{ein Stein zieht}\par\smallskip{\footnotesize\textit{Hinweis:} das Diagramm hat keinen legalen letzten Zug, kann also in keiner Partie entstehen.}\par}

\clearpage
\phantomsection\label{hm:KQQBNvk}
\begin{tcolorbox}[enhanced,colback=ink,colframe=ink,arc=0pt,boxrule=0pt,left=4mm,right=4mm,top=2.5mm,bottom=2.5mm]
\color{white}\sffamily{\LARGE\bfseries KQQBNvk}\hfill{\Large h\#3}\\[0.5mm]
{\small 6 Steine\quad Schwarz am Zug}\hfill{\small Nr.~838\qquad\href{https://helpman.komtera.lt/?fen=7B/6N1/8/8/8/8/3Q4/KQk5&moves=3}{\color{white}\bfseries lösen\,\ensuremath{\nearrow}}}
\end{tcolorbox}
\vspace{2mm}
\noindent\begin{minipage}[t]{36mm}\vspace{0pt}\sffamily\small\raggedright
\hmfigure{längste eindeutige}{{\Large h\#3}}\hmfigure{eindeutig bei dieser Länge}{{\Large 1\,128} Stellungen}\hmfigure{längstes Matt}{{\Large h\#4}}\hmfigure{in der Klasse}{15,5 Milliarden Stellungen,\\ 4 Milliarden mit Hilfsmatt}
\end{minipage}\hfill
\begin{minipage}[t]{80mm}\vspace{0pt}\begin{tikzpicture}\begin{axis}[width=72mm,height=42mm,ybar,bar width=1.6pt,xmin=-1,xmax=9,ymin=0,ymax=1.662,axis y line*=left,axis x line*=bottom,xlabel={\scriptsize Halbzüge bis Matt},ylabel={\scriptsize Stellungen (Mrd.)},tick label style={font=\tiny},label style={font=\scriptsize},enlarge x limits=false,ylabel shift=-3pt,xlabel shift=-3pt]\addplot[fill=ink!55,draw=none] coordinates {(0,0.1864) (1,0.4002) (2,1.511) (3,0.6566) (4,1.207) (5,0.006102) (6,0.02974) (8,6.574e-06)};\draw[densely dotted] (axis cs:6,0) -- (axis cs:6,1.662);\end{axis}\begin{axis}[width=72mm,height=42mm,axis y line*=right,axis x line=none,ymode=log,xmin=-1,xmax=9,ylabel={\scriptsize eindeutig (log)},tick label style={font=\tiny},label style={font=\scriptsize},ylabel shift=-3pt]\addplot[mark=*,mark size=0.7pt,draw=black,thin] coordinates {(0,186375974) (1,169506295) (2,458774013) (3,5232612) (4,17979372) (6,1128)};\end{axis}\end{tikzpicture}\end{minipage}
\vspace{3mm}

\noindent\begin{minipage}[t]{60mm}\vspace{0pt}\sbox\hmboard{\chessboard[setfen=7B/6N1/8/8/8/8/3Q4/KQk5 b - - 0 1]}\usebox\hmboard\\[1mm]
\makebox[\wd\hmboard]{\sffamily\hspace{3mm}\textbf{h\#3}\hfill(5+1)}\end{minipage}\hfill
\begin{minipage}[t]{52mm}\vspace{0pt}
\textcolor{ink}{\sffamily\textsc{Lösung}}\par\smallskip\small
\begin{tabular}{@{}r@{\enspace}l@{\enspace}l@{}}
\textbf{1.} & \symking{}xd2 & \symking{}a2\\
\textbf{2.} & \symking{}c3 & \symknight{}e6+\\
\textbf{3.} & \symking{}c4 & \symqueen{}b3\#\\
\end{tabular}\par\vspace{3mm}
\qrcode[height=12mm]{https://helpman.komtera.lt/?fen=7B/6N1/8/8/8/8/3Q4/KQk5&moves=3}\\{\sffamily\tiny\color{ink}scannen zum Lösen}
\end{minipage}
\vspace{2mm}

{\raggedright\noindent\textcolor{ink}{\sffamily\textsc{Themen}}\enspace \themetag{ein Stein zieht}\par\smallskip{\footnotesize\textit{Hinweis:} das Diagramm hat keinen legalen letzten Zug, kann also in keiner Partie entstehen.}\par}

\clearpage
\phantomsection\label{hm:KQQBPvk}
\begin{tcolorbox}[enhanced,colback=ink,colframe=ink,arc=0pt,boxrule=0pt,left=4mm,right=4mm,top=2.5mm,bottom=2.5mm]
\color{white}\sffamily{\LARGE\bfseries KQQBPvk}\hfill{\Large h\#3}\\[0.5mm]
{\small 6 Steine\quad Schwarz am Zug}\hfill{\small Nr.~839\qquad\href{https://helpman.komtera.lt/?fen=7B/8/8/4P3/8/8/3Q4/KQk5&moves=3}{\color{white}\bfseries lösen\,\ensuremath{\nearrow}}}
\end{tcolorbox}
\vspace{2mm}
\noindent\begin{minipage}[t]{36mm}\vspace{0pt}\sffamily\small\raggedright
\hmfigure{längste eindeutige}{{\Large h\#3}}\hmfigure{eindeutig bei dieser Länge}{{\Large 49\,428} Stellungen}\hmfigure{längstes Matt}{{\Large h\#5}}\hmfigure{in der Klasse}{45,4 Milliarden Stellungen,\\ 12,2 Milliarden mit Hilfsmatt}
\end{minipage}\hfill
\begin{minipage}[t]{80mm}\vspace{0pt}\begin{tikzpicture}\begin{axis}[width=72mm,height=42mm,ybar,bar width=1.6pt,xmin=-1,xmax=11,ymin=0,ymax=4.615,axis y line*=left,axis x line*=bottom,xlabel={\scriptsize Halbzüge bis Matt},ylabel={\scriptsize Stellungen (Mrd.)},tick label style={font=\tiny},label style={font=\scriptsize},enlarge x limits=false,ylabel shift=-3pt,xlabel shift=-3pt]\addplot[fill=ink!55,draw=none] coordinates {(0,0.4893) (1,1.11) (2,4.196) (3,2.221) (4,4.028) (5,0.03786) (6,0.1194) (8,0.0002253) (10,4e-09)};\draw[densely dotted] (axis cs:6,0) -- (axis cs:6,4.615);\end{axis}\begin{axis}[width=72mm,height=42mm,axis y line*=right,axis x line=none,ymode=log,xmin=-1,xmax=11,ylabel={\scriptsize eindeutig (log)},tick label style={font=\tiny},label style={font=\scriptsize},ylabel shift=-3pt]\addplot[mark=*,mark size=0.7pt,draw=black,thin] coordinates {(0,489279892) (1,482772804) (2,1364861504) (3,28893766) (4,60330712) (6,49428)};\end{axis}\end{tikzpicture}\end{minipage}
\vspace{3mm}

\noindent\begin{minipage}[t]{60mm}\vspace{0pt}\sbox\hmboard{\chessboard[setfen=7B/8/8/4P3/8/8/3Q4/KQk5 b - - 0 1]}\usebox\hmboard\\[1mm]
\makebox[\wd\hmboard]{\sffamily\hspace{3mm}\textbf{h\#3}\hfill(5+1)}\end{minipage}\hfill
\begin{minipage}[t]{52mm}\vspace{0pt}
\textcolor{ink}{\sffamily\textsc{Lösung}}\par\smallskip\small
\begin{tabular}{@{}r@{\enspace}l@{\enspace}l@{}}
\textbf{1.} & \symking{}xd2 & \symbishop{}f6\\
\textbf{2.} & \symking{}e2 & \symqueen{}d3+\\
\textbf{3.} & \symking{}e1 & \symbishop{}h4\#\\
\end{tabular}\par\vspace{3mm}
\qrcode[height=12mm]{https://helpman.komtera.lt/?fen=7B/8/8/4P3/8/8/3Q4/KQk5&moves=3}\\{\sffamily\tiny\color{ink}scannen zum Lösen}
\end{minipage}
\vspace{2mm}

{\raggedright\noindent\textcolor{ink}{\sffamily\textsc{Themen}}\enspace \themetag{reines Matt} \themetag{Mustermatt} \themetag{Spiegelmatt} \themetag{ein Stein zieht}\par\smallskip{\footnotesize\textit{Hinweis:} das Diagramm hat keinen legalen letzten Zug, kann also in keiner Partie entstehen.}\par}

\clearpage
\phantomsection\label{hm:KQQBvkb}
\begin{tcolorbox}[enhanced,colback=ink,colframe=ink,arc=0pt,boxrule=0pt,left=4mm,right=4mm,top=2.5mm,bottom=2.5mm]
\color{white}\sffamily{\LARGE\bfseries KQQBvkb}\hfill{\Large h\#3}\\[0.5mm]
{\small 6 Steine\quad Schwarz am Zug}\hfill{\small Nr.~840\qquad\href{https://helpman.komtera.lt/?fen=8/8/8/8/8/b7/BQ2Q3/K1k5&moves=3}{\color{white}\bfseries lösen\,\ensuremath{\nearrow}}}
\end{tcolorbox}
\vspace{2mm}
\noindent\begin{minipage}[t]{36mm}\vspace{0pt}\sffamily\small\raggedright
\hmfigure{längste eindeutige}{{\Large h\#3}}\hmfigure{eindeutig bei dieser Länge}{{\Large 16\,298} Stellungen}\hmfigure{längstes Matt}{{\Large h\#4}}\hmfigure{in der Klasse}{15,5 Milliarden Stellungen,\\ 3,78 Milliarden mit Hilfsmatt}
\end{minipage}\hfill
\begin{minipage}[t]{80mm}\vspace{0pt}\begin{tikzpicture}\begin{axis}[width=72mm,height=42mm,ybar,bar width=1.6pt,xmin=-1,xmax=9,ymin=0,ymax=1.468,axis y line*=left,axis x line*=bottom,xlabel={\scriptsize Halbzüge bis Matt},ylabel={\scriptsize Stellungen (Mrd.)},tick label style={font=\tiny},label style={font=\scriptsize},enlarge x limits=false,ylabel shift=-3pt,xlabel shift=-3pt]\addplot[fill=ink!55,draw=none] coordinates {(0,0.1135) (1,0.2939) (2,1.335) (3,0.8097) (4,1.148) (5,0.04452) (6,0.0356) (7,1.192e-06) (8,5.133e-05)};\draw[densely dotted] (axis cs:6,0) -- (axis cs:6,1.468);\end{axis}\begin{axis}[width=72mm,height=42mm,axis y line*=right,axis x line=none,ymode=log,xmin=-1,xmax=9,ylabel={\scriptsize eindeutig (log)},tick label style={font=\tiny},label style={font=\scriptsize},ylabel shift=-3pt]\addplot[mark=*,mark size=0.7pt,draw=black,thin] coordinates {(0,113482707) (1,145066956) (2,374189383) (3,11853833) (4,15082334) (5,178) (6,16298)};\end{axis}\end{tikzpicture}\end{minipage}
\vspace{3mm}

\noindent\begin{minipage}[t]{60mm}\vspace{0pt}\sbox\hmboard{\chessboard[setfen=8/8/8/8/8/b7/BQ2Q3/K1k5 b - - 0 1]}\usebox\hmboard\\[1mm]
\makebox[\wd\hmboard]{\sffamily\hspace{3mm}\textbf{h\#3}\hfill(4+2)}\end{minipage}\hfill
\begin{minipage}[t]{52mm}\vspace{0pt}
\textcolor{ink}{\sffamily\textsc{Lösung}}\par\smallskip\small
\begin{tabular}{@{}r@{\enspace}l@{\enspace}l@{}}
\textbf{1.} & \symbishop{}xb2+ & \symqueen{}xb2+\\
\textbf{2.} & \symking{}d1 & \symbishop{}c4\\
\textbf{3.} & \symking{}e1 & \symqueen{}e2\#\\
\end{tabular}\par\vspace{3mm}
\qrcode[height=12mm]{https://helpman.komtera.lt/?fen=8/8/8/8/8/b7/BQ2Q3/K1k5&moves=3}\\{\sffamily\tiny\color{ink}scannen zum Lösen}
\end{minipage}
\vspace{2mm}

{\raggedright\noindent\textcolor{ink}{\sffamily\textsc{Themen}}\enspace \themetag{reines Matt} \themetag{Mustermatt} \themetag{Switchback}\par\smallskip{\footnotesize\textit{Hinweis:} die Partei am Zug steht im Diagramm im Schach; die Lösung beginnt mit einem Schlag.}\par}

\clearpage
\phantomsection\label{hm:KQQBvkn}
\begin{tcolorbox}[enhanced,colback=ink,colframe=ink,arc=0pt,boxrule=0pt,left=4mm,right=4mm,top=2.5mm,bottom=2.5mm]
\color{white}\sffamily{\LARGE\bfseries KQQBvkn}\hfill{\Large h\#3}\\[0.5mm]
{\small 6 Steine\quad Schwarz am Zug}\hfill{\small Nr.~841\qquad\href{https://helpman.komtera.lt/?fen=8/8/8/8/8/2Q2Q2/8/K1k1n2B&moves=3}{\color{white}\bfseries lösen\,\ensuremath{\nearrow}}}
\end{tcolorbox}
\vspace{2mm}
\noindent\begin{minipage}[t]{36mm}\vspace{0pt}\sffamily\small\raggedright
\hmfigure{längste eindeutige}{{\Large h\#3}}\hmfigure{eindeutig bei dieser Länge}{{\Large 30\,769} Stellungen}\hmfigure{längstes Matt}{{\Large h\#5}}\hmfigure{in der Klasse}{15,5 Milliarden Stellungen,\\ 3,92 Milliarden mit Hilfsmatt}
\end{minipage}\hfill
\begin{minipage}[t]{80mm}\vspace{0pt}\begin{tikzpicture}\begin{axis}[width=72mm,height=42mm,ybar,bar width=1.6pt,xmin=-1,xmax=11,ymin=0,ymax=1.497,axis y line*=left,axis x line*=bottom,xlabel={\scriptsize Halbzüge bis Matt},ylabel={\scriptsize Stellungen (Mrd.)},tick label style={font=\tiny},label style={font=\scriptsize},enlarge x limits=false,ylabel shift=-3pt,xlabel shift=-3pt]\addplot[fill=ink!55,draw=none] coordinates {(0,0.1235) (1,0.3085) (2,1.361) (3,0.8028) (4,1.243) (5,0.03674) (6,0.04109) (7,5.19e-07) (8,0.0001145) (10,1.6e-08)};\draw[densely dotted] (axis cs:6,0) -- (axis cs:6,1.497);\end{axis}\begin{axis}[width=72mm,height=42mm,axis y line*=right,axis x line=none,ymode=log,xmin=-1,xmax=11,ylabel={\scriptsize eindeutig (log)},tick label style={font=\tiny},label style={font=\scriptsize},ylabel shift=-3pt]\addplot[mark=*,mark size=0.7pt,draw=black,thin] coordinates {(0,123478984) (1,148427996) (2,376626774) (3,9244162) (4,16978288) (5,924) (6,30769)};\end{axis}\end{tikzpicture}\end{minipage}
\vspace{3mm}

\noindent\begin{minipage}[t]{60mm}\vspace{0pt}\sbox\hmboard{\chessboard[setfen=8/8/8/8/8/2Q2Q2/8/K1k1n2B b - - 0 1]}\usebox\hmboard\\[1mm]
\makebox[\wd\hmboard]{\sffamily\hspace{3mm}\textbf{h\#3}\hfill(4+2)}\end{minipage}\hfill
\begin{minipage}[t]{52mm}\vspace{0pt}
\textcolor{ink}{\sffamily\textsc{Lösung}}\par\smallskip\small
\begin{tabular}{@{}r@{\enspace}l@{\enspace}l@{}}
\textbf{1.} & \symknight{}c2+ & \symqueen{}xc2+\\
\textbf{2.} & \symking{}xc2 & \symqueen{}c3+\\
\textbf{3.} & \symking{}d1 & \symbishop{}f3\#\\
\end{tabular}\par\vspace{3mm}
\qrcode[height=12mm]{https://helpman.komtera.lt/?fen=8/8/8/8/8/2Q2Q2/8/K1k1n2B&moves=3}\\{\sffamily\tiny\color{ink}scannen zum Lösen}
\end{minipage}
\vspace{2mm}

{\raggedright\noindent\textcolor{ink}{\sffamily\textsc{Themen}}\enspace \themetag{reines Matt} \themetag{Mustermatt} \themetag{Spiegelmatt}\par\smallskip{\footnotesize\textit{Hinweis:} die Partei am Zug steht im Diagramm im Schach.}\par}

\clearpage
\phantomsection\label{hm:KQQNNvk}
\begin{tcolorbox}[enhanced,colback=ink,colframe=ink,arc=0pt,boxrule=0pt,left=4mm,right=4mm,top=2.5mm,bottom=2.5mm]
\color{white}\sffamily{\LARGE\bfseries KQQNNvk}\hfill{\Large h\#3}\\[0.5mm]
{\small 6 Steine\quad Schwarz am Zug}\hfill{\small Nr.~842\qquad\href{https://helpman.komtera.lt/?fen=1N6/N7/8/8/8/8/3Q4/KQk5&moves=3}{\color{white}\bfseries lösen\,\ensuremath{\nearrow}}}
\end{tcolorbox}
\vspace{2mm}
\noindent\begin{minipage}[t]{36mm}\vspace{0pt}\sffamily\small\raggedright
\hmfigure{längste eindeutige}{{\Large h\#3}}\hmfigure{eindeutig bei dieser Länge}{{\Large 76\,571} Stellungen}\hmfigure{längstes Matt}{{\Large h\#4}}\hmfigure{in der Klasse}{15,5 Milliarden Stellungen,\\ 2,04 Milliarden mit Hilfsmatt}
\end{minipage}\hfill
\begin{minipage}[t]{80mm}\vspace{0pt}\begin{tikzpicture}\begin{axis}[width=72mm,height=42mm,ybar,bar width=1.6pt,xmin=-1,xmax=9,ymin=0,ymax=794.4,axis y line*=left,axis x line*=bottom,xlabel={\scriptsize Halbzüge bis Matt},ylabel={\scriptsize Stellungen (Mio.)},tick label style={font=\tiny},label style={font=\scriptsize},enlarge x limits=false,ylabel shift=-3pt,xlabel shift=-3pt]\addplot[fill=ink!55,draw=none] coordinates {(0,76.81) (1,194.2) (2,722.1) (3,371.3) (4,646.8) (5,4.296) (6,24.19) (8,0.6443)};\draw[densely dotted] (axis cs:6,0) -- (axis cs:6,794.4);\end{axis}\begin{axis}[width=72mm,height=42mm,axis y line*=right,axis x line=none,ymode=log,xmin=-1,xmax=9,ylabel={\scriptsize eindeutig (log)},tick label style={font=\tiny},label style={font=\scriptsize},ylabel shift=-3pt]\addplot[mark=*,mark size=0.7pt,draw=black,thin] coordinates {(0,76809581) (1,89757945) (2,239870144) (3,3525815) (4,6042081) (6,76571)};\end{axis}\end{tikzpicture}\end{minipage}
\vspace{3mm}

\noindent\begin{minipage}[t]{60mm}\vspace{0pt}\sbox\hmboard{\chessboard[setfen=1N6/N7/8/8/8/8/3Q4/KQk5 b - - 0 1]}\usebox\hmboard\\[1mm]
\makebox[\wd\hmboard]{\sffamily\hspace{3mm}\textbf{h\#3}\hfill(5+1)}\end{minipage}\hfill
\begin{minipage}[t]{52mm}\vspace{0pt}
\textcolor{ink}{\sffamily\textsc{Lösung}}\par\smallskip\small
\begin{tabular}{@{}r@{\enspace}l@{\enspace}l@{}}
\textbf{1.} & \symking{}xd2 & \symking{}a2\\
\textbf{2.} & \symking{}c3 & \symqueen{}c2+\\
\textbf{3.} & \symking{}b4 & \symknight{}bc6\#\\
\end{tabular}\par\vspace{3mm}
\qrcode[height=12mm]{https://helpman.komtera.lt/?fen=1N6/N7/8/8/8/8/3Q4/KQk5&moves=3}\\{\sffamily\tiny\color{ink}scannen zum Lösen}
\end{minipage}
\vspace{2mm}

{\raggedright\noindent\textcolor{ink}{\sffamily\textsc{Themen}}\enspace \themetag{Spiegelmatt} \themetag{ein Stein zieht}\par\smallskip{\footnotesize\textit{Hinweis:} das Diagramm hat keinen legalen letzten Zug, kann also in keiner Partie entstehen.}\par}

\clearpage
\phantomsection\label{hm:KQQQvkb}
\begin{tcolorbox}[enhanced,colback=ink,colframe=ink,arc=0pt,boxrule=0pt,left=4mm,right=4mm,top=2.5mm,bottom=2.5mm]
\color{white}\sffamily{\LARGE\bfseries KQQQvkb}\hfill{\Large h\#3}\\[0.5mm]
{\small 6 Steine\quad Schwarz am Zug}\hfill{\small Nr.~843\qquad\href{https://helpman.komtera.lt/?fen=8/8/8/8/3K4/8/Q6b/Q5Qk&moves=3}{\color{white}\bfseries lösen\,\ensuremath{\nearrow}}}
\end{tcolorbox}
\vspace{2mm}
\noindent\begin{minipage}[t]{36mm}\vspace{0pt}\sffamily\small\raggedright
\hmfigure{längste eindeutige}{{\Large h\#3}}\hmfigure{eindeutig bei dieser Länge}{{\Large 9} Stellungen}\hmfigure{längstes Matt}{{\Large h\#5}}\hmfigure{in der Klasse}{15,5 Milliarden Stellungen,\\ 1,15 Milliarden mit Hilfsmatt}
\end{minipage}\hfill
\begin{minipage}[t]{80mm}\vspace{0pt}\begin{tikzpicture}\begin{axis}[width=72mm,height=42mm,ybar,bar width=1.6pt,xmin=-1,xmax=11,ymin=0,ymax=622.1,axis y line*=left,axis x line*=bottom,xlabel={\scriptsize Halbzüge bis Matt},ylabel={\scriptsize Stellungen (Mio.)},tick label style={font=\tiny},label style={font=\scriptsize},enlarge x limits=false,ylabel shift=-3pt,xlabel shift=-3pt]\addplot[fill=ink!55,draw=none] coordinates {(0,91.05) (1,125) (2,565.6) (3,151.9) (4,215) (5,3.699) (6,0.2348) (7,0.000102) (8,5.7e-05) (10,1e-06)};\draw[densely dotted] (axis cs:6,0) -- (axis cs:6,622.1);\end{axis}\begin{axis}[width=72mm,height=42mm,axis y line*=right,axis x line=none,ymode=log,xmin=-1,xmax=11,ylabel={\scriptsize eindeutig (log)},tick label style={font=\tiny},label style={font=\scriptsize},ylabel shift=-3pt]\addplot[mark=*,mark size=0.7pt,draw=black,thin] coordinates {(0,91048310) (1,33492670) (2,75346808) (3,851123) (4,163301) (5,17) (6,9)};\end{axis}\end{tikzpicture}\end{minipage}
\vspace{3mm}

\noindent\begin{minipage}[t]{60mm}\vspace{0pt}\sbox\hmboard{\chessboard[setfen=8/8/8/8/3K4/8/Q6b/Q5Qk b - - 0 1]}\usebox\hmboard\\[1mm]
\makebox[\wd\hmboard]{\sffamily\hspace{3mm}\textbf{h\#3}\hfill(4+2)}\end{minipage}\hfill
\begin{minipage}[t]{52mm}\vspace{0pt}
\textcolor{ink}{\sffamily\textsc{Lösung}}\par\smallskip\small
\begin{tabular}{@{}r@{\enspace}l@{\enspace}l@{}}
\textbf{1.} & \symbishop{}xg1+ & \symqueen{}xg1+\\
\textbf{2.} & \symking{}xg1 & \symking{}e3\\
\textbf{3.} & \symking{}f1 & \symqueen{}f2\#\\
\end{tabular}\par\vspace{3mm}
\qrcode[height=12mm]{https://helpman.komtera.lt/?fen=8/8/8/8/3K4/8/Q6b/Q5Qk&moves=3}\\{\sffamily\tiny\color{ink}scannen zum Lösen}
\end{minipage}
\vspace{2mm}

{\raggedright\noindent\textcolor{ink}{\sffamily\textsc{Themen}}\enspace \textit{keine erkannt}\par\smallskip{\footnotesize\textit{Hinweis:} die Partei am Zug steht im Diagramm im Schach; die Lösung beginnt mit einem Schlag.}\par}

\clearpage
\phantomsection\label{hm:KQQRNvk}
\begin{tcolorbox}[enhanced,colback=ink,colframe=ink,arc=0pt,boxrule=0pt,left=4mm,right=4mm,top=2.5mm,bottom=2.5mm]
\color{white}\sffamily{\LARGE\bfseries KQQRNvk}\hfill{\Large h\#3}\\[0.5mm]
{\small 6 Steine\quad Schwarz am Zug}\hfill{\small Nr.~844\qquad\href{https://helpman.komtera.lt/?fen=N7/8/8/3R4/4Q3/1Q1k4/8/K7&moves=3}{\color{white}\bfseries lösen\,\ensuremath{\nearrow}}}
\end{tcolorbox}
\vspace{2mm}
\noindent\begin{minipage}[t]{36mm}\vspace{0pt}\sffamily\small\raggedright
\hmfigure{längste eindeutige}{{\Large h\#3}}\hmfigure{eindeutig bei dieser Länge}{{\Large 71} Stellungen}\hmfigure{längstes Matt}{{\Large h\#4}}\hmfigure{in der Klasse}{15,5 Milliarden Stellungen,\\ 3,87 Milliarden mit Hilfsmatt}
\end{minipage}\hfill
\begin{minipage}[t]{80mm}\vspace{0pt}\begin{tikzpicture}\begin{axis}[width=72mm,height=42mm,ybar,bar width=1.6pt,xmin=-1,xmax=9,ymin=0,ymax=1.757,axis y line*=left,axis x line*=bottom,xlabel={\scriptsize Halbzüge bis Matt},ylabel={\scriptsize Stellungen (Mrd.)},tick label style={font=\tiny},label style={font=\scriptsize},enlarge x limits=false,ylabel shift=-3pt,xlabel shift=-3pt]\addplot[fill=ink!55,draw=none] coordinates {(0,0.2771) (1,0.4322) (2,1.597) (3,0.5182) (4,1.038) (5,0.001136) (6,0.008967) (8,1.69e-07)};\draw[densely dotted] (axis cs:6,0) -- (axis cs:6,1.757);\end{axis}\begin{axis}[width=72mm,height=42mm,axis y line*=right,axis x line=none,ymode=log,xmin=-1,xmax=9,ylabel={\scriptsize eindeutig (log)},tick label style={font=\tiny},label style={font=\scriptsize},ylabel shift=-3pt]\addplot[mark=*,mark size=0.7pt,draw=black,thin] coordinates {(0,277147559) (1,131089263) (2,340349999) (3,1126051) (4,7292861) (6,71)};\end{axis}\end{tikzpicture}\end{minipage}
\vspace{3mm}

\noindent\begin{minipage}[t]{60mm}\vspace{0pt}\sbox\hmboard{\chessboard[setfen=N7/8/8/3R4/4Q3/1Q1k4/8/K7 b - - 0 1]}\usebox\hmboard\\[1mm]
\makebox[\wd\hmboard]{\sffamily\hspace{3mm}\textbf{h\#3}\hfill(5+1)}\end{minipage}\hfill
\begin{minipage}[t]{52mm}\vspace{0pt}
\textcolor{ink}{\sffamily\textsc{Lösung}}\par\smallskip\small
\begin{tabular}{@{}r@{\enspace}l@{\enspace}l@{}}
\textbf{1.} & \symking{}xe4 & \symknight{}c7\\
\textbf{2.} & \symking{}f4 & \symknight{}e6+\\
\textbf{3.} & \symking{}e4 & \symqueen{}d3\#\\
\end{tabular}\par\vspace{3mm}
\qrcode[height=12mm]{https://helpman.komtera.lt/?fen=N7/8/8/3R4/4Q3/1Q1k4/8/K7&moves=3}\\{\sffamily\tiny\color{ink}scannen zum Lösen}
\end{minipage}
\vspace{2mm}

{\raggedright\noindent\textcolor{ink}{\sffamily\textsc{Themen}}\enspace \themetag{Switchback} \themetag{ein Stein zieht} \themetag{Kniest}\par\smallskip{\footnotesize\textit{Hinweis:} das Diagramm hat keinen legalen letzten Zug, kann also in keiner Partie entstehen.}\par}

\clearpage
\phantomsection\label{hm:KQQRPvk}
\begin{tcolorbox}[enhanced,colback=ink,colframe=ink,arc=0pt,boxrule=0pt,left=4mm,right=4mm,top=2.5mm,bottom=2.5mm]
\color{white}\sffamily{\LARGE\bfseries KQQRPvk}\hfill{\Large h\#3}\\[0.5mm]
{\small 6 Steine\quad Schwarz am Zug}\hfill{\small Nr.~845\qquad\href{https://helpman.komtera.lt/?fen=8/8/8/P7/8/4Q3/1RQk4/K7&moves=3}{\color{white}\bfseries lösen\,\ensuremath{\nearrow}}}
\end{tcolorbox}
\vspace{2mm}
\noindent\begin{minipage}[t]{36mm}\vspace{0pt}\sffamily\small\raggedright
\hmfigure{längste eindeutige}{{\Large h\#3}}\hmfigure{eindeutig bei dieser Länge}{{\Large 16\,073} Stellungen}\hmfigure{längstes Matt}{{\Large h\#4}}\hmfigure{in der Klasse}{45,4 Milliarden Stellungen,\\ 11,8 Milliarden mit Hilfsmatt}
\end{minipage}\hfill
\begin{minipage}[t]{80mm}\vspace{0pt}\begin{tikzpicture}\begin{axis}[width=72mm,height=42mm,ybar,bar width=1.6pt,xmin=-1,xmax=9,ymin=0,ymax=4.964,axis y line*=left,axis x line*=bottom,xlabel={\scriptsize Halbzüge bis Matt},ylabel={\scriptsize Stellungen (Mrd.)},tick label style={font=\tiny},label style={font=\scriptsize},enlarge x limits=false,ylabel shift=-3pt,xlabel shift=-3pt]\addplot[fill=ink!55,draw=none] coordinates {(0,0.746) (1,1.23) (2,4.513) (3,1.783) (4,3.492) (5,0.007595) (6,0.04992) (8,4.509e-05)};\draw[densely dotted] (axis cs:6,0) -- (axis cs:6,4.964);\end{axis}\begin{axis}[width=72mm,height=42mm,axis y line*=right,axis x line=none,ymode=log,xmin=-1,xmax=9,ylabel={\scriptsize eindeutig (log)},tick label style={font=\tiny},label style={font=\scriptsize},ylabel shift=-3pt]\addplot[mark=*,mark size=0.7pt,draw=black,thin] coordinates {(0,745980123) (1,379075433) (2,1011813933) (3,7429934) (4,26318581) (6,16073)};\end{axis}\end{tikzpicture}\end{minipage}
\vspace{3mm}

\noindent\begin{minipage}[t]{60mm}\vspace{0pt}\sbox\hmboard{\chessboard[setfen=8/8/8/P7/8/4Q3/1RQk4/K7 b - - 0 1]}\usebox\hmboard\\[1mm]
\makebox[\wd\hmboard]{\sffamily\hspace{3mm}\textbf{h\#3}\hfill(5+1)}\end{minipage}\hfill
\begin{minipage}[t]{52mm}\vspace{0pt}
\textcolor{ink}{\sffamily\textsc{Lösung}}\par\smallskip\small
\begin{tabular}{@{}r@{\enspace}l@{\enspace}l@{}}
\textbf{1.} & \symking{}xe3 & \symking{}a2\\
\textbf{2.} & \symking{}d4 & \symqueen{}e4+\\
\textbf{3.} & \symking{}c3 & \symrook{}c2\#\\
\end{tabular}\par\vspace{3mm}
\qrcode[height=12mm]{https://helpman.komtera.lt/?fen=8/8/8/P7/8/4Q3/1RQk4/K7&moves=3}\\{\sffamily\tiny\color{ink}scannen zum Lösen}
\end{minipage}
\vspace{2mm}

{\raggedright\noindent\textcolor{ink}{\sffamily\textsc{Themen}}\enspace \themetag{ein Stein zieht}\par\smallskip{\footnotesize\textit{Hinweis:} das Diagramm hat keinen legalen letzten Zug, kann also in keiner Partie entstehen.}\par}

\clearpage
\phantomsection\label{hm:KQQRvkp}
\begin{tcolorbox}[enhanced,colback=ink,colframe=ink,arc=0pt,boxrule=0pt,left=4mm,right=4mm,top=2.5mm,bottom=2.5mm]
\color{white}\sffamily{\LARGE\bfseries KQQRvkp}\hfill{\Large h\#3}\\[0.5mm]
{\small 6 Steine\quad Schwarz am Zug}\hfill{\small Nr.~846\qquad\href{https://helpman.komtera.lt/?fen=8/8/8/1p6/8/4Q3/R1Qk4/K7&moves=3}{\color{white}\bfseries lösen\,\ensuremath{\nearrow}}}
\end{tcolorbox}
\vspace{2mm}
\noindent\begin{minipage}[t]{36mm}\vspace{0pt}\sffamily\small\raggedright
\hmfigure{längste eindeutige}{{\Large h\#3}}\hmfigure{eindeutig bei dieser Länge}{{\Large 48\,719} Stellungen}\hmfigure{längstes Matt}{{\Large h\#6}}\hmfigure{in der Klasse}{45,4 Milliarden Stellungen,\\ 11,9 Milliarden mit Hilfsmatt}
\end{minipage}\hfill
\begin{minipage}[t]{80mm}\vspace{0pt}\begin{tikzpicture}\begin{axis}[width=72mm,height=42mm,ybar,bar width=1.6pt,xmin=-1,xmax=13,ymin=0,ymax=5.204,axis y line*=left,axis x line*=bottom,xlabel={\scriptsize Halbzüge bis Matt},ylabel={\scriptsize Stellungen (Mrd.)},tick label style={font=\tiny},label style={font=\scriptsize},enlarge x limits=false,ylabel shift=-3pt,xlabel shift=-3pt]\addplot[fill=ink!55,draw=none] coordinates {(0,0.6735) (1,1.199) (2,4.731) (3,1.885) (4,3.313) (5,0.0282) (6,0.05284) (7,1.01e-07) (8,0.0001465) (10,2e-09) (12,1e-09)};\draw[densely dotted] (axis cs:6,0) -- (axis cs:6,5.204);\end{axis}\begin{axis}[width=72mm,height=42mm,axis y line*=right,axis x line=none,ymode=log,xmin=-1,xmax=13,ylabel={\scriptsize eindeutig (log)},tick label style={font=\tiny},label style={font=\scriptsize},ylabel shift=-3pt]\addplot[mark=*,mark size=0.7pt,draw=black,thin] coordinates {(0,673452182) (1,386016205) (2,929233054) (3,6508987) (4,20645490) (5,9) (6,48719)};\end{axis}\end{tikzpicture}\end{minipage}
\vspace{3mm}

\noindent\begin{minipage}[t]{60mm}\vspace{0pt}\sbox\hmboard{\chessboard[setfen=8/8/8/1p6/8/4Q3/R1Qk4/K7 b - - 0 1]}\usebox\hmboard\\[1mm]
\makebox[\wd\hmboard]{\sffamily\hspace{3mm}\textbf{h\#3}\hfill(4+2)}\end{minipage}\hfill
\begin{minipage}[t]{52mm}\vspace{0pt}
\textcolor{ink}{\sffamily\textsc{Lösung}}\par\smallskip\small
\begin{tabular}{@{}r@{\enspace}l@{\enspace}l@{}}
\textbf{1.} & \symking{}xe3 & \symrook{}a6\\
\textbf{2.} & \symking{}d4 & \symqueen{}e4+\\
\textbf{3.} & \symking{}c5 & \symrook{}c6\#\\
\end{tabular}\par\vspace{3mm}
\qrcode[height=12mm]{https://helpman.komtera.lt/?fen=8/8/8/1p6/8/4Q3/R1Qk4/K7&moves=3}\\{\sffamily\tiny\color{ink}scannen zum Lösen}
\end{minipage}
\vspace{2mm}

{\raggedright\noindent\textcolor{ink}{\sffamily\textsc{Themen}}\enspace \themetag{ein Stein zieht}\par\smallskip{\footnotesize\textit{Hinweis:} das Diagramm hat keinen legalen letzten Zug, kann also in keiner Partie entstehen.}\par}

\clearpage
\phantomsection\label{hm:KQQRvkr}
\begin{tcolorbox}[enhanced,colback=ink,colframe=ink,arc=0pt,boxrule=0pt,left=4mm,right=4mm,top=2.5mm,bottom=2.5mm]
\color{white}\sffamily{\LARGE\bfseries KQQRvkr}\hfill{\Large h\#3}\\[0.5mm]
{\small 6 Steine\quad Schwarz am Zug}\hfill{\small Nr.~847\qquad\href{https://helpman.komtera.lt/?fen=8/8/8/R7/8/1r6/Q7/KQk5&moves=3}{\color{white}\bfseries lösen\,\ensuremath{\nearrow}}}
\end{tcolorbox}
\vspace{2mm}
\noindent\begin{minipage}[t]{36mm}\vspace{0pt}\sffamily\small\raggedright
\hmfigure{längste eindeutige}{{\Large h\#3}}\hmfigure{eindeutig bei dieser Länge}{{\Large 217} Stellungen}\hmfigure{längstes Matt}{{\Large h\#5}}\hmfigure{in der Klasse}{15,5 Milliarden Stellungen,\\ 3,45 Milliarden mit Hilfsmatt}
\end{minipage}\hfill
\begin{minipage}[t]{80mm}\vspace{0pt}\begin{tikzpicture}\begin{axis}[width=72mm,height=42mm,ybar,bar width=1.6pt,xmin=-1,xmax=11,ymin=0,ymax=1.468,axis y line*=left,axis x line*=bottom,xlabel={\scriptsize Halbzüge bis Matt},ylabel={\scriptsize Stellungen (Mrd.)},tick label style={font=\tiny},label style={font=\scriptsize},enlarge x limits=false,ylabel shift=-3pt,xlabel shift=-3pt]\addplot[fill=ink!55,draw=none] coordinates {(0,0.1501) (1,0.2971) (2,1.334) (3,0.6819) (4,0.9199) (5,0.04888) (6,0.01386) (7,1.29e-07) (8,3.85e-07) (9,1e-09) (10,3.9e-08)};\draw[densely dotted] (axis cs:6,0) -- (axis cs:6,1.468);\end{axis}\begin{axis}[width=72mm,height=42mm,axis y line*=right,axis x line=none,ymode=log,xmin=-1,xmax=11,ylabel={\scriptsize eindeutig (log)},tick label style={font=\tiny},label style={font=\scriptsize},ylabel shift=-3pt]\addplot[mark=*,mark size=0.7pt,draw=black,thin] coordinates {(0,150066047) (1,120027833) (2,284153305) (3,8850682) (4,6886768) (5,68) (6,217)};\end{axis}\end{tikzpicture}\end{minipage}
\vspace{3mm}

\noindent\begin{minipage}[t]{60mm}\vspace{0pt}\sbox\hmboard{\chessboard[setfen=8/8/8/R7/8/1r6/Q7/KQk5 b - - 0 1]}\usebox\hmboard\\[1mm]
\makebox[\wd\hmboard]{\sffamily\hspace{3mm}\textbf{h\#3}\hfill(4+2)}\end{minipage}\hfill
\begin{minipage}[t]{52mm}\vspace{0pt}
\textcolor{ink}{\sffamily\textsc{Lösung}}\par\smallskip\small
\begin{tabular}{@{}r@{\enspace}l@{\enspace}l@{}}
\textbf{1.} & \symrook{}xb1+ & \symqueen{}xb1+\\
\textbf{2.} & \symking{}d2 & \symqueen{}d3+\\
\textbf{3.} & \symking{}c1 & \symrook{}c5\#\\
\end{tabular}\par\vspace{3mm}
\qrcode[height=12mm]{https://helpman.komtera.lt/?fen=8/8/8/R7/8/1r6/Q7/KQk5&moves=3}\\{\sffamily\tiny\color{ink}scannen zum Lösen}
\end{minipage}
\vspace{2mm}

{\raggedright\noindent\textcolor{ink}{\sffamily\textsc{Themen}}\enspace \themetag{Spiegelmatt} \themetag{Switchback}\par\smallskip{\footnotesize\textit{Hinweis:} die Partei am Zug steht im Diagramm im Schach; die Lösung beginnt mit einem Schlag.}\par}

\clearpage
\phantomsection\label{hm:KRRRBvk}
\begin{tcolorbox}[enhanced,colback=ink,colframe=ink,arc=0pt,boxrule=0pt,left=4mm,right=4mm,top=2.5mm,bottom=2.5mm]
\color{white}\sffamily{\LARGE\bfseries KRRRBvk}\hfill{\Large h\#3}\\[0.5mm]
{\small 6 Steine\quad Schwarz am Zug}\hfill{\small Nr.~848\qquad\href{https://helpman.komtera.lt/?fen=8/8/8/8/8/4R3/1B1R4/K1k1R3&moves=3}{\color{white}\bfseries lösen\,\ensuremath{\nearrow}}}
\end{tcolorbox}
\vspace{2mm}
\noindent\begin{minipage}[t]{36mm}\vspace{0pt}\sffamily\small\raggedright
\hmfigure{längste eindeutige}{{\Large h\#3}}\hmfigure{eindeutig bei dieser Länge}{{\Large 40\,594} Stellungen}\hmfigure{längstes Matt}{{\Large h\#5}}\hmfigure{in der Klasse}{15,5 Milliarden Stellungen,\\ 1,42 Milliarden mit Hilfsmatt}
\end{minipage}\hfill
\begin{minipage}[t]{80mm}\vspace{0pt}\begin{tikzpicture}\begin{axis}[width=72mm,height=42mm,ybar,bar width=1.6pt,xmin=-1,xmax=11,ymin=0,ymax=705.4,axis y line*=left,axis x line*=bottom,xlabel={\scriptsize Halbzüge bis Matt},ylabel={\scriptsize Stellungen (Mio.)},tick label style={font=\tiny},label style={font=\scriptsize},enlarge x limits=false,ylabel shift=-3pt,xlabel shift=-3pt]\addplot[fill=ink!55,draw=none] coordinates {(0,38.98) (1,92) (2,295.4) (3,316.7) (4,641.3) (5,23.23) (6,15.64) (7,0.000239) (8,0.1761) (10,1.4e-05)};\draw[densely dotted] (axis cs:6,0) -- (axis cs:6,705.4);\end{axis}\begin{axis}[width=72mm,height=42mm,axis y line*=right,axis x line=none,ymode=log,xmin=-1,xmax=11,ylabel={\scriptsize eindeutig (log)},tick label style={font=\tiny},label style={font=\scriptsize},ylabel shift=-3pt]\addplot[mark=*,mark size=0.7pt,draw=black,thin] coordinates {(0,38982361) (1,54106721) (2,120889955) (3,11182865) (4,8041043) (5,78) (6,40594)};\end{axis}\end{tikzpicture}\end{minipage}
\vspace{3mm}

\noindent\begin{minipage}[t]{60mm}\vspace{0pt}\sbox\hmboard{\chessboard[setfen=8/8/8/8/8/4R3/1B1R4/K1k1R3 b - - 0 1]}\usebox\hmboard\\[1mm]
\makebox[\wd\hmboard]{\sffamily\hspace{3mm}\textbf{h\#3}\hfill(5+1)}\end{minipage}\hfill
\begin{minipage}[t]{52mm}\vspace{0pt}
\textcolor{ink}{\sffamily\textsc{Lösung}}\par\smallskip\small
\begin{tabular}{@{}r@{\enspace}l@{\enspace}l@{}}
\textbf{1.} & \symking{}xd2 & \symrook{}e4\\
\textbf{2.} & \symking{}c2 & \symrook{}c1+\\
\textbf{3.} & \symking{}b3 & \symrook{}c3\#\\
\end{tabular}\par\vspace{3mm}
\qrcode[height=12mm]{https://helpman.komtera.lt/?fen=8/8/8/8/8/4R3/1B1R4/K1k1R3&moves=3}\\{\sffamily\tiny\color{ink}scannen zum Lösen}
\end{minipage}
\vspace{2mm}

{\raggedright\noindent\textcolor{ink}{\sffamily\textsc{Themen}}\enspace \themetag{ein Stein zieht}\par\smallskip{\footnotesize\textit{Hinweis:} das Diagramm hat keinen legalen letzten Zug, kann also in keiner Partie entstehen.}\par}

\clearpage
\phantomsection\label{hm:KRRRPvk}
\begin{tcolorbox}[enhanced,colback=ink,colframe=ink,arc=0pt,boxrule=0pt,left=4mm,right=4mm,top=2.5mm,bottom=2.5mm]
\color{white}\sffamily{\LARGE\bfseries KRRRPvk}\hfill{\Large h\#3}\\[0.5mm]
{\small 6 Steine\quad Schwarz am Zug}\hfill{\small Nr.~849\qquad\href{https://helpman.komtera.lt/?fen=8/8/8/8/8/2R5/3R1P2/K1k2R2&moves=3}{\color{white}\bfseries lösen\,\ensuremath{\nearrow}}}
\end{tcolorbox}
\vspace{2mm}
\noindent\begin{minipage}[t]{36mm}\vspace{0pt}\sffamily\small\raggedright
\hmfigure{längste eindeutige}{{\Large h\#3}}\hmfigure{eindeutig bei dieser Länge}{{\Large 54\,154} Stellungen}\hmfigure{längstes Matt}{{\Large h\#7}}\hmfigure{in der Klasse}{45,4 Milliarden Stellungen,\\ 4,44 Milliarden mit Hilfsmatt}
\end{minipage}\hfill
\begin{minipage}[t]{80mm}\vspace{0pt}\begin{tikzpicture}\begin{axis}[width=72mm,height=42mm,ybar,bar width=1.6pt,xmin=-1,xmax=15,ymin=0,ymax=2.178,axis y line*=left,axis x line*=bottom,xlabel={\scriptsize Halbzüge bis Matt},ylabel={\scriptsize Stellungen (Mrd.)},tick label style={font=\tiny},label style={font=\scriptsize},enlarge x limits=false,ylabel shift=-3pt,xlabel shift=-3pt]\addplot[fill=ink!55,draw=none] coordinates {(0,0.1001) (1,0.2763) (2,0.8598) (3,1.07) (4,1.98) (5,0.09441) (6,0.05472) (7,1.033e-05) (8,0.0004924) (10,1.96e-07) (12,9e-09) (14,2e-09)};\draw[densely dotted] (axis cs:6,0) -- (axis cs:6,2.178);\end{axis}\begin{axis}[width=72mm,height=42mm,axis y line*=right,axis x line=none,ymode=log,xmin=-1,xmax=15,ylabel={\scriptsize eindeutig (log)},tick label style={font=\tiny},label style={font=\scriptsize},ylabel shift=-3pt]\addplot[mark=*,mark size=0.7pt,draw=black,thin] coordinates {(0,100126004) (1,172602414) (2,369284563) (3,40898042) (4,24369356) (5,1639) (6,54154)};\end{axis}\end{tikzpicture}\end{minipage}
\vspace{3mm}

\noindent\begin{minipage}[t]{60mm}\vspace{0pt}\sbox\hmboard{\chessboard[setfen=8/8/8/8/8/2R5/3R1P2/K1k2R2 b - - 0 1]}\usebox\hmboard\\[1mm]
\makebox[\wd\hmboard]{\sffamily\hspace{3mm}\textbf{h\#3}\hfill(5+1)}\end{minipage}\hfill
\begin{minipage}[t]{52mm}\vspace{0pt}
\textcolor{ink}{\sffamily\textsc{Lösung}}\par\smallskip\small
\begin{tabular}{@{}r@{\enspace}l@{\enspace}l@{}}
\textbf{1.} & \symking{}xd2 & \symrook{}b1\\
\textbf{2.} & \symking{}e2 & \symrook{}b2+\\
\textbf{3.} & \symking{}e1 & \symrook{}c1\#\\
\end{tabular}\par\vspace{3mm}
\qrcode[height=12mm]{https://helpman.komtera.lt/?fen=8/8/8/8/8/2R5/3R1P2/K1k2R2&moves=3}\\{\sffamily\tiny\color{ink}scannen zum Lösen}
\end{minipage}
\vspace{2mm}

{\raggedright\noindent\textcolor{ink}{\sffamily\textsc{Themen}}\enspace \themetag{reines Matt} \themetag{Mustermatt} \themetag{ein Stein zieht}\par\smallskip{\footnotesize\textit{Hinweis:} das Diagramm hat keinen legalen letzten Zug, kann also in keiner Partie entstehen.}\par}

\clearpage
\phantomsection\label{hm:KQQQvkp}
\begin{tcolorbox}[enhanced,colback=ink,colframe=ink,arc=0pt,boxrule=0pt,left=4mm,right=4mm,top=2.5mm,bottom=2.5mm]
\color{white}\sffamily{\LARGE\bfseries KQQQvkp}\hfill{\Large h\#2.5}\\[0.5mm]
{\small 6 Steine\quad Weiß am Zug}\hfill{\small Nr.~850\qquad\href{https://helpman.komtera.lt/?fen=8/8/8/8/Qp6/K5Q1/1Q6/7k&moves=2.5}{\color{white}\bfseries lösen\,\ensuremath{\nearrow}}}
\end{tcolorbox}
\vspace{2mm}
\noindent\begin{minipage}[t]{36mm}\vspace{0pt}\sffamily\small\raggedright
\hmfigure{längste eindeutige}{{\Large h\#2.5}}\hmfigure{eindeutig bei dieser Länge}{{\Large 2} Stellungen}\hmfigure{längstes Matt}{{\Large h\#3.5}}\hmfigure{in der Klasse}{45,4 Milliarden Stellungen,\\ 3,76 Milliarden mit Hilfsmatt}
\end{minipage}\hfill
\begin{minipage}[t]{80mm}\vspace{0pt}\begin{tikzpicture}\begin{axis}[width=72mm,height=42mm,ybar,bar width=1.6pt,xmin=-1,xmax=8,ymin=0,ymax=2.034,axis y line*=left,axis x line*=bottom,xlabel={\scriptsize Halbzüge bis Matt},ylabel={\scriptsize Stellungen (Mrd.)},tick label style={font=\tiny},label style={font=\scriptsize},enlarge x limits=false,ylabel shift=-3pt,xlabel shift=-3pt]\addplot[fill=ink!55,draw=none] coordinates {(0,0.3503) (1,0.4423) (2,1.849) (3,0.4059) (4,0.7101) (5,0.003588) (6,0.0005292) (7,2e-09)};\draw[densely dotted] (axis cs:5,0) -- (axis cs:5,2.034);\end{axis}\begin{axis}[width=72mm,height=42mm,axis y line*=right,axis x line=none,ymode=log,xmin=-1,xmax=8,ylabel={\scriptsize eindeutig (log)},tick label style={font=\tiny},label style={font=\scriptsize},ylabel shift=-3pt]\addplot[mark=*,mark size=0.7pt,draw=black,thin] coordinates {(0,350259654) (1,100016197) (2,240294533) (3,801155) (4,519786) (5,2)};\end{axis}\end{tikzpicture}\end{minipage}
\vspace{3mm}

\noindent\begin{minipage}[t]{60mm}\vspace{0pt}\sbox\hmboard{\chessboard[setfen=8/8/8/8/Qp6/K5Q1/1Q6/7k w - - 0 1]}\usebox\hmboard\\[1mm]
\makebox[\wd\hmboard]{\sffamily\hspace{3mm}\textbf{h\#2.5}\hfill(4+2)}\end{minipage}\hfill
\begin{minipage}[t]{52mm}\vspace{0pt}
\textcolor{ink}{\sffamily\textsc{Lösung}}\par\smallskip\small
\begin{tabular}{@{}r@{\enspace}l@{\enspace}l@{}}
\textbf{1.} & \ldots & \symking{}a2\\
\textbf{2.} & b3+ & \symqueen{}gxb3\\
\textbf{3.} & \symking{}g1 & \symqueen{}d1\#\\
\end{tabular}\par\vspace{3mm}
\qrcode[height=12mm]{https://helpman.komtera.lt/?fen=8/8/8/8/Qp6/K5Q1/1Q6/7k&moves=2.5}\\{\sffamily\tiny\color{ink}scannen zum Lösen}
\end{minipage}
\vspace{2mm}

{\raggedright\noindent\textcolor{ink}{\sffamily\textsc{Themen}}\enspace \themetag{reines Matt} \themetag{Spiegelmatt}\par\smallskip{\footnotesize\textit{Hinweis:} die Partei am Zug steht im Diagramm im Schach.}\par}

\clearpage
\phantomsection\label{hm:KQQQvkq}
\begin{tcolorbox}[enhanced,colback=ink,colframe=ink,arc=0pt,boxrule=0pt,left=4mm,right=4mm,top=2.5mm,bottom=2.5mm]
\color{white}\sffamily{\LARGE\bfseries KQQQvkq}\hfill{\Large h\#2.5}\\[0.5mm]
{\small 6 Steine\quad Weiß am Zug}\hfill{\small Nr.~851\qquad\href{https://helpman.komtera.lt/?fen=8/8/8/4Q3/2Q5/k7/2Q5/1K1q4&moves=2.5}{\color{white}\bfseries lösen\,\ensuremath{\nearrow}}}
\end{tcolorbox}
\vspace{2mm}
\noindent\begin{minipage}[t]{36mm}\vspace{0pt}\sffamily\small\raggedright
\hmfigure{längste eindeutige}{{\Large h\#2.5}}\hmfigure{eindeutig bei dieser Länge}{{\Large 68} Stellungen}\hmfigure{längstes Matt}{{\Large h\#5}}\hmfigure{in der Klasse}{15,5 Milliarden Stellungen,\\ 955 Millionen mit Hilfsmatt}
\end{minipage}\hfill
\begin{minipage}[t]{80mm}\vspace{0pt}\begin{tikzpicture}\begin{axis}[width=72mm,height=42mm,ybar,bar width=1.6pt,xmin=-1,xmax=11,ymin=0,ymax=474.2,axis y line*=left,axis x line*=bottom,xlabel={\scriptsize Halbzüge bis Matt},ylabel={\scriptsize Stellungen (Mio.)},tick label style={font=\tiny},label style={font=\scriptsize},enlarge x limits=false,ylabel shift=-3pt,xlabel shift=-3pt]\addplot[fill=ink!55,draw=none] coordinates {(0,58.49) (1,88.75) (2,431.1) (3,181.5) (4,184.5) (5,10.03) (6,0.5796) (7,0.000125) (8,1.9e-05) (10,2.5e-05)};\draw[densely dotted] (axis cs:5,0) -- (axis cs:5,474.2);\end{axis}\begin{axis}[width=72mm,height=42mm,axis y line*=right,axis x line=none,ymode=log,xmin=-1,xmax=11,ylabel={\scriptsize eindeutig (log)},tick label style={font=\tiny},label style={font=\scriptsize},ylabel shift=-3pt]\addplot[mark=*,mark size=0.7pt,draw=black,thin] coordinates {(0,58491751) (1,29939682) (2,70507266) (3,2627358) (4,307145) (5,68)};\end{axis}\end{tikzpicture}\end{minipage}
\vspace{3mm}

\noindent\begin{minipage}[t]{60mm}\vspace{0pt}\sbox\hmboard{\chessboard[setfen=8/8/8/4Q3/2Q5/k7/2Q5/1K1q4 w - - 0 1]}\usebox\hmboard\\[1mm]
\makebox[\wd\hmboard]{\sffamily\hspace{3mm}\textbf{h\#2.5}\hfill(4+2)}\end{minipage}\hfill
\begin{minipage}[t]{52mm}\vspace{0pt}
\textcolor{ink}{\sffamily\textsc{Lösung}}\par\smallskip\small
\begin{tabular}{@{}r@{\enspace}l@{\enspace}l@{}}
\textbf{1.} & \ldots & \symqueen{}c1+\\
\textbf{2.} & \symqueen{}xc1+ & \symqueen{}xc1+\\
\textbf{3.} & \symking{}b3 & \symqueen{}b5\#\\
\end{tabular}\par\vspace{3mm}
\qrcode[height=12mm]{https://helpman.komtera.lt/?fen=8/8/8/4Q3/2Q5/k7/2Q5/1K1q4&moves=2.5}\\{\sffamily\tiny\color{ink}scannen zum Lösen}
\end{minipage}
\vspace{2mm}

{\raggedright\noindent\textcolor{ink}{\sffamily\textsc{Themen}}\enspace \themetag{Spiegelmatt}\par\smallskip{\footnotesize\textit{Hinweis:} die Partei am Zug steht im Diagramm im Schach.}\par}

\clearpage
\phantomsection\label{hm:KQQQvkr}
\begin{tcolorbox}[enhanced,colback=ink,colframe=ink,arc=0pt,boxrule=0pt,left=4mm,right=4mm,top=2.5mm,bottom=2.5mm]
\color{white}\sffamily{\LARGE\bfseries KQQQvkr}\hfill{\Large h\#2.5}\\[0.5mm]
{\small 6 Steine\quad Weiß am Zug}\hfill{\small Nr.~852\qquad\href{https://helpman.komtera.lt/?fen=8/8/8/3Q4/8/r5Q1/2k5/K4Q2&moves=2.5}{\color{white}\bfseries lösen\,\ensuremath{\nearrow}}}
\end{tcolorbox}
\vspace{2mm}
\noindent\begin{minipage}[t]{36mm}\vspace{0pt}\sffamily\small\raggedright
\hmfigure{längste eindeutige}{{\Large h\#2.5}}\hmfigure{eindeutig bei dieser Länge}{{\Large 24} Stellungen}\hmfigure{längstes Matt}{{\Large h\#4}}\hmfigure{in der Klasse}{15,5 Milliarden Stellungen,\\ 1,08 Milliarden mit Hilfsmatt}
\end{minipage}\hfill
\begin{minipage}[t]{80mm}\vspace{0pt}\begin{tikzpicture}\begin{axis}[width=72mm,height=42mm,ybar,bar width=1.6pt,xmin=-1,xmax=9,ymin=0,ymax=578.4,axis y line*=left,axis x line*=bottom,xlabel={\scriptsize Halbzüge bis Matt},ylabel={\scriptsize Stellungen (Mio.)},tick label style={font=\tiny},label style={font=\scriptsize},enlarge x limits=false,ylabel shift=-3pt,xlabel shift=-3pt]\addplot[fill=ink!55,draw=none] coordinates {(0,79.91) (1,113.7) (2,525.8) (3,161.4) (4,196.6) (5,5.454) (6,0.285) (7,0.000103) (8,8.5e-05)};\draw[densely dotted] (axis cs:5,0) -- (axis cs:5,578.4);\end{axis}\begin{axis}[width=72mm,height=42mm,axis y line*=right,axis x line=none,ymode=log,xmin=-1,xmax=9,ylabel={\scriptsize eindeutig (log)},tick label style={font=\tiny},label style={font=\scriptsize},ylabel shift=-3pt]\addplot[mark=*,mark size=0.7pt,draw=black,thin] coordinates {(0,79909058) (1,33023639) (2,75205086) (3,1285234) (4,175927) (5,24)};\end{axis}\end{tikzpicture}\end{minipage}
\vspace{3mm}

\noindent\begin{minipage}[t]{60mm}\vspace{0pt}\sbox\hmboard{\chessboard[setfen=8/8/8/3Q4/8/r5Q1/2k5/K4Q2 w - - 0 1]}\usebox\hmboard\\[1mm]
\makebox[\wd\hmboard]{\sffamily\hspace{3mm}\textbf{h\#2.5}\hfill(4+2)}\end{minipage}\hfill
\begin{minipage}[t]{52mm}\vspace{0pt}
\textcolor{ink}{\sffamily\textsc{Lösung}}\par\smallskip\small
\begin{tabular}{@{}r@{\enspace}l@{\enspace}l@{}}
\textbf{1.} & \ldots & \symqueen{}a2+\\
\textbf{2.} & \symrook{}xa2+ & \symking{}xa2\\
\textbf{3.} & \symking{}d2 & \symqueen{}gd3\#\\
\end{tabular}\par\vspace{3mm}
\qrcode[height=12mm]{https://helpman.komtera.lt/?fen=8/8/8/3Q4/8/r5Q1/2k5/K4Q2&moves=2.5}\\{\sffamily\tiny\color{ink}scannen zum Lösen}
\end{minipage}
\vspace{2mm}

{\raggedright\noindent\textcolor{ink}{\sffamily\textsc{Themen}}\enspace \textit{keine erkannt}\par\smallskip{\footnotesize\textit{Hinweis:} die Partei am Zug steht im Diagramm im Schach.}\par}

\clearpage
\phantomsection\label{hm:KQQQBvk}
\begin{tcolorbox}[enhanced,colback=ink,colframe=ink,arc=0pt,boxrule=0pt,left=4mm,right=4mm,top=2.5mm,bottom=2.5mm]
\color{white}\sffamily{\LARGE\bfseries KQQQBvk}\hfill{\Large h\#2}\\[0.5mm]
{\small 6 Steine\quad Schwarz am Zug}\hfill{\small Nr.~853\qquad\href{https://helpman.komtera.lt/?fen=1Q6/8/8/8/8/8/1Q1B4/K1kQ4&moves=2}{\color{white}\bfseries lösen\,\ensuremath{\nearrow}}}
\end{tcolorbox}
\vspace{2mm}
\noindent\begin{minipage}[t]{36mm}\vspace{0pt}\sffamily\small\raggedright
\hmfigure{längste eindeutige}{{\Large h\#2}}\hmfigure{eindeutig bei dieser Länge}{{\Large 142\,565} Stellungen}\hmfigure{längstes Matt}{{\Large h\#3}}\hmfigure{in der Klasse}{15,5 Milliarden Stellungen,\\ 1,2 Milliarden mit Hilfsmatt}
\end{minipage}\hfill
\begin{minipage}[t]{80mm}\vspace{0pt}\begin{tikzpicture}\begin{axis}[width=72mm,height=42mm,ybar,bar width=1.6pt,xmin=-1,xmax=7,ymin=0,ymax=659.3,axis y line*=left,axis x line*=bottom,xlabel={\scriptsize Halbzüge bis Matt},ylabel={\scriptsize Stellungen (Mio.)},tick label style={font=\tiny},label style={font=\scriptsize},enlarge x limits=false,ylabel shift=-3pt,xlabel shift=-3pt]\addplot[fill=ink!55,draw=none] coordinates {(0,166.5) (1,150.2) (2,599.3) (3,89.12) (4,198.1) (5,0.02607) (6,0.1226)};\draw[densely dotted] (axis cs:4,0) -- (axis cs:4,659.3);\end{axis}\begin{axis}[width=72mm,height=42mm,axis y line*=right,axis x line=none,ymode=log,xmin=-1,xmax=7,ylabel={\scriptsize eindeutig (log)},tick label style={font=\tiny},label style={font=\scriptsize},ylabel shift=-3pt]\addplot[mark=*,mark size=0.7pt,draw=black,thin] coordinates {(0,166481629) (1,29291606) (2,77344387) (3,28842) (4,142565)};\end{axis}\end{tikzpicture}\end{minipage}
\vspace{3mm}

\noindent\begin{minipage}[t]{60mm}\vspace{0pt}\sbox\hmboard{\chessboard[setfen=1Q6/8/8/8/8/8/1Q1B4/K1kQ4 b - - 0 1]}\usebox\hmboard\\[1mm]
\makebox[\wd\hmboard]{\sffamily\hspace{3mm}\textbf{h\#2}\hfill(5+1)}\end{minipage}\hfill
\begin{minipage}[t]{52mm}\vspace{0pt}
\textcolor{ink}{\sffamily\textsc{Lösung}}\par\smallskip\small
\begin{tabular}{@{}r@{\enspace}l@{\enspace}l@{}}
\textbf{1.} & \symking{}xd1 & \symqueen{}d4\\
\textbf{2.} & \symking{}c2 & \symqueen{}b1\#\\
\end{tabular}\par\vspace{3mm}
\qrcode[height=12mm]{https://helpman.komtera.lt/?fen=1Q6/8/8/8/8/8/1Q1B4/K1kQ4&moves=2}\\{\sffamily\tiny\color{ink}scannen zum Lösen}
\end{minipage}
\vspace{2mm}

{\raggedright\noindent\textcolor{ink}{\sffamily\textsc{Themen}}\enspace \themetag{ein Stein zieht} \themetag{ohne Schach}\par\smallskip{\footnotesize\textit{Hinweis:} das Diagramm hat keinen legalen letzten Zug, kann also in keiner Partie entstehen.}\par}

\clearpage
\phantomsection\label{hm:KQQQNvk}
\begin{tcolorbox}[enhanced,colback=ink,colframe=ink,arc=0pt,boxrule=0pt,left=4mm,right=4mm,top=2.5mm,bottom=2.5mm]
\color{white}\sffamily{\LARGE\bfseries KQQQNvk}\hfill{\Large h\#2}\\[0.5mm]
{\small 6 Steine\quad Schwarz am Zug}\hfill{\small Nr.~854\qquad\href{https://helpman.komtera.lt/?fen=8/2Q5/8/2Q5/8/8/8/K1kQ1N2&moves=2}{\color{white}\bfseries lösen\,\ensuremath{\nearrow}}}
\end{tcolorbox}
\vspace{2mm}
\noindent\begin{minipage}[t]{36mm}\vspace{0pt}\sffamily\small\raggedright
\hmfigure{längste eindeutige}{{\Large h\#2}}\hmfigure{eindeutig bei dieser Länge}{{\Large 179\,748} Stellungen}\hmfigure{längstes Matt}{{\Large h\#3}}\hmfigure{in der Klasse}{15,5 Milliarden Stellungen,\\ 1,23 Milliarden mit Hilfsmatt}
\end{minipage}\hfill
\begin{minipage}[t]{80mm}\vspace{0pt}\begin{tikzpicture}\begin{axis}[width=72mm,height=42mm,ybar,bar width=1.6pt,xmin=-1,xmax=7,ymin=0,ymax=661.7,axis y line*=left,axis x line*=bottom,xlabel={\scriptsize Halbzüge bis Matt},ylabel={\scriptsize Stellungen (Mio.)},tick label style={font=\tiny},label style={font=\scriptsize},enlarge x limits=false,ylabel shift=-3pt,xlabel shift=-3pt]\addplot[fill=ink!55,draw=none] coordinates {(0,144.4) (1,155.2) (2,601.5) (3,105.1) (4,219.6) (5,0.03748) (6,0.1595)};\draw[densely dotted] (axis cs:4,0) -- (axis cs:4,661.7);\end{axis}\begin{axis}[width=72mm,height=42mm,axis y line*=right,axis x line=none,ymode=log,xmin=-1,xmax=7,ylabel={\scriptsize eindeutig (log)},tick label style={font=\tiny},label style={font=\scriptsize},ylabel shift=-3pt]\addplot[mark=*,mark size=0.7pt,draw=black,thin] coordinates {(0,144425432) (1,33723827) (2,85462286) (3,39638) (4,179748)};\end{axis}\end{tikzpicture}\end{minipage}
\vspace{3mm}

\noindent\begin{minipage}[t]{60mm}\vspace{0pt}\sbox\hmboard{\chessboard[setfen=8/2Q5/8/2Q5/8/8/8/K1kQ1N2 b - - 0 1]}\usebox\hmboard\\[1mm]
\makebox[\wd\hmboard]{\sffamily\hspace{3mm}\textbf{h\#2}\hfill(5+1)}\end{minipage}\hfill
\begin{minipage}[t]{52mm}\vspace{0pt}
\textcolor{ink}{\sffamily\textsc{Lösung}}\par\smallskip\small
\begin{tabular}{@{}r@{\enspace}l@{\enspace}l@{}}
\textbf{1.} & \symking{}xd1 & \symqueen{}h2\\
\textbf{2.} & \symking{}e1 & \symqueen{}c1\#\\
\end{tabular}\par\vspace{3mm}
\qrcode[height=12mm]{https://helpman.komtera.lt/?fen=8/2Q5/8/2Q5/8/8/8/K1kQ1N2&moves=2}\\{\sffamily\tiny\color{ink}scannen zum Lösen}
\end{minipage}
\vspace{2mm}

{\raggedright\noindent\textcolor{ink}{\sffamily\textsc{Themen}}\enspace \themetag{ein Stein zieht} \themetag{ohne Schach}\par\smallskip{\footnotesize\textit{Hinweis:} das Diagramm hat keinen legalen letzten Zug, kann also in keiner Partie entstehen.}\par}

\clearpage
\phantomsection\label{hm:KQQQPvk}
\begin{tcolorbox}[enhanced,colback=ink,colframe=ink,arc=0pt,boxrule=0pt,left=4mm,right=4mm,top=2.5mm,bottom=2.5mm]
\color{white}\sffamily{\LARGE\bfseries KQQQPvk}\hfill{\Large h\#2}\\[0.5mm]
{\small 6 Steine\quad Schwarz am Zug}\hfill{\small Nr.~855\qquad\href{https://helpman.komtera.lt/?fen=2Q5/8/8/8/8/2Q2P2/8/K1kQ4&moves=2}{\color{white}\bfseries lösen\,\ensuremath{\nearrow}}}
\end{tcolorbox}
\vspace{2mm}
\noindent\begin{minipage}[t]{36mm}\vspace{0pt}\sffamily\small\raggedright
\hmfigure{längste eindeutige}{{\Large h\#2}}\hmfigure{eindeutig bei dieser Länge}{{\Large 1\,041\,105} Stellungen}\hmfigure{längstes Matt}{{\Large h\#3}}\hmfigure{in der Klasse}{45,4 Milliarden Stellungen,\\ 3,73 Milliarden mit Hilfsmatt}
\end{minipage}\hfill
\begin{minipage}[t]{80mm}\vspace{0pt}\begin{tikzpicture}\begin{axis}[width=72mm,height=42mm,ybar,bar width=1.6pt,xmin=-1,xmax=7,ymin=0,ymax=1.919,axis y line*=left,axis x line*=bottom,xlabel={\scriptsize Halbzüge bis Matt},ylabel={\scriptsize Stellungen (Mrd.)},tick label style={font=\tiny},label style={font=\scriptsize},enlarge x limits=false,ylabel shift=-3pt,xlabel shift=-3pt]\addplot[fill=ink!55,draw=none] coordinates {(0,0.3958) (1,0.4492) (2,1.745) (3,0.3759) (4,0.7661) (5,0.0003417) (6,0.001054)};\draw[densely dotted] (axis cs:4,0) -- (axis cs:4,1.919);\end{axis}\begin{axis}[width=72mm,height=42mm,axis y line*=right,axis x line=none,ymode=log,xmin=-1,xmax=7,ylabel={\scriptsize eindeutig (log)},tick label style={font=\tiny},label style={font=\scriptsize},ylabel shift=-3pt]\addplot[mark=*,mark size=0.7pt,draw=black,thin] coordinates {(0,395779575) (1,98784806) (2,262140335) (3,362658) (4,1041105)};\end{axis}\end{tikzpicture}\end{minipage}
\vspace{3mm}

\noindent\begin{minipage}[t]{60mm}\vspace{0pt}\sbox\hmboard{\chessboard[setfen=2Q5/8/8/8/8/2Q2P2/8/K1kQ4 b - - 0 1]}\usebox\hmboard\\[1mm]
\makebox[\wd\hmboard]{\sffamily\hspace{3mm}\textbf{h\#2}\hfill(5+1)}\end{minipage}\hfill
\begin{minipage}[t]{52mm}\vspace{0pt}
\textcolor{ink}{\sffamily\textsc{Lösung}}\par\smallskip\small
\begin{tabular}{@{}r@{\enspace}l@{\enspace}l@{}}
\textbf{1.} & \symking{}xd1 & \symqueen{}b2\\
\textbf{2.} & \symking{}e1 & \symqueen{}cc1\#\\
\end{tabular}\par\vspace{3mm}
\qrcode[height=12mm]{https://helpman.komtera.lt/?fen=2Q5/8/8/8/8/2Q2P2/8/K1kQ4&moves=2}\\{\sffamily\tiny\color{ink}scannen zum Lösen}
\end{minipage}
\vspace{2mm}

{\raggedright\noindent\textcolor{ink}{\sffamily\textsc{Themen}}\enspace \themetag{Spiegelmatt} \themetag{ein Stein zieht} \themetag{ohne Schach}\par\smallskip{\footnotesize\textit{Hinweis:} das Diagramm hat keinen legalen letzten Zug, kann also in keiner Partie entstehen.}\par}

\clearpage
\phantomsection\label{hm:KQQQQvk}
\begin{tcolorbox}[enhanced,colback=ink,colframe=ink,arc=0pt,boxrule=0pt,left=4mm,right=4mm,top=2.5mm,bottom=2.5mm]
\color{white}\sffamily{\LARGE\bfseries KQQQQvk}\hfill{\Large h\#2}\\[0.5mm]
{\small 6 Steine\quad Schwarz am Zug}\hfill{\small Nr.~856\qquad\href{https://helpman.komtera.lt/?fen=8/Q1Q5/8/8/8/8/8/QK4kQ&moves=2}{\color{white}\bfseries lösen\,\ensuremath{\nearrow}}}
\end{tcolorbox}
\vspace{2mm}
\noindent\begin{minipage}[t]{36mm}\vspace{0pt}\sffamily\small\raggedright
\hmfigure{längste eindeutige}{{\Large h\#2}}\hmfigure{eindeutig bei dieser Länge}{{\Large 251} Stellungen}\hmfigure{längstes Matt}{{\Large h\#3}}\hmfigure{in der Klasse}{15,5 Milliarden Stellungen,\\ 282 Millionen mit Hilfsmatt}
\end{minipage}\hfill
\begin{minipage}[t]{80mm}\vspace{0pt}\begin{tikzpicture}\begin{axis}[width=72mm,height=42mm,ybar,bar width=1.6pt,xmin=-1,xmax=7,ymin=0,ymax=157.1,axis y line*=left,axis x line*=bottom,xlabel={\scriptsize Halbzüge bis Matt},ylabel={\scriptsize Stellungen (Mio.)},tick label style={font=\tiny},label style={font=\scriptsize},enlarge x limits=false,ylabel shift=-3pt,xlabel shift=-3pt]\addplot[fill=ink!55,draw=none] coordinates {(0,71.98) (1,33.07) (2,142.8) (3,9.178) (4,24.86) (5,1e-05) (6,8.7e-05)};\draw[densely dotted] (axis cs:4,0) -- (axis cs:4,157.1);\end{axis}\begin{axis}[width=72mm,height=42mm,axis y line*=right,axis x line=none,ymode=log,xmin=-1,xmax=7,ylabel={\scriptsize eindeutig (log)},tick label style={font=\tiny},label style={font=\scriptsize},ylabel shift=-3pt]\addplot[mark=*,mark size=0.7pt,draw=black,thin] coordinates {(0,71982891) (1,3168119) (2,8950846) (3,30) (4,251)};\end{axis}\end{tikzpicture}\end{minipage}
\vspace{3mm}

\noindent\begin{minipage}[t]{60mm}\vspace{0pt}\sbox\hmboard{\chessboard[setfen=8/Q1Q5/8/8/8/8/8/QK4kQ b - - 0 1]}\usebox\hmboard\\[1mm]
\makebox[\wd\hmboard]{\sffamily\hspace{3mm}\textbf{h\#2}\hfill(5+1)}\end{minipage}\hfill
\begin{minipage}[t]{52mm}\vspace{0pt}
\textcolor{ink}{\sffamily\textsc{Lösung}}\par\smallskip\small
\begin{tabular}{@{}r@{\enspace}l@{\enspace}l@{}}
\textbf{1.} & \symking{}xh1 & \symqueen{}7a2\\
\textbf{2.} & \symking{}g1 & \symqueen{}c1\#\\
\end{tabular}\par\vspace{3mm}
\qrcode[height=12mm]{https://helpman.komtera.lt/?fen=8/Q1Q5/8/8/8/8/8/QK4kQ&moves=2}\\{\sffamily\tiny\color{ink}scannen zum Lösen}
\end{minipage}
\vspace{2mm}

{\raggedright\noindent\textcolor{ink}{\sffamily\textsc{Themen}}\enspace \themetag{reines Matt} \themetag{Spiegelmatt} \themetag{Switchback} \themetag{ein Stein zieht} \themetag{ohne Schach}\par\smallskip{\footnotesize\textit{Hinweis:} das Diagramm hat keinen legalen letzten Zug, kann also in keiner Partie entstehen.}\par}

\clearpage
\phantomsection\label{hm:KQQQRvk}
\begin{tcolorbox}[enhanced,colback=ink,colframe=ink,arc=0pt,boxrule=0pt,left=4mm,right=4mm,top=2.5mm,bottom=2.5mm]
\color{white}\sffamily{\LARGE\bfseries KQQQRvk}\hfill{\Large h\#2}\\[0.5mm]
{\small 6 Steine\quad Schwarz am Zug}\hfill{\small Nr.~857\qquad\href{https://helpman.komtera.lt/?fen=8/8/8/7Q/5R1Q/8/8/K1Qk4&moves=2}{\color{white}\bfseries lösen\,\ensuremath{\nearrow}}}
\end{tcolorbox}
\vspace{2mm}
\noindent\begin{minipage}[t]{36mm}\vspace{0pt}\sffamily\small\raggedright
\hmfigure{längste eindeutige}{{\Large h\#2}}\hmfigure{eindeutig bei dieser Länge}{{\Large 14\,587} Stellungen}\hmfigure{längstes Matt}{{\Large h\#3}}\hmfigure{in der Klasse}{15,5 Milliarden Stellungen,\\ 1,17 Milliarden mit Hilfsmatt}
\end{minipage}\hfill
\begin{minipage}[t]{80mm}\vspace{0pt}\begin{tikzpicture}\begin{axis}[width=72mm,height=42mm,ybar,bar width=1.6pt,xmin=-1,xmax=7,ymin=0,ymax=639.7,axis y line*=left,axis x line*=bottom,xlabel={\scriptsize Halbzüge bis Matt},ylabel={\scriptsize Stellungen (Mio.)},tick label style={font=\tiny},label style={font=\scriptsize},enlarge x limits=false,ylabel shift=-3pt,xlabel shift=-3pt]\addplot[fill=ink!55,draw=none] coordinates {(0,214.6) (1,142.5) (2,581.6) (3,67.72) (4,165.9) (5,0.001733) (6,0.006979)};\draw[densely dotted] (axis cs:4,0) -- (axis cs:4,639.7);\end{axis}\begin{axis}[width=72mm,height=42mm,axis y line*=right,axis x line=none,ymode=log,xmin=-1,xmax=7,ylabel={\scriptsize eindeutig (log)},tick label style={font=\tiny},label style={font=\scriptsize},ylabel shift=-3pt]\addplot[mark=*,mark size=0.7pt,draw=black,thin] coordinates {(0,214553967) (1,17796775) (2,51379465) (3,3961) (4,14587)};\end{axis}\end{tikzpicture}\end{minipage}
\vspace{3mm}

\noindent\begin{minipage}[t]{60mm}\vspace{0pt}\sbox\hmboard{\chessboard[setfen=8/8/8/7Q/5R1Q/8/8/K1Qk4 b - - 0 1]}\usebox\hmboard\\[1mm]
\makebox[\wd\hmboard]{\sffamily\hspace{3mm}\textbf{h\#2}\hfill(5+1)}\end{minipage}\hfill
\begin{minipage}[t]{52mm}\vspace{0pt}
\textcolor{ink}{\sffamily\textsc{Lösung}}\par\smallskip\small
\begin{tabular}{@{}r@{\enspace}l@{\enspace}l@{}}
\textbf{1.} & \symking{}xc1 & \symqueen{}e1+\\
\textbf{2.} & \symking{}c2 & \symqueen{}hd1\#\\
\end{tabular}\par\vspace{3mm}
\qrcode[height=12mm]{https://helpman.komtera.lt/?fen=8/8/8/7Q/5R1Q/8/8/K1Qk4&moves=2}\\{\sffamily\tiny\color{ink}scannen zum Lösen}
\end{minipage}
\vspace{2mm}

{\raggedright\noindent\textcolor{ink}{\sffamily\textsc{Themen}}\enspace \themetag{ein Stein zieht}\par\smallskip{\footnotesize\textit{Hinweis:} das Diagramm hat keinen legalen letzten Zug, kann also in keiner Partie entstehen.}\par}

\clearpage
\phantomsection\label{hm:KQQRBvk}
\begin{tcolorbox}[enhanced,colback=ink,colframe=ink,arc=0pt,boxrule=0pt,left=4mm,right=4mm,top=2.5mm,bottom=2.5mm]
\color{white}\sffamily{\LARGE\bfseries KQQRBvk}\hfill{\Large h\#2}\\[0.5mm]
{\small 6 Steine\quad Schwarz am Zug}\hfill{\small Nr.~858\qquad\href{https://helpman.komtera.lt/?fen=8/8/8/8/8/8/3Q3B/KQk1R3&moves=2}{\color{white}\bfseries lösen\,\ensuremath{\nearrow}}}
\end{tcolorbox}
\vspace{2mm}
\noindent\begin{minipage}[t]{36mm}\vspace{0pt}\sffamily\small\raggedright
\hmfigure{längste eindeutige}{{\Large h\#2}}\hmfigure{eindeutig bei dieser Länge}{{\Large 5\,705\,691} Stellungen}\hmfigure{längstes Matt}{{\Large h\#4}}\hmfigure{in der Klasse}{15,5 Milliarden Stellungen,\\ 3,8 Milliarden mit Hilfsmatt}
\end{minipage}\hfill
\begin{minipage}[t]{80mm}\vspace{0pt}\begin{tikzpicture}\begin{axis}[width=72mm,height=42mm,ybar,bar width=1.6pt,xmin=-1,xmax=9,ymin=0,ymax=1.8,axis y line*=left,axis x line*=bottom,xlabel={\scriptsize Halbzüge bis Matt},ylabel={\scriptsize Stellungen (Mrd.)},tick label style={font=\tiny},label style={font=\scriptsize},enlarge x limits=false,ylabel shift=-3pt,xlabel shift=-3pt]\addplot[fill=ink!55,draw=none] coordinates {(0,0.3193) (1,0.4311) (2,1.637) (3,0.4563) (4,0.9529) (5,0.0007822) (6,0.004994) (8,1e-09)};\draw[densely dotted] (axis cs:4,0) -- (axis cs:4,1.8);\end{axis}\begin{axis}[width=72mm,height=42mm,axis y line*=right,axis x line=none,ymode=log,xmin=-1,xmax=9,ylabel={\scriptsize eindeutig (log)},tick label style={font=\tiny},label style={font=\scriptsize},ylabel shift=-3pt]\addplot[mark=*,mark size=0.7pt,draw=black,thin] coordinates {(0,319299098) (1,121390864) (2,326745800) (3,873681) (4,5705691)};\end{axis}\end{tikzpicture}\end{minipage}
\vspace{3mm}

\noindent\begin{minipage}[t]{60mm}\vspace{0pt}\sbox\hmboard{\chessboard[setfen=8/8/8/8/8/8/3Q3B/KQk1R3 b - - 0 1]}\usebox\hmboard\\[1mm]
\makebox[\wd\hmboard]{\sffamily\hspace{3mm}\textbf{h\#2}\hfill(5+1)}\end{minipage}\hfill
\begin{minipage}[t]{52mm}\vspace{0pt}
\textcolor{ink}{\sffamily\textsc{Lösung}}\par\smallskip\small
\begin{tabular}{@{}r@{\enspace}l@{\enspace}l@{}}
\textbf{1.} & \symking{}xd2 & \symqueen{}d3+\\
\textbf{2.} & \symking{}xe1 & \symbishop{}g3\#\\
\end{tabular}\par\vspace{3mm}
\qrcode[height=12mm]{https://helpman.komtera.lt/?fen=8/8/8/8/8/8/3Q3B/KQk1R3&moves=2}\\{\sffamily\tiny\color{ink}scannen zum Lösen}
\end{minipage}
\vspace{2mm}

{\raggedright\noindent\textcolor{ink}{\sffamily\textsc{Themen}}\enspace \themetag{reines Matt} \themetag{Mustermatt} \themetag{Spiegelmatt} \themetag{ein Stein zieht} \themetag{Kniest}\par\smallskip{\footnotesize\textit{Hinweis:} das Diagramm hat keinen legalen letzten Zug, kann also in keiner Partie entstehen.}\par}

\clearpage
\phantomsection\label{hm:KQQRRvk}
\begin{tcolorbox}[enhanced,colback=ink,colframe=ink,arc=0pt,boxrule=0pt,left=4mm,right=4mm,top=2.5mm,bottom=2.5mm]
\color{white}\sffamily{\LARGE\bfseries KQQRRvk}\hfill{\Large h\#2}\\[0.5mm]
{\small 6 Steine\quad Schwarz am Zug}\hfill{\small Nr.~859\qquad\href{https://helpman.komtera.lt/?fen=8/8/8/8/R1R5/2Q5/8/K1kQ4&moves=2}{\color{white}\bfseries lösen\,\ensuremath{\nearrow}}}
\end{tcolorbox}
\vspace{2mm}
\noindent\begin{minipage}[t]{36mm}\vspace{0pt}\sffamily\small\raggedright
\hmfigure{längste eindeutige}{{\Large h\#2}}\hmfigure{eindeutig bei dieser Länge}{{\Large 174\,709} Stellungen}\hmfigure{längstes Matt}{{\Large h\#3}}\hmfigure{in der Klasse}{15,5 Milliarden Stellungen,\\ 1,84 Milliarden mit Hilfsmatt}
\end{minipage}\hfill
\begin{minipage}[t]{80mm}\vspace{0pt}\begin{tikzpicture}\begin{axis}[width=72mm,height=42mm,ybar,bar width=1.6pt,xmin=-1,xmax=7,ymin=0,ymax=907.6,axis y line*=left,axis x line*=bottom,xlabel={\scriptsize Halbzüge bis Matt},ylabel={\scriptsize Stellungen (Mio.)},tick label style={font=\tiny},label style={font=\scriptsize},enlarge x limits=false,ylabel shift=-3pt,xlabel shift=-3pt]\addplot[fill=ink!55,draw=none] coordinates {(0,231.9) (1,215.5) (2,825) (3,169.7) (4,395.9) (5,0.04608) (6,0.2111)};\draw[densely dotted] (axis cs:4,0) -- (axis cs:4,907.6);\end{axis}\begin{axis}[width=72mm,height=42mm,axis y line*=right,axis x line=none,ymode=log,xmin=-1,xmax=7,ylabel={\scriptsize eindeutig (log)},tick label style={font=\tiny},label style={font=\scriptsize},ylabel shift=-3pt]\addplot[mark=*,mark size=0.7pt,draw=black,thin] coordinates {(0,231934531) (1,33914205) (2,99349501) (3,52665) (4,174709)};\end{axis}\end{tikzpicture}\end{minipage}
\vspace{3mm}

\noindent\begin{minipage}[t]{60mm}\vspace{0pt}\sbox\hmboard{\chessboard[setfen=8/8/8/8/R1R5/2Q5/8/K1kQ4 b - - 0 1]}\usebox\hmboard\\[1mm]
\makebox[\wd\hmboard]{\sffamily\hspace{3mm}\textbf{h\#2}\hfill(5+1)}\end{minipage}\hfill
\begin{minipage}[t]{52mm}\vspace{0pt}
\textcolor{ink}{\sffamily\textsc{Lösung}}\par\smallskip\small
\begin{tabular}{@{}r@{\enspace}l@{\enspace}l@{}}
\textbf{1.} & \symking{}xd1 & \symqueen{}b2\\
\textbf{2.} & \symking{}e1 & \symrook{}c1\#\\
\end{tabular}\par\vspace{3mm}
\qrcode[height=12mm]{https://helpman.komtera.lt/?fen=8/8/8/8/R1R5/2Q5/8/K1kQ4&moves=2}\\{\sffamily\tiny\color{ink}scannen zum Lösen}
\end{minipage}
\vspace{2mm}

{\raggedright\noindent\textcolor{ink}{\sffamily\textsc{Themen}}\enspace \themetag{reines Matt} \themetag{Spiegelmatt} \themetag{ein Stein zieht} \themetag{ohne Schach}\par\smallskip{\footnotesize\textit{Hinweis:} das Diagramm hat keinen legalen letzten Zug, kann also in keiner Partie entstehen.}\par}

\clearpage
\phantomsection\label{hm:KQRRRvk}
\begin{tcolorbox}[enhanced,colback=ink,colframe=ink,arc=0pt,boxrule=0pt,left=4mm,right=4mm,top=2.5mm,bottom=2.5mm]
\color{white}\sffamily{\LARGE\bfseries KQRRRvk}\hfill{\Large h\#2}\\[0.5mm]
{\small 6 Steine\quad Schwarz am Zug}\hfill{\small Nr.~860\qquad\href{https://helpman.komtera.lt/?fen=8/8/8/8/8/2R1R3/4R3/K1kQ4&moves=2}{\color{white}\bfseries lösen\,\ensuremath{\nearrow}}}
\end{tcolorbox}
\vspace{2mm}
\noindent\begin{minipage}[t]{36mm}\vspace{0pt}\sffamily\small\raggedright
\hmfigure{längste eindeutige}{{\Large h\#2}}\hmfigure{eindeutig bei dieser Länge}{{\Large 272\,140} Stellungen}\hmfigure{längstes Matt}{{\Large h\#3}}\hmfigure{in der Klasse}{15,5 Milliarden Stellungen,\\ 1,29 Milliarden mit Hilfsmatt}
\end{minipage}\hfill
\begin{minipage}[t]{80mm}\vspace{0pt}\begin{tikzpicture}\begin{axis}[width=72mm,height=42mm,ybar,bar width=1.6pt,xmin=-1,xmax=7,ymin=0,ymax=521,axis y line*=left,axis x line*=bottom,xlabel={\scriptsize Halbzüge bis Matt},ylabel={\scriptsize Stellungen (Mio.)},tick label style={font=\tiny},label style={font=\scriptsize},enlarge x limits=false,ylabel shift=-3pt,xlabel shift=-3pt]\addplot[fill=ink!55,draw=none] coordinates {(0,109.8) (1,136.2) (2,473.6) (3,172.8) (4,394.4) (5,0.261) (6,0.4673)};\draw[densely dotted] (axis cs:4,0) -- (axis cs:4,521);\end{axis}\begin{axis}[width=72mm,height=42mm,axis y line*=right,axis x line=none,ymode=log,xmin=-1,xmax=7,ylabel={\scriptsize eindeutig (log)},tick label style={font=\tiny},label style={font=\scriptsize},ylabel shift=-3pt]\addplot[mark=*,mark size=0.7pt,draw=black,thin] coordinates {(0,109793150) (1,27636764) (2,65557853) (3,178210) (4,272140)};\end{axis}\end{tikzpicture}\end{minipage}
\vspace{3mm}

\noindent\begin{minipage}[t]{60mm}\vspace{0pt}\sbox\hmboard{\chessboard[setfen=8/8/8/8/8/2R1R3/4R3/K1kQ4 b - - 0 1]}\usebox\hmboard\\[1mm]
\makebox[\wd\hmboard]{\sffamily\hspace{3mm}\textbf{h\#2}\hfill(5+1)}\end{minipage}\hfill
\begin{minipage}[t]{52mm}\vspace{0pt}
\textcolor{ink}{\sffamily\textsc{Lösung}}\par\smallskip\small
\begin{tabular}{@{}r@{\enspace}l@{\enspace}l@{}}
\textbf{1.} & \symking{}xd1 & \symrook{}b3\\
\textbf{2.} & \symking{}c1 & \symrook{}b1\#\\
\end{tabular}\par\vspace{3mm}
\qrcode[height=12mm]{https://helpman.komtera.lt/?fen=8/8/8/8/8/2R1R3/4R3/K1kQ4&moves=2}\\{\sffamily\tiny\color{ink}scannen zum Lösen}
\end{minipage}
\vspace{2mm}

{\raggedright\noindent\textcolor{ink}{\sffamily\textsc{Themen}}\enspace \themetag{Switchback} \themetag{ein Stein zieht} \themetag{ohne Schach}\par\smallskip{\footnotesize\textit{Hinweis:} das Diagramm hat keinen legalen letzten Zug, kann also in keiner Partie entstehen.}\par}

\clearpage
\phantomsection\label{hm:KRRRRvk}
\begin{tcolorbox}[enhanced,colback=ink,colframe=ink,arc=0pt,boxrule=0pt,left=4mm,right=4mm,top=2.5mm,bottom=2.5mm]
\color{white}\sffamily{\LARGE\bfseries KRRRRvk}\hfill{\Large h\#2}\\[0.5mm]
{\small 6 Steine\quad Schwarz am Zug}\hfill{\small Nr.~861\qquad\href{https://helpman.komtera.lt/?fen=8/8/8/3R4/8/8/8/K1kRRR2&moves=2}{\color{white}\bfseries lösen\,\ensuremath{\nearrow}}}
\end{tcolorbox}
\vspace{2mm}
\noindent\begin{minipage}[t]{36mm}\vspace{0pt}\sffamily\small\raggedright
\hmfigure{längste eindeutige}{{\Large h\#2}}\hmfigure{eindeutig bei dieser Länge}{{\Large 102\,517} Stellungen}\hmfigure{längstes Matt}{{\Large h\#3}}\hmfigure{in der Klasse}{15,5 Milliarden Stellungen,\\ 340 Millionen mit Hilfsmatt}
\end{minipage}\hfill
\begin{minipage}[t]{80mm}\vspace{0pt}\begin{tikzpicture}\begin{axis}[width=72mm,height=42mm,ybar,bar width=1.6pt,xmin=-1,xmax=7,ymin=0,ymax=125.3,axis y line*=left,axis x line*=bottom,xlabel={\scriptsize Halbzüge bis Matt},ylabel={\scriptsize Stellungen (Mio.)},tick label style={font=\tiny},label style={font=\scriptsize},enlarge x limits=false,ylabel shift=-3pt,xlabel shift=-3pt]\addplot[fill=ink!55,draw=none] coordinates {(0,20.41) (1,34.34) (2,113.9) (3,57.52) (4,113) (5,0.1981) (6,0.1452)};\draw[densely dotted] (axis cs:4,0) -- (axis cs:4,125.3);\end{axis}\begin{axis}[width=72mm,height=42mm,axis y line*=right,axis x line=none,ymode=log,xmin=-1,xmax=7,ylabel={\scriptsize eindeutig (log)},tick label style={font=\tiny},label style={font=\scriptsize},ylabel shift=-3pt]\addplot[mark=*,mark size=0.7pt,draw=black,thin] coordinates {(0,20412026) (1,11860484) (2,27100639) (3,68953) (4,102517)};\end{axis}\end{tikzpicture}\end{minipage}
\vspace{3mm}

\noindent\begin{minipage}[t]{60mm}\vspace{0pt}\sbox\hmboard{\chessboard[setfen=8/8/8/3R4/8/8/8/K1kRRR2 b - - 0 1]}\usebox\hmboard\\[1mm]
\makebox[\wd\hmboard]{\sffamily\hspace{3mm}\textbf{h\#2}\hfill(5+1)}\end{minipage}\hfill
\begin{minipage}[t]{52mm}\vspace{0pt}
\textcolor{ink}{\sffamily\textsc{Lösung}}\par\smallskip\small
\begin{tabular}{@{}r@{\enspace}l@{\enspace}l@{}}
\textbf{1.} & \symking{}c2 & \symrook{}b1\\
\textbf{2.} & \symking{}c3 & \symrook{}ec1\#\\
\end{tabular}\par\vspace{3mm}
\qrcode[height=12mm]{https://helpman.komtera.lt/?fen=8/8/8/3R4/8/8/8/K1kRRR2&moves=2}\\{\sffamily\tiny\color{ink}scannen zum Lösen}
\end{minipage}
\vspace{2mm}

{\raggedright\noindent\textcolor{ink}{\sffamily\textsc{Themen}}\enspace \themetag{Spiegelmatt} \themetag{ein Stein zieht} \themetag{ohne Schlag} \themetag{ohne Schach}\par\smallskip{\footnotesize\textit{Hinweis:} die Partei am Zug steht im Diagramm im Schach.}\par}

\clearpage

\appendix
\chapter{Verzeichnis der Materialklassen}
\markboth{Verzeichnis der Materialklassen}{}

Sortiert nach der Länge der eindeutigen Aufgabe. \emph{Klassenmax.} ist das längste Hilfsmatt in diesem Material bei beliebiger Lösungszahl.

{\footnotesize
\begin{longtable}{lrrrr>{\raggedright\arraybackslash}p{22mm}r}
\toprule
Material & Steine & eindeutig & Klassenmax. & Abstand & veröffentlicht von & Seite \\
\midrule
\endfirsthead
\toprule
Material & Steine & eindeutig & Klassenmax. & Abstand & veröffentlicht von & Seite \\
\midrule
\endhead
\bottomrule
\endfoot
KBvkbbp & 6 & \textbf{h\#13} & h\#16.5 & 7 &  & \pageref{hm:KBvkbbp} \\
KBvkrp & 5 & \textbf{h\#13} & h\#16.5 & 7 &  & \pageref{hm:KBvkrp} \\
KQBvkrp & 6 & \textbf{h\#13} & h\#17 & 8 &  & \pageref{hm:KQBvkrp} \\
KBvkqbp & 6 & \textbf{h\#12.5} & h\#16.5 & 8 &  & \pageref{hm:KBvkqbp} \\
KBvkqqp & 6 & \textbf{h\#12.5} & h\#17.5 & 10 &  & \pageref{hm:KBvkqqp} \\
KBvkqrp & 6 & \textbf{h\#12.5} & h\#16.5 & 8 &  & \pageref{hm:KBvkqrp} \\
KBvkrpp & 6 & \textbf{h\#12.5} & h\#16.5 & 8 &  & \pageref{hm:KBvkrpp} \\
KBvkrrp & 6 & \textbf{h\#12.5} & h\#16.5 & 8 &  & \pageref{hm:KBvkrrp} \\
KPvkbpp & 6 & \textbf{h\#12.5} & h\#15.5 & 6 &  & \pageref{hm:KPvkbpp} \\
KBBPvkp & 6 & \textbf{h\#12} & h\#15 & 6 &  & \pageref{hm:KBBPvkp} \\
KBvkqp & 5 & \textbf{h\#12} & h\#17 & 10 &  & \pageref{hm:KBvkqp} \\
KBvkrbp & 6 & \textbf{h\#12} & h\#16.5 & 9 &  & \pageref{hm:KBvkrbp} \\
KQBvkqp & 6 & \textbf{h\#12} & h\#17 & 10 &  & \pageref{hm:KQBvkqp} \\
KRBvkrp & 6 & \textbf{h\#12} & h\#17 & 10 &  & \pageref{hm:KRBvkrp} \\
KBvkbp & 5 & \textbf{h\#11.5} & h\#16 & 9 &  & \pageref{hm:KBvkbp} \\
KBvkbpp & 6 & \textbf{h\#11.5} & h\#16 & 9 &  & \pageref{hm:KBvkbpp} \\
KBvkqnp & 6 & \textbf{h\#11.5} & h\#16.5 & 10 &  & \pageref{hm:KBvkqnp} \\
KBvkqpp & 6 & \textbf{h\#11.5} & h\#16.5 & 10 &  & \pageref{hm:KBvkqpp} \\
KBvkrnp & 6 & \textbf{h\#11.5} & h\#16.5 & 10 &  & \pageref{hm:KBvkrnp} \\
KPvkppp & 6 & \textbf{h\#11.5} & h\#16.5 & 10 &  & \pageref{hm:KPvkppp} \\
KBBvkrp & 6 & \textbf{h\#11} & h\#17 & 12 &  & \pageref{hm:KBBvkrp} \\
KBPPvkp & 6 & \textbf{h\#11} & h\#14 & 6 &  & \pageref{hm:KBPPvkp} \\
KBPvkrp & 6 & \textbf{h\#11} & h\#16 & 10 &  & \pageref{hm:KBPvkrp} \\
KBvkppp & 6 & \textbf{h\#11} & h\#15.5 & 9 &  & \pageref{hm:KBvkppp} \\
KQBPvkp & 6 & \textbf{h\#11} & h\#15 & 8 &  & \pageref{hm:KQBPvkp} \\
KBBBvkp & 6 & \textbf{h\#10.5} & h\#15 & 9 &  & \pageref{hm:KBBBvkp} \\
KBPvkp & 5 & \textbf{h\#10.5} & h\#16 & 11 & Szentai (1958) & \pageref{hm:KBPvkp} \\
KBPvkpp & 6 & \textbf{h\#10.5} & h\#16 & 11 &  & \pageref{hm:KBPvkpp} \\
KBPvkqp & 6 & \textbf{h\#10.5} & h\#17 & 13 &  & \pageref{hm:KBPvkqp} \\
KBvkbnp & 6 & \textbf{h\#10.5} & h\#16.5 & 12 &  & \pageref{hm:KBvkbnp} \\
KBvknpp & 6 & \textbf{h\#10.5} & h\#14.5 & 8 &  & \pageref{hm:KBvknpp} \\
KNvkqqp & 6 & \textbf{h\#10.5} & h\#12.5 & 4 &  & \pageref{hm:KNvkqqp} \\
KPPvkpp & 6 & \textbf{h\#10.5} & h\#15.5 & 10 &  & \pageref{hm:KPPvkpp} \\
KPvkbbp & 6 & \textbf{h\#10.5} & h\#15 & 9 &  & \pageref{hm:KPvkbbp} \\
KBBvkbp & 6 & \textbf{h\#10} & h\#16 & 12 &  & \pageref{hm:KBBvkbp} \\
KBBvkqp & 6 & \textbf{h\#10} & h\#17 & 14 &  & \pageref{hm:KBBvkqp} \\
KBNPvkp & 6 & \textbf{h\#10} & h\#15 & 10 &  & \pageref{hm:KBNPvkp} \\
KBvkbbn & 6 & \textbf{h\#10} & h\#12 & 4 &  & \pageref{hm:KBvkbbn} \\
KBvkqnn & 6 & \textbf{h\#10} & h\#13 & 6 &  & \pageref{hm:KBvkqnn} \\
KNPvkpp & 6 & \textbf{h\#10} & h\#15 & 10 &  & \pageref{hm:KNPvkpp} \\
KPvkpp & 5 & \textbf{h\#10} & h\#15.5 & 11 & Ban (1967) & \pageref{hm:KPvkpp} \\
KQBvkbp & 6 & \textbf{h\#10} & h\#16.5 & 13 &  & \pageref{hm:KQBvkbp} \\
KQBvknp & 6 & \textbf{h\#10} & h\#16 & 12 &  & \pageref{hm:KQBvknp} \\
KQBvkqn & 6 & \textbf{h\#10} & h\#13 & 6 &  & \pageref{hm:KQBvkqn} \\
KQBvkrn & 6 & \textbf{h\#10} & h\#13 & 6 &  & \pageref{hm:KQBvkrn} \\
KQNvkpp & 6 & \textbf{h\#10} & h\#11 & 2 &  & \pageref{hm:KQNvkpp} \\
KQNvkqp & 6 & \textbf{h\#10} & h\#12 & 4 &  & \pageref{hm:KQNvkqp} \\
KQPvkbp & 6 & \textbf{h\#10} & h\#15 & 10 &  & \pageref{hm:KQPvkbp} \\
KQPvkpp & 6 & \textbf{h\#10} & h\#16 & 12 &  & \pageref{hm:KQPvkpp} \\
KRBPvkp & 6 & \textbf{h\#10} & h\#15 & 10 &  & \pageref{hm:KRBPvkp} \\
KRBvkbp & 6 & \textbf{h\#10} & h\#16.5 & 13 &  & \pageref{hm:KRBvkbp} \\
KRBvknp & 6 & \textbf{h\#10} & h\#15.5 & 11 &  & \pageref{hm:KRBvknp} \\
KRBvkqn & 6 & \textbf{h\#10} & h\#13 & 6 &  & \pageref{hm:KRBvkqn} \\
KRBvkqp & 6 & \textbf{h\#10} & h\#17 & 14 &  & \pageref{hm:KRBvkqp} \\
KRBvkrn & 6 & \textbf{h\#10} & h\#13 & 6 &  & \pageref{hm:KRBvkrn} \\
KBBvknp & 6 & \textbf{h\#9.5} & h\#14.5 & 10 &  & \pageref{hm:KBBvknp} \\
KBBvkpp & 6 & \textbf{h\#9.5} & h\#16 & 13 &  & \pageref{hm:KBBvkpp} \\
KBBvkqn & 6 & \textbf{h\#9.5} & h\#13 & 7 &  & \pageref{hm:KBBvkqn} \\
KBPvkbp & 6 & \textbf{h\#9.5} & h\#16 & 13 &  & \pageref{hm:KBPvkbp} \\
KBPvknp & 6 & \textbf{h\#9.5} & h\#14.5 & 10 &  & \pageref{hm:KBPvknp} \\
KBvkbnn & 6 & \textbf{h\#9.5} & h\#12.5 & 6 &  & \pageref{hm:KBvkbnn} \\
KBvknnn & 6 & \textbf{h\#9.5} & h\#12 & 5 &  & \pageref{hm:KBvknnn} \\
KBvknnp & 6 & \textbf{h\#9.5} & h\#12.5 & 6 &  & \pageref{hm:KBvknnp} \\
KBvknp & 5 & \textbf{h\#9.5} & h\#15.5 & 12 &  & \pageref{hm:KBvknp} \\
KBvkqbn & 6 & \textbf{h\#9.5} & h\#12.5 & 6 &  & \pageref{hm:KBvkqbn} \\
KBvkqn & 5 & \textbf{h\#9.5} & h\#13 & 7 &  & \pageref{hm:KBvkqn} \\
KBvkqqn & 6 & \textbf{h\#9.5} & h\#12.5 & 6 &  & \pageref{hm:KBvkqqn} \\
KBvkqrn & 6 & \textbf{h\#9.5} & h\#12.5 & 6 &  & \pageref{hm:KBvkqrn} \\
KBvkrbn & 6 & \textbf{h\#9.5} & h\#12.5 & 6 &  & \pageref{hm:KBvkrbn} \\
KBvkrn & 5 & \textbf{h\#9.5} & h\#12.5 & 6 &  & \pageref{hm:KBvkrn} \\
KBvkrnn & 6 & \textbf{h\#9.5} & h\#13 & 7 &  & \pageref{hm:KBvkrnn} \\
KBvkrrn & 6 & \textbf{h\#9.5} & h\#12.5 & 6 &  & \pageref{hm:KBvkrrn} \\
KNPvkqp & 6 & \textbf{h\#9.5} & h\#14 & 9 &  & \pageref{hm:KNPvkqp} \\
KNPvkrp & 6 & \textbf{h\#9.5} & h\#12.5 & 6 &  & \pageref{hm:KNPvkrp} \\
KNvkpp & 5 & \textbf{h\#9.5} & h\#11.5 & 4 &  & \pageref{hm:KNvkpp} \\
KNvkppp & 6 & \textbf{h\#9.5} & h\#11.5 & 4 &  & \pageref{hm:KNvkppp} \\
KNvkqb & 5 & \textbf{h\#9.5} & h\#11.5 & 4 & Sheglow (1998) & \pageref{hm:KNvkqb} \\
KNvkqbb & 6 & \textbf{h\#9.5} & h\#11.5 & 4 &  & \pageref{hm:KNvkqbb} \\
KNvkqbn & 6 & \textbf{h\#9.5} & h\#11.5 & 4 &  & \pageref{hm:KNvkqbn} \\
KNvkqbp & 6 & \textbf{h\#9.5} & h\#10.5 & 2 &  & \pageref{hm:KNvkqbp} \\
KNvkqnp & 6 & \textbf{h\#9.5} & h\#11.5 & 4 &  & \pageref{hm:KNvkqnp} \\
KNvkqp & 5 & \textbf{h\#9.5} & h\#11.5 & 4 &  & \pageref{hm:KNvkqp} \\
KNvkqpp & 6 & \textbf{h\#9.5} & h\#11.5 & 4 &  & \pageref{hm:KNvkqpp} \\
KNvkqqb & 6 & \textbf{h\#9.5} & h\#11.5 & 4 &  & \pageref{hm:KNvkqqb} \\
KNvkqrb & 6 & \textbf{h\#9.5} & h\#11.5 & 4 &  & \pageref{hm:KNvkqrb} \\
KNvkqrp & 6 & \textbf{h\#9.5} & h\#11.5 & 4 &  & \pageref{hm:KNvkqrp} \\
KNvkrpp & 6 & \textbf{h\#9.5} & h\#11.5 & 4 &  & \pageref{hm:KNvkrpp} \\
KPvkbnp & 6 & \textbf{h\#9.5} & h\#14.5 & 10 &  & \pageref{hm:KPvkbnp} \\
KPvkbp & 5 & \textbf{h\#9.5} & h\#14.5 & 10 & Locker (1966) & \pageref{hm:KPvkbp} \\
KPvknnp & 6 & \textbf{h\#9.5} & h\#11.5 & 4 &  & \pageref{hm:KPvknnp} \\
KPvknpp & 6 & \textbf{h\#9.5} & h\#15.5 & 12 &  & \pageref{hm:KPvknpp} \\
KPvkqbp & 6 & \textbf{h\#9.5} & h\#14.5 & 10 &  & \pageref{hm:KPvkqbp} \\
KPvkqnp & 6 & \textbf{h\#9.5} & h\#11.5 & 4 &  & \pageref{hm:KPvkqnp} \\
KPvkqpp & 6 & \textbf{h\#9.5} & h\#15.5 & 12 &  & \pageref{hm:KPvkqpp} \\
KPvkrbp & 6 & \textbf{h\#9.5} & h\#15 & 11 &  & \pageref{hm:KPvkrbp} \\
KPvkrpp & 6 & \textbf{h\#9.5} & h\#15.5 & 12 &  & \pageref{hm:KPvkrpp} \\
KBBBvkn & 6 & \textbf{h\#9} & h\#12 & 6 &  & \pageref{hm:KBBBvkn} \\
KBBvkbn & 6 & \textbf{h\#9} & h\#12 & 6 &  & \pageref{hm:KBBvkbn} \\
KBBvkn & 5 & \textbf{h\#9} & h\#12 & 6 &  & \pageref{hm:KBBvkn} \\
KBBvkp & 5 & \textbf{h\#9} & h\#16 & 14 &  & \pageref{hm:KBBvkp} \\
KBBvkrn & 6 & \textbf{h\#9} & h\#13 & 8 &  & \pageref{hm:KBBvkrn} \\
KBNvkqp & 6 & \textbf{h\#9} & h\#17 & 16 &  & \pageref{hm:KBNvkqp} \\
KBPvkbn & 6 & \textbf{h\#9} & h\#12 & 6 &  & \pageref{hm:KBPvkbn} \\
KBPvkqn & 6 & \textbf{h\#9} & h\#13 & 8 &  & \pageref{hm:KBPvkqn} \\
KBPvkrn & 6 & \textbf{h\#9} & h\#13 & 8 &  & \pageref{hm:KBPvkrn} \\
KBvkbn & 5 & \textbf{h\#9} & h\#12 & 6 &  & \pageref{hm:KBvkbn} \\
KBvknn & 5 & \textbf{h\#9} & h\#12.5 & 7 &  & \pageref{hm:KBvknn} \\
KBvkpp & 5 & \textbf{h\#9} & h\#16 & 14 &  & \pageref{hm:KBvkpp} \\
KNNvkqb & 6 & \textbf{h\#9} & h\#11 & 4 &  & \pageref{hm:KNNvkqb} \\
KNPPvkp & 6 & \textbf{h\#9} & h\#14 & 10 &  & \pageref{hm:KNPPvkp} \\
KNPvkp & 5 & \textbf{h\#9} & h\#14 & 10 &  & \pageref{hm:KNPvkp} \\
KNvkbpp & 6 & \textbf{h\#9} & h\#10.5 & 3 &  & \pageref{hm:KNvkbpp} \\
KNvknpp & 6 & \textbf{h\#9} & h\#10.5 & 3 &  & \pageref{hm:KNvknpp} \\
KPvknp & 5 & \textbf{h\#9} & h\#14.5 & 11 &  & \pageref{hm:KPvknp} \\
KQBBvkn & 6 & \textbf{h\#9} & h\#12 & 6 &  & \pageref{hm:KQBBvkn} \\
KQBBvkp & 6 & \textbf{h\#9} & h\#15 & 12 &  & \pageref{hm:KQBBvkp} \\
KQBvkbn & 6 & \textbf{h\#9} & h\#12.5 & 7 &  & \pageref{hm:KQBvkbn} \\
KQBvknn & 6 & \textbf{h\#9} & h\#13 & 8 &  & \pageref{hm:KQBvknn} \\
KQBvkpp & 6 & \textbf{h\#9} & h\#16 & 14 &  & \pageref{hm:KQBvkpp} \\
KQNPvkp & 6 & \textbf{h\#9} & h\#12 & 6 &  & \pageref{hm:KQNPvkp} \\
KQNvkbb & 6 & \textbf{h\#9} & h\#10 & 2 &  & \pageref{hm:KQNvkbb} \\
KQNvknp & 6 & \textbf{h\#9} & h\#10.5 & 3 &  & \pageref{hm:KQNvknp} \\
KQNvkqb & 6 & \textbf{h\#9} & h\#12 & 6 &  & \pageref{hm:KQNvkqb} \\
KQNvkrp & 6 & \textbf{h\#9} & h\#10.5 & 3 &  & \pageref{hm:KQNvkrp} \\
KQPvknp & 6 & \textbf{h\#9} & h\#15.5 & 13 &  & \pageref{hm:KQPvknp} \\
KQPvkrp & 6 & \textbf{h\#9} & h\#14 & 10 &  & \pageref{hm:KQPvkrp} \\
KRBBvkp & 6 & \textbf{h\#9} & h\#16 & 14 &  & \pageref{hm:KRBBvkp} \\
KRBvkbn & 6 & \textbf{h\#9} & h\#12 & 6 &  & \pageref{hm:KRBvkbn} \\
KRBvknn & 6 & \textbf{h\#9} & h\#12 & 6 &  & \pageref{hm:KRBvknn} \\
KRBvkpp & 6 & \textbf{h\#9} & h\#16 & 14 &  & \pageref{hm:KRBvkpp} \\
KRNPvkp & 6 & \textbf{h\#9} & h\#12 & 6 &  & \pageref{hm:KRNPvkp} \\
KRNvkqb & 6 & \textbf{h\#9} & h\#11 & 4 &  & \pageref{hm:KRNvkqb} \\
KRNvkqp & 6 & \textbf{h\#9} & h\#12 & 6 &  & \pageref{hm:KRNvkqp} \\
KRPvkbp & 6 & \textbf{h\#9} & h\#15 & 12 &  & \pageref{hm:KRPvkbp} \\
KRPvkpp & 6 & \textbf{h\#9} & h\#15 & 12 &  & \pageref{hm:KRPvkpp} \\
KBBvknn & 6 & \textbf{h\#8.5} & h\#13 & 9 &  & \pageref{hm:KBBvknn} \\
KBvkbbb & 6 & \textbf{h\#8.5} & h\#10 & 3 &  & \pageref{hm:KBvkbbb} \\
KBvkqbb & 6 & \textbf{h\#8.5} & h\#10.5 & 4 &  & \pageref{hm:KBvkqbb} \\
KBvkrbb & 6 & \textbf{h\#8.5} & h\#9.5 & 2 &  & \pageref{hm:KBvkrbb} \\
KNPvkbp & 6 & \textbf{h\#8.5} & h\#14 & 11 &  & \pageref{hm:KNPvkbp} \\
KNPvknp & 6 & \textbf{h\#8.5} & h\#10.5 & 4 &  & \pageref{hm:KNPvknp} \\
KNvkbb & 5 & \textbf{h\#8.5} & h\#10.5 & 4 &  & \pageref{hm:KNvkbb} \\
KNvkbbb & 6 & \textbf{h\#8.5} & h\#10.5 & 4 &  & \pageref{hm:KNvkbbb} \\
KNvkbbn & 6 & \textbf{h\#8.5} & h\#10.5 & 4 &  & \pageref{hm:KNvkbbn} \\
KNvkbbp & 6 & \textbf{h\#8.5} & h\#10.5 & 4 &  & \pageref{hm:KNvkbbp} \\
KNvkbnn & 6 & \textbf{h\#8.5} & h\#9.5 & 2 &  & \pageref{hm:KNvkbnn} \\
KNvkbnp & 6 & \textbf{h\#8.5} & h\#10 & 3 &  & \pageref{hm:KNvkbnp} \\
KNvkbp & 5 & \textbf{h\#8.5} & h\#10.5 & 4 &  & \pageref{hm:KNvkbp} \\
KNvknn & 5 & \textbf{h\#8.5} & h\#9.5 & 2 &  & \pageref{hm:KNvknn} \\
KNvknnn & 6 & \textbf{h\#8.5} & h\#9.5 & 2 &  & \pageref{hm:KNvknnn} \\
KNvknnp & 6 & \textbf{h\#8.5} & h\#10 & 3 &  & \pageref{hm:KNvknnp} \\
KNvknp & 5 & \textbf{h\#8.5} & h\#10 & 3 & Sheglow (1998) & \pageref{hm:KNvknp} \\
KNvkqnn & 6 & \textbf{h\#8.5} & h\#9.5 & 2 &  & \pageref{hm:KNvkqnn} \\
KNvkqqn & 6 & \textbf{h\#8.5} & h\#9.5 & 2 &  & \pageref{hm:KNvkqqn} \\
KNvkqqr & 6 & \textbf{h\#8.5} & h\#9.5 & 2 &  & \pageref{hm:KNvkqqr} \\
KNvkqr & 5 & \textbf{h\#8.5} & h\#9.5 & 2 &  & \pageref{hm:KNvkqr} \\
KNvkqrn & 6 & \textbf{h\#8.5} & h\#9.5 & 2 &  & \pageref{hm:KNvkqrn} \\
KNvkrbb & 6 & \textbf{h\#8.5} & h\#10.5 & 4 &  & \pageref{hm:KNvkrbb} \\
KNvkrbn & 6 & \textbf{h\#8.5} & h\#9.5 & 2 &  & \pageref{hm:KNvkrbn} \\
KNvkrbp & 6 & \textbf{h\#8.5} & h\#10.5 & 4 &  & \pageref{hm:KNvkrbp} \\
KNvkrn & 5 & \textbf{h\#8.5} & h\#10.5 & 4 &  & \pageref{hm:KNvkrn} \\
KNvkrnn & 6 & \textbf{h\#8.5} & h\#9.5 & 2 &  & \pageref{hm:KNvkrnn} \\
KNvkrnp & 6 & \textbf{h\#8.5} & h\#10.5 & 4 &  & \pageref{hm:KNvkrnp} \\
KNvkrp & 5 & \textbf{h\#8.5} & h\#11.5 & 6 &  & \pageref{hm:KNvkrp} \\
KNvkrr & 5 & \textbf{h\#8.5} & h\#9.5 & 2 &  & \pageref{hm:KNvkrr} \\
KNvkrrn & 6 & \textbf{h\#8.5} & h\#9.5 & 2 &  & \pageref{hm:KNvkrrn} \\
KNvkrrp & 6 & \textbf{h\#8.5} & h\#9.5 & 2 &  & \pageref{hm:KNvkrrp} \\
KPPPvkp & 6 & \textbf{h\#8.5} & h\#13.5 & 10 &  & \pageref{hm:KPPPvkp} \\
KPPvkbp & 6 & \textbf{h\#8.5} & h\#14.5 & 12 &  & \pageref{hm:KPPvkbp} \\
KPPvknp & 6 & \textbf{h\#8.5} & h\#14.5 & 12 &  & \pageref{hm:KPPvknp} \\
KPPvkp & 5 & \textbf{h\#8.5} & h\#14.5 & 12 &  & \pageref{hm:KPPvkp} \\
KPPvkqp & 6 & \textbf{h\#8.5} & h\#14.5 & 12 &  & \pageref{hm:KPPvkqp} \\
KPvkbbn & 6 & \textbf{h\#8.5} & h\#9.5 & 2 &  & \pageref{hm:KPvkbbn} \\
KPvkbnn & 6 & \textbf{h\#8.5} & h\#9.5 & 2 &  & \pageref{hm:KPvkbnn} \\
KPvknnn & 6 & \textbf{h\#8.5} & h\#9.5 & 2 &  & \pageref{hm:KPvknnn} \\
KPvkp & 4 & \textbf{h\#8.5} & h\#14.5 & 12 &  & \pageref{hm:KPvkp} \\
KPvkqbn & 6 & \textbf{h\#8.5} & h\#11.5 & 6 &  & \pageref{hm:KPvkqbn} \\
KPvkqn & 5 & \textbf{h\#8.5} & h\#9.5 & 2 &  & \pageref{hm:KPvkqn} \\
KPvkqnn & 6 & \textbf{h\#8.5} & h\#10.5 & 4 &  & \pageref{hm:KPvkqnn} \\
KPvkqp & 5 & \textbf{h\#8.5} & h\#14.5 & 12 &  & \pageref{hm:KPvkqp} \\
KPvkqqn & 6 & \textbf{h\#8.5} & h\#9.5 & 2 &  & \pageref{hm:KPvkqqn} \\
KPvkqqp & 6 & \textbf{h\#8.5} & h\#10.5 & 4 &  & \pageref{hm:KPvkqqp} \\
KPvkqrn & 6 & \textbf{h\#8.5} & h\#9.5 & 2 &  & \pageref{hm:KPvkqrn} \\
KPvkqrp & 6 & \textbf{h\#8.5} & h\#12.5 & 8 &  & \pageref{hm:KPvkqrp} \\
KPvkrnp & 6 & \textbf{h\#8.5} & h\#12.5 & 8 &  & \pageref{hm:KPvkrnp} \\
KPvkrp & 5 & \textbf{h\#8.5} & h\#14.5 & 12 &  & \pageref{hm:KPvkrp} \\
KPvkrrp & 6 & \textbf{h\#8.5} & h\#10.5 & 4 &  & \pageref{hm:KPvkrrp} \\
KBBBPvk & 6 & \textbf{h\#8} & h\#10 & 4 &  & \pageref{hm:KBBBPvk} \\
KBBPvkn & 6 & \textbf{h\#8} & h\#12 & 8 &  & \pageref{hm:KBBPvkn} \\
KBBvkbb & 6 & \textbf{h\#8} & h\#10 & 4 &  & \pageref{hm:KBBvkbb} \\
KBNPvkq & 6 & \textbf{h\#8} & h\#9 & 2 &  & \pageref{hm:KBNPvkq} \\
KBNvkbp & 6 & \textbf{h\#8} & h\#16 & 16 &  & \pageref{hm:KBNvkbp} \\
KBNvkqb & 6 & \textbf{h\#8} & h\#11 & 6 &  & \pageref{hm:KBNvkqb} \\
KBNvkrp & 6 & \textbf{h\#8} & h\#16.5 & 17 &  & \pageref{hm:KBNvkrp} \\
KBPvknn & 6 & \textbf{h\#8} & h\#13 & 10 &  & \pageref{hm:KBPvknn} \\
KBvkbb & 5 & \textbf{h\#8} & h\#10 & 4 &  & \pageref{hm:KBvkbb} \\
KBvkn & 4 & \textbf{h\#8} & h\#12 & 8 &  & \pageref{hm:KBvkn} \\
KBvkp & 4 & \textbf{h\#8} & h\#16 & 16 &  & \pageref{hm:KBvkp} \\
KNNPvkp & 6 & \textbf{h\#8} & h\#11 & 6 &  & \pageref{hm:KNNPvkp} \\
KNNvkpp & 6 & \textbf{h\#8} & h\#11 & 6 &  & \pageref{hm:KNNvkpp} \\
KNNvkqp & 6 & \textbf{h\#8} & h\#12 & 8 &  & \pageref{hm:KNNvkqp} \\
KNPvkqb & 6 & \textbf{h\#8} & h\#11.5 & 7 &  & \pageref{hm:KNPvkqb} \\
KNvkb & 4 & \textbf{h\#8} & h\#10.5 & 5 &  & \pageref{hm:KNvkb} \\
KNvkbn & 5 & \textbf{h\#8} & h\#10 & 4 &  & \pageref{hm:KNvkbn} \\
KNvkp & 4 & \textbf{h\#8} & h\#11.5 & 7 &  & \pageref{hm:KNvkp} \\
KNvkqn & 5 & \textbf{h\#8} & h\#10.5 & 5 & Sheglow (1998) & \pageref{hm:KNvkqn} \\
KNvkqrr & 6 & \textbf{h\#8} & h\#9.5 & 3 &  & \pageref{hm:KNvkqrr} \\
KNvkr & 4 & \textbf{h\#8} & h\#9.5 & 3 & Elkies (1995) & \pageref{hm:KNvkr} \\
KPPvkrp & 6 & \textbf{h\#8} & h\#14.5 & 13 &  & \pageref{hm:KPPvkrp} \\
KPvkqbb & 6 & \textbf{h\#8} & h\#11.5 & 7 &  & \pageref{hm:KPvkqbb} \\
KPvkqqb & 6 & \textbf{h\#8} & h\#11.5 & 7 &  & \pageref{hm:KPvkqqb} \\
KPvkqqq & 6 & \textbf{h\#8} & h\#9.5 & 3 &  & \pageref{hm:KPvkqqq} \\
KQBvkbb & 6 & \textbf{h\#8} & h\#10.5 & 5 &  & \pageref{hm:KQBvkbb} \\
KQBvkn & 5 & \textbf{h\#8} & h\#12 & 8 &  & \pageref{hm:KQBvkn} \\
KQNPvkq & 6 & \textbf{h\#8} & h\#9 & 2 &  & \pageref{hm:KQNPvkq} \\
KQNvkb & 5 & \textbf{h\#8} & h\#11 & 6 &  & \pageref{hm:KQNvkb} \\
KQNvkbn & 6 & \textbf{h\#8} & h\#10 & 4 &  & \pageref{hm:KQNvkbn} \\
KQNvkbp & 6 & \textbf{h\#8} & h\#10 & 4 &  & \pageref{hm:KQNvkbp} \\
KQNvkn & 5 & \textbf{h\#8} & h\#10 & 4 &  & \pageref{hm:KQNvkn} \\
KQNvknn & 6 & \textbf{h\#8} & h\#10 & 4 &  & \pageref{hm:KQNvknn} \\
KQNvkp & 5 & \textbf{h\#8} & h\#11 & 6 &  & \pageref{hm:KQNvkp} \\
KQNvkqn & 6 & \textbf{h\#8} & h\#9 & 2 &  & \pageref{hm:KQNvkqn} \\
KQNvkr & 5 & \textbf{h\#8} & h\#10 & 4 &  & \pageref{hm:KQNvkr} \\
KQNvkrn & 6 & \textbf{h\#8} & h\#9 & 2 &  & \pageref{hm:KQNvkrn} \\
KQNvkrr & 6 & \textbf{h\#8} & h\#9.5 & 3 &  & \pageref{hm:KQNvkrr} \\
KQPPvkp & 6 & \textbf{h\#8} & h\#15 & 14 &  & \pageref{hm:KQPPvkp} \\
KQPvkb & 5 & \textbf{h\#8} & h\#9 & 2 &  & \pageref{hm:KQPvkb} \\
KQPvkbb & 6 & \textbf{h\#8} & h\#9.5 & 3 &  & \pageref{hm:KQPvkbb} \\
KQPvkbn & 6 & \textbf{h\#8} & h\#9 & 2 &  & \pageref{hm:KQPvkbn} \\
KQPvkn & 5 & \textbf{h\#8} & h\#9 & 2 &  & \pageref{hm:KQPvkn} \\
KQPvknn & 6 & \textbf{h\#8} & h\#10 & 4 &  & \pageref{hm:KQPvknn} \\
KQPvkp & 5 & \textbf{h\#8} & h\#15 & 14 &  & \pageref{hm:KQPvkp} \\
KQPvkqb & 6 & \textbf{h\#8} & h\#12 & 8 &  & \pageref{hm:KQPvkqb} \\
KQPvkqn & 6 & \textbf{h\#8} & h\#10 & 4 &  & \pageref{hm:KQPvkqn} \\
KQPvkqp & 6 & \textbf{h\#8} & h\#13 & 10 &  & \pageref{hm:KQPvkqp} \\
KQPvkqq & 6 & \textbf{h\#8} & h\#9.5 & 3 &  & \pageref{hm:KQPvkqq} \\
KQPvkqr & 6 & \textbf{h\#8} & h\#9 & 2 &  & \pageref{hm:KQPvkqr} \\
KQPvkrb & 6 & \textbf{h\#8} & h\#9 & 2 &  & \pageref{hm:KQPvkrb} \\
KQPvkrn & 6 & \textbf{h\#8} & h\#9 & 2 &  & \pageref{hm:KQPvkrn} \\
KRBBvkn & 6 & \textbf{h\#8} & h\#12 & 8 &  & \pageref{hm:KRBBvkn} \\
KRBvkbb & 6 & \textbf{h\#8} & h\#10 & 4 &  & \pageref{hm:KRBvkbb} \\
KRBvkn & 5 & \textbf{h\#8} & h\#12 & 8 &  & \pageref{hm:KRBvkn} \\
KRNvkbb & 6 & \textbf{h\#8} & h\#10 & 4 &  & \pageref{hm:KRNvkbb} \\
KRNvkbp & 6 & \textbf{h\#8} & h\#10 & 4 &  & \pageref{hm:KRNvkbp} \\
KRNvkn & 5 & \textbf{h\#8} & h\#10 & 4 &  & \pageref{hm:KRNvkn} \\
KRNvknp & 6 & \textbf{h\#8} & h\#10.5 & 5 &  & \pageref{hm:KRNvknp} \\
KRNvkpp & 6 & \textbf{h\#8} & h\#11.5 & 7 &  & \pageref{hm:KRNvkpp} \\
KRNvkrp & 6 & \textbf{h\#8} & h\#11.5 & 7 &  & \pageref{hm:KRNvkrp} \\
KRPvkbb & 6 & \textbf{h\#8} & h\#9 & 2 &  & \pageref{hm:KRPvkbb} \\
KRPvkbn & 6 & \textbf{h\#8} & h\#9 & 2 &  & \pageref{hm:KRPvkbn} \\
KRPvknn & 6 & \textbf{h\#8} & h\#9 & 2 &  & \pageref{hm:KRPvknn} \\
KRPvknp & 6 & \textbf{h\#8} & h\#12.5 & 9 &  & \pageref{hm:KRPvknp} \\
KRPvkqb & 6 & \textbf{h\#8} & h\#12 & 8 &  & \pageref{hm:KRPvkqb} \\
KRPvkqn & 6 & \textbf{h\#8} & h\#9 & 2 &  & \pageref{hm:KRPvkqn} \\
KRPvkqp & 6 & \textbf{h\#8} & h\#14.5 & 13 &  & \pageref{hm:KRPvkqp} \\
KRPvkrp & 6 & \textbf{h\#8} & h\#12.5 & 9 &  & \pageref{hm:KRPvkrp} \\
KBBNvkq & 6 & \textbf{h\#7.5} & h\#8.5 & 2 &  & \pageref{hm:KBBNvkq} \\
KBBPvkq & 6 & \textbf{h\#7.5} & h\#9.5 & 4 &  & \pageref{hm:KBBPvkq} \\
KBBPvkr & 6 & \textbf{h\#7.5} & h\#9.5 & 4 &  & \pageref{hm:KBBPvkr} \\
KBBvkqb & 6 & \textbf{h\#7.5} & h\#10.5 & 6 &  & \pageref{hm:KBBvkqb} \\
KBNvkq & 5 & \textbf{h\#7.5} & h\#8.5 & 2 &  & \pageref{hm:KBNvkq} \\
KBNvkqn & 6 & \textbf{h\#7.5} & h\#12.5 & 10 &  & \pageref{hm:KBNvkqn} \\
KBNvkqq & 6 & \textbf{h\#7.5} & h\#8.5 & 2 &  & \pageref{hm:KBNvkqq} \\
KBNvkqr & 6 & \textbf{h\#7.5} & h\#8.5 & 2 &  & \pageref{hm:KBNvkqr} \\
KBNvkrr & 6 & \textbf{h\#7.5} & h\#8.5 & 2 &  & \pageref{hm:KBNvkrr} \\
KBPvkqb & 6 & \textbf{h\#7.5} & h\#11 & 7 &  & \pageref{hm:KBPvkqb} \\
KBPvkqq & 6 & \textbf{h\#7.5} & h\#9 & 3 &  & \pageref{hm:KBPvkqq} \\
KBPvkqr & 6 & \textbf{h\#7.5} & h\#9 & 3 &  & \pageref{hm:KBPvkqr} \\
KBvkqqb & 6 & \textbf{h\#7.5} & h\#8.5 & 2 &  & \pageref{hm:KBvkqqb} \\
KBvkqrb & 6 & \textbf{h\#7.5} & h\#8.5 & 2 &  & \pageref{hm:KBvkqrb} \\
KBvkrrb & 6 & \textbf{h\#7.5} & h\#8.5 & 2 &  & \pageref{hm:KBvkrrb} \\
KNNPvkq & 6 & \textbf{h\#7.5} & h\#9 & 3 &  & \pageref{hm:KNNPvkq} \\
KNPPvkq & 6 & \textbf{h\#7.5} & h\#8.5 & 2 &  & \pageref{hm:KNPPvkq} \\
KNPvkb & 5 & \textbf{h\#7.5} & h\#11 & 7 &  & \pageref{hm:KNPvkb} \\
KNPvkbb & 6 & \textbf{h\#7.5} & h\#10 & 5 &  & \pageref{hm:KNPvkbb} \\
KNPvkbn & 6 & \textbf{h\#7.5} & h\#9 & 3 &  & \pageref{hm:KNPvkbn} \\
KNPvkn & 5 & \textbf{h\#7.5} & h\#9 & 3 & Paliulionis (2024) & \pageref{hm:KNPvkn} \\
KNPvknn & 6 & \textbf{h\#7.5} & h\#9 & 3 &  & \pageref{hm:KNPvknn} \\
KNPvkq & 5 & \textbf{h\#7.5} & h\#8.5 & 2 &  & \pageref{hm:KNPvkq} \\
KNPvkqn & 6 & \textbf{h\#7.5} & h\#8.5 & 2 &  & \pageref{hm:KNPvkqn} \\
KNPvkqq & 6 & \textbf{h\#7.5} & h\#8.5 & 2 &  & \pageref{hm:KNPvkqq} \\
KNPvkqr & 6 & \textbf{h\#7.5} & h\#8.5 & 2 &  & \pageref{hm:KNPvkqr} \\
KNPvkr & 5 & \textbf{h\#7.5} & h\#9 & 3 &  & \pageref{hm:KNPvkr} \\
KNPvkrb & 6 & \textbf{h\#7.5} & h\#9.5 & 4 &  & \pageref{hm:KNPvkrb} \\
KNPvkrn & 6 & \textbf{h\#7.5} & h\#9 & 3 &  & \pageref{hm:KNPvkrn} \\
KNPvkrr & 6 & \textbf{h\#7.5} & h\#8.5 & 2 &  & \pageref{hm:KNPvkrr} \\
KNvkn & 4 & \textbf{h\#7.5} & h\#10 & 5 & Kotesovec (1987) & \pageref{hm:KNvkn} \\
KNvkrb & 5 & \textbf{h\#7.5} & h\#10.5 & 6 &  & \pageref{hm:KNvkrb} \\
KNvkrrb & 6 & \textbf{h\#7.5} & h\#8.5 & 2 &  & \pageref{hm:KNvkrrb} \\
KNvkrrr & 6 & \textbf{h\#7.5} & h\#9 & 3 &  & \pageref{hm:KNvkrrr} \\
KPPvkb & 5 & \textbf{h\#7.5} & h\#8.5 & 2 &  & \pageref{hm:KPPvkb} \\
KPPvkbb & 6 & \textbf{h\#7.5} & h\#8.5 & 2 &  & \pageref{hm:KPPvkbb} \\
KPPvkbn & 6 & \textbf{h\#7.5} & h\#8.5 & 2 &  & \pageref{hm:KPPvkbn} \\
KPPvkn & 5 & \textbf{h\#7.5} & h\#8.5 & 2 &  & \pageref{hm:KPPvkn} \\
KPPvknn & 6 & \textbf{h\#7.5} & h\#9 & 3 &  & \pageref{hm:KPPvknn} \\
KPPvkq & 5 & \textbf{h\#7.5} & h\#8.5 & 2 &  & \pageref{hm:KPPvkq} \\
KPPvkqb & 6 & \textbf{h\#7.5} & h\#11 & 7 &  & \pageref{hm:KPPvkqb} \\
KPPvkqn & 6 & \textbf{h\#7.5} & h\#9 & 3 &  & \pageref{hm:KPPvkqn} \\
KPPvkqq & 6 & \textbf{h\#7.5} & h\#9 & 3 &  & \pageref{hm:KPPvkqq} \\
KPPvkqr & 6 & \textbf{h\#7.5} & h\#8.5 & 2 &  & \pageref{hm:KPPvkqr} \\
KPPvkrb & 6 & \textbf{h\#7.5} & h\#8.5 & 2 &  & \pageref{hm:KPPvkrb} \\
KPPvkrn & 6 & \textbf{h\#7.5} & h\#8.5 & 2 &  & \pageref{hm:KPPvkrn} \\
KPvkb & 4 & \textbf{h\#7.5} & h\#8.5 & 2 &  & \pageref{hm:KPvkb} \\
KPvkbb & 5 & \textbf{h\#7.5} & h\#9.5 & 4 &  & \pageref{hm:KPvkbb} \\
KPvkbbb & 6 & \textbf{h\#7.5} & h\#9.5 & 4 &  & \pageref{hm:KPvkbbb} \\
KPvkbn & 5 & \textbf{h\#7.5} & h\#9.5 & 4 &  & \pageref{hm:KPvkbn} \\
KPvkn & 4 & \textbf{h\#7.5} & h\#8.5 & 2 &  & \pageref{hm:KPvkn} \\
KPvknn & 5 & \textbf{h\#7.5} & h\#9.5 & 4 &  & \pageref{hm:KPvknn} \\
KPvkq & 4 & \textbf{h\#7.5} & h\#9 & 3 &  & \pageref{hm:KPvkq} \\
KPvkqb & 5 & \textbf{h\#7.5} & h\#11.5 & 8 &  & \pageref{hm:KPvkqb} \\
KPvkqq & 5 & \textbf{h\#7.5} & h\#9.5 & 4 &  & \pageref{hm:KPvkqq} \\
KPvkqqr & 6 & \textbf{h\#7.5} & h\#9.5 & 4 &  & \pageref{hm:KPvkqqr} \\
KPvkqr & 5 & \textbf{h\#7.5} & h\#8.5 & 2 &  & \pageref{hm:KPvkqr} \\
KPvkqrb & 6 & \textbf{h\#7.5} & h\#11.5 & 8 &  & \pageref{hm:KPvkqrb} \\
KPvkqrr & 6 & \textbf{h\#7.5} & h\#8.5 & 2 &  & \pageref{hm:KPvkqrr} \\
KPvkr & 4 & \textbf{h\#7.5} & h\#8.5 & 2 & Maslar & \pageref{hm:KPvkr} \\
KPvkrb & 5 & \textbf{h\#7.5} & h\#8.5 & 2 &  & \pageref{hm:KPvkrb} \\
KPvkrbb & 6 & \textbf{h\#7.5} & h\#9.5 & 4 &  & \pageref{hm:KPvkrbb} \\
KPvkrbn & 6 & \textbf{h\#7.5} & h\#9.5 & 4 &  & \pageref{hm:KPvkrbn} \\
KPvkrn & 5 & \textbf{h\#7.5} & h\#9.5 & 4 &  & \pageref{hm:KPvkrn} \\
KPvkrnn & 6 & \textbf{h\#7.5} & h\#9.5 & 4 &  & \pageref{hm:KPvkrnn} \\
KPvkrrb & 6 & \textbf{h\#7.5} & h\#8.5 & 2 &  & \pageref{hm:KPvkrrb} \\
KPvkrrn & 6 & \textbf{h\#7.5} & h\#9.5 & 4 &  & \pageref{hm:KPvkrrn} \\
KPvkrrr & 6 & \textbf{h\#7.5} & h\#8.5 & 2 &  & \pageref{hm:KPvkrrr} \\
KQNPvkn & 6 & \textbf{h\#7.5} & h\#9 & 3 &  & \pageref{hm:KQNPvkn} \\
KQNvkqr & 6 & \textbf{h\#7.5} & h\#9.5 & 4 &  & \pageref{hm:KQNvkqr} \\
KRNPvkq & 6 & \textbf{h\#7.5} & h\#8.5 & 2 &  & \pageref{hm:KRNPvkq} \\
KRNPvkr & 6 & \textbf{h\#7.5} & h\#8.5 & 2 &  & \pageref{hm:KRNPvkr} \\
KRPPvkq & 6 & \textbf{h\#7.5} & h\#8.5 & 2 &  & \pageref{hm:KRPPvkq} \\
KBBBBvk & 6 & \textbf{h\#7} & h\#8 & 2 &  & \pageref{hm:KBBBBvk} \\
KBBBNvk & 6 & \textbf{h\#7} & h\#8 & 2 &  & \pageref{hm:KBBBNvk} \\
KBBBvk & 5 & \textbf{h\#7} & h\#8 & 2 &  & \pageref{hm:KBBBvk} \\
KBBBvkb & 6 & \textbf{h\#7} & h\#11 & 8 &  & \pageref{hm:KBBBvkb} \\
KBBNPvk & 6 & \textbf{h\#7} & h\#9 & 4 &  & \pageref{hm:KBBNPvk} \\
KBBNvk & 5 & \textbf{h\#7} & h\#8 & 2 &  & \pageref{hm:KBBNvk} \\
KBBNvkn & 6 & \textbf{h\#7} & h\#12 & 10 &  & \pageref{hm:KBBNvkn} \\
KBBNvkp & 6 & \textbf{h\#7} & h\#15 & 16 &  & \pageref{hm:KBBNvkp} \\
KBBPPvk & 6 & \textbf{h\#7} & h\#9 & 4 &  & \pageref{hm:KBBPPvk} \\
KBBPvk & 5 & \textbf{h\#7} & h\#9 & 4 &  & \pageref{hm:KBBPvk} \\
KBBPvkb & 6 & \textbf{h\#7} & h\#10 & 6 &  & \pageref{hm:KBBPvkb} \\
KBBvkb & 5 & \textbf{h\#7} & h\#10.5 & 7 &  & \pageref{hm:KBBvkb} \\
KBBvkrb & 6 & \textbf{h\#7} & h\#10.5 & 7 &  & \pageref{hm:KBBvkrb} \\
KBNNPvk & 6 & \textbf{h\#7} & h\#8 & 2 &  & \pageref{hm:KBNNPvk} \\
KBNNvkp & 6 & \textbf{h\#7} & h\#9 & 4 &  & \pageref{hm:KBNNvkp} \\
KBNNvkq & 6 & \textbf{h\#7} & h\#8.5 & 3 &  & \pageref{hm:KBNNvkq} \\
KBNPPvk & 6 & \textbf{h\#7} & h\#8 & 2 &  & \pageref{hm:KBNPPvk} \\
KBNPvk & 5 & \textbf{h\#7} & h\#8 & 2 &  & \pageref{hm:KBNPvk} \\
KBNPvkn & 6 & \textbf{h\#7} & h\#8 & 2 &  & \pageref{hm:KBNPvkn} \\
KBNPvkr & 6 & \textbf{h\#7} & h\#9 & 4 &  & \pageref{hm:KBNPvkr} \\
KBNvkb & 5 & \textbf{h\#7} & h\#11 & 8 &  & \pageref{hm:KBNvkb} \\
KBNvkbb & 6 & \textbf{h\#7} & h\#10 & 6 &  & \pageref{hm:KBNvkbb} \\
KBNvkbn & 6 & \textbf{h\#7} & h\#12 & 10 &  & \pageref{hm:KBNvkbn} \\
KBNvkn & 5 & \textbf{h\#7} & h\#12 & 10 &  & \pageref{hm:KBNvkn} \\
KBNvknn & 6 & \textbf{h\#7} & h\#11 & 8 &  & \pageref{hm:KBNvknn} \\
KBNvknp & 6 & \textbf{h\#7} & h\#12 & 10 &  & \pageref{hm:KBNvknp} \\
KBNvkp & 5 & \textbf{h\#7} & h\#16 & 18 &  & \pageref{hm:KBNvkp} \\
KBNvkpp & 6 & \textbf{h\#7} & h\#16 & 18 &  & \pageref{hm:KBNvkpp} \\
KBNvkr & 5 & \textbf{h\#7} & h\#9 & 4 &  & \pageref{hm:KBNvkr} \\
KBNvkrn & 6 & \textbf{h\#7} & h\#12 & 10 &  & \pageref{hm:KBNvkrn} \\
KBPPPvk & 6 & \textbf{h\#7} & h\#8 & 2 &  & \pageref{hm:KBPPPvk} \\
KBPPvk & 5 & \textbf{h\#7} & h\#8 & 2 &  & \pageref{hm:KBPPvk} \\
KBPPvkb & 6 & \textbf{h\#7} & h\#8 & 2 &  & \pageref{hm:KBPPvkb} \\
KBPPvkn & 6 & \textbf{h\#7} & h\#8 & 2 &  & \pageref{hm:KBPPvkn} \\
KBPPvkq & 6 & \textbf{h\#7} & h\#8.5 & 3 &  & \pageref{hm:KBPPvkq} \\
KBPPvkr & 6 & \textbf{h\#7} & h\#8.5 & 3 &  & \pageref{hm:KBPPvkr} \\
KBPvk & 4 & \textbf{h\#7} & h\#8 & 2 & Elkies (2004) & \pageref{hm:KBPvk} \\
KBPvkb & 5 & \textbf{h\#7} & h\#10 & 6 &  & \pageref{hm:KBPvkb} \\
KBPvkbb & 6 & \textbf{h\#7} & h\#10 & 6 &  & \pageref{hm:KBPvkbb} \\
KBPvkn & 5 & \textbf{h\#7} & h\#12 & 10 &  & \pageref{hm:KBPvkn} \\
KBPvkq & 5 & \textbf{h\#7} & h\#8.5 & 3 &  & \pageref{hm:KBPvkq} \\
KBPvkr & 5 & \textbf{h\#7} & h\#8.5 & 3 &  & \pageref{hm:KBPvkr} \\
KBPvkrb & 6 & \textbf{h\#7} & h\#9.5 & 5 &  & \pageref{hm:KBPvkrb} \\
KBPvkrr & 6 & \textbf{h\#7} & h\#8.5 & 3 &  & \pageref{hm:KBPvkrr} \\
KNNNPvk & 6 & \textbf{h\#7} & h\#8 & 2 &  & \pageref{hm:KNNNPvk} \\
KNNNvkp & 6 & \textbf{h\#7} & h\#8 & 2 &  & \pageref{hm:KNNNvkp} \\
KNNPPvk & 6 & \textbf{h\#7} & h\#8 & 2 &  & \pageref{hm:KNNPPvk} \\
KNNPvk & 5 & \textbf{h\#7} & h\#8 & 2 &  & \pageref{hm:KNNPvk} \\
KNNPvkb & 6 & \textbf{h\#7} & h\#8 & 2 &  & \pageref{hm:KNNPvkb} \\
KNNvkbp & 6 & \textbf{h\#7} & h\#9 & 4 &  & \pageref{hm:KNNvkbp} \\
KNNvkp & 5 & \textbf{h\#7} & h\#11 & 8 &  & \pageref{hm:KNNvkp} \\
KNNvkq & 5 & \textbf{h\#7} & h\#8.5 & 3 & Sheglow (1998) & \pageref{hm:KNNvkq} \\
KNNvkqr & 6 & \textbf{h\#7} & h\#8.5 & 3 &  & \pageref{hm:KNNvkqr} \\
KNNvkr & 5 & \textbf{h\#7} & h\#10 & 6 &  & \pageref{hm:KNNvkr} \\
KNNvkrp & 6 & \textbf{h\#7} & h\#12 & 10 &  & \pageref{hm:KNNvkrp} \\
KNPPvk & 5 & \textbf{h\#7} & h\#8.5 & 3 &  & \pageref{hm:KNPPvk} \\
KNPPvkb & 6 & \textbf{h\#7} & h\#8.5 & 3 &  & \pageref{hm:KNPPvkb} \\
KNPPvkn & 6 & \textbf{h\#7} & h\#8.5 & 3 &  & \pageref{hm:KNPPvkn} \\
KNPPvkr & 6 & \textbf{h\#7} & h\#8.5 & 3 &  & \pageref{hm:KNPPvkr} \\
KNPvk & 4 & \textbf{h\#7} & h\#8.5 & 3 & Elkies (2004) & \pageref{hm:KNPvk} \\
KPPPvk & 5 & \textbf{h\#7} & h\#8 & 2 &  & \pageref{hm:KPPPvk} \\
KPPPvkb & 6 & \textbf{h\#7} & h\#8 & 2 &  & \pageref{hm:KPPPvkb} \\
KPPPvkn & 6 & \textbf{h\#7} & h\#8.5 & 3 &  & \pageref{hm:KPPPvkn} \\
KPPvk & 4 & \textbf{h\#7} & h\#8 & 2 & Salazar Lopez (1973) & \pageref{hm:KPPvk} \\
KPvkrr & 5 & \textbf{h\#7} & h\#8.5 & 3 &  & \pageref{hm:KPvkrr} \\
KQBBBvk & 6 & \textbf{h\#7} & h\#8 & 2 &  & \pageref{hm:KQBBBvk} \\
KQBBNvk & 6 & \textbf{h\#7} & h\#8 & 2 &  & \pageref{hm:KQBBNvk} \\
KQBBPvk & 6 & \textbf{h\#7} & h\#9 & 4 &  & \pageref{hm:KQBBPvk} \\
KQBBvk & 5 & \textbf{h\#7} & h\#8 & 2 &  & \pageref{hm:KQBBvk} \\
KQBBvkb & 6 & \textbf{h\#7} & h\#11 & 8 &  & \pageref{hm:KQBBvkb} \\
KQBNNvk & 6 & \textbf{h\#7} & h\#8 & 2 &  & \pageref{hm:KQBNNvk} \\
KQBNPvk & 6 & \textbf{h\#7} & h\#8 & 2 &  & \pageref{hm:KQBNPvk} \\
KQBNvkb & 6 & \textbf{h\#7} & h\#9 & 4 &  & \pageref{hm:KQBNvkb} \\
KQBNvkn & 6 & \textbf{h\#7} & h\#9 & 4 &  & \pageref{hm:KQBNvkn} \\
KQBNvkp & 6 & \textbf{h\#7} & h\#12 & 10 &  & \pageref{hm:KQBNvkp} \\
KQBNvkq & 6 & \textbf{h\#7} & h\#9 & 4 &  & \pageref{hm:KQBNvkq} \\
KQBNvkr & 6 & \textbf{h\#7} & h\#9 & 4 &  & \pageref{hm:KQBNvkr} \\
KQBPPvk & 6 & \textbf{h\#7} & h\#8 & 2 &  & \pageref{hm:KQBPPvk} \\
KQBPvk & 5 & \textbf{h\#7} & h\#8 & 2 &  & \pageref{hm:KQBPvk} \\
KQBPvkb & 6 & \textbf{h\#7} & h\#9 & 4 &  & \pageref{hm:KQBPvkb} \\
KQBPvkn & 6 & \textbf{h\#7} & h\#9 & 4 &  & \pageref{hm:KQBPvkn} \\
KQBPvkq & 6 & \textbf{h\#7} & h\#8 & 2 &  & \pageref{hm:KQBPvkq} \\
KQBPvkr & 6 & \textbf{h\#7} & h\#8 & 2 &  & \pageref{hm:KQBPvkr} \\
KQBvkb & 5 & \textbf{h\#7} & h\#10 & 6 &  & \pageref{hm:KQBvkb} \\
KQBvkp & 5 & \textbf{h\#7} & h\#16 & 18 &  & \pageref{hm:KQBvkp} \\
KQBvkqb & 6 & \textbf{h\#7} & h\#10.5 & 7 &  & \pageref{hm:KQBvkqb} \\
KQBvkrb & 6 & \textbf{h\#7} & h\#10 & 6 &  & \pageref{hm:KQBvkrb} \\
KQNNPvk & 6 & \textbf{h\#7} & h\#8 & 2 &  & \pageref{hm:KQNNPvk} \\
KQNNvk & 5 & \textbf{h\#7} & h\#8 & 2 &  & \pageref{hm:KQNNvk} \\
KQNNvkb & 6 & \textbf{h\#7} & h\#9 & 4 &  & \pageref{hm:KQNNvkb} \\
KQNNvkp & 6 & \textbf{h\#7} & h\#9 & 4 &  & \pageref{hm:KQNNvkp} \\
KQNNvkq & 6 & \textbf{h\#7} & h\#9 & 4 &  & \pageref{hm:KQNNvkq} \\
KQNNvkr & 6 & \textbf{h\#7} & h\#8 & 2 &  & \pageref{hm:KQNNvkr} \\
KQNPPvk & 6 & \textbf{h\#7} & h\#8 & 2 &  & \pageref{hm:KQNPPvk} \\
KQNPvk & 5 & \textbf{h\#7} & h\#9 & 4 &  & \pageref{hm:KQNPvk} \\
KQNPvkb & 6 & \textbf{h\#7} & h\#9 & 4 &  & \pageref{hm:KQNPvkb} \\
KQNPvkr & 6 & \textbf{h\#7} & h\#8.5 & 3 &  & \pageref{hm:KQNPvkr} \\
KQNvkrb & 6 & \textbf{h\#7} & h\#10 & 6 &  & \pageref{hm:KQNvkrb} \\
KQPPPvk & 6 & \textbf{h\#7} & h\#8 & 2 &  & \pageref{hm:KQPPPvk} \\
KQPPvk & 5 & \textbf{h\#7} & h\#8 & 2 &  & \pageref{hm:KQPPvk} \\
KQPPvkb & 6 & \textbf{h\#7} & h\#9 & 4 &  & \pageref{hm:KQPPvkb} \\
KQPPvkn & 6 & \textbf{h\#7} & h\#9 & 4 &  & \pageref{hm:KQPPvkn} \\
KQPPvkq & 6 & \textbf{h\#7} & h\#9 & 4 &  & \pageref{hm:KQPPvkq} \\
KQPPvkr & 6 & \textbf{h\#7} & h\#8 & 2 &  & \pageref{hm:KQPPvkr} \\
KQPvk & 4 & \textbf{h\#7} & h\#8 & 2 &  & \pageref{hm:KQPvk} \\
KQPvkq & 5 & \textbf{h\#7} & h\#9 & 4 &  & \pageref{hm:KQPvkq} \\
KQPvkr & 5 & \textbf{h\#7} & h\#9 & 4 &  & \pageref{hm:KQPvkr} \\
KQPvkrr & 6 & \textbf{h\#7} & h\#9 & 4 &  & \pageref{hm:KQPvkrr} \\
KRBBBvk & 6 & \textbf{h\#7} & h\#8 & 2 &  & \pageref{hm:KRBBBvk} \\
KRBBPvk & 6 & \textbf{h\#7} & h\#9 & 4 &  & \pageref{hm:KRBBPvk} \\
KRBBvk & 5 & \textbf{h\#7} & h\#8 & 2 &  & \pageref{hm:KRBBvk} \\
KRBBvkb & 6 & \textbf{h\#7} & h\#10 & 6 &  & \pageref{hm:KRBBvkb} \\
KRBNNvk & 6 & \textbf{h\#7} & h\#8 & 2 &  & \pageref{hm:KRBNNvk} \\
KRBNPvk & 6 & \textbf{h\#7} & h\#8 & 2 &  & \pageref{hm:KRBNPvk} \\
KRBNvkb & 6 & \textbf{h\#7} & h\#8 & 2 &  & \pageref{hm:KRBNvkb} \\
KRBNvkn & 6 & \textbf{h\#7} & h\#8 & 2 &  & \pageref{hm:KRBNvkn} \\
KRBNvkp & 6 & \textbf{h\#7} & h\#13 & 12 &  & \pageref{hm:KRBNvkp} \\
KRBNvkq & 6 & \textbf{h\#7} & h\#8.5 & 3 &  & \pageref{hm:KRBNvkq} \\
KRBNvkr & 6 & \textbf{h\#7} & h\#8 & 2 &  & \pageref{hm:KRBNvkr} \\
KRBPPvk & 6 & \textbf{h\#7} & h\#8 & 2 &  & \pageref{hm:KRBPPvk} \\
KRBPvk & 5 & \textbf{h\#7} & h\#8 & 2 & Abdurahmanovic \& Becker (2006) & \pageref{hm:KRBPvk} \\
KRBPvkb & 6 & \textbf{h\#7} & h\#8 & 2 &  & \pageref{hm:KRBPvkb} \\
KRBPvkn & 6 & \textbf{h\#7} & h\#8.5 & 3 &  & \pageref{hm:KRBPvkn} \\
KRBPvkq & 6 & \textbf{h\#7} & h\#8.5 & 3 &  & \pageref{hm:KRBPvkq} \\
KRBPvkr & 6 & \textbf{h\#7} & h\#8.5 & 3 &  & \pageref{hm:KRBPvkr} \\
KRBvkrb & 6 & \textbf{h\#7} & h\#8.5 & 3 &  & \pageref{hm:KRBvkrb} \\
KRNNPvk & 6 & \textbf{h\#7} & h\#8 & 2 &  & \pageref{hm:KRNNPvk} \\
KRNNvkb & 6 & \textbf{h\#7} & h\#8 & 2 &  & \pageref{hm:KRNNvkb} \\
KRNNvkp & 6 & \textbf{h\#7} & h\#11 & 8 &  & \pageref{hm:KRNNvkp} \\
KRNNvkq & 6 & \textbf{h\#7} & h\#8 & 2 &  & \pageref{hm:KRNNvkq} \\
KRNPPvk & 6 & \textbf{h\#7} & h\#8 & 2 &  & \pageref{hm:KRNPPvk} \\
KRNPvk & 5 & \textbf{h\#7} & h\#8 & 2 &  & \pageref{hm:KRNPvk} \\
KRNPvkb & 6 & \textbf{h\#7} & h\#8 & 2 &  & \pageref{hm:KRNPvkb} \\
KRNPvkn & 6 & \textbf{h\#7} & h\#9 & 4 &  & \pageref{hm:KRNPvkn} \\
KRNvkbn & 6 & \textbf{h\#7} & h\#9.5 & 5 &  & \pageref{hm:KRNvkbn} \\
KRNvknn & 6 & \textbf{h\#7} & h\#9 & 4 &  & \pageref{hm:KRNvknn} \\
KRNvkp & 5 & \textbf{h\#7} & h\#11.5 & 9 &  & \pageref{hm:KRNvkp} \\
KRNvkqn & 6 & \textbf{h\#7} & h\#9 & 4 &  & \pageref{hm:KRNvkqn} \\
KRNvkqr & 6 & \textbf{h\#7} & h\#9.5 & 5 &  & \pageref{hm:KRNvkqr} \\
KRNvkr & 5 & \textbf{h\#7} & h\#10 & 6 &  & \pageref{hm:KRNvkr} \\
KRNvkrb & 6 & \textbf{h\#7} & h\#10 & 6 &  & \pageref{hm:KRNvkrb} \\
KRNvkrn & 6 & \textbf{h\#7} & h\#9 & 4 &  & \pageref{hm:KRNvkrn} \\
KRNvkrr & 6 & \textbf{h\#7} & h\#9.5 & 5 &  & \pageref{hm:KRNvkrr} \\
KRPPPvk & 6 & \textbf{h\#7} & h\#8 & 2 &  & \pageref{hm:KRPPPvk} \\
KRPPvk & 5 & \textbf{h\#7} & h\#8 & 2 &  & \pageref{hm:KRPPvk} \\
KRPPvkb & 6 & \textbf{h\#7} & h\#8 & 2 &  & \pageref{hm:KRPPvkb} \\
KRPPvkn & 6 & \textbf{h\#7} & h\#8.5 & 3 &  & \pageref{hm:KRPPvkn} \\
KRPPvkp & 6 & \textbf{h\#7} & h\#14 & 14 &  & \pageref{hm:KRPPvkp} \\
KRPPvkr & 6 & \textbf{h\#7} & h\#8.5 & 3 &  & \pageref{hm:KRPPvkr} \\
KRPvkb & 5 & \textbf{h\#7} & h\#8 & 2 &  & \pageref{hm:KRPvkb} \\
KRPvkn & 5 & \textbf{h\#7} & h\#8.5 & 3 &  & \pageref{hm:KRPvkn} \\
KRPvkp & 5 & \textbf{h\#7} & h\#14 & 14 &  & \pageref{hm:KRPvkp} \\
KRPvkq & 5 & \textbf{h\#7} & h\#9 & 4 &  & \pageref{hm:KRPvkq} \\
KRPvkqq & 6 & \textbf{h\#7} & h\#9 & 4 &  & \pageref{hm:KRPvkqq} \\
KRPvkqr & 6 & \textbf{h\#7} & h\#8.5 & 3 &  & \pageref{hm:KRPvkqr} \\
KRPvkr & 5 & \textbf{h\#7} & h\#8.5 & 3 &  & \pageref{hm:KRPvkr} \\
KRPvkrb & 6 & \textbf{h\#7} & h\#9 & 4 &  & \pageref{hm:KRPvkrb} \\
KRPvkrn & 6 & \textbf{h\#7} & h\#9 & 4 &  & \pageref{hm:KRPvkrn} \\
KRPvkrr & 6 & \textbf{h\#7} & h\#8.5 & 3 &  & \pageref{hm:KRPvkrr} \\
KRvkqbn & 6 & \textbf{h\#7} & h\#8 & 2 &  & \pageref{hm:KRvkqbn} \\
KRvkqnn & 6 & \textbf{h\#7} & h\#8 & 2 &  & \pageref{hm:KRvkqnn} \\
KRvkqnp & 6 & \textbf{h\#7} & h\#8 & 2 &  & \pageref{hm:KRvkqnp} \\
KRvkqqb & 6 & \textbf{h\#7} & h\#9 & 4 &  & \pageref{hm:KRvkqqb} \\
KRvkqqn & 6 & \textbf{h\#7} & h\#8 & 2 &  & \pageref{hm:KRvkqqn} \\
KRvkqqp & 6 & \textbf{h\#7} & h\#8 & 2 &  & \pageref{hm:KRvkqqp} \\
KRvkqqq & 6 & \textbf{h\#7} & h\#8.5 & 3 &  & \pageref{hm:KRvkqqq} \\
KRvkqqr & 6 & \textbf{h\#7} & h\#8.5 & 3 &  & \pageref{hm:KRvkqqr} \\
KRvkqrb & 6 & \textbf{h\#7} & h\#8 & 2 &  & \pageref{hm:KRvkqrb} \\
KRvkqrp & 6 & \textbf{h\#7} & h\#8 & 2 &  & \pageref{hm:KRvkqrp} \\
KRvkqrr & 6 & \textbf{h\#7} & h\#8 & 2 &  & \pageref{hm:KRvkqrr} \\
KRvkrpp & 6 & \textbf{h\#7} & h\#7.5 & 1 &  & \pageref{hm:KRvkrpp} \\
KRvkrrp & 6 & \textbf{h\#7} & h\#8 & 2 &  & \pageref{hm:KRvkrrp} \\
KBBBvkq & 6 & \textbf{h\#6.5} & h\#8 & 3 &  & \pageref{hm:KBBBvkq} \\
KBBNNvk & 6 & \textbf{h\#6.5} & h\#8 & 3 &  & \pageref{hm:KBBNNvk} \\
KBBNvkb & 6 & \textbf{h\#6.5} & h\#10 & 7 &  & \pageref{hm:KBBNvkb} \\
KBBvk & 4 & \textbf{h\#6.5} & h\#8 & 3 & Maslar (1978) & \pageref{hm:KBBvk} \\
KBBvkq & 5 & \textbf{h\#6.5} & h\#7.5 & 2 &  & \pageref{hm:KBBvkq} \\
KBBvkqq & 6 & \textbf{h\#6.5} & h\#7.5 & 2 &  & \pageref{hm:KBBvkqq} \\
KBBvkqr & 6 & \textbf{h\#6.5} & h\#7.5 & 2 &  & \pageref{hm:KBBvkqr} \\
KBBvkrr & 6 & \textbf{h\#6.5} & h\#7.5 & 2 &  & \pageref{hm:KBBvkrr} \\
KBNNvk & 5 & \textbf{h\#6.5} & h\#8 & 3 & Sheglow (1998) & \pageref{hm:KBNNvk} \\
KBNPvkb & 6 & \textbf{h\#6.5} & h\#8 & 3 &  & \pageref{hm:KBNPvkb} \\
KBNvkrb & 6 & \textbf{h\#6.5} & h\#10 & 7 &  & \pageref{hm:KBNvkrb} \\
KBvkb & 4 & \textbf{h\#6.5} & h\#10 & 7 &  & \pageref{hm:KBvkb} \\
KBvkqb & 5 & \textbf{h\#6.5} & h\#9.5 & 6 & Sheglow (1998) & \pageref{hm:KBvkqb} \\
KBvkrb & 5 & \textbf{h\#6.5} & h\#9.5 & 6 &  & \pageref{hm:KBvkrb} \\
KNNPvkn & 6 & \textbf{h\#6.5} & h\#8 & 3 &  & \pageref{hm:KNNPvkn} \\
KNNPvkr & 6 & \textbf{h\#6.5} & h\#9 & 5 &  & \pageref{hm:KNNPvkr} \\
KNNvk & 4 & \textbf{h\#6.5} & h\#8 & 3 & Bebesi (1910) & \pageref{hm:KNNvk} \\
KNNvkb & 5 & \textbf{h\#6.5} & h\#10 & 7 &  & \pageref{hm:KNNvkb} \\
KNNvkbb & 6 & \textbf{h\#6.5} & h\#10 & 7 &  & \pageref{hm:KNNvkbb} \\
KNNvkbn & 6 & \textbf{h\#6.5} & h\#9 & 5 &  & \pageref{hm:KNNvkbn} \\
KNNvkn & 5 & \textbf{h\#6.5} & h\#7.5 & 2 &  & \pageref{hm:KNNvkn} \\
KNNvknn & 6 & \textbf{h\#6.5} & h\#8 & 3 &  & \pageref{hm:KNNvknn} \\
KNNvknp & 6 & \textbf{h\#6.5} & h\#9 & 5 &  & \pageref{hm:KNNvknp} \\
KNNvkqn & 6 & \textbf{h\#6.5} & h\#9 & 5 &  & \pageref{hm:KNNvkqn} \\
KNNvkqq & 6 & \textbf{h\#6.5} & h\#8.5 & 4 &  & \pageref{hm:KNNvkqq} \\
KNNvkrb & 6 & \textbf{h\#6.5} & h\#10 & 7 &  & \pageref{hm:KNNvkrb} \\
KNNvkrn & 6 & \textbf{h\#6.5} & h\#8 & 3 &  & \pageref{hm:KNNvkrn} \\
KNNvkrr & 6 & \textbf{h\#6.5} & h\#8.5 & 4 &  & \pageref{hm:KNNvkrr} \\
KNPPPvk & 6 & \textbf{h\#6.5} & h\#8 & 3 &  & \pageref{hm:KNPPPvk} \\
KPPPPvk & 6 & \textbf{h\#6.5} & h\#7.5 & 2 &  & \pageref{hm:KPPPPvk} \\
KPPPvkq & 6 & \textbf{h\#6.5} & h\#8.5 & 4 &  & \pageref{hm:KPPPvkq} \\
KPPPvkr & 6 & \textbf{h\#6.5} & h\#8.5 & 4 &  & \pageref{hm:KPPPvkr} \\
KPPvkr & 5 & \textbf{h\#6.5} & h\#8.5 & 4 &  & \pageref{hm:KPPvkr} \\
KPPvkrr & 6 & \textbf{h\#6.5} & h\#8.5 & 4 &  & \pageref{hm:KPPvkrr} \\
KPvk & 3 & \textbf{h\#6.5} & h\#8.5 & 4 & Niemann (1947) & \pageref{hm:KPvk} \\
KQvkqnn & 6 & \textbf{h\#6.5} & h\#8 & 3 &  & \pageref{hm:KQvkqnn} \\
KRBBvkq & 6 & \textbf{h\#6.5} & h\#7.5 & 2 &  & \pageref{hm:KRBBvkq} \\
KRBvkq & 5 & \textbf{h\#6.5} & h\#7.5 & 2 &  & \pageref{hm:KRBvkq} \\
KRvkbbp & 6 & \textbf{h\#6.5} & h\#7.5 & 2 &  & \pageref{hm:KRvkbbp} \\
KRvkbnp & 6 & \textbf{h\#6.5} & h\#7.5 & 2 &  & \pageref{hm:KRvkbnp} \\
KRvkbpp & 6 & \textbf{h\#6.5} & h\#8 & 3 &  & \pageref{hm:KRvkbpp} \\
KRvknnn & 6 & \textbf{h\#6.5} & h\#7.5 & 2 &  & \pageref{hm:KRvknnn} \\
KRvknpp & 6 & \textbf{h\#6.5} & h\#7.5 & 2 &  & \pageref{hm:KRvknpp} \\
KRvkqb & 5 & \textbf{h\#6.5} & h\#7.5 & 2 & Sheglow (1998) & \pageref{hm:KRvkqb} \\
KRvkqbb & 6 & \textbf{h\#6.5} & h\#8 & 3 &  & \pageref{hm:KRvkqbb} \\
KRvkqbp & 6 & \textbf{h\#6.5} & h\#8 & 3 &  & \pageref{hm:KRvkqbp} \\
KRvkqp & 5 & \textbf{h\#6.5} & h\#7.5 & 2 &  & \pageref{hm:KRvkqp} \\
KRvkqpp & 6 & \textbf{h\#6.5} & h\#7.5 & 2 &  & \pageref{hm:KRvkqpp} \\
KRvkqq & 5 & \textbf{h\#6.5} & h\#8 & 3 &  & \pageref{hm:KRvkqq} \\
KRvkqr & 5 & \textbf{h\#6.5} & h\#8 & 3 &  & \pageref{hm:KRvkqr} \\
KRvkqrn & 6 & \textbf{h\#6.5} & h\#8 & 3 &  & \pageref{hm:KRvkqrn} \\
KRvkrbb & 6 & \textbf{h\#6.5} & h\#7.5 & 2 &  & \pageref{hm:KRvkrbb} \\
KRvkrbn & 6 & \textbf{h\#6.5} & h\#8 & 3 &  & \pageref{hm:KRvkrbn} \\
KRvkrbp & 6 & \textbf{h\#6.5} & h\#7.5 & 2 &  & \pageref{hm:KRvkrbp} \\
KRvkrnp & 6 & \textbf{h\#6.5} & h\#8 & 3 &  & \pageref{hm:KRvkrnp} \\
KRvkrrb & 6 & \textbf{h\#6.5} & h\#8 & 3 &  & \pageref{hm:KRvkrrb} \\
KRvkrrn & 6 & \textbf{h\#6.5} & h\#8 & 3 &  & \pageref{hm:KRvkrrn} \\
KRvkrrr & 6 & \textbf{h\#6.5} & h\#8 & 3 &  & \pageref{hm:KRvkrrr} \\
KBBBvkr & 6 & \textbf{h\#6} & h\#8 & 4 &  & \pageref{hm:KBBBvkr} \\
KBBNvkr & 6 & \textbf{h\#6} & h\#8.5 & 5 &  & \pageref{hm:KBBNvkr} \\
KBBvkr & 5 & \textbf{h\#6} & h\#7.5 & 3 &  & \pageref{hm:KBBvkr} \\
KBNNNvk & 6 & \textbf{h\#6} & h\#8 & 4 &  & \pageref{hm:KBNNNvk} \\
KBNNvkb & 6 & \textbf{h\#6} & h\#8 & 4 &  & \pageref{hm:KBNNvkb} \\
KBNNvkn & 6 & \textbf{h\#6} & h\#8 & 4 &  & \pageref{hm:KBNNvkn} \\
KBNNvkr & 6 & \textbf{h\#6} & h\#8 & 4 &  & \pageref{hm:KBNNvkr} \\
KBNvk & 4 & \textbf{h\#6} & h\#8 & 4 &  & \pageref{hm:KBNvk} \\
KNNNvk & 5 & \textbf{h\#6} & h\#8 & 4 &  & \pageref{hm:KNNNvk} \\
KNNNvkb & 6 & \textbf{h\#6} & h\#8 & 4 &  & \pageref{hm:KNNNvkb} \\
KNNNvkq & 6 & \textbf{h\#6} & h\#8.5 & 5 &  & \pageref{hm:KNNNvkq} \\
KNNNvkr & 6 & \textbf{h\#6} & h\#8 & 4 &  & \pageref{hm:KNNNvkr} \\
KQBBvkq & 6 & \textbf{h\#6} & h\#7 & 2 &  & \pageref{hm:KQBBvkq} \\
KQBBvkr & 6 & \textbf{h\#6} & h\#7 & 2 &  & \pageref{hm:KQBBvkr} \\
KQBNvk & 5 & \textbf{h\#6} & h\#8 & 4 &  & \pageref{hm:KQBNvk} \\
KQBvk & 4 & \textbf{h\#6} & h\#7 & 2 &  & \pageref{hm:KQBvk} \\
KQNNNvk & 6 & \textbf{h\#6} & h\#8 & 4 &  & \pageref{hm:KQNNNvk} \\
KQNNvkn & 6 & \textbf{h\#6} & h\#8 & 4 &  & \pageref{hm:KQNNvkn} \\
KQNvk & 4 & \textbf{h\#6} & h\#7 & 2 &  & \pageref{hm:KQNvk} \\
KQQPPvk & 6 & \textbf{h\#6} & h\#7 & 2 &  & \pageref{hm:KQQPPvk} \\
KQQPvk & 5 & \textbf{h\#6} & h\#7 & 2 &  & \pageref{hm:KQQPvk} \\
KQQPvkp & 6 & \textbf{h\#6} & h\#7 & 2 &  & \pageref{hm:KQQPvkp} \\
KQQPvkq & 6 & \textbf{h\#6} & h\#7 & 2 &  & \pageref{hm:KQQPvkq} \\
KQQvk & 4 & \textbf{h\#6} & h\#7 & 2 &  & \pageref{hm:KQQvk} \\
KQQvkbn & 6 & \textbf{h\#6} & h\#7 & 2 &  & \pageref{hm:KQQvkbn} \\
KQQvkn & 5 & \textbf{h\#6} & h\#7 & 2 &  & \pageref{hm:KQQvkn} \\
KQQvknn & 6 & \textbf{h\#6} & h\#7 & 2 &  & \pageref{hm:KQQvknn} \\
KQQvknp & 6 & \textbf{h\#6} & h\#7 & 2 &  & \pageref{hm:KQQvknp} \\
KQQvkp & 5 & \textbf{h\#6} & h\#7 & 2 &  & \pageref{hm:KQQvkp} \\
KQQvkpp & 6 & \textbf{h\#6} & h\#7 & 2 &  & \pageref{hm:KQQvkpp} \\
KQQvkqn & 6 & \textbf{h\#6} & h\#7 & 2 &  & \pageref{hm:KQQvkqn} \\
KQQvkqp & 6 & \textbf{h\#6} & h\#7 & 2 &  & \pageref{hm:KQQvkqp} \\
KQQvkqq & 6 & \textbf{h\#6} & h\#8 & 4 &  & \pageref{hm:KQQvkqq} \\
KQRBBvk & 6 & \textbf{h\#6} & h\#8 & 4 &  & \pageref{hm:KQRBBvk} \\
KQRBPvk & 6 & \textbf{h\#6} & h\#7 & 2 &  & \pageref{hm:KQRBPvk} \\
KQRBvk & 5 & \textbf{h\#6} & h\#7 & 2 &  & \pageref{hm:KQRBvk} \\
KQRBvkb & 6 & \textbf{h\#6} & h\#7 & 2 &  & \pageref{hm:KQRBvkb} \\
KQRBvkn & 6 & \textbf{h\#6} & h\#7 & 2 &  & \pageref{hm:KQRBvkn} \\
KQRBvkp & 6 & \textbf{h\#6} & h\#14 & 16 &  & \pageref{hm:KQRBvkp} \\
KQRPPvk & 6 & \textbf{h\#6} & h\#7 & 2 &  & \pageref{hm:KQRPPvk} \\
KQRPvk & 5 & \textbf{h\#6} & h\#7 & 2 &  & \pageref{hm:KQRPvk} \\
KQRPvkb & 6 & \textbf{h\#6} & h\#7 & 2 &  & \pageref{hm:KQRPvkb} \\
KQRPvkn & 6 & \textbf{h\#6} & h\#8 & 4 &  & \pageref{hm:KQRPvkn} \\
KQRPvkp & 6 & \textbf{h\#6} & h\#10 & 8 &  & \pageref{hm:KQRPvkp} \\
KQRPvkq & 6 & \textbf{h\#6} & h\#8 & 4 &  & \pageref{hm:KQRPvkq} \\
KQRPvkr & 6 & \textbf{h\#6} & h\#7 & 2 &  & \pageref{hm:KQRPvkr} \\
KQRvk & 4 & \textbf{h\#6} & h\#7 & 2 &  & \pageref{hm:KQRvk} \\
KQRvkbb & 6 & \textbf{h\#6} & h\#7 & 2 &  & \pageref{hm:KQRvkbb} \\
KQRvkbn & 6 & \textbf{h\#6} & h\#7 & 2 &  & \pageref{hm:KQRvkbn} \\
KQRvkbp & 6 & \textbf{h\#6} & h\#7 & 2 &  & \pageref{hm:KQRvkbp} \\
KQRvkn & 5 & \textbf{h\#6} & h\#7 & 2 &  & \pageref{hm:KQRvkn} \\
KQRvknn & 6 & \textbf{h\#6} & h\#7 & 2 &  & \pageref{hm:KQRvknn} \\
KQRvknp & 6 & \textbf{h\#6} & h\#7 & 2 &  & \pageref{hm:KQRvknp} \\
KQRvkp & 5 & \textbf{h\#6} & h\#7 & 2 &  & \pageref{hm:KQRvkp} \\
KQRvkpp & 6 & \textbf{h\#6} & h\#7 & 2 &  & \pageref{hm:KQRvkpp} \\
KQRvkq & 5 & \textbf{h\#6} & h\#7 & 2 &  & \pageref{hm:KQRvkq} \\
KQRvkqb & 6 & \textbf{h\#6} & h\#8.5 & 5 &  & \pageref{hm:KQRvkqb} \\
KQRvkqn & 6 & \textbf{h\#6} & h\#7 & 2 &  & \pageref{hm:KQRvkqn} \\
KQRvkqp & 6 & \textbf{h\#6} & h\#7 & 2 &  & \pageref{hm:KQRvkqp} \\
KQRvkqq & 6 & \textbf{h\#6} & h\#8.5 & 5 &  & \pageref{hm:KQRvkqq} \\
KQRvkqr & 6 & \textbf{h\#6} & h\#8 & 4 &  & \pageref{hm:KQRvkqr} \\
KQRvkr & 5 & \textbf{h\#6} & h\#7 & 2 &  & \pageref{hm:KQRvkr} \\
KQRvkrb & 6 & \textbf{h\#6} & h\#7 & 2 &  & \pageref{hm:KQRvkrb} \\
KQRvkrn & 6 & \textbf{h\#6} & h\#7 & 2 &  & \pageref{hm:KQRvkrn} \\
KQRvkrp & 6 & \textbf{h\#6} & h\#7 & 2 &  & \pageref{hm:KQRvkrp} \\
KQRvkrr & 6 & \textbf{h\#6} & h\#8 & 4 &  & \pageref{hm:KQRvkrr} \\
KQvk & 3 & \textbf{h\#6} & h\#7 & 2 & Csaszar (1945) & \pageref{hm:KQvk} \\
KQvkbbn & 6 & \textbf{h\#6} & h\#7.5 & 3 &  & \pageref{hm:KQvkbbn} \\
KQvkbbp & 6 & \textbf{h\#6} & h\#7 & 2 &  & \pageref{hm:KQvkbbp} \\
KQvkbn & 5 & \textbf{h\#6} & h\#7 & 2 &  & \pageref{hm:KQvkbn} \\
KQvkbnn & 6 & \textbf{h\#6} & h\#7.5 & 3 &  & \pageref{hm:KQvkbnn} \\
KQvkbnp & 6 & \textbf{h\#6} & h\#7.5 & 3 &  & \pageref{hm:KQvkbnp} \\
KQvkbpp & 6 & \textbf{h\#6} & h\#7 & 2 &  & \pageref{hm:KQvkbpp} \\
KQvkn & 4 & \textbf{h\#6} & h\#7 & 2 &  & \pageref{hm:KQvkn} \\
KQvknn & 5 & \textbf{h\#6} & h\#7 & 2 &  & \pageref{hm:KQvknn} \\
KQvknnn & 6 & \textbf{h\#6} & h\#7.5 & 3 &  & \pageref{hm:KQvknnn} \\
KQvknnp & 6 & \textbf{h\#6} & h\#7.5 & 3 &  & \pageref{hm:KQvknnp} \\
KQvknp & 5 & \textbf{h\#6} & h\#7 & 2 &  & \pageref{hm:KQvknp} \\
KQvknpp & 6 & \textbf{h\#6} & h\#7 & 2 &  & \pageref{hm:KQvknpp} \\
KQvkp & 4 & \textbf{h\#6} & h\#7 & 2 &  & \pageref{hm:KQvkp} \\
KQvkpp & 5 & \textbf{h\#6} & h\#7 & 2 &  & \pageref{hm:KQvkpp} \\
KQvkppp & 6 & \textbf{h\#6} & h\#7 & 2 &  & \pageref{hm:KQvkppp} \\
KQvkqbb & 6 & \textbf{h\#6} & h\#7 & 2 &  & \pageref{hm:KQvkqbb} \\
KQvkqbn & 6 & \textbf{h\#6} & h\#7.5 & 3 &  & \pageref{hm:KQvkqbn} \\
KQvkqbp & 6 & \textbf{h\#6} & h\#7.5 & 3 &  & \pageref{hm:KQvkqbp} \\
KQvkqn & 5 & \textbf{h\#6} & h\#7 & 2 &  & \pageref{hm:KQvkqn} \\
KQvkqnp & 6 & \textbf{h\#6} & h\#7.5 & 3 &  & \pageref{hm:KQvkqnp} \\
KQvkqp & 5 & \textbf{h\#6} & h\#7 & 2 &  & \pageref{hm:KQvkqp} \\
KQvkqpp & 6 & \textbf{h\#6} & h\#7.5 & 3 &  & \pageref{hm:KQvkqpp} \\
KQvkqqb & 6 & \textbf{h\#6} & h\#7 & 2 &  & \pageref{hm:KQvkqqb} \\
KQvkqqn & 6 & \textbf{h\#6} & h\#7.5 & 3 &  & \pageref{hm:KQvkqqn} \\
KQvkqqp & 6 & \textbf{h\#6} & h\#7.5 & 3 &  & \pageref{hm:KQvkqqp} \\
KQvkqqq & 6 & \textbf{h\#6} & h\#7 & 2 &  & \pageref{hm:KQvkqqq} \\
KRBBNvk & 6 & \textbf{h\#6} & h\#8 & 4 &  & \pageref{hm:KRBBNvk} \\
KRBBvkr & 6 & \textbf{h\#6} & h\#7.5 & 3 &  & \pageref{hm:KRBBvkr} \\
KRBNvk & 5 & \textbf{h\#6} & h\#8 & 4 &  & \pageref{hm:KRBNvk} \\
KRBvk & 4 & \textbf{h\#6} & h\#7 & 2 &  & \pageref{hm:KRBvk} \\
KRBvkb & 5 & \textbf{h\#6} & h\#10 & 8 &  & \pageref{hm:KRBvkb} \\
KRBvkp & 5 & \textbf{h\#6} & h\#16 & 20 &  & \pageref{hm:KRBvkp} \\
KRBvkqb & 6 & \textbf{h\#6} & h\#10.5 & 9 &  & \pageref{hm:KRBvkqb} \\
KRBvkqq & 6 & \textbf{h\#6} & h\#8 & 4 &  & \pageref{hm:KRBvkqq} \\
KRBvkqr & 6 & \textbf{h\#6} & h\#8 & 4 &  & \pageref{hm:KRBvkqr} \\
KRBvkrr & 6 & \textbf{h\#6} & h\#7 & 2 &  & \pageref{hm:KRBvkrr} \\
KRNNNvk & 6 & \textbf{h\#6} & h\#7 & 2 &  & \pageref{hm:KRNNNvk} \\
KRNNvk & 5 & \textbf{h\#6} & h\#8 & 4 &  & \pageref{hm:KRNNvk} \\
KRNNvkn & 6 & \textbf{h\#6} & h\#7 & 2 &  & \pageref{hm:KRNNvkn} \\
KRNNvkr & 6 & \textbf{h\#6} & h\#8 & 4 &  & \pageref{hm:KRNNvkr} \\
KRNvkb & 5 & \textbf{h\#6} & h\#10 & 8 &  & \pageref{hm:KRNvkb} \\
KRNvkqq & 6 & \textbf{h\#6} & h\#8 & 4 &  & \pageref{hm:KRNvkqq} \\
KRPvk & 4 & \textbf{h\#6} & h\#8 & 4 &  & \pageref{hm:KRPvk} \\
KRRBBvk & 6 & \textbf{h\#6} & h\#8 & 4 &  & \pageref{hm:KRRBBvk} \\
KRRBPvk & 6 & \textbf{h\#6} & h\#7 & 2 &  & \pageref{hm:KRRBPvk} \\
KRRBvk & 5 & \textbf{h\#6} & h\#7 & 2 &  & \pageref{hm:KRRBvk} \\
KRRBvkb & 6 & \textbf{h\#6} & h\#8 & 4 &  & \pageref{hm:KRRBvkb} \\
KRRBvkn & 6 & \textbf{h\#6} & h\#7 & 2 &  & \pageref{hm:KRRBvkn} \\
KRRBvkp & 6 & \textbf{h\#6} & h\#8 & 4 &  & \pageref{hm:KRRBvkp} \\
KRRPPvk & 6 & \textbf{h\#6} & h\#7 & 2 &  & \pageref{hm:KRRPPvk} \\
KRRPvk & 5 & \textbf{h\#6} & h\#7 & 2 &  & \pageref{hm:KRRPvk} \\
KRRPvkb & 6 & \textbf{h\#6} & h\#8 & 4 &  & \pageref{hm:KRRPvkb} \\
KRRPvkn & 6 & \textbf{h\#6} & h\#8 & 4 &  & \pageref{hm:KRRPvkn} \\
KRRPvkp & 6 & \textbf{h\#6} & h\#10 & 8 &  & \pageref{hm:KRRPvkp} \\
KRRPvkq & 6 & \textbf{h\#6} & h\#8 & 4 &  & \pageref{hm:KRRPvkq} \\
KRRPvkr & 6 & \textbf{h\#6} & h\#7 & 2 &  & \pageref{hm:KRRPvkr} \\
KRRvkbb & 6 & \textbf{h\#6} & h\#7 & 2 &  & \pageref{hm:KRRvkbb} \\
KRRvkbn & 6 & \textbf{h\#6} & h\#7 & 2 &  & \pageref{hm:KRRvkbn} \\
KRRvkbp & 6 & \textbf{h\#6} & h\#7 & 2 &  & \pageref{hm:KRRvkbp} \\
KRRvknn & 6 & \textbf{h\#6} & h\#7 & 2 &  & \pageref{hm:KRRvknn} \\
KRRvknp & 6 & \textbf{h\#6} & h\#7 & 2 &  & \pageref{hm:KRRvknp} \\
KRRvkp & 5 & \textbf{h\#6} & h\#7 & 2 &  & \pageref{hm:KRRvkp} \\
KRRvkpp & 6 & \textbf{h\#6} & h\#7 & 2 &  & \pageref{hm:KRRvkpp} \\
KRRvkqb & 6 & \textbf{h\#6} & h\#8 & 4 &  & \pageref{hm:KRRvkqb} \\
KRRvkqn & 6 & \textbf{h\#6} & h\#7 & 2 &  & \pageref{hm:KRRvkqn} \\
KRRvkqp & 6 & \textbf{h\#6} & h\#7.5 & 3 &  & \pageref{hm:KRRvkqp} \\
KRRvkqq & 6 & \textbf{h\#6} & h\#8 & 4 &  & \pageref{hm:KRRvkqq} \\
KRRvkqr & 6 & \textbf{h\#6} & h\#8 & 4 &  & \pageref{hm:KRRvkqr} \\
KRRvkr & 5 & \textbf{h\#6} & h\#7 & 2 &  & \pageref{hm:KRRvkr} \\
KRRvkrb & 6 & \textbf{h\#6} & h\#7 & 2 &  & \pageref{hm:KRRvkrb} \\
KRRvkrn & 6 & \textbf{h\#6} & h\#7 & 2 &  & \pageref{hm:KRRvkrn} \\
KRRvkrp & 6 & \textbf{h\#6} & h\#7.5 & 3 &  & \pageref{hm:KRRvkrp} \\
KRRvkrr & 6 & \textbf{h\#6} & h\#8 & 4 &  & \pageref{hm:KRRvkrr} \\
KRvkbb & 5 & \textbf{h\#6} & h\#7 & 2 &  & \pageref{hm:KRvkbb} \\
KRvkbbb & 6 & \textbf{h\#6} & h\#7.5 & 3 &  & \pageref{hm:KRvkbbb} \\
KRvkbbn & 6 & \textbf{h\#6} & h\#7.5 & 3 &  & \pageref{hm:KRvkbbn} \\
KRvkbn & 5 & \textbf{h\#6} & h\#7 & 2 & Sheglow (1998) & \pageref{hm:KRvkbn} \\
KRvkbnn & 6 & \textbf{h\#6} & h\#7.5 & 3 &  & \pageref{hm:KRvkbnn} \\
KRvkbp & 5 & \textbf{h\#6} & h\#7 & 2 &  & \pageref{hm:KRvkbp} \\
KRvknn & 5 & \textbf{h\#6} & h\#7 & 2 &  & \pageref{hm:KRvknn} \\
KRvknnp & 6 & \textbf{h\#6} & h\#7.5 & 3 &  & \pageref{hm:KRvknnp} \\
KRvknp & 5 & \textbf{h\#6} & h\#7 & 2 &  & \pageref{hm:KRvknp} \\
KRvkp & 4 & \textbf{h\#6} & h\#7 & 2 &  & \pageref{hm:KRvkp} \\
KRvkpp & 5 & \textbf{h\#6} & h\#7 & 2 &  & \pageref{hm:KRvkpp} \\
KRvkppp & 6 & \textbf{h\#6} & h\#7.5 & 3 &  & \pageref{hm:KRvkppp} \\
KRvkqn & 5 & \textbf{h\#6} & h\#7.5 & 3 &  & \pageref{hm:KRvkqn} \\
KRvkrb & 5 & \textbf{h\#6} & h\#7 & 2 &  & \pageref{hm:KRvkrb} \\
KRvkrn & 5 & \textbf{h\#6} & h\#7.5 & 3 &  & \pageref{hm:KRvkrn} \\
KRvkrnn & 6 & \textbf{h\#6} & h\#7.5 & 3 &  & \pageref{hm:KRvkrnn} \\
KRvkrp & 5 & \textbf{h\#6} & h\#7.5 & 3 &  & \pageref{hm:KRvkrp} \\
KRvkrr & 5 & \textbf{h\#6} & h\#7.5 & 3 &  & \pageref{hm:KRvkrr} \\
KNNNvkn & 6 & \textbf{h\#5.5} & h\#7 & 3 &  & \pageref{hm:KNNNvkn} \\
KQBvkq & 5 & \textbf{h\#5.5} & h\#7 & 3 &  & \pageref{hm:KQBvkq} \\
KQNvkq & 5 & \textbf{h\#5.5} & h\#7 & 3 &  & \pageref{hm:KQNvkq} \\
KQQvkqb & 6 & \textbf{h\#5.5} & h\#8 & 5 &  & \pageref{hm:KQQvkqb} \\
KQvkbb & 5 & \textbf{h\#5.5} & h\#7 & 3 &  & \pageref{hm:KQvkbb} \\
KQvkbbb & 6 & \textbf{h\#5.5} & h\#7 & 3 &  & \pageref{hm:KQvkbbb} \\
KQvkbp & 5 & \textbf{h\#5.5} & h\#7 & 3 &  & \pageref{hm:KQvkbp} \\
KQvkq & 4 & \textbf{h\#5.5} & h\#7 & 3 & Sheglow (1998) & \pageref{hm:KQvkq} \\
KQvkqb & 5 & \textbf{h\#5.5} & h\#7.5 & 4 &  & \pageref{hm:KQvkqb} \\
KQvkqq & 5 & \textbf{h\#5.5} & h\#7.5 & 4 &  & \pageref{hm:KQvkqq} \\
KQvkqqr & 6 & \textbf{h\#5.5} & h\#6.5 & 2 &  & \pageref{hm:KQvkqqr} \\
KQvkqr & 5 & \textbf{h\#5.5} & h\#6.5 & 2 &  & \pageref{hm:KQvkqr} \\
KQvkqrb & 6 & \textbf{h\#5.5} & h\#6.5 & 2 &  & \pageref{hm:KQvkqrb} \\
KQvkqrn & 6 & \textbf{h\#5.5} & h\#7.5 & 4 &  & \pageref{hm:KQvkqrn} \\
KQvkqrp & 6 & \textbf{h\#5.5} & h\#7.5 & 4 &  & \pageref{hm:KQvkqrp} \\
KQvkrbb & 6 & \textbf{h\#5.5} & h\#6.5 & 2 &  & \pageref{hm:KQvkrbb} \\
KQvkrbn & 6 & \textbf{h\#5.5} & h\#7.5 & 4 &  & \pageref{hm:KQvkrbn} \\
KQvkrbp & 6 & \textbf{h\#5.5} & h\#6.5 & 2 &  & \pageref{hm:KQvkrbp} \\
KQvkrn & 5 & \textbf{h\#5.5} & h\#6.5 & 2 &  & \pageref{hm:KQvkrn} \\
KQvkrnn & 6 & \textbf{h\#5.5} & h\#7.5 & 4 &  & \pageref{hm:KQvkrnn} \\
KQvkrnp & 6 & \textbf{h\#5.5} & h\#7.5 & 4 &  & \pageref{hm:KQvkrnp} \\
KQvkrpp & 6 & \textbf{h\#5.5} & h\#6.5 & 2 &  & \pageref{hm:KQvkrpp} \\
KRNvkq & 5 & \textbf{h\#5.5} & h\#7 & 3 &  & \pageref{hm:KRNvkq} \\
KRRBvkq & 6 & \textbf{h\#5.5} & h\#7.5 & 4 &  & \pageref{hm:KRRBvkq} \\
KRRNvkq & 6 & \textbf{h\#5.5} & h\#6.5 & 2 &  & \pageref{hm:KRRNvkq} \\
KRvkq & 4 & \textbf{h\#5.5} & h\#7.5 & 4 &  & \pageref{hm:KRvkq} \\
KRvkr & 4 & \textbf{h\#5.5} & h\#7 & 3 &  & \pageref{hm:KRvkr} \\
KNNNNvk & 6 & \textbf{h\#5} & h\#7 & 4 &  & \pageref{hm:KNNNNvk} \\
KQBvkqq & 6 & \textbf{h\#5} & h\#8 & 6 &  & \pageref{hm:KQBvkqq} \\
KQBvkr & 5 & \textbf{h\#5} & h\#6 & 2 &  & \pageref{hm:KQBvkr} \\
KQBvkrr & 6 & \textbf{h\#5} & h\#6 & 2 &  & \pageref{hm:KQBvkrr} \\
KQNvkqq & 6 & \textbf{h\#5} & h\#7 & 4 &  & \pageref{hm:KQNvkqq} \\
KQQPvkb & 6 & \textbf{h\#5} & h\#7 & 4 &  & \pageref{hm:KQQPvkb} \\
KQQPvkn & 6 & \textbf{h\#5} & h\#9 & 8 &  & \pageref{hm:KQQPvkn} \\
KQQPvkr & 6 & \textbf{h\#5} & h\#6 & 2 &  & \pageref{hm:KQQPvkr} \\
KQQvkb & 5 & \textbf{h\#5} & h\#6 & 2 &  & \pageref{hm:KQQvkb} \\
KQQvkbb & 6 & \textbf{h\#5} & h\#7 & 4 &  & \pageref{hm:KQQvkbb} \\
KQQvkbp & 6 & \textbf{h\#5} & h\#7 & 4 &  & \pageref{hm:KQQvkbp} \\
KQQvkq & 5 & \textbf{h\#5} & h\#7 & 4 &  & \pageref{hm:KQQvkq} \\
KQQvkqr & 6 & \textbf{h\#5} & h\#7 & 4 &  & \pageref{hm:KQQvkqr} \\
KQQvkr & 5 & \textbf{h\#5} & h\#6 & 2 &  & \pageref{hm:KQQvkr} \\
KQQvkrb & 6 & \textbf{h\#5} & h\#6.5 & 3 &  & \pageref{hm:KQQvkrb} \\
KQQvkrn & 6 & \textbf{h\#5} & h\#6 & 2 &  & \pageref{hm:KQQvkrn} \\
KQQvkrp & 6 & \textbf{h\#5} & h\#6 & 2 &  & \pageref{hm:KQQvkrp} \\
KQQvkrr & 6 & \textbf{h\#5} & h\#6 & 2 &  & \pageref{hm:KQQvkrr} \\
KQRBNvk & 6 & \textbf{h\#5} & h\#7 & 4 &  & \pageref{hm:KQRBNvk} \\
KQRBvkq & 6 & \textbf{h\#5} & h\#7 & 4 &  & \pageref{hm:KQRBvkq} \\
KQRNNvk & 6 & \textbf{h\#5} & h\#6 & 2 &  & \pageref{hm:KQRNNvk} \\
KQRNPvk & 6 & \textbf{h\#5} & h\#7 & 4 &  & \pageref{hm:KQRNPvk} \\
KQRNvk & 5 & \textbf{h\#5} & h\#7 & 4 &  & \pageref{hm:KQRNvk} \\
KQRNvkb & 6 & \textbf{h\#5} & h\#6 & 2 &  & \pageref{hm:KQRNvkb} \\
KQRNvkn & 6 & \textbf{h\#5} & h\#6 & 2 &  & \pageref{hm:KQRNvkn} \\
KQRNvkp & 6 & \textbf{h\#5} & h\#8 & 6 &  & \pageref{hm:KQRNvkp} \\
KQRNvkq & 6 & \textbf{h\#5} & h\#6 & 2 &  & \pageref{hm:KQRNvkq} \\
KQRNvkr & 6 & \textbf{h\#5} & h\#6 & 2 &  & \pageref{hm:KQRNvkr} \\
KQRvkb & 5 & \textbf{h\#5} & h\#7 & 4 &  & \pageref{hm:KQRvkb} \\
KQvkb & 4 & \textbf{h\#5} & h\#6 & 2 &  & \pageref{hm:KQvkb} \\
KQvkqrr & 6 & \textbf{h\#5} & h\#6 & 2 &  & \pageref{hm:KQvkqrr} \\
KQvkr & 4 & \textbf{h\#5} & h\#6 & 2 &  & \pageref{hm:KQvkr} \\
KQvkrb & 5 & \textbf{h\#5} & h\#6.5 & 3 &  & \pageref{hm:KQvkrb} \\
KQvkrp & 5 & \textbf{h\#5} & h\#6.5 & 3 &  & \pageref{hm:KQvkrp} \\
KQvkrrb & 6 & \textbf{h\#5} & h\#6 & 2 &  & \pageref{hm:KQvkrrb} \\
KQvkrrn & 6 & \textbf{h\#5} & h\#6 & 2 &  & \pageref{hm:KQvkrrn} \\
KQvkrrp & 6 & \textbf{h\#5} & h\#6 & 2 &  & \pageref{hm:KQvkrrp} \\
KRBvkr & 5 & \textbf{h\#5} & h\#7 & 4 &  & \pageref{hm:KRBvkr} \\
KRNvk & 4 & \textbf{h\#5} & h\#7 & 4 & Maslar (1981) & \pageref{hm:KRNvk} \\
KRRBNvk & 6 & \textbf{h\#5} & h\#6 & 2 &  & \pageref{hm:KRRBNvk} \\
KRRNNvk & 6 & \textbf{h\#5} & h\#6 & 2 &  & \pageref{hm:KRRNNvk} \\
KRRNPvk & 6 & \textbf{h\#5} & h\#7 & 4 &  & \pageref{hm:KRRNPvk} \\
KRRNvk & 5 & \textbf{h\#5} & h\#7 & 4 &  & \pageref{hm:KRRNvk} \\
KRRNvkb & 6 & \textbf{h\#5} & h\#7 & 4 &  & \pageref{hm:KRRNvkb} \\
KRRNvkn & 6 & \textbf{h\#5} & h\#6 & 2 &  & \pageref{hm:KRRNvkn} \\
KRRNvkp & 6 & \textbf{h\#5} & h\#8 & 6 &  & \pageref{hm:KRRNvkp} \\
KRRNvkr & 6 & \textbf{h\#5} & h\#6.5 & 3 &  & \pageref{hm:KRRNvkr} \\
KRRvkb & 5 & \textbf{h\#5} & h\#7 & 4 &  & \pageref{hm:KRRvkb} \\
KRRvkn & 5 & \textbf{h\#5} & h\#7 & 4 &  & \pageref{hm:KRRvkn} \\
KRRvkq & 5 & \textbf{h\#5} & h\#7 & 4 &  & \pageref{hm:KRRvkq} \\
KRvkb & 4 & \textbf{h\#5} & h\#7 & 4 &  & \pageref{hm:KRvkb} \\
KRvkn & 4 & \textbf{h\#5} & h\#7 & 4 &  & \pageref{hm:KRvkn} \\
KQBvkqr & 6 & \textbf{h\#4.5} & h\#6.5 & 4 &  & \pageref{hm:KQBvkqr} \\
KQvkrr & 5 & \textbf{h\#4.5} & h\#6 & 3 &  & \pageref{hm:KQvkrr} \\
KQvkrrr & 6 & \textbf{h\#4.5} & h\#6 & 3 &  & \pageref{hm:KQvkrrr} \\
KRRBvkr & 6 & \textbf{h\#4.5} & h\#7.5 & 6 &  & \pageref{hm:KRRBvkr} \\
KRvk & 3 & \textbf{h\#4.5} & h\#7 & 5 & Ban (1965) & \pageref{hm:KRvk} \\
KQQBvkp & 6 & \textbf{h\#4} & h\#5 & 2 &  & \pageref{hm:KQQBvkp} \\
KQQBvkq & 6 & \textbf{h\#4} & h\#5 & 2 &  & \pageref{hm:KQQBvkq} \\
KQQNPvk & 6 & \textbf{h\#4} & h\#5 & 2 &  & \pageref{hm:KQQNPvk} \\
KQQNvk & 5 & \textbf{h\#4} & h\#5 & 2 &  & \pageref{hm:KQQNvk} \\
KQQNvkb & 6 & \textbf{h\#4} & h\#5 & 2 &  & \pageref{hm:KQQNvkb} \\
KQQNvkn & 6 & \textbf{h\#4} & h\#5 & 2 &  & \pageref{hm:KQQNvkn} \\
KQQNvkp & 6 & \textbf{h\#4} & h\#5 & 2 &  & \pageref{hm:KQQNvkp} \\
KQQNvkq & 6 & \textbf{h\#4} & h\#5 & 2 &  & \pageref{hm:KQQNvkq} \\
KQQNvkr & 6 & \textbf{h\#4} & h\#5 & 2 &  & \pageref{hm:KQQNvkr} \\
KQQQvkn & 6 & \textbf{h\#4} & h\#5 & 2 &  & \pageref{hm:KQQQvkn} \\
KQQRvkb & 6 & \textbf{h\#4} & h\#6 & 4 &  & \pageref{hm:KQQRvkb} \\
KQQRvkn & 6 & \textbf{h\#4} & h\#5 & 2 &  & \pageref{hm:KQQRvkn} \\
KQQRvkq & 6 & \textbf{h\#4} & h\#7 & 6 &  & \pageref{hm:KQQRvkq} \\
KQRBvkr & 6 & \textbf{h\#4} & h\#6 & 4 &  & \pageref{hm:KQRBvkr} \\
KQRRBvk & 6 & \textbf{h\#4} & h\#6 & 4 &  & \pageref{hm:KQRRBvk} \\
KQRRNvk & 6 & \textbf{h\#4} & h\#5 & 2 &  & \pageref{hm:KQRRNvk} \\
KQRRPvk & 6 & \textbf{h\#4} & h\#7 & 6 &  & \pageref{hm:KQRRPvk} \\
KQRRvkb & 6 & \textbf{h\#4} & h\#7 & 6 &  & \pageref{hm:KQRRvkb} \\
KQRRvkn & 6 & \textbf{h\#4} & h\#5 & 2 &  & \pageref{hm:KQRRvkn} \\
KQRRvkp & 6 & \textbf{h\#4} & h\#7 & 6 &  & \pageref{hm:KQRRvkp} \\
KQRRvkq & 6 & \textbf{h\#4} & h\#6 & 4 &  & \pageref{hm:KQRRvkq} \\
KRRRNvk & 6 & \textbf{h\#4} & h\#5 & 2 &  & \pageref{hm:KRRRNvk} \\
KRRRvkb & 6 & \textbf{h\#4} & h\#7 & 6 &  & \pageref{hm:KRRRvkb} \\
KRRRvkn & 6 & \textbf{h\#4} & h\#5 & 2 &  & \pageref{hm:KRRRvkn} \\
KRRRvkp & 6 & \textbf{h\#4} & h\#7 & 6 &  & \pageref{hm:KRRRvkp} \\
KRRvk & 4 & \textbf{h\#4} & h\#7 & 6 &  & \pageref{hm:KRRvk} \\
KQQBvkr & 6 & \textbf{h\#3.5} & h\#5 & 3 &  & \pageref{hm:KQQBvkr} \\
KQRRvkr & 6 & \textbf{h\#3.5} & h\#5 & 3 &  & \pageref{hm:KQRRvkr} \\
KRRRvkq & 6 & \textbf{h\#3.5} & h\#6 & 5 &  & \pageref{hm:KRRRvkq} \\
KRRRvkr & 6 & \textbf{h\#3.5} & h\#4.5 & 2 &  & \pageref{hm:KRRRvkr} \\
KQQBBvk & 6 & \textbf{h\#3} & h\#4 & 2 &  & \pageref{hm:KQQBBvk} \\
KQQBNvk & 6 & \textbf{h\#3} & h\#4 & 2 &  & \pageref{hm:KQQBNvk} \\
KQQBPvk & 6 & \textbf{h\#3} & h\#5 & 4 &  & \pageref{hm:KQQBPvk} \\
KQQBvk & 5 & \textbf{h\#3} & h\#4 & 2 &  & \pageref{hm:KQQBvk} \\
KQQBvkb & 6 & \textbf{h\#3} & h\#4 & 2 &  & \pageref{hm:KQQBvkb} \\
KQQBvkn & 6 & \textbf{h\#3} & h\#5 & 4 &  & \pageref{hm:KQQBvkn} \\
KQQNNvk & 6 & \textbf{h\#3} & h\#4 & 2 &  & \pageref{hm:KQQNNvk} \\
KQQQvkb & 6 & \textbf{h\#3} & h\#5 & 4 &  & \pageref{hm:KQQQvkb} \\
KQQRNvk & 6 & \textbf{h\#3} & h\#4 & 2 &  & \pageref{hm:KQQRNvk} \\
KQQRPvk & 6 & \textbf{h\#3} & h\#4 & 2 &  & \pageref{hm:KQQRPvk} \\
KQQRvk & 5 & \textbf{h\#3} & h\#4 & 2 &  & \pageref{hm:KQQRvk} \\
KQQRvkp & 6 & \textbf{h\#3} & h\#6 & 6 &  & \pageref{hm:KQQRvkp} \\
KQQRvkr & 6 & \textbf{h\#3} & h\#5 & 4 &  & \pageref{hm:KQQRvkr} \\
KQRRvk & 5 & \textbf{h\#3} & h\#5 & 4 &  & \pageref{hm:KQRRvk} \\
KRRRBvk & 6 & \textbf{h\#3} & h\#5 & 4 &  & \pageref{hm:KRRRBvk} \\
KRRRPvk & 6 & \textbf{h\#3} & h\#7 & 8 &  & \pageref{hm:KRRRPvk} \\
KRRRvk & 5 & \textbf{h\#3} & h\#5 & 4 &  & \pageref{hm:KRRRvk} \\
KQQQvkp & 6 & \textbf{h\#2.5} & h\#3.5 & 2 &  & \pageref{hm:KQQQvkp} \\
KQQQvkq & 6 & \textbf{h\#2.5} & h\#5 & 5 &  & \pageref{hm:KQQQvkq} \\
KQQQvkr & 6 & \textbf{h\#2.5} & h\#4 & 3 &  & \pageref{hm:KQQQvkr} \\
KQQQBvk & 6 & \textbf{h\#2} & h\#3 & 2 &  & \pageref{hm:KQQQBvk} \\
KQQQNvk & 6 & \textbf{h\#2} & h\#3 & 2 &  & \pageref{hm:KQQQNvk} \\
KQQQPvk & 6 & \textbf{h\#2} & h\#3 & 2 &  & \pageref{hm:KQQQPvk} \\
KQQQQvk & 6 & \textbf{h\#2} & h\#3 & 2 &  & \pageref{hm:KQQQQvk} \\
KQQQRvk & 6 & \textbf{h\#2} & h\#3 & 2 &  & \pageref{hm:KQQQRvk} \\
KQQQvk & 5 & \textbf{h\#2} & h\#3 & 2 &  & \pageref{hm:KQQQvk} \\
KQQRBvk & 6 & \textbf{h\#2} & h\#4 & 4 &  & \pageref{hm:KQQRBvk} \\
KQQRRvk & 6 & \textbf{h\#2} & h\#3 & 2 &  & \pageref{hm:KQQRRvk} \\
KQRRRvk & 6 & \textbf{h\#2} & h\#3 & 2 &  & \pageref{hm:KQRRRvk} \\
KRRRRvk & 6 & \textbf{h\#2} & h\#3 & 2 &  & \pageref{hm:KRRRRvk} \\
\end{longtable}
}

\end{document}
